\import{naprochelibrary.tex}

% from §12 Item 1: topology on a set
\begin{definition}\label{topology_on}
    $T$ is a topology on $X$ iff
    $T\subseteq\pow{X}$ and
    $\emptyset\in T$ and
    $X\in T$ and
    $T$ is closed under arbitrary unions and
    $T$ is closed under binary intersections.
\end{definition}

% from §12 Item 2: topological space with topology
\begin{definition}\label{topological_space}
    $X$ is a topological space with topology $T$ iff $T$ is a topology on $X$.
\end{definition}

% from §12 Item 3: open set
\begin{definition}\label{open_in}
    $U$ is open in $X$ with respect to $T$ iff $T$ is a topology on $X$ and $U\in T$.
\end{definition}

% from §12 Item 4: discrete topology
\begin{definition}\label{discrete_topology}
    $\discrete{X} = \pow{X}$.
\end{definition}

% from §12 Item 5: indiscrete topology
\begin{definition}\label{indiscrete_topology}
    $\indiscrete{X} = \{\emptyset, X\}$.
\end{definition}

% from §12 Item 6: finite complement topology
\begin{definition}\label{finite_complement_topology}
    $\finiteComplement{X} = \{U\in\pow{X} \mid \text{$X\setminus U$ is finite or $X\setminus U = X$}\}$.
\end{definition}

% from §12 Item 7: countable complement topology
\begin{definition}\label{countable_complement_topology}
    $\countableComplement{X} = \{U\in\pow{X} \mid \text{$X\setminus U$ is countable or $X\setminus U = X$}\}$.
\end{definition}

% from §12 Item 8: finer topology
\begin{definition}\label{finer_topology}
    $T_1$ is finer than $T$ on $X$ iff
    $T$ is a topology on $X$ and
    $T_1$ is a topology on $X$ and
    $T\subseteq T_1$.
\end{definition}

% from §12 Item 9: strictly finer topology
\begin{definition}\label{strictly_finer_topology}
    $T_1$ is strictly finer than $T$ on $X$ iff $T_1$ is finer than $T$ on $X$ and $T\neq T_1$.
\end{definition}

% from §12 Item 10: coarser topology
\begin{definition}\label{coarser_topology}
    $T$ is coarser than $T_1$ on $X$ iff $T_1$ is finer than $T$ on $X$.
\end{definition}

% from §12 Item 11: strictly coarser topology
\begin{definition}\label{strictly_coarser_topology}
    $T$ is strictly coarser than $T_1$ on $X$ iff $T$ is coarser than $T_1$ on $X$ and $T\neq T_1$.
\end{definition}

% from §12 Item 12: comparable topologies
\begin{definition}\label{comparable_topologies}
    $T$ is comparable with $T_1$ on $X$ iff
    $T$ is a topology on $X$ and
    $T_1$ is a topology on $X$ and
    either $T\subseteq T_1$ or $T_1\subseteq T$.
\end{definition}

% from §13 Item 1: basis on a set
\begin{definition}\label{basis_on}
    $B$ is a basis on $X$ iff
    $B\subseteq\pow{X}$ and
    for all $x\in X$ there exists $B_1\in B$ such that $x\in B_1$ and
    for all $B_1,B_2,x$ such that $B_1\in B$ and $B_2\in B$ and $x\in B_1\inter B_2$
    there exists $B_3\in B$ such that $x\in B_3$ and $B_3\subseteq B_1\inter B_2$.
\end{definition}

% from §13 Item 2: topology generated by a basis
\begin{definition}\label{topology_generated_by_basis}
    $\basisTopology{B}{X} = \{U\in\pow{X} \mid \text{for all $x\in U$ there exists $B_1\in B$ such that $x\in B_1$ and $B_1\subseteq U$}\}$.
\end{definition}

% from §13 Item 3: basis for a topology
\begin{definition}\label{basis_for_topology}
    $B$ is a basis for $T$ on $X$ iff $B$ is a basis on $X$ and $T = \basisTopology{B}{X}$.
\end{definition}

% from §13 Item 4: topology generated by a basis is a topology
\begin{lemma}\label{basis_topology_is_topology}
    Assume $B$ is a basis on $X$.
    Then $\basisTopology{B}{X}$ is a topology on $X$.
\end{lemma}
\begin{proof}
    $\basisTopology{B}{X}\subseteq\pow{X}$.
    $\emptyset\in\pow{X}$ by \cref{powerset_bottom}.
    For all $x\in\emptyset$ there exists $B_1\in B$ such that $x\in B_1$ and $B_1\subseteq\emptyset$.
    Thus $\emptyset\in\basisTopology{B}{X}$.
    For all $x\in X$ there exists $B_1\in B$ such that $x\in B_1$.
    Thus $X\in\basisTopology{B}{X}$.
    Show that for all $M$ such that $M\subseteq\basisTopology{B}{X}$ we have $\unions{M}\in\basisTopology{B}{X}$.
    \begin{subproof}
        Fix $M$.
        Assume $M\subseteq\basisTopology{B}{X}$.
        $M\subseteq\pow{X}$.
        Then $\unions{M}\in\pow{X}$ by \cref{powerset_closed_unions}.
        Show that for all $x$ such that $x\in\unions{M}$ there exists $B_1\in B$ such that $x\in B_1$ and $B_1\subseteq\unions{M}$.
        \begin{subproof}
            Fix $x$.
            Assume $x\in\unions{M}$.
            Take $U\in M$ such that $x\in U$ by \cref{unions_iff}.
            Then $U\in\basisTopology{B}{X}$ by \cref{subseteq}.
            Then there exists $B_1\in B$ such that $x\in B_1$ and $B_1\subseteq U$.
            Then $B_1\subseteq\unions{M}$.
        \end{subproof}
        Thus $\unions{M}\in\basisTopology{B}{X}$.
    \end{subproof}
    Show that for all $U,V$ such that $U\in\basisTopology{B}{X}$ and $V\in\basisTopology{B}{X}$ we have $U\inter V\in\basisTopology{B}{X}$.
    \begin{subproof}
        Fix $U$.
        Fix $V$.
        Assume $U\in\basisTopology{B}{X}$ and $V\in\basisTopology{B}{X}$.
        Then $U\in\pow{X}$ by \cref{subseteq}.
        Then $V\in\pow{X}$ by \cref{subseteq}.
        $U\subseteq X$ and $V\subseteq X$ by \cref{pow_iff}.
        Thus $U\inter V\subseteq X$ by \cref{inter_subseteq}.
        Hence $U\inter V\in\pow{X}$ by \cref{pow_iff}.
        Show that for all $x$ such that $x\in U\inter V$ there exists $B_3\in B$ such that $x\in B_3$ and $B_3\subseteq U\inter V$.
        \begin{subproof}
            Fix $x$.
            Assume $x\in U\inter V$.
            There exists $B_1\in B$ such that $x\in B_1$ and $B_1\subseteq U$.
            There exists $B_2\in B$ such that $x\in B_2$ and $B_2\subseteq V$.
            Then $x\in B_1\inter B_2$.
            There exists $B_3\in B$ such that $x\in B_3$ and $B_3\subseteq B_1\inter B_2$.
            Then $B_3\subseteq U\inter V$.
        \end{subproof}
        Thus $U\inter V\in\basisTopology{B}{X}$.
    \end{subproof}
    Thus $\basisTopology{B}{X}$ is closed under arbitrary unions.
    Thus $\basisTopology{B}{X}$ is closed under binary intersections.
    Follows by \cref{topology_on}.
\end{proof}

% from §13 helper lemma 1: basis elements are open in the generated topology
\begin{lemma}\label{basis_elem_open_in_generated}
    Assume $B$ is a basis on $X$.
    Suppose $B_1\in B$.
    Then $B_1\in\basisTopology{B}{X}$.
\end{lemma}
\begin{proof}
    $B\subseteq\pow{X}$ by \cref{basis_on}.
    Thus $B_1\in\pow{X}$ by \cref{subseteq}.
    Then $B_1\subseteq X$ by \cref{pow_iff}.
    Show that for all $x$ such that $x\in B_1$ there exists $B_2\in B$ such that $x\in B_2$ and $B_2\subseteq B_1$.
    \begin{subproof}
        Fix $x$.
        Assume $x\in B_1$.
        Let $B_2 = B_1$.
        $B_2\in B$.
        $B_2\subseteq B_1$ by \cref{subseteq_refl}.
    \end{subproof}
    Thus $B_1\in\basisTopology{B}{X}$.
\end{proof}

% from §13 Item 5: unions of basis elements
\begin{lemma}\label{basis_topology_unions}
    Assume $B$ is a basis for $T$ on $X$.
    Then $T = \{U\in\pow{X} \mid \exists M\subseteq B. U = \unions{M}\}$.
\end{lemma}
\begin{proof}
    $T = \basisTopology{B}{X}$ by \cref{basis_for_topology}.
    It suffices to show that for all $U$ we have
    $U\in\basisTopology{B}{X}$ iff there exists $M\subseteq B$ such that $U = \unions{M}$.
    Fix $U$.
        Show that if $U\in\basisTopology{B}{X}$ then there exists $M\subseteq B$ such that $U = \unions{M}$.
        \begin{subproof}
            Assume $U\in\basisTopology{B}{X}$.
            Let $M = \{B_1\in B \mid B_1\subseteq U\}$.
            Then $M\subseteq B$.
            Show that $U = \unions{M}$.
            \begin{subproof}
                Show that for all $x$ such that $x\in U$ we have $x\in\unions{M}$.
                \begin{subproof}
                    Fix $x$.
                    Assume $x\in U$.
                    Then there exists $B_1\in B$ such that $x\in B_1$ and $B_1\subseteq U$.
                    Then $B_1\in M$.
                    Thus $x\in\unions{M}$ by \cref{unions_iff}.
                \end{subproof}
                Thus $U\subseteq\unions{M}$ by \cref{subseteq}.
                Show that for all $x$ such that $x\in\unions{M}$ we have $x\in U$.
                \begin{subproof}
                    Fix $x$.
                    Assume $x\in\unions{M}$.
                    Take $B_1\in M$ such that $x\in B_1$ by \cref{unions_iff}.
                    Then $B_1\subseteq U$.
                    Thus $x\in U$.
                \end{subproof}
                Thus $\unions{M}\subseteq U$ by \cref{subseteq}.
                Thus $U = \unions{M}$ by \cref{subseteq_antisymmetric}.
            \end{subproof}
        \end{subproof}
        Show that if there exists $M\subseteq B$ such that $U = \unions{M}$ then $U\in\basisTopology{B}{X}$.
        \begin{subproof}
            Assume there exists $M\subseteq B$ such that $U = \unions{M}$.
            Take $M$ such that $M\subseteq B$ and $U = \unions{M}$.
            $\basisTopology{B}{X}$ is a topology on $X$ by \cref{basis_topology_is_topology,basis_for_topology}.
            Show that for all $B_1$ such that $B_1\in M$ we have $B_1\in\basisTopology{B}{X}$.
            \begin{subproof}
                Fix $B_1$.
                Assume $B_1\in M$.
                Then $B_1\in B$ by \cref{subseteq}.
                Then $B_1\in\basisTopology{B}{X}$ by \cref{basis_elem_open_in_generated,basis_for_topology}.
            \end{subproof}
            Thus $M\subseteq\basisTopology{B}{X}$ by \cref{subseteq}.
            Thus $\unions{M}\in\basisTopology{B}{X}$ by \cref{topology_on}.
            Thus $U\in\basisTopology{B}{X}$.
        \end{subproof}
\end{proof}

% from §13 Item 6: criterion for a collection of open sets to be a basis
\begin{lemma}\label{open_collection_basis}
    Assume $T$ is a topology on $X$.
    Assume for all $C_1\in C$ we have $C_1$ is open in $X$ with respect to $T$.
    Assume for all $U,x$ such that $U\in T$ and $x\in U$ there exists $C_1\in C$ such that $x\in C_1$ and $C_1\subseteq U$.
    Then $C$ is a basis for $T$ on $X$.
\end{lemma}
\begin{proof}
    Show that for all $x$ such that $x\in X$ there exists $C_1\in C$ such that $x\in C_1$ and $C_1\subseteq X$.
    \begin{subproof}
        Fix $x$.
        Assume $x\in X$.
        $X\in T$ by \cref{topology_on}.
        There exists $C_1\in C$ such that $x\in C_1$ and $C_1\subseteq X$.
    \end{subproof}
    Show that for all $C_1,C_2,x$ such that $C_1\in C$ and $C_2\in C$ and $x\in C_1\inter C_2$
        there exists $C_3\in C$ such that $x\in C_3$ and $C_3\subseteq C_1\inter C_2$.
    \begin{subproof}
        Fix $C_1$.
        Fix $C_2$.
        Fix $x$.
        Assume $C_1\in C$ and $C_2\in C$ and $x\in C_1\inter C_2$.
        $C_1$ is open in $X$ with respect to $T$.
        $C_2$ is open in $X$ with respect to $T$.
        $C_1\in T$ and $C_2\in T$ by \cref{open_in}.
        $C_1\inter C_2\in T$ by \cref{topology_on}.
        There exists $C_3\in C$ such that $x\in C_3$ and $C_3\subseteq C_1\inter C_2$.
    \end{subproof}
    Show that for all $C_1$ such that $C_1\in C$ we have $C_1\in\pow{X}$.
    \begin{subproof}
        Fix $C_1$.
        Assume $C_1\in C$.
        $C_1$ is open in $X$ with respect to $T$.
        Thus $C_1\in T$ by \cref{open_in}.
        $T\subseteq\pow{X}$ by \cref{topology_on}.
        Thus $C_1\in\pow{X}$ by \cref{subseteq}.
    \end{subproof}
    Thus $C\subseteq\pow{X}$ by \cref{subseteq}.
    Thus $C$ is a basis on $X$ by \cref{basis_on}.
    Show that for all $U$ such that $U\in T$ we have $U\in\basisTopology{C}{X}$.
    \begin{subproof}
        Fix $U$.
        Assume $U\in T$.
        $T\subseteq\pow{X}$ by \cref{topology_on}.
        Thus $U\in\pow{X}$ by \cref{subseteq}.
        For all $x$ such that $x\in U$ there exists $C_1\in C$ such that $x\in C_1$ and $C_1\subseteq U$.
        Thus $U\in\basisTopology{C}{X}$ by \cref{topology_generated_by_basis}.
    \end{subproof}
    Thus $T\subseteq\basisTopology{C}{X}$ by \cref{subseteq}.
    Show that for all $U$ such that $U\in\basisTopology{C}{X}$ we have $U\in T$.
    \begin{subproof}
        Fix $U$.
        Assume $U\in\basisTopology{C}{X}$.
        Let $M = \{C_1\in C \mid C_1\subseteq U\}$.
        Show that $U = \unions{M}$.
        \begin{subproof}
            Show that for all $x$ such that $x\in U$ we have $x\in\unions{M}$.
            \begin{subproof}
                Fix $x$.
                Assume $x\in U$.
                Then there exists $C_1\in C$ such that $x\in C_1$ and $C_1\subseteq U$.
                Then $C_1\in M$.
                Thus $x\in\unions{M}$ by \cref{unions_iff}.
            \end{subproof}
            Thus $U\subseteq\unions{M}$ by \cref{subseteq}.
            Show that for all $x$ such that $x\in\unions{M}$ we have $x\in U$.
            \begin{subproof}
                Fix $x$.
                Assume $x\in\unions{M}$.
                Take $C_1\in M$ such that $x\in C_1$ by \cref{unions_iff}.
                Then $C_1\subseteq U$.
                Thus $x\in U$.
            \end{subproof}
            Thus $\unions{M}\subseteq U$ by \cref{subseteq}.
            Thus $U = \unions{M}$ by \cref{subseteq_antisymmetric}.
        \end{subproof}
        Show that for all $C_1$ such that $C_1\in M$ we have $C_1\in T$.
        \begin{subproof}
            Fix $C_1$.
            Assume $C_1\in M$.
            Then $C_1\in C$ by \cref{subseteq}.
            $C_1$ is open in $X$ with respect to $T$.
            Thus $C_1\in T$ by \cref{open_in}.
        \end{subproof}
        Thus $M\subseteq T$ by \cref{subseteq}.
        Thus $\unions{M}\in T$ by \cref{topology_on}.
        Thus $U\in T$.
    \end{subproof}
    Thus $\basisTopology{C}{X}\subseteq T$ by \cref{subseteq}.
    Thus $T = \basisTopology{C}{X}$ by \cref{subseteq_antisymmetric}.
    Thus $C$ is a basis for $T$ on $X$ by \cref{basis_for_topology}.
\end{proof}

% from §13 Item 7: bases and finer topologies
\begin{lemma}\label{basis_finer_criterion}
    Assume $B$ is a basis for $T$ on $X$.
    Assume $B_1$ is a basis for $T_1$ on $X$.
    Then $T_1$ is finer than $T$ on $X$ iff
    for all $x,B_2$ such that $x\in X$ and $B_2\in B$ and $x\in B_2$
    there exists $B_3\in B_1$ such that $x\in B_3$ and $B_3\subseteq B_2$.
\end{lemma}
\begin{proof}
    Show that if for all $x,B_2$ such that $x\in X$ and $B_2\in B$ and $x\in B_2$
        there exists $B_3\in B_1$ such that $x\in B_3$ and $B_3\subseteq B_2$
        then $T_1$ is finer than $T$ on $X$.
    \begin{subproof}
        Assume for all $x,B_2$ such that $x\in X$ and $B_2\in B$ and $x\in B_2$
            there exists $B_3\in B_1$ such that $x\in B_3$ and $B_3\subseteq B_2$.
        Show that for all $U$ such that $U\in T$ we have $U\in T_1$.
        \begin{subproof}
            Fix $U$.
            Assume $U\in T$.
            For all $x$ such that $x\in U$ there exists $B_2\in B$ such that $x\in B_2$ and $B_2\subseteq U$
                by \cref{basis_for_topology,topology_generated_by_basis}.
            Show that for all $x$ such that $x\in U$ there exists $B_3\in B_1$ such that $x\in B_3$ and $B_3\subseteq U$.
            \begin{subproof}
                Fix $x$.
                Assume $x\in U$.
                There exists $B_2\in B$ such that $x\in B_2$ and $B_2\subseteq U$
                    by \cref{basis_for_topology,topology_generated_by_basis}.
                Take $B_2$ such that $B_2\in B$ and $x\in B_2$ and $B_2\subseteq U$.
                $B\subseteq\pow{X}$ by \cref{basis_for_topology,basis_on}.
                $B_2\in\pow{X}$ by \cref{subseteq}.
                $B_2\subseteq X$ by \cref{powerset_elim}.
                $x\in X$ by \cref{subseteq}.
                There exists $B_3\in B_1$ such that $x\in B_3$ and $B_3\subseteq B_2$ by assumption.
                $B_3\subseteq U$ by \cref{subseteq_transitive}.
            \end{subproof}
            $U\in\pow{X}$ by \cref{basis_for_topology,topology_generated_by_basis}.
            Thus $U\in\basisTopology{B_1}{X}$ by \cref{topology_generated_by_basis}.
            Thus $U\in T_1$ by \cref{basis_for_topology}.
        \end{subproof}
        Thus $T\subseteq T_1$ by \cref{subseteq}.
        $T$ is a topology on $X$ by \cref{basis_for_topology,basis_topology_is_topology}.
        $T_1$ is a topology on $X$ by \cref{basis_for_topology,basis_topology_is_topology}.
        Thus $T_1$ is finer than $T$ on $X$ by \cref{finer_topology}.
    \end{subproof}
    Show that if $T_1$ is finer than $T$ on $X$ then
        for all $x,B_2$ such that $x\in X$ and $B_2\in B$ and $x\in B_2$
        there exists $B_3\in B_1$ such that $x\in B_3$ and $B_3\subseteq B_2$.
    \begin{subproof}
        Assume $T_1$ is finer than $T$ on $X$.
        Show that for all $x,B_2$ such that $x\in X$ and $B_2\in B$ and $x\in B_2$
            there exists $B_3\in B_1$ such that $x\in B_3$ and $B_3\subseteq B_2$.
        \begin{subproof}
            Fix $x$.
            Fix $B_2$.
            Assume $x\in X$ and $B_2\in B$ and $x\in B_2$.
            $B_2\in T$ by \cref{basis_elem_open_in_generated,basis_for_topology}.
            $T\subseteq T_1$ by \cref{finer_topology}.
            Thus $B_2\in T_1$ by \cref{subseteq}.
            There exists $B_3\in B_1$ such that $x\in B_3$ and $B_3\subseteq B_2$
                by \cref{basis_for_topology,topology_generated_by_basis}.
        \end{subproof}
    \end{subproof}
\end{proof}

% from §13 Item 8: half-open interval on R
\begin{definition}\label{interval_closed_open}
    $\intervalclosedopen{a}{b} = \{x\in\reals \mid a \leq x \land x < b\}$.
\end{definition}

% from §13 Item 9: standard basis on R
\begin{definition}\label{standard_basis_reals}
    $\standardBasisR = \{U\in\pow{\reals} \mid \exists a,b\in\reals. a < b \land U = \intervalopen{a}{b}\}$.
\end{definition}

% from §13 Item 10: standard topology on R
\begin{definition}\label{standard_topology_reals}
    $\standardTopologyR = \basisTopology{\standardBasisR}{\reals}$.
\end{definition}

% from §13 Item 11: lower limit basis on R
\begin{definition}\label{lower_limit_basis_reals}
    $\lowerLimitBasisR = \{U\in\pow{\reals} \mid \exists a,b\in\reals. a < b \land U = \intervalclosedopen{a}{b}\}$.
\end{definition}

% from §13 Item 12: lower limit topology on R
\begin{definition}\label{lower_limit_topology_reals}
    $\lowerLimitTopologyR = \basisTopology{\lowerLimitBasisR}{\reals}$.
\end{definition}

% from §13 Item 14: K set
\begin{definition}\label{k_set}
    $\kset = \{x\in\reals \mid \exists n\in\naturalsPlus. x = 1 / n\}$.
\end{definition}

% from §13 Item 15: K basis on R
\begin{definition}\label{k_basis_reals}
    $\kBasisR = \{U\in\pow{\reals} \mid \exists a,b\in\reals. a < b \land (U = \intervalopen{a}{b} \lor U = \intervalopen{a}{b}\setminus \kset)\}$.
\end{definition}

% from §13 Item 16: K-topology on R
\begin{definition}\label{k_topology_reals}
    $\kTopologyR = \basisTopology{\kBasisR}{\reals}$.
\end{definition}

% from §13 helper lemma 2: split a weak inequality
\begin{lemma}\label{real_leq_split}
    Suppose $a \leq b$.
    Then $a < b$ or $a = b$.
\end{lemma}
\begin{proof}
    Thus $a < b$ or $a = b$.
\end{proof}

% from §13 helper lemma 3: substitute equality on the left
\begin{lemma}\label{real_lt_subst_left}
    Suppose $a = b$ and $b < c$.
    Then $a < c$.
\end{lemma}
\begin{proof}
    Thus $a < c$.
\end{proof}

% from §13 helper lemma 4: substitute equality on the right
\begin{lemma}\label{real_lt_subst_right}
    Suppose $a < b$ and $b = c$.
    Then $a < c$.
\end{lemma}
\begin{proof}
    Thus $a < c$.
\end{proof}

% from §13 helper lemma 5: negative zero
\begin{lemma}\label{real_neg_zero}
    $\neg{0} = 0$.
\end{lemma}
\begin{proof}
    $0\in\reals$ by \cref{real_add_ident}.
    $0 + 0 = 0$ by \cref{real_add_ident}.
    Thus $0 = \neg{0}$ by \cref{real_add_inverse_unique}.
    Thus $\neg{0} = 0$.
\end{proof}

% from §13 helper lemma 6: refine open intervals
\begin{lemma}\label{open_interval_refinement}
    Suppose $a,b,c,d,x\in\reals$.
    Suppose $a < x$ and $x < b$ and $c < x$ and $x < d$.
    Then there exist $e,f\in\reals$ such that $e < x$ and $x < f$ and
    $\intervalopen{e}{f} \subseteq \intervalopen{a}{b}\inter \intervalopen{c}{d}$.
\end{lemma}
\begin{proof}
    Let $A = \{a,c\}$.
    $A$ is finite by \cref{pair_is_finite}.
    $a\in A$ by \cref{upair_intro_left}.
    Thus $A\neq\emptyset$ by \cref{emptyset}.
    Show that for all $y$ such that $y\in A$ we have $y\in\reals$.
    \begin{subproof}
        Fix $y$.
        Assume $y\in A$.
        $y = a$ or $y = c$ by \cref{upair_elim}.
        Thus $y\in\reals$.
    \end{subproof}
    Thus $A\subseteq\reals$ by \cref{subseteq}.
    Take $m\in A$ such that $m$ is a maximum of $A$ by \cref{finite_real_maximum}.
    Let $B = \{b,d\}$.
    $B$ is finite by \cref{pair_is_finite}.
    $b\in B$ by \cref{upair_intro_left}.
    Thus $B\neq\emptyset$ by \cref{emptyset}.
    Show that for all $y$ such that $y\in B$ we have $y\in\reals$.
    \begin{subproof}
        Fix $y$.
        Assume $y\in B$.
        $y = b$ or $y = d$ by \cref{upair_elim}.
        Thus $y\in\reals$.
    \end{subproof}
    Thus $B\subseteq\reals$ by \cref{subseteq}.
    Take $n\in B$ such that $n$ is a minimum of $B$ by \cref{finite_real_minimum}.
    $m\in\reals$ by \cref{subseteq}.
    $n\in\reals$ by \cref{subseteq}.
    Show that $m < x$.
    \begin{subproof}
        $m = a$ or $m = c$ by \cref{upair_elim}.
        \begin{byCase}
        \caseOf{$m = a$.} Thus $m < x$.
        \caseOf{$m = c$.} Thus $m < x$.
        \end{byCase}
    \end{subproof}
    Show that $x < n$.
    \begin{subproof}
        $n = b$ or $n = d$ by \cref{upair_elim}.
        \begin{byCase}
        \caseOf{$n = b$.} Thus $x < n$.
        \caseOf{$n = d$.} Thus $x < n$.
        \end{byCase}
    \end{subproof}
    Let $e = m$.
    Let $f = n$.
    $e\in\reals$ and $f\in\reals$.
    $e < x$ and $x < f$.
    Show that for all $y$ such that $y\in\intervalopen{e}{f}$ we have
        $y\in \intervalopen{a}{b}\inter \intervalopen{c}{d}$.
    \begin{subproof}
        Fix $y$.
        Assume $y\in\intervalopen{e}{f}$.
        $e < y$ and $y < f$ by \cref{intervalopen}.
        $y\in\reals$ by \cref{intervalopen}.
        $e = m$ and $f = n$.
        Thus $m < y$ and $y < n$ by \cref{intervalopen}.
        $m$ is a maximum of $A$.
        Thus for all $t\in A$ we have $t \leq m$ by \cref{real_maximum}.
        $a\in A$ by \cref{upair_intro_left}.
        $c\in A$ by \cref{upair_intro_right}.
        Thus $a \leq m$ and $c \leq m$.
        $n$ is a minimum of $B$.
        Thus for all $t\in B$ we have $n \leq t$ by \cref{real_minimum}.
        $b\in B$ by \cref{upair_intro_left}.
        $d\in B$ by \cref{upair_intro_right}.
        Thus $n \leq b$ and $n \leq d$.
        Show that $a < y$.
        \begin{subproof}
            $a < m$ or $a = m$ by \cref{real_leq_split}.
            \begin{byCase}
            \caseOf{$a = m$.} Thus $a < y$ by \cref{real_lt_subst_left}.
            \caseOf{$a < m$.}
                Thus $a < y$ by \cref{real_lt_trans}.
            \end{byCase}
        \end{subproof}
        Show that $c < y$.
        \begin{subproof}
            $c < m$ or $c = m$ by \cref{real_leq_split}.
            \begin{byCase}
            \caseOf{$c = m$.} Thus $c < y$ by \cref{real_lt_subst_left}.
            \caseOf{$c < m$.}
                Thus $c < y$ by \cref{real_lt_trans}.
            \end{byCase}
        \end{subproof}
        Show that $y < b$.
        \begin{subproof}
            $n < b$ or $n = b$ by \cref{real_leq_split}.
            \begin{byCase}
            \caseOf{$n = b$.} Thus $y < b$ by \cref{real_lt_subst_right}.
            \caseOf{$n < b$.}
                Thus $y < b$ by \cref{real_lt_trans}.
            \end{byCase}
        \end{subproof}
        Show that $y < d$.
        \begin{subproof}
            $n < d$ or $n = d$ by \cref{real_leq_split}.
            \begin{byCase}
            \caseOf{$n = d$.} Thus $y < d$ by \cref{real_lt_subst_right}.
            \caseOf{$n < d$.}
                Thus $y < d$ by \cref{real_lt_trans}.
            \end{byCase}
        \end{subproof}
        Thus $y\in\intervalopen{a}{b}$ by \cref{intervalopen}.
        Thus $y\in\intervalopen{c}{d}$ by \cref{intervalopen}.
        Thus $y\in\intervalopen{a}{b}\inter\intervalopen{c}{d}$ by \cref{inter_intro}.
    \end{subproof}
    Thus $\intervalopen{e}{f}\subseteq \intervalopen{a}{b}\inter \intervalopen{c}{d}$ by \cref{subseteq}.
\end{proof}

% from §13 helper lemma 7: refine half-open intervals
\begin{lemma}\label{closedopen_interval_refinement}
    Suppose $a,b,c,d,x\in\reals$.
    Suppose $a \leq x$ and $x < b$ and $c \leq x$ and $x < d$.
    Then there exist $e,f\in\reals$ such that $e \leq x$ and $x < f$ and
    $\intervalclosedopen{e}{f} \subseteq \intervalclosedopen{a}{b}\inter \intervalclosedopen{c}{d}$.
\end{lemma}
\begin{proof}
    Let $A = \{a,c\}$.
    $A$ is finite by \cref{pair_is_finite}.
    $a\in A$ by \cref{upair_intro_left}.
    Thus $A\neq\emptyset$ by \cref{emptyset}.
    Show that for all $y$ such that $y\in A$ we have $y\in\reals$.
    \begin{subproof}
        Fix $y$.
        Assume $y\in A$.
        $y = a$ or $y = c$ by \cref{upair_elim}.
        Thus $y\in\reals$.
    \end{subproof}
    Thus $A\subseteq\reals$ by \cref{subseteq}.
    Take $m\in A$ such that $m$ is a maximum of $A$ by \cref{finite_real_maximum}.
    Let $B = \{b,d\}$.
    $B$ is finite by \cref{pair_is_finite}.
    $b\in B$ by \cref{upair_intro_left}.
    Thus $B\neq\emptyset$ by \cref{emptyset}.
    Show that for all $y$ such that $y\in B$ we have $y\in\reals$.
    \begin{subproof}
        Fix $y$.
        Assume $y\in B$.
        $y = b$ or $y = d$ by \cref{upair_elim}.
        Thus $y\in\reals$.
    \end{subproof}
    Thus $B\subseteq\reals$ by \cref{subseteq}.
    Take $n\in B$ such that $n$ is a minimum of $B$ by \cref{finite_real_minimum}.
    $m\in\reals$ by \cref{subseteq}.
    $n\in\reals$ by \cref{subseteq}.
    Show that $m \leq x$.
    \begin{subproof}
        $m = a$ or $m = c$ by \cref{upair_elim}.
        \begin{byCase}
        \caseOf{$m = a$.} Thus $m \leq x$.
        \caseOf{$m = c$.} Thus $m \leq x$.
        \end{byCase}
    \end{subproof}
    Show that $x < n$.
    \begin{subproof}
        $n = b$ or $n = d$ by \cref{upair_elim}.
        \begin{byCase}
        \caseOf{$n = b$.} Thus $x < n$.
        \caseOf{$n = d$.} Thus $x < n$.
        \end{byCase}
    \end{subproof}
    Let $e = m$.
    Let $f = n$.
    $e\in\reals$ and $f\in\reals$.
    $e \leq x$ and $x < f$.
    Show that for all $y$ such that $y\in\intervalclosedopen{e}{f}$ we have
        $y\in \intervalclosedopen{a}{b}\inter \intervalclosedopen{c}{d}$.
    \begin{subproof}
        Fix $y$.
        Assume $y\in\intervalclosedopen{e}{f}$.
        $e \leq y$ and $y < f$ by \cref{interval_closed_open}.
        $y\in\reals$ by \cref{interval_closed_open}.
        $m$ is a maximum of $A$.
        Thus for all $t\in A$ we have $t \leq m$ by \cref{real_maximum}.
        $a\in A$ by \cref{upair_intro_left}.
        $c\in A$ by \cref{upair_intro_right}.
        Thus $a \leq m$ and $c \leq m$.
        $n$ is a minimum of $B$.
        Thus for all $t\in B$ we have $n \leq t$ by \cref{real_minimum}.
        $b\in B$ by \cref{upair_intro_left}.
        $d\in B$ by \cref{upair_intro_right}.
        Thus $n \leq b$ and $n \leq d$.
        $a \leq y$ by \cref{real_leq_trans}.
        $c \leq y$ by \cref{real_leq_trans}.
        Show that $y < b$.
        \begin{subproof}
            $n < b$ or $n = b$ by \cref{real_leq_split}.
            \begin{byCase}
            \caseOf{$n = b$.} Thus $y < b$ by \cref{real_lt_subst_right}.
            \caseOf{$n < b$.}
                Thus $y < b$ by \cref{real_lt_trans}.
            \end{byCase}
        \end{subproof}
        Show that $y < d$.
        \begin{subproof}
            $n < d$ or $n = d$ by \cref{real_leq_split}.
            \begin{byCase}
            \caseOf{$n = d$.} Thus $y < d$ by \cref{real_lt_subst_right}.
            \caseOf{$n < d$.}
                Thus $y < d$ by \cref{real_lt_trans}.
            \end{byCase}
        \end{subproof}
        Thus $y\in\intervalclosedopen{a}{b}$ by \cref{interval_closed_open}.
        Thus $y\in\intervalclosedopen{c}{d}$ by \cref{interval_closed_open}.
        Thus $y\in\intervalclosedopen{a}{b}\inter\intervalclosedopen{c}{d}$ by \cref{inter_intro}.
    \end{subproof}
    Thus $\intervalclosedopen{e}{f}\subseteq \intervalclosedopen{a}{b}\inter \intervalclosedopen{c}{d}$ by \cref{subseteq}.
\end{proof}

% from §13 helper lemma 8: open intervals are subsets of R
\begin{lemma}\label{intervalopen_subset_reals}
    Suppose $a\in\reals$ and $b\in\reals$.
    Then $\intervalopen{a}{b}\subseteq\reals$.
\end{lemma}
\begin{proof}
    Show that for all $x$ such that $x\in\intervalopen{a}{b}$ we have $x\in\reals$.
    \begin{subproof}
        Fix $x$.
        Assume $x\in\intervalopen{a}{b}$.
        Thus $x\in\reals$ by \cref{intervalopen}.
    \end{subproof}
    Thus $\intervalopen{a}{b}\subseteq\reals$ by \cref{subseteq}.
\end{proof}

% from §13 helper lemma 9: half-open intervals are subsets of R
\begin{lemma}\label{intervalclosedopen_subset_reals}
    Suppose $a\in\reals$ and $b\in\reals$.
    Then $\intervalclosedopen{a}{b}\subseteq\reals$.
\end{lemma}
\begin{proof}
    Show that for all $x$ such that $x\in\intervalclosedopen{a}{b}$ we have $x\in\reals$.
    \begin{subproof}
        Fix $x$.
        Assume $x\in\intervalclosedopen{a}{b}$.
        Thus $x\in\reals$ by \cref{interval_closed_open}.
    \end{subproof}
    Thus $\intervalclosedopen{a}{b}\subseteq\reals$ by \cref{subseteq}.
\end{proof}

% from §13 helper lemma 10: K-set elements are positive
\begin{lemma}\label{k_set_positive}
    Suppose $x\in\kset$.
    Then $0 < x$.
\end{lemma}
\begin{proof}
    Take $n\in\naturalsPlus$ such that $x = 1 / n$ by \cref{k_set}.
    $1\in\naturalsPlus$ by \cref{real_one_in_naturals}.
    Thus $1 \leq n$ by \cref{real_one_leq_naturalsplus}.
    $0 < 1$ by \cref{real_zero_lt_one}.
    $0\in\reals$ by \cref{real_add_ident}.
    $1\in\reals$ by \cref{real_naturals_subset_reals,real_one_in_naturals,subseteq}.
    $n\in\reals$ by \cref{real_naturals_subset_reals,subseteq}.
    Show that $0 < n$.
    \begin{subproof}
        $1 < n$ or $1 = n$ by \cref{real_leq_split}.
        \begin{byCase}
        \caseOf{$1 = n$.} Thus $0 < n$ by \cref{real_lt_subst_right}.
        \caseOf{$1 < n$.} Thus $0 < n$ by \cref{real_lt_trans}.
        \end{byCase}
    \end{subproof}
    Thus $n$ is positive.
    Thus $0 < \inv{n}$ by \cref{real_inv_pos}.
    Thus $n \neq 0$ by \cref{real_pos_nonzero}.
    Thus $\inv{n}\in\reals$ by \cref{real_mul_inverse}.
    $1 / n = 1 \rmul \inv{n}$ by \cref{real_div}.
    $1 \rmul \inv{n} = \inv{n}$ by \cref{real_mul_ident}.
    Thus $\inv{n} = 1 / n$.
    Thus $0 < 1 / n$ by \cref{real_lt_subst_right}.
    Thus $0 < x$.
\end{proof}

% from §13 helper lemma 11: rational between a negative real and zero
\begin{lemma}\label{rational_between_negative_and_zero}
    Suppose $a\in\reals$ and $a < 0$.
    Then there exists $p\in\rationals$ such that $p\in\intervalopen{a}{0}$.
\end{lemma}
\begin{proof}
    $0\in\reals$ by \cref{real_add_ident}.
    $\neg{a}\in\reals$ by \cref{real_add_inverse}.
    Thus $\neg{0} < \neg{a}$ by \cref{real_lt_neg}.
    Thus $\neg{0} = 0$ by \cref{real_neg_zero}.
    Thus $0 < \neg{a}$ by \cref{real_lt_subst_left}.
    Thus $\neg{a}$ is positive.
    Take $N\in\naturalsPlus$ such that $0 < 1 / N$ and $1 / N < \neg{a}$ by \cref{real_small_reciprocal}.
    $N\in\reals$ by \cref{real_naturals_subset_reals,subseteq}.
    Show that $0 < N$.
    \begin{subproof}
        $1\in\reals$ by \cref{real_mul_ident}.
        $0\in\reals$ by \cref{real_add_ident}.
        $0 < 1$ by \cref{real_zero_lt_one}.
        $1 < N$ or $1 = N$ by \cref{real_one_leq_naturalsplus}.
        \begin{byCase}
        \caseOf{$1 = N$.} Thus $0 < N$.
        \caseOf{$1 < N$.} Thus $0 < N$ by \cref{real_lt_trans}.
        \end{byCase}
    \end{subproof}
    Thus $N\neq 0$ by \cref{real_pos_nonzero}.
    Thus $\inv{N}\in\reals$ by \cref{real_mul_inverse}.
    $1\in\reals$ by \cref{real_mul_ident}.
    $1 / N = 1 \rmul \inv{N}$ by \cref{real_div}.
    Thus $1 / N\in\reals$ by \cref{real_mul_closed}.
    Let $p = \neg{(1 / N)}$.
    $p\in\reals$ by \cref{real_add_inverse}.
    Thus $\neg{(1 / N)} < \neg{0}$ by \cref{real_lt_neg}.
    Thus $\neg{(1 / N)} < 0$ by \cref{real_neg_zero,real_lt_subst_right}.
    Thus $p < 0$ by \cref{real_lt_subst_left}.
    Thus $\neg{(\neg{a})} < \neg{(1 / N)}$ by \cref{real_lt_neg}.
    $\neg{(\neg{a})} = a$ by \cref{real_neg_neg}.
    Thus $a < \neg{(1 / N)}$ by \cref{real_lt_subst_left}.
    Thus $a < p$ by \cref{real_lt_subst_right}.
    Thus $p\in\intervalopen{a}{0}$ by \cref{intervalopen}.
    Show that $p\in\rationals$.
    \begin{subproof}
        $1\in\naturalsPlus$ by \cref{real_one_in_naturals}.
        $1\in\reals$ by \cref{real_mul_ident}.
        $\neg{1}\in\reals$ by \cref{real_add_inverse}.
        $1 + \neg{1} = 0$ by \cref{real_add_inverse}.
        $\neg{1} + 1 = 0$ by \cref{real_add_comm}.
        Thus $\neg{1}\in\integers$ by \cref{real_integers}.
        $N\in\integers$ by \cref{real_integers}.
        Show that $N\notin\{0\}$.
        \begin{subproof}
            Suppose not.
            Thus $N\in\{0\}$.
            Thus $N = 0$ by \cref{singleton_iff}.
            Contradiction.
        \end{subproof}
        Thus $N\in\integers\setminus\{0\}$ by \cref{setminus_intro}.
        $\inv{N} \rmul \neg{1} = \neg{(\inv{N} \rmul 1)}$ by \cref{real_mul_neg_right}.
        $\inv{N} \rmul 1 = 1 \rmul \inv{N}$ by \cref{real_mul_comm}.
        Thus $\inv{N} \rmul \neg{1} = \neg{(1 \rmul \inv{N})}$.
        $\neg{1} \rmul \inv{N} = \inv{N} \rmul \neg{1}$ by \cref{real_mul_comm}.
        Thus $\neg{1} \rmul \inv{N} = \neg{(1 \rmul \inv{N})}$.
        Thus $p = \neg{(1 \rmul \inv{N})}$.
        Thus $p = \neg{1} \rmul \inv{N}$.
        Thus $p = (\neg{1}) / N$ by \cref{real_div}.
        Thus $p\in\rationals$ by \cref{real_rationals}.
    \end{subproof}
\end{proof}

% from §13 Item 17: standard basis is a basis on R
\begin{lemma}\label{standard_basis_is_basis}
    $\standardBasisR$ is a basis on $\reals$.
\end{lemma}
\begin{proof}
    Show that for all $U$ such that $U\in\standardBasisR$ we have $U\in\pow{\reals}$.
    \begin{subproof}
        Fix $U$.
        Assume $U\in\standardBasisR$.
        Thus $U\in\pow{\reals}$ by \cref{standard_basis_reals}.
    \end{subproof}
    Thus $\standardBasisR\subseteq\pow{\reals}$ by \cref{subseteq}.
    Show that for all $x$ such that $x\in\reals$ there exists $B_1\in\standardBasisR$ such that $x\in B_1$.
    \begin{subproof}
        Fix $x$.
        Assume $x\in\reals$.
        $1\in\reals$ by \cref{real_naturals_subset_reals,real_one_in_naturals,subseteq}.
        $\neg{1}\in\reals$ by \cref{real_add_inverse}.
        Let $a = x + \neg{1}$.
        Let $b = x + 1$.
        $a\in\reals$ and $b\in\reals$ by \cref{real_add_closed}.
        $0 < 1$ by \cref{real_zero_lt_one}.
        $0\in\reals$ by \cref{real_add_ident}.
        Thus $\neg{1} < \neg{0}$ by \cref{real_lt_neg}.
        Thus $\neg{0} = 0$ by \cref{real_neg_zero}.
        Thus $\neg{1} < 0$ by \cref{real_lt_subst_right}.
        Thus $\neg{1} + x < 0 + x$ by \cref{real_lt_add}.
        $\neg{1} + x = x + \neg{1}$ by \cref{real_add_comm}.
        $0 + x = x + 0$ by \cref{real_add_comm}.
        $x + 0 = x$ by \cref{real_add_ident,real_add_comm}.
        Thus $a < x$.
        Thus $0 + x < 1 + x$ by \cref{real_lt_add}.
        $0 + x = x$ by \cref{real_add_ident,real_add_comm}.
        $1 + x = x + 1$ by \cref{real_add_comm}.
        Thus $x < b$.
        Thus $a < b$ by \cref{real_lt_trans}.
        Thus $x\in\intervalopen{a}{b}$ by \cref{intervalopen}.
        $\intervalopen{a}{b}\subseteq\reals$ by \cref{intervalopen_subset_reals}.
        Thus $\intervalopen{a}{b}\in\pow{\reals}$ by \cref{powerset_intro}.
        Thus $\intervalopen{a}{b}\in\standardBasisR$ by \cref{standard_basis_reals}.
    \end{subproof}
    Show that for all $B_1,B_2,x$ such that $B_1\in\standardBasisR$ and $B_2\in\standardBasisR$ and $x\in B_1\inter B_2$
        there exists $B_3\in\standardBasisR$ such that $x\in B_3$ and $B_3\subseteq B_1\inter B_2$.
    \begin{subproof}
        Fix $B_1$.
        Fix $B_2$.
        Fix $x$.
        Assume $B_1\in\standardBasisR$ and $B_2\in\standardBasisR$ and $x\in B_1\inter B_2$.
        Take $a,b\in\reals$ such that $a < b$ and $B_1 = \intervalopen{a}{b}$ by \cref{standard_basis_reals}.
        Take $c,d\in\reals$ such that $c < d$ and $B_2 = \intervalopen{c}{d}$ by \cref{standard_basis_reals}.
        $x\in B_1$ by \cref{inter_elim_left}.
        $x\in B_2$ by \cref{inter_elim_right}.
        Thus $a < x$ and $x < b$ by \cref{intervalopen}.
        Thus $x\in\reals$ by \cref{intervalopen}.
        Thus $c < x$ and $x < d$ by \cref{intervalopen}.
        Take $e,f\in\reals$ such that $e < x$ and $x < f$ and
            $\intervalopen{e}{f} \subseteq \intervalopen{a}{b}\inter \intervalopen{c}{d}$
            by \cref{open_interval_refinement}.
        Thus $e < f$ by \cref{real_lt_trans}.
        Let $B_3 = \intervalopen{e}{f}$.
        $B_3\subseteq\reals$ by \cref{intervalopen_subset_reals}.
        Thus $B_3\in\pow{\reals}$ by \cref{powerset_intro}.
        Thus $B_3\in\standardBasisR$ by \cref{standard_basis_reals}.
        Thus $x\in B_3$ by \cref{intervalopen}.
        Thus $B_3\subseteq B_1\inter B_2$.
    \end{subproof}
    Thus $\standardBasisR$ is a basis on $\reals$ by \cref{basis_on}.
\end{proof}

% from §13 Item 18: lower limit basis is a basis on R
\begin{lemma}\label{lower_limit_basis_is_basis}
    $\lowerLimitBasisR$ is a basis on $\reals$.
\end{lemma}
\begin{proof}
    Show that for all $U$ such that $U\in\lowerLimitBasisR$ we have $U\in\pow{\reals}$.
    \begin{subproof}
        Fix $U$.
        Assume $U\in\lowerLimitBasisR$.
        Thus $U\in\pow{\reals}$ by \cref{lower_limit_basis_reals}.
    \end{subproof}
    Thus $\lowerLimitBasisR\subseteq\pow{\reals}$ by \cref{subseteq}.
    Show that for all $x$ such that $x\in\reals$ there exists $B_1\in\lowerLimitBasisR$ such that $x\in B_1$.
    \begin{subproof}
        Fix $x$.
        Assume $x\in\reals$.
        $1\in\reals$ by \cref{real_naturals_subset_reals,real_one_in_naturals,subseteq}.
        Let $a = x$.
        Let $b = x + 1$.
        $b\in\reals$ by \cref{real_add_closed}.
        $0 < 1$ by \cref{real_zero_lt_one}.
        $0\in\reals$ by \cref{real_add_ident}.
        Thus $0 + x < 1 + x$ by \cref{real_lt_add}.
        $0 + x = x$ by \cref{real_add_ident,real_add_comm}.
        $1 + x = x + 1$ by \cref{real_add_comm}.
        Thus $x < b$.
        Thus $a < b$ by \cref{real_lt_subst_left}.
        Thus $x\in\intervalclosedopen{a}{b}$ by \cref{interval_closed_open}.
        $\intervalclosedopen{a}{b}\subseteq\reals$ by \cref{intervalclosedopen_subset_reals}.
        Thus $\intervalclosedopen{a}{b}\in\pow{\reals}$ by \cref{powerset_intro}.
        Thus $\intervalclosedopen{a}{b}\in\lowerLimitBasisR$ by \cref{lower_limit_basis_reals}.
    \end{subproof}
    Show that for all $B_1,B_2,x$ such that $B_1\in\lowerLimitBasisR$ and $B_2\in\lowerLimitBasisR$ and $x\in B_1\inter B_2$
        there exists $B_3\in\lowerLimitBasisR$ such that $x\in B_3$ and $B_3\subseteq B_1\inter B_2$.
    \begin{subproof}
        Fix $B_1$.
        Fix $B_2$.
        Fix $x$.
        Assume $B_1\in\lowerLimitBasisR$ and $B_2\in\lowerLimitBasisR$ and $x\in B_1\inter B_2$.
        Take $a,b\in\reals$ such that $a < b$ and $B_1 = \intervalclosedopen{a}{b}$ by \cref{lower_limit_basis_reals}.
        Take $c,d\in\reals$ such that $c < d$ and $B_2 = \intervalclosedopen{c}{d}$ by \cref{lower_limit_basis_reals}.
        $x\in B_1$ by \cref{inter_elim_left}.
        $x\in B_2$ by \cref{inter_elim_right}.
        Thus $a \leq x$ and $x < b$ by \cref{interval_closed_open}.
        Thus $x\in\reals$ by \cref{interval_closed_open}.
        Thus $c \leq x$ and $x < d$ by \cref{interval_closed_open}.
        Take $e,f\in\reals$ such that $e \leq x$ and $x < f$ and
            $\intervalclosedopen{e}{f} \subseteq \intervalclosedopen{a}{b}\inter \intervalclosedopen{c}{d}$
            by \cref{closedopen_interval_refinement}.
        Show that $e < f$.
        \begin{subproof}
            $e < x$ or $e = x$ by \cref{real_leq_split}.
            \begin{byCase}
            \caseOf{$e = x$.} Thus $e < f$ by \cref{real_lt_subst_left}.
            \caseOf{$e < x$.} Thus $e < f$ by \cref{real_lt_trans}.
            \end{byCase}
        \end{subproof}
        Let $B_3 = \intervalclosedopen{e}{f}$.
        $B_3\subseteq\reals$ by \cref{intervalclosedopen_subset_reals}.
        Thus $B_3\in\pow{\reals}$ by \cref{powerset_intro}.
        Thus $B_3\in\lowerLimitBasisR$ by \cref{lower_limit_basis_reals}.
        Thus $x\in B_3$ by \cref{interval_closed_open}.
        Thus $B_3\subseteq B_1\inter B_2$.
    \end{subproof}
    Thus $\lowerLimitBasisR$ is a basis on $\reals$ by \cref{basis_on}.
\end{proof}

% from §13 Item 19: K basis is a basis on R
\begin{lemma}\label{k_basis_is_basis}
    $\kBasisR$ is a basis on $\reals$.
\end{lemma}
\begin{proof}
    Show that for all $U$ such that $U\in\kBasisR$ we have $U\in\pow{\reals}$.
    \begin{subproof}
        Fix $U$.
        Assume $U\in\kBasisR$.
        Thus $U\in\pow{\reals}$ by \cref{k_basis_reals}.
    \end{subproof}
    Thus $\kBasisR\subseteq\pow{\reals}$ by \cref{subseteq}.
    Show that for all $x$ such that $x\in\reals$ there exists $B_1\in\kBasisR$ such that $x\in B_1$.
    \begin{subproof}
        Fix $x$.
        Assume $x\in\reals$.
        $1\in\reals$ by \cref{real_naturals_subset_reals,real_one_in_naturals,subseteq}.
        $\neg{1}\in\reals$ by \cref{real_add_inverse}.
        Let $a = x + \neg{1}$.
        Let $b = x + 1$.
        $a\in\reals$ and $b\in\reals$ by \cref{real_add_closed}.
        $0 < 1$ by \cref{real_zero_lt_one}.
        $0\in\reals$ by \cref{real_add_ident}.
        Thus $\neg{1} < \neg{0}$ by \cref{real_lt_neg}.
        Thus $\neg{0} = 0$ by \cref{real_neg_zero}.
        Thus $\neg{1} < 0$ by \cref{real_lt_subst_right}.
        Thus $\neg{1} + x < 0 + x$ by \cref{real_lt_add}.
        $\neg{1} + x = x + \neg{1}$ by \cref{real_add_comm}.
        $0 + x = x + 0$ by \cref{real_add_comm}.
        $x + 0 = x$ by \cref{real_add_ident,real_add_comm}.
        Thus $a < x$.
        Thus $0 + x < 1 + x$ by \cref{real_lt_add}.
        $0 + x = x$ by \cref{real_add_ident,real_add_comm}.
        $1 + x = x + 1$ by \cref{real_add_comm}.
        Thus $x < b$.
        Thus $a < b$ by \cref{real_lt_trans}.
        Thus $x\in\intervalopen{a}{b}$ by \cref{intervalopen}.
        $\intervalopen{a}{b}\subseteq\reals$ by \cref{intervalopen_subset_reals}.
        Thus $\intervalopen{a}{b}\in\pow{\reals}$ by \cref{powerset_intro}.
        Thus $\intervalopen{a}{b}\in\kBasisR$ by \cref{k_basis_reals}.
    \end{subproof}
    Show that for all $B_1,B_2,x$ such that $B_1\in\kBasisR$ and $B_2\in\kBasisR$ and $x\in B_1\inter B_2$
        there exists $B_3\in\kBasisR$ such that $x\in B_3$ and $B_3\subseteq B_1\inter B_2$.
    \begin{subproof}
        Fix $B_1$.
        Fix $B_2$.
        Fix $x$.
        Assume $B_1\in\kBasisR$ and $B_2\in\kBasisR$ and $x\in B_1\inter B_2$.
        Take $a,b\in\reals$ such that $a < b$ and
            $(B_1 = \intervalopen{a}{b} \lor B_1 = \intervalopen{a}{b}\setminus \kset)$
            by \cref{k_basis_reals}.
        Take $c,d\in\reals$ such that $c < d$ and
            $(B_2 = \intervalopen{c}{d} \lor B_2 = \intervalopen{c}{d}\setminus \kset)$
            by \cref{k_basis_reals}.
        $x\in B_1$ by \cref{inter_elim_left}.
        $x\in B_2$ by \cref{inter_elim_right}.
        \begin{byCase}
        \caseOf{$B_1 = \intervalopen{a}{b}$ and $B_2 = \intervalopen{c}{d}$.}
            Thus $x\in\intervalopen{a}{b}$ and $x\in\intervalopen{c}{d}$.
            Thus $x\in\reals$ by \cref{intervalopen}.
            Thus $a < x$ and $x < b$ by \cref{intervalopen}.
            Thus $c < x$ and $x < d$ by \cref{intervalopen}.
            Take $e,f\in\reals$ such that $e < x$ and $x < f$ and
                $\intervalopen{e}{f} \subseteq \intervalopen{a}{b}\inter \intervalopen{c}{d}$
                by \cref{open_interval_refinement}.
            Let $B_3 = \intervalopen{e}{f}$.
            $B_3\subseteq\reals$ by \cref{intervalopen_subset_reals}.
            Thus $B_3\in\pow{\reals}$ by \cref{powerset_intro}.
            Thus $e < f$ by \cref{real_lt_trans}.
            Let $a_1 = e$.
            Let $b_1 = f$.
            Thus $a_1 < b_1$ and $B_3 = \intervalopen{a_1}{b_1}$.
            Thus $B_3\in\kBasisR$ by \cref{k_basis_reals}.
            Thus $x\in B_3$ by \cref{intervalopen}.
            Thus $B_3\subseteq B_1\inter B_2$.
        \caseOf{$B_1 = \intervalopen{a}{b}\setminus \kset$ and $B_2 = \intervalopen{c}{d}$.}
            Thus $x\in\intervalopen{a}{b}$ and $x\notin\kset$ by \cref{setminus_elim_left,setminus_elim_right}.
            Thus $x\in\intervalopen{c}{d}$.
            Thus $x\in\reals$ by \cref{intervalopen}.
            Thus $a < x$ and $x < b$ by \cref{intervalopen}.
            Thus $c < x$ and $x < d$ by \cref{intervalopen}.
            Take $e,f\in\reals$ such that $e < x$ and $x < f$ and
                $\intervalopen{e}{f} \subseteq \intervalopen{a}{b}\inter \intervalopen{c}{d}$
                by \cref{open_interval_refinement}.
            Let $B_3 = \intervalopen{e}{f}\setminus \kset$.
            Show that for all $y$ such that $y\in B_3$ we have $y\in\reals$.
            \begin{subproof}
                Fix $y$.
                Assume $y\in B_3$.
                Thus $y\in\intervalopen{e}{f}$ by \cref{setminus_elim_left}.
                Thus $y\in\reals$ by \cref{intervalopen}.
            \end{subproof}
            Thus $B_3\subseteq\reals$ by \cref{subseteq}.
            Thus $B_3\in\pow{\reals}$ by \cref{powerset_intro}.
            Thus $e < f$ by \cref{real_lt_trans}.
            Let $a_1 = e$.
            Let $b_1 = f$.
            Thus $a_1 < b_1$ and $B_3 = \intervalopen{a_1}{b_1}\setminus \kset$.
            Thus $B_3\in\kBasisR$ by \cref{k_basis_reals}.
            $x\in\intervalopen{e}{f}$ by \cref{intervalopen}.
            Thus $x\in B_3$ by \cref{setminus_intro}.
            Show that for all $y$ such that $y\in B_3$ we have $y\in B_1\inter B_2$.
            \begin{subproof}
                Fix $y$.
                Assume $y\in B_3$.
                $y\in\intervalopen{e}{f}$ and $y\notin\kset$ by \cref{setminus_elim_left,setminus_elim_right}.
                Thus $y\in\intervalopen{a}{b}\inter\intervalopen{c}{d}$ by \cref{subseteq}.
                Thus $y\in\intervalopen{a}{b}$ by \cref{inter_elim_left}.
                Thus $y\in\intervalopen{c}{d}$ by \cref{inter_elim_right}.
                Thus $y\in B_1$ by \cref{setminus_intro}.
                Thus $y\in B_2$.
                Thus $y\in B_1\inter B_2$ by \cref{inter_intro}.
            \end{subproof}
            Thus $B_3\subseteq B_1\inter B_2$ by \cref{subseteq}.
        \caseOf{$B_1 = \intervalopen{a}{b}$ and $B_2 = \intervalopen{c}{d}\setminus \kset$.}
            Thus $x\in\intervalopen{a}{b}$.
            Thus $x\in\intervalopen{c}{d}$ and $x\notin\kset$ by \cref{setminus_elim_left,setminus_elim_right}.
            Thus $x\in\reals$ by \cref{intervalopen}.
            Thus $a < x$ and $x < b$ by \cref{intervalopen}.
            Thus $c < x$ and $x < d$ by \cref{intervalopen}.
            Take $e,f\in\reals$ such that $e < x$ and $x < f$ and
                $\intervalopen{e}{f} \subseteq \intervalopen{a}{b}\inter \intervalopen{c}{d}$
                by \cref{open_interval_refinement}.
            Let $B_3 = \intervalopen{e}{f}\setminus \kset$.
            Show that for all $y$ such that $y\in B_3$ we have $y\in\reals$.
            \begin{subproof}
                Fix $y$.
                Assume $y\in B_3$.
                Thus $y\in\intervalopen{e}{f}$ by \cref{setminus_elim_left}.
                Thus $y\in\reals$ by \cref{intervalopen}.
            \end{subproof}
            Thus $B_3\subseteq\reals$ by \cref{subseteq}.
            Thus $B_3\in\pow{\reals}$ by \cref{powerset_intro}.
            Thus $e < f$ by \cref{real_lt_trans}.
            Let $a_1 = e$.
            Let $b_1 = f$.
            Thus $a_1 < b_1$ and $B_3 = \intervalopen{a_1}{b_1}\setminus \kset$.
            Thus $B_3\in\kBasisR$ by \cref{k_basis_reals}.
            $x\in\intervalopen{e}{f}$ by \cref{intervalopen}.
            Thus $x\in B_3$ by \cref{setminus_intro}.
            Show that for all $y$ such that $y\in B_3$ we have $y\in B_1\inter B_2$.
            \begin{subproof}
                Fix $y$.
                Assume $y\in B_3$.
                $y\in\intervalopen{e}{f}$ and $y\notin\kset$ by \cref{setminus_elim_left,setminus_elim_right}.
                Thus $y\in\intervalopen{a}{b}\inter\intervalopen{c}{d}$ by \cref{subseteq}.
                Thus $y\in\intervalopen{a}{b}$ by \cref{inter_elim_left}.
                Thus $y\in\intervalopen{c}{d}$ by \cref{inter_elim_right}.
                Thus $y\in B_2$ by \cref{setminus_intro}.
                Thus $y\in B_1$.
                Thus $y\in B_1\inter B_2$ by \cref{inter_intro}.
            \end{subproof}
            Thus $B_3\subseteq B_1\inter B_2$ by \cref{subseteq}.
        \caseOf{$B_1 = \intervalopen{a}{b}\setminus \kset$ and $B_2 = \intervalopen{c}{d}\setminus \kset$.}
            Thus $x\in\intervalopen{a}{b}$ and $x\notin\kset$ by \cref{setminus_elim_left,setminus_elim_right}.
            Thus $x\in\intervalopen{c}{d}$ and $x\notin\kset$ by \cref{setminus_elim_left,setminus_elim_right}.
            Thus $x\in\reals$ by \cref{intervalopen}.
            Thus $a < x$ and $x < b$ by \cref{intervalopen}.
            Thus $c < x$ and $x < d$ by \cref{intervalopen}.
            Take $e,f\in\reals$ such that $e < x$ and $x < f$ and
                $\intervalopen{e}{f} \subseteq \intervalopen{a}{b}\inter \intervalopen{c}{d}$
                by \cref{open_interval_refinement}.
            Let $B_3 = \intervalopen{e}{f}\setminus \kset$.
            Show that for all $y$ such that $y\in B_3$ we have $y\in\reals$.
            \begin{subproof}
                Fix $y$.
                Assume $y\in B_3$.
                Thus $y\in\intervalopen{e}{f}$ by \cref{setminus_elim_left}.
                Thus $y\in\reals$ by \cref{intervalopen}.
            \end{subproof}
            Thus $B_3\subseteq\reals$ by \cref{subseteq}.
            Thus $B_3\in\pow{\reals}$ by \cref{powerset_intro}.
            Thus $e < f$ by \cref{real_lt_trans}.
            Let $a_1 = e$.
            Let $b_1 = f$.
            Thus $a_1 < b_1$ and $B_3 = \intervalopen{a_1}{b_1}\setminus \kset$.
            Thus $B_3\in\kBasisR$ by \cref{k_basis_reals}.
            $x\in\intervalopen{e}{f}$ by \cref{intervalopen}.
            Thus $x\in B_3$ by \cref{setminus_intro}.
            Show that for all $y$ such that $y\in B_3$ we have $y\in B_1\inter B_2$.
            \begin{subproof}
                Fix $y$.
                Assume $y\in B_3$.
                $y\in\intervalopen{e}{f}$ and $y\notin\kset$ by \cref{setminus_elim_left,setminus_elim_right}.
                Thus $y\in\intervalopen{a}{b}\inter\intervalopen{c}{d}$ by \cref{subseteq}.
                Thus $y\in\intervalopen{a}{b}$ by \cref{inter_elim_left}.
                Thus $y\in\intervalopen{c}{d}$ by \cref{inter_elim_right}.
                Thus $y\in B_1$ by \cref{setminus_intro}.
                Thus $y\in B_2$ by \cref{setminus_intro}.
                Thus $y\in B_1\inter B_2$ by \cref{inter_intro}.
            \end{subproof}
            Thus $B_3\subseteq B_1\inter B_2$ by \cref{subseteq}.
        \end{byCase}
    \end{subproof}
    Thus $\kBasisR$ is a basis on $\reals$ by \cref{basis_on}.
\end{proof}

% from §13 Item 20: lower limit topology strictly finer than standard topology
\begin{lemma}\label{lower_limit_strictly_finer_standard}
    $\lowerLimitTopologyR$ is strictly finer than $\standardTopologyR$ on $\reals$.
\end{lemma}
\begin{proof}
    Let $T = \standardTopologyR$.
    Let $T_1 = \lowerLimitTopologyR$.
    Show that $T_1$ is finer than $T$ on $\reals$.
    \begin{subproof}
        $\standardBasisR$ is a basis for $T$ on $\reals$ by \cref{standard_basis_is_basis,basis_for_topology,standard_topology_reals}.
        $\lowerLimitBasisR$ is a basis for $T_1$ on $\reals$ by \cref{lower_limit_basis_is_basis,basis_for_topology,lower_limit_topology_reals}.
        Show that for all $x,B_2$ such that $x\in\reals$ and $B_2\in\standardBasisR$ and $x\in B_2$
            there exists $B_3\in\lowerLimitBasisR$ such that $x\in B_3$ and $B_3\subseteq B_2$.
        \begin{subproof}
            Fix $x$.
            Fix $B_2$.
            Assume $x\in\reals$ and $B_2\in\standardBasisR$ and $x\in B_2$.
            Take $a,b\in\reals$ such that $a < b$ and $B_2 = \intervalopen{a}{b}$ by \cref{standard_basis_reals}.
            Thus $a < x$ and $x < b$ by \cref{intervalopen}.
            Let $B_3 = \intervalclosedopen{x}{b}$.
            $B_3\subseteq\reals$ by \cref{intervalclosedopen_subset_reals}.
            Thus $B_3\in\pow{\reals}$ by \cref{powerset_intro}.
            Thus $B_3\in\lowerLimitBasisR$ by \cref{lower_limit_basis_reals}.
            Thus $x\in B_3$ by \cref{interval_closed_open}.
            Show that for all $y$ such that $y\in B_3$ we have $y\in B_2$.
            \begin{subproof}
                Fix $y$.
                Assume $y\in B_3$.
                $x \leq y$ and $y < b$ by \cref{interval_closed_open}.
                $y\in\reals$ by \cref{interval_closed_open}.
                Show that $a < y$.
                \begin{subproof}
                    $x < y$ or $x = y$ by \cref{real_leq_split}.
                    \begin{byCase}
                    \caseOf{$x = y$.} Thus $a < y$ by \cref{real_lt_subst_right}.
                    \caseOf{$x < y$.} Thus $a < y$ by \cref{real_lt_trans}.
                    \end{byCase}
                \end{subproof}
                Thus $y\in\intervalopen{a}{b}$ by \cref{intervalopen}.
            \end{subproof}
            Thus $B_3\subseteq B_2$ by \cref{subseteq}.
        \end{subproof}
        Thus $T_1$ is finer than $T$ on $\reals$ by \cref{basis_finer_criterion}.
    \end{subproof}
    Show that $T \neq T_1$.
    \begin{subproof}
        Let $U = \intervalclosedopen{0}{1}$.
        $0\in\reals$ by \cref{real_add_ident}.
        $1\in\reals$ by \cref{real_naturals_subset_reals,real_one_in_naturals,subseteq}.
        $0 < 1$ by \cref{real_zero_lt_one}.
        $U\subseteq\reals$ by \cref{intervalclosedopen_subset_reals}.
        Thus $U\in\pow{\reals}$ by \cref{powerset_intro}.
        Thus $U\in\lowerLimitBasisR$ by \cref{lower_limit_basis_reals}.
        Thus $U\in T_1$ by \cref{basis_elem_open_in_generated,basis_for_topology,lower_limit_basis_is_basis,lower_limit_topology_reals}.
        Show that $U\notin T$.
        \begin{subproof}
            Suppose not.
            Then $U\in T$.
            $0\leq 0$.
            Thus $0\in U$ by \cref{interval_closed_open}.
            There exists $B_2\in\standardBasisR$ such that $0\in B_2$ and $B_2\subseteq U$
                by \cref{topology_generated_by_basis,standard_topology_reals}.
            Take $a,b\in\reals$ such that $a < b$ and $B_2 = \intervalopen{a}{b}$ by \cref{standard_basis_reals}.
            Thus $0\in\intervalopen{a}{b}$.
            Thus $a < 0$ by \cref{intervalopen}.
            Thus $0 < b$ by \cref{intervalopen}.
            $a\in\reals$.
            $0\in\reals$ by \cref{real_add_ident}.
            Thus there exists $p\in\rationals$ such that $p\in\intervalopen{a}{0}$ by \cref{rational_between_negative_and_zero}.
            Take $p\in\rationals$ such that $p\in\intervalopen{a}{0}$.
            Thus $p\in\intervalopen{a}{b}$ by \cref{intervalopen,real_lt_trans}.
            Thus $p\in B_2$.
            $p < 0$ by \cref{intervalopen}.
            Show that $p\notin U$.
            \begin{subproof}
                Suppose not.
                Thus $p\in U$.
                Thus $p\in\reals$ and $(0 < p \lor 0 = p)$ and $p < 1$ by \cref{interval_closed_open}.
                $0\in\reals$ by \cref{real_add_ident}.
                \begin{byCase}
                \caseOf{$0 < p$.}
                    Thus $0 < 0$ by \cref{real_lt_trans}.
                    Contradiction by \cref{real_lt_irreflexive}.
                \caseOf{$0 = p$.}
                    Thus $0 < 0$ by \cref{real_lt_subst_left}.
                    Contradiction by \cref{real_lt_irreflexive}.
                \end{byCase}
            \end{subproof}
            Contradiction by \cref{subseteq}.
        \end{subproof}
    \end{subproof}
    Thus $T_1$ is strictly finer than $T$ on $\reals$ by \cref{strictly_finer_topology}.
\end{proof}

% from §13 Item 21: K-topology strictly finer than standard topology
\begin{lemma}\label{k_strictly_finer_standard}
    $\kTopologyR$ is strictly finer than $\standardTopologyR$ on $\reals$.
\end{lemma}
\begin{proof}
    Let $T = \standardTopologyR$.
    Let $T_2 = \kTopologyR$.
    Show that $T_2$ is finer than $T$ on $\reals$.
    \begin{subproof}
        $\standardBasisR$ is a basis for $T$ on $\reals$ by \cref{standard_basis_is_basis,basis_for_topology,standard_topology_reals}.
        $\kBasisR$ is a basis for $T_2$ on $\reals$ by \cref{k_basis_is_basis,basis_for_topology,k_topology_reals}.
        Show that for all $x,B_2$ such that $x\in\reals$ and $B_2\in\standardBasisR$ and $x\in B_2$
            there exists $B_3\in\kBasisR$ such that $x\in B_3$ and $B_3\subseteq B_2$.
        \begin{subproof}
            Fix $x$.
            Fix $B_2$.
            Assume $x\in\reals$ and $B_2\in\standardBasisR$ and $x\in B_2$.
            Thus $B_2\in\kBasisR$ by \cref{k_basis_reals,standard_basis_reals}.
            Thus $x\in B_2$ and $B_2\subseteq B_2$.
        \end{subproof}
        Thus $T_2$ is finer than $T$ on $\reals$ by \cref{basis_finer_criterion}.
    \end{subproof}
    Show that $T \neq T_2$.
    \begin{subproof}
        Let $U = \intervalopen{\neg{1}}{1}\setminus \kset$.
        $1\in\reals$ by \cref{real_naturals_subset_reals,real_one_in_naturals,subseteq}.
        $\neg{1}\in\reals$ by \cref{real_add_inverse}.
        $0 < 1$ by \cref{real_zero_lt_one}.
        $0\in\reals$ by \cref{real_add_ident}.
        Thus $\neg{1} < \neg{0}$ by \cref{real_lt_neg}.
        Thus $\neg{0} = 0$ by \cref{real_neg_zero}.
        Thus $\neg{1} < 0$ by \cref{real_lt_subst_right}.
        Thus $0 < 1$.
        Thus $\neg{1} < 1$ by \cref{real_lt_trans}.
        Thus $0\in\intervalopen{\neg{1}}{1}$ by \cref{intervalopen}.
        Show that $0\notin\kset$.
        \begin{subproof}
            Suppose not.
            Then $0\in\kset$.
            Thus $0 < 0$ by \cref{k_set_positive}.
            Contradiction by \cref{real_lt_irreflexive}.
        \end{subproof}
        Thus $0\in U$ by \cref{setminus_intro}.
        Show that for all $x$ such that $x\in U$ we have $x\in\reals$.
        \begin{subproof}
            Fix $x$.
            Assume $x\in U$.
            Thus $x\in\intervalopen{\neg{1}}{1}$ by \cref{setminus_elim_left}.
            Thus $x\in\reals$ by \cref{intervalopen}.
        \end{subproof}
        Thus $U\subseteq\reals$ by \cref{subseteq}.
        Thus $U\in\pow{\reals}$ by \cref{powerset_intro}.
        Thus $U\in\kBasisR$ by \cref{k_basis_reals}.
        Thus $U\in T_2$ by \cref{basis_elem_open_in_generated,basis_for_topology,k_basis_is_basis,k_topology_reals}.
        Show that $U\notin T$.
        \begin{subproof}
            Suppose not.
            Then $U\in T$.
            There exists $B_2\in\standardBasisR$ such that $0\in B_2$ and $B_2\subseteq U$
                by \cref{topology_generated_by_basis,standard_topology_reals}.
            Take $a,b\in\reals$ such that $a < b$ and $B_2 = \intervalopen{a}{b}$ by \cref{standard_basis_reals}.
            Thus $a < 0$ and $0 < b$ by \cref{intervalopen}.
            Take $N\in\naturalsPlus$ such that $0 < 1 / N$ and $1 / N < b$ by \cref{real_small_reciprocal}.
            $N\in\reals$ by \cref{real_naturals_subset_reals,subseteq}.
            Show that $0 < N$.
            \begin{subproof}
                $1\in\reals$ by \cref{real_mul_ident}.
                $0\in\reals$ by \cref{real_add_ident}.
                $0 < 1$ by \cref{real_zero_lt_one}.
                $1 < N$ or $1 = N$ by \cref{real_one_leq_naturalsplus}.
                \begin{byCase}
                \caseOf{$1 = N$.} Thus $0 < N$.
                \caseOf{$1 < N$.} Thus $0 < N$ by \cref{real_lt_trans}.
                \end{byCase}
            \end{subproof}
            Thus $N\neq 0$ by \cref{real_pos_nonzero}.
            Thus $\inv{N}\in\reals$ by \cref{real_mul_inverse}.
            $1\in\reals$ by \cref{real_mul_ident}.
            $1 / N = 1 \rmul \inv{N}$ by \cref{real_div}.
            Thus $1 / N\in\reals$ by \cref{real_mul_closed}.
            Thus $a < 1 / N$ by \cref{real_lt_trans}.
            Thus $1 / N\in\intervalopen{a}{b}$ by \cref{intervalopen}.
            Thus $1 / N\in B_2$.
            Thus $1 / N\in\kset$ by \cref{k_set}.
            Thus $1 / N\notin U$ by \cref{setminus_elim_right}.
            Contradiction by \cref{subseteq}.
        \end{subproof}
    \end{subproof}
    Thus $T_2$ is strictly finer than $T$ on $\reals$ by \cref{strictly_finer_topology}.
\end{proof}

% from §13 Item 22: lower limit and K topologies are not comparable
\begin{lemma}\label{lower_limit_k_not_comparable}
    $\lowerLimitTopologyR$ is not comparable with $\kTopologyR$ on $\reals$.
\end{lemma}
\begin{proof}
    Let $T_1 = \lowerLimitTopologyR$.
    Let $T_2 = \kTopologyR$.
    Let $U_1 = \intervalclosedopen{0}{1}$.
    Let $U_2 = \intervalopen{\neg{1}}{1}\setminus \kset$.
    $0\in\reals$ by \cref{real_add_ident}.
    $1\in\reals$ by \cref{real_naturals_subset_reals,real_one_in_naturals,subseteq}.
    $0 < 1$ by \cref{real_zero_lt_one}.
    $U_1\subseteq\reals$ by \cref{intervalclosedopen_subset_reals}.
    Thus $U_1\in\pow{\reals}$ by \cref{powerset_intro}.
    Thus $U_1\in\lowerLimitBasisR$ by \cref{lower_limit_basis_reals}.
    Thus $U_1\in T_1$ by \cref{basis_elem_open_in_generated,basis_for_topology,lower_limit_basis_is_basis,lower_limit_topology_reals}.
    $\neg{1}\in\reals$ by \cref{real_add_inverse}.
    Thus $\neg{1} < \neg{0}$ by \cref{real_lt_neg}.
    Thus $\neg{0} = 0$ by \cref{real_neg_zero}.
    Thus $\neg{1} < 0$ by \cref{real_lt_subst_right}.
    Thus $\neg{1} < 1$ by \cref{real_lt_trans}.
    Show that for all $x$ such that $x\in U_2$ we have $x\in\reals$.
    \begin{subproof}
        Fix $x$.
        Assume $x\in U_2$.
        Thus $x\in\intervalopen{\neg{1}}{1}$ by \cref{setminus_elim_left}.
        Thus $x\in\reals$ by \cref{intervalopen}.
    \end{subproof}
    Thus $U_2\subseteq\reals$ by \cref{subseteq}.
    Thus $U_2\in\pow{\reals}$ by \cref{powerset_intro}.
    Thus $U_2\in\kBasisR$ by \cref{k_basis_reals}.
    Thus $U_2\in T_2$ by \cref{basis_elem_open_in_generated,basis_for_topology,k_basis_is_basis,k_topology_reals}.
    Show that $U_1\notin T_2$.
    \begin{subproof}
        Suppose not.
        Then $U_1\in T_2$.
        $0\leq 0$.
        Thus $0\in U_1$ by \cref{interval_closed_open}.
        There exists $B_2\in\kBasisR$ such that $0\in B_2$ and $B_2\subseteq U_1$
            by \cref{topology_generated_by_basis,k_topology_reals}.
        Take $a,b\in\reals$ such that $a < b$ and
            $(B_2 = \intervalopen{a}{b} \lor B_2 = \intervalopen{a}{b}\setminus \kset)$
            by \cref{k_basis_reals}.
        Thus $0\in\intervalopen{a}{b}$ by \cref{intervalopen,setminus_elim_left}.
        Thus $a < 0$ and $0 < b$ by \cref{intervalopen}.
        Thus there exists $p\in\rationals$ such that $p\in\intervalopen{a}{0}$ by \cref{rational_between_negative_and_zero}.
        Take $p\in\rationals$ such that $p\in\intervalopen{a}{0}$.
        Thus $p\in\reals$ by \cref{intervalopen}.
        Thus $p\in\intervalopen{a}{b}$ by \cref{intervalopen,real_lt_trans}.
        Show that $p\notin\kset$.
        \begin{subproof}
            Suppose not.
            Then $p\in\kset$.
            Thus $0 < p$ by \cref{k_set_positive}.
            Thus $p < 0$ by \cref{intervalopen}.
            Thus $p < p$ by \cref{real_lt_trans}.
            Contradiction by \cref{real_lt_irreflexive}.
        \end{subproof}
        Thus $p\in B_2$ by \cref{setminus_intro}.
        Thus $p < 0$ by \cref{intervalopen}.
        Show that $p\notin U_1$.
        \begin{subproof}
            Suppose not.
            Thus $p\in U_1$.
            Thus $p\in\reals$ and $(0 < p \lor 0 = p)$ and $p < 1$ by \cref{interval_closed_open}.
            $0\in\reals$ by \cref{real_add_ident}.
            \begin{byCase}
            \caseOf{$0 < p$.}
                Thus $0 < 0$ by \cref{real_lt_trans}.
                Contradiction by \cref{real_lt_irreflexive}.
            \caseOf{$0 = p$.}
                Thus $0 < 0$ by \cref{real_lt_subst_left}.
                Contradiction by \cref{real_lt_irreflexive}.
            \end{byCase}
        \end{subproof}
        Contradiction by \cref{subseteq}.
    \end{subproof}
    Show that $U_2\notin T_1$.
    \begin{subproof}
        Suppose not.
        Then $U_2\in T_1$.
        $0 < 1$ by \cref{real_zero_lt_one}.
        Thus $\neg{1} < \neg{0}$ by \cref{real_lt_neg}.
        Thus $\neg{0} = 0$ by \cref{real_neg_zero}.
        Thus $\neg{1} < 0$ by \cref{real_lt_subst_right}.
        Thus $0\in\intervalopen{\neg{1}}{1}$ by \cref{intervalopen}.
        Show that $0\notin\kset$.
        \begin{subproof}
            Suppose not.
            Then $0\in\kset$.
            Thus $0 < 0$ by \cref{k_set_positive}.
            Contradiction by \cref{real_lt_irreflexive}.
        \end{subproof}
        Thus $0\in U_2$ by \cref{setminus_intro}.
        There exists $B_1\in\lowerLimitBasisR$ such that $0\in B_1$ and $B_1\subseteq U_2$
            by \cref{topology_generated_by_basis,lower_limit_topology_reals}.
        Take $a,b\in\reals$ such that $a < b$ and $B_1 = \intervalclosedopen{a}{b}$ by \cref{lower_limit_basis_reals}.
        Thus $a \leq 0$ and $0 < b$ by \cref{interval_closed_open}.
        Take $N\in\naturalsPlus$ such that $0 < 1 / N$ and $1 / N < b$ by \cref{real_small_reciprocal}.
        $N\in\reals$ by \cref{real_naturals_subset_reals,subseteq}.
        Show that $0 < N$.
        \begin{subproof}
            $1\in\reals$ by \cref{real_mul_ident}.
            $0\in\reals$ by \cref{real_add_ident}.
            $0 < 1$ by \cref{real_zero_lt_one}.
            $1 < N$ or $1 = N$ by \cref{real_one_leq_naturalsplus}.
            \begin{byCase}
            \caseOf{$1 = N$.} Thus $0 < N$.
            \caseOf{$1 < N$.} Thus $0 < N$ by \cref{real_lt_trans}.
            \end{byCase}
        \end{subproof}
        Thus $N\neq 0$ by \cref{real_pos_nonzero}.
        Thus $\inv{N}\in\reals$ by \cref{real_mul_inverse}.
        $1\in\reals$ by \cref{real_mul_ident}.
        $1 / N = 1 \rmul \inv{N}$ by \cref{real_div}.
        Thus $1 / N\in\reals$ by \cref{real_mul_closed}.
        Show that $a < 1 / N$.
        \begin{subproof}
            $a < 0$ or $a = 0$ by \cref{real_leq_split}.
            \begin{byCase}
            \caseOf{$a = 0$.} Thus $a < 1 / N$ by \cref{real_lt_subst_left}.
            \caseOf{$a < 0$.} Thus $a < 1 / N$ by \cref{real_lt_trans}.
            \end{byCase}
        \end{subproof}
        Thus $1 / N\in\intervalclosedopen{a}{b}$ by \cref{interval_closed_open}.
        Thus $1 / N\in B_1$.
        Thus $1 / N\in\kset$ by \cref{k_set}.
        Thus $1 / N\notin U_2$ by \cref{setminus_elim_right}.
        Contradiction by \cref{subseteq}.
    \end{subproof}
    Show that $T_1$ is not comparable with $T_2$ on $\reals$.
    \begin{subproof}
        Suppose not.
        Thus $T_1$ is comparable with $T_2$ on $\reals$.
        Thus $T_1\subseteq T_2$ or $T_2\subseteq T_1$ by \cref{comparable_topologies}.
        \begin{byCase}
        \caseOf{$T_1\subseteq T_2$.}
            Contradiction by \cref{subseteq}.
        \caseOf{$T_2\subseteq T_1$.}
            Contradiction by \cref{subseteq}.
        \end{byCase}
    \end{subproof}
\end{proof}

% from §13 Item 23: subbasis on a set
\begin{definition}\label{subbasis_on}
    $S$ is a subbasis on $X$ iff $S\subseteq\pow{X}$ and $\unions{S} = X$.
\end{definition}

% from §13 Item 24: finite intersections of subbasis elements
\begin{definition}\label{finite_intersections}
    $\finiteIntersections{S} = \{B\in\pow{\unions{S}} \mid \text{there exists $F\subseteq S$ such that $F$ is finite and inhabited and $B = \inters{F}$}\}$.
\end{definition}

% from §13 Item 25: topology generated by a subbasis
\begin{definition}\label{topology_generated_by_subbasis}
    $\subbasisTopology{S}{X} = \basisTopology{\finiteIntersections{S}}{X}$.
\end{definition}

% from §13 Item 26: finite intersections of a subbasis form a basis
\begin{lemma}\label{subbasis_finite_intersections_basis}
    Assume $S$ is a subbasis on $X$.
    Then $\finiteIntersections{S}$ is a basis on $X$.
\end{lemma}
\begin{proof}
    $S\subseteq\pow{X}$ by \cref{subbasis_on}.
    $\unions{S} = X$ by \cref{subbasis_on}.
    Show that for all $B_1,B_2,x$ such that $B_1\in\finiteIntersections{S}$ and $B_2\in\finiteIntersections{S}$ and $x\in B_1\inter B_2$
        there exists $B_3\in\finiteIntersections{S}$ such that $x\in B_3$ and $B_3\subseteq B_1\inter B_2$.
    \begin{subproof}
        Fix $B_1$.
        Fix $B_2$.
        Fix $x$.
        Assume $B_1\in\finiteIntersections{S}$ and $B_2\in\finiteIntersections{S}$ and $x\in B_1\inter B_2$.
        Take $F_1\subseteq S$ such that $F_1$ is finite and inhabited and $B_1 = \inters{F_1}$ by \cref{finite_intersections}.
        Take $F_2\subseteq S$ such that $F_2$ is finite and inhabited and $B_2 = \inters{F_2}$ by \cref{finite_intersections}.
        Let $F = F_1\union F_2$.
        Thus $F\subseteq S$ by \cref{union_subsets_is_subset}.
        $F$ is finite by \cref{union_of_finite_is_finite}.
        Show that $F$ is inhabited.
        \begin{subproof}
            Take $y$ such that $y\in F_1$.
            Thus $y\in F$ by \cref{union_intro_left}.
        \end{subproof}
        $B_1\in\pow{\unions{S}}$ by \cref{finite_intersections}.
        $B_2\in\pow{\unions{S}}$ by \cref{finite_intersections}.
        Then $B_1\subseteq\unions{S}$ and $B_2\subseteq\unions{S}$ by \cref{pow_iff}.
        Thus $B_1\inter B_2\subseteq\unions{S}$ by \cref{inter_subseteq}.
        Thus $B_1\inter B_2\in\pow{\unions{S}}$ by \cref{powerset_intro}.
        Then $B_1\inter B_2 = \inters{F}$ by \cref{inters_distrib_union}.
        Thus $B_1\inter B_2\in\finiteIntersections{S}$ by \cref{finite_intersections}.
    \end{subproof}
    Show that for all $x$ such that $x\in X$ there exists $B\in\finiteIntersections{S}$ such that $x\in B$ and
        for all $B_1,B_2,x$ such that $B_1\in\finiteIntersections{S}$ and $B_2\in\finiteIntersections{S}$ and $x\in B_1\inter B_2$
        there exists $B_3\in\finiteIntersections{S}$ such that $x\in B_3$ and $B_3\subseteq B_1\inter B_2$.
    \begin{subproof}
        Fix $x$.
        Assume $x\in X$.
        Take $S_1\in S$ such that $x\in S_1$ by \cref{unions_iff}.
        $S_1\in\pow{X}$ by \cref{subseteq}.
        Then $S_1\subseteq X$ by \cref{pow_iff}.
        Thus $S_1\subseteq\unions{S}$.
        Thus $S_1\in\pow{\unions{S}}$ by \cref{powerset_intro}.
        Let $F = \{S_1\}$.
        Then $F\subseteq S$ by \cref{singleton_subset_intro}.
        Show that $F$ is finite.
        \begin{subproof}
            Let $G = \{S_1,S_1\}$.
            $G$ is finite by \cref{pair_is_finite}.
            $S_1\in G$ by \cref{upair_intro_left}.
            Thus $F\subseteq G$ by \cref{singleton_subset_intro}.
            Thus $F$ is finite by \cref{subseteq_finite_is_finite}.
        \end{subproof}
        Thus $S_1\in F$ by \cref{singleton_intro}.
        Thus $F$ is inhabited.
        Thus $\inters{F} = S_1$ by \cref{inters_singleton}.
        Thus $S_1\in\finiteIntersections{S}$ by \cref{finite_intersections}.
        Thus there exists $B\in\finiteIntersections{S}$ such that $x\in B$ and
            for all $B_1,B_2,x$ such that $B_1\in\finiteIntersections{S}$ and $B_2\in\finiteIntersections{S}$ and $x\in B_1\inter B_2$
            there exists $B_3\in\finiteIntersections{S}$ such that $x\in B_3$ and $B_3\subseteq B_1\inter B_2$.
    \end{subproof}
    Show that for all $B$ such that $B\in\finiteIntersections{S}$ we have $B\in\pow{X}$.
    \begin{subproof}
        Fix $B$.
        Assume $B\in\finiteIntersections{S}$.
        Then $B\in\pow{\unions{S}}$ by \cref{finite_intersections}.
        Then $B\subseteq\unions{S}$ by \cref{pow_iff}.
        Thus $B\subseteq X$.
        Thus $B\in\pow{X}$ by \cref{powerset_intro}.
    \end{subproof}
    Thus $\finiteIntersections{S}\subseteq\pow{X}$ by \cref{subseteq}.
    Thus $\finiteIntersections{S}$ is a basis on $X$ by \cref{basis_on}.
\end{proof}

% from §14 Item 1: more than one element
\begin{definition}\label{more_than_one_element}
    $X$ has more than one element iff there exist $x, y\in X$ such that $x\neq y$.
\end{definition}

% from §14 Item 2: simple order relation
\begin{definition}\label{simple_order_relation}
    $R$ is a simple order relation on $X$ iff
    $R$ is a binary relation on $X$ and
    $R$ is transitive and
    $R$ is asymmetric and
    $R$ is connex on $X$.
\end{definition}

% from §14 Item 3: smallest element
\begin{definition}\label{smallest_element}
    $a$ is a smallest element of $X$ with respect to $R$ iff
    $a\in X$ and for all $x\in X$ we have $a\mathrel{R}x$ or $a = x$.
\end{definition}

% from §14 Item 4: largest element
\begin{definition}\label{largest_element}
    $b$ is a largest element of $X$ with respect to $R$ iff
    $b\in X$ and for all $x\in X$ we have $x\mathrel{R}b$ or $x = b$.
\end{definition}

% from §14 Item 5: open interval in an ordered set
\begin{definition}\label{intervalopen_rel}
    $\intervalopenrel{R}{X}{a}{b} = \{x\in X \mid a\mathrel{R}x \land x\mathrel{R}b\}$.
\end{definition}

% from §14 Item 6: half-open interval in an ordered set
\begin{definition}\label{intervalopenclosed_rel}
    $\intervalopenclosedrel{R}{X}{a}{b} = \{x\in X \mid a\mathrel{R}x \land (x\mathrel{R}b \lor x = b)\}$.
\end{definition}

% from §14 Item 7: half-open interval in an ordered set
\begin{definition}\label{intervalclosedopen_rel}
    $\intervalclosedopenrel{R}{X}{a}{b} = \{x\in X \mid (a\mathrel{R}x \lor x = a) \land x\mathrel{R}b\}$.
\end{definition}

% from §14 Item 8: closed interval in an ordered set
\begin{definition}\label{intervalclosed_rel}
    $\intervalclosedrel{R}{X}{a}{b} = \{x\in X \mid (a\mathrel{R}x \lor x = a) \land (x\mathrel{R}b \lor x = b)\}$.
\end{definition}

% from §14 Item 9: order basis
\begin{definition}\label{order_basis}
    $\orderBasis{R}{X} = \{B\in\pow{X} \mid
        (\exists a, b\in X. a\mathrel{R}b \land B = \intervalopenrel{R}{X}{a}{b}) \lor
        (\exists b\in X. \exists a_0\in X. (\forall x\in X. a_0\mathrel{R}x \lor a_0 = x) \land B = \intervalclosedopenrel{R}{X}{a_0}{b}) \lor
        (\exists a\in X. \exists b_0\in X. (\forall x\in X. x\mathrel{R}b_0 \lor x = b_0) \land B = \intervalopenclosedrel{R}{X}{a}{b_0})\}$.
\end{definition}

% from §14 Item 10: order topology
\begin{definition}\label{order_topology}
    $T$ is the order topology on $X$ with respect to $R$ iff
    $R$ is a simple order relation on $X$ and
    $X$ has more than one element and
    $T = \basisTopology{\orderBasis{R}{X}}{X}$.
\end{definition}

% from §14 helper lemma 1: distinct element in a set with more than one element
\begin{lemma}\label{more_than_one_element_has_other}
    Assume $X$ has more than one element.
    Suppose $x\in X$.
    Then there exists $y\in X$ such that $y\neq x$.
\end{lemma}
\begin{proof}
    Take $u$ such that $u\in X$ and there exists $v\in X$ such that $v\neq u$ by \cref{more_than_one_element}.
    Take $v$ such that $v\in X$ and $v\neq u$.
    \begin{byCase}
    \caseOf{$x = u$.}
        Let $y = v$.
        Thus $y\in X$ and $y\neq x$.
        Thus there exists $y\in X$ such that $y\neq x$.
    \caseOf{$x \neq u$.}
        Let $y = u$.
        Thus $y\in X$ and $y\neq x$.
        Thus there exists $y\in X$ such that $y\neq x$.
    \end{byCase}
\end{proof}

% from §14 helper lemma 2: non-smallest has a predecessor
\begin{lemma}\label{non_smallest_has_predecessor}
    Assume $R$ is a simple order relation on $X$.
    Suppose $x\in X$.
    Suppose $x$ is not a smallest element of $X$ with respect to $R$.
    Then there exists $y\in X$ such that $y\mathrel{R}x$.
\end{lemma}
\begin{proof}
    Suppose not.
    Then for all $y\in X$ we have $y\not\mathrel{R}x$.
    $R$ is connex on $X$ by \cref{simple_order_relation}.
    Show that $x$ is a smallest element of $X$ with respect to $R$.
    \begin{subproof}
        Show that for all $y\in X$ we have $x\mathrel{R}y$ or $x = y$.
        \begin{subproof}
            Fix $y$.
            Assume $y\in X$.
            \begin{byCase}
            \caseOf{$y = x$.} Thus $x = y$.
            \caseOf{$y \neq x$.}
                $x\mathrel{R}y$ or $y\mathrel{R}x$ by \cref{connex}.
                Show that $x\mathrel{R}y$.
                \begin{subproof}
                    Suppose not.
                    Then $y\mathrel{R}x$.
                    Contradiction.
                \end{subproof}
            \end{byCase}
        \end{subproof}
        $x\in X$.
        Thus $x$ is a smallest element of $X$ with respect to $R$ by \cref{smallest_element}.
    \end{subproof}
    Contradiction.
\end{proof}

% from §14 helper lemma 3: non-largest has a successor
\begin{lemma}\label{non_largest_has_successor}
    Assume $R$ is a simple order relation on $X$.
    Suppose $x\in X$.
    Suppose $x$ is not a largest element of $X$ with respect to $R$.
    Then there exists $y\in X$ such that $x\mathrel{R}y$.
\end{lemma}
\begin{proof}
    Suppose not.
    Then for all $y\in X$ we have $x\not\mathrel{R}y$.
    $R$ is connex on $X$ by \cref{simple_order_relation}.
    Show that $x$ is a largest element of $X$ with respect to $R$.
    \begin{subproof}
        Show that for all $y\in X$ we have $y\mathrel{R}x$ or $y = x$.
        \begin{subproof}
            Fix $y$.
            Assume $y\in X$.
            \begin{byCase}
            \caseOf{$y = x$.} Thus $y = x$.
            \caseOf{$y \neq x$.}
                $y\mathrel{R}x$ or $x\mathrel{R}y$ by \cref{connex}.
                Show that $y\mathrel{R}x$.
                \begin{subproof}
                    Suppose not.
                    Then $x\mathrel{R}y$.
                    Contradiction.
                \end{subproof}
            \end{byCase}
        \end{subproof}
        $x\in X$.
        Thus $x$ is a largest element of $X$ with respect to $R$ by \cref{largest_element}.
    \end{subproof}
    Contradiction.
\end{proof}

% from §14 helper lemma 4: open interval inclusion by endpoints
\begin{lemma}\label{intervalopenrel_subset_by_endpoints}
    Assume $R$ is transitive.
    Suppose $a,a_1,b,b_1\in X$.
    Suppose $a_1\mathrel{R}a$ or $a_1 = a$.
    Suppose $b\mathrel{R}b_1$ or $b = b_1$.
    Then $\intervalopenrel{R}{X}{a}{b} \subseteq \intervalopenrel{R}{X}{a_1}{b_1}$.
\end{lemma}
\begin{proof}
    Show that for all $x$ such that $x\in \intervalopenrel{R}{X}{a}{b}$ we have $x\in \intervalopenrel{R}{X}{a_1}{b_1}$.
    \begin{subproof}
        Fix $x$.
        Assume $x\in \intervalopenrel{R}{X}{a}{b}$.
        Thus $x\in X$ and $a\mathrel{R}x$ and $x\mathrel{R}b$ by \cref{intervalopen_rel}.
        Show that $a_1\mathrel{R}x$.
        \begin{subproof}
            \begin{byCase}
            \caseOf{$a_1 = a$.} Thus $a_1\mathrel{R}x$.
            \caseOf{$a_1\mathrel{R}a$.}
                Thus $a_1\mathrel{R}x$ by \hyperref[transitive]{transitivity}.
            \end{byCase}
        \end{subproof}
        Show that $x\mathrel{R}b_1$.
        \begin{subproof}
            \begin{byCase}
            \caseOf{$b = b_1$.} Thus $x\mathrel{R}b_1$.
            \caseOf{$b\mathrel{R}b_1$.}
                Thus $x\mathrel{R}b_1$ by \hyperref[transitive]{transitivity}.
            \end{byCase}
        \end{subproof}
        Thus $x\in \intervalopenrel{R}{X}{a_1}{b_1}$ by \cref{intervalopen_rel}.
    \end{subproof}
    Thus $\intervalopenrel{R}{X}{a}{b} \subseteq \intervalopenrel{R}{X}{a_1}{b_1}$ by \cref{subseteq}.
\end{proof}

% from §14 helper lemma 5: closed-open interval inclusion by right endpoint
\begin{lemma}\label{intervalclosedopenrel_subset_by_right}
    Assume $R$ is transitive.
    Suppose $a,b,b_1\in X$.
    Suppose $b\mathrel{R}b_1$ or $b = b_1$.
    Then $\intervalclosedopenrel{R}{X}{a}{b} \subseteq \intervalclosedopenrel{R}{X}{a}{b_1}$.
\end{lemma}
\begin{proof}
    Show that for all $x$ such that $x\in \intervalclosedopenrel{R}{X}{a}{b}$ we have $x\in \intervalclosedopenrel{R}{X}{a}{b_1}$.
    \begin{subproof}
        Fix $x$.
        Assume $x\in \intervalclosedopenrel{R}{X}{a}{b}$.
        Thus $x\in X$ and $(a\mathrel{R}x \lor x = a)$ and $x\mathrel{R}b$ by \cref{intervalclosedopen_rel}.
        Show that $x\mathrel{R}b_1$.
        \begin{subproof}
            \begin{byCase}
            \caseOf{$b = b_1$.} Thus $x\mathrel{R}b_1$.
            \caseOf{$b\mathrel{R}b_1$.}
                Thus $x\mathrel{R}b_1$ by \hyperref[transitive]{transitivity}.
            \end{byCase}
        \end{subproof}
        Thus $x\in \intervalclosedopenrel{R}{X}{a}{b_1}$ by \cref{intervalclosedopen_rel}.
    \end{subproof}
    Thus $\intervalclosedopenrel{R}{X}{a}{b} \subseteq \intervalclosedopenrel{R}{X}{a}{b_1}$ by \cref{subseteq}.
\end{proof}

% from §14 helper lemma 6: open-closed interval inclusion by left endpoint
\begin{lemma}\label{intervalopenclosedrel_subset_by_left}
    Assume $R$ is transitive.
    Suppose $a,a_1,b\in X$.
    Suppose $a_1\mathrel{R}a$ or $a_1 = a$.
    Then $\intervalopenclosedrel{R}{X}{a}{b} \subseteq \intervalopenclosedrel{R}{X}{a_1}{b}$.
\end{lemma}
\begin{proof}
    Show that for all $x$ such that $x\in \intervalopenclosedrel{R}{X}{a}{b}$ we have $x\in \intervalopenclosedrel{R}{X}{a_1}{b}$.
    \begin{subproof}
        Fix $x$.
        Assume $x\in \intervalopenclosedrel{R}{X}{a}{b}$.
        Thus $x\in X$ and $a\mathrel{R}x$ and $(x\mathrel{R}b \lor x = b)$ by \cref{intervalopenclosed_rel}.
        Show that $a_1\mathrel{R}x$.
        \begin{subproof}
            \begin{byCase}
            \caseOf{$a_1 = a$.} Thus $a_1\mathrel{R}x$.
            \caseOf{$a_1\mathrel{R}a$.}
                Thus $a_1\mathrel{R}x$ by \hyperref[transitive]{transitivity}.
            \end{byCase}
        \end{subproof}
        Thus $x\in \intervalopenclosedrel{R}{X}{a_1}{b}$ by \cref{intervalopenclosed_rel}.
    \end{subproof}
    Thus $\intervalopenclosedrel{R}{X}{a}{b} \subseteq \intervalopenclosedrel{R}{X}{a_1}{b}$ by \cref{subseteq}.
\end{proof}

% from §14 helper lemma 7: basis element at a smallest element
\begin{lemma}\label{orderbasis_elem_at_smallest}
    Assume $R$ is a simple order relation on $X$.
    Suppose $x$ is a smallest element of $X$ with respect to $R$.
    Suppose $B\in\orderBasis{R}{X}$.
    Suppose $x\in B$.
    Then there exists $b\in X$ such that $x\mathrel{R}b$ and $B = \intervalclosedopenrel{R}{X}{x}{b}$.
\end{lemma}
\begin{proof}
    $R$ is asymmetric by \cref{simple_order_relation}.
    $B\in\pow{X}$ and
        $((\exists a, b\in X. a\mathrel{R}b \land B = \intervalopenrel{R}{X}{a}{b}) \lor
        (\exists b\in X. \exists a_0\in X. (\forall x\in X. a_0\mathrel{R}x \lor a_0 = x) \land B = \intervalclosedopenrel{R}{X}{a_0}{b}) \lor
        (\exists a\in X. \exists b_0\in X. (\forall x\in X. x\mathrel{R}b_0 \lor x = b_0) \land B = \intervalopenclosedrel{R}{X}{a}{b_0}))$
        by \cref{order_basis}.
    \begin{byCase}
    \caseOf{$\exists a, b\in X. a\mathrel{R}b \land B = \intervalopenrel{R}{X}{a}{b}$.}
        Take $a, b\in X$ such that $a\mathrel{R}b$ and $B = \intervalopenrel{R}{X}{a}{b}$.
        Thus $a\mathrel{R}x$ by \cref{intervalopen_rel}.
        $x\in X$ by \cref{smallest_element}.
        $x\mathrel{R}a$ or $x = a$ by \cref{smallest_element}.
        \begin{byCase}
        \caseOf{$x = a$.}
            Suppose not.
            Thus $a\mathrel{R}a$ by \cref{intervalopen_rel}.
            Contradiction by \cref{asymmetric}.
        \caseOf{$x\mathrel{R}a$.}
            Suppose not.
            Contradiction by \cref{asymmetric}.
        \end{byCase}
    \caseOf{$\exists b\in X. \exists a_0\in X. (\forall x\in X. a_0\mathrel{R}x \lor a_0 = x) \land B = \intervalclosedopenrel{R}{X}{a_0}{b}$.}
        Take $a_0, b\in X$ such that
            $(\forall x\in X. a_0\mathrel{R}x \lor a_0 = x)$ and $B = \intervalclosedopenrel{R}{X}{a_0}{b}$.
        Thus $x\in X$ and $(a_0\mathrel{R}x \lor x = a_0)$ and $x\mathrel{R}b$ by \cref{intervalclosedopen_rel}.
        Show that $a_0 = x$.
        \begin{subproof}
            Suppose not.
            Then $a_0 \neq x$.
            $a_0\mathrel{R}x$ by \cref{intervalclosedopen_rel}.
            $x\mathrel{R}a_0$ by \cref{smallest_element}.
            Contradiction by \cref{asymmetric}.
        \end{subproof}
        Thus $B = \intervalclosedopenrel{R}{X}{x}{b}$.
        Thus there exists $b\in X$ such that $x\mathrel{R}b$ and $B = \intervalclosedopenrel{R}{X}{x}{b}$.
    \caseOf{$\exists a\in X. \exists b_0\in X. (\forall x\in X. x\mathrel{R}b_0 \lor x = b_0) \land B = \intervalopenclosedrel{R}{X}{a}{b_0}$.}
        Take $a, b_0\in X$ such that
            $(\forall x\in X. x\mathrel{R}b_0 \lor x = b_0)$ and $B = \intervalopenclosedrel{R}{X}{a}{b_0}$.
        Thus $a\mathrel{R}x$ by \cref{intervalopenclosed_rel}.
        $x\in X$ by \cref{smallest_element}.
        $x\mathrel{R}a$ or $x = a$ by \cref{smallest_element}.
        \begin{byCase}
        \caseOf{$x = a$.}
            Suppose not.
            Thus $a\mathrel{R}a$ by \cref{intervalopenclosed_rel}.
            Contradiction by \cref{asymmetric}.
        \caseOf{$x\mathrel{R}a$.}
            Suppose not.
            Contradiction by \cref{asymmetric}.
        \end{byCase}
    \end{byCase}
\end{proof}

% from §14 helper lemma 8: basis element at a largest element
\begin{lemma}\label{orderbasis_elem_at_largest}
    Assume $R$ is a simple order relation on $X$.
    Suppose $x$ is a largest element of $X$ with respect to $R$.
    Suppose $B\in\orderBasis{R}{X}$.
    Suppose $x\in B$.
    Then there exists $a\in X$ such that $a\mathrel{R}x$ and $B = \intervalopenclosedrel{R}{X}{a}{x}$.
\end{lemma}
\begin{proof}
    $R$ is asymmetric by \cref{simple_order_relation}.
    $B\in\pow{X}$ and
        $((\exists a, b\in X. a\mathrel{R}b \land B = \intervalopenrel{R}{X}{a}{b}) \lor
        (\exists b\in X. \exists a_0\in X. (\forall x\in X. a_0\mathrel{R}x \lor a_0 = x) \land B = \intervalclosedopenrel{R}{X}{a_0}{b}) \lor
        (\exists a\in X. \exists b_0\in X. (\forall x\in X. x\mathrel{R}b_0 \lor x = b_0) \land B = \intervalopenclosedrel{R}{X}{a}{b_0}))$
        by \cref{order_basis}.
    \begin{byCase}
    \caseOf{$\exists a, b\in X. a\mathrel{R}b \land B = \intervalopenrel{R}{X}{a}{b}$.}
        Take $a, b\in X$ such that $a\mathrel{R}b$ and $B = \intervalopenrel{R}{X}{a}{b}$.
        Thus $x\mathrel{R}b$ by \cref{intervalopen_rel}.
        $x\in X$ by \cref{largest_element}.
        $b\mathrel{R}x$ or $b = x$ by \cref{largest_element}.
        \begin{byCase}
        \caseOf{$b = x$.}
            Suppose not.
            Thus $x\mathrel{R}x$ by \cref{intervalopen_rel}.
            Contradiction by \cref{asymmetric}.
        \caseOf{$b\mathrel{R}x$.}
            Suppose not.
            Contradiction by \cref{asymmetric}.
        \end{byCase}
    \caseOf{$\exists b\in X. \exists a_0\in X. (\forall x\in X. a_0\mathrel{R}x \lor a_0 = x) \land B = \intervalclosedopenrel{R}{X}{a_0}{b}$.}
        Take $a_0, b\in X$ such that
            $(\forall x\in X. a_0\mathrel{R}x \lor a_0 = x)$ and $B = \intervalclosedopenrel{R}{X}{a_0}{b}$.
        Thus $x\mathrel{R}b$ by \cref{intervalclosedopen_rel}.
        $x\in X$ by \cref{largest_element}.
        $b\mathrel{R}x$ or $b = x$ by \cref{largest_element}.
        \begin{byCase}
        \caseOf{$b = x$.}
            Suppose not.
            Thus $x\mathrel{R}x$ by \cref{intervalclosedopen_rel}.
            Contradiction by \cref{asymmetric}.
        \caseOf{$b\mathrel{R}x$.}
            Suppose not.
            Contradiction by \cref{asymmetric}.
        \end{byCase}
    \caseOf{$\exists a\in X. \exists b_0\in X. (\forall x\in X. x\mathrel{R}b_0 \lor x = b_0) \land B = \intervalopenclosedrel{R}{X}{a}{b_0}$.}
        Take $a, b_0\in X$ such that
            $(\forall x\in X. x\mathrel{R}b_0 \lor x = b_0)$ and $B = \intervalopenclosedrel{R}{X}{a}{b_0}$.
        Thus $x\in X$ and $a\mathrel{R}x$ and $(x\mathrel{R}b_0 \lor x = b_0)$ by \cref{intervalopenclosed_rel}.
        Show that $b_0 = x$.
        \begin{subproof}
            Suppose not.
            Then $b_0 \neq x$.
            $x\mathrel{R}b_0$ by \cref{intervalopenclosed_rel}.
            $b_0\mathrel{R}x$ by \cref{largest_element}.
            Contradiction by \cref{asymmetric}.
        \end{subproof}
        Thus $B = \intervalopenclosedrel{R}{X}{a}{x}$.
        Thus there exists $a\in X$ such that $a\mathrel{R}x$ and $B = \intervalopenclosedrel{R}{X}{a}{x}$.
    \end{byCase}
\end{proof}

% from §14 helper lemma 9: open interval inside a basis element
\begin{lemma}\label{orderbasis_elem_contains_open_interval}
    Assume $R$ is a simple order relation on $X$.
    Suppose $x\in X$.
    Suppose $x$ is not a smallest element of $X$ with respect to $R$.
    Suppose $x$ is not a largest element of $X$ with respect to $R$.
    Suppose $B\in\orderBasis{R}{X}$.
    Suppose $x\in B$.
    Then there exist $a,b\in X$ such that $a\mathrel{R}x$ and $x\mathrel{R}b$ and
        $\intervalopenrel{R}{X}{a}{b} \subseteq B$.
\end{lemma}
\begin{proof}
    $B\in\pow{X}$ and
        $((\exists a, b\in X. a\mathrel{R}b \land B = \intervalopenrel{R}{X}{a}{b}) \lor
        (\exists b\in X. \exists a_0\in X. (\forall x\in X. a_0\mathrel{R}x \lor a_0 = x) \land B = \intervalclosedopenrel{R}{X}{a_0}{b}) \lor
        (\exists a\in X. \exists b_0\in X. (\forall x\in X. x\mathrel{R}b_0 \lor x = b_0) \land B = \intervalopenclosedrel{R}{X}{a}{b_0}))$
        by \cref{order_basis}.
    \begin{byCase}
    \caseOf{$\exists a, b\in X. a\mathrel{R}b \land B = \intervalopenrel{R}{X}{a}{b}$.}
        Take $a, b\in X$ such that $a\mathrel{R}b$ and $B = \intervalopenrel{R}{X}{a}{b}$.
        Thus $a\mathrel{R}x$ and $x\mathrel{R}b$ by \cref{intervalopen_rel}.
        Thus $\intervalopenrel{R}{X}{a}{b} \subseteq B$ by \cref{subseteq_refl}.
    \caseOf{$\exists b\in X. \exists a_0\in X. (\forall x\in X. a_0\mathrel{R}x \lor a_0 = x) \land B = \intervalclosedopenrel{R}{X}{a_0}{b}$.}
        Take $a_0, b\in X$ such that
            $(\forall x\in X. a_0\mathrel{R}x \lor a_0 = x)$ and $B = \intervalclosedopenrel{R}{X}{a_0}{b}$.
        Thus $x\in X$ and $(a_0\mathrel{R}x \lor x = a_0)$ and $x\mathrel{R}b$ by \cref{intervalclosedopen_rel}.
        Thus $a_0\mathrel{R}x$ by \cref{smallest_element}.
        Show that $\intervalopenrel{R}{X}{a_0}{b} \subseteq B$.
        \begin{subproof}
            Show that for all $y$ such that $y\in \intervalopenrel{R}{X}{a_0}{b}$ we have $y\in B$.
            \begin{subproof}
                Fix $y$.
                Assume $y\in \intervalopenrel{R}{X}{a_0}{b}$.
                Thus $y\in X$ and $a_0\mathrel{R}y$ and $y\mathrel{R}b$ by \cref{intervalopen_rel}.
                Thus $y\in \intervalclosedopenrel{R}{X}{a_0}{b}$ by \cref{intervalclosedopen_rel}.
            \end{subproof}
            Thus $\intervalopenrel{R}{X}{a_0}{b} \subseteq B$ by \cref{subseteq}.
        \end{subproof}
        Thus $a_0\mathrel{R}x$ and $x\mathrel{R}b$ and $\intervalopenrel{R}{X}{a_0}{b} \subseteq B$.
    \caseOf{$\exists a\in X. \exists b_0\in X. (\forall x\in X. x\mathrel{R}b_0 \lor x = b_0) \land B = \intervalopenclosedrel{R}{X}{a}{b_0}$.}
        Take $a, b_0\in X$ such that
            $(\forall x\in X. x\mathrel{R}b_0 \lor x = b_0)$ and $B = \intervalopenclosedrel{R}{X}{a}{b_0}$.
        Thus $x\in X$ and $a\mathrel{R}x$ and $(x\mathrel{R}b_0 \lor x = b_0)$ by \cref{intervalopenclosed_rel}.
        Thus $x\mathrel{R}b_0$ by \cref{largest_element}.
        Show that $\intervalopenrel{R}{X}{a}{b_0} \subseteq B$.
        \begin{subproof}
            Show that for all $y$ such that $y\in \intervalopenrel{R}{X}{a}{b_0}$ we have $y\in B$.
            \begin{subproof}
                Fix $y$.
                Assume $y\in \intervalopenrel{R}{X}{a}{b_0}$.
                Thus $y\in X$ and $a\mathrel{R}y$ and $y\mathrel{R}b_0$ by \cref{intervalopen_rel}.
                Thus $y\in \intervalopenclosedrel{R}{X}{a}{b_0}$ by \cref{intervalopenclosed_rel}.
            \end{subproof}
            Thus $\intervalopenrel{R}{X}{a}{b_0} \subseteq B$ by \cref{subseteq}.
        \end{subproof}
        Thus $a\mathrel{R}x$ and $x\mathrel{R}b_0$ and $\intervalopenrel{R}{X}{a}{b_0} \subseteq B$.
    \end{byCase}
\end{proof}

% from §14 helper lemma 10: maximum of two elements in a simple order
\begin{lemma}\label{max_of_two_elements}
    Assume $R$ is a simple order relation on $X$.
    Suppose $a_1, a_2\in X$.
    Then there exists $a\in X$ such that
        $(a = a_1 \lor a = a_2) \land
        (a_1\mathrel{R}a \lor a_1 = a) \land
        (a_2\mathrel{R}a \lor a_2 = a)$.
\end{lemma}
\begin{proof}
    \begin{byCase}
    \caseOf{$a_1 = a_2$.}
        Let $a = a_1$.
        Thus $a\in X$.
        Thus $a = a_1$ or $a = a_2$.
        Thus $a_1\mathrel{R}a$ or $a_1 = a$.
        Thus $a_2\mathrel{R}a$ or $a_2 = a$.
        Thus there exists $a\in X$ such that
            $(a = a_1 \lor a = a_2) \land
            (a_1\mathrel{R}a \lor a_1 = a) \land
            (a_2\mathrel{R}a \lor a_2 = a)$.
    \caseOf{$a_1 \neq a_2$ and $a_1\mathrel{R}a_2$.}
        Let $a = a_2$.
        Thus $a\in X$.
        Thus $a = a_1$ or $a = a_2$.
        Thus $a_1\mathrel{R}a$ or $a_1 = a$.
        Thus $a_2\mathrel{R}a$ or $a_2 = a$.
        Thus there exists $a\in X$ such that
            $(a = a_1 \lor a = a_2) \land
            (a_1\mathrel{R}a \lor a_1 = a) \land
            (a_2\mathrel{R}a \lor a_2 = a)$.
    \caseOf{$a_1 \neq a_2$ and $a_2\mathrel{R}a_1$.}
        Let $a = a_1$.
        Thus $a\in X$.
        Thus $a = a_1$ or $a = a_2$.
        Thus $a_1\mathrel{R}a$ or $a_1 = a$.
        Thus $a_2\mathrel{R}a$ or $a_2 = a$.
        Thus there exists $a\in X$ such that
            $(a = a_1 \lor a = a_2) \land
            (a_1\mathrel{R}a \lor a_1 = a) \land
            (a_2\mathrel{R}a \lor a_2 = a)$.
    \end{byCase}
\end{proof}

% from §14 helper lemma 11: minimum of two elements in a simple order
\begin{lemma}\label{min_of_two_elements}
    Assume $R$ is a simple order relation on $X$.
    Suppose $b_1, b_2\in X$.
    Then there exists $b\in X$ such that
        $(b = b_1 \lor b = b_2) \land
        (b\mathrel{R}b_1 \lor b = b_1) \land
        (b\mathrel{R}b_2 \lor b = b_2)$.
\end{lemma}
\begin{proof}
    \begin{byCase}
    \caseOf{$b_1 = b_2$.}
        Let $b = b_1$.
        Thus $b\in X$.
        Thus $b = b_1$ or $b = b_2$.
        Thus $b\mathrel{R}b_1$ or $b = b_1$.
        Thus $b\mathrel{R}b_2$ or $b = b_2$.
        Thus there exists $b\in X$ such that
            $(b = b_1 \lor b = b_2) \land
            (b\mathrel{R}b_1 \lor b = b_1) \land
            (b\mathrel{R}b_2 \lor b = b_2)$.
    \caseOf{$b_1 \neq b_2$ and $b_1\mathrel{R}b_2$.}
        Let $b = b_1$.
        Thus $b\in X$.
        Thus $b = b_1$ or $b = b_2$.
        Thus $b\mathrel{R}b_1$ or $b = b_1$.
        Thus $b\mathrel{R}b_2$ or $b = b_2$.
        Thus there exists $b\in X$ such that
            $(b = b_1 \lor b = b_2) \land
            (b\mathrel{R}b_1 \lor b = b_1) \land
            (b\mathrel{R}b_2 \lor b = b_2)$.
    \caseOf{$b_1 \neq b_2$ and $b_2\mathrel{R}b_1$.}
        Let $b = b_2$.
        Thus $b\in X$.
        Thus $b = b_1$ or $b = b_2$.
        Thus $b\mathrel{R}b_1$ or $b = b_1$.
        Thus $b\mathrel{R}b_2$ or $b = b_2$.
        Thus there exists $b\in X$ such that
            $(b = b_1 \lor b = b_2) \land
            (b\mathrel{R}b_1 \lor b = b_1) \land
            (b\mathrel{R}b_2 \lor b = b_2)$.
    \end{byCase}
\end{proof}

% from §14 helper lemma 12: left endpoint from a disjunction
\begin{lemma}\label{rel_from_or_left}
    Suppose $a = a_1$ or $a = a_2$.
    Suppose $a_1\mathrel{R}x$ and $a_2\mathrel{R}x$.
    Then $a\mathrel{R}x$.
\end{lemma}
\begin{proof}
    \begin{byCase}
    \caseOf{$a = a_1$.} Thus $a\mathrel{R}x$.
    \caseOf{$a = a_2$.} Thus $a\mathrel{R}x$.
    \end{byCase}
\end{proof}

% from §14 helper lemma 13: right endpoint from a disjunction
\begin{lemma}\label{rel_from_or_right}
    Suppose $b = b_1$ or $b = b_2$.
    Suppose $x\mathrel{R}b_1$ and $x\mathrel{R}b_2$.
    Then $x\mathrel{R}b$.
\end{lemma}
\begin{proof}
    \begin{byCase}
    \caseOf{$b = b_1$.} Thus $x\mathrel{R}b$.
    \caseOf{$b = b_2$.} Thus $x\mathrel{R}b$.
    \end{byCase}
\end{proof}

% from §14 Item 11: order basis is a basis on X
\begin{lemma}\label{order_basis_is_basis}
    Assume $R$ is a simple order relation on $X$.
    Assume $X$ has more than one element.
    Then $\orderBasis{R}{X}$ is a basis on $X$.
\end{lemma}
\begin{proof}
    Show that $\orderBasis{R}{X}\subseteq\pow{X}$.
    \begin{subproof}
        Show that for all $B$ such that $B\in\orderBasis{R}{X}$ we have $B\in\pow{X}$.
        \begin{subproof}
            Fix $B$.
            Assume $B\in\orderBasis{R}{X}$.
            Thus $B\in\pow{X}$ by \cref{order_basis}.
        \end{subproof}
        Thus $\orderBasis{R}{X}\subseteq\pow{X}$ by \cref{subseteq}.
    \end{subproof}
    Show that for all $x$ such that $x\in X$ there exists $B_1\in\orderBasis{R}{X}$ such that $x\in B_1$.
    \begin{subproof}
        Fix $x$.
        Assume $x\in X$.
        \begin{byCase}
        \caseOf{$x$ is a smallest element of $X$ with respect to $R$.}
            Take $y\in X$ such that $y\neq x$ by \cref{more_than_one_element_has_other}.
            $x\mathrel{R}y$ by \cref{smallest_element}.
            Let $B_1 = \intervalclosedopenrel{R}{X}{x}{y}$.
            Show that $B_1\in\orderBasis{R}{X}$.
            \begin{subproof}
                Show that for all $z$ such that $z\in B_1$ we have $z\in X$.
                \begin{subproof}
                    Fix $z$.
                    Assume $z\in B_1$.
                    Thus $z\in X$ by \cref{intervalclosedopen_rel}.
                \end{subproof}
                Thus $B_1\subseteq X$ by \cref{subseteq}.
                Thus $B_1\in\pow{X}$ by \cref{powerset_intro}.
                Let $a_0 = x$.
                Let $b = y$.
                Thus $a_0\in X$ and $b\in X$.
                Thus for all $z\in X$ we have $a_0\mathrel{R}z$ or $a_0 = z$ by \cref{smallest_element}.
                Thus $B_1 = \intervalclosedopenrel{R}{X}{a_0}{b}$.
                Thus there exists $b\in X$ such that there exists $a_0\in X$ such that
                    $(\forall z\in X. a_0\mathrel{R}z \lor a_0 = z) \land
                    B_1 = \intervalclosedopenrel{R}{X}{a_0}{b}$.
                Thus $B_1\in\orderBasis{R}{X}$ by \cref{order_basis}.
            \end{subproof}
            Thus $x\in B_1$ by \cref{intervalclosedopen_rel}.
        \caseOf{$x$ is not a smallest element of $X$ with respect to $R$.}
            \begin{byCase}
            \caseOf{$x$ is a largest element of $X$ with respect to $R$.}
                Take $y\in X$ such that $y\neq x$ by \cref{more_than_one_element_has_other}.
                $y\mathrel{R}x$ by \cref{largest_element}.
                Let $B_1 = \intervalopenclosedrel{R}{X}{y}{x}$.
                Show that $B_1\in\orderBasis{R}{X}$.
                \begin{subproof}
                    Show that for all $z$ such that $z\in B_1$ we have $z\in X$.
                    \begin{subproof}
                        Fix $z$.
                        Assume $z\in B_1$.
                        Thus $z\in X$ by \cref{intervalopenclosed_rel}.
                    \end{subproof}
                    Thus $B_1\subseteq X$ by \cref{subseteq}.
                    Thus $B_1\in\pow{X}$ by \cref{powerset_intro}.
                    Let $a = y$.
                    Let $b_0 = x$.
                    Thus $a\in X$ and $b_0\in X$.
                    Thus for all $z\in X$ we have $z\mathrel{R}b_0$ or $z = b_0$ by \cref{largest_element}.
                    Thus $B_1 = \intervalopenclosedrel{R}{X}{a}{b_0}$.
                    Thus there exists $a\in X$ such that there exists $b_0\in X$ such that
                        $(\forall z\in X. z\mathrel{R}b_0 \lor z = b_0) \land
                        B_1 = \intervalopenclosedrel{R}{X}{a}{b_0}$.
                    Thus $B_1\in\orderBasis{R}{X}$ by \cref{order_basis}.
                \end{subproof}
                Thus $x\in B_1$ by \cref{intervalopenclosed_rel}.
            \caseOf{$x$ is not a largest element of $X$ with respect to $R$.}
                Take $a\in X$ such that $a\mathrel{R}x$ by \cref{non_smallest_has_predecessor}.
                Take $b\in X$ such that $x\mathrel{R}b$ by \cref{non_largest_has_successor}.
                Let $B_1 = \intervalopenrel{R}{X}{a}{b}$.
                Show that $B_1\in\orderBasis{R}{X}$.
                \begin{subproof}
                    Show that for all $z$ such that $z\in B_1$ we have $z\in X$.
                    \begin{subproof}
                        Fix $z$.
                        Assume $z\in B_1$.
                        Thus $z\in X$ by \cref{intervalopen_rel}.
                    \end{subproof}
                    Thus $B_1\subseteq X$ by \cref{subseteq}.
                    Thus $B_1\in\pow{X}$ by \cref{powerset_intro}.
                    $R$ is transitive by \cref{simple_order_relation}.
                    Thus $a\mathrel{R}b$ by \hyperref[transitive]{transitivity}.
                    Thus there exist $a,b\in X$ such that $a\mathrel{R}b$ and $B_1 = \intervalopenrel{R}{X}{a}{b}$.
                    Thus $B_1\in\orderBasis{R}{X}$ by \cref{order_basis}.
                \end{subproof}
                Thus $x\in B_1$ by \cref{intervalopen_rel}.
            \end{byCase}
        \end{byCase}
    \end{subproof}
    Show that for all $B_1,B_2,x$ such that $B_1\in\orderBasis{R}{X}$ and $B_2\in\orderBasis{R}{X}$ and $x\in B_1\inter B_2$
        there exists $B_3\in\orderBasis{R}{X}$ such that $x\in B_3$ and $B_3\subseteq B_1\inter B_2$.
    \begin{subproof}
        Fix $B_1$.
        Fix $B_2$.
        Fix $x$.
        Assume $B_1\in\orderBasis{R}{X}$ and $B_2\in\orderBasis{R}{X}$ and $x\in B_1\inter B_2$.
        $x\in B_1$ and $x\in B_2$ by \cref{inter_elim_left,inter_elim_right}.
        Thus $B_1\in\pow{X}$ by \cref{order_basis}.
        Thus $B_1\subseteq X$ by \cref{powerset_elim}.
        Thus $x\in X$ by \cref{subseteq}.
        \begin{byCase}
        \caseOf{$x$ is a smallest element of $X$ with respect to $R$.}
            Take $b_1\in X$ such that $x\mathrel{R}b_1$ and $B_1 = \intervalclosedopenrel{R}{X}{x}{b_1}$ by \cref{orderbasis_elem_at_smallest}.
            Take $b_2\in X$ such that $x\mathrel{R}b_2$ and $B_2 = \intervalclosedopenrel{R}{X}{x}{b_2}$ by \cref{orderbasis_elem_at_smallest}.
            $R$ is connex on $X$ by \cref{simple_order_relation}.
            \begin{byCase}
            \caseOf{$b_1 = b_2$.}
                Let $b = b_1$.
                Let $B_3 = \intervalclosedopenrel{R}{X}{x}{b}$.
                Thus $B_3\subseteq B_1$ by \cref{subseteq_refl}.
                Thus $B_3\subseteq B_2$ by \cref{subseteq_refl}.
                Thus $B_3\subseteq B_1\inter B_2$ by \cref{subseteq_inter_iff}.
                Thus $B_3\in\orderBasis{R}{X}$ by \cref{order_basis}.
                Thus $x\in B_3$ by \cref{intervalclosedopen_rel}.
            \caseOf{$b_1 \neq b_2$.}
                $b_1\mathrel{R}b_2$ or $b_2\mathrel{R}b_1$ by \cref{connex}.
                \begin{byCase}
                \caseOf{$b_1\mathrel{R}b_2$.}
                    Let $b = b_1$.
                    Let $B_3 = \intervalclosedopenrel{R}{X}{x}{b}$.
                    $B_3\subseteq B_1$ by \cref{subseteq_refl}.
                    $R$ is transitive by \cref{simple_order_relation}.
                    $x\in X$ by \cref{smallest_element}.
                    Thus $b\in X$ and $b_2\in X$.
                    Thus $b\mathrel{R}b_2$.
                    Thus $b\mathrel{R}b_2$ or $b = b_2$.
                    $B_3\subseteq B_2$ by \cref{intervalclosedopenrel_subset_by_right}.
                    Thus $B_3\subseteq B_1\inter B_2$ by \cref{subseteq_inter_iff}.
                    Thus $B_3\in\orderBasis{R}{X}$ by \cref{order_basis}.
                    Thus $x\in B_3$ by \cref{intervalclosedopen_rel}.
                \caseOf{$b_2\mathrel{R}b_1$.}
                    Let $b = b_2$.
                    Let $B_3 = \intervalclosedopenrel{R}{X}{x}{b}$.
                    $B_3\subseteq B_2$ by \cref{subseteq_refl}.
                    $R$ is transitive by \cref{simple_order_relation}.
                    $x\in X$ by \cref{smallest_element}.
                    Thus $b\in X$ and $b_1\in X$.
                    Thus $b\mathrel{R}b_1$.
                    Thus $b\mathrel{R}b_1$ or $b = b_1$.
                    $B_3\subseteq B_1$ by \cref{intervalclosedopenrel_subset_by_right}.
                    Thus $B_3\subseteq B_1\inter B_2$ by \cref{subseteq_inter_iff}.
                    Thus $B_3\in\orderBasis{R}{X}$ by \cref{order_basis}.
                    Thus $x\in B_3$ by \cref{intervalclosedopen_rel}.
                \end{byCase}
            \end{byCase}
        \caseOf{$x$ is not a smallest element of $X$ with respect to $R$.}
            \begin{byCase}
            \caseOf{$x$ is a largest element of $X$ with respect to $R$.}
                Take $a_1\in X$ such that $a_1\mathrel{R}x$ and $B_1 = \intervalopenclosedrel{R}{X}{a_1}{x}$ by \cref{orderbasis_elem_at_largest}.
                Take $a_2\in X$ such that $a_2\mathrel{R}x$ and $B_2 = \intervalopenclosedrel{R}{X}{a_2}{x}$ by \cref{orderbasis_elem_at_largest}.
                $R$ is connex on $X$ by \cref{simple_order_relation}.
                \begin{byCase}
                \caseOf{$a_1 = a_2$.}
                    Let $a = a_1$.
                    Let $B_3 = \intervalopenclosedrel{R}{X}{a}{x}$.
                    Thus $B_3\subseteq B_1$ by \cref{subseteq_refl}.
                    Thus $B_3\subseteq B_2$ by \cref{subseteq_refl}.
                    Thus $B_3\subseteq B_1\inter B_2$ by \cref{subseteq_inter_iff}.
                    Thus $B_3\in\orderBasis{R}{X}$ by \cref{order_basis}.
                    Thus $x\in B_3$ by \cref{intervalopenclosed_rel}.
                \caseOf{$a_1 \neq a_2$.}
                    $a_1\mathrel{R}a_2$ or $a_2\mathrel{R}a_1$ by \cref{connex}.
                    \begin{byCase}
                    \caseOf{$a_1\mathrel{R}a_2$.}
                        Let $a = a_2$.
                        Let $B_3 = \intervalopenclosedrel{R}{X}{a}{x}$.
                        $B_3\subseteq B_2$ by \cref{subseteq_refl}.
                        $R$ is transitive by \cref{simple_order_relation}.
                        $x\in X$ by \cref{largest_element}.
                        Thus $a\in X$ and $a_1\in X$.
                        Thus $a_1\mathrel{R}a$.
                        Thus $a_1\mathrel{R}a$ or $a_1 = a$.
                        $B_3\subseteq B_1$ by \cref{intervalopenclosedrel_subset_by_left}.
                        Thus $B_3\subseteq B_1\inter B_2$ by \cref{subseteq_inter_iff}.
                        Thus $B_3\in\orderBasis{R}{X}$ by \cref{order_basis}.
                        Thus $x\in B_3$ by \cref{intervalopenclosed_rel}.
                    \caseOf{$a_2\mathrel{R}a_1$.}
                        Let $a = a_1$.
                        Let $B_3 = \intervalopenclosedrel{R}{X}{a}{x}$.
                        $B_3\subseteq B_1$ by \cref{subseteq_refl}.
                        $R$ is transitive by \cref{simple_order_relation}.
                        $x\in X$ by \cref{largest_element}.
                        Thus $a\in X$ and $a_2\in X$.
                        Thus $a_2\mathrel{R}a$.
                        Thus $a_2\mathrel{R}a$ or $a_2 = a$.
                        $B_3\subseteq B_2$ by \cref{intervalopenclosedrel_subset_by_left}.
                        Thus $B_3\subseteq B_1\inter B_2$ by \cref{subseteq_inter_iff}.
                        Thus $B_3\in\orderBasis{R}{X}$ by \cref{order_basis}.
                        Thus $x\in B_3$ by \cref{intervalopenclosed_rel}.
                    \end{byCase}
                \end{byCase}
            \caseOf{$x$ is not a largest element of $X$ with respect to $R$.}
                Take $a_1,b_1\in X$ such that $a_1\mathrel{R}x$ and $x\mathrel{R}b_1$ and
                    $\intervalopenrel{R}{X}{a_1}{b_1} \subseteq B_1$ by \cref{orderbasis_elem_contains_open_interval}.
                Take $a_2,b_2\in X$ such that $a_2\mathrel{R}x$ and $x\mathrel{R}b_2$ and
                    $\intervalopenrel{R}{X}{a_2}{b_2} \subseteq B_2$ by \cref{orderbasis_elem_contains_open_interval}.
                Take $a\in X$ such that
                    $(a = a_1 \lor a = a_2) \land
                    (a_1\mathrel{R}a \lor a_1 = a) \land
                    (a_2\mathrel{R}a \lor a_2 = a)$
                    by \cref{max_of_two_elements}.
                Take $b\in X$ such that
                    $(b = b_1 \lor b = b_2) \land
                    (b\mathrel{R}b_1 \lor b = b_1) \land
                    (b\mathrel{R}b_2 \lor b = b_2)$
                    by \cref{min_of_two_elements}.
                Let $B_3 = \intervalopenrel{R}{X}{a}{b}$.
                Show that $B_3\in\orderBasis{R}{X}$.
                \begin{subproof}
                    Show that for all $z$ such that $z\in B_3$ we have $z\in X$.
                    \begin{subproof}
                        Fix $z$.
                        Assume $z\in B_3$.
                        Thus $z\in X$ by \cref{intervalopen_rel}.
                    \end{subproof}
                    Thus $B_3\subseteq X$ by \cref{subseteq}.
                    Thus $B_3\in\pow{X}$ by \cref{powerset_intro}.
                    $R$ is transitive by \cref{simple_order_relation}.
                    Thus $a\mathrel{R}b$ by \hyperref[transitive]{transitivity}.
                    Thus $B_3\in\orderBasis{R}{X}$ by \cref{order_basis}.
                \end{subproof}
                Thus $a = a_1$ or $a = a_2$.
                Thus $a\mathrel{R}x$ by \cref{rel_from_or_left}.
                Thus $b = b_1$ or $b = b_2$.
                Thus $x\mathrel{R}b$ by \cref{rel_from_or_right}.
                Thus $x\in B_3$ by \cref{intervalopen_rel}.
                Show that $B_3\subseteq B_1\inter B_2$.
                \begin{subproof}
                    Show that $B_3\subseteq B_1$.
                    \begin{subproof}
                        $R$ is transitive by \cref{simple_order_relation}.
                        Thus $a\in X$.
                        Thus $a_1\in X$.
                        Thus $b\in X$.
                        Thus $b_1\in X$.
                        Thus $a_1\mathrel{R}a$ or $a_1 = a$.
                        Thus $b\mathrel{R}b_1$ or $b = b_1$.
                        $B_3\subseteq \intervalopenrel{R}{X}{a_1}{b_1}$ by \cref{intervalopenrel_subset_by_endpoints}.
                        Thus $B_3\subseteq B_1$ by \cref{subseteq_transitive}.
                    \end{subproof}
                    Show that $B_3\subseteq B_2$.
                    \begin{subproof}
                        $R$ is transitive by \cref{simple_order_relation}.
                        Thus $a\in X$.
                        Thus $a_2\in X$.
                        Thus $b\in X$.
                        Thus $b_2\in X$.
                        Thus $a_2\mathrel{R}a$ or $a_2 = a$.
                        Thus $b\mathrel{R}b_2$ or $b = b_2$.
                        $B_3\subseteq \intervalopenrel{R}{X}{a_2}{b_2}$ by \cref{intervalopenrel_subset_by_endpoints}.
                        Thus $B_3\subseteq B_2$ by \cref{subseteq_transitive}.
                    \end{subproof}
                    Thus $B_3\subseteq B_1\inter B_2$ by \cref{subseteq_inter_iff}.
                \end{subproof}
            \end{byCase}
        \end{byCase}
    \end{subproof}
    Thus $\orderBasis{R}{X}$ is a basis on $X$ by \cref{basis_on}.
\end{proof}

% from §14 Item 12: open right ray
\begin{definition}\label{openray_right}
    $\openrayright{R}{X}{a} = \{x\in X \mid a\mathrel{R}x\}$.
\end{definition}

% from §14 Item 13: open left ray
\begin{definition}\label{openray_left}
    $\openrayleft{R}{X}{a} = \{x\in X \mid x\mathrel{R}a\}$.
\end{definition}

% from §14 Item 14: closed right ray
\begin{definition}\label{closedray_right}
    $\closedrayright{R}{X}{a} = \{x\in X \mid a\mathrel{R}x \lor x = a\}$.
\end{definition}

% from §14 Item 15: closed left ray
\begin{definition}\label{closedray_left}
    $\closedrayleft{R}{X}{a} = \{x\in X \mid x\mathrel{R}a \lor x = a\}$.
\end{definition}

% from §14 Item 16: open rays
\begin{definition}\label{open_rays}
    $\openRays{R}{X} = \{U\in\pow{X} \mid \exists a\in X. (U = \openrayright{R}{X}{a} \lor U = \openrayleft{R}{X}{a})\}$.
\end{definition}

% from §14 helper lemma 14: open interval as intersection of open rays
\begin{lemma}\label{intervalopenrel_as_intersection_openrays}
    $\intervalopenrel{R}{X}{a}{b} = \openrayright{R}{X}{a} \inter \openrayleft{R}{X}{b}$.
\end{lemma}
\begin{proof}
    Show that for all $x$ if $x\in \intervalopenrel{R}{X}{a}{b}$ then $x\in \openrayright{R}{X}{a} \inter \openrayleft{R}{X}{b}$.
    \begin{subproof}
        Fix $x$.
        Assume $x\in \intervalopenrel{R}{X}{a}{b}$.
        Thus $x\in X$ and $a\mathrel{R}x$ and $x\mathrel{R}b$ by \cref{intervalopen_rel}.
        Thus $x\in \openrayright{R}{X}{a}$ by \cref{openray_right}.
        Thus $x\in \openrayleft{R}{X}{b}$ by \cref{openray_left}.
        Thus $x\in \openrayright{R}{X}{a} \inter \openrayleft{R}{X}{b}$ by \cref{inter}.
    \end{subproof}
    Show that for all $x$ if $x\in \openrayright{R}{X}{a} \inter \openrayleft{R}{X}{b}$ then $x\in \intervalopenrel{R}{X}{a}{b}$.
    \begin{subproof}
        Fix $x$.
        Assume $x\in \openrayright{R}{X}{a} \inter \openrayleft{R}{X}{b}$.
        Thus $x\in \openrayright{R}{X}{a}$ and $x\in \openrayleft{R}{X}{b}$ by \cref{inter}.
        Thus $x\in X$ and $a\mathrel{R}x$ by \cref{openray_right}.
        Thus $x\in X$ and $x\mathrel{R}b$ by \cref{openray_left}.
        Thus $x\in \intervalopenrel{R}{X}{a}{b}$ by \cref{intervalopen_rel}.
    \end{subproof}
    Follows by set extensionality.
\end{proof}

% from §14 helper lemma 15: closed-open interval at a smallest element is a left ray
\begin{lemma}\label{intervalclosedopenrel_smallest_eq_openrayleft}
    Assume $a_0$ is a smallest element of $X$ with respect to $R$.
    Then $\intervalclosedopenrel{R}{X}{a_0}{b} = \openrayleft{R}{X}{b}$.
\end{lemma}
\begin{proof}
    Show that for all $x$ if $x\in \intervalclosedopenrel{R}{X}{a_0}{b}$ then $x\in \openrayleft{R}{X}{b}$.
    \begin{subproof}
        Fix $x$.
        Assume $x\in \intervalclosedopenrel{R}{X}{a_0}{b}$.
        Thus $x\in X$ and $x\mathrel{R}b$ by \cref{intervalclosedopen_rel}.
        Thus $x\in \openrayleft{R}{X}{b}$ by \cref{openray_left}.
    \end{subproof}
    Show that for all $x$ if $x\in \openrayleft{R}{X}{b}$ then $x\in \intervalclosedopenrel{R}{X}{a_0}{b}$.
    \begin{subproof}
        Fix $x$.
        Assume $x\in \openrayleft{R}{X}{b}$.
        Thus $x\in X$ and $x\mathrel{R}b$ by \cref{openray_left}.
        $a_0\mathrel{R}x$ or $x = a_0$ by \cref{smallest_element}.
        Thus $x\in \intervalclosedopenrel{R}{X}{a_0}{b}$ by \cref{intervalclosedopen_rel}.
    \end{subproof}
    Follows by set extensionality.
\end{proof}

% from §14 helper lemma 16: open-closed interval at a largest element is a right ray
\begin{lemma}\label{intervalopenclosedrel_largest_eq_openrayright}
    Assume $b_0$ is a largest element of $X$ with respect to $R$.
    Then $\intervalopenclosedrel{R}{X}{a}{b_0} = \openrayright{R}{X}{a}$.
\end{lemma}
\begin{proof}
    Show that for all $x$ if $x\in \intervalopenclosedrel{R}{X}{a}{b_0}$ then $x\in \openrayright{R}{X}{a}$.
    \begin{subproof}
        Fix $x$.
        Assume $x\in \intervalopenclosedrel{R}{X}{a}{b_0}$.
        Thus $x\in X$ and $a\mathrel{R}x$ by \cref{intervalopenclosed_rel}.
        Thus $x\in \openrayright{R}{X}{a}$ by \cref{openray_right}.
    \end{subproof}
    Show that for all $x$ if $x\in \openrayright{R}{X}{a}$ then $x\in \intervalopenclosedrel{R}{X}{a}{b_0}$.
    \begin{subproof}
        Fix $x$.
        Assume $x\in \openrayright{R}{X}{a}$.
        Thus $x\in X$ and $a\mathrel{R}x$ by \cref{openray_right}.
        $x\mathrel{R}b_0$ or $x = b_0$ by \cref{largest_element}.
        Thus $x\in \intervalopenclosedrel{R}{X}{a}{b_0}$ by \cref{intervalopenclosed_rel}.
    \end{subproof}
    Follows by set extensionality.
\end{proof}

% from §14 helper lemma 17: open rays form a subbasis on X
\begin{lemma}\label{open_rays_subbasis_on}
    Assume $R$ is a simple order relation on $X$.
    Assume $X$ has more than one element.
    Then $\openRays{R}{X}$ is a subbasis on $X$.
\end{lemma}
\begin{proof}
    Let $S = \openRays{R}{X}$.
    Show that for all $U$ such that $U\in S$ we have $U\in\pow{X}$.
    \begin{subproof}
        Fix $U$.
        Assume $U\in S$.
        Take $a\in X$ such that $U = \openrayright{R}{X}{a}$ or $U = \openrayleft{R}{X}{a}$ by \cref{open_rays}.
        \begin{byCase}
        \caseOf{$U = \openrayright{R}{X}{a}$.}
            Show that for all $x$ such that $x\in U$ we have $x\in X$.
            \begin{subproof}
                Fix $x$.
                Assume $x\in U$.
                Thus $x\in X$ by \cref{openray_right}.
            \end{subproof}
            Thus $U\subseteq X$ by \cref{subseteq}.
            Thus $U\in\pow{X}$ by \cref{powerset_intro}.
        \caseOf{$U = \openrayleft{R}{X}{a}$.}
            Show that for all $x$ such that $x\in U$ we have $x\in X$.
            \begin{subproof}
                Fix $x$.
                Assume $x\in U$.
                Thus $x\in X$ by \cref{openray_left}.
            \end{subproof}
            Thus $U\subseteq X$ by \cref{subseteq}.
            Thus $U\in\pow{X}$ by \cref{powerset_intro}.
        \end{byCase}
    \end{subproof}
    Thus $S\subseteq\pow{X}$ by \cref{subseteq}.
    Thus $\unions{S}\subseteq X$ by \cref{unions_subseteq_of_powerset_is_subseteq}.
    Show that for all $x$ such that $x\in X$ we have $x\in \unions{S}$.
    \begin{subproof}
        Fix $x$.
        Assume $x\in X$.
        Take $y\in X$ such that $y\neq x$ by \cref{more_than_one_element_has_other}.
        $R$ is connex on $X$ by \cref{simple_order_relation}.
        $x\mathrel{R}y$ or $y\mathrel{R}x$ by \cref{connex}.
        \begin{byCase}
        \caseOf{$x\mathrel{R}y$.}
            Thus $x\in \openrayleft{R}{X}{y}$ by \cref{openray_left}.
            Show that for all $z$ such that $z\in \openrayleft{R}{X}{y}$ we have $z\in X$.
            \begin{subproof}
                Fix $z$.
                Assume $z\in \openrayleft{R}{X}{y}$.
                Thus $z\in X$ by \cref{openray_left}.
            \end{subproof}
            Thus $\openrayleft{R}{X}{y}\subseteq X$ by \cref{subseteq}.
            Thus $\openrayleft{R}{X}{y}\in\pow{X}$ by \cref{powerset_intro}.
            Thus $\openrayleft{R}{X}{y}\in S$ by \cref{open_rays}.
            Thus $x\in\unions{S}$ by \cref{unions_iff}.
        \caseOf{$y\mathrel{R}x$.}
            Thus $x\in \openrayright{R}{X}{y}$ by \cref{openray_right}.
            Show that for all $z$ such that $z\in \openrayright{R}{X}{y}$ we have $z\in X$.
            \begin{subproof}
                Fix $z$.
                Assume $z\in \openrayright{R}{X}{y}$.
                Thus $z\in X$ by \cref{openray_right}.
            \end{subproof}
            Thus $\openrayright{R}{X}{y}\subseteq X$ by \cref{subseteq}.
            Thus $\openrayright{R}{X}{y}\in\pow{X}$ by \cref{powerset_intro}.
            Thus $\openrayright{R}{X}{y}\in S$ by \cref{open_rays}.
            Thus $x\in\unions{S}$ by \cref{unions_iff}.
        \end{byCase}
    \end{subproof}
    Thus $X\subseteq\unions{S}$ by \cref{subseteq}.
    Thus $\unions{S} = X$ by \cref{subseteq_antisymmetric}.
    Thus $S$ is a subbasis on $X$ by \cref{subbasis_on}.
\end{proof}

% from §14 helper lemma 18: order basis elements are finite intersections of open rays
\begin{lemma}\label{order_basis_in_finite_intersections}
    Assume $R$ is a simple order relation on $X$.
    Assume $X$ has more than one element.
    Suppose $B\in\orderBasis{R}{X}$.
    Then $B\in\finiteIntersections{\openRays{R}{X}}$.
\end{lemma}
\begin{proof}
    Let $S = \openRays{R}{X}$.
    $S$ is a subbasis on $X$ by \cref{open_rays_subbasis_on}.
    Thus $\unions{S} = X$ by \cref{subbasis_on}.
    $B\in\pow{X}$ by \cref{order_basis}.
    Thus $B\subseteq X$ by \cref{pow_iff}.
    Thus $B\subseteq\unions{S}$.
    Thus $B\in\pow{\unions{S}}$ by \cref{powerset_intro}.
    $B\in\pow{X}$ and
        $((\exists a, b\in X. a\mathrel{R}b \land B = \intervalopenrel{R}{X}{a}{b}) \lor
        (\exists b\in X. \exists a_0\in X. (\forall x\in X. a_0\mathrel{R}x \lor a_0 = x) \land B = \intervalclosedopenrel{R}{X}{a_0}{b}) \lor
        (\exists a\in X. \exists b_0\in X. (\forall x\in X. x\mathrel{R}b_0 \lor x = b_0) \land B = \intervalopenclosedrel{R}{X}{a}{b_0}))$
        by \cref{order_basis}.
    \begin{byCase}
    \caseOf{$\exists a, b\in X. a\mathrel{R}b \land B = \intervalopenrel{R}{X}{a}{b}$.}
        Take $a, b\in X$ such that $a\mathrel{R}b$ and $B = \intervalopenrel{R}{X}{a}{b}$.
        Thus $B = \openrayright{R}{X}{a} \inter \openrayleft{R}{X}{b}$ by \cref{intervalopenrel_as_intersection_openrays}.
        Let $F = \{\openrayright{R}{X}{a}, \openrayleft{R}{X}{b}\}$.
        Show that $F\subseteq S$.
        \begin{subproof}
            Show that $\openrayright{R}{X}{a}\in S$.
            \begin{subproof}
                Show that for all $x$ such that $x\in \openrayright{R}{X}{a}$ we have $x\in X$.
                \begin{subproof}
                    Fix $x$.
                    Assume $x\in \openrayright{R}{X}{a}$.
                    Thus $x\in X$ by \cref{openray_right}.
                \end{subproof}
                Thus $\openrayright{R}{X}{a}\subseteq X$ by \cref{subseteq}.
                Thus $\openrayright{R}{X}{a}\in\pow{X}$ by \cref{powerset_intro}.
                Thus $\openrayright{R}{X}{a}\in S$ by \cref{open_rays}.
            \end{subproof}
            Show that $\openrayleft{R}{X}{b}\in S$.
            \begin{subproof}
                Show that for all $x$ such that $x\in \openrayleft{R}{X}{b}$ we have $x\in X$.
                \begin{subproof}
                    Fix $x$.
                    Assume $x\in \openrayleft{R}{X}{b}$.
                    Thus $x\in X$ by \cref{openray_left}.
                \end{subproof}
                Thus $\openrayleft{R}{X}{b}\subseteq X$ by \cref{subseteq}.
                Thus $\openrayleft{R}{X}{b}\in\pow{X}$ by \cref{powerset_intro}.
                Thus $\openrayleft{R}{X}{b}\in S$ by \cref{open_rays}.
            \end{subproof}
            Thus $F\subseteq S$ by \cref{subseteq,upair_iff}.
        \end{subproof}
        $F$ is finite by \cref{pair_is_finite}.
        $\openrayright{R}{X}{a}\in F$ by \cref{upair_intro_left}.
        Thus $F$ is inhabited.
        Thus $\inters{F} = B$ by \cref{inter_as_inters}.
        Thus $B\in\finiteIntersections{S}$ by \cref{finite_intersections}.
    \caseOf{$\exists b\in X. \exists a_0\in X. (\forall x\in X. a_0\mathrel{R}x \lor a_0 = x) \land B = \intervalclosedopenrel{R}{X}{a_0}{b}$.}
        Take $a_0, b\in X$ such that
            $(\forall x\in X. a_0\mathrel{R}x \lor a_0 = x)$ and $B = \intervalclosedopenrel{R}{X}{a_0}{b}$.
        Thus $a_0$ is a smallest element of $X$ with respect to $R$ by \cref{smallest_element}.
        Thus $B = \openrayleft{R}{X}{b}$ by \cref{intervalclosedopenrel_smallest_eq_openrayleft}.
        Let $F = \{\openrayleft{R}{X}{b}\}$.
        Show that $\openrayleft{R}{X}{b}\in S$.
        \begin{subproof}
            Show that for all $x$ such that $x\in \openrayleft{R}{X}{b}$ we have $x\in X$.
            \begin{subproof}
                Fix $x$.
                Assume $x\in \openrayleft{R}{X}{b}$.
                Thus $x\in X$ by \cref{openray_left}.
            \end{subproof}
            Thus $\openrayleft{R}{X}{b}\subseteq X$ by \cref{subseteq}.
            Thus $\openrayleft{R}{X}{b}\in\pow{X}$ by \cref{powerset_intro}.
            Thus $\openrayleft{R}{X}{b}\in S$ by \cref{open_rays}.
        \end{subproof}
        Thus $F\subseteq S$ by \cref{singleton_subset_intro}.
        Show that $F$ is finite.
        \begin{subproof}
            Let $G = \{\openrayleft{R}{X}{b},\openrayleft{R}{X}{b}\}$.
            $G$ is finite by \cref{pair_is_finite}.
            $\openrayleft{R}{X}{b}\in G$ by \cref{upair_intro_left}.
            Thus $F\subseteq G$ by \cref{singleton_subset_intro}.
            Thus $F$ is finite by \cref{subseteq_finite_is_finite}.
        \end{subproof}
        Thus $F$ is inhabited.
        Thus $\inters{F} = B$ by \cref{inters_singleton}.
        Thus $B\in\finiteIntersections{S}$ by \cref{finite_intersections}.
    \caseOf{$\exists a\in X. \exists b_0\in X. (\forall x\in X. x\mathrel{R}b_0 \lor x = b_0) \land B = \intervalopenclosedrel{R}{X}{a}{b_0}$.}
        Take $a, b_0\in X$ such that
            $(\forall x\in X. x\mathrel{R}b_0 \lor x = b_0)$ and $B = \intervalopenclosedrel{R}{X}{a}{b_0}$.
        Thus $b_0$ is a largest element of $X$ with respect to $R$ by \cref{largest_element}.
        Thus $B = \openrayright{R}{X}{a}$ by \cref{intervalopenclosedrel_largest_eq_openrayright}.
        Let $F = \{\openrayright{R}{X}{a}\}$.
        Show that $\openrayright{R}{X}{a}\in S$.
        \begin{subproof}
            Show that for all $x$ such that $x\in \openrayright{R}{X}{a}$ we have $x\in X$.
            \begin{subproof}
                Fix $x$.
                Assume $x\in \openrayright{R}{X}{a}$.
                Thus $x\in X$ by \cref{openray_right}.
            \end{subproof}
            Thus $\openrayright{R}{X}{a}\subseteq X$ by \cref{subseteq}.
            Thus $\openrayright{R}{X}{a}\in\pow{X}$ by \cref{powerset_intro}.
            Thus $\openrayright{R}{X}{a}\in S$ by \cref{open_rays}.
        \end{subproof}
        Thus $F\subseteq S$ by \cref{singleton_subset_intro}.
        Show that $F$ is finite.
        \begin{subproof}
            Let $G = \{\openrayright{R}{X}{a},\openrayright{R}{X}{a}\}$.
            $G$ is finite by \cref{pair_is_finite}.
            $\openrayright{R}{X}{a}\in G$ by \cref{upair_intro_left}.
            Thus $F\subseteq G$ by \cref{singleton_subset_intro}.
            Thus $F$ is finite by \cref{subseteq_finite_is_finite}.
        \end{subproof}
        Thus $F$ is inhabited.
        Thus $\inters{F} = B$ by \cref{inters_singleton}.
        Thus $B\in\finiteIntersections{S}$ by \cref{finite_intersections}.
    \end{byCase}
\end{proof}

% from §14 helper lemma 19: finite intersections in a topology are open
\begin{lemma}\label{finite_intersections_open_in_topology}
    Assume $T$ is a topology on $X$.
    Suppose $F\subseteq T$.
    Suppose $F$ is finite and inhabited.
    Then $\inters{F}\in T$.
\end{lemma}
\begin{proof}
    Take $k$ such that $k$ is a natural number and there exists a bijection from $k$ to $F$ by \cref{finite}.
    Take $f$ such that $f$ is a bijection from $k$ to $F$.
    Let $A = \{ n\in\naturals \mid \text{for all $G,g$ such that $g$ is a bijection from $n$ to $G$ and $G\subseteq T$ and $G$ is inhabited we have $\inters{G}\in T$} \}$.
    Show that $A$ is an inductive set.
    \begin{subproof}
        Show that $\emptyset\in A$.
        \begin{subproof}
            Show that for all $G,g$ such that $g$ is a bijection from $\emptyset$ to $G$ and $G\subseteq T$ and $G$ is inhabited we have $\inters{G}\in T$.
            \begin{subproof}
                Fix $G$.
                Fix $g$.
                Assume $g$ is a bijection from $\emptyset$ to $G$.
                Assume $G\subseteq T$ and $G$ is inhabited.
                Show that $\inters{G}\in T$.
                \begin{subproof}
                    Suppose not.
                    Take $x$ such that $x\in G$.
                    $g\in\surj{\emptyset}{G}$ by \cref{bijections}.
                    Thus there exists $y\in\emptyset$ such that $g(y)=x$ by \cref{surj}.
                    Contradiction by \cref{emptyset}.
                \end{subproof}
            \end{subproof}
        \end{subproof}
        Show that for all $n\in A$ we have $\suc{n}\in A$.
        \begin{subproof}
            Fix $n$.
            Assume $n\in A$.
            Show that for all $G,g$ such that $g$ is a bijection from $\suc{n}$ to $G$ and $G\subseteq T$ and $G$ is inhabited we have $\inters{G}\in T$.
            \begin{subproof}
                Fix $G$.
                Fix $g$.
                Assume $g$ is a bijection from $\suc{n}$ to $G$.
                Assume $G\subseteq T$ and $G$ is inhabited.
                $g$ is a function by \cref{bijection_is_function}.
                $\dom{g} = \suc{n}$ by \cref{bijections_dom}.
                $n\in\suc{n}$ by \cref{suc_intro_self}.
                Thus $n\in\dom{g}$.
                Let $a = g(n)$.
                Let $G_1 = \img{g}{n}$.
                Show that for all $x$ such that $x\in n$ we have $x\in\suc{n}$.
                \begin{subproof}
                    Fix $x$.
                    Assume $x\in n$.
                    Thus $x\in\suc{n}$ by \cref{suc}.
                \end{subproof}
                Thus $n\subseteq \suc{n}$ by \cref{subseteq}.
                Thus $G_1\subseteq \img{g}{\suc{n}}$ by \cref{img_subseteq}.
                $\img{g}{\suc{n}} = \ran{g}$ by \cref{img_dom}.
                $g\in\surj{\suc{n}}{G}$ by \cref{bijections}.
                Thus $\ran{g} = G$ by \cref{ran_of_surj}.
                Thus $G_1\subseteq G$.
                Thus $G_1\subseteq T$ by \cref{subseteq_transitive}.
                $(n,a)\in g$ by \cref{function_apply_elim}.
                Thus $a\in\ran{g}$ by \cref{ran_iff}.
                Thus $a\in G$.
                Thus $a\in T$ by \cref{elem_subseteq}.
                Show that $G = G_1\union\{a\}$.
                \begin{subproof}
                    Show that $G\subseteq G_1\union\{a\}$.
                    \begin{subproof}
                        Show that for all $z$ such that $z\in G$ we have $z\in G_1\union\{a\}$.
                        \begin{subproof}
                            Fix $z$.
                            Assume $z\in G$.
                            $g\in\surj{\suc{n}}{G}$ by \cref{bijections}.
                            Thus there exists $x\in\suc{n}$ such that $g(x)=z$ by \cref{surj}.
                            $x\in n$ or $x = n$ by \cref{suc}.
                            \begin{byCase}
                            \caseOf{$x\in n$.}
                                $(x,z)\in g$ by \cref{function_apply_elim}.
                                Thus $z\in G_1$ by \cref{img_iff}.
                                Thus $z\in G_1\union\{a\}$ by \cref{union_intro_left}.
                            \caseOf{$x = n$.}
                                Thus $z = g(n)$.
                                Thus $z = a$.
                                Thus $z\in G_1\union\{a\}$ by \cref{union_intro_right,singleton_intro}.
                            \end{byCase}
                        \end{subproof}
                        Thus $G\subseteq G_1\union\{a\}$ by \cref{subseteq}.
                    \end{subproof}
                    $G_1\subseteq G$.
                    Thus $\{a\}\subseteq G$ by \cref{singleton_subset_intro}.
                    Thus $G_1\union\{a\}\subseteq G$ by \cref{union_subsets_is_subset}.
                    Thus $G = G_1\union\{a\}$ by \cref{subseteq_antisymmetric}.
                \end{subproof}
                \begin{byCase}
                \caseOf{$n = \emptyset$.}
                    Thus $G_1 = \emptyset$ by \cref{img_emptyset}.
                    Thus $G = \{a\}$ by \cref{union_emptyset,union_comm}.
                    Thus $\inters{G} = a$ by \cref{inters_singleton}.
                    Thus $\inters{G}\in T$.
                \caseOf{$n\neq\emptyset$.}
                    Thus $n$ is inhabited by \cref{nonempty_inhabited}.
                    Take $m$ such that $m\in n$.
                    $(m,g(m))\in g$ by \cref{function_apply_elim}.
                    Thus $g(m)\in G_1$ by \cref{img_iff}.
                    Thus $G_1$ is inhabited.
                    Let $h = \restrl{g}{n}$.
                    $h$ is a function by \cref{restrl_of_function_is_function}.
                    Show that for all $x$ such that $x\in n$ we have $x\in\dom{g}$.
                    \begin{subproof}
                        Fix $x$.
                        Assume $x\in n$.
                        Thus $x\in\suc{n}$ by \cref{suc}.
                        Thus $x\in\dom{g}$ by \cref{bijections_dom}.
                    \end{subproof}
                    Thus $n\subseteq\dom{g}$ by \cref{subseteq}.
                    Thus $\dom{h} = \dom{g}\inter n$ by \cref{dom_of_restrl_eq_inter}.
                    Thus $\dom{h} = n$ by \cref{inter_absorb_supseteq_left}.
                    Show that $h$ is a function to $G_1$.
                    \begin{subproof}
                        Show that for all $x\in\dom{h}$ we have $h(x)\in G_1$.
                        \begin{subproof}
                        Fix $x$.
                        Assume $x\in\dom{h}$.
                        Thus $x\in n$.
                        $(x,h(x))\in h$ by \cref{function_apply_elim}.
                        $h\subseteq g$ by \cref{restrl_subseteq}.
                        Thus $(x,h(x))\in g$ by \cref{elem_subseteq}.
                        Thus $h(x)\in G_1$ by \cref{img_iff}.
                        \end{subproof}
                        Thus $h$ is a function to $G_1$.
                    \end{subproof}
                    Thus $h\in\funs{n}{G_1}$ by \cref{funs_intro}.
                    Show that for all $y\in G_1$ there exists $x\in n$ such that $h(x) = y$.
                    \begin{subproof}
                        Fix $y$.
                        Assume $y\in G_1$.
                        Take $x\in n$ such that $(x,y)\in g$ by \cref{img_iff}.
                        Thus $x\in\dom{g}$ and $g(x) = y$ by \cref{function_apply_iff}.
                        Thus $h(x) = y$ by \cref{restrl_of_function_apply}.
                    \end{subproof}
                    Thus $h\in\surj{n}{G_1}$ by \cref{surj}.
                    Show that $h$ is injective.
                    \begin{subproof}
                        Show that for all $x,y$ such that $x\in\dom{h}$ and $y\in\dom{h}$ and $h(x) = h(y)$ we have $x = y$.
                        \begin{subproof}
                            Fix $x$.
                            Fix $y$.
                            Assume $x\in\dom{h}$ and $y\in\dom{h}$ and $h(x) = h(y)$.
                            Thus $x\in n$ and $y\in n$.
                            Thus $h(x) = g(x)$ and $h(y) = g(y)$ by \cref{restrl_of_function_apply}.
                            Thus $g(x) = g(y)$.
                            $g$ is injective by \cref{bijections}.
                            Thus $x = y$ by \cref{injective_function}.
                        \end{subproof}
                        Thus $h$ is injective by \cref{injective_function}.
                    \end{subproof}
                    Thus $h$ is a bijection from $n$ to $G_1$ by \cref{bijections}.
                    Thus $\inters{G_1}\in T$.
                    Thus $G = \{a\}\union G_1$ by \cref{union_comm}.
                    Thus $G = \cons{a}{G_1}$ by \cref{union_eq_cons}.
                    Thus $\inters{G} = a\inter\inters{G_1}$ by \cref{inters_cons}.
                    $T$ is closed under binary intersections by \cref{topology_on}.
                    Thus $a\inter\inters{G_1}\in T$.
                    Thus $\inters{G}\in T$.
                \end{byCase}
            \end{subproof}
        \end{subproof}
    \end{subproof}
    Thus $A$ is an inductive set.
    Thus $\naturals\subseteq A$ by \cref{num_naturals_smallest_inductive_set}.
    Thus $k\in\naturals$.
    Thus $k\in A$ by \cref{subseteq}.
    Thus for all $G,g$ such that $g$ is a bijection from $k$ to $G$ and $G\subseteq T$ and $G$ is inhabited we have $\inters{G}\in T$.
    Thus $\inters{F}\in T$.
\end{proof}

% from §14 Item 17: open rays form a subbasis
\begin{lemma}\label{open_rays_subbasis}
    Assume $R$ is a simple order relation on $X$.
    Assume $X$ has more than one element.
    Assume $T$ is the order topology on $X$ with respect to $R$.
    Then $\subbasisTopology{\openRays{R}{X}}{X} = T$.
\end{lemma}
\begin{proof}
    Let $S = \openRays{R}{X}$.
    Let $B = \orderBasis{R}{X}$.
    Let $T_1 = \subbasisTopology{S}{X}$.
    $T = \basisTopology{B}{X}$ by \cref{order_topology}.
    $B$ is a basis on $X$ by \cref{order_basis_is_basis}.
    Thus $B$ is a basis for $T$ on $X$ by \cref{basis_for_topology,order_topology}.
    $S$ is a subbasis on $X$ by \cref{open_rays_subbasis_on}.
    Thus $\finiteIntersections{S}$ is a basis on $X$ by \cref{subbasis_finite_intersections_basis}.
    Thus $\finiteIntersections{S}$ is a basis for $T_1$ on $X$ by \cref{basis_for_topology,topology_generated_by_subbasis}.
    Show that $T_1$ is finer than $T$ on $X$.
    \begin{subproof}
        Show that for all $x,B_2$ such that $x\in X$ and $B_2\in B$ and $x\in B_2$
            there exists $B_3\in\finiteIntersections{S}$ such that $x\in B_3$ and $B_3\subseteq B_2$.
        \begin{subproof}
            Fix $x$.
            Fix $B_2$.
            Assume $x\in X$ and $B_2\in B$ and $x\in B_2$.
            Thus $B_2\in\finiteIntersections{S}$ by \cref{order_basis_in_finite_intersections}.
            Let $B_3 = B_2$.
            Thus $x\in B_3$.
            Thus $B_3\subseteq B_2$ by \cref{subseteq_refl}.
        \end{subproof}
        Thus $T_1$ is finer than $T$ on $X$ by \cref{basis_finer_criterion}.
    \end{subproof}
    Thus $T\subseteq T_1$ by \cref{finer_topology}.
    Show that for all $U$ such that $U\in S$ we have $U\in T$.
    \begin{subproof}
        Fix $U$.
        Assume $U\in S$.
        Take $a\in X$ such that $U = \openrayright{R}{X}{a}$ or $U = \openrayleft{R}{X}{a}$ by \cref{open_rays}.
        \begin{byCase}
        \caseOf{$U = \openrayright{R}{X}{a}$.}
            \begin{byCase}
            \caseOf{$\exists b_0\in X. (\forall x\in X. x\mathrel{R}b_0 \lor x = b_0)$.}
                Take $b_0$ such that $b_0$ is a largest element of $X$ with respect to $R$.
                Thus $U = \intervalopenclosedrel{R}{X}{a}{b_0}$ by \cref{intervalopenclosedrel_largest_eq_openrayright}.
                Show that $U\in B$.
                \begin{subproof}
                    Show that for all $x$ such that $x\in U$ we have $x\in X$.
                    \begin{subproof}
                        Fix $x$.
                        Assume $x\in U$.
                        Thus $x\in X$ by \cref{intervalopenclosed_rel}.
                    \end{subproof}
                    Thus $U\subseteq X$ by \cref{subseteq}.
                    Thus $U\in\pow{X}$ by \cref{powerset_intro}.
                    Thus $a\in X$.
                    Thus $b_0\in X$ and $(\forall x\in X. x\mathrel{R}b_0 \lor x = b_0)$ by \cref{largest_element}.
                    Thus $\exists a\in X. \exists b_0\in X. (\forall x\in X. x\mathrel{R}b_0 \lor x = b_0) \land U = \intervalopenclosedrel{R}{X}{a}{b_0}$.
                    Thus $U\in B$ by \cref{order_basis}.
                \end{subproof}
                Thus $U\in T$ by \cref{basis_elem_open_in_generated,order_basis_is_basis}.
            \caseOf{it is wrong that there exists $b_0\in X$ such that for all $x\in X$ we have $x\mathrel{R}b_0$ or $x = b_0$.}
                Show that $U\in T$.
                \begin{subproof}
                    Show that $U\in\pow{X}$.
                    \begin{subproof}
                    Show that for all $x$ such that $x\in U$ we have $x\in X$.
                    \begin{subproof}
                        Fix $x$.
                        Assume $x\in U$.
                        Thus $x\in X$ by \cref{openray_right}.
                    \end{subproof}
                    Thus $U\subseteq X$ by \cref{subseteq}.
                    Thus $U\in\pow{X}$ by \cref{powerset_intro}.
                    \end{subproof}
                    Show that for all $x$ such that $x\in U$ there exists $B_1\in B$ such that $x\in B_1$ and $B_1\subseteq U$.
                    \begin{subproof}
                        Fix $x$.
                        Assume $x\in U$.
                        Thus $x\in X$ and $a\mathrel{R}x$ by \cref{openray_right}.
                        Thus $x$ is not a largest element of $X$ with respect to $R$.
                        Take $b\in X$ such that $x\mathrel{R}b$ by \cref{non_largest_has_successor}.
                        Let $B_1 = \intervalopenrel{R}{X}{a}{b}$.
                        Show that $B_1\in B$.
                        \begin{subproof}
                            Show that for all $z$ such that $z\in B_1$ we have $z\in X$.
                            \begin{subproof}
                                Fix $z$.
                                Assume $z\in B_1$.
                                Thus $z\in X$ by \cref{intervalopen_rel}.
                            \end{subproof}
                            Thus $B_1\subseteq X$ by \cref{subseteq}.
                            Thus $B_1\in\pow{X}$ by \cref{powerset_intro}.
                            $R$ is transitive by \cref{simple_order_relation}.
                            Thus $a\mathrel{R}b$ by \hyperref[transitive]{transitivity}.
                            Thus $B_1\in B$ by \cref{order_basis}.
                        \end{subproof}
                        Thus $x\in B_1$ by \cref{intervalopen_rel}.
                        Show that $B_1\subseteq U$.
                        \begin{subproof}
                            Show that for all $z$ such that $z\in B_1$ we have $z\in U$.
                            \begin{subproof}
                                Fix $z$.
                                Assume $z\in B_1$.
                                Thus $z\in X$ and $a\mathrel{R}z$ by \cref{intervalopen_rel}.
                                Thus $z\in U$ by \cref{openray_right}.
                            \end{subproof}
                            Thus $B_1\subseteq U$ by \cref{subseteq}.
                        \end{subproof}
                    \end{subproof}
                    Thus $U\in T$ by \cref{topology_generated_by_basis}.
                \end{subproof}
            \end{byCase}
        \caseOf{$U = \openrayleft{R}{X}{a}$.}
            \begin{byCase}
            \caseOf{$\exists a_0\in X. (\forall x\in X. a_0\mathrel{R}x \lor a_0 = x)$.}
                Take $a_0$ such that $a_0$ is a smallest element of $X$ with respect to $R$.
                Thus $U = \intervalclosedopenrel{R}{X}{a_0}{a}$ by \cref{intervalclosedopenrel_smallest_eq_openrayleft}.
                Show that $U\in B$.
                \begin{subproof}
                    Show that for all $x$ such that $x\in U$ we have $x\in X$.
                    \begin{subproof}
                        Fix $x$.
                        Assume $x\in U$.
                        Thus $x\in X$ by \cref{intervalclosedopen_rel}.
                    \end{subproof}
                    Thus $U\subseteq X$ by \cref{subseteq}.
                    Thus $U\in\pow{X}$ by \cref{powerset_intro}.
                    Thus $a\in X$.
                    Thus $a_0\in X$ and $(\forall x\in X. a_0\mathrel{R}x \lor a_0 = x)$ by \cref{smallest_element}.
                    Thus $\exists b\in X. \exists a_0\in X. (\forall x\in X. a_0\mathrel{R}x \lor a_0 = x) \land U = \intervalclosedopenrel{R}{X}{a_0}{b}$.
                    Thus $U\in B$ by \cref{order_basis}.
                \end{subproof}
                Thus $U\in T$ by \cref{basis_elem_open_in_generated,order_basis_is_basis}.
            \caseOf{it is wrong that there exists $a_0\in X$ such that for all $x\in X$ we have $a_0\mathrel{R}x$ or $a_0 = x$.}
                Show that $U\in T$.
                \begin{subproof}
                    Show that $U\in\pow{X}$.
                    \begin{subproof}
                        Show that for all $x$ such that $x\in U$ we have $x\in X$.
                        \begin{subproof}
                            Fix $x$.
                            Assume $x\in U$.
                            Thus $x\in X$ by \cref{openray_left}.
                        \end{subproof}
                        Thus $U\subseteq X$ by \cref{subseteq}.
                        Thus $U\in\pow{X}$ by \cref{powerset_intro}.
                    \end{subproof}
                    Show that for all $x$ such that $x\in U$ there exists $B_1\in B$ such that $x\in B_1$ and $B_1\subseteq U$.
                    \begin{subproof}
                        Fix $x$.
                        Assume $x\in U$.
                        Thus $x\in X$ and $x\mathrel{R}a$ by \cref{openray_left}.
                        Thus $x$ is not a smallest element of $X$ with respect to $R$.
                        Take $b\in X$ such that $b\mathrel{R}x$ by \cref{non_smallest_has_predecessor}.
                        Let $B_1 = \intervalopenrel{R}{X}{b}{a}$.
                        Show that $B_1\in B$.
                        \begin{subproof}
                            Show that for all $z$ such that $z\in B_1$ we have $z\in X$.
                            \begin{subproof}
                                Fix $z$.
                                Assume $z\in B_1$.
                                Thus $z\in X$ by \cref{intervalopen_rel}.
                            \end{subproof}
                            Thus $B_1\subseteq X$ by \cref{subseteq}.
                            Thus $B_1\in\pow{X}$ by \cref{powerset_intro}.
                            $R$ is transitive by \cref{simple_order_relation}.
                            Thus $b\mathrel{R}a$ by \hyperref[transitive]{transitivity}.
                            Thus $B_1\in B$ by \cref{order_basis}.
                        \end{subproof}
                        Thus $x\in B_1$ by \cref{intervalopen_rel}.
                        Show that $B_1\subseteq U$.
                        \begin{subproof}
                            Show that for all $z$ such that $z\in B_1$ we have $z\in U$.
                            \begin{subproof}
                                Fix $z$.
                                Assume $z\in B_1$.
                                Thus $z\in X$ and $z\mathrel{R}a$ by \cref{intervalopen_rel}.
                                Thus $z\in U$ by \cref{openray_left}.
                            \end{subproof}
                            Thus $B_1\subseteq U$ by \cref{subseteq}.
                        \end{subproof}
                    \end{subproof}
                    Thus $U\in T$ by \cref{topology_generated_by_basis}.
                \end{subproof}
            \end{byCase}
        \end{byCase}
    \end{subproof}
    Show that $T$ is finer than $T_1$ on $X$.
    \begin{subproof}
        Show that for all $x,B_2$ such that $x\in X$ and $B_2\in\finiteIntersections{S}$ and $x\in B_2$
            there exists $B_3\in B$ such that $x\in B_3$ and $B_3\subseteq B_2$.
        \begin{subproof}
            Fix $x$.
            Fix $B_2$.
            Assume $x\in X$ and $B_2\in\finiteIntersections{S}$ and $x\in B_2$.
            Take $F\subseteq S$ such that $F$ is finite and inhabited and $B_2 = \inters{F}$ by \cref{finite_intersections}.
            Show that for all $U$ such that $U\in F$ we have $U\in T$.
            \begin{subproof}
                Fix $U$.
                Assume $U\in F$.
                Then $U\in S$ by \cref{subseteq}.
                Thus $U\in T$ by \cref{subseteq}.
            \end{subproof}
            Thus $F\subseteq T$ by \cref{subseteq}.
            $T$ is a topology on $X$ by \cref{order_basis_is_basis,basis_topology_is_topology,order_topology}.
            Thus $B_2\in T$ by \cref{finite_intersections_open_in_topology}.
            Then there exists $B_3\in B$ such that $x\in B_3$ and $B_3\subseteq B_2$ by \cref{topology_generated_by_basis}.
        \end{subproof}
        Thus $T$ is finer than $T_1$ on $X$ by \cref{basis_finer_criterion}.
    \end{subproof}
    Thus $T_1\subseteq T$ by \cref{finer_topology}.
    Thus $T_1 = T$ by \cref{subseteq_antisymmetric}.
\end{proof}

% from §15 Item 1: product basis on X×Y
\begin{definition}\label{product_basis}
    $\productBasis{T}{T_1}{X}{Y} = \{B\in\pow{X\times Y} \mid \exists U\in T. \exists V\in T_1. B = U\times V\}$.
\end{definition}

% from §15 Item 2: product topology on X×Y
\begin{definition}\label{product_topology}
    $\productTopology{T}{T_1}{X}{Y} = \basisTopology{\productBasis{T}{T_1}{X}{Y}}{X\times Y}$.
\end{definition}

% from §15 Item 3: product basis is a basis
\begin{lemma}\label{product_basis_is_basis}
    Assume $T$ is a topology on $X$.
    Assume $T_1$ is a topology on $Y$.
    Then $\productBasis{T}{T_1}{X}{Y}$ is a basis on $X\times Y$.
\end{lemma}
\begin{proof}
    Let $B = \productBasis{T}{T_1}{X}{Y}$.
    Show that for all $B_1$ such that $B_1\in B$ we have $B_1\in\pow{X\times Y}$.
    \begin{subproof}
        Fix $B_1$.
        Assume $B_1\in B$.
        Thus $B_1\in\pow{X\times Y}$ by \cref{product_basis}.
    \end{subproof}
    Thus $B\subseteq\pow{X\times Y}$ by \cref{subseteq}.
    Show that for all $p$ such that $p\in X\times Y$ there exists $B_1\in B$ such that $p\in B_1$.
    \begin{subproof}
        Fix $p$.
        Assume $p\in X\times Y$.
        $X\in T$ by \cref{topology_on}.
        $Y\in T_1$ by \cref{topology_on}.
        $X$ is open in $X$ with respect to $T$ by \cref{open_in}.
        $Y$ is open in $Y$ with respect to $T_1$ by \cref{open_in}.
        Let $U = X$.
        Let $V = Y$.
        Let $B_1 = U\times V$.
        Thus $U\in T$ and $V\in T_1$.
        $U\subseteq X$ and $V\subseteq Y$ by \cref{subseteq_refl}.
        $B_1\subseteq X\times V$ by \cref{times_subseteq_left}.
        $X\times V\subseteq X\times Y$ by \cref{times_subseteq_right}.
        Thus $B_1\subseteq X\times Y$ by \cref{subseteq_transitive}.
        Thus $B_1\in\pow{X\times Y}$ by \cref{powerset_intro}.
        Thus $B_1\in B$ by \cref{product_basis}.
        Thus $p\in B_1$ by \cref{subseteq_refl}.
    \end{subproof}
    Show that for all $B_1,B_2,p$ such that $B_1\in B$ and $B_2\in B$ and $p\in B_1\inter B_2$
        there exists $B_3\in B$ such that $p\in B_3$ and $B_3\subseteq B_1\inter B_2$.
    \begin{subproof}
        Fix $B_1$.
        Fix $B_2$.
        Fix $p$.
        Assume $B_1\in B$ and $B_2\in B$ and $p\in B_1\inter B_2$.
        Take $U_1,V_1$ such that $U_1\in T$ and $V_1\in T_1$ and $B_1 = U_1\times V_1$ by \cref{product_basis}.
        Take $U_2,V_2$ such that $U_2\in T$ and $V_2\in T_1$ and $B_2 = U_2\times V_2$ by \cref{product_basis}.
        $U_1$ is open in $X$ with respect to $T$ by \cref{open_in}.
        $U_2$ is open in $X$ with respect to $T$ by \cref{open_in}.
        $V_1$ is open in $Y$ with respect to $T_1$ by \cref{open_in}.
        $V_2$ is open in $Y$ with respect to $T_1$ by \cref{open_in}.
        $U_1\inter U_2\in T$ by \cref{topology_on}.
        $V_1\inter V_2\in T_1$ by \cref{topology_on}.
        Thus $U_1\inter U_2$ is open in $X$ with respect to $T$ by \cref{open_in}.
        Thus $V_1\inter V_2$ is open in $Y$ with respect to $T_1$ by \cref{open_in}.
        Let $B_3 = (U_1\inter U_2)\times (V_1\inter V_2)$.
        $T\subseteq\pow{X}$ by \cref{topology_on}.
        $T_1\subseteq\pow{Y}$ by \cref{topology_on}.
        Thus $U_1\subseteq X$ and $U_2\subseteq X$ by \cref{powerset_elim,subseteq}.
        Thus $V_1\subseteq Y$ and $V_2\subseteq Y$ by \cref{powerset_elim,subseteq}.
        Thus $U_1\inter U_2\subseteq X$ by \cref{inter_subseteq}.
        Thus $V_1\inter V_2\subseteq Y$ by \cref{inter_subseteq}.
        $B_3\subseteq X\times (V_1\inter V_2)$ by \cref{times_subseteq_left}.
        $X\times (V_1\inter V_2)\subseteq X\times Y$ by \cref{times_subseteq_right}.
        Thus $B_3\subseteq X\times Y$ by \cref{subseteq_transitive}.
        Thus $B_3\in\pow{X\times Y}$ by \cref{powerset_intro}.
        Thus $B_3\in B$ by \cref{product_basis}.
        Thus $p\in B_3$ by \cref{inter_times_elim}.
        $B_3\subseteq B_1\inter B_2$ by \cref{subseteq,inter_times_intro}.
    \end{subproof}
    Thus $B$ is a basis on $X\times Y$ by \cref{basis_on}.
\end{proof}

% from §15 Item 4: bases for X and Y give a basis for X×Y
\begin{lemma}\label{product_basis_from_bases}
    Assume $B$ is a basis for $T$ on $X$.
    Assume $C$ is a basis for $T_1$ on $Y$.
    Let $D = \{D_1\in\pow{X\times Y} \mid \exists B_1\in B. \exists C_1\in C. D_1 = B_1\times C_1\}$.
    Then $D$ is a basis for $\productTopology{T}{T_1}{X}{Y}$ on $X\times Y$.
\end{lemma}
\begin{proof}
    Let $T_2 = \productTopology{T}{T_1}{X}{Y}$.
    $B$ is a basis on $X$ by \cref{basis_for_topology}.
    $C$ is a basis on $Y$ by \cref{basis_for_topology}.
    $T = \basisTopology{B}{X}$ by \cref{basis_for_topology}.
    $T_1 = \basisTopology{C}{Y}$ by \cref{basis_for_topology}.
    Thus $T$ is a topology on $X$ by \cref{basis_topology_is_topology}.
    Thus $T_1$ is a topology on $Y$ by \cref{basis_topology_is_topology}.
    $\productBasis{T}{T_1}{X}{Y}$ is a basis on $X\times Y$ by \cref{product_basis_is_basis}.
    Thus $T_2$ is a topology on $X\times Y$ by \cref{basis_topology_is_topology,product_topology}.
    Show that for all $D_1$ such that $D_1\in D$ we have $D_1$ is open in $X\times Y$ with respect to $T_2$.
    \begin{subproof}
        Fix $D_1$.
        Assume $D_1\in D$.
        Take $B_1,C_1$ such that $B_1\in B$ and $C_1\in C$ and $D_1 = B_1\times C_1$.
        $B_1\in\basisTopology{B}{X}$ by \cref{basis_elem_open_in_generated}.
        $C_1\in\basisTopology{C}{Y}$ by \cref{basis_elem_open_in_generated}.
        Thus $B_1\in T$ and $C_1\in T_1$ by \cref{basis_for_topology}.
        Thus $B_1$ is open in $X$ with respect to $T$ by \cref{open_in}.
        Thus $C_1$ is open in $Y$ with respect to $T_1$ by \cref{open_in}.
        Thus $D_1\in\productBasis{T}{T_1}{X}{Y}$ by \cref{product_basis}.
        Thus $D_1\in T_2$ by \cref{basis_elem_open_in_generated,product_topology}.
        Thus $D_1$ is open in $X\times Y$ with respect to $T_2$ by \cref{open_in}.
    \end{subproof}
    Show that for all $U,p$ such that $U\in T_2$ and $p\in U$ there exists $D_1\in D$ such that $p\in D_1$ and $D_1\subseteq U$.
    \begin{subproof}
        Fix $U$.
        Fix $p$.
        Assume $U\in T_2$ and $p\in U$.
        $U\in\basisTopology{\productBasis{T}{T_1}{X}{Y}}{X\times Y}$ by \cref{product_topology}.
        Take $B_0\in\productBasis{T}{T_1}{X}{Y}$ such that $p\in B_0$ and $B_0\subseteq U$ by \cref{topology_generated_by_basis}.
        Take $U_1,V_1$ such that $U_1\in T$ and $V_1\in T_1$ and $B_0 = U_1\times V_1$ by \cref{product_basis}.
        Thus $U_1$ is open in $X$ with respect to $T$ by \cref{open_in}.
        Thus $V_1$ is open in $Y$ with respect to $T_1$ by \cref{open_in}.
        Take $x,y$ such that $p = (x,y)$ and $x\in U_1$ and $y\in V_1$ by \cref{times_elem_is_tuple}.
        Take $B_1\in B$ such that $x\in B_1$ and $B_1\subseteq U_1$ by \cref{basis_for_topology,topology_generated_by_basis,open_in}.
        Take $C_1\in C$ such that $y\in C_1$ and $C_1\subseteq V_1$ by \cref{basis_for_topology,topology_generated_by_basis,open_in}.
        Let $D_1 = B_1\times C_1$.
        $B\subseteq\pow{X}$ by \cref{basis_on}.
        $C\subseteq\pow{Y}$ by \cref{basis_on}.
        Thus $B_1\in\pow{X}$ and $C_1\in\pow{Y}$ by \cref{subseteq}.
        Thus $B_1\subseteq X$ and $C_1\subseteq Y$ by \cref{powerset_elim}.
        $D_1\subseteq X\times C_1$ by \cref{times_subseteq_left}.
        $X\times C_1\subseteq X\times Y$ by \cref{times_subseteq_right}.
        Thus $D_1\subseteq X\times Y$ by \cref{subseteq_transitive}.
        Thus $D_1\in\pow{X\times Y}$ by \cref{powerset_intro}.
        Thus $D_1\in D$.
        $p\in D_1$ by \cref{times,inter_intro}.
        $D_1\subseteq U_1\times C_1$ by \cref{times_subseteq_left}.
        $U_1\times C_1\subseteq U_1\times V_1$ by \cref{times_subseteq_right}.
        Thus $D_1\subseteq B_0$ by \cref{subseteq_transitive}.
        Thus $D_1\subseteq U$ by \cref{subseteq_transitive}.
    \end{subproof}
    Thus $D$ is a basis for $T_2$ on $X\times Y$ by \cref{open_collection_basis}.
\end{proof}

% from §15 Item 5: projection onto the first factor
\begin{definition}\label{projection_first}
    $\piOne{X}{Y} = \{((x,y),x) \mid x\in X, y\in Y\}$.
\end{definition}

% from §15 Item 6: projection onto the second factor
\begin{definition}\label{projection_second}
    $\piTwo{X}{Y} = \{((x,y),y) \mid x\in X, y\in Y\}$.
\end{definition}

% from §15 helper lemma 1: preimage of the first projection
\begin{lemma}\label{preimg_pi1}
    Suppose $U\subseteq X$.
    Then $\preimg{\piOne{X}{Y}}{U} = U\times Y$.
\end{lemma}
\begin{proof}
    Show that for all $p$ such that $p\in\preimg{\piOne{X}{Y}}{U}$ we have $p\in U\times Y$.
    \begin{subproof}
        Fix $p$.
        Assume $p\in\preimg{\piOne{X}{Y}}{U}$.
        Take $x\in U$ such that $p\mathrel{\piOne{X}{Y}} x$ by \cref{preimg_iff}.
        Thus $(p,x)\in\piOne{X}{Y}$.
        Take $a,b$ such that $a\in X$ and $b\in Y$ and $(p,x) = ((a,b),a)$ by \cref{projection_first}.
        Thus $p = (a,b)$ and $x = a$ by \cref{pair_eq_iff}.
        Thus $a\in U$.
        Thus $p\in U\times Y$ by \cref{times}.
    \end{subproof}
    Thus $\preimg{\piOne{X}{Y}}{U}\subseteq U\times Y$ by \cref{subseteq}.
    Show that for all $p$ such that $p\in U\times Y$ we have $p\in\preimg{\piOne{X}{Y}}{U}$.
    \begin{subproof}
        Fix $p$.
        Assume $p\in U\times Y$.
        Take $x,y$ such that $p = (x,y)$ and $x\in U$ and $y\in Y$ by \cref{times_elem_is_tuple}.
        $x\in X$ by \cref{subseteq}.
        Thus $p\mathrel{\piOne{X}{Y}} x$ by \cref{projection_first}.
        Thus $p\in\preimg{\piOne{X}{Y}}{U}$ by \cref{preimg_iff}.
    \end{subproof}
    Thus $U\times Y\subseteq\preimg{\piOne{X}{Y}}{U}$ by \cref{subseteq}.
    Thus $\preimg{\piOne{X}{Y}}{U} = U\times Y$ by \cref{subseteq_antisymmetric}.
\end{proof}

% from §15 helper lemma 2: preimage of the second projection
\begin{lemma}\label{preimg_pi2}
    Suppose $V\subseteq Y$.
    Then $\preimg{\piTwo{X}{Y}}{V} = X\times V$.
\end{lemma}
\begin{proof}
    Show that for all $p$ such that $p\in\preimg{\piTwo{X}{Y}}{V}$ we have $p\in X\times V$.
    \begin{subproof}
        Fix $p$.
        Assume $p\in\preimg{\piTwo{X}{Y}}{V}$.
        Take $y\in V$ such that $p\mathrel{\piTwo{X}{Y}} y$ by \cref{preimg_iff}.
        Thus $(p,y)\in\piTwo{X}{Y}$.
        Take $a,b$ such that $a\in X$ and $b\in Y$ and $(p,y) = ((a,b),b)$ by \cref{projection_second}.
        Thus $p = (a,b)$ and $y = b$ by \cref{pair_eq_iff}.
        Thus $b\in V$.
        Thus $p\in X\times V$ by \cref{times}.
    \end{subproof}
    Thus $\preimg{\piTwo{X}{Y}}{V}\subseteq X\times V$ by \cref{subseteq}.
    Show that for all $p$ such that $p\in X\times V$ we have $p\in\preimg{\piTwo{X}{Y}}{V}$.
    \begin{subproof}
        Fix $p$.
        Assume $p\in X\times V$.
        Take $x,y$ such that $p = (x,y)$ and $x\in X$ and $y\in V$ by \cref{times_elem_is_tuple}.
        $y\in Y$ by \cref{subseteq}.
        Thus $p\mathrel{\piTwo{X}{Y}} y$ by \cref{projection_second}.
        Thus $p\in\preimg{\piTwo{X}{Y}}{V}$ by \cref{preimg_iff}.
    \end{subproof}
    Thus $X\times V\subseteq\preimg{\piTwo{X}{Y}}{V}$ by \cref{subseteq}.
    Thus $\preimg{\piTwo{X}{Y}}{V} = X\times V$ by \cref{subseteq_antisymmetric}.
\end{proof}

% from §15 Item 7: left projection subbasis collection
\begin{definition}\label{product_subbasis_left}
    $\productSubbasisLeft{T}{X}{Y} = \{A \mid \exists U\in T. A = \preimg{\piOne{X}{Y}}{U}\}$.
\end{definition}

% from §15 Item 8: right projection subbasis collection
\begin{definition}\label{product_subbasis_right}
    $\productSubbasisRight{T_1}{X}{Y} = \{A \mid \exists V\in T_1. A = \preimg{\piTwo{X}{Y}}{V}\}$.
\end{definition}

% from §15 Item 9: projection subbasis for the product topology
\begin{lemma}\label{product_topology_subbasis}
    Assume $T$ is a topology on $X$.
    Assume $T_1$ is a topology on $Y$.
    Let $S = \productSubbasisLeft{T}{X}{Y}\union \productSubbasisRight{T_1}{X}{Y}$.
    Then $\subbasisTopology{S}{X\times Y} = \productTopology{T}{T_1}{X}{Y}$.
\end{lemma}
\begin{proof}
    Let $T_2 = \productTopology{T}{T_1}{X}{Y}$.
    Let $T_3 = \subbasisTopology{S}{X\times Y}$.
    Let $S_1 = \productSubbasisLeft{T}{X}{Y}$.
    Let $S_2 = \productSubbasisRight{T_1}{X}{Y}$.
    $\productBasis{T}{T_1}{X}{Y}$ is a basis on $X\times Y$ by \cref{product_basis_is_basis}.
    Thus $T_2$ is a topology on $X\times Y$ by \cref{basis_topology_is_topology,product_topology}.
    Show that $S$ is a subbasis on $X\times Y$.
    \begin{subproof}
        Show that for all $A$ such that $A\in S$ we have $A\in\pow{X\times Y}$.
        \begin{subproof}
            Fix $A$.
            Assume $A\in S$.
            Thus $A\in S_1$ or $A\in S_2$ by \cref{union_iff}.
            \begin{byCase}
            \caseOf{$A\in S_1$.}
                Take $U\in T$ such that $A = \preimg{\piOne{X}{Y}}{U}$.
                $T\subseteq\pow{X}$ by \cref{topology_on}.
                Thus $U\in\pow{X}$ by \cref{subseteq}.
                Thus $U\subseteq X$ by \cref{powerset_elim}.
                Thus $A = U\times Y$ by \cref{preimg_pi1}.
                $A\subseteq X\times Y$ by \cref{times_subseteq_left}.
                Thus $A\in\pow{X\times Y}$ by \cref{powerset_intro}.
            \caseOf{$A\in S_2$.}
                Take $V\in T_1$ such that $A = \preimg{\piTwo{X}{Y}}{V}$.
                $T_1\subseteq\pow{Y}$ by \cref{topology_on}.
                Thus $V\in\pow{Y}$ by \cref{subseteq}.
                Thus $V\subseteq Y$ by \cref{powerset_elim}.
                Thus $A = X\times V$ by \cref{preimg_pi2}.
                $A\subseteq X\times Y$ by \cref{times_subseteq_right}.
                Thus $A\in\pow{X\times Y}$ by \cref{powerset_intro}.
            \end{byCase}
        \end{subproof}
        Thus $S\subseteq\pow{X\times Y}$ by \cref{subseteq}.
        Show that $\unions{S} = X\times Y$.
        \begin{subproof}
            $X\in T$ by \cref{topology_on}.
            Thus $X$ is open in $X$ with respect to $T$ by \cref{open_in}.
            Let $A = \preimg{\piOne{X}{Y}}{X}$.
            Thus $A\in S_1$.
            Thus $A\in S$ by \cref{union_intro_left}.
            $A = X\times Y$ by \cref{preimg_pi1,subseteq_refl}.
            Show that for all $z$ such that $z\in X\times Y$ we have $z\in\unions{S}$.
            \begin{subproof}
                Fix $z$.
                Assume $z\in X\times Y$.
                Thus $z\in A$.
                Thus $z\in\unions{S}$ by \cref{unions_iff}.
            \end{subproof}
            Thus $X\times Y\subseteq\unions{S}$ by \cref{subseteq}.
            $\unions{S}\subseteq X\times Y$ by \cref{unions_subseteq_of_powerset_is_subseteq}.
            Thus $\unions{S} = X\times Y$ by \cref{subseteq_antisymmetric}.
        \end{subproof}
        Thus $S$ is a subbasis on $X\times Y$ by \cref{subbasis_on}.
    \end{subproof}
    $\finiteIntersections{S}$ is a basis on $X\times Y$ by \cref{subbasis_finite_intersections_basis}.
    Thus $\finiteIntersections{S}$ is a basis for $T_3$ on $X\times Y$ by \cref{basis_for_topology,topology_generated_by_subbasis}.
    Show that $T_3\subseteq T_2$.
    \begin{subproof}
        Show that for all $B$ such that $B\in\finiteIntersections{S}$ we have $B\in T_2$.
        \begin{subproof}
            Fix $B$.
            Assume $B\in\finiteIntersections{S}$.
            Take $F\subseteq S$ such that $F$ is finite and inhabited and $B = \inters{F}$ by \cref{finite_intersections}.
            Show that for all $U$ such that $U\in F$ we have $U\in T_2$.
            \begin{subproof}
                Fix $U$.
                Assume $U\in F$.
                Then $U\in S$ by \cref{subseteq}.
                Thus $U\in S_1$ or $U\in S_2$ by \cref{union_iff}.
                \begin{byCase}
                \caseOf{$U\in S_1$.}
                    Take $U_0\in T$ such that $U = \preimg{\piOne{X}{Y}}{U_0}$.
                    $T\subseteq\pow{X}$ by \cref{topology_on}.
                    Thus $U_0\in\pow{X}$ by \cref{subseteq}.
                    Thus $U_0\subseteq X$ by \cref{powerset_elim}.
                    $Y\in T_1$ by \cref{topology_on}.
                    Thus $Y$ is open in $Y$ with respect to $T_1$ by \cref{open_in}.
                    Thus $U = U_0\times Y$ by \cref{preimg_pi1}.
                    $U_0\subseteq X$ by \cref{powerset_elim}.
                    $Y\subseteq Y$ by \cref{subseteq_refl}.
                    $U\subseteq X\times Y$ by \cref{times_subseteq_left,times_subseteq_right,subseteq_transitive}.
                    Thus $U\in\pow{X\times Y}$ by \cref{powerset_intro}.
                    Thus $U\in\productBasis{T}{T_1}{X}{Y}$ by \cref{product_basis}.
                    Thus $U\in T_2$ by \cref{basis_elem_open_in_generated,product_topology}.
                \caseOf{$U\in S_2$.}
                    Take $V_0\in T_1$ such that $U = \preimg{\piTwo{X}{Y}}{V_0}$.
                    $T_1\subseteq\pow{Y}$ by \cref{topology_on}.
                    Thus $V_0\in\pow{Y}$ by \cref{subseteq}.
                    Thus $V_0\subseteq Y$ by \cref{powerset_elim}.
                    $X\in T$ by \cref{topology_on}.
                    Thus $X$ is open in $X$ with respect to $T$ by \cref{open_in}.
                    Thus $U = X\times V_0$ by \cref{preimg_pi2}.
                    $X\subseteq X$ by \cref{subseteq_refl}.
                    $U\subseteq X\times Y$ by \cref{times_subseteq_left,times_subseteq_right,subseteq_transitive}.
                    Thus $U\in\pow{X\times Y}$ by \cref{powerset_intro}.
                    Thus $U\in\productBasis{T}{T_1}{X}{Y}$ by \cref{product_basis}.
                    Thus $U\in T_2$ by \cref{basis_elem_open_in_generated,product_topology}.
                \end{byCase}
            \end{subproof}
            Thus $F\subseteq T_2$ by \cref{subseteq}.
            Thus $B\in T_2$ by \cref{finite_intersections_open_in_topology}.
        \end{subproof}
        Thus $\finiteIntersections{S}\subseteq T_2$ by \cref{subseteq}.
        Show that for all $U$ such that $U\in T_3$ we have $U\in T_2$.
        \begin{subproof}
            Fix $U$.
            Assume $U\in T_3$.
            Thus there exists $M\subseteq\finiteIntersections{S}$ such that $U = \unions{M}$ by \cref{basis_topology_unions,topology_generated_by_subbasis}.
            Thus $M\subseteq T_2$ by \cref{subseteq}.
            Thus $\unions{M}\in T_2$ by \cref{topology_on}.
            Thus $U\in T_2$.
        \end{subproof}
    \end{subproof}
    Show that $T_2\subseteq T_3$.
    \begin{subproof}
        Show that for all $B$ such that $B\in\productBasis{T}{T_1}{X}{Y}$ we have $B\in T_3$.
        \begin{subproof}
            Fix $B$.
            Assume $B\in\productBasis{T}{T_1}{X}{Y}$.
            Take $U,V$ such that $U\in T$ and $V\in T_1$ and $B = U\times V$ by \cref{product_basis}.
            $T\subseteq\pow{X}$ by \cref{topology_on}.
            $T_1\subseteq\pow{Y}$ by \cref{topology_on}.
            Thus $U\in\pow{X}$ and $V\in\pow{Y}$ by \cref{subseteq}.
            Thus $U\subseteq X$ and $V\subseteq Y$ by \cref{powerset_elim}.
            Let $A_1 = \preimg{\piOne{X}{Y}}{U}$.
            Let $A_2 = \preimg{\piTwo{X}{Y}}{V}$.
            $A_1\in S_1$ and $A_2\in S_2$.
            Thus $A_1\in S$ and $A_2\in S$ by \cref{union_intro_left,union_intro_right}.
            Let $F = \{A_1, A_2\}$.
            $F\subseteq S$ by \cref{subseteq,upair_elim}.
            $F$ is finite by \cref{pair_is_finite}.
            $F$ is inhabited by \cref{upair_intro_left}.
            $A_1 = U\times Y$ by \cref{preimg_pi1}.
            $A_2 = X\times V$ by \cref{preimg_pi2}.
            $U\inter X = U$ by \cref{inter_absorb_supseteq_right}.
            $Y\inter V = V$ by \cref{inter_absorb_supseteq_left}.
            Thus $B = (U\inter X)\times (Y\inter V)$.
            Thus $B = A_1\inter A_2$ by \cref{inter_times}.
            Thus $B = \inters{F}$ by \cref{inter_as_inters}.
            $B\subseteq X\times Y$ by \cref{times_subseteq_left,times_subseteq_right,subseteq_transitive}.
            $\unions{S} = X\times Y$ by \cref{subbasis_on}.
            Thus $B\subseteq\unions{S}$ by \cref{subseteq_transitive}.
            Thus $B\in\pow{\unions{S}}$ by \cref{powerset_intro}.
            Thus $B\in\finiteIntersections{S}$ by \cref{finite_intersections}.
            Thus $B\in T_3$ by \cref{basis_elem_open_in_generated,topology_generated_by_subbasis}.
        \end{subproof}
        Show that for all $U$ such that $U\in T_2$ we have $U\in T_3$.
        \begin{subproof}
            Fix $U$.
            Assume $U\in T_2$.
            $\productBasis{T}{T_1}{X}{Y}$ is a basis on $X\times Y$ by \cref{product_basis_is_basis}.
            Thus $\productBasis{T}{T_1}{X}{Y}$ is a basis for $T_2$ on $X\times Y$ by \cref{basis_for_topology,product_topology}.
            Thus there exists $M\subseteq\productBasis{T}{T_1}{X}{Y}$ such that $U = \unions{M}$ by \cref{basis_topology_unions}.
            Thus $M\subseteq T_3$ by \cref{subseteq}.
            $T_3$ is a topology on $X\times Y$ by \cref{basis_topology_is_topology,basis_for_topology,topology_generated_by_subbasis}.
            Thus $\unions{M}\in T_3$ by \cref{topology_on}.
            Thus $U\in T_3$.
        \end{subproof}
    \end{subproof}
    Thus $T_2 = T_3$ by \cref{subseteq_antisymmetric}.
\end{proof}

% from §16 Item 1: subspace topology
\begin{definition}\label{subspace_topology}
    $\subspaceTopology{T}{Y} = \{A \mid \exists U\in T. A = Y\inter U\}$.
\end{definition}

% from §16 Item 2: subspace of a space
\begin{definition}\label{subspace_of}
    $Y$ is a subspace of $X$ with respect to $T$ iff $T$ is a topology on $X$ and $Y\subseteq X$.
\end{definition}

% from §16 helper lemma 1: subspace topology is a topology
\begin{lemma}\label{subspace_topology_is_topology}
    Assume $T$ is a topology on $X$.
    Assume $Y\subseteq X$.
    Then $\subspaceTopology{T}{Y}$ is a topology on $Y$.
\end{lemma}
\begin{proof}
    Let $T_1 = \subspaceTopology{T}{Y}$.
    Show that $T_1\subseteq\pow{Y}$.
    \begin{subproof}
        Show that for all $A$ such that $A\in T_1$ we have $A\in\pow{Y}$.
        \begin{subproof}
            Fix $A$.
            Assume $A\in T_1$.
            Take $U\in T$ such that $A = Y\inter U$ by \cref{subspace_topology}.
            Thus $A\subseteq Y$ by \cref{inter_lower_left}.
            Thus $A\in\pow{Y}$ by \cref{powerset_intro}.
        \end{subproof}
        Thus $T_1\subseteq\pow{Y}$ by \cref{subseteq}.
    \end{subproof}
    $\emptyset\in T$ by \cref{topology_on}.
    $X\in T$ by \cref{topology_on}.
    $\emptyset = Y\inter\emptyset$ by \cref{inter_emptyset}.
    Thus $\emptyset\in T_1$ by \cref{subspace_topology}.
    $Y\inter X = Y$ by \cref{inter_absorb_supseteq_right}.
    Thus $Y\in T_1$ by \cref{subspace_topology}.
    Show that for all $M$ such that $M\subseteq T_1$ we have $\unions{M}\in T_1$.
    \begin{subproof}
        Fix $M$.
        Assume $M\subseteq T_1$.
        Let $N = \{U\in T \mid \exists A\in M. A = Y\inter U\}$.
        Show that $N\subseteq T$.
        \begin{subproof}
            Show that for all $U$ such that $U\in N$ we have $U\in T$.
            \begin{subproof}
                Fix $U$.
                Assume $U\in N$.
                Thus $U\in T$ by \cref{subseteq}.
            \end{subproof}
            Thus $N\subseteq T$ by \cref{subseteq}.
        \end{subproof}
        Show that $\unions{M} = Y\inter\unions{N}$.
        \begin{subproof}
            Show that for all $x$ such that $x\in\unions{M}$ we have $x\in Y\inter\unions{N}$.
            \begin{subproof}
                Fix $x$.
                Assume $x\in\unions{M}$.
                Take $A\in M$ such that $x\in A$ by \cref{unions_iff}.
                $A\in T_1$ by \cref{subseteq}.
                Take $U\in T$ such that $A = Y\inter U$ by \cref{subspace_topology}.
                Thus $x\in Y$ and $x\in U$ by \cref{inter}.
                Thus $U\in N$.
                Thus $x\in\unions{N}$ by \cref{unions_iff}.
                Thus $x\in Y\inter\unions{N}$ by \cref{inter_intro}.
            \end{subproof}
            Thus $\unions{M}\subseteq Y\inter\unions{N}$ by \cref{subseteq}.
            Show that for all $x$ such that $x\in Y\inter\unions{N}$ we have $x\in\unions{M}$.
            \begin{subproof}
                Fix $x$.
                Assume $x\in Y\inter\unions{N}$.
                Thus $x\in Y$ and $x\in\unions{N}$ by \cref{inter}.
                Take $U\in N$ such that $x\in U$ by \cref{unions_iff}.
                Take $A\in M$ such that $A = Y\inter U$.
                Thus $x\in A$ by \cref{inter_intro}.
                Thus $x\in\unions{M}$ by \cref{unions_iff}.
            \end{subproof}
            Thus $Y\inter\unions{N}\subseteq\unions{M}$ by \cref{subseteq}.
            Thus $\unions{M} = Y\inter\unions{N}$ by \cref{subseteq_antisymmetric}.
        \end{subproof}
        $\unions{N}\in T$ by \cref{topology_on}.
        Thus $\unions{M}\in T_1$ by \cref{subspace_topology}.
    \end{subproof}
    Show that for all $A,B$ such that $A\in T_1$ and $B\in T_1$ we have $A\inter B\in T_1$.
    \begin{subproof}
        Fix $A$.
        Fix $B$.
        Assume $A\in T_1$ and $B\in T_1$.
        Take $U\in T$ such that $A = Y\inter U$ by \cref{subspace_topology}.
        Take $V\in T$ such that $B = Y\inter V$ by \cref{subspace_topology}.
        $U\inter V\in T$ by \cref{topology_on}.
        Show that $A\inter B = Y\inter (U\inter V)$.
        \begin{subproof}
            Show that for all $x$ such that $x\in A\inter B$ we have $x\in Y\inter (U\inter V)$.
            \begin{subproof}
                Fix $x$.
                Assume $x\in A\inter B$.
                Thus $x\in A$ and $x\in B$ by \cref{inter}.
                Thus $x\in Y$ and $x\in U$ by \cref{inter}.
                Thus $x\in Y$ and $x\in V$ by \cref{inter}.
                Thus $x\in U\inter V$ by \cref{inter_intro}.
                Thus $x\in Y\inter (U\inter V)$ by \cref{inter_intro}.
            \end{subproof}
            Thus $A\inter B\subseteq Y\inter (U\inter V)$ by \cref{subseteq}.
            Show that for all $x$ such that $x\in Y\inter (U\inter V)$ we have $x\in A\inter B$.
            \begin{subproof}
                Fix $x$.
                Assume $x\in Y\inter (U\inter V)$.
                Thus $x\in Y$ and $x\in U\inter V$ by \cref{inter}.
                Thus $x\in U$ and $x\in V$ by \cref{inter}.
                Thus $x\in A$ and $x\in B$ by \cref{inter_intro}.
                Thus $x\in A\inter B$ by \cref{inter_intro}.
            \end{subproof}
            Thus $Y\inter (U\inter V)\subseteq A\inter B$ by \cref{subseteq}.
            Thus $A\inter B = Y\inter (U\inter V)$ by \cref{subseteq_antisymmetric}.
        \end{subproof}
        Thus $A\inter B\in T_1$ by \cref{subspace_topology}.
    \end{subproof}
    Thus $T_1$ is closed under arbitrary unions.
    Thus $T_1$ is closed under binary intersections.
    Thus $T_1$ is a topology on $Y$ by \cref{topology_on}.
\end{proof}

% from §16 Item 3: basis for the subspace topology
\begin{lemma}\label{subspace_basis}
    Assume $B$ is a basis for $T$ on $X$.
    Assume $Y\subseteq X$.
    Let $C = \{A\in\pow{Y} \mid \exists B_1\in B. A = B_1\inter Y\}$.
    Then $C$ is a basis for $\subspaceTopology{T}{Y}$ on $Y$.
\end{lemma}
\begin{proof}
    Let $T_1 = \subspaceTopology{T}{Y}$.
    $B$ is a basis on $X$ by \cref{basis_for_topology}.
    $T = \basisTopology{B}{X}$ by \cref{basis_for_topology}.
    Thus $T$ is a topology on $X$ by \cref{basis_topology_is_topology}.
    Thus $T_1$ is a topology on $Y$ by \cref{subspace_topology_is_topology}.
    Show that for all $C_1$ such that $C_1\in C$ we have $C_1$ is open in $Y$ with respect to $T_1$.
    \begin{subproof}
        Fix $C_1$.
        Assume $C_1\in C$.
        Take $B_1\in B$ such that $C_1 = B_1\inter Y$.
        $B_1\in\basisTopology{B}{X}$ by \cref{basis_elem_open_in_generated}.
        Thus $B_1\in T$ by \cref{basis_for_topology}.
        Thus $C_1 = Y\inter B_1$ by \cref{inter_comm}.
        Thus $C_1\in T_1$ by \cref{subspace_topology}.
        Thus $C_1$ is open in $Y$ with respect to $T_1$ by \cref{open_in}.
    \end{subproof}
    Show that for all $U,x$ such that $U\in T_1$ and $x\in U$
        there exists $C_1\in C$ such that $x\in C_1$ and $C_1\subseteq U$.
    \begin{subproof}
        Fix $U$.
        Fix $x$.
        Assume $U\in T_1$ and $x\in U$.
        Take $U_1\in T$ such that $U = Y\inter U_1$ by \cref{subspace_topology}.
        Thus $x\in Y$ and $x\in U_1$ by \cref{inter}.
        Take $B_1\in B$ such that $x\in B_1$ and $B_1\subseteq U_1$
            by \cref{basis_for_topology,topology_generated_by_basis}.
        Let $C_1 = B_1\inter Y$.
        Thus $C_1\in C$.
        Thus $x\in C_1$ by \cref{inter_intro}.
        $C_1\subseteq B_1$ by \cref{inter_lower_left}.
        Thus $C_1\subseteq U_1$ by \cref{subseteq_transitive}.
        $C_1\subseteq Y$ by \cref{inter_lower_right}.
        Thus $C_1\subseteq U$ by \cref{subseteq_inter_iff}.
    \end{subproof}
    Thus $C$ is a basis for $T_1$ on $Y$ by \cref{open_collection_basis}.
\end{proof}

% from §16 Item 4: open in a subspace implies open in the space if the subspace is open
\begin{lemma}\label{open_in_subspace_open_in_space}
    Assume $Y$ is a subspace of $X$ with respect to $T$.
    Assume $U$ is open in $Y$ with respect to $\subspaceTopology{T}{Y}$.
    Assume $Y$ is open in $X$ with respect to $T$.
    Then $U$ is open in $X$ with respect to $T$.
\end{lemma}
\begin{proof}
    Let $T_1 = \subspaceTopology{T}{Y}$.
    $T$ is a topology on $X$ by \cref{subspace_of}.
    $Y\subseteq X$ by \cref{subspace_of}.
    Thus $T_1$ is a topology on $Y$ by \cref{subspace_topology_is_topology}.
    $U\in T_1$ by \cref{open_in}.
    Take $V\in T$ such that $U = Y\inter V$ by \cref{subspace_topology}.
    $Y\in T$ by \cref{open_in}.
    Thus $U\in T$ by \cref{topology_on}.
    Thus $U$ is open in $X$ with respect to $T$ by \cref{open_in}.
\end{proof}

% from §16 Item 5: product topology and subspace topology
\begin{theorem}\label{subspace_product_topology}
    Assume $A$ is a subspace of $X$ with respect to $T$.
    Assume $B$ is a subspace of $Y$ with respect to $T_1$.
    Then $\productTopology{\subspaceTopology{T}{A}}{\subspaceTopology{T_1}{B}}{A}{B}
        = \subspaceTopology{\productTopology{T}{T_1}{X}{Y}}{A\times B}$.
\end{theorem}
\begin{proof}
    Let $T_A = \subspaceTopology{T}{A}$.
    Let $T_B = \subspaceTopology{T_1}{B}$.
    Let $T_2 = \productTopology{T}{T_1}{X}{Y}$.
    Let $T_3 = \subspaceTopology{T_2}{A\times B}$.
    Let $C = \productBasis{T_A}{T_B}{A}{B}$.
    $T$ is a topology on $X$ by \cref{subspace_of}.
    $T_1$ is a topology on $Y$ by \cref{subspace_of}.
    $A\subseteq X$ by \cref{subspace_of}.
    $B\subseteq Y$ by \cref{subspace_of}.
    Thus $T_A$ is a topology on $A$ by \cref{subspace_topology_is_topology}.
    Thus $T_B$ is a topology on $B$ by \cref{subspace_topology_is_topology}.
    $\productBasis{T}{T_1}{X}{Y}$ is a basis on $X\times Y$ by \cref{product_basis_is_basis}.
    Thus $T_2$ is a topology on $X\times Y$ by \cref{basis_topology_is_topology,product_topology}.
    $A\times B\subseteq X\times B$ by \cref{times_subseteq_left}.
    $X\times B\subseteq X\times Y$ by \cref{times_subseteq_right}.
    Thus $A\times B\subseteq X\times Y$ by \cref{subseteq_transitive}.
    Thus $T_3$ is a topology on $A\times B$ by \cref{subspace_topology_is_topology}.
    Show that for all $C_1$ such that $C_1\in C$ we have $C_1$ is open in $A\times B$ with respect to $T_3$.
    \begin{subproof}
        Fix $C_1$.
        Assume $C_1\in C$.
        Take $U_A, V_B$ such that $U_A\in T_A$ and $V_B\in T_B$ and $C_1 = U_A\times V_B$ by \cref{product_basis}.
        Take $U\in T$ such that $U_A = A\inter U$ by \cref{subspace_topology}.
        Take $V\in T_1$ such that $V_B = B\inter V$ by \cref{subspace_topology}.
        Let $W = U\times V$.
        $U$ is open in $X$ with respect to $T$ by \cref{open_in}.
        $V$ is open in $Y$ with respect to $T_1$ by \cref{open_in}.
        Thus $U\in T$ by \cref{open_in}.
        Thus $V\in T_1$ by \cref{open_in}.
        $T\subseteq\pow{X}$ by \cref{topology_on}.
        $T_1\subseteq\pow{Y}$ by \cref{topology_on}.
        Thus $U\in\pow{X}$ by \cref{subseteq}.
        Thus $V\in\pow{Y}$ by \cref{subseteq}.
        Thus $U\subseteq X$ and $V\subseteq Y$ by \cref{powerset_elim}.
        $W\subseteq X\times V$ by \cref{times_subseteq_left}.
        $X\times V\subseteq X\times Y$ by \cref{times_subseteq_right}.
        Thus $W\subseteq X\times Y$ by \cref{subseteq_transitive}.
        Thus $W\in\pow{X\times Y}$ by \cref{powerset_intro}.
        Thus $W\in\productBasis{T}{T_1}{X}{Y}$ by \cref{product_basis}.
        Thus $W\in T_2$ by \cref{basis_elem_open_in_generated,product_topology}.
        $C_1 = (A\inter U)\times (B\inter V)$.
        Thus $C_1 = (A\times B)\inter W$ by \cref{inter_times}.
        Thus $C_1\in T_3$ by \cref{subspace_topology}.
        Thus $C_1$ is open in $A\times B$ with respect to $T_3$ by \cref{open_in}.
    \end{subproof}
    Show that for all $U,p$ such that $U\in T_3$ and $p\in U$
        there exists $C_1\in C$ such that $p\in C_1$ and $C_1\subseteq U$.
    \begin{subproof}
        Fix $U$.
        Fix $p$.
        Assume $U\in T_3$ and $p\in U$.
        Take $W\in T_2$ such that $U = (A\times B)\inter W$ by \cref{subspace_topology}.
        Thus $p\in A\times B$ and $p\in W$ by \cref{inter}.
        $\productBasis{T}{T_1}{X}{Y}$ is a basis on $X\times Y$ by \cref{product_basis_is_basis}.
        Thus $\productBasis{T}{T_1}{X}{Y}$ is a basis for $T_2$ on $X\times Y$ by \cref{basis_for_topology,product_topology}.
        Thus $W\in\basisTopology{\productBasis{T}{T_1}{X}{Y}}{X\times Y}$ by \cref{product_topology}.
        Take $W_1\in\productBasis{T}{T_1}{X}{Y}$ such that $p\in W_1$ and $W_1\subseteq W$
            by \cref{topology_generated_by_basis}.
        Take $U_1, V_1$ such that $U_1\in T$ and $V_1\in T_1$ and $W_1 = U_1\times V_1$ by \cref{product_basis}.
        Let $U_A = A\inter U_1$.
        Let $V_B = B\inter V_1$.
        Thus $U_A\in T_A$ and $V_B\in T_B$ by \cref{subspace_topology}.
        Let $C_1 = U_A\times V_B$.
        $T_A\subseteq\pow{A}$ by \cref{topology_on}.
        $T_B\subseteq\pow{B}$ by \cref{topology_on}.
        Thus $U_A\in\pow{A}$ and $V_B\in\pow{B}$ by \cref{subseteq}.
        Thus $U_A\subseteq A$ and $V_B\subseteq B$ by \cref{powerset_elim}.
        $C_1\subseteq A\times V_B$ by \cref{times_subseteq_left}.
        $A\times V_B\subseteq A\times B$ by \cref{times_subseteq_right}.
        Thus $C_1\subseteq A\times B$ by \cref{subseteq_transitive}.
        Thus $C_1\in\pow{A\times B}$ by \cref{powerset_intro}.
        Thus $C_1\in C$ by \cref{product_basis}.
        Take $a,b$ such that $p = (a,b)$ and $a\in A$ and $b\in B$ by \cref{times_elem_is_tuple}.
        $p\in W_1$.
        Take $a_1,b_1$ such that $p = (a_1,b_1)$ and $a_1\in U_1$ and $b_1\in V_1$ by \cref{times_elem_is_tuple}.
        Thus $a = a_1$ and $b = b_1$ by \cref{pair_eq_iff}.
        Thus $a\in U_A$ and $b\in V_B$ by \cref{inter_intro}.
        Thus $p\in C_1$ by \cref{times}.
        $C_1 = (A\inter U_1)\times (B\inter V_1)$.
        Thus $C_1 = (A\times B)\inter W_1$ by \cref{inter_times}.
        $C_1\subseteq A\times B$ by \cref{inter_lower_left}.
        $C_1\subseteq W_1$ by \cref{inter_lower_right}.
        Thus $C_1\subseteq W$ by \cref{subseteq_transitive}.
        Thus $C_1\subseteq (A\times B)\inter W$ by \cref{subseteq_inter_iff}.
        Thus $C_1\subseteq U$.
    \end{subproof}
    Thus $C$ is a basis for $T_3$ on $A\times B$ by \cref{open_collection_basis}.
    Thus $T_3 = \basisTopology{C}{A\times B}$ by \cref{basis_for_topology}.
    Thus $\productTopology{T_A}{T_B}{A}{B} = T_3$ by \cref{product_topology}.
\end{proof}

% from §16 Item 6: convex subset of an ordered set
\begin{definition}\label{convex_in}
    $Y$ is convex in $X$ with respect to $R$ iff
    $Y\subseteq X$ and
    for all $a,b$ such that $a\in Y$ and $b\in Y$ and $a\mathrel{R}b$
        for all $x$ such that $x\in X$ and $a\mathrel{R}x$ and $x\mathrel{R}b$ we have $x\in Y$.
\end{definition}

% from §16 Item 7: convex subspace order topology equals subspace topology
\begin{theorem}\label{convex_order_subspace_topology}
    Assume $T$ is the order topology on $X$ with respect to $R$.
    Assume $Y$ is convex in $X$ with respect to $R$.
    Assume $R$ is a simple order relation on $Y$.
    Assume $Y$ has more than one element.
    Then $\basisTopology{\orderBasis{R}{Y}}{Y} = \subspaceTopology{T}{Y}$.
\end{theorem}
\begin{proof}
    Let $T_Y = \basisTopology{\orderBasis{R}{Y}}{Y}$.
    Let $T_S = \subspaceTopology{T}{Y}$.
    Let $S_Y = \openRays{R}{Y}$.
    Let $S_X = \openRays{R}{X}$.
    $R$ is a simple order relation on $X$ and $X$ has more than one element and
        $T = \basisTopology{\orderBasis{R}{X}}{X}$ by \cref{order_topology}.
    $Y\subseteq X$ by \cref{convex_in}.
    $R$ is a simple order relation on $Y$ and $Y$ has more than one element.
    Thus $\orderBasis{R}{X}$ is a basis on $X$ by \cref{order_basis_is_basis}.
    Thus $\orderBasis{R}{X}$ is a basis for $T$ on $X$ by \cref{basis_for_topology,order_topology}.
    Thus $T$ is a topology on $X$ by \cref{basis_topology_is_topology}.
    Thus $T_S$ is a topology on $Y$ by \cref{subspace_topology_is_topology}.
    Thus $\orderBasis{R}{Y}$ is a basis on $Y$ by \cref{order_basis_is_basis}.
    Thus $T_Y$ is a topology on $Y$ by \cref{basis_topology_is_topology}.
    Thus $T_Y$ is the order topology on $Y$ with respect to $R$ by \cref{order_topology}.
    Thus $\subbasisTopology{S_Y}{Y} = T_Y$ by \cref{open_rays_subbasis}.
    Thus $\subbasisTopology{S_X}{X} = T$ by \cref{open_rays_subbasis}.

    Show that for all $U$ such that $U\in S_X$ we have $U\in T$.
    \begin{subproof}
        Fix $U$.
        Assume $U\in S_X$.
        Let $F = \{U\}$.
        $F\subseteq S_X$ by \cref{singleton_subset_intro}.
        Show that $F$ is finite.
        \begin{subproof}
            Let $G = \{U,U\}$.
            $G$ is finite by \cref{pair_is_finite}.
            $U\in G$ by \cref{upair_intro_left}.
            Thus $F\subseteq G$ by \cref{singleton_subset_intro}.
            Thus $F$ is finite by \cref{subseteq_finite_is_finite}.
        \end{subproof}
        $F$ is inhabited by \cref{singleton_intro}.
        Thus $\inters{F} = U$ by \cref{inters_singleton}.
        $S_X$ is a subbasis on $X$ by \cref{open_rays_subbasis_on}.
        Thus $\unions{S_X} = X$ by \cref{subbasis_on}.
        Take $a\in X$ such that $U = \openrayright{R}{X}{a}$ or $U = \openrayleft{R}{X}{a}$ by \cref{open_rays}.
        Show that $U\subseteq X$.
        \begin{subproof}
            Show that for all $x$ such that $x\in U$ we have $x\in X$.
            \begin{subproof}
                Fix $x$.
                Assume $x\in U$.
                \begin{byCase}
                \caseOf{$U = \openrayright{R}{X}{a}$.}
                    Thus $x\in X$ by \cref{openray_right}.
                \caseOf{$U = \openrayleft{R}{X}{a}$.}
                    Thus $x\in X$ by \cref{openray_left}.
                \end{byCase}
            \end{subproof}
            Thus $U\subseteq X$ by \cref{subseteq}.
        \end{subproof}
        Thus $U\subseteq\unions{S_X}$ by \cref{subseteq_transitive}.
        Thus $U\in\pow{\unions{S_X}}$ by \cref{powerset_intro}.
        Thus $U\in\finiteIntersections{S_X}$ by \cref{finite_intersections}.
        Thus $\finiteIntersections{S_X}$ is a basis on $X$ by \cref{subbasis_finite_intersections_basis}.
        Thus $U\in\basisTopology{\finiteIntersections{S_X}}{X}$ by \cref{basis_elem_open_in_generated}.
        Thus $U\in\subbasisTopology{S_X}{X}$ by \cref{topology_generated_by_subbasis}.
        Thus $U\in T$.
    \end{subproof}

    Show that for all $U$ such that $U\in S_Y$ we have $U\in T_Y$.
    \begin{subproof}
        Fix $U$.
        Assume $U\in S_Y$.
        Let $F = \{U\}$.
        $F\subseteq S_Y$ by \cref{singleton_subset_intro}.
        Show that $F$ is finite.
        \begin{subproof}
            Let $G = \{U,U\}$.
            $G$ is finite by \cref{pair_is_finite}.
            $U\in G$ by \cref{upair_intro_left}.
            Thus $F\subseteq G$ by \cref{singleton_subset_intro}.
            Thus $F$ is finite by \cref{subseteq_finite_is_finite}.
        \end{subproof}
        $F$ is inhabited by \cref{singleton_intro}.
        Thus $\inters{F} = U$ by \cref{inters_singleton}.
        $S_Y$ is a subbasis on $Y$ by \cref{open_rays_subbasis_on}.
        Thus $\unions{S_Y} = Y$ by \cref{subbasis_on}.
        Take $a\in Y$ such that $U = \openrayright{R}{Y}{a}$ or $U = \openrayleft{R}{Y}{a}$ by \cref{open_rays}.
        Show that $U\subseteq Y$.
        \begin{subproof}
            Show that for all $x$ such that $x\in U$ we have $x\in Y$.
            \begin{subproof}
                Fix $x$.
                Assume $x\in U$.
                \begin{byCase}
                \caseOf{$U = \openrayright{R}{Y}{a}$.}
                    Thus $x\in Y$ by \cref{openray_right}.
                \caseOf{$U = \openrayleft{R}{Y}{a}$.}
                    Thus $x\in Y$ by \cref{openray_left}.
                \end{byCase}
            \end{subproof}
            Thus $U\subseteq Y$ by \cref{subseteq}.
        \end{subproof}
        Thus $U\subseteq\unions{S_Y}$ by \cref{subseteq_transitive}.
        Thus $U\in\pow{\unions{S_Y}}$ by \cref{powerset_intro}.
        Thus $U\in\finiteIntersections{S_Y}$ by \cref{finite_intersections}.
        Thus $\finiteIntersections{S_Y}$ is a basis on $Y$ by \cref{subbasis_finite_intersections_basis}.
        Thus $U\in\basisTopology{\finiteIntersections{S_Y}}{Y}$ by \cref{basis_elem_open_in_generated}.
        Thus $U\in\subbasisTopology{S_Y}{Y}$ by \cref{topology_generated_by_subbasis}.
        Thus $U\in T_Y$.
    \end{subproof}

    Show that for all $U$ such that $U\in S_Y$ we have $U\in T_S$.
    \begin{subproof}
        Fix $U$.
        Assume $U\in S_Y$.
        Take $a\in Y$ such that $U = \openrayright{R}{Y}{a}$ or $U = \openrayleft{R}{Y}{a}$ by \cref{open_rays}.
        $a\in X$ by \cref{subseteq}.
        \begin{byCase}
        \caseOf{$U = \openrayright{R}{Y}{a}$.}
            Let $U_1 = \openrayright{R}{X}{a}$.
            Show that for all $x$ such that $x\in U_1$ we have $x\in X$.
            \begin{subproof}
                Fix $x$.
                Assume $x\in U_1$.
                Thus $x\in X$ by \cref{openray_right}.
            \end{subproof}
            Thus $U_1\subseteq X$ by \cref{subseteq}.
            Thus $U_1\in\pow{X}$ by \cref{powerset_intro}.
            $U_1\in S_X$ by \cref{open_rays}.
            Thus $U_1\in T$.
            Show that $U = Y\inter U_1$.
            \begin{subproof}
                Show that for all $x$ such that $x\in U$ we have $x\in Y\inter U_1$.
                \begin{subproof}
                    Fix $x$.
                    Assume $x\in U$.
                    Thus $x\in Y$ and $a\mathrel{R}x$ by \cref{openray_right}.
                    Thus $x\in X$ by \cref{subseteq}.
                    Thus $x\in U_1$ by \cref{openray_right}.
                    Thus $x\in Y\inter U_1$ by \cref{inter_intro}.
                \end{subproof}
                Thus $U\subseteq Y\inter U_1$ by \cref{subseteq}.
                Show that for all $x$ such that $x\in Y\inter U_1$ we have $x\in U$.
                \begin{subproof}
                    Fix $x$.
                    Assume $x\in Y\inter U_1$.
                    Thus $x\in Y$ and $x\in U_1$ by \cref{inter}.
                    Thus $a\mathrel{R}x$ by \cref{openray_right}.
                    Thus $x\in U$ by \cref{openray_right}.
                \end{subproof}
                Thus $Y\inter U_1\subseteq U$ by \cref{subseteq}.
                Thus $U = Y\inter U_1$ by \cref{subseteq_antisymmetric}.
            \end{subproof}
            Thus $U\in T_S$ by \cref{subspace_topology}.
        \caseOf{$U = \openrayleft{R}{Y}{a}$.}
            Let $U_1 = \openrayleft{R}{X}{a}$.
            Show that for all $x$ such that $x\in U_1$ we have $x\in X$.
            \begin{subproof}
                Fix $x$.
                Assume $x\in U_1$.
                Thus $x\in X$ by \cref{openray_left}.
            \end{subproof}
            Thus $U_1\subseteq X$ by \cref{subseteq}.
            Thus $U_1\in\pow{X}$ by \cref{powerset_intro}.
            $U_1\in S_X$ by \cref{open_rays}.
            Thus $U_1\in T$.
            Show that $U = Y\inter U_1$.
            \begin{subproof}
                Show that for all $x$ such that $x\in U$ we have $x\in Y\inter U_1$.
                \begin{subproof}
                    Fix $x$.
                    Assume $x\in U$.
                    Thus $x\in Y$ and $x\mathrel{R}a$ by \cref{openray_left}.
                    Thus $x\in X$ by \cref{subseteq}.
                    Thus $x\in U_1$ by \cref{openray_left}.
                    Thus $x\in Y\inter U_1$ by \cref{inter_intro}.
                \end{subproof}
                Thus $U\subseteq Y\inter U_1$ by \cref{subseteq}.
                Show that for all $x$ such that $x\in Y\inter U_1$ we have $x\in U$.
                \begin{subproof}
                    Fix $x$.
                    Assume $x\in Y\inter U_1$.
                    Thus $x\in Y$ and $x\in U_1$ by \cref{inter}.
                    Thus $x\mathrel{R}a$ by \cref{openray_left}.
                    Thus $x\in U$ by \cref{openray_left}.
                \end{subproof}
                Thus $Y\inter U_1\subseteq U$ by \cref{subseteq}.
                Thus $U = Y\inter U_1$ by \cref{subseteq_antisymmetric}.
            \end{subproof}
            Thus $U\in T_S$ by \cref{subspace_topology}.
        \end{byCase}
    \end{subproof}

    Show that $T_Y\subseteq T_S$.
    \begin{subproof}
        Show that for all $U$ such that $U\in T_Y$ we have $U\in T_S$.
        \begin{subproof}
            Fix $U$.
            Assume $U\in T_Y$.
            Thus $U\in\basisTopology{\finiteIntersections{S_Y}}{Y}$ by \cref{topology_generated_by_subbasis}.
            $S_Y$ is a subbasis on $Y$ by \cref{open_rays_subbasis_on}.
            Thus $\finiteIntersections{S_Y}$ is a basis on $Y$ by \cref{subbasis_finite_intersections_basis}.
            Thus $\finiteIntersections{S_Y}$ is a basis for $T_Y$ on $Y$ by \cref{basis_for_topology,topology_generated_by_subbasis}.
            Thus there exists $M\subseteq\finiteIntersections{S_Y}$ such that $U = \unions{M}$ by \cref{basis_topology_unions}.
            Show that $\finiteIntersections{S_Y}\subseteq T_S$.
            \begin{subproof}
                Show that for all $B$ such that $B\in\finiteIntersections{S_Y}$ we have $B\in T_S$.
                \begin{subproof}
                    Fix $B$.
                    Assume $B\in\finiteIntersections{S_Y}$.
                    Take $F\subseteq S_Y$ such that $F$ is finite and inhabited and $B = \inters{F}$ by \cref{finite_intersections}.
                    Show that $F\subseteq T_S$.
                    \begin{subproof}
                        Show that for all $V$ such that $V\in F$ we have $V\in T_S$.
                        \begin{subproof}
                            Fix $V$.
                            Assume $V\in F$.
                            Thus $V\in S_Y$ by \cref{subseteq}.
                            Thus $V\in T_S$.
                        \end{subproof}
                        Thus $F\subseteq T_S$ by \cref{subseteq}.
                    \end{subproof}
                    Thus $\inters{F}\in T_S$ by \cref{finite_intersections_open_in_topology}.
                    Thus $B\in T_S$.
                \end{subproof}
                Thus $\finiteIntersections{S_Y}\subseteq T_S$ by \cref{subseteq}.
            \end{subproof}
            Thus $M\subseteq T_S$ by \cref{subseteq}.
            Thus $\unions{M}\in T_S$ by \cref{topology_on}.
            Thus $U\in T_S$.
        \end{subproof}
        Thus $T_Y\subseteq T_S$ by \cref{subseteq}.
    \end{subproof}

    Show that for all $c\in X$ we have $Y\inter\openrayright{R}{X}{c}\in T_Y$.
    \begin{subproof}
        Fix $c$.
        Assume $c\in X$.
        \begin{byCase}
        \caseOf{$c\in Y$.}
            Show that $Y\inter\openrayright{R}{X}{c} = \openrayright{R}{Y}{c}$.
            \begin{subproof}
                Show that for all $x$ such that $x\in Y\inter\openrayright{R}{X}{c}$ we have $x\in \openrayright{R}{Y}{c}$.
                \begin{subproof}
                    Fix $x$.
                    Assume $x\in Y\inter\openrayright{R}{X}{c}$.
                    Thus $x\in Y$ and $c\mathrel{R}x$ by \cref{inter,openray_right}.
                    Thus $x\in \openrayright{R}{Y}{c}$ by \cref{openray_right}.
                \end{subproof}
                Thus $Y\inter\openrayright{R}{X}{c}\subseteq \openrayright{R}{Y}{c}$ by \cref{subseteq}.
                Show that for all $x$ such that $x\in \openrayright{R}{Y}{c}$ we have $x\in Y\inter\openrayright{R}{X}{c}$.
                \begin{subproof}
                    Fix $x$.
                    Assume $x\in \openrayright{R}{Y}{c}$.
                    Thus $x\in Y$ and $c\mathrel{R}x$ by \cref{openray_right}.
                    Thus $x\in X$ by \cref{subseteq}.
                    Thus $x\in \openrayright{R}{X}{c}$ by \cref{openray_right}.
                    Thus $x\in Y\inter\openrayright{R}{X}{c}$ by \cref{inter_intro}.
                \end{subproof}
                Thus $\openrayright{R}{Y}{c}\subseteq Y\inter\openrayright{R}{X}{c}$ by \cref{subseteq}.
                Thus $Y\inter\openrayright{R}{X}{c} = \openrayright{R}{Y}{c}$ by \cref{subseteq_antisymmetric}.
            \end{subproof}
            Show that for all $x$ such that $x\in \openrayright{R}{Y}{c}$ we have $x\in Y$.
            \begin{subproof}
                Fix $x$.
                Assume $x\in \openrayright{R}{Y}{c}$.
                Thus $x\in Y$ by \cref{openray_right}.
            \end{subproof}
            Thus $\openrayright{R}{Y}{c}\subseteq Y$ by \cref{subseteq}.
            Thus $\openrayright{R}{Y}{c}\in\pow{Y}$ by \cref{powerset_intro}.
            $\openrayright{R}{Y}{c}\in S_Y$ by \cref{open_rays}.
            Thus $Y\inter\openrayright{R}{X}{c}\in T_Y$.
        \caseOf{$c\notin Y$.}
            Take $y_0, y_1\in Y$ such that $y_0\neq y_1$ by \cref{more_than_one_element}.
            $y_0\in X$ by \cref{subseteq}.
            Show that $y_0\neq c$.
            \begin{subproof}
                Suppose not.
                Thus $y_0 = c$.
                Thus $c\in Y$.
                Contradiction.
            \end{subproof}
            Thus $y_0\neq c$.
            Thus $c\neq y_0$.
            $c\in X$.
            $R$ is connex on $X$ by \cref{simple_order_relation}.
            Thus $c\mathrel{R}y_0$ or $y_0\mathrel{R}c$ by \cref{connex}.
            \begin{byCase}
            \caseOf{$c\mathrel{R}y_0$.}
                Show that for all $y$ such that $y\in Y$ we have $c\mathrel{R}y$.
                \begin{subproof}
                    Fix $y$.
                    Assume $y\in Y$.
                    Suppose not.
                    $y\in X$ by \cref{subseteq}.
                    Show that $y\neq c$.
                    \begin{subproof}
                        Suppose not.
                        Thus $y = c$.
                        Thus $c\in Y$.
                        Contradiction.
                    \end{subproof}
                    Thus $c\mathrel{R}y$ or $y\mathrel{R}c$ by \cref{connex}.
                    Thus $y\mathrel{R}c$.
                    $R$ is transitive by \cref{simple_order_relation}.
                    Thus $y\mathrel{R}y_0$ by \hyperref[transitive]{transitivity}.
                    Thus $c\in Y$ by \cref{convex_in}.
                    Contradiction.
                \end{subproof}
                Show that $Y\subseteq \openrayright{R}{X}{c}$.
                \begin{subproof}
                    Show that for all $y$ such that $y\in Y$ we have $y\in \openrayright{R}{X}{c}$.
                    \begin{subproof}
                        Fix $y$.
                        Assume $y\in Y$.
                        Thus $c\mathrel{R}y$.
                        Thus $y\in X$ by \cref{elem_subseteq}.
                        Thus $y\in \openrayright{R}{X}{c}$ by \cref{openray_right}.
                    \end{subproof}
                    Thus $Y\subseteq \openrayright{R}{X}{c}$ by \cref{subseteq}.
                \end{subproof}
                Thus $Y\inter\openrayright{R}{X}{c} = Y$ by \cref{inter_absorb_supseteq_right}.
                Thus $Y\inter\openrayright{R}{X}{c}\in T_Y$.
            \caseOf{$y_0\mathrel{R}c$.}
                Show that for all $y$ such that $y\in Y$ we have $y\mathrel{R}c$.
                \begin{subproof}
                    Fix $y$.
                    Assume $y\in Y$.
                    Suppose not.
                    $y\in X$ by \cref{subseteq}.
                    Show that $y\neq c$.
                    \begin{subproof}
                        Suppose not.
                        Thus $y = c$.
                        Thus $c\in Y$.
                        Contradiction.
                    \end{subproof}
                    Thus $c\mathrel{R}y$ or $y\mathrel{R}c$ by \cref{connex}.
                    Thus $c\mathrel{R}y$.
                    $R$ is transitive by \cref{simple_order_relation}.
                    Thus $y_0\mathrel{R}y$ by \hyperref[transitive]{transitivity}.
                    Thus $c\in Y$ by \cref{convex_in}.
                    Contradiction.
                \end{subproof}
                Show that $Y\inter\openrayright{R}{X}{c} = \emptyset$.
                \begin{subproof}
                    Suppose not.
                    Thus $Y\inter\openrayright{R}{X}{c}\neq\emptyset$.
                    Thus $Y\inter\openrayright{R}{X}{c}$ is inhabited by \cref{nonempty_inhabited}.
                    Take $x$ such that $x\in Y\inter\openrayright{R}{X}{c}$.
                    Thus $x\in Y$ and $c\mathrel{R}x$ by \cref{inter,openray_right}.
                    Thus $x\mathrel{R}c$.
                    $R$ is asymmetric by \cref{simple_order_relation}.
                    Contradiction by \cref{asymmetric}.
                \end{subproof}
                Thus $Y\inter\openrayright{R}{X}{c}\in T_Y$.
            \end{byCase}
        \end{byCase}
    \end{subproof}

    Show that for all $c\in X$ we have $Y\inter\openrayleft{R}{X}{c}\in T_Y$.
    \begin{subproof}
        Fix $c$.
        Assume $c\in X$.
        \begin{byCase}
        \caseOf{$c\in Y$.}
            Show that $Y\inter\openrayleft{R}{X}{c} = \openrayleft{R}{Y}{c}$.
            \begin{subproof}
                Show that for all $x$ such that $x\in Y\inter\openrayleft{R}{X}{c}$ we have $x\in \openrayleft{R}{Y}{c}$.
                \begin{subproof}
                    Fix $x$.
                    Assume $x\in Y\inter\openrayleft{R}{X}{c}$.
                    Thus $x\in Y$ and $x\mathrel{R}c$ by \cref{inter,openray_left}.
                    Thus $x\in \openrayleft{R}{Y}{c}$ by \cref{openray_left}.
                \end{subproof}
                Thus $Y\inter\openrayleft{R}{X}{c}\subseteq \openrayleft{R}{Y}{c}$ by \cref{subseteq}.
                Show that for all $x$ such that $x\in \openrayleft{R}{Y}{c}$ we have $x\in Y\inter\openrayleft{R}{X}{c}$.
                \begin{subproof}
                    Fix $x$.
                    Assume $x\in \openrayleft{R}{Y}{c}$.
                    Thus $x\in Y$ and $x\mathrel{R}c$ by \cref{openray_left}.
                    Thus $x\in X$ by \cref{subseteq}.
                    Thus $x\in \openrayleft{R}{X}{c}$ by \cref{openray_left}.
                    Thus $x\in Y\inter\openrayleft{R}{X}{c}$ by \cref{inter_intro}.
                \end{subproof}
                Thus $\openrayleft{R}{Y}{c}\subseteq Y\inter\openrayleft{R}{X}{c}$ by \cref{subseteq}.
                Thus $Y\inter\openrayleft{R}{X}{c} = \openrayleft{R}{Y}{c}$ by \cref{subseteq_antisymmetric}.
            \end{subproof}
            Show that for all $x$ such that $x\in \openrayleft{R}{Y}{c}$ we have $x\in Y$.
            \begin{subproof}
                Fix $x$.
                Assume $x\in \openrayleft{R}{Y}{c}$.
                Thus $x\in Y$ by \cref{openray_left}.
            \end{subproof}
            Thus $\openrayleft{R}{Y}{c}\subseteq Y$ by \cref{subseteq}.
            Thus $\openrayleft{R}{Y}{c}\in\pow{Y}$ by \cref{powerset_intro}.
            $\openrayleft{R}{Y}{c}\in S_Y$ by \cref{open_rays}.
            Thus $Y\inter\openrayleft{R}{X}{c}\in T_Y$.
        \caseOf{$c\notin Y$.}
            Take $y_0, y_1\in Y$ such that $y_0\neq y_1$ by \cref{more_than_one_element}.
            $y_0\in X$ by \cref{subseteq}.
            Show that $y_0\neq c$.
            \begin{subproof}
                Suppose not.
                Thus $y_0 = c$.
                Thus $c\in Y$.
                Contradiction.
            \end{subproof}
            Thus $y_0\neq c$.
            Thus $c\neq y_0$.
            $c\in X$.
            $R$ is connex on $X$ by \cref{simple_order_relation}.
            Thus $c\mathrel{R}y_0$ or $y_0\mathrel{R}c$ by \cref{connex}.
            \begin{byCase}
            \caseOf{$y_0\mathrel{R}c$.}
                Show that for all $y$ such that $y\in Y$ we have $y\mathrel{R}c$.
                \begin{subproof}
                    Fix $y$.
                    Assume $y\in Y$.
                    Suppose not.
                    $y\in X$ by \cref{subseteq}.
                    Show that $y\neq c$.
                    \begin{subproof}
                        Suppose not.
                        Thus $y = c$.
                        Thus $c\in Y$.
                        Contradiction.
                    \end{subproof}
                    Thus $c\mathrel{R}y$ or $y\mathrel{R}c$ by \cref{connex}.
                    Thus $c\mathrel{R}y$.
                    $R$ is transitive by \cref{simple_order_relation}.
                    Thus $y_0\mathrel{R}y$ by \hyperref[transitive]{transitivity}.
                    Thus $c\in Y$ by \cref{convex_in}.
                    Contradiction.
                \end{subproof}
                Show that $Y\subseteq \openrayleft{R}{X}{c}$.
                \begin{subproof}
                    Show that for all $y$ such that $y\in Y$ we have $y\in \openrayleft{R}{X}{c}$.
                    \begin{subproof}
                        Fix $y$.
                        Assume $y\in Y$.
                        Thus $y\mathrel{R}c$.
                        Thus $y\in X$ by \cref{elem_subseteq}.
                        Thus $y\in \openrayleft{R}{X}{c}$ by \cref{openray_left}.
                    \end{subproof}
                    Thus $Y\subseteq \openrayleft{R}{X}{c}$ by \cref{subseteq}.
                \end{subproof}
                Thus $Y\inter\openrayleft{R}{X}{c} = Y$ by \cref{inter_absorb_supseteq_right}.
                Thus $Y\inter\openrayleft{R}{X}{c}\in T_Y$.
            \caseOf{$c\mathrel{R}y_0$.}
                Show that for all $y$ such that $y\in Y$ we have $c\mathrel{R}y$.
                \begin{subproof}
                    Fix $y$.
                    Assume $y\in Y$.
                    Suppose not.
                    $y\in X$ by \cref{subseteq}.
                    Show that $y\neq c$.
                    \begin{subproof}
                        Suppose not.
                        Thus $y = c$.
                        Thus $c\in Y$.
                        Contradiction.
                    \end{subproof}
                    Thus $c\mathrel{R}y$ or $y\mathrel{R}c$ by \cref{connex}.
                    Thus $y\mathrel{R}c$.
                    $R$ is transitive by \cref{simple_order_relation}.
                    Thus $y\mathrel{R}y_0$ by \hyperref[transitive]{transitivity}.
                    Thus $c\in Y$ by \cref{convex_in}.
                    Contradiction.
                \end{subproof}
                Show that $Y\inter\openrayleft{R}{X}{c} = \emptyset$.
                \begin{subproof}
                    Suppose not.
                    Thus $Y\inter\openrayleft{R}{X}{c}\neq\emptyset$.
                    Thus $Y\inter\openrayleft{R}{X}{c}$ is inhabited by \cref{nonempty_inhabited}.
                    Take $x$ such that $x\in Y\inter\openrayleft{R}{X}{c}$.
                    Thus $x\in Y$ and $x\mathrel{R}c$ by \cref{inter,openray_left}.
                    Thus $c\mathrel{R}x$.
                    $R$ is asymmetric by \cref{simple_order_relation}.
                    Contradiction by \cref{asymmetric}.
                \end{subproof}
                Thus $Y\inter\openrayleft{R}{X}{c}\in T_Y$.
            \end{byCase}
        \end{byCase}
    \end{subproof}

    Show that for all $B$ such that $B\in\orderBasis{R}{X}$ we have $Y\inter B\in T_Y$.
    \begin{subproof}
        Fix $B$.
        Assume $B\in\orderBasis{R}{X}$.
        $B\in\pow{X}$ and
            $((\exists a, b\in X. a\mathrel{R}b \land B = \intervalopenrel{R}{X}{a}{b}) \lor
            (\exists b\in X. \exists a_0\in X. (\forall x\in X. a_0\mathrel{R}x \lor a_0 = x) \land B = \intervalclosedopenrel{R}{X}{a_0}{b}) \lor
            (\exists a\in X. \exists b_0\in X. (\forall x\in X. x\mathrel{R}b_0 \lor x = b_0) \land B = \intervalopenclosedrel{R}{X}{a}{b_0}))$
            by \cref{order_basis}.
        \begin{byCase}
        \caseOf{$\exists a, b\in X. a\mathrel{R}b \land B = \intervalopenrel{R}{X}{a}{b}$.}
            Take $a, b\in X$ such that $a\mathrel{R}b$ and $B = \intervalopenrel{R}{X}{a}{b}$.
            Thus $B = \openrayright{R}{X}{a} \inter \openrayleft{R}{X}{b}$ by \cref{intervalopenrel_as_intersection_openrays}.
            $Y\inter\openrayright{R}{X}{a}\in T_Y$.
            $Y\inter\openrayleft{R}{X}{b}\in T_Y$.
            $T_Y$ is closed under binary intersections by \cref{topology_on}.
            Show that $Y\inter B = (Y\inter\openrayright{R}{X}{a}) \inter (Y\inter\openrayleft{R}{X}{b})$.
            \begin{subproof}
                Show that for all $x$ such that $x\in Y\inter B$
                    we have $x\in (Y\inter\openrayright{R}{X}{a}) \inter (Y\inter\openrayleft{R}{X}{b})$.
                \begin{subproof}
                    Fix $x$.
                    Assume $x\in Y\inter B$.
                    Thus $x\in Y$ and $x\in B$ by \cref{inter}.
                    Thus $x\in \openrayright{R}{X}{a}$ and $x\in \openrayleft{R}{X}{b}$ by \cref{inter}.
                    Thus $x\in Y\inter\openrayright{R}{X}{a}$ by \cref{inter_intro}.
                    Thus $x\in Y\inter\openrayleft{R}{X}{b}$ by \cref{inter_intro}.
                    Thus $x\in (Y\inter\openrayright{R}{X}{a}) \inter (Y\inter\openrayleft{R}{X}{b})$ by \cref{inter_intro}.
                \end{subproof}
                Thus $Y\inter B\subseteq (Y\inter\openrayright{R}{X}{a}) \inter (Y\inter\openrayleft{R}{X}{b})$ by \cref{subseteq}.
                Show that for all $x$ such that $x\in (Y\inter\openrayright{R}{X}{a}) \inter (Y\inter\openrayleft{R}{X}{b})$
                    we have $x\in Y\inter B$.
                \begin{subproof}
                    Fix $x$.
                    Assume $x\in (Y\inter\openrayright{R}{X}{a}) \inter (Y\inter\openrayleft{R}{X}{b})$.
                    Thus $x\in Y\inter\openrayright{R}{X}{a}$ and $x\in Y\inter\openrayleft{R}{X}{b}$ by \cref{inter}.
                    Thus $x\in Y$ and $x\in \openrayright{R}{X}{a}$ by \cref{inter}.
                    Thus $x\in Y$ and $x\in \openrayleft{R}{X}{b}$ by \cref{inter}.
                    Thus $x\in \openrayright{R}{X}{a} \inter \openrayleft{R}{X}{b}$ by \cref{inter_intro}.
                    Thus $x\in B$.
                    Thus $x\in Y\inter B$ by \cref{inter_intro}.
                \end{subproof}
                Thus $(Y\inter\openrayright{R}{X}{a}) \inter (Y\inter\openrayleft{R}{X}{b})\subseteq Y\inter B$ by \cref{subseteq}.
                Thus $Y\inter B = (Y\inter\openrayright{R}{X}{a}) \inter (Y\inter\openrayleft{R}{X}{b})$
                    by \cref{subseteq_antisymmetric}.
            \end{subproof}
            Thus $Y\inter B\in T_Y$.
        \caseOf{$\exists b\in X. \exists a_0\in X. (\forall x\in X. a_0\mathrel{R}x \lor a_0 = x) \land B = \intervalclosedopenrel{R}{X}{a_0}{b}$.}
            Take $a_0, b\in X$ such that
                $(\forall x\in X. a_0\mathrel{R}x \lor a_0 = x)$ and $B = \intervalclosedopenrel{R}{X}{a_0}{b}$.
            Thus $a_0$ is a smallest element of $X$ with respect to $R$ by \cref{smallest_element}.
            Thus $B = \openrayleft{R}{X}{b}$ by \cref{intervalclosedopenrel_smallest_eq_openrayleft}.
            Thus $Y\inter B = Y\inter\openrayleft{R}{X}{b}$.
            Thus $Y\inter B\in T_Y$.
        \caseOf{$\exists a\in X. \exists b_0\in X. (\forall x\in X. x\mathrel{R}b_0 \lor x = b_0) \land B = \intervalopenclosedrel{R}{X}{a}{b_0}$.}
            Take $a, b_0\in X$ such that
                $(\forall x\in X. x\mathrel{R}b_0 \lor x = b_0)$ and $B = \intervalopenclosedrel{R}{X}{a}{b_0}$.
            Thus $b_0$ is a largest element of $X$ with respect to $R$ by \cref{largest_element}.
            Thus $B = \openrayright{R}{X}{a}$ by \cref{intervalopenclosedrel_largest_eq_openrayright}.
            Thus $Y\inter B = Y\inter\openrayright{R}{X}{a}$.
            Thus $Y\inter B\in T_Y$.
        \end{byCase}
    \end{subproof}

    Show that $T_S\subseteq T_Y$.
    \begin{subproof}
        Show that for all $U$ such that $U\in T_S$ we have $U\in T_Y$.
        \begin{subproof}
            Fix $U$.
            Assume $U\in T_S$.
            Take $V\in T$ such that $U = Y\inter V$ by \cref{subspace_topology}.
            Thus there exists $M\subseteq\orderBasis{R}{X}$ such that $V = \unions{M}$ by \cref{basis_topology_unions}.
            Let $N = \{A\in\pow{Y} \mid \exists B\in M. A = Y\inter B\}$.
            Show that $U = \unions{N}$.
            \begin{subproof}
                Show that for all $x$ such that $x\in U$ we have $x\in\unions{N}$.
                \begin{subproof}
                    Fix $x$.
                    Assume $x\in U$.
                    Thus $x\in Y$ and $x\in V$ by \cref{inter}.
                    Take $B\in M$ such that $x\in B$ by \cref{unions_iff}.
                    Thus $x\in Y\inter B$ by \cref{inter_intro}.
                    Thus $Y\inter B\in N$.
                    Thus $x\in\unions{N}$ by \cref{unions_iff}.
                \end{subproof}
                Thus $U\subseteq\unions{N}$ by \cref{subseteq}.
                Show that for all $x$ such that $x\in\unions{N}$ we have $x\in U$.
                \begin{subproof}
                    Fix $x$.
                    Assume $x\in\unions{N}$.
                    Take $A\in N$ such that $x\in A$ by \cref{unions_iff}.
                    Take $B\in M$ such that $A = Y\inter B$.
                    Thus $x\in Y$ and $x\in B$ by \cref{inter}.
                    Thus $x\in V$ by \cref{unions_iff}.
                    Thus $x\in U$ by \cref{inter_intro}.
                \end{subproof}
                Thus $\unions{N}\subseteq U$ by \cref{subseteq}.
                Thus $U = \unions{N}$ by \cref{subseteq_antisymmetric}.
            \end{subproof}
            Show that $N\subseteq T_Y$.
            \begin{subproof}
                Show that for all $A$ such that $A\in N$ we have $A\in T_Y$.
                \begin{subproof}
                    Fix $A$.
                    Assume $A\in N$.
                    Take $B\in M$ such that $A = Y\inter B$.
                    Thus $B\in\orderBasis{R}{X}$ by \cref{subseteq}.
                    Thus $A\in T_Y$.
                \end{subproof}
                Thus $N\subseteq T_Y$ by \cref{subseteq}.
            \end{subproof}
            Thus $\unions{N}\in T_Y$ by \cref{topology_on}.
            Thus $U\in T_Y$.
        \end{subproof}
        Thus $T_S\subseteq T_Y$ by \cref{subseteq}.
    \end{subproof}

    Thus $T_Y = T_S$ by \cref{subseteq_antisymmetric}.
\end{proof}

% from §17 Item 1: closed set
\begin{definition}\label{closed_in}
    $A$ is closed in $X$ with respect to $T$ iff
    $A\subseteq X$ and $X\setminus A$ is open in $X$ with respect to $T$.
\end{definition}

% from §17 Item 2: empty set is closed
\begin{lemma}\label{closed_emptyset}
    Assume $T$ is a topology on $X$.
    Then $\emptyset$ is closed in $X$ with respect to $T$.
\end{lemma}
\begin{proof}
    $\emptyset\subseteq X$ by \cref{emptyset_subseteq}.
    $X\setminus \emptyset = X$ by \cref{setminus_emptyset}.
    $X\in T$ by \cref{topology_on}.
    Thus $X$ is open in $X$ with respect to $T$ by \cref{open_in}.
    Thus $\emptyset$ is closed in $X$ with respect to $T$ by \cref{closed_in}.
\end{proof}

% from §17 Item 3: the whole space is closed
\begin{lemma}\label{closed_whole_space}
    Assume $T$ is a topology on $X$.
    Then $X$ is closed in $X$ with respect to $T$.
\end{lemma}
\begin{proof}
    $X\subseteq X$ by \cref{subseteq_refl}.
    $X\setminus X = \emptyset$ by \cref{setminus_self}.
    $\emptyset\in T$ by \cref{topology_on}.
    Thus $\emptyset$ is open in $X$ with respect to $T$ by \cref{open_in}.
    Thus $X$ is closed in $X$ with respect to $T$ by \cref{closed_in}.
\end{proof}

% from §17 Item 4: intersections of closed sets are closed
\begin{lemma}\label{closed_inters}
    Assume $T$ is a topology on $X$.
    Suppose $C\subseteq\pow{X}$.
    Suppose for all $A$ such that $A\in C$ we have $A$ is closed in $X$ with respect to $T$.
    Then $\inters{C}$ is closed in $X$ with respect to $T$.
\end{lemma}
\begin{proof}
    Show that $\inters{C}\subseteq X$.
    \begin{subproof}
        Show that for all $x$ such that $x\in\inters{C}$ we have $x\in X$.
        \begin{subproof}
            Fix $x$.
            Assume $x\in\inters{C}$.
            Thus $x\in\unions{C}$ by \cref{inters}.
            $\unions{C}\subseteq X$ by \cref{unions_subseteq_of_powerset_is_subseteq}.
            Thus $x\in X$ by \cref{elem_subseteq}.
        \end{subproof}
        Thus $\inters{C}\subseteq X$ by \cref{subseteq}.
    \end{subproof}
    \begin{byCase}
    \caseOf{$C = \emptyset$.}
        $\unions{C} = \emptyset$ by \cref{unions_emptyset}.
        Show that $\inters{C} = \emptyset$.
        \begin{subproof}
            Suppose not.
            Then $\inters{C}\neq \emptyset$.
            Thus $\inters{C}$ is inhabited by \cref{nonempty_inhabited}.
            Take $x$ such that $x\in\inters{C}$.
            Thus $x\in\unions{C}$ by \cref{inters}.
            Thus $x\in\emptyset$ by \cref{unions_emptyset}.
            Contradiction by \cref{emptyset}.
        \end{subproof}
        Thus $\inters{C}$ is closed in $X$ with respect to $T$ by \cref{closed_emptyset}.
    \caseOf{$C \neq \emptyset$.}
        $C$ is inhabited by \cref{nonempty_inhabited}.
        Let $D = \{U\in T \mid \exists A\in C. U = X\setminus A\}$.
        Show that $X\setminus\inters{C} = \unions{D}$.
        \begin{subproof}
            Show that for all $x$ such that $x\in X\setminus\inters{C}$ we have $x\in\unions{D}$.
            \begin{subproof}
                Fix $x$.
                Assume $x\in X\setminus\inters{C}$.
                Suppose not.
                Then $x\notin\unions{D}$.
                Thus $x\in X$ by \cref{setminus_elim_left}.
                Show that for all $A$ such that $A\in C$ we have $x\in A$.
                \begin{subproof}
                    Fix $A$.
                    Assume $A\in C$.
                    Suppose not.
                    Thus $x\notin A$.
                    Thus $x\in X\setminus A$ by \cref{setminus_intro}.
                    $A$ is closed in $X$ with respect to $T$.
                    Thus $X\setminus A$ is open in $X$ with respect to $T$ by \cref{closed_in}.
                    Thus $X\setminus A\in T$ by \cref{open_in}.
                    Thus $X\setminus A\in D$.
                    Thus $x\in\unions{D}$ by \cref{unions_iff}.
                    Contradiction.
                \end{subproof}
                Thus $x\in\inters{C}$ by \cref{inters_intro}.
                Thus $x\notin\inters{C}$ by \cref{setminus_elim_right}.
                Contradiction.
            \end{subproof}
            Thus $X\setminus\inters{C}\subseteq\unions{D}$ by \cref{subseteq}.
            Show that for all $x$ such that $x\in\unions{D}$ we have $x\in X\setminus\inters{C}$.
            \begin{subproof}
                Fix $x$.
                Assume $x\in\unions{D}$.
                Take $U\in D$ such that $x\in U$ by \cref{unions_iff}.
                Take $A\in C$ such that $U = X\setminus A$.
                Thus $x\in X\setminus A$.
                Thus $x\in X$ by \cref{setminus_elim_left}.
                Thus $x\notin A$ by \cref{setminus_elim_right}.
                Show that $x\notin\inters{C}$.
                \begin{subproof}
                    Suppose not.
                    Then $x\in\inters{C}$.
                    Then $x\in A$ by \cref{inters_destr}.
                    Contradiction.
                \end{subproof}
                Thus $x\in X\setminus\inters{C}$ by \cref{setminus_intro}.
            \end{subproof}
            Thus $\unions{D}\subseteq X\setminus\inters{C}$ by \cref{subseteq}.
            Thus $X\setminus\inters{C} = \unions{D}$ by \cref{subseteq_antisymmetric}.
        \end{subproof}
        Show that $D\subseteq T$.
        \begin{subproof}
            Show that for all $U$ such that $U\in D$ we have $U\in T$.
            \begin{subproof}
                Fix $U$.
                Assume $U\in D$.
                Thus $U\in T$ by \cref{subseteq}.
            \end{subproof}
            Thus $D\subseteq T$ by \cref{subseteq}.
        \end{subproof}
        Thus $\unions{D}\in T$ by \cref{topology_on}.
        Thus $X\setminus\inters{C}$ is open in $X$ with respect to $T$ by \cref{open_in}.
        Thus $\inters{C}$ is closed in $X$ with respect to $T$ by \cref{closed_in}.
    \end{byCase}
\end{proof}

% from §17 helper lemma 1: complements of subsets are injective
\begin{lemma}\label{complement_eq_implies_eq}
    Suppose $A\subseteq X$ and $B\subseteq X$.
    Suppose $X\setminus A = X\setminus B$.
    Then $A = B$.
\end{lemma}
\begin{proof}
    Show that $A\subseteq B$.
    \begin{subproof}
        Show that for all $x$ such that $x\in A$ we have $x\in B$.
        \begin{subproof}
            Fix $x$.
            Assume $x\in A$.
            Thus $x\in X$ by \cref{elem_subseteq}.
            Suppose not.
            Then $x\notin B$.
            Thus $x\in X\setminus B$ by \cref{setminus_intro}.
            Thus $x\in X\setminus A$.
            Thus $x\notin A$ by \cref{setminus_elim_right}.
            Contradiction.
        \end{subproof}
        Thus $A\subseteq B$ by \cref{subseteq}.
    \end{subproof}
    Show that $B\subseteq A$.
    \begin{subproof}
        Show that for all $x$ such that $x\in B$ we have $x\in A$.
        \begin{subproof}
            Fix $x$.
            Assume $x\in B$.
            Thus $x\in X$ by \cref{elem_subseteq}.
            Suppose not.
            Then $x\notin A$.
            Thus $x\in X\setminus A$ by \cref{setminus_intro}.
            Thus $x\in X\setminus B$.
            Thus $x\notin B$ by \cref{setminus_elim_right}.
            Contradiction.
        \end{subproof}
        Thus $B\subseteq A$ by \cref{subseteq}.
    \end{subproof}
    Thus $A = B$ by \cref{subseteq_antisymmetric}.
\end{proof}

% from §17 Item 5: finite unions of closed sets are closed
\begin{lemma}\label{closed_finite_unions}
    Assume $T$ is a topology on $X$.
    Suppose $F\subseteq\pow{X}$.
    Suppose $F$ is finite.
    Suppose for all $A$ such that $A\in F$ we have $A$ is closed in $X$ with respect to $T$.
    Then $\unions{F}$ is closed in $X$ with respect to $T$.
\end{lemma}
\begin{proof}
    \begin{byCase}
    \caseOf{$F = \emptyset$.}
        $\unions{F} = \emptyset$ by \cref{unions_emptyset}.
        Thus $\unions{F}$ is closed in $X$ with respect to $T$ by \cref{closed_emptyset}.
    \caseOf{$F \neq \emptyset$.}
        $F$ is inhabited by \cref{nonempty_inhabited}.
        Let $G = \{U\in T \mid \exists A\in F. U = X\setminus A\}$.
        Show that $G$ is inhabited.
        \begin{subproof}
            Take $A$ such that $A\in F$.
            $A$ is closed in $X$ with respect to $T$.
            Thus $X\setminus A$ is open in $X$ with respect to $T$ by \cref{closed_in}.
            Thus $X\setminus A\in T$ by \cref{open_in}.
            Thus $X\setminus A\in G$.
        \end{subproof}
        Show that $X\setminus\unions{F} = \inters{G}$.
        \begin{subproof}
            Show that for all $x$ such that $x\in X\setminus\unions{F}$ we have $x\in\inters{G}$.
            \begin{subproof}
                Fix $x$.
                Assume $x\in X\setminus\unions{F}$.
                Thus $x\in X$ by \cref{setminus_elim_left}.
                Thus $x\notin\unions{F}$ by \cref{setminus_elim_right}.
                Show that for all $U$ such that $U\in G$ we have $x\in U$.
                \begin{subproof}
                    Fix $U$.
                    Assume $U\in G$.
                    Take $A\in F$ such that $U = X\setminus A$.
                    Suppose not.
                    Then $x\notin U$.
                    Thus $x\notin X\setminus A$.
                    Show that $x\in A$.
                    \begin{subproof}
                        Suppose not.
                        Then $x\notin A$.
                        Thus $x\in X\setminus A$ by \cref{setminus_intro}.
                        Contradiction.
                    \end{subproof}
                    Thus $x\in\unions{F}$ by \cref{unions_iff}.
                    Contradiction.
                    Thus $x\in U$.
                \end{subproof}
                Thus $x\in\inters{G}$ by \cref{inters_intro}.
            \end{subproof}
            Thus $X\setminus\unions{F}\subseteq\inters{G}$ by \cref{subseteq}.
            Show that for all $x$ such that $x\in\inters{G}$ we have $x\in X\setminus\unions{F}$.
            \begin{subproof}
                Fix $x$.
                Assume $x\in\inters{G}$.
                Take $A$ such that $A\in F$.
                $A$ is closed in $X$ with respect to $T$.
                Thus $X\setminus A$ is open in $X$ with respect to $T$ by \cref{closed_in}.
                Thus $X\setminus A\in T$ by \cref{open_in}.
                Thus $X\setminus A\in G$.
                Thus $x\in X\setminus A$ by \cref{inters_destr}.
                Thus $x\in X$ by \cref{setminus_elim_left}.
                Show that $x\notin\unions{F}$.
                \begin{subproof}
                    Suppose not.
                    Then $x\in\unions{F}$.
                    Take $A_1\in F$ such that $x\in A_1$ by \cref{unions_iff}.
                    $A_1$ is closed in $X$ with respect to $T$.
                    Thus $X\setminus A_1$ is open in $X$ with respect to $T$ by \cref{closed_in}.
                    Thus $X\setminus A_1\in T$ by \cref{open_in}.
                    Thus $X\setminus A_1\in G$.
                    Thus $x\in X\setminus A_1$ by \cref{inters_destr}.
                    Thus $x\notin A_1$ by \cref{setminus_elim_right}.
                    Contradiction.
                \end{subproof}
                Thus $x\in X\setminus\unions{F}$ by \cref{setminus_intro}.
            \end{subproof}
            Thus $\inters{G}\subseteq X\setminus\unions{F}$ by \cref{subseteq}.
            Thus $X\setminus\unions{F} = \inters{G}$ by \cref{subseteq_antisymmetric}.
        \end{subproof}
        Show that $G\subseteq T$.
        \begin{subproof}
            Show that for all $U$ such that $U\in G$ we have $U\in T$.
            \begin{subproof}
                Fix $U$.
                Assume $U\in G$.
                Thus $U\in T$ by \cref{subseteq}.
            \end{subproof}
            Thus $G\subseteq T$ by \cref{subseteq}.
        \end{subproof}
        Show that $G$ is finite.
        \begin{subproof}
            Take $k$ such that $k$ is a natural number and there exists a bijection from $k$ to $F$ by \cref{finite}.
            Take $f$ such that $f$ is a bijection from $k$ to $F$.
            Let $g = \{(n, X\setminus f(n)) \mid n\in k\}$.
            Show that $g$ is a bijection from $k$ to $G$.
            \begin{subproof}
                Show that $g$ is a function to $G$ and $\dom{g} = k$.
                \begin{subproof}
                    Show that $g$ is a function.
                    \begin{subproof}
                        Show that $g$ is a relation.
                        \begin{subproof}
                            Show that for all $w$ such that $w\in g$ there exist $x,y$ such that $w = (x,y)$.
                            \begin{subproof}
                                Fix $w$.
                                Assume $w\in g$.
                                Take $n\in k$ such that $w = (n, X\setminus f(n))$.
                            \end{subproof}
                        \end{subproof}
                        Show that $g$ is right-unique.
                        \begin{subproof}
                            Show that for all $n,y,y'$ such that $(n,y)\in g$ and $(n,y')\in g$ we have $y = y'$.
                            \begin{subproof}
                                Fix $n$.
                                Fix $y$.
                                Fix $y'$.
                                Assume $(n,y)\in g$ and $(n,y')\in g$.
                                Take $n_1\in k$ such that $(n,y) = (n_1, X\setminus f(n_1))$.
                                Take $n_2\in k$ such that $(n,y') = (n_2, X\setminus f(n_2))$.
                                Thus $n = n_1$ and $y = X\setminus f(n_1)$ by \cref{pair_eq_iff}.
                                Thus $n = n_2$ and $y' = X\setminus f(n_2)$ by \cref{pair_eq_iff}.
                                Thus $y = y'$.
                            \end{subproof}
                        \end{subproof}
                    \end{subproof}
                    Thus $g$ is a function.
                    Show that $k\subseteq\dom{g}$.
                    \begin{subproof}
                        Show that for all $n$ such that $n\in k$ we have $n\in\dom{g}$.
                        \begin{subproof}
                            Fix $n$.
                            Assume $n\in k$.
                            Thus $(n, X\setminus f(n))\in g$.
                            Thus $n\in\dom{g}$ by \cref{dom_intro}.
                        \end{subproof}
                        Thus $k\subseteq\dom{g}$ by \cref{subseteq}.
                    \end{subproof}
                    Show that $\dom{g}\subseteq k$.
                    \begin{subproof}
                        Show that for all $n$ such that $n\in\dom{g}$ we have $n\in k$.
                        \begin{subproof}
                            Fix $n$.
                            Assume $n\in\dom{g}$.
                            Take $y$ such that $(n,y)\in g$ by \cref{dom_iff}.
                            Take $m\in k$ such that $(n,y) = (m, X\setminus f(m))$.
                            Thus $n = m$ by \cref{pair_eq_iff}.
                            Thus $n\in k$.
                        \end{subproof}
                        Thus $\dom{g}\subseteq k$ by \cref{subseteq}.
                    \end{subproof}
                    Thus $\dom{g} = k$ by \cref{subseteq_antisymmetric}.
                    Show that for all $n$ such that $n\in\dom{g}$ we have $g(n)\in G$.
                    \begin{subproof}
                        Fix $n$.
                        Assume $n\in\dom{g}$.
                        Then $n\in k$.
                        Thus $f\in\bijections{k}{F}$ by \cref{bijections}.
                        Thus $f\in\funs{k}{F}$ by \cref{bijections_to_funs}.
                        Thus $f$ is a function to $F$ and $\dom{f} = k$ by \cref{funs_elim}.
                        Thus $n\in\dom{f}$.
                        Thus $f(n)\in F$.
                        $f(n)$ is closed in $X$ with respect to $T$.
                        Thus $X\setminus f(n)$ is open in $X$ with respect to $T$ by \cref{closed_in}.
                        Thus $X\setminus f(n)\in T$ by \cref{open_in}.
                        Thus $X\setminus f(n)\in G$.
                        Thus $(n, X\setminus f(n))\in g$.
                        Thus $g(n) = X\setminus f(n)$ by \cref{function_apply_iff}.
                        Thus $g(n)\in G$.
                    \end{subproof}
                    Thus $g$ is a function to $G$.
                \end{subproof}
                Thus $g\in\funs{k}{G}$ by \cref{funs_intro}.
                Show that $g\in\surj{k}{G}$.
                \begin{subproof}
                    Show that for all $U$ such that $U\in G$ there exists $n\in k$ such that $g(n) = U$.
                    \begin{subproof}
                        Fix $U$.
                        Assume $U\in G$.
                        Take $A\in F$ such that $U = X\setminus A$.
                        $f\in\surj{k}{F}$ by \cref{bijections}.
                        Thus there exists $n\in k$ such that $f(n) = A$ by \cref{surj}.
                        Thus $(n, X\setminus f(n))\in g$.
                        Thus $g(n) = X\setminus f(n)$ by \cref{function_apply_iff}.
                        Thus $g(n) = U$.
                    \end{subproof}
                    Thus $g\in\surj{k}{G}$ by \cref{surj}.
                \end{subproof}
                Show that $g$ is injective.
                \begin{subproof}
                    Show that for all $n,m$ such that $n\in k$ and $m\in k$ and $g(n) = g(m)$ we have $n = m$.
                    \begin{subproof}
                        Fix $n$.
                        Fix $m$.
                        Assume $n\in k$ and $m\in k$ and $g(n) = g(m)$.
                        Thus $(n, X\setminus f(n))\in g$ and $(m, X\setminus f(m))\in g$.
                        Thus $g(n) = X\setminus f(n)$ by \cref{function_apply_iff}.
                        Thus $g(m) = X\setminus f(m)$ by \cref{function_apply_iff}.
                        Thus $X\setminus f(n) = X\setminus f(m)$.
                        Thus $f\in\bijections{k}{F}$ by \cref{bijections}.
                        Thus $f\in\funs{k}{F}$ by \cref{bijections_to_funs}.
                        Thus $f$ is a function to $F$ and $\dom{f} = k$ by \cref{funs_elim}.
                        Thus $n\in\dom{f}$ and $m\in\dom{f}$.
                        Thus $f(n)\in F$ and $f(m)\in F$.
                        Thus $f(n)\in\pow{X}$ and $f(m)\in\pow{X}$ by \cref{subseteq}.
                        Thus $f(n)\subseteq X$ and $f(m)\subseteq X$ by \cref{pow_iff}.
                        Thus $f(n) = f(m)$ by \cref{complement_eq_implies_eq}.
                        $f$ is injective by \cref{bijections}.
                        Thus $n = m$ by \cref{injective_function}.
                    \end{subproof}
                    Thus $g$ is injective by \cref{injective_function}.
                \end{subproof}
                Thus $g$ is a bijection from $k$ to $G$ by \cref{bijections}.
            \end{subproof}
            Thus $G$ is finite by \cref{finite}.
        \end{subproof}
        Thus $\inters{G}\in T$ by \cref{finite_intersections_open_in_topology}.
        Thus $X\setminus\unions{F}$ is open in $X$ with respect to $T$ by \cref{open_in}.
        Show that $\unions{F}\subseteq X$.
        \begin{subproof}
            $\unions{F}\subseteq X$ by \cref{unions_subseteq_of_powerset_is_subseteq}.
        \end{subproof}
        Thus $\unions{F}$ is closed in $X$ with respect to $T$ by \cref{closed_in}.
    \end{byCase}
\end{proof}

% from §17 Item 6: closed sets in subspaces
\begin{theorem}\label{closed_in_subspace_iff}
    Assume $Y$ is a subspace of $X$ with respect to $T$.
    Then for all $A$ we have
    $A$ is closed in $Y$ with respect to $\subspaceTopology{T}{Y}$
    iff there exists $C$ such that $C$ is closed in $X$ with respect to $T$ and $A = C\inter Y$.
\end{theorem}
\begin{proof}
    $T$ is a topology on $X$ by \cref{subspace_of}.
    $Y\subseteq X$ by \cref{subspace_of}.
    Let $T_Y = \subspaceTopology{T}{Y}$.
    Fix $A$.
    Show that if $A$ is closed in $Y$ with respect to $T_Y$ then
        there exists $C$ such that $C$ is closed in $X$ with respect to $T$ and $A = C\inter Y$.
    \begin{subproof}
        Assume $A$ is closed in $Y$ with respect to $T_Y$.
        Thus $Y\setminus A$ is open in $Y$ with respect to $T_Y$ by \cref{closed_in}.
        Thus $Y\setminus A\in T_Y$ by \cref{open_in}.
        Take $U\in T$ such that $Y\setminus A = Y\inter U$ by \cref{subspace_topology}.
        Let $C = X\setminus U$.
        $C\subseteq X$ by \cref{setminus_subseteq}.
        $U$ is open in $X$ with respect to $T$ by \cref{open_in}.
        $U\in T$ by \cref{open_in}.
        $T\subseteq\pow{X}$ by \cref{topology_on}.
        Thus $U\in\pow{X}$ by \cref{subseteq}.
        Thus $U\subseteq X$ by \cref{pow_iff}.
        $X\setminus C = X\inter U$ by \cref{setminus_setminus}.
        Thus $X\setminus C = U\inter X$ by \cref{inter_comm}.
        Thus $X\setminus C = U$ by \cref{inter_absorb_supseteq_right}.
        Thus $C$ is closed in $X$ with respect to $T$ by \cref{closed_in}.
        Show that $A = C\inter Y$.
        \begin{subproof}
            Show that $A\subseteq Y\setminus U$.
            \begin{subproof}
                Show that for all $x$ such that $x\in A$ we have $x\in Y\setminus U$.
                \begin{subproof}
                    Fix $x$.
                    Assume $x\in A$.
                    $A\subseteq Y$ by \cref{closed_in}.
                    Thus $x\in Y$ by \cref{elem_subseteq}.
                    Suppose not.
                    Then $x\in U$.
                    Thus $x\in Y\inter U$ by \cref{inter_intro}.
                    Thus $x\in Y\setminus A$.
                    Thus $x\notin A$ by \cref{setminus_elim_right}.
                    Contradiction.
                \end{subproof}
                Thus $A\subseteq Y\setminus U$ by \cref{subseteq}.
            \end{subproof}
            Show that $Y\setminus U\subseteq A$.
            \begin{subproof}
                Show that for all $x$ such that $x\in Y\setminus U$ we have $x\in A$.
                \begin{subproof}
                    Fix $x$.
                    Assume $x\in Y\setminus U$.
                    Suppose not.
                    Then $x\notin A$.
                    Thus $x\in Y$ by \cref{setminus_elim_left}.
                    Thus $x\in Y\setminus A$ by \cref{setminus_intro}.
                    Thus $x\in Y\inter U$.
                    Thus $x\in U$ by \cref{inter}.
                    Thus $x\notin U$ by \cref{setminus_elim_right}.
                    Contradiction.
                \end{subproof}
                Thus $Y\setminus U\subseteq A$ by \cref{subseteq}.
            \end{subproof}
            Thus $A = Y\setminus U$ by \cref{subseteq_antisymmetric}.
            Thus $A = Y\inter (X\setminus U)$ by \cref{setminus_eq_inter_complement}.
            Thus $A = C\inter Y$ by \cref{inter_comm}.
        \end{subproof}
    \end{subproof}
    Show that if there exists $C$ such that $C$ is closed in $X$ with respect to $T$ and $A = C\inter Y$ then
        $A$ is closed in $Y$ with respect to $T_Y$.
    \begin{subproof}
        Assume there exists $C$ such that $C$ is closed in $X$ with respect to $T$ and $A = C\inter Y$.
        Take $C$ such that $C$ is closed in $X$ with respect to $T$ and $A = C\inter Y$.
        Show that $Y\setminus A = Y\setminus C$.
        \begin{subproof}
            Show that $Y\setminus A\subseteq Y\setminus C$.
            \begin{subproof}
                Show that for all $x$ such that $x\in Y\setminus A$ we have $x\in Y\setminus C$.
                \begin{subproof}
                    Fix $x$.
                    Assume $x\in Y\setminus A$.
                    Suppose not.
                    Then $x\in C$.
                    Thus $x\in C\inter Y$ by \cref{inter_intro,setminus_elim_left}.
                    Thus $x\in A$.
                    Thus $x\notin A$ by \cref{setminus_elim_right}.
                    Contradiction.
                \end{subproof}
                Thus $Y\setminus A\subseteq Y\setminus C$ by \cref{subseteq}.
            \end{subproof}
            Show that $Y\setminus C\subseteq Y\setminus A$.
            \begin{subproof}
                Show that for all $x$ such that $x\in Y\setminus C$ we have $x\in Y\setminus A$.
                \begin{subproof}
                    Fix $x$.
                    Assume $x\in Y\setminus C$.
                    Suppose not.
                    Then $x\in A$.
                    Thus $x\in C$ by \cref{inter}.
                    Thus $x\notin C$ by \cref{setminus_elim_right}.
                    Contradiction.
                \end{subproof}
                Thus $Y\setminus C\subseteq Y\setminus A$ by \cref{subseteq}.
            \end{subproof}
            Thus $Y\setminus A = Y\setminus C$ by \cref{subseteq_antisymmetric}.
        \end{subproof}
        $C\subseteq X$ by \cref{closed_in}.
        Thus $Y\setminus C = Y\inter (X\setminus C)$ by \cref{setminus_eq_inter_complement}.
        $X\setminus C$ is open in $X$ with respect to $T$ by \cref{closed_in}.
        Thus $X\setminus C\in T$ by \cref{open_in}.
        Thus $Y\setminus C\in T_Y$ by \cref{subspace_topology}.
        Thus $Y\setminus A\in T_Y$.
        Thus $Y\setminus A$ is open in $Y$ with respect to $T_Y$ by \cref{open_in,subspace_topology_is_topology}.
        $A\subseteq Y$ by \cref{inter_lower_right}.
        Thus $A$ is closed in $Y$ with respect to $T_Y$ by \cref{closed_in}.
    \end{subproof}
\end{proof}

% from §17 Item 7: closed in a closed subspace
\begin{theorem}\label{closed_in_closed_subspace}
    Assume $T$ is a topology on $X$.
    Suppose $Y\subseteq X$.
    Suppose $Y$ is closed in $X$ with respect to $T$.
    Suppose $A$ is closed in $Y$ with respect to $\subspaceTopology{T}{Y}$.
    Then $A$ is closed in $X$ with respect to $T$.
\end{theorem}
\begin{proof}
    Take $C$ such that $C$ is closed in $X$ with respect to $T$ and $A = C\inter Y$
        by \cref{closed_in_subspace_iff,subspace_of}.
    Let $D = \{C, Y\}$.
    Show that $D\subseteq\pow{X}$.
    \begin{subproof}
        Show that for all $B$ such that $B\in D$ we have $B\in\pow{X}$.
        \begin{subproof}
            Fix $B$.
            Assume $B\in D$.
            Then $B = C$ or $B = Y$ by \cref{upair_iff}.
            Thus $B\subseteq X$ by \cref{closed_in}.
            Thus $B\in\pow{X}$ by \cref{pow_iff}.
        \end{subproof}
        Thus $D\subseteq\pow{X}$ by \cref{subseteq}.
    \end{subproof}
    Show that for all $B$ such that $B\in D$ we have $B$ is closed in $X$ with respect to $T$.
    \begin{subproof}
        Fix $B$.
        Assume $B\in D$.
        Then $B = C$ or $B = Y$ by \cref{upair_iff}.
        Thus $B$ is closed in $X$ with respect to $T$.
    \end{subproof}
    Thus $\inters{D}$ is closed in $X$ with respect to $T$ by \cref{closed_inters}.
    Thus $C\inter Y$ is closed in $X$ with respect to $T$ by \cref{inter_as_inters}.
    Thus $A$ is closed in $X$ with respect to $T$.
\end{proof}

% from §17 Item 8: interior of a set
\begin{definition}\label{interior}
    $\interior{A}{X}{T} = \{x\in X \mid \exists U\in T. x\in U \land U\subseteq A\}$.
\end{definition}

% from §17 Item 9: closure of a set
\begin{definition}\label{closure}
    $\closure{A}{X}{T} = \{x\in X \mid \text{for all $C$ such that $C$ is closed in $X$ with respect to $T$ and for all $y\in A$ we have $y\in C$ we have $x\in C$}\}$.
\end{definition}

% from §17 Item 10: interior is open
\begin{lemma}\label{interior_open_in}
    Assume $T$ is a topology on $X$.
    Then $\interior{A}{X}{T}$ is open in $X$ with respect to $T$.
\end{lemma}
\begin{proof}
    Let $S = \{U\in T \mid U\subseteq A\}$.
    Show that $\interior{A}{X}{T}\subseteq\unions{S}$.
    \begin{subproof}
        Show that for all $x$ such that $x\in\interior{A}{X}{T}$ we have $x\in\unions{S}$.
        \begin{subproof}
            Fix $x$.
            Assume $x\in\interior{A}{X}{T}$.
            Take $U\in T$ such that $x\in U$ and $U\subseteq A$ by \cref{interior}.
            Thus $U\in S$.
            Thus $x\in\unions{S}$ by \cref{unions_iff}.
        \end{subproof}
        Thus $\interior{A}{X}{T}\subseteq\unions{S}$ by \cref{subseteq}.
    \end{subproof}
    Show that $\unions{S}\subseteq\interior{A}{X}{T}$.
    \begin{subproof}
        Show that for all $x$ such that $x\in\unions{S}$ we have $x\in\interior{A}{X}{T}$.
        \begin{subproof}
            Fix $x$.
            Assume $x\in\unions{S}$.
            Take $U\in S$ such that $x\in U$ by \cref{unions_iff}.
            Thus $U\in T$ and $U\subseteq A$.
            $T\subseteq\pow{X}$ by \cref{topology_on}.
            Thus $U\in\pow{X}$ by \cref{subseteq}.
            Thus $U\subseteq X$ by \cref{pow_iff}.
            Thus $x\in X$ by \cref{elem_subseteq}.
            Thus $x\in\interior{A}{X}{T}$ by \cref{interior}.
        \end{subproof}
        Thus $\unions{S}\subseteq\interior{A}{X}{T}$ by \cref{subseteq}.
    \end{subproof}
    Thus $\interior{A}{X}{T} = \unions{S}$ by \cref{subseteq_antisymmetric}.
    Show that $S\subseteq T$.
    \begin{subproof}
        Show that for all $U$ such that $U\in S$ we have $U\in T$.
        \begin{subproof}
            Fix $U$.
            Assume $U\in S$.
            Thus $U\in T$.
        \end{subproof}
        Thus $S\subseteq T$ by \cref{subseteq}.
    \end{subproof}
    Thus $\unions{S}\in T$ by \cref{topology_on}.
    Thus $\interior{A}{X}{T}$ is open in $X$ with respect to $T$ by \cref{open_in}.
\end{proof}

% from §17 Item 11: closure is closed
\begin{lemma}\label{closure_closed_in}
    Assume $T$ is a topology on $X$.
    Suppose $A\subseteq X$.
    Then $\closure{A}{X}{T}$ is closed in $X$ with respect to $T$.
\end{lemma}
\begin{proof}
    Let $F = \{C\in\pow{X} \mid \text{$C$ is closed in $X$ with respect to $T$ and for all $x\in A$ we have $x\in C$}\}$.
    Show that $F$ is inhabited.
    \begin{subproof}
        $X$ is closed in $X$ with respect to $T$ by \cref{closed_whole_space}.
        $X\in\pow{X}$ by \cref{powerset_top}.
        Thus for all $x\in A$ we have $x\in X$ by \cref{subseteq}.
        Thus $X\in F$.
    \end{subproof}
    Show that $\closure{A}{X}{T} = \inters{F}$.
    \begin{subproof}
        Show that $\closure{A}{X}{T}\subseteq\inters{F}$.
        \begin{subproof}
            Show that for all $x$ such that $x\in\closure{A}{X}{T}$ we have $x\in\inters{F}$.
            \begin{subproof}
                Fix $x$.
                Assume $x\in\closure{A}{X}{T}$.
                Thus for all $C$ such that $C\in F$ we have $x\in C$ by \cref{closure}.
                Thus $x\in\inters{F}$ by \cref{inters_iff_forall}.
            \end{subproof}
            Thus $\closure{A}{X}{T}\subseteq\inters{F}$ by \cref{subseteq}.
        \end{subproof}
        Show that $\inters{F}\subseteq\closure{A}{X}{T}$.
        \begin{subproof}
            Show that for all $x$ such that $x\in\inters{F}$ we have $x\in\closure{A}{X}{T}$.
            \begin{subproof}
                Fix $x$.
                Assume $x\in\inters{F}$.
                Thus for all $C$ such that $C\in F$ we have $x\in C$ by \cref{inters_iff_forall}.
                Take $D$ such that $D\in F$.
                Thus $D\subseteq X$ by \cref{pow_iff}.
                Thus $x\in X$ by \cref{elem_subseteq}.
                Show that for all $C$ such that $C$ is closed in $X$ with respect to $T$ and
                    for all $y\in A$ we have $y\in C$ we have $x\in C$.
                \begin{subproof}
                    Fix $C$.
                    Assume $C$ is closed in $X$ with respect to $T$ and
                        for all $y\in A$ we have $y\in C$.
                    Thus $C\in F$.
                    Thus $x\in C$.
                \end{subproof}
                Thus $x\in\closure{A}{X}{T}$ by \cref{closure}.
            \end{subproof}
            Thus $\inters{F}\subseteq\closure{A}{X}{T}$ by \cref{subseteq}.
        \end{subproof}
        Thus $\closure{A}{X}{T} = \inters{F}$ by \cref{subseteq_antisymmetric}.
    \end{subproof}
    Show that $F\subseteq\pow{X}$.
    \begin{subproof}
        Show that for all $C$ such that $C\in F$ we have $C\in\pow{X}$.
        \begin{subproof}
            Fix $C$.
            Assume $C\in F$.
            Thus $C\in\pow{X}$.
        \end{subproof}
        Thus $F\subseteq\pow{X}$ by \cref{subseteq}.
    \end{subproof}
    Show that for all $C$ such that $C\in F$ we have $C$ is closed in $X$ with respect to $T$.
    \begin{subproof}
        Fix $C$.
        Assume $C\in F$.
        Thus $C$ is closed in $X$ with respect to $T$.
    \end{subproof}
    Thus $\inters{F}$ is closed in $X$ with respect to $T$ by \cref{closed_inters}.
    Thus $\closure{A}{X}{T}$ is closed in $X$ with respect to $T$.
\end{proof}

% from §17 Item 12: interior subset
\begin{lemma}\label{interior_subseteq}
    Assume $T$ is a topology on $X$.
    Then $\interior{A}{X}{T}\subseteq A$.
\end{lemma}
\begin{proof}
    Show that for all $x$ such that $x\in\interior{A}{X}{T}$ we have $x\in A$.
    \begin{subproof}
        Fix $x$.
        Assume $x\in\interior{A}{X}{T}$.
        Take $U\in T$ such that $x\in U$ and $U\subseteq A$ by \cref{interior}.
        Thus $x\in A$ by \cref{elem_subseteq}.
    \end{subproof}
    Thus $\interior{A}{X}{T}\subseteq A$ by \cref{subseteq}.
\end{proof}

% from §17 Item 13: subset of closure
\begin{lemma}\label{subseteq_closure}
    Assume $T$ is a topology on $X$.
    Suppose $A\subseteq X$.
    Then $A\subseteq\closure{A}{X}{T}$.
\end{lemma}
\begin{proof}
    Show that for all $x$ such that $x\in A$ we have $x\in\closure{A}{X}{T}$.
    \begin{subproof}
        Fix $x$.
        Assume $x\in A$.
        Thus $x\in X$ by \cref{elem_subseteq}.
        Show that for all $C$ such that $C$ is closed in $X$ with respect to $T$ and
            for all $y\in A$ we have $y\in C$ we have $x\in C$.
        \begin{subproof}
            Fix $C$.
            Assume $C$ is closed in $X$ with respect to $T$ and
                for all $y\in A$ we have $y\in C$.
            Thus $x\in C$.
        \end{subproof}
        Thus $x\in\closure{A}{X}{T}$ by \cref{closure}.
    \end{subproof}
    Thus $A\subseteq\closure{A}{X}{T}$ by \cref{subseteq}.
\end{proof}

% from §17 helper lemma 2: closure contained in closed superset
\begin{lemma}\label{closure_subseteq_closed}
    Assume $T$ is a topology on $X$.
    Suppose $A\subseteq X$.
    Suppose $C$ is closed in $X$ with respect to $T$.
    Suppose $A\subseteq C$.
    Then $\closure{A}{X}{T}\subseteq C$.
\end{lemma}
\begin{proof}
    Show that for all $x$ such that $x\in\closure{A}{X}{T}$ we have $x\in C$.
    \begin{subproof}
        Fix $x$.
        Assume $x\in\closure{A}{X}{T}$.
        Thus for all $B$ such that $B$ is closed in $X$ with respect to $T$ and
            for all $y\in A$ we have $y\in B$ we have $x\in B$ by \cref{closure}.
        Thus for all $y\in A$ we have $y\in C$ by \cref{subseteq}.
        Thus $x\in C$.
    \end{subproof}
    Thus $\closure{A}{X}{T}\subseteq C$ by \cref{subseteq}.
\end{proof}

% from §17 Item 14: open set equals its interior
\begin{lemma}\label{open_eq_interior}
    Assume $T$ is a topology on $X$.
    Suppose $A$ is open in $X$ with respect to $T$.
    Then $A = \interior{A}{X}{T}$.
\end{lemma}
\begin{proof}
    Show that $A\subseteq\interior{A}{X}{T}$.
    \begin{subproof}
        Show that for all $x$ such that $x\in A$ we have $x\in\interior{A}{X}{T}$.
        \begin{subproof}
            Fix $x$.
            Assume $x\in A$.
            $A\in T$ by \cref{open_in}.
            $A\subseteq A$ by \cref{subseteq_refl}.
            $T\subseteq\pow{X}$ by \cref{topology_on}.
            Thus $A\in\pow{X}$ by \cref{subseteq}.
            Thus $A\subseteq X$ by \cref{pow_iff}.
            Thus $x\in X$ by \cref{elem_subseteq}.
            Thus $x\in\interior{A}{X}{T}$ by \cref{interior}.
        \end{subproof}
        Thus $A\subseteq\interior{A}{X}{T}$ by \cref{subseteq}.
    \end{subproof}
    $\interior{A}{X}{T}\subseteq A$ by \cref{interior_subseteq}.
    Thus $A = \interior{A}{X}{T}$ by \cref{subseteq_antisymmetric}.
\end{proof}

% from §17 Item 15: closed set equals its closure
\begin{lemma}\label{closed_eq_closure}
    Assume $T$ is a topology on $X$.
    Suppose $A$ is closed in $X$ with respect to $T$.
    Then $A = \closure{A}{X}{T}$.
\end{lemma}
\begin{proof}
    $A\subseteq X$ by \cref{closed_in}.
    Thus $A\subseteq\closure{A}{X}{T}$ by \cref{subseteq_closure}.
    $A\subseteq A$ by \cref{subseteq_refl}.
    Thus $\closure{A}{X}{T}\subseteq A$ by \cref{closure_subseteq_closed}.
    Thus $A = \closure{A}{X}{T}$ by \cref{subseteq_antisymmetric}.
\end{proof}

% from §17 Item 16: closure in subspace
\begin{theorem}\label{closure_in_subspace}
    Assume $Y$ is a subspace of $X$ with respect to $T$.
    Suppose $A\subseteq Y$.
    Let $T_Y = \subspaceTopology{T}{Y}$.
    Then $\closure{A}{Y}{T_Y} = \closure{A}{X}{T}\inter Y$.
\end{theorem}
\begin{proof}
    $T$ is a topology on $X$ by \cref{subspace_of}.
    $Y\subseteq X$ by \cref{subspace_of}.
    Thus $A\subseteq X$ by \cref{subseteq_transitive}.
    Let $B = \closure{A}{Y}{T_Y}$.
    Let $C = \closure{A}{X}{T}$.
    Show that $B\subseteq C\inter Y$.
    \begin{subproof}
        $C$ is closed in $X$ with respect to $T$ by \cref{closure_closed_in}.
        Thus $C\inter Y$ is closed in $Y$ with respect to $T_Y$
            by \cref{closed_in_subspace_iff,subspace_of}.
        $A\subseteq C$ by \cref{subseteq_closure}.
        Thus $A\subseteq C\inter Y$ by \cref{subseteq_inter_iff}.
        $T_Y$ is a topology on $Y$ by \cref{subspace_topology_is_topology}.
        Thus $B\subseteq C\inter Y$ by \cref{closure_subseteq_closed}.
    \end{subproof}
    Show that $C\inter Y\subseteq B$.
    \begin{subproof}
        $T_Y$ is a topology on $Y$ by \cref{subspace_topology_is_topology}.
        $B$ is closed in $Y$ with respect to $T_Y$ by \cref{closure_closed_in}.
        Take $D$ such that $D$ is closed in $X$ with respect to $T$ and $B = D\inter Y$
            by \cref{closed_in_subspace_iff,subspace_of}.
        $A\subseteq B$ by \cref{subseteq_closure}.
        Thus $B\subseteq D$ by \cref{inter_lower_left}.
        Thus $A\subseteq D$ by \cref{subseteq_transitive}.
        Thus $C\subseteq D$ by \cref{closure_subseteq_closed}.
        Show that $C\inter Y\subseteq D\inter Y$.
        \begin{subproof}
            Show that for all $x$ such that $x\in C\inter Y$ we have $x\in D\inter Y$.
            \begin{subproof}
                Fix $x$.
                Assume $x\in C\inter Y$.
                Thus $x\in C$ and $x\in Y$ by \cref{inter}.
                Thus $x\in D$ by \cref{elem_subseteq}.
                Thus $x\in D\inter Y$ by \cref{inter_intro}.
            \end{subproof}
            Thus $C\inter Y\subseteq D\inter Y$ by \cref{subseteq}.
        \end{subproof}
        Thus $C\inter Y\subseteq B$ by \cref{subseteq}.
    \end{subproof}
    Thus $B = C\inter Y$ by \cref{subseteq_antisymmetric}.
\end{proof}

% from §17 Item 17: intersects
\begin{definition}\label{intersects}
    $A$ intersects $B$ iff $A\inter B \neq \emptyset$.
\end{definition}

% from §17 Item 18: closure and open sets
\begin{theorem}\label{closure_iff_intersects_open}
    Assume $T$ is a topology on $X$.
    Suppose $A\subseteq X$.
    Suppose $x\in X$.
    Then $x\in\closure{A}{X}{T}$ iff for all $U$ such that
        $U$ is open in $X$ with respect to $T$ and $x\in U$ we have $U$ intersects $A$.
\end{theorem}
\begin{proof}
    Show that if $x\in\closure{A}{X}{T}$ then for all $U$ such that
        $U$ is open in $X$ with respect to $T$ and $x\in U$ we have $U$ intersects $A$.
    \begin{subproof}
        Assume $x\in\closure{A}{X}{T}$.
        Fix $U$.
        Assume $U$ is open in $X$ with respect to $T$ and $x\in U$.
        Suppose not.
        Then $U\inter A = \emptyset$ by \cref{intersects}.
        Thus $A\inter U = \emptyset$ by \cref{inter_comm}.
        Thus $A\subseteq X\setminus U$ by \cref{subseteq_setminus}.
        Show that $X\setminus U$ is closed in $X$ with respect to $T$.
        \begin{subproof}
            $X\setminus U\subseteq X$ by \cref{setminus_subseteq}.
            $T\subseteq\pow{X}$ by \cref{topology_on}.
            Thus $U\in T$ by \cref{open_in}.
            Thus $U\in\pow{X}$ by \cref{subseteq}.
            Thus $U\subseteq X$ by \cref{pow_iff}.
            $X\setminus (X\setminus U) = X\inter U$ by \cref{setminus_setminus}.
            Thus $X\setminus (X\setminus U) = U\inter X$ by \cref{inter_comm}.
            Thus $X\setminus (X\setminus U) = U$ by \cref{inter_absorb_supseteq_right}.
            Thus $X\setminus U$ is closed in $X$ with respect to $T$ by \cref{closed_in}.
        \end{subproof}
        Thus $\closure{A}{X}{T}\subseteq X\setminus U$ by \cref{closure_subseteq_closed}.
        Thus $x\in X\setminus U$ by \cref{elem_subseteq}.
        Thus $x\notin U$ by \cref{setminus_elim_right}.
        Contradiction.
    \end{subproof}
    Show that if for all $U$ such that $U$ is open in $X$ with respect to $T$ and $x\in U$ we have $U$ intersects $A$
        then $x\in\closure{A}{X}{T}$.
    \begin{subproof}
        Assume for all $U$ such that $U$ is open in $X$ with respect to $T$ and $x\in U$ we have $U$ intersects $A$.
        Show that for all $C$ such that $C$ is closed in $X$ with respect to $T$ and
            for all $y\in A$ we have $y\in C$ we have $x\in C$.
        \begin{subproof}
            Fix $C$.
            Assume $C$ is closed in $X$ with respect to $T$ and
                for all $y\in A$ we have $y\in C$.
            Suppose not.
            Then $x\notin C$.
            Let $U = X\setminus C$.
            $U$ is open in $X$ with respect to $T$ by \cref{closed_in}.
            Thus $x\in U$ by \cref{setminus_intro}.
            Show that $U$ does not intersect $A$.
            \begin{subproof}
                Suppose not.
                Then $U\inter A \neq \emptyset$ by \cref{intersects}.
                Take $y$ such that $y\in U\inter A$ by \cref{nonempty_inhabited}.
                Thus $y\in A$ and $y\in U$ by \cref{inter}.
                Thus $y\in C$.
                Thus $y\notin C$ by \cref{setminus_elim_right}.
                Contradiction.
            \end{subproof}
            Thus $U\inter A = \emptyset$.
            Thus $U$ does not intersect $A$ by \cref{intersects}.
            Contradiction.
        \end{subproof}
        Thus $x\in\closure{A}{X}{T}$ by \cref{closure}.
    \end{subproof}
\end{proof}

% from §17 Item 19: closure and basis elements
\begin{theorem}\label{closure_iff_intersects_basis}
    Assume $T$ is a topology on $X$.
    Suppose $B$ is a basis for $T$ on $X$.
    Suppose $A\subseteq X$.
    Suppose $x\in X$.
    Then $x\in\closure{A}{X}{T}$ iff for all $B_1$ such that $B_1\in B$ and $x\in B_1$ we have $B_1$ intersects $A$.
\end{theorem}
\begin{proof}
    Show that if $x\in\closure{A}{X}{T}$ then for all $B_1$ such that $B_1\in B$ and $x\in B_1$ we have $B_1$ intersects $A$.
    \begin{subproof}
        Assume $x\in\closure{A}{X}{T}$.
        Fix $B_1$.
        Assume $B_1\in B$ and $x\in B_1$.
        $B$ is a basis on $X$ by \cref{basis_for_topology}.
        Thus $B_1\in\basisTopology{B}{X}$ by \cref{basis_elem_open_in_generated}.
        Thus $T = \basisTopology{B}{X}$ by \cref{basis_for_topology}.
        Thus $B_1\in T$.
        Thus $B_1$ is open in $X$ with respect to $T$ by \cref{open_in}.
        Thus $B_1$ intersects $A$ by \cref{closure_iff_intersects_open}.
    \end{subproof}
    Show that if for all $B_1$ such that $B_1\in B$ and $x\in B_1$ we have $B_1$ intersects $A$ then
        $x\in\closure{A}{X}{T}$.
    \begin{subproof}
        Assume for all $B_1$ such that $B_1\in B$ and $x\in B_1$ we have $B_1$ intersects $A$.
        Show that for all $U$ such that $U$ is open in $X$ with respect to $T$ and $x\in U$ we have $U$ intersects $A$.
        \begin{subproof}
            Fix $U$.
            Assume $U$ is open in $X$ with respect to $T$ and $x\in U$.
            Thus $U\in T$ by \cref{open_in}.
            Thus $T = \basisTopology{B}{X}$ by \cref{basis_for_topology}.
            Take $B_1\in B$ such that $x\in B_1$ and $B_1\subseteq U$ by \cref{topology_generated_by_basis}.
            Thus $B_1$ intersects $A$.
            Take $y$ such that $y\in B_1\inter A$ by \cref{intersects,nonempty_inhabited}.
            Thus $y\in U\inter A$ by \cref{inter,elem_subseteq,inter_intro}.
            Show that $U\inter A \neq \emptyset$.
            \begin{subproof}
                Suppose not.
                Then $U\inter A = \emptyset$.
                Thus $y\in\emptyset$.
                Contradiction by \cref{emptyset}.
            \end{subproof}
            Thus $U$ intersects $A$ by \cref{intersects}.
        \end{subproof}
        Thus $x\in\closure{A}{X}{T}$ by \cref{closure_iff_intersects_open}.
    \end{subproof}
\end{proof}

% from §17 Item 20: neighborhood
\begin{definition}\label{neighborhood}
    $U$ is a neighborhood of $x$ in $X$ with respect to $T$ iff
        $U$ is open in $X$ with respect to $T$ and $x\in U$.
\end{definition}

% from §17 Item 21: limit point
\begin{definition}\label{limit_point}
    $x$ is a limit point of $A$ in $X$ with respect to $T$ iff
        for all $U$ such that $U$ is a neighborhood of $x$ in $X$ with respect to $T$
        we have $U$ intersects $A\setminus\{x\}$.
\end{definition}

% from §17 Item 22: set of limit points
\begin{definition}\label{limit_points}
    $\limitPoints{A}{X}{T} = \{x\in X \mid \text{$x$ is a limit point of $A$ in $X$ with respect to $T$}\}$.
\end{definition}

% from §17 Item 23: closure equals set plus limit points
\begin{theorem}\label{closure_union_limit_points}
    Assume $T$ is a topology on $X$.
    Suppose $A\subseteq X$.
    Then $\closure{A}{X}{T} = A\union\limitPoints{A}{X}{T}$.
\end{theorem}
\begin{proof}
    Show that $A\union\limitPoints{A}{X}{T}\subseteq\closure{A}{X}{T}$.
    \begin{subproof}
        Show that for all $x$ such that $x\in A\union\limitPoints{A}{X}{T}$ we have $x\in\closure{A}{X}{T}$.
        \begin{subproof}
            Fix $x$.
            Assume $x\in A\union\limitPoints{A}{X}{T}$.
            Then $x\in A$ or $x\in\limitPoints{A}{X}{T}$ by \cref{union_iff}.
            \begin{byCase}
            \caseOf{$x\in A$.}
                $A\subseteq\closure{A}{X}{T}$ by \cref{subseteq_closure}.
                Thus $x\in\closure{A}{X}{T}$ by \cref{elem_subseteq}.
            \caseOf{$x\in\limitPoints{A}{X}{T}$.}
                Show that for all $U$ such that $U$ is open in $X$ with respect to $T$ and $x\in U$ we have $U$ intersects $A$.
                \begin{subproof}
                    Fix $U$.
                    Assume $U$ is open in $X$ with respect to $T$ and $x\in U$.
                    Thus $U$ is a neighborhood of $x$ in $X$ with respect to $T$ by \cref{neighborhood}.
                    Thus $x$ is a limit point of $A$ in $X$ with respect to $T$ by \cref{limit_points}.
                    Thus $U$ intersects $A\setminus\{x\}$ by \cref{limit_point}.
                    Take $y$ such that $y\in U\inter (A\setminus\{x\})$ by \cref{intersects,nonempty_inhabited}.
                    Thus $y\in U$ and $y\in A\setminus\{x\}$ by \cref{inter}.
                    Thus $y\in A$ by \cref{setminus_elim_left}.
                    Thus $y\in U\inter A$ by \cref{inter_intro}.
                    Show that $U\inter A \neq \emptyset$.
                    \begin{subproof}
                        Suppose not.
                        Then $U\inter A = \emptyset$.
                        Thus $y\in\emptyset$.
                        Contradiction by \cref{emptyset}.
                    \end{subproof}
                    Thus $U$ intersects $A$ by \cref{intersects}.
                \end{subproof}
                Thus $x\in X$ by \cref{limit_points}.
                Thus $x\in\closure{A}{X}{T}$ by \cref{closure_iff_intersects_open}.
            \end{byCase}
        \end{subproof}
        Thus $A\union\limitPoints{A}{X}{T}\subseteq\closure{A}{X}{T}$ by \cref{subseteq}.
    \end{subproof}
    Show that $\closure{A}{X}{T}\subseteq A\union\limitPoints{A}{X}{T}$.
    \begin{subproof}
        Show that for all $x$ such that $x\in\closure{A}{X}{T}$ we have $x\in A\union\limitPoints{A}{X}{T}$.
        \begin{subproof}
            Fix $x$.
            Assume $x\in\closure{A}{X}{T}$.
            \begin{byCase}
            \caseOf{$x\in A$.}
                Thus $x\in A\union\limitPoints{A}{X}{T}$ by \cref{union_intro_left}.
            \caseOf{$x\notin A$.}
                Show that $x\in\limitPoints{A}{X}{T}$.
                \begin{subproof}
                    Show that for all $U$ such that $U$ is a neighborhood of $x$ in $X$ with respect to $T$ we have
                        $U$ intersects $A\setminus\{x\}$.
                    \begin{subproof}
                    Fix $U$.
                    Assume $U$ is a neighborhood of $x$ in $X$ with respect to $T$.
                    Thus $U$ is open in $X$ with respect to $T$ and $x\in U$ by \cref{neighborhood}.
                    Thus $x\in X$ by \cref{closure}.
                    Thus $U$ intersects $A$ by \cref{closure_iff_intersects_open}.
                        Take $y$ such that $y\in U\inter A$ by \cref{intersects,nonempty_inhabited}.
                        Thus $y\in U$ and $y\in A$ by \cref{inter}.
                        Show that $y\neq x$.
                        \begin{subproof}
                            Suppose not.
                            Then $y = x$.
                            Thus $x\in A$.
                            Contradiction.
                        \end{subproof}
                        Thus $y\notin\{x\}$ by \cref{singleton_iff}.
                        Thus $y\in A\setminus\{x\}$ by \cref{setminus_intro}.
                        Thus $y\in U\inter (A\setminus\{x\})$ by \cref{inter_intro}.
                        Show that $U\inter (A\setminus\{x\}) \neq \emptyset$.
                        \begin{subproof}
                            Suppose not.
                            Then $U\inter (A\setminus\{x\}) = \emptyset$.
                            Thus $y\in\emptyset$.
                            Contradiction by \cref{emptyset}.
                        \end{subproof}
                        Thus $U$ intersects $A\setminus\{x\}$ by \cref{intersects}.
                    \end{subproof}
                    Thus $x\in X$ by \cref{closure}.
                    Thus $x\in\limitPoints{A}{X}{T}$ by \cref{limit_point,limit_points}.
                \end{subproof}
                Thus $x\in A\union\limitPoints{A}{X}{T}$ by \cref{union_intro_right}.
            \end{byCase}
        \end{subproof}
        Thus $\closure{A}{X}{T}\subseteq A\union\limitPoints{A}{X}{T}$ by \cref{subseteq}.
    \end{subproof}
    Thus $\closure{A}{X}{T} = A\union\limitPoints{A}{X}{T}$ by \cref{subseteq_antisymmetric}.
\end{proof}

% from §17 Item 24: closed iff contains limit points
\begin{corollary}\label{closed_iff_contains_limit_points}
    Assume $T$ is a topology on $X$.
    Suppose $A\subseteq X$.
    Then $A$ is closed in $X$ with respect to $T$ iff $\limitPoints{A}{X}{T}\subseteq A$.
\end{corollary}
\begin{proof}
    Show that if $A$ is closed in $X$ with respect to $T$ then $\limitPoints{A}{X}{T}\subseteq A$.
    \begin{subproof}
        Assume $A$ is closed in $X$ with respect to $T$.
        Thus $A = \closure{A}{X}{T}$ by \cref{closed_eq_closure}.
        Thus $A = A\union\limitPoints{A}{X}{T}$ by \cref{closure_union_limit_points}.
        Thus $A\union\limitPoints{A}{X}{T} = A$.
        Thus $\limitPoints{A}{X}{T}\union A = A$ by \cref{union_comm}.
        Thus $\limitPoints{A}{X}{T}\subseteq A$ by \cref{union_eq_self_implies_subseteq}.
    \end{subproof}
    Show that if $\limitPoints{A}{X}{T}\subseteq A$ then $A$ is closed in $X$ with respect to $T$.
    \begin{subproof}
        Assume $\limitPoints{A}{X}{T}\subseteq A$.
        Thus $A\union\limitPoints{A}{X}{T} = A$ by \cref{union_absorb_subseteq_right}.
        Thus $\closure{A}{X}{T} = A$ by \cref{closure_union_limit_points}.
        $\closure{A}{X}{T}$ is closed in $X$ with respect to $T$ by \cref{closure_closed_in}.
        Thus $A$ is closed in $X$ with respect to $T$.
    \end{subproof}
\end{proof}

% from §17 Item 25: Hausdorff space
\begin{definition}\label{hausdorff}
    $X$ is Hausdorff with respect to $T$ iff
        $T$ is a topology on $X$ and
        for all $x_1,x_2$ such that $x_1\in X$ and $x_2\in X$ and $x_1\neq x_2$
        there exist $U_1,U_2$ such that
            $U_1$ is a neighborhood of $x_1$ in $X$ with respect to $T$ and
            $U_2$ is a neighborhood of $x_2$ in $X$ with respect to $T$ and
            $U_1$ is disjoint from $U_2$.
\end{definition}

% from §17 helper lemma 3: finite preserved by bijection
\begin{lemma}\label{finite_bijection}
    Suppose $A$ is finite.
    Suppose $f$ is a bijection from $A$ to $B$.
    Then $B$ is finite.
\end{lemma}
\begin{proof}
    Take $k$ such that $k$ is a natural number and there exists a bijection from $k$ to $A$ by \cref{finite}.
    Take $g$ such that $g$ is a bijection from $k$ to $A$.
    Thus $f\circ g$ is a bijection from $k$ to $B$ by \cref{bijection_circ}.
    Thus $B$ is finite by \cref{finite}.
\end{proof}

% from §17 helper lemma 4: bijection to singleton family
\begin{lemma}\label{singleton_family_bijection}
    Let $S = \{\{x\} \mid x\in A\}$.
    Then there exists a bijection from $A$ to $S$.
\end{lemma}
\begin{proof}
    Let $f = \{(x,\{x\}) \mid x\in A\}$.
    Show that $f$ is a bijection from $A$ to $S$.
    \begin{subproof}
        Show that $f$ is a function.
        \begin{subproof}
            Show that $f$ is a relation.
            \begin{subproof}
                Show that for all $w$ such that $w\in f$ there exist $x,y$ such that $w = (x,y)$.
                \begin{subproof}
                    Fix $w$.
                    Assume $w\in f$.
                    Take $x\in A$ such that $w = (x,\{x\})$.
                    Thus there exists $y$ such that $w = (x,y)$.
                \end{subproof}
                Thus $f$ is a relation by \cref{relation}.
            \end{subproof}
            Show that $f$ is right-unique.
            \begin{subproof}
                Show that for all $x,y,z$ such that $(x,y)\in f$ and $(x,z)\in f$ we have $y = z$.
                \begin{subproof}
                    Fix $x$.
                    Fix $y$.
                    Fix $z$.
                    Assume $(x,y)\in f$ and $(x,z)\in f$.
                    Take $x_1\in A$ such that $(x,y) = (x_1,\{x_1\})$.
                    Take $x_2\in A$ such that $(x,z) = (x_2,\{x_2\})$.
                    Thus $x = x_1$ and $y = \{x_1\}$ by \cref{pair_eq_iff}.
                    Thus $x = x_2$ and $z = \{x_2\}$ by \cref{pair_eq_iff}.
                    Thus $y = z$.
                \end{subproof}
                Thus $f$ is right-unique by \cref{rightunique}.
            \end{subproof}
            Thus $f$ is a function.
        \end{subproof}
        Show that $\dom{f} = A$.
        \begin{subproof}
            Show that $A\subseteq\dom{f}$.
            \begin{subproof}
                Show that for all $x$ such that $x\in A$ we have $x\in\dom{f}$.
                \begin{subproof}
                    Fix $x$.
                    Assume $x\in A$.
                    Thus $(x,\{x\})\in f$.
                    Thus $x\in\dom{f}$ by \cref{dom_iff}.
                \end{subproof}
                Thus $A\subseteq\dom{f}$ by \cref{subseteq}.
            \end{subproof}
            Show that $\dom{f}\subseteq A$.
            \begin{subproof}
                Show that for all $x$ such that $x\in\dom{f}$ we have $x\in A$.
                \begin{subproof}
                    Fix $x$.
                    Assume $x\in\dom{f}$.
                    Take $y$ such that $(x,y)\in f$ by \cref{dom_iff}.
                    Take $x_1\in A$ such that $(x,y) = (x_1,\{x_1\})$.
                    Thus $x = x_1$ by \cref{pair_eq_iff}.
                    Thus $x\in A$.
                \end{subproof}
                Thus $\dom{f}\subseteq A$ by \cref{subseteq}.
            \end{subproof}
            Thus $\dom{f} = A$ by \cref{subseteq_antisymmetric}.
        \end{subproof}
        Show that $f$ is a function to $S$.
        \begin{subproof}
            Show that for all $x$ such that $x\in\dom{f}$ we have $f(x)\in S$.
            \begin{subproof}
                Fix $x$.
                Assume $x\in\dom{f}$.
                Thus $x\in A$.
                Thus $(x,\{x\})\in f$.
                Thus $f(x) = \{x\}$ by \cref{function_apply_iff}.
                Thus $\{x\}\in S$.
                Thus $f(x)\in S$.
            \end{subproof}
            Thus $f$ is a function to $S$.
        \end{subproof}
        Thus $f\in\funs{A}{S}$ by \cref{funs_intro}.
        Show that $\ran{f} = S$.
        \begin{subproof}
            Show that $\ran{f}\subseteq S$.
            \begin{subproof}
                Show that for all $y$ such that $y\in\ran{f}$ we have $y\in S$.
                \begin{subproof}
                    Fix $y$.
                    Assume $y\in\ran{f}$.
                    Take $x$ such that $(x,y)\in f$ by \cref{ran_iff}.
                    Take $x_1\in A$ such that $(x,y) = (x_1,\{x_1\})$.
                    Thus $y = \{x_1\}$ by \cref{pair_eq_iff}.
                    Thus $y\in S$.
                \end{subproof}
                Thus $\ran{f}\subseteq S$ by \cref{subseteq}.
            \end{subproof}
            Show that $S\subseteq\ran{f}$.
            \begin{subproof}
                Show that for all $y$ such that $y\in S$ we have $y\in\ran{f}$.
                \begin{subproof}
                    Fix $y$.
                    Assume $y\in S$.
                    Take $x\in A$ such that $y = \{x\}$.
                    Thus $(x,y)\in f$.
                    Thus $y\in\ran{f}$ by \cref{ran_iff}.
                \end{subproof}
                Thus $S\subseteq\ran{f}$ by \cref{subseteq}.
            \end{subproof}
            Thus $\ran{f} = S$ by \cref{subseteq_antisymmetric}.
        \end{subproof}
        Thus $f\in\surj{A}{S}$ by \cref{funs_surj_iff}.
        Show that $f$ is injective.
        \begin{subproof}
            Show that for all $x,x',y$ such that $(x,y)\in f$ and $(x',y)\in f$ we have $x = x'$.
            \begin{subproof}
                Fix $x$.
                Fix $x'$.
                Fix $y$.
                Assume $(x,y)\in f$ and $(x',y)\in f$.
                Take $x_1\in A$ such that $(x,y) = (x_1,\{x_1\})$.
                Take $x_2\in A$ such that $(x',y) = (x_2,\{x_2\})$.
                Thus $x = x_1$ and $y = \{x_1\}$ by \cref{pair_eq_iff}.
                Thus $x' = x_2$ and $y = \{x_2\}$ by \cref{pair_eq_iff}.
                Thus $\{x_1\} = \{x_2\}$.
                Thus $x_1 = x_2$ by \cref{singleton_elim,singleton_intro}.
                Thus $x = x'$.
            \end{subproof}
            Thus $f$ is injective by \cref{injective}.
        \end{subproof}
        Thus $f\in\bijections{A}{S}$ by \cref{bijections}.
    \end{subproof}
    Thus $f$ is a bijection from $A$ to $S$.
    Thus there exists a bijection from $A$ to $S$.
\end{proof}

% from §17 Item 26: finite point sets are closed in Hausdorff spaces
\begin{theorem}\label{finite_point_set_closed_hausdorff}
    Assume $X$ is Hausdorff with respect to $T$.
    Suppose $A\subseteq X$.
    Suppose $A$ is finite.
    Then $A$ is closed in $X$ with respect to $T$.
\end{theorem}
\begin{proof}
    $T$ is a topology on $X$ by \cref{hausdorff}.
    Show that for all $x$ such that $x\in X$ we have $\{x\}$ is closed in $X$ with respect to $T$.
    \begin{subproof}
        Fix $x$.
        Assume $x\in X$.
        Show that $\limitPoints{\{x\}}{X}{T}\subseteq \{x\}$.
        \begin{subproof}
            Show that for all $y$ such that $y\in\limitPoints{\{x\}}{X}{T}$ we have $y\in\{x\}$.
            \begin{subproof}
                Fix $y$.
                Assume $y\in\limitPoints{\{x\}}{X}{T}$.
                Thus $y\in X$ and $y$ is a limit point of $\{x\}$ in $X$ with respect to $T$ by \cref{limit_points}.
                Show that $y\in\{x\}$.
                \begin{subproof}
                    Suppose not.
                    Then $y\notin\{x\}$.
                    Thus $y\neq x$ by \cref{singleton_iff}.
                    Take $U_1,U_2$ such that
                        $U_1$ is a neighborhood of $y$ in $X$ with respect to $T$ and
                        $U_2$ is a neighborhood of $x$ in $X$ with respect to $T$ and
                        $U_1$ is disjoint from $U_2$ by \cref{hausdorff}.
                    $U_1$ is a neighborhood of $y$ in $X$ with respect to $T$.
                    Show that $U_1$ does not intersect $\{x\}\setminus\{y\}$.
                    \begin{subproof}
                        Suppose not.
                        Then $U_1$ intersects $\{x\}\setminus\{y\}$.
                        Take $z$ such that $z\in U_1\inter (\{x\}\setminus\{y\})$ by \cref{intersects,nonempty_inhabited}.
                        Thus $z\in U_1$ and $z\in \{x\}\setminus\{y\}$ by \cref{inter}.
                        Thus $z\in \{x\}$ by \cref{setminus_elim_left}.
                        Thus $z = x$ by \cref{singleton_elim}.
                        Thus $x\in U_1$.
                        $x\in U_2$ by \cref{neighborhood}.
                        Thus there exists $a$ such that $a\in U_1$ and $a\in U_2$.
                        Contradiction by \cref{disjoint}.
                    \end{subproof}
                    Thus $U_1$ does not intersect $\{x\}\setminus\{y\}$.
                    Thus $y$ is not a limit point of $\{x\}$ in $X$ with respect to $T$ by \cref{limit_point}.
                    Contradiction.
                \end{subproof}
            \end{subproof}
            Thus $\limitPoints{\{x\}}{X}{T}\subseteq \{x\}$ by \cref{subseteq}.
        \end{subproof}
        $\{x\}\subseteq X$ by \cref{singleton_subset_intro}.
        Thus $\{x\}$ is closed in $X$ with respect to $T$ by \cref{closed_iff_contains_limit_points}.
    \end{subproof}
    Let $F = \{\{x\} \mid x\in A\}$.
    Take $f$ such that $f$ is a bijection from $A$ to $F$ by \cref{singleton_family_bijection}.
    Thus $F$ is finite by \cref{finite_bijection}.
    Show that $F\subseteq\pow{X}$.
    \begin{subproof}
        Show that for all $B$ such that $B\in F$ we have $B\in\pow{X}$.
        \begin{subproof}
            Fix $B$.
            Assume $B\in F$.
            Take $x\in A$ such that $B = \{x\}$.
            Thus $x\in X$ by \cref{elem_subseteq}.
            Thus $\{x\}\subseteq X$ by \cref{singleton_subset_intro}.
            Thus $B\subseteq X$.
            Thus $B\in\pow{X}$ by \cref{powerset_intro}.
        \end{subproof}
        Thus $F\subseteq\pow{X}$ by \cref{subseteq}.
    \end{subproof}
    Show that for all $B$ such that $B\in F$ we have $B$ is closed in $X$ with respect to $T$.
    \begin{subproof}
        Fix $B$.
        Assume $B\in F$.
        Take $x\in A$ such that $B = \{x\}$.
        Thus $x\in X$ by \cref{elem_subseteq}.
        Thus $\{x\}$ is closed in $X$ with respect to $T$.
        Thus $B$ is closed in $X$ with respect to $T$.
    \end{subproof}
    Thus $\unions{F}$ is closed in $X$ with respect to $T$ by \cref{closed_finite_unions}.
    Show that $\unions{F} = A$.
    \begin{subproof}
        Show that $A\subseteq\unions{F}$.
        \begin{subproof}
            Show that for all $a$ such that $a\in A$ we have $a\in\unions{F}$.
            \begin{subproof}
                Fix $a$.
                Assume $a\in A$.
                Thus $\{a\}\in F$.
                Thus $a\in\{a\}$ by \cref{singleton_intro}.
                Thus $a\in\unions{F}$ by \cref{unions_iff}.
            \end{subproof}
            Thus $A\subseteq\unions{F}$ by \cref{subseteq}.
        \end{subproof}
        Show that $\unions{F}\subseteq A$.
        \begin{subproof}
            Show that for all $a$ such that $a\in\unions{F}$ we have $a\in A$.
            \begin{subproof}
                Fix $a$.
                Assume $a\in\unions{F}$.
                Take $B\in F$ such that $a\in B$ by \cref{unions_iff}.
                Take $x\in A$ such that $B = \{x\}$.
                Thus $a\in\{x\}$.
                Thus $a = x$ by \cref{singleton_elim}.
                Thus $a\in A$.
            \end{subproof}
            Thus $\unions{F}\subseteq A$ by \cref{subseteq}.
        \end{subproof}
        Thus $\unions{F} = A$ by \cref{subseteq_antisymmetric}.
    \end{subproof}
    Thus $A$ is closed in $X$ with respect to $T$.
\end{proof}

% from §17 Item 27: T1 axiom
\begin{definition}\label{t1_axiom}
    $X$ satisfies the T-one axiom with respect to $T$ iff
        $T$ is a topology on $X$ and
        for all $A$ such that $A\subseteq X$ and $A$ is finite
        we have $A$ is closed in $X$ with respect to $T$.
\end{definition}

% from §17 Item 28: limit points in $T_1$ spaces
\begin{theorem}\label{limit_point_infinite_neighborhood}
    Assume $X$ satisfies the T-one axiom with respect to $T$.
    Suppose $A\subseteq X$.
    Suppose $x\in X$.
    Then $x$ is a limit point of $A$ in $X$ with respect to $T$ iff
        for all $U$ such that $U$ is a neighborhood of $x$ in $X$ with respect to $T$ we have $U\inter A$ is infinite.
\end{theorem}
\begin{proof}
    Show that if $x$ is a limit point of $A$ in $X$ with respect to $T$ then
        for all $U$ such that $U$ is a neighborhood of $x$ in $X$ with respect to $T$ we have $U\inter A$ is infinite.
    \begin{subproof}
        Assume $x$ is a limit point of $A$ in $X$ with respect to $T$.
        Show that for all $U$ such that $U$ is a neighborhood of $x$ in $X$ with respect to $T$ we have $U\inter A$ is infinite.
        \begin{subproof}
            Fix $U$.
            Assume $U$ is a neighborhood of $x$ in $X$ with respect to $T$.
            Suppose not.
            Then $U\inter A$ is finite.
            Let $B = U\inter (A\setminus\{x\})$.
            Show that $B$ is finite.
            \begin{subproof}
                Show that $B\subseteq U\inter A$.
                \begin{subproof}
                    Show that $B\subseteq U$.
                    \begin{subproof}
                        Show that for all $y$ such that $y\in B$ we have $y\in U$.
                        \begin{subproof}
                            Fix $y$.
                            Assume $y\in B$.
                            Thus $y\in U$ by \cref{inter}.
                        \end{subproof}
                        Thus $B\subseteq U$ by \cref{subseteq}.
                    \end{subproof}
                    Show that $B\subseteq A$.
                    \begin{subproof}
                        Show that for all $y$ such that $y\in B$ we have $y\in A$.
                        \begin{subproof}
                            Fix $y$.
                            Assume $y\in B$.
                            Thus $y\in A\setminus\{x\}$ by \cref{inter}.
                            Thus $y\in A$ by \cref{setminus_elim_left}.
                        \end{subproof}
                        Thus $B\subseteq A$ by \cref{subseteq}.
                    \end{subproof}
                    Thus $B\subseteq U\inter A$ by \cref{subseteq_inter_iff}.
                \end{subproof}
                Thus $B$ is finite by \cref{subseteq_finite_is_finite}.
            \end{subproof}
            $B\subseteq A$.
            Thus $B\subseteq X$ by \cref{subseteq_transitive}.
            $T$ is a topology on $X$ by \cref{t1_axiom}.
            Thus $B$ is closed in $X$ with respect to $T$ by \cref{t1_axiom}.
            Thus $X\setminus B$ is open in $X$ with respect to $T$ by \cref{closed_in}.
            $U$ is open in $X$ with respect to $T$ by \cref{neighborhood}.
            $U\in T$ by \cref{open_in,neighborhood}.
            $X\setminus B\in T$ by \cref{open_in,closed_in}.
            $T$ is closed under binary intersections by \cref{topology_on,t1_axiom}.
            Thus $U\inter (X\setminus B)\in T$.
            Thus $U\inter (X\setminus B)$ is open in $X$ with respect to $T$ by \cref{open_in}.
            Show that $x\in U\inter (X\setminus B)$.
            \begin{subproof}
                $x\in U$ by \cref{neighborhood}.
                Show that $x\notin B$.
                \begin{subproof}
                    Suppose not.
                    Then $x\in B$.
                    Thus $x\in A\setminus\{x\}$ by \cref{inter}.
                    Thus $x\notin\{x\}$ by \cref{setminus_elim_right}.
                    Thus $x\in\{x\}$ by \cref{singleton_intro}.
                    Contradiction.
                \end{subproof}
                Thus $x\in X\setminus B$ by \cref{setminus_intro}.
                Thus $x\in U\inter (X\setminus B)$ by \cref{inter_intro}.
            \end{subproof}
            Thus $U\inter (X\setminus B)$ is a neighborhood of $x$ in $X$ with respect to $T$ by \cref{neighborhood}.
            Show that $U\inter (X\setminus B)$ does not intersect $A\setminus\{x\}$.
            \begin{subproof}
                Suppose not.
                Then $U\inter (X\setminus B)$ intersects $A\setminus\{x\}$.
                Take $y$ such that $y\in (U\inter (X\setminus B)) \inter (A\setminus\{x\})$
                    by \cref{intersects,nonempty_inhabited}.
                Thus $y\in U$ and $y\in X\setminus B$ and $y\in A\setminus\{x\}$ by \cref{inter}.
                Thus $y\in U\inter (A\setminus\{x\})$ by \cref{inter_intro}.
                Thus $y\in B$.
                Thus $y\notin X\setminus B$ by \cref{setminus_elim_right}.
                Contradiction.
            \end{subproof}
            Thus $U\inter (X\setminus B)$ is a neighborhood of $x$ in $X$ with respect to $T$.
            Thus $U\inter (X\setminus B)$ does not intersect $A\setminus\{x\}$.
            Contradiction by \cref{limit_point}.
        \end{subproof}
    \end{subproof}
    Show that if for all $U$ such that $U$ is a neighborhood of $x$ in $X$ with respect to $T$ we have $U\inter A$ is infinite then
        $x$ is a limit point of $A$ in $X$ with respect to $T$.
    \begin{subproof}
        Assume for all $U$ such that $U$ is a neighborhood of $x$ in $X$ with respect to $T$ we have $U\inter A$ is infinite.
        Show that for all $U$ such that $U$ is a neighborhood of $x$ in $X$ with respect to $T$ we have
            $U$ intersects $A\setminus\{x\}$.
        \begin{subproof}
            Fix $U$.
            Assume $U$ is a neighborhood of $x$ in $X$ with respect to $T$.
            Suppose not.
            Then $U\inter (A\setminus\{x\}) = \emptyset$.
            Show that $U\inter A\subseteq \{x\}$.
            \begin{subproof}
                Show that for all $y$ such that $y\in U\inter A$ we have $y\in\{x\}$.
                \begin{subproof}
                    Fix $y$.
                    Assume $y\in U\inter A$.
                    Show that $y\in\{x\}$.
                    \begin{subproof}
                        Suppose not.
                        Then $y\notin\{x\}$.
                        Thus $y\in A$ by \cref{inter}.
                        Thus $y\in A\setminus\{x\}$ by \cref{setminus_intro}.
                        Thus $y\in U$ by \cref{inter}.
                        Thus $y\in U\inter (A\setminus\{x\})$ by \cref{inter_intro}.
                        Thus $y\in\emptyset$.
                        Contradiction by \cref{emptyset}.
                    \end{subproof}
                \end{subproof}
                Thus $U\inter A\subseteq \{x\}$ by \cref{subseteq}.
            \end{subproof}
            Show that $\{x\}$ is finite.
            \begin{subproof}
                $x\in\{x,x\}$ by \cref{upair_iff}.
                Thus $\{x\}\subseteq\{x,x\}$ by \cref{singleton_subset_intro}.
                $\{x,x\}$ is finite by \cref{pair_is_finite}.
                Thus $\{x\}$ is finite by \cref{subseteq_finite_is_finite}.
            \end{subproof}
            Thus $U\inter A$ is finite by \cref{subseteq_finite_is_finite}.
            Thus $U\inter A$ is not finite.
            Contradiction.
        \end{subproof}
        Thus $x$ is a limit point of $A$ in $X$ with respect to $T$ by \cref{limit_point}.
    \end{subproof}
\end{proof}

% from §17 Item 29: sequence in a set
\begin{definition}\label{sequence_in}
    $s$ is a sequence in $X$ iff $s\in\funs{\naturalsPlus}{X}$.
\end{definition}

% from §17 Item 30: convergence of a sequence
\begin{definition}\label{converges_to}
    $s$ converges to $x$ in $X$ with respect to $T$ iff
        $s$ is a sequence in $X$ and
        for all $U$ such that $U$ is a neighborhood of $x$ in $X$ with respect to $T$
        there exists $N\in\naturalsPlus$ such that for all $n$ such that $n\in\naturalsPlus$ if $N\leq n$ then $s(n)\in U$.
\end{definition}

% from §17 Item 31: limits of sequences are unique in Hausdorff spaces
\begin{theorem}\label{sequence_limit_unique}
    Assume $X$ is Hausdorff with respect to $T$.
    Suppose $s$ is a sequence in $X$.
    Suppose $x\in X$.
    Suppose $y\in X$.
    Suppose $s$ converges to $x$ in $X$ with respect to $T$.
    Suppose $s$ converges to $y$ in $X$ with respect to $T$.
    Then $x = y$.
\end{theorem}
\begin{proof}
    Suppose not.
    Then $x\neq y$.
    Take $U_1,U_2$ such that
        $U_1$ is a neighborhood of $x$ in $X$ with respect to $T$ and
        $U_2$ is a neighborhood of $y$ in $X$ with respect to $T$ and
        $U_1$ is disjoint from $U_2$ by \cref{hausdorff}.
    Take $N_1\in\naturalsPlus$ such that for all $n$ such that $n\in\naturalsPlus$ if $N_1\leq n$ then $s(n)\in U_1$ by \cref{converges_to}.
    Take $N_2\in\naturalsPlus$ such that for all $n$ such that $n\in\naturalsPlus$ if $N_2\leq n$ then $s(n)\in U_2$ by \cref{converges_to}.
    Let $A = \{N_1,N_2\}$.
    Show that $A\subseteq\naturalsPlus$.
    \begin{subproof}
        Show that for all $a$ such that $a\in A$ we have $a\in\naturalsPlus$.
        \begin{subproof}
            Fix $a$.
            Assume $a\in A$.
            Then $a = N_1$ or $a = N_2$ by \cref{upair_iff}.
            Thus $a\in\naturalsPlus$.
        \end{subproof}
        Thus $A\subseteq\naturalsPlus$ by \cref{subseteq}.
    \end{subproof}
    $A$ is finite by \cref{pair_is_finite}.
    Show that $A\neq\emptyset$.
    \begin{subproof}
        Suppose not.
        Then $A = \emptyset$.
        $N_1\in A$ by \cref{upair_iff}.
        Thus $N_1\in\emptyset$.
        Contradiction by \cref{emptyset}.
    \end{subproof}
    Take $m\in A$ such that for all $a\in A$ we have $a\leq m$ by \cref{finite_naturals_maximum}.
    $N_1\in A$ by \cref{upair_iff}.
    $N_2\in A$ by \cref{upair_iff}.
    Thus $N_1\leq m$ and $N_2\leq m$.
    Thus $s(m)\in U_1$ and $s(m)\in U_2$.
    Thus there exists $a$ such that $a\in U_1$ and $a\in U_2$.
    Contradiction by \cref{disjoint}.
\end{proof}

% from §17 helper lemma 5: open right ray is open in order topology
\begin{lemma}\label{openrayright_open_in_order_topology}
    Assume $R$ is a simple order relation on $X$.
    Assume $X$ has more than one element.
    Assume $T$ is the order topology on $X$ with respect to $R$.
    Suppose $a\in X$.
    Then $\openrayright{R}{X}{a}$ is open in $X$ with respect to $T$.
\end{lemma}
\begin{proof}
    Let $S = \openRays{R}{X}$.
    $S$ is a subbasis on $X$ by \cref{open_rays_subbasis_on}.
    Thus $\unions{S} = X$ by \cref{subbasis_on}.
    Let $U = \openrayright{R}{X}{a}$.
    Show that $U\subseteq X$.
    \begin{subproof}
        Show that for all $x$ such that $x\in U$ we have $x\in X$.
        \begin{subproof}
            Fix $x$.
            Assume $x\in U$.
            Thus $x\in X$ by \cref{openray_right}.
        \end{subproof}
        Thus $U\subseteq X$ by \cref{subseteq}.
    \end{subproof}
    Thus $U\in\pow{X}$ by \cref{powerset_intro}.
    Thus $U\in S$ by \cref{open_rays}.
    Thus $U\subseteq\unions{S}$ by \cref{subseteq_transitive}.
    Thus $U\in\pow{\unions{S}}$ by \cref{powerset_intro}.
    Let $F = \{U\}$.
    $F\subseteq S$ by \cref{singleton_subset_intro}.
    Show that $F$ is finite.
    \begin{subproof}
        $U\in\{U,U\}$ by \cref{upair_iff}.
        Thus $F\subseteq\{U,U\}$ by \cref{singleton_subset_intro}.
        $\{U,U\}$ is finite by \cref{pair_is_finite}.
        Thus $F$ is finite by \cref{subseteq_finite_is_finite}.
    \end{subproof}
    Show that $F$ is inhabited.
    \begin{subproof}
        Show that $F\neq\emptyset$.
        \begin{subproof}
            Suppose not.
            Then $F = \emptyset$.
            $U\in F$ by \cref{singleton_intro}.
            Thus $U\in\emptyset$.
            Contradiction by \cref{emptyset}.
        \end{subproof}
        Thus $F$ is inhabited by \cref{nonempty_inhabited}.
    \end{subproof}
    $\inters{F} = U$ by \cref{inters_singleton}.
    Thus $U\in\finiteIntersections{S}$ by \cref{finite_intersections}.
    Thus $\finiteIntersections{S}$ is a basis on $X$ by \cref{subbasis_finite_intersections_basis}.
    Thus $U\in\basisTopology{\finiteIntersections{S}}{X}$ by \cref{basis_elem_open_in_generated}.
    Thus $U\in\subbasisTopology{S}{X}$ by \cref{topology_generated_by_subbasis}.
    $\subbasisTopology{S}{X} = T$ by \cref{open_rays_subbasis}.
    Thus $U\in T$.
    $\orderBasis{R}{X}$ is a basis on $X$ by \cref{order_basis_is_basis}.
    Thus $T$ is a topology on $X$ by \cref{basis_topology_is_topology,order_topology}.
    Thus $U$ is open in $X$ with respect to $T$ by \cref{open_in}.
\end{proof}

% from §17 helper lemma 6: open left ray is open in order topology
\begin{lemma}\label{openrayleft_open_in_order_topology}
    Assume $R$ is a simple order relation on $X$.
    Assume $X$ has more than one element.
    Assume $T$ is the order topology on $X$ with respect to $R$.
    Suppose $a\in X$.
    Then $\openrayleft{R}{X}{a}$ is open in $X$ with respect to $T$.
\end{lemma}
\begin{proof}
    Let $S = \openRays{R}{X}$.
    $S$ is a subbasis on $X$ by \cref{open_rays_subbasis_on}.
    Thus $\unions{S} = X$ by \cref{subbasis_on}.
    Let $U = \openrayleft{R}{X}{a}$.
    Show that $U\subseteq X$.
    \begin{subproof}
        Show that for all $x$ such that $x\in U$ we have $x\in X$.
        \begin{subproof}
            Fix $x$.
            Assume $x\in U$.
            Thus $x\in X$ by \cref{openray_left}.
        \end{subproof}
        Thus $U\subseteq X$ by \cref{subseteq}.
    \end{subproof}
    Thus $U\in\pow{X}$ by \cref{powerset_intro}.
    Thus $U\in S$ by \cref{open_rays}.
    Thus $U\subseteq\unions{S}$ by \cref{subseteq_transitive}.
    Thus $U\in\pow{\unions{S}}$ by \cref{powerset_intro}.
    Let $F = \{U\}$.
    $F\subseteq S$ by \cref{singleton_subset_intro}.
    Show that $F$ is finite.
    \begin{subproof}
        $U\in\{U,U\}$ by \cref{upair_iff}.
        Thus $F\subseteq\{U,U\}$ by \cref{singleton_subset_intro}.
        $\{U,U\}$ is finite by \cref{pair_is_finite}.
        Thus $F$ is finite by \cref{subseteq_finite_is_finite}.
    \end{subproof}
    Show that $F$ is inhabited.
    \begin{subproof}
        Show that $F\neq\emptyset$.
        \begin{subproof}
            Suppose not.
            Then $F = \emptyset$.
            $U\in F$ by \cref{singleton_intro}.
            Thus $U\in\emptyset$.
            Contradiction by \cref{emptyset}.
        \end{subproof}
        Thus $F$ is inhabited by \cref{nonempty_inhabited}.
    \end{subproof}
    $\inters{F} = U$ by \cref{inters_singleton}.
    Thus $U\in\finiteIntersections{S}$ by \cref{finite_intersections}.
    Thus $\finiteIntersections{S}$ is a basis on $X$ by \cref{subbasis_finite_intersections_basis}.
    Thus $U\in\basisTopology{\finiteIntersections{S}}{X}$ by \cref{basis_elem_open_in_generated}.
    Thus $U\in\subbasisTopology{S}{X}$ by \cref{topology_generated_by_subbasis}.
    $\subbasisTopology{S}{X} = T$ by \cref{open_rays_subbasis}.
    Thus $U\in T$.
    $\orderBasis{R}{X}$ is a basis on $X$ by \cref{order_basis_is_basis}.
    Thus $T$ is a topology on $X$ by \cref{basis_topology_is_topology,order_topology}.
    Thus $U$ is open in $X$ with respect to $T$ by \cref{open_in}.
\end{proof}

% from §17 Item 32: order topology is Hausdorff
\begin{theorem}\label{order_topology_hausdorff}
    Assume $R$ is a simple order relation on $X$.
    Assume $X$ has more than one element.
    Assume $T$ is the order topology on $X$ with respect to $R$.
    Then $X$ is Hausdorff with respect to $T$.
\end{theorem}
\begin{proof}
    Show that $T$ is a topology on $X$.
    \begin{subproof}
        $\orderBasis{R}{X}$ is a basis on $X$ by \cref{order_basis_is_basis}.
        Thus $T$ is a topology on $X$ by \cref{basis_topology_is_topology,order_topology}.
    \end{subproof}
    Show that for all $x_1,x_2$ such that $x_1\in X$ and $x_2\in X$ and $x_1\neq x_2$ there exist $U_1,U_2$ such that
        $U_1$ is a neighborhood of $x_1$ in $X$ with respect to $T$ and
        $U_2$ is a neighborhood of $x_2$ in $X$ with respect to $T$ and
        $U_1$ is disjoint from $U_2$.
    \begin{subproof}
        Fix $x_1$.
        Fix $x_2$.
        Assume $x_1\in X$ and $x_2\in X$ and $x_1\neq x_2$.
        $R$ is connex on $X$ by \cref{simple_order_relation}.
        Thus $x_1\mathrel{R}x_2$ or $x_2\mathrel{R}x_1$ by \cref{connex}.
        \begin{byCase}
        \caseOf{$x_1\mathrel{R}x_2$.}
            \begin{byCase}
            \caseOf{there exists $z\in X$ such that $x_1\mathrel{R}z$ and $z\mathrel{R}x_2$.}
                Take $z\in X$ such that $x_1\mathrel{R}z$ and $z\mathrel{R}x_2$.
                Let $U_1 = \openrayleft{R}{X}{z}$.
                Let $U_2 = \openrayright{R}{X}{z}$.
                $U_1$ is open in $X$ with respect to $T$ by \cref{openrayleft_open_in_order_topology}.
                $U_2$ is open in $X$ with respect to $T$ by \cref{openrayright_open_in_order_topology}.
                $x_1\in U_1$ by \cref{openray_left}.
                $x_2\in U_2$ by \cref{openray_right}.
                Thus $U_1$ is a neighborhood of $x_1$ in $X$ with respect to $T$ by \cref{neighborhood}.
                Thus $U_2$ is a neighborhood of $x_2$ in $X$ with respect to $T$ by \cref{neighborhood}.
                Show that $U_1$ is disjoint from $U_2$.
                \begin{subproof}
                    Suppose not.
                    Then there exists $y$ such that $y\in U_1$ and $y\in U_2$.
                    Thus $y\mathrel{R}z$ by \cref{openray_left}.
                    Thus $z\mathrel{R}y$ by \cref{openray_right}.
                    $R$ is asymmetric by \cref{simple_order_relation}.
                    Contradiction by \cref{asymmetric}.
                \end{subproof}
            \caseOf{there exists no $z\in X$ such that $x_1\mathrel{R}z$ and $z\mathrel{R}x_2$.}
                Let $U_1 = \openrayleft{R}{X}{x_2}$.
                Let $U_2 = \openrayright{R}{X}{x_1}$.
                $U_1$ is open in $X$ with respect to $T$ by \cref{openrayleft_open_in_order_topology}.
                $U_2$ is open in $X$ with respect to $T$ by \cref{openrayright_open_in_order_topology}.
                $x_1\in U_1$ by \cref{openray_left}.
                $x_2\in U_2$ by \cref{openray_right}.
                Thus $U_1$ is a neighborhood of $x_1$ in $X$ with respect to $T$ by \cref{neighborhood}.
                Thus $U_2$ is a neighborhood of $x_2$ in $X$ with respect to $T$ by \cref{neighborhood}.
                Show that $U_1$ is disjoint from $U_2$.
                \begin{subproof}
                    Suppose not.
                    Then there exists $y$ such that $y\in U_1$ and $y\in U_2$.
                    Thus $x_1\mathrel{R}y$ by \cref{openray_right}.
                    Thus $y\mathrel{R}x_2$ by \cref{openray_left}.
                    Thus there exists $z\in X$ such that $x_1\mathrel{R}z$ and $z\mathrel{R}x_2$.
                    Contradiction.
                \end{subproof}
            \end{byCase}
        \caseOf{$x_2\mathrel{R}x_1$.}
            \begin{byCase}
            \caseOf{there exists $z\in X$ such that $x_2\mathrel{R}z$ and $z\mathrel{R}x_1$.}
                Take $z\in X$ such that $x_2\mathrel{R}z$ and $z\mathrel{R}x_1$.
                Let $U_2 = \openrayleft{R}{X}{z}$.
                Let $U_1 = \openrayright{R}{X}{z}$.
                $U_1$ is open in $X$ with respect to $T$ by \cref{openrayright_open_in_order_topology}.
                $U_2$ is open in $X$ with respect to $T$ by \cref{openrayleft_open_in_order_topology}.
                $x_2\in U_2$ by \cref{openray_left}.
                $x_1\in U_1$ by \cref{openray_right}.
                Thus $U_1$ is a neighborhood of $x_1$ in $X$ with respect to $T$ by \cref{neighborhood}.
                Thus $U_2$ is a neighborhood of $x_2$ in $X$ with respect to $T$ by \cref{neighborhood}.
                Show that $U_1$ is disjoint from $U_2$.
                \begin{subproof}
                    Suppose not.
                    Then there exists $y$ such that $y\in U_1$ and $y\in U_2$.
                    Thus $z\mathrel{R}y$ by \cref{openray_right}.
                    Thus $y\mathrel{R}z$ by \cref{openray_left}.
                    $R$ is asymmetric by \cref{simple_order_relation}.
                    Contradiction by \cref{asymmetric}.
                \end{subproof}
            \caseOf{there exists no $z\in X$ such that $x_2\mathrel{R}z$ and $z\mathrel{R}x_1$.}
                Let $U_2 = \openrayleft{R}{X}{x_1}$.
                Let $U_1 = \openrayright{R}{X}{x_2}$.
                $U_1$ is open in $X$ with respect to $T$ by \cref{openrayright_open_in_order_topology}.
                $U_2$ is open in $X$ with respect to $T$ by \cref{openrayleft_open_in_order_topology}.
                $x_2\in U_2$ by \cref{openray_left}.
                $x_1\in U_1$ by \cref{openray_right}.
                Thus $U_1$ is a neighborhood of $x_1$ in $X$ with respect to $T$ by \cref{neighborhood}.
                Thus $U_2$ is a neighborhood of $x_2$ in $X$ with respect to $T$ by \cref{neighborhood}.
                Show that $U_1$ is disjoint from $U_2$.
                \begin{subproof}
                    Suppose not.
                    Then there exists $y$ such that $y\in U_1$ and $y\in U_2$.
                    Thus $x_2\mathrel{R}y$ by \cref{openray_right}.
                    Thus $y\mathrel{R}x_1$ by \cref{openray_left}.
                    Thus there exists $z\in X$ such that $x_2\mathrel{R}z$ and $z\mathrel{R}x_1$.
                    Contradiction.
                \end{subproof}
            \end{byCase}
        \end{byCase}
    \end{subproof}
    Thus $X$ is Hausdorff with respect to $T$ by \cref{hausdorff}.
\end{proof}

% from §17 Item 33: product of Hausdorff spaces is Hausdorff
\begin{theorem}\label{product_hausdorff}
    Assume $X$ is Hausdorff with respect to $T$.
    Assume $Y$ is Hausdorff with respect to $T_1$.
    Let $T_2 = \productTopology{T}{T_1}{X}{Y}$.
    Then $X\times Y$ is Hausdorff with respect to $T_2$.
\end{theorem}
\begin{proof}
    $T$ is a topology on $X$ by \cref{hausdorff}.
    $T_1$ is a topology on $Y$ by \cref{hausdorff}.
    $\productBasis{T}{T_1}{X}{Y}$ is a basis on $X\times Y$ by \cref{product_basis_is_basis}.
    Thus $T_2$ is a topology on $X\times Y$ by \cref{basis_topology_is_topology,product_topology}.
    Show that for all $p_1,p_2$ such that $p_1\in X\times Y$ and $p_2\in X\times Y$ and $p_1\neq p_2$
        there exist $U_1,U_2$ such that
            $U_1$ is a neighborhood of $p_1$ in $X\times Y$ with respect to $T_2$ and
            $U_2$ is a neighborhood of $p_2$ in $X\times Y$ with respect to $T_2$ and
            $U_1$ is disjoint from $U_2$.
    \begin{subproof}
        Fix $p_1$.
        Fix $p_2$.
        Assume $p_1\in X\times Y$ and $p_2\in X\times Y$ and $p_1\neq p_2$.
        Take $x_1,y_1$ such that $x_1\in X$ and $y_1\in Y$ and $p_1 = (x_1,y_1)$ by \cref{times_elem_is_tuple}.
        Take $x_2,y_2$ such that $x_2\in X$ and $y_2\in Y$ and $p_2 = (x_2,y_2)$ by \cref{times_elem_is_tuple}.
        Show that $x_1\neq x_2$ or $y_1\neq y_2$.
        \begin{subproof}
            Suppose not.
            Then $x_1 = x_2$ and $y_1 = y_2$.
            Then $p_1 = p_2$ by \cref{pair_eq_iff}.
            Contradiction.
        \end{subproof}
        \begin{byCase}
        \caseOf{$x_1\neq x_2$.}
            Take $U_1,U_2$ such that
                $U_1$ is a neighborhood of $x_1$ in $X$ with respect to $T$ and
                $U_2$ is a neighborhood of $x_2$ in $X$ with respect to $T$ and
                $U_1$ is disjoint from $U_2$ by \cref{hausdorff}.
            Let $V_1 = Y$.
            Let $V_2 = Y$.
            $Y\in T_1$ by \cref{topology_on}.
            Thus $V_1$ is open in $Y$ with respect to $T_1$ by \cref{open_in}.
            Thus $V_2$ is open in $Y$ with respect to $T_1$ by \cref{open_in}.
            $U_1$ is open in $X$ with respect to $T$ by \cref{neighborhood}.
            $U_2$ is open in $X$ with respect to $T$ by \cref{neighborhood}.
            Let $W_1 = U_1\times V_1$.
            Let $W_2 = U_2\times V_2$.
            $U_1\in T$ by \cref{open_in,neighborhood}.
            $U_2\in T$ by \cref{open_in,neighborhood}.
            $V_1\in T_1$ by \cref{open_in}.
            $V_2\in T_1$ by \cref{open_in}.
            $T\subseteq\pow{X}$ by \cref{topology_on}.
            $T_1\subseteq\pow{Y}$ by \cref{topology_on}.
            Thus $U_1\in\pow{X}$ and $U_2\in\pow{X}$ by \cref{subseteq}.
            Thus $V_1\in\pow{Y}$ and $V_2\in\pow{Y}$ by \cref{subseteq}.
            Thus $U_1\subseteq X$ and $U_2\subseteq X$ by \cref{pow_iff}.
            Thus $V_1\subseteq Y$ and $V_2\subseteq Y$ by \cref{pow_iff}.
            $W_1\subseteq X\times V_1$ by \cref{times_subseteq_left}.
            $W_2\subseteq X\times V_2$ by \cref{times_subseteq_left}.
            $X\times V_1\subseteq X\times Y$ by \cref{times_subseteq_right}.
            $X\times V_2\subseteq X\times Y$ by \cref{times_subseteq_right}.
            Thus $W_1\subseteq X\times Y$ by \cref{subseteq_transitive}.
            Thus $W_2\subseteq X\times Y$ by \cref{subseteq_transitive}.
            Thus $W_1\in\pow{X\times Y}$ by \cref{powerset_intro}.
            Thus $W_2\in\pow{X\times Y}$ by \cref{powerset_intro}.
            Thus $W_1\in\productBasis{T}{T_1}{X}{Y}$ by \cref{product_basis}.
            Thus $W_2\in\productBasis{T}{T_1}{X}{Y}$ by \cref{product_basis}.
            Thus $W_1\in\basisTopology{\productBasis{T}{T_1}{X}{Y}}{X\times Y}$ by \cref{basis_elem_open_in_generated}.
            Thus $W_2\in\basisTopology{\productBasis{T}{T_1}{X}{Y}}{X\times Y}$ by \cref{basis_elem_open_in_generated}.
            Thus $W_1\in T_2$ by \cref{product_topology}.
            Thus $W_2\in T_2$ by \cref{product_topology}.
            Thus $W_1$ is open in $X\times Y$ with respect to $T_2$ by \cref{open_in}.
            Thus $W_2$ is open in $X\times Y$ with respect to $T_2$ by \cref{open_in}.
            $x_1\in U_1$ by \cref{neighborhood}.
            $x_2\in U_2$ by \cref{neighborhood}.
            Thus $(x_1,y_1)\in W_1$ by \cref{times_tuple_intro}.
            Thus $(x_2,y_2)\in W_2$ by \cref{times_tuple_intro}.
            Thus $W_1$ is a neighborhood of $p_1$ in $X\times Y$ with respect to $T_2$ by \cref{neighborhood}.
            Thus $W_2$ is a neighborhood of $p_2$ in $X\times Y$ with respect to $T_2$ by \cref{neighborhood}.
            Show that $W_1$ is disjoint from $W_2$.
            \begin{subproof}
                Suppose not.
                Then there exists $p$ such that $p\in W_1$ and $p\in W_2$.
                Take $x,y$ such that $p = (x,y)$ and $x\in U_1$ and $y\in V_1$ by \cref{times_elem_is_tuple}.
                Thus $x\in U_2$ by \cref{times_tuple_elim}.
                Thus there exists $a$ such that $a\in U_1$ and $a\in U_2$.
                Contradiction by \cref{disjoint}.
            \end{subproof}
        \caseOf{$y_1\neq y_2$.}
            Take $V_1,V_2$ such that
                $V_1$ is a neighborhood of $y_1$ in $Y$ with respect to $T_1$ and
                $V_2$ is a neighborhood of $y_2$ in $Y$ with respect to $T_1$ and
                $V_1$ is disjoint from $V_2$ by \cref{hausdorff}.
            Let $U_1 = X$.
            Let $U_2 = X$.
            $X\in T$ by \cref{topology_on}.
            Thus $U_1$ is open in $X$ with respect to $T$ by \cref{open_in}.
            Thus $U_2$ is open in $X$ with respect to $T$ by \cref{open_in}.
            $V_1$ is open in $Y$ with respect to $T_1$ by \cref{neighborhood}.
            $V_2$ is open in $Y$ with respect to $T_1$ by \cref{neighborhood}.
            Let $W_1 = U_1\times V_1$.
            Let $W_2 = U_2\times V_2$.
            $U_1\in T$ by \cref{open_in}.
            $U_2\in T$ by \cref{open_in}.
            $V_1\in T_1$ by \cref{open_in,neighborhood}.
            $V_2\in T_1$ by \cref{open_in,neighborhood}.
            $T\subseteq\pow{X}$ by \cref{topology_on}.
            $T_1\subseteq\pow{Y}$ by \cref{topology_on}.
            Thus $U_1\in\pow{X}$ and $U_2\in\pow{X}$ by \cref{subseteq}.
            Thus $V_1\in\pow{Y}$ and $V_2\in\pow{Y}$ by \cref{subseteq}.
            Thus $U_1\subseteq X$ and $U_2\subseteq X$ by \cref{pow_iff}.
            Thus $V_1\subseteq Y$ and $V_2\subseteq Y$ by \cref{pow_iff}.
            $W_1\subseteq X\times V_1$ by \cref{times_subseteq_left}.
            $W_2\subseteq X\times V_2$ by \cref{times_subseteq_left}.
            $X\times V_1\subseteq X\times Y$ by \cref{times_subseteq_right}.
            $X\times V_2\subseteq X\times Y$ by \cref{times_subseteq_right}.
            Thus $W_1\subseteq X\times Y$ by \cref{subseteq_transitive}.
            Thus $W_2\subseteq X\times Y$ by \cref{subseteq_transitive}.
            Thus $W_1\in\pow{X\times Y}$ by \cref{powerset_intro}.
            Thus $W_2\in\pow{X\times Y}$ by \cref{powerset_intro}.
            Thus $W_1\in\productBasis{T}{T_1}{X}{Y}$ by \cref{product_basis}.
            Thus $W_2\in\productBasis{T}{T_1}{X}{Y}$ by \cref{product_basis}.
            Thus $W_1\in\basisTopology{\productBasis{T}{T_1}{X}{Y}}{X\times Y}$ by \cref{basis_elem_open_in_generated}.
            Thus $W_2\in\basisTopology{\productBasis{T}{T_1}{X}{Y}}{X\times Y}$ by \cref{basis_elem_open_in_generated}.
            Thus $W_1\in T_2$ by \cref{product_topology}.
            Thus $W_2\in T_2$ by \cref{product_topology}.
            Thus $W_1$ is open in $X\times Y$ with respect to $T_2$ by \cref{open_in}.
            Thus $W_2$ is open in $X\times Y$ with respect to $T_2$ by \cref{open_in}.
            $y_1\in V_1$ by \cref{neighborhood}.
            $y_2\in V_2$ by \cref{neighborhood}.
            Thus $(x_1,y_1)\in W_1$ by \cref{times_tuple_intro}.
            Thus $(x_2,y_2)\in W_2$ by \cref{times_tuple_intro}.
            Thus $W_1$ is a neighborhood of $p_1$ in $X\times Y$ with respect to $T_2$ by \cref{neighborhood}.
            Thus $W_2$ is a neighborhood of $p_2$ in $X\times Y$ with respect to $T_2$ by \cref{neighborhood}.
            Show that $W_1$ is disjoint from $W_2$.
            \begin{subproof}
                Suppose not.
                Then there exists $p$ such that $p\in W_1$ and $p\in W_2$.
                Take $x,y$ such that $p = (x,y)$ and $x\in U_1$ and $y\in V_1$ by \cref{times_elem_is_tuple}.
                Thus $y\in V_2$ by \cref{times_tuple_elim}.
                Thus there exists $b$ such that $b\in V_1$ and $b\in V_2$.
                Contradiction by \cref{disjoint}.
            \end{subproof}
        \end{byCase}
    \end{subproof}
    Thus $X\times Y$ is Hausdorff with respect to $T_2$ by \cref{hausdorff}.
\end{proof}

% from §17 Item 34: subspace of a Hausdorff space is Hausdorff
\begin{theorem}\label{subspace_hausdorff}
    Assume $X$ is Hausdorff with respect to $T$.
    Suppose $Y\subseteq X$.
    Suppose $T_1 = \subspaceTopology{T}{Y}$.
    Then $Y$ is Hausdorff with respect to $T_1$.
\end{theorem}
\begin{proof}
    $T$ is a topology on $X$ by \cref{hausdorff}.
    Thus $\subspaceTopology{T}{Y}$ is a topology on $Y$ by \cref{subspace_topology_is_topology}.
    Thus $T_1$ is a topology on $Y$.
    Show that for all $y_1,y_2$ such that $y_1\in Y$ and $y_2\in Y$ and $y_1\neq y_2$
        there exist $V_1,V_2$ such that
            $V_1$ is a neighborhood of $y_1$ in $Y$ with respect to $T_1$ and
            $V_2$ is a neighborhood of $y_2$ in $Y$ with respect to $T_1$ and
            $V_1$ is disjoint from $V_2$.
    \begin{subproof}
        Fix $y_1$.
        Fix $y_2$.
        Assume $y_1\in Y$ and $y_2\in Y$ and $y_1\neq y_2$.
        Thus $y_1\in X$ and $y_2\in X$ by \cref{elem_subseteq}.
        Take $U_1,U_2$ such that
            $U_1$ is a neighborhood of $y_1$ in $X$ with respect to $T$ and
            $U_2$ is a neighborhood of $y_2$ in $X$ with respect to $T$ and
            $U_1$ is disjoint from $U_2$ by \cref{hausdorff}.
        Let $V_1 = Y\inter U_1$.
        Let $V_2 = Y\inter U_2$.
        $U_1$ is open in $X$ with respect to $T$ by \cref{neighborhood}.
        $U_2$ is open in $X$ with respect to $T$ by \cref{neighborhood}.
        Thus $U_1\in T$ by \cref{open_in,neighborhood}.
        Thus $U_2\in T$ by \cref{open_in,neighborhood}.
        Thus $V_1\in\subspaceTopology{T}{Y}$ by \cref{subspace_topology}.
        Thus $V_2\in\subspaceTopology{T}{Y}$ by \cref{subspace_topology}.
        Thus $V_1\in T_1$.
        Thus $V_2\in T_1$.
        Thus $V_1$ is open in $Y$ with respect to $T_1$ by \cref{open_in}.
        Thus $V_2$ is open in $Y$ with respect to $T_1$ by \cref{open_in}.
        $y_1\in U_1$ by \cref{neighborhood}.
        $y_2\in U_2$ by \cref{neighborhood}.
        Thus $y_1\in V_1$ by \cref{inter_intro}.
        Thus $y_2\in V_2$ by \cref{inter_intro}.
        Thus $V_1$ is a neighborhood of $y_1$ in $Y$ with respect to $T_1$ by \cref{neighborhood}.
        Thus $V_2$ is a neighborhood of $y_2$ in $Y$ with respect to $T_1$ by \cref{neighborhood}.
        Show that $V_1$ is disjoint from $V_2$.
        \begin{subproof}
            Suppose not.
            Then there exists $y$ such that $y\in V_1$ and $y\in V_2$.
            Thus $y\in U_1$ and $y\in U_2$ by \cref{inter}.
            Thus there exists $a$ such that $a\in U_1$ and $a\in U_2$.
            Contradiction by \cref{disjoint}.
        \end{subproof}
    \end{subproof}
    Thus $Y$ is Hausdorff with respect to $T_1$ by \cref{hausdorff}.
\end{proof}

% from §18 Item 1: continuous function
\begin{definition}\label{continuous}
    $f$ is continuous from $X$ to $Y$ with respect to $T$ and $T_1$ iff
        $f$ is a function from $X$ to $Y$ and
        $T$ is a topology on $X$ and
        $T_1$ is a topology on $Y$ and
        for all $V$ such that $V$ is open in $Y$ with respect to $T_1$
        we have $\preimg{f}{V}$ is open in $X$ with respect to $T$.
\end{definition}

% from §18 helper lemma 1: preimage of an intersection under a function
\begin{lemma}\label{preimg_inter_function}
    Suppose $f$ is a function.
    Then $\preimg{f}{A\inter B} = \preimg{f}{A}\inter \preimg{f}{B}$.
\end{lemma}
\begin{proof}
    Show that for all $x$ such that $x\in\preimg{f}{A\inter B}$ we have
        $x\in\preimg{f}{A}\inter\preimg{f}{B}$.
    \begin{subproof}
        Fix $x$.
        Assume $x\in\preimg{f}{A\inter B}$.
        Take $y\in A\inter B$ such that $x\mathrel{f} y$ by \cref{preimg_iff}.
        Thus $y\in A$ and $y\in B$ by \cref{inter}.
        Thus $x\in\preimg{f}{A}$ and $x\in\preimg{f}{B}$ by \cref{preimg_iff}.
        Thus $x\in\preimg{f}{A}\inter\preimg{f}{B}$ by \cref{inter_intro}.
    \end{subproof}
    Thus $\preimg{f}{A\inter B}\subseteq \preimg{f}{A}\inter\preimg{f}{B}$ by \cref{subseteq}.
    Show that for all $x$ such that $x\in\preimg{f}{A}\inter\preimg{f}{B}$ we have
        $x\in\preimg{f}{A\inter B}$.
    \begin{subproof}
        Fix $x$.
        Assume $x\in\preimg{f}{A}\inter\preimg{f}{B}$.
        Thus $x\in\preimg{f}{A}$ and $x\in\preimg{f}{B}$ by \cref{inter}.
        Take $y_1\in A$ such that $x\mathrel{f} y_1$ by \cref{preimg_iff}.
        Take $y_2\in B$ such that $x\mathrel{f} y_2$ by \cref{preimg_iff}.
        $f$ is right-unique by \cref{function_rightunique}.
        Thus $y_1 = y_2$ by \hyperref[rightunique]{right-uniqueness}.
        Let $y = y_1$.
        Thus $y\in A$ and $y\in B$.
        Thus $y\in A\inter B$ by \cref{inter_intro}.
        Thus $x\in\preimg{f}{A\inter B}$ by \cref{preimg_iff}.
    \end{subproof}
    Thus $\preimg{f}{A}\inter\preimg{f}{B}\subseteq \preimg{f}{A\inter B}$ by \cref{subseteq}.
    Thus $\preimg{f}{A\inter B} = \preimg{f}{A}\inter\preimg{f}{B}$ by \cref{subseteq_antisymmetric}.
\end{proof}

% from §18 helper lemma 2: image of a preimage is contained in the set
\begin{lemma}\label{img_preimg_subseteq}
    Suppose $f$ is a function.
    Then $\img{f}{\preimg{f}{B}}\subseteq B$.
\end{lemma}
\begin{proof}
    Show that for all $y$ such that $y\in\img{f}{\preimg{f}{B}}$ we have $y\in B$.
    \begin{subproof}
        Fix $y$.
        Assume $y\in\img{f}{\preimg{f}{B}}$.
        Take $x\in\dom{f}\inter \preimg{f}{B}$ such that $y = f(x)$ by \cref{img_of_function_elim}.
        Thus $x\in\preimg{f}{B}$ by \cref{inter_elim_right}.
        Take $b\in B$ such that $x\mathrel{f} b$ by \cref{preimg_iff}.
        $(x,f(x))\in f$ by \cref{function_apply_elim,inter_elim_left}.
        $f$ is right-unique by \cref{function_rightunique}.
        Thus $b = f(x)$ by \hyperref[rightunique]{right-uniqueness}.
        Thus $y\in B$.
    \end{subproof}
    Thus $\img{f}{\preimg{f}{B}}\subseteq B$ by \cref{subseteq}.
\end{proof}

% from §18 helper lemma 3: preimage of a complement
\begin{lemma}\label{preimg_setminus_function}
    Assume $f$ is a function from $X$ to $Y$.
    Suppose $B\subseteq Y$.
    Then $\preimg{f}{Y\setminus B} = X\setminus \preimg{f}{B}$.
\end{lemma}
\begin{proof}
    Show that for all $x$ such that $x\in\preimg{f}{Y\setminus B}$ we have
        $x\in X\setminus \preimg{f}{B}$.
    \begin{subproof}
        Fix $x$.
        Assume $x\in\preimg{f}{Y\setminus B}$.
        Take $y\in Y\setminus B$ such that $x\mathrel{f} y$ by \cref{preimg_iff}.
        $f$ is a function from $X$ to $Y$.
        Thus $f\in\funs{X}{Y}$.
        Thus $f$ is a function by \cref{funs_is_function}.
        Thus $\dom{f} = X$ by \cref{funs}.
        Thus $x\in\dom{f}$ by \cref{subseteq,preimg_subseteq_dom}.
        Thus $x\in X$ by \cref{funs}.
        Show that $x\notin\preimg{f}{B}$.
        \begin{subproof}
            Suppose not.
            Then $x\in\preimg{f}{B}$.
            Then there exists $b\in B$ such that $x\mathrel{f} b$ by \cref{preimg_iff}.
            $f$ is right-unique by \cref{function_rightunique}.
            Thus $b = y$ by \hyperref[rightunique]{right-uniqueness}.
            Thus $y\in B$ by \cref{subseteq}.
            Contradiction by \cref{setminus_elim_right}.
        \end{subproof}
        Thus $x\in X\setminus \preimg{f}{B}$ by \cref{setminus_intro}.
    \end{subproof}
    Thus $\preimg{f}{Y\setminus B}\subseteq X\setminus \preimg{f}{B}$ by \cref{subseteq}.
    Show that for all $x$ such that $x\in X\setminus \preimg{f}{B}$ we have
        $x\in\preimg{f}{Y\setminus B}$.
    \begin{subproof}
        Fix $x$.
        Assume $x\in X\setminus \preimg{f}{B}$.
        Thus $x\in X$ and $x\notin\preimg{f}{B}$ by \cref{setminus_elim_left,setminus_elim_right}.
        $f$ is a function from $X$ to $Y$.
        Thus $f\in\funs{X}{Y}$.
        Thus $f$ is a function by \cref{funs_is_function}.
        Thus $\dom{f} = X$ by \cref{funs}.
        Thus $x\in\dom{f}$.
        Let $y = f(x)$.
        Thus $y\in Y$.
        Show that $y\notin B$.
        \begin{subproof}
            Suppose not.
            Then $y\in B$.
            Thus $x\mathrel{f} y$ by \cref{function_apply_iff}.
            Thus $x\in\preimg{f}{B}$ by \cref{preimg_iff}.
            Contradiction.
        \end{subproof}
        Thus $y\in Y\setminus B$ by \cref{setminus_intro}.
        Thus $x\mathrel{f} y$ by \cref{function_apply_iff}.
        Thus $x\in\preimg{f}{Y\setminus B}$ by \cref{preimg_iff}.
    \end{subproof}
    Thus $X\setminus \preimg{f}{B}\subseteq \preimg{f}{Y\setminus B}$ by \cref{subseteq}.
    Thus $\preimg{f}{Y\setminus B} = X\setminus \preimg{f}{B}$ by \cref{subseteq_antisymmetric}.
\end{proof}

% from §18 Item 2: continuity implies image of closure is contained in closure of image
\begin{lemma}\label{continuous_closure_image_subset}
    Assume $f$ is continuous from $X$ to $Y$ with respect to $T$ and $T_1$.
    Suppose $A\subseteq X$.
    Then $\img{f}{\closure{A}{X}{T}}\subseteq \closure{\img{f}{A}}{Y}{T_1}$.
\end{lemma}
\begin{proof}
    Thus $f\in\funs{X}{Y}$ by \cref{continuous}.
    Thus $f$ is a function by \cref{funs_is_function}.
    Thus $T$ is a topology on $X$ by \cref{continuous}.
    Thus $T_1$ is a topology on $Y$ by \cref{continuous}.
    Thus $f\in\rels{X}{Y}$ by \cref{funs_subseteq_rels,subseteq}.
    Thus $\ran{f}\subseteq Y$ by \cref{rels_ran_subseteq}.
    Thus $\img{f}{A}\subseteq \ran{f}$ by \cref{img_subseteq_ran}.
    Thus $\img{f}{A}\subseteq Y$ by \cref{subseteq_transitive}.
    Show that for all $y$ such that $y\in\img{f}{\closure{A}{X}{T}}$ we have
        $y\in\closure{\img{f}{A}}{Y}{T_1}$.
    \begin{subproof}
        Fix $y$.
        Assume $y\in\img{f}{\closure{A}{X}{T}}$.
        Take $x\in\dom{f}\inter \closure{A}{X}{T}$ such that $y = f(x)$ by \cref{img_of_function_elim}.
        Thus $x\in\dom{f}$ by \cref{inter_elim_left}.
        Thus $x\mathrel{f} y$ by \cref{function_apply_elim}.
        Thus $y\in Y$ by \cref{rels_type_ran}.
        Thus $x\in\closure{A}{X}{T}$ by \cref{inter_elim_right}.
        Thus $x\in X$ by \cref{closure}.
        Show that for all $V$ such that $V$ is open in $Y$ with respect to $T_1$ and $y\in V$
            we have $V$ intersects $\img{f}{A}$.
        \begin{subproof}
            Fix $V$.
            Assume $V$ is open in $Y$ with respect to $T_1$ and $y\in V$.
            Let $U = \preimg{f}{V}$.
            $U$ is open in $X$ with respect to $T$ by \cref{continuous}.
            $x\mathrel{f} y$ by \cref{function_apply_elim,inter_elim_left}.
            Thus $x\in U$ by \cref{preimg_iff}.
            Show that $U$ intersects $A$.
            \begin{subproof}
                Suppose not.
                Then $U\inter A = \emptyset$ by \cref{intersects}.
                Thus $A\inter U = \emptyset$ by \cref{inter_comm}.
                Thus $A\subseteq X$.
                Thus $A\subseteq X\setminus U$ by \cref{subseteq_setminus}.
                $X\setminus U\subseteq X$ by \cref{setminus_subseteq}.
                $X\setminus (X\setminus U) = X\inter U$ by \cref{setminus_setminus}.
                $U\in T$ by \cref{open_in}.
                $T\subseteq\pow{X}$ by \cref{topology_on}.
                Thus $U\in\pow{X}$ by \cref{subseteq}.
                Thus $U\subseteq X$ by \cref{pow_iff}.
                Thus $X\inter U = U$ by \cref{inter_absorb_supseteq_left}.
                Thus $X\setminus (X\setminus U) = U$ by \cref{setminus_setminus,inter_absorb_supseteq_left}.
                Thus $X\setminus U$ is closed in $X$ with respect to $T$ by \cref{closed_in}.
                Show that for all $y$ such that $y\in A$ we have $y\in X\setminus U$.
                \begin{subproof}
                    Fix $y$.
                    Assume $y\in A$.
                    Thus $y\in X\setminus U$ by \cref{subseteq}.
                \end{subproof}
                Thus $x\in X\setminus U$ by \cref{closure}.
                Contradiction by \cref{setminus_elim_right}.
            \end{subproof}
            Thus $U\inter A \neq \emptyset$ by \cref{intersects}.
            Take $z$ such that $z\in U\inter A$ by \cref{nonempty_inhabited}.
            Thus $z\in U$ and $z\in A$ by \cref{inter}.
            Take $w\in V$ such that $z\mathrel{f} w$ by \cref{preimg_iff}.
            Thus $z\in\dom{f}$ and $w = f(z)$ by \cref{function_apply_iff}.
            Thus $z\in\dom{f}\inter A$ by \cref{inter_intro}.
            Thus $w\in\img{f}{A}$ by \cref{img_of_function_intro}.
            Thus $w\in V\inter\img{f}{A}$ by \cref{inter_intro}.
            Show that $V\inter\img{f}{A} \neq \emptyset$.
            \begin{subproof}
                Suppose not.
                Then $V\inter\img{f}{A} = \emptyset$.
                Thus $w\in\emptyset$.
                Contradiction by \cref{emptyset}.
            \end{subproof}
            Thus $V$ intersects $\img{f}{A}$ by \cref{intersects}.
        \end{subproof}
        Thus $y\in\closure{\img{f}{A}}{Y}{T_1}$ by \cref{closure_iff_intersects_open}.
    \end{subproof}
    Thus $\img{f}{\closure{A}{X}{T}}\subseteq \closure{\img{f}{A}}{Y}{T_1}$ by \cref{subseteq}.
\end{proof}

% from §18 Item 3: closure condition implies preimages of closed sets are closed
\begin{lemma}\label{closure_image_subset_implies_preimg_closed}
    Assume $f$ is a function from $X$ to $Y$.
    Assume $T$ is a topology on $X$.
    Assume $T_1$ is a topology on $Y$.
    Suppose for all $A$ such that $A\subseteq X$ we have
        $\img{f}{\closure{A}{X}{T}}\subseteq \closure{\img{f}{A}}{Y}{T_1}$.
    Then for all $B$ such that $B$ is closed in $Y$ with respect to $T_1$
        we have $\preimg{f}{B}$ is closed in $X$ with respect to $T$.
\end{lemma}
\begin{proof}
    Fix $B$.
    Assume $B$ is closed in $Y$ with respect to $T_1$.
    Let $A = \preimg{f}{B}$.
    Thus $f\in\funs{X}{Y}$.
    Thus $f$ is a function by \cref{funs_is_function}.
    Thus $f$ is right-unique by \cref{function_rightunique}.
    Thus $f$ is a relation by \cref{funs_is_relation}.
    Thus $f\in\rels{X}{Y}$ by \cref{funs_subseteq_rels,subseteq}.
    Thus $\ran{f}\subseteq Y$ by \cref{rels_ran_subseteq}.
    Thus $\dom{f} = X$ by \cref{funs}.
    $A\subseteq X$ by \cref{preimg_subseteq_dom}.
    Show that $\closure{A}{X}{T}\subseteq A$.
    \begin{subproof}
        Show that for all $x$ such that $x\in\closure{A}{X}{T}$ we have $x\in A$.
        \begin{subproof}
            Fix $x$.
            Assume $x\in\closure{A}{X}{T}$.
            Thus $x\in X$ by \cref{closure}.
            Thus $x\in\dom{f}$ by \cref{funs}.
            Thus $x\in\dom{f}\inter\closure{A}{X}{T}$ by \cref{inter_intro}.
            Thus $f(x)\in\img{f}{\closure{A}{X}{T}}$ by \cref{img_of_function_intro}.
            Thus $f(x)\in\closure{\img{f}{A}}{Y}{T_1}$ by \cref{subseteq}.
            $\img{f}{A}\subseteq B$ by \cref{img_preimg_subseteq}.
            Thus $\img{f}{A}\subseteq \ran{f}$ by \cref{img_subseteq_ran}.
            Thus $\img{f}{A}\subseteq Y$ by \cref{subseteq_transitive}.
            Thus $\closure{\img{f}{A}}{Y}{T_1}\subseteq B$ by \cref{closure_subseteq_closed}.
            Thus $f(x)\in B$ by \cref{subseteq}.
            Thus $x\mathrel{f} f(x)$ by \cref{function_apply_elim}.
            Thus $x\in\preimg{f}{B}$ by \cref{preimg_iff}.
            Thus $x\in A$.
        \end{subproof}
        Thus $\closure{A}{X}{T}\subseteq A$ by \cref{subseteq}.
    \end{subproof}
    $A\subseteq\closure{A}{X}{T}$ by \cref{subseteq_closure}.
    Thus $A = \closure{A}{X}{T}$ by \cref{subseteq_antisymmetric}.
    $\closure{A}{X}{T}$ is closed in $X$ with respect to $T$ by \cref{closure_closed_in}.
    Thus $A$ is closed in $X$ with respect to $T$.
\end{proof}

% from §18 Item 4: preimages of closed sets imply continuity
\begin{lemma}\label{preimg_closed_implies_continuous}
    Assume $f$ is a function from $X$ to $Y$.
    Assume $T$ is a topology on $X$.
    Assume $T_1$ is a topology on $Y$.
    Suppose for all $B$ such that $B$ is closed in $Y$ with respect to $T_1$
        we have $\preimg{f}{B}$ is closed in $X$ with respect to $T$.
    Then $f$ is continuous from $X$ to $Y$ with respect to $T$ and $T_1$.
\end{lemma}
\begin{proof}
    Thus $f\in\funs{X}{Y}$.
    Thus $\dom{f} = X$ by \cref{funs}.
    Show that for all $V$ such that $V$ is open in $Y$ with respect to $T_1$
        we have $\preimg{f}{V}$ is open in $X$ with respect to $T$.
    \begin{subproof}
        Fix $V$.
        Assume $V$ is open in $Y$ with respect to $T_1$.
        Let $B = Y\setminus V$.
        $B\subseteq Y$ by \cref{setminus_subseteq}.
        $Y\setminus B = Y\inter V$ by \cref{setminus_setminus}.
        $V\in T_1$ by \cref{open_in}.
        $T_1\subseteq\pow{Y}$ by \cref{topology_on}.
        Thus $V\in\pow{Y}$ by \cref{subseteq}.
        Thus $V\subseteq Y$ by \cref{pow_iff}.
        Thus $Y\inter V = V$ by \cref{inter_absorb_supseteq_left}.
        Thus $Y\setminus B = V$.
        Thus $B$ is closed in $Y$ with respect to $T_1$ by \cref{closed_in}.
        Thus $\preimg{f}{B}$ is closed in $X$ with respect to $T$ by \cref{subseteq}.
        $X\setminus \preimg{f}{B}$ is open in $X$ with respect to $T$ by \cref{closed_in}.
        Thus $V = Y\setminus B$.
        $f$ is a function by \cref{funs_is_function}.
        Thus $\preimg{f}{V} = X\setminus \preimg{f}{B}$ by \cref{preimg_setminus_function,subseteq}.
        Thus $\preimg{f}{V}$ is open in $X$ with respect to $T$ by \cref{subseteq}.
    \end{subproof}
    Thus $f$ is continuous from $X$ to $Y$ with respect to $T$ and $T_1$ by \cref{continuous}.
\end{proof}

% from §18 Item 5: continuity at a point
\begin{definition}\label{continuous_at}
    $f$ is continuous at $x$ from $X$ to $Y$ with respect to $T$ and $T_1$ iff
        $f$ is a function from $X$ to $Y$ and
        $T$ is a topology on $X$ and
        $T_1$ is a topology on $Y$ and
        $x\in X$ and
        for all $V$ such that $V$ is a neighborhood of $f(x)$ in $Y$ with respect to $T_1$
            there exists $U$ such that
                $U$ is a neighborhood of $x$ in $X$ with respect to $T$ and
                $\img{f}{U}\subseteq V$.
\end{definition}

% from §18 Item 6: continuity implies continuity at a point
\begin{lemma}\label{continuous_implies_continuous_at}
    Assume $f$ is continuous from $X$ to $Y$ with respect to $T$ and $T_1$.
    Suppose $x\in X$.
    Then $f$ is continuous at $x$ from $X$ to $Y$ with respect to $T$ and $T_1$.
\end{lemma}
\begin{proof}
    Show that for all $V$ such that $V$ is a neighborhood of $f(x)$ in $Y$ with respect to $T_1$
        there exists $U$ such that
            $U$ is a neighborhood of $x$ in $X$ with respect to $T$ and
            $\img{f}{U}\subseteq V$.
    \begin{subproof}
        Fix $V$.
        Assume $V$ is a neighborhood of $f(x)$ in $Y$ with respect to $T_1$.
        $V$ is open in $Y$ with respect to $T_1$ by \cref{neighborhood}.
        Let $U = \preimg{f}{V}$.
        $U$ is open in $X$ with respect to $T$ by \cref{continuous}.
        $f(x)\in V$ by \cref{neighborhood}.
        Thus $f\in\funs{X}{Y}$ by \cref{continuous}.
        Thus $f$ is a function by \cref{funs_is_function}.
        Thus $f$ is right-unique by \cref{function_rightunique}.
        Thus $f$ is a relation by \cref{funs_is_relation}.
        Thus $x\in\dom{f}$ by \cref{funs}.
        Thus $x\mathrel{f} f(x)$ by \cref{function_apply_elim}.
        Thus $x\in U$ by \cref{preimg_iff}.
        Thus $U$ is a neighborhood of $x$ in $X$ with respect to $T$ by \cref{neighborhood}.
        $\img{f}{U}\subseteq V$ by \cref{img_preimg_subseteq}.
    \end{subproof}
    Thus $f\in\funs{X}{Y}$ by \cref{continuous}.
    Thus $T$ is a topology on $X$ by \cref{continuous}.
    Thus $T_1$ is a topology on $Y$ by \cref{continuous}.
    Thus $f$ is continuous at $x$ from $X$ to $Y$ with respect to $T$ and $T_1$ by \cref{continuous_at}.
\end{proof}

% from §18 Item 7: continuity at each point implies continuity
\begin{lemma}\label{continuous_at_implies_continuous}
    Assume $f$ is a function from $X$ to $Y$.
    Assume $T$ is a topology on $X$.
    Assume $T_1$ is a topology on $Y$.
    Suppose for all $x\in X$ we have $f$ is continuous at $x$ from $X$ to $Y$ with respect to $T$ and $T_1$.
    Then $f$ is continuous from $X$ to $Y$ with respect to $T$ and $T_1$.
\end{lemma}
\begin{proof}
    Thus $f\in\funs{X}{Y}$.
    Thus $\dom{f} = X$ by \cref{funs}.
    Show that for all $V$ such that $V$ is open in $Y$ with respect to $T_1$
        we have $\preimg{f}{V}$ is open in $X$ with respect to $T$.
    \begin{subproof}
        Fix $V$.
        Assume $V$ is open in $Y$ with respect to $T_1$.
        Let $M = \{U\in T \mid U\subseteq \preimg{f}{V}\}$.
        Show that for all $x$ such that $x\in\preimg{f}{V}$ we have $x\in\unions{M}$.
        \begin{subproof}
            Fix $x$.
            Assume $x\in\preimg{f}{V}$.
            $\preimg{f}{V}\subseteq \dom{f}$ by \cref{preimg_subseteq_dom}.
            Thus $x\in\dom{f}$ by \cref{subseteq}.
            Thus $x\in X$ by \cref{funs}.
            $f$ is a function by \cref{funs_is_function}.
            Take $y\in V$ such that $x\mathrel{f} y$ by \cref{preimg_iff}.
            Thus $y = f(x)$ by \cref{function_apply_iff}.
            Thus $f(x)\in V$.
            Thus $V$ is a neighborhood of $f(x)$ in $Y$ with respect to $T_1$ by \cref{neighborhood}.
            $f$ is continuous at $x$ from $X$ to $Y$ with respect to $T$ and $T_1$ by \cref{subseteq}.
            Take $U$ such that
                $U$ is a neighborhood of $x$ in $X$ with respect to $T$ and
                $\img{f}{U}\subseteq V$ by \cref{continuous_at}.
            $U$ is open in $X$ with respect to $T$ by \cref{neighborhood}.
            Thus $U\in T$ by \cref{open_in}.
            Show that $U\subseteq \preimg{f}{V}$.
            \begin{subproof}
                Show that for all $z$ such that $z\in U$ we have $z\in\preimg{f}{V}$.
                \begin{subproof}
                Fix $z$.
                Assume $z\in U$.
                $T\subseteq\pow{X}$ by \cref{topology_on}.
                Thus $U\in\pow{X}$ by \cref{subseteq}.
                Thus $U\subseteq X$ by \cref{pow_iff}.
                Thus $z\in X$ by \cref{subseteq}.
                Thus $z\in\dom{f}$ by \cref{funs}.
                Thus $z\in\dom{f}\inter U$ by \cref{inter_intro}.
                Thus $f(z)\in\img{f}{U}$ by \cref{img_of_function_intro}.
                Thus $f(z)\in V$ by \cref{subseteq}.
                Thus $z\mathrel{f} f(z)$ by \cref{function_apply_elim}.
                Thus $z\in\preimg{f}{V}$ by \cref{preimg_iff}.
                \end{subproof}
                Thus $U\subseteq \preimg{f}{V}$ by \cref{subseteq}.
            \end{subproof}
            Thus $U\in M$.
            $x\in U$ by \cref{neighborhood}.
            Thus $x\in\unions{M}$ by \cref{unions_iff}.
        \end{subproof}
        Thus $\preimg{f}{V}\subseteq \unions{M}$ by \cref{subseteq}.
        $M\subseteq T$.
        Thus $\unions{M}\in T$ by \cref{topology_on}.
        Show that $M\subseteq\pow{\preimg{f}{V}}$.
        \begin{subproof}
            Show that for all $U$ such that $U\in M$ we have $U\in\pow{\preimg{f}{V}}$.
            \begin{subproof}
                Fix $U$.
                Assume $U\in M$.
                Thus $U\subseteq \preimg{f}{V}$ by \cref{subseteq}.
                Thus $U\in\pow{\preimg{f}{V}}$ by \cref{pow_iff}.
            \end{subproof}
            Thus $M\subseteq\pow{\preimg{f}{V}}$ by \cref{subseteq}.
        \end{subproof}
        $\unions{M}\subseteq \preimg{f}{V}$ by \cref{unions_subseteq_of_powerset_is_subseteq}.
        Thus $\preimg{f}{V} = \unions{M}$ by \cref{subseteq_antisymmetric}.
        Thus $\preimg{f}{V}$ is open in $X$ with respect to $T$ by \cref{open_in}.
    \end{subproof}
    Thus $f$ is continuous from $X$ to $Y$ with respect to $T$ and $T_1$ by \cref{continuous}.
\end{proof}

% from §18 Item 8: homeomorphism
\begin{definition}\label{homeomorphism}
    $f$ is a homeomorphism between $X$ and $Y$ with respect to $T$ and $T_1$ iff
        $f$ is a bijection from $X$ to $Y$ and
        $f$ is continuous from $X$ to $Y$ with respect to $T$ and $T_1$ and
        $\converse{f}$ is continuous from $Y$ to $X$ with respect to $T_1$ and $T$.
\end{definition}

% from §18 Item 9: topological property
\begin{definition}\label{topological_property}
    $P$ is a topological property iff
        for all $X,Y,T,T_1,f$ such that
            $f$ is a homeomorphism between $X$ and $Y$ with respect to $T$ and $T_1$
        we have $(X,T)\in P$ iff $(Y,T_1)\in P$.
\end{definition}

% from §18 Item 10: topological embedding
\begin{definition}\label{topological_embedding}
    $f$ is a topological embedding from $X$ to $Y$ with respect to $T$ and $T_1$ iff
        $f$ is a function from $X$ to $Y$ and
        $f$ is injective and
        $f$ is continuous from $X$ to $Y$ with respect to $T$ and $T_1$ and
        $f$ is a homeomorphism between $X$ and $\img{f}{X}$ with respect to
            $T$ and $\subspaceTopology{T_1}{\img{f}{X}}$.
\end{definition}

% from §18 Item 11: constant function is continuous
\begin{lemma}\label{constant_continuous}
    Assume $T$ is a topology on $X$.
    Assume $T_1$ is a topology on $Y$.
    Suppose $f$ is a function from $X$ to $Y$.
    Suppose there exists $y_0\in Y$ such that for all $x\in X$ we have $f(x) = y_0$.
    Then $f$ is continuous from $X$ to $Y$ with respect to $T$ and $T_1$.
\end{lemma}
\begin{proof}
    Show that for all $V$ such that $V$ is open in $Y$ with respect to $T_1$
        we have $\preimg{f}{V}$ is open in $X$ with respect to $T$.
    \begin{subproof}
        Fix $V$.
        Assume $V$ is open in $Y$ with respect to $T_1$.
        Take $y_0\in Y$ such that for all $x\in X$ we have $f(x) = y_0$.
        \begin{byCase}
        \caseOf{$y_0\in V$.}
            Show that $\preimg{f}{V} = X$.
            \begin{subproof}
                Show that for all $x$ such that $x\in X$ we have $x\in\preimg{f}{V}$.
                \begin{subproof}
                    Fix $x$.
                    Assume $x\in X$.
                    Thus $f(x) = y_0$.
                    Thus $f(x)\in V$.
                    Thus $f\in\funs{X}{Y}$.
                    Thus $f$ is a function by \cref{funs_is_function}.
                    Thus $f$ is right-unique by \cref{function_rightunique}.
                    Thus $f$ is a relation by \cref{funs_is_relation}.
                    Thus $x\in\dom{f}$ by \cref{funs}.
                    Thus $x\mathrel{f} f(x)$ by \cref{function_apply_elim}.
                    Thus $x\in\preimg{f}{V}$ by \cref{preimg_iff}.
                \end{subproof}
                Thus $X\subseteq \preimg{f}{V}$ by \cref{subseteq}.
                $\preimg{f}{V}\subseteq \dom{f}$ by \cref{preimg_subseteq_dom}.
                Thus $\dom{f} = X$ by \cref{funs}.
                Thus $\preimg{f}{V}\subseteq X$ by \cref{subseteq_transitive}.
                Thus $\preimg{f}{V} = X$ by \cref{subseteq_antisymmetric}.
            \end{subproof}
            Thus $\preimg{f}{V}$ is open in $X$ with respect to $T$ by \cref{open_in,topology_on}.
        \caseOf{$y_0\notin V$.}
            Show that $\preimg{f}{V} = \emptyset$.
            \begin{subproof}
                Show that for all $x$ we have $x\notin\preimg{f}{V}$.
                \begin{subproof}
                    Fix $x$.
                    Suppose not.
                    Then $x\in\preimg{f}{V}$.
                    Take $y\in V$ such that $x\mathrel{f} y$ by \cref{preimg_iff}.
                    Thus $f\in\funs{X}{Y}$.
                    Thus $f$ is a function by \cref{funs_is_function}.
                    Thus $f$ is right-unique by \cref{function_rightunique}.
                    Thus $f$ is a relation by \cref{funs_is_relation}.
                    $\preimg{f}{V}\subseteq \dom{f}$ by \cref{preimg_subseteq_dom}.
                    Thus $x\in\dom{f}$ by \cref{subseteq}.
                    Thus $y = f(x)$ by \cref{function_apply_iff}.
                    Thus $y = y_0$.
                    Thus $y_0\in V$.
                    Contradiction.
                \end{subproof}
                Thus $\preimg{f}{V}\subseteq \emptyset$ by \cref{subseteq}.
                Thus $\preimg{f}{V} = \emptyset$ by \cref{emptyset_subseteq,subseteq_antisymmetric}.
            \end{subproof}
            Thus $\preimg{f}{V}$ is open in $X$ with respect to $T$ by \cref{open_in,topology_on}.
        \end{byCase}
    \end{subproof}
    Thus $f$ is continuous from $X$ to $Y$ with respect to $T$ and $T_1$ by \cref{continuous}.
\end{proof}

% from §18 Item 12: inclusion is continuous
\begin{lemma}\label{inclusion_continuous}
    Assume $A$ is a subspace of $X$ with respect to $T$.
    Let $T_A = \subspaceTopology{T}{A}$.
    Let $j = \identity{A}$.
    Then $j$ is continuous from $A$ to $X$ with respect to $T_A$ and $T$.
\end{lemma}
\begin{proof}
    Thus $T_A$ is a topology on $A$ by \cref{subspace_topology_is_topology,subspace_of}.
    Thus $T$ is a topology on $X$ by \cref{subspace_of}.
    Thus $A\subseteq X$ by \cref{subspace_of}.
    Thus $j\in\funs{A}{A}$ by \cref{id_elem_funs}.
    Thus $j\in\funs{A}{X}$ by \cref{funs_weaken_codom}.
    Show that for all $U$ such that $U$ is open in $X$ with respect to $T$
        we have $\preimg{j}{U}$ is open in $A$ with respect to $T_A$.
    \begin{subproof}
        Fix $U$.
        Assume $U$ is open in $X$ with respect to $T$.
        Show that $\preimg{j}{U} = A\inter U$.
        \begin{subproof}
            Show that for all $x$ such that $x\in\preimg{j}{U}$ we have $x\in A\inter U$.
            \begin{subproof}
                Fix $x$.
                Assume $x\in\preimg{j}{U}$.
                Take $y\in U$ such that $x\mathrel{j} y$ by \cref{preimg_iff}.
                Thus $x = y$ and $x\in A$ by \cref{id_iff}.
                Thus $x\in A\inter U$ by \cref{inter_intro}.
            \end{subproof}
            Thus $\preimg{j}{U}\subseteq A\inter U$ by \cref{subseteq}.
            Show that for all $x$ such that $x\in A\inter U$ we have $x\in\preimg{j}{U}$.
            \begin{subproof}
                Fix $x$.
                Assume $x\in A\inter U$.
                Thus $x\in A$ and $x\in U$ by \cref{inter}.
                Thus $x\mathrel{j} x$ by \cref{id_iff}.
                Thus $x\in\preimg{j}{U}$ by \cref{preimg_iff}.
            \end{subproof}
            Thus $A\inter U\subseteq \preimg{j}{U}$ by \cref{subseteq}.
            Thus $\preimg{j}{U} = A\inter U$ by \cref{subseteq_antisymmetric}.
        \end{subproof}
        Thus $\preimg{j}{U}\in T_A$ by \cref{subspace_topology,open_in}.
        Thus $\preimg{j}{U}$ is open in $A$ with respect to $T_A$ by \cref{open_in}.
    \end{subproof}
    Thus $j$ is continuous from $A$ to $X$ with respect to $T_A$ and $T$ by \cref{continuous}.
\end{proof}

% from §18 Item 13: composites of continuous functions are continuous
\begin{lemma}\label{continuous_composite}
    Assume $f$ is continuous from $X$ to $Y$ with respect to $T$ and $T_1$.
    Assume $g$ is continuous from $Y$ to $Z$ with respect to $T_1$ and $T_2$.
    Then $g\circ f$ is continuous from $X$ to $Z$ with respect to $T$ and $T_2$.
\end{lemma}
\begin{proof}
    Show that for all $U$ such that $U$ is open in $Z$ with respect to $T_2$
        we have $\preimg{(g\circ f)}{U}$ is open in $X$ with respect to $T$.
    \begin{subproof}
        Fix $U$.
        Assume $U$ is open in $Z$ with respect to $T_2$.
        $\preimg{g}{U}$ is open in $Y$ with respect to $T_1$ by \cref{continuous}.
        $\preimg{f}{\preimg{g}{U}}$ is open in $X$ with respect to $T$ by \cref{continuous}.
        Show that $\preimg{(g\circ f)}{U} = \preimg{f}{\preimg{g}{U}}$.
        \begin{subproof}
            Show that for all $x$ such that $x\in\preimg{(g\circ f)}{U}$ we have
                $x\in\preimg{f}{\preimg{g}{U}}$.
            \begin{subproof}
                Fix $x$.
                Assume $x\in\preimg{(g\circ f)}{U}$.
                Take $z\in U$ such that $x\mathrel{(g\circ f)} z$ by \cref{preimg_iff}.
                Take $y$ such that $x\mathrel{f} y$ and $y\mathrel{g} z$ by \cref{circ_iff}.
                Thus $y\in\preimg{g}{U}$ by \cref{preimg_iff}.
                Thus $x\in\preimg{f}{\preimg{g}{U}}$ by \cref{preimg_iff}.
            \end{subproof}
            Thus $\preimg{(g\circ f)}{U}\subseteq \preimg{f}{\preimg{g}{U}}$ by \cref{subseteq}.
            Show that for all $x$ such that $x\in\preimg{f}{\preimg{g}{U}}$ we have
                $x\in\preimg{(g\circ f)}{U}$.
            \begin{subproof}
                Fix $x$.
                Assume $x\in\preimg{f}{\preimg{g}{U}}$.
                Take $y\in\preimg{g}{U}$ such that $x\mathrel{f} y$ by \cref{preimg_iff}.
                Take $z\in U$ such that $y\mathrel{g} z$ by \cref{preimg_iff}.
                Thus $x\mathrel{(g\circ f)} z$ by \cref{circ_iff}.
                Thus $x\in\preimg{(g\circ f)}{U}$ by \cref{preimg_iff}.
            \end{subproof}
            Thus $\preimg{f}{\preimg{g}{U}}\subseteq \preimg{(g\circ f)}{U}$ by \cref{subseteq}.
            Thus $\preimg{(g\circ f)}{U} = \preimg{f}{\preimg{g}{U}}$ by \cref{subseteq_antisymmetric}.
        \end{subproof}
        Thus $\preimg{(g\circ f)}{U}$ is open in $X$ with respect to $T$.
    \end{subproof}
    Thus $f\in\funs{X}{Y}$ by \cref{continuous}.
    Thus $g\in\funs{Y}{Z}$ by \cref{continuous}.
    Thus $g\circ f\in\funs{X}{Z}$ by \cref{funs_circ}.
    Thus $T$ is a topology on $X$ by \cref{continuous}.
    Thus $T_2$ is a topology on $Z$ by \cref{continuous}.
    Thus $g\circ f$ is continuous from $X$ to $Z$ with respect to $T$ and $T_2$ by \cref{continuous}.
\end{proof}

% from §18 Item 14: restricting the domain
\begin{lemma}\label{continuous_restrict_domain}
    Assume $f$ is continuous from $X$ to $Y$ with respect to $T$ and $T_1$.
    Suppose $A$ is a subspace of $X$ with respect to $T$.
    Let $T_A = \subspaceTopology{T}{A}$.
    Let $g = \restrl{f}{A}$.
    Then $g$ is continuous from $A$ to $Y$ with respect to $T_A$ and $T_1$.
\end{lemma}
\begin{proof}
    Thus $T_A$ is a topology on $A$ by \cref{subspace_topology_is_topology,subspace_of}.
    Show that for all $V$ such that $V$ is open in $Y$ with respect to $T_1$
        we have $\preimg{g}{V}$ is open in $A$ with respect to $T_A$.
    \begin{subproof}
        Fix $V$.
        Assume $V$ is open in $Y$ with respect to $T_1$.
        Show that $\preimg{g}{V} = A\inter \preimg{f}{V}$.
        \begin{subproof}
            Show that for all $x$ such that $x\in\preimg{g}{V}$ we have
                $x\in A\inter \preimg{f}{V}$.
            \begin{subproof}
                Fix $x$.
                Assume $x\in\preimg{g}{V}$.
                Take $y\in V$ such that $x\mathrel{g} y$ by \cref{preimg_iff}.
                Thus $x\in A$ and $x\mathrel{f} y$ by \cref{restrl_iff}.
                Thus $x\in\preimg{f}{V}$ by \cref{preimg_iff}.
                Thus $x\in A\inter \preimg{f}{V}$ by \cref{inter_intro}.
            \end{subproof}
            Thus $\preimg{g}{V}\subseteq A\inter \preimg{f}{V}$ by \cref{subseteq}.
            Show that for all $x$ such that $x\in A\inter \preimg{f}{V}$ we have
                $x\in\preimg{g}{V}$.
            \begin{subproof}
                Fix $x$.
                Assume $x\in A\inter \preimg{f}{V}$.
                Thus $x\in A$ and $x\in\preimg{f}{V}$ by \cref{inter}.
                Take $y\in V$ such that $x\mathrel{f} y$ by \cref{preimg_iff}.
                Thus $x\mathrel{g} y$ by \cref{restrl_iff}.
                Thus $x\in\preimg{g}{V}$ by \cref{preimg_iff}.
            \end{subproof}
            Thus $A\inter \preimg{f}{V}\subseteq \preimg{g}{V}$ by \cref{subseteq}.
            Thus $\preimg{g}{V} = A\inter \preimg{f}{V}$ by \cref{subseteq_antisymmetric}.
        \end{subproof}
        $\preimg{f}{V}$ is open in $X$ with respect to $T$ by \cref{continuous}.
        Thus $\preimg{f}{V}\in T$ by \cref{open_in}.
        Thus $A\inter \preimg{f}{V}\in T_A$ by \cref{subspace_topology}.
        Thus $A\inter \preimg{f}{V}$ is open in $A$ with respect to $T_A$ by \cref{open_in}.
        Thus $\preimg{g}{V}$ is open in $A$ with respect to $T_A$.
    \end{subproof}
    Thus $f\in\funs{X}{Y}$ by \cref{continuous}.
    Thus $f$ is a function by \cref{funs_is_function}.
    Thus $g$ is a function by \cref{restrl_of_function_is_function}.
    Thus $g$ is right-unique and $g$ is a relation.
    Thus $\dom{g} = \dom{f}\inter A$ by \cref{dom_of_restrl_eq_inter}.
    Thus $\dom{g} = A$ by \cref{funs,inter_absorb_supseteq_left,subspace_of}.
    Thus $A\subseteq X$ by \cref{subspace_of}.
    Thus $\dom{f} = X$ by \cref{funs}.
    Thus $A\subseteq \dom{f}$ by \cref{subseteq_transitive}.
    $\ran{g} = \img{f}{A}$ by \cref{restrl_ran}.
    $\img{f}{A}\subseteq \ran{f}$ by \cref{img_subseteq_ran}.
    Thus $f\in\rels{X}{Y}$ by \cref{funs_subseteq_rels,subseteq}.
    Thus $\ran{f}\subseteq Y$ by \cref{rels_ran_subseteq}.
    Thus $\ran{g}\subseteq Y$ by \cref{subseteq_transitive}.
    Show that for all $x\in\dom{g}$ we have $g(x)\in Y$.
    \begin{subproof}
        Fix $x\in\dom{g}$.
        Thus $(x,g(x))\in g$ by \cref{function_apply_elim}.
        Thus $g(x)\in\ran{g}$ by \cref{ran_iff}.
        Thus $g(x)\in Y$ by \cref{subseteq}.
    \end{subproof}
    Thus $g\in\funs{A}{Y}$ by \cref{funs_intro}.
    Thus $T_1$ is a topology on $Y$ by \cref{continuous}.
    Thus $g$ is continuous from $A$ to $Y$ with respect to $T_A$ and $T_1$ by \cref{continuous}.
\end{proof}

% from §18 Item 15: restricting the range
\begin{lemma}\label{continuous_restrict_range}
    Assume $f$ is continuous from $X$ to $Y$ with respect to $T$ and $T_1$.
    Suppose $Z$ is a subspace of $Y$ with respect to $T_1$.
    Suppose $\img{f}{X}\subseteq Z$.
    Let $T_Z = \subspaceTopology{T_1}{Z}$.
    Then $f$ is continuous from $X$ to $Z$ with respect to $T$ and $T_Z$.
\end{lemma}
\begin{proof}
    Thus $f\in\funs{X}{Y}$.
    Thus $\dom{f} = X$ by \cref{funs}.
    Thus $T_1$ is a topology on $Y$ by \cref{continuous}.
    Thus $T_Z$ is a topology on $Z$ by \cref{subspace_topology_is_topology,subspace_of}.
    Show that for all $B$ such that $B$ is open in $Z$ with respect to $T_Z$
        we have $\preimg{f}{B}$ is open in $X$ with respect to $T$.
    \begin{subproof}
        Fix $B$.
        Assume $B$ is open in $Z$ with respect to $T_Z$.
        Take $U\in T_1$ such that $B = Z\inter U$ by \cref{subspace_topology,open_in}.
        $U$ is open in $Y$ with respect to $T_1$ by \cref{open_in}.
        $\preimg{f}{U}$ is open in $X$ with respect to $T$ by \cref{continuous}.
        Show that $\preimg{f}{B} = \preimg{f}{U}$.
        \begin{subproof}
            Show that for all $x$ such that $x\in\preimg{f}{B}$ we have $x\in\preimg{f}{U}$.
            \begin{subproof}
                Fix $x$.
                Assume $x\in\preimg{f}{B}$.
                Take $y\in B$ such that $x\mathrel{f} y$ by \cref{preimg_iff}.
                Thus $y\in U$ by \cref{inter_elim_right}.
                Thus $x\in\preimg{f}{U}$ by \cref{preimg_iff}.
            \end{subproof}
            Thus $\preimg{f}{B}\subseteq \preimg{f}{U}$ by \cref{subseteq}.
            Show that for all $x$ such that $x\in\preimg{f}{U}$ we have $x\in\preimg{f}{B}$.
            \begin{subproof}
                Fix $x$.
                Assume $x\in\preimg{f}{U}$.
                Take $y\in U$ such that $x\mathrel{f} y$ by \cref{preimg_iff}.
                $\preimg{f}{U}\subseteq \dom{f}$ by \cref{preimg_subseteq_dom}.
                Thus $x\in\dom{f}$ by \cref{subseteq}.
                Thus $x\in X$ by \cref{funs}.
                Thus $f$ is a function by \cref{funs_is_function}.
                Thus $f$ is right-unique by \cref{function_rightunique}.
                Thus $f$ is a relation by \cref{funs_is_relation}.
                Thus $y = f(x)$ by \cref{function_apply_iff}.
                Thus $y\in\img{f}{X}$ by \cref{img_of_function_intro,inter_intro}.
                Thus $y\in Z$ by \cref{subseteq}.
                Thus $y\in Z\inter U$ by \cref{inter_intro}.
                Thus $y\in B$.
                Thus $x\in\preimg{f}{B}$ by \cref{preimg_iff}.
            \end{subproof}
            Thus $\preimg{f}{U}\subseteq \preimg{f}{B}$ by \cref{subseteq}.
            Thus $\preimg{f}{B} = \preimg{f}{U}$ by \cref{subseteq_antisymmetric}.
        \end{subproof}
        Thus $\preimg{f}{B}$ is open in $X$ with respect to $T$.
    \end{subproof}
    Thus $f$ is a function by \cref{funs_is_function}.
    Thus $\img{f}{\dom{f}} = \ran{f}$ by \cref{img_dom}.
    Thus $\ran{f} = \img{f}{X}$ by \cref{funs}.
    Thus $\ran{f}\subseteq Z$ by \cref{subseteq}.
    Show that for all $x\in\dom{f}$ we have $f(x)\in Z$.
    \begin{subproof}
        Fix $x\in\dom{f}$.
        Thus $(x,f(x))\in f$ by \cref{function_apply_elim}.
        Thus $f(x)\in\ran{f}$ by \cref{ran_iff}.
        Thus $f(x)\in Z$ by \cref{subseteq}.
    \end{subproof}
    Thus $f\in\funs{X}{Z}$ by \cref{funs_intro}.
    Thus $f$ is continuous from $X$ to $Z$ with respect to $T$ and $T_Z$ by \cref{continuous}.
\end{proof}

% from §18 Item 16: expanding the range
\begin{lemma}\label{continuous_expand_range}
    Assume $f$ is continuous from $X$ to $Y$ with respect to $T$ and $T_1$.
    Suppose $Y$ is a subspace of $Z$ with respect to $T_2$.
    Let $T_Y = \subspaceTopology{T_2}{Y}$.
    Suppose $T_1 = T_Y$.
    Let $j = \identity{Y}$.
    Then $j\circ f$ is continuous from $X$ to $Z$ with respect to $T$ and $T_2$.
\end{lemma}
\begin{proof}
    $j$ is continuous from $Y$ to $Z$ with respect to $T_Y$ and $T_2$ by \cref{inclusion_continuous}.
    Thus $f$ is continuous from $X$ to $Y$ with respect to $T$ and $T_Y$.
    Thus $j\circ f$ is continuous from $X$ to $Z$ with respect to $T$ and $T_2$ by \cref{continuous_composite}.
\end{proof}

% from §18 Item 17: local formulation of continuity
\begin{lemma}\label{continuous_local}
    Assume $T$ is a topology on $X$.
    Assume $T_1$ is a topology on $Y$.
    Assume $f$ is a function from $X$ to $Y$.
    Suppose $M\subseteq T$.
    Suppose $\unions{M} = X$.
    Suppose for all $U$ such that $U\in M$ we have
        $\restrl{f}{U}$ is continuous from $U$ to $Y$ with respect to
            $\subspaceTopology{T}{U}$ and $T_1$.
    Then $f$ is continuous from $X$ to $Y$ with respect to $T$ and $T_1$.
\end{lemma}
\begin{proof}
    Thus $f\in\funs{X}{Y}$.
    Thus $\dom{f} = X$ by \cref{funs}.
    Show that for all $V$ such that $V$ is open in $Y$ with respect to $T_1$
        we have $\preimg{f}{V}$ is open in $X$ with respect to $T$.
    \begin{subproof}
        Fix $V$.
        Assume $V$ is open in $Y$ with respect to $T_1$.
        Let $N = \{W\in T \mid W\subseteq \preimg{f}{V}\}$.
        Show that for all $x$ such that $x\in\preimg{f}{V}$ we have $x\in\unions{N}$.
        \begin{subproof}
            Fix $x$.
            Assume $x\in\preimg{f}{V}$.
            $\preimg{f}{V}\subseteq \dom{f}$ by \cref{preimg_subseteq_dom}.
            Thus $x\in\dom{f}$ by \cref{subseteq}.
            Thus $x\in X$ by \cref{funs}.
            Take $U\in M$ such that $x\in U$ by \cref{unions_iff,subseteq}.
            Let $g = \restrl{f}{U}$.
            $g$ is continuous from $U$ to $Y$ with respect to $\subspaceTopology{T}{U}$ and $T_1$ by \cref{subseteq}.
            $\preimg{g}{V}$ is open in $U$ with respect to $\subspaceTopology{T}{U}$ by \cref{continuous}.
            $U$ is open in $X$ with respect to $T$ by \cref{open_in,subseteq}.
            $T\subseteq\pow{X}$ by \cref{topology_on}.
            Thus $U\in T$ by \cref{open_in}.
            Thus $U\in\pow{X}$ by \cref{subseteq}.
            Thus $U\subseteq X$ by \cref{pow_iff}.
            Thus $U$ is a subspace of $X$ with respect to $T$ by \cref{subspace_of}.
            Thus $\preimg{g}{V}$ is open in $X$ with respect to $T$ by \cref{open_in_subspace_open_in_space}.
            Show that for all $z$ such that $z\in\preimg{g}{V}$ we have $z\in\preimg{f}{V}$.
            \begin{subproof}
                Fix $z$.
                Assume $z\in\preimg{g}{V}$.
                Take $y\in V$ such that $z\mathrel{g} y$ by \cref{preimg_iff}.
                Thus $z\mathrel{f} y$ by \cref{restrl_iff}.
                Thus $z\in\preimg{f}{V}$ by \cref{preimg_iff}.
            \end{subproof}
            Thus $\preimg{g}{V}\subseteq \preimg{f}{V}$ by \cref{subseteq}.
            Thus $\preimg{g}{V}\in N$.
            Take $y\in V$ such that $x\mathrel{f} y$ by \cref{preimg_iff}.
            Thus $x\mathrel{g} y$ by \cref{restrl_iff}.
            Thus $x\in\preimg{g}{V}$ by \cref{preimg_iff}.
            Thus $x\in\unions{N}$ by \cref{unions_iff}.
        \end{subproof}
        Thus $\preimg{f}{V}\subseteq \unions{N}$ by \cref{subseteq}.
        $N\subseteq T$.
        Thus $\unions{N}\in T$ by \cref{topology_on}.
        Show that $N\subseteq\pow{\preimg{f}{V}}$.
        \begin{subproof}
            Show that for all $W$ such that $W\in N$ we have $W\in\pow{\preimg{f}{V}}$.
            \begin{subproof}
                Fix $W$.
                Assume $W\in N$.
                Thus $W\subseteq \preimg{f}{V}$ by \cref{subseteq}.
                Thus $W\in\pow{\preimg{f}{V}}$ by \cref{pow_iff}.
            \end{subproof}
            Thus $N\subseteq\pow{\preimg{f}{V}}$ by \cref{subseteq}.
        \end{subproof}
        $\unions{N}\subseteq \preimg{f}{V}$ by \cref{unions_subseteq_of_powerset_is_subseteq}.
        Thus $\preimg{f}{V} = \unions{N}$ by \cref{subseteq_antisymmetric}.
        Thus $\preimg{f}{V}$ is open in $X$ with respect to $T$ by \cref{open_in}.
    \end{subproof}
    Thus $f$ is continuous from $X$ to $Y$ with respect to $T$ and $T_1$ by \cref{continuous}.
\end{proof}

% from §18 helper lemma 4: union of two compatible functions is a function
\begin{lemma}\label{union_compatible_functions}
    Suppose $f$ is a function.
    Suppose $g$ is a function.
    Suppose for all $x$ such that $x\in\dom{f}\inter\dom{g}$ we have $f(x) = g(x)$.
    Then $f\union g$ is a function.
\end{lemma}
\begin{proof}
    Show that $f\union g$ is a relation.
    \begin{subproof}
        Show that for all $w$ such that $w\in f\union g$ there exist $x,y$ such that $w = (x,y)$.
        \begin{subproof}
            Fix $w$.
            Assume $w\in f\union g$.
            Thus $w\in f$ or $w\in g$ by \cref{union_iff}.
            \begin{byCase}
            \caseOf{$w\in f$.}
                Thus $f$ is a relation.
                Take $x,y$ such that $w = (x,y)$ by \cref{relation}.
            \caseOf{$w\in g$.}
                Thus $g$ is a relation.
                Take $x,y$ such that $w = (x,y)$ by \cref{relation}.
            \end{byCase}
        \end{subproof}
        Thus $f\union g$ is a relation by \cref{relation}.
    \end{subproof}
    Show that $f\union g$ is right-unique.
    \begin{subproof}
        Show that for all $x,y,z$ such that $(x,y)\in f\union g$ and $(x,z)\in f\union g$ we have $y = z$.
        \begin{subproof}
            Fix $x$.
            Fix $y$.
            Fix $z$.
            Assume $(x,y)\in f\union g$ and $(x,z)\in f\union g$.
            \begin{byCase}
            \caseOf{$(x,y)\in f$ and $(x,z)\in f$.}
                $f$ is right-unique by \cref{function_rightunique}.
                Thus $y = z$ by \hyperref[rightunique]{right-uniqueness}.
            \caseOf{$(x,y)\in g$ and $(x,z)\in g$.}
                $g$ is right-unique by \cref{function_rightunique}.
                Thus $y = z$ by \hyperref[rightunique]{right-uniqueness}.
            \caseOf{$(x,y)\in f$ and $(x,z)\in g$.}
                Thus $x\in\dom{f}$ and $x\in\dom{g}$ by \cref{dom_intro}.
                Thus $x\in\dom{f}\inter\dom{g}$ by \cref{inter_intro}.
                Thus $f(x) = g(x)$ by \cref{subseteq}.
                Thus $y = f(x)$ and $z = g(x)$ by \cref{function_apply_iff}.
                Thus $y = z$.
            \caseOf{$(x,y)\in g$ and $(x,z)\in f$.}
                Thus $x\in\dom{f}$ and $x\in\dom{g}$ by \cref{dom_intro}.
                Thus $x\in\dom{f}\inter\dom{g}$ by \cref{inter_intro}.
                Thus $f(x) = g(x)$ by \cref{subseteq}.
                Thus $y = g(x)$ and $z = f(x)$ by \cref{function_apply_iff}.
                Thus $y = z$.
            \end{byCase}
        \end{subproof}
        Thus $f\union g$ is right-unique by \cref{rightunique}.
    \end{subproof}
    Thus $f\union g$ is a function.
\end{proof}

% from §18 Item 18: pasting lemma
\begin{theorem}\label{pasting_lemma}
    Assume $X = A\union B$.
    Assume $A$ is closed in $X$ with respect to $T$.
    Assume $B$ is closed in $X$ with respect to $T$.
    Assume $f$ is continuous from $A$ to $Y$ with respect to $\subspaceTopology{T}{A}$ and $T_1$.
    Assume $g$ is continuous from $B$ to $Y$ with respect to $\subspaceTopology{T}{B}$ and $T_1$.
    Suppose for all $x$ such that $x\in A\inter B$ we have $f(x) = g(x)$.
    Let $h = f\union g$.
    Then $h$ is continuous from $X$ to $Y$ with respect to $T$ and $T_1$.
\end{theorem}
\begin{proof}
    Thus $f\in\funs{A}{Y}$ by \cref{continuous}.
    Thus $f$ is a function by \cref{funs_is_function}.
    Thus $g\in\funs{B}{Y}$ by \cref{continuous}.
    Thus $g$ is a function by \cref{funs_is_function}.
    Thus $\dom{f} = A$ by \cref{funs}.
    Thus $\dom{g} = B$ by \cref{funs}.
    Show that for all $x$ such that $x\in\dom{f}\inter\dom{g}$ we have $f(x) = g(x)$.
    \begin{subproof}
        Fix $x$.
        Assume $x\in\dom{f}\inter\dom{g}$.
        Thus $x\in A$ and $x\in B$ by \cref{inter}.
        Thus $x\in A\inter B$ by \cref{inter_intro}.
        Thus $f(x) = g(x)$.
    \end{subproof}
    $h$ is a function by \cref{union_compatible_functions}.
    Show that for all $C$ such that $C$ is closed in $Y$ with respect to $T_1$
        we have $\preimg{h}{C}$ is closed in $X$ with respect to $T$.
    \begin{subproof}
        Fix $C$.
        Assume $C$ is closed in $Y$ with respect to $T_1$.
        Let $A_1 = \preimg{f}{C}$.
        Let $B_1 = \preimg{g}{C}$.
        Let $D = Y\setminus C$.
        $D$ is open in $Y$ with respect to $T_1$ by \cref{closed_in}.
        Thus $D\subseteq Y$ by \cref{setminus_subseteq}.
        Thus $f\in\funs{A}{Y}$ by \cref{continuous}.
        Thus $g\in\funs{B}{Y}$ by \cref{continuous}.
        $\preimg{f}{D}$ is open in $A$ with respect to $\subspaceTopology{T}{A}$ by \cref{continuous}.
        $\preimg{g}{D}$ is open in $B$ with respect to $\subspaceTopology{T}{B}$ by \cref{continuous}.
        Thus $C\subseteq Y$ by \cref{closed_in}.
        $\preimg{f}{D} = A\setminus A_1$ by \cref{preimg_setminus_function}.
        $\preimg{g}{D} = B\setminus B_1$ by \cref{preimg_setminus_function}.
        $A_1\subseteq \dom{f}$ by \cref{preimg_subseteq_dom}.
        Thus $A_1\subseteq A$ by \cref{funs}.
        $B_1\subseteq \dom{g}$ by \cref{preimg_subseteq_dom}.
        Thus $B_1\subseteq B$ by \cref{funs}.
        $A\setminus \preimg{f}{D} = A\setminus (A\setminus A_1)$ by \cref{subseteq}.
        $B\setminus \preimg{g}{D} = B\setminus (B\setminus B_1)$ by \cref{subseteq}.
        $A\setminus (A\setminus A_1) = A\inter A_1$ by \cref{setminus_setminus}.
        $B\setminus (B\setminus B_1) = B\inter B_1$ by \cref{setminus_setminus}.
        Thus $A\inter A_1 = A_1$ by \cref{inter_absorb_supseteq_left}.
        Thus $B\inter B_1 = B_1$ by \cref{inter_absorb_supseteq_left}.
        Thus $A_1 = A\setminus \preimg{f}{D}$ by \cref{subseteq}.
        Thus $B_1 = B\setminus \preimg{g}{D}$ by \cref{subseteq}.
        Thus $A_1$ is closed in $A$ with respect to $\subspaceTopology{T}{A}$ by \cref{closed_in}.
        Thus $B_1$ is closed in $B$ with respect to $\subspaceTopology{T}{B}$ by \cref{closed_in}.
        Thus $X\setminus A$ is open in $X$ with respect to $T$ by \cref{closed_in}.
        Thus $T$ is a topology on $X$ by \cref{open_in}.
        Thus $A\subseteq X$ by \cref{closed_in}.
        Thus $B\subseteq X$ by \cref{closed_in}.
        $A_1$ is closed in $X$ with respect to $T$ by \cref{closed_in_closed_subspace}.
        $B_1$ is closed in $X$ with respect to $T$ by \cref{closed_in_closed_subspace}.
        Show that $\preimg{h}{C} = A_1\union B_1$.
        \begin{subproof}
            Show that for all $x$ such that $x\in\preimg{h}{C}$ we have $x\in A_1\union B_1$.
            \begin{subproof}
                Fix $x$.
                Assume $x\in\preimg{h}{C}$.
                Take $y\in C$ such that $x\mathrel{h} y$ by \cref{preimg_iff}.
                Thus $x\mathrel{f} y$ or $x\mathrel{g} y$ by \cref{union_iff}.
                \begin{byCase}
                \caseOf{$x\mathrel{f} y$.}
                    Thus $x\in A_1$ by \cref{preimg_iff}.
                    Thus $x\in A_1\union B_1$ by \cref{union_iff}.
                \caseOf{$x\mathrel{g} y$.}
                    Thus $x\in B_1$ by \cref{preimg_iff}.
                    Thus $x\in A_1\union B_1$ by \cref{union_iff}.
                \end{byCase}
            \end{subproof}
            Thus $\preimg{h}{C}\subseteq A_1\union B_1$ by \cref{subseteq}.
            Show that for all $x$ such that $x\in A_1\union B_1$ we have $x\in\preimg{h}{C}$.
            \begin{subproof}
                Fix $x$.
                Assume $x\in A_1\union B_1$.
                Thus $x\in A_1$ or $x\in B_1$ by \cref{union_iff}.
                \begin{byCase}
                \caseOf{$x\in A_1$.}
                    Take $y\in C$ such that $x\mathrel{f} y$ by \cref{preimg_iff}.
                    Thus $x\mathrel{h} y$ by \cref{union_intro_left}.
                    Thus $x\in\preimg{h}{C}$ by \cref{preimg_iff}.
                \caseOf{$x\in B_1$.}
                    Take $y\in C$ such that $x\mathrel{g} y$ by \cref{preimg_iff}.
                    Thus $x\mathrel{h} y$ by \cref{union_intro_right}.
                    Thus $x\in\preimg{h}{C}$ by \cref{preimg_iff}.
                \end{byCase}
            \end{subproof}
            Thus $A_1\union B_1\subseteq \preimg{h}{C}$ by \cref{subseteq}.
            Thus $\preimg{h}{C} = A_1\union B_1$ by \cref{subseteq_antisymmetric}.
        \end{subproof}
        Show that $A_1\union B_1$ is closed in $X$ with respect to $T$.
        \begin{subproof}
            $X\setminus A_1$ is open in $X$ with respect to $T$ by \cref{closed_in}.
            $X\setminus B_1$ is open in $X$ with respect to $T$ by \cref{closed_in}.
            Thus $X\setminus A_1\in T$ by \cref{open_in}.
            Thus $X\setminus B_1\in T$ by \cref{open_in}.
            Thus $(X\setminus A_1)\inter (X\setminus B_1)\in T$ by \cref{topology_on}.
            Thus $(X\setminus A_1)\inter (X\setminus B_1)$ is open in $X$ with respect to $T$ by \cref{open_in}.
            $X\setminus (A_1\union B_1) = (X\setminus A_1)\inter (X\setminus B_1)$ by \cref{setminus_union}.
            Thus $X\setminus (A_1\union B_1)$ is open in $X$ with respect to $T$ by \cref{subseteq}.
            $A_1\subseteq \dom{f}$ by \cref{preimg_subseteq_dom}.
            Thus $A_1\subseteq A$ by \cref{funs}.
            Thus $A_1\subseteq X$ by \cref{subseteq_transitive,closed_in}.
            $B_1\subseteq \dom{g}$ by \cref{preimg_subseteq_dom}.
            Thus $B_1\subseteq B$ by \cref{funs}.
            Thus $B_1\subseteq X$ by \cref{subseteq_transitive,closed_in}.
            Thus $A_1\union B_1\subseteq X$ by \cref{union_subsets_is_subset}.
            Thus $A_1\union B_1$ is closed in $X$ with respect to $T$ by \cref{closed_in}.
        \end{subproof}
        Thus $\preimg{h}{C}$ is closed in $X$ with respect to $T$ by \cref{subseteq}.
    \end{subproof}
    Thus $\dom{h} = \dom{f}\union\dom{g}$ by \cref{dom_union}.
    Thus $\dom{h} = A\union B$ by \cref{funs}.
    Thus $\dom{h} = X$ by \cref{subseteq}.
    Thus $f\in\rels{A}{Y}$ by \cref{funs_subseteq_rels,subseteq}.
    Thus $g\in\rels{B}{Y}$ by \cref{funs_subseteq_rels,subseteq}.
    Thus $\ran{f}\subseteq Y$ by \cref{rels_ran_subseteq}.
    Thus $\ran{g}\subseteq Y$ by \cref{rels_ran_subseteq}.
    Thus $\ran{h} = \ran{f}\union\ran{g}$ by \cref{ran_union}.
    Thus $\ran{h}\subseteq Y$ by \cref{union_subsets_is_subset}.
    Show that for all $x\in\dom{h}$ we have $h(x)\in Y$.
    \begin{subproof}
        Fix $x\in\dom{h}$.
        Thus $(x,h(x))\in h$ by \cref{function_apply_elim}.
        Thus $h(x)\in\ran{h}$ by \cref{ran_iff}.
        Thus $h(x)\in Y$ by \cref{subseteq}.
    \end{subproof}
    Thus $h\in\funs{X}{Y}$ by \cref{funs_intro}.
    Thus $T_1$ is a topology on $Y$ by \cref{continuous}.
    Thus $X\setminus A$ is open in $X$ with respect to $T$ by \cref{closed_in}.
    Thus $T$ is a topology on $X$ by \cref{open_in}.
    Thus $h$ is continuous from $X$ to $Y$ with respect to $T$ and $T_1$ by \cref{preimg_closed_implies_continuous}.
\end{proof}

% from §18 helper lemma 5: first projection is continuous
\begin{lemma}\label{projection_one_continuous}
    Assume $T$ is a topology on $X$.
    Assume $T_1$ is a topology on $Y$.
    Let $T_2 = \productTopology{T}{T_1}{X}{Y}$.
    Then $\piOne{X}{Y}$ is continuous from $X\times Y$ to $X$ with respect to $T_2$ and $T$.
\end{lemma}
\begin{proof}
    $\productBasis{T}{T_1}{X}{Y}$ is a basis on $X\times Y$ by \cref{product_basis_is_basis}.
    Thus $T_2$ is a topology on $X\times Y$ by \cref{basis_topology_is_topology,product_topology}.
    Show that $\piOne{X}{Y}$ is a function to $X$.
    \begin{subproof}
        Show that for all $p$ we have for all $x$ we have for all $x'$ such that
            $p\mathrel{\piOne{X}{Y}} x$ and $p\mathrel{\piOne{X}{Y}} x'$ we have $x = x'$.
        \begin{subproof}
            Fix $p$.
            Fix $x$.
            Fix $x'$.
            Assume $p\mathrel{\piOne{X}{Y}} x$ and $p\mathrel{\piOne{X}{Y}} x'$.
            Thus $(p,x)\in\piOne{X}{Y}$ and $(p,x')\in\piOne{X}{Y}$.
            Take $a,b$ such that $a\in X$ and $b\in Y$ and $(p,x)=((a,b),a)$ by \cref{projection_first}.
            Take $a',b'$ such that $a'\in X$ and $b'\in Y$ and $(p,x')=((a',b'),a')$ by \cref{projection_first}.
            Thus $p=(a,b)$ and $x=a$ by \cref{pair_eq_iff}.
            Thus $p=(a',b')$ and $x'=a'$ by \cref{pair_eq_iff}.
            Thus $(a,b) = (a',b')$.
            Thus $a = a'$ by \cref{pair_eq_iff}.
            Thus $x = x'$.
        \end{subproof}
        Thus $\piOne{X}{Y}$ is right-unique by \cref{rightunique}.
        Show that for all $w\in\piOne{X}{Y}$ there exists $x,y$ such that $w = (x,y)$.
        \begin{subproof}
            Fix $w$.
            Assume $w\in\piOne{X}{Y}$.
            Take $a,b$ such that $a\in X$ and $b\in Y$ and $w = ((a,b),a)$ by \cref{projection_first}.
            Let $x = (a,b)$.
            Let $y = a$.
            Thus $w = (x,y)$.
        \end{subproof}
        Thus $\piOne{X}{Y}$ is a relation by \cref{relation}.
        Thus $\piOne{X}{Y}$ is a function.
        Show that for all $p\in\dom{\piOne{X}{Y}}$ we have $\apply{\piOne{X}{Y}}{p}\in X$.
        \begin{subproof}
            Fix $p\in\dom{\piOne{X}{Y}}$.
            Take $x$ such that $p\mathrel{\piOne{X}{Y}} x$ by \cref{dom_iff}.
            Thus $(p,x)\in\piOne{X}{Y}$.
            Take $a,b$ such that $a\in X$ and $b\in Y$ and $(p,x)=((a,b),a)$ by \cref{projection_first}.
            Thus $x = a$ by \cref{pair_eq_iff}.
            Thus $\apply{\piOne{X}{Y}}{p} = x$ by \cref{function_apply_intro}.
            Thus $\apply{\piOne{X}{Y}}{p}\in X$.
        \end{subproof}
        Thus $\piOne{X}{Y}$ is a function to $X$.
    \end{subproof}
    Show that $\dom{\piOne{X}{Y}} = X\times Y$.
    \begin{subproof}
        Show that for all $p\in\dom{\piOne{X}{Y}}$ we have $p\in X\times Y$.
        \begin{subproof}
            Fix $p\in\dom{\piOne{X}{Y}}$.
            Take $x$ such that $p\mathrel{\piOne{X}{Y}} x$ by \cref{dom_iff}.
            Thus $(p,x)\in\piOne{X}{Y}$.
            Take $a,b$ such that $a\in X$ and $b\in Y$ and $(p,x)=((a,b),a)$ by \cref{projection_first}.
            Thus $p = (a,b)$ by \cref{pair_eq_iff}.
            Thus $p\in X\times Y$ by \cref{times}.
        \end{subproof}
        Thus $\dom{\piOne{X}{Y}}\subseteq X\times Y$ by \cref{subseteq}.
        Show that for all $p\in X\times Y$ we have $p\in\dom{\piOne{X}{Y}}$.
        \begin{subproof}
            Fix $p\in X\times Y$.
            Take $a,b$ such that $p = (a,b)$ and $a\in X$ and $b\in Y$ by \cref{times_elem_is_tuple}.
            Thus $p\mathrel{\piOne{X}{Y}} a$ by \cref{projection_first}.
            Thus $p\in\dom{\piOne{X}{Y}}$ by \cref{dom_intro}.
        \end{subproof}
        Thus $X\times Y\subseteq\dom{\piOne{X}{Y}}$ by \cref{subseteq}.
        Thus $\dom{\piOne{X}{Y}} = X\times Y$ by \cref{subseteq_antisymmetric}.
    \end{subproof}
    Thus $\piOne{X}{Y}\in\funs{X\times Y}{X}$ by \cref{funs_intro}.
    Show that for all $U$ such that $U$ is open in $X$ with respect to $T$
        we have $\preimg{\piOne{X}{Y}}{U}$ is open in $X\times Y$ with respect to $T_2$.
    \begin{subproof}
        Fix $U$.
        Assume $U$ is open in $X$ with respect to $T$.
        $U\in T$ by \cref{open_in}.
        $Y\in T_1$ by \cref{topology_on}.
        $Y$ is open in $Y$ with respect to $T_1$ by \cref{open_in}.
        $T\subseteq\pow{X}$ by \cref{topology_on}.
        Thus $U\subseteq X$ by \cref{pow_iff,subseteq}.
        Thus $U\times Y\subseteq X\times Y$ by \cref{times_subseteq_left}.
        Thus $U\times Y\in\pow{X\times Y}$ by \cref{pow_iff}.
        Thus $U\times Y\in \productBasis{T}{T_1}{X}{Y}$ by \cref{product_basis}.
        $\productBasis{T}{T_1}{X}{Y}$ is a basis on $X\times Y$ by \cref{product_basis_is_basis}.
        Thus $U\times Y\in \productTopology{T}{T_1}{X}{Y}$ by \cref{basis_elem_open_in_generated,product_topology}.
        Thus $U\times Y$ is open in $X\times Y$ with respect to $T_2$ by \cref{open_in}.
        $\preimg{\piOne{X}{Y}}{U} = U\times Y$ by \cref{preimg_pi1}.
        Thus $\preimg{\piOne{X}{Y}}{U}$ is open in $X\times Y$ with respect to $T_2$.
    \end{subproof}
    Thus $\piOne{X}{Y}$ is continuous from $X\times Y$ to $X$ with respect to $T_2$ and $T$ by \cref{continuous}.
\end{proof}

% from §18 helper lemma 6: second projection is continuous
\begin{lemma}\label{projection_two_continuous}
    Assume $T$ is a topology on $X$.
    Assume $T_1$ is a topology on $Y$.
    Let $T_2 = \productTopology{T}{T_1}{X}{Y}$.
    Then $\piTwo{X}{Y}$ is continuous from $X\times Y$ to $Y$ with respect to $T_2$ and $T_1$.
\end{lemma}
\begin{proof}
    $\productBasis{T}{T_1}{X}{Y}$ is a basis on $X\times Y$ by \cref{product_basis_is_basis}.
    Thus $T_2$ is a topology on $X\times Y$ by \cref{basis_topology_is_topology,product_topology}.
    Show that $\piTwo{X}{Y}$ is a function to $Y$.
    \begin{subproof}
        Show that for all $p$ we have for all $y$ we have for all $y'$ such that
            $p\mathrel{\piTwo{X}{Y}} y$ and $p\mathrel{\piTwo{X}{Y}} y'$ we have $y = y'$.
        \begin{subproof}
            Fix $p$.
            Fix $y$.
            Fix $y'$.
            Assume $p\mathrel{\piTwo{X}{Y}} y$ and $p\mathrel{\piTwo{X}{Y}} y'$.
            Thus $(p,y)\in\piTwo{X}{Y}$ and $(p,y')\in\piTwo{X}{Y}$.
            Take $a,b$ such that $a\in X$ and $b\in Y$ and $(p,y)=((a,b),b)$ by \cref{projection_second}.
            Take $a',b'$ such that $a'\in X$ and $b'\in Y$ and $(p,y')=((a',b'),b')$ by \cref{projection_second}.
            Thus $p=(a,b)$ and $y=b$ by \cref{pair_eq_iff}.
            Thus $p=(a',b')$ and $y'=b'$ by \cref{pair_eq_iff}.
            Thus $(a,b) = (a',b')$.
            Thus $b = b'$ by \cref{pair_eq_iff}.
            Thus $y = y'$.
        \end{subproof}
        Thus $\piTwo{X}{Y}$ is right-unique by \cref{rightunique}.
        Show that for all $w\in\piTwo{X}{Y}$ there exists $x,y$ such that $w = (x,y)$.
        \begin{subproof}
            Fix $w$.
            Assume $w\in\piTwo{X}{Y}$.
            Take $a,b$ such that $a\in X$ and $b\in Y$ and $w = ((a,b),b)$ by \cref{projection_second}.
            Let $x = (a,b)$.
            Let $y = b$.
            Thus $w = (x,y)$.
        \end{subproof}
        Thus $\piTwo{X}{Y}$ is a relation by \cref{relation}.
        Thus $\piTwo{X}{Y}$ is a function.
        Show that for all $p\in\dom{\piTwo{X}{Y}}$ we have $\apply{\piTwo{X}{Y}}{p}\in Y$.
        \begin{subproof}
            Fix $p\in\dom{\piTwo{X}{Y}}$.
            Take $y$ such that $p\mathrel{\piTwo{X}{Y}} y$ by \cref{dom_iff}.
            Thus $(p,y)\in\piTwo{X}{Y}$.
            Take $a,b$ such that $a\in X$ and $b\in Y$ and $(p,y)=((a,b),b)$ by \cref{projection_second}.
            Thus $y = b$ by \cref{pair_eq_iff}.
            Thus $\apply{\piTwo{X}{Y}}{p} = y$ by \cref{function_apply_intro}.
            Thus $\apply{\piTwo{X}{Y}}{p}\in Y$.
        \end{subproof}
        Thus $\piTwo{X}{Y}$ is a function to $Y$.
    \end{subproof}
    Show that $\dom{\piTwo{X}{Y}} = X\times Y$.
    \begin{subproof}
        Show that for all $p\in\dom{\piTwo{X}{Y}}$ we have $p\in X\times Y$.
        \begin{subproof}
            Fix $p\in\dom{\piTwo{X}{Y}}$.
            Take $y$ such that $p\mathrel{\piTwo{X}{Y}} y$ by \cref{dom_iff}.
            Thus $(p,y)\in\piTwo{X}{Y}$.
            Take $a,b$ such that $a\in X$ and $b\in Y$ and $(p,y)=((a,b),b)$ by \cref{projection_second}.
            Thus $p = (a,b)$ by \cref{pair_eq_iff}.
            Thus $p\in X\times Y$ by \cref{times}.
        \end{subproof}
        Thus $\dom{\piTwo{X}{Y}}\subseteq X\times Y$ by \cref{subseteq}.
        Show that for all $p\in X\times Y$ we have $p\in\dom{\piTwo{X}{Y}}$.
        \begin{subproof}
            Fix $p\in X\times Y$.
            Take $a,b$ such that $p = (a,b)$ and $a\in X$ and $b\in Y$ by \cref{times_elem_is_tuple}.
            Thus $p\mathrel{\piTwo{X}{Y}} b$ by \cref{projection_second}.
            Thus $p\in\dom{\piTwo{X}{Y}}$ by \cref{dom_intro}.
        \end{subproof}
        Thus $X\times Y\subseteq\dom{\piTwo{X}{Y}}$ by \cref{subseteq}.
        Thus $\dom{\piTwo{X}{Y}} = X\times Y$ by \cref{subseteq_antisymmetric}.
    \end{subproof}
    Thus $\piTwo{X}{Y}\in\funs{X\times Y}{Y}$ by \cref{funs_intro}.
    Show that for all $V$ such that $V$ is open in $Y$ with respect to $T_1$
        we have $\preimg{\piTwo{X}{Y}}{V}$ is open in $X\times Y$ with respect to $T_2$.
    \begin{subproof}
        Fix $V$.
        Assume $V$ is open in $Y$ with respect to $T_1$.
        $V\in T_1$ by \cref{open_in}.
        $X\in T$ by \cref{topology_on}.
        $X$ is open in $X$ with respect to $T$ by \cref{open_in}.
        $T_1\subseteq\pow{Y}$ by \cref{topology_on}.
        Thus $V\subseteq Y$ by \cref{pow_iff,subseteq}.
        Thus $X\times V\subseteq X\times Y$ by \cref{times_subseteq_right}.
        Thus $X\times V\in\pow{X\times Y}$ by \cref{pow_iff}.
        Thus $X\times V\in \productBasis{T}{T_1}{X}{Y}$ by \cref{product_basis}.
        Thus $X\times V$ is open in $X\times Y$ with respect to $T_2$ by \cref{product_topology,basis_elem_open_in_generated,open_in}.
        $\preimg{\piTwo{X}{Y}}{V} = X\times V$ by \cref{preimg_pi2}.
        Thus $\preimg{\piTwo{X}{Y}}{V}$ is open in $X\times Y$ with respect to $T_2$.
    \end{subproof}
    Thus $\piTwo{X}{Y}$ is continuous from $X\times Y$ to $Y$ with respect to $T_2$ and $T_1$ by \cref{continuous}.
\end{proof}

% from §18 Item 19: maps into products (forward)
\begin{lemma}\label{continuous_map_into_product_forward}
    Assume $T_1$ is a topology on $X$.
    Assume $T_2$ is a topology on $Y$.
    Assume $f$ is continuous from $A$ to $X\times Y$ with respect to $T$ and
        $\productTopology{T_1}{T_2}{X}{Y}$.
    Let $f_1 = \piOne{X}{Y}\circ f$.
    Let $f_2 = \piTwo{X}{Y}\circ f$.
    Then $f_1$ is continuous from $A$ to $X$ with respect to $T$ and $T_1$ and
        $f_2$ is continuous from $A$ to $Y$ with respect to $T$ and $T_2$.
\end{lemma}
\begin{proof}
    Let $T_3 = \productTopology{T_1}{T_2}{X}{Y}$.
    $\piOne{X}{Y}$ is continuous from $X\times Y$ to $X$ with respect to $T_3$ and $T_1$
        by \cref{projection_one_continuous}.
    $\piTwo{X}{Y}$ is continuous from $X\times Y$ to $Y$ with respect to $T_3$ and $T_2$
        by \cref{projection_two_continuous}.
    Thus $f_1$ is continuous from $A$ to $X$ with respect to $T$ and $T_1$ by \cref{continuous_composite}.
    Thus $f_2$ is continuous from $A$ to $Y$ with respect to $T$ and $T_2$ by \cref{continuous_composite}.
\end{proof}

% from §18 Item 20: maps into products (backward)
\begin{lemma}\label{continuous_map_into_product_backward}
    Assume $f_1$ is continuous from $A$ to $X$ with respect to $T$ and $T_1$.
    Assume $f_2$ is continuous from $A$ to $Y$ with respect to $T$ and $T_2$.
    Let $f = \{(a,(f_1(a),f_2(a))) \mid a\in A\}$.
    Then $f$ is continuous from $A$ to $X\times Y$ with respect to $T$ and
        $\productTopology{T_1}{T_2}{X}{Y}$.
\end{lemma}
\begin{proof}
    Let $T_3 = \productTopology{T_1}{T_2}{X}{Y}$.
    Show that for all $U$ such that $U$ is open in $X\times Y$ with respect to $T_3$
        we have $\preimg{f}{U}$ is open in $A$ with respect to $T$.
    \begin{subproof}
        Fix $U$.
        Assume $U$ is open in $X\times Y$ with respect to $T_3$.
        $U\in T_3$ by \cref{open_in}.
        $T_1$ is a topology on $X$ by \cref{continuous}.
        $T_2$ is a topology on $Y$ by \cref{continuous}.
        $\productBasis{T_1}{T_2}{X}{Y}$ is a basis on $X\times Y$ by \cref{product_basis_is_basis}.
        $\productBasis{T_1}{T_2}{X}{Y}$ is a basis for $T_3$ on $X\times Y$ by \cref{basis_for_topology,product_topology}.
        Take $M\subseteq \productBasis{T_1}{T_2}{X}{Y}$ such that $U = \unions{M}$ by \cref{basis_topology_unions}.
        Let $N = \{\preimg{f}{B} \mid B\in M\}$.
        Show that for all $C$ such that $C\in N$ we have $C\in T$.
        \begin{subproof}
            Fix $C$.
            Assume $C\in N$.
            Take $B\in M$ such that $C = \preimg{f}{B}$.
            Thus $B\in \productBasis{T_1}{T_2}{X}{Y}$ by \cref{subseteq}.
            Take $U_1,V_1$ such that $U_1\in T_1$ and $V_1\in T_2$ and $B = U_1\times V_1$
                by \cref{product_basis}.
            $U_1$ is open in $X$ with respect to $T_1$ by \cref{open_in}.
            $V_1$ is open in $Y$ with respect to $T_2$ by \cref{open_in}.
            Show that $\preimg{f}{B} = \preimg{f_1}{U_1}\inter \preimg{f_2}{V_1}$.
            \begin{subproof}
                Show that for all $a$ such that $a\in\preimg{f}{B}$ we have
                    $a\in\preimg{f_1}{U_1}\inter \preimg{f_2}{V_1}$.
                \begin{subproof}
                    Fix $a$.
                    Assume $a\in\preimg{f}{B}$.
                    Take $p\in B$ such that $a\mathrel{f} p$ by \cref{preimg_iff}.
                    Take $x,y$ such that $p = (x,y)$ and $x\in U_1$ and $y\in V_1$ by \cref{times_elem_is_tuple}.
                    Take $a_0\in A$ such that $(a,p) = (a_0,(f_1(a_0),f_2(a_0)))$.
                    Thus $a = a_0$ and $p = (f_1(a),f_2(a))$ by \cref{pair_eq_iff}.
                    Thus $f_1(a) = x$ and $f_2(a) = y$ by \cref{pair_eq_iff}.
                    Thus $a\in A$.
                    $f_1\in\funs{A}{X}$ and $f_2\in\funs{A}{Y}$ by \cref{continuous}.
                    Thus $f_1$ is a function and $\dom{f_1} = A$ by \cref{funs_elim}.
                    Thus $f_2$ is a function and $\dom{f_2} = A$ by \cref{funs_elim}.
                    Thus $a\in\dom{f_1}$ and $a\in\dom{f_2}$.
                    Thus $(a,f_1(a))\in f_1$ and $(a,f_2(a))\in f_2$ by \cref{function_apply_elim}.
                    Thus $(a,x)\in f_1$ and $(a,y)\in f_2$ by \cref{pair_eq_iff}.
                    Thus $a\mathrel{f_1} x$ and $a\mathrel{f_2} y$.
                    Thus $a\in\preimg{f_1}{U_1}$ and $a\in\preimg{f_2}{V_1}$ by \cref{preimg_iff}.
                    Thus $a\in\preimg{f_1}{U_1}\inter \preimg{f_2}{V_1}$ by \cref{inter_intro}.
                \end{subproof}
                Thus $\preimg{f}{B}\subseteq \preimg{f_1}{U_1}\inter \preimg{f_2}{V_1}$ by \cref{subseteq}.
                Show that for all $a$ such that $a\in\preimg{f_1}{U_1}\inter \preimg{f_2}{V_1}$ we have
                    $a\in\preimg{f}{B}$.
                \begin{subproof}
                    Fix $a$.
                    Assume $a\in\preimg{f_1}{U_1}\inter \preimg{f_2}{V_1}$.
                    Thus $a\in\preimg{f_1}{U_1}$ and $a\in\preimg{f_2}{V_1}$ by \cref{inter}.
                    Take $x\in U_1$ such that $a\mathrel{f_1} x$ by \cref{preimg_iff}.
                    Take $y\in V_1$ such that $a\mathrel{f_2} y$ by \cref{preimg_iff}.
                    $f_1\in\funs{A}{X}$ and $f_2\in\funs{A}{Y}$ by \cref{continuous}.
                    Thus $f_1$ is a function and $\dom{f_1} = A$ by \cref{funs_elim}.
                    Thus $f_2$ is a function and $\dom{f_2} = A$ by \cref{funs_elim}.
                    Thus $f_1(a) = x$ and $f_2(a) = y$ by \cref{function_apply_intro}.
                    Thus $a\in\dom{f_1}$ by \cref{dom_intro}.
                    Thus $a\in A$.
                    Let $p = (x,y)$.
                    Thus $p\in U_1\times V_1$ by \cref{times_tuple_intro}.
                    Thus $p\in B$.
                    Thus $p = (f_1(a),f_2(a))$ by \cref{pair_eq_iff}.
                    Thus $a\mathrel{f} p$.
                    Thus $a\in\preimg{f}{B}$ by \cref{preimg_iff}.
                \end{subproof}
                Thus $\preimg{f_1}{U_1}\inter \preimg{f_2}{V_1}\subseteq \preimg{f}{B}$ by \cref{subseteq}.
                Thus $\preimg{f}{B} = \preimg{f_1}{U_1}\inter \preimg{f_2}{V_1}$ by \cref{subseteq_antisymmetric}.
            \end{subproof}
            $\preimg{f_1}{U_1}$ is open in $A$ with respect to $T$ by \cref{continuous}.
            $\preimg{f_2}{V_1}$ is open in $A$ with respect to $T$ by \cref{continuous}.
            Thus $\preimg{f_1}{U_1}\in T$ and $\preimg{f_2}{V_1}\in T$ by \cref{open_in}.
            Thus $T$ is a topology on $A$ by \cref{continuous}.
            Thus $\preimg{f_1}{U_1}\inter \preimg{f_2}{V_1}\in T$ by \cref{topology_on}.
            Thus $\preimg{f}{B}\in T$.
            Thus $\preimg{f}{B}$ is open in $A$ with respect to $T$ by \cref{open_in}.
            Thus $C\in T$.
        \end{subproof}
        Thus $N\subseteq T$ by \cref{subseteq}.
        $T$ is a topology on $A$ by \cref{continuous}.
        $\unions{N}\in T$ by \cref{topology_on}.
        Show that $\preimg{f}{U} = \unions{N}$.
        \begin{subproof}
            Show that for all $a$ such that $a\in\preimg{f}{U}$ we have $a\in\unions{N}$.
            \begin{subproof}
                Fix $a$.
                Assume $a\in\preimg{f}{U}$.
                Take $p\in U$ such that $a\mathrel{f} p$ by \cref{preimg_iff}.
                Take $B\in M$ such that $p\in B$ by \cref{unions_iff}.
                Thus $a\in\preimg{f}{B}$ by \cref{preimg_iff}.
                Thus $a\in\unions{N}$ by \cref{unions_iff}.
            \end{subproof}
            Thus $\preimg{f}{U}\subseteq \unions{N}$ by \cref{subseteq}.
            Show that for all $a$ such that $a\in\unions{N}$ we have $a\in\preimg{f}{U}$.
            \begin{subproof}
                Fix $a$.
                Assume $a\in\unions{N}$.
                Take $C\in N$ such that $a\in C$ by \cref{unions_iff}.
                    Take $B\in M$ such that $C = \preimg{f}{B}$.
                    Thus $a\in\preimg{f}{B}$.
                    Take $p\in B$ such that $a\mathrel{f} p$ by \cref{preimg_iff}.
                    Show that for all $q\in B$ we have $q\in U$.
                    \begin{subproof}
                        Fix $q\in B$.
                        Thus $q\in\unions{M}$ by \cref{unions_iff}.
                        Thus $q\in U$.
                    \end{subproof}
                    Thus $B\subseteq U$ by \cref{subseteq}.
                    Thus $p\in U$ by \cref{subseteq}.
                    Thus $a\in\preimg{f}{U}$ by \cref{preimg_iff}.
                \end{subproof}
            Thus $\unions{N}\subseteq \preimg{f}{U}$ by \cref{subseteq}.
            Thus $\preimg{f}{U} = \unions{N}$ by \cref{subseteq_antisymmetric}.
        \end{subproof}
        Thus $\preimg{f}{U}$ is open in $A$ with respect to $T$ by \cref{open_in}.
    \end{subproof}
    Show that $f\in\funs{A}{X\times Y}$.
    \begin{subproof}
        Show that $f$ is a function.
        \begin{subproof}
            Show that for all $w\in f$ there exist $x,y$ such that $w = (x,y)$.
            \begin{subproof}
                Fix $w$.
                Assume $w\in f$.
                Take $a\in A$ such that $w = (a,(f_1(a),f_2(a)))$.
                Let $x = a$.
                Let $y = (f_1(a),f_2(a))$.
                Thus $w = (x,y)$.
            \end{subproof}
            Thus $f$ is a relation by \cref{relation}.
            Show that for all $x,y,z$ such that $(x,y)\in f$ and $(x,z)\in f$ we have $y = z$.
            \begin{subproof}
                Fix $x$.
                Fix $y$.
                Fix $z$.
                Assume $(x,y)\in f$ and $(x,z)\in f$.
                Take $a_1\in A$ such that $(x,y) = (a_1,(f_1(a_1),f_2(a_1)))$.
                Take $a_2\in A$ such that $(x,z) = (a_2,(f_1(a_2),f_2(a_2)))$.
                Thus $x = a_1$ and $y = (f_1(a_1),f_2(a_1))$ by \cref{pair_eq_iff}.
                Thus $x = a_2$ and $z = (f_1(a_2),f_2(a_2))$ by \cref{pair_eq_iff}.
                Thus $a_1 = a_2$.
                Thus $y = z$.
            \end{subproof}
            Thus $f$ is right-unique by \cref{rightunique}.
            Thus $f$ is a function.
        \end{subproof}
        Show that $\dom{f} = A$.
        \begin{subproof}
            Show that for all $a\in\dom{f}$ we have $a\in A$.
            \begin{subproof}
                Fix $a\in\dom{f}$.
                Take $p$ such that $a\mathrel{f} p$ by \cref{dom_iff}.
                Take $a_0\in A$ such that $(a,p) = (a_0,(f_1(a_0),f_2(a_0)))$.
                Thus $a = a_0$ by \cref{pair_eq_iff}.
                Thus $a\in A$.
            \end{subproof}
            Thus $\dom{f}\subseteq A$ by \cref{subseteq}.
            Show that for all $a\in A$ we have $a\in\dom{f}$.
            \begin{subproof}
                Fix $a\in A$.
                Thus $(a,(f_1(a),f_2(a)))\in f$.
                Thus $a\in\dom{f}$ by \cref{dom_intro}.
            \end{subproof}
            Thus $A\subseteq\dom{f}$ by \cref{subseteq}.
            Thus $\dom{f} = A$ by \cref{subseteq_antisymmetric}.
        \end{subproof}
        Show that for all $a\in\dom{f}$ we have $f(a)\in X\times Y$.
        \begin{subproof}
            Fix $a\in\dom{f}$.
            Thus $a\in A$.
            $f_1\in\funs{A}{X}$ and $f_2\in\funs{A}{Y}$ by \cref{continuous}.
            Thus $f_1$ is a function to $X$ and $\dom{f_1} = A$ by \cref{funs_elim}.
            Thus $f_2$ is a function to $Y$ and $\dom{f_2} = A$ by \cref{funs_elim}.
            Thus $a\in\dom{f_1}$ and $a\in\dom{f_2}$.
            Thus $f_1(a)\in X$ and $f_2(a)\in Y$.
            Thus $(f_1(a),f_2(a))\in X\times Y$ by \cref{times_tuple_intro}.
            Thus $(a,(f_1(a),f_2(a)))\in f$.
            Thus $f(a) = (f_1(a),f_2(a))$ by \cref{function_apply_intro}.
            Thus $f(a)\in X\times Y$.
        \end{subproof}
        Thus $f$ is a function to $X\times Y$.
        Thus $f\in\funs{A}{X\times Y}$ by \cref{funs_intro}.
    \end{subproof}
    $T$ is a topology on $A$ by \cref{continuous}.
    $T_1$ is a topology on $X$ by \cref{continuous}.
    $T_2$ is a topology on $Y$ by \cref{continuous}.
    $\productBasis{T_1}{T_2}{X}{Y}$ is a basis on $X\times Y$ by \cref{product_basis_is_basis}.
    Thus $T_3$ is a topology on $X\times Y$ by \cref{basis_topology_is_topology,product_topology}.
    Thus $f$ is continuous from $A$ to $X\times Y$ with respect to $T$ and $T_3$ by \cref{continuous}.
\end{proof}

% from §19 Item 1: J-tuples
\begin{definition}\label{j_tuple}
    $f$ is a $J$-tuple of elements of $X$ iff $f\in\funs{J}{X}$.
\end{definition}

% from §19 Item 2: the set of J-tuples
\begin{definition}\label{tuple_power}
    $\tuplePower{X}{J} = \funs{J}{X}$.
\end{definition}

% from §19 Item 3: product of an indexed family
\begin{definition}\label{indexed_product}
    Let $A$ be a function.
    $\productFamily{A} = \{f\in\funs{\dom{A}}{\unions{\ran{A}}} \mid \forall \alpha\in\dom{A}. \apply{f}{\alpha}\in\apply{A}{\alpha}\}$.
\end{definition}

% from §19 Item 4: box basis
\begin{definition}\label{box_basis}
    $\boxBasis{T}{X} = \{B\in\pow{\productFamily{X}} \mid \exists U.
        U\in\funs{\dom{X}}{\pow{\unions{\ran{X}}}}\land
        (\forall \alpha\in\dom{X}. \apply{U}{\alpha}\in\apply{T}{\alpha})\land
        B = \productFamily{U}\}$.
\end{definition}

% from §19 Item 5: box topology
\begin{definition}\label{box_topology}
    $\boxTopology{T}{X} = \basisTopology{\boxBasis{T}{X}}{\productFamily{X}}$.
\end{definition}

% from §19 Item 6: projection mapping
\begin{definition}\label{projection_map_indexed}
    $\piIndex{X}{\beta} = \{(x,\apply{x}{\beta}) \mid x\in\productFamily{X}\}$.
\end{definition}

% from §19 Item 7: product subbasis for a fixed index
\begin{definition}\label{product_subbasis_indexed}
    $\productSubbasisAt{T}{X}{\beta} = \{A\in\pow{\productFamily{X}} \mid
        \exists U\in\apply{T}{\beta}. A = \preimg{\piIndex{X}{\beta}}{U}\}$.
\end{definition}

% from §19 Item 8: product subbasis
\begin{definition}\label{product_subbasis_union_indexed}
    $\productSubbasis{T}{X} = \{A\in\pow{\productFamily{X}} \mid
        \exists \beta\in\dom{X}. A\in\productSubbasisAt{T}{X}{\beta}\}$.
\end{definition}

% from §19 Item 9: product topology on indexed products
\begin{definition}\label{product_topology_indexed}
    $\productTopologyIndexed{T}{X} = \subbasisTopology{\productSubbasis{T}{X}}{\productFamily{X}}$.
\end{definition}

% from §19 Item 10: box basis is a basis on the indexed product
\begin{lemma}\label{box_basis_is_basis}
    Assume $X$ is a function.
    Assume $T$ is a function.
    Assume $\dom{T} = \dom{X}$.
    Assume that for all $\alpha\in\dom{X}$ we have $\apply{T}{\alpha}$ is a topology on $\apply{X}{\alpha}$.
    Then $\boxBasis{T}{X}$ is a basis on $\productFamily{X}$.
\end{lemma}
\begin{proof}
    Let $B = \boxBasis{T}{X}$.
    Show that for all $B_1$ such that $B_1\in B$ we have $B_1\in\pow{\productFamily{X}}$.
    \begin{subproof}
        Fix $B_1$.
        Assume $B_1\in B$.
        Thus $B_1\in\pow{\productFamily{X}}$ by \cref{box_basis}.
    \end{subproof}
    Thus $B\subseteq\pow{\productFamily{X}}$ by \cref{subseteq}.
    Show that for all $x$ such that $x\in\productFamily{X}$ there exists $B_1\in B$ such that $x\in B_1$.
    \begin{subproof}
        Fix $x$.
        Assume $x\in\productFamily{X}$.
        Show that for all $\alpha\in\dom{X}$ we have $\apply{X}{\alpha}\in\pow{\unions{\ran{X}}}$.
        \begin{subproof}
            Fix $\alpha$.
            Assume $\alpha\in\dom{X}$.
            Thus $\alpha\mathrel{X}\apply{X}{\alpha}$ by \cref{function_apply_elim}.
            Thus $\apply{X}{\alpha}\in\ran{X}$ by \cref{ran_intro}.
            $\ran{X}\subseteq\pow{\unions{\ran{X}}}$ by \cref{subseteq_pow_unions}.
            Thus $\apply{X}{\alpha}\in\pow{\unions{\ran{X}}}$ by \cref{subseteq}.
        \end{subproof}
        Thus $X$ is a function to $\pow{\unions{\ran{X}}}$.
        Thus $X\in\funs{\dom{X}}{\pow{\unions{\ran{X}}}}$ by \cref{funs_intro}.
        Show that for all $\alpha\in\dom{X}$ we have $\apply{X}{\alpha}\in\apply{T}{\alpha}$.
        \begin{subproof}
            Fix $\alpha$.
            Assume $\alpha\in\dom{X}$.
            $\apply{T}{\alpha}$ is a topology on $\apply{X}{\alpha}$.
            Thus $\apply{X}{\alpha}\in\apply{T}{\alpha}$ by \cref{topology_on}.
        \end{subproof}
        Let $B_1 = \productFamily{X}$.
        Thus $B_1\subseteq\productFamily{X}$ by \cref{subseteq_refl}.
        Thus $B_1\in\pow{\productFamily{X}}$ by \cref{powerset_intro}.
        Thus $B_1\in B$ by \cref{box_basis}.
        Thus $x\in B_1$.
    \end{subproof}
    Show that for all $B_1,B_2,x$ such that $B_1\in B$ and $B_2\in B$ and $x\in B_1\inter B_2$
        there exists $B_3\in B$ such that $x\in B_3$ and $B_3\subseteq B_1\inter B_2$.
    \begin{subproof}
        Fix $B_1$.
        Fix $B_2$.
        Fix $x$.
        Assume $B_1\in B$ and $B_2\in B$ and $x\in B_1\inter B_2$.
        Take $U_1$ such that
            $U_1\in\funs{\dom{X}}{\pow{\unions{\ran{X}}}}$ and
            $B_1 = \productFamily{U_1}$ and
            for all $\alpha\in\dom{X}$ we have $\apply{U_1}{\alpha}\in\apply{T}{\alpha}$
            by \cref{box_basis}.
        $\dom{U_1} = \dom{X}$ by \cref{funs_elim}.
        Take $U_2$ such that
            $U_2\in\funs{\dom{X}}{\pow{\unions{\ran{X}}}}$ and
            $B_2 = \productFamily{U_2}$ and
            for all $\alpha\in\dom{X}$ we have $\apply{U_2}{\alpha}\in\apply{T}{\alpha}$
            by \cref{box_basis}.
        $\dom{U_2} = \dom{X}$ by \cref{funs_elim}.
        Let $R = \{(\alpha,\apply{U_1}{\alpha}\inter\apply{U_2}{\alpha}) \mid \alpha\in\dom{X}\}$.
        Show that $R$ is a relation.
        \begin{subproof}
            Show that for all $w$ such that $w\in R$ there exist $x,y$ such that $w = (x,y)$.
            \begin{subproof}
                Fix $w$.
                Assume $w\in R$.
                Take $\alpha\in\dom{X}$ such that $w = (\alpha,\apply{U_1}{\alpha}\inter\apply{U_2}{\alpha})$.
            \end{subproof}
            Thus $R$ is a relation by \cref{relation}.
        \end{subproof}
        Let $U_3 = R$.
        Show that for all $a$ we have for all $y$ we have for all $z$ if $(a,y)\in U_3$ and $(a,z)\in U_3$ then $y = z$.
        \begin{subproof}
            Fix $a$.
            Fix $y$.
            Fix $z$.
            Assume $(a,y)\in U_3$ and $(a,z)\in U_3$.
            Take $\alpha_1\in\dom{X}$ such that $(a,y) = (\alpha_1,\apply{U_1}{\alpha_1}\inter\apply{U_2}{\alpha_1})$.
            Take $\alpha_2\in\dom{X}$ such that $(a,z) = (\alpha_2,\apply{U_1}{\alpha_2}\inter\apply{U_2}{\alpha_2})$.
            Thus $a = \alpha_1$ and $y = \apply{U_1}{\alpha_1}\inter\apply{U_2}{\alpha_1}$ by \cref{pair_eq_iff}.
            Thus $a = \alpha_2$ and $z = \apply{U_1}{\alpha_2}\inter\apply{U_2}{\alpha_2}$ by \cref{pair_eq_iff}.
            Thus $y = \apply{U_1}{a}\inter\apply{U_2}{a}$.
            Thus $z = \apply{U_1}{a}\inter\apply{U_2}{a}$.
            Thus $y = z$.
        \end{subproof}
        Thus $U_3$ is right-unique.
        Thus $U_3$ is a function.
        Show that $\dom{U_3} = \dom{X}$.
        \begin{subproof}
            Show that for all $\alpha\in\dom{X}$ we have $\alpha\in\dom{U_3}$.
            \begin{subproof}
                Fix $\alpha$.
                Assume $\alpha\in\dom{X}$.
                Thus $(\alpha,\apply{U_1}{\alpha}\inter\apply{U_2}{\alpha})\in U_3$.
                Thus $\alpha\in\dom{U_3}$ by \cref{dom}.
            \end{subproof}
            Thus $\dom{X}\subseteq\dom{U_3}$ by \cref{subseteq}.
            Show that for all $\alpha\in\dom{U_3}$ we have $\alpha\in\dom{X}$.
            \begin{subproof}
                Fix $\alpha$.
                Assume $\alpha\in\dom{U_3}$.
                Take $y$ such that $(\alpha,y)\in U_3$ by \cref{dom}.
                Thus $\alpha\in\dom{X}$.
            \end{subproof}
            Thus $\dom{U_3}\subseteq\dom{X}$ by \cref{subseteq}.
            Thus $\dom{U_3} = \dom{X}$ by \cref{subseteq_antisymmetric}.
        \end{subproof}
        Show that for all $\alpha\in\dom{X}$ we have $\apply{U_3}{\alpha} = \apply{U_1}{\alpha}\inter\apply{U_2}{\alpha}$.
        \begin{subproof}
            Fix $\alpha$.
            Assume $\alpha\in\dom{X}$.
            Thus $\alpha\in\dom{U_3}$.
            Thus $(\alpha,\apply{U_1}{\alpha}\inter\apply{U_2}{\alpha})\in U_3$.
            Thus $\apply{U_3}{\alpha} = \apply{U_1}{\alpha}\inter\apply{U_2}{\alpha}$ by \cref{function_apply_intro}.
        \end{subproof}
        Show that $U_3\in\funs{\dom{X}}{\pow{\unions{\ran{X}}}}$.
        \begin{subproof}
            Show that for all $\alpha\in\dom{X}$ we have $\apply{U_3}{\alpha}\in\pow{\unions{\ran{X}}}$.
            \begin{subproof}
                Fix $\alpha$.
                Assume $\alpha\in\dom{X}$.
                $U_1$ is a function to $\pow{\unions{\ran{X}}}$ by \cref{funs_elim}.
                $U_2$ is a function to $\pow{\unions{\ran{X}}}$ by \cref{funs_elim}.
                Thus $\apply{U_1}{\alpha}\in\pow{\unions{\ran{X}}}$.
                Thus $\apply{U_2}{\alpha}\in\pow{\unions{\ran{X}}}$.
                $\apply{U_1}{\alpha}\subseteq\unions{\ran{X}}$ and $\apply{U_2}{\alpha}\subseteq\unions{\ran{X}}$ by \cref{pow_iff}.
                Thus $\apply{U_3}{\alpha}\subseteq\unions{\ran{X}}$ by \cref{inter_subseteq}.
                Thus $\apply{U_3}{\alpha}\in\pow{\unions{\ran{X}}}$ by \cref{powerset_intro}.
            \end{subproof}
            Thus $U_3$ is a function to $\pow{\unions{\ran{X}}}$.
            $\dom{U_3} = \dom{X}$.
            Thus $U_3\in\funs{\dom{X}}{\pow{\unions{\ran{X}}}}$ by \cref{funs_intro}.
        \end{subproof}
        Show that for all $\alpha\in\dom{X}$ we have $\apply{U_3}{\alpha}\in\apply{T}{\alpha}$.
        \begin{subproof}
            Fix $\alpha$.
            Assume $\alpha\in\dom{X}$.
            $\apply{U_1}{\alpha}\in\apply{T}{\alpha}$.
            $\apply{U_2}{\alpha}\in\apply{T}{\alpha}$.
            Thus $\apply{U_1}{\alpha}\inter\apply{U_2}{\alpha}\in\apply{T}{\alpha}$ by \cref{topology_on}.
            Thus $\apply{U_3}{\alpha}\in\apply{T}{\alpha}$.
        \end{subproof}
        Let $B_3 = \productFamily{U_3}$.
        Show that $B_3\in\pow{\productFamily{X}}$.
        \begin{subproof}
            Show that for all $z$ such that $z\in B_3$ we have $z\in\productFamily{X}$.
            \begin{subproof}
                Fix $z$.
                Assume $z\in B_3$.
                Show that for all $\alpha\in\dom{X}$ we have $\apply{z}{\alpha}\in\apply{X}{\alpha}$.
                \begin{subproof}
                    Fix $\alpha$.
                    Assume $\alpha\in\dom{X}$.
                    Thus $\apply{z}{\alpha}\in\apply{U_3}{\alpha}$ by \cref{indexed_product}.
                    $\apply{T}{\alpha}$ is a topology on $\apply{X}{\alpha}$.
                    Thus $\apply{T}{\alpha}\subseteq\pow{\apply{X}{\alpha}}$ by \cref{topology_on}.
                    Thus $\apply{U_3}{\alpha}\in\pow{\apply{X}{\alpha}}$ by \cref{subseteq}.
                    Thus $\apply{U_3}{\alpha}\subseteq\apply{X}{\alpha}$ by \cref{pow_iff}.
                    Thus $\apply{z}{\alpha}\in\apply{X}{\alpha}$ by \cref{subseteq}.
                \end{subproof}
                Thus $z\in\funs{\dom{U_3}}{\unions{\ran{U_3}}}$ by \cref{indexed_product}.
                Thus $z$ is a function by \cref{funs_elim}.
                $\dom{z} = \dom{U_3}$ by \cref{funs_elim}.
                $\dom{U_3} = \dom{X}$.
                Show that for all $\alpha\in\dom{X}$ we have $\apply{z}{\alpha}\in\unions{\ran{X}}$.
                \begin{subproof}
                    Fix $\alpha$.
                    Assume $\alpha\in\dom{X}$.
                    Thus $\alpha\mathrel{X}\apply{X}{\alpha}$ by \cref{function_apply_elim}.
                    Thus $\apply{X}{\alpha}\in\ran{X}$ by \cref{ran_intro}.
                    Thus $\apply{z}{\alpha}\in\unions{\ran{X}}$ by \cref{unions_iff}.
                \end{subproof}
                Thus $z$ is a function to $\unions{\ran{X}}$.
                Thus $z\in\funs{\dom{X}}{\unions{\ran{X}}}$ by \cref{funs_intro}.
                Thus $z\in\productFamily{X}$ by \cref{indexed_product}.
            \end{subproof}
            Thus $B_3\subseteq\productFamily{X}$ by \cref{subseteq}.
            Thus $B_3\in\pow{\productFamily{X}}$ by \cref{powerset_intro}.
        \end{subproof}
        Thus $B_3\in B$ by \cref{box_basis}.
        Show that $x\in B_3$.
        \begin{subproof}
            $x\in B_1$ by \cref{inter_elim_left}.
            $x\in B_2$ by \cref{inter_elim_right}.
            Thus $B_1\subseteq\productFamily{U_1}$ by \cref{subseteq_refl}.
            Thus $x\in\productFamily{U_1}$ by \cref{elem_subseteq}.
            Thus $B_2\subseteq\productFamily{U_2}$ by \cref{subseteq_refl}.
            Thus $x\in\productFamily{U_2}$ by \cref{elem_subseteq}.
            Show that for all $\alpha\in\dom{X}$ we have $\apply{x}{\alpha}\in\apply{U_3}{\alpha}$.
            \begin{subproof}
                Fix $\alpha$.
                Assume $\alpha\in\dom{X}$.
                Thus $\alpha\in\dom{U_1}$.
                Thus $\alpha\in\dom{U_2}$.
                Thus $\apply{x}{\alpha}\in\apply{U_1}{\alpha}$ by \cref{indexed_product}.
                Thus $\apply{x}{\alpha}\in\apply{U_2}{\alpha}$ by \cref{indexed_product}.
                Thus $\apply{x}{\alpha}\in\apply{U_3}{\alpha}$ by \cref{inter_intro}.
            \end{subproof}
            Thus $x\in\funs{\dom{U_1}}{\unions{\ran{U_1}}}$ by \cref{indexed_product}.
            Thus $x$ is a function by \cref{funs_elim}.
            $\dom{x} = \dom{U_1}$ by \cref{funs_elim}.
            $\dom{U_1} = \dom{X}$.
            Show that for all $\alpha\in\dom{X}$ we have $\apply{x}{\alpha}\in\unions{\ran{U_3}}$.
            \begin{subproof}
                Fix $\alpha$.
                Assume $\alpha\in\dom{X}$.
                $\dom{U_3} = \dom{X}$.
                Thus $\alpha\in\dom{U_3}$.
                $U_3$ is a function.
                Thus $\alpha\mathrel{U_3}\apply{U_3}{\alpha}$ by \cref{function_apply_elim}.
                Thus $\apply{U_3}{\alpha}\in\ran{U_3}$ by \cref{ran_intro}.
                Thus $\apply{x}{\alpha}\in\unions{\ran{U_3}}$ by \cref{unions_iff}.
            \end{subproof}
            Thus $x$ is a function to $\unions{\ran{U_3}}$.
            Thus $x\in\funs{\dom{X}}{\unions{\ran{U_3}}}$ by \cref{funs_intro}.
            Thus $x\in B_3$ by \cref{indexed_product}.
        \end{subproof}
        Show that $B_3\subseteq B_1\inter B_2$.
        \begin{subproof}
            Show that for all $z$ such that $z\in B_3$ we have $z\in B_1\inter B_2$.
            \begin{subproof}
                Fix $z$.
                Assume $z\in B_3$.
                Show that $z\in B_1$.
                \begin{subproof}
                    Show that for all $\alpha\in\dom{X}$ we have $\apply{z}{\alpha}\in\apply{U_1}{\alpha}$.
                    \begin{subproof}
                        Fix $\alpha$.
                        Assume $\alpha\in\dom{X}$.
                        Thus $\apply{z}{\alpha}\in\apply{U_3}{\alpha}$ by \cref{indexed_product}.
                        Thus $\apply{z}{\alpha}\in\apply{U_1}{\alpha}$ by \cref{inter_elim_left}.
                    \end{subproof}
                    Thus $z\in\funs{\dom{U_3}}{\unions{\ran{U_3}}}$ by \cref{indexed_product}.
                    Thus $z$ is a function by \cref{funs_elim}.
                    $\dom{z} = \dom{U_3}$ by \cref{funs_elim}.
                    $\dom{U_3} = \dom{X}$.
                    Show that for all $\alpha\in\dom{X}$ we have $\apply{z}{\alpha}\in\unions{\ran{U_1}}$.
                    \begin{subproof}
                        Fix $\alpha$.
                        Assume $\alpha\in\dom{X}$.
                        $\alpha\mathrel{U_1}\apply{U_1}{\alpha}$ by \cref{function_apply_elim,funs_elim}.
                        Thus $\apply{U_1}{\alpha}\in\ran{U_1}$ by \cref{ran_intro}.
                        Thus $\apply{z}{\alpha}\in\unions{\ran{U_1}}$ by \cref{unions_iff}.
                    \end{subproof}
                    Thus $z$ is a function to $\unions{\ran{U_1}}$.
                    Thus $z\in\funs{\dom{X}}{\unions{\ran{U_1}}}$ by \cref{funs_intro}.
                    Thus $z\in B_1$ by \cref{indexed_product}.
                \end{subproof}
                Show that $z\in B_2$.
                \begin{subproof}
                    Show that for all $\alpha\in\dom{X}$ we have $\apply{z}{\alpha}\in\apply{U_2}{\alpha}$.
                    \begin{subproof}
                        Fix $\alpha$.
                        Assume $\alpha\in\dom{X}$.
                        Thus $\apply{z}{\alpha}\in\apply{U_3}{\alpha}$ by \cref{indexed_product}.
                        Thus $\apply{z}{\alpha}\in\apply{U_2}{\alpha}$ by \cref{inter_elim_right}.
                    \end{subproof}
                    Thus $z\in\funs{\dom{U_3}}{\unions{\ran{U_3}}}$ by \cref{indexed_product}.
                    Thus $z$ is a function by \cref{funs_elim}.
                    $\dom{z} = \dom{U_3}$ by \cref{funs_elim}.
                    $\dom{U_3} = \dom{X}$.
                    Show that for all $\alpha\in\dom{X}$ we have $\apply{z}{\alpha}\in\unions{\ran{U_2}}$.
                    \begin{subproof}
                        Fix $\alpha$.
                        Assume $\alpha\in\dom{X}$.
                        $\alpha\mathrel{U_2}\apply{U_2}{\alpha}$ by \cref{function_apply_elim,funs_elim}.
                        Thus $\apply{U_2}{\alpha}\in\ran{U_2}$ by \cref{ran_intro}.
                        Thus $\apply{z}{\alpha}\in\unions{\ran{U_2}}$ by \cref{unions_iff}.
                    \end{subproof}
                    Thus $z$ is a function to $\unions{\ran{U_2}}$.
                    Thus $z\in\funs{\dom{X}}{\unions{\ran{U_2}}}$ by \cref{funs_intro}.
                    Thus $z\in B_2$ by \cref{indexed_product}.
                \end{subproof}
                Thus $z\in B_1\inter B_2$ by \cref{inter_intro}.
            \end{subproof}
            Thus $B_3\subseteq B_1\inter B_2$ by \cref{subseteq}.
        \end{subproof}
    \end{subproof}
    Thus $B$ is a basis on $\productFamily{X}$ by \cref{basis_on}.
\end{proof}

% from §19 Item 11: product subbasis is a subbasis on the indexed product
\begin{lemma}\label{product_subbasis_is_subbasis}
    Assume $X$ is a function.
    Assume $T$ is a function.
    Assume $\dom{T} = \dom{X}$.
    Assume $\dom{X}$ is inhabited.
    Assume that for all $\alpha\in\dom{X}$ we have $\apply{T}{\alpha}$ is a topology on $\apply{X}{\alpha}$.
    Then $\productSubbasis{T}{X}$ is a subbasis on $\productFamily{X}$.
\end{lemma}
\begin{proof}
    Let $S = \productSubbasis{T}{X}$.
    Show that for all $A$ such that $A\in S$ we have $A\in\pow{\productFamily{X}}$.
    \begin{subproof}
        Fix $A$.
        Assume $A\in S$.
        Take $\beta\in\dom{X}$ such that $A\in\productSubbasisAt{T}{X}{\beta}$ by \cref{product_subbasis_union_indexed}.
        Thus $A\in\pow{\productFamily{X}}$ by \cref{product_subbasis_indexed}.
    \end{subproof}
    Thus $S\subseteq\pow{\productFamily{X}}$ by \cref{subseteq}.
    Take $\beta$ such that $\beta\in\dom{X}$.
    Show that for all $x$ such that $x\in\productFamily{X}$ we have $x\in\unions{S}$.
    \begin{subproof}
        Fix $x$.
        Assume $x\in\productFamily{X}$.
        $\apply{T}{\beta}$ is a topology on $\apply{X}{\beta}$.
        Thus $\apply{X}{\beta}\in\apply{T}{\beta}$ by \cref{topology_on}.
        Let $A = \preimg{\piIndex{X}{\beta}}{\apply{X}{\beta}}$.
        Show that for all $z$ such that $z\in A$ we have $z\in\productFamily{X}$.
        \begin{subproof}
            Fix $z$.
            Assume $z\in A$.
            Take $y\in\apply{X}{\beta}$ such that $z\mathrel{\piIndex{X}{\beta}} y$ by \cref{preimg_iff}.
            Thus $(z,y)\in\piIndex{X}{\beta}$.
            Take $x\in\productFamily{X}$ such that $(z,y) = (x,\apply{x}{\beta})$ by \cref{projection_map_indexed}.
            Thus $z = x$ by \cref{pair_eq_iff}.
            Thus $z\in\productFamily{X}$.
        \end{subproof}
        Thus $A\subseteq\productFamily{X}$ by \cref{subseteq}.
        Thus $A\in\pow{\productFamily{X}}$ by \cref{powerset_intro}.
        $A\in\productSubbasisAt{T}{X}{\beta}$ by \cref{product_subbasis_indexed}.
        Thus $A\in S$ by \cref{product_subbasis_union_indexed}.
        $\apply{x}{\beta}\in\apply{X}{\beta}$ by \cref{indexed_product}.
        Thus $(x,\apply{x}{\beta})\in\piIndex{X}{\beta}$ by \cref{projection_map_indexed}.
        Thus $x\mathrel{\piIndex{X}{\beta}}\apply{x}{\beta}$.
        Thus $x\in A$ by \cref{preimg_iff}.
        Thus $x\in\unions{S}$ by \cref{unions_iff}.
    \end{subproof}
    Thus $\productFamily{X}\subseteq\unions{S}$ by \cref{subseteq}.
    $\unions{S}\subseteq\productFamily{X}$ by \cref{unions_subseteq_of_powerset_is_subseteq,subseteq}.
    Thus $\unions{S} = \productFamily{X}$ by \cref{subseteq_antisymmetric}.
    Thus $S$ is a subbasis on $\productFamily{X}$ by \cref{subbasis_on}.
\end{proof}

% from §19 Item 12: finite intersections of product subbasis elements
\begin{lemma}\label{product_subbasis_finite_intersections_basis}
    Assume $X$ is a function.
    Assume $T$ is a function.
    Assume $\dom{T} = \dom{X}$.
    Assume $\dom{X}$ is inhabited.
    Assume that for all $\alpha\in\dom{X}$ we have $\apply{T}{\alpha}$ is a topology on $\apply{X}{\alpha}$.
    Then $\finiteIntersections{\productSubbasis{T}{X}}$ is a basis for $\productTopologyIndexed{T}{X}$ on $\productFamily{X}$.
\end{lemma}
\begin{proof}
    Let $S = \productSubbasis{T}{X}$.
    $S$ is a subbasis on $\productFamily{X}$ by \cref{product_subbasis_is_subbasis}.
    Thus $\finiteIntersections{S}$ is a basis on $\productFamily{X}$ by \cref{subbasis_finite_intersections_basis}.
    Let $T_1 = \productTopologyIndexed{T}{X}$.
    $T_1 = \subbasisTopology{S}{\productFamily{X}}$ by \cref{product_topology_indexed}.
    Thus $\finiteIntersections{S}$ is a basis for $T_1$ on $\productFamily{X}$ by \cref{basis_for_topology,topology_generated_by_subbasis}.
\end{proof}

% from §19 helper lemma 1: choice function on a relation
\begin{axiom}\label{choice_function}
    Suppose $R$ is a relation.
    Suppose for all $x\in A$ there exists $y$ such that $x\mathrel{R} y$.
    Then there exists $f$ such that $f$ is a function on $A$ and for all $x\in A$ we have $x\mathrel{R}\apply{f}{x}$.
\end{axiom}

% from §19 Item 13: basis for the box topology from bases
\begin{lemma}\label{box_basis_from_bases}
    Assume $X$ is a function.
    Assume $T$ is a function.
    Assume $B$ is a function.
    Assume $\dom{T} = \dom{X}$.
    Assume $\dom{B} = \dom{X}$.
    Assume that for all $\alpha\in\dom{X}$ we have $\apply{B}{\alpha}$ is a basis for $\apply{T}{\alpha}$ on $\apply{X}{\alpha}$.
    Let $C = \{A\in\pow{\productFamily{X}} \mid
        \exists U. U\in\funs{\dom{X}}{\pow{\unions{\ran{X}}}}\land
        (\forall \alpha\in\dom{X}. \apply{U}{\alpha}\in\apply{B}{\alpha})\land
        A = \productFamily{U}\}$.
    Then $C$ is a basis for $\boxTopology{T}{X}$ on $\productFamily{X}$.
\end{lemma}
\begin{proof}
    Let $T_0 = \boxTopology{T}{X}$.
    Let $B_0 = \boxBasis{T}{X}$.
    Show that for all $\alpha\in\dom{X}$ we have $\apply{T}{\alpha}$ is a topology on $\apply{X}{\alpha}$.
    \begin{subproof}
        Fix $\alpha$.
        Assume $\alpha\in\dom{X}$.
        $\apply{B}{\alpha}$ is a basis for $\apply{T}{\alpha}$ on $\apply{X}{\alpha}$.
        Thus $\apply{B}{\alpha}$ is a basis on $\apply{X}{\alpha}$ by \cref{basis_for_topology}.
        Thus $\apply{T}{\alpha} = \basisTopology{\apply{B}{\alpha}}{\apply{X}{\alpha}}$ by \cref{basis_for_topology}.
        Thus $\apply{T}{\alpha}$ is a topology on $\apply{X}{\alpha}$ by \cref{basis_topology_is_topology}.
    \end{subproof}
    $B_0$ is a basis on $\productFamily{X}$ by \cref{box_basis_is_basis}.
    Thus $T_0$ is a topology on $\productFamily{X}$ by \cref{basis_topology_is_topology,box_topology}.
    Show that for all $A$ such that $A\in C$ we have $A$ is open in $\productFamily{X}$ with respect to $T_0$.
    \begin{subproof}
        Fix $A$.
        Assume $A\in C$.
        Take $U$ such that
            $U\in\funs{\dom{X}}{\pow{\unions{\ran{X}}}}$ and
            for all $\alpha\in\dom{X}$ we have $\apply{U}{\alpha}\in\apply{B}{\alpha}$ and
            $A = \productFamily{U}$.
        Show that for all $\alpha\in\dom{X}$ we have $\apply{U}{\alpha}\in\apply{T}{\alpha}$.
        \begin{subproof}
        Fix $\alpha$.
        Assume $\alpha\in\dom{X}$.
        $\apply{B}{\alpha}$ is a basis for $\apply{T}{\alpha}$ on $\apply{X}{\alpha}$.
        Thus $\apply{B}{\alpha}$ is a basis on $\apply{X}{\alpha}$ by \cref{basis_for_topology}.
        Thus $\apply{T}{\alpha} = \basisTopology{\apply{B}{\alpha}}{\apply{X}{\alpha}}$ by \cref{basis_for_topology}.
        Thus $\apply{U}{\alpha}\in\basisTopology{\apply{B}{\alpha}}{\apply{X}{\alpha}}$ by \cref{basis_elem_open_in_generated}.
            Thus $\apply{U}{\alpha}\in\apply{T}{\alpha}$ by \cref{basis_for_topology}.
        \end{subproof}
        Thus $A\in B_0$ by \cref{box_basis}.
        Thus $A\in T_0$ by \cref{basis_elem_open_in_generated,box_topology}.
        Thus $A$ is open in $\productFamily{X}$ with respect to $T_0$ by \cref{open_in}.
    \end{subproof}
    Show that for all $U,x$ such that $U\in T_0$ and $x\in U$ there exists $C_1\in C$ such that $x\in C_1$ and $C_1\subseteq U$.
    \begin{subproof}
        Fix $U$.
        Fix $x$.
        Assume $U\in T_0$ and $x\in U$.
        $U\in\basisTopology{B_0}{\productFamily{X}}$ by \cref{box_topology}.
        Take $B_1\in B_0$ such that $x\in B_1$ and $B_1\subseteq U$ by \cref{topology_generated_by_basis}.
        Take $V$ such that
            $V\in\funs{\dom{X}}{\pow{\unions{\ran{X}}}}$ and
            $B_1 = \productFamily{V}$ and
            for all $\alpha\in\dom{X}$ we have $\apply{V}{\alpha}\in\apply{T}{\alpha}$
            by \cref{box_basis}.
        $\dom{V} = \dom{X}$ by \cref{funs_elim}.
        Thus $x\in\productFamily{V}$.
        Show that for all $\alpha\in\dom{X}$ there exists $W$ such that
            $W\in\apply{B}{\alpha}$ and $\apply{x}{\alpha}\in W$ and $W\subseteq\apply{V}{\alpha}$.
        \begin{subproof}
            Fix $\alpha$.
            Assume $\alpha\in\dom{X}$.
            Thus $\alpha\in\dom{V}$.
            Thus $\apply{x}{\alpha}\in\apply{V}{\alpha}$ by \cref{indexed_product}.
            $\apply{V}{\alpha}\in\apply{T}{\alpha}$.
            Thus $\apply{V}{\alpha}$ is open in $\apply{X}{\alpha}$ with respect to $\apply{T}{\alpha}$ by \cref{open_in}.
            $\apply{B}{\alpha}$ is a basis for $\apply{T}{\alpha}$ on $\apply{X}{\alpha}$.
            Thus there exists $W\in\apply{B}{\alpha}$ such that $\apply{x}{\alpha}\in W$ and $W\subseteq\apply{V}{\alpha}$
                by \cref{basis_for_topology,topology_generated_by_basis,open_in}.
        \end{subproof}
        Let $R = \{r\in\dom{X}\times\pow{\unions{\ran{X}}} \mid \apply{x}{\fst{r}}\in\snd{r} \land \snd{r}\subseteq\apply{V}{\fst{r}} \land \snd{r}\in\apply{B}{\fst{r}}\}$.
        Show that $R$ is a relation.
        \begin{subproof}
            Show that for all $w$ such that $w\in R$ there exist $y,z$ such that $w = (y,z)$.
            \begin{subproof}
                Fix $w$.
                Assume $w\in R$.
                Thus $w\in\dom{X}\times\pow{\unions{\ran{X}}}$.
                Take $y,z$ such that $y\in\dom{X}$ and $z\in\pow{\unions{\ran{X}}}$ and $w = (y,z)$ by \cref{times_elem_is_tuple}.
            \end{subproof}
            Thus $R$ is a relation by \cref{relation}.
        \end{subproof}
        Let $A = \dom{X}$.
        Show that for all $\alpha\in A$ there exists $W$ such that $\alpha\mathrel{R} W$.
        \begin{subproof}
            Fix $\alpha$.
            Assume $\alpha\in A$.
            Take $W$ such that $W\in\apply{B}{\alpha}$ and $\apply{x}{\alpha}\in W$ and $W\subseteq\apply{V}{\alpha}$.
            $\apply{V}{\alpha}\in\apply{T}{\alpha}$.
            $\apply{T}{\alpha}$ is a topology on $\apply{X}{\alpha}$.
            Thus $\apply{T}{\alpha}\subseteq\pow{\apply{X}{\alpha}}$ by \cref{topology_on}.
            Thus $\apply{V}{\alpha}\in\pow{\apply{X}{\alpha}}$ by \cref{subseteq}.
            Thus $\apply{V}{\alpha}\subseteq\apply{X}{\alpha}$ by \cref{pow_iff}.
            Thus $W\subseteq\apply{X}{\alpha}$ by \cref{subseteq}.
            Thus $\alpha\mathrel{X}\apply{X}{\alpha}$ by \cref{function_apply_elim}.
            Thus $\apply{X}{\alpha}\in\ran{X}$ by \cref{ran_intro}.
            Thus $W\subseteq\unions{\ran{X}}$ by \cref{unions_iff,subseteq}.
            Thus $W\in\pow{\unions{\ran{X}}}$ by \cref{powerset_intro}.
            Thus $(\alpha,W)\in\dom{X}\times\pow{\unions{\ran{X}}}$ by \cref{times_tuple_intro}.
            Thus $\apply{x}{\fst{(\alpha,W)}}\in\snd{(\alpha,W)}$ by \cref{fst_eq,snd_eq}.
            Thus $\snd{(\alpha,W)}\subseteq\apply{V}{\fst{(\alpha,W)}}$ by \cref{fst_eq,snd_eq}.
            Thus $\snd{(\alpha,W)}\in\apply{B}{\fst{(\alpha,W)}}$ by \cref{fst_eq,snd_eq}.
            Thus $(\alpha,W)\in R$.
            Thus $\alpha\mathrel{R} W$.
        \end{subproof}
        There exists $U_1$ such that $U_1$ is a function on $A$ and for all $\alpha\in A$ we have $\alpha\mathrel{R}\apply{U_1}{\alpha}$
            by \cref{choice_function}.
        Take $U_1$ such that $U_1$ is a function on $A$ and for all $\alpha\in A$ we have $\alpha\mathrel{R}\apply{U_1}{\alpha}$.
        Thus $\dom{U_1} = A$.
        Show that for all $\alpha\in\dom{X}$ we have $\apply{U_1}{\alpha}\in\apply{B}{\alpha}$ and
            $\apply{x}{\alpha}\in\apply{U_1}{\alpha}$ and $\apply{U_1}{\alpha}\subseteq\apply{V}{\alpha}$.
        \begin{subproof}
            Fix $\alpha$.
            Assume $\alpha\in\dom{X}$.
            Thus $(\alpha,\apply{U_1}{\alpha})\in R$.
            Thus $\apply{x}{\fst{(\alpha,\apply{U_1}{\alpha})}}\in\snd{(\alpha,\apply{U_1}{\alpha})}$.
            Thus $\snd{(\alpha,\apply{U_1}{\alpha})}\subseteq\apply{V}{\fst{(\alpha,\apply{U_1}{\alpha})}}$.
            Thus $\snd{(\alpha,\apply{U_1}{\alpha})}\in\apply{B}{\fst{(\alpha,\apply{U_1}{\alpha})}}$.
            Thus $\apply{x}{\alpha}\in\apply{U_1}{\alpha}$ by \cref{fst_eq,snd_eq}.
            Thus $\apply{U_1}{\alpha}\subseteq\apply{V}{\alpha}$ by \cref{fst_eq,snd_eq}.
            Thus $\apply{U_1}{\alpha}\in\apply{B}{\alpha}$ by \cref{fst_eq,snd_eq}.
        \end{subproof}
        Show that $U_1\in\funs{\dom{X}}{\pow{\unions{\ran{X}}}}$.
        \begin{subproof}
            Show that for all $\alpha\in\dom{X}$ we have $\apply{U_1}{\alpha}\in\pow{\unions{\ran{X}}}$.
            \begin{subproof}
                Fix $\alpha$.
                Assume $\alpha\in\dom{X}$.
                $\apply{V}{\alpha}\in\apply{T}{\alpha}$.
                $\apply{T}{\alpha}$ is a topology on $\apply{X}{\alpha}$.
                Thus $\apply{T}{\alpha}\subseteq\pow{\apply{X}{\alpha}}$ by \cref{topology_on}.
                Thus $\apply{V}{\alpha}\in\pow{\apply{X}{\alpha}}$ by \cref{subseteq}.
                Thus $\apply{V}{\alpha}\subseteq\apply{X}{\alpha}$ by \cref{pow_iff}.
                Thus $\apply{U_1}{\alpha}\subseteq\apply{X}{\alpha}$ by \cref{subseteq}.
                Thus $\alpha\mathrel{X}\apply{X}{\alpha}$ by \cref{function_apply_elim}.
                Thus $\apply{X}{\alpha}\in\ran{X}$ by \cref{ran_intro}.
                Thus $\apply{U_1}{\alpha}\subseteq\unions{\ran{X}}$ by \cref{unions_iff,subseteq}.
                Thus $\apply{U_1}{\alpha}\in\pow{\unions{\ran{X}}}$ by \cref{powerset_intro}.
            \end{subproof}
            Thus $U_1$ is a function to $\pow{\unions{\ran{X}}}$.
            Thus $U_1\in\funs{\dom{X}}{\pow{\unions{\ran{X}}}}$ by \cref{funs_intro}.
        \end{subproof}
        Let $C_1 = \productFamily{U_1}$.
        Show that $C_1\in\pow{\productFamily{X}}$.
        \begin{subproof}
            Show that for all $z$ such that $z\in C_1$ we have $z\in\productFamily{X}$.
            \begin{subproof}
                Fix $z$.
                Assume $z\in C_1$.
                Show that for all $\alpha\in\dom{X}$ we have $\apply{z}{\alpha}\in\apply{X}{\alpha}$.
                \begin{subproof}
                    Fix $\alpha$.
                    Assume $\alpha\in\dom{X}$.
                    Thus $\alpha\in\dom{U_1}$.
                    Thus $\apply{z}{\alpha}\in\apply{U_1}{\alpha}$ by \cref{indexed_product}.
                    $\apply{V}{\alpha}\in\apply{T}{\alpha}$.
                    $\apply{T}{\alpha}$ is a topology on $\apply{X}{\alpha}$.
                    Thus $\apply{T}{\alpha}\subseteq\pow{\apply{X}{\alpha}}$ by \cref{topology_on}.
                    Thus $\apply{V}{\alpha}\in\pow{\apply{X}{\alpha}}$ by \cref{subseteq}.
                    Thus $\apply{V}{\alpha}\subseteq\apply{X}{\alpha}$ by \cref{pow_iff}.
                    Thus $\apply{U_1}{\alpha}\subseteq\apply{X}{\alpha}$ by \cref{subseteq_transitive,subseteq}.
                    Thus $\apply{z}{\alpha}\in\apply{X}{\alpha}$ by \cref{subseteq}.
                \end{subproof}
                Thus $z\in\funs{\dom{U_1}}{\unions{\ran{U_1}}}$ by \cref{indexed_product}.
                Thus $z$ is a function by \cref{funs_elim}.
                $\dom{z} = \dom{U_1}$ by \cref{funs_elim}.
                $\dom{U_1} = \dom{X}$.
                Show that for all $\alpha\in\dom{X}$ we have $\apply{z}{\alpha}\in\unions{\ran{X}}$.
                \begin{subproof}
                    Fix $\alpha$.
                    Assume $\alpha\in\dom{X}$.
                    Thus $\alpha\mathrel{X}\apply{X}{\alpha}$ by \cref{function_apply_elim}.
                    Thus $\apply{X}{\alpha}\in\ran{X}$ by \cref{ran_intro}.
                    Thus $\apply{z}{\alpha}\in\unions{\ran{X}}$ by \cref{unions_iff}.
                \end{subproof}
                Thus $z$ is a function to $\unions{\ran{X}}$.
                Thus $z\in\funs{\dom{X}}{\unions{\ran{X}}}$ by \cref{funs_intro}.
                Thus $z\in\productFamily{X}$ by \cref{indexed_product}.
            \end{subproof}
            Thus $C_1\subseteq\productFamily{X}$ by \cref{subseteq}.
            Thus $C_1\in\pow{\productFamily{X}}$ by \cref{powerset_intro}.
        \end{subproof}
        Thus $C_1\in C$.
        Show that $x\in C_1$.
        \begin{subproof}
            Show that for all $\alpha\in\dom{X}$ we have $\apply{x}{\alpha}\in\apply{U_1}{\alpha}$.
            \begin{subproof}
                Fix $\alpha$.
                Assume $\alpha\in\dom{X}$.
                Thus $\apply{x}{\alpha}\in\apply{U_1}{\alpha}$.
            \end{subproof}
            Thus $x\in\funs{\dom{V}}{\unions{\ran{V}}}$ by \cref{indexed_product}.
            Thus $x$ is a function by \cref{funs_elim}.
            $\dom{x} = \dom{V}$ by \cref{funs_elim}.
            $\dom{V} = \dom{X}$.
            $\dom{U_1} = \dom{X}$.
            Show that for all $\alpha\in\dom{U_1}$ we have $\apply{x}{\alpha}\in\unions{\ran{U_1}}$.
            \begin{subproof}
                Fix $\alpha$.
                Assume $\alpha\in\dom{U_1}$.
                Thus $\alpha\mathrel{U_1}\apply{U_1}{\alpha}$ by \cref{function_apply_elim}.
                Thus $\apply{U_1}{\alpha}\in\ran{U_1}$ by \cref{ran_intro}.
                Thus $\apply{x}{\alpha}\in\unions{\ran{U_1}}$ by \cref{unions_iff}.
            \end{subproof}
            Thus $x$ is a function to $\unions{\ran{U_1}}$.
            Thus $x\in\funs{\dom{U_1}}{\unions{\ran{U_1}}}$ by \cref{funs_intro}.
            Thus $x\in C_1$ by \cref{indexed_product}.
        \end{subproof}
        Show that $C_1\subseteq U$.
        \begin{subproof}
            Show that for all $z$ such that $z\in C_1$ we have $z\in U$.
            \begin{subproof}
                Fix $z$.
                Assume $z\in C_1$.
                Show that $z\in B_1$.
                \begin{subproof}
                    Show that for all $\alpha\in\dom{X}$ we have $\apply{z}{\alpha}\in\apply{V}{\alpha}$.
                    \begin{subproof}
                        Fix $\alpha$.
                        Assume $\alpha\in\dom{X}$.
                        Thus $\alpha\in\dom{U_1}$.
                        Thus $\apply{z}{\alpha}\in\apply{U_1}{\alpha}$ by \cref{indexed_product}.
                        Thus $\apply{z}{\alpha}\in\apply{V}{\alpha}$ by \cref{subseteq}.
                    \end{subproof}
                    Thus $z\in\funs{\dom{U_1}}{\unions{\ran{U_1}}}$ by \cref{indexed_product}.
                    Thus $z$ is a function by \cref{funs_elim}.
                    $\dom{z} = \dom{U_1}$ by \cref{funs_elim}.
                    $\dom{U_1} = \dom{X}$.
                    Show that for all $\alpha\in\dom{X}$ we have $\apply{z}{\alpha}\in\unions{\ran{V}}$.
                    \begin{subproof}
                    Fix $\alpha$.
                    Assume $\alpha\in\dom{X}$.
                    $V$ is a function by \cref{funs_elim}.
                    $\dom{V} = \dom{X}$ by \cref{funs_elim}.
                    Thus $\alpha\in\dom{V}$.
                    Thus $\alpha\mathrel{V}\apply{V}{\alpha}$ by \cref{function_apply_elim}.
                        Thus $\apply{V}{\alpha}\in\ran{V}$ by \cref{ran_intro}.
                        Thus $\apply{z}{\alpha}\in\unions{\ran{V}}$ by \cref{unions_iff}.
                    \end{subproof}
                    Thus $z$ is a function to $\unions{\ran{V}}$.
                    Thus $z\in\funs{\dom{X}}{\unions{\ran{V}}}$ by \cref{funs_intro}.
                    Thus $z\in B_1$ by \cref{indexed_product}.
                \end{subproof}
                Thus $z\in U$ by \cref{subseteq}.
            \end{subproof}
            Thus $C_1\subseteq U$ by \cref{subseteq}.
        \end{subproof}
    \end{subproof}
    Thus $C$ is a basis for $T_0$ on $\productFamily{X}$ by \cref{open_collection_basis}.
\end{proof}

% from §19 Item 14: basis for the product topology from bases
\begin{lemma}\label{product_basis_from_bases_indexed}
    Assume $X$ is a function.
    Assume $T$ is a function.
    Assume $B$ is a function.
    Assume $\dom{T} = \dom{X}$.
    Assume $\dom{B} = \dom{X}$.
    Assume $\dom{X}$ is inhabited.
    Assume that for all $\alpha\in\dom{X}$ we have $\apply{B}{\alpha}$ is a basis for $\apply{T}{\alpha}$ on $\apply{X}{\alpha}$.
    Let $S$ be a set such that for all $A$ we have
        $A\in S$ iff there exists $\beta\in\dom{X}$ such that there exists $U\in\apply{B}{\beta}$ such that
            $A = \preimg{\piIndex{X}{\beta}}{U}$.
    Then $\finiteIntersections{S}$ is a basis for $\productTopologyIndexed{T}{X}$ on $\productFamily{X}$.
\end{lemma}
\begin{proof}
    Let $T_0 = \productTopologyIndexed{T}{X}$.
    Let $T_1 = \subbasisTopology{S}{\productFamily{X}}$.
    Show that $S$ is a subbasis on $\productFamily{X}$.
    \begin{subproof}
        Show that for all $A$ such that $A\in S$ we have $A\in\pow{\productFamily{X}}$.
        \begin{subproof}
            Fix $A$.
            Assume $A\in S$.
            Take $\beta$ such that $\beta\in\dom{X}$ and there exists $U\in\apply{B}{\beta}$ such that
                $A = \preimg{\piIndex{X}{\beta}}{U}$.
            Take $U\in\apply{B}{\beta}$ such that $A = \preimg{\piIndex{X}{\beta}}{U}$.
            $\apply{B}{\beta}$ is a basis for $\apply{T}{\beta}$ on $\apply{X}{\beta}$.
            Thus $\apply{B}{\beta}$ is a basis on $\apply{X}{\beta}$ by \cref{basis_for_topology}.
            Thus $U\in\basisTopology{\apply{B}{\beta}}{\apply{X}{\beta}}$ by \cref{basis_elem_open_in_generated}.
            Thus $U\in\apply{T}{\beta}$ by \cref{basis_for_topology}.
            Show that $A\subseteq\productFamily{X}$.
            \begin{subproof}
                Show that for all $z$ such that $z\in A$ we have $z\in\productFamily{X}$.
                \begin{subproof}
                    Fix $z$.
                    Assume $z\in A$.
                    Take $y\in U$ such that $z\mathrel{\piIndex{X}{\beta}} y$ by \cref{preimg_iff}.
                    Thus $(z,y)\in\piIndex{X}{\beta}$.
                    Take $x\in\productFamily{X}$ such that $(z,y) = (x,\apply{x}{\beta})$ by \cref{projection_map_indexed}.
                    Thus $z = x$ by \cref{pair_eq_iff}.
                    Thus $z\in\productFamily{X}$.
                \end{subproof}
                Thus $A\subseteq\productFamily{X}$ by \cref{subseteq}.
            \end{subproof}
            Thus $A\in\pow{\productFamily{X}}$ by \cref{powerset_intro}.
        \end{subproof}
        Thus $S\subseteq\pow{\productFamily{X}}$ by \cref{subseteq}.
        Show that $\productFamily{X}\subseteq\unions{S}$.
        \begin{subproof}
            Show that for all $x$ such that $x\in\productFamily{X}$ we have $x\in\unions{S}$.
            \begin{subproof}
                Fix $x$.
                Assume $x\in\productFamily{X}$.
                Take $\beta$ such that $\beta\in\dom{X}$.
                $\apply{x}{\beta}\in\apply{X}{\beta}$ by \cref{indexed_product}.
                $\apply{B}{\beta}$ is a basis for $\apply{T}{\beta}$ on $\apply{X}{\beta}$.
                Thus $\apply{B}{\beta}$ is a basis on $\apply{X}{\beta}$ by \cref{basis_for_topology}.
                Take $U\in\apply{B}{\beta}$ such that $\apply{x}{\beta}\in U$ by \cref{basis_on}.
                Let $A = \preimg{\piIndex{X}{\beta}}{U}$.
                $A\in S$.
                Thus $(x,\apply{x}{\beta})\in\piIndex{X}{\beta}$ by \cref{projection_map_indexed}.
                Thus $x\mathrel{\piIndex{X}{\beta}}\apply{x}{\beta}$.
                Thus $x\in A$ by \cref{preimg_iff}.
                Thus $x\in\unions{S}$ by \cref{unions_iff}.
            \end{subproof}
            Thus $\productFamily{X}\subseteq\unions{S}$ by \cref{subseteq}.
        \end{subproof}
        $\unions{S}\subseteq\productFamily{X}$ by \cref{unions_subseteq_of_powerset_is_subseteq,subseteq}.
        Thus $\unions{S} = \productFamily{X}$ by \cref{subseteq_antisymmetric}.
    \end{subproof}
    Thus $S$ is a subbasis on $\productFamily{X}$ by \cref{subbasis_on}.
    Thus $\finiteIntersections{S}$ is a basis on $\productFamily{X}$ by \cref{subbasis_finite_intersections_basis}.
    Thus $\finiteIntersections{S}$ is a basis for $T_1$ on $\productFamily{X}$ by \cref{basis_for_topology,topology_generated_by_subbasis}.
    Show that for all $C$ such that $C\in S$ we have $C\in T_1$.
    \begin{subproof}
        Fix $C$.
        Assume $C\in S$.
        $S$ is a subbasis on $\productFamily{X}$ by \cref{subbasis_on}.
        Thus $S\subseteq\pow{\productFamily{X}}$ by \cref{subbasis_on}.
        Thus $C\in\pow{\productFamily{X}}$ by \cref{subseteq}.
        Thus $C\subseteq\productFamily{X}$ by \cref{pow_iff}.
        $\unions{S} = \productFamily{X}$ by \cref{subbasis_on}.
        Thus $C\subseteq\unions{S}$ by \cref{subseteq}.
        Thus $C\in\pow{\unions{S}}$ by \cref{powerset_intro}.
        Let $F = \{C\}$.
        $F\subseteq S$ by \cref{singleton_subset_intro}.
        Let $G = \{C,C\}$.
        $G$ is finite by \cref{pair_is_finite}.
        Show that $F\subseteq G$.
        \begin{subproof}
            Show that for all $x$ such that $x\in F$ we have $x\in G$.
            \begin{subproof}
                Fix $x$.
                Assume $x\in F$.
                Thus $x = C$ by \cref{singleton_elim}.
                Thus $x\in G$ by \cref{upair_intro_left}.
            \end{subproof}
            Thus $F\subseteq G$ by \cref{subseteq}.
        \end{subproof}
        Thus $F$ is finite by \cref{subseteq_finite_is_finite}.
        $F$ is inhabited by \cref{singleton_intro}.
        $\inters{F} = C$ by \cref{inters_singleton}.
        Thus $C\in\finiteIntersections{S}$ by \cref{finite_intersections}.
        Thus $C\in T_1$ by \cref{basis_elem_open_in_generated,basis_for_topology}.
    \end{subproof}
    Thus $S\subseteq T_1$ by \cref{subseteq}.
    Show that for all $\alpha\in\dom{X}$ we have $\apply{T}{\alpha}$ is a topology on $\apply{X}{\alpha}$.
    \begin{subproof}
        Fix $\alpha$.
        Assume $\alpha\in\dom{X}$.
        $\apply{B}{\alpha}$ is a basis for $\apply{T}{\alpha}$ on $\apply{X}{\alpha}$.
        Thus $\apply{B}{\alpha}$ is a basis on $\apply{X}{\alpha}$ by \cref{basis_for_topology}.
        Thus $\apply{T}{\alpha} = \basisTopology{\apply{B}{\alpha}}{\apply{X}{\alpha}}$ by \cref{basis_for_topology}.
        Thus $\apply{T}{\alpha}$ is a topology on $\apply{X}{\alpha}$ by \cref{basis_topology_is_topology}.
    \end{subproof}
    Let $S_0 = \productSubbasis{T}{X}$.
    $T_0 = \subbasisTopology{S_0}{\productFamily{X}}$ by \cref{product_topology_indexed}.
    $S_0$ is a subbasis on $\productFamily{X}$ by \cref{product_subbasis_is_subbasis}.
    Thus $\finiteIntersections{S_0}$ is a basis for $T_0$ on $\productFamily{X}$ by \cref{product_subbasis_finite_intersections_basis}.
    Thus $T_0$ is a topology on $\productFamily{X}$ by \cref{basis_topology_is_topology,basis_for_topology}.
    Thus $T_1$ is a topology on $\productFamily{X}$ by \cref{basis_topology_is_topology,basis_for_topology}.
    Show that $T_1\subseteq T_0$.
    \begin{subproof}
        Show that for all $A$ such that $A\in S$ we have $A\in T_0$.
        \begin{subproof}
            Fix $A$.
            Assume $A\in S$.
            Take $\beta$ such that $\beta\in\dom{X}$ and there exists $U\in\apply{B}{\beta}$ such that
                $A = \preimg{\piIndex{X}{\beta}}{U}$.
            Take $U\in\apply{B}{\beta}$ such that $A = \preimg{\piIndex{X}{\beta}}{U}$.
            $\apply{B}{\beta}$ is a basis for $\apply{T}{\beta}$ on $\apply{X}{\beta}$.
            Thus $\apply{B}{\beta}$ is a basis on $\apply{X}{\beta}$ by \cref{basis_for_topology}.
            Thus $U\in\basisTopology{\apply{B}{\beta}}{\apply{X}{\beta}}$ by \cref{basis_elem_open_in_generated}.
            Thus $U\in\apply{T}{\beta}$ by \cref{basis_for_topology}.
            Show that $A\subseteq\productFamily{X}$.
            \begin{subproof}
                Show that for all $z$ such that $z\in A$ we have $z\in\productFamily{X}$.
                \begin{subproof}
                    Fix $z$.
                    Assume $z\in A$.
                    Take $y\in U$ such that $z\mathrel{\piIndex{X}{\beta}} y$ by \cref{preimg_iff}.
                    Thus $(z,y)\in\piIndex{X}{\beta}$.
                    Take $x\in\productFamily{X}$ such that $(z,y) = (x,\apply{x}{\beta})$ by \cref{projection_map_indexed}.
                    Thus $z = x$ by \cref{pair_eq_iff}.
                    Thus $z\in\productFamily{X}$.
                \end{subproof}
                Thus $A\subseteq\productFamily{X}$ by \cref{subseteq}.
            \end{subproof}
            Thus $A\in\pow{\productFamily{X}}$ by \cref{powerset_intro}.
            Thus $A\in\productSubbasisAt{T}{X}{\beta}$ by \cref{product_subbasis_indexed}.
            Thus $A\in S_0$ by \cref{product_subbasis_union_indexed}.
            $S_0$ is a subbasis on $\productFamily{X}$ by \cref{product_subbasis_is_subbasis}.
            Thus $\unions{S_0} = \productFamily{X}$ by \cref{subbasis_on}.
            Thus $A\subseteq\unions{S_0}$ by \cref{subseteq}.
            Thus $A\in\pow{\unions{S_0}}$ by \cref{powerset_intro}.
            Let $F = \{A\}$.
            $F\subseteq S_0$ by \cref{singleton_subset_intro}.
            Let $G = \{A,A\}$.
            $G$ is finite by \cref{pair_is_finite}.
            Show that $F\subseteq G$.
            \begin{subproof}
                Show that for all $x$ such that $x\in F$ we have $x\in G$.
                \begin{subproof}
                    Fix $x$.
                    Assume $x\in F$.
                    Thus $x = A$ by \cref{singleton_elim}.
                    Thus $x\in G$ by \cref{upair_intro_left}.
                \end{subproof}
                Thus $F\subseteq G$ by \cref{subseteq}.
            \end{subproof}
            Thus $F$ is finite by \cref{subseteq_finite_is_finite}.
            $F$ is inhabited by \cref{singleton_intro}.
            $\inters{F} = A$ by \cref{inters_singleton}.
            Thus $A\in\finiteIntersections{S_0}$ by \cref{finite_intersections}.
            Thus $A\in T_0$ by \cref{basis_elem_open_in_generated,basis_for_topology}.
        \end{subproof}
        Show that $\finiteIntersections{S}\subseteq T_0$.
        \begin{subproof}
            Show that for all $B$ such that $B\in\finiteIntersections{S}$ we have $B\in T_0$.
            \begin{subproof}
                Fix $B$.
                Assume $B\in\finiteIntersections{S}$.
                Take $F\subseteq S$ such that $F$ is finite and inhabited and $B = \inters{F}$ by \cref{finite_intersections}.
                Show that for all $C$ such that $C\in F$ we have $C\in T_0$.
                \begin{subproof}
                    Fix $C$.
                    Assume $C\in F$.
                    Thus $C\in S$ by \cref{subseteq}.
                    Thus $C\in T_0$ by \cref{subseteq}.
                \end{subproof}
                Thus $F\subseteq T_0$ by \cref{subseteq}.
                Thus $\inters{F}\in T_0$ by \cref{finite_intersections_open_in_topology}.
                Thus $B\in T_0$.
            \end{subproof}
            Thus $\finiteIntersections{S}\subseteq T_0$ by \cref{subseteq}.
        \end{subproof}
        Show that for all $U$ such that $U\in T_1$ we have $U\in T_0$.
        \begin{subproof}
            Fix $U$.
            Assume $U\in T_1$.
            Take $M\subseteq\finiteIntersections{S}$ such that $U = \unions{M}$ by \cref{basis_topology_unions}.
            Thus $M\subseteq T_0$ by \cref{subseteq}.
            Thus $\unions{M}\in T_0$ by \cref{topology_on}.
            Thus $U\in T_0$.
        \end{subproof}
    \end{subproof}
    Show that $T_0\subseteq T_1$.
    \begin{subproof}
        Show that for all $A$ such that $A\in S_0$ we have $A\in T_1$.
        \begin{subproof}
            Fix $A$.
            Assume $A\in S_0$.
            Take $\beta\in\dom{X}$ such that $A\in\productSubbasisAt{T}{X}{\beta}$ by \cref{product_subbasis_union_indexed}.
            Take $U\in\apply{T}{\beta}$ such that $A = \preimg{\piIndex{X}{\beta}}{U}$ by \cref{product_subbasis_indexed}.
            $\apply{B}{\beta}$ is a basis for $\apply{T}{\beta}$ on $\apply{X}{\beta}$.
            Take $M\subseteq\apply{B}{\beta}$ such that $U = \unions{M}$ by \cref{basis_topology_unions}.
            Let $N = \{\preimg{\piIndex{X}{\beta}}{V} \mid V\in M\}$.
            Show that for all $C$ such that $C\in N$ we have $C\in S$.
            \begin{subproof}
                Fix $C$.
                Assume $C\in N$.
                Take $V\in M$ such that $C = \preimg{\piIndex{X}{\beta}}{V}$.
                Thus $V\in\apply{B}{\beta}$ by \cref{subseteq}.
                Thus $C\in S$.
            \end{subproof}
            Thus $N\subseteq S$ by \cref{subseteq}.
            Thus $N\subseteq T_1$ by \cref{subseteq_transitive}.
            Show that $A = \unions{N}$.
            \begin{subproof}
                Show that for all $x$ such that $x\in A$ we have $x\in\unions{N}$.
                \begin{subproof}
                    Fix $x$.
                    Assume $x\in A$.
                    Take $y\in U$ such that $x\mathrel{\piIndex{X}{\beta}} y$ by \cref{preimg_iff}.
                    $y\in\unions{M}$.
                    Take $V\in M$ such that $y\in V$ by \cref{unions_iff}.
                    Let $C = \preimg{\piIndex{X}{\beta}}{V}$.
                    $C\in N$.
                    Thus $x\in C$ by \cref{preimg_iff}.
                    Thus $x\in\unions{N}$ by \cref{unions_iff}.
                \end{subproof}
                Thus $A\subseteq\unions{N}$ by \cref{subseteq}.
                Show that for all $x$ such that $x\in\unions{N}$ we have $x\in A$.
                \begin{subproof}
                    Fix $x$.
                    Assume $x\in\unions{N}$.
                    Take $C\in N$ such that $x\in C$ by \cref{unions_iff}.
                    Take $V\in M$ such that $C = \preimg{\piIndex{X}{\beta}}{V}$.
                    Show that $V\subseteq U$.
                    \begin{subproof}
                        Show that for all $z$ such that $z\in V$ we have $z\in U$.
                        \begin{subproof}
                            Fix $z$.
                            Assume $z\in V$.
                            Thus $z\in\unions{M}$ by \cref{unions_iff}.
                            Thus $z\in U$.
                        \end{subproof}
                        Thus $V\subseteq U$ by \cref{subseteq}.
                    \end{subproof}
                    Thus $\preimg{\piIndex{X}{\beta}}{V}\subseteq \preimg{\piIndex{X}{\beta}}{U}$ by \cref{preimg_subseteq}.
                    Thus $x\in A$ by \cref{subseteq}.
                \end{subproof}
                Thus $\unions{N}\subseteq A$ by \cref{subseteq}.
                Thus $A = \unions{N}$ by \cref{subseteq_antisymmetric}.
            \end{subproof}
            Thus $A\in T_1$ by \cref{topology_on,subseteq}.
        \end{subproof}
        Show that $\finiteIntersections{S_0}\subseteq T_1$.
        \begin{subproof}
            Show that for all $B$ such that $B\in\finiteIntersections{S_0}$ we have $B\in T_1$.
            \begin{subproof}
                Fix $B$.
                Assume $B\in\finiteIntersections{S_0}$.
                Take $F\subseteq S_0$ such that $F$ is finite and inhabited and $B = \inters{F}$ by \cref{finite_intersections}.
                Show that for all $C$ such that $C\in F$ we have $C\in T_1$.
                \begin{subproof}
                    Fix $C$.
                    Assume $C\in F$.
                    Thus $C\in S_0$ by \cref{subseteq}.
                    Thus $C\in T_1$ by \cref{subseteq}.
                \end{subproof}
                Thus $F\subseteq T_1$ by \cref{subseteq}.
                Thus $\inters{F}\in T_1$ by \cref{finite_intersections_open_in_topology}.
                Thus $B\in T_1$.
            \end{subproof}
            Thus $\finiteIntersections{S_0}\subseteq T_1$ by \cref{subseteq}.
        \end{subproof}
        Show that for all $U$ such that $U\in T_0$ we have $U\in T_1$.
        \begin{subproof}
            Fix $U$.
            Assume $U\in T_0$.
            Take $M\subseteq\finiteIntersections{S_0}$ such that $U = \unions{M}$ by \cref{basis_topology_unions}.
            Thus $M\subseteq T_1$ by \cref{subseteq}.
            Thus $\unions{M}\in T_1$ by \cref{topology_on}.
            Thus $U\in T_1$.
        \end{subproof}
    \end{subproof}
    Thus $T_0 = T_1$ by \cref{subseteq_antisymmetric}.
    Thus $\finiteIntersections{S}$ is a basis for $T_0$ on $\productFamily{X}$ by \cref{basis_for_topology}.
\end{proof}

% from §19 helper lemma 2: replace one coordinate in a function
\begin{lemma}\label{replace_coordinate_function}
    Assume $X$ is a function.
    Assume $\alpha\in\dom{X}$.
    Assume $U\subseteq\unions{\ran{X}}$.
    Then there exists $V$ such that
        $V\in\funs{\dom{X}}{\pow{\unions{\ran{X}}}}$ and
        $\apply{V}{\alpha} = U$ and
        for all $\beta\in\dom{X}$ such that $\beta\neq\alpha$ we have $\apply{V}{\beta} = \apply{X}{\beta}$.
\end{lemma}
\begin{proof}
    Let $R_1 = \{(\alpha,U)\}$.
    Let $R_2 = \{(\beta,\apply{X}{\beta}) \mid \beta\in\dom{X}\setminus\{\alpha\}\}$.
    Let $R = R_1\union R_2$.
    Show that $R$ is a relation.
    \begin{subproof}
        Show that for all $w$ such that $w\in R$ there exist $x,y$ such that $w = (x,y)$.
        \begin{subproof}
            Fix $w$.
            Assume $w\in R$.
            Thus $w\in R_1$ or $w\in R_2$ by \cref{union_iff}.
            \begin{byCase}
            \caseOf{$w\in R_1$.}
                Thus $w = (\alpha,U)$ by \cref{singleton_elim}.
            \caseOf{$w\in R_2$.}
                Take $\beta\in\dom{X}\setminus\{\alpha\}$ such that $w = (\beta,\apply{X}{\beta})$.
            \end{byCase}
        \end{subproof}
        Thus $R$ is a relation by \cref{relation}.
    \end{subproof}
    Show that for all $\beta\in\dom{X}$ there exists $y$ such that $\beta\mathrel{R} y$.
    \begin{subproof}
        Fix $\beta$.
        Assume $\beta\in\dom{X}$.
        \begin{byCase}
        \caseOf{$\beta = \alpha$.}
            Thus $(\beta,U)\in R_1$ by \cref{singleton_intro}.
            Thus $\beta\mathrel{R} U$ by \cref{union_iff}.
        \caseOf{$\beta\neq\alpha$.}
            Thus $\beta\notin\{\alpha\}$ by \cref{singleton_iff}.
            Thus $\beta\in\dom{X}\setminus\{\alpha\}$ by \cref{setminus_intro}.
            Thus $(\beta,\apply{X}{\beta})\in R_2$.
            Thus $\beta\mathrel{R}\apply{X}{\beta}$ by \cref{union_iff}.
        \end{byCase}
    \end{subproof}
    Take $V$ such that $V$ is a function on $\dom{X}$ and for all $\beta\in\dom{X}$ we have $\beta\mathrel{R}\apply{V}{\beta}$
        by \cref{choice_function}.
    Thus $V$ is a function and $\dom{V} = \dom{X}$.
    Show that for all $\beta\in\dom{X}$ we have $\apply{V}{\beta}\in\pow{\unions{\ran{X}}}$.
    \begin{subproof}
        Fix $\beta$.
        Assume $\beta\in\dom{X}$.
        Thus $\beta\mathrel{R}\apply{V}{\beta}$.
        Thus $(\beta,\apply{V}{\beta})\in R$.
        Thus $(\beta,\apply{V}{\beta})\in R_1$ or $(\beta,\apply{V}{\beta})\in R_2$ by \cref{union_iff}.
        \begin{byCase}
        \caseOf{$(\beta,\apply{V}{\beta})\in R_1$.}
            Thus $(\beta,\apply{V}{\beta}) = (\alpha,U)$ by \cref{singleton_elim}.
            Thus $\apply{V}{\beta} = U$ by \cref{pair_eq_iff}.
            Thus $\apply{V}{\beta}\in\pow{\unions{\ran{X}}}$ by \cref{powerset_intro,subseteq}.
        \caseOf{$(\beta,\apply{V}{\beta})\in R_2$.}
            Take $\gamma\in\dom{X}\setminus\{\alpha\}$ such that $(\beta,\apply{V}{\beta}) = (\gamma,\apply{X}{\gamma})$.
            Thus $\beta = \gamma$ and $\apply{V}{\beta} = \apply{X}{\beta}$ by \cref{pair_eq_iff}.
            Thus $\beta\mathrel{X}\apply{X}{\beta}$ by \cref{function_apply_elim}.
            Thus $\apply{X}{\beta}\in\ran{X}$ by \cref{ran_intro}.
            Show that $\apply{V}{\beta}\subseteq\unions{\ran{X}}$.
            \begin{subproof}
                Show that for all $z$ such that $z\in\apply{V}{\beta}$ we have $z\in\unions{\ran{X}}$.
                \begin{subproof}
                    Fix $z$.
                    Assume $z\in\apply{V}{\beta}$.
                    Thus $z\in\apply{X}{\beta}$.
                    Thus $z\in\unions{\ran{X}}$ by \cref{unions_iff}.
                \end{subproof}
                Thus $\apply{V}{\beta}\subseteq\unions{\ran{X}}$ by \cref{subseteq}.
            \end{subproof}
            Thus $\apply{V}{\beta}\in\pow{\unions{\ran{X}}}$ by \cref{powerset_intro}.
        \end{byCase}
    \end{subproof}
    Thus $V$ is a function to $\pow{\unions{\ran{X}}}$.
    Thus $V\in\funs{\dom{X}}{\pow{\unions{\ran{X}}}}$ by \cref{funs_intro}.
    Show that $\apply{V}{\alpha} = U$.
    \begin{subproof}
        $\alpha\mathrel{R}\apply{V}{\alpha}$.
        Thus $(\alpha,\apply{V}{\alpha})\in R$.
        Show that $(\alpha,\apply{V}{\alpha})\notin R_2$.
        \begin{subproof}
            Suppose not.
            Take $\gamma\in\dom{X}\setminus\{\alpha\}$ such that $(\alpha,\apply{V}{\alpha}) = (\gamma,\apply{X}{\gamma})$.
            Thus $\alpha = \gamma$ by \cref{pair_eq_iff}.
            Thus $\alpha\in\dom{X}\setminus\{\alpha\}$ by \cref{setminus_intro}.
            Thus $\alpha\notin\{\alpha\}$ by \cref{setminus_elim_right}.
            $\alpha\in\{\alpha\}$ by \cref{singleton_intro}.
            Contradiction.
        \end{subproof}
        Thus $(\alpha,\apply{V}{\alpha})\in R_1$ by \cref{union_iff}.
        Thus $(\alpha,\apply{V}{\alpha}) = (\alpha,U)$ by \cref{singleton_elim}.
        Thus $\apply{V}{\alpha} = U$ by \cref{pair_eq_iff}.
    \end{subproof}
    Show that for all $\beta\in\dom{X}$ such that $\beta\neq\alpha$ we have $\apply{V}{\beta} = \apply{X}{\beta}$.
    \begin{subproof}
        Fix $\beta$.
        Assume $\beta\in\dom{X}$ and $\beta\neq\alpha$.
        $\beta\mathrel{R}\apply{V}{\beta}$.
        Thus $(\beta,\apply{V}{\beta})\in R$.
        Show that $(\beta,\apply{V}{\beta})\notin R_1$.
        \begin{subproof}
            Suppose not.
            Thus $(\beta,\apply{V}{\beta}) = (\alpha,U)$ by \cref{singleton_elim}.
            Thus $\beta = \alpha$ by \cref{pair_eq_iff}.
            Contradiction.
        \end{subproof}
        Thus $(\beta,\apply{V}{\beta})\in R_2$ by \cref{union_iff}.
        Take $\gamma\in\dom{X}\setminus\{\alpha\}$ such that $(\beta,\apply{V}{\beta}) = (\gamma,\apply{X}{\gamma})$.
        Thus $\beta = \gamma$ and $\apply{V}{\beta} = \apply{X}{\beta}$ by \cref{pair_eq_iff}.
    \end{subproof}
\end{proof}

% from §19 Item 15: products of subspaces for the box topology
\begin{lemma}\label{product_subspace_box}
    Assume $X$ is a function.
    Assume $A$ is a function.
    Assume $T$ is a function.
    Assume $\dom{A} = \dom{X}$.
    Assume $\dom{T} = \dom{X}$.
    Assume that for all $\alpha\in\dom{X}$ we have $\apply{A}{\alpha}\subseteq\apply{X}{\alpha}$.
    Assume that for all $\alpha\in\dom{X}$ we have $\apply{T}{\alpha}$ is a topology on $\apply{X}{\alpha}$.
    Then $\productFamily{A}$ is a subspace of $\productFamily{X}$ with respect to $\boxTopology{T}{X}$.
\end{lemma}
\begin{proof}
    Let $T_0 = \boxTopology{T}{X}$.
    $\boxBasis{T}{X}$ is a basis on $\productFamily{X}$ by \cref{box_basis_is_basis}.
    Thus $T_0$ is a topology on $\productFamily{X}$ by \cref{basis_topology_is_topology,box_topology}.
    Show that $\productFamily{A}\subseteq\productFamily{X}$.
    \begin{subproof}
        Show that for all $f$ such that $f\in\productFamily{A}$ we have $f\in\productFamily{X}$.
        \begin{subproof}
            Fix $f$.
            Assume $f\in\productFamily{A}$.
            Show that for all $\alpha\in\dom{X}$ we have $\apply{f}{\alpha}\in\apply{X}{\alpha}$.
            \begin{subproof}
                Fix $\alpha$.
                Assume $\alpha\in\dom{X}$.
                $\dom{A} = \dom{X}$.
                Thus $\alpha\in\dom{A}$.
                Thus $\apply{f}{\alpha}\in\apply{A}{\alpha}$ by \cref{indexed_product}.
                $\apply{A}{\alpha}\subseteq\apply{X}{\alpha}$.
                Thus $\apply{f}{\alpha}\in\apply{X}{\alpha}$ by \cref{subseteq}.
            \end{subproof}
            Thus $f\in\funs{\dom{A}}{\unions{\ran{A}}}$ by \cref{indexed_product}.
            Thus $f$ is a function by \cref{funs_elim}.
            $\dom{f} = \dom{A}$ by \cref{funs_elim}.
            $\dom{A} = \dom{X}$.
            Show that for all $\alpha\in\dom{X}$ we have $\apply{f}{\alpha}\in\unions{\ran{X}}$.
            \begin{subproof}
                Fix $\alpha$.
                Assume $\alpha\in\dom{X}$.
                Thus $\alpha\mathrel{X}\apply{X}{\alpha}$ by \cref{function_apply_elim}.
                Thus $\apply{X}{\alpha}\in\ran{X}$ by \cref{ran_intro}.
                Thus $\apply{f}{\alpha}\in\unions{\ran{X}}$ by \cref{unions_iff}.
            \end{subproof}
            Thus $f$ is a function to $\unions{\ran{X}}$.
            Thus $f\in\funs{\dom{X}}{\unions{\ran{X}}}$ by \cref{funs_intro}.
            Thus $f\in\productFamily{X}$ by \cref{indexed_product}.
        \end{subproof}
        Thus $\productFamily{A}\subseteq\productFamily{X}$ by \cref{subseteq}.
    \end{subproof}
    Thus $\productFamily{A}$ is a subspace of $\productFamily{X}$ with respect to $T_0$ by \cref{subspace_of}.
\end{proof}

% from §19 Item 16: products of subspaces for the product topology
\begin{lemma}\label{product_subspace_product}
    Assume $X$ is a function.
    Assume $A$ is a function.
    Assume $T$ is a function.
    Assume $\dom{A} = \dom{X}$.
    Assume $\dom{T} = \dom{X}$.
    Assume $\dom{X}$ is inhabited.
    Assume that for all $\alpha\in\dom{X}$ we have $\apply{A}{\alpha}\subseteq\apply{X}{\alpha}$.
    Assume that for all $\alpha\in\dom{X}$ we have $\apply{T}{\alpha}$ is a topology on $\apply{X}{\alpha}$.
    Then $\productFamily{A}$ is a subspace of $\productFamily{X}$ with respect to $\productTopologyIndexed{T}{X}$.
\end{lemma}
\begin{proof}
    Let $T_0 = \productTopologyIndexed{T}{X}$.
    Let $S_0 = \productSubbasis{T}{X}$.
    Thus $\finiteIntersections{S_0}$ is a basis for $T_0$ on $\productFamily{X}$ by \cref{product_subbasis_finite_intersections_basis}.
    Thus $T_0$ is a topology on $\productFamily{X}$ by \cref{basis_topology_is_topology,basis_for_topology}.
    Show that $\productFamily{A}\subseteq\productFamily{X}$.
    \begin{subproof}
        Show that for all $f$ such that $f\in\productFamily{A}$ we have $f\in\productFamily{X}$.
        \begin{subproof}
            Fix $f$.
            Assume $f\in\productFamily{A}$.
            Show that for all $\alpha\in\dom{X}$ we have $\apply{f}{\alpha}\in\apply{X}{\alpha}$.
            \begin{subproof}
                Fix $\alpha$.
                Assume $\alpha\in\dom{X}$.
                $\dom{A} = \dom{X}$.
                Thus $\alpha\in\dom{A}$.
                Thus $\apply{f}{\alpha}\in\apply{A}{\alpha}$ by \cref{indexed_product}.
                $\apply{A}{\alpha}\subseteq\apply{X}{\alpha}$.
                Thus $\apply{f}{\alpha}\in\apply{X}{\alpha}$ by \cref{subseteq}.
            \end{subproof}
            Thus $f\in\funs{\dom{A}}{\unions{\ran{A}}}$ by \cref{indexed_product}.
            Thus $f$ is a function by \cref{funs_elim}.
            $\dom{f} = \dom{A}$ by \cref{funs_elim}.
            $\dom{A} = \dom{X}$.
            Show that for all $\alpha\in\dom{X}$ we have $\apply{f}{\alpha}\in\unions{\ran{X}}$.
            \begin{subproof}
                Fix $\alpha$.
                Assume $\alpha\in\dom{X}$.
                Thus $\alpha\mathrel{X}\apply{X}{\alpha}$ by \cref{function_apply_elim}.
                Thus $\apply{X}{\alpha}\in\ran{X}$ by \cref{ran_intro}.
                Thus $\apply{f}{\alpha}\in\unions{\ran{X}}$ by \cref{unions_iff}.
            \end{subproof}
            Thus $f$ is a function to $\unions{\ran{X}}$.
            Thus $f\in\funs{\dom{X}}{\unions{\ran{X}}}$ by \cref{funs_intro}.
            Thus $f\in\productFamily{X}$ by \cref{indexed_product}.
        \end{subproof}
        Thus $\productFamily{A}\subseteq\productFamily{X}$ by \cref{subseteq}.
    \end{subproof}
    Thus $\productFamily{A}$ is a subspace of $\productFamily{X}$ with respect to $T_0$ by \cref{subspace_of}.
\end{proof}

% from §19 Item 17: product of Hausdorff spaces with the box topology
\begin{lemma}\label{product_hausdorff_box}
    Assume $X$ is a function.
    Assume $T$ is a function.
    Assume $\dom{T} = \dom{X}$.
    Assume that for all $\alpha\in\dom{X}$ we have $\apply{X}{\alpha}$ is Hausdorff with respect to $\apply{T}{\alpha}$.
    Then $\productFamily{X}$ is Hausdorff with respect to $\boxTopology{T}{X}$.
\end{lemma}
\begin{proof}
    Let $T_0 = \boxTopology{T}{X}$.
    Let $B_0 = \boxBasis{T}{X}$.
    Show that for all $\alpha\in\dom{X}$ we have $\apply{T}{\alpha}$ is a topology on $\apply{X}{\alpha}$.
    \begin{subproof}
        Fix $\alpha$.
        Assume $\alpha\in\dom{X}$.
        $\apply{X}{\alpha}$ is Hausdorff with respect to $\apply{T}{\alpha}$.
        Thus $\apply{T}{\alpha}$ is a topology on $\apply{X}{\alpha}$ by \cref{hausdorff}.
    \end{subproof}
    $B_0$ is a basis on $\productFamily{X}$ by \cref{box_basis_is_basis}.
    Thus $T_0$ is a topology on $\productFamily{X}$ by \cref{basis_topology_is_topology,box_topology}.
    Show that for all $x_1,x_2$ such that $x_1\in\productFamily{X}$ and $x_2\in\productFamily{X}$ and $x_1\neq x_2$
        there exist $U_1,U_2$ such that
            $U_1$ is a neighborhood of $x_1$ in $\productFamily{X}$ with respect to $T_0$ and
            $U_2$ is a neighborhood of $x_2$ in $\productFamily{X}$ with respect to $T_0$ and
            $U_1$ is disjoint from $U_2$.
    \begin{subproof}
        Fix $x_1$.
        Fix $x_2$.
        Assume $x_1\in\productFamily{X}$ and $x_2\in\productFamily{X}$ and $x_1\neq x_2$.
        Take $w$ such that either $w\in x_1$ and $w\notin x_2$ or $w\notin x_1$ and $w\in x_2$
            by \cref{neq_witness}.
        \begin{byCase}
        \caseOf{$w\in x_1$ and $w\notin x_2$.}
            Thus $x_1\in\funs{\dom{X}}{\unions{\ran{X}}}$ by \cref{indexed_product}.
            Thus $x_1$ is a function and $\dom{x_1} = \dom{X}$ by \cref{funs_elim}.
            Thus $x_2\in\funs{\dom{X}}{\unions{\ran{X}}}$ by \cref{indexed_product}.
            Thus $x_2$ is a function and $\dom{x_2} = \dom{X}$ by \cref{funs_elim}.
            Take $\alpha\in\dom{x_1}$ such that $w = (\alpha,\apply{x_1}{\alpha})$ by \cref{function_member_elim}.
            Thus $\alpha\in\dom{X}$.
            Show that $\apply{x_1}{\alpha}\neq\apply{x_2}{\alpha}$.
            \begin{subproof}
                Suppose not.
                Thus $\apply{x_1}{\alpha} = \apply{x_2}{\alpha}$.
                Thus $\alpha\in\dom{x_2}$.
                Thus $(\alpha,\apply{x_2}{\alpha})\in x_2$ by \cref{function_apply_elim}.
                Thus $w = (\alpha,\apply{x_2}{\alpha})$ by \cref{pair_eq_iff}.
                Thus $w\in x_2$.
                Contradiction.
            \end{subproof}
            $\apply{x_1}{\alpha}\in\apply{X}{\alpha}$ by \cref{indexed_product}.
            $\apply{x_2}{\alpha}\in\apply{X}{\alpha}$ by \cref{indexed_product}.
            Take $U_1,U_2$ such that
                $U_1$ is a neighborhood of $\apply{x_1}{\alpha}$ in $\apply{X}{\alpha}$ with respect to $\apply{T}{\alpha}$ and
                $U_2$ is a neighborhood of $\apply{x_2}{\alpha}$ in $\apply{X}{\alpha}$ with respect to $\apply{T}{\alpha}$ and
                $U_1$ is disjoint from $U_2$ by \cref{hausdorff}.
            Show that $U_1\subseteq\unions{\ran{X}}$.
            \begin{subproof}
                $U_1$ is open in $\apply{X}{\alpha}$ with respect to $\apply{T}{\alpha}$ by \cref{neighborhood}.
                Thus $U_1\in\apply{T}{\alpha}$ by \cref{open_in}.
                $\apply{T}{\alpha}$ is a topology on $\apply{X}{\alpha}$ by \cref{hausdorff}.
                Thus $\apply{T}{\alpha}\subseteq\pow{\apply{X}{\alpha}}$ by \cref{topology_on}.
                Thus $U_1\in\pow{\apply{X}{\alpha}}$ by \cref{subseteq}.
                Thus $U_1\subseteq\apply{X}{\alpha}$ by \cref{pow_iff}.
                Thus $\alpha\mathrel{X}\apply{X}{\alpha}$ by \cref{function_apply_elim}.
                Thus $\apply{X}{\alpha}\in\ran{X}$ by \cref{ran_intro}.
                Show that $\apply{X}{\alpha}\subseteq\unions{\ran{X}}$.
                \begin{subproof}
                    Show that for all $z$ such that $z\in\apply{X}{\alpha}$ we have $z\in\unions{\ran{X}}$.
                    \begin{subproof}
                        Fix $z$.
                        Assume $z\in\apply{X}{\alpha}$.
                        Thus $z\in\unions{\ran{X}}$ by \cref{unions_iff}.
                    \end{subproof}
                    Thus $\apply{X}{\alpha}\subseteq\unions{\ran{X}}$ by \cref{subseteq}.
                \end{subproof}
                Thus $U_1\subseteq\unions{\ran{X}}$ by \cref{subseteq_transitive}.
            \end{subproof}
            Show that $U_2\subseteq\unions{\ran{X}}$.
            \begin{subproof}
                $U_2$ is open in $\apply{X}{\alpha}$ with respect to $\apply{T}{\alpha}$ by \cref{neighborhood}.
                Thus $U_2\in\apply{T}{\alpha}$ by \cref{open_in}.
                $\apply{T}{\alpha}$ is a topology on $\apply{X}{\alpha}$ by \cref{hausdorff}.
                Thus $\apply{T}{\alpha}\subseteq\pow{\apply{X}{\alpha}}$ by \cref{topology_on}.
                Thus $U_2\in\pow{\apply{X}{\alpha}}$ by \cref{subseteq}.
                Thus $U_2\subseteq\apply{X}{\alpha}$ by \cref{pow_iff}.
                Thus $\alpha\mathrel{X}\apply{X}{\alpha}$ by \cref{function_apply_elim}.
                Thus $\apply{X}{\alpha}\in\ran{X}$ by \cref{ran_intro}.
                Show that $\apply{X}{\alpha}\subseteq\unions{\ran{X}}$.
                \begin{subproof}
                    Show that for all $z$ such that $z\in\apply{X}{\alpha}$ we have $z\in\unions{\ran{X}}$.
                    \begin{subproof}
                        Fix $z$.
                        Assume $z\in\apply{X}{\alpha}$.
                        Thus $z\in\unions{\ran{X}}$ by \cref{unions_iff}.
                    \end{subproof}
                    Thus $\apply{X}{\alpha}\subseteq\unions{\ran{X}}$ by \cref{subseteq}.
                \end{subproof}
                Thus $U_2\subseteq\unions{\ran{X}}$ by \cref{subseteq_transitive}.
            \end{subproof}
            Take $V_1$ such that
                $V_1\in\funs{\dom{X}}{\pow{\unions{\ran{X}}}}$ and
                $\apply{V_1}{\alpha} = U_1$ and
                for all $\beta\in\dom{X}$ such that $\beta\neq\alpha$ we have $\apply{V_1}{\beta} = \apply{X}{\beta}$
                by \cref{replace_coordinate_function}.
            Take $V_2$ such that
                $V_2\in\funs{\dom{X}}{\pow{\unions{\ran{X}}}}$ and
                $\apply{V_2}{\alpha} = U_2$ and
                for all $\beta\in\dom{X}$ such that $\beta\neq\alpha$ we have $\apply{V_2}{\beta} = \apply{X}{\beta}$
                by \cref{replace_coordinate_function}.
            Show that for all $\beta\in\dom{X}$ we have $\apply{V_1}{\beta}\in\apply{T}{\beta}$.
            \begin{subproof}
                Fix $\beta$.
                Assume $\beta\in\dom{X}$.
                \begin{byCase}
                \caseOf{$\beta = \alpha$.}
                    Thus $\apply{V_1}{\beta} = U_1$.
                    $U_1$ is open in $\apply{X}{\alpha}$ with respect to $\apply{T}{\alpha}$ by \cref{neighborhood}.
                    Thus $U_1\in\apply{T}{\alpha}$ by \cref{open_in}.
                \caseOf{$\beta\neq\alpha$.}
                    Thus $\apply{V_1}{\beta} = \apply{X}{\beta}$.
                    Thus $\apply{V_1}{\beta}\in\apply{T}{\beta}$ by \cref{topology_on}.
                \end{byCase}
            \end{subproof}
            Show that for all $\beta\in\dom{X}$ we have $\apply{V_2}{\beta}\in\apply{T}{\beta}$.
            \begin{subproof}
                Fix $\beta$.
                Assume $\beta\in\dom{X}$.
                \begin{byCase}
                \caseOf{$\beta = \alpha$.}
                    Thus $\apply{V_2}{\beta} = U_2$.
                    $U_2$ is open in $\apply{X}{\alpha}$ with respect to $\apply{T}{\alpha}$ by \cref{neighborhood}.
                    Thus $U_2\in\apply{T}{\alpha}$ by \cref{open_in}.
                \caseOf{$\beta\neq\alpha$.}
                    Thus $\apply{V_2}{\beta} = \apply{X}{\beta}$.
                    Thus $\apply{V_2}{\beta}\in\apply{T}{\beta}$ by \cref{topology_on}.
                \end{byCase}
            \end{subproof}
            Let $W_1 = \productFamily{V_1}$.
            Let $W_2 = \productFamily{V_2}$.
            Show that $W_1\in B_0$.
            \begin{subproof}
                Show that $W_1\in\pow{\productFamily{X}}$.
                \begin{subproof}
                    Show that for all $x$ such that $x\in W_1$ we have $x\in\productFamily{X}$.
                    \begin{subproof}
                        Fix $x$.
                        Assume $x\in W_1$.
                        Thus $x\in\productFamily{V_1}$.
                        Thus $x\in\funs{\dom{V_1}}{\unions{\ran{V_1}}}$ by \cref{indexed_product}.
                        Thus $x$ is a function and $\dom{x} = \dom{V_1}$ by \cref{funs_elim}.
                        $V_1\in\funs{\dom{X}}{\pow{\unions{\ran{X}}}}$.
                        Thus $\dom{V_1} = \dom{X}$ by \cref{funs_elim}.
                        Thus $\dom{x} = \dom{X}$.
                        Show that for all $\beta\in\dom{X}$ we have $\apply{x}{\beta}\in\apply{X}{\beta}$.
                        \begin{subproof}
                            Fix $\beta$.
                            Assume $\beta\in\dom{X}$.
                            Thus $\beta\in\dom{V_1}$.
                            Thus $\apply{x}{\beta}\in\apply{V_1}{\beta}$ by \cref{indexed_product}.
                            Show that $\apply{V_1}{\beta}\subseteq\apply{X}{\beta}$.
                            \begin{subproof}
                                \begin{byCase}
                                \caseOf{$\beta = \alpha$.}
                                    Thus $\apply{V_1}{\beta} = U_1$.
                                    $U_1$ is open in $\apply{X}{\alpha}$ with respect to $\apply{T}{\alpha}$ by \cref{neighborhood}.
                                    Thus $U_1\in\apply{T}{\alpha}$ by \cref{open_in}.
                                    $\apply{T}{\alpha}$ is a topology on $\apply{X}{\alpha}$ by \cref{hausdorff}.
                                    Thus $\apply{T}{\alpha}\subseteq\pow{\apply{X}{\alpha}}$ by \cref{topology_on}.
                                    Thus $U_1\in\pow{\apply{X}{\alpha}}$ by \cref{subseteq}.
                                    Thus $U_1\subseteq\apply{X}{\alpha}$ by \cref{pow_iff}.
                                    Thus $\apply{V_1}{\beta}\subseteq\apply{X}{\beta}$.
                                \caseOf{$\beta\neq\alpha$.}
                                    Thus $\apply{V_1}{\beta} = \apply{X}{\beta}$.
                                    Thus $\apply{V_1}{\beta}\subseteq\apply{X}{\beta}$ by \cref{subseteq_refl}.
                                \end{byCase}
                            \end{subproof}
                            Thus $\apply{x}{\beta}\in\apply{X}{\beta}$ by \cref{elem_subseteq}.
                        \end{subproof}
                        Show that for all $\beta\in\dom{x}$ we have $\apply{x}{\beta}\in\unions{\ran{X}}$.
                        \begin{subproof}
                            Fix $\beta$.
                            Assume $\beta\in\dom{x}$.
                            Thus $\beta\in\dom{X}$.
                            Thus $\apply{x}{\beta}\in\apply{X}{\beta}$.
                            Thus $\beta\mathrel{X}\apply{X}{\beta}$ by \cref{function_apply_elim}.
                            Thus $\apply{X}{\beta}\in\ran{X}$ by \cref{ran_intro}.
                            Thus $\apply{x}{\beta}\in\unions{\ran{X}}$ by \cref{unions_iff}.
                        \end{subproof}
                        Thus $x$ is a function to $\unions{\ran{X}}$.
                        Thus $x\in\funs{\dom{X}}{\unions{\ran{X}}}$ by \cref{funs_intro}.
                        Thus $x\in\productFamily{X}$ by \cref{indexed_product}.
                    \end{subproof}
                    Thus $W_1\subseteq\productFamily{X}$ by \cref{subseteq}.
                    Thus $W_1\in\pow{\productFamily{X}}$ by \cref{pow_iff}.
                \end{subproof}
                Show that there exists $U$ such that
                    $U\in\funs{\dom{X}}{\pow{\unions{\ran{X}}}}$ and
                    for all $\beta\in\dom{X}$ we have $\apply{U}{\beta}\in\apply{T}{\beta}$ and
                    $W_1 = \productFamily{U}$.
                \begin{subproof}
                    Take $U$ such that $U = V_1$.
                \end{subproof}
                Thus $W_1\in B_0$ by \cref{box_basis}.
            \end{subproof}
            Show that $W_2\in B_0$.
            \begin{subproof}
                Show that $W_2\in\pow{\productFamily{X}}$.
                \begin{subproof}
                    Show that for all $x$ such that $x\in W_2$ we have $x\in\productFamily{X}$.
                    \begin{subproof}
                        Fix $x$.
                        Assume $x\in W_2$.
                        Thus $x\in\productFamily{V_2}$.
                        Thus $x\in\funs{\dom{V_2}}{\unions{\ran{V_2}}}$ by \cref{indexed_product}.
                        Thus $x$ is a function and $\dom{x} = \dom{V_2}$ by \cref{funs_elim}.
                        $V_2\in\funs{\dom{X}}{\pow{\unions{\ran{X}}}}$.
                        Thus $\dom{V_2} = \dom{X}$ by \cref{funs_elim}.
                        Thus $\dom{x} = \dom{X}$.
                        Show that for all $\beta\in\dom{X}$ we have $\apply{x}{\beta}\in\apply{X}{\beta}$.
                        \begin{subproof}
                            Fix $\beta$.
                            Assume $\beta\in\dom{X}$.
                            Thus $\beta\in\dom{V_2}$.
                            Thus $\apply{x}{\beta}\in\apply{V_2}{\beta}$ by \cref{indexed_product}.
                            Show that $\apply{V_2}{\beta}\subseteq\apply{X}{\beta}$.
                            \begin{subproof}
                                \begin{byCase}
                                \caseOf{$\beta = \alpha$.}
                                    Thus $\apply{V_2}{\beta} = U_2$.
                                    $U_2$ is open in $\apply{X}{\alpha}$ with respect to $\apply{T}{\alpha}$ by \cref{neighborhood}.
                                    Thus $U_2\in\apply{T}{\alpha}$ by \cref{open_in}.
                                    $\apply{T}{\alpha}$ is a topology on $\apply{X}{\alpha}$ by \cref{hausdorff}.
                                    Thus $\apply{T}{\alpha}\subseteq\pow{\apply{X}{\alpha}}$ by \cref{topology_on}.
                                    Thus $U_2\in\pow{\apply{X}{\alpha}}$ by \cref{subseteq}.
                                    Thus $U_2\subseteq\apply{X}{\alpha}$ by \cref{pow_iff}.
                                    Thus $\apply{V_2}{\beta}\subseteq\apply{X}{\beta}$.
                                \caseOf{$\beta\neq\alpha$.}
                                    Thus $\apply{V_2}{\beta} = \apply{X}{\beta}$.
                                    Thus $\apply{V_2}{\beta}\subseteq\apply{X}{\beta}$ by \cref{subseteq_refl}.
                                \end{byCase}
                            \end{subproof}
                            Thus $\apply{x}{\beta}\in\apply{X}{\beta}$ by \cref{elem_subseteq}.
                        \end{subproof}
                        Show that for all $\beta\in\dom{x}$ we have $\apply{x}{\beta}\in\unions{\ran{X}}$.
                        \begin{subproof}
                            Fix $\beta$.
                            Assume $\beta\in\dom{x}$.
                            Thus $\beta\in\dom{X}$.
                            Thus $\apply{x}{\beta}\in\apply{X}{\beta}$.
                            Thus $\beta\mathrel{X}\apply{X}{\beta}$ by \cref{function_apply_elim}.
                            Thus $\apply{X}{\beta}\in\ran{X}$ by \cref{ran_intro}.
                            Thus $\apply{x}{\beta}\in\unions{\ran{X}}$ by \cref{unions_iff}.
                        \end{subproof}
                        Thus $x$ is a function to $\unions{\ran{X}}$.
                        Thus $x\in\funs{\dom{X}}{\unions{\ran{X}}}$ by \cref{funs_intro}.
                        Thus $x\in\productFamily{X}$ by \cref{indexed_product}.
                    \end{subproof}
                    Thus $W_2\subseteq\productFamily{X}$ by \cref{subseteq}.
                    Thus $W_2\in\pow{\productFamily{X}}$ by \cref{pow_iff}.
                \end{subproof}
                Show that there exists $U$ such that
                    $U\in\funs{\dom{X}}{\pow{\unions{\ran{X}}}}$ and
                    for all $\beta\in\dom{X}$ we have $\apply{U}{\beta}\in\apply{T}{\beta}$ and
                    $W_2 = \productFamily{U}$.
                \begin{subproof}
                    Take $U$ such that $U = V_2$.
                \end{subproof}
                Thus $W_2\in B_0$ by \cref{box_basis}.
            \end{subproof}
            Thus $W_1\in T_0$ by \cref{basis_elem_open_in_generated,box_topology}.
            Thus $W_2\in T_0$ by \cref{basis_elem_open_in_generated,box_topology}.
            Thus $W_1$ is open in $\productFamily{X}$ with respect to $T_0$ by \cref{open_in}.
            Thus $W_2$ is open in $\productFamily{X}$ with respect to $T_0$ by \cref{open_in}.
            Show that $x_1\in W_1$.
            \begin{subproof}
                Show that for all $\beta\in\dom{X}$ we have $\apply{x_1}{\beta}\in\apply{V_1}{\beta}$.
                \begin{subproof}
                    Fix $\beta$.
                    Assume $\beta\in\dom{X}$.
                    \begin{byCase}
                    \caseOf{$\beta = \alpha$.}
                        Thus $\apply{x_1}{\beta}\in U_1$ by \cref{neighborhood}.
                        Thus $\apply{x_1}{\beta}\in\apply{V_1}{\beta}$.
                    \caseOf{$\beta\neq\alpha$.}
                        Thus $\apply{x_1}{\beta}\in\apply{X}{\beta}$ by \cref{indexed_product}.
                        Thus $\apply{x_1}{\beta}\in\apply{V_1}{\beta}$ by \cref{subseteq}.
                    \end{byCase}
                \end{subproof}
                Thus $x_1$ is a function and $\dom{x_1} = \dom{X}$ by \cref{funs_elim,indexed_product}.
                $\dom{V_1} = \dom{X}$ by \cref{funs_elim}.
                Show that for all $\beta\in\dom{X}$ we have $\apply{x_1}{\beta}\in\unions{\ran{V_1}}$.
                \begin{subproof}
                    Fix $\beta$.
                    Assume $\beta\in\dom{X}$.
                    $V_1$ is a function by \cref{funs_elim}.
                    Thus $\beta\mathrel{V_1}\apply{V_1}{\beta}$ by \cref{function_apply_elim}.
                    Thus $\apply{V_1}{\beta}\in\ran{V_1}$ by \cref{ran_intro}.
                    Thus $\apply{x_1}{\beta}\in\unions{\ran{V_1}}$ by \cref{unions_iff}.
                \end{subproof}
                Thus $x_1$ is a function to $\unions{\ran{V_1}}$.
                Thus $x_1\in\funs{\dom{X}}{\unions{\ran{V_1}}}$ by \cref{funs_intro}.
                Thus $x_1\in W_1$ by \cref{indexed_product}.
            \end{subproof}
            Show that $x_2\in W_2$.
            \begin{subproof}
                Show that for all $\beta\in\dom{X}$ we have $\apply{x_2}{\beta}\in\apply{V_2}{\beta}$.
                \begin{subproof}
                    Fix $\beta$.
                    Assume $\beta\in\dom{X}$.
                    \begin{byCase}
                    \caseOf{$\beta = \alpha$.}
                        Thus $\apply{x_2}{\beta}\in U_2$ by \cref{neighborhood}.
                        Thus $\apply{x_2}{\beta}\in\apply{V_2}{\beta}$.
                    \caseOf{$\beta\neq\alpha$.}
                        Thus $\apply{x_2}{\beta}\in\apply{X}{\beta}$ by \cref{indexed_product}.
                        Thus $\apply{x_2}{\beta}\in\apply{V_2}{\beta}$ by \cref{subseteq}.
                    \end{byCase}
                \end{subproof}
                Thus $x_2$ is a function and $\dom{x_2} = \dom{X}$ by \cref{funs_elim,indexed_product}.
                $\dom{V_2} = \dom{X}$ by \cref{funs_elim}.
                Show that for all $\beta\in\dom{X}$ we have $\apply{x_2}{\beta}\in\unions{\ran{V_2}}$.
                \begin{subproof}
                    Fix $\beta$.
                    Assume $\beta\in\dom{X}$.
                    $V_2$ is a function by \cref{funs_elim}.
                    Thus $\beta\mathrel{V_2}\apply{V_2}{\beta}$ by \cref{function_apply_elim}.
                    Thus $\apply{V_2}{\beta}\in\ran{V_2}$ by \cref{ran_intro}.
                    Thus $\apply{x_2}{\beta}\in\unions{\ran{V_2}}$ by \cref{unions_iff}.
                \end{subproof}
                Thus $x_2$ is a function to $\unions{\ran{V_2}}$.
                Thus $x_2\in\funs{\dom{X}}{\unions{\ran{V_2}}}$ by \cref{funs_intro}.
                Thus $x_2\in W_2$ by \cref{indexed_product}.
            \end{subproof}
            Thus $W_1$ is a neighborhood of $x_1$ in $\productFamily{X}$ with respect to $T_0$ by \cref{neighborhood}.
            Thus $W_2$ is a neighborhood of $x_2$ in $\productFamily{X}$ with respect to $T_0$ by \cref{neighborhood}.
            Show that $W_1$ is disjoint from $W_2$.
            \begin{subproof}
                Suppose not.
                Then there exists $p$ such that $p\in W_1$ and $p\in W_2$.
                Thus $\dom{V_1} = \dom{X}$ by \cref{funs_elim}.
                Thus $\alpha\in\dom{V_1}$.
                Thus $\apply{p}{\alpha}\in\apply{V_1}{\alpha}$ by \cref{indexed_product}.
                Thus $\apply{V_1}{\alpha} = U_1$.
                Thus $\apply{p}{\alpha}\in U_1$.
                Thus $\dom{V_2} = \dom{X}$ by \cref{funs_elim}.
                Thus $\alpha\in\dom{V_2}$.
                Thus $\apply{p}{\alpha}\in\apply{V_2}{\alpha}$ by \cref{indexed_product}.
                Thus $\apply{V_2}{\alpha} = U_2$.
                Thus $\apply{p}{\alpha}\in U_2$.
                Thus there exists $a$ such that $a\in U_1$ and $a\in U_2$.
                Contradiction by \cref{disjoint}.
            \end{subproof}
        \caseOf{$w\notin x_1$ and $w\in x_2$.}
            Thus $x_1\in\funs{\dom{X}}{\unions{\ran{X}}}$ by \cref{indexed_product}.
            Thus $x_1$ is a function and $\dom{x_1} = \dom{X}$ by \cref{funs_elim}.
            Thus $x_2\in\funs{\dom{X}}{\unions{\ran{X}}}$ by \cref{indexed_product}.
            Thus $x_2$ is a function and $\dom{x_2} = \dom{X}$ by \cref{funs_elim}.
            Take $\alpha\in\dom{x_2}$ such that $w = (\alpha,\apply{x_2}{\alpha})$ by \cref{function_member_elim}.
            Thus $\alpha\in\dom{X}$.
            Show that $\apply{x_1}{\alpha}\neq\apply{x_2}{\alpha}$.
            \begin{subproof}
                Suppose not.
                Thus $\apply{x_1}{\alpha} = \apply{x_2}{\alpha}$.
                Thus $\alpha\in\dom{x_1}$.
                Thus $(\alpha,\apply{x_1}{\alpha})\in x_1$ by \cref{function_apply_elim}.
                Thus $w = (\alpha,\apply{x_1}{\alpha})$ by \cref{pair_eq_iff}.
                Thus $w\in x_1$.
                Contradiction.
            \end{subproof}
            $\apply{x_1}{\alpha}\in\apply{X}{\alpha}$ by \cref{indexed_product}.
            $\apply{x_2}{\alpha}\in\apply{X}{\alpha}$ by \cref{indexed_product}.
            Take $U_1,U_2$ such that
                $U_1$ is a neighborhood of $\apply{x_1}{\alpha}$ in $\apply{X}{\alpha}$ with respect to $\apply{T}{\alpha}$ and
                $U_2$ is a neighborhood of $\apply{x_2}{\alpha}$ in $\apply{X}{\alpha}$ with respect to $\apply{T}{\alpha}$ and
                $U_1$ is disjoint from $U_2$ by \cref{hausdorff}.
            Show that $U_1\subseteq\unions{\ran{X}}$.
            \begin{subproof}
                $U_1$ is open in $\apply{X}{\alpha}$ with respect to $\apply{T}{\alpha}$ by \cref{neighborhood}.
                Thus $U_1\in\apply{T}{\alpha}$ by \cref{open_in}.
                $\apply{T}{\alpha}$ is a topology on $\apply{X}{\alpha}$ by \cref{hausdorff}.
                Thus $\apply{T}{\alpha}\subseteq\pow{\apply{X}{\alpha}}$ by \cref{topology_on}.
                Thus $U_1\in\pow{\apply{X}{\alpha}}$ by \cref{subseteq}.
                Thus $U_1\subseteq\apply{X}{\alpha}$ by \cref{pow_iff}.
                Thus $\alpha\mathrel{X}\apply{X}{\alpha}$ by \cref{function_apply_elim}.
                Thus $\apply{X}{\alpha}\in\ran{X}$ by \cref{ran_intro}.
                Show that $\apply{X}{\alpha}\subseteq\unions{\ran{X}}$.
                \begin{subproof}
                    Show that for all $z$ such that $z\in\apply{X}{\alpha}$ we have $z\in\unions{\ran{X}}$.
                    \begin{subproof}
                        Fix $z$.
                        Assume $z\in\apply{X}{\alpha}$.
                        Thus $z\in\unions{\ran{X}}$ by \cref{unions_iff}.
                    \end{subproof}
                    Thus $\apply{X}{\alpha}\subseteq\unions{\ran{X}}$ by \cref{subseteq}.
                \end{subproof}
                Thus $U_1\subseteq\unions{\ran{X}}$ by \cref{subseteq_transitive}.
            \end{subproof}
            Show that $U_2\subseteq\unions{\ran{X}}$.
            \begin{subproof}
                $U_2$ is open in $\apply{X}{\alpha}$ with respect to $\apply{T}{\alpha}$ by \cref{neighborhood}.
                Thus $U_2\in\apply{T}{\alpha}$ by \cref{open_in}.
                $\apply{T}{\alpha}$ is a topology on $\apply{X}{\alpha}$ by \cref{hausdorff}.
                Thus $\apply{T}{\alpha}\subseteq\pow{\apply{X}{\alpha}}$ by \cref{topology_on}.
                Thus $U_2\in\pow{\apply{X}{\alpha}}$ by \cref{subseteq}.
                Thus $U_2\subseteq\apply{X}{\alpha}$ by \cref{pow_iff}.
                Thus $\alpha\mathrel{X}\apply{X}{\alpha}$ by \cref{function_apply_elim}.
                Thus $\apply{X}{\alpha}\in\ran{X}$ by \cref{ran_intro}.
                Show that $\apply{X}{\alpha}\subseteq\unions{\ran{X}}$.
                \begin{subproof}
                    Show that for all $z$ such that $z\in\apply{X}{\alpha}$ we have $z\in\unions{\ran{X}}$.
                    \begin{subproof}
                        Fix $z$.
                        Assume $z\in\apply{X}{\alpha}$.
                        Thus $z\in\unions{\ran{X}}$ by \cref{unions_iff}.
                    \end{subproof}
                    Thus $\apply{X}{\alpha}\subseteq\unions{\ran{X}}$ by \cref{subseteq}.
                \end{subproof}
                Thus $U_2\subseteq\unions{\ran{X}}$ by \cref{subseteq_transitive}.
            \end{subproof}
            Take $V_1$ such that
                $V_1\in\funs{\dom{X}}{\pow{\unions{\ran{X}}}}$ and
                $\apply{V_1}{\alpha} = U_1$ and
                for all $\beta\in\dom{X}$ such that $\beta\neq\alpha$ we have $\apply{V_1}{\beta} = \apply{X}{\beta}$
                by \cref{replace_coordinate_function}.
            Take $V_2$ such that
                $V_2\in\funs{\dom{X}}{\pow{\unions{\ran{X}}}}$ and
                $\apply{V_2}{\alpha} = U_2$ and
                for all $\beta\in\dom{X}$ such that $\beta\neq\alpha$ we have $\apply{V_2}{\beta} = \apply{X}{\beta}$
                by \cref{replace_coordinate_function}.
            Show that for all $\beta\in\dom{X}$ we have $\apply{V_1}{\beta}\in\apply{T}{\beta}$.
            \begin{subproof}
                Fix $\beta$.
                Assume $\beta\in\dom{X}$.
                \begin{byCase}
                \caseOf{$\beta = \alpha$.}
                    Thus $\apply{V_1}{\beta} = U_1$.
                    $U_1$ is open in $\apply{X}{\alpha}$ with respect to $\apply{T}{\alpha}$ by \cref{neighborhood}.
                    Thus $U_1\in\apply{T}{\alpha}$ by \cref{open_in}.
                \caseOf{$\beta\neq\alpha$.}
                    Thus $\apply{V_1}{\beta} = \apply{X}{\beta}$.
                    Thus $\apply{V_1}{\beta}\in\apply{T}{\beta}$ by \cref{topology_on}.
                \end{byCase}
            \end{subproof}
            Show that for all $\beta\in\dom{X}$ we have $\apply{V_2}{\beta}\in\apply{T}{\beta}$.
            \begin{subproof}
                Fix $\beta$.
                Assume $\beta\in\dom{X}$.
                \begin{byCase}
                \caseOf{$\beta = \alpha$.}
                    Thus $\apply{V_2}{\beta} = U_2$.
                    $U_2$ is open in $\apply{X}{\alpha}$ with respect to $\apply{T}{\alpha}$ by \cref{neighborhood}.
                    Thus $U_2\in\apply{T}{\alpha}$ by \cref{open_in}.
                \caseOf{$\beta\neq\alpha$.}
                    Thus $\apply{V_2}{\beta} = \apply{X}{\beta}$.
                    Thus $\apply{V_2}{\beta}\in\apply{T}{\beta}$ by \cref{topology_on}.
                \end{byCase}
            \end{subproof}
            Let $W_1 = \productFamily{V_1}$.
            Let $W_2 = \productFamily{V_2}$.
            Show that $W_1\in B_0$.
            \begin{subproof}
                Show that $W_1\in\pow{\productFamily{X}}$.
                \begin{subproof}
                    Show that for all $x$ such that $x\in W_1$ we have $x\in\productFamily{X}$.
                    \begin{subproof}
                        Fix $x$.
                        Assume $x\in W_1$.
                        Thus $x\in\productFamily{V_1}$.
                        Thus $x\in\funs{\dom{V_1}}{\unions{\ran{V_1}}}$ by \cref{indexed_product}.
                        Thus $x$ is a function and $\dom{x} = \dom{V_1}$ by \cref{funs_elim}.
                        $V_1\in\funs{\dom{X}}{\pow{\unions{\ran{X}}}}$.
                        Thus $\dom{V_1} = \dom{X}$ by \cref{funs_elim}.
                        Thus $\dom{x} = \dom{X}$.
                        Show that for all $\beta\in\dom{X}$ we have $\apply{x}{\beta}\in\apply{X}{\beta}$.
                        \begin{subproof}
                            Fix $\beta$.
                            Assume $\beta\in\dom{X}$.
                            Thus $\beta\in\dom{V_1}$.
                            Thus $\apply{x}{\beta}\in\apply{V_1}{\beta}$ by \cref{indexed_product}.
                            Show that $\apply{V_1}{\beta}\subseteq\apply{X}{\beta}$.
                            \begin{subproof}
                                \begin{byCase}
                                \caseOf{$\beta = \alpha$.}
                                    Thus $\apply{V_1}{\beta} = U_1$.
                                    $U_1$ is open in $\apply{X}{\alpha}$ with respect to $\apply{T}{\alpha}$ by \cref{neighborhood}.
                                    Thus $U_1\in\apply{T}{\alpha}$ by \cref{open_in}.
                                    $\apply{T}{\alpha}$ is a topology on $\apply{X}{\alpha}$ by \cref{hausdorff}.
                                    Thus $\apply{T}{\alpha}\subseteq\pow{\apply{X}{\alpha}}$ by \cref{topology_on}.
                                    Thus $U_1\in\pow{\apply{X}{\alpha}}$ by \cref{subseteq}.
                                    Thus $U_1\subseteq\apply{X}{\alpha}$ by \cref{pow_iff}.
                                    Thus $\apply{V_1}{\beta}\subseteq\apply{X}{\beta}$.
                                \caseOf{$\beta\neq\alpha$.}
                                    Thus $\apply{V_1}{\beta} = \apply{X}{\beta}$.
                                    Thus $\apply{V_1}{\beta}\subseteq\apply{X}{\beta}$ by \cref{subseteq_refl}.
                                \end{byCase}
                            \end{subproof}
                            Thus $\apply{x}{\beta}\in\apply{X}{\beta}$ by \cref{elem_subseteq}.
                        \end{subproof}
                        Show that for all $\beta\in\dom{x}$ we have $\apply{x}{\beta}\in\unions{\ran{X}}$.
                        \begin{subproof}
                            Fix $\beta$.
                            Assume $\beta\in\dom{x}$.
                            Thus $\beta\in\dom{X}$.
                            Thus $\apply{x}{\beta}\in\apply{X}{\beta}$.
                            Thus $\beta\mathrel{X}\apply{X}{\beta}$ by \cref{function_apply_elim}.
                            Thus $\apply{X}{\beta}\in\ran{X}$ by \cref{ran_intro}.
                            Thus $\apply{x}{\beta}\in\unions{\ran{X}}$ by \cref{unions_iff}.
                        \end{subproof}
                        Thus $x$ is a function to $\unions{\ran{X}}$.
                        Thus $x\in\funs{\dom{X}}{\unions{\ran{X}}}$ by \cref{funs_intro}.
                        Thus $x\in\productFamily{X}$ by \cref{indexed_product}.
                    \end{subproof}
                    Thus $W_1\subseteq\productFamily{X}$ by \cref{subseteq}.
                    Thus $W_1\in\pow{\productFamily{X}}$ by \cref{pow_iff}.
                \end{subproof}
                Show that there exists $U$ such that
                    $U\in\funs{\dom{X}}{\pow{\unions{\ran{X}}}}$ and
                    for all $\beta\in\dom{X}$ we have $\apply{U}{\beta}\in\apply{T}{\beta}$ and
                    $W_1 = \productFamily{U}$.
                \begin{subproof}
                    Take $U$ such that $U = V_1$.
                \end{subproof}
                Thus $W_1\in B_0$ by \cref{box_basis}.
            \end{subproof}
            Show that $W_2\in B_0$.
            \begin{subproof}
                Show that $W_2\in\pow{\productFamily{X}}$.
                \begin{subproof}
                    Show that for all $x$ such that $x\in W_2$ we have $x\in\productFamily{X}$.
                    \begin{subproof}
                        Fix $x$.
                        Assume $x\in W_2$.
                        Thus $x\in\productFamily{V_2}$.
                        Thus $x\in\funs{\dom{V_2}}{\unions{\ran{V_2}}}$ by \cref{indexed_product}.
                        Thus $x$ is a function and $\dom{x} = \dom{V_2}$ by \cref{funs_elim}.
                        $V_2\in\funs{\dom{X}}{\pow{\unions{\ran{X}}}}$.
                        Thus $\dom{V_2} = \dom{X}$ by \cref{funs_elim}.
                        Thus $\dom{x} = \dom{X}$.
                        Show that for all $\beta\in\dom{X}$ we have $\apply{x}{\beta}\in\apply{X}{\beta}$.
                        \begin{subproof}
                            Fix $\beta$.
                            Assume $\beta\in\dom{X}$.
                            Thus $\beta\in\dom{V_2}$.
                            Thus $\apply{x}{\beta}\in\apply{V_2}{\beta}$ by \cref{indexed_product}.
                            Show that $\apply{V_2}{\beta}\subseteq\apply{X}{\beta}$.
                            \begin{subproof}
                                \begin{byCase}
                                \caseOf{$\beta = \alpha$.}
                                    Thus $\apply{V_2}{\beta} = U_2$.
                                    $U_2$ is open in $\apply{X}{\alpha}$ with respect to $\apply{T}{\alpha}$ by \cref{neighborhood}.
                                    Thus $U_2\in\apply{T}{\alpha}$ by \cref{open_in}.
                                    $\apply{T}{\alpha}$ is a topology on $\apply{X}{\alpha}$ by \cref{hausdorff}.
                                    Thus $\apply{T}{\alpha}\subseteq\pow{\apply{X}{\alpha}}$ by \cref{topology_on}.
                                    Thus $U_2\in\pow{\apply{X}{\alpha}}$ by \cref{subseteq}.
                                    Thus $U_2\subseteq\apply{X}{\alpha}$ by \cref{pow_iff}.
                                    Thus $\apply{V_2}{\beta}\subseteq\apply{X}{\beta}$.
                                \caseOf{$\beta\neq\alpha$.}
                                    Thus $\apply{V_2}{\beta} = \apply{X}{\beta}$.
                                    Thus $\apply{V_2}{\beta}\subseteq\apply{X}{\beta}$ by \cref{subseteq_refl}.
                                \end{byCase}
                            \end{subproof}
                            Thus $\apply{x}{\beta}\in\apply{X}{\beta}$ by \cref{elem_subseteq}.
                        \end{subproof}
                        Show that for all $\beta\in\dom{x}$ we have $\apply{x}{\beta}\in\unions{\ran{X}}$.
                        \begin{subproof}
                            Fix $\beta$.
                            Assume $\beta\in\dom{x}$.
                            Thus $\beta\in\dom{X}$.
                            Thus $\apply{x}{\beta}\in\apply{X}{\beta}$.
                            Thus $\beta\mathrel{X}\apply{X}{\beta}$ by \cref{function_apply_elim}.
                            Thus $\apply{X}{\beta}\in\ran{X}$ by \cref{ran_intro}.
                            Thus $\apply{x}{\beta}\in\unions{\ran{X}}$ by \cref{unions_iff}.
                        \end{subproof}
                        Thus $x$ is a function to $\unions{\ran{X}}$.
                        Thus $x\in\funs{\dom{X}}{\unions{\ran{X}}}$ by \cref{funs_intro}.
                        Thus $x\in\productFamily{X}$ by \cref{indexed_product}.
                    \end{subproof}
                    Thus $W_2\subseteq\productFamily{X}$ by \cref{subseteq}.
                    Thus $W_2\in\pow{\productFamily{X}}$ by \cref{pow_iff}.
                \end{subproof}
                Show that there exists $U$ such that
                    $U\in\funs{\dom{X}}{\pow{\unions{\ran{X}}}}$ and
                    for all $\beta\in\dom{X}$ we have $\apply{U}{\beta}\in\apply{T}{\beta}$ and
                    $W_2 = \productFamily{U}$.
                \begin{subproof}
                    Take $U$ such that $U = V_2$.
                \end{subproof}
                Thus $W_2\in B_0$ by \cref{box_basis}.
            \end{subproof}
            Thus $W_1\in T_0$ by \cref{basis_elem_open_in_generated,box_topology}.
            Thus $W_2\in T_0$ by \cref{basis_elem_open_in_generated,box_topology}.
            Thus $W_1$ is open in $\productFamily{X}$ with respect to $T_0$ by \cref{open_in}.
            Thus $W_2$ is open in $\productFamily{X}$ with respect to $T_0$ by \cref{open_in}.
            Show that $x_1\in W_1$.
            \begin{subproof}
                Show that for all $\beta\in\dom{X}$ we have $\apply{x_1}{\beta}\in\apply{V_1}{\beta}$.
                \begin{subproof}
                    Fix $\beta$.
                    Assume $\beta\in\dom{X}$.
                    \begin{byCase}
                    \caseOf{$\beta = \alpha$.}
                        Thus $\apply{x_1}{\beta}\in U_1$ by \cref{neighborhood}.
                        Thus $\apply{x_1}{\beta}\in\apply{V_1}{\beta}$.
                    \caseOf{$\beta\neq\alpha$.}
                        Thus $\apply{x_1}{\beta}\in\apply{X}{\beta}$ by \cref{indexed_product}.
                        Thus $\apply{x_1}{\beta}\in\apply{V_1}{\beta}$ by \cref{subseteq}.
                    \end{byCase}
                \end{subproof}
                Thus $x_1$ is a function and $\dom{x_1} = \dom{X}$ by \cref{funs_elim,indexed_product}.
                $\dom{V_1} = \dom{X}$ by \cref{funs_elim}.
                Show that for all $\beta\in\dom{X}$ we have $\apply{x_1}{\beta}\in\unions{\ran{V_1}}$.
                \begin{subproof}
                    Fix $\beta$.
                    Assume $\beta\in\dom{X}$.
                    $V_1$ is a function by \cref{funs_elim}.
                    Thus $\beta\mathrel{V_1}\apply{V_1}{\beta}$ by \cref{function_apply_elim}.
                    Thus $\apply{V_1}{\beta}\in\ran{V_1}$ by \cref{ran_intro}.
                    Thus $\apply{x_1}{\beta}\in\unions{\ran{V_1}}$ by \cref{unions_iff}.
                \end{subproof}
                Thus $x_1$ is a function to $\unions{\ran{V_1}}$.
                Thus $x_1\in\funs{\dom{X}}{\unions{\ran{V_1}}}$ by \cref{funs_intro}.
                Thus $x_1\in W_1$ by \cref{indexed_product}.
            \end{subproof}
            Show that $x_2\in W_2$.
            \begin{subproof}
                Show that for all $\beta\in\dom{X}$ we have $\apply{x_2}{\beta}\in\apply{V_2}{\beta}$.
                \begin{subproof}
                    Fix $\beta$.
                    Assume $\beta\in\dom{X}$.
                    \begin{byCase}
                    \caseOf{$\beta = \alpha$.}
                        Thus $\apply{x_2}{\beta}\in U_2$ by \cref{neighborhood}.
                        Thus $\apply{x_2}{\beta}\in\apply{V_2}{\beta}$.
                    \caseOf{$\beta\neq\alpha$.}
                        Thus $\apply{x_2}{\beta}\in\apply{X}{\beta}$ by \cref{indexed_product}.
                        Thus $\apply{x_2}{\beta}\in\apply{V_2}{\beta}$ by \cref{subseteq}.
                    \end{byCase}
                \end{subproof}
                Thus $x_2$ is a function and $\dom{x_2} = \dom{X}$ by \cref{funs_elim,indexed_product}.
                $\dom{V_2} = \dom{X}$ by \cref{funs_elim}.
                Show that for all $\beta\in\dom{X}$ we have $\apply{x_2}{\beta}\in\unions{\ran{V_2}}$.
                \begin{subproof}
                    Fix $\beta$.
                    Assume $\beta\in\dom{X}$.
                    $V_2$ is a function by \cref{funs_elim}.
                    Thus $\beta\mathrel{V_2}\apply{V_2}{\beta}$ by \cref{function_apply_elim}.
                    Thus $\apply{V_2}{\beta}\in\ran{V_2}$ by \cref{ran_intro}.
                    Thus $\apply{x_2}{\beta}\in\unions{\ran{V_2}}$ by \cref{unions_iff}.
                \end{subproof}
                Thus $x_2$ is a function to $\unions{\ran{V_2}}$.
                Thus $x_2\in\funs{\dom{X}}{\unions{\ran{V_2}}}$ by \cref{funs_intro}.
                Thus $x_2\in W_2$ by \cref{indexed_product}.
            \end{subproof}
            Thus $W_1$ is a neighborhood of $x_1$ in $\productFamily{X}$ with respect to $T_0$ by \cref{neighborhood}.
            Thus $W_2$ is a neighborhood of $x_2$ in $\productFamily{X}$ with respect to $T_0$ by \cref{neighborhood}.
            Show that $W_1$ is disjoint from $W_2$.
            \begin{subproof}
                Suppose not.
                Then there exists $p$ such that $p\in W_1$ and $p\in W_2$.
                Thus $\dom{V_1} = \dom{X}$ by \cref{funs_elim}.
                Thus $\alpha\in\dom{V_1}$.
                Thus $\apply{p}{\alpha}\in\apply{V_1}{\alpha}$ by \cref{indexed_product}.
                Thus $\apply{V_1}{\alpha} = U_1$.
                Thus $\apply{p}{\alpha}\in U_1$.
                Thus $\dom{V_2} = \dom{X}$ by \cref{funs_elim}.
                Thus $\alpha\in\dom{V_2}$.
                Thus $\apply{p}{\alpha}\in\apply{V_2}{\alpha}$ by \cref{indexed_product}.
                Thus $\apply{V_2}{\alpha} = U_2$.
                Thus $\apply{p}{\alpha}\in U_2$.
                Thus there exists $a$ such that $a\in U_1$ and $a\in U_2$.
                Contradiction by \cref{disjoint}.
            \end{subproof}
        \end{byCase}
    \end{subproof}
    Thus $\productFamily{X}$ is Hausdorff with respect to $T_0$ by \cref{hausdorff}.
\end{proof}

% from §19 Item 18: product of Hausdorff spaces with the product topology
\begin{lemma}\label{product_hausdorff_product}
    Assume $X$ is a function.
    Assume $T$ is a function.
    Assume $\dom{T} = \dom{X}$.
    Assume $\dom{X}$ is inhabited.
    Assume that for all $\alpha\in\dom{X}$ we have $\apply{X}{\alpha}$ is Hausdorff with respect to $\apply{T}{\alpha}$.
    Then $\productFamily{X}$ is Hausdorff with respect to $\productTopologyIndexed{T}{X}$.
\end{lemma}
\begin{proof}
    Let $T_0 = \productTopologyIndexed{T}{X}$.
    Let $S = \productSubbasis{T}{X}$.
    Show that for all $\alpha\in\dom{X}$ we have $\apply{T}{\alpha}$ is a topology on $\apply{X}{\alpha}$.
    \begin{subproof}
        Fix $\alpha$.
        Assume $\alpha\in\dom{X}$.
        $\apply{X}{\alpha}$ is Hausdorff with respect to $\apply{T}{\alpha}$.
        Thus $\apply{T}{\alpha}$ is a topology on $\apply{X}{\alpha}$ by \cref{hausdorff}.
    \end{subproof}
    $S$ is a subbasis on $\productFamily{X}$ by \cref{product_subbasis_is_subbasis}.
    Thus $\finiteIntersections{S}$ is a basis for $T_0$ on $\productFamily{X}$ by \cref{product_subbasis_finite_intersections_basis}.
    Thus $T_0$ is a topology on $\productFamily{X}$ by \cref{basis_topology_is_topology,basis_for_topology}.
    Show that for all $x_1,x_2$ such that $x_1\in\productFamily{X}$ and $x_2\in\productFamily{X}$ and $x_1\neq x_2$
        there exist $U_1,U_2$ such that
            $U_1$ is a neighborhood of $x_1$ in $\productFamily{X}$ with respect to $T_0$ and
            $U_2$ is a neighborhood of $x_2$ in $\productFamily{X}$ with respect to $T_0$ and
            $U_1$ is disjoint from $U_2$.
    \begin{subproof}
        Fix $x_1$.
        Fix $x_2$.
        Assume $x_1\in\productFamily{X}$ and $x_2\in\productFamily{X}$ and $x_1\neq x_2$.
        Take $w$ such that either $w\in x_1$ and $w\notin x_2$ or $w\notin x_1$ and $w\in x_2$
            by \cref{neq_witness}.
        \begin{byCase}
        \caseOf{$w\in x_1$ and $w\notin x_2$.}
            Thus $x_1\in\funs{\dom{X}}{\unions{\ran{X}}}$ by \cref{indexed_product}.
            Thus $x_1$ is a function and $\dom{x_1} = \dom{X}$ by \cref{funs_elim}.
            Thus $x_2\in\funs{\dom{X}}{\unions{\ran{X}}}$ by \cref{indexed_product}.
            Thus $x_2$ is a function and $\dom{x_2} = \dom{X}$ by \cref{funs_elim}.
            Take $\alpha\in\dom{x_1}$ such that $w = (\alpha,\apply{x_1}{\alpha})$ by \cref{function_member_elim}.
            Thus $\alpha\in\dom{X}$.
            Show that $\apply{x_1}{\alpha}\neq\apply{x_2}{\alpha}$.
            \begin{subproof}
                Suppose not.
                Thus $\apply{x_1}{\alpha} = \apply{x_2}{\alpha}$.
                Thus $\alpha\in\dom{x_2}$.
                Thus $(\alpha,\apply{x_2}{\alpha})\in x_2$ by \cref{function_apply_elim}.
                Thus $w = (\alpha,\apply{x_2}{\alpha})$ by \cref{pair_eq_iff}.
                Thus $w\in x_2$.
                Contradiction.
            \end{subproof}
            $\apply{x_1}{\alpha}\in\apply{X}{\alpha}$ by \cref{indexed_product}.
            $\apply{x_2}{\alpha}\in\apply{X}{\alpha}$ by \cref{indexed_product}.
            Take $U_1,U_2$ such that
                $U_1$ is a neighborhood of $\apply{x_1}{\alpha}$ in $\apply{X}{\alpha}$ with respect to $\apply{T}{\alpha}$ and
                $U_2$ is a neighborhood of $\apply{x_2}{\alpha}$ in $\apply{X}{\alpha}$ with respect to $\apply{T}{\alpha}$ and
                $U_1$ is disjoint from $U_2$ by \cref{hausdorff}.
            Let $W_1 = \preimg{\piIndex{X}{\alpha}}{U_1}$.
            Let $W_2 = \preimg{\piIndex{X}{\alpha}}{U_2}$.
            $U_1$ is open in $\apply{X}{\alpha}$ with respect to $\apply{T}{\alpha}$ by \cref{neighborhood}.
            $U_2$ is open in $\apply{X}{\alpha}$ with respect to $\apply{T}{\alpha}$ by \cref{neighborhood}.
            Show that $\dom{\piIndex{X}{\alpha}} = \productFamily{X}$.
            \begin{subproof}
                Show that for all $x$ such that $x\in\dom{\piIndex{X}{\alpha}}$ we have $x\in\productFamily{X}$.
                \begin{subproof}
                    Fix $x$.
                    Assume $x\in\dom{\piIndex{X}{\alpha}}$.
                    Take $y$ such that $(x,y)\in\piIndex{X}{\alpha}$ by \cref{dom}.
                    Take $x_1\in\productFamily{X}$ such that $(x,y) = (x_1,\apply{x_1}{\alpha})$ by \cref{projection_map_indexed}.
                    Thus $x = x_1$ by \cref{pair_eq_iff}.
                    Thus $x\in\productFamily{X}$.
                \end{subproof}
                Thus $\dom{\piIndex{X}{\alpha}}\subseteq\productFamily{X}$ by \cref{subseteq}.
                Show that for all $x$ such that $x\in\productFamily{X}$ we have $x\in\dom{\piIndex{X}{\alpha}}$.
                \begin{subproof}
                    Fix $x$.
                    Assume $x\in\productFamily{X}$.
                    Thus $(x,\apply{x}{\alpha})\in\piIndex{X}{\alpha}$ by \cref{projection_map_indexed}.
                    Thus $x\in\dom{\piIndex{X}{\alpha}}$ by \cref{dom}.
                \end{subproof}
                Thus $\productFamily{X}\subseteq\dom{\piIndex{X}{\alpha}}$ by \cref{subseteq}.
                Thus $\dom{\piIndex{X}{\alpha}} = \productFamily{X}$ by \cref{subseteq_antisymmetric}.
            \end{subproof}
            Show that $W_1\in\productSubbasisAt{T}{X}{\alpha}$.
            \begin{subproof}
                Show that $W_1\in\pow{\productFamily{X}}$.
                \begin{subproof}
                    $W_1\subseteq\dom{\piIndex{X}{\alpha}}$ by \cref{preimg_subseteq_dom}.
                    Thus $W_1\subseteq\productFamily{X}$ by \cref{subseteq_transitive}.
                    Thus $W_1\in\pow{\productFamily{X}}$ by \cref{pow_iff}.
                \end{subproof}
                Show that there exists $U$ such that $U\in\apply{T}{\alpha}$ and $W_1 = \preimg{\piIndex{X}{\alpha}}{U}$.
                \begin{subproof}
                    $U_1\in\apply{T}{\alpha}$ by \cref{open_in}.
                    Take $U$ such that $U = U_1$.
                \end{subproof}
                Thus $W_1\in\productSubbasisAt{T}{X}{\alpha}$ by \cref{product_subbasis_indexed}.
            \end{subproof}
            Show that $W_2\in\productSubbasisAt{T}{X}{\alpha}$.
            \begin{subproof}
                Show that $W_2\in\pow{\productFamily{X}}$.
                \begin{subproof}
                    $W_2\subseteq\dom{\piIndex{X}{\alpha}}$ by \cref{preimg_subseteq_dom}.
                    Thus $W_2\subseteq\productFamily{X}$ by \cref{subseteq_transitive}.
                    Thus $W_2\in\pow{\productFamily{X}}$ by \cref{pow_iff}.
                \end{subproof}
                Show that there exists $U$ such that $U\in\apply{T}{\alpha}$ and $W_2 = \preimg{\piIndex{X}{\alpha}}{U}$.
                \begin{subproof}
                    $U_2\in\apply{T}{\alpha}$ by \cref{open_in}.
                    Take $U$ such that $U = U_2$.
                \end{subproof}
                Thus $W_2\in\productSubbasisAt{T}{X}{\alpha}$ by \cref{product_subbasis_indexed}.
            \end{subproof}
            Show that $W_1\in S$.
            \begin{subproof}
                Thus $W_1\in\pow{\productFamily{X}}$ by \cref{product_subbasis_indexed}.
                Show that there exists $\beta\in\dom{X}$ such that $W_1\in\productSubbasisAt{T}{X}{\beta}$.
                \begin{subproof}
                    Take $\beta$ such that $\beta = \alpha$.
                \end{subproof}
                Thus $W_1\in S$ by \cref{product_subbasis_union_indexed}.
            \end{subproof}
            Show that $W_2\in S$.
            \begin{subproof}
                Thus $W_2\in\pow{\productFamily{X}}$ by \cref{product_subbasis_indexed}.
                Show that there exists $\beta\in\dom{X}$ such that $W_2\in\productSubbasisAt{T}{X}{\beta}$.
                \begin{subproof}
                    Take $\beta$ such that $\beta = \alpha$.
                \end{subproof}
                Thus $W_2\in S$ by \cref{product_subbasis_union_indexed}.
            \end{subproof}
            Let $F_1 = \{W_1\}$.
            Let $F_2 = \{W_2\}$.
            Show that $F_1\subseteq S$.
            \begin{subproof}
                Show that for all $z$ such that $z\in F_1$ we have $z\in S$.
                \begin{subproof}
                    Fix $z$.
                    Assume $z\in F_1$.
                    Thus $z = W_1$ by \cref{singleton_elim}.
                    Thus $z\in S$.
                \end{subproof}
                Thus $F_1\subseteq S$ by \cref{subseteq}.
            \end{subproof}
            Show that $F_2\subseteq S$.
            \begin{subproof}
                Show that for all $z$ such that $z\in F_2$ we have $z\in S$.
                \begin{subproof}
                    Fix $z$.
                    Assume $z\in F_2$.
                    Thus $z = W_2$ by \cref{singleton_elim}.
                    Thus $z\in S$.
                \end{subproof}
                Thus $F_2\subseteq S$ by \cref{subseteq}.
            \end{subproof}
            Let $G_1 = \{W_1,W_1\}$.
            Let $G_2 = \{W_2,W_2\}$.
            $G_1$ is finite by \cref{pair_is_finite}.
            $G_2$ is finite by \cref{pair_is_finite}.
            Show that $F_1\subseteq G_1$.
            \begin{subproof}
                Show that for all $z$ such that $z\in F_1$ we have $z\in G_1$.
                \begin{subproof}
                    Fix $z$.
                    Assume $z\in F_1$.
                    Thus $z = W_1$ by \cref{singleton_elim}.
                    Thus $z\in G_1$ by \cref{upair_intro_left}.
                \end{subproof}
                Thus $F_1\subseteq G_1$ by \cref{subseteq}.
            \end{subproof}
            Show that $F_2\subseteq G_2$.
            \begin{subproof}
                Show that for all $z$ such that $z\in F_2$ we have $z\in G_2$.
                \begin{subproof}
                    Fix $z$.
                    Assume $z\in F_2$.
                    Thus $z = W_2$ by \cref{singleton_elim}.
                    Thus $z\in G_2$ by \cref{upair_intro_left}.
                \end{subproof}
                Thus $F_2\subseteq G_2$ by \cref{subseteq}.
            \end{subproof}
            Thus $F_1$ is finite by \cref{subseteq_finite_is_finite}.
            Thus $F_2$ is finite by \cref{subseteq_finite_is_finite}.
            $F_1$ is inhabited by \cref{singleton_intro}.
            $F_2$ is inhabited by \cref{singleton_intro}.
            $\inters{F_1} = W_1$ by \cref{inters_singleton}.
            $\inters{F_2} = W_2$ by \cref{inters_singleton}.
            Show that $W_1\in\pow{\unions{S}}$.
            \begin{subproof}
                Show that for all $z$ such that $z\in W_1$ we have $z\in\unions{S}$.
                \begin{subproof}
                    Fix $z$.
                    Assume $z\in W_1$.
                    Take $A$ such that $A = W_1$.
                    Thus $A\in S$.
                    Thus $z\in A$.
                    Thus $z\in\unions{S}$ by \cref{unions_iff}.
                \end{subproof}
                Thus $W_1\subseteq\unions{S}$ by \cref{subseteq}.
                Thus $W_1\in\pow{\unions{S}}$ by \cref{pow_iff}.
            \end{subproof}
            Show that $W_2\in\pow{\unions{S}}$.
            \begin{subproof}
                Show that for all $z$ such that $z\in W_2$ we have $z\in\unions{S}$.
                \begin{subproof}
                    Fix $z$.
                    Assume $z\in W_2$.
                    Take $A$ such that $A = W_2$.
                    Thus $A\in S$.
                    Thus $z\in A$.
                    Thus $z\in\unions{S}$ by \cref{unions_iff}.
                \end{subproof}
                Thus $W_2\subseteq\unions{S}$ by \cref{subseteq}.
                Thus $W_2\in\pow{\unions{S}}$ by \cref{pow_iff}.
            \end{subproof}
            Thus $W_1\in\finiteIntersections{S}$ by \cref{finite_intersections}.
            Thus $W_2\in\finiteIntersections{S}$ by \cref{finite_intersections}.
            Thus $W_1\in T_0$ and $W_2\in T_0$ by \cref{basis_elem_open_in_generated,basis_for_topology}.
            Thus $W_1$ is open in $\productFamily{X}$ with respect to $T_0$ and
                $W_2$ is open in $\productFamily{X}$ with respect to $T_0$ by \cref{open_in}.
            $\apply{x_1}{\alpha}\in U_1$ by \cref{neighborhood}.
            $\apply{x_2}{\alpha}\in U_2$ by \cref{neighborhood}.
            Thus $(x_1,\apply{x_1}{\alpha})\in\piIndex{X}{\alpha}$ by \cref{projection_map_indexed}.
            Thus $(x_2,\apply{x_2}{\alpha})\in\piIndex{X}{\alpha}$ by \cref{projection_map_indexed}.
            Thus $x_1\in W_1$ and $x_2\in W_2$ by \cref{preimg_iff}.
            Thus $W_1$ is a neighborhood of $x_1$ in $\productFamily{X}$ with respect to $T_0$ and
                $W_2$ is a neighborhood of $x_2$ in $\productFamily{X}$ with respect to $T_0$ by \cref{neighborhood}.
            Show that $W_1$ is disjoint from $W_2$.
            \begin{subproof}
                Suppose not.
                Then there exists $p$ such that $p\in W_1$ and $p\in W_2$.
                Take $y_1\in U_1$ such that $p\mathrel{\piIndex{X}{\alpha}} y_1$ by \cref{preimg_iff}.
                Take $y_2\in U_2$ such that $p\mathrel{\piIndex{X}{\alpha}} y_2$ by \cref{preimg_iff}.
                Thus $(p,y_1)\in\piIndex{X}{\alpha}$ and $(p,y_2)\in\piIndex{X}{\alpha}$.
                Take $x\in\productFamily{X}$ such that $(p,y_1) = (x,\apply{x}{\alpha})$ by \cref{projection_map_indexed}.
                Take $x'\in\productFamily{X}$ such that $(p,y_2) = (x',\apply{x'}{\alpha})$ by \cref{projection_map_indexed}.
                Thus $p = x$ and $p = x'$ by \cref{pair_eq_iff}.
                Thus $y_1 = \apply{p}{\alpha}$ and $y_2 = \apply{p}{\alpha}$ by \cref{pair_eq_iff}.
                Thus $y_1 = y_2$.
                Thus there exists $a$ such that $a\in U_1$ and $a\in U_2$.
                Contradiction by \cref{disjoint}.
            \end{subproof}
        \caseOf{$w\notin x_1$ and $w\in x_2$.}
            Thus $x_1\in\funs{\dom{X}}{\unions{\ran{X}}}$ by \cref{indexed_product}.
            Thus $x_1$ is a function and $\dom{x_1} = \dom{X}$ by \cref{funs_elim}.
            Thus $x_2\in\funs{\dom{X}}{\unions{\ran{X}}}$ by \cref{indexed_product}.
            Thus $x_2$ is a function and $\dom{x_2} = \dom{X}$ by \cref{funs_elim}.
            Take $\alpha\in\dom{x_2}$ such that $w = (\alpha,\apply{x_2}{\alpha})$ by \cref{function_member_elim}.
            Thus $\alpha\in\dom{X}$.
            Show that $\apply{x_1}{\alpha}\neq\apply{x_2}{\alpha}$.
            \begin{subproof}
                Suppose not.
                Thus $\apply{x_1}{\alpha} = \apply{x_2}{\alpha}$.
                Thus $\alpha\in\dom{x_1}$.
                Thus $(\alpha,\apply{x_1}{\alpha})\in x_1$ by \cref{function_apply_elim}.
                Thus $w = (\alpha,\apply{x_1}{\alpha})$ by \cref{pair_eq_iff}.
                Thus $w\in x_1$.
                Contradiction.
            \end{subproof}
            $\apply{x_1}{\alpha}\in\apply{X}{\alpha}$ by \cref{indexed_product}.
            $\apply{x_2}{\alpha}\in\apply{X}{\alpha}$ by \cref{indexed_product}.
            Take $U_1,U_2$ such that
                $U_1$ is a neighborhood of $\apply{x_1}{\alpha}$ in $\apply{X}{\alpha}$ with respect to $\apply{T}{\alpha}$ and
                $U_2$ is a neighborhood of $\apply{x_2}{\alpha}$ in $\apply{X}{\alpha}$ with respect to $\apply{T}{\alpha}$ and
                $U_1$ is disjoint from $U_2$ by \cref{hausdorff}.
            Let $W_1 = \preimg{\piIndex{X}{\alpha}}{U_1}$.
            Let $W_2 = \preimg{\piIndex{X}{\alpha}}{U_2}$.
            $U_1$ is open in $\apply{X}{\alpha}$ with respect to $\apply{T}{\alpha}$ by \cref{neighborhood}.
            $U_2$ is open in $\apply{X}{\alpha}$ with respect to $\apply{T}{\alpha}$ by \cref{neighborhood}.
            Show that $\dom{\piIndex{X}{\alpha}} = \productFamily{X}$.
            \begin{subproof}
                Show that for all $x$ such that $x\in\dom{\piIndex{X}{\alpha}}$ we have $x\in\productFamily{X}$.
                \begin{subproof}
                    Fix $x$.
                    Assume $x\in\dom{\piIndex{X}{\alpha}}$.
                    Take $y$ such that $(x,y)\in\piIndex{X}{\alpha}$ by \cref{dom}.
                    Take $x_1\in\productFamily{X}$ such that $(x,y) = (x_1,\apply{x_1}{\alpha})$ by \cref{projection_map_indexed}.
                    Thus $x = x_1$ by \cref{pair_eq_iff}.
                    Thus $x\in\productFamily{X}$.
                \end{subproof}
                Thus $\dom{\piIndex{X}{\alpha}}\subseteq\productFamily{X}$ by \cref{subseteq}.
                Show that for all $x$ such that $x\in\productFamily{X}$ we have $x\in\dom{\piIndex{X}{\alpha}}$.
                \begin{subproof}
                    Fix $x$.
                    Assume $x\in\productFamily{X}$.
                    Thus $(x,\apply{x}{\alpha})\in\piIndex{X}{\alpha}$ by \cref{projection_map_indexed}.
                    Thus $x\in\dom{\piIndex{X}{\alpha}}$ by \cref{dom}.
                \end{subproof}
                Thus $\productFamily{X}\subseteq\dom{\piIndex{X}{\alpha}}$ by \cref{subseteq}.
                Thus $\dom{\piIndex{X}{\alpha}} = \productFamily{X}$ by \cref{subseteq_antisymmetric}.
            \end{subproof}
            Show that $W_1\in\productSubbasisAt{T}{X}{\alpha}$.
            \begin{subproof}
                Show that $W_1\in\pow{\productFamily{X}}$.
                \begin{subproof}
                    $W_1\subseteq\dom{\piIndex{X}{\alpha}}$ by \cref{preimg_subseteq_dom}.
                    Thus $W_1\subseteq\productFamily{X}$ by \cref{subseteq_transitive}.
                    Thus $W_1\in\pow{\productFamily{X}}$ by \cref{pow_iff}.
                \end{subproof}
                Show that there exists $U$ such that $U\in\apply{T}{\alpha}$ and $W_1 = \preimg{\piIndex{X}{\alpha}}{U}$.
                \begin{subproof}
                    $U_1\in\apply{T}{\alpha}$ by \cref{open_in}.
                    Take $U$ such that $U = U_1$.
                \end{subproof}
                Thus $W_1\in\productSubbasisAt{T}{X}{\alpha}$ by \cref{product_subbasis_indexed}.
            \end{subproof}
            Show that $W_2\in\productSubbasisAt{T}{X}{\alpha}$.
            \begin{subproof}
                Show that $W_2\in\pow{\productFamily{X}}$.
                \begin{subproof}
                    $W_2\subseteq\dom{\piIndex{X}{\alpha}}$ by \cref{preimg_subseteq_dom}.
                    Thus $W_2\subseteq\productFamily{X}$ by \cref{subseteq_transitive}.
                    Thus $W_2\in\pow{\productFamily{X}}$ by \cref{pow_iff}.
                \end{subproof}
                Show that there exists $U$ such that $U\in\apply{T}{\alpha}$ and $W_2 = \preimg{\piIndex{X}{\alpha}}{U}$.
                \begin{subproof}
                    $U_2\in\apply{T}{\alpha}$ by \cref{open_in}.
                    Take $U$ such that $U = U_2$.
                \end{subproof}
                Thus $W_2\in\productSubbasisAt{T}{X}{\alpha}$ by \cref{product_subbasis_indexed}.
            \end{subproof}
            Show that $W_1\in S$.
            \begin{subproof}
                Thus $W_1\in\pow{\productFamily{X}}$ by \cref{product_subbasis_indexed}.
                Show that there exists $\beta\in\dom{X}$ such that $W_1\in\productSubbasisAt{T}{X}{\beta}$.
                \begin{subproof}
                    Take $\beta$ such that $\beta = \alpha$.
                \end{subproof}
                Thus $W_1\in S$ by \cref{product_subbasis_union_indexed}.
            \end{subproof}
            Show that $W_2\in S$.
            \begin{subproof}
                Thus $W_2\in\pow{\productFamily{X}}$ by \cref{product_subbasis_indexed}.
                Show that there exists $\beta\in\dom{X}$ such that $W_2\in\productSubbasisAt{T}{X}{\beta}$.
                \begin{subproof}
                    Take $\beta$ such that $\beta = \alpha$.
                \end{subproof}
                Thus $W_2\in S$ by \cref{product_subbasis_union_indexed}.
            \end{subproof}
            Let $F_1 = \{W_1\}$.
            Let $F_2 = \{W_2\}$.
            Show that $F_1\subseteq S$.
            \begin{subproof}
                Show that for all $z$ such that $z\in F_1$ we have $z\in S$.
                \begin{subproof}
                    Fix $z$.
                    Assume $z\in F_1$.
                    Thus $z = W_1$ by \cref{singleton_elim}.
                    Thus $z\in S$.
                \end{subproof}
                Thus $F_1\subseteq S$ by \cref{subseteq}.
            \end{subproof}
            Show that $F_2\subseteq S$.
            \begin{subproof}
                Show that for all $z$ such that $z\in F_2$ we have $z\in S$.
                \begin{subproof}
                    Fix $z$.
                    Assume $z\in F_2$.
                    Thus $z = W_2$ by \cref{singleton_elim}.
                    Thus $z\in S$.
                \end{subproof}
                Thus $F_2\subseteq S$ by \cref{subseteq}.
            \end{subproof}
            Let $G_1 = \{W_1,W_1\}$.
            Let $G_2 = \{W_2,W_2\}$.
            $G_1$ is finite by \cref{pair_is_finite}.
            $G_2$ is finite by \cref{pair_is_finite}.
            Show that $F_1\subseteq G_1$.
            \begin{subproof}
                Show that for all $z$ such that $z\in F_1$ we have $z\in G_1$.
                \begin{subproof}
                    Fix $z$.
                    Assume $z\in F_1$.
                    Thus $z = W_1$ by \cref{singleton_elim}.
                    Thus $z\in G_1$ by \cref{upair_intro_left}.
                \end{subproof}
                Thus $F_1\subseteq G_1$ by \cref{subseteq}.
            \end{subproof}
            Show that $F_2\subseteq G_2$.
            \begin{subproof}
                Show that for all $z$ such that $z\in F_2$ we have $z\in G_2$.
                \begin{subproof}
                    Fix $z$.
                    Assume $z\in F_2$.
                    Thus $z = W_2$ by \cref{singleton_elim}.
                    Thus $z\in G_2$ by \cref{upair_intro_left}.
                \end{subproof}
                Thus $F_2\subseteq G_2$ by \cref{subseteq}.
            \end{subproof}
            Thus $F_1$ is finite by \cref{subseteq_finite_is_finite}.
            Thus $F_2$ is finite by \cref{subseteq_finite_is_finite}.
            $F_1$ is inhabited by \cref{singleton_intro}.
            $F_2$ is inhabited by \cref{singleton_intro}.
            $\inters{F_1} = W_1$ by \cref{inters_singleton}.
            $\inters{F_2} = W_2$ by \cref{inters_singleton}.
            Show that $W_1\in\pow{\unions{S}}$.
            \begin{subproof}
                Show that for all $z$ such that $z\in W_1$ we have $z\in\unions{S}$.
                \begin{subproof}
                    Fix $z$.
                    Assume $z\in W_1$.
                    Take $A$ such that $A = W_1$.
                    Thus $A\in S$.
                    Thus $z\in A$.
                    Thus $z\in\unions{S}$ by \cref{unions_iff}.
                \end{subproof}
                Thus $W_1\subseteq\unions{S}$ by \cref{subseteq}.
                Thus $W_1\in\pow{\unions{S}}$ by \cref{pow_iff}.
            \end{subproof}
            Show that $W_2\in\pow{\unions{S}}$.
            \begin{subproof}
                Show that for all $z$ such that $z\in W_2$ we have $z\in\unions{S}$.
                \begin{subproof}
                    Fix $z$.
                    Assume $z\in W_2$.
                    Take $A$ such that $A = W_2$.
                    Thus $A\in S$.
                    Thus $z\in A$.
                    Thus $z\in\unions{S}$ by \cref{unions_iff}.
                \end{subproof}
                Thus $W_2\subseteq\unions{S}$ by \cref{subseteq}.
                Thus $W_2\in\pow{\unions{S}}$ by \cref{pow_iff}.
            \end{subproof}
            Thus $W_1\in\finiteIntersections{S}$ by \cref{finite_intersections}.
            Thus $W_2\in\finiteIntersections{S}$ by \cref{finite_intersections}.
            Thus $W_1\in T_0$ and $W_2\in T_0$ by \cref{basis_elem_open_in_generated,basis_for_topology}.
            Thus $W_1$ is open in $\productFamily{X}$ with respect to $T_0$ and
                $W_2$ is open in $\productFamily{X}$ with respect to $T_0$ by \cref{open_in}.
            $\apply{x_1}{\alpha}\in U_1$ by \cref{neighborhood}.
            $\apply{x_2}{\alpha}\in U_2$ by \cref{neighborhood}.
            Thus $(x_1,\apply{x_1}{\alpha})\in\piIndex{X}{\alpha}$ by \cref{projection_map_indexed}.
            Thus $(x_2,\apply{x_2}{\alpha})\in\piIndex{X}{\alpha}$ by \cref{projection_map_indexed}.
            Thus $x_1\in W_1$ and $x_2\in W_2$ by \cref{preimg_iff}.
            Thus $W_1$ is a neighborhood of $x_1$ in $\productFamily{X}$ with respect to $T_0$ and
                $W_2$ is a neighborhood of $x_2$ in $\productFamily{X}$ with respect to $T_0$ by \cref{neighborhood}.
            Show that $W_1$ is disjoint from $W_2$.
            \begin{subproof}
                Suppose not.
                Then there exists $p$ such that $p\in W_1$ and $p\in W_2$.
                Take $y_1\in U_1$ such that $p\mathrel{\piIndex{X}{\alpha}} y_1$ by \cref{preimg_iff}.
                Take $y_2\in U_2$ such that $p\mathrel{\piIndex{X}{\alpha}} y_2$ by \cref{preimg_iff}.
                Thus $(p,y_1)\in\piIndex{X}{\alpha}$ and $(p,y_2)\in\piIndex{X}{\alpha}$.
                Take $x\in\productFamily{X}$ such that $(p,y_1) = (x,\apply{x}{\alpha})$ by \cref{projection_map_indexed}.
                Take $x'\in\productFamily{X}$ such that $(p,y_2) = (x',\apply{x'}{\alpha})$ by \cref{projection_map_indexed}.
                Thus $p = x$ and $p = x'$ by \cref{pair_eq_iff}.
                Thus $y_1 = \apply{p}{\alpha}$ and $y_2 = \apply{p}{\alpha}$ by \cref{pair_eq_iff}.
                Thus $y_1 = y_2$.
                Thus there exists $a$ such that $a\in U_1$ and $a\in U_2$.
                Contradiction by \cref{disjoint}.
            \end{subproof}
        \end{byCase}
    \end{subproof}
    Thus $\productFamily{X}$ is Hausdorff with respect to $T_0$ by \cref{hausdorff}.
\end{proof}

% from §19 helper lemma 5: closure family
\begin{definition}\label{closure_family}
    $\closureFamily{A}{X}{T} =
        \{(\alpha,\closure{\apply{A}{\alpha}}{\apply{X}{\alpha}}{\apply{T}{\alpha}}) \mid \alpha\in\dom{X}\}$.
\end{definition}

% from §19 Item 19: closure of products with the box topology
\begin{lemma}\label{product_closure_box}
    Assume $X$ is a function.
    Assume $A$ is a function.
    Assume $T$ is a function.
    Assume $\dom{A} = \dom{X}$.
    Assume $\dom{T} = \dom{X}$.
    Assume that for all $\alpha\in\dom{X}$ we have $\apply{A}{\alpha}\subseteq\apply{X}{\alpha}$.
    Assume that for all $\alpha\in\dom{X}$ we have $\apply{T}{\alpha}$ is a topology on $\apply{X}{\alpha}$.
    Let $C = \closureFamily{A}{X}{T}$.
    Then $\productFamily{C} = \closure{\productFamily{A}}{\productFamily{X}}{\boxTopology{T}{X}}$.
\end{lemma}
\begin{proof}
    Let $T_0 = \boxTopology{T}{X}$.
    Let $B_0 = \boxBasis{T}{X}$.
    Show that for all $\alpha\in\dom{X}$ we have $\apply{T}{\alpha}$ is a topology on $\apply{X}{\alpha}$.
    \begin{subproof}
        Fix $\alpha$.
        Assume $\alpha\in\dom{X}$.
        $\apply{T}{\alpha}$ is a topology on $\apply{X}{\alpha}$ by \cref{subseteq}.
    \end{subproof}
    $B_0$ is a basis on $\productFamily{X}$ by \cref{box_basis_is_basis}.
    Thus $T_0$ is a topology on $\productFamily{X}$ by \cref{basis_topology_is_topology,box_topology}.
    $\productFamily{A}$ is a subspace of $\productFamily{X}$ with respect to $T_0$ by \cref{product_subspace_box}.
    Thus $\productFamily{A}\subseteq\productFamily{X}$ by \cref{subspace_of}.
    Show that $C$ is a function.
    \begin{subproof}
        Show that $C$ is a relation.
        \begin{subproof}
            Show that for all $w$ such that $w\in C$ there exist $a,b$ such that $w = (a,b)$.
            \begin{subproof}
                Fix $w$.
                Assume $w\in C$.
                Take $\alpha\in\dom{X}$ such that
                    $w = (\alpha,\closure{\apply{A}{\alpha}}{\apply{X}{\alpha}}{\apply{T}{\alpha}})$
                    by \cref{closure_family}.
            \end{subproof}
            Thus $C$ is a relation by \cref{relation}.
        \end{subproof}
        Show that $C$ is right-unique.
        \begin{subproof}
            Show that for all $a,b_1,b_2$ such that $(a,b_1)\in C$ and $(a,b_2)\in C$ we have $b_1 = b_2$.
            \begin{subproof}
                Fix $a$.
                Fix $b_1$.
                Fix $b_2$.
                Assume $(a,b_1)\in C$ and $(a,b_2)\in C$.
                Take $\alpha_1\in\dom{X}$ such that $(a,b_1) =
                    (\alpha_1,\closure{\apply{A}{\alpha_1}}{\apply{X}{\alpha_1}}{\apply{T}{\alpha_1}})$
                    by \cref{closure_family}.
                Take $\alpha_2\in\dom{X}$ such that $(a,b_2) =
                    (\alpha_2,\closure{\apply{A}{\alpha_2}}{\apply{X}{\alpha_2}}{\apply{T}{\alpha_2}})$
                    by \cref{closure_family}.
                Thus $a = \alpha_1$ and $b_1 = \closure{\apply{A}{\alpha_1}}{\apply{X}{\alpha_1}}{\apply{T}{\alpha_1}}$ by \cref{pair_eq_iff}.
                Thus $a = \alpha_2$ and $b_2 = \closure{\apply{A}{\alpha_2}}{\apply{X}{\alpha_2}}{\apply{T}{\alpha_2}}$ by \cref{pair_eq_iff}.
                Thus $\alpha_1 = \alpha_2$.
                Thus $b_1 = b_2$.
            \end{subproof}
            Thus $C$ is right-unique.
        \end{subproof}
        Thus $C$ is a function.
    \end{subproof}
    Show that $\productFamily{C}\subseteq \closure{\productFamily{A}}{\productFamily{X}}{T_0}$.
    \begin{subproof}
        Show that for all $x$ such that $x\in\productFamily{C}$ we have
            $x\in\closure{\productFamily{A}}{\productFamily{X}}{T_0}$.
        \begin{subproof}
            Fix $x$.
            Assume $x\in\productFamily{C}$.
            Show that $\dom{C} = \dom{X}$.
            \begin{subproof}
                Show that for all $\alpha\in\dom{X}$ we have $\alpha\in\dom{C}$.
                \begin{subproof}
                    Fix $\alpha$.
                    Assume $\alpha\in\dom{X}$.
                    Thus $(\alpha,\closure{\apply{A}{\alpha}}{\apply{X}{\alpha}}{\apply{T}{\alpha}})\in C$ by \cref{closure_family}.
                    Thus $\alpha\in\dom{C}$ by \cref{dom}.
                \end{subproof}
                Thus $\dom{X}\subseteq\dom{C}$ by \cref{subseteq}.
                Show that for all $\alpha\in\dom{C}$ we have $\alpha\in\dom{X}$.
                \begin{subproof}
                    Fix $\alpha$.
                    Assume $\alpha\in\dom{C}$.
                    Take $y$ such that $(\alpha,y)\in C$ by \cref{dom}.
                    Take $\beta\in\dom{X}$ such that $(\alpha,y) =
                        (\beta,\closure{\apply{A}{\beta}}{\apply{X}{\beta}}{\apply{T}{\beta}})$
                        by \cref{closure_family}.
                    Thus $\alpha = \beta$ by \cref{pair_eq_iff}.
                    Thus $\alpha\in\dom{X}$.
                \end{subproof}
                Thus $\dom{C}\subseteq\dom{X}$ by \cref{subseteq}.
                Thus $\dom{C} = \dom{X}$ by \cref{subseteq_antisymmetric}.
            \end{subproof}
            Thus $x\in\funs{\dom{C}}{\unions{\ran{C}}}$ by \cref{indexed_product}.
            Thus $x$ is a function and $\dom{x} = \dom{C}$ by \cref{funs_elim}.
            Show that for all $\alpha\in\dom{X}$ we have $\apply{x}{\alpha}\in\apply{X}{\alpha}$.
            \begin{subproof}
                Fix $\alpha$.
                Assume $\alpha\in\dom{X}$.
                Thus $\alpha\in\dom{C}$.
                Thus $\apply{x}{\alpha}\in\apply{C}{\alpha}$ by \cref{indexed_product}.
                Thus $(\alpha,\closure{\apply{A}{\alpha}}{\apply{X}{\alpha}}{\apply{T}{\alpha}})\in C$ by \cref{closure_family}.
                Thus $\apply{C}{\alpha} = \closure{\apply{A}{\alpha}}{\apply{X}{\alpha}}{\apply{T}{\alpha}}$ by \cref{function_apply_intro}.
                Thus $\apply{x}{\alpha}\in\closure{\apply{A}{\alpha}}{\apply{X}{\alpha}}{\apply{T}{\alpha}}$.
                Thus $\apply{x}{\alpha}\in\apply{X}{\alpha}$ by \cref{closure}.
            \end{subproof}
            Show that for all $\alpha\in\dom{X}$ we have $\apply{x}{\alpha}\in\unions{\ran{X}}$.
            \begin{subproof}
                Fix $\alpha$.
                Assume $\alpha\in\dom{X}$.
                Thus $\alpha\mathrel{X}\apply{X}{\alpha}$ by \cref{function_apply_elim}.
                Thus $\apply{X}{\alpha}\in\ran{X}$ by \cref{ran_intro}.
                Thus $\apply{x}{\alpha}\in\unions{\ran{X}}$ by \cref{unions_iff}.
            \end{subproof}
            Thus $x$ is a function to $\unions{\ran{X}}$.
            Thus $x\in\funs{\dom{X}}{\unions{\ran{X}}}$ by \cref{funs_intro}.
            Thus $x\in\productFamily{X}$ by \cref{indexed_product}.
            Show that for all $B_1$ such that $B_1\in B_0$ and $x\in B_1$ we have $B_1$ intersects $\productFamily{A}$.
            \begin{subproof}
                Fix $B_1$.
                Assume $B_1\in B_0$ and $x\in B_1$.
                Take $U$ such that
                    $U\in\funs{\dom{X}}{\pow{\unions{\ran{X}}}}$ and
                    $B_1 = \productFamily{U}$ and
                    for all $\alpha\in\dom{X}$ we have $\apply{U}{\alpha}\in\apply{T}{\alpha}$
                    by \cref{box_basis}.
                $\dom{U} = \dom{X}$ by \cref{funs_elim}.
                Show that for all $\alpha\in\dom{X}$ we have $\apply{x}{\alpha}\in\apply{U}{\alpha}$.
                \begin{subproof}
                    Fix $\alpha$.
                    Assume $\alpha\in\dom{X}$.
                    Thus $\apply{x}{\alpha}\in\apply{U}{\alpha}$ by \cref{indexed_product}.
                \end{subproof}
                Show that for all $\alpha\in\dom{X}$ we have $\apply{U}{\alpha}$ intersects $\apply{A}{\alpha}$.
                \begin{subproof}
                    Fix $\alpha$.
                    Assume $\alpha\in\dom{X}$.
                    $\apply{U}{\alpha}\in\apply{T}{\alpha}$.
                    Thus $\apply{U}{\alpha}$ is open in $\apply{X}{\alpha}$ with respect to $\apply{T}{\alpha}$ by \cref{open_in}.
                    Thus $(\alpha,\closure{\apply{A}{\alpha}}{\apply{X}{\alpha}}{\apply{T}{\alpha}})\in C$ by \cref{closure_family}.
                    Thus $\alpha\in\dom{C}$ by \cref{dom}.
                    Thus $\apply{x}{\alpha}\in\apply{C}{\alpha}$ by \cref{indexed_product}.
                    Thus $\apply{C}{\alpha} = \closure{\apply{A}{\alpha}}{\apply{X}{\alpha}}{\apply{T}{\alpha}}$ by \cref{function_apply_intro}.
                    Thus $\apply{x}{\alpha}\in\closure{\apply{A}{\alpha}}{\apply{X}{\alpha}}{\apply{T}{\alpha}}$.
                    $\apply{A}{\alpha}\subseteq\apply{X}{\alpha}$.
                    Thus $\apply{U}{\alpha}$ intersects $\apply{A}{\alpha}$ by \cref{closure_iff_intersects_open}.
                \end{subproof}
                Let $R = \{w\in\dom{X}\times\unions{\ran{X}} \mid
                    \exists \alpha\in\dom{X}.\exists a\in\apply{A}{\alpha}\inter\apply{U}{\alpha}. w = (\alpha,a)\}$.
                Show that $R$ is a relation.
                \begin{subproof}
                    Show that for all $w$ such that $w\in R$ there exist $y,z$ such that $w = (y,z)$.
                    \begin{subproof}
                        Fix $w$.
                        Assume $w\in R$.
                        Take $\alpha$ such that $\alpha\in\dom{X}$ and there exists $a\in\apply{A}{\alpha}\inter\apply{U}{\alpha}$
                            such that $w = (\alpha,a)$.
                        Take $a\in\apply{A}{\alpha}\inter\apply{U}{\alpha}$ such that $w = (\alpha,a)$.
                    \end{subproof}
                    Thus $R$ is a relation by \cref{relation}.
                \end{subproof}
                Show that for all $\alpha\in\dom{X}$ there exists $a$ such that $\alpha\mathrel{R} a$.
                \begin{subproof}
                    Fix $\alpha$.
                    Assume $\alpha\in\dom{X}$.
                    $\apply{U}{\alpha}$ intersects $\apply{A}{\alpha}$.
                    Thus $\apply{U}{\alpha}\inter\apply{A}{\alpha}\neq\emptyset$ by \cref{intersects}.
                    Take $a$ such that $a\in\apply{U}{\alpha}\inter\apply{A}{\alpha}$ by \cref{nonempty_inhabited}.
                    Thus $a\in\apply{A}{\alpha}$ and $a\in\apply{U}{\alpha}$ by \cref{inter}.
                    $\apply{A}{\alpha}\subseteq\apply{X}{\alpha}$.
                    Thus $a\in\apply{X}{\alpha}$ by \cref{elem_subseteq}.
                    Thus $\alpha\mathrel{X}\apply{X}{\alpha}$ by \cref{function_apply_elim}.
                    Thus $\apply{X}{\alpha}\in\ran{X}$ by \cref{ran_intro}.
                    Thus $a\in\unions{\ran{X}}$ by \cref{unions_iff}.
                    Thus $(\alpha,a)\in\dom{X}\times\unions{\ran{X}}$ by \cref{times_tuple_intro}.
                    Thus $(\alpha,a)\in R$.
                    Thus $\alpha\mathrel{R} a$.
                \end{subproof}
                Take $f$ such that $f$ is a function on $\dom{X}$ and for all $\alpha\in\dom{X}$ we have
                    $\alpha\mathrel{R}\apply{f}{\alpha}$ by \cref{choice_function}.
                Thus $f$ is a function and $\dom{f} = \dom{X}$.
                Show that for all $\alpha\in\dom{X}$ we have $\apply{f}{\alpha}\in\apply{A}{\alpha}$ and
                    $\apply{f}{\alpha}\in\apply{U}{\alpha}$.
                \begin{subproof}
                    Fix $\alpha$.
                    Assume $\alpha\in\dom{X}$.
                    Thus $(\alpha,\apply{f}{\alpha})\in R$.
                    Take $\beta$ such that $\beta\in\dom{X}$ and there exists $a\in\apply{A}{\beta}\inter\apply{U}{\beta}$ such that
                        $(\alpha,\apply{f}{\alpha}) = (\beta,a)$.
                    Take $a\in\apply{A}{\beta}\inter\apply{U}{\beta}$ such that
                        $(\alpha,\apply{f}{\alpha}) = (\beta,a)$.
                    Thus $\alpha = \beta$ and $\apply{f}{\alpha} = a$ by \cref{pair_eq_iff}.
                    Thus $a\in\apply{A}{\beta}$ and $a\in\apply{U}{\beta}$ by \cref{inter}.
                    Thus $\apply{f}{\alpha}\in\apply{A}{\alpha}$ and $\apply{f}{\alpha}\in\apply{U}{\alpha}$.
                \end{subproof}
                Show that $f\in\productFamily{A}$.
                \begin{subproof}
                    Show that for all $\alpha\in\dom{X}$ we have $\apply{f}{\alpha}\in\unions{\ran{A}}$.
                    \begin{subproof}
                        Fix $\alpha$.
                        Assume $\alpha\in\dom{X}$.
                        $\dom{A} = \dom{X}$.
                        Thus $\alpha\in\dom{A}$.
                        Thus $\alpha\mathrel{A}\apply{A}{\alpha}$ by \cref{function_apply_elim}.
                        Thus $\apply{A}{\alpha}\in\ran{A}$ by \cref{ran_intro}.
                        Thus $\apply{f}{\alpha}\in\unions{\ran{A}}$ by \cref{unions_iff}.
                    \end{subproof}
                    Thus $f$ is a function to $\unions{\ran{A}}$.
                    Thus $f\in\funs{\dom{X}}{\unions{\ran{A}}}$ by \cref{funs_intro}.
                    Thus $f\in\productFamily{A}$ by \cref{indexed_product}.
                \end{subproof}
                Show that $f\in B_1$.
                \begin{subproof}
                    Show that for all $\alpha\in\dom{X}$ we have $\apply{f}{\alpha}\in\unions{\ran{U}}$.
                    \begin{subproof}
                        Fix $\alpha$.
                        Assume $\alpha\in\dom{X}$.
                        $U$ is a function by \cref{funs_elim}.
                        Thus $\alpha\mathrel{U}\apply{U}{\alpha}$ by \cref{function_apply_elim}.
                        Thus $\apply{U}{\alpha}\in\ran{U}$ by \cref{ran_intro}.
                        Thus $\apply{f}{\alpha}\in\unions{\ran{U}}$ by \cref{unions_iff}.
                    \end{subproof}
                    Thus $f$ is a function to $\unions{\ran{U}}$.
                    Thus $f\in\funs{\dom{X}}{\unions{\ran{U}}}$ by \cref{funs_intro}.
                    Thus $f\in B_1$ by \cref{indexed_product}.
                \end{subproof}
                Thus $f\in B_1$ and $f\in\productFamily{A}$.
                Thus $f\in B_1\inter\productFamily{A}$ by \cref{inter_intro}.
                Show that $B_1\inter\productFamily{A}\neq\emptyset$.
                \begin{subproof}
                    Suppose not.
                    Then $B_1\inter\productFamily{A} = \emptyset$.
                    Contradiction by \cref{emptyset}.
                \end{subproof}
                Thus $B_1$ intersects $\productFamily{A}$ by \cref{intersects}.
            \end{subproof}
            $B_0$ is a basis on $\productFamily{X}$ by \cref{box_basis_is_basis}.
            Thus $B_0$ is a basis for $T_0$ on $\productFamily{X}$ by \cref{basis_for_topology,box_topology}.
            Thus $x\in\closure{\productFamily{A}}{\productFamily{X}}{T_0}$ by \cref{closure_iff_intersects_basis}.
        \end{subproof}
        Thus $\productFamily{C}\subseteq \closure{\productFamily{A}}{\productFamily{X}}{T_0}$ by \cref{subseteq}.
    \end{subproof}
    Show that $\closure{\productFamily{A}}{\productFamily{X}}{T_0}\subseteq\productFamily{C}$.
    \begin{subproof}
        Show that for all $x$ such that $x\in\closure{\productFamily{A}}{\productFamily{X}}{T_0}$ we have
            $x\in\productFamily{C}$.
        \begin{subproof}
            Fix $x$.
            Assume $x\in\closure{\productFamily{A}}{\productFamily{X}}{T_0}$.
            Thus $x\in\productFamily{X}$ by \cref{closure}.
            Show that for all $\alpha\in\dom{X}$ we have
                $\apply{x}{\alpha}\in\closure{\apply{A}{\alpha}}{\apply{X}{\alpha}}{\apply{T}{\alpha}}$.
            \begin{subproof}
                Fix $\alpha$.
                Assume $\alpha\in\dom{X}$.
                Show that for all $V$ such that
                    $V$ is open in $\apply{X}{\alpha}$ with respect to $\apply{T}{\alpha}$ and $\apply{x}{\alpha}\in V$
                    we have $V$ intersects $\apply{A}{\alpha}$.
                \begin{subproof}
                    Fix $V$.
                    Assume $V$ is open in $\apply{X}{\alpha}$ with respect to $\apply{T}{\alpha}$ and $\apply{x}{\alpha}\in V$.
                    Show that $V\subseteq\unions{\ran{X}}$.
                    \begin{subproof}
                        Thus $V\in\apply{T}{\alpha}$ by \cref{open_in}.
                        $\apply{T}{\alpha}\subseteq\pow{\apply{X}{\alpha}}$ by \cref{topology_on}.
                        Thus $V\in\pow{\apply{X}{\alpha}}$ by \cref{subseteq}.
                        Thus $V\subseteq\apply{X}{\alpha}$ by \cref{pow_iff}.
                        Thus $\alpha\mathrel{X}\apply{X}{\alpha}$ by \cref{function_apply_elim}.
                        Thus $\apply{X}{\alpha}\in\ran{X}$ by \cref{ran_intro}.
                        Show that $\apply{X}{\alpha}\subseteq\unions{\ran{X}}$.
                        \begin{subproof}
                            Show that for all $z$ such that $z\in\apply{X}{\alpha}$ we have $z\in\unions{\ran{X}}$.
                            \begin{subproof}
                                Fix $z$.
                                Assume $z\in\apply{X}{\alpha}$.
                                Thus $z\in\unions{\ran{X}}$ by \cref{unions_iff}.
                            \end{subproof}
                            Thus $\apply{X}{\alpha}\subseteq\unions{\ran{X}}$ by \cref{subseteq}.
                        \end{subproof}
                        Thus $V\subseteq\unions{\ran{X}}$ by \cref{subseteq_transitive}.
                    \end{subproof}
                    Take $U$ such that
                        $U\in\funs{\dom{X}}{\pow{\unions{\ran{X}}}}$ and
                        $\apply{U}{\alpha} = V$ and
                        for all $\beta\in\dom{X}$ such that $\beta\neq\alpha$ we have $\apply{U}{\beta} = \apply{X}{\beta}$
                        by \cref{replace_coordinate_function}.
                    Show that for all $\beta\in\dom{X}$ we have $\apply{U}{\beta}\in\apply{T}{\beta}$.
                    \begin{subproof}
                        Fix $\beta$.
                        Assume $\beta\in\dom{X}$.
                        \begin{byCase}
                        \caseOf{$\beta = \alpha$.}
                            Thus $\apply{U}{\beta} = V$.
                            Thus $\apply{U}{\beta}\in\apply{T}{\beta}$ by \cref{open_in}.
                        \caseOf{$\beta\neq\alpha$.}
                            Thus $\apply{U}{\beta} = \apply{X}{\beta}$.
                            Thus $\apply{U}{\beta}\in\apply{T}{\beta}$ by \cref{topology_on}.
                        \end{byCase}
                    \end{subproof}
                    Let $B_1 = \productFamily{U}$.
                    Show that $B_1\in\pow{\productFamily{X}}$.
                    \begin{subproof}
                        Show that for all $y$ such that $y\in B_1$ we have $y\in\productFamily{X}$.
                        \begin{subproof}
                            Fix $y$.
                            Assume $y\in B_1$.
                            Thus $y\in\productFamily{U}$.
                            Thus $y\in\funs{\dom{U}}{\unions{\ran{U}}}$ by \cref{indexed_product}.
                            Thus $y$ is a function and $\dom{y} = \dom{U}$ by \cref{funs_elim}.
                            $\dom{U} = \dom{X}$ by \cref{funs_elim}.
                            Show that for all $\beta\in\dom{X}$ we have $\apply{y}{\beta}\in\apply{X}{\beta}$.
                            \begin{subproof}
                                Fix $\beta$.
                                Assume $\beta\in\dom{X}$.
                                Thus $\beta\in\dom{U}$.
                                Thus $\apply{y}{\beta}\in\apply{U}{\beta}$ by \cref{indexed_product}.
                                $\apply{U}{\beta}\in\apply{T}{\beta}$.
                                $\apply{T}{\beta}\subseteq\pow{\apply{X}{\beta}}$ by \cref{topology_on}.
                                Thus $\apply{U}{\beta}\in\pow{\apply{X}{\beta}}$ by \cref{elem_subseteq}.
                                Thus $\apply{U}{\beta}\subseteq\apply{X}{\beta}$ by \cref{pow_iff}.
                                Thus $\apply{y}{\beta}\in\apply{X}{\beta}$ by \cref{elem_subseteq}.
                            \end{subproof}
                            Show that for all $\beta\in\dom{X}$ we have $\apply{y}{\beta}\in\unions{\ran{X}}$.
                            \begin{subproof}
                                Fix $\beta$.
                                Assume $\beta\in\dom{X}$.
                                Thus $\beta\mathrel{X}\apply{X}{\beta}$ by \cref{function_apply_elim}.
                                Thus $\apply{X}{\beta}\in\ran{X}$ by \cref{ran_intro}.
                                Thus $\apply{y}{\beta}\in\unions{\ran{X}}$ by \cref{unions_iff}.
                            \end{subproof}
                            Thus $y$ is a function to $\unions{\ran{X}}$.
                            Thus $y\in\funs{\dom{X}}{\unions{\ran{X}}}$ by \cref{funs_intro}.
                            Thus $y\in\productFamily{X}$ by \cref{indexed_product}.
                        \end{subproof}
                        Thus $B_1\subseteq\productFamily{X}$ by \cref{subseteq}.
                        Thus $B_1\in\pow{\productFamily{X}}$ by \cref{pow_iff}.
                    \end{subproof}
                    $B_1\in B_0$ by \cref{box_basis}.
                    Thus $B_1\in T_0$ by \cref{basis_elem_open_in_generated,box_topology}.
                    Thus $B_1$ is open in $\productFamily{X}$ with respect to $T_0$ by \cref{open_in}.
                    Show that $x\in B_1$.
                    \begin{subproof}
                        Show that for all $\beta\in\dom{X}$ we have $\apply{x}{\beta}\in\apply{U}{\beta}$.
                        \begin{subproof}
                            Fix $\beta$.
                            Assume $\beta\in\dom{X}$.
                            \begin{byCase}
                            \caseOf{$\beta = \alpha$.}
                                Thus $\apply{x}{\beta}\in V$.
                                Thus $\apply{x}{\beta}\in\apply{U}{\beta}$.
                            \caseOf{$\beta\neq\alpha$.}
                                Thus $\apply{x}{\beta}\in\apply{X}{\beta}$ by \cref{indexed_product}.
                                Thus $\apply{x}{\beta}\in\apply{U}{\beta}$ by \cref{subseteq}.
                            \end{byCase}
                        \end{subproof}
                        Thus $x\in\funs{\dom{X}}{\unions{\ran{X}}}$ by \cref{indexed_product}.
                        Thus $x$ is a function and $\dom{x} = \dom{X}$ by \cref{funs_elim}.
                        $\dom{U} = \dom{X}$ by \cref{funs_elim}.
                        Thus $\dom{x} = \dom{U}$.
                        Show that for all $\beta\in\dom{U}$ we have $\apply{x}{\beta}\in\unions{\ran{U}}$.
                        \begin{subproof}
                            Fix $\beta$.
                            Assume $\beta\in\dom{U}$.
                            $U$ is a function by \cref{funs_elim}.
                            Thus $\beta\mathrel{U}\apply{U}{\beta}$ by \cref{function_apply_elim}.
                            Thus $\apply{U}{\beta}\in\ran{U}$ by \cref{ran_intro}.
                            Thus $\apply{x}{\beta}\in\unions{\ran{U}}$ by \cref{unions_iff}.
                        \end{subproof}
                        Thus $x$ is a function to $\unions{\ran{U}}$.
                        Thus $x\in\funs{\dom{U}}{\unions{\ran{U}}}$ by \cref{funs_intro}.
                        Thus $x\in B_1$ by \cref{indexed_product}.
                    \end{subproof}
                    $T_0$ is a topology on $\productFamily{X}$.
                    $\productFamily{A}\subseteq\productFamily{X}$.
                    Thus for all $U$ such that $U$ is open in $\productFamily{X}$ with respect to $T_0$ and $x\in U$
                        we have $U$ intersects $\productFamily{A}$ by \cref{closure_iff_intersects_open}.
                    Thus $\productFamily{A}$ intersects $B_1$.
                    Take $f$ such that $f\in\productFamily{A}\inter B_1$ by \cref{intersects,nonempty_inhabited}.
                    Thus $f\in\productFamily{A}$ and $f\in B_1$ by \cref{inter}.
                    Thus $f\in\productFamily{U}$.
                    Thus $\apply{f}{\alpha}\in\apply{A}{\alpha}$ by \cref{indexed_product}.
                    $\dom{U} = \dom{X}$ by \cref{funs_elim}.
                    Thus $\alpha\in\dom{U}$.
                    Thus $\apply{f}{\alpha}\in\apply{U}{\alpha}$ by \cref{indexed_product}.
                    Thus $\apply{U}{\alpha} = V$.
                    Thus $\apply{f}{\alpha}\in V$.
                    Thus $\apply{f}{\alpha}\in V$ and $\apply{f}{\alpha}\in\apply{A}{\alpha}$.
                    Thus $\apply{f}{\alpha}\in V\inter\apply{A}{\alpha}$ by \cref{inter_intro}.
                    Show that $V\inter\apply{A}{\alpha}\neq\emptyset$.
                    \begin{subproof}
                        Suppose not.
                        Then $V\inter\apply{A}{\alpha} = \emptyset$.
                        Contradiction by \cref{emptyset}.
                    \end{subproof}
                    Thus $V$ intersects $\apply{A}{\alpha}$ by \cref{intersects}.
                \end{subproof}
                $\apply{T}{\alpha}$ is a topology on $\apply{X}{\alpha}$.
                $\apply{A}{\alpha}\subseteq\apply{X}{\alpha}$.
                Thus $\apply{x}{\alpha}\in\apply{X}{\alpha}$ by \cref{indexed_product}.
                Thus $\apply{x}{\alpha}\in\closure{\apply{A}{\alpha}}{\apply{X}{\alpha}}{\apply{T}{\alpha}}$ by \cref{closure_iff_intersects_open}.
            \end{subproof}
            Show that $x\in\productFamily{C}$.
            \begin{subproof}
                Show that for all $\alpha\in\dom{X}$ we have $\apply{x}{\alpha}\in\apply{C}{\alpha}$.
                \begin{subproof}
                    Fix $\alpha$.
                    Assume $\alpha\in\dom{X}$.
                    Thus $(\alpha,\closure{\apply{A}{\alpha}}{\apply{X}{\alpha}}{\apply{T}{\alpha}})\in C$ by \cref{closure_family}.
                    Thus $\apply{C}{\alpha} = \closure{\apply{A}{\alpha}}{\apply{X}{\alpha}}{\apply{T}{\alpha}}$ by \cref{function_apply_intro}.
                    Thus $\apply{x}{\alpha}\in\closure{\apply{A}{\alpha}}{\apply{X}{\alpha}}{\apply{T}{\alpha}}$.
                    Thus $\apply{x}{\alpha}\in\apply{C}{\alpha}$.
                \end{subproof}
                Thus $x\in\funs{\dom{X}}{\unions{\ran{X}}}$ by \cref{indexed_product}.
                Thus $x$ is a function and $\dom{x} = \dom{X}$ by \cref{funs_elim}.
                Show that $\dom{C} = \dom{X}$.
                \begin{subproof}
                    Show that for all $\alpha\in\dom{X}$ we have $\alpha\in\dom{C}$.
                    \begin{subproof}
                        Fix $\alpha$.
                        Assume $\alpha\in\dom{X}$.
                        Thus $(\alpha,\closure{\apply{A}{\alpha}}{\apply{X}{\alpha}}{\apply{T}{\alpha}})\in C$ by \cref{closure_family}.
                        Thus $\alpha\in\dom{C}$ by \cref{dom}.
                    \end{subproof}
                    Thus $\dom{X}\subseteq\dom{C}$ by \cref{subseteq}.
                    Show that for all $\alpha\in\dom{C}$ we have $\alpha\in\dom{X}$.
                    \begin{subproof}
                        Fix $\alpha$.
                        Assume $\alpha\in\dom{C}$.
                        Take $y$ such that $(\alpha,y)\in C$ by \cref{dom}.
                        Take $\beta\in\dom{X}$ such that $(\alpha,y) =
                            (\beta,\closure{\apply{A}{\beta}}{\apply{X}{\beta}}{\apply{T}{\beta}})$
                            by \cref{closure_family}.
                        Thus $\alpha = \beta$ by \cref{pair_eq_iff}.
                        Thus $\alpha\in\dom{X}$.
                    \end{subproof}
                    Thus $\dom{C}\subseteq\dom{X}$ by \cref{subseteq}.
                    Thus $\dom{C} = \dom{X}$ by \cref{subseteq_antisymmetric}.
                \end{subproof}
                Thus $\dom{x} = \dom{C}$.
                Show that for all $\alpha\in\dom{C}$ we have $\apply{x}{\alpha}\in\unions{\ran{C}}$.
                \begin{subproof}
                    Fix $\alpha$.
                    Assume $\alpha\in\dom{C}$.
                    Thus $\alpha\in\dom{X}$.
                    Thus $\apply{x}{\alpha}\in\apply{C}{\alpha}$.
                    Thus $\alpha\mathrel{C}\apply{C}{\alpha}$ by \cref{function_apply_elim}.
                    Thus $\apply{C}{\alpha}\in\ran{C}$ by \cref{ran_intro}.
                    Thus $\apply{x}{\alpha}\in\unions{\ran{C}}$ by \cref{unions_iff}.
                \end{subproof}
                Thus $x$ is a function to $\unions{\ran{C}}$.
                Thus $x\in\funs{\dom{C}}{\unions{\ran{C}}}$ by \cref{funs_intro}.
                Thus $x\in\productFamily{C}$ by \cref{indexed_product}.
            \end{subproof}
        \end{subproof}
        Thus $\closure{\productFamily{A}}{\productFamily{X}}{T_0}\subseteq\productFamily{C}$ by \cref{subseteq}.
    \end{subproof}
    Thus $\productFamily{C} = \closure{\productFamily{A}}{\productFamily{X}}{T_0}$ by \cref{subseteq_antisymmetric}.
\end{proof}

% from §19 Item 20: closure of products with the product topology
\begin{lemma}\label{product_closure_product}
    Assume $X$ is a function.
    Assume $A$ is a function.
    Assume $T$ is a function.
    Assume $\dom{A} = \dom{X}$.
    Assume $\dom{T} = \dom{X}$.
    Assume $\dom{X}$ is inhabited.
    Assume that for all $\alpha\in\dom{X}$ we have $\apply{A}{\alpha}\subseteq\apply{X}{\alpha}$.
    Assume that for all $\alpha\in\dom{X}$ we have $\apply{T}{\alpha}$ is a topology on $\apply{X}{\alpha}$.
    Let $C = \closureFamily{A}{X}{T}$.
    Then $\productFamily{C} = \closure{\productFamily{A}}{\productFamily{X}}{\productTopologyIndexed{T}{X}}$.
\end{lemma}
\begin{proof}
    Let $T_0 = \productTopologyIndexed{T}{X}$.
    Let $T_1 = \boxTopology{T}{X}$.
    Let $S = \productSubbasis{T}{X}$.
    Show that for all $\alpha\in\dom{X}$ we have $\apply{T}{\alpha}$ is a topology on $\apply{X}{\alpha}$.
    \begin{subproof}
        Fix $\alpha$.
        Assume $\alpha\in\dom{X}$.
        Thus $\apply{T}{\alpha}$ is a topology on $\apply{X}{\alpha}$ by \cref{subseteq}.
    \end{subproof}
    Show that $C$ is a function.
    \begin{subproof}
        Show that $C$ is a relation.
        \begin{subproof}
            Show that for all $w$ such that $w\in C$ there exist $a,b$ such that $w = (a,b)$.
            \begin{subproof}
                Fix $w$.
                Assume $w\in C$.
                Take $\alpha\in\dom{X}$ such that
                    $w = (\alpha,\closure{\apply{A}{\alpha}}{\apply{X}{\alpha}}{\apply{T}{\alpha}})$
                    by \cref{closure_family}.
            \end{subproof}
            Thus $C$ is a relation by \cref{relation}.
        \end{subproof}
        Show that $C$ is right-unique.
        \begin{subproof}
            Show that for all $a,b_1,b_2$ such that $(a,b_1)\in C$ and $(a,b_2)\in C$ we have $b_1 = b_2$.
            \begin{subproof}
                Fix $a$.
                Fix $b_1$.
                Fix $b_2$.
                Assume $(a,b_1)\in C$ and $(a,b_2)\in C$.
                Take $\alpha_1\in\dom{X}$ such that $(a,b_1) =
                    (\alpha_1,\closure{\apply{A}{\alpha_1}}{\apply{X}{\alpha_1}}{\apply{T}{\alpha_1}})$
                    by \cref{closure_family}.
                Take $\alpha_2\in\dom{X}$ such that $(a,b_2) =
                    (\alpha_2,\closure{\apply{A}{\alpha_2}}{\apply{X}{\alpha_2}}{\apply{T}{\alpha_2}})$
                    by \cref{closure_family}.
                Thus $a = \alpha_1$ and $b_1 = \closure{\apply{A}{\alpha_1}}{\apply{X}{\alpha_1}}{\apply{T}{\alpha_1}}$ by \cref{pair_eq_iff}.
                Thus $a = \alpha_2$ and $b_2 = \closure{\apply{A}{\alpha_2}}{\apply{X}{\alpha_2}}{\apply{T}{\alpha_2}}$ by \cref{pair_eq_iff}.
                Thus $\alpha_1 = \alpha_2$.
                Thus $b_1 = b_2$.
            \end{subproof}
            Thus $C$ is right-unique.
        \end{subproof}
        Thus $C$ is a function.
    \end{subproof}
    $\productFamily{A}$ is a subspace of $\productFamily{X}$ with respect to $T_0$ by \cref{product_subspace_product}.
    Thus $\productFamily{A}\subseteq\productFamily{X}$ by \cref{subspace_of}.
    Show that $\closure{\productFamily{A}}{\productFamily{X}}{T_0}\subseteq\productFamily{C}$.
    \begin{subproof}
        Show that for all $x$ such that $x\in\closure{\productFamily{A}}{\productFamily{X}}{T_0}$ we have
            $x\in\productFamily{C}$.
        \begin{subproof}
            Fix $x$.
            Assume $x\in\closure{\productFamily{A}}{\productFamily{X}}{T_0}$.
            Thus $x\in\productFamily{X}$ by \cref{closure}.
            Show that $\dom{C} = \dom{X}$.
            \begin{subproof}
                Show that for all $\alpha\in\dom{X}$ we have $\alpha\in\dom{C}$.
                \begin{subproof}
                    Fix $\alpha$.
                    Assume $\alpha\in\dom{X}$.
                    Thus $(\alpha,\closure{\apply{A}{\alpha}}{\apply{X}{\alpha}}{\apply{T}{\alpha}})\in C$ by \cref{closure_family}.
                    Thus $\alpha\in\dom{C}$ by \cref{dom}.
                \end{subproof}
                Thus $\dom{X}\subseteq\dom{C}$ by \cref{subseteq}.
                Show that for all $\alpha\in\dom{C}$ we have $\alpha\in\dom{X}$.
                \begin{subproof}
                    Fix $\alpha$.
                    Assume $\alpha\in\dom{C}$.
                    Take $y$ such that $(\alpha,y)\in C$ by \cref{dom}.
                    Take $\beta\in\dom{X}$ such that $(\alpha,y) =
                        (\beta,\closure{\apply{A}{\beta}}{\apply{X}{\beta}}{\apply{T}{\beta}})$
                        by \cref{closure_family}.
                    Thus $\alpha = \beta$ by \cref{pair_eq_iff}.
                    Thus $\alpha\in\dom{X}$.
                \end{subproof}
                Thus $\dom{C}\subseteq\dom{X}$ by \cref{subseteq}.
                Thus $\dom{C} = \dom{X}$ by \cref{subseteq_antisymmetric}.
            \end{subproof}
            Show that for all $\alpha\in\dom{X}$ we have
                $\apply{x}{\alpha}\in\closure{\apply{A}{\alpha}}{\apply{X}{\alpha}}{\apply{T}{\alpha}}$.
            \begin{subproof}
                Fix $\alpha$.
                Assume $\alpha\in\dom{X}$.
                Show that for all $V$ such that
                    $V$ is open in $\apply{X}{\alpha}$ with respect to $\apply{T}{\alpha}$ and $\apply{x}{\alpha}\in V$
                    we have $V$ intersects $\apply{A}{\alpha}$.
                \begin{subproof}
                    Fix $V$.
                    Assume $V$ is open in $\apply{X}{\alpha}$ with respect to $\apply{T}{\alpha}$ and $\apply{x}{\alpha}\in V$.
                    Let $W = \preimg{\piIndex{X}{\alpha}}{V}$.
                    Show that $W\in\productSubbasisAt{T}{X}{\alpha}$.
                    \begin{subproof}
                        Show that $W\in\pow{\productFamily{X}}$.
                        \begin{subproof}
                            Show that for all $z$ such that $z\in W$ we have $z\in\productFamily{X}$.
                            \begin{subproof}
                                Fix $z$.
                                Assume $z\in W$.
                                Take $y\in V$ such that $z\mathrel{\piIndex{X}{\alpha}} y$ by \cref{preimg_iff}.
                                Thus $(z,y)\in\piIndex{X}{\alpha}$.
                                Take $p\in\productFamily{X}$ such that $(z,y) = (p,\apply{p}{\alpha})$ by \cref{projection_map_indexed}.
                                Thus $z = p$ by \cref{pair_eq_iff}.
                                Thus $z\in\productFamily{X}$.
                            \end{subproof}
                            Thus $W\subseteq\productFamily{X}$ by \cref{subseteq}.
                            Thus $W\in\pow{\productFamily{X}}$ by \cref{pow_iff}.
                        \end{subproof}
                        $V\in\apply{T}{\alpha}$ by \cref{open_in}.
                        Thus $W\in\productSubbasisAt{T}{X}{\alpha}$ by \cref{product_subbasis_indexed}.
                    \end{subproof}
                    Show that $W\in S$.
                    \begin{subproof}
                        $W\in\productSubbasisAt{T}{X}{\alpha}$.
                        Thus $W\in\pow{\productFamily{X}}$ by \cref{product_subbasis_indexed}.
                        Thus there exists $\beta\in\dom{X}$ such that $W\in\productSubbasisAt{T}{X}{\beta}$.
                        Thus $W\in S$ by \cref{product_subbasis_union_indexed}.
                    \end{subproof}
                    Let $F = \{W\}$.
                    Show that $F\subseteq S$.
                    \begin{subproof}
                        Show that for all $z$ such that $z\in F$ we have $z\in S$.
                        \begin{subproof}
                            Fix $z$.
                            Assume $z\in F$.
                            Thus $z = W$ by \cref{singleton_elim}.
                            Thus $z\in S$.
                        \end{subproof}
                        Thus $F\subseteq S$ by \cref{subseteq}.
                    \end{subproof}
                    Let $G = \{W,W\}$.
                    $G$ is finite by \cref{pair_is_finite}.
                    Show that $F\subseteq G$.
                    \begin{subproof}
                        Show that for all $z$ such that $z\in F$ we have $z\in G$.
                        \begin{subproof}
                            Fix $z$.
                            Assume $z\in F$.
                            Thus $z = W$ by \cref{singleton_elim}.
                            Thus $z\in G$ by \cref{upair_intro_left}.
                        \end{subproof}
                        Thus $F\subseteq G$ by \cref{subseteq}.
                    \end{subproof}
                    Thus $F$ is finite by \cref{subseteq_finite_is_finite}.
                    $F$ is inhabited by \cref{singleton_intro}.
                    $\inters{F} = W$ by \cref{inters_singleton}.
                    Show that $W\in\pow{\unions{S}}$.
                    \begin{subproof}
                        Show that for all $z$ such that $z\in W$ we have $z\in\unions{S}$.
                        \begin{subproof}
                            Fix $z$.
                            Assume $z\in W$.
                            Thus $z\in\unions{S}$ by \cref{unions_iff}.
                        \end{subproof}
                        Thus $W\subseteq\unions{S}$ by \cref{subseteq}.
                        Thus $W\in\pow{\unions{S}}$ by \cref{pow_iff}.
                    \end{subproof}
                    Thus $W\in\finiteIntersections{S}$ by \cref{finite_intersections}.
                    $S$ is a subbasis on $\productFamily{X}$ by \cref{product_subbasis_is_subbasis}.
                    Thus $\finiteIntersections{S}$ is a basis for $T_0$ on $\productFamily{X}$ by \cref{product_subbasis_finite_intersections_basis}.
                    Thus $W\in T_0$ by \cref{basis_elem_open_in_generated,basis_for_topology}.
                    $T_0$ is a topology on $\productFamily{X}$ by \cref{product_subbasis_finite_intersections_basis,basis_for_topology,basis_topology_is_topology}.
                    Thus $W$ is open in $\productFamily{X}$ with respect to $T_0$ by \cref{open_in}.
                    $\apply{x}{\alpha}\in V$.
                    Thus $(x,\apply{x}{\alpha})\in\piIndex{X}{\alpha}$ by \cref{projection_map_indexed}.
                    Thus $x\mathrel{\piIndex{X}{\alpha}}\apply{x}{\alpha}$.
                    Thus $x\in W$ by \cref{preimg_iff}.
                    $T_0$ is a topology on $\productFamily{X}$.
                    $\productFamily{A}\subseteq\productFamily{X}$ by \cref{subspace_of}.
                    Thus for all $U$ such that $U$ is open in $\productFamily{X}$ with respect to $T_0$ and $x\in U$
                        we have $U$ intersects $\productFamily{A}$ by \cref{closure_iff_intersects_open}.
                    Thus $W$ intersects $\productFamily{A}$.
                    Take $f$ such that $f\in W\inter\productFamily{A}$ by \cref{intersects,nonempty_inhabited}.
                    Thus $f\in W$ and $f\in\productFamily{A}$ by \cref{inter}.
                    Take $y\in V$ such that $f\mathrel{\piIndex{X}{\alpha}} y$ by \cref{preimg_iff}.
                    Thus $(f,y)\in\piIndex{X}{\alpha}$.
                    Take $p\in\productFamily{X}$ such that $(f,y) = (p,\apply{p}{\alpha})$ by \cref{projection_map_indexed}.
                    Thus $f = p$ and $y = \apply{p}{\alpha}$ by \cref{pair_eq_iff}.
                    Thus $\apply{f}{\alpha}\in V$.
                    Thus $\apply{f}{\alpha}\in\apply{A}{\alpha}$ by \cref{indexed_product}.
                    Thus $\apply{f}{\alpha}\in V$ and $\apply{f}{\alpha}\in\apply{A}{\alpha}$.
                    Thus $\apply{f}{\alpha}\in V\inter\apply{A}{\alpha}$ by \cref{inter_intro}.
                    Show that $V\inter\apply{A}{\alpha}\neq\emptyset$.
                    \begin{subproof}
                        Suppose not.
                        Then $V\inter\apply{A}{\alpha} = \emptyset$.
                        Contradiction by \cref{emptyset}.
                    \end{subproof}
                    Thus $V$ intersects $\apply{A}{\alpha}$ by \cref{intersects}.
                \end{subproof}
                $\apply{T}{\alpha}$ is a topology on $\apply{X}{\alpha}$.
                $\apply{A}{\alpha}\subseteq\apply{X}{\alpha}$.
                Thus $\apply{x}{\alpha}\in\apply{X}{\alpha}$ by \cref{indexed_product}.
                Thus $\apply{x}{\alpha}\in\closure{\apply{A}{\alpha}}{\apply{X}{\alpha}}{\apply{T}{\alpha}}$ by \cref{closure_iff_intersects_open}.
            \end{subproof}
            Show that $x\in\productFamily{C}$.
            \begin{subproof}
                Show that for all $\alpha\in\dom{X}$ we have $\apply{x}{\alpha}\in\apply{C}{\alpha}$.
                \begin{subproof}
                    Fix $\alpha$.
                    Assume $\alpha\in\dom{X}$.
                    Thus $(\alpha,\closure{\apply{A}{\alpha}}{\apply{X}{\alpha}}{\apply{T}{\alpha}})\in C$ by \cref{closure_family}.
                    Thus $\apply{C}{\alpha} = \closure{\apply{A}{\alpha}}{\apply{X}{\alpha}}{\apply{T}{\alpha}}$ by \cref{function_apply_intro}.
                    Thus $\apply{x}{\alpha}\in\closure{\apply{A}{\alpha}}{\apply{X}{\alpha}}{\apply{T}{\alpha}}$.
                    Thus $\apply{x}{\alpha}\in\apply{C}{\alpha}$.
                \end{subproof}
                Thus $x\in\funs{\dom{X}}{\unions{\ran{X}}}$ by \cref{indexed_product}.
                Thus $x$ is a function and $\dom{x} = \dom{X}$ by \cref{funs_elim}.
                Thus $\dom{x} = \dom{C}$.
                Show that for all $\alpha\in\dom{C}$ we have $\apply{x}{\alpha}\in\unions{\ran{C}}$.
                \begin{subproof}
                    Fix $\alpha$.
                    Assume $\alpha\in\dom{C}$.
                    Thus $\apply{x}{\alpha}\in\apply{C}{\alpha}$.
                    Thus $\alpha\mathrel{C}\apply{C}{\alpha}$ by \cref{function_apply_elim}.
                    Thus $\apply{C}{\alpha}\in\ran{C}$ by \cref{ran_intro}.
                    Thus $\apply{x}{\alpha}\in\unions{\ran{C}}$ by \cref{unions_iff}.
                \end{subproof}
                Thus $x$ is a function to $\unions{\ran{C}}$.
                Thus $x\in\funs{\dom{C}}{\unions{\ran{C}}}$ by \cref{funs_intro}.
                Thus $x\in\productFamily{C}$ by \cref{indexed_product}.
            \end{subproof}
        \end{subproof}
        Thus $\closure{\productFamily{A}}{\productFamily{X}}{T_0}\subseteq\productFamily{C}$ by \cref{subseteq}.
    \end{subproof}
    Show that $\productFamily{C}\subseteq\closure{\productFamily{A}}{\productFamily{X}}{T_0}$.
    \begin{subproof}
        $\productFamily{C} = \closure{\productFamily{A}}{\productFamily{X}}{T_1}$ by \cref{product_closure_box}.
        Show that $T_0\subseteq T_1$.
        \begin{subproof}
            Show that for all $U$ such that $U\in T_0$ we have $U\in T_1$.
            \begin{subproof}
                Fix $U$.
                Assume $U\in T_0$.
                $S$ is a subbasis on $\productFamily{X}$ by \cref{product_subbasis_is_subbasis}.
                Thus $\finiteIntersections{S}$ is a basis on $\productFamily{X}$ by \cref{subbasis_finite_intersections_basis}.
                $T_0 = \subbasisTopology{S}{\productFamily{X}}$ by \cref{product_topology_indexed}.
                Thus $U\in\basisTopology{\finiteIntersections{S}}{\productFamily{X}}$ by \cref{topology_generated_by_subbasis}.
                Show that for all $x$ such that $x\in U$ there exists $B\in\finiteIntersections{S}$ such that $x\in B$ and $B\subseteq U$.
                \begin{subproof}
                    Fix $x$.
                    Assume $x\in U$.
                    Take $B\in\finiteIntersections{S}$ such that $x\in B$ and $B\subseteq U$ by \cref{topology_generated_by_basis}.
                \end{subproof}
                Let $M = \{B\in\finiteIntersections{S} \mid B\subseteq U\}$.
                Show that for all $x$ such that $x\in U$ we have $x\in\unions{M}$.
                \begin{subproof}
                    Fix $x$.
                    Assume $x\in U$.
                    Take $B\in\finiteIntersections{S}$ such that $x\in B$ and $B\subseteq U$ by \cref{topology_generated_by_basis}.
                    Thus $B\in M$.
                    Thus $x\in\unions{M}$ by \cref{unions_iff}.
                \end{subproof}
                Thus $U\subseteq\unions{M}$ by \cref{subseteq}.
                Show that $\unions{M}\subseteq U$.
                \begin{subproof}
                    Show that for all $x$ such that $x\in\unions{M}$ we have $x\in U$.
                    \begin{subproof}
                        Fix $x$.
                        Assume $x\in\unions{M}$.
                        Take $B\in M$ such that $x\in B$ by \cref{unions_iff}.
                        Thus $B\subseteq U$.
                        Thus $x\in U$ by \cref{subseteq}.
                    \end{subproof}
                    Thus $\unions{M}\subseteq U$ by \cref{subseteq}.
                \end{subproof}
                Thus $U = \unions{M}$ by \cref{subseteq_antisymmetric}.
                Show that for all $B$ such that $B\in\finiteIntersections{S}$ we have $B\in T_1$.
                \begin{subproof}
                    Fix $B$.
                    Assume $B\in\finiteIntersections{S}$.
                    Take $F\subseteq S$ such that $F$ is finite and inhabited and $B = \inters{F}$ by \cref{finite_intersections}.
                    Show that for all $A$ such that $A\in F$ we have $A\in T_1$.
                    \begin{subproof}
                        Fix $A$.
                        Assume $A\in F$.
                        Thus $A\in S$ by \cref{subseteq}.
                        Take $\beta\in\dom{X}$ such that $A\in\productSubbasisAt{T}{X}{\beta}$ by \cref{product_subbasis_union_indexed}.
                        Take $V\in\apply{T}{\beta}$ such that $A = \preimg{\piIndex{X}{\beta}}{V}$ by \cref{product_subbasis_indexed}.
                        Show that for all $x$ such that $x\in A$ there exists $B_1\in\boxBasis{T}{X}$ such that $x\in B_1$ and $B_1\subseteq A$.
                        \begin{subproof}
                            Fix $x$.
                            Assume $x\in A$.
                            $V$ is open in $\apply{X}{\beta}$ with respect to $\apply{T}{\beta}$ by \cref{open_in}.
                            Show that $V\subseteq\unions{\ran{X}}$.
                            \begin{subproof}
                                $\apply{T}{\beta}\subseteq\pow{\apply{X}{\beta}}$ by \cref{topology_on}.
                                Thus $V\in\pow{\apply{X}{\beta}}$ by \cref{subseteq}.
                                Thus $V\subseteq\apply{X}{\beta}$ by \cref{pow_iff}.
                                Thus $\beta\mathrel{X}\apply{X}{\beta}$ by \cref{function_apply_elim}.
                                Thus $\apply{X}{\beta}\in\ran{X}$ by \cref{ran_intro}.
                                Show that $\apply{X}{\beta}\subseteq\unions{\ran{X}}$.
                                \begin{subproof}
                                    Show that for all $z$ such that $z\in\apply{X}{\beta}$ we have $z\in\unions{\ran{X}}$.
                                    \begin{subproof}
                                        Fix $z$.
                                        Assume $z\in\apply{X}{\beta}$.
                                        Thus $z\in\unions{\ran{X}}$ by \cref{unions_iff}.
                                    \end{subproof}
                                    Thus $\apply{X}{\beta}\subseteq\unions{\ran{X}}$ by \cref{subseteq}.
                                \end{subproof}
                                Thus $V\subseteq\unions{\ran{X}}$ by \cref{subseteq_transitive}.
                            \end{subproof}
                            Take $U$ such that
                                $U\in\funs{\dom{X}}{\pow{\unions{\ran{X}}}}$ and
                                $\apply{U}{\beta} = V$ and
                                for all $\gamma\in\dom{X}$ such that $\gamma\neq\beta$ we have $\apply{U}{\gamma} = \apply{X}{\gamma}$
                                by \cref{replace_coordinate_function}.
                            Show that for all $\gamma\in\dom{X}$ we have $\apply{U}{\gamma}\in\apply{T}{\gamma}$.
                            \begin{subproof}
                                Fix $\gamma$.
                                Assume $\gamma\in\dom{X}$.
                                \begin{byCase}
                                \caseOf{$\gamma = \beta$.}
                                    Thus $\apply{U}{\gamma} = V$.
                                    Thus $\apply{U}{\gamma}\in\apply{T}{\gamma}$ by \cref{open_in}.
                                \caseOf{$\gamma\neq\beta$.}
                                    Thus $\apply{U}{\gamma} = \apply{X}{\gamma}$.
                                    Thus $\apply{U}{\gamma}\in\apply{T}{\gamma}$ by \cref{topology_on}.
                                \end{byCase}
                            \end{subproof}
                            Let $B_1 = \productFamily{U}$.
                            Show that $B_1\in\pow{\productFamily{X}}$.
                            \begin{subproof}
                                Show that for all $z$ such that $z\in B_1$ we have $z\in\productFamily{X}$.
                                \begin{subproof}
                                    Fix $z$.
                                    Assume $z\in B_1$.
                                    Thus $z\in\productFamily{U}$.
                                    Thus $z\in\funs{\dom{U}}{\unions{\ran{U}}}$ by \cref{indexed_product}.
                                    Thus $z$ is a function and $\dom{z} = \dom{U}$ by \cref{funs_elim}.
                                    $\dom{U} = \dom{X}$ by \cref{funs_elim}.
                                    Show that for all $\gamma\in\dom{X}$ we have $\apply{z}{\gamma}\in\apply{X}{\gamma}$.
                                    \begin{subproof}
                                        Fix $\gamma$.
                                        Assume $\gamma\in\dom{X}$.
                                        Thus $\gamma\in\dom{U}$.
                                        Thus $\apply{z}{\gamma}\in\apply{U}{\gamma}$ by \cref{indexed_product}.
                                        $\apply{U}{\gamma}\in\apply{T}{\gamma}$.
                                        $\apply{T}{\gamma}\subseteq\pow{\apply{X}{\gamma}}$ by \cref{topology_on}.
                                        Thus $\apply{U}{\gamma}\in\pow{\apply{X}{\gamma}}$ by \cref{elem_subseteq}.
                                        Thus $\apply{U}{\gamma}\subseteq\apply{X}{\gamma}$ by \cref{pow_iff}.
                                        Thus $\apply{z}{\gamma}\in\apply{X}{\gamma}$ by \cref{elem_subseteq}.
                                    \end{subproof}
                                    Show that for all $\gamma\in\dom{X}$ we have $\apply{z}{\gamma}\in\unions{\ran{X}}$.
                                    \begin{subproof}
                                        Fix $\gamma$.
                                        Assume $\gamma\in\dom{X}$.
                                        Thus $\gamma\mathrel{X}\apply{X}{\gamma}$ by \cref{function_apply_elim}.
                                        Thus $\apply{X}{\gamma}\in\ran{X}$ by \cref{ran_intro}.
                                        Thus $\apply{z}{\gamma}\in\unions{\ran{X}}$ by \cref{unions_iff}.
                                    \end{subproof}
                                    Thus $z$ is a function to $\unions{\ran{X}}$.
                                    Thus $z\in\funs{\dom{X}}{\unions{\ran{X}}}$ by \cref{funs_intro}.
                                    Thus $z\in\productFamily{X}$ by \cref{indexed_product}.
                                \end{subproof}
                                Thus $B_1\subseteq\productFamily{X}$ by \cref{subseteq}.
                                Thus $B_1\in\pow{\productFamily{X}}$ by \cref{pow_iff}.
                            \end{subproof}
                            $B_1\in\boxBasis{T}{X}$ by \cref{box_basis}.
                            Show that $x\in B_1$.
                            \begin{subproof}
                                Take $y\in V$ such that $x\mathrel{\piIndex{X}{\beta}} y$ by \cref{preimg_iff}.
                                Thus $x\in\productFamily{X}$ by \cref{projection_map_indexed,pair_eq_iff}.
                                Thus $(x,y)\in\piIndex{X}{\beta}$.
                                Take $p\in\productFamily{X}$ such that $(x,y) = (p,\apply{p}{\beta})$ by \cref{projection_map_indexed}.
                                Thus $x = p$ and $y = \apply{p}{\beta}$ by \cref{pair_eq_iff}.
                                Thus $\apply{x}{\beta}\in V$.
                                Show that for all $\gamma\in\dom{X}$ we have $\apply{x}{\gamma}\in\apply{U}{\gamma}$.
                                \begin{subproof}
                                    Fix $\gamma$.
                                    Assume $\gamma\in\dom{X}$.
                                    \begin{byCase}
                                    \caseOf{$\gamma = \beta$.}
                                        Thus $\apply{x}{\gamma}\in V$.
                                        Thus $\apply{x}{\gamma}\in\apply{U}{\gamma}$.
                                    \caseOf{$\gamma\neq\beta$.}
                                        Thus $\apply{x}{\gamma}\in\apply{X}{\gamma}$ by \cref{indexed_product}.
                                        Thus $\apply{x}{\gamma}\in\apply{U}{\gamma}$ by \cref{subseteq}.
                                    \end{byCase}
                                \end{subproof}
                                Thus $x\in\funs{\dom{X}}{\unions{\ran{X}}}$ by \cref{indexed_product}.
                                Thus $x$ is a function and $\dom{x} = \dom{X}$ by \cref{funs_elim}.
                                $\dom{U} = \dom{X}$ by \cref{funs_elim}.
                                Thus $\dom{x} = \dom{U}$.
                                Show that for all $\gamma\in\dom{U}$ we have $\apply{x}{\gamma}\in\unions{\ran{U}}$.
                                \begin{subproof}
                                    Fix $\gamma$.
                                    Assume $\gamma\in\dom{U}$.
                                    $U$ is a function by \cref{funs_elim}.
                                    Thus $\gamma\mathrel{U}\apply{U}{\gamma}$ by \cref{function_apply_elim}.
                                    Thus $\apply{U}{\gamma}\in\ran{U}$ by \cref{ran_intro}.
                                    Thus $\apply{x}{\gamma}\in\unions{\ran{U}}$ by \cref{unions_iff}.
                                \end{subproof}
                                Thus $x$ is a function to $\unions{\ran{U}}$.
                                Thus $x\in\funs{\dom{U}}{\unions{\ran{U}}}$ by \cref{funs_intro}.
                                Thus $x\in B_1$ by \cref{indexed_product}.
                            \end{subproof}
                            Show that $B_1\subseteq A$.
                            \begin{subproof}
                                Show that for all $z$ such that $z\in B_1$ we have $z\in A$.
                                \begin{subproof}
                                    Fix $z$.
                                    Assume $z\in B_1$.
                                    Thus $z\in\productFamily{U}$.
                                    Thus $\beta\in\dom{X}$.
                                    Thus $\beta\in\dom{U}$ by \cref{funs_elim}.
                                    Thus $\apply{z}{\beta}\in\apply{U}{\beta}$ by \cref{indexed_product}.
                                    Thus $\apply{U}{\beta} = V$.
                                    Thus $\apply{z}{\beta}\in V$.
                                    $B_1\in\pow{\productFamily{X}}$.
                                    Thus $B_1\subseteq\productFamily{X}$ by \cref{pow_iff}.
                                    Thus $z\in\productFamily{X}$ by \cref{elem_subseteq}.
                                    Thus $(z,\apply{z}{\beta})\in\piIndex{X}{\beta}$ by \cref{projection_map_indexed}.
                                    Thus $z\in A$ by \cref{preimg_iff}.
                                \end{subproof}
                                Thus $B_1\subseteq A$ by \cref{subseteq}.
                            \end{subproof}
                        \end{subproof}
                        Show that $A\in\pow{\productFamily{X}}$.
                        \begin{subproof}
                            Show that for all $x$ such that $x\in A$ we have $x\in\productFamily{X}$.
                            \begin{subproof}
                                Fix $x$.
                                Assume $x\in A$.
                                Take $y\in V$ such that $x\mathrel{\piIndex{X}{\beta}} y$ by \cref{preimg_iff}.
                                Thus $x\in\productFamily{X}$ by \cref{projection_map_indexed,pair_eq_iff}.
                            \end{subproof}
                            Thus $A\subseteq\productFamily{X}$ by \cref{subseteq}.
                            Thus $A\in\pow{\productFamily{X}}$ by \cref{pow_iff}.
                        \end{subproof}
                        Thus $A\in T_1$ by \cref{topology_generated_by_basis,box_topology}.
                    \end{subproof}
                    Thus $F\subseteq T_1$ by \cref{subseteq}.
                    $\boxBasis{T}{X}$ is a basis on $\productFamily{X}$ by \cref{box_basis_is_basis}.
                    $T_1$ is a topology on $\productFamily{X}$ by \cref{basis_topology_is_topology,box_topology}.
                    Thus $B\in T_1$ by \cref{finite_intersections_open_in_topology}.
                \end{subproof}
                Thus $M\subseteq T_1$ by \cref{subseteq}.
                $\boxBasis{T}{X}$ is a basis on $\productFamily{X}$ by \cref{box_basis_is_basis}.
                $T_1$ is a topology on $\productFamily{X}$ by \cref{basis_topology_is_topology,box_topology}.
                Thus $\unions{M}\in T_1$ by \cref{topology_on}.
                Thus $U\in T_1$.
            \end{subproof}
        \end{subproof}
        Show that for all $x$ such that $x\in\closure{\productFamily{A}}{\productFamily{X}}{T_1}$ we have
            $x\in\closure{\productFamily{A}}{\productFamily{X}}{T_0}$.
        \begin{subproof}
            Fix $x$.
            Assume $x\in\closure{\productFamily{A}}{\productFamily{X}}{T_1}$.
            Show that for all $U$ such that
                $U$ is open in $\productFamily{X}$ with respect to $T_0$ and $x\in U$
                we have $U$ intersects $\productFamily{A}$.
            \begin{subproof}
                Fix $U$.
                Assume $U$ is open in $\productFamily{X}$ with respect to $T_0$ and $x\in U$.
                Thus $U\in T_0$ by \cref{open_in}.
                Thus $U\in T_1$ by \cref{subseteq}.
                $\boxBasis{T}{X}$ is a basis on $\productFamily{X}$ by \cref{box_basis_is_basis}.
                $T_1$ is a topology on $\productFamily{X}$ by \cref{basis_topology_is_topology,box_topology}.
                Thus $U$ is open in $\productFamily{X}$ with respect to $T_1$ by \cref{open_in}.
                $T_0$ is a topology on $\productFamily{X}$ by \cref{open_in}.
                $T_0\subseteq\pow{\productFamily{X}}$ by \cref{topology_on}.
                Thus $U\in\pow{\productFamily{X}}$ by \cref{subseteq}.
                Thus $U\subseteq\productFamily{X}$ by \cref{pow_iff}.
                Thus $x\in\productFamily{X}$ by \cref{elem_subseteq}.
                $\productFamily{A}\subseteq\productFamily{X}$ by \cref{subspace_of}.
                Thus $U$ intersects $\productFamily{A}$ by \cref{closure_iff_intersects_open}.
            \end{subproof}
            $S$ is a subbasis on $\productFamily{X}$ by \cref{product_subbasis_is_subbasis}.
            Thus $\finiteIntersections{S}$ is a basis for $T_0$ on $\productFamily{X}$ by \cref{product_subbasis_finite_intersections_basis}.
            Thus $T_0$ is a topology on $\productFamily{X}$ by \cref{basis_topology_is_topology,basis_for_topology}.
            $\productFamily{A}\subseteq\productFamily{X}$ by \cref{subspace_of}.
            $\boxBasis{T}{X}$ is a basis on $\productFamily{X}$ by \cref{box_basis_is_basis}.
            $T_1$ is a topology on $\productFamily{X}$ by \cref{basis_topology_is_topology,box_topology}.
            $\productFamily{X}$ is closed in $\productFamily{X}$ with respect to $T_1$ by \cref{closed_whole_space}.
            Thus $\closure{\productFamily{A}}{\productFamily{X}}{T_1}\subseteq\productFamily{X}$ by \cref{closure_subseteq_closed}.
            Thus $x\in\productFamily{X}$ by \cref{elem_subseteq}.
            Thus $x\in\closure{\productFamily{A}}{\productFamily{X}}{T_0}$ by \cref{closure_iff_intersects_open}.
        \end{subproof}
        Thus $\productFamily{C}\subseteq\closure{\productFamily{A}}{\productFamily{X}}{T_0}$ by \cref{subseteq}.
    \end{subproof}
    Thus $\productFamily{C} = \closure{\productFamily{A}}{\productFamily{X}}{T_0}$ by \cref{subseteq_antisymmetric}.
\end{proof}

% from §19 helper lemma 3: projection map is a function
\begin{lemma}\label{projection_map_indexed_function}
    Assume $X$ is a function.
    Assume $\beta\in\dom{X}$.
    Then $\piIndex{X}{\beta}\in\funs{\productFamily{X}}{\apply{X}{\beta}}$.
\end{lemma}
\begin{proof}
    Let $P = \piIndex{X}{\beta}$.
    Show that $P$ is a relation.
    \begin{subproof}
        Show that for all $w$ such that $w\in P$ there exist $x,y$ such that $w = (x,y)$.
        \begin{subproof}
            Fix $w$.
            Assume $w\in P$.
            Take $x\in\productFamily{X}$ such that $w = (x,\apply{x}{\beta})$ by \cref{projection_map_indexed}.
        \end{subproof}
        Thus $P$ is a relation by \cref{relation}.
    \end{subproof}
    Show that $P$ is right-unique.
    \begin{subproof}
        Show that for all $x,y_1,y_2$ such that $(x,y_1)\in P$ and $(x,y_2)\in P$ we have $y_1 = y_2$.
        \begin{subproof}
            Fix $x$.
            Fix $y_1$.
            Fix $y_2$.
            Assume $(x,y_1)\in P$ and $(x,y_2)\in P$.
            Take $x_1\in\productFamily{X}$ such that $(x,y_1) = (x_1,\apply{x_1}{\beta})$ by \cref{projection_map_indexed}.
            Take $x_2\in\productFamily{X}$ such that $(x,y_2) = (x_2,\apply{x_2}{\beta})$ by \cref{projection_map_indexed}.
            Thus $x = x_1$ and $y_1 = \apply{x_1}{\beta}$ by \cref{pair_eq_iff}.
            Thus $x = x_2$ and $y_2 = \apply{x_2}{\beta}$ by \cref{pair_eq_iff}.
            Thus $y_1 = \apply{x}{\beta}$ and $y_2 = \apply{x}{\beta}$.
            Thus $y_1 = y_2$.
        \end{subproof}
        Thus $P$ is right-unique.
    \end{subproof}
    Thus $P$ is a function.
    Show that $\dom{P} = \productFamily{X}$.
    \begin{subproof}
        Show that for all $x\in\productFamily{X}$ we have $x\in\dom{P}$.
        \begin{subproof}
            Fix $x$.
            Assume $x\in\productFamily{X}$.
            Thus $(x,\apply{x}{\beta})\in P$ by \cref{projection_map_indexed}.
            Thus $x\in\dom{P}$ by \cref{dom}.
        \end{subproof}
        Thus $\productFamily{X}\subseteq\dom{P}$ by \cref{subseteq}.
        Show that for all $x\in\dom{P}$ we have $x\in\productFamily{X}$.
        \begin{subproof}
            Fix $x$.
            Assume $x\in\dom{P}$.
            Take $y$ such that $(x,y)\in P$ by \cref{dom}.
            Take $x_0\in\productFamily{X}$ such that $(x,y) = (x_0,\apply{x_0}{\beta})$ by \cref{projection_map_indexed}.
            Thus $x = x_0$ by \cref{pair_eq_iff}.
            Thus $x\in\productFamily{X}$.
        \end{subproof}
        Thus $\dom{P}\subseteq\productFamily{X}$ by \cref{subseteq}.
        Thus $\dom{P} = \productFamily{X}$ by \cref{subseteq_antisymmetric}.
    \end{subproof}
    Show that for all $x\in\dom{P}$ we have $\apply{P}{x}\in\apply{X}{\beta}$.
    \begin{subproof}
        Fix $x$.
        Assume $x\in\dom{P}$.
        Thus $(x,\apply{P}{x})\in P$ by \cref{function_apply_elim}.
        Take $x_0\in\productFamily{X}$ such that $(x,\apply{P}{x}) = (x_0,\apply{x_0}{\beta})$ by \cref{projection_map_indexed}.
        Thus $x = x_0$ and $\apply{P}{x} = \apply{x}{\beta}$ by \cref{pair_eq_iff}.
        Thus $\apply{x}{\beta}\in\apply{X}{\beta}$ by \cref{indexed_product}.
        Thus $\apply{P}{x}\in\apply{X}{\beta}$.
    \end{subproof}
    Thus $P$ is a function to $\apply{X}{\beta}$.
    Thus $P\in\funs{\productFamily{X}}{\apply{X}{\beta}}$ by \cref{funs_intro}.
\end{proof}

% from §19 helper lemma 4: image of a finite set under a function is finite
\begin{lemma}\label{finite_image_function}
    Assume $F$ is finite.
    Assume $g$ is a function on $F$.
    Then $\img{g}{F}$ is finite.
\end{lemma}
\begin{proof}
    Let $B = \img{g}{F}$.
    Let $R = \{w\in B\times F \mid \exists y\in B.\exists x\in F. w = (y,x) \land \apply{g}{x} = y\}$.
    Show that $R$ is a relation.
    \begin{subproof}
        Show that for all $w$ such that $w\in R$ there exist $y,x$ such that $w = (y,x)$.
        \begin{subproof}
            Fix $w$.
            Assume $w\in R$.
            Take $y\in B$ such that there exists $x\in F$ such that $w = (y,x)$ and $\apply{g}{x} = y$.
            Take $x\in F$ such that $w = (y,x)$ and $\apply{g}{x} = y$.
        \end{subproof}
        Thus $R$ is a relation by \cref{relation}.
    \end{subproof}
    Show that for all $y\in B$ there exists $x$ such that $y\mathrel{R} x$.
    \begin{subproof}
        Fix $y$.
        Assume $y\in B$.
        Take $x\in F$ such that $(x,y)\in g$ by \cref{img_iff}.
        Thus $y = \apply{g}{x}$ by \cref{function_apply_intro}.
        Thus $(y,x)\in R$.
        Thus $y\mathrel{R} x$.
    \end{subproof}
    Take $s$ such that $s$ is a function on $B$ and for all $y\in B$ we have $y\mathrel{R}\apply{s}{y}$ by \cref{choice_function}.
    Thus $s$ is a function and $\dom{s} = B$.
    Show that for all $y\in B$ we have $\apply{s}{y}\in F$ and $\apply{g}{\apply{s}{y}} = y$.
    \begin{subproof}
        Fix $y$.
        Assume $y\in B$.
        Thus $(y,\apply{s}{y})\in R$.
        Take $y_0$ such that $y_0\in B$ and
            there exists $x_0\in F$ such that
            $(y,\apply{s}{y}) = (y_0,x_0)$ and $\apply{g}{x_0} = y_0$.
        Take $x_0\in F$ such that
            $(y,\apply{s}{y}) = (y_0,x_0)$ and $\apply{g}{x_0} = y_0$.
        Thus $y = y_0$ and $\apply{s}{y} = x_0$ by \cref{pair_eq_iff}.
        Thus $\apply{s}{y}\in F$.
        Thus $\apply{g}{\apply{s}{y}} = y$.
    \end{subproof}
    Let $C = \ran{s}$.
    Show that $s$ is injective.
    \begin{subproof}
        Show that for all $y_1,y_2$ such that $y_1\in B$ and $y_2\in B$ and $\apply{s}{y_1} = \apply{s}{y_2}$ we have $y_1 = y_2$.
        \begin{subproof}
            Fix $y_1$.
            Fix $y_2$.
            Assume $y_1\in B$ and $y_2\in B$ and $\apply{s}{y_1} = \apply{s}{y_2}$.
            Thus $\apply{g}{\apply{s}{y_1}} = y_1$ and $\apply{g}{\apply{s}{y_2}} = y_2$ by \cref{subseteq}.
            Thus $y_1 = y_2$.
        \end{subproof}
        Thus $s$ is injective.
    \end{subproof}
    Show that $C\subseteq F$.
    \begin{subproof}
        Show that for all $x$ such that $x\in C$ we have $x\in F$.
        \begin{subproof}
            Fix $x$.
            Assume $x\in C$.
            Take $y$ such that $(y,x)\in s$ by \cref{ran_iff}.
            Thus $y\in\dom{s}$ by \cref{dom_iff}.
            Thus $y\in B$.
            Thus $x = \apply{s}{y}$ by \cref{function_apply_intro}.
            Thus $\apply{s}{y}\in F$.
            Thus $x\in F$.
        \end{subproof}
        Thus $C\subseteq F$ by \cref{subseteq}.
    \end{subproof}
    Thus $C$ is finite by \cref{subseteq_finite_is_finite}.
    Show that $s\in\funs{B}{C}$.
    \begin{subproof}
        Show that for all $y\in B$ we have $\apply{s}{y}\in C$.
        \begin{subproof}
            Fix $y$.
            Assume $y\in B$.
            Thus $y\in\dom{s}$.
            Thus $(y,\apply{s}{y})\in s$ by \cref{function_apply_elim}.
            Thus $\apply{s}{y}\in\ran{s}$ by \cref{ran_intro}.
            Thus $\apply{s}{y}\in C$.
        \end{subproof}
        Thus $s$ is a function to $C$.
        Thus $s\in\funs{B}{C}$ by \cref{funs_intro}.
    \end{subproof}
    Thus $s\in\surj{B}{C}$ by \cref{funs_surj_iff}.
    Thus $s\in\bijections{B}{C}$ by \cref{bijections}.
    Thus $\converse{s}$ is a bijection from $C$ to $B$ by \cref{bijection_converse_is_bijection}.
    Thus $B$ is finite by \cref{finite_bijection}.
\end{proof}

% from §19 Item 21: continuity into an indexed product
\begin{lemma}\label{continuous_map_into_indexed_product}
    Assume $X$ is a function.
    Assume $T$ is a function.
    Assume $\dom{T} = \dom{X}$.
    Assume $\dom{X}$ is inhabited.
    Assume that for all $\alpha\in\dom{X}$ we have $\apply{T}{\alpha}$ is a topology on $\apply{X}{\alpha}$.
    Assume $T_A$ is a topology on $A$.
    Assume $f\in\funs{A}{\productFamily{X}}$.
    Then $f$ is continuous from $A$ to $\productFamily{X}$ with respect to $T_A$ and $\productTopologyIndexed{T}{X}$ iff
        for all $\alpha\in\dom{X}$ we have $\piIndex{X}{\alpha}\circ f$ is continuous from $A$ to $\apply{X}{\alpha}$ with respect to $T_A$ and $\apply{T}{\alpha}$.
\end{lemma}
\begin{proof}
    Let $T_0 = \productTopologyIndexed{T}{X}$.
    Let $S = \productSubbasis{T}{X}$.
    Show that for all $\alpha\in\dom{X}$ we have $\apply{T}{\alpha}$ is a topology on $\apply{X}{\alpha}$.
    \begin{subproof}
        Fix $\alpha$.
        Assume $\alpha\in\dom{X}$.
        Thus $\apply{T}{\alpha}$ is a topology on $\apply{X}{\alpha}$ by \cref{subseteq}.
    \end{subproof}
    Show that if $f$ is continuous from $A$ to $\productFamily{X}$ with respect to $T_A$ and $T_0$ then
        for all $\alpha\in\dom{X}$ we have $\piIndex{X}{\alpha}\circ f$ is continuous from $A$ to $\apply{X}{\alpha}$
        with respect to $T_A$ and $\apply{T}{\alpha}$.
    \begin{subproof}
        Assume $f$ is continuous from $A$ to $\productFamily{X}$ with respect to $T_A$ and $T_0$.
        Fix $\alpha$.
        Assume $\alpha\in\dom{X}$.
        Let $g = \piIndex{X}{\alpha}\circ f$.
        $\piIndex{X}{\alpha}\in\funs{\productFamily{X}}{\apply{X}{\alpha}}$ by \cref{projection_map_indexed_function}.
        Thus $g\in\funs{A}{\apply{X}{\alpha}}$ by \cref{funs_circ,continuous}.
        Show that for all $V$ such that $V$ is open in $\apply{X}{\alpha}$ with respect to $\apply{T}{\alpha}$
            we have $\preimg{g}{V}$ is open in $A$ with respect to $T_A$.
        \begin{subproof}
            Fix $V$.
            Assume $V$ is open in $\apply{X}{\alpha}$ with respect to $\apply{T}{\alpha}$.
            Let $W = \preimg{\piIndex{X}{\alpha}}{V}$.
            Show that $W\in\productSubbasisAt{T}{X}{\alpha}$.
            \begin{subproof}
                Show that $W\in\pow{\productFamily{X}}$.
                \begin{subproof}
                    $\piIndex{X}{\alpha}\in\funs{\productFamily{X}}{\apply{X}{\alpha}}$ by \cref{projection_map_indexed_function}.
                    Thus $\dom{\piIndex{X}{\alpha}} = \productFamily{X}$ by \cref{funs_elim}.
                    $\preimg{\piIndex{X}{\alpha}}{V}\subseteq\dom{\piIndex{X}{\alpha}}$ by \cref{preimg_subseteq_dom}.
                    Thus $W\subseteq\productFamily{X}$ by \cref{subseteq_transitive}.
                    Thus $W\in\pow{\productFamily{X}}$ by \cref{pow_iff}.
                \end{subproof}
                $V\in\apply{T}{\alpha}$ by \cref{open_in}.
                Thus $W\in\productSubbasisAt{T}{X}{\alpha}$ by \cref{product_subbasis_indexed}.
            \end{subproof}
            Show that $W\in\pow{\productFamily{X}}$.
            \begin{subproof}
                $\piIndex{X}{\alpha}\in\funs{\productFamily{X}}{\apply{X}{\alpha}}$ by \cref{projection_map_indexed_function}.
                Thus $\dom{\piIndex{X}{\alpha}} = \productFamily{X}$ by \cref{funs_elim}.
                $\preimg{\piIndex{X}{\alpha}}{V}\subseteq\dom{\piIndex{X}{\alpha}}$ by \cref{preimg_subseteq_dom}.
                Thus $W\subseteq\productFamily{X}$ by \cref{subseteq_transitive}.
                Thus $W\in\pow{\productFamily{X}}$ by \cref{pow_iff}.
            \end{subproof}
            Thus $W\in S$ by \cref{product_subbasis_union_indexed}.
            Let $F = \{W\}$.
            Let $G = \{W,W\}$.
            $G$ is finite by \cref{pair_is_finite}.
            Show that $F\subseteq G$.
            \begin{subproof}
                Show that for all $z$ such that $z\in F$ we have $z\in G$.
                \begin{subproof}
                    Fix $z$.
                    Assume $z\in F$.
                    Thus $z = W$ by \cref{singleton_elim}.
                    Thus $z\in G$ by \cref{upair_intro_left}.
                \end{subproof}
                Thus $F\subseteq G$ by \cref{subseteq}.
            \end{subproof}
            Thus $F$ is finite by \cref{subseteq_finite_is_finite}.
            $F$ is inhabited by \cref{singleton_intro}.
            $\inters{F} = W$ by \cref{inters_singleton}.
            Show that $F\subseteq S$.
            \begin{subproof}
                Show that for all $z$ such that $z\in F$ we have $z\in S$.
                \begin{subproof}
                    Fix $z$.
                    Assume $z\in F$.
                    Thus $z = W$ by \cref{singleton_elim}.
                    Thus $z\in S$.
                \end{subproof}
                Thus $F\subseteq S$ by \cref{subseteq}.
            \end{subproof}
            Show that $W\in\pow{\unions{S}}$.
            \begin{subproof}
                Show that for all $z$ such that $z\in W$ we have $z\in\unions{S}$.
                \begin{subproof}
                    Fix $z$.
                    Assume $z\in W$.
                    Thus $z\in\unions{S}$ by \cref{unions_iff}.
                \end{subproof}
                Thus $W\subseteq\unions{S}$ by \cref{subseteq}.
                Thus $W\in\pow{\unions{S}}$ by \cref{pow_iff}.
            \end{subproof}
            Thus $W\in\finiteIntersections{S}$ by \cref{finite_intersections}.
            $S$ is a subbasis on $\productFamily{X}$ by \cref{product_subbasis_is_subbasis}.
            Thus $\finiteIntersections{S}$ is a basis for $T_0$ on $\productFamily{X}$ by \cref{product_subbasis_finite_intersections_basis}.
            Thus $W\in T_0$ by \cref{basis_elem_open_in_generated,basis_for_topology}.
            Thus $T_0$ is a topology on $\productFamily{X}$ by \cref{basis_topology_is_topology,basis_for_topology}.
            Thus $W$ is open in $\productFamily{X}$ with respect to $T_0$ by \cref{open_in}.
            $\preimg{f}{W}$ is open in $A$ with respect to $T_A$ by \cref{continuous}.
            Show that $\preimg{g}{V} = \preimg{f}{W}$.
            \begin{subproof}
                Show that for all $a$ such that $a\in\preimg{g}{V}$ we have $a\in\preimg{f}{W}$.
                \begin{subproof}
                    Fix $a$.
                    Assume $a\in\preimg{g}{V}$.
                    Take $y\in V$ such that $a\mathrel{g} y$ by \cref{preimg_iff}.
                    Take $z$ such that $a\mathrel{f} z$ and $z\mathrel{\piIndex{X}{\alpha}} y$ by \cref{circ_iff}.
                    Thus $z\in W$ by \cref{preimg_iff}.
                    Thus $a\in\preimg{f}{W}$ by \cref{preimg_iff}.
                \end{subproof}
                Thus $\preimg{g}{V}\subseteq \preimg{f}{W}$ by \cref{subseteq}.
                Show that for all $a$ such that $a\in\preimg{f}{W}$ we have $a\in\preimg{g}{V}$.
                \begin{subproof}
                    Fix $a$.
                    Assume $a\in\preimg{f}{W}$.
                    Take $z\in W$ such that $a\mathrel{f} z$ by \cref{preimg_iff}.
                    Take $y\in V$ such that $z\mathrel{\piIndex{X}{\alpha}} y$ by \cref{preimg_iff}.
                    Thus $a\mathrel{g} y$ by \cref{circ_iff}.
                    Thus $a\in\preimg{g}{V}$ by \cref{preimg_iff}.
                \end{subproof}
                Thus $\preimg{f}{W}\subseteq \preimg{g}{V}$ by \cref{subseteq}.
                Thus $\preimg{g}{V} = \preimg{f}{W}$ by \cref{subseteq_antisymmetric}.
            \end{subproof}
            Thus $\preimg{g}{V}$ is open in $A$ with respect to $T_A$.
        \end{subproof}
        Thus $g$ is continuous from $A$ to $\apply{X}{\alpha}$ with respect to $T_A$ and $\apply{T}{\alpha}$ by \cref{continuous}.
    \end{subproof}
    Show that if for all $\alpha\in\dom{X}$ we have $\piIndex{X}{\alpha}\circ f$ is continuous from $A$ to $\apply{X}{\alpha}$
        with respect to $T_A$ and $\apply{T}{\alpha}$ then $f$ is continuous from $A$ to $\productFamily{X}$ with respect to $T_A$ and $T_0$.
    \begin{subproof}
        Assume for all $\alpha\in\dom{X}$ we have $\piIndex{X}{\alpha}\circ f$ is continuous from $A$ to $\apply{X}{\alpha}$
            with respect to $T_A$ and $\apply{T}{\alpha}$.
        Show that for all $U$ such that $U$ is open in $\productFamily{X}$ with respect to $T_0$
            we have $\preimg{f}{U}$ is open in $A$ with respect to $T_A$.
        \begin{subproof}
            Fix $U$.
            Assume $U$ is open in $\productFamily{X}$ with respect to $T_0$.
            Thus $U\in T_0$ by \cref{open_in}.
            $S$ is a subbasis on $\productFamily{X}$ by \cref{product_subbasis_is_subbasis}.
            Thus $\finiteIntersections{S}$ is a basis for $T_0$ on $\productFamily{X}$ by \cref{product_subbasis_finite_intersections_basis}.
            Take $M\subseteq \finiteIntersections{S}$ such that $U = \unions{M}$ by \cref{basis_topology_unions}.
            Let $N = \{\preimg{f}{B} \mid B\in M\}$.
            Show that for all $C$ such that $C\in N$ we have $C$ is open in $A$ with respect to $T_A$.
            \begin{subproof}
                Fix $C$.
                Assume $C\in N$.
                Take $B\in M$ such that $C = \preimg{f}{B}$.
                Thus $B\in\finiteIntersections{S}$ by \cref{subseteq}.
                Take $F\subseteq S$ such that $F$ is finite and inhabited and $B = \inters{F}$ by \cref{finite_intersections}.
                Let $H = \{\preimg{f}{A} \mid A\in F\}$.
                Show that for all $D$ such that $D\in H$ we have $D$ is open in $A$ with respect to $T_A$.
                \begin{subproof}
                    Fix $D$.
                    Assume $D\in H$.
                    Take $A\in F$ such that $D = \preimg{f}{A}$.
                    Thus $A\in S$ by \cref{subseteq}.
                    Take $\beta\in\dom{X}$ such that $A\in\productSubbasisAt{T}{X}{\beta}$ by \cref{product_subbasis_union_indexed}.
                    Take $V\in\apply{T}{\beta}$ such that $A = \preimg{\piIndex{X}{\beta}}{V}$ by \cref{product_subbasis_indexed}.
                    $\apply{T}{\beta}$ is a topology on $\apply{X}{\beta}$ by \cref{subseteq}.
                    Thus $V$ is open in $\apply{X}{\beta}$ with respect to $\apply{T}{\beta}$ by \cref{open_in}.
                    Let $g = \piIndex{X}{\beta}\circ f$.
                    Thus $g$ is continuous from $A$ to $\apply{X}{\beta}$ with respect to $T_A$ and $\apply{T}{\beta}$ by \cref{subseteq}.
                    Thus $\preimg{g}{V}$ is open in $A$ with respect to $T_A$ by \cref{continuous}.
                    Show that $\preimg{f}{A} = \preimg{g}{V}$.
                    \begin{subproof}
                        Show that for all $a$ such that $a\in\preimg{f}{A}$ we have $a\in\preimg{g}{V}$.
                        \begin{subproof}
                            Fix $a$.
                            Assume $a\in\preimg{f}{A}$.
                            Take $z\in A$ such that $a\mathrel{f} z$ by \cref{preimg_iff}.
                            Take $y\in V$ such that $z\mathrel{\piIndex{X}{\beta}} y$ by \cref{preimg_iff}.
                            Thus $a\mathrel{g} y$ by \cref{circ_iff}.
                            Thus $a\in\preimg{g}{V}$ by \cref{preimg_iff}.
                        \end{subproof}
                        Thus $\preimg{f}{A}\subseteq \preimg{g}{V}$ by \cref{subseteq}.
                        Show that for all $a$ such that $a\in\preimg{g}{V}$ we have $a\in\preimg{f}{A}$.
                        \begin{subproof}
                            Fix $a$.
                            Assume $a\in\preimg{g}{V}$.
                            Take $y\in V$ such that $a\mathrel{g} y$ by \cref{preimg_iff}.
                            Take $z$ such that $a\mathrel{f} z$ and $z\mathrel{\piIndex{X}{\beta}} y$ by \cref{circ_iff}.
                            Thus $z\in A$ by \cref{preimg_iff}.
                            Thus $a\in\preimg{f}{A}$ by \cref{preimg_iff}.
                        \end{subproof}
                        Thus $\preimg{g}{V}\subseteq \preimg{f}{A}$ by \cref{subseteq}.
                        Thus $\preimg{f}{A} = \preimg{g}{V}$ by \cref{subseteq_antisymmetric}.
                    \end{subproof}
                    Thus $D$ is open in $A$ with respect to $T_A$.
                \end{subproof}
                Show that $H$ is finite.
                \begin{subproof}
                    Let $h = \{(A,\preimg{f}{A}) \mid A\in F\}$.
                    Show that $h$ is a function on $F$.
                    \begin{subproof}
                        Show that $h$ is a relation.
                        \begin{subproof}
                            Show that for all $w$ such that $w\in h$ there exist $x,y$ such that $w = (x,y)$.
                            \begin{subproof}
                                Fix $w$.
                                Assume $w\in h$.
                                Take $A\in F$ such that $w = (A,\preimg{f}{A})$.
                            \end{subproof}
                            Thus $h$ is a relation by \cref{relation}.
                        \end{subproof}
                        Show that for all $A\in F$ there exists $y$ such that $A\mathrel{h} y$.
                        \begin{subproof}
                            Fix $A$.
                            Assume $A\in F$.
                            Thus $(A,\preimg{f}{A})\in h$.
                            Thus $A\mathrel{h}\preimg{f}{A}$.
                        \end{subproof}
                        Show that for all $x,y_1,y_2$ such that $(x,y_1)\in h$ and $(x,y_2)\in h$ we have $y_1 = y_2$.
                        \begin{subproof}
                            Fix $x$.
                            Fix $y_1$.
                            Fix $y_2$.
                            Assume $(x,y_1)\in h$ and $(x,y_2)\in h$.
                            Take $A\in F$ such that $(x,y_1) = (A,\preimg{f}{A})$.
                            Take $A'\in F$ such that $(x,y_2) = (A',\preimg{f}{A'})$.
                            Thus $x = A$ and $y_1 = \preimg{f}{A}$ by \cref{pair_eq_iff}.
                            Thus $x = A'$ and $y_2 = \preimg{f}{A'}$ by \cref{pair_eq_iff}.
                            Thus $A = A'$.
                            Thus $y_1 = y_2$.
                        \end{subproof}
                        Thus $h$ is right-unique.
                        Show that $\dom{h} = F$.
                        \begin{subproof}
                            Show that $\dom{h}\subseteq F$.
                            \begin{subproof}
                                Show that for all $x$ such that $x\in\dom{h}$ we have $x\in F$.
                                \begin{subproof}
                                    Fix $x$.
                                    Assume $x\in\dom{h}$.
                                    Take $y$ such that $(x,y)\in h$ by \cref{dom_iff}.
                                    Take $A\in F$ such that $(x,y) = (A,\preimg{f}{A})$.
                                    Thus $x = A$ by \cref{pair_eq_iff}.
                                    Thus $x\in F$.
                                \end{subproof}
                                Thus $\dom{h}\subseteq F$ by \cref{subseteq}.
                            \end{subproof}
                            Show that $F\subseteq\dom{h}$.
                            \begin{subproof}
                                Show that for all $A$ such that $A\in F$ we have $A\in\dom{h}$.
                                \begin{subproof}
                                    Fix $A$.
                                    Assume $A\in F$.
                                    Thus $(A,\preimg{f}{A})\in h$.
                                    Thus $A\in\dom{h}$ by \cref{dom_iff}.
                                \end{subproof}
                                Thus $F\subseteq\dom{h}$ by \cref{subseteq}.
                            \end{subproof}
                            Thus $\dom{h} = F$ by \cref{subseteq_antisymmetric}.
                        \end{subproof}
                        Thus $h$ is a function.
                        Thus $h$ is a function on $F$.
                    \end{subproof}
                    Thus $\img{h}{F}$ is finite by \cref{finite_image_function}.
                    Show that $H\subseteq\img{h}{F}$.
                    \begin{subproof}
                        Show that for all $D$ such that $D\in H$ we have $D\in\img{h}{F}$.
                        \begin{subproof}
                            Fix $D$.
                            Assume $D\in H$.
                            Take $A\in F$ such that $D = \preimg{f}{A}$.
                            Thus $(A,\preimg{f}{A})\in h$.
                            Thus $D\in\img{h}{F}$ by \cref{img_iff}.
                        \end{subproof}
                        Thus $H\subseteq\img{h}{F}$ by \cref{subseteq}.
                    \end{subproof}
                    Thus $H$ is finite by \cref{subseteq_finite_is_finite}.
                \end{subproof}
                Show that $H$ is inhabited.
                \begin{subproof}
                    Take $A$ such that $A\in F$.
                    Thus $\preimg{f}{A}\in H$.
                    Thus $H$ is inhabited.
                \end{subproof}
                Show that $H\subseteq T_A$.
                \begin{subproof}
                    Show that for all $D$ such that $D\in H$ we have $D\in T_A$.
                    \begin{subproof}
                        Fix $D$.
                        Assume $D\in H$.
                        Thus $D$ is open in $A$ with respect to $T_A$.
                        Thus $D\in T_A$ by \cref{open_in}.
                    \end{subproof}
                    Thus $H\subseteq T_A$ by \cref{subseteq}.
                \end{subproof}
                $T_A$ is a topology on $A$ by \cref{topology_on}.
                Thus $\inters{H}\in T_A$ by \cref{finite_intersections_open_in_topology}.
                Thus $\inters{H}$ is open in $A$ with respect to $T_A$ by \cref{open_in}.
                Show that $C = \inters{H}$.
                \begin{subproof}
                    Show that for all $a$ such that $a\in C$ we have $a\in\inters{H}$.
                    \begin{subproof}
                        Fix $a$.
                        Assume $a\in C$.
                        Thus $a\in\preimg{f}{B}$.
                        Take $z\in B$ such that $a\mathrel{f} z$ by \cref{preimg_iff}.
                        Thus $z\in\inters{F}$ by \cref{subseteq}.
                        Thus for all $A\in F$ we have $z\in A$ by \cref{inters_iff_forall}.
                        Show that for all $D\in H$ we have $a\in D$.
                        \begin{subproof}
                            Fix $D$.
                            Assume $D\in H$.
                            Take $A\in F$ such that $D = \preimg{f}{A}$.
                            Thus $a\in D$ by \cref{preimg_iff}.
                        \end{subproof}
                        Thus $a\in\inters{H}$ by \cref{inters_iff_forall}.
                    \end{subproof}
                    Thus $C\subseteq\inters{H}$ by \cref{subseteq}.
                    Show that for all $a$ such that $a\in\inters{H}$ we have $a\in C$.
                    \begin{subproof}
                        Fix $a$.
                        Assume $a\in\inters{H}$.
                        Show that for all $A\in F$ we have $a\in\preimg{f}{A}$.
                        \begin{subproof}
                            Fix $A$.
                            Assume $A\in F$.
                            Thus $\preimg{f}{A}\in H$ by \cref{img_iff}.
                            Thus $a\in\preimg{f}{A}$ by \cref{inters_iff_forall}.
                        \end{subproof}
                        Show that $a\in\preimg{f}{B}$.
                        \begin{subproof}
                            Take $y$ such that $a\mathrel{f} y$ by \cref{preimg_iff,preimg_subseteq_dom}.
                            Show that for all $A\in F$ we have $y\in A$.
                            \begin{subproof}
                                Fix $A$.
                                Assume $A\in F$.
                                Thus $a\in\preimg{f}{A}$.
                                Take $y_A\in A$ such that $a\mathrel{f} y_A$ by \cref{preimg_iff}.
                                Thus $f\in\funs{A}{\productFamily{X}}$ by \cref{continuous}.
                                Thus $f$ is a function by \cref{funs_is_function}.
                                Thus $f$ is right-unique by \cref{function_rightunique}.
                                Thus $y = y_A$ by \hyperref[rightunique]{right-uniqueness}.
                                Thus $y\in A$.
                            \end{subproof}
                            Thus $y\in\inters{F}$ by \cref{inters_iff_forall}.
                            Thus $y\in B$ by \cref{subseteq}.
                            Thus $a\in\preimg{f}{B}$ by \cref{preimg_iff}.
                        \end{subproof}
                        Thus $a\in C$.
                    \end{subproof}
                    Thus $\inters{H}\subseteq C$ by \cref{subseteq}.
                    Thus $C = \inters{H}$ by \cref{subseteq_antisymmetric}.
                \end{subproof}
                Thus $C$ is open in $A$ with respect to $T_A$.
            \end{subproof}
            Thus $N\subseteq T_A$ by \cref{subseteq,open_in}.
            $T_A$ is a topology on $A$ by \cref{continuous}.
            Thus $\unions{N}\in T_A$ by \cref{topology_on}.
            Show that $\preimg{f}{U} = \unions{N}$.
            \begin{subproof}
                Show that for all $a$ such that $a\in\preimg{f}{U}$ we have $a\in\unions{N}$.
                \begin{subproof}
                    Fix $a$.
                    Assume $a\in\preimg{f}{U}$.
                    Take $p\in U$ such that $a\mathrel{f} p$ by \cref{preimg_iff}.
                    Take $B\in M$ such that $p\in B$ by \cref{unions_iff}.
                    Thus $a\in\preimg{f}{B}$ by \cref{preimg_iff}.
                    Thus $a\in\unions{N}$ by \cref{unions_iff}.
                \end{subproof}
                Thus $\preimg{f}{U}\subseteq \unions{N}$ by \cref{subseteq}.
                Show that for all $a$ such that $a\in\unions{N}$ we have $a\in\preimg{f}{U}$.
                \begin{subproof}
                    Fix $a$.
                    Assume $a\in\unions{N}$.
                    Take $C\in N$ such that $a\in C$ by \cref{unions_iff}.
                    Take $B\in M$ such that $C = \preimg{f}{B}$.
                    Thus $a\in\preimg{f}{B}$.
                    Take $p\in B$ such that $a\mathrel{f} p$ by \cref{preimg_iff}.
                    Show that for all $q\in B$ we have $q\in U$.
                    \begin{subproof}
                        Fix $q\in B$.
                        Thus $q\in\unions{M}$ by \cref{unions_iff}.
                        Thus $q\in U$.
                    \end{subproof}
                    Thus $B\subseteq U$ by \cref{subseteq}.
                    Thus $p\in U$ by \cref{subseteq}.
                    Thus $a\in\preimg{f}{U}$ by \cref{preimg_iff}.
                \end{subproof}
                Thus $\unions{N}\subseteq \preimg{f}{U}$ by \cref{subseteq}.
                Thus $\preimg{f}{U} = \unions{N}$ by \cref{subseteq_antisymmetric}.
            \end{subproof}
            Thus $\preimg{f}{U}$ is open in $A$ with respect to $T_A$ by \cref{open_in}.
        \end{subproof}
        $S$ is a subbasis on $\productFamily{X}$ by \cref{product_subbasis_is_subbasis}.
        Thus $\finiteIntersections{S}$ is a basis for $T_0$ on $\productFamily{X}$ by \cref{product_subbasis_finite_intersections_basis}.
        Thus $T_0$ is a topology on $\productFamily{X}$ by \cref{basis_topology_is_topology,basis_for_topology}.
        Thus $f$ is continuous from $A$ to $\productFamily{X}$ with respect to $T_A$ and $T_0$ by \cref{continuous}.
    \end{subproof}
\end{proof}

% from §20 helper lemma 1: metric nonnegative
\begin{definition}\label{metric_nonnegative}
    $d$ is nonnegative on $X$ iff
    for all $x,y\in X$ we have $d((x,y)) \geq 0$.
\end{definition}

% from §20 helper lemma 2: metric zero characterization
\begin{definition}\label{metric_zero_characterization}
    $d$ is zero-characterized on $X$ iff
    for all $x,y\in X$ we have $d((x,y)) = 0$ iff $x = y$.
\end{definition}

% from §20 helper lemma 3: metric symmetry
\begin{definition}\label{metric_symmetric}
    $d$ is symmetric on $X$ iff
    for all $x,y\in X$ we have $d((x,y)) = d((y,x))$.
\end{definition}

% from §20 helper lemma 4: metric triangle inequality
\begin{definition}\label{metric_triangle_inequality}
    $d$ satisfies the triangle inequality on $X$ iff
    for all $x,y,z\in X$ we have $d((x,y)) + d((y,z)) \geq d((x,z))$.
\end{definition}

% from §20 Item 1: metric on a set
\begin{definition}\label{metric_on}
    $d$ is a metric on $X$ iff
    $d\in\funs{X\times X}{\reals}$ and
    $d$ is nonnegative on $X$ and
    $d$ is zero-characterized on $X$ and
    $d$ is symmetric on $X$ and
    $d$ satisfies the triangle inequality on $X$.
\end{definition}

% from §20 Item 2: metric ball
\begin{definition}\label{metric_ball}
    $\metricBall{d}{X}{x}{\epsilon} = \{y\in X \mid d((x,y)) < \epsilon\}$.
\end{definition}

% from §20 Item 3: metric basis
\begin{definition}\label{metric_basis}
    $\metricBasis{d}{X} = \{B\in\pow{X} \mid \exists x\in X. \exists \epsilon\in\reals. 0 < \epsilon \land B = \metricBall{d}{X}{x}{\epsilon}\}$.
\end{definition}

% from §20 Item 4: metric topology induced by a metric
\begin{definition}\label{metric_topology}
    $\metricTopology{d}{X} = \basisTopology{\metricBasis{d}{X}}{X}$.
\end{definition}

% from §20 Item 5: metric basis is a basis
\begin{lemma}\label{metric_basis_is_basis}
    Assume $d$ is a metric on $X$.
    Then $\metricBasis{d}{X}$ is a basis on $X$.
\end{lemma}
\begin{proof}
    Let $B = \metricBasis{d}{X}$.
    Show that for all $B_1$ such that $B_1\in B$ we have $B_1\in\pow{X}$.
    \begin{subproof}
        Fix $B_1$.
        Assume $B_1\in B$.
        Take $x,\epsilon$ such that $x\in X$ and $\epsilon\in\reals$ and $0 < \epsilon$ and $B_1 = \metricBall{d}{X}{x}{\epsilon}$
            by \cref{metric_basis}.
        Show that for all $y$ such that $y\in B_1$ we have $y\in X$.
        \begin{subproof}
            Fix $y$.
            Assume $y\in B_1$.
            Thus $y\in X$ by \cref{metric_ball}.
        \end{subproof}
        Thus $B_1\subseteq X$ by \cref{subseteq}.
        Thus $B_1\in\pow{X}$ by \cref{powerset_intro}.
    \end{subproof}
    Thus $B\subseteq\pow{X}$ by \cref{subseteq}.
    Show that for all $x$ such that $x\in X$ there exists $B_1\in B$ such that $x\in B_1$.
    \begin{subproof}
        Fix $x$.
        Assume $x\in X$.
        $1\in\reals$ by \cref{real_naturals_subset_reals,real_one_in_naturals,subseteq}.
        $0 < 1$ by \cref{real_zero_lt_one}.
        Show that $d((x,x)) = 0$.
        \begin{subproof}
            Show that for all $u,v$ such that $u\in X$ and $v\in X$ and $u = v$ we have $d((u,v)) = 0$.
            \begin{subproof}
                Fix $u$.
                Fix $v$.
                Assume $u\in X$ and $v\in X$ and $u = v$.
                $d$ is zero-characterized on $X$ by \cref{metric_on}.
                Thus $d((u,v)) = 0$ iff $u = v$ by \cref{metric_zero_characterization}.
                Thus $d((u,v)) = 0$.
            \end{subproof}
            Thus $d((x,x)) = 0$.
        \end{subproof}
        Thus $d((x,x)) < 1$ by \cref{real_lt_subst_left}.
        Let $B_1 = \metricBall{d}{X}{x}{1}$.
        Then $x\in B_1$ by \cref{metric_ball}.
        Show that for all $y$ such that $y\in B_1$ we have $y\in X$.
        \begin{subproof}
            Fix $y$.
            Assume $y\in B_1$.
            Thus $y\in X$ by \cref{metric_ball}.
        \end{subproof}
        Thus $B_1\subseteq X$ by \cref{subseteq}.
        Thus $B_1\in\pow{X}$ by \cref{powerset_intro}.
        Thus $B_1\in B$ by \cref{metric_basis}.
    \end{subproof}
    Show that for all $B_1,B_2,x$ such that $B_1\in B$ and $B_2\in B$ and $x\in B_1\inter B_2$
        there exists $B_3\in B$ such that $x\in B_3$ and $B_3\subseteq B_1\inter B_2$.
    \begin{subproof}
        Fix $B_1$.
        Fix $B_2$.
        Fix $x$.
        Assume $B_1\in B$ and $B_2\in B$ and $x\in B_1\inter B_2$.
        Take $x_1,\epsilon_1$ such that $x_1\in X$ and $\epsilon_1\in\reals$ and $0 < \epsilon_1$ and $B_1 = \metricBall{d}{X}{x_1}{\epsilon_1}$
            by \cref{metric_basis}.
        Take $x_2,\epsilon_2$ such that $x_2\in X$ and $\epsilon_2\in\reals$ and $0 < \epsilon_2$ and $B_2 = \metricBall{d}{X}{x_2}{\epsilon_2}$
            by \cref{metric_basis}.
        $x\in B_1$ and $x\in B_2$ by \cref{inter_elim_left,inter_elim_right}.
        Thus $x\in X$ and $d((x_1,x)) < \epsilon_1$ by \cref{metric_ball}.
        Thus $x\in X$ and $d((x_2,x)) < \epsilon_2$ by \cref{metric_ball}.
        $d\in\funs{X\times X}{\reals}$ by \cref{metric_on}.
        $(x_1,x)\in X\times X$ by \cref{times_tuple_intro}.
        $(x_2,x)\in X\times X$ by \cref{times_tuple_intro}.
        Thus $d((x_1,x))\in\reals$ by \cref{funs_type_apply}.
        Thus $d((x_2,x))\in\reals$ by \cref{funs_type_apply}.
        Let $r_1 = \epsilon_1 - d((x_1,x))$.
        Let $r_2 = \epsilon_2 - d((x_2,x))$.
        $0 < r_1$ by \cref{real_lt_imp_pos_diff}.
        $0 < r_2$ by \cref{real_lt_imp_pos_diff}.
        Let $A = \{r_1, r_2\}$.
        $A$ is finite by \cref{pair_is_finite}.
        $r_1\in A$ by \cref{upair_intro_left}.
        Thus $A\neq\emptyset$ by \cref{emptyset}.
        Show that for all $y$ such that $y\in A$ we have $y\in\reals$.
        \begin{subproof}
            Fix $y$.
            Assume $y\in A$.
            $y = r_1$ or $y = r_2$ by \cref{upair_elim}.
            Thus $y\in\reals$.
        \end{subproof}
        Thus $A\subseteq\reals$ by \cref{subseteq}.
        Take $m\in A$ such that $m$ is a minimum of $A$ by \cref{finite_real_minimum}.
        Show that $0 < m$.
        \begin{subproof}
            $m = r_1$ or $m = r_2$ by \cref{upair_elim}.
            \begin{byCase}
            \caseOf{$m = r_1$.} Thus $0 < m$.
            \caseOf{$m = r_2$.} Thus $0 < m$.
            \end{byCase}
        \end{subproof}
        Let $B_3 = \metricBall{d}{X}{x}{m}$.
        Show that $B_3\in B$.
        \begin{subproof}
            $m\in\reals$ by \cref{subseteq}.
            Show that for all $y$ such that $y\in B_3$ we have $y\in X$.
            \begin{subproof}
                Fix $y$.
                Assume $y\in B_3$.
                Thus $y\in X$ by \cref{metric_ball}.
            \end{subproof}
            Thus $B_3\subseteq X$ by \cref{subseteq}.
            Thus $B_3\in\pow{X}$ by \cref{powerset_intro}.
            Thus $B_3\in B$ by \cref{metric_basis}.
        \end{subproof}
        Show that $x\in B_3$.
        \begin{subproof}
            $d$ is zero-characterized on $X$ by \cref{metric_on}.
            Thus $d((x,x)) = 0$ iff $x = x$ by \cref{metric_zero_characterization}.
            Thus $d((x,x)) = 0$.
            Thus $d((x,x)) < m$ by \cref{real_lt_subst_left}.
        \end{subproof}
        Show that $B_3\subseteq B_1\inter B_2$.
        \begin{subproof}
            Show that for all $z$ such that $z\in B_3$ we have $z\in B_1\inter B_2$.
            \begin{subproof}
                Fix $z$.
                Assume $z\in B_3$.
                Thus $z\in X$ and $d((x,z)) < m$ by \cref{metric_ball}.
                $m$ is a minimum of $A$.
                Thus for all $t\in A$ we have $m \leq t$ by \cref{real_minimum}.
                $r_1\in A$ by \cref{upair_intro_left}.
                $r_2\in A$ by \cref{upair_intro_right}.
                Thus $m \leq r_1$ and $m \leq r_2$.
                Show that $z\in B_1$.
                \begin{subproof}
                    Show that $d((x,z)) < r_1$.
                    \begin{subproof}
                        $d((x,z)) < m$.
                        $(x,z)\in X\times X$ by \cref{times_tuple_intro}.
                        Thus $d((x,z))\in\reals$ by \cref{funs_type_apply}.
                        $m\in\reals$ by \cref{subseteq}.
                        $r_1 = \epsilon_1 - d((x_1,x))$.
                        Thus $r_1\in\reals$ by \cref{real_sub,real_add_closed}.
                        $m < r_1$ or $m = r_1$ by \cref{real_leq_split}.
                        \begin{byCase}
                        \caseOf{$m < r_1$.} Thus $d((x,z)) < r_1$ by \cref{real_lt_trans}.
                        \caseOf{$m = r_1$.} Thus $d((x,z)) < r_1$ by \cref{real_lt_subst_right}.
                        \end{byCase}
                    \end{subproof}
                    $r_1 = \epsilon_1 - d((x_1,x))$.
                    Thus $d((x,z)) < \epsilon_1 - d((x_1,x))$ by \cref{real_lt_subst_right}.
                    $(x,z)\in X\times X$ by \cref{times_tuple_intro}.
                    Thus $d((x,z))\in\reals$ by \cref{funs_type_apply}.
                    Thus $\epsilon_1 - d((x_1,x))\in\reals$ by \cref{real_sub,real_add_closed}.
                    $d((x,z)) + d((x_1,x)) < (\epsilon_1 - d((x_1,x))) + d((x_1,x))$ by \cref{real_lt_add}.
                    $\epsilon_1 - d((x_1,x)) = \epsilon_1 + \neg{(d((x_1,x)))}$ by \cref{real_sub}.
                    Thus $(\epsilon_1 - d((x_1,x))) + d((x_1,x)) = (\epsilon_1 + \neg{(d((x_1,x)))}) + d((x_1,x))$.
                    $\epsilon_1\in\reals$.
                    $d((x_1,x))\in\reals$.
                    $\neg{(d((x_1,x)))}\in\reals$ by \cref{real_add_inverse}.
                    $(\epsilon_1 + \neg{(d((x_1,x)))}) + d((x_1,x)) = \epsilon_1 + (\neg{(d((x_1,x)))} + d((x_1,x)))$ by \cref{real_add_assoc}.
                    $\neg{(d((x_1,x)))} + d((x_1,x)) = d((x_1,x)) + \neg{(d((x_1,x)))}$ by \cref{real_add_comm}.
                    $d((x_1,x)) + \neg{(d((x_1,x)))} = 0$ by \cref{real_add_inverse}.
                    Thus $(\epsilon_1 - d((x_1,x))) + d((x_1,x)) = \epsilon_1 + 0$.
                    $0 + \epsilon_1 = \epsilon_1$ by \cref{real_add_ident}.
                    $0\in\reals$ by \cref{real_add_ident}.
                    $\epsilon_1 + 0 = 0 + \epsilon_1$ by \cref{real_add_comm}.
                    Thus $\epsilon_1 + 0 = \epsilon_1$.
                    Thus $d((x,z)) + d((x_1,x)) < \epsilon_1$ by \cref{real_lt_subst_right}.
                    $d((x_1,x)) + d((x,z)) = d((x,z)) + d((x_1,x))$ by \cref{real_add_comm}.
                    Thus $d((x_1,x)) + d((x,z)) < \epsilon_1$ by \cref{real_lt_subst_left}.
                    $d$ satisfies the triangle inequality on $X$ by \cref{metric_on}.
                    Thus $d((x_1,x)) + d((x,z)) \geq d((x_1,z))$ by \cref{metric_triangle_inequality}.
                    Thus $d((x_1,z)) \leq d((x_1,x)) + d((x,z))$.
                    Show that $d((x_1,z)) < \epsilon_1$.
                    \begin{subproof}
                        $(x_1,z)\in X\times X$ by \cref{times_tuple_intro}.
                        Thus $d((x_1,z))\in\reals$ by \cref{funs_type_apply}.
                        $d((x_1,x))\in\reals$.
                        $d((x,z))\in\reals$.
                        Thus $d((x_1,x)) + d((x,z))\in\reals$ by \cref{real_add_closed}.
                        $d((x_1,z)) < d((x_1,x)) + d((x,z))$ or $d((x_1,z)) = d((x_1,x)) + d((x,z))$ by \cref{real_leq_split}.
                        \begin{byCase}
                        \caseOf{$d((x_1,z)) < d((x_1,x)) + d((x,z))$.}
                            Thus $d((x_1,z)) < \epsilon_1$ by \cref{real_lt_trans}.
                        \caseOf{$d((x_1,z)) = d((x_1,x)) + d((x,z))$.}
                            Thus $d((x_1,z)) < \epsilon_1$ by \cref{real_lt_subst_left}.
                        \end{byCase}
                    \end{subproof}
                    Thus $z\in B_1$ by \cref{metric_ball}.
                \end{subproof}
                Show that $z\in B_2$.
                \begin{subproof}
                    Show that $d((x,z)) < r_2$.
                    \begin{subproof}
                        $d((x,z)) < m$.
                        $(x,z)\in X\times X$ by \cref{times_tuple_intro}.
                        Thus $d((x,z))\in\reals$ by \cref{funs_type_apply}.
                        $m\in\reals$ by \cref{subseteq}.
                        $r_2 = \epsilon_2 - d((x_2,x))$.
                        Thus $r_2\in\reals$ by \cref{real_sub,real_add_closed}.
                        $m < r_2$ or $m = r_2$ by \cref{real_leq_split}.
                        \begin{byCase}
                        \caseOf{$m < r_2$.} Thus $d((x,z)) < r_2$ by \cref{real_lt_trans}.
                        \caseOf{$m = r_2$.} Thus $d((x,z)) < r_2$ by \cref{real_lt_subst_right}.
                        \end{byCase}
                    \end{subproof}
                    $r_2 = \epsilon_2 - d((x_2,x))$.
                    Thus $d((x,z)) < \epsilon_2 - d((x_2,x))$ by \cref{real_lt_subst_right}.
                    $(x,z)\in X\times X$ by \cref{times_tuple_intro}.
                    Thus $d((x,z))\in\reals$ by \cref{funs_type_apply}.
                    Thus $\epsilon_2 - d((x_2,x))\in\reals$ by \cref{real_sub,real_add_closed}.
                    $d((x,z)) + d((x_2,x)) < (\epsilon_2 - d((x_2,x))) + d((x_2,x))$ by \cref{real_lt_add}.
                    $\epsilon_2 - d((x_2,x)) = \epsilon_2 + \neg{(d((x_2,x)))}$ by \cref{real_sub}.
                    Thus $(\epsilon_2 - d((x_2,x))) + d((x_2,x)) = (\epsilon_2 + \neg{(d((x_2,x)))}) + d((x_2,x))$.
                    $\epsilon_2\in\reals$.
                    $d((x_2,x))\in\reals$.
                    $\neg{(d((x_2,x)))}\in\reals$ by \cref{real_add_inverse}.
                    $(\epsilon_2 + \neg{(d((x_2,x)))}) + d((x_2,x)) = \epsilon_2 + (\neg{(d((x_2,x)))} + d((x_2,x)))$ by \cref{real_add_assoc}.
                    $\neg{(d((x_2,x)))} + d((x_2,x)) = d((x_2,x)) + \neg{(d((x_2,x)))}$ by \cref{real_add_comm}.
                    $d((x_2,x)) + \neg{(d((x_2,x)))} = 0$ by \cref{real_add_inverse}.
                    Thus $(\epsilon_2 - d((x_2,x))) + d((x_2,x)) = \epsilon_2 + 0$.
                    $0 + \epsilon_2 = \epsilon_2$ by \cref{real_add_ident}.
                    $0\in\reals$ by \cref{real_add_ident}.
                    $\epsilon_2 + 0 = 0 + \epsilon_2$ by \cref{real_add_comm}.
                    Thus $\epsilon_2 + 0 = \epsilon_2$.
                    Thus $d((x,z)) + d((x_2,x)) < \epsilon_2$ by \cref{real_lt_subst_right}.
                    $d((x_2,x)) + d((x,z)) = d((x,z)) + d((x_2,x))$ by \cref{real_add_comm}.
                    Thus $d((x_2,x)) + d((x,z)) < \epsilon_2$ by \cref{real_lt_subst_left}.
                    $d$ satisfies the triangle inequality on $X$ by \cref{metric_on}.
                    Thus $d((x_2,x)) + d((x,z)) \geq d((x_2,z))$ by \cref{metric_triangle_inequality}.
                    Thus $d((x_2,z)) \leq d((x_2,x)) + d((x,z))$.
                    Show that $d((x_2,z)) < \epsilon_2$.
                    \begin{subproof}
                        $(x_2,z)\in X\times X$ by \cref{times_tuple_intro}.
                        Thus $d((x_2,z))\in\reals$ by \cref{funs_type_apply}.
                        $d((x_2,x))\in\reals$.
                        $d((x,z))\in\reals$.
                        Thus $d((x_2,x)) + d((x,z))\in\reals$ by \cref{real_add_closed}.
                        $d((x_2,z)) < d((x_2,x)) + d((x,z))$ or $d((x_2,z)) = d((x_2,x)) + d((x,z))$ by \cref{real_leq_split}.
                        \begin{byCase}
                        \caseOf{$d((x_2,z)) < d((x_2,x)) + d((x,z))$.}
                            Thus $d((x_2,z)) < \epsilon_2$ by \cref{real_lt_trans}.
                        \caseOf{$d((x_2,z)) = d((x_2,x)) + d((x,z))$.}
                            Thus $d((x_2,z)) < \epsilon_2$ by \cref{real_lt_subst_left}.
                        \end{byCase}
                    \end{subproof}
                    Thus $z\in B_2$ by \cref{metric_ball}.
                \end{subproof}
                Thus $z\in B_1\inter B_2$ by \cref{inter_intro}.
            \end{subproof}
            Thus $B_3\subseteq B_1\inter B_2$ by \cref{subseteq}.
        \end{subproof}
    \end{subproof}
    Thus $B$ is a basis on $X$ by \cref{basis_on}.
\end{proof}

% from §20 Item 6: metric open sets
\begin{lemma}\label{metric_open_iff}
    Assume $d$ is a metric on $X$.
    Let $T = \metricTopology{d}{X}$.
    Then for all $U$ such that $U\subseteq X$ we have
    $U$ is open in $X$ with respect to $T$ iff
    for all $y\in U$ there exists $\delta\in\reals$ such that $0 < \delta$ and $\metricBall{d}{X}{y}{\delta}\subseteq U$.
\end{lemma}
\begin{proof}
    Let $B = \metricBasis{d}{X}$.
    $T = \basisTopology{B}{X}$ by \cref{metric_topology}.
    Fix $U$.
    Assume $U\subseteq X$.
    Show that if $U$ is open in $X$ with respect to $T$ then
        for all $y\in U$ there exists $\delta\in\reals$ such that $0 < \delta$ and $\metricBall{d}{X}{y}{\delta}\subseteq U$.
    \begin{subproof}
        Assume $U$ is open in $X$ with respect to $T$.
        $T$ is a topology on $X$ and $U\in T$ by \cref{open_in}.
        Thus $U\in\basisTopology{B}{X}$.
        Show that for all $y\in U$ there exists $\delta\in\reals$ such that $0 < \delta$ and $\metricBall{d}{X}{y}{\delta}\subseteq U$.
        \begin{subproof}
            Fix $y$.
            Assume $y\in U$.
            There exists $B_1\in B$ such that $y\in B_1$ and $B_1\subseteq U$ by \cref{topology_generated_by_basis}.
            Take $x,\epsilon$ such that $x\in X$ and $\epsilon\in\reals$ and $0 < \epsilon$ and
                $B_1 = \metricBall{d}{X}{x}{\epsilon}$ by \cref{metric_basis}.
            Thus $y\in X$ and $d((x,y)) < \epsilon$ by \cref{metric_ball}.
            $d\in\funs{X\times X}{\reals}$ by \cref{metric_on}.
            $(x,y)\in X\times X$ by \cref{times_tuple_intro}.
            Thus $d((x,y))\in\reals$ by \cref{funs_type_apply}.
            Let $\delta = \epsilon - d((x,y))$.
            $0 < \delta$ by \cref{real_lt_imp_pos_diff}.
            Show that $\metricBall{d}{X}{y}{\delta}\subseteq U$.
            \begin{subproof}
                Show that $\metricBall{d}{X}{y}{\delta}\subseteq B_1$.
                \begin{subproof}
                    Show that for all $z$ such that $z\in\metricBall{d}{X}{y}{\delta}$ we have $z\in B_1$.
                    \begin{subproof}
                        Fix $z$.
                        Assume $z\in\metricBall{d}{X}{y}{\delta}$.
                        Thus $z\in X$ and $d((y,z)) < \delta$ by \cref{metric_ball}.
                        $(y,z)\in X\times X$ by \cref{times_tuple_intro}.
                        Thus $d((y,z))\in\reals$ by \cref{funs_type_apply}.
                        $\delta = \epsilon - d((x,y))$.
                        Thus $d((y,z)) < \epsilon - d((x,y))$ by \cref{real_lt_subst_right}.
                        $\neg{(d((x,y)))}\in\reals$ by \cref{real_add_inverse}.
                        $\epsilon - d((x,y)) = \epsilon + \neg{(d((x,y)))}$ by \cref{real_sub}.
                        Thus $\epsilon - d((x,y))\in\reals$ by \cref{real_add_closed}.
                        $d((y,z)) + d((x,y)) < (\epsilon - d((x,y))) + d((x,y))$ by \cref{real_lt_add}.
                        $\epsilon - d((x,y)) = \epsilon + \neg{(d((x,y)))}$ by \cref{real_sub}.
                        Thus $(\epsilon - d((x,y))) + d((x,y)) = (\epsilon + \neg{(d((x,y)))}) + d((x,y))$.
                        $\epsilon\in\reals$.
                        $d((x,y))\in\reals$.
                        $\neg{(d((x,y)))}\in\reals$ by \cref{real_add_inverse}.
                        $(\epsilon + \neg{(d((x,y)))}) + d((x,y)) = \epsilon + (\neg{(d((x,y)))} + d((x,y)))$ by \cref{real_add_assoc}.
                        $\neg{(d((x,y)))} + d((x,y)) = d((x,y)) + \neg{(d((x,y)))}$ by \cref{real_add_comm}.
                        $d((x,y)) + \neg{(d((x,y)))} = 0$ by \cref{real_add_inverse}.
                        Thus $(\epsilon - d((x,y))) + d((x,y)) = \epsilon + 0$.
                        $0 + \epsilon = \epsilon$ by \cref{real_add_ident}.
                        $0\in\reals$ by \cref{real_add_ident}.
                        $\epsilon + 0 = 0 + \epsilon$ by \cref{real_add_comm}.
                        Thus $\epsilon + 0 = \epsilon$.
                        Thus $d((y,z)) + d((x,y)) < \epsilon$ by \cref{real_lt_subst_right}.
                        $d((x,y)) + d((y,z)) = d((y,z)) + d((x,y))$ by \cref{real_add_comm}.
                        Thus $d((x,y)) + d((y,z)) < \epsilon$ by \cref{real_lt_subst_left}.
                        $d$ satisfies the triangle inequality on $X$ by \cref{metric_on}.
                        Thus $d((x,y)) + d((y,z)) \geq d((x,z))$ by \cref{metric_triangle_inequality}.
                        Thus $d((x,z)) \leq d((x,y)) + d((y,z))$.
                        Show that $d((x,z)) < \epsilon$.
                        \begin{subproof}
                            $(x,z)\in X\times X$ by \cref{times_tuple_intro}.
                            Thus $d((x,z))\in\reals$ by \cref{funs_type_apply}.
                            $d((x,y))\in\reals$.
                            $d((y,z))\in\reals$.
                            Thus $d((x,y)) + d((y,z))\in\reals$ by \cref{real_add_closed}.
                            $d((x,z)) < d((x,y)) + d((y,z))$ or $d((x,z)) = d((x,y)) + d((y,z))$ by \cref{real_leq_split}.
                            \begin{byCase}
                            \caseOf{$d((x,z)) < d((x,y)) + d((y,z))$.}
                                Thus $d((x,z)) < \epsilon$ by \cref{real_lt_trans}.
                            \caseOf{$d((x,z)) = d((x,y)) + d((y,z))$.}
                                Thus $d((x,z)) < \epsilon$ by \cref{real_lt_subst_left}.
                            \end{byCase}
                        \end{subproof}
                        Thus $z\in B_1$ by \cref{metric_ball}.
                    \end{subproof}
                    Thus $\metricBall{d}{X}{y}{\delta}\subseteq B_1$ by \cref{subseteq}.
                \end{subproof}
                Thus $\metricBall{d}{X}{y}{\delta}\subseteq U$ by \cref{subseteq_transitive}.
            \end{subproof}
        \end{subproof}
    \end{subproof}
    Show that if for all $y\in U$ there exists $\delta\in\reals$ such that $0 < \delta$ and $\metricBall{d}{X}{y}{\delta}\subseteq U$
        then $U$ is open in $X$ with respect to $T$.
    \begin{subproof}
        Assume for all $y\in U$ there exists $\delta\in\reals$ such that $0 < \delta$ and $\metricBall{d}{X}{y}{\delta}\subseteq U$.
        $B$ is a basis on $X$ by \cref{metric_basis_is_basis}.
        Thus $T$ is a topology on $X$ by \cref{basis_topology_is_topology}.
        Show that $U\in T$.
        \begin{subproof}
            Show that for all $y\in U$ there exists $B_1\in B$ such that $y\in B_1$ and $B_1\subseteq U$.
            \begin{subproof}
                Fix $y$.
                Assume $y\in U$.
                Thus $y\in X$ by \cref{subseteq}.
                Take $\delta\in\reals$ such that $0 < \delta$ and $\metricBall{d}{X}{y}{\delta}\subseteq U$.
                Let $B_1 = \metricBall{d}{X}{y}{\delta}$.
                Show that $y\in B_1$.
                \begin{subproof}
                    $d$ is zero-characterized on $X$ by \cref{metric_on}.
                    Thus $d((y,y)) = 0$ iff $y = y$ by \cref{metric_zero_characterization}.
                    Thus $d((y,y)) = 0$.
                    Thus $d((y,y)) < \delta$ by \cref{real_lt_subst_left}.
                    Thus $y\in B_1$ by \cref{metric_ball}.
                \end{subproof}
                Show that $B_1\in B$.
                \begin{subproof}
                    Show that for all $z$ such that $z\in B_1$ we have $z\in X$.
                    \begin{subproof}
                        Fix $z$.
                        Assume $z\in B_1$.
                        Thus $z\in X$ by \cref{metric_ball}.
                    \end{subproof}
                    Thus $B_1\subseteq X$ by \cref{subseteq}.
                    Thus $B_1\in\pow{X}$ by \cref{powerset_intro}.
                    Thus $B_1\in B$ by \cref{metric_basis}.
                \end{subproof}
            \end{subproof}
            Thus $U\in\pow{X}$ by \cref{powerset_intro}.
            Thus $U\in\basisTopology{B}{X}$ by \cref{topology_generated_by_basis}.
        \end{subproof}
        Thus $U$ is open in $X$ with respect to $T$ by \cref{open_in}.
    \end{subproof}
\end{proof}

% from §20 Item 7: discrete metric
\begin{definition}\label{discrete_metric}
    $\discreteMetric{X} = \{w \mid \exists p\in X\times X. \exists z\in\reals.
        w = (p,z) \land
        ((\fst{p} = \snd{p} \land z = 0) \lor (\fst{p}\neq \snd{p} \land z = 1))\}$.
\end{definition}

% from §20 Item 8: discrete metric is a metric
\begin{lemma}\label{discrete_metric_is_metric}
    Then $\discreteMetric{X}$ is a metric on $X$.
\end{lemma}
\begin{proof}
    Let $d = \discreteMetric{X}$.
    Show that $d\in\funs{X\times X}{\reals}$.
    \begin{subproof}
        Show that $d$ is a function to $\reals$.
        \begin{subproof}
            Show that $d$ is a function.
            \begin{subproof}
                Show that for all $p,z_1,z_2$ such that $(p,z_1)\in d$ and $(p,z_2)\in d$ we have $z_1 = z_2$.
                \begin{subproof}
                    Fix $p$.
                    Fix $z_1$.
                    Fix $z_2$.
                    Assume $(p,z_1)\in d$ and $(p,z_2)\in d$.
                    Take $p_1, w_1$ such that
                        $p_1\in X\times X$ and $w_1\in\reals$ and
                        $(p,z_1) = (p_1,w_1)$ and
                        $((\fst{p_1} = \snd{p_1} \land w_1 = 0) \lor (\fst{p_1}\neq \snd{p_1} \land w_1 = 1))$
                        by \cref{discrete_metric}.
                    Take $p_2, w_2$ such that
                        $p_2\in X\times X$ and $w_2\in\reals$ and
                        $(p,z_2) = (p_2,w_2)$ and
                        $((\fst{p_2} = \snd{p_2} \land w_2 = 0) \lor (\fst{p_2}\neq \snd{p_2} \land w_2 = 1))$
                        by \cref{discrete_metric}.
                    $p = p_1$ and $z_1 = w_1$ by \cref{pair_eq_iff}.
                    $p = p_2$ and $z_2 = w_2$ by \cref{pair_eq_iff}.
                    Thus $((\fst{p} = \snd{p} \land z_1 = 0) \lor (\fst{p}\neq \snd{p} \land z_1 = 1))$.
                    Thus $((\fst{p} = \snd{p} \land z_2 = 0) \lor (\fst{p}\neq \snd{p} \land z_2 = 1))$.
                    \begin{byCase}
                    \caseOf{$\fst{p} = \snd{p}$.}
                        Thus $z_1 = 0$ and $z_2 = 0$.
                        Thus $z_1 = z_2$.
                    \caseOf{$\fst{p} \neq \snd{p}$.}
                        Thus $z_1 = 1$ and $z_2 = 1$.
                        Thus $z_1 = z_2$.
                    \end{byCase}
                \end{subproof}
                Thus $d$ is right-unique by \cref{rightunique}.
                Show that for all $w$ such that $w\in d$ there exist $x,y$ such that $w = (x,y)$.
                \begin{subproof}
                    Fix $w$.
                    Assume $w\in d$.
                    Take $p, z$ such that
                        $p\in X\times X$ and $z\in\reals$ and
                        $w = (p,z)$ and
                        $((\fst{p} = \snd{p} \land z = 0) \lor (\fst{p}\neq \snd{p} \land z = 1))$
                        by \cref{discrete_metric}.
                    Thus there exist $x,y$ such that $w = (x,y)$.
                \end{subproof}
                Thus $d$ is a relation by \cref{relation}.
                Thus $d$ is a function.
            \end{subproof}
            Show that for all $p$ such that $p\in\dom{d}$ we have $d(p)\in\reals$.
            \begin{subproof}
                Fix $p$.
                Assume $p\in\dom{d}$.
                Take $z$ such that $(p,z)\in d$ by \cref{dom_iff}.
                Take $p_0, w$ such that
                    $p_0\in X\times X$ and $w\in\reals$ and
                    $(p,z) = (p_0,w)$ and
                    $((\fst{p_0} = \snd{p_0} \land w = 0) \lor (\fst{p_0}\neq \snd{p_0} \land w = 1))$
                    by \cref{discrete_metric}.
                $p = p_0$ and $z = w$ by \cref{pair_eq_iff}.
                Thus $z\in\reals$.
                $d$ is a function.
                Thus $d(p) = z$ by \cref{function_apply_intro}.
                Thus $d(p)\in\reals$.
            \end{subproof}
            Thus $d$ is a function to $\reals$.
        \end{subproof}
        Show that $\dom{d} = X\times X$.
        \begin{subproof}
            Show that for all $p$ such that $p\in\dom{d}$ we have $p\in X\times X$.
            \begin{subproof}
                Fix $p$.
                Assume $p\in\dom{d}$.
                Take $z$ such that $(p,z)\in d$ by \cref{dom_iff}.
                Take $p_0, w$ such that
                    $p_0\in X\times X$ and $w\in\reals$ and
                    $(p,z) = (p_0,w)$ and
                    $((\fst{p_0} = \snd{p_0} \land w = 0) \lor (\fst{p_0}\neq \snd{p_0} \land w = 1))$
                    by \cref{discrete_metric}.
                $p = p_0$ and $z = w$ by \cref{pair_eq_iff}.
                Thus $p\in X\times X$.
            \end{subproof}
            Thus $\dom{d}\subseteq X\times X$ by \cref{subseteq}.
            Show that for all $p$ such that $p\in X\times X$ we have $p\in\dom{d}$.
            \begin{subproof}
                Fix $p$.
                Assume $p\in X\times X$.
                \begin{byCase}
                \caseOf{$\fst{p} = \snd{p}$.}
                    $0\in\reals$ by \cref{real_add_ident}.
                    Let $w = (p,0)$.
                    $w\in d$ by \cref{discrete_metric}.
                    Thus $p\in\dom{d}$ by \cref{dom_intro}.
                \caseOf{$\fst{p} \neq \snd{p}$.}
                    $1\in\reals$ by \cref{real_naturals_subset_reals,real_one_in_naturals,subseteq}.
                    Let $w = (p,1)$.
                    $w\in d$ by \cref{discrete_metric}.
                    Thus $p\in\dom{d}$ by \cref{dom_intro}.
                \end{byCase}
            \end{subproof}
            Thus $X\times X\subseteq \dom{d}$ by \cref{subseteq}.
            Thus $\dom{d} = X\times X$ by \cref{subseteq_antisymmetric}.
        \end{subproof}
        Thus $d\in\funs{X\times X}{\reals}$ by \cref{funs_intro}.
    \end{subproof}
    Show that for all $x,y$ such that $x\in X$ and $y\in X$ we have $d((x,y)) \geq 0$.
    \begin{subproof}
        Fix $x$.
        Fix $y$.
        Assume $x\in X$ and $y\in X$.
        \begin{byCase}
        \caseOf{$x = y$.}
            $(x,y)\in X\times X$ by \cref{times_tuple_intro}.
            $0\in\reals$ by \cref{real_add_ident}.
            $\fst{(x,y)} = x$ and $\snd{(x,y)} = y$ by \cref{fst_eq,snd_eq}.
            Thus $\fst{(x,y)} = \snd{(x,y)}$.
            Thus $(\fst{(x,y)} = \snd{(x,y)} \land 0 = 0)$.
            Thus $((\fst{(x,y)} = \snd{(x,y)} \land 0 = 0) \lor
                (\fst{(x,y)} \neq \snd{(x,y)} \land 0 = 1))$.
            Let $w = ((x,y),0)$.
            $w\in d$ by \cref{discrete_metric}.
            $d$ is a function by \cref{funs_elim}.
            Thus $d((x,y)) = 0$ by \cref{function_apply_intro}.
            Thus $0 = d((x,y))$.
            Thus $0 \leq d((x,y))$.
            Thus $d((x,y)) \geq 0$.
        \caseOf{$x \neq y$.}
            $(x,y)\in X\times X$ by \cref{times_tuple_intro}.
            $1\in\reals$ by \cref{real_naturals_subset_reals,real_one_in_naturals,subseteq}.
            $\fst{(x,y)} = x$ and $\snd{(x,y)} = y$ by \cref{fst_eq,snd_eq}.
            Thus $\fst{(x,y)} \neq \snd{(x,y)}$.
            Thus $(\fst{(x,y)} \neq \snd{(x,y)} \land 1 = 1)$.
            Thus $((\fst{(x,y)} = \snd{(x,y)} \land 1 = 0) \lor
                (\fst{(x,y)} \neq \snd{(x,y)} \land 1 = 1))$.
            Let $w = ((x,y),1)$.
            $w\in d$ by \cref{discrete_metric}.
            $d$ is a function by \cref{funs_elim}.
            Thus $d((x,y)) = 1$ by \cref{function_apply_intro}.
            $0 < 1$ by \cref{real_zero_lt_one}.
            Thus $0 < d((x,y))$ by \cref{real_lt_subst_right}.
            Thus $0 \leq d((x,y))$.
            Thus $d((x,y)) \geq 0$.
        \end{byCase}
    \end{subproof}
    Thus $d$ is nonnegative on $X$ by \cref{metric_nonnegative}.
    Show that for all $x,y$ such that $x\in X$ and $y\in X$ we have $d((x,y)) = 0$ iff $x = y$.
    \begin{subproof}
        Fix $x$.
        Fix $y$.
        Assume $x\in X$ and $y\in X$.
        Show that if $d((x,y)) = 0$ then $x = y$.
        \begin{subproof}
            Assume $d((x,y)) = 0$.
            $d\in\funs{X\times X}{\reals}$.
            Thus $X\times X = \dom{d}$ by \cref{funs_elim}.
            $(x,y)\in X\times X$ by \cref{times_tuple_intro}.
            Thus $(x,y)\in\dom{d}$ by \cref{subseteq}.
            $d$ is a function by \cref{funs_elim}.
            Thus $((x,y),0)\in d$ by \cref{function_apply_iff}.
            Take $p, z$ such that
                $p\in X\times X$ and $z\in\reals$ and
                $((x,y),0) = (p,z)$ and
                $((\fst{p} = \snd{p} \land z = 0) \lor (\fst{p}\neq \snd{p} \land z = 1))$
                by \cref{discrete_metric}.
            $p = (x,y)$ and $z = 0$ by \cref{pair_eq_iff}.
            Thus $((\fst{p} = \snd{p} \land 0 = 0) \lor (\fst{p} \neq \snd{p} \land 0 = 1))$.
            Show that $\fst{p} = \snd{p}$.
            \begin{subproof}
                Suppose not.
                Thus $\fst{p} \neq \snd{p}$.
                Thus $0 = 1$.
                $0 < 1$ by \cref{real_zero_lt_one}.
                Thus $1 < 1$ by \cref{real_lt_subst_left}.
                Contradiction by \cref{real_lt_irreflexive}.
            \end{subproof}
            Thus $x = y$ by \cref{fst_eq,snd_eq}.
        \end{subproof}
        Show that if $x = y$ then $d((x,y)) = 0$.
        \begin{subproof}
            Assume $x = y$.
            $(x,y)\in X\times X$ by \cref{times_tuple_intro}.
            $0\in\reals$ by \cref{real_add_ident}.
            $\fst{(x,y)} = x$ and $\snd{(x,y)} = y$ by \cref{fst_eq,snd_eq}.
            Thus $\fst{(x,y)} = \snd{(x,y)}$.
            Thus $(\fst{(x,y)} = \snd{(x,y)} \land 0 = 0)$.
            Thus $((\fst{(x,y)} = \snd{(x,y)} \land 0 = 0) \lor
                (\fst{(x,y)} \neq \snd{(x,y)} \land 0 = 1))$.
            Let $w = ((x,y),0)$.
            $w\in d$ by \cref{discrete_metric}.
            $d$ is a function by \cref{funs_elim}.
            Thus $d((x,y)) = 0$ by \cref{function_apply_intro}.
        \end{subproof}
    \end{subproof}
    Thus $d$ is zero-characterized on $X$ by \cref{metric_zero_characterization}.
    Show that for all $x,y$ such that $x\in X$ and $y\in X$ we have $d((x,y)) = d((y,x))$.
    \begin{subproof}
        Fix $x$.
        Fix $y$.
        Assume $x\in X$ and $y\in X$.
        \begin{byCase}
        \caseOf{$x = y$.}
            $(x,y)\in X\times X$ by \cref{times_tuple_intro}.
            $(y,x)\in X\times X$ by \cref{times_tuple_intro}.
            $0\in\reals$ by \cref{real_add_ident}.
            $\fst{(x,y)} = x$ and $\snd{(x,y)} = y$ by \cref{fst_eq,snd_eq}.
            Thus $\fst{(x,y)} = \snd{(x,y)}$.
            Thus $(\fst{(x,y)} = \snd{(x,y)} \land 0 = 0)$.
            Thus $((\fst{(x,y)} = \snd{(x,y)} \land 0 = 0) \lor
                (\fst{(x,y)} \neq \snd{(x,y)} \land 0 = 1))$.
            Let $w_1 = ((x,y),0)$.
            $w_1\in d$ by \cref{discrete_metric}.
            $\fst{(y,x)} = y$ and $\snd{(y,x)} = x$ by \cref{fst_eq,snd_eq}.
            Thus $\fst{(y,x)} = \snd{(y,x)}$.
            Thus $(\fst{(y,x)} = \snd{(y,x)} \land 0 = 0)$.
            Thus $((\fst{(y,x)} = \snd{(y,x)} \land 0 = 0) \lor
                (\fst{(y,x)} \neq \snd{(y,x)} \land 0 = 1))$.
            Let $w_2 = ((y,x),0)$.
            $w_2\in d$ by \cref{discrete_metric}.
            $d$ is a function by \cref{funs_elim}.
            Thus $d((x,y)) = 0$ by \cref{function_apply_intro}.
            Thus $d((y,x)) = 0$ by \cref{function_apply_intro}.
            Thus $d((x,y)) = d((y,x))$.
        \caseOf{$x \neq y$.}
            $(x,y)\in X\times X$ by \cref{times_tuple_intro}.
            $(y,x)\in X\times X$ by \cref{times_tuple_intro}.
            $1\in\reals$ by \cref{real_naturals_subset_reals,real_one_in_naturals,subseteq}.
            $\fst{(x,y)} = x$ and $\snd{(x,y)} = y$ by \cref{fst_eq,snd_eq}.
            Thus $\fst{(x,y)} \neq \snd{(x,y)}$.
            Thus $(\fst{(x,y)} \neq \snd{(x,y)} \land 1 = 1)$.
            Thus $((\fst{(x,y)} = \snd{(x,y)} \land 1 = 0) \lor
                (\fst{(x,y)} \neq \snd{(x,y)} \land 1 = 1))$.
            Let $w_1 = ((x,y),1)$.
            $w_1\in d$ by \cref{discrete_metric}.
            $d$ is a function by \cref{funs_elim}.
            Thus $d((x,y)) = 1$ by \cref{function_apply_intro}.
            Show that $y \neq x$.
            \begin{subproof}
                Suppose not.
                Thus $y = x$.
                Thus $x = y$.
                Contradiction.
            \end{subproof}
            $\fst{(y,x)} = y$ and $\snd{(y,x)} = x$ by \cref{fst_eq,snd_eq}.
            Thus $\fst{(y,x)} \neq \snd{(y,x)}$.
            Thus $(\fst{(y,x)} \neq \snd{(y,x)} \land 1 = 1)$.
            Thus $((\fst{(y,x)} = \snd{(y,x)} \land 1 = 0) \lor
                (\fst{(y,x)} \neq \snd{(y,x)} \land 1 = 1))$.
            Let $w_2 = ((y,x),1)$.
            $w_2\in d$ by \cref{discrete_metric}.
            Thus $d((y,x)) = 1$ by \cref{function_apply_intro}.
            Thus $d((x,y)) = d((y,x))$.
        \end{byCase}
    \end{subproof}
    Thus $d$ is symmetric on $X$ by \cref{metric_symmetric}.
    Show that for all $x,y,z$ such that $x\in X$ and $y\in X$ and $z\in X$ we have
        $d((x,y)) + d((y,z)) \geq d((x,z))$.
    \begin{subproof}
        Fix $x$.
        Fix $y$.
        Fix $z$.
        Assume $x\in X$ and $y\in X$ and $z\in X$.
        \begin{byCase}
        \caseOf{$x = z$.}
            $(x,z)\in X\times X$ by \cref{times_tuple_intro}.
            $0\in\reals$ by \cref{real_add_ident}.
            $\fst{(x,z)} = x$ and $\snd{(x,z)} = z$ by \cref{fst_eq,snd_eq}.
            Thus $\fst{(x,z)} = \snd{(x,z)}$.
            Thus $(\fst{(x,z)} = \snd{(x,z)} \land 0 = 0)$.
            Thus $((\fst{(x,z)} = \snd{(x,z)} \land 0 = 0) \lor
                (\fst{(x,z)} \neq \snd{(x,z)} \land 0 = 1))$.
            Let $w = ((x,z),0)$.
            $w\in d$ by \cref{discrete_metric}.
            $d$ is a function by \cref{funs_elim}.
            Thus $d((x,z)) = 0$ by \cref{function_apply_intro}.
            $d\in\funs{X\times X}{\reals}$.
            Thus $d$ is a function to $\reals$ and $X\times X = \dom{d}$ by \cref{funs_elim}.
            $(x,y)\in X\times X$ and $(y,z)\in X\times X$ by \cref{times_tuple_intro}.
            Thus $(x,y)\in\dom{d}$ and $(y,z)\in\dom{d}$ by \cref{subseteq}.
            Thus $d((x,y))\in\reals$ and $d((y,z))\in\reals$.
            Thus $d((x,y)) + d((y,z))\in\reals$ by \cref{real_add_closed}.
            $d$ is nonnegative on $X$.
            Thus $d((x,y)) \geq 0$ by \cref{metric_nonnegative}.
            Thus $0 \leq d((x,y))$.
            Thus $d((y,z)) \geq 0$ by \cref{metric_nonnegative}.
            Thus $0 \leq d((y,z))$.
            Thus $0 + d((y,z)) \leq d((x,y)) + d((y,z))$ by \cref{real_leq_add}.
            $0 + d((y,z)) = d((y,z))$ by \cref{real_add_ident}.
            Thus $0 \leq d((x,y)) + d((y,z))$ by \cref{real_leq_trans}.
            Thus $d((x,z)) \leq 0$.
            Thus $d((x,z)) \leq d((x,y)) + d((y,z))$ by \cref{real_leq_trans}.
            Thus $d((x,y)) + d((y,z)) \geq d((x,z))$.
        \caseOf{$x \neq z$.}
            $(x,z)\in X\times X$ by \cref{times_tuple_intro}.
            $1\in\reals$ by \cref{real_naturals_subset_reals,real_one_in_naturals,subseteq}.
            $0\in\reals$ by \cref{real_add_ident}.
            $\fst{(x,z)} = x$ and $\snd{(x,z)} = z$ by \cref{fst_eq,snd_eq}.
            Thus $\fst{(x,z)} \neq \snd{(x,z)}$.
            Thus $(\fst{(x,z)} \neq \snd{(x,z)} \land 1 = 1)$.
            Thus $((\fst{(x,z)} = \snd{(x,z)} \land 1 = 0) \lor
                (\fst{(x,z)} \neq \snd{(x,z)} \land 1 = 1))$.
            Let $w = ((x,z),1)$.
            $w\in d$ by \cref{discrete_metric}.
            $d$ is a function by \cref{funs_elim}.
            Thus $d((x,z)) = 1$ by \cref{function_apply_intro}.
            Show that $d((x,y)) + d((y,z)) \geq 1$.
            \begin{subproof}
                \begin{byCase}
                \caseOf{$x = y$.}
                    Show that $y \neq z$.
                    \begin{subproof}
                        Suppose not.
                        Thus $y = z$.
                        Thus $x = z$.
                        Contradiction.
                    \end{subproof}
                    $(y,z)\in X\times X$ by \cref{times_tuple_intro}.
                    $\fst{(y,z)} = y$ and $\snd{(y,z)} = z$ by \cref{fst_eq,snd_eq}.
                    Thus $\fst{(y,z)} \neq \snd{(y,z)}$.
                    Thus $(\fst{(y,z)} \neq \snd{(y,z)} \land 1 = 1)$.
                    Thus $((\fst{(y,z)} = \snd{(y,z)} \land 1 = 0) \lor
                        (\fst{(y,z)} \neq \snd{(y,z)} \land 1 = 1))$.
                    Let $w_1 = ((y,z),1)$.
                    $w_1\in d$ by \cref{discrete_metric}.
                    $d$ is a function by \cref{funs_elim}.
                    Thus $d((y,z)) = 1$ by \cref{function_apply_intro}.
                    $d\in\funs{X\times X}{\reals}$.
                    Thus $d$ is a function to $\reals$ and $X\times X = \dom{d}$ by \cref{funs_elim}.
                    $(x,y)\in X\times X$ by \cref{times_tuple_intro}.
                    Thus $(x,y)\in\dom{d}$ by \cref{subseteq}.
                    Thus $d((x,y))\in\reals$.
                    $d$ is nonnegative on $X$.
                    Thus $d((x,y)) \geq 0$ by \cref{metric_nonnegative}.
                    Thus $0 \leq d((x,y))$.
                    Thus $0 + 1 \leq d((x,y)) + 1$ by \cref{real_leq_add}.
                    $0 + 1 = 1$ by \cref{real_add_ident}.
                    $d((x,y)) + 1 = 1 + d((x,y))$ by \cref{real_add_comm}.
                    Thus $1 \leq d((x,y)) + d((y,z))$ by \cref{real_leq_trans}.
                    Thus $d((x,y)) + d((y,z)) \geq 1$.
                \caseOf{$x \neq y$.}
                    $(x,y)\in X\times X$ by \cref{times_tuple_intro}.
                    $\fst{(x,y)} = x$ and $\snd{(x,y)} = y$ by \cref{fst_eq,snd_eq}.
                    Thus $\fst{(x,y)} \neq \snd{(x,y)}$.
                    Thus $(\fst{(x,y)} \neq \snd{(x,y)} \land 1 = 1)$.
                    Thus $((\fst{(x,y)} = \snd{(x,y)} \land 1 = 0) \lor
                        (\fst{(x,y)} \neq \snd{(x,y)} \land 1 = 1))$.
                    Let $w_2 = ((x,y),1)$.
                    $w_2\in d$ by \cref{discrete_metric}.
                    $d$ is a function by \cref{funs_elim}.
                    Thus $d((x,y)) = 1$ by \cref{function_apply_intro}.
                    $d\in\funs{X\times X}{\reals}$.
                    Thus $d$ is a function to $\reals$ and $X\times X = \dom{d}$ by \cref{funs_elim}.
                    $(y,z)\in X\times X$ by \cref{times_tuple_intro}.
                    Thus $(y,z)\in\dom{d}$ by \cref{subseteq}.
                    Thus $d((y,z))\in\reals$.
                    $d$ is nonnegative on $X$.
                    Thus $d((y,z)) \geq 0$ by \cref{metric_nonnegative}.
                    Thus $0 \leq d((y,z))$.
                    Thus $0 + 1 \leq d((y,z)) + 1$ by \cref{real_leq_add}.
                    $0 + 1 = 1$ by \cref{real_add_ident}.
                    $d((y,z)) + 1 = 1 + d((y,z))$ by \cref{real_add_comm}.
                    Thus $1 \leq d((x,y)) + d((y,z))$ by \cref{real_leq_trans}.
                    Thus $d((x,y)) + d((y,z)) \geq 1$.
                \end{byCase}
            \end{subproof}
            Thus $1 \leq d((x,y)) + d((y,z))$.
            Thus $d((x,z)) \leq 1$.
            Thus $d((x,z)) \leq d((x,y)) + d((y,z))$ by \cref{real_leq_trans}.
            Thus $d((x,y)) + d((y,z)) \geq d((x,z))$.
        \end{byCase}
    \end{subproof}
    Thus $d$ satisfies the triangle inequality on $X$ by \cref{metric_triangle_inequality}.
    Thus $d$ is a metric on $X$ by \cref{metric_on}.
\end{proof}

% from §20 Item 9: discrete metric topology
\begin{lemma}\label{discrete_metric_topology}
    Then $\metricTopology{\discreteMetric{X}}{X} = \discrete{X}$.
\end{lemma}
\begin{proof}
    Let $d = \discreteMetric{X}$.
    Let $T = \metricTopology{d}{X}$.
    $d$ is a metric on $X$ by \cref{discrete_metric_is_metric}.
    Thus $\metricBasis{d}{X}$ is a basis on $X$ by \cref{metric_basis_is_basis}.
    Thus $T$ is a topology on $X$ by \cref{basis_topology_is_topology,metric_topology}.
    Show that $T\subseteq\discrete{X}$.
    \begin{subproof}
        $T\subseteq\pow{X}$ by \cref{topology_on}.
        Thus $T\subseteq\discrete{X}$ by \cref{discrete_topology}.
    \end{subproof}
    Show that $\discrete{X}\subseteq T$.
    \begin{subproof}
        Show that for all $U$ such that $U\in\discrete{X}$ we have $U\in T$.
        \begin{subproof}
            Fix $U$.
            Assume $U\in\discrete{X}$.
            Thus $U\in\pow{X}$ by \cref{discrete_topology}.
            Thus $U\subseteq X$ by \cref{powerset_elim}.
            Show that $U$ is open in $X$ with respect to $T$.
            \begin{subproof}
                Show that for all $y\in U$ there exists $\delta\in\reals$ such that
                    $0 < \delta$ and $\metricBall{d}{X}{y}{\delta}\subseteq U$.
                \begin{subproof}
                    Fix $y$.
                    Assume $y\in U$.
                    Thus $y\in X$ by \cref{subseteq}.
                    $1\in\reals$ by \cref{real_naturals_subset_reals,real_one_in_naturals,subseteq}.
                    $0 < 1$ by \cref{real_zero_lt_one}.
                    Let $\delta = 1$.
                    Show that $\metricBall{d}{X}{y}{\delta}\subseteq U$.
                    \begin{subproof}
                        Show that for all $z$ such that $z\in\metricBall{d}{X}{y}{\delta}$ we have $z\in U$.
                        \begin{subproof}
                            Fix $z$.
                            Assume $z\in\metricBall{d}{X}{y}{\delta}$.
                            Thus $z\in X$ and $d((y,z)) < 1$ by \cref{metric_ball}.
                            Show that $z = y$.
                            \begin{subproof}
                                Suppose not.
                                Thus $z \neq y$.
                                $(y,z)\in X\times X$ by \cref{times_tuple_intro}.
                                $\fst{(y,z)} = y$ and $\snd{(y,z)} = z$ by \cref{fst_eq,snd_eq}.
                                Thus $\fst{(y,z)} \neq \snd{(y,z)}$.
                                Thus $(\fst{(y,z)} \neq \snd{(y,z)} \land 1 = 1)$.
                                Thus $((\fst{(y,z)} = \snd{(y,z)} \land 1 = 0) \lor
                                    (\fst{(y,z)} \neq \snd{(y,z)} \land 1 = 1))$.
                                Let $w = ((y,z),1)$.
                                $w\in d$ by \cref{discrete_metric}.
                                $d$ is a function by \cref{funs_elim,metric_on}.
                                Thus $d((y,z)) = 1$ by \cref{function_apply_intro}.
                                Thus $1 < 1$ by \cref{real_lt_subst_right}.
                                Contradiction by \cref{real_lt_irreflexive}.
                            \end{subproof}
                            Thus $z\in U$.
                        \end{subproof}
                        Thus $\metricBall{d}{X}{y}{\delta}\subseteq U$ by \cref{subseteq}.
                    \end{subproof}
                \end{subproof}
                Thus $U$ is open in $X$ with respect to $T$ by \cref{metric_open_iff}.
            \end{subproof}
            Thus $U\in T$ by \cref{open_in}.
        \end{subproof}
        Thus $\discrete{X}\subseteq T$ by \cref{subseteq}.
    \end{subproof}
    Thus $T = \discrete{X}$ by \cref{subseteq_antisymmetric}.
\end{proof}

% from §20 Item 10: standard metric on $\reals$
\begin{definition}\label{standard_metric_r}
    $\standardMetricR = \{w \mid \exists p\in\reals\times\reals. \exists z\in\reals.
        w = (p,z) \land z = \abs{\fst{p} - \snd{p}}\}$.
\end{definition}

% from §20 Item 11: standard metric is a metric
\begin{lemma}\label{standard_metric_is_metric}
    Then $\standardMetricR$ is a metric on $\reals$.
\end{lemma}
\begin{proof}
    Let $d = \standardMetricR$.
    Show that $d\in\funs{\reals\times\reals}{\reals}$.
    \begin{subproof}
        Show that $d$ is a function to $\reals$.
        \begin{subproof}
            Show that $d$ is a function.
            \begin{subproof}
                Show that for all $p,z_1,z_2$ such that $(p,z_1)\in d$ and $(p,z_2)\in d$ we have $z_1 = z_2$.
                \begin{subproof}
                    Fix $p$.
                    Fix $z_1$.
                    Fix $z_2$.
                    Assume $(p,z_1)\in d$ and $(p,z_2)\in d$.
                    Take $p_1, w_1$ such that
                        $p_1\in\reals\times\reals$ and $w_1\in\reals$ and
                        $(p,z_1) = (p_1,w_1)$ and $w_1 = \abs{\fst{p_1} - \snd{p_1}}$
                        by \cref{standard_metric_r}.
                    Take $p_2, w_2$ such that
                        $p_2\in\reals\times\reals$ and $w_2\in\reals$ and
                        $(p,z_2) = (p_2,w_2)$ and $w_2 = \abs{\fst{p_2} - \snd{p_2}}$
                        by \cref{standard_metric_r}.
                    $p = p_1$ and $z_1 = w_1$ by \cref{pair_eq_iff}.
                    $p = p_2$ and $z_2 = w_2$ by \cref{pair_eq_iff}.
                    Thus $z_1 = \abs{\fst{p} - \snd{p}}$.
                    Thus $z_2 = \abs{\fst{p} - \snd{p}}$.
                    Thus $z_1 = z_2$.
                \end{subproof}
                Thus $d$ is right-unique by \cref{rightunique}.
                Show that for all $w$ such that $w\in d$ there exist $x,y$ such that $w = (x,y)$.
                \begin{subproof}
                    Fix $w$.
                    Assume $w\in d$.
                    Take $p, z$ such that
                        $p\in\reals\times\reals$ and $z\in\reals$ and
                        $w = (p,z)$ and $z = \abs{\fst{p} - \snd{p}}$
                        by \cref{standard_metric_r}.
                    Thus there exist $x,y$ such that $w = (x,y)$.
                \end{subproof}
                Thus $d$ is a relation by \cref{relation}.
                Thus $d$ is a function.
            \end{subproof}
            Show that for all $p$ such that $p\in\dom{d}$ we have $d(p)\in\reals$.
            \begin{subproof}
                Fix $p$.
                Assume $p\in\dom{d}$.
                Take $z$ such that $(p,z)\in d$ by \cref{dom_iff}.
                Take $p_0, w$ such that
                    $p_0\in\reals\times\reals$ and $w\in\reals$ and
                    $(p,z) = (p_0,w)$ and $w = \abs{\fst{p_0} - \snd{p_0}}$
                    by \cref{standard_metric_r}.
                $p = p_0$ and $z = w$ by \cref{pair_eq_iff}.
                Thus $z\in\reals$.
                $d$ is a function.
                Thus $d(p) = z$ by \cref{function_apply_intro}.
                Thus $d(p)\in\reals$.
            \end{subproof}
            Thus $d$ is a function to $\reals$.
        \end{subproof}
        Show that $\dom{d} = \reals\times\reals$.
        \begin{subproof}
            Show that for all $p$ such that $p\in\dom{d}$ we have $p\in\reals\times\reals$.
            \begin{subproof}
                Fix $p$.
                Assume $p\in\dom{d}$.
                Take $z$ such that $(p,z)\in d$ by \cref{dom_iff}.
                Take $p_0, w$ such that
                    $p_0\in\reals\times\reals$ and $w\in\reals$ and
                    $(p,z) = (p_0,w)$ and $w = \abs{\fst{p_0} - \snd{p_0}}$
                    by \cref{standard_metric_r}.
                $p = p_0$ and $z = w$ by \cref{pair_eq_iff}.
                Thus $p\in\reals\times\reals$.
            \end{subproof}
            Thus $\dom{d}\subseteq\reals\times\reals$ by \cref{subseteq}.
            Show that for all $p$ such that $p\in\reals\times\reals$ we have $p\in\dom{d}$.
            \begin{subproof}
                Fix $p$.
                Assume $p\in\reals\times\reals$.
                Take $x,y$ such that $x\in\reals$ and $y\in\reals$ and $p = (x,y)$ by \cref{times_elem_is_tuple}.
                $\neg{y}\in\reals$ by \cref{real_add_inverse}.
                $x - y = x + \neg{y}$ by \cref{real_sub}.
                Thus $x - y\in\reals$ by \cref{real_add_closed}.
                Thus $\abs{(x - y)}\in\reals$ by \cref{real_abs_in_reals}.
                $\fst{p} = x$ and $\snd{p} = y$ by \cref{fst_eq,snd_eq}.
                Thus $\fst{p} - \snd{p} = x - y$.
                Thus $\fst{p} - \snd{p}\in\reals$.
                Let $z = \abs{\fst{p} - \snd{p}}$.
                Thus $z\in\reals$ by \cref{real_abs_in_reals}.
                Thus $z = \abs{(x - y)}$.
                Let $w = (p,z)$.
                $w\in d$ by \cref{standard_metric_r}.
                Thus $p\in\dom{d}$ by \cref{dom_intro}.
            \end{subproof}
            Thus $\reals\times\reals\subseteq \dom{d}$ by \cref{subseteq}.
            Thus $\dom{d} = \reals\times\reals$ by \cref{subseteq_antisymmetric}.
        \end{subproof}
        Thus $d\in\funs{\reals\times\reals}{\reals}$ by \cref{funs_intro}.
    \end{subproof}
    Show that for all $x,y$ such that $x\in\reals$ and $y\in\reals$ we have $d((x,y)) \geq 0$.
    \begin{subproof}
        Fix $x$.
        Fix $y$.
        Assume $x\in\reals$ and $y\in\reals$.
        $\neg{y}\in\reals$ by \cref{real_add_inverse}.
        $x - y = x + \neg{y}$ by \cref{real_sub}.
        Thus $x - y\in\reals$ by \cref{real_add_closed}.
        Thus $\abs{(x - y)}\in\reals$ by \cref{real_abs_in_reals}.
        Let $p = (x,y)$.
        $p\in\reals\times\reals$ by \cref{times_tuple_intro}.
        $\fst{p} = x$ and $\snd{p} = y$ by \cref{fst_eq,snd_eq}.
        Thus $\fst{p} - \snd{p} = x - y$.
        Thus $\fst{p} - \snd{p}\in\reals$.
        Let $z = \abs{\fst{p} - \snd{p}}$.
        Thus $z\in\reals$ by \cref{real_abs_in_reals}.
        Thus $z = \abs{(x - y)}$.
        Let $w = (p,z)$.
        $w\in d$ by \cref{standard_metric_r}.
        $d$ is a function by \cref{funs_elim}.
        Thus $d((x,y)) = z$ by \cref{function_apply_intro}.
        Thus $d((x,y)) = \abs{(x - y)}$.
        Thus $d((x,y)) \geq 0$ by \cref{real_abs_nonnegative}.
    \end{subproof}
    Thus $d$ is nonnegative on $\reals$ by \cref{metric_nonnegative}.
    Show that for all $x,y$ such that $x\in\reals$ and $y\in\reals$ we have $d((x,y)) = 0$ iff $x = y$.
    \begin{subproof}
        Fix $x$.
        Fix $y$.
        Assume $x\in\reals$ and $y\in\reals$.
        Show that if $d((x,y)) = 0$ then $x = y$.
        \begin{subproof}
            Assume $d((x,y)) = 0$.
            $(x,y)\in\reals\times\reals$ by \cref{times_tuple_intro}.
            $\neg{y}\in\reals$ by \cref{real_add_inverse}.
            $x - y = x + \neg{y}$ by \cref{real_sub}.
            Thus $x - y\in\reals$ by \cref{real_add_closed}.
            Thus $\abs{(x - y)}\in\reals$ by \cref{real_abs_in_reals}.
            Let $p = (x,y)$.
            $p\in\reals\times\reals$ by \cref{times_tuple_intro}.
            $\fst{p} = x$ and $\snd{p} = y$ by \cref{fst_eq,snd_eq}.
            Thus $\fst{p} - \snd{p} = x - y$.
            Thus $\fst{p} - \snd{p}\in\reals$.
            Let $z = \abs{\fst{p} - \snd{p}}$.
            Thus $z\in\reals$ by \cref{real_abs_in_reals}.
            Thus $z = \abs{(x - y)}$.
            Let $w = (p,z)$.
            $w\in d$ by \cref{standard_metric_r}.
            $d$ is a function by \cref{funs_elim}.
            Thus $d((x,y)) = z$ by \cref{function_apply_intro}.
            Thus $d((x,y)) = \abs{(x - y)}$.
            Thus $\abs{(x - y)} = 0$.
            Thus $x - y = 0$ by \cref{real_abs_eq_zero}.
            $x - y = x + \neg{y}$ by \cref{real_sub}.
            $y + \neg{y} = 0$ by \cref{real_add_inverse}.
            Thus $x + \neg{y} = y + \neg{y}$.
            Thus $x = y$ by \cref{real_add_cancel}.
        \end{subproof}
        Show that if $x = y$ then $d((x,y)) = 0$.
        \begin{subproof}
            Assume $x = y$.
            $(x,y)\in\reals\times\reals$ by \cref{times_tuple_intro}.
            $x - y = x + \neg{y}$ by \cref{real_sub}.
            $y + \neg{y} = 0$ by \cref{real_add_inverse}.
            Thus $x - y = 0$.
            $0\in\reals$ by \cref{real_add_ident}.
            Thus $0 \geq 0$.
            Thus $\abs{(x - y)} = 0$ by \cref{real_abs_nonneg}.
            Let $p = (x,y)$.
            $p\in\reals\times\reals$ by \cref{times_tuple_intro}.
            $\fst{p} = x$ and $\snd{p} = y$ by \cref{fst_eq,snd_eq}.
            Thus $\fst{p} - \snd{p} = x - y$.
            Thus $\fst{p} - \snd{p}\in\reals$.
            Let $z = \abs{\fst{p} - \snd{p}}$.
            Thus $z\in\reals$ by \cref{real_abs_in_reals}.
            Thus $z = \abs{(x - y)}$.
            Let $w = (p,z)$.
            $w\in d$ by \cref{standard_metric_r}.
            $d$ is a function by \cref{funs_elim}.
            Thus $d((x,y)) = z$ by \cref{function_apply_intro}.
            Thus $d((x,y)) = \abs{(x - y)}$.
            Thus $d((x,y)) = 0$.
        \end{subproof}
    \end{subproof}
    Thus $d$ is zero-characterized on $\reals$ by \cref{metric_zero_characterization}.
    Show that for all $x,y$ such that $x\in\reals$ and $y\in\reals$ we have $d((x,y)) = d((y,x))$.
    \begin{subproof}
        Fix $x$.
        Fix $y$.
        Assume $x\in\reals$ and $y\in\reals$.
        $(x,y)\in\reals\times\reals$ by \cref{times_tuple_intro}.
        $(y,x)\in\reals\times\reals$ by \cref{times_tuple_intro}.
        $\neg{y}\in\reals$ by \cref{real_add_inverse}.
        $x - y = x + \neg{y}$ by \cref{real_sub}.
        Thus $x - y\in\reals$ by \cref{real_add_closed}.
        $\neg{x}\in\reals$ by \cref{real_add_inverse}.
        $y - x = y + \neg{x}$ by \cref{real_sub}.
        Thus $y - x\in\reals$ by \cref{real_add_closed}.
        Thus $\abs{(x - y)}\in\reals$ by \cref{real_abs_in_reals}.
        Thus $\abs{(y - x)}\in\reals$ by \cref{real_abs_in_reals}.
        Let $p_1 = (x,y)$.
        Let $p_2 = (y,x)$.
        $p_1\in\reals\times\reals$ by \cref{times_tuple_intro}.
        $p_2\in\reals\times\reals$ by \cref{times_tuple_intro}.
        $\fst{p_1} = x$ and $\snd{p_1} = y$ by \cref{fst_eq,snd_eq}.
        $\fst{p_2} = y$ and $\snd{p_2} = x$ by \cref{fst_eq,snd_eq}.
        Thus $\fst{p_1} - \snd{p_1} = x - y$.
        Thus $\fst{p_2} - \snd{p_2} = y - x$.
        Thus $\fst{p_1} - \snd{p_1}\in\reals$.
        Thus $\fst{p_2} - \snd{p_2}\in\reals$.
        Let $z_1 = \abs{\fst{p_1} - \snd{p_1}}$.
        Let $z_2 = \abs{\fst{p_2} - \snd{p_2}}$.
        Thus $z_1\in\reals$ by \cref{real_abs_in_reals}.
        Thus $z_2\in\reals$ by \cref{real_abs_in_reals}.
        Thus $z_1 = \abs{(x - y)}$.
        Thus $z_2 = \abs{(y - x)}$.
        Let $w_1 = (p_1,z_1)$.
        Let $w_2 = (p_2,z_2)$.
        $w_1\in d$ by \cref{standard_metric_r}.
        $w_2\in d$ by \cref{standard_metric_r}.
        $d$ is a function by \cref{funs_elim}.
        Thus $d((x,y)) = z_1$ by \cref{function_apply_intro}.
        Thus $d((y,x)) = z_2$ by \cref{function_apply_intro}.
        Thus $d((x,y)) = \abs{(x - y)}$.
        Thus $d((y,x)) = \abs{(y - x)}$.
        Thus $d((x,y)) = d((y,x))$ by \cref{real_abs_sub_swap}.
    \end{subproof}
    Thus $d$ is symmetric on $\reals$ by \cref{metric_symmetric}.
    Show that for all $x,y,z$ such that $x\in\reals$ and $y\in\reals$ and $z\in\reals$ we have
        $d((x,y)) + d((y,z)) \geq d((x,z))$.
    \begin{subproof}
        Fix $x$.
        Fix $y$.
        Fix $z$.
        Assume $x\in\reals$ and $y\in\reals$ and $z\in\reals$.
        $(x,y)\in\reals\times\reals$ by \cref{times_tuple_intro}.
        $(y,z)\in\reals\times\reals$ by \cref{times_tuple_intro}.
        $(x,z)\in\reals\times\reals$ by \cref{times_tuple_intro}.
        $\neg{y}\in\reals$ by \cref{real_add_inverse}.
        $x - y = x + \neg{y}$ by \cref{real_sub}.
        Thus $x - y\in\reals$ by \cref{real_add_closed}.
        $\neg{z}\in\reals$ by \cref{real_add_inverse}.
        $y - z = y + \neg{z}$ by \cref{real_sub}.
        Thus $y - z\in\reals$ by \cref{real_add_closed}.
        $x - z = x + \neg{z}$ by \cref{real_sub}.
        Thus $x - z\in\reals$ by \cref{real_add_closed}.
        Thus $\abs{(x - y)}\in\reals$ by \cref{real_abs_in_reals}.
        Thus $\abs{(y - z)}\in\reals$ by \cref{real_abs_in_reals}.
        Thus $\abs{(x - z)}\in\reals$ by \cref{real_abs_in_reals}.
        Let $p_1 = (x,y)$.
        Let $p_2 = (y,z)$.
        Let $p_3 = (x,z)$.
        $p_1\in\reals\times\reals$ by \cref{times_tuple_intro}.
        $p_2\in\reals\times\reals$ by \cref{times_tuple_intro}.
        $p_3\in\reals\times\reals$ by \cref{times_tuple_intro}.
        $\fst{p_1} = x$ and $\snd{p_1} = y$ by \cref{fst_eq,snd_eq}.
        $\fst{p_2} = y$ and $\snd{p_2} = z$ by \cref{fst_eq,snd_eq}.
        $\fst{p_3} = x$ and $\snd{p_3} = z$ by \cref{fst_eq,snd_eq}.
        Thus $\fst{p_1} - \snd{p_1} = x - y$.
        Thus $\fst{p_2} - \snd{p_2} = y - z$.
        Thus $\fst{p_3} - \snd{p_3} = x - z$.
        Thus $\fst{p_1} - \snd{p_1}\in\reals$.
        Thus $\fst{p_2} - \snd{p_2}\in\reals$.
        Thus $\fst{p_3} - \snd{p_3}\in\reals$.
        Let $z_1 = \abs{\fst{p_1} - \snd{p_1}}$.
        Let $z_2 = \abs{\fst{p_2} - \snd{p_2}}$.
        Let $z_3 = \abs{\fst{p_3} - \snd{p_3}}$.
        Thus $z_1\in\reals$ by \cref{real_abs_in_reals}.
        Thus $z_2\in\reals$ by \cref{real_abs_in_reals}.
        Thus $z_3\in\reals$ by \cref{real_abs_in_reals}.
        Thus $z_1 = \abs{(x - y)}$.
        Thus $z_2 = \abs{(y - z)}$.
        Thus $z_3 = \abs{(x - z)}$.
        Let $w_1 = (p_1,z_1)$.
        Let $w_2 = (p_2,z_2)$.
        Let $w_3 = (p_3,z_3)$.
        $w_1\in d$ by \cref{standard_metric_r}.
        $w_2\in d$ by \cref{standard_metric_r}.
        $w_3\in d$ by \cref{standard_metric_r}.
        $d$ is a function by \cref{funs_elim}.
        Thus $d((x,y)) = z_1$ by \cref{function_apply_intro}.
        Thus $d((y,z)) = z_2$ by \cref{function_apply_intro}.
        Thus $d((x,z)) = z_3$ by \cref{function_apply_intro}.
        Thus $d((x,y)) = \abs{(x - y)}$.
        Thus $d((y,z)) = \abs{(y - z)}$.
        Thus $d((x,z)) = \abs{(x - z)}$.
        Show that $\abs{(x - z)} \leq \abs{(x - y)} + \abs{(y - z)}$.
        \begin{subproof}
            $x - y = x + \neg{y}$ by \cref{real_sub}.
            $y - z = y + \neg{z}$ by \cref{real_sub}.
            $(x - y) + (y - z) = (x + \neg{y}) + (y + \neg{z})$ by \cref{real_sub}.
            $(x + \neg{y}) + (y + \neg{z}) = x + (\neg{y} + (y + \neg{z}))$ by \cref{real_add_assoc}.
            $\neg{y} + (y + \neg{z}) = (\neg{y} + y) + \neg{z}$ by \cref{real_add_assoc}.
            $\neg{y} + y = y + \neg{y}$ by \cref{real_add_comm}.
            $y + \neg{y} = 0$ by \cref{real_add_inverse}.
            Thus $\neg{y} + (y + \neg{z}) = 0 + \neg{z}$.
            $0 + \neg{z} = \neg{z}$ by \cref{real_add_ident}.
            Thus $(x - y) + (y - z) = x + \neg{z}$.
            $x + \neg{z} = x - z$ by \cref{real_sub}.
            Thus $\abs{((x - y) + (y - z))} = \abs{(x - z)}$.
            Thus $\abs{(x - z)} \leq \abs{(x - y)} + \abs{(y - z)}$ by \cref{real_triangle_inequality}.
        \end{subproof}
        Thus $d((x,z)) \leq \abs{(x - y)} + \abs{(y - z)}$.
        Thus $d((x,z)) \leq d((x,y)) + d((y,z))$.
        Thus $d((x,y)) + d((y,z)) \geq d((x,z))$.
    \end{subproof}
    Thus $d$ satisfies the triangle inequality on $\reals$ by \cref{metric_triangle_inequality}.
    Thus $d$ is a metric on $\reals$ by \cref{metric_on}.
\end{proof}

% from §20 Item 12: standard metric topology
\begin{lemma}\label{standard_metric_topology}
    Then $\metricTopology{\standardMetricR}{\reals} = \standardTopologyR$.
\end{lemma}
\begin{proof}
    Let $d = \standardMetricR$.
    Let $T = \metricTopology{d}{\reals}$.
    Let $T_1 = \standardTopologyR$.
    Let $B = \metricBasis{d}{\reals}$.
    Let $B_1 = \standardBasisR$.
    $d$ is a metric on $\reals$ by \cref{standard_metric_is_metric}.
    $d\in\funs{\reals\times\reals}{\reals}$ by \cref{metric_on}.
    Thus $B$ is a basis on $\reals$ by \cref{metric_basis_is_basis}.
    Thus $B$ is a basis for $T$ on $\reals$ by \cref{basis_for_topology,metric_topology}.
    $B_1$ is a basis on $\reals$ by \cref{standard_basis_is_basis}.
    Thus $B_1$ is a basis for $T_1$ on $\reals$ by \cref{basis_for_topology,standard_topology_reals}.
    Show that $T$ is finer than $T_1$ on $\reals$.
    \begin{subproof}
        Show that for all $x,B_2$ such that $x\in\reals$ and $B_2\in B_1$ and $x\in B_2$
            there exists $B_3\in B$ such that $x\in B_3$ and $B_3\subseteq B_2$.
        \begin{subproof}
            Fix $x$.
            Fix $B_2$.
            Assume $x\in\reals$ and $B_2\in B_1$ and $x\in B_2$.
            Take $a,b\in\reals$ such that $a < b$ and $B_2 = \intervalopen{a}{b}$ by \cref{standard_basis_reals}.
            Thus $a < x$ and $x < b$ by \cref{intervalopen}.
            Let $r_1 = x - a$.
            Let $r_2 = b - x$.
            $0 < r_1$ by \cref{real_lt_imp_pos_diff}.
            $0 < r_2$ by \cref{real_lt_imp_pos_diff}.
            Let $A = \{r_1, r_2\}$.
            $A$ is finite by \cref{pair_is_finite}.
            $r_1\in A$ by \cref{upair_intro_left}.
            Thus $A\neq\emptyset$ by \cref{emptyset}.
            Show that for all $y$ such that $y\in A$ we have $y\in\reals$.
            \begin{subproof}
                Fix $y$.
                Assume $y\in A$.
                $y = r_1$ or $y = r_2$ by \cref{upair_elim}.
                Thus $y\in\reals$.
            \end{subproof}
            Thus $A\subseteq\reals$ by \cref{subseteq}.
            Take $m\in A$ such that $m$ is a minimum of $A$ by \cref{finite_real_minimum}.
            Show that $0 < m$.
            \begin{subproof}
                $m = r_1$ or $m = r_2$ by \cref{upair_elim}.
                \begin{byCase}
                \caseOf{$m = r_1$.} Thus $0 < m$.
                \caseOf{$m = r_2$.} Thus $0 < m$.
                \end{byCase}
            \end{subproof}
            Let $B_3 = \metricBall{d}{\reals}{x}{m}$.
            Show that $B_3\in B$.
            \begin{subproof}
                $m\in\reals$ by \cref{subseteq}.
                Show that $B_3\in\pow{\reals}$.
                \begin{subproof}
                    Show that for all $y$ such that $y\in B_3$ we have $y\in\reals$.
                    \begin{subproof}
                        Fix $y$.
                        Assume $y\in B_3$.
                        Thus $y\in\reals$ by \cref{metric_ball}.
                    \end{subproof}
                    Thus $B_3\subseteq\reals$ by \cref{subseteq}.
                    Thus $B_3\in\pow{\reals}$ by \cref{powerset_intro}.
                \end{subproof}
                Thus $B_3\in B$ by \cref{metric_basis}.
            \end{subproof}
            Show that $x\in B_3$.
            \begin{subproof}
                $d$ is zero-characterized on $\reals$ by \cref{metric_on}.
                Thus $d((x,x)) = 0$ by \cref{metric_zero_characterization}.
                Thus $d((x,x)) < m$ by \cref{real_lt_subst_left}.
                Thus $x\in B_3$ by \cref{metric_ball}.
            \end{subproof}
            Show that $B_3\subseteq B_2$.
            \begin{subproof}
                Show that for all $z$ such that $z\in B_3$ we have $z\in B_2$.
                \begin{subproof}
                    Fix $z$.
                    Assume $z\in B_3$.
                    Thus $z\in\reals$ and $d((x,z)) < m$ by \cref{metric_ball}.
                    $d$ is a function by \cref{funs_elim}.
                    Show that $d((x,z)) = \abs{(x - z)}$.
                    \begin{subproof}
                        $\neg{z}\in\reals$ by \cref{real_add_inverse}.
                        $x - z = x + \neg{z}$ by \cref{real_sub}.
                        Thus $x - z\in\reals$ by \cref{real_add_closed}.
                        Thus $\abs{(x - z)}\in\reals$ by \cref{real_abs_in_reals}.
                        Let $p = (x,z)$.
                        $p\in\reals\times\reals$ by \cref{times_tuple_intro}.
                        $\fst{p} = x$ and $\snd{p} = z$ by \cref{fst_eq,snd_eq}.
                        Thus $\fst{p} - \snd{p} = x - z$.
                        Thus $\fst{p} - \snd{p}\in\reals$.
                        Let $z_0 = \abs{\fst{p} - \snd{p}}$.
                        $z_0\in\reals$ by \cref{real_abs_in_reals}.
                        Thus $\abs{\fst{p} - \snd{p}}\in\reals$ by \cref{real_abs_in_reals}.
                        Let $w = (p,z_0)$.
                        $w\in d$ by \cref{standard_metric_r}.
                        Thus $d((x,z)) = \abs{\fst{p} - \snd{p}}$ by \cref{function_apply_intro}.
                        Thus $d((x,z)) = \abs{(x - z)}$.
                    \end{subproof}
                    Thus $\abs{(x - z)} < m$ by \cref{real_lt_subst_left}.
                    $m$ is a minimum of $A$.
                    Thus for all $t\in A$ we have $m \leq t$ by \cref{real_minimum}.
                    $r_1\in A$ by \cref{upair_intro_left}.
                    $r_2\in A$ by \cref{upair_intro_right}.
                    Thus $m \leq r_1$ and $m \leq r_2$.
                Show that $a < z$.
                \begin{subproof}
                    $\neg{z}\in\reals$ by \cref{real_add_inverse}.
                    $x - z = x + \neg{z}$ by \cref{real_sub}.
                    Thus $x - z\in\reals$ by \cref{real_add_closed}.
                    $\abs{(x - z)} < m$.
                    Show that $\abs{(x - z)} < r_1$.
                    \begin{subproof}
                        Thus $\abs{(x - z)}\in\reals$ by \cref{real_abs_in_reals}.
                        $m\in\reals$ by \cref{subseteq}.
                        $r_1\in\reals$ by \cref{subseteq}.
                            $m < r_1$ or $m = r_1$ by \cref{real_leq_split}.
                            \begin{byCase}
                            \caseOf{$m < r_1$.}
                                Thus $\abs{(x - z)} < r_1$ by \cref{real_lt_trans}.
                            \caseOf{$m = r_1$.}
                                Thus $\abs{(x - z)} < r_1$ by \cref{real_lt_subst_right}.
                            \end{byCase}
                        \end{subproof}
                    Show that $x - z < r_1$.
                        \begin{subproof}
                            $\abs{(x - z)}\in\reals$ by \cref{real_abs_in_reals}.
                            $r_1\in\reals$ by \cref{subseteq}.
                            $x - z \leq \abs{(x - z)}$ by \cref{real_leq_abs}.
                            Thus $x - z < \abs{(x - z)}$ or $x - z = \abs{(x - z)}$ by \cref{real_leq_split}.
                            \begin{byCase}
                            \caseOf{$x - z < \abs{(x - z)}$.}
                                Thus $x - z < r_1$ by \cref{real_lt_trans}.
                            \caseOf{$x - z = \abs{(x - z)}$.}
                                Thus $x - z < r_1$ by \cref{real_lt_subst_left}.
                            \end{byCase}
                        \end{subproof}
                        $r_1 = x - a$.
                        Thus $x - z < x - a$ by \cref{real_lt_subst_right}.
                        $x - z\in\reals$.
                        $a\in\reals$.
                        $x\in\reals$.
                        $x - a\in\reals$ by \cref{real_sub,real_add_closed}.
                        $\neg{(x - a)} < \neg{(x - z)}$ by \cref{real_lt_neg}.
                        $\neg{(x - a)} = a - x$ by \cref{real_sub_swap_neg}.
                        $\neg{(x - z)} = z - x$ by \cref{real_sub_swap_neg}.
                        Thus $a - x < z - x$ by \cref{real_lt_subst_left,real_lt_subst_right}.
                        $\neg{x}\in\reals$ by \cref{real_add_inverse}.
                        $a - x = a + \neg{x}$ by \cref{real_sub}.
                        Thus $a - x\in\reals$ by \cref{real_add_closed}.
                        $z - x = z + \neg{x}$ by \cref{real_sub}.
                        Thus $z - x\in\reals$ by \cref{real_add_closed}.
                        $(a - x) + x < (z - x) + x$ by \cref{real_lt_add}.
                        Show that $(a - x) + x = a$.
                        \begin{subproof}
                            $a + \neg{x}\in\reals$ by \cref{real_add_closed}.
                            $a - x = a + \neg{x}$ by \cref{real_sub}.
                            Thus $(a - x) + x = (a + \neg{x}) + x$ by \cref{real_add_eq_subst}.
                            $(a + \neg{x}) + x = a + (\neg{x} + x)$ by \cref{real_add_assoc}.
                            $\neg{x} + x = x + \neg{x}$ by \cref{real_add_comm}.
                            $x + \neg{x} = 0$ by \cref{real_add_inverse}.
                            Thus $\neg{x} + x = 0$ by \cref{real_add_comm,real_add_inverse}.
                            $\neg{x} + x\in\reals$ by \cref{real_add_closed}.
                            $0\in\reals$ by \cref{real_add_ident}.
                            Thus $a + (\neg{x} + x) = a + 0$ by \cref{real_add_eq_subst}.
                            $a + 0 = a$ by \cref{real_add_ident,real_add_comm}.
                            Thus $(a - x) + x = a$.
                        \end{subproof}
                        Show that $(z - x) + x = z$.
                        \begin{subproof}
                            $z\in\reals$ by \cref{metric_ball}.
                            $x\in\reals$.
                            $\neg{x}\in\reals$ by \cref{real_add_inverse}.
                            $z + \neg{x}\in\reals$ by \cref{real_add_closed}.
                            $z - x = z + \neg{x}$ by \cref{real_sub}.
                            Thus $(z - x) + x = (z + \neg{x}) + x$ by \cref{real_add_eq_subst}.
                            $(z + \neg{x}) + x = z + (\neg{x} + x)$ by \cref{real_add_assoc}.
                            $\neg{x} + x = x + \neg{x}$ by \cref{real_add_comm}.
                            $x + \neg{x} = 0$ by \cref{real_add_inverse}.
                            Thus $\neg{x} + x = 0$ by \cref{real_add_comm,real_add_inverse}.
                            $\neg{x} + x\in\reals$ by \cref{real_add_closed}.
                            $0\in\reals$ by \cref{real_add_ident}.
                            Thus $z + (\neg{x} + x) = z + 0$ by \cref{real_add_eq_subst}.
                            $z + 0 = z$ by \cref{real_add_ident,real_add_comm}.
                            Thus $(z - x) + x = z$.
                        \end{subproof}
                        Thus $a < z$.
                    \end{subproof}
                Show that $z < b$.
                \begin{subproof}
                    $\neg{z}\in\reals$ by \cref{real_add_inverse}.
                    $x - z = x + \neg{z}$ by \cref{real_sub}.
                    Thus $x - z\in\reals$ by \cref{real_add_closed}.
                    $\abs{(x - z)} < m$.
                    Show that $\abs{(x - z)} < r_2$.
                    \begin{subproof}
                        Thus $\abs{(x - z)}\in\reals$ by \cref{real_abs_in_reals}.
                        $m\in\reals$ by \cref{subseteq}.
                        $r_2\in\reals$ by \cref{subseteq}.
                            $m < r_2$ or $m = r_2$ by \cref{real_leq_split}.
                            \begin{byCase}
                            \caseOf{$m < r_2$.}
                                Thus $\abs{(x - z)} < r_2$ by \cref{real_lt_trans}.
                            \caseOf{$m = r_2$.}
                                Thus $\abs{(x - z)} < r_2$ by \cref{real_lt_subst_right}.
                            \end{byCase}
                        \end{subproof}
                        $\neg{z}\in\reals$ by \cref{real_add_inverse}.
                        $x - z = x + \neg{z}$ by \cref{real_sub}.
                        Thus $x - z\in\reals$ by \cref{real_add_closed}.
                        $\neg{(x - z)}\in\reals$ by \cref{real_add_inverse}.
                        $\abs{(x - z)}\in\reals$ by \cref{real_abs_in_reals}.
                        $r_2\in\reals$ by \cref{subseteq}.
                        Show that $\neg{(x - z)} < r_2$.
                        \begin{subproof}
                            $\neg{(x - z)} \leq \abs{(x - z)}$ by \cref{real_neg_leq_abs}.
                            Thus $\neg{(x - z)} < \abs{(x - z)}$ or $\neg{(x - z)} = \abs{(x - z)}$ by \cref{real_leq_split}.
                            \begin{byCase}
                            \caseOf{$\neg{(x - z)} < \abs{(x - z)}$.}
                                Thus $\neg{(x - z)} < r_2$ by \cref{real_lt_trans}.
                            \caseOf{$\neg{(x - z)} = \abs{(x - z)}$.}
                                Thus $\neg{(x - z)} < r_2$ by \cref{real_lt_subst_left}.
                            \end{byCase}
                        \end{subproof}
                        $\neg{(x - z)} = z - x$ by \cref{real_sub_swap_neg}.
                        Thus $z - x < r_2$ by \cref{real_lt_subst_left}.
                        $r_2 = b - x$.
                        Thus $z - x < b - x$ by \cref{real_lt_subst_right}.
                        $z - x\in\reals$ by \cref{real_sub,real_add_closed}.
                        $b - x\in\reals$ by \cref{real_sub,real_add_closed}.
                        $(z - x) + x < (b - x) + x$ by \cref{real_lt_add}.
                        Show that $(z - x) + x = z$.
                        \begin{subproof}
                            $z\in\reals$ by \cref{metric_ball}.
                            $x\in\reals$.
                            $\neg{x}\in\reals$ by \cref{real_add_inverse}.
                            $z + \neg{x}\in\reals$ by \cref{real_add_closed}.
                            $z - x = z + \neg{x}$ by \cref{real_sub}.
                            Thus $(z - x) + x = (z + \neg{x}) + x$ by \cref{real_add_eq_subst}.
                            $(z + \neg{x}) + x = z + (\neg{x} + x)$ by \cref{real_add_assoc}.
                            $\neg{x} + x = x + \neg{x}$ by \cref{real_add_comm}.
                            $x + \neg{x} = 0$ by \cref{real_add_inverse}.
                            Thus $\neg{x} + x = 0$ by \cref{real_add_comm,real_add_inverse}.
                            $\neg{x} + x\in\reals$ by \cref{real_add_closed}.
                            $0\in\reals$ by \cref{real_add_ident}.
                            Thus $z + (\neg{x} + x) = z + 0$ by \cref{real_add_eq_subst}.
                            $z + 0 = z$ by \cref{real_add_ident,real_add_comm}.
                            Thus $(z - x) + x = z$.
                        \end{subproof}
                        Show that $(b - x) + x = b$.
                        \begin{subproof}
                            $b\in\reals$.
                            $x\in\reals$.
                            $\neg{x}\in\reals$ by \cref{real_add_inverse}.
                            $b + \neg{x}\in\reals$ by \cref{real_add_closed}.
                            $b - x = b + \neg{x}$ by \cref{real_sub}.
                            Thus $(b - x) + x = (b + \neg{x}) + x$ by \cref{real_add_eq_subst}.
                            $(b + \neg{x}) + x = b + (\neg{x} + x)$ by \cref{real_add_assoc}.
                            $\neg{x} + x = x + \neg{x}$ by \cref{real_add_comm}.
                            $x + \neg{x} = 0$ by \cref{real_add_inverse}.
                            Thus $\neg{x} + x = 0$ by \cref{real_add_comm,real_add_inverse}.
                            $\neg{x} + x\in\reals$ by \cref{real_add_closed}.
                            $0\in\reals$ by \cref{real_add_ident}.
                            Thus $b + (\neg{x} + x) = b + 0$ by \cref{real_add_eq_subst}.
                            $b + 0 = b$ by \cref{real_add_ident,real_add_comm}.
                            Thus $(b - x) + x = b$.
                        \end{subproof}
                        Thus $z < b$.
                    \end{subproof}
                    Thus $z\in\intervalopen{a}{b}$ by \cref{intervalopen}.
                    Thus $z\in B_2$.
                \end{subproof}
                Thus $B_3\subseteq B_2$ by \cref{subseteq}.
            \end{subproof}
        \end{subproof}
        Thus $T$ is finer than $T_1$ on $\reals$ by \cref{basis_finer_criterion}.
    \end{subproof}
    Show that $T_1$ is finer than $T$ on $\reals$.
    \begin{subproof}
        Show that for all $x,B_2$ such that $x\in\reals$ and $B_2\in B$ and $x\in B_2$
            there exists $B_3\in B_1$ such that $x\in B_3$ and $B_3\subseteq B_2$.
        \begin{subproof}
            Fix $x$.
            Fix $B_2$.
            Assume $x\in\reals$ and $B_2\in B$ and $x\in B_2$.
            Take $x_0,\epsilon\in\reals$ such that $0 < \epsilon$ and $B_2 = \metricBall{d}{\reals}{x_0}{\epsilon}$
                by \cref{metric_basis}.
            Thus $x\in\reals$ and $d((x_0,x)) < \epsilon$ by \cref{metric_ball}.
            $d\in\funs{\reals\times\reals}{\reals}$ by \cref{metric_on}.
            $(x_0,x)\in\reals\times\reals$ by \cref{times_tuple_intro}.
            Thus $d((x_0,x))\in\reals$ by \cref{funs_type_apply}.
            Let $\delta = \epsilon - d((x_0,x))$.
            $\neg{d((x_0,x))}\in\reals$ by \cref{real_add_inverse}.
            $\delta = \epsilon + \neg{d((x_0,x))}$ by \cref{real_sub}.
            Thus $\delta\in\reals$ by \cref{real_add_closed}.
            $0 < \delta$ by \cref{real_lt_imp_pos_diff}.
            Show that $x - \delta < x$.
            \begin{subproof}
                $\neg{\delta}\in\reals$ by \cref{real_add_inverse}.
                $x - \delta = x + \neg{\delta}$ by \cref{real_sub}.
                $0\in\reals$ by \cref{real_add_ident}.
                $0 < \delta$.
                Thus $\neg{\delta} < \neg{0}$ by \cref{real_lt_neg}.
                $\neg{0} = 0$ by \cref{real_add_inverse_unique,real_add_ident}.
                Thus $\neg{\delta} < 0$ by \cref{real_lt_subst_right}.
                $\neg{\delta} + x < 0 + x$ by \cref{real_lt_add}.
                $\neg{\delta} + x = x + \neg{\delta}$ by \cref{real_add_comm}.
                Thus $x + \neg{\delta} < 0 + x$ by \cref{real_lt_subst_left}.
                Thus $x - \delta < 0 + x$ by \cref{real_lt_subst_left,real_sub}.
                $0 + x = x$ by \cref{real_add_ident}.
                Thus $x - \delta < x$ by \cref{real_lt_subst_right}.
            \end{subproof}
            Show that $x < x + \delta$.
            \begin{subproof}
                $0\in\reals$ by \cref{real_add_ident}.
                $0 + x < \delta + x$ by \cref{real_lt_add}.
                $0 + x = x$ by \cref{real_add_ident}.
                $\delta + x = x + \delta$ by \cref{real_add_comm}.
                Thus $x < x + \delta$ by \cref{real_lt_subst_left,real_lt_subst_right}.
            \end{subproof}
            Let $B_3 = \intervalopen{x - \delta}{x + \delta}$.
            Show that $B_3\in B_1$.
            \begin{subproof}
                $\neg{\delta}\in\reals$ by \cref{real_add_inverse}.
                $x - \delta = x + \neg{\delta}$ by \cref{real_sub}.
                Thus $x - \delta\in\reals$ by \cref{real_add_closed}.
                $x + \delta\in\reals$ by \cref{real_add_closed}.
                Thus $x - \delta < x + \delta$ by \cref{real_lt_trans}.
                $B_3\subseteq\reals$ by \cref{intervalopen_subset_reals}.
                Thus $B_3\in\pow{\reals}$ by \cref{powerset_intro}.
                Thus $B_3\in B_1$ by \cref{standard_basis_reals}.
            \end{subproof}
            Show that $x\in B_3$.
            \begin{subproof}
                Thus $x\in\intervalopen{x - \delta}{x + \delta}$ by \cref{intervalopen}.
                Thus $x\in B_3$.
            \end{subproof}
            Show that $B_3\subseteq B_2$.
            \begin{subproof}
                Show that for all $z$ such that $z\in B_3$ we have $z\in B_2$.
                \begin{subproof}
                    Fix $z$.
                    Assume $z\in B_3$.
                    Thus $x - \delta < z$ and $z < x + \delta$ by \cref{intervalopen}.
                    Thus $z\in\reals$ by \cref{intervalopen}.
                    Show that $d((x,z)) < \delta$.
                    \begin{subproof}
                        $\neg{z}\in\reals$ by \cref{real_add_inverse}.
                        $x - z = x + \neg{z}$ by \cref{real_sub}.
                        Thus $x - z\in\reals$ by \cref{real_add_closed}.
                        $\neg{x}\in\reals$ by \cref{real_add_inverse}.
                        $z - x = z + \neg{x}$ by \cref{real_sub}.
                        Thus $z - x\in\reals$ by \cref{real_add_closed}.
                        Show that $z - x < \delta$.
                        \begin{subproof}
                            $z < x + \delta$.
                            $z\in\reals$ by \cref{intervalopen}.
                            $x + \delta\in\reals$ by \cref{real_add_closed}.
                            $\neg{x}\in\reals$ by \cref{real_add_inverse}.
                            $z + \neg{x} < (x + \delta) + \neg{x}$ by \cref{real_lt_add}.
                            $z + \neg{x} = z - x$ by \cref{real_sub}.
                            $(x + \delta) + \neg{x} = x + (\delta + \neg{x})$ by \cref{real_add_assoc}.
                            $\delta + \neg{x} = \neg{x} + \delta$ by \cref{real_add_comm}.
                            $x + (\neg{x} + \delta) = (x + \neg{x}) + \delta$ by \cref{real_add_assoc}.
                            $x + \neg{x} = 0$ by \cref{real_add_inverse}.
                            Thus $(x + \delta) + \neg{x} = 0 + \delta$ by \cref{real_add_assoc,real_add_comm,real_add_inverse}.
                            $0 + \delta = \delta$ by \cref{real_add_ident}.
                            Thus $z - x < \delta$ by \cref{real_lt_subst_left,real_lt_subst_right}.
                        \end{subproof}
                        Show that $x - z < \delta$.
                        \begin{subproof}
                            $x - \delta < z$.
                            $\neg{\delta}\in\reals$ by \cref{real_add_inverse}.
                            $x - \delta = x + \neg{\delta}$ by \cref{real_sub}.
                            Thus $x - \delta\in\reals$ by \cref{real_add_closed}.
                            $\neg{x}\in\reals$ by \cref{real_add_inverse}.
                            $(x - \delta) + \neg{x} < z + \neg{x}$ by \cref{real_lt_add}.
                            $(x - \delta) + \neg{x} = \neg{\delta}$ by \cref{real_sub,real_add_assoc,real_add_inverse,real_add_comm,real_add_ident}.
                            $z + \neg{x} = z - x$ by \cref{real_sub}.
                            Thus $\neg{\delta} < z - x$ by \cref{real_lt_subst_left,real_lt_subst_right}.
                            Thus $\neg{(z - x)} < \neg{(\neg{\delta})}$ by \cref{real_lt_neg}.
                            $\neg{(z - x)} = x - z$ by \cref{real_sub_swap_neg}.
                            $\delta + \neg{\delta} = 0$ by \cref{real_add_inverse}.
                            $\neg{\delta} + \delta = \delta + \neg{\delta}$ by \cref{real_add_comm}.
                            Thus $\neg{\delta} + \delta = 0$ by \cref{real_add_inverse}.
                            Thus $\delta = \neg{(\neg{\delta})}$ by \cref{real_add_inverse_unique}.
                            Thus $x - z < \delta$ by \cref{real_lt_subst_left,real_lt_subst_right}.
                        \end{subproof}
                        Thus $x - z \leq \delta$ and $\neg{(x - z)} \leq \delta$.
                        Thus $\abs{(x - z)} \leq \delta$ by \cref{real_abs_leq_if_pm_leq}.
                        Thus $\abs{(x - z)} < \delta$.
                        $d$ is a function by \cref{funs_elim}.
                        Let $p = (x,z)$.
                        $p\in\reals\times\reals$ by \cref{times_tuple_intro}.
                        $\fst{p} = x$ and $\snd{p} = z$ by \cref{fst_eq,snd_eq}.
                        Thus $\fst{p} - \snd{p} = x - z$.
                        Thus $\fst{p} - \snd{p}\in\reals$.
                        Thus $\abs{\fst{p} - \snd{p}}\in\reals$ by \cref{real_abs_in_reals}.
                        Let $w = (p,\abs{\fst{p} - \snd{p}})$.
                        $w\in d$ by \cref{standard_metric_r}.
                        Thus $d((x,z)) = \abs{\fst{p} - \snd{p}}$ by \cref{function_apply_intro}.
                        Thus $d((x,z)) = \abs{(x - z)}$.
                        Thus $d((x,z)) < \delta$.
                    \end{subproof}
                    $d$ satisfies the triangle inequality on $\reals$ by \cref{metric_on}.
                    Thus $d((x_0,z)) \leq d((x_0,x)) + d((x,z))$ by \cref{metric_triangle_inequality}.
                    $d((x_0,x))\in\reals$ by \cref{funs_type_apply}.
                    $(x,z)\in\reals\times\reals$ by \cref{times_tuple_intro}.
                    $d((x,z))\in\reals$ by \cref{funs_type_apply}.
                    Thus $d((x_0,x)) + d((x,z))\in\reals$ by \cref{real_add_closed}.
                    Thus $d((x_0,x)) + \delta\in\reals$ by \cref{real_add_closed}.
                    $d((x,z)) + d((x_0,x)) < \delta + d((x_0,x))$ by \cref{real_lt_add}.
                    $d((x,z)) + d((x_0,x)) = d((x_0,x)) + d((x,z))$ by \cref{real_add_comm}.
                    $\delta + d((x_0,x)) = d((x_0,x)) + \delta$ by \cref{real_add_comm}.
                    Thus $d((x_0,x)) + d((x,z)) < d((x_0,x)) + \delta$
                        by \cref{real_lt_subst_left,real_lt_subst_right}.
                    $(x_0,z)\in\reals\times\reals$ by \cref{times_tuple_intro}.
                    Thus $d((x_0,z))\in\reals$ by \cref{funs_type_apply}.
                    Show that $d((x_0,z)) < d((x_0,x)) + \delta$.
                    \begin{subproof}
                        Thus $d((x_0,z)) < d((x_0,x)) + d((x,z))$ or
                            $d((x_0,z)) = d((x_0,x)) + d((x,z))$ by \cref{real_leq_split}.
                        \begin{byCase}
                        \caseOf{$d((x_0,z)) < d((x_0,x)) + d((x,z))$.}
                            Thus $d((x_0,z)) < d((x_0,x)) + \delta$ by \cref{real_lt_trans}.
                        \caseOf{$d((x_0,z)) = d((x_0,x)) + d((x,z))$.}
                            Thus $d((x_0,z)) < d((x_0,x)) + \delta$ by \cref{real_lt_subst_left}.
                        \end{byCase}
                    \end{subproof}
                    Show that $d((x_0,x)) + \delta = \epsilon$.
                    \begin{subproof}
                        $\delta = \epsilon + \neg{d((x_0,x))}$ by \cref{real_sub}.
                        $d((x_0,x)) + \delta = \delta + d((x_0,x))$ by \cref{real_add_comm}.
                        Thus $d((x_0,x)) + \delta = (\epsilon + \neg{d((x_0,x))}) + d((x_0,x))$
                            by \cref{real_add_eq_subst}.
                        $(\epsilon + \neg{d((x_0,x))}) + d((x_0,x)) =
                            \epsilon + (\neg{d((x_0,x))} + d((x_0,x)))$ by \cref{real_add_assoc}.
                        $\neg{d((x_0,x))} + d((x_0,x)) = 0$ by \cref{real_add_comm,real_add_inverse}.
                        Thus $d((x_0,x)) + \delta = \epsilon + 0$ by \cref{real_add_eq_subst}.
                        $\epsilon + 0 = \epsilon$ by \cref{real_add_ident,real_add_comm}.
                        Thus $d((x_0,x)) + \delta = \epsilon$.
                    \end{subproof}
                    Thus $d((x_0,z)) < \epsilon$ by \cref{real_lt_subst_right}.
                    Thus $z\in B_2$ by \cref{metric_ball}.
                \end{subproof}
                Thus $B_3\subseteq B_2$ by \cref{subseteq}.
            \end{subproof}
        \end{subproof}
        Thus $T_1$ is finer than $T$ on $\reals$ by \cref{basis_finer_criterion}.
    \end{subproof}
    Thus $T\subseteq T_1$ by \cref{finer_topology}.
    Thus $T_1\subseteq T$ by \cref{finer_topology}.
    Thus $T = T_1$ by \cref{subseteq_antisymmetric}.
\end{proof}

% from §20 Item 13: metrizable space
\begin{definition}\label{metrizable}
    $X$ is metrizable with respect to $T$ iff
    $T$ is a topology on $X$ and
    there exists $d$ such that $d$ is a metric on $X$ and $T = \metricTopology{d}{X}$.
\end{definition}

% from §20 Item 14: metric space
\begin{definition}\label{metric_space}
    $X$ is a metric space with metric $d$ and topology $T$ iff
    $d$ is a metric on $X$ and $T = \metricTopology{d}{X}$.
\end{definition}

% from §20 Item 15: bounded subset
\begin{definition}\label{metric_bounded}
    $A$ is bounded in $X$ with respect to $d$ iff
    $d$ is a metric on $X$ and
    $A\subseteq X$ and
    there exists $M\in\reals$ such that for all $a_1,a_2\in A$ we have $d((a_1,a_2)) \leq M$.
\end{definition}

% from §20 helper lemma 5: distance set of a subset
\begin{definition}\label{metric_distance_set}
    $\metricDistanceSet{A}{d} = \{d((a_1,a_2)) \mid a_1\in A, a_2\in A\}$.
\end{definition}

% from §20 Item 16: diameter of a set
\begin{definition}\label{metric_diameter}
    $r$ is a diameter of $A$ in $X$ with respect to $d$ iff
    $A$ is bounded in $X$ with respect to $d$ and
    $A\neq\emptyset$ and
    $r$ is a supremum of $\metricDistanceSet{A}{d}$.
\end{definition}

% from §20 Item 17: standard bounded metric
\begin{definition}\label{bounded_metric}
    $\boundedMetric{d} = \{w \mid \exists p\in\dom{d}. \exists z\in\reals.
        w = (p,z) \land ((z = d(p) \land d(p) < 1) \lor (z = 1 \land 1 \leq d(p)))\}$.
\end{definition}

% from §20 Item 18: standard bounded metric is a metric
\begin{theorem}\label{bounded_metric_is_metric}
    Assume $d$ is a metric on $X$.
    Then $\boundedMetric{d}$ is a metric on $X$.
\end{theorem}
\begin{proof}
    Let $e = \boundedMetric{d}$.
    Show that $e\in\funs{X\times X}{\reals}$.
    \begin{subproof}
        Show that $e$ is a function to $\reals$.
        \begin{subproof}
            Show that $e$ is a function.
            \begin{subproof}
                Show that for all $p,z_1,z_2$ such that $(p,z_1)\in e$ and $(p,z_2)\in e$ we have $z_1 = z_2$.
                \begin{subproof}
                    Fix $p$.
                    Fix $z_1$.
                    Fix $z_2$.
                    Assume $(p,z_1)\in e$ and $(p,z_2)\in e$.
                    Take $p_1, w_1$ such that
                        $p_1\in\dom{d}$ and $w_1\in\reals$ and
                        $(p,z_1) = (p_1,w_1)$ and
                        $((w_1 = d(p_1) \land d(p_1) < 1) \lor (w_1 = 1 \land 1 \leq d(p_1)))$
                        by \cref{bounded_metric}.
                    Take $p_2, w_2$ such that
                        $p_2\in\dom{d}$ and $w_2\in\reals$ and
                        $(p,z_2) = (p_2,w_2)$ and
                        $((w_2 = d(p_2) \land d(p_2) < 1) \lor (w_2 = 1 \land 1 \leq d(p_2)))$
                        by \cref{bounded_metric}.
                    $p = p_1$ and $z_1 = w_1$ by \cref{pair_eq_iff}.
                    $p = p_2$ and $z_2 = w_2$ by \cref{pair_eq_iff}.
                    Thus $((z_1 = d(p) \land d(p) < 1) \lor (z_1 = 1 \land 1 \leq d(p)))$.
                    Thus $((z_2 = d(p) \land d(p) < 1) \lor (z_2 = 1 \land 1 \leq d(p)))$.
                    $d\in\funs{X\times X}{\reals}$ by \cref{metric_on}.
                    Thus $d(p)\in\reals$ by \cref{funs_elim}.
                    $1\in\reals$ by \cref{real_naturals_subset_reals,real_one_in_naturals,subseteq}.
                    \begin{byCase}
                    \caseOf{$d(p) < 1$.}
                        Show that $z_1 = d(p)$.
                        \begin{subproof}
                            Suppose not.
                            \begin{byCase}
                            \caseOf{$z_1 = d(p) \land d(p) < 1$.}
                                Thus $z_1 = d(p)$.
                            \caseOf{$z_1 = 1 \land 1 \leq d(p)$.}
                                Thus $1 < d(p)$ or $1 = d(p)$ by \cref{real_leq_split}.
                                \begin{byCase}
                                \caseOf{$1 < d(p)$.}
                                    Thus $1 < 1$ by \cref{real_lt_trans}.
                                    Contradiction by \cref{real_lt_irreflexive}.
                                \caseOf{$1 = d(p)$.}
                                    Thus $1 < 1$ by \cref{real_lt_subst_left}.
                                    Contradiction by \cref{real_lt_irreflexive}.
                                \end{byCase}
                            \end{byCase}
                        \end{subproof}
                        Show that $z_2 = d(p)$.
                        \begin{subproof}
                            Suppose not.
                            \begin{byCase}
                            \caseOf{$z_2 = d(p) \land d(p) < 1$.}
                                Thus $z_2 = d(p)$.
                            \caseOf{$z_2 = 1 \land 1 \leq d(p)$.}
                                Thus $1 < d(p)$ or $1 = d(p)$ by \cref{real_leq_split}.
                                \begin{byCase}
                                \caseOf{$1 < d(p)$.}
                                    Thus $1 < 1$ by \cref{real_lt_trans}.
                                    Contradiction by \cref{real_lt_irreflexive}.
                                \caseOf{$1 = d(p)$.}
                                    Thus $1 < 1$ by \cref{real_lt_subst_left}.
                                    Contradiction by \cref{real_lt_irreflexive}.
                                \end{byCase}
                            \end{byCase}
                        \end{subproof}
                        Thus $z_1 = z_2$.
                    \caseOf{it is wrong that $d(p) < 1$.}
                        Show that $z_1 = 1$.
                        \begin{subproof}
                            Suppose not.
                            \begin{byCase}
                            \caseOf{$z_1 = d(p) \land d(p) < 1$.}
                                Contradiction.
                            \caseOf{$z_1 = 1 \land 1 \leq d(p)$.}
                                Thus $z_1 = 1$.
                            \end{byCase}
                        \end{subproof}
                        Show that $z_2 = 1$.
                        \begin{subproof}
                            Suppose not.
                            \begin{byCase}
                            \caseOf{$z_2 = d(p) \land d(p) < 1$.}
                                Contradiction.
                            \caseOf{$z_2 = 1 \land 1 \leq d(p)$.}
                                Thus $z_2 = 1$.
                            \end{byCase}
                        \end{subproof}
                        Thus $z_1 = z_2$.
                    \end{byCase}
                \end{subproof}
                Thus $e$ is right-unique by \cref{rightunique}.
                Show that for all $w$ such that $w\in e$ there exist $x,y$ such that $w = (x,y)$.
                \begin{subproof}
                    Fix $w$.
                    Assume $w\in e$.
                    Take $p,z$ such that
                        $p\in\dom{d}$ and $z\in\reals$ and
                        $w = (p,z)$ and
                        $((z = d(p) \land d(p) < 1) \lor (z = 1 \land 1 \leq d(p)))$
                        by \cref{bounded_metric}.
                    Thus there exist $x,y$ such that $w = (x,y)$.
                \end{subproof}
                Thus $e$ is a relation by \cref{relation}.
                Thus $e$ is a function.
            \end{subproof}
            Show that for all $p\in\dom{e}$ we have $e(p)\in\reals$.
            \begin{subproof}
                Fix $p$.
                Assume $p\in\dom{e}$.
                Take $z$ such that $(p,z)\in e$ by \cref{dom_iff}.
                Take $p_0, w$ such that
                    $p_0\in\dom{d}$ and $w\in\reals$ and
                    $(p,z) = (p_0,w)$ and
                    $((w = d(p_0) \land d(p_0) < 1) \lor (w = 1 \land 1 \leq d(p_0)))$
                    by \cref{bounded_metric}.
                $p = p_0$ and $z = w$ by \cref{pair_eq_iff}.
                Thus $z\in\reals$.
                $e$ is a function.
                Thus $e(p) = z$ by \cref{function_apply_intro}.
                Thus $e(p)\in\reals$.
            \end{subproof}
            Thus $e$ is a function to $\reals$.
        \end{subproof}
        Show that $\dom{e} = X\times X$.
        \begin{subproof}
            Show that for all $p$ such that $p\in\dom{e}$ we have $p\in X\times X$.
            \begin{subproof}
                Fix $p$.
                Assume $p\in\dom{e}$.
                Take $z$ such that $(p,z)\in e$ by \cref{dom_iff}.
                Take $p_0, w$ such that
                    $p_0\in\dom{d}$ and $w\in\reals$ and
                    $(p,z) = (p_0,w)$ and
                    $((w = d(p_0) \land d(p_0) < 1) \lor (w = 1 \land 1 \leq d(p_0)))$
                    by \cref{bounded_metric}.
                $p = p_0$ by \cref{pair_eq_iff}.
                $d\in\funs{X\times X}{\reals}$ by \cref{metric_on}.
                Thus $\dom{d} = X\times X$ by \cref{funs_elim}.
                Thus $p\in X\times X$ by \cref{subseteq}.
            \end{subproof}
            Thus $\dom{e}\subseteq X\times X$ by \cref{subseteq}.
            Show that for all $p$ such that $p\in X\times X$ we have $p\in\dom{e}$.
            \begin{subproof}
                Fix $p$.
                Assume $p\in X\times X$.
                $d\in\funs{X\times X}{\reals}$ by \cref{metric_on}.
                Thus $\dom{d} = X\times X$ by \cref{funs_elim}.
                Thus $p\in\dom{d}$ by \cref{subseteq}.
                Thus $d(p)\in\reals$ by \cref{funs_elim}.
                $1\in\reals$ by \cref{real_naturals_subset_reals,real_one_in_naturals,subseteq}.
                \begin{byCase}
                \caseOf{$d(p) < 1$.}
                    Let $w = (p,d(p))$.
                    $w\in e$ by \cref{bounded_metric}.
                    Thus $p\in\dom{e}$ by \cref{dom_intro}.
                \caseOf{it is wrong that $d(p) < 1$.}
                    Show that $1 \leq d(p)$.
                    \begin{subproof}
                        Suppose not.
                        \begin{byCase}
                        \caseOf{$d(p) = 1$.}
                            Thus $1 < d(p)$ or $1 = d(p)$.
                            Thus $1 \leq d(p)$.
                        \caseOf{$d(p) \neq 1$.}
                            Thus $d(p) < 1$ or $1 < d(p)$ by \cref{real_lt_trichotomy}.
                            \begin{byCase}
                            \caseOf{$d(p) < 1$.}
                                Contradiction.
                            \caseOf{$1 < d(p)$.}
                                Thus $1 < d(p)$ or $1 = d(p)$.
                                Thus $1 \leq d(p)$.
                            \end{byCase}
                        \end{byCase}
                    \end{subproof}
                    Let $w = (p,1)$.
                    $w\in e$ by \cref{bounded_metric}.
                    Thus $p\in\dom{e}$ by \cref{dom_intro}.
                \end{byCase}
            \end{subproof}
            Thus $X\times X\subseteq \dom{e}$ by \cref{subseteq}.
            Thus $\dom{e} = X\times X$ by \cref{subseteq_antisymmetric}.
        \end{subproof}
        Thus $e\in\funs{X\times X}{\reals}$ by \cref{funs_intro}.
    \end{subproof}
    Show that for all $x,y$ such that $x\in X$ and $y\in X$ we have $e((x,y)) \geq 0$.
    \begin{subproof}
        Fix $x$.
        Fix $y$.
        Assume $x\in X$ and $y\in X$.
        Let $p = (x,y)$.
        $p\in X\times X$ by \cref{times_tuple_intro}.
        $d\in\funs{X\times X}{\reals}$ by \cref{metric_on}.
        Thus $\dom{d} = X\times X$ by \cref{funs_elim}.
        Thus $p\in\dom{d}$ by \cref{subseteq}.
        Thus $d(p)\in\reals$ by \cref{funs_elim}.
        $1\in\reals$ by \cref{real_naturals_subset_reals,real_one_in_naturals,subseteq}.
        \begin{byCase}
        \caseOf{$d(p) < 1$.}
            Let $w = (p,d(p))$.
            $w\in e$ by \cref{bounded_metric}.
            $e$ is a function by \cref{funs_elim}.
            Thus $e(p) = d(p)$ by \cref{function_apply_intro}.
            $d$ is nonnegative on $X$ by \cref{metric_on}.
            Thus $d((x,y)) \geq 0$ by \cref{metric_nonnegative}.
            Thus $e((x,y)) \geq 0$.
        \caseOf{it is wrong that $d(p) < 1$.}
            Show that $1 \leq d(p)$.
            \begin{subproof}
                Suppose not.
                \begin{byCase}
                \caseOf{$d(p) = 1$.}
                    Thus $1 < d(p)$ or $1 = d(p)$.
                    Thus $1 \leq d(p)$.
                \caseOf{$d(p) \neq 1$.}
                    Thus $d(p) < 1$ or $1 < d(p)$ by \cref{real_lt_trichotomy}.
                    \begin{byCase}
                    \caseOf{$d(p) < 1$.}
                        Contradiction.
                    \caseOf{$1 < d(p)$.}
                        Thus $1 < d(p)$ or $1 = d(p)$.
                        Thus $1 \leq d(p)$.
                    \end{byCase}
                \end{byCase}
            \end{subproof}
            Let $w = (p,1)$.
            $w\in e$ by \cref{bounded_metric}.
            $e$ is a function by \cref{funs_elim}.
            Thus $e(p) = 1$ by \cref{function_apply_intro}.
            $0 < 1$ by \cref{real_zero_lt_one}.
            Thus $0 < 1$ or $0 = 1$.
            Thus $0 \leq 1$.
            Thus $e((x,y)) \geq 0$.
        \end{byCase}
    \end{subproof}
    Thus $e$ is nonnegative on $X$ by \cref{metric_nonnegative}.
    Show that for all $x,y$ such that $x\in X$ and $y\in X$ we have $e((x,y)) = 0$ iff $x = y$.
    \begin{subproof}
        Fix $x$.
        Fix $y$.
        Assume $x\in X$ and $y\in X$.
        Show that if $e((x,y)) = 0$ then $x = y$.
        \begin{subproof}
            Suppose not.
            Thus $e((x,y)) = 0$ and $x \neq y$.
            Let $p = (x,y)$.
            $p\in X\times X$ by \cref{times_tuple_intro}.
            $e\in\funs{X\times X}{\reals}$.
            Thus $\dom{e} = X\times X$ by \cref{funs_elim}.
            Thus $p\in\dom{e}$ by \cref{subseteq}.
            $e$ is a function by \cref{funs_elim}.
            Thus $(p,0)\in e$ by \cref{function_apply_iff}.
            Take $p_0, w$ such that
                $p_0\in\dom{d}$ and $w\in\reals$ and
                $(p,0) = (p_0,w)$ and
                $((w = d(p_0) \land d(p_0) < 1) \lor (w = 1 \land 1 \leq d(p_0)))$
                by \cref{bounded_metric}.
            $p = p_0$ and $0 = w$ by \cref{pair_eq_iff}.
            Thus $((0 = d(p) \land d(p) < 1) \lor (0 = 1 \land 1 \leq d(p)))$.
            \begin{byCase}
            \caseOf{$0 = d(p) \land d(p) < 1$.}
                Thus $0 = d(p)$.
                $d$ is zero-characterized on $X$ by \cref{metric_on}.
                Thus $x = y$ by \cref{metric_zero_characterization}.
            \caseOf{$0 = 1 \land 1 \leq d(p)$.}
                $0 < 1$ by \cref{real_zero_lt_one}.
                Thus $0 < 0$ by \cref{real_lt_subst_right}.
                Contradiction by \cref{real_lt_irreflexive}.
            \end{byCase}
        \end{subproof}
        Show that if $x = y$ then $e((x,y)) = 0$.
        \begin{subproof}
            Assume $x = y$.
            Let $p = (x,y)$.
            $p\in X\times X$ by \cref{times_tuple_intro}.
            $d\in\funs{X\times X}{\reals}$ by \cref{metric_on}.
            Thus $\dom{d} = X\times X$ by \cref{funs_elim}.
            Thus $p\in\dom{d}$ by \cref{subseteq}.
            Thus $d(p)\in\reals$ by \cref{funs_elim}.
            $d$ is zero-characterized on $X$ by \cref{metric_on}.
            Thus $d((x,y)) = 0$ by \cref{metric_zero_characterization}.
            $0 < 1$ by \cref{real_zero_lt_one}.
            Thus $d((x,y)) < 1$ by \cref{real_lt_subst_left}.
            Let $w = (p,d(p))$.
            $w\in e$ by \cref{bounded_metric}.
            $e$ is a function by \cref{funs_elim}.
            Thus $e(p) = d(p)$ by \cref{function_apply_intro}.
            Thus $e((x,y)) = 0$.
        \end{subproof}
    \end{subproof}
    Thus $e$ is zero-characterized on $X$ by \cref{metric_zero_characterization}.
    Show that for all $x,y$ such that $x\in X$ and $y\in X$ we have $e((x,y)) = e((y,x))$.
    \begin{subproof}
        Fix $x$.
        Fix $y$.
        Assume $x\in X$ and $y\in X$.
        Let $p_1 = (x,y)$.
        Let $p_2 = (y,x)$.
        $p_1\in X\times X$ by \cref{times_tuple_intro}.
        $p_2\in X\times X$ by \cref{times_tuple_intro}.
        $d$ is symmetric on $X$ by \cref{metric_on}.
        Thus $d((x,y)) = d((y,x))$ by \cref{metric_symmetric}.
        $d\in\funs{X\times X}{\reals}$ by \cref{metric_on}.
        Thus $\dom{d} = X\times X$ by \cref{funs_elim}.
        Thus $p_1\in\dom{d}$ and $p_2\in\dom{d}$ by \cref{subseteq}.
        Thus $d(p_1)\in\reals$ and $d(p_2)\in\reals$ by \cref{funs_elim}.
        $1\in\reals$ by \cref{real_naturals_subset_reals,real_one_in_naturals,subseteq}.
        \begin{byCase}
        \caseOf{$d(p_1) < 1$.}
            Thus $d(p_2) = d(p_1)$.
            Thus $d(p_2) < 1$ by \cref{real_lt_subst_left}.
            Let $w_1 = (p_1,d(p_1))$.
            Let $w_2 = (p_2,d(p_2))$.
            $w_1\in e$ by \cref{bounded_metric}.
            $w_2\in e$ by \cref{bounded_metric}.
            $e$ is a function by \cref{funs_elim}.
            Thus $e(p_1) = d(p_1)$ by \cref{function_apply_intro}.
            Thus $e(p_2) = d(p_2)$ by \cref{function_apply_intro}.
            Thus $e((x,y)) = e((y,x))$.
        \caseOf{it is wrong that $d(p_1) < 1$.}
            Show that $1 \leq d(p_1)$.
            \begin{subproof}
                Suppose not.
                \begin{byCase}
                \caseOf{$d(p_1) = 1$.}
                    Thus $1 < d(p_1)$ or $1 = d(p_1)$.
                    Thus $1 \leq d(p_1)$.
                \caseOf{$d(p_1) \neq 1$.}
                    Thus $d(p_1) < 1$ or $1 < d(p_1)$ by \cref{real_lt_trichotomy}.
                    \begin{byCase}
                    \caseOf{$d(p_1) < 1$.}
                        Contradiction.
                    \caseOf{$1 < d(p_1)$.}
                        Thus $1 < d(p_1)$ or $1 = d(p_1)$.
                        Thus $1 \leq d(p_1)$.
                    \end{byCase}
                \end{byCase}
            \end{subproof}
            Show that $1 \leq d(p_2)$.
            \begin{subproof}
                Thus $1 < d(p_1)$ or $1 = d(p_1)$ by \cref{real_leq_split}.
                \begin{byCase}
                \caseOf{$1 < d(p_1)$.}
                    Thus $1 < d(p_2)$ by \cref{real_lt_subst_right}.
                    Thus $1 < d(p_2)$ or $1 = d(p_2)$.
                    Thus $1 \leq d(p_2)$.
                \caseOf{$1 = d(p_1)$.}
                    Thus $1 = d(p_2)$.
                    Thus $1 < d(p_2)$ or $1 = d(p_2)$.
                    Thus $1 \leq d(p_2)$.
                \end{byCase}
            \end{subproof}
            Let $w_1 = (p_1,1)$.
            Let $w_2 = (p_2,1)$.
            $w_1\in e$ by \cref{bounded_metric}.
            $w_2\in e$ by \cref{bounded_metric}.
            $e$ is a function by \cref{funs_elim}.
            Thus $e(p_1) = 1$ by \cref{function_apply_intro}.
            Thus $e(p_2) = 1$ by \cref{function_apply_intro}.
            Thus $e((x,y)) = e((y,x))$.
        \end{byCase}
    \end{subproof}
    Thus $e$ is symmetric on $X$ by \cref{metric_symmetric}.
    Show that for all $x,y,z$ such that $x\in X$ and $y\in X$ and $z\in X$ we have
        $e((x,y)) + e((y,z)) \geq e((x,z))$.
    \begin{subproof}
        Fix $x$.
        Fix $y$.
        Fix $z$.
        Assume $x\in X$ and $y\in X$ and $z\in X$.
        Let $p_1 = (x,y)$.
        Let $p_2 = (y,z)$.
        Let $p_3 = (x,z)$.
        $p_1\in X\times X$ by \cref{times_tuple_intro}.
        $p_2\in X\times X$ by \cref{times_tuple_intro}.
        $p_3\in X\times X$ by \cref{times_tuple_intro}.
        $d\in\funs{X\times X}{\reals}$ by \cref{metric_on}.
        Thus $\dom{d} = X\times X$ by \cref{funs_elim}.
        Thus $p_1\in\dom{d}$ and $p_2\in\dom{d}$ and $p_3\in\dom{d}$ by \cref{subseteq}.
        Thus $d(p_1)\in\reals$ and $d(p_2)\in\reals$ and $d(p_3)\in\reals$ by \cref{funs_elim}.
        $1\in\reals$ by \cref{real_naturals_subset_reals,real_one_in_naturals,subseteq}.
        $d$ satisfies the triangle inequality on $X$ by \cref{metric_on}.
        Thus $d(p_3) \leq d(p_1) + d(p_2)$ by \cref{metric_triangle_inequality}.
        \begin{byCase}
        \caseOf{$d(p_1) < 1$.}
            \begin{byCase}
            \caseOf{$d(p_2) < 1$.}
                Let $w_1 = (p_1,d(p_1))$.
                Let $w_2 = (p_2,d(p_2))$.
                $w_1\in e$ by \cref{bounded_metric}.
                $w_2\in e$ by \cref{bounded_metric}.
                $e$ is a function by \cref{funs_elim}.
                Thus $e(p_1) = d(p_1)$ by \cref{function_apply_intro}.
                Thus $e(p_2) = d(p_2)$ by \cref{function_apply_intro}.
                Show that $e(p_3) \leq d(p_3)$.
                \begin{subproof}
                    \begin{byCase}
                    \caseOf{$d(p_3) < 1$.}
                        Let $w_3 = (p_3,d(p_3))$.
                        $w_3\in e$ by \cref{bounded_metric}.
                        Thus $e(p_3) = d(p_3)$ by \cref{function_apply_intro}.
                        Thus $e(p_3) \leq d(p_3)$.
                    \caseOf{it is wrong that $d(p_3) < 1$.}
                        Show that $1 \leq d(p_3)$.
                        \begin{subproof}
                            Suppose not.
                            \begin{byCase}
                            \caseOf{$d(p_3) = 1$.}
                                Thus $1 < d(p_3)$ or $1 = d(p_3)$.
                                Thus $1 \leq d(p_3)$.
                            \caseOf{$d(p_3) \neq 1$.}
                                Thus $d(p_3) < 1$ or $1 < d(p_3)$ by \cref{real_lt_trichotomy}.
                                \begin{byCase}
                                \caseOf{$d(p_3) < 1$.}
                                    Contradiction.
                                \caseOf{$1 < d(p_3)$.}
                                    Thus $1 < d(p_3)$ or $1 = d(p_3)$.
                                    Thus $1 \leq d(p_3)$.
                                \end{byCase}
                            \end{byCase}
                        \end{subproof}
                        Let $w_3 = (p_3,1)$.
                        $w_3\in e$ by \cref{bounded_metric}.
                        Thus $e(p_3) = 1$ by \cref{function_apply_intro}.
                        Thus $e(p_3) \leq d(p_3)$.
                    \end{byCase}
                \end{subproof}
                $e\in\funs{X\times X}{\reals}$.
                Thus $\dom{e} = X\times X$ by \cref{funs_elim}.
                Thus $p_3\in\dom{e}$ by \cref{subseteq}.
                Thus $e(p_3)\in\reals$ by \cref{funs_elim}.
                $d(p_1) + d(p_2)\in\reals$ by \cref{real_add_closed}.
                Thus $e(p_3) \leq d(p_1) + d(p_2)$ by \cref{real_leq_trans}.
                Thus $e(p_1)\in\reals$ and $e(p_2)\in\reals$ by \cref{funs_elim}.
                Thus $e(p_1) + e(p_2) = d(p_1) + e(p_2)$ by \cref{real_add_eq_subst}.
                Thus $d(p_1) + e(p_2) = d(p_1) + d(p_2)$ by \cref{real_add_eq_subst}.
                Thus $e(p_1) + e(p_2) = d(p_1) + d(p_2)$.
                Thus $d(p_1) + d(p_2) \leq e(p_1) + e(p_2)$.
                Thus $e(p_3) \leq e(p_1) + e(p_2)$ by \cref{real_leq_trans}.
            \caseOf{it is wrong that $d(p_2) < 1$.}
                Show that $1 \leq d(p_2)$.
                \begin{subproof}
                    Suppose not.
                    \begin{byCase}
                \caseOf{$d(p_2) = 1$.}
                    Thus $1 < d(p_2)$ or $1 = d(p_2)$.
                    Thus $1 \leq d(p_2)$.
                \caseOf{$d(p_2) \neq 1$.}
                    Thus $d(p_2) < 1$ or $1 < d(p_2)$ by \cref{real_lt_trichotomy}.
                        \begin{byCase}
                        \caseOf{$d(p_2) < 1$.}
                            Contradiction.
                        \caseOf{$1 < d(p_2)$.}
                            Thus $1 < d(p_2)$ or $1 = d(p_2)$.
                            Thus $1 \leq d(p_2)$.
                        \end{byCase}
                    \end{byCase}
                \end{subproof}
                Let $w_1 = (p_1,d(p_1))$.
                Let $w_2 = (p_2,1)$.
                $w_1\in e$ by \cref{bounded_metric}.
                $w_2\in e$ by \cref{bounded_metric}.
                $e$ is a function by \cref{funs_elim}.
                Thus $e(p_1) = d(p_1)$ by \cref{function_apply_intro}.
                Thus $e(p_2) = 1$ by \cref{function_apply_intro}.
                Show that $e(p_3) \leq 1$.
                \begin{subproof}
                    \begin{byCase}
                    \caseOf{$d(p_3) < 1$.}
                        Let $w_3 = (p_3,d(p_3))$.
                        $w_3\in e$ by \cref{bounded_metric}.
                        Thus $e(p_3) = d(p_3)$ by \cref{function_apply_intro}.
                        Thus $e(p_3) < 1$ by \cref{real_lt_subst_left}.
                        Thus $e(p_3) < 1$ or $e(p_3) = 1$.
                        Thus $e(p_3) \leq 1$.
                    \caseOf{it is wrong that $d(p_3) < 1$.}
                        Show that $1 \leq d(p_3)$.
                        \begin{subproof}
                            Suppose not.
                            \begin{byCase}
                            \caseOf{$d(p_3) = 1$.}
                                Thus $1 < d(p_3)$ or $1 = d(p_3)$.
                                Thus $1 \leq d(p_3)$.
                            \caseOf{$d(p_3) \neq 1$.}
                                Thus $d(p_3) < 1$ or $1 < d(p_3)$ by \cref{real_lt_trichotomy}.
                                \begin{byCase}
                                \caseOf{$d(p_3) < 1$.}
                                    Contradiction.
                                \caseOf{$1 < d(p_3)$.}
                                    Thus $1 < d(p_3)$ or $1 = d(p_3)$.
                                    Thus $1 \leq d(p_3)$.
                                \end{byCase}
                            \end{byCase}
                        \end{subproof}
                        Let $w_3 = (p_3,1)$.
                        $w_3\in e$ by \cref{bounded_metric}.
                        Thus $e(p_3) = 1$ by \cref{function_apply_intro}.
                        Thus $e(p_3) \leq 1$.
                    \end{byCase}
                \end{subproof}
                $e$ is nonnegative on $X$ by \cref{metric_nonnegative}.
                Thus $e(p_1) \geq 0$ by \cref{metric_nonnegative}.
                Thus $0 \leq e(p_1)$.
                $e\in\funs{X\times X}{\reals}$.
                Thus $e(p_1)\in\reals$ by \cref{funs_elim}.
                Thus $e(p_2)\in\reals$ and $e(p_3)\in\reals$ by \cref{funs_elim}.
                Thus $e(p_1) + e(p_2)\in\reals$ by \cref{real_add_closed}.
                $0\in\reals$ by \cref{real_add_ident}.
                $1\in\reals$ by \cref{real_naturals_subset_reals,real_one_in_naturals,subseteq}.
                Thus $0 + 1 \leq e(p_1) + 1$ by \cref{real_leq_add}.
                $0 + 1 = 1$ by \cref{real_add_ident}.
                $e(p_1) + 1 = 1 + e(p_1)$ by \cref{real_add_comm}.
                Thus $1 \leq e(p_1) + e(p_2)$ by \cref{real_leq_trans}.
                Thus $e(p_3) \leq e(p_1) + e(p_2)$ by \cref{real_leq_trans}.
            \end{byCase}
        \caseOf{it is wrong that $d(p_1) < 1$.}
            Show that $1 \leq d(p_1)$.
            \begin{subproof}
                Suppose not.
                \begin{byCase}
                \caseOf{$d(p_1) = 1$.}
                    Thus $1 < d(p_1)$ or $1 = d(p_1)$.
                    Thus $1 \leq d(p_1)$.
                \caseOf{$d(p_1) \neq 1$.}
                    Thus $d(p_1) < 1$ or $1 < d(p_1)$ by \cref{real_lt_trichotomy}.
                    \begin{byCase}
                    \caseOf{$d(p_1) < 1$.}
                        Contradiction.
                    \caseOf{$1 < d(p_1)$.}
                        Thus $1 < d(p_1)$ or $1 = d(p_1)$.
                        Thus $1 \leq d(p_1)$.
                    \end{byCase}
                \end{byCase}
            \end{subproof}
            Let $w_1 = (p_1,1)$.
            $w_1\in e$ by \cref{bounded_metric}.
            $e$ is a function by \cref{funs_elim}.
            Thus $e(p_1) = 1$ by \cref{function_apply_intro}.
            Show that $e(p_3) \leq 1$.
            \begin{subproof}
                \begin{byCase}
                \caseOf{$d(p_3) < 1$.}
                    Let $w_3 = (p_3,d(p_3))$.
                    $w_3\in e$ by \cref{bounded_metric}.
                    Thus $e(p_3) = d(p_3)$ by \cref{function_apply_intro}.
                    Thus $e(p_3) < 1$ by \cref{real_lt_subst_left}.
                    Thus $e(p_3) < 1$ or $e(p_3) = 1$.
                    Thus $e(p_3) \leq 1$.
                \caseOf{it is wrong that $d(p_3) < 1$.}
                    Show that $1 \leq d(p_3)$.
                    \begin{subproof}
                        Suppose not.
                        \begin{byCase}
                        \caseOf{$d(p_3) = 1$.}
                            Thus $1 < d(p_3)$ or $1 = d(p_3)$.
                            Thus $1 \leq d(p_3)$.
                        \caseOf{$d(p_3) \neq 1$.}
                            Thus $d(p_3) < 1$ or $1 < d(p_3)$ by \cref{real_lt_trichotomy}.
                            \begin{byCase}
                            \caseOf{$d(p_3) < 1$.}
                                Contradiction.
                            \caseOf{$1 < d(p_3)$.}
                                Thus $1 < d(p_3)$ or $1 = d(p_3)$.
                                Thus $1 \leq d(p_3)$.
                            \end{byCase}
                        \end{byCase}
                    \end{subproof}
                    Let $w_3 = (p_3,1)$.
                    $w_3\in e$ by \cref{bounded_metric}.
                    Thus $e(p_3) = 1$ by \cref{function_apply_intro}.
                    Thus $e(p_3) \leq 1$.
                \end{byCase}
            \end{subproof}
            $e$ is nonnegative on $X$ by \cref{metric_nonnegative}.
            Thus $e(p_2) \geq 0$ by \cref{metric_nonnegative}.
            Thus $0 \leq e(p_2)$.
            $e\in\funs{X\times X}{\reals}$.
            Thus $e(p_1)\in\reals$ and $e(p_2)\in\reals$ and $e(p_3)\in\reals$ by \cref{funs_elim}.
            Thus $e(p_1) + e(p_2)\in\reals$ by \cref{real_add_closed}.
            $0\in\reals$ by \cref{real_add_ident}.
            $1\in\reals$ by \cref{real_naturals_subset_reals,real_one_in_naturals,subseteq}.
            Thus $0 + 1 \leq e(p_2) + 1$ by \cref{real_leq_add}.
            $0 + 1 = 1$ by \cref{real_add_ident}.
            $e(p_2) + 1 = 1 + e(p_2)$ by \cref{real_add_comm}.
            Thus $1 \leq e(p_1) + e(p_2)$ by \cref{real_leq_trans}.
            Thus $e(p_3) \leq e(p_1) + e(p_2)$ by \cref{real_leq_trans}.
        \end{byCase}
    \end{subproof}
    Thus $e$ satisfies the triangle inequality on $X$ by \cref{metric_triangle_inequality}.
    Thus $e$ is a metric on $X$ by \cref{metric_on}.
\end{proof}

% from §20 Item 19: standard bounded metric induces the same topology
\begin{theorem}\label{bounded_metric_topology}
    Assume $d$ is a metric on $X$.
    Then $\metricTopology{\boundedMetric{d}}{X} = \metricTopology{d}{X}$.
\end{theorem}
\begin{proof}
    Let $e = \boundedMetric{d}$.
    $e$ is a metric on $X$ by \cref{bounded_metric_is_metric}.
    Let $T = \metricTopology{d}{X}$.
    Let $T_1 = \metricTopology{e}{X}$.
    Show that for all $U$ such that $U\in T_1$ we have $U\in T$.
    \begin{subproof}
        Fix $U$.
        Assume $U\in T_1$.
        Let $B_1 = \metricBasis{e}{X}$.
        $B_1$ is a basis on $X$ by \cref{metric_basis_is_basis}.
        $T_1 = \basisTopology{B_1}{X}$ by \cref{metric_topology}.
        Thus $T_1$ is a topology on $X$ by \cref{basis_topology_is_topology,basis_for_topology}.
        $U$ is open in $X$ with respect to $T_1$ by \cref{open_in}.
        Thus $T_1\subseteq\pow{X}$ by \cref{topology_on}.
        Thus $U\in\pow{X}$ by \cref{subseteq}.
        Thus $U\subseteq X$ by \cref{powerset_elim}.
        Show that $U$ is open in $X$ with respect to $T$.
        \begin{subproof}
            Show that for all $y\in U$ there exists $\delta\in\reals$ such that $0 < \delta$ and $\metricBall{d}{X}{y}{\delta}\subseteq U$.
            \begin{subproof}
                Fix $y$.
                Assume $y\in U$.
                There exists $\epsilon\in\reals$ such that $0 < \epsilon$ and $\metricBall{e}{X}{y}{\epsilon}\subseteq U$
                    by \cref{metric_open_iff}.
                Take $\epsilon$ such that $\epsilon\in\reals$ and $0 < \epsilon$ and $\metricBall{e}{X}{y}{\epsilon}\subseteq U$.
                $1\in\reals$ by \cref{real_naturals_subset_reals,real_one_in_naturals,subseteq}.
                Thus $\epsilon < 1$ or $\epsilon = 1$ or $1 < \epsilon$ by \cref{real_lt_trichotomy}.
                \begin{byCase}
                \caseOf{$\epsilon < 1$.}
                    Let $\delta = \epsilon$.
                    Show that $\metricBall{d}{X}{y}{\delta}\subseteq U$.
                    \begin{subproof}
                        Show that for all $z$ such that $z\in\metricBall{d}{X}{y}{\delta}$ we have $z\in U$.
                        \begin{subproof}
                            Fix $z$.
                            Assume $z\in\metricBall{d}{X}{y}{\delta}$.
                            Thus $z\in X$ and $d((y,z)) < \delta$ by \cref{metric_ball}.
                            Thus $y\in X$ by \cref{subseteq}.
                            $(y,z)\in X\times X$ by \cref{times_tuple_intro}.
                            $d\in\funs{X\times X}{\reals}$ by \cref{metric_on}.
                            Thus $d((y,z))\in\reals$ by \cref{funs_type_apply}.
                            Thus $d((y,z)) < 1$ by \cref{real_lt_trans}.
                            Let $p = (y,z)$.
                            Thus $y\in X$ by \cref{subseteq}.
                            $p\in X\times X$ by \cref{times_tuple_intro}.
                            $d\in\funs{X\times X}{\reals}$ by \cref{metric_on}.
                            Thus $p\in\dom{d}$ by \cref{funs_elim,subseteq}.
                            Thus $d(p)\in\reals$ by \cref{funs_type_apply}.
                            Thus $d(p) = d((y,z))$.
                            Thus $d(p) < 1$ by \cref{real_lt_subst_left}.
                            Thus $d(p)\in\reals$ by \cref{funs_type_apply}.
                            Thus $d(p) < 1$.
                            Thus $d(p)\in\reals$ by \cref{funs_type_apply}.
                            Thus $d(p) < 1$.
                            $e\in\funs{X\times X}{\reals}$ by \cref{metric_on}.
                            Thus $\dom{e} = X\times X$ by \cref{funs_elim}.
                            Thus $p\in\dom{e}$ by \cref{subseteq}.
                            $e$ is a function by \cref{funs_elim}.
                            Thus $(p,e(p))\in e$ by \cref{function_apply_iff}.
                            Take $p_1,w$ such that
                                $p_1\in\dom{d}$ and $w\in\reals$ and
                                $(p,e(p)) = (p_1,w)$ and
                                $((w = d(p_1) \land d(p_1) < 1) \lor (w = 1 \land 1 \leq d(p_1)))$
                                by \cref{bounded_metric}.
                            Thus $p = p_1$ by \cref{pair_eq_iff}.
                            Show that $e(p) = d(p)$.
                            \begin{subproof}
                                Suppose not.
                                Thus $e(p) \neq d(p)$.
                                Thus $(w = d(p) \land d(p) < 1) \lor (w = 1 \land 1 \leq d(p))$.
                                \begin{byCase}
                                \caseOf{$w = d(p) \land d(p) < 1$.}
                                    Thus $e(p) = d(p)$.
                                    Contradiction.
                                \caseOf{$w = 1 \land 1 \leq d(p)$.}
                                    Thus $1 \leq d(p)$.
                                    Thus $1 < d(p)$ or $1 = d(p)$.
                                    \begin{byCase}
                                    \caseOf{$1 < d(p)$.}
                                        Thus $d(p) = d((y,z))$.
                                        Thus $d(p) < 1$ by \cref{real_lt_subst_left}.
                                        Thus $1 < 1$ by \cref{real_lt_trans}.
                                        Contradiction by \cref{real_lt_irreflexive}.
                                    \caseOf{$1 = d(p)$.}
                                        Thus $1 < 1$ by \cref{real_lt_subst_right}.
                                        Contradiction by \cref{real_lt_irreflexive}.
                                    \end{byCase}
                                \end{byCase}
                            \end{subproof}
                            Thus $e((y,z)) < \delta$ by \cref{real_lt_subst_left}.
                            Thus $z\in\metricBall{e}{X}{y}{\delta}$ by \cref{metric_ball}.
                            Thus $z\in U$ by \cref{subseteq}.
                        \end{subproof}
                        Thus $\metricBall{d}{X}{y}{\delta}\subseteq U$ by \cref{subseteq}.
                    \end{subproof}
                \caseOf{$\epsilon = 1$.}
                    Let $\delta = \epsilon$.
                    Show that $\metricBall{d}{X}{y}{\delta}\subseteq U$.
                    \begin{subproof}
                        Show that for all $z$ such that $z\in\metricBall{d}{X}{y}{\delta}$ we have $z\in U$.
                        \begin{subproof}
                            Fix $z$.
                            Assume $z\in\metricBall{d}{X}{y}{\delta}$.
                            Thus $z\in X$ and $d((y,z)) < \delta$ by \cref{metric_ball}.
                            Thus $d((y,z)) < 1$ by \cref{real_lt_subst_right}.
                            Let $p = (y,z)$.
                            Thus $y\in X$ by \cref{subseteq}.
                            $p\in X\times X$ by \cref{times_tuple_intro}.
                            $d\in\funs{X\times X}{\reals}$ by \cref{metric_on}.
                            Thus $p\in\dom{d}$ by \cref{funs_elim,subseteq}.
                            Thus $d(p)\in\reals$ by \cref{funs_type_apply}.
                            Thus $d(p) = d((y,z))$.
                            Thus $d(p) < 1$ by \cref{real_lt_subst_left}.
                            $e\in\funs{X\times X}{\reals}$ by \cref{metric_on}.
                            Thus $\dom{e} = X\times X$ by \cref{funs_elim}.
                            Thus $p\in\dom{e}$ by \cref{subseteq}.
                            $e$ is a function by \cref{funs_elim}.
                            Thus $(p,e(p))\in e$ by \cref{function_apply_iff}.
                            Take $p_1,w$ such that
                                $p_1\in\dom{d}$ and $w\in\reals$ and
                                $(p,e(p)) = (p_1,w)$ and
                                $((w = d(p_1) \land d(p_1) < 1) \lor (w = 1 \land 1 \leq d(p_1)))$
                                by \cref{bounded_metric}.
                            Thus $p = p_1$ by \cref{pair_eq_iff}.
                            Show that $e(p) = d(p)$.
                            \begin{subproof}
                                Suppose not.
                                Thus $e(p) \neq d(p)$.
                                Thus $(w = d(p) \land d(p) < 1) \lor (w = 1 \land 1 \leq d(p))$.
                                \begin{byCase}
                                \caseOf{$w = d(p) \land d(p) < 1$.}
                                    Thus $e(p) = d(p)$.
                                    Contradiction.
                                \caseOf{$w = 1 \land 1 \leq d(p)$.}
                                    Thus $1 \leq d(p)$.
                                    Thus $1 < d(p)$ or $1 = d(p)$.
                                    \begin{byCase}
                                    \caseOf{$1 < d(p)$.}
                                        Thus $d(p) = d((y,z))$.
                                        Thus $d(p) < 1$ by \cref{real_lt_subst_left}.
                                        Thus $1 < 1$ by \cref{real_lt_trans}.
                                        Contradiction by \cref{real_lt_irreflexive}.
                                    \caseOf{$1 = d(p)$.}
                                        Thus $1 < 1$ by \cref{real_lt_subst_right}.
                                        Contradiction by \cref{real_lt_irreflexive}.
                                    \end{byCase}
                                \end{byCase}
                            \end{subproof}
                            Thus $e((y,z)) < \delta$ by \cref{real_lt_subst_left}.
                            Thus $z\in\metricBall{e}{X}{y}{\delta}$ by \cref{metric_ball}.
                            Thus $z\in U$ by \cref{subseteq}.
                        \end{subproof}
                        Thus $\metricBall{d}{X}{y}{\delta}\subseteq U$ by \cref{subseteq}.
                    \end{subproof}
                \caseOf{$1 < \epsilon$.}
                    Let $\delta = 1$.
                    Show that $\metricBall{d}{X}{y}{\delta}\subseteq U$.
                    \begin{subproof}
                        Show that for all $z$ such that $z\in\metricBall{d}{X}{y}{\delta}$ we have $z\in U$.
                        \begin{subproof}
                            Fix $z$.
                            Assume $z\in\metricBall{d}{X}{y}{\delta}$.
                            Thus $z\in X$ and $d((y,z)) < \delta$ by \cref{metric_ball}.
                            Thus $d((y,z)) < 1$ by \cref{real_lt_subst_right}.
                            Let $p = (y,z)$.
                            Thus $y\in X$ by \cref{subseteq}.
                            $p\in X\times X$ by \cref{times_tuple_intro}.
                            $d\in\funs{X\times X}{\reals}$ by \cref{metric_on}.
                            Thus $p\in\dom{d}$ by \cref{funs_elim,subseteq}.
                            Thus $d(p)\in\reals$ by \cref{funs_type_apply}.
                            Thus $d(p) = d((y,z))$.
                            Thus $d(p) < 1$ by \cref{real_lt_subst_left}.
                            $e\in\funs{X\times X}{\reals}$ by \cref{metric_on}.
                            Thus $\dom{e} = X\times X$ by \cref{funs_elim}.
                            Thus $p\in\dom{e}$ by \cref{subseteq}.
                            $e$ is a function by \cref{funs_elim}.
                            Thus $(p,e(p))\in e$ by \cref{function_apply_iff}.
                            Take $p_1,w$ such that
                                $p_1\in\dom{d}$ and $w\in\reals$ and
                                $(p,e(p)) = (p_1,w)$ and
                                $((w = d(p_1) \land d(p_1) < 1) \lor (w = 1 \land 1 \leq d(p_1)))$
                                by \cref{bounded_metric}.
                            Thus $p = p_1$ by \cref{pair_eq_iff}.
                            Show that $e(p) = d(p)$.
                            \begin{subproof}
                                Suppose not.
                                Thus $e(p) \neq d(p)$.
                                Thus $(w = d(p) \land d(p) < 1) \lor (w = 1 \land 1 \leq d(p))$.
                                \begin{byCase}
                                \caseOf{$w = d(p) \land d(p) < 1$.}
                                    Thus $e(p) = d(p)$.
                                    Contradiction.
                                \caseOf{$w = 1 \land 1 \leq d(p)$.}
                                    Thus $1 \leq d(p)$.
                                    Thus $1 < d(p)$ or $1 = d(p)$.
                                    \begin{byCase}
                                    \caseOf{$1 < d(p)$.}
                                        Thus $d(p) = d((y,z))$.
                                        Thus $d(p) < 1$ by \cref{real_lt_subst_left}.
                                        Thus $1 < 1$ by \cref{real_lt_trans}.
                                        Contradiction by \cref{real_lt_irreflexive}.
                                    \caseOf{$1 = d(p)$.}
                                        Thus $1 < 1$ by \cref{real_lt_subst_right}.
                                        Contradiction by \cref{real_lt_irreflexive}.
                                    \end{byCase}
                                \end{byCase}
                            \end{subproof}
                            Thus $e((y,z)) < 1$ by \cref{real_lt_subst_left}.
                            Thus $e((y,z)) < \epsilon$ by \cref{real_lt_trans}.
                            Thus $z\in\metricBall{e}{X}{y}{\epsilon}$ by \cref{metric_ball}.
                            Thus $z\in U$ by \cref{subseteq}.
                        \end{subproof}
                        Thus $\metricBall{d}{X}{y}{\delta}\subseteq U$ by \cref{subseteq}.
                    \end{subproof}
                \end{byCase}
            \end{subproof}
            Thus $U$ is open in $X$ with respect to $T$ by \cref{metric_open_iff}.
        \end{subproof}
        Thus $U\in T$ by \cref{open_in}.
    \end{subproof}
    Thus $T_1\subseteq T$ by \cref{subseteq}.
    Show that for all $U$ such that $U\in T$ we have $U\in T_1$.
    \begin{subproof}
        Fix $U$.
        Assume $U\in T$.
        Let $B = \metricBasis{d}{X}$.
        $B$ is a basis on $X$ by \cref{metric_basis_is_basis}.
        $T = \basisTopology{B}{X}$ by \cref{metric_topology}.
        Thus $T$ is a topology on $X$ by \cref{basis_topology_is_topology,basis_for_topology}.
        $U$ is open in $X$ with respect to $T$ by \cref{open_in}.
        Thus $T\subseteq\pow{X}$ by \cref{topology_on}.
        Thus $U\in\pow{X}$ by \cref{subseteq}.
        Thus $U\subseteq X$ by \cref{powerset_elim}.
        Show that $U$ is open in $X$ with respect to $T_1$.
        \begin{subproof}
            Show that for all $y\in U$ there exists $\delta\in\reals$ such that $0 < \delta$ and $\metricBall{e}{X}{y}{\delta}\subseteq U$.
            \begin{subproof}
                Fix $y$.
                Assume $y\in U$.
                There exists $\epsilon\in\reals$ such that $0 < \epsilon$ and $\metricBall{d}{X}{y}{\epsilon}\subseteq U$
                    by \cref{metric_open_iff}.
                Take $\epsilon$ such that $\epsilon\in\reals$ and $0 < \epsilon$ and $\metricBall{d}{X}{y}{\epsilon}\subseteq U$.
                $1\in\reals$ by \cref{real_naturals_subset_reals,real_one_in_naturals,subseteq}.
                Let $A = \{\epsilon,1\}$.
                $A$ is finite by \cref{pair_is_finite}.
                $\epsilon\in A$ by \cref{upair_intro_left}.
                $1\in A$ by \cref{upair_intro_right}.
                Thus $A\neq\emptyset$ by \cref{emptyset}.
                Show that for all $t$ such that $t\in A$ we have $t\in\reals$.
                \begin{subproof}
                    Fix $t$.
                    Assume $t\in A$.
                    $t = \epsilon$ or $t = 1$ by \cref{upair_elim}.
                    Thus $t\in\reals$.
                \end{subproof}
                Thus $A\subseteq\reals$ by \cref{subseteq}.
                Take $m\in A$ such that $m$ is a minimum of $A$ by \cref{finite_real_minimum}.
                Thus $m\in\reals$ by \cref{subseteq}.
                Show that $0 < m$.
                \begin{subproof}
                    $m = \epsilon$ or $m = 1$ by \cref{upair_elim}.
                    \begin{byCase}
                    \caseOf{$m = \epsilon$.} Thus $0 < m$.
                    \caseOf{$m = 1$.} $0 < 1$ by \cref{real_zero_lt_one}. Thus $0 < m$.
                    \end{byCase}
                \end{subproof}
                Show that $\metricBall{e}{X}{y}{m}\subseteq U$.
                \begin{subproof}
                    Show that for all $z$ such that $z\in\metricBall{e}{X}{y}{m}$ we have $z\in U$.
                    \begin{subproof}
                        Fix $z$.
                        Assume $z\in\metricBall{e}{X}{y}{m}$.
                        Thus $z\in X$ and $e((y,z)) < m$ by \cref{metric_ball}.
                        Thus $y\in X$ by \cref{subseteq}.
                        Let $p = (y,z)$.
                        $p\in X\times X$ by \cref{times_tuple_intro}.
                        $d\in\funs{X\times X}{\reals}$ by \cref{metric_on}.
                        Thus $p\in\dom{d}$ by \cref{funs_elim,subseteq}.
                        $e\in\funs{X\times X}{\reals}$ by \cref{metric_on}.
                        Thus $\dom{e} = X\times X$ by \cref{funs_elim}.
                        Thus $p\in\dom{e}$ by \cref{subseteq}.
                        Thus $e(p)\in\reals$ by \cref{funs_type_apply}.
                        Thus $e(p) = e((y,z))$.
                        Thus $e(p) < m$ by \cref{real_lt_subst_left}.
                        Thus $m \leq \epsilon$ by \cref{real_minimum}.
                        Show that $e(p) < \epsilon$.
                        \begin{subproof}
                            Thus $m < \epsilon$ or $m = \epsilon$.
                            \begin{byCase}
                            \caseOf{$m < \epsilon$.}
                                Thus $e(p) < \epsilon$ by \cref{real_lt_trans}.
                            \caseOf{$m = \epsilon$.}
                                Thus $e(p) < \epsilon$ by \cref{real_lt_subst_right}.
                            \end{byCase}
                        \end{subproof}
                        $e$ is a function by \cref{funs_elim}.
                        Thus $(p,e(p))\in e$ by \cref{function_apply_iff}.
                        Take $p_1,w$ such that
                            $p_1\in\dom{d}$ and $w\in\reals$ and
                            $(p,e(p)) = (p_1,w)$ and
                            $((w = d(p_1) \land d(p_1) < 1) \lor (w = 1 \land 1 \leq d(p_1)))$
                            by \cref{bounded_metric}.
                        Thus $p = p_1$ by \cref{pair_eq_iff}.
                        Show that $e(p) = d(p)$.
                        \begin{subproof}
                            Suppose not.
                            Thus $e(p) \neq d(p)$.
                            Show that $e(p) < 1$.
                            \begin{subproof}
                                Thus $m \leq 1$ by \cref{real_minimum}.
                                Thus $m < 1$ or $m = 1$.
                                \begin{byCase}
                                \caseOf{$m < 1$.}
                                    Thus $e(p) < 1$ by \cref{real_lt_trans}.
                                \caseOf{$m = 1$.}
                                    Thus $e(p) < 1$ by \cref{real_lt_subst_right}.
                                \end{byCase}
                            \end{subproof}
                            Thus $(w = d(p) \land d(p) < 1) \lor (w = 1 \land 1 \leq d(p))$.
                            \begin{byCase}
                            \caseOf{$w = d(p) \land d(p) < 1$.}
                                Thus $e(p) = d(p)$.
                                Contradiction.
                            \caseOf{$w = 1 \land 1 \leq d(p)$.}
                                Thus $e(p) = 1$.
                                Thus $1 < 1$ by \cref{real_lt_subst_left}.
                                Contradiction by \cref{real_lt_irreflexive}.
                            \end{byCase}
                        \end{subproof}
                        Thus $d(p) < \epsilon$ by \cref{real_lt_subst_left}.
                        Thus $z\in\metricBall{d}{X}{y}{\epsilon}$ by \cref{metric_ball}.
                        Thus $z\in U$ by \cref{subseteq}.
                    \end{subproof}
                    Thus $\metricBall{e}{X}{y}{m}\subseteq U$ by \cref{subseteq}.
                \end{subproof}
            \end{subproof}
            Thus $U$ is open in $X$ with respect to $T_1$ by \cref{metric_open_iff}.
        \end{subproof}
        Thus $U\in T_1$ by \cref{open_in}.
    \end{subproof}
    Thus $T\subseteq T_1$ by \cref{subseteq}.
    Thus $T_1 = T$ by \cref{subseteq_antisymmetric}.
\end{proof}

% from §20 Item 20: $\reals^n$
\begin{definition}\label{reals_power}
    $\realsPower{J} = \tuplePower{\reals}{J}$.
\end{definition}

% from §20 Item 21: euclidean norm
\begin{definition}\label{euclidean_norm}
    $r$ is a euclidean norm of $x$ with respect to $n$ iff
    $x\in\realsPower{\initialSegment{n}}$ and
    $r\in\reals$ and
    there exists $a\in\funs{\initialSegment{n}}{\reals}$ such that
        for all $i\in\initialSegment{n}$ we have $a(i) = \exp{x(i)}{1+1}$ and
        there exists $s$ such that $s$ is the sum of $a$ and $r = \sqrt{s}$.
\end{definition}

% from §20 helper lemma 6: sum of a finite sequence predicate
\begin{definition}\label{sum_finite_sequence_pred}
    $\sumFiniteSequence{a}{s}$ iff
    there exists $n\in\naturalsPlus$ such that
    $a\in\funs{\initialSegment{n}}{\reals}$ and
    there exists $t\in\funs{\initialSegment{n}}{\reals}$ such that
    $t(1) = a(1)$ and
    $s = t(n)$ and
    for all $k\in\naturalsPlus$ such that $k + 1 \leq n$ we have
    $t(k + 1) = t(k) + a(k + 1)$.
\end{definition}

% from §20 helper lemma 7: sum of a finite sequence is unique
\begin{lemma}\label{sum_finite_sequence_unique}
    Suppose $\sumFiniteSequence{a}{s_1}$.
    Suppose $\sumFiniteSequence{a}{s_2}$.
    Then $s_1 = s_2$.
\end{lemma}
\begin{proof}
    Take $n_1$ such that $n_1\in\naturalsPlus$ and $a\in\funs{\initialSegment{n_1}}{\reals}$ and
        there exists $t_1\in\funs{\initialSegment{n_1}}{\reals}$ such that
            $t_1(1) = a(1)$ and
            $s_1 = t_1(n_1)$ and
            for all $k\in\naturalsPlus$ such that $k + 1 \leq n_1$ we have $t_1(k + 1) = t_1(k) + a(k + 1)$
        by \cref{sum_finite_sequence_pred}.
    Take $n_2$ such that $n_2\in\naturalsPlus$ and $a\in\funs{\initialSegment{n_2}}{\reals}$ and
        there exists $t_2\in\funs{\initialSegment{n_2}}{\reals}$ such that
            $t_2(1) = a(1)$ and
            $s_2 = t_2(n_2)$ and
            for all $k\in\naturalsPlus$ such that $k + 1 \leq n_2$ we have $t_2(k + 1) = t_2(k) + a(k + 1)$
        by \cref{sum_finite_sequence_pred}.
    Take $t_1$ such that $t_1\in\funs{\initialSegment{n_1}}{\reals}$ and
        $t_1(1) = a(1)$ and
        $s_1 = t_1(n_1)$ and
        for all $k\in\naturalsPlus$ such that $k + 1 \leq n_1$ we have $t_1(k + 1) = t_1(k) + a(k + 1)$.
    Take $t_2$ such that $t_2\in\funs{\initialSegment{n_2}}{\reals}$ and
        $t_2(1) = a(1)$ and
        $s_2 = t_2(n_2)$ and
        for all $k\in\naturalsPlus$ such that $k + 1 \leq n_2$ we have $t_2(k + 1) = t_2(k) + a(k + 1)$.
    Thus $\dom{a} = \initialSegment{n_1}$ by \cref{funs_elim}.
    Thus $\dom{a} = \initialSegment{n_2}$ by \cref{funs_elim}.
    Thus $\initialSegment{n_1} = \initialSegment{n_2}$ by \cref{subseteq_antisymmetric,subseteq}.
    $n_1\in\naturalsPlus$ and $n_2\in\naturalsPlus$.
    Thus $n_1 \leq n_1$ and $n_2 \leq n_2$.
    Thus $n_1\in\initialSegment{n_1}$ and $n_2\in\initialSegment{n_2}$.
    Thus $n_1\in\initialSegment{n_2}$ and $n_2\in\initialSegment{n_1}$ by \cref{subseteq}.
    Thus $n_1 \leq n_2$ and $n_2 \leq n_1$.
    $n_1\in\reals$ and $n_2\in\reals$ by \cref{real_naturals_subset_reals,subseteq}.
    Thus $n_1 = n_2$ by \cref{real_leq_antisym}.
    Let $n = n_1$.
    Thus $t_1\in\funs{\initialSegment{n}}{\reals}$ and $t_2\in\funs{\initialSegment{n}}{\reals}$.
    Let $A = \{k\in\naturalsPlus \mid \text{if $k \leq n$ then $t_1(k) = t_2(k)$}\}$.
    $A\subseteq\naturalsPlus$ by \cref{subseteq}.
    Show that $1\in A$.
    \begin{subproof}
        $1\in\naturalsPlus$ by \cref{real_one_in_naturals}.
        $1 \leq n$ by \cref{real_one_leq_naturalsplus}.
        Thus $t_1(1) = a(1)$ and $t_2(1) = a(1)$.
        Thus $t_1(1) = t_2(1)$.
        Thus $1\in A$.
    \end{subproof}
    Show that for all $k\in A$ we have $k + 1\in A$.
    \begin{subproof}
        Fix $k$.
        Assume $k\in A$.
        Thus $k\in\naturalsPlus$ by \cref{subseteq}.
        Show that $k + 1\in A$.
        \begin{subproof}
            Show that if $k + 1 \leq n$ then $t_1(k + 1) = t_2(k + 1)$.
            \begin{subproof}
                Assume $k + 1 \leq n$.
                $k\in\reals$ and $n\in\reals$ by \cref{real_naturals_subset_reals,subseteq}.
                $0\in\reals$ by \cref{real_add_ident}.
                $1\in\reals$ by \cref{real_naturals_subset_reals,real_one_in_naturals,subseteq}.
                $0 < 1$ by \cref{real_zero_lt_one}.
                $0 + k < 1 + k$ by \cref{real_lt_add}.
                $0 + k = k + 0$ by \cref{real_add_comm}.
                $1 + k = k + 1$ by \cref{real_add_comm}.
                Thus $k + 0 < k + 1$ by \cref{real_lt_subst_left}.
                $k + 0 = k$ by \cref{real_add_ident}.
                Thus $k < k + 1$ by \cref{real_lt_subst_left}.
                Thus $k < k + 1$ or $k = k + 1$.
                Thus $k \leq k + 1$.
                $k + 1\in\reals$ by \cref{real_add_closed}.
                Thus $k \leq n$ by \cref{real_leq_trans}.
                Thus $k\in\initialSegment{n}$ by \cref{initial_segment_naturalsplus}.
                Thus $t_1(k) = t_2(k)$.
                $k + 1\in\naturalsPlus$ by \cref{real_naturals_closed_succ}.
                $a\in\funs{\initialSegment{n}}{\reals}$.
                Thus $k + 1\in\initialSegment{n}$.
                Thus $a(k + 1)\in\reals$ by \cref{funs_type_apply}.
                Thus $t_1(k)\in\reals$ and $t_2(k)\in\reals$ by \cref{funs_type_apply}.
                Thus $t_1(k + 1) = t_1(k) + a(k + 1)$.
                Thus $t_2(k + 1) = t_2(k) + a(k + 1)$.
                Thus $t_1(k + 1) = t_2(k + 1)$ by \cref{real_add_eq_subst}.
            \end{subproof}
            Thus $k + 1\in A$.
        \end{subproof}
    \end{subproof}
    Thus for all $k\in\naturalsPlus$ we have $k\in A$ by \cref{real_naturals_induction}.
    Show that for all $k\in\naturalsPlus$ we have if $k \leq n$ then $t_1(k) = t_2(k)$.
    \begin{subproof}
        Fix $k\in\naturalsPlus$.
        Thus $k\in A$.
        Thus if $k \leq n$ then $t_1(k) = t_2(k)$.
    \end{subproof}
    Thus $\dom{t_1} = \initialSegment{n}$ and $\dom{t_2} = \initialSegment{n}$ by \cref{funs_elim}.
    Thus $t_1$ is a function and $t_2$ is a function by \cref{funs_elim}.
    Thus $\dom{t_1} = \dom{t_2}$.
    Thus $\dom{t_1}\subseteq \dom{t_2}$ and $\dom{t_2}\subseteq \dom{t_1}$.
    Show that for all $x\in\dom{t_1}$ we have $t_1(x) = t_2(x)$.
    \begin{subproof}
        Fix $x\in\dom{t_1}$.
        Thus $x\in\initialSegment{n}$.
        Thus $x\in\naturalsPlus$ and $x \leq n$ by \cref{initial_segment_naturalsplus}.
        Thus $t_1(x) = t_2(x)$.
    \end{subproof}
    Thus $t_1\subseteq t_2$ by \cref{fun_subseteq}.
    Show that for all $x\in\dom{t_2}$ we have $t_1(x) = t_2(x)$.
    \begin{subproof}
        Fix $x\in\dom{t_2}$.
        Thus $x\in\initialSegment{n}$.
        Thus $x\in\naturalsPlus$ and $x \leq n$ by \cref{initial_segment_naturalsplus}.
        Thus $t_1(x) = t_2(x)$.
    \end{subproof}
    Thus $t_2\subseteq t_1$ by \cref{fun_subseteq}.
    Thus $t_1 = t_2$ by \cref{subseteq_antisymmetric}.
    Thus $s_1 = t_1(n_1)$.
    Thus $s_2 = t_2(n_2)$.
    Thus $t_1(n_1) = t_2(n_1)$.
    Thus $t_2(n_1) = t_2(n_2)$.
    Thus $s_1 = s_2$.
\end{proof}

% from §20 helper lemma 8: squared coordinate sequence
\begin{definition}\label{coord_square_sequence}
    $\coordSquareSequence{n}{p} = \{w \mid \exists i\in\initialSegment{n}.
        w = (i, \exp{\apply{\fst{p}}{i} - \apply{\snd{p}}{i}}{1+1})\}$.
\end{definition}

% from §20 Item 22: euclidean metric
\begin{definition}\label{euclidean_metric}
    $\euclideanMetric{n} = \{w \mid \exists p\in\realsPower{\initialSegment{n}}\times\realsPower{\initialSegment{n}}. \exists r\in\reals.
        \exists s\in\reals.
        w = (p,r) \land
        \sumFiniteSequence{\coordSquareSequence{n}{p}}{s} \land
        r = \sqrt{s}\}$.
\end{definition}

% from §20 helper lemma 9: coordinate difference set
\begin{definition}\label{coord_diff_set}
    $\coordDiffSet{n}{x}{y} = \{\abs{x(i) - y(i)} \mid i\in\initialSegment{n}\}$.
\end{definition}

% from §20 Item 23: square metric
\begin{definition}\label{square_metric}
    $\squareMetric{n} = \{w \mid \exists p\in\realsPower{\initialSegment{n}}\times\realsPower{\initialSegment{n}}.
        \exists x\in\realsPower{\initialSegment{n}}. \exists y\in\realsPower{\initialSegment{n}}. \exists r\in\reals.
        w = (p,r) \land p = (x,y) \land
        r\in\coordDiffSet{n}{x}{y} \land
        (\forall t\in\coordDiffSet{n}{x}{y}. t \leq r)\}$.
\end{definition}

% from §20 Item 24: comparison of metric topologies
\begin{lemma}\label{metric_topology_finer_iff}
    Suppose $d$ is a metric on $X$.
    Suppose $d_1$ is a metric on $X$.
    Let $T = \metricTopology{d}{X}$.
    Let $T_1 = \metricTopology{d_1}{X}$.
    Then $T_1$ is finer than $T$ on $X$ iff
    for all $x\in X$ we have
    for all $\epsilon\in\reals$ such that $0 < \epsilon$ there exists $\delta\in\reals$
        such that $0 < \delta$ and
        $\metricBall{d_1}{X}{x}{\delta} \subseteq \metricBall{d}{X}{x}{\epsilon}$.
\end{lemma}
\begin{proof}
    Show that if $T_1$ is finer than $T$ on $X$ then
        for all $x\in X$ we have
        for all $\epsilon\in\reals$ such that $0 < \epsilon$ there exists $\delta\in\reals$
            such that $0 < \delta$ and
            $\metricBall{d_1}{X}{x}{\delta} \subseteq \metricBall{d}{X}{x}{\epsilon}$.
    \begin{subproof}
        Assume $T_1$ is finer than $T$ on $X$.
        Fix $x\in X$.
        Fix $\epsilon\in\reals$.
        Assume $0 < \epsilon$.
        Let $U = \metricBall{d}{X}{x}{\epsilon}$.
        Show that for all $y$ such that $y\in U$ we have $y\in X$.
        \begin{subproof}
            Fix $y$.
            Assume $y\in U$.
            Thus $y\in X$ by \cref{metric_ball}.
        \end{subproof}
        Thus $U\subseteq X$ by \cref{subseteq}.
        $U\in\pow{X}$ by \cref{powerset_intro}.
        $U\in\metricBasis{d}{X}$ by \cref{metric_basis}.
        $\metricBasis{d}{X}$ is a basis on $X$ by \cref{metric_basis_is_basis}.
        Thus $U\in T$ by \cref{basis_elem_open_in_generated,basis_for_topology,metric_topology}.
        $T\subseteq T_1$ by \cref{finer_topology}.
        Thus $U\in T_1$ by \cref{subseteq}.
        $T_1$ is a topology on $X$ by \cref{basis_topology_is_topology,metric_topology,metric_basis_is_basis,basis_for_topology}.
        Thus $U$ is open in $X$ with respect to $T_1$ by \cref{open_in}.
        $d_1$ is a metric on $X$.
        Thus $x\in U$ by \cref{metric_ball,metric_zero_characterization,real_lt_subst_left,metric_on}.
        There exists $\delta\in\reals$ such that $0 < \delta$ and
            $\metricBall{d_1}{X}{x}{\delta} \subseteq U$
            by \cref{metric_open_iff}.
        Thus there exists $\delta\in\reals$ such that $0 < \delta$ and
            $\metricBall{d_1}{X}{x}{\delta} \subseteq \metricBall{d}{X}{x}{\epsilon}$.
    \end{subproof}
    Show that if
        for all $x\in X$ we have
        for all $\epsilon\in\reals$ such that $0 < \epsilon$ there exists $\delta\in\reals$
            such that $0 < \delta$ and
            $\metricBall{d_1}{X}{x}{\delta} \subseteq \metricBall{d}{X}{x}{\epsilon}$
        then $T_1$ is finer than $T$ on $X$.
    \begin{subproof}
        Assume for all $x\in X$ we have
            for all $\epsilon\in\reals$ such that $0 < \epsilon$ there exists $\delta\in\reals$
                such that $0 < \delta$ and
                $\metricBall{d_1}{X}{x}{\delta} \subseteq \metricBall{d}{X}{x}{\epsilon}$.
        Show that for all $U$ such that $U\in T$ we have $U\in T_1$.
        \begin{subproof}
            Fix $U$.
            Assume $U\in T$.
            $T$ is a topology on $X$ by \cref{basis_topology_is_topology,metric_topology,metric_basis_is_basis,basis_for_topology}.
            $U\in\pow{X}$ by \cref{topology_on,subseteq}.
            $U\subseteq X$ by \cref{powerset_elim}.
            $d$ is a metric on $X$.
            $d_1$ is a metric on $X$.
            Show that $U$ is open in $X$ with respect to $T_1$.
            \begin{subproof}
                Show that for all $x\in U$ there exists $\delta\in\reals$ such that $0 < \delta$ and
                    $\metricBall{d_1}{X}{x}{\delta} \subseteq U$.
                \begin{subproof}
                    Fix $x\in U$.
                    $U$ is open in $X$ with respect to $T$ by \cref{open_in}.
                    There exists $\epsilon\in\reals$ such that $0 < \epsilon$ and
                        $\metricBall{d}{X}{x}{\epsilon} \subseteq U$
                        by \cref{metric_open_iff}.
                    Take $\epsilon\in\reals$ such that $0 < \epsilon$ and
                        $\metricBall{d}{X}{x}{\epsilon} \subseteq U$.
                    $x\in X$ by \cref{subseteq}.
                    There exists $\delta\in\reals$ such that $0 < \delta$ and
                        $\metricBall{d_1}{X}{x}{\delta} \subseteq \metricBall{d}{X}{x}{\epsilon}$ by assumption.
                    Thus there exists $\delta\in\reals$ such that $0 < \delta$ and
                        $\metricBall{d_1}{X}{x}{\delta} \subseteq U$ by \cref{subseteq_transitive}.
                \end{subproof}
                Thus $U$ is open in $X$ with respect to $T_1$ by \cref{metric_open_iff}.
            \end{subproof}
            Thus $U\in T_1$ by \cref{open_in}.
        \end{subproof}
        Thus $T\subseteq T_1$ by \cref{subseteq}.
        $T$ is a topology on $X$ by \cref{basis_topology_is_topology,metric_topology,metric_basis_is_basis,basis_for_topology}.
        $T_1$ is a topology on $X$ by \cref{basis_topology_is_topology,metric_topology,metric_basis_is_basis,basis_for_topology}.
        Thus $T_1$ is finer than $T$ on $X$ by \cref{finer_topology}.
    \end{subproof}
\end{proof}

% from §20 helper lemma 13: initial segment is finite
\begin{lemma}\label{initial_segment_finite}
    Suppose $n\in\naturalsPlus$.
    Then $\initialSegment{n}$ is finite.
\end{lemma}
\begin{proof}
    Let $A = \{n\in\naturalsPlus \mid \text{$\initialSegment{n}$ is finite}\}$.
    Show that for all $n\in\naturalsPlus$ we have $n\in A$.
    \begin{subproof}
        Show that $A\subseteq\naturalsPlus$.
        \begin{subproof}
            Show that for all $n$ such that $n\in A$ we have $n\in\naturalsPlus$.
            \begin{subproof}
                Fix $n$.
                Assume $n\in A$.
                Then $n\in\naturalsPlus$.
            \end{subproof}
            Thus $A\subseteq\naturalsPlus$ by \cref{subseteq}.
        \end{subproof}
        $1\in\naturalsPlus$ by \cref{real_one_in_naturals}.
        Show that $1\in A$.
        \begin{subproof}
            Let $B = \initialSegment{1}$.
            Show that $B\subseteq\{1,1\}$.
            \begin{subproof}
                Show that for all $k$ such that $k\in B$ we have $k\in\{1,1\}$.
                \begin{subproof}
                    Fix $k$.
                    Assume $k\in B$.
                    Thus $k\in\naturalsPlus$ and $k \leq 1$ by \cref{initial_segment_naturalsplus}.
                    $1 \rless k$ or $1 = k$ by \cref{real_one_leq_naturalsplus}.
                    Thus $k = 1$.
                    Thus $k\in\{1,1\}$ by \cref{upair_intro_left}.
                \end{subproof}
                Thus $B\subseteq\{1,1\}$ by \cref{subseteq}.
            \end{subproof}
            $\{1,1\}$ is finite by \cref{pair_is_finite}.
            Thus $B$ is finite by \cref{subseteq_finite_is_finite}.
            Thus $1\in A$.
        \end{subproof}
        Show that for all $n\in A$ we have $n + 1\in A$.
        \begin{subproof}
            Fix $n$.
            Assume $n\in A$.
            Thus $\initialSegment{n}$ is finite.
            $n\in\naturalsPlus$ by \cref{subseteq}.
            $n + 1\in\naturalsPlus$ by \cref{real_naturals_closed_succ}.
            Let $B = \initialSegment{n + 1}$.
            Show that $B\subseteq \initialSegment{n}\union\{n + 1,n + 1\}$.
            \begin{subproof}
                Show that for all $k$ such that $k\in B$ we have $k\in \initialSegment{n}\union\{n + 1,n + 1\}$.
                \begin{subproof}
                    Fix $k$.
                    Assume $k\in B$.
                    Thus $k\in\naturalsPlus$ and $k \leq n + 1$ by \cref{initial_segment_naturalsplus}.
                    $k = n + 1$ or $k \neq n + 1$.
                    \begin{byCase}
                    \caseOf{$k = n + 1$.} Thus $k\in\{n + 1,n + 1\}$ by \cref{upair_intro_left}.
                    Thus $k\in \initialSegment{n}\union\{n + 1,n + 1\}$ by \cref{union_iff}.
                    \caseOf{$k \neq n + 1$.} Thus $k\in\initialSegment{n}$ by \cref{real_naturals_leq_succ_elim,initial_segment_naturalsplus}.
                    Thus $k\in \initialSegment{n}\union\{n + 1,n + 1\}$ by \cref{union_iff}.
                    \end{byCase}
                \end{subproof}
                Thus $B\subseteq \initialSegment{n}\union\{n + 1,n + 1\}$ by \cref{subseteq}.
            \end{subproof}
            $\{n + 1,n + 1\}$ is finite by \cref{pair_is_finite}.
            Thus $\initialSegment{n}\union\{n + 1,n + 1\}$ is finite by \cref{union_of_finite_is_finite}.
            Thus $B$ is finite by \cref{subseteq_finite_is_finite}.
            Thus $n + 1\in A$.
        \end{subproof}
        Thus for all $n\in\naturalsPlus$ we have $n\in A$ by \cref{real_naturals_induction}.
    \end{subproof}
    Thus $n\in A$ by \cref{subseteq}.
    Thus $\initialSegment{n}$ is finite.
\end{proof}

% from §20 helper lemma 14: product family of constant reals
\begin{lemma}\label{product_family_reals_power}
    Suppose $J$ is inhabited.
    Let $X = \{(i,\reals) \mid i\in J\}$.
    Then $\productFamily{X} = \realsPower{J}$.
\end{lemma}
\begin{proof}
    Show that $X$ is a function.
    \begin{subproof}
        Show that for all $w$ such that $w\in X$ there exist $x,y$ such that $w = (x,y)$.
        \begin{subproof}
            Fix $w$.
            Assume $w\in X$.
            Take $i\in J$ such that $w = (i,\reals)$.
            Thus there exist $x,y$ such that $w = (x,y)$.
        \end{subproof}
        Thus $X$ is a relation by \cref{relation}.
        Show that for all $i,y_1,y_2$ such that $(i,y_1)\in X$ and $(i,y_2)\in X$ we have $y_1 = y_2$.
        \begin{subproof}
            Fix $i$.
            Fix $y_1$.
            Fix $y_2$.
            Assume $(i,y_1)\in X$ and $(i,y_2)\in X$.
            Then $y_1 = \reals$ and $y_2 = \reals$.
            Thus $y_1 = y_2$.
        \end{subproof}
        Thus $X$ is right-unique by \cref{rightunique}.
        Thus $X$ is a function.
    \end{subproof}
    Show that $\dom{X} = J$.
    \begin{subproof}
        Show that for all $i$ such that $i\in J$ we have $i\in\dom{X}$.
        \begin{subproof}
            Fix $i$.
            Assume $i\in J$.
            Thus $(i,\reals)\in X$.
            Thus $i\in\dom{X}$ by \cref{dom_intro}.
        \end{subproof}
        Thus $J\subseteq\dom{X}$ by \cref{subseteq}.
        Show that for all $i$ such that $i\in\dom{X}$ we have $i\in J$.
        \begin{subproof}
            Fix $i$.
            Assume $i\in\dom{X}$.
            Take $y$ such that $(i,y)\in X$ by \cref{dom_iff}.
            Thus $y = \reals$ and $i\in J$.
            Thus $i\in J$.
        \end{subproof}
        Thus $\dom{X}\subseteq J$ by \cref{subseteq}.
        Thus $\dom{X} = J$ by \cref{subseteq_antisymmetric}.
    \end{subproof}
    Show that $\unions{\ran{X}} = \reals$.
    \begin{subproof}
        Show that for all $y$ such that $y\in\ran{X}$ we have $y = \reals$.
        \begin{subproof}
            Fix $y$.
            Assume $y\in\ran{X}$.
            Take $i$ such that $(i,y)\in X$ by \cref{ran_iff}.
            Thus $y = \reals$.
        \end{subproof}
        Show that for all $z$ such that $z\in\unions{\ran{X}}$ we have $z\in\reals$.
        \begin{subproof}
            Fix $z$.
            Assume $z\in\unions{\ran{X}}$.
            Take $y\in\ran{X}$ such that $z\in y$ by \cref{unions_iff}.
            Thus $y = \reals$.
            Thus $z\in\reals$.
        \end{subproof}
        Thus $\unions{\ran{X}}\subseteq\reals$ by \cref{subseteq}.
        Show that for all $z$ such that $z\in\reals$ we have $z\in\unions{\ran{X}}$.
        \begin{subproof}
            Fix $z$.
            Assume $z\in\reals$.
            Take $i$ such that $i\in J$.
            Thus $(i,\reals)\in X$.
            Thus $\reals\in\ran{X}$ by \cref{ran_intro}.
            Thus $z\in\unions{\ran{X}}$ by \cref{unions_iff}.
        \end{subproof}
        Thus $\reals\subseteq\unions{\ran{X}}$ by \cref{subseteq}.
        Thus $\unions{\ran{X}} = \reals$ by \cref{subseteq_antisymmetric}.
    \end{subproof}
    Show that for all $f$ such that $f\in\productFamily{X}$ we have $f\in\realsPower{J}$.
    \begin{subproof}
        Fix $f$.
        Assume $f\in\productFamily{X}$.
        Thus $f\in\funs{\dom{X}}{\unions{\ran{X}}}$ by \cref{indexed_product}.
        Thus $\dom{f} = \dom{X}$ by \cref{funs_elim}.
        Thus $f$ is a function to $\unions{\ran{X}}$ by \cref{funs_elim}.
        Thus $f$ is a function to $\reals$ by \cref{function_on_weaken_codom}.
        Thus $f\in\funs{J}{\reals}$ by \cref{funs_intro}.
        Thus $f\in\tuplePower{\reals}{J}$ by \cref{tuple_power}.
        Thus $f\in\realsPower{J}$ by \cref{reals_power}.
    \end{subproof}
    Thus $\productFamily{X}\subseteq\realsPower{J}$ by \cref{subseteq}.
    Show that for all $f$ such that $f\in\realsPower{J}$ we have $f\in\productFamily{X}$.
    \begin{subproof}
        Fix $f$.
        Assume $f\in\realsPower{J}$.
        Thus $f\in\tuplePower{\reals}{J}$ by \cref{reals_power}.
        Thus $f\in\funs{J}{\reals}$ by \cref{tuple_power}.
        Thus $\dom{f} = J$ and $f$ is a function to $\reals$ by \cref{funs_elim}.
        Thus $f$ is a function to $\unions{\ran{X}}$ by \cref{function_on_weaken_codom}.
        Thus $f\in\funs{\dom{X}}{\unions{\ran{X}}}$ by \cref{funs_intro}.
        Show that for all $\alpha\in\dom{X}$ we have $f(\alpha)\in\apply{X}{\alpha}$.
        \begin{subproof}
            Fix $\alpha$.
            Assume $\alpha\in\dom{X}$.
            Thus $\alpha\in J$ by \cref{subseteq}.
            Thus $(\alpha,\reals)\in X$.
            Thus $\apply{X}{\alpha} = \reals$ by \cref{function_apply_intro}.
            Thus $f(\alpha)\in\reals$ by \cref{funs_elim}.
            Thus $f(\alpha)\in\apply{X}{\alpha}$.
        \end{subproof}
        Thus $f\in\productFamily{X}$ by \cref{indexed_product}.
    \end{subproof}
    Thus $\realsPower{J}\subseteq\productFamily{X}$ by \cref{subseteq}.
    Thus $\productFamily{X} = \realsPower{J}$ by \cref{subseteq_antisymmetric}.
\end{proof}

% from §20 helper lemma 15: coordinate difference set is finite
\begin{lemma}\label{coord_diff_set_finite}
    Suppose $n\in\naturalsPlus$.
    Let $J = \initialSegment{n}$.
    Suppose $x\in\realsPower{J}$.
    Suppose $y\in\realsPower{J}$.
    Then $\coordDiffSet{n}{x}{y}\subseteq\reals$ and
         $\coordDiffSet{n}{x}{y}$ is finite and
         $\coordDiffSet{n}{x}{y}$ is inhabited.
\end{lemma}
\begin{proof}
    Let $D = \coordDiffSet{n}{x}{y}$.
    Show that for all $t$ such that $t\in D$ we have $t\in\reals$.
    \begin{subproof}
        Fix $t$.
        Assume $t\in D$.
        Take $i\in J$ such that $t = \abs{x(i) - y(i)}$ by \cref{coord_diff_set}.
        $x\in\tuplePower{\reals}{J}$ by \cref{reals_power}.
        $y\in\tuplePower{\reals}{J}$ by \cref{reals_power}.
        Thus $x\in\funs{J}{\reals}$ and $y\in\funs{J}{\reals}$ by \cref{tuple_power}.
        Thus $\dom{x} = J$ and $\dom{y} = J$ by \cref{funs_elim}.
        Thus $x(i)\in\reals$ and $y(i)\in\reals$ by \cref{funs_elim}.
        Thus $\neg{y(i)}\in\reals$ by \cref{real_add_inverse}.
        Thus $x(i) + \neg{y(i)}\in\reals$ by \cref{real_add_closed}.
        Thus $x(i) - y(i)\in\reals$ by \cref{real_sub}.
        Thus $t\in\reals$ by \cref{real_abs_in_reals}.
    \end{subproof}
    Thus $D\subseteq\reals$ by \cref{subseteq}.
    Let $g = \{(i,\abs{x(i) - y(i)}) \mid i\in J\}$.
    Show that $g$ is a function on $J$.
    \begin{subproof}
        Show that $g$ is a relation.
        \begin{subproof}
            Show that for all $w$ such that $w\in g$ there exist $i,t$ such that $w = (i,t)$.
            \begin{subproof}
                Fix $w$.
                Assume $w\in g$.
                Take $i\in J$ such that $w = (i,\abs{x(i) - y(i)})$.
                Thus there exist $i_0,t$ such that $w = (i_0,t)$.
            \end{subproof}
            Thus $g$ is a relation by \cref{relation}.
        \end{subproof}
        Show that for all $i,t_1,t_2$ such that $(i,t_1)\in g$ and $(i,t_2)\in g$ we have $t_1 = t_2$.
        \begin{subproof}
            Fix $i$.
            Fix $t_1$.
            Fix $t_2$.
            Assume $(i,t_1)\in g$ and $(i,t_2)\in g$.
            Then $t_1 = \abs{x(i) - y(i)}$ and $t_2 = \abs{x(i) - y(i)}$.
            Thus $t_1 = t_2$.
        \end{subproof}
        Thus $g$ is right-unique by \cref{rightunique}.
        Thus $g$ is a function.
        Show that $\dom{g} = J$.
        \begin{subproof}
            Show that for all $i$ such that $i\in J$ we have $i\in\dom{g}$.
            \begin{subproof}
                Fix $i$.
                Assume $i\in J$.
                Thus $(i,\abs{x(i) - y(i)})\in g$.
                Thus $i\in\dom{g}$ by \cref{dom_intro}.
            \end{subproof}
            Thus $J\subseteq\dom{g}$ by \cref{subseteq}.
            Show that for all $i$ such that $i\in\dom{g}$ we have $i\in J$.
            \begin{subproof}
                Fix $i$.
                Assume $i\in\dom{g}$.
                Take $t$ such that $(i,t)\in g$ by \cref{dom_iff}.
                Thus $i\in J$.
            \end{subproof}
            Thus $\dom{g}\subseteq J$ by \cref{subseteq}.
            Thus $\dom{g} = J$ by \cref{subseteq_antisymmetric}.
        \end{subproof}
        Thus $g$ is a function on $J$.
    \end{subproof}
    Show that $D = \img{g}{J}$.
    \begin{subproof}
        Show that for all $t$ we have $t\in D$ iff $t\in\img{g}{J}$.
        \begin{subproof}
            Fix $t$.
            Show that if $t\in D$ then $t\in\img{g}{J}$.
            \begin{subproof}
                Assume $t\in D$.
                Take $i\in J$ such that $t = \abs{x(i) - y(i)}$ by \cref{coord_diff_set}.
                Thus $(i,t)\in g$.
                Thus $t\in\img{g}{J}$ by \cref{img_iff}.
            \end{subproof}
            Show that if $t\in\img{g}{J}$ then $t\in D$.
            \begin{subproof}
                Assume $t\in\img{g}{J}$.
                Take $i\in J$ such that $(i,t)\in g$ by \cref{img_iff}.
                Thus $t = \abs{x(i) - y(i)}$.
                Thus $t\in D$ by \cref{coord_diff_set}.
            \end{subproof}
        \end{subproof}
        Thus $D = \img{g}{J}$ by \cref{setext}.
    \end{subproof}
    $J$ is finite by \cref{initial_segment_finite}.
    Thus $\img{g}{J}$ is finite by \cref{finite_image_function}.
    Thus $D$ is finite.
    $1\in\naturalsPlus$ by \cref{real_one_in_naturals}.
    $1 \leq n$ by \cref{real_one_leq_naturalsplus}.
    Thus $1\in J$ by \cref{initial_segment_naturalsplus}.
    Thus $(1,\abs{x(1) - y(1)})\in g$.
    Thus $\abs{x(1) - y(1)}\in D$ by \cref{coord_diff_set}.
    Thus $D$ is inhabited.
\end{proof}

% from §20 helper lemma 16: square metric is a function
\begin{lemma}\label{square_metric_function}
    Suppose $n\in\naturalsPlus$.
    Then $\squareMetric{n}$ is a function.
\end{lemma}
\begin{proof}
    Let $d = \squareMetric{n}$.
    Show that for all $w$ such that $w\in d$ there exist $p,r$ such that $w = (p,r)$.
    \begin{subproof}
        Fix $w$.
        Assume $w\in d$.
        Take $p,x,y,r$ such that
            $p\in\realsPower{\initialSegment{n}}\times\realsPower{\initialSegment{n}}$ and
            $x\in\realsPower{\initialSegment{n}}$ and
            $y\in\realsPower{\initialSegment{n}}$ and
            $r\in\reals$ and
            $w = (p,r)$ and
            $p = (x,y)$ and
            $r\in\coordDiffSet{n}{x}{y}$ and
            for all $t\in\coordDiffSet{n}{x}{y}$ we have $t \leq r$
            by \cref{square_metric}.
        Thus there exist $p_0,r_0$ such that $w = (p_0,r_0)$.
    \end{subproof}
    Thus $d$ is a relation by \cref{relation}.
    Show that for all $p,r_1,r_2$ such that $(p,r_1)\in d$ and $(p,r_2)\in d$ we have $r_1 = r_2$.
    \begin{subproof}
        Fix $p$.
        Fix $r_1$.
        Fix $r_2$.
        Assume $(p,r_1)\in d$ and $(p,r_2)\in d$.
        Take $x_1,y_1\in\realsPower{\initialSegment{n}}$ such that
            $p = (x_1,y_1)$ and
            $r_1\in\coordDiffSet{n}{x_1}{y_1}$ and
            for all $t\in\coordDiffSet{n}{x_1}{y_1}$ we have $t \leq r_1$
            by \cref{square_metric,pair_eq_iff}.
        Take $x_2,y_2\in\realsPower{\initialSegment{n}}$ such that
            $p = (x_2,y_2)$ and
            $r_2\in\coordDiffSet{n}{x_2}{y_2}$ and
            for all $t\in\coordDiffSet{n}{x_2}{y_2}$ we have $t \leq r_2$
            by \cref{square_metric,pair_eq_iff}.
        Thus $x_1 = x_2$ and $y_1 = y_2$ by \cref{pair_eq_iff}.
        Thus $r_1\in\coordDiffSet{n}{x_2}{y_2}$.
        Thus $r_1\in\reals$ by \cref{coord_diff_set_finite,subseteq}.
        Thus $r_2\in\reals$ by \cref{coord_diff_set_finite,subseteq}.
        Thus $r_2 \leq r_1$.
        Thus $r_1 \leq r_2$.
        Thus $r_1 = r_2$ by \cref{real_leq_antisym}.
    \end{subproof}
    Thus $d$ is right-unique by \cref{rightunique}.
    Thus $d$ is a function.
\end{proof}

% from §20 helper lemma 17: absolute value bound from interval
\begin{lemma}\label{real_abs_lt_from_interval}
    Suppose $a\in\reals$.
    Suppose $b\in\reals$.
    Suppose $\epsilon\in\reals$.
    Suppose $0 < \epsilon$.
    Suppose $a - \epsilon < b$ and $b < a + \epsilon$.
    Then $\abs{a - b} < \epsilon$.
\end{lemma}
\begin{proof}
    $\neg{a}\in\reals$ by \cref{real_add_inverse}.
    $b - a = b + \neg{a}$ by \cref{real_sub}.
    Thus $b - a\in\reals$ by \cref{real_add_closed}.
    Show that $b - a < \epsilon$.
    \begin{subproof}
        $b < a + \epsilon$.
        $a + \epsilon\in\reals$ by \cref{real_add_closed}.
        Thus $b + \neg{a} < (a + \epsilon) + \neg{a}$ by \cref{real_lt_add}.
        $(a + \epsilon) + \neg{a} = a + (\epsilon + \neg{a})$ by \cref{real_add_assoc}.
        $\epsilon + \neg{a} = \neg{a} + \epsilon$ by \cref{real_add_comm}.
        $a + (\neg{a} + \epsilon) = (a + \neg{a}) + \epsilon$ by \cref{real_add_assoc}.
        $a + \neg{a} = 0$ by \cref{real_add_inverse}.
        Thus $(a + \epsilon) + \neg{a} = 0 + \epsilon$.
        $0 + \epsilon = \epsilon$ by \cref{real_add_ident}.
        Thus $b - a < \epsilon$ by \cref{real_sub}.
    \end{subproof}
    Show that $\neg{\epsilon} < b - a$.
    \begin{subproof}
        $a - \epsilon < b$.
        $\neg{\epsilon}\in\reals$ by \cref{real_add_inverse}.
        $a + \neg{\epsilon}\in\reals$ by \cref{real_add_closed}.
        $a - \epsilon = a + \neg{\epsilon}$ by \cref{real_sub}.
        Thus $a + \neg{\epsilon} < b$ by \cref{real_lt_subst_left}.
        Thus $(a + \neg{\epsilon}) + \neg{a} < b + \neg{a}$ by \cref{real_lt_add}.
        Thus $(a - \epsilon) + \neg{a} = (a + \neg{\epsilon}) + \neg{a}$ by \cref{real_sub}.
        $(a + \neg{\epsilon}) + \neg{a} = a + (\neg{\epsilon} + \neg{a})$ by \cref{real_add_assoc}.
        $\neg{\epsilon} + \neg{a} = \neg{a} + \neg{\epsilon}$ by \cref{real_add_comm}.
        $a + (\neg{a} + \neg{\epsilon}) = (a + \neg{a}) + \neg{\epsilon}$ by \cref{real_add_assoc}.
        $a + \neg{a} = 0$ by \cref{real_add_inverse}.
        Thus $(a - \epsilon) + \neg{a} = 0 + \neg{\epsilon}$.
        $0 + \neg{\epsilon} = \neg{\epsilon}$ by \cref{real_add_ident}.
        Thus $\neg{\epsilon} < b - a$ by \cref{real_sub}.
    \end{subproof}
    Show that $\abs{b - a} < \epsilon$.
    \begin{subproof}
        \begin{byCase}
        \caseOf{$b - a = 0$.}
            Thus $0 \leq b - a$.
            Thus $\abs{b - a} = b - a$ by \cref{real_abs_nonneg}.
            Thus $\abs{b - a} = 0$.
            Thus $\abs{b - a} < \epsilon$ by \cref{real_lt_subst_left}.
        \caseOf{$b - a \neq 0$.}
            $0\in\reals$ by \cref{real_add_ident}.
            Thus $b - a < 0$ or $0 < b - a$ by \cref{real_lt_trichotomy}.
            \begin{byCase}
            \caseOf{$b - a < 0$.}
                Thus $\abs{b - a} = \neg{(b - a)}$ by \cref{real_abs_neg}.
                $\neg{(\neg{\epsilon})} = \epsilon$ by \cref{real_neg_neg}.
                $\neg{\epsilon}\in\reals$ by \cref{real_add_inverse}.
                Thus $\neg{\epsilon} < b - a$.
                Thus $\neg{(b - a)} < \neg{(\neg{\epsilon})}$ by \cref{real_lt_neg}.
                Thus $\neg{(b - a)} < \epsilon$ by \cref{real_lt_subst_right}.
                Thus $\abs{b - a} < \epsilon$.
            \caseOf{$0 < b - a$.}
                Thus $\abs{b - a} = b - a$ by \cref{real_abs_nonneg}.
                Thus $\abs{b - a} < \epsilon$ by \cref{real_lt_subst_left}.
            \end{byCase}
        \end{byCase}
    \end{subproof}
    Thus $\abs{a - b} < \epsilon$ by \cref{real_abs_sub_swap}.
\end{proof}

% from §20 helper lemma 18: sum of nonnegative finite sequence is nonnegative
\begin{lemma}\label{sum_finite_sequence_nonneg}
    Suppose $\sumFiniteSequence{a}{s}$.
    Suppose for all $i\in\dom{a}$ we have $0 \leq a(i)$.
    Then $0 \leq s$.
\end{lemma}
\begin{proof}
    Take $n$ such that $n\in\naturalsPlus$ and $a\in\funs{\initialSegment{n}}{\reals}$ and
        there exists $t\in\funs{\initialSegment{n}}{\reals}$ such that
            $t(1) = a(1)$ and $s = t(n)$ and
            for all $k\in\naturalsPlus$ such that $k + 1 \leq n$ we have $t(k + 1) = t(k) + a(k + 1)$
        by \cref{sum_finite_sequence_pred}.
    Take $t$ such that $t\in\funs{\initialSegment{n}}{\reals}$ and
            $t(1) = a(1)$ and $s = t(n)$ and
            for all $k\in\naturalsPlus$ such that $k + 1 \leq n$ we have $t(k + 1) = t(k) + a(k + 1)$.
    Thus $\dom{a} = \initialSegment{n}$ by \cref{funs_elim}.
    Let $A = \{k\in\naturalsPlus \mid \text{if $k \leq n$ then $0 \leq t(k)$}\}$.
    $A\subseteq\naturalsPlus$ by \cref{subseteq}.
    Show that $1\in A$.
    \begin{subproof}
        $1\in\naturalsPlus$ by \cref{real_one_in_naturals}.
        $1 \leq n$ by \cref{real_one_leq_naturalsplus}.
        Thus $1\in\dom{a}$.
        Thus $0 \leq a(1)$.
        Thus $0 \leq t(1)$.
        Thus $1\in A$.
    \end{subproof}
    Show that for all $k\in A$ we have $k + 1\in A$.
    \begin{subproof}
        Fix $k$.
        Assume $k\in A$.
        Show that $k + 1\in A$.
        \begin{subproof}
            Show that if $k + 1 \leq n$ then $0 \leq t(k + 1)$.
            \begin{subproof}
                Assume $k + 1 \leq n$.
                $k\in\naturalsPlus$ by \cref{subseteq}.
                $k\in\reals$ by \cref{real_naturals_subset_reals,subseteq}.
                $0\in\reals$ by \cref{real_add_ident}.
                $1\in\reals$ by \cref{real_naturals_subset_reals,real_one_in_naturals,subseteq}.
                $0 < 1$ by \cref{real_zero_lt_one}.
                $0 + k < 1 + k$ by \cref{real_lt_add}.
                $0 + k = k$ by \cref{real_add_ident}.
                $1 + k = k + 1$ by \cref{real_add_comm}.
                Thus $k < k + 1$ by \cref{real_lt_subst_left}.
                Show that $k \leq n$.
                \begin{subproof}
                    Thus $k + 1 < n$ or $k + 1 = n$.
                    $k + 1\in\reals$ by \cref{real_add_closed}.
                    $n\in\reals$ by \cref{real_naturals_subset_reals,subseteq}.
                    \begin{byCase}
                    \caseOf{$k + 1 = n$.}
                        Thus $k < n$ by \cref{real_lt_subst_right}.
                        Thus $k \leq n$.
                    \caseOf{$k + 1 < n$.}
                        Thus $k < n$ by \cref{real_lt_trans}.
                        Thus $k \leq n$.
                    \end{byCase}
                \end{subproof}
                Thus $0 \leq t(k)$.
                $k + 1\in\naturalsPlus$ by \cref{real_naturals_closed_succ}.
                Thus $k + 1\in\dom{a}$.
                Thus $0 \leq a(k + 1)$.
                Thus $k\in\initialSegment{n}$ by \cref{initial_segment_naturalsplus}.
                Thus $k\in\dom{t}$ by \cref{funs_elim}.
                Thus $t(k)\in\reals$ by \cref{funs_type_apply}.
                Thus $a(k + 1)\in\reals$ by \cref{funs_type_apply}.
                Thus $t(k) + a(k + 1)\in\reals$ by \cref{real_add_closed}.
                $0\in\reals$ by \cref{real_add_ident}.
                $0 + a(k + 1) \leq t(k) + a(k + 1)$ by \cref{real_leq_add}.
                $0 + a(k + 1) = a(k + 1)$ by \cref{real_add_ident}.
                Thus $0 \leq t(k) + a(k + 1)$ by \cref{real_leq_trans}.
                Thus $0 \leq t(k + 1)$.
            \end{subproof}
            Thus $k + 1\in A$.
        \end{subproof}
        Thus $k + 1\in A$.
    \end{subproof}
    Thus for all $k\in\naturalsPlus$ we have $k\in A$ by \cref{real_naturals_induction}.
    Thus $n\in A$ by \cref{subseteq}.
    Thus $0 \leq t(n)$.
    Thus $0 \leq s$.
\end{proof}

% from §20 helper lemma 19: terms are bounded by the sum
\begin{lemma}\label{sum_finite_sequence_term_leq}
    Suppose $\sumFiniteSequence{a}{s}$.
    Suppose for all $i\in\dom{a}$ we have $0 \leq a(i)$.
    Then for all $i\in\dom{a}$ we have $a(i) \leq s$.
\end{lemma}
\begin{proof}
    Take $n$ such that $n\in\naturalsPlus$ and $a\in\funs{\initialSegment{n}}{\reals}$ and
        there exists $t\in\funs{\initialSegment{n}}{\reals}$ such that
            $t(1) = a(1)$ and $s = t(n)$ and
            for all $k\in\naturalsPlus$ such that $k + 1 \leq n$ we have $t(k + 1) = t(k) + a(k + 1)$
        by \cref{sum_finite_sequence_pred}.
    Take $t$ such that $t\in\funs{\initialSegment{n}}{\reals}$ and
            $t(1) = a(1)$ and $s = t(n)$ and
            for all $k\in\naturalsPlus$ such that $k + 1 \leq n$ we have $t(k + 1) = t(k) + a(k + 1)$.
    Thus $\dom{a} = \initialSegment{n}$ and $\dom{t} = \initialSegment{n}$ by \cref{funs_elim}.
    Let $A = \{k\in\naturalsPlus \mid \text{if $k \leq n$ then
        $0 \leq t(k)$ and $a(k) \leq t(k)$ and
        for all $j\in\initialSegment{k}$ we have $t(j) \leq t(k)$}\}$.
    $A\subseteq\naturalsPlus$ by \cref{subseteq}.
    Show that $1\in A$.
    \begin{subproof}
        $1\in\naturalsPlus$ by \cref{real_one_in_naturals}.
        $1 \leq n$ by \cref{real_one_leq_naturalsplus}.
        Thus $1\in\dom{a}$.
        Thus $0 \leq a(1)$.
        Thus $t(1) = a(1)$.
        Thus $0 \leq t(1)$.
        Thus $a(1) \leq t(1)$.
        Show that for all $j\in\initialSegment{1}$ we have $t(j) \leq t(1)$.
        \begin{subproof}
            Fix $j$.
            Assume $j\in\initialSegment{1}$.
            Thus $j\in\naturalsPlus$ and $j \leq 1$ by \cref{initial_segment_naturalsplus}.
            $1 \rless j$ or $1 = j$ by \cref{real_one_leq_naturalsplus}.
            Thus $j = 1$.
            Thus $t(j) = t(1)$.
            Thus $t(j) \leq t(1)$.
        \end{subproof}
        Thus $1\in A$.
    \end{subproof}
    Show that for all $k\in A$ we have $k + 1\in A$.
    \begin{subproof}
        Fix $k$.
        Assume $k\in A$.
        Show that $k + 1\in A$.
        \begin{subproof}
            Show that if $k + 1 \leq n$ then
                $0 \leq t(k + 1)$ and $a(k + 1) \leq t(k + 1)$ and
                for all $j\in\initialSegment{k + 1}$ we have $t(j) \leq t(k + 1)$.
            \begin{subproof}
                Assume $k + 1 \leq n$.
                $k\in\naturalsPlus$ by \cref{subseteq}.
                $k + 1\in\naturalsPlus$ by \cref{real_naturals_closed_succ}.
                $k\in\reals$ by \cref{real_naturals_subset_reals,subseteq}.
                $1\in\reals$ by \cref{real_naturals_subset_reals,real_one_in_naturals,subseteq}.
                $0\in\reals$ by \cref{real_add_ident}.
                $0 < 1$ by \cref{real_zero_lt_one}.
                $0 + k < 1 + k$ by \cref{real_lt_add}.
                $0 + k = k$ by \cref{real_add_ident}.
                $1 + k = k + 1$ by \cref{real_add_comm}.
                Thus $k < k + 1$ by \cref{real_lt_subst_left}.
                Show that $k \leq n$.
                \begin{subproof}
                    Thus $k + 1 < n$ or $k + 1 = n$.
                    $k + 1\in\reals$ by \cref{real_add_closed}.
                    $n\in\reals$ by \cref{real_naturals_subset_reals,subseteq}.
                    \begin{byCase}
                    \caseOf{$k + 1 = n$.}
                        Thus $k < n$ by \cref{real_lt_subst_right}.
                        Thus $k \leq n$.
                    \caseOf{$k + 1 < n$.}
                        Thus $k < n$ by \cref{real_lt_trans}.
                        Thus $k \leq n$.
                    \end{byCase}
                \end{subproof}
                Thus $k\in\dom{a}$ and $k\in\dom{t}$ by \cref{initial_segment_naturalsplus}.
                Thus $0 \leq t(k)$ and $a(k) \leq t(k)$ and
                    for all $j\in\initialSegment{k}$ we have $t(j) \leq t(k)$.
                Thus $k + 1\in\dom{a}$.
                Thus $0 \leq a(k + 1)$.
                Thus $a(k + 1)\in\reals$ by \cref{funs_type_apply}.
                Thus $t(k)\in\reals$ by \cref{funs_type_apply}.
                Thus $t(k) + a(k + 1)\in\reals$ by \cref{real_add_closed}.
                $0 + a(k + 1) \leq t(k) + a(k + 1)$ by \cref{real_leq_add}.
                $0 + a(k + 1) = a(k + 1)$ by \cref{real_add_ident}.
                Thus $0 \leq t(k) + a(k + 1)$ by \cref{real_leq_trans}.
                Thus $0 \leq t(k + 1)$.
                $t(k + 1) = t(k) + a(k + 1)$.
                $0 + t(k) \leq a(k + 1) + t(k)$ by \cref{real_leq_add}.
                $0 + t(k) = t(k)$ by \cref{real_add_ident}.
                $a(k + 1) + t(k) = t(k) + a(k + 1)$ by \cref{real_add_comm}.
                Thus $t(k) \leq t(k + 1)$ by \cref{real_leq_trans}.
                Thus $a(k + 1) \leq t(k + 1)$ by \cref{real_leq_trans}.
                Show that for all $j\in\initialSegment{k + 1}$ we have $t(j) \leq t(k + 1)$.
                \begin{subproof}
                    Fix $j$.
                    Assume $j\in\initialSegment{k + 1}$.
                    Thus $j\in\naturalsPlus$ and $j \leq k + 1$ by \cref{initial_segment_naturalsplus}.
                    $j = k + 1$ or $j \neq k + 1$.
                    \begin{byCase}
                    \caseOf{$j = k + 1$.}
                        Thus $t(j) = t(k + 1)$.
                        Thus $t(j) \leq t(k + 1)$.
                    \caseOf{$j \neq k + 1$.}
                        Thus $j \leq k$ by \cref{real_naturals_leq_succ_elim}.
                        Thus $k \leq n$.
                        $j\in\reals$ by \cref{real_naturals_subset_reals,subseteq}.
                        $k\in\reals$ by \cref{real_naturals_subset_reals,subseteq}.
                        $n\in\reals$ by \cref{real_naturals_subset_reals,subseteq}.
                        Thus $j \leq n$ by \cref{real_leq_trans}.
                        Thus $j\in\initialSegment{k}$ by \cref{initial_segment_naturalsplus}.
                        Thus $t(j) \leq t(k)$.
                        Thus $k \leq n$.
                        Thus $j \leq n$ by \cref{real_leq_trans}.
                        Thus $j\in\initialSegment{n}$ by \cref{initial_segment_naturalsplus}.
                        Thus $j\in\dom{t}$ by \cref{funs_elim}.
                        Thus $t(j)\in\reals$ by \cref{funs_type_apply}.
                        Thus $t(k + 1)\in\reals$ by \cref{funs_type_apply}.
                        Thus $t(j) \leq t(k + 1)$ by \cref{real_leq_trans}.
                    \end{byCase}
                \end{subproof}
            \end{subproof}
            Thus $k + 1\in A$.
        \end{subproof}
        Thus $k + 1\in A$.
    \end{subproof}
    Thus for all $k\in\naturalsPlus$ we have $k\in A$ by \cref{real_naturals_induction}.
    Show that for all $i\in\dom{a}$ we have $a(i) \leq s$.
    \begin{subproof}
        Fix $i$.
        Assume $i\in\dom{a}$.
        Thus $i\in\initialSegment{n}$.
        Thus $i\in\naturalsPlus$ and $i \leq n$ by \cref{initial_segment_naturalsplus}.
        Thus $i\in A$ by \cref{subseteq}.
        Thus $a(i) \leq t(i)$.
        Thus $n\in A$ by \cref{subseteq}.
        Thus for all $j\in\initialSegment{n}$ we have $t(j) \leq t(n)$.
        Thus $t(i) \leq t(n)$ by \cref{initial_segment_naturalsplus}.
        Thus $a(i)\in\reals$ by \cref{funs_type_apply}.
        Thus $t(i)\in\reals$ by \cref{funs_type_apply}.
        Thus $n = n$.
        Thus $n < n$ or $n = n$.
        Thus $n \leq n$.
        Thus $n\in\initialSegment{n}$ by \cref{initial_segment_naturalsplus}.
        Thus $n\in\dom{t}$ by \cref{funs_elim}.
        Thus $t(n)\in\reals$ by \cref{funs_type_apply}.
        Thus $a(i) \leq t(n)$ by \cref{real_leq_trans}.
        Thus $a(i) \leq s$.
    \end{subproof}
\end{proof}

% from §20 helper lemma 20: square metric value exists
\begin{lemma}\label{square_metric_value_exists}
    Suppose $n\in\naturalsPlus$.
    Let $J = \initialSegment{n}$.
    Suppose $x\in\realsPower{J}$.
    Suppose $y\in\realsPower{J}$.
    Then there exists $r\in\reals$ such that $((x,y),r)\in\squareMetric{n}$ and
        for all $i\in J$ we have $\abs{x(i) - y(i)} \leq r$.
\end{lemma}
\begin{proof}
    Let $D = \coordDiffSet{n}{x}{y}$.
    $D\subseteq\reals$ and $D$ is finite and $D$ is inhabited by \cref{coord_diff_set_finite}.
    Take $d_0$ such that $d_0\in D$.
    Thus $D\neq\emptyset$ by \cref{emptyset}.
    Take $r\in D$ such that $r$ is a maximum of $D$ by \cref{finite_real_maximum}.
    Show that $((x,y),r)\in\squareMetric{n}$.
    \begin{subproof}
        Let $p = (x,y)$.
        $p\in\realsPower{J}\times\realsPower{J}$ by \cref{times_tuple_intro}.
        $r\in\reals$ by \cref{subseteq}.
        Thus $r\in\coordDiffSet{n}{x}{y}$.
        For all $t\in\coordDiffSet{n}{x}{y}$ we have $t \leq r$ by \cref{real_maximum}.
        Thus $((x,y),r)\in\squareMetric{n}$ by \cref{square_metric}.
    \end{subproof}
    Show that for all $i\in J$ we have $\abs{x(i) - y(i)} \leq r$.
    \begin{subproof}
        Fix $i\in J$.
        Thus $\abs{x(i) - y(i)}\in D$ by \cref{coord_diff_set}.
        Thus $\abs{x(i) - y(i)} \leq r$ by \cref{real_maximum}.
    \end{subproof}
\end{proof}

% from §20 helper lemma 21: bounded finite sum
\begin{lemma}\label{sum_finite_sequence_leq_n_times}
    Suppose $n\in\naturalsPlus$.
    Suppose $a\in\funs{\initialSegment{n}}{\reals}$.
    Suppose $\sumFiniteSequence{a}{s}$.
    Suppose $M\in\reals$.
    Suppose for all $i\in\initialSegment{n}$ we have $0 \leq a(i)$ and $a(i) \leq M$.
    Then $s \leq n \rmul M$.
\end{lemma}
\begin{proof}
    Take $n_0$ such that $n_0\in\naturalsPlus$ and $a\in\funs{\initialSegment{n_0}}{\reals}$ and
        there exists $t\in\funs{\initialSegment{n_0}}{\reals}$ such that
            $t(1) = a(1)$ and $s = t(n_0)$ and
            for all $k\in\naturalsPlus$ such that $k + 1 \leq n_0$ we have $t(k + 1) = t(k) + a(k + 1)$
        by \cref{sum_finite_sequence_pred}.
    Take $t$ such that $t\in\funs{\initialSegment{n_0}}{\reals}$ and
        $t(1) = a(1)$ and $s = t(n_0)$ and
        for all $k\in\naturalsPlus$ such that $k + 1 \leq n_0$ we have $t(k + 1) = t(k) + a(k + 1)$.
    Thus $\dom{a} = \initialSegment{n}$ by \cref{funs_elim}.
    Thus $\dom{a} = \initialSegment{n_0}$ by \cref{funs_elim}.
    Thus $\initialSegment{n} = \initialSegment{n_0}$ by \cref{subseteq_antisymmetric,subseteq}.
    $n\in\naturalsPlus$ and $n_0\in\naturalsPlus$.
    Thus $n \leq n$ and $n_0 \leq n_0$.
    Thus $n\in\initialSegment{n}$ and $n_0\in\initialSegment{n_0}$.
    Thus $n\in\initialSegment{n_0}$ and $n_0\in\initialSegment{n}$ by \cref{subseteq}.
    Thus $n \leq n_0$ and $n_0 \leq n$.
    $n\in\reals$ and $n_0\in\reals$ by \cref{real_naturals_subset_reals,subseteq}.
    Thus $n = n_0$ by \cref{real_leq_antisym}.
    Thus $s = t(n)$.
    Thus $t\in\funs{\initialSegment{n}}{\reals}$ by \cref{funs_elim}.
    Thus $\dom{t} = \initialSegment{n}$ by \cref{funs_elim}.
    Let $A = \{k\in\naturalsPlus \mid \text{if $k \leq n$ then $t(k) \leq k \rmul M$}\}$.
    $A\subseteq\naturalsPlus$ by \cref{subseteq}.
    Show that $1\in A$.
    \begin{subproof}
        $1\in\naturalsPlus$ by \cref{real_one_in_naturals}.
        $1 \leq n$ by \cref{real_one_leq_naturalsplus}.
        Thus $1\in\dom{a}$.
        Thus $0 \leq a(1)$ and $a(1) \leq M$.
        Thus $t(1) \leq M$.
        $1\in\reals$ by \cref{real_naturals_subset_reals,real_one_in_naturals,subseteq}.
        $1 \rmul M = M$ by \cref{real_mul_ident}.
        Thus $t(1) \leq 1 \rmul M$.
        Thus $1\in A$.
    \end{subproof}
    Show that for all $k\in A$ we have $k + 1\in A$.
    \begin{subproof}
        Fix $k$.
        Assume $k\in A$.
        Show that $k + 1\in A$.
        \begin{subproof}
            Show that if $k + 1 \leq n$ then $t(k + 1) \leq (k + 1) \rmul M$.
            \begin{subproof}
                Assume $k + 1 \leq n$.
                $k\in\naturalsPlus$ by \cref{subseteq}.
                $k + 1\in\naturalsPlus$ by \cref{real_naturals_closed_succ}.
                $k\in\reals$ by \cref{real_naturals_subset_reals,subseteq}.
                $1\in\reals$ by \cref{real_naturals_subset_reals,real_one_in_naturals,subseteq}.
                $0\in\reals$ by \cref{real_add_ident}.
                $0 < 1$ by \cref{real_zero_lt_one}.
                $0 + k < 1 + k$ by \cref{real_lt_add}.
                $0 + k = k$ by \cref{real_add_ident}.
                $1 + k = k + 1$ by \cref{real_add_comm}.
                Thus $k < k + 1$ by \cref{real_lt_subst_left}.
                Show that $k \leq n$.
                \begin{subproof}
                    Thus $k + 1 < n$ or $k + 1 = n$.
                    $k + 1\in\reals$ by \cref{real_add_closed}.
                    $n\in\reals$ by \cref{real_naturals_subset_reals,subseteq}.
                    \begin{byCase}
                    \caseOf{$k + 1 = n$.}
                        Thus $k < n$ by \cref{real_lt_subst_right}.
                        Thus $k \leq n$.
                    \caseOf{$k + 1 < n$.}
                        Thus $k < n$ by \cref{real_lt_trans}.
                        Thus $k \leq n$.
                    \end{byCase}
                \end{subproof}
                Thus $t(k) \leq k \rmul M$.
                Thus $k\in\initialSegment{n}$ by \cref{initial_segment_naturalsplus}.
                Thus $k\in\dom{t}$ by \cref{funs_elim}.
                Thus $k + 1\in\dom{a}$.
                Thus $0 \leq a(k + 1)$ and $a(k + 1) \leq M$.
                Thus $a(k + 1)\in\reals$ by \cref{funs_type_apply}.
                Thus $t(k)\in\reals$ by \cref{funs_type_apply}.
                Thus $k \rmul M\in\reals$ by \cref{real_mul_closed}.
                Thus $t(k) + a(k + 1)\in\reals$ by \cref{real_add_closed}.
                $t(k) + a(k + 1) \leq (k \rmul M) + a(k + 1)$ by \cref{real_leq_add}.
                $a(k + 1) + (k \rmul M) \leq M + (k \rmul M)$ by \cref{real_leq_add}.
                $a(k + 1) + (k \rmul M) = (k \rmul M) + a(k + 1)$ by \cref{real_add_comm}.
                $M + (k \rmul M) = (k \rmul M) + M$ by \cref{real_add_comm}.
                $(k \rmul M) + a(k + 1)\in\reals$ by \cref{real_add_closed}.
                $(k \rmul M) + M\in\reals$ by \cref{real_add_closed}.
                Thus $(k \rmul M) + a(k + 1) \leq (k \rmul M) + M$ by \cref{real_leq_trans}.
                Thus $t(k) + a(k + 1) \leq (k \rmul M) + M$ by \cref{real_leq_trans}.
                $t(k + 1) = t(k) + a(k + 1)$.
                Show that $(k \rmul M) + M = (k + 1) \rmul M$.
                \begin{subproof}
                    $M\in\reals$.
                    $1 \rmul M = M$ by \cref{real_mul_ident}.
                    $k + 1\in\reals$ by \cref{real_add_closed}.
                    $(k + 1) \rmul M = M \rmul (k + 1)$ by \cref{real_mul_comm}.
                    $M \rmul (k + 1) = (M \rmul k) + (M \rmul 1)$ by \cref{real_distrib}.
                    $M \rmul k = k \rmul M$ by \cref{real_mul_comm}.
                    $M \rmul 1 = 1 \rmul M$ by \cref{real_mul_comm}.
                    Thus $(k + 1) \rmul M = (k \rmul M) + (1 \rmul M)$.
                    Thus $(k + 1) \rmul M = (k \rmul M) + M$.
                \end{subproof}
                Thus $t(k + 1) \leq (k + 1) \rmul M$ by \cref{real_leq_trans}.
            \end{subproof}
            Thus $k + 1\in A$.
        \end{subproof}
        Thus $k + 1\in A$.
    \end{subproof}
    Thus for all $k\in\naturalsPlus$ we have $k\in A$ by \cref{real_naturals_induction}.
    Thus $n\in A$ by \cref{subseteq}.
    Thus $t(n) \leq n \rmul M$.
    Thus $s \leq n \rmul M$.
\end{proof}

% from §20 helper lemma 22: squares preserve order on nonnegative reals
\begin{lemma}\label{real_nonneg_leq_iff_sq_leq}
    Suppose $a\in\reals$.
    Suppose $b\in\reals$.
    Suppose $0 \leq a$.
    Suppose $0 \leq b$.
    Then $a \leq b$ iff $\exp{a}{1+1} \leq \exp{b}{1+1}$.
\end{lemma}
\begin{proof}
    Show that if $a \leq b$ then $\exp{a}{1+1} \leq \exp{b}{1+1}$.
    \begin{subproof}
        Assume $a \leq b$.
        Thus $a < b$ or $a = b$ by \cref{real_leq_split}.
        \begin{byCase}
        \caseOf{$a = b$.}
            Thus $\exp{a}{1+1} = \exp{b}{1+1}$.
            Thus $\exp{a}{1+1} \leq \exp{b}{1+1}$.
        \caseOf{$a < b$.}
            $0 < a$ or $a = 0$ by \cref{real_leq_split}.
            \begin{byCase}
            \caseOf{$a = 0$.}
                Thus $0 < b$ by \cref{real_lt_subst_left}.
                Thus $b$ is positive.
                $\exp{a}{1+1} = \exp{0}{1+1}$.
                $\exp{0}{1} = 0$ by \cref{real_pow_one}.
                $1\in\naturalsPlus$ by \cref{real_one_in_naturals}.
                $\exp{0}{1+1} = \exp{0}{1} \rmul 0$ by \cref{real_pow_succ}.
                $0 \rmul 0 = 0$ by \cref{real_mul_zero}.
                Thus $\exp{a}{1+1} = 0$.
                $\exp{b}{1+1} \geq 0$ by \cref{real_sq_nonnegative}.
                Thus $\exp{a}{1+1} \leq \exp{b}{1+1}$.
            \caseOf{$0 < a$.}
                Thus $a$ is positive.
                Thus $0 < b$ by \cref{real_lt_trans}.
                Thus $b$ is positive.
                Thus $\exp{a}{1+1} < \exp{b}{1+1}$ by \cref{real_pos_lt_iff_sq_lt}.
                Thus $\exp{a}{1+1} \leq \exp{b}{1+1}$.
            \end{byCase}
        \end{byCase}
    \end{subproof}
    Show that if $\exp{a}{1+1} \leq \exp{b}{1+1}$ then $a \leq b$.
    \begin{subproof}
        Assume $\exp{a}{1+1} \leq \exp{b}{1+1}$.
        Suppose not.
        Thus $b < a$.
        $0 < b$ or $b = 0$ by \cref{real_leq_split}.
        \begin{byCase}
        \caseOf{$b = 0$.}
            Thus $0 < a$ by \cref{real_lt_subst_left}.
            Thus $a$ is positive.
            $1\in\naturalsPlus$ by \cref{real_one_in_naturals}.
            $\exp{a}{1+1} = \exp{a}{1} \rmul a$ by \cref{real_pow_succ}.
            $\exp{a}{1} = a$ by \cref{real_pow_one}.
            Thus $\exp{a}{1+1} = a \rmul a$.
            Thus $\exp{a}{1+1}\in\reals$ by \cref{real_mul_closed}.
            Thus $\exp{a}{1+1} > 0$ by \cref{real_mul_pos_pos_gt}.
            $\exp{b}{1+1} = \exp{0}{1+1}$.
            $\exp{0}{1} = 0$ by \cref{real_pow_one}.
            $1\in\naturalsPlus$ by \cref{real_one_in_naturals}.
            $\exp{0}{1+1} = \exp{0}{1} \rmul 0$ by \cref{real_pow_succ}.
            $0 \rmul 0 = 0$ by \cref{real_mul_zero}.
            Thus $\exp{b}{1+1} = 0$.
            $0\in\reals$ by \cref{real_add_ident}.
            Thus $\exp{b}{1+1}\in\reals$.
            Thus $\exp{b}{1+1} < \exp{a}{1+1}$ by \cref{real_lt_subst_left}.
            Thus $\exp{a}{1+1} < \exp{b}{1+1}$ or $\exp{a}{1+1} = \exp{b}{1+1}$ by \cref{real_leq_split}.
            \begin{byCase}
            \caseOf{$\exp{a}{1+1} < \exp{b}{1+1}$.}
                Thus $\exp{a}{1+1} < \exp{a}{1+1}$ by \cref{real_lt_trans}.
                Contradiction by \cref{real_lt_irreflexive}.
            \caseOf{$\exp{a}{1+1} = \exp{b}{1+1}$.}
                Thus $\exp{b}{1+1} < \exp{b}{1+1}$ by \cref{real_lt_subst_right}.
                Contradiction by \cref{real_lt_irreflexive}.
            \end{byCase}
        \caseOf{$0 < b$.}
            Thus $b$ is positive.
            Thus $0 < a$ by \cref{real_lt_trans}.
            Thus $a$ is positive.
            Thus $\exp{b}{1+1} < \exp{a}{1+1}$ by \cref{real_pos_lt_iff_sq_lt}.
            $1\in\naturalsPlus$ by \cref{real_one_in_naturals}.
            $\exp{a}{1+1} = \exp{a}{1} \rmul a$ by \cref{real_pow_succ}.
            $\exp{a}{1} = a$ by \cref{real_pow_one}.
            Thus $\exp{a}{1+1} = a \rmul a$.
            Thus $\exp{a}{1+1}\in\reals$ by \cref{real_mul_closed}.
            $\exp{b}{1+1} = \exp{b}{1} \rmul b$ by \cref{real_pow_succ}.
            $\exp{b}{1} = b$ by \cref{real_pow_one}.
            Thus $\exp{b}{1+1} = b \rmul b$.
            Thus $\exp{b}{1+1}\in\reals$ by \cref{real_mul_closed}.
            Thus $\exp{a}{1+1} < \exp{b}{1+1}$ or $\exp{a}{1+1} = \exp{b}{1+1}$ by \cref{real_leq_split}.
            \begin{byCase}
            \caseOf{$\exp{a}{1+1} < \exp{b}{1+1}$.}
                Thus $\exp{a}{1+1} < \exp{a}{1+1}$ by \cref{real_lt_trans}.
                Contradiction by \cref{real_lt_irreflexive}.
            \caseOf{$\exp{a}{1+1} = \exp{b}{1+1}$.}
                Thus $\exp{b}{1+1} < \exp{b}{1+1}$ by \cref{real_lt_subst_right}.
                Contradiction by \cref{real_lt_irreflexive}.
            \end{byCase}
        \end{byCase}
    \end{subproof}
\end{proof}

% from §20 helper lemma 23: coordinate square sequence element
\begin{lemma}\label{coord_square_sequence_elem}
    Suppose $n\in\naturalsPlus$.
    Let $J = \initialSegment{n}$.
    Suppose $x\in\realsPower{J}$.
    Suppose $y\in\realsPower{J}$.
    Let $a = \coordSquareSequence{n}{(x,y)}$.
    Suppose $i\in J$.
    Then $(i,\exp{x(i) - y(i)}{1+1})\in a$.
\end{lemma}
\begin{proof}
    Thus $i\in\initialSegment{n}$.
    Thus $(i,\exp{x(i) - y(i)}{1+1}) =
        (i,\exp{\apply{\fst{(x,y)}}{i} - \apply{\snd{(x,y)}}{i}}{1+1})$ by \cref{fst_eq,snd_eq}.
    Thus $(i,\exp{x(i) - y(i)}{1+1})\in a$ by \cref{coord_square_sequence}.
\end{proof}

% from §20 helper lemma 24: euclidean and square metric bounds
\begin{lemma}\label{euclidean_square_metric_bounds}
    Suppose $n\in\naturalsPlus$.
    Let $J = \initialSegment{n}$.
    Suppose $x\in\realsPower{J}$.
    Suppose $y\in\realsPower{J}$.
    Suppose $((x,y),r)\in\euclideanMetric{n}$.
    Suppose $((x,y),m)\in\squareMetric{n}$.
    Then $m \leq r$ and $r \leq \sqrt{n} \rmul m$.
\end{lemma}
\begin{proof}
    $1\in\naturalsPlus$ by \cref{real_one_in_naturals}.
    Let $D = \coordDiffSet{n}{x}{y}$.
    Take $s\in\reals$ such that
        $\sumFiniteSequence{\coordSquareSequence{n}{(x,y)}}{s}$ and $r = \sqrt{s}$
        by \cref{euclidean_metric,pair_eq_iff}.
    Let $a = \coordSquareSequence{n}{(x,y)}$.
    Take $n_1$ such that $n_1\in\naturalsPlus$ and $a\in\funs{\initialSegment{n_1}}{\reals}$ and
        there exists $t\in\funs{\initialSegment{n_1}}{\reals}$ such that
            $t(1) = a(1)$ and $s = t(n_1)$ and
            for all $k\in\naturalsPlus$ such that $k + 1 \leq n_1$ we have $t(k + 1) = t(k) + a(k + 1)$
        by \cref{sum_finite_sequence_pred}.
    Thus $a$ is a function and $\dom{a} = \initialSegment{n_1}$ by \cref{funs_elim}.
    Take $x_0,y_0\in\realsPower{J}$ such that
        $(x,y) = (x_0,y_0)$ and
        $m\in\coordDiffSet{n}{x_0}{y_0}$ and
        for all $t\in\coordDiffSet{n}{x_0}{y_0}$ we have $t \leq m$
        by \cref{square_metric,pair_eq_iff}.
    Thus $x = x_0$ and $y = y_0$ by \cref{pair_eq_iff}.
    Thus $m\in D$ and for all $t\in D$ we have $t \leq m$.
    Take $i_0\in J$ such that $m = \abs{x(i_0) - y(i_0)}$ by \cref{coord_diff_set}.
    Show that for all $i\in J$ we have $0 \leq a(i)$.
    \begin{subproof}
        Fix $i$.
        Assume $i\in J$.
        Thus $(i,\exp{x(i) - y(i)}{1+1})\in a$ by \cref{coord_square_sequence_elem}.
        $x\in\tuplePower{\reals}{J}$ and $y\in\tuplePower{\reals}{J}$ by \cref{reals_power}.
        Thus $x\in\funs{J}{\reals}$ and $y\in\funs{J}{\reals}$ by \cref{tuple_power}.
        Thus $x(i)\in\reals$ and $y(i)\in\reals$ by \cref{funs_type_apply}.
        Thus $\neg{y(i)}\in\reals$ by \cref{real_add_inverse}.
        Thus $x(i) + \neg{y(i)}\in\reals$ by \cref{real_add_closed}.
        Thus $x(i) - y(i)\in\reals$ by \cref{real_sub}.
        Thus $\exp{x(i) - y(i)}{1+1} \geq 0$ by \cref{real_sq_nonnegative}.
        Thus $a(i) = \exp{x(i) - y(i)}{1+1}$ by \cref{function_apply_intro}.
        Thus $a(i) \geq 0$.
    \end{subproof}
    Show that $\dom{a} = J$.
    \begin{subproof}
        Show that for all $i$ such that $i\in J$ we have $i\in\dom{a}$.
        \begin{subproof}
            Fix $i$.
            Assume $i\in J$.
            Thus $(i,\exp{x(i) - y(i)}{1+1})\in a$ by \cref{coord_square_sequence_elem}.
            Thus $i\in\dom{a}$ by \cref{dom_intro}.
        \end{subproof}
        Thus $J\subseteq\dom{a}$ by \cref{subseteq}.
        Show that for all $i$ such that $i\in\dom{a}$ we have $i\in J$.
        \begin{subproof}
            Fix $i$.
            Assume $i\in\dom{a}$.
            Take $w$ such that $(i,w)\in a$ by \cref{dom_iff}.
            Take $i_1\in\initialSegment{n}$ such that
                $(i,w) = (i_1,\exp{\apply{\fst{(x,y)}}{i_1} - \apply{\snd{(x,y)}}{i_1}}{1+1})$
                by \cref{coord_square_sequence}.
            Thus $i = i_1$ by \cref{pair_eq_iff}.
            Thus $i\in\initialSegment{n}$.
            Thus $i\in J$.
        \end{subproof}
        Thus $\dom{a}\subseteq J$ by \cref{subseteq}.
        Thus $\dom{a} = J$ by \cref{subseteq_antisymmetric}.
    \end{subproof}
    Thus $a$ is a function to $\reals$ by \cref{funs_elim}.
    Thus $a\in\funs{J}{\reals}$ by \cref{funs_intro}.
    Thus $0 \leq s$ by \cref{sum_finite_sequence_nonneg}.
    $r\in\reals$ and $r \geq 0$ and $r \rmul r = s$ by \cref{positive_square_root}.
    Show that $m \leq r$.
    \begin{subproof}
        Thus $i_0\in J$.
        Thus $(i_0,\exp{x(i_0) - y(i_0)}{1+1})\in a$ by \cref{coord_square_sequence_elem}.
        Thus $i_0\in\dom{a}$ by \cref{dom_intro}.
        Thus $a(i_0) \leq s$ by \cref{sum_finite_sequence_term_leq}.
        $x\in\tuplePower{\reals}{J}$ and $y\in\tuplePower{\reals}{J}$ by \cref{reals_power}.
        Thus $x\in\funs{J}{\reals}$ and $y\in\funs{J}{\reals}$ by \cref{tuple_power}.
        Thus $x(i_0)\in\reals$ and $y(i_0)\in\reals$ by \cref{funs_type_apply}.
        Thus $\neg{y(i_0)}\in\reals$ by \cref{real_add_inverse}.
        Thus $x(i_0) + \neg{y(i_0)}\in\reals$ by \cref{real_add_closed}.
        Thus $x(i_0) - y(i_0)\in\reals$ by \cref{real_sub}.
        $\exp{\abs{x(i_0) - y(i_0)}}{1+1} = \exp{x(i_0) - y(i_0)}{1+1}$ by \cref{real_abs_square}.
        Thus $\exp{m}{1+1} = \exp{x(i_0) - y(i_0)}{1+1}$.
        Thus $\exp{m}{1+1} = a(i_0)$ by \cref{function_apply_intro}.
        Thus $\exp{m}{1+1} \leq s$.
        $\exp{r}{1+1} = r \rmul r$ by \cref{real_pow_succ,real_pow_one}.
        Thus $\exp{r}{1+1} = s$.
        $m\in\reals$ by \cref{coord_diff_set_finite,subseteq}.
        $m \geq 0$ by \cref{real_abs_nonnegative}.
        Thus $m \leq r$ by \cref{real_nonneg_leq_iff_sq_leq}.
    \end{subproof}
    Show that $r \leq \sqrt{n} \rmul m$.
    \begin{subproof}
        Show that for all $i\in\initialSegment{n}$ we have $0 \leq a(i)$ and $a(i) \leq \exp{m}{1+1}$.
        \begin{subproof}
            Fix $i$.
            Assume $i\in\initialSegment{n}$.
            Thus $i\in J$.
            Thus $0 \leq a(i)$.
            Thus $i\in\dom{a}$ by \cref{subseteq}.
            Thus $\abs{x(i) - y(i)}\in D$ by \cref{coord_diff_set}.
            Thus $\abs{x(i) - y(i)} \leq m$.
            $m\in\reals$ by \cref{coord_diff_set_finite,subseteq}.
            $x\in\tuplePower{\reals}{J}$ and $y\in\tuplePower{\reals}{J}$ by \cref{reals_power}.
            Thus $x\in\funs{J}{\reals}$ and $y\in\funs{J}{\reals}$ by \cref{tuple_power}.
            Thus $x(i)\in\reals$ and $y(i)\in\reals$ by \cref{funs_type_apply}.
            Thus $\neg{y(i)}\in\reals$ by \cref{real_add_inverse}.
            Thus $x(i) + \neg{y(i)}\in\reals$ by \cref{real_add_closed}.
            Thus $x(i) - y(i)\in\reals$ by \cref{real_sub}.
            Thus $\abs{x(i) - y(i)}\in\reals$ by \cref{real_abs_in_reals}.
            $\abs{x(i) - y(i)} \geq 0$ by \cref{real_abs_nonnegative}.
            $0\in\reals$ by \cref{real_add_ident}.
            Thus $0 \leq m$ by \cref{real_leq_trans}.
            Thus $\exp{\abs{x(i) - y(i)}}{1+1} \leq \exp{m}{1+1}$ by \cref{real_nonneg_leq_iff_sq_leq}.
            $\exp{\abs{x(i) - y(i)}}{1+1} = \exp{x(i) - y(i)}{1+1}$ by \cref{real_abs_square}.
            Thus $(i,\exp{x(i) - y(i)}{1+1})\in a$ by \cref{coord_square_sequence_elem}.
            Thus $a(i) = \exp{x(i) - y(i)}{1+1}$ by \cref{function_apply_intro}.
            Thus $a(i) \leq \exp{m}{1+1}$.
            Thus $0 \leq a(i)$ and $a(i) \leq \exp{m}{1+1}$.
        \end{subproof}
        Thus $a\in\funs{\initialSegment{n}}{\reals}$.
        Thus $\sumFiniteSequence{a}{s}$.
        $m\in\reals$ by \cref{coord_diff_set_finite,subseteq}.
        $\exp{m}{1+1} = \exp{m}{1} \rmul m$ by \cref{real_pow_succ}.
        $\exp{m}{1} = m$ by \cref{real_pow_one}.
        Thus $\exp{m}{1+1} = m \rmul m$.
        Thus $\exp{m}{1+1}\in\reals$ by \cref{real_mul_closed}.
        Thus $s \leq n \rmul \exp{m}{1+1}$ by \cref{sum_finite_sequence_leq_n_times}.
        $n\in\reals$ by \cref{real_naturals_subset_reals,subseteq}.
        Show that $0 \leq n$.
        \begin{subproof}
            $0 < 1$ by \cref{real_zero_lt_one}.
            $0\in\reals$ by \cref{real_add_ident}.
            $1\in\reals$ by \cref{real_naturals_subset_reals,subseteq}.
            $1 \leq n$ by \cref{real_one_leq_naturalsplus}.
            Thus $1 < n$ or $1 = n$ by \cref{real_leq_split}.
            \begin{byCase}
            \caseOf{$1 < n$.}
                Thus $0 < n$ by \cref{real_lt_trans}.
                Thus $0 \leq n$.
            \caseOf{$1 = n$.}
                Thus $0 < n$ by \cref{real_lt_subst_right}.
                Thus $0 \leq n$.
            \end{byCase}
        \end{subproof}
        Thus $0 < n$.
        $\sqrt{n}\in\reals$ and $\sqrt{n} \geq 0$ and $\sqrt{n} \rmul \sqrt{n} = n$ by \cref{positive_square_root}.
        Show that $\exp{\sqrt{n} \rmul m}{1+1} = n \rmul \exp{m}{1+1}$.
        \begin{subproof}
            $\sqrt{n} \rmul m\in\reals$ by \cref{real_mul_closed}.
            $\exp{\sqrt{n} \rmul m}{1+1} = (\sqrt{n} \rmul m) \rmul (\sqrt{n} \rmul m)$ by \cref{real_pow_succ,real_pow_one}.
            $(\sqrt{n} \rmul m) \rmul (\sqrt{n} \rmul m) = (\sqrt{n} \rmul \sqrt{n}) \rmul (m \rmul m)$ by \cref{real_mul_assoc,real_mul_comm}.
            Thus $\exp{\sqrt{n} \rmul m}{1+1} = n \rmul (m \rmul m)$.
            $\exp{m}{1+1} = m \rmul m$ by \cref{real_pow_succ,real_pow_one}.
            Thus $\exp{\sqrt{n} \rmul m}{1+1} = n \rmul \exp{m}{1+1}$.
        \end{subproof}
        $\exp{r}{1+1} = r \rmul r$ by \cref{real_pow_succ,real_pow_one}.
        Thus $\exp{r}{1+1} = s$.
        Thus $\exp{r}{1+1} \leq \exp{\sqrt{n} \rmul m}{1+1}$ by \cref{real_leq_trans}.
        $r\in\reals$ and $r \geq 0$.
        $\sqrt{n} \rmul m\in\reals$ by \cref{real_mul_closed}.
        Show that $\sqrt{n} \rmul m \geq 0$.
        \begin{subproof}
            $x\in\tuplePower{\reals}{J}$ and $y\in\tuplePower{\reals}{J}$ by \cref{reals_power}.
            Thus $x\in\funs{J}{\reals}$ and $y\in\funs{J}{\reals}$ by \cref{tuple_power}.
            Thus $x(i_0)\in\reals$ and $y(i_0)\in\reals$ by \cref{funs_type_apply}.
            Thus $\neg{y(i_0)}\in\reals$ by \cref{real_add_inverse}.
            Thus $x(i_0) + \neg{y(i_0)}\in\reals$ by \cref{real_add_closed}.
            Thus $x(i_0) - y(i_0)\in\reals$ by \cref{real_sub}.
            Thus $\abs{x(i_0) - y(i_0)} \geq 0$ by \cref{real_abs_nonnegative}.
            Thus $m \geq 0$.
            $0 < m$ or $m = 0$ by \cref{real_leq_split}.
            \begin{byCase}
            \caseOf{$m = 0$.}
                Thus $\sqrt{n} \rmul m = \sqrt{n} \rmul 0$.
                Thus $\sqrt{n} \rmul 0 = 0 \rmul \sqrt{n}$ by \cref{real_mul_comm}.
                Thus $0 \rmul \sqrt{n} = 0$ by \cref{real_mul_zero}.
                Thus $\sqrt{n} \rmul m = 0$.
                Thus $\sqrt{n} \rmul m \geq 0$.
            \caseOf{$0 < m$.}
                Thus $m$ is positive.
                Show that $\sqrt{n}$ is positive.
                \begin{subproof}
                    Suppose not.
                    Thus $\sqrt{n} = 0$.
                    Thus $n = 0$ by \cref{positive_square_root,real_mul_zero}.
                    Thus $0 < 0$ by \cref{real_lt_subst_right}.
                    Contradiction by \cref{real_lt_irreflexive}.
                \end{subproof}
                Thus $\sqrt{n}$ is positive.
                Thus $\sqrt{n} \rmul m > 0$ by \cref{real_mul_pos_pos_gt}.
                Thus $\sqrt{n} \rmul m \geq 0$.
            \end{byCase}
        \end{subproof}
        Thus $r \leq \sqrt{n} \rmul m$ by \cref{real_nonneg_leq_iff_sq_leq}.
    \end{subproof}
    Thus $m \leq r$ and $r \leq \sqrt{n} \rmul m$.
\end{proof}

% from §20 helper lemma 25: interval bound from absolute value
\begin{lemma}\label{real_interval_from_abs_lt}
    Suppose $a\in\reals$.
    Suppose $b\in\reals$.
    Suppose $\epsilon\in\reals$.
    Suppose $0 < \epsilon$.
    Suppose $\abs{a - b} < \epsilon$.
    Then $a - \epsilon < b$ and $b < a + \epsilon$.
\end{lemma}
\begin{proof}
    Thus $\abs{a - b} \leq \epsilon$.
    Thus $\epsilon$ is nonnegative.
    Thus $b - \epsilon \leq a$ by \cref{real_abs_sub_leq_imp_left_bound}.
    Thus $a \leq b + \epsilon$ by \cref{real_abs_sub_leq_imp_right_bound}.
    $\neg{\epsilon}\in\reals$ by \cref{real_add_inverse}.
    Thus $b - \epsilon\in\reals$ by \cref{real_sub,real_add_closed}.
    $a + \epsilon\in\reals$ by \cref{real_add_closed}.
    $b - \epsilon + \epsilon \leq a + \epsilon$ by \cref{real_leq_add}.
    $b - \epsilon = b + \neg{\epsilon}$ by \cref{real_sub}.
    $(b + \neg{\epsilon}) + \epsilon = b + (\neg{\epsilon} + \epsilon)$ by \cref{real_add_assoc}.
    $\epsilon + \neg{\epsilon} = 0$ by \cref{real_add_inverse}.
    Thus $\neg{\epsilon} + \epsilon = 0$ by \cref{real_add_comm}.
    Thus $b - \epsilon + \epsilon = b + 0$.
    $0\in\reals$ by \cref{real_add_ident}.
    $b + 0 = 0 + b$ by \cref{real_add_comm}.
    $0 + b = b$ by \cref{real_add_ident}.
    Thus $b + 0 = b$.
    Thus $b \leq a + \epsilon$ by \cref{real_leq_trans}.
    $b + \epsilon\in\reals$ by \cref{real_add_closed}.
    $a + \neg{\epsilon}\in\reals$ by \cref{real_add_closed}.
    $a + \neg{\epsilon} \leq (b + \epsilon) + \neg{\epsilon}$ by \cref{real_leq_add}.
    $a + \neg{\epsilon} = a - \epsilon$ by \cref{real_sub}.
    $(b + \epsilon) + \neg{\epsilon} = b + (\epsilon + \neg{\epsilon})$ by \cref{real_add_assoc}.
    $\epsilon + \neg{\epsilon} = 0$ by \cref{real_add_inverse}.
    Thus $(b + \epsilon) + \neg{\epsilon} = b + 0$.
    $0\in\reals$ by \cref{real_add_ident}.
    $b + 0 = 0 + b$ by \cref{real_add_comm}.
    $0 + b = b$ by \cref{real_add_ident}.
    Thus $b + 0 = b$.
    Thus $a - \epsilon \leq b$.
    Show that $b \neq a + \epsilon$.
    \begin{subproof}
        Suppose not.
        Thus $b = a + \epsilon$.
        $b - a = b + \neg{a}$ by \cref{real_sub}.
        $\neg{a}\in\reals$ by \cref{real_add_inverse}.
        $(a + \epsilon) + \neg{a} = (a + \neg{a}) + \epsilon$ by \cref{real_add_assoc,real_add_comm}.
        $a + \neg{a} = 0$ by \cref{real_add_inverse}.
        Thus $b - a = 0 + \epsilon$.
        $0 + \epsilon = \epsilon$ by \cref{real_add_ident}.
        Thus $b - a = \epsilon$.
        $\epsilon \geq 0$.
        Thus $\abs{b - a} = \epsilon$ by \cref{real_abs_nonneg}.
        Thus $\abs{a - b} = \epsilon$ by \cref{real_abs_sub_swap}.
        Thus $\epsilon < \epsilon$ by \cref{real_lt_subst_left}.
        Contradiction by \cref{real_lt_irreflexive}.
    \end{subproof}
    Thus $b < a + \epsilon$ by \cref{real_leq_split}.
    Show that $a - \epsilon \neq b$.
    \begin{subproof}
        Suppose not.
        Thus $a - \epsilon = b$.
        $a - b = a + \neg{b}$ by \cref{real_sub}.
        $a - \epsilon = a + \neg{\epsilon}$ by \cref{real_sub}.
        Thus $b = a + \neg{\epsilon}$.
        Thus $\neg{b} = \neg{(a + \neg{\epsilon})}$.
        $\neg{(a + \neg{\epsilon})} = \neg{a} + \neg{(\neg{\epsilon})}$ by \cref{real_neg_addition}.
        $\neg{(\neg{\epsilon})} = \epsilon$ by \cref{real_neg_neg}.
        Thus $\neg{b} = \neg{a} + \epsilon$.
        Thus $a - b = a + (\neg{a} + \epsilon)$.
        $\neg{a}\in\reals$ by \cref{real_add_inverse}.
        $\epsilon\in\reals$.
        $a + (\neg{a} + \epsilon) = (a + \neg{a}) + \epsilon$ by \cref{real_add_assoc}.
        $a + \neg{a} = 0$ by \cref{real_add_inverse}.
        Thus $a - b = 0 + \epsilon$.
        $0 + \epsilon = \epsilon$ by \cref{real_add_ident}.
        Thus $a - b = \epsilon$.
        $\epsilon \geq 0$.
        Thus $\abs{a - b} = \epsilon$ by \cref{real_abs_nonneg}.
        Thus $\epsilon < \epsilon$ by \cref{real_lt_subst_left}.
        Contradiction by \cref{real_lt_irreflexive}.
    \end{subproof}
    Thus $a - \epsilon < b$ by \cref{real_leq_split}.
    Thus $a - \epsilon < b$ and $b < a + \epsilon$.
\end{proof}

% from §20 helper lemma 26: coordinate square sequence is a function
\begin{lemma}\label{coord_square_sequence_function}
    Suppose $n\in\naturalsPlus$.
    Let $J = \initialSegment{n}$.
    Suppose $x\in\realsPower{J}$.
    Suppose $y\in\realsPower{J}$.
    Then $\coordSquareSequence{n}{(x,y)}\in\funs{J}{\reals}$.
\end{lemma}
\begin{proof}
    Let $a = \coordSquareSequence{n}{(x,y)}$.
    Show that $a$ is a function.
    \begin{subproof}
        Show that $a$ is a relation.
        \begin{subproof}
            Show that for all $w$ such that $w\in a$ there exist $i,t$ such that $w = (i,t)$.
            \begin{subproof}
                Fix $w$.
                Assume $w\in a$.
                Take $i\in\initialSegment{n}$ such that
                    $w = (i,\exp{\apply{\fst{(x,y)}}{i} - \apply{\snd{(x,y)}}{i}}{1+1})$
                    by \cref{coord_square_sequence}.
                Thus there exist $i_0,t$ such that $w = (i_0,t)$.
            \end{subproof}
            Thus $a$ is a relation by \cref{relation}.
        \end{subproof}
        Show that for all $i,t_1,t_2$ such that $(i,t_1)\in a$ and $(i,t_2)\in a$ we have $t_1 = t_2$.
        \begin{subproof}
            Fix $i$.
            Fix $t_1$.
            Fix $t_2$.
            Assume $(i,t_1)\in a$ and $(i,t_2)\in a$.
            Take $i_1\in\initialSegment{n}$ such that
                $(i,t_1) = (i_1,\exp{\apply{\fst{(x,y)}}{i_1} - \apply{\snd{(x,y)}}{i_1}}{1+1})$
                by \cref{coord_square_sequence}.
            Take $i_2\in\initialSegment{n}$ such that
                $(i,t_2) = (i_2,\exp{\apply{\fst{(x,y)}}{i_2} - \apply{\snd{(x,y)}}{i_2}}{1+1})$
                by \cref{coord_square_sequence}.
            Thus $i = i_1$ and $t_1 = \exp{\apply{\fst{(x,y)}}{i_1} - \apply{\snd{(x,y)}}{i_1}}{1+1}$ by \cref{pair_eq_iff}.
            Thus $i = i_2$ and $t_2 = \exp{\apply{\fst{(x,y)}}{i_2} - \apply{\snd{(x,y)}}{i_2}}{1+1}$ by \cref{pair_eq_iff}.
            Thus $t_1 = \exp{\apply{\fst{(x,y)}}{i} - \apply{\snd{(x,y)}}{i}}{1+1}$.
            Thus $t_2 = \exp{\apply{\fst{(x,y)}}{i} - \apply{\snd{(x,y)}}{i}}{1+1}$.
            Thus $t_1 = t_2$.
        \end{subproof}
        Thus $a$ is right-unique by \cref{rightunique}.
        Thus $a$ is a function.
    \end{subproof}
    Show that $\dom{a} = J$.
    \begin{subproof}
        Show that for all $i$ such that $i\in J$ we have $i\in\dom{a}$.
        \begin{subproof}
            Fix $i$.
            Assume $i\in J$.
            Thus $(i,\exp{x(i) - y(i)}{1+1})\in a$ by \cref{coord_square_sequence_elem}.
            Thus $i\in\dom{a}$ by \cref{dom_intro}.
        \end{subproof}
        Thus $J\subseteq\dom{a}$ by \cref{subseteq}.
        Show that for all $i$ such that $i\in\dom{a}$ we have $i\in J$.
        \begin{subproof}
            Fix $i$.
            Assume $i\in\dom{a}$.
            Take $w$ such that $(i,w)\in a$ by \cref{dom_iff}.
            Take $i_0\in\initialSegment{n}$ such that
                $(i,w) = (i_0,\exp{\apply{\fst{(x,y)}}{i_0} - \apply{\snd{(x,y)}}{i_0}}{1+1})$
                by \cref{coord_square_sequence}.
            Thus $i = i_0$ by \cref{pair_eq_iff}.
            Thus $i\in J$.
        \end{subproof}
        Thus $\dom{a}\subseteq J$ by \cref{subseteq}.
        Thus $\dom{a} = J$ by \cref{subseteq_antisymmetric}.
    \end{subproof}
    Show that for all $i\in\dom{a}$ we have $a(i)\in\reals$.
    \begin{subproof}
        Fix $i$.
        Assume $i\in\dom{a}$.
        Thus $i\in J$ by \cref{subseteq}.
        $x\in\tuplePower{\reals}{J}$ and $y\in\tuplePower{\reals}{J}$ by \cref{reals_power}.
        Thus $x\in\funs{J}{\reals}$ and $y\in\funs{J}{\reals}$ by \cref{tuple_power}.
        Thus $x(i)\in\reals$ and $y(i)\in\reals$ by \cref{funs_type_apply}.
        Thus $\neg{y(i)}\in\reals$ by \cref{real_add_inverse}.
        Thus $x(i) + \neg{y(i)}\in\reals$ by \cref{real_add_closed}.
        Thus $x(i) - y(i)\in\reals$ by \cref{real_sub}.
        $1\in\naturalsPlus$ by \cref{real_one_in_naturals}.
        $\exp{x(i) - y(i)}{1+1} = \exp{x(i) - y(i)}{1} \rmul (x(i) - y(i))$ by \cref{real_pow_succ}.
        $\exp{x(i) - y(i)}{1} = x(i) - y(i)$ by \cref{real_pow_one}.
        Thus $\exp{x(i) - y(i)}{1+1} = (x(i) - y(i)) \rmul (x(i) - y(i))$.
        Thus $\exp{x(i) - y(i)}{1+1}\in\reals$ by \cref{real_mul_closed}.
        $(i,\exp{x(i) - y(i)}{1+1})\in a$ by \cref{coord_square_sequence_elem}.
        Thus $a(i) = \exp{x(i) - y(i)}{1+1}$ by \cref{function_apply_intro}.
        Thus $a(i)\in\reals$.
    \end{subproof}
    Thus $a\in\funs{J}{\reals}$ by \cref{funs_intro}.
\end{proof}

% from §20 helper lemma 27: sum finite sequence successor closure
\begin{lemma}\label{sum_finite_sequence_exists_succ}
    Let $A = \{n\in\naturalsPlus \mid \text{for all $a\in\funs{\initialSegment{n}}{\reals}$ there exists $s\in\reals$ such that $\sumFiniteSequence{a}{s}$}\}$.
    Then for all $n\in A$ we have $n + 1\in A$.
\end{lemma}
\begin{proof}
        Show that for all $n\in A$ we have $n + 1\in A$.
        \begin{subproof}
        Fix $n$.
        Assume $n\in A$.
        Thus for all $a\in\funs{\initialSegment{n}}{\reals}$ there exists $s\in\reals$
            such that $\sumFiniteSequence{a}{s}$.
        Then $n\in\naturalsPlus$.
        $n + 1\in\naturalsPlus$ by \cref{real_naturals_closed_succ}.
        Show that $n + 1\in A$.
            \begin{subproof}
                Show that for all $a\in\funs{\initialSegment{n + 1}}{\reals}$ there exists $s\in\reals$ such that $\sumFiniteSequence{a}{s}$.
                \begin{subproof}
                Fix $a$.
                Assume $a\in\funs{\initialSegment{n + 1}}{\reals}$.
                Let $a_0 = \{(i,a(i)) \mid i\in\initialSegment{n}\}$.
                Show that $a_0\in\funs{\initialSegment{n}}{\reals}$.
                \begin{subproof}
                    Show that $a_0$ is a function.
                    \begin{subproof}
                        Show that $a_0$ is a relation.
                        \begin{subproof}
                            Show that for all $w$ such that $w\in a_0$ there exist $i,t$ such that $w = (i,t)$.
                            \begin{subproof}
                                Fix $w$.
                                Assume $w\in a_0$.
                                Take $i\in\initialSegment{n}$ such that $w = (i,a(i))$.
                                Thus there exist $i_0,t$ such that $w = (i_0,t)$.
                            \end{subproof}
                            Thus $a_0$ is a relation by \cref{relation}.
                        \end{subproof}
                        Show that for all $i,t_1,t_2$ such that $(i,t_1)\in a_0$ and $(i,t_2)\in a_0$ we have $t_1 = t_2$.
                        \begin{subproof}
                            Fix $i$.
                            Fix $t_1$.
                            Fix $t_2$.
                            Assume $(i,t_1)\in a_0$ and $(i,t_2)\in a_0$.
                            Then $t_1 = a(i)$ and $t_2 = a(i)$.
                            Thus $t_1 = t_2$.
                        \end{subproof}
                        Thus $a_0$ is right-unique by \cref{rightunique}.
                        Thus $a_0$ is a function.
                    \end{subproof}
                    Show that $\dom{a_0} = \initialSegment{n}$.
                    \begin{subproof}
                        Show that for all $i$ such that $i\in\initialSegment{n}$ we have $i\in\dom{a_0}$.
                        \begin{subproof}
                            Fix $i$.
                            Assume $i\in\initialSegment{n}$.
                            Thus $(i,a(i))\in a_0$.
                            Thus $i\in\dom{a_0}$ by \cref{dom_intro}.
                        \end{subproof}
                        Thus $\initialSegment{n}\subseteq\dom{a_0}$ by \cref{subseteq}.
                        Show that for all $i$ such that $i\in\dom{a_0}$ we have $i\in\initialSegment{n}$.
                        \begin{subproof}
                            Fix $i$.
                            Assume $i\in\dom{a_0}$.
                            Take $w$ such that $(i,w)\in a_0$ by \cref{dom_iff}.
                            Thus $i\in\initialSegment{n}$.
                        \end{subproof}
                        Thus $\dom{a_0}\subseteq\initialSegment{n}$ by \cref{subseteq}.
                        Thus $\dom{a_0} = \initialSegment{n}$ by \cref{subseteq_antisymmetric}.
                    \end{subproof}
                    Show that for all $i\in\dom{a_0}$ we have $a_0(i)\in\reals$.
                    \begin{subproof}
                        Fix $i$.
                        Assume $i\in\dom{a_0}$.
                        Thus $i\in\initialSegment{n}$ by \cref{subseteq}.
                        Thus $i\in\naturalsPlus$ by \cref{initial_segment_naturalsplus}.
                        Thus $i < n$ or $i = n$ by \cref{initial_segment_naturalsplus}.
                        $i\in\reals$ by \cref{real_naturals_subset_reals,subseteq}.
                        $n\in\reals$ by \cref{real_naturals_subset_reals,subseteq}.
                        $1\in\reals$ by \cref{real_naturals_subset_reals,real_one_in_naturals,subseteq}.
                        $n + 1\in\reals$ by \cref{real_add_closed}.
                        $0\in\reals$ by \cref{real_add_ident}.
                        $0 < 1$ by \cref{real_zero_lt_one}.
                        $0 + n < 1 + n$ by \cref{real_lt_add}.
                        $0 + n = n$ by \cref{real_add_ident}.
                        $1 + n = n + 1$ by \cref{real_add_comm}.
                        Thus $n < n + 1$ by \cref{real_lt_subst_left}.
                        Show that $i < n + 1$.
                        \begin{subproof}
                            \begin{byCase}
                            \caseOf{$i < n$.}
                                Thus $i < n + 1$ by \cref{real_lt_trans}.
                            \caseOf{$i = n$.}
                                Thus $i < n + 1$ by \cref{real_lt_subst_left}.
                            \end{byCase}
                        \end{subproof}
                        Thus $i\in\initialSegment{n + 1}$ by \cref{initial_segment_naturalsplus}.
                        Thus $i\in\dom{a}$ by \cref{funs_elim}.
                        Thus $a(i)\in\reals$ by \cref{funs_type_apply}.
                        $(i,a(i))\in a_0$.
                        Thus $a_0(i) = a(i)$ by \cref{function_apply_intro}.
                        Thus $a_0(i)\in\reals$.
                    \end{subproof}
                    Thus $a_0\in\funs{\initialSegment{n}}{\reals}$ by \cref{funs_intro}.
                \end{subproof}
                Thus there exists $s_0\in\reals$ such that $\sumFiniteSequence{a_0}{s_0}$.
                Take $s_0\in\reals$ such that $\sumFiniteSequence{a_0}{s_0}$.
                Take $n_0$ such that $n_0\in\naturalsPlus$ and $a_0\in\funs{\initialSegment{n_0}}{\reals}$ and
                    there exists $t_0\in\funs{\initialSegment{n_0}}{\reals}$ such that
                        $t_0(1) = a_0(1)$ and $s_0 = t_0(n_0)$ and
                        for all $k\in\naturalsPlus$ such that $k + 1 \leq n_0$ we have $t_0(k + 1) = t_0(k) + a_0(k + 1)$
                    by \cref{sum_finite_sequence_pred}.
                Take $t_0$ such that $t_0\in\funs{\initialSegment{n_0}}{\reals}$ and
                    $t_0(1) = a_0(1)$ and $s_0 = t_0(n_0)$ and
                    for all $k\in\naturalsPlus$ such that $k + 1 \leq n_0$ we have $t_0(k + 1) = t_0(k) + a_0(k + 1)$.
                Thus $\dom{a_0} = \initialSegment{n_0}$ by \cref{funs_elim}.
                Show that $\initialSegment{n_0} = \initialSegment{n}$.
                \begin{subproof}
                    Show that $\initialSegment{n_0}\subseteq\initialSegment{n}$.
                    \begin{subproof}
                        Show that for all $i$ such that $i\in\initialSegment{n_0}$ we have $i\in\initialSegment{n}$.
                        \begin{subproof}
                            Fix $i$.
                            Assume $i\in\initialSegment{n_0}$.
                            Thus $i\in\dom{a_0}$ by \cref{funs_elim}.
                            Thus $i\in\initialSegment{n}$.
                        \end{subproof}
                        Thus $\initialSegment{n_0}\subseteq\initialSegment{n}$ by \cref{subseteq}.
                    \end{subproof}
                    Show that $\initialSegment{n}\subseteq\initialSegment{n_0}$.
                    \begin{subproof}
                        Show that for all $i$ such that $i\in\initialSegment{n}$ we have $i\in\initialSegment{n_0}$.
                        \begin{subproof}
                            Fix $i$.
                            Assume $i\in\initialSegment{n}$.
                            Thus $i\in\dom{a_0}$.
                            Thus $i\in\initialSegment{n_0}$ by \cref{funs_elim}.
                        \end{subproof}
                        Thus $\initialSegment{n}\subseteq\initialSegment{n_0}$ by \cref{subseteq}.
                    \end{subproof}
                    Thus $\initialSegment{n_0} = \initialSegment{n}$ by \cref{subseteq_antisymmetric}.
                \end{subproof}
                Thus $\dom{t_0} = \initialSegment{n}$ by \cref{funs_elim}.
                $t_0$ is a function by \cref{funs_elim}.
                $t_0\in\funs{\initialSegment{n}}{\reals}$ by \cref{funs_intro}.
                $n\in\initialSegment{n}$ by \cref{initial_segment_naturalsplus}.
                Thus $t_0(n)\in\reals$ by \cref{funs_type_apply}.
                $n + 1\in\initialSegment{n + 1}$ by \cref{initial_segment_naturalsplus}.
                Thus $a(n + 1)\in\reals$ by \cref{funs_type_apply}.
                Let $v = t_0(n) + a(n + 1)$.
                $v\in\reals$ by \cref{real_add_closed}.
                Let $g = \{(n + 1,v)\}$.
                Show that $g$ is a function.
                \begin{subproof}
                    Show that $g$ is a relation.
                    \begin{subproof}
                        Show that for all $w$ such that $w\in g$ there exist $x,y$ such that $w = (x,y)$.
                        \begin{subproof}
                            Fix $w$.
                            Assume $w\in g$.
                            Thus $w = (n + 1,v)$ by \cref{singleton_elim}.
                            Thus there exist $x,y$ such that $w = (x,y)$.
                        \end{subproof}
                        Thus $g$ is a relation by \cref{relation}.
                    \end{subproof}
                    Show that for all $x,y_1,y_2$ such that $(x,y_1)\in g$ and $(x,y_2)\in g$ we have $y_1 = y_2$.
                    \begin{subproof}
                        Fix $x$.
                        Fix $y_1$.
                        Fix $y_2$.
                        Assume $(x,y_1)\in g$ and $(x,y_2)\in g$.
                        Thus $(x,y_1) = (n + 1,v)$ and $(x,y_2) = (n + 1,v)$ by \cref{singleton_elim}.
                        Thus $y_1 = v$ and $y_2 = v$ by \cref{pair_eq_iff}.
                        Thus $y_1 = y_2$.
                    \end{subproof}
                    Thus $g$ is right-unique by \cref{rightunique}.
                    Thus $g$ is a function.
                \end{subproof}
                Show that for all $x$ such that $x\in\dom{t_0}\inter\dom{g}$ we have $t_0(x) = g(x)$.
                \begin{subproof}
                    Fix $x$.
                    Assume $x\in\dom{t_0}\inter\dom{g}$.
                    Suppose not.
                    Thus $x\in\dom{g}$.
                    Take $y$ such that $(x,y)\in g$ by \cref{dom_iff}.
                    Thus $(x,y) = (n + 1,v)$ by \cref{singleton_elim}.
                    Thus $x = n + 1$ by \cref{pair_eq_iff}.
                    Thus $x\in\dom{t_0}$.
                    Thus $x\in\initialSegment{n}$ by \cref{funs_elim}.
                    Thus $x\in\naturalsPlus$ by \cref{initial_segment_naturalsplus}.
                    Thus $x < n$ or $x = n$ by \cref{initial_segment_naturalsplus}.
                    $n\in\naturalsPlus$.
                    $x\in\reals$ by \cref{real_naturals_subset_reals,subseteq}.
                    $n\in\reals$ by \cref{real_naturals_subset_reals,subseteq}.
                    $1\in\reals$ by \cref{real_naturals_subset_reals,real_one_in_naturals,subseteq}.
                    $0\in\reals$ by \cref{real_add_ident}.
                    $0 < 1$ by \cref{real_zero_lt_one}.
                    $0 + n < 1 + n$ by \cref{real_lt_add}.
                    $0 + n = n$ by \cref{real_add_ident}.
                    $1 + n = n + 1$ by \cref{real_add_comm}.
                    Thus $n < n + 1$ by \cref{real_lt_subst_left}.
                    Thus $n < x$ by \cref{real_lt_subst_right}.
                    Show that $n < n$.
                    \begin{subproof}
                        \begin{byCase}
                        \caseOf{$x < n$.}
                            Thus $n < n$ by \cref{real_lt_trans}.
                        \caseOf{$x = n$.}
                            Thus $n < n$ by \cref{real_lt_subst_right}.
                        \end{byCase}
                    \end{subproof}
                    Contradiction by \cref{real_lt_irreflexive}.
                \end{subproof}
                Let $t = t_0\union g$.
                $t$ is a function by \cref{union_compatible_functions}.
                Show that $\dom{t} = \initialSegment{n}\union\{n + 1\}$.
                \begin{subproof}
                    Show that for all $i$ such that $i\in\dom{t_0}\union\dom{g}$ we have $i\in\dom{t}$.
                    \begin{subproof}
                        Fix $i$.
                        Assume $i\in\dom{t_0}\union\dom{g}$.
                        Thus $i\in\dom{t_0}$ or $i\in\dom{g}$ by \cref{union_iff}.
                        \begin{byCase}
                        \caseOf{$i\in\dom{t_0}$.}
                            Take $w$ such that $(i,w)\in t_0$ by \cref{dom_iff}.
                            Thus $(i,w)\in t$ by \cref{union_iff}.
                            Thus $i\in\dom{t}$ by \cref{dom_intro}.
                        \caseOf{$i\in\dom{g}$.}
                            Take $w$ such that $(i,w)\in g$ by \cref{dom_iff}.
                            Thus $(i,w)\in t$ by \cref{union_iff}.
                            Thus $i\in\dom{t}$ by \cref{dom_intro}.
                        \end{byCase}
                    \end{subproof}
                    Thus $\dom{t_0}\union\dom{g}\subseteq\dom{t}$ by \cref{subseteq}.
                    Show that for all $i$ such that $i\in\dom{t}$ we have $i\in\dom{t_0}\union\dom{g}$.
                    \begin{subproof}
                        Fix $i$.
                        Assume $i\in\dom{t}$.
                        Take $w$ such that $(i,w)\in t$ by \cref{dom_iff}.
                        Thus $(i,w)\in t_0$ or $(i,w)\in g$ by \cref{union_iff}.
                        \begin{byCase}
                        \caseOf{$(i,w)\in t_0$.} Thus $i\in\dom{t_0}$ by \cref{dom_intro}.
                        Thus $i\in\dom{t_0}\union\dom{g}$ by \cref{union_iff}.
                        \caseOf{$(i,w)\in g$.} Thus $i\in\dom{g}$ by \cref{dom_intro}.
                        Thus $i\in\dom{t_0}\union\dom{g}$ by \cref{union_iff}.
                        \end{byCase}
                    \end{subproof}
                    Thus $\dom{t}\subseteq\dom{t_0}\union\dom{g}$ by \cref{subseteq}.
                    Thus $\dom{t} = \dom{t_0}\union\dom{g}$ by \cref{subseteq_antisymmetric}.
                    Thus $\dom{t_0} = \initialSegment{n}$ by \cref{funs_elim}.
                    Show that $\dom{g} = \{n + 1\}$.
                    \begin{subproof}
                        Show that for all $i$ such that $i\in\dom{g}$ we have $i\in\{n + 1\}$.
                        \begin{subproof}
                            Fix $i$.
                            Assume $i\in\dom{g}$.
                            Take $y$ such that $(i,y)\in g$ by \cref{dom_iff}.
                            Thus $(i,y) = (n + 1,v)$ by \cref{singleton_elim}.
                            Thus $i = n + 1$ by \cref{pair_eq_iff}.
                            Thus $n + 1\in\{n + 1\}$ by \cref{singleton_intro}.
                            Thus $i\in\{n + 1\}$.
                        \end{subproof}
                        Thus $\dom{g}\subseteq\{n + 1\}$ by \cref{subseteq}.
                        Show that for all $i$ such that $i\in\{n + 1\}$ we have $i\in\dom{g}$.
                        \begin{subproof}
                            Fix $i$.
                            Assume $i\in\{n + 1\}$.
                            Thus $i = n + 1$ by \cref{singleton_elim}.
                            Thus $(n + 1,v)\in g$ by \cref{singleton_intro}.
                            Thus $(i,v) = (n + 1,v)$ by \cref{pair_eq_iff}.
                            Thus $(i,v)\in g$.
                            Thus $i\in\dom{g}$ by \cref{dom_intro}.
                        \end{subproof}
                        Thus $\{n + 1\}\subseteq\dom{g}$ by \cref{subseteq}.
                        Thus $\dom{g} = \{n + 1\}$ by \cref{subseteq_antisymmetric}.
                    \end{subproof}
                    Thus $\dom{t} = \initialSegment{n}\union\{n + 1\}$.
                \end{subproof}
                Show that for all $i\in\dom{t}$ we have $t(i)\in\reals$.
                \begin{subproof}
                    Fix $i$.
                    Assume $i\in\dom{t}$.
                    Thus $i\in\dom{t_0}\union\dom{g}$.
                    Thus $i\in\dom{t_0}$ or $i\in\dom{g}$ by \cref{union_iff}.
                    \begin{byCase}
                    \caseOf{$i\in\dom{t_0}$.}
                        Thus $t_0(i)\in\reals$ by \cref{funs_type_apply}.
                        $(i,t_0(i))\in t_0$ by \cref{function_apply_elim}.
                        Thus $(i,t_0(i))\in t$ by \cref{union_iff}.
                        Thus $t(i) = t_0(i)$ by \cref{function_apply_intro}.
                        Thus $t(i)\in\reals$.
                    \caseOf{$i\in\dom{g}$.}
                        Take $y$ such that $(i,y)\in g$ by \cref{dom_iff}.
                        Thus $(i,y) = (n + 1,v)$ by \cref{singleton_elim}.
                        Thus $i = n + 1$ and $y = v$ by \cref{pair_eq_iff}.
                        Thus $(i,v)\in g$.
                        Thus $(i,v)\in t$ by \cref{union_iff}.
                        Thus $t(i) = v$ by \cref{function_apply_intro}.
                        Thus $t(i)\in\reals$.
                    \end{byCase}
                \end{subproof}
                Show that $t(1) = a(1)$.
                \begin{subproof}
                    $1\in\naturalsPlus$ by \cref{real_one_in_naturals}.
                    $1 \leq n$ by \cref{real_one_leq_naturalsplus}.
                    Thus $1\in\initialSegment{n}$ by \cref{initial_segment_naturalsplus}.
                    Thus $1\in\dom{t_0}$ by \cref{funs_elim}.
                    $(1,t_0(1))\in t_0$ by \cref{function_apply_elim}.
                    Thus $(1,t_0(1))\in t$ by \cref{union_iff}.
                    Thus $t(1) = t_0(1)$ by \cref{function_apply_intro}.
                    $(1,a(1))\in a_0$.
                    $a_0$ is a function by \cref{funs_elim}.
                    Thus $a_0(1) = a(1)$ by \cref{function_apply_intro}.
                    Thus $t(1) = a(1)$.
                \end{subproof}
                Show that for all $k\in\naturalsPlus$ such that $k + 1 \leq n + 1$ we have
                    $t(k + 1) = t(k) + a(k + 1)$.
                \begin{subproof}
                    Fix $k$.
                    Assume $k\in\naturalsPlus$.
                    Assume $k + 1 \leq n + 1$.
                    $k + 1\in\naturalsPlus$ by \cref{real_naturals_closed_succ}.
                    $k + 1 = n + 1$ or $k + 1 \neq n + 1$.
                    \begin{byCase}
                    \caseOf{$k + 1 = n + 1$.}
                        $k\in\reals$ and $n\in\reals$ and $1\in\reals$ by \cref{real_naturals_subset_reals,subseteq,real_one_in_naturals}.
                        Thus $k = n$ by \cref{real_add_cancel}.
                        Thus $k\in\initialSegment{n}$ by \cref{initial_segment_naturalsplus}.
                        Thus $k\in\dom{t_0}$ by \cref{funs_elim}.
                        $(k,t_0(k))\in t_0$ by \cref{function_apply_elim}.
                        Thus $(k,t_0(k))\in t$ by \cref{union_iff}.
                        Thus $t(k) = t_0(k)$ by \cref{function_apply_intro}.
                        $n\in\dom{t_0}$ by \cref{funs_elim,initial_segment_naturalsplus}.
                        $t_0(n)\in\reals$ by \cref{funs_type_apply}.
                        Thus $t(k) = t_0(n)$.
                        $(n + 1,v)\in g$ by \cref{singleton_intro}.
                        Thus $(k + 1,v) = (n + 1,v)$ by \cref{pair_eq_iff}.
                        Thus $(k + 1,v)\in g$.
                        Thus $(k + 1,v)\in t$ by \cref{union_iff}.
                        Thus $t(k + 1) = v$ by \cref{function_apply_intro}.
                        Thus $t(k + 1) = t(k) + a(k + 1)$.
                    \caseOf{$k + 1 \neq n + 1$.}
                        Thus $k + 1 \leq n$ by \cref{real_naturals_leq_succ_elim}.
                        Thus $k + 1\in\initialSegment{n}$ by \cref{initial_segment_naturalsplus}.
                        $k\in\reals$ by \cref{real_naturals_subset_reals,subseteq}.
                        $n\in\reals$ by \cref{real_naturals_subset_reals,subseteq}.
                        $1\in\reals$ by \cref{real_naturals_subset_reals,real_one_in_naturals,subseteq}.
                        $k + 1\in\reals$ by \cref{real_add_closed}.
                        $0\in\reals$ by \cref{real_add_ident}.
                        $0 < 1$ by \cref{real_zero_lt_one}.
                        $0 + k < 1 + k$ by \cref{real_lt_add}.
                        $0 + k = k$ by \cref{real_add_ident}.
                        $1 + k = k + 1$ by \cref{real_add_comm}.
                        Thus $k < k + 1$ by \cref{real_lt_subst_left}.
                        Thus $k + 1 < n$ or $k + 1 = n$ by \cref{initial_segment_naturalsplus}.
                        Show that $k < n$.
                        \begin{subproof}
                            \begin{byCase}
                            \caseOf{$k + 1 < n$.}
                                Thus $k < n$ by \cref{real_lt_trans}.
                            \caseOf{$k + 1 = n$.}
                                Thus $k < n$ by \cref{real_lt_subst_right}.
                            \end{byCase}
                        \end{subproof}
                        Thus $k\in\initialSegment{n}$ by \cref{initial_segment_naturalsplus}.
                        $k\in\dom{t_0}$ and $k + 1\in\dom{t_0}$ by \cref{funs_elim,subseteq}.
                        $(k,t_0(k))\in t_0$ and $(k + 1,t_0(k + 1))\in t_0$ by \cref{function_apply_elim}.
                        Thus $(k,t_0(k))\in t$ and $(k + 1,t_0(k + 1))\in t$ by \cref{union_iff}.
                        Thus $t(k) = t_0(k)$ and $t(k + 1) = t_0(k + 1)$ by \cref{function_apply_intro}.
                        Thus $k + 1\in\initialSegment{n_0}$.
                        Thus $k + 1 < n_0$ or $k + 1 = n_0$ by \cref{initial_segment_naturalsplus}.
                        Thus $t_0(k + 1) = t_0(k) + a_0(k + 1)$.
                        $(k + 1,a(k + 1))\in a_0$.
                        $a_0$ is a function by \cref{funs_elim}.
                        Thus $a_0(k + 1) = a(k + 1)$ by \cref{function_apply_intro}.
                        Thus $t(k + 1) = t(k) + a(k + 1)$.
                    \end{byCase}
                \end{subproof}
                Let $s = t(n + 1)$.
                $(n + 1,v)\in g$ by \cref{singleton_intro}.
                Thus $(n + 1,v)\in t$ by \cref{union_iff}.
                Thus $n + 1\in\dom{t}$ by \cref{dom_intro}.
                Thus $s\in\reals$ by \cref{funs_type_apply}.
                $n + 1\in\naturalsPlus$ by \cref{real_naturals_closed_succ}.
                Show that $\dom{t} = \initialSegment{n + 1}$.
                \begin{subproof}
                    Show that $\initialSegment{n}\union\{n + 1\} = \initialSegment{n + 1}$.
                    \begin{subproof}
                        Show that $\initialSegment{n}\union\{n + 1\}\subseteq\initialSegment{n + 1}$.
                        \begin{subproof}
                            Show that for all $i$ such that $i\in\initialSegment{n}\union\{n + 1\}$ we have $i\in\initialSegment{n + 1}$.
                            \begin{subproof}
                                Fix $i$.
                                Assume $i\in\initialSegment{n}\union\{n + 1\}$.
                                Thus $i\in\initialSegment{n}$ or $i\in\{n + 1\}$ by \cref{union_iff}.
                                \begin{byCase}
                                \caseOf{$i\in\initialSegment{n}$.}
                                    Thus $i\in\naturalsPlus$ by \cref{initial_segment_naturalsplus}.
                                    Thus $i < n$ or $i = n$ by \cref{initial_segment_naturalsplus}.
                                    $i\in\reals$ by \cref{real_naturals_subset_reals,subseteq}.
                                    $n\in\reals$ by \cref{real_naturals_subset_reals,subseteq}.
                                    $1\in\reals$ by \cref{real_naturals_subset_reals,real_one_in_naturals,subseteq}.
                                    $0\in\reals$ by \cref{real_add_ident}.
                                    $0 < 1$ by \cref{real_zero_lt_one}.
                                    $0 + n < 1 + n$ by \cref{real_lt_add}.
                                    $0 + n = n$ by \cref{real_add_ident}.
                                    $1 + n = n + 1$ by \cref{real_add_comm}.
                                    $n + 1\in\reals$ by \cref{real_add_closed}.
                                    Thus $n < n + 1$ by \cref{real_lt_subst_left}.
                                    Show that $i < n + 1$.
                                    \begin{subproof}
                                        \begin{byCase}
                                        \caseOf{$i < n$.}
                                            Thus $i < n + 1$ by \cref{real_lt_trans}.
                                        \caseOf{$i = n$.}
                                            Thus $i < n + 1$ by \cref{real_lt_subst_left}.
                                        \end{byCase}
                                    \end{subproof}
                                    Thus $i\in\initialSegment{n + 1}$ by \cref{initial_segment_naturalsplus}.
                                \caseOf{$i\in\{n + 1\}$.}
                                    Thus $i = n + 1$ by \cref{singleton_elim}.
                                    Thus $i\in\initialSegment{n + 1}$ by \cref{initial_segment_naturalsplus}.
                                \end{byCase}
                            \end{subproof}
                            Thus $\initialSegment{n}\union\{n + 1\}\subseteq\initialSegment{n + 1}$ by \cref{subseteq}.
                        \end{subproof}
                        Show that $\initialSegment{n + 1}\subseteq\initialSegment{n}\union\{n + 1\}$.
                        \begin{subproof}
                            Show that for all $i$ such that $i\in\initialSegment{n + 1}$ we have $i\in\initialSegment{n}\union\{n + 1\}$.
                            \begin{subproof}
                                Fix $i$.
                                Assume $i\in\initialSegment{n + 1}$.
                                Thus $i\in\naturalsPlus$ by \cref{initial_segment_naturalsplus}.
                                Thus $i < n + 1$ or $i = n + 1$ by \cref{initial_segment_naturalsplus}.
                                \begin{byCase}
                                \caseOf{$i = n + 1$.}
                                    Thus $i\in\{n + 1\}$ by \cref{singleton_intro}.
                                    Thus $i\in\initialSegment{n}\union\{n + 1\}$ by \cref{union_iff}.
                                \caseOf{$i < n + 1$.}
                                    Thus $i \leq n + 1$.
                                    Thus $i \neq n + 1$.
                                    Thus $i \leq n$ by \cref{real_naturals_leq_succ_elim}.
                                    Thus $i\in\initialSegment{n}$ by \cref{initial_segment_naturalsplus}.
                                    Thus $i\in\initialSegment{n}\union\{n + 1\}$ by \cref{union_iff}.
                                \end{byCase}
                            \end{subproof}
                            Thus $\initialSegment{n + 1}\subseteq\initialSegment{n}\union\{n + 1\}$ by \cref{subseteq}.
                        \end{subproof}
                        Thus $\initialSegment{n}\union\{n + 1\} = \initialSegment{n + 1}$ by \cref{subseteq_antisymmetric}.
                    \end{subproof}
                    Thus $\dom{t} = \initialSegment{n + 1}$.
                \end{subproof}
                Thus $t\in\funs{\initialSegment{n + 1}}{\reals}$ by \cref{funs_intro}.
                Thus $\sumFiniteSequence{a}{s}$ by \cref{sum_finite_sequence_pred}.
            \end{subproof}
            Thus $n + 1\in A$.
        \end{subproof}
        \end{subproof}
\end{proof}

% from §20 helper lemma 28: sum of a finite sequence exists
\begin{lemma}\label{sum_finite_sequence_exists}
    Suppose $m\in\naturalsPlus$.
    Suppose $f\in\funs{\initialSegment{m}}{\reals}$.
    Then there exists $s\in\reals$ such that $\sumFiniteSequence{f}{s}$.
\end{lemma}
\begin{proof}
    Let $A = \{n\in\naturalsPlus \mid \text{for all $a\in\funs{\initialSegment{n}}{\reals}$ there exists $s\in\reals$ such that $\sumFiniteSequence{a}{s}$}\}$.
        Show that $A\subseteq\naturalsPlus$.
        \begin{subproof}
            Show that for all $n$ such that $n\in A$ we have $n\in\naturalsPlus$.
            \begin{subproof}
                Fix $n$.
                Assume $n\in A$.
                Then $n\in\naturalsPlus$.
            \end{subproof}
            Thus $A\subseteq\naturalsPlus$ by \cref{subseteq}.
        \end{subproof}
        $1\in\naturalsPlus$ by \cref{real_one_in_naturals}.
        Show that $1\in A$.
        \begin{subproof}
            Show that for all $a\in\funs{\initialSegment{1}}{\reals}$ there exists $s\in\reals$ such that $\sumFiniteSequence{a}{s}$.
            \begin{subproof}
                Fix $a$.
                Assume $a\in\funs{\initialSegment{1}}{\reals}$.
                Let $t = a$.
                $t\in\funs{\initialSegment{1}}{\reals}$.
                Thus $\dom{t} = \initialSegment{1}$ by \cref{funs_elim}.
                $1 \leq 1$.
                Thus $1\in\initialSegment{1}$ by \cref{initial_segment_naturalsplus,real_one_in_naturals}.
                Thus $t(1)\in\reals$ by \cref{funs_type_apply}.
                Let $s = t(1)$.
                Show that $\sumFiniteSequence{a}{s}$.
                \begin{subproof}
                    Show that for all $k\in\naturalsPlus$ such that $k + 1 \leq 1$ we have $t(k + 1) = t(k) + a(k + 1)$.
                    \begin{subproof}
                        Fix $k$.
                        Assume $k\in\naturalsPlus$.
                        Assume $k + 1 \leq 1$.
                        Suppose not.
                        $1 \rless k$ or $1 = k$ by \cref{real_one_leq_naturalsplus}.
                        Thus $1 \leq k$.
                        $k\in\reals$ by \cref{real_naturals_subset_reals,subseteq}.
                        $1\in\reals$ by \cref{real_naturals_subset_reals,real_one_in_naturals,subseteq}.
                        $k + 1\in\reals$ by \cref{real_add_closed}.
                        $1 + 1\in\reals$ by \cref{real_add_closed}.
                        $1 + 1 \leq k + 1$ by \cref{real_leq_add}.
                        $0\in\reals$ by \cref{real_add_ident}.
                        $0 < 1$ by \cref{real_zero_lt_one}.
                        $0 + 1 < 1 + 1$ by \cref{real_lt_add}.
                        $0 + 1 = 1$ by \cref{real_add_ident}.
                        Thus $1 < 1 + 1$ by \cref{real_lt_subst_left}.
                        Thus $1 < k + 1$ by \cref{real_lt_trans}.
                        Thus $1 < 1$ by \cref{real_lt_trans,real_lt_subst_right}.
                        Contradiction by \cref{real_lt_irreflexive}.
                    \end{subproof}
                    Thus $\sumFiniteSequence{a}{s}$ by \cref{sum_finite_sequence_pred}.
                \end{subproof}
                Thus $\sumFiniteSequence{a}{s}$ by \cref{sum_finite_sequence_pred}.
            \end{subproof}
            Thus $1\in A$.
        \end{subproof}
        Thus for all $n\in A$ we have $n + 1\in A$ by \cref{sum_finite_sequence_exists_succ}.
        Thus for all $n\in\naturalsPlus$ we have $n\in A$ by \cref{real_naturals_induction}.
    Thus $m\in A$ by \cref{subseteq}.
    Thus there exists $s\in\reals$ such that $\sumFiniteSequence{f}{s}$.
\end{proof}

% from §20 helper definition 1: coordinate sum sequence
\begin{definition}\label{coord_sum_sequence}
    $\coordSumSequence{n}{a}{b} = \{w \mid \exists i\in\initialSegment{n}.
        w = (i, a(i) + b(i))\}$.
\end{definition}

% from §20 helper definition 2: coordinate product sequence
\begin{definition}\label{coord_product_sequence}
    $\coordProductSequence{n}{a}{b} = \{w \mid \exists i\in\initialSegment{n}.
        w = (i, a(i) \rmul b(i))\}$.
\end{definition}

% from §20 helper definition 3: coordinate square sequence
\begin{definition}\label{coord_square_sequence_mul}
    $\coordSquareSequenceMul{n}{a} = \{w \mid \exists i\in\initialSegment{n}.
        w = (i, a(i) \rmul a(i))\}$.
\end{definition}

% from §20 helper definition 4: coordinate difference sequence
\begin{definition}\label{coord_diff_sequence}
    $\coordDiffSequence{n}{x}{y} = \{w \mid \exists i\in\initialSegment{n}.
        w = (i, x(i) - y(i))\}$.
\end{definition}

% from §20 helper definition 5: coordinate scalar sequence
\begin{definition}\label{coord_scalar_sequence}
    $\coordScalarSequence{n}{c}{a} = \{w \mid \exists i\in\initialSegment{n}.
        w = (i, c \rmul a(i))\}$.
\end{definition}

% from §20 helper lemma 29: coordinate sum sequence is a function
\begin{lemma}\label{coord_sum_sequence_function}
    Suppose $n\in\naturalsPlus$.
    Let $J = \initialSegment{n}$.
    Suppose $a\in\funs{J}{\reals}$.
    Suppose $b\in\funs{J}{\reals}$.
    Then $\coordSumSequence{n}{a}{b}\in\funs{J}{\reals}$.
\end{lemma}
\begin{proof}
    Let $c = \coordSumSequence{n}{a}{b}$.
    Show that $c$ is a function.
    \begin{subproof}
        Show that $c$ is a relation.
        \begin{subproof}
            Show that for all $w$ such that $w\in c$ there exist $i,t$ such that $w = (i,t)$.
            \begin{subproof}
                Fix $w$.
                Assume $w\in c$.
                Take $i\in\initialSegment{n}$ such that $w = (i, a(i) + b(i))$ by \cref{coord_sum_sequence}.
                Thus there exist $i_0,t$ such that $w = (i_0,t)$.
            \end{subproof}
            Thus $c$ is a relation by \cref{relation}.
        \end{subproof}
        Show that for all $i,t_1,t_2$ such that $(i,t_1)\in c$ and $(i,t_2)\in c$ we have $t_1 = t_2$.
        \begin{subproof}
            Fix $i$.
            Fix $t_1$.
            Fix $t_2$.
            Assume $(i,t_1)\in c$ and $(i,t_2)\in c$.
            Take $i_1\in\initialSegment{n}$ such that $(i,t_1) = (i_1, a(i_1) + b(i_1))$ by \cref{coord_sum_sequence}.
            Take $i_2\in\initialSegment{n}$ such that $(i,t_2) = (i_2, a(i_2) + b(i_2))$ by \cref{coord_sum_sequence}.
            Thus $i = i_1$ and $t_1 = a(i_1) + b(i_1)$ by \cref{pair_eq_iff}.
            Thus $i = i_2$ and $t_2 = a(i_2) + b(i_2)$ by \cref{pair_eq_iff}.
            Thus $t_1 = a(i) + b(i)$.
            Thus $t_2 = a(i) + b(i)$.
            Thus $t_1 = t_2$.
        \end{subproof}
        Thus $c$ is right-unique by \cref{rightunique}.
        Thus $c$ is a function.
    \end{subproof}
    Show that $\dom{c} = J$.
    \begin{subproof}
        Show that for all $i$ such that $i\in J$ we have $i\in\dom{c}$.
        \begin{subproof}
            Fix $i$.
            Assume $i\in J$.
            Thus $(i, a(i) + b(i))\in c$ by \cref{coord_sum_sequence}.
            Thus $i\in\dom{c}$ by \cref{dom_intro}.
        \end{subproof}
        Thus $J\subseteq\dom{c}$ by \cref{subseteq}.
        Show that for all $i$ such that $i\in\dom{c}$ we have $i\in J$.
        \begin{subproof}
            Fix $i$.
            Assume $i\in\dom{c}$.
            Take $w$ such that $(i,w)\in c$ by \cref{dom_iff}.
            Take $i_0\in\initialSegment{n}$ such that $(i,w) = (i_0, a(i_0) + b(i_0))$ by \cref{coord_sum_sequence}.
            Thus $i = i_0$ by \cref{pair_eq_iff}.
            Thus $i\in J$.
        \end{subproof}
        Thus $\dom{c}\subseteq J$ by \cref{subseteq}.
        Thus $\dom{c} = J$ by \cref{subseteq_antisymmetric}.
    \end{subproof}
    Show that for all $i\in\dom{c}$ we have $c(i)\in\reals$.
    \begin{subproof}
        Fix $i$.
        Assume $i\in\dom{c}$.
        Thus $i\in J$ by \cref{subseteq}.
        $a\in\funs{J}{\reals}$ and $b\in\funs{J}{\reals}$.
        Thus $a(i)\in\reals$ and $b(i)\in\reals$ by \cref{funs_type_apply}.
        Thus $a(i) + b(i)\in\reals$ by \cref{real_add_closed}.
        $(i, a(i) + b(i))\in c$ by \cref{coord_sum_sequence}.
        Thus $c(i) = a(i) + b(i)$ by \cref{function_apply_intro}.
        Thus $c(i)\in\reals$.
    \end{subproof}
    Thus $c\in\funs{J}{\reals}$ by \cref{funs_intro}.
\end{proof}

% from §20 helper lemma 30: coordinate product sequence is a function
\begin{lemma}\label{coord_product_sequence_function}
    Suppose $n\in\naturalsPlus$.
    Let $J = \initialSegment{n}$.
    Suppose $a\in\funs{J}{\reals}$.
    Suppose $b\in\funs{J}{\reals}$.
    Then $\coordProductSequence{n}{a}{b}\in\funs{J}{\reals}$.
\end{lemma}
\begin{proof}
    Let $c = \coordProductSequence{n}{a}{b}$.
    Show that $c$ is a function.
    \begin{subproof}
        Show that $c$ is a relation.
        \begin{subproof}
            Show that for all $w$ such that $w\in c$ there exist $i,t$ such that $w = (i,t)$.
            \begin{subproof}
                Fix $w$.
                Assume $w\in c$.
                Take $i\in\initialSegment{n}$ such that $w = (i, a(i) \rmul b(i))$ by \cref{coord_product_sequence}.
                Thus there exist $i_0,t$ such that $w = (i_0,t)$.
            \end{subproof}
            Thus $c$ is a relation by \cref{relation}.
        \end{subproof}
        Show that for all $i,t_1,t_2$ such that $(i,t_1)\in c$ and $(i,t_2)\in c$ we have $t_1 = t_2$.
        \begin{subproof}
            Fix $i$.
            Fix $t_1$.
            Fix $t_2$.
            Assume $(i,t_1)\in c$ and $(i,t_2)\in c$.
            Take $i_1\in\initialSegment{n}$ such that $(i,t_1) = (i_1, a(i_1) \rmul b(i_1))$ by \cref{coord_product_sequence}.
            Take $i_2\in\initialSegment{n}$ such that $(i,t_2) = (i_2, a(i_2) \rmul b(i_2))$ by \cref{coord_product_sequence}.
            Thus $i = i_1$ and $t_1 = a(i_1) \rmul b(i_1)$ by \cref{pair_eq_iff}.
            Thus $i = i_2$ and $t_2 = a(i_2) \rmul b(i_2)$ by \cref{pair_eq_iff}.
            Thus $t_1 = a(i) \rmul b(i)$.
            Thus $t_2 = a(i) \rmul b(i)$.
            Thus $t_1 = t_2$.
        \end{subproof}
        Thus $c$ is right-unique by \cref{rightunique}.
        Thus $c$ is a function.
    \end{subproof}
    Show that $\dom{c} = J$.
    \begin{subproof}
        Show that for all $i$ such that $i\in J$ we have $i\in\dom{c}$.
        \begin{subproof}
            Fix $i$.
            Assume $i\in J$.
            Thus $(i, a(i) \rmul b(i))\in c$ by \cref{coord_product_sequence}.
            Thus $i\in\dom{c}$ by \cref{dom_intro}.
        \end{subproof}
        Thus $J\subseteq\dom{c}$ by \cref{subseteq}.
        Show that for all $i$ such that $i\in\dom{c}$ we have $i\in J$.
        \begin{subproof}
            Fix $i$.
            Assume $i\in\dom{c}$.
            Take $w$ such that $(i,w)\in c$ by \cref{dom_iff}.
            Take $i_0\in\initialSegment{n}$ such that $(i,w) = (i_0, a(i_0) \rmul b(i_0))$ by \cref{coord_product_sequence}.
            Thus $i = i_0$ by \cref{pair_eq_iff}.
            Thus $i\in J$.
        \end{subproof}
        Thus $\dom{c}\subseteq J$ by \cref{subseteq}.
        Thus $\dom{c} = J$ by \cref{subseteq_antisymmetric}.
    \end{subproof}
    Show that for all $i\in\dom{c}$ we have $c(i)\in\reals$.
    \begin{subproof}
        Fix $i$.
        Assume $i\in\dom{c}$.
        Thus $i\in J$ by \cref{subseteq}.
        $a\in\funs{J}{\reals}$ and $b\in\funs{J}{\reals}$.
        Thus $a(i)\in\reals$ and $b(i)\in\reals$ by \cref{funs_type_apply}.
        Thus $a(i) \rmul b(i)\in\reals$ by \cref{real_mul_closed}.
        $(i, a(i) \rmul b(i))\in c$ by \cref{coord_product_sequence}.
        Thus $c(i) = a(i) \rmul b(i)$ by \cref{function_apply_intro}.
        Thus $c(i)\in\reals$.
    \end{subproof}
    Thus $c\in\funs{J}{\reals}$ by \cref{funs_intro}.
\end{proof}

% from §20 helper lemma 31: coordinate square sequence is a function
\begin{lemma}\label{coord_square_sequence_mul_function}
    Suppose $n\in\naturalsPlus$.
    Let $J = \initialSegment{n}$.
    Suppose $a\in\funs{J}{\reals}$.
    Then $\coordSquareSequenceMul{n}{a}\in\funs{J}{\reals}$.
\end{lemma}
\begin{proof}
    Let $c = \coordSquareSequenceMul{n}{a}$.
    Show that $c$ is a function.
    \begin{subproof}
        Show that $c$ is a relation.
        \begin{subproof}
            Show that for all $w$ such that $w\in c$ there exist $i,t$ such that $w = (i,t)$.
            \begin{subproof}
                Fix $w$.
                Assume $w\in c$.
                Take $i\in\initialSegment{n}$ such that $w = (i, a(i) \rmul a(i))$ by \cref{coord_square_sequence_mul}.
                Thus there exist $i_0,t$ such that $w = (i_0,t)$.
            \end{subproof}
            Thus $c$ is a relation by \cref{relation}.
        \end{subproof}
        Show that for all $i,t_1,t_2$ such that $(i,t_1)\in c$ and $(i,t_2)\in c$ we have $t_1 = t_2$.
        \begin{subproof}
            Fix $i$.
            Fix $t_1$.
            Fix $t_2$.
            Assume $(i,t_1)\in c$ and $(i,t_2)\in c$.
            Take $i_1\in\initialSegment{n}$ such that $(i,t_1) = (i_1, a(i_1) \rmul a(i_1))$ by \cref{coord_square_sequence_mul}.
            Take $i_2\in\initialSegment{n}$ such that $(i,t_2) = (i_2, a(i_2) \rmul a(i_2))$ by \cref{coord_square_sequence_mul}.
            Thus $i = i_1$ and $t_1 = a(i_1) \rmul a(i_1)$ by \cref{pair_eq_iff}.
            Thus $i = i_2$ and $t_2 = a(i_2) \rmul a(i_2)$ by \cref{pair_eq_iff}.
            Thus $t_1 = a(i) \rmul a(i)$.
            Thus $t_2 = a(i) \rmul a(i)$.
            Thus $t_1 = t_2$.
        \end{subproof}
        Thus $c$ is right-unique by \cref{rightunique}.
        Thus $c$ is a function.
    \end{subproof}
    Show that $\dom{c} = J$.
    \begin{subproof}
        Show that for all $i$ such that $i\in J$ we have $i\in\dom{c}$.
        \begin{subproof}
            Fix $i$.
            Assume $i\in J$.
            Thus $(i, a(i) \rmul a(i))\in c$ by \cref{coord_square_sequence_mul}.
            Thus $i\in\dom{c}$ by \cref{dom_intro}.
        \end{subproof}
        Thus $J\subseteq\dom{c}$ by \cref{subseteq}.
        Show that for all $i$ such that $i\in\dom{c}$ we have $i\in J$.
        \begin{subproof}
            Fix $i$.
            Assume $i\in\dom{c}$.
            Take $w$ such that $(i,w)\in c$ by \cref{dom_iff}.
            Take $i_0\in\initialSegment{n}$ such that $(i,w) = (i_0, a(i_0) \rmul a(i_0))$ by \cref{coord_square_sequence_mul}.
            Thus $i = i_0$ by \cref{pair_eq_iff}.
            Thus $i\in J$.
        \end{subproof}
        Thus $\dom{c}\subseteq J$ by \cref{subseteq}.
        Thus $\dom{c} = J$ by \cref{subseteq_antisymmetric}.
    \end{subproof}
    Show that for all $i\in\dom{c}$ we have $c(i)\in\reals$.
    \begin{subproof}
        Fix $i$.
        Assume $i\in\dom{c}$.
        Thus $i\in J$ by \cref{subseteq}.
        $a\in\funs{J}{\reals}$.
        Thus $a(i)\in\reals$ by \cref{funs_type_apply}.
        Thus $a(i) \rmul a(i)\in\reals$ by \cref{real_mul_closed}.
        $(i, a(i) \rmul a(i))\in c$ by \cref{coord_square_sequence_mul}.
        Thus $c(i) = a(i) \rmul a(i)$ by \cref{function_apply_intro}.
        Thus $c(i)\in\reals$.
    \end{subproof}
    Thus $c\in\funs{J}{\reals}$ by \cref{funs_intro}.
\end{proof}

% from §20 helper lemma 32: coordinate difference sequence is a function
\begin{lemma}\label{coord_diff_sequence_function}
    Suppose $n\in\naturalsPlus$.
    Let $J = \initialSegment{n}$.
    Suppose $x\in\realsPower{J}$.
    Suppose $y\in\realsPower{J}$.
    Then $\coordDiffSequence{n}{x}{y}\in\funs{J}{\reals}$.
\end{lemma}
\begin{proof}
    Let $c = \coordDiffSequence{n}{x}{y}$.
    Show that $c$ is a function.
    \begin{subproof}
        Show that $c$ is a relation.
        \begin{subproof}
            Show that for all $w$ such that $w\in c$ there exist $i,t$ such that $w = (i,t)$.
            \begin{subproof}
                Fix $w$.
                Assume $w\in c$.
                Take $i\in\initialSegment{n}$ such that $w = (i, x(i) - y(i))$ by \cref{coord_diff_sequence}.
                Thus there exist $i_0,t$ such that $w = (i_0,t)$.
            \end{subproof}
            Thus $c$ is a relation by \cref{relation}.
        \end{subproof}
        Show that for all $i,t_1,t_2$ such that $(i,t_1)\in c$ and $(i,t_2)\in c$ we have $t_1 = t_2$.
        \begin{subproof}
            Fix $i$.
            Fix $t_1$.
            Fix $t_2$.
            Assume $(i,t_1)\in c$ and $(i,t_2)\in c$.
            Take $i_1\in\initialSegment{n}$ such that $(i,t_1) = (i_1, x(i_1) - y(i_1))$ by \cref{coord_diff_sequence}.
            Take $i_2\in\initialSegment{n}$ such that $(i,t_2) = (i_2, x(i_2) - y(i_2))$ by \cref{coord_diff_sequence}.
            Thus $i = i_1$ and $t_1 = x(i_1) - y(i_1)$ by \cref{pair_eq_iff}.
            Thus $i = i_2$ and $t_2 = x(i_2) - y(i_2)$ by \cref{pair_eq_iff}.
            Thus $t_1 = x(i) - y(i)$.
            Thus $t_2 = x(i) - y(i)$.
            Thus $t_1 = t_2$.
        \end{subproof}
        Thus $c$ is right-unique by \cref{rightunique}.
        Thus $c$ is a function.
    \end{subproof}
    Show that $\dom{c} = J$.
    \begin{subproof}
        Show that for all $i$ such that $i\in J$ we have $i\in\dom{c}$.
        \begin{subproof}
            Fix $i$.
            Assume $i\in J$.
            Thus $(i, x(i) - y(i))\in c$ by \cref{coord_diff_sequence}.
            Thus $i\in\dom{c}$ by \cref{dom_intro}.
        \end{subproof}
        Thus $J\subseteq\dom{c}$ by \cref{subseteq}.
        Show that for all $i$ such that $i\in\dom{c}$ we have $i\in J$.
        \begin{subproof}
            Fix $i$.
            Assume $i\in\dom{c}$.
            Take $w$ such that $(i,w)\in c$ by \cref{dom_iff}.
            Take $i_0\in\initialSegment{n}$ such that $(i,w) = (i_0, x(i_0) - y(i_0))$ by \cref{coord_diff_sequence}.
            Thus $i = i_0$ by \cref{pair_eq_iff}.
            Thus $i\in J$.
        \end{subproof}
        Thus $\dom{c}\subseteq J$ by \cref{subseteq}.
        Thus $\dom{c} = J$ by \cref{subseteq_antisymmetric}.
    \end{subproof}
    Show that for all $i\in\dom{c}$ we have $c(i)\in\reals$.
    \begin{subproof}
        Fix $i$.
        Assume $i\in\dom{c}$.
        Thus $i\in J$ by \cref{subseteq}.
        $x\in\tuplePower{\reals}{J}$ and $y\in\tuplePower{\reals}{J}$ by \cref{reals_power}.
        Thus $x\in\funs{J}{\reals}$ and $y\in\funs{J}{\reals}$ by \cref{tuple_power}.
        Thus $x(i)\in\reals$ and $y(i)\in\reals$ by \cref{funs_type_apply}.
        Thus $\neg{y(i)}\in\reals$ by \cref{real_add_inverse}.
        Thus $x(i) + \neg{y(i)}\in\reals$ by \cref{real_add_closed}.
        Thus $x(i) - y(i)\in\reals$ by \cref{real_sub}.
        $(i, x(i) - y(i))\in c$ by \cref{coord_diff_sequence}.
        Thus $c(i) = x(i) - y(i)$ by \cref{function_apply_intro}.
        Thus $c(i)\in\reals$.
    \end{subproof}
    Thus $c\in\funs{J}{\reals}$ by \cref{funs_intro}.
\end{proof}

% from §20 helper lemma 33: coordinate scalar sequence is a function
\begin{lemma}\label{coord_scalar_sequence_function}
    Suppose $n\in\naturalsPlus$.
    Let $J = \initialSegment{n}$.
    Suppose $c\in\reals$.
    Suppose $a\in\funs{J}{\reals}$.
    Then $\coordScalarSequence{n}{c}{a}\in\funs{J}{\reals}$.
\end{lemma}
\begin{proof}
    Let $b = \coordScalarSequence{n}{c}{a}$.
    Show that $b$ is a function.
    \begin{subproof}
        Show that $b$ is a relation.
        \begin{subproof}
            Show that for all $w$ such that $w\in b$ there exist $i,t$ such that $w = (i,t)$.
            \begin{subproof}
                Fix $w$.
                Assume $w\in b$.
                Take $i\in\initialSegment{n}$ such that $w = (i, c \rmul a(i))$ by \cref{coord_scalar_sequence}.
                Thus there exist $i_0,t$ such that $w = (i_0,t)$.
            \end{subproof}
            Thus $b$ is a relation by \cref{relation}.
        \end{subproof}
        Show that for all $i,t_1,t_2$ such that $(i,t_1)\in b$ and $(i,t_2)\in b$ we have $t_1 = t_2$.
        \begin{subproof}
            Fix $i$.
            Fix $t_1$.
            Fix $t_2$.
            Assume $(i,t_1)\in b$ and $(i,t_2)\in b$.
            Take $i_1\in\initialSegment{n}$ such that $(i,t_1) = (i_1, c \rmul a(i_1))$ by \cref{coord_scalar_sequence}.
            Take $i_2\in\initialSegment{n}$ such that $(i,t_2) = (i_2, c \rmul a(i_2))$ by \cref{coord_scalar_sequence}.
            Thus $i = i_1$ and $t_1 = c \rmul a(i_1)$ by \cref{pair_eq_iff}.
            Thus $i = i_2$ and $t_2 = c \rmul a(i_2)$ by \cref{pair_eq_iff}.
            Thus $t_1 = c \rmul a(i)$.
            Thus $t_2 = c \rmul a(i)$.
            Thus $t_1 = t_2$.
        \end{subproof}
        Thus $b$ is right-unique by \cref{rightunique}.
        Thus $b$ is a function.
    \end{subproof}
    Show that $\dom{b} = J$.
    \begin{subproof}
        Show that for all $i$ such that $i\in J$ we have $i\in\dom{b}$.
        \begin{subproof}
            Fix $i$.
            Assume $i\in J$.
            Thus $(i, c \rmul a(i))\in b$ by \cref{coord_scalar_sequence}.
            Thus $i\in\dom{b}$ by \cref{dom_intro}.
        \end{subproof}
        Thus $J\subseteq\dom{b}$ by \cref{subseteq}.
        Show that for all $i$ such that $i\in\dom{b}$ we have $i\in J$.
        \begin{subproof}
            Fix $i$.
            Assume $i\in\dom{b}$.
            Take $w$ such that $(i,w)\in b$ by \cref{dom_iff}.
            Take $i_0\in\initialSegment{n}$ such that $(i,w) = (i_0, c \rmul a(i_0))$ by \cref{coord_scalar_sequence}.
            Thus $i = i_0$ by \cref{pair_eq_iff}.
            Thus $i\in J$.
        \end{subproof}
        Thus $\dom{b}\subseteq J$ by \cref{subseteq}.
        Thus $\dom{b} = J$ by \cref{subseteq_antisymmetric}.
    \end{subproof}
    Show that for all $i\in\dom{b}$ we have $b(i)\in\reals$.
    \begin{subproof}
        Fix $i$.
        Assume $i\in\dom{b}$.
        Thus $i\in J$ by \cref{subseteq}.
        $a\in\funs{J}{\reals}$.
        Thus $a(i)\in\reals$ by \cref{funs_type_apply}.
        Thus $c \rmul a(i)\in\reals$ by \cref{real_mul_closed}.
        $(i, c \rmul a(i))\in b$ by \cref{coord_scalar_sequence}.
        Thus $b(i) = c \rmul a(i)$ by \cref{function_apply_intro}.
        Thus $b(i)\in\reals$.
    \end{subproof}
    Thus $b\in\funs{J}{\reals}$ by \cref{funs_intro}.
\end{proof}

% from §20 helper lemma 34: strict inequality implies weak inequality
\begin{lemma}\label{real_lt_imp_leq}
    Suppose $a < b$.
    Then $a \leq b$.
\end{lemma}
\begin{proof}
    Thus $a \leq b$.
\end{proof}

% from §20 helper lemma 35: sum of a coordinate sum sequence
\begin{lemma}\label{sum_finite_sequence_add}
    Suppose $n\in\naturalsPlus$.
    Let $J = \initialSegment{n}$.
    Suppose $a\in\funs{J}{\reals}$.
    Suppose $b\in\funs{J}{\reals}$.
    Suppose $\sumFiniteSequence{a}{s}$.
    Suppose $\sumFiniteSequence{b}{t}$.
    Let $c = \coordSumSequence{n}{a}{b}$.
    Then $\sumFiniteSequence{c}{s + t}$.
\end{lemma}
\begin{proof}
    Take $n_1$ such that $n_1\in\naturalsPlus$ and $a\in\funs{\initialSegment{n_1}}{\reals}$ and
        there exists $u\in\funs{\initialSegment{n_1}}{\reals}$ such that
            $u(1) = a(1)$ and $s = u(n_1)$ and
            for all $k\in\naturalsPlus$ such that $k + 1 \leq n_1$ we have $u(k + 1) = u(k) + a(k + 1)$
        by \cref{sum_finite_sequence_pred}.
    Take $n_2$ such that $n_2\in\naturalsPlus$ and $b\in\funs{\initialSegment{n_2}}{\reals}$ and
        there exists $v\in\funs{\initialSegment{n_2}}{\reals}$ such that
            $v(1) = b(1)$ and $t = v(n_2)$ and
            for all $k\in\naturalsPlus$ such that $k + 1 \leq n_2$ we have $v(k + 1) = v(k) + b(k + 1)$
        by \cref{sum_finite_sequence_pred}.
    Take $u$ such that $u\in\funs{\initialSegment{n_1}}{\reals}$ and
        $u(1) = a(1)$ and $s = u(n_1)$ and
        for all $k\in\naturalsPlus$ such that $k + 1 \leq n_1$ we have $u(k + 1) = u(k) + a(k + 1)$.
    Take $v$ such that $v\in\funs{\initialSegment{n_2}}{\reals}$ and
        $v(1) = b(1)$ and $t = v(n_2)$ and
        for all $k\in\naturalsPlus$ such that $k + 1 \leq n_2$ we have $v(k + 1) = v(k) + b(k + 1)$.
    Thus $\dom{a} = \initialSegment{n_1}$ and $\dom{b} = \initialSegment{n_2}$ by \cref{funs_elim}.
    Thus $\initialSegment{n_1} = J$ and $\initialSegment{n_2} = J$ by \cref{subseteq_antisymmetric,subseteq,funs_elim}.
    $n_1\in\naturalsPlus$ and $n_2\in\naturalsPlus$.
    Thus $n_1\leq n_1$ and $n_2\leq n_2$.
    Thus $n_1\in\initialSegment{n_1}$ and $n_2\in\initialSegment{n_2}$.
    Thus $n_1\in J$ and $n_2\in J$ by \cref{subseteq}.
    Thus $n_1 \leq n$ and $n_2 \leq n$ by \cref{initial_segment_naturalsplus}.
    $n\in\naturalsPlus$.
    Thus $n\leq n$.
    Thus $n\in J$ by \cref{initial_segment_naturalsplus}.
    Thus $n\in\initialSegment{n_1}$ and $n\in\initialSegment{n_2}$ by \cref{subseteq}.
    Thus $n \leq n_1$ and $n \leq n_2$ by \cref{initial_segment_naturalsplus}.
    $n_1\in\reals$ and $n_2\in\reals$ and $n\in\reals$ by \cref{real_naturals_subset_reals,subseteq}.
    Thus $n_1 = n$ and $n_2 = n$ by \cref{real_leq_antisym}.
    Let $w = \coordSumSequence{n}{u}{v}$.
    $u\in\funs{J}{\reals}$ and $v\in\funs{J}{\reals}$ by \cref{subseteq_antisymmetric,funs_elim}.
    Thus $w\in\funs{J}{\reals}$ by \cref{coord_sum_sequence_function}.
    Show that $w(1) = c(1)$.
    \begin{subproof}
        $1\in\naturalsPlus$ by \cref{real_one_in_naturals}.
        $1 \leq n$ by \cref{real_one_leq_naturalsplus}.
        Thus $1\in J$ by \cref{initial_segment_naturalsplus}.
        Thus $(1, u(1) + v(1))\in w$ by \cref{coord_sum_sequence}.
        $w$ is a function by \cref{funs_elim}.
        Thus $w(1) = u(1) + v(1)$ by \cref{function_apply_intro}.
        Thus $w(1) = a(1) + b(1)$.
        $(1, a(1) + b(1))\in c$ by \cref{coord_sum_sequence}.
        $c\in\funs{J}{\reals}$ by \cref{coord_sum_sequence_function}.
        Thus $c$ is a function by \cref{funs_elim}.
        Thus $c(1) = a(1) + b(1)$ by \cref{function_apply_intro}.
        Thus $w(1) = c(1)$.
    \end{subproof}
    Show that for all $k\in\naturalsPlus$ such that $k + 1 \leq n$ we have $w(k + 1) = w(k) + c(k + 1)$.
    \begin{subproof}
        Fix $k$.
        Assume $k\in\naturalsPlus$.
        Assume $k + 1 \leq n$.
        Thus $k + 1\in\naturalsPlus$ by \cref{real_naturals_closed_succ}.
        Thus $k + 1\in J$ by \cref{initial_segment_naturalsplus}.
        Show that $k\in J$.
        \begin{subproof}
            $k\in\reals$ by \cref{real_naturals_subset_reals,subseteq}.
            $0\in\reals$ by \cref{real_add_ident}.
            $1\in\reals$ by \cref{real_naturals_subset_reals,real_one_in_naturals,subseteq}.
            $0 < 1$ by \cref{real_zero_lt_one}.
            $0 + k < 1 + k$ by \cref{real_lt_add}.
            $0 + k = k$ by \cref{real_add_ident}.
            $1 + k = k + 1$ by \cref{real_add_comm}.
            Thus $k < k + 1$ by \cref{real_lt_subst_left}.
            Thus $k \leq k + 1$ by \cref{real_lt_imp_leq}.
            $k + 1\in\reals$ by \cref{real_add_closed}.
            $n\in\naturalsPlus$ by assumption.
            Thus $n\in\reals$ by \cref{real_naturals_subset_reals,subseteq}.
            Thus $k \leq n$ by \cref{real_leq_trans}.
            Thus $k\in J$ by \cref{initial_segment_naturalsplus}.
        \end{subproof}
        Thus $k\in J$.
        Thus $(k + 1, u(k + 1) + v(k + 1))\in w$ by \cref{coord_sum_sequence}.
        Thus $(k, u(k) + v(k))\in w$ by \cref{coord_sum_sequence}.
        $w$ is a function by \cref{funs_elim}.
        Thus $w(k + 1) = u(k + 1) + v(k + 1)$ by \cref{function_apply_intro}.
        Thus $w(k) = u(k) + v(k)$ by \cref{function_apply_intro}.
        $a\in\funs{J}{\reals}$ and $b\in\funs{J}{\reals}$.
        Thus $a(k + 1)\in\reals$ and $b(k + 1)\in\reals$ by \cref{funs_type_apply}.
        Thus $u(k)\in\reals$ and $v(k)\in\reals$ by \cref{funs_type_apply}.
        Thus $u(k) + a(k + 1)\in\reals$ by \cref{real_add_closed}.
        Thus $v(k) + b(k + 1)\in\reals$ by \cref{real_add_closed}.
        Thus $u(k) + v(k)\in\reals$ by \cref{real_add_closed}.
        Thus $a(k + 1) + b(k + 1)\in\reals$ by \cref{real_add_closed}.
        Thus $v(k) + a(k + 1)\in\reals$ by \cref{real_add_closed}.
        Thus $u(k + 1) = u(k) + a(k + 1)$.
        Thus $v(k + 1) = v(k) + b(k + 1)$.
        Thus $u(k + 1) + v(k + 1) = (u(k) + a(k + 1)) + v(k + 1)$ by \cref{real_add_eq_subst}.
        Thus $u(k + 1) + v(k + 1) = (u(k) + a(k + 1)) + (v(k) + b(k + 1))$ by \cref{real_add_eq_subst}.
        Thus $(u(k) + a(k + 1)) + (v(k) + b(k + 1)) = ((u(k) + a(k + 1)) + v(k)) + b(k + 1)$ by \cref{real_add_assoc}.
        Thus $(u(k) + a(k + 1)) + v(k) = u(k) + (a(k + 1) + v(k))$ by \cref{real_add_assoc}.
        Thus $a(k + 1) + v(k) = v(k) + a(k + 1)$ by \cref{real_add_comm}.
        Thus $u(k) + (a(k + 1) + v(k)) = u(k) + (v(k) + a(k + 1))$ by \cref{real_add_eq_subst}.
        Thus $(u(k) + a(k + 1)) + v(k) = u(k) + (v(k) + a(k + 1))$ by \cref{real_add_eq_subst}.
        Thus $((u(k) + a(k + 1)) + v(k)) + b(k + 1) = (u(k) + (v(k) + a(k + 1))) + b(k + 1)$ by \cref{real_add_eq_subst}.
        Thus $(u(k) + (v(k) + a(k + 1))) + b(k + 1) = u(k) + ((v(k) + a(k + 1)) + b(k + 1))$ by \cref{real_add_assoc}.
        Thus $(v(k) + a(k + 1)) + b(k + 1) = v(k) + (a(k + 1) + b(k + 1))$ by \cref{real_add_assoc}.
        Thus $u(k) + ((v(k) + a(k + 1)) + b(k + 1)) = u(k) + (v(k) + (a(k + 1) + b(k + 1)))$ by \cref{real_add_eq_subst}.
        Thus $u(k) + (v(k) + (a(k + 1) + b(k + 1))) = (u(k) + v(k)) + (a(k + 1) + b(k + 1))$ by \cref{real_add_assoc}.
        Thus $u(k + 1) + v(k + 1) = (u(k) + v(k)) + (a(k + 1) + b(k + 1))$ by \cref{real_add_eq_subst}.
        Thus $w(k + 1) = (u(k) + v(k)) + (a(k + 1) + b(k + 1))$ by \cref{real_add_eq_subst}.
        Thus $w(k + 1) = w(k) + (a(k + 1) + b(k + 1))$ by \cref{real_add_eq_subst}.
        $(k + 1, a(k + 1) + b(k + 1))\in c$ by \cref{coord_sum_sequence}.
        $c\in\funs{J}{\reals}$ by \cref{coord_sum_sequence_function}.
        Thus $c$ is a function by \cref{funs_elim}.
        Thus $c(k + 1) = a(k + 1) + b(k + 1)$ by \cref{function_apply_intro}.
        Thus $w(k + 1) = w(k) + c(k + 1)$ by \cref{real_add_eq_subst}.
    \end{subproof}
    Show that $s + t = w(n)$.
    \begin{subproof}
        $n\in J$ by \cref{initial_segment_naturalsplus}.
        Thus $(n, u(n) + v(n))\in w$ by \cref{coord_sum_sequence}.
        $w$ is a function by \cref{funs_elim}.
        Thus $w(n) = u(n) + v(n)$ by \cref{function_apply_intro}.
        Thus $w(n) = s + t$.
    \end{subproof}
    Show that $\sumFiniteSequence{c}{s + t}$.
    \begin{subproof}
        $c\in\funs{J}{\reals}$ by \cref{coord_sum_sequence_function}.
        Thus $c\in\funs{\initialSegment{n}}{\reals}$.
        $w\in\funs{J}{\reals}$.
        Thus $w\in\funs{\initialSegment{n}}{\reals}$.
        Thus $\sumFiniteSequence{c}{s + t}$ by \cref{sum_finite_sequence_pred}.
    \end{subproof}
\end{proof}

% from §20 helper lemma 36: sum of a coordinate scalar sequence
\begin{lemma}\label{sum_finite_sequence_scalar}
    Suppose $n\in\naturalsPlus$.
    Let $J = \initialSegment{n}$.
    Suppose $c\in\reals$.
    Suppose $a\in\funs{J}{\reals}$.
    Suppose $\sumFiniteSequence{a}{s}$.
    Let $b = \coordScalarSequence{n}{c}{a}$.
    Then $\sumFiniteSequence{b}{c \rmul s}$.
\end{lemma}
\begin{proof}
    Take $n_1$ such that $n_1\in\naturalsPlus$ and $a\in\funs{\initialSegment{n_1}}{\reals}$ and
        there exists $u\in\funs{\initialSegment{n_1}}{\reals}$ such that
            $u(1) = a(1)$ and $s = u(n_1)$ and
            for all $k\in\naturalsPlus$ such that $k + 1 \leq n_1$ we have $u(k + 1) = u(k) + a(k + 1)$
        by \cref{sum_finite_sequence_pred}.
    Take $u$ such that $u\in\funs{\initialSegment{n_1}}{\reals}$ and
        $u(1) = a(1)$ and $s = u(n_1)$ and
        for all $k\in\naturalsPlus$ such that $k + 1 \leq n_1$ we have $u(k + 1) = u(k) + a(k + 1)$.
    Thus $\dom{a} = \initialSegment{n_1}$ by \cref{funs_elim}.
    Thus $\initialSegment{n_1} = J$ by \cref{subseteq_antisymmetric,subseteq,funs_elim}.
    $n_1\in\naturalsPlus$.
    Thus $n_1\leq n_1$.
    Thus $n_1\in\initialSegment{n_1}$.
    Thus $n_1\in J$ by \cref{subseteq}.
    Thus $n_1 \leq n$ by \cref{initial_segment_naturalsplus}.
    $n\in\naturalsPlus$.
    Thus $n\leq n$.
    Thus $n\in J$ by \cref{initial_segment_naturalsplus}.
    Thus $n\in\initialSegment{n_1}$ by \cref{subseteq}.
    Thus $n \leq n_1$ by \cref{initial_segment_naturalsplus}.
    $n_1\in\reals$ and $n\in\reals$ by \cref{real_naturals_subset_reals,subseteq}.
    Thus $n_1 = n$ by \cref{real_leq_antisym}.
    Let $w = \coordScalarSequence{n}{c}{u}$.
    $u\in\funs{J}{\reals}$ by \cref{subseteq_antisymmetric,funs_elim}.
    Thus $w\in\funs{J}{\reals}$ by \cref{coord_scalar_sequence_function}.
    Show that $w(1) = b(1)$.
    \begin{subproof}
        $1\in\naturalsPlus$ by \cref{real_one_in_naturals}.
        $1 \leq n$ by \cref{real_one_leq_naturalsplus}.
        Thus $1\in J$ by \cref{initial_segment_naturalsplus}.
        Thus $(1, c \rmul u(1))\in w$ by \cref{coord_scalar_sequence}.
        $w$ is a function by \cref{funs_elim}.
        Thus $w(1) = c \rmul u(1)$ by \cref{function_apply_intro}.
        Thus $w(1) = c \rmul a(1)$.
        $(1, c \rmul a(1))\in b$ by \cref{coord_scalar_sequence}.
        $b\in\funs{J}{\reals}$ by \cref{coord_scalar_sequence_function}.
        Thus $b$ is a function by \cref{funs_elim}.
        Thus $b(1) = c \rmul a(1)$ by \cref{function_apply_intro}.
        Thus $w(1) = b(1)$.
    \end{subproof}
    Show that for all $k\in\naturalsPlus$ such that $k + 1 \leq n$ we have $w(k + 1) = w(k) + b(k + 1)$.
    \begin{subproof}
        Fix $k$.
        Assume $k\in\naturalsPlus$.
        Assume $k + 1 \leq n$.
        Thus $k + 1\in\naturalsPlus$ by \cref{real_naturals_closed_succ}.
        Thus $k + 1\in J$ by \cref{initial_segment_naturalsplus}.
        Show that $k\in J$.
        \begin{subproof}
            $k\in\reals$ by \cref{real_naturals_subset_reals,subseteq}.
            $0\in\reals$ by \cref{real_add_ident}.
            $1\in\reals$ by \cref{real_naturals_subset_reals,real_one_in_naturals,subseteq}.
            $0 < 1$ by \cref{real_zero_lt_one}.
            $0 + k < 1 + k$ by \cref{real_lt_add}.
            $0 + k = k$ by \cref{real_add_ident}.
            $1 + k = k + 1$ by \cref{real_add_comm}.
            Thus $k < k + 1$ by \cref{real_lt_subst_left}.
            Thus $k \leq k + 1$ by \cref{real_lt_imp_leq}.
            $k + 1\in\reals$ by \cref{real_add_closed}.
            $n\in\reals$ by \cref{real_naturals_subset_reals,subseteq}.
            Thus $k \leq n$ by \cref{real_leq_trans}.
            Thus $k\in J$ by \cref{initial_segment_naturalsplus}.
        \end{subproof}
        Thus $k\in J$.
        Thus $(k + 1, c \rmul u(k + 1))\in w$ by \cref{coord_scalar_sequence}.
        Thus $(k, c \rmul u(k))\in w$ by \cref{coord_scalar_sequence}.
        $w$ is a function by \cref{funs_elim}.
        Thus $w(k + 1) = c \rmul u(k + 1)$ by \cref{function_apply_intro}.
        Thus $w(k) = c \rmul u(k)$ by \cref{function_apply_intro}.
        $a\in\funs{J}{\reals}$.
        Thus $a(k + 1)\in\reals$ by \cref{funs_type_apply}.
        Thus $u(k)\in\reals$ by \cref{funs_type_apply}.
        Thus $u(k + 1) = u(k) + a(k + 1)$.
        Thus $c \rmul u(k + 1) = c \rmul (u(k) + a(k + 1))$.
        Thus $c \rmul u(k + 1) = (c \rmul u(k)) + (c \rmul a(k + 1))$ by \cref{real_distrib}.
        Thus $w(k + 1) = (c \rmul u(k)) + (c \rmul a(k + 1))$ by \cref{real_add_eq_subst}.
        Thus $w(k + 1) = w(k) + (c \rmul a(k + 1))$ by \cref{real_add_eq_subst}.
        $(k + 1, c \rmul a(k + 1))\in b$ by \cref{coord_scalar_sequence}.
        $b\in\funs{J}{\reals}$ by \cref{coord_scalar_sequence_function}.
        Thus $b$ is a function by \cref{funs_elim}.
        Thus $b(k + 1) = c \rmul a(k + 1)$ by \cref{function_apply_intro}.
        Thus $w(k + 1) = w(k) + b(k + 1)$ by \cref{real_add_eq_subst}.
    \end{subproof}
    Show that $c \rmul s = w(n)$.
    \begin{subproof}
        $n\in J$ by \cref{initial_segment_naturalsplus}.
        Thus $(n, c \rmul u(n))\in w$ by \cref{coord_scalar_sequence}.
        $w$ is a function by \cref{funs_elim}.
        Thus $w(n) = c \rmul u(n)$ by \cref{function_apply_intro}.
        Thus $w(n) = c \rmul s$.
    \end{subproof}
    Show that $\sumFiniteSequence{b}{c \rmul s}$.
    \begin{subproof}
        $b\in\funs{J}{\reals}$ by \cref{coord_scalar_sequence_function}.
        Thus $b\in\funs{\initialSegment{n}}{\reals}$.
        $w\in\funs{J}{\reals}$.
        Thus $w\in\funs{\initialSegment{n}}{\reals}$.
        Thus $\sumFiniteSequence{b}{c \rmul s}$ by \cref{sum_finite_sequence_pred}.
    \end{subproof}
\end{proof}

% from §20 helper lemma 37: split sum at the last term
\begin{lemma}\label{sum_finite_sequence_split_last}
    Suppose $n\in\naturalsPlus$.
    Let $J = \initialSegment{n}$.
    Let $J_1 = \initialSegment{n + 1}$.
    Suppose $a\in\funs{J_1}{\reals}$.
    Suppose $\sumFiniteSequence{a}{s}$.
    Let $a_0 = \restrl{a}{J}$.
    Then $a_0\in\funs{J}{\reals}$ and
        there exists $s_0\in\reals$ such that $\sumFiniteSequence{a_0}{s_0}$ and $s = s_0 + a(n + 1)$.
\end{lemma}
\begin{proof}
    $a$ is a function by \cref{funs_elim}.
    Thus $a_0$ is a function by \cref{restrl_of_function_is_function}.
    $\dom{a} = J_1$ by \cref{funs_elim}.
    Thus $\dom{a_0} = J_1\inter J$ by \cref{dom_of_restrl_eq_inter}.
    Show that $J\subseteq J_1$.
    \begin{subproof}
        Show that for all $i$ such that $i\in J$ we have $i\in J_1$.
        \begin{subproof}
            Fix $i$.
            Assume $i\in J$.
            Thus $i\in\naturalsPlus$ and $i \leq n$ by \cref{initial_segment_naturalsplus}.
            $n\in\reals$ and $i\in\reals$ by \cref{real_naturals_subset_reals,subseteq}.
            $1\in\reals$ by \cref{real_naturals_subset_reals,real_one_in_naturals,subseteq}.
            $0\in\reals$ by \cref{real_add_ident}.
            $0 < 1$ by \cref{real_zero_lt_one}.
            $0 + n < 1 + n$ by \cref{real_lt_add}.
            $0 + n = n$ by \cref{real_add_ident}.
            $1 + n = n + 1$ by \cref{real_add_comm}.
            $n + 1\in\reals$ by \cref{real_add_closed}.
            Thus $n < n + 1$ by \cref{real_lt_subst_left}.
            Thus $n \leq n + 1$ by \cref{real_lt_imp_leq}.
            Thus $i \leq n + 1$ by \cref{real_leq_trans}.
            Thus $i\in J_1$ by \cref{initial_segment_naturalsplus}.
        \end{subproof}
        Thus $J\subseteq J_1$ by \cref{subseteq}.
    \end{subproof}
    Thus $J_1\inter J = J$ by \cref{inter_absorb_supseteq_left,subseteq}.
    Thus $\dom{a_0} = J$.
    Show that for all $i\in\dom{a_0}$ we have $a_0(i)\in\reals$.
    \begin{subproof}
        Fix $i$.
        Assume $i\in\dom{a_0}$.
        Thus $i\in J$.
        $i\in J_1$ by \cref{subseteq}.
        Thus $i\in\dom{a}$ by \cref{funs_elim}.
        Thus $a(i)\in\reals$ by \cref{funs_type_apply}.
        Thus $a_0(i) = a(i)$ by \cref{restrl_of_function_apply,subseteq}.
        Thus $a_0(i)\in\reals$.
    \end{subproof}
    Thus $a_0\in\funs{J}{\reals}$ by \cref{funs_intro}.
    Take $n_0$ such that $n_0\in\naturalsPlus$ and $a\in\funs{\initialSegment{n_0}}{\reals}$ and
        there exists $t\in\funs{\initialSegment{n_0}}{\reals}$ such that
            $t(1) = a(1)$ and $s = t(n_0)$ and
            for all $k\in\naturalsPlus$ such that $k + 1 \leq n_0$ we have $t(k + 1) = t(k) + a(k + 1)$
        by \cref{sum_finite_sequence_pred}.
    Take $t$ such that $t\in\funs{\initialSegment{n_0}}{\reals}$ and
        $t(1) = a(1)$ and $s = t(n_0)$ and
        for all $k\in\naturalsPlus$ such that $k + 1 \leq n_0$ we have $t(k + 1) = t(k) + a(k + 1)$.
    Thus $\dom{a} = \initialSegment{n_0}$ by \cref{funs_elim}.
    Thus $\initialSegment{n_0} = J_1$ by \cref{subseteq_antisymmetric,subseteq,funs_elim}.
    $n_0\in\naturalsPlus$.
    Thus $n_0\leq n_0$.
    Thus $n_0\in\initialSegment{n_0}$.
    Thus $n_0\in J_1$ by \cref{subseteq}.
    Thus $n_0 \leq n + 1$ by \cref{initial_segment_naturalsplus}.
    $n + 1\in\naturalsPlus$ by \cref{real_naturals_closed_succ}.
    Thus $n + 1\leq n + 1$.
    Thus $n + 1\in J_1$ by \cref{initial_segment_naturalsplus}.
    Thus $n + 1\in\initialSegment{n_0}$ by \cref{subseteq}.
    Thus $n + 1 \leq n_0$ by \cref{initial_segment_naturalsplus}.
    $n_0\in\reals$ and $n + 1\in\reals$ by \cref{real_naturals_subset_reals,subseteq}.
    Thus $n_0 = n + 1$ by \cref{real_leq_antisym}.
    Let $t_0 = \restrl{t}{J}$.
    $t$ is a function by \cref{funs_elim}.
    Thus $t_0$ is a function by \cref{restrl_of_function_is_function}.
    $\dom{t} = J_1$ by \cref{funs_elim}.
    Thus $\dom{t_0} = J_1\inter J$ by \cref{dom_of_restrl_eq_inter}.
    Thus $\dom{t_0} = J$ by \cref{inter_absorb_supseteq_left,subseteq}.
    Show that for all $i\in\dom{t_0}$ we have $t_0(i)\in\reals$.
    \begin{subproof}
        Fix $i$.
        Assume $i\in\dom{t_0}$.
        Thus $i\in J$ by \cref{subseteq}.
        Thus $i\in J_1$ by \cref{subseteq}.
        Thus $i\in\dom{t}$ by \cref{funs_elim}.
        Thus $t(i)\in\reals$ by \cref{funs_type_apply}.
        Thus $t_0(i) = t(i)$ by \cref{restrl_of_function_apply,subseteq}.
        Thus $t_0(i)\in\reals$.
    \end{subproof}
    Thus $t_0\in\funs{J}{\reals}$ by \cref{funs_intro}.
    Show that $t_0(1) = a_0(1)$.
    \begin{subproof}
        $1\in\naturalsPlus$ by \cref{real_one_in_naturals}.
        $1 \leq n$ by \cref{real_one_leq_naturalsplus}.
        Thus $1\in J$ by \cref{initial_segment_naturalsplus}.
        Thus $t_0(1) = t(1)$ by \cref{restrl_of_function_apply,subseteq}.
        Thus $a_0(1) = a(1)$ by \cref{restrl_of_function_apply,subseteq}.
        Thus $t_0(1) = a_0(1)$.
    \end{subproof}
    Show that for all $k\in\naturalsPlus$ such that $k + 1 \leq n$ we have
        $t_0(k + 1) = t_0(k) + a_0(k + 1)$.
    \begin{subproof}
        Fix $k$.
        Assume $k\in\naturalsPlus$.
        Assume $k + 1 \leq n$.
        Thus $k + 1\in\naturalsPlus$ by \cref{real_naturals_closed_succ}.
        Thus $k + 1\in J$ by \cref{initial_segment_naturalsplus}.
        Show that $k\in J$.
        \begin{subproof}
            $k\in\reals$ by \cref{real_naturals_subset_reals,subseteq}.
            $0\in\reals$ by \cref{real_add_ident}.
            $1\in\reals$ by \cref{real_naturals_subset_reals,real_one_in_naturals,subseteq}.
            $0 < 1$ by \cref{real_zero_lt_one}.
            $0 + k < 1 + k$ by \cref{real_lt_add}.
            $0 + k = k$ by \cref{real_add_ident}.
            $1 + k = k + 1$ by \cref{real_add_comm}.
            Thus $k < k + 1$ by \cref{real_lt_subst_left}.
            Thus $k \leq k + 1$ by \cref{real_lt_imp_leq}.
            $k + 1\in\reals$ by \cref{real_add_closed}.
            $n\in\naturalsPlus$ by assumption.
            Thus $n\in\reals$ by \cref{real_naturals_subset_reals,subseteq}.
            Thus $k \leq n$ by \cref{real_leq_trans}.
            Thus $k\in J$ by \cref{initial_segment_naturalsplus}.
        \end{subproof}
        Thus $k\in J$.
        Thus $t_0(k + 1) = t(k + 1)$ by \cref{restrl_of_function_apply,subseteq}.
        Thus $t_0(k) = t(k)$ by \cref{restrl_of_function_apply,subseteq}.
        Thus $a_0(k + 1) = a(k + 1)$ by \cref{restrl_of_function_apply,subseteq}.
        Show that $n \leq n + 1$.
        \begin{subproof}
            $n\in\naturalsPlus$ by assumption.
            Thus $n\in\reals$ by \cref{real_naturals_subset_reals,subseteq}.
            $0\in\reals$ by \cref{real_add_ident}.
            $1\in\reals$ by \cref{real_naturals_subset_reals,real_one_in_naturals,subseteq}.
            $0 < 1$ by \cref{real_zero_lt_one}.
            $0 + n < 1 + n$ by \cref{real_lt_add}.
            $0 + n = n$ by \cref{real_add_ident}.
            $1 + n = n + 1$ by \cref{real_add_comm}.
            Thus $n < n + 1$ by \cref{real_lt_subst_left}.
            Thus $n \leq n + 1$ by \cref{real_lt_imp_leq}.
        \end{subproof}
        $n\in\reals$ by \cref{real_naturals_subset_reals,subseteq}.
        $n + 1\in\reals$ by \cref{real_add_closed}.
        $k\in\reals$ by \cref{real_naturals_subset_reals,subseteq}.
        $1\in\reals$ by \cref{real_naturals_subset_reals,real_one_in_naturals,subseteq}.
        $k + 1\in\reals$ by \cref{real_add_closed}.
        Thus $k + 1 \leq n + 1$ by \cref{real_leq_trans}.
        Thus $t(k + 1) = t(k) + a(k + 1)$.
        Thus $t_0(k + 1) = t_0(k) + a_0(k + 1)$ by \cref{real_add_eq_subst}.
    \end{subproof}
    Let $s_0 = t_0(n)$.
    Show that $\sumFiniteSequence{a_0}{s_0}$.
    \begin{subproof}
        $n\in J$ by \cref{initial_segment_naturalsplus}.
        Thus $t_0(n)\in\reals$ by \cref{funs_type_apply}.
        Thus $\sumFiniteSequence{a_0}{s_0}$ by \cref{sum_finite_sequence_pred}.
    \end{subproof}
    Show that $s = s_0 + a(n + 1)$.
    \begin{subproof}
        $n + 1\in J_1$ by \cref{initial_segment_naturalsplus,real_naturals_closed_succ}.
        Thus $t(n + 1) = s$.
        Thus $t(n + 1) = t(n) + a(n + 1)$.
        Thus $t(n) + a(n + 1) = s$.
        $n\in J$ by \cref{initial_segment_naturalsplus}.
        Thus $t_0(n) = t(n)$ by \cref{restrl_of_function_apply,subseteq}.
        Thus $t_0(n) + a(n + 1) = t(n) + a(n + 1)$ by \cref{real_add_eq_subst}.
        Thus $t_0(n) + a(n + 1) = s$.
        Thus $s = s_0 + a(n + 1)$ by \cref{real_add_comm}.
    \end{subproof}
    Thus there exists $s_0\in\reals$ such that $\sumFiniteSequence{a_0}{s_0}$ and $s = s_0 + a(n + 1)$.
\end{proof}

% from §20 helper lemma 38: zero sum of a nonnegative finite sequence
\begin{lemma}\label{sum_finite_sequence_zero_terms}
    Suppose $n\in\naturalsPlus$.
    Let $J = \initialSegment{n}$.
    Suppose $a\in\funs{J}{\reals}$.
    Suppose $\sumFiniteSequence{a}{s}$.
    Suppose $s = 0$.
    Suppose for all $i\in J$ we have $0 \leq a(i)$.
    Then for all $i\in J$ we have $a(i) = 0$.
\end{lemma}
\begin{proof}
    Let $A = \{m\in\naturalsPlus \mid \text{for all $b\in\funs{\initialSegment{m}}{\reals}$ for all $t\in\reals$
        if $\sumFiniteSequence{b}{t}$ and $t = 0$ and for all $i\in\initialSegment{m}$ we have $0 \leq b(i)$ then
        for all $i\in\initialSegment{m}$ we have $b(i) = 0$}\}$.
    $A\subseteq\naturalsPlus$ by \cref{subseteq}.
    Show that $1\in A$.
    \begin{subproof}
        Show that for all $b\in\funs{\initialSegment{1}}{\reals}$ for all $t\in\reals$
            if $\sumFiniteSequence{b}{t}$ and $t = 0$ and for all $i\in\initialSegment{1}$ we have $0 \leq b(i)$ then
            for all $i\in\initialSegment{1}$ we have $b(i) = 0$.
        \begin{subproof}
            Fix $b\in\funs{\initialSegment{1}}{\reals}$.
            Fix $t\in\reals$.
            Assume $\sumFiniteSequence{b}{t}$.
            Assume $t = 0$.
            Assume for all $i\in\initialSegment{1}$ we have $0 \leq b(i)$.
            Take $n_0$ such that $n_0\in\naturalsPlus$ and $b\in\funs{\initialSegment{n_0}}{\reals}$ and
                there exists $u\in\funs{\initialSegment{n_0}}{\reals}$ such that
                    $u(1) = b(1)$ and $t = u(n_0)$ and
                    for all $k\in\naturalsPlus$ such that $k + 1 \leq n_0$ we have $u(k + 1) = u(k) + b(k + 1)$
                by \cref{sum_finite_sequence_pred}.
            Thus $\dom{b} = \initialSegment{n_0}$ by \cref{funs_elim}.
            Thus $\initialSegment{n_0} = \initialSegment{1}$ by \cref{subseteq_antisymmetric,subseteq,funs_elim}.
            $n_0\in\naturalsPlus$.
            Thus $n_0\leq n_0$.
            Thus $n_0\in\initialSegment{n_0}$.
            Thus $n_0\in\initialSegment{1}$ by \cref{subseteq}.
            Thus $n_0 \leq 1$ by \cref{initial_segment_naturalsplus}.
            $1\in\naturalsPlus$ by \cref{real_one_in_naturals}.
            Thus $1\leq 1$.
            Thus $1\in\initialSegment{1}$ by \cref{initial_segment_naturalsplus}.
            Thus $1\in\initialSegment{n_0}$ by \cref{subseteq}.
            Thus $1 \leq n_0$ by \cref{initial_segment_naturalsplus}.
            $n_0\in\reals$ and $1\in\reals$ by \cref{real_naturals_subset_reals,subseteq}.
            Thus $n_0 = 1$ by \cref{real_leq_antisym}.
            Take $u$ such that $u\in\funs{\initialSegment{n_0}}{\reals}$ and
                $u(1) = b(1)$ and $t = u(n_0)$ and
                for all $k\in\naturalsPlus$ such that $k + 1 \leq n_0$ we have $u(k + 1) = u(k) + b(k + 1)$.
            Thus $t = u(1)$.
            Thus $b(1) = 0$.
            Show that for all $i\in\initialSegment{1}$ we have $b(i) = 0$.
            \begin{subproof}
                Fix $i$.
                Assume $i\in\initialSegment{1}$.
                Thus $i\in\naturalsPlus$ and $i \leq 1$ by \cref{initial_segment_naturalsplus}.
                Thus $1 \leq i$ by \cref{real_one_leq_naturalsplus}.
                $i\in\reals$ and $1\in\reals$ by \cref{real_naturals_subset_reals,subseteq}.
                Thus $i = 1$ by \cref{real_leq_antisym}.
                Thus $b(i) = 0$.
            \end{subproof}
        \end{subproof}
        Thus $1\in A$.
    \end{subproof}
    Show that for all $k\in A$ we have $k + 1\in A$.
    \begin{subproof}
        Fix $k$.
        Assume $k\in A$.
        Show that $k + 1\in A$.
        \begin{subproof}
            Show that for all $b\in\funs{\initialSegment{k + 1}}{\reals}$ for all $t\in\reals$
                if $\sumFiniteSequence{b}{t}$ and $t = 0$ and for all $i\in\initialSegment{k + 1}$ we have $0 \leq b(i)$ then
                for all $i\in\initialSegment{k + 1}$ we have $b(i) = 0$.
            \begin{subproof}
                Fix $b\in\funs{\initialSegment{k + 1}}{\reals}$.
                Fix $t\in\reals$.
                Assume $\sumFiniteSequence{b}{t}$.
                Assume $t = 0$.
                Assume for all $i\in\initialSegment{k + 1}$ we have $0 \leq b(i)$.
                Let $J = \initialSegment{k}$.
                Let $J_1 = \initialSegment{k + 1}$.
                Let $b_0 = \restrl{b}{J}$.
                $k\in\naturalsPlus$ by \cref{subseteq}.
                Thus $b_0\in\funs{J}{\reals}$ and
                    there exists $s_0\in\reals$ such that $\sumFiniteSequence{b_0}{s_0}$ and $t = s_0 + b(k + 1)$
                    by \cref{sum_finite_sequence_split_last}.
                Take $s_0\in\reals$ such that $\sumFiniteSequence{b_0}{s_0}$ and $t = s_0 + b(k + 1)$.
                Show that $J\subseteq J_1$.
                \begin{subproof}
                    Show that for all $i$ such that $i\in J$ we have $i\in J_1$.
                    \begin{subproof}
                        Fix $i$.
                        Assume $i\in J$.
                        Thus $i\in\naturalsPlus$ and $i \leq k$ by \cref{initial_segment_naturalsplus}.
                        $k\in\naturalsPlus$ by \cref{subseteq}.
                        Thus $k\in\reals$ by \cref{real_naturals_subset_reals,subseteq}.
                        $0\in\reals$ by \cref{real_add_ident}.
                        $1\in\reals$ by \cref{real_naturals_subset_reals,real_one_in_naturals,subseteq}.
                        $0 < 1$ by \cref{real_zero_lt_one}.
                        $0 + k < 1 + k$ by \cref{real_lt_add}.
                        $0 + k = k$ by \cref{real_add_ident}.
                        $1 + k = k + 1$ by \cref{real_add_comm}.
                        Thus $k < k + 1$ by \cref{real_lt_subst_left}.
                        Thus $k \leq k + 1$ by \cref{real_lt_imp_leq}.
                        $k + 1\in\reals$ by \cref{real_add_closed}.
                        $i\in\reals$ by \cref{real_naturals_subset_reals,subseteq}.
                        Thus $i \leq k + 1$ by \cref{real_leq_trans}.
                        Thus $i\in J_1$ by \cref{initial_segment_naturalsplus}.
                    \end{subproof}
                    Thus $J\subseteq J_1$ by \cref{subseteq}.
                \end{subproof}
                Show that for all $i\in J$ we have $0 \leq b_0(i)$.
                \begin{subproof}
                    Fix $i$.
                    Assume $i\in J$.
                    Thus $i\in J_1$ by \cref{subseteq}.
                    Thus $\dom{b} = J_1$ by \cref{funs_elim}.
                    Thus $J\subseteq\dom{b}$ by \cref{subseteq}.
                    $b$ is a function by \cref{funs_elim}.
                    Thus $b_0(i) = b(i)$ by \cref{restrl_of_function_apply}.
                    Thus $0 \leq b_0(i)$.
                \end{subproof}
                Thus $\dom{b_0} = J$ by \cref{funs_elim}.
                Thus for all $i\in\dom{b_0}$ we have $0 \leq b_0(i)$ by \cref{subseteq}.
                Thus $0 \leq s_0$ by \cref{sum_finite_sequence_nonneg}.
                $k + 1\in\naturalsPlus$ by \cref{real_naturals_closed_succ}.
                Thus $k + 1\in J_1$ by \cref{initial_segment_naturalsplus}.
                Thus $0 \leq b(k + 1)$.
                $t = 0$.
                Thus $s_0 + b(k + 1) = 0$.
                Show that $s_0 = 0$.
                \begin{subproof}
                    $0 \leq b(k + 1)$.
                    $s_0\in\reals$ and $b(k + 1)\in\reals$ by \cref{funs_type_apply,restrl_of_function_apply,subseteq}.
                    $0\in\reals$ by \cref{real_add_ident}.
                    $0 + s_0 \leq b(k + 1) + s_0$ by \cref{real_leq_add}.
                    $0 + s_0 = s_0$ by \cref{real_add_ident}.
                    $b(k + 1) + s_0 = s_0 + b(k + 1)$ by \cref{real_add_comm}.
                    Thus $s_0 \leq s_0 + b(k + 1)$.
                    Thus $s_0 + b(k + 1) \leq 0$.
                    Thus $s_0 \leq 0$ by \cref{real_leq_trans}.
                    Thus $s_0 = 0$ by \cref{real_leq_antisym}.
                \end{subproof}
                Show that $b(k + 1) = 0$.
                \begin{subproof}
                    $0 \leq s_0$.
                    $b(k + 1)\in\reals$ and $s_0\in\reals$ by \cref{funs_type_apply,restrl_of_function_apply,subseteq}.
                    $0\in\reals$ by \cref{real_add_ident}.
                    $0 + b(k + 1) \leq s_0 + b(k + 1)$ by \cref{real_leq_add}.
                    $0 + b(k + 1) = b(k + 1)$ by \cref{real_add_ident}.
                    Thus $b(k + 1) \leq s_0 + b(k + 1)$.
                    Thus $s_0 + b(k + 1) \leq 0$.
                    Thus $b(k + 1) \leq 0$ by \cref{real_leq_trans}.
                    Thus $b(k + 1) = 0$ by \cref{real_leq_antisym}.
                \end{subproof}
                Show that for all $i\in J$ we have $b_0(i) = 0$.
                \begin{subproof}
                    $k\in A$.
                    Thus for all $b_1\in\funs{J}{\reals}$ for all $t_1\in\reals$
                        if $\sumFiniteSequence{b_1}{t_1}$ and $t_1 = 0$ and for all $i\in J$ we have $0 \leq b_1(i)$ then
                        for all $i\in J$ we have $b_1(i) = 0$.
                    Thus for all $i\in J$ we have $b_0(i) = 0$.
                \end{subproof}
                Show that for all $i\in J_1$ we have $b(i) = 0$.
                \begin{subproof}
                    Fix $i$.
                    Assume $i\in J_1$.
                    Thus $i\in\naturalsPlus$ and $i \leq k + 1$ by \cref{initial_segment_naturalsplus}.
                    $k + 1\in\reals$ by \cref{real_naturals_subset_reals,subseteq}.
                    $i\in\reals$ by \cref{real_naturals_subset_reals,subseteq}.
                    Thus $i < k + 1$ or $i = k + 1$ by \cref{real_leq_split}.
                    \begin{byCase}
                    \caseOf{$i = k + 1$.}
                        Thus $b(i) = 0$.
                    \caseOf{$i < k + 1$.}
                        Thus $i \leq k + 1$.
                        Thus $i \neq k + 1$.
                        $i\in\naturalsPlus$ by \cref{initial_segment_naturalsplus}.
                        $k\in\naturalsPlus$ by \cref{subseteq}.
                        Thus $i \leq k$ by \cref{real_naturals_leq_succ_elim}.
                        Thus $i\in J$ by \cref{initial_segment_naturalsplus}.
                        Thus $\dom{b} = J_1$ by \cref{funs_elim}.
                        Thus $J\subseteq\dom{b}$ by \cref{subseteq}.
                        $b$ is a function by \cref{funs_elim}.
                        Thus $b(i) = b_0(i)$ by \cref{restrl_of_function_apply}.
                        Thus $b(i) = 0$.
                    \end{byCase}
                \end{subproof}
            \end{subproof}
            Thus $k + 1\in A$.
        \end{subproof}
    \end{subproof}
    Thus for all $m\in\naturalsPlus$ we have $m\in A$ by \cref{real_naturals_induction}.
    Thus $n\in A$ by \cref{subseteq}.
    Thus for all $b_1\in\funs{J}{\reals}$ for all $t_1\in\reals$
        if $\sumFiniteSequence{b_1}{t_1}$ and $t_1 = 0$ and for all $i\in J$ we have $0 \leq b_1(i)$ then
        for all $i\in J$ we have $b_1(i) = 0$.
    Thus for all $i\in J$ we have $a(i) = 0$.
\end{proof}

% from §20 helper lemma 39: sum of squares of a scalar sequence
\begin{lemma}\label{sum_finite_sequence_square_scalar}
    Suppose $n\in\naturalsPlus$.
    Let $J = \initialSegment{n}$.
    Suppose $c\in\reals$.
    Suppose $a\in\funs{J}{\reals}$.
    Suppose $\sumFiniteSequence{\coordSquareSequenceMul{n}{a}}{s}$.
    Let $b = \coordScalarSequence{n}{c}{a}$.
    Then $\sumFiniteSequence{\coordSquareSequenceMul{n}{b}}{(c \rmul c) \rmul s}$.
\end{lemma}
\begin{proof}
    Let $u = \coordSquareSequenceMul{n}{a}$.
    Let $v = \coordSquareSequenceMul{n}{b}$.
    Let $w = \coordScalarSequence{n}{c \rmul c}{u}$.
    $a\in\funs{J}{\reals}$.
    Thus $u\in\funs{J}{\reals}$ by \cref{coord_square_sequence_mul_function}.
    $b\in\funs{J}{\reals}$ by \cref{coord_scalar_sequence_function}.
    Thus $v\in\funs{J}{\reals}$ by \cref{coord_square_sequence_mul_function}.
    $c \rmul c\in\reals$ by \cref{real_mul_closed}.
    Thus $w\in\funs{J}{\reals}$ by \cref{coord_scalar_sequence_function}.
    Show that $v = w$.
    \begin{subproof}
        Show that for all $i\in J$ we have $v(i) = w(i)$.
        \begin{subproof}
            Fix $i$.
            Assume $i\in J$.
            Thus $(i, c \rmul a(i))\in b$ by \cref{coord_scalar_sequence}.
            $b$ is a function by \cref{funs_elim}.
            Thus $b(i) = c \rmul a(i)$ by \cref{function_apply_intro}.
            Thus $(i, b(i) \rmul b(i))\in v$ by \cref{coord_square_sequence_mul}.
            $v$ is a function by \cref{funs_elim}.
            Thus $v(i) = b(i) \rmul b(i)$ by \cref{function_apply_intro}.
            Thus $v(i) = (c \rmul a(i)) \rmul (c \rmul a(i))$.
            $a(i)\in\reals$ by \cref{funs_type_apply}.
            Thus $c \rmul a(i)\in\reals$ by \cref{real_mul_closed}.
            Thus $c \rmul c\in\reals$ by \cref{real_mul_closed}.
            Thus $a(i) \rmul a(i)\in\reals$ by \cref{real_mul_closed}.
            Thus $(c \rmul a(i)) \rmul (c \rmul a(i)) =
                ((c \rmul a(i)) \rmul c) \rmul a(i)$ by \cref{real_mul_assoc}.
            Thus $(c \rmul a(i)) \rmul c = c \rmul (a(i) \rmul c)$ by \cref{real_mul_assoc}.
            Thus $((c \rmul a(i)) \rmul c) \rmul a(i) =
                (c \rmul (a(i) \rmul c)) \rmul a(i)$ by \cref{real_mul_assoc}.
            Thus $a(i) \rmul c = c \rmul a(i)$ by \cref{real_mul_comm}.
            Thus $(c \rmul (a(i) \rmul c)) \rmul a(i) =
                (c \rmul (c \rmul a(i))) \rmul a(i)$ by \cref{real_mul_assoc}.
            Thus $c \rmul (c \rmul a(i)) = (c \rmul c) \rmul a(i)$ by \cref{real_mul_assoc}.
            Thus $(c \rmul (c \rmul a(i))) \rmul a(i) =
                ((c \rmul c) \rmul a(i)) \rmul a(i)$ by \cref{real_mul_assoc}.
            Thus $((c \rmul c) \rmul a(i)) \rmul a(i) =
                (c \rmul c) \rmul (a(i) \rmul a(i))$ by \cref{real_mul_assoc}.
            Thus $(c \rmul a(i)) \rmul (c \rmul a(i)) =
                (c \rmul c) \rmul (a(i) \rmul a(i))$ by \cref{real_mul_assoc}.
            Thus $v(i) = (c \rmul c) \rmul (a(i) \rmul a(i))$.
            Thus $(i, a(i) \rmul a(i))\in u$ by \cref{coord_square_sequence_mul}.
            $u$ is a function by \cref{funs_elim}.
            Thus $u(i) = a(i) \rmul a(i)$ by \cref{function_apply_intro}.
            Thus $(i, (c \rmul c) \rmul u(i))\in w$ by \cref{coord_scalar_sequence}.
            $w$ is a function by \cref{funs_elim}.
            Thus $w(i) = (c \rmul c) \rmul u(i)$ by \cref{function_apply_intro}.
            Thus $w(i) = (c \rmul c) \rmul (a(i) \rmul a(i))$.
            Thus $v(i) = w(i)$.
        \end{subproof}
        $v$ is a function and $w$ is a function by \cref{funs_elim}.
        $\dom{v} = J$ and $\dom{w} = J$ by \cref{funs_elim}.
        Show that $v\subseteq w$.
        \begin{subproof}
            Thus $\dom{v}\subseteq \dom{w}$ by \cref{subseteq}.
            Show that for all $x\in\dom{v}$ we have $v(x) = w(x)$.
            \begin{subproof}
                Fix $x$.
                Assume $x\in\dom{v}$.
                Thus $x\in J$ by \cref{subseteq}.
                Thus $v(x) = w(x)$.
            \end{subproof}
            Thus $v\subseteq w$ by \cref{fun_subseteq}.
        \end{subproof}
        Show that $w\subseteq v$.
        \begin{subproof}
            Thus $\dom{w}\subseteq \dom{v}$ by \cref{subseteq}.
            Show that for all $x\in\dom{w}$ we have $w(x) = v(x)$.
            \begin{subproof}
                Fix $x$.
                Assume $x\in\dom{w}$.
                Thus $x\in J$ by \cref{subseteq}.
                Thus $w(x) = v(x)$.
            \end{subproof}
            Thus $w\subseteq v$ by \cref{fun_subseteq}.
        \end{subproof}
        Thus $v = w$ by \cref{subseteq_antisymmetric}.
    \end{subproof}
    Thus $\sumFiniteSequence{w}{(c \rmul c) \rmul s}$ by \cref{sum_finite_sequence_scalar}.
    Thus $\sumFiniteSequence{v}{(c \rmul c) \rmul s}$.
\end{proof}

% from §20 helper lemma 40: sum of squares of a coordinate sum sequence
\begin{lemma}\label{sum_finite_sequence_square_sum}
    Suppose $n\in\naturalsPlus$.
    Let $J = \initialSegment{n}$.
    Suppose $a\in\funs{J}{\reals}$.
    Suppose $b\in\funs{J}{\reals}$.
    Suppose $\sumFiniteSequence{\coordSquareSequenceMul{n}{a}}{s_1}$.
    Suppose $\sumFiniteSequence{\coordSquareSequenceMul{n}{b}}{s_2}$.
    Suppose $\sumFiniteSequence{\coordProductSequence{n}{a}{b}}{s}$.
    Let $c = \coordSumSequence{n}{a}{b}$.
    Then $\sumFiniteSequence{\coordSquareSequenceMul{n}{c}}{s_1 + s_2 + (1+1) \rmul s}$.
\end{lemma}
\begin{proof}
    Let $u = \coordSquareSequenceMul{n}{a}$.
    Let $v = \coordSquareSequenceMul{n}{b}$.
    Let $p = \coordProductSequence{n}{a}{b}$.
    $1\in\naturalsPlus$ by \cref{real_one_in_naturals}.
    $1\in\reals$ by \cref{real_naturals_subset_reals,subseteq}.
    $1 + 1\in\reals$ by \cref{real_add_closed}.
    Let $p_2 = \coordScalarSequence{n}{1+1}{p}$.
    Let $w = \coordSumSequence{n}{u}{v}$.
    Let $d = \coordSumSequence{n}{w}{p_2}$.
    Let $q = \coordSquareSequenceMul{n}{c}$.
    $u\in\funs{J}{\reals}$ by \cref{coord_square_sequence_mul_function}.
    $v\in\funs{J}{\reals}$ by \cref{coord_square_sequence_mul_function}.
    $p\in\funs{J}{\reals}$ by \cref{coord_product_sequence_function}.
    Thus $p_2\in\funs{J}{\reals}$ by \cref{coord_scalar_sequence_function}.
    Thus $w\in\funs{J}{\reals}$ by \cref{coord_sum_sequence_function}.
    Thus $d\in\funs{J}{\reals}$ by \cref{coord_sum_sequence_function}.
    $c\in\funs{J}{\reals}$ by \cref{coord_sum_sequence_function}.
    Thus $q\in\funs{J}{\reals}$ by \cref{coord_square_sequence_mul_function}.
    Show that $q = d$.
    \begin{subproof}
        Show that for all $i\in J$ we have $q(i) = d(i)$.
        \begin{subproof}
            Fix $i$.
            Assume $i\in J$.
            Thus $a(i)\in\reals$ and $b(i)\in\reals$ by \cref{funs_type_apply}.
            Thus $a(i) + b(i)\in\reals$ by \cref{real_add_closed}.
            Thus $(i, a(i) + b(i))\in c$ by \cref{coord_sum_sequence}.
            $c$ is a function by \cref{funs_elim}.
            Thus $c(i) = a(i) + b(i)$ by \cref{function_apply_intro}.
            Thus $(i, c(i) \rmul c(i))\in q$ by \cref{coord_square_sequence_mul}.
            $q$ is a function by \cref{funs_elim}.
            Thus $q(i) = c(i) \rmul c(i)$ by \cref{function_apply_intro}.
            Thus $q(i) = (a(i) + b(i)) \rmul (a(i) + b(i))$.
            Thus $(a(i) + b(i)) \rmul (a(i) + b(i)) =
                (a(i) + b(i)) \rmul a(i) + (a(i) + b(i)) \rmul b(i)$ by \cref{real_distrib}.
            Thus $(a(i) + b(i)) \rmul a(i) = a(i) \rmul (a(i) + b(i))$ by \cref{real_mul_comm}.
            Thus $a(i) \rmul (a(i) + b(i)) = (a(i) \rmul a(i)) + (a(i) \rmul b(i))$ by \cref{real_distrib}.
            Thus $(a(i) + b(i)) \rmul a(i) = (a(i) \rmul a(i)) + (a(i) \rmul b(i))$.
            Thus $(a(i) + b(i)) \rmul b(i) = b(i) \rmul (a(i) + b(i))$ by \cref{real_mul_comm}.
            Thus $b(i) \rmul (a(i) + b(i)) = (b(i) \rmul a(i)) + (b(i) \rmul b(i))$ by \cref{real_distrib}.
            Thus $(a(i) + b(i)) \rmul b(i) = (b(i) \rmul a(i)) + (b(i) \rmul b(i))$.
            Thus $q(i) =
                ((a(i) \rmul a(i)) + (a(i) \rmul b(i))) + ((b(i) \rmul a(i)) + (b(i) \rmul b(i)))$ by \cref{real_add_eq_subst}.
            Thus $a(i) \rmul a(i)\in\reals$ and $a(i) \rmul b(i)\in\reals$ and $b(i) \rmul b(i)\in\reals$
                by \cref{real_mul_closed}.
            Thus $b(i) \rmul a(i) = a(i) \rmul b(i)$ by \cref{real_mul_comm}.
            Thus $q(i) =
                ((a(i) \rmul a(i)) + (a(i) \rmul b(i))) + ((a(i) \rmul b(i)) + (b(i) \rmul b(i)))$ by \cref{real_add_eq_subst}.
            Thus $(a(i) \rmul b(i)) + (b(i) \rmul b(i))\in\reals$ by \cref{real_add_closed}.
            Thus $q(i) =
                (a(i) \rmul a(i)) + ((a(i) \rmul b(i)) + ((a(i) \rmul b(i)) + (b(i) \rmul b(i))))$ by \cref{real_add_assoc}.
            Thus $(a(i) \rmul b(i)) + ((a(i) \rmul b(i)) + (b(i) \rmul b(i))) =
                ((a(i) \rmul b(i)) + (a(i) \rmul b(i))) + (b(i) \rmul b(i))$ by \cref{real_add_assoc}.
            Thus $q(i) =
                (a(i) \rmul a(i)) + (((a(i) \rmul b(i)) + (a(i) \rmul b(i))) + (b(i) \rmul b(i)))$ by \cref{real_add_eq_subst}.
            Thus $(a(i) \rmul b(i)) + (a(i) \rmul b(i))\in\reals$ by \cref{real_add_closed}.
            Thus $q(i) =
                ((a(i) \rmul a(i)) + ((a(i) \rmul b(i)) + (a(i) \rmul b(i)))) + (b(i) \rmul b(i))$ by \cref{real_add_assoc}.
            Thus $(1 + 1) \rmul (a(i) \rmul b(i)) = (a(i) \rmul b(i)) \rmul (1 + 1)$ by \cref{real_mul_comm}.
            Thus $(a(i) \rmul b(i)) \rmul (1 + 1) =
                (a(i) \rmul b(i)) \rmul 1 + (a(i) \rmul b(i)) \rmul 1$ by \cref{real_distrib}.
            Thus $(a(i) \rmul b(i)) \rmul 1 = 1 \rmul (a(i) \rmul b(i))$ by \cref{real_mul_comm}.
            Thus $1 \rmul (a(i) \rmul b(i)) = a(i) \rmul b(i)$ by \cref{real_mul_ident}.
            Thus $(a(i) \rmul b(i)) \rmul 1 = a(i) \rmul b(i)$.
            Thus $(1 + 1) \rmul (a(i) \rmul b(i)) =
                (a(i) \rmul b(i)) + (a(i) \rmul b(i))$ by \cref{real_add_eq_subst}.
            Thus $q(i) =
                ((a(i) \rmul a(i)) + ((1 + 1) \rmul (a(i) \rmul b(i)))) + (b(i) \rmul b(i))$ by \cref{real_add_eq_subst}.
            Thus $q(i) =
                (a(i) \rmul a(i)) + ((1 + 1) \rmul (a(i) \rmul b(i))) + (b(i) \rmul b(i))$ by \cref{real_add_assoc}.
            Thus $(i, a(i) \rmul a(i))\in u$ by \cref{coord_square_sequence_mul}.
            Thus $(i, b(i) \rmul b(i))\in v$ by \cref{coord_square_sequence_mul}.
            Thus $(i, a(i) \rmul b(i))\in p$ by \cref{coord_product_sequence}.
            $u$ is a function by \cref{funs_elim}.
            $v$ is a function by \cref{funs_elim}.
            $p$ is a function by \cref{funs_elim}.
            Thus $u(i) = a(i) \rmul a(i)$ and $v(i) = b(i) \rmul b(i)$ and $p(i) = a(i) \rmul b(i)$
                by \cref{function_apply_intro}.
            Thus $(i, (1 + 1) \rmul p(i))\in p_2$ by \cref{coord_scalar_sequence}.
            $p_2$ is a function by \cref{funs_elim}.
            Thus $p_2(i) = (1 + 1) \rmul p(i)$ by \cref{function_apply_intro}.
            Thus $p_2(i) = (1 + 1) \rmul (a(i) \rmul b(i))$.
            Thus $(i, u(i) + v(i))\in w$ by \cref{coord_sum_sequence}.
            $w$ is a function by \cref{funs_elim}.
            Thus $w(i) = u(i) + v(i)$ by \cref{function_apply_intro}.
            Thus $(i, w(i) + p_2(i))\in d$ by \cref{coord_sum_sequence}.
            $d$ is a function by \cref{funs_elim}.
            Thus $d(i) = w(i) + p_2(i)$ by \cref{function_apply_intro}.
            Thus $w(i) = (a(i) \rmul a(i)) + (b(i) \rmul b(i))$ by \cref{real_add_eq_subst}.
            Thus $d(i) =
                ((a(i) \rmul a(i)) + (b(i) \rmul b(i))) + ((1 + 1) \rmul (a(i) \rmul b(i)))$ by \cref{real_add_eq_subst}.
            Thus $(1 + 1) \rmul (a(i) \rmul b(i))\in\reals$ by \cref{real_mul_closed}.
            Thus $d(i) =
                (a(i) \rmul a(i)) + ((1 + 1) \rmul (a(i) \rmul b(i))) + (b(i) \rmul b(i))$ by \cref{real_add_assoc,real_add_comm}.
            Thus $q(i) = d(i)$.
        \end{subproof}
        $q$ is a function and $d$ is a function by \cref{funs_elim}.
        $\dom{q} = J$ and $\dom{d} = J$ by \cref{funs_elim}.
        Show that $q\subseteq d$.
        \begin{subproof}
            Thus $\dom{q}\subseteq \dom{d}$ by \cref{subseteq}.
            Show that for all $x\in\dom{q}$ we have $q(x) = d(x)$.
            \begin{subproof}
                Fix $x$.
                Assume $x\in\dom{q}$.
                Thus $x\in J$ by \cref{subseteq}.
                Thus $q(x) = d(x)$.
            \end{subproof}
            Thus $q\subseteq d$ by \cref{fun_subseteq}.
        \end{subproof}
        Show that $d\subseteq q$.
        \begin{subproof}
            Thus $\dom{d}\subseteq \dom{q}$ by \cref{subseteq}.
            Show that for all $x\in\dom{d}$ we have $d(x) = q(x)$.
            \begin{subproof}
                Fix $x$.
                Assume $x\in\dom{d}$.
                Thus $x\in J$ by \cref{subseteq}.
                Thus $d(x) = q(x)$.
            \end{subproof}
            Thus $d\subseteq q$ by \cref{fun_subseteq}.
        \end{subproof}
        Thus $q = d$ by \cref{subseteq_antisymmetric}.
    \end{subproof}
    Thus $\sumFiniteSequence{w}{s_1 + s_2}$ by \cref{sum_finite_sequence_add}.
    Thus $\sumFiniteSequence{p_2}{(1+1) \rmul s}$ by \cref{sum_finite_sequence_scalar}.
    Thus $\sumFiniteSequence{d}{(s_1 + s_2) + (1+1) \rmul s}$ by \cref{sum_finite_sequence_add}.
    Thus $\sumFiniteSequence{q}{(s_1 + s_2) + (1+1) \rmul s}$.
    Thus $\sumFiniteSequence{\coordSquareSequenceMul{n}{c}}{s_1 + s_2 + (1+1) \rmul s}$
        by \cref{real_add_assoc}.
\end{proof}

% from §20 helper lemma 41: Cauchy-Schwarz for finite sequences
\begin{lemma}\label{cauchy_schwarz_finite_sequences}
    Suppose $n\in\naturalsPlus$.
    Let $J = \initialSegment{n}$.
    Suppose $a\in\funs{J}{\reals}$.
    Suppose $b\in\funs{J}{\reals}$.
    Suppose $\sumFiniteSequence{\coordProductSequence{n}{a}{b}}{s}$.
    Suppose $\sumFiniteSequence{\coordSquareSequenceMul{n}{a}}{s_1}$.
    Suppose $\sumFiniteSequence{\coordSquareSequenceMul{n}{b}}{s_2}$.
    Then $s \rmul s \leq s_1 \rmul s_2$.
\end{lemma}
\begin{proof}
    Let $a_1 = \coordScalarSequence{n}{s_2}{a}$.
    Let $b_1 = \coordScalarSequence{n}{\neg{s}}{b}$.
    Let $c = \coordSumSequence{n}{a_1}{b_1}$.
    Let $p = \coordProductSequence{n}{a}{b}$.
    Let $p_1 = \coordProductSequence{n}{a_1}{b_1}$.
    Show that $s_2\in\reals$.
    \begin{subproof}
        Take $n_0$ such that $n_0\in\naturalsPlus$ and
            $\coordSquareSequenceMul{n}{b}\in\funs{\initialSegment{n_0}}{\reals}$ and
            there exists $t\in\funs{\initialSegment{n_0}}{\reals}$ such that
                $t(1) = \apply{\coordSquareSequenceMul{n}{b}}{1}$ and
                $s_2 = t(n_0)$ and
                for all $k\in\naturalsPlus$ such that $k + 1 \leq n_0$ we have
                    $t(k + 1) = t(k) + \apply{\coordSquareSequenceMul{n}{b}}{k + 1}$
            by \cref{sum_finite_sequence_pred}.
        Take $t$ such that $t\in\funs{\initialSegment{n_0}}{\reals}$ and
            $t(1) = \apply{\coordSquareSequenceMul{n}{b}}{1}$ and
            $s_2 = t(n_0)$ and
            for all $k\in\naturalsPlus$ such that $k + 1 \leq n_0$ we have
                $t(k + 1) = t(k) + \apply{\coordSquareSequenceMul{n}{b}}{k + 1}$.
        $n_0\in\naturalsPlus$.
        Thus $n_0\leq n_0$.
        Thus $n_0\in\initialSegment{n_0}$ by \cref{initial_segment_naturalsplus}.
        Thus $t(n_0)\in\reals$ by \cref{funs_type_apply}.
        Thus $s_2\in\reals$.
    \end{subproof}
    Show that $s\in\reals$.
    \begin{subproof}
        Take $n_1$ such that $n_1\in\naturalsPlus$ and
            $\coordProductSequence{n}{a}{b}\in\funs{\initialSegment{n_1}}{\reals}$ and
            there exists $u\in\funs{\initialSegment{n_1}}{\reals}$ such that
                $u(1) = \apply{\coordProductSequence{n}{a}{b}}{1}$ and
                $s = u(n_1)$ and
                for all $k\in\naturalsPlus$ such that $k + 1 \leq n_1$ we have
                    $u(k + 1) = u(k) + \apply{\coordProductSequence{n}{a}{b}}{k + 1}$
            by \cref{sum_finite_sequence_pred}.
        Take $u$ such that $u\in\funs{\initialSegment{n_1}}{\reals}$ and
            $u(1) = \apply{\coordProductSequence{n}{a}{b}}{1}$ and
            $s = u(n_1)$ and
            for all $k\in\naturalsPlus$ such that $k + 1 \leq n_1$ we have
                $u(k + 1) = u(k) + \apply{\coordProductSequence{n}{a}{b}}{k + 1}$.
        $n_1\in\naturalsPlus$.
        Thus $n_1\leq n_1$.
        Thus $n_1\in\initialSegment{n_1}$ by \cref{initial_segment_naturalsplus}.
        Thus $u(n_1)\in\reals$ by \cref{funs_type_apply}.
        Thus $s\in\reals$.
    \end{subproof}
    Show that $s_1\in\reals$.
    \begin{subproof}
        Take $n_2$ such that $n_2\in\naturalsPlus$ and
            $\coordSquareSequenceMul{n}{a}\in\funs{\initialSegment{n_2}}{\reals}$ and
            there exists $v\in\funs{\initialSegment{n_2}}{\reals}$ such that
                $v(1) = \apply{\coordSquareSequenceMul{n}{a}}{1}$ and
                $s_1 = v(n_2)$ and
                for all $k\in\naturalsPlus$ such that $k + 1 \leq n_2$ we have
                    $v(k + 1) = v(k) + \apply{\coordSquareSequenceMul{n}{a}}{k + 1}$
            by \cref{sum_finite_sequence_pred}.
        Take $v$ such that $v\in\funs{\initialSegment{n_2}}{\reals}$ and
            $v(1) = \apply{\coordSquareSequenceMul{n}{a}}{1}$ and
            $s_1 = v(n_2)$ and
            for all $k\in\naturalsPlus$ such that $k + 1 \leq n_2$ we have
                $v(k + 1) = v(k) + \apply{\coordSquareSequenceMul{n}{a}}{k + 1}$.
        $n_2\in\naturalsPlus$.
        Thus $n_2\leq n_2$.
        Thus $n_2\in\initialSegment{n_2}$ by \cref{initial_segment_naturalsplus}.
        Thus $v(n_2)\in\reals$ by \cref{funs_type_apply}.
        Thus $s_1\in\reals$.
    \end{subproof}
    Thus $\neg{s}\in\reals$ by \cref{real_add_inverse}.
    $a_1\in\funs{J}{\reals}$ by \cref{coord_scalar_sequence_function}.
    $b_1\in\funs{J}{\reals}$ by \cref{coord_scalar_sequence_function}.
    $p\in\funs{J}{\reals}$ by \cref{coord_product_sequence_function}.
    Thus $p_1\in\funs{J}{\reals}$ by \cref{coord_product_sequence_function}.
    Show that $\sumFiniteSequence{\coordSquareSequenceMul{n}{a_1}}{(s_2 \rmul s_2) \rmul s_1}$.
    \begin{subproof}
        Thus $\sumFiniteSequence{\coordSquareSequenceMul{n}{a_1}}{(s_2 \rmul s_2) \rmul s_1}$ by \cref{sum_finite_sequence_square_scalar}.
    \end{subproof}
    Show that $\sumFiniteSequence{\coordSquareSequenceMul{n}{b_1}}{(s \rmul s) \rmul s_2}$.
    \begin{subproof}
        Thus $\sumFiniteSequence{\coordSquareSequenceMul{n}{b_1}}{((\neg{s}) \rmul (\neg{s})) \rmul s_2}$ by \cref{sum_finite_sequence_square_scalar}.
        Show that $(\neg{s}) \rmul (\neg{s}) = s \rmul s$.
        \begin{subproof}
            $(\neg{s}) \rmul (\neg{s}) = \neg{((\neg{s}) \rmul s)}$ by \cref{real_mul_neg_right}.
            $(\neg{s}) \rmul s = s \rmul (\neg{s})$ by \cref{real_mul_comm}.
            $s \rmul (\neg{s}) = \neg{(s \rmul s)}$ by \cref{real_mul_neg_right}.
            Thus $(\neg{s}) \rmul (\neg{s}) = \neg{\neg{(s \rmul s)}}$.
            Thus $s \rmul s\in\reals$ by \cref{real_mul_closed}.
            Thus $(\neg{s}) \rmul (\neg{s}) = s \rmul s$ by \cref{real_neg_neg}.
        \end{subproof}
        Thus $\sumFiniteSequence{\coordSquareSequenceMul{n}{b_1}}{(s \rmul s) \rmul s_2}$.
    \end{subproof}
    Show that $\sumFiniteSequence{p_1}{(s_2 \rmul (\neg{s})) \rmul s}$.
    \begin{subproof}
        Let $r = \coordScalarSequence{n}{s_2 \rmul (\neg{s})}{p}$.
        $s_2 \rmul (\neg{s})\in\reals$ by \cref{real_mul_closed}.
        Thus $r\in\funs{J}{\reals}$ by \cref{coord_scalar_sequence_function}.
        Show that $p_1 = r$.
        \begin{subproof}
            Show that for all $i\in J$ we have $p_1(i) = r(i)$.
            \begin{subproof}
                Fix $i$.
                Assume $i\in J$.
                Thus $(i, s_2 \rmul a(i))\in a_1$ by \cref{coord_scalar_sequence}.
                Thus $(i, \neg{s} \rmul b(i))\in b_1$ by \cref{coord_scalar_sequence}.
                $a_1$ is a function by \cref{funs_elim}.
                $b_1$ is a function by \cref{funs_elim}.
                Thus $a_1(i) = s_2 \rmul a(i)$ and $b_1(i) = \neg{s} \rmul b(i)$
                    by \cref{function_apply_intro}.
                Thus $(i, a_1(i) \rmul b_1(i))\in p_1$ by \cref{coord_product_sequence}.
                $p_1$ is a function by \cref{funs_elim}.
                Thus $p_1(i) = a_1(i) \rmul b_1(i)$ by \cref{function_apply_intro}.
                Thus $p_1(i) = (s_2 \rmul a(i)) \rmul (\neg{s} \rmul b(i))$.
                Thus $a(i)\in\reals$ and $b(i)\in\reals$ by \cref{funs_type_apply}.
                Thus $s_2 \rmul a(i)\in\reals$ and $\neg{s} \rmul b(i)\in\reals$ by \cref{real_mul_closed}.
                Thus $(s_2 \rmul a(i)) \rmul (\neg{s} \rmul b(i)) =
                    s_2 \rmul (a(i) \rmul (\neg{s} \rmul b(i)))$ by \cref{real_mul_assoc}.
                Thus $a(i) \rmul (\neg{s} \rmul b(i)) = (a(i) \rmul \neg{s}) \rmul b(i)$ by \cref{real_mul_assoc}.
                Thus $p_1(i) = s_2 \rmul ((a(i) \rmul \neg{s}) \rmul b(i))$ by \cref{real_mul_assoc}.
                Thus $p_1(i) = (s_2 \rmul (a(i) \rmul \neg{s})) \rmul b(i)$ by \cref{real_mul_assoc}.
                Thus $a(i) \rmul \neg{s} = \neg{s} \rmul a(i)$ by \cref{real_mul_comm}.
                Thus $p_1(i) = (s_2 \rmul (\neg{s} \rmul a(i))) \rmul b(i)$ by \cref{real_mul_assoc}.
                Thus $p_1(i) = ((s_2 \rmul \neg{s}) \rmul a(i)) \rmul b(i)$ by \cref{real_mul_assoc}.
                Thus $p_1(i) = (s_2 \rmul \neg{s}) \rmul (a(i) \rmul b(i))$ by \cref{real_mul_assoc}.
                Thus $(i, a(i) \rmul b(i))\in p$ by \cref{coord_product_sequence}.
                $p$ is a function by \cref{funs_elim}.
                Thus $p(i) = a(i) \rmul b(i)$ by \cref{function_apply_intro}.
                Thus $(i, (s_2 \rmul (\neg{s})) \rmul p(i))\in r$ by \cref{coord_scalar_sequence}.
                $r$ is a function by \cref{funs_elim}.
                Thus $r(i) = (s_2 \rmul (\neg{s})) \rmul p(i)$ by \cref{function_apply_intro}.
                Thus $r(i) = (s_2 \rmul (\neg{s})) \rmul (a(i) \rmul b(i))$.
                Thus $p_1(i) = r(i)$.
            \end{subproof}
            $p_1$ is a function and $r$ is a function by \cref{funs_elim}.
            $\dom{p_1} = J$ and $\dom{r} = J$ by \cref{funs_elim}.
            Show that $p_1\subseteq r$.
            \begin{subproof}
                Thus $\dom{p_1}\subseteq \dom{r}$ by \cref{subseteq}.
                Show that for all $x\in\dom{p_1}$ we have $p_1(x) = r(x)$.
                \begin{subproof}
                    Fix $x$.
                    Assume $x\in\dom{p_1}$.
                    Thus $x\in J$ by \cref{subseteq}.
                    Thus $p_1(x) = r(x)$.
                \end{subproof}
                Thus $p_1\subseteq r$ by \cref{fun_subseteq}.
            \end{subproof}
            Show that $r\subseteq p_1$.
            \begin{subproof}
                Thus $\dom{r}\subseteq \dom{p_1}$ by \cref{subseteq}.
                Show that for all $x\in\dom{r}$ we have $r(x) = p_1(x)$.
                \begin{subproof}
                    Fix $x$.
                    Assume $x\in\dom{r}$.
                    Thus $x\in J$ by \cref{subseteq}.
                    Thus $r(x) = p_1(x)$.
                \end{subproof}
                Thus $r\subseteq p_1$ by \cref{fun_subseteq}.
            \end{subproof}
            Thus $p_1 = r$ by \cref{subseteq_antisymmetric}.
        \end{subproof}
        Thus $\sumFiniteSequence{r}{(s_2 \rmul (\neg{s})) \rmul s}$ by \cref{sum_finite_sequence_scalar}.
        Thus $\sumFiniteSequence{p_1}{(s_2 \rmul (\neg{s})) \rmul s}$.
    \end{subproof}
    $c\in\funs{J}{\reals}$ by \cref{coord_sum_sequence_function}.
    Thus $\coordSquareSequenceMul{n}{c}\in\funs{J}{\reals}$ by \cref{coord_square_sequence_mul_function}.
    Take $S\in\reals$ such that $\sumFiniteSequence{\coordSquareSequenceMul{n}{c}}{S}$ by \cref{sum_finite_sequence_exists}.
    Thus $\sumFiniteSequence{\coordSquareSequenceMul{n}{c}}{(s_2 \rmul s_2) \rmul s_1 + (s \rmul s) \rmul s_2 + (1+1) \rmul ((s_2 \rmul (\neg{s})) \rmul s)}$
        by \cref{sum_finite_sequence_square_sum}.
    Thus $S = (s_2 \rmul s_2) \rmul s_1 + (s \rmul s) \rmul s_2 + (1+1) \rmul ((s_2 \rmul (\neg{s})) \rmul s)$
        by \cref{sum_finite_sequence_unique}.
    Show that $0 \leq S$.
    \begin{subproof}
        Show that for all $i\in\dom{\coordSquareSequenceMul{n}{c}}$ we have
            $0 \leq \apply{\coordSquareSequenceMul{n}{c}}{i}$.
        \begin{subproof}
            Fix $i$.
            Assume $i\in\dom{\coordSquareSequenceMul{n}{c}}$.
            Thus $i\in J$ by \cref{funs_elim}.
            Thus $(i, c(i) \rmul c(i))\in \coordSquareSequenceMul{n}{c}$ by \cref{coord_square_sequence_mul}.
            Thus $\apply{\coordSquareSequenceMul{n}{c}}{i} = c(i) \rmul c(i)$ by \cref{function_apply_intro,funs_elim}.
            Thus $c(i)\in\reals$ by \cref{funs_type_apply}.
            Thus $\exp{c(i)}{1+1} \geq 0$ by \cref{real_sq_nonnegative}.
            $1\in\naturalsPlus$ by \cref{real_one_in_naturals}.
            $\exp{c(i)}{1+1} = \exp{c(i)}{1} \rmul c(i)$ by \cref{real_pow_succ}.
            $\exp{c(i)}{1} = c(i)$ by \cref{real_pow_one}.
            Thus $\exp{c(i)}{1+1} = c(i) \rmul c(i)$.
            Thus $0 \leq c(i) \rmul c(i)$.
            Thus $0 \leq \apply{\coordSquareSequenceMul{n}{c}}{i}$.
        \end{subproof}
        Thus $0 \leq S$ by \cref{sum_finite_sequence_nonneg}.
    \end{subproof}
    Show that $S = s_2 \rmul (s_2 \rmul s_1 - s \rmul s)$.
    \begin{subproof}
        $(s_2 \rmul (\neg{s})) \rmul s = s_2 \rmul ((\neg{s}) \rmul s)$ by \cref{real_mul_assoc}.
        $(\neg{s}) \rmul s = s \rmul (\neg{s})$ by \cref{real_mul_comm}.
        $s \rmul (\neg{s}) = \neg{(s \rmul s)}$ by \cref{real_mul_neg_right}.
        Thus $(s_2 \rmul (\neg{s})) \rmul s = s_2 \rmul \neg{(s \rmul s)}$ by \cref{real_mul_assoc}.
        Thus $s \rmul s\in\reals$ by \cref{real_mul_closed}.
        Thus $(s_2 \rmul (\neg{s})) \rmul s = \neg{(s_2 \rmul (s \rmul s))}$ by \cref{real_mul_neg_right}.
        Thus $(1+1) \rmul ((s_2 \rmul (\neg{s})) \rmul s) =
            (1+1) \rmul \neg{(s_2 \rmul (s \rmul s))}$.
        $1\in\naturalsPlus$ by \cref{real_one_in_naturals}.
        $1\in\reals$ by \cref{real_naturals_subset_reals,subseteq}.
        $1 + 1\in\reals$ by \cref{real_add_closed}.
        Thus $s_2 \rmul (s \rmul s)\in\reals$ by \cref{real_mul_closed}.
        Thus $(1+1) \rmul \neg{(s_2 \rmul (s \rmul s))} =
            \neg{((1+1) \rmul (s_2 \rmul (s \rmul s)))}$ by \cref{real_mul_neg_right}.
        Thus $S = (s_2 \rmul s_2) \rmul s_1 + (s \rmul s) \rmul s_2 +
            \neg{((1+1) \rmul (s_2 \rmul (s \rmul s)))}$ by \cref{real_add_eq_subst}.
        $(s \rmul s) \rmul s_2 = s_2 \rmul (s \rmul s)$ by \cref{real_mul_comm}.
        Thus $S = s_2 \rmul (s_2 \rmul s_1) + s_2 \rmul (s \rmul s) +
            \neg{((1+1) \rmul (s_2 \rmul (s \rmul s)))}$ by \cref{real_mul_assoc}.
        Thus $s_2 \rmul (s_2 \rmul s_1)\in\reals$ by \cref{real_mul_closed}.
        Thus $s_2 \rmul (s \rmul s)\in\reals$ by \cref{real_mul_closed}.
        Thus $\neg{((1+1) \rmul (s_2 \rmul (s \rmul s)))}\in\reals$ by \cref{real_add_inverse,real_mul_closed}.
        Thus $S = s_2 \rmul (s_2 \rmul s_1) +
            (s_2 \rmul (s \rmul s) + \neg{((1+1) \rmul (s_2 \rmul (s \rmul s)))})$ by \cref{real_add_assoc}.
        Thus $(1+1) \rmul (s_2 \rmul (s \rmul s)) =
            (s_2 \rmul (s \rmul s)) \rmul (1+1)$ by \cref{real_mul_comm}.
        Thus $(s_2 \rmul (s \rmul s)) \rmul (1+1) =
            (s_2 \rmul (s \rmul s)) \rmul 1 + (s_2 \rmul (s \rmul s)) \rmul 1$ by \cref{real_distrib}.
        Thus $(s_2 \rmul (s \rmul s)) \rmul 1 = 1 \rmul (s_2 \rmul (s \rmul s))$ by \cref{real_mul_comm}.
        Thus $1 \rmul (s_2 \rmul (s \rmul s)) = s_2 \rmul (s \rmul s)$ by \cref{real_mul_ident}.
        Thus $(1+1) \rmul (s_2 \rmul (s \rmul s)) =
            (s_2 \rmul (s \rmul s)) + (s_2 \rmul (s \rmul s))$ by \cref{real_add_eq_subst}.
        Thus $\neg{((1+1) \rmul (s_2 \rmul (s \rmul s)))} =
            \neg{(s_2 \rmul (s \rmul s) + s_2 \rmul (s \rmul s))}$ by \cref{real_add_eq_subst}.
        Thus $\neg{(s_2 \rmul (s \rmul s) + s_2 \rmul (s \rmul s))} =
            \neg{(s_2 \rmul (s \rmul s))} + \neg{(s_2 \rmul (s \rmul s))}$ by \cref{real_neg_addition}.
        Thus $s_2 \rmul (s \rmul s) + \neg{((1+1) \rmul (s_2 \rmul (s \rmul s)))} =
            s_2 \rmul (s \rmul s) + (\neg{(s_2 \rmul (s \rmul s))} + \neg{(s_2 \rmul (s \rmul s))})$ by \cref{real_add_eq_subst}.
        Thus $\neg{(s_2 \rmul (s \rmul s))}\in\reals$ by \cref{real_add_inverse}.
        Thus $s_2 \rmul (s \rmul s) + \neg{((1+1) \rmul (s_2 \rmul (s \rmul s)))} =
            (s_2 \rmul (s \rmul s) + \neg{(s_2 \rmul (s \rmul s))}) + \neg{(s_2 \rmul (s \rmul s))}$ by \cref{real_add_assoc}.
        Thus $s_2 \rmul (s \rmul s) + \neg{(s_2 \rmul (s \rmul s))} = 0$ by \cref{real_add_inverse}.
        Thus $s_2 \rmul (s \rmul s) + \neg{((1+1) \rmul (s_2 \rmul (s \rmul s)))} =
            0 + \neg{(s_2 \rmul (s \rmul s))}$ by \cref{real_add_eq_subst}.
        Thus $s_2 \rmul (s \rmul s) + \neg{((1+1) \rmul (s_2 \rmul (s \rmul s)))} =
            \neg{(s_2 \rmul (s \rmul s))}$ by \cref{real_add_ident}.
        Thus $S = s_2 \rmul (s_2 \rmul s_1) + \neg{(s_2 \rmul (s \rmul s))}$ by \cref{real_add_eq_subst}.
        Thus $s_2 \rmul s_1\in\reals$ by \cref{real_mul_closed}.
        Thus $s \rmul s\in\reals$ by \cref{real_mul_closed}.
        Thus $\neg{(s \rmul s)}\in\reals$ by \cref{real_add_inverse}.
        Thus $s_2 \rmul (s_2 \rmul s_1 + \neg{(s \rmul s)}) =
            s_2 \rmul (s_2 \rmul s_1) + s_2 \rmul \neg{(s \rmul s)}$ by \cref{real_distrib}.
        Thus $s_2 \rmul \neg{(s \rmul s)} = \neg{(s_2 \rmul (s \rmul s))}$ by \cref{real_mul_neg_right}.
        Thus $s_2 \rmul (s_2 \rmul s_1 + \neg{(s \rmul s)}) =
            s_2 \rmul (s_2 \rmul s_1) + \neg{(s_2 \rmul (s \rmul s))}$ by \cref{real_add_eq_subst}.
        Thus $S = s_2 \rmul (s_2 \rmul s_1 + \neg{(s \rmul s)})$.
        Thus $S = s_2 \rmul (s_2 \rmul s_1 - s \rmul s)$ by \cref{real_sub}.
    \end{subproof}
    Show that $s \rmul s \leq s_1 \rmul s_2$.
    \begin{subproof}
        Show that $0 \leq s_2$.
        \begin{subproof}
            Show that for all $i\in\dom{\coordSquareSequenceMul{n}{b}}$ we have $0 \leq \apply{\coordSquareSequenceMul{n}{b}}{i}$.
            \begin{subproof}
                Fix $i$.
                Assume $i\in\dom{\coordSquareSequenceMul{n}{b}}$.
                Thus $\coordSquareSequenceMul{n}{b}\in\funs{J}{\reals}$ by \cref{coord_square_sequence_mul_function}.
                Thus $i\in J$ by \cref{funs_elim}.
                Thus $(i, b(i) \rmul b(i))\in \coordSquareSequenceMul{n}{b}$ by \cref{coord_square_sequence_mul}.
                Thus $\apply{\coordSquareSequenceMul{n}{b}}{i} = b(i) \rmul b(i)$ by \cref{function_apply_intro,funs_elim}.
                Thus $b(i)\in\reals$ by \cref{funs_type_apply}.
                Thus $\exp{b(i)}{1+1} \geq 0$ by \cref{real_sq_nonnegative}.
                $1\in\naturalsPlus$ by \cref{real_one_in_naturals}.
                $\exp{b(i)}{1+1} = \exp{b(i)}{1} \rmul b(i)$ by \cref{real_pow_succ}.
                $\exp{b(i)}{1} = b(i)$ by \cref{real_pow_one}.
                Thus $\exp{b(i)}{1+1} = b(i) \rmul b(i)$.
                Thus $0 \leq b(i) \rmul b(i)$.
                Thus $0 \leq \apply{\coordSquareSequenceMul{n}{b}}{i}$.
            \end{subproof}
            Thus $0 \leq s_2$ by \cref{sum_finite_sequence_nonneg}.
        \end{subproof}
        $0 < s_2$ or $s_2 = 0$ by \cref{real_leq_split}.
        \begin{byCase}
        \caseOf{$s_2 = 0$.}
            Show that $s = 0$.
            \begin{subproof}
                Show that for all $i\in J$ we have $b(i) = 0$.
                \begin{subproof}
                    Show that for all $i\in J$ we have $0 \leq \apply{\coordSquareSequenceMul{n}{b}}{i}$.
                    \begin{subproof}
                        Fix $i$.
                        Assume $i\in J$.
                        Thus $\coordSquareSequenceMul{n}{b}\in\funs{J}{\reals}$ by \cref{coord_square_sequence_mul_function}.
                        Thus $(i, b(i) \rmul b(i))\in \coordSquareSequenceMul{n}{b}$ by \cref{coord_square_sequence_mul}.
                        Thus $\apply{\coordSquareSequenceMul{n}{b}}{i} = b(i) \rmul b(i)$ by \cref{function_apply_intro,funs_elim}.
                        Thus $b(i)\in\reals$ by \cref{funs_type_apply}.
                        Thus $\exp{b(i)}{1+1} \geq 0$ by \cref{real_sq_nonnegative}.
                        $1\in\naturalsPlus$ by \cref{real_one_in_naturals}.
                        $\exp{b(i)}{1+1} = \exp{b(i)}{1} \rmul b(i)$ by \cref{real_pow_succ}.
                        $\exp{b(i)}{1} = b(i)$ by \cref{real_pow_one}.
                        Thus $\exp{b(i)}{1+1} = b(i) \rmul b(i)$.
                        Thus $0 \leq b(i) \rmul b(i)$.
                        Thus $0 \leq \apply{\coordSquareSequenceMul{n}{b}}{i}$.
                    \end{subproof}
                    Thus $n\in\naturalsPlus$.
                    Thus $J = \initialSegment{n}$.
                    Thus $\coordSquareSequenceMul{n}{b}\in\funs{J}{\reals}$ by \cref{coord_square_sequence_mul_function}.
                    Thus $\sumFiniteSequence{\coordSquareSequenceMul{n}{b}}{s_2}$.
                    Thus $s_2 = 0$.
                    Thus for all $i\in J$ we have $\apply{\coordSquareSequenceMul{n}{b}}{i} = 0$ by \cref{sum_finite_sequence_zero_terms}.
                    Fix $i$.
                    Assume $i\in J$.
                    Thus $(i, b(i) \rmul b(i))\in \coordSquareSequenceMul{n}{b}$ by \cref{coord_square_sequence_mul}.
                    Thus $\apply{\coordSquareSequenceMul{n}{b}}{i} = b(i) \rmul b(i)$ by \cref{function_apply_intro,funs_elim}.
                    Thus $b(i) \rmul b(i) = 0$.
                    Thus $b(i)\in\reals$ by \cref{funs_type_apply}.
                    Thus $b(i) = 0$ or $b(i) = 0$ by \cref{real_mul_zero_product}.
                    Thus $b(i) = 0$.
                \end{subproof}
                Show that for all $i\in J$ we have $p(i) = 0$.
                \begin{subproof}
                    Fix $i$.
                    Assume $i\in J$.
                    Thus $p(i) = a(i) \rmul b(i)$ by \cref{coord_product_sequence,funs_elim,function_apply_intro}.
                    Thus $a(i)\in\reals$ and $b(i)\in\reals$ by \cref{funs_type_apply}.
                    Thus $a(i) \rmul b(i) = b(i) \rmul a(i)$ by \cref{real_mul_comm}.
                    Thus $b(i) \rmul a(i) = 0$ by \cref{real_mul_zero}.
                    Thus $p(i) = 0$.
                \end{subproof}
                Show that $s \leq 0$.
                \begin{subproof}
                    $0\in\reals$ by \cref{real_add_ident}.
                    Thus $\sumFiniteSequence{p}{s}$.
                    Thus $s \leq n \rmul 0$ by \cref{sum_finite_sequence_leq_n_times}.
                    $n\in\reals$ by \cref{real_naturals_subset_reals,subseteq}.
                    $n \rmul 0 = 0 \rmul n$ by \cref{real_mul_comm}.
                    $0 \rmul n = 0$ by \cref{real_mul_zero}.
                    Thus $n \rmul 0 = 0$.
                    Thus $s \leq 0$.
                \end{subproof}
                Show that $0 \leq s$.
                \begin{subproof}
                    Show that for all $i\in\dom{p}$ we have $0 \leq p(i)$.
                    \begin{subproof}
                        Fix $i$.
                        Assume $i\in\dom{p}$.
                        Thus $i\in J$ by \cref{funs_elim}.
                        Thus $p(i) = 0$.
                        Thus $0 \leq p(i)$.
                    \end{subproof}
                    Thus $0 \leq s$ by \cref{sum_finite_sequence_nonneg}.
                \end{subproof}
                Thus $s = 0$ by \cref{real_leq_antisym}.
            \end{subproof}
            Thus $s \rmul s = s \rmul 0$.
            Thus $s \rmul 0 = 0 \rmul s$ by \cref{real_mul_comm}.
            Thus $0 \rmul s = 0$ by \cref{real_mul_zero}.
            Thus $s \rmul s = 0$.
            Thus $s_1 \rmul s_2 = s_1 \rmul 0$.
            Thus $s_1 \rmul 0 = 0 \rmul s_1$ by \cref{real_mul_comm}.
            Thus $0 \rmul s_1 = 0$ by \cref{real_mul_zero}.
            Thus $s_1 \rmul s_2 = 0$.
            Thus $s \rmul s = s_1 \rmul s_2$.
            Thus $s \rmul s \leq s_1 \rmul s_2$.
        \caseOf{$0 < s_2$.}
            Show that $s_2 \rmul s_1 - s \rmul s \geq 0$.
            \begin{subproof}
                Suppose not.
                Thus $s_2 \rmul s_1 - s \rmul s < 0$.
                Thus $0\in\reals$ by \cref{real_add_ident}.
                Thus $s_2 \rmul s_1\in\reals$ by \cref{real_mul_closed}.
                Thus $s \rmul s\in\reals$ by \cref{real_mul_closed}.
                Thus $\neg{(s \rmul s)}\in\reals$ by \cref{real_add_inverse}.
                Thus $s_2 \rmul s_1 - s \rmul s\in\reals$ by \cref{real_sub,real_add_closed}.
                Thus $(s_2 \rmul s_1 - s \rmul s) \rmul s_2 < 0 \rmul s_2$ by \cref{real_lt_mul_pos}.
                Thus $0 \rmul s_2 = 0$ by \cref{real_mul_zero}.
                Thus $(s_2 \rmul s_1 - s \rmul s) \rmul s_2 < 0$ by \cref{real_lt_subst_right}.
                Thus $(s_2 \rmul s_1 - s \rmul s) \rmul s_2 = s_2 \rmul (s_2 \rmul s_1 - s \rmul s)$ by \cref{real_mul_comm}.
                Thus $s_2 \rmul (s_2 \rmul s_1 - s \rmul s) < 0$ by \cref{real_lt_subst_left}.
                Thus $S < 0$ by \cref{real_lt_subst_left}.
                Thus $0 < 0$ by \cref{real_lt_trans,real_lt_subst_right}.
                Contradiction by \cref{real_lt_irreflexive}.
            \end{subproof}
            $0\in\reals$ by \cref{real_add_ident}.
            $s_2 \rmul s_1\in\reals$ by \cref{real_mul_closed}.
            $s \rmul s\in\reals$ by \cref{real_mul_closed}.
            Thus $\neg{(s \rmul s)}\in\reals$ by \cref{real_add_inverse}.
            $s_2 \rmul s_1 - s \rmul s \in\reals$ by \cref{real_sub,real_add_closed}.
            Thus $0 + (s \rmul s) \leq (s_2 \rmul s_1 - s \rmul s) + (s \rmul s)$ by \cref{real_leq_add}.
            $0 + (s \rmul s) = s \rmul s$ by \cref{real_add_ident}.
            Thus $(s_2 \rmul s_1 - s \rmul s) + (s \rmul s) =
                (s_2 \rmul s_1 + \neg{(s \rmul s)}) + (s \rmul s)$ by \cref{real_sub}.
            Thus $(s_2 \rmul s_1 + \neg{(s \rmul s)}) + (s \rmul s) =
                s_2 \rmul s_1 + (\neg{(s \rmul s)} + (s \rmul s))$ by \cref{real_add_assoc}.
            Thus $\neg{(s \rmul s)} + (s \rmul s) = (s \rmul s) + \neg{(s \rmul s)}$ by \cref{real_add_comm}.
            Thus $s_2 \rmul s_1 + (\neg{(s \rmul s)} + (s \rmul s)) =
                s_2 \rmul s_1 + ((s \rmul s) + \neg{(s \rmul s)})$ by \cref{real_add_eq_subst}.
            Thus $(s \rmul s) + \neg{(s \rmul s)} = 0$ by \cref{real_add_inverse}.
            Thus $s_2 \rmul s_1 + ((s \rmul s) + \neg{(s \rmul s)}) =
                s_2 \rmul s_1 + 0$ by \cref{real_add_eq_subst}.
            Thus $s_2 \rmul s_1 + 0 = s_2 \rmul s_1$ by \cref{real_add_comm,real_add_ident}.
            Thus $(s_2 \rmul s_1 - s \rmul s) + (s \rmul s) = s_2 \rmul s_1$.
            Thus $s \rmul s \leq s_2 \rmul s_1$ by \cref{real_leq_trans}.
            Thus $s \rmul s \leq s_1 \rmul s_2$ by \cref{real_mul_comm}.
        \end{byCase}
    \end{subproof}
\end{proof}

% from §20 helper lemma 42: Minkowski for finite sequences
\begin{lemma}\label{minkowski_finite_sequences}
    Suppose $n\in\naturalsPlus$.
    Let $J = \initialSegment{n}$.
    Suppose $a\in\funs{J}{\reals}$.
    Suppose $b\in\funs{J}{\reals}$.
    Suppose $\sumFiniteSequence{\coordSquareSequenceMul{n}{a}}{s_1}$.
    Suppose $\sumFiniteSequence{\coordSquareSequenceMul{n}{b}}{s_2}$.
    Let $c = \coordSumSequence{n}{a}{b}$.
    Suppose $\sumFiniteSequence{\coordSquareSequenceMul{n}{c}}{s}$.
    Then $\sqrt{s} \leq \sqrt{s_1} + \sqrt{s_2}$.
\end{lemma}
\begin{proof}
    Let $p = \coordProductSequence{n}{a}{b}$.
    Thus $p\in\funs{J}{\reals}$ by \cref{coord_product_sequence_function}.
    Take $t\in\reals$ such that $\sumFiniteSequence{p}{t}$ by \cref{sum_finite_sequence_exists}.
    Thus $\sumFiniteSequence{\coordSquareSequenceMul{n}{c}}{s_1 + s_2 + (1+1) \rmul t}$ by \cref{sum_finite_sequence_square_sum}.
    Thus $s = s_1 + s_2 + (1+1) \rmul t$ by \cref{sum_finite_sequence_unique}.
    Show that $s_1\in\reals$.
    \begin{subproof}
        Take $n_1$ such that $n_1\in\naturalsPlus$ and
            $\coordSquareSequenceMul{n}{a}\in\funs{\initialSegment{n_1}}{\reals}$ and
            there exists $u\in\funs{\initialSegment{n_1}}{\reals}$ such that
                $u(1) = \apply{\coordSquareSequenceMul{n}{a}}{1}$ and
                $s_1 = u(n_1)$ and
                for all $k\in\naturalsPlus$ such that $k + 1 \leq n_1$ we have
                    $u(k + 1) = u(k) + \apply{\coordSquareSequenceMul{n}{a}}{k + 1}$
            by \cref{sum_finite_sequence_pred}.
        Take $u$ such that $u\in\funs{\initialSegment{n_1}}{\reals}$ and
            $u(1) = \apply{\coordSquareSequenceMul{n}{a}}{1}$ and
            $s_1 = u(n_1)$ and
            for all $k\in\naturalsPlus$ such that $k + 1 \leq n_1$ we have
                $u(k + 1) = u(k) + \apply{\coordSquareSequenceMul{n}{a}}{k + 1}$.
        $n_1\in\naturalsPlus$.
        Thus $n_1\leq n_1$.
        Thus $n_1\in\initialSegment{n_1}$ by \cref{initial_segment_naturalsplus}.
        Thus $u(n_1)\in\reals$ by \cref{funs_type_apply}.
        Thus $s_1\in\reals$.
    \end{subproof}
    Show that $s_2\in\reals$.
    \begin{subproof}
        Take $n_2$ such that $n_2\in\naturalsPlus$ and
            $\coordSquareSequenceMul{n}{b}\in\funs{\initialSegment{n_2}}{\reals}$ and
            there exists $v\in\funs{\initialSegment{n_2}}{\reals}$ such that
                $v(1) = \apply{\coordSquareSequenceMul{n}{b}}{1}$ and
                $s_2 = v(n_2)$ and
                for all $k\in\naturalsPlus$ such that $k + 1 \leq n_2$ we have
                    $v(k + 1) = v(k) + \apply{\coordSquareSequenceMul{n}{b}}{k + 1}$
            by \cref{sum_finite_sequence_pred}.
        Take $v$ such that $v\in\funs{\initialSegment{n_2}}{\reals}$ and
            $v(1) = \apply{\coordSquareSequenceMul{n}{b}}{1}$ and
            $s_2 = v(n_2)$ and
            for all $k\in\naturalsPlus$ such that $k + 1 \leq n_2$ we have
                $v(k + 1) = v(k) + \apply{\coordSquareSequenceMul{n}{b}}{k + 1}$.
        $n_2\in\naturalsPlus$.
        Thus $n_2\leq n_2$.
        Thus $n_2\in\initialSegment{n_2}$ by \cref{initial_segment_naturalsplus}.
        Thus $v(n_2)\in\reals$ by \cref{funs_type_apply}.
        Thus $s_2\in\reals$.
    \end{subproof}
    $1\in\naturalsPlus$ by \cref{real_one_in_naturals}.
    $1\in\reals$ by \cref{real_naturals_subset_reals,subseteq}.
    $1 + 1\in\reals$ by \cref{real_add_closed}.
    Thus $(1+1) \rmul t\in\reals$ by \cref{real_mul_closed}.
    Thus $s_1 + s_2\in\reals$ by \cref{real_add_closed}.
    Thus $s\in\reals$ by \cref{real_add_closed}.
    Show that $0 \leq s_1$.
    \begin{subproof}
        Show that for all $i\in\dom{\coordSquareSequenceMul{n}{a}}$ we have $0 \leq \apply{\coordSquareSequenceMul{n}{a}}{i}$.
        \begin{subproof}
            Fix $i$.
            Assume $i\in\dom{\coordSquareSequenceMul{n}{a}}$.
            Thus $\coordSquareSequenceMul{n}{a}\in\funs{J}{\reals}$ by \cref{coord_square_sequence_mul_function}.
            Thus $i\in J$ by \cref{funs_elim}.
            Thus $(i, a(i) \rmul a(i))\in \coordSquareSequenceMul{n}{a}$ by \cref{coord_square_sequence_mul}.
            Thus $\apply{\coordSquareSequenceMul{n}{a}}{i} = a(i) \rmul a(i)$ by \cref{function_apply_intro,funs_elim}.
            Thus $a(i)\in\reals$ by \cref{funs_type_apply}.
            Thus $\exp{a(i)}{1+1} \geq 0$ by \cref{real_sq_nonnegative}.
            $1\in\naturalsPlus$ by \cref{real_one_in_naturals}.
            $\exp{a(i)}{1+1} = \exp{a(i)}{1} \rmul a(i)$ by \cref{real_pow_succ}.
            $\exp{a(i)}{1} = a(i)$ by \cref{real_pow_one}.
            Thus $\exp{a(i)}{1+1} = a(i) \rmul a(i)$.
            Thus $0 \leq a(i) \rmul a(i)$.
            Thus $0 \leq \apply{\coordSquareSequenceMul{n}{a}}{i}$.
        \end{subproof}
        Thus $0 \leq s_1$ by \cref{sum_finite_sequence_nonneg}.
    \end{subproof}
    Show that $0 \leq s_2$.
    \begin{subproof}
        Show that for all $i\in\dom{\coordSquareSequenceMul{n}{b}}$ we have $0 \leq \apply{\coordSquareSequenceMul{n}{b}}{i}$.
        \begin{subproof}
            Fix $i$.
            Assume $i\in\dom{\coordSquareSequenceMul{n}{b}}$.
            Thus $\coordSquareSequenceMul{n}{b}\in\funs{J}{\reals}$ by \cref{coord_square_sequence_mul_function}.
            Thus $i\in J$ by \cref{funs_elim}.
            Thus $(i, b(i) \rmul b(i))\in \coordSquareSequenceMul{n}{b}$ by \cref{coord_square_sequence_mul}.
            Thus $\apply{\coordSquareSequenceMul{n}{b}}{i} = b(i) \rmul b(i)$ by \cref{function_apply_intro,funs_elim}.
            Thus $b(i)\in\reals$ by \cref{funs_type_apply}.
            Thus $\exp{b(i)}{1+1} \geq 0$ by \cref{real_sq_nonnegative}.
            $1\in\naturalsPlus$ by \cref{real_one_in_naturals}.
            $\exp{b(i)}{1+1} = \exp{b(i)}{1} \rmul b(i)$ by \cref{real_pow_succ}.
            $\exp{b(i)}{1} = b(i)$ by \cref{real_pow_one}.
            Thus $\exp{b(i)}{1+1} = b(i) \rmul b(i)$.
            Thus $0 \leq b(i) \rmul b(i)$.
            Thus $0 \leq \apply{\coordSquareSequenceMul{n}{b}}{i}$.
        \end{subproof}
        Thus $0 \leq s_2$ by \cref{sum_finite_sequence_nonneg}.
    \end{subproof}
    Show that $0 \leq s$.
    \begin{subproof}
        Show that for all $i\in\dom{\coordSquareSequenceMul{n}{c}}$ we have $0 \leq \apply{\coordSquareSequenceMul{n}{c}}{i}$.
        \begin{subproof}
            Fix $i$.
            Assume $i\in\dom{\coordSquareSequenceMul{n}{c}}$.
            $c\in\funs{J}{\reals}$ by \cref{coord_sum_sequence_function}.
            Thus $\coordSquareSequenceMul{n}{c}\in\funs{J}{\reals}$ by \cref{coord_square_sequence_mul_function}.
            Thus $i\in J$ by \cref{funs_elim}.
            Thus $(i, c(i) \rmul c(i))\in \coordSquareSequenceMul{n}{c}$ by \cref{coord_square_sequence_mul}.
            Thus $\apply{\coordSquareSequenceMul{n}{c}}{i} = c(i) \rmul c(i)$ by \cref{function_apply_intro,funs_elim}.
            Thus $c(i)\in\reals$ by \cref{funs_type_apply}.
            Thus $\exp{c(i)}{1+1} \geq 0$ by \cref{real_sq_nonnegative}.
            $1\in\naturalsPlus$ by \cref{real_one_in_naturals}.
            $\exp{c(i)}{1+1} = \exp{c(i)}{1} \rmul c(i)$ by \cref{real_pow_succ}.
            $\exp{c(i)}{1} = c(i)$ by \cref{real_pow_one}.
            Thus $\exp{c(i)}{1+1} = c(i) \rmul c(i)$.
            Thus $0 \leq c(i) \rmul c(i)$.
            Thus $0 \leq \apply{\coordSquareSequenceMul{n}{c}}{i}$.
        \end{subproof}
        Thus $0 \leq s$ by \cref{sum_finite_sequence_nonneg}.
    \end{subproof}
    Let $r = \sqrt{s}$.
    Let $r_1 = \sqrt{s_1}$.
    Let $r_2 = \sqrt{s_2}$.
    $r\in\reals$ and $r \geq 0$ and $r \rmul r = s$ by \cref{positive_square_root}.
    $r_1\in\reals$ and $r_1 \geq 0$ and $r_1 \rmul r_1 = s_1$ by \cref{positive_square_root}.
    $r_2\in\reals$ and $r_2 \geq 0$ and $r_2 \rmul r_2 = s_2$ by \cref{positive_square_root}.
    Show that $0 \leq r_1 \rmul r_2$.
    \begin{subproof}
        $0 < r_1$ or $r_1 = 0$ by \cref{real_leq_split}.
        \begin{byCase}
        \caseOf{$r_1 = 0$.}
            Thus $r_1 \rmul r_2 = 0 \rmul r_2$.
            Thus $0 \rmul r_2 = 0$ by \cref{real_mul_zero}.
            Thus $r_1 \rmul r_2 = 0$.
            Thus $r_1 \rmul r_2 \geq 0$.
        \caseOf{$0 < r_1$.}
            $0 < r_2$ or $r_2 = 0$ by \cref{real_leq_split}.
            \begin{byCase}
            \caseOf{$r_2 = 0$.}
                Thus $r_1 \rmul r_2 = r_1 \rmul 0$.
                Thus $r_1 \rmul 0 = 0 \rmul r_1$ by \cref{real_mul_comm}.
                Thus $0 \rmul r_1 = 0$ by \cref{real_mul_zero}.
                Thus $r_1 \rmul r_2 = 0$.
                Thus $r_1 \rmul r_2 \geq 0$.
            \caseOf{$0 < r_2$.}
                Thus $r_1 \rmul r_2 > 0$ by \cref{real_mul_pos_pos_gt}.
                Thus $r_1 \rmul r_2 \geq 0$.
            \end{byCase}
        \end{byCase}
    \end{subproof}
    Show that $\abs{t} \leq r_1 \rmul r_2$.
    \begin{subproof}
        Show that $\abs{t}\in\reals$.
        \begin{subproof}
            $t = 0$ or $t \neq 0$.
            \begin{byCase}
            \caseOf{$t = 0$.}
                Thus $t \geq 0$.
                Thus $\abs{t} = t$ by \cref{real_abs_nonneg}.
                Thus $\abs{t}\in\reals$.
            \caseOf{$t \neq 0$.}
                $0\in\reals$ by \cref{real_add_ident}.
                Thus $t < 0$ or $0 < t$ by \cref{real_lt_trichotomy}.
                \begin{byCase}
                \caseOf{$t < 0$.}
                    Thus $\abs{t} = \neg{t}$ by \cref{real_abs_neg}.
                    Thus $\neg{t}\in\reals$ by \cref{real_add_inverse}.
                    Thus $\abs{t}\in\reals$.
                \caseOf{$0 < t$.}
                    Thus $t \geq 0$.
                    Thus $\abs{t} = t$ by \cref{real_abs_nonneg}.
                    Thus $\abs{t}\in\reals$.
                \end{byCase}
            \end{byCase}
        \end{subproof}
        Thus $\abs{t} \geq 0$ by \cref{real_abs_nonnegative}.
        Thus $r_1 \rmul r_2\in\reals$ by \cref{real_mul_closed}.
        Thus $\abs{t} \leq r_1 \rmul r_2$ iff
            $\exp{\abs{t}}{1+1} \leq \exp{r_1 \rmul r_2}{1+1}$ by \cref{real_nonneg_leq_iff_sq_leq}.
        $\exp{\abs{t}}{1+1} = \exp{t}{1+1}$ by \cref{real_abs_square}.
        $\exp{t}{1+1} = \exp{t}{1} \rmul t$ by \cref{real_pow_succ}.
        $\exp{t}{1} = t$ by \cref{real_pow_one}.
        Thus $\exp{t}{1+1} = t \rmul t$.
        Thus $t \rmul t \leq s_1 \rmul s_2$ by \cref{cauchy_schwarz_finite_sequences}.
        $\exp{r_1 \rmul r_2}{1+1} = (r_1 \rmul r_2) \rmul (r_1 \rmul r_2)$ by \cref{real_pow_succ,real_pow_one}.
        Thus $(r_1 \rmul r_2) \rmul (r_1 \rmul r_2) =
            (r_1 \rmul r_1) \rmul (r_2 \rmul r_2)$ by \cref{real_mul_assoc,real_mul_comm}.
        Thus $(r_1 \rmul r_1) \rmul (r_2 \rmul r_2) = s_1 \rmul s_2$.
        Thus $\exp{r_1 \rmul r_2}{1+1} = s_1 \rmul s_2$.
        Thus $\exp{\abs{t}}{1+1} \leq \exp{r_1 \rmul r_2}{1+1}$ by \cref{real_leq_trans}.
        Thus $\abs{t} \leq r_1 \rmul r_2$ by \cref{real_nonneg_leq_iff_sq_leq}.
    \end{subproof}
    Thus $\abs{t}\in\reals$.
    Thus $r_1 \rmul r_2\in\reals$ by \cref{real_mul_closed}.
    Thus $t \leq \abs{t}$ by \cref{real_leq_abs}.
    Thus $t \leq r_1 \rmul r_2$ by \cref{real_leq_trans}.
    Show that $(1+1) \rmul t \leq (1+1) \rmul (r_1 \rmul r_2)$.
    \begin{subproof}
        Thus $t < r_1 \rmul r_2$ or $t = r_1 \rmul r_2$ by \cref{real_leq_split}.
        $0\in\reals$ by \cref{real_add_ident}.
        $1\in\reals$ by \cref{real_naturals_subset_reals,real_one_in_naturals,subseteq}.
        $0 < 1$ by \cref{real_zero_lt_one}.
        $0 + 1 < 1 + 1$ by \cref{real_lt_add}.
        $0 + 1 = 1$ by \cref{real_add_ident}.
        Thus $1 < 1 + 1$ by \cref{real_lt_subst_left}.
        Thus $0 < 1 + 1$ by \cref{real_lt_trans}.
        \begin{byCase}
        \caseOf{$t = r_1 \rmul r_2$.}
            Thus $(1+1) \rmul t = (1+1) \rmul (r_1 \rmul r_2)$.
        \caseOf{$t < r_1 \rmul r_2$.}
            $t\in\reals$.
            Thus $r_1 \rmul r_2\in\reals$ by \cref{real_mul_closed}.
            Thus $1 + 1\in\reals$ by \cref{real_add_closed}.
            Thus $t \rmul (1+1) < (r_1 \rmul r_2) \rmul (1+1)$ by \cref{real_lt_mul_pos}.
            Thus $t \rmul (1+1) = (1+1) \rmul t$ by \cref{real_mul_comm}.
            Thus $(1+1) \rmul t < (r_1 \rmul r_2) \rmul (1+1)$ by \cref{real_lt_subst_left}.
            Thus $(r_1 \rmul r_2) \rmul (1+1) = (1+1) \rmul (r_1 \rmul r_2)$ by \cref{real_mul_comm}.
            Thus $(1+1) \rmul t < (1+1) \rmul (r_1 \rmul r_2)$ by \cref{real_lt_subst_right}.
            Thus $(1+1) \rmul t \leq (1+1) \rmul (r_1 \rmul r_2)$.
        \end{byCase}
    \end{subproof}
    Thus $(1+1) \rmul t\in\reals$ by \cref{real_mul_closed}.
    Thus $(1+1) \rmul (r_1 \rmul r_2)\in\reals$ by \cref{real_mul_closed}.
    Thus $(1+1) \rmul t + s_1 \leq (1+1) \rmul (r_1 \rmul r_2) + s_1$ by \cref{real_leq_add}.
    Thus $(1+1) \rmul t + s_1 = s_1 + (1+1) \rmul t$ by \cref{real_add_comm}.
    Thus $(1+1) \rmul (r_1 \rmul r_2) + s_1 = s_1 + (1+1) \rmul (r_1 \rmul r_2)$ by \cref{real_add_comm}.
    Thus $s_1 + (1+1) \rmul t \leq s_1 + (1+1) \rmul (r_1 \rmul r_2)$ by \cref{real_leq_trans}.
    Thus $s_1 + (1+1) \rmul t\in\reals$ by \cref{real_add_closed}.
    Thus $s_1 + (1+1) \rmul (r_1 \rmul r_2)\in\reals$ by \cref{real_add_closed}.
    Thus $(s_1 + (1+1) \rmul t) + s_2 \leq (s_1 + (1+1) \rmul (r_1 \rmul r_2)) + s_2$ by \cref{real_leq_add}.
    Thus $(s_1 + (1+1) \rmul t) + s_2 = s_1 + ((1+1) \rmul t + s_2)$ by \cref{real_add_assoc}.
    Thus $(1+1) \rmul t + s_2 = s_2 + (1+1) \rmul t$ by \cref{real_add_comm}.
    Thus $s_1 + ((1+1) \rmul t + s_2) = s_1 + (s_2 + (1+1) \rmul t)$ by \cref{real_add_eq_subst}.
    Thus $s_1 + (s_2 + (1+1) \rmul t) = (s_1 + s_2) + (1+1) \rmul t$ by \cref{real_add_assoc}.
    Thus $(s_1 + (1+1) \rmul t) + s_2 = (s_1 + s_2) + (1+1) \rmul t$.
    Thus $(s_1 + (1+1) \rmul (r_1 \rmul r_2)) + s_2 = s_1 + ((1+1) \rmul (r_1 \rmul r_2) + s_2)$ by \cref{real_add_assoc}.
    Thus $(1+1) \rmul (r_1 \rmul r_2) + s_2 = s_2 + (1+1) \rmul (r_1 \rmul r_2)$ by \cref{real_add_comm}.
    Thus $s_1 + ((1+1) \rmul (r_1 \rmul r_2) + s_2) = s_1 + (s_2 + (1+1) \rmul (r_1 \rmul r_2))$ by \cref{real_add_eq_subst}.
    Thus $s_1 + (s_2 + (1+1) \rmul (r_1 \rmul r_2)) = (s_1 + s_2) + (1+1) \rmul (r_1 \rmul r_2)$ by \cref{real_add_assoc}.
    Thus $(s_1 + (1+1) \rmul (r_1 \rmul r_2)) + s_2 = (s_1 + s_2) + (1+1) \rmul (r_1 \rmul r_2)$.
    Thus $(s_1 + s_2) + (1+1) \rmul t \leq (s_1 + s_2) + (1+1) \rmul (r_1 \rmul r_2)$ by \cref{real_leq_trans}.
    Thus $s \leq (s_1 + s_2) + (1+1) \rmul (r_1 \rmul r_2)$ by \cref{real_leq_trans}.
    Show that $(r_1 + r_2) \rmul (r_1 + r_2) = s_1 + s_2 + (1+1) \rmul (r_1 \rmul r_2)$.
    \begin{subproof}
        Thus $r_1 + r_2\in\reals$ by \cref{real_add_closed}.
        Thus $(r_1 + r_2) \rmul (r_1 + r_2) =
            (r_1 + r_2) \rmul r_1 + (r_1 + r_2) \rmul r_2$ by \cref{real_distrib}.
        Thus $(r_1 + r_2) \rmul r_1 = r_1 \rmul (r_1 + r_2)$ by \cref{real_mul_comm}.
        Thus $r_1 \rmul (r_1 + r_2) = (r_1 \rmul r_1) + (r_1 \rmul r_2)$ by \cref{real_distrib}.
        Thus $(r_1 + r_2) \rmul r_1 = (r_1 \rmul r_1) + (r_1 \rmul r_2)$.
        Thus $(r_1 + r_2) \rmul r_2 = r_2 \rmul (r_1 + r_2)$ by \cref{real_mul_comm}.
        Thus $r_2 \rmul (r_1 + r_2) = (r_2 \rmul r_1) + (r_2 \rmul r_2)$ by \cref{real_distrib}.
        Thus $(r_1 + r_2) \rmul r_2 = (r_2 \rmul r_1) + (r_2 \rmul r_2)$.
        Thus $(r_1 + r_2) \rmul (r_1 + r_2) =
            ((r_1 \rmul r_1) + (r_1 \rmul r_2)) + ((r_2 \rmul r_1) + (r_2 \rmul r_2))$ by \cref{real_add_eq_subst}.
        Thus $r_2 \rmul r_1 = r_1 \rmul r_2$ by \cref{real_mul_comm}.
        Thus $(r_1 + r_2) \rmul (r_1 + r_2) =
            ((r_1 \rmul r_1) + (r_1 \rmul r_2)) + ((r_1 \rmul r_2) + (r_2 \rmul r_2))$ by \cref{real_add_eq_subst}.
        Thus $(r_1 \rmul r_2) + (r_2 \rmul r_2)\in\reals$ by \cref{real_add_closed}.
        Thus $(r_1 + r_2) \rmul (r_1 + r_2) =
            (r_1 \rmul r_1) + ((r_1 \rmul r_2) + ((r_1 \rmul r_2) + (r_2 \rmul r_2)))$ by \cref{real_add_assoc}.
        Thus $(r_1 \rmul r_2) + ((r_1 \rmul r_2) + (r_2 \rmul r_2)) =
            ((r_1 \rmul r_2) + (r_1 \rmul r_2)) + (r_2 \rmul r_2)$ by \cref{real_add_assoc}.
        Thus $(r_1 + r_2) \rmul (r_1 + r_2) =
            (r_1 \rmul r_1) + (((r_1 \rmul r_2) + (r_1 \rmul r_2)) + (r_2 \rmul r_2))$ by \cref{real_add_eq_subst}.
        Thus $(r_1 \rmul r_2) + (r_1 \rmul r_2)\in\reals$ by \cref{real_add_closed}.
        Thus $(r_1 + r_2) \rmul (r_1 + r_2) =
            ((r_1 \rmul r_1) + ((r_1 \rmul r_2) + (r_1 \rmul r_2))) + (r_2 \rmul r_2)$ by \cref{real_add_assoc}.
        Thus $(1 + 1) \rmul (r_1 \rmul r_2) = (r_1 \rmul r_2) \rmul (1 + 1)$ by \cref{real_mul_comm}.
        Thus $(r_1 \rmul r_2) \rmul (1 + 1) =
            (r_1 \rmul r_2) \rmul 1 + (r_1 \rmul r_2) \rmul 1$ by \cref{real_distrib}.
        Thus $(r_1 \rmul r_2) \rmul 1 = 1 \rmul (r_1 \rmul r_2)$ by \cref{real_mul_comm}.
        Thus $1 \rmul (r_1 \rmul r_2) = r_1 \rmul r_2$ by \cref{real_mul_ident}.
        Thus $(r_1 \rmul r_2) \rmul 1 = r_1 \rmul r_2$.
        Thus $(1 + 1) \rmul (r_1 \rmul r_2) =
            (r_1 \rmul r_2) + (r_1 \rmul r_2)$ by \cref{real_add_eq_subst}.
        Thus $(r_1 + r_2) \rmul (r_1 + r_2) =
            ((r_1 \rmul r_1) + ((1 + 1) \rmul (r_1 \rmul r_2))) + (r_2 \rmul r_2)$ by \cref{real_add_eq_subst}.
        Thus $(r_1 + r_2) \rmul (r_1 + r_2) =
            (r_1 \rmul r_1) + ((1 + 1) \rmul (r_1 \rmul r_2)) + (r_2 \rmul r_2)$ by \cref{real_add_assoc}.
        Thus $r_1 \rmul r_1 = s_1$ and $r_2 \rmul r_2 = s_2$.
        Thus $(r_1 + r_2) \rmul (r_1 + r_2) =
            s_1 + ((1 + 1) \rmul (r_1 \rmul r_2)) + s_2$ by \cref{real_add_eq_subst}.
        Thus $(r_1 + r_2) \rmul (r_1 + r_2) =
            s_1 + s_2 + (1 + 1) \rmul (r_1 \rmul r_2)$ by \cref{real_add_assoc,real_add_comm}.
    \end{subproof}
    Thus $r_1 + r_2\in\reals$ by \cref{real_add_closed}.
    $0\in\reals$ by \cref{real_add_ident}.
    Thus $0 \leq r_1$.
    Thus $0 + r_2 \leq r_1 + r_2$ by \cref{real_leq_add}.
    $0 + r_2 = r_2$ by \cref{real_add_ident}.
    Thus $r_2 \leq r_1 + r_2$ by \cref{real_leq_trans}.
    Thus $0 \leq r_2$.
    Thus $0 \leq r_1 + r_2$ by \cref{real_leq_trans}.
    Thus $\exp{r}{1+1} = r \rmul r$ by \cref{real_pow_succ,real_pow_one}.
    Thus $\exp{r}{1+1} = s$.
    Thus $\exp{r_1 + r_2}{1+1} = (r_1 + r_2) \rmul (r_1 + r_2)$ by \cref{real_pow_succ,real_pow_one}.
    Thus $\exp{r_1 + r_2}{1+1} = s_1 + s_2 + (1 + 1) \rmul (r_1 \rmul r_2)$ by \cref{real_add_eq_subst}.
    Thus $\exp{r}{1+1} \leq \exp{r_1 + r_2}{1+1}$ by \cref{real_leq_trans}.
    Thus $r\in\reals$ and $r \geq 0$.
    Thus $r_1 + r_2\in\reals$ and $0 \leq r_1 + r_2$.
    Thus $r \leq r_1 + r_2$ by \cref{real_nonneg_leq_iff_sq_leq}.
    Thus $\sqrt{s} \leq \sqrt{s_1} + \sqrt{s_2}$.
\end{proof}

% from §20 helper lemma 43: euclidean metric value exists
\begin{lemma}\label{euclidean_metric_value_exists}
    Suppose $n\in\naturalsPlus$.
    Let $J = \initialSegment{n}$.
    Suppose $x\in\realsPower{J}$.
    Suppose $y\in\realsPower{J}$.
    Then there exists $r\in\reals$ such that $((x,y),r)\in\euclideanMetric{n}$.
\end{lemma}
\begin{proof}
    Let $a = \coordSquareSequence{n}{(x,y)}$.
    $a\in\funs{J}{\reals}$ by \cref{coord_square_sequence_function}.
    Take $s\in\reals$ such that $\sumFiniteSequence{a}{s}$ by \cref{sum_finite_sequence_exists}.
    Show that $0 \leq s$.
    \begin{subproof}
        Show that for all $i\in\dom{a}$ we have $0 \leq a(i)$.
        \begin{subproof}
            Fix $i$.
            Assume $i\in\dom{a}$.
            Thus $i\in J$ by \cref{funs_elim}.
            $x\in\tuplePower{\reals}{J}$ and $y\in\tuplePower{\reals}{J}$ by \cref{reals_power}.
            Thus $x\in\funs{J}{\reals}$ and $y\in\funs{J}{\reals}$ by \cref{tuple_power}.
            Thus $x(i)\in\reals$ and $y(i)\in\reals$ by \cref{funs_type_apply}.
            Thus $\neg{y(i)}\in\reals$ by \cref{real_add_inverse}.
            Thus $x(i) + \neg{y(i)}\in\reals$ by \cref{real_add_closed}.
            Thus $x(i) - y(i)\in\reals$ by \cref{real_sub}.
            Thus $\exp{x(i) - y(i)}{1+1} \geq 0$ by \cref{real_sq_nonnegative}.
            $(i,\exp{x(i) - y(i)}{1+1})\in a$ by \cref{coord_square_sequence_elem}.
            $a$ is a function by \cref{funs_elim}.
            Thus $a(i) = \exp{x(i) - y(i)}{1+1}$ by \cref{function_apply_intro}.
            Thus $0 \leq a(i)$.
        \end{subproof}
        Thus $0 \leq s$ by \cref{sum_finite_sequence_nonneg}.
    \end{subproof}
    Thus $\sqrt{s}\in\reals$ and $\sqrt{s} \geq 0$ and $\sqrt{s} \rmul \sqrt{s} = s$ by \cref{positive_square_root}.
    Let $r = \sqrt{s}$.
    Show that $((x,y),r)\in\euclideanMetric{n}$.
    \begin{subproof}
        Let $p = (x,y)$.
        $p\in\realsPower{J}\times\realsPower{J}$ by \cref{times_tuple_intro}.
        $r\in\reals$.
        Thus $((x,y),r)\in\euclideanMetric{n}$ by \cref{euclidean_metric}.
    \end{subproof}
    Thus there exists $r\in\reals$ such that $((x,y),r)\in\euclideanMetric{n}$.
\end{proof}

% from §20 helper lemma 44: euclidean metric is a function
\begin{lemma}\label{euclidean_metric_function}
    Suppose $n\in\naturalsPlus$.
    Then $\euclideanMetric{n}$ is a function.
\end{lemma}
\begin{proof}
    Let $d = \euclideanMetric{n}$.
    Show that for all $w$ such that $w\in d$ there exist $p,r$ such that $w = (p,r)$.
    \begin{subproof}
        Fix $w$.
        Assume $w\in d$.
        Take $p,r,s$ such that
            $p\in\realsPower{\initialSegment{n}}\times\realsPower{\initialSegment{n}}$ and
            $r\in\reals$ and $s\in\reals$ and
            $w = (p,r)$ and
            $\sumFiniteSequence{\coordSquareSequence{n}{p}}{s}$ and
            $r = \sqrt{s}$
            by \cref{euclidean_metric}.
        Thus there exist $p_0,r_0$ such that $w = (p_0,r_0)$.
    \end{subproof}
    Thus $d$ is a relation by \cref{relation}.
    Show that for all $p,r_1,r_2$ such that $(p,r_1)\in d$ and $(p,r_2)\in d$ we have $r_1 = r_2$.
    \begin{subproof}
        Fix $p$.
        Fix $r_1$.
        Fix $r_2$.
        Assume $(p,r_1)\in d$ and $(p,r_2)\in d$.
        Take $s_1\in\reals$ such that
            $\sumFiniteSequence{\coordSquareSequence{n}{p}}{s_1}$ and $r_1 = \sqrt{s_1}$
            by \cref{euclidean_metric,pair_eq_iff}.
        Take $s_2\in\reals$ such that
            $\sumFiniteSequence{\coordSquareSequence{n}{p}}{s_2}$ and $r_2 = \sqrt{s_2}$
            by \cref{euclidean_metric,pair_eq_iff}.
        Thus $s_1 = s_2$ by \cref{sum_finite_sequence_unique}.
        Thus $r_1 = r_2$.
    \end{subproof}
    Thus $d$ is right-unique by \cref{rightunique}.
    Thus $d$ is a function.
\end{proof}

% from §20 helper lemma 45: square metric is a metric
\begin{lemma}\label{square_metric_is_metric}
    Suppose $n\in\naturalsPlus$.
    Let $J = \initialSegment{n}$.
    Let $X = \realsPower{J}$.
    Then $\squareMetric{n}$ is a metric on $X$.
\end{lemma}
\begin{proof}
    Let $d = \squareMetric{n}$.
    Show that $d\in\funs{X\times X}{\reals}$.
    \begin{subproof}
        $d$ is a function by \cref{square_metric_function}.
        Show that $\dom{d} = X\times X$.
        \begin{subproof}
            Show that $X\times X\subseteq\dom{d}$.
            \begin{subproof}
                Show that for all $p$ such that $p\in X\times X$ we have $p\in\dom{d}$.
                \begin{subproof}
                    Fix $p$.
                    Assume $p\in X\times X$.
                    Take $x,y$ such that $x\in X$ and $y\in X$ and $p = (x,y)$ by \cref{times_elem_is_tuple}.
                    Take $r\in\reals$ such that $((x,y),r)\in d$ and
                        for all $i\in J$ we have $\abs{x(i) - y(i)} \leq r$
                        by \cref{square_metric_value_exists}.
                    Thus $(p,r)\in d$ by \cref{pair_eq_iff}.
                    Thus $p\in\dom{d}$ by \cref{dom_intro}.
                \end{subproof}
                Thus $X\times X\subseteq\dom{d}$ by \cref{subseteq}.
            \end{subproof}
            Show that $\dom{d}\subseteq X\times X$.
            \begin{subproof}
                Show that for all $p$ such that $p\in\dom{d}$ we have $p\in X\times X$.
                \begin{subproof}
                    Fix $p$.
                    Assume $p\in\dom{d}$.
                    Take $r$ such that $(p,r)\in d$ by \cref{dom_iff}.
                    Take $x,y\in\realsPower{J}$ such that
                        $p = (x,y)$ and
                        $r\in\coordDiffSet{n}{x}{y}$ and
                        for all $t\in\coordDiffSet{n}{x}{y}$ we have $t \leq r$
                        by \cref{square_metric,pair_eq_iff}.
                    Thus $p\in X\times X$ by \cref{times_tuple_intro}.
                \end{subproof}
                Thus $\dom{d}\subseteq X\times X$ by \cref{subseteq}.
            \end{subproof}
            Thus $\dom{d} = X\times X$ by \cref{subseteq_antisymmetric}.
        \end{subproof}
        Show that for all $p\in\dom{d}$ we have $d(p)\in\reals$.
        \begin{subproof}
            Fix $p$.
            Assume $p\in\dom{d}$.
            Take $r$ such that $(p,r)\in d$ by \cref{dom_iff}.
            Take $x,y\in\realsPower{J}$ such that
                $p = (x,y)$ and
                $r\in\coordDiffSet{n}{x}{y}$ and
                for all $t\in\coordDiffSet{n}{x}{y}$ we have $t \leq r$
                by \cref{square_metric,pair_eq_iff}.
            $r\in\reals$ by \cref{coord_diff_set_finite,subseteq}.
            Thus $d(p) = r$ by \cref{function_apply_intro}.
            Thus $d(p)\in\reals$.
        \end{subproof}
        Thus $d\in\funs{X\times X}{\reals}$ by \cref{funs_intro}.
    \end{subproof}
    Thus $d$ is a function by \cref{funs_elim}.
    Show that for all $x,y$ such that $x\in X$ and $y\in X$ we have $d((x,y)) \geq 0$.
    \begin{subproof}
        Fix $x$.
        Fix $y$.
        Assume $x\in X$ and $y\in X$.
        Take $r\in\reals$ such that $((x,y),r)\in d$ and
            for all $i\in J$ we have $\abs{x(i) - y(i)} \leq r$
            by \cref{square_metric_value_exists}.
        Thus $r\in\coordDiffSet{n}{x}{y}$ by \cref{square_metric,pair_eq_iff}.
        Take $i_0\in J$ such that $r = \abs{x(i_0) - y(i_0)}$ by \cref{coord_diff_set}.
        $x\in\tuplePower{\reals}{J}$ and $y\in\tuplePower{\reals}{J}$ by \cref{reals_power}.
        Thus $x\in\funs{J}{\reals}$ and $y\in\funs{J}{\reals}$ by \cref{tuple_power}.
        Thus $x(i_0)\in\reals$ and $y(i_0)\in\reals$ by \cref{funs_type_apply}.
        Thus $x(i_0) - y(i_0)\in\reals$ by \cref{real_sub,real_add_closed,real_add_inverse}.
        Thus $r \geq 0$ by \cref{real_abs_nonnegative}.
        Thus $d((x,y)) = r$ by \cref{function_apply_intro}.
        Thus $d((x,y)) \geq 0$.
    \end{subproof}
    Thus $d$ is nonnegative on $X$ by \cref{metric_nonnegative}.
    Show that $d$ is zero-characterized on $X$.
    \begin{subproof}
        Show that for all $x,y$ such that $x\in X$ and $y\in X$ we have $d((x,y)) = 0$ iff $x = y$.
        \begin{subproof}
            Fix $x$.
            Fix $y$.
            Assume $x\in X$ and $y\in X$.
            Show that if $d((x,y)) = 0$ then $x = y$.
            \begin{subproof}
                Assume $d((x,y)) = 0$.
                Take $r\in\reals$ such that $((x,y),r)\in d$ and
                    for all $i\in J$ we have $\abs{x(i) - y(i)} \leq r$
                    by \cref{square_metric_value_exists}.
                Thus $d((x,y)) = r$ by \cref{function_apply_intro}.
                Thus $r = 0$.
                Show that for all $i\in J$ we have $x(i) = y(i)$.
                \begin{subproof}
                    Fix $i$.
                    Assume $i\in J$.
                    Thus $\abs{x(i) - y(i)} \leq 0$.
                    $x\in\tuplePower{\reals}{J}$ and $y\in\tuplePower{\reals}{J}$ by \cref{reals_power}.
                    Thus $x\in\funs{J}{\reals}$ and $y\in\funs{J}{\reals}$ by \cref{tuple_power}.
                    Thus $x(i)\in\reals$ and $y(i)\in\reals$ by \cref{funs_type_apply}.
                    Thus $x(i) - y(i)\in\reals$ by \cref{real_sub,real_add_closed,real_add_inverse}.
                    Thus $\abs{x(i) - y(i)}\in\reals$ by \cref{real_abs_in_reals}.
                    $\abs{x(i) - y(i)} \geq 0$ by \cref{real_abs_nonnegative}.
                    Thus $\abs{x(i) - y(i)} = 0$ by \cref{real_leq_antisym}.
                    Thus $x(i) - y(i) = 0$ by \cref{real_abs_eq_zero}.
                    Show that $x(i) = y(i)$.
                    \begin{subproof}
                        $x(i) - y(i) = x(i) + \neg{y(i)}$ by \cref{real_sub}.
                        Thus $x(i) + \neg{y(i)} = 0$.
                        $\neg{y(i)}\in\reals$ by \cref{real_add_inverse}.
                        $(x(i) + \neg{y(i)}) + y(i) = 0 + y(i)$ by \cref{real_add_eq_subst}.
                        $(x(i) + \neg{y(i)}) + y(i) = x(i) + (\neg{y(i)} + y(i))$ by \cref{real_add_assoc}.
                        $\neg{y(i)} + y(i) = y(i) + \neg{y(i)}$ by \cref{real_add_comm}.
                        $y(i) + \neg{y(i)} = 0$ by \cref{real_add_inverse}.
                        Thus $x(i) + (\neg{y(i)} + y(i)) = x(i) + 0$ by \cref{real_add_eq_subst}.
                        $x(i)\in\reals$ and $y(i)\in\reals$.
                        $x(i) + 0 = 0 + x(i)$ by \cref{real_add_comm}.
                        $0 + x(i) = x(i)$ by \cref{real_add_ident}.
                        $0 + y(i) = y(i)$ by \cref{real_add_ident}.
                        Thus $x(i) = y(i)$.
                    \end{subproof}
                \end{subproof}
                $x\in\funs{J}{\reals}$ and $y\in\funs{J}{\reals}$ by \cref{reals_power,tuple_power}.
                Thus $\dom{x} = J$ and $\dom{y} = J$ by \cref{funs_elim}.
                Show that $x\subseteq y$.
                \begin{subproof}
                    $x$ is a function and $y$ is a function by \cref{funs_elim}.
                    Thus $\dom{x}\subseteq\dom{y}$ by \cref{subseteq}.
                    Show that for all $i\in\dom{x}$ we have $x(i) = y(i)$.
                    \begin{subproof}
                        Fix $i$.
                        Assume $i\in\dom{x}$.
                        Thus $i\in J$ by \cref{subseteq}.
                        Thus $x(i) = y(i)$.
                    \end{subproof}
                    Thus $x\subseteq y$ by \cref{fun_subseteq}.
                \end{subproof}
                Show that $y\subseteq x$.
                \begin{subproof}
                    $x$ is a function and $y$ is a function by \cref{funs_elim}.
                    Thus $\dom{y}\subseteq\dom{x}$ by \cref{subseteq}.
                    Show that for all $i\in\dom{y}$ we have $y(i) = x(i)$.
                    \begin{subproof}
                        Fix $i$.
                        Assume $i\in\dom{y}$.
                        Thus $i\in J$ by \cref{subseteq}.
                        Thus $y(i) = x(i)$.
                    \end{subproof}
                    Thus $y\subseteq x$ by \cref{fun_subseteq}.
                \end{subproof}
                Thus $x = y$ by \cref{subseteq_antisymmetric}.
            \end{subproof}
            Show that if $x = y$ then $d((x,y)) = 0$.
            \begin{subproof}
                Assume $x = y$.
                Take $r\in\reals$ such that $((x,y),r)\in d$ and
                    for all $i\in J$ we have $\abs{x(i) - y(i)} \leq r$
                    by \cref{square_metric_value_exists}.
                Show that for all $t\in\coordDiffSet{n}{x}{y}$ we have $t = 0$.
                \begin{subproof}
                    Fix $t$.
                    Assume $t\in\coordDiffSet{n}{x}{y}$.
                    Take $i\in J$ such that $t = \abs{x(i) - y(i)}$ by \cref{coord_diff_set}.
                    Thus $t = \abs{x(i) - x(i)}$.
                    $x(i)\in\reals$ by \cref{reals_power,tuple_power,funs_type_apply}.
                    $x(i) - x(i) = x(i) + \neg{x(i)}$ by \cref{real_sub}.
                    $\neg{x(i)}\in\reals$ by \cref{real_add_inverse}.
                    Thus $x(i) + \neg{x(i)} = 0$ by \cref{real_add_inverse}.
                    Thus $x(i) - x(i) = 0$ by \cref{real_add_eq_subst}.
                    Thus $0 < x(i) - x(i)$ or $0 = x(i) - x(i)$.
                    Thus $0 \leq x(i) - x(i)$.
                    Thus $x(i) - x(i) \geq 0$.
                    $x(i) - x(i)\in\reals$ by \cref{real_sub,real_add_closed}.
                    Thus $\abs{x(i) - x(i)} = x(i) - x(i)$ by \cref{real_abs_nonneg}.
                    Thus $t = 0$.
                \end{subproof}
                Thus $r\in\coordDiffSet{n}{x}{y}$ by \cref{square_metric,pair_eq_iff}.
                Thus $r = 0$.
                Thus $d((x,y)) = r$ by \cref{function_apply_intro}.
                Thus $d((x,y)) = 0$.
            \end{subproof}
        \end{subproof}
        Thus $d$ is zero-characterized on $X$ by \cref{metric_zero_characterization}.
    \end{subproof}
    Show that $d$ is symmetric on $X$.
    \begin{subproof}
        Show that for all $x,y$ such that $x\in X$ and $y\in X$ we have $d((x,y)) = d((y,x))$.
        \begin{subproof}
            Fix $x$.
            Fix $y$.
            Assume $x\in X$ and $y\in X$.
            Take $r\in\reals$ such that $((x,y),r)\in d$ and
                for all $i\in J$ we have $\abs{x(i) - y(i)} \leq r$
                by \cref{square_metric_value_exists}.
            Take $s\in\reals$ such that $((y,x),s)\in d$ and
                for all $i\in J$ we have $\abs{y(i) - x(i)} \leq s$
                by \cref{square_metric_value_exists}.
            Show that $r \leq s$.
            \begin{subproof}
                Show that for all $t\in\coordDiffSet{n}{x}{y}$ we have $t \leq s$.
                \begin{subproof}
                    Fix $t$.
                    Assume $t\in\coordDiffSet{n}{x}{y}$.
                    Take $i\in J$ such that $t = \abs{x(i) - y(i)}$ by \cref{coord_diff_set}.
                    $x(i)\in\reals$ by \cref{reals_power,tuple_power,funs_type_apply}.
                    $y(i)\in\reals$ by \cref{reals_power,tuple_power,funs_type_apply}.
                    Thus $t = \abs{y(i) - x(i)}$ by \cref{real_abs_sub_swap}.
                    Thus $t \leq s$.
                \end{subproof}
                Thus $r\in\coordDiffSet{n}{x}{y}$ by \cref{square_metric,pair_eq_iff}.
                Thus $r \leq s$.
            \end{subproof}
            Show that $s \leq r$.
            \begin{subproof}
                Show that for all $t\in\coordDiffSet{n}{y}{x}$ we have $t \leq r$.
                \begin{subproof}
                    Fix $t$.
                    Assume $t\in\coordDiffSet{n}{y}{x}$.
                    Take $i\in J$ such that $t = \abs{y(i) - x(i)}$ by \cref{coord_diff_set}.
                    $x(i)\in\reals$ by \cref{reals_power,tuple_power,funs_type_apply}.
                    $y(i)\in\reals$ by \cref{reals_power,tuple_power,funs_type_apply}.
                    Thus $t = \abs{x(i) - y(i)}$ by \cref{real_abs_sub_swap}.
                    Thus $t \leq r$.
                \end{subproof}
                Thus $s\in\coordDiffSet{n}{y}{x}$ by \cref{square_metric,pair_eq_iff}.
                Thus $s \leq r$.
            \end{subproof}
            Thus $r = s$ by \cref{real_leq_antisym}.
            Thus $d((x,y)) = r$ by \cref{function_apply_intro}.
            Thus $d((y,x)) = s$ by \cref{function_apply_intro}.
            Thus $d((x,y)) = d((y,x))$.
        \end{subproof}
        Thus $d$ is symmetric on $X$ by \cref{metric_symmetric}.
    \end{subproof}
    Show that $d$ satisfies the triangle inequality on $X$.
    \begin{subproof}
        Show that for all $x,y,z$ such that $x\in X$ and $y\in X$ and $z\in X$ we have
            $d((x,y)) + d((y,z)) \geq d((x,z))$.
        \begin{subproof}
            Fix $x$.
            Fix $y$.
            Fix $z$.
            Assume $x\in X$ and $y\in X$ and $z\in X$.
            Take $r_1\in\reals$ such that $((x,y),r_1)\in d$ and
                for all $i\in J$ we have $\abs{x(i) - y(i)} \leq r_1$
                by \cref{square_metric_value_exists}.
            Take $r_2\in\reals$ such that $((y,z),r_2)\in d$ and
                for all $i\in J$ we have $\abs{y(i) - z(i)} \leq r_2$
                by \cref{square_metric_value_exists}.
            Take $r_3\in\reals$ such that $((x,z),r_3)\in d$ and
                for all $i\in J$ we have $\abs{x(i) - z(i)} \leq r_3$
                by \cref{square_metric_value_exists}.
            Thus $r_3\in\coordDiffSet{n}{x}{z}$ by \cref{square_metric,pair_eq_iff}.
            Take $i_0\in J$ such that $r_3 = \abs{x(i_0) - z(i_0)}$ by \cref{coord_diff_set}.
            $x(i_0)\in\reals$ by \cref{reals_power,tuple_power,funs_type_apply}.
            $y(i_0)\in\reals$ by \cref{reals_power,tuple_power,funs_type_apply}.
            $z(i_0)\in\reals$ by \cref{reals_power,tuple_power,funs_type_apply}.
            $\neg{y(i_0)}\in\reals$ by \cref{real_add_inverse}.
            $\neg{z(i_0)}\in\reals$ by \cref{real_add_inverse}.
            $x(i_0) - y(i_0)\in\reals$ by \cref{real_sub,real_add_closed}.
            $y(i_0) - z(i_0)\in\reals$ by \cref{real_sub,real_add_closed}.
            $x(i_0) - y(i_0) = x(i_0) + \neg{y(i_0)}$ by \cref{real_sub}.
            $y(i_0) - z(i_0) = y(i_0) + \neg{z(i_0)}$ by \cref{real_sub}.
            $(x(i_0) - y(i_0)) + (y(i_0) - z(i_0)) = (x(i_0) + \neg{y(i_0)}) + (y(i_0) + \neg{z(i_0)})$.
            $(x(i_0) + \neg{y(i_0)}) + (y(i_0) + \neg{z(i_0)}) = ((x(i_0) + \neg{y(i_0)}) + y(i_0)) + \neg{z(i_0)}$ by \cref{real_add_assoc}.
            $(x(i_0) + \neg{y(i_0)}) + y(i_0) = x(i_0) + (\neg{y(i_0)} + y(i_0))$ by \cref{real_add_assoc}.
            $\neg{y(i_0)} + y(i_0) = 0$ by \cref{real_add_inverse,real_add_comm}.
            Thus $x(i_0) + (\neg{y(i_0)} + y(i_0)) = x(i_0) + 0$.
            $0\in\reals$ by \cref{real_add_ident}.
            $x(i_0) + 0 = 0 + x(i_0)$ by \cref{real_add_comm}.
            $0 + x(i_0) = x(i_0)$ by \cref{real_add_ident}.
            Thus $x(i_0) + 0 = x(i_0)$.
            Thus $(x(i_0) + \neg{y(i_0)}) + y(i_0) = x(i_0)$.
            Thus $(x(i_0) + \neg{y(i_0)}) + (y(i_0) + \neg{z(i_0)}) = x(i_0) + \neg{z(i_0)}$ by \cref{real_add_assoc}.
            Thus $(x(i_0) - y(i_0)) + (y(i_0) - z(i_0)) = x(i_0) + \neg{z(i_0)}$.
            $x(i_0) + \neg{z(i_0)} = x(i_0) - z(i_0)$ by \cref{real_sub}.
            Thus $(x(i_0) - y(i_0)) + (y(i_0) - z(i_0)) = x(i_0) - z(i_0)$.
            Thus $\abs{(x(i_0) - y(i_0)) + (y(i_0) - z(i_0))} \leq \abs{x(i_0) - y(i_0)} + \abs{y(i_0) - z(i_0)}$ by \cref{real_triangle_inequality}.
            Thus $\abs{x(i_0) - z(i_0)} \leq \abs{x(i_0) - y(i_0)} + \abs{y(i_0) - z(i_0)}$.
            Thus $r_3 \leq \abs{x(i_0) - y(i_0)} + \abs{y(i_0) - z(i_0)}$.
            Thus $\abs{x(i_0) - y(i_0)} \leq r_1$.
            $\abs{x(i_0) - y(i_0)}\in\reals$ by \cref{real_abs_in_reals}.
            $\abs{y(i_0) - z(i_0)}\in\reals$ by \cref{real_abs_in_reals}.
            Thus $\abs{x(i_0) - y(i_0)} + \abs{y(i_0) - z(i_0)} \leq r_1 + \abs{y(i_0) - z(i_0)}$ by \cref{real_leq_add}.
            $\abs{x(i_0) - y(i_0)} + \abs{y(i_0) - z(i_0)}\in\reals$ by \cref{real_add_closed}.
            $r_1 + \abs{y(i_0) - z(i_0)}\in\reals$ by \cref{real_add_closed}.
            Thus $r_3 \leq r_1 + \abs{y(i_0) - z(i_0)}$ by \cref{real_leq_trans}.
            Thus $\abs{y(i_0) - z(i_0)} \leq r_2$.
            Thus $r_1\in\reals$.
            Thus $r_2\in\reals$.
            Thus $\abs{y(i_0) - z(i_0)} + r_1 \leq r_2 + r_1$ by \cref{real_leq_add}.
            $r_1 + \abs{y(i_0) - z(i_0)} = \abs{y(i_0) - z(i_0)} + r_1$ by \cref{real_add_comm}.
            $r_1 + r_2 = r_2 + r_1$ by \cref{real_add_comm}.
            Thus $r_1 + \abs{y(i_0) - z(i_0)} \leq r_1 + r_2$.
            $r_1 + r_2\in\reals$ by \cref{real_add_closed}.
            Thus $r_3 \leq r_1 + r_2$ by \cref{real_leq_trans}.
            Thus $d((x,y)) = r_1$ by \cref{function_apply_intro}.
            Thus $d((y,z)) = r_2$ by \cref{function_apply_intro}.
            Thus $d((x,z)) = r_3$ by \cref{function_apply_intro}.
            Thus $d((x,z)) \leq d((x,y)) + d((y,z))$.
            Thus $d((x,y)) + d((y,z)) \geq d((x,z))$.
        \end{subproof}
    Thus $d$ satisfies the triangle inequality on $X$ by \cref{metric_triangle_inequality}.
    \end{subproof}
    Thus $d$ is a metric on $X$ by \cref{metric_on}.
\end{proof}

% from §20 helper lemma 46: square sequence of a difference sequence
\begin{lemma}\label{coord_square_sequence_diff}
    Suppose $n\in\naturalsPlus$.
    Let $J = \initialSegment{n}$.
    Suppose $x\in\realsPower{J}$.
    Suppose $y\in\realsPower{J}$.
    Let $a = \coordDiffSequence{n}{x}{y}$.
    Then $\coordSquareSequenceMul{n}{a} = \coordSquareSequence{n}{(x,y)}$.
\end{lemma}
\begin{proof}
    Let $u = \coordSquareSequenceMul{n}{a}$.
    Let $v = \coordSquareSequence{n}{(x,y)}$.
    Show that $u\subseteq v$.
    \begin{subproof}
        Show that for all $w$ such that $w\in u$ we have $w\in v$.
        \begin{subproof}
            Fix $w$.
            Assume $w\in u$.
            Take $i\in J$ such that $w = (i, a(i) \rmul a(i))$ by \cref{coord_square_sequence_mul}.
            Thus $i\in\initialSegment{n}$.
            $a\in\funs{J}{\reals}$ by \cref{coord_diff_sequence_function}.
            Thus $a$ is a function by \cref{funs_elim}.
            Thus $(i, x(i) - y(i))\in a$ by \cref{coord_diff_sequence}.
            Thus $a(i) = x(i) - y(i)$ by \cref{function_apply_intro}.
            $x\in\tuplePower{\reals}{J}$ and $y\in\tuplePower{\reals}{J}$ by \cref{reals_power}.
            Thus $x\in\funs{J}{\reals}$ and $y\in\funs{J}{\reals}$ by \cref{tuple_power}.
            Thus $x(i)\in\reals$ and $y(i)\in\reals$ by \cref{funs_type_apply}.
            Thus $\neg{y(i)}\in\reals$ by \cref{real_add_inverse}.
            Thus $x(i) + \neg{y(i)}\in\reals$ by \cref{real_add_closed}.
            Thus $x(i) - y(i)\in\reals$ by \cref{real_sub}.
            $1\in\naturalsPlus$ by \cref{real_one_in_naturals}.
            Thus $\exp{x(i) - y(i)}{1+1} = (x(i) - y(i)) \rmul (x(i) - y(i))$ by \cref{real_pow_succ,real_pow_one}.
            Thus $w = (i,\exp{x(i) - y(i)}{1+1})$.
            Let $p = (x,y)$.
            $\fst{p} = x$ and $\snd{p} = y$ by \cref{fst_eq,snd_eq}.
            Thus $\exp{x(i) - y(i)}{1+1} = \exp{\apply{\fst{p}}{i} - \apply{\snd{p}}{i}}{1+1}$.
            Thus $w = (i,\exp{\apply{\fst{p}}{i} - \apply{\snd{p}}{i}}{1+1})$.
            Thus $w\in v$ by \cref{coord_square_sequence}.
        \end{subproof}
        Thus $u\subseteq v$ by \cref{subseteq}.
    \end{subproof}
    Show that $v\subseteq u$.
    \begin{subproof}
        Show that for all $w$ such that $w\in v$ we have $w\in u$.
        \begin{subproof}
            Fix $w$.
            Assume $w\in v$.
            Let $p = (x,y)$.
            Take $i\in\initialSegment{n}$ such that
                $w = (i,\exp{\apply{\fst{p}}{i} - \apply{\snd{p}}{i}}{1+1})$ by \cref{coord_square_sequence}.
            Thus $i\in J$.
            $\fst{p} = x$ and $\snd{p} = y$ by \cref{fst_eq,snd_eq}.
            Thus $w = (i,\exp{x(i) - y(i)}{1+1})$.
            $a\in\funs{J}{\reals}$ by \cref{coord_diff_sequence_function}.
            Thus $a$ is a function by \cref{funs_elim}.
            Thus $(i, x(i) - y(i))\in a$ by \cref{coord_diff_sequence}.
            Thus $a(i) = x(i) - y(i)$ by \cref{function_apply_intro}.
            $x\in\tuplePower{\reals}{J}$ and $y\in\tuplePower{\reals}{J}$ by \cref{reals_power}.
            Thus $x\in\funs{J}{\reals}$ and $y\in\funs{J}{\reals}$ by \cref{tuple_power}.
            Thus $x(i)\in\reals$ and $y(i)\in\reals$ by \cref{funs_type_apply}.
            Thus $\neg{y(i)}\in\reals$ by \cref{real_add_inverse}.
            Thus $x(i) + \neg{y(i)}\in\reals$ by \cref{real_add_closed}.
            Thus $x(i) - y(i)\in\reals$ by \cref{real_sub}.
            $1\in\naturalsPlus$ by \cref{real_one_in_naturals}.
            Thus $\exp{x(i) - y(i)}{1+1} = (x(i) - y(i)) \rmul (x(i) - y(i))$ by \cref{real_pow_succ,real_pow_one}.
            Thus $w = (i, a(i) \rmul a(i))$.
            Thus $w\in u$ by \cref{coord_square_sequence_mul}.
        \end{subproof}
        Thus $v\subseteq u$ by \cref{subseteq}.
    \end{subproof}
    Thus $u = v$ by \cref{subseteq_antisymmetric}.
\end{proof}

% from §20 helper lemma 47: difference sequence of a sum
\begin{lemma}\label{coord_diff_sequence_sum}
    Suppose $n\in\naturalsPlus$.
    Let $J = \initialSegment{n}$.
    Suppose $x\in\realsPower{J}$.
    Suppose $y\in\realsPower{J}$.
    Suppose $z\in\realsPower{J}$.
    Let $a = \coordDiffSequence{n}{x}{y}$.
    Let $b = \coordDiffSequence{n}{y}{z}$.
    Let $c = \coordSumSequence{n}{a}{b}$.
    Then $c = \coordDiffSequence{n}{x}{z}$.
\end{lemma}
\begin{proof}
    Let $d = \coordDiffSequence{n}{x}{z}$.
    Show that $c\subseteq d$.
    \begin{subproof}
        Show that for all $w$ such that $w\in c$ we have $w\in d$.
        \begin{subproof}
            Fix $w$.
            Assume $w\in c$.
            Take $i\in J$ such that $w = (i, a(i) + b(i))$ by \cref{coord_sum_sequence}.
            $a\in\funs{J}{\reals}$ by \cref{coord_diff_sequence_function}.
            $b\in\funs{J}{\reals}$ by \cref{coord_diff_sequence_function}.
            Thus $a$ is a function by \cref{funs_elim}.
            Thus $b$ is a function by \cref{funs_elim}.
            Thus $(i, x(i) - y(i))\in a$ by \cref{coord_diff_sequence}.
            Thus $(i, y(i) - z(i))\in b$ by \cref{coord_diff_sequence}.
            Thus $a(i) = x(i) - y(i)$ by \cref{function_apply_intro}.
            Thus $b(i) = y(i) - z(i)$ by \cref{function_apply_intro}.
            $x\in\tuplePower{\reals}{J}$ and $y\in\tuplePower{\reals}{J}$ by \cref{reals_power}.
            Thus $x\in\funs{J}{\reals}$ and $y\in\funs{J}{\reals}$ by \cref{tuple_power}.
            Thus $x(i)\in\reals$ and $y(i)\in\reals$ by \cref{funs_type_apply}.
            $z\in\tuplePower{\reals}{J}$ by \cref{reals_power}.
            Thus $z\in\funs{J}{\reals}$ by \cref{tuple_power}.
            Thus $z(i)\in\reals$ by \cref{funs_type_apply}.
            $\neg{y(i)}\in\reals$ by \cref{real_add_inverse}.
            $\neg{z(i)}\in\reals$ by \cref{real_add_inverse}.
            $x(i) + \neg{y(i)}\in\reals$ by \cref{real_add_closed}.
            $y(i) + \neg{z(i)}\in\reals$ by \cref{real_add_closed}.
            $x(i) - y(i) = x(i) + \neg{y(i)}$ by \cref{real_sub}.
            $y(i) - z(i) = y(i) + \neg{z(i)}$ by \cref{real_sub}.
            $(x(i) - y(i)) + (y(i) - z(i)) = (x(i) + \neg{y(i)}) + (y(i) + \neg{z(i)})$.
            $(x(i) + \neg{y(i)}) + (y(i) + \neg{z(i)}) = ((x(i) + \neg{y(i)}) + y(i)) + \neg{z(i)}$ by \cref{real_add_assoc}.
            $(x(i) + \neg{y(i)}) + y(i) = x(i) + (\neg{y(i)} + y(i))$ by \cref{real_add_assoc}.
            $\neg{y(i)} + y(i) = 0$ by \cref{real_add_inverse,real_add_comm}.
            Thus $x(i) + (\neg{y(i)} + y(i)) = x(i) + 0$.
            $0\in\reals$ by \cref{real_add_ident}.
            $x(i) + 0 = 0 + x(i)$ by \cref{real_add_comm}.
            $0 + x(i) = x(i)$ by \cref{real_add_ident}.
            Thus $x(i) + 0 = x(i)$.
            Thus $(x(i) + \neg{y(i)}) + y(i) = x(i)$.
            Thus $(x(i) + \neg{y(i)}) + (y(i) + \neg{z(i)}) = x(i) + \neg{z(i)}$ by \cref{real_add_assoc}.
            Thus $(x(i) - y(i)) + (y(i) - z(i)) = x(i) + \neg{z(i)}$.
            $x(i) + \neg{z(i)} = x(i) - z(i)$ by \cref{real_sub}.
            Thus $a(i) + b(i) = x(i) - z(i)$.
            Thus $w = (i, x(i) - z(i))$.
            Thus $w\in d$ by \cref{coord_diff_sequence}.
        \end{subproof}
        Thus $c\subseteq d$ by \cref{subseteq}.
    \end{subproof}
    Show that $d\subseteq c$.
    \begin{subproof}
        Show that for all $w$ such that $w\in d$ we have $w\in c$.
        \begin{subproof}
            Fix $w$.
            Assume $w\in d$.
            Take $i\in J$ such that $w = (i, x(i) - z(i))$ by \cref{coord_diff_sequence}.
            $a\in\funs{J}{\reals}$ by \cref{coord_diff_sequence_function}.
            $b\in\funs{J}{\reals}$ by \cref{coord_diff_sequence_function}.
            Thus $a$ is a function by \cref{funs_elim}.
            Thus $b$ is a function by \cref{funs_elim}.
            Thus $(i, x(i) - y(i))\in a$ by \cref{coord_diff_sequence}.
            Thus $(i, y(i) - z(i))\in b$ by \cref{coord_diff_sequence}.
            Thus $a(i) = x(i) - y(i)$ by \cref{function_apply_intro}.
            Thus $b(i) = y(i) - z(i)$ by \cref{function_apply_intro}.
            $x\in\tuplePower{\reals}{J}$ and $y\in\tuplePower{\reals}{J}$ by \cref{reals_power}.
            Thus $x\in\funs{J}{\reals}$ and $y\in\funs{J}{\reals}$ by \cref{tuple_power}.
            Thus $x(i)\in\reals$ and $y(i)\in\reals$ by \cref{funs_type_apply}.
            $z\in\tuplePower{\reals}{J}$ by \cref{reals_power}.
            Thus $z\in\funs{J}{\reals}$ by \cref{tuple_power}.
            Thus $z(i)\in\reals$ by \cref{funs_type_apply}.
            $\neg{y(i)}\in\reals$ by \cref{real_add_inverse}.
            $\neg{z(i)}\in\reals$ by \cref{real_add_inverse}.
            $x(i) + \neg{y(i)}\in\reals$ by \cref{real_add_closed}.
            $y(i) + \neg{z(i)}\in\reals$ by \cref{real_add_closed}.
            $x(i) - y(i) = x(i) + \neg{y(i)}$ by \cref{real_sub}.
            $y(i) - z(i) = y(i) + \neg{z(i)}$ by \cref{real_sub}.
            $(x(i) + \neg{y(i)}) + (y(i) + \neg{z(i)}) = ((x(i) + \neg{y(i)}) + y(i)) + \neg{z(i)}$ by \cref{real_add_assoc}.
            $(x(i) + \neg{y(i)}) + y(i) = x(i) + (\neg{y(i)} + y(i))$ by \cref{real_add_assoc}.
            $\neg{y(i)} + y(i) = 0$ by \cref{real_add_inverse,real_add_comm}.
            Thus $x(i) + (\neg{y(i)} + y(i)) = x(i) + 0$.
            $0\in\reals$ by \cref{real_add_ident}.
            $x(i) + 0 = 0 + x(i)$ by \cref{real_add_comm}.
            $0 + x(i) = x(i)$ by \cref{real_add_ident}.
            Thus $x(i) + 0 = x(i)$.
            Thus $(x(i) + \neg{y(i)}) + y(i) = x(i)$.
            Thus $(x(i) + \neg{y(i)}) + (y(i) + \neg{z(i)}) = x(i) + \neg{z(i)}$ by \cref{real_add_assoc}.
            $x(i) + \neg{z(i)} = x(i) - z(i)$ by \cref{real_sub}.
            Thus $(x(i) - y(i)) + (y(i) - z(i)) = x(i) - z(i)$.
            Thus $a(i) + b(i) = x(i) - z(i)$.
            Thus $w = (i, a(i) + b(i))$.
            Thus $w\in c$ by \cref{coord_sum_sequence}.
        \end{subproof}
        Thus $d\subseteq c$ by \cref{subseteq}.
    \end{subproof}
    Thus $c = d$ by \cref{subseteq_antisymmetric}.
\end{proof}

% from §20 Item 25: euclidean metric is a metric
\begin{theorem}\label{euclidean_metric_is_metric}
    Suppose $n\in\naturalsPlus$.
    Let $J = \initialSegment{n}$.
    Let $X = \realsPower{J}$.
    Then $\euclideanMetric{n}$ is a metric on $X$.
\end{theorem}
\begin{proof}
    Let $d = \euclideanMetric{n}$.
    Show that $d\in\funs{X\times X}{\reals}$.
    \begin{subproof}
        $d$ is a function by \cref{euclidean_metric_function}.
        Show that $\dom{d} = X\times X$.
        \begin{subproof}
            Show that $X\times X\subseteq\dom{d}$.
            \begin{subproof}
                Show that for all $p$ such that $p\in X\times X$ we have $p\in\dom{d}$.
                \begin{subproof}
                    Fix $p$.
                    Assume $p\in X\times X$.
                    Take $x,y$ such that $x\in X$ and $y\in X$ and $p = (x,y)$ by \cref{times_elem_is_tuple}.
                    Take $r\in\reals$ such that $((x,y),r)\in d$ by \cref{euclidean_metric_value_exists}.
                    Thus $(p,r)\in d$ by \cref{pair_eq_iff}.
                    Thus $p\in\dom{d}$ by \cref{dom_intro}.
                \end{subproof}
                Thus $X\times X\subseteq\dom{d}$ by \cref{subseteq}.
            \end{subproof}
            Show that $\dom{d}\subseteq X\times X$.
            \begin{subproof}
                Show that for all $p$ such that $p\in\dom{d}$ we have $p\in X\times X$.
                \begin{subproof}
                    Fix $p$.
                    Assume $p\in\dom{d}$.
                    Take $r$ such that $(p,r)\in d$ by \cref{dom_iff}.
                    Take $p_0,r_0,s_0$ such that
                        $p_0\in\realsPower{J}\times\realsPower{J}$ and
                        $r_0\in\reals$ and
                        $s_0\in\reals$ and
                        $(p,r) = (p_0,r_0)$ and
                        $\sumFiniteSequence{\coordSquareSequence{n}{p_0}}{s_0}$ and
                        $r_0 = \sqrt{s_0}$
                        by \cref{euclidean_metric}.
                    Thus $p = p_0$ by \cref{pair_eq_iff}.
                    Take $x,y$ such that $x\in X$ and $y\in X$ and $p_0 = (x,y)$ by \cref{times_elem_is_tuple}.
                    Thus $p = (x,y)$ by \cref{pair_eq_iff}.
                    Thus $p\in X\times X$ by \cref{times_tuple_intro}.
                \end{subproof}
                Thus $\dom{d}\subseteq X\times X$ by \cref{subseteq}.
            \end{subproof}
            Thus $\dom{d} = X\times X$ by \cref{subseteq_antisymmetric}.
        \end{subproof}
        Show that for all $p\in\dom{d}$ we have $d(p)\in\reals$.
        \begin{subproof}
            Fix $p$.
            Assume $p\in\dom{d}$.
            Take $r$ such that $(p,r)\in d$ by \cref{dom_iff}.
            Thus $r\in\reals$ by \cref{euclidean_metric,pair_eq_iff}.
            Thus $d(p) = r$ by \cref{function_apply_intro}.
            Thus $d(p)\in\reals$.
        \end{subproof}
        Thus $d\in\funs{X\times X}{\reals}$ by \cref{funs_intro}.
    \end{subproof}
    Thus $d$ is a function by \cref{funs_elim}.
    Show that for all $x,y$ such that $x\in X$ and $y\in X$ we have $d((x,y)) \geq 0$.
    \begin{subproof}
        Fix $x$.
        Fix $y$.
        Assume $x\in X$ and $y\in X$.
        Take $r\in\reals$ such that $((x,y),r)\in d$ by \cref{euclidean_metric_value_exists}.
        Take $s\in\reals$ such that
            $\sumFiniteSequence{\coordSquareSequence{n}{(x,y)}}{s}$ and $r = \sqrt{s}$
            by \cref{euclidean_metric,pair_eq_iff}.
        Let $a = \coordSquareSequence{n}{(x,y)}$.
        Show that $0 \leq s$.
        \begin{subproof}
            Show that for all $i\in\dom{a}$ we have $0 \leq a(i)$.
            \begin{subproof}
                Fix $i$.
                Assume $i\in\dom{a}$.
                $a\in\funs{J}{\reals}$ by \cref{coord_square_sequence_function}.
                Thus $i\in J$ by \cref{funs_elim}.
                Thus $(i,\exp{x(i) - y(i)}{1+1})\in a$ by \cref{coord_square_sequence_elem}.
                $x\in\tuplePower{\reals}{J}$ and $y\in\tuplePower{\reals}{J}$ by \cref{reals_power}.
                Thus $x\in\funs{J}{\reals}$ and $y\in\funs{J}{\reals}$ by \cref{tuple_power}.
                Thus $x(i)\in\reals$ and $y(i)\in\reals$ by \cref{funs_type_apply}.
                Thus $\neg{y(i)}\in\reals$ by \cref{real_add_inverse}.
                Thus $x(i) + \neg{y(i)}\in\reals$ by \cref{real_add_closed}.
                Thus $x(i) - y(i)\in\reals$ by \cref{real_sub}.
                Thus $\exp{x(i) - y(i)}{1+1} \geq 0$ by \cref{real_sq_nonnegative}.
                $a$ is a function by \cref{funs_elim}.
                Thus $a(i) = \exp{x(i) - y(i)}{1+1}$ by \cref{function_apply_intro}.
                Thus $0 \leq a(i)$.
            \end{subproof}
            Thus $0 \leq s$ by \cref{sum_finite_sequence_nonneg}.
        \end{subproof}
        Thus $\sqrt{s}\in\reals$ and $\sqrt{s} \geq 0$ by \cref{positive_square_root}.
        Thus $d((x,y)) = r$ by \cref{function_apply_intro}.
        Thus $d((x,y)) \geq 0$.
    \end{subproof}
    Thus $d$ is nonnegative on $X$ by \cref{metric_nonnegative}.
    Show that $d$ is zero-characterized on $X$.
    \begin{subproof}
        Show that for all $x,y$ such that $x\in X$ and $y\in X$ we have $d((x,y)) = 0$ iff $x = y$.
        \begin{subproof}
            Fix $x$.
            Fix $y$.
            Assume $x\in X$ and $y\in X$.
            Show that if $d((x,y)) = 0$ then $x = y$.
            \begin{subproof}
                Assume $d((x,y)) = 0$.
                Take $r\in\reals$ such that $((x,y),r)\in d$ by \cref{euclidean_metric_value_exists}.
                Thus $d((x,y)) = r$ by \cref{function_apply_intro}.
                Thus $r = 0$.
                Let $d_1 = \squareMetric{n}$.
                Take $m\in\reals$ such that $((x,y),m)\in d_1$ by \cref{square_metric_value_exists}.
                Thus $m \leq r$ by \cref{euclidean_square_metric_bounds}.
                $d_1$ is a metric on $X$ by \cref{square_metric_is_metric}.
                Thus $d_1$ is nonnegative on $X$ by \cref{metric_on}.
                Thus $d_1((x,y)) \geq 0$ by \cref{metric_nonnegative}.
                Thus $d_1\in\funs{X\times X}{\reals}$ by \cref{metric_on}.
                Thus $d_1$ is a function by \cref{funs_elim}.
                Thus $d_1((x,y)) = m$ by \cref{function_apply_intro}.
                Thus $m \geq 0$.
                Thus $m = 0$ by \cref{real_leq_antisym}.
                Thus $d_1$ is zero-characterized on $X$ by \cref{metric_on}.
                Thus $d_1((x,y)) = 0$.
                Thus $x = y$ by \cref{metric_zero_characterization}.
            \end{subproof}
            Show that if $x = y$ then $d((x,y)) = 0$.
            \begin{subproof}
                Assume $x = y$.
                Take $r\in\reals$ such that $((x,y),r)\in d$ by \cref{euclidean_metric_value_exists}.
                Let $d_1 = \squareMetric{n}$.
                Take $m\in\reals$ such that $((x,y),m)\in d_1$ by \cref{square_metric_value_exists}.
                $d_1$ is a metric on $X$ by \cref{square_metric_is_metric}.
                Thus $d_1$ is zero-characterized on $X$ by \cref{metric_on}.
                Thus $d_1((x,y)) = 0$ by \cref{metric_zero_characterization}.
                Thus $d_1\in\funs{X\times X}{\reals}$ by \cref{metric_on}.
                Thus $d_1$ is a function by \cref{funs_elim}.
                Thus $d_1((x,y)) = m$ by \cref{function_apply_intro}.
                Thus $m = 0$.
                $0\in\reals$ by \cref{real_add_ident}.
                $0 < 1$ by \cref{real_zero_lt_one}.
                Thus $0 \leq 1$ by \cref{real_lt_imp_leq}.
                $1 \leq n$ by \cref{real_one_leq_naturalsplus}.
                $1\in\reals$ by \cref{real_naturals_subset_reals,real_one_in_naturals,subseteq}.
                $n\in\reals$ by \cref{real_naturals_subset_reals,subseteq}.
                Thus $0 \leq n$ by \cref{real_leq_trans}.
                Thus $\sqrt{n}\in\reals$ by \cref{positive_square_root}.
                Thus $\sqrt{n} \rmul m = \sqrt{n} \rmul 0$.
                Thus $\sqrt{n} \rmul 0 = 0 \rmul \sqrt{n}$ by \cref{real_mul_comm}.
                Thus $0 \rmul \sqrt{n} = 0$ by \cref{real_mul_zero}.
                Thus $\sqrt{n} \rmul 0 = 0$.
                Thus $r \leq \sqrt{n} \rmul m$ by \cref{euclidean_square_metric_bounds}.
                Thus $r \leq 0$.
                Thus $d((x,y)) = r$ by \cref{function_apply_intro}.
                Thus $d((x,y)) \geq 0$ by \cref{metric_nonnegative}.
                Thus $r \geq 0$.
                Thus $r = 0$ by \cref{real_leq_antisym}.
                Thus $d((x,y)) = 0$.
            \end{subproof}
        \end{subproof}
        Thus $d$ is zero-characterized on $X$ by \cref{metric_zero_characterization}.
    \end{subproof}
    Show that $d$ is symmetric on $X$.
    \begin{subproof}
        Show that for all $x,y$ such that $x\in X$ and $y\in X$ we have $d((x,y)) = d((y,x))$.
        \begin{subproof}
            Fix $x$.
            Fix $y$.
            Assume $x\in X$ and $y\in X$.
            Take $r\in\reals$ such that $((x,y),r)\in d$ by \cref{euclidean_metric_value_exists}.
            Take $s\in\reals$ such that $((y,x),s)\in d$ by \cref{euclidean_metric_value_exists}.
            Take $s_1\in\reals$ such that
                $\sumFiniteSequence{\coordSquareSequence{n}{(x,y)}}{s_1}$ and $r = \sqrt{s_1}$
                by \cref{euclidean_metric,pair_eq_iff}.
            Take $s_2\in\reals$ such that
                $\sumFiniteSequence{\coordSquareSequence{n}{(y,x)}}{s_2}$ and $s = \sqrt{s_2}$
                by \cref{euclidean_metric,pair_eq_iff}.
            Let $a = \coordSquareSequence{n}{(x,y)}$.
            Let $b = \coordSquareSequence{n}{(y,x)}$.
            Show that $a = b$.
            \begin{subproof}
                Show that for all $w$ we have $w\in a$ iff $w\in b$.
                \begin{subproof}
                    Fix $w$.
                    Show that if $w\in a$ then $w\in b$.
                    \begin{subproof}
                        Assume $w\in a$.
                        Let $p = (x,y)$.
                        Take $i\in\initialSegment{n}$ such that
                            $w = (i,\exp{\apply{\fst{p}}{i} - \apply{\snd{p}}{i}}{1+1})$ by \cref{coord_square_sequence}.
                        Thus $i\in J$.
                        $\fst{p} = x$ and $\snd{p} = y$ by \cref{fst_eq,snd_eq}.
                        Thus $w = (i,\exp{x(i) - y(i)}{1+1})$.
                        $x\in\tuplePower{\reals}{J}$ and $y\in\tuplePower{\reals}{J}$ by \cref{reals_power}.
                        Thus $x\in\funs{J}{\reals}$ and $y\in\funs{J}{\reals}$ by \cref{tuple_power}.
                        Thus $x(i)\in\reals$ and $y(i)\in\reals$ by \cref{funs_type_apply}.
                        Thus $\neg{y(i)}\in\reals$ by \cref{real_add_inverse}.
                        Thus $x(i) + \neg{y(i)}\in\reals$ by \cref{real_add_closed}.
                        Thus $x(i) - y(i)\in\reals$ by \cref{real_sub}.
                        Thus $y(i) - x(i) = \neg{(x(i) - y(i))}$ by \cref{real_sub_swap_neg}.
                        Let $u = x(i) - y(i)$.
                        Thus $y(i) - x(i) = \neg{u}$.
                        $1\in\naturalsPlus$ by \cref{real_one_in_naturals}.
                        $\neg{u}\in\reals$ by \cref{real_add_inverse}.
                        $\exp{u}{1+1} = u \rmul u$ by \cref{real_pow_succ,real_pow_one}.
                        $\exp{\neg{u}}{1+1} = \neg{u} \rmul \neg{u}$ by \cref{real_pow_succ,real_pow_one}.
                        $\neg{u} \rmul \neg{u} = \neg{(\neg{u} \rmul u)}$ by \cref{real_mul_neg_right}.
                        $\neg{u} \rmul u = u \rmul \neg{u}$ by \cref{real_mul_comm}.
                        $u \rmul \neg{u} = \neg{(u \rmul u)}$ by \cref{real_mul_neg_right}.
                        Thus $\neg{u} \rmul \neg{u} = \neg{(\neg{(u \rmul u)})}$.
                        $u \rmul u\in\reals$ by \cref{real_mul_closed}.
                        Thus $\neg{(\neg{(u \rmul u)})} = u \rmul u$ by \cref{real_neg_neg}.
                        Thus $\exp{\neg{u}}{1+1} = \exp{u}{1+1}$.
                        Thus $w = (i,\exp{y(i) - x(i)}{1+1})$.
                        Thus $i\in J$.
                        Thus $i\in\initialSegment{n}$.
                        Let $q = (y,x)$.
                        $\fst{q} = y$ and $\snd{q} = x$ by \cref{fst_eq,snd_eq}.
                        Thus $\exp{y(i) - x(i)}{1+1} = \exp{\apply{\fst{q}}{i} - \apply{\snd{q}}{i}}{1+1}$.
                        Thus $w = (i,\exp{\apply{\fst{q}}{i} - \apply{\snd{q}}{i}}{1+1})$.
                        Thus $w\in b$ by \cref{coord_square_sequence}.
                    \end{subproof}
                    Show that if $w\in b$ then $w\in a$.
                    \begin{subproof}
                        Assume $w\in b$.
                        Let $p = (y,x)$.
                        Take $i\in\initialSegment{n}$ such that
                            $w = (i,\exp{\apply{\fst{p}}{i} - \apply{\snd{p}}{i}}{1+1})$ by \cref{coord_square_sequence}.
                        Thus $i\in J$.
                        $\fst{p} = y$ and $\snd{p} = x$ by \cref{fst_eq,snd_eq}.
                        Thus $w = (i,\exp{y(i) - x(i)}{1+1})$.
                        $x\in\tuplePower{\reals}{J}$ and $y\in\tuplePower{\reals}{J}$ by \cref{reals_power}.
                        Thus $x\in\funs{J}{\reals}$ and $y\in\funs{J}{\reals}$ by \cref{tuple_power}.
                        Thus $x(i)\in\reals$ and $y(i)\in\reals$ by \cref{funs_type_apply}.
                        Thus $\neg{y(i)}\in\reals$ by \cref{real_add_inverse}.
                        Thus $x(i) + \neg{y(i)}\in\reals$ by \cref{real_add_closed}.
                        Thus $x(i) - y(i)\in\reals$ by \cref{real_sub}.
                        Thus $y(i) - x(i) = \neg{(x(i) - y(i))}$ by \cref{real_sub_swap_neg}.
                        Let $u = x(i) - y(i)$.
                        Thus $y(i) - x(i) = \neg{u}$.
                        $1\in\naturalsPlus$ by \cref{real_one_in_naturals}.
                        $\neg{u}\in\reals$ by \cref{real_add_inverse}.
                        $\exp{u}{1+1} = u \rmul u$ by \cref{real_pow_succ,real_pow_one}.
                        $\exp{\neg{u}}{1+1} = \neg{u} \rmul \neg{u}$ by \cref{real_pow_succ,real_pow_one}.
                        $\neg{u} \rmul \neg{u} = \neg{(\neg{u} \rmul u)}$ by \cref{real_mul_neg_right}.
                        $\neg{u} \rmul u = u \rmul \neg{u}$ by \cref{real_mul_comm}.
                        $u \rmul \neg{u} = \neg{(u \rmul u)}$ by \cref{real_mul_neg_right}.
                        Thus $\neg{u} \rmul \neg{u} = \neg{(\neg{(u \rmul u)})}$.
                        $u \rmul u\in\reals$ by \cref{real_mul_closed}.
                        Thus $\neg{(\neg{(u \rmul u)})} = u \rmul u$ by \cref{real_neg_neg}.
                        Thus $\exp{\neg{u}}{1+1} = \exp{u}{1+1}$.
                        Thus $w = (i,\exp{x(i) - y(i)}{1+1})$.
                        Thus $i\in J$.
                        Thus $i\in\initialSegment{n}$.
                        Let $q = (x,y)$.
                        $\fst{q} = x$ and $\snd{q} = y$ by \cref{fst_eq,snd_eq}.
                        Thus $\exp{x(i) - y(i)}{1+1} = \exp{\apply{\fst{q}}{i} - \apply{\snd{q}}{i}}{1+1}$.
                        Thus $w = (i,\exp{\apply{\fst{q}}{i} - \apply{\snd{q}}{i}}{1+1})$.
                        Thus $w\in a$ by \cref{coord_square_sequence}.
                    \end{subproof}
                \end{subproof}
                Thus $a = b$ by \cref{setext}.
            \end{subproof}
            Thus $s_1 = s_2$ by \cref{sum_finite_sequence_unique}.
            Thus $r = s$.
            Thus $d((x,y)) = r$ by \cref{function_apply_intro}.
            Thus $d((y,x)) = s$ by \cref{function_apply_intro}.
            Thus $d((x,y)) = d((y,x))$.
        \end{subproof}
        Thus $d$ is symmetric on $X$ by \cref{metric_symmetric}.
    \end{subproof}
    Show that $d$ satisfies the triangle inequality on $X$.
    \begin{subproof}
        Show that for all $x,y,z$ such that $x\in X$ and $y\in X$ and $z\in X$ we have
            $d((x,y)) + d((y,z)) \geq d((x,z))$.
        \begin{subproof}
            Fix $x$.
            Fix $y$.
            Fix $z$.
            Assume $x\in X$ and $y\in X$ and $z\in X$.
            Take $r_1\in\reals$ such that $((x,y),r_1)\in d$ by \cref{euclidean_metric_value_exists}.
            Take $r_2\in\reals$ such that $((y,z),r_2)\in d$ by \cref{euclidean_metric_value_exists}.
            Take $r_3\in\reals$ such that $((x,z),r_3)\in d$ by \cref{euclidean_metric_value_exists}.
            Take $s_1\in\reals$ such that
                $\sumFiniteSequence{\coordSquareSequence{n}{(x,y)}}{s_1}$ and $r_1 = \sqrt{s_1}$
                by \cref{euclidean_metric,pair_eq_iff}.
            Take $s_2\in\reals$ such that
                $\sumFiniteSequence{\coordSquareSequence{n}{(y,z)}}{s_2}$ and $r_2 = \sqrt{s_2}$
                by \cref{euclidean_metric,pair_eq_iff}.
            Take $s_3\in\reals$ such that
                $\sumFiniteSequence{\coordSquareSequence{n}{(x,z)}}{s_3}$ and $r_3 = \sqrt{s_3}$
                by \cref{euclidean_metric,pair_eq_iff}.
            Let $a = \coordDiffSequence{n}{x}{y}$.
            Let $b = \coordDiffSequence{n}{y}{z}$.
            Let $c = \coordSumSequence{n}{a}{b}$.
            Thus $a\in\funs{J}{\reals}$ by \cref{coord_diff_sequence_function}.
            Thus $b\in\funs{J}{\reals}$ by \cref{coord_diff_sequence_function}.
            Thus $c\in\funs{J}{\reals}$ by \cref{coord_sum_sequence_function}.
            Thus $\coordSquareSequenceMul{n}{a} = \coordSquareSequence{n}{(x,y)}$ by \cref{coord_square_sequence_diff}.
            Thus $\coordSquareSequenceMul{n}{b} = \coordSquareSequence{n}{(y,z)}$ by \cref{coord_square_sequence_diff}.
            Thus $c = \coordDiffSequence{n}{x}{z}$ by \cref{coord_diff_sequence_sum}.
            Thus $\coordSquareSequenceMul{n}{c} = \coordSquareSequence{n}{(x,z)}$ by \cref{coord_square_sequence_diff}.
            Thus $\sumFiniteSequence{\coordSquareSequenceMul{n}{a}}{s_1}$.
            Thus $\sumFiniteSequence{\coordSquareSequenceMul{n}{b}}{s_2}$.
            Thus $\sumFiniteSequence{\coordSquareSequenceMul{n}{c}}{s_3}$.
            Thus $\sqrt{s_3} \leq \sqrt{s_1} + \sqrt{s_2}$ by \cref{minkowski_finite_sequences}.
            Thus $r_3 \leq r_1 + r_2$.
            Thus $d((x,y)) = r_1$ by \cref{function_apply_intro}.
            Thus $d((y,z)) = r_2$ by \cref{function_apply_intro}.
            Thus $d((x,z)) = r_3$ by \cref{function_apply_intro}.
            Thus $d((x,z)) \leq d((x,y)) + d((y,z))$.
            Thus $d((x,y)) + d((y,z)) \geq d((x,z))$.
        \end{subproof}
        Thus $d$ satisfies the triangle inequality on $X$ by \cref{metric_triangle_inequality}.
    \end{subproof}
    Thus $d$ is a metric on $X$ by \cref{metric_on}.
\end{proof}

% from §20 Item 26: euclidean metric topology equals product topology
\begin{theorem}\label{euclidean_metric_product_topology}
    Suppose $n\in\naturalsPlus$.
    Let $J = \initialSegment{n}$.
    Let $X = \{(i,\reals) \mid i\in J\}$.
    Let $T = \{(i,\standardTopologyR) \mid i\in J\}$.
    Let $T_0 = \productTopologyIndexed{T}{X}$.
    Then $\metricTopology{\euclideanMetric{n}}{\realsPower{J}} = T_0$.
\end{theorem}
\begin{proof}
    Let $X_0 = \realsPower{J}$.
    Let $d = \euclideanMetric{n}$.
    Let $d_1 = \squareMetric{n}$.
    Let $T_E = \metricTopology{d}{X_0}$.
    Let $T_S = \metricTopology{d_1}{X_0}$.
    $1\in\naturalsPlus$ by \cref{real_one_in_naturals}.
    $1 \leq n$ by \cref{real_one_leq_naturalsplus}.
    Thus $1\in J$ by \cref{initial_segment_naturalsplus}.
    Thus $J$ is inhabited.
    Show that $X$ is a function.
    \begin{subproof}
        Show that for all $w$ such that $w\in X$ there exist $x,y$ such that $w = (x,y)$.
        \begin{subproof}
            Fix $w$.
            Assume $w\in X$.
            Take $i\in J$ such that $w = (i,\reals)$.
            Thus there exist $x,y$ such that $w = (x,y)$.
        \end{subproof}
        Thus $X$ is a relation by \cref{relation}.
        Show that for all $i,y_1,y_2$ such that $(i,y_1)\in X$ and $(i,y_2)\in X$ we have $y_1 = y_2$.
        \begin{subproof}
            Fix $i$.
            Fix $y_1$.
            Fix $y_2$.
            Assume $(i,y_1)\in X$ and $(i,y_2)\in X$.
            Then $y_1 = \reals$ and $y_2 = \reals$.
            Thus $y_1 = y_2$.
        \end{subproof}
        Thus $X$ is right-unique by \cref{rightunique}.
        Thus $X$ is a function.
    \end{subproof}
    Show that $T$ is a function.
    \begin{subproof}
        Show that for all $w$ such that $w\in T$ there exist $x,y$ such that $w = (x,y)$.
        \begin{subproof}
            Fix $w$.
            Assume $w\in T$.
            Take $i\in J$ such that $w = (i,\standardTopologyR)$.
            Thus there exist $x,y$ such that $w = (x,y)$.
        \end{subproof}
        Thus $T$ is a relation by \cref{relation}.
        Show that for all $i,y_1,y_2$ such that $(i,y_1)\in T$ and $(i,y_2)\in T$ we have $y_1 = y_2$.
        \begin{subproof}
            Fix $i$.
            Fix $y_1$.
            Fix $y_2$.
            Assume $(i,y_1)\in T$ and $(i,y_2)\in T$.
            Then $y_1 = \standardTopologyR$ and $y_2 = \standardTopologyR$.
            Thus $y_1 = y_2$.
        \end{subproof}
        Thus $T$ is right-unique by \cref{rightunique}.
        Thus $T$ is a function.
    \end{subproof}
    Show that $\dom{X} = J$.
    \begin{subproof}
        Show that for all $i$ such that $i\in J$ we have $i\in\dom{X}$.
        \begin{subproof}
            Fix $i$.
            Assume $i\in J$.
            Thus $(i,\reals)\in X$.
            Thus $i\in\dom{X}$ by \cref{dom_intro}.
        \end{subproof}
        Thus $J\subseteq\dom{X}$ by \cref{subseteq}.
        Show that for all $i$ such that $i\in\dom{X}$ we have $i\in J$.
        \begin{subproof}
            Fix $i$.
            Assume $i\in\dom{X}$.
            Take $y$ such that $(i,y)\in X$ by \cref{dom_iff}.
            Thus $i\in J$.
        \end{subproof}
        Thus $\dom{X}\subseteq J$ by \cref{subseteq}.
        Thus $\dom{X} = J$ by \cref{subseteq_antisymmetric}.
    \end{subproof}
    Show that $\dom{T} = J$.
    \begin{subproof}
        Show that for all $i$ such that $i\in J$ we have $i\in\dom{T}$.
        \begin{subproof}
            Fix $i$.
            Assume $i\in J$.
            Thus $(i,\standardTopologyR)\in T$.
            Thus $i\in\dom{T}$ by \cref{dom_intro}.
        \end{subproof}
        Thus $J\subseteq\dom{T}$ by \cref{subseteq}.
        Show that for all $i$ such that $i\in\dom{T}$ we have $i\in J$.
        \begin{subproof}
            Fix $i$.
            Assume $i\in\dom{T}$.
            Take $y$ such that $(i,y)\in T$ by \cref{dom_iff}.
            Thus $i\in J$.
        \end{subproof}
        Thus $\dom{T}\subseteq J$ by \cref{subseteq}.
        Thus $\dom{T} = J$ by \cref{subseteq_antisymmetric}.
    \end{subproof}
    Show that for all $\alpha\in\dom{X}$ we have $\apply{T}{\alpha}$ is a topology on $\apply{X}{\alpha}$.
    \begin{subproof}
        Fix $\alpha$.
        Assume $\alpha\in\dom{X}$.
        Thus $\alpha\in J$ by \cref{subseteq}.
        Thus $(\alpha,\reals)\in X$.
        Thus $\apply{X}{\alpha} = \reals$ by \cref{function_apply_intro}.
        Thus $(\alpha,\standardTopologyR)\in T$.
        Thus $\apply{T}{\alpha} = \standardTopologyR$ by \cref{function_apply_intro}.
        $\standardBasisR$ is a basis on $\reals$ by \cref{standard_basis_is_basis}.
        Thus $\standardTopologyR$ is a topology on $\reals$ by \cref{basis_topology_is_topology,standard_topology_reals}.
        Thus $\apply{T}{\alpha}$ is a topology on $\apply{X}{\alpha}$.
    \end{subproof}
    Let $S_0 = \productSubbasis{T}{X}$.
    Let $B_0 = \finiteIntersections{S_0}$.
    $B_0$ is a basis for $T_0$ on $\productFamily{X}$ by \cref{product_subbasis_finite_intersections_basis}.
    Thus $\productFamily{X} = X_0$ by \cref{product_family_reals_power}.
    Thus $B_0$ is a basis for $T_0$ on $X_0$.
    $d$ is a metric on $X_0$ by \cref{euclidean_metric_is_metric}.
    $d_1$ is a metric on $X_0$ by \cref{square_metric_is_metric}.
    Let $B_S = \metricBasis{d_1}{X_0}$.
    $B_S$ is a basis on $X_0$ by \cref{metric_basis_is_basis}.
    Thus $B_S$ is a basis for $T_S$ on $X_0$ by \cref{basis_for_topology,metric_topology}.
    Show that $T_E = T_S$.
    \begin{subproof}
        Show that $T_E$ is finer than $T_S$ on $X_0$.
        \begin{subproof}
            Show that for all $x\in X_0$ we have
                for all $\epsilon\in\reals$ such that $0 < \epsilon$ there exists $\delta\in\reals$
                    such that $0 < \delta$ and
                    $\metricBall{d}{X_0}{x}{\delta} \subseteq \metricBall{d_1}{X_0}{x}{\epsilon}$.
            \begin{subproof}
                Fix $x$.
                Assume $x\in X_0$.
                Fix $\epsilon\in\reals$.
                Assume $0 < \epsilon$.
                Let $\delta = \epsilon$.
                $0 < \delta$.
                Show that for all $y$ such that $y\in\metricBall{d}{X_0}{x}{\delta}$ we have $y\in\metricBall{d_1}{X_0}{x}{\epsilon}$.
                \begin{subproof}
                    Fix $y$.
                    Assume $y\in\metricBall{d}{X_0}{x}{\delta}$.
                    Thus $y\in X_0$ and $d((x,y)) < \delta$ by \cref{metric_ball}.
                    Take $r\in\reals$ such that $((x,y),r)\in d$ by \cref{euclidean_metric_value_exists}.
                    Take $m\in\reals$ such that $((x,y),m)\in d_1$ by \cref{square_metric_value_exists}.
                    Thus $m \leq r$ by \cref{euclidean_square_metric_bounds}.
                    $d$ is a function by \cref{euclidean_metric_function}.
                    $d_1$ is a function by \cref{square_metric_function}.
                    Thus $d((x,y)) = r$ by \cref{function_apply_intro}.
                    Thus $r < \epsilon$ by \cref{real_lt_subst_left}.
                    Show that $m < \epsilon$.
                    \begin{subproof}
                        $m < r$ or $m = r$ by \cref{real_leq_split}.
                        \begin{byCase}
                        \caseOf{$m < r$.} Thus $m < \epsilon$ by \cref{real_lt_trans}.
                        \caseOf{$m = r$.} Thus $m < \epsilon$ by \cref{real_lt_subst_left}.
                        \end{byCase}
                    \end{subproof}
                    Thus $d_1((x,y)) = m$ by \cref{function_apply_intro}.
                    Thus $y\in\metricBall{d_1}{X_0}{x}{\epsilon}$ by \cref{metric_ball}.
                \end{subproof}
                Thus $\metricBall{d}{X_0}{x}{\delta} \subseteq \metricBall{d_1}{X_0}{x}{\epsilon}$ by \cref{subseteq}.
            \end{subproof}
            Thus $T_E$ is finer than $T_S$ on $X_0$ by \cref{metric_topology_finer_iff}.
        \end{subproof}
        Show that $T_S$ is finer than $T_E$ on $X_0$.
        \begin{subproof}
            Show that for all $x\in X_0$ we have
                for all $\epsilon\in\reals$ such that $0 < \epsilon$ there exists $\delta\in\reals$
                    such that $0 < \delta$ and
                    $\metricBall{d_1}{X_0}{x}{\delta} \subseteq \metricBall{d}{X_0}{x}{\epsilon}$.
            \begin{subproof}
                Fix $x$.
                Assume $x\in X_0$.
                Fix $\epsilon\in\reals$.
                Assume $0 < \epsilon$.
                $n\in\reals$ by \cref{real_naturals_subset_reals,subseteq}.
                Show that $0 \leq n$.
                \begin{subproof}
                    $0 < 1$ by \cref{real_zero_lt_one}.
                    $0\in\reals$ by \cref{real_add_ident}.
                    $1\in\reals$ by \cref{real_naturals_subset_reals,subseteq}.
                    $1 \leq n$ by \cref{real_one_leq_naturalsplus}.
                    Thus $1 < n$ or $1 = n$ by \cref{real_leq_split}.
                    \begin{byCase}
                    \caseOf{$1 < n$.}
                        Thus $0 < n$ by \cref{real_lt_trans}.
                        Thus $0 \leq n$.
                    \caseOf{$1 = n$.}
                        Thus $0 < n$ by \cref{real_lt_subst_right}.
                        Thus $0 \leq n$.
                    \end{byCase}
                \end{subproof}
                Thus $\sqrt{n}\in\reals$ and $\sqrt{n} \geq 0$ and $\sqrt{n} \rmul \sqrt{n} = n$ by \cref{positive_square_root}.
                Show that $\sqrt{n}$ is positive.
                \begin{subproof}
                    Show that $\sqrt{n} \neq 0$.
                    \begin{subproof}
                        Suppose not.
                        Thus $\sqrt{n} = 0$.
                        Thus $n = 0$ by \cref{positive_square_root,real_mul_zero}.
                        $0 < 1$ by \cref{real_zero_lt_one}.
                        $1 \leq n$ by \cref{real_one_leq_naturalsplus}.
                        Show that $0 < n$.
                        \begin{subproof}
                            $0\in\reals$ by \cref{real_add_ident}.
                            $1\in\reals$ by \cref{real_naturals_subset_reals,subseteq}.
                            $n\in\reals$ by \cref{real_naturals_subset_reals,subseteq}.
                            Thus $1 < n$ or $1 = n$ by \cref{real_leq_split}.
                            \begin{byCase}
                            \caseOf{$1 < n$.}
                                Thus $0 < n$ by \cref{real_lt_trans}.
                            \caseOf{$1 = n$.}
                                Thus $0 < n$ by \cref{real_lt_subst_right}.
                            \end{byCase}
                        \end{subproof}
                        Thus $0 < 0$ by \cref{real_lt_subst_right}.
                        Contradiction by \cref{real_lt_irreflexive}.
                    \end{subproof}
                    Thus $\sqrt{n} > 0$.
                \end{subproof}
                Let $\delta = \epsilon \rmul \inv{\sqrt{n}}$.
                Show that $0 < \delta$.
                \begin{subproof}
                    Show that $\sqrt{n} \neq 0$.
                    \begin{subproof}
                        Suppose not.
                        Thus $\sqrt{n} = 0$.
                        Thus $0 < 0$ by \cref{real_lt_subst_left}.
                        Contradiction by \cref{real_lt_irreflexive}.
                    \end{subproof}
                    $\inv{\sqrt{n}}\in\reals$ by \cref{real_mul_inverse}.
                    $\inv{\sqrt{n}} > 0$ by \cref{real_inv_pos}.
                    $\epsilon \rmul \inv{\sqrt{n}}\in\reals$ by \cref{real_mul_closed}.
                    Thus $\delta > 0$ by \cref{real_mul_pos_pos_gt}.
                \end{subproof}
                Show that for all $y$ such that $y\in\metricBall{d_1}{X_0}{x}{\delta}$ we have $y\in\metricBall{d}{X_0}{x}{\epsilon}$.
                \begin{subproof}
                    Fix $y$.
                    Assume $y\in\metricBall{d_1}{X_0}{x}{\delta}$.
                    Thus $y\in X_0$ and $d_1((x,y)) < \delta$ by \cref{metric_ball}.
                    Take $m\in\reals$ such that $((x,y),m)\in d_1$ by \cref{square_metric_value_exists}.
                    Take $r\in\reals$ such that $((x,y),r)\in d$ by \cref{euclidean_metric_value_exists}.
                    $d$ is a function by \cref{euclidean_metric_function}.
                    $d_1$ is a function by \cref{square_metric_function}.
                    Thus $d_1((x,y)) = m$ by \cref{function_apply_intro}.
                    Thus $m < \delta$ by \cref{real_lt_subst_left}.
                    Thus $r \leq \sqrt{n} \rmul m$ by \cref{euclidean_square_metric_bounds}.
                    Show that $\sqrt{n} \rmul m < \epsilon$.
                    \begin{subproof}
                        $\sqrt{n}\in\reals$.
                        $m\in\reals$.
                        $\delta\in\reals$.
                        $m \rmul \sqrt{n} < \delta \rmul \sqrt{n}$ by \cref{real_lt_mul_pos}.
                        Thus $\sqrt{n} \rmul m < \sqrt{n} \rmul \delta$ by \cref{real_mul_comm}.
                        Show that $\sqrt{n} \rmul \delta = \epsilon$.
                        \begin{subproof}
                            $\sqrt{n}\in\reals$.
                            $\epsilon\in\reals$.
                            $\inv{\sqrt{n}}\in\reals$.
                            $\sqrt{n} \rmul \delta = \sqrt{n} \rmul (\epsilon \rmul \inv{\sqrt{n}})$.
                            $\sqrt{n} \rmul (\epsilon \rmul \inv{\sqrt{n}}) = (\sqrt{n} \rmul \epsilon) \rmul \inv{\sqrt{n}}$ by \cref{real_mul_assoc}.
                            $(\sqrt{n} \rmul \epsilon) \rmul \inv{\sqrt{n}} = (\epsilon \rmul \sqrt{n}) \rmul \inv{\sqrt{n}}$ by \cref{real_mul_comm}.
                            $(\epsilon \rmul \sqrt{n}) \rmul \inv{\sqrt{n}} = \epsilon \rmul (\sqrt{n} \rmul \inv{\sqrt{n}})$ by \cref{real_mul_assoc}.
                            Show that $\sqrt{n} \neq 0$.
                            \begin{subproof}
                                Suppose not.
                                Thus $\sqrt{n} = 0$.
                                Thus $0 < 0$ by \cref{real_lt_subst_left}.
                                Contradiction by \cref{real_lt_irreflexive}.
                            \end{subproof}
                            $\sqrt{n} \rmul \inv{\sqrt{n}} = 1$ by \cref{real_mul_inverse,real_mul_comm}.
                            Thus $\sqrt{n} \rmul \delta = \epsilon \rmul 1$.
                            $1\in\reals$ by \cref{real_mul_ident}.
                            $1 \rmul \epsilon = \epsilon$ by \cref{real_mul_ident}.
                            Thus $\epsilon \rmul 1 = \epsilon$ by \cref{real_mul_comm}.
                        \end{subproof}
                        Thus $\sqrt{n} \rmul m < \epsilon$ by \cref{real_lt_subst_right}.
                    \end{subproof}
                    Show that $r < \epsilon$.
                    \begin{subproof}
                        $r\in\reals$.
                        $\sqrt{n} \rmul m\in\reals$ by \cref{real_mul_closed}.
                        $\epsilon\in\reals$.
                        $r < \sqrt{n} \rmul m$ or $r = \sqrt{n} \rmul m$ by \cref{real_leq_split}.
                        \begin{byCase}
                        \caseOf{$r < \sqrt{n} \rmul m$.} Thus $r < \epsilon$ by \cref{real_lt_trans}.
                        \caseOf{$r = \sqrt{n} \rmul m$.} Thus $r < \epsilon$ by \cref{real_lt_subst_left}.
                        \end{byCase}
                    \end{subproof}
                    Thus $d((x,y)) = r$ by \cref{function_apply_intro}.
                    Thus $y\in\metricBall{d}{X_0}{x}{\epsilon}$ by \cref{metric_ball}.
                \end{subproof}
                Thus $\metricBall{d_1}{X_0}{x}{\delta} \subseteq \metricBall{d}{X_0}{x}{\epsilon}$ by \cref{subseteq}.
            \end{subproof}
            Thus $T_S$ is finer than $T_E$ on $X_0$ by \cref{metric_topology_finer_iff}.
        \end{subproof}
        Thus $T_S\subseteq T_E$ by \cref{finer_topology}.
        Thus $T_E\subseteq T_S$ by \cref{finer_topology}.
        Thus $T_E = T_S$ by \cref{subseteq_antisymmetric}.
    \end{subproof}
    Show that $T_S = T_0$.
    \begin{subproof}
        Show that $T_0$ is finer than $T_S$ on $X_0$.
        \begin{subproof}
            Show that for all $x,B_2$ such that $x\in X_0$ and $B_2\in B_S$ and $x\in B_2$
                there exists $B_3\in B_0$ such that $x\in B_3$ and $B_3\subseteq B_2$.
            \begin{subproof}
                Fix $x$.
                Fix $B_2$.
                Assume $x\in X_0$ and $B_2\in B_S$ and $x\in B_2$.
                Take $x_0,\epsilon$ such that $x_0\in X_0$ and $\epsilon\in\reals$ and $0 < \epsilon$ and
                    $B_2 = \metricBall{d_1}{X_0}{x_0}{\epsilon}$ by \cref{metric_basis}.
                Let $R_g = \{w\in J\times\pow{\productFamily{X}} \mid \exists i\in J.
                    w = (i,\preimg{\piIndex{X}{i}}{\intervalopen{x_0(i) - \epsilon}{x_0(i) + \epsilon}})\}$.
                Show that $R_g$ is a relation.
                \begin{subproof}
                    Show that for all $w$ such that $w\in R_g$ there exist $i,A$ such that $w = (i,A)$.
                    \begin{subproof}
                        Fix $w$.
                        Assume $w\in R_g$.
                        Thus $w\in J\times\pow{\productFamily{X}}$.
                        Take $i,A$ such that $i\in J$ and $A\in\pow{\productFamily{X}}$ and $w = (i,A)$ by \cref{times_elem_is_tuple}.
                    \end{subproof}
                    Thus $R_g$ is a relation by \cref{relation}.
                \end{subproof}
                Show that for all $i\in J$ there exists $A$ such that $i\mathrel{R_g} A$.
                \begin{subproof}
                    Fix $i$.
                    Assume $i\in J$.
                    Let $A = \preimg{\piIndex{X}{i}}{\intervalopen{x_0(i) - \epsilon}{x_0(i) + \epsilon}}$.
                    Show that $A\subseteq\productFamily{X}$.
                    \begin{subproof}
                        Show that for all $z$ such that $z\in A$ we have $z\in\productFamily{X}$.
                        \begin{subproof}
                            Fix $z$.
                            Assume $z\in A$.
                            Take $y\in\intervalopen{x_0(i) - \epsilon}{x_0(i) + \epsilon}$ such that $z\mathrel{\piIndex{X}{i}} y$ by \cref{preimg_iff}.
                            Thus $(z,y)\in\piIndex{X}{i}$.
                            Take $x_1\in\productFamily{X}$ such that $(z,y) = (x_1,\apply{x_1}{i})$ by \cref{projection_map_indexed}.
                            Thus $z = x_1$ by \cref{pair_eq_iff}.
                            Thus $z\in\productFamily{X}$.
                        \end{subproof}
                        Thus $A\subseteq\productFamily{X}$ by \cref{subseteq}.
                    \end{subproof}
                    Thus $A\in\pow{\productFamily{X}}$ by \cref{powerset_intro}.
                    Thus $(i,A)\in J\times\pow{\productFamily{X}}$ by \cref{times_tuple_intro}.
                    Thus $(i,A)\in R_g$.
                    Thus $i\mathrel{R_g} A$.
                \end{subproof}
                There exists $g$ such that $g$ is a function on $J$ and for all $i\in J$ we have $i\mathrel{R_g}\apply{g}{i}$ by \cref{choice_function}.
                Take $g$ such that $g$ is a function on $J$ and for all $i\in J$ we have $i\mathrel{R_g}\apply{g}{i}$.
                Let $F = \img{g}{J}$.
                $J$ is finite by \cref{initial_segment_finite}.
                Thus $F$ is finite by \cref{finite_image_function}.
                Show that $F$ is inhabited.
                \begin{subproof}
                    Take $i$ such that $i\in J$.
                    Thus $(i,\apply{g}{i})\in g$ by \cref{function_apply_elim}.
                    Thus $\apply{g}{i}\in F$ by \cref{img_iff}.
                    Thus $F\neq\emptyset$ by \cref{emptyset}.
                \end{subproof}
                Show that for all $A$ such that $A\in F$ we have $A\in S_0$.
                \begin{subproof}
                    Fix $A$.
                    Assume $A\in F$.
                    Take $i\in J$ such that $(i,A)\in g$ by \cref{img_iff}.
                    Thus $\apply{g}{i} = A$ by \cref{function_apply_intro}.
                    Thus $i\mathrel{R_g}\apply{g}{i}$.
                    Take $i_0\in J$ such that
                        $(i,\apply{g}{i}) = (i_0,\preimg{\piIndex{X}{i_0}}{\intervalopen{x_0(i_0) - \epsilon}{x_0(i_0) + \epsilon}})$.
                    Thus $\apply{g}{i} = \preimg{\piIndex{X}{i}}{\intervalopen{x_0(i) - \epsilon}{x_0(i) + \epsilon}}$ by \cref{pair_eq_iff}.
                    Thus $A = \preimg{\piIndex{X}{i}}{\intervalopen{x_0(i) - \epsilon}{x_0(i) + \epsilon}}$.
                    Show that $\intervalopen{x_0(i) - \epsilon}{x_0(i) + \epsilon}\in\standardTopologyR$.
                    \begin{subproof}
                        $x_0\in X_0$.
                        Thus $x_0\in\realsPower{J}$.
                        Thus $x_0\in\tuplePower{\reals}{J}$ by \cref{reals_power}.
                        Thus $x_0\in\funs{J}{\reals}$ by \cref{tuple_power}.
                        Thus $x_0(i)\in\reals$ by \cref{funs_type_apply}.
                        $\epsilon\in\reals$.
                        $\neg{\epsilon}\in\reals$ by \cref{real_add_inverse}.
                        Thus $x_0(i) - \epsilon\in\reals$ by \cref{real_sub,real_add_closed}.
                        Thus $x_0(i) + \epsilon\in\reals$ by \cref{real_add_closed}.
                        Show that $x_0(i) - \epsilon < x_0(i) + \epsilon$.
                        \begin{subproof}
                            $0 < \epsilon$.
                            $0\in\reals$ by \cref{real_add_ident}.
                            Thus $\neg{\epsilon} < 0$ by \cref{real_lt_neg,real_neg_zero,real_lt_subst_right}.
                            $\neg{\epsilon}\in\reals$ by \cref{real_add_inverse}.
                            $x_0(i)\in\reals$ by \cref{funs_type_apply}.
                            Thus $\neg{\epsilon} + x_0(i) < 0 + x_0(i)$ by \cref{real_lt_add}.
                            Thus $x_0(i) + \neg{\epsilon} < x_0(i) + 0$ by \cref{real_add_comm}.
                            $x_0(i) + \neg{\epsilon} = x_0(i) - \epsilon$ by \cref{real_sub}.
                            $0 + x_0(i) = x_0(i)$ by \cref{real_add_ident}.
                            Thus $x_0(i) + 0 = x_0(i)$ by \cref{real_add_comm}.
                            Thus $x_0(i) - \epsilon < x_0(i)$.
                            Thus $0 + x_0(i) < \epsilon + x_0(i)$ by \cref{real_lt_add}.
                            Thus $x_0(i) < x_0(i) + \epsilon$ by \cref{real_add_comm,real_add_ident}.
                            Thus $x_0(i) - \epsilon < x_0(i) + \epsilon$ by \cref{real_lt_trans}.
                        \end{subproof}
                        $\intervalopen{x_0(i) - \epsilon}{x_0(i) + \epsilon}\subseteq\reals$ by \cref{intervalopen_subset_reals}.
                        Thus $\intervalopen{x_0(i) - \epsilon}{x_0(i) + \epsilon}\in\pow{\reals}$ by \cref{powerset_intro}.
                        Thus $\intervalopen{x_0(i) - \epsilon}{x_0(i) + \epsilon}\in\standardBasisR$ by \cref{standard_basis_reals}.
                        $\standardBasisR$ is a basis on $\reals$ by \cref{standard_basis_is_basis}.
                        Thus $\intervalopen{x_0(i) - \epsilon}{x_0(i) + \epsilon}\in\standardTopologyR$
                            by \cref{basis_elem_open_in_generated,standard_basis_is_basis,standard_topology_reals}.
                    \end{subproof}
                    Let $U = \intervalopen{x_0(i) - \epsilon}{x_0(i) + \epsilon}$.
                    Thus $U\in\standardTopologyR$.
                    Thus $(i,\standardTopologyR)\in T$.
                    $T$ is a function.
                    Thus $\apply{T}{i} = \standardTopologyR$ by \cref{function_apply_intro}.
                    Thus $U\in\apply{T}{i}$.
                    Show that $A\subseteq\productFamily{X}$.
                    \begin{subproof}
                        Show that for all $z$ such that $z\in A$ we have $z\in\productFamily{X}$.
                        \begin{subproof}
                            Fix $z$.
                            Assume $z\in A$.
                            Take $y\in U$ such that $z\mathrel{\piIndex{X}{i}} y$ by \cref{preimg_iff}.
                            Thus $(z,y)\in\piIndex{X}{i}$.
                            Take $x_1\in\productFamily{X}$ such that $(z,y) = (x_1,\apply{x_1}{i})$ by \cref{projection_map_indexed}.
                            Thus $z = x_1$ by \cref{pair_eq_iff}.
                            Thus $z\in\productFamily{X}$.
                        \end{subproof}
                        Thus $A\subseteq\productFamily{X}$ by \cref{subseteq}.
                    \end{subproof}
                    Thus $A\in\pow{\productFamily{X}}$ by \cref{powerset_intro}.
                    Thus $A\in\productSubbasisAt{T}{X}{i}$ by \cref{product_subbasis_indexed}.
                    Thus $A\in S_0$ by \cref{product_subbasis_union_indexed}.
                \end{subproof}
                Thus $F\subseteq S_0$ by \cref{subseteq}.
                Show that $B_2 = \inters{F}$.
                \begin{subproof}
                    Show that $B_2\subseteq\inters{F}$.
                    \begin{subproof}
                        Show that for all $y$ such that $y\in B_2$ we have $y\in\inters{F}$.
                        \begin{subproof}
                            Fix $y$.
                            Assume $y\in B_2$.
                            Thus $y\in X_0$ and $d_1((x_0,y)) < \epsilon$ by \cref{metric_ball}.
                            Thus $x_0\in\realsPower{J}$.
                            Thus $y\in\realsPower{J}$.
                            Take $m\in\reals$ such that $((x_0,y),m)\in d_1$ and
                                for all $j\in J$ we have $\abs{x_0(j) - y(j)} \leq m$ by \cref{square_metric_value_exists}.
                            $d_1$ is a function by \cref{square_metric_function}.
                            Thus $d_1((x_0,y)) = m$ by \cref{function_apply_intro}.
                            Thus $m < \epsilon$ by \cref{real_lt_subst_left}.
                            Show that for all $A$ such that $A\in F$ we have $y\in A$.
                            \begin{subproof}
                                Fix $A$.
                                Assume $A\in F$.
                                Take $i\in J$ such that $(i,A)\in g$ by \cref{img_iff}.
                                Thus $\apply{g}{i} = A$ by \cref{function_apply_intro}.
                                Thus $i\mathrel{R_g}\apply{g}{i}$.
                                Take $i_0\in J$ such that
                                    $(i,\apply{g}{i}) = (i_0,\preimg{\piIndex{X}{i_0}}{\intervalopen{x_0(i_0) - \epsilon}{x_0(i_0) + \epsilon}})$.
                                Thus $\apply{g}{i} = \preimg{\piIndex{X}{i}}{\intervalopen{x_0(i) - \epsilon}{x_0(i) + \epsilon}}$ by \cref{pair_eq_iff}.
                                Thus $A = \preimg{\piIndex{X}{i}}{\intervalopen{x_0(i) - \epsilon}{x_0(i) + \epsilon}}$.
                                Show that $y\in A$.
                                \begin{subproof}
                                    Thus $y\in\productFamily{X}$ by \cref{product_family_reals_power,subseteq}.
                                    Thus $(y,\apply{y}{i})\in\piIndex{X}{i}$ by \cref{projection_map_indexed}.
                                    Show that $\apply{y}{i}\in\intervalopen{x_0(i) - \epsilon}{x_0(i) + \epsilon}$.
                                    \begin{subproof}
                                        $x_0\in\realsPower{J}$.
                                        $y\in\realsPower{J}$.
                                        Thus $x_0\in\funs{J}{\reals}$ and $y\in\funs{J}{\reals}$ by \cref{reals_power,tuple_power}.
                                        Thus $x_0(i)\in\reals$ and $\apply{y}{i}\in\reals$ by \cref{funs_type_apply}.
                                        $\neg{\apply{y}{i}}\in\reals$ by \cref{real_add_inverse}.
                                        $x_0(i) - \apply{y}{i} = x_0(i) + \neg{\apply{y}{i}}$ by \cref{real_sub}.
                                        Thus $x_0(i) - \apply{y}{i}\in\reals$ by \cref{real_add_closed}.
                                        Thus $\abs{x_0(i) - \apply{y}{i}}\in\reals$ by \cref{real_abs_in_reals}.
                                        Thus $\abs{x_0(i) - \apply{y}{i}} \leq m$.
                                        Show that $\abs{x_0(i) - \apply{y}{i}} < \epsilon$.
                                        \begin{subproof}
                                            $\abs{x_0(i) - \apply{y}{i}} < m$ or $\abs{x_0(i) - \apply{y}{i}} = m$ by \cref{real_leq_split}.
                                            \begin{byCase}
                                            \caseOf{$\abs{x_0(i) - \apply{y}{i}} < m$.}
                                                Thus $\abs{x_0(i) - \apply{y}{i}} < \epsilon$ by \cref{real_lt_trans}.
                                            \caseOf{$\abs{x_0(i) - \apply{y}{i}} = m$.}
                                                Thus $\abs{x_0(i) - \apply{y}{i}} < \epsilon$ by \cref{real_lt_subst_left}.
                                            \end{byCase}
                                        \end{subproof}
                                        Thus $x_0(i) - \epsilon < \apply{y}{i}$ and $\apply{y}{i} < x_0(i) + \epsilon$
                                            by \cref{real_interval_from_abs_lt}.
                                        Thus $\apply{y}{i}\in\intervalopen{x_0(i) - \epsilon}{x_0(i) + \epsilon}$ by \cref{intervalopen}.
                                    \end{subproof}
                                    Thus $y\in A$ by \cref{preimg_iff}.
                                \end{subproof}
                            \end{subproof}
                            Thus $y\in\inters{F}$ by \cref{inters_iff_forall}.
                        \end{subproof}
                        Thus $B_2\subseteq\inters{F}$ by \cref{subseteq}.
                    \end{subproof}
                    Show that $\inters{F}\subseteq B_2$.
                    \begin{subproof}
                        Show that for all $y$ such that $y\in\inters{F}$ we have $y\in B_2$.
                        \begin{subproof}
                            Fix $y$.
                            Assume $y\in\inters{F}$.
                            Take $i$ such that $i\in J$.
                            Thus $(i,\apply{g}{i})\in g$ by \cref{function_apply_elim}.
                            Thus $\apply{g}{i}\in F$ by \cref{img_iff}.
                            Thus $y\in\apply{g}{i}$ by \cref{inters_iff_forall}.
                            Thus $i\mathrel{R_g}\apply{g}{i}$.
                            Take $i_0\in J$ such that
                                $(i,\apply{g}{i}) = (i_0,\preimg{\piIndex{X}{i_0}}{\intervalopen{x_0(i_0) - \epsilon}{x_0(i_0) + \epsilon}})$.
                            Thus $\apply{g}{i} = \preimg{\piIndex{X}{i}}{\intervalopen{x_0(i) - \epsilon}{x_0(i) + \epsilon}}$ by \cref{pair_eq_iff}.
                            Thus $y\in\preimg{\piIndex{X}{i}}{\intervalopen{x_0(i) - \epsilon}{x_0(i) + \epsilon}}$.
                            Take $z\in\intervalopen{x_0(i) - \epsilon}{x_0(i) + \epsilon}$ such that
                                $y\mathrel{\piIndex{X}{i}} z$ by \cref{preimg_iff}.
                            Thus $(y,z)\in\piIndex{X}{i}$.
                            Take $y_0\in\productFamily{X}$ such that $(y,z) = (y_0,\apply{y_0}{i})$ by \cref{projection_map_indexed}.
                            Thus $y = y_0$ by \cref{pair_eq_iff}.
                            Thus $y\in\productFamily{X}$.
                            Thus $y\in X_0$ by \cref{subseteq}.
                            Take $m\in\reals$ such that $((x_0,y),m)\in d_1$ and
                                for all $j\in J$ we have $\abs{x_0(j) - y(j)} \leq m$ by \cref{square_metric_value_exists}.
                            Show that $m < \epsilon$.
                            \begin{subproof}
                                Show that for all $t\in\coordDiffSet{n}{x_0}{y}$ we have $t < \epsilon$.
                                \begin{subproof}
                                    Fix $t$.
                                    Assume $t\in\coordDiffSet{n}{x_0}{y}$.
                                    Take $i_0\in J$ such that $t = \abs{x_0(i_0) - y(i_0)}$ by \cref{coord_diff_set}.
                                    Thus $(i_0,\apply{g}{i_0})\in g$ by \cref{function_apply_elim}.
                                    Thus $\apply{g}{i_0}\in F$ by \cref{img_iff}.
                                    Thus $y\in\apply{g}{i_0}$ by \cref{inters_iff_forall}.
                                    Thus $y\in\preimg{\piIndex{X}{i_0}}{\intervalopen{x_0(i_0) - \epsilon}{x_0(i_0) + \epsilon}}$
                                        by \cref{function_apply_intro,pair_eq_iff}.
                                    Take $z\in\intervalopen{x_0(i_0) - \epsilon}{x_0(i_0) + \epsilon}$ such that
                                        $y\mathrel{\piIndex{X}{i_0}} z$ by \cref{preimg_iff}.
                                    Thus $(y,z)\in\piIndex{X}{i_0}$.
                                    Take $y_0\in\productFamily{X}$ such that $(y,z) = (y_0,\apply{y_0}{i_0})$ by \cref{projection_map_indexed}.
                                    Thus $y = y_0$ and $z = \apply{y}{i_0}$ by \cref{pair_eq_iff}.
                                    Thus $x_0(i_0) - \epsilon < y(i_0)$ and $y(i_0) < x_0(i_0) + \epsilon$ by \cref{intervalopen}.
                                    $x_0\in\realsPower{J}$.
                                    $y\in\realsPower{J}$.
                                    Thus $x_0\in\funs{J}{\reals}$ and $y\in\funs{J}{\reals}$ by \cref{reals_power,tuple_power}.
                                    Thus $x_0(i_0)\in\reals$ and $y(i_0)\in\reals$ by \cref{funs_type_apply}.
                                    $\epsilon\in\reals$.
                                    $0 < \epsilon$.
                                    Thus $\abs{x_0(i_0) - y(i_0)} < \epsilon$ by \cref{real_abs_lt_from_interval}.
                                    Thus $t < \epsilon$.
                                \end{subproof}
                                Thus $m\in\coordDiffSet{n}{x_0}{y}$ by \cref{square_metric,pair_eq_iff}.
                                Thus $m < \epsilon$.
                            \end{subproof}
                            $d_1$ is a function by \cref{square_metric_function}.
                            Thus $d_1((x_0,y)) = m$ by \cref{function_apply_intro}.
                            Thus $d_1((x_0,y)) < \epsilon$ by \cref{real_lt_subst_left}.
                            Thus $y\in B_2$ by \cref{metric_ball}.
                        \end{subproof}
                        Thus $\inters{F}\subseteq B_2$ by \cref{subseteq}.
                    \end{subproof}
                    Thus $B_2 = \inters{F}$ by \cref{subseteq_antisymmetric}.
                \end{subproof}
                Show that $B_2\in B_0$.
                \begin{subproof}
                    Show that $B_2\in\pow{\unions{S_0}}$.
                    \begin{subproof}
                        Show that for all $y$ such that $y\in B_2$ we have $y\in\unions{S_0}$.
                        \begin{subproof}
                            Fix $y$.
                            Assume $y\in B_2$.
                            Thus $y\in\inters{F}$.
                            Thus there exists $A\in F$ such that $y\in A$ by \cref{inters_iff_forall}.
                            Thus $y\in\unions{F}$ by \cref{unions_iff}.
                            Thus $y\in\unions{S_0}$ by \cref{unions_iff,subseteq}.
                        \end{subproof}
                        Thus $B_2\subseteq\unions{S_0}$ by \cref{subseteq}.
                        Thus $B_2\in\pow{\unions{S_0}}$ by \cref{powerset_intro}.
                    \end{subproof}
                    Thus $B_2\in B_0$ by \cref{finite_intersections}.
                \end{subproof}
                Thus $x\in B_2$ and $B_2\subseteq B_2$.
            \end{subproof}
            Thus $T_0$ is finer than $T_S$ on $X_0$ by \cref{basis_finer_criterion}.
        \end{subproof}
        Show that $T_S$ is finer than $T_0$ on $X_0$.
        \begin{subproof}
            Show that for all $x,B_2$ such that $x\in X_0$ and $B_2\in B_0$ and $x\in B_2$
                there exists $B_3\in B_S$ such that $x\in B_3$ and $B_3\subseteq B_2$.
            \begin{subproof}
                Fix $x$.
                Fix $B_2$.
                Assume $x\in X_0$ and $B_2\in B_0$ and $x\in B_2$.
                Take $F\subseteq S_0$ such that $F$ is finite and inhabited and $B_2 = \inters{F}$ by \cref{finite_intersections}.
                Let $R = \{w\in F\times\reals \mid \exists \beta\in J.\exists U\in\standardTopologyR.
                    \fst{w} = \preimg{\piIndex{X}{\beta}}{U} \land \apply{x}{\beta}\in U \land
                    0 < \snd{w} \land \metricBall{\standardMetricR}{\reals}{\apply{x}{\beta}}{\snd{w}} \subseteq U\}$.
                Show that $R$ is a relation.
                \begin{subproof}
                    Show that for all $w$ such that $w\in R$ there exist $y,z$ such that $w = (y,z)$.
                    \begin{subproof}
                        Fix $w$.
                        Assume $w\in R$.
                        Thus $w\in F\times\reals$.
                        Take $y,z$ such that $y\in F$ and $z\in\reals$ and $w = (y,z)$ by \cref{times_elem_is_tuple}.
                    \end{subproof}
                    Thus $R$ is a relation by \cref{relation}.
                \end{subproof}
                Show that for all $A\in F$ there exists $\epsilon\in\reals$ such that $A\mathrel{R}\epsilon$.
                \begin{subproof}
                    Fix $A$.
                    Assume $A\in F$.
                    Thus $A\in S_0$ by \cref{subseteq}.
                    Thus $x\in B_2$.
                    Thus $x\in\inters{F}$.
                    Thus $x\in A$ by \cref{inters_iff_forall}.
                    Take $\beta\in\dom{X}$ such that $A\in\productSubbasisAt{T}{X}{\beta}$ by \cref{product_subbasis_union_indexed}.
                    Take $U\in\apply{T}{\beta}$ such that $A = \preimg{\piIndex{X}{\beta}}{U}$ by \cref{product_subbasis_indexed}.
                    Thus $\beta\in J$ by \cref{subseteq}.
                    Thus $x\in\preimg{\piIndex{X}{\beta}}{U}$.
                    Take $y\in U$ such that $x\mathrel{\piIndex{X}{\beta}} y$ by \cref{preimg_iff}.
                    Thus $(x,y)\in\piIndex{X}{\beta}$.
                    Take $x_1\in\productFamily{X}$ such that $(x,y) = (x_1,\apply{x_1}{\beta})$ by \cref{projection_map_indexed}.
                    Thus $y = \apply{x}{\beta}$ by \cref{pair_eq_iff}.
                    Thus $\apply{x}{\beta}\in U$.
                    $\apply{T}{\beta} = \standardTopologyR$ by \cref{function_apply_intro}.
                    Thus $U\in\standardTopologyR$.
                    $\standardBasisR$ is a basis on $\reals$ by \cref{standard_basis_is_basis}.
                    Thus $\standardTopologyR$ is a topology on $\reals$ by \cref{basis_topology_is_topology,standard_topology_reals}.
                    Thus $\standardTopologyR\subseteq\pow{\reals}$ by \cref{topology_on}.
                    Thus $U\in\pow{\reals}$ by \cref{subseteq}.
                    Thus $U\subseteq\reals$ by \cref{pow_iff}.
                    Thus $U\in\metricTopology{\standardMetricR}{\reals}$ by \cref{standard_metric_topology}.
                    $\standardMetricR$ is a metric on $\reals$ by \cref{standard_metric_is_metric}.
                    Thus $U$ is open in $\reals$ with respect to $\metricTopology{\standardMetricR}{\reals}$ by \cref{open_in,basis_topology_is_topology,metric_basis_is_basis,metric_topology}.
                    Thus there exists $\epsilon\in\reals$ such that $0 < \epsilon$ and
                        $\metricBall{\standardMetricR}{\reals}{\apply{x}{\beta}}{\epsilon}\subseteq U$ by \cref{metric_open_iff}.
                    Take $\epsilon\in\reals$ such that $0 < \epsilon$ and
                        $\metricBall{\standardMetricR}{\reals}{\apply{x}{\beta}}{\epsilon}\subseteq U$.
                    Let $w = (A,\epsilon)$.
                    $w\in F\times\reals$ by \cref{times_tuple_intro}.
                    $\fst{w} = A$ and $\snd{w} = \epsilon$ by \cref{fst_eq,snd_eq}.
                    Thus $\fst{w} = \preimg{\piIndex{X}{\beta}}{U}$.
                    Thus $\apply{x}{\beta}\in U$.
                    Thus $0 < \snd{w}$.
                    Thus $\metricBall{\standardMetricR}{\reals}{\apply{x}{\beta}}{\snd{w}}\subseteq U$.
                    Thus $w\in R$.
                    Thus $(A,\epsilon)\in R$ by \cref{pair_eq_iff}.
                    Thus $A\mathrel{R}\epsilon$.
                \end{subproof}
                There exists $e$ such that $e$ is a function on $F$ and for all $A\in F$ we have $A\mathrel{R}\apply{e}{A}$ by \cref{choice_function}.
                Take $e$ such that $e$ is a function on $F$ and for all $A\in F$ we have $A\mathrel{R}\apply{e}{A}$.
                Let $E = \img{e}{F}$.
                Thus $E$ is finite by \cref{finite_image_function}.
                Show that $E$ is inhabited.
                \begin{subproof}
                    Take $A$ such that $A\in F$.
                    Thus $(A,\apply{e}{A})\in e$ by \cref{function_apply_elim}.
                    Thus $\apply{e}{A}\in E$ by \cref{img_iff}.
                    Thus $E\neq\emptyset$ by \cref{emptyset}.
                \end{subproof}
                Thus $E\neq\emptyset$ by \cref{emptyset}.
                Show that for all $t$ such that $t\in E$ we have $t\in\reals$ and $0 < t$.
                \begin{subproof}
                    Fix $t$.
                    Assume $t\in E$.
                    Take $A\in F$ such that $(A,t)\in e$ by \cref{img_iff}.
                    Thus $A\mathrel{R}t$.
                    Thus $t\in\reals$ and $0 < t$.
                \end{subproof}
                Thus $E\subseteq\reals$ by \cref{subseteq}.
                Take $\delta\in E$ such that $\delta$ is a minimum of $E$ by \cref{finite_real_minimum}.
                Show that $0 < \delta$.
                \begin{subproof}
                    Thus $\delta\in E$.
                    Thus $0 < \delta$ by \cref{subseteq}.
                \end{subproof}
                Let $B_3 = \metricBall{d_1}{X_0}{x}{\delta}$.
                Show that $B_3\in B_S$.
                \begin{subproof}
                    $x\in X_0$.
                    Thus $\delta\in\reals$ by \cref{subseteq}.
                    Thus $0 < \delta$ by \cref{subseteq}.
                    Show that $B_3\subseteq X_0$.
                    \begin{subproof}
                        Show that for all $y$ such that $y\in B_3$ we have $y\in X_0$.
                        \begin{subproof}
                            Fix $y$.
                            Assume $y\in B_3$.
                            Thus $y\in X_0$ by \cref{metric_ball}.
                        \end{subproof}
                        Thus $B_3\subseteq X_0$ by \cref{subseteq}.
                    \end{subproof}
                    Thus $B_3\in\pow{X_0}$ by \cref{powerset_intro}.
                    Thus $B_3\in B_S$ by \cref{metric_basis}.
                \end{subproof}
                Show that $B_3\subseteq B_2$.
                \begin{subproof}
                    Show that for all $y$ such that $y\in B_3$ we have $y\in B_2$.
                    \begin{subproof}
                        Fix $y$.
                        Assume $y\in B_3$.
                        Thus $y\in X_0$ and $d_1((x,y)) < \delta$ by \cref{metric_ball}.
                        Thus $x\in\realsPower{J}$.
                        Thus $y\in\realsPower{J}$.
                        Take $m\in\reals$ such that $((x,y),m)\in d_1$ and
                            for all $j\in J$ we have $\abs{x(j) - y(j)} \leq m$ by \cref{square_metric_value_exists}.
                        $d_1$ is a function by \cref{square_metric_function}.
                        Thus $d_1((x,y)) = m$ by \cref{function_apply_intro}.
                        Thus $m < \delta$ by \cref{real_lt_subst_left}.
                        Show that for all $A\in F$ we have $y\in A$.
                        \begin{subproof}
                            Fix $A$.
                            Assume $A\in F$.
                            Thus $(A,\apply{e}{A})\in R$.
                            Let $w = (A,\apply{e}{A})$.
                            Thus $w\in R$.
                            Take $\beta,U$ such that $\beta\in J$ and $U\in\standardTopologyR$ and
                                $\fst{w} = \preimg{\piIndex{X}{\beta}}{U}$ and $\apply{x}{\beta}\in U$ and
                                $0 < \snd{w}$ and
                                $\metricBall{\standardMetricR}{\reals}{\apply{x}{\beta}}{\snd{w}} \subseteq U$.
                            $\fst{w} = A$ and $\snd{w} = \apply{e}{A}$ by \cref{fst_eq,snd_eq}.
                            Thus $A = \preimg{\piIndex{X}{\beta}}{U}$.
                            Thus $0 < \apply{e}{A}$.
                            Thus $\metricBall{\standardMetricR}{\reals}{\apply{x}{\beta}}{\apply{e}{A}} \subseteq U$.
                            Thus $(A,\apply{e}{A})\in e$ by \cref{function_apply_elim}.
                            Thus $\apply{e}{A}\in E$ by \cref{img_iff}.
                            Thus $\delta \leq \apply{e}{A}$ by \cref{real_minimum}.
                            Show that $m < \apply{e}{A}$.
                            \begin{subproof}
                                $m < \delta$.
                                $\delta < \apply{e}{A}$ or $\delta = \apply{e}{A}$ by \cref{real_leq_split}.
                                \begin{byCase}
                                \caseOf{$\delta < \apply{e}{A}$.} Thus $m < \apply{e}{A}$ by \cref{real_lt_trans}.
                                \caseOf{$\delta = \apply{e}{A}$.} Thus $m < \apply{e}{A}$ by \cref{real_lt_subst_right}.
                                \end{byCase}
                            \end{subproof}
                            Show that $y\in A$.
                            \begin{subproof}
                                $x\in\realsPower{J}$ and $y\in\realsPower{J}$.
                                Thus $x\in\funs{J}{\reals}$ and $y\in\funs{J}{\reals}$ by \cref{reals_power,tuple_power}.
                                Thus $x(\beta)\in\reals$ and $\apply{y}{\beta}\in\reals$ by \cref{funs_type_apply}.
                                $\neg{\apply{y}{\beta}}\in\reals$ by \cref{real_add_inverse}.
                                $x(\beta) - \apply{y}{\beta} = x(\beta) + \neg{\apply{y}{\beta}}$ by \cref{real_sub}.
                                Thus $x(\beta) - \apply{y}{\beta}\in\reals$ by \cref{real_add_closed}.
                                Thus $\abs{x(\beta) - \apply{y}{\beta}}\in\reals$ by \cref{real_abs_in_reals}.
                                Thus $\abs{x(\beta) - \apply{y}{\beta}} \leq m$.
                                Show that $\abs{x(\beta) - \apply{y}{\beta}} < \apply{e}{A}$.
                                \begin{subproof}
                                    $\abs{x(\beta) - \apply{y}{\beta}} < m$ or $\abs{x(\beta) - \apply{y}{\beta}} = m$ by \cref{real_leq_split}.
                                    \begin{byCase}
                                    \caseOf{$\abs{x(\beta) - \apply{y}{\beta}} < m$.}
                                        Thus $\abs{x(\beta) - \apply{y}{\beta}} < \apply{e}{A}$ by \cref{real_lt_trans}.
                                    \caseOf{$\abs{x(\beta) - \apply{y}{\beta}} = m$.}
                                        Thus $\abs{x(\beta) - \apply{y}{\beta}} < \apply{e}{A}$ by \cref{real_lt_subst_left}.
                                    \end{byCase}
                                \end{subproof}
                                Let $p = (\apply{x}{\beta},\apply{y}{\beta})$.
                                $p\in\reals\times\reals$ by \cref{times_tuple_intro}.
                                $\fst{p} = \apply{x}{\beta}$ and $\snd{p} = \apply{y}{\beta}$ by \cref{fst_eq,snd_eq}.
                                Thus $\abs{\fst{p} - \snd{p}}\in\reals$ by \cref{real_abs_in_reals}.
                                Let $w = (p,\abs{\fst{p} - \snd{p}})$.
                                $w\in\standardMetricR$ by \cref{standard_metric_r}.
                                $\standardMetricR$ is a metric on $\reals$ by \cref{standard_metric_is_metric}.
                                Thus $\standardMetricR\in\funs{\reals\times\reals}{\reals}$ by \cref{metric_on}.
                                Thus $\standardMetricR$ is a function by \cref{funs_elim}.
                                Thus $\apply{\standardMetricR}{p} = \abs{\apply{x}{\beta} - \apply{y}{\beta}}$ by \cref{function_apply_intro}.
                                Thus $\apply{\standardMetricR}{p} < \apply{e}{A}$ by \cref{real_lt_subst_left}.
                                Thus $\apply{y}{\beta}\in\metricBall{\standardMetricR}{\reals}{\apply{x}{\beta}}{\apply{e}{A}}$ by \cref{metric_ball}.
                                Thus $\apply{y}{\beta}\in U$ by \cref{subseteq}.
                                Thus $y\in X_0$ by \cref{metric_ball}.
                                Thus $\productFamily{X} = \realsPower{J}$ by \cref{product_family_reals_power}.
                                Thus $y\in\productFamily{X}$ by \cref{subseteq}.
                                Thus $(y,\apply{y}{\beta})\in\piIndex{X}{\beta}$ by \cref{projection_map_indexed}.
                                Thus $y\in\preimg{\piIndex{X}{\beta}}{U}$ by \cref{preimg_iff}.
                                Thus $y\in A$.
                            \end{subproof}
                        \end{subproof}
                        Thus $y\in\inters{F}$ by \cref{inters_iff_forall}.
                        Thus $y\in B_2$ by \cref{subseteq_antisymmetric}.
                    \end{subproof}
                    Thus $B_3\subseteq B_2$ by \cref{subseteq}.
                \end{subproof}
                Thus $x\in B_3$ and $B_3\subseteq B_2$ by \cref{metric_ball,metric_zero_characterization,metric_on}.
            \end{subproof}
            Thus $T_S$ is finer than $T_0$ on $X_0$ by \cref{basis_finer_criterion}.
        \end{subproof}
        Thus $T_0\subseteq T_S$ by \cref{finer_topology}.
        Thus $T_S\subseteq T_0$ by \cref{finer_topology}.
        Thus $T_S = T_0$ by \cref{subseteq_antisymmetric}.
    \end{subproof}
    Thus $T_E = T_0$.
\end{proof}

% from §20 helper lemma 10: uniform coordinate metric set
\begin{definition}\label{uniform_metric_set}
    $\uniformMetricSet{J}{x}{y} = \{\apply{\boundedMetric{\standardMetricR}}{(x(\alpha),y(\alpha))} \mid \alpha\in J\}$.
\end{definition}

% from §20 helper lemma 11: supremum predicate
\begin{definition}\label{real_supremum_pred}
    $\realSupremum{p}{A}$ iff $p$ is a supremum of $A$.
\end{definition}

% from §20 Item 27: uniform metric
\begin{definition}\label{uniform_metric}
    $\uniformMetric{J} = \{w \mid \exists p\in\realsPower{J}\times\realsPower{J}.
        \exists x\in\realsPower{J}. \exists y\in\realsPower{J}. \exists r\in\reals.
        w = (p,r) \land p = (x,y) \land
        \realSupremum{r}{\uniformMetricSet{J}{x}{y}}\}$.
\end{definition}

% from §20 Item 28: uniform topology
\begin{definition}\label{uniform_topology}
    $\uniformTopology{J} = \metricTopology{\uniformMetric{J}}{\realsPower{J}}$.
\end{definition}

% from §20 Item 29: uniform topology is finer than product topology
\begin{theorem}\label{uniform_finer_product}
    Let $X = \{(\alpha,\reals) \mid \alpha\in J\}$.
    Let $T = \{(\alpha,\standardTopologyR) \mid \alpha\in J\}$.
    Let $T_0 = \productTopologyIndexed{T}{X}$.
    Then $\uniformTopology{J}$ is finer than $T_0$ on $\realsPower{J}$.
\end{theorem}
\begin{proof}
    Omitted. % do not touch
\end{proof}


% from §20 Item 30: uniform topology is coarser than box topology
\begin{theorem}\label{uniform_coarser_box}
    Let $X = \{(\alpha,\reals) \mid \alpha\in J\}$.
    Let $T = \{(\alpha,\standardTopologyR) \mid \alpha\in J\}$.
    Let $T_0 = \boxTopology{T}{X}$.
    Then $\uniformTopology{J}$ is coarser than $T_0$ on $\realsPower{J}$.
\end{theorem}
\begin{proof}
    Omitted. % do not touch
\end{proof}

% from §20 helper lemma 12: omega coordinate metric set
\begin{definition}\label{omega_metric_set}
    $\omegaMetricSet{x}{y} = \{\apply{\boundedMetric{\standardMetricR}}{(x(i),y(i))} / i \mid i\in\naturalsPlus\}$.
\end{definition}

% from §20 Item 31: metric on $\reals^{\omega}$
\begin{definition}\label{omega_metric}
    $\omegaMetric = \{w \mid \exists p\in\realsPower{\naturalsPlus}\times\realsPower{\naturalsPlus}.
        \exists x\in\realsPower{\naturalsPlus}. \exists y\in\realsPower{\naturalsPlus}. \exists r\in\reals.
        w = (p,r) \land p = (x,y) \land
        \realSupremum{r}{\omegaMetricSet{x}{y}}\}$.
\end{definition}

% from §20 Item 32: $\omega$-metric is a metric
\begin{theorem}\label{omega_metric_is_metric}
    Then $\omegaMetric$ is a metric on $\realsPower{\naturalsPlus}$.
\end{theorem}
\begin{proof}
    Omitted. % do not touch
\end{proof}

% from §20 Item 33: $\omega$-metric induces the product topology
\begin{theorem}\label{omega_metric_product_topology}
    Let $X = \{(i,\reals) \mid i\in\naturalsPlus\}$.
    Let $T = \{(i,\standardTopologyR) \mid i\in\naturalsPlus\}$.
    Let $T_0 = \productTopologyIndexed{T}{X}$.
    Then $\metricTopology{\omegaMetric}{\realsPower{\naturalsPlus}} = T_0$.
\end{theorem}
\begin{proof}
    Omitted.  % do not touch
\end{proof}
 
% from §21 helper lemma 1: positive reciprocal of naturalsplus
\begin{lemma}\label{positive_reciprocal_naturalsplus}
    Suppose $n\in\naturalsPlus$.
    Then $0 < 1 / n$.
\end{lemma}
\begin{proof}
    $1\in\naturalsPlus$ by \cref{real_one_in_naturals}.
    Thus $1 \leq n$ by \cref{real_one_leq_naturalsplus}.
    $0 < 1$ by \cref{real_zero_lt_one}.
    $0\in\reals$ by \cref{real_add_ident}.
    $1\in\reals$ by \cref{real_naturals_subset_reals,real_one_in_naturals,subseteq}.
    $n\in\reals$ by \cref{real_naturals_subset_reals,subseteq}.
    Show that $0 < n$.
    \begin{subproof}
        $1 < n$ or $1 = n$ by \cref{real_leq_split}.
        \begin{byCase}
        \caseOf{$1 = n$.} Thus $0 < n$ by \cref{real_lt_subst_right}.
        \caseOf{$1 < n$.} Thus $0 < n$ by \cref{real_lt_trans}.
        \end{byCase}
    \end{subproof}
    Thus $n$ is positive.
    Thus $0 < \inv{n}$ by \cref{real_inv_pos}.
    Thus $n \neq 0$ by \cref{real_pos_nonzero}.
    Thus $\inv{n}\in\reals$ by \cref{real_mul_inverse}.
    $1 / n = 1 \rmul \inv{n}$ by \cref{real_div}.
    $1 \rmul \inv{n} = \inv{n}$ by \cref{real_mul_ident}.
    Thus $\inv{n} = 1 / n$.
    Thus $0 < 1 / n$ by \cref{real_lt_subst_right}.
\end{proof}

% from §21 helper lemma 2: metric ball monotonic in radius
\begin{lemma}\label{metric_ball_mono}
    Assume $d$ is a metric on $X$.
    Suppose $x\in X$.
    Suppose $r_1\in\reals$.
    Suppose $r_2\in\reals$.
    Suppose $r_1 < r_2$.
    Then $\metricBall{d}{X}{x}{r_1} \subseteq \metricBall{d}{X}{x}{r_2}$.
\end{lemma}
\begin{proof}
    Show that for all $y$ such that $y\in\metricBall{d}{X}{x}{r_1}$ we have $y\in\metricBall{d}{X}{x}{r_2}$.
    \begin{subproof}
        Fix $y$.
        Assume $y\in\metricBall{d}{X}{x}{r_1}$.
        Thus $y\in X$ and $d((x,y)) < r_1$ by \cref{metric_ball}.
        $d\in\funs{X\times X}{\reals}$ by \cref{metric_on}.
        $(x,y)\in X\times X$ by \cref{times_tuple_intro}.
        Thus $d((x,y))\in\reals$ by \cref{funs_type_apply}.
        Thus $d((x,y)) < r_2$ by \cref{real_lt_trans}.
        Thus $y\in\metricBall{d}{X}{x}{r_2}$ by \cref{metric_ball}.
    \end{subproof}
    Thus $\metricBall{d}{X}{x}{r_1} \subseteq \metricBall{d}{X}{x}{r_2}$ by \cref{subseteq}.
\end{proof}

% from §21 Item 1: epsilon-delta characterization of continuity in metric spaces
\begin{theorem}\label{metric_continuity_eps_delta}
    Suppose $d_X$ is a metric on $X$.
    Suppose $d_Y$ is a metric on $Y$.
    Let $T_X = \metricTopology{d_X}{X}$.
    Let $T_Y = \metricTopology{d_Y}{Y}$.
    Suppose $f$ is a function from $X$ to $Y$.
    Then $f$ is continuous from $X$ to $Y$ with respect to $T_X$ and $T_Y$ iff
        for all $x\in X$ we have
        for all $\epsilon\in\reals$ such that $0 < \epsilon$
        there exists $\delta\in\reals$ such that $0 < \delta$ and
        for all $y\in X$ if $d_X((x,y)) < \delta$ then $d_Y((f(x),f(y))) < \epsilon$.
\end{theorem}
\begin{proof}
    Show that if $f$ is continuous from $X$ to $Y$ with respect to $T_X$ and $T_Y$ then
        for all $x\in X$ we have
        for all $\epsilon\in\reals$ such that $0 < \epsilon$
        there exists $\delta\in\reals$ such that $0 < \delta$ and
        for all $y\in X$ if $d_X((x,y)) < \delta$ then $d_Y((f(x),f(y))) < \epsilon$.
    \begin{subproof}
        Assume $f$ is continuous from $X$ to $Y$ with respect to $T_X$ and $T_Y$.
        Fix $x\in X$.
        Fix $\epsilon\in\reals$.
        Assume $0 < \epsilon$.
        Let $V = \metricBall{d_Y}{Y}{f(x)}{\epsilon}$.
        $f\in\funs{X}{Y}$.
        Show that $V$ is open in $Y$ with respect to $T_Y$.
        \begin{subproof}
            Show that $V\subseteq Y$.
            \begin{subproof}
                Show that for all $y$ such that $y\in V$ we have $y\in Y$.
                \begin{subproof}
                    Fix $y$.
                    Assume $y\in V$.
                    Thus $y\in Y$ by \cref{metric_ball}.
                \end{subproof}
            Thus $V\subseteq Y$ by \cref{subseteq}.
        \end{subproof}
        Thus $V\in\pow{Y}$ by \cref{powerset_intro}.
        $f(x)\in Y$ by \cref{funs_type_apply}.
        $V\in\metricBasis{d_Y}{Y}$ by \cref{metric_basis}.
            $\metricBasis{d_Y}{Y}$ is a basis on $Y$ by \cref{metric_basis_is_basis}.
            Thus $V\in T_Y$ by \cref{basis_elem_open_in_generated,metric_topology,basis_for_topology}.
            $T_Y$ is a topology on $Y$ by \cref{basis_topology_is_topology,metric_topology,metric_basis_is_basis,basis_for_topology}.
            Thus $V$ is open in $Y$ with respect to $T_Y$ by \cref{open_in}.
        \end{subproof}
        $d_Y$ is a metric on $Y$.
        $f\in\funs{X}{Y}$.
        $f(x)\in Y$ by \cref{funs_type_apply}.
        $d_Y((f(x),f(x))) = 0$ by \cref{metric_zero_characterization,metric_on}.
        Thus $d_Y((f(x),f(x))) < \epsilon$ by \cref{real_lt_subst_left}.
        Thus $f(x)\in V$ by \cref{metric_ball}.
        $x\in\dom{f}$ by \cref{funs}.
        $f$ is a function by \cref{funs_is_function}.
        Thus $(x,f(x))\in f$ by \cref{function_apply_iff}.
        Thus $x\in\preimg{f}{V}$ by \cref{preimg_iff}.
        $\preimg{f}{V}$ is open in $X$ with respect to $T_X$ by \cref{continuous}.
        $\preimg{f}{V}\subseteq X$ by \cref{preimg_subseteq_dom,funs}.
        There exists $\delta\in\reals$ such that $0 < \delta$ and
            $\metricBall{d_X}{X}{x}{\delta} \subseteq \preimg{f}{V}$
            by \cref{metric_open_iff}.
        Take $\delta\in\reals$ such that $0 < \delta$ and
            $\metricBall{d_X}{X}{x}{\delta} \subseteq \preimg{f}{V}$.
        Show that for all $y\in X$ if $d_X((x,y)) < \delta$ then $d_Y((f(x),f(y))) < \epsilon$.
        \begin{subproof}
            Fix $y\in X$.
            Assume $d_X((x,y)) < \delta$.
            Thus $y\in\metricBall{d_X}{X}{x}{\delta}$ by \cref{metric_ball}.
            Thus $y\in\preimg{f}{V}$ by \cref{subseteq}.
            Take $z\in V$ such that $(y,z)\in f$ by \cref{preimg_iff}.
            $f$ is a function by \cref{funs_is_function}.
            Thus $z = f(y)$ by \cref{function_apply_iff}.
            Thus $f(y)\in V$.
            Thus $d_Y((f(x),f(y))) < \epsilon$ by \cref{metric_ball}.
        \end{subproof}
        Thus there exists $\delta\in\reals$ such that $0 < \delta$ and
        for all $y\in X$ if $d_X((x,y)) < \delta$ then $d_Y((f(x),f(y))) < \epsilon$.
    \end{subproof}
    Show that if
        for all $x\in X$ we have
        for all $\epsilon\in\reals$ such that $0 < \epsilon$
        there exists $\delta\in\reals$ such that $0 < \delta$ and
        for all $y\in X$ if $d_X((x,y)) < \delta$ then $d_Y((f(x),f(y))) < \epsilon$
        then $f$ is continuous from $X$ to $Y$ with respect to $T_X$ and $T_Y$.
    \begin{subproof}
        Assume for all $x\in X$ we have
            for all $\epsilon\in\reals$ such that $0 < \epsilon$
            there exists $\delta\in\reals$ such that $0 < \delta$ and
            for all $y\in X$ if $d_X((x,y)) < \delta$ then $d_Y((f(x),f(y))) < \epsilon$.
        Show that for all $V$ such that $V$ is open in $Y$ with respect to $T_Y$ we have
            $\preimg{f}{V}$ is open in $X$ with respect to $T_X$.
        \begin{subproof}
            Fix $V$.
            Assume $V$ is open in $Y$ with respect to $T_Y$.
            $T_Y$ is a topology on $Y$ by \cref{basis_topology_is_topology,metric_topology,metric_basis_is_basis,basis_for_topology}.
            $V\in T_Y$ by \cref{open_in}.
            $T_Y\subseteq\pow{Y}$ by \cref{topology_on}.
            Thus $V\in\pow{Y}$ by \cref{subseteq}.
            Thus $V\subseteq Y$ by \cref{pow_iff}.
            Show that for all $x\in\preimg{f}{V}$ there exists $\delta\in\reals$ such that
                $0 < \delta$ and $\metricBall{d_X}{X}{x}{\delta} \subseteq \preimg{f}{V}$.
            \begin{subproof}
                Fix $x\in\preimg{f}{V}$.
                Take $z\in V$ such that $(x,z)\in f$ by \cref{preimg_iff}.
                $f$ is a function by \cref{funs_is_function}.
                Thus $x\in\dom{f}$ by \cref{preimg_subseteq_dom,elem_subseteq}.
                Thus $z = f(x)$ by \cref{function_apply_iff}.
                Thus $f(x)\in V$.
                Thus $f(x)\in Y$ by \cref{subseteq}.
                There exists $\epsilon\in\reals$ such that $0 < \epsilon$ and
                    $\metricBall{d_Y}{Y}{f(x)}{\epsilon} \subseteq V$ by \cref{metric_open_iff}.
                Take $\epsilon\in\reals$ such that $0 < \epsilon$ and
                    $\metricBall{d_Y}{Y}{f(x)}{\epsilon} \subseteq V$.
                There exists $\delta\in\reals$ such that $0 < \delta$ and
                    for all $y\in X$ if $d_X((x,y)) < \delta$ then $d_Y((f(x),f(y))) < \epsilon$.
                Take $\delta\in\reals$ such that $0 < \delta$ and
                    for all $y\in X$ if $d_X((x,y)) < \delta$ then $d_Y((f(x),f(y))) < \epsilon$.
                Show that $\metricBall{d_X}{X}{x}{\delta} \subseteq \preimg{f}{V}$.
                \begin{subproof}
                    Show that for all $y$ such that $y\in\metricBall{d_X}{X}{x}{\delta}$ we have $y\in\preimg{f}{V}$.
                    \begin{subproof}
                        Fix $y$.
                        Assume $y\in\metricBall{d_X}{X}{x}{\delta}$.
                        Thus $y\in X$ and $d_X((x,y)) < \delta$ by \cref{metric_ball}.
                        Thus $d_Y((f(x),f(y))) < \epsilon$.
                        $f(y)\in Y$ by \cref{funs_type_apply}.
                        Thus $f(y)\in\metricBall{d_Y}{Y}{f(x)}{\epsilon}$ by \cref{metric_ball}.
                        Thus $f(y)\in V$ by \cref{subseteq}.
                        $f$ is a function by \cref{funs_is_function}.
                        $f$ is a relation by \cref{funs_is_relation}.
                        $f$ is right-unique by \cref{function_rightunique}.
                        $y\in\dom{f}$ by \cref{funs}.
                        Thus $(y,f(y))\in f$ by \cref{function_apply_iff}.
                        Thus $y\in\preimg{f}{V}$ by \cref{preimg_iff}.
                    \end{subproof}
                    Thus $\metricBall{d_X}{X}{x}{\delta} \subseteq \preimg{f}{V}$ by \cref{subseteq}.
                \end{subproof}
            \end{subproof}
            $\preimg{f}{V}\subseteq X$ by \cref{preimg_subseteq_dom,funs}.
            Thus $\preimg{f}{V}$ is open in $X$ with respect to $T_X$ by \cref{metric_open_iff}.
        \end{subproof}
        $f\in\funs{X}{Y}$.
        $T_X$ is a topology on $X$ by \cref{basis_topology_is_topology,metric_topology,metric_basis_is_basis,basis_for_topology}.
        $T_Y$ is a topology on $Y$ by \cref{basis_topology_is_topology,metric_topology,metric_basis_is_basis,basis_for_topology}.
        Thus $f$ is continuous from $X$ to $Y$ with respect to $T_X$ and $T_Y$ by \cref{continuous}.
    \end{subproof}
\end{proof}

% from §21 Item 2: sequence implies closure
\begin{lemma}\label{sequence_implies_closure}
    Assume $T$ is a topology on $X$.
    Suppose $A\subseteq X$.
    Suppose $s$ is a sequence in $A$.
    Suppose $x\in X$.
    Suppose $s$ converges to $x$ in $X$ with respect to $T$.
    Then $x\in\closure{A}{X}{T}$.
\end{lemma}
\begin{proof}
    Show that for all $U$ such that $U$ is open in $X$ with respect to $T$ and $x\in U$ we have $U$ intersects $A$.
    \begin{subproof}
        Fix $U$.
        Assume $U$ is open in $X$ with respect to $T$ and $x\in U$.
        Thus $U$ is a neighborhood of $x$ in $X$ with respect to $T$ by \cref{neighborhood}.
        Take $N\in\naturalsPlus$ such that
            for all $n$ such that $n\in\naturalsPlus$ if $N \leq n$ then $s(n)\in U$ by \cref{converges_to}.
        $N + 1\in\naturalsPlus$ by \cref{real_naturals_closed_succ}.
        $0\in\reals$ by \cref{real_add_ident}.
        $1\in\reals$ by \cref{real_naturals_subset_reals,real_one_in_naturals,subseteq}.
        $N\in\reals$ by \cref{real_naturals_subset_reals,subseteq}.
        $0 < 1$ by \cref{real_zero_lt_one}.
        $0 + N < 1 + N$ by \cref{real_lt_add}.
        $0 + N = N$ by \cref{real_add_ident,real_add_comm}.
        $1 + N = N + 1$ by \cref{real_add_comm}.
        Thus $N < N + 1$ by \cref{real_lt_subst_left,real_lt_subst_right}.
        Thus $N \leq N + 1$ by \cref{real_lt_imp_leq}.
        Thus $s(N + 1)\in U$.
        $s\in\funs{\naturalsPlus}{A}$ by \cref{sequence_in}.
        Thus $s(N + 1)\in A$ by \cref{funs_type_apply}.
        Thus $s(N + 1)\in U\inter A$ by \cref{inter_intro}.
        Show that $U\inter A \neq \emptyset$.
        \begin{subproof}
            Suppose not.
            Then $U\inter A = \emptyset$.
            Thus $s(N + 1)\in\emptyset$.
            Contradiction by \cref{emptyset}.
        \end{subproof}
        Thus $U$ intersects $A$ by \cref{intersects}.
    \end{subproof}
    Thus $x\in\closure{A}{X}{T}$ by \cref{closure_iff_intersects_open}.
\end{proof}

% from §21 Item 3: closure implies sequence in a metrizable space
\begin{lemma}\label{closure_implies_sequence_metrizable}
    Assume $T$ is a topology on $X$.
    Suppose $A\subseteq X$.
    Suppose $x\in X$.
    Suppose $X$ is metrizable with respect to $T$.
    Suppose $x\in\closure{A}{X}{T}$.
    Then there exists $s$ such that $s$ is a sequence in $A$ and $s$ converges to $x$ in $X$ with respect to $T$.
\end{lemma}
\begin{proof}
    Take $d$ such that $d$ is a metric on $X$ and $T = \metricTopology{d}{X}$ by \cref{metrizable}.
    Let $b = \{(k,\metricBall{d}{X}{x}{1 / k}) \mid k\in\naturalsPlus\}$.
    Show that $b$ is a function on $\naturalsPlus$.
    \begin{subproof}
        Show that $b$ is a function.
        \begin{subproof}
            Show that for all $k,U_1,U_2$ such that $(k,U_1)\in b$ and $(k,U_2)\in b$ we have $U_1 = U_2$.
            \begin{subproof}
                Fix $k$.
                Fix $U_1$.
                Fix $U_2$.
                Assume $(k,U_1)\in b$ and $(k,U_2)\in b$.
                Thus $U_1 = \metricBall{d}{X}{x}{1 / k}$ and $U_2 = \metricBall{d}{X}{x}{1 / k}$ by \cref{pair_eq_iff}.
                Thus $U_1 = U_2$.
            \end{subproof}
            Thus $b$ is right-unique by \cref{rightunique}.
            Thus $b$ is a function.
        \end{subproof}
        Show that $\dom{b} = \naturalsPlus$.
        \begin{subproof}
            Show that $\naturalsPlus\subseteq\dom{b}$.
            \begin{subproof}
                Show that for all $k$ such that $k\in\naturalsPlus$ we have $k\in\dom{b}$.
                \begin{subproof}
                    Fix $k$.
                    Assume $k\in\naturalsPlus$.
                    Let $U = \metricBall{d}{X}{x}{1 / k}$.
                    $(k,U)\in b$ by \cref{pair_eq_iff}.
                    Thus $k\in\dom{b}$ by \cref{dom_intro}.
                \end{subproof}
                Thus $\naturalsPlus\subseteq\dom{b}$ by \cref{subseteq}.
            \end{subproof}
            Show that $\dom{b}\subseteq\naturalsPlus$.
            \begin{subproof}
                Show that for all $k$ such that $k\in\dom{b}$ we have $k\in\naturalsPlus$.
                \begin{subproof}
                    Fix $k$.
                    Assume $k\in\dom{b}$.
                    Take $U$ such that $(k,U)\in b$ by \cref{dom_iff}.
                    Thus $k\in\naturalsPlus$ by \cref{pair_eq_iff}.
                \end{subproof}
                Thus $\dom{b}\subseteq\naturalsPlus$ by \cref{subseteq}.
            \end{subproof}
            Thus $\dom{b} = \naturalsPlus$ by \cref{subseteq_antisymmetric}.
        \end{subproof}
        Thus $b$ is a function on $\naturalsPlus$.
    \end{subproof}

    Let $F = \{(n,\img{\restrl{b}{\initialSegment{n}}}{\initialSegment{n}}) \mid n\in\naturalsPlus\}$.
    Let $B = \{(n,\inters{\apply{F}{n}}) \mid n\in\naturalsPlus\}$.

    Show that $F$ is a function on $\naturalsPlus$.
    \begin{subproof}
        Show that $F$ is a function.
        \begin{subproof}
            Show that for all $n,U_1,U_2$ such that $(n,U_1)\in F$ and $(n,U_2)\in F$ we have $U_1 = U_2$.
            \begin{subproof}
                Fix $n$.
                Fix $U_1$.
                Fix $U_2$.
                Assume $(n,U_1)\in F$ and $(n,U_2)\in F$.
                Thus $U_1 = \img{\restrl{b}{\initialSegment{n}}}{\initialSegment{n}}$ and
                    $U_2 = \img{\restrl{b}{\initialSegment{n}}}{\initialSegment{n}}$ by \cref{pair_eq_iff}.
                Thus $U_1 = U_2$.
            \end{subproof}
            Thus $F$ is right-unique by \cref{rightunique}.
            Thus $F$ is a function.
        \end{subproof}
        Show that $\dom{F} = \naturalsPlus$.
        \begin{subproof}
            Show that $\naturalsPlus\subseteq\dom{F}$.
            \begin{subproof}
                Show that for all $n$ such that $n\in\naturalsPlus$ we have $n\in\dom{F}$.
                \begin{subproof}
                    Fix $n$.
                    Assume $n\in\naturalsPlus$.
                    Let $U = \img{\restrl{b}{\initialSegment{n}}}{\initialSegment{n}}$.
                    $(n,U)\in F$ by \cref{pair_eq_iff}.
                    Thus $n\in\dom{F}$ by \cref{dom_intro}.
                \end{subproof}
                Thus $\naturalsPlus\subseteq\dom{F}$ by \cref{subseteq}.
            \end{subproof}
            Show that $\dom{F}\subseteq\naturalsPlus$.
            \begin{subproof}
                Show that for all $n$ such that $n\in\dom{F}$ we have $n\in\naturalsPlus$.
                \begin{subproof}
                    Fix $n$.
                    Assume $n\in\dom{F}$.
                    Take $U$ such that $(n,U)\in F$ by \cref{dom_iff}.
                    Thus $n\in\naturalsPlus$ by \cref{pair_eq_iff}.
                \end{subproof}
                Thus $\dom{F}\subseteq\naturalsPlus$ by \cref{subseteq}.
            \end{subproof}
            Thus $\dom{F} = \naturalsPlus$ by \cref{subseteq_antisymmetric}.
        \end{subproof}
        Thus $F$ is a function on $\naturalsPlus$.
    \end{subproof}

    Show that for all $n\in\naturalsPlus$ we have $F(n) = \img{\restrl{b}{\initialSegment{n}}}{\initialSegment{n}}$.
    \begin{subproof}
        Fix $n$.
        Assume $n\in\naturalsPlus$.
        Let $U = \img{\restrl{b}{\initialSegment{n}}}{\initialSegment{n}}$.
        $(n,U)\in F$ by \cref{pair_eq_iff}.
        Thus $F(n) = U$ by \cref{function_apply_intro}.
    \end{subproof}

    Show that $B$ is a function on $\naturalsPlus$.
    \begin{subproof}
        Show that $B$ is a function.
        \begin{subproof}
            Show that for all $n,U_1,U_2$ such that $(n,U_1)\in B$ and $(n,U_2)\in B$ we have $U_1 = U_2$.
            \begin{subproof}
                Fix $n$.
                Fix $U_1$.
                Fix $U_2$.
                Assume $(n,U_1)\in B$ and $(n,U_2)\in B$.
                Thus $U_1 = \inters{F(n)}$ and $U_2 = \inters{F(n)}$ by \cref{pair_eq_iff}.
                Thus $U_1 = U_2$.
            \end{subproof}
            Thus $B$ is right-unique by \cref{rightunique}.
            Thus $B$ is a function.
        \end{subproof}
        Show that $\dom{B} = \naturalsPlus$.
        \begin{subproof}
            Show that $\naturalsPlus\subseteq\dom{B}$.
            \begin{subproof}
                Show that for all $n$ such that $n\in\naturalsPlus$ we have $n\in\dom{B}$.
                \begin{subproof}
                    Fix $n$.
                    Assume $n\in\naturalsPlus$.
                    Let $U = \inters{F(n)}$.
                    $(n,U)\in B$ by \cref{pair_eq_iff}.
                    Thus $n\in\dom{B}$ by \cref{dom_intro}.
                \end{subproof}
                Thus $\naturalsPlus\subseteq\dom{B}$ by \cref{subseteq}.
            \end{subproof}
            Show that $\dom{B}\subseteq\naturalsPlus$.
            \begin{subproof}
                Show that for all $n$ such that $n\in\dom{B}$ we have $n\in\naturalsPlus$.
                \begin{subproof}
                    Fix $n$.
                    Assume $n\in\dom{B}$.
                    Take $U$ such that $(n,U)\in B$ by \cref{dom_iff}.
                    Thus $n\in\naturalsPlus$ by \cref{pair_eq_iff}.
                \end{subproof}
                Thus $\dom{B}\subseteq\naturalsPlus$ by \cref{subseteq}.
            \end{subproof}
            Thus $\dom{B} = \naturalsPlus$ by \cref{subseteq_antisymmetric}.
        \end{subproof}
        Thus $B$ is a function on $\naturalsPlus$.
    \end{subproof}

    Show that for all $n\in\naturalsPlus$ we have $B(n) = \inters{F(n)}$.
    \begin{subproof}
        Fix $n$.
        Assume $n\in\naturalsPlus$.
        Let $U = \inters{F(n)}$.
        $(n,U)\in B$ by \cref{pair_eq_iff}.
        Thus $B(n) = U$ by \cref{function_apply_intro}.
    \end{subproof}

    Show that for all $n\in\naturalsPlus$ we have $B(n)$ is a neighborhood of $x$ in $X$ with respect to $T$.
    \begin{subproof}
        Fix $n\in\naturalsPlus$.
        Show that $F(n)$ is finite and inhabited.
        \begin{subproof}
            $\initialSegment{n}$ is finite by \cref{initial_segment_finite}.
            Show that $\initialSegment{n}\subseteq\dom{b}$.
            \begin{subproof}
                Show that for all $i$ such that $i\in\initialSegment{n}$ we have $i\in\dom{b}$.
                \begin{subproof}
                    Fix $i$.
                    Assume $i\in\initialSegment{n}$.
                    $i\in\naturalsPlus$ by \cref{initial_segment_naturalsplus}.
                    Let $U = \metricBall{d}{X}{x}{1 / i}$.
                    $(i,U)\in b$ by \cref{pair_eq_iff}.
                    Thus $i\in\dom{b}$ by \cref{dom_intro}.
                \end{subproof}
                Thus $\initialSegment{n}\subseteq\dom{b}$ by \cref{subseteq}.
            \end{subproof}
            Let $g = \restrl{b}{\initialSegment{n}}$.
            $g$ is a function by \cref{restrl_of_function_is_function}.
            Show that $\dom{g} = \initialSegment{n}$.
            \begin{subproof}
                Thus $\dom{g}\subseteq\initialSegment{n}$ by \cref{restrl_dom}.
                Show that $\initialSegment{n}\subseteq\dom{g}$.
                \begin{subproof}
                    Show that for all $i$ such that $i\in\initialSegment{n}$ we have $i\in\dom{g}$.
                    \begin{subproof}
                        Fix $i$.
                        Assume $i\in\initialSegment{n}$.
                        $i\in\naturalsPlus$ by \cref{initial_segment_naturalsplus}.
                        Let $U = \metricBall{d}{X}{x}{1 / i}$.
                        $(i,U)\in b$ by \cref{pair_eq_iff}.
                        Thus $(i,U)\in g$ by \cref{restrl_iff}.
                        Thus $i\in\dom{g}$ by \cref{dom_intro}.
                    \end{subproof}
                    Thus $\initialSegment{n}\subseteq\dom{g}$ by \cref{subseteq}.
                \end{subproof}
                Thus $\dom{g} = \initialSegment{n}$ by \cref{subseteq_antisymmetric}.
            \end{subproof}
            Thus $g$ is a function on $\initialSegment{n}$.
            Thus $F(n) = \img{g}{\initialSegment{n}}$.
            Thus $F(n)$ is finite by \cref{finite_image_function}.
            $1\in\naturalsPlus$ by \cref{real_one_in_naturals}.
            $1 \leq n$ by \cref{real_one_leq_naturalsplus}.
            Thus $1\in\initialSegment{n}$ by \cref{initial_segment_naturalsplus}.
            Let $U_1 = \metricBall{d}{X}{x}{1 / 1}$.
            $(1,U_1)\in b$ by \cref{pair_eq_iff}.
            Thus $(1,U_1)\in g$ by \cref{restrl_iff}.
            Thus $1\in\dom{g}$ by \cref{dom_intro}.
            Thus $1\in\dom{g}\inter\initialSegment{n}$ by \cref{inter_intro}.
            $(\restrl{b}{\initialSegment{n}})(1) = b(1)$ by \cref{restrl_of_function_apply}.
            Thus $b(1)\in F(n)$ by \cref{img_of_function_intro}.
            Thus $F(n)$ is inhabited.
        \end{subproof}
        Show that $F(n)\subseteq T$.
        \begin{subproof}
            Show that for all $U$ such that $U\in F(n)$ we have $U\in T$.
            \begin{subproof}
                Fix $U$.
                Assume $U\in F(n)$.
                Thus $U\in\img{\restrl{b}{\initialSegment{n}}}{\initialSegment{n}}$.
                Take $k\in\initialSegment{n}$ such that $k\mathrel{\restrl{b}{\initialSegment{n}}} U$ by \cref{img_iff}.
                Thus $k\mathrel{b} U$ by \cref{restrl_iff}.
                Thus $U = \metricBall{d}{X}{x}{1 / k}$ by \cref{pair_eq_iff}.
                $k\in\naturalsPlus$ by \cref{initial_segment_naturalsplus}.
                $0 < 1 / k$ by \cref{positive_reciprocal_naturalsplus}.
                Show that for all $y$ such that $y\in U$ we have $y\in X$.
                \begin{subproof}
                    Fix $y$.
                    Assume $y\in U$.
                    Thus $y\in X$ by \cref{metric_ball}.
                \end{subproof}
                Thus $U\subseteq X$ by \cref{subseteq}.
                Thus $U\in\pow{X}$ by \cref{powerset_intro}.
                $k\in\reals$ by \cref{real_naturals_subset_reals,subseteq}.
                Show that $0 < k$.
                \begin{subproof}
                    $0\in\reals$ by \cref{real_add_ident}.
                    $1\in\reals$ by \cref{real_naturals_subset_reals,real_one_in_naturals,subseteq}.
                    $1 < k$ or $1 = k$ by \cref{real_one_leq_naturalsplus}.
                    \begin{byCase}
                    \caseOf{$1 = k$.}
                        $0 < 1$ by \cref{real_zero_lt_one}.
                        Thus $0 < k$ by \cref{real_lt_subst_right}.
                    \caseOf{$1 < k$.}
                        $0 < 1$ by \cref{real_zero_lt_one}.
                        Thus $0 < k$ by \cref{real_lt_trans}.
                    \end{byCase}
                \end{subproof}
                Thus $k\neq 0$ by \cref{real_pos_nonzero}.
                Thus $\inv{k}\in\reals$ by \cref{real_mul_inverse}.
                $1\in\reals$ by \cref{real_naturals_subset_reals,real_one_in_naturals,subseteq}.
                Thus $1 / k\in\reals$ by \cref{real_div,real_mul_closed}.
                Thus $U\in\metricBasis{d}{X}$ by \cref{metric_basis}.
                $\metricBasis{d}{X}$ is a basis on $X$ by \cref{metric_basis_is_basis}.
                Thus $U\in T$ by \cref{basis_elem_open_in_generated,metric_topology,basis_for_topology}.
            \end{subproof}
            Thus $F(n)\subseteq T$ by \cref{subseteq}.
        \end{subproof}
        Thus $B(n)\in T$ by \cref{finite_intersections_open_in_topology}.
        Show that $x\in B(n)$.
        \begin{subproof}
            Show that for all $U$ such that $U\in F(n)$ we have $x\in U$.
            \begin{subproof}
                Fix $U$.
                Assume $U\in F(n)$.
                Thus $U\in\img{\restrl{b}{\initialSegment{n}}}{\initialSegment{n}}$.
                Take $k\in\initialSegment{n}$ such that $k\mathrel{\restrl{b}{\initialSegment{n}}} U$ by \cref{img_iff}.
                Thus $k\mathrel{b} U$ by \cref{restrl_iff}.
                Thus $U = \metricBall{d}{X}{x}{1 / k}$ by \cref{pair_eq_iff}.
                $k\in\naturalsPlus$ by \cref{initial_segment_naturalsplus}.
                $0 < 1 / k$ by \cref{positive_reciprocal_naturalsplus}.
                $d((x,x)) = 0$ by \cref{metric_zero_characterization,metric_on}.
                Thus $d((x,x)) < 1 / k$ by \cref{real_lt_subst_left}.
                Thus $x\in U$ by \cref{metric_ball}.
            \end{subproof}
            Thus $x\in B(n)$ by \cref{inters_iff_forall}.
        \end{subproof}
        Thus $B(n)$ is a neighborhood of $x$ in $X$ with respect to $T$ by \cref{neighborhood,open_in}.
    \end{subproof}

    Show that for all $n\in\naturalsPlus$ we have $B(n)$ intersects $A$.
    \begin{subproof}
        Fix $n\in\naturalsPlus$.
        Thus $B(n)$ is a neighborhood of $x$ in $X$ with respect to $T$.
        Thus $B(n)$ is open in $X$ with respect to $T$ and $x\in B(n)$ by \cref{neighborhood}.
        Thus $B(n)$ intersects $A$ by \cref{closure_iff_intersects_open}.
    \end{subproof}

    Let $R = \{w\in\naturalsPlus\times X \mid \snd{w}\in A\inter \apply{B}{\fst{w}}\}$.
    Show that $R$ is a relation.
    \begin{subproof}
        Show that for all $w$ such that $w\in R$ there exist $n,y$ such that $w = (n,y)$.
        \begin{subproof}
            Fix $w$.
            Assume $w\in R$.
            Thus $w\in\naturalsPlus\times X$.
            Take $n,y$ such that $n\in\naturalsPlus$ and $y\in X$ and $w = (n,y)$ by \cref{times_elem_is_tuple}.
        \end{subproof}
        Thus $R$ is a relation by \cref{relation}.
    \end{subproof}
    Show that for all $n\in\naturalsPlus$ there exists $y$ such that $n\mathrel{R} y$.
    \begin{subproof}
        Fix $n\in\naturalsPlus$.
        $B(n)$ intersects $A$.
        Thus $B(n)\inter A \neq \emptyset$ by \cref{intersects}.
        Thus $A\inter B(n) = B(n)\inter A$ by \cref{inter_comm}.
        Thus $A\inter B(n) \neq \emptyset$.
        Take $y$ such that $y\in A\inter B(n)$ by \cref{nonempty_inhabited}.
        Thus $y\in A$ by \cref{inter_elim_left}.
        Thus $y\in X$ by \cref{subseteq}.
        Thus $(n,y)\in\naturalsPlus\times X$ by \cref{times_tuple_intro}.
        $\fst{(n,y)} = n$ by \cref{fst_eq}.
        $\snd{(n,y)} = y$ by \cref{snd_eq}.
        Thus $\snd{(n,y)}\in A\inter \apply{B}{\fst{(n,y)}}$ by \cref{inter_intro}.
        Thus $(n,y)\in R$.
        Thus $n\mathrel{R} y$ by \cref{pair_eq_iff}.
    \end{subproof}
    Take $s$ such that $s$ is a function on $\naturalsPlus$ and
        for all $n\in\naturalsPlus$ we have $n\mathrel{R} s(n)$ by \cref{axiom_of_choice}.

    Show that $s$ is a sequence in $A$.
    \begin{subproof}
        Show that $s$ is a function to $A$.
        \begin{subproof}
            Show that for all $n\in\dom{s}$ we have $s(n)\in A$.
            \begin{subproof}
                Fix $n$.
                Assume $n\in\dom{s}$.
                $\dom{s} = \naturalsPlus$.
                Thus $n\in\naturalsPlus$ by \cref{subseteq}.
                Thus $n\mathrel{R} s(n)$.
                Thus $(n,s(n))\in R$ by \cref{pair_eq_iff}.
                Thus $\snd{(n,s(n))}\in A\inter \apply{B}{\fst{(n,s(n))}}$.
                $\fst{(n,s(n))} = n$ by \cref{fst_eq}.
                $\snd{(n,s(n))} = s(n)$ by \cref{snd_eq}.
                Thus $s(n)\in A\inter B(n)$.
                Thus $s(n)\in A$ by \cref{inter}.
            \end{subproof}
            Thus $s$ is a function to $A$.
        \end{subproof}
        Thus $s\in\funs{\naturalsPlus}{A}$ by \cref{funs_intro}.
        Thus $s$ is a sequence in $A$ by \cref{sequence_in}.
    \end{subproof}

    Show that $s$ converges to $x$ in $X$ with respect to $T$.
    \begin{subproof}
        Show that $s$ is a sequence in $X$.
        \begin{subproof}
            $s$ is a function to $A$.
            Thus $s$ is a function to $X$ by \cref{function_on_weaken_codom,subseteq}.
            $\dom{s} = \naturalsPlus$ by \cref{funs_elim}.
            Thus $s\in\funs{\naturalsPlus}{X}$ by \cref{funs_intro}.
            Thus $s$ is a sequence in $X$ by \cref{sequence_in}.
        \end{subproof}
        Show that for all $U$ such that $U$ is a neighborhood of $x$ in $X$ with respect to $T$
            there exists $N\in\naturalsPlus$ such that for all $n\in\naturalsPlus$ if $N \leq n$ then $s(n)\in U$.
        \begin{subproof}
            Fix $U$.
            Assume $U$ is a neighborhood of $x$ in $X$ with respect to $T$.
            Thus $U$ is open in $X$ with respect to $T$ and $x\in U$ by \cref{neighborhood}.
            Thus $U\in T$ by \cref{open_in}.
            $T\subseteq\pow{X}$ by \cref{topology_on}.
            Thus $U\in\pow{X}$ by \cref{subseteq}.
            Thus $U\subseteq X$ by \cref{powerset_elim}.
            There exists $\epsilon\in\reals$ such that $0 < \epsilon$ and
                $\metricBall{d}{X}{x}{\epsilon} \subseteq U$ by \cref{metric_open_iff}.
            Take $\epsilon\in\reals$ such that $0 < \epsilon$ and
                $\metricBall{d}{X}{x}{\epsilon} \subseteq U$.
            Take $N\in\naturalsPlus$ such that $0 < 1 / N$ and $1 / N < \epsilon$ by \cref{real_small_reciprocal}.
            $N\in\reals$ by \cref{real_naturals_subset_reals,subseteq}.
            Show that $0 < N$.
            \begin{subproof}
                $0\in\reals$ by \cref{real_add_ident}.
                $1\in\reals$ by \cref{real_naturals_subset_reals,real_one_in_naturals,subseteq}.
                $1 < N$ or $1 = N$ by \cref{real_one_leq_naturalsplus}.
                \begin{byCase}
                \caseOf{$1 = N$.}
                    $0 < 1$ by \cref{real_zero_lt_one}.
                    Thus $0 < N$ by \cref{real_lt_subst_right}.
                \caseOf{$1 < N$.}
                    $0 < 1$ by \cref{real_zero_lt_one}.
                    Thus $0 < N$ by \cref{real_lt_trans}.
                \end{byCase}
            \end{subproof}
            Thus $N\neq 0$ by \cref{real_pos_nonzero}.
            Thus $\inv{N}\in\reals$ by \cref{real_mul_inverse}.
            $1\in\reals$ by \cref{real_naturals_subset_reals,real_one_in_naturals,subseteq}.
            Thus $1 / N\in\reals$ by \cref{real_div,real_mul_closed}.
            $\metricBall{d}{X}{x}{1 / N} \subseteq \metricBall{d}{X}{x}{\epsilon}$ by \cref{metric_ball_mono}.
            Thus $\metricBall{d}{X}{x}{1 / N} \subseteq U$ by \cref{subseteq_transitive}.
            Show that for all $n\in\naturalsPlus$ if $N \leq n$ then $s(n)\in U$.
            \begin{subproof}
                Fix $n\in\naturalsPlus$.
                Assume $N \leq n$.
                Show that $B(n)\subseteq \metricBall{d}{X}{x}{1 / N}$.
                \begin{subproof}
                    Show that $F(N)\subseteq F(n)$.
                    \begin{subproof}
                        Show that for all $U_0$ such that $U_0\in F(N)$ we have $U_0\in F(n)$.
                        \begin{subproof}
                            Fix $U_0$.
                            Assume $U_0\in F(N)$.
                            Thus $U_0\in\img{\restrl{b}{\initialSegment{N}}}{\initialSegment{N}}$.
                            Take $k\in\initialSegment{N}$ such that $k\mathrel{\restrl{b}{\initialSegment{N}}} U_0$ by \cref{img_iff}.
                            Thus $k\mathrel{b} U_0$ by \cref{restrl_iff}.
                            $k\in\naturalsPlus$ and $k \leq N$ by \cref{initial_segment_naturalsplus}.
                            $k\in\reals$ by \cref{real_naturals_subset_reals,subseteq}.
                            $N\in\reals$ by \cref{real_naturals_subset_reals,subseteq}.
                            $n\in\reals$ by \cref{real_naturals_subset_reals,subseteq}.
                            Thus $k \leq n$ by \cref{real_leq_trans}.
                            Thus $k\in\initialSegment{n}$ by \cref{initial_segment_naturalsplus}.
                            Thus $(k,U_0)\in\restrl{b}{\initialSegment{n}}$ by \cref{restrl_iff}.
                            Thus $U_0\in F(n)$ by \cref{img_iff}.
                        \end{subproof}
                        Thus $F(N)\subseteq F(n)$ by \cref{subseteq}.
                    \end{subproof}
                    Show that $B(n)\subseteq B(N)$.
                    \begin{subproof}
                        Show that for all $U_0$ such that $U_0\in F(N)$ we have $B(n)\subseteq U_0$.
                        \begin{subproof}
                            Fix $U_0$.
                            Assume $U_0\in F(N)$.
                            Thus $U_0\in F(n)$ by \cref{subseteq}.
                            Thus $B(n)\subseteq U_0$ by \cref{inters_subseteq_elem}.
                        \end{subproof}
                        Thus $F(N)$ is inhabited.
                        Thus $B(n)\subseteq B(N)$ by \cref{subseteq_inters_iff}.
                    \end{subproof}
                    Show that $B(N)\subseteq \metricBall{d}{X}{x}{1 / N}$.
                    \begin{subproof}
                        $N\in\initialSegment{N}$ by \cref{initial_segment_naturalsplus}.
                        Show that $\initialSegment{N}\subseteq\dom{b}$.
                        \begin{subproof}
                            Show that for all $i$ such that $i\in\initialSegment{N}$ we have $i\in\dom{b}$.
                            \begin{subproof}
                                Fix $i$.
                                Assume $i\in\initialSegment{N}$.
                                $i\in\naturalsPlus$ by \cref{initial_segment_naturalsplus}.
                                Let $V = \metricBall{d}{X}{x}{1 / i}$.
                                $(i,V)\in b$ by \cref{pair_eq_iff}.
                                Thus $i\in\dom{b}$ by \cref{dom_intro}.
                            \end{subproof}
                            Thus $\initialSegment{N}\subseteq\dom{b}$ by \cref{subseteq}.
                        \end{subproof}
                        Thus $N\in\dom{b}$ by \cref{subseteq}.
                        Thus $(N,b(N))\in b$ by \cref{function_apply_elim}.
                        Thus $(N,b(N))\in\restrl{b}{\initialSegment{N}}$ by \cref{restrl_iff}.
                        Thus $b(N)\in F(N)$ by \cref{img_iff}.
                        Thus $B(N)\subseteq b(N)$ by \cref{inters_subseteq_elem}.
                        Thus $B(N)\subseteq \metricBall{d}{X}{x}{1 / N}$.
                    \end{subproof}
                    Thus $B(n)\subseteq \metricBall{d}{X}{x}{1 / N}$ by \cref{subseteq_transitive}.
                \end{subproof}
                $n\mathrel{R} s(n)$.
                Thus $(n,s(n))\in R$ by \cref{pair_eq_iff}.
                Thus $\snd{(n,s(n))}\in A\inter \apply{B}{\fst{(n,s(n))}}$.
                $\fst{(n,s(n))} = n$ by \cref{fst_eq}.
                $\snd{(n,s(n))} = s(n)$ by \cref{snd_eq}.
                Thus $s(n)\in A\inter B(n)$.
                Thus $s(n)\in B(n)$ by \cref{inter}.
                Thus $s(n)\in U$ by \cref{subseteq}.
            \end{subproof}
            Thus there exists $N\in\naturalsPlus$ such that for all $n\in\naturalsPlus$ if $N \leq n$ then $s(n)\in U$.
        \end{subproof}
        Thus $s$ converges to $x$ in $X$ with respect to $T$ by \cref{converges_to}.
    \end{subproof}

    Thus there exists $s$ such that $s$ is a sequence in $A$ and $s$ converges to $x$ in $X$ with respect to $T$.
\end{proof}

% from §21 Item 4: continuous implies sequential continuity
\begin{theorem}\label{continuous_implies_sequential}
    Assume $T$ is a topology on $X$.
    Assume $T_1$ is a topology on $Y$.
    Suppose $f$ is a function from $X$ to $Y$.
    Suppose $f$ is continuous from $X$ to $Y$ with respect to $T$ and $T_1$.
    Then for all $s$ such that $s$ is a sequence in $X$ we have
        for all $x\in X$ if $s$ converges to $x$ in $X$ with respect to $T$ then
        $f\circ s$ converges to $f(x)$ in $Y$ with respect to $T_1$.
\end{theorem}
\begin{proof}
    Fix $s$ such that $s$ is a sequence in $X$.
    Fix $x\in X$.
    Assume $s$ converges to $x$ in $X$ with respect to $T$.
    $f\in\funs{X}{Y}$ by \cref{continuous}.
    Thus $f\circ s\in\funs{\naturalsPlus}{Y}$ by \cref{funs_circ,sequence_in}.
    Thus $f\circ s$ is a sequence in $Y$ by \cref{sequence_in}.
    Show that for all $V$ such that $V$ is a neighborhood of $f(x)$ in $Y$ with respect to $T_1$
        there exists $N\in\naturalsPlus$ such that for all $n\in\naturalsPlus$ if $N\leq n$ then
        $(f\circ s)(n)\in V$.
    \begin{subproof}
        Fix $V$.
        Assume $V$ is a neighborhood of $f(x)$ in $Y$ with respect to $T_1$.
        Thus $V$ is open in $Y$ with respect to $T_1$ and $f(x)\in V$ by \cref{neighborhood}.
        Thus $\preimg{f}{V}$ is open in $X$ with respect to $T$ by \cref{continuous}.
        $f\in\funs{X}{Y}$ by \cref{continuous}.
        Thus $f$ is a function by \cref{funs_is_function}.
        $x\in\dom{f}$ by \cref{funs}.
        Thus $(x,f(x))\in f$ by \cref{function_apply_iff}.
        Thus $x\in\preimg{f}{V}$ by \cref{preimg_iff}.
        Thus $\preimg{f}{V}$ is a neighborhood of $x$ in $X$ with respect to $T$ by \cref{neighborhood}.
        Take $N\in\naturalsPlus$ such that for all $n\in\naturalsPlus$ if $N\leq n$ then
            $s(n)\in\preimg{f}{V}$ by \cref{converges_to}.
        Show that for all $n\in\naturalsPlus$ if $N\leq n$ then $(f\circ s)(n)\in V$.
        \begin{subproof}
            Fix $n\in\naturalsPlus$.
            Assume $N\leq n$.
            Thus $s(n)\in\preimg{f}{V}$.
            Take $y\in V$ such that $(s(n),y)\in f$ by \cref{preimg_iff}.
            $f$ is a function by \cref{funs_is_function,continuous}.
            Thus $y = f(s(n))$ by \cref{function_apply_iff}.
            Thus $f(s(n))\in V$.
            $s\in\funs{\naturalsPlus}{X}$ by \cref{sequence_in}.
            Thus $s$ is a function by \cref{funs_is_function}.
            Thus $s\in\rels{\naturalsPlus}{X}$ by \cref{funs_subseteq_rels,subseteq}.
            Thus $\ran{s}\subseteq X$ by \cref{rels_ran_subseteq}.
            $f\in\funs{X}{Y}$ by \cref{continuous}.
            Thus $\dom{f} = X$ by \cref{funs}.
            Thus $\ran{s}\subseteq\dom{f}$ by \cref{subseteq}.
            Thus $f$ is composable with $s$.
            Thus $n\in\dom{s}$ by \cref{funs,subseteq}.
            Thus $(f\circ s)(n) = f(s(n))$ by \cref{circ_apply}.
            Thus $(f\circ s)(n)\in V$.
        \end{subproof}
        Thus there exists $N\in\naturalsPlus$ such that for all $n\in\naturalsPlus$ if $N\leq n$ then
            $(f\circ s)(n)\in V$.
    \end{subproof}
    Thus $f\circ s$ converges to $f(x)$ in $Y$ with respect to $T_1$ by \cref{converges_to}.
\end{proof}

% from §21 Item 5: sequential continuity implies continuity in metrizable spaces
\begin{theorem}\label{sequential_implies_continuous_metrizable}
    Assume $T$ is a topology on $X$.
    Assume $T_1$ is a topology on $Y$.
    Suppose $f$ is a function from $X$ to $Y$.
    Suppose $X$ is metrizable with respect to $T$.
    Suppose for all $s$ such that $s$ is a sequence in $X$ we have
        for all $x\in X$ if $s$ converges to $x$ in $X$ with respect to $T$ then
        $f\circ s$ converges to $f(x)$ in $Y$ with respect to $T_1$.
    Then $f$ is continuous from $X$ to $Y$ with respect to $T$ and $T_1$.
\end{theorem}
\begin{proof}
    Show that for all $A$ such that $A\subseteq X$ we have
        $\img{f}{\closure{A}{X}{T}}\subseteq \closure{\img{f}{A}}{Y}{T_1}$.
    \begin{subproof}
        Fix $A$.
        Assume $A\subseteq X$.
        Show that for all $y$ such that $y\in\img{f}{\closure{A}{X}{T}}$ we have
            $y\in\closure{\img{f}{A}}{Y}{T_1}$.
        \begin{subproof}
            Fix $y$.
            Assume $y\in\img{f}{\closure{A}{X}{T}}$.
            $f$ is a function from $X$ to $Y$.
            Thus $f\in\funs{X}{Y}$.
            Thus $f$ is a function by \cref{funs_is_function}.
            Thus $f$ is right-unique by \cref{function_rightunique}.
            Thus $f$ is a relation by \cref{funs_is_relation}.
            Take $x\in\dom{f}\inter\closure{A}{X}{T}$ such that $y = f(x)$ by \cref{img_of_function_elim}.
            Thus $x\in\closure{A}{X}{T}$ by \cref{inter_elim_right}.
            Thus $x\in X$ by \cref{closure}.
            Take $s$ such that $s$ is a sequence in $A$ and
                $s$ converges to $x$ in $X$ with respect to $T$
                by \cref{closure_implies_sequence_metrizable}.
            $A\subseteq X$.
            Thus $s$ is a sequence in $X$ by \cref{sequence_in,funs_weaken_codom}.
            Thus $f\circ s$ converges to $f(x)$ in $Y$ with respect to $T_1$.
            Show that $f\circ s$ is a sequence in $\img{f}{A}$.
            \begin{subproof}
                $f\in\funs{X}{Y}$.
                Thus $f\circ s\in\funs{\naturalsPlus}{Y}$ by \cref{funs_circ,sequence_in}.
                Thus $f\circ s$ is a function by \cref{funs_is_function}.
                Show that for all $n\in\dom{(f\circ s)}$ we have $(f\circ s)(n)\in\img{f}{A}$.
                \begin{subproof}
                    Fix $n\in\dom{(f\circ s)}$.
                    Thus $n\in\naturalsPlus$ by \cref{funs}.
                    $s\in\funs{\naturalsPlus}{A}$ by \cref{sequence_in}.
                    Thus $s(n)\in A$ by \cref{funs_type_apply}.
                    $s(n)\in\dom{f}$ by \cref{funs,subseteq}.
                    Thus $s(n)\in\dom{f}\inter A$ by \cref{inter_intro}.
                    Thus $f(s(n))\in\img{f}{A}$ by \cref{img_of_function_intro}.
                    $s\in\funs{\naturalsPlus}{A}$ by \cref{sequence_in}.
                    Thus $s$ is a function by \cref{funs_is_function}.
                    Thus $s\in\rels{\naturalsPlus}{A}$ by \cref{funs_subseteq_rels,subseteq}.
                    Thus $\ran{s}\subseteq A$ by \cref{rels_ran_subseteq}.
                    Thus $\ran{s}\subseteq X$ by \cref{subseteq_transitive}.
                    $f\in\funs{X}{Y}$.
                    Thus $\dom{f} = X$ by \cref{funs}.
                    Thus $\ran{s}\subseteq\dom{f}$ by \cref{subseteq}.
                    Thus $f$ is composable with $s$.
                    Thus $n\in\dom{s}$ by \cref{funs,subseteq}.
                    Thus $(f\circ s)(n) = f(s(n))$ by \cref{circ_apply}.
                    Thus $(f\circ s)(n)\in\img{f}{A}$.
            \end{subproof}
            Thus $f\circ s$ is a function to $\img{f}{A}$.
            $\dom{(f\circ s)} = \naturalsPlus$ by \cref{funs}.
            Thus $f\circ s\in\funs{\naturalsPlus}{\img{f}{A}}$ by \cref{funs_intro}.
            Thus $f\circ s$ is a sequence in $\img{f}{A}$ by \cref{sequence_in}.
        \end{subproof}
        $f\in\funs{X}{Y}$.
        Thus $f\in\rels{X}{Y}$ by \cref{funs_subseteq_rels,subseteq}.
        Thus $\ran{f}\subseteq Y$ by \cref{rels_ran_subseteq}.
        Thus $\img{f}{A}\subseteq \ran{f}$ by \cref{img_subseteq_ran}.
        Thus $\img{f}{A}\subseteq Y$ by \cref{subseteq_transitive}.
        $x\in X$.
        Thus $f(x)\in Y$ by \cref{funs_type_apply}.
        Thus $f(x)\in\closure{\img{f}{A}}{Y}{T_1}$ by \cref{sequence_implies_closure}.
        Thus $y\in\closure{\img{f}{A}}{Y}{T_1}$.
        \end{subproof}
        Thus $\img{f}{\closure{A}{X}{T}}\subseteq \closure{\img{f}{A}}{Y}{T_1}$ by \cref{subseteq}.
    \end{subproof}
    Thus for all $B$ such that $B$ is closed in $Y$ with respect to $T_1$
        we have $\preimg{f}{B}$ is closed in $X$ with respect to $T$ by
        \cref{closure_image_subset_implies_preimg_closed}.
    Thus $f$ is continuous from $X$ to $Y$ with respect to $T$ and $T_1$ by
        \cref{preimg_closed_implies_continuous}.
\end{proof}

% from §21 Item 6: countable basis at a point
\begin{definition}\label{countable_basis_at_point}
    $b$ is a countable basis at $x$ in $X$ with respect to $T$ iff
        $b$ is a sequence in $\pow{X}$ and
        for all $n\in\naturalsPlus$ we have
            $b(n)$ is a neighborhood of $x$ in $X$ with respect to $T$ and
        for all $U$ such that $U$ is a neighborhood of $x$ in $X$ with respect to $T$
            there exists $n\in\naturalsPlus$ such that $b(n)\subseteq U$.
\end{definition}

% from §21 Item 7: first countability axiom
\begin{definition}\label{first_countability_axiom}
    $X$ satisfies the first countability axiom with respect to $T$ iff
        for all $x\in X$ there exists $b$ such that
            $b$ is a countable basis at $x$ in $X$ with respect to $T$.
\end{definition}

% from §21 helper definition 1: addition operation on reals
\begin{definition}\label{real_addition_operation}
    $\realAddition = \{w \mid \exists p\in\reals\times\reals. \exists z\in\reals.
        w = (p,z) \land z = \fst{p} + \snd{p}\}$.
\end{definition}

% from §21 helper definition 2: subtraction operation on reals
\begin{definition}\label{real_subtraction_operation}
    $\realSubtraction = \{w \mid \exists p\in\reals\times\reals. \exists z\in\reals.
        w = (p,z) \land z = \fst{p} - \snd{p}\}$.
\end{definition}

% from §21 helper definition 3: multiplication operation on reals
\begin{definition}\label{real_multiplication_operation}
    $\realMultiplication = \{w \mid \exists p\in\reals\times\reals. \exists z\in\reals.
        w = (p,z) \land z = \fst{p} \rmul \snd{p}\}$.
\end{definition}

% from §21 helper definition 4: quotient operation on reals
\begin{definition}\label{real_division_operation}
    $\realDivision = \{w \mid \exists p\in\reals\times(\reals\setminus\{0\}). \exists z\in\reals.
        w = (p,z) \land z = \fst{p} \rmul \inv{\snd{p}}\}$.
\end{definition}

% from §21 helper lemma 3: standard metric value
\begin{lemma}\label{standard_metric_value}
    Suppose $x\in\reals$.
    Suppose $y\in\reals$.
    Then $\apply{\standardMetricR}{(x,y)} = \abs{x - y}$.
\end{lemma}
\begin{proof}
    Let $d = \standardMetricR$.
    $d$ is a metric on $\reals$ by \cref{standard_metric_is_metric}.
    Thus $d\in\funs{\reals\times\reals}{\reals}$ by \cref{metric_on}.
    Let $p = (x,y)$.
    $p\in\reals\times\reals$ by \cref{times_tuple_intro}.
    $\fst{p} = x$ and $\snd{p} = y$ by \cref{fst_eq,snd_eq}.
    $\neg{y}\in\reals$ by \cref{real_add_inverse}.
    $x - y = x + \neg{y}$ by \cref{real_sub}.
    Thus $x - y\in\reals$ by \cref{real_add_closed}.
    Thus $\abs{x - y}\in\reals$ by \cref{real_abs_in_reals}.
    Let $z = \abs{x - y}$.
    Let $w = (p,z)$.
    $w\in d$ by \cref{standard_metric_r}.
    $d$ is a function by \cref{funs_elim}.
    Thus $d(p) = z$ by \cref{function_apply_intro}.
    Thus $\apply{\standardMetricR}{(x,y)} = \abs{x - y}$.
\end{proof}

% from §21 Item 8: addition operation is continuous
\begin{lemma}\label{real_addition_continuous}
    Let $T = \productTopology{\standardTopologyR}{\standardTopologyR}{\reals}{\reals}$.
    Then $\realAddition$ is continuous from $\reals\times\reals$ to $\reals$ with respect to $T$ and $\standardTopologyR$.
\end{lemma}
\begin{proof}
    Let $f = \realAddition$.
    Show that $f$ is a function from $\reals\times\reals$ to $\reals$.
    \begin{subproof}
        Show that $f$ is a function.
        \begin{subproof}
            Show that for all $p,z_1,z_2$ such that $(p,z_1)\in f$ and $(p,z_2)\in f$ we have $z_1 = z_2$.
            \begin{subproof}
                Fix $p$.
                Fix $z_1$.
                Fix $z_2$.
                Assume $(p,z_1)\in f$ and $(p,z_2)\in f$.
                Take $p_1,w_1$ such that
                    $p_1\in\reals\times\reals$ and $w_1\in\reals$ and
                    $(p,z_1) = (p_1,w_1)$ and $w_1 = \fst{p_1} + \snd{p_1}$
                    by \cref{real_addition_operation}.
                Take $p_2,w_2$ such that
                    $p_2\in\reals\times\reals$ and $w_2\in\reals$ and
                    $(p,z_2) = (p_2,w_2)$ and $w_2 = \fst{p_2} + \snd{p_2}$
                    by \cref{real_addition_operation}.
                $p = p_1$ and $z_1 = w_1$ by \cref{pair_eq_iff}.
                $p = p_2$ and $z_2 = w_2$ by \cref{pair_eq_iff}.
                Thus $z_1 = \fst{p} + \snd{p}$.
                Thus $z_2 = \fst{p} + \snd{p}$.
                Thus $z_1 = z_2$.
            \end{subproof}
            Thus $f$ is right-unique by \cref{rightunique}.
            Show that for all $w$ such that $w\in f$ there exist $x,y$ such that $w = (x,y)$.
            \begin{subproof}
                Fix $w$.
                Assume $w\in f$.
                Take $p,z$ such that $p\in\reals\times\reals$ and $z\in\reals$ and $w = (p,z)$ and
                    $z = \fst{p} + \snd{p}$ by \cref{real_addition_operation}.
                Thus there exist $x,y$ such that $w = (x,y)$.
            \end{subproof}
            Thus $f$ is a relation by \cref{relation}.
            Thus $f$ is a function.
        \end{subproof}
        Show that for all $p\in\dom{f}$ we have $f(p)\in\reals$.
        \begin{subproof}
            Fix $p\in\dom{f}$.
            Take $z$ such that $(p,z)\in f$ by \cref{dom_iff}.
            Take $p_0,w$ such that
                $p_0\in\reals\times\reals$ and $w\in\reals$ and
                $(p,z) = (p_0,w)$ and $w = \fst{p_0} + \snd{p_0}$
                by \cref{real_addition_operation}.
            $z = w$ by \cref{pair_eq_iff}.
            Thus $z\in\reals$.
            $f$ is a function.
            Thus $f(p) = z$ by \cref{function_apply_intro}.
            Thus $f(p)\in\reals$.
        \end{subproof}
        Thus $f$ is a function to $\reals$.
        Show that $\dom{f} = \reals\times\reals$.
        \begin{subproof}
            Show that for all $p\in\dom{f}$ we have $p\in\reals\times\reals$.
            \begin{subproof}
                Fix $p\in\dom{f}$.
                Take $z$ such that $(p,z)\in f$ by \cref{dom_iff}.
                Take $p_0,w$ such that
                    $p_0\in\reals\times\reals$ and $w\in\reals$ and
                    $(p,z) = (p_0,w)$ and $w = \fst{p_0} + \snd{p_0}$
                    by \cref{real_addition_operation}.
                Thus $p = p_0$ by \cref{pair_eq_iff}.
                Thus $p\in\reals\times\reals$.
            \end{subproof}
            Thus $\dom{f}\subseteq\reals\times\reals$ by \cref{subseteq}.
            Show that for all $p\in\reals\times\reals$ we have $p\in\dom{f}$.
            \begin{subproof}
                Fix $p\in\reals\times\reals$.
                Take $x,y$ such that $p = (x,y)$ and $x\in\reals$ and $y\in\reals$ by \cref{times_elem_is_tuple}.
                $x + y\in\reals$ by \cref{real_add_closed}.
                $\fst{p} = x$ and $\snd{p} = y$ by \cref{fst_eq,snd_eq}.
                Let $z = x + y$.
                Let $w = (p,z)$.
                $w\in f$ by \cref{real_addition_operation}.
                Thus $p\in\dom{f}$ by \cref{dom_intro}.
            \end{subproof}
            Thus $\reals\times\reals\subseteq\dom{f}$ by \cref{subseteq}.
            Thus $\dom{f} = \reals\times\reals$ by \cref{subseteq_antisymmetric}.
        \end{subproof}
        Thus $f\in\funs{\reals\times\reals}{\reals}$ by \cref{funs_intro}.
    \end{subproof}
    Show that for all $p\in\reals\times\reals$ we have
        $f$ is continuous at $p$ from $\reals\times\reals$ to $\reals$ with respect to $T$ and $\standardTopologyR$.
    \begin{subproof}
        Fix $p$.
        Assume $p\in\reals\times\reals$.
        Take $x_0,y_0$ such that $p = (x_0,y_0)$ and $x_0\in\reals$ and $y_0\in\reals$ by \cref{times_elem_is_tuple}.
        $x_0 + y_0\in\reals$ by \cref{real_add_closed}.
        Let $z_0 = x_0 + y_0$.
        $\fst{p} = x_0$ and $\snd{p} = y_0$ by \cref{fst_eq,snd_eq}.
        Let $w_0 = (p,z_0)$.
        $w_0\in f$ by \cref{real_addition_operation}.
        $f$ is a function.
        Thus $f(p) = z_0$ by \cref{function_apply_intro}.
        Show that for all $V$ such that $V$ is a neighborhood of $f(p)$ in $\reals$ with respect to $\standardTopologyR$
            there exists $U$ such that $U$ is a neighborhood of $p$ in $\reals\times\reals$ with respect to $T$ and $\img{f}{U}\subseteq V$.
        \begin{subproof}
            Fix $V$.
            Assume $V$ is a neighborhood of $f(p)$ in $\reals$ with respect to $\standardTopologyR$.
            $V$ is open in $\reals$ with respect to $\standardTopologyR$ by \cref{neighborhood}.
            Let $d = \standardMetricR$.
            Let $T_d = \metricTopology{d}{\reals}$.
            $T_d = \standardTopologyR$ by \cref{standard_metric_topology}.
            $\standardTopologyR\subseteq T_d$ by \cref{subseteq}.
            Thus $V\in T_d$ by \cref{open_in,subseteq}.
            $T_d\subseteq\standardTopologyR$ by \cref{subseteq}.
            $f(p)\in V$ by \cref{neighborhood}.
            $V\in\standardTopologyR$ by \cref{open_in,neighborhood}.
            $\standardBasisR$ is a basis on $\reals$ by \cref{standard_basis_is_basis}.
            Thus $\standardTopologyR$ is a topology on $\reals$ by \cref{basis_topology_is_topology,standard_topology_reals}.
            Thus $V\subseteq\reals$ by \cref{topology_on,subseteq,pow_iff}.
            $d$ is a metric on $\reals$ by \cref{standard_metric_is_metric}.
            Take $\epsilon\in\reals$ such that $0 < \epsilon$ and
                $\metricBall{d}{\reals}{f(p)}{\epsilon} \subseteq V$ by \cref{metric_open_iff}.
            Let $\delta = \epsilon / (1 + 1)$.
            $1\in\reals$ by \cref{real_mul_ident}.
            $0\in\reals$ by \cref{real_add_ident}.
            $0 < 1$ by \cref{real_zero_lt_one}.
            $0 + 1 < 1 + 1$ by \cref{real_lt_add}.
            $0 + 1 = 1$ by \cref{real_add_ident}.
            $1 + 1\in\reals$ by \cref{real_add_closed}.
            Thus $1 < 1 + 1$ by \cref{real_lt_subst_left}.
            Thus $0 < 1 + 1$ by \cref{real_lt_trans}.
            Thus $(1 + 1)\neq 0$ by \cref{real_pos_nonzero}.
            Thus $\inv{(1 + 1)}\in\reals$ by \cref{real_mul_inverse}.
            $\epsilon\in\reals$.
            $\delta = \epsilon \rmul \inv{(1 + 1)}$ by \cref{real_div}.
            $\delta\in\reals$ by \cref{real_mul_closed}.
            $0\rmul \inv{(1 + 1)} = 0$ by \cref{real_mul_zero}.
            $0 < \inv{(1 + 1)}$ by \cref{real_inv_pos}.
            Thus $0 < \delta$ by \cref{real_lt_mul_pos}.
            Let $U_1 = \metricBall{d}{\reals}{x_0}{\delta}$.
            Let $U_2 = \metricBall{d}{\reals}{y_0}{\delta}$.
            Show that $U_1$ is open in $\reals$ with respect to $\standardTopologyR$.
            \begin{subproof}
                $x_0\in\reals$ and $\delta\in\reals$ and $0 < \delta$.
                Show that $U_1\in\pow{\reals}$.
                \begin{subproof}
                    Show that for all $x$ such that $x\in U_1$ we have $x\in\reals$.
                    \begin{subproof}
                        Fix $x$.
                        Assume $x\in U_1$.
                        Thus $x\in\reals$ by \cref{metric_ball}.
                    \end{subproof}
                    Thus $U_1\subseteq\reals$ by \cref{subseteq}.
                    Thus $U_1\in\pow{\reals}$ by \cref{powerset_intro}.
                \end{subproof}
                Thus $U_1\in\metricBasis{d}{\reals}$ by \cref{metric_basis}.
                $\metricBasis{d}{\reals}$ is a basis on $\reals$ by \cref{metric_basis_is_basis}.
                Thus $U_1\in T_d$ by \cref{basis_elem_open_in_generated,metric_topology,basis_for_topology}.
                Thus $U_1\in\standardTopologyR$ by \cref{subseteq}.
            \end{subproof}
            Show that $U_2$ is open in $\reals$ with respect to $\standardTopologyR$.
            \begin{subproof}
                $y_0\in\reals$ and $\delta\in\reals$ and $0 < \delta$.
                Show that $U_2\in\pow{\reals}$.
                \begin{subproof}
                    Show that for all $y$ such that $y\in U_2$ we have $y\in\reals$.
                    \begin{subproof}
                        Fix $y$.
                        Assume $y\in U_2$.
                        Thus $y\in\reals$ by \cref{metric_ball}.
                    \end{subproof}
                    Thus $U_2\subseteq\reals$ by \cref{subseteq}.
                    Thus $U_2\in\pow{\reals}$ by \cref{powerset_intro}.
                \end{subproof}
                Thus $U_2\in\metricBasis{d}{\reals}$ by \cref{metric_basis}.
                $\metricBasis{d}{\reals}$ is a basis on $\reals$ by \cref{metric_basis_is_basis}.
                Thus $U_2\in T_d$ by \cref{basis_elem_open_in_generated,metric_topology,basis_for_topology}.
                Thus $U_2\in\standardTopologyR$ by \cref{subseteq}.
            \end{subproof}
            Let $U = U_1\times U_2$.
            $U_1\in\standardTopologyR$ and $U_2\in\standardTopologyR$ by \cref{open_in}.
            $\standardTopologyR\subseteq\pow{\reals}$ by \cref{topology_on}.
            Thus $U_1\in\pow{\reals}$ and $U_2\in\pow{\reals}$ by \cref{subseteq}.
            Thus $U_1\subseteq\reals$ and $U_2\subseteq\reals$ by \cref{pow_iff}.
            $U_1\times U_2\subseteq\reals\times U_2$ by \cref{times_subseteq_left}.
            $\reals\times U_2\subseteq\reals\times\reals$ by \cref{times_subseteq_right}.
            Thus $U\subseteq\reals\times\reals$ by \cref{subseteq_transitive}.
            $U\in\pow{\reals\times\reals}$ by \cref{pow_iff}.
            Thus $U\in\productBasis{\standardTopologyR}{\standardTopologyR}{\reals}{\reals}$ by \cref{product_basis}.
            $\productBasis{\standardTopologyR}{\standardTopologyR}{\reals}{\reals}$ is a basis on $\reals\times\reals$ by \cref{product_basis_is_basis}.
            Thus $T$ is a topology on $\reals\times\reals$ by \cref{basis_topology_is_topology,product_topology}.
            Thus $U\in T$ by \cref{basis_elem_open_in_generated,product_topology}.
            Thus $U$ is open in $\reals\times\reals$ with respect to $T$ by \cref{open_in}.
            Show that $p\in U$.
            \begin{subproof}
                $d((x_0,x_0)) = 0$ by \cref{metric_zero_characterization,standard_metric_is_metric,metric_on}.
                Thus $d((x_0,x_0)) < \delta$ by \cref{real_lt_subst_left}.
                Thus $x_0\in U_1$ by \cref{metric_ball}.
                $d((y_0,y_0)) = 0$ by \cref{metric_zero_characterization,standard_metric_is_metric,metric_on}.
                Thus $d((y_0,y_0)) < \delta$ by \cref{real_lt_subst_left}.
                Thus $y_0\in U_2$ by \cref{metric_ball}.
                Thus $p\in U$ by \cref{times_tuple_intro}.
            \end{subproof}
            Show that $\img{f}{U}\subseteq V$.
            \begin{subproof}
                Show that for all $q$ such that $q\in U$ we have
                    $f(q)\in\metricBall{d}{\reals}{f(p)}{\epsilon}$.
                \begin{subproof}
                    Fix $q$.
                    Assume $q\in U$.
                    Take $x,y$ such that $q = (x,y)$ and $x\in U_1$ and $y\in U_2$ by \cref{times_elem_is_tuple}.
                    $x\in\reals$ and $y\in\reals$ by \cref{metric_ball}.
                    $d((x_0,x)) < \delta$ by \cref{metric_ball}.
                    $d((y_0,y)) < \delta$ by \cref{metric_ball}.
                    $d((x_0,x)) = \abs{x_0 - x}$ by \cref{standard_metric_value}.
                    $d((y_0,y)) = \abs{y_0 - y}$ by \cref{standard_metric_value}.
                    Thus $\abs{x_0 - x} < \delta$ and $\abs{y_0 - y} < \delta$.
                    $\abs{x - x_0} = \abs{x_0 - x}$ by \cref{real_abs_sub_swap}.
                    $\abs{y - y_0} = \abs{y_0 - y}$ by \cref{real_abs_sub_swap}.
                    Thus $\abs{x - x_0} < \delta$ by \cref{real_lt_subst_left}.
                    Thus $\abs{y - y_0} < \delta$ by \cref{real_lt_subst_left}.
                    Show that $f(q) = x + y$.
                    \begin{subproof}
                        $x + y\in\reals$ by \cref{real_add_closed}.
                        $q\in\reals\times\reals$ by \cref{times_tuple_intro}.
                        $\fst{q} = x$ and $\snd{q} = y$ by \cref{fst_eq,snd_eq}.
                        Let $w = (q, x + y)$.
                        $w\in f$ by \cref{real_addition_operation}.
                        $f$ is a function.
                        Thus $f(q) = x + y$ by \cref{function_apply_intro}.
                    \end{subproof}
                    $f(p) = x_0 + y_0$.
                    Show that $\abs{f(q) - f(p)} < \epsilon$.
                    \begin{subproof}
                        $x - x_0 = x + \neg{x_0}$ by \cref{real_sub}.
                        $y - y_0 = y + \neg{y_0}$ by \cref{real_sub}.
                        $(x + y) - (x_0 + y_0) = (x + y) + \neg{(x_0 + y_0)}$ by \cref{real_sub}.
                        $\neg{(x_0 + y_0)} = \neg{x_0} + \neg{y_0}$ by \cref{real_neg_addition}.
                        Show that $(x + y) + (\neg{x_0} + \neg{y_0}) = (x + \neg{x_0}) + (y + \neg{y_0})$.
                        \begin{subproof}
                            $\neg{x_0}\in\reals$ and $\neg{y_0}\in\reals$ by \cref{real_add_inverse}.
                            $\neg{x_0} + \neg{y_0}\in\reals$ by \cref{real_add_closed}.
                            $(x + y) + (\neg{x_0} + \neg{y_0}) = x + (y + (\neg{x_0} + \neg{y_0}))$ by \cref{real_add_assoc}.
                            $y + (\neg{x_0} + \neg{y_0}) = (y + \neg{x_0}) + \neg{y_0}$ by \cref{real_add_assoc}.
                            $y + \neg{x_0} = \neg{x_0} + y$ by \cref{real_add_comm}.
                            Thus $y + (\neg{x_0} + \neg{y_0}) = (\neg{x_0} + y) + \neg{y_0}$.
                            $(\neg{x_0} + y) + \neg{y_0} = \neg{x_0} + (y + \neg{y_0})$ by \cref{real_add_assoc}.
                            Thus $(x + y) + (\neg{x_0} + \neg{y_0}) = x + (\neg{x_0} + (y + \neg{y_0}))$ by \cref{real_add_assoc}.
                            $y + \neg{y_0}\in\reals$ by \cref{real_add_closed}.
                            $\neg{x_0} + (y + \neg{y_0})\in\reals$ by \cref{real_add_closed}.
                            $x + (\neg{x_0} + (y + \neg{y_0})) = (x + \neg{x_0}) + (y + \neg{y_0})$ by \cref{real_add_assoc}.
                        \end{subproof}
                        Thus $(x + y) - (x_0 + y_0) = (x - x_0) + (y - y_0)$.
                        $\neg{x_0}\in\reals$ by \cref{real_add_inverse}.
                        $\neg{y_0}\in\reals$ by \cref{real_add_inverse}.
                        Thus $x - x_0\in\reals$ by \cref{real_sub,real_add_closed}.
                        Thus $y - y_0\in\reals$ by \cref{real_sub,real_add_closed}.
                        $x + y\in\reals$ by \cref{real_add_closed}.
                        $x_0 + y_0\in\reals$ by \cref{real_add_closed}.
                        Thus $(x + y) - (x_0 + y_0)\in\reals$ by \cref{real_sub,real_add_closed}.
                        Thus $\abs{(x + y) - (x_0 + y_0)}\in\reals$ by \cref{real_abs_in_reals}.
                        Thus $\abs{(x + y) - (x_0 + y_0)} \leq \abs{x - x_0} + \abs{y - y_0}$ by \cref{real_triangle_inequality}.
                        $\abs{x - x_0}\in\reals$ by \cref{real_abs_in_reals}.
                        $\abs{y - y_0}\in\reals$ by \cref{real_abs_in_reals}.
                        $\abs{x - x_0} + \abs{y - y_0}\in\reals$ by \cref{real_add_closed}.
                        $\delta\in\reals$ by \cref{real_mul_closed}.
                        $\delta + \delta\in\reals$ by \cref{real_add_closed}.
                        $\abs{x - x_0} + \abs{y - y_0} < \delta + \delta$ by \cref{real_lt_add_two}.
                        Show that $\abs{(x + y) - (x_0 + y_0)} < \delta + \delta$.
                        \begin{subproof}
                            Thus $\abs{(x + y) - (x_0 + y_0)} < \abs{x - x_0} + \abs{y - y_0}$ or
                                $\abs{(x + y) - (x_0 + y_0)} = \abs{x - x_0} + \abs{y - y_0}$ by \cref{real_leq_split}.
                            \begin{byCase}
                            \caseOf{$\abs{(x + y) - (x_0 + y_0)} < \abs{x - x_0} + \abs{y - y_0}$.}
                                Thus $\abs{(x + y) - (x_0 + y_0)} < \delta + \delta$ by \cref{real_lt_trans}.
                            \caseOf{$\abs{(x + y) - (x_0 + y_0)} = \abs{x - x_0} + \abs{y - y_0}$.}
                                Thus $\abs{(x + y) - (x_0 + y_0)} < \delta + \delta$ by \cref{real_lt_subst_left}.
                            \end{byCase}
                        \end{subproof}
                        $1\rmul \delta = \delta$ by \cref{real_mul_ident}.
                        $1 \rmul \delta\in\reals$ by \cref{real_mul_closed}.
                        $(1 + 1) \rmul \delta = (1 \rmul \delta) + (1 \rmul \delta)$ by \cref{real_right_distrib}.
                        Thus $(1 \rmul \delta) + (1 \rmul \delta) = \delta + (1 \rmul \delta)$ by \cref{real_add_eq_subst}.
                        Thus $(1 \rmul \delta) + (1 \rmul \delta) = \delta + \delta$ by \cref{real_add_eq_subst}.
                        Thus $(1 + 1) \rmul \delta = \delta + \delta$.
                        $\delta = \epsilon \rmul \inv{(1 + 1)}$ by \cref{real_div}.
                        Thus $(1 + 1) \rmul \delta = (1 + 1) \rmul (\epsilon \rmul \inv{(1 + 1)})$.
                        $(1 + 1) \rmul (\epsilon \rmul \inv{(1 + 1)}) = \epsilon \rmul ((1 + 1) \rmul \inv{(1 + 1)})$ by \cref{real_mul_assoc,real_mul_comm}.
                        $(1 + 1) \rmul \inv{(1 + 1)} = 1$ by \cref{real_mul_inverse}.
                        $\epsilon \rmul 1 = 1 \rmul \epsilon$ by \cref{real_mul_comm}.
                        $1 \rmul \epsilon = \epsilon$ by \cref{real_mul_ident}.
                        Thus $\epsilon \rmul 1 = \epsilon$.
                        Thus $\delta + \delta = \epsilon$.
                        Thus $\abs{(x + y) - (x_0 + y_0)} < \epsilon$ by \cref{real_lt_subst_right}.
                        Thus $\abs{f(q) - f(p)} = \abs{(x + y) - (x_0 + y_0)}$.
                    Thus $\abs{f(q) - f(p)} < \epsilon$ by \cref{real_lt_subst_left}.
                \end{subproof}
                $f$ is a function to $\reals$ and $\dom{f} = \reals\times\reals$ by \cref{funs_elim}.
                Thus $f(p)\in\reals$.
                Thus $f(q)\in\reals$.
                $d((f(p),f(q))) = \abs{f(p) - f(q)}$ by \cref{standard_metric_value}.
                $\abs{f(p) - f(q)} = \abs{f(q) - f(p)}$ by \cref{real_abs_sub_swap}.
                Thus $\abs{f(p) - f(q)} < \epsilon$ by \cref{real_lt_subst_left}.
                Thus $d((f(p),f(q))) < \epsilon$ by \cref{real_lt_subst_left}.
                Thus $f(q)\in\metricBall{d}{\reals}{f(p)}{\epsilon}$ by \cref{metric_ball}.
                \end{subproof}
                Show that for all $y$ such that $y\in\img{f}{U}$ we have
                    $y\in\metricBall{d}{\reals}{f(p)}{\epsilon}$.
                \begin{subproof}
                    Fix $y$.
                    Assume $y\in\img{f}{U}$.
                    Take $q\in U$ such that $q\mathrel{f} y$ by \cref{img_iff}.
                    $f$ is a function.
                    Thus $y = f(q)$ by \cref{function_apply_intro}.
                    Thus $y\in\metricBall{d}{\reals}{f(p)}{\epsilon}$.
                \end{subproof}
                Thus $\img{f}{U}\subseteq \metricBall{d}{\reals}{f(p)}{\epsilon}$ by \cref{subseteq}.
                $\metricBall{d}{\reals}{f(p)}{\epsilon}\subseteq V$ by \cref{subseteq}.
                Thus $\img{f}{U}\subseteq V$ by \cref{subseteq_transitive}.
            \end{subproof}
            Thus $U$ is a neighborhood of $p$ in $\reals\times\reals$ with respect to $T$ by \cref{neighborhood}.
        \end{subproof}
        $\standardBasisR$ is a basis on $\reals$ by \cref{standard_basis_is_basis}.
        Thus $\standardTopologyR$ is a topology on $\reals$ by \cref{basis_topology_is_topology,standard_topology_reals}.
        $\productBasis{\standardTopologyR}{\standardTopologyR}{\reals}{\reals}$ is a basis on $\reals\times\reals$ by \cref{product_basis_is_basis}.
        Thus $T$ is a topology on $\reals\times\reals$ by \cref{basis_topology_is_topology,product_topology}.
        Thus $f$ is continuous at $p$ from $\reals\times\reals$ to $\reals$ with respect to $T$ and $\standardTopologyR$ by \cref{continuous_at}.
    \end{subproof}
    $\standardBasisR$ is a basis on $\reals$ by \cref{standard_basis_is_basis}.
    Thus $\standardTopologyR$ is a topology on $\reals$ by \cref{basis_topology_is_topology,standard_topology_reals}.
    $\productBasis{\standardTopologyR}{\standardTopologyR}{\reals}{\reals}$ is a basis on $\reals\times\reals$ by \cref{product_basis_is_basis}.
    Thus $T$ is a topology on $\reals\times\reals$ by \cref{basis_topology_is_topology,product_topology}.
    Thus $f$ is continuous from $\reals\times\reals$ to $\reals$ with respect to $T$ and $\standardTopologyR$ by \cref{continuous_at_implies_continuous}.
\end{proof}

% from §21 Item 9: subtraction operation is continuous
\begin{lemma}\label{real_subtraction_continuous}
    Let $T = \productTopology{\standardTopologyR}{\standardTopologyR}{\reals}{\reals}$.
    Then $\realSubtraction$ is continuous from $\reals\times\reals$ to $\reals$ with respect to $T$ and $\standardTopologyR$.
\end{lemma}
\begin{proof}
    Let $f = \realSubtraction$.
    Show that $f$ is a function from $\reals\times\reals$ to $\reals$.
    \begin{subproof}
        Show that $f$ is a function.
        \begin{subproof}
            Show that for all $p,z_1,z_2$ such that $(p,z_1)\in f$ and $(p,z_2)\in f$ we have $z_1 = z_2$.
            \begin{subproof}
                Fix $p$.
                Fix $z_1$.
                Fix $z_2$.
                Assume $(p,z_1)\in f$ and $(p,z_2)\in f$.
                Take $p_1,w_1$ such that
                    $p_1\in\reals\times\reals$ and $w_1\in\reals$ and
                    $(p,z_1) = (p_1,w_1)$ and $w_1 = \fst{p_1} - \snd{p_1}$
                    by \cref{real_subtraction_operation}.
                Take $p_2,w_2$ such that
                    $p_2\in\reals\times\reals$ and $w_2\in\reals$ and
                    $(p,z_2) = (p_2,w_2)$ and $w_2 = \fst{p_2} - \snd{p_2}$
                    by \cref{real_subtraction_operation}.
                $p = p_1$ and $z_1 = w_1$ by \cref{pair_eq_iff}.
                $p = p_2$ and $z_2 = w_2$ by \cref{pair_eq_iff}.
                Thus $z_1 = \fst{p} - \snd{p}$.
                Thus $z_2 = \fst{p} - \snd{p}$.
                Thus $z_1 = z_2$.
            \end{subproof}
            Thus $f$ is right-unique by \cref{rightunique}.
            Show that for all $w$ such that $w\in f$ there exist $x,y$ such that $w = (x,y)$.
            \begin{subproof}
                Fix $w$.
                Assume $w\in f$.
                Take $p,z$ such that $p\in\reals\times\reals$ and $z\in\reals$ and $w = (p,z)$ and
                    $z = \fst{p} - \snd{p}$ by \cref{real_subtraction_operation}.
                Thus there exist $x,y$ such that $w = (x,y)$.
            \end{subproof}
            Thus $f$ is a relation by \cref{relation}.
            Thus $f$ is a function.
        \end{subproof}
        Show that for all $p\in\dom{f}$ we have $f(p)\in\reals$.
        \begin{subproof}
            Fix $p\in\dom{f}$.
            Take $z$ such that $(p,z)\in f$ by \cref{dom_iff}.
            Take $p_0,w$ such that
                $p_0\in\reals\times\reals$ and $w\in\reals$ and
                $(p,z) = (p_0,w)$ and $w = \fst{p_0} - \snd{p_0}$
                by \cref{real_subtraction_operation}.
            $z = w$ by \cref{pair_eq_iff}.
            Thus $z\in\reals$.
            $f$ is a function.
            Thus $f(p) = z$ by \cref{function_apply_intro}.
            Thus $f(p)\in\reals$.
        \end{subproof}
        Thus $f$ is a function to $\reals$.
        Show that $\dom{f} = \reals\times\reals$.
        \begin{subproof}
            Show that for all $p\in\dom{f}$ we have $p\in\reals\times\reals$.
            \begin{subproof}
                Fix $p\in\dom{f}$.
                Take $z$ such that $(p,z)\in f$ by \cref{dom_iff}.
                Take $p_0,w$ such that
                    $p_0\in\reals\times\reals$ and $w\in\reals$ and
                    $(p,z) = (p_0,w)$ and $w = \fst{p_0} - \snd{p_0}$
                    by \cref{real_subtraction_operation}.
                Thus $p = p_0$ by \cref{pair_eq_iff}.
                Thus $p\in\reals\times\reals$.
            \end{subproof}
            Thus $\dom{f}\subseteq\reals\times\reals$ by \cref{subseteq}.
            Show that for all $p\in\reals\times\reals$ we have $p\in\dom{f}$.
            \begin{subproof}
                Fix $p\in\reals\times\reals$.
                Take $x,y$ such that $p = (x,y)$ and $x\in\reals$ and $y\in\reals$ by \cref{times_elem_is_tuple}.
                $x - y\in\reals$ by \cref{real_sub,real_add_closed,real_add_inverse}.
                $\fst{p} = x$ and $\snd{p} = y$ by \cref{fst_eq,snd_eq}.
                Let $z = x - y$.
                Let $w = (p,z)$.
                $w\in f$ by \cref{real_subtraction_operation}.
                Thus $p\in\dom{f}$ by \cref{dom_intro}.
            \end{subproof}
            Thus $\reals\times\reals\subseteq\dom{f}$ by \cref{subseteq}.
            Thus $\dom{f} = \reals\times\reals$ by \cref{subseteq_antisymmetric}.
        \end{subproof}
        Thus $f\in\funs{\reals\times\reals}{\reals}$ by \cref{funs_intro}.
    \end{subproof}
    Show that for all $p\in\reals\times\reals$ we have
        $f$ is continuous at $p$ from $\reals\times\reals$ to $\reals$ with respect to $T$ and $\standardTopologyR$.
    \begin{subproof}
        Fix $p$.
        Assume $p\in\reals\times\reals$.
        Take $x_0,y_0$ such that $p = (x_0,y_0)$ and $x_0\in\reals$ and $y_0\in\reals$ by \cref{times_elem_is_tuple}.
        $x_0 - y_0\in\reals$ by \cref{real_sub,real_add_closed,real_add_inverse}.
        Let $z_0 = x_0 - y_0$.
        $\fst{p} = x_0$ and $\snd{p} = y_0$ by \cref{fst_eq,snd_eq}.
        Let $w_0 = (p,z_0)$.
        $w_0\in f$ by \cref{real_subtraction_operation}.
        $f$ is a function.
        Thus $f(p) = z_0$ by \cref{function_apply_intro}.
        Show that for all $V$ such that $V$ is a neighborhood of $f(p)$ in $\reals$ with respect to $\standardTopologyR$
            there exists $U$ such that $U$ is a neighborhood of $p$ in $\reals\times\reals$ with respect to $T$ and $\img{f}{U}\subseteq V$.
        \begin{subproof}
            Fix $V$.
            Assume $V$ is a neighborhood of $f(p)$ in $\reals$ with respect to $\standardTopologyR$.
            $V$ is open in $\reals$ with respect to $\standardTopologyR$ by \cref{neighborhood}.
            Let $d = \standardMetricR$.
            Let $T_d = \metricTopology{d}{\reals}$.
            $T_d = \standardTopologyR$ by \cref{standard_metric_topology}.
            $\standardTopologyR\subseteq T_d$ by \cref{subseteq}.
            Thus $V\in T_d$ by \cref{open_in,subseteq}.
            $T_d\subseteq\standardTopologyR$ by \cref{subseteq}.
            $f(p)\in V$ by \cref{neighborhood}.
            $V\in\standardTopologyR$ by \cref{open_in,neighborhood}.
            $\standardBasisR$ is a basis on $\reals$ by \cref{standard_basis_is_basis}.
            Thus $\standardTopologyR$ is a topology on $\reals$ by \cref{basis_topology_is_topology,standard_topology_reals}.
            Thus $V\subseteq\reals$ by \cref{topology_on,subseteq,pow_iff}.
            $d$ is a metric on $\reals$ by \cref{standard_metric_is_metric}.
            Take $\epsilon\in\reals$ such that $0 < \epsilon$ and
                $\metricBall{d}{\reals}{f(p)}{\epsilon} \subseteq V$ by \cref{metric_open_iff}.
            Let $\delta = \epsilon / (1 + 1)$.
            $1\in\reals$ by \cref{real_mul_ident}.
            $0\in\reals$ by \cref{real_add_ident}.
            $0 < 1$ by \cref{real_zero_lt_one}.
            $0 + 1 < 1 + 1$ by \cref{real_lt_add}.
            $0 + 1 = 1$ by \cref{real_add_ident}.
            $1 + 1\in\reals$ by \cref{real_add_closed}.
            Thus $1 < 1 + 1$ by \cref{real_lt_subst_left}.
            Thus $0 < 1 + 1$ by \cref{real_lt_trans}.
            Thus $(1 + 1)\neq 0$ by \cref{real_pos_nonzero}.
            Thus $\inv{(1 + 1)}\in\reals$ by \cref{real_mul_inverse}.
            $\epsilon\in\reals$.
            $\delta = \epsilon \rmul \inv{(1 + 1)}$ by \cref{real_div}.
            $0\rmul \inv{(1 + 1)} = 0$ by \cref{real_mul_zero}.
            $0 < \inv{(1 + 1)}$ by \cref{real_inv_pos}.
            Thus $0 < \delta$ by \cref{real_lt_mul_pos}.
            Let $U_1 = \metricBall{d}{\reals}{x_0}{\delta}$.
            Let $U_2 = \metricBall{d}{\reals}{y_0}{\delta}$.
            Show that $U_1$ is open in $\reals$ with respect to $\standardTopologyR$.
            \begin{subproof}
                $x_0\in\reals$ and $\delta\in\reals$ and $0 < \delta$.
                Show that $U_1\in\pow{\reals}$.
                \begin{subproof}
                    Show that for all $x$ such that $x\in U_1$ we have $x\in\reals$.
                    \begin{subproof}
                        Fix $x$.
                        Assume $x\in U_1$.
                        Thus $x\in\reals$ by \cref{metric_ball}.
                    \end{subproof}
                    Thus $U_1\subseteq\reals$ by \cref{subseteq}.
                    Thus $U_1\in\pow{\reals}$ by \cref{powerset_intro}.
                \end{subproof}
                Thus $U_1\in\metricBasis{d}{\reals}$ by \cref{metric_basis}.
                $\metricBasis{d}{\reals}$ is a basis on $\reals$ by \cref{metric_basis_is_basis}.
                Thus $U_1\in T_d$ by \cref{basis_elem_open_in_generated,metric_topology,basis_for_topology}.
                Thus $U_1\in\standardTopologyR$ by \cref{subseteq}.
            \end{subproof}
            Show that $U_2$ is open in $\reals$ with respect to $\standardTopologyR$.
            \begin{subproof}
                $y_0\in\reals$ and $\delta\in\reals$ and $0 < \delta$.
                Show that $U_2\in\pow{\reals}$.
                \begin{subproof}
                    Show that for all $y$ such that $y\in U_2$ we have $y\in\reals$.
                    \begin{subproof}
                        Fix $y$.
                        Assume $y\in U_2$.
                        Thus $y\in\reals$ by \cref{metric_ball}.
                    \end{subproof}
                    Thus $U_2\subseteq\reals$ by \cref{subseteq}.
                    Thus $U_2\in\pow{\reals}$ by \cref{powerset_intro}.
                \end{subproof}
                Thus $U_2\in\metricBasis{d}{\reals}$ by \cref{metric_basis}.
                $\metricBasis{d}{\reals}$ is a basis on $\reals$ by \cref{metric_basis_is_basis}.
                Thus $U_2\in T_d$ by \cref{basis_elem_open_in_generated,metric_topology,basis_for_topology}.
                Thus $U_2\in\standardTopologyR$ by \cref{subseteq}.
            \end{subproof}
            Let $U = U_1\times U_2$.
            $U_1\in\standardTopologyR$ and $U_2\in\standardTopologyR$ by \cref{open_in}.
            $\standardTopologyR\subseteq\pow{\reals}$ by \cref{topology_on}.
            Thus $U_1\in\pow{\reals}$ and $U_2\in\pow{\reals}$ by \cref{subseteq}.
            Thus $U_1\subseteq\reals$ and $U_2\subseteq\reals$ by \cref{pow_iff}.
            $U_1\times U_2\subseteq\reals\times U_2$ by \cref{times_subseteq_left}.
            $\reals\times U_2\subseteq\reals\times\reals$ by \cref{times_subseteq_right}.
            Thus $U\subseteq\reals\times\reals$ by \cref{subseteq_transitive}.
            $U\in\pow{\reals\times\reals}$ by \cref{pow_iff}.
            Thus $U\in\productBasis{\standardTopologyR}{\standardTopologyR}{\reals}{\reals}$ by \cref{product_basis}.
            $\productBasis{\standardTopologyR}{\standardTopologyR}{\reals}{\reals}$ is a basis on $\reals\times\reals$ by \cref{product_basis_is_basis}.
            Thus $T$ is a topology on $\reals\times\reals$ by \cref{basis_topology_is_topology,product_topology}.
            Thus $U\in T$ by \cref{basis_elem_open_in_generated,product_topology}.
            Thus $U$ is open in $\reals\times\reals$ with respect to $T$ by \cref{open_in}.
            Show that $p\in U$.
            \begin{subproof}
                $d((x_0,x_0)) = 0$ by \cref{metric_zero_characterization,standard_metric_is_metric,metric_on}.
                Thus $d((x_0,x_0)) < \delta$ by \cref{real_lt_subst_left}.
                Thus $x_0\in U_1$ by \cref{metric_ball}.
                $d((y_0,y_0)) = 0$ by \cref{metric_zero_characterization,standard_metric_is_metric,metric_on}.
                Thus $d((y_0,y_0)) < \delta$ by \cref{real_lt_subst_left}.
                Thus $y_0\in U_2$ by \cref{metric_ball}.
                Thus $p\in U$ by \cref{times_tuple_intro}.
            \end{subproof}
            Show that $\img{f}{U}\subseteq V$.
            \begin{subproof}
                Show that for all $q$ such that $q\in U$ we have
                    $f(q)\in\metricBall{d}{\reals}{f(p)}{\epsilon}$.
                \begin{subproof}
                    Fix $q$.
                    Assume $q\in U$.
                    Take $x,y$ such that $q = (x,y)$ and $x\in U_1$ and $y\in U_2$ by \cref{times_elem_is_tuple}.
                    $x\in\reals$ and $y\in\reals$ by \cref{metric_ball}.
                    $d((x_0,x)) < \delta$ by \cref{metric_ball}.
                    $d((y_0,y)) < \delta$ by \cref{metric_ball}.
                    $d((x_0,x)) = \abs{x_0 - x}$ by \cref{standard_metric_value}.
                    $d((y_0,y)) = \abs{y_0 - y}$ by \cref{standard_metric_value}.
                    Thus $\abs{x_0 - x} < \delta$ and $\abs{y_0 - y} < \delta$.
                    $\abs{x - x_0} = \abs{x_0 - x}$ by \cref{real_abs_sub_swap}.
                    Thus $\abs{x - x_0} < \delta$ by \cref{real_lt_subst_left}.
                    $\abs{y - y_0} = \abs{y_0 - y}$ by \cref{real_abs_sub_swap}.
                    Thus $\abs{y - y_0} < \delta$ by \cref{real_lt_subst_left}.
                    Show that $f(q) = x - y$.
                    \begin{subproof}
                        $x - y\in\reals$ by \cref{real_sub,real_add_closed,real_add_inverse}.
                        $q\in\reals\times\reals$ by \cref{times_tuple_intro}.
                        $\fst{q} = x$ and $\snd{q} = y$ by \cref{fst_eq,snd_eq}.
                        Let $w = (q, x - y)$.
                        $w\in f$ by \cref{real_subtraction_operation}.
                        $f$ is a function.
                        Thus $f(q) = x - y$ by \cref{function_apply_intro}.
                    \end{subproof}
                    $f(p) = x_0 - y_0$.
                    Show that $\abs{f(q) - f(p)} < \epsilon$.
                    \begin{subproof}
                        $x - y = x + \neg{y}$ by \cref{real_sub}.
                        $x_0 - y_0 = x_0 + \neg{y_0}$ by \cref{real_sub}.
                        $(x - y) - (x_0 - y_0) = (x + \neg{y}) + \neg{(x_0 + \neg{y_0})}$ by \cref{real_sub}.
                        $\neg{y_0}\in\reals$ by \cref{real_add_inverse}.
                        $\neg{(x_0 + \neg{y_0})} = \neg{x_0} + \neg{(\neg{y_0})}$ by \cref{real_neg_addition}.
                        $\neg{(\neg{y_0})} = y_0$ by \cref{real_neg_neg}.
                        Thus $(x - y) - (x_0 - y_0) = (x + \neg{y}) + (\neg{x_0} + y_0)$.
                        $\neg{y}\in\reals$ by \cref{real_add_inverse}.
                        $\neg{x_0}\in\reals$ by \cref{real_add_inverse}.
                        $\neg{x_0} + y_0\in\reals$ by \cref{real_add_closed}.
                        $\neg{y} + \neg{x_0}\in\reals$ by \cref{real_add_closed}.
                        $\neg{y} + y_0\in\reals$ by \cref{real_add_closed}.
                        $(x + \neg{y}) + (\neg{x_0} + y_0) = x + (\neg{y} + (\neg{x_0} + y_0))$ by \cref{real_add_assoc}.
                        $\neg{y} + (\neg{x_0} + y_0) = (\neg{y} + \neg{x_0}) + y_0$ by \cref{real_add_assoc}.
                        $\neg{y} + \neg{x_0} = \neg{x_0} + \neg{y}$ by \cref{real_add_comm}.
                        Thus $\neg{y} + (\neg{x_0} + y_0) = (\neg{x_0} + \neg{y}) + y_0$.
                        $(\neg{x_0} + \neg{y}) + y_0 = \neg{x_0} + (\neg{y} + y_0)$ by \cref{real_add_assoc}.
                        Thus $(x - y) - (x_0 - y_0) = x + (\neg{x_0} + (\neg{y} + y_0))$ by \cref{real_add_assoc}.
                        $\neg{y} + y_0 = y_0 + \neg{y}$ by \cref{real_add_comm}.
                        Thus $(x - y) - (x_0 - y_0) = x + (\neg{x_0} + (y_0 + \neg{y}))$.
                        $y_0 + \neg{y} = y_0 - y$ by \cref{real_sub}.
                        $x + \neg{x_0} = x - x_0$ by \cref{real_sub}.
                        Thus $(x - y) - (x_0 - y_0) = (x - x_0) + (y_0 - y)$ by \cref{real_add_assoc}.
                        $x - x_0\in\reals$ by \cref{real_sub,real_add_closed}.
                        $y_0 - y\in\reals$ by \cref{real_sub,real_add_closed}.
                        $\abs{(x - x_0) + (y_0 - y)} \leq \abs{x - x_0} + \abs{y_0 - y}$ by \cref{real_triangle_inequality}.
                        $\abs{y_0 - y} = \abs{y - y_0}$ by \cref{real_abs_sub_swap}.
                        Thus $\abs{y_0 - y} < \delta$ by \cref{real_lt_subst_left}.
                        $\abs{x - x_0}\in\reals$ by \cref{real_abs_in_reals}.
                        $\abs{y_0 - y}\in\reals$ by \cref{real_abs_in_reals}.
                        $\delta\in\reals$ by \cref{real_mul_closed}.
                        $\abs{x - x_0} + \abs{y_0 - y} < \delta + \delta$ by \cref{real_lt_add_two}.
                        $\abs{x - x_0} + \abs{y_0 - y}\in\reals$ by \cref{real_add_closed}.
                        $\delta + \delta\in\reals$ by \cref{real_add_closed}.
                        $(x - x_0) + (y_0 - y)\in\reals$ by \cref{real_add_closed}.
                        $\abs{(x - x_0) + (y_0 - y)}\in\reals$ by \cref{real_abs_in_reals}.
                        Show that $\abs{(x - x_0) + (y_0 - y)} < \delta + \delta$.
                        \begin{subproof}
                            Thus $\abs{(x - x_0) + (y_0 - y)} < \abs{x - x_0} + \abs{y_0 - y}$ or
                                $\abs{(x - x_0) + (y_0 - y)} = \abs{x - x_0} + \abs{y_0 - y}$ by \cref{real_leq_split}.
                            \begin{byCase}
                            \caseOf{$\abs{(x - x_0) + (y_0 - y)} < \abs{x - x_0} + \abs{y_0 - y}$.}
                                $\abs{(x - x_0) + (y_0 - y)}\in\reals$ by \cref{real_abs_in_reals}.
                                $\abs{x - x_0} + \abs{y_0 - y}\in\reals$ by \cref{real_add_closed}.
                                $\delta + \delta\in\reals$ by \cref{real_add_closed}.
                                Thus $\abs{(x - x_0) + (y_0 - y)} < \delta + \delta$ by \cref{real_lt_trans}.
                            \caseOf{$\abs{(x - x_0) + (y_0 - y)} = \abs{x - x_0} + \abs{y_0 - y}$.}
                                Thus $\abs{(x - x_0) + (y_0 - y)} < \delta + \delta$ by \cref{real_lt_subst_left}.
                            \end{byCase}
                        \end{subproof}
                        $1\rmul \delta = \delta$ by \cref{real_mul_ident}.
                        $1 \rmul \delta\in\reals$ by \cref{real_mul_closed}.
                        $(1 + 1) \rmul \delta = (1 \rmul \delta) + (1 \rmul \delta)$ by \cref{real_right_distrib}.
                        Thus $(1 \rmul \delta) + (1 \rmul \delta) = \delta + (1 \rmul \delta)$ by \cref{real_add_eq_subst}.
                        Thus $(1 \rmul \delta) + (1 \rmul \delta) = \delta + \delta$ by \cref{real_add_eq_subst}.
                        Thus $(1 + 1) \rmul \delta = \delta + \delta$.
                        $\delta = \epsilon \rmul \inv{(1 + 1)}$ by \cref{real_div}.
                        Thus $(1 + 1) \rmul \delta = (1 + 1) \rmul (\epsilon \rmul \inv{(1 + 1)})$.
                        $(1 + 1) \rmul (\epsilon \rmul \inv{(1 + 1)}) = \epsilon \rmul ((1 + 1) \rmul \inv{(1 + 1)})$ by \cref{real_mul_assoc,real_mul_comm}.
                        $(1 + 1) \rmul \inv{(1 + 1)} = 1$ by \cref{real_mul_inverse}.
                        $\epsilon \rmul 1 = 1 \rmul \epsilon$ by \cref{real_mul_comm}.
                        $1 \rmul \epsilon = \epsilon$ by \cref{real_mul_ident}.
                        Thus $\epsilon \rmul 1 = \epsilon$.
                        Thus $\delta + \delta = \epsilon$.
                        Thus $\abs{(x - x_0) + (y_0 - y)} < \epsilon$ by \cref{real_lt_subst_right}.
                        Thus $\abs{f(q) - f(p)} = \abs{(x - y) - (x_0 - y_0)}$.
                    Thus $\abs{f(q) - f(p)} < \epsilon$ by \cref{real_lt_subst_left}.
                \end{subproof}
                $f$ is a function to $\reals$ and $\dom{f} = \reals\times\reals$ by \cref{funs_elim}.
                Thus $f(p)\in\reals$.
                Thus $f(q)\in\reals$.
                $d((f(p),f(q))) = \abs{f(p) - f(q)}$ by \cref{standard_metric_value}.
                $\abs{f(p) - f(q)} = \abs{f(q) - f(p)}$ by \cref{real_abs_sub_swap}.
                Thus $\abs{f(p) - f(q)} < \epsilon$ by \cref{real_lt_subst_left}.
                Thus $d((f(p),f(q))) < \epsilon$ by \cref{real_lt_subst_left}.
                Thus $f(q)\in\metricBall{d}{\reals}{f(p)}{\epsilon}$ by \cref{metric_ball}.
                \end{subproof}
                Show that for all $y$ such that $y\in\img{f}{U}$ we have
                    $y\in\metricBall{d}{\reals}{f(p)}{\epsilon}$.
                \begin{subproof}
                    Fix $y$.
                    Assume $y\in\img{f}{U}$.
                    Take $q\in U$ such that $q\mathrel{f} y$ by \cref{img_iff}.
                    $f$ is a function.
                    Thus $y = f(q)$ by \cref{function_apply_intro}.
                    Thus $y\in\metricBall{d}{\reals}{f(p)}{\epsilon}$.
                \end{subproof}
                Thus $\img{f}{U}\subseteq \metricBall{d}{\reals}{f(p)}{\epsilon}$ by \cref{subseteq}.
                $\metricBall{d}{\reals}{f(p)}{\epsilon}\subseteq V$ by \cref{subseteq}.
                Thus $\img{f}{U}\subseteq V$ by \cref{subseteq_transitive}.
            \end{subproof}
            Thus $U$ is a neighborhood of $p$ in $\reals\times\reals$ with respect to $T$ by \cref{neighborhood}.
        \end{subproof}
        $\standardBasisR$ is a basis on $\reals$ by \cref{standard_basis_is_basis}.
        Thus $\standardTopologyR$ is a topology on $\reals$ by \cref{basis_topology_is_topology,standard_topology_reals}.
        $\productBasis{\standardTopologyR}{\standardTopologyR}{\reals}{\reals}$ is a basis on $\reals\times\reals$ by \cref{product_basis_is_basis}.
        Thus $T$ is a topology on $\reals\times\reals$ by \cref{basis_topology_is_topology,product_topology}.
        Thus $f$ is continuous at $p$ from $\reals\times\reals$ to $\reals$ with respect to $T$ and $\standardTopologyR$ by \cref{continuous_at}.
    \end{subproof}
    $\standardBasisR$ is a basis on $\reals$ by \cref{standard_basis_is_basis}.
    Thus $\standardTopologyR$ is a topology on $\reals$ by \cref{basis_topology_is_topology,standard_topology_reals}.
    $\productBasis{\standardTopologyR}{\standardTopologyR}{\reals}{\reals}$ is a basis on $\reals\times\reals$ by \cref{product_basis_is_basis}.
    Thus $T$ is a topology on $\reals\times\reals$ by \cref{basis_topology_is_topology,product_topology}.
    Thus $f$ is continuous from $\reals\times\reals$ to $\reals$ with respect to $T$ and $\standardTopologyR$ by \cref{continuous_at_implies_continuous}.
\end{proof}

% from §21 helper lemma 4: absolute value bound by nearby point
\begin{lemma}\label{real_abs_sub_lt_bound}
    Suppose $x\in\reals$ and $y\in\reals$ and $\delta\in\reals$.
    Suppose $\abs{x - y} < \delta$.
    Then $\abs{x} < \abs{y} + \delta$.
\end{lemma}
\begin{proof}
    $x - y = x + \neg{y}$ by \cref{real_sub}.
    $\neg{y}\in\reals$ by \cref{real_add_inverse}.
    Thus $x - y\in\reals$ by \cref{real_add_closed}.
    $(x - y) + y = x + (\neg{y} + y)$ by \cref{real_add_assoc}.
    $\neg{y} + y = 0$ by \cref{real_add_inverse,real_add_comm}.
    $0 + x = x$ by \cref{real_add_ident}.
    $0\in\reals$ by \cref{real_add_ident}.
    $x + 0 = 0 + x$ by \cref{real_add_comm}.
    Thus $x + 0 = x$.
    Thus $(x - y) + y = x$.
    $\abs{(x - y) + y} \leq \abs{x - y} + \abs{y}$ by \cref{real_triangle_inequality}.
    Thus $\abs{x} \leq \abs{x - y} + \abs{y}$ by \cref{real_lt_subst_left}.
    $\abs{x - y}\in\reals$ by \cref{real_abs_in_reals}.
    $\abs{y}\in\reals$ by \cref{real_abs_in_reals}.
    Thus $\abs{x - y} + \abs{y}\in\reals$ by \cref{real_add_closed}.
    $\abs{x}\in\reals$ by \cref{real_abs_in_reals}.
    $\delta + \abs{y}\in\reals$ by \cref{real_add_closed}.
    $\delta\in\reals$.
    Thus $\abs{x - y} + \abs{y} < \delta + \abs{y}$ by \cref{real_lt_add}.
    Show that $\abs{x} < \delta + \abs{y}$.
    \begin{subproof}
        Thus $\abs{x} < \abs{x - y} + \abs{y}$ or
            $\abs{x} = \abs{x - y} + \abs{y}$ by \cref{real_leq_split}.
        \begin{byCase}
        \caseOf{$\abs{x} < \abs{x - y} + \abs{y}$.}
            Thus $\abs{x} < \delta + \abs{y}$ by \cref{real_lt_trans}.
        \caseOf{$\abs{x} = \abs{x - y} + \abs{y}$.}
            Thus $\abs{x} < \delta + \abs{y}$ by \cref{real_lt_subst_left}.
        \end{byCase}
    \end{subproof}
    $\delta + \abs{y} = \abs{y} + \delta$ by \cref{real_add_comm}.
    Thus $\abs{x} < \abs{y} + \delta$ by \cref{real_lt_subst_right}.
\end{proof}

% from §21 helper lemma 5: positive sum with nonnegative term
\begin{lemma}\label{real_pos_add_nonneg}
    Suppose $a\in\reals$ and $b\in\reals$.
    Suppose $a \geq 0$ and $0 < b$.
    Then $0 < a + b$.
\end{lemma}
\begin{proof}
    $0\in\reals$ by \cref{real_add_ident}.
    $a + b\in\reals$ by \cref{real_add_closed}.
    Thus $0 < a$ or $0 = a$ by \cref{real_leq_split}.
    \begin{byCase}
    \caseOf{$0 = a$.}
        Thus $a = 0$.
        $0 + b = b$ by \cref{real_add_ident}.
        Thus $b = 0 + b$.
        Thus $0 < 0 + b$ by \cref{real_lt_subst_right}.
        Thus $a + b = 0 + b$ by \cref{real_add_eq_subst}.
        Thus $0 + b = a + b$.
        Thus $0 < a + b$ by \cref{real_lt_subst_right}.
    \caseOf{$0 < a$.}
        $0 + b < a + b$ by \cref{real_lt_add}.
        $0 + b = b$ by \cref{real_add_ident}.
        Thus $b < a + b$ by \cref{real_lt_subst_left}.
        Thus $0 < a + b$ by \cref{real_lt_trans}.
    \end{byCase}
\end{proof}

% from §21 helper lemma 6: substitute equality in multiplication
\begin{lemma}\label{real_mul_eq_subst}
    Suppose $a\in\reals$ and $b\in\reals$ and $c\in\reals$.
    Suppose $a = b$.
    Then $a \rmul c = b \rmul c$.
\end{lemma}
\begin{proof}
    Thus $a \rmul c = b \rmul c$.
\end{proof}

% from §21 Item 10: multiplication operation is continuous
\begin{lemma}\label{real_multiplication_continuous}
    Let $T = \productTopology{\standardTopologyR}{\standardTopologyR}{\reals}{\reals}$.
    Then $\realMultiplication$ is continuous from $\reals\times\reals$ to $\reals$ with respect to $T$ and $\standardTopologyR$.
\end{lemma}
\begin{proof}
    Let $f = \realMultiplication$.
    Show that $f$ is a function from $\reals\times\reals$ to $\reals$.
    \begin{subproof}
        Show that $f$ is a function.
        \begin{subproof}
            Show that for all $p,z_1,z_2$ such that $(p,z_1)\in f$ and $(p,z_2)\in f$ we have $z_1 = z_2$.
            \begin{subproof}
                Fix $p$.
                Fix $z_1$.
                Fix $z_2$.
                Assume $(p,z_1)\in f$ and $(p,z_2)\in f$.
                Take $p_1,w_1$ such that
                    $p_1\in\reals\times\reals$ and $w_1\in\reals$ and
                    $(p,z_1) = (p_1,w_1)$ and $w_1 = \fst{p_1} \rmul \snd{p_1}$
                    by \cref{real_multiplication_operation}.
                Take $p_2,w_2$ such that
                    $p_2\in\reals\times\reals$ and $w_2\in\reals$ and
                    $(p,z_2) = (p_2,w_2)$ and $w_2 = \fst{p_2} \rmul \snd{p_2}$
                    by \cref{real_multiplication_operation}.
                $p = p_1$ and $z_1 = w_1$ by \cref{pair_eq_iff}.
                $p = p_2$ and $z_2 = w_2$ by \cref{pair_eq_iff}.
                Thus $z_1 = \fst{p} \rmul \snd{p}$.
                Thus $z_2 = \fst{p} \rmul \snd{p}$.
                Thus $z_1 = z_2$.
            \end{subproof}
            Thus $f$ is right-unique by \cref{rightunique}.
            Show that for all $w$ such that $w\in f$ there exist $x,y$ such that $w = (x,y)$.
            \begin{subproof}
                Fix $w$.
                Assume $w\in f$.
                Take $p,z$ such that $p\in\reals\times\reals$ and $z\in\reals$ and $w = (p,z)$ and
                    $z = \fst{p} \rmul \snd{p}$ by \cref{real_multiplication_operation}.
                Thus there exist $x,y$ such that $w = (x,y)$.
            \end{subproof}
            Thus $f$ is a relation by \cref{relation}.
            Thus $f$ is a function.
        \end{subproof}
        Show that for all $p\in\dom{f}$ we have $f(p)\in\reals$.
        \begin{subproof}
            Fix $p\in\dom{f}$.
            Take $z$ such that $(p,z)\in f$ by \cref{dom_iff}.
            Take $p_0,w$ such that
                $p_0\in\reals\times\reals$ and $w\in\reals$ and
                $(p,z) = (p_0,w)$ and $w = \fst{p_0} \rmul \snd{p_0}$
                by \cref{real_multiplication_operation}.
            $z = w$ by \cref{pair_eq_iff}.
            Thus $z\in\reals$.
            $f$ is a function.
            Thus $f(p) = z$ by \cref{function_apply_intro}.
            Thus $f(p)\in\reals$.
        \end{subproof}
        Thus $f$ is a function to $\reals$.
        Show that $\dom{f} = \reals\times\reals$.
        \begin{subproof}
            Show that for all $p\in\dom{f}$ we have $p\in\reals\times\reals$.
            \begin{subproof}
                Fix $p\in\dom{f}$.
                Take $z$ such that $(p,z)\in f$ by \cref{dom_iff}.
                Take $p_0,w$ such that
                    $p_0\in\reals\times\reals$ and $w\in\reals$ and
                    $(p,z) = (p_0,w)$ and $w = \fst{p_0} \rmul \snd{p_0}$
                    by \cref{real_multiplication_operation}.
                Thus $p = p_0$ by \cref{pair_eq_iff}.
                Thus $p\in\reals\times\reals$.
            \end{subproof}
            Thus $\dom{f}\subseteq\reals\times\reals$ by \cref{subseteq}.
            Show that for all $p\in\reals\times\reals$ we have $p\in\dom{f}$.
            \begin{subproof}
                Fix $p\in\reals\times\reals$.
                Take $x,y$ such that $p = (x,y)$ and $x\in\reals$ and $y\in\reals$ by \cref{times_elem_is_tuple}.
                $x \rmul y\in\reals$ by \cref{real_mul_closed}.
                $\fst{p} = x$ and $\snd{p} = y$ by \cref{fst_eq,snd_eq}.
                Let $z = x \rmul y$.
                Let $w = (p,z)$.
                $w\in f$ by \cref{real_multiplication_operation}.
                Thus $p\in\dom{f}$ by \cref{dom_intro}.
            \end{subproof}
            Thus $\reals\times\reals\subseteq\dom{f}$ by \cref{subseteq}.
            Thus $\dom{f} = \reals\times\reals$ by \cref{subseteq_antisymmetric}.
        \end{subproof}
        Thus $f\in\funs{\reals\times\reals}{\reals}$ by \cref{funs_intro}.
    \end{subproof}
    Show that for all $p\in\reals\times\reals$ we have
        $f$ is continuous at $p$ from $\reals\times\reals$ to $\reals$ with respect to $T$ and $\standardTopologyR$.
    \begin{subproof}
        Fix $p$.
        Assume $p\in\reals\times\reals$.
        Take $x_0,y_0$ such that $p = (x_0,y_0)$ and $x_0\in\reals$ and $y_0\in\reals$ by \cref{times_elem_is_tuple}.
        $x_0 \rmul y_0\in\reals$ by \cref{real_mul_closed}.
        Let $z_0 = x_0 \rmul y_0$.
        $\fst{p} = x_0$ and $\snd{p} = y_0$ by \cref{fst_eq,snd_eq}.
        Let $w_0 = (p,z_0)$.
        $w_0\in f$ by \cref{real_multiplication_operation}.
        $f$ is a function.
        Thus $f(p) = z_0$ by \cref{function_apply_intro}.
        Show that for all $V$ such that $V$ is a neighborhood of $f(p)$ in $\reals$ with respect to $\standardTopologyR$
            there exists $U$ such that $U$ is a neighborhood of $p$ in $\reals\times\reals$ with respect to $T$ and $\img{f}{U}\subseteq V$.
        \begin{subproof}
            Fix $V$.
            Assume $V$ is a neighborhood of $f(p)$ in $\reals$ with respect to $\standardTopologyR$.
            $V$ is open in $\reals$ with respect to $\standardTopologyR$ by \cref{neighborhood}.
            Let $d = \standardMetricR$.
            Let $T_d = \metricTopology{d}{\reals}$.
            $T_d = \standardTopologyR$ by \cref{standard_metric_topology}.
            $\standardTopologyR\subseteq T_d$ by \cref{subseteq}.
            Thus $V\in T_d$ by \cref{open_in,subseteq}.
            $T_d\subseteq\standardTopologyR$ by \cref{subseteq}.
            $f(p)\in V$ by \cref{neighborhood}.
            $V\in\standardTopologyR$ by \cref{open_in,neighborhood}.
            $\standardBasisR$ is a basis on $\reals$ by \cref{standard_basis_is_basis}.
            Thus $\standardTopologyR$ is a topology on $\reals$ by \cref{basis_topology_is_topology,standard_topology_reals}.
            Thus $V\subseteq\reals$ by \cref{topology_on,subseteq,pow_iff}.
            $d$ is a metric on $\reals$ by \cref{standard_metric_is_metric}.
            Take $\epsilon\in\reals$ such that $0 < \epsilon$ and
                $\metricBall{d}{\reals}{f(p)}{\epsilon} \subseteq V$ by \cref{metric_open_iff}.
            Let $m = \abs{x_0} + \abs{y_0} + 1$.
            $\abs{x_0}\in\reals$ by \cref{real_abs_in_reals}.
            $\abs{y_0}\in\reals$ by \cref{real_abs_in_reals}.
            $\abs{x_0} + \abs{y_0}\in\reals$ by \cref{real_add_closed}.
            $1\in\reals$ by \cref{real_mul_ident}.
            Thus $m\in\reals$ by \cref{real_add_closed}.
            $\abs{y_0} \geq 0$ by \cref{real_abs_nonnegative}.
            $0 < 1$ by \cref{real_zero_lt_one}.
            $\abs{y_0} + 1\in\reals$ by \cref{real_add_closed}.
            Thus $0 < \abs{y_0} + 1$ by \cref{real_pos_add_nonneg}.
            $\abs{x_0} \geq 0$ by \cref{real_abs_nonnegative}.
            Thus $0 < \abs{x_0} + (\abs{y_0} + 1)$ by \cref{real_pos_add_nonneg}.
            $(\abs{x_0} + \abs{y_0}) + 1 = \abs{x_0} + (\abs{y_0} + 1)$ by \cref{real_add_assoc}.
            Thus $0 < m$ by \cref{real_lt_subst_right}.
            Thus $m \geq 0$ by \cref{real_lt_imp_leq}.
            Let $\delta = \epsilon / (m + \epsilon)$.
            $m + \epsilon\in\reals$ by \cref{real_add_closed}.
            $0 < m + \epsilon$ by \cref{real_pos_add_nonneg}.
            Thus $(m + \epsilon)\neq 0$ by \cref{real_pos_nonzero}.
            Thus $\inv{(m + \epsilon)}\in\reals$ by \cref{real_mul_inverse}.
            $0 < \inv{(m + \epsilon)}$ by \cref{real_inv_pos}.
            $0\in\reals$ by \cref{real_add_ident}.
            $0 \rmul \inv{(m + \epsilon)} = 0$ by \cref{real_mul_zero}.
            Thus $0 \rmul \inv{(m + \epsilon)} < \epsilon \rmul \inv{(m + \epsilon)}$ by \cref{real_lt_mul_pos}.
            Thus $0 < \epsilon \rmul \inv{(m + \epsilon)}$ by \cref{real_lt_subst_left}.
            $\delta = \epsilon \rmul \inv{(m + \epsilon)}$ by \cref{real_div}.
            Thus $0 < \delta$ by \cref{real_lt_subst_left}.
            $0 + \epsilon < m + \epsilon$ by \cref{real_lt_add}.
            $0 + \epsilon = \epsilon$ by \cref{real_add_ident}.
            Thus $\epsilon < m + \epsilon$ by \cref{real_lt_subst_left}.
            $\epsilon \rmul \inv{(m + \epsilon)} < (m + \epsilon) \rmul \inv{(m + \epsilon)}$ by \cref{real_lt_mul_pos}.
            $(m + \epsilon) \rmul \inv{(m + \epsilon)} = 1$ by \cref{real_mul_inverse}.
            Thus $\epsilon \rmul \inv{(m + \epsilon)} < 1$ by \cref{real_lt_subst_right}.
            $\delta = \epsilon \rmul \inv{(m + \epsilon)}$ by \cref{real_div}.
            Thus $\delta < 1$ by \cref{real_lt_subst_left}.
            Let $U_1 = \metricBall{d}{\reals}{x_0}{\delta}$.
            Let $U_2 = \metricBall{d}{\reals}{y_0}{\delta}$.
            Show that $U_1$ is open in $\reals$ with respect to $\standardTopologyR$.
            \begin{subproof}
                $x_0\in\reals$ and $\delta\in\reals$ and $0 < \delta$.
                Show that $U_1\in\pow{\reals}$.
                \begin{subproof}
                    Show that for all $x$ such that $x\in U_1$ we have $x\in\reals$.
                    \begin{subproof}
                        Fix $x$.
                        Assume $x\in U_1$.
                        Thus $x\in\reals$ by \cref{metric_ball}.
                    \end{subproof}
                    Thus $U_1\subseteq\reals$ by \cref{subseteq}.
                \end{subproof}
                Thus $U_1\in\pow{\reals}$ by \cref{pow_iff}.
                Thus $U_1\in\metricBasis{d}{\reals}$ by \cref{metric_basis}.
                $\metricBasis{d}{\reals}$ is a basis on $\reals$ by \cref{metric_basis_is_basis}.
                Thus $U_1\in T_d$ by \cref{basis_elem_open_in_generated,metric_topology,basis_for_topology}.
                Thus $U_1\in\standardTopologyR$ by \cref{subseteq}.
                Thus $U_1$ is open in $\reals$ with respect to $\standardTopologyR$ by \cref{open_in}.
            \end{subproof}
            Show that $U_2$ is open in $\reals$ with respect to $\standardTopologyR$.
            \begin{subproof}
                $y_0\in\reals$ and $\delta\in\reals$ and $0 < \delta$.
                Show that $U_2\in\pow{\reals}$.
                \begin{subproof}
                    Show that for all $y$ such that $y\in U_2$ we have $y\in\reals$.
                    \begin{subproof}
                        Fix $y$.
                        Assume $y\in U_2$.
                        Thus $y\in\reals$ by \cref{metric_ball}.
                    \end{subproof}
                    Thus $U_2\subseteq\reals$ by \cref{subseteq}.
                \end{subproof}
                Thus $U_2\in\pow{\reals}$ by \cref{pow_iff}.
                Thus $U_2\in\metricBasis{d}{\reals}$ by \cref{metric_basis}.
                $\metricBasis{d}{\reals}$ is a basis on $\reals$ by \cref{metric_basis_is_basis}.
                Thus $U_2\in T_d$ by \cref{basis_elem_open_in_generated,metric_topology,basis_for_topology}.
                Thus $U_2\in\standardTopologyR$ by \cref{subseteq}.
                Thus $U_2$ is open in $\reals$ with respect to $\standardTopologyR$ by \cref{open_in}.
            \end{subproof}
            Let $U = U_1\times U_2$.
            Show that $U$ is open in $\reals\times\reals$ with respect to $T$.
            \begin{subproof}
                $U_1\in\standardTopologyR$ by \cref{open_in}.
                $U_2\in\standardTopologyR$ by \cref{open_in}.
                $\standardTopologyR\subseteq\pow{\reals}$ by \cref{topology_on}.
                Thus $U_1\in\pow{\reals}$ and $U_2\in\pow{\reals}$ by \cref{subseteq}.
                Thus $U_1\subseteq\reals$ and $U_2\subseteq\reals$ by \cref{pow_iff}.
                $U_1\times U_2\subseteq\reals\times U_2$ by \cref{times_subseteq_left}.
                $\reals\times U_2\subseteq\reals\times\reals$ by \cref{times_subseteq_right}.
                Thus $U\subseteq\reals\times\reals$ by \cref{subseteq_transitive}.
                $U\in\pow{\reals\times\reals}$ by \cref{pow_iff}.
                Thus $U\in\productBasis{\standardTopologyR}{\standardTopologyR}{\reals}{\reals}$ by \cref{product_basis}.
                $\productBasis{\standardTopologyR}{\standardTopologyR}{\reals}{\reals}$ is a basis on $\reals\times\reals$ by \cref{product_basis_is_basis}.
                Thus $T$ is a topology on $\reals\times\reals$ by \cref{basis_topology_is_topology,product_topology}.
                Thus $U\in T$ by \cref{basis_elem_open_in_generated,product_topology}.
                Thus $U$ is open in $\reals\times\reals$ with respect to $T$ by \cref{open_in}.
            \end{subproof}
            Show that $p\in U$.
            \begin{subproof}
                $d((x_0,x_0)) = 0$ by \cref{metric_zero_characterization,standard_metric_is_metric,metric_on}.
                Thus $d((x_0,x_0)) < \delta$ by \cref{real_lt_subst_left}.
                Thus $x_0\in U_1$ by \cref{metric_ball}.
                $d((y_0,y_0)) = 0$ by \cref{metric_zero_characterization,standard_metric_is_metric,metric_on}.
                Thus $d((y_0,y_0)) < \delta$ by \cref{real_lt_subst_left}.
                Thus $y_0\in U_2$ by \cref{metric_ball}.
                Thus $p\in U$ by \cref{times_tuple_intro}.
            \end{subproof}
            Show that $\img{f}{U}\subseteq V$.
            \begin{subproof}
                Show that for all $q$ such that $q\in U$ we have
                    $f(q)\in\metricBall{d}{\reals}{f(p)}{\epsilon}$.
                \begin{subproof}
                    Fix $q$.
                    Assume $q\in U$.
                    Take $x,y$ such that $q = (x,y)$ and $x\in U_1$ and $y\in U_2$ by \cref{times_elem_is_tuple}.
                    $x\in\reals$ and $y\in\reals$ by \cref{metric_ball}.
                    $d((x_0,x)) < \delta$ by \cref{metric_ball}.
                    $d((y_0,y)) < \delta$ by \cref{metric_ball}.
                    $d((x_0,x)) = \abs{x_0 - x}$ by \cref{standard_metric_value}.
                    $d((y_0,y)) = \abs{y_0 - y}$ by \cref{standard_metric_value}.
                    Thus $\abs{x_0 - x} < \delta$ and $\abs{y_0 - y} < \delta$.
                    $\abs{x - x_0} = \abs{x_0 - x}$ by \cref{real_abs_sub_swap}.
                    Thus $\abs{x - x_0} < \delta$ by \cref{real_lt_subst_left}.
                    $\abs{y - y_0} = \abs{y_0 - y}$ by \cref{real_abs_sub_swap}.
                    Thus $\abs{y - y_0} < \delta$ by \cref{real_lt_subst_left}.
                    Show that $f(q) = x \rmul y$.
                    \begin{subproof}
                        $x \rmul y\in\reals$ by \cref{real_mul_closed}.
                        $q\in\reals\times\reals$ by \cref{times_tuple_intro}.
                        $\fst{q} = x$ and $\snd{q} = y$ by \cref{fst_eq,snd_eq}.
                        Let $w = (q, x \rmul y)$.
                        $w\in f$ by \cref{real_multiplication_operation}.
                        $f$ is a function.
                        Thus $f(q) = x \rmul y$ by \cref{function_apply_intro}.
                    \end{subproof}
                    $f(p) = x_0 \rmul y_0$.
                    Show that $\abs{f(q) - f(p)} < \epsilon$.
                    \begin{subproof}
                        $x - x_0 = x + \neg{x_0}$ by \cref{real_sub}.
                        $y - y_0 = y + \neg{y_0}$ by \cref{real_sub}.
                        $\neg{x_0}\in\reals$ by \cref{real_add_inverse}.
                        $\neg{y_0}\in\reals$ by \cref{real_add_inverse}.
                        Thus $(x - x_0) \rmul y = (x + \neg{x_0}) \rmul y$.
                        Thus $x_0 \rmul (y - y_0) = x_0 \rmul (y + \neg{y_0})$.
                        $y + \neg{y_0}\in\reals$ by \cref{real_add_closed}.
                        $(x + \neg{x_0}) \rmul y = (x \rmul y) + (\neg{x_0} \rmul y)$ by \cref{real_right_distrib}.
                        $x_0 \rmul (y + \neg{y_0}) = (y + \neg{y_0}) \rmul x_0$ by \cref{real_mul_comm}.
                        $(y + \neg{y_0}) \rmul x_0 = (y \rmul x_0) + (\neg{y_0} \rmul x_0)$ by \cref{real_right_distrib}.
                        $y \rmul x_0 = x_0 \rmul y$ by \cref{real_mul_comm}.
                        $\neg{y_0} \rmul x_0 = x_0 \rmul \neg{y_0}$ by \cref{real_mul_comm}.
                        $x_0 \rmul \neg{y_0} = \neg{(x_0 \rmul y_0)}$ by \cref{real_mul_neg_right}.
                        Thus $(x - x_0) \rmul y + x_0 \rmul (y - y_0) = (x \rmul y) + (\neg{x_0} \rmul y) + (x_0 \rmul y) + \neg{(x_0 \rmul y_0)}$.
                        $(\neg{x_0} \rmul y) + (x_0 \rmul y) = (\neg{x_0} + x_0) \rmul y$ by \cref{real_right_distrib}.
                        $\neg{x_0} + x_0 = 0$ by \cref{real_add_inverse,real_add_comm}.
                        $0 \rmul y = 0$ by \cref{real_mul_zero}.
                        Thus $(\neg{x_0} \rmul y) + (x_0 \rmul y) = 0$.
                        $0\in\reals$ by \cref{real_add_ident}.
                        $x \rmul y\in\reals$ by \cref{real_mul_closed}.
                        $\neg{x_0} \rmul y\in\reals$ by \cref{real_mul_closed}.
                        $x_0 \rmul y\in\reals$ by \cref{real_mul_closed}.
                        $(\neg{x_0} \rmul y) + (x_0 \rmul y)\in\reals$ by \cref{real_add_closed}.
                        Show that $(x \rmul y) + (\neg{x_0} \rmul y) + (x_0 \rmul y) = x \rmul y$.
                        \begin{subproof}
                            $((x \rmul y) + (\neg{x_0} \rmul y)) + (x_0 \rmul y) = (x \rmul y) + ((\neg{x_0} \rmul y) + (x_0 \rmul y))$ by \cref{real_add_assoc}.
                            $(x \rmul y) + ((\neg{x_0} \rmul y) + (x_0 \rmul y)) = ((\neg{x_0} \rmul y) + (x_0 \rmul y)) + (x \rmul y)$ by \cref{real_add_comm}.
                            $((\neg{x_0} \rmul y) + (x_0 \rmul y)) + (x \rmul y) = 0 + (x \rmul y)$ by \cref{real_add_eq_subst}.
                            $0 + (x \rmul y) = x \rmul y$ by \cref{real_add_ident}.
                            Thus $((x \rmul y) + (\neg{x_0} \rmul y)) + (x_0 \rmul y) = x \rmul y$.
                            Thus $(x \rmul y) + (\neg{x_0} \rmul y) + (x_0 \rmul y) = x \rmul y$.
                        \end{subproof}
                        $(x \rmul y) + (\neg{x_0} \rmul y)\in\reals$ by \cref{real_add_closed}.
                        $(x \rmul y) + (\neg{x_0} \rmul y) + (x_0 \rmul y)\in\reals$ by \cref{real_add_closed}.
                        $\neg{(x_0 \rmul y_0)}\in\reals$ by \cref{real_add_inverse}.
                        Thus $(x \rmul y) + (\neg{x_0} \rmul y) + (x_0 \rmul y) + \neg{(x_0 \rmul y_0)} = (x \rmul y) + \neg{(x_0 \rmul y_0)}$ by \cref{real_add_eq_subst}.
                        Thus $(x - x_0) \rmul y + x_0 \rmul (y - y_0) = (x \rmul y) + \neg{(x_0 \rmul y_0)}$.
                        $(x \rmul y) + \neg{(x_0 \rmul y_0)} = (x \rmul y) - (x_0 \rmul y_0)$ by \cref{real_sub}.
                        Thus $(x \rmul y) - (x_0 \rmul y_0) = (x - x_0) \rmul y + x_0 \rmul (y - y_0)$.
                        $x - x_0\in\reals$ by \cref{real_sub,real_add_closed}.
                        $y - y_0\in\reals$ by \cref{real_sub,real_add_closed}.
                        $(x - x_0) \rmul y\in\reals$ by \cref{real_mul_closed}.
                        $x_0 \rmul (y - y_0)\in\reals$ by \cref{real_mul_closed}.
                        $\abs{(x - x_0) \rmul y + x_0 \rmul (y - y_0)} \leq \abs{(x - x_0) \rmul y} + \abs{x_0 \rmul (y - y_0)}$ by \cref{real_triangle_inequality}.
                        $\abs{(x - x_0) \rmul y} = \abs{x - x_0} \rmul \abs{y}$ by \cref{real_abs_mul}.
                        $\abs{x_0 \rmul (y - y_0)} = \abs{x_0} \rmul \abs{y - y_0}$ by \cref{real_abs_mul}.
                        Thus $\abs{(x - x_0) \rmul y + x_0 \rmul (y - y_0)} \leq \abs{x - x_0} \rmul \abs{y} + \abs{x_0} \rmul \abs{y - y_0}$.
                        $\delta\in\reals$ by \cref{real_mul_closed}.
                        $\abs{y} < \abs{y_0} + \delta$ by \cref{real_abs_sub_lt_bound}.
                        $\abs{y_0} \geq 0$ by \cref{real_abs_nonnegative}.
                        Thus $0 < \abs{y_0} + \delta$ by \cref{real_pos_add_nonneg}.
                        $\abs{x - x_0} \geq 0$ by \cref{real_abs_nonnegative}.
                        Thus $0 < \abs{x - x_0}$ or $0 = \abs{x - x_0}$ by \cref{real_leq_split}.
                        Show that $\abs{x - x_0} \rmul \abs{y} < \delta \rmul (\abs{y_0} + \delta)$.
                        \begin{subproof}
                            \begin{byCase}
                            \caseOf{$0 = \abs{x - x_0}$.}
                                Thus $\abs{x - x_0} = 0$ by assumption.
                                $\abs{x - x_0}\in\reals$ by \cref{real_abs_in_reals}.
                                $\abs{y}\in\reals$ by \cref{real_abs_in_reals}.
                                $0\in\reals$ by \cref{real_add_ident}.
                                $\abs{x - x_0} \rmul \abs{y} = 0 \rmul \abs{y}$ by \cref{real_mul_eq_subst}.
                                $0 \rmul \abs{y} = 0$ by \cref{real_mul_zero}.
                                Thus $\abs{x - x_0} \rmul \abs{y} = 0$.
                                $\delta$ is positive.
                                $\abs{y_0} + \delta$ is positive.
                                $\abs{y_0}\in\reals$ by \cref{real_abs_in_reals}.
                                $\abs{y_0} + \delta\in\reals$ by \cref{real_add_closed}.
                                Thus $\delta \rmul (\abs{y_0} + \delta)$ is positive by \cref{real_mul_pos_pos}.
                                Thus $0 < \delta \rmul (\abs{y_0} + \delta)$.
                                Thus $\abs{x - x_0} \rmul \abs{y} < \delta \rmul (\abs{y_0} + \delta)$ by \cref{real_lt_subst_left}.
                            \caseOf{$0 < \abs{x - x_0}$.}
                                $\abs{y}\in\reals$ by \cref{real_abs_in_reals}.
                                $\abs{y_0}\in\reals$ by \cref{real_abs_in_reals}.
                                $\abs{x - x_0}\in\reals$ by \cref{real_abs_in_reals}.
                                $\abs{y_0} + \delta\in\reals$ by \cref{real_add_closed}.
                                $\abs{x - x_0} \rmul \abs{y}\in\reals$ by \cref{real_mul_closed}.
                                $\abs{x - x_0} \rmul (\abs{y_0} + \delta)\in\reals$ by \cref{real_mul_closed}.
                                $\delta \rmul (\abs{y_0} + \delta)\in\reals$ by \cref{real_mul_closed}.
                                $\abs{y} \rmul \abs{x - x_0} < (\abs{y_0} + \delta) \rmul \abs{x - x_0}$ by \cref{real_lt_mul_pos}.
                                $\abs{y} \rmul \abs{x - x_0} = \abs{x - x_0} \rmul \abs{y}$ by \cref{real_mul_comm}.
                                $(\abs{y_0} + \delta) \rmul \abs{x - x_0} = \abs{x - x_0} \rmul (\abs{y_0} + \delta)$ by \cref{real_mul_comm}.
                                Thus $\abs{x - x_0} \rmul \abs{y} < \abs{x - x_0} \rmul (\abs{y_0} + \delta)$.
                                $\abs{x - x_0} \rmul (\abs{y_0} + \delta) < \delta \rmul (\abs{y_0} + \delta)$ by \cref{real_lt_mul_pos}.
                                Thus $\abs{x - x_0} \rmul \abs{y} < \delta \rmul (\abs{y_0} + \delta)$ by \cref{real_lt_trans}.
                            \end{byCase}
                        \end{subproof}
                        $\abs{x_0} \geq 0$ by \cref{real_abs_nonnegative}.
                        Thus $0 < \abs{x_0}$ or $0 = \abs{x_0}$ by \cref{real_leq_split}.
                        $\abs{x_0} \geq 0$ by \cref{real_abs_nonnegative}.
                        Thus $0 < \abs{x_0}$ or $0 = \abs{x_0}$ by \cref{real_leq_split}.
                        Show that $\abs{x - x_0} \rmul \abs{y} + \abs{x_0} \rmul \abs{y - y_0} < \delta \rmul m$.
                        \begin{subproof}
                            \begin{byCase}
                            \caseOf{$0 = \abs{x_0}$.}
                                Thus $\abs{x_0} = 0$ by assumption.
                                $\abs{x_0}\in\reals$ by \cref{real_abs_in_reals}.
                                $\abs{y - y_0}\in\reals$ by \cref{real_abs_in_reals}.
                                $0\in\reals$ by \cref{real_add_ident}.
                                $\abs{x_0} \rmul \abs{y - y_0} = 0 \rmul \abs{y - y_0}$ by \cref{real_mul_eq_subst}.
                                $0 \rmul \abs{y - y_0} = 0$ by \cref{real_mul_zero}.
                                Thus $\abs{x_0} \rmul \abs{y - y_0} = 0$.
                                $\abs{x - x_0}\in\reals$ by \cref{real_abs_in_reals}.
                                $\abs{y}\in\reals$ by \cref{real_abs_in_reals}.
                                $\abs{x - x_0} \rmul \abs{y}\in\reals$ by \cref{real_mul_closed}.
                                Thus $\abs{x - x_0} \rmul \abs{y} + \abs{x_0} \rmul \abs{y - y_0} = \abs{x - x_0} \rmul \abs{y} + 0$ by \cref{real_add_eq_subst}.
                                $\abs{x - x_0} \rmul \abs{y} + 0 = 0 + \abs{x - x_0} \rmul \abs{y}$ by \cref{real_add_comm}.
                                $0 + \abs{x - x_0} \rmul \abs{y} = \abs{x - x_0} \rmul \abs{y}$ by \cref{real_add_ident}.
                                Thus $\abs{x - x_0} \rmul \abs{y} + \abs{x_0} \rmul \abs{y - y_0} = \abs{x - x_0} \rmul \abs{y}$.
                                Thus $\abs{x - x_0} \rmul \abs{y} < \delta \rmul (\abs{y_0} + \delta)$.
                                $\delta + \abs{y_0} < 1 + \abs{y_0}$ by \cref{real_lt_add}.
                                Thus $\abs{y_0} + \delta < \abs{y_0} + 1$ by \cref{real_add_comm,real_add_assoc}.
                                $\abs{y_0} + \delta\in\reals$ by \cref{real_add_closed}.
                                $\abs{y_0} + 1\in\reals$ by \cref{real_add_closed}.
                                $\delta\in\reals$ by \cref{real_mul_closed}.
                                $\delta$ is positive.
                                $\delta \rmul (\abs{y_0} + \delta) = (\abs{y_0} + \delta) \rmul \delta$ by \cref{real_mul_comm}.
                                $\delta \rmul (\abs{y_0} + 1) = (\abs{y_0} + 1) \rmul \delta$ by \cref{real_mul_comm}.
                                $(\abs{y_0} + \delta) \rmul \delta < (\abs{y_0} + 1) \rmul \delta$ by \cref{real_lt_mul_pos}.
                                Thus $\delta \rmul (\abs{y_0} + \delta) < (\abs{y_0} + 1) \rmul \delta$ by \cref{real_lt_subst_left}.
                                Thus $\delta \rmul (\abs{y_0} + \delta) < \delta \rmul (\abs{y_0} + 1)$ by \cref{real_lt_subst_right}.
                                $\delta \rmul (\abs{y_0} + \delta)\in\reals$ by \cref{real_mul_closed}.
                                $\delta \rmul (\abs{y_0} + 1)\in\reals$ by \cref{real_mul_closed}.
                                Thus $\abs{x - x_0} \rmul \abs{y} < \delta \rmul (\abs{y_0} + 1)$ by \cref{real_lt_trans}.
                                Thus $\abs{x_0} = 0$.
                                Thus $\abs{x_0} + \abs{y_0} + 1 = \abs{y_0} + 1$ by \cref{real_add_ident,real_add_comm,real_add_assoc,real_add_eq_subst}.
                                $\delta \rmul (\abs{y_0} + 1) = (\abs{y_0} + 1) \rmul \delta$ by \cref{real_mul_comm}.
                                $(\abs{y_0} + 1) \rmul \delta = (\abs{x_0} + \abs{y_0} + 1) \rmul \delta$ by \cref{real_mul_eq_subst}.
                                Thus $\delta \rmul (\abs{y_0} + 1) = \delta \rmul (\abs{x_0} + \abs{y_0} + 1)$ by \cref{real_mul_comm}.
                                Thus $\abs{x - x_0} \rmul \abs{y} + \abs{x_0} \rmul \abs{y - y_0} < \delta \rmul (\abs{x_0} + \abs{y_0} + 1)$ by \cref{real_lt_subst_right}.
                                Thus $\abs{x - x_0} \rmul \abs{y} + \abs{x_0} \rmul \abs{y - y_0} < \delta \rmul m$ by \cref{real_lt_subst_right}.
                            \caseOf{$0 < \abs{x_0}$.}
                                $\abs{y - y_0}\in\reals$ by \cref{real_abs_in_reals}.
                                $\abs{x_0}\in\reals$ by \cref{real_abs_in_reals}.
                                $\delta\in\reals$ by \cref{real_mul_closed}.
                                $\abs{y - y_0} \rmul \abs{x_0} < \delta \rmul \abs{x_0}$ by \cref{real_lt_mul_pos}.
                                $\abs{y - y_0} \rmul \abs{x_0} = \abs{x_0} \rmul \abs{y - y_0}$ by \cref{real_mul_comm}.
                                $\delta \rmul \abs{x_0} = \abs{x_0} \rmul \delta$ by \cref{real_mul_comm}.
                                Thus $\abs{x_0} \rmul \abs{y - y_0} < \delta \rmul \abs{x_0}$ by \cref{real_lt_subst_left}.
                                Thus $\abs{x_0} \rmul \abs{y - y_0} < \abs{x_0} \rmul \delta$ by \cref{real_lt_subst_right}.
                                $\abs{x - x_0}\in\reals$ by \cref{real_abs_in_reals}.
                                $\abs{y}\in\reals$ by \cref{real_abs_in_reals}.
                                $\abs{x - x_0} \rmul \abs{y}\in\reals$ by \cref{real_mul_closed}.
                                $\abs{y_0}\in\reals$ by \cref{real_abs_in_reals}.
                                $\abs{y_0} + \delta\in\reals$ by \cref{real_add_closed}.
                                $\delta \rmul (\abs{y_0} + \delta)\in\reals$ by \cref{real_mul_closed}.
                                $\abs{y - y_0}\in\reals$ by \cref{real_abs_in_reals}.
                                $\abs{x_0}\in\reals$ by \cref{real_abs_in_reals}.
                                $\abs{x_0} \rmul \abs{y - y_0}\in\reals$ by \cref{real_mul_closed}.
                                $\abs{x_0} \rmul \delta\in\reals$ by \cref{real_mul_closed}.
                                Thus $\abs{x - x_0} \rmul \abs{y} + \abs{x_0} \rmul \abs{y - y_0} < \delta \rmul (\abs{y_0} + \delta) + \abs{x_0} \rmul \delta$ by \cref{real_lt_add_two}.
                                $(\abs{x_0} + \abs{y_0}) + \delta = \abs{x_0} + (\abs{y_0} + \delta)$ by \cref{real_add_assoc}.
                                $\delta \rmul (\abs{y_0} + \delta) = (\abs{y_0} + \delta) \rmul \delta$ by \cref{real_mul_comm}.
                                $(\abs{y_0} + \delta) \rmul \delta = (\abs{y_0} \rmul \delta) + (\delta \rmul \delta)$ by \cref{real_right_distrib}.
                                $\abs{x_0} \rmul \delta = \delta \rmul \abs{x_0}$ by \cref{real_mul_comm}.
                                $\abs{y_0} \rmul \delta = \delta \rmul \abs{y_0}$ by \cref{real_mul_comm}.
                                Show that $\delta \rmul (\abs{y_0} + \delta) + \abs{x_0} \rmul \delta = (\delta \rmul \abs{x_0}) + (\delta \rmul \abs{y_0}) + (\delta \rmul \delta)$.
                                \begin{subproof}
                                    Thus $\delta \rmul (\abs{y_0} + \delta) + \abs{x_0} \rmul \delta = (\abs{y_0} \rmul \delta + \delta \rmul \delta) + \abs{x_0} \rmul \delta$ by \cref{real_add_eq_subst,real_right_distrib,real_mul_comm}.
                                    $\abs{y_0} \rmul \delta\in\reals$ by \cref{real_mul_closed}.
                                    $\delta \rmul \delta\in\reals$ by \cref{real_mul_closed}.
                                    $\abs{x_0} \rmul \delta\in\reals$ by \cref{real_mul_closed}.
                                    $(\abs{y_0} \rmul \delta + \delta \rmul \delta) + \abs{x_0} \rmul \delta = \abs{y_0} \rmul \delta + (\delta \rmul \delta + \abs{x_0} \rmul \delta)$ by \cref{real_add_assoc}.
                                    $\delta \rmul \delta + \abs{x_0} \rmul \delta = \abs{x_0} \rmul \delta + \delta \rmul \delta$ by \cref{real_add_comm}.
                                    Thus $\abs{y_0} \rmul \delta + (\delta \rmul \delta + \abs{x_0} \rmul \delta) = \abs{y_0} \rmul \delta + (\abs{x_0} \rmul \delta + \delta \rmul \delta)$ by \cref{real_add_eq_subst}.
                                    $\abs{y_0} \rmul \delta + (\abs{x_0} \rmul \delta + \delta \rmul \delta) = (\abs{y_0} \rmul \delta + \abs{x_0} \rmul \delta) + \delta \rmul \delta$ by \cref{real_add_assoc}.
                                    $\abs{y_0} \rmul \delta + \abs{x_0} \rmul \delta = \abs{x_0} \rmul \delta + \abs{y_0} \rmul \delta$ by \cref{real_add_comm}.
                                    Thus $(\abs{y_0} \rmul \delta + \abs{x_0} \rmul \delta) + \delta \rmul \delta = (\abs{x_0} \rmul \delta + \abs{y_0} \rmul \delta) + \delta \rmul \delta$ by \cref{real_add_eq_subst}.
                                    Thus $\delta \rmul (\abs{y_0} + \delta) + \abs{x_0} \rmul \delta = (\abs{x_0} \rmul \delta + \abs{y_0} \rmul \delta) + \delta \rmul \delta$.
                                    $\abs{x_0} \rmul \delta = \delta \rmul \abs{x_0}$ by \cref{real_mul_comm}.
                                    $\abs{y_0} \rmul \delta = \delta \rmul \abs{y_0}$ by \cref{real_mul_comm}.
                                    Thus $\abs{x_0} \rmul \delta + \abs{y_0} \rmul \delta = \delta \rmul \abs{x_0} + \delta \rmul \abs{y_0}$ by \cref{real_add_eq_subst}.
                                    Thus $\delta \rmul (\abs{y_0} + \delta) + \abs{x_0} \rmul \delta = (\delta \rmul \abs{x_0} + \delta \rmul \abs{y_0}) + \delta \rmul \delta$ by \cref{real_add_eq_subst}.
                                    Thus $\delta \rmul (\abs{y_0} + \delta) + \abs{x_0} \rmul \delta = (\delta \rmul \abs{x_0}) + (\delta \rmul \abs{y_0}) + (\delta \rmul \delta)$ by \cref{real_add_assoc}.
                                \end{subproof}
                                Show that $(\delta \rmul \abs{x_0}) + (\delta \rmul \abs{y_0}) = \delta \rmul (\abs{x_0} + \abs{y_0})$.
                                \begin{subproof}
                                    $\delta \rmul \abs{x_0} = \abs{x_0} \rmul \delta$ by \cref{real_mul_comm}.
                                    $\delta \rmul \abs{y_0} = \abs{y_0} \rmul \delta$ by \cref{real_mul_comm}.
                                    Thus $(\delta \rmul \abs{x_0}) + (\delta \rmul \abs{y_0}) = \abs{x_0} \rmul \delta + \abs{y_0} \rmul \delta$ by \cref{real_add_eq_subst}.
                                    $\abs{x_0} \rmul \delta + \abs{y_0} \rmul \delta = (\abs{x_0} + \abs{y_0}) \rmul \delta$ by \cref{real_right_distrib}.
                                    $(\abs{x_0} + \abs{y_0}) \rmul \delta = \delta \rmul (\abs{x_0} + \abs{y_0})$ by \cref{real_mul_comm}.
                                    Thus $(\delta \rmul \abs{x_0}) + (\delta \rmul \abs{y_0}) = \delta \rmul (\abs{x_0} + \abs{y_0})$.
                                \end{subproof}
                                Thus $\delta \rmul (\abs{y_0} + \delta) + \abs{x_0} \rmul \delta = (\delta \rmul (\abs{x_0} + \abs{y_0})) + (\delta \rmul \delta)$ by \cref{real_add_assoc}.
                                Show that $(\delta \rmul (\abs{x_0} + \abs{y_0})) + (\delta \rmul \delta) = \delta \rmul ((\abs{x_0} + \abs{y_0}) + \delta)$.
                                \begin{subproof}
                                    $\abs{x_0}\in\reals$ by \cref{real_abs_in_reals}.
                                    $\abs{y_0}\in\reals$ by \cref{real_abs_in_reals}.
                                    $\abs{x_0} + \abs{y_0}\in\reals$ by \cref{real_add_closed}.
                                    $\delta\in\reals$ by \cref{real_mul_closed}.
                                    $(\abs{x_0} + \abs{y_0}) + \delta\in\reals$ by \cref{real_add_closed}.
                                    $\delta \rmul ((\abs{x_0} + \abs{y_0}) + \delta) = ((\abs{x_0} + \abs{y_0}) + \delta) \rmul \delta$ by \cref{real_mul_comm}.
                                    $((\abs{x_0} + \abs{y_0}) + \delta) \rmul \delta = (\abs{x_0} + \abs{y_0}) \rmul \delta + \delta \rmul \delta$ by \cref{real_right_distrib}.
                                    $(\abs{x_0} + \abs{y_0}) \rmul \delta = \delta \rmul (\abs{x_0} + \abs{y_0})$ by \cref{real_mul_comm}.
                                    Thus $\delta \rmul ((\abs{x_0} + \abs{y_0}) + \delta) = \delta \rmul (\abs{x_0} + \abs{y_0}) + \delta \rmul \delta$ by \cref{real_add_eq_subst}.
                                    Thus $(\delta \rmul (\abs{x_0} + \abs{y_0})) + (\delta \rmul \delta) = \delta \rmul ((\abs{x_0} + \abs{y_0}) + \delta)$.
                                \end{subproof}
                                Thus $\delta \rmul (\abs{y_0} + \delta) + \abs{x_0} \rmul \delta = \delta \rmul ((\abs{x_0} + \abs{y_0}) + \delta)$.
                                $\abs{x_0} + \abs{y_0} = \abs{y_0} + \abs{x_0}$ by \cref{real_add_comm}.
                                Thus $(\abs{x_0} + \abs{y_0}) + \delta = (\abs{y_0} + \abs{x_0}) + \delta$ by \cref{real_add_eq_subst}.
                                Thus $\delta \rmul ((\abs{x_0} + \abs{y_0}) + \delta) = \delta \rmul (\abs{x_0} + \abs{y_0} + \delta)$ by \cref{real_add_assoc}.
                                Thus $\abs{x - x_0} \rmul \abs{y} + \abs{x_0} \rmul \abs{y - y_0} < \delta \rmul (\abs{x_0} + \abs{y_0} + \delta)$ by \cref{real_lt_subst_right}.
                                $\delta + (\abs{x_0} + \abs{y_0}) < 1 + (\abs{x_0} + \abs{y_0})$ by \cref{real_lt_add}.
                                Thus $\abs{x_0} + \abs{y_0} + \delta < \abs{x_0} + \abs{y_0} + 1$ by \cref{real_add_comm,real_add_assoc}.
                                $\abs{x_0}\in\reals$ by \cref{real_abs_in_reals}.
                                $\abs{y_0}\in\reals$ by \cref{real_abs_in_reals}.
                                $\abs{x_0} + \abs{y_0}\in\reals$ by \cref{real_add_closed}.
                                $\abs{x_0} + \abs{y_0} + \delta\in\reals$ by \cref{real_add_closed}.
                                $1\in\reals$ by \cref{real_mul_ident}.
                                $\abs{x_0} + \abs{y_0} + 1\in\reals$ by \cref{real_add_closed}.
                                $(\abs{x_0} + \abs{y_0} + \delta) \rmul \delta < (\abs{x_0} + \abs{y_0} + 1) \rmul \delta$ by \cref{real_lt_mul_pos}.
                                $(\abs{x_0} + \abs{y_0} + \delta) \rmul \delta = \delta \rmul (\abs{x_0} + \abs{y_0} + \delta)$ by \cref{real_mul_comm}.
                                Thus $\delta \rmul (\abs{x_0} + \abs{y_0} + \delta) < (\abs{x_0} + \abs{y_0} + 1) \rmul \delta$ by \cref{real_lt_subst_left}.
                                $(\abs{x_0} + \abs{y_0} + 1) \rmul \delta = \delta \rmul (\abs{x_0} + \abs{y_0} + 1)$ by \cref{real_mul_comm}.
                                Thus $\delta \rmul (\abs{x_0} + \abs{y_0} + \delta) < \delta \rmul (\abs{x_0} + \abs{y_0} + 1)$ by \cref{real_lt_subst_right}.
                                $\abs{x - x_0} \rmul \abs{y}\in\reals$ by \cref{real_mul_closed}.
                                $\abs{x_0} \rmul \abs{y - y_0}\in\reals$ by \cref{real_mul_closed}.
                                $\abs{x - x_0} \rmul \abs{y} + \abs{x_0} \rmul \abs{y - y_0}\in\reals$ by \cref{real_add_closed}.
                                $\delta \rmul (\abs{x_0} + \abs{y_0} + \delta)\in\reals$ by \cref{real_mul_closed}.
                                $\delta \rmul (\abs{x_0} + \abs{y_0} + 1)\in\reals$ by \cref{real_mul_closed}.
                                Thus $\abs{x - x_0} \rmul \abs{y} + \abs{x_0} \rmul \abs{y - y_0} < \delta \rmul (\abs{x_0} + \abs{y_0} + 1)$ by \cref{real_lt_trans}.
                                Thus $\abs{x - x_0} \rmul \abs{y} + \abs{x_0} \rmul \abs{y - y_0} < \delta \rmul m$ by \cref{real_lt_subst_right}.
                            \end{byCase}
                        \end{subproof}
                        $0 + m < \epsilon + m$ by \cref{real_lt_add}.
                        $0 + m = m$ by \cref{real_add_ident}.
                        Thus $m < \epsilon + m$ by \cref{real_lt_subst_left}.
                        $\epsilon + m = m + \epsilon$ by \cref{real_add_comm}.
                        Thus $m < m + \epsilon$ by \cref{real_lt_subst_right}.
                        $m \rmul \inv{(m + \epsilon)} < (m + \epsilon) \rmul \inv{(m + \epsilon)}$ by \cref{real_lt_mul_pos}.
                        $(m + \epsilon) \rmul \inv{(m + \epsilon)} = 1$ by \cref{real_mul_inverse}.
                        Thus $m \rmul \inv{(m + \epsilon)} < 1$ by \cref{real_lt_subst_right}.
                        $m \rmul \inv{(m + \epsilon)}\in\reals$ by \cref{real_mul_closed}.
                        $1\in\reals$ by \cref{real_mul_ident}.
                        $(m \rmul \inv{(m + \epsilon)}) \rmul \epsilon < 1 \rmul \epsilon$ by \cref{real_lt_mul_pos}.
                        $(m \rmul \inv{(m + \epsilon)}) \rmul \epsilon = \epsilon \rmul (m \rmul \inv{(m + \epsilon)})$ by \cref{real_mul_comm}.
                        Thus $\epsilon \rmul (m \rmul \inv{(m + \epsilon)}) < 1 \rmul \epsilon$ by \cref{real_lt_subst_left}.
                        $1 \rmul \epsilon = \epsilon \rmul 1$ by \cref{real_mul_comm}.
                        Thus $\epsilon \rmul (m \rmul \inv{(m + \epsilon)}) < \epsilon \rmul 1$ by \cref{real_lt_subst_right}.
                        $\epsilon \rmul 1 = 1 \rmul \epsilon$ by \cref{real_mul_comm}.
                        $1 \rmul \epsilon = \epsilon$ by \cref{real_mul_ident}.
                        Thus $\epsilon \rmul 1 = \epsilon$.
                        $\delta = \epsilon \rmul \inv{(m + \epsilon)}$ by \cref{real_div}.
                        $(\epsilon \rmul \inv{(m + \epsilon)}) \rmul m = \epsilon \rmul (\inv{(m + \epsilon)} \rmul m)$ by \cref{real_mul_assoc}.
                        $\inv{(m + \epsilon)} \rmul m = m \rmul \inv{(m + \epsilon)}$ by \cref{real_mul_comm}.
                        Thus $\delta \rmul m = \epsilon \rmul (m \rmul \inv{(m + \epsilon)})$ by \cref{real_div,real_mul_assoc,real_mul_comm}.
                        Thus $\delta \rmul m < \epsilon$ by \cref{real_lt_subst_left}.
                        $x - x_0\in\reals$ by \cref{real_sub,real_add_closed}.
                        $y - y_0\in\reals$ by \cref{real_sub,real_add_closed}.
                        $\abs{x - x_0}\in\reals$ by \cref{real_abs_in_reals}.
                        $\abs{y}\in\reals$ by \cref{real_abs_in_reals}.
                        $\abs{y - y_0}\in\reals$ by \cref{real_abs_in_reals}.
                        $\abs{x - x_0} \rmul \abs{y}\in\reals$ by \cref{real_mul_closed}.
                        $\abs{x_0} \rmul \abs{y - y_0}\in\reals$ by \cref{real_mul_closed}.
                        $\abs{x - x_0} \rmul \abs{y} + \abs{x_0} \rmul \abs{y - y_0}\in\reals$ by \cref{real_add_closed}.
                        $\delta \rmul m\in\reals$ by \cref{real_mul_closed}.
                        Thus $\abs{x - x_0} \rmul \abs{y} + \abs{x_0} \rmul \abs{y - y_0} < \epsilon$ by \cref{real_lt_trans}.
                        Show that $\abs{(x - x_0) \rmul y + x_0 \rmul (y - y_0)} < \epsilon$.
                        \begin{subproof}
                            $(x - x_0) \rmul y\in\reals$ by \cref{real_mul_closed}.
                            $x_0 \rmul (y - y_0)\in\reals$ by \cref{real_mul_closed}.
                            $(x - x_0) \rmul y + x_0 \rmul (y - y_0)\in\reals$ by \cref{real_add_closed}.
                            $\abs{(x - x_0) \rmul y + x_0 \rmul (y - y_0)}\in\reals$ by \cref{real_abs_in_reals}.
                            Thus $\abs{(x - x_0) \rmul y + x_0 \rmul (y - y_0)} < \abs{x - x_0} \rmul \abs{y} + \abs{x_0} \rmul \abs{y - y_0}$ or
                                $\abs{(x - x_0) \rmul y + x_0 \rmul (y - y_0)} = \abs{x - x_0} \rmul \abs{y} + \abs{x_0} \rmul \abs{y - y_0}$ by \cref{real_leq_split}.
                            \begin{byCase}
                            \caseOf{$\abs{(x - x_0) \rmul y + x_0 \rmul (y - y_0)} < \abs{x - x_0} \rmul \abs{y} + \abs{x_0} \rmul \abs{y - y_0}$.}
                                Thus $\abs{(x - x_0) \rmul y + x_0 \rmul (y - y_0)} < \epsilon$ by \cref{real_lt_trans}.
                            \caseOf{$\abs{(x - x_0) \rmul y + x_0 \rmul (y - y_0)} = \abs{x - x_0} \rmul \abs{y} + \abs{x_0} \rmul \abs{y - y_0}$.}
                                Thus $\abs{(x - x_0) \rmul y + x_0 \rmul (y - y_0)} < \epsilon$ by \cref{real_lt_subst_left}.
                            \end{byCase}
                        \end{subproof}
                        Thus $\abs{(x \rmul y) - (x_0 \rmul y_0)} < \epsilon$ by \cref{real_lt_subst_left}.
                        Thus $\abs{f(q) - f(p)} < \epsilon$ by \cref{real_lt_subst_left}.
                    \end{subproof}
                    $f$ is a function to $\reals$ and $\dom{f} = \reals\times\reals$ by \cref{funs_elim}.
                    Thus $f(p)\in\reals$.
                    Thus $f(q)\in\reals$.
                    $d((f(p),f(q))) = \abs{f(p) - f(q)}$ by \cref{standard_metric_value}.
                    $\abs{f(p) - f(q)} = \abs{f(q) - f(p)}$ by \cref{real_abs_sub_swap}.
                    Thus $\abs{f(p) - f(q)} < \epsilon$ by \cref{real_lt_subst_left}.
                    Thus $d((f(p),f(q))) < \epsilon$ by \cref{real_lt_subst_left}.
                    Thus $f(q)\in\metricBall{d}{\reals}{f(p)}{\epsilon}$ by \cref{metric_ball}.
                \end{subproof}
                Show that for all $y$ such that $y\in\img{f}{U}$ we have
                    $y\in\metricBall{d}{\reals}{f(p)}{\epsilon}$.
                \begin{subproof}
                    Fix $y$.
                    Assume $y\in\img{f}{U}$.
                    Take $q\in U$ such that $q\mathrel{f} y$ by \cref{img_iff}.
                    $f$ is a function.
                    Thus $y = f(q)$ by \cref{function_apply_intro}.
                    Thus $y\in\metricBall{d}{\reals}{f(p)}{\epsilon}$.
                \end{subproof}
                Thus $\img{f}{U}\subseteq \metricBall{d}{\reals}{f(p)}{\epsilon}$ by \cref{subseteq}.
                $\metricBall{d}{\reals}{f(p)}{\epsilon}\subseteq V$ by \cref{subseteq}.
                Thus $\img{f}{U}\subseteq V$ by \cref{subseteq_transitive}.
            \end{subproof}
            Thus $U$ is a neighborhood of $p$ in $\reals\times\reals$ with respect to $T$ by \cref{neighborhood}.
        \end{subproof}
        $\standardBasisR$ is a basis on $\reals$ by \cref{standard_basis_is_basis}.
        Thus $\standardTopologyR$ is a topology on $\reals$ by \cref{basis_topology_is_topology,standard_topology_reals}.
        $\productBasis{\standardTopologyR}{\standardTopologyR}{\reals}{\reals}$ is a basis on $\reals\times\reals$ by \cref{product_basis_is_basis}.
        Thus $T$ is a topology on $\reals\times\reals$ by \cref{basis_topology_is_topology,product_topology}.
        Thus $f$ is continuous at $p$ from $\reals\times\reals$ to $\reals$ with respect to $T$ and $\standardTopologyR$ by \cref{continuous_at}.
    \end{subproof}
    $\standardBasisR$ is a basis on $\reals$ by \cref{standard_basis_is_basis}.
    Thus $\standardTopologyR$ is a topology on $\reals$ by \cref{basis_topology_is_topology,standard_topology_reals}.
    $\productBasis{\standardTopologyR}{\standardTopologyR}{\reals}{\reals}$ is a basis on $\reals\times\reals$ by \cref{product_basis_is_basis}.
    Thus $T$ is a topology on $\reals\times\reals$ by \cref{basis_topology_is_topology,product_topology}.
    Thus $f$ is continuous from $\reals\times\reals$ to $\reals$ with respect to $T$ and $\standardTopologyR$ by \cref{continuous_at_implies_continuous}.
\end{proof}

% from §21 helper definition 5: reciprocal operation on reals
\begin{definition}\label{real_inverse_operation}
    $\realInverse = \{w \mid \exists y\in\reals\setminus\{0\}. \exists z\in\reals.
        w = (y,z) \land z = \inv{y}\}$.
\end{definition}

% from §21 helper lemma 7: reciprocal operation is continuous
\begin{lemma}\label{real_inverse_continuous}
    Let $T_1 = \standardTopologyR$.
    Let $T_2 = \subspaceTopology{T_1}{\reals\setminus\{0\}}$.
    Then $\realInverse$ is continuous from $\reals\setminus\{0\}$ to $\reals$
        with respect to $T_2$ and $T_1$.
\end{lemma}
\begin{proof}
    Let $f = \realInverse$.
    Show that $f$ is a function from $\reals\setminus\{0\}$ to $\reals$.
    \begin{subproof}
        Show that $f$ is a function.
        \begin{subproof}
            Show that for all $p,z_1,z_2$ such that $(p,z_1)\in f$ and $(p,z_2)\in f$ we have $z_1 = z_2$.
            \begin{subproof}
                Fix $p$.
                Fix $z_1$.
                Fix $z_2$.
                Assume $(p,z_1)\in f$ and $(p,z_2)\in f$.
                Take $p_1,w_1$ such that
                    $p_1\in\reals\setminus\{0\}$ and $w_1\in\reals$ and
                    $(p,z_1) = (p_1,w_1)$ and $w_1 = \inv{p_1}$
                    by \cref{real_inverse_operation}.
                Take $p_2,w_2$ such that
                    $p_2\in\reals\setminus\{0\}$ and $w_2\in\reals$ and
                    $(p,z_2) = (p_2,w_2)$ and $w_2 = \inv{p_2}$
                    by \cref{real_inverse_operation}.
                $p = p_1$ and $z_1 = w_1$ by \cref{pair_eq_iff}.
                $p = p_2$ and $z_2 = w_2$ by \cref{pair_eq_iff}.
                Thus $z_1 = \inv{p}$.
                Thus $z_2 = \inv{p}$.
                Thus $z_1 = z_2$.
            \end{subproof}
            Thus $f$ is right-unique by \cref{rightunique}.
            Show that for all $w$ such that $w\in f$ there exist $x,y$ such that $w = (x,y)$.
            \begin{subproof}
                Fix $w$.
                Assume $w\in f$.
                Take $p,z$ such that $p\in\reals\setminus\{0\}$ and $z\in\reals$ and $w = (p,z)$ and
                    $z = \inv{p}$ by \cref{real_inverse_operation}.
                Thus there exist $x,y$ such that $w = (x,y)$.
            \end{subproof}
            Thus $f$ is a relation by \cref{relation}.
            Thus $f$ is a function.
        \end{subproof}
        Show that for all $p\in\dom{f}$ we have $f(p)\in\reals$.
        \begin{subproof}
            Fix $p\in\dom{f}$.
            Take $z$ such that $(p,z)\in f$ by \cref{dom_iff}.
            Take $p_0,w$ such that
                $p_0\in\reals\setminus\{0\}$ and $w\in\reals$ and
                $(p,z) = (p_0,w)$ and $w = \inv{p_0}$
                by \cref{real_inverse_operation}.
            $z = w$ by \cref{pair_eq_iff}.
            Thus $z\in\reals$.
            $f$ is a function.
            Thus $f(p) = z$ by \cref{function_apply_intro}.
            Thus $f(p)\in\reals$.
        \end{subproof}
        Thus $f$ is a function to $\reals$.
        Show that $\dom{f} = \reals\setminus\{0\}$.
        \begin{subproof}
            Show that for all $p\in\dom{f}$ we have $p\in\reals\setminus\{0\}$.
            \begin{subproof}
                Fix $p\in\dom{f}$.
                Take $z$ such that $(p,z)\in f$ by \cref{dom_iff}.
                Take $p_0,w$ such that
                    $p_0\in\reals\setminus\{0\}$ and $w\in\reals$ and
                    $(p,z) = (p_0,w)$ and $w = \inv{p_0}$
                    by \cref{real_inverse_operation}.
                Thus $p = p_0$ by \cref{pair_eq_iff}.
                Thus $p\in\reals\setminus\{0\}$.
            \end{subproof}
            Thus $\dom{f}\subseteq\reals\setminus\{0\}$ by \cref{subseteq}.
            Show that for all $p\in\reals\setminus\{0\}$ we have $p\in\dom{f}$.
            \begin{subproof}
                Fix $p\in\reals\setminus\{0\}$.
                $p\in\reals$ by \cref{setminus_elim_left}.
                $p\notin\{0\}$ by \cref{setminus_elim_right}.
                Thus $p\neq 0$ by \cref{singleton_iff}.
                Thus $\inv{p}\in\reals$ by \cref{real_mul_inverse}.
                Let $z = \inv{p}$.
                Let $w = (p,z)$.
                $w\in f$ by \cref{real_inverse_operation}.
                Thus $p\in\dom{f}$ by \cref{dom_intro}.
            \end{subproof}
            Thus $\reals\setminus\{0\}\subseteq\dom{f}$ by \cref{subseteq}.
            Thus $\dom{f} = \reals\setminus\{0\}$ by \cref{subseteq_antisymmetric}.
        \end{subproof}
        Thus $f\in\funs{\reals\setminus\{0\}}{\reals}$ by \cref{funs_intro}.
    \end{subproof}
    Show that for all $p\in\reals\setminus\{0\}$ we have
        $f$ is continuous at $p$ from $\reals\setminus\{0\}$ to $\reals$ with respect to $T_2$ and $T_1$.
    \begin{subproof}
        Fix $p$.
        Assume $p\in\reals\setminus\{0\}$.
        $p\in\reals$ by \cref{setminus_elim_left}.
        $p\notin\{0\}$ by \cref{setminus_elim_right}.
        Thus $p\neq 0$ by \cref{singleton_iff}.
        $\inv{p}\in\reals$ by \cref{real_mul_inverse}.
        Let $z_0 = \inv{p}$.
        Let $w_0 = (p,z_0)$.
        $w_0\in f$ by \cref{real_inverse_operation}.
        $f$ is a function.
        Thus $f(p) = z_0$ by \cref{function_apply_intro}.
        Show that for all $V$ such that $V$ is a neighborhood of $f(p)$ in $\reals$ with respect to $T_1$
            there exists $U$ such that $U$ is a neighborhood of $p$ in $\reals\setminus\{0\}$ with respect to $T_2$
            and $\img{f}{U}\subseteq V$.
        \begin{subproof}
            Fix $V$.
            Assume $V$ is a neighborhood of $f(p)$ in $\reals$ with respect to $T_1$.
            $V$ is open in $\reals$ with respect to $T_1$ by \cref{neighborhood}.
            Let $d = \standardMetricR$.
            Let $T_d = \metricTopology{d}{\reals}$.
            $T_d = \standardTopologyR$ by \cref{standard_metric_topology}.
            $T_1 = \standardTopologyR$.
            $\standardTopologyR\subseteq T_d$ by \cref{subseteq}.
            Thus $V\in T_d$ by \cref{open_in,subseteq}.
            $T_d\subseteq\standardTopologyR$ by \cref{subseteq}.
            $f(p)\in V$ by \cref{neighborhood}.
            $\standardBasisR$ is a basis on $\reals$ by \cref{standard_basis_is_basis}.
            Thus $\standardTopologyR$ is a topology on $\reals$ by \cref{basis_topology_is_topology,standard_topology_reals}.
            Thus $V\subseteq\reals$ by \cref{topology_on,subseteq,pow_iff}.
            $d$ is a metric on $\reals$ by \cref{standard_metric_is_metric}.
            Take $\epsilon\in\reals$ such that $0 < \epsilon$ and
                $\metricBall{d}{\reals}{f(p)}{\epsilon} \subseteq V$ by \cref{metric_open_iff}.
            $\abs{p}\in\reals$ by \cref{real_abs_in_reals}.
            $\abs{p} \geq 0$ by \cref{real_abs_nonnegative}.
            $\abs{p}\neq 0$ by \cref{real_abs_eq_zero}.
            Thus $\abs{p} > 0$ by \cref{real_lt_trichotomy,real_abs_nonnegative}.
            Let $a = \abs{p} \rmul \inv{(1 + 1)}$.
            $1\in\reals$ by \cref{real_mul_ident}.
            $0\in\reals$ by \cref{real_add_ident}.
            $0 < 1$ by \cref{real_zero_lt_one}.
            $1 + 1\in\reals$ by \cref{real_add_closed}.
            $0 + 1 < 1 + 1$ by \cref{real_lt_add}.
            $0 + 1 = 1$ by \cref{real_add_ident}.
            Thus $1 < 1 + 1$ by \cref{real_lt_subst_left}.
            Thus $0 < 1 + 1$ by \cref{real_lt_trans}.
            Thus $(1 + 1)\neq 0$ by \cref{real_pos_nonzero}.
            Thus $\inv{(1 + 1)}\in\reals$ by \cref{real_mul_inverse}.
            $0 < \inv{(1 + 1)}$ by \cref{real_inv_pos}.
            $0 \rmul \inv{(1 + 1)} < \abs{p} \rmul \inv{(1 + 1)}$ by \cref{real_lt_mul_pos}.
            Thus $0 \rmul \inv{(1 + 1)} < a$ by \cref{real_lt_subst_right}.
            $0 \rmul \inv{(1 + 1)} = 0$ by \cref{real_mul_zero}.
            Thus $0 < a$ by \cref{real_lt_subst_left}.
            Let $b = (\epsilon \rmul (\abs{p} \rmul \abs{p})) \rmul \inv{(1 + 1)}$.
            $\abs{p} \rmul \abs{p}\in\reals$ by \cref{real_mul_closed}.
            $\epsilon \rmul (\abs{p} \rmul \abs{p})\in\reals$ by \cref{real_mul_closed}.
            $(1 + 1)\in\reals$ by \cref{real_add_closed}.
            $0 < (1 + 1)$ by \cref{real_lt_add}.
            Thus $\inv{(1 + 1)}\in\reals$ by \cref{real_mul_inverse}.
            $0 < \inv{(1 + 1)}$ by \cref{real_inv_pos}.
            Thus $\abs{p} \rmul \abs{p} > 0$ by \cref{real_mul_pos_pos_gt}.
            $0 \rmul (\abs{p} \rmul \abs{p}) < \epsilon \rmul (\abs{p} \rmul \abs{p})$ by \cref{real_lt_mul_pos}.
            $0 \rmul (\abs{p} \rmul \abs{p}) = 0$ by \cref{real_mul_zero}.
            Thus $0 < \epsilon \rmul (\abs{p} \rmul \abs{p})$ by \cref{real_lt_subst_left}.
            $0 \rmul \inv{(1 + 1)} < (\epsilon \rmul (\abs{p} \rmul \abs{p})) \rmul \inv{(1 + 1)}$ by \cref{real_lt_mul_pos}.
            Thus $0 \rmul \inv{(1 + 1)} < b$ by \cref{real_lt_subst_right}.
            $0 \rmul \inv{(1 + 1)} = 0$ by \cref{real_mul_zero}.
            Thus $0 < b$ by \cref{real_lt_subst_left}.
            $a\in\reals$ by \cref{real_mul_closed}.
            $b\in\reals$ by \cref{real_mul_closed}.
            Let $E = \{a,b\}$.
            Show that for all $x$ such that $x\in E$ we have $x\in\reals$.
            \begin{subproof}
                Fix $x$.
                Assume $x\in E$.
                $x = a$ or $x = b$ by \cref{upair_elim}.
                \begin{byCase}
                \caseOf{$x = a$.}
                    Thus $x\in\reals$.
                \caseOf{$x = b$.}
                    Thus $x\in\reals$.
                \end{byCase}
            \end{subproof}
            Thus $E\subseteq\reals$ by \cref{subseteq}.
            $E$ is finite by \cref{pair_is_finite}.
            Show that $E\neq\emptyset$.
            \begin{subproof}
                Suppose not.
                Then $E = \emptyset$.
                $a\in E$ by \cref{upair_intro_left}.
                Thus $a\in\emptyset$.
                Contradiction by \cref{emptyset}.
            \end{subproof}
            Take $\delta\in E$ such that $\delta$ is a minimum of $E$ by \cref{finite_real_minimum}.
            $\delta\in E$.
            Thus $\delta\in\reals$ by \cref{subseteq}.
            Show that $0 < \delta$.
            \begin{subproof}
                $\delta = a$ or $\delta = b$ by \cref{upair_elim}.
                \begin{byCase}
                \caseOf{$\delta = a$.}
                    Thus $0 < \delta$ by \cref{real_lt_subst_right}.
                \caseOf{$\delta = b$.}
                    Thus $0 < \delta$ by \cref{real_lt_subst_right}.
                \end{byCase}
            \end{subproof}
            $a\in E$ by \cref{upair_intro_left}.
            $b\in E$ by \cref{upair_intro_right}.
            Thus $\delta \leq a$ by \cref{real_minimum}.
            Thus $\delta \leq b$ by \cref{real_minimum}.
            Let $U = \metricBall{d}{\reals}{p}{\delta}$.
            Show that $U\subseteq\reals\setminus\{0\}$.
            \begin{subproof}
                Show that for all $y$ such that $y\in U$ we have $y\in\reals\setminus\{0\}$.
                \begin{subproof}
                    Fix $y$.
                    Assume $y\in U$.
                    Thus $y\in\reals$ by \cref{metric_ball}.
                    $d((p,y)) < \delta$ by \cref{metric_ball}.
                    $d((p,y)) = \abs{p - y}$ by \cref{standard_metric_value}.
                    Thus $\abs{p - y} < \delta$.
                    Show that $y\neq 0$.
                    \begin{subproof}
                        Show that $\abs{p - y} < a$.
                        \begin{subproof}
                            $\neg{y}\in\reals$ by \cref{real_add_inverse}.
                            $p - y\in\reals$ by \cref{real_sub,real_add_closed}.
                            Thus $\abs{p - y}\in\reals$ by \cref{real_abs_in_reals}.
                            $\delta < a$ or $\delta = a$ by \cref{real_leq_split}.
                            \begin{byCase}
                            \caseOf{$\delta < a$.}
                                Thus $\abs{p - y} < a$ by \cref{real_lt_trans}.
                            \caseOf{$\delta = a$.}
                                Thus $\abs{p - y} < a$ by \cref{real_lt_subst_right}.
                            \end{byCase}
                        \end{subproof}
                        Show that $a + a = \abs{p}$.
                        \begin{subproof}
                            $a + a = (\abs{p} + \abs{p}) \rmul \inv{(1 + 1)}$ by \cref{real_right_distrib}.
                            $(1 + 1) \rmul \abs{p} = \abs{p} + \abs{p}$ by \cref{real_right_distrib,real_mul_ident}.
                            Thus $a + a = ((1 + 1) \rmul \abs{p}) \rmul \inv{(1 + 1)}$ by \cref{real_mul_eq_subst}.
                            $((1 + 1) \rmul \abs{p}) \rmul \inv{(1 + 1)} = \abs{p} \rmul ((1 + 1) \rmul \inv{(1 + 1)})$ by \cref{real_mul_assoc,real_mul_comm}.
                            $(1 + 1) \rmul \inv{(1 + 1)} = 1$ by \cref{real_mul_inverse}.
                            Thus $a + a = \abs{p} \rmul 1$ by \cref{real_mul_eq_subst}.
                            $\abs{p} \rmul 1 = 1 \rmul \abs{p}$ by \cref{real_mul_comm}.
                            Thus $a + a = 1 \rmul \abs{p}$ by \cref{real_mul_eq_subst}.
                            Thus $a + a = \abs{p}$ by \cref{real_mul_ident}.
                        \end{subproof}
                        $\abs{p} < \abs{y} + a$ by \cref{real_abs_sub_lt_bound}.
                        Show that it is wrong that $\abs{y} \leq a$.
                        \begin{subproof}
                            Suppose not.
                            Then $\abs{y} \leq a$.
                            $\abs{y}\in\reals$ by \cref{real_abs_in_reals}.
                            Then $\abs{y} < a$ or $\abs{y} = a$ by \cref{real_leq_split}.
                            \begin{byCase}
                            \caseOf{$\abs{y} < a$.}
                                $\abs{y} + a < a + a$ by \cref{real_lt_add}.
                                Thus $\abs{y} + a < \abs{p}$ by \cref{real_lt_subst_right}.
                                $\abs{y} + a\in\reals$ by \cref{real_add_closed}.
                                Thus $\abs{p} < \abs{p}$ by \cref{real_lt_trans}.
                                Contradiction by \cref{real_lt_irreflexive}.
                            \caseOf{$\abs{y} = a$.}
                                $\abs{y} + a = a + a$ by \cref{real_add_eq_subst}.
                                Thus $\abs{y} + a = \abs{p}$.
                                Thus $\abs{p} < \abs{p}$ by \cref{real_lt_subst_right}.
                                Contradiction by \cref{real_lt_irreflexive}.
                            \end{byCase}
                        \end{subproof}
                        Thus $\abs{y} > a$.
                        Thus $\abs{y}\neq 0$ by \cref{real_lt_trichotomy,real_abs_nonnegative}.
                        Show that $y\neq 0$.
                        \begin{subproof}
                            Suppose not.
                            Then $y = 0$.
                            Thus $0 = y$.
                            Thus $0 < y$ or $0 = y$.
                            Thus $0 \leq y$.
                            Thus $\abs{y} = y$ by \cref{real_abs_nonneg}.
                            Thus $\abs{y} = 0$.
                            Contradiction.
                        \end{subproof}
                    \end{subproof}
                    Thus $y\notin\{0\}$ by \cref{singleton_iff}.
                    Thus $y\in\reals\setminus\{0\}$ by \cref{setminus_intro}.
                \end{subproof}
                Thus $U\subseteq\reals\setminus\{0\}$ by \cref{subseteq}.
            \end{subproof}
            Show that $U$ is a neighborhood of $p$ in $\reals\setminus\{0\}$ with respect to $T_2$.
            \begin{subproof}
                $p\in\reals$ and $\delta\in\reals$ and $0 < \delta$.
                Show that $U$ is open in $\reals$ with respect to $\standardTopologyR$.
                \begin{subproof}
                    Show that for all $y$ such that $y\in U$ we have $y\in\reals$.
                    \begin{subproof}
                        Fix $y$.
                        Assume $y\in U$.
                        Thus $y\in\reals$ by \cref{metric_ball}.
                    \end{subproof}
                    Thus $U\subseteq\reals$ by \cref{subseteq}.
                    $U\in\pow{\reals}$ by \cref{powerset_intro}.
                    $U\in\metricBasis{d}{\reals}$ by \cref{metric_basis}.
                    $\metricBasis{d}{\reals}$ is a basis on $\reals$ by \cref{metric_basis_is_basis}.
                    Thus $U\in T_d$ by \cref{basis_elem_open_in_generated,basis_for_topology,metric_topology}.
                    $T_d\subseteq\standardTopologyR$ by \cref{subseteq}.
                    Thus $U\in\standardTopologyR$ by \cref{subseteq}.
                    $\standardBasisR$ is a basis on $\reals$ by \cref{standard_basis_is_basis}.
                    Thus $\standardTopologyR$ is a topology on $\reals$ by \cref{basis_topology_is_topology,standard_topology_reals}.
                    Thus $U$ is open in $\reals$ with respect to $\standardTopologyR$ by \cref{open_in}.
                \end{subproof}
                $U\subseteq\reals\setminus\{0\}$.
                $U\in T_1$ by \cref{open_in}.
                $U\subseteq U$ by \cref{subseteq}.
                Thus $U\subseteq(\reals\setminus\{0\})\inter U$ by \cref{subseteq_inter_iff}.
                $(\reals\setminus\{0\})\inter U\subseteq U$ by \cref{inter_lower_right}.
                Thus $U = (\reals\setminus\{0\})\inter U$ by \cref{subseteq_antisymmetric}.
                Thus $U\in T_2$ by \cref{subspace_topology}.
                $\standardBasisR$ is a basis on $\reals$ by \cref{standard_basis_is_basis}.
                Thus $T_1$ is a topology on $\reals$ by \cref{basis_topology_is_topology,standard_topology_reals}.
                Show that for all $y$ such that $y\in\reals\setminus\{0\}$ we have $y\in\reals$.
                \begin{subproof}
                    Fix $y$.
                    Assume $y\in\reals\setminus\{0\}$.
                    Thus $y\in\reals$ by \cref{setminus_elim_left}.
                \end{subproof}
                Thus $\reals\setminus\{0\}\subseteq\reals$ by \cref{subseteq}.
                Thus $T_2$ is a topology on $\reals\setminus\{0\}$ by \cref{subspace_topology_is_topology}.
                Thus $U$ is open in $\reals\setminus\{0\}$ with respect to $T_2$ by \cref{open_in}.
                $d((p,p)) = 0$ by \cref{metric_zero_characterization,standard_metric_is_metric,metric_on}.
                Thus $d((p,p)) < \delta$ by \cref{real_lt_subst_left}.
                Thus $p\in U$ by \cref{metric_ball}.
                Thus $U$ is a neighborhood of $p$ in $\reals\setminus\{0\}$ with respect to $T_2$ by \cref{neighborhood}.
            \end{subproof}
            Show that $\img{f}{U}\subseteq V$.
            \begin{subproof}
                Show that for all $q$ such that $q\in U$ we have
                    $f(q)\in\metricBall{d}{\reals}{f(p)}{\epsilon}$.
                \begin{subproof}
                    Fix $q$.
                    Assume $q\in U$.
                    $q\in\reals$ by \cref{metric_ball}.
                    $d((p,q)) < \delta$ by \cref{metric_ball}.
                    $d((p,q)) = \abs{p - q}$ by \cref{standard_metric_value}.
                    Thus $\abs{p - q} < \delta$.
                    Show that $f(q) = \inv{q}$.
                    \begin{subproof}
                        $q\notin\{0\}$ by \cref{subseteq,setminus_elim_right,metric_ball}.
                        Thus $q\neq 0$ by \cref{singleton_iff}.
                        Thus $q\in\reals\setminus\{0\}$ by \cref{setminus_intro}.
                        $\inv{q}\in\reals$ by \cref{real_mul_inverse}.
                        Let $w = (q, \inv{q})$.
                        $w\in f$ by \cref{real_inverse_operation}.
                        $f$ is a function.
                        Thus $f(q) = \inv{q}$ by \cref{function_apply_intro}.
                    \end{subproof}
                    Show that $\abs{f(q) - f(p)} < \epsilon$.
                    \begin{subproof}
                        $q\neq 0$ and $p\neq 0$.
                        Show that $q \rmul p \neq 0$.
                        \begin{subproof}
                            Suppose not.
                            Then $q \rmul p = 0$.
                            $\inv{q}\in\reals$ by \cref{real_mul_inverse}.
                            $\inv{q} \rmul (q \rmul p) = \inv{q} \rmul 0$ by \cref{real_mul_zero}.
                            $\inv{q} \rmul (q \rmul p) = (\inv{q} \rmul q) \rmul p$ by \cref{real_mul_assoc}.
                            $\inv{q} \rmul q = 1$ by \cref{real_mul_inverse,real_mul_comm}.
                            Thus $p = 0$ by \cref{real_mul_ident,real_mul_zero,real_mul_comm}.
                            Contradiction.
                        \end{subproof}
                        Thus $q \rmul p \neq 0$.
                        $q \rmul p\in\reals$ by \cref{real_mul_closed}.
                        $\inv{(q \rmul p)}\in\reals$ by \cref{real_mul_inverse}.
                        Show that $\inv{q} = p \rmul \inv{(q \rmul p)}$.
                        \begin{subproof}
                            $q \rmul (p \rmul \inv{(q \rmul p)}) = (q \rmul p) \rmul \inv{(q \rmul p)}$ by \cref{real_mul_assoc}.
                            $(q \rmul p) \rmul \inv{(q \rmul p)} = 1$ by \cref{real_mul_inverse}.
                            $p \rmul \inv{(q \rmul p)}\in\reals$ by \cref{real_mul_closed}.
                            Thus $p \rmul \inv{(q \rmul p)} = \inv{q}$ by \cref{real_mul_inverse_unique}.
                        \end{subproof}
                        Show that $\inv{p} = q \rmul \inv{(q \rmul p)}$.
                        \begin{subproof}
                            $p \rmul (q \rmul \inv{(q \rmul p)}) = (p \rmul q) \rmul \inv{(q \rmul p)}$ by \cref{real_mul_assoc}.
                            $p \rmul q = q \rmul p$ by \cref{real_mul_comm}.
                            Thus $p \rmul (q \rmul \inv{(q \rmul p)}) = (q \rmul p) \rmul \inv{(q \rmul p)}$ by \cref{real_mul_assoc}.
                            $(q \rmul p) \rmul \inv{(q \rmul p)} = 1$ by \cref{real_mul_inverse}.
                            $q \rmul \inv{(q \rmul p)}\in\reals$ by \cref{real_mul_closed}.
                            Thus $q \rmul \inv{(q \rmul p)} = \inv{p}$ by \cref{real_mul_inverse_unique}.
                        \end{subproof}
                        $\inv{q} - \inv{p} = (p \rmul \inv{(q \rmul p)}) - (q \rmul \inv{(q \rmul p)})$ by \cref{real_lt_subst_left}.
                        Show that $(p \rmul \inv{(q \rmul p)}) - (q \rmul \inv{(q \rmul p)}) = (p - q) \rmul \inv{(q \rmul p)}$.
                        \begin{subproof}
                            $(p \rmul \inv{(q \rmul p)}) - (q \rmul \inv{(q \rmul p)}) =
                                (p \rmul \inv{(q \rmul p)}) + \neg{(q \rmul \inv{(q \rmul p)})}$ by \cref{real_sub}.
                            Show that $\neg{(q \rmul \inv{(q \rmul p)})} = (\neg{q}) \rmul \inv{(q \rmul p)}$.
                            \begin{subproof}
                                $\neg{(q \rmul \inv{(q \rmul p)})} = \neg{1} \rmul (q \rmul \inv{(q \rmul p)})$ by \cref{real_neg_as_mul}.
                                $q \rmul \inv{(q \rmul p)} = \inv{(q \rmul p)} \rmul q$ by \cref{real_mul_comm}.
                                Thus $\neg{1} \rmul (q \rmul \inv{(q \rmul p)}) = \neg{1} \rmul (\inv{(q \rmul p)} \rmul q)$ by \cref{real_mul_eq_subst}.
                                Thus $\neg{(q \rmul \inv{(q \rmul p)})} = \neg{1} \rmul (\inv{(q \rmul p)} \rmul q)$ by \cref{real_mul_eq_subst}.
                                $\neg{1} \rmul (\inv{(q \rmul p)} \rmul q) = \neg{(\inv{(q \rmul p)} \rmul q)}$ by \cref{real_neg_as_mul}.
                                $\inv{(q \rmul p)} \rmul \neg{q} = \neg{(\inv{(q \rmul p)} \rmul q)}$ by \cref{real_mul_neg_right}.
                                $\neg{q}\in\reals$ by \cref{real_add_inverse}.
                                $\inv{(q \rmul p)} \rmul \neg{q} = (\neg{q}) \rmul \inv{(q \rmul p)}$ by \cref{real_mul_comm}.
                                Thus $\neg{(q \rmul \inv{(q \rmul p)})} = (\neg{q}) \rmul \inv{(q \rmul p)}$.
                            \end{subproof}
                            Thus $(p \rmul \inv{(q \rmul p)}) + \neg{(q \rmul \inv{(q \rmul p)})} =
                                (p \rmul \inv{(q \rmul p)}) + ((\neg{q}) \rmul \inv{(q \rmul p)})$ by \cref{real_add_eq_subst}.
                            $\neg{q}\in\reals$ by \cref{real_add_inverse}.
                            $(p \rmul \inv{(q \rmul p)}) + ((\neg{q}) \rmul \inv{(q \rmul p)}) =
                                (p + \neg{q}) \rmul \inv{(q \rmul p)}$ by \cref{real_right_distrib}.
                            $p + \neg{q} = p - q$ by \cref{real_sub}.
                            Thus $(p \rmul \inv{(q \rmul p)}) + ((\neg{q}) \rmul \inv{(q \rmul p)}) =
                                (p - q) \rmul \inv{(q \rmul p)}$ by \cref{real_mul_eq_subst}.
                            Thus $(p \rmul \inv{(q \rmul p)}) - (q \rmul \inv{(q \rmul p)}) = (p - q) \rmul \inv{(q \rmul p)}$
                                by \cref{real_add_eq_subst}.
                        \end{subproof}
                        Thus $\abs{\inv{q} - \inv{p}} = \abs{(p - q) \rmul \inv{(q \rmul p)}}$ by \cref{real_lt_subst_left}.
                        $\neg{q}\in\reals$ by \cref{real_add_inverse}.
                        $p - q\in\reals$ by \cref{real_sub,real_add_closed}.
                        $\abs{(p - q) \rmul \inv{(q \rmul p)}} = \abs{p - q} \rmul \abs{\inv{(q \rmul p)}}$ by \cref{real_abs_mul}.
                        $\abs{p - q} = \abs{q - p}$ by \cref{real_abs_sub_swap}.
                        Show that $\abs{p - q} < \epsilon \rmul \abs{q \rmul p}$.
                        \begin{subproof}
                            $\neg{q}\in\reals$ by \cref{real_add_inverse}.
                            $p - q\in\reals$ by \cref{real_sub,real_add_closed}.
                            Thus $\abs{p - q}\in\reals$ by \cref{real_abs_in_reals}.
                            Show that $a + a = \abs{p}$.
                            \begin{subproof}
                                $a + a = (\abs{p} + \abs{p}) \rmul \inv{(1 + 1)}$ by \cref{real_right_distrib}.
                                $(1 + 1) \rmul \abs{p} = \abs{p} + \abs{p}$ by \cref{real_right_distrib,real_mul_ident}.
                                Thus $a + a = ((1 + 1) \rmul \abs{p}) \rmul \inv{(1 + 1)}$ by \cref{real_mul_eq_subst}.
                                $((1 + 1) \rmul \abs{p}) \rmul \inv{(1 + 1)} = \abs{p} \rmul ((1 + 1) \rmul \inv{(1 + 1)})$ by \cref{real_mul_assoc,real_mul_comm}.
                                $(1 + 1) \rmul \inv{(1 + 1)} = 1$ by \cref{real_mul_inverse}.
                                Thus $a + a = \abs{p} \rmul 1$ by \cref{real_mul_eq_subst}.
                                $\abs{p} \rmul 1 = 1 \rmul \abs{p}$ by \cref{real_mul_comm}.
                                Thus $a + a = 1 \rmul \abs{p}$ by \cref{real_mul_eq_subst}.
                                Thus $a + a = \abs{p}$ by \cref{real_mul_ident}.
                            \end{subproof}
                            Show that $\abs{q} > a$.
                            \begin{subproof}
                                Show that it is wrong that $\abs{q} \leq a$.
                                \begin{subproof}
                                    Suppose not.
                                    Then $\abs{q} \leq a$.
                                    $\abs{q}\in\reals$ by \cref{real_abs_in_reals}.
                                    Show that $\abs{p - q} < a$.
                                    \begin{subproof}
                                        $\delta < a$ or $\delta = a$ by \cref{real_leq_split}.
                                        \begin{byCase}
                                        \caseOf{$\delta < a$.}
                                            Thus $\abs{p - q} < a$ by \cref{real_lt_trans}.
                                        \caseOf{$\delta = a$.}
                                            Thus $\abs{p - q} < a$ by \cref{real_lt_subst_right}.
                                        \end{byCase}
                                    \end{subproof}
                                    Thus $\abs{p - q} < a$.
                                    $\abs{p} < \abs{q} + a$ by \cref{real_abs_sub_lt_bound}.
                                    Then $\abs{q} < a$ or $\abs{q} = a$ by \cref{real_leq_split}.
                                    \begin{byCase}
                                    \caseOf{$\abs{q} < a$.}
                                        $\abs{q} + a < a + a$ by \cref{real_lt_add}.
                                        Thus $\abs{q} + a < \abs{p}$ by \cref{real_lt_subst_right}.
                                        $\abs{q} + a\in\reals$ by \cref{real_add_closed}.
                                        Thus $\abs{p} < \abs{p}$ by \cref{real_lt_trans}.
                                        Contradiction by \cref{real_lt_irreflexive}.
                                    \caseOf{$\abs{q} = a$.}
                                        $\abs{q} + a = a + a$ by \cref{real_add_eq_subst}.
                                        Thus $\abs{q} + a = \abs{p}$.
                                        Thus $\abs{p} < \abs{p}$ by \cref{real_lt_subst_right}.
                                        Contradiction by \cref{real_lt_irreflexive}.
                                    \end{byCase}
                                \end{subproof}
                                Thus $\abs{q} > a$.
                            \end{subproof}
                            $\abs{p}\in\reals$ by \cref{real_abs_in_reals}.
                            Show that $\abs{p}\neq 0$.
                            \begin{subproof}
                                Suppose not.
                                Then $\abs{p} = 0$.
                                Thus $p = 0$ by \cref{real_abs_eq_zero}.
                                Contradiction.
                            \end{subproof}
                            $\abs{p} \geq 0$ by \cref{real_abs_nonnegative}.
                            Thus $0 < \abs{p}$ by \cref{real_lt_trichotomy}.
                            $a\in\reals$ by \cref{real_mul_closed}.
                            $\abs{q}\in\reals$ by \cref{real_abs_in_reals}.
                            $a \rmul \abs{p} < \abs{q} \rmul \abs{p}$ by \cref{real_lt_mul_pos}.
                            $\abs{q \rmul p} = \abs{q} \rmul \abs{p}$ by \cref{real_abs_mul}.
                            Thus $\abs{q \rmul p} > a \rmul \abs{p}$ by \cref{real_lt_subst_left}.
                            Show that $\abs{p - q} < b$.
                            \begin{subproof}
                                $\delta < b$ or $\delta = b$ by \cref{real_leq_split}.
                                \begin{byCase}
                                \caseOf{$\delta < b$.}
                                    Thus $\abs{p - q} < b$ by \cref{real_lt_trans}.
                                \caseOf{$\delta = b$.}
                                    Thus $\abs{p - q} < b$ by \cref{real_lt_subst_right}.
                                \end{byCase}
                            \end{subproof}
                            Thus $\abs{p - q} < b$.
                            Show that $b = \epsilon \rmul (a \rmul \abs{p})$.
                            \begin{subproof}
                                $a \rmul \abs{p} = (\abs{p} \rmul \inv{(1 + 1)}) \rmul \abs{p}$ by \cref{real_mul_eq_subst}.
                                $(\abs{p} \rmul \inv{(1 + 1)}) \rmul \abs{p} = (\abs{p} \rmul \abs{p}) \rmul \inv{(1 + 1)}$ by \cref{real_mul_assoc,real_mul_comm}.
                                Thus $\epsilon \rmul (a \rmul \abs{p}) = (\epsilon \rmul (\abs{p} \rmul \abs{p})) \rmul \inv{(1 + 1)}$ by \cref{real_mul_assoc,real_mul_comm}.
                                Thus $b = \epsilon \rmul (a \rmul \abs{p})$.
                            \end{subproof}
                            $\abs{q \rmul p} = \abs{q} \rmul \abs{p}$ by \cref{real_abs_mul}.
                            Thus $a \rmul \abs{p} < \abs{q \rmul p}$ by \cref{real_lt_subst_right}.
                            $a \rmul \abs{p}\in\reals$ by \cref{real_mul_closed}.
                            $\abs{q \rmul p}\in\reals$ by \cref{real_abs_in_reals}.
                            $(a \rmul \abs{p}) \rmul \epsilon < \abs{q \rmul p} \rmul \epsilon$ by \cref{real_lt_mul_pos}.
                            $\epsilon \rmul (a \rmul \abs{p}) = (a \rmul \abs{p}) \rmul \epsilon$ by \cref{real_mul_comm}.
                            $\epsilon \rmul \abs{q \rmul p} = \abs{q \rmul p} \rmul \epsilon$ by \cref{real_mul_comm}.
                            Thus $\epsilon \rmul (a \rmul \abs{p}) < \abs{q \rmul p} \rmul \epsilon$ by \cref{real_lt_subst_left}.
                            Thus $\epsilon \rmul (a \rmul \abs{p}) < \epsilon \rmul \abs{q \rmul p}$ by \cref{real_lt_subst_right}.
                            Thus $b < \epsilon \rmul \abs{q \rmul p}$ by \cref{real_lt_subst_left}.
                            $\epsilon \rmul \abs{q \rmul p}\in\reals$ by \cref{real_mul_closed}.
                            Thus $\abs{p - q} < \epsilon \rmul \abs{q \rmul p}$ by \cref{real_lt_trans}.
                        \end{subproof}
                        Show that $0 < \abs{\inv{(q \rmul p)}}$.
                        \begin{subproof}
                            $\inv{(q \rmul p)} \rmul (q \rmul p) = 1$ by \cref{real_mul_inverse,real_mul_comm}.
                            Thus $\inv{(q \rmul p)} \neq 0$ by \cref{real_mul_unit_left_nonzero}.
                            Thus $\abs{\inv{(q \rmul p)}}\neq 0$ by \cref{real_abs_eq_zero}.
                            $\abs{\inv{(q \rmul p)}} \geq 0$ by \cref{real_abs_nonnegative}.
                            Thus $0 < \abs{\inv{(q \rmul p)}}$ by \cref{real_lt_trichotomy}.
                        \end{subproof}
                        $\abs{p - q}\in\reals$ by \cref{real_abs_in_reals}.
                        $\abs{\inv{(q \rmul p)}}\in\reals$ by \cref{real_abs_in_reals}.
                        $\epsilon\in\reals$.
                        $\abs{q \rmul p}\in\reals$ by \cref{real_abs_in_reals}.
                        $\epsilon \rmul \abs{q \rmul p}\in\reals$ by \cref{real_mul_closed}.
                        Thus $\abs{p - q} \rmul \abs{\inv{(q \rmul p)}} < \epsilon \rmul \abs{q \rmul p} \rmul \abs{\inv{(q \rmul p)}}$
                            by \cref{real_lt_mul_pos}.
                        $\abs{q \rmul p} \rmul \abs{\inv{(q \rmul p)}} = \abs{(q \rmul p) \rmul \inv{(q \rmul p)}}$ by \cref{real_abs_mul,real_mul_assoc}.
                        $(q \rmul p) \rmul \inv{(q \rmul p)} = 1$ by \cref{real_mul_inverse}.
                        $\abs{1} = 1$ by \cref{real_abs_nonneg,real_zero_lt_one,real_lt_imp_leq}.
                        Thus $\abs{q \rmul p} \rmul \abs{\inv{(q \rmul p)}} = 1$ by \cref{real_lt_subst_right}.
                        $\epsilon \rmul \abs{q \rmul p} \rmul \abs{\inv{(q \rmul p)}} =
                            \epsilon \rmul (\abs{q \rmul p} \rmul \abs{\inv{(q \rmul p)}})$ by \cref{real_mul_assoc}.
                        Thus $\epsilon \rmul \abs{q \rmul p} \rmul \abs{\inv{(q \rmul p)}} = \epsilon \rmul 1$ by \cref{real_mul_eq_subst}.
                        $\epsilon \rmul 1 = 1 \rmul \epsilon$ by \cref{real_mul_comm}.
                        Thus $\epsilon \rmul 1 = \epsilon$ by \cref{real_mul_ident}.
                        Thus $\epsilon \rmul \abs{q \rmul p} \rmul \abs{\inv{(q \rmul p)}} = \epsilon$ by \cref{real_mul_eq_subst}.
                        Thus $\abs{p - q} \rmul \abs{\inv{(q \rmul p)}} < \epsilon$ by \cref{real_lt_subst_right}.
                        Thus $\abs{\inv{q} - \inv{p}} < \epsilon$ by \cref{real_lt_subst_left}.
                    \end{subproof}
                    Thus $\abs{f(q) - f(p)} < \epsilon$ by \cref{real_lt_subst_left,real_div}.
                $f(p)\in\reals$ by \cref{funs_type_apply}.
                $q\in\reals\setminus\{0\}$ by \cref{subseteq}.
                $f(q)\in\reals$ by \cref{funs_type_apply}.
                Thus $d((f(p),f(q))) < \epsilon$ by \cref{standard_metric_value,real_abs_sub_swap,real_lt_subst_left}.
                Thus $f(q)\in\metricBall{d}{\reals}{f(p)}{\epsilon}$ by \cref{metric_ball}.
            \end{subproof}
            Show that for all $y$ such that $y\in\img{f}{U}$ we have
                $y\in\metricBall{d}{\reals}{f(p)}{\epsilon}$.
            \begin{subproof}
                Fix $y$.
                Assume $y\in\img{f}{U}$.
                Take $q\in U$ such that $q\mathrel{f} y$ by \cref{img_iff}.
                $f$ is a function.
                Thus $y = f(q)$ by \cref{function_apply_intro}.
                Thus $y\in\metricBall{d}{\reals}{f(p)}{\epsilon}$.
            \end{subproof}
            Thus $\img{f}{U}\subseteq \metricBall{d}{\reals}{f(p)}{\epsilon}$ by \cref{subseteq}.
            $\metricBall{d}{\reals}{f(p)}{\epsilon}\subseteq V$ by \cref{subseteq}.
            Thus $\img{f}{U}\subseteq V$ by \cref{subseteq_transitive}.
        \end{subproof}
        \end{subproof}
        $\standardBasisR$ is a basis on $\reals$ by \cref{standard_basis_is_basis}.
        Thus $T_1$ is a topology on $\reals$ by \cref{basis_topology_is_topology,standard_topology_reals}.
        Show that for all $y$ such that $y\in\reals\setminus\{0\}$ we have $y\in\reals$.
        \begin{subproof}
            Fix $y$.
            Assume $y\in\reals\setminus\{0\}$.
            Thus $y\in\reals$ by \cref{setminus_elim_left}.
        \end{subproof}
        Thus $\reals\setminus\{0\}\subseteq\reals$ by \cref{subseteq}.
        Thus $T_2$ is a topology on $\reals\setminus\{0\}$ by \cref{subspace_topology_is_topology}.
        Thus $f$ is continuous at $p$ from $\reals\setminus\{0\}$ to $\reals$ with respect to $T_2$ and $T_1$ by \cref{continuous_at}.
    \end{subproof}
    $\standardBasisR$ is a basis on $\reals$ by \cref{standard_basis_is_basis}.
    Thus $T_1$ is a topology on $\reals$ by \cref{basis_topology_is_topology,standard_topology_reals}.
    Show that for all $y$ such that $y\in\reals\setminus\{0\}$ we have $y\in\reals$.
    \begin{subproof}
        Fix $y$.
        Assume $y\in\reals\setminus\{0\}$.
        Thus $y\in\reals$ by \cref{setminus_elim_left}.
    \end{subproof}
    Thus $\reals\setminus\{0\}\subseteq\reals$ by \cref{subseteq}.
    Thus $T_2$ is a topology on $\reals\setminus\{0\}$ by \cref{subspace_topology_is_topology}.
    Thus $f$ is continuous from $\reals\setminus\{0\}$ to $\reals$ with respect to $T_2$ and $T_1$ by \cref{continuous_at_implies_continuous}.
\end{proof}

% from §21 Item 11: quotient operation is continuous
\begin{lemma}\label{real_division_continuous}
    Let $T_1 = \standardTopologyR$.
    Let $T_2 = \subspaceTopology{T_1}{\reals\setminus\{0\}}$.
    Let $T = \productTopology{T_1}{T_2}{\reals}{\reals\setminus\{0\}}$.
    Then $\realDivision$ is continuous from $\reals\times(\reals\setminus\{0\})$ to $\reals$
        with respect to $T$ and $\standardTopologyR$.
\end{lemma}
\begin{proof}
    Let $T_3 = \productTopology{T_1}{T_1}{\reals}{\reals}$.
    $\standardBasisR$ is a basis on $\reals$ by \cref{standard_basis_is_basis}.
    Thus $T_1$ is a topology on $\reals$ by \cref{basis_topology_is_topology,standard_topology_reals}.
    Show that for all $y$ such that $y\in\reals\setminus\{0\}$ we have $y\in\reals$.
    \begin{subproof}
        Fix $y$.
        Assume $y\in\reals\setminus\{0\}$.
        Thus $y\in\reals$ by \cref{setminus_elim_left}.
    \end{subproof}
    Thus $\reals\setminus\{0\}\subseteq\reals$ by \cref{subseteq}.
    Thus $T_2$ is a topology on $\reals\setminus\{0\}$ by \cref{subspace_topology_is_topology}.
    Let $f_1 = \piOne{\reals}{\reals\setminus\{0\}}$.
    Let $f_2 = \piTwo{\reals}{\reals\setminus\{0\}}$.
    $f_1$ is continuous from $\reals\times(\reals\setminus\{0\})$ to $\reals$ with respect to $T$ and $T_1$
        by \cref{projection_one_continuous}.
    $f_2$ is continuous from $\reals\times(\reals\setminus\{0\})$ to $\reals\setminus\{0\}$ with respect to $T$ and $T_2$
        by \cref{projection_two_continuous}.
    $\realInverse$ is continuous from $\reals\setminus\{0\}$ to $\reals$ with respect to $T_2$ and $T_1$
        by \cref{real_inverse_continuous}.
    Let $g = \realInverse\circ f_2$.
    Thus $g$ is continuous from $\reals\times(\reals\setminus\{0\})$ to $\reals$ with respect to $T$ and $T_1$
        by \cref{continuous_composite}.
    Let $k = \{(p,(f_1(p), g(p))) \mid p\in\reals\times(\reals\setminus\{0\})\}$.
    Thus $k$ is continuous from $\reals\times(\reals\setminus\{0\})$ to $\reals\times\reals$
        with respect to $T$ and $T_3$ by \cref{continuous_map_into_product_backward}.
    $\realMultiplication$ is continuous from $\reals\times\reals$ to $\reals$ with respect to $T_3$ and $T_1$
        by \cref{real_multiplication_continuous}.
    Let $m = \realMultiplication\circ k$.
    Thus $m$ is continuous from $\reals\times(\reals\setminus\{0\})$ to $\reals$ with respect to $T$ and $T_1$
        by \cref{continuous_composite}.
    Show that $m = \realDivision$.
    \begin{subproof}
        $m\in\funs{\reals\times(\reals\setminus\{0\})}{\reals}$ by \cref{continuous}.
        Thus $\dom{m} = \reals\times(\reals\setminus\{0\})$ by \cref{funs}.
        Show that $\realDivision$ is a function from $\reals\times(\reals\setminus\{0\})$ to $\reals$.
        \begin{subproof}
            Show that $\realDivision$ is a function.
            \begin{subproof}
                Show that for all $p,z_1,z_2$ such that $(p,z_1)\in\realDivision$ and $(p,z_2)\in\realDivision$ we have $z_1 = z_2$.
                \begin{subproof}
                    Fix $p$.
                    Fix $z_1$.
                    Fix $z_2$.
                    Assume $(p,z_1)\in\realDivision$ and $(p,z_2)\in\realDivision$.
                    Take $p_1,w_1$ such that
                        $p_1\in\reals\times(\reals\setminus\{0\})$ and $w_1\in\reals$ and
                        $(p,z_1) = (p_1,w_1)$ and $w_1 = \fst{p_1} \rmul \inv{\snd{p_1}}$
                        by \cref{real_division_operation}.
                    Take $p_2,w_2$ such that
                        $p_2\in\reals\times(\reals\setminus\{0\})$ and $w_2\in\reals$ and
                        $(p,z_2) = (p_2,w_2)$ and $w_2 = \fst{p_2} \rmul \inv{\snd{p_2}}$
                        by \cref{real_division_operation}.
                    $p = p_1$ and $z_1 = w_1$ by \cref{pair_eq_iff}.
                    $p = p_2$ and $z_2 = w_2$ by \cref{pair_eq_iff}.
                    Thus $z_1 = \fst{p} \rmul \inv{\snd{p}}$.
                    Thus $z_2 = \fst{p} \rmul \inv{\snd{p}}$.
                    Thus $z_1 = z_2$.
                \end{subproof}
                Thus $\realDivision$ is right-unique by \cref{rightunique}.
                Show that for all $w$ such that $w\in\realDivision$ there exist $x,y$ such that $w = (x,y)$.
                \begin{subproof}
                    Fix $w$.
                    Assume $w\in\realDivision$.
                    Take $p,z$ such that $p\in\reals\times(\reals\setminus\{0\})$ and $z\in\reals$ and
                        $w = (p,z)$ and $z = \fst{p} \rmul \inv{\snd{p}}$ by \cref{real_division_operation}.
                    Thus there exist $x,y$ such that $w = (x,y)$.
                \end{subproof}
                Thus $\realDivision$ is a relation by \cref{relation}.
                Thus $\realDivision$ is a function.
            \end{subproof}
            Show that for all $p\in\dom{\realDivision}$ we have $\apply{\realDivision}{p}\in\reals$.
            \begin{subproof}
                Fix $p\in\dom{\realDivision}$.
                Take $z$ such that $(p,z)\in\realDivision$ by \cref{dom_iff}.
                Take $p_0,w$ such that
                    $p_0\in\reals\times(\reals\setminus\{0\})$ and $w\in\reals$ and
                    $(p,z) = (p_0,w)$ and $w = \fst{p_0} \rmul \inv{\snd{p_0}}$
                    by \cref{real_division_operation}.
                $z = w$ by \cref{pair_eq_iff}.
                Thus $z\in\reals$.
                $\realDivision$ is a function.
                Thus $\apply{\realDivision}{p} = z$ by \cref{function_apply_intro}.
                Thus $\apply{\realDivision}{p}\in\reals$.
            \end{subproof}
            Thus $\realDivision$ is a function to $\reals$.
            Show that $\dom{\realDivision} = \reals\times(\reals\setminus\{0\})$.
            \begin{subproof}
                Show that for all $p\in\dom{\realDivision}$ we have $p\in\reals\times(\reals\setminus\{0\})$.
                \begin{subproof}
                    Fix $p\in\dom{\realDivision}$.
                    Take $z$ such that $(p,z)\in\realDivision$ by \cref{dom_iff}.
                    Take $p_0,w$ such that
                        $p_0\in\reals\times(\reals\setminus\{0\})$ and $w\in\reals$ and
                        $(p,z) = (p_0,w)$ and $w = \fst{p_0} \rmul \inv{\snd{p_0}}$
                        by \cref{real_division_operation}.
                    Thus $p = p_0$ by \cref{pair_eq_iff}.
                    Thus $p\in\reals\times(\reals\setminus\{0\})$.
                \end{subproof}
                Thus $\dom{\realDivision}\subseteq\reals\times(\reals\setminus\{0\})$ by \cref{subseteq}.
                Show that for all $p\in\reals\times(\reals\setminus\{0\})$ we have $p\in\dom{\realDivision}$.
                \begin{subproof}
                    Fix $p\in\reals\times(\reals\setminus\{0\})$.
                    Take $a,b$ such that $p = (a,b)$ and $a\in\reals$ and $b\in\reals\setminus\{0\}$ by \cref{times_elem_is_tuple}.
                    Thus $\fst{p} = a$ and $\snd{p} = b$ by \cref{fst_eq,snd_eq}.
                    Thus $\fst{p}\in\reals$ and $\snd{p}\in\reals\setminus\{0\}$.
                    $\snd{p}\in\reals$ by \cref{setminus_elim_left}.
                    $\snd{p}\notin\{0\}$ by \cref{setminus_elim_right}.
                    Thus $\snd{p}\neq 0$ by \cref{singleton_iff}.
                    $\inv{\snd{p}}\in\reals$ by \cref{real_mul_inverse}.
                    $\fst{p} \rmul \inv{\snd{p}}\in\reals$ by \cref{real_mul_closed}.
                    Let $z = \fst{p} \rmul \inv{\snd{p}}$.
                    Let $w = (p,z)$.
                    $w\in\realDivision$ by \cref{real_division_operation}.
                    Thus $p\in\dom{\realDivision}$ by \cref{dom_intro}.
                \end{subproof}
                Thus $\reals\times(\reals\setminus\{0\})\subseteq\dom{\realDivision}$ by \cref{subseteq}.
                Thus $\dom{\realDivision} = \reals\times(\reals\setminus\{0\})$ by \cref{subseteq_antisymmetric}.
            \end{subproof}
            Thus $\realDivision\in\funs{\reals\times(\reals\setminus\{0\})}{\reals}$ by \cref{funs_intro}.
        \end{subproof}
        Show that for all $p\in\reals\times(\reals\setminus\{0\})$ we have $m(p) = \apply{\realDivision}{p}$.
        \begin{subproof}
            Fix $p\in\reals\times(\reals\setminus\{0\})$.
            $k\in\funs{\reals\times(\reals\setminus\{0\})}{\reals\times\reals}$ by \cref{continuous}.
            Thus $\dom{k} = \reals\times(\reals\setminus\{0\})$ by \cref{funs}.
            $(p,(f_1(p), g(p)))\in k$.
            $k$ is a function.
            Thus $k$ is right-unique and $k$ is a relation.
            Thus $k(p) = (f_1(p), g(p))$ by \cref{function_apply_intro}.
            $\realMultiplication\in\funs{\reals\times\reals}{\reals}$ by \cref{continuous}.
            Thus $\dom{\realMultiplication} = \reals\times\reals$ by \cref{funs}.
            $\ran{k}\subseteq\reals\times\reals$ by \cref{funs_ran}.
            Thus $\realMultiplication$ is a function by \cref{funs_is_function}.
            Thus $\realMultiplication$ is right-unique and $\realMultiplication$ is a relation.
            Thus $\realMultiplication$ is composable with $k$.
            Thus $p\in\dom{k}$ by \cref{funs}.
            Thus $m(p) = \apply{\realMultiplication\circ k}{p}$.
            Thus $m(p) = \apply{\realMultiplication}{k(p)}$ by \cref{circ_apply}.
            $f_1\in\funs{\reals\times(\reals\setminus\{0\})}{\reals}$ by \cref{continuous}.
            $g\in\funs{\reals\times(\reals\setminus\{0\})}{\reals}$ by \cref{continuous}.
            $f_1(p)\in\reals$ by \cref{funs_type_apply}.
            $g(p)\in\reals$ by \cref{funs_type_apply}.
            $(f_1(p), g(p))\in\reals\times\reals$ by \cref{times_tuple_intro}.
            Thus $k(p)\in\reals\times\reals$.
            $f_1(p) \rmul g(p)\in\reals$ by \cref{real_mul_closed}.
            Let $z_0 = f_1(p) \rmul g(p)$.
            $\fst{k(p)} = f_1(p)$ and $\snd{k(p)} = g(p)$ by \cref{fst_eq,snd_eq}.
            Thus $z_0 = \fst{k(p)} \rmul \snd{k(p)}$.
            $(k(p), z_0)\in\realMultiplication$ by \cref{real_multiplication_operation}.
            $\realMultiplication$ is a function.
            Thus $\apply{\realMultiplication}{k(p)} = z_0$ by \cref{function_apply_intro}.
            Thus $m(p) = f_1(p) \rmul g(p)$.
            Take $a,b$ such that $p = (a,b)$ and $a\in\reals$ and $b\in\reals\setminus\{0\}$ by \cref{times_elem_is_tuple}.
            $p\mathrel{f_1} a$ by \cref{projection_first}.
            $f_1$ is a function.
            Thus $f_1(p) = a$ by \cref{function_apply_intro}.
            $p\mathrel{f_2} b$ by \cref{projection_second}.
            $f_2$ is a function.
            Thus $f_2(p) = b$ by \cref{function_apply_intro}.
            $f_2\in\funs{\reals\times(\reals\setminus\{0\})}{\reals\setminus\{0\}}$ by \cref{continuous}.
            Thus $\ran{f_2}\subseteq\reals\setminus\{0\}$ by \cref{funs_ran}.
            $\realInverse\in\funs{\reals\setminus\{0\}}{\reals}$ by \cref{continuous}.
            Thus $\dom{\realInverse} = \reals\setminus\{0\}$ by \cref{funs}.
            $f_2$ is right-unique and $f_2$ is a relation.
            $\realInverse$ is right-unique and $\realInverse$ is a relation.
            Thus $\ran{f_2}\subseteq\dom{\realInverse}$ by \cref{subseteq}.
            Thus $\realInverse$ is composable with $f_2$.
            Thus $p\in\dom{f_2}$ by \cref{funs}.
            Thus $g(p) = \apply{\realInverse}{f_2(p)}$ by \cref{circ_apply}.
            $\fst{p} = a$ and $\snd{p} = b$ by \cref{fst_eq,snd_eq}.
            Thus $f_1(p) = \fst{p}$.
            Thus $f_2(p) = \snd{p}$.
            Thus $g(p) = \apply{\realInverse}{\snd{p}}$.
            $\snd{p}\in\reals\setminus\{0\}$.
            $\snd{p}\notin\{0\}$ by \cref{setminus_elim_right}.
            Thus $\snd{p}\neq 0$ by \cref{singleton_iff}.
            $\snd{p}\in\reals$ by \cref{setminus_elim_left}.
            $\inv{\snd{p}}\in\reals$ by \cref{real_mul_inverse}.
            $(\snd{p},\inv{\snd{p}})\in\realInverse$ by \cref{real_inverse_operation}.
            $\realInverse$ is a function.
            Thus $\apply{\realInverse}{\snd{p}} = \inv{\snd{p}}$ by \cref{function_apply_intro}.
            Thus $g(p) = \inv{\snd{p}}$.
            Thus $m(p) = \fst{p} \rmul \inv{\snd{p}}$.
            $\realDivision$ is a function.
            $\fst{p} \rmul \inv{\snd{p}}\in\reals$ by \cref{real_mul_closed,real_mul_inverse}.
            $(p,\fst{p} \rmul \inv{\snd{p}})\in\realDivision$ by \cref{real_division_operation}.
            Thus $\apply{\realDivision}{p} = \fst{p} \rmul \inv{\snd{p}}$ by \cref{function_apply_intro}.
            Thus $m(p) = \apply{\realDivision}{p}$.
        \end{subproof}
        Show that $m\subseteq\realDivision$.
        \begin{subproof}
            $m$ is a function.
            $\realDivision$ is a function.
            $\dom{m} = \reals\times(\reals\setminus\{0\})$ by \cref{funs}.
            $\dom{\realDivision} = \reals\times(\reals\setminus\{0\})$ by \cref{funs}.
            Thus $\dom{m}\subseteq\dom{\realDivision}$ by \cref{subseteq}.
            Show that for all $x\in\dom{m}$ we have $m(x) = \apply{\realDivision}{x}$.
            \begin{subproof}
                Fix $x\in\dom{m}$.
                Thus $x\in\reals\times(\reals\setminus\{0\})$ by \cref{funs}.
                Thus $m(x) = \apply{\realDivision}{x}$.
            \end{subproof}
            Thus $m\subseteq\realDivision$ by \cref{fun_subseteq}.
        \end{subproof}
        Show that $\realDivision\subseteq m$.
        \begin{subproof}
            $m$ is a function.
            $\realDivision$ is a function.
            $\dom{m} = \reals\times(\reals\setminus\{0\})$ by \cref{funs}.
            $\dom{\realDivision} = \reals\times(\reals\setminus\{0\})$ by \cref{funs}.
            Thus $\dom{\realDivision}\subseteq\dom{m}$ by \cref{subseteq}.
            Show that for all $x\in\dom{\realDivision}$ we have $m(x) = \apply{\realDivision}{x}$.
            \begin{subproof}
                Fix $x\in\dom{\realDivision}$.
                Thus $x\in\reals\times(\reals\setminus\{0\})$ by \cref{funs}.
                Thus $m(x) = \apply{\realDivision}{x}$.
            \end{subproof}
            Thus $\realDivision\subseteq m$ by \cref{fun_subseteq}.
        \end{subproof}
        Thus $m = \realDivision$ by \cref{subseteq_antisymmetric}.
    \end{subproof}
    Thus $\realDivision$ is continuous from $\reals\times(\reals\setminus\{0\})$ to $\reals$
        with respect to $T$ and $\standardTopologyR$ by \cref{continuous_composite}.
\end{proof}

% from §21 Item 12: sum of continuous functions
\begin{theorem}\label{real_function_addition_continuous}
    Suppose $f$ is continuous from $X$ to $\reals$ with respect to $T$ and $\standardTopologyR$.
    Suppose $g$ is continuous from $X$ to $\reals$ with respect to $T$ and $\standardTopologyR$.
    Let $h = \{(x, f(x) + g(x)) \mid x\in X\}$.
    Then $h$ is continuous from $X$ to $\reals$ with respect to $T$ and $\standardTopologyR$.
\end{theorem}
\begin{proof}
    Let $T_1 = \standardTopologyR$.
    Let $T_2 = \productTopology{T_1}{T_1}{\reals}{\reals}$.
    Let $k = \{(x,(f(x), g(x))) \mid x\in X\}$.
    Thus $k$ is continuous from $X$ to $\reals\times\reals$ with respect to $T$ and $T_2$
        by \cref{continuous_map_into_product_backward}.
    $\realAddition$ is continuous from $\reals\times\reals$ to $\reals$ with respect to $T_2$ and $T_1$
        by \cref{real_addition_continuous}.
    Let $m = \realAddition\circ k$.
    Thus $m$ is continuous from $X$ to $\reals$ with respect to $T$ and $T_1$ by \cref{continuous_composite}.
    Show that $h$ is a function from $X$ to $\reals$.
    \begin{subproof}
        Show that $h$ is a function.
        \begin{subproof}
            Show that for all $x,z_1,z_2$ such that $(x,z_1)\in h$ and $(x,z_2)\in h$ we have $z_1 = z_2$.
            \begin{subproof}
                Fix $x$.
                Fix $z_1$.
                Fix $z_2$.
                Assume $(x,z_1)\in h$ and $(x,z_2)\in h$.
                Take $x_1\in X$ such that $(x,z_1) = (x_1, f(x_1) + g(x_1))$.
                Thus $x = x_1$ and $z_1 = f(x_1) + g(x_1)$ by \cref{pair_eq_iff}.
                Thus $x\in X$ and $z_1 = f(x) + g(x)$.
                Take $x_2\in X$ such that $(x,z_2) = (x_2, f(x_2) + g(x_2))$.
                Thus $x = x_2$ and $z_2 = f(x_2) + g(x_2)$ by \cref{pair_eq_iff}.
                Thus $x\in X$ and $z_2 = f(x) + g(x)$.
                Thus $z_1 = z_2$.
            \end{subproof}
            Thus $h$ is right-unique by \cref{rightunique}.
            Show that for all $w$ such that $w\in h$ there exist $x,y$ such that $w = (x,y)$.
            \begin{subproof}
                Fix $w$.
                Assume $w\in h$.
                Take $x\in X$ such that $w = (x, f(x) + g(x))$.
                Thus there exist $x_0,y_0$ such that $w = (x_0,y_0)$.
            \end{subproof}
            Thus $h$ is a relation by \cref{relation}.
            Thus $h$ is a function.
        \end{subproof}
        Show that for all $x\in\dom{h}$ we have $h(x)\in\reals$.
        \begin{subproof}
            Fix $x\in\dom{h}$.
            Take $z$ such that $(x,z)\in h$ by \cref{dom_iff}.
            Thus $x\in X$ and $z = f(x) + g(x)$.
            $f\in\funs{X}{\reals}$ and $g\in\funs{X}{\reals}$ by \cref{continuous}.
            Thus $f(x)\in\reals$ and $g(x)\in\reals$ by \cref{funs_type_apply}.
            Thus $z\in\reals$ by \cref{real_add_closed}.
            Thus $h(x) = z$ by \cref{function_apply_intro}.
            Thus $h(x)\in\reals$.
        \end{subproof}
        Thus $h$ is a function to $\reals$.
        Show that $\dom{h} = X$.
        \begin{subproof}
            Show that for all $x\in\dom{h}$ we have $x\in X$.
            \begin{subproof}
                Fix $x\in\dom{h}$.
                Take $z$ such that $(x,z)\in h$ by \cref{dom_iff}.
                Thus $x\in X$.
            \end{subproof}
            Thus $\dom{h}\subseteq X$ by \cref{subseteq}.
            Show that for all $x\in X$ we have $x\in\dom{h}$.
            \begin{subproof}
                Fix $x\in X$.
                Let $z = f(x) + g(x)$.
                $(x,z)\in h$.
                Thus $x\in\dom{h}$ by \cref{dom_intro}.
            \end{subproof}
            Thus $X\subseteq\dom{h}$ by \cref{subseteq}.
            Thus $\dom{h} = X$ by \cref{subseteq_antisymmetric}.
        \end{subproof}
        Thus $h\in\funs{X}{\reals}$ by \cref{funs_intro}.
    \end{subproof}
    Show that for all $x\in X$ we have $m(x) = \apply{h}{x}$.
    \begin{subproof}
        Fix $x\in X$.
        $k\in\funs{X}{\reals\times\reals}$ by \cref{continuous}.
        $(x,(f(x), g(x)))\in k$.
        $k$ is a function.
        Thus $k$ is right-unique and $k$ is a relation.
        Thus $k(x) = (f(x), g(x))$ by \cref{function_apply_intro}.
        $f\in\funs{X}{\reals}$ and $g\in\funs{X}{\reals}$ by \cref{continuous}.
        Thus $f(x)\in\reals$ and $g(x)\in\reals$ by \cref{funs_type_apply}.
        $(f(x), g(x))\in\reals\times\reals$ by \cref{times_tuple_intro}.
        Thus $k(x)\in\reals\times\reals$.
        $f(x) + g(x)\in\reals$ by \cref{real_add_closed}.
        Let $p_0 = (f(x), g(x))$.
        Let $z_0 = f(x) + g(x)$.
        $\fst{p_0} = f(x)$ and $\snd{p_0} = g(x)$ by \cref{fst_eq,snd_eq}.
        Thus $z_0 = \fst{p_0} + \snd{p_0}$.
        Thus $k(x) = p_0$ by \cref{function_apply_intro}.
        Thus $(k(x), z_0) = (p_0, z_0)$.
        Thus $(k(x), z_0)\in\realAddition$ by \cref{real_addition_operation}.
        $\realAddition\in\funs{\reals\times\reals}{\reals}$ by \cref{continuous,real_addition_continuous}.
        Thus $\realAddition$ is a function by \cref{funs_is_function}.
        Thus $\realAddition$ is right-unique and $\realAddition$ is a relation.
        Thus $\apply{\realAddition}{k(x)} = f(x) + g(x)$ by \cref{function_apply_intro}.
        Thus $\ran{k}\subseteq\reals\times\reals$ by \cref{funs_ran}.
        Thus $\dom{\realAddition} = \reals\times\reals$ by \cref{funs}.
        Thus $\ran{k}\subseteq\dom{\realAddition}$ by \cref{subseteq}.
        Thus $\realAddition$ is composable with $k$.
        Thus $x\in\dom{k}$ by \cref{funs}.
        Thus $m(x) = \apply{\realAddition}{k(x)}$ by \cref{circ_apply}.
        $(x, f(x) + g(x))\in h$.
        $h$ is right-unique and $h$ is a relation.
        Thus $\apply{h}{x} = f(x) + g(x)$ by \cref{function_apply_intro}.
        Thus $m(x) = \apply{h}{x}$.
    \end{subproof}
    Show that $m\subseteq h$.
    \begin{subproof}
        $m$ is a function.
        $h$ is a function.
        $\dom{m} = X$ by \cref{funs,continuous}.
        $\dom{h} = X$ by \cref{funs}.
        Thus $\dom{m}\subseteq\dom{h}$ by \cref{subseteq}.
        Show that for all $x\in\dom{m}$ we have $m(x) = \apply{h}{x}$.
        \begin{subproof}
            Fix $x\in\dom{m}$.
            Thus $x\in X$ by \cref{funs,continuous}.
            Thus $m(x) = \apply{h}{x}$.
        \end{subproof}
        Thus $m\subseteq h$ by \cref{fun_subseteq}.
    \end{subproof}
    Show that $h\subseteq m$.
    \begin{subproof}
        $m$ is a function.
        $h$ is a function.
        $\dom{m} = X$ by \cref{funs,continuous}.
        $\dom{h} = X$ by \cref{funs}.
        Thus $\dom{h}\subseteq\dom{m}$ by \cref{subseteq}.
        Show that for all $x\in\dom{h}$ we have $m(x) = \apply{h}{x}$.
        \begin{subproof}
            Fix $x\in\dom{h}$.
            Thus $x\in X$ by \cref{funs}.
            Thus $m(x) = \apply{h}{x}$.
        \end{subproof}
        Thus $h\subseteq m$ by \cref{fun_subseteq}.
    \end{subproof}
    Thus $m = h$ by \cref{subseteq_antisymmetric}.
    Thus $h$ is continuous from $X$ to $\reals$ with respect to $T$ and $\standardTopologyR$ by \cref{continuous_composite}.
\end{proof}

% from §21 Item 13: difference of continuous functions
\begin{theorem}\label{real_function_subtraction_continuous}
    Suppose $f$ is continuous from $X$ to $\reals$ with respect to $T$ and $\standardTopologyR$.
    Suppose $g$ is continuous from $X$ to $\reals$ with respect to $T$ and $\standardTopologyR$.
    Let $h = \{(x, f(x) - g(x)) \mid x\in X\}$.
    Then $h$ is continuous from $X$ to $\reals$ with respect to $T$ and $\standardTopologyR$.
\end{theorem}
\begin{proof}
    Let $T_1 = \standardTopologyR$.
    Let $T_2 = \productTopology{T_1}{T_1}{\reals}{\reals}$.
    Let $k = \{(x,(f(x), g(x))) \mid x\in X\}$.
    Thus $k$ is continuous from $X$ to $\reals\times\reals$ with respect to $T$ and $T_2$
        by \cref{continuous_map_into_product_backward}.
    $\realSubtraction$ is continuous from $\reals\times\reals$ to $\reals$ with respect to $T_2$ and $T_1$
        by \cref{real_subtraction_continuous}.
    Let $m = \realSubtraction\circ k$.
    Thus $m$ is continuous from $X$ to $\reals$ with respect to $T$ and $T_1$ by \cref{continuous_composite}.
    Show that $h$ is a function from $X$ to $\reals$.
    \begin{subproof}
        Show that $h$ is a function.
        \begin{subproof}
            Show that for all $x,z_1,z_2$ such that $(x,z_1)\in h$ and $(x,z_2)\in h$ we have $z_1 = z_2$.
            \begin{subproof}
                Fix $x$.
                Fix $z_1$.
                Fix $z_2$.
                Assume $(x,z_1)\in h$ and $(x,z_2)\in h$.
                Take $x_1\in X$ such that $(x,z_1) = (x_1, f(x_1) - g(x_1))$.
                Thus $x = x_1$ and $z_1 = f(x_1) - g(x_1)$ by \cref{pair_eq_iff}.
                Thus $x\in X$ and $z_1 = f(x) - g(x)$.
                Take $x_2\in X$ such that $(x,z_2) = (x_2, f(x_2) - g(x_2))$.
                Thus $x = x_2$ and $z_2 = f(x_2) - g(x_2)$ by \cref{pair_eq_iff}.
                Thus $x\in X$ and $z_2 = f(x) - g(x)$.
                Thus $z_1 = z_2$.
            \end{subproof}
            Thus $h$ is right-unique by \cref{rightunique}.
            Show that for all $w$ such that $w\in h$ there exist $x,y$ such that $w = (x,y)$.
            \begin{subproof}
                Fix $w$.
                Assume $w\in h$.
                Take $x\in X$ such that $w = (x, f(x) - g(x))$.
                Thus there exist $x_0,y_0$ such that $w = (x_0,y_0)$.
            \end{subproof}
            Thus $h$ is a relation by \cref{relation}.
            Thus $h$ is a function.
        \end{subproof}
        Show that for all $x\in\dom{h}$ we have $h(x)\in\reals$.
        \begin{subproof}
            Fix $x\in\dom{h}$.
            Take $z$ such that $(x,z)\in h$ by \cref{dom_iff}.
            Thus $x\in X$ and $z = f(x) - g(x)$.
            $f\in\funs{X}{\reals}$ and $g\in\funs{X}{\reals}$ by \cref{continuous}.
            Thus $f(x)\in\reals$ and $g(x)\in\reals$ by \cref{funs_type_apply}.
            $\neg{g(x)}\in\reals$ by \cref{real_add_inverse}.
            Thus $z\in\reals$ by \cref{real_sub,real_add_closed}.
            Thus $h(x) = z$ by \cref{function_apply_intro}.
            Thus $h(x)\in\reals$.
        \end{subproof}
        Thus $h$ is a function to $\reals$.
        Show that $\dom{h} = X$.
        \begin{subproof}
            Show that for all $x\in\dom{h}$ we have $x\in X$.
            \begin{subproof}
                Fix $x\in\dom{h}$.
                Take $z$ such that $(x,z)\in h$ by \cref{dom_iff}.
                Thus $x\in X$.
            \end{subproof}
            Thus $\dom{h}\subseteq X$ by \cref{subseteq}.
            Show that for all $x\in X$ we have $x\in\dom{h}$.
            \begin{subproof}
                Fix $x\in X$.
                Let $z = f(x) - g(x)$.
                $(x,z)\in h$.
                Thus $x\in\dom{h}$ by \cref{dom_intro}.
            \end{subproof}
            Thus $X\subseteq\dom{h}$ by \cref{subseteq}.
            Thus $\dom{h} = X$ by \cref{subseteq_antisymmetric}.
        \end{subproof}
        Thus $h\in\funs{X}{\reals}$ by \cref{funs_intro}.
    \end{subproof}
    Show that for all $x\in X$ we have $m(x) = \apply{h}{x}$.
    \begin{subproof}
        Fix $x\in X$.
        $k\in\funs{X}{\reals\times\reals}$ by \cref{continuous}.
        $(x,(f(x), g(x)))\in k$.
        $k$ is a function.
        Thus $k(x) = (f(x), g(x))$ by \cref{function_apply_intro}.
        $f\in\funs{X}{\reals}$ and $g\in\funs{X}{\reals}$ by \cref{continuous}.
        Thus $f(x)\in\reals$ and $g(x)\in\reals$ by \cref{funs_type_apply}.
        $\neg{g(x)}\in\reals$ by \cref{real_add_inverse}.
        $(f(x), g(x))\in\reals\times\reals$ by \cref{times_tuple_intro}.
        Thus $k(x)\in\reals\times\reals$.
        $f(x) - g(x)\in\reals$ by \cref{real_sub,real_add_closed}.
        Let $p_0 = (f(x), g(x))$.
        Let $z_0 = f(x) - g(x)$.
        $\fst{p_0} = f(x)$ and $\snd{p_0} = g(x)$ by \cref{fst_eq,snd_eq}.
        Thus $z_0 = \fst{p_0} - \snd{p_0}$.
        Thus $k(x) = p_0$ by \cref{function_apply_intro}.
        Thus $(k(x), z_0) = (p_0, z_0)$.
        Thus $(k(x), z_0)\in\realSubtraction$ by \cref{real_subtraction_operation}.
        $\realSubtraction\in\funs{\reals\times\reals}{\reals}$ by \cref{continuous,real_subtraction_continuous}.
        Thus $\realSubtraction$ is a function by \cref{funs_is_function}.
        Thus $\apply{\realSubtraction}{k(x)} = f(x) - g(x)$ by \cref{function_apply_intro}.
        Thus $\ran{k}\subseteq\reals\times\reals$ by \cref{funs_ran}.
        Thus $\dom{\realSubtraction} = \reals\times\reals$ by \cref{funs}.
        Thus $\ran{k}\subseteq\dom{\realSubtraction}$ by \cref{subseteq}.
        Thus $\realSubtraction$ is composable with $k$.
        Thus $x\in\dom{k}$ by \cref{funs}.
        Thus $m(x) = \apply{\realSubtraction}{k(x)}$ by \cref{circ_apply}.
        $(x, f(x) - g(x))\in h$.
        $h$ is right-unique and $h$ is a relation.
        Thus $\apply{h}{x} = f(x) - g(x)$ by \cref{function_apply_intro}.
        Thus $m(x) = \apply{h}{x}$.
    \end{subproof}
    Show that $m\subseteq h$.
    \begin{subproof}
        $m$ is a function.
        $h$ is a function.
        $\dom{m} = X$ by \cref{funs,continuous}.
        $\dom{h} = X$ by \cref{funs}.
        Thus $\dom{m}\subseteq\dom{h}$ by \cref{subseteq}.
        Show that for all $x\in\dom{m}$ we have $m(x) = \apply{h}{x}$.
        \begin{subproof}
            Fix $x\in\dom{m}$.
            Thus $x\in X$ by \cref{funs,continuous}.
            Thus $m(x) = \apply{h}{x}$.
        \end{subproof}
        Thus $m\subseteq h$ by \cref{fun_subseteq}.
    \end{subproof}
    Show that $h\subseteq m$.
    \begin{subproof}
        $m$ is a function.
        $h$ is a function.
        $\dom{m} = X$ by \cref{funs,continuous}.
        $\dom{h} = X$ by \cref{funs}.
        Thus $\dom{h}\subseteq\dom{m}$ by \cref{subseteq}.
        Show that for all $x\in\dom{h}$ we have $m(x) = \apply{h}{x}$.
        \begin{subproof}
            Fix $x\in\dom{h}$.
            Thus $x\in X$ by \cref{funs}.
            Thus $m(x) = \apply{h}{x}$.
        \end{subproof}
        Thus $h\subseteq m$ by \cref{fun_subseteq}.
    \end{subproof}
    Thus $m = h$ by \cref{subseteq_antisymmetric}.
    Thus $h$ is continuous from $X$ to $\reals$ with respect to $T$ and $\standardTopologyR$ by \cref{continuous_composite}.
\end{proof}

% from §21 Item 14: product of continuous functions
\begin{theorem}\label{real_function_multiplication_continuous}
    Suppose $f$ is continuous from $X$ to $\reals$ with respect to $T$ and $\standardTopologyR$.
    Suppose $g$ is continuous from $X$ to $\reals$ with respect to $T$ and $\standardTopologyR$.
    Let $h = \{(x, f(x) \rmul g(x)) \mid x\in X\}$.
    Then $h$ is continuous from $X$ to $\reals$ with respect to $T$ and $\standardTopologyR$.
\end{theorem}
\begin{proof}
    Let $T_1 = \standardTopologyR$.
    Let $T_2 = \productTopology{T_1}{T_1}{\reals}{\reals}$.
    Let $k = \{(x,(f(x), g(x))) \mid x\in X\}$.
    Thus $k$ is continuous from $X$ to $\reals\times\reals$ with respect to $T$ and $T_2$
        by \cref{continuous_map_into_product_backward}.
    $\realMultiplication$ is continuous from $\reals\times\reals$ to $\reals$ with respect to $T_2$ and $T_1$
        by \cref{real_multiplication_continuous}.
    Let $m = \realMultiplication\circ k$.
    Thus $m$ is continuous from $X$ to $\reals$ with respect to $T$ and $T_1$ by \cref{continuous_composite}.
    Show that $h$ is a function from $X$ to $\reals$.
    \begin{subproof}
        Show that $h$ is a function.
        \begin{subproof}
            Show that for all $x,z_1,z_2$ such that $(x,z_1)\in h$ and $(x,z_2)\in h$ we have $z_1 = z_2$.
            \begin{subproof}
                Fix $x$.
                Fix $z_1$.
                Fix $z_2$.
                Assume $(x,z_1)\in h$ and $(x,z_2)\in h$.
                Take $x_1\in X$ such that $(x,z_1) = (x_1, f(x_1) \rmul g(x_1))$.
                Thus $x = x_1$ and $z_1 = f(x_1) \rmul g(x_1)$ by \cref{pair_eq_iff}.
                Thus $x\in X$ and $z_1 = f(x) \rmul g(x)$.
                Take $x_2\in X$ such that $(x,z_2) = (x_2, f(x_2) \rmul g(x_2))$.
                Thus $x = x_2$ and $z_2 = f(x_2) \rmul g(x_2)$ by \cref{pair_eq_iff}.
                Thus $x\in X$ and $z_2 = f(x) \rmul g(x)$.
                Thus $z_1 = z_2$.
            \end{subproof}
            Thus $h$ is right-unique by \cref{rightunique}.
            Show that for all $w$ such that $w\in h$ there exist $x,y$ such that $w = (x,y)$.
            \begin{subproof}
                Fix $w$.
                Assume $w\in h$.
                Take $x\in X$ such that $w = (x, f(x) \rmul g(x))$.
                Thus there exist $x_0,y_0$ such that $w = (x_0,y_0)$.
            \end{subproof}
            Thus $h$ is a relation by \cref{relation}.
            Thus $h$ is a function.
        \end{subproof}
        Show that for all $x\in\dom{h}$ we have $h(x)\in\reals$.
        \begin{subproof}
            Fix $x\in\dom{h}$.
            Take $z$ such that $(x,z)\in h$ by \cref{dom_iff}.
            Thus $x\in X$ and $z = f(x) \rmul g(x)$.
            $f\in\funs{X}{\reals}$ and $g\in\funs{X}{\reals}$ by \cref{continuous}.
            Thus $f(x)\in\reals$ and $g(x)\in\reals$ by \cref{funs_type_apply}.
            Thus $z\in\reals$ by \cref{real_mul_closed}.
            Thus $h(x) = z$ by \cref{function_apply_intro}.
            Thus $h(x)\in\reals$.
        \end{subproof}
        Thus $h$ is a function to $\reals$.
        Show that $\dom{h} = X$.
        \begin{subproof}
            Show that for all $x\in\dom{h}$ we have $x\in X$.
            \begin{subproof}
                Fix $x\in\dom{h}$.
                Take $z$ such that $(x,z)\in h$ by \cref{dom_iff}.
                Thus $x\in X$.
            \end{subproof}
            Thus $\dom{h}\subseteq X$ by \cref{subseteq}.
            Show that for all $x\in X$ we have $x\in\dom{h}$.
            \begin{subproof}
                Fix $x\in X$.
                Let $z = f(x) \rmul g(x)$.
                $(x,z)\in h$.
                Thus $x\in\dom{h}$ by \cref{dom_intro}.
            \end{subproof}
            Thus $X\subseteq\dom{h}$ by \cref{subseteq}.
            Thus $\dom{h} = X$ by \cref{subseteq_antisymmetric}.
        \end{subproof}
        Thus $h\in\funs{X}{\reals}$ by \cref{funs_intro}.
    \end{subproof}
    Show that for all $x\in X$ we have $m(x) = \apply{h}{x}$.
    \begin{subproof}
        Fix $x\in X$.
        $k\in\funs{X}{\reals\times\reals}$ by \cref{continuous}.
        $(x,(f(x), g(x)))\in k$.
        $k$ is a function.
        Thus $k(x) = (f(x), g(x))$ by \cref{function_apply_intro}.
        $f\in\funs{X}{\reals}$ and $g\in\funs{X}{\reals}$ by \cref{continuous}.
        Thus $f(x)\in\reals$ and $g(x)\in\reals$ by \cref{funs_type_apply}.
        $(f(x), g(x))\in\reals\times\reals$ by \cref{times_tuple_intro}.
        Thus $k(x)\in\reals\times\reals$.
        $f(x) \rmul g(x)\in\reals$ by \cref{real_mul_closed}.
        Let $p_0 = (f(x), g(x))$.
        Let $z_0 = f(x) \rmul g(x)$.
        $\fst{p_0} = f(x)$ and $\snd{p_0} = g(x)$ by \cref{fst_eq,snd_eq}.
        Thus $z_0 = \fst{p_0} \rmul \snd{p_0}$.
        Thus $k(x) = p_0$ by \cref{function_apply_intro}.
        Thus $(k(x), z_0) = (p_0, z_0)$.
        Thus $(k(x), z_0)\in\realMultiplication$ by \cref{real_multiplication_operation}.
        $\realMultiplication\in\funs{\reals\times\reals}{\reals}$ by \cref{continuous,real_multiplication_continuous}.
        Thus $\realMultiplication$ is a function by \cref{funs_is_function}.
        Thus $\apply{\realMultiplication}{k(x)} = f(x) \rmul g(x)$ by \cref{function_apply_intro}.
        Thus $\ran{k}\subseteq\reals\times\reals$ by \cref{funs_ran}.
        Thus $\dom{\realMultiplication} = \reals\times\reals$ by \cref{funs}.
        Thus $\ran{k}\subseteq\dom{\realMultiplication}$ by \cref{subseteq}.
        Thus $\realMultiplication$ is composable with $k$.
        Thus $x\in\dom{k}$ by \cref{funs}.
        Thus $m(x) = \apply{\realMultiplication}{k(x)}$ by \cref{circ_apply}.
        $(x, f(x) \rmul g(x))\in h$.
        $h$ is right-unique and $h$ is a relation.
        Thus $\apply{h}{x} = f(x) \rmul g(x)$ by \cref{function_apply_intro}.
        Thus $m(x) = \apply{h}{x}$.
    \end{subproof}
    Show that $m\subseteq h$.
    \begin{subproof}
        $m$ is a function.
        $h$ is a function.
        $\dom{m} = X$ by \cref{funs,continuous}.
        $\dom{h} = X$ by \cref{funs}.
        Thus $\dom{m}\subseteq\dom{h}$ by \cref{subseteq}.
        Show that for all $x\in\dom{m}$ we have $m(x) = \apply{h}{x}$.
        \begin{subproof}
            Fix $x\in\dom{m}$.
            Thus $x\in X$ by \cref{funs,continuous}.
            Thus $m(x) = \apply{h}{x}$.
        \end{subproof}
        Thus $m\subseteq h$ by \cref{fun_subseteq}.
    \end{subproof}
    Show that $h\subseteq m$.
    \begin{subproof}
        $m$ is a function.
        $h$ is a function.
        $\dom{m} = X$ by \cref{funs,continuous}.
        $\dom{h} = X$ by \cref{funs}.
        Thus $\dom{h}\subseteq\dom{m}$ by \cref{subseteq}.
        Show that for all $x\in\dom{h}$ we have $m(x) = \apply{h}{x}$.
        \begin{subproof}
            Fix $x\in\dom{h}$.
            Thus $x\in X$ by \cref{funs}.
            Thus $m(x) = \apply{h}{x}$.
        \end{subproof}
        Thus $h\subseteq m$ by \cref{fun_subseteq}.
    \end{subproof}
    Thus $m = h$ by \cref{subseteq_antisymmetric}.
    Thus $h$ is continuous from $X$ to $\reals$ with respect to $T$ and $\standardTopologyR$ by \cref{continuous_composite}.
\end{proof}

% from §21 Item 15: quotient of continuous functions
\begin{theorem}\label{real_function_division_continuous}
    Suppose $f$ is continuous from $X$ to $\reals$ with respect to $T$ and $\standardTopologyR$.
    Suppose $g$ is continuous from $X$ to $\reals$ with respect to $T$ and $\standardTopologyR$.
    Suppose for all $x\in X$ we have $g(x)\neq 0$.
    Let $h = \{(x, f(x) \rmul \inv{(g(x))}) \mid x\in X\}$.
    Then $h$ is continuous from $X$ to $\reals$ with respect to $T$ and $\standardTopologyR$.
\end{theorem}
\begin{proof}
    Let $T_1 = \standardTopologyR$.
    Let $T_2 = \subspaceTopology{T_1}{\reals\setminus\{0\}}$.
    Show that $\img{g}{X}\subseteq \reals\setminus\{0\}$.
    \begin{subproof}
        Show that for all $y$ such that $y\in\img{g}{X}$ we have $y\in\reals\setminus\{0\}$.
        \begin{subproof}
            Fix $y$.
            Assume $y\in\img{g}{X}$.
            Take $x\in X$ such that $x\mathrel{g} y$ by \cref{img_iff}.
            $g$ is a function.
            Thus $y = g(x)$ by \cref{function_apply_intro}.
            $g\in\funs{X}{\reals}$ by \cref{continuous}.
            Thus $g(x)\in\reals$ by \cref{funs_type_apply}.
            $g(x)\neq 0$.
            Show that $y\notin\{0\}$.
            \begin{subproof}
                Suppose not.
                Thus $y\in\{0\}$.
                Thus $y = 0$ by \cref{singleton_iff}.
                Contradiction.
            \end{subproof}
            Thus $y\in\reals\setminus\{0\}$ by \cref{setminus_intro}.
        \end{subproof}
        Thus $\img{g}{X}\subseteq \reals\setminus\{0\}$ by \cref{subseteq}.
    \end{subproof}
    $\standardBasisR$ is a basis on $\reals$ by \cref{standard_basis_is_basis}.
    Thus $T_1$ is a topology on $\reals$ by \cref{basis_topology_is_topology,standard_topology_reals}.
    Show that for all $y$ such that $y\in\reals\setminus\{0\}$ we have $y\in\reals$.
    \begin{subproof}
        Fix $y$.
        Assume $y\in\reals\setminus\{0\}$.
        Thus $y\in\reals$ by \cref{setminus_elim_left}.
    \end{subproof}
    Thus $\reals\setminus\{0\}\subseteq\reals$ by \cref{subseteq}.
    Thus $\reals\setminus\{0\}$ is a subspace of $\reals$ with respect to $T_1$ by \cref{subspace_of}.
    Thus $g$ is continuous from $X$ to $\reals\setminus\{0\}$ with respect to $T$ and $T_2$
        by \cref{continuous_restrict_range}.
    $\realInverse$ is continuous from $\reals\setminus\{0\}$ to $\reals$ with respect to $T_2$ and $T_1$
        by \cref{real_inverse_continuous}.
    Let $j = \realInverse\circ g$.
    Thus $j$ is continuous from $X$ to $\reals$ with respect to $T$ and $T_1$ by \cref{continuous_composite}.
    Let $k = \{(x,(f(x), j(x))) \mid x\in X\}$.
    Thus $k$ is continuous from $X$ to $\reals\times\reals$ with respect to $T$ and
        $\productTopology{T_1}{T_1}{\reals}{\reals}$ by \cref{continuous_map_into_product_backward}.
    $\realMultiplication$ is continuous from $\reals\times\reals$ to $\reals$ with respect to
        $\productTopology{T_1}{T_1}{\reals}{\reals}$ and $T_1$ by \cref{real_multiplication_continuous}.
    Let $m = \realMultiplication\circ k$.
    Thus $m$ is continuous from $X$ to $\reals$ with respect to $T$ and $T_1$ by \cref{continuous_composite}.
    Show that $h$ is a function from $X$ to $\reals$.
    \begin{subproof}
        Show that $h$ is a function.
        \begin{subproof}
            Show that for all $x,z_1,z_2$ such that $(x,z_1)\in h$ and $(x,z_2)\in h$ we have $z_1 = z_2$.
            \begin{subproof}
                Fix $x$.
                Fix $z_1$.
                Fix $z_2$.
                Assume $(x,z_1)\in h$ and $(x,z_2)\in h$.
                Take $x_1\in X$ such that $(x,z_1) = (x_1, f(x_1) \rmul \inv{(g(x_1))})$.
                Thus $x = x_1$ and $z_1 = f(x_1) \rmul \inv{(g(x_1))}$ by \cref{pair_eq_iff}.
                Thus $x\in X$ and $z_1 = f(x) \rmul \inv{(g(x))}$.
                Take $x_2\in X$ such that $(x,z_2) = (x_2, f(x_2) \rmul \inv{(g(x_2))})$.
                Thus $x = x_2$ and $z_2 = f(x_2) \rmul \inv{(g(x_2))}$ by \cref{pair_eq_iff}.
                Thus $x\in X$ and $z_2 = f(x) \rmul \inv{(g(x))}$.
                Thus $z_1 = z_2$.
            \end{subproof}
            Thus $h$ is right-unique by \cref{rightunique}.
            Show that for all $w$ such that $w\in h$ there exist $x,y$ such that $w = (x,y)$.
            \begin{subproof}
                Fix $w$.
                Assume $w\in h$.
                Take $x\in X$ such that $w = (x, f(x) \rmul \inv{(g(x))})$.
                Thus there exist $x_0,y_0$ such that $w = (x_0,y_0)$.
            \end{subproof}
            Thus $h$ is a relation by \cref{relation}.
            Thus $h$ is a function.
        \end{subproof}
        Show that for all $x\in\dom{h}$ we have $h(x)\in\reals$.
        \begin{subproof}
            Fix $x\in\dom{h}$.
            Take $z$ such that $(x,z)\in h$ by \cref{dom_iff}.
            Take $x_0\in X$ such that $(x,z) = (x_0, f(x_0) \rmul \inv{(g(x_0))})$.
            Thus $x = x_0$ and $z = f(x_0) \rmul \inv{(g(x_0))}$ by \cref{pair_eq_iff}.
            Thus $x\in X$ and $z = f(x) \rmul \inv{(g(x))}$.
            $f\in\funs{X}{\reals}$ and $g\in\funs{X}{\reals}$ by \cref{continuous}.
            Thus $f(x)\in\reals$ and $g(x)\in\reals$ by \cref{funs_type_apply}.
            $g(x)\neq 0$.
            $\inv{(g(x))}\in\reals$ by \cref{real_mul_inverse}.
            Thus $z\in\reals$ by \cref{real_mul_closed}.
            Thus $h(x) = z$ by \cref{function_apply_intro}.
            Thus $h(x)\in\reals$.
        \end{subproof}
        Thus $h$ is a function to $\reals$.
        Show that $\dom{h} = X$.
        \begin{subproof}
            Show that for all $x\in\dom{h}$ we have $x\in X$.
            \begin{subproof}
                Fix $x\in\dom{h}$.
                Take $z$ such that $(x,z)\in h$ by \cref{dom_iff}.
                Thus $x\in X$.
            \end{subproof}
            Thus $\dom{h}\subseteq X$ by \cref{subseteq}.
            Show that for all $x\in X$ we have $x\in\dom{h}$.
            \begin{subproof}
                Fix $x\in X$.
                Let $z = f(x) \rmul \inv{(g(x))}$.
                $(x,z)\in h$.
                Thus $x\in\dom{h}$ by \cref{dom_intro}.
            \end{subproof}
            Thus $X\subseteq\dom{h}$ by \cref{subseteq}.
            Thus $\dom{h} = X$ by \cref{subseteq_antisymmetric}.
        \end{subproof}
        Thus $h\in\funs{X}{\reals}$ by \cref{funs_intro}.
    \end{subproof}
    Show that for all $x\in X$ we have $m(x) = \apply{h}{x}$.
    \begin{subproof}
        Fix $x\in X$.
        $k\in\funs{X}{\reals\times\reals}$ by \cref{continuous}.
        $(x,(f(x), j(x)))\in k$.
        $k$ is a function.
        Thus $k(x) = (f(x), j(x))$ by \cref{function_apply_intro}.
        $f\in\funs{X}{\reals}$ and $j\in\funs{X}{\reals}$ by \cref{continuous}.
        Thus $f(x)\in\reals$ and $j(x)\in\reals$ by \cref{funs_type_apply}.
        $(f(x), j(x))\in\reals\times\reals$ by \cref{times_tuple_intro}.
        Thus $k(x)\in\reals\times\reals$.
        $f(x) \rmul j(x)\in\reals$ by \cref{real_mul_closed}.
        Let $p_0 = (f(x), j(x))$.
        Let $z_0 = f(x) \rmul j(x)$.
        $\fst{p_0} = f(x)$ and $\snd{p_0} = j(x)$ by \cref{fst_eq,snd_eq}.
        Thus $z_0 = \fst{p_0} \rmul \snd{p_0}$.
        Thus $k(x) = p_0$ by \cref{function_apply_intro}.
        Thus $(k(x), z_0) = (p_0, z_0)$.
        Thus $(k(x), z_0)\in\realMultiplication$ by \cref{real_multiplication_operation}.
        $\realMultiplication\in\funs{\reals\times\reals}{\reals}$ by \cref{continuous,real_multiplication_continuous}.
        Thus $\realMultiplication$ is a function by \cref{funs_is_function}.
        Thus $\apply{\realMultiplication}{k(x)} = f(x) \rmul j(x)$ by \cref{function_apply_intro}.
        Thus $\ran{k}\subseteq\reals\times\reals$ by \cref{funs_ran}.
        Thus $\dom{\realMultiplication} = \reals\times\reals$ by \cref{funs}.
        Thus $\ran{k}\subseteq\dom{\realMultiplication}$ by \cref{subseteq}.
        Thus $\realMultiplication$ is composable with $k$.
        Thus $x\in\dom{k}$ by \cref{funs}.
        Thus $m(x) = \apply{\realMultiplication}{k(x)}$ by \cref{circ_apply}.
        $g\in\funs{X}{\reals\setminus\{0\}}$ by \cref{continuous}.
        Thus $g(x)\in\reals\setminus\{0\}$ by \cref{funs_type_apply}.
        Thus $\ran{g}\subseteq\reals\setminus\{0\}$ by \cref{funs_ran}.
        $\realInverse\in\funs{\reals\setminus\{0\}}{\reals}$ by \cref{continuous,real_inverse_continuous}.
        Thus $\dom{\realInverse} = \reals\setminus\{0\}$ by \cref{funs}.
        Thus $\ran{g}\subseteq\dom{\realInverse}$ by \cref{subseteq}.
        Thus $\realInverse$ is composable with $g$.
        $g$ is a function by \cref{funs_is_function}.
        $\realInverse$ is a function by \cref{funs_is_function}.
        Thus $x\in\dom{g}$ by \cref{funs}.
        Thus $j(x) = \apply{\realInverse}{g(x)}$ by \cref{circ_apply}.
        $g(x)\notin\{0\}$ by \cref{setminus_elim_right}.
        Thus $g(x)\neq 0$ by \cref{singleton_iff}.
        $g(x)\in\reals$ by \cref{setminus_elim_left}.
        $\inv{(g(x))}\in\reals$ by \cref{real_mul_inverse}.
        $(g(x),\inv{(g(x))})\in\realInverse$ by \cref{real_inverse_operation}.
        $\realInverse$ is a function.
        Thus $\apply{\realInverse}{g(x)} = \inv{(g(x))}$ by \cref{function_apply_intro}.
        Thus $m(x) = f(x) \rmul \inv{(g(x))}$.
        $(x, f(x) \rmul \inv{(g(x))})\in h$.
        $h$ is right-unique and $h$ is a relation.
        Thus $\apply{h}{x} = f(x) \rmul \inv{(g(x))}$ by \cref{function_apply_intro}.
        Thus $m(x) = \apply{h}{x}$.
    \end{subproof}
    Show that $m\subseteq h$.
    \begin{subproof}
        $m$ is a function.
        $h$ is a function.
        $\dom{m} = X$ by \cref{funs,continuous}.
        $\dom{h} = X$ by \cref{funs}.
        Thus $\dom{m}\subseteq\dom{h}$ by \cref{subseteq}.
        Show that for all $x\in\dom{m}$ we have $m(x) = \apply{h}{x}$.
        \begin{subproof}
            Fix $x\in\dom{m}$.
            Thus $x\in X$ by \cref{funs,continuous}.
            Thus $m(x) = \apply{h}{x}$.
        \end{subproof}
        Thus $m\subseteq h$ by \cref{fun_subseteq}.
    \end{subproof}
    Show that $h\subseteq m$.
    \begin{subproof}
        $m$ is a function.
        $h$ is a function.
        $\dom{m} = X$ by \cref{funs,continuous}.
        $\dom{h} = X$ by \cref{funs}.
        Thus $\dom{h}\subseteq\dom{m}$ by \cref{subseteq}.
        Show that for all $x\in\dom{h}$ we have $m(x) = \apply{h}{x}$.
        \begin{subproof}
            Fix $x\in\dom{h}$.
            Thus $x\in X$ by \cref{funs}.
            Thus $m(x) = \apply{h}{x}$.
        \end{subproof}
        Thus $h\subseteq m$ by \cref{fun_subseteq}.
    \end{subproof}
    Thus $m = h$ by \cref{subseteq_antisymmetric}.
    Thus $h$ is continuous from $X$ to $\reals$ with respect to $T$ and $\standardTopologyR$ by \cref{continuous_composite}.
\end{proof}

% from §21 Item 16: uniform convergence
\begin{definition}\label{uniform_converges_to}
    $s$ converges uniformly to $f$ from $X$ to $Y$ with respect to $d$ iff
        $s$ is a sequence in $\funs{X}{Y}$ and
        $f\in\funs{X}{Y}$ and
        for all $\epsilon\in\reals$ such that $0 < \epsilon$
        there exists $N\in\naturalsPlus$ such that for all $n\in\naturalsPlus$ if $N\leq n$ then
            for all $x\in X$ we have $d((\apply{\apply{s}{n}}{x}, f(x))) < \epsilon$.
\end{definition}

% from §21 Item 17: uniform limit theorem
\begin{theorem}\label{uniform_limit_theorem}
    Suppose $T_X$ is a topology on $X$.
    Suppose $s$ is a sequence in $\funs{X}{Y}$.
    Suppose $f\in\funs{X}{Y}$.
    Suppose $d$ is a metric on $Y$.
    Let $T = \metricTopology{d}{Y}$.
    Suppose for all $n\in\naturalsPlus$ we have
        $s(n)$ is continuous from $X$ to $Y$ with respect to $T_X$ and $T$.
    Suppose $s$ converges uniformly to $f$ from $X$ to $Y$ with respect to $d$.
    Then $f$ is continuous from $X$ to $Y$ with respect to $T_X$ and $T$.
\end{theorem}
\begin{proof}
    $f$ is a function from $X$ to $Y$ by \cref{funs}.
    $d$ is a metric on $Y$.
    $\metricBasis{d}{Y}$ is a basis on $Y$ by \cref{metric_basis_is_basis}.
    Thus $T$ is a topology on $Y$ by \cref{basis_topology_is_topology,metric_topology,basis_for_topology}.
    Show that for all $V$ such that $V$ is open in $Y$ with respect to $T$ we have
        $\preimg{f}{V}$ is open in $X$ with respect to $T_X$.
    \begin{subproof}
        Fix $V$.
        Assume $V$ is open in $Y$ with respect to $T$.
        $V\in T$ by \cref{open_in}.
        $T\subseteq\pow{Y}$ by \cref{topology_on}.
        Thus $V\in\pow{Y}$ by \cref{subseteq}.
        Thus $V\subseteq Y$ by \cref{pow_iff}.
        Show that for all $x\in\preimg{f}{V}$ there exists $U$ such that
            $U$ is open in $X$ with respect to $T_X$ and $x\in U$ and $U\subseteq\preimg{f}{V}$.
        \begin{subproof}
            Fix $x$.
            Assume $x\in\preimg{f}{V}$.
            Take $y\in V$ such that $(x,y)\in f$ by \cref{preimg_iff}.
            $f$ is a function.
            Thus $y = f(x)$ by \cref{function_apply_iff}.
            Thus $f(x)\in V$.
            There exists $\epsilon\in\reals$ such that $0 < \epsilon$ and
                $\metricBall{d}{Y}{f(x)}{\epsilon} \subseteq V$ by \cref{metric_open_iff}.
            Take $\epsilon\in\reals$ such that $0 < \epsilon$ and
                $\metricBall{d}{Y}{f(x)}{\epsilon} \subseteq V$.
            Let $\epsilon_1 = \epsilon / (1 + 1)$.
            Show that $0 < \epsilon_1$.
            \begin{subproof}
                $0 < 1$ by \cref{real_zero_lt_one}.
                $1\in\reals$ by \cref{real_naturals_subset_reals,real_one_in_naturals,subseteq}.
                Thus $1$ is positive.
                Thus $1 + 1$ is positive by \cref{real_pos_add}.
                Thus $1 + 1\in\reals$ by \cref{real_add_closed}.
                Thus $1 + 1 \neq 0$ by \cref{real_pos_nonzero}.
                Thus $\inv{(1 + 1)}\in\reals$ by \cref{real_mul_inverse}.
                Thus $0 < \inv{(1 + 1)}$ by \cref{real_inv_pos}.
                $\epsilon\in\reals$.
                Thus $\epsilon \rmul \inv{(1 + 1)}$ is positive by \cref{real_mul_pos_pos}.
                $\epsilon / (1 + 1) = \epsilon \rmul \inv{(1 + 1)}$ by \cref{real_div}.
                Thus $0 < \epsilon_1$ by \cref{real_lt_subst_right}.
            \end{subproof}
            $\epsilon\in\reals$.
            $0 < 1$ by \cref{real_zero_lt_one}.
            $1\in\reals$ by \cref{real_naturals_subset_reals,real_one_in_naturals,subseteq}.
            Thus $1$ is positive.
            Thus $1 + 1$ is positive by \cref{real_pos_add}.
            Thus $1 + 1\in\reals$ by \cref{real_add_closed}.
            Thus $1 + 1 \neq 0$ by \cref{real_pos_nonzero}.
            Thus $\inv{(1 + 1)}\in\reals$ by \cref{real_mul_inverse}.
            Thus $\epsilon_1\in\reals$ by \cref{real_div,real_mul_closed}.
            $s$ is a sequence in $\funs{X}{Y}$.
            Thus $s\in\funs{\naturalsPlus}{\funs{X}{Y}}$ by \cref{sequence_in}.
            There exists $N\in\naturalsPlus$ such that for all $n\in\naturalsPlus$ if $N\leq n$ then
                for all $z\in X$ we have $d((\apply{\apply{s}{n}}{z}, f(z))) < \epsilon_1$
                by \cref{uniform_converges_to}.
            Take $N\in\naturalsPlus$ such that for all $n\in\naturalsPlus$ if $N\leq n$ then
                for all $z\in X$ we have $d((\apply{\apply{s}{n}}{z}, f(z))) < \epsilon_1$.
            Let $g = s(N)$.
            Thus $g\in\funs{X}{Y}$ by \cref{funs_type_apply}.
            Thus $g$ is a function by \cref{funs_is_function}.
            $f\in\funs{X}{Y}$.
            Thus $f$ is a function by \cref{funs_is_function}.
            $N\leq N$.
            Thus for all $z\in X$ we have $d((g(z), f(z))) < \epsilon_1$.
            $N\in\naturalsPlus$.
            Thus $g$ is continuous from $X$ to $Y$ with respect to $T_X$ and $T$.
            Let $W = \metricBall{d}{Y}{f(x)}{\epsilon_1}$.
            Show that $W$ is open in $Y$ with respect to $T$.
            \begin{subproof}
                Show that $W\subseteq Y$.
                \begin{subproof}
                    Show that for all $y$ such that $y\in W$ we have $y\in Y$.
                    \begin{subproof}
                        Fix $y$.
                        Assume $y\in W$.
                        Thus $y\in Y$ by \cref{metric_ball}.
                \end{subproof}
                Thus $W\subseteq Y$ by \cref{subseteq}.
            \end{subproof}
            Thus $W\in\pow{Y}$ by \cref{powerset_intro}.
            $x\in\dom{f}$ by \cref{preimg_subseteq_dom,elem_subseteq}.
            Thus $x\in X$ by \cref{funs}.
            $f(x)\in Y$ by \cref{funs_type_apply}.
            $W\in\metricBasis{d}{Y}$ by \cref{metric_basis}.
                $\metricBasis{d}{Y}$ is a basis on $Y$ by \cref{metric_basis_is_basis}.
                Thus $W\in T$ by \cref{basis_elem_open_in_generated,metric_topology,basis_for_topology}.
                Thus $W$ is open in $Y$ with respect to $T$ by \cref{open_in}.
            \end{subproof}
            Let $U = \preimg{g}{W}$.
            Thus $U$ is open in $X$ with respect to $T_X$ by \cref{continuous}.
            Show that $x\in U$.
            \begin{subproof}
                $x\in\dom{f}$ by \cref{preimg_subseteq_dom,elem_subseteq}.
                Thus $x\in X$ by \cref{funs}.
                $g(x)\in Y$ by \cref{funs_type_apply}.
                $f(x)\in Y$ by \cref{funs_type_apply}.
                $d$ is symmetric on $Y$ by \cref{metric_on}.
                Thus $d((f(x), g(x))) = d((g(x), f(x)))$ by \cref{metric_symmetric}.
                Thus $d((f(x), g(x))) < \epsilon_1$ by \cref{real_lt_subst_left}.
                Thus $g(x)\in W$ by \cref{metric_ball}.
                $x\in\dom{g}$ by \cref{funs}.
                Thus $(x,g(x))\in g$ by \cref{function_apply_iff}.
                Thus $x\in U$ by \cref{preimg_iff}.
            \end{subproof}
            Show that $U\subseteq\preimg{f}{V}$.
            \begin{subproof}
                Show that for all $y$ such that $y\in U$ we have $y\in\preimg{f}{V}$.
                \begin{subproof}
                    Fix $y$.
                    Assume $y\in U$.
                    Take $z\in W$ such that $(y,z)\in g$ by \cref{preimg_iff}.
                    Thus $z = g(y)$ by \cref{function_apply_iff}.
                    Thus $g(y)\in W$.
                    Thus $g(y)\in Y$ and $d((f(x), g(y))) < \epsilon_1$ by \cref{metric_ball}.
                    Thus $f(x)\in Y$ by \cref{subseteq}.
                    $d$ is symmetric on $Y$ by \cref{metric_on}.
                    Thus $d((g(y), f(x))) = d((f(x), g(y)))$ by \cref{metric_symmetric}.
                    Thus $d((g(y), f(x))) < \epsilon_1$ by \cref{real_lt_subst_left}.
                    $y\in\dom{g}$ by \cref{preimg_subseteq_dom,elem_subseteq}.
                    Thus $y\in X$ by \cref{funs}.
                    Thus $d((g(y), f(y))) < \epsilon_1$.
                    $f(y)\in Y$ by \cref{funs_type_apply}.
                    $d$ is symmetric on $Y$ by \cref{metric_on}.
                    Thus $d((f(y), g(y))) = d((g(y), f(y)))$ by \cref{metric_symmetric}.
                    Thus $d((f(y), g(y))) < \epsilon_1$ by \cref{real_lt_subst_left}.
                    $d\in\funs{Y\times Y}{\reals}$ by \cref{metric_on}.
                    $(f(y), g(y))\in Y\times Y$ by \cref{times_tuple_intro}.
                    $(g(y), f(x))\in Y\times Y$ by \cref{times_tuple_intro}.
                    $(f(y), f(x))\in Y\times Y$ by \cref{times_tuple_intro}.
                    Thus $d((f(y), g(y)))\in\reals$ by \cref{funs_type_apply}.
                    Thus $d((g(y), f(x)))\in\reals$ by \cref{funs_type_apply}.
                    Thus $d((f(y), f(x)))\in\reals$ by \cref{funs_type_apply}.
                    $d$ satisfies the triangle inequality on $Y$ by \cref{metric_on}.
                    Thus $d((f(y), g(y))) + d((g(y), f(x))) \geq d((f(y), f(x)))$ by \cref{metric_triangle_inequality}.
                    Thus $d((f(y), f(x))) \leq d((f(y), g(y))) + d((g(y), f(x)))$.
                    $d((f(y), g(y))) + d((g(y), f(x)))\in\reals$ by \cref{real_add_closed}.
                    $d((f(y), g(y))) < \epsilon_1$ and $d((g(y), f(x))) < \epsilon_1$.
                    Thus $d((f(y), g(y))) + d((g(y), f(x))) < \epsilon_1 + \epsilon_1$ by \cref{real_lt_add_two}.
                    Show that $\epsilon_1 + \epsilon_1 = \epsilon$.
                    \begin{subproof}
                        $\epsilon_1 = \epsilon / (1 + 1)$.
                        $\epsilon / (1 + 1) = \epsilon \rmul \inv{(1 + 1)}$ by \cref{real_div}.
                        Thus $\epsilon_1 = \epsilon \rmul \inv{(1 + 1)}$.
                        $\epsilon \rmul \inv{(1 + 1)} = \inv{(1 + 1)} \rmul \epsilon$ by \cref{real_mul_comm}.
                        Thus $\epsilon_1 = \inv{(1 + 1)} \rmul \epsilon$.
                        Thus $\epsilon_1 + \epsilon_1 = (\inv{(1 + 1)} \rmul \epsilon) + (\inv{(1 + 1)} \rmul \epsilon)$.
                        $\inv{(1 + 1)} \rmul (\epsilon + \epsilon) = (\inv{(1 + 1)} \rmul \epsilon) + (\inv{(1 + 1)} \rmul \epsilon)$ by \cref{real_distrib}.
                        Thus $\epsilon_1 + \epsilon_1 = \inv{(1 + 1)} \rmul (\epsilon + \epsilon)$.
                        Show that $\epsilon + \epsilon = (1 + 1) \rmul \epsilon$.
                        \begin{subproof}
                            $\epsilon \rmul (1 + 1) = (\epsilon \rmul 1) + (\epsilon \rmul 1)$ by \cref{real_distrib}.
                            $\epsilon \rmul 1 = 1 \rmul \epsilon$ by \cref{real_mul_comm}.
                            $1 \rmul \epsilon = \epsilon$ by \cref{real_mul_ident}.
                            Thus $\epsilon \rmul 1 = \epsilon$.
                            Thus $\epsilon \rmul (1 + 1) = \epsilon + \epsilon$.
                            $(1 + 1) \rmul \epsilon = \epsilon \rmul (1 + 1)$ by \cref{real_mul_comm}.
                            Thus $(1 + 1) \rmul \epsilon = \epsilon + \epsilon$.
                        \end{subproof}
                        Thus $\epsilon_1 + \epsilon_1 = \inv{(1 + 1)} \rmul ((1 + 1) \rmul \epsilon)$.
                        $\inv{(1 + 1)} \rmul ((1 + 1) \rmul \epsilon) = (\inv{(1 + 1)} \rmul (1 + 1)) \rmul \epsilon$ by \cref{real_mul_assoc}.
                        $(1 + 1) \rmul \inv{(1 + 1)} = 1$ by \cref{real_mul_inverse}.
                        Thus $\inv{(1 + 1)} \rmul (1 + 1) = 1$ by \cref{real_mul_comm}.
                        Thus $\epsilon_1 + \epsilon_1 = 1 \rmul \epsilon$.
                        $1 \rmul \epsilon = \epsilon$ by \cref{real_mul_ident}.
                        Thus $\epsilon_1 + \epsilon_1 = \epsilon$.
                    \end{subproof}
                    Thus $d((f(y), g(y))) + d((g(y), f(x))) < \epsilon$ by \cref{real_lt_subst_right}.
                    Show that $d((f(y), f(x))) < \epsilon$.
                    \begin{subproof}
                        $d((f(y), f(x))) < d((f(y), g(y))) + d((g(y), f(x)))$ or
                            $d((f(y), f(x))) = d((f(y), g(y))) + d((g(y), f(x)))$ by \cref{real_leq_split}.
                        \begin{byCase}
                        \caseOf{$d((f(y), f(x))) < d((f(y), g(y))) + d((g(y), f(x)))$.}
                            Thus $d((f(y), f(x))) < \epsilon$ by \cref{real_lt_trans}.
                        \caseOf{$d((f(y), f(x))) = d((f(y), g(y))) + d((g(y), f(x)))$.}
                            Thus $d((f(y), f(x))) < \epsilon$ by \cref{real_lt_subst_right}.
                        \end{byCase}
                    \end{subproof}
                    $d$ is symmetric on $Y$ by \cref{metric_on}.
                    Thus $d((f(x), f(y))) = d((f(y), f(x)))$ by \cref{metric_symmetric}.
                    Thus $d((f(x), f(y))) < \epsilon$ by \cref{real_lt_subst_left}.
                    Thus $f(y)\in\metricBall{d}{Y}{f(x)}{\epsilon}$ by \cref{metric_ball}.
                    Thus $f(y)\in V$ by \cref{subseteq}.
                    $y\in\dom{f}$ by \cref{funs}.
                    Thus $(y,f(y))\in f$ by \cref{function_apply_iff}.
                    Thus $y\in\preimg{f}{V}$ by \cref{preimg_iff}.
                \end{subproof}
                Thus $U\subseteq\preimg{f}{V}$ by \cref{subseteq}.
            \end{subproof}
            Thus there exists $U$ such that
                $U$ is open in $X$ with respect to $T_X$ and $x\in U$ and $U\subseteq\preimg{f}{V}$.
        \end{subproof}
        Let $M = \{U\in T_X \mid U\subseteq\preimg{f}{V}\}$.
        Show that $M\subseteq T_X$.
        \begin{subproof}
            Show that for all $U$ such that $U\in M$ we have $U\in T_X$.
            \begin{subproof}
                Fix $U$.
                Assume $U\in M$.
                Thus $U\in T_X$.
            \end{subproof}
            Thus $M\subseteq T_X$ by \cref{subseteq}.
        \end{subproof}
        Thus $\unions{M}\in T_X$ by \cref{topology_on}.
        Show that $\preimg{f}{V} = \unions{M}$.
        \begin{subproof}
            Show that for all $x$ such that $x\in\preimg{f}{V}$ we have $x\in\unions{M}$.
            \begin{subproof}
                Fix $x$.
                Assume $x\in\preimg{f}{V}$.
                Take $U$ such that $U$ is open in $X$ with respect to $T_X$ and $x\in U$ and $U\subseteq\preimg{f}{V}$.
                Thus $U\in T_X$ by \cref{open_in}.
                Thus $U\in M$.
                Thus $x\in\unions{M}$ by \cref{unions_iff}.
            \end{subproof}
            Thus $\preimg{f}{V}\subseteq\unions{M}$ by \cref{subseteq}.
            Show that for all $x$ such that $x\in\unions{M}$ we have $x\in\preimg{f}{V}$.
            \begin{subproof}
                Fix $x$.
                Assume $x\in\unions{M}$.
                Take $U\in M$ such that $x\in U$ by \cref{unions_iff}.
                Thus $U\subseteq\preimg{f}{V}$.
                Thus $x\in\preimg{f}{V}$ by \cref{subseteq}.
            \end{subproof}
            Thus $\unions{M}\subseteq\preimg{f}{V}$ by \cref{subseteq}.
            Thus $\preimg{f}{V} = \unions{M}$ by \cref{subseteq_antisymmetric}.
        \end{subproof}
        Thus $\preimg{f}{V}\in T_X$.
        Thus $\preimg{f}{V}$ is open in $X$ with respect to $T_X$ by \cref{open_in}.
    \end{subproof}
    Thus $f$ is continuous from $X$ to $Y$ with respect to $T_X$ and $T$ by \cref{continuous}.
\end{proof}
% from §22 Item 1: quotient map
\begin{definition}\label{quotient_map}
    $p$ is a quotient map from $X$ to $Y$ with respect to $T_X$ and $T_Y$ iff
        $p$ is a surjection from $X$ to $Y$ and
        $p$ is continuous from $X$ to $Y$ with respect to $T_X$ and $T_Y$ and
        for all $U$ such that $U\subseteq Y$ if $\preimg{p}{U}$ is open in $X$ with respect to $T_X$ then
            $U$ is open in $Y$ with respect to $T_Y$.
\end{definition}

% from §22 Item 2: quotient map closed set characterization
\begin{lemma}\label{quotient_map_closed_iff}
    Suppose $p$ is a surjection from $X$ to $Y$.
    Suppose $T_X$ is a topology on $X$.
    Suppose $T_Y$ is a topology on $Y$.
    Then $p$ is a quotient map from $X$ to $Y$ with respect to $T_X$ and $T_Y$ iff
        for all $A$ such that $A\subseteq Y$ we have
            $A$ is closed in $Y$ with respect to $T_Y$ iff
            $\preimg{p}{A}$ is closed in $X$ with respect to $T_X$.
\end{lemma}
\begin{proof}
    $p\in\surj{X}{Y}$.
    Thus $p\in\funs{X}{Y}$ by \cref{surj}.
    Thus $p$ is a function from $X$ to $Y$ by \cref{funs}.
    Show that if $p$ is a quotient map from $X$ to $Y$ with respect to $T_X$ and $T_Y$ then
        for all $A$ such that $A\subseteq Y$ we have
            $A$ is closed in $Y$ with respect to $T_Y$ iff
            $\preimg{p}{A}$ is closed in $X$ with respect to $T_X$.
    \begin{subproof}
        Assume $p$ is a quotient map from $X$ to $Y$ with respect to $T_X$ and $T_Y$.
        Thus $p$ is continuous from $X$ to $Y$ with respect to $T_X$ and $T_Y$ by \cref{quotient_map}.
        Fix $A$.
        Assume $A\subseteq Y$.
        Show that if $A$ is closed in $Y$ with respect to $T_Y$ then
            $\preimg{p}{A}$ is closed in $X$ with respect to $T_X$.
        \begin{subproof}
            Assume $A$ is closed in $Y$ with respect to $T_Y$.
            $Y\setminus A$ is open in $Y$ with respect to $T_Y$ by \cref{closed_in}.
            $Y\setminus A\subseteq Y$ by \cref{setminus_subseteq}.
            $\preimg{p}{Y\setminus A}$ is open in $X$ with respect to $T_X$ by \cref{continuous}.
            $\preimg{p}{Y\setminus A} = X\setminus \preimg{p}{A}$ by \cref{preimg_setminus_function}.
            Thus $X\setminus \preimg{p}{A}$ is open in $X$ with respect to $T_X$ by \cref{open_in}.
            $\preimg{p}{A}\subseteq \dom{p}$ by \cref{preimg_subseteq_dom}.
            Thus $\preimg{p}{A}\subseteq X$ by \cref{funs,subseteq}.
            Thus $\preimg{p}{A}$ is closed in $X$ with respect to $T_X$ by \cref{closed_in}.
        \end{subproof}
        Show that if $\preimg{p}{A}$ is closed in $X$ with respect to $T_X$ then
            $A$ is closed in $Y$ with respect to $T_Y$.
        \begin{subproof}
            Assume $\preimg{p}{A}$ is closed in $X$ with respect to $T_X$.
            $X\setminus \preimg{p}{A}$ is open in $X$ with respect to $T_X$ by \cref{closed_in}.
            $\preimg{p}{Y\setminus A} = X\setminus \preimg{p}{A}$ by \cref{preimg_setminus_function}.
            Thus $\preimg{p}{Y\setminus A}$ is open in $X$ with respect to $T_X$ by \cref{open_in}.
            $Y\setminus A\subseteq Y$ by \cref{setminus_subseteq}.
            Thus $Y\setminus A$ is open in $Y$ with respect to $T_Y$ by \cref{quotient_map}.
            Thus $A$ is closed in $Y$ with respect to $T_Y$ by \cref{closed_in}.
        \end{subproof}
        Thus $A$ is closed in $Y$ with respect to $T_Y$ iff
            $\preimg{p}{A}$ is closed in $X$ with respect to $T_X$.
    \end{subproof}
    Show that if for all $A$ such that $A\subseteq Y$ we have
            $A$ is closed in $Y$ with respect to $T_Y$ iff
            $\preimg{p}{A}$ is closed in $X$ with respect to $T_X$
        then $p$ is a quotient map from $X$ to $Y$ with respect to $T_X$ and $T_Y$.
    \begin{subproof}
        Assume that for all $A$ such that $A\subseteq Y$ we have
            $A$ is closed in $Y$ with respect to $T_Y$ iff
            $\preimg{p}{A}$ is closed in $X$ with respect to $T_X$.
        Show that for all $B$ such that $B$ is closed in $Y$ with respect to $T_Y$ we have
            $\preimg{p}{B}$ is closed in $X$ with respect to $T_X$.
        \begin{subproof}
            Fix $B$.
            Assume $B$ is closed in $Y$ with respect to $T_Y$.
            Thus $B\subseteq Y$ by \cref{closed_in}.
            Thus $\preimg{p}{B}$ is closed in $X$ with respect to $T_X$.
        \end{subproof}
        Thus $p$ is continuous from $X$ to $Y$ with respect to $T_X$ and $T_Y$
            by \cref{preimg_closed_implies_continuous}.
        Show that for all $U$ such that $U\subseteq Y$ we have
            $U$ is open in $Y$ with respect to $T_Y$ iff
            $\preimg{p}{U}$ is open in $X$ with respect to $T_X$.
        \begin{subproof}
            Fix $U$.
            Assume $U\subseteq Y$.
            Show that if $U$ is open in $Y$ with respect to $T_Y$ then
                $\preimg{p}{U}$ is open in $X$ with respect to $T_X$.
            \begin{subproof}
                Assume $U$ is open in $Y$ with respect to $T_Y$.
                $Y\setminus U\subseteq Y$ by \cref{setminus_subseteq}.
                $Y\setminus (Y\setminus U) = U$ by \cref{double_relative_complement,subseteq}.
                Thus $Y\setminus (Y\setminus U)$ is open in $Y$ with respect to $T_Y$ by \cref{open_in}.
                Thus $Y\setminus U$ is closed in $Y$ with respect to $T_Y$ by \cref{closed_in}.
                Thus $\preimg{p}{Y\setminus U}$ is closed in $X$ with respect to $T_X$.
                $\preimg{p}{Y\setminus U} = X\setminus \preimg{p}{U}$ by \cref{preimg_setminus_function}.
                Thus $X\setminus \preimg{p}{U}$ is closed in $X$ with respect to $T_X$ by \cref{closed_in}.
                $\preimg{p}{U}\subseteq \dom{p}$ by \cref{preimg_subseteq_dom}.
                Thus $\preimg{p}{U}\subseteq X$ by \cref{funs,subseteq}.
                $X\setminus (X\setminus \preimg{p}{U})$ is open in $X$ with respect to $T_X$ by \cref{closed_in}.
                $X\setminus (X\setminus \preimg{p}{U}) = \preimg{p}{U}$ by \cref{double_relative_complement,subseteq}.
                Thus $\preimg{p}{U}$ is open in $X$ with respect to $T_X$ by \cref{open_in}.
            \end{subproof}
            Show that if $\preimg{p}{U}$ is open in $X$ with respect to $T_X$ then
                $U$ is open in $Y$ with respect to $T_Y$.
            \begin{subproof}
                Assume $\preimg{p}{U}$ is open in $X$ with respect to $T_X$.
                $\preimg{p}{U}\subseteq \dom{p}$ by \cref{preimg_subseteq_dom}.
                Thus $\preimg{p}{U}\subseteq X$ by \cref{funs,subseteq}.
                $X\setminus \preimg{p}{U}\subseteq X$ by \cref{setminus_subseteq}.
                $X\setminus (X\setminus \preimg{p}{U}) = \preimg{p}{U}$ by \cref{double_relative_complement,subseteq}.
                Thus $X\setminus (X\setminus \preimg{p}{U})$ is open in $X$ with respect to $T_X$ by \cref{open_in}.
                Thus $X\setminus \preimg{p}{U}$ is closed in $X$ with respect to $T_X$ by \cref{closed_in}.
                $\preimg{p}{Y\setminus U} = X\setminus \preimg{p}{U}$ by \cref{preimg_setminus_function}.
                Thus $\preimg{p}{Y\setminus U}$ is closed in $X$ with respect to $T_X$ by \cref{closed_in}.
                $Y\setminus U\subseteq Y$ by \cref{setminus_subseteq}.
                Thus $Y\setminus U$ is closed in $Y$ with respect to $T_Y$.
                $Y\setminus (Y\setminus U)$ is open in $Y$ with respect to $T_Y$ by \cref{closed_in}.
                $Y\setminus (Y\setminus U) = U$ by \cref{double_relative_complement,subseteq}.
                Thus $U$ is open in $Y$ with respect to $T_Y$ by \cref{open_in}.
            \end{subproof}
            Thus $U$ is open in $Y$ with respect to $T_Y$ iff
                $\preimg{p}{U}$ is open in $X$ with respect to $T_X$.
        \end{subproof}
        Thus $p$ is a quotient map from $X$ to $Y$ with respect to $T_X$ and $T_Y$ by \cref{quotient_map}.
    \end{subproof}
\end{proof}

% from §22 Item 3: saturated subset
\begin{definition}\label{saturated}
    $C$ is saturated with respect to $p$ iff
        $C\subseteq\dom{p}$ and
        for all $y\in\ran{p}$ if $C$ intersects $\preimg{p}{\{y\}}$ then
            $\preimg{p}{\{y\}}\subseteq C$.
\end{definition}

% from §22 Item 4: saturated sets are complete inverse images
\begin{lemma}\label{saturated_iff_preimg}
    Suppose $p$ is a function.
    Then $C$ is saturated with respect to $p$ iff there exists $B$ such that
        $C = \preimg{p}{B}$.
\end{lemma}
\begin{proof}
    Show that if $C$ is saturated with respect to $p$ then there exists $B$ such that
        $C = \preimg{p}{B}$.
    \begin{subproof}
        Assume $C$ is saturated with respect to $p$.
        Let $B = \img{p}{C}$.
        Show that $C\subseteq \preimg{p}{B}$.
        \begin{subproof}
            Show that for all $x$ such that $x\in C$ we have $x\in\preimg{p}{B}$.
            \begin{subproof}
                Fix $x$.
                Assume $x\in C$.
                $C\subseteq \dom{p}$ by \cref{saturated}.
                Thus $x\in\dom{p}$ by \cref{subseteq}.
                Thus $p(x)\in \img{p}{C}$ by \cref{img_of_function_intro,inter_intro}.
                Thus $p(x)\in B$.
                Thus $(x,p(x))\in p$ by \cref{function_apply_elim}.
                Thus $x\in\preimg{p}{B}$ by \cref{preimg_iff}.
            \end{subproof}
            Thus $C\subseteq \preimg{p}{B}$ by \cref{subseteq}.
        \end{subproof}
        Show that $\preimg{p}{B}\subseteq C$.
        \begin{subproof}
            Show that for all $x$ such that $x\in\preimg{p}{B}$ we have $x\in C$.
            \begin{subproof}
                Fix $x$.
                Assume $x\in\preimg{p}{B}$.
                Take $y\in B$ such that $x\mathrel{p} y$ by \cref{preimg_iff}.
                Take $c\in\dom{p}\inter C$ such that $y = p(c)$ by \cref{img_of_function_elim}.
                Thus $c\in C$ by \cref{inter_elim_right}.
                $c\in\dom{p}$ by \cref{inter_elim_left}.
                Thus $p(c) = y$.
                Thus $c\mathrel{p} y$ by \cref{function_apply_iff}.
                Thus $c\in\preimg{p}{\{y\}}$ by \cref{preimg_iff,singleton_iff}.
                Thus $c\in C\inter \preimg{p}{\{y\}}$ by \cref{inter_intro}.
                Show that $C\inter \preimg{p}{\{y\}}\neq \emptyset$.
                \begin{subproof}
                    Suppose not.
                    Thus $C\inter \preimg{p}{\{y\}} = \emptyset$.
                    Thus $c\notin C\inter \preimg{p}{\{y\}}$ by \cref{emptyset}.
                    Contradiction.
                \end{subproof}
                Thus $C$ intersects $\preimg{p}{\{y\}}$ by \cref{intersects}.
                Thus $y\in \img{p}{C}$.
                $\img{p}{C}\subseteq \ran{p}$ by \cref{img_subseteq_ran}.
                Thus $y\in \ran{p}$ by \cref{subseteq}.
                Thus $\preimg{p}{\{y\}}\subseteq C$ by \cref{saturated}.
                Thus $x\in\preimg{p}{\{y\}}$ by \cref{preimg_iff,singleton_iff}.
                Thus $x\in C$ by \cref{subseteq}.
            \end{subproof}
            Thus $\preimg{p}{B}\subseteq C$ by \cref{subseteq}.
        \end{subproof}
        Thus $C = \preimg{p}{B}$ by \cref{subseteq_antisymmetric}.
    \end{subproof}
    Show that if there exists $B$ such that $C = \preimg{p}{B}$ then
        $C$ is saturated with respect to $p$.
    \begin{subproof}
        Assume that there exists $B$ such that $C = \preimg{p}{B}$.
        Take $B$ such that $C = \preimg{p}{B}$.
        Show that $C\subseteq\dom{p}$.
        \begin{subproof}
            Show that for all $x$ such that $x\in C$ we have $x\in\dom{p}$.
            \begin{subproof}
                Fix $x$.
                Assume $x\in C$.
                Thus $x\in\preimg{p}{B}$.
                Take $y\in B$ such that $x\mathrel{p} y$ by \cref{preimg_iff}.
                Thus $x\in\dom{p}$ by \cref{function_apply_iff}.
            \end{subproof}
            Thus $C\subseteq\dom{p}$ by \cref{subseteq}.
        \end{subproof}
        Show that for all $y\in\ran{p}$ if $C$ intersects $\preimg{p}{\{y\}}$ then
            $\preimg{p}{\{y\}}\subseteq C$.
        \begin{subproof}
            Fix $y\in\ran{p}$.
            Assume $C$ intersects $\preimg{p}{\{y\}}$.
            Take $c$ such that $c\in C\inter \preimg{p}{\{y\}}$ by \cref{intersects,nonempty_inhabited}.
            Thus $c\in C$ and $c\in\preimg{p}{\{y\}}$ by \cref{inter_elim_left,inter_elim_right}.
            Take $b\in B$ such that $c\mathrel{p} b$ by \cref{preimg_iff,subseteq}.
            Thus $c\mathrel{p} y$ by \cref{preimg_iff,singleton_iff,inter_elim_right}.
            $p$ is right-unique by \cref{function_rightunique}.
            Thus $b = y$ by \hyperref[rightunique]{right-uniqueness}.
            Thus $y\in B$.
            Show that for all $x$ such that $x\in\preimg{p}{\{y\}}$ we have $x\in C$.
            \begin{subproof}
                Fix $x$.
                Assume $x\in\preimg{p}{\{y\}}$.
                Thus $x\mathrel{p} y$ by \cref{preimg_iff,singleton_iff}.
                Thus $x\in\preimg{p}{B}$ by \cref{preimg_iff}.
                Thus $x\in C$ by \cref{subseteq}.
            \end{subproof}
            Thus $\preimg{p}{\{y\}}\subseteq C$ by \cref{subseteq}.
        \end{subproof}
        Thus $C$ is saturated with respect to $p$ by \cref{saturated}.
    \end{subproof}
\end{proof}

% from §22 helper lemma 1: image of a preimage for surjections
\begin{lemma}\label{img_preimg_surjection}
    Suppose $p$ is a surjection from $X$ to $Y$.
    Suppose $B\subseteq Y$.
    Then $\img{p}{\preimg{p}{B}} = B$.
\end{lemma}
\begin{proof}
    $p\in\surj{X}{Y}$.
    Thus $p\in\funs{X}{Y}$ by \cref{surj}.
    Thus $p$ is a function by \cref{funs_is_function}.
    Show that $\img{p}{\preimg{p}{B}}\subseteq B$.
    \begin{subproof}
        Thus $\img{p}{\preimg{p}{B}}\subseteq B$ by \cref{img_preimg_subseteq}.
    \end{subproof}
    Show that $B\subseteq \img{p}{\preimg{p}{B}}$.
    \begin{subproof}
        Show that for all $y$ such that $y\in B$ we have $y\in\img{p}{\preimg{p}{B}}$.
        \begin{subproof}
            Fix $y$.
            Assume $y\in B$.
            Thus $y\in Y$ by \cref{subseteq}.
            Take $x\in X$ such that $p(x) = y$ by \cref{surj}.
            Thus $x\in\dom{p}$ by \cref{funs}.
            Thus $x\mathrel{p} y$ by \cref{function_apply_iff}.
            Thus $x\in\preimg{p}{B}$ by \cref{preimg_iff}.
            Thus $y\in\img{p}{\preimg{p}{B}}$ by \cref{img_iff}.
        \end{subproof}
        Thus $B\subseteq \img{p}{\preimg{p}{B}}$ by \cref{subseteq}.
    \end{subproof}
    Thus $\img{p}{\preimg{p}{B}} = B$ by \cref{subseteq_antisymmetric}.
\end{proof}

% from §22 Item 5: quotient maps and saturated open sets
\begin{theorem}\label{quotient_map_saturated_open_iff}
    Suppose $p$ is a surjection from $X$ to $Y$.
    Suppose $T_X$ is a topology on $X$.
    Suppose $T_Y$ is a topology on $Y$.
    Then $p$ is a quotient map from $X$ to $Y$ with respect to $T_X$ and $T_Y$ iff
        $p$ is continuous from $X$ to $Y$ with respect to $T_X$ and $T_Y$ and
        for all $C$ such that
            $C$ is saturated with respect to $p$ and
            $C$ is open in $X$ with respect to $T_X$
        we have $\img{p}{C}$ is open in $Y$ with respect to $T_Y$.
\end{theorem}
\begin{proof}
    Show that if $p$ is a quotient map from $X$ to $Y$ with respect to $T_X$ and $T_Y$ then
        $p$ is continuous from $X$ to $Y$ with respect to $T_X$ and $T_Y$ and
        for all $C$ such that
            $C$ is saturated with respect to $p$ and
            $C$ is open in $X$ with respect to $T_X$
        we have $\img{p}{C}$ is open in $Y$ with respect to $T_Y$.
    \begin{subproof}
        Assume $p$ is a quotient map from $X$ to $Y$ with respect to $T_X$ and $T_Y$.
        $p$ is continuous from $X$ to $Y$ with respect to $T_X$ and $T_Y$ by \cref{continuous,quotient_map}.
        Show that for all $C$ such that
            $C$ is saturated with respect to $p$ and
            $C$ is open in $X$ with respect to $T_X$
        we have $\img{p}{C}$ is open in $Y$ with respect to $T_Y$.
        \begin{subproof}
            Fix $C$.
            Assume $C$ is saturated with respect to $p$.
            Assume $C$ is open in $X$ with respect to $T_X$.
            $p\in\surj{X}{Y}$.
            Thus $p\in\funs{X}{Y}$ by \cref{surj}.
            Thus $p$ is a function by \cref{funs_is_function}.
            Take $B$ such that $C = \preimg{p}{B}$ by \cref{saturated_iff_preimg}.
            Let $B_1 = B\inter Y$.
            Show that $\preimg{p}{B_1} = C$.
            \begin{subproof}
                $B_1\subseteq B$ by \cref{inter_lower_left}.
                Thus $\preimg{p}{B_1}\subseteq \preimg{p}{B}$ by \cref{preimg_subseteq}.
                Show that $\preimg{p}{B}\subseteq \preimg{p}{B_1}$.
                \begin{subproof}
                    Show that for all $x$ such that $x\in\preimg{p}{B}$ we have $x\in\preimg{p}{B_1}$.
                    \begin{subproof}
                        Fix $x$.
                        Assume $x\in\preimg{p}{B}$.
                        Take $y\in B$ such that $x\mathrel{p} y$ by \cref{preimg_iff}.
                        $p\in\surj{X}{Y}$.
                        Thus $p\in\funs{X}{Y}$ by \cref{surj}.
                        Thus $\ran{p}\subseteq Y$ by \cref{funs_ran}.
                        Thus $y\in Y$ by \cref{ran_iff,subseteq}.
                        Thus $y\in B_1$ by \cref{inter_intro}.
                        Thus $x\in\preimg{p}{B_1}$ by \cref{preimg_iff}.
                    \end{subproof}
                    Thus $\preimg{p}{B}\subseteq \preimg{p}{B_1}$ by \cref{subseteq}.
                \end{subproof}
                Thus $\preimg{p}{B_1} = \preimg{p}{B}$ by \cref{subseteq_antisymmetric}.
                Thus $\preimg{p}{B_1} = C$ by \cref{subseteq}.
            \end{subproof}
            Thus $\preimg{p}{B_1}$ is open in $X$ with respect to $T_X$ by \cref{open_in}.
            Thus $B_1$ is open in $Y$ with respect to $T_Y$ by \cref{quotient_map,inter_lower_right}.
            $\img{p}{C} = \img{p}{\preimg{p}{B_1}}$ by \cref{subseteq}.
            Thus $\img{p}{C} = B_1$ by \cref{img_preimg_surjection,inter_lower_right}.
            Thus $\img{p}{C}$ is open in $Y$ with respect to $T_Y$ by \cref{open_in}.
        \end{subproof}
    \end{subproof}
    Show that if
        $p$ is continuous from $X$ to $Y$ with respect to $T_X$ and $T_Y$ and
        for all $C$ such that
            $C$ is saturated with respect to $p$ and
            $C$ is open in $X$ with respect to $T_X$
        we have $\img{p}{C}$ is open in $Y$ with respect to $T_Y$
    then $p$ is a quotient map from $X$ to $Y$ with respect to $T_X$ and $T_Y$.
    \begin{subproof}
        Assume $p$ is continuous from $X$ to $Y$ with respect to $T_X$ and $T_Y$.
        Assume that for all $C$ such that
            $C$ is saturated with respect to $p$ and
            $C$ is open in $X$ with respect to $T_X$
        we have $\img{p}{C}$ is open in $Y$ with respect to $T_Y$.
        Show that for all $U$ such that $U\subseteq Y$ we have
            $U$ is open in $Y$ with respect to $T_Y$ iff
            $\preimg{p}{U}$ is open in $X$ with respect to $T_X$.
        \begin{subproof}
            Fix $U$.
            Assume $U\subseteq Y$.
            Show that if $U$ is open in $Y$ with respect to $T_Y$ then
                $\preimg{p}{U}$ is open in $X$ with respect to $T_X$.
            \begin{subproof}
                Assume $U$ is open in $Y$ with respect to $T_Y$.
                Thus $\preimg{p}{U}$ is open in $X$ with respect to $T_X$ by \cref{continuous}.
            \end{subproof}
            Show that if $\preimg{p}{U}$ is open in $X$ with respect to $T_X$ then
                $U$ is open in $Y$ with respect to $T_Y$.
            \begin{subproof}
                Assume $\preimg{p}{U}$ is open in $X$ with respect to $T_X$.
                Thus $p\in\funs{X}{Y}$ by \cref{continuous}.
                Thus $p$ is a function by \cref{funs_is_function}.
                $\preimg{p}{U}$ is saturated with respect to $p$ by \cref{saturated_iff_preimg}.
                Thus $\img{p}{\preimg{p}{U}}$ is open in $Y$ with respect to $T_Y$.
                Thus $U$ is open in $Y$ with respect to $T_Y$ by \cref{img_preimg_surjection}.
            \end{subproof}
            Thus $U$ is open in $Y$ with respect to $T_Y$ iff
                $\preimg{p}{U}$ is open in $X$ with respect to $T_X$.
        \end{subproof}
        Thus $p$ is a quotient map from $X$ to $Y$ with respect to $T_X$ and $T_Y$ by \cref{quotient_map}.
    \end{subproof}
\end{proof}

% from §22 Item 6: open map
\begin{definition}\label{open_map}
    $f$ is an open map from $X$ to $Y$ with respect to $T$ and $T_1$ iff
        $f$ is a function from $X$ to $Y$ and
        for all $U$ such that $U$ is open in $X$ with respect to $T$ we have
            $\img{f}{U}$ is open in $Y$ with respect to $T_1$.
\end{definition}

% from §22 Item 7: closed map
\begin{definition}\label{closed_map}
    $f$ is a closed map from $X$ to $Y$ with respect to $T$ and $T_1$ iff
        $f$ is a function from $X$ to $Y$ and
        for all $A$ such that $A$ is closed in $X$ with respect to $T$ we have
            $\img{f}{A}$ is closed in $Y$ with respect to $T_1$.
\end{definition}

% from §22 Item 8: open or closed surjections are quotient maps
\begin{lemma}\label{open_or_closed_surjection_quotient}
    Suppose $p$ is a surjection from $X$ to $Y$.
    Suppose $T_X$ is a topology on $X$.
    Suppose $T_Y$ is a topology on $Y$.
    Suppose $p$ is continuous from $X$ to $Y$ with respect to $T_X$ and $T_Y$.
    Suppose $p$ is an open map from $X$ to $Y$ with respect to $T_X$ and $T_Y$
        or $p$ is a closed map from $X$ to $Y$ with respect to $T_X$ and $T_Y$.
    Then $p$ is a quotient map from $X$ to $Y$ with respect to $T_X$ and $T_Y$.
\end{lemma}
\begin{proof}
    \begin{byCase}
    \caseOf{$p$ is an open map from $X$ to $Y$ with respect to $T_X$ and $T_Y$.}
        Show that for all $U$ such that $U\subseteq Y$ we have
            $U$ is open in $Y$ with respect to $T_Y$ iff
            $\preimg{p}{U}$ is open in $X$ with respect to $T_X$.
        \begin{subproof}
            Fix $U$.
            Assume $U\subseteq Y$.
            Show that if $U$ is open in $Y$ with respect to $T_Y$ then
                $\preimg{p}{U}$ is open in $X$ with respect to $T_X$.
            \begin{subproof}
                Assume $U$ is open in $Y$ with respect to $T_Y$.
                Thus $\preimg{p}{U}$ is open in $X$ with respect to $T_X$ by \cref{continuous}.
            \end{subproof}
            Show that if $\preimg{p}{U}$ is open in $X$ with respect to $T_X$ then
                $U$ is open in $Y$ with respect to $T_Y$.
            \begin{subproof}
                Assume $\preimg{p}{U}$ is open in $X$ with respect to $T_X$.
                $\img{p}{\preimg{p}{U}}$ is open in $Y$ with respect to $T_Y$ by \cref{open_map}.
                Thus $U$ is open in $Y$ with respect to $T_Y$ by \cref{img_preimg_surjection}.
            \end{subproof}
            Thus $U$ is open in $Y$ with respect to $T_Y$ iff
                $\preimg{p}{U}$ is open in $X$ with respect to $T_X$.
        \end{subproof}
        Thus $p$ is a quotient map from $X$ to $Y$ with respect to $T_X$ and $T_Y$ by \cref{quotient_map}.
    \caseOf{$p$ is a closed map from $X$ to $Y$ with respect to $T_X$ and $T_Y$.}
        Show that for all $A$ such that $A\subseteq Y$ we have
            $A$ is closed in $Y$ with respect to $T_Y$ iff
            $\preimg{p}{A}$ is closed in $X$ with respect to $T_X$.
        \begin{subproof}
            Fix $A$.
            Assume $A\subseteq Y$.
            Show that if $A$ is closed in $Y$ with respect to $T_Y$ then
                $\preimg{p}{A}$ is closed in $X$ with respect to $T_X$.
            \begin{subproof}
                Assume $A$ is closed in $Y$ with respect to $T_Y$.
                $Y\setminus A$ is open in $Y$ with respect to $T_Y$ by \cref{closed_in}.
                Thus $\preimg{p}{Y\setminus A}$ is open in $X$ with respect to $T_X$ by \cref{continuous}.
                Thus $p$ is a function from $X$ to $Y$ by \cref{continuous}.
                $\preimg{p}{Y\setminus A} = X\setminus \preimg{p}{A}$ by \cref{preimg_setminus_function}.
                $\preimg{p}{A}\subseteq \dom{p}$ by \cref{preimg_subseteq_dom}.
                Thus $\preimg{p}{A}\subseteq X$ by \cref{funs,subseteq}.
                Thus $\preimg{p}{A}$ is closed in $X$ with respect to $T_X$ by \cref{closed_in}.
            \end{subproof}
            Show that if $\preimg{p}{A}$ is closed in $X$ with respect to $T_X$ then
                $A$ is closed in $Y$ with respect to $T_Y$.
            \begin{subproof}
                Assume $\preimg{p}{A}$ is closed in $X$ with respect to $T_X$.
                $\img{p}{\preimg{p}{A}}$ is closed in $Y$ with respect to $T_Y$ by \cref{closed_map}.
                Thus $A$ is closed in $Y$ with respect to $T_Y$ by \cref{img_preimg_surjection}.
            \end{subproof}
            Thus $A$ is closed in $Y$ with respect to $T_Y$ iff
                $\preimg{p}{A}$ is closed in $X$ with respect to $T_X$.
        \end{subproof}
        Thus $p$ is a quotient map from $X$ to $Y$ with respect to $T_X$ and $T_Y$
            by \cref{quotient_map_closed_iff}.
    \end{byCase}
\end{proof}

% from §22 Item 9: quotient topology
\begin{definition}\label{quotient_topology}
    $\quotientTopology{p}{T}{X}{A}
        = \{U\in\pow{A} \mid \text{$\preimg{p}{U}$ is open in $X$ with respect to $T$}\}$.
\end{definition}

% from §22 Item 10: quotient topology is a topology
\begin{lemma}\label{quotient_topology_is_topology}
    Assume $p$ is a function from $X$ to $A$.
    Assume $T$ is a topology on $X$.
    Then $\quotientTopology{p}{T}{X}{A}$ is a topology on $A$.
\end{lemma}
\begin{proof}
    Let $T_Q = \quotientTopology{p}{T}{X}{A}$.
    $p\in\funs{X}{A}$ by \cref{funs}.
    Thus $p$ is a function by \cref{funs_is_function}.
    $\dom{p} = X$ by \cref{funs}.
    Show that $T_Q\subseteq\pow{A}$.
    \begin{subproof}
        Show that for all $U$ such that $U\in T_Q$ we have $U\in\pow{A}$.
        \begin{subproof}
            Fix $U$.
            Assume $U\in T_Q$.
            Thus $U\in\pow{A}$ by \cref{quotient_topology}.
        \end{subproof}
        Thus $T_Q\subseteq\pow{A}$ by \cref{subseteq}.
    \end{subproof}
    Show that $\emptyset\in T_Q$.
    \begin{subproof}
        Show that $\preimg{p}{\emptyset} = \emptyset$.
        \begin{subproof}
            Show that $\preimg{p}{\emptyset}\subseteq\emptyset$.
            \begin{subproof}
                Show that for all $x$ such that $x\in\preimg{p}{\emptyset}$ we have $x\in\emptyset$.
                \begin{subproof}
                Fix $x$.
                Assume $x\in\preimg{p}{\emptyset}$.
                Suppose not.
                Thus $x\notin\emptyset$ by \cref{emptyset}.
                Take $y\in\emptyset$ such that $x\mathrel{p} y$ by \cref{preimg_iff}.
                Contradiction by \cref{emptyset}.
                \end{subproof}
                Thus $\preimg{p}{\emptyset}\subseteq\emptyset$ by \cref{subseteq}.
            \end{subproof}
            $\emptyset\subseteq\preimg{p}{\emptyset}$ by \cref{emptyset_subseteq}.
            Thus $\preimg{p}{\emptyset} = \emptyset$ by \cref{subseteq_antisymmetric}.
        \end{subproof}
        $\emptyset\in T$ by \cref{topology_on}.
        Thus $\emptyset$ is open in $X$ with respect to $T$ by \cref{open_in}.
        Thus $\preimg{p}{\emptyset}$ is open in $X$ with respect to $T$ by \cref{open_in}.
        $\emptyset\subseteq A$ by \cref{emptyset_subseteq}.
        Thus $\emptyset\in\pow{A}$ by \cref{pow_iff}.
        Thus $\emptyset\in T_Q$ by \cref{quotient_topology}.
    \end{subproof}
    Show that $A\in T_Q$.
    \begin{subproof}
        Show that $\preimg{p}{A} = X$.
        \begin{subproof}
            $\preimg{p}{A}\subseteq\dom{p}$ by \cref{preimg_subseteq_dom}.
            $\dom{p} = X$ by \cref{funs}.
            Thus $\preimg{p}{A}\subseteq X$ by \cref{subseteq}.
            Show that $X\subseteq \preimg{p}{A}$.
            \begin{subproof}
                Show that for all $x$ such that $x\in X$ we have $x\in\preimg{p}{A}$.
                \begin{subproof}
                    Fix $x$.
                    Assume $x\in X$.
                    $x\in\dom{p}$ by \cref{funs}.
                    Take $y$ such that $x\mathrel{p} y$ by \cref{dom_iff}.
                    $y\in\ran{p}$ by \cref{ran_iff}.
                    $\ran{p}\subseteq A$ by \cref{funs_ran}.
                    Thus $y\in A$ by \cref{subseteq}.
                    Thus $x\in\preimg{p}{A}$ by \cref{preimg_iff}.
                \end{subproof}
                Thus $X\subseteq \preimg{p}{A}$ by \cref{subseteq}.
            \end{subproof}
            Thus $\preimg{p}{A} = X$ by \cref{subseteq_antisymmetric}.
        \end{subproof}
        $X\in T$ by \cref{topology_on}.
        Thus $\preimg{p}{A}$ is open in $X$ with respect to $T$ by \cref{open_in}.
        $A\subseteq A$ by \cref{subseteq_refl}.
        Thus $A\in\pow{A}$ by \cref{pow_iff}.
        Thus $A\in T_Q$ by \cref{quotient_topology}.
    \end{subproof}
    Show that for all $M$ such that $M\subseteq T_Q$ we have $\unions{M}\in T_Q$.
    \begin{subproof}
        Fix $M$.
        Assume $M\subseteq T_Q$.
        Let $N = \{V\in T \mid \text{there exists $U\in M$ such that $V = \preimg{p}{U}$}\}$.
        Show that $N\subseteq T$.
        \begin{subproof}
            Show that for all $V$ such that $V\in N$ we have $V\in T$.
            \begin{subproof}
                Fix $V$.
                Assume $V\in N$.
                Thus $V\in T$.
            \end{subproof}
            Thus $N\subseteq T$ by \cref{subseteq}.
        \end{subproof}
        Thus $\unions{N}\in T$ by \cref{topology_on}.
        Show that $\preimg{p}{\unions{M}} = \unions{N}$.
        \begin{subproof}
            Show that for all $x$ such that $x\in\preimg{p}{\unions{M}}$ we have $x\in\unions{N}$.
            \begin{subproof}
                Fix $x$.
                Assume $x\in\preimg{p}{\unions{M}}$.
                Take $y\in\unions{M}$ such that $x\mathrel{p} y$ by \cref{preimg_iff}.
                Take $U\in M$ such that $y\in U$ by \cref{unions_iff}.
                Thus $U\in T_Q$ by \cref{subseteq}.
                Thus $\preimg{p}{U}$ is open in $\dom{p}$ with respect to $T$ by \cref{quotient_topology}.
                Thus $\preimg{p}{U}\in T$ by \cref{open_in}.
                Thus $\preimg{p}{U}\in N$.
                Thus $x\in\preimg{p}{U}$ by \cref{preimg_iff}.
                Thus $x\in\unions{N}$ by \cref{unions_iff}.
            \end{subproof}
            Thus $\preimg{p}{\unions{M}}\subseteq \unions{N}$ by \cref{subseteq}.
            Show that for all $x$ such that $x\in\unions{N}$ we have $x\in\preimg{p}{\unions{M}}$.
            \begin{subproof}
                Fix $x$.
                Assume $x\in\unions{N}$.
                Take $V\in N$ such that $x\in V$ by \cref{unions_iff}.
                Take $U\in M$ such that $V = \preimg{p}{U}$.
                Thus $x\in\preimg{p}{U}$ by \cref{subseteq}.
                Take $y\in U$ such that $x\mathrel{p} y$ by \cref{preimg_iff}.
                Thus $y\in\unions{M}$ by \cref{unions_iff}.
                Thus $x\in\preimg{p}{\unions{M}}$ by \cref{preimg_iff}.
            \end{subproof}
            Thus $\unions{N}\subseteq \preimg{p}{\unions{M}}$ by \cref{subseteq}.
            Thus $\preimg{p}{\unions{M}} = \unions{N}$ by \cref{subseteq_antisymmetric}.
        \end{subproof}
        Thus $\preimg{p}{\unions{M}}$ is open in $\dom{p}$ with respect to $T$ by \cref{open_in}.
        Show that $\unions{M}\in\pow{A}$.
        \begin{subproof}
            Show that $M\subseteq \pow{A}$.
            \begin{subproof}
                Show that for all $U$ such that $U\in M$ we have $U\in\pow{A}$.
                \begin{subproof}
                    Fix $U$.
                    Assume $U\in M$.
                    Thus $U\in T_Q$ by \cref{subseteq}.
                    Thus $U\in\pow{A}$ by \cref{quotient_topology}.
                \end{subproof}
                Thus $M\subseteq \pow{A}$ by \cref{subseteq}.
            \end{subproof}
            Thus $\unions{M}\subseteq A$ by \cref{unions_subseteq_of_powerset_is_subseteq}.
            Thus $\unions{M}\in\pow{A}$ by \cref{pow_iff}.
        \end{subproof}
        Thus $\unions{M}\in T_Q$ by \cref{quotient_topology}.
    \end{subproof}
    Show that for all $U,V$ such that $U\in T_Q$ and $V\in T_Q$ we have $U\inter V\in T_Q$.
    \begin{subproof}
        Fix $U$.
        Fix $V$.
        Assume $U\in T_Q$ and $V\in T_Q$.
        $\preimg{p}{U}$ is open in $X$ with respect to $T$ by \cref{quotient_topology}.
        Thus $\preimg{p}{U}\in T$ by \cref{open_in}.
        $\preimg{p}{V}$ is open in $X$ with respect to $T$ by \cref{quotient_topology}.
        Thus $\preimg{p}{V}\in T$ by \cref{open_in}.
        Thus $\preimg{p}{U}\inter \preimg{p}{V}\in T$ by \cref{topology_on}.
        $\preimg{p}{U\inter V} = \preimg{p}{U}\inter \preimg{p}{V}$ by \cref{preimg_inter_function}.
        Thus $\preimg{p}{U\inter V}\in T$ by \cref{preimg_inter_function}.
        Thus $\preimg{p}{U\inter V}$ is open in $X$ with respect to $T$ by \cref{open_in}.
        $U\in\pow{A}$ by \cref{quotient_topology}.
        $V\in\pow{A}$ by \cref{quotient_topology}.
        Thus $U\subseteq A$ and $V\subseteq A$ by \cref{pow_iff}.
        Thus $U\inter V\subseteq A$ by \cref{inter_subseteq}.
        Thus $U\inter V\in\pow{A}$ by \cref{pow_iff}.
        Thus $U\inter V\in T_Q$ by \cref{quotient_topology}.
    \end{subproof}
    Thus $T_Q$ is closed under arbitrary unions.
    Thus $T_Q$ is closed under binary intersections.
    Thus $\quotientTopology{p}{T}{X}{A}$ is a topology on $A$ by \cref{topology_on}.
\end{proof}

% from §22 Item 11: quotient topology induces a quotient map
\begin{lemma}\label{quotient_topology_quotient_map}
    Suppose $p$ is a surjection from $X$ to $A$.
    Suppose $T$ is a topology on $X$.
    Let $T_Q = \quotientTopology{p}{T}{X}{A}$.
    Then $p$ is a quotient map from $X$ to $A$ with respect to $T$ and $T_Q$.
\end{lemma}
\begin{proof}
    $p\in\funs{X}{A}$ by \cref{surj}.
    Thus $p$ is a function by \cref{funs_is_function}.
    Thus $T_Q$ is a topology on $A$ by \cref{quotient_topology_is_topology}.
    Show that $p$ is continuous from $X$ to $A$ with respect to $T$ and $T_Q$.
    \begin{subproof}
        Show that for all $V$ such that $V$ is open in $A$ with respect to $T_Q$
            we have $\preimg{p}{V}$ is open in $X$ with respect to $T$.
        \begin{subproof}
            Fix $V$.
            Assume $V$ is open in $A$ with respect to $T_Q$.
            Thus $V\in T_Q$ by \cref{open_in}.
            Thus $\preimg{p}{V}$ is open in $X$ with respect to $T$ by \cref{quotient_topology}.
        \end{subproof}
        Thus $p$ is continuous from $X$ to $A$ with respect to $T$ and $T_Q$ by \cref{continuous}.
    \end{subproof}
    Show that for all $U$ such that $U\subseteq A$ if $\preimg{p}{U}$ is open in $X$ with respect to $T$
        then $U$ is open in $A$ with respect to $T_Q$.
    \begin{subproof}
        Fix $U$.
        Assume $U\subseteq A$.
        Assume $\preimg{p}{U}$ is open in $X$ with respect to $T$.
        Thus $U\in\pow{A}$ by \cref{pow_iff}.
        Thus $U\in T_Q$ by \cref{quotient_topology}.
        Thus $U$ is open in $A$ with respect to $T_Q$ by \cref{open_in}.
    \end{subproof}
    Thus $p$ is a quotient map from $X$ to $A$ with respect to $T$ and $T_Q$
        by \cref{quotient_map}.
\end{proof}

% from §22 Item 12: quotient topology is unique
\begin{lemma}\label{quotient_topology_unique}
    Suppose $p$ is a surjection from $X$ to $A$.
    Suppose $T$ is a topology on $X$.
    Suppose $T_A$ is a topology on $A$.
    Suppose $p$ is a quotient map from $X$ to $A$ with respect to $T$ and $T_A$.
    Then $T_A = \quotientTopology{p}{T}{X}{A}$.
\end{lemma}
\begin{proof}
    Let $T_Q = \quotientTopology{p}{T}{X}{A}$.
    Show that $T_A\subseteq T_Q$.
    \begin{subproof}
        Show that for all $U$ such that $U\in T_A$ we have $U\in T_Q$.
        \begin{subproof}
            Fix $U$.
            Assume $U\in T_A$.
            Thus $U$ is open in $A$ with respect to $T_A$ by \cref{open_in}.
            $T_A\subseteq \pow{A}$ by \cref{topology_on}.
            Thus $U\in\pow{A}$ by \cref{subseteq}.
            Thus $\preimg{p}{U}$ is open in $X$ with respect to $T$ by \cref{continuous,quotient_map}.
            Thus $U\in T_Q$ by \cref{quotient_topology}.
        \end{subproof}
        Thus $T_A\subseteq T_Q$ by \cref{subseteq}.
    \end{subproof}
    Show that $T_Q\subseteq T_A$.
    \begin{subproof}
        Show that for all $U$ such that $U\in T_Q$ we have $U\in T_A$.
        \begin{subproof}
            Fix $U$.
            Assume $U\in T_Q$.
            Thus $U\in\pow{A}$ by \cref{quotient_topology}.
            Thus $U\subseteq A$ by \cref{pow_iff}.
            Thus $\preimg{p}{U}$ is open in $X$ with respect to $T$ by \cref{quotient_topology}.
            Thus $U$ is open in $A$ with respect to $T_A$ by \cref{quotient_map}.
            Thus $U\in T_A$ by \cref{open_in}.
        \end{subproof}
        Thus $T_Q\subseteq T_A$ by \cref{subseteq}.
    \end{subproof}
    Thus $T_A = T_Q$ by \cref{subseteq_antisymmetric}.
\end{proof}

% from §22 helper definition 1: partition map
\begin{definition}\label{partition_map}
    $\partitionMap{P}{X} = \{w\in X\times P \mid \fst{w}\in \snd{w}\}$.
\end{definition}

% from §22 Item 13: quotient space
\begin{definition}\label{quotient_space}
    $P$ is a quotient space of $X$ with respect to $T$ and $T_P$ iff
        $P$ is a partition of $X$ and
        $T$ is a topology on $X$ and
        $T_P = \quotientTopology{\partitionMap{P}{X}}{T}{X}{P}$.
\end{definition}

% from §22 helper lemma 2: restriction preimage for saturated subsets
\begin{lemma}\label{preimg_restrl_saturated}
    Suppose $p$ is a surjection from $X$ to $Y$.
    Suppose $A\subseteq X$.
    Suppose $A$ is saturated with respect to $p$.
    Let $q = \restrl{p}{A}$.
    Suppose $V\subseteq \img{p}{A}$.
    Then $\preimg{q}{V} = \preimg{p}{V}$.
\end{lemma}
\begin{proof}
    Show that $\preimg{q}{V}\subseteq \preimg{p}{V}$.
    \begin{subproof}
        Show that for all $x$ such that $x\in\preimg{q}{V}$ we have $x\in\preimg{p}{V}$.
        \begin{subproof}
            Fix $x$.
            Assume $x\in\preimg{q}{V}$.
            Take $y\in V$ such that $x\mathrel{q} y$ by \cref{preimg_iff}.
            Thus $x\mathrel{p} y$ by \cref{restrl_iff}.
            Thus $x\in\preimg{p}{V}$ by \cref{preimg_iff}.
        \end{subproof}
        Thus $\preimg{q}{V}\subseteq \preimg{p}{V}$ by \cref{subseteq}.
    \end{subproof}
    Show that $\preimg{p}{V}\subseteq \preimg{q}{V}$.
    \begin{subproof}
        Show that for all $x$ such that $x\in\preimg{p}{V}$ we have $x\in\preimg{q}{V}$.
        \begin{subproof}
            Fix $x$.
            Assume $x\in\preimg{p}{V}$.
            Take $y\in V$ such that $x\mathrel{p} y$ by \cref{preimg_iff}.
            Take $a\in A$ such that $a\mathrel{p} y$ by \cref{img_iff,subseteq}.
            Thus $a\in\preimg{p}{\{y\}}$ by \cref{preimg_iff,singleton_iff}.
            Thus $a\in A\inter \preimg{p}{\{y\}}$ by \cref{inter_intro}.
            Show that $A\inter \preimg{p}{\{y\}}\neq \emptyset$.
            \begin{subproof}
                Suppose not.
                Thus $A\inter \preimg{p}{\{y\}} = \emptyset$.
                Thus $a\notin A\inter \preimg{p}{\{y\}}$ by \cref{emptyset}.
                Contradiction.
            \end{subproof}
            Thus $A$ intersects $\preimg{p}{\{y\}}$ by \cref{intersects}.
            $p\in\funs{X}{Y}$ by \cref{surj}.
            Thus $\ran{p}\subseteq Y$ by \cref{funs_ran}.
            Thus $y\in\ran{p}$ by \cref{ran_iff,subseteq}.
            Thus $\preimg{p}{\{y\}}\subseteq A$ by \cref{saturated}.
            Thus $x\in A$ by \cref{preimg_iff,subseteq,singleton_iff}.
            Thus $x\mathrel{q} y$ by \cref{restrl_iff}.
            Thus $x\in\preimg{q}{V}$ by \cref{preimg_iff}.
        \end{subproof}
        Thus $\preimg{p}{V}\subseteq \preimg{q}{V}$ by \cref{subseteq}.
    \end{subproof}
    Thus $\preimg{q}{V} = \preimg{p}{V}$ by \cref{subseteq_antisymmetric}.
\end{proof}

% from §22 helper lemma 3: image of intersection with saturated set
\begin{lemma}\label{img_inter_saturated}
    Suppose $p$ is a function.
    Suppose $A\subseteq \dom{p}$.
    Suppose $A$ is saturated with respect to $p$.
    Suppose $U\subseteq \dom{p}$.
    Then $\img{p}{U\inter A} = \img{p}{U}\inter \img{p}{A}$.
\end{lemma}
\begin{proof}
    Show that $\img{p}{U\inter A}\subseteq \img{p}{U}\inter \img{p}{A}$.
    \begin{subproof}
        Show that for all $y$ such that $y\in\img{p}{U\inter A}$ we have
            $y\in\img{p}{U}\inter \img{p}{A}$.
        \begin{subproof}
            Fix $y$.
            Assume $y\in\img{p}{U\inter A}$.
            Take $x\in U\inter A$ such that $x\mathrel{p} y$ by \cref{img_iff}.
            Thus $x\in U$ and $x\in A$ by \cref{inter}.
            Thus $y\in\img{p}{U}$ and $y\in\img{p}{A}$ by \cref{img_iff}.
            Thus $y\in\img{p}{U}\inter \img{p}{A}$ by \cref{inter_intro}.
        \end{subproof}
        Thus $\img{p}{U\inter A}\subseteq \img{p}{U}\inter \img{p}{A}$ by \cref{subseteq}.
    \end{subproof}
    Show that $\img{p}{U}\inter \img{p}{A}\subseteq \img{p}{U\inter A}$.
    \begin{subproof}
        Show that for all $y$ such that $y\in\img{p}{U}\inter \img{p}{A}$ we have
            $y\in\img{p}{U\inter A}$.
        \begin{subproof}
            Fix $y$.
            Assume $y\in\img{p}{U}\inter \img{p}{A}$.
            Thus $y\in\img{p}{U}$ and $y\in\img{p}{A}$ by \cref{inter}.
            Take $u\in U$ such that $u\mathrel{p} y$ by \cref{img_iff}.
            Take $a\in A$ such that $a\mathrel{p} y$ by \cref{img_iff}.
            Thus $a\in\preimg{p}{\{y\}}$ by \cref{preimg_iff,singleton_iff}.
            Thus $a\in A\inter \preimg{p}{\{y\}}$ by \cref{inter_intro}.
            Show that $A\inter \preimg{p}{\{y\}}\neq \emptyset$.
            \begin{subproof}
                Suppose not.
                Thus $A\inter \preimg{p}{\{y\}} = \emptyset$.
                Thus $a\notin A\inter \preimg{p}{\{y\}}$ by \cref{emptyset}.
                Contradiction.
            \end{subproof}
            Thus $A$ intersects $\preimg{p}{\{y\}}$ by \cref{intersects}.
            $\img{p}{A}\subseteq \ran{p}$ by \cref{img_subseteq_ran}.
            Thus $y\in\ran{p}$ by \cref{subseteq}.
            Thus $\preimg{p}{\{y\}}\subseteq A$ by \cref{saturated}.
            Thus $u\in A$ by \cref{preimg_iff,subseteq,singleton_iff}.
            Thus $u\in U\inter A$ by \cref{inter_intro}.
            Thus $y\in\img{p}{U\inter A}$ by \cref{img_iff}.
        \end{subproof}
        Thus $\img{p}{U}\inter \img{p}{A}\subseteq \img{p}{U\inter A}$ by \cref{subseteq}.
    \end{subproof}
    Thus $\img{p}{U\inter A} = \img{p}{U}\inter \img{p}{A}$ by \cref{subseteq_antisymmetric}.
\end{proof}

% from §22 helper lemma 4: preimage of image of a saturated set
\begin{lemma}\label{preimg_img_saturated}
    Suppose $p$ is a function.
    Suppose $A\subseteq \dom{p}$.
    Suppose $A$ is saturated with respect to $p$.
    Then $\preimg{p}{\img{p}{A}} = A$.
\end{lemma}
\begin{proof}
    Show that $A\subseteq \preimg{p}{\img{p}{A}}$.
    \begin{subproof}
        Show that for all $x$ such that $x\in A$ we have
            $x\in\preimg{p}{\img{p}{A}}$.
        \begin{subproof}
            Fix $x$.
            Assume $x\in A$.
            Thus $x\in\dom{p}$ by \cref{subseteq}.
            Thus $(x,p(x))\in p$ by \cref{function_apply_elim}.
            Thus $p(x)\in\img{p}{A}$ by \cref{img_iff}.
            Thus $x\in\preimg{p}{\img{p}{A}}$ by \cref{preimg_iff}.
        \end{subproof}
        Thus $A\subseteq \preimg{p}{\img{p}{A}}$ by \cref{subseteq}.
    \end{subproof}
    Show that $\preimg{p}{\img{p}{A}}\subseteq A$.
    \begin{subproof}
        Show that for all $x$ such that $x\in\preimg{p}{\img{p}{A}}$ we have $x\in A$.
        \begin{subproof}
            Fix $x$.
            Assume $x\in\preimg{p}{\img{p}{A}}$.
            Take $y\in\img{p}{A}$ such that $x\mathrel{p} y$ by \cref{preimg_iff}.
            Take $a\in A$ such that $a\mathrel{p} y$ by \cref{img_iff}.
            Thus $a\in\preimg{p}{\{y\}}$ by \cref{preimg_iff,singleton_iff}.
            Thus $a\in A\inter \preimg{p}{\{y\}}$ by \cref{inter_intro}.
            Show that $A\inter \preimg{p}{\{y\}}\neq \emptyset$.
            \begin{subproof}
                Suppose not.
                Thus $A\inter \preimg{p}{\{y\}} = \emptyset$.
                Thus $a\notin A\inter \preimg{p}{\{y\}}$ by \cref{emptyset}.
                Contradiction.
            \end{subproof}
            Thus $A$ intersects $\preimg{p}{\{y\}}$ by \cref{intersects}.
            $\img{p}{A}\subseteq \ran{p}$ by \cref{img_subseteq_ran}.
            Thus $y\in\ran{p}$ by \cref{subseteq}.
            Thus $\preimg{p}{\{y\}}\subseteq A$ by \cref{saturated}.
            Thus $x\in A$ by \cref{preimg_iff,subseteq,singleton_iff}.
        \end{subproof}
        Thus $\preimg{p}{\img{p}{A}}\subseteq A$ by \cref{subseteq}.
    \end{subproof}
    Thus $\preimg{p}{\img{p}{A}} = A$ by \cref{subseteq_antisymmetric}.
\end{proof}

% from §22 Item 14: restriction to a saturated subspace
\begin{theorem}\label{quotient_map_restriction_open_or_closed}
    Suppose $p$ is a quotient map from $X$ to $Y$ with respect to $T_X$ and $T_Y$.
    Suppose $A$ is a subspace of $X$ with respect to $T_X$.
    Suppose $A$ is saturated with respect to $p$.
    Let $q = \restrl{p}{A}$.
    Let $T_A = \subspaceTopology{T_X}{A}$.
    Let $T_B = \subspaceTopology{T_Y}{\img{p}{A}}$.
    Suppose $A$ is open in $X$ with respect to $T_X$
        or $A$ is closed in $X$ with respect to $T_X$.
    Then $q$ is a quotient map from $A$ to $\img{p}{A}$ with respect to $T_A$ and $T_B$.
\end{theorem}
\begin{proof}
    $p\in\surj{X}{Y}$ by \cref{quotient_map}.
    Thus $p\in\funs{X}{Y}$ by \cref{surj}.
    Thus $p$ is a function by \cref{funs_is_function}.
    $p$ is continuous from $X$ to $Y$ with respect to $T_X$ and $T_Y$ by \cref{quotient_map}.
    Thus $T_X$ is a topology on $X$ by \cref{continuous}.
    Thus $T_Y$ is a topology on $Y$ by \cref{continuous}.
    $A\subseteq X$ by \cref{subspace_of}.
    Thus $A$ is a subspace of $X$ with respect to $T_X$ by \cref{subspace_of}.
    $\img{p}{A}\subseteq \ran{p}$ by \cref{img_subseteq_ran}.
    Thus $\img{p}{A}\subseteq Y$ by \cref{funs_ran,subseteq_transitive}.
    Thus $\img{p}{A}$ is a subspace of $Y$ with respect to $T_Y$ by \cref{subspace_of}.
    Thus $T_A$ is a topology on $A$ by \cref{subspace_topology_is_topology,subspace_of}.
    Thus $T_B$ is a topology on $\img{p}{A}$ by \cref{subspace_topology_is_topology,subspace_of}.
    Show that $q$ is a surjection from $A$ to $\img{p}{A}$.
    \begin{subproof}
        Thus $q$ is a function by \cref{restrl_of_function_is_function}.
        Thus $q$ is right-unique and $q$ is a relation.
        $\dom{q} = \dom{p}\inter A$ by \cref{dom_of_restrl_eq_inter}.
        Thus $\dom{q} = A$ by \cref{funs,inter_absorb_supseteq_left,subspace_of}.
        $A\subseteq \dom{p}$ by \cref{funs,subseteq_transitive,subspace_of}.
        $\ran{q} = \img{p}{A}$ by \cref{restrl_ran}.
        Thus $\dom{q}\subseteq A$ by \cref{subseteq_refl}.
        Thus $\ran{q}\subseteq \img{p}{A}$ by \cref{subseteq_refl}.
        Thus $q\in\rels{A}{\img{p}{A}}$ by \cref{rels_intro_dom_and_ran}.
        Thus $q\in\funs{A}{\img{p}{A}}$ by \cref{funs}.
        Thus $q\in\surj{A}{\img{p}{A}}$ by \cref{funs_surj_iff,restrl_ran}.
    \end{subproof}
    Show that $q$ is continuous from $A$ to $\img{p}{A}$ with respect to $T_A$ and $T_B$.
    \begin{subproof}
        $p$ is continuous from $X$ to $Y$ with respect to $T_X$ and $T_Y$ by \cref{quotient_map}.
        Thus $q$ is continuous from $A$ to $Y$ with respect to $T_A$ and $T_Y$
            by \cref{continuous_restrict_domain}.
        $A\subseteq \dom{p}$ by \cref{funs,subseteq_transitive,subspace_of}.
        $\img{q}{A} = \img{p}{A\inter A}$ by \cref{restrl_img}.
        $A\inter A = A$ by \cref{inter_absorb_supseteq_left,subseteq_refl}.
        Thus $\img{q}{A} = \img{p}{A}$ by \cref{subseteq}.
        Thus $\img{q}{A}\subseteq \img{p}{A}$ by \cref{subseteq_refl}.
        Thus $q$ is continuous from $A$ to $\img{p}{A}$ with respect to $T_A$ and $T_B$
            by \cref{continuous_restrict_range}.
    \end{subproof}
    \begin{byCase}
    \caseOf{$A$ is open in $X$ with respect to $T_X$.}
        Show that for all $V$ such that $V\subseteq \img{p}{A}$ if $\preimg{q}{V}$ is open in $A$ with respect to $T_A$ then
            $V$ is open in $\img{p}{A}$ with respect to $T_B$.
        \begin{subproof}
            Fix $V$.
            Assume $V\subseteq \img{p}{A}$.
            Assume $\preimg{q}{V}$ is open in $A$ with respect to $T_A$.
            Thus $\preimg{q}{V}$ is open in $X$ with respect to $T_X$ by \cref{open_in_subspace_open_in_space,subspace_of}.
            $\preimg{q}{V} = \preimg{p}{V}$ by \cref{preimg_restrl_saturated}.
            Thus $\preimg{p}{V}$ is open in $X$ with respect to $T_X$ by \cref{open_in}.
            Thus $V\subseteq Y$ by \cref{subseteq_transitive}.
            Thus $V$ is open in $Y$ with respect to $T_Y$ by \cref{quotient_map}.
            Thus $V\in T_Y$ by \cref{open_in}.
            Thus $\img{p}{A}\inter V = V$ by \cref{inter_absorb_supseteq_left}.
            Thus $V = \img{p}{A}\inter V$.
            Thus $V\in T_B$ by \cref{subspace_topology}.
            Thus $V$ is open in $\img{p}{A}$ with respect to $T_B$ by \cref{open_in}.
        \end{subproof}
        Thus $q$ is a quotient map from $A$ to $\img{p}{A}$ with respect to $T_A$ and $T_B$ by \cref{quotient_map}.
    \caseOf{$A$ is closed in $X$ with respect to $T_X$.}
        Show that for all $B$ such that $B\subseteq \img{p}{A}$ we have
            $B$ is closed in $\img{p}{A}$ with respect to $T_B$ iff
            $\preimg{q}{B}$ is closed in $A$ with respect to $T_A$.
        \begin{subproof}
            Fix $B$.
            Assume $B\subseteq \img{p}{A}$.
            Show that if $B$ is closed in $\img{p}{A}$ with respect to $T_B$ then
                $\preimg{q}{B}$ is closed in $A$ with respect to $T_A$.
            \begin{subproof}
                Assume $B$ is closed in $\img{p}{A}$ with respect to $T_B$.
                Take $C$ such that $C$ is closed in $Y$ with respect to $T_Y$ and
                    $B = C\inter \img{p}{A}$ by \cref{closed_in_subspace_iff,subspace_of}.
                $Y\setminus C$ is open in $Y$ with respect to $T_Y$ by \cref{closed_in}.
                Thus $\preimg{p}{Y\setminus C}$ is open in $X$ with respect to $T_X$ by \cref{continuous}.
                $C\subseteq Y$ by \cref{closed_in}.
                $\preimg{p}{Y\setminus C} = X\setminus \preimg{p}{C}$ by \cref{preimg_setminus_function}.
                $\preimg{p}{C}\subseteq \dom{p}$ by \cref{preimg_subseteq_dom}.
                Thus $\preimg{p}{C}\subseteq X$ by \cref{funs,subseteq}.
                Thus $\preimg{p}{C}$ is closed in $X$ with respect to $T_X$ by \cref{closed_in}.
                $\preimg{q}{B} = \preimg{p}{B}$ by \cref{preimg_restrl_saturated}.
                $\preimg{p}{B} = \preimg{p}{C}\inter \preimg{p}{\img{p}{A}}$ by \cref{preimg_inter_function}.
                $A\subseteq \dom{p}$ by \cref{funs,subseteq_transitive,subspace_of}.
                $\preimg{p}{\img{p}{A}} = A$ by \cref{preimg_img_saturated}.
                Thus $\preimg{q}{B} = \preimg{p}{C}\inter A$ by \cref{subseteq}.
                Thus $\preimg{q}{B}$ is closed in $A$ with respect to $T_A$ by \cref{closed_in_subspace_iff,subspace_of}.
            \end{subproof}
            Show that if $\preimg{q}{B}$ is closed in $A$ with respect to $T_A$ then
                $B$ is closed in $\img{p}{A}$ with respect to $T_B$.
            \begin{subproof}
                Assume $\preimg{q}{B}$ is closed in $A$ with respect to $T_A$.
                Take $C$ such that $C$ is closed in $X$ with respect to $T_X$ and
                    $\preimg{q}{B} = C\inter A$ by \cref{closed_in_subspace_iff,subspace_of}.
                $\preimg{q}{B} = \preimg{p}{B}$ by \cref{preimg_restrl_saturated}.
                Show that $\preimg{p}{B}$ is closed in $X$ with respect to $T_X$.
                \begin{subproof}
                    $X\setminus C$ is open in $X$ with respect to $T_X$ by \cref{closed_in}.
                    $X\setminus A$ is open in $X$ with respect to $T_X$ by \cref{closed_in}.
                    Let $M = \{X\setminus C, X\setminus A\}$.
                    Show that $M\subseteq T_X$.
                    \begin{subproof}
                        Show that for all $U$ such that $U\in M$ we have $U\in T_X$.
                        \begin{subproof}
                            Fix $U$.
                            Assume $U\in M$.
                            Thus $U = X\setminus C$ or $U = X\setminus A$ by \cref{upair_elim}.
                            \begin{byCase}
                            \caseOf{$U = X\setminus C$.}
                                Thus $U\in T_X$ by \cref{open_in}.
                            \caseOf{$U = X\setminus A$.}
                                Thus $U\in T_X$ by \cref{open_in}.
                            \end{byCase}
                        \end{subproof}
                        Thus $M\subseteq T_X$ by \cref{subseteq}.
                    \end{subproof}
                    Thus $\unions{M}\in T_X$ by \cref{topology_on}.
                    $\unions{M} = (X\setminus C)\union (X\setminus A)$ by \cref{union_as_unions}.
                    Thus $(X\setminus C)\union (X\setminus A)$ is open in $X$ with respect to $T_X$ by \cref{open_in}.
                    $X\setminus (C\inter A) = (X\setminus C)\union (X\setminus A)$ by \cref{setminus_inter}.
                    Thus $X\setminus (C\inter A)$ is open in $X$ with respect to $T_X$ by \cref{open_in}.
                    $C\subseteq X$ by \cref{closed_in}.
                    $C\inter A\subseteq C$ by \cref{inter_lower_left}.
                    Thus $C\inter A\subseteq X$ by \cref{subseteq_transitive}.
                    Thus $C\inter A$ is closed in $X$ with respect to $T_X$ by \cref{closed_in}.
                    Thus $\preimg{p}{B}$ is closed in $X$ with respect to $T_X$ by \cref{subseteq}.
                \end{subproof}
                Thus $B\subseteq Y$ by \cref{subseteq_transitive}.
                Thus $B$ is closed in $Y$ with respect to $T_Y$ by \cref{quotient_map_closed_iff,quotient_map}.
                Thus $\img{p}{A}\inter B = B$ by \cref{inter_absorb_supseteq_left}.
                Thus $B = \img{p}{A}\inter B$.
                Let $D = B$.
                Thus $D$ is closed in $Y$ with respect to $T_Y$.
                Thus $B = D\inter \img{p}{A}$.
                Thus $B$ is closed in $\img{p}{A}$ with respect to $T_B$ by \cref{closed_in_subspace_iff,subspace_of}.
            \end{subproof}
            Thus $B$ is closed in $\img{p}{A}$ with respect to $T_B$ iff
                $\preimg{q}{B}$ is closed in $A$ with respect to $T_A$.
        \end{subproof}
        Thus $q$ is a quotient map from $A$ to $\img{p}{A}$ with respect to $T_A$ and $T_B$
            by \cref{quotient_map_closed_iff}.
    \end{byCase}
\end{proof}

% from §22 Item 15: restriction of a quotient map with open or closed map
\begin{theorem}\label{quotient_map_restriction_open_or_closed_map}
    Suppose $p$ is a quotient map from $X$ to $Y$ with respect to $T_X$ and $T_Y$.
    Suppose $A$ is a subspace of $X$ with respect to $T_X$.
    Suppose $A$ is saturated with respect to $p$.
    Let $q = \restrl{p}{A}$.
    Let $T_A = \subspaceTopology{T_X}{A}$.
    Let $T_B = \subspaceTopology{T_Y}{\img{p}{A}}$.
    Suppose $p$ is an open map from $X$ to $Y$ with respect to $T_X$ and $T_Y$
        or $p$ is a closed map from $X$ to $Y$ with respect to $T_X$ and $T_Y$.
    Then $q$ is a quotient map from $A$ to $\img{p}{A}$ with respect to $T_A$ and $T_B$.
\end{theorem}
\begin{proof}
    $p\in\surj{X}{Y}$ by \cref{quotient_map}.
    Thus $p\in\funs{X}{Y}$ by \cref{surj}.
    Thus $p$ is a function by \cref{funs_is_function}.
    $p$ is continuous from $X$ to $Y$ with respect to $T_X$ and $T_Y$ by \cref{quotient_map}.
    Thus $T_X$ is a topology on $X$ by \cref{continuous}.
    Thus $T_Y$ is a topology on $Y$ by \cref{continuous}.
    $A\subseteq X$ by \cref{subspace_of}.
    Thus $A$ is a subspace of $X$ with respect to $T_X$ by \cref{subspace_of}.
    $\img{p}{A}\subseteq \ran{p}$ by \cref{img_subseteq_ran}.
    Thus $\img{p}{A}\subseteq Y$ by \cref{funs_ran,subseteq_transitive}.
    Thus $\img{p}{A}$ is a subspace of $Y$ with respect to $T_Y$ by \cref{subspace_of}.
    Thus $T_A$ is a topology on $A$ by \cref{subspace_topology_is_topology,subspace_of}.
    Thus $T_B$ is a topology on $\img{p}{A}$ by \cref{subspace_topology_is_topology,subspace_of}.
    Show that $q$ is a surjection from $A$ to $\img{p}{A}$.
    \begin{subproof}
        Thus $q$ is a function by \cref{restrl_of_function_is_function}.
        Thus $q$ is right-unique and $q$ is a relation.
        $\dom{q} = \dom{p}\inter A$ by \cref{dom_of_restrl_eq_inter}.
        Thus $\dom{q} = A$ by \cref{funs,inter_absorb_supseteq_left,subspace_of}.
        $A\subseteq \dom{p}$ by \cref{funs,subseteq_transitive,subspace_of}.
        $\ran{q} = \img{p}{A}$ by \cref{restrl_ran}.
        Thus $\dom{q}\subseteq A$ by \cref{subseteq_refl}.
        Thus $\ran{q}\subseteq \img{p}{A}$ by \cref{subseteq_refl}.
        Thus $q\in\rels{A}{\img{p}{A}}$ by \cref{rels_intro_dom_and_ran}.
        Thus $q\in\funs{A}{\img{p}{A}}$ by \cref{funs}.
        Thus $q\in\surj{A}{\img{p}{A}}$ by \cref{funs_surj_iff,restrl_ran}.
    \end{subproof}
    Show that $q$ is continuous from $A$ to $\img{p}{A}$ with respect to $T_A$ and $T_B$.
    \begin{subproof}
        Thus $q$ is continuous from $A$ to $Y$ with respect to $T_A$ and $T_Y$
            by \cref{continuous_restrict_domain}.
        $A\subseteq \dom{p}$ by \cref{funs,subseteq_transitive,subspace_of}.
        $\img{q}{A} = \img{p}{A\inter A}$ by \cref{restrl_img}.
        $A\inter A = A$ by \cref{inter_absorb_supseteq_left,subseteq_refl}.
        Thus $\img{q}{A} = \img{p}{A}$ by \cref{subseteq}.
        Thus $\img{q}{A}\subseteq \img{p}{A}$ by \cref{subseteq_refl}.
        Thus $q$ is continuous from $A$ to $\img{p}{A}$ with respect to $T_A$ and $T_B$
            by \cref{continuous_restrict_range}.
    \end{subproof}
    \begin{byCase}
    \caseOf{$p$ is an open map from $X$ to $Y$ with respect to $T_X$ and $T_Y$.}
        Show that $q$ is an open map from $A$ to $\img{p}{A}$ with respect to $T_A$ and $T_B$.
        \begin{subproof}
            Show that for all $U$ such that $U$ is open in $A$ with respect to $T_A$ we have
                $\img{q}{U}$ is open in $\img{p}{A}$ with respect to $T_B$.
            \begin{subproof}
                Fix $U$.
                Assume $U$ is open in $A$ with respect to $T_A$.
                Thus $U\in T_A$ by \cref{open_in}.
                $T_A\subseteq \pow{A}$ by \cref{topology_on}.
                Thus $U\in\pow{A}$ by \cref{subseteq}.
                Thus $U\subseteq A$ by \cref{pow_iff}.
                Take $W\in T_X$ such that $U = A\inter W$ by \cref{subspace_topology}.
                $W$ is open in $X$ with respect to $T_X$ by \cref{open_in}.
                $A\subseteq \dom{p}$ by \cref{funs,subseteq_transitive,subspace_of}.
                $\img{q}{U} = \img{p}{A\inter U}$ by \cref{restrl_img}.
                $A\inter U = U$ by \cref{inter_absorb_supseteq_left}.
                Thus $\img{q}{U} = \img{p}{U}$ by \cref{subseteq}.
                Thus $\img{q}{U} = \img{p}{A\inter W}$ by \cref{subseteq}.
                $T_X\subseteq \pow{X}$ by \cref{topology_on}.
                Thus $W\in\pow{X}$ by \cref{subseteq}.
                Thus $W\subseteq X$ by \cref{pow_iff}.
                $\dom{p} = X$ by \cref{funs}.
                Thus $W\subseteq \dom{p}$ by \cref{subseteq_transitive}.
                Thus $\img{q}{U} = \img{p}{W\inter A}$ by \cref{inter_comm,subseteq}.
                Thus $\img{q}{U} = \img{p}{W}\inter \img{p}{A}$ by \cref{img_inter_saturated}.
                Thus $\img{q}{U} = \img{p}{A}\inter \img{p}{W}$ by \cref{inter_comm}.
                Thus $\img{p}{W}$ is open in $Y$ with respect to $T_Y$ by \cref{open_map}.
                Thus $\img{p}{W}\in T_Y$ by \cref{open_in}.
                Thus $\img{q}{U}\in T_B$ by \cref{subspace_topology}.
                Thus $\img{q}{U}$ is open in $\img{p}{A}$ with respect to $T_B$ by \cref{open_in}.
            \end{subproof}
            $q\in\funs{A}{\img{p}{A}}$ by \cref{surj_to_fun}.
            Thus $q$ is a function from $A$ to $\img{p}{A}$ by \cref{funs}.
            Thus $q$ is an open map from $A$ to $\img{p}{A}$ with respect to $T_A$ and $T_B$ by \cref{open_map}.
        \end{subproof}
        Thus $q$ is a quotient map from $A$ to $\img{p}{A}$ with respect to $T_A$ and $T_B$
            by \cref{open_or_closed_surjection_quotient}.
    \caseOf{$p$ is a closed map from $X$ to $Y$ with respect to $T_X$ and $T_Y$.}
        Show that $q$ is a closed map from $A$ to $\img{p}{A}$ with respect to $T_A$ and $T_B$.
        \begin{subproof}
            Show that for all $U$ such that $U$ is closed in $A$ with respect to $T_A$ we have
                $\img{q}{U}$ is closed in $\img{p}{A}$ with respect to $T_B$.
            \begin{subproof}
                Fix $U$.
                Assume $U$ is closed in $A$ with respect to $T_A$.
                Take $C$ such that $C$ is closed in $X$ with respect to $T_X$ and $U = C\inter A$
                    by \cref{closed_in_subspace_iff,subspace_of}.
                $U\subseteq A$ by \cref{closed_in}.
                $A\subseteq \dom{p}$ by \cref{funs,subseteq_transitive,subspace_of}.
                $\img{q}{U} = \img{p}{A\inter U}$ by \cref{restrl_img}.
                $A\inter U = U$ by \cref{inter_absorb_supseteq_left}.
                Thus $\img{q}{U} = \img{p}{U}$ by \cref{subseteq}.
                Thus $\img{q}{U} = \img{p}{C\inter A}$ by \cref{subseteq}.
                $C\subseteq X$ by \cref{closed_in}.
                $\dom{p} = X$ by \cref{funs}.
                Thus $C\subseteq \dom{p}$ by \cref{subseteq_transitive}.
                Thus $\img{q}{U} = \img{p}{C}\inter \img{p}{A}$ by \cref{img_inter_saturated}.
                Thus $\img{p}{C}$ is closed in $Y$ with respect to $T_Y$ by \cref{closed_map}.
                Thus $\img{p}{A}\inter \img{p}{C} = \img{p}{C}\inter \img{p}{A}$ by \cref{inter_comm}.
                Thus $\img{q}{U}$ is closed in $\img{p}{A}$ with respect to $T_B$ by \cref{closed_in_subspace_iff,subspace_of}.
            \end{subproof}
            $q\in\funs{A}{\img{p}{A}}$ by \cref{surj_to_fun}.
            Thus $q$ is a function from $A$ to $\img{p}{A}$ by \cref{funs}.
            Thus $q$ is a closed map from $A$ to $\img{p}{A}$ with respect to $T_A$ and $T_B$ by \cref{closed_map}.
        \end{subproof}
        Thus $q$ is a quotient map from $A$ to $\img{p}{A}$ with respect to $T_A$ and $T_B$
            by \cref{open_or_closed_surjection_quotient}.
    \end{byCase}
\end{proof}

% from §22 Item 16: composite of quotient maps
\begin{lemma}\label{quotient_map_composite}
    Suppose $p$ is a quotient map from $X$ to $Y$ with respect to $T_X$ and $T_Y$.
    Suppose $q$ is a quotient map from $Y$ to $Z$ with respect to $T_Y$ and $T_Z$.
    Then $q\circ p$ is a quotient map from $X$ to $Z$ with respect to $T_X$ and $T_Z$.
\end{lemma}
\begin{proof}
    $p\in\surj{X}{Y}$ by \cref{quotient_map}.
    $q\in\surj{Y}{Z}$ by \cref{quotient_map}.
    Thus $q\circ p\in\surj{X}{Z}$ by \cref{surjections_circ}.
    $p$ is continuous from $X$ to $Y$ with respect to $T_X$ and $T_Y$ by \cref{quotient_map}.
    $q$ is continuous from $Y$ to $Z$ with respect to $T_Y$ and $T_Z$ by \cref{quotient_map}.
    Thus $q\circ p$ is continuous from $X$ to $Z$ with respect to $T_X$ and $T_Z$
        by \cref{continuous_composite}.
    Show that for all $U$ such that $U\subseteq Z$ if $\preimg{(q\circ p)}{U}$ is open in $X$ with respect to $T_X$
        then $U$ is open in $Z$ with respect to $T_Z$.
    \begin{subproof}
        Fix $U$.
        Assume $U\subseteq Z$.
        Assume $\preimg{(q\circ p)}{U}$ is open in $X$ with respect to $T_X$.
        Show that $\preimg{(q\circ p)}{U} = \preimg{p}{\preimg{q}{U}}$.
        \begin{subproof}
            Show that for all $x$ such that $x\in\preimg{(q\circ p)}{U}$ we have
                $x\in\preimg{p}{\preimg{q}{U}}$.
            \begin{subproof}
                Fix $x$.
                Assume $x\in\preimg{(q\circ p)}{U}$.
                Take $z\in U$ such that $x\mathrel{(q\circ p)} z$ by \cref{preimg_iff}.
                Take $y$ such that $x\mathrel{p} y$ and $y\mathrel{q} z$ by \cref{circ_iff}.
                Thus $y\in\preimg{q}{U}$ by \cref{preimg_iff}.
                Thus $x\in\preimg{p}{\preimg{q}{U}}$ by \cref{preimg_iff}.
            \end{subproof}
            Thus $\preimg{(q\circ p)}{U}\subseteq \preimg{p}{\preimg{q}{U}}$ by \cref{subseteq}.
            Show that for all $x$ such that $x\in\preimg{p}{\preimg{q}{U}}$ we have
                $x\in\preimg{(q\circ p)}{U}$.
            \begin{subproof}
                Fix $x$.
                Assume $x\in\preimg{p}{\preimg{q}{U}}$.
                Take $y\in\preimg{q}{U}$ such that $x\mathrel{p} y$ by \cref{preimg_iff}.
                Take $z\in U$ such that $y\mathrel{q} z$ by \cref{preimg_iff}.
                Thus $x\mathrel{(q\circ p)} z$ by \cref{circ_iff}.
                Thus $x\in\preimg{(q\circ p)}{U}$ by \cref{preimg_iff}.
            \end{subproof}
            Thus $\preimg{p}{\preimg{q}{U}}\subseteq \preimg{(q\circ p)}{U}$ by \cref{subseteq}.
            Thus $\preimg{(q\circ p)}{U} = \preimg{p}{\preimg{q}{U}}$ by \cref{subseteq_antisymmetric}.
        \end{subproof}
        Thus $\preimg{p}{\preimg{q}{U}}$ is open in $X$ with respect to $T_X$ by \cref{open_in}.
        $q\in\surj{Y}{Z}$ by \cref{quotient_map}.
        Thus $q\in\funs{Y}{Z}$ by \cref{surj_to_fun}.
        Thus $\dom{q} = Y$ by \cref{funs}.
        $\preimg{q}{U}\subseteq \dom{q}$ by \cref{preimg_subseteq_dom}.
        Thus $\preimg{q}{U}\subseteq Y$ by \cref{subseteq}.
        Thus $\preimg{q}{U}$ is open in $Y$ with respect to $T_Y$ by \cref{quotient_map}.
        Thus $U$ is open in $Z$ with respect to $T_Z$ by \cref{quotient_map}.
    \end{subproof}
    Thus $q\circ p$ is a quotient map from $X$ to $Z$ with respect to $T_X$ and $T_Z$ by \cref{quotient_map}.
\end{proof}

% from §22 helper definition 2: induced map
\begin{definition}\label{induced_map}
    $\inducedMap{p}{g}{X}{Y}{Z} =
        \{w\in Y\times Z \mid \snd{w}\in \img{g}{\preimg{p}{\{\fst{w}\}}}\}$.
\end{definition}

% from §22 Item 17: induced map exists
\begin{lemma}\label{quotient_map_induced_map}
    Suppose $p$ is a quotient map from $X$ to $Y$ with respect to $T_X$ and $T_Y$.
    Suppose $g$ is a function from $X$ to $Z$.
    Suppose for all $y\in Y$ there exists $z$ such that for all $x\in\preimg{p}{\{y\}}$ we have $g(x) = z$.
    Then there exists $f$ such that
        $f$ is a function from $Y$ to $Z$ and
        $f\circ p = g$.
\end{lemma}
\begin{proof}
    $p\in\surj{X}{Y}$ by \cref{quotient_map}.
    Thus $p\in\funs{X}{Y}$ by \cref{surj}.
    Thus $p$ is a function by \cref{funs_is_function}.
    Let $f = \inducedMap{p}{g}{X}{Y}{Z}$.
    Show that $f$ is a function from $Y$ to $Z$.
    \begin{subproof}
        Show that $f$ is right-unique.
        \begin{subproof}
            Show that for all $y,z_1,z_2$ such that $y\mathrel{f} z_1$ and $y\mathrel{f} z_2$ we have $z_1 = z_2$.
            \begin{subproof}
                Fix $y$.
                Fix $z_1$.
                Fix $z_2$.
                Assume $y\mathrel{f} z_1$ and $y\mathrel{f} z_2$.
                Thus $(y,z_1)\in Y\times Z$ by \cref{induced_map}.
                Take $y_0,z_0$ such that $(y,z_1) = (y_0,z_0)$ and $y_0\in Y$ and $z_0\in Z$ by \cref{times_elem_is_tuple}.
                Thus $y = y_0$ by \cref{pair_eq_iff}.
                Thus $y\in Y$.
                Thus $\snd{(y,z_1)}\in \img{g}{\preimg{p}{\{\fst{(y,z_1)}\}}}$ by \cref{induced_map}.
                Thus $\snd{(y,z_2)}\in \img{g}{\preimg{p}{\{\fst{(y,z_2)}\}}}$ by \cref{induced_map}.
                Thus $z_1\in \img{g}{\preimg{p}{\{y\}}}$ and $z_2\in \img{g}{\preimg{p}{\{y\}}}$ by \cref{fst_eq,snd_eq}.
                Take $x_1\in \preimg{p}{\{y\}}$ such that $x_1\mathrel{g} z_1$ by \cref{img_iff}.
                Take $x_2\in \preimg{p}{\{y\}}$ such that $x_2\mathrel{g} z_2$ by \cref{img_iff}.
                $x_1\in \dom{p}$ and $x_2\in \dom{p}$ by \cref{preimg_subseteq_dom,subseteq}.
                $\dom{p} = X$ by \cref{funs}.
                Thus $x_1\in X$ and $x_2\in X$ by \cref{subseteq}.
                $g\in\funs{X}{Z}$ by \cref{funs}.
                Thus $g$ is a function by \cref{funs_is_function}.
                Thus $g(x_1) = z_1$ and $g(x_2) = z_2$ by \cref{function_apply_intro}.
                Take $z$ such that for all $x\in\preimg{p}{\{y\}}$ we have $g(x) = z$.
                Thus $g(x_1) = z$ and $g(x_2) = z$.
                Thus $z_1 = z$ and $z_2 = z$.
                Thus $z_1 = z_2$.
            \end{subproof}
        \end{subproof}
        Show that $\dom{f} = Y$.
        \begin{subproof}
            Show that $\dom{f}\subseteq Y$.
            \begin{subproof}
                Show that for all $y$ such that $y\in\dom{f}$ we have $y\in Y$.
                \begin{subproof}
                    Fix $y$.
                    Assume $y\in\dom{f}$.
                    Take $z$ such that $(y,z)\in f$ by \cref{dom}.
                    Thus $(y,z)\in Y\times Z$ by \cref{induced_map}.
                    Take $y_0,z_0$ such that $(y,z) = (y_0,z_0)$ and $y_0\in Y$ and $z_0\in Z$ by \cref{times_elem_is_tuple}.
                    Thus $y = y_0$ by \cref{pair_eq_iff}.
                    Thus $y\in Y$.
                \end{subproof}
                Thus $\dom{f}\subseteq Y$ by \cref{subseteq}.
            \end{subproof}
            Show that $Y\subseteq \dom{f}$.
            \begin{subproof}
                Show that for all $y$ such that $y\in Y$ we have $y\in\dom{f}$.
                \begin{subproof}
                    Fix $y$.
                    Assume $y\in Y$.
                    Take $x\in X$ such that $p(x) = y$ by \cref{surj}.
                    Thus $x\in\dom{p}$ by \cref{funs}.
                    Thus $x\mathrel{p} y$ by \cref{function_apply_iff}.
                    $g\in\funs{X}{Z}$ by \cref{funs}.
                    Thus $g$ is a function to $Z$ and $X = \dom{g}$ by \cref{funs_elim}.
                    Thus $x\in\dom{g}$ by \cref{subseteq}.
                    Thus $g(x)\in Z$ by \cref{funs_elim}.
                    Thus $(y,g(x))\in Y\times Z$ by \cref{times_tuple_intro}.
                    $y\in\{y\}$ by \cref{singleton_intro}.
                    Thus $x\in\preimg{p}{\{y\}}$ by \cref{preimg_iff}.
                    Thus $x\mathrel{g} g(x)$ by \cref{function_apply_elim}.
                    Thus $g(x)\in \img{g}{\preimg{p}{\{y\}}}$ by \cref{img_iff}.
                    Thus $\snd{(y,g(x))}\in \img{g}{\preimg{p}{\{\fst{(y,g(x))}\}}}$ by \cref{fst_eq,snd_eq}.
                    Thus $(y,g(x))\in f$ by \cref{induced_map}.
                    Thus $y\in\dom{f}$ by \cref{dom}.
                \end{subproof}
                Thus $Y\subseteq \dom{f}$ by \cref{subseteq}.
            \end{subproof}
            Thus $\dom{f} = Y$ by \cref{subseteq_antisymmetric}.
        \end{subproof}
        Show that $f\subseteq Y\times Z$.
        \begin{subproof}
            Show that for all $w$ such that $w\in f$ we have $w\in Y\times Z$.
            \begin{subproof}
                Fix $w$.
                Assume $w\in f$.
                Thus $w\in Y\times Z$ by \cref{induced_map}.
                Take $y,z$ such that $w = (y,z)$ and $y\in Y$ and $z\in Z$ by \cref{times_elem_is_tuple}.
                Thus $\snd{(y,z)}\in \img{g}{\preimg{p}{\{\fst{(y,z)}\}}}$ by \cref{induced_map}.
                Thus $z\in \img{g}{\preimg{p}{\{y\}}}$ by \cref{fst_eq,snd_eq}.
                $\img{g}{\preimg{p}{\{y\}}}\subseteq \ran{g}$ by \cref{img_subseteq_ran}.
                Thus $z\in\ran{g}$ by \cref{subseteq}.
                $g\in\funs{X}{Z}$ by \cref{funs}.
                Thus $\ran{g}\subseteq Z$ by \cref{funs_ran}.
                Thus $z\in Z$ by \cref{subseteq}.
                Thus $w\in Y\times Z$ by \cref{times_tuple_intro}.
            \end{subproof}
            Thus $f\subseteq Y\times Z$ by \cref{subseteq}.
        \end{subproof}
        Thus $f\in\rels{Y}{Z}$ by \cref{rels_intro}.
        Thus $f$ is a relation by \cref{rels_is_relation}.
        Thus $f\in\funs{Y}{Z}$ by \cref{funs}.
        Thus $f$ is a function from $Y$ to $Z$ by \cref{funs}.
    \end{subproof}
    Show that $f\circ p = g$.
    \begin{subproof}
        Show that $g\subseteq f\circ p$.
        \begin{subproof}
            Show that for all $w$ such that $w\in g$ we have $w\in f\circ p$.
            \begin{subproof}
                Fix $w$.
                Assume $w\in g$.
                $g\in\funs{X}{Z}$ by \cref{funs}.
                Thus $g\in\rels{X}{Z}$ by \cref{funs}.
                Thus $g\subseteq X\times Z$ by \cref{rels_elim}.
                Thus $w\in X\times Z$ by \cref{subseteq}.
                Take $x,z$ such that $w = (x,z)$ and $x\in X$ and $z\in Z$ by \cref{times_elem_is_tuple}.
                Thus $(x,z)\in g$.
                Thus $g$ is a function by \cref{funs_is_function}.
                Thus $z = g(x)$ by \cref{function_apply_intro}.
                $\dom{p} = X$ by \cref{funs}.
                Thus $x\in\dom{p}$ by \cref{subseteq}.
                Take $y$ such that $x\mathrel{p} y$ by \cref{dom_iff}.
                Thus $y\in\ran{p}$ by \cref{ran_iff}.
                Thus $y\in Y$ by \cref{funs_ran,subseteq}.
                $y\in\{y\}$ by \cref{singleton_intro}.
                Thus $x\in\preimg{p}{\{y\}}$ by \cref{preimg_iff}.
                Thus $z\in \img{g}{\preimg{p}{\{y\}}}$ by \cref{img_iff}.
                Thus $(y,z)\in Y\times Z$ by \cref{times_tuple_intro}.
                Thus $\snd{(y,z)}\in \img{g}{\preimg{p}{\{\fst{(y,z)}\}}}$ by \cref{fst_eq,snd_eq}.
                Thus $(y,z)\in f$ by \cref{induced_map}.
                Thus $(x,z)\in f\circ p$ by \cref{circ_iff}.
                Thus $w\in f\circ p$.
            \end{subproof}
            Thus $g\subseteq f\circ p$ by \cref{subseteq}.
        \end{subproof}
        Show that $f\circ p\subseteq g$.
        \begin{subproof}
            Show that for all $w$ such that $w\in f\circ p$ we have $w\in g$.
            \begin{subproof}
                Fix $w$.
                Assume $w\in f\circ p$.
                $f\circ p$ is a relation by \cref{circ_is_relation}.
                Take $x,z$ such that $w = (x,z)$ by \cref{relation}.
                Thus $(x,z)\in f\circ p$.
                Take $y$ such that $x\mathrel{p} y$ and $y\mathrel{f} z$ by \cref{circ_iff}.
                Thus $(y,z)\in Y\times Z$ by \cref{induced_map}.
                Thus $\snd{(y,z)}\in \img{g}{\preimg{p}{\{\fst{(y,z)}\}}}$ by \cref{induced_map}.
                Thus $z\in \img{g}{\preimg{p}{\{y\}}}$ by \cref{fst_eq,snd_eq}.
                $g\in\funs{X}{Z}$ by \cref{funs}.
                Thus $g$ is a function by \cref{funs_is_function}.
                Take $x_1\in \preimg{p}{\{y\}}$ such that $x_1\mathrel{g} z$ by \cref{img_iff}.
                Thus $g(x_1) = z$ by \cref{function_apply_intro}.
                $y\in\{y\}$ by \cref{singleton_intro}.
                Thus $x\in\preimg{p}{\{y\}}$ by \cref{preimg_iff}.
                Thus $y\in Y$ by \cref{funs_ran,ran_iff,subseteq}.
                Take $z_0$ such that for all $x\in\preimg{p}{\{y\}}$ we have $g(x) = z_0$.
                Thus $g(x) = z_0$ and $g(x_1) = z_0$.
                Thus $g(x) = z$.
                $\dom{p} = X$ by \cref{funs}.
                Thus $x\in\dom{p}$ by \cref{preimg_subseteq_dom,subseteq}.
                Thus $x\in X$ by \cref{subseteq}.
                $\dom{g} = X$ by \cref{funs}.
                Thus $x\in\dom{g}$ by \cref{subseteq}.
                Thus $(x,z)\in g$ by \cref{function_apply_elim}.
                Thus $w\in g$.
            \end{subproof}
            Thus $f\circ p\subseteq g$ by \cref{subseteq}.
        \end{subproof}
        Thus $f\circ p = g$ by \cref{subseteq_antisymmetric}.
    \end{subproof}
\end{proof}

% from §22 Item 18: induced map continuity
\begin{lemma}\label{quotient_map_induced_map_continuous}
    Suppose $p$ is a quotient map from $X$ to $Y$ with respect to $T_X$ and $T_Y$.
    Suppose $g$ is a function from $X$ to $Z$.
    Suppose for all $y\in Y$ there exists $z$ such that for all $x\in\preimg{p}{\{y\}}$ we have $g(x) = z$.
    Suppose $f$ is a function from $Y$ to $Z$.
    Suppose $f\circ p = g$.
    Then $f$ is continuous from $Y$ to $Z$ with respect to $T_Y$ and $T_Z$
        iff $g$ is continuous from $X$ to $Z$ with respect to $T_X$ and $T_Z$.
\end{lemma}
\begin{proof}
    Show that if $f$ is continuous from $Y$ to $Z$ with respect to $T_Y$ and $T_Z$
        then $g$ is continuous from $X$ to $Z$ with respect to $T_X$ and $T_Z$.
    \begin{subproof}
        Assume $f$ is continuous from $Y$ to $Z$ with respect to $T_Y$ and $T_Z$.
        Thus $p$ is continuous from $X$ to $Y$ with respect to $T_X$ and $T_Y$ by \cref{quotient_map}.
        Thus $T_X$ is a topology on $X$ by \cref{continuous}.
        Thus $f\circ p$ is continuous from $X$ to $Z$ with respect to $T_X$ and $T_Z$ by \cref{continuous_composite}.
        Thus $g$ is continuous from $X$ to $Z$ with respect to $T_X$ and $T_Z$.
    \end{subproof}
    Show that if $g$ is continuous from $X$ to $Z$ with respect to $T_X$ and $T_Z$
        then $f$ is continuous from $Y$ to $Z$ with respect to $T_Y$ and $T_Z$.
    \begin{subproof}
        Assume $g$ is continuous from $X$ to $Z$ with respect to $T_X$ and $T_Z$.
        Thus $p$ is continuous from $X$ to $Y$ with respect to $T_X$ and $T_Y$ by \cref{quotient_map}.
        Thus $T_Y$ is a topology on $Y$ by \cref{continuous}.
        $T_Z$ is a topology on $Z$ by \cref{continuous}.
        Show that for all $V$ such that $V$ is open in $Z$ with respect to $T_Z$
            we have $\preimg{f}{V}$ is open in $Y$ with respect to $T_Y$.
        \begin{subproof}
            Fix $V$.
            Assume $V$ is open in $Z$ with respect to $T_Z$.
            Thus $\preimg{g}{V}$ is open in $X$ with respect to $T_X$ by \cref{continuous}.
            Show that $\preimg{g}{V} = \preimg{p}{\preimg{f}{V}}$.
            \begin{subproof}
                Show that $\preimg{g}{V}\subseteq \preimg{p}{\preimg{f}{V}}$.
                \begin{subproof}
                    Show that for all $x$ such that $x\in\preimg{g}{V}$ we have $x\in\preimg{p}{\preimg{f}{V}}$.
                    \begin{subproof}
                        Fix $x$.
                        Assume $x\in\preimg{g}{V}$.
                        Take $z\in V$ such that $x\mathrel{g} z$ by \cref{preimg_iff}.
                        Thus $x\mathrel{(f\circ p)} z$.
                        Take $y$ such that $x\mathrel{p} y$ and $y\mathrel{f} z$ by \cref{circ_iff}.
                        Thus $y\in\preimg{f}{V}$ by \cref{preimg_iff}.
                        Thus $x\in\preimg{p}{\preimg{f}{V}}$ by \cref{preimg_iff}.
                    \end{subproof}
                    Thus $\preimg{g}{V}\subseteq \preimg{p}{\preimg{f}{V}}$ by \cref{subseteq}.
                \end{subproof}
                Show that $\preimg{p}{\preimg{f}{V}}\subseteq \preimg{g}{V}$.
                \begin{subproof}
                    Show that for all $x$ such that $x\in\preimg{p}{\preimg{f}{V}}$ we have $x\in\preimg{g}{V}$.
                    \begin{subproof}
                        Fix $x$.
                        Assume $x\in\preimg{p}{\preimg{f}{V}}$.
                        Take $y\in\preimg{f}{V}$ such that $x\mathrel{p} y$ by \cref{preimg_iff}.
                        Take $z\in V$ such that $y\mathrel{f} z$ by \cref{preimg_iff}.
                        Thus $x\mathrel{(f\circ p)} z$ by \cref{circ_iff}.
                        Thus $x\mathrel{g} z$.
                        Thus $x\in\preimg{g}{V}$ by \cref{preimg_iff}.
                    \end{subproof}
                    Thus $\preimg{p}{\preimg{f}{V}}\subseteq \preimg{g}{V}$ by \cref{subseteq}.
                \end{subproof}
                Thus $\preimg{g}{V} = \preimg{p}{\preimg{f}{V}}$ by \cref{subseteq_antisymmetric}.
            \end{subproof}
            Thus $\preimg{p}{\preimg{f}{V}}$ is open in $X$ with respect to $T_X$.
            $\preimg{f}{V}\subseteq \dom{f}$ by \cref{preimg_subseteq_dom}.
            $\dom{f} = Y$ by \cref{funs}.
            Thus $\preimg{f}{V}\subseteq Y$ by \cref{subseteq}.
            Thus $\preimg{f}{V}$ is open in $Y$ with respect to $T_Y$ by \cref{quotient_map}.
        \end{subproof}
        Thus $f$ is continuous from $Y$ to $Z$ with respect to $T_Y$ and $T_Z$ by \cref{continuous}.
    \end{subproof}
\end{proof}

% from §22 Item 19: induced map and quotient maps
\begin{lemma}\label{quotient_map_induced_map_quotient}
    Suppose $p$ is a quotient map from $X$ to $Y$ with respect to $T_X$ and $T_Y$.
    Suppose $g$ is a function from $X$ to $Z$.
    Suppose for all $y\in Y$ there exists $z$ such that for all $x\in\preimg{p}{\{y\}}$ we have $g(x) = z$.
    Suppose $f$ is a function from $Y$ to $Z$.
    Suppose $f\circ p = g$.
    Then $f$ is a quotient map from $Y$ to $Z$ with respect to $T_Y$ and $T_Z$
        iff $g$ is a quotient map from $X$ to $Z$ with respect to $T_X$ and $T_Z$.
\end{lemma}
\begin{proof}
    Show that if $f$ is a quotient map from $Y$ to $Z$ with respect to $T_Y$ and $T_Z$
        then $g$ is a quotient map from $X$ to $Z$ with respect to $T_X$ and $T_Z$.
    \begin{subproof}
        Assume $f$ is a quotient map from $Y$ to $Z$ with respect to $T_Y$ and $T_Z$.
        Thus $f\circ p$ is a quotient map from $X$ to $Z$ with respect to $T_X$ and $T_Z$
            by \cref{quotient_map_composite}.
        Thus $g$ is a quotient map from $X$ to $Z$ with respect to $T_X$ and $T_Z$.
    \end{subproof}
    Show that if $g$ is a quotient map from $X$ to $Z$ with respect to $T_X$ and $T_Z$
        then $f$ is a quotient map from $Y$ to $Z$ with respect to $T_Y$ and $T_Z$.
    \begin{subproof}
        Assume $g$ is a quotient map from $X$ to $Z$ with respect to $T_X$ and $T_Z$.
        Show that $f$ is a surjection from $Y$ to $Z$.
        \begin{subproof}
            Show that for all $z\in Z$ there exists $y\in Y$ such that $f(y) = z$.
            \begin{subproof}
                Fix $z\in Z$.
                $g\in\surj{X}{Z}$ by \cref{quotient_map}.
                Take $x\in X$ such that $g(x) = z$ by \cref{surj}.
                Let $y = p(x)$.
                $p\in\surj{X}{Y}$ by \cref{quotient_map}.
                Thus $p\in\funs{X}{Y}$ by \cref{surj}.
                Thus $p$ is a function by \cref{funs_is_function}.
                Thus $y\in Y$ by \cref{funs_type_apply}.
                $f\in\funs{Y}{Z}$.
                Thus $f$ is a function by \cref{funs_is_function}.
                $g\in\funs{X}{Z}$ by \cref{surj}.
                Thus $g$ is a function by \cref{funs_is_function}.
                $x\in\dom{p}$ by \cref{funs}.
                Thus $x\mathrel{p} y$ by \cref{function_apply_elim}.
                $x\in\dom{g}$ by \cref{funs}.
                Thus $x\mathrel{g} z$ by \cref{function_apply_iff}.
                Thus $x\mathrel{(f\circ p)} z$.
                Take $y_0$ such that $x\mathrel{p} y_0$ and $y_0\mathrel{f} z$ by \cref{circ_iff}.
                Thus $y_0 = y$ by \cref{function_rightunique}.
                Thus $y\mathrel{f} z$.
                Thus $f(y) = z$ by \cref{function_apply_iff}.
            \end{subproof}
            Thus $f\in\surj{Y}{Z}$ by \cref{surj}.
        \end{subproof}
        Thus $g$ is continuous from $X$ to $Z$ with respect to $T_X$ and $T_Z$ by \cref{quotient_map}.
        Thus $f$ is continuous from $Y$ to $Z$ with respect to $T_Y$ and $T_Z$
            by \cref{quotient_map_induced_map_continuous}.
        Show that for all $V$ such that $V\subseteq Z$ if $\preimg{f}{V}$ is open in $Y$ with respect to $T_Y$
            then $V$ is open in $Z$ with respect to $T_Z$.
        \begin{subproof}
            Fix $V$.
            Assume $V\subseteq Z$.
            Assume $\preimg{f}{V}$ is open in $Y$ with respect to $T_Y$.
            Thus $p$ is continuous from $X$ to $Y$ with respect to $T_X$ and $T_Y$ by \cref{quotient_map}.
            Thus $\preimg{p}{\preimg{f}{V}}$ is open in $X$ with respect to $T_X$ by \cref{continuous}.
            Show that $\preimg{p}{\preimg{f}{V}} = \preimg{g}{V}$.
            \begin{subproof}
                Show that $\preimg{p}{\preimg{f}{V}}\subseteq \preimg{g}{V}$.
                \begin{subproof}
                    Show that for all $x$ such that $x\in\preimg{p}{\preimg{f}{V}}$ we have $x\in\preimg{g}{V}$.
                    \begin{subproof}
                        Fix $x$.
                        Assume $x\in\preimg{p}{\preimg{f}{V}}$.
                        Take $y\in\preimg{f}{V}$ such that $x\mathrel{p} y$ by \cref{preimg_iff}.
                        Take $z\in V$ such that $y\mathrel{f} z$ by \cref{preimg_iff}.
                        Thus $x\mathrel{(f\circ p)} z$ by \cref{circ_iff}.
                        Thus $x\mathrel{g} z$.
                        Thus $x\in\preimg{g}{V}$ by \cref{preimg_iff}.
                    \end{subproof}
                    Thus $\preimg{p}{\preimg{f}{V}}\subseteq \preimg{g}{V}$ by \cref{subseteq}.
                \end{subproof}
                Show that $\preimg{g}{V}\subseteq \preimg{p}{\preimg{f}{V}}$.
                \begin{subproof}
                    Show that for all $x$ such that $x\in\preimg{g}{V}$ we have $x\in\preimg{p}{\preimg{f}{V}}$.
                    \begin{subproof}
                        Fix $x$.
                        Assume $x\in\preimg{g}{V}$.
                        Take $z\in V$ such that $x\mathrel{g} z$ by \cref{preimg_iff}.
                        Thus $x\mathrel{(f\circ p)} z$.
                        Take $y$ such that $x\mathrel{p} y$ and $y\mathrel{f} z$ by \cref{circ_iff}.
                        Thus $y\in\preimg{f}{V}$ by \cref{preimg_iff}.
                        Thus $x\in\preimg{p}{\preimg{f}{V}}$ by \cref{preimg_iff}.
                    \end{subproof}
                    Thus $\preimg{g}{V}\subseteq \preimg{p}{\preimg{f}{V}}$ by \cref{subseteq}.
                \end{subproof}
                Thus $\preimg{p}{\preimg{f}{V}} = \preimg{g}{V}$ by \cref{subseteq_antisymmetric}.
            \end{subproof}
            Thus $\preimg{g}{V}$ is open in $X$ with respect to $T_X$.
            Thus $V$ is open in $Z$ with respect to $T_Z$ by \cref{quotient_map}.
        \end{subproof}
        Thus $f$ is a quotient map from $Y$ to $Z$ with respect to $T_Y$ and $T_Z$ by \cref{quotient_map}.
    \end{subproof}
\end{proof}

% from §22 helper definition 3: fiber partition
\begin{definition}\label{fiber_partition}
    $\fiberPartition{g}{Z} = \{\preimg{g}{\{z\}} \mid z\in Z\}$.
\end{definition}

% from §22 helper lemma 5: fibers form a partition
\begin{lemma}\label{fiber_partition_is_partition}
    Suppose $g$ is a surjection from $X$ to $Z$.
    Let $X_s = \fiberPartition{g}{Z}$.
    Then $X_s$ is a partition of $X$.
\end{lemma}
\begin{proof}
    $g\in\surj{X}{Z}$.
    Thus $g\in\funs{X}{Z}$ by \cref{surj}.
    Thus $g$ is a function by \cref{funs_is_function}.
    Show that $\emptyset\notin X_s$.
    \begin{subproof}
        Suppose not.
        Thus $\emptyset\in X_s$.
        Take $z\in Z$ such that $\emptyset = \preimg{g}{\{z\}}$ by \cref{fiber_partition}.
        Take $x\in X$ such that $g(x) = z$ by \cref{surj}.
        Thus $x\in\dom{g}$ by \cref{funs}.
        Thus $x\mathrel{g} z$ by \cref{function_apply_iff}.
        $z\in\{z\}$ by \cref{singleton_intro}.
        Thus $x\in\preimg{g}{\{z\}}$ by \cref{preimg_iff}.
        Contradiction by \cref{emptyset}.
    \end{subproof}
    Show that for all $A,B$ such that $A\in X_s$ and $B\in X_s$ and $A\neq B$ we have $A$ is disjoint from $B$.
    \begin{subproof}
        Fix $A$.
        Fix $B$.
        Assume $A\in X_s$ and $B\in X_s$ and $A\neq B$.
        Take $z_1\in Z$ such that $A = \preimg{g}{\{z_1\}}$ by \cref{fiber_partition}.
        Take $z_2\in Z$ such that $B = \preimg{g}{\{z_2\}}$ by \cref{fiber_partition}.
        Show that $A$ is disjoint from $B$.
        \begin{subproof}
            Suppose not.
            Then there exists $x$ such that $x\in A$ and $x\in B$.
            Thus $x\in\preimg{g}{\{z_1\}}$ and $x\in\preimg{g}{\{z_2\}}$.
            Take $b_1\in\{z_1\}$ such that $x\mathrel{g} b_1$ by \cref{preimg_iff}.
            Take $b_2\in\{z_2\}$ such that $x\mathrel{g} b_2$ by \cref{preimg_iff}.
            Thus $b_1 = z_1$ by \cref{singleton_elim}.
            Thus $b_2 = z_2$ by \cref{singleton_elim}.
            Thus $x\mathrel{g} z_1$ and $x\mathrel{g} z_2$.
            Thus $z_1 = z_2$ by \cref{function_rightunique}.
            Thus $A = B$.
            Contradiction.
        \end{subproof}
        Thus $A$ is disjoint from $B$ by \cref{disjoint}.
    \end{subproof}
    Thus $X_s$ is a partition by \cref{partition}.
    Show that $\unions{X_s} = X$.
    \begin{subproof}
        Show that $X\subseteq \unions{X_s}$.
        \begin{subproof}
            Show that for all $x$ such that $x\in X$ we have $x\in\unions{X_s}$.
            \begin{subproof}
                Fix $x$.
                Assume $x\in X$.
                $g(x)\in Z$ by \cref{funs_type_apply}.
                Thus $\preimg{g}{\{g(x)\}}\in X_s$ by \cref{fiber_partition}.
                $x\in\dom{g}$ by \cref{funs}.
                Thus $x\mathrel{g} g(x)$ by \cref{function_apply_elim}.
                Thus $x\in\preimg{g}{\{g(x)\}}$ by \cref{preimg_iff,singleton_intro}.
                Thus $x\in\unions{X_s}$ by \cref{unions_intro}.
            \end{subproof}
            Thus $X\subseteq \unions{X_s}$ by \cref{subseteq}.
        \end{subproof}
        Show that $\unions{X_s}\subseteq X$.
        \begin{subproof}
            Show that for all $x$ such that $x\in\unions{X_s}$ we have $x\in X$.
            \begin{subproof}
                Fix $x$.
                Assume $x\in\unions{X_s}$.
                Take $A\in X_s$ such that $x\in A$ by \cref{unions_iff}.
                Take $z\in Z$ such that $A = \preimg{g}{\{z\}}$ by \cref{fiber_partition}.
                Thus $x\in\preimg{g}{\{z\}}$.
                $\preimg{g}{\{z\}}\subseteq \dom{g}$ by \cref{preimg_subseteq_dom}.
                Thus $x\in\dom{g}$ by \cref{subseteq}.
                $\dom{g} = X$ by \cref{funs}.
                Thus $x\in X$ by \cref{subseteq}.
            \end{subproof}
            Thus $\unions{X_s}\subseteq X$ by \cref{subseteq}.
        \end{subproof}
        Thus $\unions{X_s} = X$ by \cref{subseteq_antisymmetric}.
    \end{subproof}
    Thus $X_s$ is a partition of $X$.
\end{proof}

% from §22 helper lemma 6: partition map is a surjection
\begin{lemma}\label{partition_map_surjection}
    Suppose $P$ is a partition of $X$.
    Let $p = \partitionMap{P}{X}$.
    Then $p$ is a surjection from $X$ to $P$.
\end{lemma}
\begin{proof}
    $P$ is a partition.
    $\unions{P} = X$.
    Show that $p\in\funs{X}{P}$.
    \begin{subproof}
        Show that $p\subseteq X\times P$.
        \begin{subproof}
            Show that for all $w$ such that $w\in p$ we have $w\in X\times P$.
            \begin{subproof}
                Fix $w$.
                Assume $w\in p$.
                Thus $w\in X\times P$ by \cref{partition_map}.
            \end{subproof}
            Thus $p\subseteq X\times P$ by \cref{subseteq}.
        \end{subproof}
        Thus $p\in\rels{X}{P}$ by \cref{rels_intro}.
        Show that $\dom{p} = X$.
        \begin{subproof}
            $\dom{p}\subseteq X$ by \cref{rels_dom_subseteq}.
            Show that $X\subseteq \dom{p}$.
            \begin{subproof}
                Show that for all $x$ such that $x\in X$ we have $x\in\dom{p}$.
                \begin{subproof}
                    Fix $x$.
                    Assume $x\in X$.
                    Thus $x\in\unions{P}$ by \cref{subseteq}.
                    Take $C\in P$ such that $x\in C$ by \cref{unions_iff}.
                    Thus $(x,C)\in X\times P$ by \cref{times_tuple_intro}.
                    Thus $\fst{(x,C)}\in \snd{(x,C)}$ by \cref{fst_eq,snd_eq}.
                    Thus $(x,C)\in p$ by \cref{partition_map}.
                    Thus $x\in\dom{p}$ by \cref{dom}.
                \end{subproof}
                Thus $X\subseteq \dom{p}$ by \cref{subseteq}.
            \end{subproof}
            Thus $\dom{p} = X$ by \cref{subseteq_antisymmetric}.
        \end{subproof}
        Show that $p$ is right-unique.
        \begin{subproof}
            Show that for all $x,C,D$ such that $x\mathrel{p} C$ and $x\mathrel{p} D$ we have $C = D$.
            \begin{subproof}
                Fix $x$.
                Fix $C$.
                Fix $D$.
                Assume $x\mathrel{p} C$ and $x\mathrel{p} D$.
                Thus $(x,C)\in p$ and $(x,D)\in p$.
                Thus $(x,C)\in X\times P$ and $(x,D)\in X\times P$ by \cref{partition_map}.
                Thus $C\in P$ and $D\in P$ by \cref{times_tuple_elim}.
                Thus $x\in C$ and $x\in D$ by \cref{partition_map,fst_eq,snd_eq}.
                Show that $C = D$.
                \begin{subproof}
                    Suppose not.
                    Then $C\neq D$.
                    Then $C$ is disjoint from $D$ by \cref{partition}.
                    Contradiction by \cref{disjoint}.
                \end{subproof}
            \end{subproof}
        \end{subproof}
        Thus $p\in\funs{X}{P}$ by \cref{funs}.
    \end{subproof}
    Show that for all $C\in P$ there exists $x\in X$ such that $p(x) = C$.
    \begin{subproof}
        Fix $C\in P$.
        $\emptyset\notin P$ by \cref{partition}.
        Thus $C\neq\emptyset$.
        Take $x$ such that $x\in C$ by \cref{nonempty_inhabited}.
        Thus $x\in\unions{P}$ by \cref{unions_intro}.
        Thus $x\in X$ by \cref{subseteq}.
        Thus $(x,C)\in X\times P$ by \cref{times_tuple_intro}.
        Thus $\fst{(x,C)}\in \snd{(x,C)}$ by \cref{fst_eq,snd_eq}.
        Thus $(x,C)\in p$ by \cref{partition_map}.
        Thus $p$ is a function by \cref{funs_is_function}.
        Thus $p(x) = C$ by \cref{function_apply_intro}.
    \end{subproof}
    Thus $p\in\surj{X}{P}$ by \cref{surj}.
\end{proof}

% from §22 Item 20: Corollary 22.3(a) induced bijective map
\begin{lemma}\label{corollary_quotient_space_induced_bijection}
    Suppose $g$ is a surjection from $X$ to $Z$.
    Suppose $g$ is continuous from $X$ to $Z$ with respect to $T_X$ and $T_Z$.
    Let $X_s = \fiberPartition{g}{Z}$.
    Let $p = \partitionMap{X_s}{X}$.
    Let $T_s = \quotientTopology{p}{T_X}{X}{X_s}$.
    Then there exists $f$ such that
        $f$ is a bijection from $X_s$ to $Z$ and
        $f$ is continuous from $X_s$ to $Z$ with respect to $T_s$ and $T_Z$ and
        $f\circ p = g$.
\end{lemma}
\begin{proof}
    $T_X$ is a topology on $X$ by \cref{continuous}.
    $g\in\surj{X}{Z}$.
    Thus $g\in\funs{X}{Z}$ by \cref{surj}.
    Thus $g$ is a function by \cref{funs_is_function}.
    $X_s$ is a partition of $X$ by \cref{fiber_partition_is_partition}.
    Thus $p$ is a surjection from $X$ to $X_s$ by \cref{partition_map_surjection}.
    Thus $p$ is a quotient map from $X$ to $X_s$ with respect to $T_X$ and $T_s$ by \cref{quotient_topology_quotient_map}.
    $g\in\surj{X}{Z}$.
    Thus $g\in\funs{X}{Z}$ by \cref{surj}.
    Thus $g$ is a function by \cref{funs_is_function}.
    Show that for all $y\in X_s$ there exists $z$ such that for all $x\in\preimg{p}{\{y\}}$ we have $g(x) = z$.
    \begin{subproof}
        Fix $y$.
        Assume $y\in X_s$.
        Take $z\in Z$ such that $y = \preimg{g}{\{z\}}$ by \cref{fiber_partition}.
        Show that for all $x\in\preimg{p}{\{y\}}$ we have $g(x) = z$.
        \begin{subproof}
            Fix $x$.
            Assume $x\in\preimg{p}{\{y\}}$.
            Take $y_0\in\{y\}$ such that $x\mathrel{p} y_0$ by \cref{preimg_iff}.
            Thus $y_0 = y$ by \cref{singleton_elim}.
            Thus $(x,y)\in p$.
            Thus $x\in y$ by \cref{partition_map,fst_eq,snd_eq}.
            Thus $x\in\preimg{g}{\{z\}}$.
            Take $b\in\{z\}$ such that $x\mathrel{g} b$ by \cref{preimg_iff}.
            Thus $b = z$ by \cref{singleton_elim}.
            Thus $x\mathrel{g} z$.
            Thus $g(x) = z$ by \cref{function_apply_intro}.
        \end{subproof}
    \end{subproof}
    Take $f$ such that $f$ is a function from $X_s$ to $Z$ and $f\circ p = g$ by \cref{quotient_map_induced_map}.
    Show that $f$ is a bijection from $X_s$ to $Z$.
    \begin{subproof}
        Show that $f$ is a surjection from $X_s$ to $Z$.
        \begin{subproof}
            Show that for all $z\in Z$ there exists $y\in X_s$ such that $f(y) = z$.
            \begin{subproof}
                Fix $z\in Z$.
                Take $x\in X$ such that $g(x) = z$ by \cref{surj}.
                $p\in\funs{X}{X_s}$ by \cref{surj}.
                Thus $p$ is a function by \cref{funs_is_function}.
                Thus $p(x)\in X_s$ by \cref{funs_type_apply}.
                $f\in\funs{X_s}{Z}$.
                Thus $f$ is a function by \cref{funs_is_function}.
                $x\in\dom{p}$ by \cref{funs}.
                Thus $x\mathrel{p} p(x)$ by \cref{function_apply_elim}.
                $x\in\dom{g}$ by \cref{funs}.
                Thus $x\mathrel{g} z$ by \cref{function_apply_iff}.
                Thus $x\mathrel{(f\circ p)} z$.
                Take $y$ such that $x\mathrel{p} y$ and $y\mathrel{f} z$ by \cref{circ_iff}.
                Thus $y = p(x)$ by \cref{function_rightunique}.
                Thus $p(x)\mathrel{f} z$.
                Thus $f(p(x)) = z$ by \cref{function_apply_iff}.
            \end{subproof}
            Thus $f\in\surj{X_s}{Z}$ by \cref{surj}.
        \end{subproof}
        Show that $f$ is injective.
        \begin{subproof}
            Show that for all $y_1,y_2,z$ such that $y_1\mathrel{f} z$ and $y_2\mathrel{f} z$ we have $y_1 = y_2$.
            \begin{subproof}
                Fix $y_1$.
                Fix $y_2$.
                Fix $z$.
                Assume $y_1\mathrel{f} z$ and $y_2\mathrel{f} z$.
                $f\in\funs{X_s}{Z}$.
                Thus $f$ is a function by \cref{funs_is_function}.
                Thus $f(y_1) = z$ and $f(y_2) = z$ by \cref{function_apply_intro}.
                Thus $y_1\in\dom{f}$ and $y_2\in\dom{f}$ by \cref{function_apply_iff}.
                $\dom{f} = X_s$ by \cref{funs}.
                Thus $y_1\in X_s$ and $y_2\in X_s$.
                Take $z_1\in Z$ such that $y_1 = \preimg{g}{\{z_1\}}$ by \cref{fiber_partition}.
                Take $z_2\in Z$ such that $y_2 = \preimg{g}{\{z_2\}}$ by \cref{fiber_partition}.
                $\emptyset\notin X_s$ by \cref{partition}.
                Thus $y_1\neq\emptyset$ and $y_2\neq\emptyset$.
                Take $x_1$ such that $x_1\in y_1$ by \cref{nonempty_inhabited}.
                Take $x_2$ such that $x_2\in y_2$ by \cref{nonempty_inhabited}.
                Thus $x_1\in\unions{X_s}$ and $x_2\in\unions{X_s}$ by \cref{unions_intro}.
                Thus $x_1\in X$ and $x_2\in X$ by \cref{subseteq}.
                Thus $(x_1,y_1)\in X\times X_s$ and $(x_2,y_2)\in X\times X_s$ by \cref{times_tuple_intro}.
                Thus $\fst{(x_1,y_1)}\in \snd{(x_1,y_1)}$ and $\fst{(x_2,y_2)}\in \snd{(x_2,y_2)}$ by \cref{fst_eq,snd_eq}.
                Thus $(x_1,y_1)\in p$ and $(x_2,y_2)\in p$ by \cref{partition_map}.
                $p\in\funs{X}{X_s}$ by \cref{surj}.
                Thus $p$ is a function by \cref{funs_is_function}.
                Thus $x_1\mathrel{p} y_1$ and $x_2\mathrel{p} y_2$.
                Thus $x_1\mathrel{(f\circ p)} z$ and $x_2\mathrel{(f\circ p)} z$ by \cref{circ_iff}.
                Thus $x_1\mathrel{g} z$ and $x_2\mathrel{g} z$.
                Thus $g(x_1) = z$ and $g(x_2) = z$ by \cref{function_apply_intro}.
                Thus $x_1\in\preimg{g}{\{z_1\}}$ and $x_2\in\preimg{g}{\{z_2\}}$.
                Take $b_1\in\{z_1\}$ such that $x_1\mathrel{g} b_1$ by \cref{preimg_iff}.
                Take $b_2\in\{z_2\}$ such that $x_2\mathrel{g} b_2$ by \cref{preimg_iff}.
                Thus $b_1 = z_1$ and $b_2 = z_2$ by \cref{singleton_elim}.
                Thus $x_1\mathrel{g} z_1$ and $x_2\mathrel{g} z_2$.
                Thus $g(x_1) = z_1$ and $g(x_2) = z_2$ by \cref{function_apply_intro}.
                Thus $z_1 = z$ and $z_2 = z$.
                Thus $y_1 = \preimg{g}{\{z\}}$ and $y_2 = \preimg{g}{\{z\}}$.
                Thus $y_1 = y_2$.
            \end{subproof}
            Thus $f$ is injective by \cref{injective}.
        \end{subproof}
        Thus $f\in\bijections{X_s}{Z}$ by \cref{bijections}.
    \end{subproof}
    Thus $f$ is a bijection from $X_s$ to $Z$.
    Thus $f$ is continuous from $X_s$ to $Z$ with respect to $T_s$ and $T_Z$
        by \cref{quotient_map_induced_map_continuous}.
\end{proof}

% from §22 Item 21: Corollary 22.3(a) homeomorphism criterion
\begin{lemma}\label{corollary_quotient_space_homeomorphism}
    Suppose $g$ is a surjection from $X$ to $Z$.
    Suppose $g$ is continuous from $X$ to $Z$ with respect to $T_X$ and $T_Z$.
    Let $X_s = \fiberPartition{g}{Z}$.
    Let $p = \partitionMap{X_s}{X}$.
    Let $T_s = \quotientTopology{p}{T_X}{X}{X_s}$.
    Suppose $f$ is a function from $X_s$ to $Z$.
    Suppose $f\circ p = g$.
    Suppose $f$ is a bijection from $X_s$ to $Z$.
    Then $f$ is a homeomorphism between $X_s$ and $Z$ with respect to $T_s$ and $T_Z$
        iff $g$ is a quotient map from $X$ to $Z$ with respect to $T_X$ and $T_Z$.
\end{lemma}
\begin{proof}
    $T_X$ is a topology on $X$ by \cref{continuous}.
    $X_s$ is a partition of $X$ by \cref{fiber_partition_is_partition}.
    Thus $p$ is a surjection from $X$ to $X_s$ by \cref{partition_map_surjection}.
    Thus $p$ is a quotient map from $X$ to $X_s$ with respect to $T_X$ and $T_s$ by \cref{quotient_topology_quotient_map}.
    Show that if $f$ is a homeomorphism between $X_s$ and $Z$ with respect to $T_s$ and $T_Z$
        then $g$ is a quotient map from $X$ to $Z$ with respect to $T_X$ and $T_Z$.
    \begin{subproof}
        Assume $f$ is a homeomorphism between $X_s$ and $Z$ with respect to $T_s$ and $T_Z$.
        Thus $f$ is continuous from $X_s$ to $Z$ with respect to $T_s$ and $T_Z$ by \cref{homeomorphism}.
        Show that $f$ is an open map from $X_s$ to $Z$ with respect to $T_s$ and $T_Z$.
        \begin{subproof}
            Show that for all $U$ such that $U$ is open in $X_s$ with respect to $T_s$
                we have $\img{f}{U}$ is open in $Z$ with respect to $T_Z$.
            \begin{subproof}
                Fix $U$.
                Assume $U$ is open in $X_s$ with respect to $T_s$.
                Thus $\converse{f}$ is continuous from $Z$ to $X_s$ with respect to $T_Z$ and $T_s$ by \cref{homeomorphism}.
                Thus $\preimg{\converse{f}}{U}$ is open in $Z$ with respect to $T_Z$ by \cref{continuous}.
                Thus $f$ is a relation by \cref{funs_is_relation,funs}.
                Thus $\preimg{\converse{f}}{U} = \img{f}{U}$ by \cref{preim_eq_img_of_converse,converse_converse_eq}.
            \end{subproof}
            Thus $f$ is an open map from $X_s$ to $Z$ with respect to $T_s$ and $T_Z$ by \cref{open_map}.
        \end{subproof}
        Thus $f$ is a surjection from $X_s$ to $Z$ by \cref{bijections}.
        $p\in\funs{X}{X_s}$ by \cref{surj}.
        $T_s$ is a topology on $X_s$ by \cref{quotient_topology_is_topology}.
        $T_Z$ is a topology on $Z$ by \cref{continuous}.
        Thus $f$ is a quotient map from $X_s$ to $Z$ with respect to $T_s$ and $T_Z$ by \cref{open_or_closed_surjection_quotient}.
        Thus $f\circ p$ is a quotient map from $X$ to $Z$ with respect to $T_X$ and $T_Z$
            by \cref{quotient_map_composite}.
        Thus $g$ is a quotient map from $X$ to $Z$ with respect to $T_X$ and $T_Z$.
    \end{subproof}
    Show that if $g$ is a quotient map from $X$ to $Z$ with respect to $T_X$ and $T_Z$
        then $f$ is a homeomorphism between $X_s$ and $Z$ with respect to $T_s$ and $T_Z$.
    \begin{subproof}
        Assume $g$ is a quotient map from $X$ to $Z$ with respect to $T_X$ and $T_Z$.
        Thus $g\in\surj{X}{Z}$ by \cref{quotient_map}.
        Thus $g\in\funs{X}{Z}$ by \cref{surj}.
        Thus $g$ is a function by \cref{funs_is_function}.
        Show that for all $y\in X_s$ there exists $z$ such that for all $x\in\preimg{p}{\{y\}}$ we have $g(x) = z$.
        \begin{subproof}
            Fix $y$.
            Assume $y\in X_s$.
            Take $z\in Z$ such that $y = \preimg{g}{\{z\}}$ by \cref{fiber_partition}.
            Show that for all $x\in\preimg{p}{\{y\}}$ we have $g(x) = z$.
            \begin{subproof}
                Fix $x$.
                Assume $x\in\preimg{p}{\{y\}}$.
                Take $y_0\in\{y\}$ such that $x\mathrel{p} y_0$ by \cref{preimg_iff}.
                Thus $y_0 = y$ by \cref{singleton_elim}.
                Thus $(x,y)\in p$.
                Thus $x\in y$ by \cref{partition_map,fst_eq,snd_eq}.
                Thus $x\in\preimg{g}{\{z\}}$.
                Take $b\in\{z\}$ such that $x\mathrel{g} b$ by \cref{preimg_iff}.
                Thus $b = z$ by \cref{singleton_elim}.
                Thus $x\mathrel{g} z$.
                Thus $g(x) = z$ by \cref{function_apply_intro}.
            \end{subproof}
        \end{subproof}
        Thus $f$ is a quotient map from $X_s$ to $Z$ with respect to $T_s$ and $T_Z$
            by \cref{quotient_map_induced_map_quotient}.
        Thus $f$ is continuous from $X_s$ to $Z$ with respect to $T_s$ and $T_Z$ by \cref{quotient_map}.
        Show that $\converse{f}$ is continuous from $Z$ to $X_s$ with respect to $T_Z$ and $T_s$.
        \begin{subproof}
            Show that for all $U$ such that $U$ is open in $X_s$ with respect to $T_s$
                we have $\preimg{\converse{f}}{U}$ is open in $Z$ with respect to $T_Z$.
            \begin{subproof}
                Fix $U$.
                Assume $U$ is open in $X_s$ with respect to $T_s$.
                $f$ is injective by \cref{bijections}.
                Show that $\preimg{f}{\img{f}{U}} = U$.
                \begin{subproof}
                    Show that $\preimg{f}{\img{f}{U}}\subseteq U$.
                    \begin{subproof}
                        Show that for all $x$ such that $x\in\preimg{f}{\img{f}{U}}$ we have $x\in U$.
                        \begin{subproof}
                            Fix $x$.
                            Assume $x\in\preimg{f}{\img{f}{U}}$.
                            Take $z\in\img{f}{U}$ such that $x\mathrel{f} z$ by \cref{preimg_iff}.
                            Take $u\in U$ such that $u\mathrel{f} z$ by \cref{img_iff}.
                            Thus $x = u$ by \cref{injective}.
                            Thus $x\in U$.
                        \end{subproof}
                        Thus $\preimg{f}{\img{f}{U}}\subseteq U$ by \cref{subseteq}.
                    \end{subproof}
                    Show that $U\subseteq \preimg{f}{\img{f}{U}}$.
                    \begin{subproof}
                        Show that for all $x$ such that $x\in U$ we have $x\in\preimg{f}{\img{f}{U}}$.
                        \begin{subproof}
                            Fix $x$.
                            Assume $x\in U$.
                            $T_s$ is a topology on $X_s$ by \cref{open_in}.
                            $U\in T_s$ by \cref{open_in}.
                            Thus $T_s\subseteq \pow{X_s}$ by \cref{topology_on}.
                            Thus $U\in\pow{X_s}$ by \cref{subseteq}.
                            Thus $U\subseteq X_s$ by \cref{powerset_elim}.
                            Thus $x\in X_s$ by \cref{subseteq}.
                            $\dom{f} = X_s$ by \cref{funs}.
                            Thus $x\in\dom{f}$ by \cref{subseteq}.
                            $f\in\funs{X_s}{Z}$.
                            Thus $f$ is a function by \cref{funs_is_function}.
                            Thus $x\mathrel{f} f(x)$ by \cref{function_apply_elim}.
                            Thus $f(x)\in\img{f}{U}$ by \cref{img_iff}.
                            Thus $x\in\preimg{f}{\img{f}{U}}$ by \cref{preimg_iff}.
                        \end{subproof}
                        Thus $U\subseteq \preimg{f}{\img{f}{U}}$ by \cref{subseteq}.
                    \end{subproof}
                    Thus $\preimg{f}{\img{f}{U}} = U$ by \cref{subseteq_antisymmetric}.
                \end{subproof}
                $f\in\funs{X_s}{Z}$.
                Thus $f\in\rels{X_s}{Z}$ by \cref{funs}.
                Thus $\ran{f}\subseteq Z$ by \cref{rels_ran_subseteq}.
                $\img{f}{U}\subseteq \ran{f}$ by \cref{img_subseteq_ran}.
                Thus $\img{f}{U}\subseteq Z$ by \cref{subseteq_transitive}.
                Thus $\img{f}{U}$ is open in $Z$ with respect to $T_Z$ by \cref{quotient_map}.
                Thus $f$ is a relation by \cref{funs_is_relation,funs}.
                Thus $\preimg{\converse{f}}{U} = \img{f}{U}$ by \cref{preim_eq_img_of_converse,converse_converse_eq}.
            \end{subproof}
            $\converse{f}\in\funs{Z}{X_s}$ by \cref{bijective_converse_are_funs}.
            $T_Z$ is a topology on $Z$ by \cref{continuous}.
            $T_s$ is a topology on $X_s$ by \cref{continuous}.
            Thus $\converse{f}$ is continuous from $Z$ to $X_s$ with respect to $T_Z$ and $T_s$ by \cref{continuous}.
        \end{subproof}
        Thus $f$ is a homeomorphism between $X_s$ and $Z$ with respect to $T_s$ and $T_Z$ by \cref{homeomorphism}.
    \end{subproof}
\end{proof}

% from §22 Item 22: Corollary 22.3(b) Hausdorff
\begin{lemma}\label{corollary_quotient_space_hausdorff}
    Suppose $g$ is a surjection from $X$ to $Z$.
    Suppose $g$ is continuous from $X$ to $Z$ with respect to $T_X$ and $T_Z$.
    Suppose $Z$ is Hausdorff with respect to $T_Z$.
    Let $X_s = \fiberPartition{g}{Z}$.
    Let $p = \partitionMap{X_s}{X}$.
    Let $T_s = \quotientTopology{p}{T_X}{X}{X_s}$.
    Suppose $f$ is a bijection from $X_s$ to $Z$.
    Suppose $f$ is continuous from $X_s$ to $Z$ with respect to $T_s$ and $T_Z$.
    Then $X_s$ is Hausdorff with respect to $T_s$.
\end{lemma}
\begin{proof}
    $T_X$ is a topology on $X$ by \cref{continuous}.
    $X_s$ is a partition of $X$ by \cref{fiber_partition_is_partition}.
    Thus $p$ is a surjection from $X$ to $X_s$ by \cref{partition_map_surjection}.
    Thus $p\in\funs{X}{X_s}$ by \cref{surj}.
    Thus $T_s$ is a topology on $X_s$ by \cref{quotient_topology_is_topology}.
    Show that for all $x_1,x_2$ such that $x_1\in X_s$ and $x_2\in X_s$ and $x_1\neq x_2$ there exist $U_1,U_2$ such that
        $U_1$ is a neighborhood of $x_1$ in $X_s$ with respect to $T_s$ and
        $U_2$ is a neighborhood of $x_2$ in $X_s$ with respect to $T_s$ and
        $U_1$ is disjoint from $U_2$.
    \begin{subproof}
        Fix $x_1$.
        Fix $x_2$.
        Assume $x_1\in X_s$.
        Assume $x_2\in X_s$.
        Assume $x_1\neq x_2$.
        $f$ is injective by \cref{bijections}.
        $f\in\funs{X_s}{Z}$ by \cref{bijections_to_funs}.
        Thus $f$ is a function by \cref{funs_is_function}.
        $\dom{f} = X_s$ by \cref{funs}.
        Thus $x_1\in\dom{f}$ and $x_2\in\dom{f}$ by \cref{subseteq}.
        Show that $f(x_1) \neq f(x_2)$.
        \begin{subproof}
            Suppose not.
            Thus $f(x_1) = f(x_2)$.
            Thus $x_1 = x_2$ by \cref{injective_function}.
            Contradiction.
        \end{subproof}
        Thus $f(x_1)\in Z$ and $f(x_2)\in Z$ by \cref{funs_type_apply}.
        Thus there exist $V_1,V_2$ such that
            $V_1$ is a neighborhood of $f(x_1)$ in $Z$ with respect to $T_Z$ and
            $V_2$ is a neighborhood of $f(x_2)$ in $Z$ with respect to $T_Z$ and
            $V_1$ is disjoint from $V_2$ by \cref{hausdorff}.
        Thus $V_1$ and $V_2$ are open in $Z$ with respect to $T_Z$ and
            $f(x_1)\in V_1$ and $f(x_2)\in V_2$ by \cref{neighborhood}.
        Let $U_1 = \preimg{f}{V_1}$.
        Let $U_2 = \preimg{f}{V_2}$.
        Thus $U_1$ and $U_2$ are open in $X_s$ with respect to $T_s$ by \cref{continuous}.
        Thus $x_1\in U_1$ and $x_2\in U_2$ by \cref{preimg_iff,function_apply_elim}.
        Thus $U_1$ is a neighborhood of $x_1$ in $X_s$ with respect to $T_s$ and
            $U_2$ is a neighborhood of $x_2$ in $X_s$ with respect to $T_s$ by \cref{neighborhood}.
        Show that $U_1$ is disjoint from $U_2$.
        \begin{subproof}
            Suppose not.
            Then there exists $x$ such that $x\in U_1$ and $x\in U_2$.
            Thus $x\in\preimg{f}{V_1}$ and $x\in\preimg{f}{V_2}$.
            Take $b_1\in V_1$ such that $x\mathrel{f} b_1$ by \cref{preimg_iff}.
            Take $b_2\in V_2$ such that $x\mathrel{f} b_2$ by \cref{preimg_iff}.
            $f\in\funs{X_s}{Z}$ by \cref{bijections_to_funs}.
            Thus $f$ is a function by \cref{funs_is_function}.
            $T_s$ is a topology on $X_s$ by \cref{open_in}.
            $U_1\in T_s$ by \cref{open_in}.
            Thus $T_s\subseteq \pow{X_s}$ by \cref{topology_on}.
            Thus $U_1\in\pow{X_s}$ by \cref{subseteq}.
            Thus $U_1\subseteq X_s$ by \cref{powerset_elim}.
            Thus $x\in X_s$ by \cref{subseteq}.
            $\dom{f} = X_s$ by \cref{funs}.
            Thus $x\in\dom{f}$ by \cref{subseteq}.
            Thus $f(x) = b_1$ and $f(x) = b_2$ by \cref{function_apply_iff}.
            Thus $f(x)\in V_1$ and $f(x)\in V_2$.
            Contradiction by \cref{disjoint}.
        \end{subproof}
        Thus $U_1$ is disjoint from $U_2$ by \cref{disjoint}.
    \end{subproof}
    Thus $X_s$ is Hausdorff with respect to $T_s$ by \cref{hausdorff}.
\end{proof}

% from §23 Item 1: separation
\begin{definition}\label{separation}
    $X$ is separated by $U$ and $V$ with respect to $T$ iff
        $U$ is open in $X$ with respect to $T$ and
        $V$ is open in $X$ with respect to $T$ and
        $U\neq\emptyset$ and $V\neq\emptyset$ and
        $U$ is disjoint from $V$ and
        $U\union V = X$.
\end{definition}

% from §23 Item 2: connected
\begin{definition}\label{connected}
    $X$ is connected with respect to $T$ iff
        there exists no $U$ such that there exists $V$ such that
            $X$ is separated by $U$ and $V$ with respect to $T$.
\end{definition}

% from §23 Item 3: connected iff only clopen sets are empty or whole space
\begin{lemma}\label{connected_iff_clopen}
    Assume $T$ is a topology on $X$.
    Then $X$ is connected with respect to $T$ iff
        for all $A$ such that $A\subseteq X$ we have
            if $A$ is open in $X$ with respect to $T$ and $A$ is closed in $X$ with respect to $T$ then
                $A = \emptyset$ or $A = X$.
\end{lemma}
\begin{proof}
    Show that if $X$ is connected with respect to $T$ then
        for all $A$ such that $A\subseteq X$ we have
            if $A$ is open in $X$ with respect to $T$ and $A$ is closed in $X$ with respect to $T$ then
                $A = \emptyset$ or $A = X$.
    \begin{subproof}
        Assume $X$ is connected with respect to $T$.
        Fix $A$.
        Assume $A\subseteq X$.
        Assume $A$ is open in $X$ with respect to $T$.
        Assume $A$ is closed in $X$ with respect to $T$.
        \begin{byCase}
        \caseOf{$A = \emptyset$.}
            Thus $A = \emptyset$ or $A = X$.
        \caseOf{$A\neq \emptyset$.}
            Show that $A = X$.
            \begin{subproof}
                Suppose not.
                Then $A\neq X$.
                Let $U = A$.
                Let $V = X\setminus A$.
                $U$ is open in $X$ with respect to $T$.
                $V$ is open in $X$ with respect to $T$ by \cref{closed_in}.
                $U\neq\emptyset$.
                Show that $V\neq\emptyset$.
                \begin{subproof}
                    Suppose not.
                    Then $V = \emptyset$.
                    Thus $X\setminus A = \emptyset$.
                    Thus $X\subseteq A$ by \cref{setminus_eq_emptyset_iff_subseteq}.
                    Thus $A = X$ by \cref{subseteq_antisymmetric,subseteq}.
                    Contradiction.
                \end{subproof}
                Show that $U$ is disjoint from $V$.
                \begin{subproof}
                    Show that there exists no $x$ such that $x\in U$ and $x\in V$.
                    \begin{subproof}
                        Suppose not.
                        Then there exists $x$ such that $x\in U$ and $x\in V$.
                        Thus $x\in A$ and $x\in X\setminus A$.
                        Thus $x\notin A$ by \cref{setminus_elim_right}.
                        Contradiction.
                    \end{subproof}
                    Thus $U$ is disjoint from $V$ by \cref{disjoint}.
                \end{subproof}
                $U\union V = X$ by \cref{setminus_partition}.
                Thus $X$ is separated by $U$ and $V$ with respect to $T$ by \cref{separation}.
                Thus there exists $U$ such that there exists $V$ such that
                    $X$ is separated by $U$ and $V$ with respect to $T$.
                Contradiction by \cref{connected}.
            \end{subproof}
            Thus $A = \emptyset$ or $A = X$.
        \end{byCase}
    \end{subproof}
    Show that if
        for all $A$ such that $A\subseteq X$ we have
            if $A$ is open in $X$ with respect to $T$ and $A$ is closed in $X$ with respect to $T$ then
                $A = \emptyset$ or $A = X$
        then $X$ is connected with respect to $T$.
    \begin{subproof}
        Assume that for all $A$ such that $A\subseteq X$ we have
            if $A$ is open in $X$ with respect to $T$ and $A$ is closed in $X$ with respect to $T$ then
                $A = \emptyset$ or $A = X$.
        Show that there exists no $U$ such that there exists $V$ such that
            $X$ is separated by $U$ and $V$ with respect to $T$.
        \begin{subproof}
            Suppose not.
            Then there exists $U$ such that there exists $V$ such that
                $X$ is separated by $U$ and $V$ with respect to $T$.
            Take $U$ such that there exists $V$ such that
                $X$ is separated by $U$ and $V$ with respect to $T$.
            Take $V$ such that $X$ is separated by $U$ and $V$ with respect to $T$.
            $U$ is open in $X$ with respect to $T$ by \cref{separation}.
            $V$ is open in $X$ with respect to $T$ by \cref{separation}.
            $U\neq\emptyset$ by \cref{separation}.
            $V\neq\emptyset$ by \cref{separation}.
            $U$ is disjoint from $V$ by \cref{separation}.
            $U\union V = X$ by \cref{separation}.
            Show that $U\subseteq X$.
            \begin{subproof}
                Show that for all $x$ such that $x\in U$ we have $x\in X$.
                \begin{subproof}
                    Fix $x$.
                    Assume $x\in U$.
                    Thus $x\in U\union V$ by \cref{union_intro_left}.
                    Thus $x\in X$ by \cref{subseteq}.
                \end{subproof}
                Thus $U\subseteq X$ by \cref{subseteq}.
            \end{subproof}
            Show that $X\setminus U = V$.
            \begin{subproof}
                Show that $V\subseteq X\setminus U$.
                \begin{subproof}
                    Show that for all $x$ such that $x\in V$ we have $x\in X\setminus U$.
                    \begin{subproof}
                        Fix $x$.
                        Assume $x\in V$.
                        Thus $x\in U\union V$ by \cref{union_intro_right}.
                        Thus $x\in X$ by \cref{subseteq}.
                        Show that $x\notin U$.
                        \begin{subproof}
                            Suppose not.
                            Then $x\in U$.
                            Thus there exists $x$ such that $x\in U$ and $x\in V$.
                            Contradiction by \cref{disjoint}.
                        \end{subproof}
                        Thus $x\in X\setminus U$ by \cref{setminus_intro}.
                    \end{subproof}
                    Thus $V\subseteq X\setminus U$ by \cref{subseteq}.
                \end{subproof}
                Show that $X\setminus U\subseteq V$.
                \begin{subproof}
                    Show that for all $x$ such that $x\in X\setminus U$ we have $x\in V$.
                    \begin{subproof}
                        Fix $x$.
                        Assume $x\in X\setminus U$.
                        Thus $x\in X$ by \cref{setminus_elim_left}.
                        Thus $x\in U\union V$ by \cref{subseteq}.
                        Thus $x\in U$ or $x\in V$ by \cref{union_iff}.
                        Show that $x\in V$.
                        \begin{subproof}
                            \begin{byCase}
                            \caseOf{$x\in V$.}
                                Thus $x\in V$.
                            \caseOf{$x\in U$.}
                                Thus $x\notin U$ by \cref{setminus_elim_right}.
                                Contradiction.
                            \end{byCase}
                        \end{subproof}
                        Thus $x\in V$.
                    \end{subproof}
                    Thus $X\setminus U\subseteq V$ by \cref{subseteq}.
                \end{subproof}
                Thus $X\setminus U = V$ by \cref{subseteq_antisymmetric}.
            \end{subproof}
            $V$ is open in $X$ with respect to $T$.
            Thus $X\setminus U$ is open in $X$ with respect to $T$ by \cref{open_in}.
            Thus $U$ is closed in $X$ with respect to $T$ by \cref{closed_in}.
            $U$ is open in $X$ with respect to $T$.
            Show that $U\neq X$.
            \begin{subproof}
                Take $y$ such that $y\in V$ by \cref{nonempty_inhabited}.
                Thus $y\in U\union V$ by \cref{union_intro_right}.
                Thus $y\in X$ by \cref{subseteq}.
                Show that $y\notin U$.
                \begin{subproof}
                    Suppose not.
                    Then $y\in U$.
                    Thus there exists $y$ such that $y\in U$ and $y\in V$.
                    Contradiction by \cref{disjoint}.
                \end{subproof}
                Suppose not.
                Then $U = X$.
                Thus $y\in U$.
                Contradiction.
            \end{subproof}
            Thus $U = \emptyset$ or $U = X$.
            Contradiction.
        \end{subproof}
        Thus $X$ is connected with respect to $T$ by \cref{connected}.
    \end{subproof}
\end{proof}

% from §23 Item 4: separation of a subspace
\begin{lemma}\label{separation_subspace_limit_points}
    Assume $Y$ is a subspace of $X$ with respect to $T$.
    Let $T_Y = \subspaceTopology{T}{Y}$.
    Suppose $A\subseteq Y$.
    Suppose $B\subseteq Y$.
    Then $Y$ is separated by $A$ and $B$ with respect to $T_Y$ iff
        $A\neq\emptyset$ and $B\neq\emptyset$ and
        $A$ is disjoint from $B$ and
        $A\union B = Y$ and
        $\limitPoints{A}{X}{T}$ is disjoint from $B$ and
        $\limitPoints{B}{X}{T}$ is disjoint from $A$.
\end{lemma}
\begin{proof}
    Show that if $Y$ is separated by $A$ and $B$ with respect to $T_Y$ then
        $A\neq\emptyset$ and $B\neq\emptyset$ and
        $A$ is disjoint from $B$ and
        $A\union B = Y$ and
        $\limitPoints{A}{X}{T}$ is disjoint from $B$ and
        $\limitPoints{B}{X}{T}$ is disjoint from $A$.
    \begin{subproof}
        Assume $Y$ is separated by $A$ and $B$ with respect to $T_Y$.
        $A$ is open in $Y$ with respect to $T_Y$ by \cref{separation}.
        $B$ is open in $Y$ with respect to $T_Y$ by \cref{separation}.
        $A\neq\emptyset$ by \cref{separation}.
        $B\neq\emptyset$ by \cref{separation}.
        $A$ is disjoint from $B$ by \cref{separation}.
        $A\union B = Y$ by \cref{separation}.
        $T$ is a topology on $X$ by \cref{subspace_of}.
        $Y\subseteq X$ by \cref{subspace_of}.
        Thus $A\subseteq X$ by \cref{subseteq_transitive}.
        Thus $B\subseteq X$ by \cref{subseteq_transitive}.
        $T_Y$ is a topology on $Y$ by \cref{subspace_topology_is_topology,subspace_of}.
        Show that $B = Y\setminus A$.
        \begin{subproof}
            Show that $B\subseteq Y\setminus A$.
            \begin{subproof}
                Show that for all $x$ such that $x\in B$ we have $x\in Y\setminus A$.
                \begin{subproof}
                    Fix $x$.
                    Assume $x\in B$.
                    Thus $x\in Y$ by \cref{subseteq}.
                    Show that $x\notin A$.
                    \begin{subproof}
                        Suppose not.
                        Then $x\in A$.
                        Thus there exists $x$ such that $x\in A$ and $x\in B$.
                        Contradiction by \cref{disjoint}.
                    \end{subproof}
                    Thus $x\in Y\setminus A$ by \cref{setminus_intro}.
                \end{subproof}
                Thus $B\subseteq Y\setminus A$ by \cref{subseteq}.
            \end{subproof}
            Show that $Y\setminus A\subseteq B$.
            \begin{subproof}
                Show that for all $x$ such that $x\in Y\setminus A$ we have $x\in B$.
                \begin{subproof}
                    Fix $x$.
                    Assume $x\in Y\setminus A$.
                    Thus $x\in Y$ by \cref{setminus_elim_left}.
                    Thus $x\in A\union B$ by \cref{subseteq}.
                    Thus $x\in A$ or $x\in B$ by \cref{union_iff}.
                    Show that $x\in B$.
                    \begin{subproof}
                        Suppose not.
                        Then $x\in A$.
                        Thus $x\notin A$ by \cref{setminus_elim_right}.
                        Contradiction.
                    \end{subproof}
                    Thus $x\in B$.
                \end{subproof}
                Thus $Y\setminus A\subseteq B$ by \cref{subseteq}.
            \end{subproof}
            Thus $B = Y\setminus A$ by \cref{subseteq_antisymmetric}.
        \end{subproof}
        $B$ is open in $Y$ with respect to $T_Y$.
        Thus $Y\setminus A$ is open in $Y$ with respect to $T_Y$ by \cref{open_in}.
        Thus $A$ is closed in $Y$ with respect to $T_Y$ by \cref{closed_in}.
        $A\subseteq Y$.
        Thus $Y\setminus (Y\setminus A) = A$ by \cref{double_complement,subseteq_refl}.
        Thus $A = Y\setminus B$.
        $A$ is open in $Y$ with respect to $T_Y$.
        Thus $Y\setminus B$ is open in $Y$ with respect to $T_Y$ by \cref{open_in}.
        Thus $B$ is closed in $Y$ with respect to $T_Y$ by \cref{closed_in}.
        Thus $\closure{A}{Y}{T_Y} = A$ by \cref{closed_eq_closure}.
        Thus $\closure{B}{Y}{T_Y} = B$ by \cref{closed_eq_closure}.
        Thus $\closure{A}{X}{T}\inter Y = A$ by \cref{closure_in_subspace,subspace_of}.
        Thus $\closure{B}{X}{T}\inter Y = B$ by \cref{closure_in_subspace,subspace_of}.
        Show that $\limitPoints{A}{X}{T}$ is disjoint from $B$.
        \begin{subproof}
            Show that there exists no $x$ such that $x\in\limitPoints{A}{X}{T}$ and $x\in B$.
            \begin{subproof}
                Suppose not.
                Then there exists $x$ such that $x\in\limitPoints{A}{X}{T}$ and $x\in B$.
                Thus $x\in B$.
                Thus $x\in Y$ by \cref{subseteq}.
                $\closure{A}{X}{T} = A\union\limitPoints{A}{X}{T}$ by \cref{closure_union_limit_points}.
                Thus $x\in A\union\limitPoints{A}{X}{T}$ by \cref{union_intro_right}.
                Thus $x\in\closure{A}{X}{T}$ by \cref{subseteq}.
                Thus $x\in\closure{A}{X}{T}\inter Y$ by \cref{inter_intro}.
                Thus $x\in A$ by \cref{subseteq}.
                Thus there exists $x$ such that $x\in A$ and $x\in B$.
                Contradiction by \cref{disjoint}.
            \end{subproof}
            Thus $\limitPoints{A}{X}{T}$ is disjoint from $B$ by \cref{disjoint}.
        \end{subproof}
        Show that $\limitPoints{B}{X}{T}$ is disjoint from $A$.
        \begin{subproof}
            Show that there exists no $x$ such that $x\in\limitPoints{B}{X}{T}$ and $x\in A$.
            \begin{subproof}
                Suppose not.
                Then there exists $x$ such that $x\in\limitPoints{B}{X}{T}$ and $x\in A$.
                Thus $x\in A$.
                Thus $x\in Y$ by \cref{subseteq}.
                $\closure{B}{X}{T} = B\union\limitPoints{B}{X}{T}$ by \cref{closure_union_limit_points}.
                Thus $x\in B\union\limitPoints{B}{X}{T}$ by \cref{union_intro_right}.
                Thus $x\in\closure{B}{X}{T}$ by \cref{subseteq}.
                Thus $x\in\closure{B}{X}{T}\inter Y$ by \cref{inter_intro}.
                Thus $x\in B$ by \cref{subseteq}.
                Thus there exists $x$ such that $x\in A$ and $x\in B$.
                Contradiction by \cref{disjoint}.
            \end{subproof}
            Thus $\limitPoints{B}{X}{T}$ is disjoint from $A$ by \cref{disjoint}.
        \end{subproof}
    \end{subproof}
    Show that if
        $A\neq\emptyset$ and $B\neq\emptyset$ and
        $A$ is disjoint from $B$ and
        $A\union B = Y$ and
        $\limitPoints{A}{X}{T}$ is disjoint from $B$ and
        $\limitPoints{B}{X}{T}$ is disjoint from $A$
        then $Y$ is separated by $A$ and $B$ with respect to $T_Y$.
    \begin{subproof}
        Assume $A\neq\emptyset$.
        Assume $B\neq\emptyset$.
        Assume $A$ is disjoint from $B$.
        Assume $A\union B = Y$.
        Assume $\limitPoints{A}{X}{T}$ is disjoint from $B$.
        Assume $\limitPoints{B}{X}{T}$ is disjoint from $A$.
        $T$ is a topology on $X$ by \cref{subspace_of}.
        $Y\subseteq X$ by \cref{subspace_of}.
        Thus $A\subseteq X$ by \cref{subseteq_transitive}.
        Thus $B\subseteq X$ by \cref{subseteq_transitive}.
        $T_Y$ is a topology on $Y$ by \cref{subspace_topology_is_topology,subspace_of}.
        Show that $B = Y\setminus A$.
        \begin{subproof}
            Show that $B\subseteq Y\setminus A$.
            \begin{subproof}
                Show that for all $x$ such that $x\in B$ we have $x\in Y\setminus A$.
                \begin{subproof}
                    Fix $x$.
                    Assume $x\in B$.
                    Thus $x\in Y$ by \cref{subseteq}.
                    Show that $x\notin A$.
                    \begin{subproof}
                        Suppose not.
                        Then $x\in A$.
                        Thus there exists $x$ such that $x\in A$ and $x\in B$.
                        Contradiction by \cref{disjoint}.
                    \end{subproof}
                    Thus $x\in Y\setminus A$ by \cref{setminus_intro}.
                \end{subproof}
                Thus $B\subseteq Y\setminus A$ by \cref{subseteq}.
            \end{subproof}
            Show that $Y\setminus A\subseteq B$.
            \begin{subproof}
                Show that for all $x$ such that $x\in Y\setminus A$ we have $x\in B$.
                \begin{subproof}
                    Fix $x$.
                    Assume $x\in Y\setminus A$.
                    Thus $x\in Y$ by \cref{setminus_elim_left}.
                    Thus $x\in A\union B$ by \cref{subseteq}.
                    Thus $x\in A$ or $x\in B$ by \cref{union_iff}.
                    Show that $x\in B$.
                    \begin{subproof}
                        Suppose not.
                        Then $x\in A$.
                        Thus $x\notin A$ by \cref{setminus_elim_right}.
                        Contradiction.
                    \end{subproof}
                    Thus $x\in B$.
                \end{subproof}
                Thus $Y\setminus A\subseteq B$ by \cref{subseteq}.
            \end{subproof}
            Thus $B = Y\setminus A$ by \cref{subseteq_antisymmetric}.
        \end{subproof}
        $A\subseteq Y$.
        Thus $Y\setminus (Y\setminus A) = A$ by \cref{double_complement,subseteq_refl}.
        Thus $A = Y\setminus B$.
        Show that $\closure{A}{X}{T}\inter Y\subseteq A$.
        \begin{subproof}
            Show that for all $x$ such that $x\in\closure{A}{X}{T}\inter Y$ we have $x\in A$.
            \begin{subproof}
                Fix $x$.
                Assume $x\in\closure{A}{X}{T}\inter Y$.
                Thus $x\in\closure{A}{X}{T}$ and $x\in Y$ by \cref{inter}.
                $\closure{A}{X}{T} = A\union\limitPoints{A}{X}{T}$ by \cref{closure_union_limit_points}.
                Thus $x\in A\union\limitPoints{A}{X}{T}$ by \cref{subseteq}.
                Thus $x\in A$ or $x\in\limitPoints{A}{X}{T}$ by \cref{union_iff}.
                \begin{byCase}
                \caseOf{$x\in A$.}
                    Thus $x\in A$.
                \caseOf{$x\in\limitPoints{A}{X}{T}$.}
                    Thus $x\in A$ or $x\in B$ by \cref{union_iff}.
                    Show that $x\in A$.
                    \begin{subproof}
                        Suppose not.
                        Then $x\in B$.
                        Thus there exists $x$ such that $x\in\limitPoints{A}{X}{T}$ and $x\in B$.
                        Contradiction by \cref{disjoint}.
                    \end{subproof}
                    Thus $x\in A$.
                \end{byCase}
            \end{subproof}
            Thus $\closure{A}{X}{T}\inter Y\subseteq A$ by \cref{subseteq}.
        \end{subproof}
        $A\subseteq \closure{A}{X}{T}$ by \cref{subseteq_closure}.
        Thus $A\subseteq \closure{A}{X}{T}\inter Y$ by \cref{subseteq_inter_iff}.
        Thus $\closure{A}{X}{T}\inter Y = A$ by \cref{subseteq_antisymmetric}.
        Thus $\closure{A}{Y}{T_Y} = A$ by \cref{closure_in_subspace,subspace_of}.
        $\closure{A}{Y}{T_Y}$ is closed in $Y$ with respect to $T_Y$ by \cref{closure_closed_in}.
        Thus $Y\setminus\closure{A}{Y}{T_Y}$ is open in $Y$ with respect to $T_Y$ by \cref{closed_in}.
        Thus $Y\setminus A$ is open in $Y$ with respect to $T_Y$ by \cref{open_in}.
        Thus $A$ is closed in $Y$ with respect to $T_Y$ by \cref{closed_in}.
        Show that $\closure{B}{X}{T}\inter Y\subseteq B$.
        \begin{subproof}
            Show that for all $x$ such that $x\in\closure{B}{X}{T}\inter Y$ we have $x\in B$.
            \begin{subproof}
                Fix $x$.
                Assume $x\in\closure{B}{X}{T}\inter Y$.
                Thus $x\in\closure{B}{X}{T}$ and $x\in Y$ by \cref{inter}.
                $\closure{B}{X}{T} = B\union\limitPoints{B}{X}{T}$ by \cref{closure_union_limit_points}.
                Thus $x\in B\union\limitPoints{B}{X}{T}$ by \cref{subseteq}.
                Thus $x\in B$ or $x\in\limitPoints{B}{X}{T}$ by \cref{union_iff}.
                \begin{byCase}
                \caseOf{$x\in B$.}
                    Thus $x\in B$.
                \caseOf{$x\in\limitPoints{B}{X}{T}$.}
                    Thus $x\in A$ or $x\in B$ by \cref{union_iff}.
                    Show that $x\in B$.
                    \begin{subproof}
                        Suppose not.
                        Then $x\in A$.
                        Thus there exists $x$ such that $x\in\limitPoints{B}{X}{T}$ and $x\in A$.
                        Contradiction by \cref{disjoint}.
                    \end{subproof}
                    Thus $x\in B$.
                \end{byCase}
            \end{subproof}
            Thus $\closure{B}{X}{T}\inter Y\subseteq B$ by \cref{subseteq}.
        \end{subproof}
        $B\subseteq \closure{B}{X}{T}$ by \cref{subseteq_closure}.
        Thus $B\subseteq \closure{B}{X}{T}\inter Y$ by \cref{subseteq_inter_iff}.
        Thus $\closure{B}{X}{T}\inter Y = B$ by \cref{subseteq_antisymmetric}.
        Thus $\closure{B}{Y}{T_Y} = B$ by \cref{closure_in_subspace,subspace_of}.
        $\closure{B}{Y}{T_Y}$ is closed in $Y$ with respect to $T_Y$ by \cref{closure_closed_in}.
        Thus $Y\setminus\closure{B}{Y}{T_Y}$ is open in $Y$ with respect to $T_Y$ by \cref{closed_in}.
        Thus $Y\setminus B$ is open in $Y$ with respect to $T_Y$ by \cref{open_in}.
        Thus $B$ is closed in $Y$ with respect to $T_Y$ by \cref{closed_in}.
        $Y\setminus A$ is open in $Y$ with respect to $T_Y$ by \cref{closed_in}.
        $Y\setminus B$ is open in $Y$ with respect to $T_Y$ by \cref{closed_in}.
        Thus $A$ is open in $Y$ with respect to $T_Y$ by \cref{open_in}.
        Thus $B$ is open in $Y$ with respect to $T_Y$ by \cref{open_in}.
        Thus $Y$ is separated by $A$ and $B$ with respect to $T_Y$ by \cref{separation}.
    \end{subproof}
\end{proof}

% from §23 Item 5: connected subspace lies in a separation set
\begin{lemma}\label{connected_subspace_in_separation}
    Assume $X$ is separated by $C$ and $D$ with respect to $T$.
    Assume $Y$ is a subspace of $X$ with respect to $T$.
    Let $T_Y = \subspaceTopology{T}{Y}$.
    Assume $Y$ is connected with respect to $T_Y$.
    Then $Y\subseteq C$ or $Y\subseteq D$.
\end{lemma}
\begin{proof}
    $C$ is open in $X$ with respect to $T$ by \cref{separation}.
    $D$ is open in $X$ with respect to $T$ by \cref{separation}.
    $C$ is disjoint from $D$ by \cref{separation}.
    $C\union D = X$ by \cref{separation}.
    $T$ is a topology on $X$ by \cref{subspace_of}.
    $Y\subseteq X$ by \cref{subspace_of}.
    Thus $T_Y$ is a topology on $Y$ by \cref{subspace_topology_is_topology}.
    Show that $C\inter Y$ is open in $Y$ with respect to $T_Y$.
    \begin{subproof}
        $C\in T$ by \cref{open_in}.
        $C\inter Y = Y\inter C$ by \cref{inter_comm}.
        Thus $C\inter Y\in T_Y$ by \cref{subspace_topology}.
        Thus $C\inter Y$ is open in $Y$ with respect to $T_Y$ by \cref{open_in}.
    \end{subproof}
    Show that $D\inter Y$ is open in $Y$ with respect to $T_Y$.
    \begin{subproof}
        $D\in T$ by \cref{open_in}.
        $D\inter Y = Y\inter D$ by \cref{inter_comm}.
        Thus $D\inter Y\in T_Y$ by \cref{subspace_topology}.
        Thus $D\inter Y$ is open in $Y$ with respect to $T_Y$ by \cref{open_in}.
    \end{subproof}
    Show that $C\inter Y$ is disjoint from $D\inter Y$.
    \begin{subproof}
        Show that there exists no $x$ such that $x\in C\inter Y$ and $x\in D\inter Y$.
        \begin{subproof}
            Suppose not.
            Then there exists $x$ such that $x\in C\inter Y$ and $x\in D\inter Y$.
            Thus $x\in C$ and $x\in D$ by \cref{inter}.
            Thus there exists $x$ such that $x\in C$ and $x\in D$.
            Contradiction by \cref{disjoint}.
        \end{subproof}
        Thus $C\inter Y$ is disjoint from $D\inter Y$ by \cref{disjoint}.
    \end{subproof}
    Show that $(C\inter Y)\union(D\inter Y) = Y$.
    \begin{subproof}
        Show that $(C\inter Y)\union(D\inter Y)\subseteq Y$.
        \begin{subproof}
            Show that for all $x$ such that $x\in (C\inter Y)\union(D\inter Y)$ we have $x\in Y$.
            \begin{subproof}
                Fix $x$.
                Assume $x\in (C\inter Y)\union(D\inter Y)$.
                Thus $x\in C\inter Y$ or $x\in D\inter Y$ by \cref{union_iff}.
                \begin{byCase}
                \caseOf{$x\in C\inter Y$.}
                    Thus $x\in Y$ by \cref{inter}.
                \caseOf{$x\in D\inter Y$.}
                    Thus $x\in Y$ by \cref{inter}.
                \end{byCase}
            \end{subproof}
            Thus $(C\inter Y)\union(D\inter Y)\subseteq Y$ by \cref{subseteq}.
        \end{subproof}
        Show that $Y\subseteq (C\inter Y)\union(D\inter Y)$.
        \begin{subproof}
            Show that for all $x$ such that $x\in Y$ we have $x\in (C\inter Y)\union(D\inter Y)$.
            \begin{subproof}
                Fix $x$.
                Assume $x\in Y$.
                Thus $x\in X$ by \cref{subseteq}.
                Thus $x\in C\union D$ by \cref{subseteq}.
                Thus $x\in C$ or $x\in D$ by \cref{union_iff}.
                \begin{byCase}
                \caseOf{$x\in C$.}
                    Thus $x\in C\inter Y$ by \cref{inter_intro}.
                    Thus $x\in (C\inter Y)\union(D\inter Y)$ by \cref{union_intro_left}.
                \caseOf{$x\in D$.}
                    Thus $x\in D\inter Y$ by \cref{inter_intro}.
                    Thus $x\in (C\inter Y)\union(D\inter Y)$ by \cref{union_intro_right}.
                \end{byCase}
            \end{subproof}
            Thus $Y\subseteq (C\inter Y)\union(D\inter Y)$ by \cref{subseteq}.
        \end{subproof}
        Thus $(C\inter Y)\union(D\inter Y) = Y$ by \cref{subseteq_antisymmetric}.
    \end{subproof}
    Show that $Y\subseteq C$ or $Y\subseteq D$.
    \begin{subproof}
        \begin{byCase}
        \caseOf{$C\inter Y = \emptyset$.}
            $(C\inter Y)\union(D\inter Y) = (\emptyset)\union(D\inter Y)$.
            $(\emptyset)\union(D\inter Y) = (D\inter Y)\union\emptyset$ by \cref{union_comm}.
            Thus $(C\inter Y)\union(D\inter Y) = D\inter Y$ by \cref{union_emptyset}.
            Thus $Y = D\inter Y$.
            Thus $Y\subseteq D$ by \cref{inter_lower_left,subseteq}.
        \caseOf{$C\inter Y\neq \emptyset$.}
            Show that $D\inter Y = \emptyset$.
            \begin{subproof}
                Suppose not.
                Thus $(C\inter Y)\neq\emptyset$ and $(D\inter Y)\neq\emptyset$.
                Thus $Y$ is separated by $C\inter Y$ and $D\inter Y$ with respect to $T_Y$ by \cref{separation}.
                Thus there exists $U$ such that there exists $V$ such that
                    $Y$ is separated by $U$ and $V$ with respect to $T_Y$.
                Contradiction by \cref{connected}.
            \end{subproof}
            $(C\inter Y)\union(D\inter Y) = (C\inter Y)\union\emptyset$.
            Thus $(C\inter Y)\union(D\inter Y) = C\inter Y$ by \cref{union_emptyset}.
            Thus $Y = C\inter Y$.
            Thus $Y\subseteq C$ by \cref{inter_lower_left,subseteq}.
        \end{byCase}
    \end{subproof}
\end{proof}

% from §23 helper lemma 1: subspace topology of a subspace
\begin{lemma}\label{subspace_subspace_topology}
    Assume $T$ is a topology on $X$.
    Suppose $A\subseteq Y$.
    Suppose $Y\subseteq X$.
    Let $T_Y = \subspaceTopology{T}{Y}$.
    Then $\subspaceTopology{T_Y}{A} = \subspaceTopology{T}{A}$.
\end{lemma}
\begin{proof}
    Show that $\subspaceTopology{T_Y}{A}\subseteq \subspaceTopology{T}{A}$.
    \begin{subproof}
        Show that for all $U$ such that $U\in\subspaceTopology{T_Y}{A}$ we have $U\in\subspaceTopology{T}{A}$.
        \begin{subproof}
            Fix $U$.
            Assume $U\in\subspaceTopology{T_Y}{A}$.
            Take $V\in T_Y$ such that $U = A\inter V$ by \cref{subspace_topology}.
            Take $W\in T$ such that $V = Y\inter W$ by \cref{subspace_topology}.
            $A\inter V = A\inter (Y\inter W)$.
            $A\inter (Y\inter W) = (A\inter Y)\inter W$ by \cref{inter_assoc}.
            $A\inter Y = A$ by \cref{inter_absorb_supseteq_right,subseteq}.
            Thus $U = A\inter W$.
            Thus $U\in\subspaceTopology{T}{A}$ by \cref{subspace_topology}.
        \end{subproof}
        Thus $\subspaceTopology{T_Y}{A}\subseteq \subspaceTopology{T}{A}$ by \cref{subseteq}.
    \end{subproof}
    Show that $\subspaceTopology{T}{A}\subseteq \subspaceTopology{T_Y}{A}$.
    \begin{subproof}
        Show that for all $U$ such that $U\in\subspaceTopology{T}{A}$ we have $U\in\subspaceTopology{T_Y}{A}$.
        \begin{subproof}
            Fix $U$.
            Assume $U\in\subspaceTopology{T}{A}$.
            Take $W\in T$ such that $U = A\inter W$ by \cref{subspace_topology}.
            Let $V = Y\inter W$.
            Thus $V\in T_Y$ by \cref{subspace_topology}.
            $A\inter W = A\inter (Y\inter W)$.
            $A\inter (Y\inter W) = (A\inter Y)\inter W$ by \cref{inter_assoc}.
            $A\inter Y = A$ by \cref{inter_absorb_supseteq_right,subseteq}.
            Thus $U = A\inter V$.
            Thus $U\in\subspaceTopology{T_Y}{A}$ by \cref{subspace_topology}.
        \end{subproof}
        Thus $\subspaceTopology{T}{A}\subseteq \subspaceTopology{T_Y}{A}$ by \cref{subseteq}.
    \end{subproof}
    Thus $\subspaceTopology{T_Y}{A} = \subspaceTopology{T}{A}$ by \cref{subseteq_antisymmetric}.
\end{proof}

% from §23 Item 6: union of connected subspaces with a common point
\begin{theorem}\label{union_connected_common_point}
    Assume $M\neq\emptyset$.
    Assume $T$ is a topology on $X$.
    Assume for all $A$ such that $A\in M$ we have $A\subseteq X$.
    Assume for all $A$ such that $A\in M$ we have
        $A$ is connected with respect to $\subspaceTopology{T}{A}$.
    Assume $\inters{M}\neq\emptyset$.
    Then $\unions{M}$ is connected with respect to $\subspaceTopology{T}{\unions{M}}$.
\end{theorem}
\begin{proof}
    Let $Y = \unions{M}$.
    Let $T_Y = \subspaceTopology{T}{Y}$.
    Show that $Y\subseteq X$.
    \begin{subproof}
        Show that for all $x$ such that $x\in Y$ we have $x\in X$.
        \begin{subproof}
            Fix $x$.
            Assume $x\in Y$.
            Take $A\in M$ such that $x\in A$ by \cref{unions_iff}.
            Thus $A\subseteq X$ by \cref{subseteq}.
            Thus $x\in X$ by \cref{subseteq}.
        \end{subproof}
        Thus $Y\subseteq X$ by \cref{subseteq}.
    \end{subproof}
    Thus $Y$ is a subspace of $X$ with respect to $T$ by \cref{subspace_of}.
    Thus $T_Y$ is a topology on $Y$ by \cref{subspace_topology_is_topology,subspace_of}.
    Show that there exists no $C$ such that there exists $D$ such that
        $Y$ is separated by $C$ and $D$ with respect to $T_Y$.
    \begin{subproof}
        Suppose not.
        Then there exists $C$ such that there exists $D$ such that
            $Y$ is separated by $C$ and $D$ with respect to $T_Y$.
        Take $C$ such that there exists $D$ such that
            $Y$ is separated by $C$ and $D$ with respect to $T_Y$.
        Take $D$ such that $Y$ is separated by $C$ and $D$ with respect to $T_Y$.
        $C$ is open in $Y$ with respect to $T_Y$ by \cref{separation}.
        $D$ is open in $Y$ with respect to $T_Y$ by \cref{separation}.
        $C\neq\emptyset$ by \cref{separation}.
        $D\neq\emptyset$ by \cref{separation}.
        $C$ is disjoint from $D$ by \cref{separation}.
        $C\union D = Y$ by \cref{separation}.
        Take $p$ such that $p\in\inters{M}$ by \cref{nonempty_inhabited}.
        Take $A_0$ such that $A_0\in M$ by \cref{nonempty_inhabited}.
        Thus $p\in A_0$ by \cref{inters_iff_forall}.
        Thus $p\in Y$ by \cref{unions_iff}.
        Thus $p\in C$ or $p\in D$ by \cref{union_iff}.
        \begin{byCase}
        \caseOf{$p\in C$.}
            Show that for all $A$ such that $A\in M$ we have $A\subseteq C$.
            \begin{subproof}
                Fix $A$.
                Assume $A\in M$.
                Thus $A\subseteq X$ by \cref{subseteq}.
                Show that $A\subseteq Y$.
                \begin{subproof}
                    Show that for all $x$ such that $x\in A$ we have $x\in Y$.
                    \begin{subproof}
                        Fix $x$.
                        Assume $x\in A$.
                        Thus $x\in Y$ by \cref{unions_iff}.
                    \end{subproof}
                    Thus $A\subseteq Y$ by \cref{subseteq}.
                \end{subproof}
                Let $T_A = \subspaceTopology{T}{A}$.
                Let $T_{YA} = \subspaceTopology{T_Y}{A}$.
                $A$ is connected with respect to $T_A$.
                $T_{YA} = T_A$ by \cref{subspace_subspace_topology,subspace_of}.
                Thus $A$ is connected with respect to $T_{YA}$.
                $A$ is a subspace of $Y$ with respect to $T_Y$ by \cref{subspace_of}.
                Thus $A\subseteq C$ or $A\subseteq D$ by \cref{connected_subspace_in_separation}.
                Show that $A\subseteq C$.
                \begin{subproof}
                    Suppose not.
                    Then $A\subseteq D$.
                    Thus $p\in A$ by \cref{inters_iff_forall}.
                    Thus $p\in D$ by \cref{subseteq}.
                    Thus there exists $p$ such that $p\in C$ and $p\in D$.
                    Contradiction by \cref{disjoint}.
                \end{subproof}
            \end{subproof}
            Show that $Y\subseteq C$.
            \begin{subproof}
                Show that for all $x$ such that $x\in Y$ we have $x\in C$.
                \begin{subproof}
                    Fix $x$.
                    Assume $x\in Y$.
                    Take $A\in M$ such that $x\in A$ by \cref{unions_iff}.
                    Thus $A\subseteq C$ by \cref{subseteq}.
                    Thus $x\in C$ by \cref{subseteq}.
                \end{subproof}
                Thus $Y\subseteq C$ by \cref{subseteq}.
            \end{subproof}
            $C\in T_Y$ by \cref{open_in}.
            Thus $T_Y\subseteq\pow{Y}$ by \cref{topology_on}.
            Thus $C\in\pow{Y}$ by \cref{subseteq}.
            Thus $C\subseteq Y$ by \cref{powerset_elim}.
            Thus $C = Y$ by \cref{subseteq_antisymmetric}.
            $D\in T_Y$ by \cref{open_in}.
            Thus $D\in\pow{Y}$ by \cref{subseteq}.
            Thus $D\subseteq Y$ by \cref{powerset_elim}.
            Show that $D = \emptyset$.
            \begin{subproof}
                Suppose not.
                Take $d$ such that $d\in D$ by \cref{nonempty_inhabited}.
                Thus $d\in Y$ by \cref{subseteq}.
                Thus $d\in C$ by \cref{subseteq}.
                Thus there exists $d$ such that $d\in C$ and $d\in D$.
                Contradiction by \cref{disjoint}.
            \end{subproof}
            Contradiction.
        \caseOf{$p\in D$.}
            Show that for all $A$ such that $A\in M$ we have $A\subseteq D$.
            \begin{subproof}
                Fix $A$.
                Assume $A\in M$.
                Thus $A\subseteq X$ by \cref{subseteq}.
                Show that $A\subseteq Y$.
                \begin{subproof}
                    Show that for all $x$ such that $x\in A$ we have $x\in Y$.
                    \begin{subproof}
                        Fix $x$.
                        Assume $x\in A$.
                        Thus $x\in Y$ by \cref{unions_iff}.
                    \end{subproof}
                    Thus $A\subseteq Y$ by \cref{subseteq}.
                \end{subproof}
                Let $T_A = \subspaceTopology{T}{A}$.
                Let $T_{YA} = \subspaceTopology{T_Y}{A}$.
                $A$ is connected with respect to $T_A$.
                $T_{YA} = T_A$ by \cref{subspace_subspace_topology,subspace_of}.
                Thus $A$ is connected with respect to $T_{YA}$.
                $A$ is a subspace of $Y$ with respect to $T_Y$ by \cref{subspace_of}.
                Thus $A\subseteq C$ or $A\subseteq D$ by \cref{connected_subspace_in_separation}.
                Show that $A\subseteq D$.
                \begin{subproof}
                    Suppose not.
                    Then $A\subseteq C$.
                    Thus $p\in A$ by \cref{inters_iff_forall}.
                    Thus $p\in C$ by \cref{subseteq}.
                    Thus there exists $p$ such that $p\in C$ and $p\in D$.
                    Contradiction by \cref{disjoint}.
                \end{subproof}
            \end{subproof}
            Show that $Y\subseteq D$.
            \begin{subproof}
                Show that for all $x$ such that $x\in Y$ we have $x\in D$.
                \begin{subproof}
                    Fix $x$.
                    Assume $x\in Y$.
                    Take $A\in M$ such that $x\in A$ by \cref{unions_iff}.
                    Thus $A\subseteq D$ by \cref{subseteq}.
                    Thus $x\in D$ by \cref{subseteq}.
                \end{subproof}
                Thus $Y\subseteq D$ by \cref{subseteq}.
            \end{subproof}
            $D\in T_Y$ by \cref{open_in}.
            Thus $T_Y\subseteq\pow{Y}$ by \cref{topology_on}.
            Thus $D\in\pow{Y}$ by \cref{subseteq}.
            Thus $D\subseteq Y$ by \cref{powerset_elim}.
            Thus $D = Y$ by \cref{subseteq_antisymmetric}.
            $C\in T_Y$ by \cref{open_in}.
            Thus $C\in\pow{Y}$ by \cref{subseteq}.
            Thus $C\subseteq Y$ by \cref{powerset_elim}.
            Show that $C = \emptyset$.
            \begin{subproof}
                Suppose not.
                Take $c$ such that $c\in C$ by \cref{nonempty_inhabited}.
                Thus $c\in Y$ by \cref{subseteq}.
                Thus $c\in D$ by \cref{subseteq}.
                Thus there exists $c$ such that $c\in C$ and $c\in D$.
                Contradiction by \cref{disjoint}.
            \end{subproof}
            Contradiction.
        \end{byCase}
    \end{subproof}
    Thus $Y$ is connected with respect to $T_Y$ by \cref{connected}.
\end{proof}

% from §23 Item 7: connectedness preserved under adjoining limit points
\begin{theorem}\label{connected_superset_closure}
    Assume $T$ is a topology on $X$.
    Suppose $A\subseteq X$.
    Suppose $B\subseteq X$.
    Assume $A$ is connected with respect to $\subspaceTopology{T}{A}$.
    Suppose $A\subseteq B$.
    Suppose $B\subseteq \closure{A}{X}{T}$.
    Then $B$ is connected with respect to $\subspaceTopology{T}{B}$.
\end{theorem}
\begin{proof}
    Let $T_B = \subspaceTopology{T}{B}$.
    $T$ is a topology on $X$.
    $A\subseteq X$.
    $B\subseteq X$.
    Thus $T_B$ is a topology on $B$ by \cref{subspace_topology_is_topology}.
    Show that there exists no $C$ such that there exists $D$ such that
        $B$ is separated by $C$ and $D$ with respect to $T_B$.
    \begin{subproof}
        Suppose not.
        Then there exists $C$ such that there exists $D$ such that
            $B$ is separated by $C$ and $D$ with respect to $T_B$.
        Take $C$ such that there exists $D$ such that
            $B$ is separated by $C$ and $D$ with respect to $T_B$.
        Take $D$ such that $B$ is separated by $C$ and $D$ with respect to $T_B$.
        $C$ is open in $B$ with respect to $T_B$ by \cref{separation}.
        $D$ is open in $B$ with respect to $T_B$ by \cref{separation}.
        $C\neq\emptyset$ by \cref{separation}.
        $D\neq\emptyset$ by \cref{separation}.
        $C$ is disjoint from $D$ by \cref{separation}.
        $C\union D = B$ by \cref{separation}.
        Show that $C\subseteq B$.
        \begin{subproof}
            $C\in T_B$ by \cref{open_in}.
            Thus $T_B\subseteq\pow{B}$ by \cref{topology_on}.
            Thus $C\in\pow{B}$ by \cref{subseteq}.
            Thus $C\subseteq B$ by \cref{powerset_elim}.
        \end{subproof}
        Show that $D\subseteq B$.
        \begin{subproof}
            $D\in T_B$ by \cref{open_in}.
            Thus $T_B\subseteq\pow{B}$ by \cref{topology_on}.
            Thus $D\in\pow{B}$ by \cref{subseteq}.
            Thus $D\subseteq B$ by \cref{powerset_elim}.
        \end{subproof}
        $A\subseteq B$.
        Thus $A$ is a subspace of $B$ with respect to $T_B$ by \cref{subspace_of}.
        Let $T_A = \subspaceTopology{T}{A}$.
        Let $T_{BA} = \subspaceTopology{T_B}{A}$.
        $T_{BA} = T_A$ by \cref{subspace_subspace_topology,subspace_of}.
        $A$ is connected with respect to $T_A$.
        Thus $A$ is connected with respect to $T_{BA}$.
        Thus $A\subseteq C$ or $A\subseteq D$ by \cref{connected_subspace_in_separation}.
        \begin{byCase}
        \caseOf{$A\subseteq C$.}
            $B$ is a subspace of $X$ with respect to $T$ by \cref{subspace_of}.
            Thus $\limitPoints{C}{X}{T}$ is disjoint from $D$ by \cref{separation_subspace_limit_points}.
            Show that $\closure{C}{X}{T}$ is disjoint from $D$.
            \begin{subproof}
                Show that there exists no $x$ such that $x\in\closure{C}{X}{T}$ and $x\in D$.
                \begin{subproof}
                    Suppose not.
                    Then there exists $x$ such that $x\in\closure{C}{X}{T}$ and $x\in D$.
                    Thus $C\subseteq X$ by \cref{subseteq_transitive,subseteq}.
                    $\closure{C}{X}{T} = C\union\limitPoints{C}{X}{T}$ by \cref{closure_union_limit_points}.
                    Thus $x\in C$ or $x\in\limitPoints{C}{X}{T}$ by \cref{union_iff,subseteq}.
                    \begin{byCase}
                    \caseOf{$x\in C$.}
                        Thus there exists $x$ such that $x\in C$ and $x\in D$.
                        Contradiction by \cref{disjoint}.
                    \caseOf{$x\in\limitPoints{C}{X}{T}$.}
                        Thus there exists $x$ such that $x\in\limitPoints{C}{X}{T}$ and $x\in D$.
                        Contradiction by \cref{disjoint}.
                    \end{byCase}
                \end{subproof}
                Thus $\closure{C}{X}{T}$ is disjoint from $D$ by \cref{disjoint}.
            \end{subproof}
            Thus $C\subseteq X$ by \cref{subseteq_transitive,subseteq}.
            $C\subseteq \closure{C}{X}{T}$ by \cref{subseteq_closure}.
            Thus $A\subseteq \closure{C}{X}{T}$ by \cref{subseteq_transitive}.
            $\closure{C}{X}{T}$ is closed in $X$ with respect to $T$ by \cref{closure_closed_in}.
            Thus $\closure{A}{X}{T}\subseteq \closure{C}{X}{T}$ by \cref{closure_subseteq_closed}.
            Thus $B\subseteq \closure{C}{X}{T}$ by \cref{subseteq_transitive,subseteq}.
            Thus $D\subseteq \closure{C}{X}{T}$ by \cref{subseteq_transitive,subseteq}.
            Show that $D = \emptyset$.
            \begin{subproof}
                Suppose not.
                Take $d$ such that $d\in D$ by \cref{nonempty_inhabited}.
                Thus $d\in\closure{C}{X}{T}$ by \cref{subseteq}.
                Thus there exists $d$ such that $d\in\closure{C}{X}{T}$ and $d\in D$.
                Contradiction by \cref{disjoint}.
            \end{subproof}
            Contradiction by \cref{emptyset}.
        \caseOf{$A\subseteq D$.}
            $B$ is a subspace of $X$ with respect to $T$ by \cref{subspace_of}.
            Thus $\limitPoints{D}{X}{T}$ is disjoint from $C$ by \cref{separation_subspace_limit_points}.
            Show that $\closure{D}{X}{T}$ is disjoint from $C$.
            \begin{subproof}
                Show that there exists no $x$ such that $x\in\closure{D}{X}{T}$ and $x\in C$.
                \begin{subproof}
                    Suppose not.
                    Then there exists $x$ such that $x\in\closure{D}{X}{T}$ and $x\in C$.
                    Thus $D\subseteq X$ by \cref{subseteq_transitive,subseteq}.
                    $\closure{D}{X}{T} = D\union\limitPoints{D}{X}{T}$ by \cref{closure_union_limit_points}.
                    Thus $x\in D$ or $x\in\limitPoints{D}{X}{T}$ by \cref{union_iff,subseteq}.
                    \begin{byCase}
                    \caseOf{$x\in D$.}
                        Thus there exists $x$ such that $x\in C$ and $x\in D$.
                        Contradiction by \cref{disjoint}.
                    \caseOf{$x\in\limitPoints{D}{X}{T}$.}
                        Thus there exists $x$ such that $x\in\limitPoints{D}{X}{T}$ and $x\in C$.
                        Contradiction by \cref{disjoint}.
                    \end{byCase}
                \end{subproof}
                Thus $\closure{D}{X}{T}$ is disjoint from $C$ by \cref{disjoint}.
            \end{subproof}
            Thus $D\subseteq X$ by \cref{subseteq_transitive,subseteq}.
            $D\subseteq \closure{D}{X}{T}$ by \cref{subseteq_closure}.
            Thus $A\subseteq \closure{D}{X}{T}$ by \cref{subseteq_transitive}.
            $\closure{D}{X}{T}$ is closed in $X$ with respect to $T$ by \cref{closure_closed_in}.
            Thus $\closure{A}{X}{T}\subseteq \closure{D}{X}{T}$ by \cref{closure_subseteq_closed}.
            Thus $B\subseteq \closure{D}{X}{T}$ by \cref{subseteq_transitive,subseteq}.
            Thus $C\subseteq \closure{D}{X}{T}$ by \cref{subseteq_transitive,subseteq}.
            Show that $C = \emptyset$.
            \begin{subproof}
                Suppose not.
                Take $c$ such that $c\in C$ by \cref{nonempty_inhabited}.
                Thus $c\in\closure{D}{X}{T}$ by \cref{subseteq}.
                Thus there exists $c$ such that $c\in\closure{D}{X}{T}$ and $c\in C$.
                Contradiction by \cref{disjoint}.
            \end{subproof}
            Contradiction by \cref{emptyset}.
        \end{byCase}
    \end{subproof}
    Thus $B$ is connected with respect to $T_B$ by \cref{connected}.
\end{proof}

% from §23 Item 8: continuous image of connected space
\begin{theorem}\label{connected_image_continuous}
    Assume $T$ is a topology on $X$.
    Assume $T_1$ is a topology on $Y$.
    Suppose $f$ is continuous from $X$ to $Y$ with respect to $T$ and $T_1$.
    Assume $X$ is connected with respect to $T$.
    Let $Z = \img{f}{X}$.
    Let $T_Z = \subspaceTopology{T_1}{Z}$.
    Then $Z$ is connected with respect to $T_Z$.
\end{theorem}
\begin{proof}
    Show that there exists no $A$ such that there exists $B$ such that
        $Z$ is separated by $A$ and $B$ with respect to $T_Z$.
    \begin{subproof}
        Suppose not.
        Then there exists $A$ such that there exists $B$ such that
            $Z$ is separated by $A$ and $B$ with respect to $T_Z$.
        Take $A$ such that there exists $B$ such that
            $Z$ is separated by $A$ and $B$ with respect to $T_Z$.
        Take $B$ such that $Z$ is separated by $A$ and $B$ with respect to $T_Z$.
        $A$ is open in $Z$ with respect to $T_Z$ by \cref{separation}.
        $B$ is open in $Z$ with respect to $T_Z$ by \cref{separation}.
        $f\in\funs{X}{Y}$ by \cref{continuous}.
        Thus $\ran{f}\subseteq Y$ by \cref{funs_ran}.
        $\img{f}{X}\subseteq \ran{f}$ by \cref{img_subseteq_ran}.
        Thus $Z\subseteq Y$ by \cref{subseteq_transitive}.
        $Z$ is a subspace of $Y$ with respect to $T_1$ by \cref{subspace_of}.
        Thus $T_Z$ is a topology on $Z$ by \cref{subspace_topology_is_topology,subspace_of}.
        Show that $A\subseteq Z$.
        \begin{subproof}
            $A\in T_Z$ by \cref{open_in}.
            Thus $T_Z\subseteq\pow{Z}$ by \cref{topology_on}.
            Thus $A\in\pow{Z}$ by \cref{subseteq}.
            Thus $A\subseteq Z$ by \cref{powerset_elim}.
        \end{subproof}
        Show that $B\subseteq Z$.
        \begin{subproof}
            $B\in T_Z$ by \cref{open_in}.
            Thus $T_Z\subseteq\pow{Z}$ by \cref{topology_on}.
            Thus $B\in\pow{Z}$ by \cref{subseteq}.
            Thus $B\subseteq Z$ by \cref{powerset_elim}.
        \end{subproof}
        $A\neq\emptyset$ by \cref{separation}.
        $B\neq\emptyset$ by \cref{separation}.
        $A$ is disjoint from $B$ by \cref{separation}.
        $A\union B = Z$ by \cref{separation}.
        Thus $\img{f}{X}\subseteq Z$ by \cref{subseteq_refl}.
        Thus $f$ is continuous from $X$ to $Z$ with respect to $T$ and $T_Z$
            by \cref{continuous_restrict_range}.
        Thus $f\in\funs{X}{Z}$ by \cref{continuous}.
        Thus $\dom{f} = X$ by \cref{funs}.
        Show that $\preimg{f}{A}$ is open in $X$ with respect to $T$.
        \begin{subproof}
            Thus $\preimg{f}{A}$ is open in $X$ with respect to $T$ by \cref{continuous}.
        \end{subproof}
        Show that $\preimg{f}{B}$ is open in $X$ with respect to $T$.
        \begin{subproof}
            Thus $\preimg{f}{B}$ is open in $X$ with respect to $T$ by \cref{continuous}.
        \end{subproof}
        Show that $\preimg{f}{A}\neq\emptyset$.
        \begin{subproof}
            Take $a$ such that $a\in A$ by \cref{nonempty_inhabited}.
            Thus $a\in Z$ by \cref{subseteq}.
            Take $x\in X$ such that $x\mathrel{f} a$ by \cref{img_iff}.
            Thus $x\in\preimg{f}{A}$ by \cref{preimg_iff}.
            Show that $\preimg{f}{A}\neq\emptyset$.
            \begin{subproof}
                Suppose not.
                Then $\preimg{f}{A} = \emptyset$.
                Thus $x\in\emptyset$.
                Contradiction by \cref{emptyset}.
            \end{subproof}
        \end{subproof}
        Show that $\preimg{f}{B}\neq\emptyset$.
        \begin{subproof}
            Take $b$ such that $b\in B$ by \cref{nonempty_inhabited}.
            Thus $b\in Z$ by \cref{subseteq}.
            Take $x\in X$ such that $x\mathrel{f} b$ by \cref{img_iff}.
            Thus $x\in\preimg{f}{B}$ by \cref{preimg_iff}.
            Show that $\preimg{f}{B}\neq\emptyset$.
            \begin{subproof}
                Suppose not.
                Then $\preimg{f}{B} = \emptyset$.
                Thus $x\in\emptyset$.
                Contradiction by \cref{emptyset}.
            \end{subproof}
        \end{subproof}
        Show that $\preimg{f}{A}$ is disjoint from $\preimg{f}{B}$.
        \begin{subproof}
            Show that there exists no $x$ such that $x\in\preimg{f}{A}$ and $x\in\preimg{f}{B}$.
            \begin{subproof}
                Suppose not.
                Then there exists $x$ such that $x\in\preimg{f}{A}$ and $x\in\preimg{f}{B}$.
                Take $a\in A$ such that $x\mathrel{f} a$ by \cref{preimg_iff}.
                Take $b\in B$ such that $x\mathrel{f} b$ by \cref{preimg_iff}.
                $f$ is a function by \cref{funs_is_function}.
                Thus $f(x) = a$ and $f(x) = b$ by \cref{function_apply_iff}.
                Thus $a = b$.
                Thus there exists $a$ such that $a\in A$ and $a\in B$.
                Contradiction by \cref{disjoint}.
            \end{subproof}
            Thus $\preimg{f}{A}$ is disjoint from $\preimg{f}{B}$ by \cref{disjoint}.
        \end{subproof}
        Show that $\preimg{f}{A}\union\preimg{f}{B} = X$.
        \begin{subproof}
            Show that $\preimg{f}{Z} = X$.
            \begin{subproof}
                Show that for all $x$ such that $x\in X$ we have $x\in\preimg{f}{Z}$.
                \begin{subproof}
                    Fix $x$.
                    Assume $x\in X$.
                    Thus $x\in\dom{f}$ by \cref{subseteq}.
                    Thus $f(x)\in Z$ by \cref{funs_type_apply}.
                    $f$ is a function by \cref{funs_is_function}.
                    Thus $x\mathrel{f} f(x)$ by \cref{function_apply_elim}.
                    Thus $x\in\preimg{f}{Z}$ by \cref{preimg_iff}.
                \end{subproof}
                Thus $X\subseteq\preimg{f}{Z}$ by \cref{subseteq}.
                $\preimg{f}{Z}\subseteq \dom{f}$ by \cref{preimg_subseteq_dom}.
                Thus $\preimg{f}{Z}\subseteq X$ by \cref{subseteq}.
                Thus $\preimg{f}{Z} = X$ by \cref{subseteq_antisymmetric}.
            \end{subproof}
            $A\union B = Z$.
            Thus $\preimg{f}{A\union B} = X$.
            $\preimg{f}{A\union B} = \preimg{f}{A}\union\preimg{f}{B}$ by \cref{preimg_union}.
            Thus $\preimg{f}{A}\union\preimg{f}{B} = X$ by \cref{subseteq_antisymmetric,subseteq}.
        \end{subproof}
        Thus $X$ is separated by $\preimg{f}{A}$ and $\preimg{f}{B}$ with respect to $T$ by \cref{separation}.
        Thus there exists $U$ such that there exists $V$ such that
            $X$ is separated by $U$ and $V$ with respect to $T$.
        Contradiction by \cref{connected}.
    \end{subproof}
    Thus $Z$ is connected with respect to $T_Z$ by \cref{connected}.
\end{proof}

% from §23 helper lemma 2: identity is continuous
\begin{lemma}\label{identity_continuous}
    Assume $T$ is a topology on $X$.
    Let $j = \identity{X}$.
    Then $j$ is continuous from $X$ to $X$ with respect to $T$ and $T$.
\end{lemma}
\begin{proof}
    $j\in\funs{X}{X}$ by \cref{id_elem_funs}.
    Show that for all $U$ such that $U$ is open in $X$ with respect to $T$
        we have $\preimg{j}{U}$ is open in $X$ with respect to $T$.
    \begin{subproof}
        Fix $U$.
        Assume $U$ is open in $X$ with respect to $T$.
        Show that $\preimg{j}{U} = U$.
        \begin{subproof}
            Show that $\preimg{j}{U}\subseteq U$.
            \begin{subproof}
                Show that for all $x$ such that $x\in\preimg{j}{U}$ we have $x\in U$.
                \begin{subproof}
                    Fix $x$.
                    Assume $x\in\preimg{j}{U}$.
                    Take $y\in U$ such that $x\mathrel{j} y$ by \cref{preimg_iff}.
                    Thus $x = y$ and $x\in X$ by \cref{id_iff}.
                    Thus $x\in U$.
                \end{subproof}
                Thus $\preimg{j}{U}\subseteq U$ by \cref{subseteq}.
            \end{subproof}
            Show that $U\subseteq\preimg{j}{U}$.
            \begin{subproof}
                $U\in T$ by \cref{open_in}.
                $T\subseteq\pow{X}$ by \cref{topology_on}.
                Thus $U\subseteq X$ by \cref{pow_iff,subseteq}.
                Show that for all $x$ such that $x\in U$ we have $x\in\preimg{j}{U}$.
                \begin{subproof}
                    Fix $x$.
                    Assume $x\in U$.
                    Thus $x\in X$ by \cref{subseteq}.
                    Thus $x\mathrel{j} x$ by \cref{id_iff}.
                    Thus $x\in\preimg{j}{U}$ by \cref{preimg_iff}.
                \end{subproof}
                Thus $U\subseteq\preimg{j}{U}$ by \cref{subseteq}.
            \end{subproof}
            Thus $\preimg{j}{U} = U$ by \cref{subseteq_antisymmetric}.
        \end{subproof}
        Thus $\preimg{j}{U}$ is open in $X$ with respect to $T$ by \cref{open_in}.
    \end{subproof}
    Thus $j$ is continuous from $X$ to $X$ with respect to $T$ and $T$ by \cref{continuous}.
\end{proof}

% from §23 helper lemma 3: constant map is a function
\begin{lemma}\label{constant_map_function}
    Suppose $c\in B$.
    Let $f = \{(a,c) \mid a\in A\}$.
    Then $f$ is a function from $A$ to $B$.
\end{lemma}
\begin{proof}
    Show that $f$ is right-unique.
    \begin{subproof}
        Show that for all $a,b_1,b_2$ such that $(a,b_1)\in f$ and $(a,b_2)\in f$ we have $b_1 = b_2$.
        \begin{subproof}
            Fix $a$.
            Fix $b_1$.
            Fix $b_2$.
            Assume $(a,b_1)\in f$ and $(a,b_2)\in f$.
            Take $a_1$ such that $a_1\in A$ and $(a,b_1) = (a_1,c)$.
            Take $a_2$ such that $a_2\in A$ and $(a,b_2) = (a_2,c)$.
            Thus $b_1 = c$ and $b_2 = c$ by \cref{pair_eq_iff}.
            Thus $b_1 = b_2$.
        \end{subproof}
        Thus $f$ is right-unique by \cref{rightunique}.
    \end{subproof}
    Show that $f$ is a relation.
    \begin{subproof}
        Show that for all $w\in f$ there exist $x,y$ such that $w = (x,y)$.
        \begin{subproof}
            Fix $w$.
            Assume $w\in f$.
            Take $a$ such that $a\in A$ and $w = (a,c)$.
            Thus there exist $x,y$ such that $w = (x,y)$.
        \end{subproof}
        Thus $f$ is a relation by \cref{relation}.
    \end{subproof}
    Thus $f$ is a function.
    Show that for all $x\in\dom{f}$ we have $f(x)\in B$.
    \begin{subproof}
        Fix $x$.
        Assume $x\in\dom{f}$.
        Take $y$ such that $(x,y)\in f$ by \cref{dom_iff}.
        Take $a$ such that $a\in A$ and $(x,y) = (a,c)$.
        Thus $y = c$ by \cref{pair_eq_iff}.
        Thus $f(x) = y$ by \cref{function_apply_intro}.
        Thus $f(x) = c$.
        Thus $f(x)\in B$.
    \end{subproof}
    Thus $f$ is a function to $B$.
    Show that $\dom{f} = A$.
    \begin{subproof}
        Show that for all $x$ such that $x\in\dom{f}$ we have $x\in A$.
        \begin{subproof}
            Fix $x$.
            Assume $x\in\dom{f}$.
            Take $y$ such that $(x,y)\in f$ by \cref{dom_iff}.
            Take $a$ such that $a\in A$ and $(x,y) = (a,c)$.
            Thus $x = a$ by \cref{pair_eq_iff}.
            Thus $x\in A$.
        \end{subproof}
        Thus $\dom{f}\subseteq A$ by \cref{subseteq}.
        Show that for all $a$ such that $a\in A$ we have $a\in\dom{f}$.
        \begin{subproof}
            Fix $a$.
            Assume $a\in A$.
            Thus $(a,c)\in f$.
            Thus $a\in\dom{f}$ by \cref{dom_intro}.
        \end{subproof}
        Thus $A\subseteq\dom{f}$ by \cref{subseteq}.
        Thus $\dom{f} = A$ by \cref{subseteq_antisymmetric}.
    \end{subproof}
    Thus $f\in\funs{A}{B}$ by \cref{funs_intro}.
    Thus $f$ is a function from $A$ to $B$.
\end{proof}

% from §23 Item 9: finite products of connected spaces
\begin{theorem}\label{connected_product_finite}
    Assume $T_1$ is a topology on $X$.
    Assume $T_2$ is a topology on $Y$.
    Assume $X$ is connected with respect to $T_1$.
    Assume $Y$ is connected with respect to $T_2$.
    Let $T = \productTopology{T_1}{T_2}{X}{Y}$.
    Then $X\times Y$ is connected with respect to $T$.
\end{theorem}
\begin{proof}
    $T_1$ is a topology on $X$.
    $T_2$ is a topology on $Y$.
    $\productBasis{T_1}{T_2}{X}{Y}$ is a basis on $X\times Y$ by \cref{product_basis_is_basis}.
    Thus $T$ is a topology on $X\times Y$ by \cref{basis_topology_is_topology,product_topology}.
    \begin{byCase}
    \caseOf{$X = \emptyset$.}
        Thus $X\times Y = \emptyset$ by \cref{times_empty_left}.
        Show that $X\times Y$ is connected with respect to $T$.
        \begin{subproof}
            Show that there exists no $U$ such that there exists $V$ such that
                $X\times Y$ is separated by $U$ and $V$ with respect to $T$.
            \begin{subproof}
                Suppose not.
                Then there exists $U$ such that there exists $V$ such that
                    $X\times Y$ is separated by $U$ and $V$ with respect to $T$.
                Take $U$ such that there exists $V$ such that
                    $X\times Y$ is separated by $U$ and $V$ with respect to $T$.
                Take $V$ such that $X\times Y$ is separated by $U$ and $V$ with respect to $T$.
                $U$ is open in $X\times Y$ with respect to $T$ by \cref{separation}.
                $U\in T$ by \cref{open_in}.
                $T\subseteq\pow{X\times Y}$ by \cref{topology_on}.
                Thus $U\subseteq X\times Y$ by \cref{pow_iff,subseteq}.
                Thus $U\subseteq \emptyset$.
                Thus $U = \emptyset$ by \cref{subseteq_emptyset_iff}.
                Contradiction by \cref{separation}.
            \end{subproof}
            Thus $X\times Y$ is connected with respect to $T$ by \cref{connected}.
        \end{subproof}
    \caseOf{$X\neq\emptyset$.}
        \begin{byCase}
        \caseOf{$Y = \emptyset$.}
            Thus $X\times Y = \emptyset$ by \cref{times_empty_right}.
            Show that $X\times Y$ is connected with respect to $T$.
            \begin{subproof}
                Show that there exists no $U$ such that there exists $V$ such that
                    $X\times Y$ is separated by $U$ and $V$ with respect to $T$.
                \begin{subproof}
                    Suppose not.
                    Then there exists $U$ such that there exists $V$ such that
                        $X\times Y$ is separated by $U$ and $V$ with respect to $T$.
                    Take $U$ such that there exists $V$ such that
                        $X\times Y$ is separated by $U$ and $V$ with respect to $T$.
                    Take $V$ such that $X\times Y$ is separated by $U$ and $V$ with respect to $T$.
                    $U$ is open in $X\times Y$ with respect to $T$ by \cref{separation}.
                    $U\in T$ by \cref{open_in}.
                    $T\subseteq\pow{X\times Y}$ by \cref{topology_on}.
                    Thus $U\subseteq X\times Y$ by \cref{pow_iff,subseteq}.
                    Thus $U\subseteq \emptyset$.
                    Thus $U = \emptyset$ by \cref{subseteq_emptyset_iff}.
                    Contradiction by \cref{separation}.
                \end{subproof}
                Thus $X\times Y$ is connected with respect to $T$ by \cref{connected}.
            \end{subproof}
        \caseOf{$Y\neq\emptyset$.}
            Take $a$ such that $a\in X$ by \cref{nonempty_inhabited}.
            Take $b$ such that $b\in Y$ by \cref{nonempty_inhabited}.
            Let $j_X = \identity{X}$.
            Thus $j_X$ is continuous from $X$ to $X$ with respect to $T_1$ and $T_1$
                by \cref{identity_continuous}.
            Show that for all $x$ such that $x\in X$ we have $j_X(x) = x$.
            \begin{subproof}
                Fix $x$.
                Assume $x\in X$.
                Thus $x\mathrel{j_X} x$ by \cref{id_iff}.
                $j_X\in\funs{X}{X}$ by \cref{id_elem_funs}.
                Thus $j_X$ is a function by \cref{funs_is_function}.
                Thus $j_X(x) = x$ by \cref{function_apply_intro}.
            \end{subproof}
            Let $k = \{(x,b) \mid x\in X\}$.
            Thus $k$ is a function from $X$ to $Y$ by \cref{constant_map_function}.
            Show that for all $x$ such that $x\in X$ we have $k(x) = b$.
            \begin{subproof}
                Fix $x$.
                Assume $x\in X$.
                Thus $(x,b)\in k$.
                Thus $k$ is a function.
                Thus $k(x) = b$ by \cref{function_apply_intro}.
            \end{subproof}
            Thus $k$ is continuous from $X$ to $Y$ with respect to $T_1$ and $T_2$
                by \cref{constant_continuous}.
            Let $f = \{(x,(j_X(x),k(x))) \mid x\in X\}$.
            Thus $f$ is continuous from $X$ to $X\times Y$ with respect to $T_1$ and $T$
                by \cref{continuous_map_into_product_backward}.
            Let $Z = \img{f}{X}$.
            Show that $Z = X\times\{b\}$.
            \begin{subproof}
                Show that $Z\subseteq X\times\{b\}$.
                \begin{subproof}
                    Show that for all $p$ such that $p\in Z$ we have $p\in X\times\{b\}$.
                    \begin{subproof}
                        Fix $p$.
                        Assume $p\in Z$.
                        Take $x\in X$ such that $x\mathrel{f} p$ by \cref{img_iff}.
                        Take $x_0\in X$ such that $(x,p) = (x_0,(j_X(x_0),k(x_0)))$.
                        Thus $p = (j_X(x),k(x))$ by \cref{pair_eq_iff}.
                        Thus $j_X(x) = x$.
                        Thus $k(x) = b$.
                        Thus $p = (x,b)$.
                        Thus $b\in\{b\}$ by \cref{singleton_intro}.
                        Thus $p\in X\times\{b\}$ by \cref{times_tuple_intro}.
                    \end{subproof}
                    Thus $Z\subseteq X\times\{b\}$ by \cref{subseteq}.
                \end{subproof}
                Show that $X\times\{b\}\subseteq Z$.
                \begin{subproof}
                    Show that for all $p$ such that $p\in X\times\{b\}$ we have $p\in Z$.
                    \begin{subproof}
                        Fix $p$.
                        Assume $p\in X\times\{b\}$.
                        Take $x,y$ such that $x\in X$ and $y\in\{b\}$ and $p = (x,y)$
                            by \cref{times_elem_is_tuple}.
                        Thus $y = b$ by \cref{singleton_iff}.
                        Thus $p = (x,b)$ by \cref{pair_eq_iff}.
                        Thus $(x,b)\in k$.
                        Thus $k$ is a function.
                        Thus $k(x) = b$ by \cref{function_apply_intro}.
                        Thus $j_X(x) = x$.
                        Thus $(x,(j_X(x),k(x))) = (x,p)$.
                        Thus $(x,p)\in f$.
                        Thus $x\mathrel{f} p$.
                        Thus $p\in Z$ by \cref{img_iff}.
                    \end{subproof}
                    Thus $X\times\{b\}\subseteq Z$ by \cref{subseteq}.
                \end{subproof}
                Thus $Z = X\times\{b\}$ by \cref{subseteq_antisymmetric}.
            \end{subproof}
            Let $T_{Xb} = \subspaceTopology{T}{X\times\{b\}}$.
            Thus $X\times\{b\}$ is connected with respect to $T_{Xb}$ by \cref{connected_image_continuous}.
            Show that for all $x$ such that $x\in X$ we have
                $\{x\}\times Y$ is connected with respect to $\subspaceTopology{T}{\{x\}\times Y}$.
            \begin{subproof}
                Fix $x$.
                Assume $x\in X$.
                Let $j_Y = \identity{Y}$.
                Thus $j_Y$ is continuous from $Y$ to $Y$ with respect to $T_2$ and $T_2$
                    by \cref{identity_continuous}.
                Show that for all $y$ such that $y\in Y$ we have $j_Y(y) = y$.
                \begin{subproof}
                    Fix $y$.
                    Assume $y\in Y$.
                    Thus $y\mathrel{j_Y} y$ by \cref{id_iff}.
                    $j_Y\in\funs{Y}{Y}$ by \cref{id_elem_funs}.
                    Thus $j_Y$ is a function by \cref{funs_is_function}.
                    Thus $j_Y(y) = y$ by \cref{function_apply_intro}.
                \end{subproof}
                Let $h = \{(y,x) \mid y\in Y\}$.
                Thus $h$ is a function from $Y$ to $X$ by \cref{constant_map_function}.
                Show that for all $y$ such that $y\in Y$ we have $h(y) = x$.
                \begin{subproof}
                Fix $y$.
                Assume $y\in Y$.
                Thus $(y,x)\in h$.
                Thus $h$ is a function.
                Thus $h(y) = x$ by \cref{function_apply_intro}.
            \end{subproof}
                Thus $h$ is continuous from $Y$ to $X$ with respect to $T_2$ and $T_1$
                    by \cref{constant_continuous}.
                Let $g = \{(y,(h(y),j_Y(y))) \mid y\in Y\}$.
                Thus $g$ is continuous from $Y$ to $X\times Y$ with respect to $T_2$ and $T$
                    by \cref{continuous_map_into_product_backward}.
                Let $Z_x = \img{g}{Y}$.
                Show that $Z_x = \{x\}\times Y$.
                \begin{subproof}
                    Show that $Z_x\subseteq \{x\}\times Y$.
                    \begin{subproof}
                        Show that for all $p$ such that $p\in Z_x$ we have $p\in \{x\}\times Y$.
                        \begin{subproof}
                            Fix $p$.
                            Assume $p\in Z_x$.
                            Take $y\in Y$ such that $y\mathrel{g} p$ by \cref{img_iff}.
                            Take $y_0\in Y$ such that $(y,p) = (y_0,(h(y_0),j_Y(y_0)))$.
                            Thus $p = (h(y),j_Y(y))$ by \cref{pair_eq_iff}.
                            Thus $j_Y(y) = y$.
                            Thus $h(y) = x$.
                            Thus $p = (x,y)$.
                            Thus $x\in\{x\}$ by \cref{singleton_intro}.
                            Thus $p\in \{x\}\times Y$ by \cref{times_tuple_intro}.
                        \end{subproof}
                        Thus $Z_x\subseteq \{x\}\times Y$ by \cref{subseteq}.
                    \end{subproof}
                    Show that $\{x\}\times Y\subseteq Z_x$.
                    \begin{subproof}
                        Show that for all $p$ such that $p\in \{x\}\times Y$ we have $p\in Z_x$.
                        \begin{subproof}
                            Fix $p$.
                            Assume $p\in \{x\}\times Y$.
                            Take $y_0,y$ such that $y_0\in\{x\}$ and $y\in Y$ and $p = (y_0,y)$
                                by \cref{times_elem_is_tuple}.
                            Thus $y_0 = x$ by \cref{singleton_iff}.
                            Thus $p = (x,y)$ by \cref{pair_eq_iff}.
                            Thus $(y,x)\in h$.
                            Thus $h$ is a function.
                            Thus $h(y) = x$ by \cref{function_apply_intro}.
                            Thus $j_Y(y) = y$.
                            Thus $(y,(h(y),j_Y(y))) = (y,p)$.
                            Thus $(y,p)\in g$.
                            Thus $y\mathrel{g} p$.
                            Thus $p\in Z_x$ by \cref{img_iff}.
                        \end{subproof}
                        Thus $\{x\}\times Y\subseteq Z_x$ by \cref{subseteq}.
                    \end{subproof}
                    Thus $Z_x = \{x\}\times Y$ by \cref{subseteq_antisymmetric}.
                \end{subproof}
                Thus $\{x\}\times Y$ is connected with respect to $\subspaceTopology{T}{\{x\}\times Y}$
                    by \cref{connected_image_continuous}.
            \end{subproof}
            Show that for all $x$ such that $x\in X$ we have
                $(X\times\{b\})\union(\{x\}\times Y)$ is connected with respect to
                $\subspaceTopology{T}{(X\times\{b\})\union(\{x\}\times Y)}$.
            \begin{subproof}
                Fix $x$.
                Assume $x\in X$.
                Let $M = \{X\times\{b\},\{x\}\times Y\}$.
                Show that $M\neq\emptyset$.
                \begin{subproof}
                    Suppose not.
                    Then $M = \emptyset$.
                    $X\times\{b\}\in M$ by \cref{upair_intro_left}.
                    Contradiction by \cref{emptyset}.
                \end{subproof}
                Show that for all $A$ such that $A\in M$ we have $A\subseteq X\times Y$.
                \begin{subproof}
                    Fix $A$.
                    Assume $A\in M$.
                    Thus $A = X\times\{b\}$ or $A = \{x\}\times Y$ by \cref{upair_iff}.
                    \begin{byCase}
                    \caseOf{$A = X\times\{b\}$.}
                        Show that $\{b\}\subseteq Y$.
                        \begin{subproof}
                            Show that for all $y$ such that $y\in\{b\}$ we have $y\in Y$.
                            \begin{subproof}
                                Fix $y$.
                                Assume $y\in\{b\}$.
                                Thus $y = b$ by \cref{singleton_iff}.
                                Thus $y\in Y$.
                            \end{subproof}
                            Thus $\{b\}\subseteq Y$ by \cref{subseteq}.
                        \end{subproof}
                        Thus $A\subseteq X\times Y$ by \cref{times_subseteq_right}.
                    \caseOf{$A = \{x\}\times Y$.}
                        Show that $\{x\}\subseteq X$.
                        \begin{subproof}
                            Show that for all $y$ such that $y\in\{x\}$ we have $y\in X$.
                            \begin{subproof}
                                Fix $y$.
                                Assume $y\in\{x\}$.
                                Thus $y = x$ by \cref{singleton_iff}.
                                Thus $y\in X$.
                            \end{subproof}
                            Thus $\{x\}\subseteq X$ by \cref{subseteq}.
                        \end{subproof}
                        Thus $A\subseteq X\times Y$ by \cref{times_subseteq_left}.
                    \end{byCase}
                \end{subproof}
                Show that for all $A$ such that $A\in M$ we have
                    $A$ is connected with respect to $\subspaceTopology{T}{A}$.
                \begin{subproof}
                    Fix $A$.
                    Assume $A\in M$.
                    Thus $A = X\times\{b\}$ or $A = \{x\}\times Y$ by \cref{upair_iff}.
                    \begin{byCase}
                    \caseOf{$A = X\times\{b\}$.}
                        Thus $A$ is connected with respect to $\subspaceTopology{T}{A}$.
                    \caseOf{$A = \{x\}\times Y$.}
                        Thus $A$ is connected with respect to $\subspaceTopology{T}{A}$.
                    \end{byCase}
                \end{subproof}
                Show that $\inters{M}\neq\emptyset$.
                \begin{subproof}
                    Show that $(x,b)\in\inters{M}$.
                    \begin{subproof}
                        Show that there exists $A$ such that $A\in M$.
                        \begin{subproof}
                            Let $A = X\times\{b\}$.
                            Thus $A\in M$ by \cref{upair_intro_left}.
                        \end{subproof}
                        Show that for all $A$ such that $A\in M$ we have $(x,b)\in A$.
                        \begin{subproof}
                            Fix $A$.
                            Assume $A\in M$.
                            Thus $A = X\times\{b\}$ or $A = \{x\}\times Y$ by \cref{upair_iff}.
                            \begin{byCase}
                            \caseOf{$A = X\times\{b\}$.}
                                Thus $x\in X$.
                                Thus $b\in\{b\}$ by \cref{singleton_intro}.
                                Thus $(x,b)\in A$ by \cref{times_tuple_intro}.
                            \caseOf{$A = \{x\}\times Y$.}
                                Thus $x\in\{x\}$ by \cref{singleton_intro}.
                                Thus $b\in Y$.
                                Thus $(x,b)\in A$ by \cref{times_tuple_intro}.
                            \end{byCase}
                        \end{subproof}
                        Thus $(x,b)\in\inters{M}$ by \cref{inters_iff_forall}.
                    \end{subproof}
                    Suppose not.
                    Then $\inters{M} = \emptyset$.
                    Thus $(x,b)\in\emptyset$.
                    Contradiction by \cref{emptyset}.
                \end{subproof}
                Let $T_x = (X\times\{b\})\union(\{x\}\times Y)$.
                Thus $\unions{M} = T_x$ by \cref{union_as_unions}.
                Thus $T_x$ is connected with respect to $\subspaceTopology{T}{T_x}$
                    by \cref{union_connected_common_point}.
            \end{subproof}
            Show that there exists no $C$ such that there exists $D$ such that
                $X\times Y$ is separated by $C$ and $D$ with respect to $T$.
            \begin{subproof}
                Suppose not.
                Then there exists $C$ such that there exists $D$ such that
                    $X\times Y$ is separated by $C$ and $D$ with respect to $T$.
                Take $C$ such that there exists $D$ such that
                    $X\times Y$ is separated by $C$ and $D$ with respect to $T$.
                Take $D$ such that $X\times Y$ is separated by $C$ and $D$ with respect to $T$.
                Show that for all $x$ such that $x\in X$ we have
                    $(X\times\{b\})\union(\{x\}\times Y)\subseteq C$ or
                    $(X\times\{b\})\union(\{x\}\times Y)\subseteq D$.
                \begin{subproof}
                    Fix $x$.
                    Assume $x\in X$.
                    Let $T_x = (X\times\{b\})\union(\{x\}\times Y)$.
                    Show that $T_x\subseteq X\times Y$.
                    \begin{subproof}
                        Show that for all $p$ such that $p\in T_x$ we have $p\in X\times Y$.
                        \begin{subproof}
                            Fix $p$.
                            Assume $p\in T_x$.
                            Thus $p\in X\times\{b\}$ or $p\in\{x\}\times Y$ by \cref{union_iff}.
                            \begin{byCase}
                            \caseOf{$p\in X\times\{b\}$.}
                                Show that $\{b\}\subseteq Y$.
                                \begin{subproof}
                                    Show that for all $y$ such that $y\in\{b\}$ we have $y\in Y$.
                                    \begin{subproof}
                                        Fix $y$.
                                        Assume $y\in\{b\}$.
                                        Thus $y = b$ by \cref{singleton_iff}.
                                        Thus $y\in Y$.
                                    \end{subproof}
                                    Thus $\{b\}\subseteq Y$ by \cref{subseteq}.
                                \end{subproof}
                                Thus $p\in X\times Y$ by \cref{times_subseteq_right,subseteq}.
                            \caseOf{$p\in\{x\}\times Y$.}
                                Show that $\{x\}\subseteq X$.
                                \begin{subproof}
                                    Show that for all $y$ such that $y\in\{x\}$ we have $y\in X$.
                                    \begin{subproof}
                                        Fix $y$.
                                        Assume $y\in\{x\}$.
                                        Thus $y = x$ by \cref{singleton_iff}.
                                        Thus $y\in X$.
                                    \end{subproof}
                                    Thus $\{x\}\subseteq X$ by \cref{subseteq}.
                                \end{subproof}
                                Thus $p\in X\times Y$ by \cref{times_subseteq_left,subseteq}.
                            \end{byCase}
                        \end{subproof}
                        Thus $T_x\subseteq X\times Y$ by \cref{subseteq}.
                    \end{subproof}
                    Thus $T_x$ is a subspace of $X\times Y$ with respect to $T$ by \cref{subspace_of}.
                    Thus $T_x$ is connected with respect to $\subspaceTopology{T}{T_x}$.
                    Thus $T_x\subseteq C$ or $T_x\subseteq D$ by \cref{connected_subspace_in_separation}.
                \end{subproof}
                Thus $(a,b)\in X\times Y$ by \cref{times_tuple_intro}.
                $C\union D = X\times Y$ by \cref{separation}.
                Thus $(a,b)\in C\union D$ by \cref{subseteq}.
                Thus $(a,b)\in C$ or $(a,b)\in D$ by \cref{union_iff}.
                \begin{byCase}
                \caseOf{$(a,b)\in C$.}
                    Show that for all $x$ such that $x\in X$ we have
                        $(X\times\{b\})\union(\{x\}\times Y)\subseteq C$.
                    \begin{subproof}
                        Fix $x$.
                        Assume $x\in X$.
                        Let $T_x = (X\times\{b\})\union(\{x\}\times Y)$.
                        Show that $(a,b)\in T_x$.
                        \begin{subproof}
                            Thus $b\in\{b\}$ by \cref{singleton_intro}.
                            Thus $(a,b)\in X\times\{b\}$ by \cref{times_tuple_intro}.
                            Thus $(a,b)\in T_x$ by \cref{union_intro_left}.
                        \end{subproof}
                        Show that $T_x\subseteq C$.
                        \begin{subproof}
                            Thus $T_x\subseteq C$ or $T_x\subseteq D$.
                            \begin{byCase}
                            \caseOf{$T_x\subseteq C$.}
                                Thus $T_x\subseteq C$.
                            \caseOf{$T_x\subseteq D$.}
                                Thus $(a,b)\in D$ by \cref{subseteq}.
                                Contradiction by \cref{separation,disjoint}.
                            \end{byCase}
                        \end{subproof}
                        Thus $T_x\subseteq C$.
                    \end{subproof}
                    Show that $X\times Y\subseteq C$.
                    \begin{subproof}
                        Show that for all $p$ such that $p\in X\times Y$ we have $p\in C$.
                        \begin{subproof}
                            Fix $p$.
                            Assume $p\in X\times Y$.
                            Take $x,y$ such that $x\in X$ and $y\in Y$ and $p = (x,y)$ by \cref{times_elem_is_tuple}.
                            Let $T_x = (X\times\{b\})\union(\{x\}\times Y)$.
                            Thus $x\in\{x\}$ by \cref{singleton_intro}.
                            Thus $p\in \{x\}\times Y$ by \cref{times_tuple_intro}.
                            Thus $p\in T_x$ by \cref{union_intro_right}.
                            Thus $T_x\subseteq C$.
                            Thus $p\in C$ by \cref{subseteq}.
                        \end{subproof}
                        Thus $X\times Y\subseteq C$ by \cref{subseteq}.
                    \end{subproof}
                    Show that $D = \emptyset$.
                    \begin{subproof}
                        Take $d$ such that $d\in D$ by \cref{nonempty_inhabited,separation}.
                        $D$ is open in $X\times Y$ with respect to $T$ by \cref{separation}.
                        $D\in T$ by \cref{open_in}.
                        $T\subseteq\pow{X\times Y}$ by \cref{topology_on}.
                        Thus $D\subseteq X\times Y$ by \cref{pow_iff,subseteq}.
                        Thus $d\in X\times Y$ by \cref{subseteq}.
                        Thus $d\in C$ by \cref{subseteq}.
                        Contradiction by \cref{separation,disjoint}.
                    \end{subproof}
                    Contradiction by \cref{separation}.
                \caseOf{$(a,b)\in D$.}
                    Show that for all $x$ such that $x\in X$ we have
                        $(X\times\{b\})\union(\{x\}\times Y)\subseteq D$.
                    \begin{subproof}
                        Fix $x$.
                        Assume $x\in X$.
                        Let $T_x = (X\times\{b\})\union(\{x\}\times Y)$.
                        Show that $(a,b)\in T_x$.
                        \begin{subproof}
                            Thus $b\in\{b\}$ by \cref{singleton_intro}.
                            Thus $(a,b)\in X\times\{b\}$ by \cref{times_tuple_intro}.
                            Thus $(a,b)\in T_x$ by \cref{union_intro_left}.
                        \end{subproof}
                        Show that $T_x\subseteq D$.
                        \begin{subproof}
                            Thus $T_x\subseteq C$ or $T_x\subseteq D$.
                            \begin{byCase}
                            \caseOf{$T_x\subseteq D$.}
                                Thus $T_x\subseteq D$.
                            \caseOf{$T_x\subseteq C$.}
                                Thus $(a,b)\in C$ by \cref{subseteq}.
                                Contradiction by \cref{separation,disjoint}.
                            \end{byCase}
                        \end{subproof}
                        Thus $T_x\subseteq D$.
                    \end{subproof}
                    Show that $X\times Y\subseteq D$.
                    \begin{subproof}
                        Show that for all $p$ such that $p\in X\times Y$ we have $p\in D$.
                        \begin{subproof}
                            Fix $p$.
                            Assume $p\in X\times Y$.
                            Take $x,y$ such that $x\in X$ and $y\in Y$ and $p = (x,y)$ by \cref{times_elem_is_tuple}.
                            Let $T_x = (X\times\{b\})\union(\{x\}\times Y)$.
                            Thus $x\in\{x\}$ by \cref{singleton_intro}.
                            Thus $p\in \{x\}\times Y$ by \cref{times_tuple_intro}.
                            Thus $p\in T_x$ by \cref{union_intro_right}.
                            Thus $T_x\subseteq D$.
                            Thus $p\in D$ by \cref{subseteq}.
                        \end{subproof}
                        Thus $X\times Y\subseteq D$ by \cref{subseteq}.
                    \end{subproof}
                    Show that $C = \emptyset$.
                    \begin{subproof}
                        Take $c$ such that $c\in C$ by \cref{nonempty_inhabited,separation}.
                        $C$ is open in $X\times Y$ with respect to $T$ by \cref{separation}.
                        $C\in T$ by \cref{open_in}.
                        $T\subseteq\pow{X\times Y}$ by \cref{topology_on}.
                        Thus $C\subseteq X\times Y$ by \cref{pow_iff,subseteq}.
                        Thus $c\in X\times Y$ by \cref{subseteq}.
                        Thus $c\in D$ by \cref{subseteq}.
                        Contradiction by \cref{separation,disjoint}.
                    \end{subproof}
                    Contradiction by \cref{separation}.
                \end{byCase}
            \end{subproof}
            Thus $X\times Y$ is connected with respect to $T$ by \cref{connected}.
        \end{byCase}
    \end{byCase}
\end{proof}

% from §24 helper definition 1: upper bound in an ordered set
\begin{definition}\label{upper_bound_rel}
    $b$ is an upper bound of $A$ in $X$ with respect to $R$ iff
        $b\in X$ and $A\subseteq X$ and
        for all $a\in A$ we have $a\mathrel{R}b$ or $a = b$.
\end{definition}

% from §24 helper definition 2: least upper bound in an ordered set
\begin{definition}\label{least_upper_bound_rel}
    $b$ is a least upper bound of $A$ in $X$ with respect to $R$ iff
        $b$ is an upper bound of $A$ in $X$ with respect to $R$ and
        for all $c$ such that $c$ is an upper bound of $A$ in $X$ with respect to $R$
            we have $b\mathrel{R}c$ or $b = c$.
\end{definition}

% from §24 helper definition 3: least upper bound property
\begin{definition}\label{least_upper_bound_property}
    $X$ has the least upper bound property with respect to $R$ iff
        for all $A$ such that $A\subseteq X$ and $A$ is inhabited and
            there exists $b$ such that $b$ is an upper bound of $A$ in $X$ with respect to $R$
        there exists $c$ such that $c$ is a least upper bound of $A$ in $X$ with respect to $R$.
\end{definition}

% from §24 Item 1: linear continuum
\begin{definition}\label{linear_continuum}
    $X$ is a linear continuum with respect to $R$ iff
        $R$ is a simple order relation on $X$ and
        $X$ has more than one element and
        $X$ has the least upper bound property with respect to $R$ and
        for all $x,y\in X$ if $x\mathrel{R}y$ then
            there exists $z\in X$ such that $x\mathrel{R}z$ and $z\mathrel{R}y$.
\end{definition}

% from §24 helper lemma 1: open interval is convex
\begin{lemma}\label{intervalopenrel_convex}
    Assume $R$ is transitive.
    Suppose $a,b\in X$.
    Then $\intervalopenrel{R}{X}{a}{b}$ is convex in $X$ with respect to $R$.
\end{lemma}
\begin{proof}
    Let $Y = \intervalopenrel{R}{X}{a}{b}$.
    Show that $Y\subseteq X$.
    \begin{subproof}
        Show that for all $y$ such that $y\in Y$ we have $y\in X$.
        \begin{subproof}
            Fix $y$.
            Assume $y\in Y$.
            Thus $y\in X$ by \cref{intervalopen_rel}.
        \end{subproof}
        Thus $Y\subseteq X$ by \cref{subseteq}.
    \end{subproof}
    Show that for all $u,v$ such that $u\in Y$ and $v\in Y$ and $u\mathrel{R}v$
        for all $x$ such that $x\in X$ and $u\mathrel{R}x$ and $x\mathrel{R}v$ we have $x\in Y$.
    \begin{subproof}
        Fix $u$.
        Fix $v$.
        Assume $u\in Y$ and $v\in Y$ and $u\mathrel{R}v$.
        Fix $x$.
        Assume $x\in X$ and $u\mathrel{R}x$ and $x\mathrel{R}v$.
        Thus $a\mathrel{R}u$ and $u\mathrel{R}b$ by \cref{intervalopen_rel}.
        Thus $a\mathrel{R}x$ by \hyperref[transitive]{transitivity}.
        Thus $v\mathrel{R}b$ by \cref{intervalopen_rel}.
        Thus $x\mathrel{R}b$ by \hyperref[transitive]{transitivity}.
        Thus $x\in Y$ by \cref{intervalopen_rel}.
    \end{subproof}
    Thus $Y$ is convex in $X$ with respect to $R$ by \cref{convex_in}.
\end{proof}

% from §24 helper lemma 2: open-closed interval is convex
\begin{lemma}\label{intervalopenclosedrel_convex}
    Assume $R$ is transitive.
    Suppose $a,b\in X$.
    Then $\intervalopenclosedrel{R}{X}{a}{b}$ is convex in $X$ with respect to $R$.
\end{lemma}
\begin{proof}
    Let $Y = \intervalopenclosedrel{R}{X}{a}{b}$.
    Show that $Y\subseteq X$.
    \begin{subproof}
        Show that for all $y$ such that $y\in Y$ we have $y\in X$.
        \begin{subproof}
            Fix $y$.
            Assume $y\in Y$.
            Thus $y\in X$ by \cref{intervalopenclosed_rel}.
        \end{subproof}
        Thus $Y\subseteq X$ by \cref{subseteq}.
    \end{subproof}
    Show that for all $u,v$ such that $u\in Y$ and $v\in Y$ and $u\mathrel{R}v$
        for all $x$ such that $x\in X$ and $u\mathrel{R}x$ and $x\mathrel{R}v$ we have $x\in Y$.
    \begin{subproof}
        Fix $u$.
        Fix $v$.
        Assume $u\in Y$ and $v\in Y$ and $u\mathrel{R}v$.
        Fix $x$.
        Assume $x\in X$ and $u\mathrel{R}x$ and $x\mathrel{R}v$.
        Thus $a\mathrel{R}u$ by \cref{intervalopenclosed_rel}.
        Thus $a\mathrel{R}x$ by \hyperref[transitive]{transitivity}.
        Thus $v\mathrel{R}b$ or $v = b$ by \cref{intervalopenclosed_rel}.
        \begin{byCase}
        \caseOf{$v\mathrel{R}b$.}
            Thus $x\mathrel{R}b$ by \hyperref[transitive]{transitivity}.
            Thus $x\in Y$ by \cref{intervalopenclosed_rel}.
        \caseOf{$v = b$.}
            Thus $x\mathrel{R}b$.
            Thus $x\in Y$ by \cref{intervalopenclosed_rel}.
        \end{byCase}
    \end{subproof}
    Thus $Y$ is convex in $X$ with respect to $R$ by \cref{convex_in}.
\end{proof}

% from §24 helper lemma 3: closed-open interval is convex
\begin{lemma}\label{intervalclosedopenrel_convex}
    Assume $R$ is transitive.
    Suppose $a,b\in X$.
    Then $\intervalclosedopenrel{R}{X}{a}{b}$ is convex in $X$ with respect to $R$.
\end{lemma}
\begin{proof}
    Let $Y = \intervalclosedopenrel{R}{X}{a}{b}$.
    Show that $Y\subseteq X$.
    \begin{subproof}
        Show that for all $y$ such that $y\in Y$ we have $y\in X$.
        \begin{subproof}
            Fix $y$.
            Assume $y\in Y$.
            Thus $y\in X$ by \cref{intervalclosedopen_rel}.
        \end{subproof}
        Thus $Y\subseteq X$ by \cref{subseteq}.
    \end{subproof}
    Show that for all $u,v$ such that $u\in Y$ and $v\in Y$ and $u\mathrel{R}v$
        for all $x$ such that $x\in X$ and $u\mathrel{R}x$ and $x\mathrel{R}v$ we have $x\in Y$.
    \begin{subproof}
        Fix $u$.
        Fix $v$.
        Assume $u\in Y$ and $v\in Y$ and $u\mathrel{R}v$.
        Fix $x$.
        Assume $x\in X$ and $u\mathrel{R}x$ and $x\mathrel{R}v$.
        Thus $a\mathrel{R}u$ or $u = a$ by \cref{intervalclosedopen_rel}.
        Thus $v\mathrel{R}b$ by \cref{intervalclosedopen_rel}.
        Thus $x\mathrel{R}b$ by \hyperref[transitive]{transitivity}.
        \begin{byCase}
        \caseOf{$a\mathrel{R}u$.}
            Thus $a\mathrel{R}x$ by \hyperref[transitive]{transitivity}.
            Thus $x\in Y$ by \cref{intervalclosedopen_rel}.
        \caseOf{$u = a$.}
            Thus $a\mathrel{R}x$.
            Thus $x\in Y$ by \cref{intervalclosedopen_rel}.
        \end{byCase}
    \end{subproof}
    Thus $Y$ is convex in $X$ with respect to $R$ by \cref{convex_in}.
\end{proof}

% from §24 helper lemma 4: closed interval is convex
\begin{lemma}\label{intervalclosedrel_convex}
    Assume $R$ is transitive.
    Suppose $a,b\in X$.
    Then $\intervalclosedrel{R}{X}{a}{b}$ is convex in $X$ with respect to $R$.
\end{lemma}
\begin{proof}
    Let $Y = \intervalclosedrel{R}{X}{a}{b}$.
    Show that $Y\subseteq X$.
    \begin{subproof}
        Show that for all $y$ such that $y\in Y$ we have $y\in X$.
        \begin{subproof}
            Fix $y$.
            Assume $y\in Y$.
            Thus $y\in X$ by \cref{intervalclosed_rel}.
        \end{subproof}
        Thus $Y\subseteq X$ by \cref{subseteq}.
    \end{subproof}
    Show that for all $u,v$ such that $u\in Y$ and $v\in Y$ and $u\mathrel{R}v$
        for all $x$ such that $x\in X$ and $u\mathrel{R}x$ and $x\mathrel{R}v$ we have $x\in Y$.
    \begin{subproof}
        Fix $u$.
        Fix $v$.
        Assume $u\in Y$ and $v\in Y$ and $u\mathrel{R}v$.
        Fix $x$.
        Assume $x\in X$ and $u\mathrel{R}x$ and $x\mathrel{R}v$.
        Thus $a\mathrel{R}u$ or $u = a$ by \cref{intervalclosed_rel}.
        Show that $a\mathrel{R}x$.
        \begin{subproof}
            \begin{byCase}
            \caseOf{$a\mathrel{R}u$.}
                Thus $a\mathrel{R}x$ by \hyperref[transitive]{transitivity}.
            \caseOf{$u = a$.}
                Thus $a\mathrel{R}x$.
            \end{byCase}
        \end{subproof}
        Thus $v\mathrel{R}b$ or $v = b$ by \cref{intervalclosed_rel}.
        Show that $x\mathrel{R}b$.
        \begin{subproof}
            \begin{byCase}
            \caseOf{$v\mathrel{R}b$.}
                Thus $x\mathrel{R}b$ by \hyperref[transitive]{transitivity}.
            \caseOf{$v = b$.}
                Thus $x\mathrel{R}b$.
            \end{byCase}
        \end{subproof}
        Thus $x\in Y$ by \cref{intervalclosed_rel}.
    \end{subproof}
    Thus $Y$ is convex in $X$ with respect to $R$ by \cref{convex_in}.
\end{proof}

% from §24 helper lemma 5: open right ray is convex
\begin{lemma}\label{openrayright_convex}
    Assume $R$ is transitive.
    Suppose $a\in X$.
    Then $\openrayright{R}{X}{a}$ is convex in $X$ with respect to $R$.
\end{lemma}
\begin{proof}
    Let $Y = \openrayright{R}{X}{a}$.
    Show that $Y\subseteq X$.
    \begin{subproof}
        Show that for all $y$ such that $y\in Y$ we have $y\in X$.
        \begin{subproof}
            Fix $y$.
            Assume $y\in Y$.
            Thus $y\in X$ by \cref{openray_right}.
        \end{subproof}
        Thus $Y\subseteq X$ by \cref{subseteq}.
    \end{subproof}
    Show that for all $u,v$ such that $u\in Y$ and $v\in Y$ and $u\mathrel{R}v$
        for all $x$ such that $x\in X$ and $u\mathrel{R}x$ and $x\mathrel{R}v$ we have $x\in Y$.
    \begin{subproof}
        Fix $u$.
        Fix $v$.
        Assume $u\in Y$ and $v\in Y$ and $u\mathrel{R}v$.
        Fix $x$.
        Assume $x\in X$ and $u\mathrel{R}x$ and $x\mathrel{R}v$.
        Thus $a\mathrel{R}u$ by \cref{openray_right}.
        Thus $a\mathrel{R}x$ by \hyperref[transitive]{transitivity}.
        Thus $x\in Y$ by \cref{openray_right}.
    \end{subproof}
    Thus $Y$ is convex in $X$ with respect to $R$ by \cref{convex_in}.
\end{proof}

% from §24 helper lemma 6: open left ray is convex
\begin{lemma}\label{openrayleft_convex}
    Assume $R$ is transitive.
    Suppose $a\in X$.
    Then $\openrayleft{R}{X}{a}$ is convex in $X$ with respect to $R$.
\end{lemma}
\begin{proof}
    Let $Y = \openrayleft{R}{X}{a}$.
    Show that $Y\subseteq X$.
    \begin{subproof}
        Show that for all $y$ such that $y\in Y$ we have $y\in X$.
        \begin{subproof}
            Fix $y$.
            Assume $y\in Y$.
            Thus $y\in X$ by \cref{openray_left}.
        \end{subproof}
        Thus $Y\subseteq X$ by \cref{subseteq}.
    \end{subproof}
    Show that for all $u,v$ such that $u\in Y$ and $v\in Y$ and $u\mathrel{R}v$
        for all $x$ such that $x\in X$ and $u\mathrel{R}x$ and $x\mathrel{R}v$ we have $x\in Y$.
    \begin{subproof}
        Fix $u$.
        Fix $v$.
        Assume $u\in Y$ and $v\in Y$ and $u\mathrel{R}v$.
        Fix $x$.
        Assume $x\in X$ and $u\mathrel{R}x$ and $x\mathrel{R}v$.
        Thus $v\mathrel{R}a$ by \cref{openray_left}.
        Thus $x\mathrel{R}a$ by \hyperref[transitive]{transitivity}.
        Thus $x\in Y$ by \cref{openray_left}.
    \end{subproof}
    Thus $Y$ is convex in $X$ with respect to $R$ by \cref{convex_in}.
\end{proof}

% from §24 helper lemma 7: closed right ray is convex
\begin{lemma}\label{closedrayright_convex}
    Assume $R$ is transitive.
    Suppose $a\in X$.
    Then $\closedrayright{R}{X}{a}$ is convex in $X$ with respect to $R$.
\end{lemma}
\begin{proof}
    Let $Y = \closedrayright{R}{X}{a}$.
    Show that $Y\subseteq X$.
    \begin{subproof}
        Show that for all $y$ such that $y\in Y$ we have $y\in X$.
        \begin{subproof}
            Fix $y$.
            Assume $y\in Y$.
            Thus $y\in X$ by \cref{closedray_right}.
        \end{subproof}
        Thus $Y\subseteq X$ by \cref{subseteq}.
    \end{subproof}
    Show that for all $u,v$ such that $u\in Y$ and $v\in Y$ and $u\mathrel{R}v$
        for all $x$ such that $x\in X$ and $u\mathrel{R}x$ and $x\mathrel{R}v$ we have $x\in Y$.
    \begin{subproof}
        Fix $u$.
        Fix $v$.
        Assume $u\in Y$ and $v\in Y$ and $u\mathrel{R}v$.
        Fix $x$.
        Assume $x\in X$ and $u\mathrel{R}x$ and $x\mathrel{R}v$.
        Thus $a\mathrel{R}u$ or $u = a$ by \cref{closedray_right}.
        Show that $a\mathrel{R}x$.
        \begin{subproof}
            \begin{byCase}
            \caseOf{$a\mathrel{R}u$.}
                Thus $a\mathrel{R}x$ by \hyperref[transitive]{transitivity}.
            \caseOf{$u = a$.}
                Thus $a\mathrel{R}x$.
            \end{byCase}
        \end{subproof}
        Thus $x\in Y$ by \cref{closedray_right}.
    \end{subproof}
    Thus $Y$ is convex in $X$ with respect to $R$ by \cref{convex_in}.
\end{proof}

% from §24 helper lemma 8: closed left ray is convex
\begin{lemma}\label{closedrayleft_convex}
    Assume $R$ is transitive.
    Suppose $a\in X$.
    Then $\closedrayleft{R}{X}{a}$ is convex in $X$ with respect to $R$.
\end{lemma}
\begin{proof}
    Let $Y = \closedrayleft{R}{X}{a}$.
    Show that $Y\subseteq X$.
    \begin{subproof}
        Show that for all $y$ such that $y\in Y$ we have $y\in X$.
        \begin{subproof}
            Fix $y$.
            Assume $y\in Y$.
            Thus $y\in X$ by \cref{closedray_left}.
        \end{subproof}
        Thus $Y\subseteq X$ by \cref{subseteq}.
    \end{subproof}
    Show that for all $u,v$ such that $u\in Y$ and $v\in Y$ and $u\mathrel{R}v$
        for all $x$ such that $x\in X$ and $u\mathrel{R}x$ and $x\mathrel{R}v$ we have $x\in Y$.
    \begin{subproof}
        Fix $u$.
        Fix $v$.
        Assume $u\in Y$ and $v\in Y$ and $u\mathrel{R}v$.
        Fix $x$.
        Assume $x\in X$ and $u\mathrel{R}x$ and $x\mathrel{R}v$.
        Thus $v\mathrel{R}a$ or $v = a$ by \cref{closedray_left}.
        Show that $x\mathrel{R}a$.
        \begin{subproof}
            \begin{byCase}
            \caseOf{$v\mathrel{R}a$.}
                Thus $x\mathrel{R}a$ by \hyperref[transitive]{transitivity}.
            \caseOf{$v = a$.}
                Thus $x\mathrel{R}a$.
            \end{byCase}
        \end{subproof}
        Thus $x\in Y$ by \cref{closedray_left}.
    \end{subproof}
    Thus $Y$ is convex in $X$ with respect to $R$ by \cref{convex_in}.
\end{proof}

% from §24 helper lemma 9: convex separation contradiction
\begin{lemma}\label{convex_separation_contradiction}
    Assume $T$ is the order topology on $L$ with respect to $R$.
    Assume $L$ is a linear continuum with respect to $R$.
    Suppose $Y$ is convex in $L$ with respect to $R$.
    Let $T_Y = \subspaceTopology{T}{Y}$.
    Suppose $Y$ is separated by $A$ and $B$ with respect to $T_Y$.
    Suppose $a\in A$ and $b\in B$ and $a\mathrel{R}b$.
    Contradiction.
\end{lemma}
\begin{proof}
    $R$ is a simple order relation on $L$ by \cref{linear_continuum}.
    $R$ is transitive by \cref{simple_order_relation}.
    $L$ has more than one element by \cref{linear_continuum}.
    $R$ is transitive by \cref{simple_order_relation}.
    $R$ is asymmetric by \cref{simple_order_relation}.
    $R$ is connex on $L$ by \cref{simple_order_relation}.
    $L$ has the least upper bound property with respect to $R$ by \cref{linear_continuum}.
    $Y\subseteq L$ by \cref{convex_in}.

    $R$ is a simple order relation on $L$ by \cref{linear_continuum}.
    $L$ has more than one element by \cref{linear_continuum}.
    $T = \basisTopology{\orderBasis{R}{L}}{L}$ by \cref{order_topology}.
    $\orderBasis{R}{L}$ is a basis on $L$ by \cref{order_basis_is_basis}.
    Thus $T$ is a topology on $L$ by \cref{basis_topology_is_topology}.
    Thus $T_Y$ is a topology on $Y$ by \cref{subspace_topology_is_topology,subspace_of,convex_in}.
    $T_Y\subseteq\pow{Y}$ by \cref{topology_on}.
    $A$ is open in $Y$ with respect to $T_Y$ by \cref{separation}.
    $B$ is open in $Y$ with respect to $T_Y$ by \cref{separation}.
    $A\in T_Y$ by \cref{open_in}.
    $B\in T_Y$ by \cref{open_in}.
    Thus $A\in\pow{Y}$ by \cref{subseteq}.
    Thus $B\in\pow{Y}$ by \cref{subseteq}.
    Thus $A\subseteq Y$ by \cref{powerset_elim}.
    Thus $B\subseteq Y$ by \cref{powerset_elim}.
    Thus $a\in Y$ and $b\in Y$ by \cref{subseteq}.
    Thus $a\in L$ and $b\in L$ by \cref{subseteq}.

    Let $I = \intervalclosedrel{R}{L}{a}{b}$.
    Show that $I\subseteq Y$.
    \begin{subproof}
        Show that for all $x$ such that $x\in I$ we have $x\in Y$.
        \begin{subproof}
            Fix $x$.
            Assume $x\in I$.
            Thus $x\in L$ and $(a\mathrel{R}x \lor x = a)$ and $(x\mathrel{R}b \lor x = b)$
                by \cref{intervalclosed_rel}.
            \begin{byCase}
            \caseOf{$x = a$.} Thus $x\in Y$. 
            \caseOf{$x = b$.} Thus $x\in Y$. 
            \caseOf{$a\mathrel{R}x$ and $x\mathrel{R}b$.}
                Thus $x\in Y$ by \cref{convex_in}.
            \end{byCase}
        \end{subproof}
        Thus $I\subseteq Y$ by \cref{subseteq}.
    \end{subproof}

    Show that $I\subseteq L$.
    \begin{subproof}
        Show that for all $x$ such that $x\in I$ we have $x\in L$.
        \begin{subproof}
            Fix $x$.
            Assume $x\in I$.
            Thus $x\in L$ by \cref{intervalclosed_rel}.
        \end{subproof}
        Thus $I\subseteq L$ by \cref{subseteq}.
    \end{subproof}

    Let $A_0 = A\inter I$.
    Let $B_0 = B\inter I$.
    Let $T_I = \subspaceTopology{T}{I}$.

    Show that $A_0$ is open in $I$ with respect to $T_I$.
    \begin{subproof}
        Take $U\in T$ such that $A = Y\inter U$ by \cref{subspace_topology,open_in}.
        Show that $A_0 = I\inter U$.
        \begin{subproof}
            Show that for all $x$ such that $x\in A_0$ we have $x\in I\inter U$.
            \begin{subproof}
                Fix $x$.
                Assume $x\in A_0$.
                Thus $x\in A$ and $x\in I$ by \cref{inter_elim_left,inter_elim_right}.
                Thus $x\in U$ by \cref{inter_elim_right,subspace_topology}.
                Thus $x\in I\inter U$ by \cref{inter_intro}.
            \end{subproof}
            Thus $A_0\subseteq I\inter U$ by \cref{subseteq}.
            Show that for all $x$ such that $x\in I\inter U$ we have $x\in A_0$.
            \begin{subproof}
                Fix $x$.
                Assume $x\in I\inter U$.
                Thus $x\in I$ and $x\in U$ by \cref{inter_elim_left,inter_elim_right}.
                $I\subseteq Y$.
                Thus $x\in Y$ by \cref{subseteq}.
                Thus $x\in Y\inter U$ by \cref{inter_intro}.
                Thus $x\in A$ by \cref{subspace_topology}.
                Thus $x\in A_0$ by \cref{inter_intro}.
            \end{subproof}
            Thus $I\inter U\subseteq A_0$ by \cref{subseteq}.
            Thus $A_0 = I\inter U$ by \cref{subseteq_antisymmetric}.
        \end{subproof}
        Thus $A_0\in T_I$ by \cref{subspace_topology}.
        $I\subseteq L$.
        Thus $T_I$ is a topology on $I$ by \cref{subspace_topology_is_topology,subspace_of}.
        Thus $A_0$ is open in $I$ with respect to $T_I$ by \cref{open_in}.
    \end{subproof}

    Show that $B_0$ is open in $I$ with respect to $T_I$.
    \begin{subproof}
        Take $U\in T$ such that $B = Y\inter U$ by \cref{subspace_topology,open_in}.
        Show that $B_0 = I\inter U$.
        \begin{subproof}
            Show that for all $x$ such that $x\in B_0$ we have $x\in I\inter U$.
            \begin{subproof}
                Fix $x$.
                Assume $x\in B_0$.
                Thus $x\in B$ and $x\in I$ by \cref{inter_elim_left,inter_elim_right}.
                Thus $x\in U$ by \cref{inter_elim_right,subspace_topology}.
                Thus $x\in I\inter U$ by \cref{inter_intro}.
            \end{subproof}
            Thus $B_0\subseteq I\inter U$ by \cref{subseteq}.
            Show that for all $x$ such that $x\in I\inter U$ we have $x\in B_0$.
            \begin{subproof}
                Fix $x$.
                Assume $x\in I\inter U$.
                Thus $x\in I$ and $x\in U$ by \cref{inter_elim_left,inter_elim_right}.
                $I\subseteq Y$.
                Thus $x\in Y$ by \cref{subseteq}.
                Thus $x\in Y\inter U$ by \cref{inter_intro}.
                Thus $x\in B$ by \cref{subspace_topology}.
                Thus $x\in B_0$ by \cref{inter_intro}.
            \end{subproof}
            Thus $I\inter U\subseteq B_0$ by \cref{subseteq}.
            Thus $B_0 = I\inter U$ by \cref{subseteq_antisymmetric}.
        \end{subproof}
        Thus $B_0\in T_I$ by \cref{subspace_topology}.
        $I\subseteq L$.
        Thus $T_I$ is a topology on $I$ by \cref{subspace_topology_is_topology,subspace_of}.
        Thus $B_0$ is open in $I$ with respect to $T_I$ by \cref{open_in}.
    \end{subproof}

    $a\in I$ by \cref{intervalclosed_rel}.
    $b\in I$ by \cref{intervalclosed_rel}.
    Thus $a\in A_0$ and $b\in B_0$ by \cref{inter_intro}.
    Show that $A_0\neq\emptyset$.
    \begin{subproof}
        Suppose not.
        Thus $a\notin A_0$ by \cref{emptyset}.
        Contradiction.
    \end{subproof}
    Show that $B_0\neq\emptyset$.
    \begin{subproof}
        Suppose not.
        Thus $b\notin B_0$ by \cref{emptyset}.
        Contradiction.
    \end{subproof}
    Thus $A_0\neq\emptyset$ and $B_0\neq\emptyset$.

    Show that $A_0$ is disjoint from $B_0$.
    \begin{subproof}
        $A$ is disjoint from $B$ by \cref{separation}.
        Show that there exists no $x$ such that $x\in A_0$ and $x\in B_0$.
        \begin{subproof}
            Suppose not.
            Take $x$ such that $x\in A_0$ and $x\in B_0$.
            Thus $x\in A$ and $x\in B$ by \cref{inter_elim_left}.
            Contradiction by \cref{disjoint}.
        \end{subproof}
        Thus $A_0$ is disjoint from $B_0$ by \cref{disjoint}.
    \end{subproof}

    Show that $A_0\union B_0 = I$.
    \begin{subproof}
        Show that $A_0\union B_0\subseteq I$.
        \begin{subproof}
            $A_0\subseteq I$ and $B_0\subseteq I$ by \cref{inter_elim_right,subseteq}.
            Thus $A_0\union B_0\subseteq I$ by \cref{union_subsets_is_subset}.
        \end{subproof}
        Show that for all $x$ such that $x\in I$ we have $x\in A_0\union B_0$.
        \begin{subproof}
            Fix $x$.
            Assume $x\in I$.
            Thus $x\in Y$ by \cref{subseteq}.
            $A\union B = Y$ by \cref{separation}.
            Thus $x\in A\union B$.
            Thus $x\in A$ or $x\in B$ by \cref{union_iff}.
            \begin{byCase}
            \caseOf{$x\in A$.}
                Thus $x\in A_0$ by \cref{inter_intro}.
                Thus $x\in A_0\union B_0$ by \cref{union_intro_left}.
            \caseOf{$x\in B$.}
                Thus $x\in B_0$ by \cref{inter_intro}.
                Thus $x\in A_0\union B_0$ by \cref{union_intro_right}.
            \end{byCase}
        \end{subproof}
        Thus $I\subseteq A_0\union B_0$ by \cref{subseteq}.
        Thus $A_0\union B_0 = I$ by \cref{subseteq_antisymmetric}.
    \end{subproof}

    $I$ is separated by $A_0$ and $B_0$ with respect to $T_I$ by \cref{separation}.

    Show that for all $x$ such that $x\in A_0$ we have $x\in I$.
    \begin{subproof}
        Fix $x$.
        Assume $x\in A_0$.
        Thus $x\in I$ by \cref{inter_elim_right}.
    \end{subproof}
    Thus $A_0\subseteq I$ by \cref{subseteq}.
    $I\subseteq L$.
    Thus $A_0\subseteq L$ by \cref{subseteq_transitive}.
    Thus $A_0$ is inhabited by \cref{nonempty_inhabited}.

    Show that there exists $u$ such that $u$ is an upper bound of $A_0$ in $L$ with respect to $R$.
    \begin{subproof}
        Show that $b$ is an upper bound of $A_0$ in $L$ with respect to $R$.
        \begin{subproof}
            $b\in L$.
            Show that for all $x\in A_0$ we have $x\mathrel{R}b$ or $x = b$.
            \begin{subproof}
                Fix $x$.
                Assume $x\in A_0$.
                Thus $x\in I$ by \cref{inter_elim_right}.
                Thus $x\mathrel{R}b$ or $x = b$ by \cref{intervalclosed_rel}.
            \end{subproof}
            Thus $b$ is an upper bound of $A_0$ in $L$ with respect to $R$ by \cref{upper_bound_rel}.
        \end{subproof}
        Thus there exists $u$ such that $u$ is an upper bound of $A_0$ in $L$ with respect to $R$.
    \end{subproof}

    Take $c$ such that $c$ is a least upper bound of $A_0$ in $L$ with respect to $R$ by
        \cref{least_upper_bound_property,linear_continuum}.

    Show that $c\in I$.
    \begin{subproof}
        $c$ is an upper bound of $A_0$ in $L$ with respect to $R$ by \cref{least_upper_bound_rel}.
        Thus $c\in L$ by \cref{upper_bound_rel}.
        Thus $a\mathrel{R}c$ or $a = c$ by \cref{upper_bound_rel,least_upper_bound_rel}.
        Show that $b$ is an upper bound of $A_0$ in $L$ with respect to $R$.
        \begin{subproof}
            $b\in L$.
            Show that for all $x\in A_0$ we have $x\mathrel{R}b$ or $x = b$.
            \begin{subproof}
                Fix $x$.
                Assume $x\in A_0$.
                Thus $x\in I$ by \cref{subseteq}.
                Thus $x\mathrel{R}b$ or $x = b$ by \cref{intervalclosed_rel}.
            \end{subproof}
            Thus $b$ is an upper bound of $A_0$ in $L$ with respect to $R$ by \cref{upper_bound_rel}.
        \end{subproof}
        Thus $c\mathrel{R}b$ or $c = b$ by \cref{least_upper_bound_rel}.
        Thus $c\in I$ by \cref{intervalclosed_rel}.
    \end{subproof}

    Thus $c\in A_0\union B_0$.
    Thus $c\in A_0$ or $c\in B_0$ by \cref{union_iff}.
    \begin{byCase}
    \caseOf{$c\in B_0$.}
        Suppose not.
        $c\notin A_0$ by \cref{disjoint}.
        Thus $c\neq a$.
        Thus $a\mathrel{R}c$ by \cref{intervalclosed_rel}.

        Take $U\in T$ such that $B_0 = I\inter U$ by \cref{subspace_topology,open_in}.
        Thus $c\in U$ by \cref{inter_elim_right}.
        Take $B_1\in\orderBasis{R}{L}$ such that $c\in B_1$ and $B_1\subseteq U$
            by \cref{topology_generated_by_basis}.
        Thus $B_1\inter I \subseteq B_0$ by \cref{subseteq_inter_iff,subseteq}.

        Show that there exists $d\in I$ such that $d\mathrel{R}c$ and
            $\intervalopenclosedrel{R}{I}{d}{c} \subseteq B_0$.
        \begin{subproof}
            $B_1\in\pow{L}$ and
                $((\exists u, v\in L. u\mathrel{R}v \land B_1 = \intervalopenrel{R}{L}{u}{v}) \lor
                (\exists v\in L. \exists u_0\in L. (\forall x\in L. u_0\mathrel{R}x \lor u_0 = x) \land B_1 = \intervalclosedopenrel{R}{L}{u_0}{v}) \lor
                (\exists u\in L. \exists v_0\in L. (\forall x\in L. x\mathrel{R}v_0 \lor x = v_0) \land B_1 = \intervalopenclosedrel{R}{L}{u}{v_0}))$
                by \cref{order_basis}.
            \begin{byCase}
            \caseOf{$\exists u, v\in L. u\mathrel{R}v \land B_1 = \intervalopenrel{R}{L}{u}{v}$.}
                Take $u, v\in L$ such that $u\mathrel{R}v$ and $B_1 = \intervalopenrel{R}{L}{u}{v}$.
                Thus $u\mathrel{R}c$ and $c\mathrel{R}v$ by \cref{intervalopen_rel}.
                Take $d\in L$ such that
                    $(d = u \lor d = a) \land (u\mathrel{R}d \lor u = d) \land (a\mathrel{R}d \lor a = d)$
                    by \cref{max_of_two_elements}.
                Show that $d\in I$.
                \begin{subproof}
                    $d\in L$.
                    $a\mathrel{R}d$ or $d = a$.
                    Show that $d\mathrel{R}b$ or $d = b$.
                    \begin{subproof}
                        $d\mathrel{R}c$ by \cref{intervalopen_rel}.
                        $c\mathrel{R}b$ or $c = b$ by \cref{intervalclosed_rel}.
                        \begin{byCase}
                        \caseOf{$c\mathrel{R}b$.}
                            Thus $d\mathrel{R}b$ by \hyperref[transitive]{transitivity}.
                        \caseOf{$c = b$.}
                            Thus $d\mathrel{R}b$.
                        \end{byCase}
                    \end{subproof}
                    Thus $d\in I$ by \cref{intervalclosed_rel}.
                \end{subproof}
                Show that $\intervalopenclosedrel{R}{I}{d}{c} \subseteq B_0$.
                \begin{subproof}
                    Show that for all $x$ such that $x\in \intervalopenclosedrel{R}{I}{d}{c}$ we have $x\in B_0$.
                    \begin{subproof}
                        Fix $x$.
                        Assume $x\in \intervalopenclosedrel{R}{I}{d}{c}$.
                        Thus $x\in I$ and $d\mathrel{R}x$ and $(x\mathrel{R}c \lor x = c)$ by \cref{intervalopenclosed_rel}.
                        Show that $u\mathrel{R}x$.
                        \begin{subproof}
                            \begin{byCase}
                            \caseOf{$d = u$.} Thus $u\mathrel{R}x$.
                            \caseOf{$d = a$.}
                                $u\mathrel{R}d$ or $u = d$.
                                \begin{byCase}
                                \caseOf{$u = d$.} Thus $u\mathrel{R}x$.
                                \caseOf{$u\mathrel{R}d$.}
                                    Thus $u\mathrel{R}x$ by \hyperref[transitive]{transitivity}.
                                \end{byCase}
                            \end{byCase}
                        \end{subproof}
                        Show that $x\mathrel{R}v$.
                        \begin{subproof}
                            \begin{byCase}
                            \caseOf{$x = c$.} Thus $x\mathrel{R}v$ by \cref{intervalopen_rel}.
                            \caseOf{$x\mathrel{R}c$.}
                                Thus $x\mathrel{R}v$ by \hyperref[transitive]{transitivity}.
                            \end{byCase}
                        \end{subproof}
                        Thus $x\in L$ by \cref{subseteq}.
                        Thus $x\in B_1$ by \cref{intervalopen_rel}.
                        Thus $x\in B_1\inter I$ by \cref{inter_intro}.
                        Thus $x\in B_0$ by \cref{subseteq}.
                    \end{subproof}
                    Thus $\intervalopenclosedrel{R}{I}{d}{c} \subseteq B_0$ by \cref{subseteq}.
                \end{subproof}
                Thus there exists $d\in I$ such that $d\mathrel{R}c$ and
                    $\intervalopenclosedrel{R}{I}{d}{c} \subseteq B_0$.
            \caseOf{$\exists v\in L. \exists u_0\in L. (\forall x\in L. u_0\mathrel{R}x \lor u_0 = x) \land B_1 = \intervalclosedopenrel{R}{L}{u_0}{v}$.}
                Take $u_0, v\in L$ such that
                    $(\forall x\in L. u_0\mathrel{R}x \lor u_0 = x)$ and $B_1 = \intervalclosedopenrel{R}{L}{u_0}{v}$.
                Thus $c\mathrel{R}v$ by \cref{intervalclosedopen_rel}.
                Let $d = a$.
                $d\in I$ by \cref{intervalclosed_rel}.
                Thus $d\mathrel{R}c$ by \cref{intervalclosed_rel}.
                Show that $\intervalopenclosedrel{R}{I}{d}{c} \subseteq B_0$.
                \begin{subproof}
                    Show that for all $x$ such that $x\in \intervalopenclosedrel{R}{I}{d}{c}$ we have $x\in B_0$.
                    \begin{subproof}
                        Fix $x$.
                        Assume $x\in \intervalopenclosedrel{R}{I}{d}{c}$.
                        Thus $x\in I$ and $d\mathrel{R}x$ and $(x\mathrel{R}c \lor x = c)$ by \cref{intervalopenclosed_rel}.
                        Thus $u_0\mathrel{R}x$ or $u_0 = x$ by \cref{subseteq}.
                        Thus $x\mathrel{R}v$ by \hyperref[transitive]{transitivity}.
                        Thus $x\in L$ by \cref{subseteq}.
                        Thus $x\in B_1$ by \cref{intervalclosedopen_rel}.
                        Thus $x\in B_1\inter I$ by \cref{inter_intro}.
                        Thus $x\in B_0$ by \cref{subseteq}.
                    \end{subproof}
                    Thus $\intervalopenclosedrel{R}{I}{d}{c} \subseteq B_0$ by \cref{subseteq}.
                \end{subproof}
                Thus there exists $d\in I$ such that $d\mathrel{R}c$ and
                    $\intervalopenclosedrel{R}{I}{d}{c} \subseteq B_0$.
            \caseOf{$\exists u\in L. \exists v_0\in L. (\forall x\in L. x\mathrel{R}v_0 \lor x = v_0) \land B_1 = \intervalopenclosedrel{R}{L}{u}{v_0}$.}
                Take $u, v_0\in L$ such that
                    $(\forall x\in L. x\mathrel{R}v_0 \lor x = v_0)$ and $B_1 = \intervalopenclosedrel{R}{L}{u}{v_0}$.
                Thus $u\mathrel{R}c$ by \cref{intervalopenclosed_rel}.
                Take $d\in L$ such that
                    $(d = u \lor d = a) \land (u\mathrel{R}d \lor u = d) \land (a\mathrel{R}d \lor a = d)$
                    by \cref{max_of_two_elements}.
                Show that $d\in I$.
                \begin{subproof}
                    $d\in L$.
                    $a\mathrel{R}d$ or $d = a$.
                    Show that $d\mathrel{R}b$ or $d = b$.
                    \begin{subproof}
                        $d\mathrel{R}c$ by \cref{intervalopenclosed_rel}.
                        $c\mathrel{R}b$ or $c = b$ by \cref{intervalclosed_rel}.
                        \begin{byCase}
                        \caseOf{$c\mathrel{R}b$.}
                            Thus $d\mathrel{R}b$ by \hyperref[transitive]{transitivity}.
                        \caseOf{$c = b$.}
                            Thus $d\mathrel{R}b$.
                        \end{byCase}
                    \end{subproof}
                    Thus $d\in I$ by \cref{intervalclosed_rel}.
                \end{subproof}
                Show that $\intervalopenclosedrel{R}{I}{d}{c} \subseteq B_0$.
                \begin{subproof}
                    Show that for all $x$ such that $x\in \intervalopenclosedrel{R}{I}{d}{c}$ we have $x\in B_0$.
                    \begin{subproof}
                        Fix $x$.
                        Assume $x\in \intervalopenclosedrel{R}{I}{d}{c}$.
                        Thus $x\in I$ and $d\mathrel{R}x$ and $(x\mathrel{R}c \lor x = c)$ by \cref{intervalopenclosed_rel}.
                        Show that $u\mathrel{R}x$.
                        \begin{subproof}
                            \begin{byCase}
                            \caseOf{$d = u$.} Thus $u\mathrel{R}x$.
                            \caseOf{$d = a$.}
                                $u\mathrel{R}d$ or $u = d$.
                                \begin{byCase}
                                \caseOf{$u = d$.} Thus $u\mathrel{R}x$.
                                \caseOf{$u\mathrel{R}d$.}
                                    Thus $u\mathrel{R}x$ by \hyperref[transitive]{transitivity}.
                                \end{byCase}
                            \end{byCase}
                        \end{subproof}
                        Thus $x\in L$ by \cref{subseteq}.
                        Thus $x\mathrel{R}v_0$ or $x = v_0$.
                        Thus $x\in B_1$ by \cref{intervalopenclosed_rel}.
                        Thus $x\in B_1\inter I$ by \cref{inter_intro}.
                        Thus $x\in B_0$ by \cref{subseteq}.
                    \end{subproof}
                    Thus $\intervalopenclosedrel{R}{I}{d}{c} \subseteq B_0$ by \cref{subseteq}.
                \end{subproof}
                Thus there exists $d\in I$ such that $d\mathrel{R}c$ and
                    $\intervalopenclosedrel{R}{I}{d}{c} \subseteq B_0$.
            \end{byCase}
        \end{subproof}

        Take $d\in I$ such that $d\mathrel{R}c$ and $\intervalopenclosedrel{R}{I}{d}{c} \subseteq B_0$.
        Show that $d$ is an upper bound of $A_0$ in $L$ with respect to $R$.
        \begin{subproof}
            $d\in L$ by \cref{intervalclosed_rel}.
            Show that for all $x\in A_0$ we have $x\mathrel{R}d$ or $x = d$.
            \begin{subproof}
                Fix $x$.
                Assume $x\in A_0$.
                $c$ is an upper bound of $A_0$ in $L$ with respect to $R$ by \cref{least_upper_bound_rel}.
                Thus $x\mathrel{R}c$ or $x = c$ by \cref{upper_bound_rel}.
                Show that it is wrong that $d\mathrel{R}x$.
                \begin{subproof}
                    Suppose not.
                    Then $d\mathrel{R}x$.
                    \begin{byCase}
                    \caseOf{$x = c$.}
                        Thus $c\in A_0$.
                        Contradiction by \cref{disjoint}.
                    \caseOf{$x\mathrel{R}c$.}
                        Thus $x\in \intervalopenclosedrel{R}{I}{d}{c}$ by \cref{intervalopenclosed_rel}.
                        Thus $x\in B_0$ by \cref{subseteq}.
                        Contradiction by \cref{disjoint}.
                    \end{byCase}
                \end{subproof}
                Thus $x\in L$ by \cref{subseteq}.
                \begin{byCase}
                \caseOf{$x = d$.}
                    Thus $x\mathrel{R}d$ or $x = d$.
                \caseOf{$x \neq d$.}
                    $x\mathrel{R}d$ or $d\mathrel{R}x$ by \cref{connex}.
                    Show that $x\mathrel{R}d$.
                    \begin{subproof}
                        Suppose not.
                        Then $d\mathrel{R}x$.
                        Contradiction.
                    \end{subproof}
                    Thus $x\mathrel{R}d$ or $x = d$.
                \end{byCase}
            \end{subproof}
            Thus $d$ is an upper bound of $A_0$ in $L$ with respect to $R$ by \cref{upper_bound_rel}.
        \end{subproof}
        Thus $c\mathrel{R}d$ or $c = d$ by \cref{least_upper_bound_rel}.
        Contradiction by \cref{asymmetric}.
    \caseOf{$c\in A_0$.}
        Suppose not.
        $c\notin B_0$ by \cref{disjoint}.
        Thus $c\neq b$.

        Take $U\in T$ such that $A_0 = I\inter U$ by \cref{subspace_topology,open_in}.
        Thus $c\in U$ by \cref{inter_elim_right}.
        Take $B_1\in\orderBasis{R}{L}$ such that $c\in B_1$ and $B_1\subseteq U$
            by \cref{topology_generated_by_basis}.
        Thus $B_1\inter I \subseteq A_0$ by \cref{subseteq_inter_iff,subseteq}.

        Show that there exists $e\in I$ such that $c\mathrel{R}e$ and
            $\intervalclosedopenrel{R}{I}{c}{e} \subseteq A_0$.
        \begin{subproof}
            $B_1\in\pow{L}$ and
                $((\exists u, v\in L. u\mathrel{R}v \land B_1 = \intervalopenrel{R}{L}{u}{v}) \lor
                (\exists v\in L. \exists u_0\in L. (\forall x\in L. u_0\mathrel{R}x \lor u_0 = x) \land B_1 = \intervalclosedopenrel{R}{L}{u_0}{v}) \lor
                (\exists u\in L. \exists v_0\in L. (\forall x\in L. x\mathrel{R}v_0 \lor x = v_0) \land B_1 = \intervalopenclosedrel{R}{L}{u}{v_0}))$
                by \cref{order_basis}.
            \begin{byCase}
            \caseOf{$\exists u, v\in L. u\mathrel{R}v \land B_1 = \intervalopenrel{R}{L}{u}{v}$.}
                Take $u, v\in L$ such that $u\mathrel{R}v$ and $B_1 = \intervalopenrel{R}{L}{u}{v}$.
                Thus $u\mathrel{R}c$ and $c\mathrel{R}v$ by \cref{intervalopen_rel}.
                Take $e\in L$ such that
                    $(e = v \lor e = b) \land (e\mathrel{R}v \lor e = v) \land (e\mathrel{R}b \lor e = b)$
                    by \cref{min_of_two_elements}.
                Show that $e\in I$.
                \begin{subproof}
                    $e\in L$.
                    $e\mathrel{R}b$ or $e = b$.
                    $c\mathrel{R}b$ by \cref{intervalclosed_rel}.
                    Show that $c\mathrel{R}e$.
                    \begin{subproof}
                        \begin{byCase}
                        \caseOf{$e = v$.} Thus $c\mathrel{R}e$ by \cref{intervalopen_rel}.
                        \caseOf{$e = b$.} Thus $c\mathrel{R}e$.
                        \end{byCase}
                    \end{subproof}
                    $a\mathrel{R}c$ or $a = c$ by \cref{intervalclosed_rel}.
                    Show that $a\mathrel{R}e$ or $e = a$.
                    \begin{subproof}
                        \begin{byCase}
                        \caseOf{$a = c$.} Thus $a\mathrel{R}e$.
                        \caseOf{$a\mathrel{R}c$.}
                            Thus $a\mathrel{R}e$ by \hyperref[transitive]{transitivity}.
                        \end{byCase}
                    \end{subproof}
                    Thus $e\in I$ by \cref{intervalclosed_rel}.
                \end{subproof}
                Show that $\intervalclosedopenrel{R}{I}{c}{e} \subseteq A_0$.
                \begin{subproof}
                    Show that for all $x$ such that $x\in \intervalclosedopenrel{R}{I}{c}{e}$ we have $x\in A_0$.
                    \begin{subproof}
                        Fix $x$.
                        Assume $x\in \intervalclosedopenrel{R}{I}{c}{e}$.
                        Thus $x\in I$ and $(c\mathrel{R}x \lor x = c)$ and $x\mathrel{R}e$ by \cref{intervalclosedopen_rel}.
                        Show that $x\mathrel{R}v$.
                        \begin{subproof}
                            \begin{byCase}
                            \caseOf{$e = v$.} Thus $x\mathrel{R}v$.
                            \caseOf{$e = b$.}
                                $e\mathrel{R}v$ or $e = v$.
                                \begin{byCase}
                                \caseOf{$e = v$.} Thus $x\mathrel{R}v$.
                                \caseOf{$e\mathrel{R}v$.}
                                    Thus $x\mathrel{R}v$ by \hyperref[transitive]{transitivity}.
                                \end{byCase}
                            \end{byCase}
                        \end{subproof}
                        Show that $u\mathrel{R}x$.
                        \begin{subproof}
                            \begin{byCase}
                            \caseOf{$x = c$.} Thus $u\mathrel{R}x$ by \cref{intervalopen_rel}.
                            \caseOf{$c\mathrel{R}x$.}
                                Thus $u\mathrel{R}x$ by \hyperref[transitive]{transitivity}.
                            \end{byCase}
                        \end{subproof}
                        Thus $x\in L$ by \cref{subseteq}.
                        Thus $x\in B_1$ by \cref{intervalopen_rel}.
                        Thus $x\in B_1\inter I$ by \cref{inter_intro}.
                        Thus $x\in A_0$ by \cref{subseteq}.
                    \end{subproof}
                    Thus $\intervalclosedopenrel{R}{I}{c}{e} \subseteq A_0$ by \cref{subseteq}.
                \end{subproof}
                Thus there exists $e\in I$ such that $c\mathrel{R}e$ and
                    $\intervalclosedopenrel{R}{I}{c}{e} \subseteq A_0$.
            \caseOf{$\exists v\in L. \exists u_0\in L. (\forall x\in L. u_0\mathrel{R}x \lor u_0 = x) \land B_1 = \intervalclosedopenrel{R}{L}{u_0}{v}$.}
                Take $u_0, v\in L$ such that
                    $(\forall x\in L. u_0\mathrel{R}x \lor u_0 = x)$ and $B_1 = \intervalclosedopenrel{R}{L}{u_0}{v}$.
                Thus $c\mathrel{R}v$ by \cref{intervalclosedopen_rel}.
                Take $e\in L$ such that
                    $(e = v \lor e = b) \land (e\mathrel{R}v \lor e = v) \land (e\mathrel{R}b \lor e = b)$
                    by \cref{min_of_two_elements}.
                Show that $e\in I$.
                \begin{subproof}
                    $e\in L$.
                    $e\mathrel{R}b$ or $e = b$.
                    $c\mathrel{R}b$ by \cref{intervalclosed_rel}.
                    Show that $c\mathrel{R}e$.
                    \begin{subproof}
                        \begin{byCase}
                        \caseOf{$e = v$.} Thus $c\mathrel{R}e$ by \cref{intervalclosedopen_rel}.
                        \caseOf{$e = b$.} Thus $c\mathrel{R}e$.
                        \end{byCase}
                    \end{subproof}
                    $a\mathrel{R}c$ or $a = c$ by \cref{intervalclosed_rel}.
                    Show that $a\mathrel{R}e$ or $e = a$.
                    \begin{subproof}
                        \begin{byCase}
                        \caseOf{$a = c$.} Thus $a\mathrel{R}e$.
                        \caseOf{$a\mathrel{R}c$.}
                            Thus $a\mathrel{R}e$ by \hyperref[transitive]{transitivity}.
                        \end{byCase}
                    \end{subproof}
                    Thus $e\in I$ by \cref{intervalclosed_rel}.
                \end{subproof}
                Show that $\intervalclosedopenrel{R}{I}{c}{e} \subseteq A_0$.
                \begin{subproof}
                    Show that for all $x$ such that $x\in \intervalclosedopenrel{R}{I}{c}{e}$ we have $x\in A_0$.
                    \begin{subproof}
                        Fix $x$.
                        Assume $x\in \intervalclosedopenrel{R}{I}{c}{e}$.
                        Thus $x\in I$ and $(c\mathrel{R}x \lor x = c)$ and $x\mathrel{R}e$ by \cref{intervalclosedopen_rel}.
                        Thus $x\mathrel{R}v$ by \hyperref[transitive]{transitivity}.
                        Thus $x\in L$ by \cref{subseteq}.
                        Thus $u_0\mathrel{R}x$ or $u_0 = x$.
                        Thus $x\in B_1$ by \cref{intervalclosedopen_rel}.
                        Thus $x\in B_1\inter I$ by \cref{inter_intro}.
                        Thus $x\in A_0$ by \cref{subseteq}.
                    \end{subproof}
                    Thus $\intervalclosedopenrel{R}{I}{c}{e} \subseteq A_0$ by \cref{subseteq}.
                \end{subproof}
                Thus there exists $e\in I$ such that $c\mathrel{R}e$ and
                    $\intervalclosedopenrel{R}{I}{c}{e} \subseteq A_0$.
            \caseOf{$\exists u\in L. \exists v_0\in L. (\forall x\in L. x\mathrel{R}v_0 \lor x = v_0) \land B_1 = \intervalopenclosedrel{R}{L}{u}{v_0}$.}
                Take $u, v_0\in L$ such that
                    $(\forall x\in L. x\mathrel{R}v_0 \lor x = v_0)$ and $B_1 = \intervalopenclosedrel{R}{L}{u}{v_0}$.
                Thus $u\mathrel{R}c$ by \cref{intervalopenclosed_rel}.
                Take $e\in L$ such that
                    $(e = v_0 \lor e = b) \land (e\mathrel{R}v_0 \lor e = v_0) \land (e\mathrel{R}b \lor e = b)$
                    by \cref{min_of_two_elements}.
                Show that $c\mathrel{R}e$.
                \begin{subproof}
                    \begin{byCase}
                    \caseOf{$e = b$.}
                        $c\mathrel{R}b$ or $c = b$ by \cref{intervalclosed_rel}.
                        $c \neq b$ by \cref{disjoint}.
                        Thus $c\mathrel{R}e$.
                    \caseOf{$e = v_0$.}
                        $c\in L$ by \cref{subseteq}.
                        Thus $c\mathrel{R}v_0$ or $c = v_0$.
                        Show that $c\mathrel{R}v_0$.
                        \begin{subproof}
                            Suppose not.
                            Then $c = v_0$.
                            $b\mathrel{R}v_0$ or $b = v_0$.
                            \begin{byCase}
                            \caseOf{$b = v_0$.} Contradiction by \cref{disjoint}.
                            \caseOf{$b\mathrel{R}v_0$.}
                                $c\mathrel{R}b$ or $c = b$ by \cref{intervalclosed_rel}.
                                \begin{byCase}
                                \caseOf{$c = b$.} Contradiction by \cref{disjoint}.
                                \caseOf{$c\mathrel{R}b$.}
                                    Thus $b\mathrel{R}c$.
                                    Contradiction by \cref{asymmetric}.
                                \end{byCase}
                            \end{byCase}
                        \end{subproof}
                        Thus $c\mathrel{R}e$.
                    \end{byCase}
                \end{subproof}
                Show that $e\in I$.
                \begin{subproof}
                    $e\in L$.
                    $e\mathrel{R}b$ or $e = b$.
                    $a\mathrel{R}c$ or $a = c$ by \cref{intervalclosed_rel}.
                    Show that $a\mathrel{R}e$ or $e = a$.
                    \begin{subproof}
                        \begin{byCase}
                        \caseOf{$a = c$.} Thus $a\mathrel{R}e$.
                        \caseOf{$a\mathrel{R}c$.}
                            Thus $a\mathrel{R}e$ by \hyperref[transitive]{transitivity}.
                        \end{byCase}
                    \end{subproof}
                    Thus $e\in I$ by \cref{intervalclosed_rel}.
                \end{subproof}
                Show that $\intervalclosedopenrel{R}{I}{c}{e} \subseteq A_0$.
                \begin{subproof}
                    Show that for all $x$ such that $x\in \intervalclosedopenrel{R}{I}{c}{e}$ we have $x\in A_0$.
                    \begin{subproof}
                        Fix $x$.
                        Assume $x\in \intervalclosedopenrel{R}{I}{c}{e}$.
                        Thus $x\in I$ and $(c\mathrel{R}x \lor x = c)$ and $x\mathrel{R}e$ by \cref{intervalclosedopen_rel}.
                        Show that $u\mathrel{R}x$.
                        \begin{subproof}
                            \begin{byCase}
                            \caseOf{$x = c$.} Thus $u\mathrel{R}x$ by \cref{intervalopenclosed_rel}.
                            \caseOf{$c\mathrel{R}x$.}
                                Thus $u\mathrel{R}x$ by \hyperref[transitive]{transitivity}.
                            \end{byCase}
                        \end{subproof}
                        Show that $x\mathrel{R}v_0$ or $x = v_0$.
                        \begin{subproof}
                            $e\mathrel{R}v_0$ or $e = v_0$.
                            \begin{byCase}
                            \caseOf{$e = v_0$.} Thus $x\mathrel{R}v_0$.
                            \caseOf{$e\mathrel{R}v_0$.}
                                Thus $x\mathrel{R}v_0$ by \hyperref[transitive]{transitivity}.
                            \end{byCase}
                        \end{subproof}
                        Thus $x\in L$ by \cref{subseteq}.
                        Thus $x\in B_1$ by \cref{intervalopenclosed_rel}.
                        Thus $x\in B_1\inter I$ by \cref{inter_intro}.
                        Thus $x\in A_0$ by \cref{subseteq}.
                    \end{subproof}
                    Thus $\intervalclosedopenrel{R}{I}{c}{e} \subseteq A_0$ by \cref{subseteq}.
                \end{subproof}
                Thus there exists $e\in I$ such that $c\mathrel{R}e$ and
                    $\intervalclosedopenrel{R}{I}{c}{e} \subseteq A_0$.
            \end{byCase}
        \end{subproof}

        Take $e\in I$ such that $c\mathrel{R}e$ and $\intervalclosedopenrel{R}{I}{c}{e} \subseteq A_0$.
        $c\in L$ by \cref{subseteq}.
        $e\in L$ by \cref{subseteq}.
        Take $z\in L$ such that $c\mathrel{R}z$ and $z\mathrel{R}e$ by \cref{linear_continuum}.
        Show that $z\in I$.
        \begin{subproof}
            $a\mathrel{R}c$ or $a = c$ by \cref{intervalclosed_rel}.
            Show that $a\mathrel{R}z$ or $z = a$.
            \begin{subproof}
                \begin{byCase}
                \caseOf{$a = c$.} Thus $a\mathrel{R}z$.
                \caseOf{$a\mathrel{R}c$.}
                    Thus $a\mathrel{R}z$ by \hyperref[transitive]{transitivity}.
                \end{byCase}
            \end{subproof}
            $e\mathrel{R}b$ or $e = b$ by \cref{intervalclosed_rel}.
            Show that $z\mathrel{R}b$ or $z = b$.
            \begin{subproof}
                \begin{byCase}
                \caseOf{$e = b$.} Thus $z\mathrel{R}b$.
                \caseOf{$e\mathrel{R}b$.}
                    Thus $z\mathrel{R}b$ by \hyperref[transitive]{transitivity}.
                \end{byCase}
            \end{subproof}
            Thus $z\in I$ by \cref{intervalclosed_rel}.
        \end{subproof}
        Thus $z\in \intervalclosedopenrel{R}{I}{c}{e}$ by \cref{intervalclosedopen_rel}.
        Thus $z\in A_0$ by \cref{subseteq}.
        $c$ is an upper bound of $A_0$ in $L$ with respect to $R$ by \cref{least_upper_bound_rel}.
        Thus $z\mathrel{R}c$ or $z = c$ by \cref{upper_bound_rel}.
        Contradiction by \cref{asymmetric}.
    \end{byCase}
\end{proof}

% from §24 helper lemma 10: convex subset of a linear continuum is connected
\begin{lemma}\label{convex_connected_linear_continuum}
    Assume $T$ is the order topology on $L$ with respect to $R$.
    Assume $L$ is a linear continuum with respect to $R$.
    Suppose $Y$ is convex in $L$ with respect to $R$.
    Let $T_Y = \subspaceTopology{T}{Y}$.
    Then $Y$ is connected with respect to $T_Y$.
\end{lemma}
\begin{proof}
    Show that there exists no $A$ such that there exists $B$ such that
        $Y$ is separated by $A$ and $B$ with respect to $T_Y$.
    \begin{subproof}
        Suppose not.
        Then there exists $A$ such that there exists $B$ such that
            $Y$ is separated by $A$ and $B$ with respect to $T_Y$.
        Take $A$ such that there exists $B$ such that
            $Y$ is separated by $A$ and $B$ with respect to $T_Y$.
        Take $B$ such that $Y$ is separated by $A$ and $B$ with respect to $T_Y$.
        Take $a$ such that $a\in A$ by \cref{nonempty_inhabited,separation}.
        Take $b$ such that $b\in B$ by \cref{nonempty_inhabited,separation}.
        $R$ is a simple order relation on $L$ by \cref{linear_continuum}.
        $R$ is connex on $L$ by \cref{simple_order_relation}.
        $Y\subseteq L$ by \cref{convex_in}.
        $A\union B = Y$ by \cref{separation}.
        $A\subseteq A\union B$ and $B\subseteq A\union B$ by \cref{union_upper_left,union_upper_right}.
        Thus $A\subseteq Y$ and $B\subseteq Y$ by \cref{subseteq_transitive}.
        Thus $A\subseteq L$ and $B\subseteq L$ by \cref{subseteq_transitive}.
        Thus $a\in L$ and $b\in L$ by \cref{subseteq}.
        $A$ is disjoint from $B$ by \cref{separation}.
        Thus $a\neq b$ by \cref{disjoint}.
        Thus $a\mathrel{R}b$ or $b\mathrel{R}a$ by \cref{connex}.
        \begin{byCase}
        \caseOf{$a\mathrel{R}b$.}
            Contradiction by \cref{convex_separation_contradiction}.
        \caseOf{$b\mathrel{R}a$.}
            Show that $Y$ is separated by $B$ and $A$ with respect to $T_Y$.
            \begin{subproof}
                $A$ is open in $Y$ with respect to $T_Y$ by \cref{separation}.
                $B$ is open in $Y$ with respect to $T_Y$ by \cref{separation}.
                $A\neq\emptyset$ and $B\neq\emptyset$ by \cref{separation}.
                $A$ is disjoint from $B$ by \cref{separation}.
                $B$ is disjoint from $A$ by \cref{disjoint_symmetric}.
                $A\union B = Y$ by \cref{separation}.
                $B\union A = Y$ by \cref{union_comm}.
                Thus $Y$ is separated by $B$ and $A$ with respect to $T_Y$ by \cref{separation}.
            \end{subproof}
            Contradiction by \cref{convex_separation_contradiction}.
        \end{byCase}
    \end{subproof}
    Thus $Y$ is connected with respect to $T_Y$ by \cref{connected}.
\end{proof}

% from §24 Item 2: linear continuum connected, intervals and rays connected
\begin{theorem}\label{linear_continuum_connected}
    Assume $T$ is the order topology on $L$ with respect to $R$.
    Assume $L$ is a linear continuum with respect to $R$.
    Then $L$ is connected with respect to $T$ and
        for all $a,b\in L$ if $a\mathrel{R}b$ then
            $\intervalopenrel{R}{L}{a}{b}$ is connected with respect to $\subspaceTopology{T}{\intervalopenrel{R}{L}{a}{b}}$ and
            $\intervalopenclosedrel{R}{L}{a}{b}$ is connected with respect to $\subspaceTopology{T}{\intervalopenclosedrel{R}{L}{a}{b}}$ and
            $\intervalclosedopenrel{R}{L}{a}{b}$ is connected with respect to $\subspaceTopology{T}{\intervalclosedopenrel{R}{L}{a}{b}}$ and
            $\intervalclosedrel{R}{L}{a}{b}$ is connected with respect to $\subspaceTopology{T}{\intervalclosedrel{R}{L}{a}{b}}$ and
        for all $a\in L$ we have
            $\openrayright{R}{L}{a}$ is connected with respect to $\subspaceTopology{T}{\openrayright{R}{L}{a}}$ and
            $\openrayleft{R}{L}{a}$ is connected with respect to $\subspaceTopology{T}{\openrayleft{R}{L}{a}}$ and
            $\closedrayright{R}{L}{a}$ is connected with respect to $\subspaceTopology{T}{\closedrayright{R}{L}{a}}$ and
            $\closedrayleft{R}{L}{a}$ is connected with respect to $\subspaceTopology{T}{\closedrayleft{R}{L}{a}}$.
\end{theorem}
\begin{proof}
    $R$ is a simple order relation on $L$ by \cref{linear_continuum}.
    $L$ has more than one element by \cref{linear_continuum}.
    $T = \basisTopology{\orderBasis{R}{L}}{L}$ by \cref{order_topology}.
    $\orderBasis{R}{L}$ is a basis on $L$ by \cref{order_basis_is_basis}.
    Thus $T$ is a topology on $L$ by \cref{basis_topology_is_topology}.
    Show that $L$ is connected with respect to $T$.
    \begin{subproof}
        $L$ is convex in $L$ with respect to $R$ by \cref{convex_in,subseteq}.
        Let $T_L = \subspaceTopology{T}{L}$.
        Thus $L$ is connected with respect to $T_L$ by \cref{convex_connected_linear_continuum}.
        Show that $T_L = T$.
        \begin{subproof}
            Show that $T_L\subseteq T$.
            \begin{subproof}
                Show that for all $U$ such that $U\in T_L$ we have $U\in T$.
                \begin{subproof}
                    Fix $U$.
                    Assume $U\in T_L$.
                    Take $V\in T$ such that $U = L\inter V$ by \cref{subspace_topology}.
                    $T\subseteq\pow{L}$ by \cref{topology_on}.
                    Thus $V\in\pow{L}$ by \cref{subseteq}.
                    Thus $V\subseteq L$ by \cref{powerset_elim}.
                    Thus $U = V$ by \cref{inter_absorb_supseteq_left}.
                    Thus $U\in T$.
                \end{subproof}
                Thus $T_L\subseteq T$ by \cref{subseteq}.
            \end{subproof}
            Show that $T\subseteq T_L$.
            \begin{subproof}
                Show that for all $U$ such that $U\in T$ we have $U\in T_L$.
                \begin{subproof}
                    Fix $U$.
                    Assume $U\in T$.
                    $T\subseteq\pow{L}$ by \cref{topology_on}.
                    Thus $U\in\pow{L}$ by \cref{subseteq}.
                    Thus $U\subseteq L$ by \cref{powerset_elim}.
                    Thus $U = L\inter U$ by \cref{inter_absorb_supseteq_left}.
                    Thus $U\in T_L$ by \cref{subspace_topology}.
                \end{subproof}
                Thus $T\subseteq T_L$ by \cref{subseteq}.
            \end{subproof}
            Thus $T_L = T$ by \cref{subseteq_antisymmetric}.
        \end{subproof}
        Thus $L$ is connected with respect to $T$ by \cref{subseteq}.
    \end{subproof}

    Show that for all $a,b\in L$ if $a\mathrel{R}b$ then
        $\intervalopenrel{R}{L}{a}{b}$ is connected with respect to $\subspaceTopology{T}{\intervalopenrel{R}{L}{a}{b}}$ and
        $\intervalopenclosedrel{R}{L}{a}{b}$ is connected with respect to $\subspaceTopology{T}{\intervalopenclosedrel{R}{L}{a}{b}}$ and
        $\intervalclosedopenrel{R}{L}{a}{b}$ is connected with respect to $\subspaceTopology{T}{\intervalclosedopenrel{R}{L}{a}{b}}$ and
        $\intervalclosedrel{R}{L}{a}{b}$ is connected with respect to $\subspaceTopology{T}{\intervalclosedrel{R}{L}{a}{b}}$.
    \begin{subproof}
        Fix $a$.
        Fix $b$.
        Assume $a\in L$ and $b\in L$.
        $R$ is transitive by \cref{simple_order_relation}.
        Assume $a\mathrel{R}b$.
        Let $I_1 = \intervalopenrel{R}{L}{a}{b}$.
        Let $I_2 = \intervalopenclosedrel{R}{L}{a}{b}$.
        Let $I_3 = \intervalclosedopenrel{R}{L}{a}{b}$.
        Let $I_4 = \intervalclosedrel{R}{L}{a}{b}$.
        $I_1$ is convex in $L$ with respect to $R$ by \cref{intervalopenrel_convex}.
        $I_2$ is convex in $L$ with respect to $R$ by \cref{intervalopenclosedrel_convex}.
        $I_3$ is convex in $L$ with respect to $R$ by \cref{intervalclosedopenrel_convex}.
        $I_4$ is convex in $L$ with respect to $R$ by \cref{intervalclosedrel_convex}.
        Thus $I_1$ is connected with respect to $\subspaceTopology{T}{I_1}$ by \cref{convex_connected_linear_continuum}.
        Thus $I_2$ is connected with respect to $\subspaceTopology{T}{I_2}$ by \cref{convex_connected_linear_continuum}.
        Thus $I_3$ is connected with respect to $\subspaceTopology{T}{I_3}$ by \cref{convex_connected_linear_continuum}.
        Thus $I_4$ is connected with respect to $\subspaceTopology{T}{I_4}$ by \cref{convex_connected_linear_continuum}.
    \end{subproof}

    Show that for all $a\in L$ we have
        $\openrayright{R}{L}{a}$ is connected with respect to $\subspaceTopology{T}{\openrayright{R}{L}{a}}$ and
        $\openrayleft{R}{L}{a}$ is connected with respect to $\subspaceTopology{T}{\openrayleft{R}{L}{a}}$ and
        $\closedrayright{R}{L}{a}$ is connected with respect to $\subspaceTopology{T}{\closedrayright{R}{L}{a}}$ and
        $\closedrayleft{R}{L}{a}$ is connected with respect to $\subspaceTopology{T}{\closedrayleft{R}{L}{a}}$.
    \begin{subproof}
        Fix $a\in L$.
        $R$ is transitive by \cref{simple_order_relation}.
        Let $R_1 = \openrayright{R}{L}{a}$.
        Let $R_2 = \openrayleft{R}{L}{a}$.
        Let $R_3 = \closedrayright{R}{L}{a}$.
        Let $R_4 = \closedrayleft{R}{L}{a}$.
        $R_1$ is convex in $L$ with respect to $R$ by \cref{openrayright_convex}.
        $R_2$ is convex in $L$ with respect to $R$ by \cref{openrayleft_convex}.
        $R_3$ is convex in $L$ with respect to $R$ by \cref{closedrayright_convex}.
        $R_4$ is convex in $L$ with respect to $R$ by \cref{closedrayleft_convex}.
        Thus $R_1$ is connected with respect to $\subspaceTopology{T}{R_1}$ by \cref{convex_connected_linear_continuum}.
        Thus $R_2$ is connected with respect to $\subspaceTopology{T}{R_2}$ by \cref{convex_connected_linear_continuum}.
        Thus $R_3$ is connected with respect to $\subspaceTopology{T}{R_3}$ by \cref{convex_connected_linear_continuum}.
        Thus $R_4$ is connected with respect to $\subspaceTopology{T}{R_4}$ by \cref{convex_connected_linear_continuum}.
    \end{subproof}
\end{proof}

% from §24 helper definition 4: real order relation
\begin{definition}\label{real_order_relation}
    $\realOrderRel = \{w\in\reals\times\reals \mid \fst{w} < \snd{w}\}$.
\end{definition}

% from §24 helper definition 5: order topology on reals
\begin{definition}\label{real_order_topology}
    $\realOrderTopology = \basisTopology{\orderBasis{\realOrderRel}{\reals}}{\reals}$.
\end{definition}

% from §24 helper definition 6: between set in an ordered set
\begin{definition}\label{between_rel}
    $\betweenrel{R}{Y}{x}{y} = \intervalclosedrel{R}{Y}{x}{y} \union \intervalclosedrel{R}{Y}{y}{x}$.
\end{definition}

% from §24 helper lemma 11: real order relation intro
\begin{lemma}\label{real_order_relation_intro}
    Suppose $a\in\reals$ and $b\in\reals$ and $a < b$.
    Then $a\mathrel{\realOrderRel} b$.
\end{lemma}
\begin{proof}
    Thus $(a,b)\in\reals\times\reals$ by \cref{times_tuple_intro}.
    $\fst{(a,b)} = a$ and $\snd{(a,b)} = b$ by \cref{fst_eq,snd_eq}.
    Thus $\fst{(a,b)} < b$ by \cref{real_lt_subst_left}.
    Thus $b = \snd{(a,b)}$ by \cref{snd_eq}.
    Thus $\fst{(a,b)} < \snd{(a,b)}$ by \cref{real_lt_subst_right}.
    Thus $(a,b)\in \realOrderRel$ by \cref{real_order_relation}.
    Thus $a\mathrel{\realOrderRel} b$.
\end{proof}

% from §24 helper lemma 12: real order relation elim
\begin{lemma}\label{real_order_relation_elim}
    Suppose $a\mathrel{\realOrderRel} b$.
    Then $a\in\reals$ and $b\in\reals$ and $a < b$.
\end{lemma}
\begin{proof}
    Thus $(a,b)\in \realOrderRel$.
    Thus $(a,b)\in\reals\times\reals$ and $\fst{(a,b)} < \snd{(a,b)}$ by \cref{real_order_relation}.
    Thus $a\in\reals$ and $b\in\reals$ by \cref{times_tuple_elim}.
    $\fst{(a,b)} = a$ and $\snd{(a,b)} = b$ by \cref{fst_eq,snd_eq}.
    Thus $a = \fst{(a,b)}$ by \cref{fst_eq}.
    Thus $a < \snd{(a,b)}$ by \cref{real_lt_subst_left}.
    Thus $a < b$ by \cref{real_lt_subst_right}.
\end{proof}

% from §24 helper lemma 13: real order relation is simple
\begin{lemma}\label{real_order_relation_simple}
    $\realOrderRel$ is a simple order relation on $\reals$.
\end{lemma}
\begin{proof}
    Show that $\realOrderRel$ is a binary relation on $\reals$.
    \begin{subproof}
        Show that for all $w$ such that $w\in \realOrderRel$ we have $w\in\reals\times\reals$.
        \begin{subproof}
            Fix $w$.
            Assume $w\in \realOrderRel$.
            Thus $w\in\reals\times\reals$ by \cref{real_order_relation}.
        \end{subproof}
        Thus $\realOrderRel\subseteq\reals\times\reals$ by \cref{subseteq}.
        Thus $\realOrderRel$ is a binary relation on $\reals$.
    \end{subproof}
    Show that $\realOrderRel$ is transitive.
    \begin{subproof}
        Show that for all $a,b,c$ such that $a\mathrel{\realOrderRel} b$ and $b\mathrel{\realOrderRel} c$ we have
            $a\mathrel{\realOrderRel} c$.
        \begin{subproof}
            Fix $a$.
            Fix $b$.
            Fix $c$.
            Assume $a\mathrel{\realOrderRel} b$ and $b\mathrel{\realOrderRel} c$.
            Thus $a\in\reals$ and $b\in\reals$ and $a < b$ by \cref{real_order_relation_elim}.
            Thus $c\in\reals$ and $b < c$ by \cref{real_order_relation_elim}.
            Thus $a < c$ by \cref{real_lt_trans}.
            Thus $a\mathrel{\realOrderRel} c$ by \cref{real_order_relation_intro}.
        \end{subproof}
        Thus $\realOrderRel$ is transitive by \cref{transitive}.
    \end{subproof}
    Show that $\realOrderRel$ is asymmetric.
    \begin{subproof}
        Show that for all $a,b$ such that $a\mathrel{\realOrderRel} b$ we have it is wrong that $b\mathrel{\realOrderRel} a$.
        \begin{subproof}
            Fix $a$.
            Fix $b$.
            Assume $a\mathrel{\realOrderRel} b$.
            Thus $a\in\reals$ and $b\in\reals$ and $a < b$ by \cref{real_order_relation_elim}.
            Show that it is wrong that $b\mathrel{\realOrderRel} a$.
            \begin{subproof}
                Suppose not.
                Then $b\mathrel{\realOrderRel} a$.
                Thus $b < a$ by \cref{real_order_relation_elim}.
                Show that $a \neq b$.
                \begin{subproof}
                    Suppose not.
                    Then $a = b$.
                    Thus $a < a$.
                    Contradiction by \cref{real_lt_irreflexive}.
                \end{subproof}
                Contradiction by \cref{real_lt_asym}.
            \end{subproof}
        \end{subproof}
        Thus $\realOrderRel$ is asymmetric by \cref{asymmetric}.
    \end{subproof}
    Show that $\realOrderRel$ is connex on $\reals$.
    \begin{subproof}
        Show that for all $a,b$ such that $a\in\reals$ and $b\in\reals$ and $a\neq b$ we have
            $a\mathrel{\realOrderRel} b$ or $b\mathrel{\realOrderRel} a$.
        \begin{subproof}
            Fix $a$.
            Fix $b$.
            Assume $a\in\reals$ and $b\in\reals$ and $a\neq b$.
            $a < b$ or $b < a$ by \cref{real_lt_trichotomy}.
            Thus $a\mathrel{\realOrderRel} b$ or $b\mathrel{\realOrderRel} a$ by \cref{real_order_relation_intro}.
        \end{subproof}
        Thus $\realOrderRel$ is connex on $\reals$ by \cref{connex}.
    \end{subproof}
    Thus $\realOrderRel$ is a simple order relation on $\reals$ by \cref{simple_order_relation}.
\end{proof}

% from §24 helper lemma 14: inhabited implies nonempty
\begin{lemma}\label{inhabited_nonempty}
    Suppose $A$ is inhabited.
    Then $A\neq\emptyset$.
\end{lemma}
\begin{proof}
    Take $a$ such that $a\in A$.
    Suppose not.
    Then $A = \emptyset$.
    Thus $a\notin A$ by \cref{emptyset}.
    Contradiction.
\end{proof}

% from §24 helper lemma 15: real line is a linear continuum
\begin{lemma}\label{reals_linear_continuum}
    $\reals$ is a linear continuum with respect to $\realOrderRel$.
\end{lemma}
\begin{proof}
    $\realOrderRel$ is a simple order relation on $\reals$ by \cref{real_order_relation_simple}.
    Show that $\reals$ has more than one element.
    \begin{subproof}
        $0\in\reals$ by \cref{real_add_ident}.
        $1\in\reals$ by \cref{real_mul_ident}.
        $0 < 1$ by \cref{real_zero_lt_one}.
        Show that $0 \neq 1$.
        \begin{subproof}
            Suppose not.
            Then $0 = 1$.
            Thus $0 < 0$.
            Contradiction by \cref{real_lt_irreflexive}.
        \end{subproof}
        Thus $\reals$ has more than one element by \cref{more_than_one_element}.
    \end{subproof}
    Show that $\reals$ has the least upper bound property with respect to $\realOrderRel$.
    \begin{subproof}
        Show that for all $A$ such that $A\subseteq\reals$ and $A$ is inhabited and
            there exists $u$ such that $u$ is an upper bound of $A$ in $\reals$ with respect to $\realOrderRel$
            there exists $c$ such that $c$ is a least upper bound of $A$ in $\reals$ with respect to $\realOrderRel$.
        \begin{subproof}
            Fix $A$.
            Assume $A\subseteq\reals$ and $A$ is inhabited and
                there exists $u$ such that $u$ is an upper bound of $A$ in $\reals$ with respect to $\realOrderRel$.
            Take $u$ such that $u$ is an upper bound of $A$ in $\reals$ with respect to $\realOrderRel$.
        Show that $A$ is bounded above.
        \begin{subproof}
            $u\in\reals$ by \cref{upper_bound_rel}.
            Show that for all $x\in A$ we have $x \leq u$.
            \begin{subproof}
                Fix $x$.
                Assume $x\in A$.
                $x\mathrel{\realOrderRel} u$ or $x = u$ by \cref{upper_bound_rel}.
                \begin{byCase}
                \caseOf{$x = u$.} Thus $x \leq u$.
                \caseOf{$x\mathrel{\realOrderRel} u$.}
                    Thus $x < u$ by \cref{real_order_relation_elim}.
                    Thus $x \leq u$.
                \end{byCase}
            \end{subproof}
            Thus $u$ is an upper bound of $A$ by \cref{real_upper_bound}.
            Thus $A$ is bounded above by \cref{real_bounded_above}.
        \end{subproof}
            Thus $A\neq\emptyset$ by \cref{inhabited_nonempty}.
            Take $p$ such that $p$ is a supremum of $A$ by \cref{real_completeness_axiom}.
            Show that $p$ is a least upper bound of $A$ in $\reals$ with respect to $\realOrderRel$.
            \begin{subproof}
                Show that $p$ is an upper bound of $A$ in $\reals$ with respect to $\realOrderRel$.
                \begin{subproof}
                    $p$ is an upper bound of $A$ by \cref{real_supremum}.
                    $p\in\reals$ by \cref{real_upper_bound}.
                    Show that for all $x\in A$ we have $x\mathrel{\realOrderRel} p$ or $x = p$.
                    \begin{subproof}
                        Fix $x$.
                        Assume $x\in A$.
                        Thus $x \leq p$ by \cref{real_upper_bound}.
                        \begin{byCase}
                        \caseOf{$x = p$.} Thus $x = p$.
                        \caseOf{$x \neq p$.}
                            Thus $x < p$.
                            Thus $x\in\reals$ by \cref{subseteq}.
                            Thus $x\mathrel{\realOrderRel} p$ by \cref{real_order_relation_intro}.
                        \end{byCase}
                    \end{subproof}
                    Thus $p$ is an upper bound of $A$ in $\reals$ with respect to $\realOrderRel$ by \cref{upper_bound_rel}.
                \end{subproof}
                Show that for all $c$ such that $c$ is an upper bound of $A$ in $\reals$ with respect to $\realOrderRel$ we have
                    $p\mathrel{\realOrderRel} c$ or $p = c$.
                \begin{subproof}
                    Fix $c$.
                    Assume $c$ is an upper bound of $A$ in $\reals$ with respect to $\realOrderRel$.
                    $c\in\reals$ by \cref{upper_bound_rel}.
                    $p$ is an upper bound of $A$ by \cref{real_supremum}.
                    $p\in\reals$ by \cref{real_upper_bound}.
                    Show that $c$ is an upper bound of $A$.
                    \begin{subproof}
                        Show that for all $x\in A$ we have $x \leq c$.
                        \begin{subproof}
                            Fix $x$.
                            Assume $x\in A$.
                            $x\mathrel{\realOrderRel} c$ or $x = c$ by \cref{upper_bound_rel}.
                            \begin{byCase}
                            \caseOf{$x = c$.} Thus $x \leq c$.
                            \caseOf{$x\mathrel{\realOrderRel} c$.}
                                Thus $x < c$ by \cref{real_order_relation_elim}.
                                Thus $x \leq c$.
                            \end{byCase}
                        \end{subproof}
                        Thus $c$ is an upper bound of $A$ by \cref{real_upper_bound}.
                    \end{subproof}
                    Thus $p \leq c$ by \cref{real_supremum}.
                    \begin{byCase}
                    \caseOf{$p = c$.} Thus $p = c$.
                    \caseOf{$p \neq c$.}
                        Thus $p < c$.
                        Thus $p\mathrel{\realOrderRel} c$ by \cref{real_order_relation_intro}.
                    \end{byCase}
                \end{subproof}
                Thus $p$ is a least upper bound of $A$ in $\reals$ with respect to $\realOrderRel$ by \cref{least_upper_bound_rel}.
            \end{subproof}
            Thus there exists $c$ such that $c$ is a least upper bound of $A$ in $\reals$ with respect to $\realOrderRel$.
        \end{subproof}
        Thus $\reals$ has the least upper bound property with respect to $\realOrderRel$ by \cref{least_upper_bound_property}.
    \end{subproof}
    Show that for all $a,b\in\reals$ if $a\mathrel{\realOrderRel} b$ then
        there exists $z\in\reals$ such that $a\mathrel{\realOrderRel} z$ and $z\mathrel{\realOrderRel} b$.
    \begin{subproof}
        Fix $a$.
        Fix $b$.
        Assume $a\in\reals$ and $b\in\reals$.
        Assume $a\mathrel{\realOrderRel} b$.
        Thus $a < b$ by \cref{real_order_relation_elim}.
        Take $p\in\rationals$ such that $p\in\intervalopen{a}{b}$ by \cref{real_rational_open_interval}.
        Thus $p\in\reals$ and $a < p$ and $p < b$ by \cref{intervalopen}.
        Thus $a\mathrel{\realOrderRel} p$ and $p\mathrel{\realOrderRel} b$ by \cref{real_order_relation_intro}.
        Thus there exists $z\in\reals$ such that $a\mathrel{\realOrderRel} z$ and $z\mathrel{\realOrderRel} b$.
    \end{subproof}
    Thus $\reals$ is a linear continuum with respect to $\realOrderRel$ by \cref{linear_continuum}.
\end{proof}

% from §24 Item 3: real line connected, intervals and rays connected
\begin{corollary}\label{reals_connected}
    $\reals$ is connected with respect to $\realOrderTopology$ and
        for all $a,b\in\reals$ if $a\mathrel{\realOrderRel} b$ then
            $\intervalopenrel{\realOrderRel}{\reals}{a}{b}$ is connected with respect to $\subspaceTopology{\realOrderTopology}{\intervalopenrel{\realOrderRel}{\reals}{a}{b}}$ and
            $\intervalopenclosedrel{\realOrderRel}{\reals}{a}{b}$ is connected with respect to $\subspaceTopology{\realOrderTopology}{\intervalopenclosedrel{\realOrderRel}{\reals}{a}{b}}$ and
            $\intervalclosedopenrel{\realOrderRel}{\reals}{a}{b}$ is connected with respect to $\subspaceTopology{\realOrderTopology}{\intervalclosedopenrel{\realOrderRel}{\reals}{a}{b}}$ and
            $\intervalclosedrel{\realOrderRel}{\reals}{a}{b}$ is connected with respect to $\subspaceTopology{\realOrderTopology}{\intervalclosedrel{\realOrderRel}{\reals}{a}{b}}$ and
        for all $a\in\reals$ we have
            $\openrayright{\realOrderRel}{\reals}{a}$ is connected with respect to $\subspaceTopology{\realOrderTopology}{\openrayright{\realOrderRel}{\reals}{a}}$ and
            $\openrayleft{\realOrderRel}{\reals}{a}$ is connected with respect to $\subspaceTopology{\realOrderTopology}{\openrayleft{\realOrderRel}{\reals}{a}}$ and
            $\closedrayright{\realOrderRel}{\reals}{a}$ is connected with respect to $\subspaceTopology{\realOrderTopology}{\closedrayright{\realOrderRel}{\reals}{a}}$ and
            $\closedrayleft{\realOrderRel}{\reals}{a}$ is connected with respect to $\subspaceTopology{\realOrderTopology}{\closedrayleft{\realOrderRel}{\reals}{a}}$.
\end{corollary}
\begin{proof}
    $\realOrderRel$ is a simple order relation on $\reals$ and $\reals$ has more than one element by
        \cref{reals_linear_continuum,linear_continuum}.
    $\realOrderTopology = \basisTopology{\orderBasis{\realOrderRel}{\reals}}{\reals}$ by \cref{real_order_topology}.
    Thus $\realOrderTopology$ is the order topology on $\reals$ with respect to $\realOrderRel$ by \cref{order_topology}.
    Thus $\reals$ is connected with respect to $\realOrderTopology$ and
        for all $a,b\in\reals$ if $a\mathrel{\realOrderRel} b$ then
            $\intervalopenrel{\realOrderRel}{\reals}{a}{b}$ is connected with respect to $\subspaceTopology{\realOrderTopology}{\intervalopenrel{\realOrderRel}{\reals}{a}{b}}$ and
            $\intervalopenclosedrel{\realOrderRel}{\reals}{a}{b}$ is connected with respect to $\subspaceTopology{\realOrderTopology}{\intervalopenclosedrel{\realOrderRel}{\reals}{a}{b}}$ and
            $\intervalclosedopenrel{\realOrderRel}{\reals}{a}{b}$ is connected with respect to $\subspaceTopology{\realOrderTopology}{\intervalclosedopenrel{\realOrderRel}{\reals}{a}{b}}$ and
            $\intervalclosedrel{\realOrderRel}{\reals}{a}{b}$ is connected with respect to $\subspaceTopology{\realOrderTopology}{\intervalclosedrel{\realOrderRel}{\reals}{a}{b}}$ and
        for all $a\in\reals$ we have
            $\openrayright{\realOrderRel}{\reals}{a}$ is connected with respect to $\subspaceTopology{\realOrderTopology}{\openrayright{\realOrderRel}{\reals}{a}}$ and
            $\openrayleft{\realOrderRel}{\reals}{a}$ is connected with respect to $\subspaceTopology{\realOrderTopology}{\openrayleft{\realOrderRel}{\reals}{a}}$ and
            $\closedrayright{\realOrderRel}{\reals}{a}$ is connected with respect to $\subspaceTopology{\realOrderTopology}{\closedrayright{\realOrderRel}{\reals}{a}}$ and
            $\closedrayleft{\realOrderRel}{\reals}{a}$ is connected with respect to $\subspaceTopology{\realOrderTopology}{\closedrayleft{\realOrderRel}{\reals}{a}}$
        by \cref{linear_continuum_connected,reals_linear_continuum}.
\end{proof}

% from §24 Item 4: intermediate value theorem
\begin{theorem}\label{intermediate_value_theorem}
    Suppose $f$ is continuous from $X$ to $Y$ with respect to $T$ and $T_1$.
    Assume $X$ is connected with respect to $T$.
    Assume $R$ is a simple order relation on $Y$.
    Assume $Y$ has more than one element.
    Assume $T_1$ is the order topology on $Y$ with respect to $R$.
    Suppose $a\in X$.
    Suppose $b\in X$.
    Suppose $r\in Y$.
    Suppose $r\in\betweenrel{R}{Y}{f(a)}{f(b)}$.
    Then there exists $c\in X$ such that $f(c) = r$.
\end{theorem}
\begin{proof}
    $f\in\funs{X}{Y}$ by \cref{continuous}.
    $f$ is a function by \cref{funs_elim}.
    $T$ is a topology on $X$ by \cref{continuous}.
    $T_1$ is a topology on $Y$ by \cref{continuous}.
    Let $Z = \img{f}{X}$.
    Let $T_Z = \subspaceTopology{T_1}{Z}$.
    $Z$ is connected with respect to $T_Z$ by \cref{connected_image_continuous}.
    $\ran{f}\subseteq Y$ by \cref{funs_ran}.
    $Z\subseteq \ran{f}$ by \cref{img_subseteq_ran}.
    Thus $Z\subseteq Y$ by \cref{subseteq_transitive}.
    Thus $Z$ is a subspace of $Y$ with respect to $T_1$ by \cref{subspace_of}.
    Thus $T_Z$ is a topology on $Z$ by \cref{subspace_topology_is_topology,subspace_of}.
    Suppose not.
    Then there exists no $c\in X$ such that $f(c) = r$.
    Show that $r\notin Z$.
    \begin{subproof}
        Suppose not.
        Then $r\in Z$.
        Take $c\in X$ such that $c\mathrel{f} r$ by \cref{img_iff}.
        Thus $f(c) = r$ by \cref{function_apply_iff}.
        Contradiction.
    \end{subproof}
    Show that $f(a)\in Z$.
    \begin{subproof}
        $a\in X$.
        Thus $a\in\dom{f}$ by \cref{funs}.
        Thus $(a,f(a))\in f$ by \cref{function_apply_elim}.
        Thus $f(a)\in Z$ by \cref{img_elem_intro}.
    \end{subproof}
    Show that $f(b)\in Z$.
    \begin{subproof}
        $b\in X$.
        Thus $b\in\dom{f}$ by \cref{funs}.
        Thus $(b,f(b))\in f$ by \cref{function_apply_elim}.
        Thus $f(b)\in Z$ by \cref{img_elem_intro}.
    \end{subproof}
    Show that $r\neq f(a)$.
    \begin{subproof}
        Suppose not.
        Then $r = f(a)$.
        Thus $r\in Z$.
        Contradiction.
    \end{subproof}
    Show that $r\neq f(b)$.
    \begin{subproof}
        Suppose not.
        Then $r = f(b)$.
        Thus $r\in Z$.
        Contradiction.
    \end{subproof}
    $f(a)\in Y$ by \cref{funs_type_apply}.
    $f(b)\in Y$ by \cref{funs_type_apply}.
    Let $A = Z\inter\openrayleft{R}{Y}{r}$.
    Let $B = Z\inter\openrayright{R}{Y}{r}$.
    $\openrayleft{R}{Y}{r}$ is open in $Y$ with respect to $T_1$ by \cref{openrayleft_open_in_order_topology}.
    $\openrayright{R}{Y}{r}$ is open in $Y$ with respect to $T_1$ by \cref{openrayright_open_in_order_topology}.
    Thus $\openrayleft{R}{Y}{r}\in T_1$ by \cref{open_in}.
    Thus $\openrayright{R}{Y}{r}\in T_1$ by \cref{open_in}.
    Thus $A\in T_Z$ by \cref{subspace_topology}.
    Thus $B\in T_Z$ by \cref{subspace_topology}.
    Thus $A$ is open in $Z$ with respect to $T_Z$ by \cref{open_in}.
    Thus $B$ is open in $Z$ with respect to $T_Z$ by \cref{open_in}.
    Show that $A\neq\emptyset$ and $B\neq\emptyset$.
    \begin{subproof}
        $r\in \intervalclosedrel{R}{Y}{f(a)}{f(b)} \union \intervalclosedrel{R}{Y}{f(b)}{f(a)}$ by \cref{between_rel}.
        Thus $r\in \intervalclosedrel{R}{Y}{f(a)}{f(b)}$ or $r\in \intervalclosedrel{R}{Y}{f(b)}{f(a)}$ by \cref{union_iff}.
        \begin{byCase}
    \caseOf{$r\in \intervalclosedrel{R}{Y}{f(a)}{f(b)}$.}
        Thus $f(a)\mathrel{R} r$ or $r = f(a)$ by \cref{intervalclosed_rel}.
        Thus $r\mathrel{R} f(b)$ or $r = f(b)$ by \cref{intervalclosed_rel}.
        Show that $f(a)\mathrel{R} r$.
        \begin{subproof}
            \begin{byCase}
            \caseOf{$r = f(a)$.} Contradiction.
            \caseOf{$r\neq f(a)$.} Thus $f(a)\mathrel{R} r$.
            \end{byCase}
        \end{subproof}
        Show that $r\mathrel{R} f(b)$.
        \begin{subproof}
            \begin{byCase}
            \caseOf{$r = f(b)$.} Contradiction.
            \caseOf{$r\neq f(b)$.} Thus $r\mathrel{R} f(b)$.
            \end{byCase}
        \end{subproof}
        Thus $f(a)\in \openrayleft{R}{Y}{r}$ by \cref{openray_left}.
        Thus $f(b)\in \openrayright{R}{Y}{r}$ by \cref{openray_right}.
        Thus $f(a)\in A$ by \cref{inter_intro}.
        Thus $f(b)\in B$ by \cref{inter_intro}.
        Thus $A\neq\emptyset$ by \cref{emptyset}.
        Thus $B\neq\emptyset$ by \cref{emptyset}.
    \caseOf{$r\in \intervalclosedrel{R}{Y}{f(b)}{f(a)}$.}
        Thus $f(b)\mathrel{R} r$ or $r = f(b)$ by \cref{intervalclosed_rel}.
        Thus $r\mathrel{R} f(a)$ or $r = f(a)$ by \cref{intervalclosed_rel}.
        Show that $f(b)\mathrel{R} r$.
        \begin{subproof}
            \begin{byCase}
            \caseOf{$r = f(b)$.} Contradiction.
            \caseOf{$r\neq f(b)$.} Thus $f(b)\mathrel{R} r$.
            \end{byCase}
        \end{subproof}
        Show that $r\mathrel{R} f(a)$.
        \begin{subproof}
            \begin{byCase}
            \caseOf{$r = f(a)$.} Contradiction.
            \caseOf{$r\neq f(a)$.} Thus $r\mathrel{R} f(a)$.
            \end{byCase}
        \end{subproof}
        Thus $f(b)\in \openrayleft{R}{Y}{r}$ by \cref{openray_left}.
        Thus $f(a)\in \openrayright{R}{Y}{r}$ by \cref{openray_right}.
        Thus $f(b)\in A$ by \cref{inter_intro}.
        Thus $f(a)\in B$ by \cref{inter_intro}.
        Thus $A\neq\emptyset$ by \cref{emptyset}.
        Thus $B\neq\emptyset$ by \cref{emptyset}.
        \end{byCase}
    \end{subproof}
    Show that $A$ is disjoint from $B$.
    \begin{subproof}
        Suppose not.
        Then there exists $y$ such that $y\in A$ and $y\in B$.
        Thus $y\in \openrayleft{R}{Y}{r}$ and $y\in \openrayright{R}{Y}{r}$ by \cref{inter}.
        Thus $y\mathrel{R} r$ and $r\mathrel{R} y$ by \cref{openray_left,openray_right}.
        $R$ is asymmetric by \cref{simple_order_relation}.
        Contradiction by \cref{asymmetric}.
    \end{subproof}
    Show that $A\union B = Z$.
    \begin{subproof}
        Show that $A\union B\subseteq Z$.
        \begin{subproof}
            Show that for all $y$ such that $y\in A\union B$ we have $y\in Z$.
            \begin{subproof}
                Fix $y$.
                Assume $y\in A\union B$.
                Thus $y\in A$ or $y\in B$ by \cref{union_iff}.
                \begin{byCase}
                \caseOf{$y\in A$.}
                    Thus $y\in Z$ by \cref{inter}.
                \caseOf{$y\in B$.}
                    Thus $y\in Z$ by \cref{inter}.
                \end{byCase}
            \end{subproof}
            Thus $A\union B\subseteq Z$ by \cref{subseteq}.
        \end{subproof}
        Show that $Z\subseteq A\union B$.
        \begin{subproof}
            Show that for all $y$ such that $y\in Z$ we have $y\in A\union B$.
            \begin{subproof}
                Fix $y$.
                Assume $y\in Z$.
                Thus $y\in Y$ by \cref{subseteq}.
                Show that $y\neq r$.
                \begin{subproof}
                    Suppose not.
                    Then $y = r$.
                    Thus $r\in Z$.
                    Contradiction.
                \end{subproof}
                $R$ is connex on $Y$ by \cref{simple_order_relation}.
                Thus $y\mathrel{R} r$ or $r\mathrel{R} y$ by \cref{connex}.
                \begin{byCase}
                \caseOf{$y\mathrel{R} r$.}
                    Thus $y\in \openrayleft{R}{Y}{r}$ by \cref{openray_left}.
                    Thus $y\in A$ by \cref{inter_intro}.
                    Thus $y\in A\union B$ by \cref{union_iff}.
                \caseOf{$r\mathrel{R} y$.}
                    Thus $y\in \openrayright{R}{Y}{r}$ by \cref{openray_right}.
                    Thus $y\in B$ by \cref{inter_intro}.
                    Thus $y\in A\union B$ by \cref{union_iff}.
                \end{byCase}
            \end{subproof}
            Thus $Z\subseteq A\union B$ by \cref{subseteq}.
        \end{subproof}
        Thus $A\union B = Z$ by \cref{subseteq_antisymmetric}.
    \end{subproof}
    Thus $Z$ is separated by $A$ and $B$ with respect to $T_Z$ by \cref{separation}.
    Contradiction by \cref{connected}.
\end{proof}

% from §24 Item 5: path in a space
\begin{definition}\label{path_in_space}
    $f$ is a path in $X$ from $x$ to $y$ with respect to $T$ iff
        there exist $a,b\in\reals$ such that
            $a\mathrel{\realOrderRel} b$ and
            $f$ is continuous from $\intervalclosedrel{\realOrderRel}{\reals}{a}{b}$ to $X$ with respect to
                $\subspaceTopology{\realOrderTopology}{\intervalclosedrel{\realOrderRel}{\reals}{a}{b}}$ and $T$ and
            $f(a) = x$ and $f(b) = y$.
\end{definition}

% from §24 Item 6: path connected
\begin{definition}\label{path_connected}
    $X$ is path connected with respect to $T$ iff
        for all $x,y\in X$ there exists $f$ such that
            $f$ is a path in $X$ from $x$ to $y$ with respect to $T$.
\end{definition}

% from §24 Item 7: path-connected implies connected
\begin{theorem}\label{path_connected_implies_connected}
    Assume $T$ is a topology on $X$.
    Assume $X$ is path connected with respect to $T$.
    Then $X$ is connected with respect to $T$.
\end{theorem}
\begin{proof}
    Suppose not.
    Then there exists $A$ such that there exists $B$ such that
        $X$ is separated by $A$ and $B$ with respect to $T$.
    Take $A$ such that there exists $B$ such that
        $X$ is separated by $A$ and $B$ with respect to $T$.
    Take $B$ such that $X$ is separated by $A$ and $B$ with respect to $T$.
    $A\neq\emptyset$ by \cref{separation}.
    $B\neq\emptyset$ by \cref{separation}.
    $A$ is open in $X$ with respect to $T$ by \cref{separation}.
    $B$ is open in $X$ with respect to $T$ by \cref{separation}.
    Thus $A\in T$ by \cref{open_in}.
    Thus $B\in T$ by \cref{open_in}.
    $T\subseteq\pow{X}$ by \cref{topology_on}.
    Thus $A\subseteq X$ by \cref{powerset_elim,subseteq}.
    Thus $B\subseteq X$ by \cref{powerset_elim,subseteq}.
    Take $x$ such that $x\in A$ by \cref{nonempty_inhabited}.
    Take $y$ such that $y\in B$ by \cref{nonempty_inhabited}.
    Thus $x\in X$ and $y\in X$ by \cref{subseteq}.
    $A$ is disjoint from $B$ by \cref{separation}.
    Take $f$ such that $f$ is a path in $X$ from $x$ to $y$ with respect to $T$ by \cref{path_connected}.
    Take $a,b\in\reals$ such that
        $a\mathrel{\realOrderRel} b$ and
        $f$ is continuous from $\intervalclosedrel{\realOrderRel}{\reals}{a}{b}$ to $X$ with respect to
            $\subspaceTopology{\realOrderTopology}{\intervalclosedrel{\realOrderRel}{\reals}{a}{b}}$ and $T$ and
        $f(a) = x$ and $f(b) = y$
        by \cref{path_in_space}.
    Let $I = \intervalclosedrel{\realOrderRel}{\reals}{a}{b}$.
    Let $T_I = \subspaceTopology{\realOrderTopology}{I}$.
    $T_I$ is a topology on $I$ by \cref{continuous}.
    $I$ is connected with respect to $T_I$ by \cref{reals_connected}.
    Let $Z = \img{f}{I}$.
    Let $T_Z = \subspaceTopology{T}{Z}$.
    $Z$ is connected with respect to $T_Z$ by \cref{connected_image_continuous}.
    $f\in\funs{I}{X}$ by \cref{continuous}.
    $f$ is a function by \cref{funs_elim}.
    $\ran{f}\subseteq X$ by \cref{funs_ran}.
    $Z\subseteq \ran{f}$ by \cref{img_subseteq_ran}.
    Thus $Z\subseteq X$ by \cref{subseteq_transitive}.
    Thus $Z$ is a subspace of $X$ with respect to $T$ by \cref{subspace_of}.
    Thus $Z\subseteq A$ or $Z\subseteq B$ by \cref{connected_subspace_in_separation}.
    Show that $a\in I$.
    \begin{subproof}
        $a\in\reals$.
        $a = a$.
        Thus $a\mathrel{\realOrderRel} a$ or $a = a$.
        Thus $a\mathrel{\realOrderRel} b$ or $a = b$.
        Thus $a\in I$ by \cref{intervalclosed_rel}.
    \end{subproof}
    Show that $b\in I$.
    \begin{subproof}
        $b\in\reals$.
        $b = b$.
        Thus $b\mathrel{\realOrderRel} b$ or $b = b$.
        Thus $a\mathrel{\realOrderRel} b$ or $b = a$.
        Thus $b\in I$ by \cref{intervalclosed_rel}.
    \end{subproof}
    Show that $x\in Z$.
    \begin{subproof}
        $a\in\dom{f}$ by \cref{funs}.
        Thus $(a,f(a))\in f$ by \cref{function_apply_elim}.
        Thus $f(a)\in Z$ by \cref{img_elem_intro}.
        Thus $x\in Z$.
    \end{subproof}
    Show that $y\in Z$.
    \begin{subproof}
        $b\in\dom{f}$ by \cref{funs}.
        Thus $(b,f(b))\in f$ by \cref{function_apply_elim}.
        Thus $f(b)\in Z$ by \cref{img_elem_intro}.
        Thus $y\in Z$.
    \end{subproof}
    Show that there exists $z$ such that $z\in A$ and $z\in B$.
    \begin{subproof}
        \begin{byCase}
        \caseOf{$Z\subseteq A$.}
            Thus $y\in A$ by \cref{subseteq}.
            Thus there exists $z$ such that $z\in A$ and $z\in B$.
        \caseOf{$Z\subseteq B$.}
            Thus $x\in B$ by \cref{subseteq}.
            Thus there exists $z$ such that $z\in A$ and $z\in B$.
        \end{byCase}
    \end{subproof}
    Contradiction by \cref{disjoint}.
\end{proof}

% from §25 Item 1: component relation
\begin{definition}\label{component_relation}
    $\componentRelation{T}{X} = \{w\in X\times X \mid
        \text{there exists $A$ such that
            $A$ is a subspace of $X$ with respect to $T$ and
            $\fst{w}\in A$ and $\snd{w}\in A$ and
            $A$ is connected with respect to $\subspaceTopology{T}{A}$}\}$.
\end{definition}

% from §25 Item 2: component
\begin{definition}\label{component}
    Assume $T$ is a topology on $X$.
    $C$ is a component of $X$ with respect to $T$ iff
        there exists $x\in X$ such that
            $C = \equivalenceClass{\componentRelation{T}{X}}{x}$.
\end{definition}

% from §25 helper lemma 1: component relation is a binary relation
\begin{lemma}\label{component_relation_binary}
    Assume $T$ is a topology on $X$.
    Then $\componentRelation{T}{X}$ is a binary relation on $X$.
\end{lemma}
\begin{proof}
    Show that for all $w$ such that $w\in \componentRelation{T}{X}$ we have $w\in X\times X$.
    \begin{subproof}
        Fix $w$.
        Assume $w\in \componentRelation{T}{X}$.
        Thus $w\in X\times X$ and there exists $A$ such that
            $A$ is a subspace of $X$ with respect to $T$ and
            $\fst{w}\in A$ and $\snd{w}\in A$ and
            $A$ is connected with respect to $\subspaceTopology{T}{A}$
            by \cref{component_relation}.
        Thus $w\in X\times X$.
    \end{subproof}
    Thus $\componentRelation{T}{X}\subseteq X\times X$ by \cref{subseteq}.
    Thus $\componentRelation{T}{X}$ is a binary relation on $X$.
\end{proof}

% from §25 helper lemma 2: component relation is reflexive
\begin{lemma}\label{component_relation_reflexive}
    Assume $T$ is a topology on $X$.
    Then $\componentRelation{T}{X}$ is reflexive on $X$.
\end{lemma}
\begin{proof}
    Show that for all $x\in X$ we have $x\mathrel{\componentRelation{T}{X}}x$.
    \begin{subproof}
        Fix $x\in X$.
        Let $A = \{x\}$.
        Thus $A\subseteq X$ by \cref{singleton_subset_intro}.
        Thus $A$ is a subspace of $X$ with respect to $T$ by \cref{subspace_of}.
        Show that $A$ is connected with respect to $\subspaceTopology{T}{A}$.
        \begin{subproof}
            Suppose not.
            Then there exists $U$ such that there exists $V$ such that
                $A$ is separated by $U$ and $V$ with respect to $\subspaceTopology{T}{A}$.
            Take $U$ such that there exists $V$ such that
                $A$ is separated by $U$ and $V$ with respect to $\subspaceTopology{T}{A}$.
            Take $V$ such that $A$ is separated by $U$ and $V$ with respect to $\subspaceTopology{T}{A}$.
            $U\neq\emptyset$ and $V\neq\emptyset$ by \cref{separation}.
            Take $u$ such that $u\in U$ by \cref{nonempty_inhabited}.
            Take $v$ such that $v\in V$ by \cref{nonempty_inhabited}.
            $U$ is open in $A$ with respect to $\subspaceTopology{T}{A}$ by \cref{separation}.
            $V$ is open in $A$ with respect to $\subspaceTopology{T}{A}$ by \cref{separation}.
            $\subspaceTopology{T}{A}$ is a topology on $A$ by \cref{subspace_topology_is_topology,subspace_of}.
            Thus $U\in\subspaceTopology{T}{A}$ and $V\in\subspaceTopology{T}{A}$ by \cref{open_in}.
            Take $U_1\in T$ such that $U = A\inter U_1$ by \cref{subspace_topology}.
            Take $V_1\in T$ such that $V = A\inter V_1$ by \cref{subspace_topology}.
            Thus $U\subseteq A$ and $V\subseteq A$ by \cref{inter_lower_left,subseteq}.
            Thus $u\in A$ and $v\in A$ by \cref{subseteq}.
            Thus $u = x$ and $v = x$ by \cref{singleton_iff}.
            Thus there exists $z$ such that $z\in U$ and $z\in V$.
            Contradiction by \cref{disjoint,separation}.
        \end{subproof}
        $x\in A$ by \cref{singleton_intro}.
        Thus $\fst{(x,x)}\in A$ and $\snd{(x,x)}\in A$ by \cref{fst_eq,snd_eq}.
        Thus $(x,x)\in X\times X$ by \cref{times_tuple_intro}.
        Thus there exists $A$ such that
            $A$ is a subspace of $X$ with respect to $T$ and
            $\fst{(x,x)}\in A$ and $\snd{(x,x)}\in A$ and
            $A$ is connected with respect to $\subspaceTopology{T}{A}$.
        Thus $(x,x)\in \componentRelation{T}{X}$ by \cref{component_relation}.
        Thus $x\mathrel{\componentRelation{T}{X}}x$.
    \end{subproof}
    Thus $\componentRelation{T}{X}$ is reflexive on $X$ by \cref{reflexive_on}.
\end{proof}

% from §25 helper lemma 3: component relation is symmetric
\begin{lemma}\label{component_relation_symmetric}
    Assume $T$ is a topology on $X$.
    Then $\componentRelation{T}{X}$ is symmetric.
\end{lemma}
\begin{proof}
    Show that for all $x,y$ such that $x\mathrel{\componentRelation{T}{X}}y$ we have
        $y\mathrel{\componentRelation{T}{X}}x$.
    \begin{subproof}
        Fix $x$.
        Fix $y$.
        Assume $x\mathrel{\componentRelation{T}{X}}y$.
        Thus $(x,y)\in X\times X$ and there exists $A$ such that
            $A$ is a subspace of $X$ with respect to $T$ and
            $\fst{(x,y)}\in A$ and $\snd{(x,y)}\in A$ and
            $A$ is connected with respect to $\subspaceTopology{T}{A}$
            by \cref{component_relation}.
        Take $A$ such that
            $A$ is a subspace of $X$ with respect to $T$ and
            $\fst{(x,y)}\in A$ and $\snd{(x,y)}\in A$ and
            $A$ is connected with respect to $\subspaceTopology{T}{A}$.
        Thus $x\in X$ and $y\in X$ by \cref{times_tuple_elim}.
        Thus $x\in A$ and $y\in A$ by \cref{fst_eq,snd_eq}.
        Thus $(y,x)\in X\times X$ by \cref{times_tuple_intro}.
        Thus $\fst{(y,x)}\in A$ and $\snd{(y,x)}\in A$ by \cref{fst_eq,snd_eq}.
        Thus there exists $A$ such that
            $A$ is a subspace of $X$ with respect to $T$ and
            $\fst{(y,x)}\in A$ and $\snd{(y,x)}\in A$ and
            $A$ is connected with respect to $\subspaceTopology{T}{A}$.
        Thus $(y,x)\in \componentRelation{T}{X}$ by \cref{component_relation}.
        Thus $y\mathrel{\componentRelation{T}{X}}x$.
    \end{subproof}
    Thus $\componentRelation{T}{X}$ is symmetric by \cref{symmetric}.
\end{proof}

% from §25 helper lemma 4: component relation is transitive
\begin{lemma}\label{component_relation_transitive}
    Assume $T$ is a topology on $X$.
    Then $\componentRelation{T}{X}$ is transitive.
\end{lemma}
\begin{proof}
    Show that for all $x,y,z$ such that
        $x\mathrel{\componentRelation{T}{X}}y$ and $y\mathrel{\componentRelation{T}{X}}z$
        we have $x\mathrel{\componentRelation{T}{X}}z$.
    \begin{subproof}
        Fix $x$.
        Fix $y$.
        Fix $z$.
        Assume $x\mathrel{\componentRelation{T}{X}}y$ and $y\mathrel{\componentRelation{T}{X}}z$.
        Thus $(x,y)\in X\times X$ and there exists $A$ such that
            $A$ is a subspace of $X$ with respect to $T$ and
            $\fst{(x,y)}\in A$ and $\snd{(x,y)}\in A$ and
            $A$ is connected with respect to $\subspaceTopology{T}{A}$
            by \cref{component_relation}.
        Take $A$ such that
            $A$ is a subspace of $X$ with respect to $T$ and
            $\fst{(x,y)}\in A$ and $\snd{(x,y)}\in A$ and
            $A$ is connected with respect to $\subspaceTopology{T}{A}$.
        Thus $x\in X$ and $y\in X$ by \cref{times_tuple_elim}.
        Thus $x\in A$ and $y\in A$ by \cref{fst_eq,snd_eq}.
        Thus $(y,z)\in X\times X$ and there exists $B$ such that
            $B$ is a subspace of $X$ with respect to $T$ and
            $\fst{(y,z)}\in B$ and $\snd{(y,z)}\in B$ and
            $B$ is connected with respect to $\subspaceTopology{T}{B}$
            by \cref{component_relation}.
        Take $B$ such that
            $B$ is a subspace of $X$ with respect to $T$ and
            $\fst{(y,z)}\in B$ and $\snd{(y,z)}\in B$ and
            $B$ is connected with respect to $\subspaceTopology{T}{B}$.
        Thus $y\in X$ and $z\in X$ by \cref{times_tuple_elim}.
        Thus $y\in B$ and $z\in B$ by \cref{fst_eq,snd_eq}.
        Let $M = \{A,B\}$.
        $A\in M$ by \cref{upair_intro_left}.
        Thus $M\neq\emptyset$ by \cref{emptyset}.
        Show that for all $C$ such that $C\in M$ we have $C\subseteq X$.
        \begin{subproof}
            Fix $C$.
            Assume $C\in M$.
            $C = A$ or $C = B$ by \cref{upair_iff}.
            \begin{byCase}
            \caseOf{$C = A$.} Thus $C\subseteq X$ by \cref{subspace_of}.
            \caseOf{$C = B$.} Thus $C\subseteq X$ by \cref{subspace_of}.
            \end{byCase}
        \end{subproof}
        Show that for all $C$ such that $C\in M$ we have
            $C$ is connected with respect to $\subspaceTopology{T}{C}$.
        \begin{subproof}
            Fix $C$.
            Assume $C\in M$.
            $C = A$ or $C = B$ by \cref{upair_iff}.
            \begin{byCase}
            \caseOf{$C = A$.} Thus $C$ is connected with respect to $\subspaceTopology{T}{C}$.
            \caseOf{$C = B$.} Thus $C$ is connected with respect to $\subspaceTopology{T}{C}$.
            \end{byCase}
        \end{subproof}
        Show that $\inters{M}\neq\emptyset$.
        \begin{subproof}
            Show that $y\in \inters{M}$.
            \begin{subproof}
                Show that for all $C$ such that $C\in M$ we have $y\in C$.
                \begin{subproof}
                    Fix $C$.
                    Assume $C\in M$.
                    $C = A$ or $C = B$ by \cref{upair_iff}.
                    \begin{byCase}
                    \caseOf{$C = A$.} Thus $y\in C$.
                    \caseOf{$C = B$.} Thus $y\in C$.
                    \end{byCase}
                \end{subproof}
                Thus $y\in \inters{M}$ by \cref{inters_iff_forall}.
            \end{subproof}
            Thus $\inters{M}\neq\emptyset$ by \cref{emptyset}.
        \end{subproof}
        Thus $\unions{M}$ is connected with respect to $\subspaceTopology{T}{\unions{M}}$ by \cref{union_connected_common_point}.
        Thus $\unions{M} = A\union B$ by \cref{union_as_unions}.
        Thus $A\union B$ is connected with respect to $\subspaceTopology{T}{A\union B}$.
        Show that $A\union B$ is a subspace of $X$ with respect to $T$.
        \begin{subproof}
            Show that $A\union B\subseteq X$.
            \begin{subproof}
                Show that for all $u$ such that $u\in A\union B$ we have $u\in X$.
                \begin{subproof}
                    Fix $u$.
                    Assume $u\in A\union B$.
                    Thus $u\in A$ or $u\in B$ by \cref{union_iff}.
                    \begin{byCase}
                    \caseOf{$u\in A$.} Thus $u\in X$ by \cref{subspace_of,subseteq}.
                    \caseOf{$u\in B$.} Thus $u\in X$ by \cref{subspace_of,subseteq}.
                    \end{byCase}
                \end{subproof}
                Thus $A\union B\subseteq X$ by \cref{subseteq}.
            \end{subproof}
            Thus $A\union B$ is a subspace of $X$ with respect to $T$ by \cref{subspace_of}.
        \end{subproof}
        Thus $x\in A\union B$ by \cref{union_iff}.
        Thus $z\in A\union B$ by \cref{union_iff}.
        Thus $(x,z)\in X\times X$ by \cref{times_tuple_intro}.
        Thus $\fst{(x,z)}\in A\union B$ and $\snd{(x,z)}\in A\union B$ by \cref{fst_eq,snd_eq}.
        Let $C = A\union B$.
        Thus $\fst{(x,z)}\in C$ and $\snd{(x,z)}\in C$.
        Thus $C$ is a subspace of $X$ with respect to $T$.
        Thus $C$ is connected with respect to $\subspaceTopology{T}{C}$.
        Thus there exists $C$ such that
            $C$ is a subspace of $X$ with respect to $T$ and
            $\fst{(x,z)}\in C$ and $\snd{(x,z)}\in C$ and
            $C$ is connected with respect to $\subspaceTopology{T}{C}$.
        Thus $(x,z)\in \componentRelation{T}{X}$ by \cref{component_relation}.
        Thus $x\mathrel{\componentRelation{T}{X}}z$.
    \end{subproof}
    Thus $\componentRelation{T}{X}$ is transitive by \cref{transitive}.
\end{proof}

% from §25 helper lemma 5: component relation is an equivalence
\begin{lemma}\label{component_relation_equivalence}
    Assume $T$ is a topology on $X$.
    Then $\componentRelation{T}{X}$ is an equivalence on $X$.
\end{lemma}
\begin{proof}
    $\componentRelation{T}{X}$ is a binary relation on $X$ by \cref{component_relation_binary}.
    $\componentRelation{T}{X}$ is reflexive on $X$ by \cref{component_relation_reflexive}.
    $\componentRelation{T}{X}$ is transitive by \cref{component_relation_transitive}.
    $\componentRelation{T}{X}$ is symmetric by \cref{component_relation_symmetric}.
    Thus $\componentRelation{T}{X}$ is an equivalence on $X$.
\end{proof}

% from §25 Item 3: components are connected
\begin{theorem}\label{component_connected}
    Assume $T$ is a topology on $X$.
    Suppose $C$ is a component of $X$ with respect to $T$.
    Then $C$ is connected with respect to $\subspaceTopology{T}{C}$.
\end{theorem}
\begin{proof}
    Take $x\in X$ such that $C = \equivalenceClass{\componentRelation{T}{X}}{x}$ by \cref{component}.
    Let $M = \{A\in\pow{X} \mid \text{$x\in A$ and $A$ is connected with respect to $\subspaceTopology{T}{A}$}\}$.
    Show that $M\neq\emptyset$.
    \begin{subproof}
        $\componentRelation{T}{X}$ is reflexive on $X$ by \cref{component_relation_reflexive}.
        Thus $x\mathrel{\componentRelation{T}{X}}x$ by \cref{reflexive_on}.
        Thus $(x,x)\in X\times X$ and there exists $A$ such that
            $A$ is a subspace of $X$ with respect to $T$ and
            $\fst{(x,x)}\in A$ and $\snd{(x,x)}\in A$ and
            $A$ is connected with respect to $\subspaceTopology{T}{A}$
            by \cref{component_relation}.
        Take $A$ such that
            $A$ is a subspace of $X$ with respect to $T$ and
            $\fst{(x,x)}\in A$ and $\snd{(x,x)}\in A$ and
            $A$ is connected with respect to $\subspaceTopology{T}{A}$.
        Thus $x\in A$ by \cref{fst_eq,snd_eq}.
        Thus $A\subseteq X$ by \cref{subspace_of}.
        Thus $A\in\pow{X}$ by \cref{powerset_intro}.
        Thus $A\in M$.
        Thus $M\neq\emptyset$ by \cref{emptyset}.
    \end{subproof}
    Show that for all $A$ such that $A\in M$ we have $A\subseteq X$.
    \begin{subproof}
        Fix $A$.
        Assume $A\in M$.
        Thus $A\in\pow{X}$.
        Thus $A\subseteq X$ by \cref{powerset_elim}.
    \end{subproof}
    Show that for all $A$ such that $A\in M$ we have
        $A$ is connected with respect to $\subspaceTopology{T}{A}$.
    \begin{subproof}
        Fix $A$.
        Assume $A\in M$.
        Thus $A$ is connected with respect to $\subspaceTopology{T}{A}$.
    \end{subproof}
    Show that $\inters{M}\neq\emptyset$.
    \begin{subproof}
        Show that $x\in \inters{M}$.
        \begin{subproof}
            Show that there exists $A$ such that $A\in M$.
            \begin{subproof}
                Thus there exists $A$ such that $A\in M$ by \cref{nonempty_inhabited}.
            \end{subproof}
            Show that for all $A$ such that $A\in M$ we have $x\in A$.
            \begin{subproof}
                Fix $A$.
                Assume $A\in M$.
                Thus $x\in A$.
            \end{subproof}
            Thus $x\in \inters{M}$ by \cref{inters_iff_forall}.
        \end{subproof}
        Thus $\inters{M}\neq\emptyset$ by \cref{emptyset}.
    \end{subproof}
    Let $U = \unions{M}$.
    $U$ is connected with respect to $\subspaceTopology{T}{U}$ by \cref{union_connected_common_point}.
    Show that $U = C$.
    \begin{subproof}
        Show that $U\subseteq C$.
        \begin{subproof}
            Show that for all $y$ such that $y\in U$ we have $y\in C$.
            \begin{subproof}
                Fix $y$.
                Assume $y\in U$.
                Take $A\in M$ such that $y\in A$ by \cref{unions_iff}.
                Thus $A\in\pow{X}$.
                Thus $x\in A$.
                Thus $A$ is connected with respect to $\subspaceTopology{T}{A}$.
                Thus $A\subseteq X$ by \cref{powerset_elim}.
                Thus $A$ is a subspace of $X$ with respect to $T$ by \cref{subspace_of}.
                Thus $y\in X$ by \cref{subseteq}.
                Thus $(x,y)\in X\times X$ by \cref{times_tuple_intro}.
                Thus $\fst{(x,y)}\in A$ and $\snd{(x,y)}\in A$ by \cref{fst_eq,snd_eq}.
                Thus there exists $A$ such that
                    $A$ is a subspace of $X$ with respect to $T$ and
                    $\fst{(x,y)}\in A$ and $\snd{(x,y)}\in A$ and
                    $A$ is connected with respect to $\subspaceTopology{T}{A}$.
                Thus $(x,y)\in \componentRelation{T}{X}$ by \cref{component_relation}.
                $\componentRelation{T}{X}$ is symmetric by \cref{component_relation_symmetric}.
                Thus $y\mathrel{\componentRelation{T}{X}}x$ by \cref{symmetric}.
                Thus $y\in \downward{\componentRelation{T}{X}}{x}$ by \cref{downward_closure_iff}.
                Thus $y\in C$.
            \end{subproof}
            Thus $U\subseteq C$ by \cref{subseteq}.
        \end{subproof}
        Show that $C\subseteq U$.
        \begin{subproof}
            Show that for all $y$ such that $y\in C$ we have $y\in U$.
            \begin{subproof}
                Fix $y$.
                Assume $y\in C$.
                Thus $y\in \downward{\componentRelation{T}{X}}{x}$.
                Thus $y\mathrel{\componentRelation{T}{X}}x$ by \cref{downward_closure_iff}.
                Thus $(y,x)\in X\times X$ and there exists $A$ such that
                    $A$ is a subspace of $X$ with respect to $T$ and
                    $\fst{(y,x)}\in A$ and $\snd{(y,x)}\in A$ and
                    $A$ is connected with respect to $\subspaceTopology{T}{A}$
                    by \cref{component_relation}.
                Take $A$ such that
                    $A$ is a subspace of $X$ with respect to $T$ and
                    $\fst{(y,x)}\in A$ and $\snd{(y,x)}\in A$ and
                    $A$ is connected with respect to $\subspaceTopology{T}{A}$.
                Thus $y\in A$ and $x\in A$ by \cref{fst_eq,snd_eq}.
                Thus $A\subseteq X$ by \cref{subspace_of}.
                Thus $A\in\pow{X}$ by \cref{powerset_intro}.
                Thus $A\in M$.
                Thus $y\in U$ by \cref{unions_iff}.
            \end{subproof}
            Thus $C\subseteq U$ by \cref{subseteq}.
        \end{subproof}
        Thus $U = C$ by \cref{subseteq_antisymmetric}.
    \end{subproof}
    Thus $C$ is connected with respect to $\subspaceTopology{T}{C}$.
\end{proof}

% from §25 Item 4: distinct components are disjoint
\begin{theorem}\label{components_disjoint}
    Assume $T$ is a topology on $X$.
    Suppose $C$ is a component of $X$ with respect to $T$.
    Suppose $D$ is a component of $X$ with respect to $T$.
    Suppose $C\neq D$.
    Then $C$ is disjoint from $D$.
\end{theorem}
\begin{proof}
    Suppose not.
    Then there exists $z$ such that $z\in C$ and $z\in D$.
    Take $x\in X$ such that $C = \equivalenceClass{\componentRelation{T}{X}}{x}$ by \cref{component}.
    Take $y\in X$ such that $D = \equivalenceClass{\componentRelation{T}{X}}{y}$ by \cref{component}.
    Thus $z\in \downward{\componentRelation{T}{X}}{x}$.
    Thus $z\mathrel{\componentRelation{T}{X}}x$ by \cref{downward_closure_iff}.
    Thus $z\in \downward{\componentRelation{T}{X}}{y}$.
    Thus $z\mathrel{\componentRelation{T}{X}}y$ by \cref{downward_closure_iff}.
    $\componentRelation{T}{X}$ is symmetric by \cref{component_relation_symmetric}.
    Thus $x\mathrel{\componentRelation{T}{X}}z$ by \cref{symmetric}.
    $\componentRelation{T}{X}$ is transitive by \cref{component_relation_transitive}.
    Thus $x\mathrel{\componentRelation{T}{X}}y$ by \cref{transitive}.
    $\componentRelation{T}{X}$ is an equivalence on $X$ by \cref{component_relation_equivalence}.
    Thus $\equivalenceClass{\componentRelation{T}{X}}{x} = \equivalenceClass{\componentRelation{T}{X}}{y}$ by \cref{equiv_iff_equivclasses_eq}.
    Thus $C = D$.
    Contradiction.
\end{proof}

% from §25 Item 5: components cover the space
\begin{theorem}\label{components_cover}
    Assume $T$ is a topology on $X$.
    Then for all $x\in X$ there exists $C$ such that
        $C$ is a component of $X$ with respect to $T$ and $x\in C$.
\end{theorem}
\begin{proof}
    Fix $x\in X$.
    Let $C = \equivalenceClass{\componentRelation{T}{X}}{x}$.
    $\componentRelation{T}{X}$ is an equivalence on $X$ by \cref{component_relation_equivalence}.
    Thus $x\in C$ by \cref{equivclasses_inhabited_equivalenceon}.
    Thus $C$ is a component of $X$ with respect to $T$ by \cref{component}.
    Thus there exists $C$ such that $C$ is a component of $X$ with respect to $T$ and $x\in C$.
\end{proof}

% from §25 Item 6: connected subspaces intersect only one component
\begin{theorem}\label{connected_subspace_single_component}
    Assume $T$ is a topology on $X$.
    Suppose $A$ is a subspace of $X$ with respect to $T$.
    Suppose $A$ is connected with respect to $\subspaceTopology{T}{A}$.
    Suppose $A\neq\emptyset$.
    Then for all $C,D$ such that
        $C$ is a component of $X$ with respect to $T$ and
        $D$ is a component of $X$ with respect to $T$ and
        $A$ intersects $C$ and $A$ intersects $D$
        we have $C = D$.
\end{theorem}
\begin{proof}
    Fix $C$.
    Fix $D$.
    Assume $C$ is a component of $X$ with respect to $T$.
    Assume $D$ is a component of $X$ with respect to $T$.
    Assume $A$ intersects $C$ and $A$ intersects $D$.
    $A\inter C\neq\emptyset$ by \cref{intersects}.
    Take $x$ such that $x\in A\inter C$ by \cref{nonempty_inhabited}.
    Thus $x\in A$ and $x\in C$ by \cref{inter}.
    $A\inter D\neq\emptyset$ by \cref{intersects}.
    Take $y$ such that $y\in A\inter D$ by \cref{nonempty_inhabited}.
    Thus $y\in A$ and $y\in D$ by \cref{inter}.
    Thus $x\in A$ and $y\in A$ by \cref{inter}.
    Thus $x\in X$ and $y\in X$ by \cref{subspace_of,subseteq}.
    Thus $(x,y)\in X\times X$ by \cref{times_tuple_intro}.
    Thus $\fst{(x,y)}\in A$ and $\snd{(x,y)}\in A$ by \cref{fst_eq,snd_eq}.
    Thus there exists $A$ such that
        $A$ is a subspace of $X$ with respect to $T$ and
        $\fst{(x,y)}\in A$ and $\snd{(x,y)}\in A$ and
        $A$ is connected with respect to $\subspaceTopology{T}{A}$.
    Thus $(x,y)\in \componentRelation{T}{X}$ by \cref{component_relation}.
    Take $c\in X$ such that $C = \equivalenceClass{\componentRelation{T}{X}}{c}$ by \cref{component}.
    Take $d\in X$ such that $D = \equivalenceClass{\componentRelation{T}{X}}{d}$ by \cref{component}.
    $\componentRelation{T}{X}$ is an equivalence on $X$ by \cref{component_relation_equivalence}.
    Thus $x\in \downward{\componentRelation{T}{X}}{c}$.
    Thus $x\mathrel{\componentRelation{T}{X}}c$ by \cref{downward_closure_iff}.
    Thus $y\in \downward{\componentRelation{T}{X}}{d}$.
    Thus $y\mathrel{\componentRelation{T}{X}}d$ by \cref{downward_closure_iff}.
    $\componentRelation{T}{X}$ is symmetric by \cref{component_relation_symmetric}.
    Thus $c\mathrel{\componentRelation{T}{X}}x$ by \cref{symmetric}.
    $\componentRelation{T}{X}$ is transitive by \cref{component_relation_transitive}.
    Thus $c\mathrel{\componentRelation{T}{X}}y$ by \cref{transitive}.
    Thus $c\mathrel{\componentRelation{T}{X}}d$ by \cref{transitive,symmetric}.
    Thus $\equivalenceClass{\componentRelation{T}{X}}{c} = \equivalenceClass{\componentRelation{T}{X}}{d}$ by \cref{equiv_iff_equivclasses_eq}.
    Thus $C = D$.
\end{proof}

% from §25 Item 7: path relation
\begin{definition}\label{path_relation}
    $\pathRelation{T}{X} = \{w\in X\times X \mid
        \text{there exists $f$ such that
            $f$ is a path in $X$ from $\fst{w}$ to $\snd{w}$ with respect to $T$}\}$.
\end{definition}

% from §25 Item 8: path component
\begin{definition}\label{path_component}
    Assume $T$ is a topology on $X$.
    $C$ is a path component of $X$ with respect to $T$ iff
        there exists $x\in X$ such that
            $C = \equivalenceClass{\pathRelation{T}{X}}{x}$.
\end{definition}

% from §25 helper lemma 6: real line has smaller elements
\begin{lemma}\label{reals_has_smaller}
    For all $a\in\reals$ there exists $b\in\reals$ such that $b < a$.
\end{lemma}
\begin{proof}
    Fix $a\in\reals$.
    $0\in\reals$ by \cref{real_add_ident}.
    $1\in\reals$ by \cref{real_mul_ident}.
    $\neg{1}\in\reals$ by \cref{real_add_inverse}.
    $0 < 1$ by \cref{real_zero_lt_one}.
    Thus $\neg{1} < \neg{0}$ by \cref{real_lt_neg}.
    $\neg{0} = 0$ by \cref{real_neg_zero}.
    Thus $\neg{1} < 0$ by \cref{real_lt_subst_right}.
    Thus $\neg{1} + a < 0 + a$ by \cref{real_lt_add}.
    $\neg{1} + a = a + \neg{1}$ by \cref{real_add_comm}.
    Thus $a + \neg{1} < 0 + a$ by \cref{real_lt_subst_left}.
    $0 + a = a$ by \cref{real_add_ident}.
    Thus $a + \neg{1} < a$ by \cref{real_lt_subst_right}.
    Let $b = a + \neg{1}$.
    Thus $b\in\reals$ by \cref{real_add_closed}.
    Thus $b < a$.
    Thus there exists $b\in\reals$ such that $b < a$.
\end{proof}

% from §25 helper lemma 7: real line has larger elements
\begin{lemma}\label{reals_has_larger}
    For all $a\in\reals$ there exists $b\in\reals$ such that $a < b$.
\end{lemma}
\begin{proof}
    Fix $a\in\reals$.
    $0\in\reals$ by \cref{real_add_ident}.
    $1\in\reals$ by \cref{real_mul_ident}.
    $0 < 1$ by \cref{real_zero_lt_one}.
    Thus $0 + a < 1 + a$ by \cref{real_lt_add}.
    $0 + a = a$ by \cref{real_add_ident}.
    $1 + a = a + 1$ by \cref{real_add_comm}.
    Thus $a < a + 1$ by \cref{real_lt_subst_left,real_lt_subst_right}.
    Let $b = a + 1$.
    $1\in\reals$ by \cref{real_mul_ident}.
    Thus $b\in\reals$ by \cref{real_add_closed}.
    Thus $a < b$.
    Thus there exists $b\in\reals$ such that $a < b$.
\end{proof}

% from §25 helper lemma 8: order interval equals real open interval
\begin{lemma}\label{intervalopenrel_real_eq_intervalopen}
    Suppose $a\in\reals$ and $b\in\reals$.
    Then $\intervalopenrel{\realOrderRel}{\reals}{a}{b} = \intervalopen{a}{b}$.
\end{lemma}
\begin{proof}
    Show that $\intervalopenrel{\realOrderRel}{\reals}{a}{b} \subseteq \intervalopen{a}{b}$.
    \begin{subproof}
        Show that for all $x$ such that $x\in\intervalopenrel{\realOrderRel}{\reals}{a}{b}$ we have
            $x\in\intervalopen{a}{b}$.
        \begin{subproof}
            Fix $x$.
            Assume $x\in\intervalopenrel{\realOrderRel}{\reals}{a}{b}$.
            Thus $x\in\reals$ by \cref{intervalopen_rel}.
            Thus $a\mathrel{\realOrderRel}x$ and $x\mathrel{\realOrderRel}b$ by \cref{intervalopen_rel}.
            Thus $a < x$ and $x < b$ by \cref{real_order_relation_elim}.
            Thus $x\in\intervalopen{a}{b}$ by \cref{intervalopen}.
        \end{subproof}
        Thus $\intervalopenrel{\realOrderRel}{\reals}{a}{b} \subseteq \intervalopen{a}{b}$ by \cref{subseteq}.
    \end{subproof}
    Show that $\intervalopen{a}{b} \subseteq \intervalopenrel{\realOrderRel}{\reals}{a}{b}$.
    \begin{subproof}
        Show that for all $x$ such that $x\in\intervalopen{a}{b}$ we have
            $x\in\intervalopenrel{\realOrderRel}{\reals}{a}{b}$.
        \begin{subproof}
            Fix $x$.
            Assume $x\in\intervalopen{a}{b}$.
            Thus $x\in\reals$ and $a < x$ and $x < b$ by \cref{intervalopen}.
            Thus $a\mathrel{\realOrderRel}x$ by \cref{real_order_relation_intro}.
            Thus $x\mathrel{\realOrderRel}b$ by \cref{real_order_relation_intro}.
            Thus $x\in\intervalopenrel{\realOrderRel}{\reals}{a}{b}$ by \cref{intervalopen_rel}.
        \end{subproof}
        Thus $\intervalopen{a}{b} \subseteq \intervalopenrel{\realOrderRel}{\reals}{a}{b}$ by \cref{subseteq}.
    \end{subproof}
    Thus $\intervalopenrel{\realOrderRel}{\reals}{a}{b} = \intervalopen{a}{b}$ by \cref{subseteq_antisymmetric}.
\end{proof}

% from §25 helper lemma 9: order basis on reals equals the standard basis
\begin{lemma}\label{order_basis_reals_standard}
    $\orderBasis{\realOrderRel}{\reals} = \standardBasisR$.
\end{lemma}
\begin{proof}
    Show that $\orderBasis{\realOrderRel}{\reals} \subseteq \standardBasisR$.
    \begin{subproof}
        Show that for all $B$ such that $B\in\orderBasis{\realOrderRel}{\reals}$ we have
            $B\in\standardBasisR$.
        \begin{subproof}
            Fix $B$.
            Assume $B\in\orderBasis{\realOrderRel}{\reals}$.
            Thus
                $(\exists a, b\in\reals. a\mathrel{\realOrderRel}b \land B = \intervalopenrel{\realOrderRel}{\reals}{a}{b}) \lor
                (\exists b\in\reals. \exists a_0\in\reals. (\forall x\in\reals. a_0\mathrel{\realOrderRel}x \lor a_0 = x) \land B = \intervalclosedopenrel{\realOrderRel}{\reals}{a_0}{b}) \lor
                (\exists a\in\reals. \exists b_0\in\reals. (\forall x\in\reals. x\mathrel{\realOrderRel}b_0 \lor x = b_0) \land B = \intervalopenclosedrel{\realOrderRel}{\reals}{a}{b_0})$
                by \cref{order_basis}.
            \begin{byCase}
            \caseOf{$\exists a, b\in\reals. a\mathrel{\realOrderRel}b \land B = \intervalopenrel{\realOrderRel}{\reals}{a}{b}$.}
                Take $a,b\in\reals$ such that $a\mathrel{\realOrderRel}b$ and
                    $B = \intervalopenrel{\realOrderRel}{\reals}{a}{b}$.
                Thus $a < b$ by \cref{real_order_relation_elim}.
                Thus $B = \intervalopen{a}{b}$ by \cref{intervalopenrel_real_eq_intervalopen}.
                $\intervalopen{a}{b}\subseteq\reals$ by \cref{intervalopen_subset_reals}.
                Thus $B\subseteq\reals$ by \cref{subseteq}.
                Thus $B\in\pow{\reals}$ by \cref{powerset_intro}.
                Thus $B\in\standardBasisR$ by \cref{standard_basis_reals}.
            \caseOf{$\exists b\in\reals. \exists a_0\in\reals. (\forall x\in\reals. a_0\mathrel{\realOrderRel}x \lor a_0 = x) \land B = \intervalclosedopenrel{\realOrderRel}{\reals}{a_0}{b}$.}
                Suppose not.
                Take $b\in\reals$ such that there exists $a_0\in\reals$ such that
                    $(\forall x\in\reals. a_0\mathrel{\realOrderRel}x \lor a_0 = x)$ and
                    $B = \intervalclosedopenrel{\realOrderRel}{\reals}{a_0}{b}$.
                Take $a_0\in\reals$ such that
                    $(\forall x\in\reals. a_0\mathrel{\realOrderRel}x \lor a_0 = x)$ and
                    $B = \intervalclosedopenrel{\realOrderRel}{\reals}{a_0}{b}$.
                Take $c\in\reals$ such that $c < a_0$ by \cref{reals_has_smaller}.
                Thus $a_0\mathrel{\realOrderRel}c$ or $a_0 = c$.
                \begin{byCase}
                \caseOf{$a_0\mathrel{\realOrderRel}c$.}
                    Thus $a_0 < c$ by \cref{real_order_relation_elim}.
                    Show that $a_0 \neq c$.
                    \begin{subproof}
                        Suppose not.
                        Then $a_0 = c$.
                        Thus $c < c$ by \cref{real_lt_subst_right}.
                        Contradiction by \cref{real_lt_irreflexive}.
                    \end{subproof}
                    Contradiction by \cref{real_lt_asym}.
                \caseOf{$a_0 = c$.}
                    Thus $c < c$ by \cref{real_lt_subst_right}.
                    Contradiction by \cref{real_lt_irreflexive}.
                \end{byCase}
            \caseOf{$\exists a\in\reals. \exists b_0\in\reals. (\forall x\in\reals. x\mathrel{\realOrderRel}b_0 \lor x = b_0) \land B = \intervalopenclosedrel{\realOrderRel}{\reals}{a}{b_0}$.}
                Suppose not.
                Take $a\in\reals$ such that there exists $b_0\in\reals$ such that
                    $(\forall x\in\reals. x\mathrel{\realOrderRel}b_0 \lor x = b_0)$ and
                    $B = \intervalopenclosedrel{\realOrderRel}{\reals}{a}{b_0}$.
                Take $b_0\in\reals$ such that
                    $(\forall x\in\reals. x\mathrel{\realOrderRel}b_0 \lor x = b_0)$ and
                    $B = \intervalopenclosedrel{\realOrderRel}{\reals}{a}{b_0}$.
                Take $c\in\reals$ such that $b_0 < c$ by \cref{reals_has_larger}.
                Thus $c\mathrel{\realOrderRel}b_0$ or $c = b_0$.
                \begin{byCase}
                \caseOf{$c\mathrel{\realOrderRel}b_0$.}
                    Thus $c < b_0$ by \cref{real_order_relation_elim}.
                    Show that $c \neq b_0$.
                    \begin{subproof}
                        Suppose not.
                        Then $c = b_0$.
                        Thus $c < c$ by \cref{real_lt_subst_right}.
                        Contradiction by \cref{real_lt_irreflexive}.
                    \end{subproof}
                    Contradiction by \cref{real_lt_asym}.
                \caseOf{$c = b_0$.}
                    Thus $c < c$ by \cref{real_lt_subst_right}.
                    Contradiction by \cref{real_lt_irreflexive}.
                \end{byCase}
            \end{byCase}
        \end{subproof}
        Thus $\orderBasis{\realOrderRel}{\reals} \subseteq \standardBasisR$ by \cref{subseteq}.
    \end{subproof}
    Show that $\standardBasisR \subseteq \orderBasis{\realOrderRel}{\reals}$.
    \begin{subproof}
        Show that for all $B$ such that $B\in\standardBasisR$ we have
            $B\in\orderBasis{\realOrderRel}{\reals}$.
        \begin{subproof}
            Fix $B$.
            Assume $B\in\standardBasisR$.
            Take $a,b\in\reals$ such that $a < b$ and $B = \intervalopen{a}{b}$ by \cref{standard_basis_reals}.
            Thus $a\mathrel{\realOrderRel}b$ by \cref{real_order_relation_intro}.
            Thus $B = \intervalopenrel{\realOrderRel}{\reals}{a}{b}$ by \cref{intervalopenrel_real_eq_intervalopen}.
            $\intervalopen{a}{b}\subseteq\reals$ by \cref{intervalopen_subset_reals}.
            Thus $B\subseteq\reals$ by \cref{subseteq}.
            Thus $B\in\pow{\reals}$ by \cref{powerset_intro}.
            Thus $\exists a,b\in\reals. a\mathrel{\realOrderRel}b \land B = \intervalopenrel{\realOrderRel}{\reals}{a}{b}$.
            Thus $B\in\orderBasis{\realOrderRel}{\reals}$ by \cref{order_basis}.
        \end{subproof}
        Thus $\standardBasisR \subseteq \orderBasis{\realOrderRel}{\reals}$ by \cref{subseteq}.
    \end{subproof}
    Thus $\orderBasis{\realOrderRel}{\reals} = \standardBasisR$ by \cref{subseteq_antisymmetric}.
\end{proof}

% from §25 helper lemma 10: real order topology equals standard topology
\begin{lemma}\label{real_order_topology_eq_standard}
    $\realOrderTopology = \standardTopologyR$.
\end{lemma}
\begin{proof}
    $\realOrderTopology = \basisTopology{\orderBasis{\realOrderRel}{\reals}}{\reals}$ by \cref{real_order_topology}.
    $\standardTopologyR = \basisTopology{\standardBasisR}{\reals}$ by \cref{standard_topology_reals}.
    $\orderBasis{\realOrderRel}{\reals} = \standardBasisR$ by \cref{order_basis_reals_standard}.
    Thus $\realOrderTopology = \standardTopologyR$.
\end{proof}

% from §25 helper lemma 11: closed right ray is complement of open left ray
\begin{lemma}\label{closedrayright_complement_openrayleft}
    Assume $R$ is a simple order relation on $X$.
    Suppose $a\in X$.
    Then $\closedrayright{R}{X}{a} = X\setminus \openrayleft{R}{X}{a}$.
\end{lemma}
\begin{proof}
    Show that $\closedrayright{R}{X}{a} \subseteq X\setminus \openrayleft{R}{X}{a}$.
    \begin{subproof}
        Show that for all $x$ such that $x\in\closedrayright{R}{X}{a}$ we have
            $x\in X\setminus \openrayleft{R}{X}{a}$.
        \begin{subproof}
            Fix $x$.
            Assume $x\in\closedrayright{R}{X}{a}$.
            Thus $x\in X$ and $a\mathrel{R}x$ or $x = a$ by \cref{closedray_right}.
            $R$ is asymmetric by \cref{simple_order_relation}.
            Show that $x\notin \openrayleft{R}{X}{a}$.
            \begin{subproof}
                Suppose not.
                Then $x\in\openrayleft{R}{X}{a}$.
                Thus $x\mathrel{R}a$ by \cref{openray_left}.
                \begin{byCase}
                \caseOf{$a\mathrel{R}x$.}
                    Contradiction by \cref{asymmetric}.
                \caseOf{$x = a$.}
                    Thus $a\mathrel{R}a$ by \cref{real_lt_subst_left}.
                    Contradiction by \cref{asymmetric}.
                \end{byCase}
            \end{subproof}
            Thus $x\in X\setminus \openrayleft{R}{X}{a}$ by \cref{setminus_intro}.
        \end{subproof}
        Thus $\closedrayright{R}{X}{a} \subseteq X\setminus \openrayleft{R}{X}{a}$ by \cref{subseteq}.
    \end{subproof}
    Show that $X\setminus \openrayleft{R}{X}{a} \subseteq \closedrayright{R}{X}{a}$.
    \begin{subproof}
        Show that for all $x$ such that $x\in X\setminus \openrayleft{R}{X}{a}$ we have
            $x\in\closedrayright{R}{X}{a}$.
        \begin{subproof}
            Fix $x$.
            Assume $x\in X\setminus \openrayleft{R}{X}{a}$.
            Thus $x\in X$ and $x\notin\openrayleft{R}{X}{a}$ by \cref{setminus_elim_left,setminus_elim_right}.
            Show that it is wrong that $x\mathrel{R}a$.
            \begin{subproof}
                Suppose not.
                Then $x\mathrel{R}a$.
                Thus $x\in\openrayleft{R}{X}{a}$ by \cref{openray_left}.
                Contradiction by \cref{setminus_elim_right}.
            \end{subproof}
            \begin{byCase}
            \caseOf{$x = a$.}
                Thus $x\in\closedrayright{R}{X}{a}$ by \cref{closedray_right}.
            \caseOf{$x \neq a$.}
                $R$ is connex on $X$ by \cref{simple_order_relation}.
                Thus $a\mathrel{R}x$ or $x\mathrel{R}a$ by \cref{connex}.
                Thus $a\mathrel{R}x$.
                Thus $x\in\closedrayright{R}{X}{a}$ by \cref{closedray_right}.
            \end{byCase}
        \end{subproof}
        Thus $X\setminus \openrayleft{R}{X}{a} \subseteq \closedrayright{R}{X}{a}$ by \cref{subseteq}.
    \end{subproof}
    Thus $\closedrayright{R}{X}{a} = X\setminus \openrayleft{R}{X}{a}$ by \cref{subseteq_antisymmetric}.
\end{proof}

% from §25 helper lemma 12: closed left ray is complement of open right ray
\begin{lemma}\label{closedrayleft_complement_openrayright}
    Assume $R$ is a simple order relation on $X$.
    Suppose $a\in X$.
    Then $\closedrayleft{R}{X}{a} = X\setminus \openrayright{R}{X}{a}$.
\end{lemma}
\begin{proof}
    Show that $\closedrayleft{R}{X}{a} \subseteq X\setminus \openrayright{R}{X}{a}$.
    \begin{subproof}
        Show that for all $x$ such that $x\in\closedrayleft{R}{X}{a}$ we have
            $x\in X\setminus \openrayright{R}{X}{a}$.
        \begin{subproof}
            Fix $x$.
            Assume $x\in\closedrayleft{R}{X}{a}$.
            Thus $x\in X$ and $x\mathrel{R}a$ or $x = a$ by \cref{closedray_left}.
            $R$ is asymmetric by \cref{simple_order_relation}.
            Show that $x\notin \openrayright{R}{X}{a}$.
            \begin{subproof}
                Suppose not.
                Then $x\in\openrayright{R}{X}{a}$.
                Thus $a\mathrel{R}x$ by \cref{openray_right}.
                \begin{byCase}
                \caseOf{$x\mathrel{R}a$.}
                    Contradiction by \cref{asymmetric}.
                \caseOf{$x = a$.}
                    Thus $a\mathrel{R}a$ by \cref{real_lt_subst_left}.
                    Contradiction by \cref{asymmetric}.
                \end{byCase}
            \end{subproof}
            Thus $x\in X\setminus \openrayright{R}{X}{a}$ by \cref{setminus_intro}.
        \end{subproof}
        Thus $\closedrayleft{R}{X}{a} \subseteq X\setminus \openrayright{R}{X}{a}$ by \cref{subseteq}.
    \end{subproof}
    Show that $X\setminus \openrayright{R}{X}{a} \subseteq \closedrayleft{R}{X}{a}$.
    \begin{subproof}
        Show that for all $x$ such that $x\in X\setminus \openrayright{R}{X}{a}$ we have
            $x\in\closedrayleft{R}{X}{a}$.
        \begin{subproof}
            Fix $x$.
            Assume $x\in X\setminus \openrayright{R}{X}{a}$.
            Thus $x\in X$ and $x\notin\openrayright{R}{X}{a}$ by \cref{setminus_elim_left,setminus_elim_right}.
            Show that it is wrong that $a\mathrel{R}x$.
            \begin{subproof}
                Suppose not.
                Then $a\mathrel{R}x$.
                Thus $x\in\openrayright{R}{X}{a}$ by \cref{openray_right}.
                Contradiction by \cref{setminus_elim_right}.
            \end{subproof}
            \begin{byCase}
            \caseOf{$x = a$.}
                Thus $x\in\closedrayleft{R}{X}{a}$ by \cref{closedray_left}.
            \caseOf{$x \neq a$.}
                $R$ is connex on $X$ by \cref{simple_order_relation}.
                Thus $x\mathrel{R}a$ or $a\mathrel{R}x$ by \cref{connex}.
                Thus $x\mathrel{R}a$.
                Thus $x\in\closedrayleft{R}{X}{a}$ by \cref{closedray_left}.
            \end{byCase}
        \end{subproof}
        Thus $X\setminus \openrayright{R}{X}{a} \subseteq \closedrayleft{R}{X}{a}$ by \cref{subseteq}.
    \end{subproof}
    Thus $\closedrayleft{R}{X}{a} = X\setminus \openrayright{R}{X}{a}$ by \cref{subseteq_antisymmetric}.
\end{proof}

% from §25 helper lemma 13: closed right ray is closed in the order topology
\begin{lemma}\label{closedrayright_closed_in_order_topology}
    Assume $R$ is a simple order relation on $X$.
    Assume $X$ has more than one element.
    Assume $T$ is the order topology on $X$ with respect to $R$.
    Suppose $a\in X$.
    Then $\closedrayright{R}{X}{a}$ is closed in $X$ with respect to $T$.
\end{lemma}
\begin{proof}
    Let $U = \openrayleft{R}{X}{a}$.
    $U$ is open in $X$ with respect to $T$ by \cref{openrayleft_open_in_order_topology}.
    $\closedrayright{R}{X}{a} = X\setminus U$ by \cref{closedrayright_complement_openrayleft}.
    Show that $\closedrayright{R}{X}{a}\subseteq X$.
    \begin{subproof}
        Show that for all $x$ such that $x\in\closedrayright{R}{X}{a}$ we have $x\in X$.
        \begin{subproof}
            Fix $x$.
            Assume $x\in\closedrayright{R}{X}{a}$.
            Thus $x\in X$ by \cref{closedray_right}.
        \end{subproof}
        Thus $\closedrayright{R}{X}{a}\subseteq X$ by \cref{subseteq}.
    \end{subproof}
    Show that $X\setminus \closedrayright{R}{X}{a} = U$.
    \begin{subproof}
        Show that $U\subseteq X$.
        \begin{subproof}
            Show that for all $x$ such that $x\in U$ we have $x\in X$.
            \begin{subproof}
                Fix $x$.
                Assume $x\in U$.
                Thus $x\in X$ by \cref{openray_left}.
            \end{subproof}
            Thus $U\subseteq X$ by \cref{subseteq}.
        \end{subproof}
        Thus $X\setminus \closedrayright{R}{X}{a} = X\setminus (X\setminus U)$ by \cref{subseteq}.
        Thus $X\setminus \closedrayright{R}{X}{a} = X\inter U$ by \cref{setminus_setminus}.
        Thus $X\inter U = U$ by \cref{inter_absorb_supseteq_left}.
        Thus $X\setminus \closedrayright{R}{X}{a} = U$ by \cref{subseteq}.
    \end{subproof}
    Thus $X\setminus \closedrayright{R}{X}{a}$ is open in $X$ with respect to $T$ by \cref{open_in}.
    Thus $\closedrayright{R}{X}{a}$ is closed in $X$ with respect to $T$ by \cref{closed_in}.
\end{proof}

% from §25 helper lemma 14: closed left ray is closed in the order topology
\begin{lemma}\label{closedrayleft_closed_in_order_topology}
    Assume $R$ is a simple order relation on $X$.
    Assume $X$ has more than one element.
    Assume $T$ is the order topology on $X$ with respect to $R$.
    Suppose $a\in X$.
    Then $\closedrayleft{R}{X}{a}$ is closed in $X$ with respect to $T$.
\end{lemma}
\begin{proof}
    Let $U = \openrayright{R}{X}{a}$.
    $U$ is open in $X$ with respect to $T$ by \cref{openrayright_open_in_order_topology}.
    $\closedrayleft{R}{X}{a} = X\setminus U$ by \cref{closedrayleft_complement_openrayright}.
    Show that $\closedrayleft{R}{X}{a}\subseteq X$.
    \begin{subproof}
        Show that for all $x$ such that $x\in\closedrayleft{R}{X}{a}$ we have $x\in X$.
        \begin{subproof}
            Fix $x$.
            Assume $x\in\closedrayleft{R}{X}{a}$.
            Thus $x\in X$ by \cref{closedray_left}.
        \end{subproof}
        Thus $\closedrayleft{R}{X}{a}\subseteq X$ by \cref{subseteq}.
    \end{subproof}
    Show that $X\setminus \closedrayleft{R}{X}{a} = U$.
    \begin{subproof}
        Show that $U\subseteq X$.
        \begin{subproof}
            Show that for all $x$ such that $x\in U$ we have $x\in X$.
            \begin{subproof}
                Fix $x$.
                Assume $x\in U$.
                Thus $x\in X$ by \cref{openray_right}.
            \end{subproof}
            Thus $U\subseteq X$ by \cref{subseteq}.
        \end{subproof}
        Thus $X\setminus \closedrayleft{R}{X}{a} = X\setminus (X\setminus U)$ by \cref{subseteq}.
        Thus $X\setminus \closedrayleft{R}{X}{a} = X\inter U$ by \cref{setminus_setminus}.
        Thus $X\inter U = U$ by \cref{inter_absorb_supseteq_left}.
        Thus $X\setminus \closedrayleft{R}{X}{a} = U$ by \cref{subseteq}.
    \end{subproof}
    Thus $X\setminus \closedrayleft{R}{X}{a}$ is open in $X$ with respect to $T$ by \cref{open_in}.
    Thus $\closedrayleft{R}{X}{a}$ is closed in $X$ with respect to $T$ by \cref{closed_in}.
\end{proof}

% from §25 helper lemma 15: closed interval is closed in the order topology
\begin{lemma}\label{intervalclosedrel_closed_in_order_topology}
    Assume $R$ is a simple order relation on $X$.
    Assume $X$ has more than one element.
    Assume $T$ is the order topology on $X$ with respect to $R$.
    Suppose $a\in X$.
    Suppose $b\in X$.
    Then $\intervalclosedrel{R}{X}{a}{b}$ is closed in $X$ with respect to $T$.
\end{lemma}
\begin{proof}
    Let $A = \closedrayright{R}{X}{a}$.
    Let $B = \closedrayleft{R}{X}{b}$.
    $\orderBasis{R}{X}$ is a basis on $X$ by \cref{order_basis_is_basis}.
    Thus $T$ is a topology on $X$ by \cref{basis_topology_is_topology,order_topology}.
    $A$ is closed in $X$ with respect to $T$ by \cref{closedrayright_closed_in_order_topology}.
    $B$ is closed in $X$ with respect to $T$ by \cref{closedrayleft_closed_in_order_topology}.
    Let $C = \{A,B\}$.
    Show that $C\subseteq\pow{X}$.
    \begin{subproof}
        Show that for all $D$ such that $D\in C$ we have $D\in\pow{X}$.
        \begin{subproof}
            Fix $D$.
            Assume $D\in C$.
            Then $D = A$ or $D = B$ by \cref{upair_iff}.
            \begin{byCase}
            \caseOf{$D = A$.}
                Show that $D\subseteq X$.
                \begin{subproof}
                    Show that for all $x$ such that $x\in D$ we have $x\in X$.
                    \begin{subproof}
                        Fix $x$.
                        Assume $x\in D$.
                        Thus $x\in A$.
                        Thus $x\in X$ by \cref{closedray_right}.
                    \end{subproof}
                    Thus $D\subseteq X$ by \cref{subseteq}.
                \end{subproof}
                Thus $D\in\pow{X}$ by \cref{powerset_intro}.
            \caseOf{$D = B$.}
                Show that $D\subseteq X$.
                \begin{subproof}
                    Show that for all $x$ such that $x\in D$ we have $x\in X$.
                    \begin{subproof}
                        Fix $x$.
                        Assume $x\in D$.
                        Thus $x\in B$.
                        Thus $x\in X$ by \cref{closedray_left}.
                    \end{subproof}
                    Thus $D\subseteq X$ by \cref{subseteq}.
                \end{subproof}
                Thus $D\in\pow{X}$ by \cref{powerset_intro}.
            \end{byCase}
        \end{subproof}
        Thus $C\subseteq\pow{X}$ by \cref{subseteq}.
    \end{subproof}
    Show that for all $D$ such that $D\in C$ we have $D$ is closed in $X$ with respect to $T$.
    \begin{subproof}
        Fix $D$.
        Assume $D\in C$.
        Then $D = A$ or $D = B$ by \cref{upair_iff}.
        Thus $D$ is closed in $X$ with respect to $T$ by \cref{closedrayright_closed_in_order_topology,closedrayleft_closed_in_order_topology}.
    \end{subproof}
    Thus $\inters{C}$ is closed in $X$ with respect to $T$ by \cref{closed_inters}.
    Show that $\intervalclosedrel{R}{X}{a}{b} = \inters{C}$.
    \begin{subproof}
        Show that $\intervalclosedrel{R}{X}{a}{b} \subseteq \inters{C}$.
        \begin{subproof}
            Show that for all $x$ such that $x\in\intervalclosedrel{R}{X}{a}{b}$ we have $x\in\inters{C}$.
            \begin{subproof}
                Fix $x$.
                Assume $x\in\intervalclosedrel{R}{X}{a}{b}$.
                Thus $x\in X$ by \cref{intervalclosed_rel}.
                Thus $a\mathrel{R}x$ or $x = a$ by \cref{intervalclosed_rel}.
                Thus $x\mathrel{R}b$ or $x = b$ by \cref{intervalclosed_rel}.
                Thus $x\in A$ and $x\in B$ by \cref{closedray_right,closedray_left}.
                Show that there exists $D$ such that $D\in C$.
                \begin{subproof}
                    $A\in C$ by \cref{upair_iff}.
                    Thus there exists $D$ such that $D\in C$.
                \end{subproof}
                Show that for all $D$ such that $D\in C$ we have $x\in D$.
                \begin{subproof}
                    Fix $D$.
                    Assume $D\in C$.
                    Then $D = A$ or $D = B$ by \cref{upair_iff}.
                    Thus $x\in D$.
                \end{subproof}
                Thus $x\in\inters{C}$ by \cref{inters_iff_forall}.
            \end{subproof}
            Thus $\intervalclosedrel{R}{X}{a}{b} \subseteq \inters{C}$ by \cref{subseteq}.
        \end{subproof}
        Show that $\inters{C} \subseteq \intervalclosedrel{R}{X}{a}{b}$.
        \begin{subproof}
            Show that for all $x$ such that $x\in\inters{C}$ we have $x\in\intervalclosedrel{R}{X}{a}{b}$.
            \begin{subproof}
                Fix $x$.
                Assume $x\in\inters{C}$.
                Thus $x\in A$ and $x\in B$ by \cref{inters_iff_forall,upair_iff}.
                Thus $x\in X$ by \cref{closedray_right}.
                Thus $a\mathrel{R}x$ or $x = a$ by \cref{closedray_right}.
                Thus $x\mathrel{R}b$ or $x = b$ by \cref{closedray_left}.
                Thus $x\in\intervalclosedrel{R}{X}{a}{b}$ by \cref{intervalclosed_rel}.
            \end{subproof}
            Thus $\inters{C} \subseteq \intervalclosedrel{R}{X}{a}{b}$ by \cref{subseteq}.
        \end{subproof}
        Thus $\intervalclosedrel{R}{X}{a}{b} = \inters{C}$ by \cref{subseteq_antisymmetric}.
    \end{subproof}
    Thus $\intervalclosedrel{R}{X}{a}{b}$ is closed in $X$ with respect to $T$ by \cref{subseteq}.
\end{proof}

% from §25 helper lemma 16: reverse a path
\begin{lemma}\label{path_reverse}
    Assume $T$ is a topology on $X$.
    Suppose $f$ is a path in $X$ from $x$ to $y$ with respect to $T$.
    Then there exists $g$ such that $g$ is a path in $X$ from $y$ to $x$ with respect to $T$.
\end{lemma}
\begin{proof}
    Take $a,b\in\reals$ such that
        $a\mathrel{\realOrderRel} b$ and
        $f$ is continuous from $\intervalclosedrel{\realOrderRel}{\reals}{a}{b}$ to $X$ with respect to
            $\subspaceTopology{\realOrderTopology}{\intervalclosedrel{\realOrderRel}{\reals}{a}{b}}$ and $T$ and
        $f(a) = x$ and $f(b) = y$
        by \cref{path_in_space}.
    Let $I = \intervalclosedrel{\realOrderRel}{\reals}{a}{b}$.
    Let $T_I = \subspaceTopology{\realOrderTopology}{I}$.
    Show that $I\subseteq\reals$.
    \begin{subproof}
        Show that for all $t$ such that $t\in I$ we have $t\in\reals$.
        \begin{subproof}
            Fix $t$.
            Assume $t\in I$.
            Thus $t\in\reals$ by \cref{intervalclosed_rel}.
        \end{subproof}
        Thus $I\subseteq\reals$ by \cref{subseteq}.
    \end{subproof}
    $\standardBasisR$ is a basis on $\reals$ by \cref{standard_basis_is_basis}.
    $\standardTopologyR = \basisTopology{\standardBasisR}{\reals}$ by \cref{standard_topology_reals}.
    Thus $\standardTopologyR$ is a topology on $\reals$ by \cref{basis_topology_is_topology}.
    $\realOrderTopology = \standardTopologyR$ by \cref{real_order_topology_eq_standard}.
    Thus $\realOrderTopology$ is a topology on $\reals$.
    Thus $I$ is a subspace of $\reals$ with respect to $\realOrderTopology$ by \cref{subspace_of}.
    Thus $T_I$ is a topology on $I$ by \cref{subspace_topology_is_topology,subspace_of}.
    Let $c_0 = a + b$.
    $c_0\in\reals$ by \cref{real_add_closed}.
    Let $k = \{(t,c_0) \mid t\in I\}$.
    $k$ is a function from $I$ to $\reals$ by \cref{constant_map_function}.
    Show that for all $t\in I$ we have $k(t) = c_0$.
    \begin{subproof}
        Fix $t$.
        Assume $t\in I$.
        Thus $(t,c_0)\in k$.
        $k$ is a function.
        Thus $k(t) = c_0$ by \cref{function_apply_iff}.
    \end{subproof}
    Thus $k$ is continuous from $I$ to $\reals$ with respect to $T_I$ and $\realOrderTopology$
        by \cref{constant_continuous}.
    Thus $k$ is continuous from $I$ to $\reals$ with respect to $T_I$ and $\standardTopologyR$
        by \cref{real_order_topology_eq_standard}.
    Let $j = \identity{I}$.
    $j$ is continuous from $I$ to $\reals$ with respect to $T_I$ and $\realOrderTopology$ by \cref{inclusion_continuous}.
    Thus $j$ is continuous from $I$ to $\reals$ with respect to $T_I$ and $\standardTopologyR$
        by \cref{real_order_topology_eq_standard}.
    Let $r = \{(t, k(t) - j(t)) \mid t\in I\}$.
    $r$ is continuous from $I$ to $\reals$ with respect to $T_I$ and $\standardTopologyR$
        by \cref{real_function_subtraction_continuous}.
    Thus $r$ is continuous from $I$ to $\reals$ with respect to $T_I$ and $\realOrderTopology$
        by \cref{real_order_topology_eq_standard}.
    Show that $\img{r}{I}\subseteq I$.
    \begin{subproof}
        Show that for all $t$ such that $t\in I$ we have $r(t)\in I$.
        \begin{subproof}
            Fix $t$.
            Assume $t\in I$.
            Thus $t\in\reals$ by \cref{intervalclosed_rel}.
            Thus $a\mathrel{\realOrderRel}t$ or $t = a$ by \cref{intervalclosed_rel}.
            Thus $t\mathrel{\realOrderRel}b$ or $t = b$ by \cref{intervalclosed_rel}.
            Show that $j(t) = t$.
            \begin{subproof}
                $t\in I$.
                Thus $t\mathrel{j} t$ by \cref{id_iff}.
                $j\in\funs{I}{I}$ by \cref{id_elem_funs}.
                Thus $j$ is a function by \cref{funs_is_function}.
                Thus $j(t) = t$ by \cref{function_apply_iff}.
            \end{subproof}
            Show that $r(t) = c_0 - t$.
            \begin{subproof}
                $(t, k(t) - j(t))\in r$.
                $r\in\funs{I}{\reals}$ by \cref{continuous}.
                Thus $r$ is a function by \cref{funs_is_function}.
                Thus $r(t) = k(t) - j(t)$ by \cref{function_apply_iff}.
                Thus $k(t) = c_0$.
                Thus $j(t) = t$.
                Thus $r(t) = c_0 - t$.
            \end{subproof}
            Show that $r(t)\in\reals$.
            \begin{subproof}
                $\neg{t}\in\reals$ by \cref{real_add_inverse}.
                $c_0 + \neg{t}\in\reals$ by \cref{real_add_closed}.
                Thus $c_0 - t\in\reals$ by \cref{real_sub}.
                Thus $r(t)\in\reals$.
            \end{subproof}
            Show that $a\mathrel{\realOrderRel}r(t)$ or $r(t) = a$.
            \begin{subproof}
                \begin{byCase}
                \caseOf{$t = b$.}
                    $b + \neg{b} = 0$ by \cref{real_add_inverse}.
                    $c_0 - b = c_0 + \neg{b}$ by \cref{real_sub}.
                    $c_0 = a + b$.
                    Thus $c_0 + \neg{b} = (a + b) + \neg{b}$ by \cref{real_add_eq_subst}.
                    $a\in\reals$.
                    $b\in\reals$.
                    $\neg{b}\in\reals$ by \cref{real_add_inverse}.
                    $(a + b) + \neg{b} = a + (b + \neg{b})$ by \cref{real_add_assoc}.
                    Thus $r(t) = a + 0$ by \cref{real_add_eq_subst}.
                    Thus $r(t) = a$ by \cref{real_add_ident,real_add_comm}.
                \caseOf{$t\mathrel{\realOrderRel}b$.}
                    Thus $t < b$ by \cref{real_order_relation_elim}.
                    $b\in\reals$.
                    $\neg{t}\in\reals$ by \cref{real_add_inverse}.
                    Thus $t + \neg{t} < b + \neg{t}$ by \cref{real_lt_add}.
                    $t + \neg{t} = 0$ by \cref{real_add_inverse}.
                    Thus $0 < b + \neg{t}$ by \cref{real_lt_subst_left}.
                    $b + \neg{t} = b - t$ by \cref{real_sub}.
                    Thus $0 < b - t$ by \cref{real_lt_subst_right}.
                    $0\in\reals$ by \cref{real_add_ident}.
                    $a\in\reals$.
                    $b - t = b + \neg{t}$ by \cref{real_sub}.
                    Thus $b - t\in\reals$ by \cref{real_add_closed}.
                    Thus $0 + a < (b - t) + a$ by \cref{real_lt_add}.
                    $0 + a = a$ by \cref{real_add_ident}.
                    $(b - t) + a = a + (b - t)$ by \cref{real_add_comm}.
                    Thus $a < a + (b - t)$ by \cref{real_lt_subst_left,real_lt_subst_right}.
                    $b - t = b + \neg{t}$ by \cref{real_sub}.
                    $a + (b - t) = a + (b + \neg{t})$ by \cref{real_add_eq_subst}.
                    $a + (b + \neg{t}) = (a + b) + \neg{t}$ by \cref{real_add_assoc}.
                    $(a + b) + \neg{t} = (a + b) - t$ by \cref{real_sub}.
                    Thus $a + (b - t) = c_0 - t$.
                    Thus $a < c_0 - t$ by \cref{real_lt_subst_right}.
                    Thus $a\mathrel{\realOrderRel}r(t)$ by \cref{real_order_relation_intro}.
                \end{byCase}
            \end{subproof}
            Show that $r(t)\mathrel{\realOrderRel}b$ or $r(t) = b$.
            \begin{subproof}
                \begin{byCase}
                \caseOf{$t = a$.}
                    $a + b = b + a$ by \cref{real_add_comm}.
                    $c_0 = a + b$.
                    Thus $c_0 = b + a$ by \cref{real_add_eq_subst}.
                    $c_0 - a = c_0 + \neg{a}$ by \cref{real_sub}.
                    Thus $c_0 - a = (b + a) + \neg{a}$ by \cref{real_add_eq_subst}.
                    $a\in\reals$.
                    $b\in\reals$.
                    $\neg{a}\in\reals$ by \cref{real_add_inverse}.
                    $(b + a) + \neg{a} = b + (a + \neg{a})$ by \cref{real_add_assoc}.
                    $a + \neg{a} = 0$ by \cref{real_add_inverse}.
                    Thus $r(t) = b + 0$ by \cref{real_add_eq_subst}.
                    Thus $r(t) = b$ by \cref{real_add_ident,real_add_comm}.
                \caseOf{$a\mathrel{\realOrderRel}t$.}
                    Thus $a < t$ by \cref{real_order_relation_elim}.
                    $a\in\reals$.
                    $t\in\reals$.
                    $b\in\reals$.
                    $\neg{t}\in\reals$ by \cref{real_add_inverse}.
                    $b - t = b + \neg{t}$ by \cref{real_sub}.
                    Thus $b - t\in\reals$ by \cref{real_add_closed}.
                    Thus $a + (b - t) < t + (b - t)$ by \cref{real_lt_add}.
                    $b - t = b + \neg{t}$ by \cref{real_sub}.
                    $t + (b - t) = t + (b + \neg{t})$ by \cref{real_add_eq_subst}.
                    $t + (b + \neg{t}) = (t + b) + \neg{t}$ by \cref{real_add_assoc}.
                    $(t + b) + \neg{t} = (b + t) + \neg{t}$ by \cref{real_add_comm}.
                    $(b + t) + \neg{t} = b + (t + \neg{t})$ by \cref{real_add_assoc}.
                    $t + \neg{t} = 0$ by \cref{real_add_inverse}.
                    Thus $t + (b - t) = b + 0$ by \cref{real_add_eq_subst}.
                    Thus $t + (b - t) = b$ by \cref{real_add_ident,real_add_comm}.
                    Thus $a + (b - t) < b$ by \cref{real_lt_subst_right}.
                    $b - t = b + \neg{t}$ by \cref{real_sub}.
                    $a + (b - t) = a + (b + \neg{t})$ by \cref{real_add_eq_subst}.
                    $a + (b + \neg{t}) = (a + b) + \neg{t}$ by \cref{real_add_assoc}.
                    $(a + b) + \neg{t} = (a + b) - t$ by \cref{real_sub}.
                    Thus $a + (b - t) = c_0 - t$.
                    Thus $c_0 - t < b$ by \cref{real_lt_subst_left}.
                    Thus $r(t)\mathrel{\realOrderRel}b$ by \cref{real_order_relation_intro}.
                \end{byCase}
            \end{subproof}
            Thus $r(t)\in I$ by \cref{intervalclosed_rel}.
        \end{subproof}
        Show that for all $y$ such that $y\in\img{r}{I}$ we have $y\in I$.
        \begin{subproof}
            Fix $y$.
            Assume $y\in\img{r}{I}$.
            Take $t\in I$ such that $t\mathrel{r} y$ by \cref{img_iff}.
            $r\in\funs{I}{\reals}$ by \cref{continuous}.
            Thus $r$ is a function by \cref{funs_is_function}.
            Thus $y = r(t)$ by \cref{function_apply_iff}.
            Thus $y\in I$.
        \end{subproof}
        Thus $\img{r}{I}\subseteq I$ by \cref{subseteq}.
    \end{subproof}
    Thus $r$ is continuous from $I$ to $I$ with respect to $T_I$ and $T_I$ by \cref{continuous_restrict_range}.
    Let $g = f\circ r$.
    Show that $f$ is composable with $r$.
    \begin{subproof}
        $r\in\funs{I}{\reals}$ by \cref{continuous}.
        Thus $r$ is a function by \cref{funs_is_function}.
        Thus $\dom{r} = I$ by \cref{funs}.
        Thus $\img{r}{I} = \ran{r}$ by \cref{img_dom}.
        Thus $\ran{r}\subseteq I$ by \cref{subseteq}.
        $f\in\funs{I}{X}$ by \cref{continuous}.
        Thus $\dom{f} = I$ by \cref{funs}.
        Thus $\ran{r}\subseteq \dom{f}$ by \cref{subseteq}.
        Thus $f$ is composable with $r$.
    \end{subproof}
    Thus $g$ is continuous from $I$ to $X$ with respect to $T_I$ and $T$ by \cref{continuous_composite}.
    Show that $r(a) = b$ and $r(b) = a$.
    \begin{subproof}
        Show that $a\in I$ and $b\in I$.
        \begin{subproof}
            $a = a$.
            $b = b$.
            Thus $a\mathrel{\realOrderRel}a$ or $a = a$.
            Thus $b\mathrel{\realOrderRel}b$ or $b = b$.
            Thus $a\mathrel{\realOrderRel}b$ or $a = b$.
            Thus $a\in I$ and $b\in I$ by \cref{intervalclosed_rel}.
        \end{subproof}
        Show that $r(a) = c_0 - a$.
        \begin{subproof}
            $a\in I$.
            $(a, k(a) - j(a))\in r$.
            $r\in\funs{I}{\reals}$ by \cref{continuous}.
            Thus $r$ is a function by \cref{funs_is_function}.
            Thus $r(a) = k(a) - j(a)$ by \cref{function_apply_iff}.
            $(a,c_0)\in k$.
            $k$ is a function by \cref{funs_is_function,constant_map_function}.
            Thus $k(a) = c_0$ by \cref{function_apply_iff}.
            $a\mathrel{j} a$ by \cref{id_iff}.
            $j\in\funs{I}{I}$ by \cref{id_elem_funs}.
            Thus $j$ is a function by \cref{funs_is_function}.
            Thus $j(a) = a$ by \cref{function_apply_iff}.
            Thus $r(a) = c_0 - a$.
        \end{subproof}
        Show that $r(b) = c_0 - b$.
        \begin{subproof}
            $b\in I$.
            $(b, k(b) - j(b))\in r$.
            $r\in\funs{I}{\reals}$ by \cref{continuous}.
            Thus $r$ is a function by \cref{funs_is_function}.
            Thus $r(b) = k(b) - j(b)$ by \cref{function_apply_iff}.
            $(b,c_0)\in k$.
            $k$ is a function by \cref{funs_is_function,constant_map_function}.
            Thus $k(b) = c_0$ by \cref{function_apply_iff}.
            $b\mathrel{j} b$ by \cref{id_iff}.
            $j\in\funs{I}{I}$ by \cref{id_elem_funs}.
            Thus $j$ is a function by \cref{funs_is_function}.
            Thus $j(b) = b$ by \cref{function_apply_iff}.
            Thus $r(b) = c_0 - b$.
        \end{subproof}
        $a + b = b + a$ by \cref{real_add_comm}.
        $c_0 = a + b$.
        Thus $c_0 = b + a$ by \cref{real_add_eq_subst}.
        $c_0 - a = c_0 + \neg{a}$ by \cref{real_sub}.
        Thus $c_0 - a = (b + a) + \neg{a}$ by \cref{real_add_eq_subst}.
        $a\in\reals$.
        $b\in\reals$.
        $\neg{a}\in\reals$ by \cref{real_add_inverse}.
        $(b + a) + \neg{a} = b + (a + \neg{a})$ by \cref{real_add_assoc}.
        $a + \neg{a} = 0$ by \cref{real_add_inverse}.
        Thus $r(a) = b + 0$ by \cref{real_add_eq_subst}.
        Thus $r(a) = b$ by \cref{real_add_ident,real_add_comm}.
        $c_0 - b = c_0 + \neg{b}$ by \cref{real_sub}.
        $c_0 + \neg{b} = (a + b) + \neg{b}$.
        $a\in\reals$.
        $b\in\reals$.
        $\neg{b}\in\reals$ by \cref{real_add_inverse}.
        $(a + b) + \neg{b} = a + (b + \neg{b})$ by \cref{real_add_assoc}.
        $b + \neg{b} = 0$ by \cref{real_add_inverse}.
        Thus $r(b) = a + 0$ by \cref{real_add_eq_subst}.
        Thus $r(b) = a$ by \cref{real_add_ident,real_add_comm}.
    \end{subproof}
    Show that $g(a) = y$ and $g(b) = x$.
    \begin{subproof}
        $a = a$.
        $b = b$.
        Thus $a\mathrel{\realOrderRel}a$ or $a = a$.
        Thus $b\mathrel{\realOrderRel}b$ or $b = b$.
        Thus $a\mathrel{\realOrderRel}b$ or $a = b$.
        Thus $a\in I$ and $b\in I$ by \cref{intervalclosed_rel}.
        $r\in\funs{I}{\reals}$ by \cref{continuous}.
        Thus $r$ is a function by \cref{funs_is_function}.
        Thus $\dom{r} = I$ by \cref{funs}.
        Thus $a\in\dom{r}$ and $b\in\dom{r}$ by \cref{subseteq}.
        Show that $g(a) = f(r(a))$.
        \begin{subproof}
            $r\in\funs{I}{\reals}$ by \cref{continuous}.
            Thus $r$ is a function by \cref{funs_is_function}.
            $f\in\funs{I}{X}$ by \cref{continuous}.
            Thus $f$ is a function by \cref{funs_is_function}.
            Thus $g(a) = f(r(a))$ by \cref{circ_apply}.
        \end{subproof}
        Show that $g(b) = f(r(b))$.
        \begin{subproof}
            $r\in\funs{I}{\reals}$ by \cref{continuous}.
            Thus $r$ is a function by \cref{funs_is_function}.
            $f\in\funs{I}{X}$ by \cref{continuous}.
            Thus $f$ is a function by \cref{funs_is_function}.
            Thus $g(b) = f(r(b))$ by \cref{circ_apply}.
        \end{subproof}
        Thus $g(a) = y$ and $g(b) = x$.
    \end{subproof}
    Thus $g$ is a path in $X$ from $y$ to $x$ with respect to $T$ by \cref{path_in_space}.
    Thus there exists $g$ such that $g$ is a path in $X$ from $y$ to $x$ with respect to $T$.
\end{proof}

% from §25 helper lemma 17: concatenate paths
\begin{lemma}\label{path_concatenation}
    Assume $T$ is a topology on $X$.
    Suppose $f$ is a path in $X$ from $x$ to $y$ with respect to $T$.
    Suppose $g$ is a path in $X$ from $y$ to $z$ with respect to $T$.
    Then there exists $h$ such that $h$ is a path in $X$ from $x$ to $z$ with respect to $T$.
\end{lemma}
\begin{proof}
    Take $a,b\in\reals$ such that
        $a\mathrel{\realOrderRel} b$ and
        $f$ is continuous from $\intervalclosedrel{\realOrderRel}{\reals}{a}{b}$ to $X$ with respect to
            $\subspaceTopology{\realOrderTopology}{\intervalclosedrel{\realOrderRel}{\reals}{a}{b}}$ and $T$ and
        $f(a) = x$ and $f(b) = y$
        by \cref{path_in_space}.
    Let $I_1 = \intervalclosedrel{\realOrderRel}{\reals}{a}{b}$.
    Let $T_1 = \subspaceTopology{\realOrderTopology}{I_1}$.
    Take $c,d\in\reals$ such that
        $c\mathrel{\realOrderRel} d$ and
        $g$ is continuous from $\intervalclosedrel{\realOrderRel}{\reals}{c}{d}$ to $X$ with respect to
            $\subspaceTopology{\realOrderTopology}{\intervalclosedrel{\realOrderRel}{\reals}{c}{d}}$ and $T$ and
        $g(c) = y$ and $g(d) = z$
        by \cref{path_in_space}.
    Let $I_2 = \intervalclosedrel{\realOrderRel}{\reals}{c}{d}$.
    Let $T_2 = \subspaceTopology{\realOrderTopology}{I_2}$.
    Let $s = c - b$.
    $b\in\reals$.
    $c\in\reals$.
    $\neg{b}\in\reals$ by \cref{real_add_inverse}.
    $s\in\reals$ by \cref{real_sub,real_add_closed}.
    Let $e = b + (d - c)$.
    $d - c = d + \neg{c}$ by \cref{real_sub}.
    $\neg{c}\in\reals$ by \cref{real_add_inverse}.
    Thus $d - c\in\reals$ by \cref{real_add_closed}.
    Thus $e\in\reals$ by \cref{real_add_closed}.
    Let $I_3 = \intervalclosedrel{\realOrderRel}{\reals}{b}{e}$.
    Let $I = \intervalclosedrel{\realOrderRel}{\reals}{a}{e}$.
    Let $T_I = \subspaceTopology{\realOrderTopology}{I}$.
    Let $T_3 = \subspaceTopology{\realOrderTopology}{I_3}$.
    $\standardBasisR$ is a basis on $\reals$ by \cref{standard_basis_is_basis}.
    $\standardTopologyR = \basisTopology{\standardBasisR}{\reals}$ by \cref{standard_topology_reals}.
    Thus $\standardTopologyR$ is a topology on $\reals$ by \cref{basis_topology_is_topology}.
    $\realOrderTopology = \standardTopologyR$ by \cref{real_order_topology_eq_standard}.
    Thus $\realOrderTopology$ is a topology on $\reals$.
    Show that $I_1\subseteq\reals$.
    \begin{subproof}
        Show that for all $t$ such that $t\in I_1$ we have $t\in\reals$.
        \begin{subproof}
            Fix $t$.
            Assume $t\in I_1$.
            Thus $t\in\reals$ by \cref{intervalclosed_rel}.
        \end{subproof}
        Thus $I_1\subseteq\reals$ by \cref{subseteq}.
    \end{subproof}
    Show that $I_2\subseteq\reals$.
    \begin{subproof}
        Show that for all $t$ such that $t\in I_2$ we have $t\in\reals$.
        \begin{subproof}
            Fix $t$.
            Assume $t\in I_2$.
            Thus $t\in\reals$ by \cref{intervalclosed_rel}.
        \end{subproof}
        Thus $I_2\subseteq\reals$ by \cref{subseteq}.
    \end{subproof}
    Show that $I_3\subseteq\reals$.
    \begin{subproof}
        Show that for all $t$ such that $t\in I_3$ we have $t\in\reals$.
        \begin{subproof}
            Fix $t$.
            Assume $t\in I_3$.
            Thus $t\in\reals$ by \cref{intervalclosed_rel}.
        \end{subproof}
        Thus $I_3\subseteq\reals$ by \cref{subseteq}.
    \end{subproof}
    Show that $I\subseteq\reals$.
    \begin{subproof}
        Show that for all $t$ such that $t\in I$ we have $t\in\reals$.
        \begin{subproof}
            Fix $t$.
            Assume $t\in I$.
            Thus $t\in\reals$ by \cref{intervalclosed_rel}.
        \end{subproof}
        Thus $I\subseteq\reals$ by \cref{subseteq}.
    \end{subproof}
    Thus $I$ is a subspace of $\reals$ with respect to $\realOrderTopology$ by \cref{subspace_of}.
    Thus $T_I$ is a topology on $I$ by \cref{subspace_topology_is_topology,subspace_of}.
    Thus $I_1$ is a subspace of $\reals$ with respect to $\realOrderTopology$ by \cref{subspace_of}.
    Thus $I_2$ is a subspace of $\reals$ with respect to $\realOrderTopology$ by \cref{subspace_of}.
    Thus $I_3$ is a subspace of $\reals$ with respect to $\realOrderTopology$ by \cref{subspace_of}.
    Thus $T_1$ is a topology on $I_1$ by \cref{subspace_topology_is_topology,subspace_of}.
    Thus $T_2$ is a topology on $I_2$ by \cref{subspace_topology_is_topology,subspace_of}.
    Thus $T_3$ is a topology on $I_3$ by \cref{subspace_topology_is_topology,subspace_of}.
    Let $k = \{(t,s) \mid t\in I_3\}$.
    $k$ is a function from $I_3$ to $\reals$ by \cref{constant_map_function}.
    Show that for all $t\in I_3$ we have $k(t) = s$.
    \begin{subproof}
        Fix $t$.
        Assume $t\in I_3$.
        Thus $(t,s)\in k$.
        $k$ is a function.
        Thus $k(t) = s$ by \cref{function_apply_iff}.
    \end{subproof}
    Thus $k$ is continuous from $I_3$ to $\reals$ with respect to $T_3$ and $\realOrderTopology$
        by \cref{constant_continuous}.
    Thus $k$ is continuous from $I_3$ to $\reals$ with respect to $T_3$ and $\standardTopologyR$
        by \cref{real_order_topology_eq_standard}.
    Let $j = \identity{I_3}$.
    $j$ is continuous from $I_3$ to $\reals$ with respect to $T_3$ and $\realOrderTopology$
        by \cref{inclusion_continuous}.
    Thus $j$ is continuous from $I_3$ to $\reals$ with respect to $T_3$ and $\standardTopologyR$
        by \cref{real_order_topology_eq_standard}.
    Let $u = \{(t, j(t) + k(t)) \mid t\in I_3\}$.
    $u$ is continuous from $I_3$ to $\reals$ with respect to $T_3$ and $\standardTopologyR$
        by \cref{real_function_addition_continuous}.
    Thus $u$ is continuous from $I_3$ to $\reals$ with respect to $T_3$ and $\realOrderTopology$
        by \cref{real_order_topology_eq_standard}.
    Show that $\img{u}{I_3}\subseteq I_2$.
    \begin{subproof}
        Show that for all $t$ such that $t\in I_3$ we have $u(t)\in I_2$.
        \begin{subproof}
            Fix $t$.
            Assume $t\in I_3$.
            Thus $t\in\reals$ by \cref{intervalclosed_rel}.
            Thus $b\mathrel{\realOrderRel}t$ or $t = b$ by \cref{intervalclosed_rel}.
            Thus $t\mathrel{\realOrderRel}e$ or $t = e$ by \cref{intervalclosed_rel}.
            Show that $j(t) = t$.
            \begin{subproof}
                $t\mathrel{j} t$ by \cref{id_iff}.
                $j\in\funs{I_3}{I_3}$ by \cref{id_elem_funs}.
                Thus $j$ is a function by \cref{funs_is_function}.
                Thus $j(t) = t$ by \cref{function_apply_iff}.
            \end{subproof}
            Show that $u(t) = t + s$.
            \begin{subproof}
                $(t, j(t) + k(t))\in u$.
                $u\in\funs{I_3}{\reals}$ by \cref{continuous}.
                Thus $u$ is a function by \cref{funs_is_function}.
                Thus $u(t) = j(t) + k(t)$ by \cref{function_apply_iff}.
                Thus $k(t) = s$.
                Thus $u(t) = t + s$.
            \end{subproof}
            Show that $u(t)\in\reals$.
            \begin{subproof}
                $s\in\reals$.
                Thus $t + s\in\reals$ by \cref{real_add_closed}.
                Thus $u(t)\in\reals$.
            \end{subproof}
            Show that $c\mathrel{\realOrderRel}u(t)$ or $u(t) = c$.
            \begin{subproof}
                \begin{byCase}
                \caseOf{$t = b$.}
                    $s = c - b$.
                    $c - b = c + \neg{b}$ by \cref{real_sub}.
                    Thus $b + s = b + (c + \neg{b})$ by \cref{real_add_eq_subst}.
                    $b + (c + \neg{b}) = (b + c) + \neg{b}$ by \cref{real_add_assoc}.
                    $(b + c) + \neg{b} = (c + b) + \neg{b}$ by \cref{real_add_comm}.
                    $(c + b) + \neg{b} = c + (b + \neg{b})$ by \cref{real_add_assoc}.
                    $b + \neg{b} = 0$ by \cref{real_add_inverse}.
                    Thus $b + s = c + 0$ by \cref{real_add_eq_subst}.
                    Thus $b + s = c$ by \cref{real_add_ident,real_add_comm}.
                    Thus $u(t) = c$.
                \caseOf{$b\mathrel{\realOrderRel}t$.}
                    Thus $b < t$ by \cref{real_order_relation_elim}.
                    $b\in\reals$.
                    $t\in\reals$.
                    $s\in\reals$.
                    Thus $b + s < t + s$ by \cref{real_lt_add}.
                    $s = c - b$.
                    $c - b = c + \neg{b}$ by \cref{real_sub}.
                    Thus $b + s = b + (c + \neg{b})$ by \cref{real_add_eq_subst}.
                    $b + (c + \neg{b}) = (b + c) + \neg{b}$ by \cref{real_add_assoc}.
                    $(b + c) + \neg{b} = (c + b) + \neg{b}$ by \cref{real_add_comm}.
                    $(c + b) + \neg{b} = c + (b + \neg{b})$ by \cref{real_add_assoc}.
                    $b + \neg{b} = 0$ by \cref{real_add_inverse}.
                    Thus $b + s = c + 0$ by \cref{real_add_eq_subst}.
                    Thus $b + s = c$ by \cref{real_add_ident,real_add_comm}.
                    Thus $c < u(t)$ by \cref{real_lt_subst_left,real_lt_subst_right}.
                    Thus $c\mathrel{\realOrderRel}u(t)$ by \cref{real_order_relation_intro}.
                \end{byCase}
            \end{subproof}
            Show that $u(t)\mathrel{\realOrderRel}d$ or $u(t) = d$.
            \begin{subproof}
                \begin{byCase}
                \caseOf{$t = e$.}
                    $s = c - b$.
                    $e = b + (d - c)$.
                    $d - c = d + \neg{c}$ by \cref{real_sub}.
                    $c - b = c + \neg{b}$ by \cref{real_sub}.
                    Thus $e + s = (b + (d + \neg{c})) + (c + \neg{b})$ by \cref{real_add_eq_subst}.
                    $(b + (d + \neg{c})) + (c + \neg{b}) = b + ((d + \neg{c}) + (c + \neg{b}))$ by \cref{real_add_assoc}.
                    $(d + \neg{c}) + (c + \neg{b}) = ((d + \neg{c}) + c) + \neg{b}$ by \cref{real_add_assoc}.
                    $(d + \neg{c}) + c = d + (\neg{c} + c)$ by \cref{real_add_assoc}.
                    $\neg{c} + c = 0$ by \cref{real_add_comm,real_add_inverse}.
                    Thus $d + (\neg{c} + c) = d + 0$ by \cref{real_add_eq_subst}.
                    Thus $(d + \neg{c}) + c = d + 0$ by \cref{real_add_eq_subst}.
                    Thus $((d + \neg{c}) + c) + \neg{b} = (d + 0) + \neg{b}$ by \cref{real_add_eq_subst}.
                    $d + 0 = d$ by \cref{real_add_ident,real_add_comm}.
                    Thus $(d + 0) + \neg{b} = d + \neg{b}$ by \cref{real_add_eq_subst}.
                    Thus $(d + \neg{c}) + (c + \neg{b}) = d + \neg{b}$.
                    Thus $e + s = b + (d + \neg{b})$ by \cref{real_add_eq_subst}.
                    $b + (d + \neg{b}) = (b + d) + \neg{b}$ by \cref{real_add_assoc}.
                    $(b + d) + \neg{b} = (d + b) + \neg{b}$ by \cref{real_add_comm}.
                    $(d + b) + \neg{b} = d + (b + \neg{b})$ by \cref{real_add_assoc}.
                    $b + \neg{b} = 0$ by \cref{real_add_inverse}.
                    Thus $e + s = d + 0$ by \cref{real_add_eq_subst}.
                    Thus $e + s = d$ by \cref{real_add_ident,real_add_comm}.
                    Thus $u(t) = d$.
                \caseOf{$t\mathrel{\realOrderRel}e$.}
                    Thus $t < e$ by \cref{real_order_relation_elim}.
                    $t\in\reals$.
                    $e\in\reals$.
                    $s\in\reals$.
                    Thus $t + s < e + s$ by \cref{real_lt_add}.
                    $s = c - b$.
                    $e = b + (d - c)$.
                    $d - c = d + \neg{c}$ by \cref{real_sub}.
                    $c - b = c + \neg{b}$ by \cref{real_sub}.
                    Thus $e + s = (b + (d + \neg{c})) + (c + \neg{b})$ by \cref{real_add_eq_subst}.
                    $(b + (d + \neg{c})) + (c + \neg{b}) = b + ((d + \neg{c}) + (c + \neg{b}))$ by \cref{real_add_assoc}.
                    $(d + \neg{c}) + (c + \neg{b}) = ((d + \neg{c}) + c) + \neg{b}$ by \cref{real_add_assoc}.
                    $(d + \neg{c}) + c = d + (\neg{c} + c)$ by \cref{real_add_assoc}.
                    $\neg{c} + c = 0$ by \cref{real_add_comm,real_add_inverse}.
                    Thus $d + (\neg{c} + c) = d + 0$ by \cref{real_add_eq_subst}.
                    Thus $(d + \neg{c}) + c = d + 0$ by \cref{real_add_eq_subst}.
                    Thus $((d + \neg{c}) + c) + \neg{b} = (d + 0) + \neg{b}$ by \cref{real_add_eq_subst}.
                    $d + 0 = d$ by \cref{real_add_ident,real_add_comm}.
                    Thus $(d + 0) + \neg{b} = d + \neg{b}$ by \cref{real_add_eq_subst}.
                    Thus $(d + \neg{c}) + (c + \neg{b}) = d + \neg{b}$.
                    Thus $e + s = b + (d + \neg{b})$ by \cref{real_add_eq_subst}.
                    $b + (d + \neg{b}) = (b + d) + \neg{b}$ by \cref{real_add_assoc}.
                    $(b + d) + \neg{b} = (d + b) + \neg{b}$ by \cref{real_add_comm}.
                    $(d + b) + \neg{b} = d + (b + \neg{b})$ by \cref{real_add_assoc}.
                    $b + \neg{b} = 0$ by \cref{real_add_inverse}.
                    Thus $e + s = d + 0$ by \cref{real_add_eq_subst}.
                    Thus $e + s = d$ by \cref{real_add_ident,real_add_comm}.
                    Thus $u(t) < d$ by \cref{real_lt_subst_right}.
                    Thus $u(t)\mathrel{\realOrderRel}d$ by \cref{real_order_relation_intro}.
                \end{byCase}
            \end{subproof}
            Thus $u(t)\in I_2$ by \cref{intervalclosed_rel}.
        \end{subproof}
        Show that for all $y$ such that $y\in\img{u}{I_3}$ we have $y\in I_2$.
        \begin{subproof}
            Fix $y$.
            Assume $y\in\img{u}{I_3}$.
            Take $t\in I_3$ such that $t\mathrel{u} y$ by \cref{img_iff}.
            $u\in\funs{I_3}{\reals}$ by \cref{continuous}.
            Thus $u$ is a function by \cref{funs_is_function}.
            Thus $y = u(t)$ by \cref{function_apply_iff}.
            Thus $y\in I_2$.
        \end{subproof}
        Thus $\img{u}{I_3}\subseteq I_2$ by \cref{subseteq}.
    \end{subproof}
    Thus $u$ is continuous from $I_3$ to $I_2$ with respect to $T_3$ and $T_2$
        by \cref{continuous_restrict_range}.
    Let $g_1 = g\circ u$.
    Show that $g$ is composable with $u$.
    \begin{subproof}
        $u\in\funs{I_3}{I_2}$ by \cref{continuous}.
        Thus $u$ is a function by \cref{funs_is_function}.
        Thus $\dom{u} = I_3$ by \cref{funs}.
        Thus $\img{u}{I_3} = \ran{u}$ by \cref{img_dom}.
        Thus $\ran{u}\subseteq I_2$ by \cref{subseteq}.
        $g\in\funs{I_2}{X}$ by \cref{continuous}.
        Thus $\dom{g} = I_2$ by \cref{funs}.
        Thus $\ran{u}\subseteq \dom{g}$ by \cref{subseteq}.
        Thus $g$ is composable with $u$.
    \end{subproof}
    Thus $g_1$ is continuous from $I_3$ to $X$ with respect to $T_3$ and $T$ by \cref{continuous_composite}.
    Show that $b\in I_3$.
    \begin{subproof}
        $c\mathrel{\realOrderRel} d$.
        Thus $c < d$ by \cref{real_order_relation_elim}.
        $c\in\reals$.
        $d\in\reals$.
        Thus $c + \neg{c} < d + \neg{c}$ by \cref{real_lt_add}.
        $c + \neg{c} = 0$ by \cref{real_add_inverse}.
        Thus $0 < d + \neg{c}$ by \cref{real_lt_subst_left}.
        $d + \neg{c} = d - c$ by \cref{real_sub}.
        Thus $0 < d - c$ by \cref{real_lt_subst_right}.
        $b\in\reals$.
        $d - c\in\reals$.
        $0\in\reals$ by \cref{real_add_ident}.
        Thus $0 + b < (d - c) + b$ by \cref{real_lt_add}.
        $0 + b = b$ by \cref{real_add_comm,real_add_ident}.
        Thus $b < (d - c) + b$ by \cref{real_lt_subst_left}.
        $(d - c) + b = b + (d - c)$ by \cref{real_add_comm}.
        Thus $b < b + (d - c)$ by \cref{real_lt_subst_right}.
        Thus $b < e$ by \cref{real_lt_subst_right}.
        Thus $b\mathrel{\realOrderRel}e$ by \cref{real_order_relation_intro}.
        Thus $b = b$.
        Thus $b\mathrel{\realOrderRel}b$ or $b = b$.
        Thus $b\mathrel{\realOrderRel}e$ or $b = e$.
        Thus $b\in I_3$ by \cref{intervalclosed_rel}.
    \end{subproof}
    Show that $e\in I_3$.
    \begin{subproof}
        $b\in\reals$.
        $e\in\reals$.
        Thus $e = e$.
        Thus $e\mathrel{\realOrderRel}e$ or $e = e$.
        Show that $b\mathrel{\realOrderRel}e$ or $b = e$.
        \begin{subproof}
            $c\mathrel{\realOrderRel} d$.
            Thus $c < d$ by \cref{real_order_relation_elim}.
            $c\in\reals$.
            $d\in\reals$.
            Thus $c + \neg{c} < d + \neg{c}$ by \cref{real_lt_add}.
            $c + \neg{c} = 0$ by \cref{real_add_inverse}.
            Thus $0 < d + \neg{c}$ by \cref{real_lt_subst_left}.
            $d + \neg{c} = d - c$ by \cref{real_sub}.
            Thus $0 < d - c$ by \cref{real_lt_subst_right}.
            $b\in\reals$.
            $d - c\in\reals$.
            $0\in\reals$ by \cref{real_add_ident}.
            Thus $0 + b < (d - c) + b$ by \cref{real_lt_add}.
            $0 + b = b$ by \cref{real_add_comm,real_add_ident}.
            Thus $b < (d - c) + b$ by \cref{real_lt_subst_left}.
            $(d - c) + b = b + (d - c)$ by \cref{real_add_comm}.
            Thus $b < b + (d - c)$ by \cref{real_lt_subst_right}.
            Thus $b < e$ by \cref{real_lt_subst_right}.
            Thus $b\mathrel{\realOrderRel}e$ by \cref{real_order_relation_intro}.
            Thus $b\mathrel{\realOrderRel}e$ or $b = e$.
        \end{subproof}
        Thus $e\in I_3$ by \cref{intervalclosed_rel}.
    \end{subproof}
    Show that $u(b) = c$ and $u(e) = d$.
    \begin{subproof}
        Show that $u(b) = b + s$.
        \begin{subproof}
            $(b, j(b) + k(b))\in u$.
            $u\in\funs{I_3}{\reals}$ by \cref{continuous}.
            Thus $u$ is a function by \cref{funs_is_function}.
            Thus $u(b) = j(b) + k(b)$ by \cref{function_apply_iff}.
            $b\mathrel{j} b$ by \cref{id_iff}.
            $j\in\funs{I_3}{I_3}$ by \cref{id_elem_funs}.
            Thus $j$ is a function by \cref{funs_is_function}.
            Thus $j(b) = b$ by \cref{function_apply_iff}.
            Thus $k(b) = s$.
            Thus $u(b) = b + s$.
        \end{subproof}
        Show that $u(e) = e + s$.
        \begin{subproof}
            $(e, j(e) + k(e))\in u$.
            $u\in\funs{I_3}{\reals}$ by \cref{continuous}.
            Thus $u$ is a function by \cref{funs_is_function}.
            Thus $u(e) = j(e) + k(e)$ by \cref{function_apply_iff}.
            $e\mathrel{j} e$ by \cref{id_iff}.
            $j\in\funs{I_3}{I_3}$ by \cref{id_elem_funs}.
            Thus $j$ is a function by \cref{funs_is_function}.
            Thus $j(e) = e$ by \cref{function_apply_iff}.
            Thus $k(e) = s$.
            Thus $u(e) = e + s$.
        \end{subproof}
        $s = c - b$.
        $c - b = c + \neg{b}$ by \cref{real_sub}.
        Thus $b + s = b + (c + \neg{b})$ by \cref{real_add_eq_subst}.
        $b + (c + \neg{b}) = (b + c) + \neg{b}$ by \cref{real_add_assoc}.
        $(b + c) + \neg{b} = (c + b) + \neg{b}$ by \cref{real_add_comm}.
        $(c + b) + \neg{b} = c + (b + \neg{b})$ by \cref{real_add_assoc}.
        $b + \neg{b} = 0$ by \cref{real_add_inverse}.
        Thus $b + s = c + 0$ by \cref{real_add_eq_subst}.
        Thus $u(b) = c$ by \cref{real_add_ident,real_add_comm}.
        $s = c - b$.
        $e = b + (d - c)$.
        $d - c = d + \neg{c}$ by \cref{real_sub}.
        $c - b = c + \neg{b}$ by \cref{real_sub}.
        Thus $e + s = (b + (d + \neg{c})) + (c + \neg{b})$ by \cref{real_add_eq_subst}.
        $(b + (d + \neg{c})) + (c + \neg{b}) = b + ((d + \neg{c}) + (c + \neg{b}))$ by \cref{real_add_assoc}.
        $(d + \neg{c}) + (c + \neg{b}) = (d + (\neg{c} + c)) + \neg{b}$ by \cref{real_add_assoc}.
        $\neg{c} + c = 0$ by \cref{real_add_comm,real_add_inverse}.
        Thus $(d + \neg{c}) + (c + \neg{b}) = (d + 0) + \neg{b}$ by \cref{real_add_eq_subst}.
        $d + 0 = d$ by \cref{real_add_ident,real_add_comm}.
        Thus $(d + 0) + \neg{b} = d + \neg{b}$ by \cref{real_add_eq_subst}.
        Thus $(d + \neg{c}) + (c + \neg{b}) = d + \neg{b}$ by \cref{real_add_eq_subst}.
        Thus $e + s = b + (d + \neg{b})$ by \cref{real_add_eq_subst}.
        $b + (d + \neg{b}) = (b + d) + \neg{b}$ by \cref{real_add_assoc}.
        $(b + d) + \neg{b} = (d + b) + \neg{b}$ by \cref{real_add_comm}.
        $(d + b) + \neg{b} = d + (b + \neg{b})$ by \cref{real_add_assoc}.
        $b + \neg{b} = 0$ by \cref{real_add_inverse}.
        Thus $e + s = d + 0$ by \cref{real_add_eq_subst}.
        Thus $u(e) = d$ by \cref{real_add_ident,real_add_comm}.
    \end{subproof}
    Show that $g_1(b) = y$ and $g_1(e) = z$.
    \begin{subproof}
        Show that $g_1(b) = g(u(b))$.
        \begin{subproof}
            $u\in\funs{I_3}{I_2}$ by \cref{continuous}.
            Thus $u$ is a function by \cref{funs_is_function}.
            Show that $b\in\dom{u}$.
            \begin{subproof}
                Take $y\in I_2$ such that $(b,y)\in u$ by \cref{funs_tuple_intro}.
                Thus $b\in\dom{u}$ by \cref{dom_intro}.
            \end{subproof}
            $g\in\funs{I_2}{X}$ by \cref{continuous}.
            Thus $g$ is a function by \cref{funs_is_function}.
            Thus $g_1(b) = g(u(b))$ by \cref{circ_apply}.
        \end{subproof}
        Show that $g_1(e) = g(u(e))$.
        \begin{subproof}
            $u\in\funs{I_3}{I_2}$ by \cref{continuous}.
            Thus $u$ is a function by \cref{funs_is_function}.
            Show that $e\in\dom{u}$.
            \begin{subproof}
                Take $y\in I_2$ such that $(e,y)\in u$ by \cref{funs_tuple_intro}.
                Thus $e\in\dom{u}$ by \cref{dom_intro}.
            \end{subproof}
            $g\in\funs{I_2}{X}$ by \cref{continuous}.
            Thus $g$ is a function by \cref{funs_is_function}.
            Thus $g_1(e) = g(u(e))$ by \cref{circ_apply}.
        \end{subproof}
        Thus $g_1(b) = y$ and $g_1(e) = z$.
    \end{subproof}
    Show that $I = I_1\union I_3$.
    \begin{subproof}
        Show that $I_1\union I_3\subseteq I$.
        \begin{subproof}
            Show that for all $t$ such that $t\in I_1\union I_3$ we have $t\in I$.
            \begin{subproof}
                Fix $t$.
                Assume $t\in I_1\union I_3$.
                Thus $t\in I_1$ or $t\in I_3$ by \cref{union_iff}.
                \begin{byCase}
                \caseOf{$t\in I_1$.}
                    Thus $a\mathrel{\realOrderRel}t$ or $t = a$ by \cref{intervalclosed_rel}.
                    Thus $t\mathrel{\realOrderRel}b$ or $t = b$ by \cref{intervalclosed_rel}.
                    $c\mathrel{\realOrderRel} d$.
                    Thus $c < d$ by \cref{real_order_relation_elim}.
                    $c\in\reals$.
                    $d\in\reals$.
                    Thus $c + \neg{c} < d + \neg{c}$ by \cref{real_lt_add}.
                    $c + \neg{c} = 0$ by \cref{real_add_inverse}.
                    Thus $0 < d + \neg{c}$ by \cref{real_lt_subst_left}.
                    $d + \neg{c} = d - c$ by \cref{real_sub}.
                    Thus $0 < d - c$ by \cref{real_lt_subst_right}.
                    $b\in\reals$.
                    $d - c\in\reals$.
                    $0\in\reals$ by \cref{real_add_ident}.
                    Thus $0 + b < (d - c) + b$ by \cref{real_lt_add}.
                    $0 + b = b$ by \cref{real_add_comm,real_add_ident}.
                    Thus $b < (d - c) + b$ by \cref{real_lt_subst_left}.
                    $(d - c) + b = b + (d - c)$ by \cref{real_add_comm}.
                    Thus $b < b + (d - c)$ by \cref{real_lt_subst_right}.
                    Thus $b < e$ by \cref{real_lt_subst_right}.
                    Thus $b\mathrel{\realOrderRel}e$ by \cref{real_order_relation_intro}.
                    \begin{byCase}
                    \caseOf{$t = b$.}
                        Thus $t\mathrel{\realOrderRel}e$ by \cref{real_order_relation_intro,real_order_relation_elim}.
                        Thus $t\in I$ by \cref{intervalclosed_rel}.
                    \caseOf{$t\mathrel{\realOrderRel}b$.}
                        Thus $t < b$ by \cref{real_order_relation_elim}.
                        Thus $b < e$ by \cref{real_order_relation_elim}.
                        Thus $t\in\reals$ by \cref{intervalclosed_rel}.
                        $b\in\reals$.
                        $e\in\reals$.
                        Thus $t < e$ by \cref{real_lt_trans}.
                        Thus $t\mathrel{\realOrderRel}e$ by \cref{real_order_relation_intro}.
                        Thus $t\in I$ by \cref{intervalclosed_rel}.
                    \end{byCase}
                \caseOf{$t\in I_3$.}
                    Thus $b\mathrel{\realOrderRel}t$ or $t = b$ by \cref{intervalclosed_rel}.
                    Thus $t\mathrel{\realOrderRel}e$ or $t = e$ by \cref{intervalclosed_rel}.
                    \begin{byCase}
                    \caseOf{$t = b$.}
                        Thus $a\mathrel{\realOrderRel}b$.
                        Thus $a\mathrel{\realOrderRel}t$ by \cref{real_lt_subst_right}.
                        Thus $t\in I$ by \cref{intervalclosed_rel}.
                    \caseOf{$b\mathrel{\realOrderRel}t$.}
                        Thus $a\mathrel{\realOrderRel}b$.
                        Thus $a < b$ by \cref{real_order_relation_elim}.
                        Thus $b < t$ by \cref{real_order_relation_elim}.
                        $a\in\reals$.
                        $b\in\reals$.
                        Thus $t\in\reals$ by \cref{intervalclosed_rel}.
                        Thus $a < t$ by \cref{real_lt_trans}.
                        Thus $a\mathrel{\realOrderRel}t$ by \cref{real_order_relation_intro}.
                        Thus $t\in I$ by \cref{intervalclosed_rel}.
                    \end{byCase}
                \end{byCase}
            \end{subproof}
            Thus $I_1\union I_3\subseteq I$ by \cref{subseteq}.
        \end{subproof}
        Show that $I\subseteq I_1\union I_3$.
        \begin{subproof}
            Show that for all $t$ such that $t\in I$ we have $t\in I_1\union I_3$.
            \begin{subproof}
                Fix $t$.
                Assume $t\in I$.
                Thus $t\in\reals$ by \cref{intervalclosed_rel}.
                Thus $a\mathrel{\realOrderRel}t$ or $t = a$ by \cref{intervalclosed_rel}.
                Thus $t\mathrel{\realOrderRel}e$ or $t = e$ by \cref{intervalclosed_rel}.
                \begin{byCase}
                \caseOf{$t = b$.}
                    Thus $t\in I_1$ by \cref{intervalclosed_rel}.
                    Thus $t\in I_1\union I_3$ by \cref{union_iff}.
                \caseOf{$t\neq b$.}
                    $b\in\reals$.
                    $t\in\reals$.
                    $t\neq b$.
                    Thus $t < b$ or $b < t$ by \cref{real_lt_trichotomy}.
                    \begin{byCase}
                    \caseOf{$t < b$.}
                        Thus $t\mathrel{\realOrderRel}b$ by \cref{real_order_relation_intro}.
                        Thus $t\in I_1$ by \cref{intervalclosed_rel}.
                        Thus $t\in I_1\union I_3$ by \cref{union_iff}.
                    \caseOf{$b < t$.}
                        Thus $b\mathrel{\realOrderRel}t$ by \cref{real_order_relation_intro}.
                        Thus $t\in I_3$ by \cref{intervalclosed_rel}.
                        Thus $t\in I_1\union I_3$ by \cref{union_iff}.
                    \end{byCase}
                \end{byCase}
        \end{subproof}
        Thus $I\subseteq I_1\union I_3$ by \cref{subseteq}.
    \end{subproof}
    Thus $I = I_1\union I_3$ by \cref{subseteq_antisymmetric}.
    Show that $I_1\subseteq I$.
    \begin{subproof}
        Show that for all $x$ such that $x\in I_1$ we have $x\in I$.
        \begin{subproof}
            Fix $x$.
            Assume $x\in I_1$.
            Thus $x\in I_1\union I_3$ by \cref{union_iff}.
            Thus $x\in I$.
        \end{subproof}
        Thus $I_1\subseteq I$ by \cref{subseteq}.
    \end{subproof}
    Show that $I_3\subseteq I$.
    \begin{subproof}
        Show that for all $x$ such that $x\in I_3$ we have $x\in I$.
        \begin{subproof}
            Fix $x$.
            Assume $x\in I_3$.
            Thus $x\in I_1\union I_3$ by \cref{union_iff}.
            Thus $x\in I$.
        \end{subproof}
        Thus $I_3\subseteq I$ by \cref{subseteq}.
    \end{subproof}
    \end{subproof}
    $\realOrderRel$ is a simple order relation on $\reals$ and $\reals$ has more than one element
        by \cref{reals_linear_continuum,linear_continuum}.
    $\realOrderTopology = \basisTopology{\orderBasis{\realOrderRel}{\reals}}{\reals}$ by \cref{real_order_topology}.
    Thus $\realOrderTopology$ is the order topology on $\reals$ with respect to $\realOrderRel$ by \cref{order_topology}.
    Show that $I_1$ is closed in $I$ with respect to $T_I$.
    \begin{subproof}
        $I$ is a subspace of $\reals$ with respect to $\realOrderTopology$ by \cref{subspace_of}.
        $I_1$ is closed in $\reals$ with respect to $\realOrderTopology$ by \cref{intervalclosedrel_closed_in_order_topology}.
        Thus there exists $C$ such that $C$ is closed in $\reals$ with respect to $\realOrderTopology$ and $I_1 = C\inter I$.
        Thus $I_1$ is closed in $I$ with respect to $T_I$ by \cref{closed_in_subspace_iff,subspace_of}.
    \end{subproof}
    Show that $I_3$ is closed in $I$ with respect to $T_I$.
    \begin{subproof}
        $I$ is a subspace of $\reals$ with respect to $\realOrderTopology$ by \cref{subspace_of}.
        $I_3$ is closed in $\reals$ with respect to $\realOrderTopology$ by \cref{intervalclosedrel_closed_in_order_topology}.
        Thus there exists $C$ such that $C$ is closed in $\reals$ with respect to $\realOrderTopology$ and $I_3 = C\inter I$.
        Thus $I_3$ is closed in $I$ with respect to $T_I$ by \cref{closed_in_subspace_iff,subspace_of}.
    \end{subproof}
    Show that for all $t$ such that $t\in I_1\inter I_3$ we have $f(t) = g_1(t)$.
    \begin{subproof}
        Fix $t$.
        Assume $t\in I_1\inter I_3$.
        Thus $t\in I_1$ and $t\in I_3$ by \cref{inter}.
        Thus $t\mathrel{\realOrderRel}b$ or $t = b$ by \cref{intervalclosed_rel}.
        Thus $b\mathrel{\realOrderRel}t$ or $t = b$ by \cref{intervalclosed_rel}.
        Show that $t = b$.
        \begin{subproof}
            Suppose not.
            Then $t\neq b$.
            Thus $t\mathrel{\realOrderRel}b$ and $b\mathrel{\realOrderRel}t$ by \cref{intervalclosed_rel}.
            Thus $t < b$ and $b < t$ by \cref{real_order_relation_elim}.
            Thus $t\in\reals$ by \cref{intervalclosed_rel}.
            $b\in\reals$.
            Contradiction by \cref{real_lt_asym}.
        \end{subproof}
        Thus $f(t) = y$.
        Thus $g_1(t) = y$ by \cref{real_lt_subst_right}.
    \end{subproof}
    Let $T_{I1} = \subspaceTopology{T_I}{I_1}$.
    Let $T_{I3} = \subspaceTopology{T_I}{I_3}$.
    Show that $T_{I1} = T_1$.
    \begin{subproof}
        Show that $I_1\subseteq I$.
        \begin{subproof}
            Show that for all $x$ such that $x\in I_1$ we have $x\in I$.
            \begin{subproof}
                Fix $x$.
                Assume $x\in I_1$.
                Thus $x\in I_1\union I_3$ by \cref{union_iff}.
                Thus $x\in I$.
            \end{subproof}
            Thus $I_1\subseteq I$ by \cref{subseteq}.
        \end{subproof}
        $I\subseteq \reals$.
        Thus $T_{I1} = \subspaceTopology{\realOrderTopology}{I_1}$
            by \cref{subspace_subspace_topology,subspace_of,subseteq}.
        Thus $T_{I1} = T_1$ by \cref{subspace_topology}.
    \end{subproof}
    Show that $T_{I3} = T_3$.
    \begin{subproof}
        Show that $I_3\subseteq I$.
        \begin{subproof}
            Show that for all $x$ such that $x\in I_3$ we have $x\in I$.
            \begin{subproof}
                Fix $x$.
                Assume $x\in I_3$.
                Thus $x\in I_1\union I_3$ by \cref{union_iff}.
                Thus $x\in I$.
            \end{subproof}
            Thus $I_3\subseteq I$ by \cref{subseteq}.
        \end{subproof}
        $I\subseteq \reals$.
        Thus $T_{I3} = \subspaceTopology{\realOrderTopology}{I_3}$
            by \cref{subspace_subspace_topology,subspace_of,subseteq}.
        Thus $T_{I3} = T_3$ by \cref{subspace_topology}.
    \end{subproof}
    Thus $f$ is continuous from $I_1$ to $X$ with respect to $T_{I1}$ and $T$.
    Thus $g_1$ is continuous from $I_3$ to $X$ with respect to $T_{I3}$ and $T$.
    Let $h = f\union g_1$.
    Thus $h$ is continuous from $I$ to $X$ with respect to $T_I$ and $T$ by \cref{pasting_lemma}.
    Show that $h(a) = x$ and $h(e) = z$.
    \begin{subproof}
        $a\in I_1$ by \cref{intervalclosed_rel,real_order_relation_elim}.
        $e\in I_3$ by \cref{intervalclosed_rel,real_order_relation_elim}.
        $f\in\funs{I_1}{X}$ by \cref{continuous}.
        $g_1\in\funs{I_3}{X}$ by \cref{continuous}.
        $(a, f(a))\in f$ by \cref{funs_elim,function_apply_iff}.
        $(e, g_1(e))\in g_1$ by \cref{funs_elim,function_apply_iff}.
        $h\in\funs{I}{X}$ by \cref{continuous}.
        Thus $h$ is a function by \cref{funs_is_function}.
        Thus $h(a) = f(a)$ and $h(e) = g_1(e)$ by \cref{function_apply_iff,union_iff}.
        Thus $h(a) = x$ and $h(e) = z$.
    \end{subproof}
    Show that $a\mathrel{\realOrderRel}e$.
    \begin{subproof}
        $a\mathrel{\realOrderRel}b$.
        Thus $a < b$ by \cref{real_order_relation_elim}.
        Show that $b < e$.
        \begin{subproof}
            $c\mathrel{\realOrderRel} d$.
            Thus $c < d$ by \cref{real_order_relation_elim}.
            $c\in\reals$.
            $d\in\reals$.
            Thus $c + \neg{c} < d + \neg{c}$ by \cref{real_lt_add}.
            $c + \neg{c} = 0$ by \cref{real_add_inverse}.
            Thus $0 < d + \neg{c}$ by \cref{real_lt_subst_left}.
            $d + \neg{c} = d - c$ by \cref{real_sub}.
            Thus $0 < d - c$ by \cref{real_lt_subst_right}.
            $b\in\reals$.
            $d - c\in\reals$.
            $0\in\reals$ by \cref{real_add_ident}.
            Thus $0 + b < (d - c) + b$ by \cref{real_lt_add}.
            $0 + b = b$ by \cref{real_add_comm,real_add_ident}.
            Thus $b < (d - c) + b$ by \cref{real_lt_subst_left}.
            $(d - c) + b = b + (d - c)$ by \cref{real_add_comm}.
            Thus $b < b + (d - c)$ by \cref{real_lt_subst_right}.
            Thus $b < e$ by \cref{real_lt_subst_right}.
        \end{subproof}
        $a\in\reals$.
        $b\in\reals$.
        $e\in\reals$.
        Thus $a < e$ by \cref{real_lt_trans}.
        Thus $a\mathrel{\realOrderRel}e$ by \cref{real_order_relation_intro}.
    \end{subproof}
    Thus $h$ is continuous from $\intervalclosedrel{\realOrderRel}{\reals}{a}{e}$ to $X$ with respect to
        $\subspaceTopology{\realOrderTopology}{\intervalclosedrel{\realOrderRel}{\reals}{a}{e}}$ and $T$.
    Thus $h$ is a path in $X$ from $x$ to $z$ with respect to $T$ by \cref{path_in_space}.
    Thus there exists $h$ such that $h$ is a path in $X$ from $x$ to $z$ with respect to $T$.
\end{proof}

% from §25 helper lemma 18: path relation is a binary relation
\begin{lemma}\label{path_relation_binary}
    Assume $T$ is a topology on $X$.
    Then $\pathRelation{T}{X}$ is a binary relation on $X$.
\end{lemma}
\begin{proof}
    Show that for all $w$ such that $w\in \pathRelation{T}{X}$ we have $w\in X\times X$.
    \begin{subproof}
        Fix $w$.
        Assume $w\in \pathRelation{T}{X}$.
        Thus $w\in X\times X$ and there exists $f$ such that
            $f$ is a path in $X$ from $\fst{w}$ to $\snd{w}$ with respect to $T$
            by \cref{path_relation}.
        Thus $w\in X\times X$.
    \end{subproof}
    Thus $\pathRelation{T}{X}\subseteq X\times X$ by \cref{subseteq}.
    Thus $\pathRelation{T}{X}$ is a binary relation on $X$.
\end{proof}

% from §25 helper lemma 19: path relation is reflexive
\begin{lemma}\label{path_relation_reflexive}
    Assume $T$ is a topology on $X$.
    Then $\pathRelation{T}{X}$ is reflexive on $X$.
\end{lemma}
\begin{proof}
    Show that for all $x\in X$ we have $x\mathrel{\pathRelation{T}{X}}x$.
    \begin{subproof}
        Fix $x\in X$.
        $0\in\reals$ by \cref{real_add_ident}.
        $1\in\reals$ by \cref{real_mul_ident}.
        $0 < 1$ by \cref{real_zero_lt_one}.
        Thus $0\mathrel{\realOrderRel}1$ by \cref{real_order_relation_intro}.
        Let $I = \intervalclosedrel{\realOrderRel}{\reals}{0}{1}$.
        Let $T_I = \subspaceTopology{\realOrderTopology}{I}$.
        Show that $I\subseteq\reals$.
        \begin{subproof}
            Show that for all $t$ such that $t\in I$ we have $t\in\reals$.
            \begin{subproof}
                Fix $t$.
                Assume $t\in I$.
                Thus $t\in\reals$ by \cref{intervalclosed_rel}.
            \end{subproof}
            Thus $I\subseteq\reals$ by \cref{subseteq}.
        \end{subproof}
        $\standardBasisR$ is a basis on $\reals$ by \cref{standard_basis_is_basis}.
        $\standardTopologyR = \basisTopology{\standardBasisR}{\reals}$ by \cref{standard_topology_reals}.
        Thus $\standardTopologyR$ is a topology on $\reals$ by \cref{basis_topology_is_topology}.
        $\realOrderTopology = \standardTopologyR$ by \cref{real_order_topology_eq_standard}.
        Thus $\realOrderTopology$ is a topology on $\reals$.
        Thus $I$ is a subspace of $\reals$ with respect to $\realOrderTopology$ by \cref{subspace_of}.
        Thus $T_I$ is a topology on $I$ by \cref{subspace_topology_is_topology,subspace_of}.
        Let $f = \{(t,x) \mid t\in I\}$.
        $f$ is a function from $I$ to $X$ by \cref{constant_map_function}.
        Show that for all $t\in I$ we have $f(t) = x$.
        \begin{subproof}
            Fix $t$.
            Assume $t\in I$.
            Thus $(t,x)\in f$.
            $f$ is a function.
            Thus $f(t) = x$ by \cref{function_apply_iff}.
        \end{subproof}
        Thus $f$ is continuous from $I$ to $X$ with respect to $T_I$ and $T$ by \cref{constant_continuous}.
        Show that $0\in I$ and $1\in I$.
        \begin{subproof}
            $0\in I$ by \cref{intervalclosed_rel,real_order_relation_elim}.
            $1\in I$ by \cref{intervalclosed_rel,real_order_relation_elim}.
        \end{subproof}
        Thus $f(0) = x$ and $f(1) = x$.
        Thus $f$ is a path in $X$ from $x$ to $x$ with respect to $T$ by \cref{path_in_space}.
        Thus there exists $f$ such that $f$ is a path in $X$ from $x$ to $x$ with respect to $T$.
        Thus $(x,x)\in X\times X$ by \cref{times_tuple_intro}.
        Thus $\fst{(x,x)} = x$ and $\snd{(x,x)} = x$ by \cref{fst_eq,snd_eq}.
        Thus there exists $f$ such that
            $f$ is a path in $X$ from $\fst{(x,x)}$ to $\snd{(x,x)}$ with respect to $T$.
        Thus $(x,x)\in \pathRelation{T}{X}$ by \cref{path_relation}.
        Thus $x\mathrel{\pathRelation{T}{X}}x$.
    \end{subproof}
    Thus $\pathRelation{T}{X}$ is reflexive on $X$ by \cref{reflexive_on}.
\end{proof}

% from §25 helper lemma 20: path relation is symmetric
\begin{lemma}\label{path_relation_symmetric}
    Assume $T$ is a topology on $X$.
    Then $\pathRelation{T}{X}$ is symmetric.
\end{lemma}
\begin{proof}
    Show that for all $x,y$ such that $x\mathrel{\pathRelation{T}{X}}y$ we have
        $y\mathrel{\pathRelation{T}{X}}x$.
    \begin{subproof}
        Fix $x$.
        Fix $y$.
        Assume $x\mathrel{\pathRelation{T}{X}}y$.
        Thus $(x,y)\in X\times X$ and there exists $f$ such that
            $f$ is a path in $X$ from $\fst{(x,y)}$ to $\snd{(x,y)}$ with respect to $T$
            by \cref{path_relation}.
        Thus $x\in X$ and $y\in X$ by \cref{times_tuple_elim}.
        Thus there exists $f$ such that $f$ is a path in $X$ from $x$ to $y$ with respect to $T$
            by \cref{fst_eq,snd_eq}.
        Take $f$ such that $f$ is a path in $X$ from $x$ to $y$ with respect to $T$.
        Thus there exists $g$ such that $g$ is a path in $X$ from $y$ to $x$ with respect to $T$
            by \cref{path_reverse}.
        Take $g$ such that $g$ is a path in $X$ from $y$ to $x$ with respect to $T$.
        Thus $(y,x)\in X\times X$ by \cref{times_tuple_intro}.
        Thus $\fst{(y,x)} = y$ and $\snd{(y,x)} = x$ by \cref{fst_eq,snd_eq}.
        Thus there exists $g$ such that
            $g$ is a path in $X$ from $\fst{(y,x)}$ to $\snd{(y,x)}$ with respect to $T$.
        Thus $(y,x)\in \pathRelation{T}{X}$ by \cref{path_relation}.
        Thus $y\mathrel{\pathRelation{T}{X}}x$.
    \end{subproof}
    Thus $\pathRelation{T}{X}$ is symmetric by \cref{symmetric}.
\end{proof}

% from §25 helper lemma 21: path relation is transitive
\begin{lemma}\label{path_relation_transitive}
    Assume $T$ is a topology on $X$.
    Then $\pathRelation{T}{X}$ is transitive.
\end{lemma}
\begin{proof}
    Show that for all $x,y,z$ such that
        $x\mathrel{\pathRelation{T}{X}}y$ and $y\mathrel{\pathRelation{T}{X}}z$
        we have $x\mathrel{\pathRelation{T}{X}}z$.
    \begin{subproof}
        Fix $x$.
        Fix $y$.
        Fix $z$.
        Assume $x\mathrel{\pathRelation{T}{X}}y$ and $y\mathrel{\pathRelation{T}{X}}z$.
        Thus $(x,y)\in X\times X$ and there exists $f$ such that
            $f$ is a path in $X$ from $\fst{(x,y)}$ to $\snd{(x,y)}$ with respect to $T$
            by \cref{path_relation}.
        Thus $(y,z)\in X\times X$ and there exists $g$ such that
            $g$ is a path in $X$ from $\fst{(y,z)}$ to $\snd{(y,z)}$ with respect to $T$
            by \cref{path_relation}.
        Thus $x\in X$ and $y\in X$ by \cref{times_tuple_elim}.
        Thus $z\in X$ by \cref{times_tuple_elim}.
        Thus there exists $f$ such that $f$ is a path in $X$ from $x$ to $y$ with respect to $T$
            by \cref{fst_eq,snd_eq}.
        Thus there exists $g$ such that $g$ is a path in $X$ from $y$ to $z$ with respect to $T$
            by \cref{fst_eq,snd_eq}.
        Take $f$ such that $f$ is a path in $X$ from $x$ to $y$ with respect to $T$.
        Take $g$ such that $g$ is a path in $X$ from $y$ to $z$ with respect to $T$.
        Thus there exists $h$ such that $h$ is a path in $X$ from $x$ to $z$ with respect to $T$
            by \cref{path_concatenation}.
        Take $h$ such that $h$ is a path in $X$ from $x$ to $z$ with respect to $T$.
        Thus $(x,z)\in X\times X$ by \cref{times_tuple_intro}.
        Thus $\fst{(x,z)} = x$ and $\snd{(x,z)} = z$ by \cref{fst_eq,snd_eq}.
        Thus there exists $h$ such that
            $h$ is a path in $X$ from $\fst{(x,z)}$ to $\snd{(x,z)}$ with respect to $T$.
        Thus $(x,z)\in \pathRelation{T}{X}$ by \cref{path_relation}.
        Thus $x\mathrel{\pathRelation{T}{X}}z$.
    \end{subproof}
    Thus $\pathRelation{T}{X}$ is transitive by \cref{transitive}.
\end{proof}

% from §25 helper lemma 22: path relation is an equivalence
\begin{lemma}\label{path_relation_equivalence}
    Assume $T$ is a topology on $X$.
    Then $\pathRelation{T}{X}$ is an equivalence on $X$.
\end{lemma}
\begin{proof}
    $\pathRelation{T}{X}$ is a binary relation on $X$ by \cref{path_relation_binary}.
    $\pathRelation{T}{X}$ is reflexive on $X$ by \cref{path_relation_reflexive}.
    $\pathRelation{T}{X}$ is transitive by \cref{path_relation_transitive}.
    $\pathRelation{T}{X}$ is symmetric by \cref{path_relation_symmetric}.
    Thus $\pathRelation{T}{X}$ is an equivalence on $X$.
\end{proof}

% from §25 helper lemma 23: union of path-connected subspaces with a common point
\begin{lemma}\label{union_path_connected_common_point}
    Assume $M\neq\emptyset$.
    Assume $T$ is a topology on $X$.
    Assume for all $A$ such that $A\in M$ we have $A\subseteq X$.
    Assume for all $A$ such that $A\in M$ we have
        $A$ is path connected with respect to $\subspaceTopology{T}{A}$.
    Assume $\inters{M}\neq\emptyset$.
    Then $\unions{M}$ is path connected with respect to $\subspaceTopology{T}{\unions{M}}$.
\end{lemma}
\begin{proof}
    Let $U = \unions{M}$.
    Let $T_U = \subspaceTopology{T}{U}$.
    Show that $U\subseteq X$.
    \begin{subproof}
        Show that for all $x$ such that $x\in U$ we have $x\in X$.
        \begin{subproof}
            Fix $x$.
            Assume $x\in U$.
            Take $A\in M$ such that $x\in A$ by \cref{unions_iff}.
            Thus $A\subseteq X$ by \cref{subseteq}.
            Thus $x\in X$ by \cref{subseteq}.
        \end{subproof}
        Thus $U\subseteq X$ by \cref{subseteq}.
    \end{subproof}
    Thus $U$ is a subspace of $X$ with respect to $T$ by \cref{subspace_of}.
    Thus $T_U$ is a topology on $U$ by \cref{subspace_topology_is_topology,subspace_of}.
    Show that for all $x,y\in U$ there exists $h$ such that
        $h$ is a path in $U$ from $x$ to $y$ with respect to $T_U$.
    \begin{subproof}
        Fix $x$.
        Fix $y$.
        Assume $x\in U$ and $y\in U$.
        Take $A\in M$ such that $x\in A$ by \cref{unions_iff}.
        Take $B\in M$ such that $y\in B$ by \cref{unions_iff}.
        Take $p$ such that $p\in\inters{M}$ by \cref{nonempty_inhabited}.
        Thus $p\in A$ and $p\in B$ by \cref{inters_iff_forall}.
        Let $T_A = \subspaceTopology{T}{A}$.
        Let $T_B = \subspaceTopology{T}{B}$.
        $A$ is path connected with respect to $T_A$.
        $B$ is path connected with respect to $T_B$.
        Thus there exists $f$ such that $f$ is a path in $A$ from $x$ to $p$ with respect to $T_A$
            by \cref{path_connected}.
        Thus there exists $g$ such that $g$ is a path in $B$ from $p$ to $y$ with respect to $T_B$
            by \cref{path_connected}.
        Take $f$ such that $f$ is a path in $A$ from $x$ to $p$ with respect to $T_A$.
        Take $g$ such that $g$ is a path in $B$ from $p$ to $y$ with respect to $T_B$.
        Show that $A\subseteq U$.
        \begin{subproof}
            Show that for all $u$ such that $u\in A$ we have $u\in U$.
            \begin{subproof}
                Fix $u$.
                Assume $u\in A$.
                Thus $u\in U$ by \cref{unions_iff}.
            \end{subproof}
            Thus $A\subseteq U$ by \cref{subseteq}.
        \end{subproof}
        Show that $B\subseteq U$.
        \begin{subproof}
            Show that for all $u$ such that $u\in B$ we have $u\in U$.
            \begin{subproof}
                Fix $u$.
                Assume $u\in B$.
                Thus $u\in U$ by \cref{unions_iff}.
            \end{subproof}
            Thus $B\subseteq U$ by \cref{subseteq}.
        \end{subproof}
        Thus $A$ is a subspace of $U$ with respect to $T_U$ by \cref{subspace_of}.
        Thus $B$ is a subspace of $U$ with respect to $T_U$ by \cref{subspace_of}.
        Let $T_{UA} = \subspaceTopology{T_U}{A}$.
        Let $T_{UB} = \subspaceTopology{T_U}{B}$.
        Thus $T_{UA} = T_A$ by \cref{subspace_subspace_topology,subspace_of,subseteq}.
        Thus $T_{UB} = T_B$ by \cref{subspace_subspace_topology,subspace_of,subseteq}.
        Let $j_A = \identity{A}$.
        Let $j_B = \identity{B}$.
        Thus $j_A$ is continuous from $A$ to $U$ with respect to $T_{UA}$ and $T_U$
            by \cref{inclusion_continuous}.
        Thus $j_B$ is continuous from $B$ to $U$ with respect to $T_{UB}$ and $T_U$
            by \cref{inclusion_continuous}.
        Thus $j_A$ is continuous from $A$ to $U$ with respect to $T_A$ and $T_U$.
        Thus $j_B$ is continuous from $B$ to $U$ with respect to $T_B$ and $T_U$.
        Let $f_1 = j_A\circ f$.
        Let $g_1 = j_B\circ g$.
        Show that $f_1$ is a path in $U$ from $x$ to $p$ with respect to $T_U$.
        \begin{subproof}
            Take $a,b\in\reals$ such that
                $a\mathrel{\realOrderRel} b$ and
                $f$ is continuous from $\intervalclosedrel{\realOrderRel}{\reals}{a}{b}$ to $A$ with respect to
                    $\subspaceTopology{\realOrderTopology}{\intervalclosedrel{\realOrderRel}{\reals}{a}{b}}$ and $T_A$ and
                $f(a) = x$ and $f(b) = p$
                by \cref{path_in_space}.
            Let $I = \intervalclosedrel{\realOrderRel}{\reals}{a}{b}$.
            $f\in\funs{I}{A}$ by \cref{continuous}.
            Thus $f$ is a function by \cref{funs_is_function}.
            $j_A\in\funs{A}{U}$ by \cref{continuous}.
            Thus $j_A$ is a function by \cref{funs_is_function}.
            Thus $\dom{j_A} = A$ by \cref{funs}.
            Thus $\ran{f}\subseteq A$ by \cref{funs_ran}.
            Thus $\ran{f}\subseteq \dom{j_A}$ by \cref{subseteq}.
            Thus $j_A$ is composable with $f$.
            Thus $f_1$ is continuous from $\intervalclosedrel{\realOrderRel}{\reals}{a}{b}$ to $U$ with respect to
                $\subspaceTopology{\realOrderTopology}{\intervalclosedrel{\realOrderRel}{\reals}{a}{b}}$ and $T_U$
                by \cref{continuous_composite}.
            Thus $b\mathrel{\realOrderRel} b$ or $b = b$.
            Thus $a\mathrel{\realOrderRel} b$ or $a = b$.
            Thus $a\in I$ and $b\in I$ by \cref{intervalclosed_rel}.
            Thus $\dom{f} = I$ by \cref{funs}.
            Thus $a\in\dom{f}$ and $b\in\dom{f}$ by \cref{subseteq}.
            Thus $(a,f(a))\in f$ and $(b,f(b))\in f$ by \cref{function_apply_elim}.
            Thus $f(a)\in\ran{f}$ and $f(b)\in\ran{f}$ by \cref{ran_intro}.
            Thus $\ran{f}\subseteq A$ by \cref{funs_ran}.
            Thus $f(a)\in A$ and $f(b)\in A$ by \cref{subseteq}.
            Thus $f_1(a) = j_A(f(a))$ and $f_1(b) = j_A(f(b))$ by \cref{circ_apply}.
            Thus $j_A(f(a)) = f(a)$ and $j_A(f(b)) = f(b)$ by \cref{id_apply}.
            Thus $f_1(a) = x$ and $f_1(b) = p$.
            Thus $f_1$ is a path in $U$ from $x$ to $p$ with respect to $T_U$ by \cref{path_in_space}.
        \end{subproof}
        Show that $g_1$ is a path in $U$ from $p$ to $y$ with respect to $T_U$.
        \begin{subproof}
            Take $c,d\in\reals$ such that
                $c\mathrel{\realOrderRel} d$ and
                $g$ is continuous from $\intervalclosedrel{\realOrderRel}{\reals}{c}{d}$ to $B$ with respect to
                    $\subspaceTopology{\realOrderTopology}{\intervalclosedrel{\realOrderRel}{\reals}{c}{d}}$ and $T_B$ and
                $g(c) = p$ and $g(d) = y$
                by \cref{path_in_space}.
            Let $J = \intervalclosedrel{\realOrderRel}{\reals}{c}{d}$.
            $g\in\funs{J}{B}$ by \cref{continuous}.
            Thus $g$ is a function by \cref{funs_is_function}.
            $j_B\in\funs{B}{U}$ by \cref{continuous}.
            Thus $j_B$ is a function by \cref{funs_is_function}.
            Thus $\dom{j_B} = B$ by \cref{funs}.
            Thus $\ran{g}\subseteq B$ by \cref{funs_ran}.
            Thus $\ran{g}\subseteq \dom{j_B}$ by \cref{subseteq}.
            Thus $j_B$ is composable with $g$.
            Thus $g_1$ is continuous from $\intervalclosedrel{\realOrderRel}{\reals}{c}{d}$ to $U$ with respect to
                $\subspaceTopology{\realOrderTopology}{\intervalclosedrel{\realOrderRel}{\reals}{c}{d}}$ and $T_U$
                by \cref{continuous_composite}.
            Thus $d\mathrel{\realOrderRel} d$ or $d = d$.
            Thus $c\mathrel{\realOrderRel} d$ or $c = d$.
            Thus $c\in J$ and $d\in J$ by \cref{intervalclosed_rel}.
            Thus $\dom{g} = J$ by \cref{funs}.
            Thus $c\in\dom{g}$ and $d\in\dom{g}$ by \cref{subseteq}.
            Thus $(c,g(c))\in g$ and $(d,g(d))\in g$ by \cref{function_apply_elim}.
            Thus $g(c)\in\ran{g}$ and $g(d)\in\ran{g}$ by \cref{ran_intro}.
            Thus $\ran{g}\subseteq B$ by \cref{funs_ran}.
            Thus $g(c)\in B$ and $g(d)\in B$ by \cref{subseteq}.
            Thus $g_1(c) = j_B(g(c))$ and $g_1(d) = j_B(g(d))$ by \cref{circ_apply}.
            Thus $j_B(g(c)) = g(c)$ and $j_B(g(d)) = g(d)$ by \cref{id_apply}.
            Thus $g_1(c) = p$ and $g_1(d) = y$.
            Thus $g_1$ is a path in $U$ from $p$ to $y$ with respect to $T_U$ by \cref{path_in_space}.
        \end{subproof}
        Thus there exists $h$ such that $h$ is a path in $U$ from $x$ to $y$ with respect to $T_U$
            by \cref{path_concatenation}.
    \end{subproof}
    Thus $U$ is path connected with respect to $T_U$ by \cref{path_connected}.
\end{proof}


% from §25 helper lemma 24: initial segment of a path
\begin{lemma}\label{path_initial_segment}
    Assume $T$ is a topology on $X$.
    Suppose $f$ is a path in $X$ from $x$ to $y$ with respect to $T$.
    Suppose $t\in\dom{f}$.
    Then there exists $g$ such that $g$ is a path in $X$ from $x$ to $f(t)$ with respect to $T$.
\end{lemma}
\begin{proof}
    Take $a,b\in\reals$ such that
        $a\mathrel{\realOrderRel} b$ and
        $f$ is continuous from $\intervalclosedrel{\realOrderRel}{\reals}{a}{b}$ to $X$ with respect to
            $\subspaceTopology{\realOrderTopology}{\intervalclosedrel{\realOrderRel}{\reals}{a}{b}}$ and $T$ and
        $f(a) = x$ and $f(b) = y$
        by \cref{path_in_space}.
    Let $I = \intervalclosedrel{\realOrderRel}{\reals}{a}{b}$.
    Let $T_I = \subspaceTopology{\realOrderTopology}{I}$.
    $f\in\funs{I}{X}$ by \cref{continuous}.
    Thus $\dom{f} = I$ by \cref{funs}.
    Thus $t\in I$ by \cref{subseteq}.
    Show that $a\in I$.
    \begin{subproof}
        $a\in\reals$.
        $a = a$.
        Thus $a\mathrel{\realOrderRel} a$ or $a = a$.
        Thus $a\mathrel{\realOrderRel} b$ or $a = b$.
        Thus $a\in I$ by \cref{intervalclosed_rel}.
    \end{subproof}
    Thus $f$ is a function by \cref{funs_is_function}.
    Thus $\ran{f}\subseteq X$ by \cref{funs_ran}.
    Thus $a\in\dom{f}$ by \cref{funs,subseteq}.
    Thus $(a,f(a))\in f$ by \cref{function_apply_elim}.
    Thus $f(a)\in\ran{f}$ by \cref{ran_intro}.
    Thus $x\in X$ by \cref{subseteq}.
    \begin{byCase}
    \caseOf{$t = a$.}
        Thus $f(t) = x$.
        $\pathRelation{T}{X}$ is reflexive on $X$ by \cref{path_relation_reflexive}.
        Thus $x\mathrel{\pathRelation{T}{X}}x$ by \cref{reflexive_on}.
        Thus $(x,x)\in X\times X$ and there exists $g$ such that
            $g$ is a path in $X$ from $\fst{(x,x)}$ to $\snd{(x,x)}$ with respect to $T$
            by \cref{path_relation}.
        Thus there exists $g$ such that $g$ is a path in $X$ from $x$ to $f(t)$ with respect to $T$
            by \cref{fst_eq,snd_eq}.
    \caseOf{$t\neq a$.}
        Thus $a\mathrel{\realOrderRel} t$ or $t = a$ by \cref{intervalclosed_rel}.
        Thus $a\mathrel{\realOrderRel} t$.
        Thus $t\in\reals$ by \cref{real_order_relation_elim}.
        Let $J = \intervalclosedrel{\realOrderRel}{\reals}{a}{t}$.
        $\realOrderRel$ is a simple order relation on $\reals$ by \cref{real_order_relation_simple}.
        Thus $\realOrderRel$ is transitive by \cref{simple_order_relation}.
        Thus $I$ is convex in $\reals$ with respect to $\realOrderRel$ by \cref{intervalclosedrel_convex}.
        Show that $J\subseteq I$.
        \begin{subproof}
            Show that for all $s$ such that $s\in J$ we have $s\in I$.
            \begin{subproof}
                Fix $s$.
                Assume $s\in J$.
                Thus $s\in\reals$ by \cref{intervalclosed_rel}.
                Thus $a\mathrel{\realOrderRel} s$ or $s = a$ by \cref{intervalclosed_rel}.
                Thus $s\mathrel{\realOrderRel} t$ or $s = t$ by \cref{intervalclosed_rel}.
                \begin{byCase}
                \caseOf{$s = a$.}
                    Thus $s\in I$.
                \caseOf{$s\neq a$.}
                    Thus $a\mathrel{\realOrderRel} s$ by \cref{intervalclosed_rel}.
                    \begin{byCase}
                    \caseOf{$s = t$.}
                        Thus $s\in I$.
                    \caseOf{$s\neq t$.}
                        Thus $s\mathrel{\realOrderRel} t$ by \cref{intervalclosed_rel}.
                        Thus $s\in I$ by \cref{convex_in}.
                    \end{byCase}
                \end{byCase}
            \end{subproof}
            Thus $J\subseteq I$ by \cref{subseteq}.
        \end{subproof}
        Thus $T_I$ is a topology on $I$ by \cref{continuous}.
        Thus $J$ is a subspace of $I$ with respect to $T_I$ by \cref{subspace_of}.
        Let $T_J = \subspaceTopology{T_I}{J}$.
        Let $g = \restrl{f}{J}$.
        Thus $g$ is continuous from $J$ to $X$ with respect to $T_J$ and $T$
            by \cref{continuous_restrict_domain}.
        Show that $I\subseteq\reals$.
        \begin{subproof}
            Show that for all $u$ such that $u\in I$ we have $u\in\reals$.
            \begin{subproof}
                Fix $u$.
                Assume $u\in I$.
                Thus $u\in\reals$ by \cref{intervalclosed_rel}.
            \end{subproof}
            Thus $I\subseteq\reals$ by \cref{subseteq}.
        \end{subproof}
        $\standardBasisR$ is a basis on $\reals$ by \cref{standard_basis_is_basis}.
        $\standardTopologyR = \basisTopology{\standardBasisR}{\reals}$ by \cref{standard_topology_reals}.
        Thus $\standardTopologyR$ is a topology on $\reals$ by \cref{basis_topology_is_topology}.
        $\realOrderTopology = \standardTopologyR$ by \cref{real_order_topology_eq_standard}.
        Thus $\realOrderTopology$ is a topology on $\reals$.
        Thus $T_J = \subspaceTopology{\realOrderTopology}{J}$ by \cref{subspace_subspace_topology,subseteq}.
        Thus $g$ is continuous from $J$ to $X$ with respect to $\subspaceTopology{\realOrderTopology}{J}$ and $T$.
        $f$ is a function by \cref{funs_is_function}.
        Thus $a\in J$ and $t\in J$ by \cref{intervalclosed_rel}.
        Thus $g(a) = f(a)$ and $g(t) = f(t)$ by \cref{restrl_of_function_apply}.
        Thus $g(a) = x$ and $g(t) = f(t)$.
        Thus $g$ is a path in $X$ from $x$ to $f(t)$ with respect to $T$ by \cref{path_in_space}.
        Thus there exists $g$ such that $g$ is a path in $X$ from $x$ to $f(t)$ with respect to $T$.
    \end{byCase}
\end{proof}

% from §25 Item 9: path components are path connected
\begin{theorem}\label{path_component_path_connected}
    Assume $T$ is a topology on $X$.
    Suppose $C$ is a path component of $X$ with respect to $T$.
    Then $C$ is path connected with respect to $\subspaceTopology{T}{C}$.
\end{theorem}
\begin{proof}
    Take $x\in X$ such that $C = \equivalenceClass{\pathRelation{T}{X}}{x}$ by \cref{path_component}.
    Show that $C\subseteq X$.
    \begin{subproof}
        Show that for all $u$ such that $u\in C$ we have $u\in X$.
        \begin{subproof}
            Fix $u$.
            Assume $u\in C$.
            Thus $u\mathrel{\pathRelation{T}{X}}x$ by \cref{downward_closure_iff}.
            Thus $(u,x)\in X\times X$ and there exists $f$ such that
                $f$ is a path in $X$ from $\fst{(u,x)}$ to $\snd{(u,x)}$ with respect to $T$
                by \cref{path_relation}.
            Thus $u\in X$ by \cref{times_tuple_elim}.
        \end{subproof}
        Thus $C\subseteq X$ by \cref{subseteq}.
    \end{subproof}
    Thus $C$ is a subspace of $X$ with respect to $T$ by \cref{subspace_of}.
    Let $T_C = \subspaceTopology{T}{C}$.
    Thus $T_C$ is a topology on $C$ by \cref{subspace_topology_is_topology,subspace_of}.
    Show that for all $y,z\in C$ there exists $h$ such that
        $h$ is a path in $C$ from $y$ to $z$ with respect to $T_C$.
    \begin{subproof}
        Fix $y$.
        Fix $z$.
        Assume $y\in C$ and $z\in C$.
        Thus $y\mathrel{\pathRelation{T}{X}}x$ by \cref{downward_closure_iff}.
        Thus $z\mathrel{\pathRelation{T}{X}}x$ by \cref{downward_closure_iff}.
        $\pathRelation{T}{X}$ is symmetric by \cref{path_relation_symmetric}.
        Thus $x\mathrel{\pathRelation{T}{X}}y$ by \cref{symmetric}.
        Thus $x\mathrel{\pathRelation{T}{X}}z$ by \cref{symmetric}.
        Thus $(x,y)\in X\times X$ and there exists $f$ such that
            $f$ is a path in $X$ from $\fst{(x,y)}$ to $\snd{(x,y)}$ with respect to $T$
            by \cref{path_relation}.
        Thus $(x,z)\in X\times X$ and there exists $g$ such that
            $g$ is a path in $X$ from $\fst{(x,z)}$ to $\snd{(x,z)}$ with respect to $T$
            by \cref{path_relation}.
        Thus there exists $f$ such that $f$ is a path in $X$ from $x$ to $y$ with respect to $T$
            by \cref{fst_eq,snd_eq}.
        Thus there exists $g$ such that $g$ is a path in $X$ from $x$ to $z$ with respect to $T$
            by \cref{fst_eq,snd_eq}.
        Take $f$ such that $f$ is a path in $X$ from $x$ to $y$ with respect to $T$.
        Take $g$ such that $g$ is a path in $X$ from $x$ to $z$ with respect to $T$.
        Take $a,b\in\reals$ such that
            $a\mathrel{\realOrderRel} b$ and
            $f$ is continuous from $\intervalclosedrel{\realOrderRel}{\reals}{a}{b}$ to $X$ with respect to
                $\subspaceTopology{\realOrderTopology}{\intervalclosedrel{\realOrderRel}{\reals}{a}{b}}$ and $T$ and
            $f(a) = x$ and $f(b) = y$
            by \cref{path_in_space}.
        Let $I = \intervalclosedrel{\realOrderRel}{\reals}{a}{b}$.
        Let $T_I = \subspaceTopology{\realOrderTopology}{I}$.
        Take $c,d\in\reals$ such that
            $c\mathrel{\realOrderRel} d$ and
            $g$ is continuous from $\intervalclosedrel{\realOrderRel}{\reals}{c}{d}$ to $X$ with respect to
                $\subspaceTopology{\realOrderTopology}{\intervalclosedrel{\realOrderRel}{\reals}{c}{d}}$ and $T$ and
            $g(c) = x$ and $g(d) = z$
            by \cref{path_in_space}.
        Let $J = \intervalclosedrel{\realOrderRel}{\reals}{c}{d}$.
        Let $T_J = \subspaceTopology{\realOrderTopology}{J}$.
        $f\in\funs{I}{X}$ by \cref{continuous}.
        $g\in\funs{J}{X}$ by \cref{continuous}.
        Thus $f$ is a function by \cref{funs_is_function}.
        Thus $g$ is a function by \cref{funs_is_function}.
        Show that $\ran{f}\subseteq C$.
        \begin{subproof}
            Show that for all $u$ such that $u\in\ran{f}$ we have $u\in C$.
            \begin{subproof}
                Fix $u$.
                Assume $u\in\ran{f}$.
                Take $t\in\dom{f}$ such that $f(t) = u$ by \cref{fun_ran_iff}.
                Thus $(t,f(t))\in f$ by \cref{function_apply_elim}.
                Thus $f(t)\in\ran{f}$ by \cref{ran_intro}.
                Thus $\ran{f}\subseteq X$ by \cref{funs_ran}.
                Thus $f(t)\in X$ by \cref{subseteq}.
                Thus $(x,f(t))\in X\times X$ by \cref{times_tuple_intro}.
                Thus there exists $h$ such that $h$ is a path in $X$ from $x$ to $f(t)$ with respect to $T$
                    by \cref{path_initial_segment}.
                Thus $\fst{(x,f(t))} = x$ and $\snd{(x,f(t))} = f(t)$ by \cref{fst_eq,snd_eq}.
                Thus there exists $h$ such that
                    $h$ is a path in $X$ from $\fst{(x,f(t))}$ to $\snd{(x,f(t))}$ with respect to $T$.
                Thus $(x,f(t))\in \pathRelation{T}{X}$ by \cref{path_relation}.
                Thus $x\mathrel{\pathRelation{T}{X}}f(t)$.
                Thus $f(t)\mathrel{\pathRelation{T}{X}}x$ by \cref{path_relation_symmetric,symmetric}.
                Thus $u\mathrel{\pathRelation{T}{X}}x$.
                Thus $u\in C$ by \cref{downward_closure_iff}.
            \end{subproof}
            Thus $\ran{f}\subseteq C$ by \cref{subseteq}.
        \end{subproof}
        Show that $\ran{g}\subseteq C$.
        \begin{subproof}
            Show that for all $u$ such that $u\in\ran{g}$ we have $u\in C$.
            \begin{subproof}
                Fix $u$.
                Assume $u\in\ran{g}$.
                Take $t\in\dom{g}$ such that $g(t) = u$ by \cref{fun_ran_iff}.
                Thus $(t,g(t))\in g$ by \cref{function_apply_elim}.
                Thus $g(t)\in\ran{g}$ by \cref{ran_intro}.
                Thus $\ran{g}\subseteq X$ by \cref{funs_ran}.
                Thus $g(t)\in X$ by \cref{subseteq}.
                Thus $(x,g(t))\in X\times X$ by \cref{times_tuple_intro}.
                Thus there exists $h$ such that $h$ is a path in $X$ from $x$ to $g(t)$ with respect to $T$
                    by \cref{path_initial_segment}.
                Thus $\fst{(x,g(t))} = x$ and $\snd{(x,g(t))} = g(t)$ by \cref{fst_eq,snd_eq}.
                Thus there exists $h$ such that
                    $h$ is a path in $X$ from $\fst{(x,g(t))}$ to $\snd{(x,g(t))}$ with respect to $T$.
                Thus $(x,g(t))\in \pathRelation{T}{X}$ by \cref{path_relation}.
                Thus $x\mathrel{\pathRelation{T}{X}}g(t)$.
                Thus $g(t)\mathrel{\pathRelation{T}{X}}x$ by \cref{path_relation_symmetric,symmetric}.
                Thus $u\mathrel{\pathRelation{T}{X}}x$.
                Thus $u\in C$ by \cref{downward_closure_iff}.
            \end{subproof}
            Thus $\ran{g}\subseteq C$ by \cref{subseteq}.
        \end{subproof}
        $\img{f}{I}\subseteq\ran{f}$ by \cref{img_subseteq_ran}.
        $\img{g}{J}\subseteq\ran{g}$ by \cref{img_subseteq_ran}.
        Thus $\img{f}{I}\subseteq C$ by \cref{subseteq}.
        Thus $\img{g}{J}\subseteq C$ by \cref{subseteq}.
        Thus $f$ is continuous from $I$ to $C$ with respect to $T_I$ and $T_C$
            by \cref{continuous_restrict_range}.
        Thus $g$ is continuous from $J$ to $C$ with respect to $T_J$ and $T_C$
            by \cref{continuous_restrict_range}.
        Thus $f$ is a path in $C$ from $x$ to $y$ with respect to $T_C$ by \cref{path_in_space}.
        Thus $g$ is a path in $C$ from $x$ to $z$ with respect to $T_C$ by \cref{path_in_space}.
        Thus there exists $f_1$ such that $f_1$ is a path in $C$ from $y$ to $x$ with respect to $T_C$
            by \cref{path_reverse}.
        Take $f_1$ such that $f_1$ is a path in $C$ from $y$ to $x$ with respect to $T_C$.
        Thus there exists $h$ such that $h$ is a path in $C$ from $y$ to $z$ with respect to $T_C$
            by \cref{path_concatenation}.
    \end{subproof}
    Thus $C$ is path connected with respect to $T_C$ by \cref{path_connected}.
\end{proof}

% from §25 Item 10: distinct path components are disjoint
\begin{theorem}\label{path_components_disjoint}
    Assume $T$ is a topology on $X$.
    Suppose $C$ is a path component of $X$ with respect to $T$.
    Suppose $D$ is a path component of $X$ with respect to $T$.
    Suppose $C\neq D$.
    Then $C$ is disjoint from $D$.
\end{theorem}
\begin{proof}
    Suppose not.
    Then there exists $z$ such that $z\in C$ and $z\in D$.
    Take $x\in X$ such that $C = \equivalenceClass{\pathRelation{T}{X}}{x}$ by \cref{path_component}.
    Take $y\in X$ such that $D = \equivalenceClass{\pathRelation{T}{X}}{y}$ by \cref{path_component}.
    Thus $z\in \downward{\pathRelation{T}{X}}{x}$.
    Thus $z\mathrel{\pathRelation{T}{X}}x$ by \cref{downward_closure_iff}.
    Thus $z\in \downward{\pathRelation{T}{X}}{y}$.
    Thus $z\mathrel{\pathRelation{T}{X}}y$ by \cref{downward_closure_iff}.
    $\pathRelation{T}{X}$ is symmetric by \cref{path_relation_symmetric}.
    Thus $x\mathrel{\pathRelation{T}{X}}z$ by \cref{symmetric}.
    $\pathRelation{T}{X}$ is transitive by \cref{path_relation_transitive}.
    Thus $x\mathrel{\pathRelation{T}{X}}y$ by \cref{transitive}.
    $\pathRelation{T}{X}$ is an equivalence on $X$ by \cref{path_relation_equivalence}.
    Thus $\equivalenceClass{\pathRelation{T}{X}}{x} = \equivalenceClass{\pathRelation{T}{X}}{y}$ by \cref{equiv_iff_equivclasses_eq}.
    Thus $C = D$.
    Contradiction.
\end{proof}

% from §25 Item 11: path components cover the space
\begin{theorem}\label{path_components_cover}
    Assume $T$ is a topology on $X$.
    Then for all $x\in X$ there exists $C$ such that
        $C$ is a path component of $X$ with respect to $T$ and $x\in C$.
\end{theorem}
\begin{proof}
    Fix $x\in X$.
    Let $C = \equivalenceClass{\pathRelation{T}{X}}{x}$.
    $\pathRelation{T}{X}$ is an equivalence on $X$ by \cref{path_relation_equivalence}.
    Thus $x\in C$ by \cref{equivclasses_inhabited_equivalenceon}.
    Thus $C$ is a path component of $X$ with respect to $T$ by \cref{path_component}.
    Thus there exists $C$ such that $C$ is a path component of $X$ with respect to $T$ and $x\in C$.
\end{proof}

% from §25 Item 12: path-connected subspaces intersect only one path component
\begin{theorem}\label{path_connected_subspace_single_component}
    Assume $T$ is a topology on $X$.
    Suppose $A$ is a subspace of $X$ with respect to $T$.
    Suppose $A$ is path connected with respect to $\subspaceTopology{T}{A}$.
    Suppose $A\neq\emptyset$.
    Then for all $C,D$ such that
        $C$ is a path component of $X$ with respect to $T$ and
        $D$ is a path component of $X$ with respect to $T$ and
        $A$ intersects $C$ and $A$ intersects $D$
        we have $C = D$.
\end{theorem}
\begin{proof}
    Fix $C$.
    Fix $D$.
    Assume $C$ is a path component of $X$ with respect to $T$.
    Assume $D$ is a path component of $X$ with respect to $T$.
    Assume $A$ intersects $C$ and $A$ intersects $D$.
    $A\inter C\neq\emptyset$ by \cref{intersects}.
    Take $x$ such that $x\in A\inter C$ by \cref{nonempty_inhabited}.
    Thus $x\in A$ and $x\in C$ by \cref{inter}.
    $A\inter D\neq\emptyset$ by \cref{intersects}.
    Take $y$ such that $y\in A\inter D$ by \cref{nonempty_inhabited}.
    Thus $y\in A$ and $y\in D$ by \cref{inter}.
    Thus $x\in A$ and $y\in A$ by \cref{inter}.
    Let $T_A = \subspaceTopology{T}{A}$.
    Thus $T_A$ is a topology on $A$ by \cref{subspace_topology_is_topology,subspace_of}.
    Thus there exists $f$ such that $f$ is a path in $A$ from $x$ to $y$ with respect to $T_A$
        by \cref{path_connected}.
    Take $f$ such that $f$ is a path in $A$ from $x$ to $y$ with respect to $T_A$.
    Take $a,b\in\reals$ such that
        $a\mathrel{\realOrderRel} b$ and
        $f$ is continuous from $\intervalclosedrel{\realOrderRel}{\reals}{a}{b}$ to $A$ with respect to
            $\subspaceTopology{\realOrderTopology}{\intervalclosedrel{\realOrderRel}{\reals}{a}{b}}$ and $T_A$ and
        $f(a) = x$ and $f(b) = y$
        by \cref{path_in_space}.
    Let $I = \intervalclosedrel{\realOrderRel}{\reals}{a}{b}$.
    Let $T_I = \subspaceTopology{\realOrderTopology}{I}$.
    Let $j = \identity{A}$.
    Thus $j$ is continuous from $A$ to $X$ with respect to $T_A$ and $T$ by \cref{inclusion_continuous}.
    Let $f_1 = j\circ f$.
    $f\in\funs{I}{A}$ by \cref{continuous}.
    Thus $f$ is a function by \cref{funs_is_function}.
    Thus $\ran{f}\subseteq A$ by \cref{funs_ran}.
    $j\in\funs{A}{X}$ by \cref{continuous}.
    Thus $j$ is a function by \cref{funs_is_function}.
    Thus $\dom{j} = A$ by \cref{funs}.
    Thus $\ran{f}\subseteq \dom{j}$ by \cref{subseteq}.
    Thus $j$ is composable with $f$.
    Thus $f_1$ is continuous from $I$ to $X$ with respect to $T_I$ and $T$ by \cref{continuous_composite}.
    Thus $a\in I$ and $b\in I$ by \cref{intervalclosed_rel}.
    Thus $\dom{f} = I$ by \cref{funs}.
    Thus $a\in\dom{f}$ and $b\in\dom{f}$ by \cref{subseteq}.
    Thus $(a,f(a))\in f$ and $(b,f(b))\in f$ by \cref{function_apply_elim}.
    Thus $f(a)\in\ran{f}$ and $f(b)\in\ran{f}$ by \cref{ran_intro}.
    Thus $f(a)\in A$ and $f(b)\in A$ by \cref{subseteq}.
    Thus $f_1(a) = j(f(a))$ and $f_1(b) = j(f(b))$ by \cref{circ_apply}.
    Thus $j(f(a)) = f(a)$ and $j(f(b)) = f(b)$ by \cref{id_apply}.
    Thus $f_1(a) = x$ and $f_1(b) = y$.
    Thus $f_1$ is a path in $X$ from $x$ to $y$ with respect to $T$ by \cref{path_in_space}.
    Thus $x\in X$ and $y\in X$ by \cref{subspace_of,subseteq}.
    Thus $(x,y)\in X\times X$ by \cref{times_tuple_intro}.
    Thus $\fst{(x,y)} = x$ and $\snd{(x,y)} = y$ by \cref{fst_eq,snd_eq}.
    Thus there exists $f_1$ such that
        $f_1$ is a path in $X$ from $\fst{(x,y)}$ to $\snd{(x,y)}$ with respect to $T$.
    Thus $(x,y)\in \pathRelation{T}{X}$ by \cref{path_relation}.
    Thus $x\mathrel{\pathRelation{T}{X}}y$.
    Take $c\in X$ such that $C = \equivalenceClass{\pathRelation{T}{X}}{c}$ by \cref{path_component}.
    Take $d\in X$ such that $D = \equivalenceClass{\pathRelation{T}{X}}{d}$ by \cref{path_component}.
    $\pathRelation{T}{X}$ is an equivalence on $X$ by \cref{path_relation_equivalence}.
    Thus $x\in \downward{\pathRelation{T}{X}}{c}$.
    Thus $x\mathrel{\pathRelation{T}{X}}c$ by \cref{downward_closure_iff}.
    Thus $y\in \downward{\pathRelation{T}{X}}{d}$.
    Thus $y\mathrel{\pathRelation{T}{X}}d$ by \cref{downward_closure_iff}.
    $\pathRelation{T}{X}$ is symmetric by \cref{path_relation_symmetric}.
    Thus $c\mathrel{\pathRelation{T}{X}}x$ by \cref{symmetric}.
    $\pathRelation{T}{X}$ is transitive by \cref{path_relation_transitive}.
    Thus $c\mathrel{\pathRelation{T}{X}}y$ by \cref{transitive}.
    Thus $c\mathrel{\pathRelation{T}{X}}d$ by \cref{transitive,symmetric}.
    Thus $\equivalenceClass{\pathRelation{T}{X}}{c} = \equivalenceClass{\pathRelation{T}{X}}{d}$ by \cref{equiv_iff_equivclasses_eq}.
    Thus $C = D$.
\end{proof}

% from §25 Item 13: locally connected at a point
\begin{definition}\label{locally_connected_at}
    $X$ is locally connected at $x$ with respect to $T$ iff
        for all $U$ such that $U$ is a neighborhood of $x$ in $X$ with respect to $T$
        there exists $V$ such that
            $V$ is a neighborhood of $x$ in $X$ with respect to $T$ and
            $V\subseteq U$ and
            $V$ is connected with respect to $\subspaceTopology{T}{V}$.
\end{definition}

% from §25 Item 14: locally connected
\begin{definition}\label{locally_connected}
    $X$ is locally connected with respect to $T$ iff
        for all $x\in X$ we have $X$ is locally connected at $x$ with respect to $T$.
\end{definition}

% from §25 Item 15: locally path connected at a point
\begin{definition}\label{locally_path_connected_at}
    $X$ is locally path connected at $x$ with respect to $T$ iff
        for all $U$ such that $U$ is a neighborhood of $x$ in $X$ with respect to $T$
        there exists $V$ such that
            $V$ is a neighborhood of $x$ in $X$ with respect to $T$ and
            $V\subseteq U$ and
            $V$ is path connected with respect to $\subspaceTopology{T}{V}$.
\end{definition}

% from §25 Item 16: locally path connected
\begin{definition}\label{locally_path_connected}
    $X$ is locally path connected with respect to $T$ iff
        for all $x\in X$ we have $X$ is locally path connected at $x$ with respect to $T$.
\end{definition}

% from §25 Item 17: locally connected iff components of open sets are open
\begin{theorem}\label{locally_connected_iff_components_open}
    Assume $T$ is a topology on $X$.
    Then $X$ is locally connected with respect to $T$ iff
        for all $U$ such that $U$ is open in $X$ with respect to $T$
        for all $C$ such that $C$ is a component of $U$ with respect to $\subspaceTopology{T}{U}$
        we have $C$ is open in $X$ with respect to $T$.
\end{theorem}
\begin{proof}
    Show that if $X$ is locally connected with respect to $T$ then
        for all $U$ such that $U$ is open in $X$ with respect to $T$
        for all $C$ such that $C$ is a component of $U$ with respect to $\subspaceTopology{T}{U}$
        we have $C$ is open in $X$ with respect to $T$.
    \begin{subproof}
        Assume $X$ is locally connected with respect to $T$.
        Fix $U$.
        Assume $U$ is open in $X$ with respect to $T$.
        $T\subseteq\pow{X}$ by \cref{topology_on}.
        Thus $U\in T$ by \cref{open_in}.
        Thus $U\subseteq X$ by \cref{powerset_elim,subseteq}.
        Thus $U$ is a subspace of $X$ with respect to $T$ by \cref{subspace_of}.
        Let $T_U = \subspaceTopology{T}{U}$.
        Thus $T_U$ is a topology on $U$ by \cref{subspace_topology_is_topology,subspace_of}.
        Fix $C$.
        Assume $C$ is a component of $U$ with respect to $\subspaceTopology{T}{U}$.
        Thus $C$ is a component of $U$ with respect to $T_U$.
        Show that $C$ is open in $X$ with respect to $T$.
        \begin{subproof}
            Show that for all $x$ such that $x\in C$ we have
                there exists $V$ such that $V$ is open in $X$ with respect to $T$ and $x\in V$ and $V\subseteq C$.
            \begin{subproof}
                Fix $x$.
                Assume $x\in C$.
                Show that $C\subseteq U$.
                \begin{subproof}
                    Take $x_0\in U$ such that $C = \equivalenceClass{\componentRelation{T_U}{U}}{x_0}$ by \cref{component}.
                    Show that for all $u$ such that $u\in C$ we have $u\in U$.
                    \begin{subproof}
                        Fix $u$.
                        Assume $u\in C$.
                        Thus $u\in \downward{\componentRelation{T_U}{U}}{x_0}$.
                        Thus $u\mathrel{\componentRelation{T_U}{U}}x_0$ by \cref{downward_closure_iff}.
                        Thus $(u,x_0)\in U\times U$ and there exists $A$ such that
                            $A$ is a subspace of $U$ with respect to $T_U$ and
                            $\fst{(u,x_0)}\in A$ and $\snd{(u,x_0)}\in A$ and
                            $A$ is connected with respect to $\subspaceTopology{T_U}{A}$
                            by \cref{component_relation}.
                        Thus $u\in U$ by \cref{times_tuple_elim}.
                    \end{subproof}
                    Thus $C\subseteq U$ by \cref{subseteq}.
                \end{subproof}
                Thus $x\in U$ by \cref{subseteq}.
                Thus $x\in X$ by \cref{subseteq}.
                Thus $U$ is a neighborhood of $x$ in $X$ with respect to $T$ by \cref{neighborhood}.
                $X$ is locally connected at $x$ with respect to $T$ by \cref{locally_connected}.
                Thus there exists $V$ such that
                    $V$ is a neighborhood of $x$ in $X$ with respect to $T$ and
                    $V\subseteq U$ and
                    $V$ is connected with respect to $\subspaceTopology{T}{V}$
                    by \cref{locally_connected_at}.
                Take $V$ such that
                    $V$ is a neighborhood of $x$ in $X$ with respect to $T$ and
                    $V\subseteq U$ and
                    $V$ is connected with respect to $\subspaceTopology{T}{V}$.
                Thus $V$ is open in $X$ with respect to $T$ and $x\in V$ by \cref{neighborhood}.
                Show that $V\subseteq C$.
                \begin{subproof}
                    Show that for all $y$ such that $y\in V$ we have $y\in C$.
                    \begin{subproof}
                        Fix $y$.
                        Assume $y\in V$.
                        Thus $y\in U$ by \cref{subseteq}.
                        Thus $V$ is a subspace of $U$ with respect to $T_U$ by \cref{subspace_of}.
                        Let $T_{UV} = \subspaceTopology{T_U}{V}$.
                        Let $T_V = \subspaceTopology{T}{V}$.
                        Thus $T_{UV} = T_V$ by \cref{subspace_subspace_topology,subseteq}.
                        Thus $V$ is connected with respect to $T_{UV}$.
                        Thus there exists $D$ such that
                            $D$ is a component of $U$ with respect to $T_U$ and $y\in D$
                            by \cref{components_cover}.
                        Thus $y\in V\inter D$ by \cref{inter}.
                        Thus $V\inter D\neq\emptyset$ by \cref{emptyset}.
                        Thus $V$ intersects $D$ by \cref{intersects}.
                        Thus $x\in V\inter C$ by \cref{inter}.
                        Thus $V\inter C\neq\emptyset$ by \cref{emptyset}.
                        Thus $V$ intersects $C$ by \cref{intersects}.
                        Thus $V\neq\emptyset$ by \cref{emptyset}.
                        Thus $C = D$ by \cref{connected_subspace_single_component,subspace_of}.
                        Thus $y\in C$.
                    \end{subproof}
                    Thus $V\subseteq C$ by \cref{subseteq}.
                \end{subproof}
                Thus there exists $V$ such that $V$ is open in $X$ with respect to $T$ and $x\in V$ and $V\subseteq C$.
            \end{subproof}
            Let $M = \{V\in T \mid V\subseteq C\}$.
            Show that $M\subseteq T$.
            \begin{subproof}
                Show that for all $V$ such that $V\in M$ we have $V\in T$.
                \begin{subproof}
                    Fix $V$.
                    Assume $V\in M$.
                    Thus $V\in T$.
                \end{subproof}
                Thus $M\subseteq T$ by \cref{subseteq}.
            \end{subproof}
            Show that $C = \unions{M}$.
            \begin{subproof}
                Show that $C\subseteq \unions{M}$.
                \begin{subproof}
                    Show that for all $x$ such that $x\in C$ we have $x\in \unions{M}$.
                    \begin{subproof}
                        Fix $x$.
                        Assume $x\in C$.
                        Take $V$ such that $V$ is open in $X$ with respect to $T$ and $x\in V$ and $V\subseteq C$.
                        Thus $V\in T$ by \cref{open_in}.
                        Thus $V\in M$.
                        Thus $x\in \unions{M}$ by \cref{unions_iff}.
                    \end{subproof}
                    Thus $C\subseteq \unions{M}$ by \cref{subseteq}.
                \end{subproof}
                Show that $\unions{M}\subseteq C$.
                \begin{subproof}
                    Show that for all $x$ such that $x\in \unions{M}$ we have $x\in C$.
                    \begin{subproof}
                        Fix $x$.
                        Assume $x\in \unions{M}$.
                        Take $V\in M$ such that $x\in V$ by \cref{unions_iff}.
                        Thus $V\subseteq C$.
                        Thus $x\in C$ by \cref{subseteq}.
                    \end{subproof}
                    Thus $\unions{M}\subseteq C$ by \cref{subseteq}.
                \end{subproof}
                Thus $C = \unions{M}$ by \cref{subseteq_antisymmetric}.
            \end{subproof}
            Thus $\unions{M}\in T$ by \cref{topology_on,subseteq}.
            Thus $C\in T$.
            Thus $C$ is open in $X$ with respect to $T$ by \cref{open_in}.
        \end{subproof}
    \end{subproof}
    Show that if for all $U$ such that $U$ is open in $X$ with respect to $T$
        for all $C$ such that $C$ is a component of $U$ with respect to $\subspaceTopology{T}{U}$
        we have $C$ is open in $X$ with respect to $T$ then
        $X$ is locally connected with respect to $T$.
    \begin{subproof}
        Assume for all $U$ such that $U$ is open in $X$ with respect to $T$
            for all $C$ such that $C$ is a component of $U$ with respect to $\subspaceTopology{T}{U}$
            we have $C$ is open in $X$ with respect to $T$.
        Show that for all $x\in X$ we have $X$ is locally connected at $x$ with respect to $T$.
        \begin{subproof}
            Fix $x\in X$.
            Show that for all $U$ such that $U$ is a neighborhood of $x$ in $X$ with respect to $T$
                there exists $V$ such that
                    $V$ is a neighborhood of $x$ in $X$ with respect to $T$ and
                    $V\subseteq U$ and
                    $V$ is connected with respect to $\subspaceTopology{T}{V}$.
            \begin{subproof}
                Fix $U$.
                Assume $U$ is a neighborhood of $x$ in $X$ with respect to $T$.
                Thus $U$ is open in $X$ with respect to $T$ and $x\in U$ by \cref{neighborhood}.
                $T\subseteq\pow{X}$ by \cref{topology_on}.
                Thus $U\in T$ by \cref{open_in}.
                Thus $U\subseteq X$ by \cref{powerset_elim,subseteq}.
                Thus $U$ is a subspace of $X$ with respect to $T$ by \cref{subspace_of}.
                Let $T_U = \subspaceTopology{T}{U}$.
                Thus $T_U$ is a topology on $U$ by \cref{subspace_topology_is_topology,subspace_of}.
                Thus there exists $C$ such that
                    $C$ is a component of $U$ with respect to $T_U$ and $x\in C$
                    by \cref{components_cover}.
                Take $C$ such that
                    $C$ is a component of $U$ with respect to $T_U$ and $x\in C$.
                Thus $C$ is open in $X$ with respect to $T$.
                Thus $C$ is connected with respect to $\subspaceTopology{T_U}{C}$
                    by \cref{component_connected}.
                Show that $C\subseteq U$.
                \begin{subproof}
                    Take $x_0\in U$ such that $C = \equivalenceClass{\componentRelation{T_U}{U}}{x_0}$ by \cref{component}.
                    Show that for all $u$ such that $u\in C$ we have $u\in U$.
                    \begin{subproof}
                        Fix $u$.
                        Assume $u\in C$.
                        Thus $u\in \downward{\componentRelation{T_U}{U}}{x_0}$.
                        Thus $u\mathrel{\componentRelation{T_U}{U}}x_0$ by \cref{downward_closure_iff}.
                        Thus $(u,x_0)\in U\times U$ and there exists $A$ such that
                            $A$ is a subspace of $U$ with respect to $T_U$ and
                            $\fst{(u,x_0)}\in A$ and $\snd{(u,x_0)}\in A$ and
                            $A$ is connected with respect to $\subspaceTopology{T_U}{A}$
                            by \cref{component_relation}.
                        Thus $u\in U$ by \cref{times_tuple_elim}.
                    \end{subproof}
                    Thus $C\subseteq U$ by \cref{subseteq}.
                \end{subproof}
                Let $T_C = \subspaceTopology{T}{C}$.
                Thus $T_C = \subspaceTopology{T_U}{C}$ by \cref{subspace_subspace_topology,subseteq}.
                Thus $C$ is connected with respect to $T_C$.
                Thus $C$ is a neighborhood of $x$ in $X$ with respect to $T$ by \cref{neighborhood,open_in}.
                Thus $C\subseteq U$ by \cref{subseteq}.
                Thus there exists $V$ such that
                    $V$ is a neighborhood of $x$ in $X$ with respect to $T$ and
                    $V\subseteq U$ and
                    $V$ is connected with respect to $\subspaceTopology{T}{V}$.
            \end{subproof}
            Thus $X$ is locally connected at $x$ with respect to $T$ by \cref{locally_connected_at}.
        \end{subproof}
        Thus $X$ is locally connected with respect to $T$ by \cref{locally_connected}.
    \end{subproof}
\end{proof}

% from §25 Item 18: locally path connected iff path components of open sets are open
\begin{theorem}\label{locally_path_connected_iff_components_open}
    Assume $T$ is a topology on $X$.
    Then $X$ is locally path connected with respect to $T$ iff
        for all $U$ such that $U$ is open in $X$ with respect to $T$
        for all $C$ such that $C$ is a path component of $U$ with respect to $\subspaceTopology{T}{U}$
        we have $C$ is open in $X$ with respect to $T$.
\end{theorem}
\begin{proof}
    Show that if $X$ is locally path connected with respect to $T$ then
        for all $U$ such that $U$ is open in $X$ with respect to $T$
        for all $C$ such that $C$ is a path component of $U$ with respect to $\subspaceTopology{T}{U}$
        we have $C$ is open in $X$ with respect to $T$.
    \begin{subproof}
        Assume $X$ is locally path connected with respect to $T$.
        Fix $U$.
        Assume $U$ is open in $X$ with respect to $T$.
        $T\subseteq\pow{X}$ by \cref{topology_on}.
        Thus $U\in T$ by \cref{open_in}.
        Thus $U\subseteq X$ by \cref{powerset_elim,subseteq}.
        Thus $U$ is a subspace of $X$ with respect to $T$ by \cref{subspace_of}.
        Let $T_U = \subspaceTopology{T}{U}$.
        Thus $T_U$ is a topology on $U$ by \cref{subspace_topology_is_topology,subspace_of}.
        Fix $C$.
        Assume $C$ is a path component of $U$ with respect to $\subspaceTopology{T}{U}$.
        Thus $C$ is a path component of $U$ with respect to $T_U$.
        Show that $C$ is open in $X$ with respect to $T$.
        \begin{subproof}
            Show that for all $x$ such that $x\in C$ we have
                there exists $V$ such that $V$ is open in $X$ with respect to $T$ and $x\in V$ and $V\subseteq C$.
            \begin{subproof}
                Fix $x$.
                Assume $x\in C$.
                Show that $C\subseteq U$.
                \begin{subproof}
                    Take $x_0\in U$ such that $C = \equivalenceClass{\pathRelation{T_U}{U}}{x_0}$ by \cref{path_component}.
                    Show that for all $u$ such that $u\in C$ we have $u\in U$.
                    \begin{subproof}
                        Fix $u$.
                        Assume $u\in C$.
                        Thus $u\in \downward{\pathRelation{T_U}{U}}{x_0}$.
                        Thus $u\mathrel{\pathRelation{T_U}{U}}x_0$ by \cref{downward_closure_iff}.
                        Thus $(u,x_0)\in U\times U$ and there exists $f$ such that
                            $f$ is a path in $U$ from $\fst{(u,x_0)}$ to $\snd{(u,x_0)}$ with respect to $T_U$
                            by \cref{path_relation}.
                        Thus $u\in U$ by \cref{times_tuple_elim}.
                    \end{subproof}
                    Thus $C\subseteq U$ by \cref{subseteq}.
                \end{subproof}
                Thus $x\in U$ by \cref{subseteq}.
                Thus $U$ is a neighborhood of $x$ in $X$ with respect to $T$ by \cref{neighborhood}.
                Thus $x\in X$ by \cref{subseteq}.
                $X$ is locally path connected at $x$ with respect to $T$ by \cref{locally_path_connected}.
                Thus there exists $V$ such that
                    $V$ is a neighborhood of $x$ in $X$ with respect to $T$ and
                    $V\subseteq U$ and
                    $V$ is path connected with respect to $\subspaceTopology{T}{V}$
                    by \cref{locally_path_connected_at}.
                Take $V$ such that
                    $V$ is a neighborhood of $x$ in $X$ with respect to $T$ and
                    $V\subseteq U$ and
                    $V$ is path connected with respect to $\subspaceTopology{T}{V}$.
                Thus $V$ is open in $X$ with respect to $T$ and $x\in V$ by \cref{neighborhood}.
                Show that $V\subseteq C$.
                \begin{subproof}
                    Show that for all $y$ such that $y\in V$ we have $y\in C$.
                    \begin{subproof}
                        Fix $y$.
                        Assume $y\in V$.
                        Thus $y\in U$ by \cref{subseteq}.
                        Thus $V$ is a subspace of $U$ with respect to $T_U$ by \cref{subspace_of}.
                        Let $T_{UV} = \subspaceTopology{T_U}{V}$.
                        Let $T_V = \subspaceTopology{T}{V}$.
                        Thus $T_{UV} = T_V$ by \cref{subspace_subspace_topology,subseteq}.
                        Thus $V$ is path connected with respect to $T_{UV}$.
                        Thus there exists $D$ such that
                            $D$ is a path component of $U$ with respect to $T_U$ and $y\in D$
                            by \cref{path_components_cover}.
                        Thus $y\in V\inter D$ by \cref{inter}.
                        Thus $V\inter D\neq\emptyset$ by \cref{emptyset}.
                        Thus $V$ intersects $D$ by \cref{intersects}.
                        Thus $x\in V\inter C$ by \cref{inter}.
                        Thus $V\inter C\neq\emptyset$ by \cref{emptyset}.
                        Thus $V$ intersects $C$ by \cref{intersects}.
                        Thus $V\neq\emptyset$ by \cref{emptyset}.
                        Thus $C = D$ by \cref{path_connected_subspace_single_component,subspace_of}.
                        Thus $y\in C$.
                    \end{subproof}
                    Thus $V\subseteq C$ by \cref{subseteq}.
                \end{subproof}
                Thus there exists $V$ such that $V$ is open in $X$ with respect to $T$ and $x\in V$ and $V\subseteq C$.
            \end{subproof}
            Let $M = \{V\in T \mid V\subseteq C\}$.
            Show that $M\subseteq T$.
            \begin{subproof}
                Show that for all $V$ such that $V\in M$ we have $V\in T$.
                \begin{subproof}
                    Fix $V$.
                    Assume $V\in M$.
                    Thus $V\in T$.
                \end{subproof}
                Thus $M\subseteq T$ by \cref{subseteq}.
            \end{subproof}
            Show that $C = \unions{M}$.
            \begin{subproof}
                Show that $C\subseteq \unions{M}$.
                \begin{subproof}
                    Show that for all $x$ such that $x\in C$ we have $x\in \unions{M}$.
                    \begin{subproof}
                        Fix $x$.
                        Assume $x\in C$.
                        Take $V$ such that $V$ is open in $X$ with respect to $T$ and $x\in V$ and $V\subseteq C$.
                        Thus $V\in T$ by \cref{open_in}.
                        Thus $V\in M$.
                        Thus $x\in \unions{M}$ by \cref{unions_iff}.
                    \end{subproof}
                    Thus $C\subseteq \unions{M}$ by \cref{subseteq}.
                \end{subproof}
                Show that $\unions{M}\subseteq C$.
                \begin{subproof}
                    Show that for all $x$ such that $x\in \unions{M}$ we have $x\in C$.
                    \begin{subproof}
                        Fix $x$.
                        Assume $x\in \unions{M}$.
                        Take $V\in M$ such that $x\in V$ by \cref{unions_iff}.
                        Thus $V\subseteq C$.
                        Thus $x\in C$ by \cref{subseteq}.
                    \end{subproof}
                    Thus $\unions{M}\subseteq C$ by \cref{subseteq}.
                \end{subproof}
                Thus $C = \unions{M}$ by \cref{subseteq_antisymmetric}.
            \end{subproof}
            Thus $\unions{M}\in T$ by \cref{topology_on,subseteq}.
            Thus $C\in T$.
            Thus $C$ is open in $X$ with respect to $T$ by \cref{open_in}.
        \end{subproof}
    \end{subproof}
    Show that if for all $U$ such that $U$ is open in $X$ with respect to $T$
        for all $C$ such that $C$ is a path component of $U$ with respect to $\subspaceTopology{T}{U}$
        we have $C$ is open in $X$ with respect to $T$ then
        $X$ is locally path connected with respect to $T$.
    \begin{subproof}
        Assume for all $U$ such that $U$ is open in $X$ with respect to $T$
            for all $C$ such that $C$ is a path component of $U$ with respect to $\subspaceTopology{T}{U}$
            we have $C$ is open in $X$ with respect to $T$.
        Show that for all $x\in X$ we have $X$ is locally path connected at $x$ with respect to $T$.
        \begin{subproof}
            Fix $x\in X$.
            Show that for all $U$ such that $U$ is a neighborhood of $x$ in $X$ with respect to $T$
                there exists $V$ such that
                    $V$ is a neighborhood of $x$ in $X$ with respect to $T$ and
                    $V\subseteq U$ and
                    $V$ is path connected with respect to $\subspaceTopology{T}{V}$.
            \begin{subproof}
                Fix $U$.
                Assume $U$ is a neighborhood of $x$ in $X$ with respect to $T$.
                Thus $U$ is open in $X$ with respect to $T$ and $x\in U$ by \cref{neighborhood}.
                $T\subseteq\pow{X}$ by \cref{topology_on}.
                Thus $U\in T$ by \cref{open_in}.
                Thus $U\subseteq X$ by \cref{powerset_elim,subseteq}.
                Thus $U$ is a subspace of $X$ with respect to $T$ by \cref{subspace_of}.
                Let $T_U = \subspaceTopology{T}{U}$.
                Thus $T_U$ is a topology on $U$ by \cref{subspace_topology_is_topology,subspace_of}.
                Thus there exists $C$ such that
                    $C$ is a path component of $U$ with respect to $T_U$ and $x\in C$
                    by \cref{path_components_cover}.
                Take $C$ such that
                    $C$ is a path component of $U$ with respect to $T_U$ and $x\in C$.
                Thus $C$ is open in $X$ with respect to $T$.
                Thus $C$ is path connected with respect to $\subspaceTopology{T_U}{C}$
                    by \cref{path_component_path_connected}.
                Show that $C\subseteq U$.
                \begin{subproof}
                    Take $x_0\in U$ such that $C = \equivalenceClass{\pathRelation{T_U}{U}}{x_0}$ by \cref{path_component}.
                    Show that for all $u$ such that $u\in C$ we have $u\in U$.
                    \begin{subproof}
                        Fix $u$.
                        Assume $u\in C$.
                        Thus $u\in \downward{\pathRelation{T_U}{U}}{x_0}$.
                        Thus $u\mathrel{\pathRelation{T_U}{U}}x_0$ by \cref{downward_closure_iff}.
                        Thus $(u,x_0)\in U\times U$ and there exists $f$ such that
                            $f$ is a path in $U$ from $\fst{(u,x_0)}$ to $\snd{(u,x_0)}$ with respect to $T_U$
                            by \cref{path_relation}.
                        Thus $u\in U$ by \cref{times_tuple_elim}.
                    \end{subproof}
                    Thus $C\subseteq U$ by \cref{subseteq}.
                \end{subproof}
                Let $T_C = \subspaceTopology{T}{C}$.
                Thus $T_C = \subspaceTopology{T_U}{C}$ by \cref{subspace_subspace_topology,subseteq}.
                Thus $C$ is path connected with respect to $T_C$.
                Thus $C$ is a neighborhood of $x$ in $X$ with respect to $T$ by \cref{neighborhood,open_in}.
                Thus $C\subseteq U$ by \cref{subseteq}.
                Thus there exists $V$ such that
                    $V$ is a neighborhood of $x$ in $X$ with respect to $T$ and
                    $V\subseteq U$ and
                    $V$ is path connected with respect to $\subspaceTopology{T}{V}$.
            \end{subproof}
            Thus $X$ is locally path connected at $x$ with respect to $T$ by \cref{locally_path_connected_at}.
        \end{subproof}
        Thus $X$ is locally path connected with respect to $T$ by \cref{locally_path_connected}.
    \end{subproof}
\end{proof}

% from §25 Item 19: path components lie in components
\begin{theorem}\label{path_component_subset_component}
    Assume $T$ is a topology on $X$.
    Suppose $P$ is a path component of $X$ with respect to $T$.
    Then there exists $C$ such that $C$ is a component of $X$ with respect to $T$ and $P\subseteq C$.
\end{theorem}
\begin{proof}
    Take $x\in X$ such that $P = \equivalenceClass{\pathRelation{T}{X}}{x}$ by \cref{path_component}.
    Show that $P\subseteq X$.
    \begin{subproof}
        Show that for all $y$ such that $y\in P$ we have $y\in X$.
        \begin{subproof}
            Fix $y$.
            Assume $y\in P$.
            Thus $y\mathrel{\pathRelation{T}{X}}x$ by \cref{downward_closure_iff}.
            Thus $(y,x)\in X\times X$ and there exists $f$ such that
                $f$ is a path in $X$ from $\fst{(y,x)}$ to $\snd{(y,x)}$ with respect to $T$
                by \cref{path_relation}.
            Thus $y\in X$ by \cref{times_tuple_elim}.
        \end{subproof}
        Thus $P\subseteq X$ by \cref{subseteq}.
    \end{subproof}
    Thus $P$ is a subspace of $X$ with respect to $T$ by \cref{subspace_of}.
    Let $T_P = \subspaceTopology{T}{P}$.
    Thus $T_P$ is a topology on $P$ by \cref{subspace_topology_is_topology,subspace_of}.
    Thus $P$ is path connected with respect to $T_P$ by \cref{path_component_path_connected}.
    Thus $P$ is connected with respect to $T_P$ by \cref{path_connected_implies_connected}.
    Thus there exists $C$ such that
        $C$ is a component of $X$ with respect to $T$ and $x\in C$
        by \cref{components_cover}.
    Take $C$ such that
        $C$ is a component of $X$ with respect to $T$ and $x\in C$.
    Show that $P\subseteq C$.
    \begin{subproof}
        Show that for all $y$ such that $y\in P$ we have $y\in C$.
        \begin{subproof}
            Fix $y$.
            Assume $y\in P$.
            Thus $y\in X$ by \cref{subseteq}.
            Thus there exists $D$ such that
                $D$ is a component of $X$ with respect to $T$ and $y\in D$
                by \cref{components_cover}.
            Take $D$ such that
                $D$ is a component of $X$ with respect to $T$ and $y\in D$.
            Thus $x\in P$ by \cref{path_relation_equivalence,equivclasses_inhabited_equivalenceon}.
            Thus $x\in P\inter C$ by \cref{inter}.
            Thus $P\inter C\neq\emptyset$ by \cref{emptyset}.
            Thus $P$ intersects $C$ by \cref{intersects}.
            Thus $y\in P\inter D$ by \cref{inter}.
            Thus $P\inter D\neq\emptyset$ by \cref{emptyset}.
            Thus $P$ intersects $D$ by \cref{intersects}.
            Thus $P\neq\emptyset$ by \cref{emptyset}.
            Thus $C = D$ by \cref{connected_subspace_single_component,subspace_of}.
            Thus $y\in C$.
        \end{subproof}
        Thus $P\subseteq C$ by \cref{subseteq}.
    \end{subproof}
    Thus there exists $C$ such that $C$ is a component of $X$ with respect to $T$ and $P\subseteq C$.
\end{proof}


% from §25 Item 20: components equal path components in locally path connected spaces
\begin{theorem}\label{components_equal_path_components_locally_path_connected}
    Assume $T$ is a topology on $X$.
    Suppose $X$ is locally path connected with respect to $T$.
    Then for all $C$ such that $C$ is a component of $X$ with respect to $T$
        there exists $P$ such that $P$ is a path component of $X$ with respect to $T$ and $C = P$.
\end{theorem}
\begin{proof}
    Fix $C$.
    Assume $C$ is a component of $X$ with respect to $T$.
    Take $x\in X$ such that $C = \equivalenceClass{\componentRelation{T}{X}}{x}$ by \cref{component}.
    $\componentRelation{T}{X}$ is an equivalence on $X$ by \cref{component_relation_equivalence}.
    Thus $x\in C$ by \cref{equivclasses_inhabited_equivalenceon}.
    Show that $C\subseteq X$.
    \begin{subproof}
        Show that for all $y$ such that $y\in C$ we have $y\in X$.
        \begin{subproof}
            Fix $y$.
            Assume $y\in C$.
            Thus $y\in \downward{\componentRelation{T}{X}}{x}$.
            Thus $y\mathrel{\componentRelation{T}{X}}x$ by \cref{downward_closure_iff}.
            Thus $(y,x)\in X\times X$ and there exists $A$ such that
                $A$ is a subspace of $X$ with respect to $T$ and
                $\fst{(y,x)}\in A$ and $\snd{(y,x)}\in A$ and
                $A$ is connected with respect to $\subspaceTopology{T}{A}$
                by \cref{component_relation}.
            Thus $y\in X$ by \cref{times_tuple_elim}.
        \end{subproof}
        Thus $C\subseteq X$ by \cref{subseteq}.
    \end{subproof}
    Thus $C$ is a subspace of $X$ with respect to $T$ by \cref{subspace_of}.
    $X\subseteq X$ by \cref{subseteq_refl}.
    Thus $X$ is a subspace of $X$ with respect to $T$ by \cref{subspace_of}.
    Let $T_X = \subspaceTopology{T}{X}$.
    Thus $T_X$ is a topology on $X$ by \cref{subspace_topology_is_topology,subspace_of}.
    Show that $T_X = T$.
    \begin{subproof}
        Show that $T\subseteq T_X$.
        \begin{subproof}
            Show that for all $U$ such that $U\in T$ we have $U\in T_X$.
            \begin{subproof}
                Fix $U$.
                Assume $U\in T$.
                $T\subseteq\pow{X}$ by \cref{topology_on}.
                Thus $U\subseteq X$ by \cref{powerset_elim,subseteq}.
                Thus $U = X\inter U$ by \cref{inter_absorb_supseteq_left}.
                Thus $U\in T_X$ by \cref{subspace_topology}.
            \end{subproof}
            Thus $T\subseteq T_X$ by \cref{subseteq}.
        \end{subproof}
        Show that $T_X\subseteq T$.
        \begin{subproof}
            Show that for all $U$ such that $U\in T_X$ we have $U\in T$.
            \begin{subproof}
                Fix $U$.
                Assume $U\in T_X$.
                Take $V\in T$ such that $U = X\inter V$ by \cref{subspace_topology}.
                $T\subseteq\pow{X}$ by \cref{topology_on}.
                Thus $V\subseteq X$ by \cref{powerset_elim,subseteq}.
                Thus $U = V$ by \cref{inter_absorb_supseteq_left}.
                Thus $U\in T$.
            \end{subproof}
            Thus $T_X\subseteq T$ by \cref{subseteq}.
        \end{subproof}
        Thus $T_X = T$ by \cref{subseteq_antisymmetric}.
    \end{subproof}
    Show that for all $D$ such that $D$ is a path component of $X$ with respect to $T$
        we have $D$ is open in $X$ with respect to $T$.
    \begin{subproof}
        Fix $D$.
        Assume $D$ is a path component of $X$ with respect to $T$.
        Thus $D$ is a path component of $X$ with respect to $T_X$.
        $X\in T$ by \cref{topology_on}.
        Thus $X$ is open in $X$ with respect to $T$ by \cref{open_in}.
        Thus $D$ is open in $X$ with respect to $T$ by \cref{locally_path_connected_iff_components_open}.
    \end{subproof}
    Thus there exists $P$ such that
        $P$ is a path component of $X$ with respect to $T$ and $x\in P$
        by \cref{path_components_cover}.
    Take $P$ such that $P$ is a path component of $X$ with respect to $T$ and $x\in P$.
    Show that $C = P$.
    \begin{subproof}
        Suppose not.
        Thus $C\neq P$.
        Show that $P\subseteq C$.
        \begin{subproof}
            Thus there exists $C_1$ such that
                $C_1$ is a component of $X$ with respect to $T$ and $P\subseteq C_1$
                by \cref{path_component_subset_component}.
            Take $C_1$ such that
                $C_1$ is a component of $X$ with respect to $T$ and $P\subseteq C_1$.
            Thus $x\in C_1$ by \cref{subseteq}.
            Show that $C = C_1$.
            \begin{subproof}
                Suppose not.
                Then $C\neq C_1$.
                Then $C$ is disjoint from $C_1$ by \cref{components_disjoint}.
                Thus $x\in C$ and $x\in C_1$.
                Contradiction by \cref{disjoint}.
            \end{subproof}
            Thus $P\subseteq C$ by \cref{subseteq}.
        \end{subproof}
        Thus $P\subset C$ by \cref{subset}.
        Thus $C\setminus P$ is inhabited by \cref{difference_with_proper_subset_is_inhabited}.
        Take $y$ such that $y\in C\setminus P$ by \cref{nonempty_inhabited}.
        Thus $y\in C$ and $y\notin P$ by \cref{setminus_elim_left,setminus_elim_right}.
        Thus $y\in X$ by \cref{subseteq}.
        Let $M = \{D\in\pow{X} \mid
            \text{$D$ is a path component of $X$ with respect to $T$ and $D\subseteq C$ and $D\neq P$}\}$.
        Let $Q = \unions{M}$.
        Show that $Q = C\setminus P$.
        \begin{subproof}
            Show that $Q\subseteq C\setminus P$.
            \begin{subproof}
                Show that for all $z$ such that $z\in Q$ we have $z\in C\setminus P$.
                \begin{subproof}
                    Fix $z$.
                    Assume $z\in Q$.
                    Take $D\in M$ such that $z\in D$ by \cref{unions_iff}.
                    Thus $D$ is a path component of $X$ with respect to $T$ and $D\subseteq C$ and $D\neq P$.
                    Thus $P$ is a path component of $X$ with respect to $T$.
                    Thus $D$ is disjoint from $P$ by \cref{path_components_disjoint}.
                    Thus $z\in C$ by \cref{subseteq}.
                    Show that $z\notin P$.
                    \begin{subproof}
                        Suppose not.
                        Then $z\in P$.
                        Thus $z\in D$ and $z\in P$.
                        Contradiction by \cref{disjoint}.
                    \end{subproof}
                    Thus $z\in C\setminus P$ by \cref{setminus_intro}.
                \end{subproof}
                Thus $Q\subseteq C\setminus P$ by \cref{subseteq}.
            \end{subproof}
            Show that $C\setminus P\subseteq Q$.
            \begin{subproof}
                Show that for all $z$ such that $z\in C\setminus P$ we have $z\in Q$.
                \begin{subproof}
                    Fix $z$.
                    Assume $z\in C\setminus P$.
                    Thus $z\in C$ and $z\notin P$ by \cref{setminus_elim_left,setminus_elim_right}.
                    Thus $z\in X$ by \cref{subseteq}.
                    Thus there exists $D$ such that
                        $D$ is a path component of $X$ with respect to $T$ and $z\in D$
                        by \cref{path_components_cover}.
                    Take $D$ such that $D$ is a path component of $X$ with respect to $T$ and $z\in D$.
                    Thus there exists $C_1$ such that
                        $C_1$ is a component of $X$ with respect to $T$ and $D\subseteq C_1$
                        by \cref{path_component_subset_component}.
                    Take $C_1$ such that
                        $C_1$ is a component of $X$ with respect to $T$ and $D\subseteq C_1$.
                    Thus $z\in C_1$ by \cref{subseteq}.
                    Show that $C = C_1$.
                    \begin{subproof}
                        Suppose not.
                        Then $C\neq C_1$.
                        Then $C$ is disjoint from $C_1$ by \cref{components_disjoint}.
                        Thus $z\in C$ and $z\in C_1$.
                        Contradiction by \cref{disjoint}.
                    \end{subproof}
                    Thus $D\subseteq C$.
                    Show that $D\neq P$.
                    \begin{subproof}
                        Suppose not.
                        Then $D = P$.
                        Then $z\in P$.
                        Contradiction.
                    \end{subproof}
                    Thus $D\subseteq X$ by \cref{subseteq_transitive}.
                    Thus $D\in\pow{X}$ by \cref{powerset_intro}.
                    Thus $D\in M$.
                    Thus $z\in Q$ by \cref{unions_iff}.
                \end{subproof}
                Thus $C\setminus P\subseteq Q$ by \cref{subseteq}.
            \end{subproof}
            Thus $Q = C\setminus P$ by \cref{subseteq_antisymmetric}.
        \end{subproof}
        Show that $M\subseteq T$.
        \begin{subproof}
            Show that for all $D$ such that $D\in M$ we have $D\in T$.
            \begin{subproof}
                Fix $D$.
                Assume $D\in M$.
                Thus $D$ is a path component of $X$ with respect to $T$.
                Thus $D$ is open in $X$ with respect to $T$.
                Thus $D\in T$ by \cref{open_in}.
            \end{subproof}
            Thus $M\subseteq T$ by \cref{subseteq}.
        \end{subproof}
        Thus $Q\in T$ by \cref{topology_on,subseteq}.
        Thus $Q$ is open in $X$ with respect to $T$ by \cref{open_in}.
        Thus $P$ is open in $X$ with respect to $T$.
        Let $T_C = \subspaceTopology{T}{C}$.
        Thus $T_C$ is a topology on $C$ by \cref{subspace_topology_is_topology,subspace_of}.
        Show that $P$ is open in $C$ with respect to $T_C$.
        \begin{subproof}
            Thus $P\in T$ by \cref{open_in}.
            Thus $P = C\inter P$ by \cref{inter_absorb_supseteq_left}.
            Thus $P\in T_C$ by \cref{subspace_topology}.
            Thus $P$ is open in $C$ with respect to $T_C$ by \cref{open_in}.
        \end{subproof}
        Show that $Q$ is open in $C$ with respect to $T_C$.
        \begin{subproof}
            Thus $Q\in T$ by \cref{open_in}.
            Thus $Q = C\inter Q$ by \cref{inter_absorb_supseteq_left,setminus_subseteq}.
            Thus $Q\in T_C$ by \cref{subspace_topology}.
            Thus $Q$ is open in $C$ with respect to $T_C$ by \cref{open_in}.
        \end{subproof}
        Thus $x\in P$.
        Thus $P\neq\emptyset$ by \cref{emptyset}.
        Thus $y\in Q$.
        Thus $Q\neq\emptyset$ by \cref{emptyset}.
        Thus $P\inter Q = \emptyset$ by \cref{setminus_disjoint}.
        Show that $P$ is disjoint from $Q$.
        \begin{subproof}
            Suppose not.
            Then there exists $z$ such that $z\in P$ and $z\in Q$.
            Thus $z\in P\inter Q$ by \cref{inter}.
            Thus $P\inter Q\neq\emptyset$ by \cref{emptyset}.
            Contradiction by \cref{setminus_disjoint}.
        \end{subproof}
        Thus $P\union Q = C$ by \cref{setminus_partition,subseteq}.
        Thus $C$ is separated by $P$ and $Q$ with respect to $T_C$ by \cref{separation}.
        Thus $C$ is not connected with respect to $T_C$ by \cref{connected}.
        Contradiction by \cref{component_connected}.
        Thus $C = P$.
    \end{subproof}
    Thus there exists $P$ such that $P$ is a path component of $X$ with respect to $T$ and $C = P$.
\end{proof}


% from §26 Item 1: covering
\begin{definition}\label{covering}
    $A$ is a covering of $X$ iff $X\subseteq\unions{A}$.
\end{definition}

% from §26 Item 2: open covering
\begin{definition}\label{open_covering}
    $A$ is an open covering of $X$ with respect to $T$ iff
        $A$ is a covering of $X$ and
        for all $U$ such that $U\in A$ we have $U$ is open in $X$ with respect to $T$.
\end{definition}

% from §26 Item 3: compact space
\begin{definition}\label{compact}
    $X$ is compact with respect to $T$ iff
        for all $A$ such that $A$ is an open covering of $X$ with respect to $T$
        there exists $B$ such that
            $B\subseteq A$ and
            $B$ is finite and
            $B$ is a covering of $X$.
\end{definition}

% from §26 Item 4: compactness via coverings open in the ambient space
\begin{lemma}\label{compact_subspace_open_covering}
    Assume $Y$ is a subspace of $X$ with respect to $T$.
    Let $T_Y = \subspaceTopology{T}{Y}$.
    Then $Y$ is compact with respect to $T_Y$ iff
        for all $A$ such that
            $A$ is a covering of $Y$ and
            for all $U$ such that $U\in A$ we have $U$ is open in $X$ with respect to $T$
        there exists $B$ such that
            $B\subseteq A$ and
            $B$ is finite and
            $B$ is a covering of $Y$.
\end{lemma}
\begin{proof}
    Let $T_Y = \subspaceTopology{T}{Y}$.
    $T$ is a topology on $X$ by \cref{subspace_of}.
    Thus $Y\subseteq X$ by \cref{subspace_of}.
    Thus $T_Y$ is a topology on $Y$ by \cref{subspace_topology_is_topology}.
    Show that if $Y$ is compact with respect to $T_Y$ then
        for all $A$ such that
            $A$ is a covering of $Y$ and
            for all $U$ such that $U\in A$ we have $U$ is open in $X$ with respect to $T$
        there exists $B$ such that
            $B\subseteq A$ and
            $B$ is finite and
            $B$ is a covering of $Y$.
    \begin{subproof}
        Assume $Y$ is compact with respect to $T_Y$.
        Fix $A$.
        Assume $A$ is a covering of $Y$.
        Assume for all $U$ such that $U\in A$ we have $U$ is open in $X$ with respect to $T$.
        Let $A_Y = \{V \mid \exists U\in A. V = Y\inter U\}$.
        Show that $A_Y$ is an open covering of $Y$ with respect to $T_Y$.
        \begin{subproof}
            Show that for all $V$ such that $V\in A_Y$ we have $V$ is open in $Y$ with respect to $T_Y$.
            \begin{subproof}
                Fix $V$.
                Assume $V\in A_Y$.
                Take $U\in A$ such that $V = Y\inter U$.
                $U$ is open in $X$ with respect to $T$.
                Thus $U\in T$ by \cref{open_in}.
                Thus $V\in T_Y$ by \cref{subspace_topology}.
                Thus $V$ is open in $Y$ with respect to $T_Y$ by \cref{open_in}.
            \end{subproof}
            Show that $A_Y$ is a covering of $Y$.
            \begin{subproof}
                Show that for all $y$ such that $y\in Y$ we have $y\in\unions{A_Y}$.
                \begin{subproof}
                    Fix $y$.
                    Assume $y\in Y$.
                    $Y\subseteq\unions{A}$ by \cref{covering}.
                    Thus $y\in\unions{A}$ by \cref{subseteq}.
                    Take $U\in A$ such that $y\in U$ by \cref{unions_iff}.
                    Thus $y\in Y\inter U$ by \cref{inter_intro}.
                    Thus $Y\inter U\in A_Y$.
                    Thus $y\in\unions{A_Y}$ by \cref{unions_iff}.
                \end{subproof}
                Thus $Y\subseteq\unions{A_Y}$ by \cref{subseteq}.
                Thus $A_Y$ is a covering of $Y$ by \cref{covering}.
            \end{subproof}
            Thus $A_Y$ is an open covering of $Y$ with respect to $T_Y$ by \cref{open_covering}.
        \end{subproof}
        Thus there exists $B_Y$ such that
            $B_Y\subseteq A_Y$ and $B_Y$ is finite and $B_Y$ is a covering of $Y$
            by \cref{compact}.
        Take $B_Y$ such that
            $B_Y\subseteq A_Y$ and $B_Y$ is finite and $B_Y$ is a covering of $Y$.
        Let $R = \{w\in B_Y\times A \mid \exists V\in B_Y.\exists U\in A. w = (V,U) \land V = Y\inter U\}$.
        Show that $R$ is a relation.
        \begin{subproof}
            Show that for all $w$ such that $w\in R$ there exist $V,U$ such that $w = (V,U)$.
            \begin{subproof}
                Fix $w$.
                Assume $w\in R$.
                Take $V\in B_Y$ such that there exists $U\in A$ such that $w = (V,U)$ and $V = Y\inter U$.
                Take $U\in A$ such that $w = (V,U)$ and $V = Y\inter U$.
            \end{subproof}
            Thus $R$ is a relation by \cref{relation}.
        \end{subproof}
        Show that for all $V\in B_Y$ there exists $U$ such that $V\mathrel{R} U$.
        \begin{subproof}
            Fix $V$.
            Assume $V\in B_Y$.
            Thus $V\in A_Y$ by \cref{subseteq}.
            Take $U\in A$ such that $V = Y\inter U$.
            Thus $(V,U)\in R$.
            Thus $V\mathrel{R} U$.
        \end{subproof}
        Take $g$ such that $g$ is a function on $B_Y$ and for all $V\in B_Y$ we have $V\mathrel{R}\apply{g}{V}$ by \cref{choice_function}.
        Let $B = \img{g}{B_Y}$.
        Show that $B\subseteq A$.
        \begin{subproof}
            Show that for all $U$ such that $U\in B$ we have $U\in A$.
            \begin{subproof}
                Fix $U$.
                Assume $U\in B$.
                Take $V\in B_Y$ such that $(V,U)\in g$ by \cref{img_iff}.
                Thus $V\mathrel{R} U$ by \cref{function_apply_iff}.
                Take $V_1\in B_Y$ such that there exists $U_1\in A$ such that
                    $(V,U) = (V_1,U_1)$ and $V_1 = Y\inter U_1$.
                Take $U_1\in A$ such that $(V,U) = (V_1,U_1)$ and $V_1 = Y\inter U_1$.
                Thus $U = U_1$ by \cref{pair_eq_iff}.
                Thus $U\in A$.
            \end{subproof}
            Thus $B\subseteq A$ by \cref{subseteq}.
        \end{subproof}
        $B$ is finite by \cref{finite_image_function}.
        Show that $B$ is a covering of $Y$.
        \begin{subproof}
            Show that for all $y$ such that $y\in Y$ we have $y\in\unions{B}$.
            \begin{subproof}
                Fix $y$.
                Assume $y\in Y$.
                $Y\subseteq\unions{B_Y}$ by \cref{covering}.
                Thus $y\in\unions{B_Y}$ by \cref{subseteq}.
                Take $V\in B_Y$ such that $y\in V$ by \cref{unions_iff}.
                Thus $V\mathrel{R}\apply{g}{V}$ by \cref{subseteq}.
                Take $V_1\in B_Y$ such that there exists $U_1\in A$ such that
                    $(V,\apply{g}{V}) = (V_1,U_1)$ and $V_1 = Y\inter U_1$.
                Take $U_1\in A$ such that $(V,\apply{g}{V}) = (V_1,U_1)$ and $V_1 = Y\inter U_1$.
                Thus $V = V_1$ and $\apply{g}{V} = U_1$ by \cref{pair_eq_iff}.
                Thus $y\in Y\inter \apply{g}{V}$ by \cref{inter_intro,subseteq}.
                Thus $y\in \apply{g}{V}$ by \cref{inter_elim_right}.
                Thus $\dom{g} = B_Y$.
                Thus $V\in\dom{g}$ by \cref{subseteq}.
                Thus $(V,\apply{g}{V})\in g$ by \cref{function_apply_elim}.
                Thus $\apply{g}{V}\in B$ by \cref{img_iff}.
                Thus $y\in\unions{B}$ by \cref{unions_iff}.
            \end{subproof}
            Thus $Y\subseteq\unions{B}$ by \cref{subseteq}.
            Thus $B$ is a covering of $Y$ by \cref{covering}.
        \end{subproof}
    \end{subproof}
    Show that if for all $A$ such that
            $A$ is a covering of $Y$ and
            for all $U$ such that $U\in A$ we have $U$ is open in $X$ with respect to $T$
        there exists $B$ such that
            $B\subseteq A$ and
            $B$ is finite and
            $B$ is a covering of $Y$
        then $Y$ is compact with respect to $T_Y$.
    \begin{subproof}
        Assume for all $A$ such that
                $A$ is a covering of $Y$ and
                for all $U$ such that $U\in A$ we have $U$ is open in $X$ with respect to $T$
            there exists $B$ such that
                $B\subseteq A$ and
                $B$ is finite and
                $B$ is a covering of $Y$.
        Show that for all $A$ such that $A$ is an open covering of $Y$ with respect to $T_Y$
            there exists $B$ such that $B\subseteq A$ and $B$ is finite and $B$ is a covering of $Y$.
        \begin{subproof}
            Fix $A$.
            Assume $A$ is an open covering of $Y$ with respect to $T_Y$.
            Let $R = \{w\in A\times T \mid \exists V\in A.\exists U\in T. w = (V,U) \land V = Y\inter U\}$.
            Show that $R$ is a relation.
            \begin{subproof}
                Show that for all $w$ such that $w\in R$ there exist $V,U$ such that $w = (V,U)$.
                \begin{subproof}
                    Fix $w$.
                    Assume $w\in R$.
                    Take $V\in A$ such that there exists $U\in T$ such that $w = (V,U)$ and $V = Y\inter U$.
                    Take $U\in T$ such that $w = (V,U)$ and $V = Y\inter U$.
                \end{subproof}
                Thus $R$ is a relation by \cref{relation}.
            \end{subproof}
            Show that for all $V\in A$ there exists $U$ such that $V\mathrel{R} U$.
            \begin{subproof}
                Fix $V$.
                Assume $V\in A$.
                $V$ is open in $Y$ with respect to $T_Y$ by \cref{open_covering}.
                Thus $V\in T_Y$ by \cref{open_in}.
                Take $U\in T$ such that $V = Y\inter U$ by \cref{subspace_topology}.
                Thus $(V,U)\in R$.
                Thus $V\mathrel{R} U$.
            \end{subproof}
            Take $g$ such that $g$ is a function on $A$ and for all $V\in A$ we have $V\mathrel{R}\apply{g}{V}$ by \cref{choice_function}.
            Let $C = \img{g}{A}$.
            Show that for all $U$ such that $U\in C$ we have $U$ is open in $X$ with respect to $T$.
            \begin{subproof}
                Fix $U$.
                Assume $U\in C$.
                Take $V\in A$ such that $(V,U)\in g$ by \cref{img_iff}.
                Thus $V\mathrel{R} U$ by \cref{function_apply_iff}.
                Take $V_1\in A$ such that there exists $U_1\in T$ such that
                    $(V,U) = (V_1,U_1)$ and $V_1 = Y\inter U_1$.
                Take $U_1\in T$ such that $(V,U) = (V_1,U_1)$ and $V_1 = Y\inter U_1$.
                Thus $U = U_1$ by \cref{pair_eq_iff}.
                Thus $U\in T$.
                Thus $U$ is open in $X$ with respect to $T$ by \cref{open_in}.
            \end{subproof}
            Show that $C$ is a covering of $Y$.
            \begin{subproof}
                Show that for all $y$ such that $y\in Y$ we have $y\in\unions{C}$.
                \begin{subproof}
                    Fix $y$.
                    Assume $y\in Y$.
                    $Y\subseteq\unions{A}$ by \cref{open_covering,covering}.
                    Thus $y\in\unions{A}$ by \cref{subseteq}.
                    Take $V\in A$ such that $y\in V$ by \cref{unions_iff}.
                    Thus $V\mathrel{R}\apply{g}{V}$ by \cref{subseteq}.
                    Take $V_1\in A$ such that there exists $U_1\in T$ such that
                        $(V,\apply{g}{V}) = (V_1,U_1)$ and $V_1 = Y\inter U_1$.
                    Take $U_1\in T$ such that $(V,\apply{g}{V}) = (V_1,U_1)$ and $V_1 = Y\inter U_1$.
                    Thus $V = V_1$ and $\apply{g}{V} = U_1$ by \cref{pair_eq_iff}.
                    Thus $y\in Y\inter \apply{g}{V}$ by \cref{inter_intro,subseteq}.
                    Thus $y\in \apply{g}{V}$ by \cref{inter_elim_right}.
                    Thus $\dom{g} = A$.
                    Thus $V\in\dom{g}$ by \cref{subseteq}.
                    Thus $(V,\apply{g}{V})\in g$ by \cref{function_apply_elim}.
                    Thus $\apply{g}{V}\in C$ by \cref{img_iff}.
                    Thus $y\in\unions{C}$ by \cref{unions_iff}.
                \end{subproof}
                Thus $Y\subseteq\unions{C}$ by \cref{subseteq}.
                Thus $C$ is a covering of $Y$ by \cref{covering}.
            \end{subproof}
            Thus there exists $B$ such that $B\subseteq C$ and $B$ is finite and $B$ is a covering of $Y$.
            Take $B$ such that $B\subseteq C$ and $B$ is finite and $B$ is a covering of $Y$.
            Let $S = \{w\in B\times A \mid \exists U\in B.\exists V\in A. w = (U,V) \land \apply{g}{V} = U\}$.
            Show that $S$ is a relation.
            \begin{subproof}
                Show that for all $w$ such that $w\in S$ there exist $U,V$ such that $w = (U,V)$.
                \begin{subproof}
                    Fix $w$.
                    Assume $w\in S$.
                    Take $U\in B$ such that there exists $V\in A$ such that $w = (U,V)$ and $\apply{g}{V} = U$.
                    Take $V\in A$ such that $w = (U,V)$ and $\apply{g}{V} = U$.
                \end{subproof}
                Thus $S$ is a relation by \cref{relation}.
            \end{subproof}
            Show that for all $U\in B$ there exists $V$ such that $U\mathrel{S} V$.
            \begin{subproof}
                Fix $U$.
                Assume $U\in B$.
                Thus $U\in C$ by \cref{subseteq}.
                Take $V\in A$ such that $(V,U)\in g$ by \cref{img_iff}.
                Thus $U = \apply{g}{V}$ by \cref{function_apply_intro}.
                Thus $(U,V)\in S$.
                Thus $U\mathrel{S} V$.
            \end{subproof}
            Take $h$ such that $h$ is a function on $B$ and for all $U\in B$ we have $U\mathrel{S}\apply{h}{U}$ by \cref{choice_function}.
            Let $B_1 = \img{h}{B}$.
            Show that $B_1\subseteq A$.
            \begin{subproof}
                Show that for all $V$ such that $V\in B_1$ we have $V\in A$.
                \begin{subproof}
                    Fix $V$.
                    Assume $V\in B_1$.
                    Take $U\in B$ such that $(U,V)\in h$ by \cref{img_iff}.
                    Thus $U\mathrel{S} V$ by \cref{function_apply_iff}.
                    Take $U_1\in B$ such that there exists $V_1\in A$ such that
                        $(U,V) = (U_1,V_1)$ and $\apply{g}{V_1} = U_1$.
                    Take $V_1\in A$ such that $(U,V) = (U_1,V_1)$ and $\apply{g}{V_1} = U_1$.
                    Thus $V = V_1$ by \cref{pair_eq_iff}.
                    Thus $V\in A$.
                \end{subproof}
                Thus $B_1\subseteq A$ by \cref{subseteq}.
            \end{subproof}
            $B_1$ is finite by \cref{finite_image_function}.
            Show that $B_1$ is a covering of $Y$.
            \begin{subproof}
                Show that for all $y$ such that $y\in Y$ we have $y\in\unions{B_1}$.
                \begin{subproof}
                    Fix $y$.
                    Assume $y\in Y$.
                    $Y\subseteq\unions{B}$ by \cref{covering}.
                    Thus $y\in\unions{B}$ by \cref{subseteq}.
                    Take $U\in B$ such that $y\in U$ by \cref{unions_iff}.
                    Thus $U\mathrel{S}\apply{h}{U}$ by \cref{subseteq}.
                    Take $U_1\in B$ such that there exists $V_1\in A$ such that
                        $(U,\apply{h}{U}) = (U_1,V_1)$ and $\apply{g}{V_1} = U_1$.
                    Take $V_1\in A$ such that $(U,\apply{h}{U}) = (U_1,V_1)$ and $\apply{g}{V_1} = U_1$.
                    Thus $U = U_1$ and $\apply{h}{U} = V_1$ by \cref{pair_eq_iff}.
                    Thus $y\in Y\inter \apply{g}{V_1}$ by \cref{inter_intro,subseteq}.
                    Thus $V_1\in A$.
                    Thus $V_1\mathrel{R}\apply{g}{V_1}$ by \cref{subseteq}.
                    Take $V_2\in A$ such that there exists $U_2\in T$ such that
                        $(V_1,\apply{g}{V_1}) = (V_2,U_2)$ and $V_2 = Y\inter U_2$.
                    Take $U_2\in T$ such that $(V_1,\apply{g}{V_1}) = (V_2,U_2)$ and $V_2 = Y\inter U_2$.
                    Thus $V_1 = V_2$ and $\apply{g}{V_1} = U_2$ by \cref{pair_eq_iff}.
                    Thus $V_1 = Y\inter \apply{g}{V_1}$.
                    Thus $y\in V_1$.
                    Thus $\dom{h} = B$.
                    Thus $U\in\dom{h}$ by \cref{subseteq}.
                    Thus $(U,\apply{h}{U})\in h$ by \cref{function_apply_elim}.
                    Thus $V_1\in B_1$ by \cref{img_iff}.
                    Thus $y\in\unions{B_1}$ by \cref{unions_iff}.
                \end{subproof}
                Thus $Y\subseteq\unions{B_1}$ by \cref{subseteq}.
                Thus $B_1$ is a covering of $Y$ by \cref{covering}.
            \end{subproof}
            Thus there exists $B$ such that $B\subseteq A$ and $B$ is finite and $B$ is a covering of $Y$.
        \end{subproof}
        Thus $Y$ is compact with respect to $T_Y$ by \cref{compact}.
    \end{subproof}
\end{proof}

% from §26 Item 5: closed subspace of a compact space is compact
\begin{theorem}\label{closed_subspace_compact}
    Assume $T$ is a topology on $X$.
    Assume $X$ is compact with respect to $T$.
    Suppose $Y\subseteq X$.
    Suppose $Y$ is closed in $X$ with respect to $T$.
    Let $T_Y = \subspaceTopology{T}{Y}$.
    Then $Y$ is compact with respect to $T_Y$.
\end{theorem}
\begin{proof}
    $Y$ is a subspace of $X$ with respect to $T$ by \cref{subspace_of}.
    Let $T_Y = \subspaceTopology{T}{Y}$.
    Show that for all $A$ such that
        $A$ is a covering of $Y$ and
        for all $U$ such that $U\in A$ we have $U$ is open in $X$ with respect to $T$
        there exists $B$ such that
            $B\subseteq A$ and
            $B$ is finite and
            $B$ is a covering of $Y$.
    \begin{subproof}
        Fix $A$.
        Assume $A$ is a covering of $Y$.
        Assume for all $U$ such that $U\in A$ we have $U$ is open in $X$ with respect to $T$.
        Let $C = A\union \{X\setminus Y\}$.
        Show that $C$ is an open covering of $X$ with respect to $T$.
        \begin{subproof}
            Show that for all $U$ such that $U\in C$ we have $U$ is open in $X$ with respect to $T$.
            \begin{subproof}
                Fix $U$.
                Assume $U\in C$.
                \begin{byCase}
                \caseOf{$U\in A$.}
                    Thus $U$ is open in $X$ with respect to $T$.
                \caseOf{$U\notin A$.}
                    Thus $U\in \{X\setminus Y\}$ by \cref{union_iff}.
                    Thus $U = X\setminus Y$ by \cref{singleton_iff}.
                    $Y$ is closed in $X$ with respect to $T$.
                    Thus $X\setminus Y$ is open in $X$ with respect to $T$ by \cref{closed_in}.
                    Thus $U$ is open in $X$ with respect to $T$.
                \end{byCase}
            \end{subproof}
            Show that $C$ is a covering of $X$.
            \begin{subproof}
                Show that for all $x$ such that $x\in X$ we have $x\in\unions{C}$.
                \begin{subproof}
                    Fix $x$.
                    Assume $x\in X$.
                    \begin{byCase}
                    \caseOf{$x\in Y$.}
                        $Y\subseteq\unions{A}$ by \cref{covering}.
                        Thus $x\in\unions{A}$ by \cref{subseteq}.
                        Take $U\in A$ such that $x\in U$ by \cref{unions_iff}.
                        Thus $U\in C$ by \cref{union_iff}.
                        Thus $x\in\unions{C}$ by \cref{unions_iff}.
                    \caseOf{$x\notin Y$.}
                        Thus $x\in X\setminus Y$ by \cref{setminus_intro}.
                        Thus $X\setminus Y\in \{X\setminus Y\}$ by \cref{singleton_iff}.
                        Thus $X\setminus Y\in C$ by \cref{union_iff}.
                        Thus $x\in\unions{C}$ by \cref{unions_iff}.
                    \end{byCase}
                \end{subproof}
                Thus $X\subseteq\unions{C}$ by \cref{subseteq}.
                Thus $C$ is a covering of $X$ by \cref{covering}.
            \end{subproof}
            Thus $C$ is an open covering of $X$ with respect to $T$ by \cref{open_covering}.
        \end{subproof}
        Thus there exists $B$ such that $B\subseteq C$ and $B$ is finite and $B$ is a covering of $X$ by \cref{compact}.
        Take $B$ such that $B\subseteq C$ and $B$ is finite and $B$ is a covering of $X$.
        \begin{byCase}
        \caseOf{$X\setminus Y\in B$.}
            Let $B_1 = B\setminus \{X\setminus Y\}$.
            Show that $B_1$ is a covering of $Y$.
            \begin{subproof}
                Show that for all $y$ such that $y\in Y$ we have $y\in\unions{B_1}$.
                \begin{subproof}
                    Fix $y$.
                    Assume $y\in Y$.
                    $X\subseteq\unions{B}$ by \cref{covering}.
                    Thus $y\in\unions{B}$ by \cref{subseteq}.
                    Take $U\in B$ such that $y\in U$ by \cref{unions_iff}.
                    Show that $U\in B_1$.
                    \begin{subproof}
                        Suppose not.
                        Show that $U\in \{X\setminus Y\}$.
                        \begin{subproof}
                            Suppose not.
                            Thus $U\notin \{X\setminus Y\}$.
                            Thus $U\in B\setminus \{X\setminus Y\}$ by \cref{setminus_intro}.
                            Thus $U\in B_1$.
                            Contradiction.
                        \end{subproof}
                        Thus $U = X\setminus Y$ by \cref{singleton_iff}.
                        Thus $y\notin U$ by \cref{setminus_elim_right}.
                        Contradiction.
                    \end{subproof}
                    Thus $y\in\unions{B_1}$ by \cref{unions_iff}.
                \end{subproof}
                Thus $Y\subseteq\unions{B_1}$ by \cref{subseteq}.
                Thus $B_1$ is a covering of $Y$ by \cref{covering}.
            \end{subproof}
            $B_1\subseteq B$.
            Thus $B_1$ is finite by \cref{subseteq_finite_is_finite}.
            Show that $B_1\subseteq A$.
            \begin{subproof}
                Show that for all $U$ such that $U\in B_1$ we have $U\in A$.
                \begin{subproof}
                    Fix $U$.
                    Assume $U\in B_1$.
                    Thus $U\in B$ by \cref{subseteq}.
                    Thus $U\in C$ by \cref{subseteq}.
                    Thus $U\notin \{X\setminus Y\}$ by \cref{setminus_elim_right}.
                    Thus $U\in A$ by \cref{union_iff}.
                \end{subproof}
                Thus $B_1\subseteq A$ by \cref{subseteq}.
            \end{subproof}
            Thus there exists $B$ such that $B\subseteq A$ and $B$ is finite and $B$ is a covering of $Y$.
        \caseOf{$X\setminus Y\notin B$.}
            Show that $B\subseteq A$.
            \begin{subproof}
                Show that for all $U$ such that $U\in B$ we have $U\in A$.
                \begin{subproof}
                    Fix $U$.
                    Assume $U\in B$.
                    Thus $U\in C$ by \cref{subseteq}.
                    Show that $U\notin \{X\setminus Y\}$.
                    \begin{subproof}
                        Suppose not.
                        Then $U = X\setminus Y$ by \cref{singleton_iff}.
                        Thus $X\setminus Y\in B$.
                        Contradiction.
                    \end{subproof}
                    Thus $U\in A$ by \cref{union_iff}.
                \end{subproof}
                Thus $B\subseteq A$ by \cref{subseteq}.
            \end{subproof}
            Show that $B$ is a covering of $Y$.
            \begin{subproof}
                Show that for all $y$ such that $y\in Y$ we have $y\in\unions{B}$.
                \begin{subproof}
                    Fix $y$.
                    Assume $y\in Y$.
                    $X\subseteq\unions{B}$ by \cref{covering}.
                    Thus $y\in\unions{B}$ by \cref{subseteq}.
                \end{subproof}
                Thus $Y\subseteq\unions{B}$ by \cref{subseteq}.
                Thus $B$ is a covering of $Y$ by \cref{covering}.
            \end{subproof}
            Thus there exists $B$ such that $B\subseteq A$ and $B$ is finite and $B$ is a covering of $Y$.
        \end{byCase}
    \end{subproof}
    Thus $Y$ is compact with respect to $T_Y$ by \cref{compact_subspace_open_covering}.
\end{proof}

% from §26 Item 6: compact subspace of Hausdorff is closed
\begin{theorem}\label{compact_subspace_closed}
    Assume $T$ is a topology on $X$.
    Assume $X$ is Hausdorff with respect to $T$.
    Suppose $Y\subseteq X$.
    Let $T_Y = \subspaceTopology{T}{Y}$.
    Suppose $Y$ is compact with respect to $T_Y$.
    Then $Y$ is closed in $X$ with respect to $T$.
\end{theorem}
\begin{proof}
    Let $T_Y = \subspaceTopology{T}{Y}$.
    \begin{byCase}
    \caseOf{$Y = \emptyset$.}
        Thus $Y$ is closed in $X$ with respect to $T$ by \cref{closed_emptyset}.
    \caseOf{$Y \neq \emptyset$.}
        $Y$ is inhabited by \cref{nonempty_inhabited}.
        Show that $X\setminus Y$ is open in $X$ with respect to $T$.
        \begin{subproof}
            Let $M = \{U\in T \mid \text{$U$ is disjoint from $Y$}\}$.
            Show that $X\setminus Y = \unions{M}$.
            \begin{subproof}
                Show that $X\setminus Y\subseteq \unions{M}$.
                \begin{subproof}
                    Show that for all $x_0$ such that $x_0\in X\setminus Y$ we have $x_0\in\unions{M}$.
                    \begin{subproof}
                        Fix $x_0$.
                        Assume $x_0\in X\setminus Y$.
                        Thus $x_0\in X$ and $x_0\notin Y$ by \cref{setminus_elim_left,setminus_elim_right}.
                        Let $A = \{V\in T \mid \text{there exists $y\in Y$ such that there exists $U$ such that
                                $U$ is a neighborhood of $x_0$ in $X$ with respect to $T$ and
                                $V$ is a neighborhood of $y$ in $X$ with respect to $T$ and
                                $U$ is disjoint from $V$}\}$.
                        Show that $A$ is a covering of $Y$.
                        \begin{subproof}
                            Show that for all $y$ such that $y\in Y$ we have $y\in\unions{A}$.
                            \begin{subproof}
                                Fix $y$.
                                Assume $y\in Y$.
                                Thus $y\in X$ by \cref{subseteq}.
                                Show that $x_0 \neq y$.
                                \begin{subproof}
                                    Suppose not.
                                    Then $x_0 = y$.
                                    Thus $x_0\in Y$.
                                    Contradiction.
                                \end{subproof}
                                Thus there exist $U_1,U_2$ such that
                                    $U_1$ is a neighborhood of $x_0$ in $X$ with respect to $T$ and
                                    $U_2$ is a neighborhood of $y$ in $X$ with respect to $T$ and
                                    $U_1$ is disjoint from $U_2$
                                    by \cref{hausdorff}.
                                Take $U_1,U_2$ such that
                                    $U_1$ is a neighborhood of $x_0$ in $X$ with respect to $T$ and
                                    $U_2$ is a neighborhood of $y$ in $X$ with respect to $T$ and
                                    $U_1$ is disjoint from $U_2$.
                                Thus $U_2$ is open in $X$ with respect to $T$ by \cref{neighborhood}.
                                Thus $U_2\in T$ by \cref{open_in}.
                                Thus $U_2\in A$.
                                Thus $y\in U_2$ by \cref{neighborhood}.
                                Thus $y\in\unions{A}$ by \cref{unions_iff}.
                            \end{subproof}
                            Thus $Y\subseteq\unions{A}$ by \cref{subseteq}.
                            Thus $A$ is a covering of $Y$ by \cref{covering}.
                        \end{subproof}
                        Show that for all $U$ such that $U\in A$ we have $U$ is open in $X$ with respect to $T$.
                        \begin{subproof}
                            Fix $U$.
                            Assume $U\in A$.
                            Thus $U\in T$ by \cref{subseteq}.
                            Thus $U$ is open in $X$ with respect to $T$ by \cref{open_in}.
                        \end{subproof}
                        $Y$ is a subspace of $X$ with respect to $T$ by \cref{subspace_of}.
                        Thus there exists $B$ such that
                            $B\subseteq A$ and $B$ is finite and $B$ is a covering of $Y$
                            by \cref{compact_subspace_open_covering,subseteq}.
                        Take $B$ such that $B\subseteq A$ and $B$ is finite and $B$ is a covering of $Y$.
                        Show that $B$ is inhabited.
                        \begin{subproof}
                            Take $y$ such that $y\in Y$ by \cref{nonempty_inhabited}.
                            $Y\subseteq\unions{B}$ by \cref{covering}.
                            Thus $y\in\unions{B}$ by \cref{subseteq}.
                            Take $V\in B$ such that $y\in V$ by \cref{unions_iff}.
                            Thus $B$ is inhabited.
                        \end{subproof}
                        Let $R = \{w\in B\times T \mid \text{there exists $V\in B$ such that there exists $U\in T$ such that
                                $w = (V,U)$ and
                                $U$ is a neighborhood of $x_0$ in $X$ with respect to $T$ and
                                $U$ is disjoint from $V$}\}$.
                        Show that $R$ is a relation.
                        \begin{subproof}
                            Show that for all $w$ such that $w\in R$ there exist $V,U$ such that $w = (V,U)$.
                            \begin{subproof}
                                Fix $w$.
                                Assume $w\in R$.
                                Take $V\in B$ such that there exists $U\in T$ such that
                                    $w = (V,U)$ and
                                    $U$ is a neighborhood of $x_0$ in $X$ with respect to $T$ and
                                    $U$ is disjoint from $V$.
                                Take $U\in T$ such that
                                    $w = (V,U)$ and
                                    $U$ is a neighborhood of $x_0$ in $X$ with respect to $T$ and
                                    $U$ is disjoint from $V$.
                            \end{subproof}
                            Thus $R$ is a relation by \cref{relation}.
                        \end{subproof}
                        Show that for all $V\in B$ there exists $U$ such that $V\mathrel{R} U$.
                        \begin{subproof}
                            Fix $V$.
                            Assume $V\in B$.
                            Thus $V\in A$ by \cref{subseteq}.
                            Take $y\in Y$ such that there exists $U$ such that
                                $U$ is a neighborhood of $x_0$ in $X$ with respect to $T$ and
                                $V$ is a neighborhood of $y$ in $X$ with respect to $T$ and
                                $U$ is disjoint from $V$.
                            Take $U$ such that
                                $U$ is a neighborhood of $x_0$ in $X$ with respect to $T$ and
                                $V$ is a neighborhood of $y$ in $X$ with respect to $T$ and
                                $U$ is disjoint from $V$.
                            Thus $U$ is open in $X$ with respect to $T$ by \cref{neighborhood}.
                            Thus $U\in T$ by \cref{open_in}.
                            Thus $(V,U)\in R$.
                            Thus $V\mathrel{R} U$.
                        \end{subproof}
                        Take $g$ such that $g$ is a function on $B$ and for all $V\in B$ we have $V\mathrel{R}\apply{g}{V}$ by \cref{choice_function}.
                        Let $G = \img{g}{B}$.
                        Show that $G\subseteq T$.
                        \begin{subproof}
                            Show that for all $U$ such that $U\in G$ we have $U\in T$.
                            \begin{subproof}
                                Fix $U$.
                                Assume $U\in G$.
                                Take $V\in B$ such that $(V,U)\in g$ by \cref{img_iff}.
                                Thus $V\mathrel{R} U$ by \cref{function_apply_iff}.
                                Take $V_1\in B$ such that there exists $U_1\in T$ such that
                                    $(V,U) = (V_1,U_1)$ and
                                    $U_1$ is a neighborhood of $x_0$ in $X$ with respect to $T$ and
                                    $U_1$ is disjoint from $V_1$.
                                Take $U_1\in T$ such that
                                    $(V,U) = (V_1,U_1)$ and
                                    $U_1$ is a neighborhood of $x_0$ in $X$ with respect to $T$ and
                                    $U_1$ is disjoint from $V_1$.
                                Thus $U = U_1$ by \cref{pair_eq_iff}.
                                Thus $U\in T$.
                            \end{subproof}
                            Thus $G\subseteq T$ by \cref{subseteq}.
                        \end{subproof}
                        $G$ is finite by \cref{finite_image_function}.
                        Show that $G$ is inhabited.
                        \begin{subproof}
                            Take $y$ such that $y\in Y$ by \cref{nonempty_inhabited}.
                            $Y\subseteq\unions{B}$ by \cref{covering}.
                            Thus $y\in\unions{B}$ by \cref{subseteq}.
                            Take $V\in B$ such that $y\in V$ by \cref{unions_iff}.
                            Thus $(V,\apply{g}{V})\in g$ by \cref{function_apply_elim,subseteq}.
                            Thus $\apply{g}{V}\in G$ by \cref{img_iff}.
                            Thus $G$ is inhabited.
                        \end{subproof}
                        Let $U_0 = \inters{G}$.
                        $T$ is a topology on $X$ by \cref{hausdorff}.
                        Thus $U_0\in T$ by \cref{finite_intersections_open_in_topology}.
                        Show that $x_0\in U_0$.
                        \begin{subproof}
                            Show that for all $U$ such that $U\in G$ we have $x_0\in U$.
                            \begin{subproof}
                                Fix $U$.
                                Assume $U\in G$.
                                Take $V\in B$ such that $(V,U)\in g$ by \cref{img_iff}.
                                Thus $V\mathrel{R} U$ by \cref{function_apply_iff}.
                                Take $V_1\in B$ such that there exists $U_1\in T$ such that
                                    $(V,U) = (V_1,U_1)$ and
                                    $U_1$ is a neighborhood of $x_0$ in $X$ with respect to $T$ and
                                    $U_1$ is disjoint from $V_1$.
                                Take $U_1\in T$ such that
                                    $(V,U) = (V_1,U_1)$ and
                                    $U_1$ is a neighborhood of $x_0$ in $X$ with respect to $T$ and
                                    $U_1$ is disjoint from $V_1$.
                                Thus $U = U_1$ by \cref{pair_eq_iff}.
                                Thus $x_0\in U$ by \cref{neighborhood}.
                            \end{subproof}
                            Thus $x_0\in U_0$ by \cref{inters_intro}.
                        \end{subproof}
                        Show that $U_0$ is disjoint from $Y$.
                        \begin{subproof}
                            Suppose not.
                            Then there exists $z$ such that $z\in U_0$ and $z\in Y$.
                            $Y\subseteq\unions{B}$ by \cref{covering}.
                            Thus $z\in\unions{B}$ by \cref{subseteq}.
                            Take $V\in B$ such that $z\in V$ by \cref{unions_iff}.
                            Thus $(V,\apply{g}{V})\in g$ by \cref{function_apply_elim,subseteq}.
                            Thus $\apply{g}{V}\in G$ by \cref{img_iff}.
                            Thus $U_0\subseteq \apply{g}{V}$ by \cref{inters_subseteq_elem}.
                            Thus $z\in \apply{g}{V}$ by \cref{subseteq}.
                            Thus $V\mathrel{R}\apply{g}{V}$ by \cref{subseteq}.
                            Take $V_1\in B$ such that there exists $U_1\in T$ such that
                                $(V,\apply{g}{V}) = (V_1,U_1)$ and
                                $U_1$ is a neighborhood of $x_0$ in $X$ with respect to $T$ and
                                $U_1$ is disjoint from $V_1$.
                            Take $U_1\in T$ such that
                                $(V,\apply{g}{V}) = (V_1,U_1)$ and
                                $U_1$ is a neighborhood of $x_0$ in $X$ with respect to $T$ and
                                $U_1$ is disjoint from $V_1$.
                            Thus $V = V_1$ and $\apply{g}{V} = U_1$ by \cref{pair_eq_iff}.
                            Thus $\apply{g}{V}$ is disjoint from $V$.
                            Thus $z\notin \apply{g}{V}$ by \cref{disjoint}.
                            Contradiction.
                        \end{subproof}
                        Thus $U_0$ is disjoint from $Y$.
                        Thus $U_0\in M$.
                        Thus $x_0\in\unions{M}$ by \cref{unions_iff}.
                    \end{subproof}
                    Thus $X\setminus Y\subseteq \unions{M}$ by \cref{subseteq}.
                \end{subproof}
                Show that $\unions{M}\subseteq X\setminus Y$.
                \begin{subproof}
                    Show that for all $x$ such that $x\in\unions{M}$ we have $x\in X\setminus Y$.
                    \begin{subproof}
                        Fix $x$.
                        Assume $x\in\unions{M}$.
                        Take $U\in M$ such that $x\in U$ by \cref{unions_iff}.
                        Thus $U$ is disjoint from $Y$.
                        Thus $x\notin Y$ by \cref{disjoint}.
                        $T\subseteq\pow{X}$ by \cref{topology_on}.
                        Thus $U\in\pow{X}$ by \cref{subseteq}.
                        Thus $x\in X$ by \cref{powerset_elim,subseteq}.
                        Thus $x\in X\setminus Y$ by \cref{setminus_intro}.
                    \end{subproof}
                    Thus $\unions{M}\subseteq X\setminus Y$ by \cref{subseteq}.
                \end{subproof}
                Thus $X\setminus Y = \unions{M}$ by \cref{subseteq_antisymmetric}.
            \end{subproof}
            Thus $M\subseteq T$ by \cref{subseteq}.
            Thus $X\setminus Y\in T$ by \cref{topology_on}.
            Thus $X\setminus Y$ is open in $X$ with respect to $T$ by \cref{open_in}.
        \end{subproof}
        Thus $Y$ is closed in $X$ with respect to $T$ by \cref{closed_in}.
    \end{byCase}
\end{proof}

% from §26 Item 7: separation of a point and a compact set in a Hausdorff space
\begin{lemma}\label{compact_hausdorff_separation}
    Assume $T$ is a topology on $X$.
    Assume $X$ is Hausdorff with respect to $T$.
    Suppose $Y\subseteq X$.
    Let $T_Y = \subspaceTopology{T}{Y}$.
    Suppose $Y$ is compact with respect to $T_Y$.
    Suppose $x_0\in X\setminus Y$.
    Then there exist $U,V$ such that
        $U$ is open in $X$ with respect to $T$ and
        $V$ is open in $X$ with respect to $T$ and
        $x_0\in U$ and
        $Y\subseteq V$ and
        $U$ is disjoint from $V$.
\end{lemma}
\begin{proof}
    Let $T_Y = \subspaceTopology{T}{Y}$.
    \begin{byCase}
    \caseOf{$Y = \emptyset$.}
        Let $U = X$.
        Let $V = \emptyset$.
        $U$ is open in $X$ with respect to $T$ by \cref{open_in,topology_on}.
        $V\in T$ by \cref{topology_on}.
        $V$ is open in $X$ with respect to $T$ by \cref{open_in}.
        $x_0\in U$ by \cref{setminus_elim_left}.
        $Y\subseteq V$ by \cref{emptyset_subseteq}.
        Show that $U$ is disjoint from $V$.
        \begin{subproof}
            Suppose not.
            Then there exists $z$ such that $z\in U$ and $z\in V$.
            Thus $z\in\emptyset$.
            Contradiction by \cref{emptyset}.
        \end{subproof}
        Thus there exist $U,V$ such that
            $U$ is open in $X$ with respect to $T$ and
            $V$ is open in $X$ with respect to $T$ and
            $x_0\in U$ and
            $Y\subseteq V$ and
            $U$ is disjoint from $V$.
    \caseOf{$Y \neq \emptyset$.}
        $Y$ is inhabited by \cref{nonempty_inhabited}.
        Thus $x_0\in X$ and $x_0\notin Y$ by \cref{setminus_elim_left,setminus_elim_right}.
        Let $A = \{V\in T \mid \text{there exists $y\in Y$ such that there exists $U$ such that
                $U$ is a neighborhood of $x_0$ in $X$ with respect to $T$ and
                $V$ is a neighborhood of $y$ in $X$ with respect to $T$ and
                $U$ is disjoint from $V$}\}$.
        Show that $A$ is a covering of $Y$.
        \begin{subproof}
            Show that for all $y$ such that $y\in Y$ we have $y\in\unions{A}$.
            \begin{subproof}
                Fix $y$.
                Assume $y\in Y$.
                Thus $y\in X$ by \cref{subseteq}.
                Show that $x_0 \neq y$.
                \begin{subproof}
                    Suppose not.
                    Then $x_0 = y$.
                    Thus $x_0\in Y$.
                    Contradiction.
                \end{subproof}
                Thus there exist $U_1,U_2$ such that
                    $U_1$ is a neighborhood of $x_0$ in $X$ with respect to $T$ and
                    $U_2$ is a neighborhood of $y$ in $X$ with respect to $T$ and
                    $U_1$ is disjoint from $U_2$
                    by \cref{hausdorff}.
                Take $U_1,U_2$ such that
                    $U_1$ is a neighborhood of $x_0$ in $X$ with respect to $T$ and
                    $U_2$ is a neighborhood of $y$ in $X$ with respect to $T$ and
                    $U_1$ is disjoint from $U_2$.
                Thus $U_2$ is open in $X$ with respect to $T$ by \cref{neighborhood}.
                Thus $U_2\in T$ by \cref{open_in}.
                Thus $U_2\in A$.
                Thus $y\in U_2$ by \cref{neighborhood}.
                Thus $y\in\unions{A}$ by \cref{unions_iff}.
            \end{subproof}
            Thus $Y\subseteq\unions{A}$ by \cref{subseteq}.
            Thus $A$ is a covering of $Y$ by \cref{covering}.
        \end{subproof}
        Show that for all $U$ such that $U\in A$ we have $U$ is open in $X$ with respect to $T$.
        \begin{subproof}
            Fix $U$.
            Assume $U\in A$.
            Thus $U\in T$ by \cref{subseteq}.
            Thus $U$ is open in $X$ with respect to $T$ by \cref{open_in}.
        \end{subproof}
        $Y$ is a subspace of $X$ with respect to $T$ by \cref{subspace_of}.
        Thus there exists $B$ such that
            $B\subseteq A$ and $B$ is finite and $B$ is a covering of $Y$
            by \cref{compact_subspace_open_covering,subseteq}.
        Take $B$ such that $B\subseteq A$ and $B$ is finite and $B$ is a covering of $Y$.
        Show that $B$ is inhabited.
        \begin{subproof}
            Take $y$ such that $y\in Y$ by \cref{nonempty_inhabited}.
            $Y\subseteq\unions{B}$ by \cref{covering}.
            Thus $y\in\unions{B}$ by \cref{subseteq}.
            Take $V\in B$ such that $y\in V$ by \cref{unions_iff}.
            Thus $B$ is inhabited.
        \end{subproof}
        Let $R = \{w\in B\times T \mid \text{there exists $V\in B$ such that there exists $U\in T$ such that
                $w = (V,U)$ and
                $U$ is a neighborhood of $x_0$ in $X$ with respect to $T$ and
                $U$ is disjoint from $V$}\}$.
        Show that $R$ is a relation.
        \begin{subproof}
            Show that for all $w$ such that $w\in R$ there exist $V,U$ such that $w = (V,U)$.
            \begin{subproof}
                Fix $w$.
                Assume $w\in R$.
                Take $V\in B$ such that there exists $U\in T$ such that
                    $w = (V,U)$ and
                    $U$ is a neighborhood of $x_0$ in $X$ with respect to $T$ and
                    $U$ is disjoint from $V$.
                Take $U\in T$ such that
                    $w = (V,U)$ and
                    $U$ is a neighborhood of $x_0$ in $X$ with respect to $T$ and
                    $U$ is disjoint from $V$.
            \end{subproof}
            Thus $R$ is a relation by \cref{relation}.
        \end{subproof}
        Show that for all $V\in B$ there exists $U$ such that $V\mathrel{R} U$.
        \begin{subproof}
            Fix $V$.
            Assume $V\in B$.
            Thus $V\in A$ by \cref{subseteq}.
            Take $y\in Y$ such that there exists $U$ such that
                $U$ is a neighborhood of $x_0$ in $X$ with respect to $T$ and
                $V$ is a neighborhood of $y$ in $X$ with respect to $T$ and
                $U$ is disjoint from $V$.
            Take $U$ such that
                $U$ is a neighborhood of $x_0$ in $X$ with respect to $T$ and
                $V$ is a neighborhood of $y$ in $X$ with respect to $T$ and
                $U$ is disjoint from $V$.
            Thus $U$ is open in $X$ with respect to $T$ by \cref{neighborhood}.
            Thus $U\in T$ by \cref{open_in}.
            Thus $(V,U)\in R$.
            Thus $V\mathrel{R} U$.
        \end{subproof}
        Take $g$ such that $g$ is a function on $B$ and for all $V\in B$ we have $V\mathrel{R}\apply{g}{V}$ by \cref{choice_function}.
        Let $G = \img{g}{B}$.
        Show that $G\subseteq T$.
        \begin{subproof}
            Show that for all $U$ such that $U\in G$ we have $U\in T$.
            \begin{subproof}
                Fix $U$.
                Assume $U\in G$.
                Take $V\in B$ such that $(V,U)\in g$ by \cref{img_iff}.
                Thus $V\mathrel{R} U$ by \cref{function_apply_iff}.
                Take $V_1\in B$ such that there exists $U_1\in T$ such that
                    $(V,U) = (V_1,U_1)$ and
                    $U_1$ is a neighborhood of $x_0$ in $X$ with respect to $T$ and
                    $U_1$ is disjoint from $V_1$.
                Take $U_1\in T$ such that
                    $(V,U) = (V_1,U_1)$ and
                    $U_1$ is a neighborhood of $x_0$ in $X$ with respect to $T$ and
                    $U_1$ is disjoint from $V_1$.
                Thus $U = U_1$ by \cref{pair_eq_iff}.
                Thus $U\in T$.
            \end{subproof}
            Thus $G\subseteq T$ by \cref{subseteq}.
        \end{subproof}
        $G$ is finite by \cref{finite_image_function}.
        Show that $G$ is inhabited.
        \begin{subproof}
            Take $y$ such that $y\in Y$ by \cref{nonempty_inhabited}.
            $Y\subseteq\unions{B}$ by \cref{covering}.
            Thus $y\in\unions{B}$ by \cref{subseteq}.
            Take $V\in B$ such that $y\in V$ by \cref{unions_iff}.
            Thus $(V,\apply{g}{V})\in g$ by \cref{function_apply_elim,subseteq}.
            Thus $\apply{g}{V}\in G$ by \cref{img_iff}.
            Thus $G$ is inhabited.
        \end{subproof}
        Let $U_0 = \inters{G}$.
        $T$ is a topology on $X$ by \cref{hausdorff}.
        Thus $U_0\in T$ by \cref{finite_intersections_open_in_topology}.
        Show that $x_0\in U_0$.
        \begin{subproof}
            Show that for all $U$ such that $U\in G$ we have $x_0\in U$.
            \begin{subproof}
                Fix $U$.
                Assume $U\in G$.
                Take $V\in B$ such that $(V,U)\in g$ by \cref{img_iff}.
                Thus $V\mathrel{R} U$ by \cref{function_apply_iff}.
                Take $V_1\in B$ such that there exists $U_1\in T$ such that
                    $(V,U) = (V_1,U_1)$ and
                    $U_1$ is a neighborhood of $x_0$ in $X$ with respect to $T$ and
                    $U_1$ is disjoint from $V_1$.
                Take $U_1\in T$ such that
                    $(V,U) = (V_1,U_1)$ and
                    $U_1$ is a neighborhood of $x_0$ in $X$ with respect to $T$ and
                    $U_1$ is disjoint from $V_1$.
                Thus $U = U_1$ by \cref{pair_eq_iff}.
                Thus $x_0\in U$ by \cref{neighborhood}.
            \end{subproof}
            Thus $x_0\in U_0$ by \cref{inters_intro}.
        \end{subproof}
        Let $V_0 = \unions{B}$.
        Show that $V_0$ is open in $X$ with respect to $T$.
        \begin{subproof}
            Thus $B\subseteq A$ by \cref{subseteq}.
            Thus $B\subseteq T$ by \cref{subseteq}.
            Thus $V_0\in T$ by \cref{topology_on}.
            Thus $V_0$ is open in $X$ with respect to $T$ by \cref{open_in}.
        \end{subproof}
        $Y\subseteq V_0$ by \cref{covering}.
        Show that $U_0$ is disjoint from $V_0$.
        \begin{subproof}
            Suppose not.
            Then there exists $z$ such that $z\in U_0$ and $z\in V_0$.
            Take $V\in B$ such that $z\in V$ by \cref{unions_iff}.
            Thus $(V,\apply{g}{V})\in g$ by \cref{function_apply_elim,subseteq}.
            Thus $\apply{g}{V}\in G$ by \cref{img_iff}.
            Thus $U_0\subseteq \apply{g}{V}$ by \cref{inters_subseteq_elem}.
            Thus $z\in \apply{g}{V}$ by \cref{subseteq}.
            Thus $V\mathrel{R}\apply{g}{V}$ by \cref{subseteq}.
            Take $V_1\in B$ such that there exists $U_1\in T$ such that
                $(V,\apply{g}{V}) = (V_1,U_1)$ and
                $U_1$ is a neighborhood of $x_0$ in $X$ with respect to $T$ and
                $U_1$ is disjoint from $V_1$.
            Take $U_1\in T$ such that
                $(V,\apply{g}{V}) = (V_1,U_1)$ and
                $U_1$ is a neighborhood of $x_0$ in $X$ with respect to $T$ and
                $U_1$ is disjoint from $V_1$.
            Thus $V = V_1$ and $\apply{g}{V} = U_1$ by \cref{pair_eq_iff}.
            Thus $\apply{g}{V}$ is disjoint from $V$.
            Thus $z\notin \apply{g}{V}$ by \cref{disjoint}.
            Contradiction.
        \end{subproof}
        Thus there exist $U,V$ such that
            $U$ is open in $X$ with respect to $T$ and
            $V$ is open in $X$ with respect to $T$ and
            $x_0\in U$ and
            $Y\subseteq V$ and
            $U$ is disjoint from $V$.
    \end{byCase}
\end{proof}

% from §26 Item 8: image of a compact space under a continuous map is compact
\begin{theorem}\label{continuous_image_compact}
    Assume $f$ is continuous from $X$ to $Y$ with respect to $T$ and $T_1$.
    Assume $X$ is compact with respect to $T$.
    Let $A = \img{f}{X}$.
    Let $T_A = \subspaceTopology{T_1}{A}$.
    Then $A$ is compact with respect to $T_A$.
\end{theorem}
\begin{proof}
    Let $A = \img{f}{X}$.
    Let $T_A = \subspaceTopology{T_1}{A}$.
    Show that for all $C$ such that
        $C$ is a covering of $A$ and
        for all $U$ such that $U\in C$ we have $U$ is open in $Y$ with respect to $T_1$
        there exists $B$ such that
            $B\subseteq C$ and
            $B$ is finite and
            $B$ is a covering of $A$.
    \begin{subproof}
        Fix $C$.
        Assume $C$ is a covering of $A$.
        Assume for all $U$ such that $U\in C$ we have $U$ is open in $Y$ with respect to $T_1$.
        Let $D = \{W \mid \exists U\in C. W = \preimg{f}{U}\}$.
        Show that $D$ is a covering of $X$.
        \begin{subproof}
            Show that for all $x$ such that $x\in X$ we have $x\in\unions{D}$.
            \begin{subproof}
                Fix $x$.
                Assume $x\in X$.
                Thus $f\in\funs{X}{Y}$ by \cref{continuous}.
                Thus $f$ is a function by \cref{funs_is_function}.
                Thus $x\in\dom{f}$ by \cref{funs}.
                Thus $(x,f(x))\in f$ by \cref{function_apply_elim}.
                Thus $f(x)\in A$ by \cref{img_iff}.
                $A\subseteq\unions{C}$ by \cref{covering}.
                Thus $f(x)\in\unions{C}$ by \cref{subseteq}.
                Take $U\in C$ such that $f(x)\in U$ by \cref{unions_iff}.
                Thus $x\in\preimg{f}{U}$ by \cref{preimg_iff}.
                Thus $\preimg{f}{U}\in D$.
                Thus $x\in\unions{D}$ by \cref{unions_iff}.
            \end{subproof}
            Thus $X\subseteq\unions{D}$ by \cref{subseteq}.
            Thus $D$ is a covering of $X$ by \cref{covering}.
        \end{subproof}
        Show that for all $W$ such that $W\in D$ we have $W$ is open in $X$ with respect to $T$.
        \begin{subproof}
            Fix $W$.
            Assume $W\in D$.
            Take $U\in C$ such that $W = \preimg{f}{U}$.
            $U$ is open in $Y$ with respect to $T_1$.
            Thus $\preimg{f}{U}$ is open in $X$ with respect to $T$ by \cref{continuous}.
            Thus $W$ is open in $X$ with respect to $T$.
        \end{subproof}
        Thus $D$ is an open covering of $X$ with respect to $T$ by \cref{open_covering}.
        Thus there exists $E$ such that
            $E\subseteq D$ and $E$ is finite and $E$ is a covering of $X$
            by \cref{compact}.
        Take $E$ such that $E\subseteq D$ and $E$ is finite and $E$ is a covering of $X$.
        Let $R = \{w\in E\times C \mid \exists W\in E.\exists U\in C. w = (W,U) \land W = \preimg{f}{U}\}$.
        Show that $R$ is a relation.
        \begin{subproof}
            Show that for all $w$ such that $w\in R$ there exist $W,U$ such that $w = (W,U)$.
            \begin{subproof}
                Fix $w$.
                Assume $w\in R$.
                Take $W\in E$ such that there exists $U\in C$ such that $w = (W,U)$ and $W = \preimg{f}{U}$.
                Take $U\in C$ such that $w = (W,U)$ and $W = \preimg{f}{U}$.
            \end{subproof}
            Thus $R$ is a relation by \cref{relation}.
        \end{subproof}
        Show that for all $W\in E$ there exists $U$ such that $W\mathrel{R} U$.
        \begin{subproof}
            Fix $W$.
            Assume $W\in E$.
            Thus $W\in D$ by \cref{subseteq}.
            Take $U\in C$ such that $W = \preimg{f}{U}$.
            Thus $(W,U)\in R$.
            Thus $W\mathrel{R} U$.
        \end{subproof}
        Take $h$ such that $h$ is a function on $E$ and for all $W\in E$ we have $W\mathrel{R}\apply{h}{W}$ by \cref{choice_function}.
        Thus $h$ is a function and $\dom{h} = E$.
        Let $B = \img{h}{E}$.
        Show that $B\subseteq C$.
        \begin{subproof}
            Show that for all $U$ such that $U\in B$ we have $U\in C$.
            \begin{subproof}
                Fix $U$.
                Assume $U\in B$.
                Take $W\in E$ such that $(W,U)\in h$ by \cref{img_iff}.
                Thus $U = \apply{h}{W}$ by \cref{function_apply_intro}.
                Thus $W\mathrel{R}\apply{h}{W}$.
                Thus $W\mathrel{R} U$.
                Take $W_1\in E$ such that there exists $U_1\in C$ such that
                    $(W,U) = (W_1,U_1)$ and $W_1 = \preimg{f}{U_1}$.
                Take $U_1\in C$ such that $(W,U) = (W_1,U_1)$ and $W_1 = \preimg{f}{U_1}$.
                Thus $U = U_1$ by \cref{pair_eq_iff}.
                Thus $U\in C$.
            \end{subproof}
            Thus $B\subseteq C$ by \cref{subseteq}.
        \end{subproof}
        $B$ is finite by \cref{finite_image_function}.
        Show that $B$ is a covering of $A$.
        \begin{subproof}
            Show that for all $a$ such that $a\in A$ we have $a\in\unions{B}$.
            \begin{subproof}
                Fix $a$.
                Assume $a\in A$.
                Take $x\in X$ such that $(x,a)\in f$ by \cref{img_iff}.
                Thus $f\in\funs{X}{Y}$ by \cref{continuous}.
                Thus $f$ is a function by \cref{funs_is_function}.
                Thus $f(x) = a$ by \cref{function_apply_intro}.
                $X\subseteq\unions{E}$ by \cref{covering}.
                Thus $x\in\unions{E}$ by \cref{subseteq}.
                Take $W\in E$ such that $x\in W$ by \cref{unions_iff}.
                Thus $W\mathrel{R}\apply{h}{W}$.
                Take $W_1\in E$ such that there exists $U_1\in C$ such that
                    $(W,\apply{h}{W}) = (W_1,U_1)$ and $W_1 = \preimg{f}{U_1}$.
                Take $U_1\in C$ such that $(W,\apply{h}{W}) = (W_1,U_1)$ and $W_1 = \preimg{f}{U_1}$.
                Thus $W = W_1$ and $\apply{h}{W} = U_1$ by \cref{pair_eq_iff}.
                Thus $x\in \preimg{f}{\apply{h}{W}}$.
                Take $b\in \apply{h}{W}$ such that $x\mathrel{f} b$ by \cref{preimg_iff}.
                Thus $f(x) = b$ by \cref{function_apply_intro}.
                Thus $a = b$.
                Thus $a\in \apply{h}{W}$.
                Thus $(W,\apply{h}{W})\in h$ by \cref{function_apply_elim}.
                Thus $\apply{h}{W}\in B$ by \cref{img_iff}.
                Thus $a\in\unions{B}$ by \cref{unions_iff}.
            \end{subproof}
            Thus $A\subseteq\unions{B}$ by \cref{subseteq}.
            Thus $B$ is a covering of $A$ by \cref{covering}.
        \end{subproof}
    \end{subproof}
    $T_1$ is a topology on $Y$ by \cref{continuous}.
    $f\in\funs{X}{Y}$ by \cref{continuous}.
    Thus $\ran{f}\subseteq Y$ by \cref{funs,rels_ran_subseteq}.
    Thus $A\subseteq \ran{f}$ by \cref{img_subseteq_ran}.
    Thus $A\subseteq Y$ by \cref{subseteq_transitive}.
    Thus $A$ is a subspace of $Y$ with respect to $T_1$ by \cref{subspace_of}.
    Thus $A$ is compact with respect to $T_A$ by \cref{compact_subspace_open_covering}.
\end{proof}

% from §26 Item 9: bijective continuous map from compact to Hausdorff is a homeomorphism
\begin{theorem}\label{compact_hausdorff_bijection_homeomorphism}
    Assume $f$ is a bijection from $X$ to $Y$.
    Assume $f$ is continuous from $X$ to $Y$ with respect to $T$ and $T_1$.
    Assume $X$ is compact with respect to $T$.
    Assume $Y$ is Hausdorff with respect to $T_1$.
    Then $f$ is a homeomorphism between $X$ and $Y$ with respect to $T$ and $T_1$.
\end{theorem}
\begin{proof}
    $T$ is a topology on $X$ by \cref{continuous}.
    $T_1$ is a topology on $Y$ by \cref{continuous}.
    Show that for all $A$ such that $A$ is closed in $X$ with respect to $T$
        we have $\img{f}{A}$ is closed in $Y$ with respect to $T_1$.
    \begin{subproof}
        Fix $A$.
        Assume $A$ is closed in $X$ with respect to $T$.
        Thus $A\subseteq X$ by \cref{closed_in}.
        Thus $A$ is a subspace of $X$ with respect to $T$ by \cref{subspace_of}.
        Let $T_A = \subspaceTopology{T}{A}$.
        Thus $A$ is compact with respect to $T_A$ by \cref{closed_subspace_compact}.
        Let $g = \restrl{f}{A}$.
        Thus $g$ is continuous from $A$ to $Y$ with respect to $T_A$ and $T_1$
            by \cref{continuous_restrict_domain}.
        Let $B = \img{g}{A}$.
        Thus $f\in\funs{X}{Y}$ by \cref{bijections_to_funs}.
        Show that $A\subseteq\dom{f}$.
        \begin{subproof}
            Show that for all $a$ such that $a\in A$ we have $a\in\dom{f}$.
            \begin{subproof}
                Fix $a$.
                Assume $a\in A$.
                Thus $a\in X$ by \cref{subseteq}.
                Thus $a\in\dom{f}$ by \cref{funs}.
            \end{subproof}
            Thus $A\subseteq\dom{f}$ by \cref{subseteq}.
        \end{subproof}
        Thus $B = \img{f}{A\inter A}$ by \cref{restrl_img}.
        Thus $B = \img{f}{A}$ by \cref{inter_idempotent}.
        Let $T_B = \subspaceTopology{T_1}{B}$.
        Thus $B$ is compact with respect to $T_B$ by \cref{continuous_image_compact}.
        Thus $\ran{f}\subseteq Y$ by \cref{funs,rels_ran_subseteq}.
        Thus $B\subseteq \ran{f}$ by \cref{img_subseteq_ran}.
        Thus $B\subseteq Y$ by \cref{subseteq_transitive}.
        Thus $B$ is a subspace of $Y$ with respect to $T_1$ by \cref{subspace_of}.
        Thus $B$ is closed in $Y$ with respect to $T_1$ by \cref{compact_subspace_closed}.
        Thus $\img{f}{A}$ is closed in $Y$ with respect to $T_1$.
    \end{subproof}
    Thus $f\in\funs{X}{Y}$ by \cref{bijections_to_funs}.
    Thus $f$ is a function from $X$ to $Y$.
    Thus $f$ is a closed map from $X$ to $Y$ with respect to $T$ and $T_1$ by \cref{closed_map}.
    Show that $\converse{f}$ is continuous from $Y$ to $X$ with respect to $T_1$ and $T$.
    \begin{subproof}
        $\converse{f}\in\funs{Y}{X}$ by \cref{bijective_converse_are_funs}.
        Show that for all $C$ such that $C$ is closed in $X$ with respect to $T$
            we have $\preimg{\converse{f}}{C}$ is closed in $Y$ with respect to $T_1$.
        \begin{subproof}
            Fix $C$.
            Assume $C$ is closed in $X$ with respect to $T$.
            Thus $\img{f}{C}$ is closed in $Y$ with respect to $T_1$ by \cref{closed_map}.
            Thus $f\in\funs{X}{Y}$ by \cref{bijections_to_funs}.
            Thus $f$ is a relation by \cref{funs_is_relation}.
            $\preimg{\converse{f}}{C} = \img{\converse{\converse{f}}}{C}$ by \cref{preim_eq_img_of_converse}.
            Thus $\preimg{\converse{f}}{C} = \img{f}{C}$ by \cref{converse_converse_eq}.
            Thus $\preimg{\converse{f}}{C}$ is closed in $Y$ with respect to $T_1$.
        \end{subproof}
        Thus $\converse{f}$ is continuous from $Y$ to $X$ with respect to $T_1$ and $T$
            by \cref{preimg_closed_implies_continuous}.
    \end{subproof}
    Thus $f$ is a homeomorphism between $X$ and $Y$ with respect to $T$ and $T_1$
        by \cref{homeomorphism}.
\end{proof}

% from §26 helper lemma 1: slice is compact
\begin{lemma}\label{slice_compact}
    Assume $T$ is a topology on $X$.
    Assume $T_1$ is a topology on $Y$.
    Assume $Y$ is compact with respect to $T_1$.
    Assume $x_0\in X$.
    Let $T_2 = \productTopology{T}{T_1}{X}{Y}$.
    Let $S = \{x_0\}\times Y$.
    Let $T_S = \subspaceTopology{T_2}{S}$.
    Then $S$ is compact with respect to $T_S$.
\end{lemma}
\begin{proof}
    Let $f_1 = \{(y,x_0) \mid y\in Y\}$.
    Thus $f_1$ is a function from $Y$ to $X$ by \cref{constant_map_function}.
    Show that for all $y$ such that $y\in Y$ we have $f_1(y) = x_0$.
    \begin{subproof}
        Fix $y$.
        Assume $y\in Y$.
        Thus $(y,x_0)\in f_1$.
        $f_1\in\funs{Y}{X}$ by \cref{constant_map_function}.
        Thus $f_1$ is a function by \cref{funs_is_function}.
        Thus $f_1(y) = x_0$ by \cref{function_apply_intro}.
    \end{subproof}
    Thus $f_1$ is continuous from $Y$ to $X$ with respect to $T_1$ and $T$
        by \cref{constant_continuous}.
    Let $f_2 = \identity{Y}$.
    Thus $f_2$ is continuous from $Y$ to $Y$ with respect to $T_1$ and $T_1$
        by \cref{identity_continuous}.
    Show that for all $y$ such that $y\in Y$ we have $f_2(y) = y$.
    \begin{subproof}
        Fix $y$.
        Assume $y\in Y$.
        Thus $y\mathrel{f_2} y$ by \cref{id_iff}.
        $f_2\in\funs{Y}{Y}$ by \cref{id_elem_funs}.
        Thus $f_2$ is a function by \cref{funs_is_function}.
        Thus $f_2(y) = y$ by \cref{function_apply_intro}.
    \end{subproof}
    Let $k = \{(y,(f_1(y),f_2(y))) \mid y\in Y\}$.
    Thus $k$ is continuous from $Y$ to $X\times Y$ with respect to $T_1$ and $T_2$
        by \cref{continuous_map_into_product_backward}.
    Let $S_0 = \img{k}{Y}$.
    Thus $S_0$ is compact with respect to $\subspaceTopology{T_2}{S_0}$ by \cref{continuous_image_compact}.
    Show that $S_0 = S$.
    \begin{subproof}
        Show that $S_0\subseteq S$.
        \begin{subproof}
            Show that for all $p$ such that $p\in S_0$ we have $p\in S$.
            \begin{subproof}
                Fix $p$.
                Assume $p\in S_0$.
                Take $y\in Y$ such that $(y,p)\in k$ by \cref{img_iff}.
                Take $y_1\in Y$ such that $(y,p) = (y_1,(f_1(y_1),f_2(y_1)))$.
                Thus $y = y_1$ and $p = (f_1(y_1),f_2(y_1))$ by \cref{pair_eq_iff}.
                Thus $p = (f_1(y),f_2(y))$.
                Thus $p = (x_0,y)$.
                Thus $x_0\in\{x_0\}$ by \cref{singleton_iff}.
                Thus $p\in \{x_0\}\times Y$ by \cref{times_tuple_intro}.
                Thus $p\in S$.
            \end{subproof}
            Thus $S_0\subseteq S$ by \cref{subseteq}.
        \end{subproof}
        Show that $S\subseteq S_0$.
        \begin{subproof}
            Show that for all $p$ such that $p\in S$ we have $p\in S_0$.
            \begin{subproof}
                Fix $p$.
                Assume $p\in S$.
                Take $x,y$ such that $p = (x,y)$ and $x\in\{x_0\}$ and $y\in Y$
                    by \cref{times_elem_is_tuple}.
                Thus $x = x_0$ by \cref{singleton_iff}.
                Thus $p = (x_0,y)$.
                Thus $f_1(y) = x_0$.
                Thus $f_2(y) = y$.
                Thus $(y,p) = (y,(f_1(y),f_2(y)))$.
                Thus $(y,p)\in k$.
                Thus $p\in S_0$ by \cref{img_iff}.
            \end{subproof}
            Thus $S\subseteq S_0$ by \cref{subseteq}.
        \end{subproof}
        Thus $S_0 = S$ by \cref{subseteq_antisymmetric}.
    \end{subproof}
    Thus $S$ is compact with respect to $T_S$.
\end{proof}

% from §26 helper lemma 2: tube lemma
\begin{lemma}\label{tube_lemma_helper}
    Assume $T$ is a topology on $X$.
    Assume $T_1$ is a topology on $Y$.
    Assume $Y$ is compact with respect to $T_1$.
    Let $T_2 = \productTopology{T}{T_1}{X}{Y}$.
    Assume $N$ is open in $X\times Y$ with respect to $T_2$.
    Assume $x_0\in X$.
    Suppose $\{x_0\}\times Y\subseteq N$.
    Then there exists $W$ such that
        $W$ is open in $X$ with respect to $T$ and
        $x_0\in W$ and
        $W\times Y\subseteq N$.
\end{lemma}
\begin{proof}
    Let $S = \{x_0\}\times Y$.
    Let $T_S = \subspaceTopology{T_2}{S}$.
    \begin{byCase}
    \caseOf{$Y = \emptyset$.}
        Let $W = X$.
        $W$ is open in $X$ with respect to $T$ by \cref{open_in,topology_on}.
        $x_0\in W$ by \cref{subseteq_refl}.
        $W\times Y = \emptyset$ by \cref{times_empty_right}.
        Thus $W\times Y\subseteq N$ by \cref{emptyset_subseteq}.
        Thus there exists $W$ such that
            $W$ is open in $X$ with respect to $T$ and
            $x_0\in W$ and
            $W\times Y\subseteq N$.
    \caseOf{$Y\neq\emptyset$.}
        Take $y_0$ such that $y_0\in Y$ by \cref{nonempty_inhabited}.
        Thus $x_0\in\{x_0\}$ by \cref{singleton_iff}.
        Thus $(x_0,y_0)\in S$ by \cref{times_tuple_intro}.
        Thus $S$ is inhabited.
        Thus $S$ is compact with respect to $T_S$ by \cref{slice_compact}.
        Let $B = \productBasis{T}{T_1}{X}{Y}$.
        $B$ is a basis on $X\times Y$ by \cref{product_basis_is_basis}.
        Thus $T_2$ is a topology on $X\times Y$ by \cref{basis_topology_is_topology,product_topology}.
        Show that $S\subseteq X\times Y$.
        \begin{subproof}
            Show that for all $p$ such that $p\in S$ we have $p\in X\times Y$.
            \begin{subproof}
                Fix $p$.
                Assume $p\in S$.
                Take $x,y$ such that $p = (x,y)$ and $x\in\{x_0\}$ and $y\in Y$
                    by \cref{times_elem_is_tuple}.
                Thus $x = x_0$ by \cref{singleton_iff}.
                Thus $x\in X$.
                Thus $p\in X\times Y$ by \cref{times_tuple_intro}.
            \end{subproof}
            Thus $S\subseteq X\times Y$ by \cref{subseteq}.
        \end{subproof}
        Thus $S$ is a subspace of $X\times Y$ with respect to $T_2$ by \cref{subspace_of}.
        Let $R = \{w\in S\times B \mid \exists p\in S.\exists B_1\in B. w = (p,B_1) \land p\in B_1 \land B_1\subseteq N\}$.
        Show that $R$ is a relation.
        \begin{subproof}
            Show that for all $w$ such that $w\in R$ there exist $p,B_1$ such that $w = (p,B_1)$.
            \begin{subproof}
                Fix $w$.
                Assume $w\in R$.
                Take $p\in S$ such that there exists $B_1\in B$ such that
                    $w = (p,B_1)$ and $p\in B_1$ and $B_1\subseteq N$.
                Take $B_1\in B$ such that $w = (p,B_1)$ and $p\in B_1$ and $B_1\subseteq N$.
            \end{subproof}
            Thus $R$ is a relation by \cref{relation}.
        \end{subproof}
        Show that for all $p\in S$ there exists $B_1$ such that $p\mathrel{R} B_1$.
        \begin{subproof}
            Fix $p$.
            Assume $p\in S$.
            Thus $p\in N$ by \cref{subseteq}.
            $N\in T_2$ by \cref{open_in}.
            $T_2 = \basisTopology{B}{X\times Y}$ by \cref{product_topology}.
            Take $B_1\in B$ such that $p\in B_1$ and $B_1\subseteq N$
                by \cref{topology_generated_by_basis}.
            Thus $(p,B_1)\in R$.
            Thus $p\mathrel{R} B_1$.
        \end{subproof}
        Take $g$ such that $g$ is a function on $S$ and for all $p\in S$ we have $p\mathrel{R}\apply{g}{p}$
            by \cref{choice_function}.
        Let $C = \img{g}{S}$.
        Show that $C$ is a covering of $S$.
        \begin{subproof}
            Show that for all $p$ such that $p\in S$ we have $p\in\unions{C}$.
            \begin{subproof}
                Fix $p$.
                Assume $p\in S$.
                Thus $p\mathrel{R}\apply{g}{p}$.
                Take $p_1\in S$ such that there exists $B_1\in B$ such that
                    $(p,\apply{g}{p}) = (p_1,B_1)$ and $p_1\in B_1$ and $B_1\subseteq N$.
                Take $B_1\in B$ such that
                    $(p,\apply{g}{p}) = (p_1,B_1)$ and $p_1\in B_1$ and $B_1\subseteq N$.
                Thus $p = p_1$ and $\apply{g}{p} = B_1$ by \cref{pair_eq_iff}.
                Thus $p\in \apply{g}{p}$.
                Thus $(p,\apply{g}{p})\in g$ by \cref{function_apply_elim}.
                Thus $\apply{g}{p}\in C$ by \cref{img_iff}.
                Thus $p\in\unions{C}$ by \cref{unions_iff}.
            \end{subproof}
            Thus $S\subseteq\unions{C}$ by \cref{subseteq}.
            Thus $C$ is a covering of $S$ by \cref{covering}.
        \end{subproof}
        Show that for all $U$ such that $U\in C$ we have $U$ is open in $X\times Y$ with respect to $T_2$.
        \begin{subproof}
            Fix $U$.
            Assume $U\in C$.
            Take $p\in S$ such that $(p,U)\in g$ by \cref{img_iff}.
            Thus $p\mathrel{R} U$ by \cref{function_apply_iff}.
            Take $p_1\in S$ such that there exists $B_1\in B$ such that
                $(p,U) = (p_1,B_1)$ and $p_1\in B_1$ and $B_1\subseteq N$.
            Take $B_1\in B$ such that $(p,U) = (p_1,B_1)$ and $p_1\in B_1$ and $B_1\subseteq N$.
            Thus $U = B_1$ by \cref{pair_eq_iff}.
            Thus $U\in B$.
            Thus $U\in T_2$ by \cref{basis_elem_open_in_generated,product_topology}.
            Thus $U$ is open in $X\times Y$ with respect to $T_2$ by \cref{open_in}.
        \end{subproof}
        Thus there exists $F$ such that
            $F\subseteq C$ and $F$ is finite and $F$ is a covering of $S$
            by \cref{compact_subspace_open_covering}.
        Take $F$ such that $F\subseteq C$ and $F$ is finite and $F$ is a covering of $S$.
        Let $R_1 = \{w\in F\times T \mid \exists E\in F.\exists U\in T.\exists V\in T_1. w = (E,U) \land E = U\times V\}$.
        Show that $R_1$ is a relation.
        \begin{subproof}
            Show that for all $w$ such that $w\in R_1$ there exist $E,U$ such that $w = (E,U)$.
            \begin{subproof}
                Fix $w$.
                Assume $w\in R_1$.
                Take $E\in F$ such that there exists $U\in T$ such that there exists $V\in T_1$ such that
                    $w = (E,U)$ and $E = U\times V$.
                Take $U\in T$ such that there exists $V\in T_1$ such that $w = (E,U)$ and $E = U\times V$.
            \end{subproof}
            Thus $R_1$ is a relation by \cref{relation}.
        \end{subproof}
        Show that for all $E\in F$ there exists $U$ such that $E\mathrel{R_1} U$.
        \begin{subproof}
            Fix $E$.
            Assume $E\in F$.
            Thus $E\in C$ by \cref{subseteq}.
            Take $p\in S$ such that $(p,E)\in g$ by \cref{img_iff}.
            Thus $p\mathrel{R} E$ by \cref{function_apply_iff}.
            Take $p_1\in S$ such that there exists $B_1\in B$ such that
                $(p,E) = (p_1,B_1)$ and $p_1\in B_1$ and $B_1\subseteq N$.
            Take $B_1\in B$ such that $(p,E) = (p_1,B_1)$ and $p_1\in B_1$ and $B_1\subseteq N$.
            Thus $E = B_1$ by \cref{pair_eq_iff}.
            Take $U\in T$ such that there exists $V\in T_1$ such that $E = U\times V$
                by \cref{product_basis}.
            Take $V\in T_1$ such that $E = U\times V$.
            Thus $(E,U)\in R_1$.
            Thus $E\mathrel{R_1} U$.
        \end{subproof}
        Take $h$ such that $h$ is a function on $F$ and for all $E\in F$ we have $E\mathrel{R_1}\apply{h}{E}$
            by \cref{choice_function}.
        Let $G = \img{h}{F}$.
        Thus $G$ is finite by \cref{finite_image_function}.
        Show that $G$ is inhabited.
        \begin{subproof}
            Take $p_0$ such that $p_0\in S$ by \cref{nonempty_inhabited}.
            $S\subseteq\unions{F}$ by \cref{covering}.
            Thus $p_0\in\unions{F}$ by \cref{subseteq}.
            Take $E\in F$ such that $p_0\in E$ by \cref{unions_iff}.
            Thus $(E,\apply{h}{E})\in h$ by \cref{function_apply_elim}.
            Thus $\apply{h}{E}\in G$ by \cref{img_iff}.
            Thus $G$ is inhabited.
        \end{subproof}
        Show that for all $U$ such that $U\in G$ we have $U\in T$.
        \begin{subproof}
            Fix $U$.
            Assume $U\in G$.
            Take $E\in F$ such that $(E,U)\in h$ by \cref{img_iff}.
            Thus $E\mathrel{R_1} U$ by \cref{function_apply_iff}.
            Take $E_1\in F$ such that there exists $U_1\in T$ such that there exists $V_1\in T_1$ such that
                $(E,U) = (E_1,U_1)$ and $E_1 = U_1\times V_1$.
            Take $U_1\in T$ such that there exists $V_1\in T_1$ such that
                $(E,U) = (E_1,U_1)$ and $E_1 = U_1\times V_1$.
            Thus $U = U_1$ by \cref{pair_eq_iff}.
            Thus $U\in T$.
        \end{subproof}
        Thus $G\subseteq T$ by \cref{subseteq}.
        Let $W = \inters{G}$.
        Thus $W\in T$ by \cref{finite_intersections_open_in_topology}.
        Thus $W$ is open in $X$ with respect to $T$ by \cref{open_in}.
        Show that $x_0\in W$.
        \begin{subproof}
            Show that for all $U$ such that $U\in G$ we have $x_0\in U$.
            \begin{subproof}
                Fix $U$.
                Assume $U\in G$.
                Take $E\in F$ such that $(E,U)\in h$ by \cref{img_iff}.
                Thus $E\mathrel{R_1} U$ by \cref{function_apply_iff}.
                Take $E_1\in F$ such that there exists $U_1\in T$ such that there exists $V_1\in T_1$ such that
                    $(E,U) = (E_1,U_1)$ and $E_1 = U_1\times V_1$.
                Take $U_1\in T$ such that there exists $V_1\in T_1$ such that
                    $(E,U) = (E_1,U_1)$ and $E_1 = U_1\times V_1$.
                Take $V_1\in T_1$ such that
                    $(E,U) = (E_1,U_1)$ and $E_1 = U_1\times V_1$.
                Thus $E = E_1$ and $U = U_1$ by \cref{pair_eq_iff}.
                Thus $E\in C$ by \cref{subseteq}.
                Take $p\in S$ such that $(p,E)\in g$ by \cref{img_iff}.
                Thus $p\mathrel{R} E$ by \cref{function_apply_iff}.
                Take $p_1\in S$ such that there exists $B_1\in B$ such that
                    $(p,E) = (p_1,B_1)$ and $p_1\in B_1$ and $B_1\subseteq N$.
                Take $B_1\in B$ such that $(p,E) = (p_1,B_1)$ and $p_1\in B_1$ and $B_1\subseteq N$.
                Thus $E = B_1$ by \cref{pair_eq_iff}.
                Take $x,y$ such that $p_1 = (x,y)$ and $x\in\{x_0\}$ and $y\in Y$
                    by \cref{times_elem_is_tuple}.
                Thus $x = x_0$ by \cref{singleton_iff}.
                Thus $(x_0,y)\in E$.
                Thus $E = U_1\times V_1$.
                Thus $(x_0,y)\in U_1\times V_1$.
                Thus $x_0\in U_1$ by \cref{times_tuple_elim}.
                Thus $x_0\in U$.
            \end{subproof}
            Thus $x_0\in W$ by \cref{inters_intro}.
        \end{subproof}
        Show that $W\times Y\subseteq N$.
        \begin{subproof}
            Show that for all $p$ such that $p\in W\times Y$ we have $p\in N$.
            \begin{subproof}
                Fix $p$.
                Assume $p\in W\times Y$.
                Take $x,y$ such that $p = (x,y)$ and $x\in W$ and $y\in Y$
                    by \cref{times_elem_is_tuple}.
                Thus $(x_0,y)\in S$ by \cref{times_tuple_intro}.
                $S\subseteq\unions{F}$ by \cref{covering}.
                Thus $(x_0,y)\in\unions{F}$ by \cref{subseteq}.
                Take $E\in F$ such that $(x_0,y)\in E$ by \cref{unions_iff}.
                Thus $E\mathrel{R_1}\apply{h}{E}$.
                Take $E_1\in F$ such that there exists $U_1\in T$ such that there exists $V_1\in T_1$ such that
                    $(E,\apply{h}{E}) = (E_1,U_1)$ and $E_1 = U_1\times V_1$.
                Take $U_1\in T$ such that there exists $V_1\in T_1$ such that
                    $(E,\apply{h}{E}) = (E_1,U_1)$ and $E_1 = U_1\times V_1$.
                Take $V_1\in T_1$ such that
                    $(E,\apply{h}{E}) = (E_1,U_1)$ and $E_1 = U_1\times V_1$.
                Thus $E = E_1$ and $\apply{h}{E} = U_1$ by \cref{pair_eq_iff}.
                Thus $E = U_1\times V_1$.
                Thus $(x_0,y)\in U_1\times V_1$.
                Thus $y\in V_1$ by \cref{times_tuple_elim}.
                Thus $(E,\apply{h}{E})\in h$ by \cref{function_apply_elim}.
                Thus $\apply{h}{E}\in G$ by \cref{img_iff}.
                Thus $U_1\in G$.
                Thus $W\subseteq U_1$ by \cref{inters_subseteq_elem}.
                Thus $x\in U_1$ by \cref{subseteq}.
                Thus $p\in U_1\times V_1$ by \cref{times_tuple_intro}.
                Thus $p\in E$ by \cref{subseteq}.
                Thus $E\in C$ by \cref{subseteq}.
                Take $p_0\in S$ such that $(p_0,E)\in g$ by \cref{img_iff}.
                Thus $p_0\mathrel{R} E$ by \cref{function_apply_iff}.
                Take $p_1\in S$ such that there exists $B_1\in B$ such that
                    $(p_0,E) = (p_1,B_1)$ and $p_1\in B_1$ and $B_1\subseteq N$.
                Take $B_1\in B$ such that $(p_0,E) = (p_1,B_1)$ and $p_1\in B_1$ and $B_1\subseteq N$.
                Thus $E = B_1$ by \cref{pair_eq_iff}.
                Thus $E\subseteq N$ by \cref{subseteq}.
                Thus $p\in N$ by \cref{subseteq}.
            \end{subproof}
            Thus $W\times Y\subseteq N$ by \cref{subseteq}.
        \end{subproof}
        Thus there exists $W$ such that
            $W$ is open in $X$ with respect to $T$ and
            $x_0\in W$ and
            $W\times Y\subseteq N$.
    \end{byCase}
\end{proof}

% from §26 helper lemma 3: union of a finite family of finite sets is finite
\begin{lemma}\label{finite_union_of_finite_family}
    Suppose $F$ is finite.
    Suppose for all $A$ such that $A\in F$ we have $A$ is finite.
    Then $\unions{F}$ is finite.
\end{lemma}
\begin{proof}
    Take $k$ such that $k$ is a natural number and there exists a bijection from $k$ to $F$ by \cref{finite}.
    Take $f$ such that $f$ is a bijection from $k$ to $F$.
    Let $A = \{ n\in\naturals \mid \text{for all $G,g$ such that $g$ is a bijection from $n$ to $G$ and for all $B\in G$ we have $B$ is finite we have $\unions{G}$ is finite} \}$.
    Show that $A$ is an inductive set.
    \begin{subproof}
        Show that $\emptyset\in A$.
        \begin{subproof}
            Show that for all $G,g$ such that $g$ is a bijection from $\emptyset$ to $G$ and
                for all $B\in G$ we have $B$ is finite we have $\unions{G}$ is finite.
            \begin{subproof}
                Fix $G$.
                Fix $g$.
                Assume $g$ is a bijection from $\emptyset$ to $G$.
                Assume for all $B\in G$ we have $B$ is finite.
                $g\in\surj{\emptyset}{G}$ by \cref{bijections}.
                Show that $G = \emptyset$.
                \begin{subproof}
                    Suppose not.
                    Then $G$ is inhabited by \cref{nonempty_inhabited}.
                    Take $z$ such that $z\in G$.
                    Thus there exists $y\in\emptyset$ such that $g(y) = z$ by \cref{surj}.
                    Contradiction by \cref{emptyset}.
                \end{subproof}
                Thus $\unions{G} = \emptyset$ by \cref{unions_emptyset}.
                Thus $\unions{G}$ is finite by \cref{emptyset_is_finite}.
            \end{subproof}
        \end{subproof}
        Show that for all $n\in A$ we have $\suc{n}\in A$.
        \begin{subproof}
            Fix $n$.
            Assume $n\in A$.
            Show that for all $G,g$ such that $g$ is a bijection from $\suc{n}$ to $G$
                and for all $B\in G$ we have $B$ is finite we have $\unions{G}$ is finite.
            \begin{subproof}
                Fix $G$.
                Fix $g$.
                Assume $g$ is a bijection from $\suc{n}$ to $G$.
                Assume for all $B\in G$ we have $B$ is finite.
                $g$ is a function by \cref{bijection_is_function}.
                $\dom{g} = \suc{n}$ by \cref{bijections_dom}.
                $n\in\suc{n}$ by \cref{suc_intro_self}.
                Thus $n\in\dom{g}$.
                Let $a = g(n)$.
                Let $G_1 = \img{g}{n}$.
                Show that for all $x$ such that $x\in n$ we have $x\in\suc{n}$.
                \begin{subproof}
                    Fix $x$.
                    Assume $x\in n$.
                    Thus $x\in\suc{n}$ by \cref{suc}.
                \end{subproof}
                Thus $n\subseteq \suc{n}$ by \cref{subseteq}.
                Thus $G_1\subseteq \img{g}{\suc{n}}$ by \cref{img_subseteq}.
                $\img{g}{\suc{n}} = \ran{g}$ by \cref{img_dom}.
                $g\in\surj{\suc{n}}{G}$ by \cref{bijections}.
                Thus $\ran{g} = G$ by \cref{ran_of_surj}.
                Thus $G_1\subseteq G$.
                $(n,a)\in g$ by \cref{function_apply_elim}.
                Thus $a\in\ran{g}$ by \cref{ran_iff}.
                Thus $a\in G$.
                Show that $G = G_1\union\{a\}$.
                \begin{subproof}
                    Show that $G\subseteq G_1\union\{a\}$.
                    \begin{subproof}
                        Show that for all $z$ such that $z\in G$ we have $z\in G_1\union\{a\}$.
                        \begin{subproof}
                            Fix $z$.
                            Assume $z\in G$.
                            $g\in\surj{\suc{n}}{G}$ by \cref{bijections}.
                            Thus there exists $x\in\suc{n}$ such that $g(x)=z$ by \cref{surj}.
                            $x\in n$ or $x = n$ by \cref{suc}.
                            \begin{byCase}
                            \caseOf{$x\in n$.}
                                $(x,z)\in g$ by \cref{function_apply_elim}.
                                Thus $z\in G_1$ by \cref{img_iff}.
                                Thus $z\in G_1\union\{a\}$ by \cref{union_intro_left}.
                            \caseOf{$x = n$.}
                                Thus $z = g(n)$.
                                Thus $z = a$.
                                Thus $z\in G_1\union\{a\}$ by \cref{union_intro_right,singleton_intro}.
                            \end{byCase}
                        \end{subproof}
                        Thus $G\subseteq G_1\union\{a\}$ by \cref{subseteq}.
                    \end{subproof}
                    $G_1\subseteq G$.
                    Thus $\{a\}\subseteq G$ by \cref{singleton_subset_intro}.
                    Thus $G_1\union\{a\}\subseteq G$ by \cref{union_subsets_is_subset}.
                    Thus $G = G_1\union\{a\}$ by \cref{subseteq_antisymmetric}.
                \end{subproof}
                Show that for all $B$ such that $B\in G_1$ we have $B$ is finite.
                \begin{subproof}
                    Fix $B$.
                    Assume $B\in G_1$.
                    Thus $B\in G$ by \cref{subseteq}.
                    Thus $B$ is finite.
                \end{subproof}
                Let $h = \restrl{g}{n}$.
                $h$ is a function by \cref{restrl_of_function_is_function}.
                Show that for all $x$ such that $x\in n$ we have $x\in\dom{g}$.
                \begin{subproof}
                    Fix $x$.
                    Assume $x\in n$.
                    Thus $x\in\suc{n}$ by \cref{suc}.
                    Thus $x\in\dom{g}$ by \cref{bijections_dom}.
                \end{subproof}
                Thus $n\subseteq\dom{g}$ by \cref{subseteq}.
                Thus $\dom{h} = \dom{g}\inter n$ by \cref{dom_of_restrl_eq_inter}.
                Thus $\dom{h} = n$ by \cref{inter_absorb_supseteq_left}.
                Show that $h$ is a function to $G_1$.
                \begin{subproof}
                    Show that for all $x\in\dom{h}$ we have $h(x)\in G_1$.
                    \begin{subproof}
                        Fix $x$.
                        Assume $x\in\dom{h}$.
                        Thus $x\in n$.
                        $(x,h(x))\in h$ by \cref{function_apply_elim}.
                        $h\subseteq g$ by \cref{restrl_subseteq}.
                        Thus $(x,h(x))\in g$ by \cref{elem_subseteq}.
                        Thus $h(x)\in G_1$ by \cref{img_iff}.
                    \end{subproof}
                    Thus $h$ is a function to $G_1$.
                \end{subproof}
                Thus $h\in\funs{n}{G_1}$ by \cref{funs_intro}.
                Show that for all $y\in G_1$ there exists $x\in n$ such that $h(x) = y$.
                \begin{subproof}
                    Fix $y$.
                    Assume $y\in G_1$.
                    Take $x\in n$ such that $(x,y)\in g$ by \cref{img_iff}.
                    Thus $x\in\dom{g}$ and $g(x) = y$ by \cref{function_apply_iff}.
                    Thus $h(x) = y$ by \cref{restrl_of_function_apply}.
                \end{subproof}
                Thus $h\in\surj{n}{G_1}$ by \cref{surj}.
                Show that $h$ is injective.
                \begin{subproof}
                    Show that for all $x,y$ such that $x\in\dom{h}$ and $y\in\dom{h}$ and $h(x) = h(y)$ we have $x = y$.
                    \begin{subproof}
                        Fix $x$.
                        Fix $y$.
                        Assume $x\in\dom{h}$ and $y\in\dom{h}$ and $h(x) = h(y)$.
                        Thus $x\in n$ and $y\in n$.
                        Thus $h(x) = g(x)$ and $h(y) = g(y)$ by \cref{restrl_of_function_apply}.
                        Thus $g(x) = g(y)$.
                        $g$ is injective by \cref{bijections}.
                        Thus $x = y$ by \cref{injective_function}.
                    \end{subproof}
                    Thus $h$ is injective by \cref{injective_function}.
                \end{subproof}
                Thus $h$ is a bijection from $n$ to $G_1$ by \cref{bijections}.
                Thus $\unions{G_1}$ is finite.
                Thus $a$ is finite.
                Thus $G = \{a\}\union G_1$ by \cref{union_comm}.
                Thus $G = \cons{a}{G_1}$ by \cref{union_eq_cons}.
                Thus $\unions{G} = a\union\unions{G_1}$ by \cref{unions_cons}.
                Thus $\unions{G}$ is finite by \cref{union_of_finite_is_finite}.
            \end{subproof}
        \end{subproof}
    \end{subproof}
    Thus $A$ is an inductive set.
    Thus $\naturals\subseteq A$ by \cref{num_naturals_smallest_inductive_set}.
    Thus $k\in\naturals$.
    Thus $k\in A$ by \cref{subseteq}.
    Thus for all $G,g$ such that $g$ is a bijection from $k$ to $G$ and for all $B\in G$ we have $B$ is finite
        we have $\unions{G}$ is finite.
    Thus $\unions{F}$ is finite.
\end{proof}

% from §26 Item 10: product of finitely many compact spaces is compact
\begin{theorem}\label{product_compact}
    Assume $T$ is a topology on $X$.
    Assume $T_1$ is a topology on $Y$.
    Assume $X$ is compact with respect to $T$.
    Assume $Y$ is compact with respect to $T_1$.
    Let $T_2 = \productTopology{T}{T_1}{X}{Y}$.
    Then $X\times Y$ is compact with respect to $T_2$.
\end{theorem}
\begin{proof}
    Let $T_2 = \productTopology{T}{T_1}{X}{Y}$.
    Let $B_0 = \productBasis{T}{T_1}{X}{Y}$.
    $B_0$ is a basis on $X\times Y$ by \cref{product_basis_is_basis}.
    Thus $T_2$ is a topology on $X\times Y$ by \cref{basis_topology_is_topology,product_topology}.
    Show that for all $A$ such that $A$ is an open covering of $X\times Y$ with respect to $T_2$
        there exists $B$ such that
            $B\subseteq A$ and
            $B$ is finite and
            $B$ is a covering of $X\times Y$.
    \begin{subproof}
        Fix $A$.
        Assume $A$ is an open covering of $X\times Y$ with respect to $T_2$.
        Let $C = \{W\in T \mid \text{there exists $B\subseteq A$ such that $B$ is finite and $W\times Y\subseteq \unions{B}$}\}$.
        Show that $C$ is an open covering of $X$ with respect to $T$.
        \begin{subproof}
            Show that for all $W$ such that $W\in C$ we have $W$ is open in $X$ with respect to $T$.
            \begin{subproof}
                Fix $W$.
                Assume $W\in C$.
                Thus $W\in T$.
                Thus $W$ is open in $X$ with respect to $T$ by \cref{open_in}.
            \end{subproof}
            Show that $C$ is a covering of $X$.
            \begin{subproof}
                Show that for all $x$ such that $x\in X$ we have $x\in\unions{C}$.
                \begin{subproof}
                    Fix $x$.
                    Assume $x\in X$.
                    Let $S = \{x\}\times Y$.
                    Let $T_S = \subspaceTopology{T_2}{S}$.
                    Thus $S$ is compact with respect to $T_S$ by \cref{slice_compact}.
                    $T_2$ is a topology on $X\times Y$ by \cref{basis_topology_is_topology,product_topology}.
                    Show that $S\subseteq X\times Y$.
                    \begin{subproof}
                        Show that for all $p$ such that $p\in S$ we have $p\in X\times Y$.
                        \begin{subproof}
                            Fix $p$.
                            Assume $p\in S$.
                            Take $x_0,y$ such that $p = (x_0,y)$ and $x_0\in\{x\}$ and $y\in Y$
                                by \cref{times_elem_is_tuple}.
                            Thus $x_0 = x$ by \cref{singleton_iff}.
                            Thus $x_0\in X$.
                            Thus $p\in X\times Y$ by \cref{times_tuple_intro}.
                        \end{subproof}
                        Thus $S\subseteq X\times Y$ by \cref{subseteq}.
                    \end{subproof}
                    Thus $S$ is a subspace of $X\times Y$ with respect to $T_2$ by \cref{subspace_of}.
                    Show that $A$ is a covering of $S$.
                    \begin{subproof}
                        $A$ is a covering of $X\times Y$ by \cref{open_covering}.
                        Thus $X\times Y\subseteq\unions{A}$ by \cref{covering}.
                        Thus $S\subseteq\unions{A}$ by \cref{subseteq}.
                        Thus $A$ is a covering of $S$ by \cref{covering}.
                    \end{subproof}
                    Show that for all $U$ such that $U\in A$ we have $U$ is open in $X\times Y$ with respect to $T_2$.
                    \begin{subproof}
                        Fix $U$.
                        Assume $U\in A$.
                        Thus $U$ is open in $X\times Y$ with respect to $T_2$ by \cref{open_covering}.
                    \end{subproof}
                    Thus there exists $B$ such that
                        $B\subseteq A$ and $B$ is finite and $B$ is a covering of $S$
                        by \cref{compact_subspace_open_covering}.
                    Take $B$ such that $B\subseteq A$ and $B$ is finite and $B$ is a covering of $S$.
                    Let $N = \unions{B}$.
                    Show that $N$ is open in $X\times Y$ with respect to $T_2$.
                    \begin{subproof}
                        Show that $B\subseteq T_2$.
                        \begin{subproof}
                            Show that for all $U$ such that $U\in B$ we have $U\in T_2$.
                            \begin{subproof}
                                Fix $U$.
                                Assume $U\in B$.
                                Thus $U\in A$ by \cref{subseteq}.
                                Thus $U$ is open in $X\times Y$ with respect to $T_2$ by \cref{open_covering}.
                                Thus $U\in T_2$ by \cref{open_in}.
                            \end{subproof}
                            Thus $B\subseteq T_2$ by \cref{subseteq}.
                        \end{subproof}
                        Thus $N\in T_2$ by \cref{topology_on}.
                        Thus $N$ is open in $X\times Y$ with respect to $T_2$ by \cref{open_in}.
                    \end{subproof}
                    $S\subseteq N$ by \cref{covering}.
                    Thus $\{x\}\times Y\subseteq N$.
                    Thus there exists $W$ such that
                        $W$ is open in $X$ with respect to $T$ and
                        $x\in W$ and
                        $W\times Y\subseteq N$ by \cref{tube_lemma_helper}.
                    Take $W$ such that
                        $W$ is open in $X$ with respect to $T$ and
                        $x\in W$ and
                        $W\times Y\subseteq N$.
                    Thus $W\in T$ by \cref{open_in}.
                    Thus $B\subseteq A$ and $B$ is finite and $W\times Y\subseteq \unions{B}$.
                    Thus $W\in C$.
                    Thus $x\in\unions{C}$ by \cref{unions_iff}.
                \end{subproof}
                Thus $X\subseteq\unions{C}$ by \cref{subseteq}.
                Thus $C$ is a covering of $X$ by \cref{covering}.
            \end{subproof}
            Thus $C$ is an open covering of $X$ with respect to $T$ by \cref{open_covering}.
        \end{subproof}
        Thus there exists $F$ such that
            $F\subseteq C$ and $F$ is finite and $F$ is a covering of $X$
            by \cref{compact}.
        Take $F$ such that $F\subseteq C$ and $F$ is finite and $F$ is a covering of $X$.
        Let $R = \{w\in F\times\pow{A} \mid \text{there exists $W\in F$ such that there exists $B\subseteq A$ such that $w = (W,B)$ and $B$ is finite and $W\times Y\subseteq \unions{B}$}\}$.
        Show that $R$ is a relation.
        \begin{subproof}
            Show that for all $w$ such that $w\in R$ there exist $W,B$ such that $w = (W,B)$.
            \begin{subproof}
                Fix $w$.
                Assume $w\in R$.
                Take $W\in F$ such that there exists $B\subseteq A$ such that
                    $w = (W,B)$ and $B$ is finite and $W\times Y\subseteq \unions{B}$.
                Take $B\subseteq A$ such that $w = (W,B)$ and $B$ is finite and $W\times Y\subseteq \unions{B}$.
            \end{subproof}
            Thus $R$ is a relation by \cref{relation}.
        \end{subproof}
        Show that for all $W\in F$ there exists $B$ such that $W\mathrel{R} B$.
        \begin{subproof}
            Fix $W$.
            Assume $W\in F$.
            Thus $W\in C$ by \cref{subseteq}.
            Take $B\subseteq A$ such that $B$ is finite and $W\times Y\subseteq \unions{B}$.
            Thus $(W,B)\in R$.
            Thus $W\mathrel{R} B$.
        \end{subproof}
        Take $h$ such that $h$ is a function on $F$ and for all $W\in F$ we have $W\mathrel{R}\apply{h}{W}$
            by \cref{choice_function}.
        Let $H = \img{h}{F}$.
        Thus $H$ is finite by \cref{finite_image_function}.
        Show that for all $B$ such that $B\in H$ we have $B$ is finite.
        \begin{subproof}
            Fix $B$.
            Assume $B\in H$.
            Take $W\in F$ such that $(W,B)\in h$ by \cref{img_iff}.
            Thus $W\mathrel{R} B$ by \cref{function_apply_iff}.
            Take $W_1\in F$ such that there exists $B_1\subseteq A$ such that
                $(W,B) = (W_1,B_1)$ and $B_1$ is finite and $W_1\times Y\subseteq \unions{B_1}$.
            Take $B_1\subseteq A$ such that $(W,B) = (W_1,B_1)$ and $B_1$ is finite and $W_1\times Y\subseteq \unions{B_1}$.
            Thus $B = B_1$ by \cref{pair_eq_iff}.
            Thus $B$ is finite.
        \end{subproof}
        Let $B = \unions{H}$.
        Thus $B$ is finite by \cref{finite_union_of_finite_family}.
        Show that $B\subseteq A$.
        \begin{subproof}
            Show that $H\subseteq \pow{A}$.
            \begin{subproof}
                Show that for all $B_1$ such that $B_1\in H$ we have $B_1\in\pow{A}$.
                \begin{subproof}
                    Fix $B_1$.
                    Assume $B_1\in H$.
                    Take $W\in F$ such that $(W,B_1)\in h$ by \cref{img_iff}.
                    Thus $W\mathrel{R} B_1$ by \cref{function_apply_iff}.
                    Take $W_1\in F$ such that there exists $B_2\subseteq A$ such that
                        $(W,B_1) = (W_1,B_2)$ and $B_2$ is finite and $W_1\times Y\subseteq \unions{B_2}$.
                    Take $B_2\subseteq A$ such that
                        $(W,B_1) = (W_1,B_2)$ and $B_2$ is finite and $W_1\times Y\subseteq \unions{B_2}$.
                    Thus $B_1 = B_2$ by \cref{pair_eq_iff}.
                    Thus $B_1\subseteq A$.
                    Thus $B_1\in\pow{A}$ by \cref{powerset_intro}.
                \end{subproof}
                Thus $H\subseteq \pow{A}$ by \cref{subseteq}.
            \end{subproof}
            Thus $B\subseteq A$ by \cref{unions_subseteq_of_powerset_is_subseteq}.
        \end{subproof}
        Show that $B$ is a covering of $X\times Y$.
        \begin{subproof}
            Show that for all $p$ such that $p\in X\times Y$ we have $p\in\unions{B}$.
            \begin{subproof}
                Fix $p$.
                Assume $p\in X\times Y$.
                Take $x,y$ such that $p = (x,y)$ and $x\in X$ and $y\in Y$
                    by \cref{times_elem_is_tuple}.
                $X\subseteq\unions{F}$ by \cref{covering}.
                Thus $x\in\unions{F}$ by \cref{subseteq}.
                Take $W\in F$ such that $x\in W$ by \cref{unions_iff}.
                Thus $p\in W\times Y$ by \cref{times_tuple_intro}.
                Thus $W\mathrel{R}\apply{h}{W}$.
                Take $W_1\in F$ such that there exists $B_1\subseteq A$ such that
                    $(W,\apply{h}{W}) = (W_1,B_1)$ and $B_1$ is finite and $W_1\times Y\subseteq \unions{B_1}$.
                Take $B_1\subseteq A$ such that
                    $(W,\apply{h}{W}) = (W_1,B_1)$ and $B_1$ is finite and $W_1\times Y\subseteq \unions{B_1}$.
                Thus $W = W_1$ and $\apply{h}{W} = B_1$ by \cref{pair_eq_iff}.
                Thus $W\times Y\subseteq \unions{\apply{h}{W}}$.
                Thus $p\in\unions{\apply{h}{W}}$ by \cref{subseteq}.
                Take $U\in\apply{h}{W}$ such that $p\in U$ by \cref{unions_iff}.
                Thus $(W,\apply{h}{W})\in h$ by \cref{function_apply_elim}.
                Thus $\apply{h}{W}\in H$ by \cref{img_iff}.
                Thus $U\in B$ by \cref{unions_iff}.
                Thus $p\in\unions{B}$ by \cref{unions_iff}.
            \end{subproof}
            Thus $X\times Y\subseteq\unions{B}$ by \cref{subseteq}.
            Thus $B$ is a covering of $X\times Y$ by \cref{covering}.
        \end{subproof}
    \end{subproof}
    Thus $X\times Y$ is compact with respect to $T_2$ by \cref{compact}.
\end{proof}

% from §26 Item 11: tube lemma
\begin{lemma}\label{tube_lemma}
    Assume $T$ is a topology on $X$.
    Assume $T_1$ is a topology on $Y$.
    Assume $Y$ is compact with respect to $T_1$.
    Let $T_2 = \productTopology{T}{T_1}{X}{Y}$.
    Assume $N$ is open in $X\times Y$ with respect to $T_2$.
    Assume $x_0\in X$.
    Suppose $\{x_0\}\times Y\subseteq N$.
    Then there exists $W$ such that
        $W$ is open in $X$ with respect to $T$ and
        $x_0\in W$ and
        $W\times Y\subseteq N$.
\end{lemma}
\begin{proof}
    Follows by \cref{tube_lemma_helper}.
\end{proof}

% from §26 Item 12: finite intersection property
\begin{definition}\label{finite_intersection_property}
    $C$ has the finite intersection property iff
        $C$ is inhabited and
        for all $F$ such that $F\subseteq C$ and $F$ is finite and inhabited
        we have $\inters{F}\neq\emptyset$.
\end{definition}

% from §26 Item 13: compactness via finite intersection property
\begin{theorem}\label{compact_iff_finite_intersection_property}
    Assume $T$ is a topology on $X$.
    Then $X$ is compact with respect to $T$ iff
        for all $C$ if
            $C$ has the finite intersection property and
            for all $A$ such that $A\in C$ we have $A$ is closed in $X$ with respect to $T$
        then $\inters{C}\neq\emptyset$.
\end{theorem}
\begin{proof}
    Show that if $X$ is compact with respect to $T$ then
        for all $C$ if
            $C$ has the finite intersection property and
            for all $A$ such that $A\in C$ we have $A$ is closed in $X$ with respect to $T$
        then $\inters{C}\neq\emptyset$.
    \begin{subproof}
        Assume $X$ is compact with respect to $T$.
        Fix $C$.
        Assume $C$ has the finite intersection property and
            for all $A$ such that $A\in C$ we have $A$ is closed in $X$ with respect to $T$.
        Show that $\inters{C}\neq\emptyset$.
        \begin{subproof}
            Suppose not.
            Then $\inters{C} = \emptyset$ by \cref{emptyset}.
            $C$ is inhabited by \cref{finite_intersection_property}.
            Take $A_0$ such that $A_0\in C$.
            Let $F_0 = \{A_0\}$.
            Show that $F_0$ is finite.
            \begin{subproof}
                Let $G_0 = \{A_0,A_0\}$.
                $G_0$ is finite by \cref{pair_is_finite}.
                Thus $A_0\in G_0$ by \cref{upair_intro_left}.
                Thus $F_0\subseteq G_0$ by \cref{singleton_subset_intro}.
                Thus $F_0$ is finite by \cref{subseteq_finite_is_finite}.
            \end{subproof}
            $F_0\subseteq C$ by \cref{singleton_subset_intro}.
            $F_0$ is inhabited.
            Thus $\inters{F_0}\neq\emptyset$ by \cref{finite_intersection_property}.
            Thus $\inters{F_0} = A_0$ by \cref{inters_singleton}.
            Thus $A_0\neq\emptyset$ by \cref{emptyset}.
            Thus $A_0\subseteq X$ by \cref{closed_in}.
            Take $x_0$ such that $x_0\in A_0$ by \cref{nonempty_inhabited}.
            Thus $x_0\in X$ by \cref{subseteq}.
            Thus $X$ is inhabited.
            Let $A = \{U \mid \exists D\in C. U = X\setminus D\}$.
            Show that $A$ is an open covering of $X$ with respect to $T$.
            \begin{subproof}
                Show that for all $U$ such that $U\in A$ we have $U$ is open in $X$ with respect to $T$.
                \begin{subproof}
                    Fix $U$.
                    Assume $U\in A$.
                    Take $D\in C$ such that $U = X\setminus D$.
                    $D$ is closed in $X$ with respect to $T$.
                    Thus $X\setminus D$ is open in $X$ with respect to $T$ by \cref{closed_in}.
                    Thus $U$ is open in $X$ with respect to $T$.
                \end{subproof}
                Show that $A$ is a covering of $X$.
                \begin{subproof}
                    Show that for all $x$ such that $x\in X$ we have $x\in\unions{A}$.
                    \begin{subproof}
                        Fix $x$.
                        Assume $x\in X$.
                        Suppose not.
                        Then $x\notin\unions{A}$.
                        Show that for all $D$ such that $D\in C$ we have $x\in D$.
                        \begin{subproof}
                            Fix $D$.
                            Assume $D\in C$.
                            Suppose not.
                            Then $x\notin D$.
                            Thus $x\in X\setminus D$ by \cref{setminus_intro}.
                            Thus $X\setminus D\in A$.
                            Thus $x\in\unions{A}$ by \cref{unions_iff}.
                            Contradiction.
                        \end{subproof}
                        Thus $x\in\inters{C}$ by \cref{inters_intro}.
                        Thus $\inters{C}\neq\emptyset$ by \cref{emptyset}.
                        Contradiction.
                    \end{subproof}
                    Thus $X\subseteq\unions{A}$ by \cref{subseteq}.
                    Thus $A$ is a covering of $X$ by \cref{covering}.
                \end{subproof}
                Thus $A$ is an open covering of $X$ with respect to $T$ by \cref{open_covering}.
            \end{subproof}
            Thus there exists $B$ such that
                $B\subseteq A$ and $B$ is finite and $B$ is a covering of $X$
                by \cref{compact}.
            Take $B$ such that $B\subseteq A$ and $B$ is finite and $B$ is a covering of $X$.
            Show that $B$ is inhabited.
            \begin{subproof}
                Take $x_1$ such that $x_1\in X$ by \cref{nonempty_inhabited}.
                $X\subseteq\unions{B}$ by \cref{covering}.
                Thus $x_1\in\unions{B}$ by \cref{subseteq}.
                Take $U_1\in B$ such that $x_1\in U_1$ by \cref{unions_iff}.
                Thus $B$ is inhabited.
            \end{subproof}
            Let $R = \{w\in B\times C \mid \exists U\in B.\exists D\in C. w = (U,D) \land U = X\setminus D\}$.
            Show that $R$ is a relation.
            \begin{subproof}
                Show that for all $w$ such that $w\in R$ there exist $U,D$ such that $w = (U,D)$.
                \begin{subproof}
                    Fix $w$.
                    Assume $w\in R$.
                    Take $U\in B$ such that there exists $D\in C$ such that $w = (U,D)$ and $U = X\setminus D$.
                    Take $D\in C$ such that $w = (U,D)$ and $U = X\setminus D$.
                \end{subproof}
                Thus $R$ is a relation by \cref{relation}.
            \end{subproof}
            Show that for all $U\in B$ there exists $D$ such that $U\mathrel{R} D$.
            \begin{subproof}
                Fix $U$.
                Assume $U\in B$.
                Thus $U\in A$ by \cref{subseteq}.
                Take $D\in C$ such that $U = X\setminus D$.
                Thus $(U,D)\in R$.
                Thus $U\mathrel{R} D$.
            \end{subproof}
            Take $h$ such that $h$ is a function on $B$ and for all $U\in B$ we have $U\mathrel{R}\apply{h}{U}$
                by \cref{choice_function}.
            Let $F = \img{h}{B}$.
            Thus $F$ is finite by \cref{finite_image_function}.
            Show that $F$ is inhabited.
            \begin{subproof}
                Take $U_0$ such that $U_0\in B$ by \cref{nonempty_inhabited}.
                Thus $(U_0,\apply{h}{U_0})\in h$ by \cref{function_apply_elim}.
                Thus $\apply{h}{U_0}\in F$ by \cref{img_iff}.
                Thus $F$ is inhabited.
            \end{subproof}
            Show that $F\subseteq C$.
            \begin{subproof}
                Show that for all $D$ such that $D\in F$ we have $D\in C$.
                \begin{subproof}
                    Fix $D$.
                    Assume $D\in F$.
                    Take $U\in B$ such that $(U,D)\in h$ by \cref{img_iff}.
                    Thus $U\mathrel{R} D$ by \cref{function_apply_iff}.
                    Take $U_1\in B$ such that there exists $D_1\in C$ such that
                        $(U,D) = (U_1,D_1)$ and $U_1 = X\setminus D_1$.
                    Take $D_1\in C$ such that $(U,D) = (U_1,D_1)$ and $U_1 = X\setminus D_1$.
                    Thus $D = D_1$ by \cref{pair_eq_iff}.
                    Thus $D\in C$.
                \end{subproof}
                Thus $F\subseteq C$ by \cref{subseteq}.
            \end{subproof}
            Show that $X\setminus\unions{B} = \inters{F}$.
            \begin{subproof}
                Show that for all $x$ such that $x\in X\setminus\unions{B}$ we have $x\in\inters{F}$.
                \begin{subproof}
                    Fix $x$.
                    Assume $x\in X\setminus\unions{B}$.
                    Thus $x\in X$ by \cref{setminus_elim_left}.
                    Thus $x\notin\unions{B}$ by \cref{setminus_elim_right}.
                    Show that for all $D$ such that $D\in F$ we have $x\in D$.
                    \begin{subproof}
                        Fix $D$.
                        Assume $D\in F$.
                        Take $U\in B$ such that $(U,D)\in h$ by \cref{img_iff}.
                        Thus $U\mathrel{R} D$ by \cref{function_apply_iff}.
                        Take $U_1\in B$ such that there exists $D_1\in C$ such that
                            $(U,D) = (U_1,D_1)$ and $U_1 = X\setminus D_1$.
                        Take $D_1\in C$ such that $(U,D) = (U_1,D_1)$ and $U_1 = X\setminus D_1$.
                        Thus $U = U_1$ and $D = D_1$ by \cref{pair_eq_iff}.
                        Suppose not.
                        Then $x\notin D$.
                        Thus $x\in X\setminus D$ by \cref{setminus_intro}.
                        Thus $x\in U$.
                        Thus $x\in\unions{B}$ by \cref{unions_iff}.
                        Contradiction.
                    \end{subproof}
                    Thus $x\in\inters{F}$ by \cref{inters_intro}.
                \end{subproof}
                Thus $X\setminus\unions{B}\subseteq\inters{F}$ by \cref{subseteq}.
                Show that for all $x$ such that $x\in\inters{F}$ we have $x\in X\setminus\unions{B}$.
                \begin{subproof}
                    Fix $x$.
                    Assume $x\in\inters{F}$.
                    Take $D_0$ such that $D_0\in F$ by \cref{nonempty_inhabited}.
                    Thus $D_0\in C$ by \cref{subseteq}.
                    Thus $D_0$ is closed in $X$ with respect to $T$.
                    Thus $D_0\subseteq X$ by \cref{closed_in}.
                    Thus $x\in D_0$ by \cref{inters_destr}.
                    Thus $x\in X$ by \cref{subseteq}.
                    Suppose not.
                    Then $x\in\unions{B}$.
                    Take $U\in B$ such that $x\in U$ by \cref{unions_iff}.
                    Thus $U\mathrel{R}\apply{h}{U}$.
                    Take $U_1\in B$ such that there exists $D_1\in C$ such that
                        $(U,\apply{h}{U}) = (U_1,D_1)$ and $U_1 = X\setminus D_1$.
                    Take $D_1\in C$ such that $(U,\apply{h}{U}) = (U_1,D_1)$ and $U_1 = X\setminus D_1$.
                    Thus $U = U_1$ and $\apply{h}{U} = D_1$ by \cref{pair_eq_iff}.
                    Thus $D_1\in F$ by \cref{img_iff,function_apply_elim}.
                    Thus $x\in D_1$ by \cref{inters_destr}.
                    Thus $x\in X\setminus D_1$ by \cref{setminus_intro}.
                    Thus $x\notin D_1$ by \cref{setminus_elim_right}.
                    Contradiction.
                \end{subproof}
                Thus $\inters{F}\subseteq X\setminus\unions{B}$ by \cref{subseteq}.
                Thus $X\setminus\unions{B} = \inters{F}$ by \cref{subseteq_antisymmetric}.
            \end{subproof}
            Show that $\unions{B} = X$.
            \begin{subproof}
                $X\subseteq\unions{B}$ by \cref{covering}.
                Show that $\unions{B}\subseteq X$.
                \begin{subproof}
                    Show that for all $U$ such that $U\in B$ we have $U\subseteq X$.
                    \begin{subproof}
                        Fix $U$.
                        Assume $U\in B$.
                        Thus $U\in A$ by \cref{subseteq}.
                        Take $D\in C$ such that $U = X\setminus D$.
                        Thus $U\subseteq X$ by \cref{setminus_subseteq}.
                    \end{subproof}
                    Thus $B\subseteq\pow{X}$ by \cref{powerset_intro,subseteq}.
                    Thus $\unions{B}\subseteq X$ by \cref{unions_subseteq_of_powerset_is_subseteq}.
                \end{subproof}
                Thus $\unions{B} = X$ by \cref{subseteq_antisymmetric}.
            \end{subproof}
            Thus $X\setminus\unions{B} = \emptyset$ by \cref{setminus_self}.
            Thus $\inters{F} = \emptyset$ by \cref{subseteq_antisymmetric,subseteq_refl}.
            $F$ is inhabited and $F$ is finite.
            Thus $\inters{F}\neq\emptyset$ by \cref{finite_intersection_property}.
            Contradiction.
        \end{subproof}
    \end{subproof}
    Show that if
        for all $C$ if
            $C$ has the finite intersection property and
            for all $A$ such that $A\in C$ we have $A$ is closed in $X$ with respect to $T$
        then $\inters{C}\neq\emptyset$
        then $X$ is compact with respect to $T$.
    \begin{subproof}
        Assume for all $C$ if
            $C$ has the finite intersection property and
            for all $A$ such that $A\in C$ we have $A$ is closed in $X$ with respect to $T$
        then $\inters{C}\neq\emptyset$.
        Show that $X$ is compact with respect to $T$.
        \begin{subproof}
            Show that for all $A$ such that $A$ is an open covering of $X$ with respect to $T$
                there exists $B$ such that
                    $B\subseteq A$ and
                    $B$ is finite and
                    $B$ is a covering of $X$.
            \begin{subproof}
                Fix $A$.
                Assume $A$ is an open covering of $X$ with respect to $T$.
                \begin{byCase}
                \caseOf{$X = \emptyset$.}
                    Let $B = \emptyset$.
                    $B\subseteq A$ by \cref{emptyset_subseteq}.
                    $B$ is finite by \cref{emptyset_is_finite}.
                    $\unions{B} = \emptyset$ by \cref{unions_emptyset}.
                    Thus $X\subseteq\unions{B}$ by \cref{subseteq_refl}.
                    Thus $B$ is a covering of $X$ by \cref{covering}.
                    Thus there exists $B$ such that
                        $B\subseteq A$ and $B$ is finite and $B$ is a covering of $X$.
                \caseOf{$X \neq \emptyset$.}
                    Show that there exists $B$ such that
                        $B\subseteq A$ and $B$ is finite and $B$ is a covering of $X$.
                    \begin{subproof}
                        Suppose not.
                        Then for all $B$ such that $B\subseteq A$ and $B$ is finite we have
                            $B$ is not a covering of $X$.
                        $A$ is a covering of $X$ by \cref{open_covering}.
                        Let $C = \{D \mid \exists U\in A. D = X\setminus U\}$.
                        Show that for all $D$ such that $D\in C$ we have $D$ is closed in $X$ with respect to $T$.
                        \begin{subproof}
                            Fix $D$.
                            Assume $D\in C$.
                            Take $U\in A$ such that $D = X\setminus U$.
                            $U$ is open in $X$ with respect to $T$ by \cref{open_covering}.
                            Thus $U\in T$ by \cref{open_in}.
                            $T\subseteq\pow{X}$ by \cref{topology_on}.
                            Thus $U\in\pow{X}$ by \cref{subseteq}.
                            Thus $U\subseteq X$ by \cref{pow_iff}.
                            Thus $X\setminus (X\setminus U) = U$ by \cref{double_relative_complement}.
                            Thus $X\setminus (X\setminus U)$ is open in $X$ with respect to $T$.
                            Thus $X\setminus U\subseteq X$ by \cref{setminus_subseteq}.
                            Thus $X\setminus U$ is closed in $X$ with respect to $T$ by \cref{closed_in}.
                            Thus $D$ is closed in $X$ with respect to $T$.
                        \end{subproof}
                        Show that $C$ has the finite intersection property.
                        \begin{subproof}
                            Show that $C$ is inhabited.
                            \begin{subproof}
                                Show that $A$ is inhabited.
                                \begin{subproof}
                                    Take $x_0$ such that $x_0\in X$ by \cref{nonempty_inhabited}.
                                    $X\subseteq\unions{A}$ by \cref{covering}.
                                    Thus $x_0\in\unions{A}$ by \cref{subseteq}.
                                    Take $U_0\in A$ such that $x_0\in U_0$ by \cref{unions_iff}.
                                    Thus $A$ is inhabited.
                                \end{subproof}
                                Take $U_1$ such that $U_1\in A$ by \cref{nonempty_inhabited}.
                                Thus $X\setminus U_1\in C$.
                                Thus $C$ is inhabited.
                            \end{subproof}
                            Show that for all $F$ such that $F\subseteq C$ and $F$ is finite and inhabited
                                we have $\inters{F}\neq\emptyset$.
                            \begin{subproof}
                                Fix $F$.
                                Assume $F\subseteq C$ and $F$ is finite and inhabited.
                                Suppose not.
                                Then $\inters{F} = \emptyset$ by \cref{emptyset}.
                                Let $R_0 = \{w\in F\times A \mid \exists D\in F.\exists U\in A. w = (D,U) \land D = X\setminus U\}$.
                                Show that $R_0$ is a relation.
                                \begin{subproof}
                                    Show that for all $w$ such that $w\in R_0$ there exist $D,U$ such that $w = (D,U)$.
                                    \begin{subproof}
                                        Fix $w$.
                                        Assume $w\in R_0$.
                                        Take $D\in F$ such that there exists $U\in A$ such that $w = (D,U)$ and $D = X\setminus U$.
                                        Take $U\in A$ such that $w = (D,U)$ and $D = X\setminus U$.
                                    \end{subproof}
                                    Thus $R_0$ is a relation by \cref{relation}.
                                \end{subproof}
                                Show that for all $D\in F$ there exists $U$ such that $D\mathrel{R_0} U$.
                                \begin{subproof}
                                    Fix $D$.
                                    Assume $D\in F$.
                                    Thus $D\in C$ by \cref{subseteq}.
                                    Take $U\in A$ such that $D = X\setminus U$.
                                    Thus $(D,U)\in R_0$.
                                    Thus $D\mathrel{R_0} U$.
                                \end{subproof}
                                Take $h$ such that $h$ is a function on $F$ and for all $D\in F$ we have $D\mathrel{R_0}\apply{h}{D}$
                                    by \cref{choice_function}.
                                Let $B = \img{h}{F}$.
                                Thus $B$ is finite by \cref{finite_image_function}.
                                Show that $B\subseteq A$.
                                \begin{subproof}
                                    Show that for all $U$ such that $U\in B$ we have $U\in A$.
                                    \begin{subproof}
                                        Fix $U$.
                                        Assume $U\in B$.
                                        Take $D\in F$ such that $(D,U)\in h$ by \cref{img_iff}.
                                        Thus $D\mathrel{R_0} U$ by \cref{function_apply_iff}.
                                        Take $D_1\in F$ such that there exists $U_1\in A$ such that
                                            $(D,U) = (D_1,U_1)$ and $D_1 = X\setminus U_1$.
                                        Take $U_1\in A$ such that $(D,U) = (D_1,U_1)$ and $D_1 = X\setminus U_1$.
                                        Thus $U = U_1$ by \cref{pair_eq_iff}.
                                        Thus $U\in A$.
                                    \end{subproof}
                                    Thus $B\subseteq A$ by \cref{subseteq}.
                                \end{subproof}
                                Show that $X\setminus\unions{B} = \inters{F}$.
                                \begin{subproof}
                                    Show that for all $x$ such that $x\in X\setminus\unions{B}$ we have $x\in\inters{F}$.
                                    \begin{subproof}
                                        Fix $x$.
                                        Assume $x\in X\setminus\unions{B}$.
                                        Thus $x\in X$ by \cref{setminus_elim_left}.
                                        Thus $x\notin\unions{B}$ by \cref{setminus_elim_right}.
                                        Show that for all $D$ such that $D\in F$ we have $x\in D$.
                                        \begin{subproof}
                                            Fix $D$.
                                            Assume $D\in F$.
                                            Thus $D\mathrel{R_0}\apply{h}{D}$.
                                            Take $D_1\in F$ such that there exists $U_1\in A$ such that
                                                $(D,\apply{h}{D}) = (D_1,U_1)$ and $D_1 = X\setminus U_1$.
                                            Take $U_1\in A$ such that
                                                $(D,\apply{h}{D}) = (D_1,U_1)$ and $D_1 = X\setminus U_1$.
                                            Thus $D = D_1$ and $\apply{h}{D} = U_1$ by \cref{pair_eq_iff}.
                                            Suppose not.
                                            Then $x\notin D$.
                                            Thus $x\in X\setminus D$ by \cref{setminus_intro}.
                                            Thus $x\in U_1$.
                                            Thus $(D,\apply{h}{D})\in h$ by \cref{function_apply_elim}.
                                            Thus $\apply{h}{D}\in B$ by \cref{img_iff}.
                                            Thus $U_1\in B$ by \cref{pair_eq_iff}.
                                            Thus $x\in\unions{B}$ by \cref{unions_iff}.
                                            Contradiction.
                                        \end{subproof}
                                        Thus $x\in\inters{F}$ by \cref{inters_intro}.
                                    \end{subproof}
                                    Thus $X\setminus\unions{B}\subseteq\inters{F}$ by \cref{subseteq}.
                                    Show that for all $x$ such that $x\in\inters{F}$ we have $x\in X\setminus\unions{B}$.
                                    \begin{subproof}
                                        Fix $x$.
                                        Assume $x\in\inters{F}$.
                                        Take $D_0$ such that $D_0\in F$ by \cref{nonempty_inhabited}.
                                        Thus $x\in D_0$ by \cref{inters_destr}.
                                        Thus $D_0\in C$ by \cref{subseteq}.
                                        Thus $D_0\subseteq X$ by \cref{closed_in}.
                                        Thus $x\in X$ by \cref{subseteq}.
                                        Suppose not.
                                        Then $x\in\unions{B}$.
                                        Take $U\in B$ such that $x\in U$ by \cref{unions_iff}.
                                        Take $D\in F$ such that $(D,U)\in h$ by \cref{img_iff}.
                                        Thus $D\mathrel{R_0} U$ by \cref{function_apply_iff}.
                                        Take $D_1\in F$ such that there exists $U_1\in A$ such that
                                            $(D,U) = (D_1,U_1)$ and $D_1 = X\setminus U_1$.
                                        Take $U_1\in A$ such that $(D,U) = (D_1,U_1)$ and $D_1 = X\setminus U_1$.
                                        Thus $D = D_1$ and $U = U_1$ by \cref{pair_eq_iff}.
                                        $U$ is open in $X$ with respect to $T$ by \cref{open_covering}.
                                        Thus $U\in T$ by \cref{open_in}.
                                        $T\subseteq\pow{X}$ by \cref{topology_on}.
                                        Thus $U\in\pow{X}$ by \cref{subseteq}.
                                        Thus $U\subseteq X$ by \cref{pow_iff}.
                                        Thus $X\setminus D = U$ by \cref{double_relative_complement}.
                                        Thus $U\subseteq X\setminus D$ by \cref{subseteq_refl}.
                                        Thus $x\in X\setminus D$ by \cref{subseteq}.
                                        Thus $x\notin D$ by \cref{setminus_elim_right}.
                                        Thus $x\in D$ by \cref{inters_destr}.
                                        Contradiction.
                                    \end{subproof}
                                    Thus $\inters{F}\subseteq X\setminus\unions{B}$ by \cref{subseteq}.
                                    Thus $X\setminus\unions{B} = \inters{F}$ by \cref{subseteq_antisymmetric}.
                                \end{subproof}
                                Thus $X\setminus\unions{B} = \emptyset$.
                                Thus $X\subseteq\unions{B}$ by \cref{setminus_eq_emptyset_iff_subseteq}.
                                Thus $B$ is a covering of $X$ by \cref{covering}.
                                Contradiction.
                            \end{subproof}
                            Thus $C$ has the finite intersection property by \cref{finite_intersection_property}.
                        \end{subproof}
                        Thus $\inters{C}\neq\emptyset$.
                        Show that $X\setminus\unions{A} = \inters{C}$.
                        \begin{subproof}
                            Show that for all $x$ such that $x\in X\setminus\unions{A}$ we have $x\in\inters{C}$.
                            \begin{subproof}
                                Fix $x$.
                                Assume $x\in X\setminus\unions{A}$.
                                Thus $x\in X$ by \cref{setminus_elim_left}.
                                Thus $x\notin\unions{A}$ by \cref{setminus_elim_right}.
                                Show that for all $D$ such that $D\in C$ we have $x\in D$.
                                \begin{subproof}
                                    Fix $D$.
                                    Assume $D\in C$.
                                    Take $U\in A$ such that $D = X\setminus U$.
                                    Suppose not.
                                    Then $x\notin D$.
                                    Thus $x\in X\setminus D$ by \cref{setminus_intro}.
                                    Thus $x\in U$.
                                    Thus $x\in\unions{A}$ by \cref{unions_iff}.
                                    Contradiction.
                                \end{subproof}
                                Thus $C$ is inhabited by \cref{finite_intersection_property}.
                                Thus $x\in\inters{C}$ by \cref{inters_intro}.
                            \end{subproof}
                            Thus $X\setminus\unions{A}\subseteq\inters{C}$ by \cref{subseteq}.
                            Show that for all $x$ such that $x\in\inters{C}$ we have $x\in X\setminus\unions{A}$.
                            \begin{subproof}
                                Fix $x$.
                                Assume $x\in\inters{C}$.
                                Take $D_0$ such that $D_0\in C$ by \cref{finite_intersection_property}.
                                Thus $x\in D_0$ by \cref{inters_destr}.
                                Thus $D_0\subseteq X$ by \cref{closed_in}.
                                Thus $x\in X$ by \cref{subseteq}.
                                Suppose not.
                                Then $x\in\unions{A}$.
                                Take $U\in A$ such that $x\in U$ by \cref{unions_iff}.
                                Thus $X\setminus U\in C$.
                                Thus $x\in X\setminus U$ by \cref{inters_destr}.
                                Thus $x\notin U$ by \cref{setminus_elim_right}.
                                Contradiction.
                            \end{subproof}
                            Thus $\inters{C}\subseteq X\setminus\unions{A}$ by \cref{subseteq}.
                            Thus $X\setminus\unions{A} = \inters{C}$ by \cref{subseteq_antisymmetric}.
                        \end{subproof}
                        Show that $\unions{A} = X$.
                        \begin{subproof}
                            $X\subseteq\unions{A}$ by \cref{covering}.
                            Show that $\unions{A}\subseteq X$.
                            \begin{subproof}
                                Show that for all $U$ such that $U\in A$ we have $U\subseteq X$.
                                \begin{subproof}
                                Fix $U$.
                                Assume $U\in A$.
                                Thus $U$ is open in $X$ with respect to $T$ by \cref{open_covering}.
                                Thus $U\in T$ by \cref{open_in}.
                                $T\subseteq\pow{X}$ by \cref{topology_on}.
                                Thus $U\in\pow{X}$ by \cref{subseteq}.
                                Thus $U\subseteq X$ by \cref{pow_iff}.
                                \end{subproof}
                                Thus $A\subseteq\pow{X}$ by \cref{powerset_intro,subseteq}.
                                Thus $\unions{A}\subseteq X$ by \cref{unions_subseteq_of_powerset_is_subseteq}.
                            \end{subproof}
                            Thus $\unions{A} = X$ by \cref{subseteq_antisymmetric}.
                        \end{subproof}
                        Thus $X\setminus\unions{A} = \emptyset$ by \cref{setminus_self}.
                        Thus $\inters{C} = \emptyset$ by \cref{subseteq_antisymmetric,subseteq_refl}.
                        Contradiction.
                    \end{subproof}
                \end{byCase}
            \end{subproof}
            Thus $X$ is compact with respect to $T$ by \cref{compact}.
        \end{subproof}
    \end{subproof}
\end{proof}

% from §27 helper definition 1: immediate successor
\begin{definition}\label{immediate_successor}
    $y$ is an immediate successor of $x$ in $X$ with respect to $R$ iff
        $x\in X$ and $y\in X$ and $x\mathrel{R}y$ and
        $\intervalclosedrel{R}{X}{x}{y} = \{x,y\}$.
\end{definition}

% from §27 helper lemma 1: interval of an immediate successor
\begin{lemma}\label{immediate_successor_interval}
    Assume $R$ is asymmetric.
    Suppose $y$ is an immediate successor of $x$ in $X$ with respect to $R$.
    Let $I = \intervalclosedrel{R}{X}{x}{y}$.
    Then $I = \{x,y\}$.
\end{lemma}
\begin{proof}
    Follows by \cref{immediate_successor}.
\end{proof}

% from §27 helper lemma 2: non-immediate successor gives an intermediate point
\begin{lemma}\label{no_immediate_successor_between}
    Suppose $x\in X$.
    Suppose $y\in X$.
    Suppose $x\mathrel{R}y$.
    Suppose it is wrong that $y$ is an immediate successor of $x$ in $X$ with respect to $R$.
    Then there exists $z\in X$ such that $x\mathrel{R}z$ and $z\mathrel{R}y$.
\end{lemma}
\begin{proof}
    Let $I = \intervalclosedrel{R}{X}{x}{y}$.
    Suppose not.
    Then it is wrong that there exists $z\in X$ such that $x\mathrel{R}z$ and $z\mathrel{R}y$.
    Show that $I = \{x,y\}$.
    \begin{subproof}
        Show that $\{x,y\}\subseteq I$.
        \begin{subproof}
            $x\in I$ by \cref{intervalclosed_rel}.
            $y\in I$ by \cref{intervalclosed_rel}.
            Show that for all $z$ such that $z\in\{x,y\}$ we have $z\in I$.
            \begin{subproof}
                Fix $z$.
                Assume $z\in\{x,y\}$.
                Thus $z = x$ or $z = y$ by \cref{upair_elim}.
                \begin{byCase}
                \caseOf{$z = x$.} Thus $z\in I$ by \cref{intervalclosed_rel}.
                \caseOf{$z = y$.} Thus $z\in I$ by \cref{intervalclosed_rel}.
                \end{byCase}
            \end{subproof}
            Thus $\{x,y\}\subseteq I$ by \cref{subseteq}.
        \end{subproof}
        Show that $I\subseteq \{x,y\}$.
        \begin{subproof}
            Show that for all $z$ such that $z\in I$ we have $z\in\{x,y\}$.
            \begin{subproof}
                Fix $z$.
                Assume $z\in I$.
                Thus $z\in X$ and $(x\mathrel{R}z \lor z = x)$ and $(z\mathrel{R}y \lor z = y)$
                    by \cref{intervalclosed_rel}.
                \begin{byCase}
                \caseOf{$z = x$.}
                    Thus $z\in\{x,y\}$ by \cref{upair_intro_left}.
                \caseOf{$z = y$.}
                    Thus $z\in\{x,y\}$ by \cref{upair_intro_right}.
                \caseOf{$x\mathrel{R}z$ and $z\mathrel{R}y$.}
                    Then it is wrong that there exists $z\in X$ such that $x\mathrel{R}z$ and $z\mathrel{R}y$.
                    Contradiction.
                \end{byCase}
            \end{subproof}
            Thus $I\subseteq \{x,y\}$ by \cref{subseteq}.
        \end{subproof}
        Thus $I = \{x,y\}$ by \cref{subseteq_antisymmetric}.
    \end{subproof}
    Thus $y$ is an immediate successor of $x$ in $X$ with respect to $R$ by \cref{immediate_successor}.
    Contradiction.
\end{proof}

% from §27 Item 1: closed interval is compact in an ordered set with the least upper bound property
\begin{theorem}\label{ordered_closed_interval_compact}
    Assume $R$ is a simple order relation on $X$.
    Assume $X$ has the least upper bound property with respect to $R$.
    Assume $T$ is the order topology on $X$ with respect to $R$.
    Suppose $a\in X$.
    Suppose $b\in X$.
    Suppose $a\mathrel{R}b$ or $a = b$.
    Let $Y = \intervalclosedrel{R}{X}{a}{b}$.
    Let $T_Y = \subspaceTopology{T}{Y}$.
    Then $Y$ is compact with respect to $T_Y$.
\end{theorem}
\begin{proof}
    $R$ is transitive by \cref{simple_order_relation}.
    $R$ is asymmetric by \cref{simple_order_relation}.
    $R$ is connex on $X$ by \cref{simple_order_relation}.
    Show that for all $A$ such that $A$ is an open covering of $Y$ with respect to $T_Y$
        there exists $B$ such that $B\subseteq A$ and $B$ is finite and $B$ is a covering of $Y$.
    \begin{subproof}
        Fix $A$.
        Assume $A$ is an open covering of $Y$ with respect to $T_Y$.
        \begin{byCase}
        \caseOf{$a = b$.}
            Show that there exists $B$ such that $B\subseteq A$ and $B$ is finite and $B$ is a covering of $Y$.
            \begin{subproof}
                $Y\subseteq\unions{A}$ by \cref{open_covering,covering}.
                $a\in Y$ by \cref{intervalclosed_rel}.
                Thus $a\in\unions{A}$ by \cref{subseteq}.
                Take $U\in A$ such that $a\in U$ by \cref{unions_iff}.
                Let $B = \{U\}$.
                $B\subseteq A$ by \cref{singleton_subset_intro}.
                Show that $B$ is finite.
                \begin{subproof}
                    Let $G = \{U,U\}$.
                    $G$ is finite by \cref{pair_is_finite}.
                    $U\in G$ by \cref{upair_intro_left}.
                    Thus $B\subseteq G$ by \cref{singleton_subset_intro}.
                    Thus $B$ is finite by \cref{subseteq_finite_is_finite}.
                \end{subproof}
                Show that $B$ is a covering of $Y$.
                \begin{subproof}
                    Show that for all $x$ such that $x\in Y$ we have $x\in\unions{B}$.
                    \begin{subproof}
                        Fix $x$.
                        Assume $x\in Y$.
                        Thus $x\in X$ and $(a\mathrel{R}x \lor x = a)$ and $(x\mathrel{R}a \lor x = a)$
                            by \cref{intervalclosed_rel}.
                        Show that $x = a$.
                        \begin{subproof}
                            Suppose not.
                            Then $a\mathrel{R}x$ and $x\mathrel{R}a$ by \cref{intervalclosed_rel}.
                            Contradiction by \cref{asymmetric}.
                        \end{subproof}
                        Thus $x\in U$.
                        Thus $U\in B$ by \cref{singleton_intro}.
                        Thus $x\in\unions{B}$ by \cref{unions_iff}.
                    \end{subproof}
                    Thus $Y\subseteq\unions{B}$ by \cref{subseteq}.
                    Thus $B$ is a covering of $Y$ by \cref{covering}.
                \end{subproof}
            \end{subproof}
        \caseOf{$a\mathrel{R}b$.}
            $a\in Y$ and $b\in Y$ by \cref{intervalclosed_rel}.
            Thus $a\neq b$ by \cref{asymmetric}.
            Thus $X$ has more than one element by \cref{more_than_one_element}.
            Thus $Y$ has more than one element by \cref{more_than_one_element}.
            $Y$ is convex in $X$ with respect to $R$ by \cref{intervalclosedrel_convex}.
            Show that $Y\subseteq X$.
            \begin{subproof}
                Show that for all $y$ such that $y\in Y$ we have $y\in X$.
                \begin{subproof}
                    Fix $y$.
                    Assume $y\in Y$.
                    Thus $y\in X$ by \cref{intervalclosed_rel}.
                \end{subproof}
                Thus $Y\subseteq X$ by \cref{subseteq}.
            \end{subproof}
            Show that $a$ is a smallest element of $Y$ with respect to $R$.
            \begin{subproof}
                $a\in Y$ by \cref{intervalclosed_rel}.
                Show that for all $x\in Y$ we have $a\mathrel{R}x$ or $a = x$.
                \begin{subproof}
                    Fix $x$.
                    Assume $x\in Y$.
                    Thus $a\mathrel{R}x$ or $x = a$ by \cref{intervalclosed_rel}.
                \end{subproof}
                Thus $a$ is a smallest element of $Y$ with respect to $R$ by \cref{smallest_element}.
            \end{subproof}
            Show that $b$ is a largest element of $Y$ with respect to $R$.
            \begin{subproof}
                $b\in Y$ by \cref{intervalclosed_rel}.
                Show that for all $x\in Y$ we have $x\mathrel{R}b$ or $x = b$.
                \begin{subproof}
                    Fix $x$.
                    Assume $x\in Y$.
                    Thus $x\mathrel{R}b$ or $x = b$ by \cref{intervalclosed_rel}.
                \end{subproof}
                Thus $b$ is a largest element of $Y$ with respect to $R$ by \cref{largest_element}.
            \end{subproof}
            $T = \basisTopology{\orderBasis{R}{X}}{X}$ by \cref{order_topology}.
            $\orderBasis{R}{X}$ is a basis on $X$ by \cref{order_basis_is_basis}.
            Thus $T$ is a topology on $X$ by \cref{basis_topology_is_topology}.
            Thus $T_Y$ is a topology on $Y$ by \cref{subspace_topology_is_topology,subspace_of}.

            Show that for all $x$ such that $x\in Y$ and $x\neq b$ there exists $y\in Y$ such that
                $x\mathrel{R}y$ and
                there exists $B$ such that $B\subseteq A$ and $B$ is finite and
                    $B$ is a covering of $\intervalclosedrel{R}{X}{x}{y}$.
            \begin{subproof}
                Fix $x$.
                Assume $x\in Y$ and $x\neq b$.
                $Y\subseteq\unions{A}$ by \cref{open_covering,covering}.
                Thus $x\in\unions{A}$ by \cref{subseteq}.
                Take $U\in A$ such that $x\in U$ by \cref{unions_iff}.
                $U$ is open in $Y$ with respect to $T_Y$ by \cref{open_covering}.
                Thus $U\in T_Y$ by \cref{open_in}.
                Take $V\in T$ such that $U = Y\inter V$ by \cref{subspace_topology}.
                Thus $x\in V$ by \cref{inter_elim_right}.
                Take $B_1\in\orderBasis{R}{X}$ such that $x\in B_1$ and $B_1\subseteq V$
                    by \cref{topology_generated_by_basis}.
                Show that there exists $c\in Y$ such that $x\mathrel{R}c$ and
                    $\intervalclosedopenrel{R}{X}{x}{c} \subseteq V$.
                \begin{subproof}
                    $B_1\in\pow{X}$ and
                        $((\exists u, v\in X. u\mathrel{R}v \land B_1 = \intervalopenrel{R}{X}{u}{v}) \lor
                        (\exists v\in X. \exists a_0\in X. (\forall z\in X. a_0\mathrel{R}z \lor a_0 = z) \land
                            B_1 = \intervalclosedopenrel{R}{X}{a_0}{v}) \lor
                        (\exists u\in X. \exists b_0\in X. (\forall z\in X. z\mathrel{R}b_0 \lor z = b_0) \land
                            B_1 = \intervalopenclosedrel{R}{X}{u}{b_0}))$
                        by \cref{order_basis}.
                    \begin{byCase}
                    \caseOf{$\exists u, v\in X. u\mathrel{R}v \land B_1 = \intervalopenrel{R}{X}{u}{v}$.}
                        Take $u, v\in X$ such that $u\mathrel{R}v$ and $B_1 = \intervalopenrel{R}{X}{u}{v}$.
                        Thus $u\mathrel{R}x$ and $x\mathrel{R}v$ by \cref{intervalopen_rel}.
                        Show that there exists $c\in Y$ such that $x\mathrel{R}c$ and
                            $\intervalclosedopenrel{R}{X}{x}{c} \subseteq B_1$.
                        \begin{subproof}
                            \begin{byCase}
                            \caseOf{$v\in Y$.}
                                Let $c = v$.
                                Show that $\intervalclosedopenrel{R}{X}{x}{c} \subseteq B_1$.
                                \begin{subproof}
                                    Show that for all $z$ such that $z\in\intervalclosedopenrel{R}{X}{x}{c}$ we have $z\in B_1$.
                                    \begin{subproof}
                                        Fix $z$.
                                        Assume $z\in\intervalclosedopenrel{R}{X}{x}{c}$.
                                        Thus $z\in X$ and $(x\mathrel{R}z \lor z = x)$ and $z\mathrel{R}c$
                                            by \cref{intervalclosedopen_rel}.
                                        Show that $u\mathrel{R}z$.
                                        \begin{subproof}
                                            \begin{byCase}
                                            \caseOf{$z = x$.}
                                                Thus $u\mathrel{R}z$.
                                            \caseOf{$x\mathrel{R}z$.}
                                                Thus $u\mathrel{R}z$ by \hyperref[transitive]{transitivity}.
                                            \end{byCase}
                                        \end{subproof}
                                        Thus $u\mathrel{R}z$ and $z\mathrel{R}v$.
                                        Thus $z\in\intervalopenrel{R}{X}{u}{v}$ by \cref{intervalopen_rel}.
                                        Thus $z\in B_1$.
                                    \end{subproof}
                                    Thus $\intervalclosedopenrel{R}{X}{x}{c} \subseteq B_1$ by \cref{subseteq}.
                                \end{subproof}
                            \caseOf{it is wrong that $v\in Y$.}
                                Thus $a\mathrel{R}x$ or $x = a$ by \cref{intervalclosed_rel}.
                                Show that $a\mathrel{R}v$.
                                \begin{subproof}
                                    \begin{byCase}
                                    \caseOf{$x = a$.}
                                        Thus $a\mathrel{R}v$.
                                    \caseOf{$a\mathrel{R}x$.}
                                        Thus $a\mathrel{R}v$ by \hyperref[transitive]{transitivity}.
                                    \end{byCase}
                                \end{subproof}
                                Show that $b\mathrel{R}v$.
                                \begin{subproof}
                                    Suppose not.
                                    Then $v\mathrel{R}b$ or $v = b$ by \cref{connex}.
                                    \begin{byCase}
                                    \caseOf{$v = b$.}
                                        Thus $v\in Y$.
                                        Contradiction.
                                    \caseOf{$v\mathrel{R}b$.}
                                        Thus $v\in Y$ by \cref{intervalclosed_rel}.
                                        Contradiction.
                                    \end{byCase}
                                \end{subproof}
                                Show that $b\in B_1$.
                                \begin{subproof}
                                    Show that $u\mathrel{R}b$.
                                    \begin{subproof}
                                        Thus $x\mathrel{R}b$ by \cref{intervalclosed_rel}.
                                        Thus $u\mathrel{R}b$ by \hyperref[transitive]{transitivity}.
                                    \end{subproof}
                                    Thus $u\mathrel{R}b$ and $b\mathrel{R}v$.
                                    Thus $b\in\intervalopenrel{R}{X}{u}{v}$ by \cref{intervalopen_rel}.
                                    Thus $b\in B_1$.
                                \end{subproof}
                                Show that $\intervalclosedopenrel{R}{X}{x}{b} \subseteq B_1$.
                                \begin{subproof}
                                    Show that for all $z$ such that $z\in\intervalclosedopenrel{R}{X}{x}{b}$ we have $z\in B_1$.
                                    \begin{subproof}
                                        Fix $z$.
                                        Assume $z\in\intervalclosedopenrel{R}{X}{x}{b}$.
                                        Thus $z\in X$ and $(x\mathrel{R}z \lor z = x)$ and $z\mathrel{R}b$
                                            by \cref{intervalclosedopen_rel}.
                                        Show that $u\mathrel{R}z$.
                                        \begin{subproof}
                                            \begin{byCase}
                                            \caseOf{$z = x$.}
                                                Thus $u\mathrel{R}z$.
                                            \caseOf{$x\mathrel{R}z$.}
                                                Thus $u\mathrel{R}z$ by \hyperref[transitive]{transitivity}.
                                            \end{byCase}
                                        \end{subproof}
                                        Thus $u\mathrel{R}z$.
                                        Thus $z\mathrel{R}v$ by \hyperref[transitive]{transitivity}.
                                        Thus $z\in\intervalopenrel{R}{X}{u}{v}$ by \cref{intervalopen_rel}.
                                        Thus $z\in B_1$.
                                    \end{subproof}
                                    Thus $\intervalclosedopenrel{R}{X}{x}{b} \subseteq B_1$ by \cref{subseteq}.
                                \end{subproof}
                                Thus $x\mathrel{R}b$ by \cref{intervalclosed_rel}.
                                Thus there exists $c\in Y$ such that $x\mathrel{R}c$ and
                                    $\intervalclosedopenrel{R}{X}{x}{c} \subseteq B_1$.
                            \end{byCase}
                        \end{subproof}
                    \caseOf{$\exists v\in X. \exists a_0\in X. (\forall z\in X. a_0\mathrel{R}z \lor a_0 = z) \land
                        B_1 = \intervalclosedopenrel{R}{X}{a_0}{v}$.}
                        Take $a_0, v\in X$ such that
                            $(\forall z\in X. a_0\mathrel{R}z \lor a_0 = z)$ and
                            $B_1 = \intervalclosedopenrel{R}{X}{a_0}{v}$.
                        Thus $x\in X$ and $(a_0\mathrel{R}x \lor x = a_0)$ and $x\mathrel{R}v$
                            by \cref{intervalclosedopen_rel}.
                        Show that there exists $c\in Y$ such that $x\mathrel{R}c$ and
                            $\intervalclosedopenrel{R}{X}{x}{c} \subseteq B_1$.
                        \begin{subproof}
                            \begin{byCase}
                            \caseOf{$v\in Y$.}
                                Let $c = v$.
                                Show that $\intervalclosedopenrel{R}{X}{x}{c} \subseteq B_1$.
                                \begin{subproof}
                                    Show that for all $z$ such that $z\in\intervalclosedopenrel{R}{X}{x}{c}$ we have $z\in B_1$.
                                    \begin{subproof}
                                        Fix $z$.
                                        Assume $z\in\intervalclosedopenrel{R}{X}{x}{c}$.
                                        Thus $z\in X$ and $(x\mathrel{R}z \lor z = x)$ and $z\mathrel{R}c$
                                            by \cref{intervalclosedopen_rel}.
                                        Thus $a_0\mathrel{R}z$ or $z = a_0$ by \cref{smallest_element}.
                                        Thus $z\mathrel{R}b$.
                                        Thus $z\mathrel{R}v$ by \hyperref[transitive]{transitivity}.
                                        Thus $z\in\intervalclosedopenrel{R}{X}{a_0}{v}$ by \cref{intervalclosedopen_rel}.
                                        Thus $z\in B_1$.
                                    \end{subproof}
                                    Thus $\intervalclosedopenrel{R}{X}{x}{c} \subseteq B_1$ by \cref{subseteq}.
                                \end{subproof}
                            \caseOf{it is wrong that $v\in Y$.}
                                Thus $a\mathrel{R}x$ or $x = a$ by \cref{intervalclosed_rel}.
                                Show that $a\mathrel{R}v$.
                                \begin{subproof}
                                    \begin{byCase}
                                    \caseOf{$x = a$.}
                                        Thus $a\mathrel{R}v$.
                                    \caseOf{$a\mathrel{R}x$.}
                                        Thus $a\mathrel{R}v$ by \hyperref[transitive]{transitivity}.
                                    \end{byCase}
                                \end{subproof}
                                Show that $b\mathrel{R}v$.
                                \begin{subproof}
                                    Suppose not.
                                    Then $v\mathrel{R}b$ or $v = b$ by \cref{connex}.
                                    \begin{byCase}
                                    \caseOf{$v = b$.}
                                        Thus $v\in Y$.
                                        Contradiction.
                                    \caseOf{$v\mathrel{R}b$.}
                                        Thus $v\in Y$ by \cref{intervalclosed_rel}.
                                        Contradiction.
                                    \end{byCase}
                                \end{subproof}
                                Show that $b\in B_1$.
                                \begin{subproof}
                                    Thus $a_0\mathrel{R}b$ or $b = a_0$ by \cref{smallest_element}.
                                    Thus $b\in\intervalclosedopenrel{R}{X}{a_0}{v}$ by \cref{intervalclosedopen_rel}.
                                    Thus $b\in B_1$.
                                \end{subproof}
                                Show that $\intervalclosedopenrel{R}{X}{x}{b} \subseteq B_1$.
                                \begin{subproof}
                                    Show that for all $z$ such that $z\in\intervalclosedopenrel{R}{X}{x}{b}$ we have $z\in B_1$.
                                    \begin{subproof}
                                        Fix $z$.
                                        Assume $z\in\intervalclosedopenrel{R}{X}{x}{b}$.
                                        Thus $z\in X$ and $(x\mathrel{R}z \lor z = x)$ and $z\mathrel{R}b$
                                            by \cref{intervalclosedopen_rel}.
                                        Thus $a_0\mathrel{R}z$ or $z = a_0$ by \cref{smallest_element}.
                                        Thus $z\mathrel{R}v$ by \hyperref[transitive]{transitivity}.
                                        Thus $z\in\intervalclosedopenrel{R}{X}{a_0}{v}$ by \cref{intervalclosedopen_rel}.
                                        Thus $z\in B_1$.
                                    \end{subproof}
                                    Thus $\intervalclosedopenrel{R}{X}{x}{b} \subseteq B_1$ by \cref{subseteq}.
                                \end{subproof}
                                Thus $x\mathrel{R}b$ by \cref{intervalclosed_rel}.
                                Thus there exists $c\in Y$ such that $x\mathrel{R}c$ and
                                    $\intervalclosedopenrel{R}{X}{x}{c} \subseteq B_1$.
                            \end{byCase}
                        \end{subproof}
                    \caseOf{$\exists u\in X. \exists b_0\in X. (\forall z\in X. z\mathrel{R}b_0 \lor z = b_0) \land
                        B_1 = \intervalopenclosedrel{R}{X}{u}{b_0}$.}
                        Take $u, b_0\in X$ such that
                            $(\forall z\in X. z\mathrel{R}b_0 \lor z = b_0)$ and
                            $B_1 = \intervalopenclosedrel{R}{X}{u}{b_0}$.
                        Thus $u\mathrel{R}x$ and $(x\mathrel{R}b_0 \lor x = b_0)$ by \cref{intervalopenclosed_rel}.
                        Show that $\intervalclosedopenrel{R}{X}{x}{b} \subseteq B_1$.
                        \begin{subproof}
                            Show that for all $z$ such that $z\in\intervalclosedopenrel{R}{X}{x}{b}$ we have $z\in B_1$.
                            \begin{subproof}
                                Fix $z$.
                                Assume $z\in\intervalclosedopenrel{R}{X}{x}{b}$.
                                Thus $z\in X$ and $(x\mathrel{R}z \lor z = x)$ and $z\mathrel{R}b$
                                    by \cref{intervalclosedopen_rel}.
                                Show that $u\mathrel{R}z$.
                                \begin{subproof}
                                    \begin{byCase}
                                    \caseOf{$z = x$.}
                                        Thus $u\mathrel{R}z$.
                                    \caseOf{$x\mathrel{R}z$.}
                                        Thus $u\mathrel{R}z$ by \hyperref[transitive]{transitivity}.
                                    \end{byCase}
                                \end{subproof}
                                Thus $u\mathrel{R}z$.
                                Thus $z\mathrel{R}b_0$ by \hyperref[transitive]{transitivity}.
                                Thus $z\in\intervalopenclosedrel{R}{X}{u}{b_0}$ by \cref{intervalopenclosed_rel}.
                                Thus $z\in B_1$.
                            \end{subproof}
                            Thus $\intervalclosedopenrel{R}{X}{x}{b} \subseteq B_1$ by \cref{subseteq}.
                        \end{subproof}
                        Thus $x\mathrel{R}b$ by \cref{intervalclosed_rel}.
                        Thus there exists $c\in Y$ such that $x\mathrel{R}c$ and
                            $\intervalclosedopenrel{R}{X}{x}{c} \subseteq B_1$.
                    \end{byCase}
                \end{subproof}
                Take $c\in Y$ such that $x\mathrel{R}c$ and $\intervalclosedopenrel{R}{X}{x}{c} \subseteq V$.
                \begin{byCase}
                \caseOf{there exists $y\in X$ such that $y$ is an immediate successor of $x$ in $X$ with respect to $R$.}
                    Take $y\in X$ such that $y$ is an immediate successor of $x$ in $X$ with respect to $R$.
                    Show that $y\in Y$.
                    \begin{subproof}
                        Thus $x\mathrel{R}b$ by \cref{intervalclosed_rel}.
                        Show that $y\mathrel{R}b$ or $y = b$.
                        \begin{subproof}
                            Thus $y\mathrel{R}b$ or $y = b$ or $b\mathrel{R}y$ by \cref{connex}.
                            \begin{byCase}
                            \caseOf{$b\mathrel{R}y$.}
                                Then $b\in\intervalclosedrel{R}{X}{x}{y}$ by \cref{intervalclosed_rel}.
                                Let $I = \intervalclosedrel{R}{X}{x}{y}$.
                                Thus $I = \{x,y\}$ by \cref{immediate_successor_interval}.
                                Thus $b\in\{x,y\}$.
                                Thus $b = x$ or $b = y$ by \cref{upair_elim}.
                                Thus $y = b$.
                            \caseOf{$y\mathrel{R}b$ or $y = b$.}
                            \end{byCase}
                        \end{subproof}
                        Show that $a\mathrel{R}y$.
                        \begin{subproof}
                            Thus $x\mathrel{R}y$ by \cref{immediate_successor}.
                            Thus $a\mathrel{R}x$ or $x = a$ by \cref{intervalclosed_rel}.
                            \begin{byCase}
                            \caseOf{$x = a$.}
                                Thus $a\mathrel{R}y$.
                            \caseOf{$a\mathrel{R}x$.}
                                Thus $a\mathrel{R}y$ by \hyperref[transitive]{transitivity}.
                            \end{byCase}
                        \end{subproof}
                        Thus $y\mathrel{R}b$ or $y = b$.
                        Thus $y\in Y$ by \cref{intervalclosed_rel}.
                    \end{subproof}
                    $Y\subseteq\unions{A}$ by \cref{open_covering,covering}.
                    Thus $y\in\unions{A}$ by \cref{subseteq}.
                    Take $U_1\in A$ such that $y\in U_1$ by \cref{unions_iff}.
                    Let $B = \{U,U_1\}$.
                    Show that $B\subseteq A$.
                    \begin{subproof}
                        Show that for all $V$ such that $V\in B$ we have $V\in A$.
                        \begin{subproof}
                            Fix $V$.
                            Assume $V\in B$.
                            Thus $V = U$ or $V = U_1$ by \cref{upair_elim}.
                            \begin{byCase}
                            \caseOf{$V = U$.} Thus $V\in A$.
                            \caseOf{$V = U_1$.} Thus $V\in A$.
                            \end{byCase}
                        \end{subproof}
                        Thus $B\subseteq A$ by \cref{subseteq}.
                    \end{subproof}
                    $B$ is finite by \cref{pair_is_finite}.
                    Show that $B$ is a covering of $\intervalclosedrel{R}{X}{x}{y}$.
                    \begin{subproof}
                        Show that for all $z$ such that $z\in\intervalclosedrel{R}{X}{x}{y}$ we have $z\in\unions{B}$.
                        \begin{subproof}
                            Fix $z$.
                            Assume $z\in\intervalclosedrel{R}{X}{x}{y}$.
                            Thus $z\in\{x,y\}$ by \cref{immediate_successor_interval}.
                            Thus $z = x$ or $z = y$ by \cref{upair_elim}.
                            \begin{byCase}
                            \caseOf{$z = x$.}
                                Thus $z\in U$.
                                Thus $U\in B$ by \cref{upair_intro_left}.
                                Thus $z\in\unions{B}$ by \cref{unions_iff}.
                            \caseOf{$z = y$.}
                                Thus $z\in U_1$.
                                Thus $U_1\in B$ by \cref{upair_intro_right}.
                                Thus $z\in\unions{B}$ by \cref{unions_iff}.
                            \end{byCase}
                        \end{subproof}
                        Thus $\intervalclosedrel{R}{X}{x}{y}\subseteq\unions{B}$ by \cref{subseteq}.
                        Thus $B$ is a covering of $\intervalclosedrel{R}{X}{x}{y}$ by \cref{covering}.
                    \end{subproof}
                \caseOf{it is wrong that there exists $y\in X$ such that $y$ is an immediate successor of $x$ in $X$ with respect to $R$.}
                    Thus $x\in X$ by \cref{subseteq}.
                    Thus $c\in X$ by \cref{subseteq}.
                    Thus it is wrong that $c$ is an immediate successor of $x$ in $X$ with respect to $R$.
                    Take $y\in X$ such that $x\mathrel{R}y$ and $y\mathrel{R}c$ by \cref{no_immediate_successor_between}.
                    Thus $x\in Y$ and $c\in Y$.
                    Thus $y\in Y$ by \cref{convex_in}.
                    Let $B = \{U\}$.
                    $B\subseteq A$ by \cref{singleton_subset_intro}.
                    Show that $B$ is finite.
                    \begin{subproof}
                        Let $G = \{U,U\}$.
                        $G$ is finite by \cref{pair_is_finite}.
                        $U\in G$ by \cref{upair_intro_left}.
                        Thus $B\subseteq G$ by \cref{singleton_subset_intro}.
                        Thus $B$ is finite by \cref{subseteq_finite_is_finite}.
                    \end{subproof}
                    Show that $B$ is a covering of $\intervalclosedrel{R}{X}{x}{y}$.
                    \begin{subproof}
                        Show that for all $z$ such that $z\in\intervalclosedrel{R}{X}{x}{y}$ we have $z\in\unions{B}$.
                        \begin{subproof}
                            Fix $z$.
                            Assume $z\in\intervalclosedrel{R}{X}{x}{y}$.
                            Thus $z\in X$ and $(x\mathrel{R}z \lor z = x)$ and $(z\mathrel{R}y \lor z = y)$
                                by \cref{intervalclosed_rel}.
                            Show that $z\mathrel{R}c$.
                            \begin{subproof}
                                \begin{byCase}
                                \caseOf{$z = y$.}
                                    Thus $z\mathrel{R}c$.
                                \caseOf{$z\mathrel{R}y$.}
                                    Thus $z\mathrel{R}c$ by \hyperref[transitive]{transitivity}.
                                \end{byCase}
                            \end{subproof}
                            Thus $z\in X$ and $(x\mathrel{R}z \lor z = x)$ and $(z\mathrel{R}y \lor z = y)$
                                by \cref{intervalclosed_rel}.
                            Show that $z\mathrel{R}c$.
                            \begin{subproof}
                                \begin{byCase}
                                \caseOf{$z = y$.}
                                    Thus $z\mathrel{R}c$.
                                \caseOf{$z\mathrel{R}y$.}
                                    Thus $z\mathrel{R}c$ by \hyperref[transitive]{transitivity}.
                                \end{byCase}
                            \end{subproof}
                            Thus $z\in\intervalclosedopenrel{R}{X}{x}{c}$ by \cref{intervalclosedopen_rel}.
                            Thus $z\in V$ by \cref{subseteq}.
                            Show that $z\in Y$.
                            \begin{subproof}
                                \begin{byCase}
                                \caseOf{$z = x$.}
                                    Thus $z\in Y$.
                                \caseOf{$x\mathrel{R}z$.}
                                    Thus $x\in Y$ and $c\in Y$.
                                    Thus $z\in Y$ by \cref{convex_in}.
                                \end{byCase}
                            \end{subproof}
                            Thus $z\in Y\inter V$ by \cref{inter_intro}.
                            Thus $z\in U$.
                            Thus $U\in B$ by \cref{singleton_intro}.
                            Thus $z\in\unions{B}$ by \cref{unions_iff}.
                        \end{subproof}
                        Thus $\intervalclosedrel{R}{X}{x}{y}\subseteq\unions{B}$ by \cref{subseteq}.
                        Thus $B$ is a covering of $\intervalclosedrel{R}{X}{x}{y}$ by \cref{covering}.
                    \end{subproof}
                \end{byCase}
            \end{subproof}

            Let $C = \{y\in Y \mid \text{$a\mathrel{R}y$ and there exists $B\subseteq A$ such that
                $B$ is finite and $B$ is a covering of $\intervalclosedrel{R}{X}{a}{y}$}\}$.
            Show that $C$ is inhabited.
            \begin{subproof}
                Thus $a\in Y$.
                Take $y\in Y$ such that $a\mathrel{R}y$ and
                    there exists $B$ such that $B\subseteq A$ and $B$ is finite and
                        $B$ is a covering of $\intervalclosedrel{R}{X}{a}{y}$.
                Thus $y\in C$.
                Thus $C$ is inhabited.
            \end{subproof}
            Show that $C\subseteq X$.
            \begin{subproof}
                Show that for all $y$ such that $y\in C$ we have $y\in X$.
                \begin{subproof}
                    Fix $y$.
                    Assume $y\in C$.
                    Thus $y\in Y$.
                    Thus $y\in X$ by \cref{subseteq}.
                \end{subproof}
                Thus $C\subseteq X$ by \cref{subseteq}.
            \end{subproof}
            Show that $b$ is an upper bound of $C$ in $X$ with respect to $R$.
            \begin{subproof}
                $b\in X$.
                Show that for all $y\in C$ we have $y\mathrel{R}b$ or $y = b$.
                \begin{subproof}
                    Fix $y$.
                    Assume $y\in C$.
                    Thus $y\in Y$.
                    Thus $y\mathrel{R}b$ or $y = b$ by \cref{intervalclosed_rel}.
                \end{subproof}
                Thus $b$ is an upper bound of $C$ in $X$ with respect to $R$ by \cref{upper_bound_rel}.
            \end{subproof}
            Take $c$ such that $c$ is a least upper bound of $C$ in $X$ with respect to $R$
                by \cref{least_upper_bound_property}.
            Show that $a\mathrel{R}c$.
            \begin{subproof}
                Take $y\in C$ such that $a\mathrel{R}y$ by \cref{nonempty_inhabited}.
                Suppose not.
                \begin{byCase}
                \caseOf{$c = a$.}
                    $c$ is an upper bound of $C$ in $X$ with respect to $R$ by \cref{least_upper_bound_rel}.
                    Thus $y\mathrel{R}c$ or $y = c$ by \cref{upper_bound_rel}.
                    Show that $y\mathrel{R}c$.
                    \begin{subproof}
                        Suppose not.
                        Then $y = c$.
                        Thus $y = a$.
                        Thus $a\mathrel{R}a$.
                        Contradiction by \cref{asymmetric}.
                    \end{subproof}
                    Thus $y\mathrel{R}a$.
                    Contradiction by \cref{asymmetric}.
                \caseOf{$c\neq a$.}
                    $c$ is an upper bound of $C$ in $X$ with respect to $R$ by \cref{least_upper_bound_rel}.
                    Thus $c\in X$ by \cref{upper_bound_rel}.
                    Thus $a\in X$.
                    Thus $c\mathrel{R}a$ or $a\mathrel{R}c$ by \cref{connex}.
                    Show that $c\mathrel{R}a$.
                    \begin{subproof}
                        Suppose not.
                        Then $a\mathrel{R}c$.
                        Contradiction.
                    \end{subproof}
                    Thus $c\mathrel{R}y$ by \hyperref[transitive]{transitivity}.
                    $c$ is an upper bound of $C$ in $X$ with respect to $R$ by \cref{least_upper_bound_rel}.
                    Thus $y\mathrel{R}c$ or $y = c$ by \cref{upper_bound_rel}.
                    Contradiction by \cref{asymmetric}.
                \end{byCase}
            \end{subproof}
            Thus $c\mathrel{R}b$ or $c = b$ by \cref{least_upper_bound_rel}.
            $c$ is an upper bound of $C$ in $X$ with respect to $R$ by \cref{least_upper_bound_rel}.
            Thus $c\in X$ by \cref{upper_bound_rel}.
            Thus $c\in Y$ by \cref{intervalclosed_rel}.

            Show that $c\in C$.
            \begin{subproof}
                $Y\subseteq\unions{A}$ by \cref{open_covering,covering}.
                Thus $c\in\unions{A}$ by \cref{subseteq}.
                Take $U\in A$ such that $c\in U$ by \cref{unions_iff}.
                $U$ is open in $Y$ with respect to $T_Y$ by \cref{open_covering}.
                Thus $U\in T_Y$ by \cref{open_in}.
                Take $V\in T$ such that $U = Y\inter V$ by \cref{subspace_topology}.
                Thus $c\in Y\inter V$.
                Thus $c\in V$ by \cref{inter_elim_right}.
                Take $B_1\in\orderBasis{R}{X}$ such that $c\in B_1$ and $B_1\subseteq V$
                    by \cref{topology_generated_by_basis}.
                Show that there exists $d\in Y$ such that $d\mathrel{R}c$ and
                    $\intervalopenclosedrel{R}{X}{d}{c} \subseteq V$.
                \begin{subproof}
                    $B_1\in\pow{X}$ and
                        $((\exists u, v\in X. u\mathrel{R}v \land B_1 = \intervalopenrel{R}{X}{u}{v}) \lor
                        (\exists v\in X. \exists a_0\in X. (\forall z\in X. a_0\mathrel{R}z \lor a_0 = z) \land
                            B_1 = \intervalclosedopenrel{R}{X}{a_0}{v}) \lor
                        (\exists u\in X. \exists b_0\in X. (\forall z\in X. z\mathrel{R}b_0 \lor z = b_0) \land
                            B_1 = \intervalopenclosedrel{R}{X}{u}{b_0}))$
                        by \cref{order_basis}.
                    \begin{byCase}
                    \caseOf{$\exists u, v\in X. u\mathrel{R}v \land B_1 = \intervalopenrel{R}{X}{u}{v}$.}
                        Take $u, v\in X$ such that $u\mathrel{R}v$ and $B_1 = \intervalopenrel{R}{X}{u}{v}$.
                        Thus $u\mathrel{R}c$ and $c\mathrel{R}v$ by \cref{intervalopen_rel}.
                        Show that there exists $d\in Y$ such that $d\mathrel{R}c$ and
                            $\intervalopenclosedrel{R}{X}{d}{c} \subseteq B_1$.
                        \begin{subproof}
                            \begin{byCase}
                            \caseOf{$u\in Y$.}
                                Let $d = u$.
                                Show that $\intervalopenclosedrel{R}{X}{d}{c} \subseteq B_1$.
                                \begin{subproof}
                                    Show that for all $z$ such that $z\in\intervalopenclosedrel{R}{X}{d}{c}$ we have $z\in B_1$.
                                    \begin{subproof}
                                        Fix $z$.
                                        Assume $z\in\intervalopenclosedrel{R}{X}{d}{c}$.
                                        Thus $z\in X$ and $d\mathrel{R}z$ and $(z\mathrel{R}c \lor z = c)$
                                            by \cref{intervalopenclosed_rel}.
                                        Show that $z\mathrel{R}v$.
                                        \begin{subproof}
                                            \begin{byCase}
                                            \caseOf{$z = c$.}
                                                Thus $z\mathrel{R}v$.
                                            \caseOf{$z\mathrel{R}c$.}
                                                Thus $z\mathrel{R}v$ by \hyperref[transitive]{transitivity}.
                                            \end{byCase}
                                        \end{subproof}
                                        Thus $u\mathrel{R}z$ and $z\mathrel{R}v$.
                                        Thus $z\in\intervalopenrel{R}{X}{u}{v}$ by \cref{intervalopen_rel}.
                                        Thus $z\in B_1$.
                                    \end{subproof}
                                    Thus $\intervalopenclosedrel{R}{X}{d}{c} \subseteq B_1$ by \cref{subseteq}.
                                \end{subproof}
                            \caseOf{it is wrong that $u\in Y$.}
                                Show that $u\mathrel{R}a$.
                                \begin{subproof}
                                    Suppose not.
                                    Then $a\mathrel{R}u$ or $a = u$ by \cref{connex}.
                                    Show that $u\mathrel{R}b$ or $u = b$.
                                    \begin{subproof}
                                        Thus $c\mathrel{R}b$ or $c = b$ by \cref{least_upper_bound_rel}.
                                        \begin{byCase}
                                        \caseOf{$c = b$.}
                                            Thus $u\mathrel{R}b$.
                                        \caseOf{$c\mathrel{R}b$.}
                                            Thus $u\mathrel{R}b$ by \hyperref[transitive]{transitivity}.
                                        \end{byCase}
                                    \end{subproof}
                                    Thus $u\in Y$ by \cref{intervalclosed_rel}.
                                    Contradiction.
                                \end{subproof}
                                Let $d = a$.
                                Show that $\intervalopenclosedrel{R}{X}{d}{c} \subseteq B_1$.
                                \begin{subproof}
                                    Show that for all $z$ such that $z\in\intervalopenclosedrel{R}{X}{d}{c}$ we have $z\in B_1$.
                                    \begin{subproof}
                                        Fix $z$.
                                        Assume $z\in\intervalopenclosedrel{R}{X}{d}{c}$.
                                        Thus $z\in X$ and $d\mathrel{R}z$ and $(z\mathrel{R}c \lor z = c)$
                                            by \cref{intervalopenclosed_rel}.
                                        Show that $u\mathrel{R}z$.
                                        \begin{subproof}
                                            Thus $u\mathrel{R}a$.
                                            Thus $u\mathrel{R}z$ by \hyperref[transitive]{transitivity}.
                                        \end{subproof}
                                        Show that $z\mathrel{R}v$.
                                        \begin{subproof}
                                            \begin{byCase}
                                            \caseOf{$z = c$.}
                                                Thus $z\mathrel{R}v$.
                                            \caseOf{$z\mathrel{R}c$.}
                                                Thus $z\mathrel{R}v$ by \hyperref[transitive]{transitivity}.
                                            \end{byCase}
                                        \end{subproof}
                                        Thus $u\mathrel{R}z$ and $z\mathrel{R}v$.
                                        Thus $z\in\intervalopenrel{R}{X}{u}{v}$ by \cref{intervalopen_rel}.
                                        Thus $z\in B_1$.
                                    \end{subproof}
                                    Thus $\intervalopenclosedrel{R}{X}{d}{c} \subseteq B_1$ by \cref{subseteq}.
                                \end{subproof}
                            \end{byCase}
                        \end{subproof}
                    \caseOf{$\exists v\in X. \exists a_0\in X. (\forall z\in X. a_0\mathrel{R}z \lor a_0 = z) \land
                        B_1 = \intervalclosedopenrel{R}{X}{a_0}{v}$.}
                        Take $a_0, v\in X$ such that
                            $(\forall z\in X. a_0\mathrel{R}z \lor a_0 = z)$ and
                            $B_1 = \intervalclosedopenrel{R}{X}{a_0}{v}$.
                        Thus $a_0$ is a smallest element of $X$ with respect to $R$ by \cref{smallest_element}.
                        Thus $c\in X$ and $(a_0\mathrel{R}c \lor c = a_0)$ and $c\mathrel{R}v$
                            by \cref{intervalclosedopen_rel}.
                        Let $d = a$.
                        Show that $\intervalopenclosedrel{R}{X}{d}{c} \subseteq B_1$.
                        \begin{subproof}
                            Show that for all $z$ such that $z\in\intervalopenclosedrel{R}{X}{d}{c}$ we have $z\in B_1$.
                            \begin{subproof}
                                Fix $z$.
                                Assume $z\in\intervalopenclosedrel{R}{X}{d}{c}$.
                                Thus $z\in X$ and $d\mathrel{R}z$ and $(z\mathrel{R}c \lor z = c)$
                                    by \cref{intervalopenclosed_rel}.
                                Show that $z\mathrel{R}v$.
                                \begin{subproof}
                                    \begin{byCase}
                                    \caseOf{$z = c$.}
                                        Thus $z\mathrel{R}v$.
                                    \caseOf{$z\mathrel{R}c$.}
                                        Thus $z\mathrel{R}v$ by \hyperref[transitive]{transitivity}.
                                    \end{byCase}
                                \end{subproof}
                                Thus $a_0\mathrel{R}z$ or $z = a_0$ by \cref{smallest_element}.
                                Thus $z\in\intervalclosedopenrel{R}{X}{a_0}{v}$ by \cref{intervalclosedopen_rel}.
                                Thus $z\in B_1$.
                            \end{subproof}
                            Thus $\intervalopenclosedrel{R}{X}{d}{c} \subseteq B_1$ by \cref{subseteq}.
                        \end{subproof}
                    \caseOf{$\exists u\in X. \exists b_0\in X. (\forall z\in X. z\mathrel{R}b_0 \lor z = b_0) \land
                        B_1 = \intervalopenclosedrel{R}{X}{u}{b_0}$.}
                        Take $u, b_0\in X$ such that
                            $(\forall z\in X. z\mathrel{R}b_0 \lor z = b_0)$ and
                            $B_1 = \intervalopenclosedrel{R}{X}{u}{b_0}$.
                        Thus $b_0$ is a largest element of $X$ with respect to $R$ by \cref{largest_element}.
                        Thus $u\mathrel{R}c$ and $(c\mathrel{R}b_0 \lor c = b_0)$ by \cref{intervalopenclosed_rel}.
                        Show that there exists $d\in Y$ such that $d\mathrel{R}c$ and
                            $\intervalopenclosedrel{R}{X}{d}{c} \subseteq B_1$.
                        \begin{subproof}
                            \begin{byCase}
                            \caseOf{$u\in Y$.}
                                Let $d = u$.
                                Show that $\intervalopenclosedrel{R}{X}{d}{c} \subseteq B_1$.
                                \begin{subproof}
                                    Show that for all $z$ such that $z\in\intervalopenclosedrel{R}{X}{d}{c}$ we have $z\in B_1$.
                                    \begin{subproof}
                                        Fix $z$.
                                        Assume $z\in\intervalopenclosedrel{R}{X}{d}{c}$.
                                        Thus $z\in X$ and $d\mathrel{R}z$ and $(z\mathrel{R}c \lor z = c)$
                                            by \cref{intervalopenclosed_rel}.
                                        Show that $z\mathrel{R}b_0$ or $z = b_0$.
                                        \begin{subproof}
                                            \begin{byCase}
                                            \caseOf{$z = c$.}
                                                Thus $z\mathrel{R}b_0$ or $z = b_0$ by \cref{intervalopenclosed_rel}.
                                            \caseOf{$z\mathrel{R}c$.}
                                                Thus $z\mathrel{R}b_0$ by \hyperref[transitive]{transitivity}.
                                            \end{byCase}
                                        \end{subproof}
                                        Thus $z\in\intervalopenclosedrel{R}{X}{u}{b_0}$ by \cref{intervalopenclosed_rel}.
                                        Thus $z\in B_1$.
                                    \end{subproof}
                                    Thus $\intervalopenclosedrel{R}{X}{d}{c} \subseteq B_1$ by \cref{subseteq}.
                                \end{subproof}
                            \caseOf{it is wrong that $u\in Y$.}
                                Show that $u\mathrel{R}a$.
                                \begin{subproof}
                                    Suppose not.
                                    Then $a\mathrel{R}u$ or $a = u$ by \cref{connex}.
                                    Show that $u\mathrel{R}b$ or $u = b$.
                                    \begin{subproof}
                                        Thus $c\mathrel{R}b$ or $c = b$ by \cref{least_upper_bound_rel}.
                                        \begin{byCase}
                                        \caseOf{$c = b$.}
                                            Thus $u\mathrel{R}b$.
                                        \caseOf{$c\mathrel{R}b$.}
                                            Thus $u\mathrel{R}b$ by \hyperref[transitive]{transitivity}.
                                        \end{byCase}
                                    \end{subproof}
                                    Thus $u\in Y$ by \cref{intervalclosed_rel}.
                                    Contradiction.
                                \end{subproof}
                                Let $d = a$.
                                Show that $\intervalopenclosedrel{R}{X}{d}{c} \subseteq B_1$.
                                \begin{subproof}
                                    Show that for all $z$ such that $z\in\intervalopenclosedrel{R}{X}{d}{c}$ we have $z\in B_1$.
                                    \begin{subproof}
                                        Fix $z$.
                                        Assume $z\in\intervalopenclosedrel{R}{X}{d}{c}$.
                                        Thus $z\in X$ and $d\mathrel{R}z$ and $(z\mathrel{R}c \lor z = c)$
                                            by \cref{intervalopenclosed_rel}.
                                        Show that $u\mathrel{R}z$.
                                        \begin{subproof}
                                            Thus $u\mathrel{R}a$.
                                            Thus $u\mathrel{R}z$ by \hyperref[transitive]{transitivity}.
                                        \end{subproof}
                                        Show that $z\mathrel{R}b_0$ or $z = b_0$.
                                        \begin{subproof}
                                            \begin{byCase}
                                            \caseOf{$z = c$.}
                                                Thus $z\mathrel{R}b_0$ or $z = b_0$ by \cref{intervalopenclosed_rel}.
                                            \caseOf{$z\mathrel{R}c$.}
                                                Thus $z\mathrel{R}b_0$ by \hyperref[transitive]{transitivity}.
                                            \end{byCase}
                                        \end{subproof}
                                        Thus $z\in\intervalopenclosedrel{R}{X}{u}{b_0}$ by \cref{intervalopenclosed_rel}.
                                        Thus $z\in B_1$.
                                    \end{subproof}
                                    Thus $\intervalopenclosedrel{R}{X}{d}{c} \subseteq B_1$ by \cref{subseteq}.
                                \end{subproof}
                            \end{byCase}
                        \end{subproof}
                    \end{byCase}
                \end{subproof}
                Take $d\in Y$ such that $d\mathrel{R}c$ and $\intervalopenclosedrel{R}{X}{d}{c} \subseteq V$.
                Suppose not.
                Then $c\notin C$.
                Show that there exists $z\in C$ such that $d\mathrel{R}z$.
                \begin{subproof}
                    Suppose not.
                    Show that for all $y\in C$ we have $y\mathrel{R}d$ or $y = d$.
                    \begin{subproof}
                        Fix $y$.
                        Assume $y\in C$.
                        Thus $y\in X$ by \cref{subseteq}.
                        Thus $d\in X$ by \cref{subseteq}.
                        Thus $y\mathrel{R}d$ or $y = d$ or $d\mathrel{R}y$ by \cref{connex}.
                        Show that $y\mathrel{R}d$ or $y = d$.
                        \begin{subproof}
                            Suppose not.
                            Then $d\mathrel{R}y$.
                            Thus there exists $z\in C$ such that $d\mathrel{R}z$.
                            Contradiction.
                        \end{subproof}
                    \end{subproof}
                    Thus $d\in X$ by \cref{subseteq}.
                    Thus $d$ is an upper bound of $C$ in $X$ with respect to $R$ by \cref{upper_bound_rel}.
                    Thus $c\mathrel{R}d$ or $c = d$ by \cref{least_upper_bound_rel}.
                    Contradiction by \cref{asymmetric}.
                \end{subproof}
                Take $z\in C$ such that $d\mathrel{R}z$.
                Thus $c$ is an upper bound of $C$ in $X$ with respect to $R$ by \cref{least_upper_bound_rel}.
                Thus $z\mathrel{R}c$ or $z = c$ by \cref{upper_bound_rel}.
                Thus $z\mathrel{R}c$.
                Take $B_0$ such that $B_0\subseteq A$ and $B_0$ is finite and
                    $B_0$ is a covering of $\intervalclosedrel{R}{X}{a}{z}$.
                Let $B = B_0\union\{U\}$.
                Show that $B$ is finite.
                \begin{subproof}
                    Let $G = \{U,U\}$.
                    $G$ is finite by \cref{pair_is_finite}.
                    Thus $U\in G$ by \cref{upair_intro_left}.
                    $\{U\}\subseteq G$ by \cref{singleton_subset_intro}.
                    Thus $\{U\}$ is finite by \cref{subseteq_finite_is_finite}.
                    Thus $B$ is finite by \cref{union_of_finite_is_finite}.
                \end{subproof}
                Show that $B$ is a covering of $\intervalclosedrel{R}{X}{a}{c}$.
                \begin{subproof}
                    Show that for all $w$ such that $w\in\intervalclosedrel{R}{X}{a}{c}$ we have $w\in\unions{B}$.
                    \begin{subproof}
                        Fix $w$.
                        Assume $w\in\intervalclosedrel{R}{X}{a}{c}$.
                        Thus $w\in X$ and $(a\mathrel{R}w \lor w = a)$ and $(w\mathrel{R}c \lor w = c)$
                            by \cref{intervalclosed_rel}.
                        $R$ is connex on $X$ by \cref{connex}.
                        Thus $w\mathrel{R}z$ or $w = z$ or $z\mathrel{R}w$.
                        \begin{byCase}
                        \caseOf{$w\mathrel{R}z$ or $w = z$.}
                            Thus $z\in X$ by \cref{subseteq}.
                            Thus $w\in\intervalclosedrel{R}{X}{a}{z}$ by \cref{intervalclosed_rel}.
                            Thus $\intervalclosedrel{R}{X}{a}{z}\subseteq\unions{B_0}$ by \cref{covering}.
                            Thus $w\in\unions{B_0}$ by \cref{subseteq}.
                            Thus $w\in\unions{B}$ by \cref{unions_iff,union_intro_left}.
                        \caseOf{$z\mathrel{R}w$.}
                            Show that $d\mathrel{R}w$.
                            \begin{subproof}
                                Thus $d\mathrel{R}z$.
                                Thus $d\mathrel{R}w$ by \hyperref[transitive]{transitivity}.
                            \end{subproof}
                            Thus $w\in\intervalopenclosedrel{R}{X}{d}{c}$ by \cref{intervalopenclosed_rel}.
                            Thus $w\in V$ by \cref{subseteq}.
                            Show that $w\in Y$.
                            \begin{subproof}
                                \begin{byCase}
                                \caseOf{$w = a$.}
                                    Thus $w\in Y$.
                                \caseOf{$a\mathrel{R}w$.}
                                    Thus $a\in Y$ and $c\in Y$.
                                    Thus $w\in Y$ by \cref{convex_in}.
                                \end{byCase}
                            \end{subproof}
                            Thus $w\in Y\inter V$ by \cref{inter_intro}.
                            Thus $w\in U$.
                            Thus $U\in\{U\}$ by \cref{singleton_intro}.
                            Thus $U\in B$ by \cref{union_intro_right}.
                            Thus $w\in\unions{B}$ by \cref{unions_iff}.
                        \end{byCase}
                    \end{subproof}
                    Thus $\intervalclosedrel{R}{X}{a}{c}\subseteq\unions{B}$ by \cref{subseteq}.
                    Thus $B$ is a covering of $\intervalclosedrel{R}{X}{a}{c}$ by \cref{covering}.
                \end{subproof}
                Thus $c\in C$.
                Contradiction.
            \end{subproof}

            Show that $c = b$.
            \begin{subproof}
                Suppose not.
                Thus $c\mathrel{R}b$ or $c = b$ by \cref{least_upper_bound_rel}.
                Show that $c\mathrel{R}b$.
                \begin{subproof}
                    Suppose not.
                    Then $c = b$.
                    Contradiction.
                \end{subproof}
                Take $y\in Y$ such that $c\mathrel{R}y$ and
                    there exists $B_1$ such that $B_1\subseteq A$ and $B_1$ is finite and
                        $B_1$ is a covering of $\intervalclosedrel{R}{X}{c}{y}$.
                Take $B_1$ such that $B_1\subseteq A$ and $B_1$ is finite and
                    $B_1$ is a covering of $\intervalclosedrel{R}{X}{c}{y}$.
                Thus $a\mathrel{R}y$ by \hyperref[transitive]{transitivity}.
                Take $B_0$ such that $B_0\subseteq A$ and $B_0$ is finite and
                    $B_0$ is a covering of $\intervalclosedrel{R}{X}{a}{c}$.
                Let $B = B_0\union B_1$.
                Show that $B$ is finite.
                \begin{subproof}
                    Thus $B_0$ is finite.
                    Thus $B_1$ is finite.
                    Thus $B$ is finite by \cref{union_of_finite_is_finite}.
                \end{subproof}
                Show that $B$ is a covering of $\intervalclosedrel{R}{X}{a}{y}$.
                \begin{subproof}
                    Show that for all $w$ such that $w\in\intervalclosedrel{R}{X}{a}{y}$ we have $w\in\unions{B}$.
                    \begin{subproof}
                        Fix $w$.
                        Assume $w\in\intervalclosedrel{R}{X}{a}{y}$.
                        Thus $w\in X$ and $(a\mathrel{R}w \lor w = a)$ and $(w\mathrel{R}y \lor w = y)$
                            by \cref{intervalclosed_rel}.
                        $R$ is connex on $X$ by \cref{connex}.
                        Thus $w\mathrel{R}c$ or $w = c$ or $c\mathrel{R}w$.
                        \begin{byCase}
                        \caseOf{$w\mathrel{R}c$ or $w = c$.}
                            Thus $w\in\intervalclosedrel{R}{X}{a}{c}$ by \cref{intervalclosed_rel}.
                            Thus $\intervalclosedrel{R}{X}{a}{c}\subseteq\unions{B_0}$ by \cref{covering}.
                            Thus $w\in\unions{B_0}$ by \cref{subseteq}.
                            Thus $w\in\unions{B}$ by \cref{unions_iff,union_intro_left}.
                        \caseOf{$c\mathrel{R}w$.}
                            Thus $w\in\intervalclosedrel{R}{X}{c}{y}$ by \cref{intervalclosed_rel}.
                            Thus $\intervalclosedrel{R}{X}{c}{y}\subseteq\unions{B_1}$ by \cref{covering}.
                            Thus $w\in\unions{B_1}$ by \cref{subseteq}.
                            Thus $w\in\unions{B}$ by \cref{unions_iff,union_intro_right}.
                        \end{byCase}
                    \end{subproof}
                    Thus $\intervalclosedrel{R}{X}{a}{y}\subseteq\unions{B}$ by \cref{subseteq}.
                    Thus $B$ is a covering of $\intervalclosedrel{R}{X}{a}{y}$ by \cref{covering}.
                \end{subproof}
                Thus $y\in C$.
                Thus $c$ is an upper bound of $C$ in $X$ with respect to $R$ by \cref{least_upper_bound_rel}.
                Thus $y\mathrel{R}c$ or $y = c$ by \cref{upper_bound_rel}.
                Contradiction by \cref{asymmetric}.
            \end{subproof}
            Take $B$ such that $B\subseteq A$ and $B$ is finite and
                $B$ is a covering of $\intervalclosedrel{R}{X}{a}{b}$.
        \end{byCase}
    \end{subproof}
    Thus $Y$ is compact with respect to $T_Y$ by \cref{compact}.
\end{proof}

% from §27 Item 2: closed interval in reals is compact
\begin{corollary}\label{real_closed_interval_compact}
    Suppose $a\in\reals$.
    Suppose $b\in\reals$.
    Suppose $a < b$ or $a = b$.
    Let $R = \realOrderRel$.
    Let $T = \realOrderTopology$.
    Let $Y = \intervalclosedrel{R}{\reals}{a}{b}$.
    Let $T_Y = \subspaceTopology{T}{Y}$.
    Then $Y$ is compact with respect to $T_Y$.
\end{corollary}
\begin{proof}
    Let $X = \reals$.
    $R$ is a simple order relation on $X$ by \cref{real_order_relation_simple}.
    $\reals$ is a linear continuum with respect to $\realOrderRel$ by \cref{reals_linear_continuum}.
    Thus $X$ has the least upper bound property with respect to $R$ by \cref{linear_continuum}.
    Thus $X$ has more than one element by \cref{linear_continuum}.
    $\realOrderTopology = \basisTopology{\orderBasis{\realOrderRel}{\reals}}{\reals}$ by \cref{real_order_topology}.
    Thus $T$ is the order topology on $X$ with respect to $R$ by \cref{order_topology}.
    Show that $a\mathrel{R}b$ or $a = b$.
    \begin{subproof}
        \begin{byCase}
        \caseOf{$a = b$.}
            Thus $a = b$.
        \caseOf{$a < b$.}
            Thus $a\mathrel{R}b$ by \cref{real_order_relation_intro}.
        \end{byCase}
    \end{subproof}
    Thus $Y$ is compact with respect to $T_Y$ by \cref{ordered_closed_interval_compact}.
\end{proof}

% from §27 helper lemma 3: four-term real sum bound
\begin{lemma}\label{real_four_sum_bound}
    Suppose $r_1,r_2,r_3,r_4,M_0\in\reals$.
    Suppose $r_1 \leq 1$ and $r_4 \leq 1$ and $r_2 \leq M_0$ and $r_3 \leq M_0$.
    Then $(r_1 + r_2) + (r_3 + r_4) \leq (M_0 + 1) + (M_0 + 1)$.
\end{lemma}
\begin{proof}
    $1\in\reals$ by \cref{real_naturals_subset_reals,real_one_in_naturals,subseteq}.
    $r_2 + 1\in\reals$ by \cref{real_add_closed}.
    $r_3 + 1\in\reals$ by \cref{real_add_closed}.
    $1 + r_2\in\reals$ by \cref{real_add_closed}.
    $1 + r_3\in\reals$ by \cref{real_add_closed}.
    $M_0 + 1\in\reals$ by \cref{real_add_closed}.
    $r_1 + r_2\in\reals$ by \cref{real_add_closed}.
    $r_3 + r_4\in\reals$ by \cref{real_add_closed}.
    $r_4 + r_3\in\reals$ by \cref{real_add_closed}.
    $(r_1 + r_2) + (r_3 + r_4)\in\reals$ by \cref{real_add_closed}.
    $(M_0 + 1) + (r_3 + r_4)\in\reals$ by \cref{real_add_closed}.
    $(r_3 + r_4) + (M_0 + 1)\in\reals$ by \cref{real_add_closed}.
    $(M_0 + 1) + (M_0 + 1)\in\reals$ by \cref{real_add_closed}.
    Thus $r_1 + r_2 \leq 1 + r_2$ by \cref{real_leq_add}.
    $1 + r_2 = r_2 + 1$ by \cref{real_add_comm}.
    Thus $1 + r_2 \leq r_2 + 1$.
    Thus $r_2 + 1 \leq M_0 + 1$ by \cref{real_leq_add}.
    Thus $1 + r_2 \leq M_0 + 1$ by \cref{real_leq_trans}.
    Thus $r_1 + r_2 \leq M_0 + 1$ by \cref{real_leq_trans}.
    Thus $r_4 + r_3 \leq 1 + r_3$ by \cref{real_leq_add}.
    $r_4 + r_3 = r_3 + r_4$ by \cref{real_add_comm}.
    Thus $r_3 + r_4 \leq r_4 + r_3$.
    Thus $r_3 + r_4 \leq 1 + r_3$ by \cref{real_leq_trans}.
    $1 + r_3 = r_3 + 1$ by \cref{real_add_comm}.
    Thus $1 + r_3 \leq r_3 + 1$.
    Thus $r_3 + 1 \leq M_0 + 1$ by \cref{real_leq_add}.
    Thus $1 + r_3 \leq M_0 + 1$ by \cref{real_leq_trans}.
    Thus $r_3 + r_4 \leq M_0 + 1$ by \cref{real_leq_trans}.
    Thus $(r_1 + r_2) + (r_3 + r_4) \leq (M_0 + 1) + (r_3 + r_4)$ by \cref{real_leq_add}.
    Thus $(r_3 + r_4) + (M_0 + 1) \leq (M_0 + 1) + (M_0 + 1)$ by \cref{real_leq_add}.
    $(M_0 + 1) + (r_3 + r_4) = (r_3 + r_4) + (M_0 + 1)$ by \cref{real_add_comm}.
    Thus $(M_0 + 1) + (r_3 + r_4) \leq (r_3 + r_4) + (M_0 + 1)$.
    Thus $(M_0 + 1) + (r_3 + r_4) \leq (M_0 + 1) + (M_0 + 1)$ by \cref{real_leq_trans}.
    Thus $(r_1 + r_2) + (r_3 + r_4) \leq (M_0 + 1) + (M_0 + 1)$ by \cref{real_leq_trans}.
\end{proof}

% from §27 helper lemma 4: compact standard closed interval
\begin{lemma}\label{real_intervalclosed_compact_standard}
    Suppose $a\in\reals$.
    Suppose $b\in\reals$.
    Suppose $a \leq b$.
    Let $I = \intervalclosed{a}{b}$.
    Then $I$ is compact with respect to $\subspaceTopology{\standardTopologyR}{I}$.
\end{lemma}
\begin{proof}
    Let $R = \realOrderRel$.
    Let $T = \realOrderTopology$.
    Let $I_0 = \intervalclosedrel{R}{\reals}{a}{b}$.
    Let $T_0 = \subspaceTopology{T}{I_0}$.
    Show that for all $x$ such that $x\in I_0$ we have $x\in I$.
    \begin{subproof}
        Fix $x$.
        Assume $x\in I_0$.
        Thus $x\in\reals$ and $(a\mathrel{R}x \lor x = a)$ and $(x\mathrel{R}b \lor x = b)$ by \cref{intervalclosed_rel}.
        Thus $a < x$ or $x = a$ by \cref{real_order_relation_elim}.
        Thus $x < b$ or $x = b$ by \cref{real_order_relation_elim}.
        Thus $a \leq x$.
        Thus $x \leq b$.
        Thus $x\in I$ by \cref{intervalclosed}.
    \end{subproof}
    Thus $I_0\subseteq I$ by \cref{subseteq}.
    Show that for all $x$ such that $x\in I$ we have $x\in I_0$.
    \begin{subproof}
        Fix $x$.
        Assume $x\in I$.
        Thus $x\in\reals$ and $a \leq x$ and $x \leq b$ by \cref{intervalclosed}.
        Thus $a < x$ or $x = a$.
        Thus $x < b$ or $x = b$.
        Thus $a\mathrel{R}x$ or $x = a$ by \cref{real_order_relation_intro}.
        Thus $x\mathrel{R}b$ or $x = b$ by \cref{real_order_relation_intro}.
        Thus $x\in I_0$ by \cref{intervalclosed_rel}.
    \end{subproof}
    Thus $I\subseteq I_0$ by \cref{subseteq}.
    Thus $I_0 = I$ by \cref{subseteq_antisymmetric}.
    Thus $I_0$ is compact with respect to $T_0$ by \cref{real_closed_interval_compact}.
    Thus $I$ is compact with respect to $\subspaceTopology{\standardTopologyR}{I}$
        by \cref{real_order_topology_eq_standard}.
\end{proof}

% from §27 helper lemma 5: indexed projection map is continuous
\begin{lemma}\label{projection_map_indexed_continuous}
    Assume $X$ is a function.
    Assume $T$ is a function.
    Assume $\dom{T} = \dom{X}$.
    Assume $\dom{X}$ is inhabited.
    Assume that for all $\alpha\in\dom{X}$ we have $\apply{T}{\alpha}$ is a topology on $\apply{X}{\alpha}$.
    Let $T_0 = \productTopologyIndexed{T}{X}$.
    Then for all $\alpha\in\dom{X}$ we have $\piIndex{X}{\alpha}$ is continuous from $\productFamily{X}$ to $\apply{X}{\alpha}$
        with respect to $T_0$ and $\apply{T}{\alpha}$.
\end{lemma}
\begin{proof}
    Let $T_0 = \productTopologyIndexed{T}{X}$.
    Let $S = \productSubbasis{T}{X}$.
    $S$ is a subbasis on $\productFamily{X}$ by \cref{product_subbasis_is_subbasis}.
    Thus $\finiteIntersections{S}$ is a basis for $T_0$ on $\productFamily{X}$ by \cref{product_subbasis_finite_intersections_basis}.
    Thus $T_0$ is a topology on $\productFamily{X}$ by \cref{basis_topology_is_topology,basis_for_topology}.
    Let $j = \identity{\productFamily{X}}$.
    $j\in\funs{\productFamily{X}}{\productFamily{X}}$ by \cref{id_elem_funs}.
    $j$ is continuous from $\productFamily{X}$ to $\productFamily{X}$ with respect to $T_0$ and $T_0$
        by \cref{identity_continuous}.
    Thus for all $\alpha\in\dom{X}$ we have $\piIndex{X}{\alpha}\circ j$ is continuous from $\productFamily{X}$
        to $\apply{X}{\alpha}$ with respect to $T_0$ and $\apply{T}{\alpha}$
        by \cref{continuous_map_into_indexed_product}.
    Fix $\alpha$.
    Assume $\alpha\in\dom{X}$.
    Show that $\piIndex{X}{\alpha}\circ j = \piIndex{X}{\alpha}$.
    \begin{subproof}
        $\piIndex{X}{\alpha}\circ j$ is a relation by \cref{circ_is_relation}.
        Show that $\piIndex{X}{\alpha}$ is a relation.
        \begin{subproof}
            Show that for all $w$ such that $w\in\piIndex{X}{\alpha}$ there exist $x,y$ such that $w = (x,y)$.
            \begin{subproof}
                Fix $w$.
                Assume $w\in\piIndex{X}{\alpha}$.
                Take $x\in\productFamily{X}$ such that $w = (x,\apply{x}{\alpha})$ by \cref{projection_map_indexed}.
                Thus there exist $x_0,y$ such that $w = (x_0,y)$.
            \end{subproof}
            Thus $\piIndex{X}{\alpha}$ is a relation by \cref{relation}.
        \end{subproof}
        Show that for all $x,z$ we have $x\mathrel{\piIndex{X}{\alpha}\circ j} z$ iff $x\mathrel{\piIndex{X}{\alpha}} z$.
        \begin{subproof}
            Fix $x$.
            Fix $z$.
            Show that if $x\mathrel{\piIndex{X}{\alpha}\circ j} z$ then $x\mathrel{\piIndex{X}{\alpha}} z$.
            \begin{subproof}
                Assume $x\mathrel{\piIndex{X}{\alpha}\circ j} z$.
                Take $y$ such that $x\mathrel{j} y$ and $y\mathrel{\piIndex{X}{\alpha}} z$ by \cref{circ_iff}.
                Thus $x = y$ and $x\in\productFamily{X}$ by \cref{id_iff}.
                Thus $x\mathrel{\piIndex{X}{\alpha}} z$.
            \end{subproof}
            Show that if $x\mathrel{\piIndex{X}{\alpha}} z$ then $x\mathrel{\piIndex{X}{\alpha}\circ j} z$.
            \begin{subproof}
                Assume $x\mathrel{\piIndex{X}{\alpha}} z$.
                Take $x_0\in\productFamily{X}$ such that $(x,z) = (x_0,\apply{x_0}{\alpha})$ by \cref{projection_map_indexed}.
                Thus $x = x_0$ by \cref{pair_eq_iff}.
                Thus $x\in\productFamily{X}$.
                Thus $x\mathrel{j} x$ by \cref{id_iff}.
                Thus $x\mathrel{\piIndex{X}{\alpha}\circ j} z$ by \cref{circ_iff}.
            \end{subproof}
        \end{subproof}
        Thus $\piIndex{X}{\alpha}\circ j = \piIndex{X}{\alpha}$ by \cref{relext}.
    \end{subproof}
    Thus $\piIndex{X}{\alpha}$ is continuous from $\productFamily{X}$ to $\apply{X}{\alpha}$
        with respect to $T_0$ and $\apply{T}{\alpha}$.
\end{proof}

% from §27 helper definition 2: indexed product compact index
\begin{definition}\label{indexed_product_compact_index}
    A natural number $m$ is indexed-product-compact iff
        for all $Y_1,T_1$ such that $Y_1$ is a function and $T_1$ is a function and
            $\dom{Y_1} = \initialSegment{m}$ and $\dom{T_1} = \initialSegment{m}$,
            if for all $i\in\initialSegment{m}$ we have $\apply{T_1}{i}$ is a topology on $\apply{Y_1}{i}$ then
            if for all $i\in\initialSegment{m}$ we have $\apply{Y_1}{i}$ is compact with respect to $\apply{T_1}{i}$ then
            $\productFamily{Y_1}$ is compact with respect to $\productTopologyIndexed{T_1}{Y_1}$.
\end{definition}

% from §27 helper definition 3: finite indexed product compactness property
\begin{definition}\label{finite_indexed_product_compact_property}
    $\finiteIndexedProductCompactProp{m}$ iff
        for all $Y_1,T_1$ such that
            $Y_1$ is a function and $T_1$ is a function and
            $\dom{Y_1} = \initialSegment{m}$ and $\dom{T_1} = \initialSegment{m}$ and
            for all $i\in\dom{Y_1}$ we have $\apply{T_1}{i}$ is a topology on $\apply{Y_1}{i}$ and
                $\apply{Y_1}{i}$ is compact with respect to $\apply{T_1}{i}$,
        we have $\productFamily{Y_1}$ is compact with respect to $\productTopologyIndexed{T_1}{Y_1}$.
\end{definition}

% from §27 helper lemma 7: continuity preserved under equality
\begin{lemma}\label{continuous_eq_subst}
    Let $f,g$ be functions.
    Suppose $f = g$.
    Suppose $g$ is continuous from $X$ to $Y$ with respect to $T$ and $T_1$.
    Then $f$ is continuous from $X$ to $Y$ with respect to $T$ and $T_1$.
\end{lemma}
\begin{proof}
    Thus $f = g$.
    Thus $f$ is continuous from $X$ to $Y$ with respect to $T$ and $T_1$.
\end{proof}

% from §27 helper lemma 8: continuity respects equal codomain/topology
\begin{lemma}\label{continuous_eq_codomain}
    Suppose $Y = Y_2$.
    Suppose $T_1 = T_2$.
    Suppose $f$ is continuous from $X$ to $Y$ with respect to $T$ and $T_1$.
    Then $f$ is continuous from $X$ to $Y_2$ with respect to $T$ and $T_2$.
\end{lemma}
\begin{proof}
    Thus $Y = Y_2$.
    Thus $T_1 = T_2$.
    Thus $f$ is continuous from $X$ to $Y_2$ with respect to $T$ and $T_2$.
\end{proof}

% from §27 helper lemma 6: compact indexed product on an initial segment
\begin{lemma}\label{compact_indexed_product_initial_segment}
    Suppose $n\in\naturalsPlus$.
    Let $J = \initialSegment{n}$.
    Assume $Y$ is a function.
    Assume $T_Y$ is a function.
    Assume $\dom{Y} = J$.
    Assume $\dom{T_Y} = J$.
    Assume that for all $i\in J$ we have $\apply{T_Y}{i}$ is a topology on $\apply{Y}{i}$.
    Assume that for all $i\in J$ we have $\apply{Y}{i}$ is compact with respect to $\apply{T_Y}{i}$.
    Then $\productFamily{Y}$ is compact with respect to $\productTopologyIndexed{T_Y}{Y}$.
\end{lemma}
\begin{proof}
    Let $A = \{m\in\naturalsPlus \mid \finiteIndexedProductCompactProp{m}\}$.
    $A\subseteq\naturalsPlus$ by \cref{subseteq}.
    Show that $1\in A$.
    \begin{subproof}
        It suffices to show that $1\in\naturalsPlus$ and $\finiteIndexedProductCompactProp{1}$.
        Show that $\finiteIndexedProductCompactProp{1}$.
        \begin{subproof}
        Show that for all $Y_1,T_1$ such that $Y_1$ is a function and $T_1$ is a function and $\dom{Y_1} = \initialSegment{1}$ and $\dom{T_1} = \initialSegment{1}$,
            if for all $i\in\initialSegment{1}$ we have $\apply{T_1}{i}$ is a topology on $\apply{Y_1}{i}$ then
            if for all $i\in\initialSegment{1}$ we have $\apply{Y_1}{i}$ is compact with respect to $\apply{T_1}{i}$ then
            $\productFamily{Y_1}$ is compact with respect to $\productTopologyIndexed{T_1}{Y_1}$.
        \begin{subproof}
            Fix $Y_1$.
            Fix $T_1$.
            Assume $Y_1$ is a function.
            Assume $T_1$ is a function.
            Assume $\dom{Y_1} = \initialSegment{1}$.
            Assume $\dom{T_1} = \initialSegment{1}$.
            Assume that for all $i\in\initialSegment{1}$ we have $\apply{T_1}{i}$ is a topology on $\apply{Y_1}{i}$.
            Assume that for all $i\in\initialSegment{1}$ we have $\apply{Y_1}{i}$ is compact with respect to $\apply{T_1}{i}$.
            Let $J_1 = \initialSegment{1}$.
            $1\in\naturalsPlus$ by \cref{real_one_in_naturals}.
            $1 \leq 1$.
            Thus $1\in J_1$ by \cref{initial_segment_naturalsplus}.
            Show that for all $i\in J_1$ we have $i = 1$.
            \begin{subproof}
                Fix $i$.
                Assume $i\in J_1$.
                Thus $i\in\naturalsPlus$ by \cref{initial_segment_naturalsplus}.
                Thus $i < 1$ or $i = 1$ by \cref{initial_segment_naturalsplus}.
                $1\leq i$ by \cref{real_one_leq_naturalsplus}.
                Thus $i = 1$.
            \end{subproof}
            Thus $J_1 = \{1,1\}$ by \cref{subseteq_antisymmetric,subseteq,upair_iff}.
            Let $X = \apply{Y_1}{1}$.
            Let $T_X = \apply{T_1}{1}$.
            $T_X$ is a topology on $X$ by \cref{subseteq}.
            $X$ is compact with respect to $T_X$ by \cref{subseteq}.
            Let $g = \{(x,\{(1,x)\}) \mid x\in X\}$.
            Show that $g\in\funs{X}{\productFamily{Y_1}}$.
            \begin{subproof}
            Show that $g$ is a function.
            \begin{subproof}
                Show that $g$ is a relation.
                \begin{subproof}
                    Show that for all $w$ such that $w\in g$ there exist $x,y$ such that $w = (x,y)$.
                    \begin{subproof}
                        Fix $w$.
                        Assume $w\in g$.
                        Take $x\in X$ such that $w = (x,\{(1,x)\})$.
                        Thus there exist $x_0,y$ such that $w = (x_0,y)$.
                    \end{subproof}
                    Thus $g$ is a relation by \cref{relation}.
                \end{subproof}
                Show that for all $x,u,v$ such that $(x,u)\in g$ and $(x,v)\in g$ we have $u = v$.
                \begin{subproof}
                    Fix $x$.
                    Fix $u$.
                    Fix $v$.
                    Assume $(x,u)\in g$.
                    Assume $(x,v)\in g$.
                    Take $x_1\in X$ such that $(x,u) = (x_1,\{(1,x_1)\})$.
                    Take $x_2\in X$ such that $(x,v) = (x_2,\{(1,x_2)\})$.
                    Thus $x = x_1$ and $u = \{(1,x_1)\}$ by \cref{pair_eq_iff}.
                    Thus $x = x_2$ and $v = \{(1,x_2)\}$ by \cref{pair_eq_iff}.
                    Thus $u = v$.
                \end{subproof}
                Thus $g$ is right-unique by \cref{rightunique}.
                Thus $g$ is a function.
            \end{subproof}
            Show that $\dom{g} = X$.
            \begin{subproof}
                Show that for all $x$ such that $x\in X$ we have $x\in\dom{g}$.
                \begin{subproof}
                    Fix $x$.
                    Assume $x\in X$.
                    Thus $(x,\{(1,x)\})\in g$.
                    Thus $x\in\dom{g}$ by \cref{dom_intro}.
                \end{subproof}
                Thus $X\subseteq\dom{g}$ by \cref{subseteq}.
                Show that for all $x$ such that $x\in\dom{g}$ we have $x\in X$.
                \begin{subproof}
                    Fix $x$.
                    Assume $x\in\dom{g}$.
                    Take $y$ such that $(x,y)\in g$ by \cref{dom_iff}.
                    Take $x_0\in X$ such that $(x,y) = (x_0,\{(1,x_0)\})$.
                    Thus $x = x_0$ by \cref{pair_eq_iff}.
                    Thus $x\in X$.
                \end{subproof}
                Thus $\dom{g}\subseteq X$ by \cref{subseteq}.
                Thus $\dom{g} = X$ by \cref{subseteq_antisymmetric}.
            \end{subproof}
            Show that for all $x\in\dom{g}$ we have $g(x)\in\productFamily{Y_1}$.
            \begin{subproof}
                Fix $x$.
                Assume $x\in\dom{g}$.
                Take $y$ such that $(x,y)\in g$ by \cref{dom_iff}.
                Take $x_0\in X$ such that $(x,y) = (x_0,\{(1,x_0)\})$.
                Thus $x = x_0$ and $y = \{(1,x_0)\}$ by \cref{pair_eq_iff}.
                Show that $y\in\productFamily{Y_1}$.
                \begin{subproof}
                    Show that $y\in\funs{\dom{Y_1}}{\unions{\ran{Y_1}}}$.
                    \begin{subproof}
                        Show that $y$ is a function.
                        \begin{subproof}
                            Show that $y$ is a relation.
                            \begin{subproof}
                                Show that for all $w$ such that $w\in y$ there exist $a,b$ such that $w = (a,b)$.
                                \begin{subproof}
                                    Fix $w$.
                                    Assume $w\in y$.
                                    Thus $w = (1,x_0)$ by \cref{singleton_iff}.
                                    Thus there exist $a,b$ such that $w = (a,b)$.
                                \end{subproof}
                                Thus $y$ is a relation by \cref{relation}.
                            \end{subproof}
                            Show that for all $a,b_1,b_2$ such that $(a,b_1)\in y$ and $(a,b_2)\in y$ we have $b_1 = b_2$.
                            \begin{subproof}
                                Fix $a$.
                                Fix $b_1$.
                                Fix $b_2$.
                                Assume $(a,b_1)\in y$ and $(a,b_2)\in y$.
                                Thus $(a,b_1) = (1,x_0)$ by \cref{singleton_iff}.
                                Thus $(a,b_2) = (1,x_0)$ by \cref{singleton_iff}.
                                Thus $b_1 = x_0$ and $b_2 = x_0$ by \cref{pair_eq_iff}.
                                Thus $b_1 = b_2$.
                            \end{subproof}
                            Thus $y$ is right-unique by \cref{rightunique}.
                            Thus $y$ is a function.
                        \end{subproof}
                        Show that $\dom{y} = \dom{Y_1}$.
                        \begin{subproof}
                            Show that for all $i$ such that $i\in\dom{Y_1}$ we have $i\in\dom{y}$.
                            \begin{subproof}
                                Fix $i$.
                                Assume $i\in\dom{Y_1}$.
                                Thus $i\in J_1$ by \cref{subseteq}.
                                Thus $i = 1$.
                                Thus $(1,x_0)\in y$ by \cref{singleton_iff}.
                                Thus $i\in\dom{y}$ by \cref{dom_intro}.
                            \end{subproof}
                            Thus $\dom{Y_1}\subseteq\dom{y}$ by \cref{subseteq}.
                            Show that for all $i$ such that $i\in\dom{y}$ we have $i\in\dom{Y_1}$.
                            \begin{subproof}
                                Fix $i$.
                                Assume $i\in\dom{y}$.
                                Take $v$ such that $(i,v)\in y$ by \cref{dom_iff}.
                                Thus $(i,v) = (1,x_0)$ by \cref{singleton_iff}.
                                Thus $i = 1$ by \cref{pair_eq_iff}.
                                Thus $i\in J_1$.
                                Thus $i\in\dom{Y_1}$ by \cref{subseteq}.
                            \end{subproof}
                            Thus $\dom{y}\subseteq\dom{Y_1}$ by \cref{subseteq}.
                            Thus $\dom{y} = \dom{Y_1}$ by \cref{subseteq_antisymmetric}.
                        \end{subproof}
                        Show that for all $i\in\dom{Y_1}$ we have $y(i)\in\unions{\ran{Y_1}}$.
                        \begin{subproof}
                            Fix $i$.
                            Assume $i\in\dom{Y_1}$.
                            Thus $i\in J_1$ by \cref{subseteq}.
                            Thus $i = 1$.
                            Thus $(1,x_0)\in y$ by \cref{singleton_iff}.
                            Thus $y(1) = x_0$ by \cref{function_apply_intro}.
                            Thus $1\in\dom{Y_1}$.
                            Thus $(1,\apply{Y_1}{1})\in Y_1$ by \cref{function_apply_elim}.
                            Thus $(1,X)\in Y_1$.
                            Thus $X\in\ran{Y_1}$ by \cref{ran_intro}.
                            Thus $x_0\in\unions{\ran{Y_1}}$ by \cref{unions_iff}.
                            Thus $y(i)\in\unions{\ran{Y_1}}$.
                        \end{subproof}
                        Thus $y$ is a function to $\unions{\ran{Y_1}}$.
                        Thus $y\in\funs{\dom{Y_1}}{\unions{\ran{Y_1}}}$ by \cref{funs_intro}.
                    \end{subproof}
                    Show that for all $i\in\dom{Y_1}$ we have $y(i)\in\apply{Y_1}{i}$.
                    \begin{subproof}
                        Fix $i$.
                        Assume $i\in\dom{Y_1}$.
                        Thus $i\in J_1$ by \cref{subseteq}.
                        Thus $i = 1$.
                        Thus $(1,x_0)\in y$ by \cref{singleton_iff}.
                        Thus $y$ is a function by \cref{funs_elim}.
                        Thus $y(1) = x_0$ by \cref{function_apply_intro}.
                        Thus $x_0\in X$ by \cref{subseteq}.
                        Thus $y(i)\in\apply{Y_1}{i}$.
                    \end{subproof}
                    Thus $y\in\productFamily{Y_1}$ by \cref{indexed_product}.
                \end{subproof}
                Thus $g(x)\in\productFamily{Y_1}$ by \cref{function_apply_intro}.
            \end{subproof}
            Thus $g$ is a function to $\productFamily{Y_1}$.
            Thus $g\in\funs{X}{\productFamily{Y_1}}$ by \cref{funs_intro}.
        \end{subproof}
        Let $T_P = \productTopologyIndexed{T_1}{Y_1}$.
        Show that $g$ is continuous from $X$ to $\productFamily{Y_1}$ with respect to $T_X$ and $T_P$.
        \begin{subproof}
            $J_1$ is inhabited.
            Show that for all $i\in J_1$ we have $\piIndex{Y_1}{i}\circ g$ is continuous from $X$ to $\apply{Y_1}{i}$ with respect to $T_X$ and $\apply{T_1}{i}$.
            \begin{subproof}
                Fix $i$.
                Assume $i\in J_1$.
                Thus $i = 1$.
                Let $j = \identity{X}$.
                $j$ is continuous from $X$ to $X$ with respect to $T_X$ and $T_X$ by \cref{identity_continuous}.
                Show that $\piIndex{Y_1}{i}\circ g = j$.
                \begin{subproof}
                    Show that for all $x\in X$ we have $\apply{\piIndex{Y_1}{i}\circ g}{x} = \apply{j}{x}$.
                    \begin{subproof}
                        Fix $x$.
                        Assume $x\in X$.
                        Thus $g$ is a function to $\productFamily{Y_1}$ and $\dom{g} = X$ by \cref{funs_elim}.
                        Thus $x\in\dom{g}$.
                        Thus $(x,g(x))\in g$ by \cref{function_apply_elim}.
                        Thus $g(x)\in\ran{g}$ by \cref{ran_intro}.
                        Thus $\ran{g}\subseteq\productFamily{Y_1}$ by \cref{function_to_ran}.
                        Thus $g(x)\in\productFamily{Y_1}$ by \cref{subseteq}.
                        Thus $(x,\{(1,x)\})\in g$.
                        Thus $g(x) = \{(1,x)\}$ by \cref{function_apply_intro}.
                        Thus $g(x)\in\funs{\dom{Y_1}}{\unions{\ran{Y_1}}}$ by \cref{indexed_product}.
                        Thus $g(x)$ is a function by \cref{funs_elim}.
                        Thus $(1,x)\in g(x)$ by \cref{singleton_iff}.
                        Thus $\apply{g(x)}{1} = x$ by \cref{function_apply_intro}.
                        Thus $\apply{g(x)}{i} = x$.
                        Thus $(g(x),x) = (g(x),\apply{g(x)}{i})$ by \cref{pair_eq_iff}.
                        Thus $(g(x),x)\in\piIndex{Y_1}{i}$ by \cref{projection_map_indexed}.
                        Thus $(x,x)\in\piIndex{Y_1}{i}\circ g$ by \cref{circ_iff}.
                        Thus $\piIndex{Y_1}{i}\in\funs{\productFamily{Y_1}}{\apply{Y_1}{i}}$ by \cref{projection_map_indexed_function}.
                        Thus $\piIndex{Y_1}{i}$ is a function by \cref{funs_is_function}.
                        Thus $\piIndex{Y_1}{i}\circ g$ is a function by \cref{function_circ}.
                        Thus $\apply{\piIndex{Y_1}{i}\circ g}{x} = x$ by \cref{function_apply_intro}.
                        Thus $\apply{j}{x} = x$ by \cref{id_apply}.
                        Thus $\apply{\piIndex{Y_1}{i}\circ g}{x} = \apply{j}{x}$.
                    \end{subproof}
                    Show that $\piIndex{Y_1}{i}\circ g = j$.
                    \begin{subproof}
                        $g$ is a function to $\productFamily{Y_1}$ and $\dom{g} = X$ by \cref{funs_elim}.
                        $\piIndex{Y_1}{i}\in\funs{\productFamily{Y_1}}{\apply{Y_1}{i}}$ by \cref{projection_map_indexed_function}.
                        Thus $\piIndex{Y_1}{i}$ is a function by \cref{funs_is_function}.
                        $\piIndex{Y_1}{i}\circ g$ is a function by \cref{function_circ}.
                        $j$ is a function by \cref{id_is_function}.
                        Show that $\dom{\piIndex{Y_1}{i}\circ g} = X$.
                        \begin{subproof}
                            Show that for all $x$ such that $x\in\dom{\piIndex{Y_1}{i}\circ g}$ we have $x\in X$.
                            \begin{subproof}
                                Fix $x$.
                                Assume $x\in\dom{\piIndex{Y_1}{i}\circ g}$.
                                Take $z$ such that $(x,z)\in\piIndex{Y_1}{i}\circ g$ by \cref{dom_iff}.
                                Take $y$ such that $x\mathrel{g} y$ and $y\mathrel{\piIndex{Y_1}{i}} z$ by \cref{circ_iff}.
                                Thus $x\in\dom{g}$ by \cref{dom_intro}.
                                Thus $x\in X$ by \cref{subseteq}.
                            \end{subproof}
                            Thus $\dom{\piIndex{Y_1}{i}\circ g}\subseteq X$ by \cref{subseteq}.
                            Show that for all $x$ such that $x\in X$ we have $x\in\dom{\piIndex{Y_1}{i}\circ g}$.
                            \begin{subproof}
                                Fix $x$.
                                Assume $x\in X$.
                                Thus $x\in\dom{g}$ by \cref{subseteq}.
                                Thus $(x,g(x))\in g$ by \cref{function_apply_elim}.
                                Thus $g(x)\in\ran{g}$ by \cref{ran_intro}.
                                Thus $\ran{g}\subseteq\productFamily{Y_1}$ by \cref{function_to_ran}.
                                Thus $g(x)\in\productFamily{Y_1}$ by \cref{subseteq}.
                                Thus $(g(x),\apply{g(x)}{i})\in\piIndex{Y_1}{i}$ by \cref{projection_map_indexed}.
                                Thus $(x,\apply{g(x)}{i})\in\piIndex{Y_1}{i}\circ g$ by \cref{circ_iff}.
                                Thus $x\in\dom{\piIndex{Y_1}{i}\circ g}$ by \cref{dom_intro}.
                            \end{subproof}
                            Thus $X\subseteq\dom{\piIndex{Y_1}{i}\circ g}$ by \cref{subseteq}.
                            Thus $\dom{\piIndex{Y_1}{i}\circ g} = X$ by \cref{subseteq_antisymmetric}.
                        \end{subproof}
                        Thus $\dom{j} = X$ by \cref{id_dom}.
                        Show that for all $x$ such that $x\in\dom{\piIndex{Y_1}{i}\circ g}$ we have $x\in\dom{j}$.
                        \begin{subproof}
                            Fix $x$.
                            Assume $x\in\dom{\piIndex{Y_1}{i}\circ g}$.
                            Thus $x\in X$ by \cref{subseteq}.
                            Thus $x\in\dom{j}$ by \cref{subseteq}.
                        \end{subproof}
                        Thus $\dom{\piIndex{Y_1}{i}\circ g}\subseteq\dom{j}$ by \cref{subseteq}.
                        Thus for all $x\in\dom{\piIndex{Y_1}{i}\circ g}$ we have $\apply{\piIndex{Y_1}{i}\circ g}{x} = \apply{j}{x}$ by \cref{subseteq}.
                        Thus $\piIndex{Y_1}{i}\circ g\subseteq j$ by \cref{fun_subseteq}.
                        Show that for all $x$ such that $x\in\dom{j}$ we have $x\in\dom{\piIndex{Y_1}{i}\circ g}$.
                        \begin{subproof}
                            Fix $x$.
                            Assume $x\in\dom{j}$.
                            Thus $x\in X$ by \cref{subseteq}.
                            Thus $x\in\dom{\piIndex{Y_1}{i}\circ g}$ by \cref{subseteq}.
                        \end{subproof}
                        Thus $\dom{j}\subseteq\dom{\piIndex{Y_1}{i}\circ g}$ by \cref{subseteq}.
                        Thus for all $x\in\dom{j}$ we have $\apply{\piIndex{Y_1}{i}\circ g}{x} = \apply{j}{x}$ by \cref{subseteq}.
                        Thus $j\subseteq \piIndex{Y_1}{i}\circ g$ by \cref{fun_subseteq}.
                        Thus $\piIndex{Y_1}{i}\circ g = j$ by \cref{subseteq_antisymmetric}.
                    \end{subproof}
                \end{subproof}
                Thus $\piIndex{Y_1}{i}\circ g$ is continuous from $X$ to $\apply{Y_1}{i}$ with respect to $T_X$ and $\apply{T_1}{i}$.
            \end{subproof}
            Thus $g$ is continuous from $X$ to $\productFamily{Y_1}$ with respect to $T_X$ and $T_P$
                by \cref{continuous_map_into_indexed_product}.
        \end{subproof}
        Show that $g\in\surj{X}{\productFamily{Y_1}}$.
        \begin{subproof}
            Show that for all $y\in\productFamily{Y_1}$ there exists $x\in X$ such that $g(x) = y$.
            \begin{subproof}
                Fix $y$.
                Assume $y\in\productFamily{Y_1}$.
                Thus $y\in\funs{\dom{Y_1}}{\unions{\ran{Y_1}}}$ by \cref{indexed_product}.
                Thus $y$ is a function and $\dom{y} = \dom{Y_1}$ by \cref{funs_elim}.
                Thus $1\in\dom{y}$ by \cref{subseteq,real_one_in_naturals,initial_segment_naturalsplus}.
                Take $x\in X$ such that $x = y(1)$.
                Thus $(x,\{(1,x)\})\in g$.
                Show that $g(x) = y$.
                \begin{subproof}
                    Show that for all $i\in\dom{y}$ we have $\apply{g(x)}{i} = y(i)$.
                    \begin{subproof}
                        Fix $i$.
                        Assume $i\in\dom{y}$.
                        Thus $i\in J_1$ by \cref{subseteq}.
                        Thus $i = 1$.
                        Thus $g$ is a function to $\productFamily{Y_1}$ and $\dom{g} = X$ by \cref{funs_elim}.
                        Thus $g(x) = \{(1,x)\}$ by \cref{function_apply_intro}.
                        Thus $(1,x)\in g(x)$ by \cref{singleton_iff}.
                        Thus $x\in\dom{g}$ by \cref{subseteq}.
                        Thus $(x,g(x))\in g$ by \cref{function_apply_elim}.
                        Thus $g(x)\in\ran{g}$ by \cref{ran_intro}.
                        Thus $\ran{g}\subseteq\productFamily{Y_1}$ by \cref{function_to_ran}.
                        Thus $g(x)\in\productFamily{Y_1}$ by \cref{subseteq}.
                        Thus $g(x)\in\funs{\dom{Y_1}}{\unions{\ran{Y_1}}}$ by \cref{indexed_product}.
                        Thus $g(x)$ is a function by \cref{funs_elim}.
                        Thus $\apply{g(x)}{1} = x$ by \cref{function_apply_intro}.
                        Thus $\apply{g(x)}{i} = y(i)$.
                    \end{subproof}
                    Show that $g(x) = y$.
                    \begin{subproof}
                        Thus $y$ is a function and $\dom{y} = \dom{Y_1}$ by \cref{funs_elim}.
                        Thus $g$ is a function to $\productFamily{Y_1}$ and $\dom{g} = X$ by \cref{funs_elim}.
                        Thus $x\in\dom{g}$ by \cref{subseteq}.
                        Thus $(x,g(x))\in g$ by \cref{function_apply_elim}.
                        Thus $g(x)\in\ran{g}$ by \cref{ran_intro}.
                        Thus $\ran{g}\subseteq\productFamily{Y_1}$ by \cref{function_to_ran}.
                        Thus $g(x)\in\productFamily{Y_1}$ by \cref{subseteq}.
                        Thus $g(x)\in\funs{\dom{Y_1}}{\unions{\ran{Y_1}}}$ by \cref{indexed_product}.
                        Thus $g(x)$ is a function by \cref{funs_elim}.
                        Thus $g(x)\in\rels{\dom{Y_1}}{\unions{\ran{Y_1}}}$ and $g(x)$ is right-unique by \cref{funs}.
                        Thus $g(x)$ is a relation by \cref{rels_is_relation}.
                        Thus $\dom{g(x)} = \dom{Y_1}$ by \cref{funs_elim}.
                        Thus $\dom{g(x)}\subseteq\dom{y}$ by \cref{subseteq}.
                        Thus for all $i\in\dom{g(x)}$ we have $\apply{g(x)}{i} = y(i)$ by \cref{subseteq}.
                        Thus $g(x)\subseteq y$ by \cref{fun_subseteq}.
                        Thus $\dom{y}\subseteq\dom{g(x)}$ by \cref{subseteq}.
                        Thus $y\subseteq g(x)$ by \cref{fun_subseteq}.
                        Thus $g(x) = y$ by \cref{subseteq_antisymmetric}.
                    \end{subproof}
                \end{subproof}
                Thus there exists $x\in X$ such that $g(x) = y$.
            \end{subproof}
            Thus $g\in\surj{X}{\productFamily{Y_1}}$ by \cref{surj}.
        \end{subproof}
        Thus $\dom{g} = X$ by \cref{surj,funs_elim}.
        Thus $\img{g}{X} = \ran{g}$ by \cref{img_dom}.
        Thus $\ran{g} = \productFamily{Y_1}$ by \cref{funs_surj_iff}.
        Let $A = \img{g}{X}$.
        Let $T_A = \subspaceTopology{T_P}{A}$.
        Thus $A = \productFamily{Y_1}$.
        Thus $T_A = \subspaceTopology{T_P}{\productFamily{Y_1}}$.
        Show that $T_A = T_P$.
        \begin{subproof}
            $T_P$ is a topology on $\productFamily{Y_1}$
                by \cref{product_subbasis_finite_intersections_basis,basis_for_topology,basis_topology_is_topology}.
            Show that $T_A\subseteq T_P$.
            \begin{subproof}
                Show that for all $U$ such that $U\in T_A$ we have $U\in T_P$.
                \begin{subproof}
                    Fix $U$.
                    Assume $U\in T_A$.
                    Take $V\in T_P$ such that $U = \productFamily{Y_1}\inter V$ by \cref{subspace_topology}.
                    Thus $T_P\subseteq\pow{\productFamily{Y_1}}$ by \cref{topology_on}.
                    Thus $V\in\pow{\productFamily{Y_1}}$ by \cref{subseteq}.
                    Thus $V\subseteq\productFamily{Y_1}$ by \cref{powerset_elim}.
                    Thus $U = V$ by \cref{inter_absorb_supseteq_left}.
                    Thus $U\in T_P$.
                \end{subproof}
                Thus $T_A\subseteq T_P$ by \cref{subseteq}.
            \end{subproof}
            Show that $T_P\subseteq T_A$.
            \begin{subproof}
                Show that for all $U$ such that $U\in T_P$ we have $U\in T_A$.
                \begin{subproof}
                    Fix $U$.
                    Assume $U\in T_P$.
                    Thus $T_P\subseteq\pow{\productFamily{Y_1}}$ by \cref{topology_on}.
                    Thus $U\in\pow{\productFamily{Y_1}}$ by \cref{subseteq}.
                    Thus $U\subseteq\productFamily{Y_1}$ by \cref{powerset_elim}.
                    Thus $U = \productFamily{Y_1}\inter U$ by \cref{inter_absorb_supseteq_left}.
                    Thus $U\in T_A$ by \cref{subspace_topology}.
                \end{subproof}
                Thus $T_P\subseteq T_A$ by \cref{subseteq}.
            \end{subproof}
            Thus $T_A = T_P$ by \cref{subseteq_antisymmetric}.
        \end{subproof}
        Thus $A$ is compact with respect to $T_A$ by \cref{continuous_image_compact}.
        Thus $\productFamily{Y_1}$ is compact with respect to $T_A$.
        Thus $\productFamily{Y_1}$ is compact with respect to $T_P$.
    \end{subproof}
        Thus $\finiteIndexedProductCompactProp{1}$ by \cref{finite_indexed_product_compact_property}.
        \end{subproof}
        $1\in\naturalsPlus$ by \cref{real_one_in_naturals}.
        Thus $1\in\naturalsPlus$ and $\finiteIndexedProductCompactProp{1}$.
    \end{subproof}
    Show that for all $k\in A$ we have $k + 1\in A$.
    \begin{subproof}
        Fix $k$.
        Assume $k\in A$.
        Thus $k\in\naturalsPlus$ by \cref{subseteq}.
        Thus $\finiteIndexedProductCompactProp{k}$.
        Thus for all $Y_1,T_1$ such that $Y_1$ is a function and $T_1$ is a function and $\dom{Y_1} = \initialSegment{k}$ and $\dom{T_1} = \initialSegment{k}$,
            if for all $i\in\initialSegment{k}$ we have $\apply{T_1}{i}$ is a topology on $\apply{Y_1}{i}$ then
            if for all $i\in\initialSegment{k}$ we have $\apply{Y_1}{i}$ is compact with respect to $\apply{T_1}{i}$ then
            $\productFamily{Y_1}$ is compact with respect to $\productTopologyIndexed{T_1}{Y_1}$ by \cref{finite_indexed_product_compact_property}.
        Show that $k + 1\in A$.
        \begin{subproof}
            It suffices to show that $k + 1\in\naturalsPlus$ and $\finiteIndexedProductCompactProp{k + 1}$.
            Show that $\finiteIndexedProductCompactProp{k + 1}$.
            \begin{subproof}
            Show that for all $Y_1,T_1$ such that $Y_1$ is a function and $T_1$ is a function and $\dom{Y_1} = \initialSegment{k + 1}$ and $\dom{T_1} = \initialSegment{k + 1}$,
                if for all $i\in\initialSegment{k + 1}$ we have $\apply{T_1}{i}$ is a topology on $\apply{Y_1}{i}$ then
                if for all $i\in\initialSegment{k + 1}$ we have $\apply{Y_1}{i}$ is compact with respect to $\apply{T_1}{i}$ then
                $\productFamily{Y_1}$ is compact with respect to $\productTopologyIndexed{T_1}{Y_1}$.
            \begin{subproof}
                Fix $Y_1$.
                Fix $T_1$.
                Assume $Y_1$ is a function.
                Assume $T_1$ is a function.
                Assume $\dom{Y_1} = \initialSegment{k + 1}$.
                Assume $\dom{T_1} = \initialSegment{k + 1}$.
                Assume that for all $i\in\initialSegment{k + 1}$ we have $\apply{T_1}{i}$ is a topology on $\apply{Y_1}{i}$.
                Assume that for all $i\in\initialSegment{k + 1}$ we have $\apply{Y_1}{i}$ is compact with respect to $\apply{T_1}{i}$.
                Let $J_1 = \initialSegment{k}$.
                Let $J = \initialSegment{k + 1}$.
                Show that $J_1\subseteq J$.
                \begin{subproof}
                    Show that for all $i$ such that $i\in J_1$ we have $i\in J$.
                    \begin{subproof}
                        Fix $i$.
                        Assume $i\in J_1$.
                        Thus $i\in\naturalsPlus$ and $i \leq k$ by \cref{initial_segment_naturalsplus}.
                        Thus $i\in\reals$ by \cref{real_naturals_subset_reals,subseteq}.
                        $k\in\naturalsPlus$ by \cref{subseteq}.
                        $k\in\reals$ by \cref{real_naturals_subset_reals,subseteq}.
                        $k + 1\in\naturalsPlus$ by \cref{real_naturals_closed_succ}.
                        Thus $k + 1\in\reals$ by \cref{real_naturals_subset_reals,subseteq}.
                        $0\in\reals$ by \cref{real_add_ident}.
                        $1\in\reals$ by \cref{real_naturals_subset_reals,real_one_in_naturals,subseteq}.
                        $0 < 1$ by \cref{real_zero_lt_one}.
                        $0 + k < 1 + k$ by \cref{real_lt_add}.
                        $0 + k = k + 0$ by \cref{real_add_comm}.
                        $1 + k = k + 1$ by \cref{real_add_comm}.
                        Thus $k + 0 < k + 1$ by \cref{real_lt_subst_left}.
                        $k + 0 = k$ by \cref{real_add_ident}.
                        Thus $k < k + 1$ by \cref{real_lt_subst_left}.
                        Thus $k \leq k + 1$ by \cref{real_lt_imp_leq}.
                        Thus $i \leq k + 1$ by \cref{real_leq_trans}.
                        Thus $i\in J$ by \cref{initial_segment_naturalsplus}.
                    \end{subproof}
                    Thus $J_1\subseteq J$ by \cref{subseteq}.
                \end{subproof}
                $k\in\reals$ by \cref{real_naturals_subset_reals,subseteq}.
                $k + 1\in\naturalsPlus$ by \cref{real_naturals_closed_succ}.
                Thus $k + 1\in\reals$ by \cref{real_naturals_subset_reals,subseteq}.
                $0\in\reals$ by \cref{real_add_ident}.
                $1\in\reals$ by \cref{real_naturals_subset_reals,real_one_in_naturals,subseteq}.
                $0 < 1$ by \cref{real_zero_lt_one}.
                $0 + k < 1 + k$ by \cref{real_lt_add}.
                $0 + k = k + 0$ by \cref{real_add_comm}.
                $1 + k = k + 1$ by \cref{real_add_comm}.
                Thus $k + 0 < k + 1$ by \cref{real_lt_subst_left}.
                $k + 0 = k$ by \cref{real_add_ident}.
                Thus $k < k + 1$ by \cref{real_lt_subst_left}.
                Show that $k \neq k + 1$.
                \begin{subproof}
                    Suppose not.
                    Thus $k = k + 1$.
                    Thus $k < k$ by \cref{real_lt_subst_left}.
                    Contradiction by \cref{real_lt_irreflexive}.
                \end{subproof}
                Let $Y_0 = \{(i,\apply{Y_1}{i}) \mid i\in J_1\}$.
                Let $T_0 = \{(i,\apply{T_1}{i}) \mid i\in J_1\}$.
                Show that $\dom{Y_0} = J_1$.
                \begin{subproof}
                    Show that for all $i$ such that $i\in J_1$ we have $i\in\dom{Y_0}$.
                    \begin{subproof}
                        Fix $i$.
                        Assume $i\in J_1$.
                        Thus $(i,\apply{Y_1}{i})\in Y_0$.
                        Thus $i\in\dom{Y_0}$ by \cref{dom_intro}.
                    \end{subproof}
                    Thus $J_1\subseteq\dom{Y_0}$ by \cref{subseteq}.
                    Show that for all $i$ such that $i\in\dom{Y_0}$ we have $i\in J_1$.
                    \begin{subproof}
                        Fix $i$.
                        Assume $i\in\dom{Y_0}$.
                        Take $v$ such that $(i,v)\in Y_0$ by \cref{dom_iff}.
                        Take $j\in J_1$ such that $(i,v) = (j,\apply{Y_1}{j})$.
                        Thus $i = j$ by \cref{pair_eq_iff}.
                        Thus $i\in J_1$.
                    \end{subproof}
                    Thus $\dom{Y_0}\subseteq J_1$ by \cref{subseteq}.
                    Thus $\dom{Y_0} = J_1$ by \cref{subseteq_antisymmetric}.
                \end{subproof}
                Show that $\dom{T_0} = J_1$.
                \begin{subproof}
                    Show that for all $i$ such that $i\in J_1$ we have $i\in\dom{T_0}$.
                    \begin{subproof}
                        Fix $i$.
                        Assume $i\in J_1$.
                        Thus $(i,\apply{T_1}{i})\in T_0$.
                        Thus $i\in\dom{T_0}$ by \cref{dom_intro}.
                    \end{subproof}
                    Thus $J_1\subseteq\dom{T_0}$ by \cref{subseteq}.
                    Show that for all $i$ such that $i\in\dom{T_0}$ we have $i\in J_1$.
                    \begin{subproof}
                        Fix $i$.
                        Assume $i\in\dom{T_0}$.
                        Take $v$ such that $(i,v)\in T_0$ by \cref{dom_iff}.
                        Take $j\in J_1$ such that $(i,v) = (j,\apply{T_1}{j})$.
                        Thus $i = j$ by \cref{pair_eq_iff}.
                        Thus $i\in J_1$.
                    \end{subproof}
                    Thus $\dom{T_0}\subseteq J_1$ by \cref{subseteq}.
                    Thus $\dom{T_0} = J_1$ by \cref{subseteq_antisymmetric}.
                \end{subproof}
                Show that $Y_0$ is a function.
                \begin{subproof}
                    Show that $Y_0$ is a relation.
                    \begin{subproof}
                        Show that for all $w$ such that $w\in Y_0$ there exist $a,b$ such that $w = (a,b)$.
                        \begin{subproof}
                            Fix $w$.
                            Assume $w\in Y_0$.
                            Take $i\in J_1$ such that $w = (i,\apply{Y_1}{i})$.
                            Thus there exist $a,b$ such that $w = (a,b)$.
                        \end{subproof}
                        Thus $Y_0$ is a relation by \cref{relation}.
                    \end{subproof}
                    Show that for all $a,b_1,b_2$ such that $(a,b_1)\in Y_0$ and $(a,b_2)\in Y_0$ we have $b_1 = b_2$.
                    \begin{subproof}
                        Fix $a$.
                        Fix $b_1$.
                        Fix $b_2$.
                        Assume $(a,b_1)\in Y_0$ and $(a,b_2)\in Y_0$.
                        Take $i_1\in J_1$ such that $(a,b_1) = (i_1,\apply{Y_1}{i_1})$.
                        Take $i_2\in J_1$ such that $(a,b_2) = (i_2,\apply{Y_1}{i_2})$.
                        Thus $a = i_1$ and $b_1 = \apply{Y_1}{i_1}$ by \cref{pair_eq_iff}.
                        Thus $a = i_2$ and $b_2 = \apply{Y_1}{i_2}$ by \cref{pair_eq_iff}.
                        Thus $i_1 = i_2$.
                        Thus $b_1 = b_2$.
                    \end{subproof}
                    Thus $Y_0$ is right-unique by \cref{rightunique}.
                    Thus $Y_0$ is a function.
                \end{subproof}
                Show that $T_0$ is a function.
                \begin{subproof}
                    Show that $T_0$ is a relation.
                    \begin{subproof}
                        Show that for all $w$ such that $w\in T_0$ there exist $a,b$ such that $w = (a,b)$.
                        \begin{subproof}
                            Fix $w$.
                            Assume $w\in T_0$.
                            Take $i\in J_1$ such that $w = (i,\apply{T_1}{i})$.
                            Thus there exist $a,b$ such that $w = (a,b)$.
                        \end{subproof}
                        Thus $T_0$ is a relation by \cref{relation}.
                    \end{subproof}
                    Show that for all $a,b_1,b_2$ such that $(a,b_1)\in T_0$ and $(a,b_2)\in T_0$ we have $b_1 = b_2$.
                    \begin{subproof}
                        Fix $a$.
                        Fix $b_1$.
                        Fix $b_2$.
                        Assume $(a,b_1)\in T_0$ and $(a,b_2)\in T_0$.
                        Take $i_1\in J_1$ such that $(a,b_1) = (i_1,\apply{T_1}{i_1})$.
                        Take $i_2\in J_1$ such that $(a,b_2) = (i_2,\apply{T_1}{i_2})$.
                        Thus $a = i_1$ and $b_1 = \apply{T_1}{i_1}$ by \cref{pair_eq_iff}.
                        Thus $a = i_2$ and $b_2 = \apply{T_1}{i_2}$ by \cref{pair_eq_iff}.
                        Thus $i_1 = i_2$.
                        Thus $b_1 = b_2$.
                    \end{subproof}
                    Thus $T_0$ is right-unique by \cref{rightunique}.
                    Thus $T_0$ is a function.
                \end{subproof}
                Show that for all $i\in J_1$ we have $\apply{T_0}{i}$ is a topology on $\apply{Y_0}{i}$.
                \begin{subproof}
                    Fix $i$.
                    Assume $i\in J_1$.
                    Thus $i\in J$ by \cref{subseteq}.
                    Thus $\apply{T_1}{i}$ is a topology on $\apply{Y_1}{i}$ by \cref{subseteq}.
                    Thus $(i,\apply{Y_1}{i})\in Y_0$.
                    Thus $\apply{Y_0}{i} = \apply{Y_1}{i}$ by \cref{function_apply_intro}.
                    Thus $(i,\apply{T_1}{i})\in T_0$.
                    Thus $\apply{T_0}{i} = \apply{T_1}{i}$ by \cref{function_apply_intro}.
                    Thus $\apply{T_0}{i}$ is a topology on $\apply{Y_0}{i}$ by \cref{subseteq}.
                \end{subproof}
                Show that for all $i\in J_1$ we have $\apply{Y_0}{i}$ is compact with respect to $\apply{T_0}{i}$.
                \begin{subproof}
                    Fix $i$.
                    Assume $i\in J_1$.
                    Thus $i\in J$ by \cref{subseteq}.
                    Thus $\apply{Y_1}{i}$ is compact with respect to $\apply{T_1}{i}$ by \cref{subseteq}.
                    Thus $(i,\apply{Y_1}{i})\in Y_0$.
                    Thus $\apply{Y_0}{i} = \apply{Y_1}{i}$ by \cref{function_apply_intro}.
                    Thus $(i,\apply{T_1}{i})\in T_0$.
                    Thus $\apply{T_0}{i} = \apply{T_1}{i}$ by \cref{function_apply_intro}.
                    Thus $\apply{Y_0}{i}$ is compact with respect to $\apply{T_0}{i}$ by \cref{subseteq}.
                \end{subproof}
                Thus $\productFamily{Y_0}$ is compact with respect to $\productTopologyIndexed{T_0}{Y_0}$ by \cref{subseteq}.
                Let $X_0 = \productFamily{Y_0}$.
                Let $T_X = \productTopologyIndexed{T_0}{Y_0}$.
                Let $X_1 = \apply{Y_1}{k + 1}$.
                Let $T_{1p} = \apply{T_1}{k + 1}$.
                $k + 1\in\naturalsPlus$ by \cref{real_naturals_closed_succ}.
                $k + 1 \leq k + 1$.
                Thus $k + 1\in J$ by \cref{initial_segment_naturalsplus}.
                $T_{1p}$ is a topology on $X_1$ by \cref{subseteq}.
                $X_1$ is compact with respect to $T_{1p}$ by \cref{subseteq}.
                $1\in\naturalsPlus$ by \cref{real_one_in_naturals}.
                $1 \leq k$ by \cref{real_one_leq_naturalsplus}.
                Thus $1\in J_1$ by \cref{initial_segment_naturalsplus}.
                Thus $J_1$ is inhabited.
                Thus $T_X$ is a topology on $X_0$
                    by \cref{product_subbasis_finite_intersections_basis,basis_for_topology,basis_topology_is_topology}.
                Let $T_P = \productTopology{T_X}{T_{1p}}{X_0}{X_1}$.
                Thus $X_0\times X_1$ is compact with respect to $T_P$ by \cref{product_compact}.
                Let $g = \{(p,\fst{p}\union\{(k + 1,\snd{p})\}) \mid p\in X_0\times X_1\}$.
                Show that $g\in\funs{X_0\times X_1}{\productFamily{Y_1}}$.
                \begin{subproof}
                    Show that $g$ is a function.
                    \begin{subproof}
                        Show that $g$ is a relation.
                        \begin{subproof}
                            Show that for all $w$ such that $w\in g$ there exist $x,y$ such that $w = (x,y)$.
                            \begin{subproof}
                                Fix $w$.
                                Assume $w\in g$.
                                Take $p\in X_0\times X_1$ such that $w = (p,\fst{p}\union\{(k + 1,\snd{p})\})$.
                                Thus there exist $x,y$ such that $w = (x,y)$.
                            \end{subproof}
                            Thus $g$ is a relation by \cref{relation}.
                        \end{subproof}
                        Show that for all $p,u,v$ such that $(p,u)\in g$ and $(p,v)\in g$ we have $u = v$.
                        \begin{subproof}
                            Fix $p$.
                            Fix $u$.
                            Fix $v$.
                            Assume $(p,u)\in g$.
                            Assume $(p,v)\in g$.
                            Take $p_1\in X_0\times X_1$ such that $(p,u) = (p_1,\fst{p_1}\union\{(k + 1,\snd{p_1})\})$.
                            Take $p_2\in X_0\times X_1$ such that $(p,v) = (p_2,\fst{p_2}\union\{(k + 1,\snd{p_2})\})$.
                            Thus $p = p_1$ and $u = \fst{p_1}\union\{(k + 1,\snd{p_1})\}$ by \cref{pair_eq_iff}.
                            Thus $p = p_2$ and $v = \fst{p_2}\union\{(k + 1,\snd{p_2})\}$ by \cref{pair_eq_iff}.
                            Thus $u = v$.
                        \end{subproof}
                        Thus $g$ is right-unique by \cref{rightunique}.
                        Thus $g$ is a function.
                    \end{subproof}
                    Show that $\dom{g} = X_0\times X_1$.
                    \begin{subproof}
                        Show that for all $p$ such that $p\in X_0\times X_1$ we have $p\in\dom{g}$.
                        \begin{subproof}
                            Fix $p$.
                            Assume $p\in X_0\times X_1$.
                            Thus $(p,\fst{p}\union\{(k + 1,\snd{p})\})\in g$.
                            Thus $p\in\dom{g}$ by \cref{dom_intro}.
                        \end{subproof}
                        Thus $X_0\times X_1\subseteq\dom{g}$ by \cref{subseteq}.
                        Show that for all $p$ such that $p\in\dom{g}$ we have $p\in X_0\times X_1$.
                        \begin{subproof}
                            Fix $p$.
                            Assume $p\in\dom{g}$.
                            Take $v$ such that $(p,v)\in g$ by \cref{dom_iff}.
                            Take $p_0\in X_0\times X_1$ such that $(p,v) = (p_0,\fst{p_0}\union\{(k + 1,\snd{p_0})\})$.
                            Thus $p = p_0$ by \cref{pair_eq_iff}.
                            Thus $p\in X_0\times X_1$.
                        \end{subproof}
                        Thus $\dom{g}\subseteq X_0\times X_1$ by \cref{subseteq}.
                        Thus $\dom{g} = X_0\times X_1$ by \cref{subseteq_antisymmetric}.
                    \end{subproof}
                    Show that for all $p\in\dom{g}$ we have $g(p)\in\productFamily{Y_1}$.
                    \begin{subproof}
                        Fix $p$.
                        Assume $p\in\dom{g}$.
                        Take $v$ such that $(p,v)\in g$ by \cref{dom_iff}.
                        Take $p_0\in X_0\times X_1$ such that $(p,v) = (p_0,\fst{p_0}\union\{(k + 1,\snd{p_0})\})$.
                        Thus $p = p_0$ and $v = \fst{p}\union\{(k + 1,\snd{p})\}$ by \cref{pair_eq_iff}.
                        Show that $v\in\productFamily{Y_1}$.
                        \begin{subproof}
                            Let $x = \fst{p}$.
                            Let $y = \snd{p}$.
                            Thus $p\in X_0\times X_1$ by \cref{subseteq}.
                            $x\in X_0$ and $y\in X_1$ by \cref{times_proj_elim}.
                            Thus $x\in\productFamily{Y_0}$ by \cref{subseteq}.
                            Thus $x\in\funs{\dom{Y_0}}{\unions{\ran{Y_0}}}$ by \cref{indexed_product}.
                            Thus $x$ is a function and $\dom{x} = \dom{Y_0}$ by \cref{funs_elim}.
                            Thus $x\in\rels{\dom{Y_0}}{\unions{\ran{Y_0}}}$ by \cref{funs_subseteq_rels,subseteq}.
                            Thus $x$ is a relation by \cref{rels_is_relation}.
                            Show that $v\in\funs{\dom{Y_1}}{\unions{\ran{Y_1}}}$.
                            \begin{subproof}
                                Show that $v$ is a function.
                                \begin{subproof}
                                    Show that $v$ is a relation.
                                    \begin{subproof}
                                        Show that for all $w$ such that $w\in v$ there exist $a,b$ such that $w = (a,b)$.
                                        \begin{subproof}
                                            Fix $w$.
                                            Assume $w\in v$.
                                            Thus $w\in x$ or $w = (k + 1,y)$ by \cref{union_iff,singleton_iff}.
                                            \begin{byCase}
                                            \caseOf{$w\in x$.}
                                                Take $a,b$ such that $w = (a,b)$ by \cref{relation}.
                                                Thus there exist $a,b$ such that $w = (a,b)$.
                                            \caseOf{$w = (k + 1,y)$.}
                                                Thus there exist $a,b$ such that $w = (a,b)$.
                                            \end{byCase}
                                        \end{subproof}
                                        Thus $v$ is a relation by \cref{relation}.
                                    \end{subproof}
                                    Show that for all $a,b_1,b_2$ such that $(a,b_1)\in v$ and $(a,b_2)\in v$ we have $b_1 = b_2$.
                                    \begin{subproof}
                                        Fix $a$.
                                        Fix $b_1$.
                                        Fix $b_2$.
                                        Assume $(a,b_1)\in v$ and $(a,b_2)\in v$.
                                        \begin{byCase}
                                        \caseOf{$a = k + 1$.}
                                            Thus $(a,b_1)\in x$ or $(a,b_1) = (k + 1,y)$ by \cref{union_iff,singleton_iff}.
                                            Show that $(a,b_1)\notin x$.
                                            \begin{subproof}
                                                Suppose not.
                                                Thus $(a,b_1)\in x$.
                                                Thus $a\in\dom{x}$ by \cref{dom_intro}.
                                                Thus $\dom{x} = \dom{Y_0}$ by \cref{funs_elim}.
                                                Thus $a\in J_1$ by \cref{subseteq}.
                                                Thus $a\in\naturalsPlus$ and $a \leq k$ by \cref{initial_segment_naturalsplus}.
                                                Thus $k + 1 \leq k$.
                                                Thus $k + 1 < k$ or $k + 1 = k$ by \cref{real_leq_split}.
                                                \begin{byCase}
                                                \caseOf{$k + 1 < k$.}
                                                    Contradiction by \cref{real_lt_asym}.
                                                \caseOf{$k + 1 = k$.}
                                                    Thus $k < k$ by \cref{real_lt_subst_left}.
                                                    Contradiction by \cref{real_lt_irreflexive}.
                                                \end{byCase}
                                            \end{subproof}
                                            Thus $(a,b_1) = (k + 1,y)$.
                                            Thus $(a,b_2)\in x$ or $(a,b_2) = (k + 1,y)$ by \cref{union_iff,singleton_iff}.
                                            Show that $(a,b_2)\notin x$.
                                            \begin{subproof}
                                                Suppose not.
                                                Thus $(a,b_2)\in x$.
                                                Thus $a\in\dom{x}$ by \cref{dom_intro}.
                                                Thus $\dom{x} = \dom{Y_0}$ by \cref{funs_elim}.
                                                Thus $a\in J_1$ by \cref{subseteq}.
                                                Thus $a\in\naturalsPlus$ and $a \leq k$ by \cref{initial_segment_naturalsplus}.
                                                Thus $k + 1 \leq k$.
                                                Thus $k + 1 < k$ or $k + 1 = k$ by \cref{real_leq_split}.
                                                \begin{byCase}
                                                \caseOf{$k + 1 < k$.}
                                                    Contradiction by \cref{real_lt_asym}.
                                                \caseOf{$k + 1 = k$.}
                                                    Thus $k < k$ by \cref{real_lt_subst_left}.
                                                    Contradiction by \cref{real_lt_irreflexive}.
                                                \end{byCase}
                                            \end{subproof}
                                            Thus $(a,b_2) = (k + 1,y)$.
                                            Thus $b_1 = y$ and $b_2 = y$ by \cref{pair_eq_iff}.
                                            Thus $b_1 = b_2$.
                                        \caseOf{$a \neq k + 1$.}
                                            Show that $(a,b_1)\in x$.
                                            \begin{subproof}
                                                Suppose not.
                                                Thus $(a,b_1) = (k + 1,y)$ by \cref{union_iff,singleton_iff}.
                                                Thus $a = k + 1$ by \cref{pair_eq_iff}.
                                                Contradiction.
                                            \end{subproof}
                                            Show that $(a,b_2)\in x$.
                                            \begin{subproof}
                                                Suppose not.
                                                Thus $(a,b_2) = (k + 1,y)$ by \cref{union_iff,singleton_iff}.
                                                Thus $a = k + 1$ by \cref{pair_eq_iff}.
                                                Contradiction.
                                            \end{subproof}
                                            Thus $b_1 = b_2$ by \cref{rightunique,relation}.
                                        \end{byCase}
                                    \end{subproof}
                                    Thus $v$ is right-unique by \cref{rightunique}.
                                    Thus $v$ is a function.
                                    Thus $v$ is right-unique and $v$ is a relation.
                                \end{subproof}
                                Show that $\dom{v} = \dom{Y_1}$.
                                \begin{subproof}
                                    Show that for all $i$ such that $i\in\dom{Y_1}$ we have $i\in\dom{v}$.
                                    \begin{subproof}
                                        Fix $i$.
                                        Assume $i\in\dom{Y_1}$.
                                        Thus $i\in J$ by \cref{subseteq}.
                                        \begin{byCase}
                                        \caseOf{$i = k + 1$.}
                                            Thus $(k + 1,y)\in v$ by \cref{union_iff,singleton_iff}.
                                            Thus $i\in\dom{v}$ by \cref{dom_intro}.
                                        \caseOf{$i \neq k + 1$.}
                                            Thus $i\in J_1$ by \cref{real_naturals_leq_succ_elim,initial_segment_naturalsplus}.
                                            Thus $i\in\dom{x}$ by \cref{indexed_product,funs_elim}.
                                            Take $b$ such that $(i,b)\in x$ by \cref{dom_iff}.
                                            Thus $(i,b)\in v$ by \cref{union_iff}.
                                            Thus $i\in\dom{v}$ by \cref{dom_intro}.
                                        \end{byCase}
                                    \end{subproof}
                                    Thus $\dom{Y_1}\subseteq\dom{v}$ by \cref{subseteq}.
                                    Show that for all $i$ such that $i\in\dom{v}$ we have $i\in\dom{Y_1}$.
                                    \begin{subproof}
                                        Fix $i$.
                                        Assume $i\in\dom{v}$.
                                        Take $b$ such that $(i,b)\in v$ by \cref{dom_iff}.
                                        Thus $(i,b)\in x$ or $(i,b) = (k + 1,y)$ by \cref{union_iff,singleton_iff}.
                                        \begin{byCase}
                                        \caseOf{$(i,b) = (k + 1,y)$.}
                                            Thus $i = k + 1$ by \cref{pair_eq_iff}.
                                            Thus $i\in J$.
                                            Thus $i\in\dom{Y_1}$ by \cref{subseteq}.
                                        \caseOf{$(i,b)\in x$.}
                                            Thus $i\in\dom{x}$ by \cref{dom_intro}.
                                            Thus $i\in J_1$ by \cref{indexed_product,funs_elim}.
                                            Thus $i\in J$ by \cref{subseteq,initial_segment_naturalsplus}.
                                            Thus $i\in\dom{Y_1}$ by \cref{subseteq}.
                                        \end{byCase}
                                    \end{subproof}
                                    Thus $\dom{v}\subseteq\dom{Y_1}$ by \cref{subseteq}.
                                    Thus $\dom{v} = \dom{Y_1}$ by \cref{subseteq_antisymmetric}.
                                \end{subproof}
                                Show that for all $i\in\dom{Y_1}$ we have $v(i)\in\unions{\ran{Y_1}}$.
                                \begin{subproof}
                                    Fix $i$.
                                    Assume $i\in\dom{Y_1}$.
                                    Thus $i\in J$ by \cref{subseteq}.
                                    \begin{byCase}
                                    \caseOf{$i = k + 1$.}
                                        Thus $(k + 1,y)\in v$ by \cref{union_iff,singleton_iff}.
                                        Thus $v(k + 1) = y$ by \cref{function_apply_intro}.
                                        Thus $(k + 1,X_1)\in Y_1$ by \cref{function_apply_elim,subseteq}.
                                        Thus $X_1\in\ran{Y_1}$ by \cref{ran_intro}.
                                        Thus $y\in\unions{\ran{Y_1}}$ by \cref{unions_iff}.
                                        Thus $v(i)\in\unions{\ran{Y_1}}$.
                                    \caseOf{$i \neq k + 1$.}
                                        Thus $i\in J_1$ by \cref{real_naturals_leq_succ_elim,initial_segment_naturalsplus}.
                                        Thus $i\in J$ by \cref{subseteq}.
                                        Thus $i\in\dom{Y_1}$ by \cref{subseteq}.
                                        Thus $i\in\dom{x}$ by \cref{indexed_product,funs_elim}.
                                        Thus $x(i)\in\apply{Y_0}{i}$ by \cref{indexed_product}.
                                        Thus $(i,\apply{Y_1}{i})\in Y_0$.
                                        Thus $\apply{Y_0}{i} = \apply{Y_1}{i}$ by \cref{function_apply_intro}.
                                        Thus $x(i)\in\apply{Y_1}{i}$.
                                        Thus $(i,\apply{Y_1}{i})\in Y_1$ by \cref{function_apply_elim}.
                                        Thus $\apply{Y_1}{i}\in\ran{Y_1}$ by \cref{ran_intro}.
                                        Thus $x(i)\in\unions{\ran{Y_1}}$ by \cref{unions_iff}.
                                        Thus $(i,x(i))\in v$ by \cref{function_apply_elim,union_iff}.
                                        Thus $v(i) = x(i)$ by \cref{function_apply_intro}.
                                        Thus $v(i)\in\unions{\ran{Y_1}}$.
                                    \end{byCase}
                                \end{subproof}
                                Thus $v$ is a function to $\unions{\ran{Y_1}}$.
                                Thus $v\in\funs{\dom{Y_1}}{\unions{\ran{Y_1}}}$ by \cref{funs_intro}.
                                Thus $v\in\rels{\dom{Y_1}}{\unions{\ran{Y_1}}}$ and $v$ is right-unique by \cref{funs}.
                                Thus $v$ is a relation by \cref{rels_is_relation}.
                            \end{subproof}
                            Thus $v\in\funs{\dom{Y_1}}{\unions{\ran{Y_1}}}$.
                            Thus $v$ is right-unique by \cref{funs}.
                            Thus $v$ is a relation by \cref{funs,rels_is_relation}.
                            Show that for all $i\in\dom{Y_1}$ we have $v(i)\in\apply{Y_1}{i}$.
                            \begin{subproof}
                                Fix $i$.
                                Assume $i\in\dom{Y_1}$.
                                Thus $i\in J$ by \cref{subseteq}.
                                \begin{byCase}
                                \caseOf{$i = k + 1$.}
                                    Thus $(k + 1,y)\in v$ by \cref{union_iff,singleton_iff}.
                                    Thus $v(k + 1) = y$ by \cref{function_apply_intro}.
                                    Thus $y\in X_1$ by \cref{subseteq}.
                                    Thus $v(i)\in\apply{Y_1}{i}$.
                                \caseOf{$i \neq k + 1$.}
                                    Thus $i\in J_1$ by \cref{real_naturals_leq_succ_elim,initial_segment_naturalsplus}.
                                    Thus $x(i)\in\apply{Y_0}{i}$ by \cref{indexed_product}.
                                    Thus $\apply{Y_0}{i} = \apply{Y_1}{i}$ by \cref{function_apply_intro}.
                                    Thus $x(i)\in\apply{Y_1}{i}$.
                                    Thus $(i,x(i))\in v$ by \cref{function_apply_elim,union_iff}.
                                    Thus $v(i) = x(i)$ by \cref{function_apply_intro}.
                                    Thus $v(i)\in\apply{Y_1}{i}$.
                                \end{byCase}
                            \end{subproof}
                            Thus $v\in\productFamily{Y_1}$ by \cref{indexed_product}.
                        \end{subproof}
                        Thus $g(p)\in\productFamily{Y_1}$ by \cref{function_apply_intro}.
                    \end{subproof}
                    Thus $g$ is a function to $\productFamily{Y_1}$.
                    Thus $g\in\funs{X_0\times X_1}{\productFamily{Y_1}}$ by \cref{funs_intro}.
                \end{subproof}
                Thus $g\in\funs{X_0\times X_1}{\productFamily{Y_1}}$.
                Thus $g$ is right-unique by \cref{funs}.
                Thus $g$ is a relation by \cref{funs,rels_is_relation}.
                Show that $g$ is continuous from $X_0\times X_1$ to $\productFamily{Y_1}$ with respect to $T_P$ and $\productTopologyIndexed{T_1}{Y_1}$.
                \begin{subproof}
                    $J$ is inhabited.
                    Show that for all $i\in J$ we have $\piIndex{Y_1}{i}\circ g$ is continuous from $X_0\times X_1$
                        to $\apply{Y_1}{i}$ with respect to $T_P$ and $\apply{T_1}{i}$.
                    \begin{subproof}
                        Fix $i$.
                        Assume $i\in J$.
                        \begin{byCase}
                        \caseOf{$i = k + 1$.}
                            Let $h = \piTwo{X_0}{X_1}$.
                            $h$ is continuous from $X_0\times X_1$ to $X_1$ with respect to $T_P$ and $T_{1p}$
                                by \cref{projection_two_continuous}.
                            Show that $\piIndex{Y_1}{i}\circ g = h$.
                            \begin{subproof}
                                Show that for all $p\in X_0\times X_1$ we have $\apply{\piIndex{Y_1}{i}\circ g}{p} = \apply{h}{p}$.
                                \begin{subproof}
                                    Fix $p$.
                                    Assume $p\in X_0\times X_1$.
                                    Thus $p\in\dom{g}$ by \cref{funs}.
                                    Thus $(p,g(p))\in g$ by \cref{function_apply_elim}.
                                    Thus $g(p)\in\ran{g}$ by \cref{ran_intro}.
                                    Thus $g$ is a function to $\productFamily{Y_1}$ by \cref{funs_elim}.
                                    Thus $\ran{g}\subseteq\productFamily{Y_1}$ by \cref{function_to_ran}.
                                    Thus $g(p)\in\productFamily{Y_1}$ by \cref{subseteq}.
                                    Thus $(p,\fst{p}\union\{(k + 1,\snd{p})\})\in g$.
                                    Thus $g(p) = \fst{p}\union\{(k + 1,\snd{p})\}$ by \cref{function_apply_intro}.
                                    Thus $(i,\snd{p})\in g(p)$ by \cref{union_iff,singleton_iff}.
                                    Thus $g(p)\in\funs{\dom{Y_1}}{\unions{\ran{Y_1}}}$ by \cref{indexed_product}.
                                    Thus $g(p)$ is a function by \cref{funs_is_function}.
                                    Thus $\apply{g(p)}{i} = \snd{p}$ by \cref{function_apply_intro}.
                                    Thus $(g(p),\apply{g(p)}{i})\in\piIndex{Y_1}{i}$ by \cref{projection_map_indexed}.
                                    Thus $(g(p),\snd{p})\in\piIndex{Y_1}{i}$.
                                    Thus $(p,\snd{p})\in\piIndex{Y_1}{i}\circ g$ by \cref{circ_iff}.
                                    Thus $\piIndex{Y_1}{i}\in\funs{\productFamily{Y_1}}{\apply{Y_1}{i}}$ by \cref{projection_map_indexed_function}.
                                    Thus $\piIndex{Y_1}{i}$ is a function by \cref{funs_is_function}.
                                    Thus $g$ is a function by \cref{funs_is_function}.
                                    Thus $\piIndex{Y_1}{i}\circ g$ is a function by \cref{function_circ}.
                                    Thus $\apply{\piIndex{Y_1}{i}\circ g}{p} = \snd{p}$ by \cref{function_apply_intro}.
                                    Thus $h\in\funs{X_0\times X_1}{X_1}$ by \cref{continuous}.
                                    Thus $h$ is a function by \cref{funs_is_function}.
                                    Take $x,y$ such that $p = (x,y)$ by \cref{times_elem_is_tuple}.
                                    Thus $(x,y)\in X_0\times X_1$.
                                    Thus $x\in X_0$ and $y\in X_1$ by \cref{times_tuple_elim}.
                                    Thus $\snd{p} = y$ by \cref{snd_eq}.
                                    Thus $(p,\snd{p}) = ((x,y),y)$ by \cref{pair_eq_iff}.
                                    Thus $(p,\snd{p})\in h$ by \cref{projection_second}.
                                    Thus $\apply{h}{p} = \snd{p}$ by \cref{function_apply_intro}.
                                    Thus $\apply{\piIndex{Y_1}{i}\circ g}{p} = \apply{h}{p}$.
                                \end{subproof}
                                Show that $\piIndex{Y_1}{i}\circ g = h$.
                                \begin{subproof}
                                    $\piIndex{Y_1}{i}\in\funs{\productFamily{Y_1}}{\apply{Y_1}{i}}$ by \cref{projection_map_indexed_function}.
                                    Thus $\piIndex{Y_1}{i}$ is right-unique and $\piIndex{Y_1}{i}$ is a relation.
                                    $\piIndex{Y_1}{i}\circ g$ is a function by \cref{function_circ}.
                                    $h\in\funs{X_0\times X_1}{X_1}$ by \cref{continuous}.
                                    $h$ is a function by \cref{funs_is_function}.
                                    Show that $\dom{\piIndex{Y_1}{i}\circ g} = X_0\times X_1$.
                                    \begin{subproof}
                                        Show that for all $p$ such that $p\in\dom{\piIndex{Y_1}{i}\circ g}$ we have $p\in X_0\times X_1$.
                                        \begin{subproof}
                                            Fix $p$.
                                            Assume $p\in\dom{\piIndex{Y_1}{i}\circ g}$.
                                            Take $z$ such that $(p,z)\in\piIndex{Y_1}{i}\circ g$ by \cref{dom_iff}.
                                            Take $y$ such that $p\mathrel{g} y$ and $y\mathrel{\piIndex{Y_1}{i}} z$ by \cref{circ_iff}.
                                            Thus $p\in\dom{g}$ by \cref{dom_intro}.
                                            Thus $\dom{g} = X_0\times X_1$ by \cref{funs}.
                                            Thus $p\in X_0\times X_1$ by \cref{subseteq}.
                                        \end{subproof}
                                        Thus $\dom{\piIndex{Y_1}{i}\circ g}\subseteq X_0\times X_1$ by \cref{subseteq}.
                                        Show that for all $p$ such that $p\in X_0\times X_1$ we have $p\in\dom{\piIndex{Y_1}{i}\circ g}$.
                                        \begin{subproof}
                                            Fix $p$.
                                            Assume $p\in X_0\times X_1$.
                                            Thus $p\in\dom{g}$ by \cref{funs}.
                                            Thus $(p,g(p))\in g$ by \cref{function_apply_elim}.
                                            Thus $g(p)\in\ran{g}$ by \cref{ran_intro}.
                                            Thus $g$ is a function to $\productFamily{Y_1}$ by \cref{funs_elim}.
                                            Thus $\ran{g}\subseteq\productFamily{Y_1}$ by \cref{function_to_ran}.
                                            Thus $g(p)\in\productFamily{Y_1}$ by \cref{subseteq}.
                                            Thus $(g(p),\apply{g(p)}{i})\in\piIndex{Y_1}{i}$ by \cref{projection_map_indexed}.
                                            Thus $(p,\apply{g(p)}{i})\in\piIndex{Y_1}{i}\circ g$ by \cref{circ_iff}.
                                            Thus $p\in\dom{\piIndex{Y_1}{i}\circ g}$ by \cref{dom_intro}.
                                        \end{subproof}
                                        Thus $X_0\times X_1\subseteq\dom{\piIndex{Y_1}{i}\circ g}$ by \cref{subseteq}.
                                        Thus $\dom{\piIndex{Y_1}{i}\circ g} = X_0\times X_1$ by \cref{subseteq_antisymmetric}.
                                    \end{subproof}
                                    Thus $\dom{h} = X_0\times X_1$ by \cref{funs}.
                                    Thus $\dom{\piIndex{Y_1}{i}\circ g}\subseteq\dom{h}$ by \cref{subseteq}.
                                    Thus for all $p\in\dom{\piIndex{Y_1}{i}\circ g}$ we have $\apply{\piIndex{Y_1}{i}\circ g}{p} = \apply{h}{p}$ by \cref{subseteq}.
                                    Thus $\piIndex{Y_1}{i}\circ g\subseteq h$ by \cref{fun_subseteq}.
                                    Thus $\dom{h}\subseteq\dom{\piIndex{Y_1}{i}\circ g}$ by \cref{subseteq}.
                                    Thus for all $p\in\dom{h}$ we have $\apply{\piIndex{Y_1}{i}\circ g}{p} = \apply{h}{p}$ by \cref{subseteq}.
                                    Thus $h\subseteq \piIndex{Y_1}{i}\circ g$ by \cref{fun_subseteq}.
                                    Thus $\piIndex{Y_1}{i}\circ g = h$ by \cref{subseteq_antisymmetric}.
                                \end{subproof}
                            \end{subproof}
                            Thus $\piIndex{Y_1}{i}\circ g$ is continuous from $X_0\times X_1$ to $\apply{Y_1}{i}$ with respect to $T_P$ and $\apply{T_1}{i}$
                                by \cref{continuous_eq_subst}.
                        \caseOf{$i \neq k + 1$.}
                            Thus $i\in J_1$ by \cref{real_naturals_leq_succ_elim,initial_segment_naturalsplus}.
                            Let $h_0 = \piOne{X_0}{X_1}$.
                            $h_0$ is continuous from $X_0\times X_1$ to $X_0$ with respect to $T_P$ and $T_X$
                                by \cref{projection_one_continuous}.
                            $\piIndex{Y_0}{i}$ is continuous from $X_0$ to $\apply{Y_0}{i}$ with respect to $T_X$ and $\apply{T_0}{i}$
                                by \cref{projection_map_indexed_continuous,subseteq,subseteq_antisymmetric}.
                            Let $h = \piIndex{Y_0}{i}\circ h_0$.
                            $h$ is continuous from $X_0\times X_1$ to $\apply{Y_0}{i}$ with respect to $T_P$ and $\apply{T_0}{i}$
                                by \cref{continuous_composite}.
                            Show that $\piIndex{Y_1}{i}\circ g = h$.
                            \begin{subproof}
                                Show that for all $p\in X_0\times X_1$ we have $\apply{\piIndex{Y_1}{i}\circ g}{p} = \apply{h}{p}$.
                                \begin{subproof}
                                    Fix $p$.
                                    Assume $p\in X_0\times X_1$.
                                    Thus $(p,\fst{p}\union\{(k + 1,\snd{p})\})\in g$.
                                    Thus $g(p) = \fst{p}\union\{(k + 1,\snd{p})\}$ by \cref{function_apply_intro}.
                                    Let $x = \fst{p}$.
                                    Take $a,b$ such that $p = (a,b)$ and $a\in X_0$ and $b\in X_1$ by \cref{times_elem_is_tuple}.
                                    Thus $\fst{p} = a$ by \cref{fst_eq}.
                                    Thus $x = a$.
                                    Thus $x\in X_0$.
                                    Thus $(p,x) = ((a,b),a)$ by \cref{pair_eq_iff}.
                                    Thus $(p,x)\in h_0$ by \cref{projection_first}.
                                    Thus $(x,\apply{x}{i})\in\piIndex{Y_0}{i}$ by \cref{projection_map_indexed}.
                                    Thus $(p,\apply{x}{i})\in\piIndex{Y_0}{i}\circ h_0$ by \cref{circ_iff}.
                                    Thus $h\in\funs{X_0\times X_1}{\apply{Y_0}{i}}$ by \cref{continuous}.
                                    Thus $h$ is right-unique and $h$ is a relation by \cref{funs_is_function}.
                                    Thus $\apply{h}{p} = \apply{x}{i}$ by \cref{function_apply_intro}.
                                    Thus $p\in\dom{g}$ by \cref{funs}.
                                    Thus $(p,g(p))\in g$ by \cref{function_apply_elim}.
                                    Thus $g(p)\in\ran{g}$ by \cref{ran_intro}.
                                    Thus $g$ is a function to $\productFamily{Y_1}$ by \cref{funs_elim}.
                                    Thus $\ran{g}\subseteq\productFamily{Y_1}$ by \cref{function_to_ran}.
                                    Thus $g(p)\in\productFamily{Y_1}$ by \cref{subseteq}.
                                    Thus $g(p)\in\funs{\dom{Y_1}}{\unions{\ran{Y_1}}}$ by \cref{indexed_product}.
                                    Thus $g(p)$ is a function by \cref{funs_is_function}.
                                    Thus $x\in\funs{\dom{Y_0}}{\unions{\ran{Y_0}}}$ by \cref{indexed_product}.
                                    Thus $x$ is right-unique and $x$ is a relation by \cref{funs_is_function}.
                                    Thus $i\in\dom{x}$ by \cref{funs_elim,subseteq}.
                                    Thus $(i,\apply{x}{i})\in x$ by \cref{function_apply_elim}.
                                    Thus $(i,\apply{x}{i})\in g(p)$ by \cref{union_iff}.
                                    Thus $\apply{g(p)}{i} = \apply{x}{i}$ by \cref{function_apply_intro}.
                                    Thus $(g(p),\apply{x}{i})\in\piIndex{Y_1}{i}$ by \cref{projection_map_indexed}.
                                    Thus $(p,\apply{x}{i})\in\piIndex{Y_1}{i}\circ g$ by \cref{circ_iff}.
                                    Thus $\piIndex{Y_1}{i}\in\funs{\productFamily{Y_1}}{\apply{Y_1}{i}}$ by \cref{projection_map_indexed_function}.
                                    Thus $\piIndex{Y_1}{i}$ is right-unique and $\piIndex{Y_1}{i}$ is a relation by \cref{funs_is_function}.
                                    Thus $g$ is right-unique by \cref{funs}.
                                    Thus $\piIndex{Y_1}{i}\circ g$ is right-unique and $\piIndex{Y_1}{i}\circ g$ is a relation by \cref{function_circ}.
                                    Thus $\apply{\piIndex{Y_1}{i}\circ g}{p} = \apply{x}{i}$ by \cref{function_apply_intro}.
                                    Thus $\apply{\piIndex{Y_1}{i}\circ g}{p} = \apply{h}{p}$.
                                \end{subproof}
                                Show that $\piIndex{Y_1}{i}\circ g = h$.
                                \begin{subproof}
                                    $\piIndex{Y_1}{i}\in\funs{\productFamily{Y_1}}{\apply{Y_1}{i}}$ by \cref{projection_map_indexed_function}.
                                    Thus $\piIndex{Y_1}{i}$ is right-unique and $\piIndex{Y_1}{i}$ is a relation by \cref{funs_is_function}.
                                    Thus $g$ is right-unique by \cref{funs}.
                                    Thus $\piIndex{Y_1}{i}\circ g$ is a function by \cref{function_circ}.
                                    $h\in\funs{X_0\times X_1}{\apply{Y_0}{i}}$ by \cref{continuous}.
                                    Thus $h$ is a function by \cref{funs_is_function}.
                                    Show that $\dom{\piIndex{Y_1}{i}\circ g} = X_0\times X_1$.
                                    \begin{subproof}
                                        Show that for all $p$ such that $p\in\dom{\piIndex{Y_1}{i}\circ g}$ we have $p\in X_0\times X_1$.
                                        \begin{subproof}
                                            Fix $p$.
                                            Assume $p\in\dom{\piIndex{Y_1}{i}\circ g}$.
                                            Take $z$ such that $(p,z)\in\piIndex{Y_1}{i}\circ g$ by \cref{dom_iff}.
                                            Take $y$ such that $p\mathrel{g} y$ and $y\mathrel{\piIndex{Y_1}{i}} z$ by \cref{circ_iff}.
                                            Thus $p\in\dom{g}$ by \cref{dom_intro}.
                                            Thus $\dom{g} = X_0\times X_1$ by \cref{funs}.
                                            Thus $p\in X_0\times X_1$ by \cref{subseteq}.
                                        \end{subproof}
                                        Thus $\dom{\piIndex{Y_1}{i}\circ g}\subseteq X_0\times X_1$ by \cref{subseteq}.
                                        Show that for all $p$ such that $p\in X_0\times X_1$ we have $p\in\dom{\piIndex{Y_1}{i}\circ g}$.
                                        \begin{subproof}
                                            Fix $p$.
                                            Assume $p\in X_0\times X_1$.
                                            Thus $p\in\dom{g}$ by \cref{funs}.
                                            Thus $(p,g(p))\in g$ by \cref{function_apply_elim}.
                                            Thus $g(p)\in\ran{g}$ by \cref{ran_intro}.
                                            Thus $g$ is a function to $\productFamily{Y_1}$ by \cref{funs_elim}.
                                            Thus $\ran{g}\subseteq\productFamily{Y_1}$ by \cref{function_to_ran}.
                                            Thus $g(p)\in\productFamily{Y_1}$ by \cref{subseteq}.
                                            Thus $(g(p),\apply{g(p)}{i})\in\piIndex{Y_1}{i}$ by \cref{projection_map_indexed}.
                                            Thus $(p,\apply{g(p)}{i})\in\piIndex{Y_1}{i}\circ g$ by \cref{circ_iff}.
                                            Thus $p\in\dom{\piIndex{Y_1}{i}\circ g}$ by \cref{dom_intro}.
                                        \end{subproof}
                                        Thus $X_0\times X_1\subseteq\dom{\piIndex{Y_1}{i}\circ g}$ by \cref{subseteq}.
                                        Thus $\dom{\piIndex{Y_1}{i}\circ g} = X_0\times X_1$ by \cref{subseteq_antisymmetric}.
                                    \end{subproof}
                                    Thus $\dom{h} = X_0\times X_1$ by \cref{funs}.
                                    Thus $\dom{\piIndex{Y_1}{i}\circ g}\subseteq\dom{h}$ by \cref{subseteq}.
                                    Thus for all $p\in\dom{\piIndex{Y_1}{i}\circ g}$ we have $\apply{\piIndex{Y_1}{i}\circ g}{p} = \apply{h}{p}$ by \cref{subseteq}.
                                    Thus $\piIndex{Y_1}{i}\circ g\subseteq h$ by \cref{fun_subseteq}.
                                    Thus $\dom{h}\subseteq\dom{\piIndex{Y_1}{i}\circ g}$ by \cref{subseteq}.
                                    Thus for all $p\in\dom{h}$ we have $\apply{\piIndex{Y_1}{i}\circ g}{p} = \apply{h}{p}$ by \cref{subseteq}.
                                    Thus $h\subseteq \piIndex{Y_1}{i}\circ g$ by \cref{fun_subseteq}.
                                    Thus $\piIndex{Y_1}{i}\circ g = h$ by \cref{subseteq_antisymmetric}.
                                \end{subproof}
                            \end{subproof}
                            Thus $\apply{Y_0}{i} = \apply{Y_1}{i}$ by \cref{function_apply_intro}.
                            Thus $\apply{T_0}{i} = \apply{T_1}{i}$ by \cref{function_apply_intro}.
                            Thus $h$ is continuous from $X_0\times X_1$ to $\apply{Y_1}{i}$ with respect to $T_P$ and $\apply{T_1}{i}$ by \cref{continuous_eq_codomain}.
                            Thus $h\in\funs{X_0\times X_1}{\apply{Y_1}{i}}$ by \cref{continuous}.
                            Thus $h$ is right-unique and $h$ is a relation by \cref{funs_is_function}.
                            Thus $\piIndex{Y_1}{i}\in\funs{\productFamily{Y_1}}{\apply{Y_1}{i}}$ by \cref{projection_map_indexed_function}.
                            Thus $\piIndex{Y_1}{i}$ is right-unique and $\piIndex{Y_1}{i}$ is a relation by \cref{funs_is_function}.
                            Thus $g$ is right-unique by \cref{funs}.
                            Thus $\piIndex{Y_1}{i}\circ g$ is right-unique and $\piIndex{Y_1}{i}\circ g$ is a relation by \cref{function_circ}.
                            Thus $\piIndex{Y_1}{i}\circ g$ is continuous from $X_0\times X_1$ to $\apply{Y_1}{i}$ with respect to $T_P$ and $\apply{T_1}{i}$
                                by \cref{continuous_eq_subst}.
                        \end{byCase}
                    \end{subproof}
                    $Y_1$ is right-unique and $Y_1$ is a relation.
                    $T_1$ is right-unique and $T_1$ is a relation.
                    Thus $\dom{T_1} = \dom{Y_1}$.
                    $J$ is inhabited.
                    $T_P$ is a topology on $X_0\times X_1$.
                    Thus $g$ is continuous from $X_0\times X_1$ to $\productFamily{Y_1}$ with respect to $T_P$ and $\productTopologyIndexed{T_1}{Y_1}$
                        by \cref{continuous_map_into_indexed_product}.
                \end{subproof}
                Show that $g\in\surj{X_0\times X_1}{\productFamily{Y_1}}$.
                \begin{subproof}
                    Show that for all $z\in\productFamily{Y_1}$ there exists $p\in X_0\times X_1$ such that $g(p) = z$.
                    \begin{subproof}
                        Fix $z$.
                        Assume $z\in\productFamily{Y_1}$.
                        Thus $z\in\funs{\dom{Y_1}}{\unions{\ran{Y_1}}}$ by \cref{indexed_product}.
                        Thus $z$ is a function and $\dom{z} = \dom{Y_1}$ by \cref{funs_elim}.
                        Let $x = \{(i,z(i)) \mid i\in J_1\}$.
                        Show that $x\in X_0$.
                        \begin{subproof}
                            Show that $x\in\productFamily{Y_0}$.
                            \begin{subproof}
                                Show that $x\in\funs{\dom{Y_0}}{\unions{\ran{Y_0}}}$.
                                \begin{subproof}
                                    Show that $x$ is a function.
                                    \begin{subproof}
                                        Show that $x$ is a relation.
                                        \begin{subproof}
                                            Show that for all $w$ such that $w\in x$ there exist $a,b$ such that $w = (a,b)$.
                                            \begin{subproof}
                                                Fix $w$.
                                                Assume $w\in x$.
                                                Take $i\in J_1$ such that $w = (i,z(i))$.
                                                Thus there exist $a,b$ such that $w = (a,b)$.
                                            \end{subproof}
                                            Thus $x$ is a relation by \cref{relation}.
                                        \end{subproof}
                                        Show that for all $a,b_1,b_2$ such that $(a,b_1)\in x$ and $(a,b_2)\in x$ we have $b_1 = b_2$.
                                        \begin{subproof}
                                            Fix $a$.
                                            Fix $b_1$.
                                            Fix $b_2$.
                                            Assume $(a,b_1)\in x$ and $(a,b_2)\in x$.
                                            Take $i_1\in J_1$ such that $(a,b_1) = (i_1,z(i_1))$.
                                            Take $i_2\in J_1$ such that $(a,b_2) = (i_2,z(i_2))$.
                                            Thus $a = i_1$ and $b_1 = z(i_1)$ by \cref{pair_eq_iff}.
                                            Thus $a = i_2$ and $b_2 = z(i_2)$ by \cref{pair_eq_iff}.
                                            Thus $b_1 = b_2$.
                                        \end{subproof}
                                        Thus $x$ is right-unique by \cref{rightunique}.
                                        Thus $x$ is a function.
                                    \end{subproof}
                                    Show that $\dom{x} = \dom{Y_0}$.
                                    \begin{subproof}
                                        Show that for all $i$ such that $i\in\dom{Y_0}$ we have $i\in\dom{x}$.
                                        \begin{subproof}
                                            Fix $i$.
                                            Assume $i\in\dom{Y_0}$.
                                            Thus $i\in J_1$ by \cref{subseteq}.
                                            Thus $(i,z(i))\in x$.
                                            Thus $i\in\dom{x}$ by \cref{dom_intro}.
                                        \end{subproof}
                                        Thus $\dom{Y_0}\subseteq\dom{x}$ by \cref{subseteq}.
                                        Show that for all $i$ such that $i\in\dom{x}$ we have $i\in\dom{Y_0}$.
                                        \begin{subproof}
                                            Fix $i$.
                                            Assume $i\in\dom{x}$.
                                            Take $v$ such that $(i,v)\in x$ by \cref{dom_iff}.
                                            Take $j\in J_1$ such that $(i,v) = (j,z(j))$.
                                            Thus $i = j$ by \cref{pair_eq_iff}.
                                            Thus $i\in J_1$.
                                            Thus $i\in\dom{Y_0}$ by \cref{subseteq}.
                                        \end{subproof}
                                        Thus $\dom{x}\subseteq\dom{Y_0}$ by \cref{subseteq}.
                                        Thus $\dom{x} = \dom{Y_0}$ by \cref{subseteq_antisymmetric}.
                                    \end{subproof}
                                    Show that for all $i\in\dom{Y_0}$ we have $x(i)\in\unions{\ran{Y_0}}$.
                                    \begin{subproof}
                                        Fix $i$.
                                        Assume $i\in\dom{Y_0}$.
                                        Thus $i\in J_1$ by \cref{subseteq}.
                                        Thus $i\in J$ by \cref{subseteq}.
                                        Thus $i\in\dom{Y_1}$ by \cref{subseteq}.
                                        Thus $(i,z(i))\in x$.
                                        Thus $x$ is a function.
                                        Thus $x(i) = z(i)$ by \cref{function_apply_intro}.
                                        Thus $(i,\apply{Y_1}{i})\in Y_0$.
                                        Thus $\apply{Y_0}{i} = \apply{Y_1}{i}$ by \cref{function_apply_intro}.
                                        Thus $z(i)\in\apply{Y_1}{i}$ by \cref{indexed_product}.
                                        Thus $z(i)\in\apply{Y_0}{i}$.
                                        Thus $\apply{Y_0}{i}\in\ran{Y_0}$ by \cref{ran_intro}.
                                        Thus $x(i)\in\unions{\ran{Y_0}}$ by \cref{unions_iff}.
                                    \end{subproof}
                                    Thus $x$ is a function to $\unions{\ran{Y_0}}$.
                                    Thus $x\in\funs{\dom{Y_0}}{\unions{\ran{Y_0}}}$ by \cref{funs_intro}.
                                    Thus $x$ is right-unique and $x$ is a relation by \cref{funs_is_function}.
                                \end{subproof}
                                Show that for all $i\in\dom{Y_0}$ we have $x(i)\in\apply{Y_0}{i}$.
                                \begin{subproof}
                                Fix $i$.
                                Assume $i\in\dom{Y_0}$.
                                Thus $i\in J_1$ by \cref{subseteq}.
                                Thus $i\in J$ by \cref{subseteq}.
                                Thus $i\in\dom{Y_1}$ by \cref{subseteq}.
                                Thus $(i,z(i))\in x$.
                                Thus $x$ is a function.
                                Thus $x(i) = z(i)$ by \cref{function_apply_intro}.
                                Thus $(i,\apply{Y_1}{i})\in Y_0$.
                                Thus $\apply{Y_0}{i} = \apply{Y_1}{i}$ by \cref{function_apply_intro}.
                                Thus $z(i)\in\apply{Y_1}{i}$ by \cref{indexed_product}.
                                Thus $x(i)\in\apply{Y_0}{i}$.
                            \end{subproof}
                                Thus $x\in\productFamily{Y_0}$ by \cref{indexed_product}.
                            \end{subproof}
                            Thus $x\in X_0$.
                        \end{subproof}
                        Let $y = z(k + 1)$.
                        Thus $y\in X_1$ by \cref{indexed_product,subseteq}.
                        Let $p = (x,y)$.
                        Thus $p\in X_0\times X_1$ by \cref{times_tuple_intro}.
                        Show that $g(p) = z$.
                        \begin{subproof}
                            Show that for all $i\in\dom{z}$ we have $\apply{g(p)}{i} = z(i)$.
                            \begin{subproof}
                                Fix $i$.
                                Assume $i\in\dom{z}$.
                                Thus $i\in J$ by \cref{subseteq}.
                                Thus $p\in\dom{g}$ by \cref{funs}.
                                Thus $(p,g(p))\in g$ by \cref{function_apply_elim}.
                                Thus $g(p)\in\ran{g}$ by \cref{ran_intro}.
                                Thus $g$ is a function to $\productFamily{Y_1}$ by \cref{funs_elim}.
                                Thus $\ran{g}\subseteq\productFamily{Y_1}$ by \cref{function_to_ran}.
                                Thus $g(p)\in\productFamily{Y_1}$ by \cref{subseteq}.
                                Thus $g(p)\in\funs{\dom{Y_1}}{\unions{\ran{Y_1}}}$ by \cref{indexed_product}.
                                Thus $g(p)$ is a function by \cref{funs_is_function}.
                                Thus $(p,\fst{p}\union\{(k + 1,\snd{p})\})\in g$.
                                Thus $g(p) = \fst{p}\union\{(k + 1,\snd{p})\}$ by \cref{function_apply_intro}.
                                Thus $\fst{p} = x$ by \cref{fst_eq}.
                                Thus $\snd{p} = y$ by \cref{snd_eq}.
                                Thus $g(p) = x\union\{(k + 1,y)\}$.
                                \begin{byCase}
                                \caseOf{$i = k + 1$.}
                                    Thus $(k + 1,y)\in\{(k + 1,y)\}$ by \cref{singleton_intro}.
                                    Thus $(k + 1,y)\in g(p)$ by \cref{union_iff}.
                                    Thus $\apply{g(p)}{k + 1} = y$ by \cref{function_apply_intro}.
                                    Thus $\apply{g(p)}{i} = z(i)$.
                                \caseOf{$i \neq k + 1$.}
                                    Thus $i\in J_1$ by \cref{real_naturals_leq_succ_elim,initial_segment_naturalsplus}.
                                    Thus $(i,z(i))\in x$.
                                    Thus $(i,z(i))\in g(p)$ by \cref{union_iff}.
                                    Thus $\apply{g(p)}{i} = z(i)$ by \cref{function_apply_intro}.
                                \end{byCase}
                            \end{subproof}
                            Show that $g(p) = z$.
                            \begin{subproof}
                                Thus $z$ is a function and $\dom{z} = \dom{Y_1}$ by \cref{funs_elim}.
                                Thus $g$ is a function to $\productFamily{Y_1}$ by \cref{funs_elim}.
                                Thus $p\in\dom{g}$ by \cref{funs}.
                                Thus $(p,g(p))\in g$ by \cref{function_apply_elim}.
                                Thus $g(p)\in\ran{g}$ by \cref{ran_intro}.
                                Thus $\ran{g}\subseteq\productFamily{Y_1}$ by \cref{function_to_ran}.
                                Thus $g(p)\in\productFamily{Y_1}$ by \cref{subseteq}.
                                Thus $g(p)\in\funs{\dom{Y_1}}{\unions{\ran{Y_1}}}$ by \cref{indexed_product}.
                                Thus $g(p)$ is a function and $\dom{g(p)} = \dom{Y_1}$ by \cref{funs_elim}.
                                Show that $g(p)\subseteq z$.
                                \begin{subproof}
                                    Show that for all $i$ such that $i\in\dom{g(p)}$ we have $i\in\dom{z}$.
                                    \begin{subproof}
                                        Fix $i$.
                                        Assume $i\in\dom{g(p)}$.
                                        Thus $i\in\dom{Y_1}$ by \cref{subseteq}.
                                        Thus $i\in\dom{z}$ by \cref{subseteq}.
                                    \end{subproof}
                                    Thus $\dom{g(p)}\subseteq\dom{z}$ by \cref{subseteq}.
                                    Show that for all $i\in\dom{g(p)}$ we have $\apply{g(p)}{i} = z(i)$.
                                    \begin{subproof}
                                        Fix $i$.
                                        Assume $i\in\dom{g(p)}$.
                                        Thus $i\in\dom{z}$ by \cref{subseteq}.
                                        Thus $\apply{g(p)}{i} = z(i)$.
                                    \end{subproof}
                                    Thus $g(p)\subseteq z$ by \cref{fun_subseteq}.
                                \end{subproof}
                                Show that $z\subseteq g(p)$.
                                \begin{subproof}
                                    Show that for all $i$ such that $i\in\dom{z}$ we have $i\in\dom{g(p)}$.
                                    \begin{subproof}
                                        Fix $i$.
                                        Assume $i\in\dom{z}$.
                                        Thus $i\in\dom{Y_1}$ by \cref{subseteq}.
                                        Thus $i\in\dom{g(p)}$ by \cref{subseteq}.
                                    \end{subproof}
                                    Thus $\dom{z}\subseteq\dom{g(p)}$ by \cref{subseteq}.
                                    Thus for all $i\in\dom{z}$ we have $z(i) = \apply{g(p)}{i}$.
                                    Thus $z\subseteq g(p)$ by \cref{fun_subseteq}.
                                \end{subproof}
                                Thus $g(p) = z$ by \cref{subseteq_antisymmetric}.
                            \end{subproof}
                        \end{subproof}
                        Thus there exists $p\in X_0\times X_1$ such that $g(p) = z$.
                    \end{subproof}
                    Thus $g\in\surj{X_0\times X_1}{\productFamily{Y_1}}$ by \cref{surj}.
                \end{subproof}
                Thus $\dom{g} = X_0\times X_1$ by \cref{surj,funs_elim}.
                Thus $\img{g}{X_0\times X_1} = \ran{g}$ by \cref{img_dom}.
                Thus $\ran{g} = \productFamily{Y_1}$ by \cref{funs_surj_iff}.
                Thus $\img{g}{X_0\times X_1} = \productFamily{Y_1}$.
                Let $A = \img{g}{X_0\times X_1}$.
                Let $T_Q = \productTopologyIndexed{T_1}{Y_1}$.
                Let $T_A = \subspaceTopology{T_Q}{A}$.
                Thus $A = \productFamily{Y_1}$.
                Thus $T_A = \subspaceTopology{T_Q}{\productFamily{Y_1}}$.
                Show that $T_A = T_Q$.
                \begin{subproof}
                    $T_Q$ is a topology on $\productFamily{Y_1}$
                        by \cref{product_subbasis_finite_intersections_basis,basis_for_topology,basis_topology_is_topology}.
                    Show that $T_A\subseteq T_Q$.
                    \begin{subproof}
                        Show that for all $U$ such that $U\in T_A$ we have $U\in T_Q$.
                        \begin{subproof}
                            Fix $U$.
                            Assume $U\in T_A$.
                            Take $V\in T_Q$ such that $U = \productFamily{Y_1}\inter V$ by \cref{subspace_topology}.
                            Thus $T_Q\subseteq\pow{\productFamily{Y_1}}$ by \cref{topology_on}.
                            Thus $V\in\pow{\productFamily{Y_1}}$ by \cref{subseteq}.
                            Thus $V\subseteq\productFamily{Y_1}$ by \cref{powerset_elim}.
                            Thus $U = V$ by \cref{inter_absorb_supseteq_left}.
                            Thus $U\in T_Q$.
                        \end{subproof}
                        Thus $T_A\subseteq T_Q$ by \cref{subseteq}.
                    \end{subproof}
                    Show that $T_Q\subseteq T_A$.
                    \begin{subproof}
                        Show that for all $U$ such that $U\in T_Q$ we have $U\in T_A$.
                        \begin{subproof}
                            Fix $U$.
                            Assume $U\in T_Q$.
                            Thus $T_Q\subseteq\pow{\productFamily{Y_1}}$ by \cref{topology_on}.
                            Thus $U\in\pow{\productFamily{Y_1}}$ by \cref{subseteq}.
                            Thus $U\subseteq\productFamily{Y_1}$ by \cref{powerset_elim}.
                            Thus $U = \productFamily{Y_1}\inter U$ by \cref{inter_absorb_supseteq_left}.
                            Thus $U\in T_A$ by \cref{subspace_topology}.
                        \end{subproof}
                        Thus $T_Q\subseteq T_A$ by \cref{subseteq}.
                    \end{subproof}
                    Thus $T_A = T_Q$ by \cref{subseteq_antisymmetric}.
                \end{subproof}
                Thus $A$ is compact with respect to $T_A$ by \cref{continuous_image_compact}.
                Thus $\productFamily{Y_1}$ is compact with respect to $T_A$.
                Thus $\productFamily{Y_1}$ is compact with respect to $T_Q$.
            \end{subproof}
        \end{subproof}
        Thus for all $Y_1,T_1$ such that $Y_1$ is a function and $T_1$ is a function and $\dom{Y_1} = \initialSegment{k + 1}$ and $\dom{T_1} = \initialSegment{k + 1}$,
            if for all $i\in\initialSegment{k + 1}$ we have $\apply{T_1}{i}$ is a topology on $\apply{Y_1}{i}$ then
            if for all $i\in\initialSegment{k + 1}$ we have $\apply{Y_1}{i}$ is compact with respect to $\apply{T_1}{i}$ then
            $\productFamily{Y_1}$ is compact with respect to $\productTopologyIndexed{T_1}{Y_1}$.
        Thus $\finiteIndexedProductCompactProp{k + 1}$ by \cref{finite_indexed_product_compact_property}.
        \end{subproof}
        Thus $k + 1\in\naturalsPlus$ by \cref{real_naturals_closed_succ}.
        Thus $k + 1\in\naturalsPlus$ and $\finiteIndexedProductCompactProp{k + 1}$.
    \end{subproof}
    Thus $n\in A$ by \cref{real_naturals_induction}.
    Thus $\finiteIndexedProductCompactProp{n}$.
    Thus for all $Y_1,T_1$ such that $Y_1$ is a function and $T_1$ is a function and $\dom{Y_1} = \initialSegment{n}$ and $\dom{T_1} = \initialSegment{n}$,
        if for all $i\in\initialSegment{n}$ we have $\apply{T_1}{i}$ is a topology on $\apply{Y_1}{i}$ then
        if for all $i\in\initialSegment{n}$ we have $\apply{Y_1}{i}$ is compact with respect to $\apply{T_1}{i}$ then
        $\productFamily{Y_1}$ is compact with respect to $\productTopologyIndexed{T_1}{Y_1}$ by \cref{finite_indexed_product_compact_property}.
    Thus $\productFamily{Y}$ is compact with respect to $\productTopologyIndexed{T_Y}{Y}$.
\end{proof}

% from §27 helper lemma 9: square and euclidean metric topologies agree
\begin{lemma}\label{square_metric_topology_equals_euclidean}
    Suppose $n\in\naturalsPlus$.
    Let $J = \initialSegment{n}$.
    Let $X = \realsPower{J}$.
    Let $T_E = \metricTopology{\euclideanMetric{n}}{X}$.
    Let $T_S = \metricTopology{\squareMetric{n}}{X}$.
    Then $T_E = T_S$.
\end{lemma}
\begin{proof}
    Let $X_0 = X$.
    Let $d = \euclideanMetric{n}$.
    Let $d_1 = \squareMetric{n}$.
    $d$ is a metric on $X_0$ by \cref{euclidean_metric_is_metric}.
    $d_1$ is a metric on $X_0$ by \cref{square_metric_is_metric}.
    Show that $T_E = T_S$.
    \begin{subproof}
        Show that $T_E$ is finer than $T_S$ on $X_0$.
        \begin{subproof}
            Show that for all $x\in X_0$ we have
                for all $\epsilon\in\reals$ such that $0 < \epsilon$ there exists $\delta\in\reals$
                    such that $0 < \delta$ and
                    $\metricBall{d}{X_0}{x}{\delta} \subseteq \metricBall{d_1}{X_0}{x}{\epsilon}$.
            \begin{subproof}
                Fix $x$.
                Assume $x\in X_0$.
                Fix $\epsilon\in\reals$.
                Assume $0 < \epsilon$.
                Let $\delta = \epsilon$.
                $0 < \delta$.
                Show that for all $y$ such that $y\in\metricBall{d}{X_0}{x}{\delta}$ we have $y\in\metricBall{d_1}{X_0}{x}{\epsilon}$.
                \begin{subproof}
                    Fix $y$.
                    Assume $y\in\metricBall{d}{X_0}{x}{\delta}$.
                    Thus $y\in X_0$ and $d((x,y)) < \delta$ by \cref{metric_ball}.
                    Take $r\in\reals$ such that $((x,y),r)\in d$ by \cref{euclidean_metric_value_exists}.
                    Take $m\in\reals$ such that $((x,y),m)\in d_1$ by \cref{square_metric_value_exists}.
                    Thus $m \leq r$ by \cref{euclidean_square_metric_bounds}.
                    $d$ is a function by \cref{euclidean_metric_function}.
                    $d_1$ is a function by \cref{square_metric_function}.
                    Thus $d((x,y)) = r$ by \cref{function_apply_intro}.
                    Thus $r < \epsilon$ by \cref{real_lt_subst_left}.
                    Show that $m < \epsilon$.
                    \begin{subproof}
                        $m < r$ or $m = r$ by \cref{real_leq_split}.
                        \begin{byCase}
                        \caseOf{$m < r$.} Thus $m < \epsilon$ by \cref{real_lt_trans}.
                        \caseOf{$m = r$.} Thus $m < \epsilon$ by \cref{real_lt_subst_left}.
                        \end{byCase}
                    \end{subproof}
                    Thus $d_1((x,y)) = m$ by \cref{function_apply_intro}.
                    Thus $y\in\metricBall{d_1}{X_0}{x}{\epsilon}$ by \cref{metric_ball}.
                \end{subproof}
                Thus $\metricBall{d}{X_0}{x}{\delta} \subseteq \metricBall{d_1}{X_0}{x}{\epsilon}$ by \cref{subseteq}.
            \end{subproof}
            Thus $T_E$ is finer than $T_S$ on $X_0$ by \cref{metric_topology_finer_iff}.
        \end{subproof}
        Show that $T_S$ is finer than $T_E$ on $X_0$.
        \begin{subproof}
            Show that for all $x\in X_0$ we have
                for all $\epsilon\in\reals$ such that $0 < \epsilon$ there exists $\delta\in\reals$
                    such that $0 < \delta$ and
                    $\metricBall{d_1}{X_0}{x}{\delta} \subseteq \metricBall{d}{X_0}{x}{\epsilon}$.
            \begin{subproof}
                Fix $x$.
                Assume $x\in X_0$.
                Fix $\epsilon\in\reals$.
                Assume $0 < \epsilon$.
                $n\in\reals$ by \cref{real_naturals_subset_reals,subseteq}.
                Show that $0 \leq n$.
                \begin{subproof}
                    $0 < 1$ by \cref{real_zero_lt_one}.
                    $0\in\reals$ by \cref{real_add_ident}.
                    $1\in\naturalsPlus$ by \cref{real_one_in_naturals}.
                    $1\in\reals$ by \cref{real_naturals_subset_reals,real_one_in_naturals,subseteq}.
                    $1 \leq n$ by \cref{real_one_leq_naturalsplus}.
                    Thus $1 < n$ or $1 = n$ by \cref{real_leq_split}.
                    \begin{byCase}
                    \caseOf{$1 < n$.}
                        Thus $0 < n$ by \cref{real_lt_trans}.
                        Thus $0 \leq n$.
                    \caseOf{$1 = n$.}
                        Thus $0 < n$ by \cref{real_lt_subst_right}.
                        Thus $0 \leq n$.
                    \end{byCase}
                \end{subproof}
                Thus $\sqrt{n}\in\reals$ and $\sqrt{n} \geq 0$ and $\sqrt{n} \rmul \sqrt{n} = n$ by \cref{positive_square_root}.
                Show that $\sqrt{n}$ is positive.
                \begin{subproof}
                    Show that $\sqrt{n} \neq 0$.
                    \begin{subproof}
                        Suppose not.
                        Thus $\sqrt{n} = 0$.
                        Thus $n = 0$ by \cref{positive_square_root,real_mul_zero}.
                        $0 < 1$ by \cref{real_zero_lt_one}.
                        $1 \leq n$ by \cref{real_one_leq_naturalsplus}.
                        Show that $0 < n$.
                        \begin{subproof}
                            $0\in\reals$ by \cref{real_add_ident}.
                            $1\in\reals$ by \cref{real_naturals_subset_reals,real_one_in_naturals,subseteq}.
                            $n\in\reals$ by \cref{real_naturals_subset_reals,subseteq}.
                            Thus $1 < n$ or $1 = n$ by \cref{real_leq_split}.
                            \begin{byCase}
                            \caseOf{$1 < n$.}
                                Thus $0 < n$ by \cref{real_lt_trans}.
                            \caseOf{$1 = n$.}
                                Thus $0 < n$ by \cref{real_lt_subst_right}.
                            \end{byCase}
                        \end{subproof}
                        Thus $0 < 0$ by \cref{real_lt_subst_right}.
                        Contradiction by \cref{real_lt_irreflexive}.
                    \end{subproof}
                    Thus $\sqrt{n} > 0$.
                \end{subproof}
                Let $\delta = \epsilon \rmul \inv{\sqrt{n}}$.
                Show that $0 < \delta$.
                \begin{subproof}
                    Show that $\sqrt{n} \neq 0$.
                    \begin{subproof}
                        Suppose not.
                        Thus $\sqrt{n} = 0$.
                        Thus $0 < 0$ by \cref{real_lt_subst_left}.
                        Contradiction by \cref{real_lt_irreflexive}.
                    \end{subproof}
                    $\inv{\sqrt{n}}\in\reals$ by \cref{real_mul_inverse}.
                    $\inv{\sqrt{n}} > 0$ by \cref{real_inv_pos}.
                    $\epsilon \rmul \inv{\sqrt{n}}\in\reals$ by \cref{real_mul_closed}.
                    Thus $\delta > 0$ by \cref{real_mul_pos_pos_gt}.
                \end{subproof}
                Show that for all $y$ such that $y\in\metricBall{d_1}{X_0}{x}{\delta}$ we have
                    $y\in\metricBall{d}{X_0}{x}{\epsilon}$.
                \begin{subproof}
                    Fix $y$.
                    Assume $y\in\metricBall{d_1}{X_0}{x}{\delta}$.
                    Thus $y\in X_0$ and $d_1((x,y)) < \delta$ by \cref{metric_ball}.
                    Take $m\in\reals$ such that $((x,y),m)\in d_1$ by \cref{square_metric_value_exists}.
                    Take $r\in\reals$ such that $((x,y),r)\in d$ by \cref{euclidean_metric_value_exists}.
                    $d$ is a function by \cref{euclidean_metric_function}.
                    $d_1$ is a function by \cref{square_metric_function}.
                    Thus $d_1((x,y)) = m$ by \cref{function_apply_intro}.
                    Thus $m < \delta$ by \cref{real_lt_subst_left}.
                    Thus $r \leq \sqrt{n} \rmul m$ by \cref{euclidean_square_metric_bounds}.
                    Show that $\sqrt{n} \rmul m < \epsilon$.
                    \begin{subproof}
                        $\sqrt{n}\in\reals$.
                        $m\in\reals$.
                        $\delta\in\reals$.
                        $m \rmul \sqrt{n} < \delta \rmul \sqrt{n}$ by \cref{real_lt_mul_pos}.
                        Thus $\sqrt{n} \rmul m < \sqrt{n} \rmul \delta$ by \cref{real_mul_comm}.
                        Show that $\sqrt{n} \rmul \delta = \epsilon$.
                        \begin{subproof}
                            $\sqrt{n}\in\reals$.
                            $\epsilon\in\reals$.
                            $\inv{\sqrt{n}}\in\reals$.
                            $\sqrt{n} \rmul \delta = \sqrt{n} \rmul (\epsilon \rmul \inv{\sqrt{n}})$.
                            $\sqrt{n} \rmul (\epsilon \rmul \inv{\sqrt{n}}) = (\sqrt{n} \rmul \epsilon) \rmul \inv{\sqrt{n}}$ by \cref{real_mul_assoc}.
                            $(\sqrt{n} \rmul \epsilon) \rmul \inv{\sqrt{n}} = (\epsilon \rmul \sqrt{n}) \rmul \inv{\sqrt{n}}$ by \cref{real_mul_comm}.
                            $(\epsilon \rmul \sqrt{n}) \rmul \inv{\sqrt{n}} = \epsilon \rmul (\sqrt{n} \rmul \inv{\sqrt{n}})$ by \cref{real_mul_assoc}.
                            Show that $\sqrt{n} \neq 0$.
                            \begin{subproof}
                                Suppose not.
                                Thus $\sqrt{n} = 0$.
                                Thus $0 < 0$ by \cref{real_lt_subst_left}.
                                Contradiction by \cref{real_lt_irreflexive}.
                            \end{subproof}
                            $\sqrt{n} \rmul \inv{\sqrt{n}} = 1$ by \cref{real_mul_inverse,real_mul_comm}.
                            Thus $\sqrt{n} \rmul \delta = \epsilon \rmul 1$.
                            $1\in\reals$ by \cref{real_mul_ident}.
                            $1 \rmul \epsilon = \epsilon$ by \cref{real_mul_ident}.
                            Thus $\epsilon \rmul 1 = \epsilon$ by \cref{real_mul_comm}.
                        \end{subproof}
                        Thus $\sqrt{n} \rmul m < \epsilon$ by \cref{real_lt_subst_right}.
                    \end{subproof}
                    Show that $r < \epsilon$.
                    \begin{subproof}
                        $r\in\reals$.
                        $\sqrt{n} \rmul m\in\reals$ by \cref{real_mul_closed}.
                        $\epsilon\in\reals$.
                        $r < \sqrt{n} \rmul m$ or $r = \sqrt{n} \rmul m$ by \cref{real_leq_split}.
                        \begin{byCase}
                        \caseOf{$r < \sqrt{n} \rmul m$.} Thus $r < \epsilon$ by \cref{real_lt_trans}.
                        \caseOf{$r = \sqrt{n} \rmul m$.} Thus $r < \epsilon$ by \cref{real_lt_subst_left}.
                        \end{byCase}
                    \end{subproof}
                    Thus $d((x,y)) = r$ by \cref{function_apply_intro}.
                    Thus $y\in\metricBall{d}{X_0}{x}{\epsilon}$ by \cref{metric_ball}.
                \end{subproof}
                Thus $\metricBall{d_1}{X_0}{x}{\delta} \subseteq \metricBall{d}{X_0}{x}{\epsilon}$ by \cref{subseteq}.
            \end{subproof}
            Thus $T_S$ is finer than $T_E$ on $X_0$ by \cref{metric_topology_finer_iff}.
        \end{subproof}
        Thus $T_S\subseteq T_E$ by \cref{finer_topology}.
        Thus $T_E\subseteq T_S$ by \cref{finer_topology}.
        Thus $T_E = T_S$ by \cref{subseteq_antisymmetric}.
    \end{subproof}
    Thus $T_E = T_S$.
\end{proof}

% from §27 Item 3: compact subsets of $\reals^n$
\begin{theorem}\label{compact_reals_power_iff_closed_bounded}
    Suppose $n\in\naturalsPlus$.
    Let $J = \initialSegment{n}$.
    Let $X = \realsPower{J}$.
    Let $d = \euclideanMetric{n}$.
    Let $\rho = \squareMetric{n}$.
    Let $T = \metricTopology{\rho}{X}$.
    Suppose $A\subseteq X$.
    Let $T_A = \subspaceTopology{T}{A}$.
    Then $A$ is compact with respect to $T_A$ iff
        $A$ is closed in $X$ with respect to $T$ and
        $A$ is bounded in $X$ with respect to $d$ or $A$ is bounded in $X$ with respect to $\rho$.
\end{theorem}
\begin{proof}
    Let $X_0 = X$.
    Let $d_1 = \rho$.
    Let $T_S = T$.
    Show that if $A$ is compact with respect to $T_A$ then
        $A$ is closed in $X$ with respect to $T$ and
        $A$ is bounded in $X$ with respect to $d$ or $A$ is bounded in $X$ with respect to $\rho$.
    \begin{subproof}
        Assume $A$ is compact with respect to $T_A$.
        $d_1$ is a metric on $X$ by \cref{square_metric_is_metric}.
        $T$ is a topology on $X$ by \cref{basis_topology_is_topology,metric_topology,metric_basis_is_basis,basis_for_topology}.
        Show that $X$ is Hausdorff with respect to $T$.
        \begin{subproof}
            Let $d_E = d$.
            Let $T_E = \metricTopology{d_E}{X_0}$.
            Let $T_S = \metricTopology{d_1}{X_0}$.
            Let $X_1 = \{(i,\reals) \mid i\in J\}$.
            Let $T_1 = \{(i,\standardTopologyR) \mid i\in J\}$.
            Let $T_0 = \productTopologyIndexed{T_1}{X_1}$.
            $d_E$ is a metric on $X_0$ by \cref{euclidean_metric_is_metric}.
            $d_1$ is a metric on $X_0$ by \cref{square_metric_is_metric}.
            $1\in\naturalsPlus$ by \cref{real_one_in_naturals}.
            $1 \leq n$ by \cref{real_one_leq_naturalsplus}.
            Thus $1\in J$ by \cref{initial_segment_naturalsplus}.
            Thus $J$ is inhabited.
            Show that $X_1$ is a function.
            \begin{subproof}
                Show that for all $w$ such that $w\in X_1$ there exist $x,y$ such that $w = (x,y)$.
                \begin{subproof}
                    Fix $w$.
                    Assume $w\in X_1$.
                    Take $i\in J$ such that $w = (i,\reals)$.
                    Thus there exist $x,y$ such that $w = (x,y)$.
                \end{subproof}
                Thus $X_1$ is a relation by \cref{relation}.
                Show that for all $i,y_1,y_2$ such that $(i,y_1)\in X_1$ and $(i,y_2)\in X_1$ we have $y_1 = y_2$.
                \begin{subproof}
                    Fix $i$.
                    Fix $y_1$.
                    Fix $y_2$.
                    Assume $(i,y_1)\in X_1$ and $(i,y_2)\in X_1$.
                    Then $y_1 = \reals$ and $y_2 = \reals$.
                    Thus $y_1 = y_2$.
                \end{subproof}
                Thus $X_1$ is right-unique by \cref{rightunique}.
                Thus $X_1$ is a function.
            \end{subproof}
            Show that $T_1$ is a function.
            \begin{subproof}
                Show that for all $w$ such that $w\in T_1$ there exist $x,y$ such that $w = (x,y)$.
                \begin{subproof}
                    Fix $w$.
                    Assume $w\in T_1$.
                    Take $i\in J$ such that $w = (i,\standardTopologyR)$.
                    Thus there exist $x,y$ such that $w = (x,y)$.
                \end{subproof}
                Thus $T_1$ is a relation by \cref{relation}.
                Show that for all $i,y_1,y_2$ such that $(i,y_1)\in T_1$ and $(i,y_2)\in T_1$ we have $y_1 = y_2$.
                \begin{subproof}
                    Fix $i$.
                    Fix $y_1$.
                    Fix $y_2$.
                    Assume $(i,y_1)\in T_1$ and $(i,y_2)\in T_1$.
                    Then $y_1 = \standardTopologyR$ and $y_2 = \standardTopologyR$.
                    Thus $y_1 = y_2$.
                \end{subproof}
                Thus $T_1$ is right-unique by \cref{rightunique}.
                Thus $T_1$ is a function.
            \end{subproof}
            Show that $\dom{X_1} = J$.
            \begin{subproof}
                Show that for all $i$ such that $i\in J$ we have $i\in\dom{X_1}$.
                \begin{subproof}
                    Fix $i$.
                    Assume $i\in J$.
                    Thus $(i,\reals)\in X_1$.
                    Thus $i\in\dom{X_1}$ by \cref{dom_intro}.
                \end{subproof}
                Thus $J\subseteq\dom{X_1}$ by \cref{subseteq}.
                Show that for all $i$ such that $i\in\dom{X_1}$ we have $i\in J$.
                \begin{subproof}
                    Fix $i$.
                    Assume $i\in\dom{X_1}$.
                    Take $y$ such that $(i,y)\in X_1$ by \cref{dom_iff}.
                    Thus $i\in J$.
                \end{subproof}
                Thus $\dom{X_1}\subseteq J$ by \cref{subseteq}.
                Thus $\dom{X_1} = J$ by \cref{subseteq_antisymmetric}.
            \end{subproof}
            Show that $\dom{T_1} = J$.
            \begin{subproof}
                Show that for all $i$ such that $i\in J$ we have $i\in\dom{T_1}$.
                \begin{subproof}
                    Fix $i$.
                    Assume $i\in J$.
                    Thus $(i,\standardTopologyR)\in T_1$.
                    Thus $i\in\dom{T_1}$ by \cref{dom_intro}.
                \end{subproof}
                Thus $J\subseteq\dom{T_1}$ by \cref{subseteq}.
                Show that for all $i$ such that $i\in\dom{T_1}$ we have $i\in J$.
                \begin{subproof}
                    Fix $i$.
                    Assume $i\in\dom{T_1}$.
                    Take $y$ such that $(i,y)\in T_1$ by \cref{dom_iff}.
                    Thus $i\in J$.
                \end{subproof}
                Thus $\dom{T_1}\subseteq J$ by \cref{subseteq}.
                Thus $\dom{T_1} = J$ by \cref{subseteq_antisymmetric}.
            \end{subproof}
            Show that for all $\alpha\in\dom{X_1}$ we have $\apply{X_1}{\alpha}$ is Hausdorff with respect to $\apply{T_1}{\alpha}$.
            \begin{subproof}
                Fix $\alpha$.
                Assume $\alpha\in\dom{X_1}$.
                Thus $\alpha\in J$ by \cref{subseteq}.
                Thus $(\alpha,\reals)\in X_1$.
                Thus $\apply{X_1}{\alpha} = \reals$ by \cref{function_apply_intro}.
                Thus $(\alpha,\standardTopologyR)\in T_1$.
                Thus $\apply{T_1}{\alpha} = \standardTopologyR$ by \cref{function_apply_intro}.
                $\realOrderTopology = \standardTopologyR$ by \cref{real_order_topology_eq_standard}.
                $\realOrderTopology = \basisTopology{\orderBasis{\realOrderRel}{\reals}}{\reals}$ by \cref{real_order_topology}.
                $\realOrderRel$ is a simple order relation on $\reals$ by \cref{real_order_relation_simple}.
                $\reals$ has more than one element by \cref{linear_continuum,reals_linear_continuum}.
                Thus $\realOrderTopology$ is the order topology on $\reals$ with respect to $\realOrderRel$ by \cref{order_topology}.
                Thus $\reals$ is Hausdorff with respect to $\realOrderTopology$ by \cref{order_topology_hausdorff}.
                Thus $\reals$ is Hausdorff with respect to $\standardTopologyR$ by \cref{real_order_topology_eq_standard}.
                Thus $\apply{X_1}{\alpha}$ is Hausdorff with respect to $\apply{T_1}{\alpha}$.
            \end{subproof}
            Thus $\productFamily{X_1}$ is Hausdorff with respect to $T_0$ by \cref{product_hausdorff_product}.
            Thus $\productFamily{X_1} = X_0$ by \cref{product_family_reals_power}.
            Thus $X_0$ is Hausdorff with respect to $T_0$.
            $T_E = T_0$ by \cref{euclidean_metric_product_topology}.
            Show that $T_E = T_S$.
            \begin{subproof}
                Show that $T_E$ is finer than $T_S$ on $X_0$.
                \begin{subproof}
                    Show that for all $x\in X_0$ we have
                        for all $\epsilon\in\reals$ such that $0 < \epsilon$ there exists $\delta\in\reals$
                            such that $0 < \delta$ and
                            $\metricBall{d_E}{X_0}{x}{\delta} \subseteq \metricBall{d_1}{X_0}{x}{\epsilon}$.
                    \begin{subproof}
                        Fix $x$.
                        Assume $x\in X_0$.
                        Fix $\epsilon\in\reals$.
                        Assume $0 < \epsilon$.
                        Let $\delta = \epsilon$.
                        $0 < \delta$.
                        Show that for all $y$ such that $y\in\metricBall{d_E}{X_0}{x}{\delta}$ we have
                            $y\in\metricBall{d_1}{X_0}{x}{\epsilon}$.
                        \begin{subproof}
                            Fix $y$.
                            Assume $y\in\metricBall{d_E}{X_0}{x}{\delta}$.
                            Thus $y\in X_0$ and $d_E((x,y)) < \delta$ by \cref{metric_ball}.
                            Take $r\in\reals$ such that $((x,y),r)\in d_E$ by \cref{euclidean_metric_value_exists}.
                            Take $m\in\reals$ such that $((x,y),m)\in d_1$ by \cref{square_metric_value_exists}.
                            Thus $m \leq r$ by \cref{euclidean_square_metric_bounds}.
                            $d_E$ is a function by \cref{euclidean_metric_function}.
                            $d_1$ is a function by \cref{square_metric_function}.
                            Thus $d_E((x,y)) = r$ by \cref{function_apply_intro}.
                            Thus $r < \epsilon$ by \cref{real_lt_subst_left}.
                            Show that $m < \epsilon$.
                            \begin{subproof}
                                $m < r$ or $m = r$ by \cref{real_leq_split}.
                                \begin{byCase}
                                \caseOf{$m < r$.} Thus $m < \epsilon$ by \cref{real_lt_trans}.
                                \caseOf{$m = r$.} Thus $m < \epsilon$ by \cref{real_lt_subst_left}.
                                \end{byCase}
                            \end{subproof}
                            Thus $d_1((x,y)) = m$ by \cref{function_apply_intro}.
                            Thus $y\in\metricBall{d_1}{X_0}{x}{\epsilon}$ by \cref{metric_ball}.
                        \end{subproof}
                        Thus $\metricBall{d_E}{X_0}{x}{\delta} \subseteq \metricBall{d_1}{X_0}{x}{\epsilon}$ by \cref{subseteq}.
                    \end{subproof}
                    Thus $T_E$ is finer than $T_S$ on $X_0$ by \cref{metric_topology_finer_iff}.
                \end{subproof}
                Show that $T_S$ is finer than $T_E$ on $X_0$.
                \begin{subproof}
                    Show that for all $x\in X_0$ we have
                        for all $\epsilon\in\reals$ such that $0 < \epsilon$ there exists $\delta\in\reals$
                            such that $0 < \delta$ and
                            $\metricBall{d_1}{X_0}{x}{\delta} \subseteq \metricBall{d_E}{X_0}{x}{\epsilon}$.
                    \begin{subproof}
                        Fix $x$.
                        Assume $x\in X_0$.
                        Fix $\epsilon\in\reals$.
                        Assume $0 < \epsilon$.
                        $n\in\reals$ by \cref{real_naturals_subset_reals,subseteq}.
                        Show that $0 \leq n$.
                        \begin{subproof}
                            $0 < 1$ by \cref{real_zero_lt_one}.
                            $0\in\reals$ by \cref{real_add_ident}.
                            $1\in\reals$ by \cref{real_naturals_subset_reals,subseteq}.
                            $1 \leq n$ by \cref{real_one_leq_naturalsplus}.
                            Thus $1 < n$ or $1 = n$ by \cref{real_leq_split}.
                            \begin{byCase}
                            \caseOf{$1 < n$.}
                                Thus $0 < n$ by \cref{real_lt_trans}.
                                Thus $0 \leq n$.
                            \caseOf{$1 = n$.}
                                Thus $0 < n$ by \cref{real_lt_subst_right}.
                                Thus $0 \leq n$.
                            \end{byCase}
                        \end{subproof}
                        Thus $\sqrt{n}\in\reals$ and $\sqrt{n} \geq 0$ and $\sqrt{n} \rmul \sqrt{n} = n$ by \cref{positive_square_root}.
                        Show that $\sqrt{n}$ is positive.
                        \begin{subproof}
                            Show that $\sqrt{n} \neq 0$.
                            \begin{subproof}
                                Suppose not.
                                Thus $\sqrt{n} = 0$.
                                Thus $n = 0$ by \cref{positive_square_root,real_mul_zero}.
                                $0 < 1$ by \cref{real_zero_lt_one}.
                                $1 \leq n$ by \cref{real_one_leq_naturalsplus}.
                                Show that $0 < n$.
                                \begin{subproof}
                                    $0\in\reals$ by \cref{real_add_ident}.
                                    $1\in\reals$ by \cref{real_naturals_subset_reals,subseteq}.
                                    $n\in\reals$ by \cref{real_naturals_subset_reals,subseteq}.
                                    Thus $1 < n$ or $1 = n$ by \cref{real_leq_split}.
                                    \begin{byCase}
                                    \caseOf{$1 < n$.}
                                        Thus $0 < n$ by \cref{real_lt_trans}.
                                    \caseOf{$1 = n$.}
                                        Thus $0 < n$ by \cref{real_lt_subst_right}.
                                    \end{byCase}
                                \end{subproof}
                                Thus $0 < 0$ by \cref{real_lt_subst_right}.
                                Contradiction by \cref{real_lt_irreflexive}.
                            \end{subproof}
                            Thus $\sqrt{n} > 0$.
                        \end{subproof}
                        Let $\delta = \epsilon \rmul \inv{\sqrt{n}}$.
                        Show that $0 < \delta$.
                        \begin{subproof}
                            Show that $\sqrt{n} \neq 0$.
                            \begin{subproof}
                                Suppose not.
                                Thus $\sqrt{n} = 0$.
                                Thus $0 < 0$ by \cref{real_lt_subst_left}.
                                Contradiction by \cref{real_lt_irreflexive}.
                            \end{subproof}
                            $\inv{\sqrt{n}}\in\reals$ by \cref{real_mul_inverse}.
                            $\inv{\sqrt{n}} > 0$ by \cref{real_inv_pos}.
                            $\epsilon \rmul \inv{\sqrt{n}}\in\reals$ by \cref{real_mul_closed}.
                            Thus $\delta > 0$ by \cref{real_mul_pos_pos_gt}.
                        \end{subproof}
                        Show that for all $y$ such that $y\in\metricBall{d_1}{X_0}{x}{\delta}$ we have
                            $y\in\metricBall{d_E}{X_0}{x}{\epsilon}$.
                        \begin{subproof}
                            Fix $y$.
                            Assume $y\in\metricBall{d_1}{X_0}{x}{\delta}$.
                            Thus $y\in X_0$ and $d_1((x,y)) < \delta$ by \cref{metric_ball}.
                            Take $m\in\reals$ such that $((x,y),m)\in d_1$ by \cref{square_metric_value_exists}.
                            Take $r\in\reals$ such that $((x,y),r)\in d_E$ by \cref{euclidean_metric_value_exists}.
                            $d_E$ is a function by \cref{euclidean_metric_function}.
                            $d_1$ is a function by \cref{square_metric_function}.
                            Thus $d_1((x,y)) = m$ by \cref{function_apply_intro}.
                            Thus $m < \delta$ by \cref{real_lt_subst_left}.
                            Thus $r \leq \sqrt{n} \rmul m$ by \cref{euclidean_square_metric_bounds}.
                            Show that $\sqrt{n} \rmul m < \epsilon$.
                            \begin{subproof}
                                $\sqrt{n}\in\reals$.
                                $m\in\reals$.
                                $\delta\in\reals$.
                                $m \rmul \sqrt{n} < \delta \rmul \sqrt{n}$ by \cref{real_lt_mul_pos}.
                                Thus $\sqrt{n} \rmul m < \sqrt{n} \rmul \delta$ by \cref{real_mul_comm}.
                                Show that $\sqrt{n} \rmul \delta = \epsilon$.
                                \begin{subproof}
                                    $\sqrt{n}\in\reals$.
                                    $\epsilon\in\reals$.
                                    $\inv{\sqrt{n}}\in\reals$.
                                    $\sqrt{n} \rmul \delta = \sqrt{n} \rmul (\epsilon \rmul \inv{\sqrt{n}})$.
                                    $\sqrt{n} \rmul (\epsilon \rmul \inv{\sqrt{n}}) = (\sqrt{n} \rmul \epsilon) \rmul \inv{\sqrt{n}}$ by \cref{real_mul_assoc}.
                                    $(\sqrt{n} \rmul \epsilon) \rmul \inv{\sqrt{n}} = (\epsilon \rmul \sqrt{n}) \rmul \inv{\sqrt{n}}$ by \cref{real_mul_comm}.
                                    $(\epsilon \rmul \sqrt{n}) \rmul \inv{\sqrt{n}} = \epsilon \rmul (\sqrt{n} \rmul \inv{\sqrt{n}})$ by \cref{real_mul_assoc}.
                                    Show that $\sqrt{n} \neq 0$.
                                    \begin{subproof}
                                        Suppose not.
                                        Thus $\sqrt{n} = 0$.
                                        Thus $0 < 0$ by \cref{real_lt_subst_left}.
                                        Contradiction by \cref{real_lt_irreflexive}.
                                    \end{subproof}
                                    $\sqrt{n} \rmul \inv{\sqrt{n}} = 1$ by \cref{real_mul_inverse,real_mul_comm}.
                                    Thus $\sqrt{n} \rmul \delta = \epsilon \rmul 1$.
                                    $1\in\reals$ by \cref{real_mul_ident}.
                                    $1 \rmul \epsilon = \epsilon$ by \cref{real_mul_ident}.
                                    Thus $\epsilon \rmul 1 = \epsilon$ by \cref{real_mul_comm}.
                                \end{subproof}
                                Thus $\sqrt{n} \rmul m < \epsilon$ by \cref{real_lt_subst_right}.
                            \end{subproof}
                            Show that $r < \epsilon$.
                            \begin{subproof}
                                $r\in\reals$.
                                $\sqrt{n} \rmul m\in\reals$ by \cref{real_mul_closed}.
                                $\epsilon\in\reals$.
                                $r < \sqrt{n} \rmul m$ or $r = \sqrt{n} \rmul m$ by \cref{real_leq_split}.
                                \begin{byCase}
                                \caseOf{$r < \sqrt{n} \rmul m$.} Thus $r < \epsilon$ by \cref{real_lt_trans}.
                                \caseOf{$r = \sqrt{n} \rmul m$.} Thus $r < \epsilon$ by \cref{real_lt_subst_left}.
                                \end{byCase}
                            \end{subproof}
                            Thus $d_E((x,y)) = r$ by \cref{function_apply_intro}.
                            Thus $y\in\metricBall{d_E}{X_0}{x}{\epsilon}$ by \cref{metric_ball}.
                        \end{subproof}
                        Thus $\metricBall{d_1}{X_0}{x}{\delta} \subseteq \metricBall{d_E}{X_0}{x}{\epsilon}$ by \cref{subseteq}.
                    \end{subproof}
                    Thus $T_S$ is finer than $T_E$ on $X_0$ by \cref{metric_topology_finer_iff}.
                \end{subproof}
                Thus $T_S\subseteq T_E$ by \cref{finer_topology}.
                Thus $T_E\subseteq T_S$ by \cref{finer_topology}.
                Thus $T_E = T_S$ by \cref{subseteq_antisymmetric}.
            \end{subproof}
            Thus $T_S = T_0$ by \cref{subseteq_antisymmetric,subseteq}.
            Thus $X$ is Hausdorff with respect to $T$.
        \end{subproof}
        Thus $A$ is closed in $X$ with respect to $T$ by \cref{compact_subspace_closed}.
        Show that $A$ is bounded in $X$ with respect to $\rho$.
        \begin{subproof}
            \begin{byCase}
            \caseOf{$A = \emptyset$.}
                Let $M = 0$.
                $M\in\reals$ by \cref{real_add_ident}.
                Show that for all $a_1,a_2\in A$ we have $\rho((a_1,a_2)) \leq M$.
                \begin{subproof}
                    Fix $a_1$.
                    Fix $a_2$.
                    Assume $a_1\in A$ and $a_2\in A$.
                    Suppose not.
                    Thus $A\subseteq\emptyset$ by \cref{subseteq_emptyset_iff}.
                    Thus $a_1\in\emptyset$ by \cref{elem_subseteq}.
                    Contradiction by \cref{emptyset}.
                \end{subproof}
                Thus $A$ is bounded in $X$ with respect to $\rho$ by \cref{metric_bounded}.
            \caseOf{$A \neq \emptyset$.}
                $A$ is inhabited by \cref{nonempty_inhabited}.
                Let $C = \{\metricBall{\rho}{X}{a}{1} \mid a\in A\}$.
                Show that for all $U$ such that $U\in C$ we have $U$ is open in $X$ with respect to $T$.
                \begin{subproof}
                    Fix $U$.
                    Assume $U\in C$.
                    Take $a\in A$ such that $U = \metricBall{\rho}{X}{a}{1}$.
                    Thus $a\in X$ by \cref{subseteq}.
                    Show that $U\subseteq X$.
                    \begin{subproof}
                        Show that for all $x$ such that $x\in U$ we have $x\in X$.
                        \begin{subproof}
                            Fix $x$.
                            Assume $x\in U$.
                            Thus $x\in X$ by \cref{metric_ball}.
                        \end{subproof}
                        Thus $U\subseteq X$ by \cref{subseteq}.
                    \end{subproof}
                    Thus $U\in\pow{X}$ by \cref{powerset_intro}.
                    $1\in\reals$ by \cref{real_naturals_subset_reals,real_one_in_naturals,subseteq}.
                    $0 < 1$ by \cref{real_zero_lt_one}.
                    Thus $U\in\metricBasis{\rho}{X}$ by \cref{metric_basis}.
                    $\metricBasis{\rho}{X}$ is a basis on $X$ by \cref{metric_basis_is_basis}.
                    Thus $U\in T$ by \cref{basis_elem_open_in_generated,metric_topology,basis_for_topology}.
                    Thus $U$ is open in $X$ with respect to $T$ by \cref{open_in}.
                \end{subproof}
                Show that $C$ is a covering of $A$.
                \begin{subproof}
                    Show that for all $x$ such that $x\in A$ we have $x\in\unions{C}$.
                    \begin{subproof}
                        Fix $x$.
                        Assume $x\in A$.
                        Let $U = \metricBall{\rho}{X}{x}{1}$.
                        $U\in C$.
                        Thus $x\in X$ by \cref{subseteq}.
                        $\rho$ is a metric on $X$ by \cref{square_metric_is_metric}.
                        Thus $\rho((x,x)) = 0$ by \cref{metric_zero_characterization,metric_on}.
                        $0 < 1$ by \cref{real_zero_lt_one}.
                        Thus $\rho((x,x)) < 1$ by \cref{real_lt_subst_left}.
                        Thus $x\in U$ by \cref{metric_ball}.
                        Thus $x\in\unions{C}$ by \cref{unions_iff}.
                    \end{subproof}
                    Thus $A\subseteq\unions{C}$ by \cref{subseteq}.
                    Thus $C$ is a covering of $A$ by \cref{covering}.
                \end{subproof}
                $T$ is a topology on $X$.
                $A$ is a subspace of $X$ with respect to $T$ by \cref{subspace_of}.
                Thus there exists $B$ such that
                    $B\subseteq C$ and
                    $B$ is finite and
                    $B$ is a covering of $A$
                    by \cref{compact_subspace_open_covering}.
                Let $R = \{w\in B\times A \mid \exists U\in B.\exists a\in A. w = (U,a) \land U = \metricBall{\rho}{X}{a}{1}\}$.
                Show that $R$ is a relation.
                \begin{subproof}
                    Show that for all $w$ such that $w\in R$ there exist $x,y$ such that $w = (x,y)$.
                    \begin{subproof}
                        Fix $w$.
                        Assume $w\in R$.
                        Thus $w\in B\times A$.
                        Take $x,y$ such that $w = (x,y)$ by \cref{times_elem_is_tuple}.
                        Thus there exist $x,y$ such that $w = (x,y)$.
                    \end{subproof}
                    Thus $R$ is a relation by \cref{relation}.
                \end{subproof}
                Show that for all $U\in B$ there exists $a$ such that $U\mathrel{R} a$.
                \begin{subproof}
                    Fix $U$.
                    Assume $U\in B$.
                    Thus $U\in C$ by \cref{subseteq}.
                    Take $a\in A$ such that $U = \metricBall{\rho}{X}{a}{1}$.
                    Thus $U\mathrel{R} a$.
                \end{subproof}
                Take $g$ such that $g$ is a function on $B$ and for all $U\in B$ we have $U\mathrel{R}\apply{g}{U}$
                    by \cref{choice_function}.
                Let $F = \img{g}{B}$.
                $F$ is finite by \cref{finite_image_function}.
                Show that $F\subseteq A$.
                \begin{subproof}
                    Show that for all $f$ such that $f\in F$ we have $f\in A$.
                    \begin{subproof}
                        Fix $f$.
                        Assume $f\in F$.
                        Take $U\in B$ such that $(U,f)\in g$ by \cref{img_iff}.
                        Thus $f = g(U)$ by \cref{function_apply_intro}.
                        Thus $U\mathrel{R} f$ by \cref{function_apply_elim}.
                        Thus $f\in A$.
                    \end{subproof}
                    Thus $F\subseteq A$ by \cref{subseteq}.
                \end{subproof}
                Show that $F\neq \emptyset$.
                \begin{subproof}
                    Take $a_0$ such that $a_0\in A$.
                    Thus $A\subseteq\unions{B}$ by \cref{covering}.
                    Thus $a_0\in\unions{B}$ by \cref{subseteq}.
                    Take $U_0\in B$ such that $a_0\in U_0$ by \cref{unions_iff}.
                    Thus $(U_0,g(U_0))\in g$ by \cref{function_apply_elim}.
                    Thus $g(U_0)\in F$ by \cref{img_iff}.
                    Thus $F$ is inhabited by \cref{nonempty_inhabited}.
                    Thus $F\neq\emptyset$.
                \end{subproof}
                Take $a_0$ such that $a_0\in F$.
                Let $h = \{(f,\rho((a_0,f))) \mid f\in F\}$.
                Show that $h$ is a function on $F$.
                \begin{subproof}
                    Show that for all $w$ such that $w\in h$ there exist $x,y$ such that $w = (x,y)$.
                    \begin{subproof}
                        Fix $w$.
                        Assume $w\in h$.
                        Take $f\in F$ such that $w = (f,\rho((a_0,f)))$.
                        Thus there exist $x,y$ such that $w = (x,y)$.
                    \end{subproof}
                    Thus $h$ is a relation by \cref{relation}.
                    Show that for all $x,y_1,y_2$ such that $(x,y_1)\in h$ and $(x,y_2)\in h$ we have $y_1 = y_2$.
                    \begin{subproof}
                        Fix $x$.
                        Fix $y_1$.
                        Fix $y_2$.
                        Assume $(x,y_1)\in h$ and $(x,y_2)\in h$.
                        Take $f_1\in F$ such that $(x,y_1) = (f_1,\rho((a_0,f_1)))$.
                        Take $f_2\in F$ such that $(x,y_2) = (f_2,\rho((a_0,f_2)))$.
                        Thus $f_1 = x$ and $f_2 = x$ by \cref{pair_eq_iff}.
                        Thus $\rho((a_0,f_1)) = \rho((a_0,f_2))$.
                        Thus $y_1 = y_2$ by \cref{pair_eq_iff}.
                    \end{subproof}
                    Thus $h$ is right-unique by \cref{rightunique}.
                    Show that $\dom{h} = F$.
                    \begin{subproof}
                        Show that for all $f$ such that $f\in F$ we have $f\in\dom{h}$.
                        \begin{subproof}
                            Fix $f$.
                            Assume $f\in F$.
                            Thus $(f,\rho((a_0,f)))\in h$.
                            Thus $f\in\dom{h}$ by \cref{dom_intro}.
                        \end{subproof}
                        Thus $F\subseteq\dom{h}$ by \cref{subseteq}.
                        Show that for all $f$ such that $f\in\dom{h}$ we have $f\in F$.
                        \begin{subproof}
                            Fix $f$.
                            Assume $f\in\dom{h}$.
                            Take $y$ such that $(f,y)\in h$ by \cref{dom_iff}.
                            Take $f_0\in F$ such that $(f,y) = (f_0,\rho((a_0,f_0)))$.
                            Thus $f = f_0$ by \cref{pair_eq_iff}.
                            Thus $f\in F$.
                        \end{subproof}
                        Thus $\dom{h}\subseteq F$ by \cref{subseteq}.
                        Thus $\dom{h} = F$ by \cref{subseteq_antisymmetric}.
                    \end{subproof}
                    Thus $h$ is a function on $F$.
                \end{subproof}
                Let $D = \img{h}{F}$.
                $D$ is finite by \cref{finite_image_function}.
                Show that $D\subseteq\reals$.
                \begin{subproof}
                    Show that for all $r$ such that $r\in D$ we have $r\in\reals$.
                    \begin{subproof}
                        Fix $r$.
                        Assume $r\in D$.
                        Take $f\in F$ such that $(f,r)\in h$ by \cref{img_iff}.
                        Show that $h\subseteq F\times\reals$.
                        \begin{subproof}
                            Show that for all $w$ such that $w\in h$ we have $w\in F\times\reals$.
                            \begin{subproof}
                                Fix $w$.
                                Assume $w\in h$.
                                Take $f_0\in F$ such that $w = (f_0,\rho((a_0,f_0)))$.
                                $f_0\in A$ by \cref{subseteq}.
                                $f_0\in X$ by \cref{subseteq}.
                                $a_0\in A$ by \cref{subseteq}.
                                $a_0\in X$ by \cref{subseteq}.
                                $(a_0,f_0)\in X\times X$ by \cref{times_tuple_intro}.
                                Thus $\rho((a_0,f_0))\in\reals$ by \cref{metric_on,funs_type_apply,square_metric_is_metric}.
                                Thus $(f_0,\rho((a_0,f_0)))\in F\times\reals$ by \cref{times_tuple_intro}.
                                Thus $w\in F\times\reals$.
                            \end{subproof}
                            Thus $h\subseteq F\times\reals$ by \cref{subseteq}.
                        \end{subproof}
                        Thus $(f,r)\in F\times\reals$ by \cref{subseteq}.
                        Thus $r\in\reals$ by \cref{times_tuple_elim}.
                    \end{subproof}
                    Thus $D\subseteq\reals$ by \cref{subseteq}.
                \end{subproof}
                Show that $D\neq\emptyset$.
                \begin{subproof}
                    Take $f_0$ such that $f_0\in F$.
                    Thus $(f_0,\rho((a_0,f_0)))\in h$.
                    Thus $\rho((a_0,f_0))\in D$ by \cref{img_iff}.
                    Thus $D$ is inhabited by \cref{nonempty_inhabited}.
                    Thus $D\neq\emptyset$.
                \end{subproof}
                Take $M_0\in D$ such that $M_0$ is a maximum of $D$ by \cref{finite_real_maximum}.
                Let $M = (M_0 + 1) + (M_0 + 1)$.
                Show that for all $x,y\in A$ we have $\rho((x,y)) \leq M$.
                \begin{subproof}
                    Fix $x$.
                    Fix $y$.
                    Assume $x\in A$.
                    Assume $y\in A$.
                    Take $U_x\in B$ such that $x\in U_x$ by \cref{covering,subseteq,unions_iff}.
                    Take $U_y\in B$ such that $y\in U_y$ by \cref{covering,subseteq,unions_iff}.
                    Let $f_x = g(U_x)$.
                    Let $f_y = g(U_y)$.
                    $U_x\mathrel{R} f_x$ by \cref{function_apply_elim}.
                    $U_y\mathrel{R} f_y$ by \cref{function_apply_elim}.
                    Show that $U_x = \metricBall{\rho}{X}{f_x}{1}$.
                    \begin{subproof}
                        Thus $(U_x,f_x)\in R$.
                        Take $U_0, a_0$ such that
                            $U_0\in B$ and $a_0\in A$ and
                            $(U_x,f_x) = (U_0,a_0)$ and $U_0 = \metricBall{\rho}{X}{a_0}{1}$.
                        Thus $U_x = U_0$ by \cref{pair_eq_iff}.
                        Thus $U_x = \metricBall{\rho}{X}{a_0}{1}$.
                        Thus $f_x = a_0$ by \cref{pair_eq_iff}.
                        Thus $U_x = \metricBall{\rho}{X}{f_x}{1}$.
                    \end{subproof}
                    Show that $U_y = \metricBall{\rho}{X}{f_y}{1}$.
                    \begin{subproof}
                        Thus $(U_y,f_y)\in R$.
                        Take $U_1, a_1$ such that
                            $U_1\in B$ and $a_1\in A$ and
                            $(U_y,f_y) = (U_1,a_1)$ and $U_1 = \metricBall{\rho}{X}{a_1}{1}$.
                        Thus $U_y = U_1$ by \cref{pair_eq_iff}.
                        Thus $U_y = \metricBall{\rho}{X}{a_1}{1}$.
                        Thus $f_y = a_1$ by \cref{pair_eq_iff}.
                        Thus $U_y = \metricBall{\rho}{X}{f_y}{1}$.
                    \end{subproof}
                    Show that $\rho((x,f_x)) < 1$.
                    \begin{subproof}
                        Thus $x\in U_x$.
                        Thus $x\in \metricBall{\rho}{X}{f_x}{1}$.
                        Thus $x\in X$ and $\rho((f_x,x)) < 1$ by \cref{metric_ball}.
                        Thus $(U_x,f_x)\in R$.
                        Thus $(U_x,f_x)\in B\times A$.
                        Thus $f_x\in A$ by \cref{times_tuple_elim}.
                        $f_x\in X$ by \cref{subseteq}.
                        $\rho$ is symmetric on $X$ by \cref{metric_symmetric,metric_on,square_metric_is_metric}.
                        Thus $\rho((x,f_x)) = \rho((f_x,x))$ by \cref{metric_symmetric}.
                        Thus $\rho((x,f_x)) < 1$ by \cref{real_lt_subst_left}.
                    \end{subproof}
                    Show that $\rho((y,f_y)) < 1$.
                    \begin{subproof}
                        Thus $y\in U_y$.
                        Thus $y\in \metricBall{\rho}{X}{f_y}{1}$.
                        Thus $y\in X$ and $\rho((f_y,y)) < 1$ by \cref{metric_ball}.
                        Thus $(U_y,f_y)\in R$.
                        Thus $(U_y,f_y)\in B\times A$.
                        Thus $f_y\in A$ by \cref{times_tuple_elim}.
                        $f_y\in X$ by \cref{subseteq}.
                        $\rho$ is symmetric on $X$ by \cref{metric_symmetric,metric_on,square_metric_is_metric}.
                        Thus $\rho((y,f_y)) = \rho((f_y,y))$ by \cref{metric_symmetric}.
                        Thus $\rho((y,f_y)) < 1$ by \cref{real_lt_subst_left}.
                    \end{subproof}
                    Show that $f_x\in F$.
                    \begin{subproof}
                        $U_x\in B$.
                        $g$ is a function on $B$.
                        Thus $B = \dom{g}$.
                        Thus $U_x\in\dom{g}$ by \cref{subseteq}.
                        Thus $(U_x,f_x)\in g$ by \cref{function_apply_elim}.
                        Thus $f_x\in \img{g}{B}$ by \cref{img_iff}.
                        Thus $f_x\in F$.
                    \end{subproof}
                    Show that $f_y\in F$.
                    \begin{subproof}
                        $U_y\in B$.
                        $g$ is a function on $B$.
                        Thus $B = \dom{g}$.
                        Thus $U_y\in\dom{g}$ by \cref{subseteq}.
                        Thus $(U_y,f_y)\in g$ by \cref{function_apply_elim}.
                        Thus $f_y\in \img{g}{B}$ by \cref{img_iff}.
                        Thus $f_y\in F$.
                    \end{subproof}
                    Show that $h(f_x)\in D$.
                    \begin{subproof}
                        $h$ is a function on $F$.
                        Thus $F = \dom{h}$.
                        Thus $f_x\in\dom{h}$ by \cref{subseteq}.
                        Thus $(f_x,h(f_x))\in h$ by \cref{function_apply_elim}.
                        Thus $h(f_x)\in \img{h}{F}$ by \cref{img_iff}.
                        Thus $h(f_x)\in D$.
                    \end{subproof}
                    Show that $h(f_y)\in D$.
                    \begin{subproof}
                        $h$ is a function on $F$.
                        Thus $F = \dom{h}$.
                        Thus $f_y\in\dom{h}$ by \cref{subseteq}.
                        Thus $(f_y,h(f_y))\in h$ by \cref{function_apply_elim}.
                        Thus $h(f_y)\in \img{h}{F}$ by \cref{img_iff}.
                        Thus $h(f_y)\in D$.
                    \end{subproof}
                    Thus $h(f_x) \leq M_0$ and $h(f_y) \leq M_0$ by \cref{real_maximum}.
                    Show that $\rho((a_0,f_x)) \leq M_0$.
                    \begin{subproof}
                        Thus $(f_x,\rho((a_0,f_x)))\in h$.
                        Thus $h(f_x) = \rho((a_0,f_x))$ by \cref{function_apply_intro}.
                        Thus $\rho((a_0,f_x)) \leq M_0$.
                    \end{subproof}
                    Show that $\rho((a_0,f_y)) \leq M_0$.
                    \begin{subproof}
                        Thus $(f_y,\rho((a_0,f_y)))\in h$.
                        Thus $h(f_y) = \rho((a_0,f_y))$ by \cref{function_apply_intro}.
                        Thus $\rho((a_0,f_y)) \leq M_0$.
                    \end{subproof}
                    $\rho$ is a metric on $X$ by \cref{square_metric_is_metric}.
                    Thus $\rho$ satisfies the triangle inequality on $X$ by \cref{metric_on}.
                    Show that $\rho((x,y)) \leq \rho((x,f_x)) + \rho((f_x,a_0)) + \rho((a_0,f_y)) + \rho((f_y,y))$.
                    \begin{subproof}
                        $x\in X$ by \cref{subseteq}.
                        $y\in X$ by \cref{subseteq}.
                        $f_x\in X$ by \cref{subseteq}.
                        $f_y\in X$ by \cref{subseteq}.
                        $a_0\in X$ by \cref{subseteq}.
                        Thus $\rho((x,y)) \leq \rho((x,f_x)) + \rho((f_x,y))$
                            by \cref{metric_triangle_inequality}.
                        Thus $\rho((f_x,y)) \leq \rho((f_x,a_0)) + \rho((a_0,y))$
                            by \cref{metric_triangle_inequality}.
                        Thus $\rho((a_0,y)) \leq \rho((a_0,f_y)) + \rho((f_y,y))$
                            by \cref{metric_triangle_inequality}.
                        Thus $x\in X$ by \cref{elem_subseteq}.
                        Thus $f_x\in A$ by \cref{elem_subseteq}.
                        Thus $f_x\in X$ by \cref{elem_subseteq}.
                        Thus $f_y\in A$ by \cref{elem_subseteq}.
                        Thus $f_y\in X$ by \cref{elem_subseteq}.
                        Thus $y\in X$ by \cref{elem_subseteq}.
                        Thus $a_0\in A$ by \cref{elem_subseteq}.
                        Thus $a_0\in X$ by \cref{elem_subseteq}.
                        $(x,f_x)\in X\times X$ by \cref{times_tuple_intro}.
                        $(f_x,y)\in X\times X$ by \cref{times_tuple_intro}.
                        $(f_x,a_0)\in X\times X$ by \cref{times_tuple_intro}.
                        $(a_0,y)\in X\times X$ by \cref{times_tuple_intro}.
                        $(a_0,f_y)\in X\times X$ by \cref{times_tuple_intro}.
                        $(f_y,y)\in X\times X$ by \cref{times_tuple_intro}.
                        $(x,y)\in X\times X$ by \cref{times_tuple_intro}.
                        Thus $\rho((x,y))\in\reals$ by \cref{metric_on,funs_type_apply}.
                        Thus $\rho((x,f_x))\in\reals$ by \cref{metric_on,funs_type_apply}.
                        Thus $\rho((f_x,y))\in\reals$ by \cref{metric_on,funs_type_apply}.
                        Thus $\rho((f_x,a_0))\in\reals$ by \cref{metric_on,funs_type_apply}.
                        Thus $\rho((a_0,y))\in\reals$ by \cref{metric_on,funs_type_apply}.
                        Thus $\rho((a_0,f_y))\in\reals$ by \cref{metric_on,funs_type_apply}.
                        Thus $\rho((f_y,y))\in\reals$ by \cref{metric_on,funs_type_apply}.
                        Let $r_1 = \rho((x,f_x))$.
                        Let $r_2 = \rho((f_x,y))$.
                        Let $r_3 = \rho((f_x,a_0))$.
                        Let $r_4 = \rho((a_0,y))$.
                        Thus $r_1\in\reals$.
                        Thus $r_2\in\reals$.
                        Thus $r_3\in\reals$.
                        Thus $r_4\in\reals$.
                        Thus $r_3 + r_4\in\reals$ by \cref{real_add_closed}.
                        Thus $r_1 + r_2\in\reals$ by \cref{real_add_closed}.
                        Thus $r_1 + (r_3 + r_4)\in\reals$ by \cref{real_add_closed}.
                        Thus $r_2 \leq r_3 + r_4$.
                        Thus $r_2 + r_1 \leq (r_3 + r_4) + r_1$ by \cref{real_leq_add}.
                        Thus $r_1 + r_2 = r_2 + r_1$ by \cref{real_add_comm}.
                        Thus $(r_3 + r_4) + r_1 = r_1 + (r_3 + r_4)$ by \cref{real_add_comm}.
                        Thus $r_1 + r_2 \leq (r_3 + r_4) + r_1$ by \cref{real_leq_trans}.
                        Thus $r_1 + r_2 \leq r_1 + (r_3 + r_4)$ by \cref{real_leq_trans}.
                        Thus $\rho((x,y)) \leq r_1 + (r_3 + r_4)$ by \cref{real_leq_trans}.
                        Let $u = \rho((a_0,f_y)) + \rho((f_y,y))$.
                        Thus $u\in\reals$ by \cref{real_add_closed}.
                        Thus $r_1 + r_3\in\reals$ by \cref{real_add_closed}.
                        Thus $(r_1 + r_3) + u\in\reals$ by \cref{real_add_closed}.
                        Thus $r_4 \leq u$.
                        Thus $r_4 + (r_1 + r_3) \leq u + (r_1 + r_3)$ by \cref{real_leq_add}.
                        Thus $(r_1 + r_3) + r_4 = r_4 + (r_1 + r_3)$ by \cref{real_add_comm}.
                        Thus $u + (r_1 + r_3) = (r_1 + r_3) + u$ by \cref{real_add_comm}.
                        Thus $(r_1 + r_3) + r_4 \leq u + (r_1 + r_3)$ by \cref{real_leq_trans}.
                        Thus $(r_1 + r_3) + r_4 \leq (r_1 + r_3) + u$ by \cref{real_leq_trans}.
                        Thus $r_1 + (r_3 + r_4) = (r_1 + r_3) + r_4$ by \cref{real_add_assoc}.
                        Thus $r_1 + (r_3 + r_4) \leq (r_1 + r_3) + u$ by \cref{real_leq_trans}.
                        Thus $\rho((x,y)) \leq (r_1 + r_3) + u$ by \cref{real_leq_trans}.
                        Thus $\rho((x,y)) \leq \rho((x,f_x)) + \rho((f_x,a_0)) + \rho((a_0,f_y)) + \rho((f_y,y))$
                            by \cref{real_add_assoc,real_add_comm}.
                    \end{subproof}
                    Show that $\rho((x,y)) \leq M$.
                    \begin{subproof}
                        Let $s_1 = \rho((x,f_x))$.
                        Let $s_2 = \rho((f_x,a_0))$.
                        Let $s_3 = \rho((a_0,f_y))$.
                        Let $s_4 = \rho((f_y,y))$.
                        Thus $x\in A$.
                        Thus $y\in A$.
                        Thus $f_x\in F$.
                        Thus $f_y\in F$.
                        Thus $a_0\in F$.
                        Thus $x\in X$ by \cref{elem_subseteq}.
                        Thus $f_x\in A$ by \cref{elem_subseteq}.
                        Thus $f_x\in X$ by \cref{elem_subseteq}.
                        Thus $f_y\in A$ by \cref{elem_subseteq}.
                        Thus $f_y\in X$ by \cref{elem_subseteq}.
                        Thus $y\in X$ by \cref{elem_subseteq}.
                        Thus $a_0\in A$ by \cref{elem_subseteq}.
                        Thus $a_0\in X$ by \cref{elem_subseteq}.
                        $(x,y)\in X\times X$ by \cref{times_tuple_intro}.
                        $(x,f_x)\in X\times X$ by \cref{times_tuple_intro}.
                        $(f_x,a_0)\in X\times X$ by \cref{times_tuple_intro}.
                        $(a_0,f_y)\in X\times X$ by \cref{times_tuple_intro}.
                        $(f_y,y)\in X\times X$ by \cref{times_tuple_intro}.
                        Thus $s_1\in\reals$ by \cref{metric_on,funs_type_apply}.
                        Thus $s_2\in\reals$ by \cref{metric_on,funs_type_apply}.
                        Thus $s_3\in\reals$ by \cref{metric_on,funs_type_apply}.
                        Thus $s_4\in\reals$ by \cref{metric_on,funs_type_apply}.
                        $s_1 + s_2\in\reals$ by \cref{real_add_closed}.
                        $s_3 + s_4\in\reals$ by \cref{real_add_closed}.
                        $1\in\reals$ by \cref{real_naturals_subset_reals,real_one_in_naturals,subseteq}.
                        $\rho$ is symmetric on $X$ by \cref{metric_symmetric,metric_on,square_metric_is_metric}.
                        Thus $\rho((f_x,a_0)) = \rho((a_0,f_x))$ by \cref{metric_symmetric}.
                        Thus $s_2 = \rho((a_0,f_x))$.
                        Thus $s_1 \leq 1$ by \cref{real_lt_imp_leq}.
                        Thus $\rho((f_y,y)) = \rho((y,f_y))$ by \cref{metric_symmetric}.
                        Thus $s_4 = \rho((y,f_y))$.
                        Thus $s_4 < 1$ by \cref{real_lt_subst_left}.
                        Thus $s_4 \leq 1$ by \cref{real_lt_imp_leq}.
                        Thus $s_2 \leq M_0$.
                        Thus $s_3 \leq M_0$.
                        Thus $M_0\in D$ by \cref{real_maximum}.
                        Thus $M_0\in\reals$ by \cref{subseteq}.
                        $M_0 + 1\in\reals$ by \cref{real_add_closed}.
                        $(s_1 + s_2) + (s_3 + s_4)\in\reals$ by \cref{real_add_closed}.
                        $(M_0 + 1) + (M_0 + 1)\in\reals$ by \cref{real_add_closed}.
                        Thus $\rho((x,y))\in\reals$ by \cref{metric_on,funs_type_apply}.
                        Thus $\rho((x,y)) \leq s_1 + s_2 + s_3 + s_4$.
                        Thus $s_1 + s_2 + s_3 + s_4 = (s_1 + s_2) + (s_3 + s_4)$ by \cref{real_add_assoc}.
                        Thus $\rho((x,y)) \leq (s_1 + s_2) + (s_3 + s_4)$ by \cref{real_leq_trans}.
                        Thus $(s_1 + s_2) + (s_3 + s_4) \leq (M_0 + 1) + (M_0 + 1)$ by \cref{real_four_sum_bound}.
                        Thus $\rho((x,y)) \leq (M_0 + 1) + (M_0 + 1)$ by \cref{real_leq_trans}.
                        Thus $\rho((x,y)) \leq M$.
                    \end{subproof}
                \end{subproof}
                Thus $M_0\in D$ by \cref{real_maximum}.
                Thus $M_0\in\reals$ by \cref{subseteq}.
                $1\in\reals$ by \cref{real_naturals_subset_reals,real_one_in_naturals,subseteq}.
                $M_0 + 1\in\reals$ by \cref{real_add_closed}.
                Thus $M\in\reals$ by \cref{real_add_closed}.
                Thus $A$ is bounded in $X$ with respect to $\rho$ by \cref{metric_bounded}.
            \end{byCase}
        \end{subproof}
        Thus $A$ is bounded in $X$ with respect to $d$ or $A$ is bounded in $X$ with respect to $\rho$.
    \end{subproof}
    Show that if
        $A$ is closed in $X$ with respect to $T$ and
        $A$ is bounded in $X$ with respect to $d$ or $A$ is bounded in $X$ with respect to $\rho$
        then $A$ is compact with respect to $T_A$.
    \begin{subproof}
        Assume $A$ is closed in $X$ with respect to $T$.
        Assume $A$ is bounded in $X$ with respect to $d$ or $A$ is bounded in $X$ with respect to $\rho$.
        \begin{byCase}
        \caseOf{$A = \emptyset$.}
            Show that $A$ is compact with respect to $T_A$.
            \begin{subproof}
                Show that for all $C$ such that $C$ is an open covering of $A$ with respect to $T_A$ there exists $B$ such that
                    $B\subseteq C$ and $B$ is finite and $B$ is a covering of $A$.
                \begin{subproof}
                    Fix $C$.
                    Assume $C$ is an open covering of $A$ with respect to $T_A$.
                    Let $B = \emptyset$.
                    $B\subseteq C$.
                    $B$ is finite by \cref{emptyset_is_finite}.
                    $A\subseteq\unions{B}$ by \cref{emptyset_subseteq}.
                    Thus $B$ is a covering of $A$ by \cref{covering}.
                \end{subproof}
                Thus $A$ is compact with respect to $T_A$ by \cref{compact}.
            \end{subproof}
        \caseOf{$A \neq \emptyset$.}
            $A$ is inhabited by \cref{nonempty_inhabited}.
            Show that $A$ is bounded in $X$ with respect to $\rho$.
            \begin{subproof}
                \begin{byCase}
                \caseOf{$A$ is bounded in $X$ with respect to $\rho$.}
                    Thus $A$ is bounded in $X$ with respect to $\rho$.
                \caseOf{$A$ is bounded in $X$ with respect to $d$.}
                    Take $M\in\reals$ such that for all $a_1,a_2\in A$ we have $d((a_1,a_2)) \leq M$
                        by \cref{metric_bounded}.
                    Show that for all $a_1,a_2\in A$ we have $\rho((a_1,a_2)) \leq M$.
                    \begin{subproof}
                        Fix $a_1$.
                        Fix $a_2$.
                        Assume $a_1\in A$.
                        Assume $a_2\in A$.
                        Thus $a_1\in X$ by \cref{elem_subseteq}.
                        Thus $a_2\in X$ by \cref{elem_subseteq}.
                        Take $r\in\reals$ such that $((a_1,a_2),r)\in d$ by \cref{euclidean_metric_value_exists}.
                        Take $m\in\reals$ such that $((a_1,a_2),m)\in \rho$ by \cref{square_metric_value_exists}.
                        Thus $m \leq r$ by \cref{euclidean_square_metric_bounds}.
                        $(a_1,a_2)\in X\times X$ by \cref{times_tuple_intro}.
                        Thus $d((a_1,a_2))\in\reals$ by \cref{metric_on,funs_type_apply,euclidean_metric_is_metric}.
                        Thus $\rho((a_1,a_2))\in\reals$ by \cref{metric_on,funs_type_apply,square_metric_is_metric}.
                        $d$ is a function by \cref{euclidean_metric_function}.
                        Thus $d((a_1,a_2)) = r$ by \cref{function_apply_intro}.
                        Thus $r = d((a_1,a_2))$.
                        Thus $r \leq d((a_1,a_2))$.
                        Thus $d((a_1,a_2)) \leq M$.
                        Thus $r \leq M$ by \cref{real_leq_trans}.
                        Thus $m \leq M$ by \cref{real_leq_trans}.
                        Thus $\rho((a_1,a_2)) = m$ by \cref{square_metric_function,function_apply_intro}.
                        Thus $\rho((a_1,a_2)) \leq m$.
                        Thus $\rho((a_1,a_2)) \leq M$ by \cref{real_leq_trans}.
                    \end{subproof}
                    $\rho$ is a metric on $X$ by \cref{square_metric_is_metric}.
                    Thus $A$ is bounded in $X$ with respect to $\rho$ by \cref{metric_bounded}.
                \end{byCase}
            \end{subproof}
            Take $M\in\reals$ such that for all $a_1,a_2\in A$ we have $\rho((a_1,a_2)) \leq M$
                by \cref{metric_bounded}.
            Take $x_0$ such that $x_0\in A$.
            Let $Y = \{(i,\intervalclosed{x_0(i) - M}{x_0(i) + M}) \mid i\in J\}$.
            Show that $A\subseteq \productFamily{Y}$.
            \begin{subproof}
                Show that for all $x$ such that $x\in A$ we have $x\in \productFamily{Y}$.
                \begin{subproof}
                    Fix $x$.
                    Assume $x\in A$.
                    Show that for all $i\in J$ we have $x(i)\in \intervalclosed{x_0(i) - M}{x_0(i) + M}$.
                    \begin{subproof}
                        Fix $i$.
                        Assume $i\in J$.
                        Thus $x_0\in X$ by \cref{elem_subseteq}.
                        Thus $x\in X$ by \cref{elem_subseteq}.
                        Thus $x_0\in\realsPower{J}$.
                        Thus $x\in\realsPower{J}$.
                        Thus $x_0\in\tuplePower{\reals}{J}$ by \cref{reals_power}.
                        Thus $x\in\tuplePower{\reals}{J}$ by \cref{reals_power}.
                        Thus $x_0\in\funs{J}{\reals}$ by \cref{tuple_power}.
                        Thus $x\in\funs{J}{\reals}$ by \cref{tuple_power}.
                        Thus $x_0(i)\in\reals$ by \cref{funs_type_apply}.
                        Thus $x(i)\in\reals$ by \cref{funs_type_apply}.
                        $\neg{x(i)}\in\reals$ by \cref{real_add_inverse}.
                        $x_0(i) - x(i) = x_0(i) + \neg{x(i)}$ by \cref{real_sub}.
                        Thus $x_0(i) - x(i)\in\reals$ by \cref{real_add_closed}.
                        Thus $\abs{x_0(i) - x(i)}\in\reals$ by \cref{real_abs_in_reals}.
                        Take $m\in\reals$ such that $((x_0,x),m)\in \rho$ by \cref{square_metric_value_exists}.
                        Thus $m\in\coordDiffSet{n}{x_0}{x}$ and for all $t\in\coordDiffSet{n}{x_0}{x}$ we have $t \leq m$
                            by \cref{square_metric,pair_eq_iff}.
                        Thus $\abs{x_0(i) - x(i)} \leq m$ by \cref{coord_diff_set}.
                        $\rho$ is a function by \cref{square_metric_function}.
                        Thus $\rho((x_0,x)) = m$ by \cref{function_apply_intro}.
                        Thus $m = \rho((x_0,x))$.
                        Thus $m \leq \rho((x_0,x))$.
                        Thus $\rho((x_0,x)) \leq M$.
                        Thus $m \leq M$ by \cref{real_leq_trans}.
                        Thus $\abs{x_0(i) - x(i)} \leq M$ by \cref{real_leq_trans}.
                        Thus $\abs{x(i) - x_0(i)} = \abs{x_0(i) - x(i)}$ by \cref{real_abs_sub_swap}.
                        Thus $\abs{x(i) - x_0(i)} \leq \abs{x_0(i) - x(i)}$.
                        Thus $\abs{x(i) - x_0(i)} \leq M$ by \cref{real_leq_trans}.
                        $\abs{x(i) - x_0(i)} \geq 0$ by \cref{real_abs_nonnegative}.
                        $0\in\reals$ by \cref{real_add_ident}.
                        Thus $0 \leq \abs{x(i) - x_0(i)}$.
                        Thus $0 \leq M$ by \cref{real_leq_trans}.
                        Thus $M$ is nonnegative.
                        Thus $x_0(i) - M \leq x(i)$ by \cref{real_abs_sub_leq_imp_left_bound}.
                        Thus $x(i) \leq x_0(i) + M$ by \cref{real_abs_sub_leq_imp_right_bound}.
                        Thus $x(i)\in\intervalclosed{x_0(i) - M}{x_0(i) + M}$ by \cref{intervalclosed}.
                    \end{subproof}
                    Show that $x\in\productFamily{Y}$.
                    \begin{subproof}
                        Show that $x\in\funs{\dom{Y}}{\unions{\ran{Y}}}$.
                        \begin{subproof}
                            Thus $x\in\realsPower{J}$.
                            Thus $x\in\tuplePower{\reals}{J}$ by \cref{reals_power}.
                            Thus $x\in\funs{J}{\reals}$ by \cref{tuple_power}.
                            Thus $x$ is a function to $\reals$ and $J = \dom{x}$ by \cref{funs_elim}.
                            Show that for all $i$ such that $i\in J$ we have $i\in\dom{Y}$.
                            \begin{subproof}
                                Fix $i$.
                                Assume $i\in J$.
                                Thus $(i,\intervalclosed{x_0(i) - M}{x_0(i) + M})\in Y$.
                                Thus $i\in\dom{Y}$ by \cref{dom_intro}.
                            \end{subproof}
                            Thus $J\subseteq\dom{Y}$ by \cref{subseteq}.
                            Show that for all $i$ such that $i\in\dom{Y}$ we have $i\in J$.
                            \begin{subproof}
                                Fix $i$.
                                Assume $i\in\dom{Y}$.
                                Take $v$ such that $(i,v)\in Y$ by \cref{dom_iff}.
                                Take $j\in J$ such that $(i,v) = (j,\intervalclosed{x_0(j) - M}{x_0(j) + M})$.
                                Thus $i = j$ by \cref{pair_eq_iff}.
                                Thus $i\in J$.
                            \end{subproof}
                            Thus $\dom{Y}\subseteq J$ by \cref{subseteq}.
                            Thus $\dom{Y} = J$ by \cref{subseteq_antisymmetric}.
                            Thus $\dom{x} = \dom{Y}$.
                            Show that for all $i\in\dom{Y}$ we have $x(i)\in\unions{\ran{Y}}$.
                            \begin{subproof}
                                Fix $i$.
                                Assume $i\in\dom{Y}$.
                                Thus $i\in J$ by \cref{subseteq}.
                                Thus $x(i)\in\intervalclosed{x_0(i) - M}{x_0(i) + M}$.
                                Thus $(i,\intervalclosed{x_0(i) - M}{x_0(i) + M})\in Y$.
                                Thus $\intervalclosed{x_0(i) - M}{x_0(i) + M}\in\ran{Y}$ by \cref{ran_intro}.
                                Thus $x(i)\in\unions{\ran{Y}}$ by \cref{unions_iff}.
                            \end{subproof}
                            Thus $x$ is a function to $\unions{\ran{Y}}$.
                            Thus $x\in\funs{\dom{Y}}{\unions{\ran{Y}}}$ by \cref{funs_intro}.
                        \end{subproof}
                        Show that for all $i\in\dom{Y}$ we have $x(i)\in\apply{Y}{i}$.
                        \begin{subproof}
                            Fix $i$.
                            Assume $i\in\dom{Y}$.
                            Take $v$ such that $(i,v)\in Y$ by \cref{dom_iff}.
                            Take $j\in J$ such that $(i,v) = (j,\intervalclosed{x_0(j) - M}{x_0(j) + M})$.
                            Thus $i = j$ by \cref{pair_eq_iff}.
                            Thus $i\in J$.
                            Thus $x(i)\in\intervalclosed{x_0(i) - M}{x_0(i) + M}$.
                            Thus $(i,\intervalclosed{x_0(i) - M}{x_0(i) + M})\in Y$.
                            Show that $Y$ is a function.
                            \begin{subproof}
                                Show that for all $i,u,v$ such that $(i,u)\in Y$ and $(i,v)\in Y$ we have $u = v$.
                                \begin{subproof}
                                    Fix $i$.
                                    Fix $u$.
                                    Fix $v$.
                                    Assume $(i,u)\in Y$.
                                    Assume $(i,v)\in Y$.
                                    Take $i_1\in J$ such that $(i,u) = (i_1,\intervalclosed{x_0(i_1) - M}{x_0(i_1) + M})$.
                                    Take $i_2\in J$ such that $(i,v) = (i_2,\intervalclosed{x_0(i_2) - M}{x_0(i_2) + M})$.
                                    Thus $i = i_1$ and $u = \intervalclosed{x_0(i_1) - M}{x_0(i_1) + M}$ by \cref{pair_eq_iff}.
                                    Thus $i = i_2$ and $v = \intervalclosed{x_0(i_2) - M}{x_0(i_2) + M}$ by \cref{pair_eq_iff}.
                                    Thus $u = v$.
                                \end{subproof}
                                Thus $Y$ is right-unique by \cref{rightunique}.
                                Show that for all $w$ such that $w\in Y$ there exist $x,y$ such that $w = (x,y)$.
                                \begin{subproof}
                                    Fix $w$.
                                    Assume $w\in Y$.
                                    Take $i\in J$ such that $w = (i,\intervalclosed{x_0(i) - M}{x_0(i) + M})$.
                                    Thus there exist $x,y$ such that $w = (x,y)$.
                                \end{subproof}
                                Thus $Y$ is a relation by \cref{relation}.
                                Thus $Y$ is a function.
                            \end{subproof}
                            Thus $Y(i) = \intervalclosed{x_0(i) - M}{x_0(i) + M}$ by \cref{function_apply_intro}.
                            Thus $x(i)\in\apply{Y}{i}$.
                        \end{subproof}
                        Thus $x\in\productFamily{Y}$ by \cref{indexed_product}.
                    \end{subproof}
                \end{subproof}
                Thus $A\subseteq \productFamily{Y}$ by \cref{subseteq}.
            \end{subproof}
            Let $T_Y = \{(i,\subspaceTopology{\standardTopologyR}{\intervalclosed{x_0(i) - M}{x_0(i) + M}}) \mid i\in J\}$.
            Show that $Y$ is a function.
            \begin{subproof}
                Show that $Y$ is a relation.
                \begin{subproof}
                    Show that for all $w$ such that $w\in Y$ there exist $x,y$ such that $w = (x,y)$.
                    \begin{subproof}
                        Fix $w$.
                        Assume $w\in Y$.
                        Take $i\in J$ such that $w = (i,\intervalclosed{x_0(i) - M}{x_0(i) + M})$.
                        Thus there exist $x,y$ such that $w = (x,y)$.
                    \end{subproof}
                    Thus $Y$ is a relation by \cref{relation}.
                \end{subproof}
                Show that for all $i,u,v$ such that $(i,u)\in Y$ and $(i,v)\in Y$ we have $u = v$.
                \begin{subproof}
                    Fix $i$.
                    Fix $u$.
                    Fix $v$.
                    Assume $(i,u)\in Y$.
                    Assume $(i,v)\in Y$.
                    Take $i_1\in J$ such that $(i,u) = (i_1,\intervalclosed{x_0(i_1) - M}{x_0(i_1) + M})$.
                    Take $i_2\in J$ such that $(i,v) = (i_2,\intervalclosed{x_0(i_2) - M}{x_0(i_2) + M})$.
                    Thus $i = i_1$ and $u = \intervalclosed{x_0(i_1) - M}{x_0(i_1) + M}$ by \cref{pair_eq_iff}.
                    Thus $i = i_2$ and $v = \intervalclosed{x_0(i_2) - M}{x_0(i_2) + M}$ by \cref{pair_eq_iff}.
                    Thus $u = v$.
                \end{subproof}
                Thus $Y$ is right-unique by \cref{rightunique}.
                Thus $Y$ is a function.
            \end{subproof}
            Show that $T_Y$ is a function.
            \begin{subproof}
                Show that $T_Y$ is a relation.
                \begin{subproof}
                    Show that for all $w$ such that $w\in T_Y$ there exist $x,y$ such that $w = (x,y)$.
                    \begin{subproof}
                        Fix $w$.
                        Assume $w\in T_Y$.
                        Take $i\in J$ such that $w = (i,\subspaceTopology{\standardTopologyR}{\intervalclosed{x_0(i) - M}{x_0(i) + M}})$.
                        Thus there exist $x,y$ such that $w = (x,y)$.
                    \end{subproof}
                    Thus $T_Y$ is a relation by \cref{relation}.
                \end{subproof}
                Show that for all $i,u,v$ such that $(i,u)\in T_Y$ and $(i,v)\in T_Y$ we have $u = v$.
                \begin{subproof}
                    Fix $i$.
                    Fix $u$.
                    Fix $v$.
                    Assume $(i,u)\in T_Y$.
                    Assume $(i,v)\in T_Y$.
                    Take $i_1\in J$ such that $(i,u) = (i_1,\subspaceTopology{\standardTopologyR}{\intervalclosed{x_0(i_1) - M}{x_0(i_1) + M}})$.
                    Take $i_2\in J$ such that $(i,v) = (i_2,\subspaceTopology{\standardTopologyR}{\intervalclosed{x_0(i_2) - M}{x_0(i_2) + M}})$.
                    Thus $i = i_1$ and $u = \subspaceTopology{\standardTopologyR}{\intervalclosed{x_0(i_1) - M}{x_0(i_1) + M}}$ by \cref{pair_eq_iff}.
                    Thus $i = i_2$ and $v = \subspaceTopology{\standardTopologyR}{\intervalclosed{x_0(i_2) - M}{x_0(i_2) + M}}$ by \cref{pair_eq_iff}.
                    Thus $u = v$.
                \end{subproof}
                Thus $T_Y$ is right-unique by \cref{rightunique}.
                Thus $T_Y$ is a function.
            \end{subproof}
            Show that $\dom{Y} = J$.
            \begin{subproof}
                Show that for all $i$ such that $i\in J$ we have $i\in\dom{Y}$.
                \begin{subproof}
                    Fix $i$.
                    Assume $i\in J$.
                    Thus $(i,\intervalclosed{x_0(i) - M}{x_0(i) + M})\in Y$.
                    Thus $i\in\dom{Y}$ by \cref{dom_intro}.
                \end{subproof}
                Thus $J\subseteq\dom{Y}$ by \cref{subseteq}.
                Show that for all $i$ such that $i\in\dom{Y}$ we have $i\in J$.
                \begin{subproof}
                    Fix $i$.
                    Assume $i\in\dom{Y}$.
                    Take $v$ such that $(i,v)\in Y$ by \cref{dom_iff}.
                    Take $j\in J$ such that $(i,v) = (j,\intervalclosed{x_0(j) - M}{x_0(j) + M})$.
                    Thus $i = j$ by \cref{pair_eq_iff}.
                    Thus $i\in J$.
                \end{subproof}
                Thus $\dom{Y}\subseteq J$ by \cref{subseteq}.
                Thus $\dom{Y} = J$ by \cref{subseteq_antisymmetric}.
            \end{subproof}
            Show that $\dom{T_Y} = J$.
            \begin{subproof}
                Show that for all $i$ such that $i\in J$ we have $i\in\dom{T_Y}$.
                \begin{subproof}
                    Fix $i$.
                    Assume $i\in J$.
                    Thus $(i,\subspaceTopology{\standardTopologyR}{\intervalclosed{x_0(i) - M}{x_0(i) + M}})\in T_Y$.
                    Thus $i\in\dom{T_Y}$ by \cref{dom_intro}.
                \end{subproof}
                Thus $J\subseteq\dom{T_Y}$ by \cref{subseteq}.
                Show that for all $i$ such that $i\in\dom{T_Y}$ we have $i\in J$.
                \begin{subproof}
                    Fix $i$.
                    Assume $i\in\dom{T_Y}$.
                    Take $v$ such that $(i,v)\in T_Y$ by \cref{dom_iff}.
                    Take $j\in J$ such that $(i,v) = (j,\subspaceTopology{\standardTopologyR}{\intervalclosed{x_0(j) - M}{x_0(j) + M}})$.
                    Thus $i = j$ by \cref{pair_eq_iff}.
                    Thus $i\in J$.
                \end{subproof}
                Thus $\dom{T_Y}\subseteq J$ by \cref{subseteq}.
                Thus $\dom{T_Y} = J$ by \cref{subseteq_antisymmetric}.
            \end{subproof}
            Show that $\productFamily{Y}$ is compact with respect to $\productTopologyIndexed{T_Y}{Y}$.
            \begin{subproof}
                % product of finitely many compact intervals
                Show that for all $i\in J$ we have $\intervalclosed{x_0(i) - M}{x_0(i) + M}$ is compact
                    with respect to $\subspaceTopology{\standardTopologyR}{\intervalclosed{x_0(i) - M}{x_0(i) + M}}$.
                \begin{subproof}
                    Fix $i$.
                    Assume $i\in J$.
                    Thus $x_0\in X$ by \cref{elem_subseteq}.
                    Thus $x_0\in\realsPower{J}$.
                    Thus $x_0\in\tuplePower{\reals}{J}$ by \cref{reals_power}.
                    Thus $x_0\in\funs{J}{\reals}$ by \cref{tuple_power}.
                    Thus $x_0(i)\in\reals$ by \cref{funs_type_apply}.
                    Thus $M\in\reals$.
                    $0\in\reals$ by \cref{real_add_ident}.
                    $\rho$ is a metric on $X$ by \cref{square_metric_is_metric}.
                    Thus $\rho$ is nonnegative on $X$ by \cref{metric_on}.
                    Thus $x_0\in A$.
                    Thus $x_0\in X$ by \cref{elem_subseteq}.
                    $(x_0,x_0)\in X\times X$ by \cref{times_tuple_intro}.
                    Thus $\rho((x_0,x_0))\in\reals$ by \cref{metric_on,funs_type_apply}.
                    Thus $\rho((x_0,x_0)) \geq 0$ by \cref{metric_nonnegative}.
                    Thus $\rho((x_0,x_0)) \leq M$.
                    Thus $0 \leq M$ by \cref{real_leq_trans}.
                    Thus $M$ is nonnegative.
                    $\neg{M}\in\reals$ by \cref{real_add_inverse}.
                    $\neg{0}\in\reals$ by \cref{real_add_inverse}.
                    $0 + \neg{0} = 0$ by \cref{real_add_inverse}.
                    $0 + \neg{0} = \neg{0}$ by \cref{real_add_ident}.
                    Thus $\neg{0} = 0$.
                    Thus $\neg{M} \leq \neg{0}$ by \cref{real_leq_neg}.
                    Thus $\neg{M} \leq 0$ by \cref{real_leq_trans}.
                    Thus $\neg{M} \leq M$ by \cref{real_leq_trans}.
                    Thus $\neg{M} + x_0(i) \leq M + x_0(i)$ by \cref{real_leq_add}.
                    $x_0(i) + \neg{M} = \neg{M} + x_0(i)$ by \cref{real_add_comm}.
                    $x_0(i) + M = M + x_0(i)$ by \cref{real_add_comm}.
                    Thus $x_0(i) + \neg{M} \leq x_0(i) + M$ by \cref{real_leq_trans}.
                    Thus $x_0(i) + M\in\reals$ by \cref{real_add_closed}.
                    $x_0(i) - M = x_0(i) + \neg{M}$ by \cref{real_sub}.
                    Thus $x_0(i) - M\in\reals$ by \cref{real_add_closed}.
                    Thus $x_0(i) - M \leq x_0(i) + M$ by \cref{real_sub}.
                    Let $I = \intervalclosed{x_0(i) - M}{x_0(i) + M}$.
                    Thus $I$ is compact with respect to $\subspaceTopology{\standardTopologyR}{I}$
                        by \cref{real_intervalclosed_compact_standard}.
                    Thus $\intervalclosed{x_0(i) - M}{x_0(i) + M}$ is compact
                        with respect to $\subspaceTopology{\standardTopologyR}{\intervalclosed{x_0(i) - M}{x_0(i) + M}}$.
                \end{subproof}
                Show that for all $i\in J$ we have $\apply{T_Y}{i}$ is a topology on $\apply{Y}{i}$.
                \begin{subproof}
                    Fix $i$.
                    Assume $i\in J$.
                    Thus $(i,\intervalclosed{x_0(i) - M}{x_0(i) + M})\in Y$.
                    Thus $(i,\subspaceTopology{\standardTopologyR}{\intervalclosed{x_0(i) - M}{x_0(i) + M}})\in T_Y$.
                    Thus $\apply{Y}{i} = \intervalclosed{x_0(i) - M}{x_0(i) + M}$ by \cref{function_apply_intro}.
                    Thus $\apply{T_Y}{i} = \subspaceTopology{\standardTopologyR}{\intervalclosed{x_0(i) - M}{x_0(i) + M}}$ by \cref{function_apply_intro}.
                    $\standardBasisR$ is a basis on $\reals$ by \cref{standard_basis_is_basis}.
                    Thus $\standardTopologyR$ is a topology on $\reals$ by \cref{basis_topology_is_topology,standard_topology_reals}.
                    Show that $\intervalclosed{x_0(i) - M}{x_0(i) + M}\subseteq\reals$.
                    \begin{subproof}
                        Show that for all $x$ such that $x\in\intervalclosed{x_0(i) - M}{x_0(i) + M}$ we have $x\in\reals$.
                        \begin{subproof}
                            Fix $x$.
                            Assume $x\in\intervalclosed{x_0(i) - M}{x_0(i) + M}$.
                            Thus $x\in\reals$ by \cref{intervalclosed}.
                        \end{subproof}
                        Thus $\intervalclosed{x_0(i) - M}{x_0(i) + M}\subseteq\reals$ by \cref{subseteq}.
                    \end{subproof}
                    Thus $\apply{T_Y}{i}$ is a topology on $\apply{Y}{i}$ by \cref{subspace_topology_is_topology}.
                \end{subproof}
                Show that for all $i\in J$ we have $\apply{Y}{i}$ is compact with respect to $\apply{T_Y}{i}$.
                \begin{subproof}
                    Fix $i$.
                    Assume $i\in J$.
                    Thus $\intervalclosed{x_0(i) - M}{x_0(i) + M}$ is compact
                        with respect to $\subspaceTopology{\standardTopologyR}{\intervalclosed{x_0(i) - M}{x_0(i) + M}}$.
                    Thus $(i,\intervalclosed{x_0(i) - M}{x_0(i) + M})\in Y$.
                    Thus $(i,\subspaceTopology{\standardTopologyR}{\intervalclosed{x_0(i) - M}{x_0(i) + M}})\in T_Y$.
                    Thus $\apply{Y}{i} = \intervalclosed{x_0(i) - M}{x_0(i) + M}$ by \cref{function_apply_intro}.
                    Thus $\apply{T_Y}{i} = \subspaceTopology{\standardTopologyR}{\intervalclosed{x_0(i) - M}{x_0(i) + M}}$ by \cref{function_apply_intro}.
                    Thus $\apply{Y}{i}$ is compact with respect to $\apply{T_Y}{i}$.
                \end{subproof}
                Thus $\productFamily{Y}$ is compact with respect to $\productTopologyIndexed{T_Y}{Y}$ by \cref{compact_indexed_product_initial_segment}.
            \end{subproof}
            Let $T_{P0} = \productTopologyIndexed{T_Y}{Y}$.
            Let $T_{Y0} = \subspaceTopology{T}{\productFamily{Y}}$.
            Show that $T_{Y0}\subseteq T_{P0}$.
            \begin{subproof}
                Let $X_1 = \{(i,\reals) \mid i\in J\}$.
                Let $T_1 = \{(i,\standardTopologyR) \mid i\in J\}$.
                Let $T_0 = \productTopologyIndexed{T_1}{X_1}$.
                Let $T_E = \metricTopology{d}{X}$.
                Thus $T_E = T_0$ by \cref{euclidean_metric_product_topology}.
                Thus $T = T_E$ by \cref{square_metric_topology_equals_euclidean}.
                Thus $T = T_0$.
                Show that $X_1$ is a function.
                \begin{subproof}
                    Show that for all $w$ such that $w\in X_1$ there exist $x,y$ such that $w = (x,y)$.
                    \begin{subproof}
                        Fix $w$.
                        Assume $w\in X_1$.
                        Take $i\in J$ such that $w = (i,\reals)$.
                        Thus there exist $x,y$ such that $w = (x,y)$.
                    \end{subproof}
                    Thus $X_1$ is a relation by \cref{relation}.
                    Show that for all $i,y_1,y_2$ such that $(i,y_1)\in X_1$ and $(i,y_2)\in X_1$ we have $y_1 = y_2$.
                    \begin{subproof}
                        Fix $i$.
                        Fix $y_1$.
                        Fix $y_2$.
                        Assume $(i,y_1)\in X_1$ and $(i,y_2)\in X_1$.
                        Then $y_1 = \reals$ and $y_2 = \reals$.
                        Thus $y_1 = y_2$.
                    \end{subproof}
                    Thus $X_1$ is right-unique by \cref{rightunique}.
                    Thus $X_1$ is a function.
                \end{subproof}
                Show that $T_1$ is a function.
                \begin{subproof}
                    Show that for all $w$ such that $w\in T_1$ there exist $x,y$ such that $w = (x,y)$.
                    \begin{subproof}
                        Fix $w$.
                        Assume $w\in T_1$.
                        Take $i\in J$ such that $w = (i,\standardTopologyR)$.
                        Thus there exist $x,y$ such that $w = (x,y)$.
                    \end{subproof}
                    Thus $T_1$ is a relation by \cref{relation}.
                    Show that for all $i,y_1,y_2$ such that $(i,y_1)\in T_1$ and $(i,y_2)\in T_1$ we have $y_1 = y_2$.
                    \begin{subproof}
                        Fix $i$.
                        Fix $y_1$.
                        Fix $y_2$.
                        Assume $(i,y_1)\in T_1$ and $(i,y_2)\in T_1$.
                        Then $y_1 = \standardTopologyR$ and $y_2 = \standardTopologyR$.
                        Thus $y_1 = y_2$.
                    \end{subproof}
                    Thus $T_1$ is right-unique by \cref{rightunique}.
                    Thus $T_1$ is a function.
                \end{subproof}
                Show that $\dom{X_1} = J$.
                \begin{subproof}
                    Show that for all $i$ such that $i\in J$ we have $i\in\dom{X_1}$.
                    \begin{subproof}
                        Fix $i$.
                        Assume $i\in J$.
                        Thus $(i,\reals)\in X_1$.
                        Thus $i\in\dom{X_1}$ by \cref{dom_intro}.
                    \end{subproof}
                    Thus $J\subseteq\dom{X_1}$ by \cref{subseteq}.
                    Show that for all $i$ such that $i\in\dom{X_1}$ we have $i\in J$.
                    \begin{subproof}
                        Fix $i$.
                        Assume $i\in\dom{X_1}$.
                        Take $y$ such that $(i,y)\in X_1$ by \cref{dom_iff}.
                        Thus $i\in J$.
                    \end{subproof}
                    Thus $\dom{X_1}\subseteq J$ by \cref{subseteq}.
                    Thus $\dom{X_1} = J$ by \cref{subseteq_antisymmetric}.
                \end{subproof}
                Show that $\dom{T_1} = J$.
                \begin{subproof}
                    Show that for all $i$ such that $i\in J$ we have $i\in\dom{T_1}$.
                    \begin{subproof}
                        Fix $i$.
                        Assume $i\in J$.
                        Thus $(i,\standardTopologyR)\in T_1$.
                        Thus $i\in\dom{T_1}$ by \cref{dom_intro}.
                    \end{subproof}
                    Thus $J\subseteq\dom{T_1}$ by \cref{subseteq}.
                    Show that for all $i$ such that $i\in\dom{T_1}$ we have $i\in J$.
                    \begin{subproof}
                        Fix $i$.
                        Assume $i\in\dom{T_1}$.
                        Take $y$ such that $(i,y)\in T_1$ by \cref{dom_iff}.
                        Thus $i\in J$.
                    \end{subproof}
                    Thus $\dom{T_1}\subseteq J$ by \cref{subseteq}.
                    Thus $\dom{T_1} = J$ by \cref{subseteq_antisymmetric}.
                \end{subproof}
                Thus $\dom{T_1} = \dom{X_1}$.
                Show that $\dom{X_1}$ is inhabited.
                \begin{subproof}
                    $1\in\naturalsPlus$ by \cref{real_one_in_naturals}.
                    $1 \leq n$ by \cref{real_one_leq_naturalsplus}.
                    Thus $1\in J$ by \cref{initial_segment_naturalsplus}.
                    Thus $1\in\dom{X_1}$ by \cref{subseteq}.
                    Thus $\dom{X_1}$ is inhabited by \cref{nonempty_inhabited}.
                \end{subproof}
                Show that for all $i\in\dom{X_1}$ we have $\apply{Y}{i}\subseteq\apply{X_1}{i}$.
                \begin{subproof}
                    Fix $i$.
                    Assume $i\in\dom{X_1}$.
                    Thus $i\in J$ by \cref{subseteq}.
                    Thus $(i,\intervalclosed{x_0(i) - M}{x_0(i) + M})\in Y$.
                    Thus $(i,\reals)\in X_1$.
                    Thus $\apply{Y}{i} = \intervalclosed{x_0(i) - M}{x_0(i) + M}$ by \cref{function_apply_intro}.
                    Thus $\apply{X_1}{i} = \reals$ by \cref{function_apply_intro}.
                    Show that $\intervalclosed{x_0(i) - M}{x_0(i) + M}\subseteq\reals$.
                    \begin{subproof}
                        Show that for all $x$ such that $x\in\intervalclosed{x_0(i) - M}{x_0(i) + M}$ we have $x\in\reals$.
                        \begin{subproof}
                            Fix $x$.
                            Assume $x\in\intervalclosed{x_0(i) - M}{x_0(i) + M}$.
                            Thus $x\in\reals$ by \cref{intervalclosed}.
                        \end{subproof}
                        Thus $\intervalclosed{x_0(i) - M}{x_0(i) + M}\subseteq\reals$ by \cref{subseteq}.
                    \end{subproof}
                    Thus $\apply{Y}{i}\subseteq\apply{X_1}{i}$ by \cref{subseteq_transitive}.
                \end{subproof}
                Show that for all $i\in\dom{X_1}$ we have $\apply{T_1}{i}$ is a topology on $\apply{X_1}{i}$.
                \begin{subproof}
                    Fix $i$.
                    Assume $i\in\dom{X_1}$.
                    Thus $i\in J$ by \cref{subseteq}.
                    Thus $(i,\reals)\in X_1$.
                    Thus $(i,\standardTopologyR)\in T_1$.
                    Thus $\apply{X_1}{i} = \reals$ by \cref{function_apply_intro}.
                    Thus $\apply{T_1}{i} = \standardTopologyR$ by \cref{function_apply_intro}.
                    $\standardBasisR$ is a basis on $\reals$ by \cref{standard_basis_is_basis}.
                    Thus $\standardTopologyR$ is a topology on $\reals$
                        by \cref{basis_topology_is_topology,standard_topology_reals}.
                    Thus $\apply{T_1}{i}$ is a topology on $\apply{X_1}{i}$.
                \end{subproof}
                Thus $\productFamily{Y}$ is a subspace of $\productFamily{X_1}$ with respect to $T$ by \cref{product_subspace_product}.
                Show that $j$ is continuous from $\productFamily{Y}$ to $X$ with respect to $T_{P0}$ and $T$.
                \begin{subproof}
                    Let $j = \identity{\productFamily{Y}}$.
                    Thus $Y$ is a function.
                    Thus $T_Y$ is a function.
                    Thus $\dom{T_Y} = \dom{Y}$.
                    Show that $\dom{Y}$ is inhabited.
                    \begin{subproof}
                        $1\in\naturalsPlus$ by \cref{real_one_in_naturals}.
                        $1 \leq n$ by \cref{real_one_leq_naturalsplus}.
                        Thus $1\in J$ by \cref{initial_segment_naturalsplus}.
                        Thus $1\in\dom{Y}$ by \cref{subseteq}.
                        Thus $\dom{Y}$ is inhabited by \cref{nonempty_inhabited}.
                    \end{subproof}
                    Show that for all $i\in\dom{Y}$ we have $\apply{T_Y}{i}$ is a topology on $\apply{Y}{i}$.
                    \begin{subproof}
                        Fix $i$.
                        Assume $i\in\dom{Y}$.
                        Thus $i\in J$ by \cref{subseteq}.
                        Thus $(i,\intervalclosed{x_0(i) - M}{x_0(i) + M})\in Y$.
                        Thus $(i,\subspaceTopology{\standardTopologyR}{\intervalclosed{x_0(i) - M}{x_0(i) + M}})\in T_Y$.
                        Thus $\apply{Y}{i} = \intervalclosed{x_0(i) - M}{x_0(i) + M}$ by \cref{function_apply_intro}.
                        Thus $\apply{T_Y}{i} = \subspaceTopology{\standardTopologyR}{\intervalclosed{x_0(i) - M}{x_0(i) + M}}$ by \cref{function_apply_intro}.
                        $\standardBasisR$ is a basis on $\reals$ by \cref{standard_basis_is_basis}.
                        Thus $\standardTopologyR$ is a topology on $\reals$ by \cref{basis_topology_is_topology,standard_topology_reals}.
                        Show that $\intervalclosed{x_0(i) - M}{x_0(i) + M}\subseteq\reals$.
                        \begin{subproof}
                            Show that for all $x$ such that $x\in\intervalclosed{x_0(i) - M}{x_0(i) + M}$ we have $x\in\reals$.
                            \begin{subproof}
                                Fix $x$.
                                Assume $x\in\intervalclosed{x_0(i) - M}{x_0(i) + M}$.
                                Thus $x\in\reals$ by \cref{intervalclosed}.
                            \end{subproof}
                            Thus $\intervalclosed{x_0(i) - M}{x_0(i) + M}\subseteq\reals$ by \cref{subseteq}.
                        \end{subproof}
                        Thus $\apply{T_Y}{i}$ is a topology on $\apply{Y}{i}$ by \cref{subspace_topology_is_topology}.
                    \end{subproof}
                    Thus $T_{P0}$ is a topology on $\productFamily{Y}$
                        by \cref{product_subbasis_finite_intersections_basis,basis_for_topology,basis_topology_is_topology}.
                    Thus $\productFamily{X_1} = X$ by \cref{product_family_reals_power}.
                    Thus $\productFamily{Y}\subseteq X$ by \cref{subspace_of,subseteq_transitive}.
                    $j\in\funs{\productFamily{Y}}{\productFamily{Y}}$ by \cref{id_elem_funs}.
                    Thus $j\in\funs{\productFamily{Y}}{X}$ by \cref{funs_weaken_codom}.
                    Show that for all $i\in\dom{X_1}$ we have $\piIndex{X_1}{i}\circ j$ is continuous from $\productFamily{Y}$
                        to $\apply{X_1}{i}$ with respect to $T_{P0}$ and $\apply{T_1}{i}$.
                    \begin{subproof}
                        Fix $i$.
                        Assume $i\in\dom{X_1}$.
                        Thus $i\in J$ by \cref{subseteq}.
                        Let $g = \piIndex{Y}{i}$.
                        Let $k = \identity{\apply{Y}{i}}$.
                        Show that $k$ is continuous from $\apply{Y}{i}$ to $\apply{X_1}{i}$ with respect to $\apply{T_Y}{i}$ and $\apply{T_1}{i}$.
                        \begin{subproof}
                            Thus $\apply{Y}{i}\subseteq\apply{X_1}{i}$ by \cref{subseteq}.
                            Thus $\apply{T_1}{i}$ is a topology on $\apply{X_1}{i}$ by \cref{subseteq}.
                            Thus $\apply{Y}{i}$ is a subspace of $\apply{X_1}{i}$ with respect to $\apply{T_1}{i}$ by \cref{subspace_of}.
                            Thus $\apply{T_Y}{i} = \subspaceTopology{\apply{T_1}{i}}{\apply{Y}{i}}$ by \cref{function_apply_intro}.
                            Thus $k$ is continuous from $\apply{Y}{i}$ to $\apply{X_1}{i}$ with respect to $\apply{T_Y}{i}$ and $\apply{T_1}{i}$
                                by \cref{inclusion_continuous}.
                        \end{subproof}
                        $g$ is continuous from $\productFamily{Y}$ to $\apply{Y}{i}$ with respect to $T_{P0}$ and $\apply{T_Y}{i}$
                            by \cref{projection_map_indexed_continuous}.
                        Thus $k\circ g$ is continuous from $\productFamily{Y}$ to $\apply{X_1}{i}$ with respect to $T_{P0}$ and $\apply{T_1}{i}$
                            by \cref{continuous_composite}.
                        Show that $k\circ g = \piIndex{X_1}{i}\circ j$.
                        \begin{subproof}
                            Show that for all $y,z$ if $y\mathrel{k\circ g} z$ then $y\mathrel{\piIndex{X_1}{i}\circ j} z$.
                            \begin{subproof}
                                Fix $y$.
                                Fix $z$.
                                Assume $y\mathrel{k\circ g} z$.
                                Take $x$ such that $y\mathrel{g} x$ and $x\mathrel{k} z$ by \cref{circ_iff}.
                                Thus $x = z$ and $x\in\apply{Y}{i}$ by \cref{id_iff}.
                                Take $y_0\in\productFamily{Y}$ such that $(y,x) = (y_0,\apply{y_0}{i})$ by \cref{projection_map_indexed}.
                                Thus $y = y_0$ and $x = \apply{y_0}{i}$ by \cref{pair_eq_iff}.
                                Thus $y\in\productFamily{Y}$ and $z = y(i)$.
                                Thus $\productFamily{Y}\subseteq\productFamily{X_1}$ by \cref{subspace_of}.
                                Thus $y\in\productFamily{X_1}$ by \cref{subseteq}.
                                Thus $(y,\apply{y}{i})\in\piIndex{X_1}{i}$ by \cref{projection_map_indexed}.
                                Thus $(y,y)\in j$ by \cref{id_iff}.
                                Thus $(y,z)\in\piIndex{X_1}{i}\circ j$ by \cref{circ_iff}.
                            \end{subproof}
                            Show that for all $y,z$ if $y\mathrel{\piIndex{X_1}{i}\circ j} z$ then $y\mathrel{k\circ g} z$.
                            \begin{subproof}
                                Fix $y$.
                                Fix $z$.
                                Assume $y\mathrel{\piIndex{X_1}{i}\circ j} z$.
                                Take $x$ such that $y\mathrel{j} x$ and $x\mathrel{\piIndex{X_1}{i}} z$ by \cref{circ_iff}.
                                Thus $y = x$ and $y\in\productFamily{Y}$ by \cref{id_iff}.
                                Take $x_1\in\productFamily{X_1}$ such that $(x,z) = (x_1,\apply{x_1}{i})$ by \cref{projection_map_indexed}.
                                Thus $x = x_1$ and $z = \apply{x_1}{i}$ by \cref{pair_eq_iff}.
                                Thus $z = y(i)$.
                                Thus $(y,\apply{y}{i})\in g$ by \cref{projection_map_indexed}.
                                Thus $(y,z)\in g$.
                                Thus $i\in J$ by \cref{subseteq}.
                                Thus $i\in\dom{Y}$ by \cref{subseteq}.
                                Thus $z\in\apply{Y}{i}$ by \cref{indexed_product}.
                                Thus $(z,z)\in k$ by \cref{id_iff}.
                                Thus $y\mathrel{k\circ g} z$ by \cref{circ_iff}.
                            \end{subproof}
                            Follows by set extensionality.
                        \end{subproof}
                        Thus $\piIndex{X_1}{i}\circ j$ is continuous from $\productFamily{Y}$ to $\apply{X_1}{i}$ with respect to $T_{P0}$ and $\apply{T_1}{i}$
                            by \cref{continuous_eq_subst}.
                    \end{subproof}
                    Thus $j$ is continuous from $\productFamily{Y}$ to $\productFamily{X_1}$ with respect to $T_{P0}$ and $T_0$
                        by \cref{continuous_map_into_indexed_product}.
                    Thus $j$ is continuous from $\productFamily{Y}$ to $X$ with respect to $T_{P0}$ and $T$ by \cref{continuous_eq_codomain}.
                \end{subproof}
                Show that for all $U$ such that $U\in T_{Y0}$ we have $U\in T_{P0}$.
                \begin{subproof}
                    Fix $U$.
                    Assume $U\in T_{Y0}$.
                    Take $V\in T$ such that $U = \productFamily{Y}\inter V$ by \cref{subspace_topology}.
                    Let $j = \identity{\productFamily{Y}}$.
                    Thus $T$ is a topology on $X$ by \cref{continuous}.
                    Thus $V$ is open in $X$ with respect to $T$ by \cref{open_in}.
                    Thus $\preimg{j}{V}$ is open in $\productFamily{Y}$ with respect to $T_{P0}$ by \cref{continuous}.
                    Show that $\preimg{j}{V} = \productFamily{Y}\inter V$.
                    \begin{subproof}
                        Show that for all $y$ such that $y\in\preimg{j}{V}$ we have $y\in \productFamily{Y}\inter V$.
                        \begin{subproof}
                            Fix $y$.
                            Assume $y\in\preimg{j}{V}$.
                            Take $z\in V$ such that $y\mathrel{j} z$ by \cref{preimg_iff}.
                            Thus $y = z$ and $y\in \productFamily{Y}$ by \cref{id_iff}.
                            Thus $y\in V$.
                            Thus $y\in \productFamily{Y}$ and $y\in V$.
                            Thus $y\in \productFamily{Y}\inter V$ by \cref{inter_intro}.
                        \end{subproof}
                        Thus $\preimg{j}{V}\subseteq \productFamily{Y}\inter V$ by \cref{subseteq}.
                        Show that for all $y$ such that $y\in \productFamily{Y}\inter V$ we have $y\in\preimg{j}{V}$.
                        \begin{subproof}
                            Fix $y$.
                            Assume $y\in \productFamily{Y}\inter V$.
                            Thus $y\in \productFamily{Y}$ and $y\in V$ by \cref{inter}.
                            Thus $y\mathrel{j} y$ by \cref{id_iff}.
                            Thus $y\in\preimg{j}{V}$ by \cref{preimg_iff}.
                        \end{subproof}
                        Thus $\productFamily{Y}\inter V\subseteq \preimg{j}{V}$ by \cref{subseteq}.
                        Thus $\preimg{j}{V} = \productFamily{Y}\inter V$ by \cref{subseteq_antisymmetric}.
                    \end{subproof}
                    Thus $U\in T_{P0}$ by \cref{open_in}.
                \end{subproof}
                Thus $T_{Y0}\subseteq T_{P0}$ by \cref{subseteq}.
            \end{subproof}
            Show that $\productFamily{Y}$ is compact with respect to $T_{Y0}$.
            \begin{subproof}
                Show that for all $C$ such that $C$ is an open covering of $\productFamily{Y}$ with respect to $T_{Y0}$ there exists $B$ such that
                    $B\subseteq C$ and $B$ is finite and $B$ is a covering of $\productFamily{Y}$.
                \begin{subproof}
                    Fix $C$.
                    Assume $C$ is an open covering of $\productFamily{Y}$ with respect to $T_{Y0}$.
                    Show that $T_{P0}$ is a topology on $\productFamily{Y}$.
                    \begin{subproof}
                        Thus $Y$ is right-unique and $Y$ is a relation.
                        Thus $Y$ is a function.
                        Thus $T_Y$ is right-unique and $T_Y$ is a relation.
                        Thus $T_Y$ is a function.
                        Thus $\dom{T_Y} = \dom{Y}$.
                        Show that $\dom{Y}$ is inhabited.
                        \begin{subproof}
                            $1\in\naturalsPlus$ by \cref{real_one_in_naturals}.
                            $1 \leq n$ by \cref{real_one_leq_naturalsplus}.
                            Thus $1\in J$ by \cref{initial_segment_naturalsplus}.
                            Thus $1\in\dom{Y}$ by \cref{subseteq}.
                            Thus $\dom{Y}$ is inhabited.
                        \end{subproof}
                        Show that for all $i\in\dom{Y}$ we have $\apply{T_Y}{i}$ is a topology on $\apply{Y}{i}$.
                        \begin{subproof}
                            Fix $i$.
                            Assume $i\in\dom{Y}$.
                            Thus $i\in J$ by \cref{subseteq}.
                            Thus $(i,\intervalclosed{x_0(i) - M}{x_0(i) + M})\in Y$.
                            Thus $(i,\subspaceTopology{\standardTopologyR}{\intervalclosed{x_0(i) - M}{x_0(i) + M}})\in T_Y$.
                            Thus $\apply{Y}{i} = \intervalclosed{x_0(i) - M}{x_0(i) + M}$ by \cref{function_apply_intro}.
                            Thus $\apply{T_Y}{i} = \subspaceTopology{\standardTopologyR}{\intervalclosed{x_0(i) - M}{x_0(i) + M}}$ by \cref{function_apply_intro}.
                            $\standardBasisR$ is a basis on $\reals$ by \cref{standard_basis_is_basis}.
                            Thus $\standardTopologyR$ is a topology on $\reals$ by \cref{basis_topology_is_topology,standard_topology_reals}.
                            Show that $\intervalclosed{x_0(i) - M}{x_0(i) + M}\subseteq\reals$.
                            \begin{subproof}
                                Show that for all $x$ such that $x\in\intervalclosed{x_0(i) - M}{x_0(i) + M}$ we have $x\in\reals$.
                                \begin{subproof}
                                    Fix $x$.
                                    Assume $x\in\intervalclosed{x_0(i) - M}{x_0(i) + M}$.
                                    Thus $x\in\reals$ by \cref{intervalclosed}.
                                \end{subproof}
                                Thus $\intervalclosed{x_0(i) - M}{x_0(i) + M}\subseteq\reals$ by \cref{subseteq}.
                            \end{subproof}
                            Thus $\apply{T_Y}{i}$ is a topology on $\apply{Y}{i}$ by \cref{subspace_topology_is_topology}.
                        \end{subproof}
                        Thus $T_{P0}$ is a topology on $\productFamily{Y}$
                            by \cref{product_subbasis_finite_intersections_basis,basis_for_topology,basis_topology_is_topology}.
                    \end{subproof}
                    Show that $C$ is an open covering of $\productFamily{Y}$ with respect to $T_{P0}$.
                    \begin{subproof}
                        Show that for all $U$ such that $U\in C$ we have $U$ is open in $\productFamily{Y}$ with respect to $T_{P0}$.
                        \begin{subproof}
                            Fix $U$.
                            Assume $U\in C$.
                            Thus $U$ is open in $\productFamily{Y}$ with respect to $T_{Y0}$ by \cref{open_covering}.
                            Thus $U\in T_{Y0}$ by \cref{open_in}.
                            Thus $U\in T_{P0}$ by \cref{subseteq}.
                            Thus $T_{P0}$ is a topology on $\productFamily{Y}$ by \cref{continuous}.
                            Thus $U$ is open in $\productFamily{Y}$ with respect to $T_{P0}$ by \cref{open_in}.
                        \end{subproof}
                        Thus $C$ is an open covering of $\productFamily{Y}$ with respect to $T_{P0}$ by \cref{open_covering}.
                    \end{subproof}
                    Take $B$ such that $B\subseteq C$ and $B$ is finite and $B$ is a covering of $\productFamily{Y}$
                        by \cref{compact}.
                \end{subproof}
                Thus $\productFamily{Y}$ is compact with respect to $T_{Y0}$ by \cref{compact}.
            \end{subproof}
            Show that $\productFamily{Y}$ is a subspace of $X$ with respect to $T$.
            \begin{subproof}
                Show that $T$ is a topology on $X$.
                \begin{subproof}
                    $\rho$ is a metric on $X$ by \cref{square_metric_is_metric}.
                    Thus $\metricBasis{\rho}{X}$ is a basis on $X$ by \cref{metric_basis_is_basis}.
                    Thus $T$ is a topology on $X$ by \cref{basis_topology_is_topology,metric_topology,basis_for_topology}.
                \end{subproof}
                Show that $\productFamily{Y}\subseteq X$.
                \begin{subproof}
                    Show that for all $y$ such that $y\in\productFamily{Y}$ we have $y\in X$.
                    \begin{subproof}
                        Fix $y$.
                        Assume $y\in\productFamily{Y}$.
                        Thus $y\in\funs{\dom{Y}}{\unions{\ran{Y}}}$ and $\dom{y} = \dom{Y}$ by \cref{indexed_product,funs_elim}.
                        Thus $\dom{y} = J$ by \cref{subseteq}.
                        Show that for all $i\in J$ we have $y(i)\in\reals$.
                        \begin{subproof}
                            Fix $i$.
                            Assume $i\in J$.
                            Thus $i\in\dom{Y}$ by \cref{subseteq}.
                            Thus $y(i)\in\apply{Y}{i}$ by \cref{indexed_product}.
                            Thus $(i,\intervalclosed{x_0(i) - M}{x_0(i) + M})\in Y$.
                            Thus $\apply{Y}{i} = \intervalclosed{x_0(i) - M}{x_0(i) + M}$ by \cref{function_apply_intro}.
                            Thus $y(i)\in\intervalclosed{x_0(i) - M}{x_0(i) + M}$.
                            Thus $y(i)\in\reals$ by \cref{intervalclosed}.
                        \end{subproof}
                        Thus $y$ is a function to $\reals$.
                        Thus $y\in\funs{J}{\reals}$ by \cref{funs_intro}.
                        Thus $y\in\tuplePower{\reals}{J}$ by \cref{tuple_power}.
                        Thus $y\in\realsPower{J}$ by \cref{reals_power}.
                        Thus $y\in X$.
                    \end{subproof}
                    Thus $\productFamily{Y}\subseteq X$ by \cref{subseteq}.
                \end{subproof}
                Thus $\productFamily{Y}$ is a subspace of $X$ with respect to $T$ by \cref{subspace_of}.
            \end{subproof}
            Show that $A$ is closed in $\productFamily{Y}$ with respect to $T_{Y0}$.
            \begin{subproof}
                Thus $A\subseteq\productFamily{Y}$ by \cref{subseteq}.
                Thus $A = A\inter\productFamily{Y}$ by \cref{inter_absorb_supseteq_right}.
                Thus $A$ is closed in $\productFamily{Y}$ with respect to $T_{Y0}$ by \cref{closed_in_subspace_iff,subspace_of}.
            \end{subproof}
            Thus $T$ is a topology on $X$ and $\productFamily{Y}\subseteq X$ by \cref{subspace_of}.
            Thus $T_{Y0}$ is a topology on $\productFamily{Y}$ by \cref{subspace_topology_is_topology}.
            Let $T_{YA} = \subspaceTopology{T_{Y0}}{A}$.
            Thus $A$ is compact with respect to $T_{YA}$ by \cref{closed_subspace_compact}.
            Thus $T_{YA} = T_A$ by \cref{subspace_subspace_topology}.
            Thus $A$ is compact with respect to $T_A$.
        \end{byCase}
    \end{subproof}
\end{proof}

% from §27 helper lemma 10: finite inhabited subset has a largest element
\begin{lemma}\label{finite_inhabited_largest_element}
    Assume $R$ is a simple order relation on $X$.
    Suppose $F\subseteq X$.
    Suppose $F$ is finite.
    Suppose $F$ is inhabited.
    Then there exists $b$ such that $b$ is a largest element of $F$ with respect to $R$.
\end{lemma}
\begin{proof}
    Take $k$ such that $k$ is a natural number and there exists a bijection from $k$ to $F$ by \cref{finite}.
    Take $f$ such that $f$ is a bijection from $k$ to $F$.
    Let $A = \{ n\in\naturals \mid \text{for all $G,g$ such that $g$ is a bijection from $n$ to $G$ and $G\subseteq X$ and $G$ is inhabited
        there exists $b$ such that $b$ is a largest element of $G$ with respect to $R$}\}$.
    Show that $A$ is an inductive set.
    \begin{subproof}
        Show that $\emptyset\in A$.
        \begin{subproof}
            Show that for all $G,g$ such that $g$ is a bijection from $\emptyset$ to $G$ and $G\subseteq X$ and $G$ is inhabited
                there exists $b$ such that $b$ is a largest element of $G$ with respect to $R$.
            \begin{subproof}
                Fix $G$.
                Fix $g$.
                Assume $g$ is a bijection from $\emptyset$ to $G$.
                Assume $G\subseteq X$ and $G$ is inhabited.
                Show that there exists $b$ such that $b$ is a largest element of $G$ with respect to $R$.
                \begin{subproof}
                    Suppose not.
                    Take $x$ such that $x\in G$.
                    $g\in\surj{\emptyset}{G}$ by \cref{bijections}.
                    Thus there exists $y\in\emptyset$ such that $g(y)=x$ by \cref{surj}.
                    Contradiction by \cref{emptyset}.
                \end{subproof}
            \end{subproof}
        \end{subproof}
        Show that for all $n\in A$ we have $\suc{n}\in A$.
        \begin{subproof}
            Fix $n$.
            Assume $n\in A$.
            Show that for all $G,g$ such that $g$ is a bijection from $\suc{n}$ to $G$ and $G\subseteq X$ and $G$ is inhabited
                there exists $b$ such that $b$ is a largest element of $G$ with respect to $R$.
            \begin{subproof}
                Fix $G$.
                Fix $g$.
                Assume $g$ is a bijection from $\suc{n}$ to $G$.
                Assume $G\subseteq X$ and $G$ is inhabited.
                $g$ is a function by \cref{bijection_is_function}.
                $\dom{g} = \suc{n}$ by \cref{bijections_dom}.
                $n\in\suc{n}$ by \cref{suc_intro_self}.
                Thus $n\in\dom{g}$.
                Let $a = g(n)$.
                Let $G_1 = \img{g}{n}$.
                Show that for all $x$ such that $x\in n$ we have $x\in\suc{n}$.
                \begin{subproof}
                    Fix $x$.
                    Assume $x\in n$.
                    Thus $x\in\suc{n}$ by \cref{suc}.
                \end{subproof}
                Thus $n\subseteq \suc{n}$ by \cref{subseteq}.
                Thus $G_1\subseteq \img{g}{\suc{n}}$ by \cref{img_subseteq}.
                $\img{g}{\suc{n}} = \ran{g}$ by \cref{img_dom}.
                $g\in\surj{\suc{n}}{G}$ by \cref{bijections}.
                Thus $\ran{g} = G$ by \cref{ran_of_surj}.
                Thus $G_1\subseteq G$.
                Thus $G_1\subseteq X$ by \cref{subseteq_transitive}.
                $(n,a)\in g$ by \cref{function_apply_elim}.
                Thus $a\in\ran{g}$ by \cref{ran_iff}.
                Thus $a\in G$.
                Thus $a\in X$ by \cref{elem_subseteq}.
                Show that $G = G_1\union\{a\}$.
                \begin{subproof}
                    Show that $G\subseteq G_1\union\{a\}$.
                    \begin{subproof}
                        Show that for all $z$ such that $z\in G$ we have $z\in G_1\union\{a\}$.
                        \begin{subproof}
                            Fix $z$.
                            Assume $z\in G$.
                            $g\in\surj{\suc{n}}{G}$ by \cref{bijections}.
                            Thus there exists $x\in\suc{n}$ such that $g(x)=z$ by \cref{surj}.
                            $x\in n$ or $x = n$ by \cref{suc}.
                            \begin{byCase}
                            \caseOf{$x\in n$.}
                                $(x,z)\in g$ by \cref{function_apply_elim}.
                                Thus $z\in G_1$ by \cref{img_iff}.
                                Thus $z\in G_1\union\{a\}$ by \cref{union_intro_left}.
                            \caseOf{$x = n$.}
                                Thus $z = g(n)$.
                                Thus $z = a$.
                                Thus $z\in G_1\union\{a\}$ by \cref{union_intro_right,singleton_intro}.
                            \end{byCase}
                        \end{subproof}
                        Thus $G\subseteq G_1\union\{a\}$ by \cref{subseteq}.
                    \end{subproof}
                    $G_1\subseteq G$.
                    Thus $\{a\}\subseteq G$ by \cref{singleton_subset_intro}.
                    Thus $G_1\union\{a\}\subseteq G$ by \cref{union_subsets_is_subset}.
                    Thus $G = G_1\union\{a\}$ by \cref{subseteq_antisymmetric}.
                \end{subproof}
                Let $h = \restrl{g}{n}$.
                $h$ is a function by \cref{restrl_of_function_is_function}.
                Show that $h$ is a bijection from $n$ to $G_1$.
                \begin{subproof}
                    Show that for all $x$ such that $x\in n$ we have $x\in\dom{g}$.
                    \begin{subproof}
                        Fix $x$.
                        Assume $x\in n$.
                        Thus $x\in\suc{n}$ by \cref{suc}.
                        Thus $x\in\dom{g}$ by \cref{bijections_dom}.
                    \end{subproof}
                    Thus $n\subseteq\dom{g}$ by \cref{subseteq}.
                    Thus $\dom{h} = \dom{g}\inter n$ by \cref{dom_of_restrl_eq_inter}.
                    Thus $\dom{h} = n$ by \cref{inter_absorb_supseteq_left}.
                    Show that $h$ is a function to $G_1$.
                    \begin{subproof}
                        Show that for all $x\in\dom{h}$ we have $h(x)\in G_1$.
                        \begin{subproof}
                            Fix $x$.
                            Assume $x\in\dom{h}$.
                            Thus $x\in n$.
                            $(x,h(x))\in h$ by \cref{function_apply_elim}.
                            $h\subseteq g$ by \cref{restrl_subseteq}.
                            Thus $(x,h(x))\in g$ by \cref{elem_subseteq}.
                            Thus $h(x)\in G_1$ by \cref{img_iff}.
                        \end{subproof}
                        Thus $h$ is a function to $G_1$.
                    \end{subproof}
                    Thus $h\in\funs{n}{G_1}$ by \cref{funs_intro}.
                    Show that for all $y\in G_1$ there exists $x\in n$ such that $h(x) = y$.
                    \begin{subproof}
                        Fix $y$.
                        Assume $y\in G_1$.
                        Take $x\in n$ such that $(x,y)\in g$ by \cref{img_iff}.
                        Thus $x\in\dom{g}$ and $g(x) = y$ by \cref{function_apply_iff}.
                        Thus $h(x) = y$ by \cref{restrl_of_function_apply}.
                    \end{subproof}
                    Thus $h\in\surj{n}{G_1}$ by \cref{surj}.
                    Show that $h$ is injective.
                    \begin{subproof}
                        Show that for all $x,y$ such that $x\in\dom{h}$ and $y\in\dom{h}$ and $h(x) = h(y)$ we have $x = y$.
                        \begin{subproof}
                            Fix $x$.
                            Fix $y$.
                            Assume $x\in\dom{h}$ and $y\in\dom{h}$ and $h(x) = h(y)$.
                            Thus $x\in n$ and $y\in n$.
                            Thus $h(x) = g(x)$ and $h(y) = g(y)$ by \cref{restrl_of_function_apply}.
                            Thus $g(x) = g(y)$.
                            $g$ is injective by \cref{bijections}.
                            Thus $x = y$ by \cref{injective_function}.
                        \end{subproof}
                        Thus $h$ is injective by \cref{injective_function}.
                    \end{subproof}
                    Thus $h$ is a bijection from $n$ to $G_1$ by \cref{bijections}.
                \end{subproof}
                \begin{byCase}
                \caseOf{$n = \emptyset$.}
                    Thus $G_1 = \emptyset$ by \cref{img_emptyset}.
                    Thus $G = \{a\}$ by \cref{union_emptyset,union_comm}.
                    Show that $a$ is a largest element of $G$ with respect to $R$.
                    \begin{subproof}
                        $a\in G$ by \cref{singleton_intro}.
                        Show that for all $x\in G$ we have $x\mathrel{R}a$ or $x = a$.
                        \begin{subproof}
                            Fix $x$.
                            Assume $x\in G$.
                            Thus $x = a$ by \cref{singleton_iff}.
                        \end{subproof}
                        Thus $a$ is a largest element of $G$ with respect to $R$ by \cref{largest_element}.
                    \end{subproof}
                \caseOf{$n\neq \emptyset$.}
                    $n$ is inhabited by \cref{nonempty_inhabited}.
                    Take $x_0$ such that $x_0\in n$.
                    Thus $x_0\in\suc{n}$ by \cref{suc}.
                    Thus $x_0\in\dom{g}$ by \cref{bijections_dom}.
                    Thus $(x_0,g(x_0))\in g$ by \cref{function_apply_elim}.
                    Thus $g(x_0)\in G_1$ by \cref{img_iff}.
                    Thus $G_1$ is inhabited.
                    Thus there exists $b$ such that $b$ is a largest element of $G_1$ with respect to $R$.
                    Take $b$ such that $b$ is a largest element of $G_1$ with respect to $R$.
                    $b\in G_1$ by \cref{largest_element}.
                    Thus $b\in X$ by \cref{elem_subseteq}.
                    $R$ is transitive by \cref{simple_order_relation}.
                    $R$ is asymmetric by \cref{simple_order_relation}.
                    $R$ is connex on $X$ by \cref{simple_order_relation}.
                    Thus $b\mathrel{R}a$ or $a\mathrel{R}b$ or $b = a$ by \cref{connex}.
                    \begin{byCase}
                    \caseOf{$b\mathrel{R}a$ or $b = a$.}
                        Show that $a$ is a largest element of $G$ with respect to $R$.
                        \begin{subproof}
                            $a\in G$ by \cref{union_intro_right,singleton_intro}.
                            Show that for all $y\in G$ we have $y\mathrel{R}a$ or $y = a$.
                            \begin{subproof}
                                Fix $y$.
                                Assume $y\in G$.
                                Thus $y\in G_1$ or $y = a$ by \cref{union_iff,singleton_iff}.
                                \begin{byCase}
                                \caseOf{$y\in G_1$.}
                                    Thus $y\mathrel{R}b$ or $y = b$ by \cref{largest_element}.
                                    \begin{byCase}
                                    \caseOf{$b = a$.}
                                        Thus $y\mathrel{R}a$ or $y = a$.
                                    \caseOf{$b\mathrel{R}a$.}
                                        \begin{byCase}
                                        \caseOf{$y\mathrel{R}b$.}
                                            Thus $y\mathrel{R}a$ by \cref{transitive}.
                                        \caseOf{$y = b$.}
                                            Thus $y\mathrel{R}a$.
                                        \end{byCase}
                                    \end{byCase}
                                \caseOf{$y = a$.}
                                    Thus $y = a$.
                                \end{byCase}
                            \end{subproof}
                            Thus $a$ is a largest element of $G$ with respect to $R$ by \cref{largest_element}.
                        \end{subproof}
                    \caseOf{$a\mathrel{R}b$.}
                        Show that $b$ is a largest element of $G$ with respect to $R$.
                        \begin{subproof}
                            $b\in G$ by \cref{elem_subseteq}.
                            Show that for all $y\in G$ we have $y\mathrel{R}b$ or $y = b$.
                            \begin{subproof}
                                Fix $y$.
                                Assume $y\in G$.
                                Thus $y\in G_1$ or $y = a$ by \cref{union_iff,singleton_iff}.
                                \begin{byCase}
                                \caseOf{$y\in G_1$.}
                                    Thus $y\mathrel{R}b$ or $y = b$ by \cref{largest_element}.
                                \caseOf{$y = a$.}
                                    Thus $y\mathrel{R}b$.
                                \end{byCase}
                            \end{subproof}
                            Thus $b$ is a largest element of $G$ with respect to $R$ by \cref{largest_element}.
                        \end{subproof}
                    \end{byCase}
                \end{byCase}
            \end{subproof}
        \end{subproof}
    \end{subproof}
    Thus $A$ is an inductive set.
    Thus $\naturals\subseteq A$ by \cref{num_naturals_smallest_inductive_set}.
    Thus $k\in\naturals$.
    Thus $k\in A$ by \cref{subseteq}.
    Thus for all $G,g$ such that $g$ is a bijection from $k$ to $G$ and $G\subseteq X$ and $G$ is inhabited
        there exists $b$ such that $b$ is a largest element of $G$ with respect to $R$.
    Thus there exists $b$ such that $b$ is a largest element of $F$ with respect to $R$.
\end{proof}

% from §27 helper lemma 11: finite inhabited subset has a smallest element
\begin{lemma}\label{finite_inhabited_smallest_element}
    Assume $R$ is a simple order relation on $X$.
    Suppose $F\subseteq X$.
    Suppose $F$ is finite.
    Suppose $F$ is inhabited.
    Then there exists $a$ such that $a$ is a smallest element of $F$ with respect to $R$.
\end{lemma}
\begin{proof}
    Let $R_1 = \converse{R}$.
    Show that $R_1$ is a simple order relation on $X$.
    \begin{subproof}
        $R$ is a binary relation on $X$ by \cref{simple_order_relation}.
        Thus $R\subseteq X\times X$.
        Thus $R_1\subseteq X\times X$ by \cref{converse_type}.
        Thus $R_1$ is a binary relation on $X$.
        $R$ is transitive by \cref{simple_order_relation}.
        $R$ is asymmetric by \cref{simple_order_relation}.
        $R$ is connex on $X$ by \cref{simple_order_relation}.
        Show that $R_1$ is transitive.
        \begin{subproof}
            Show that for all $x,y,z$ if $x\mathrel{R_1}y$ and $y\mathrel{R_1}z$ then $x\mathrel{R_1}z$.
            \begin{subproof}
                Fix $x$.
                Fix $y$.
                Fix $z$.
                Assume $x\mathrel{R_1}y$ and $y\mathrel{R_1}z$.
                Thus $(x,y)\in R_1$ and $(y,z)\in R_1$.
                Thus $(x,y)\in X\times X$ and $(y,z)\in X\times X$ by \cref{elem_subseteq}.
                Thus $x\in X$ and $y\in X$ and $z\in X$ by \cref{times_tuple_elim}.
                Thus $y\mathrel{R}x$ and $z\mathrel{R}y$ by \cref{converse_iff}.
                Thus $z\mathrel{R}x$ by \cref{transitive}.
                Thus $x\mathrel{R_1}z$ by \cref{converse_iff}.
            \end{subproof}
            Thus $R_1$ is transitive by \cref{transitive}.
        \end{subproof}
        Show that $R_1$ is asymmetric.
        \begin{subproof}
            Show that for all $x,y$ if $x\mathrel{R_1}y$ then it is wrong that $y\mathrel{R_1}x$.
            \begin{subproof}
                Fix $x$.
                Fix $y$.
                Assume $x\mathrel{R_1}y$.
                Thus $(x,y)\in R_1$.
                Thus $(x,y)\in X\times X$ by \cref{elem_subseteq}.
                Thus $x\in X$ and $y\in X$ by \cref{times_tuple_elim}.
                Show that it is wrong that $y\mathrel{R_1}x$.
                \begin{subproof}
                    Suppose not.
                    Thus $y\mathrel{R_1}x$.
                    Thus $y\mathrel{R}x$ and $x\mathrel{R}y$ by \cref{converse_iff}.
                    Contradiction by \cref{asymmetric}.
                \end{subproof}
            \end{subproof}
            Thus $R_1$ is asymmetric by \cref{asymmetric}.
        \end{subproof}
        Show that $R_1$ is connex on $X$.
        \begin{subproof}
            Show that for all $x,y$ such that $x\in X$ and $y\in X$ and $x\neq y$ we have $x\mathrel{R_1}y$ or $y\mathrel{R_1}x$.
            \begin{subproof}
                Fix $x$.
                Fix $y$.
                Assume $x\in X$ and $y\in X$ and $x\neq y$.
                Thus $x\mathrel{R}y$ or $y\mathrel{R}x$ by \cref{connex}.
                Thus $y\mathrel{R_1}x$ or $x\mathrel{R_1}y$ by \cref{converse_iff}.
            \end{subproof}
            Thus $R_1$ is connex on $X$ by \cref{connex}.
        \end{subproof}
    \end{subproof}
    Thus there exists $b$ such that $b$ is a largest element of $F$ with respect to $R_1$
        by \cref{finite_inhabited_largest_element}.
    Take $b$ such that $b$ is a largest element of $F$ with respect to $R_1$.
    Show that $b$ is a smallest element of $F$ with respect to $R$.
    \begin{subproof}
        $b\in F$ by \cref{largest_element}.
        Show that for all $x\in F$ we have $b\mathrel{R}x$ or $b = x$.
        \begin{subproof}
            Fix $x$.
            Assume $x\in F$.
            Thus $x\mathrel{R_1}b$ or $x = b$ by \cref{largest_element}.
            Thus $b\mathrel{R}x$ or $b = x$ by \cref{converse_iff}.
        \end{subproof}
        Thus $b$ is a smallest element of $F$ with respect to $R$ by \cref{smallest_element}.
    \end{subproof}
\end{proof}

% from §27 Item 4: extreme value theorem
\begin{theorem}\label{extreme_value_theorem}
    Assume $T_X$ is a topology on $X$.
    Assume $T_Y$ is the order topology on $Y$ with respect to $R$.
    Suppose $f$ is continuous from $X$ to $Y$ with respect to $T_X$ and $T_Y$.
    Suppose $X$ is compact with respect to $T_X$.
    Suppose $X\neq\emptyset$.
    Then there exist $c,d\in X$ such that
        $\apply{f}{c}$ is a smallest element of $\img{f}{X}$ with respect to $R$ and
        $\apply{f}{d}$ is a largest element of $\img{f}{X}$ with respect to $R$.
\end{theorem}
\begin{proof}
    Let $A = \img{f}{X}$.
    Let $T_A = \subspaceTopology{T_Y}{A}$.
    Thus $A$ is compact with respect to $T_A$ by \cref{continuous_image_compact}.
    $T_Y = \basisTopology{\orderBasis{R}{Y}}{Y}$ by \cref{order_topology}.
    $R$ is a simple order relation on $Y$ by \cref{order_topology}.
    $Y$ has more than one element by \cref{order_topology}.
    $\orderBasis{R}{Y}$ is a basis on $Y$ by \cref{order_basis_is_basis}.
    Thus $T_Y$ is a topology on $Y$ by \cref{basis_topology_is_topology}.
    $f\in\funs{X}{Y}$ by \cref{continuous}.
    Thus $f$ is a function to $Y$ and $\dom{f} = X$ by \cref{funs_elim}.
    Thus $\img{f}{X} = \ran{f}$ by \cref{img_dom}.
    Thus $A\subseteq Y$ by \cref{function_to_ran}.
    Thus $A$ is a subspace of $Y$ with respect to $T_Y$ by \cref{subspace_of}.
    Thus $X$ is inhabited by \cref{nonempty_inhabited}.
    Take $x_0$ such that $x_0\in X$.
    Thus $f$ is a function by \cref{funs_is_function}.
    Thus $x_0\in\dom{f}$.
    Thus $(x_0,\apply{f}{x_0})\in f$ by \cref{function_apply_elim}.
    Thus $\apply{f}{x_0}\in A$ by \cref{img_iff}.
    Thus $A$ is inhabited.
    Show that there exists $M$ such that $M$ is a largest element of $A$ with respect to $R$.
    \begin{subproof}
        Suppose not.
        Let $C = \{U \mid \exists a\in A. U = \openrayleft{R}{Y}{a}\}$.
        Show that $C$ is a covering of $A$.
        \begin{subproof}
            Show that for all $y$ such that $y\in A$ we have $y\in\unions{C}$.
            \begin{subproof}
                Fix $y$.
                Assume $y\in A$.
                Thus $y$ is not a largest element of $A$ with respect to $R$.
                Take $a\in A$ such that it is wrong that $a\mathrel{R}y$ or $a = y$ by \cref{largest_element}.
                Thus $a\neq y$ and it is wrong that $a\mathrel{R}y$.
                Thus $y\in Y$ by \cref{subseteq}.
                Thus $a\in Y$ by \cref{subseteq}.
                $R$ is connex on $Y$ by \cref{order_topology,simple_order_relation}.
                Thus $y\mathrel{R}a$ or $a\mathrel{R}y$ by \cref{connex}.
                \begin{byCase}
                \caseOf{$y\mathrel{R}a$.}
                    Thus $y\in\openrayleft{R}{Y}{a}$ by \cref{openray_left}.
                    Thus $\openrayleft{R}{Y}{a}\in C$.
                    Thus $y\in\unions{C}$ by \cref{unions_iff}.
                \caseOf{$a\mathrel{R}y$.}
                    Contradiction.
                \end{byCase}
            \end{subproof}
            Thus $A\subseteq\unions{C}$ by \cref{subseteq}.
            Thus $C$ is a covering of $A$ by \cref{covering}.
        \end{subproof}
        Show that for all $U$ such that $U\in C$ we have $U$ is open in $Y$ with respect to $T_Y$.
        \begin{subproof}
            Fix $U$.
            Assume $U\in C$.
            Take $a\in A$ such that $U = \openrayleft{R}{Y}{a}$.
            Thus $a\in Y$ by \cref{subseteq}.
            Thus $U$ is open in $Y$ with respect to $T_Y$ by \cref{openrayleft_open_in_order_topology}.
        \end{subproof}
        Thus there exists $B$ such that $B\subseteq C$ and $B$ is finite and $B$ is a covering of $A$
            by \cref{compact_subspace_open_covering,compact}.
        Take $B$ such that $B\subseteq C$ and $B$ is finite and $B$ is a covering of $A$.
        Show that $B$ is inhabited.
        \begin{subproof}
            Take $y$ such that $y\in A$.
            Thus $y\in\unions{B}$ by \cref{covering,subseteq}.
            Take $U\in B$ such that $y\in U$ by \cref{unions_iff}.
        \end{subproof}
        Let $R_B = \{w\in B\times A \mid \exists U\in B.\exists a\in A. w = (U,a) \land U = \openrayleft{R}{Y}{a}\}$.
        Show that $R_B$ is a relation.
        \begin{subproof}
            Show that for all $w$ such that $w\in R_B$ there exist $U,a$ such that $w = (U,a)$.
            \begin{subproof}
                Fix $w$.
                Assume $w\in R_B$.
                Take $U\in B$ such that there exists $a\in A$ such that $w = (U,a)$ and $U = \openrayleft{R}{Y}{a}$.
                Take $a\in A$ such that $w = (U,a)$ and $U = \openrayleft{R}{Y}{a}$.
            \end{subproof}
            Thus $R_B$ is a relation by \cref{relation}.
        \end{subproof}
        Show that for all $U\in B$ there exists $a$ such that $U\mathrel{R_B}a$.
        \begin{subproof}
            Fix $U$.
            Assume $U\in B$.
            Thus $U\in C$ by \cref{subseteq}.
            Take $a\in A$ such that $U = \openrayleft{R}{Y}{a}$.
            Thus $(U,a)\in R_B$.
            Thus $U\mathrel{R_B} a$ by \cref{relation}.
        \end{subproof}
        Take $g$ such that $g$ is a function on $B$ and for all $U\in B$ we have $U\mathrel{R_B} \apply{g}{U}$
            by \cref{axiom_of_choice}.
        Let $E = \img{g}{B}$.
        Show that $E\subseteq A$.
        \begin{subproof}
            Show that for all $e$ such that $e\in E$ we have $e\in A$.
            \begin{subproof}
                Fix $e$.
                Assume $e\in E$.
                Thus $e\in\img{g}{B}$.
                Take $U\in B$ such that $(U,e)\in g$ by \cref{img_iff}.
                Thus $g$ is a function.
                Thus $g(U) = e$ by \cref{function_apply_intro}.
                Thus $U\mathrel{R_B} g(U)$.
                Thus $(U,\apply{g}{U})\in R_B$ by \cref{relation}.
                Thus $(U,\apply{g}{U})\in B\times A$.
                Thus there exists $U_1$ such that $U_1\in B$ and there exists $a$ such that
                    $a\in A$ and $(U,\apply{g}{U}) = (U_1,a)$ and $U_1 = \openrayleft{R}{Y}{a}$.
                Take $U_1$ such that $U_1\in B$ and there exists $a$ such that
                    $a\in A$ and $(U,\apply{g}{U}) = (U_1,a)$ and $U_1 = \openrayleft{R}{Y}{a}$.
                Take $a$ such that $a\in A$ and $(U,\apply{g}{U}) = (U_1,a)$ and
                    $U_1 = \openrayleft{R}{Y}{a}$.
                Thus $U = U_1$ and $g(U) = a$ by \cref{pair_eq_iff}.
                Thus $U = \openrayleft{R}{Y}{a}$ and $g(U) = a$.
                Thus $e\in A$.
            \end{subproof}
            Thus $E\subseteq A$ by \cref{subseteq}.
        \end{subproof}
        Thus $E$ is finite by \cref{finite_image_function}.
        Show that $E$ is inhabited.
        \begin{subproof}
            Take $U_0$ such that $U_0\in B$.
            Thus $g$ is a function.
            Thus $U_0\in\dom{g}$.
            Thus $(U_0,\apply{g}{U_0})\in g$ by \cref{function_apply_elim}.
            Thus $g(U_0)\in E$ by \cref{img_iff}.
        \end{subproof}
        Thus $E\subseteq Y$ by \cref{subseteq_transitive}.
        Thus there exists $a_0$ such that $a_0$ is a largest element of $E$ with respect to $R$
            by \cref{finite_inhabited_largest_element}.
        Take $a_0$ such that $a_0$ is a largest element of $E$ with respect to $R$.
        Show that for all $U\in B$ we have $a_0\notin U$.
        \begin{subproof}
            Fix $U$.
            Assume $U\in B$.
            Thus $U\mathrel{R_B} g(U)$.
            Thus $(U,\apply{g}{U})\in R_B$ by \cref{relation}.
            Thus $(U,\apply{g}{U})\in B\times A$.
            Thus there exists $U_1$ such that $U_1\in B$ and there exists $a$ such that
                $a\in A$ and $(U,\apply{g}{U}) = (U_1,a)$ and $U_1 = \openrayleft{R}{Y}{a}$.
            Take $U_1$ such that $U_1\in B$ and there exists $a$ such that
                $a\in A$ and $(U,\apply{g}{U}) = (U_1,a)$ and $U_1 = \openrayleft{R}{Y}{a}$.
            Take $a$ such that $a\in A$ and $(U,\apply{g}{U}) = (U_1,a)$ and
                $U_1 = \openrayleft{R}{Y}{a}$.
            Thus $U = U_1$ and $g(U) = a$ by \cref{pair_eq_iff}.
            Thus $U = \openrayleft{R}{Y}{a}$ and $g(U) = a$.
            Thus $g$ is a function.
            Thus $U\in\dom{g}$.
            Thus $(U,\apply{g}{U})\in g$ by \cref{function_apply_elim}.
            Thus $(U,a)\in g$.
            Thus $a\in E$ by \cref{img_iff}.
            Thus $a\mathrel{R}a_0$ or $a = a_0$ by \cref{largest_element}.
            Show that $a_0\notin U$.
            \begin{subproof}
                Suppose not.
                Thus $a_0\in U$.
                Thus $a_0\mathrel{R}a$ by \cref{openray_left}.
                $R$ is asymmetric by \cref{simple_order_relation,order_topology}.
                \begin{byCase}
                \caseOf{$a\mathrel{R}a_0$.}
                    Contradiction by \cref{asymmetric}.
                \caseOf{$a = a_0$.}
                    Thus $a_0\mathrel{R}a_0$.
                    Contradiction by \cref{asymmetric}.
                \end{byCase}
            \end{subproof}
        \end{subproof}
        Show that $a_0\notin\unions{B}$.
        \begin{subproof}
            Suppose not.
            Thus $a_0\in\unions{B}$.
            Take $U\in B$ such that $a_0\in U$ by \cref{unions_iff}.
            Contradiction.
        \end{subproof}
        $a_0\in A$ by \cref{largest_element,elem_subseteq}.
        Thus $a_0\in\unions{B}$ by \cref{covering,subseteq}.
        Contradiction.
    \end{subproof}
    Show that there exists $m$ such that $m$ is a smallest element of $A$ with respect to $R$.
    \begin{subproof}
        Suppose not.
        Let $C = \{U \mid \exists a\in A. U = \openrayright{R}{Y}{a}\}$.
        Show that $C$ is a covering of $A$.
        \begin{subproof}
            Show that for all $y$ such that $y\in A$ we have $y\in\unions{C}$.
            \begin{subproof}
                Fix $y$.
                Assume $y\in A$.
                Thus $y$ is not a smallest element of $A$ with respect to $R$.
                Take $a\in A$ such that it is wrong that $y\mathrel{R}a$ or $y = a$ by \cref{smallest_element}.
                Thus $a\neq y$ and it is wrong that $y\mathrel{R}a$.
                Thus $y\in Y$ by \cref{subseteq}.
                Thus $a\in Y$ by \cref{subseteq}.
                $R$ is connex on $Y$ by \cref{order_topology,simple_order_relation}.
                Thus $y\mathrel{R}a$ or $a\mathrel{R}y$ by \cref{connex}.
                \begin{byCase}
                \caseOf{$a\mathrel{R}y$.}
                    Thus $y\in\openrayright{R}{Y}{a}$ by \cref{openray_right}.
                    Thus $\openrayright{R}{Y}{a}\in C$.
                    Thus $y\in\unions{C}$ by \cref{unions_iff}.
                \caseOf{$y\mathrel{R}a$.}
                    Contradiction.
                \end{byCase}
            \end{subproof}
            Thus $A\subseteq\unions{C}$ by \cref{subseteq}.
            Thus $C$ is a covering of $A$ by \cref{covering}.
        \end{subproof}
        Show that for all $U$ such that $U\in C$ we have $U$ is open in $Y$ with respect to $T_Y$.
        \begin{subproof}
            Fix $U$.
            Assume $U\in C$.
            Take $a\in A$ such that $U = \openrayright{R}{Y}{a}$.
            Thus $a\in Y$ by \cref{subseteq}.
            Thus $U$ is open in $Y$ with respect to $T_Y$ by \cref{openrayright_open_in_order_topology}.
        \end{subproof}
        Thus there exists $B$ such that $B\subseteq C$ and $B$ is finite and $B$ is a covering of $A$
            by \cref{compact_subspace_open_covering,compact}.
        Take $B$ such that $B\subseteq C$ and $B$ is finite and $B$ is a covering of $A$.
        Show that $B$ is inhabited.
        \begin{subproof}
            Take $y$ such that $y\in A$.
            Thus $y\in\unions{B}$ by \cref{covering,subseteq}.
            Take $U\in B$ such that $y\in U$ by \cref{unions_iff}.
        \end{subproof}
        Let $R_B = \{w\in B\times A \mid \exists U\in B.\exists a\in A. w = (U,a) \land U = \openrayright{R}{Y}{a}\}$.
        Show that $R_B$ is a relation.
        \begin{subproof}
            Show that for all $w$ such that $w\in R_B$ there exist $U,a$ such that $w = (U,a)$.
            \begin{subproof}
                Fix $w$.
                Assume $w\in R_B$.
                Take $U\in B$ such that there exists $a\in A$ such that $w = (U,a)$ and $U = \openrayright{R}{Y}{a}$.
                Take $a\in A$ such that $w = (U,a)$ and $U = \openrayright{R}{Y}{a}$.
            \end{subproof}
            Thus $R_B$ is a relation by \cref{relation}.
        \end{subproof}
        Show that for all $U\in B$ there exists $a$ such that $U\mathrel{R_B}a$.
        \begin{subproof}
            Fix $U$.
            Assume $U\in B$.
            Thus $U\in C$ by \cref{subseteq}.
            Take $a\in A$ such that $U = \openrayright{R}{Y}{a}$.
            Thus $(U,a)\in R_B$.
            Thus $U\mathrel{R_B} a$ by \cref{relation}.
        \end{subproof}
        Take $g$ such that $g$ is a function on $B$ and for all $U\in B$ we have $U\mathrel{R_B} \apply{g}{U}$
            by \cref{axiom_of_choice}.
        Let $E = \img{g}{B}$.
        Show that $E\subseteq A$.
        \begin{subproof}
            Show that for all $e$ such that $e\in E$ we have $e\in A$.
            \begin{subproof}
                Fix $e$.
                Assume $e\in E$.
                Thus $e\in\img{g}{B}$.
                Take $U\in B$ such that $(U,e)\in g$ by \cref{img_iff}.
                Thus $g$ is a function.
                Thus $g(U) = e$ by \cref{function_apply_intro}.
                Thus $U\mathrel{R_B} g(U)$.
                Thus $(U,\apply{g}{U})\in R_B$ by \cref{relation}.
                Thus $(U,\apply{g}{U})\in B\times A$.
                Thus there exists $U_1$ such that $U_1\in B$ and there exists $a$ such that
                    $a\in A$ and $(U,\apply{g}{U}) = (U_1,a)$ and $U_1 = \openrayright{R}{Y}{a}$.
                Take $U_1$ such that $U_1\in B$ and there exists $a$ such that
                    $a\in A$ and $(U,\apply{g}{U}) = (U_1,a)$ and $U_1 = \openrayright{R}{Y}{a}$.
                Take $a$ such that $a\in A$ and $(U,\apply{g}{U}) = (U_1,a)$ and
                    $U_1 = \openrayright{R}{Y}{a}$.
                Thus $U = U_1$ and $g(U) = a$ by \cref{pair_eq_iff}.
                Thus $U = \openrayright{R}{Y}{a}$ and $g(U) = a$.
                Thus $e\in A$.
            \end{subproof}
            Thus $E\subseteq A$ by \cref{subseteq}.
        \end{subproof}
        Thus $E$ is finite by \cref{finite_image_function}.
        Show that $E$ is inhabited.
        \begin{subproof}
            Take $U_0$ such that $U_0\in B$.
            Thus $g$ is a function.
            Thus $U_0\in\dom{g}$.
            Thus $(U_0,\apply{g}{U_0})\in g$ by \cref{function_apply_elim}.
            Thus $g(U_0)\in E$ by \cref{img_iff}.
        \end{subproof}
        Thus $E\subseteq Y$ by \cref{subseteq_transitive}.
        Thus there exists $a_0$ such that $a_0$ is a smallest element of $E$ with respect to $R$
            by \cref{finite_inhabited_smallest_element}.
        Take $a_0$ such that $a_0$ is a smallest element of $E$ with respect to $R$.
        Show that for all $U\in B$ we have $a_0\notin U$.
        \begin{subproof}
            Fix $U$.
            Assume $U\in B$.
            Thus $U\mathrel{R_B} g(U)$.
            Thus $(U,\apply{g}{U})\in R_B$ by \cref{relation}.
            Thus $(U,\apply{g}{U})\in B\times A$.
            Thus there exists $U_1$ such that $U_1\in B$ and there exists $a$ such that
                $a\in A$ and $(U,\apply{g}{U}) = (U_1,a)$ and $U_1 = \openrayright{R}{Y}{a}$.
            Take $U_1$ such that $U_1\in B$ and there exists $a$ such that
                $a\in A$ and $(U,\apply{g}{U}) = (U_1,a)$ and $U_1 = \openrayright{R}{Y}{a}$.
            Take $a$ such that $a\in A$ and $(U,\apply{g}{U}) = (U_1,a)$ and
                $U_1 = \openrayright{R}{Y}{a}$.
            Thus $U = U_1$ and $g(U) = a$ by \cref{pair_eq_iff}.
            Thus $U = \openrayright{R}{Y}{a}$ and $g(U) = a$.
            Thus $g$ is a function.
            Thus $U\in\dom{g}$.
            Thus $(U,\apply{g}{U})\in g$ by \cref{function_apply_elim}.
            Thus $(U,a)\in g$.
            Thus $a\in E$ by \cref{img_iff}.
            Thus $a_0\mathrel{R}a$ or $a_0 = a$ by \cref{smallest_element}.
            Show that $a_0\notin U$.
            \begin{subproof}
                Suppose not.
                Thus $a_0\in U$.
                Thus $a\mathrel{R}a_0$ by \cref{openray_right}.
                $R$ is asymmetric by \cref{simple_order_relation,order_topology}.
                \begin{byCase}
                \caseOf{$a_0\mathrel{R}a$.}
                    Contradiction by \cref{asymmetric}.
                \caseOf{$a_0 = a$.}
                    Thus $a_0\mathrel{R}a_0$.
                    Contradiction by \cref{asymmetric}.
                \end{byCase}
            \end{subproof}
        \end{subproof}
        Show that $a_0\notin\unions{B}$.
        \begin{subproof}
            Suppose not.
            Thus $a_0\in\unions{B}$.
            Take $U\in B$ such that $a_0\in U$ by \cref{unions_iff}.
            Contradiction.
        \end{subproof}
        $a_0\in A$ by \cref{smallest_element,elem_subseteq}.
        Thus $a_0\in\unions{B}$ by \cref{covering,subseteq}.
        Contradiction.
    \end{subproof}
    Take $m$ such that $m$ is a smallest element of $A$ with respect to $R$.
    Thus $m\in A$ by \cref{smallest_element}.
    Thus $m\in\img{f}{X}$.
    Take $c\in X$ such that $(c,m)\in f$ by \cref{img_iff}.
    Thus $f$ is a function by \cref{funs_is_function}.
    Thus $f(c) = m$ by \cref{function_apply_intro}.
    Thus $\apply{f}{c}$ is a smallest element of $A$ with respect to $R$ by \cref{smallest_element}.
    Take $M$ such that $M$ is a largest element of $A$ with respect to $R$.
    Thus $M\in A$ by \cref{largest_element}.
    Thus $M\in\img{f}{X}$.
    Take $d\in X$ such that $(d,M)\in f$ by \cref{img_iff}.
    Thus $f$ is a function by \cref{funs_is_function}.
    Thus $f(d) = M$ by \cref{function_apply_intro}.
    Thus $\apply{f}{d}$ is a largest element of $A$ with respect to $R$ by \cref{largest_element}.
\end{proof}

% from §27 helper definition 4: distance set from a point to a set
\begin{definition}\label{point_distance_set}
    $\pointDistanceSet{d}{A}{x} = \{d((x,a)) \mid a\in A\}$.
\end{definition}

% from §27 Item 5: distance from a point to a set
\begin{definition}\label{point_set_distance}
    $r$ is the distance from $x$ to $A$ in $X$ with respect to $d$ iff
        $x\in X$ and $A\subseteq X$ and $r\in\reals$ and
        $r$ is an infimum of $\pointDistanceSet{d}{A}{x}$.
\end{definition}

% from §27 helper lemma 3: distance-to-set map is continuous
\begin{lemma}\label{point_set_distance_continuous}
    Suppose $d$ is a metric on $X$.
    Suppose $A\subseteq X$.
    Suppose $A\neq\emptyset$.
    Let $T = \metricTopology{d}{X}$.
    Then there exists $f$ such that
        $f$ is a function from $X$ to $\reals$ and
        for all $x\in X$ we have $f(x)$ is the distance from $x$ to $A$ in $X$ with respect to $d$ and
        $f$ is continuous from $X$ to $\reals$ with respect to $T$ and $\standardTopologyR$.
\end{lemma}
\begin{proof}
    Let $R = \{w\in X\times\reals \mid \text{there exists $x\in X$ such that there exists $r\in\reals$ such that
        $w = (x,r)$ and $r$ is the distance from $x$ to $A$ in $X$ with respect to $d$}\}$.
    Show that $R$ is a relation.
    \begin{subproof}
        Show that for all $w$ such that $w\in R$ there exist $x,r$ such that $w = (x,r)$.
        \begin{subproof}
            Fix $w$.
            Assume $w\in R$.
            Take $x\in X$ such that there exists $r\in\reals$ such that
                $w = (x,r)$ and $r$ is the distance from $x$ to $A$ in $X$ with respect to $d$.
            Take $r\in\reals$ such that $w = (x,r)$ and
                $r$ is the distance from $x$ to $A$ in $X$ with respect to $d$.
        \end{subproof}
        Thus $R$ is a relation by \cref{relation}.
    \end{subproof}
    Show that for all $x\in X$ there exists $r$ such that $x\mathrel{R} r$.
    \begin{subproof}
        Fix $x$.
        Assume $x\in X$.
        Let $S = \pointDistanceSet{d}{A}{x}$.
        Show that $S\subseteq\reals$.
        \begin{subproof}
            Show that for all $s$ such that $s\in S$ we have $s\in\reals$.
            \begin{subproof}
                Fix $s$.
                Assume $s\in S$.
                Take $a\in A$ such that $s = d((x,a))$ by \cref{point_distance_set}.
                Thus $a\in X$ by \cref{subseteq}.
                $(x,a)\in X\times X$ by \cref{times_tuple_intro}.
                $d\in\funs{X\times X}{\reals}$ by \cref{metric_on}.
                Thus $s\in\reals$ by \cref{funs_type_apply}.
            \end{subproof}
            Thus $S\subseteq\reals$ by \cref{subseteq}.
        \end{subproof}
        Show that $S\neq\emptyset$.
        \begin{subproof}
            Take $a_0$ such that $a_0\in A$ by \cref{nonempty_inhabited}.
            Thus $d((x,a_0))\in S$ by \cref{point_distance_set}.
        \end{subproof}
        Show that $S$ is bounded below.
        \begin{subproof}
            Show that $0$ is a lower bound of $S$.
            \begin{subproof}
                $0\in\reals$ by \cref{real_add_ident}.
                Show that for all $s\in S$ we have $0 \leq s$.
                \begin{subproof}
                    Fix $s$.
                    Assume $s\in S$.
                    Take $a\in A$ such that $s = d((x,a))$ by \cref{point_distance_set}.
                    Thus $a\in X$ by \cref{subseteq}.
                    $d$ is nonnegative on $X$ by \cref{metric_on}.
                    Thus $s \geq 0$ by \cref{metric_nonnegative}.
                \end{subproof}
                Thus $0$ is a lower bound of $S$ by \cref{real_lower_bound,subseteq}.
            \end{subproof}
            Thus $S$ is bounded below by \cref{real_bounded_below}.
        \end{subproof}
        Take $r$ such that $r$ is an infimum of $S$ by \cref{real_infimum_exists}.
        Thus $r$ is a lower bound of $S$ by \cref{real_infimum}.
        Thus $r\in\reals$ by \cref{real_lower_bound}.
        Thus $r$ is the distance from $x$ to $A$ in $X$ with respect to $d$ by \cref{point_set_distance}.
        Thus $(x,r)\in R$.
        Thus $x\mathrel{R} r$ by \cref{relation}.
    \end{subproof}
    Take $f$ such that $f$ is a function on $X$ and for all $x\in X$ we have $x\mathrel{R} f(x)$
        by \cref{axiom_of_choice}.
    Show that for all $x\in X$ we have $f(x)$ is the distance from $x$ to $A$ in $X$ with respect to $d$.
    \begin{subproof}
        Fix $x$.
        Assume $x\in X$.
        Thus $(x,\apply{f}{x})\in R$ by \cref{relation}.
        Thus $(x,\apply{f}{x})\in X\times\reals$ and there exists $x_1\in X$ such that
            there exists $r\in\reals$ such that $(x,\apply{f}{x}) = (x_1,r)$ and
            $r$ is the distance from $x_1$ to $A$ in $X$ with respect to $d$.
        Take $x_1\in X$ such that there exists $r\in\reals$ such that
            $(x,\apply{f}{x}) = (x_1,r)$ and $r$ is the distance from $x_1$ to $A$ in $X$ with respect to $d$.
        Take $r\in\reals$ such that $(x,\apply{f}{x}) = (x_1,r)$ and
            $r$ is the distance from $x_1$ to $A$ in $X$ with respect to $d$.
        Thus $x = x_1$ and $f(x) = r$ by \cref{pair_eq_iff}.
        Thus $f(x)$ is the distance from $x$ to $A$ in $X$ with respect to $d$.
    \end{subproof}
    Show that $f$ is a function from $X$ to $\reals$.
    \begin{subproof}
        Thus $f$ is a function.
        Thus $X = \dom{f}$.
        Show that for all $x\in\dom{f}$ we have $f(x)\in\reals$.
        \begin{subproof}
            Fix $x$.
            Assume $x\in\dom{f}$.
            Thus $x\in X$.
            Thus $f(x)$ is the distance from $x$ to $A$ in $X$ with respect to $d$.
            Thus $f(x)\in\reals$ by \cref{point_set_distance}.
        \end{subproof}
        Thus $f$ is a function to $\reals$.
        Thus $f\in\funs{X}{\reals}$ by \cref{funs_intro}.
        Thus $f$ is a function from $X$ to $\reals$.
    \end{subproof}
    Let $T_Y = \metricTopology{\standardMetricR}{\reals}$.
    Show that $f$ is continuous from $X$ to $\reals$ with respect to $T$ and $T_Y$.
    \begin{subproof}
        Show that for all $x\in X$ we have
            for all $\epsilon\in\reals$ such that $0 < \epsilon$
            there exists $\delta\in\reals$ such that $0 < \delta$ and
            for all $y\in X$ if $d((x,y)) < \delta$ then
                $\apply{\standardMetricR}{(f(x),f(y))} < \epsilon$.
        \begin{subproof}
            Fix $x$.
            Assume $x\in X$.
            Fix $\epsilon\in\reals$.
            Assume $0 < \epsilon$.
            Let $\delta = \epsilon$.
            Thus $0 < \delta$.
            Show that for all $y\in X$ if $d((x,y)) < \delta$ then
                $\apply{\standardMetricR}{(f(x),f(y))} < \epsilon$.
            \begin{subproof}
            Fix $y$.
            Assume $y\in X$.
            Assume $d((x,y)) < \delta$.
            Thus $\delta = \epsilon$.
            Thus $d((x,y)) < \epsilon$ by \cref{real_lt_subst_right}.
            Let $S_x = \pointDistanceSet{d}{A}{x}$.
            Let $S_y = \pointDistanceSet{d}{A}{y}$.
                $f(x)$ is an infimum of $S_x$ by \cref{point_set_distance}.
                $f(y)$ is an infimum of $S_y$ by \cref{point_set_distance}.
                Thus $f(x)$ is a lower bound of $S_x$ by \cref{real_infimum}.
                Thus $f(y)$ is a lower bound of $S_y$ by \cref{real_infimum}.
                Show that $f(x) \geq f(y) - d((x,y))$.
                \begin{subproof}
                    Let $b = f(y) - d((x,y))$.
                    Show that $b$ is a lower bound of $S_x$.
                    \begin{subproof}
                        $f(y)\in\reals$ by \cref{point_set_distance}.
                        $(x,y)\in X\times X$ by \cref{times_tuple_intro}.
                        $d\in\funs{X\times X}{\reals}$ by \cref{metric_on}.
                        Thus $d((x,y))\in\reals$ by \cref{funs_type_apply}.
                        $b = f(y) + \neg{(d((x,y)))}$ by \cref{real_sub}.
                        $\neg{(d((x,y)))}\in\reals$ by \cref{real_add_inverse}.
                        Thus $b\in\reals$ by \cref{real_add_closed}.
                        Show that for all $s\in S_x$ we have $b \leq s$.
                        \begin{subproof}
                            Fix $s$.
                            Assume $s\in S_x$.
                            Take $a\in A$ such that $s = d((x,a))$ by \cref{point_distance_set}.
                            Thus $a\in X$ by \cref{subseteq}.
                            $d$ is symmetric on $X$ by \cref{metric_on}.
                            Thus $d((y,x)) = d((x,y))$ by \cref{metric_symmetric}.
                            $d$ satisfies the triangle inequality on $X$ by \cref{metric_on}.
                            Thus $d((y,x)) + d((x,a)) \geq d((y,a))$ by \cref{metric_triangle_inequality}.
                            Thus $d((y,a)) \leq d((y,x)) + d((x,a))$.
                            Thus $d((y,a)) \leq d((x,y)) + d((x,a))$.
                            $(y,a)\in X\times X$ by \cref{times_tuple_intro}.
                            $(x,y)\in X\times X$ by \cref{times_tuple_intro}.
                            $(x,a)\in X\times X$ by \cref{times_tuple_intro}.
                            Thus $d((y,a))\in\reals$ by \cref{funs_type_apply}.
                            Thus $d((x,y))\in\reals$ by \cref{funs_type_apply}.
                            Thus $d((x,a))\in\reals$ by \cref{funs_type_apply}.
                            Thus $d((x,y)) + d((x,a))\in\reals$ by \cref{real_add_closed}.
                            Thus $f(y)\in\reals$ by \cref{point_set_distance}.
                            Thus $f(y) \leq d((y,a))$ by \cref{real_lower_bound,point_distance_set}.
                            Thus $f(y) \leq d((x,y)) + d((x,a))$ by \cref{real_leq_trans}.
                            Let $u = d((x,y))$.
                            Let $v = d((x,a))$.
                            Thus $f(y) \leq u + v$ by \cref{real_leq_trans}.
                            Thus $f(y) + \neg{u} \leq (u + v) + \neg{u}$ by \cref{real_leq_add}.
                            Show that $(u + v) + \neg{u} = v$.
                            \begin{subproof}
                                Thus $u\in\reals$.
                                Thus $v\in\reals$.
                                $u + \neg{u} = 0$ by \cref{real_add_inverse}.
                                $(u + v) + \neg{u} = v + (u + \neg{u})$ by \cref{real_add_assoc,real_add_comm}.
                                Thus $(u + v) + \neg{u} = v + 0$ by \cref{real_add_inverse}.
                                $0\in\reals$ by \cref{real_add_ident}.
                                $v + 0 = 0 + v$ by \cref{real_add_comm}.
                                $0 + v = v$ by \cref{real_add_ident}.
                                Thus $(u + v) + \neg{u} = v$ by \cref{real_add_comm,real_add_ident}.
                            \end{subproof}
                            Thus $f(y) + \neg{u} \leq v$ by \cref{real_leq_trans}.
                            Thus $f(y) - u \leq v$ by \cref{real_sub}.
                            Thus $f(y) - d((x,y)) \leq d((x,a))$.
                            Thus $b \leq s$.
                        \end{subproof}
                        Thus $b$ is a lower bound of $S_x$ by \cref{real_lower_bound,subseteq}.
                    \end{subproof}
                    Thus $f(x) \geq b$ by \cref{real_infimum}.
                \end{subproof}
                Show that $f(y) \geq f(x) - d((x,y))$.
                \begin{subproof}
                    Let $b = f(x) - d((x,y))$.
                    Show that $b$ is a lower bound of $S_y$.
                    \begin{subproof}
                        $f(x)\in\reals$ by \cref{point_set_distance}.
                        $(x,y)\in X\times X$ by \cref{times_tuple_intro}.
                        $d\in\funs{X\times X}{\reals}$ by \cref{metric_on}.
                        Thus $d((x,y))\in\reals$ by \cref{funs_type_apply}.
                        $b = f(x) + \neg{(d((x,y)))}$ by \cref{real_sub}.
                        $\neg{(d((x,y)))}\in\reals$ by \cref{real_add_inverse}.
                        Thus $b\in\reals$ by \cref{real_add_closed}.
                        Show that for all $s\in S_y$ we have $b \leq s$.
                        \begin{subproof}
                            Fix $s$.
                            Assume $s\in S_y$.
                            Take $a\in A$ such that $s = d((y,a))$ by \cref{point_distance_set}.
                            Thus $a\in X$ by \cref{subseteq}.
                            $d$ is symmetric on $X$ by \cref{metric_on}.
                            Thus $d((x,y)) = d((y,x))$ by \cref{metric_symmetric}.
                            $d$ satisfies the triangle inequality on $X$ by \cref{metric_on}.
                            Thus $d((x,y)) + d((y,a)) \geq d((x,a))$ by \cref{metric_triangle_inequality}.
                            Thus $d((x,a)) \leq d((x,y)) + d((y,a))$.
                            $(x,a)\in X\times X$ by \cref{times_tuple_intro}.
                            $(x,y)\in X\times X$ by \cref{times_tuple_intro}.
                            $(y,a)\in X\times X$ by \cref{times_tuple_intro}.
                            Thus $d((x,a))\in\reals$ by \cref{funs_type_apply}.
                            Thus $d((x,y))\in\reals$ by \cref{funs_type_apply}.
                            Thus $d((y,a))\in\reals$ by \cref{funs_type_apply}.
                            Thus $d((x,y)) + d((y,a))\in\reals$ by \cref{real_add_closed}.
                            Thus $f(x)\in\reals$ by \cref{point_set_distance}.
                            Thus $f(x) \leq d((x,a))$ by \cref{real_lower_bound,point_distance_set}.
                            Thus $f(x) \leq d((x,y)) + d((y,a))$ by \cref{real_leq_trans}.
                            Let $u = d((x,y))$.
                            Let $v = d((y,a))$.
                            Thus $f(x) \leq u + v$ by \cref{real_leq_trans}.
                            Thus $f(x) + \neg{u} \leq (u + v) + \neg{u}$ by \cref{real_leq_add}.
                            Show that $(u + v) + \neg{u} = v$.
                            \begin{subproof}
                                Thus $u\in\reals$.
                                Thus $v\in\reals$.
                                $u + \neg{u} = 0$ by \cref{real_add_inverse}.
                                $(u + v) + \neg{u} = v + (u + \neg{u})$ by \cref{real_add_assoc,real_add_comm}.
                                Thus $(u + v) + \neg{u} = v + 0$ by \cref{real_add_inverse}.
                                $0\in\reals$ by \cref{real_add_ident}.
                                $v + 0 = 0 + v$ by \cref{real_add_comm}.
                                $0 + v = v$ by \cref{real_add_ident}.
                                Thus $(u + v) + \neg{u} = v$ by \cref{real_add_comm,real_add_ident}.
                            \end{subproof}
                            Thus $f(x) + \neg{u} \leq v$ by \cref{real_leq_trans}.
                            Thus $f(x) - u \leq v$ by \cref{real_sub}.
                            Thus $f(x) - d((x,y)) \leq d((y,a))$.
                            Thus $b \leq s$.
                        \end{subproof}
                        Thus $b$ is a lower bound of $S_y$ by \cref{real_lower_bound,subseteq}.
                    \end{subproof}
                    Thus $f(y) \geq b$ by \cref{real_infimum}.
                \end{subproof}
                Show that $f(x) \leq f(y) + d((x,y))$.
                \begin{subproof}
                    Thus $f(x) - d((x,y)) \leq f(y)$.
                    $f(x)\in\reals$ by \cref{point_set_distance}.
                    $f(y)\in\reals$ by \cref{point_set_distance}.
                    $(x,y)\in X\times X$ by \cref{times_tuple_intro}.
                    $d\in\funs{X\times X}{\reals}$ by \cref{metric_on}.
                    Thus $d((x,y))\in\reals$ by \cref{funs_type_apply}.
                    $\neg{(d((x,y)))}\in\reals$ by \cref{real_add_inverse}.
                    Thus $f(x) + \neg{(d((x,y)))}\in\reals$ by \cref{real_add_closed}.
                    Thus $f(x) - d((x,y))\in\reals$ by \cref{real_sub}.
                    Thus $(f(x) - d((x,y))) + d((x,y)) \leq f(y) + d((x,y))$ by \cref{real_leq_add}.
                    Let $s_0 = (f(x) - d((x,y))) + d((x,y))$.
                    Show that $s_0 = f(x)$.
                    \begin{subproof}
                        $f(x) - d((x,y)) = f(x) + \neg{(d((x,y)))}$ by \cref{real_sub}.
                        Thus $s_0 = f(x) + (\neg{(d((x,y)))} + d((x,y)))$ by \cref{real_add_assoc}.
                        $\neg{(d((x,y)))} + d((x,y)) = 0$ by \cref{real_add_inverse,real_add_comm}.
                        Thus $s_0 = f(x) + 0$ by \cref{real_add_assoc}.
                        $0\in\reals$ by \cref{real_add_ident}.
                        $f(x)\in\reals$ by \cref{point_set_distance}.
                        $f(x) + 0 = 0 + f(x)$ by \cref{real_add_comm}.
                        $0 + f(x) = f(x)$ by \cref{real_add_ident}.
                        Thus $s_0 = f(x)$ by \cref{real_add_comm,real_add_ident}.
                    \end{subproof}
                    Thus $f(x) \leq s_0$.
                    Thus $f(x) \leq f(y) + d((x,y))$ by \cref{real_leq_trans}.
                \end{subproof}
                Show that $f(y) \leq f(x) + d((x,y))$.
                \begin{subproof}
                    Thus $f(y) - d((x,y)) \leq f(x)$.
                    $f(y)\in\reals$ by \cref{point_set_distance}.
                    $f(x)\in\reals$ by \cref{point_set_distance}.
                    $(x,y)\in X\times X$ by \cref{times_tuple_intro}.
                    $d\in\funs{X\times X}{\reals}$ by \cref{metric_on}.
                    Thus $d((x,y))\in\reals$ by \cref{funs_type_apply}.
                    $\neg{(d((x,y)))}\in\reals$ by \cref{real_add_inverse}.
                    Thus $f(y) + \neg{(d((x,y)))}\in\reals$ by \cref{real_add_closed}.
                    Thus $f(y) - d((x,y))\in\reals$ by \cref{real_sub}.
                    Thus $(f(y) - d((x,y))) + d((x,y)) \leq f(x) + d((x,y))$ by \cref{real_leq_add}.
                    Let $s_1 = (f(y) - d((x,y))) + d((x,y))$.
                    Show that $s_1 = f(y)$.
                    \begin{subproof}
                        $f(y) - d((x,y)) = f(y) + \neg{(d((x,y)))}$ by \cref{real_sub}.
                        Thus $s_1 = f(y) + (\neg{(d((x,y)))} + d((x,y)))$ by \cref{real_add_assoc}.
                        $\neg{(d((x,y)))} + d((x,y)) = 0$ by \cref{real_add_inverse,real_add_comm}.
                        Thus $s_1 = f(y) + 0$ by \cref{real_add_assoc}.
                        $0\in\reals$ by \cref{real_add_ident}.
                        $f(y)\in\reals$ by \cref{point_set_distance}.
                        $f(y) + 0 = 0 + f(y)$ by \cref{real_add_comm}.
                        $0 + f(y) = f(y)$ by \cref{real_add_ident}.
                        Thus $s_1 = f(y)$ by \cref{real_add_comm,real_add_ident}.
                    \end{subproof}
                    Thus $f(y) \leq s_1$.
                    Thus $f(y) \leq f(x) + d((x,y))$ by \cref{real_leq_trans}.
                \end{subproof}
                Show that $\abs{(f(x) - f(y))} \leq d((x,y))$.
                \begin{subproof}
                    $f(x)\in\reals$ by \cref{point_set_distance}.
                    $f(y)\in\reals$ by \cref{point_set_distance}.
                    $(x,y)\in X\times X$ by \cref{times_tuple_intro}.
                    $d\in\funs{X\times X}{\reals}$ by \cref{metric_on}.
                    Thus $d((x,y))\in\reals$ by \cref{funs_type_apply}.
                    $\neg{(f(y))}\in\reals$ by \cref{real_add_inverse}.
                    Thus $f(x) + \neg{(f(y))}\in\reals$ by \cref{real_add_closed}.
                    $f(x) - f(y)\in\reals$ by \cref{real_sub}.
                    $\neg{(f(x) - f(y))}\in\reals$ by \cref{real_add_inverse}.

                    Show that $f(x) - f(y) \leq d((x,y))$.
                    \begin{subproof}
                        Thus $f(y) + d((x,y))\in\reals$ by \cref{real_add_closed}.
                        Thus $f(x) + \neg{(f(y))} \leq (f(y) + d((x,y))) + \neg{(f(y))}$ by \cref{real_leq_add}.
                        $(f(y) + d((x,y))) + \neg{(f(y))} = f(y) + (d((x,y)) + \neg{(f(y))})$ by \cref{real_add_assoc}.
                        $d((x,y)) + \neg{(f(y))} = \neg{(f(y))} + d((x,y))$ by \cref{real_add_comm}.
                        $f(y) + (\neg{(f(y))} + d((x,y))) = (f(y) + \neg{(f(y))}) + d((x,y))$ by \cref{real_add_assoc}.
                        $f(y) + \neg{(f(y))} = 0$ by \cref{real_add_inverse}.
                        $0\in\reals$ by \cref{real_add_ident}.
                        Thus $(f(y) + d((x,y))) + \neg{(f(y))} = d((x,y))$ by \cref{real_add_ident,real_add_assoc,real_add_comm}.
                        Thus $f(x) + \neg{(f(y))} \leq d((x,y))$ by \cref{real_leq_trans}.
                        Thus $f(x) - f(y) \leq d((x,y))$ by \cref{real_sub}.
                    \end{subproof}

                    Show that $\neg{(f(x) - f(y))} \leq d((x,y))$.
                    \begin{subproof}
                        Thus $f(x) + d((x,y))\in\reals$ by \cref{real_add_closed}.
                        $\neg{(f(x))}\in\reals$ by \cref{real_add_inverse}.
                        Thus $f(y) + \neg{(f(x))} \leq (f(x) + d((x,y))) + \neg{(f(x))}$ by \cref{real_leq_add}.
                        $(f(x) + d((x,y))) + \neg{(f(x))} = f(x) + (d((x,y)) + \neg{(f(x))})$ by \cref{real_add_assoc}.
                        $d((x,y)) + \neg{(f(x))} = \neg{(f(x))} + d((x,y))$ by \cref{real_add_comm}.
                        $f(x) + (\neg{(f(x))} + d((x,y))) = (f(x) + \neg{(f(x))}) + d((x,y))$ by \cref{real_add_assoc}.
                        $f(x) + \neg{(f(x))} = 0$ by \cref{real_add_inverse}.
                        $0\in\reals$ by \cref{real_add_ident}.
                        Thus $(f(x) + d((x,y))) + \neg{(f(x))} = d((x,y))$ by \cref{real_add_ident,real_add_assoc,real_add_comm}.
                        Thus $f(y) + \neg{(f(x))} \leq d((x,y))$ by \cref{real_leq_trans}.
                        Thus $f(y) - f(x) \leq d((x,y))$ by \cref{real_sub}.
                        Thus $f(y) - f(x) = \neg{(f(x) - f(y))}$ by \cref{real_sub_swap_neg}.
                        Thus $\neg{(f(x) - f(y))} \leq d((x,y))$.
                    \end{subproof}

                    Thus $\abs{(f(x) - f(y))} \leq d((x,y))$ by \cref{real_abs_leq_if_pm_leq}.
                \end{subproof}
                $f(x)\in\reals$ by \cref{point_set_distance}.
                $f(y)\in\reals$ by \cref{point_set_distance}.
                $\neg{(f(y))}\in\reals$ by \cref{real_add_inverse}.
                Thus $f(x) + \neg{(f(y))}\in\reals$ by \cref{real_add_closed}.
                $f(x) - f(y)\in\reals$ by \cref{real_sub}.
                $(f(x),f(y))\in\reals\times\reals$ by \cref{times_tuple_intro}.
                $\abs{(f(x) - f(y))}\in\reals$ by \cref{real_abs_in_reals}.
                Let $p = (f(x),f(y))$.
                $\fst{p} = f(x)$ and $\snd{p} = f(y)$ by \cref{fst_eq,snd_eq}.
                Thus $\abs{(\fst{p} - \snd{p})} = \abs{(f(x) - f(y))}$.
                Thus $(p,\abs{(f(x) - f(y))})\in\standardMetricR$ by \cref{standard_metric_r}.
                $\standardMetricR$ is a metric on $\reals$ by \cref{standard_metric_is_metric}.
                Thus $\standardMetricR\in\funs{\reals\times\reals}{\reals}$ by \cref{metric_on}.
                Thus $\standardMetricR$ is a function by \cref{funs_elim}.
                Thus $\apply{\standardMetricR}{(f(x),f(y))} = \abs{(f(x) - f(y))}$ by \cref{function_apply_intro}.
                Thus $\apply{\standardMetricR}{(f(x),f(y))} \leq d((x,y))$.
                $(x,y)\in X\times X$ by \cref{times_tuple_intro}.
                $d\in\funs{X\times X}{\reals}$ by \cref{metric_on}.
                Thus $d((x,y))\in\reals$ by \cref{funs_type_apply}.
                Thus $\apply{\standardMetricR}{(f(x),f(y))}\in\reals$ by \cref{funs_type_apply}.
                Thus $\apply{\standardMetricR}{(f(x),f(y))} < d((x,y))$ or
                    $\apply{\standardMetricR}{(f(x),f(y))} = d((x,y))$ by \cref{real_leq_split}.
                \begin{byCase}
                \caseOf{$\apply{\standardMetricR}{(f(x),f(y))} < d((x,y))$.}
                    Thus $\apply{\standardMetricR}{(f(x),f(y))} < \epsilon$ by \cref{real_lt_trans}.
                \caseOf{$\apply{\standardMetricR}{(f(x),f(y))} = d((x,y))$.}
                    Thus $\apply{\standardMetricR}{(f(x),f(y))} < \epsilon$ by \cref{real_lt_subst_left}.
                \end{byCase}
            \end{subproof}
        \end{subproof}
        $\standardMetricR$ is a metric on $\reals$ by \cref{standard_metric_is_metric}.
        Thus $f$ is continuous from $X$ to $\reals$ with respect to $T$ and $T_Y$
            by \cref{metric_continuity_eps_delta}.
    \end{subproof}
    Thus $f$ is continuous from $X$ to $\reals$ with respect to $T$ and $\standardTopologyR$
        by \cref{standard_metric_topology,continuous_eq_codomain}.
\end{proof}

% from §27 Item 6: Lebesgue number lemma
\begin{lemma}\label{lebesgue_number_lemma}
    Suppose $d$ is a metric on $X$.
    Let $T = \metricTopology{d}{X}$.
    Suppose $A$ is an open covering of $X$ with respect to $T$.
    Suppose $X$ is compact with respect to $T$.
    Then there exists $\delta\in\reals$ such that $0 < \delta$ and
        for all $B$ such that $B\subseteq X$ and
            for all $x,y\in B$ we have $d((x,y)) < \delta$
        there exists $U\in A$ such that $B\subseteq U$.
\end{lemma}
\begin{proof}
    Omitted. % do not touch
\end{proof}

% from §27 Item 7: Lebesgue number
\begin{definition}\label{lebesgue_number}
    $\delta$ is a Lebesgue number for $A$ in $X$ with respect to $d$ iff
        $d$ is a metric on $X$ and
        $A$ is an open covering of $X$ with respect to $\metricTopology{d}{X}$ and
        $\delta\in\reals$ and $0 < \delta$ and
        for all $B$ such that $B\subseteq X$ and
            for all $x,y\in B$ we have $d((x,y)) < \delta$
        there exists $U\in A$ such that $B\subseteq U$.
\end{definition}

% from §27 Item 8: uniformly continuous function
\begin{definition}\label{uniformly_continuous}
    $f$ is uniformly continuous from $X$ to $Y$ with respect to $d_X$ and $d_Y$ iff
        $f\in\funs{X}{Y}$ and
        $d_X$ is a metric on $X$ and
        $d_Y$ is a metric on $Y$ and
        for all $\epsilon\in\reals$ such that $0 < \epsilon$
            there exists $\delta\in\reals$ such that $0 < \delta$ and
                for all $x_0,x_1\in X$ such that $d_X((x_0,x_1)) < \delta$
                    we have $d_Y((\apply{f}{x_0},\apply{f}{x_1})) < \epsilon$.
\end{definition}

% from §27 Item 9: uniform continuity theorem
\begin{theorem}\label{uniform_continuity_theorem}
    Suppose $d_X$ is a metric on $X$.
    Suppose $d_Y$ is a metric on $Y$.
    Let $T_X = \metricTopology{d_X}{X}$.
    Let $T_Y = \metricTopology{d_Y}{Y}$.
    Suppose $f$ is continuous from $X$ to $Y$ with respect to $T_X$ and $T_Y$.
    Suppose $X$ is compact with respect to $T_X$.
    Then $f$ is uniformly continuous from $X$ to $Y$ with respect to $d_X$ and $d_Y$.
\end{theorem}
\begin{proof}
    $f$ is a function from $X$ to $Y$ by \cref{continuous}.
    Thus $f\in\funs{X}{Y}$ by \cref{funs}.
    Show that for all $\epsilon\in\reals$ such that $0 < \epsilon$
        there exists $\delta\in\reals$ such that $0 < \delta$ and
            for all $x_0,x_1\in X$ such that $d_X((x_0,x_1)) < \delta$
                we have $d_Y((\apply{f}{x_0},\apply{f}{x_1})) < \epsilon$.
    \begin{subproof}
        Fix $\epsilon\in\reals$.
        Assume $0 < \epsilon$.
        Let $\epsilon_0 = \epsilon / (1 + 1)$.
        $1\in\reals$ by \cref{real_mul_ident}.
        $0\in\reals$ by \cref{real_add_ident}.
        $0 < 1$ by \cref{real_zero_lt_one}.
        $0 + 1 < 1 + 1$ by \cref{real_lt_add}.
        $0 + 1 = 1$ by \cref{real_add_ident}.
        Thus $1 < 1 + 1$.
        $(1 + 1)\in\reals$ by \cref{real_add_closed}.
        Thus $0 < 1 + 1$ by \cref{real_lt_trans}.
        Thus $1 + 1$ is positive.
        Thus $(1 + 1)\neq 0$ by \cref{real_pos_nonzero}.
        Thus $\inv{(1 + 1)}\in\reals$ by \cref{real_mul_inverse}.
        Thus $\inv{(1 + 1)} > 0$ by \cref{real_inv_pos}.
        $\epsilon_0 = \epsilon \rmul \inv{(1 + 1)}$ by \cref{real_div}.
        Thus $\epsilon_0\in\reals$ by \cref{real_mul_closed}.
        $0\in\reals$ by \cref{real_add_ident}.
        $0 \rmul \inv{(1 + 1)} = 0$ by \cref{real_mul_zero}.
        $0 \rmul \inv{(1 + 1)} < \epsilon \rmul \inv{(1 + 1)}$ by \cref{real_lt_mul_pos}.
        Thus $0 < \epsilon \rmul \inv{(1 + 1)}$ by \cref{real_lt_subst_left}.
        Thus $0 < \epsilon_0$ by \cref{real_lt_subst_left}.

        Let $R = \{w\in X\times\reals \mid \text{$0 < \snd{w}$ and
            for all $y\in X$ if $d_X((\fst{w},y)) < \snd{w}$ then
            $d_Y((\apply{f}{\fst{w}},\apply{f}{y})) < \epsilon_0$}\}$.
        Show that $R$ is a relation.
        \begin{subproof}
            Show that for all $w$ such that $w\in R$ there exist $x,r$ such that $w = (x,r)$.
            \begin{subproof}
                Fix $w$.
                Assume $w\in R$.
                Thus $w\in X\times\reals$.
                Thus there exist $x,r$ such that $w = (x,r)$ by \cref{times_elem_is_tuple}.
            \end{subproof}
            Thus $R$ is a relation by \cref{relation}.
        \end{subproof}
        Show that for all $x\in X$ there exists $\delta\in\reals$ such that $x\mathrel{R}\delta$.
        \begin{subproof}
            Fix $x\in X$.
            Thus for all $\eta\in\reals$ such that $0 < \eta$
                there exists $\delta\in\reals$ such that $0 < \delta$ and
                    for all $y\in X$ if $d_X((x,y)) < \delta$ then
                        $d_Y((\apply{f}{x},\apply{f}{y})) < \eta$
                by \cref{metric_continuity_eps_delta}.
            Thus there exists $\delta\in\reals$ such that $0 < \delta$ and
                for all $y\in X$ if $d_X((x,y)) < \delta$ then
                    $d_Y((\apply{f}{x},\apply{f}{y})) < \epsilon_0$.
            Take $\delta\in\reals$ such that $0 < \delta$ and
                for all $y\in X$ if $d_X((x,y)) < \delta$ then
                    $d_Y((\apply{f}{x},\apply{f}{y})) < \epsilon_0$.
            $(x,\delta)\in X\times\reals$ by \cref{times_tuple_intro}.
            $\fst{(x,\delta)} = x$ and $\snd{(x,\delta)} = \delta$ by \cref{fst_eq,snd_eq}.
            Thus $(x,\delta)\in R$.
            Thus $x\mathrel{R}\delta$.
        \end{subproof}
        Take $g$ such that $g$ is a function on $X$ and for all $x\in X$ we have $x\mathrel{R}\apply{g}{x}$
            by \cref{choice_function}.
        Let $A = \{U\in\pow{X} \mid \exists x\in X. U = \metricBall{d_X}{X}{x}{\apply{g}{x}}\}$.

        Show that $A$ is an open covering of $X$ with respect to $T_X$.
        \begin{subproof}
            Show that $A$ is a covering of $X$.
            \begin{subproof}
                Show that for all $x\in X$ we have $x\in\unions{A}$.
                \begin{subproof}
                    Fix $x\in X$.
                    $(x,\apply{g}{x})\in g$ by \cref{function_apply_elim}.
                    Thus $x\mathrel{R}\apply{g}{x}$.
                    Thus $0 < \apply{g}{x}$.
                    $d_X((x,x)) = 0$ by \cref{metric_zero_characterization,metric_on}.
                    Thus $d_X((x,x)) < \apply{g}{x}$ by \cref{real_lt_subst_left}.
                    Thus $x\in \metricBall{d_X}{X}{x}{\apply{g}{x}}$ by \cref{metric_ball}.
                    Show that $\metricBall{d_X}{X}{x}{\apply{g}{x}}\subseteq X$.
                    \begin{subproof}
                        Show that for all $y$ such that $y\in\metricBall{d_X}{X}{x}{\apply{g}{x}}$ we have $y\in X$.
                        \begin{subproof}
                            Fix $y$.
                            Assume $y\in\metricBall{d_X}{X}{x}{\apply{g}{x}}$.
                            Thus $y\in X$ by \cref{metric_ball}.
                        \end{subproof}
                        Thus $\metricBall{d_X}{X}{x}{\apply{g}{x}}\subseteq X$ by \cref{subseteq}.
                    \end{subproof}
                    Thus $\metricBall{d_X}{X}{x}{\apply{g}{x}}\in\pow{X}$ by \cref{powerset_intro}.
                    Thus $\metricBall{d_X}{X}{x}{\apply{g}{x}}\in A$.
                    Thus $x\in\unions{A}$ by \cref{unions_iff}.
                \end{subproof}
                Thus $X\subseteq\unions{A}$ by \cref{subseteq}.
            \end{subproof}
            Show that for all $U$ such that $U\in A$ we have $U$ is open in $X$ with respect to $T_X$.
            \begin{subproof}
                Fix $U$.
                Assume $U\in A$.
                Take $x\in X$ such that $U = \metricBall{d_X}{X}{x}{\apply{g}{x}}$.
                $(x,\apply{g}{x})\in g$ by \cref{function_apply_elim}.
                Thus $x\mathrel{R}\apply{g}{x}$.
                Thus $(x,\apply{g}{x})\in R$.
                Thus $(x,\apply{g}{x})\in X\times\reals$.
                Thus $\apply{g}{x}\in\reals$ by \cref{times_tuple_elim}.
                Thus $0 < \apply{g}{x}$.
                Show that $\metricBall{d_X}{X}{x}{\apply{g}{x}}\subseteq X$.
                \begin{subproof}
                    Show that for all $y$ such that $y\in\metricBall{d_X}{X}{x}{\apply{g}{x}}$ we have $y\in X$.
                    \begin{subproof}
                        Fix $y$.
                        Assume $y\in\metricBall{d_X}{X}{x}{\apply{g}{x}}$.
                        Thus $y\in X$ by \cref{metric_ball}.
                    \end{subproof}
                    Thus $\metricBall{d_X}{X}{x}{\apply{g}{x}}\subseteq X$ by \cref{subseteq}.
                \end{subproof}
                Thus $\metricBall{d_X}{X}{x}{\apply{g}{x}}\in\pow{X}$ by \cref{powerset_intro}.
                Thus $U\in\metricBasis{d_X}{X}$ by \cref{metric_basis}.
                $\metricBasis{d_X}{X}$ is a basis on $X$ by \cref{metric_basis_is_basis}.
                Thus $U\in T_X$ by \cref{basis_elem_open_in_generated,basis_for_topology,metric_topology}.
                Thus $T_X$ is a topology on $X$ by \cref{basis_topology_is_topology,metric_topology}.
                Thus $U$ is open in $X$ with respect to $T_X$ by \cref{open_in}.
            \end{subproof}
            Thus $A$ is an open covering of $X$ with respect to $T_X$ by \cref{open_covering}.
        \end{subproof}

        Take $\delta\in\reals$ such that $0 < \delta$ and
            for all $B$ such that $B\subseteq X$ and
                for all $x,y\in B$ we have $d_X((x,y)) < \delta$
            there exists $U\in A$ such that $B\subseteq U$
            by \cref{lebesgue_number_lemma}.

        Show that for all $x_0,x_1$ such that $x_0\in X$ and $x_1\in X$ and $d_X((x_0,x_1)) < \delta$
            we have $d_Y((\apply{f}{x_0},\apply{f}{x_1})) < \epsilon$.
        \begin{subproof}
            Fix $x_0$.
            Fix $x_1$.
            Assume $x_0\in X$ and $x_1\in X$ and $d_X((x_0,x_1)) < \delta$.
            Let $B = \{x_0,x_1\}$.
            Show that $B\subseteq X$.
            \begin{subproof}
                Show that for all $z$ such that $z\in B$ we have $z\in X$.
                \begin{subproof}
                    Fix $z$.
                    Assume $z\in B$.
                    Thus $z = x_0$ or $z = x_1$ by \cref{upair_iff}.
                    Thus $z\in X$.
                \end{subproof}
                Thus $B\subseteq X$ by \cref{subseteq}.
            \end{subproof}
            Show that for all $x\in B$ we have for all $y\in B$ we have $d_X((x,y)) < \delta$.
            \begin{subproof}
                Fix $x$.
                Assume $x\in B$.
                Fix $y$.
                Assume $y\in B$.
                Thus $x = x_0$ or $x = x_1$ by \cref{upair_iff}.
                Thus $y = x_0$ or $y = x_1$ by \cref{upair_iff}.
                \begin{byCase}
                \caseOf{$x = x_0$.}
                    \begin{byCase}
                    \caseOf{$y = x_0$.}
                        $d_X((x_0,x_0)) = 0$ by \cref{metric_zero_characterization,metric_on}.
                        Thus $d_X((x,y)) < \delta$ by \cref{real_lt_subst_left}.
                    \caseOf{$y \neq x_0$.}
                        Thus $y = x_1$.
                        Thus $d_X((x,y)) < \delta$.
                    \end{byCase}
                \caseOf{$x \neq x_0$.}
                    Thus $x = x_1$.
                    \begin{byCase}
                    \caseOf{$y = x_0$.}
                        $d_X((x_1,x_0)) = d_X((x_0,x_1))$ by \cref{metric_symmetric,metric_on}.
                        Thus $d_X((x,y)) < \delta$.
                    \caseOf{$y \neq x_0$.}
                        Thus $y = x_1$.
                        $d_X((x_1,x_1)) = 0$ by \cref{metric_zero_characterization,metric_on}.
                        Thus $d_X((x,y)) < \delta$ by \cref{real_lt_subst_left}.
                    \end{byCase}
                \end{byCase}
            \end{subproof}
            Take $U\in A$ such that $B\subseteq U$.
            $x_0\in B$ by \cref{upair_intro_left}.
            $x_1\in B$ by \cref{upair_intro_right}.
            Thus $x_0\in U$ by \cref{subseteq}.
            Thus $x_1\in U$ by \cref{subseteq}.
            Take $x\in X$ such that $U = \metricBall{d_X}{X}{x}{\apply{g}{x}}$.
            Thus $d_X((x,x_0)) < \apply{g}{x}$ by \cref{metric_ball}.
            Thus $d_X((x,x_1)) < \apply{g}{x}$ by \cref{metric_ball}.
            $(x,\apply{g}{x})\in g$ by \cref{function_apply_elim}.
            Thus $x\mathrel{R}\apply{g}{x}$.
            Thus $(x,\apply{g}{x})\in R$.
            Thus $(x,\apply{g}{x})\in X\times\reals$.
            Show that for all $y\in X$ if $d_X((x,y)) < \apply{g}{x}$ then
                $d_Y((\apply{f}{x},\apply{f}{y})) < \epsilon_0$.
            \begin{subproof}
                Fix $y$.
                Assume $y\in X$.
                Assume $d_X((x,y)) < \apply{g}{x}$.
                Let $w = (x,\apply{g}{x})$.
                Thus $w\in R$.
                Thus $w\in X\times\reals$ and
                    for all $y_0\in X$ if $d_X((\fst{w},y_0)) < \snd{w}$ then
                        $d_Y((\apply{f}{\fst{w}},\apply{f}{y_0})) < \epsilon_0$.
                $\fst{w} = x$ and $\snd{w} = \apply{g}{x}$ by \cref{fst_eq,snd_eq}.
                Thus for all $y_0\in X$ if $d_X((x,y_0)) < \apply{g}{x}$ then
                    $d_Y((\apply{f}{x},\apply{f}{y_0})) < \epsilon_0$.
                Thus $d_Y((\apply{f}{x},\apply{f}{y})) < \epsilon_0$.
            \end{subproof}
            Thus $x_0\in X$.
            Thus $d_Y((\apply{f}{x},\apply{f}{x_0})) < \epsilon_0$.
            Thus $x_1\in X$.
            Thus $d_Y((\apply{f}{x},\apply{f}{x_1})) < \epsilon_0$.
            $\apply{f}{x_0}\in Y$ by \cref{funs_type_apply}.
            $\apply{f}{x}\in Y$ by \cref{funs_type_apply}.
            $\apply{f}{x_1}\in Y$ by \cref{funs_type_apply}.
            $d_Y((\apply{f}{x_0},\apply{f}{x})) = d_Y((\apply{f}{x},\apply{f}{x_0}))$ by \cref{metric_symmetric,metric_on}.
            Thus $d_Y((\apply{f}{x_0},\apply{f}{x})) < \epsilon_0$ by \cref{real_lt_subst_left}.
            $d_Y$ satisfies the triangle inequality on $Y$ by \cref{metric_on}.
            Let $s = d_Y((\apply{f}{x_0},\apply{f}{x})) + d_Y((\apply{f}{x},\apply{f}{x_1}))$.
            Thus $s \geq d_Y((\apply{f}{x_0},\apply{f}{x_1}))$ by \cref{metric_triangle_inequality}.
            Thus $d_Y((\apply{f}{x_0},\apply{f}{x_1})) \leq s$.
            $d_Y\in\funs{Y\times Y}{\reals}$ by \cref{metric_on}.
            $(\apply{f}{x_0},\apply{f}{x})\in Y\times Y$ by \cref{times_tuple_intro}.
            $(\apply{f}{x},\apply{f}{x_1})\in Y\times Y$ by \cref{times_tuple_intro}.
            Thus $d_Y((\apply{f}{x_0},\apply{f}{x}))\in\reals$ by \cref{funs_type_apply}.
            Thus $d_Y((\apply{f}{x},\apply{f}{x_1}))\in\reals$ by \cref{funs_type_apply}.
            Thus $\epsilon_0\in\reals$.
            Thus $s < \epsilon_0 + \epsilon_0$ by \cref{real_lt_add_two}.
            Show that $d_Y((\apply{f}{x_0},\apply{f}{x_1})) < \epsilon_0 + \epsilon_0$.
            \begin{subproof}
                $(\apply{f}{x_0},\apply{f}{x_1})\in Y\times Y$ by \cref{times_tuple_intro}.
                Thus $d_Y((\apply{f}{x_0},\apply{f}{x_1}))\in\reals$ by \cref{funs_type_apply}.
                Thus $s\in\reals$ by \cref{real_add_closed}.
                Thus $\epsilon_0 + \epsilon_0\in\reals$ by \cref{real_add_closed}.
                Thus $d_Y((\apply{f}{x_0},\apply{f}{x_1})) < s$ or
                    $d_Y((\apply{f}{x_0},\apply{f}{x_1})) = s$ by \cref{real_leq_split}.
                \begin{byCase}
                \caseOf{$d_Y((\apply{f}{x_0},\apply{f}{x_1})) < s$.}
                    Thus $d_Y((\apply{f}{x_0},\apply{f}{x_1})) < \epsilon_0 + \epsilon_0$ by \cref{real_lt_trans}.
                \caseOf{$d_Y((\apply{f}{x_0},\apply{f}{x_1})) = s$.}
                    Thus $d_Y((\apply{f}{x_0},\apply{f}{x_1})) < \epsilon_0 + \epsilon_0$ by \cref{real_lt_subst_left}.
                \end{byCase}
            \end{subproof}

            Show that $\epsilon_0 + \epsilon_0 = \epsilon$.
            \begin{subproof}
                $\epsilon_0 = \epsilon \rmul \inv{(1 + 1)}$ by \cref{real_div}.
                Thus $\epsilon_0 + \epsilon_0 =
                    (\epsilon \rmul \inv{(1 + 1)}) + (\epsilon \rmul \inv{(1 + 1)})$.
                $(\epsilon + \epsilon) \rmul \inv{(1 + 1)}
                    = (\epsilon \rmul \inv{(1 + 1)}) + (\epsilon \rmul \inv{(1 + 1)})$
                    by \cref{real_right_distrib}.
                Thus $\epsilon_0 + \epsilon_0 = (\epsilon + \epsilon) \rmul \inv{(1 + 1)}$.
                $\epsilon \rmul (1 + 1) = (\epsilon \rmul 1) + (\epsilon \rmul 1)$ by \cref{real_distrib}.
                $1 \rmul \epsilon = \epsilon$ by \cref{real_mul_ident}.
                $\epsilon \rmul 1 = 1 \rmul \epsilon$ by \cref{real_mul_comm}.
                Thus $\epsilon \rmul 1 = \epsilon$ by \cref{real_mul_ident,real_mul_comm}.
                Thus $\epsilon \rmul (1 + 1) = \epsilon + \epsilon$.
                Thus $(\epsilon + \epsilon) \rmul \inv{(1 + 1)} = (\epsilon \rmul (1 + 1)) \rmul \inv{(1 + 1)}$.
                $(\epsilon \rmul (1 + 1)) \rmul \inv{(1 + 1)}
                    = \epsilon \rmul ((1 + 1) \rmul \inv{(1 + 1)})$ by \cref{real_mul_assoc}.
                $(1 + 1) \rmul \inv{(1 + 1)} = 1$ by \cref{real_mul_inverse}.
                Thus $\epsilon_0 + \epsilon_0 = \epsilon \rmul 1$.
                Thus $\epsilon_0 + \epsilon_0 = \epsilon$ by \cref{real_mul_ident}.
            \end{subproof}
            Thus $d_Y((\apply{f}{x_0},\apply{f}{x_1})) < \epsilon$ by \cref{real_lt_subst_right}.
        \end{subproof}
    \end{subproof}
    Thus $f$ is uniformly continuous from $X$ to $Y$ with respect to $d_X$ and $d_Y$ by \cref{uniformly_continuous}.
\end{proof}

% from §27 Item 10: isolated point
\begin{definition}\label{isolated_point}
    $x$ is an isolated point of $X$ with respect to $T$ iff
        $x\in X$ and $\{x\}$ is open in $X$ with respect to $T$.
\end{definition}

% from §27 Item 11: compact Hausdorff without isolated points is uncountable
\begin{theorem}\label{compact_hausdorff_no_isolated_uncountable}
    Assume $T$ is a topology on $X$.
    Assume $X$ is compact with respect to $T$.
    Assume $X$ is Hausdorff with respect to $T$.
    Assume $X\neq\emptyset$.
    Suppose for all $x\in X$ we have $x$ is not an isolated point of $X$ with respect to $T$.
    Then $X$ is not countable.
\end{theorem}
\begin{proof}
    Omitted. % do not touch
\end{proof}

% from §27 Item 12: closed interval in reals is uncountable
\begin{corollary}\label{real_closed_interval_uncountable}
    Suppose $a\in\reals$.
    Suppose $b\in\reals$.
    Suppose $a < b$.
    Let $R = \realOrderRel$.
    Let $T = \realOrderTopology$.
    Let $Y = \intervalclosedrel{R}{\reals}{a}{b}$.
    Let $T_Y = \subspaceTopology{T}{Y}$.
    Then $Y$ is not countable.
\end{corollary}
\begin{proof}
    $R$ is a simple order relation on $\reals$ by \cref{real_order_relation_simple}.
    $\reals$ is a linear continuum with respect to $\realOrderRel$ by \cref{reals_linear_continuum}.
    Thus $\reals$ has more than one element by \cref{linear_continuum}.
    $\realOrderTopology = \basisTopology{\orderBasis{\realOrderRel}{\reals}}{\reals}$ by \cref{real_order_topology}.
    Thus $T$ is the order topology on $\reals$ with respect to $R$ by \cref{order_topology}.
    $\orderBasis{R}{\reals}$ is a basis on $\reals$ by \cref{order_basis_is_basis}.
    Thus $T$ is a topology on $\reals$ by \cref{basis_topology_is_topology,order_topology}.

    $a\mathrel{R}b$ by \cref{real_order_relation_intro}.
    Show that $a\neq b$.
    \begin{subproof}
        Suppose not.
        Thus $a = b$.
        Thus $a < a$ by \cref{real_lt_subst_right}.
        Contradiction by \cref{real_lt_irreflexive}.
    \end{subproof}
    Thus $a\in Y$ by \cref{intervalclosed_rel}.
    Thus $b\in Y$ by \cref{intervalclosed_rel}.
    Thus $Y$ has more than one element by \cref{more_than_one_element}.
    Thus $Y\neq\emptyset$ by \cref{emptyset}.

    Show that $Y\subseteq\reals$.
    \begin{subproof}
        Show that for all $y$ such that $y\in Y$ we have $y\in\reals$.
        \begin{subproof}
            Fix $y$.
            Assume $y\in Y$.
            Thus $y\in\reals$ by \cref{intervalclosed_rel}.
        \end{subproof}
        Thus $Y\subseteq\reals$ by \cref{subseteq}.
    \end{subproof}

    Thus $\reals$ is Hausdorff with respect to $T$ by \cref{order_topology_hausdorff}.
    Thus $Y$ is Hausdorff with respect to $T_Y$ by \cref{subspace_hausdorff}.
    $T_Y$ is a topology on $Y$ by \cref{subspace_topology_is_topology}.

    Thus $a < b$ or $a = b$.
    Thus $Y$ is compact with respect to $T_Y$ by \cref{real_closed_interval_compact}.
    Thus $Y$ is connected with respect to $T_Y$ by \cref{linear_continuum_connected}.

    Show that for all $x\in Y$ we have $x$ is not an isolated point of $Y$ with respect to $T_Y$.
    \begin{subproof}
        Fix $x\in Y$.
        Suppose not.
        Thus $x$ is an isolated point of $Y$ with respect to $T_Y$.
        Thus $\{x\}$ is open in $Y$ with respect to $T_Y$ by \cref{isolated_point}.
        Take $y\in Y$ such that $y\neq x$ by \cref{more_than_one_element_has_other}.
        Show that $\{x\}\subseteq Y$.
        \begin{subproof}
            Show that for all $u$ such that $u\in\{x\}$ we have $u\in Y$.
            \begin{subproof}
                Fix $u$.
                Assume $u\in\{x\}$.
                Thus $u = x$ by \cref{singleton_iff}.
                Thus $u\in Y$.
            \end{subproof}
            Thus $\{x\}\subseteq Y$ by \cref{subseteq}.
        \end{subproof}
        $\{x,y\}$ is finite by \cref{pair_is_finite}.
        Show that $\{x\}\subseteq\{x,y\}$.
        \begin{subproof}
            Show that for all $u$ such that $u\in\{x\}$ we have $u\in\{x,y\}$.
            \begin{subproof}
                Fix $u$.
                Assume $u\in\{x\}$.
                Thus $u = x$ by \cref{singleton_iff}.
                Thus $u\in\{x,y\}$ by \cref{upair_intro_left}.
            \end{subproof}
            Thus $\{x\}\subseteq\{x,y\}$ by \cref{subseteq}.
        \end{subproof}
        Thus $\{x\}$ is finite by \cref{subseteq_finite_is_finite}.
        Thus $\{x\}$ is closed in $Y$ with respect to $T_Y$ by \cref{finite_point_set_closed_hausdorff}.
        $T_Y$ is a topology on $Y$ by \cref{subspace_topology_is_topology}.
        Thus $\{x\} = \emptyset$ or $\{x\} = Y$ by \cref{connected_iff_clopen}.
        $x\in\{x\}$ by \cref{singleton_iff}.
        Thus $\{x\}\neq\emptyset$ by \cref{emptyset}.
        Thus $\{x\} = Y$.
        Take $z\in Y$ such that $z\neq x$ by \cref{more_than_one_element_has_other}.
        Thus $Y\subseteq\{x\}$ by \cref{subseteq}.
        Thus $z\in\{x\}$ by \cref{subseteq}.
        Thus $z = x$ by \cref{singleton_iff}.
        Contradiction.
    \end{subproof}

    Thus $Y$ is not countable by \cref{compact_hausdorff_no_isolated_uncountable}.
\end{proof}


% from §28 Item 1: limit point compact
\begin{definition}\label{limit_point_compact}
    $X$ is limit point compact with respect to $T$ iff
        $T$ is a topology on $X$ and
        for all $A$ such that $A\subseteq X$ and $A$ is infinite
            there exists $x\in X$ such that $x$ is a limit point of $A$ in $X$ with respect to $T$.
\end{definition}

% from §28 Item 2: compact implies limit point compact
\begin{theorem}\label{compact_implies_limit_point_compact}
    Assume $T$ is a topology on $X$.
    Suppose $X$ is compact with respect to $T$.
    Then $X$ is limit point compact with respect to $T$.
\end{theorem}
\begin{proof}
    Show that for all $A$ such that $A\subseteq X$ and $A$ is infinite
        there exists $x\in X$ such that $x$ is a limit point of $A$ in $X$ with respect to $T$.
    \begin{subproof}
        Fix $A$.
        Assume $A\subseteq X$ and $A$ is infinite.
        Suppose not.
        Then for all $x\in X$ we have $x$ is not a limit point of $A$ in $X$ with respect to $T$.

        Let $C = \{U\in\pow{X} \mid \text{there exists $a\in A$ such that
            $U$ is a neighborhood of $a$ in $X$ with respect to $T$ and
            $U\inter (A\setminus\{a\}) = \emptyset$}\}$.

        Show that for all $U$ such that $U\in C$ we have $U$ is open in $X$ with respect to $T$.
        \begin{subproof}
            Fix $U$.
            Assume $U\in C$.
            Take $a\in A$ such that
                $U$ is a neighborhood of $a$ in $X$ with respect to $T$ and
                $U\inter (A\setminus\{a\}) = \emptyset$.
            Thus $U$ is open in $X$ with respect to $T$ by \cref{neighborhood}.
        \end{subproof}

        Show that $A\subseteq\unions{C}$.
        \begin{subproof}
            Show that for all $a$ such that $a\in A$ we have $a\in\unions{C}$.
            \begin{subproof}
                Fix $a$.
                Assume $a\in A$.
                Thus $a$ is not a limit point of $A$ in $X$ with respect to $T$.
                Thus there exists $U$ such that
                    $U$ is a neighborhood of $a$ in $X$ with respect to $T$ and
                    $U$ does not intersect $A\setminus\{a\}$ by \cref{limit_point}.
                Take $U$ such that
                    $U$ is a neighborhood of $a$ in $X$ with respect to $T$ and
                    $U$ does not intersect $A\setminus\{a\}$.
                Thus $U\inter (A\setminus\{a\}) = \emptyset$ by \cref{intersects}.
                Thus $U\in C$.
                Thus $a\in U$ by \cref{neighborhood}.
                Thus $a\in\unions{C}$ by \cref{unions_iff}.
            \end{subproof}
            Thus $A\subseteq\unions{C}$ by \cref{subseteq}.
        \end{subproof}

        Show that $\limitPoints{A}{X}{T}\subseteq\emptyset$.
        \begin{subproof}
            Show that for all $x$ such that $x\in\limitPoints{A}{X}{T}$ we have $x\in\emptyset$.
            \begin{subproof}
                Fix $x$.
                Assume $x\in\limitPoints{A}{X}{T}$.
                Thus $x$ is a limit point of $A$ in $X$ with respect to $T$ by \cref{limit_points}.
                Contradiction.
            \end{subproof}
            Thus $\limitPoints{A}{X}{T}\subseteq\emptyset$ by \cref{subseteq}.
        \end{subproof}
        $\emptyset\subseteq A$ by \cref{emptyset_subseteq}.
        Thus $\limitPoints{A}{X}{T}\subseteq A$ by \cref{subseteq_transitive}.
        Thus $A$ is closed in $X$ with respect to $T$ by \cref{closed_iff_contains_limit_points}.
        Thus $X\setminus A$ is open in $X$ with respect to $T$ by \cref{closed_in}.

        Let $D = C\union\{X\setminus A\}$.

        Show that $D$ is an open covering of $X$ with respect to $T$.
        \begin{subproof}
            Show that for all $U$ such that $U\in D$ we have $U$ is open in $X$ with respect to $T$.
            \begin{subproof}
                Fix $U$.
                Assume $U\in D$.
                \begin{byCase}
                \caseOf{$U\in C$.}
                    Thus $U$ is open in $X$ with respect to $T$.
                \caseOf{$U\notin C$.}
                    Thus $U\in\{X\setminus A\}$ by \cref{union_iff}.
                    Thus $U = X\setminus A$ by \cref{singleton_iff}.
                    Thus $U$ is open in $X$ with respect to $T$.
                \end{byCase}
            \end{subproof}
            Show that $D$ is a covering of $X$.
            \begin{subproof}
                Show that for all $x$ such that $x\in X$ we have $x\in\unions{D}$.
                \begin{subproof}
                    Fix $x$.
                    Assume $x\in X$.
                    \begin{byCase}
                    \caseOf{$x\in A$.}
                        Thus $x\in\unions{C}$ by \cref{subseteq}.
                        Thus $x\in\unions{D}$ by \cref{union_iff,unions_iff}.
                    \caseOf{$x\notin A$.}
                        Thus $x\in X\setminus A$ by \cref{setminus_intro}.
                        Show that $X\setminus A\in D$.
                        \begin{subproof}
                            $X\setminus A\in\{X\setminus A\}$ by \cref{singleton_intro}.
                            Thus $X\setminus A\in D$ by \cref{union_iff}.
                        \end{subproof}
                        Thus $x\in\unions{D}$ by \cref{unions_iff}.
                    \end{byCase}
                \end{subproof}
                Thus $X\subseteq\unions{D}$ by \cref{subseteq}.
                Thus $D$ is a covering of $X$ by \cref{covering}.
            \end{subproof}
            Thus $D$ is an open covering of $X$ with respect to $T$ by \cref{open_covering}.
        \end{subproof}

        Take $B$ such that $B\subseteq D$ and $B$ is finite and $B$ is a covering of $X$ by \cref{compact}.
        Let $B_A = \{U\in B \mid U\neq X\setminus A\}$.
        $B_A\subseteq B$ by \cref{subseteq}.
        Thus $B_A$ is finite by \cref{subseteq_finite_is_finite}.

        Show that $A\subseteq\unions{B_A}$.
        \begin{subproof}
            Show that for all $a$ such that $a\in A$ we have $a\in\unions{B_A}$.
            \begin{subproof}
                Fix $a$.
                Assume $a\in A$.
                Thus $a\in X$ by \cref{subseteq}.
                $X\subseteq\unions{B}$ by \cref{covering}.
                Thus $a\in\unions{B}$ by \cref{subseteq}.
                Take $U\in B$ such that $a\in U$ by \cref{unions_iff}.
                Show that $U\neq X\setminus A$.
                \begin{subproof}
                    Suppose not.
                    Then $U = X\setminus A$.
                    Thus $a\in X\setminus A$ by \cref{subseteq}.
                    Thus $a\notin A$ by \cref{setminus_elim_right}.
                    Contradiction.
                \end{subproof}
                Thus $U\in B_A$.
                Thus $a\in\unions{B_A}$ by \cref{unions_iff}.
            \end{subproof}
            Thus $A\subseteq\unions{B_A}$ by \cref{subseteq}.
        \end{subproof}

        Let $h = \{(U, U\inter A) \mid U\in B_A\}$.
        Show that $h$ is a function on $B_A$.
        \begin{subproof}
            Show that $h$ is a function.
            \begin{subproof}
                Show that for all $U,V_1,V_2$ such that $(U,V_1)\in h$ and $(U,V_2)\in h$ we have $V_1 = V_2$.
                \begin{subproof}
                    Fix $U$.
                    Fix $V_1$.
                    Fix $V_2$.
                    Assume $(U,V_1)\in h$ and $(U,V_2)\in h$.
                    Take $U_1\in B_A$ such that $V_1 = U_1\inter A$ and $(U,V_1) = (U_1,U_1\inter A)$ by \cref{pair_eq_iff}.
                    Take $U_2\in B_A$ such that $V_2 = U_2\inter A$ and $(U,V_2) = (U_2,U_2\inter A)$ by \cref{pair_eq_iff}.
                    Thus $U = U_1$ and $U = U_2$ by \cref{pair_eq_iff}.
                    Thus $V_1 = U\inter A$ and $V_2 = U\inter A$.
                    Thus $V_1 = V_2$.
                \end{subproof}
                Thus $h$ is right-unique by \cref{rightunique}.
                Thus $h$ is a function.
            \end{subproof}
            Show that $\dom{h} = B_A$.
            \begin{subproof}
                Show that $B_A\subseteq\dom{h}$.
                \begin{subproof}
                    Show that for all $U$ such that $U\in B_A$ we have $U\in\dom{h}$.
                    \begin{subproof}
                        Fix $U$.
                        Assume $U\in B_A$.
                        Let $V = U\inter A$.
                        $(U,V)\in h$ by \cref{pair_eq_iff}.
                        Thus $U\in\dom{h}$ by \cref{dom_intro}.
                    \end{subproof}
                    Thus $B_A\subseteq\dom{h}$ by \cref{subseteq}.
                \end{subproof}
                Show that $\dom{h}\subseteq B_A$.
                \begin{subproof}
                    Show that for all $U$ such that $U\in\dom{h}$ we have $U\in B_A$.
                    \begin{subproof}
                        Fix $U$.
                        Assume $U\in\dom{h}$.
                        Take $V$ such that $(U,V)\in h$ by \cref{dom_iff}.
                        Take $U_0\in B_A$ such that $(U,V) = (U_0,U_0\inter A)$ by \cref{pair_eq_iff}.
                        Thus $U = U_0$ by \cref{pair_eq_iff}.
                        Thus $U\in B_A$.
                    \end{subproof}
                    Thus $\dom{h}\subseteq B_A$ by \cref{subseteq}.
                \end{subproof}
                Thus $\dom{h} = B_A$ by \cref{subseteq_antisymmetric}.
            \end{subproof}
            Thus $h$ is a function on $B_A$.
        \end{subproof}

        Let $F = \img{h}{B_A}$.
        Thus $F$ is finite by \cref{finite_image_function}.

        Show that for all $V$ such that $V\in F$ we have $V$ is finite.
        \begin{subproof}
            Fix $V$.
            Assume $V\in F$.
            Take $U\in B_A$ such that $(U,V)\in h$ by \cref{img_iff}.
            Thus $V = \apply{h}{U}$ by \cref{function_apply_intro}.
            Take $U_0\in B_A$ such that $(U,V) = (U_0,U_0\inter A)$ by \cref{pair_eq_iff}.
            Thus $V = U\inter A$ by \cref{pair_eq_iff}.
            $U\in B_A$.
            Thus $U\in B$ by \cref{subseteq}.
            Thus $U\neq X\setminus A$.
            $U\in D$ by \cref{subseteq}.
            Show that $U\in C$.
            \begin{subproof}
                Suppose not.
                Then $U\in\{X\setminus A\}$ by \cref{union_iff}.
                Thus $U = X\setminus A$ by \cref{singleton_iff}.
                Contradiction.
            \end{subproof}
            Take $a\in A$ such that
                $U$ is a neighborhood of $a$ in $X$ with respect to $T$ and
                $U\inter (A\setminus\{a\}) = \emptyset$.
            Show that $V\subseteq\{a\}$.
            \begin{subproof}
                Show that for all $x$ such that $x\in V$ we have $x\in\{a\}$.
                \begin{subproof}
                    Fix $x$.
                    Assume $x\in V$.
                    Thus $x\in U$ and $x\in A$ by \cref{inter}.
                    \begin{byCase}
                    \caseOf{$x = a$.}
                        Thus $x\in\{a\}$ by \cref{singleton_intro}.
                    \caseOf{$x\neq a$.}
                        Show that $x\notin\{a\}$.
                        \begin{subproof}
                            Suppose not.
                            Thus $x\in\{a\}$.
                            Thus $x = a$ by \cref{singleton_iff}.
                            Contradiction.
                        \end{subproof}
                        Thus $x\in A\setminus\{a\}$ by \cref{setminus_intro}.
                        Thus $x\in U\inter (A\setminus\{a\})$ by \cref{inter_intro}.
                        Contradiction by \cref{emptyset}.
                    \end{byCase}
                \end{subproof}
                Thus $V\subseteq\{a\}$ by \cref{subseteq}.
            \end{subproof}
            Show that $\{a\}$ is finite.
                \begin{subproof}
                    $\{a,a\}$ is finite by \cref{pair_is_finite}.
                    $a\in\{a,a\}$ by \cref{upair_intro_left}.
                    Thus $\{a\}\subseteq\{a,a\}$ by \cref{singleton_subset_intro}.
                    Thus $\{a\}$ is finite by \cref{subseteq_finite_is_finite}.
                \end{subproof}
            Thus $V$ is finite by \cref{subseteq_finite_is_finite}.
        \end{subproof}

        Show that $\unions{F} = A$.
        \begin{subproof}
            Show that $A\subseteq\unions{F}$.
            \begin{subproof}
                Show that for all $a$ such that $a\in A$ we have $a\in\unions{F}$.
                \begin{subproof}
                    Fix $a$.
                    Assume $a\in A$.
                    Thus $a\in\unions{B_A}$ by \cref{subseteq}.
                    Take $U\in B_A$ such that $a\in U$ by \cref{unions_iff}.
                    Let $V = U\inter A$.
                    $(U,V)\in h$ by \cref{pair_eq_iff}.
                    Thus $V\in F$ by \cref{img_iff}.
                    Thus $a\in V$ by \cref{inter_intro}.
                    Thus $a\in\unions{F}$ by \cref{unions_iff}.
                \end{subproof}
                Thus $A\subseteq\unions{F}$ by \cref{subseteq}.
            \end{subproof}
            Show that $\unions{F}\subseteq A$.
            \begin{subproof}
                Show that for all $x$ such that $x\in\unions{F}$ we have $x\in A$.
                \begin{subproof}
                    Fix $x$.
                    Assume $x\in\unions{F}$.
                    Take $V\in F$ such that $x\in V$ by \cref{unions_iff}.
                    Take $U\in B_A$ such that $(U,V)\in h$ by \cref{img_iff}.
                    Take $U_0\in B_A$ such that $(U,V) = (U_0,U_0\inter A)$ by \cref{pair_eq_iff}.
                    Thus $V = U_0\inter A$ by \cref{pair_eq_iff}.
                    Thus $x\in A$ by \cref{inter}.
                \end{subproof}
                Thus $\unions{F}\subseteq A$ by \cref{subseteq}.
            \end{subproof}
            Thus $\unions{F} = A$ by \cref{subseteq_antisymmetric}.
        \end{subproof}

        Thus $\unions{F}$ is finite by \cref{finite_union_of_finite_family}.
        Thus $A$ is finite by \cref{subseteq_antisymmetric}.
        Contradiction.
    \end{subproof}
    Thus $X$ is limit point compact with respect to $T$ by \cref{limit_point_compact}.
\end{proof}

% from §28 helper lemma 1: naturalsPlus is infinite
\begin{lemma}\label{naturalsplus_infinite}
    $\naturalsPlus$ is infinite.
\end{lemma}
\begin{proof}
    Suppose not.
    Then $\naturalsPlus$ is finite.
    $1\in\naturalsPlus$ by \cref{real_one_in_naturals}.
    Thus $\naturalsPlus\neq\emptyset$ by \cref{emptyset}.
    $\naturalsPlus\subseteq\naturalsPlus$ by \cref{subseteq_refl}.
    Take $m\in\naturalsPlus$ such that for all $a\in\naturalsPlus$ we have $a \leq m$
        by \cref{finite_naturals_maximum}.
    $m + 1\in\naturalsPlus$ by \cref{real_naturals_closed_succ}.
    Thus $m + 1 \leq m$.
    $\naturalsPlus\subseteq\reals$ by \cref{real_naturals_subset_reals}.
    Thus $m\in\reals$ by \cref{subseteq}.
    Thus $m + 1\in\reals$ by \cref{real_naturals_closed_succ,subseteq}.
    Thus $1\in\reals$ by \cref{real_one_in_naturals,subseteq}.
    $0\in\reals$ by \cref{real_add_ident}.
    $0 < 1$ by \cref{real_zero_lt_one}.
    Thus $0 + m < 1 + m$ by \cref{real_lt_add}.
    $0 + m = m$ by \cref{real_add_ident}.
    $1 + m = m + 1$ by \cref{real_add_comm}.
    Thus $m < m + 1$.
    Thus $m \leq m + 1$.
    Thus $m = m + 1$ by \cref{real_leq_antisym}.
    Thus $1 + m = m$ by \cref{real_add_comm}.
    $0 + m = m$ by \cref{real_add_ident}.
    Thus $1 + m = 0 + m$.
    Thus $1 = 0$ by \cref{real_add_cancel}.
    Thus $0 < 0$ by \cref{real_zero_lt_one}.
    Contradiction by \cref{real_lt_irreflexive}.
\end{proof}

% from §28 Item 3: limit point compact does not imply compact
\begin{theorem}\label{limit_point_compact_not_compact}
    There exists $X,T$ such that
        $X$ is limit point compact with respect to $T$ and
        $X$ is not compact with respect to $T$.
\end{theorem}
\begin{proof}
    Let $Y = \{0,1\}$.
    Let $X = \naturalsPlus\times Y$.
    Let $T_0 = \discrete{\naturalsPlus}$.
    Let $T_1 = \{\emptyset, Y\}$.
    Let $T = \productTopology{T_0}{T_1}{\naturalsPlus}{Y}$.

    Show that $T_0$ is a topology on $\naturalsPlus$.
    \begin{subproof}
        $T_0 = \pow{\naturalsPlus}$ by \cref{discrete_topology}.
        Thus $T_0\subseteq\pow{\naturalsPlus}$ by \cref{subseteq_refl}.
        $\emptyset\in T_0$ by \cref{discrete_topology,powerset_bottom}.
        $\naturalsPlus\in T_0$ by \cref{discrete_topology,powerset_top}.
        Show that $T_0$ is closed under arbitrary unions.
        \begin{subproof}
            Fix $M$.
            Assume $M\subseteq T_0$.
            $T_0 = \pow{\naturalsPlus}$ by \cref{discrete_topology}.
            Thus $T_0\subseteq\pow{\naturalsPlus}$ by \cref{subseteq_refl}.
            Thus $M\subseteq\pow{\naturalsPlus}$ by \cref{subseteq_transitive}.
            Thus $\unions{M}\in\pow{\naturalsPlus}$ by \cref{powerset_closed_unions}.
            Thus $\unions{M}\in T_0$ by \cref{discrete_topology}.
        \end{subproof}
        Show that $T_0$ is closed under binary intersections.
        \begin{subproof}
            Show that for all $A$ such that $A\in T_0$ we have
                for all $B$ such that $B\in T_0$ we have $A\inter B\in T_0$.
            \begin{subproof}
                Fix $A$.
                Assume $A\in T_0$.
                Fix $B$.
                Assume $B\in T_0$.
                $T_0 = \pow{\naturalsPlus}$ by \cref{discrete_topology}.
                Thus $A\in\pow{\naturalsPlus}$ and $B\in\pow{\naturalsPlus}$ by \cref{subseteq}.
                Thus $A\inter B\in\pow{\naturalsPlus}$ by \cref{inter_powerset}.
                Thus $A\inter B\in T_0$ by \cref{discrete_topology}.
            \end{subproof}
            Thus $T_0$ is closed under binary intersections.
        \end{subproof}
        Thus $T_0$ is a topology on $\naturalsPlus$ by \cref{topology_on}.
    \end{subproof}

    Show that $T_1$ is a topology on $Y$.
    \begin{subproof}
        Show that $T_1\subseteq\pow{Y}$.
        \begin{subproof}
            Show that for all $U$ such that $U\in T_1$ we have $U\in\pow{Y}$.
            \begin{subproof}
                Fix $U$.
                Assume $U\in T_1$.
                Thus $U = \emptyset$ or $U = Y$ by \cref{upair_iff}.
                \begin{byCase}
                \caseOf{$U = \emptyset$.}
                    $U\subseteq Y$ by \cref{emptyset_subseteq}.
                    Thus $U\in\pow{Y}$ by \cref{powerset_intro}.
                \caseOf{$U = Y$.}
                    $U\subseteq Y$ by \cref{subseteq_refl}.
                    Thus $U\in\pow{Y}$ by \cref{powerset_intro}.
                \end{byCase}
            \end{subproof}
            Thus $T_1\subseteq\pow{Y}$ by \cref{subseteq}.
        \end{subproof}
        $\emptyset\in T_1$ by \cref{upair_intro_left}.
        $Y\in T_1$ by \cref{upair_intro_right}.
        Show that $T_1$ is closed under arbitrary unions.
        \begin{subproof}
            Fix $M$.
            Assume $M\subseteq T_1$.
            \begin{byCase}
            \caseOf{$Y\in M$.}
                Show that $\unions{M} = Y$.
                \begin{subproof}
                    Show that $\unions{M}\subseteq Y$.
                    \begin{subproof}
                        $T_1\subseteq\pow{Y}$.
                        Thus $M\subseteq\pow{Y}$ by \cref{subseteq_transitive}.
                        Thus $\unions{M}\in\pow{Y}$ by \cref{powerset_closed_unions}.
                        Thus $\unions{M}\subseteq Y$ by \cref{pow_iff}.
                    \end{subproof}
                    Show that $Y\subseteq\unions{M}$.
                    \begin{subproof}
                        Show that for all $y$ such that $y\in Y$ we have $y\in\unions{M}$.
                        \begin{subproof}
                            Fix $y$.
                            Assume $y\in Y$.
                            Thus $y\in Y$ and $Y\in M$.
                            Thus $y\in\unions{M}$ by \cref{unions_iff}.
                        \end{subproof}
                        Thus $Y\subseteq\unions{M}$ by \cref{subseteq}.
                    \end{subproof}
                    Thus $\unions{M} = Y$ by \cref{subseteq_antisymmetric}.
                \end{subproof}
                Thus $\unions{M}\in T_1$ by \cref{upair_intro_right}.
            \caseOf{$Y\notin M$.}
                Show that $\unions{M} = \emptyset$.
                \begin{subproof}
                    Show that for all $U$ such that $U\in M$ we have $U\subseteq\emptyset$.
                    \begin{subproof}
                        Fix $U$.
                        Assume $U\in M$.
                        $U\in T_1$ by \cref{subseteq}.
                        Thus $U = \emptyset$ or $U = Y$ by \cref{upair_iff}.
                        Show that $U = \emptyset$.
                        \begin{subproof}
                            Suppose not.
                            Then $U \neq \emptyset$.
                            \begin{byCase}
                            \caseOf{$U = \emptyset$.}
                                Contradiction.
                            \caseOf{$U = Y$.}
                                Thus $Y\in M$.
                                Contradiction.
                            \end{byCase}
                        \end{subproof}
                        Thus $U\subseteq\emptyset$ by \cref{subseteq_refl}.
                    \end{subproof}
                    Thus $\unions{M}\subseteq\emptyset$ by \cref{unions_family}.
                    Thus $\unions{M} = \emptyset$ by \cref{subseteq_emptyset_iff}.
                \end{subproof}
                Thus $\unions{M}\in T_1$ by \cref{upair_intro_left}.
            \end{byCase}
        \end{subproof}
        Show that $T_1$ is closed under binary intersections.
        \begin{subproof}
            Show that for all $A$ such that $A\in T_1$ we have
                for all $B$ such that $B\in T_1$ we have $A\inter B\in T_1$.
            \begin{subproof}
                Fix $A$.
                Assume $A\in T_1$.
                Fix $B$.
                Assume $B\in T_1$.
                Thus $A = \emptyset$ or $A = Y$ by \cref{upair_iff}.
                Thus $B = \emptyset$ or $B = Y$ by \cref{upair_iff}.
                \begin{byCase}
                \caseOf{$A = \emptyset$.}
                    Thus $A\inter B = \emptyset$ by \cref{inter_comm,inter_emptyset}.
                    Thus $A\inter B\in T_1$ by \cref{upair_intro_left}.
                \caseOf{$A = Y$.}
                    \begin{byCase}
                    \caseOf{$B = \emptyset$.}
                        Thus $A\inter B = \emptyset$ by \cref{inter_emptyset}.
                        Thus $A\inter B\in T_1$ by \cref{upair_intro_left}.
                    \caseOf{$B = Y$.}
                        Thus $A\inter B = Y$ by \cref{inter_idempotent}.
                        Thus $A\inter B\in T_1$ by \cref{upair_intro_right}.
                    \end{byCase}
                \end{byCase}
            \end{subproof}
            Thus $T_1$ is closed under binary intersections.
        \end{subproof}
        Thus $T_1$ is a topology on $Y$ by \cref{topology_on}.
    \end{subproof}

    $\productBasis{T_0}{T_1}{\naturalsPlus}{Y}$ is a basis on $X$ by \cref{product_basis_is_basis}.
    Thus $T$ is a topology on $X$ by \cref{basis_topology_is_topology,product_topology}.

    Show that $X$ is limit point compact with respect to $T$.
    \begin{subproof}

        Show that for all $A$ such that $A\subseteq X$ and $A$ is infinite
            there exists $x\in X$ such that $x$ is a limit point of $A$ in $X$ with respect to $T$.
        \begin{subproof}
            Fix $A$.
            Assume $A\subseteq X$ and $A$ is infinite.
            Show that $A\neq\emptyset$.
            \begin{subproof}
                Suppose not.
                Then $A = \emptyset$.
                Thus $A$ is finite by \cref{emptyset_is_finite}.
                Contradiction.
            \end{subproof}
            Take $a$ such that $a\in A$ by \cref{nonempty_inhabited}.
            Thus $a\in X$ by \cref{subseteq}.
            Take $n,y$ such that $a = (n,y)$ and $n\in\naturalsPlus$ and $y\in Y$ by \cref{times_elem_is_tuple}.

            Show that there exists $y'$ such that $y'\in Y$ and $y'\neq y$.
            \begin{subproof}
                Thus $y = 0$ or $y = 1$ by \cref{upair_iff}.
                \begin{byCase}
                \caseOf{$y = 0$.}
                    Let $y' = 1$.
                    $y'\in Y$ by \cref{upair_intro_right}.
                    Show that $y'\neq y$.
                    \begin{subproof}
                        Suppose not.
                        Thus $y' = y$.
                        Thus $1 = 0$.
                        Contradiction by \cref{real_mul_ident}.
                    \end{subproof}
                    Thus there exists $y'$ such that $y'\in Y$ and $y'\neq y$.
                \caseOf{$y = 1$.}
                    Let $y' = 0$.
                    $y'\in Y$ by \cref{upair_intro_left}.
                    Show that $y'\neq y$.
                    \begin{subproof}
                        Suppose not.
                        Thus $y' = y$.
                        Thus $0 = 1$.
                        Thus $1 = 0$.
                        Contradiction by \cref{real_mul_ident}.
                    \end{subproof}
                    Thus there exists $y'$ such that $y'\in Y$ and $y'\neq y$.
                \end{byCase}
            \end{subproof}
            Take $y'$ such that $y'\in Y$ and $y'\neq y$.
            Let $x = (n,y')$.
            $x\in X$ by \cref{times_tuple_intro}.

            Show that $x$ is a limit point of $A$ in $X$ with respect to $T$.
            \begin{subproof}
                Show that for all $U$ such that $U$ is a neighborhood of $x$ in $X$ with respect to $T$
                    we have $U$ intersects $A\setminus\{x\}$.
                \begin{subproof}
                    Fix $U$.
                    Assume $U$ is a neighborhood of $x$ in $X$ with respect to $T$.
                    Thus $U$ is open in $X$ with respect to $T$ and $x\in U$ by \cref{neighborhood}.
                    Thus $U\in T$ by \cref{open_in}.
                    $U\in\basisTopology{\productBasis{T_0}{T_1}{\naturalsPlus}{Y}}{X}$ by \cref{product_topology}.
                    Take $B\in\productBasis{T_0}{T_1}{\naturalsPlus}{Y}$ such that $x\in B$ and $B\subseteq U$
                        by \cref{topology_generated_by_basis}.
                    Take $U_0,V_0$ such that $U_0\in T_0$ and $V_0\in T_1$ and $B = U_0\times V_0$
                        by \cref{product_basis}.
                    Thus $(n,y')\in U_0\times V_0$.
                    Thus $n\in U_0$ and $y'\in V_0$ by \cref{times_tuple_elim}.
                    Thus $V_0 = \emptyset$ or $V_0 = Y$ by \cref{upair_iff}.
                    Show that $a\in U$.
                    \begin{subproof}
                        Show that $V_0 = Y$.
                        \begin{subproof}
                            Suppose not.
                            Then $V_0 \neq Y$.
                            \begin{byCase}
                            \caseOf{$V_0 = \emptyset$.}
                                Thus $y'\in\emptyset$.
                                Contradiction by \cref{emptyset}.
                            \caseOf{$V_0 = Y$.}
                                Contradiction.
                            \end{byCase}
                        \end{subproof}
                        Thus $B = U_0\times Y$.
                        Thus $(n,y)\in U_0\times Y$ by \cref{times_tuple_intro}.
                        Thus $(n,y)\in B$.
                        Thus $a\in U$ by \cref{subseteq}.
                    \end{subproof}
                    Show that $a\neq x$.
                    \begin{subproof}
                        Suppose not.
                        Thus $a = x$.
                        Thus $(n,y) = (n,y')$.
                        Thus $y = y'$ by \cref{pair_eq_iff}.
                        Contradiction.
                    \end{subproof}
                    Show that $a\in A\setminus\{x\}$.
                    \begin{subproof}
                        Show that $a\notin\{x\}$.
                        \begin{subproof}
                            Suppose not.
                            Thus $a\in\{x\}$.
                            Thus $a = x$ by \cref{singleton_iff}.
                            Contradiction.
                        \end{subproof}
                        Thus $a\in A\setminus\{x\}$ by \cref{setminus_intro}.
                    \end{subproof}
                    Show that $U$ intersects $A\setminus\{x\}$.
                    \begin{subproof}
                        Suppose not.
                        Then $U\inter (A\setminus\{x\}) = \emptyset$ by \cref{intersects}.
                        Thus $a\in U\inter (A\setminus\{x\})$ by \cref{inter_intro}.
                        Contradiction by \cref{emptyset}.
                    \end{subproof}
                \end{subproof}
                Thus $x$ is a limit point of $A$ in $X$ with respect to $T$ by \cref{limit_point}.
            \end{subproof}
        \end{subproof}
        Thus $X$ is limit point compact with respect to $T$ by \cref{limit_point_compact}.
    \end{subproof}

    Show that $X$ is not compact with respect to $T$.
    \begin{subproof}
        Suppose not.
        Thus $X$ is compact with respect to $T$.
        Let $\pi = \piOne{\naturalsPlus}{Y}$.
        $\pi$ is continuous from $X$ to $\naturalsPlus$ with respect to $T$ and $T_0$
            by \cref{projection_one_continuous}.
        Show that $\img{\pi}{X} = \naturalsPlus$.
        \begin{subproof}
            Show that $\img{\pi}{X}\subseteq\naturalsPlus$.
            \begin{subproof}
                Show that for all $z$ such that $z\in\img{\pi}{X}$ we have $z\in\naturalsPlus$.
                \begin{subproof}
                    Fix $z$.
                    Assume $z\in\img{\pi}{X}$.
                    Take $p\in X$ such that $(p,z)\in\pi$ by \cref{img_iff}.
                    Take $n_0,y_0$ such that $p = (n_0,y_0)$ and $n_0\in\naturalsPlus$ and $y_0\in Y$
                        by \cref{times_elem_is_tuple}.
                    Take $n_1,y_1$ such that $n_1\in\naturalsPlus$ and $y_1\in Y$ and $(p,z) = ((n_1,y_1),n_1)$
                        by \cref{projection_first}.
                    Thus $(n_0,y_0) = (n_1,y_1)$ by \cref{pair_eq_iff}.
                    Thus $z = n_1$ by \cref{pair_eq_iff}.
                    Thus $z = n_0$ by \cref{pair_eq_iff}.
                    Thus $z\in\naturalsPlus$.
                \end{subproof}
                Thus $\img{\pi}{X}\subseteq\naturalsPlus$ by \cref{subseteq}.
            \end{subproof}
            Show that $\naturalsPlus\subseteq\img{\pi}{X}$.
            \begin{subproof}
                Show that for all $n$ such that $n\in\naturalsPlus$ we have $n\in\img{\pi}{X}$.
                \begin{subproof}
                    Fix $n$.
                    Assume $n\in\naturalsPlus$.
                    $0\in Y$ by \cref{upair_intro_left}.
                    Thus $(n,0)\in X$ by \cref{times_tuple_intro}.
                    Thus $((n,0),n)\in\pi$ by \cref{projection_first}.
                    Thus $n\in\img{\pi}{X}$ by \cref{img_iff}.
                \end{subproof}
                Thus $\naturalsPlus\subseteq\img{\pi}{X}$ by \cref{subseteq}.
            \end{subproof}
            Thus $\img{\pi}{X} = \naturalsPlus$ by \cref{subseteq_antisymmetric}.
        \end{subproof}
        Let $T_A = \subspaceTopology{T_0}{\naturalsPlus}$.
        Thus $\naturalsPlus$ is compact with respect to $T_A$ by \cref{continuous_image_compact}.
        Show that $T_A = T_0$.
        \begin{subproof}
            Show that $T_A\subseteq T_0$.
            \begin{subproof}
                Show that for all $U$ such that $U\in T_A$ we have $U\in T_0$.
                \begin{subproof}
                    Fix $U$.
                    Assume $U\in T_A$.
                    Take $V\in T_0$ such that $U = \naturalsPlus\inter V$ by \cref{subspace_topology}.
                    $T_0 = \pow{\naturalsPlus}$ by \cref{discrete_topology}.
                    Thus $V\subseteq\naturalsPlus$ by \cref{pow_iff}.
                    Thus $\naturalsPlus\inter V = V$ by \cref{inter_absorb_supseteq_left,subseteq}.
                    Thus $U = V$.
                    Thus $U\in T_0$.
                \end{subproof}
                Thus $T_A\subseteq T_0$ by \cref{subseteq}.
            \end{subproof}
            Show that $T_0\subseteq T_A$.
            \begin{subproof}
                Show that for all $U$ such that $U\in T_0$ we have $U\in T_A$.
                \begin{subproof}
                    Fix $U$.
                    Assume $U\in T_0$.
                    $T_0 = \pow{\naturalsPlus}$ by \cref{discrete_topology}.
                    Thus $U\subseteq\naturalsPlus$ by \cref{pow_iff}.
                    Thus $U = \naturalsPlus\inter U$ by \cref{inter_absorb_supseteq_left,subseteq}.
                    Thus $U\in T_A$ by \cref{subspace_topology}.
                \end{subproof}
                Thus $T_0\subseteq T_A$ by \cref{subseteq}.
            \end{subproof}
            Thus $T_A = T_0$ by \cref{subseteq_antisymmetric}.
        \end{subproof}
        Thus $\naturalsPlus$ is compact with respect to $T_0$.

        Let $C = \{\cons{n}{\emptyset} \mid n\in\naturalsPlus\}$.
        Show that $C$ is a covering of $\naturalsPlus$.
        \begin{subproof}
            Show that $\naturalsPlus\subseteq\unions{C}$.
            \begin{subproof}
                Show that for all $n$ such that $n\in\naturalsPlus$ we have $n\in\unions{C}$.
                \begin{subproof}
                    Fix $n$.
                    Assume $n\in\naturalsPlus$.
                    Thus $\cons{n}{\emptyset}\in C$.
                    $n\in\cons{n}{\emptyset}$ by \cref{cons_left}.
                    Thus $n\in\unions{C}$ by \cref{unions_iff}.
                \end{subproof}
                Thus $\naturalsPlus\subseteq\unions{C}$ by \cref{subseteq}.
            \end{subproof}
            Thus $C$ is a covering of $\naturalsPlus$ by \cref{covering}.
        \end{subproof}
        Show that for all $U$ such that $U\in C$ we have $U$ is open in $\naturalsPlus$ with respect to $T_0$.
        \begin{subproof}
            Fix $U$.
            Assume $U\in C$.
            Take $n\in\naturalsPlus$ such that $U = \cons{n}{\emptyset}$.
            Show that $U\in T_0$.
            \begin{subproof}
                $n\in\naturalsPlus$.
                $\emptyset\subseteq\naturalsPlus$ by \cref{emptyset_subseteq}.
                Thus $\cons{n}{\emptyset}\subseteq\naturalsPlus$ by \cref{cons_subseteq_intro}.
                Thus $\cons{n}{\emptyset}\in\pow{\naturalsPlus}$ by \cref{powerset_intro}.
                Thus $U\in T_0$ by \cref{discrete_topology}.
            \end{subproof}
            Thus $U$ is open in $\naturalsPlus$ with respect to $T_0$ by \cref{open_in}.
        \end{subproof}
        Thus $C$ is an open covering of $\naturalsPlus$ with respect to $T_0$ by \cref{open_covering}.
        Take $B$ such that
            $B\subseteq C$ and
            $B$ is finite and
            $B$ is a covering of $\naturalsPlus$ by \cref{compact}.
        Thus $\naturalsPlus\subseteq\unions{B}$ by \cref{covering}.
        Show that $\unions{B}$ is finite.
        \begin{subproof}
            Show that for all $U$ such that $U\in B$ we have $U$ is finite.
            \begin{subproof}
                Fix $U$.
                Assume $U\in B$.
                Thus $U\in C$ by \cref{subseteq}.
                Take $n\in\naturalsPlus$ such that $U = \cons{n}{\emptyset}$.
                $\{n,n\}$ is finite by \cref{pair_is_finite}.
                $n\in\{n,n\}$ by \cref{upair_intro_left}.
                $\emptyset\subseteq\{n,n\}$ by \cref{emptyset_subseteq}.
                Thus $\cons{n}{\emptyset}\subseteq\{n,n\}$ by \cref{cons_subseteq_intro}.
                Thus $U$ is finite by \cref{subseteq_finite_is_finite}.
            \end{subproof}
            Thus $\unions{B}$ is finite by \cref{finite_union_of_finite_family}.
        \end{subproof}
        Thus $\naturalsPlus$ is finite by \cref{subseteq_finite_is_finite}.
        Contradiction by \cref{naturalsplus_infinite}.
        Contradiction.
    \end{subproof}
    Thus there exists $X,T$ such that
        $X$ is limit point compact with respect to $T$ and
        $X$ is not compact with respect to $T$.
\end{proof}

% from §28 helper definition 1: strictly increasing on a set
\begin{definition}\label{strictly_increasing_on_set}
    $u$ is strictly increasing on $A$ iff
        $u\in\funs{A}{A}$ and
        for all $m$ such that $m\in A$ we have
            for all $n$ such that $n\in A$ and $m < n$ we have $u(m) < u(n)$.
\end{definition}

% from §28 Item 4: subsequence
\begin{definition}\label{subsequence}
    $t$ is a subsequence of $s$ iff
        there exists $u$ such that
            $u$ is strictly increasing on $\naturalsPlus$ and
            for all $n\in\naturalsPlus$ we have $t(n) = \apply{s}{\apply{u}{n}}$.
\end{definition}

% from §28 Item 5: sequentially compact
\begin{definition}\label{sequentially_compact}
    $X$ is sequentially compact with respect to $T$ iff
        $T$ is a topology on $X$ and
        for all $s$ such that $s$ is a sequence in $X$
            there exist $t,x$ such that
                $t$ is a sequence in $X$ and
                $t$ is a subsequence of $s$ and
                $x\in X$ and
                $t$ converges to $x$ in $X$ with respect to $T$.
\end{definition}

% from §28 helper lemma 2: metric spaces are Hausdorff
\begin{lemma}\label{metric_space_hausdorff}
    Suppose $d$ is a metric on $X$.
    Let $T = \metricTopology{d}{X}$.
    Then $X$ is Hausdorff with respect to $T$.
\end{lemma}
\begin{proof}
    Let $B = \metricBasis{d}{X}$.
    $B$ is a basis on $X$ by \cref{metric_basis_is_basis}.
    Thus $T$ is a topology on $X$ by \cref{basis_topology_is_topology,metric_topology}.
    Show that for all $x_1,x_2$ such that $x_1\in X$ and $x_2\in X$ and $x_1\neq x_2$
        there exist $U_1,U_2$ such that
            $U_1$ is a neighborhood of $x_1$ in $X$ with respect to $T$ and
            $U_2$ is a neighborhood of $x_2$ in $X$ with respect to $T$ and
            $U_1$ is disjoint from $U_2$.
    \begin{subproof}
        Fix $x_1$.
        Fix $x_2$.
        Assume $x_1\in X$ and $x_2\in X$ and $x_1\neq x_2$.
        Let $r = d((x_1,x_2))$.
        Show that $r\in\reals$.
        \begin{subproof}
            $d\in\funs{X\times X}{\reals}$ by \cref{metric_on}.
            $(x_1,x_2)\in X\times X$ by \cref{times_tuple_intro}.
            Thus $r\in\reals$ by \cref{funs_type_apply}.
        \end{subproof}
        Show that $0 < r$.
        \begin{subproof}
            $d$ is nonnegative on $X$ by \cref{metric_on}.
            Thus $r \geq 0$ by \cref{metric_nonnegative}.
            $d$ is zero-characterized on $X$ by \cref{metric_on}.
            Thus $r = 0$ iff $x_1 = x_2$ by \cref{metric_zero_characterization}.
            Thus $r\neq 0$.
            Thus $0 < r$ by \cref{real_leq_split}.
        \end{subproof}
        Let $m_2 = 1 + 1$.
        Show that $m_2\in\reals$ and $0 < m_2$.
        \begin{subproof}
            $1\in\naturalsPlus$ by \cref{real_one_in_naturals}.
            Thus $1\in\reals$ by \cref{real_naturals_subset_reals,subseteq}.
            $m_2\in\reals$ by \cref{real_add_closed}.
            $0\in\reals$ by \cref{real_add_ident}.
            $0 < 1$ by \cref{real_zero_lt_one}.
            Thus $0 + 1 < 1 + 1$ by \cref{real_lt_add}.
            $0 + 1 = 1$ by \cref{real_add_ident}.
            Thus $1 < m_2$.
            Thus $0 < m_2$ by \cref{real_lt_trans}.
        \end{subproof}
        Show that $m_2 \neq 0$.
        \begin{subproof}
            Suppose not.
            Then $m_2 = 0$.
            Thus $0 < 0$.
            Contradiction by \cref{real_lt_irreflexive}.
        \end{subproof}
        Let $\epsilon = r / m_2$.
        Show that $\epsilon\in\reals$ and $0 < \epsilon$.
        \begin{subproof}
            $\epsilon = r \rmul \inv{m_2}$ by \cref{real_div}.
            $\inv{m_2}\in\reals$ by \cref{real_mul_inverse}.
            Thus $\epsilon\in\reals$ by \cref{real_mul_closed}.
            Thus $0 < \inv{m_2}$ by \cref{real_inv_pos}.
            $0\in\reals$ by \cref{real_add_ident}.
            Thus $0 \rmul \inv{m_2} < r \rmul \inv{m_2}$ by \cref{real_lt_mul_pos}.
            $0 \rmul \inv{m_2} = 0$ by \cref{real_mul_zero}.
            Thus $0 < r \rmul \inv{m_2}$.
            Thus $0 < \epsilon$.
        \end{subproof}
        Let $U_1 = \metricBall{d}{X}{x_1}{\epsilon}$.
        Let $U_2 = \metricBall{d}{X}{x_2}{\epsilon}$.
        Show that $U_1$ is a neighborhood of $x_1$ in $X$ with respect to $T$.
        \begin{subproof}
            Show that $U_1$ is open in $X$ with respect to $T$.
            \begin{subproof}
                Show that $U_1\subseteq X$.
                \begin{subproof}
                    Show that for all $y\in U_1$ we have $y\in X$.
                    \begin{subproof}
                        Fix $y$.
                        Assume $y\in U_1$.
                        Thus $y\in X$ by \cref{metric_ball}.
                    \end{subproof}
                \end{subproof}
                Show that for all $y\in U_1$ there exists $\delta\in\reals$ such that
                    $0 < \delta$ and $\metricBall{d}{X}{y}{\delta}\subseteq U_1$.
                \begin{subproof}
                    Fix $y$.
                    Assume $y\in U_1$.
                    Thus $y\in X$ and $d((x_1,y)) < \epsilon$ by \cref{metric_ball}.
                    Show that $d((x_1,y))\in\reals$.
                    \begin{subproof}
                        $d\in\funs{X\times X}{\reals}$ by \cref{metric_on}.
                        $(x_1,y)\in X\times X$ by \cref{times_tuple_intro}.
                        Thus $d((x_1,y))\in\reals$ by \cref{funs_type_apply}.
                    \end{subproof}
                    Let $\delta = \epsilon - d((x_1,y))$.
                    $0 < \delta$ by \cref{real_lt_imp_pos_diff}.
                    Show that $\metricBall{d}{X}{y}{\delta}\subseteq U_1$.
                    \begin{subproof}
                        Show that for all $z$ such that $z\in\metricBall{d}{X}{y}{\delta}$ we have $z\in U_1$.
                        \begin{subproof}
                            Fix $z$.
                            Assume $z\in\metricBall{d}{X}{y}{\delta}$.
                            Thus $z\in X$ and $d((y,z)) < \delta$ by \cref{metric_ball}.
                            Show that $d((y,z))\in\reals$.
                            \begin{subproof}
                                $d\in\funs{X\times X}{\reals}$ by \cref{metric_on}.
                                $(y,z)\in X\times X$ by \cref{times_tuple_intro}.
                                Thus $d((y,z))\in\reals$ by \cref{funs_type_apply}.
                            \end{subproof}
                            Show that $\delta\in\reals$.
                            \begin{subproof}
                                $\delta = \epsilon - d((x_1,y))$ by \cref{real_sub}.
                                $\epsilon\in\reals$.
                                $\neg{d((x_1,y))}\in\reals$ by \cref{real_add_inverse}.
                                Thus $\delta\in\reals$ by \cref{real_sub,real_add_closed}.
                            \end{subproof}
                            $d$ satisfies the triangle inequality on $X$ by \cref{metric_on}.
                            Thus $d((x_1,z)) \leq d((x_1,y)) + d((y,z))$ by \cref{metric_triangle_inequality}.
                            Thus $d((y,z)) + d((x_1,y)) < \delta + d((x_1,y))$ by \cref{real_lt_add}.
                            Thus $d((x_1,y)) + d((y,z)) < d((x_1,y)) + \delta$ by \cref{real_add_comm}.
                            Show that $d((x_1,y)) + \delta = \epsilon$.
                            \begin{subproof}
                                $\epsilon - d((x_1,y)) = \epsilon + \neg{d((x_1,y))}$ by \cref{real_sub}.
                                Thus $\delta = \epsilon + \neg{d((x_1,y))}$ by \cref{real_sub}.
                                Thus $d((x_1,y)) + \delta = d((x_1,y)) + (\epsilon + \neg{d((x_1,y))})$.
                                $d((x_1,y))\in\reals$.
                                $\epsilon\in\reals$.
                                $\neg{d((x_1,y))}\in\reals$ by \cref{real_add_inverse}.
                                Thus $d((x_1,y)) + (\epsilon + \neg{d((x_1,y))})
                                    = (d((x_1,y)) + \epsilon) + \neg{d((x_1,y))}$ by \cref{real_add_assoc}.
                                $d((x_1,y)) + \epsilon = \epsilon + d((x_1,y))$ by \cref{real_add_comm}.
                                Thus $(d((x_1,y)) + \epsilon) + \neg{d((x_1,y))}
                                    = (\epsilon + d((x_1,y))) + \neg{d((x_1,y))}$ by \cref{real_add_eq_subst}.
                                Thus $(\epsilon + d((x_1,y))) + \neg{d((x_1,y))}
                                    = \epsilon + (d((x_1,y)) + \neg{d((x_1,y))})$ by \cref{real_add_assoc}.
                                $d((x_1,y)) + \neg{d((x_1,y))} = 0$ by \cref{real_add_inverse}.
                                Thus $d((x_1,y)) + \delta = \epsilon + 0$.
                                Thus $d((x_1,y)) + \delta = \epsilon$ by \cref{real_add_ident,real_add_comm}.
                            \end{subproof}
                            Show that $d((x_1,z)) < \epsilon$.
                            \begin{subproof}
                                Show that $d((x_1,z))\in\reals$.
                                \begin{subproof}
                                    $d\in\funs{X\times X}{\reals}$ by \cref{metric_on}.
                                    $(x_1,z)\in X\times X$ by \cref{times_tuple_intro}.
                                    Thus $d((x_1,z))\in\reals$ by \cref{funs_type_apply}.
                                \end{subproof}
                                Show that $d((x_1,y)) + d((y,z))\in\reals$.
                                \begin{subproof}
                                    $d((x_1,y))\in\reals$.
                                    $d((y,z))\in\reals$.
                                    Thus $d((x_1,y)) + d((y,z))\in\reals$ by \cref{real_add_closed}.
                                \end{subproof}
                                $\epsilon\in\reals$.
                                Thus $d((x_1,z)) < d((x_1,y)) + d((y,z))$ or
                                    $d((x_1,z)) = d((x_1,y)) + d((y,z))$ by \cref{real_leq_split}.
                                \begin{byCase}
                                \caseOf{$d((x_1,z)) < d((x_1,y)) + d((y,z))$.}
                                    Thus $d((x_1,z)) < \epsilon$ by \cref{real_lt_trans}.
                                \caseOf{$d((x_1,z)) = d((x_1,y)) + d((y,z))$.}
                                    Thus $d((x_1,z)) < \epsilon$ by \cref{real_lt_subst_left}.
                                \end{byCase}
                            \end{subproof}
                            Thus $z\in U_1$ by \cref{metric_ball}.
                        \end{subproof}
                        Thus $\metricBall{d}{X}{y}{\delta}\subseteq U_1$ by \cref{subseteq}.
                    \end{subproof}
                \end{subproof}
                Thus $U_1$ is open in $X$ with respect to $T$ by \cref{metric_open_iff}.
            \end{subproof}
            Show that $x_1\in U_1$.
            \begin{subproof}
                Show that $d((x_1,x_1)) = 0$.
                \begin{subproof}
                    $d$ is zero-characterized on $X$ by \cref{metric_on}.
                    Thus $d((x_1,x_1)) = 0$ iff $x_1 = x_1$ by \cref{metric_zero_characterization}.
                    Thus $d((x_1,x_1)) = 0$.
                \end{subproof}
                Thus $x_1\in U_1$ by \cref{metric_ball}.
            \end{subproof}
            Thus $U_1$ is a neighborhood of $x_1$ in $X$ with respect to $T$ by \cref{neighborhood}.
        \end{subproof}
        Show that $U_2$ is a neighborhood of $x_2$ in $X$ with respect to $T$.
        \begin{subproof}
            Show that $U_2$ is open in $X$ with respect to $T$.
            \begin{subproof}
                Show that $U_2\subseteq X$.
                \begin{subproof}
                    Show that for all $y\in U_2$ we have $y\in X$.
                    \begin{subproof}
                        Fix $y$.
                        Assume $y\in U_2$.
                        Thus $y\in X$ by \cref{metric_ball}.
                    \end{subproof}
                \end{subproof}
                Show that for all $y\in U_2$ there exists $\delta\in\reals$ such that
                    $0 < \delta$ and $\metricBall{d}{X}{y}{\delta}\subseteq U_2$.
                \begin{subproof}
                Fix $y$.
                Assume $y\in U_2$.
                Thus $y\in X$ and $d((x_2,y)) < \epsilon$ by \cref{metric_ball}.
                Show that $d((x_2,y))\in\reals$.
                \begin{subproof}
                    $d\in\funs{X\times X}{\reals}$ by \cref{metric_on}.
                    $(x_2,y)\in X\times X$ by \cref{times_tuple_intro}.
                    Thus $d((x_2,y))\in\reals$ by \cref{funs_type_apply}.
                \end{subproof}
                Let $\delta = \epsilon - d((x_2,y))$.
                $0 < \delta$ by \cref{real_lt_imp_pos_diff}.
                    Show that $\metricBall{d}{X}{y}{\delta}\subseteq U_2$.
                    \begin{subproof}
                        Show that for all $z$ such that $z\in\metricBall{d}{X}{y}{\delta}$ we have $z\in U_2$.
                        \begin{subproof}
                            Fix $z$.
                            Assume $z\in\metricBall{d}{X}{y}{\delta}$.
                            Thus $z\in X$ and $d((y,z)) < \delta$ by \cref{metric_ball}.
                            Show that $d((y,z))\in\reals$.
                            \begin{subproof}
                                $d\in\funs{X\times X}{\reals}$ by \cref{metric_on}.
                                $(y,z)\in X\times X$ by \cref{times_tuple_intro}.
                                Thus $d((y,z))\in\reals$ by \cref{funs_type_apply}.
                            \end{subproof}
                            Show that $\delta\in\reals$.
                            \begin{subproof}
                                $\delta = \epsilon - d((x_2,y))$ by \cref{real_sub}.
                                $\epsilon\in\reals$.
                                $\neg{d((x_2,y))}\in\reals$ by \cref{real_add_inverse}.
                                Thus $\delta\in\reals$ by \cref{real_sub,real_add_closed}.
                            \end{subproof}
                            $d$ satisfies the triangle inequality on $X$ by \cref{metric_on}.
                            Thus $d((x_2,z)) \leq d((x_2,y)) + d((y,z))$ by \cref{metric_triangle_inequality}.
                            Thus $d((y,z)) + d((x_2,y)) < \delta + d((x_2,y))$ by \cref{real_lt_add}.
                            Thus $d((x_2,y)) + d((y,z)) < d((x_2,y)) + \delta$ by \cref{real_add_comm}.
                            Show that $d((x_2,y)) + \delta = \epsilon$.
                            \begin{subproof}
                                $\epsilon - d((x_2,y)) = \epsilon + \neg{d((x_2,y))}$ by \cref{real_sub}.
                                Thus $\delta = \epsilon + \neg{d((x_2,y))}$ by \cref{real_sub}.
                                Thus $d((x_2,y)) + \delta = d((x_2,y)) + (\epsilon + \neg{d((x_2,y))})$.
                                $d((x_2,y))\in\reals$.
                                $\epsilon\in\reals$.
                                $\neg{d((x_2,y))}\in\reals$ by \cref{real_add_inverse}.
                                Thus $d((x_2,y)) + (\epsilon + \neg{d((x_2,y))})
                                    = (d((x_2,y)) + \epsilon) + \neg{d((x_2,y))}$ by \cref{real_add_assoc}.
                                $d((x_2,y)) + \epsilon = \epsilon + d((x_2,y))$ by \cref{real_add_comm}.
                                Thus $(d((x_2,y)) + \epsilon) + \neg{d((x_2,y))}
                                    = (\epsilon + d((x_2,y))) + \neg{d((x_2,y))}$ by \cref{real_add_eq_subst}.
                                Thus $(\epsilon + d((x_2,y))) + \neg{d((x_2,y))}
                                    = \epsilon + (d((x_2,y)) + \neg{d((x_2,y))})$ by \cref{real_add_assoc}.
                                $d((x_2,y)) + \neg{d((x_2,y))} = 0$ by \cref{real_add_inverse}.
                                Thus $d((x_2,y)) + \delta = \epsilon + 0$.
                                Thus $d((x_2,y)) + \delta = \epsilon$ by \cref{real_add_ident,real_add_comm}.
                            \end{subproof}
                            Show that $d((x_2,z)) < \epsilon$.
                            \begin{subproof}
                                Show that $d((x_2,z))\in\reals$.
                                \begin{subproof}
                                    $d\in\funs{X\times X}{\reals}$ by \cref{metric_on}.
                                    $(x_2,z)\in X\times X$ by \cref{times_tuple_intro}.
                                    Thus $d((x_2,z))\in\reals$ by \cref{funs_type_apply}.
                                \end{subproof}
                                Show that $d((x_2,y)) + d((y,z))\in\reals$.
                                \begin{subproof}
                                    $d((x_2,y))\in\reals$.
                                    $d((y,z))\in\reals$.
                                    Thus $d((x_2,y)) + d((y,z))\in\reals$ by \cref{real_add_closed}.
                                \end{subproof}
                                $\epsilon\in\reals$.
                                Thus $d((x_2,z)) < d((x_2,y)) + d((y,z))$ or
                                    $d((x_2,z)) = d((x_2,y)) + d((y,z))$ by \cref{real_leq_split}.
                                \begin{byCase}
                                \caseOf{$d((x_2,z)) < d((x_2,y)) + d((y,z))$.}
                                    Thus $d((x_2,z)) < \epsilon$ by \cref{real_lt_trans}.
                                \caseOf{$d((x_2,z)) = d((x_2,y)) + d((y,z))$.}
                                    Thus $d((x_2,z)) < \epsilon$ by \cref{real_lt_subst_left}.
                                \end{byCase}
                            \end{subproof}
                            Thus $z\in U_2$ by \cref{metric_ball}.
                        \end{subproof}
                        Thus $\metricBall{d}{X}{y}{\delta}\subseteq U_2$ by \cref{subseteq}.
                    \end{subproof}
                \end{subproof}
                Thus $U_2$ is open in $X$ with respect to $T$ by \cref{metric_open_iff}.
            \end{subproof}
            Show that $x_2\in U_2$.
            \begin{subproof}
                Show that $d((x_2,x_2)) = 0$.
                \begin{subproof}
                    $d$ is zero-characterized on $X$ by \cref{metric_on}.
                    Thus $d((x_2,x_2)) = 0$ iff $x_2 = x_2$ by \cref{metric_zero_characterization}.
                    Thus $d((x_2,x_2)) = 0$.
                \end{subproof}
                Thus $x_2\in U_2$ by \cref{metric_ball}.
            \end{subproof}
            Thus $U_2$ is a neighborhood of $x_2$ in $X$ with respect to $T$ by \cref{neighborhood}.
        \end{subproof}
        Show that $U_1$ is disjoint from $U_2$.
        \begin{subproof}
            Show that for all $y$ such that $y\in U_1$ we have $y\notin U_2$.
            \begin{subproof}
                Fix $y$.
                Assume $y\in U_1$.
                Suppose not.
                Then $y\in U_2$.
                Thus $y\in X$ and $d((x_1,y)) < \epsilon$ and $d((x_2,y)) < \epsilon$ by \cref{metric_ball}.
                $d$ satisfies the triangle inequality on $X$ by \cref{metric_on}.
                Thus $r \leq d((x_1,y)) + d((y,x_2))$ by \cref{metric_triangle_inequality}.
                $d$ is symmetric on $X$ by \cref{metric_on}.
                Thus $d((y,x_2)) = d((x_2,y))$ by \cref{metric_symmetric}.
                Thus $r \leq d((x_1,y)) + d((x_2,y))$.
                Show that $d((x_1,y))\in\reals$ and $d((x_2,y))\in\reals$.
                \begin{subproof}
                    $d\in\funs{X\times X}{\reals}$ by \cref{metric_on}.
                    $(x_1,y)\in X\times X$ by \cref{times_tuple_intro}.
                    $(x_2,y)\in X\times X$ by \cref{times_tuple_intro}.
                    Thus $d((x_1,y))\in\reals$ and $d((x_2,y))\in\reals$ by \cref{funs_type_apply}.
                \end{subproof}
                $\epsilon\in\reals$.
                Thus $d((x_2,y)) + d((x_1,y)) < \epsilon + d((x_1,y))$ by \cref{real_lt_add}.
                Thus $d((x_1,y)) + d((x_2,y)) < d((x_1,y)) + \epsilon$ by \cref{real_add_comm}.
                Thus $d((x_1,y)) + \epsilon < \epsilon + \epsilon$ by \cref{real_lt_add}.
                Show that $d((x_1,y)) + \epsilon\in\reals$.
                \begin{subproof}
                    $d((x_1,y))\in\reals$.
                    $\epsilon\in\reals$.
                    Thus $d((x_1,y)) + \epsilon\in\reals$ by \cref{real_add_closed}.
                \end{subproof}
                Show that $d((x_1,y)) + d((x_2,y))\in\reals$.
                \begin{subproof}
                    $d((x_1,y))\in\reals$.
                    $d((x_2,y))\in\reals$.
                    Thus $d((x_1,y)) + d((x_2,y))\in\reals$ by \cref{real_add_closed}.
                \end{subproof}
                Show that $\epsilon + \epsilon\in\reals$.
                \begin{subproof}
                    $\epsilon\in\reals$.
                    Thus $\epsilon + \epsilon\in\reals$ by \cref{real_add_closed}.
                \end{subproof}
                Thus $d((x_1,y)) + d((x_2,y)) < \epsilon + \epsilon$ by \cref{real_lt_trans}.
                Show that $r < \epsilon + \epsilon$.
                \begin{subproof}
                    Show that $d((x_1,y)) + d((x_2,y))\in\reals$.
                    \begin{subproof}
                        $d((x_1,y))\in\reals$.
                        $d((x_2,y))\in\reals$.
                        Thus $d((x_1,y)) + d((x_2,y))\in\reals$ by \cref{real_add_closed}.
                    \end{subproof}
                    Show that $\epsilon + \epsilon\in\reals$.
                    \begin{subproof}
                        $\epsilon\in\reals$.
                        Thus $\epsilon + \epsilon\in\reals$ by \cref{real_add_closed}.
                    \end{subproof}
                    $r\in\reals$.
                    Thus $r < d((x_1,y)) + d((x_2,y))$ or
                        $r = d((x_1,y)) + d((x_2,y))$ by \cref{real_leq_split}.
                    \begin{byCase}
                    \caseOf{$r < d((x_1,y)) + d((x_2,y))$.}
                        Thus $r < \epsilon + \epsilon$ by \cref{real_lt_trans}.
                    \caseOf{$r = d((x_1,y)) + d((x_2,y))$.}
                        Thus $r < \epsilon + \epsilon$ by \cref{real_lt_subst_left}.
                    \end{byCase}
                \end{subproof}
                Show that $\epsilon + \epsilon = r$.
                \begin{subproof}
                    $\epsilon = r \rmul \inv{m_2}$ by \cref{real_div}.
                    $(1 + 1) \rmul \epsilon = \epsilon + \epsilon$ by \cref{real_right_distrib,real_mul_ident}.
                    Thus $m_2 \rmul \epsilon = \epsilon + \epsilon$.
                    $m_2 \rmul \epsilon = m_2 \rmul (r \rmul \inv{m_2})$.
                    $m_2\in\reals$.
                    $r\in\reals$.
                    $\inv{m_2}\in\reals$ by \cref{real_mul_inverse}.
                    Thus $m_2 \rmul \epsilon = (m_2 \rmul r) \rmul \inv{m_2}$ by \cref{real_mul_assoc}.
                    $m_2 \rmul r = r \rmul m_2$ by \cref{real_mul_comm}.
                    Thus $m_2 \rmul \epsilon = r \rmul (m_2 \rmul \inv{m_2})$ by \cref{real_mul_assoc}.
                    $m_2 \rmul \inv{m_2} = 1$ by \cref{real_mul_inverse,real_mul_comm}.
                    Thus $m_2 \rmul \epsilon = r \rmul 1$.
                    $1\in\reals$ by \cref{real_mul_ident}.
                    $r \rmul 1 = 1 \rmul r$ by \cref{real_mul_comm}.
                    Thus $m_2 \rmul \epsilon = 1 \rmul r$.
                    Thus $m_2 \rmul \epsilon = r$ by \cref{real_mul_ident}.
                    Thus $\epsilon + \epsilon = r$.
                \end{subproof}
                Thus $r < r$.
                Contradiction by \cref{real_lt_irreflexive}.
            \end{subproof}
            Thus $U_1$ is disjoint from $U_2$ by \cref{disjoint}.
        \end{subproof}
    \end{subproof}
    Thus $X$ is Hausdorff with respect to $T$ by \cref{hausdorff}.
\end{proof}


% from §28 helper definition 1: iteration sequence predicate
\begin{definition}\label{iteration_sequence}
    $u$ is an iteration sequence for $f$ on $n$ starting at $a_0$ iff
        $u$ is a function and
        $\dom{u} = \initialSegment{n}$ and
        $u(1) = a_0$ and
        for all $k\in\naturalsPlus$ such that $k + 1 \leq n$ we have $u(k + 1) = \apply{f}{u(k)}$.
\end{definition}

% from §28 helper definition 2: iteration set
\begin{definition}\label{iteration_set}
    $\iterationSet{f}{a_0} = \{ n\in\naturalsPlus \mid \text{there exists $u$ such that
        $u$ is an iteration sequence for $f$ on $n$ starting at $a_0$} \}$.
\end{definition}

% from §28 helper lemma 3: iteration set introduction
\begin{lemma}\label{iteration_set_intro}
    Suppose $n\in\naturalsPlus$.
    Suppose there exists $u$ such that
        $u$ is an iteration sequence for $f$ on $n$ starting at $a_0$.
    Then $n\in\iterationSet{f}{a_0}$.
\end{lemma}
\begin{proof}
    Thus $n\in\iterationSet{f}{a_0}$ by \cref{iteration_set}.
\end{proof}

% from §28 helper lemma 4: iteration set elimination
\begin{lemma}\label{iteration_set_elim}
    Suppose $n\in\iterationSet{f}{a_0}$.
    Then $n\in\naturalsPlus$ and there exists $u$ such that
        $u$ is an iteration sequence for $f$ on $n$ starting at $a_0$.
\end{lemma}
\begin{proof}
    Thus $n\in\naturalsPlus$ and there exists $u$ such that
        $u$ is an iteration sequence for $f$ on $n$ starting at $a_0$
        by \cref{iteration_set}.
\end{proof}

% from §28 helper lemma 5: iteration sequence introduction
\begin{lemma}\label{iteration_sequence_intro}
    Suppose $u$ is a function and $\dom{u} = \initialSegment{n}$ and $u(1) = a_0$.
    Suppose for all $k\in\naturalsPlus$ such that $k + 1 \leq n$ we have $u(k + 1) = \apply{f}{u(k)}$.
    Then $u$ is an iteration sequence for $f$ on $n$ starting at $a_0$.
\end{lemma}
\begin{proof}
    Thus $u$ is an iteration sequence for $f$ on $n$ starting at $a_0$ by \cref{iteration_sequence}.
\end{proof}

% from §28 helper lemma 6: iteration sequence uniqueness
\begin{lemma}\label{iteration_sequence_unique}
    Suppose $n\in\naturalsPlus$.
    Suppose $u$ and $v$ are iteration sequences for $f$ on $n$ starting at $a_0$.
    Then $u = v$.
\end{lemma}
\begin{proof}
    Thus $u$ is a function and $v$ is a function by \cref{iteration_sequence}.
    Thus $\dom{u} = \initialSegment{n}$ and $\dom{v} = \initialSegment{n}$ by \cref{iteration_sequence}.
    Show that for all $k\in\initialSegment{n}$ we have $u(k) = v(k)$.
    \begin{subproof}
        Let $A = \{k\in\naturalsPlus \mid \text{if $k\leq n$ then $u(k) = v(k)$}\}$.
        Show that for all $k\in\naturalsPlus$ we have $k\in A$.
        \begin{subproof}
            Show that $A\subseteq\naturalsPlus$.
            \begin{subproof}
                Show that for all $k$ such that $k\in A$ we have $k\in\naturalsPlus$.
                \begin{subproof}
                    Fix $k$.
                    Assume $k\in A$.
                    Then $k\in\naturalsPlus$.
                \end{subproof}
                Thus $A\subseteq\naturalsPlus$ by \cref{subseteq}.
            \end{subproof}
            $1\in\naturalsPlus$ by \cref{real_one_in_naturals}.
            Show that $1\in A$.
            \begin{subproof}
                Show that if $1\leq n$ then $u(1) = v(1)$.
                \begin{subproof}
                    Assume $1\leq n$.
                    Then $u(1) = a_0$ and $v(1) = a_0$ by \cref{iteration_sequence}.
                    Thus $u(1) = v(1)$.
                \end{subproof}
                Thus $1\in A$.
            \end{subproof}
            Show that for all $k\in A$ we have $k + 1\in A$.
            \begin{subproof}
                Fix $k$.
                Assume $k\in A$.
                Show that if $k + 1\leq n$ then $u(k + 1) = v(k + 1)$.
                \begin{subproof}
                    Assume $k + 1\leq n$.
                    Then $k \leq n$.
                    Thus $u(k) = v(k)$.
                    Thus $u(k + 1) = \apply{f}{u(k)}$ and $v(k + 1) = \apply{f}{v(k)}$ by \cref{iteration_sequence}.
                    Thus $u(k + 1) = v(k + 1)$.
                \end{subproof}
                Thus $k + 1\in A$.
            \end{subproof}
            Thus for all $k\in\naturalsPlus$ we have $k\in A$ by \cref{real_naturals_induction}.
        \end{subproof}
        Show that for all $k\in\initialSegment{n}$ we have $u(k) = v(k)$.
        \begin{subproof}
            Fix $k$.
            Assume $k\in\initialSegment{n}$.
            Then $k\in\naturalsPlus$ and $k\leq n$ by \cref{initial_segment_naturalsplus}.
            Thus $k\in A$.
            Thus if $k\leq n$ then $u(k) = v(k)$.
            Thus $u(k) = v(k)$.
        \end{subproof}
    \end{subproof}
    Show that $u\subseteq v$.
    \begin{subproof}
        Show that for all $w$ such that $w\in u$ we have $w\in v$.
        \begin{subproof}
            Fix $w$.
            Assume $w\in u$.
            Thus $u$ is a function by \cref{iteration_sequence}.
            Take $k\in\dom{u}$ such that $w = (k,u(k))$ by \cref{function_member_elim}.
            Thus $\dom{u} = \initialSegment{n}$ by \cref{iteration_sequence}.
            Thus $k\in\initialSegment{n}$.
            Thus $u(k) = v(k)$.
            Thus $w = (k,v(k))$.
            Thus $\dom{v} = \initialSegment{n}$ by \cref{iteration_sequence}.
            Thus $k\in\dom{v}$.
            Thus $v$ is a function by \cref{iteration_sequence}.
            Thus $(k,v(k))\in v$ by \cref{function_apply_elim}.
            Thus $w\in v$.
        \end{subproof}
        Thus $u\subseteq v$ by \cref{subseteq}.
    \end{subproof}
    Show that $v\subseteq u$.
    \begin{subproof}
        Show that for all $w$ such that $w\in v$ we have $w\in u$.
        \begin{subproof}
            Fix $w$.
            Assume $w\in v$.
            Thus $v$ is a function by \cref{iteration_sequence}.
            Take $k\in\dom{v}$ such that $w = (k,v(k))$ by \cref{function_member_elim}.
            Thus $\dom{v} = \initialSegment{n}$ by \cref{iteration_sequence}.
            Thus $k\in\initialSegment{n}$.
            Thus $u(k) = v(k)$.
            Thus $w = (k,u(k))$.
            Thus $\dom{u} = \initialSegment{n}$ by \cref{iteration_sequence}.
            Thus $k\in\dom{u}$.
            Thus $u$ is a function by \cref{iteration_sequence}.
            Thus $(k,u(k))\in u$ by \cref{function_apply_elim}.
            Thus $w\in u$.
        \end{subproof}
        Thus $v\subseteq u$ by \cref{subseteq}.
    \end{subproof}
    Thus $u = v$ by \cref{subseteq_antisymmetric}.
\end{proof}

% from §28 helper lemma 7: infinite subset of naturalsPlus has increasing sequence
\begin{lemma}\label{infinite_naturalsplus_has_strictly_increasing_sequence}
    Suppose $A\subseteq\naturalsPlus$.
    Suppose $A$ is infinite.
    Then there exists $u$ such that
        $u$ is strictly increasing on $\naturalsPlus$ and
        for all $n\in\naturalsPlus$ we have $u(n)\in A$.
\end{lemma}
\begin{proof}
    Show that $A\neq\emptyset$.
    \begin{subproof}
        Suppose not.
        Then $A=\emptyset$.
        Thus $A$ is finite by \cref{emptyset_is_finite}.
        Contradiction.
    \end{subproof}
    Take $a_0\in A$ such that for all $a\in A$ we have $a_0 \leq a$
        by \cref{real_naturals_least_element}.
    Show that for all $a\in A$ there exists $b\in A$ such that $a < b$.
    \begin{subproof}
        Fix $a$.
        Assume $a\in A$.
        Thus $a\in\naturalsPlus$ by \cref{subseteq}.
        Suppose not.
        Then for all $b\in A$ we have $b \leq a$.
        Show that $A\subseteq\initialSegment{a}$.
        \begin{subproof}
            Show that for all $b\in A$ we have $b\in\initialSegment{a}$.
            \begin{subproof}
                Fix $b$.
                Assume $b\in A$.
                Then $b \leq a$.
                Thus $b\in\naturalsPlus$ by \cref{subseteq}.
                Thus $b\in\initialSegment{a}$ by \cref{initial_segment_naturalsplus}.
            \end{subproof}
            Thus $A\subseteq\initialSegment{a}$ by \cref{subseteq}.
        \end{subproof}
        $\initialSegment{a}$ is finite by \cref{initial_segment_finite}.
        Thus $A$ is finite by \cref{subseteq_finite_is_finite}.
        Contradiction.
    \end{subproof}
    Let $R = \{w\in A\times A \mid \fst{w} < \snd{w}\}$.
    Show that $R$ is a relation.
    \begin{subproof}
        Show that for all $w$ such that $w\in R$ there exist $x,y$ such that $w = (x,y)$.
        \begin{subproof}
            Fix $w$.
            Assume $w\in R$.
            Then $w\in A\times A$ and $\fst{w} < \snd{w}$.
            Take $x,y$ such that $w = (x,y)$ by \cref{times_elem_is_tuple}.
        \end{subproof}
        Thus $R$ is a relation by \cref{relation}.
    \end{subproof}
    Show that for all $x\in A$ there exists $y$ such that $x\mathrel{R} y$.
    \begin{subproof}
        Fix $x\in A$.
        Take $y\in A$ such that $x < y$.
        Thus $(x,y)\in A\times A$ by \cref{times_tuple_intro}.
        Thus $(x,y)\in R$.
        Thus $x\mathrel{R} y$.
    \end{subproof}
    Take $f$ such that $f$ is a function on $A$ and for all $x\in A$ we have
        $x\mathrel{R}\apply{f}{x}$ by \cref{choice_function}.
    Show that for all $x\in A$ we have $\apply{f}{x}\in A$ and $x < \apply{f}{x}$.
    \begin{subproof}
        Fix $x\in A$.
        Thus $x\mathrel{R}\apply{f}{x}$.
        Thus $(x,\apply{f}{x})\in A\times A$ and $x < \apply{f}{x}$.
        Thus $\apply{f}{x}\in A$ by \cref{times_tuple_elim}.
    \end{subproof}
    Let $S = \iterationSet{f}{a_0}$.
    Show that for all $n\in\naturalsPlus$ we have $n\in S$.
    \begin{subproof}
        Show that $S\subseteq\naturalsPlus$.
        \begin{subproof}
            Show that for all $n$ such that $n\in S$ we have $n\in\naturalsPlus$.
            \begin{subproof}
                Fix $n$.
                Assume $n\in S$.
                Thus $n\in\naturalsPlus$ by \cref{iteration_set_elim}.
            \end{subproof}
            Thus $S\subseteq\naturalsPlus$ by \cref{subseteq}.
        \end{subproof}
        $1\in\naturalsPlus$ by \cref{real_one_in_naturals}.
        Show that $1\in S$.
        \begin{subproof}
            Let $u = \{(1,a_0)\}$.
            Show that $u$ is a function.
            \begin{subproof}
                Show that $u$ is a relation.
                \begin{subproof}
                    Show that for all $w$ such that $w\in u$ there exist $x,y$ such that $w = (x,y)$.
                    \begin{subproof}
                        Fix $w$.
                        Assume $w\in u$.
                        Then $w = (1,a_0)$.
                    \end{subproof}
                    Thus $u$ is a relation by \cref{relation}.
                \end{subproof}
                Show that $u$ is right-unique.
                \begin{subproof}
                    Show that for all $x,y,y'$ such that $(x,y)\in u$ and $(x,y')\in u$ we have $y = y'$.
                    \begin{subproof}
                        Fix $x$.
                        Fix $y$.
                        Fix $y'$.
                        Assume $(x,y)\in u$ and $(x,y')\in u$.
                        Then $(x,y) = (1,a_0)$ and $(x,y') = (1,a_0)$.
                        Thus $y = y'$.
                    \end{subproof}
                    Thus $u$ is right-unique by \cref{rightunique}.
                \end{subproof}
                Thus $u$ is a function.
            \end{subproof}
            Show that $\dom{u} = \initialSegment{1}$.
            \begin{subproof}
                Show that $\dom{u}\subseteq\initialSegment{1}$.
                \begin{subproof}
                    Show that for all $x$ such that $x\in\dom{u}$ we have $x\in\initialSegment{1}$.
                    \begin{subproof}
                        Fix $x$.
                        Assume $x\in\dom{u}$.
                        Take $y$ such that $(x,y)\in u$ by \cref{dom}.
                        Then $(x,y) = (1,a_0)$.
                        Thus $x = 1$.
                        $1\in\naturalsPlus$ by \cref{real_one_in_naturals}.
                        Thus $x\in\initialSegment{1}$ by \cref{initial_segment_naturalsplus}.
                    \end{subproof}
                    Thus $\dom{u}\subseteq\initialSegment{1}$ by \cref{subseteq}.
                \end{subproof}
                Show that $\initialSegment{1}\subseteq\dom{u}$.
                \begin{subproof}
                    Show that for all $x$ such that $x\in\initialSegment{1}$ we have $x\in\dom{u}$.
                    \begin{subproof}
                        Fix $x$.
                        Assume $x\in\initialSegment{1}$.
                        Thus $x\in\naturalsPlus$ and $x \leq 1$ by \cref{initial_segment_naturalsplus}.
                        $1\in\reals$ by \cref{real_naturals_subset_reals,real_one_in_naturals,subseteq}.
                        $x\in\reals$ by \cref{real_naturals_subset_reals,subseteq}.
                        $1 < x$ or $1 = x$ by \cref{real_one_leq_naturalsplus}.
                        \begin{byCase}
                        \caseOf{$1 = x$.}
                            Thus $(x,a_0) = (1,a_0)$.
                            Thus $(x,a_0)\in u$ by \cref{singleton_intro}.
                            Thus $x\in\dom{u}$ by \cref{dom}.
                        \caseOf{$1 < x$.}
                            Suppose not.
                            Show that $1 < 1$.
                            \begin{subproof}
                                $x < 1$ or $x = 1$ by \cref{real_leq_split}.
                                \begin{byCase}
                                \caseOf{$x < 1$.}
                                    Thus $1 < 1$ by \cref{real_lt_trans}.
                                \caseOf{$x = 1$.}
                                    Thus $1 < 1$ by \cref{real_lt_subst_right}.
                                \end{byCase}
                            \end{subproof}
                            Contradiction by \cref{real_lt_irreflexive}.
                        \end{byCase}
                    \end{subproof}
                    Thus $\initialSegment{1}\subseteq\dom{u}$ by \cref{subseteq}.
                \end{subproof}
                Thus $\dom{u} = \initialSegment{1}$ by \cref{subseteq_antisymmetric}.
            \end{subproof}
            Thus $(1,a_0)\in u$ by \cref{singleton_intro}.
            Thus $u(1) = a_0$ by \cref{function_apply_intro}.
            Show that for all $k\in\naturalsPlus$ such that $k + 1 \leq 1$ we have $u(k + 1) = \apply{f}{u(k)}$.
            \begin{subproof}
                Fix $k$.
                Assume $k\in\naturalsPlus$.
                Assume $k + 1 \leq 1$.
                Suppose not.
                $k\in\reals$ by \cref{real_naturals_subset_reals,subseteq}.
                $1\in\reals$ by \cref{real_naturals_subset_reals,real_one_in_naturals,subseteq}.
                $k + 1\in\reals$ by \cref{real_add_closed}.
                Show that $k < k + 1$.
                \begin{subproof}
                    $0\in\reals$ by \cref{real_add_ident}.
                    $0 < 1$ by \cref{real_zero_lt_one}.
                    $0 + k < 1 + k$ by \cref{real_lt_add}.
                    $0 + k = k + 0$ by \cref{real_add_comm}.
                    $1 + k = k + 1$ by \cref{real_add_comm}.
                    Thus $k + 0 < k + 1$ by \cref{real_lt_subst_left}.
                    $k + 0 = k$ by \cref{real_add_ident}.
                    Thus $k < k + 1$ by \cref{real_lt_subst_left}.
                \end{subproof}
                Show that $1 < k + 1$.
                \begin{subproof}
                    $1 < k$ or $1 = k$ by \cref{real_one_leq_naturalsplus}.
                    \begin{byCase}
                    \caseOf{$1 = k$.}
                        Thus $1 < k + 1$ by \cref{real_lt_subst_left}.
                    \caseOf{$1 < k$.}
                        Thus $1 < k + 1$ by \cref{real_lt_trans}.
                    \end{byCase}
                \end{subproof}
                $k + 1 < 1$ or $k + 1 = 1$ by \cref{real_leq_split}.
                Show that $1 < 1$.
                \begin{subproof}
                    \begin{byCase}
                    \caseOf{$k + 1 < 1$.}
                        Thus $1 < 1$ by \cref{real_lt_trans}.
                    \caseOf{$k + 1 = 1$.}
                        Thus $1 < 1$ by \cref{real_lt_subst_right}.
                    \end{byCase}
                \end{subproof}
                Contradiction by \cref{real_lt_irreflexive}.
            \end{subproof}
            Thus $1\in\naturalsPlus$.
            Thus $u$ is an iteration sequence for $f$ on $1$ starting at $a_0$ by \cref{iteration_sequence}.
            Thus there exists $u$ such that
                $u$ is an iteration sequence for $f$ on $1$ starting at $a_0$.
            Thus $1\in S$ by \cref{iteration_set_intro}.
        \end{subproof}
        Show that for all $n\in S$ we have $n + 1\in S$.
        \begin{subproof}
            Fix $n$.
            Assume $n\in S$.
            Thus $n\in\naturalsPlus$ and there exists $u$ such that
                $u$ is an iteration sequence for $f$ on $n$ starting at $a_0$
                by \cref{iteration_set_elim}.
            Take $u$ such that
                $u$ is an iteration sequence for $f$ on $n$ starting at $a_0$.
            Thus $n + 1\in\naturalsPlus$ by \cref{real_naturals_closed_succ}.
            Let $v = u\union\{(n + 1,\apply{f}{u(n)})\}$.
            Show that $v$ is a function.
            \begin{subproof}
                Show that $v$ is a relation.
                \begin{subproof}
                    Thus $u$ is a relation.
                    Show that $\{(n + 1,\apply{f}{u(n)})\}$ is a relation.
                    \begin{subproof}
                        Show that for all $w$ such that $w\in\{(n + 1,\apply{f}{u(n)})\}$ there exist $x,y$ such that $w = (x,y)$.
                        \begin{subproof}
                            Fix $w$.
                            Assume $w\in\{(n + 1,\apply{f}{u(n)})\}$.
                            Then $w = (n + 1,\apply{f}{u(n)})$.
                        \end{subproof}
                        Thus $\{(n + 1,\apply{f}{u(n)})\}$ is a relation by \cref{relation}.
                    \end{subproof}
                    Thus $v$ is a relation by \cref{union_relations_is_relation}.
                \end{subproof}
                Show that $v$ is right-unique.
                \begin{subproof}
                    Show that for all $x,y,y'$ such that $(x,y)\in v$ and $(x,y')\in v$ we have $y = y'$.
                    \begin{subproof}
                        Fix $x$.
                        Fix $y$.
                        Fix $y'$.
                        Assume $(x,y)\in v$ and $(x,y')\in v$.
                        \begin{byCase}
                        \caseOf{$x = n + 1$.}
                            $(n + 1,y)\in v$.
                            $(n + 1,y')\in v$.
                            Show that $(n + 1,y)\in\{(n + 1,\apply{f}{u(n)})\}$.
                            \begin{subproof}
                                Suppose not.
                                Thus $(n + 1,y)\in u$ by \cref{union_iff}.
                                Thus $n + 1\in\dom{u}$ by \cref{dom}.
                                Thus $\dom{u} = \initialSegment{n}$ by \cref{iteration_sequence}.
                                Thus $n + 1\in\initialSegment{n}$.
                                Thus $n + 1 \leq n$ by \cref{initial_segment_naturalsplus}.
                                Show that $n < n + 1$.
                                \begin{subproof}
                                    $n\in\reals$ by \cref{real_naturals_subset_reals,subseteq}.
                                    $0\in\reals$ by \cref{real_add_ident}.
                                    $1\in\reals$ by \cref{real_naturals_subset_reals,real_one_in_naturals,subseteq}.
                                    $0 < 1$ by \cref{real_zero_lt_one}.
                                    $0 + n < 1 + n$ by \cref{real_lt_add}.
                                    $0 + n = n + 0$ by \cref{real_add_comm}.
                                    $1 + n = n + 1$ by \cref{real_add_comm}.
                                    Thus $n + 0 < n + 1$ by \cref{real_lt_subst_left}.
                                    $n + 0 = n$ by \cref{real_add_ident}.
                                    Thus $n < n + 1$ by \cref{real_lt_subst_left}.
                                \end{subproof}
                                $n + 1 < n$ or $n + 1 = n$ by \cref{real_leq_split}.
                                \begin{byCase}
                                \caseOf{$n + 1 < n$.}
                                    $n\in\reals$ by \cref{real_naturals_subset_reals,subseteq}.
                                    $1\in\reals$ by \cref{real_naturals_subset_reals,real_one_in_naturals,subseteq}.
                                    $n + 1\in\reals$ by \cref{real_add_closed}.
                                    Thus $n < n$ by \cref{real_lt_trans}.
                                    Contradiction by \cref{real_lt_irreflexive}.
                                \caseOf{$n + 1 = n$.}
                                    $n\in\reals$ by \cref{real_naturals_subset_reals,subseteq}.
                                    Thus $n < n$ by \cref{real_lt_subst_right}.
                                    Contradiction by \cref{real_lt_irreflexive}.
                                \end{byCase}
                            \end{subproof}
                            Thus $(n + 1,y) = (n + 1,\apply{f}{u(n)})$ by \cref{singleton_iff}.
                            Thus $y = \apply{f}{u(n)}$ by \cref{pair_eq_iff}.
                            Show that $(n + 1,y')\in\{(n + 1,\apply{f}{u(n)})\}$.
                            \begin{subproof}
                                Suppose not.
                                Thus $(n + 1,y')\in u$ by \cref{union_iff}.
                                Thus $n + 1\in\dom{u}$ by \cref{dom}.
                                Thus $\dom{u} = \initialSegment{n}$ by \cref{iteration_sequence}.
                                Thus $n + 1\in\initialSegment{n}$.
                                Thus $n + 1 \leq n$ by \cref{initial_segment_naturalsplus}.
                                Show that $n < n + 1$.
                                \begin{subproof}
                                    $n\in\reals$ by \cref{real_naturals_subset_reals,subseteq}.
                                    $0\in\reals$ by \cref{real_add_ident}.
                                    $1\in\reals$ by \cref{real_naturals_subset_reals,real_one_in_naturals,subseteq}.
                                    $0 < 1$ by \cref{real_zero_lt_one}.
                                    $0 + n < 1 + n$ by \cref{real_lt_add}.
                                    $0 + n = n + 0$ by \cref{real_add_comm}.
                                    $1 + n = n + 1$ by \cref{real_add_comm}.
                                    Thus $n + 0 < n + 1$ by \cref{real_lt_subst_left}.
                                    $n + 0 = n$ by \cref{real_add_ident}.
                                    Thus $n < n + 1$ by \cref{real_lt_subst_left}.
                                \end{subproof}
                                $n + 1 < n$ or $n + 1 = n$ by \cref{real_leq_split}.
                                \begin{byCase}
                                \caseOf{$n + 1 < n$.}
                                    $n\in\reals$ by \cref{real_naturals_subset_reals,subseteq}.
                                    $1\in\reals$ by \cref{real_naturals_subset_reals,real_one_in_naturals,subseteq}.
                                    $n + 1\in\reals$ by \cref{real_add_closed}.
                                    Thus $n < n$ by \cref{real_lt_trans}.
                                    Contradiction by \cref{real_lt_irreflexive}.
                                \caseOf{$n + 1 = n$.}
                                    $n\in\reals$ by \cref{real_naturals_subset_reals,subseteq}.
                                    Thus $n < n$ by \cref{real_lt_subst_right}.
                                    Contradiction by \cref{real_lt_irreflexive}.
                                \end{byCase}
                            \end{subproof}
                            Thus $(n + 1,y') = (n + 1,\apply{f}{u(n)})$ by \cref{singleton_iff}.
                            Thus $y' = \apply{f}{u(n)}$ by \cref{pair_eq_iff}.
                            Thus $y = y'$.
                        \caseOf{$x \neq n + 1$.}
                            Show that $(x,y)\in u$.
                            \begin{subproof}
                                Suppose not.
                                Thus $(x,y)\in\{(n + 1,\apply{f}{u(n)})\}$ by \cref{union_iff}.
                                Thus $(x,y) = (n + 1,\apply{f}{u(n)})$ by \cref{singleton_iff}.
                                Thus $x = n + 1$ by \cref{pair_eq_iff}.
                                Contradiction.
                            \end{subproof}
                            Show that $(x,y')\in u$.
                            \begin{subproof}
                                Suppose not.
                                Thus $(x,y')\in\{(n + 1,\apply{f}{u(n)})\}$ by \cref{union_iff}.
                                Thus $(x,y') = (n + 1,\apply{f}{u(n)})$ by \cref{singleton_iff}.
                                Thus $x = n + 1$ by \cref{pair_eq_iff}.
                                Contradiction.
                            \end{subproof}
                            Thus $u$ is a function by \cref{iteration_sequence}.
                            Thus $y = y'$ by \cref{function_rightunique}.
                        \end{byCase}
                    \end{subproof}
                    Thus $v$ is right-unique by \cref{rightunique}.
                \end{subproof}
                Thus $v$ is a function.
            \end{subproof}
            Show that $\dom{v} = \initialSegment{n + 1}$.
            \begin{subproof}
                Show that $\dom{v}\subseteq\initialSegment{n + 1}$.
                \begin{subproof}
                    Show that for all $x$ such that $x\in\dom{v}$ we have $x\in\initialSegment{n + 1}$.
                    \begin{subproof}
                        Fix $x$.
                        Assume $x\in\dom{v}$.
                        Take $y$ such that $(x,y)\in v$ by \cref{dom}.
                        Thus $(x,y)\in u$ or $(x,y)\in\{(n + 1,\apply{f}{u(n)})\}$ by \cref{union_iff}.
                        \begin{byCase}
                        \caseOf{$(x,y)\in u$.}
                            Thus $x\in\dom{u}$ by \cref{dom}.
                            Thus $\dom{u} = \initialSegment{n}$ by \cref{iteration_sequence}.
                            Thus $x\in\initialSegment{n}$.
                            Thus $x\in\naturalsPlus$ and $x \leq n$ by \cref{initial_segment_naturalsplus}.
                            Show that $x \leq n + 1$.
                            \begin{subproof}
                                $x < n$ or $x = n$ by \cref{real_leq_split}.
                                \begin{byCase}
                                \caseOf{$x < n$.}
                                    Show that $n < n + 1$.
                                    \begin{subproof}
                                        $n\in\reals$ by \cref{real_naturals_subset_reals,subseteq}.
                                        $0\in\reals$ by \cref{real_add_ident}.
                                        $1\in\reals$ by \cref{real_naturals_subset_reals,real_one_in_naturals,subseteq}.
                                        $0 < 1$ by \cref{real_zero_lt_one}.
                                        $0 + n < 1 + n$ by \cref{real_lt_add}.
                                        $0 + n = n + 0$ by \cref{real_add_comm}.
                                        $1 + n = n + 1$ by \cref{real_add_comm}.
                                        Thus $n + 0 < n + 1$ by \cref{real_lt_subst_left}.
                                        $n + 0 = n$ by \cref{real_add_ident}.
                                        Thus $n < n + 1$ by \cref{real_lt_subst_left}.
                                    \end{subproof}
                                    $x\in\reals$ by \cref{real_naturals_subset_reals,subseteq}.
                                    $n\in\reals$ by \cref{real_naturals_subset_reals,subseteq}.
                                    $1\in\reals$ by \cref{real_naturals_subset_reals,real_one_in_naturals,subseteq}.
                                    $n + 1\in\reals$ by \cref{real_add_closed}.
                                    Thus $x < n + 1$ by \cref{real_lt_trans}.
                                    Thus $x \leq n + 1$.
                                \caseOf{$x = n$.}
                                    Show that $n < n + 1$.
                                    \begin{subproof}
                                        $n\in\reals$ by \cref{real_naturals_subset_reals,subseteq}.
                                        $0\in\reals$ by \cref{real_add_ident}.
                                        $1\in\reals$ by \cref{real_naturals_subset_reals,real_one_in_naturals,subseteq}.
                                        $0 < 1$ by \cref{real_zero_lt_one}.
                                        $0 + n < 1 + n$ by \cref{real_lt_add}.
                                        $0 + n = n + 0$ by \cref{real_add_comm}.
                                        $1 + n = n + 1$ by \cref{real_add_comm}.
                                        Thus $n + 0 < n + 1$ by \cref{real_lt_subst_left}.
                                        $n + 0 = n$ by \cref{real_add_ident}.
                                        Thus $n < n + 1$ by \cref{real_lt_subst_left}.
                                    \end{subproof}
                                    Thus $x < n + 1$ by \cref{real_lt_subst_left}.
                                    Thus $x \leq n + 1$.
                                \end{byCase}
                            \end{subproof}
                            Thus $x\in\initialSegment{n + 1}$ by \cref{initial_segment_naturalsplus}.
                        \caseOf{$(x,y)\in\{(n + 1,\apply{f}{u(n)})\}$.}
                            Thus $(x,y) = (n + 1,\apply{f}{u(n)})$ by \cref{singleton_iff}.
                            Thus $x = n + 1$ by \cref{pair_eq_iff}.
                            Thus $x\in\initialSegment{n + 1}$ by \cref{initial_segment_naturalsplus}.
                        \end{byCase}
                    \end{subproof}
                    Thus $\dom{v}\subseteq\initialSegment{n + 1}$ by \cref{subseteq}.
                \end{subproof}
                Show that $\initialSegment{n + 1}\subseteq\dom{v}$.
                \begin{subproof}
                    Show that for all $x$ such that $x\in\initialSegment{n + 1}$ we have $x\in\dom{v}$.
                    \begin{subproof}
                        Fix $x$.
                        Assume $x\in\initialSegment{n + 1}$.
                        Thus $x\in\naturalsPlus$ and $x \leq n + 1$ by \cref{initial_segment_naturalsplus}.
                        \begin{byCase}
                        \caseOf{$x = n + 1$.}
                            Thus $(x,\apply{f}{u(n)}) = (n + 1,\apply{f}{u(n)})$.
                            Thus $(x,\apply{f}{u(n)})\in\{(n + 1,\apply{f}{u(n)})\}$ by \cref{cons_left}.
                            Thus $(x,\apply{f}{u(n)})\in v$ by \cref{union_intro_right}.
                            Thus $x\in\dom{v}$ by \cref{dom}.
                        \caseOf{$x \neq n + 1$.}
                            Thus $x \leq n$ by \cref{real_naturals_leq_succ_elim}.
                            Thus $x\in\initialSegment{n}$ by \cref{initial_segment_naturalsplus}.
                            Thus $x\in\dom{u}$ by \cref{iteration_sequence}.
                            Thus $x\in\dom{v}$ by \cref{dom,union_iff}.
                        \end{byCase}
                    \end{subproof}
                    Thus $\initialSegment{n + 1}\subseteq\dom{v}$ by \cref{subseteq}.
                \end{subproof}
                Thus $\dom{v} = \initialSegment{n + 1}$ by \cref{subseteq_antisymmetric}.
            \end{subproof}
            Show that $v(1) = a_0$.
            \begin{subproof}
                Thus $u$ is a function by \cref{iteration_sequence}.
                Thus $1\in\naturalsPlus$ by \cref{real_one_in_naturals}.
                Thus $1 < n$ or $1 = n$ by \cref{real_one_leq_naturalsplus}.
                Thus $1 \leq n$.
                Thus $1\in\initialSegment{n}$ by \cref{initial_segment_naturalsplus}.
                Thus $1\in\dom{u}$ by \cref{iteration_sequence}.
                Thus $(1,u(1))\in u$ by \cref{function_apply_elim}.
                Thus $(1,u(1))\in v$ by \cref{union_intro_left}.
                Thus $v(1) = u(1)$ by \cref{function_apply_intro}.
                Thus $u(1) = a_0$ by \cref{iteration_sequence}.
                Thus $v(1) = a_0$.
            \end{subproof}
            Show that for all $k\in\naturalsPlus$ such that $k + 1 \leq n + 1$ we have $v(k + 1) = \apply{f}{v(k)}$.
            \begin{subproof}
                Fix $k$.
                Assume $k\in\naturalsPlus$.
                Assume $k + 1 \leq n + 1$.
                \begin{byCase}
                \caseOf{$k + 1 = n + 1$.}
                    Thus $k = n$.
                    Thus $(k + 1,\apply{f}{u(n)})\in\{(n + 1,\apply{f}{u(n)})\}$ by \cref{cons_left}.
                    Thus $(k + 1,\apply{f}{u(n)})\in v$ by \cref{union_intro_right}.
                    Thus $v(k + 1) = \apply{f}{u(n)}$ by \cref{function_apply_intro}.
                    Thus $v(k + 1) = \apply{f}{u(k)}$.
                    Thus $k \leq n$.
                    Thus $k\in\initialSegment{n}$ by \cref{initial_segment_naturalsplus}.
                    Thus $k\in\dom{u}$ by \cref{iteration_sequence}.
                    Thus $u$ is a function by \cref{iteration_sequence}.
                    Thus $(k,u(k))\in u$ by \cref{function_apply_elim}.
                    Thus $(k,u(k))\in v$ by \cref{union_intro_left}.
                    Thus $v(k) = u(k)$ by \cref{function_apply_intro}.
                    Thus $v(k + 1) = \apply{f}{v(k)}$.
                \caseOf{$k + 1 \neq n + 1$.}
                    Thus $k + 1\in\naturalsPlus$ by \cref{real_naturals_closed_succ}.
                    Thus $k + 1 \leq n$ by \cref{real_naturals_leq_succ_elim}.
                    Thus $k + 1\in\initialSegment{n}$ by \cref{initial_segment_naturalsplus}.
                    Thus $k + 1\in\dom{u}$ by \cref{iteration_sequence}.
                    Thus $u$ is a function by \cref{iteration_sequence}.
                    Thus $(k + 1,u(k + 1))\in u$ by \cref{function_apply_elim}.
                    Thus $(k + 1,u(k + 1))\in v$ by \cref{union_intro_left}.
                    Thus $k + 1\in\initialSegment{n + 1}$ by \cref{initial_segment_naturalsplus}.
                    Thus $k + 1\in\dom{v}$ by \cref{dom}.
                    Thus $(k + 1,v(k + 1))\in v$ by \cref{function_apply_elim}.
                    Thus $u(k + 1) = v(k + 1)$ by \cref{function_rightunique}.
                    Thus $v(k + 1) = u(k + 1)$.
                    Thus $u(k + 1) = \apply{f}{u(k)}$.
                    Thus $v(k + 1) = \apply{f}{u(k)}$.
                    Show that $k \leq n$.
                    \begin{subproof}
                        Show that $k < k + 1$.
                        \begin{subproof}
                            $k\in\reals$ by \cref{real_naturals_subset_reals,subseteq}.
                            $0\in\reals$ by \cref{real_add_ident}.
                            $1\in\reals$ by \cref{real_naturals_subset_reals,real_one_in_naturals,subseteq}.
                            $0 < 1$ by \cref{real_zero_lt_one}.
                            $0 + k < 1 + k$ by \cref{real_lt_add}.
                            $0 + k = k + 0$ by \cref{real_add_comm}.
                            $1 + k = k + 1$ by \cref{real_add_comm}.
                            Thus $k + 0 < k + 1$ by \cref{real_lt_subst_left}.
                            $k + 0 = k$ by \cref{real_add_ident}.
                            Thus $k < k + 1$ by \cref{real_lt_subst_left}.
                        \end{subproof}
                        Thus $k < k + 1$ or $k = k + 1$.
                        Thus $k \leq k + 1$.
                        $k\in\reals$ by \cref{real_naturals_subset_reals,subseteq}.
                        $k + 1\in\reals$ by \cref{real_naturals_subset_reals,subseteq}.
                        $n\in\reals$ by \cref{real_naturals_subset_reals,subseteq}.
                        Thus $k \leq n$ by \cref{real_leq_trans}.
                    \end{subproof}
                    Thus $k\in\initialSegment{n}$ by \cref{initial_segment_naturalsplus}.
                    Thus $k\in\dom{u}$ by \cref{iteration_sequence}.
                    Thus $(k,u(k))\in u$ by \cref{function_apply_elim}.
                    Thus $(k,u(k))\in v$ by \cref{union_intro_left}.
                    Thus $k\in\dom{v}$ by \cref{dom}.
                    Thus $(k,v(k))\in v$ by \cref{function_apply_elim}.
                    Thus $u(k) = v(k)$ by \cref{function_rightunique}.
                    Thus $v(k + 1) = \apply{f}{v(k)}$.
                \end{byCase}
            \end{subproof}
            Thus $v$ is an iteration sequence for $f$ on $n + 1$ starting at $a_0$ by \cref{iteration_sequence_intro}.
            Thus $n + 1\in\naturalsPlus$ by \cref{real_naturals_closed_succ}.
            Thus $n + 1\in S$ by \cref{iteration_set_intro}.
        \end{subproof}
        Thus for all $n\in\naturalsPlus$ we have $n\in S$ by \cref{real_naturals_induction}.
    \end{subproof}
    Let $F = \{u \mid \text{there exists $n\in\naturalsPlus$ such that
        $u$ is an iteration sequence for $f$ on $n$ starting at $a_0$} \}$.
    Show that for all $u\in F$ we have $u$ is a function.
    \begin{subproof}
        Fix $u$.
        Assume $u\in F$.
        Take $n\in\naturalsPlus$ such that
            $u$ is an iteration sequence for $f$ on $n$ starting at $a_0$.
        Thus $u$ is a function by \cref{iteration_sequence}.
    \end{subproof}
    Show that for all $u,v$ such that $u\in F$ and $v\in F$ we have $u\subseteq v$ or $v\subseteq u$.
    \begin{subproof}
        Fix $u$.
        Fix $v$.
        Assume $u\in F$ and $v\in F$.
        Take $m\in\naturalsPlus$ such that
            $u$ is an iteration sequence for $f$ on $m$ starting at $a_0$.
        Take $n\in\naturalsPlus$ such that
            $v$ is an iteration sequence for $f$ on $n$ starting at $a_0$.
        Thus $u$ is a function by \cref{iteration_sequence}.
        Thus $\dom{u} = \initialSegment{m}$ by \cref{iteration_sequence}.
        Thus $u(1) = a_0$ by \cref{iteration_sequence}.
        Thus for all $k\in\naturalsPlus$ such that $k + 1 \leq m$ we have $u(k + 1) = \apply{f}{u(k)}$
            by \cref{iteration_sequence}.
        Thus $v$ is a function by \cref{iteration_sequence}.
        Thus $\dom{v} = \initialSegment{n}$ by \cref{iteration_sequence}.
        Thus $v(1) = a_0$ by \cref{iteration_sequence}.
        Thus for all $k\in\naturalsPlus$ such that $k + 1 \leq n$ we have $v(k + 1) = \apply{f}{v(k)}$
            by \cref{iteration_sequence}.
        $m = n$ or $m \neq n$.
        \begin{byCase}
        \caseOf{$m = n$.}
            Thus $u = v$ by \cref{iteration_sequence_unique}.
            Thus $u\subseteq v$ or $v\subseteq u$ by \cref{subseteq_refl}.
        \caseOf{$m \neq n$.}
            $m\in\reals$ by \cref{real_naturals_subset_reals,subseteq}.
            $n\in\reals$ by \cref{real_naturals_subset_reals,subseteq}.
            Thus $m < n$ or $n < m$ by \cref{real_lt_trichotomy}.
            \begin{byCase}
        \caseOf{$m < n$.}
            Show that for all $k\in\initialSegment{m}$ we have $u(k) = v(k)$.
            \begin{subproof}
                Let $A = \{k\in\naturalsPlus \mid \text{if $k\leq m$ then $u(k) = v(k)$}\}$.
                Show that for all $k\in\naturalsPlus$ we have $k\in A$.
                \begin{subproof}
                    Show that $A\subseteq\naturalsPlus$.
                    \begin{subproof}
                        Show that for all $k$ such that $k\in A$ we have $k\in\naturalsPlus$.
                        \begin{subproof}
                            Fix $k$.
                            Assume $k\in A$.
                            Then $k\in\naturalsPlus$.
                        \end{subproof}
                        Thus $A\subseteq\naturalsPlus$ by \cref{subseteq}.
                    \end{subproof}
                    $1\in\naturalsPlus$ by \cref{real_one_in_naturals}.
                    Show that $1\in A$.
                    \begin{subproof}
                        Show that if $1\leq m$ then $u(1) = v(1)$.
                        \begin{subproof}
                            Assume $1\leq m$.
                            Then $u(1) = a_0$ and $v(1) = a_0$.
                            Thus $u(1) = v(1)$.
                        \end{subproof}
                        Thus $1\in A$.
                    \end{subproof}
                    Show that for all $k\in A$ we have $k + 1\in A$.
                    \begin{subproof}
                        Fix $k$.
                        Assume $k\in A$.
                        Show that if $k + 1\leq m$ then $u(k + 1) = v(k + 1)$.
                        \begin{subproof}
                            Assume $k + 1\leq m$.
                            Then $k \leq m$.
                            Thus $u(k) = v(k)$.
                            Thus $u(k + 1) = \apply{f}{u(k)}$ and $v(k + 1) = \apply{f}{v(k)}$.
                            Thus $u(k + 1) = v(k + 1)$.
                        \end{subproof}
                        Thus $k + 1\in A$.
                    \end{subproof}
                    Thus for all $k\in\naturalsPlus$ we have $k\in A$ by \cref{real_naturals_induction}.
                \end{subproof}
                Show that for all $k\in\initialSegment{m}$ we have $u(k) = v(k)$.
                \begin{subproof}
                    Fix $k$.
                    Assume $k\in\initialSegment{m}$.
                    Thus $k\in\naturalsPlus$ and $k \leq m$ by \cref{initial_segment_naturalsplus}.
                    Thus $k\in A$.
                    Thus $u(k) = v(k)$.
                \end{subproof}
            \end{subproof}
            Show that $u\subseteq v$.
            \begin{subproof}
                Show that for all $w$ such that $w\in u$ we have $w\in v$.
                \begin{subproof}
                    Fix $w$.
                    Assume $w\in u$.
                    Thus $u$ is a function by \cref{iteration_sequence}.
                    Take $k\in\dom{u}$ such that $w = (k,u(k))$ by \cref{function_member_elim}.
                    Thus $\dom{u} = \initialSegment{m}$ by \cref{iteration_sequence}.
                    Thus $k\in\initialSegment{m}$.
                    Thus $u(k) = v(k)$.
                    Thus $w = (k,v(k))$.
                    Thus $k\in\naturalsPlus$ and $k \leq m$ by \cref{initial_segment_naturalsplus}.
                    Thus $m \leq n$.
                    $k\in\reals$ by \cref{real_naturals_subset_reals,subseteq}.
                    $m\in\reals$ by \cref{real_naturals_subset_reals,subseteq}.
                    $n\in\reals$ by \cref{real_naturals_subset_reals,subseteq}.
                    Thus $k \leq n$ by \cref{real_leq_trans}.
                    Thus $k\in\initialSegment{n}$ by \cref{initial_segment_naturalsplus}.
                    Thus $\dom{v} = \initialSegment{n}$ by \cref{iteration_sequence}.
                    Thus $k\in\dom{v}$.
                    Thus $v$ is a function by \cref{iteration_sequence}.
                    Thus $(k,v(k))\in v$ by \cref{function_apply_elim}.
                    Thus $w\in v$.
                \end{subproof}
                Thus $u\subseteq v$ by \cref{subseteq}.
            \end{subproof}
        \caseOf{$n < m$.}
            Show that for all $k\in\initialSegment{n}$ we have $v(k) = u(k)$.
            \begin{subproof}
                Let $A = \{k\in\naturalsPlus \mid \text{if $k\leq n$ then $v(k) = u(k)$}\}$.
                Show that for all $k\in\naturalsPlus$ we have $k\in A$.
                \begin{subproof}
                    Show that $A\subseteq\naturalsPlus$.
                    \begin{subproof}
                        Show that for all $k$ such that $k\in A$ we have $k\in\naturalsPlus$.
                        \begin{subproof}
                            Fix $k$.
                            Assume $k\in A$.
                            Then $k\in\naturalsPlus$.
                        \end{subproof}
                        Thus $A\subseteq\naturalsPlus$ by \cref{subseteq}.
                    \end{subproof}
                    $1\in\naturalsPlus$ by \cref{real_one_in_naturals}.
                    Show that $1\in A$.
                    \begin{subproof}
                        Show that if $1\leq n$ then $v(1) = u(1)$.
                        \begin{subproof}
                            Assume $1\leq n$.
                            Then $v(1) = a_0$ and $u(1) = a_0$.
                            Thus $v(1) = u(1)$.
                        \end{subproof}
                        Thus $1\in A$.
                    \end{subproof}
                    Show that for all $k\in A$ we have $k + 1\in A$.
                    \begin{subproof}
                        Fix $k$.
                        Assume $k\in A$.
                        Show that if $k + 1\leq n$ then $v(k + 1) = u(k + 1)$.
                        \begin{subproof}
                            Assume $k + 1\leq n$.
                            $k\in\reals$ and $n\in\reals$ by \cref{real_naturals_subset_reals,subseteq}.
                            $0\in\reals$ by \cref{real_add_ident}.
                            $1\in\reals$ by \cref{real_naturals_subset_reals,real_one_in_naturals,subseteq}.
                            $0 < 1$ by \cref{real_zero_lt_one}.
                            $0 + k < 1 + k$ by \cref{real_lt_add}.
                            $0 + k = k + 0$ by \cref{real_add_comm}.
                            $1 + k = k + 1$ by \cref{real_add_comm}.
                            Thus $k + 0 < k + 1$ by \cref{real_lt_subst_left}.
                            $k + 0 = k$ by \cref{real_add_ident}.
                            Thus $k < k + 1$ by \cref{real_lt_subst_left}.
                            Thus $k < k + 1$ or $k = k + 1$.
                            Thus $k \leq k + 1$.
                            $k + 1\in\reals$ by \cref{real_add_closed}.
                            Thus $k \leq n$ by \cref{real_leq_trans}.
                            Thus $v(k) = u(k)$.
                            Thus $k\in\naturalsPlus$.
                            Thus $v(k + 1) = \apply{f}{v(k)}$ by \cref{iteration_sequence}.
                            Thus $n < m$ or $n = m$.
                            Thus $n \leq m$.
                            $k + 1\in\reals$ by \cref{real_add_closed}.
                            $n\in\reals$ and $m\in\reals$ by \cref{real_naturals_subset_reals,subseteq}.
                            Thus $k + 1 \leq m$ by \cref{real_leq_trans}.
                            Thus $u(k + 1) = \apply{f}{u(k)}$ by \cref{iteration_sequence}.
                            Thus $v(k + 1) = u(k + 1)$.
                        \end{subproof}
                        Thus $k + 1\in A$.
                    \end{subproof}
                    Thus for all $k\in\naturalsPlus$ we have $k\in A$ by \cref{real_naturals_induction}.
                \end{subproof}
                Show that for all $k\in\initialSegment{n}$ we have $v(k) = u(k)$.
                \begin{subproof}
                    Fix $k$.
                    Assume $k\in\initialSegment{n}$.
                    Thus $k\in\naturalsPlus$ and $k \leq n$ by \cref{initial_segment_naturalsplus}.
                    Thus $k\in A$.
                    Thus $v(k) = u(k)$.
                \end{subproof}
            \end{subproof}
            Show that $v\subseteq u$.
            \begin{subproof}
                Show that for all $w$ such that $w\in v$ we have $w\in u$.
                \begin{subproof}
                    Fix $w$.
                    Assume $w\in v$.
                    Thus $v$ is a function by \cref{iteration_sequence}.
                    Take $k\in\dom{v}$ such that $w = (k,v(k))$ by \cref{function_member_elim}.
                    Thus $\dom{v} = \initialSegment{n}$ by \cref{iteration_sequence}.
                    Thus $k\in\initialSegment{n}$.
                    Thus $v(k) = u(k)$.
                    Thus $w = (k,u(k))$.
                    Thus $k\in\naturalsPlus$ and $k \leq n$ by \cref{initial_segment_naturalsplus}.
                    Thus $n \leq m$.
                    $k\in\reals$ by \cref{real_naturals_subset_reals,subseteq}.
                    $n\in\reals$ by \cref{real_naturals_subset_reals,subseteq}.
                    $m\in\reals$ by \cref{real_naturals_subset_reals,subseteq}.
                    Thus $k \leq m$ by \cref{real_leq_trans}.
                    Thus $k\in\initialSegment{m}$ by \cref{initial_segment_naturalsplus}.
                    Thus $\dom{u} = \initialSegment{m}$ by \cref{iteration_sequence}.
                    Thus $k\in\dom{u}$.
                    Thus $u$ is a function by \cref{iteration_sequence}.
                    Thus $(k,u(k))\in u$ by \cref{function_apply_elim}.
                    Thus $w\in u$.
                \end{subproof}
                Thus $v\subseteq u$ by \cref{subseteq}.
            \end{subproof}
        \end{byCase}
        \end{byCase}
    \end{subproof}
    Let $u = \unions{F}$.
    Thus $F$ is a family of functions.
    Thus $u$ is a function by \cref{unions_of_compatible_family_of_function_is_function}.
    Show that $\dom{u} = \naturalsPlus$.
    \begin{subproof}
        Show that $\dom{u}\subseteq\naturalsPlus$.
        \begin{subproof}
            Show that for all $x$ such that $x\in\dom{u}$ we have $x\in\naturalsPlus$.
            \begin{subproof}
                Fix $x$.
                Assume $x\in\dom{u}$.
                Take $y$ such that $(x,y)\in u$ by \cref{dom}.
                Take $v\in F$ such that $(x,y)\in v$ by \cref{unions_iff}.
                Take $n\in\naturalsPlus$ such that
                    $v$ is an iteration sequence for $f$ on $n$ starting at $a_0$.
                Thus $x\in\dom{v}$ by \cref{dom}.
                Thus $\dom{v} = \initialSegment{n}$ by \cref{iteration_sequence}.
                Thus $x\in\initialSegment{n}$.
                Thus $x\in\naturalsPlus$ by \cref{initial_segment_naturalsplus}.
            \end{subproof}
            Thus $\dom{u}\subseteq\naturalsPlus$ by \cref{subseteq}.
        \end{subproof}
        Show that $\naturalsPlus\subseteq\dom{u}$.
        \begin{subproof}
            Show that for all $n\in\naturalsPlus$ we have $n\in\dom{u}$.
            \begin{subproof}
                Fix $n\in\naturalsPlus$.
                $n\in S$ by \cref{subseteq}.
                Take $v$ such that
                    $v$ is an iteration sequence for $f$ on $n$ starting at $a_0$.
                Thus $v\in F$.
                Thus $\dom{v} = \initialSegment{n}$ by \cref{iteration_sequence}.
                Thus $n\in\dom{v}$ by \cref{initial_segment_naturalsplus}.
                Thus $n\in\dom{u}$ by \cref{dom,unions_iff}.
            \end{subproof}
            Thus $\naturalsPlus\subseteq\dom{u}$ by \cref{subseteq}.
        \end{subproof}
        Thus $\dom{u} = \naturalsPlus$ by \cref{subseteq_antisymmetric}.
    \end{subproof}
    Show that for all $n\in\naturalsPlus$ we have $u(n)\in A$.
    \begin{subproof}
        Fix $n\in\naturalsPlus$.
        Thus $n\in\dom{u}$.
        Thus $(n,u(n))\in u$ by \cref{function_apply_elim}.
        Thus $(n,u(n))\in\unions{F}$.
        Take $v\in F$ such that $(n,u(n))\in v$ by \cref{unions_iff}.
        Thus $n\in\dom{v}$ by \cref{dom}.
        Take $n_0\in\naturalsPlus$ such that
            $v$ is an iteration sequence for $f$ on $n_0$ starting at $a_0$.
        Thus $\dom{v} = \initialSegment{n_0}$ by \cref{iteration_sequence}.
        Thus $n\in\initialSegment{n_0}$.
        Thus $n \leq n_0$ by \cref{initial_segment_naturalsplus}.
        Thus $v(1) = a_0$ by \cref{iteration_sequence}.
        Thus for all $k\in\naturalsPlus$ such that $k + 1 \leq n_0$ we have $v(k + 1) = \apply{f}{v(k)}$
            by \cref{iteration_sequence}.
        Show that $v(n)\in A$.
        \begin{subproof}
            Let $B = \{k\in\naturalsPlus \mid \text{if $k\leq n$ then $v(k)\in A$}\}$.
            Show that for all $k\in\naturalsPlus$ we have $k\in B$.
            \begin{subproof}
                Show that $B\subseteq\naturalsPlus$.
                \begin{subproof}
                    Show that for all $k$ such that $k\in B$ we have $k\in\naturalsPlus$.
                    \begin{subproof}
                        Fix $k$.
                        Assume $k\in B$.
                        Then $k\in\naturalsPlus$.
                    \end{subproof}
                    Thus $B\subseteq\naturalsPlus$ by \cref{subseteq}.
                \end{subproof}
                $1\in\naturalsPlus$ by \cref{real_one_in_naturals}.
                Show that $1\in B$.
                \begin{subproof}
                    Show that if $1\leq n$ then $v(1)\in A$.
                    \begin{subproof}
                        Assume $1\leq n$.
                        Then $v(1) = a_0$.
                        Thus $v(1)\in A$.
                    \end{subproof}
                    Thus $1\in B$.
                \end{subproof}
                Show that for all $k\in B$ we have $k + 1\in B$.
                \begin{subproof}
                    Fix $k$.
                    Assume $k\in B$.
                    Show that if $k + 1\leq n$ then $v(k + 1)\in A$.
                    \begin{subproof}
                        Assume $k + 1\leq n$.
                        Then $k \leq n$.
                        Thus $v(k)\in A$.
                        Thus $\apply{f}{v(k)}\in A$.
                        $k + 1\in\naturalsPlus$ by \cref{real_naturals_closed_succ}.
                        $k + 1\in\reals$ by \cref{real_naturals_subset_reals,subseteq}.
                        $n\in\reals$ by \cref{real_naturals_subset_reals,subseteq}.
                        $n_0\in\reals$ by \cref{real_naturals_subset_reals,subseteq}.
                        Thus $k + 1 \leq n_0$ by \cref{real_leq_trans}.
                        Thus $v(k + 1) = \apply{f}{v(k)}$ by \cref{iteration_sequence}.
                        Thus $v(k + 1)\in A$.
                    \end{subproof}
                    Thus $k + 1\in B$.
                \end{subproof}
                Thus for all $k\in\naturalsPlus$ we have $k\in B$ by \cref{real_naturals_induction}.
            \end{subproof}
            Thus $v(n)\in A$ by \cref{initial_segment_naturalsplus}.
        \end{subproof}
        Thus $u(n)\in A$ by \cref{function_apply_intro}.
    \end{subproof}
    Show that $u$ is strictly increasing on $\naturalsPlus$.
    \begin{subproof}
        Show that for all $n\in\naturalsPlus$ we have $u(n) < u(n + 1)$.
        \begin{subproof}
            Fix $n\in\naturalsPlus$.
            $n + 1\in\naturalsPlus$ by \cref{real_naturals_closed_succ}.
            Thus $n + 1\in\dom{u}$.
            Thus $(n + 1,u(n + 1))\in u$ by \cref{function_apply_elim}.
            Thus $(n + 1,u(n + 1))\in\unions{F}$.
            Take $v\in F$ such that $(n + 1,u(n + 1))\in v$ by \cref{unions_iff}.
            Thus $n + 1\in\dom{v}$ by \cref{dom}.
            Thus $v$ is a function.
            Thus $v(n + 1) = u(n + 1)$ by \cref{function_apply_intro}.
            Take $n_0\in\naturalsPlus$ such that
                $v$ is an iteration sequence for $f$ on $n_0$ starting at $a_0$.
            Thus $\dom{v} = \initialSegment{n_0}$ by \cref{iteration_sequence}.
            Thus $n + 1\in\initialSegment{n_0}$.
            Thus $n + 1 \leq n_0$ by \cref{initial_segment_naturalsplus}.
            Thus $v(n + 1) = \apply{f}{v(n)}$ by \cref{iteration_sequence}.
            Thus $u(n + 1) = \apply{f}{v(n)}$.
            Show that $n\in\dom{v}$.
            \begin{subproof}
                Show that $n \leq n_0$.
                \begin{subproof}
                    Show that $n < n + 1$.
                    \begin{subproof}
                        $n\in\reals$ by \cref{real_naturals_subset_reals,subseteq}.
                        $0\in\reals$ by \cref{real_add_ident}.
                        $1\in\reals$ by \cref{real_naturals_subset_reals,real_one_in_naturals,subseteq}.
                        $0 < 1$ by \cref{real_zero_lt_one}.
                        $0 + n < 1 + n$ by \cref{real_lt_add}.
                        $0 + n = n + 0$ by \cref{real_add_comm}.
                        $1 + n = n + 1$ by \cref{real_add_comm}.
                        Thus $n + 0 < n + 1$ by \cref{real_lt_subst_left}.
                        $n + 0 = n$ by \cref{real_add_ident}.
                        Thus $n < n + 1$ by \cref{real_lt_subst_left}.
                    \end{subproof}
                    Thus $n < n + 1$ or $n = n + 1$.
                    Thus $n \leq n + 1$.
                    $n\in\reals$ by \cref{real_naturals_subset_reals,subseteq}.
                    $n + 1\in\naturalsPlus$ by \cref{real_naturals_closed_succ}.
                    $n + 1\in\reals$ by \cref{real_naturals_subset_reals,subseteq}.
                    $n_0\in\reals$ by \cref{real_naturals_subset_reals,subseteq}.
                    Thus $n \leq n_0$ by \cref{real_leq_trans}.
                \end{subproof}
                Thus $n\in\initialSegment{n_0}$ by \cref{initial_segment_naturalsplus}.
                Thus $n\in\dom{v}$ by \cref{iteration_sequence}.
            \end{subproof}
            Thus $v$ is a function by \cref{iteration_sequence}.
            Thus $(n,v(n))\in v$ by \cref{function_apply_elim}.
            Thus $(n,v(n))\in\unions{F}$ by \cref{unions_iff}.
            Thus $(n,v(n))\in u$.
            Thus $u(n) = v(n)$ by \cref{function_apply_intro}.
            Thus $u(n + 1) = \apply{f}{u(n)}$.
            Thus $u(n) < u(n + 1)$ by \cref{function_apply_intro}.
        \end{subproof}
        Show that for all $m,n\in\naturalsPlus$ such that $m < n$ we have $u(m) < u(n)$.
        \begin{subproof}
            Let $T = \{n\in\naturalsPlus \mid \text{for all $m\in\naturalsPlus$ such that $m < n$ we have $u(m) < u(n)$}\}$.
            Show that for all $n\in\naturalsPlus$ we have $n\in T$.
            \begin{subproof}
                Show that $T\subseteq\naturalsPlus$.
                \begin{subproof}
                    Show that for all $n$ such that $n\in T$ we have $n\in\naturalsPlus$.
                    \begin{subproof}
                        Fix $n$.
                        Assume $n\in T$.
                        Then $n\in\naturalsPlus$.
                    \end{subproof}
                    Thus $T\subseteq\naturalsPlus$ by \cref{subseteq}.
                \end{subproof}
                $1\in\naturalsPlus$ by \cref{real_one_in_naturals}.
                Show that $1\in T$.
                \begin{subproof}
                    Show that for all $m\in\naturalsPlus$ such that $m < 1$ we have $u(m) < u(1)$.
                \begin{subproof}
                    Fix $m$.
                    Assume $m\in\naturalsPlus$ and $m < 1$.
                    Suppose not.
                    $m\in\reals$ by \cref{real_naturals_subset_reals,subseteq}.
                    $1\in\reals$ by \cref{real_naturals_subset_reals,real_one_in_naturals,subseteq}.
                    $1 < m$ or $1 = m$ by \cref{real_one_leq_naturalsplus}.
                    \begin{byCase}
                    \caseOf{$1 = m$.}
                        Thus $m < m$ by \cref{real_lt_subst_right}.
                        Contradiction by \cref{real_lt_irreflexive}.
                    \caseOf{$1 < m$.}
                        Thus $1 < 1$ by \cref{real_lt_trans}.
                        Contradiction by \cref{real_lt_irreflexive}.
                    \end{byCase}
                    \end{subproof}
                    Thus $1\in T$.
                \end{subproof}
                Show that for all $n\in T$ we have $n + 1\in T$.
                \begin{subproof}
                    Fix $n$.
                    Assume $n\in T$.
                    Show that for all $m\in\naturalsPlus$ such that $m < n + 1$ we have $u(m) < u(n + 1)$.
                    \begin{subproof}
                        Fix $m$.
                        Assume $m\in\naturalsPlus$ and $m < n + 1$.
                        Suppose not.
                        Thus $n\in\naturalsPlus$ by \cref{subseteq}.
                        Thus $n + 1\in\naturalsPlus$ by \cref{real_naturals_closed_succ}.
                        Thus $u(m)\in A$ and $u(n)\in A$ and $u(n + 1)\in A$.
                        Thus $u(m)\in\naturalsPlus$ and $u(n)\in\naturalsPlus$ and $u(n + 1)\in\naturalsPlus$
                            by \cref{subseteq}.
                        Thus $u(m)\in\reals$ and $u(n)\in\reals$ and $u(n + 1)\in\reals$
                            by \cref{real_naturals_subset_reals,subseteq}.
                        Thus $m < n$ or $m = n$ or $m = n + 1$ by \cref{real_succ_discrete}.
                        \begin{byCase}
                        \caseOf{$m = n$.}
                            Thus $u(n) < u(n + 1)$.
                            Thus $u(m) < u(n + 1)$ by \cref{real_lt_subst_left}.
                            Contradiction.
                        \caseOf{$m < n$.}
                            Thus $u(m) < u(n)$ by \cref{subseteq}.
                            Thus $u(n) < u(n + 1)$.
                            Thus $u(m) < u(n + 1)$ by \cref{real_lt_trans}.
                            Contradiction.
                        \caseOf{$m = n + 1$.}
                            Thus $n + 1 < n + 1$ by \cref{real_lt_subst_left}.
                            $n + 1\in\reals$ by \cref{real_naturals_subset_reals,subseteq}.
                            Contradiction by \cref{real_lt_irreflexive}.
                        \end{byCase}
                    \end{subproof}
                    Thus $n + 1\in T$.
                \end{subproof}
                Thus for all $n\in\naturalsPlus$ we have $n\in T$ by \cref{real_naturals_induction}.
            \end{subproof}
            Thus for all $m,n\in\naturalsPlus$ such that $m < n$ we have $u(m) < u(n)$ by \cref{subseteq}.
        \end{subproof}
        Show that $u\in\funs{\naturalsPlus}{\naturalsPlus}$.
        \begin{subproof}
            Show that $u$ is a function to $\naturalsPlus$.
            \begin{subproof}
                Show that for all $n\in\dom{u}$ we have $u(n)\in\naturalsPlus$.
                \begin{subproof}
                    Fix $n\in\dom{u}$.
                    Thus $n\in\naturalsPlus$.
                    Thus $u(n)\in A$.
                    Thus $u(n)\in\naturalsPlus$ by \cref{subseteq}.
                \end{subproof}
                Thus $u$ is a function to $\naturalsPlus$.
            \end{subproof}
            Thus $u\in\funs{\naturalsPlus}{\naturalsPlus}$ by \cref{funs_intro}.
        \end{subproof}
        Thus $u$ is strictly increasing on $\naturalsPlus$ by \cref{strictly_increasing_on_set}.
    \end{subproof}
    Thus there exists $u$ such that
        $u$ is strictly increasing on $\naturalsPlus$ and
        for all $n\in\naturalsPlus$ we have $u(n)\in A$.
\end{proof}


% from §28 helper lemma 8: constant subsequence in a finite range
\begin{lemma}\label{finite_sequence_constant_subsequence}
    Suppose $A\subseteq X$.
    Suppose $A$ is finite.
    Suppose $s$ is a sequence in $A$.
    Then there exist $t,a$ such that
        $a\in A$ and
        $t$ is a sequence in $A$ and
        $t$ is a subsequence of $s$ and
        for all $n\in\naturalsPlus$ we have $t(n) = a$.
\end{lemma}
\begin{proof}
    $s\in\funs{\naturalsPlus}{A}$ by \cref{sequence_in}.
    Thus $s\in\rels{\naturalsPlus}{A}$ and $\dom{s} = \naturalsPlus$ and $s$ is right-unique by \cref{funs}.
    Show that $s$ is a relation.
    \begin{subproof}
        Show that for all $w$ such that $w\in s$ there exist $x,y$ such that $w = (x,y)$.
        \begin{subproof}
            Fix $w$.
            Assume $w\in s$.
            Thus $s\subseteq\naturalsPlus\times A$ by \cref{rels,pow_iff}.
            Thus $w\in\naturalsPlus\times A$ by \cref{subseteq}.
            Thus there exist $x,y$ such that $w = (x,y)$ by \cref{times_elem_is_tuple}.
        \end{subproof}
        Thus $s$ is a relation by \cref{relation}.
    \end{subproof}
    Thus $s$ is a function.
    Let $g = \{ (a,\preimg{s}{\{a\}}) \mid a\in A \}$.
    Let $F = \img{g}{A}$.
    Show that $g$ is a function on $A$.
    \begin{subproof}
        Show that $g$ is a relation.
        \begin{subproof}
            Show that for all $w$ such that $w\in g$ there exist $a,B$ such that $w = (a,B)$.
            \begin{subproof}
                Fix $w$.
                Assume $w\in g$.
                Take $a\in A$ such that $w = (a,\preimg{s}{\{a\}})$.
                Thus there exist $a,B$ such that $w = (a,B)$.
            \end{subproof}
            Thus $g$ is a relation by \cref{relation}.
        \end{subproof}
        Show that $g$ is right-unique.
        \begin{subproof}
            Show that for all $a,B_1,B_2$ such that $(a,B_1)\in g$ and $(a,B_2)\in g$ we have $B_1 = B_2$.
            \begin{subproof}
                Fix $a$.
                Fix $B_1$.
                Fix $B_2$.
                Assume $(a,B_1)\in g$ and $(a,B_2)\in g$.
                Take $a_1\in A$ such that $(a,B_1) = (a_1,\preimg{s}{\{a_1\}})$.
                Take $a_2\in A$ such that $(a,B_2) = (a_2,\preimg{s}{\{a_2\}})$.
                Thus $a = a_1$ and $B_1 = \preimg{s}{\{a_1\}}$ by \cref{pair_eq_iff}.
                Thus $a = a_2$ and $B_2 = \preimg{s}{\{a_2\}}$ by \cref{pair_eq_iff}.
                Thus $B_1 = \preimg{s}{\{a\}}$ and $B_2 = \preimg{s}{\{a\}}$.
                Thus $B_1 = B_2$.
            \end{subproof}
            Thus $g$ is right-unique by \cref{rightunique}.
        \end{subproof}
        Show that $\dom{g} = A$.
        \begin{subproof}
            Show that $\dom{g}\subseteq A$.
            \begin{subproof}
                Show that for all $a$ such that $a\in\dom{g}$ we have $a\in A$.
                \begin{subproof}
                    Fix $a$.
                    Assume $a\in\dom{g}$.
                    Take $B$ such that $(a,B)\in g$ by \cref{dom}.
                    Thus $a\in A$.
                \end{subproof}
                Thus $\dom{g}\subseteq A$ by \cref{subseteq}.
            \end{subproof}
            Show that $A\subseteq\dom{g}$.
            \begin{subproof}
                Show that for all $a$ such that $a\in A$ we have $a\in\dom{g}$.
                \begin{subproof}
                    Fix $a$.
                    Assume $a\in A$.
                    Thus $(a,\preimg{s}{\{a\}})\in g$.
                    Thus $a\in\dom{g}$ by \cref{dom}.
                \end{subproof}
                Thus $A\subseteq\dom{g}$ by \cref{subseteq}.
            \end{subproof}
            Thus $\dom{g} = A$ by \cref{subseteq_antisymmetric}.
        \end{subproof}
        Thus $g$ is a function on $A$.
    \end{subproof}
    Thus $F$ is finite by \cref{finite_image_function}.
    Show that $\unions{F} = \naturalsPlus$.
    \begin{subproof}
        Show that $\naturalsPlus\subseteq\unions{F}$.
        \begin{subproof}
            Show that for all $n$ such that $n\in\naturalsPlus$ we have $n\in\unions{F}$.
            \begin{subproof}
                Fix $n$.
                Assume $n\in\naturalsPlus$.
                Thus $s(n)\in A$ by \cref{funs_type_apply}.
                Thus $(s(n),\preimg{s}{\{s(n)\}})\in g$.
                Thus $g(s(n)) = \preimg{s}{\{s(n)\}}$ by \cref{function_apply_iff}.
                Thus $\preimg{s}{\{s(n)\}}\in F$ by \cref{img_iff}.
                $s(n)\in\{s(n)\}$ by \cref{singleton_intro}.
                Thus $n\in\preimg{s}{\{s(n)\}}$ by \cref{preimg_iff,function_apply_iff}.
                Thus $n\in\unions{F}$ by \cref{unions_iff}.
            \end{subproof}
            Thus $\naturalsPlus\subseteq\unions{F}$ by \cref{subseteq}.
        \end{subproof}
        Show that $\unions{F}\subseteq\naturalsPlus$.
        \begin{subproof}
            Show that for all $n$ such that $n\in\unions{F}$ we have $n\in\naturalsPlus$.
            \begin{subproof}
                Fix $n$.
                Assume $n\in\unions{F}$.
                Take $B\in F$ such that $n\in B$ by \cref{unions_iff}.
                Take $a\in A$ such that $(a,B)\in g$ by \cref{img_iff}.
                Thus $B = g(a)$ by \cref{function_apply_iff}.
                Thus $B = \preimg{s}{\{a\}}$ by \cref{function_apply_iff}.
                Thus $n\in\preimg{s}{\{a\}}$.
                Thus $n\in\dom{s}$ by \cref{preimg_subseteq_dom,subseteq}.
                Thus $n\in\naturalsPlus$.
            \end{subproof}
            Thus $\unions{F}\subseteq\naturalsPlus$ by \cref{subseteq}.
        \end{subproof}
        Thus $\unions{F} = \naturalsPlus$ by \cref{subseteq_antisymmetric}.
    \end{subproof}
    Show that there exists $a\in A$ such that $\preimg{s}{\{a\}}$ is infinite.
    \begin{subproof}
        Suppose not.
        Then for all $a\in A$ we have $\preimg{s}{\{a\}}$ is finite.
        Show that for all $B\in F$ we have $B$ is finite.
        \begin{subproof}
            Fix $B$.
            Assume $B\in F$.
            Take $a\in A$ such that $(a,B)\in g$ by \cref{img_iff}.
            Thus $B = g(a)$ by \cref{function_apply_iff}.
            Thus $B = \preimg{s}{\{a\}}$ by \cref{function_apply_iff}.
            Thus $B$ is finite.
        \end{subproof}
        Thus $\unions{F}$ is finite by \cref{finite_union_of_finite_family}.
        Thus $\naturalsPlus$ is finite.
        Contradiction by \cref{naturalsplus_infinite}.
    \end{subproof}
    Take $a\in A$ such that $\preimg{s}{\{a\}}$ is infinite.
    Let $I = \preimg{s}{\{a\}}$.
    $I\subseteq\dom{s}$ by \cref{preimg_subseteq_dom}.
    Thus $I\subseteq\naturalsPlus$ by \cref{subseteq}.
    Take $u$ such that
        $u$ is strictly increasing on $\naturalsPlus$ and
        for all $n\in\naturalsPlus$ we have $u(n)\in I$
        by \cref{infinite_naturalsplus_has_strictly_increasing_sequence}.
    Let $t = s\circ u$.
    Show that $t$ is a sequence in $A$.
    \begin{subproof}
        Show that $t\in\funs{\naturalsPlus}{A}$.
        \begin{subproof}
            Show that $t$ is a function.
            \begin{subproof}
                Thus $u\in\funs{\naturalsPlus}{\naturalsPlus}$ by \cref{strictly_increasing_on_set}.
                Thus $s$ is right-unique and $u$ is right-unique by \cref{funs}.
                Thus $t$ is a function by \cref{function_circ}.
            \end{subproof}
            Show that $\dom{t} = \naturalsPlus$.
            \begin{subproof}
                Show that $\dom{t}\subseteq\naturalsPlus$.
                \begin{subproof}
                    Show that for all $n$ such that $n\in\dom{t}$ we have $n\in\naturalsPlus$.
                    \begin{subproof}
                        Fix $n$.
                        Assume $n\in\dom{t}$.
                        Take $y$ such that $(n,y)\in t$ by \cref{dom}.
                        Thus $(n,y)\in s\circ u$.
                        Take $z$ such that $n\mathrel{u} z$ and $z\mathrel{s} y$ by \cref{circ_elem_elim}.
                        Thus $n\in\dom{u}$ by \cref{dom}.
                        Thus $u\in\funs{\naturalsPlus}{\naturalsPlus}$ by \cref{strictly_increasing_on_set}.
                        Thus $\dom{u} = \naturalsPlus$ by \cref{funs}.
                        Thus $n\in\naturalsPlus$.
                    \end{subproof}
                    Thus $\dom{t}\subseteq\naturalsPlus$ by \cref{subseteq}.
                \end{subproof}
                Show that $\naturalsPlus\subseteq\dom{t}$.
                \begin{subproof}
                    Show that for all $n$ such that $n\in\naturalsPlus$ we have $n\in\dom{t}$.
                    \begin{subproof}
                        Fix $n$.
                        Assume $n\in\naturalsPlus$.
                        Thus $u\in\funs{\naturalsPlus}{\naturalsPlus}$ by \cref{strictly_increasing_on_set}.
                        Thus $\dom{u} = \naturalsPlus$ by \cref{funs}.
                        Thus $u\in\rels{\naturalsPlus}{\naturalsPlus}$ and $u$ is right-unique by \cref{funs}.
                        Show that $u$ is a relation.
                        \begin{subproof}
                            Show that for all $w$ such that $w\in u$ there exist $x,y$ such that $w = (x,y)$.
                            \begin{subproof}
                                Fix $w$.
                                Assume $w\in u$.
                                Thus $u\subseteq\naturalsPlus\times\naturalsPlus$ by \cref{rels,pow_iff}.
                                Thus $w\in\naturalsPlus\times\naturalsPlus$ by \cref{subseteq}.
                                Thus there exist $x,y$ such that $w = (x,y)$ by \cref{times_elem_is_tuple}.
                            \end{subproof}
                            Thus $u$ is a relation by \cref{relation}.
                        \end{subproof}
                        Thus $u$ is a function.
                        Thus $n\in\dom{u}$.
                        Thus $u(n)\in I$.
                        Thus $u(n)\in\naturalsPlus$ by \cref{subseteq}.
                        Thus $u(n)\in\dom{s}$.
                        Thus $(n,u(n))\in u$ by \cref{function_apply_elim}.
                        Thus $(u(n),s(u(n)))\in s$ by \cref{function_apply_elim}.
                        Thus $(n,s(u(n)))\in s\circ u$ by \cref{circ_elem_intro}.
                        Thus $n\in\dom{t}$ by \cref{dom}.
                    \end{subproof}
                    Thus $\naturalsPlus\subseteq\dom{t}$ by \cref{subseteq}.
                \end{subproof}
                Thus $\dom{t} = \naturalsPlus$ by \cref{subseteq_antisymmetric}.
            \end{subproof}
            Show that for all $n\in\dom{t}$ we have $t(n)\in A$.
            \begin{subproof}
                Fix $n$.
                Assume $n\in\dom{t}$.
                Thus $n\in\naturalsPlus$.
                Thus $u(n)\in I$.
                Thus $u(n)\in\naturalsPlus$ by \cref{subseteq}.
                Thus $u(n)\in\dom{s}$.
                Thus $(n,t(n))\in t$ by \cref{function_apply_elim}.
                Thus $(n,t(n))\in s\circ u$.
                Take $z$ such that $n\mathrel{u} z$ and $z\mathrel{s} t(n)$ by \cref{circ_elem_elim}.
                Thus $n\in\dom{u}$ by \cref{dom}.
                Thus $u\in\funs{\naturalsPlus}{\naturalsPlus}$ by \cref{strictly_increasing_on_set}.
                Thus $u$ is a function by \cref{funs_is_function}.
                Thus $z = u(n)$ by \cref{function_apply_iff}.
                Thus $t(n) = s(u(n))$ by \cref{function_apply_iff}.
                Thus $t(n)\in A$ by \cref{funs_type_apply}.
            \end{subproof}
            Thus $t$ is a function to $A$.
            Thus $t\in\funs{\naturalsPlus}{A}$ by \cref{funs_intro}.
        \end{subproof}
        Thus $t$ is a sequence in $A$ by \cref{sequence_in}.
    \end{subproof}
    Show that $t$ is a subsequence of $s$.
    \begin{subproof}
        Show that for all $n\in\naturalsPlus$ we have $t(n) = \apply{s}{\apply{u}{n}}$.
        \begin{subproof}
            Fix $n\in\naturalsPlus$.
            Thus $t\in\funs{\naturalsPlus}{A}$ by \cref{sequence_in}.
            Thus $\dom{t} = \naturalsPlus$ by \cref{funs}.
            Thus $n\in\dom{t}$.
            Thus $t$ is a function by \cref{funs_is_function}.
            Thus $(n,t(n))\in t$ by \cref{function_apply_elim}.
            Thus $(n,t(n))\in s\circ u$.
            Take $z$ such that $n\mathrel{u} z$ and $z\mathrel{s} t(n)$ by \cref{circ_elem_elim}.
            Thus $u\in\funs{\naturalsPlus}{\naturalsPlus}$ by \cref{strictly_increasing_on_set}.
            Thus $u$ is a function by \cref{funs_is_function}.
            Thus $z = u(n)$ by \cref{function_apply_iff}.
            Thus $t(n) = s(u(n))$ by \cref{function_apply_iff}.
        \end{subproof}
        Thus $t$ is a subsequence of $s$ by \cref{subsequence}.
    \end{subproof}
    Show that for all $n\in\naturalsPlus$ we have $t(n) = a$.
    \begin{subproof}
        Fix $n\in\naturalsPlus$.
        Thus $u(n)\in I$.
        Take $b\in\{a\}$ such that $u(n)\mathrel{s} b$ by \cref{preimg_iff}.
        Thus $b = a$ by \cref{singleton_iff}.
        Thus $s(u(n)) = a$ by \cref{function_apply_iff}.
        Thus $t\in\funs{\naturalsPlus}{A}$ by \cref{sequence_in}.
        Thus $\dom{t} = \naturalsPlus$ by \cref{funs}.
        Thus $n\in\dom{t}$.
        Thus $t$ is a function by \cref{funs_is_function}.
        Thus $(n,t(n))\in t$ by \cref{function_apply_elim}.
        Thus $(n,t(n))\in s\circ u$.
        Take $z$ such that $n\mathrel{u} z$ and $z\mathrel{s} t(n)$ by \cref{circ_elem_elim}.
        Thus $u\in\funs{\naturalsPlus}{\naturalsPlus}$ by \cref{strictly_increasing_on_set}.
        Thus $u$ is a function by \cref{funs_is_function}.
        Thus $z = u(n)$ by \cref{function_apply_iff}.
        Thus $t(n) = s(u(n))$ by \cref{function_apply_iff}.
        Thus $t(n) = a$.
    \end{subproof}
    Thus there exist $t,a$ such that
        $a\in A$ and
        $t$ is a sequence in $A$ and
        $t$ is a subsequence of $s$ and
        for all $n\in\naturalsPlus$ we have $t(n) = a$.
\end{proof}

% from §28 helper lemma 9: sequential compactness implies Lebesgue number lemma
\begin{lemma}\label{sequentially_compact_lebesgue}
    Suppose $d$ is a metric on $X$.
    Let $T = \metricTopology{d}{X}$.
    Suppose $A$ is an open covering of $X$ with respect to $T$.
    Suppose $X\neq\emptyset$.
    Suppose $X$ is sequentially compact with respect to $T$.
    Then there exists $\delta\in\reals$ such that $0 < \delta$ and
        for all $B$ such that $B\subseteq X$ and
            for all $x,y\in B$ we have $d((x,y)) < \delta$
        there exists $U\in A$ such that $B\subseteq U$.
\end{lemma}
\begin{proof}
    Suppose not.
    Then for all $\delta\in\reals$ such that $0 < \delta$ there exists $B$ such that
        $B\subseteq X$ and
        for all $U\in A$ we have $B\not\subseteq U$ and
        for all $x,y\in B$ we have $d((x,y)) < \delta$.

    $X\subseteq\unions{A}$ by \cref{open_covering,covering}.
    Take $x_0$ such that $x_0\in X$ by \cref{nonempty_inhabited}.
    Thus $x_0\in\unions{A}$ by \cref{subseteq}.
    Take $U_0\in A$ such that $x_0\in U_0$ by \cref{unions_iff}.

    Let $R = \{w\in\naturalsPlus\times X \mid
        \text{for all $U\in A$ we have $\metricBall{d}{X}{\snd{w}}{1 / (\fst{w})}\not\subseteq U$}\}$.
    Show that $R$ is a relation.
    \begin{subproof}
        Show that for all $w$ such that $w\in R$ there exist $n,x$ such that $w = (n,x)$.
        \begin{subproof}
            Fix $w$.
            Assume $w\in R$.
            Thus $w\in\naturalsPlus\times X$.
            Take $n,x$ such that $w = (n,x)$ by \cref{times_elem_is_tuple}.
        \end{subproof}
        Thus $R$ is a relation by \cref{relation}.
    \end{subproof}

    Show that for all $n\in\naturalsPlus$ there exists $x$ such that $n\mathrel{R} x$.
    \begin{subproof}
        Fix $n\in\naturalsPlus$.
        Let $\delta = 1 / n$.
        $0 < \delta$ by \cref{positive_reciprocal_naturalsplus}.
        Take $B$ such that
            $B\subseteq X$ and
            for all $U\in A$ we have $B\not\subseteq U$ and
            for all $x,y\in B$ we have $d((x,y)) < \delta$.
        Show that $B\neq\emptyset$.
        \begin{subproof}
            Suppose not.
            Then $B = \emptyset$.
            Thus $B\subseteq U_0$ by \cref{emptyset_subseteq}.
            Thus $B\not\subseteq U_0$.
            Contradiction.
        \end{subproof}
        Take $x$ such that $x\in B$.
        Show that for all $U\in A$ we have $\metricBall{d}{X}{x}{\delta}\not\subseteq U$.
        \begin{subproof}
            Fix $U$.
            Assume $U\in A$.
            Suppose not.
            Then $\metricBall{d}{X}{x}{\delta}\subseteq U$.
            Show that $B\subseteq\metricBall{d}{X}{x}{\delta}$.
            \begin{subproof}
                Show that for all $y\in B$ we have $y\in\metricBall{d}{X}{x}{\delta}$.
                \begin{subproof}
                    Fix $y$.
                    Assume $y\in B$.
                    Thus $y\in X$ by \cref{subseteq}.
                    Thus $d((x,y)) < \delta$.
                    Thus $y\in\metricBall{d}{X}{x}{\delta}$ by \cref{metric_ball}.
                \end{subproof}
                Thus $B\subseteq\metricBall{d}{X}{x}{\delta}$ by \cref{subseteq}.
            \end{subproof}
            Thus $B\subseteq U$ by \cref{subseteq_transitive}.
            Contradiction.
        \end{subproof}
        Thus $x\in X$ by \cref{subseteq}.
        Thus $(n,x)\in R$.
        Thus $n\mathrel{R} x$.
    \end{subproof}

    Take $s$ such that $s$ is a function on $\naturalsPlus$ and
        for all $n\in\naturalsPlus$ we have $n\mathrel{R} \apply{s}{n}$ by \cref{choice_function}.
    Show that $s$ is a sequence in $X$.
    \begin{subproof}
        Show that $s\in\funs{\naturalsPlus}{X}$.
        \begin{subproof}
            Show that $s$ is a function to $X$.
            \begin{subproof}
                Show that for all $n\in\dom{s}$ we have $s(n)\in X$.
                \begin{subproof}
                    Fix $n\in\dom{s}$.
                    Thus $n\in\naturalsPlus$.
                    Thus $n\mathrel{R} s(n)$.
                    Thus $(n,s(n))\in\naturalsPlus\times X$.
                    Thus $s(n)\in X$ by \cref{times_tuple_elim}.
                \end{subproof}
                Thus $s$ is a function to $X$.
            \end{subproof}
            Thus $s\in\funs{\naturalsPlus}{X}$ by \cref{funs_intro}.
        \end{subproof}
        Thus $s$ is a sequence in $X$ by \cref{sequence_in}.
    \end{subproof}

    Take $t,x$ such that
        $t$ is a sequence in $X$ and
        $t$ is a subsequence of $s$ and
        $x\in X$ and
        $t$ converges to $x$ in $X$ with respect to $T$ by \cref{sequentially_compact}.
    Take $u$ such that
        $u$ is strictly increasing on $\naturalsPlus$ and
        for all $n\in\naturalsPlus$ we have $t(n) = \apply{s}{\apply{u}{n}}$ by \cref{subsequence}.

    Show that for all $n\in\naturalsPlus$ we have $n \leq u(n)$.
    \begin{subproof}
        Let $P = \{n\in\naturalsPlus \mid n \leq u(n)\}$.
        Show that for all $n\in\naturalsPlus$ we have $n\in P$.
        \begin{subproof}
            Show that $P\subseteq\naturalsPlus$.
            \begin{subproof}
                Show that for all $n$ such that $n\in P$ we have $n\in\naturalsPlus$.
                \begin{subproof}
                    Fix $n$.
                    Assume $n\in P$.
                    Then $n\in\naturalsPlus$.
                \end{subproof}
                Thus $P\subseteq\naturalsPlus$ by \cref{subseteq}.
            \end{subproof}
            $1\in\naturalsPlus$ by \cref{real_one_in_naturals}.
            Show that $1\in P$.
            \begin{subproof}
                $u\in\funs{\naturalsPlus}{\naturalsPlus}$ by \cref{strictly_increasing_on_set}.
                Thus $u(1)\in\naturalsPlus$ by \cref{funs_type_apply}.
                Thus $1 \leq u(1)$ by \cref{real_one_leq_naturalsplus}.
                Thus $1\in P$.
            \end{subproof}
            Show that for all $n\in P$ we have $n + 1\in P$.
            \begin{subproof}
                Fix $n$.
                Assume $n\in P$.
                Thus $n \leq u(n)$.
                Thus $n\in\naturalsPlus$.
                Thus $n + 1\in\naturalsPlus$ by \cref{real_naturals_closed_succ}.
                $0 < 1$ by \cref{real_zero_lt_one}.
                $n\in\reals$ by \cref{real_naturals_subset_reals,subseteq}.
                $0\in\reals$ by \cref{real_add_ident}.
                $1\in\reals$ by \cref{real_mul_ident}.
                Thus $0 + n < 1 + n$ by \cref{real_lt_add}.
                $0 + n = n + 0$ by \cref{real_add_comm}.
                $1 + n = n + 1$ by \cref{real_add_comm}.
                Thus $n + 0 < 1 + n$ by \cref{real_lt_subst_left}.
                Thus $n + 0 < n + 1$ by \cref{real_lt_subst_right}.
                Thus $n < n + 1$ by \cref{real_add_ident}.
                $u$ is strictly increasing on $\naturalsPlus$ by \cref{strictly_increasing_on_set}.
                Thus $u(n) < u(n + 1)$ by \cref{strictly_increasing_on_set}.
                Show that it is wrong that $u(n + 1) < n + 1$.
                \begin{subproof}
                    Suppose not.
                    Then $u(n + 1) < n + 1$.
                    $u\in\funs{\naturalsPlus}{\naturalsPlus}$ by \cref{strictly_increasing_on_set}.
                    Thus $u(n + 1)\in\naturalsPlus$ by \cref{funs_type_apply}.
                    $n\in\naturalsPlus$.
                    Thus $u(n + 1) < n$ or $u(n + 1) = n$ or $u(n + 1) = n + 1$
                        by \cref{real_succ_discrete}.
                    \begin{byCase}
                    \caseOf{$u(n + 1) < n$.}
                        $u\in\funs{\naturalsPlus}{\naturalsPlus}$ by \cref{strictly_increasing_on_set}.
                        Thus $u(n)\in\naturalsPlus$ and $u(n + 1)\in\naturalsPlus$ by \cref{funs_type_apply}.
                        Thus $u(n)\in\reals$ and $u(n + 1)\in\reals$
                            by \cref{real_naturals_subset_reals,subseteq}.
                        $n < u(n)$ or $n = u(n)$ by \cref{real_leq_split}.
                        \begin{byCase}
                        \caseOf{$n < u(n)$.}
                            Thus $u(n + 1) < u(n)$ by \cref{real_lt_trans}.
                            Thus $u(n + 1) < u(n + 1)$ by \cref{real_lt_trans}.
                            Contradiction by \cref{real_lt_irreflexive}.
                        \caseOf{$n = u(n)$.}
                            Thus $u(n + 1) < u(n)$ by \cref{real_lt_subst_right}.
                            Thus $u(n + 1) < u(n + 1)$ by \cref{real_lt_trans}.
                            Contradiction by \cref{real_lt_irreflexive}.
                        \end{byCase}
                    \caseOf{$u(n + 1) = n$.}
                        $u\in\funs{\naturalsPlus}{\naturalsPlus}$ by \cref{strictly_increasing_on_set}.
                        Thus $u(n)\in\naturalsPlus$ by \cref{funs_type_apply}.
                        Thus $u(n)\in\reals$ by \cref{real_naturals_subset_reals,subseteq}.
                        Thus $u(n) < n$ by \cref{real_lt_subst_right}.
                        $n < u(n)$ or $n = u(n)$ by \cref{real_leq_split}.
                        \begin{byCase}
                        \caseOf{$n < u(n)$.}
                            Thus $u(n) < u(n)$ by \cref{real_lt_trans}.
                            Contradiction by \cref{real_lt_irreflexive}.
                        \caseOf{$n = u(n)$.}
                            Thus $u(n) < u(n)$ by \cref{real_lt_subst_left}.
                            Contradiction by \cref{real_lt_irreflexive}.
                        \end{byCase}
                    \caseOf{$u(n + 1) = n + 1$.}
                        Thus $n + 1 < n + 1$ by \cref{real_lt_subst_left}.
                        $n + 1\in\naturalsPlus$ by \cref{real_naturals_closed_succ}.
                        Thus $n + 1\in\reals$ by \cref{real_naturals_subset_reals,subseteq}.
                        Contradiction by \cref{real_lt_irreflexive}.
                    \end{byCase}
                \end{subproof}
                $u\in\funs{\naturalsPlus}{\naturalsPlus}$ by \cref{strictly_increasing_on_set}.
                Thus $u(n + 1)\in\naturalsPlus$ by \cref{funs_type_apply}.
                $n + 1\in\naturalsPlus$.
                Thus $u(n + 1)\in\reals$ and $n + 1\in\reals$
                    by \cref{real_naturals_subset_reals,subseteq}.
                Thus $n + 1 < u(n + 1)$ or $n + 1 = u(n + 1)$ or $u(n + 1) < n + 1$
                    by \cref{real_lt_trichotomy}.
                Thus $n + 1 < u(n + 1)$ or $n + 1 = u(n + 1)$.
                Thus $n + 1 \leq u(n + 1)$.
                Thus $n + 1\in P$.
            \end{subproof}
            Thus for all $n\in\naturalsPlus$ we have $n\in P$ by \cref{real_naturals_induction}.
        \end{subproof}
        Thus for all $n\in\naturalsPlus$ we have $n \leq u(n)$ by \cref{subseteq}.
    \end{subproof}

    Show that if $m,n\in\naturalsPlus$ and $m \leq n$ then $1 / n \leq 1 / m$.
    \begin{subproof}
        Assume $m\in\naturalsPlus$ and $n\in\naturalsPlus$.
        Assume $m < n$ or $m = n$.
        Show that $0 < m$ and $0 < n$.
        \begin{subproof}
            $1\in\naturalsPlus$ by \cref{real_one_in_naturals}.
            $0 < 1$ by \cref{real_zero_lt_one}.
            $m\in\reals$ and $n\in\reals$ by \cref{real_naturals_subset_reals,subseteq}.
            $0\in\reals$ by \cref{real_add_ident}.
            $1\in\reals$ by \cref{real_mul_ident}.

            Show that $0 < m$.
            \begin{subproof}
                $1 < m$ or $1 = m$ by \cref{real_one_leq_naturalsplus}.
                \begin{byCase}
                \caseOf{$1 < m$.}
                    Thus $0 < m$ by \cref{real_lt_trans}.
                \caseOf{$1 = m$.}
                    Thus $0 < m$ by \cref{real_lt_subst_left}.
                \end{byCase}
            \end{subproof}

            Show that $0 < n$.
            \begin{subproof}
                $1 < n$ or $1 = n$ by \cref{real_one_leq_naturalsplus}.
                \begin{byCase}
                \caseOf{$1 < n$.}
                    Thus $0 < n$ by \cref{real_lt_trans}.
                \caseOf{$1 = n$.}
                    Thus $0 < n$ by \cref{real_lt_subst_left}.
                \end{byCase}
            \end{subproof}
        \end{subproof}
        Suppose not.
        Then $1 / m < 1 / n$.
        $m\in\reals$ and $n\in\reals$ by \cref{real_naturals_subset_reals,subseteq}.
        $0\in\reals$ by \cref{real_add_ident}.
        Thus $0 \rmul n < m \rmul n$ by \cref{real_lt_mul_pos}.
        $0 \rmul n = 0$ by \cref{real_mul_zero}.
        Thus $0 < m \rmul n$ by \cref{real_lt_subst_left}.
        $1\in\reals$ by \cref{real_mul_ident}.
        $m\neq 0$ and $n\neq 0$ by \cref{real_pos_nonzero}.
        $\inv{m}\in\reals$ and $\inv{n}\in\reals$ by \cref{real_mul_inverse}.
        Thus $1 / m\in\reals$ and $1 / n\in\reals$ by \cref{real_div,real_mul_closed}.
        $m \rmul n\in\reals$ by \cref{real_mul_closed}.
        Thus $(1 / m) \rmul (m \rmul n) < (1 / n) \rmul (m \rmul n)$ by \cref{real_lt_mul_pos}.
        $(1 / m) \rmul (m \rmul n) = (m \rmul n) \rmul (1 / m)$ by \cref{real_mul_comm}.
        $(1 / n) \rmul (m \rmul n) = (m \rmul n) \rmul (1 / n)$ by \cref{real_mul_comm}.
        Thus $(m \rmul n) \rmul (1 / m) < (m \rmul n) \rmul (1 / n)$
            by \cref{real_lt_subst_left,real_lt_subst_right}.
        Show that $(m \rmul n) \rmul (1 / m) = n$.
        \begin{subproof}
            $1 / m = 1 \rmul \inv{m}$ by \cref{real_div}.
            Thus $(m \rmul n) \rmul (1 / m) = (m \rmul n) \rmul \inv{m}$.
            Thus $(m \rmul n) \rmul \inv{m} = n \rmul (m \rmul \inv{m})$ by \cref{real_mul_assoc,real_mul_comm}.
            $m \rmul \inv{m}\in\reals$ by \cref{real_mul_closed}.
            Thus $n \rmul (m \rmul \inv{m}) = (m \rmul \inv{m}) \rmul n$ by \cref{real_mul_comm}.
            $m\neq 0$ by \cref{real_pos_nonzero}.
            $m\rmul\inv{m} = 1$ by \cref{real_mul_inverse}.
            Thus $(m \rmul n) \rmul (1 / m) = 1 \rmul n$.
            Thus $(m \rmul n) \rmul (1 / m) = n$ by \cref{real_mul_ident}.
        \end{subproof}
        Show that $(m \rmul n) \rmul (1 / n) = m$.
        \begin{subproof}
            $1 / n = 1 \rmul \inv{n}$ by \cref{real_div}.
            Thus $(m \rmul n) \rmul (1 / n) = (m \rmul n) \rmul \inv{n}$.
            Thus $(m \rmul n) \rmul \inv{n} = m \rmul (n \rmul \inv{n})$ by \cref{real_mul_assoc}.
            $n \rmul \inv{n}\in\reals$ by \cref{real_mul_closed}.
            Thus $m \rmul (n \rmul \inv{n}) = (n \rmul \inv{n}) \rmul m$ by \cref{real_mul_comm}.
            $n\neq 0$ by \cref{real_pos_nonzero}.
            $n\rmul\inv{n} = 1$ by \cref{real_mul_inverse}.
            Thus $(m \rmul n) \rmul (1 / n) = 1 \rmul m$.
            Thus $(m \rmul n) \rmul (1 / n) = m$ by \cref{real_mul_ident}.
        \end{subproof}
        Thus $n < m$.
        \begin{byCase}
        \caseOf{$m < n$.}
            Thus $n < n$ by \cref{real_lt_trans}.
            Contradiction by \cref{real_lt_irreflexive}.
        \caseOf{$m = n$.}
            Thus $n < n$ by \cref{real_lt_subst_left}.
            Contradiction by \cref{real_lt_irreflexive}.
        \end{byCase}
    \end{subproof}

    $X\subseteq\unions{A}$ by \cref{open_covering,covering}.
    Thus $x\in\unions{A}$ by \cref{subseteq}.
    Take $U\in A$ such that $x\in U$ by \cref{unions_iff}.
    Show that $U\subseteq X$.
    \begin{subproof}
        $T$ is a topology on $X$ by \cref{metric_topology,basis_topology_is_topology,metric_basis_is_basis}.
        Thus $U\in T$ by \cref{open_in,open_covering}.
        Thus $T\subseteq\pow{X}$ by \cref{topology_on}.
        Thus $U\subseteq X$ by \cref{subseteq,powerset_elim}.
    \end{subproof}
    Thus $U$ is open in $X$ with respect to $T$ by \cref{open_covering}.
    Take $\delta\in\reals$ such that $0 < \delta$ and $\metricBall{d}{X}{x}{\delta}\subseteq U$
        by \cref{metric_open_iff}.
    Let $m_2 = 1 + 1$.
    Let $\epsilon = \delta / m_2$.
    $1\in\reals$ by \cref{real_mul_ident}.
    Thus $m_2\in\reals$ by \cref{real_add_closed}.
    $0\in\reals$ by \cref{real_add_ident}.
    $0 < 1$ by \cref{real_zero_lt_one}.
    Thus $0 + 1 < 1 + 1$ by \cref{real_lt_add}.
    $0 + 1 = 1$ by \cref{real_add_ident}.
    Thus $1 < m_2$ by \cref{real_lt_subst_left}.
    Thus $0 < m_2$ by \cref{real_lt_trans}.
    $m_2\neq 0$ by \cref{real_pos_nonzero}.
    Thus $\inv{m_2}\in\reals$ by \cref{real_mul_inverse}.
    Thus $0 < \inv{m_2}$ by \cref{real_inv_pos}.
    $\epsilon = \delta \rmul \inv{m_2}$ by \cref{real_div}.
    Thus $\epsilon\in\reals$ by \cref{real_mul_closed}.
    $0\in\reals$ by \cref{real_add_ident}.
    Thus $0 \rmul \inv{m_2} < \delta \rmul \inv{m_2}$ by \cref{real_lt_mul_pos}.
    $0 \rmul \inv{m_2} = 0$ by \cref{real_mul_zero}.
    Thus $0 < \epsilon$ by \cref{real_lt_subst_left}.
    Take $N\in\naturalsPlus$ such that $0 < 1 / N$ and $1 / N < \epsilon$ by \cref{real_small_reciprocal}.

    Show that $\metricBall{d}{X}{x}{\epsilon}$ is a neighborhood of $x$ in $X$ with respect to $T$.
    \begin{subproof}
        Show that $\metricBall{d}{X}{x}{\epsilon}$ is open in $X$ with respect to $T$.
        \begin{subproof}
            Show that $\metricBall{d}{X}{x}{\epsilon}\subseteq X$.
            \begin{subproof}
                Show that for all $y\in\metricBall{d}{X}{x}{\epsilon}$ we have $y\in X$.
                \begin{subproof}
                    Fix $y$.
                    Assume $y\in\metricBall{d}{X}{x}{\epsilon}$.
                    Thus $y\in X$ by \cref{metric_ball}.
                \end{subproof}
            \end{subproof}
            Show that for all $y\in\metricBall{d}{X}{x}{\epsilon}$ there exists $\eta\in\reals$ such that
                $0 < \eta$ and $\metricBall{d}{X}{y}{\eta}\subseteq\metricBall{d}{X}{x}{\epsilon}$.
            \begin{subproof}
                Fix $y$.
                Assume $y\in\metricBall{d}{X}{x}{\epsilon}$.
                Thus $y\in X$ and $d((x,y)) < \epsilon$ by \cref{metric_ball}.
                $(x,y)\in X\times X$ by \cref{times_tuple_intro}.
                $d\in\funs{X\times X}{\reals}$ by \cref{metric_on}.
                Thus $d((x,y))\in\reals$ by \cref{funs_type_apply}.
                $\epsilon\in\reals$.
                Let $\eta = \epsilon - d((x,y))$.
                $0 < \eta$ by \cref{real_lt_imp_pos_diff}.
                Show that $\metricBall{d}{X}{y}{\eta}\subseteq\metricBall{d}{X}{x}{\epsilon}$.
                \begin{subproof}
                    Show that for all $z$ such that $z\in\metricBall{d}{X}{y}{\eta}$ we have $z\in\metricBall{d}{X}{x}{\epsilon}$.
                    \begin{subproof}
                        Fix $z$.
                        Assume $z\in\metricBall{d}{X}{y}{\eta}$.
                        Thus $z\in X$ and $d((y,z)) < \eta$ by \cref{metric_ball}.
                        $(y,z)\in X\times X$ by \cref{times_tuple_intro}.
                        Thus $d((y,z))\in\reals$ by \cref{funs_type_apply}.
                        $\neg{(d((x,y)))}\in\reals$ by \cref{real_add_inverse}.
                        $\eta = \epsilon + \neg{(d((x,y)))}$ by \cref{real_sub}.
                        Thus $\eta\in\reals$ by \cref{real_add_closed}.
                        $d$ satisfies the triangle inequality on $X$ by \cref{metric_on}.
                        Thus $d((x,z)) \leq d((x,y)) + d((y,z))$ by \cref{metric_triangle_inequality}.
                        $(x,z)\in X\times X$ by \cref{times_tuple_intro}.
                        Thus $d((x,z))\in\reals$ by \cref{funs_type_apply}.
                        Thus $\eta + d((x,y))\in\reals$ by \cref{real_add_closed}.
                        Thus $d((x,y)) + d((y,z))\in\reals$ by \cref{real_add_closed}.
                        Thus $d((y,z)) + d((x,y)) < \eta + d((x,y))$ by \cref{real_lt_add}.
                        $d((y,z)) + d((x,y)) = d((x,y)) + d((y,z))$ by \cref{real_add_comm}.
                        Thus $d((x,y)) + d((y,z)) < \eta + d((x,y))$ by \cref{real_lt_subst_left}.
                        Show that $d((x,z)) < \eta + d((x,y))$.
                        \begin{subproof}
                            $d((x,z)) < d((x,y)) + d((y,z))$ or
                                $d((x,z)) = d((x,y)) + d((y,z))$ by \cref{real_leq_split}.
                            \begin{byCase}
                            \caseOf{$d((x,z)) < d((x,y)) + d((y,z))$.}
                                $d((x,z))\in\reals$.
                                $d((x,y)) + d((y,z))\in\reals$.
                                $\eta + d((x,y))\in\reals$.
                                Thus $d((x,z)) < \eta + d((x,y))$ by \cref{real_lt_trans}.
                            \caseOf{$d((x,z)) = d((x,y)) + d((y,z))$.}
                                Thus $d((x,z)) < \eta + d((x,y))$ by \cref{real_lt_subst_left}.
                            \end{byCase}
                        \end{subproof}
                        Show that $\eta + d((x,y)) = \epsilon$.
                        \begin{subproof}
                            $\eta = \epsilon + \neg{(d((x,y)))}$ by \cref{real_sub}.
                            Thus $\eta + d((x,y)) = (\epsilon + \neg{(d((x,y)))}) + d((x,y))$.
                            Thus $(\epsilon + \neg{(d((x,y)))}) + d((x,y)) = \epsilon + (\neg{(d((x,y)))} + d((x,y)))$
                                by \cref{real_add_assoc}.
                            $\neg{(d((x,y)))} + d((x,y)) = d((x,y)) + \neg{(d((x,y)))}$ by \cref{real_add_comm}.
                            $d((x,y)) + \neg{(d((x,y)))} = 0$ by \cref{real_add_inverse}.
                            Thus $\eta + d((x,y)) = \epsilon + 0$.
                            $\epsilon + 0 = 0 + \epsilon$ by \cref{real_add_comm}.
                            Thus $\eta + d((x,y)) = 0 + \epsilon$.
                            Thus $\eta + d((x,y)) = \epsilon$ by \cref{real_add_ident}.
                        \end{subproof}
                        Thus $d((x,z)) < \epsilon$ by \cref{real_lt_subst_right}.
                        Thus $z\in\metricBall{d}{X}{x}{\epsilon}$ by \cref{metric_ball}.
                    \end{subproof}
                    Thus $\metricBall{d}{X}{y}{\eta}\subseteq\metricBall{d}{X}{x}{\epsilon}$ by \cref{subseteq}.
                \end{subproof}
            \end{subproof}
            Thus $\metricBall{d}{X}{x}{\epsilon}$ is open in $X$ with respect to $T$ by \cref{metric_open_iff}.
        \end{subproof}
        Show that $x\in\metricBall{d}{X}{x}{\epsilon}$.
        \begin{subproof}
            $d((x,x)) = 0$ by \cref{metric_zero_characterization,metric_on}.
            Thus $x\in\metricBall{d}{X}{x}{\epsilon}$ by \cref{metric_ball,real_lt_subst_left}.
        \end{subproof}
        Thus $\metricBall{d}{X}{x}{\epsilon}$ is a neighborhood of $x$ in $X$ with respect to $T$ by \cref{neighborhood}.
    \end{subproof}

    Take $N_0\in\naturalsPlus$ such that for all $n\in\naturalsPlus$ if $N_0 \leq n$ then
        $t(n)\in\metricBall{d}{X}{x}{\epsilon}$ by \cref{converges_to}.
    Show that there exists $n\in\naturalsPlus$ such that $N_0 \leq n$ and $N \leq n$.
    \begin{subproof}
        $N_0\in\reals$ and $N\in\reals$ by \cref{real_naturals_subset_reals,subseteq}.
        \begin{byCase}
        \caseOf{$N_0 = N$.}
            Take $n$ such that $n = N$.
            Thus $N_0 \leq n$ and $N \leq n$.
        \caseOf{$N_0 \neq N$.}
            Thus $N_0 < N$ or $N < N_0$ by \cref{real_lt_trichotomy}.
            \begin{byCase}
            \caseOf{$N_0 < N$.}
                Take $n$ such that $n = N$.
                Thus $N_0 \leq n$ by \cref{real_lt_imp_leq}.
                Thus $N \leq n$.
            \caseOf{$N < N_0$.}
                Take $n$ such that $n = N_0$.
                Thus $N_0 \leq n$.
                Thus $N \leq n$ by \cref{real_lt_imp_leq}.
            \end{byCase}
        \end{byCase}
    \end{subproof}
    Take $n\in\naturalsPlus$ such that $N_0 \leq n$ and $N \leq n$.

    Thus $t(n)\in\metricBall{d}{X}{x}{\epsilon}$.
    Thus $d((x,t(n))) < \epsilon$ by \cref{metric_ball}.
    $u\in\funs{\naturalsPlus}{\naturalsPlus}$ by \cref{strictly_increasing_on_set}.
    Thus $u(n)\in\naturalsPlus$ by \cref{funs_type_apply}.
    Thus $n \leq u(n)$.
    $N\in\reals$ and $n\in\reals$ by \cref{real_naturals_subset_reals,subseteq}.
    Thus $u(n)\in\reals$ by \cref{real_naturals_subset_reals,subseteq}.
    Thus $N \leq u(n)$ by \cref{real_leq_trans}.
    Thus $1 / u(n) \leq 1 / N$.
    Show that $1 / u(n) < \epsilon$.
    \begin{subproof}
        $u(n)\in\naturalsPlus$.
        $u(n)\in\reals$ and $N\in\reals$ by \cref{real_naturals_subset_reals,subseteq}.
        $1\in\naturalsPlus$ by \cref{real_one_in_naturals}.
        $1 \leq u(n)$ by \cref{real_one_leq_naturalsplus}.
        $1 \leq N$ by \cref{real_one_leq_naturalsplus}.
        $0 < 1$ by \cref{real_zero_lt_one}.
        Thus $0 < u(n)$ and $0 < N$ by \cref{real_lt_trans}.
        $u(n)\neq 0$ and $N\neq 0$ by \cref{real_pos_nonzero}.
        Thus $1 / u(n)\in\reals$ and $1 / N\in\reals$ by \cref{real_div,real_mul_closed,real_mul_inverse}.
        $\epsilon\in\reals$.
        $1 / u(n) < 1 / N$ or $1 / u(n) = 1 / N$ by \cref{real_leq_split}.
        \begin{byCase}
        \caseOf{$1 / u(n) < 1 / N$.}
            Thus $1 / u(n) < \epsilon$ by \cref{real_lt_trans}.
        \caseOf{$1 / u(n) = 1 / N$.}
            Thus $1 / u(n) < \epsilon$ by \cref{real_lt_subst_left}.
        \end{byCase}
    \end{subproof}

    Show that $\metricBall{d}{X}{t(n)}{1 / u(n)}\subseteq U$.
    \begin{subproof}
        Show that for all $z$ such that $z\in\metricBall{d}{X}{t(n)}{1 / u(n)}$ we have $z\in U$.
        \begin{subproof}
            Fix $z$.
            Assume $z\in\metricBall{d}{X}{t(n)}{1 / u(n)}$.
            Thus $z\in X$ and $d((t(n),z)) < 1 / u(n)$ by \cref{metric_ball}.
            $t$ is a sequence in $X$.
            Thus $t\in\funs{\naturalsPlus}{X}$ by \cref{sequence_in}.
            Thus $t(n)\in X$ by \cref{funs_type_apply}.
            $d$ satisfies the triangle inequality on $X$ by \cref{metric_on}.
            Thus $d((x,z)) \leq d((x,t(n))) + d((t(n),z))$ by \cref{metric_triangle_inequality}.
            Show that $d((x,z)) < \epsilon + \epsilon$.
            \begin{subproof}
                $d((t(n),z)) < 1 / u(n)$ by \cref{metric_ball}.
                $(t(n),z)\in X\times X$ by \cref{times_tuple_intro}.
                $d\in\funs{X\times X}{\reals}$ by \cref{metric_on}.
                Thus $d((t(n),z))\in\reals$ by \cref{funs_type_apply}.
                $1 / u(n)\in\reals$.
                $\epsilon\in\reals$.
                Thus $d((t(n),z)) < \epsilon$ by \cref{real_lt_trans}.
                Thus $d((x,t(n))) < \epsilon$ by \cref{metric_ball}.
                $(x,t(n))\in X\times X$ by \cref{times_tuple_intro}.
                Thus $d((x,t(n)))\in\reals$ by \cref{funs_type_apply}.
                Thus $d((t(n),z)) + d((x,t(n))) < \epsilon + d((x,t(n)))$ by \cref{real_lt_add}.
                $d((t(n),z)) + d((x,t(n))) = d((x,t(n))) + d((t(n),z))$ by \cref{real_add_comm}.
                $\epsilon + d((x,t(n))) = d((x,t(n))) + \epsilon$ by \cref{real_add_comm}.
                Thus $d((x,t(n))) + d((t(n),z)) < d((x,t(n))) + \epsilon$
                    by \cref{real_lt_subst_left,real_lt_subst_right}.
                Thus $d((x,t(n))) + \epsilon < \epsilon + \epsilon$ by \cref{real_lt_add}.
                Thus $d((x,t(n))) + d((t(n),z))\in\reals$ by \cref{real_add_closed}.
                Thus $d((x,t(n))) + \epsilon\in\reals$ by \cref{real_add_closed}.
                Thus $\epsilon + \epsilon\in\reals$ by \cref{real_add_closed}.
                Thus $d((x,t(n))) + d((t(n),z)) < \epsilon + \epsilon$ by \cref{real_lt_trans}.
                $x\in X$.
                $z\in X$.
                Thus $(x,z)\in X\times X$ by \cref{times_tuple_intro}.
                Thus $d((x,z))\in\reals$ by \cref{funs_type_apply}.
                $d((x,z)) < d((x,t(n))) + d((t(n),z))$ or
                    $d((x,z)) = d((x,t(n))) + d((t(n),z))$ by \cref{real_leq_split}.
                \begin{byCase}
                \caseOf{$d((x,z)) < d((x,t(n))) + d((t(n),z))$.}
                    Thus $d((x,z)) < \epsilon + \epsilon$ by \cref{real_lt_trans}.
                \caseOf{$d((x,z)) = d((x,t(n))) + d((t(n),z))$.}
                    Thus $d((x,z)) < \epsilon + \epsilon$ by \cref{real_lt_subst_left}.
                \end{byCase}
            \end{subproof}
            Show that $\epsilon + \epsilon = \delta$.
            \begin{subproof}
                $\epsilon = \delta \rmul \inv{m_2}$ by \cref{real_div}.
                $m_2 \rmul \epsilon = \delta$ by \cref{real_mul_assoc,real_mul_inverse,real_mul_comm,real_mul_ident}.
                $m_2 \rmul \epsilon = \epsilon + \epsilon$ by \cref{real_right_distrib,real_mul_ident}.
                Thus $\epsilon + \epsilon = \delta$.
            \end{subproof}
            Thus $d((x,z)) < \delta$ by \cref{real_lt_subst_right}.
            Thus $z\in\metricBall{d}{X}{x}{\delta}$ by \cref{metric_ball}.
            Thus $z\in U$ by \cref{subseteq}.
        \end{subproof}
        Thus $\metricBall{d}{X}{t(n)}{1 / u(n)}\subseteq U$ by \cref{subseteq}.
    \end{subproof}

    $n\in\naturalsPlus$.
    Thus $u(n)\in\naturalsPlus$.
    Thus $u(n)\mathrel{R} \apply{s}{\apply{u}{n}}$.
    Thus $(u(n),\apply{s}{\apply{u}{n}})\in R$.
    $t(n) = \apply{s}{\apply{u}{n}}$.
    Thus $(u(n),t(n)) = (u(n),\apply{s}{\apply{u}{n}})$ by \cref{pair_eq_iff}.
    Thus $(u(n),t(n))\in R$.
    Thus for all $V\in A$ we have $\metricBall{d}{X}{t(n)}{1 / u(n)}\not\subseteq V$.
    Thus $\metricBall{d}{X}{t(n)}{1 / u(n)}\not\subseteq U$.
    Contradiction.
\end{proof}

% from §28 helper lemma 11: iteration sequence on naturalsPlus
\begin{lemma}\label{iteration_sequence_naturalsplus}
    Suppose $f\in\funs{A}{A}$.
    Suppose $a_0\in A$.
    Then there exists $u$ such that
        $u$ is a sequence in $A$ and
        $u(1) = a_0$ and
        for all $n\in\naturalsPlus$ we have $u(n + 1) = \apply{f}{u(n)}$.
\end{lemma}
\begin{proof}
    Let $S = \iterationSet{f}{a_0}$.
    Show that for all $n\in\naturalsPlus$ we have $n\in S$.
    \begin{subproof}
        Show that $S\subseteq\naturalsPlus$.
        \begin{subproof}
            Show that for all $n$ such that $n\in S$ we have $n\in\naturalsPlus$.
            \begin{subproof}
                Fix $n$.
                Assume $n\in S$.
                Thus $n\in\naturalsPlus$ by \cref{iteration_set_elim}.
            \end{subproof}
            Thus $S\subseteq\naturalsPlus$ by \cref{subseteq}.
        \end{subproof}
        $1\in\naturalsPlus$ by \cref{real_one_in_naturals}.
        Show that $1\in S$.
        \begin{subproof}
            Let $u = \{(1,a_0)\}$.
            Show that $u$ is a function.
            \begin{subproof}
                Show that $u$ is a relation.
                \begin{subproof}
                    Show that for all $w$ such that $w\in u$ there exist $x,y$ such that $w = (x,y)$.
                    \begin{subproof}
                        Fix $w$.
                        Assume $w\in u$.
                        Then $w = (1,a_0)$.
                    \end{subproof}
                    Thus $u$ is a relation by \cref{relation}.
                \end{subproof}
                Show that $u$ is right-unique.
                \begin{subproof}
                    Show that for all $x,y,y'$ such that $(x,y)\in u$ and $(x,y')\in u$ we have $y = y'$.
                    \begin{subproof}
                        Fix $x$.
                        Fix $y$.
                        Fix $y'$.
                        Assume $(x,y)\in u$ and $(x,y')\in u$.
                        Then $(x,y) = (1,a_0)$ and $(x,y') = (1,a_0)$.
                        Thus $y = y'$.
                    \end{subproof}
                    Thus $u$ is right-unique by \cref{rightunique}.
                \end{subproof}
                Thus $u$ is a function.
            \end{subproof}
            Show that $\dom{u} = \initialSegment{1}$.
            \begin{subproof}
                Show that $\dom{u}\subseteq\initialSegment{1}$.
                \begin{subproof}
                    Show that for all $x$ such that $x\in\dom{u}$ we have $x\in\initialSegment{1}$.
                    \begin{subproof}
                        Fix $x$.
                        Assume $x\in\dom{u}$.
                        Take $y$ such that $(x,y)\in u$ by \cref{dom}.
                        Then $(x,y) = (1,a_0)$.
                        Thus $x = 1$.
                        $1\in\naturalsPlus$ by \cref{real_one_in_naturals}.
                        Thus $x\in\initialSegment{1}$ by \cref{initial_segment_naturalsplus}.
                    \end{subproof}
                    Thus $\dom{u}\subseteq\initialSegment{1}$ by \cref{subseteq}.
                \end{subproof}
                Show that $\initialSegment{1}\subseteq\dom{u}$.
                \begin{subproof}
                    Show that for all $x$ such that $x\in\initialSegment{1}$ we have $x\in\dom{u}$.
                    \begin{subproof}
                        Fix $x$.
                        Assume $x\in\initialSegment{1}$.
                        Thus $x\in\naturalsPlus$ and $x \leq 1$ by \cref{initial_segment_naturalsplus}.
                        $1\in\reals$ by \cref{real_naturals_subset_reals,real_one_in_naturals,subseteq}.
                        $x\in\reals$ by \cref{real_naturals_subset_reals,subseteq}.
                        $1 < x$ or $1 = x$ by \cref{real_one_leq_naturalsplus}.
                        \begin{byCase}
                        \caseOf{$1 = x$.}
                            Thus $(x,a_0) = (1,a_0)$.
                            Thus $(x,a_0)\in u$ by \cref{singleton_intro}.
                            Thus $x\in\dom{u}$ by \cref{dom}.
                        \caseOf{$1 < x$.}
                            Suppose not.
                            Show that $1 < 1$.
                            \begin{subproof}
                                $x < 1$ or $x = 1$ by \cref{real_leq_split}.
                                \begin{byCase}
                                \caseOf{$x < 1$.}
                                    Thus $1 < 1$ by \cref{real_lt_trans}.
                                \caseOf{$x = 1$.}
                                    Thus $1 < 1$ by \cref{real_lt_subst_right}.
                                \end{byCase}
                            \end{subproof}
                            Contradiction by \cref{real_lt_irreflexive}.
                        \end{byCase}
                    \end{subproof}
                    Thus $\initialSegment{1}\subseteq\dom{u}$ by \cref{subseteq}.
                \end{subproof}
                Thus $\dom{u} = \initialSegment{1}$ by \cref{subseteq_antisymmetric}.
            \end{subproof}
            Thus $(1,a_0)\in u$ by \cref{singleton_intro}.
            Thus $u(1) = a_0$ by \cref{function_apply_intro}.
            Show that for all $k\in\naturalsPlus$ such that $k + 1 \leq 1$ we have $u(k + 1) = \apply{f}{u(k)}$.
            \begin{subproof}
                Fix $k$.
                Assume $k\in\naturalsPlus$.
                Assume $k + 1 \leq 1$.
                Suppose not.
                $k\in\reals$ by \cref{real_naturals_subset_reals,subseteq}.
                $1\in\reals$ by \cref{real_naturals_subset_reals,real_one_in_naturals,subseteq}.
                $k + 1\in\reals$ by \cref{real_add_closed}.
                Show that $k < k + 1$.
                \begin{subproof}
                    $0\in\reals$ by \cref{real_add_ident}.
                    $0 < 1$ by \cref{real_zero_lt_one}.
                    $0 + k < 1 + k$ by \cref{real_lt_add}.
                    $0 + k = k + 0$ by \cref{real_add_comm}.
                    $1 + k = k + 1$ by \cref{real_add_comm}.
                    Thus $k + 0 < k + 1$ by \cref{real_lt_subst_left}.
                    $k + 0 = k$ by \cref{real_add_ident}.
                    Thus $k < k + 1$ by \cref{real_lt_subst_left}.
                \end{subproof}
                Show that $1 < k + 1$.
                \begin{subproof}
                    $1 < k$ or $1 = k$ by \cref{real_one_leq_naturalsplus}.
                    \begin{byCase}
                    \caseOf{$1 = k$.}
                        Thus $1 < k + 1$ by \cref{real_lt_subst_left}.
                    \caseOf{$1 < k$.}
                        Thus $1 < k + 1$ by \cref{real_lt_trans}.
                    \end{byCase}
                \end{subproof}
                $k + 1 < 1$ or $k + 1 = 1$ by \cref{real_leq_split}.
                Show that $1 < 1$.
                \begin{subproof}
                    \begin{byCase}
                    \caseOf{$k + 1 < 1$.}
                        Thus $1 < 1$ by \cref{real_lt_trans}.
                    \caseOf{$k + 1 = 1$.}
                        Thus $1 < 1$ by \cref{real_lt_subst_right}.
                    \end{byCase}
                \end{subproof}
                Contradiction by \cref{real_lt_irreflexive}.
            \end{subproof}
            Thus $1\in\naturalsPlus$.
            Thus $u$ is an iteration sequence for $f$ on $1$ starting at $a_0$ by \cref{iteration_sequence}.
            Thus there exists $u$ such that
                $u$ is an iteration sequence for $f$ on $1$ starting at $a_0$.
            Thus $1\in S$ by \cref{iteration_set_intro}.
        \end{subproof}
        Show that for all $n\in S$ we have $n + 1\in S$.
        \begin{subproof}
            Fix $n$.
            Assume $n\in S$.
            Thus $n\in\naturalsPlus$ and there exists $u$ such that
                $u$ is an iteration sequence for $f$ on $n$ starting at $a_0$
                by \cref{iteration_set_elim}.
            Take $u$ such that
                $u$ is an iteration sequence for $f$ on $n$ starting at $a_0$.
            Thus $n + 1\in\naturalsPlus$ by \cref{real_naturals_closed_succ}.
            Let $v = u\union\{(n + 1,\apply{f}{u(n)})\}$.
            Show that $v$ is a function.
            \begin{subproof}
                Show that $v$ is a relation.
                \begin{subproof}
                    Thus $u$ is a relation.
                    Show that $\{(n + 1,\apply{f}{u(n)})\}$ is a relation.
                    \begin{subproof}
                        Show that for all $w$ such that $w\in\{(n + 1,\apply{f}{u(n)})\}$ there exist $x,y$ such that $w = (x,y)$.
                        \begin{subproof}
                            Fix $w$.
                            Assume $w\in\{(n + 1,\apply{f}{u(n)})\}$.
                            Then $w = (n + 1,\apply{f}{u(n)})$.
                        \end{subproof}
                        Thus $\{(n + 1,\apply{f}{u(n)})\}$ is a relation by \cref{relation}.
                    \end{subproof}
                    Thus $v$ is a relation by \cref{union_relations_is_relation}.
                \end{subproof}
                Show that $v$ is right-unique.
                \begin{subproof}
                    Show that for all $x,y,y'$ such that $(x,y)\in v$ and $(x,y')\in v$ we have $y = y'$.
                    \begin{subproof}
                        Fix $x$.
                        Fix $y$.
                        Fix $y'$.
                        Assume $(x,y)\in v$ and $(x,y')\in v$.
                        \begin{byCase}
                        \caseOf{$x = n + 1$.}
                            $(n + 1,y)\in v$.
                            $(n + 1,y')\in v$.
                            Show that $(n + 1,y)\in\{(n + 1,\apply{f}{u(n)})\}$.
                            \begin{subproof}
                                Suppose not.
                                Thus $(n + 1,y)\in u$ by \cref{union_iff}.
                                Thus $n + 1\in\dom{u}$ by \cref{dom}.
                                Thus $\dom{u} = \initialSegment{n}$ by \cref{iteration_sequence}.
                                Thus $n + 1\in\initialSegment{n}$.
                                Thus $n + 1 \leq n$ by \cref{initial_segment_naturalsplus}.
                                Show that $n < n + 1$.
                                \begin{subproof}
                                    $n\in\reals$ by \cref{real_naturals_subset_reals,subseteq}.
                                    $0\in\reals$ by \cref{real_add_ident}.
                                    $1\in\reals$ by \cref{real_naturals_subset_reals,real_one_in_naturals,subseteq}.
                                    $0 < 1$ by \cref{real_zero_lt_one}.
                                    $0 + n < 1 + n$ by \cref{real_lt_add}.
                                    $0 + n = n + 0$ by \cref{real_add_comm}.
                                    $1 + n = n + 1$ by \cref{real_add_comm}.
                                    Thus $n + 0 < n + 1$ by \cref{real_lt_subst_left}.
                                    $n + 0 = n$ by \cref{real_add_ident}.
                                    Thus $n < n + 1$ by \cref{real_lt_subst_left}.
                                \end{subproof}
                                $n + 1 < n$ or $n + 1 = n$ by \cref{real_leq_split}.
                                \begin{byCase}
                                \caseOf{$n + 1 < n$.}
                                    $n\in\reals$ by \cref{real_naturals_subset_reals,subseteq}.
                                    $1\in\reals$ by \cref{real_naturals_subset_reals,real_one_in_naturals,subseteq}.
                                    $n + 1\in\reals$ by \cref{real_add_closed}.
                                    Thus $n < n$ by \cref{real_lt_trans}.
                                    Contradiction by \cref{real_lt_irreflexive}.
                                \caseOf{$n + 1 = n$.}
                                    $n\in\reals$ by \cref{real_naturals_subset_reals,subseteq}.
                                    Thus $n < n$ by \cref{real_lt_subst_right}.
                                    Contradiction by \cref{real_lt_irreflexive}.
                                \end{byCase}
                            \end{subproof}
                            Thus $(n + 1,y) = (n + 1,\apply{f}{u(n)})$ by \cref{singleton_iff}.
                            Thus $y = \apply{f}{u(n)}$ by \cref{pair_eq_iff}.
                            Show that $(n + 1,y')\in\{(n + 1,\apply{f}{u(n)})\}$.
                            \begin{subproof}
                                Suppose not.
                                Thus $(n + 1,y')\in u$ by \cref{union_iff}.
                                Thus $n + 1\in\dom{u}$ by \cref{dom}.
                                Thus $\dom{u} = \initialSegment{n}$ by \cref{iteration_sequence}.
                                Thus $n + 1\in\initialSegment{n}$.
                                Thus $n + 1 \leq n$ by \cref{initial_segment_naturalsplus}.
                                Show that $n < n + 1$.
                                \begin{subproof}
                                    $n\in\reals$ by \cref{real_naturals_subset_reals,subseteq}.
                                    $0\in\reals$ by \cref{real_add_ident}.
                                    $1\in\reals$ by \cref{real_naturals_subset_reals,real_one_in_naturals,subseteq}.
                                    $0 < 1$ by \cref{real_zero_lt_one}.
                                    $0 + n < 1 + n$ by \cref{real_lt_add}.
                                    $0 + n = n + 0$ by \cref{real_add_comm}.
                                    $1 + n = n + 1$ by \cref{real_add_comm}.
                                    Thus $n + 0 < n + 1$ by \cref{real_lt_subst_left}.
                                    $n + 0 = n$ by \cref{real_add_ident}.
                                    Thus $n < n + 1$ by \cref{real_lt_subst_left}.
                                \end{subproof}
                                $n + 1 < n$ or $n + 1 = n$ by \cref{real_leq_split}.
                                \begin{byCase}
                                \caseOf{$n + 1 < n$.}
                                    $n\in\reals$ by \cref{real_naturals_subset_reals,subseteq}.
                                    $1\in\reals$ by \cref{real_naturals_subset_reals,real_one_in_naturals,subseteq}.
                                    $n + 1\in\reals$ by \cref{real_add_closed}.
                                    Thus $n < n$ by \cref{real_lt_trans}.
                                    Contradiction by \cref{real_lt_irreflexive}.
                                \caseOf{$n + 1 = n$.}
                                    $n\in\reals$ by \cref{real_naturals_subset_reals,subseteq}.
                                    Thus $n < n$ by \cref{real_lt_subst_right}.
                                    Contradiction by \cref{real_lt_irreflexive}.
                                \end{byCase}
                            \end{subproof}
                            Thus $(n + 1,y') = (n + 1,\apply{f}{u(n)})$ by \cref{singleton_iff}.
                            Thus $y' = \apply{f}{u(n)}$ by \cref{pair_eq_iff}.
                            Thus $y = y'$.
                        \caseOf{$x \neq n + 1$.}
                            Show that $(x,y)\in u$.
                            \begin{subproof}
                                Suppose not.
                                Thus $(x,y)\in\{(n + 1,\apply{f}{u(n)})\}$ by \cref{union_iff}.
                                Thus $(x,y) = (n + 1,\apply{f}{u(n)})$ by \cref{singleton_iff}.
                                Thus $x = n + 1$ by \cref{pair_eq_iff}.
                                Contradiction.
                            \end{subproof}
                            Show that $(x,y')\in u$.
                            \begin{subproof}
                                Suppose not.
                                Thus $(x,y')\in\{(n + 1,\apply{f}{u(n)})\}$ by \cref{union_iff}.
                                Thus $(x,y') = (n + 1,\apply{f}{u(n)})$ by \cref{singleton_iff}.
                                Thus $x = n + 1$ by \cref{pair_eq_iff}.
                                Contradiction.
                            \end{subproof}
                            Thus $u$ is a function by \cref{iteration_sequence}.
                            Thus $y = y'$ by \cref{function_rightunique}.
                        \end{byCase}
                    \end{subproof}
                    Thus $v$ is right-unique by \cref{rightunique}.
                \end{subproof}
                Thus $v$ is a function.
            \end{subproof}
            Show that $\dom{v} = \initialSegment{n + 1}$.
            \begin{subproof}
                Show that $\dom{v}\subseteq\initialSegment{n + 1}$.
                \begin{subproof}
                    Show that for all $x$ such that $x\in\dom{v}$ we have $x\in\initialSegment{n + 1}$.
                    \begin{subproof}
                        Fix $x$.
                        Assume $x\in\dom{v}$.
                        Take $y$ such that $(x,y)\in v$ by \cref{dom}.
                        Thus $(x,y)\in u$ or $(x,y)\in\{(n + 1,\apply{f}{u(n)})\}$ by \cref{union_iff}.
                        \begin{byCase}
                        \caseOf{$(x,y)\in u$.}
                            Thus $x\in\dom{u}$ by \cref{dom}.
                            Thus $\dom{u} = \initialSegment{n}$ by \cref{iteration_sequence}.
                            Thus $x\in\initialSegment{n}$.
                            Thus $x\in\naturalsPlus$ and $x \leq n$ by \cref{initial_segment_naturalsplus}.
                            Thus $x \leq n + 1$.
                            Thus $x\in\initialSegment{n + 1}$ by \cref{initial_segment_naturalsplus}.
                        \caseOf{$(x,y)\in\{(n + 1,\apply{f}{u(n)})\}$.}
                            Thus $(x,y) = (n + 1,\apply{f}{u(n)})$ by \cref{singleton_iff}.
                            Thus $x = n + 1$ by \cref{pair_eq_iff}.
                            Thus $x\in\initialSegment{n + 1}$ by \cref{initial_segment_naturalsplus}.
                        \end{byCase}
                    \end{subproof}
                    Thus $\dom{v}\subseteq\initialSegment{n + 1}$ by \cref{subseteq}.
                \end{subproof}
                Show that $\initialSegment{n + 1}\subseteq\dom{v}$.
                \begin{subproof}
                    Show that for all $x$ such that $x\in\initialSegment{n + 1}$ we have $x\in\dom{v}$.
                    \begin{subproof}
                        Fix $x$.
                        Assume $x\in\initialSegment{n + 1}$.
                        Thus $x\in\naturalsPlus$ and $x \leq n + 1$ by \cref{initial_segment_naturalsplus}.
                        \begin{byCase}
                        \caseOf{$x = n + 1$.}
                            Thus $(x,\apply{f}{u(n)}) = (n + 1,\apply{f}{u(n)})$.
                            Thus $(x,\apply{f}{u(n)})\in\{(n + 1,\apply{f}{u(n)})\}$ by \cref{cons_left}.
                            Thus $(x,\apply{f}{u(n)})\in v$ by \cref{union_intro_right}.
                            Thus $x\in\dom{v}$ by \cref{dom}.
                        \caseOf{$x \neq n + 1$.}
                            Thus $x \leq n$ by \cref{real_naturals_leq_succ_elim}.
                            Thus $x\in\initialSegment{n}$ by \cref{initial_segment_naturalsplus}.
                            Thus $x\in\dom{u}$ by \cref{iteration_sequence}.
                            Thus $x\in\dom{v}$ by \cref{dom,union_iff}.
                        \end{byCase}
                    \end{subproof}
                    Thus $\initialSegment{n + 1}\subseteq\dom{v}$ by \cref{subseteq}.
                \end{subproof}
                Thus $\dom{v} = \initialSegment{n + 1}$ by \cref{subseteq_antisymmetric}.
            \end{subproof}
            Show that $v(1) = a_0$.
            \begin{subproof}
                Thus $u$ is a function by \cref{iteration_sequence}.
                Thus $1\in\naturalsPlus$ by \cref{real_one_in_naturals}.
                Thus $1 < n$ or $1 = n$ by \cref{real_one_leq_naturalsplus}.
                Thus $1 \leq n$.
                Thus $1\in\initialSegment{n}$ by \cref{initial_segment_naturalsplus}.
                Thus $1\in\dom{u}$ by \cref{iteration_sequence}.
                Thus $(1,u(1))\in u$ by \cref{function_apply_elim}.
                Thus $(1,u(1))\in v$ by \cref{union_intro_left}.
                Thus $v(1) = u(1)$ by \cref{function_apply_intro}.
                Thus $u(1) = a_0$ by \cref{iteration_sequence}.
                Thus $v(1) = a_0$.
            \end{subproof}
            Show that for all $k\in\naturalsPlus$ such that $k + 1 \leq n + 1$ we have $v(k + 1) = \apply{f}{v(k)}$.
            \begin{subproof}
                Fix $k$.
                Assume $k\in\naturalsPlus$.
                Assume $k + 1 \leq n + 1$.
                \begin{byCase}
                \caseOf{$k + 1 = n + 1$.}
                    Thus $k = n$.
                    Thus $(k + 1,\apply{f}{u(n)})\in\{(n + 1,\apply{f}{u(n)})\}$ by \cref{cons_left}.
                    Thus $(k + 1,\apply{f}{u(n)})\in v$ by \cref{union_intro_right}.
                    Thus $v(k + 1) = \apply{f}{u(n)}$ by \cref{function_apply_intro}.
                    Thus $v(k + 1) = \apply{f}{u(k)}$.
                    Thus $k \leq n$.
                    Thus $k\in\initialSegment{n}$ by \cref{initial_segment_naturalsplus}.
                    Thus $k\in\dom{u}$ by \cref{iteration_sequence}.
                    Thus $u$ is a function by \cref{iteration_sequence}.
                    Thus $(k,u(k))\in u$ by \cref{function_apply_elim}.
                    Thus $(k,u(k))\in v$ by \cref{union_intro_left}.
                    Thus $v(k) = u(k)$ by \cref{function_apply_intro}.
                    Thus $v(k + 1) = \apply{f}{v(k)}$.
                \caseOf{$k + 1 \neq n + 1$.}
                    Thus $k + 1\in\naturalsPlus$ by \cref{real_naturals_closed_succ}.
                    Thus $k + 1 \leq n$ by \cref{real_naturals_leq_succ_elim}.
                    Thus $k + 1\in\initialSegment{n}$ by \cref{initial_segment_naturalsplus}.
                    Thus $k + 1\in\dom{u}$ by \cref{iteration_sequence}.
                    Thus $u$ is a function by \cref{iteration_sequence}.
                    Thus $(k + 1,u(k + 1))\in u$ by \cref{function_apply_elim}.
                    Thus $(k + 1,u(k + 1))\in v$ by \cref{union_intro_left}.
                    Thus $k + 1\in\initialSegment{n + 1}$ by \cref{initial_segment_naturalsplus}.
                    Thus $k + 1\in\dom{v}$ by \cref{dom}.
                    Thus $(k + 1,v(k + 1))\in v$ by \cref{function_apply_elim}.
                    Thus $u(k + 1) = v(k + 1)$ by \cref{function_rightunique}.
                    Thus $v(k + 1) = u(k + 1)$.
                    Thus $u(k + 1) = \apply{f}{u(k)}$.
                    Thus $v(k + 1) = \apply{f}{u(k)}$.
                    Show that $k \leq n$.
                    \begin{subproof}
                        Show that $k < k + 1$.
                        \begin{subproof}
                            $k\in\reals$ by \cref{real_naturals_subset_reals,subseteq}.
                            $0\in\reals$ by \cref{real_add_ident}.
                            $1\in\reals$ by \cref{real_naturals_subset_reals,real_one_in_naturals,subseteq}.
                            $0 < 1$ by \cref{real_zero_lt_one}.
                            $0 + k < 1 + k$ by \cref{real_lt_add}.
                            $0 + k = k + 0$ by \cref{real_add_comm}.
                            $1 + k = k + 1$ by \cref{real_add_comm}.
                            Thus $k + 0 < k + 1$ by \cref{real_lt_subst_left}.
                            $k + 0 = k$ by \cref{real_add_ident}.
                            Thus $k < k + 1$ by \cref{real_lt_subst_left}.
                        \end{subproof}
                        Thus $k < k + 1$ or $k = k + 1$.
                        Thus $k \leq k + 1$.
                        $k\in\reals$ by \cref{real_naturals_subset_reals,subseteq}.
                        $k + 1\in\reals$ by \cref{real_naturals_subset_reals,subseteq}.
                        $n\in\reals$ by \cref{real_naturals_subset_reals,subseteq}.
                        Thus $k \leq n$ by \cref{real_leq_trans}.
                    \end{subproof}
                    Thus $k\in\initialSegment{n}$ by \cref{initial_segment_naturalsplus}.
                    Thus $k\in\dom{u}$ by \cref{iteration_sequence}.
                    Thus $(k,u(k))\in u$ by \cref{function_apply_elim}.
                    Thus $(k,u(k))\in v$ by \cref{union_intro_left}.
                    Thus $k\in\dom{v}$ by \cref{dom}.
                    Thus $(k,v(k))\in v$ by \cref{function_apply_elim}.
                    Thus $u(k) = v(k)$ by \cref{function_rightunique}.
                    Thus $v(k + 1) = \apply{f}{v(k)}$.
                \end{byCase}
            \end{subproof}
            Thus $v$ is an iteration sequence for $f$ on $n + 1$ starting at $a_0$ by \cref{iteration_sequence_intro}.
            Thus $n + 1\in\naturalsPlus$ by \cref{real_naturals_closed_succ}.
            Thus $n + 1\in S$ by \cref{iteration_set_intro}.
        \end{subproof}
        Thus for all $n\in\naturalsPlus$ we have $n\in S$ by \cref{real_naturals_induction}.
    \end{subproof}
    Let $F = \{u \mid \text{there exists $n\in\naturalsPlus$ such that
        $u$ is an iteration sequence for $f$ on $n$ starting at $a_0$} \}$.
    Show that for all $u\in F$ we have $u$ is a function.
    \begin{subproof}
        Fix $u$.
        Assume $u\in F$.
        Take $n\in\naturalsPlus$ such that
            $u$ is an iteration sequence for $f$ on $n$ starting at $a_0$.
        Thus $u$ is a function by \cref{iteration_sequence}.
    \end{subproof}
    Show that for all $u,v$ such that $u\in F$ and $v\in F$ we have $u\subseteq v$ or $v\subseteq u$.
    \begin{subproof}
        Fix $u$.
        Fix $v$.
        Assume $u\in F$ and $v\in F$.
        Take $m\in\naturalsPlus$ such that
            $u$ is an iteration sequence for $f$ on $m$ starting at $a_0$.
        Take $n\in\naturalsPlus$ such that
            $v$ is an iteration sequence for $f$ on $n$ starting at $a_0$.
        Thus $u$ is a function by \cref{iteration_sequence}.
        Thus $\dom{u} = \initialSegment{m}$ by \cref{iteration_sequence}.
        Thus $u(1) = a_0$ by \cref{iteration_sequence}.
        Thus for all $k\in\naturalsPlus$ such that $k + 1 \leq m$ we have $u(k + 1) = \apply{f}{u(k)}$
            by \cref{iteration_sequence}.
        Thus $v$ is a function by \cref{iteration_sequence}.
        Thus $\dom{v} = \initialSegment{n}$ by \cref{iteration_sequence}.
        Thus $v(1) = a_0$ by \cref{iteration_sequence}.
        Thus for all $k\in\naturalsPlus$ such that $k + 1 \leq n$ we have $v(k + 1) = \apply{f}{v(k)}$
            by \cref{iteration_sequence}.
        $m = n$ or $m \neq n$.
        \begin{byCase}
        \caseOf{$m = n$.}
            Thus $u = v$ by \cref{iteration_sequence_unique}.
            Thus $u\subseteq v$ or $v\subseteq u$ by \cref{subseteq_refl}.
        \caseOf{$m \neq n$.}
            $m\in\reals$ by \cref{real_naturals_subset_reals,subseteq}.
            $n\in\reals$ by \cref{real_naturals_subset_reals,subseteq}.
            Thus $m < n$ or $n < m$ by \cref{real_lt_trichotomy}.
            \begin{byCase}
        \caseOf{$m < n$.}
            Show that for all $k\in\initialSegment{m}$ we have $u(k) = v(k)$.
            \begin{subproof}
                Let $A_1 = \{k\in\naturalsPlus \mid \text{if $k\leq m$ then $u(k) = v(k)$}\}$.
                Show that for all $k\in\naturalsPlus$ we have $k\in A_1$.
                \begin{subproof}
                    Show that $A_1\subseteq\naturalsPlus$.
                    \begin{subproof}
                        Show that for all $k$ such that $k\in A_1$ we have $k\in\naturalsPlus$.
                        \begin{subproof}
                            Fix $k$.
                            Assume $k\in A_1$.
                            Then $k\in\naturalsPlus$.
                        \end{subproof}
                        Thus $A_1\subseteq\naturalsPlus$ by \cref{subseteq}.
                    \end{subproof}
                    $1\in\naturalsPlus$ by \cref{real_one_in_naturals}.
                    Show that $1\in A_1$.
                    \begin{subproof}
                        Show that if $1\leq m$ then $u(1) = v(1)$.
                        \begin{subproof}
                            Assume $1\leq m$.
                            Then $u(1) = a_0$ and $v(1) = a_0$.
                            Thus $u(1) = v(1)$.
                        \end{subproof}
                        Thus $1\in A_1$.
                    \end{subproof}
                    Show that for all $k\in A_1$ we have $k + 1\in A_1$.
                    \begin{subproof}
                        Fix $k$.
                        Assume $k\in A_1$.
                        Show that if $k + 1\leq m$ then $u(k + 1) = v(k + 1)$.
                        \begin{subproof}
                            Assume $k + 1\leq m$.
                            Then $k \leq m$.
                            Thus $u(k) = v(k)$.
                            Thus $u(k + 1) = \apply{f}{u(k)}$ and $v(k + 1) = \apply{f}{v(k)}$.
                            Thus $u(k + 1) = v(k + 1)$.
                        \end{subproof}
                        Thus $k + 1\in A_1$.
                    \end{subproof}
                    Thus for all $k\in\naturalsPlus$ we have $k\in A_1$ by \cref{real_naturals_induction}.
                \end{subproof}
                Show that for all $k\in\initialSegment{m}$ we have $u(k) = v(k)$.
                \begin{subproof}
                    Fix $k$.
                    Assume $k\in\initialSegment{m}$.
                    Thus $k\in\naturalsPlus$ and $k \leq m$ by \cref{initial_segment_naturalsplus}.
                    Thus $k\in A_1$.
                    Thus $u(k) = v(k)$.
                \end{subproof}
            \end{subproof}
            Show that $u\subseteq v$.
            \begin{subproof}
                Show that for all $w$ such that $w\in u$ we have $w\in v$.
                \begin{subproof}
                    Fix $w$.
                    Assume $w\in u$.
                    Thus $u$ is a function by \cref{iteration_sequence}.
                    Take $k\in\dom{u}$ such that $w = (k,u(k))$ by \cref{function_member_elim}.
                    Thus $\dom{u} = \initialSegment{m}$ by \cref{iteration_sequence}.
                    Thus $k\in\initialSegment{m}$.
                    Thus $u(k) = v(k)$.
                    Thus $w = (k,v(k))$.
                    Thus $k\in\naturalsPlus$ and $k \leq m$ by \cref{initial_segment_naturalsplus}.
                    Thus $m \leq n$.
                    $k\in\reals$ by \cref{real_naturals_subset_reals,subseteq}.
                    $m\in\reals$ by \cref{real_naturals_subset_reals,subseteq}.
                    $n\in\reals$ by \cref{real_naturals_subset_reals,subseteq}.
                    Thus $k \leq n$ by \cref{real_leq_trans}.
                    Thus $k\in\initialSegment{n}$ by \cref{initial_segment_naturalsplus}.
                    Thus $\dom{v} = \initialSegment{n}$ by \cref{iteration_sequence}.
                    Thus $k\in\dom{v}$.
                    Thus $v$ is a function by \cref{iteration_sequence}.
                    Thus $(k,v(k))\in v$ by \cref{function_apply_elim}.
                    Thus $w\in v$.
                \end{subproof}
                Thus $u\subseteq v$ by \cref{subseteq}.
            \end{subproof}
        \caseOf{$n < m$.}
            Show that for all $k\in\initialSegment{n}$ we have $v(k) = u(k)$.
            \begin{subproof}
                Let $A_2 = \{k\in\naturalsPlus \mid \text{if $k\leq n$ then $v(k) = u(k)$}\}$.
                Show that for all $k\in\naturalsPlus$ we have $k\in A_2$.
                \begin{subproof}
                    Show that $A_2\subseteq\naturalsPlus$.
                    \begin{subproof}
                        Show that for all $k$ such that $k\in A_2$ we have $k\in\naturalsPlus$.
                        \begin{subproof}
                            Fix $k$.
                            Assume $k\in A_2$.
                            Then $k\in\naturalsPlus$.
                        \end{subproof}
                        Thus $A_2\subseteq\naturalsPlus$ by \cref{subseteq}.
                    \end{subproof}
                    $1\in\naturalsPlus$ by \cref{real_one_in_naturals}.
                    Show that $1\in A_2$.
                    \begin{subproof}
                        Show that if $1\leq n$ then $v(1) = u(1)$.
                        \begin{subproof}
                            Assume $1\leq n$.
                            Then $v(1) = a_0$ and $u(1) = a_0$.
                            Thus $v(1) = u(1)$.
                        \end{subproof}
                        Thus $1\in A_2$.
                    \end{subproof}
                    Show that for all $k\in A_2$ we have $k + 1\in A_2$.
                    \begin{subproof}
                        Fix $k$.
                        Assume $k\in A_2$.
                        Show that if $k + 1\leq n$ then $v(k + 1) = u(k + 1)$.
                        \begin{subproof}
                            Assume $k + 1\leq n$.
                            $k\in\reals$ and $n\in\reals$ by \cref{real_naturals_subset_reals,subseteq}.
                            $0\in\reals$ by \cref{real_add_ident}.
                            $1\in\reals$ by \cref{real_naturals_subset_reals,real_one_in_naturals,subseteq}.
                            $0 < 1$ by \cref{real_zero_lt_one}.
                            $0 + k < 1 + k$ by \cref{real_lt_add}.
                            $0 + k = k + 0$ by \cref{real_add_comm}.
                            $1 + k = k + 1$ by \cref{real_add_comm}.
                            Thus $k + 0 < k + 1$ by \cref{real_lt_subst_left}.
                            $k + 0 = k$ by \cref{real_add_ident}.
                            Thus $k < k + 1$ by \cref{real_lt_subst_left}.
                            Thus $k < k + 1$ or $k = k + 1$.
                            Thus $k \leq k + 1$.
                            $k + 1\in\reals$ by \cref{real_add_closed}.
                            Thus $k \leq n$ by \cref{real_leq_trans}.
                            Thus $v(k) = u(k)$.
                            Thus $k\in\naturalsPlus$.
                            Thus $v(k + 1) = \apply{f}{v(k)}$ by \cref{iteration_sequence}.
                            Thus $n < m$ or $n = m$.
                            Thus $n \leq m$.
                            $k + 1\in\reals$ by \cref{real_add_closed}.
                            $n\in\reals$ and $m\in\reals$ by \cref{real_naturals_subset_reals,subseteq}.
                            Thus $k + 1 \leq m$ by \cref{real_leq_trans}.
                            Thus $u(k + 1) = \apply{f}{u(k)}$ by \cref{iteration_sequence}.
                            Thus $v(k + 1) = u(k + 1)$.
                        \end{subproof}
                        Thus $k + 1\in A_2$.
                    \end{subproof}
                    Thus for all $k\in\naturalsPlus$ we have $k\in A_2$ by \cref{real_naturals_induction}.
                \end{subproof}
                Show that for all $k\in\initialSegment{n}$ we have $v(k) = u(k)$.
                \begin{subproof}
                    Fix $k$.
                    Assume $k\in\initialSegment{n}$.
                    Thus $k\in\naturalsPlus$ and $k \leq n$ by \cref{initial_segment_naturalsplus}.
                    Thus $k\in A_2$.
                    Thus $v(k) = u(k)$.
                \end{subproof}
            \end{subproof}
            Show that $v\subseteq u$.
            \begin{subproof}
                Show that for all $w$ such that $w\in v$ we have $w\in u$.
                \begin{subproof}
                    Fix $w$.
                    Assume $w\in v$.
                    Thus $v$ is a function by \cref{iteration_sequence}.
                    Take $k\in\dom{v}$ such that $w = (k,v(k))$ by \cref{function_member_elim}.
                    Thus $\dom{v} = \initialSegment{n}$ by \cref{iteration_sequence}.
                    Thus $k\in\initialSegment{n}$.
                    Thus $v(k) = u(k)$.
                    Thus $w = (k,u(k))$.
                    Thus $k\in\naturalsPlus$ and $k \leq n$ by \cref{initial_segment_naturalsplus}.
                    Thus $n \leq m$.
                    $k\in\reals$ by \cref{real_naturals_subset_reals,subseteq}.
                    $n\in\reals$ by \cref{real_naturals_subset_reals,subseteq}.
                    $m\in\reals$ by \cref{real_naturals_subset_reals,subseteq}.
                    Thus $k \leq m$ by \cref{real_leq_trans}.
                    Thus $k\in\initialSegment{m}$ by \cref{initial_segment_naturalsplus}.
                    Thus $\dom{u} = \initialSegment{m}$ by \cref{iteration_sequence}.
                    Thus $k\in\dom{u}$.
                    Thus $u$ is a function by \cref{iteration_sequence}.
                    Thus $(k,u(k))\in u$ by \cref{function_apply_elim}.
                    Thus $w\in u$.
                \end{subproof}
                Thus $v\subseteq u$ by \cref{subseteq}.
            \end{subproof}
        \end{byCase}
        \end{byCase}
    \end{subproof}
    Let $u = \unions{F}$.
    Thus $F$ is a family of functions.
    Thus $u$ is a function by \cref{unions_of_compatible_family_of_function_is_function}.
    Show that $\dom{u} = \naturalsPlus$.
    \begin{subproof}
        Show that $\dom{u}\subseteq\naturalsPlus$.
        \begin{subproof}
            Show that for all $x$ such that $x\in\dom{u}$ we have $x\in\naturalsPlus$.
            \begin{subproof}
                Fix $x$.
                Assume $x\in\dom{u}$.
                Take $y$ such that $(x,y)\in u$ by \cref{dom}.
                Take $v\in F$ such that $(x,y)\in v$ by \cref{unions_iff}.
                Take $n\in\naturalsPlus$ such that
                    $v$ is an iteration sequence for $f$ on $n$ starting at $a_0$.
                Thus $x\in\dom{v}$ by \cref{dom}.
                Thus $\dom{v} = \initialSegment{n}$ by \cref{iteration_sequence}.
                Thus $x\in\initialSegment{n}$.
                Thus $x\in\naturalsPlus$ by \cref{initial_segment_naturalsplus}.
            \end{subproof}
            Thus $\dom{u}\subseteq\naturalsPlus$ by \cref{subseteq}.
        \end{subproof}
        Show that $\naturalsPlus\subseteq\dom{u}$.
        \begin{subproof}
            Show that for all $n\in\naturalsPlus$ we have $n\in\dom{u}$.
            \begin{subproof}
                Fix $n\in\naturalsPlus$.
                $n\in S$ by \cref{subseteq}.
                Take $v$ such that
                    $v$ is an iteration sequence for $f$ on $n$ starting at $a_0$.
                Thus $v\in F$.
                Thus $\dom{v} = \initialSegment{n}$ by \cref{iteration_sequence}.
                Thus $n\in\dom{v}$ by \cref{initial_segment_naturalsplus}.
                Thus $n\in\dom{u}$ by \cref{dom,unions_iff}.
            \end{subproof}
            Thus $\naturalsPlus\subseteq\dom{u}$ by \cref{subseteq}.
        \end{subproof}
        Thus $\dom{u} = \naturalsPlus$ by \cref{subseteq_antisymmetric}.
    \end{subproof}
    Show that for all $n\in\naturalsPlus$ we have $u(n)\in A$.
    \begin{subproof}
        Fix $n\in\naturalsPlus$.
        Thus $n\in\dom{u}$.
        Thus $(n,u(n))\in u$ by \cref{function_apply_elim}.
        Thus $(n,u(n))\in\unions{F}$.
        Take $v\in F$ such that $(n,u(n))\in v$ by \cref{unions_iff}.
        Thus $n\in\dom{v}$ by \cref{dom}.
        Take $n_0\in\naturalsPlus$ such that
            $v$ is an iteration sequence for $f$ on $n_0$ starting at $a_0$.
        Thus $\dom{v} = \initialSegment{n_0}$ by \cref{iteration_sequence}.
        Thus $n\in\initialSegment{n_0}$.
        Thus $n \leq n_0$ by \cref{initial_segment_naturalsplus}.
        Thus $v(1) = a_0$ by \cref{iteration_sequence}.
        Thus for all $k\in\naturalsPlus$ such that $k + 1 \leq n_0$ we have $v(k + 1) = \apply{f}{v(k)}$
            by \cref{iteration_sequence}.
        Show that $v(n)\in A$.
        \begin{subproof}
            Let $B = \{k\in\naturalsPlus \mid \text{if $k\leq n$ then $v(k)\in A$}\}$.
            Show that for all $k\in\naturalsPlus$ we have $k\in B$.
            \begin{subproof}
                Show that $B\subseteq\naturalsPlus$.
                \begin{subproof}
                    Show that for all $k$ such that $k\in B$ we have $k\in\naturalsPlus$.
                    \begin{subproof}
                        Fix $k$.
                        Assume $k\in B$.
                        Then $k\in\naturalsPlus$.
                    \end{subproof}
                    Thus $B\subseteq\naturalsPlus$ by \cref{subseteq}.
                \end{subproof}
                $1\in\naturalsPlus$ by \cref{real_one_in_naturals}.
                Show that $1\in B$.
                \begin{subproof}
                    Show that if $1\leq n$ then $v(1)\in A$.
                    \begin{subproof}
                        Assume $1\leq n$.
                        Then $v(1) = a_0$.
                        Thus $v(1)\in A$.
                    \end{subproof}
                    Thus $1\in B$.
                \end{subproof}
                Show that for all $k\in B$ we have $k + 1\in B$.
                \begin{subproof}
                    Fix $k$.
                    Assume $k\in B$.
                    Show that if $k + 1\leq n$ then $v(k + 1)\in A$.
                    \begin{subproof}
                        Assume $k + 1\leq n$.
                        Then $k \leq n$.
                        Thus $v(k)\in A$.
                        $f\in\funs{A}{A}$.
                        Thus $\apply{f}{v(k)}\in A$ by \cref{funs_type_apply}.
                        $k + 1\in\naturalsPlus$ by \cref{real_naturals_closed_succ}.
                        $k + 1\in\reals$ by \cref{real_naturals_subset_reals,subseteq}.
                        $n\in\reals$ by \cref{real_naturals_subset_reals,subseteq}.
                        $n_0\in\reals$ by \cref{real_naturals_subset_reals,subseteq}.
                        Thus $k + 1 \leq n_0$ by \cref{real_leq_trans}.
                        Thus $v(k + 1) = \apply{f}{v(k)}$ by \cref{iteration_sequence}.
                        Thus $v(k + 1)\in A$.
                    \end{subproof}
                    Thus $k + 1\in B$.
                \end{subproof}
                Thus for all $k\in\naturalsPlus$ we have $k\in B$ by \cref{real_naturals_induction}.
            \end{subproof}
            Thus $v(n)\in A$ by \cref{initial_segment_naturalsplus}.
        \end{subproof}
        Thus $u(n)\in A$ by \cref{function_apply_intro}.
    \end{subproof}
    Show that $u(1) = a_0$.
    \begin{subproof}
        $1\in\naturalsPlus$ by \cref{real_one_in_naturals}.
        Thus $1\in\dom{u}$.
        Thus $(1,u(1))\in u$ by \cref{function_apply_elim}.
        Take $v\in F$ such that $(1,u(1))\in v$ by \cref{unions_iff}.
        Take $n_0\in\naturalsPlus$ such that
            $v$ is an iteration sequence for $f$ on $n_0$ starting at $a_0$.
        Thus $v$ is a function by \cref{iteration_sequence}.
        Thus $1\in\dom{v}$ by \cref{dom}.
        Thus $v(1) = a_0$ by \cref{iteration_sequence}.
        Thus $u(1) = v(1)$ by \cref{function_apply_intro}.
        Thus $u(1) = a_0$.
    \end{subproof}
    Show that for all $n\in\naturalsPlus$ we have $u(n + 1) = \apply{f}{u(n)}$.
    \begin{subproof}
        Fix $n\in\naturalsPlus$.
        Thus $n + 1\in\naturalsPlus$ by \cref{real_naturals_closed_succ}.
        Thus $n + 1\in\dom{u}$.
        Thus $(n + 1,u(n + 1))\in u$ by \cref{function_apply_elim}.
        Take $v\in F$ such that $(n + 1,u(n + 1))\in v$ by \cref{unions_iff}.
        Take $n_0\in\naturalsPlus$ such that
            $v$ is an iteration sequence for $f$ on $n_0$ starting at $a_0$.
        Thus $v$ is a function by \cref{iteration_sequence}.
        Thus $n + 1\in\dom{v}$ by \cref{dom}.
        Thus $\dom{v} = \initialSegment{n_0}$ by \cref{iteration_sequence}.
        Thus $n + 1\in\initialSegment{n_0}$ by \cref{subseteq}.
        Thus $n + 1 \leq n_0$ by \cref{initial_segment_naturalsplus}.
        Thus $v(n + 1) = \apply{f}{v(n)}$ by \cref{iteration_sequence}.
        Thus $u(n + 1) = v(n + 1)$ by \cref{function_apply_intro}.
        Thus $u(n + 1) = \apply{f}{v(n)}$.
        $0 < 1$ by \cref{real_zero_lt_one}.
        $n\in\reals$ by \cref{real_naturals_subset_reals,subseteq}.
        $0\in\reals$ by \cref{real_add_ident}.
        $1\in\reals$ by \cref{real_mul_ident}.
        Thus $0 + n < 1 + n$ by \cref{real_lt_add}.
        $0 + n = n + 0$ by \cref{real_add_comm}.
        $1 + n = n + 1$ by \cref{real_add_comm}.
        Thus $n + 0 < n + 1$ by \cref{real_lt_subst_left}.
        $n + 0 = n$ by \cref{real_add_ident}.
        Thus $n < n + 1$ by \cref{real_lt_subst_left}.
        Thus $n \leq n + 1$.
        $n + 1\in\reals$ by \cref{real_naturals_subset_reals,subseteq}.
        $n_0\in\reals$ by \cref{real_naturals_subset_reals,subseteq}.
        Thus $n \leq n_0$ by \cref{real_leq_trans}.
        Thus $n\in\initialSegment{n_0}$ by \cref{initial_segment_naturalsplus}.
        Thus $n\in\dom{v}$ by \cref{iteration_sequence,subseteq}.
        Thus $(n,v(n))\in v$ by \cref{function_apply_elim}.
        Thus $(n,v(n))\in u$ by \cref{unions_iff}.
        Thus $u(n) = v(n)$ by \cref{function_apply_intro}.
        Thus $u(n + 1) = \apply{f}{u(n)}$.
    \end{subproof}
    Thus $u\in\funs{\naturalsPlus}{A}$ by \cref{funs_intro}.
    Thus $u$ is a sequence in $A$ by \cref{sequence_in}.
    Thus there exists $u$ such that
        $u$ is a sequence in $A$ and
        $u(1) = a_0$ and
        for all $n\in\naturalsPlus$ we have $u(n + 1) = \apply{f}{u(n)}$.
\end{proof}

% from §28 helper lemma 10: sequential compactness implies finite epsilon-ball cover
\begin{lemma}\label{sequentially_compact_finite_ball_cover}
    Suppose $d$ is a metric on $X$.
    Let $T = \metricTopology{d}{X}$.
    Suppose $X$ is sequentially compact with respect to $T$.
    Suppose $\epsilon\in\reals$.
    Suppose $0 < \epsilon$.
    Then there exists $B$ such that
        $B$ is finite and
        $B$ is a covering of $X$ and
        for all $U\in B$ there exists $x\in X$ such that $U = \metricBall{d}{X}{x}{\epsilon}$.
\end{lemma}
\begin{proof}
    \begin{byCase}
    \caseOf{$X = \emptyset$.}
        Let $B = \emptyset$.
        $B$ is finite by \cref{emptyset_is_finite}.
        Show that for all $U\in B$ we have $\exists x\in X. U = \metricBall{d}{X}{x}{\epsilon}$.
        \begin{subproof}
            Fix $U$.
            Assume $U\in B$.
            Suppose not.
            Contradiction by \cref{emptyset}.
        \end{subproof}
        $X\subseteq\unions{B}$ by \cref{emptyset_subseteq,unions_emptyset}.
        Thus $B$ is a covering of $X$ by \cref{covering}.
        Thus there exists $B$ such that
            $B$ is finite and
            $B$ is a covering of $X$ and
            for all $U\in B$ there exists $x\in X$ such that $U = \metricBall{d}{X}{x}{\epsilon}$.
    \caseOf{$X \neq \emptyset$.}
        Suppose not.
        Then for all $B$ such that
            $B$ is finite and
            for all $U\in B$ we have $\exists x\in X. U = \metricBall{d}{X}{x}{\epsilon}$
        we have $B$ is not a covering of $X$.
        Let $F = \{B\in\pow{X} \mid \text{$B$ is finite}\}$.
        Show that for all $B\in F$ there exists $x\in X$ such that
            for all $y\in B$ we have $x\notin\metricBall{d}{X}{y}{\epsilon}$.
        \begin{subproof}
            Fix $B$.
            Assume $B\in F$.
            Thus $B\in\pow{X}$ and $B$ is finite.
            Thus $B\subseteq X$ by \cref{powerset_elim}.
            Let $g = \{(x,\metricBall{d}{X}{x}{\epsilon}) \mid x\in B\}$.
            Show that $g$ is a function on $B$.
            \begin{subproof}
                Show that $g$ is a relation.
                \begin{subproof}
                    Show that for all $w$ such that $w\in g$ there exist $x,y$ such that $w = (x,y)$.
                    \begin{subproof}
                        Fix $w$.
                        Assume $w\in g$.
                        Take $x\in B$ such that $w = (x,\metricBall{d}{X}{x}{\epsilon})$.
                        Thus there exist $x,y$ such that $w = (x,y)$.
                    \end{subproof}
                    Thus $g$ is a relation by \cref{relation}.
                \end{subproof}
                Show that $g$ is right-unique.
                \begin{subproof}
                    Show that for all $x,U_1,U_2$ such that $(x,U_1)\in g$ and $(x,U_2)\in g$ we have $U_1 = U_2$.
                    \begin{subproof}
                        Fix $x$.
                        Fix $U_1$.
                        Fix $U_2$.
                        Assume $(x,U_1)\in g$ and $(x,U_2)\in g$.
                        Take $x_1\in B$ such that $(x,U_1) = (x_1,\metricBall{d}{X}{x_1}{\epsilon})$.
                        Take $x_2\in B$ such that $(x,U_2) = (x_2,\metricBall{d}{X}{x_2}{\epsilon})$.
                        Thus $x = x_1$ and $U_1 = \metricBall{d}{X}{x_1}{\epsilon}$ by \cref{pair_eq_iff}.
                        Thus $x = x_2$ and $U_2 = \metricBall{d}{X}{x_2}{\epsilon}$ by \cref{pair_eq_iff}.
                        Thus $U_1 = \metricBall{d}{X}{x}{\epsilon}$ and $U_2 = \metricBall{d}{X}{x}{\epsilon}$.
                        Thus $U_1 = U_2$.
                    \end{subproof}
                    Thus $g$ is right-unique by \cref{rightunique}.
                \end{subproof}
                Show that $\dom{g} = B$.
                \begin{subproof}
                    Show that $\dom{g}\subseteq B$.
                    \begin{subproof}
                        Show that for all $x$ such that $x\in\dom{g}$ we have $x\in B$.
                        \begin{subproof}
                            Fix $x$.
                            Assume $x\in\dom{g}$.
                            Take $U$ such that $(x,U)\in g$ by \cref{dom}.
                            Take $x_1\in B$ such that $(x,U) = (x_1,\metricBall{d}{X}{x_1}{\epsilon})$.
                            Thus $x = x_1$ by \cref{pair_eq_iff}.
                            Thus $x\in B$.
                        \end{subproof}
                        Thus $\dom{g}\subseteq B$ by \cref{subseteq}.
                    \end{subproof}
                    Show that $B\subseteq\dom{g}$.
                    \begin{subproof}
                        Show that for all $x$ such that $x\in B$ we have $x\in\dom{g}$.
                        \begin{subproof}
                            Fix $x$.
                            Assume $x\in B$.
                            Thus $(x,\metricBall{d}{X}{x}{\epsilon})\in g$.
                            Thus $x\in\dom{g}$ by \cref{dom}.
                        \end{subproof}
                        Thus $B\subseteq\dom{g}$ by \cref{subseteq}.
                    \end{subproof}
                    Thus $\dom{g} = B$ by \cref{subseteq_antisymmetric}.
                \end{subproof}
                Thus $g$ is a function on $B$.
            \end{subproof}
            Let $G = \img{g}{B}$.
            Thus $G$ is finite by \cref{finite_image_function}.
            Show that for all $U\in G$ we have $\exists x\in X. U = \metricBall{d}{X}{x}{\epsilon}$.
            \begin{subproof}
                Fix $U$.
                Assume $U\in G$.
                Take $x\in B$ such that $(x,U)\in g$ by \cref{img_iff}.
                Thus $x\in X$ by \cref{subseteq}.
                Take $x_1\in B$ such that $(x,U) = (x_1,\metricBall{d}{X}{x_1}{\epsilon})$.
                Thus $U = \metricBall{d}{X}{x_1}{\epsilon}$ by \cref{pair_eq_iff}.
                Thus $x_1\in X$ by \cref{subseteq}.
                Thus there exists $x\in X$ such that $U = \metricBall{d}{X}{x}{\epsilon}$.
            \end{subproof}
            Thus $G$ is not a covering of $X$.
            Thus $X\not\subseteq\unions{G}$ by \cref{covering}.
            Show that $\unions{G}\subseteq X$.
            \begin{subproof}
                Show that for all $U\in G$ we have $U\subseteq X$.
                \begin{subproof}
                    Fix $U$.
                    Assume $U\in G$.
                    Take $x\in X$ such that $U = \metricBall{d}{X}{x}{\epsilon}$.
                    Show that $U\subseteq X$.
                    \begin{subproof}
                        Show that for all $y\in U$ we have $y\in X$.
                        \begin{subproof}
                            Fix $y$.
                            Assume $y\in U$.
                            Thus $y\in\metricBall{d}{X}{x}{\epsilon}$.
                            Thus $y\in X$ by \cref{metric_ball}.
                        \end{subproof}
                        Thus $U\subseteq X$ by \cref{subseteq}.
                    \end{subproof}
                \end{subproof}
                Thus $\unions{G}\subseteq X$ by \cref{unions_family}.
            \end{subproof}
            Show that $\unions{G}\neq X$.
            \begin{subproof}
                Suppose not.
                Then $\unions{G} = X$.
                Thus $X\subseteq\unions{G}$ by \cref{subseteq_refl}.
                Contradiction.
            \end{subproof}
            Thus $\unions{G}\subset X$ by \cref{subset}.
            Take $x\in X$ such that $x\notin\unions{G}$ by \cref{subset_witness}.
            Show that for all $y\in B$ we have $x\notin\metricBall{d}{X}{y}{\epsilon}$.
            \begin{subproof}
                Fix $y$.
                Assume $y\in B$.
                $(y,\metricBall{d}{X}{y}{\epsilon})\in g$.
                Thus $g(y) = \metricBall{d}{X}{y}{\epsilon}$ by \cref{function_apply_iff}.
                Thus $\metricBall{d}{X}{y}{\epsilon}\in G$ by \cref{img_iff}.
                Suppose not.
                Then $x\in\metricBall{d}{X}{y}{\epsilon}$.
                Thus $x\in\unions{G}$ by \cref{unions_iff}.
                Contradiction.
            \end{subproof}
            Thus there exists $x\in X$ such that for all $y\in B$ we have $x\notin\metricBall{d}{X}{y}{\epsilon}$.
        \end{subproof}
        Let $R = \{ w\in F\times X \mid
            \text{for all $y\in\fst{w}$ we have $\snd{w}\notin\metricBall{d}{X}{y}{\epsilon}$} \}$.
        Show that $R$ is a relation.
        \begin{subproof}
            Show that for all $w$ such that $w\in R$ there exist $B,x$ such that $w = (B,x)$.
            \begin{subproof}
                Fix $w$.
                Assume $w\in R$.
                Thus $w\in F\times X$.
                Take $B,x$ such that $w = (B,x)$ by \cref{times_elem_is_tuple}.
            \end{subproof}
            Thus $R$ is a relation by \cref{relation}.
        \end{subproof}
        Show that for all $B\in F$ there exists $x$ such that $B\mathrel{R} x$.
        \begin{subproof}
            Fix $B\in F$.
            Take $x\in X$ such that for all $y\in B$ we have $x\notin\metricBall{d}{X}{y}{\epsilon}$.
            Thus $(B,x)\in F\times X$ by \cref{times_tuple_intro}.
            Thus for all $y\in\fst{(B,x)}$ we have $\snd{(B,x)}\notin\metricBall{d}{X}{y}{\epsilon}$
                by \cref{fst_eq,snd_eq}.
            Thus $(B,x)\in R$.
            Thus $B\mathrel{R} x$.
            Thus there exists $x$ such that $B\mathrel{R} x$.
        \end{subproof}
        Take $c$ such that $c$ is a function on $F$ and for all $B\in F$ we have $B\mathrel{R} \apply{c}{B}$ by \cref{choice_function}.
        Show that for all $B\in F$ we have $\apply{c}{B}\in X$ and for all $y\in B$ we have $\apply{c}{B}\notin\metricBall{d}{X}{y}{\epsilon}$.
        \begin{subproof}
            Fix $B\in F$.
            Thus $B\mathrel{R}\apply{c}{B}$.
            Thus $(B,\apply{c}{B})\in F\times X$ and
                for all $y\in B$ we have $\apply{c}{B}\notin\metricBall{d}{X}{y}{\epsilon}$.
            Thus $\apply{c}{B}\in X$ by \cref{times_tuple_elim}.
        \end{subproof}
        Show that $c\in\funs{F}{X}$.
        \begin{subproof}
            Show that for all $B\in\dom{c}$ we have $c(B)\in X$.
            \begin{subproof}
                Fix $B\in\dom{c}$.
                Thus $B\in F$.
                Thus $c(B)\in X$.
            \end{subproof}
            Thus $c\in\funs{F}{X}$ by \cref{funs_intro}.
        \end{subproof}
        Let $g = \{(B,B\union\{\apply{c}{B},\apply{c}{B}\}) \mid B\in F\}$.
        Show that $g\in\funs{F}{F}$.
        \begin{subproof}
            Show that $g$ is a function on $F$.
            \begin{subproof}
                Show that $g$ is a relation.
                \begin{subproof}
                    Show that for all $w$ such that $w\in g$ there exist $B,D$ such that $w = (B,D)$.
                    \begin{subproof}
                        Fix $w$.
                        Assume $w\in g$.
                        Take $B\in F$ such that $w = (B,B\union\{\apply{c}{B},\apply{c}{B}\})$.
                        Thus there exist $B,D$ such that $w = (B,D)$.
                    \end{subproof}
                    Thus $g$ is a relation by \cref{relation}.
                \end{subproof}
                Show that $g$ is right-unique.
                \begin{subproof}
                    Show that for all $B,D_1,D_2$ such that $(B,D_1)\in g$ and $(B,D_2)\in g$ we have $D_1 = D_2$.
                    \begin{subproof}
                        Fix $B$.
                        Fix $D_1$.
                        Fix $D_2$.
                        Assume $(B,D_1)\in g$ and $(B,D_2)\in g$.
                        Take $B_1\in F$ such that $(B,D_1) = (B_1,B_1\union\{\apply{c}{B_1},\apply{c}{B_1}\})$.
                        Take $B_2\in F$ such that $(B,D_2) = (B_2,B_2\union\{\apply{c}{B_2},\apply{c}{B_2}\})$.
                        Thus $B = B_1$ and $D_1 = B_1\union\{\apply{c}{B_1},\apply{c}{B_1}\}$ by \cref{pair_eq_iff}.
                        Thus $B = B_2$ and $D_2 = B_2\union\{\apply{c}{B_2},\apply{c}{B_2}\}$ by \cref{pair_eq_iff}.
                        Thus $D_1 = B\union\{\apply{c}{B},\apply{c}{B}\}$ and
                            $D_2 = B\union\{\apply{c}{B},\apply{c}{B}\}$.
                        Thus $D_1 = D_2$.
                    \end{subproof}
                    Thus $g$ is right-unique by \cref{rightunique}.
                \end{subproof}
                Show that $\dom{g} = F$.
                \begin{subproof}
                    Show that $\dom{g}\subseteq F$.
                    \begin{subproof}
                        Show that for all $B$ such that $B\in\dom{g}$ we have $B\in F$.
                        \begin{subproof}
                            Fix $B$.
                            Assume $B\in\dom{g}$.
                            Take $D$ such that $(B,D)\in g$ by \cref{dom}.
                            Take $B_1\in F$ such that $(B,D) = (B_1,B_1\union\{\apply{c}{B_1},\apply{c}{B_1}\})$.
                            Thus $B = B_1$ by \cref{pair_eq_iff}.
                            Thus $B\in F$.
                        \end{subproof}
                        Thus $\dom{g}\subseteq F$ by \cref{subseteq}.
                    \end{subproof}
                    Show that $F\subseteq\dom{g}$.
                    \begin{subproof}
                        Show that for all $B$ such that $B\in F$ we have $B\in\dom{g}$.
                        \begin{subproof}
                            Fix $B$.
                            Assume $B\in F$.
                            Thus $(B,B\union\{\apply{c}{B},\apply{c}{B}\})\in g$.
                            Thus $B\in\dom{g}$ by \cref{dom}.
                        \end{subproof}
                        Thus $F\subseteq\dom{g}$ by \cref{subseteq}.
                    \end{subproof}
                    Thus $\dom{g} = F$ by \cref{subseteq_antisymmetric}.
                \end{subproof}
                Thus $g$ is a function on $F$.
            \end{subproof}
            Show that for all $B\in\dom{g}$ we have $g(B)\in F$.
            \begin{subproof}
                Fix $B\in\dom{g}$.
                Thus $B\in F$.
                Thus $(B,B\union\{\apply{c}{B},\apply{c}{B}\})\in g$.
                Thus $g(B) = B\union\{\apply{c}{B},\apply{c}{B}\}$ by \cref{function_apply_intro}.
                Thus $B$ is finite and $B\subseteq X$ by \cref{powerset_elim}.
                Thus $c(B)\in X$.
                Show that $\{c(B),c(B)\}\subseteq X$.
                \begin{subproof}
                    Show that for all $y$ such that $y\in\{c(B),c(B)\}$ we have $y\in X$.
                    \begin{subproof}
                        Fix $y$.
                        Assume $y\in\{c(B),c(B)\}$.
                        Thus $y = c(B)$ by \cref{upair_elim}.
                        Thus $y\in X$.
                    \end{subproof}
                    Thus $\{c(B),c(B)\}\subseteq X$ by \cref{subseteq}.
                \end{subproof}
                Thus $\{c(B),c(B)\}$ is finite by \cref{pair_is_finite}.
                Thus $B\union\{c(B),c(B)\}$ is finite by \cref{union_of_finite_is_finite}.
                Thus $B\union\{c(B),c(B)\}\subseteq X$ by \cref{union_subsets_is_subset}.
                Thus $B\union\{c(B),c(B)\}\in\pow{X}$ by \cref{powerset_intro}.
                Thus $g(B)\in F$.
            \end{subproof}
            Thus $g\in\funs{F}{F}$ by \cref{funs_intro}.
        \end{subproof}
        Take $x_0$ such that $x_0\in X$ by \cref{nonempty_inhabited}.
        Let $a_0 = \{x_0,x_0\}$.
        $a_0$ is finite by \cref{pair_is_finite}.
        Show that $a_0\subseteq X$.
        \begin{subproof}
            Show that for all $y$ such that $y\in a_0$ we have $y\in X$.
            \begin{subproof}
                Fix $y$.
                Assume $y\in a_0$.
                Thus $y = x_0$ by \cref{upair_elim}.
                Thus $y\in X$.
            \end{subproof}
            Thus $a_0\subseteq X$ by \cref{subseteq}.
        \end{subproof}
        Thus $a_0\in\pow{X}$ by \cref{powerset_intro}.
        Thus $a_0\in F$.
        Take $u$ such that
            $u$ is a sequence in $F$ and
            $u(1) = a_0$ and
            for all $n\in\naturalsPlus$ we have $u(n + 1) = \apply{g}{u(n)}$ by \cref{iteration_sequence_naturalsplus}.
        Let $s = c\circ u$.
        Show that $s$ is a sequence in $X$.
        \begin{subproof}
            $u\in\funs{\naturalsPlus}{F}$ by \cref{sequence_in}.
            Show that $\ran{u}\subseteq F$.
            \begin{subproof}
                Show that for all $z$ such that $z\in\ran{u}$ we have $z\in F$.
                \begin{subproof}
                    Fix $z$.
                    Assume $z\in\ran{u}$.
                    Take $w$ such that $w\in u$ and there exists $n$ such that $w = (n,z)$ by \cref{ran}.
                    Take $n$ such that $w = (n,z)$.
                    Thus $(n,z)\in u$.
                    Thus $n\in\dom{u}$ by \cref{dom}.
                    Thus $u\in\rels{\naturalsPlus}{F}$ by \cref{funs}.
                    Thus $u\in\pow{\naturalsPlus\times F}$ by \cref{rels}.
                    Thus $u\subseteq\naturalsPlus\times F$ by \cref{powerset_elim}.
                    Show that $u$ is a relation.
                    \begin{subproof}
                        Show that for all $w_1$ such that $w_1\in u$ there exist $x,y$ such that $w_1 = (x,y)$.
                        \begin{subproof}
                            Fix $w_1$.
                            Assume $w_1\in u$.
                            Thus $w_1\in\naturalsPlus\times F$ by \cref{subseteq}.
                            Take $x,y$ such that $w_1 = (x,y)$ by \cref{times_elem_is_tuple}.
                        \end{subproof}
                        Thus $u$ is a relation by \cref{relation}.
                    \end{subproof}
                    Thus $u$ is right-unique by \cref{funs}.
                    Thus $u(n) = z$ by \cref{function_apply_intro}.
                    Thus $\dom{u} = \naturalsPlus$ by \cref{funs}.
                    Thus $n\in\naturalsPlus$ by \cref{subseteq}.
                    Thus $u(n)\in F$ by \cref{funs_type_apply}.
                    Thus $z\in F$.
                \end{subproof}
                Thus $\ran{u}\subseteq F$ by \cref{subseteq}.
            \end{subproof}
            Thus $\dom{c} = F$ by \cref{funs}.
            Thus $\ran{u}\subseteq\dom{c}$ by \cref{subseteq}.
            Thus $u\in\rels{\naturalsPlus}{F}$ by \cref{funs}.
            Thus $u\in\pow{\naturalsPlus\times F}$ by \cref{rels}.
            Thus $u\subseteq\naturalsPlus\times F$ by \cref{powerset_elim}.
            Show that $u$ is a relation.
            \begin{subproof}
                Show that for all $w_1$ such that $w_1\in u$ there exist $x,y$ such that $w_1 = (x,y)$.
                \begin{subproof}
                    Fix $w_1$.
                    Assume $w_1\in u$.
                    Thus $w_1\in\naturalsPlus\times F$ by \cref{subseteq}.
                    Take $x,y$ such that $w_1 = (x,y)$ by \cref{times_elem_is_tuple}.
                \end{subproof}
                Thus $u$ is a relation by \cref{relation}.
            \end{subproof}
            Thus $u$ is right-unique by \cref{funs}.
            Thus $c\in\rels{F}{X}$ by \cref{funs}.
            Thus $c\in\pow{F\times X}$ by \cref{rels}.
            Thus $c\subseteq F\times X$ by \cref{powerset_elim}.
            Show that $c$ is a relation.
            \begin{subproof}
                Show that for all $w_2$ such that $w_2\in c$ there exist $x,y$ such that $w_2 = (x,y)$.
                \begin{subproof}
                    Fix $w_2$.
                    Assume $w_2\in c$.
                    Thus $w_2\in F\times X$ by \cref{subseteq}.
                    Take $x,y$ such that $w_2 = (x,y)$ by \cref{times_elem_is_tuple}.
                \end{subproof}
                Thus $c$ is a relation by \cref{relation}.
            \end{subproof}
            Thus $c$ is right-unique by \cref{funs}.
            Thus $s$ is a function by \cref{function_circ}.
            Thus $\dom{s} = \preimg{u}{\dom{c}}$ by \cref{dom_of_circ}.
            Show that $\preimg{u}{\dom{c}} = \dom{u}$.
            \begin{subproof}
                Show that $\preimg{u}{\dom{c}}\subseteq\dom{u}$.
                \begin{subproof}
                    Thus $\preimg{u}{\dom{c}}\subseteq\dom{u}$ by \cref{preimg_subseteq_dom}.
                \end{subproof}
                Show that $\dom{u}\subseteq\preimg{u}{\dom{c}}$.
                \begin{subproof}
                    Show that for all $n$ such that $n\in\dom{u}$ we have $n\in\preimg{u}{\dom{c}}$.
                    \begin{subproof}
                        Fix $n$.
                        Assume $n\in\dom{u}$.
                        Thus $\dom{u} = \naturalsPlus$ by \cref{funs}.
                        Thus $n\in\naturalsPlus$ by \cref{subseteq}.
                        Thus $u(n)\in F$ by \cref{funs_type_apply}.
                        Thus $u(n)\in\dom{c}$ by \cref{subseteq}.
                        Thus $(n,u(n))\in u$ by \cref{function_apply_elim}.
                        Thus $n\in\preimg{u}{\dom{c}}$ by \cref{preimg_iff}.
                    \end{subproof}
                    Thus $\dom{u}\subseteq\preimg{u}{\dom{c}}$ by \cref{subseteq}.
                \end{subproof}
                Thus $\preimg{u}{\dom{c}} = \dom{u}$ by \cref{subseteq_antisymmetric}.
            \end{subproof}
            Thus $\dom{s} = \dom{u}$.
            Thus $\dom{s} = \naturalsPlus$ by \cref{funs}.
            Show that for all $n\in\dom{s}$ we have $s(n)\in X$.
            \begin{subproof}
                Fix $n\in\dom{s}$.
                Thus $n\in\naturalsPlus$.
                Thus $n\in\dom{u}$.
                Thus $u(n)\in F$ by \cref{funs_type_apply}.
                Thus $c(u(n))\in X$ by \cref{funs_type_apply}.
                Thus $s(n)\in X$ by \cref{circ_apply}.
            \end{subproof}
            Thus $s\in\funs{\naturalsPlus}{X}$ by \cref{funs_intro}.
            Thus $s$ is a sequence in $X$ by \cref{sequence_in}.
        \end{subproof}
        Show that for all $n\in\naturalsPlus$ we have
            for all $y\in u(n)$ we have $s(n)\notin\metricBall{d}{X}{y}{\epsilon}$.
        \begin{subproof}
            Fix $n\in\naturalsPlus$.
            Fix $y$.
            Assume $y\in u(n)$.
            Thus $u(n)\in F$ by \cref{funs_type_apply,sequence_in}.
            Thus $u(n)\mathrel{R} c(u(n))$.
            Thus for all $z\in u(n)$ we have $c(u(n))\notin\metricBall{d}{X}{z}{\epsilon}$.
            $u\in\funs{\naturalsPlus}{F}$ by \cref{sequence_in}.
            Thus $u$ is a function by \cref{funs_is_function}.
            Thus $u$ is right-unique and $u$ is a relation.
            $c\in\funs{F}{X}$ by \cref{funs}.
            Thus $c$ is a function by \cref{funs_is_function}.
            Thus $c$ is right-unique and $c$ is a relation.
            Show that $\ran{u}\subseteq F$.
            \begin{subproof}
                Show that for all $z$ such that $z\in\ran{u}$ we have $z\in F$.
                \begin{subproof}
                    Fix $z$.
                    Assume $z\in\ran{u}$.
                    Take $w$ such that $w\in u$ and there exists $n$ such that $w = (n,z)$ by \cref{ran}.
                    Take $n$ such that $w = (n,z)$.
                    Thus $(n,z)\in u$.
                    Thus $n\in\dom{u}$ by \cref{dom}.
                    Thus $u$ is right-unique and $u$ is a relation.
                    Thus $u(n) = z$ by \cref{function_apply_intro}.
                    Thus $\dom{u} = \naturalsPlus$ by \cref{funs}.
                    Thus $n\in\naturalsPlus$ by \cref{subseteq}.
                    Thus $u(n)\in F$ by \cref{funs_type_apply}.
                    Thus $z\in F$.
                \end{subproof}
                Thus $\ran{u}\subseteq F$ by \cref{subseteq}.
            \end{subproof}
            Thus $\dom{c} = F$ by \cref{funs}.
            Thus $\ran{u}\subseteq\dom{c}$ by \cref{subseteq}.
            Thus $\dom{u} = \naturalsPlus$ by \cref{funs}.
            Thus $n\in\dom{u}$ by \cref{subseteq}.
            Thus $s(n) = c(u(n))$ by \cref{circ_apply}.
            Thus $s(n)\notin\metricBall{d}{X}{y}{\epsilon}$.
        \end{subproof}
        Show that for all $n\in\naturalsPlus$ we have $u(n)\subseteq u(n + 1)$.
        \begin{subproof}
            Fix $n\in\naturalsPlus$.
            Thus $u(n + 1) = g(u(n))$.
            Thus $u(n)\in F$ by \cref{sequence_in,funs_type_apply}.
            Thus $(u(n),u(n)\union\{c(u(n)),c(u(n))\})\in g$.
            Thus $g$ is a function.
            Thus $g(u(n)) = u(n)\union\{c(u(n)),c(u(n))\}$ by \cref{function_apply_intro}.
            Thus $u(n + 1) = u(n)\union\{c(u(n)),c(u(n))\}$.
            Thus $u(n)\subseteq u(n + 1)$ by \cref{union_upper_left}.
        \end{subproof}
        Show that for all $n\in\naturalsPlus$ we have $s(n)\in u(n + 1)$.
        \begin{subproof}
            Fix $n\in\naturalsPlus$.
            Thus $u(n + 1) = g(u(n))$.
            Thus $u(n)\in F$ by \cref{sequence_in,funs_type_apply}.
            Thus $(u(n),u(n)\union\{c(u(n)),c(u(n))\})\in g$.
            Thus $g$ is a function.
            Thus $g(u(n)) = u(n)\union\{c(u(n)),c(u(n))\}$ by \cref{function_apply_intro}.
            Thus $u(n + 1) = u(n)\union\{c(u(n)),c(u(n))\}$.
            $c(u(n))\in\{c(u(n)),c(u(n))\}$ by \cref{upair_intro_left}.
            Thus $c(u(n))\in u(n + 1)$ by \cref{union_iff}.
            $u\in\funs{\naturalsPlus}{F}$ by \cref{sequence_in}.
            $c\in\funs{F}{X}$ by \cref{funs}.
            Thus $u$ is right-unique by \cref{funs}.
            Thus $c$ is right-unique by \cref{funs}.
            Thus $s$ is a function by \cref{function_circ}.
            Thus $\dom{u} = \naturalsPlus$ by \cref{funs}.
            Thus $n\in\dom{u}$ by \cref{subseteq}.
            Thus $u\in\rels{\naturalsPlus}{F}$ by \cref{funs}.
            Thus $u\in\pow{\naturalsPlus\times F}$ by \cref{rels}.
            Thus $u\subseteq\naturalsPlus\times F$ by \cref{powerset_elim}.
            Show that $u$ is a relation.
            \begin{subproof}
                Show that for all $w_1$ such that $w_1\in u$ there exist $x,y$ such that $w_1 = (x,y)$.
                \begin{subproof}
                    Fix $w_1$.
                    Assume $w_1\in u$.
                    Thus $w_1\in\naturalsPlus\times F$ by \cref{subseteq}.
                    Take $x,y$ such that $w_1 = (x,y)$ by \cref{times_elem_is_tuple}.
                \end{subproof}
                Thus $u$ is a relation by \cref{relation}.
            \end{subproof}
            Thus $(n,u(n))\in u$ by \cref{function_apply_elim}.
            Thus $u(n)\in F$ by \cref{funs_type_apply}.
            Thus $\dom{c} = F$ by \cref{funs}.
            Thus $u(n)\in\dom{c}$ by \cref{subseteq}.
            Thus $(u(n),c(u(n)))\in c$ by \cref{function_apply_elim}.
            Thus $(n,c(u(n)))\in c\circ u$ by \cref{circ_elem_intro}.
            Thus $(n,c(u(n)))\in s$.
            Thus $s(n) = c(u(n))$ by \cref{function_apply_intro}.
            Thus $s(n)\in u(n + 1)$.
        \end{subproof}
        Show that for all $n\in\naturalsPlus$ we have
            for all $m\in\naturalsPlus$ such that $m < n$ we have $s(m)\in u(n)$.
        \begin{subproof}
            Let $P = \{n\in\naturalsPlus \mid \text{for all $m\in\naturalsPlus$ such that $m < n$ we have $s(m)\in u(n)$}\}$.
            Show that for all $n\in\naturalsPlus$ we have $n\in P$.
            \begin{subproof}
                Show that $P\subseteq\naturalsPlus$.
                \begin{subproof}
                    Show that for all $n$ such that $n\in P$ we have $n\in\naturalsPlus$.
                    \begin{subproof}
                        Fix $n$.
                        Assume $n\in P$.
                        Then $n\in\naturalsPlus$.
                    \end{subproof}
                    Thus $P\subseteq\naturalsPlus$ by \cref{subseteq}.
                \end{subproof}
                $1\in\naturalsPlus$ by \cref{real_one_in_naturals}.
                Show that $1\in P$.
                \begin{subproof}
                    Show that for all $m\in\naturalsPlus$ such that $m < 1$ we have $s(m)\in u(1)$.
                    \begin{subproof}
                        Fix $m$.
                        Assume $m\in\naturalsPlus$ and $m < 1$.
                        Suppose not.
                        $m\in\reals$ by \cref{real_naturals_subset_reals,subseteq}.
                        $1\in\reals$ by \cref{real_naturals_subset_reals,real_one_in_naturals,subseteq}.
                        $1 < m$ or $1 = m$ by \cref{real_one_leq_naturalsplus}.
                        \begin{byCase}
                        \caseOf{$1 = m$.}
                            Thus $m < m$ by \cref{real_lt_subst_right}.
                            Contradiction by \cref{real_lt_irreflexive}.
                        \caseOf{$1 < m$.}
                            Thus $1 < 1$ by \cref{real_lt_trans}.
                            Contradiction by \cref{real_lt_irreflexive}.
                        \end{byCase}
                    \end{subproof}
                    Thus $1\in P$.
                \end{subproof}
                Show that for all $n\in P$ we have $n + 1\in P$.
                \begin{subproof}
                    Fix $n$.
                    Assume $n\in P$.
                    Show that for all $m\in\naturalsPlus$ such that $m < n + 1$ we have $s(m)\in u(n + 1)$.
                    \begin{subproof}
                        Fix $m$.
                        Assume $m\in\naturalsPlus$ and $m < n + 1$.
                        Thus $m \neq n + 1$.
                        Thus $m \leq n$ by \cref{real_naturals_leq_succ_elim}.
                        $m < n$ or $m = n$ by \cref{real_leq_split}.
                        \begin{byCase}
                        \caseOf{$m < n$.}
                            Thus $s(m)\in u(n)$.
                            Thus $u(n)\subseteq u(n + 1)$.
                            Thus $s(m)\in u(n + 1)$ by \cref{subseteq}.
                        \caseOf{$m = n$.}
                            Thus $s(n)\in u(n + 1)$.
                        \end{byCase}
                    \end{subproof}
                    Thus $n + 1\in P$.
                \end{subproof}
                Thus for all $n\in\naturalsPlus$ we have $n\in P$ by \cref{real_naturals_induction}.
            \end{subproof}
            Thus for all $n\in\naturalsPlus$ we have
                for all $m\in\naturalsPlus$ such that $m < n$ we have $s(m)\in u(n)$ by \cref{subseteq}.
        \end{subproof}
        Show that for all $m\in\naturalsPlus$ we have
            for all $n\in\naturalsPlus$ such that $m < n$ we have
                it is wrong that $d((s(m),s(n))) < \epsilon$.
        \begin{subproof}
            Fix $m\in\naturalsPlus$.
            Fix $n\in\naturalsPlus$.
            Assume $m < n$.
            Thus $s(m)\in u(n)$.
            Thus $s(n)\notin\metricBall{d}{X}{s(m)}{\epsilon}$.
            Suppose not.
            Then $d((s(m),s(n))) < \epsilon$.
            $s\in\funs{\naturalsPlus}{X}$ by \cref{sequence_in}.
            Thus $s(m)\in X$ and $s(n)\in X$ by \cref{funs_type_apply}.
            Thus $s(n)\in\metricBall{d}{X}{s(m)}{\epsilon}$ by \cref{metric_ball}.
            Contradiction.
        \end{subproof}
        Take $t,x$ such that
            $t$ is a sequence in $X$ and
            $t$ is a subsequence of $s$ and
            $x\in X$ and
            $t$ converges to $x$ in $X$ with respect to $T$ by \cref{sequentially_compact}.
        Take $v$ such that
            $v$ is strictly increasing on $\naturalsPlus$ and
            for all $n\in\naturalsPlus$ we have $t(n) = \apply{s}{\apply{v}{n}}$ by \cref{subsequence}.
        Let $m_2 = 1 + 1$.
        Let $\delta = \epsilon / m_2$.
        $1\in\reals$ by \cref{real_naturals_subset_reals,real_one_in_naturals,subseteq}.
        $m_2\in\reals$ by \cref{real_add_closed}.
        $0\in\reals$ by \cref{real_add_ident}.
        $0 < 1$ by \cref{real_zero_lt_one}.
        $0 + 1 < 1 + 1$ by \cref{real_lt_add}.
        $0 + 1 = 1$ by \cref{real_add_ident}.
        Thus $1 < m_2$ by \cref{real_lt_subst_right}.
        Thus $0 < m_2$ by \cref{real_lt_trans}.
        $m_2\neq 0$ by \cref{real_pos_nonzero}.
        Thus $\inv{m_2}\in\reals$ by \cref{real_mul_inverse}.
        Thus $0 < \inv{m_2}$ by \cref{real_inv_pos}.
        $\delta = \epsilon \rmul \inv{m_2}$ by \cref{real_div}.
        $0\in\reals$ by \cref{real_add_ident}.
        $\epsilon\in\reals$.
        Thus $0 \rmul \inv{m_2} < \epsilon \rmul \inv{m_2}$ by \cref{real_lt_mul_pos}.
        $0 \rmul \inv{m_2} = 0$ by \cref{real_mul_zero}.
        Thus $0 < \delta$ by \cref{real_lt_subst_right}.
        Show that $\metricBall{d}{X}{x}{\delta}$ is a neighborhood of $x$ in $X$ with respect to $T$.
        \begin{subproof}
            Show that $\metricBall{d}{X}{x}{\delta}$ is open in $X$ with respect to $T$.
            \begin{subproof}
                Show that $\metricBall{d}{X}{x}{\delta}\in\pow{X}$.
                \begin{subproof}
                    Show that for all $y\in\metricBall{d}{X}{x}{\delta}$ we have $y\in X$.
                    \begin{subproof}
                        Fix $y$.
                        Assume $y\in\metricBall{d}{X}{x}{\delta}$.
                        Thus $y\in X$ by \cref{metric_ball}.
                    \end{subproof}
                    Thus $\metricBall{d}{X}{x}{\delta}\subseteq X$ by \cref{subseteq}.
                    Thus $\metricBall{d}{X}{x}{\delta}\in\pow{X}$ by \cref{powerset_intro}.
                \end{subproof}
                Show that there exists $x_0\in X$ such that
                    there exists $\epsilon_0\in\reals$ such that
                        $0 < \epsilon_0$ and $\metricBall{d}{X}{x}{\delta} = \metricBall{d}{X}{x_0}{\epsilon_0}$.
                \begin{subproof}
                    Let $x_0 = x$.
                    Let $\epsilon_0 = \delta$.
                    $x_0\in X$.
                    $\epsilon_0\in\reals$.
                    Thus $0 < \epsilon_0$ by \cref{real_lt_subst_right}.
                \end{subproof}
                Thus $\metricBall{d}{X}{x}{\delta}\in\metricBasis{d}{X}$ by \cref{metric_basis}.
                $\metricBasis{d}{X}$ is a basis on $X$ by \cref{metric_basis_is_basis}.
                Thus $\metricBall{d}{X}{x}{\delta}\in\basisTopology{\metricBasis{d}{X}}{X}$
                    by \cref{basis_elem_open_in_generated}.
                $T = \metricTopology{d}{X}$.
                Thus $\metricBall{d}{X}{x}{\delta}\in T$ by \cref{metric_topology}.
                $T$ is a topology on $X$ by \cref{basis_topology_is_topology,metric_topology,metric_basis_is_basis,basis_for_topology}.
                Thus $\metricBall{d}{X}{x}{\delta}$ is open in $X$ with respect to $T$ by \cref{open_in}.
            \end{subproof}
            $x\in\metricBall{d}{X}{x}{\delta}$ by \cref{metric_ball,metric_zero_characterization,real_lt_subst_left,metric_on}.
            Thus $\metricBall{d}{X}{x}{\delta}$ is a neighborhood of $x$ in $X$ with respect to $T$ by \cref{neighborhood}.
        \end{subproof}
        Take $N\in\naturalsPlus$ such that for all $n\in\naturalsPlus$ if $N\leq n$ then $t(n)\in\metricBall{d}{X}{x}{\delta}$ by \cref{converges_to}.
        $N + 1\in\naturalsPlus$ by \cref{real_naturals_closed_succ}.
        Thus $N\leq N + 1$.
        Thus $t(N)\in\metricBall{d}{X}{x}{\delta}$ and $t(N + 1)\in\metricBall{d}{X}{x}{\delta}$ by \cref{converges_to}.
        Thus $d((x,t(N))) < \delta$ and $d((x,t(N + 1))) < \delta$ by \cref{metric_ball}.
        $t\in\funs{\naturalsPlus}{X}$ by \cref{sequence_in}.
        Thus $t(N)\in X$ by \cref{funs_type_apply}.
        $d$ is symmetric on $X$ by \cref{metric_on}.
        Thus $d((t(N),x)) = d((x,t(N)))$ by \cref{metric_symmetric}.
        Thus $d((t(N),x)) < \delta$ by \cref{real_lt_subst_left}.
        Thus $t(N + 1)\in X$ by \cref{funs_type_apply}.
        $x\in X$.
        $d$ satisfies the triangle inequality on $X$ by \cref{metric_on}.
        Thus $d((t(N),t(N + 1))) \leq d((t(N),x)) + d((x,t(N + 1)))$ by \cref{metric_triangle_inequality}.
        Show that $d((t(N),t(N + 1))) < \epsilon$.
        \begin{subproof}
            $d((t(N),t(N + 1))) < d((t(N),x)) + d((x,t(N + 1)))$ or
                $d((t(N),t(N + 1))) = d((t(N),x)) + d((x,t(N + 1)))$ by \cref{real_leq_split}.
            \begin{byCase}
            \caseOf{$d((t(N),t(N + 1))) < d((t(N),x)) + d((x,t(N + 1)))$.}
                $(t(N),x)\in X\times X$ by \cref{times_tuple_intro}.
                $(x,t(N + 1))\in X\times X$ by \cref{times_tuple_intro}.
                $d\in\funs{X\times X}{\reals}$ by \cref{metric_on}.
                Thus $d((t(N),x))\in\reals$ and $d((x,t(N + 1)))\in\reals$ by \cref{funs_type_apply}.
                $\delta\in\reals$.
                Thus $d((x,t(N + 1))) + d((t(N),x)) < \delta + d((t(N),x))$ by \cref{real_lt_add}.
                $d((x,t(N + 1))) + d((t(N),x)) = d((t(N),x)) + d((x,t(N + 1)))$ by \cref{real_add_comm}.
                $\delta + d((t(N),x)) = d((t(N),x)) + \delta$ by \cref{real_add_comm}.
                Thus $d((t(N),x)) + d((x,t(N + 1))) < d((t(N),x)) + \delta$
                    by \cref{real_lt_subst_left,real_lt_subst_right}.
                Thus $d((t(N),x)) + \delta < \delta + \delta$ by \cref{real_lt_add}.
                Thus $d((t(N),x)) + d((x,t(N + 1)))\in\reals$ by \cref{real_add_closed}.
                Thus $d((t(N),x)) + \delta\in\reals$ by \cref{real_add_closed}.
                Thus $\delta + \delta\in\reals$ by \cref{real_add_closed}.
                Thus $d((t(N),x)) + d((x,t(N + 1))) < \delta + \delta$ by \cref{real_lt_trans}.
                Show that $\delta + \delta = \epsilon$.
                \begin{subproof}
                    $\delta = \epsilon \rmul \inv{m_2}$ by \cref{real_div}.
                    $m_2 \rmul \delta = \epsilon$ by \cref{real_mul_assoc,real_mul_inverse,real_mul_comm,real_mul_ident}.
                    $m_2 \rmul \delta = \delta + \delta$ by \cref{real_right_distrib,real_mul_ident}.
                Thus $\delta + \delta = \epsilon$.
            \end{subproof}
            Thus $d((t(N),x)) + d((x,t(N + 1))) < \epsilon$ by \cref{real_lt_subst_right}.
            $(t(N),t(N + 1))\in X\times X$ by \cref{times_tuple_intro}.
            $d\in\funs{X\times X}{\reals}$ by \cref{metric_on}.
            Thus $d((t(N),t(N + 1)))\in\reals$ by \cref{funs_type_apply}.
            $\epsilon\in\reals$.
            Thus $d((t(N),t(N + 1))) < \epsilon$ by \cref{real_lt_trans}.
        \caseOf{$d((t(N),t(N + 1))) = d((t(N),x)) + d((x,t(N + 1)))$.}
            Show that $d((t(N),x)) + d((x,t(N + 1))) < \epsilon$.
                \begin{subproof}
                    $(t(N),x)\in X\times X$ by \cref{times_tuple_intro}.
                    $(x,t(N + 1))\in X\times X$ by \cref{times_tuple_intro}.
                    $d\in\funs{X\times X}{\reals}$ by \cref{metric_on}.
                    Thus $d((t(N),x))\in\reals$ and $d((x,t(N + 1)))\in\reals$ by \cref{funs_type_apply}.
                    $\delta\in\reals$.
                    Thus $d((x,t(N + 1))) + d((t(N),x)) < \delta + d((t(N),x))$ by \cref{real_lt_add}.
                    $d((x,t(N + 1))) + d((t(N),x)) = d((t(N),x)) + d((x,t(N + 1)))$ by \cref{real_add_comm}.
                    $\delta + d((t(N),x)) = d((t(N),x)) + \delta$ by \cref{real_add_comm}.
                    Thus $d((t(N),x)) + d((x,t(N + 1))) < d((t(N),x)) + \delta$
                        by \cref{real_lt_subst_left,real_lt_subst_right}.
                    Thus $d((t(N),x)) + \delta < \delta + \delta$ by \cref{real_lt_add}.
                    Thus $d((t(N),x)) + d((x,t(N + 1)))\in\reals$ by \cref{real_add_closed}.
                    Thus $d((t(N),x)) + \delta\in\reals$ by \cref{real_add_closed}.
                    Thus $\delta + \delta\in\reals$ by \cref{real_add_closed}.
                    Thus $d((t(N),x)) + d((x,t(N + 1))) < \delta + \delta$ by \cref{real_lt_trans}.
                    Show that $\delta + \delta = \epsilon$.
                    \begin{subproof}
                        $\delta = \epsilon \rmul \inv{m_2}$ by \cref{real_div}.
                        $m_2 \rmul \delta = \epsilon$ by \cref{real_mul_assoc,real_mul_inverse,real_mul_comm,real_mul_ident}.
                        $m_2 \rmul \delta = \delta + \delta$ by \cref{real_right_distrib,real_mul_ident}.
                        Thus $\delta + \delta = \epsilon$.
                    \end{subproof}
                    Thus $d((t(N),x)) + d((x,t(N + 1))) < \epsilon$ by \cref{real_lt_subst_right}.
                \end{subproof}
                Thus $d((t(N),t(N + 1))) < \epsilon$ by \cref{real_lt_subst_left}.
            \end{byCase}
        \end{subproof}
        Thus $v(N)\in\naturalsPlus$ and $v(N + 1)\in\naturalsPlus$ by \cref{strictly_increasing_on_set,funs_type_apply}.
        $0 < 1$ by \cref{real_zero_lt_one}.
        $N\in\reals$ by \cref{real_naturals_subset_reals,subseteq}.
        $0\in\reals$ by \cref{real_add_ident}.
        $1\in\reals$ by \cref{real_mul_ident}.
        Thus $0 + N < 1 + N$ by \cref{real_lt_add}.
        $0 + N = N + 0$ by \cref{real_add_comm}.
        $1 + N = N + 1$ by \cref{real_add_comm}.
        Thus $N + 0 < N + 1$ by \cref{real_lt_subst_left}.
        $N + 0 = N$ by \cref{real_add_ident}.
        Thus $N < N + 1$ by \cref{real_lt_subst_left}.
        Thus $v(N) < v(N + 1)$ by \cref{strictly_increasing_on_set}.
        Thus $t(N) = \apply{s}{\apply{v}{N}}$ and $t(N + 1) = \apply{s}{\apply{v}{(N + 1)}}$.
        Show that it is wrong that $d((t(N),t(N + 1))) < \epsilon$.
        \begin{subproof}
            Suppose not.
            Then $d((t(N),t(N + 1))) < \epsilon$.
            Thus $(t(N),t(N + 1)) = (\apply{s}{\apply{v}{N}},\apply{s}{\apply{v}{(N + 1)}})$ by \cref{pair_eq_iff}.
            Thus $d((t(N),t(N + 1))) = d((\apply{s}{\apply{v}{N}},\apply{s}{\apply{v}{(N + 1)}}))$.
            Thus $d((\apply{s}{\apply{v}{N}},\apply{s}{\apply{v}{(N + 1)}})) < \epsilon$ by \cref{real_lt_subst_left}.
            Thus it is wrong that $d((\apply{s}{\apply{v}{N}},\apply{s}{\apply{v}{(N + 1)}})) < \epsilon$.
            Contradiction.
        \end{subproof}
        Contradiction.
    \end{byCase}
\end{proof}

% from §28 Item 6: compactness, limit point compactness, and sequential compactness
\begin{theorem}\label{metrizable_compact_iff_limit_point_compact_iff_sequentially_compact}
    Assume $T$ is a topology on $X$.
    Suppose $X$ is metrizable with respect to $T$.
    Then $X$ is compact with respect to $T$ iff
        $X$ is limit point compact with respect to $T$ and
        $X$ is sequentially compact with respect to $T$.
\end{theorem}
\begin{proof}
    Take $d$ such that $d$ is a metric on $X$ and $T = \metricTopology{d}{X}$.

    Show that if $X$ is compact with respect to $T$ then $X$ is limit point compact with respect to $T$.
    \begin{subproof}
        Assume $X$ is compact with respect to $T$.
        Thus $X$ is limit point compact with respect to $T$ by \cref{compact_implies_limit_point_compact}.
    \end{subproof}

    Show that if $X$ is sequentially compact with respect to $T$ then $X$ is compact with respect to $T$.
    \begin{subproof}
        Assume $X$ is sequentially compact with respect to $T$.
        Show that for all $A$ if $A$ is an open covering of $X$ with respect to $T$ then
            there exists $B$ such that $B\subseteq A$ and $B$ is finite and $B$ is a covering of $X$.
        \begin{subproof}
            Fix $A$.
            Assume $A$ is an open covering of $X$ with respect to $T$.
            \begin{byCase}
            \caseOf{$X = \emptyset$.}
                Let $B = \emptyset$.
                $B\subseteq A$ by \cref{emptyset_subseteq}.
                $B$ is finite by \cref{emptyset_is_finite}.
                $X\subseteq\unions{B}$ by \cref{emptyset_subseteq,unions_emptyset}.
                Thus $B$ is a covering of $X$ by \cref{covering}.
                Thus there exists $B$ such that $B\subseteq A$ and $B$ is finite and $B$ is a covering of $X$.
            \caseOf{$X\neq\emptyset$.}
                Take $\delta\in\reals$ such that $0 < \delta$ and
                    for all $B$ such that $B\subseteq X$ and
                        for all $x,y\in B$ we have $d((x,y)) < \delta$
                    there exists $U\in A$ such that $B\subseteq U$
                    by \cref{sequentially_compact_lebesgue}.
                Let $m_2 = 1 + 1$.
                Let $\epsilon = \delta / m_2$.
                $1\in\reals$ by \cref{real_mul_ident}.
                Thus $m_2\in\reals$ by \cref{real_add_closed}.
                $0\in\reals$ by \cref{real_add_ident}.
                $0 < 1$ by \cref{real_zero_lt_one}.
                Thus $0 + 1 < 1 + 1$ by \cref{real_lt_add}.
                $0 + 1 = 1$ by \cref{real_add_ident}.
                Thus $1 < m_2$ by \cref{real_lt_subst_left}.
                Thus $0 < m_2$ by \cref{real_lt_trans}.
                $m_2\neq 0$ by \cref{real_pos_nonzero}.
                Thus $\inv{m_2}\in\reals$ by \cref{real_mul_inverse}.
                Thus $0 < \inv{m_2}$ by \cref{real_inv_pos}.
                $\epsilon = \delta \rmul \inv{m_2}$ by \cref{real_div}.
                Thus $\epsilon\in\reals$ by \cref{real_mul_closed}.
                $0\in\reals$ by \cref{real_add_ident}.
                Thus $0 \rmul \inv{m_2} < \delta \rmul \inv{m_2}$ by \cref{real_lt_mul_pos}.
                $0 \rmul \inv{m_2} = 0$ by \cref{real_mul_zero}.
                Thus $0 < \epsilon$ by \cref{real_lt_subst_left}.
                Take $B$ such that
                    $B$ is finite and
                    $B$ is a covering of $X$ and
                    for all $U\in B$ there exists $x\in X$ such that $U = \metricBall{d}{X}{x}{\epsilon}$
                    by \cref{sequentially_compact_finite_ball_cover}.

                Show that for all $V\in B$ there exists $U\in A$ such that $V\subseteq U$.
                \begin{subproof}
                    Fix $V$.
                    Assume $V\in B$.
                    Take $x_0\in X$ such that $V = \metricBall{d}{X}{x_0}{\epsilon}$.
                    Show that $V\subseteq X$.
                    \begin{subproof}
                        Show that for all $y\in V$ we have $y\in X$.
                        \begin{subproof}
                            Fix $y$.
                            Assume $y\in V$.
                            Thus $y\in X$ by \cref{metric_ball}.
                        \end{subproof}
                        Thus $V\subseteq X$ by \cref{subseteq}.
                    \end{subproof}
                    Show that for all $x,y\in V$ we have $d((x,y)) < \delta$.
                    \begin{subproof}
                        Fix $x$.
                        Fix $y$.
                        Assume $x\in V$ and $y\in V$.
                        Thus $x\in X$ and $y\in X$ by \cref{subseteq}.
                        Thus $d((x_0,x)) < \epsilon$ and $d((x_0,y)) < \epsilon$ by \cref{metric_ball}.
                        $d$ satisfies the triangle inequality on $X$ by \cref{metric_on}.
                        Thus $d((x,y)) \leq d((x,x_0)) + d((x_0,y))$ by \cref{metric_triangle_inequality}.
                        $d$ is symmetric on $X$ by \cref{metric_on,metric_symmetric}.
                        Thus $d((x,x_0)) = d((x_0,x))$ by \cref{metric_symmetric}.
                        Thus $d((x,x_0)) < \epsilon$ by \cref{real_lt_subst_left}.
                        Show that $d((x,y)) < \epsilon + \epsilon$.
                        \begin{subproof}
                            $(x,x_0)\in X\times X$ and $(x_0,y)\in X\times X$ by \cref{times_tuple_intro}.
                            $(x,y)\in X\times X$ by \cref{times_tuple_intro}.
                            $d\in\funs{X\times X}{\reals}$ by \cref{metric_on}.
                            Thus $d((x,x_0))\in\reals$ and $d((x_0,y))\in\reals$ by \cref{funs_type_apply}.
                            Thus $d((x,y))\in\reals$ by \cref{funs_type_apply}.
                            $\epsilon\in\reals$.
                            Thus $d((x,x_0)) + d((x_0,y)) < \epsilon + d((x_0,y))$ by \cref{real_lt_add}.
                            $d((x,x_0)) + d((x_0,y)) = d((x_0,y)) + d((x,x_0))$ by \cref{real_add_comm}.
                            $\epsilon + d((x_0,y)) = d((x_0,y)) + \epsilon$ by \cref{real_add_comm}.
                            Thus $d((x,x_0)) + d((x_0,y)) < d((x_0,y)) + \epsilon$
                                by \cref{real_lt_subst_left,real_lt_subst_right}.
                            Thus $d((x_0,y)) < \epsilon$ by \cref{real_lt_subst_left}.
                            Thus $d((x_0,y)) + \epsilon < \epsilon + \epsilon$ by \cref{real_lt_add}.
                            Thus $d((x,x_0)) + d((x_0,y))\in\reals$ by \cref{real_add_closed}.
                            Thus $d((x_0,y)) + \epsilon\in\reals$ by \cref{real_add_closed}.
                            Thus $\epsilon + \epsilon\in\reals$ by \cref{real_add_closed}.
                            Thus $d((x,x_0)) + d((x_0,y)) < \epsilon + \epsilon$ by \cref{real_lt_trans}.
                            $d((x,y)) < d((x,x_0)) + d((x_0,y))$ or
                                $d((x,y)) = d((x,x_0)) + d((x_0,y))$ by \cref{real_leq_split}.
                            \begin{byCase}
                            \caseOf{$d((x,y)) < d((x,x_0)) + d((x_0,y))$.}
                                Thus $d((x,y)) < \epsilon + \epsilon$ by \cref{real_lt_trans}.
                            \caseOf{$d((x,y)) = d((x,x_0)) + d((x_0,y))$.}
                                Thus $d((x,y)) < \epsilon + \epsilon$ by \cref{real_lt_subst_left}.
                            \end{byCase}
                        \end{subproof}
                        Show that $\epsilon + \epsilon = \delta$.
                        \begin{subproof}
                            $\epsilon = \delta \rmul \inv{m_2}$ by \cref{real_div}.
                            $m_2 \rmul \epsilon = \delta$ by \cref{real_mul_assoc,real_mul_inverse,real_mul_comm,real_mul_ident}.
                            $m_2 \rmul \epsilon = \epsilon + \epsilon$ by \cref{real_right_distrib,real_mul_ident}.
                            Thus $\epsilon + \epsilon = \delta$.
                        \end{subproof}
                        Thus $d((x,y)) < \delta$ by \cref{real_lt_subst_right}.
                    \end{subproof}
                    Thus there exists $U\in A$ such that $V\subseteq U$.
                \end{subproof}

                Let $R = \{w\in B\times A \mid \fst{w}\subseteq \snd{w}\}$.
                Show that $R$ is a relation.
                \begin{subproof}
                    Show that for all $w$ such that $w\in R$ there exist $V,U$ such that $w = (V,U)$.
                    \begin{subproof}
                        Fix $w$.
                        Assume $w\in R$.
                        Thus $w\in B\times A$.
                        Take $V,U$ such that $w = (V,U)$ by \cref{times_elem_is_tuple}.
                    \end{subproof}
                    Thus $R$ is a relation by \cref{relation}.
                \end{subproof}
                Show that for all $V\in B$ there exists $U$ such that $V\mathrel{R} U$.
                \begin{subproof}
                    Fix $V$.
                    Assume $V\in B$.
                    Take $U\in A$ such that $V\subseteq U$.
                    Thus $(V,U)\in B\times A$ by \cref{times_tuple_intro}.
                    Thus $(V,U)\in R$.
                    Thus $V\mathrel{R} U$.
                \end{subproof}
                Take $h$ such that $h$ is a function on $B$ and for all $V\in B$ we have $V\mathrel{R}\apply{h}{V}$
                    by \cref{choice_function}.
                Let $C = \img{h}{B}$.
                Thus $C$ is finite by \cref{finite_image_function}.
                Show that $C\subseteq A$.
                \begin{subproof}
                    Show that for all $U$ such that $U\in C$ we have $U\in A$.
                    \begin{subproof}
                        Fix $U$.
                        Assume $U\in C$.
                        Take $V\in B$ such that $(V,U)\in h$ by \cref{img_iff}.
                        Thus $U = \apply{h}{V}$ by \cref{function_apply_intro}.
                        Thus $V\mathrel{R} U$.
                        Thus $(V,U)\in B\times A$.
                        Thus $U\in A$ by \cref{times_tuple_elim}.
                    \end{subproof}
                    Thus $C\subseteq A$ by \cref{subseteq}.
                \end{subproof}
                Show that $C$ is a covering of $X$.
                \begin{subproof}
                    Show that for all $x$ such that $x\in X$ we have $x\in\unions{C}$.
                    \begin{subproof}
                        Fix $x$.
                        Assume $x\in X$.
                        $X\subseteq\unions{B}$ by \cref{covering}.
                        Thus $x\in\unions{B}$ by \cref{subseteq}.
                        Take $V\in B$ such that $x\in V$ by \cref{unions_iff}.
                        Thus $V\mathrel{R}\apply{h}{V}$.
                        Thus $(V,\apply{h}{V})\in R$.
                        Thus $\fst{(V,\apply{h}{V})}\subseteq \snd{(V,\apply{h}{V})}$.
                        Thus $V\subseteq \apply{h}{V}$ by \cref{fst_eq,snd_eq}.
                        Thus $x\in\apply{h}{V}$ by \cref{subseteq}.
                        Thus $V\in\dom{h}$.
                        Thus $(V,\apply{h}{V})\in h$ by \cref{function_apply_elim}.
                        Thus $\apply{h}{V}\in C$ by \cref{img_iff}.
                        Thus $x\in\unions{C}$ by \cref{unions_iff}.
                    \end{subproof}
                    Thus $X\subseteq\unions{C}$ by \cref{subseteq}.
                    Thus $C$ is a covering of $X$ by \cref{covering}.
                \end{subproof}
                Thus there exists $B$ such that $B\subseteq A$ and $B$ is finite and $B$ is a covering of $X$.
            \end{byCase}
        \end{subproof}
        Thus $X$ is compact with respect to $T$ by \cref{compact}.
    \end{subproof}

    Show that if $X$ is limit point compact with respect to $T$ then $X$ is sequentially compact with respect to $T$.
    \begin{subproof}
        Assume $X$ is limit point compact with respect to $T$.
        Show that for all $s$ such that $s$ is a sequence in $X$
            there exist $t,x$ such that
                $t$ is a sequence in $X$ and
                $t$ is a subsequence of $s$ and
                $x\in X$ and
                $t$ converges to $x$ in $X$ with respect to $T$.
        \begin{subproof}
            Fix $s$.
            Assume $s$ is a sequence in $X$.
            $s\in\funs{\naturalsPlus}{X}$ by \cref{sequence_in}.
            Let $A = \ran{s}$.
            Thus $A\subseteq X$ by \cref{funs_ran}.
            \begin{byCase}
            \caseOf{$A$ is finite.}
                Show that $s$ is a sequence in $A$.
                \begin{subproof}
                    Show that $s\in\funs{\naturalsPlus}{A}$.
                    \begin{subproof}
                        Show that for all $n\in\dom{s}$ we have $s(n)\in A$.
                        \begin{subproof}
                            Fix $n\in\dom{s}$.
                            Thus $s$ is a function by \cref{funs_is_function}.
                            Thus $(n,s(n))\in s$ by \cref{function_apply_elim}.
                            Thus $s(n)\in\ran{s}$ by \cref{ran_iff}.
                            Thus $s(n)\in A$.
                        \end{subproof}
                        Thus $s$ is a function to $A$.
                    \end{subproof}
                    Thus $s$ is a sequence in $A$ by \cref{sequence_in}.
                \end{subproof}
                Take $t,a$ such that
                    $a\in A$ and
                    $t$ is a sequence in $A$ and
                    $t$ is a subsequence of $s$ and
                    for all $n\in\naturalsPlus$ we have $t(n) = a$
                    by \cref{finite_sequence_constant_subsequence}.
                Thus $a\in X$ by \cref{subseteq}.
                Show that $t$ is a sequence in $X$.
                \begin{subproof}
                    Show that $t\in\funs{\naturalsPlus}{X}$.
                    \begin{subproof}
                        Show that for all $n\in\dom{t}$ we have $t(n)\in X$.
                        \begin{subproof}
                            Fix $n\in\dom{t}$.
                            Thus $t$ is a sequence in $A$.
                            Thus $t\in\funs{\naturalsPlus}{A}$ by \cref{sequence_in}.
                            Thus $\dom{t} = \naturalsPlus$ by \cref{funs}.
                            Thus $n\in\naturalsPlus$ by \cref{subseteq}.
                            Thus $t(n)\in A$ by \cref{funs_type_apply}.
                            Thus $t(n)\in X$ by \cref{subseteq}.
                        \end{subproof}
                        Thus $t$ is a function to $X$.
                    \end{subproof}
                    Thus $t$ is a sequence in $X$ by \cref{sequence_in}.
                \end{subproof}
                Show that $t$ converges to $a$ in $X$ with respect to $T$.
                \begin{subproof}
                    Show that for all $U$ such that $U$ is a neighborhood of $a$ in $X$ with respect to $T$
                        there exists $N\in\naturalsPlus$ such that for all $n\in\naturalsPlus$ if $N\leq n$ then $t(n)\in U$.
                    \begin{subproof}
                        Fix $U$.
                        Assume $U$ is a neighborhood of $a$ in $X$ with respect to $T$.
                        Thus $a\in U$ by \cref{neighborhood}.
                        Take $N$ such that $N = 1$.
                        Show that for all $n\in\naturalsPlus$ if $N\leq n$ then $t(n)\in U$.
                        \begin{subproof}
                            Fix $n\in\naturalsPlus$.
                            Assume $N\leq n$.
                            Thus $t(n) = a$.
                            Thus $t(n)\in U$ by \cref{subseteq}.
                        \end{subproof}
                        Thus there exists $N\in\naturalsPlus$ such that for all $n\in\naturalsPlus$ if $N\leq n$ then $t(n)\in U$.
                    \end{subproof}
                    Thus $t$ converges to $a$ in $X$ with respect to $T$ by \cref{converges_to}.
                \end{subproof}
                Thus there exist $t,a$ such that
                    $t$ is a sequence in $X$ and
                    $t$ is a subsequence of $s$ and
                    $a\in X$ and
                    $t$ converges to $a$ in $X$ with respect to $T$.
            \caseOf{$A$ is infinite.}
                Take $x\in X$ such that $x$ is a limit point of $A$ in $X$ with respect to $T$
                    by \cref{limit_point_compact}.

                Show that $X$ satisfies the T-one axiom with respect to $T$.
                \begin{subproof}
                    $T$ is a topology on $X$ by \cref{metrizable}.
                    $X$ is Hausdorff with respect to $T$ by \cref{metric_space_hausdorff}.
                    Show that for all $F$ such that $F\subseteq X$ and $F$ is finite we have
                        $F$ is closed in $X$ with respect to $T$.
                    \begin{subproof}
                        Fix $F$.
                        Assume $F\subseteq X$ and $F$ is finite.
                        Thus $F$ is closed in $X$ with respect to $T$ by \cref{finite_point_set_closed_hausdorff}.
                    \end{subproof}
                    Thus $X$ satisfies the T-one axiom with respect to $T$ by \cref{t1_axiom}.
                \end{subproof}

                Show that for all $U$ such that $U$ is a neighborhood of $x$ in $X$ with respect to $T$
                    we have $U\inter A$ is infinite.
                \begin{subproof}
                    Thus for all $U$ such that $U$ is a neighborhood of $x$ in $X$ with respect to $T$
                        we have $U\inter A$ is infinite by \cref{limit_point_infinite_neighborhood}.
                \end{subproof}

                Show that for all $n\in\naturalsPlus$ we have
                    $\metricBall{d}{X}{x}{1 / n}$ is a neighborhood of $x$ in $X$ with respect to $T$.
                \begin{subproof}
                    Fix $n\in\naturalsPlus$.
                    Let $U = \metricBall{d}{X}{x}{1 / n}$.
                    Show that $U$ is open in $X$ with respect to $T$.
                    \begin{subproof}
                        Let $B = \metricBasis{d}{X}$.
                        $B$ is a basis on $X$ by \cref{metric_basis_is_basis}.
                        $0 < 1 / n$ by \cref{positive_reciprocal_naturalsplus}.
                        $n\in\reals$ by \cref{real_naturals_subset_reals,subseteq}.
                        $1\in\reals$ by \cref{real_naturals_subset_reals,real_one_in_naturals,subseteq}.
                        $0\in\reals$ by \cref{real_add_ident}.
                        $0 < 1$ by \cref{real_zero_lt_one}.
                        Show that $0 < n$.
                        \begin{subproof}
                            $1 < n$ or $1 = n$ by \cref{real_one_leq_naturalsplus}.
                            \begin{byCase}
                            \caseOf{$1 < n$.}
                                Thus $0 < n$ by \cref{real_lt_trans}.
                            \caseOf{$1 = n$.}
                                Thus $0 < n$ by \cref{real_lt_subst_right}.
                            \end{byCase}
                        \end{subproof}
                        $n\neq 0$ by \cref{real_pos_nonzero}.
                        Thus $\inv{n}\in\reals$ by \cref{real_mul_inverse}.
                        $1 / n = 1 \rmul \inv{n}$ by \cref{real_div}.
                        Thus $1 / n\in\reals$ by \cref{real_mul_closed}.
                        Show that $U\subseteq X$.
                        \begin{subproof}
                            Show that for all $y$ such that $y\in U$ we have $y\in X$.
                            \begin{subproof}
                                Fix $y$.
                                Assume $y\in U$.
                                Then $y\in\metricBall{d}{X}{x}{1 / n}$.
                                Thus $y\in X$ by \cref{metric_ball}.
                            \end{subproof}
                            Thus $U\subseteq X$ by \cref{subseteq}.
                        \end{subproof}
                        Thus $U\in\pow{X}$ by \cref{powerset_intro}.
                        Thus $U\in B$ by \cref{metric_basis}.
                        Thus $U\in\basisTopology{B}{X}$ by \cref{basis_elem_open_in_generated}.
                        Thus $U\in T$ by \cref{metric_topology}.
                        Thus $U$ is open in $X$ with respect to $T$ by \cref{open_in,basis_topology_is_topology,metric_basis_is_basis,metric_topology}.
                    \end{subproof}
                    Show that $x\in U$.
                    \begin{subproof}
                        $d((x,x)) = 0$ by \cref{metric_zero_characterization,metric_on}.
                        $0 < 1 / n$ by \cref{positive_reciprocal_naturalsplus}.
                        Thus $d((x,x)) < 1 / n$ by \cref{real_lt_subst_left}.
                        Thus $x\in U$ by \cref{metric_ball}.
                    \end{subproof}
                    Thus $U$ is a neighborhood of $x$ in $X$ with respect to $T$ by \cref{neighborhood}.
                \end{subproof}

                Let $D = \{w\in\naturalsPlus\times\naturalsPlus \mid
                    \apply{s}{\snd{w}}\in\metricBall{d}{X}{x}{1 / (\fst{w})}\}$.

                Show that there exists $a_0\in D$ such that $\fst{a_0} = 1$.
                \begin{subproof}
                    Let $U_1 = \metricBall{d}{X}{x}{1 / 1}$.
                    $U_1$ is a neighborhood of $x$ in $X$ with respect to $T$.
                    Thus $U_1\inter A$ is infinite.
                    Show that $U_1\inter A\neq \emptyset$.
                    \begin{subproof}
                        Suppose not.
                        Then $U_1\inter A = \emptyset$.
                        Then $U_1\inter A$ is finite by \cref{emptyset_is_finite}.
                        Contradiction.
                    \end{subproof}
                    Take $y$ such that $y\in U_1\inter A$ by \cref{nonempty_inhabited}.
                    Thus $y\in A$ by \cref{inter_elim_right}.
                    Thus $y\in\ran{s}$.
                    Take $k$ such that $(k,y)\in s$ by \cref{ran_iff}.
                    Thus $k\in\dom{s}$ by \cref{dom}.
                    Thus $k\in\naturalsPlus$ by \cref{funs,subseteq,sequence_in}.
                    Thus $s(k) = y$ by \cref{function_apply_iff,funs_is_function}.
                    Thus $s(k)\in U_1$.
                    Let $a_0 = (1,k)$.
                    $1\in\naturalsPlus$ by \cref{real_one_in_naturals}.
                    Thus $a_0\in\naturalsPlus\times\naturalsPlus$ by \cref{times_tuple_intro}.
                    Thus $\apply{s}{k}\in\metricBall{d}{X}{x}{1 / 1}$.
                    Thus $\apply{s}{\snd{a_0}}\in\metricBall{d}{X}{x}{1 / (\fst{a_0})}$ by \cref{fst_eq,snd_eq}.
                    Thus $a_0\in D$.
                    Thus there exists $a_0\in D$ such that $\fst{a_0} = 1$ by \cref{fst_eq}.
                \end{subproof}

                Let $R = \{w\in D\times D \mid
                    \exists n,k,n_1,k_1\in\naturalsPlus.
                        w = ((n,k),(n_1,k_1)) \land n_1 = n + 1 \land k < k_1\}$.

                Show that $R$ is a relation.
                \begin{subproof}
                    Show that for all $w$ such that $w\in R$ there exist $p,q$ such that $w = (p,q)$.
                    \begin{subproof}
                        Fix $w$.
                        Assume $w\in R$.
                        Thus $w\in D\times D$.
                        Take $p,q$ such that $w = (p,q)$ by \cref{times_elem_is_tuple}.
                    \end{subproof}
                    Thus $R$ is a relation by \cref{relation}.
                \end{subproof}

                Show that for all $p\in D$ there exists $q$ such that $p\mathrel{R} q$.
                \begin{subproof}
                    Fix $p\in D$.
                    Take $n,k$ such that $p = (n,k)$ by \cref{times_elem_is_tuple}.
                    Thus $n\in\naturalsPlus$ and $k\in\naturalsPlus$ by \cref{times_tuple_elim}.
                    Let $V = \metricBall{d}{X}{x}{1 / (n + 1)}$.
                    Show that $V$ is a neighborhood of $x$ in $X$ with respect to $T$.
                    \begin{subproof}
                        $n + 1\in\naturalsPlus$ by \cref{real_naturals_closed_succ}.
                        Thus $\metricBall{d}{X}{x}{1 / (n + 1)}$ is a neighborhood of $x$ in $X$ with respect to $T$.
                    \end{subproof}
                    Thus $V\inter A$ is infinite.
                    Let $J = \{m\in\naturalsPlus \mid \apply{s}{m}\in V\}$.
                    Show that $J$ is infinite.
                    \begin{subproof}
                        Suppose not.
                        Then $J$ is finite.
                        Show that $\img{s}{J} = V\inter A$.
                        \begin{subproof}
                            Show that $\img{s}{J}\subseteq V\inter A$.
                            \begin{subproof}
                                Show that for all $y$ such that $y\in\img{s}{J}$ we have $y\in V\inter A$.
                                \begin{subproof}
                                    Fix $y$.
                                    Assume $y\in\img{s}{J}$.
                                    Take $m\in J$ such that $(m,y)\in s$ by \cref{img_iff}.
                                    Thus $y = \apply{s}{m}$ by \cref{function_apply_intro,funs_is_function}.
                                    Thus $y\in V$ by \cref{subseteq}.
                                    Thus $y\in A$ by \cref{ran_iff}.
                                    Thus $y\in V\inter A$ by \cref{inter_intro}.
                                \end{subproof}
                                Thus $\img{s}{J}\subseteq V\inter A$ by \cref{subseteq}.
                            \end{subproof}
                            Show that $V\inter A\subseteq \img{s}{J}$.
                            \begin{subproof}
                                Show that for all $y$ such that $y\in V\inter A$ we have $y\in\img{s}{J}$.
                                \begin{subproof}
                                    Fix $y$.
                                    Assume $y\in V\inter A$.
                                    Thus $y\in A$ by \cref{inter_elim_right}.
                                    Thus $y\in\ran{s}$.
                                    Take $m$ such that $(m,y)\in s$ by \cref{ran_iff}.
                                    Thus $m\in\dom{s}$ by \cref{dom}.
                                    Thus $m\in\naturalsPlus$ by \cref{funs,subseteq,sequence_in}.
                                    Thus $y = \apply{s}{m}$ by \cref{function_apply_intro,funs_is_function}.
                                    Thus $m\in J$ by \cref{subseteq,inter_elim_left}.
                                    Thus $y\in\img{s}{J}$ by \cref{img_iff}.
                                \end{subproof}
                                Thus $V\inter A\subseteq \img{s}{J}$ by \cref{subseteq}.
                            \end{subproof}
                            Thus $\img{s}{J} = V\inter A$ by \cref{subseteq_antisymmetric}.
                        \end{subproof}
                        Let $g = \restrl{s}{J}$.
                        Show that $g$ is a function on $J$.
                        \begin{subproof}
                            $s\in\funs{\naturalsPlus}{X}$ by \cref{sequence_in}.
                            Thus $s$ is a function by \cref{funs_is_function}.
                            Thus $g$ is a function by \cref{restrl_of_function_is_function}.
                            Show that $\dom{g} = J$.
                            \begin{subproof}
                                Thus $\dom{s} = \naturalsPlus$ by \cref{funs_elim}.
                                Thus $J\subseteq\dom{s}$ by \cref{subseteq}.
                                $\dom{g} = \dom{s}\inter J$ by \cref{dom_of_restrl_eq_inter}.
                                Thus $\dom{g} = J$ by \cref{inter_absorb_supseteq_left}.
                            \end{subproof}
                            Thus $g$ is a function on $J$.
                        \end{subproof}
                        $s\in\funs{\naturalsPlus}{X}$ by \cref{sequence_in}.
                        Thus $\dom{s} = \naturalsPlus$ by \cref{funs_elim}.
                        Thus $J\subseteq\dom{s}$ by \cref{subseteq}.
                        Thus $\img{g}{J} = \img{s}{J\inter J}$ by \cref{restrl_img}.
                        $J\subseteq J$ by \cref{subseteq_refl}.
                        Thus $J\inter J = J$ by \cref{inter_absorb_supseteq_right}.
                        Thus $\img{g}{J} = \img{s}{J}$.
                        Thus $\img{g}{J} = V\inter A$.
                        Thus $\img{g}{J}$ is finite by \cref{finite_image_function}.
                        Thus $V\inter A$ is finite.
                        Contradiction.
                    \end{subproof}
                    Show that there exists $k_1\in J$ such that $k < k_1$.
                    \begin{subproof}
                        Suppose not.
                        Then for all $m\in J$ we have $m \leq k$.
                        Show that $J\subseteq\initialSegment{k}$.
                        \begin{subproof}
                            Show that for all $m\in J$ we have $m\in\initialSegment{k}$.
                            \begin{subproof}
                                Fix $m$.
                                Assume $m\in J$.
                                Thus $m\in\naturalsPlus$.
                                Thus $m\leq k$.
                                Thus $m\in\initialSegment{k}$ by \cref{initial_segment_naturalsplus}.
                            \end{subproof}
                            Thus $J\subseteq\initialSegment{k}$ by \cref{subseteq}.
                        \end{subproof}
                        $\initialSegment{k}$ is finite by \cref{initial_segment_finite}.
                        Thus $J$ is finite by \cref{subseteq_finite_is_finite}.
                        Contradiction.
                    \end{subproof}
                    Take $k_1\in J$ such that $k < k_1$.
                    Let $q = (n + 1,k_1)$.
                    Show that $q\in D$.
                    \begin{subproof}
                        $n + 1\in\naturalsPlus$ by \cref{real_naturals_closed_succ}.
                        $k_1\in\naturalsPlus$ by \cref{subseteq}.
                        Thus $q\in\naturalsPlus\times\naturalsPlus$ by \cref{times_tuple_intro}.
                        Thus $\apply{s}{\snd{q}}\in V$ by \cref{snd_eq}.
                        Thus $\apply{s}{\snd{q}}\in\metricBall{d}{X}{x}{1 / (n + 1)}$.
                        Thus $q\in D$.
                    \end{subproof}
                    Let $n_1 = n + 1$.
                    Thus $n_1\in\naturalsPlus$ by \cref{real_naturals_closed_succ}.
                    Thus $(p,q) = ((n,k),(n_1,k_1))$ by \cref{pair_eq_iff}.
                    Thus $(p,q)\in R$.
                    Thus $p\mathrel{R} q$.
                \end{subproof}

                Take $f$ such that $f$ is a function on $D$ and for all $p\in D$ we have
                    $p\mathrel{R}\apply{f}{p}$ by \cref{choice_function}.
                Show that $f\in\funs{D}{D}$.
                \begin{subproof}
                    Show that $f$ is a function to $D$.
                    \begin{subproof}
                        Show that for all $p\in\dom{f}$ we have $f(p)\in D$.
                        \begin{subproof}
                            Fix $p$.
                            Assume $p\in\dom{f}$.
                            Thus $\dom{f} = D$.
                            Thus $p\in D$.
                            Thus $p\mathrel{R}\apply{f}{p}$.
                            Thus $(p,\apply{f}{p})\in R$.
                            Thus $(p,\apply{f}{p})\in D\times D$ by \cref{subseteq}.
                            Thus $\apply{f}{p}\in D$ by \cref{times_tuple_elim}.
                        \end{subproof}
                        Thus $f$ is a function to $D$.
                    \end{subproof}
                    Thus $\dom{f} = D$.
                    Thus $f\in\funs{D}{D}$ by \cref{funs_intro}.
                \end{subproof}
                Take $a_0\in D$ such that $\fst{a_0} = 1$.
                Take $v$ such that
                    $v$ is a sequence in $D$ and
                    $v(1) = a_0$ and
                    for all $n\in\naturalsPlus$ we have $v(n + 1) = \apply{f}{v(n)}$
                    by \cref{iteration_sequence_naturalsplus}.

                Let $g = \{(p,\snd{p}) \mid p\in D\}$.
                Show that $g$ is a function on $D$.
                \begin{subproof}
                    Show that $g$ is a relation.
                    \begin{subproof}
                        Show that for all $w$ such that $w\in g$ there exist $p,y$ such that $w = (p,y)$.
                        \begin{subproof}
                            Fix $w$.
                            Assume $w\in g$.
                            Take $p\in D$ such that $w = (p,\snd{p})$.
                            Thus there exist $p,y$ such that $w = (p,y)$.
                        \end{subproof}
                        Thus $g$ is a relation by \cref{relation}.
                    \end{subproof}
                    Show that $g$ is right-unique.
                    \begin{subproof}
                        Show that for all $p,y_1,y_2$ such that $(p,y_1)\in g$ and $(p,y_2)\in g$ we have $y_1 = y_2$.
                        \begin{subproof}
                            Fix $p$.
                            Fix $y_1$.
                            Fix $y_2$.
                            Assume $(p,y_1)\in g$ and $(p,y_2)\in g$.
                            Take $p_1\in D$ such that $(p,y_1) = (p_1,\snd{p_1})$.
                            Take $p_2\in D$ such that $(p,y_2) = (p_2,\snd{p_2})$.
                            Thus $p = p_1$ and $y_1 = \snd{p_1}$ by \cref{pair_eq_iff}.
                            Thus $p = p_2$ and $y_2 = \snd{p_2}$ by \cref{pair_eq_iff}.
                            Thus $y_1 = y_2$.
                        \end{subproof}
                        Thus $g$ is right-unique by \cref{rightunique}.
                    \end{subproof}
                    Show that $\dom{g} = D$.
                    \begin{subproof}
                        Show that $\dom{g}\subseteq D$.
                        \begin{subproof}
                            Show that for all $p$ such that $p\in\dom{g}$ we have $p\in D$.
                            \begin{subproof}
                                Fix $p$.
                                Assume $p\in\dom{g}$.
                                Take $y$ such that $(p,y)\in g$ by \cref{dom}.
                                Take $p_1\in D$ such that $(p,y) = (p_1,\snd{p_1})$.
                                Thus $p = p_1$ by \cref{pair_eq_iff}.
                                Thus $p\in D$.
                            \end{subproof}
                            Thus $\dom{g}\subseteq D$ by \cref{subseteq}.
                        \end{subproof}
                        Show that $D\subseteq\dom{g}$.
                        \begin{subproof}
                            Show that for all $p$ such that $p\in D$ we have $p\in\dom{g}$.
                            \begin{subproof}
                                Fix $p$.
                                Assume $p\in D$.
                                Thus $(p,\snd{p})\in g$.
                                Thus $p\in\dom{g}$ by \cref{dom}.
                            \end{subproof}
                            Thus $D\subseteq\dom{g}$ by \cref{subseteq}.
                        \end{subproof}
                        Thus $\dom{g} = D$ by \cref{subseteq_antisymmetric}.
                    \end{subproof}
                    Show that for all $p\in\dom{g}$ we have $g(p)\in\naturalsPlus$.
                    \begin{subproof}
                        Fix $p\in\dom{g}$.
                        Take $p_1\in D$ such that $(p,g(p)) = (p_1,\snd{p_1})$ by \cref{pair_eq_iff,function_apply_intro}.
                        Thus $p = p_1$ and $g(p) = \snd{p_1}$ by \cref{pair_eq_iff}.
                        Thus $p_1\in\naturalsPlus\times\naturalsPlus$ by \cref{subseteq}.
                        Take $n_0,k_0$ such that $p_1 = (n_0,k_0)$ and $n_0\in\naturalsPlus$ and $k_0\in\naturalsPlus$
                            by \cref{times_elem_is_tuple}.
                        Thus $\snd{p_1} = k_0$ by \cref{snd_eq}.
                        Thus $g(p) = k_0$.
                        Thus $g(p)\in\naturalsPlus$.
                    \end{subproof}
                    Thus $g\in\funs{D}{\naturalsPlus}$ by \cref{funs_intro}.
                \end{subproof}

                Let $u = g\circ v$.
                Show that $u$ is a sequence in $\naturalsPlus$.
                \begin{subproof}
                    Show that $u\in\funs{\naturalsPlus}{\naturalsPlus}$.
                    \begin{subproof}
                        Show that $u$ is a function.
                        \begin{subproof}
                            $v\in\funs{\naturalsPlus}{D}$ by \cref{sequence_in}.
                            Thus $v$ is right-unique by \cref{funs}.
                            Thus $g$ is right-unique by \cref{funs}.
                            Thus $u$ is a function by \cref{function_circ}.
                        \end{subproof}
                        Show that $\dom{u} = \naturalsPlus$.
                        \begin{subproof}
                            $v\in\funs{\naturalsPlus}{D}$ by \cref{sequence_in}.
                            Thus $\dom{v} = \naturalsPlus$ by \cref{funs}.
                            Thus $\ran{v}\subseteq D$ by \cref{funs_ran}.
                            Thus $\ran{v}\subseteq\dom{g}$ by \cref{funs}.
                            Thus $v$ is right-unique by \cref{funs}.
                            Thus $v$ is a relation by \cref{funs_is_relation}.
                            Thus $g$ is right-unique by \cref{funs}.
                            Thus $g$ is a relation by \cref{funs_is_relation}.
                            Thus $\dom{u} = \preimg{v}{\dom{g}}$ by \cref{dom_of_circ}.
                            Show that $\preimg{v}{\dom{g}} = \dom{v}$.
                            \begin{subproof}
                                Show that $\preimg{v}{\dom{g}}\subseteq\dom{v}$.
                                \begin{subproof}
                                    Thus $\preimg{v}{\dom{g}}\subseteq\dom{v}$ by \cref{preimg_subseteq_dom}.
                                \end{subproof}
                                Show that $\dom{v}\subseteq\preimg{v}{\dom{g}}$.
                                \begin{subproof}
                                    Show that for all $n$ such that $n\in\dom{v}$ we have $n\in\preimg{v}{\dom{g}}$.
                                    \begin{subproof}
                                    Fix $n$.
                                    Assume $n\in\dom{v}$.
                                    $v$ is a function by \cref{funs_is_function}.
                                    Thus $(n,v(n))\in v$ by \cref{function_apply_elim}.
                                    Thus $v(n)\in\ran{v}$ by \cref{ran_iff}.
                                    Thus $v(n)\in\dom{g}$ by \cref{subseteq}.
                                    Thus $n\in\preimg{v}{\dom{g}}$ by \cref{preimg_iff}.
                                    \end{subproof}
                                    Thus $\dom{v}\subseteq\preimg{v}{\dom{g}}$ by \cref{subseteq}.
                                \end{subproof}
                                Thus $\preimg{v}{\dom{g}} = \dom{v}$ by \cref{subseteq_antisymmetric}.
                            \end{subproof}
                            Thus $\dom{u} = \naturalsPlus$ by \cref{subseteq_antisymmetric}.
                        \end{subproof}
                        Show that for all $n\in\dom{u}$ we have $u(n)\in\naturalsPlus$.
                        \begin{subproof}
                            Fix $n\in\dom{u}$.
                            Thus $n\in\naturalsPlus$.
                            Thus $n\in\dom{v}$.
                            $v\in\funs{\naturalsPlus}{D}$ by \cref{sequence_in}.
                            Thus $\ran{v}\subseteq D$ by \cref{funs_ran}.
                            Thus $\ran{v}\subseteq\dom{g}$ by \cref{funs}.
                            $v$ is a function by \cref{funs_is_function}.
                            Thus $(n,v(n))\in v$ by \cref{function_apply_elim}.
                            Thus $v(n)\in\ran{v}$ by \cref{ran_iff}.
                            Thus $v(n)\in\dom{g}$ by \cref{subseteq}.
                            Thus $\dom{g} = D$ by \cref{funs}.
                            Thus $v(n)\in D$.
                            Thus $v(n)\in\naturalsPlus\times\naturalsPlus$ by \cref{subseteq}.
                            Take $m,k$ such that $v(n) = (m,k)$ and $m\in\naturalsPlus$ and $k\in\naturalsPlus$
                                by \cref{times_elem_is_tuple}.
                            Thus $\snd{(v(n))} = k$ by \cref{snd_eq}.
                            Thus $(v(n),\snd{(v(n))})\in g$.
                            Thus $g(v(n)) = \snd{(v(n))}$ by \cref{function_apply_intro}.
                            Thus $g(v(n)) = k$.
                            Thus $g(v(n))\in\naturalsPlus$.
                            Show that $u(n) = g(v(n))$.
                            \begin{subproof}
                                $v$ is a function by \cref{funs_is_function}.
                                Thus $(n,v(n))\in v$ by \cref{function_apply_elim}.
                                $g$ is a function by \cref{funs_is_function}.
                                Thus $(v(n),g(v(n)))\in g$ by \cref{function_apply_elim}.
                                Thus $(n,g(v(n)))\in g\circ v$ by \cref{circ_iff}.
                                Thus $(n,g(v(n)))\in u$.
                                Thus $u(n) = g(v(n))$ by \cref{function_apply_intro}.
                            \end{subproof}
                            Thus $u(n)\in\naturalsPlus$.
                        \end{subproof}
                        Thus $u$ is a function to $\naturalsPlus$.
                    \end{subproof}
                    Thus $u$ is a sequence in $\naturalsPlus$ by \cref{sequence_in}.
                \end{subproof}

                Show that for all $n\in\naturalsPlus$ we have $u(n) = \snd{v(n)}$.
                \begin{subproof}
                    Fix $n\in\naturalsPlus$.
                    Thus $n\in\dom{u}$.
                    Show that $u(n) = g(v(n))$.
                    \begin{subproof}
                        $v\in\funs{\naturalsPlus}{D}$ by \cref{sequence_in}.
                        Thus $v$ is a function by \cref{funs_is_function}.
                        Thus $\dom{v} = \naturalsPlus$ by \cref{funs}.
                        Thus $n\in\dom{v}$.
                        Thus $(n,v(n))\in v$ by \cref{function_apply_elim}.
                        Thus $v(n)\in D$ by \cref{funs_type_apply}.
                        Thus $\dom{g} = D$ by \cref{funs}.
                        Thus $v(n)\in\dom{g}$.
                        Thus $g$ is a function by \cref{funs_is_function}.
                        Thus $(v(n),g(v(n)))\in g$ by \cref{function_apply_elim}.
                        Thus $(n,g(v(n)))\in g\circ v$ by \cref{circ_iff}.
                        Thus $(n,g(v(n)))\in u$.
                        $u\in\funs{\naturalsPlus}{\naturalsPlus}$ by \cref{sequence_in}.
                        Thus $u$ is a function by \cref{funs_is_function}.
                        Thus $u(n) = g(v(n))$ by \cref{function_apply_intro}.
                    \end{subproof}
                    $v\in\funs{\naturalsPlus}{D}$ by \cref{sequence_in}.
                    Thus $v(n)\in D$ by \cref{funs_type_apply}.
                    Thus $(v(n),\snd{v(n)})\in g$.
                    Thus $g(v(n)) = \snd{v(n)}$ by \cref{function_apply_intro}.
                    Thus $u(n) = \snd{v(n)}$.
                \end{subproof}

                Show that for all $n\in\naturalsPlus$ we have $\fst{v(n)} = n$.
                \begin{subproof}
                    Let $P = \{n\in\naturalsPlus \mid \fst{v(n)} = n\}$.
                    Show that for all $n\in\naturalsPlus$ we have $n\in P$.
                    \begin{subproof}
                        Show that $P\subseteq\naturalsPlus$.
                        \begin{subproof}
                            Show that for all $n$ such that $n\in P$ we have $n\in\naturalsPlus$.
                            \begin{subproof}
                                Fix $n$.
                                Assume $n\in P$.
                                Then $n\in\naturalsPlus$.
                            \end{subproof}
                            Thus $P\subseteq\naturalsPlus$ by \cref{subseteq}.
                        \end{subproof}
                        $1\in\naturalsPlus$ by \cref{real_one_in_naturals}.
                        Show that $1\in P$.
                        \begin{subproof}
                            Thus $v(1) = a_0$.
                            Thus $v(1) = (1,\snd{a_0})$ by \cref{fst_eq,pair_eq_pair_of_proj}.
                            Thus $\fst{v(1)} = 1$ by \cref{fst_eq}.
                            Thus $1\in P$.
                        \end{subproof}
                        Show that for all $n\in P$ we have $n + 1\in P$.
                        \begin{subproof}
                            Fix $n$.
                            Assume $n\in P$.
                            Thus $\fst{v(n)} = n$.
                            $n + 1\in\naturalsPlus$ by \cref{real_naturals_closed_succ}.
                            Thus $n + 1\in\dom{v}$.
                            Thus $v(n + 1) = \apply{f}{v(n)}$ by \cref{iteration_sequence_naturalsplus}.
                            Thus $v(n)\mathrel{R} v(n + 1)$.
                            Take $n_0,k_0,n_1,k_1\in\naturalsPlus$ such that
                                $(v(n),v(n + 1)) = ((n_0,k_0),(n_1,k_1))$ and
                                $n_1 = n_0 + 1$ and $k_0 < k_1$.
                            Thus $\fst{v(n)} = n_0$ and $\fst{v(n + 1)} = n_1$ by \cref{pair_eq_iff,fst_eq}.
                            Thus $\fst{v(n + 1)} = \fst{v(n)} + 1$.
                            Thus $\fst{v(n + 1)} = n + 1$ by \cref{real_add_comm,real_add_ident}.
                            Thus $n + 1\in P$.
                        \end{subproof}
                        Thus for all $n\in\naturalsPlus$ we have $n\in P$ by \cref{real_naturals_induction}.
                    \end{subproof}
                    Thus for all $n\in\naturalsPlus$ we have $\fst{v(n)} = n$ by \cref{subseteq}.
                \end{subproof}

                Show that for all $n\in\naturalsPlus$ we have $\apply{s}{u(n)}\in\metricBall{d}{X}{x}{1 / n}$.
                \begin{subproof}
                    Fix $n\in\naturalsPlus$.
                    $v\in\funs{\naturalsPlus}{D}$ by \cref{sequence_in}.
                    Thus $v(n)\in D$ by \cref{funs_type_apply}.
                    Thus $\apply{s}{\snd{v(n)}}\in\metricBall{d}{X}{x}{1 / n}$.
                    Thus $\apply{s}{u(n)}\in\metricBall{d}{X}{x}{1 / n}$ by \cref{snd_eq}.
                \end{subproof}

                Show that for all $n\in\naturalsPlus$ we have $u(n) < u(n + 1)$.
                \begin{subproof}
                    Fix $n\in\naturalsPlus$.
                    Thus $v(n + 1) = \apply{f}{v(n)}$ by \cref{iteration_sequence_naturalsplus}.
                    Thus $v(n)\mathrel{R} v(n + 1)$.
                    Take $n_0,k_0,n_1,k_1\in\naturalsPlus$ such that
                        $(v(n),v(n + 1)) = ((n_0,k_0),(n_1,k_1))$ and
                        $n_1 = n_0 + 1$ and $k_0 < k_1$.
                    Thus $\snd{v(n)} = k_0$ and $\snd{v(n + 1)} = k_1$ by \cref{pair_eq_iff,snd_eq}.
                    Thus $\snd{v(n)} < \snd{v(n + 1)}$ by \cref{real_lt_subst_left,real_lt_subst_right}.
                    Thus $u(n) = \snd{v(n)}$.
                    Thus $n + 1\in\naturalsPlus$ by \cref{real_naturals_closed_succ}.
                    Thus $u(n + 1) = \snd{v(n + 1)}$.
                    Thus $u(n) < u(n + 1)$ by \cref{real_lt_subst_left,real_lt_subst_right}.
                \end{subproof}

                Show that $u$ is strictly increasing on $\naturalsPlus$.
                \begin{subproof}
                    Show that for all $m,n\in\naturalsPlus$ such that $m < n$ we have $u(m) < u(n)$.
                    \begin{subproof}
                        Let $T_0 = \{n\in\naturalsPlus \mid \text{for all $m\in\naturalsPlus$ such that $m < n$ we have $u(m) < u(n)$}\}$.
                        Show that for all $n\in\naturalsPlus$ we have $n\in T_0$.
                        \begin{subproof}
                            Show that $T_0\subseteq\naturalsPlus$.
                            \begin{subproof}
                                Show that for all $n$ such that $n\in T_0$ we have $n\in\naturalsPlus$.
                                \begin{subproof}
                                    Fix $n$.
                                    Assume $n\in T_0$.
                                    Then $n\in\naturalsPlus$.
                                \end{subproof}
                                Thus $T_0\subseteq\naturalsPlus$ by \cref{subseteq}.
                            \end{subproof}
                            $1\in\naturalsPlus$ by \cref{real_one_in_naturals}.
                            Show that $1\in T_0$.
                            \begin{subproof}
                                Show that for all $m\in\naturalsPlus$ such that $m < 1$ we have $u(m) < u(1)$.
                                \begin{subproof}
                                    Fix $m$.
                                    Assume $m\in\naturalsPlus$ and $m < 1$.
                                    Suppose not.
                                    $m\in\reals$ by \cref{real_naturals_subset_reals,subseteq}.
                                    $1\in\reals$ by \cref{real_naturals_subset_reals,real_one_in_naturals,subseteq}.
                                    $1 < m$ or $1 = m$ by \cref{real_one_leq_naturalsplus}.
                                    \begin{byCase}
                                    \caseOf{$1 = m$.}
                                        Thus $m < m$ by \cref{real_lt_subst_right}.
                                        Contradiction by \cref{real_lt_irreflexive}.
                                    \caseOf{$1 < m$.}
                                        Thus $1 < 1$ by \cref{real_lt_trans}.
                                        Contradiction by \cref{real_lt_irreflexive}.
                                    \end{byCase}
                                \end{subproof}
                                Thus $1\in T_0$.
                            \end{subproof}
                            Show that for all $n\in T_0$ we have $n + 1\in T_0$.
                            \begin{subproof}
                                Fix $n$.
                                Assume $n\in T_0$.
                                Show that for all $m\in\naturalsPlus$ such that $m < n + 1$ we have $u(m) < u(n + 1)$.
                                \begin{subproof}
                                    Fix $m$.
                                    Assume $m\in\naturalsPlus$ and $m < n + 1$.
                                    Suppose not.
                                    Thus $m < n$ or $m = n$ or $m = n + 1$ by \cref{real_succ_discrete}.
                                    \begin{byCase}
                                    \caseOf{$m = n$.}
                                        Thus $u(n) < u(n + 1)$.
                                        Thus $u(m) < u(n + 1)$ by \cref{real_lt_subst_left}.
                                        Contradiction.
                                    \caseOf{$m < n$.}
                                        Thus $u(m) < u(n)$ by \cref{subseteq}.
                                        Thus $u(n) < u(n + 1)$.
                                        $u\in\funs{\naturalsPlus}{\naturalsPlus}$ by \cref{sequence_in}.
                                        Thus $u(m)\in\naturalsPlus$ by \cref{funs_type_apply}.
                                        Thus $u(n)\in\naturalsPlus$ by \cref{funs_type_apply}.
                                        Thus $n + 1\in\naturalsPlus$ by \cref{real_naturals_closed_succ}.
                                        Thus $u(n + 1)\in\naturalsPlus$ by \cref{funs_type_apply}.
                                        Thus $u(m)\in\reals$ by \cref{real_naturals_subset_reals,subseteq}.
                                        Thus $u(n)\in\reals$ by \cref{real_naturals_subset_reals,subseteq}.
                                        Thus $u(n + 1)\in\reals$ by \cref{real_naturals_subset_reals,subseteq}.
                                        Thus $u(m) < u(n + 1)$ by \cref{real_lt_trans}.
                                        Contradiction.
                                    \caseOf{$m = n + 1$.}
                                        Thus $n + 1 < n + 1$ by \cref{real_lt_subst_left}.
                                        $n + 1\in\reals$ by \cref{real_naturals_subset_reals,subseteq}.
                                        Contradiction by \cref{real_lt_irreflexive}.
                                    \end{byCase}
                                \end{subproof}
                                Thus $n + 1\in T_0$.
                            \end{subproof}
                            Thus for all $n\in\naturalsPlus$ we have $n\in T_0$ by \cref{real_naturals_induction}.
                        \end{subproof}
                        Thus for all $m,n\in\naturalsPlus$ such that $m < n$ we have $u(m) < u(n)$ by \cref{subseteq}.
                    \end{subproof}
                    $u\in\funs{\naturalsPlus}{\naturalsPlus}$ by \cref{sequence_in}.
                    Thus $u$ is strictly increasing on $\naturalsPlus$ by \cref{strictly_increasing_on_set}.
                \end{subproof}

                Let $t = s\circ u$.
                Show that $t$ is a sequence in $X$.
                \begin{subproof}
                    Show that $t\in\funs{\naturalsPlus}{X}$.
                    \begin{subproof}
                        Show that $t$ is a function.
                        \begin{subproof}
                            Thus $s$ is right-unique and $u$ is right-unique by \cref{funs,sequence_in}.
                            Thus $t$ is a function by \cref{function_circ}.
                        \end{subproof}
                        Show that $\dom{t} = \naturalsPlus$.
                        \begin{subproof}
                            $u\in\funs{\naturalsPlus}{\naturalsPlus}$ by \cref{sequence_in}.
                            Thus $u$ is right-unique by \cref{funs}.
                            Thus $u$ is a relation by \cref{funs_is_relation}.
                            $s\in\funs{\naturalsPlus}{X}$ by \cref{sequence_in}.
                            Thus $s$ is right-unique by \cref{funs}.
                            Thus $s$ is a relation by \cref{funs_is_relation}.
                            Thus $\ran{u}\subseteq\naturalsPlus$ by \cref{funs_ran}.
                            Thus $\dom{s} = \naturalsPlus$ by \cref{funs}.
                            Thus $\ran{u}\subseteq\dom{s}$ by \cref{subseteq}.
                            Thus $\dom{t} = \preimg{u}{\dom{s}}$ by \cref{dom_of_circ}.
                            Thus $\dom{t} = \preimg{u}{\naturalsPlus}$.
                            Show that $\preimg{u}{\naturalsPlus} = \dom{u}$.
                            \begin{subproof}
                                Show that $\preimg{u}{\naturalsPlus}\subseteq\dom{u}$.
                                \begin{subproof}
                                    Thus $\preimg{u}{\naturalsPlus}\subseteq\dom{u}$ by \cref{preimg_subseteq_dom}.
                                \end{subproof}
                                Show that $\dom{u}\subseteq\preimg{u}{\naturalsPlus}$.
                                \begin{subproof}
                                    Show that for all $n$ such that $n\in\dom{u}$ we have $n\in\preimg{u}{\naturalsPlus}$.
                                    \begin{subproof}
                                        Fix $n$.
                                        Assume $n\in\dom{u}$.
                                        $u\in\funs{\naturalsPlus}{\naturalsPlus}$ by \cref{sequence_in}.
                                        Thus $\dom{u} = \naturalsPlus$ by \cref{funs}.
                                        Thus $n\in\naturalsPlus$.
                                        Thus $u(n)\in\naturalsPlus$ by \cref{funs_type_apply}.
                                        Thus $u$ is a function by \cref{funs_is_function}.
                                        Thus $(n,u(n))\in u$ by \cref{function_apply_elim}.
                                        Thus $n\in\preimg{u}{\naturalsPlus}$ by \cref{preimg_iff}.
                                    \end{subproof}
                                    Thus $\dom{u}\subseteq\preimg{u}{\naturalsPlus}$ by \cref{subseteq}.
                                \end{subproof}
                                Thus $\preimg{u}{\naturalsPlus} = \dom{u}$ by \cref{subseteq_antisymmetric}.
                            \end{subproof}
                            Thus $\dom{t} = \dom{u}$ by \cref{subseteq_antisymmetric}.
                            $u\in\funs{\naturalsPlus}{\naturalsPlus}$ by \cref{sequence_in}.
                            Thus $\dom{u} = \naturalsPlus$ by \cref{funs}.
                            Thus $\dom{t} = \naturalsPlus$.
                        \end{subproof}
                        Show that for all $n\in\dom{t}$ we have $t(n)\in X$.
                        \begin{subproof}
                            Fix $n\in\dom{t}$.
                            Thus $n\in\naturalsPlus$.
                            Thus $u(n)\in\naturalsPlus$ by \cref{sequence_in,funs_type_apply}.
                            Thus $u(n)\in\dom{s}$ by \cref{funs,sequence_in}.
                            Thus $s(u(n))\in X$ by \cref{funs_type_apply,sequence_in}.
                            $u$ is a function by \cref{funs_is_function,sequence_in}.
                            Thus $\dom{u} = \naturalsPlus$ by \cref{funs,sequence_in}.
                            Thus $n\in\dom{u}$.
                            Thus $(n,u(n))\in u$ by \cref{function_apply_elim}.
                            $s$ is a function by \cref{funs_is_function,sequence_in}.
                            Thus $(u(n),s(u(n)))\in s$ by \cref{function_apply_elim}.
                            Thus $(n,s(u(n)))\in s\circ u$ by \cref{circ_iff}.
                            Thus $(n,s(u(n)))\in t$.
                            Thus $t(n) = s(u(n))$ by \cref{function_apply_intro}.
                            Thus $t(n)\in X$.
                        \end{subproof}
                        Thus $t$ is a function to $X$.
                    \end{subproof}
                    Thus $t$ is a sequence in $X$ by \cref{sequence_in}.
                \end{subproof}

                Show that $t$ is a subsequence of $s$.
                \begin{subproof}
                    Show that for all $n\in\naturalsPlus$ we have $t(n) = \apply{s}{\apply{u}{n}}$.
                    \begin{subproof}
                        Fix $n\in\naturalsPlus$.
                        Thus $n\in\dom{t}$.
                        $u\in\funs{\naturalsPlus}{\naturalsPlus}$ by \cref{sequence_in}.
                        Thus $\dom{u} = \naturalsPlus$ by \cref{funs}.
                        Thus $n\in\dom{u}$.
                        Thus $u$ is a function by \cref{funs_is_function}.
                        Thus $(n,u(n))\in u$ by \cref{function_apply_elim}.
                        $s\in\funs{\naturalsPlus}{X}$ by \cref{sequence_in}.
                        Thus $\dom{s} = \naturalsPlus$ by \cref{funs}.
                        Thus $u(n)\in\naturalsPlus$ by \cref{funs_type_apply}.
                        Thus $u(n)\in\dom{s}$.
                        Thus $s$ is a function by \cref{funs_is_function}.
                        Thus $(u(n),s(u(n)))\in s$ by \cref{function_apply_elim}.
                        Thus $(n,s(u(n)))\in s\circ u$ by \cref{circ_iff}.
                        Thus $(n,s(u(n)))\in t$.
                        $t\in\funs{\naturalsPlus}{X}$ by \cref{sequence_in}.
                        Thus $t$ is a function by \cref{funs_is_function}.
                        Thus $t(n) = s(u(n))$ by \cref{function_apply_intro}.
                        Thus $t(n) = \apply{s}{\apply{u}{n}}$.
                    \end{subproof}
                    Thus $t$ is a subsequence of $s$ by \cref{subsequence}.
                \end{subproof}

                Show that $t$ converges to $x$ in $X$ with respect to $T$.
                \begin{subproof}
                    Show that for all $U$ if $U$ is a neighborhood of $x$ in $X$ with respect to $T$ then
                        there exists $N\in\naturalsPlus$ such that for all $n\in\naturalsPlus$ if $N\leq n$ then $t(n)\in U$.
                    \begin{subproof}
                        Fix $U$.
                        Assume $U$ is a neighborhood of $x$ in $X$ with respect to $T$.
                        Thus $U$ is open in $X$ with respect to $T$ by \cref{neighborhood}.
                        Thus $x\in U$ by \cref{neighborhood}.
                        Thus $U\in T$ by \cref{open_in}.
                        $T$ is a topology on $X$ by \cref{topology_on}.
                        Thus $T\subseteq\pow{X}$ by \cref{topology_on}.
                        Thus $U\in\pow{X}$ by \cref{subseteq}.
                        Thus $U\subseteq X$ by \cref{powerset_elim}.
                        Take $\epsilon\in\reals$ such that $0 < \epsilon$ and $\metricBall{d}{X}{x}{\epsilon}\subseteq U$
                            by \cref{metric_open_iff}.
                        Take $N\in\naturalsPlus$ such that $0 < 1 / N$ and $1 / N < \epsilon$ by \cref{real_small_reciprocal}.

                        Show that for all $m,n$ such that $m\in\naturalsPlus$ and $n\in\naturalsPlus$ and $m \leq n$
                            we have $1 / n \leq 1 / m$.
                        \begin{subproof}
                            Fix $m$.
                            Fix $n$.
                            Assume $m\in\naturalsPlus$ and $n\in\naturalsPlus$ and $m \leq n$.
                            $m\in\reals$ and $n\in\reals$ by \cref{real_naturals_subset_reals,subseteq}.
                            Thus $m < n$ or $m = n$ by \cref{real_leq_split}.
                            Show that $0 < m$ and $0 < n$.
                            \begin{subproof}
                                $1\in\naturalsPlus$ by \cref{real_one_in_naturals}.
                                $0 < 1$ by \cref{real_zero_lt_one}.
                                $m\in\reals$ and $n\in\reals$ by \cref{real_naturals_subset_reals,subseteq}.
                                $0\in\reals$ by \cref{real_add_ident}.
                                $1\in\reals$ by \cref{real_mul_ident}.

                                Show that $0 < m$.
                                \begin{subproof}
                                    $1 < m$ or $1 = m$ by \cref{real_one_leq_naturalsplus}.
                                    \begin{byCase}
                                    \caseOf{$1 < m$.}
                                        Thus $0 < m$ by \cref{real_lt_trans}.
                                    \caseOf{$1 = m$.}
                                        Thus $0 < m$ by \cref{real_lt_subst_left}.
                                    \end{byCase}
                                \end{subproof}

                                Show that $0 < n$.
                                \begin{subproof}
                                    $1 < n$ or $1 = n$ by \cref{real_one_leq_naturalsplus}.
                                    \begin{byCase}
                                    \caseOf{$1 < n$.}
                                        Thus $0 < n$ by \cref{real_lt_trans}.
                                    \caseOf{$1 = n$.}
                                        Thus $0 < n$ by \cref{real_lt_subst_left}.
                                    \end{byCase}
                                \end{subproof}
                            \end{subproof}
                            Suppose not.
                            Then $1 / m < 1 / n$.
                            $m\in\reals$ and $n\in\reals$ by \cref{real_naturals_subset_reals,subseteq}.
                            $0\in\reals$ by \cref{real_add_ident}.
                            $1\in\reals$ by \cref{real_mul_ident}.
                            $0 < 1$ by \cref{real_zero_lt_one}.
                            Thus $0 < m$.
                            Thus $m\neq 0$ by \cref{real_pos_nonzero}.
                            Thus $\inv{m}\in\reals$ by \cref{real_mul_inverse}.
                            Thus $0 < \inv{m}$ by \cref{real_inv_pos}.
                            Thus $0 \rmul \inv{m} < 1 \rmul \inv{m}$ by \cref{real_lt_mul_pos}.
                            Thus $0 < \inv{m}$ by \cref{real_mul_zero,real_mul_ident}.
                            Thus $1 / m = 1 \rmul \inv{m}$ by \cref{real_div}.
                            Thus $1 / m\in\reals$ by \cref{real_mul_closed}.
                            Thus $0 < n$.
                            Thus $n\neq 0$ by \cref{real_pos_nonzero}.
                            Thus $\inv{n}\in\reals$ by \cref{real_mul_inverse}.
                            Thus $1 / n = 1 \rmul \inv{n}$ by \cref{real_div}.
                            Thus $1 / n\in\reals$ by \cref{real_mul_closed}.
                            Show that $1 / m \rmul (m \rmul n) < 1 / n \rmul (m \rmul n)$.
                            \begin{subproof}
                                $m \rmul n\in\reals$ by \cref{real_mul_closed}.
                                Thus $0 \rmul n < m \rmul n$ by \cref{real_lt_mul_pos}.
                                Thus $0 < m \rmul n$ by \cref{real_mul_zero}.
                                Thus $1 / m \rmul (m \rmul n) < 1 / n \rmul (m \rmul n)$ by \cref{real_lt_mul_pos}.
                            \end{subproof}
                            Show that $1 / m \rmul (m \rmul n) = n$.
                            \begin{subproof}
                                $1 / m = 1 \rmul \inv{m}$ by \cref{real_div}.
                                Thus $1 / m \rmul (m \rmul n) = (1 \rmul \inv{m}) \rmul (m \rmul n)$.
                                Thus $(1 \rmul \inv{m}) \rmul (m \rmul n) = \inv{m} \rmul (m \rmul n)$
                                    by \cref{real_mul_ident}.
                                Thus $\inv{m} \rmul (m \rmul n) = (\inv{m} \rmul m) \rmul n$ by \cref{real_mul_assoc}.
                                $\inv{m} \rmul m = 1$ by \cref{real_mul_inverse,real_mul_comm}.
                                Thus $1 / m \rmul (m \rmul n) = 1 \rmul n$.
                                Thus $1 / m \rmul (m \rmul n) = n$ by \cref{real_mul_ident}.
                            \end{subproof}
                            Show that $1 / n \rmul (m \rmul n) = m$.
                            \begin{subproof}
                                $1 / n = 1 \rmul \inv{n}$ by \cref{real_div}.
                                Thus $1 / n \rmul (m \rmul n) = (1 \rmul \inv{n}) \rmul (m \rmul n)$.
                                Thus $(1 \rmul \inv{n}) \rmul (m \rmul n) = \inv{n} \rmul (m \rmul n)$
                                    by \cref{real_mul_ident}.
                                Thus $\inv{n} \rmul (m \rmul n) = (\inv{n} \rmul n) \rmul m$ by \cref{real_mul_assoc,real_mul_comm}.
                                $\inv{n} \rmul n = 1$ by \cref{real_mul_inverse,real_mul_comm}.
                                Thus $1 / n \rmul (m \rmul n) = 1 \rmul m$.
                                Thus $1 / n \rmul (m \rmul n) = m$ by \cref{real_mul_ident}.
                            \end{subproof}
                            Thus $n < m$ by \cref{real_lt_subst_left,real_lt_subst_right}.
                            Thus $m < n$ or $m = n$ by \cref{real_leq_split}.
                            \begin{byCase}
                            \caseOf{$m < n$.}
                                Thus $m < m$ by \cref{real_lt_trans}.
                                Contradiction by \cref{real_lt_irreflexive}.
                            \caseOf{$m = n$.}
                                Thus $n < n$ by \cref{real_lt_subst_left}.
                                Contradiction by \cref{real_lt_irreflexive}.
                            \end{byCase}
                        \end{subproof}

                        Show that for all $n\in\naturalsPlus$ if $N\leq n$ then $t(n)\in U$.
                        \begin{subproof}
                            Fix $n\in\naturalsPlus$.
                            Assume $N\leq n$.
                            Thus $1 / n \leq 1 / N$.
                            Show that $\metricBall{d}{X}{x}{1 / n}\subseteq \metricBall{d}{X}{x}{\epsilon}$.
                            \begin{subproof}
                                Show that for all $y$ such that $y\in\metricBall{d}{X}{x}{1 / n}$ we have
                                    $y\in\metricBall{d}{X}{x}{\epsilon}$.
                                \begin{subproof}
                                Fix $y$.
                                Assume $y\in\metricBall{d}{X}{x}{1 / n}$.
                                Thus $d((x,y)) < 1 / n$ by \cref{metric_ball}.
                                Thus $y\in X$ by \cref{metric_ball}.
                                $d\in\funs{X\times X}{\reals}$ by \cref{metric_on}.
                                $(x,y)\in X \times X$ by \cref{times_tuple_intro}.
                                Thus $d((x,y))\in\reals$ by \cref{funs_elim}.
                                $0\in\reals$ by \cref{real_add_ident}.
                                $1\in\reals$ by \cref{real_mul_ident}.
                                Show that $1 / n\in\reals$.
                                \begin{subproof}
                                    $n\in\reals$ by \cref{real_naturals_subset_reals,subseteq}.
                                    $0 < 1$ by \cref{real_zero_lt_one}.
                                    Show that $0 < n$.
                                    \begin{subproof}
                                        $1 < n$ or $1 = n$ by \cref{real_one_leq_naturalsplus}.
                                        \begin{byCase}
                                        \caseOf{$1 < n$.}
                                            Thus $0 < n$ by \cref{real_lt_trans}.
                                        \caseOf{$1 = n$.}
                                            Thus $0 < n$ by \cref{real_lt_subst_right}.
                                        \end{byCase}
                                    \end{subproof}
                                    Thus $n \neq 0$ by \cref{real_pos_nonzero}.
                                    Thus $\inv{n}\in\reals$ by \cref{real_mul_inverse}.
                                    $1 / n = 1 \rmul \inv{n}$ by \cref{real_div}.
                                    Thus $1 / n\in\reals$ by \cref{real_mul_closed}.
                                \end{subproof}
                                Show that $1 / N\in\reals$.
                                \begin{subproof}
                                    $N\in\reals$ by \cref{real_naturals_subset_reals,subseteq}.
                                    $0 < 1$ by \cref{real_zero_lt_one}.
                                    Show that $0 < N$.
                                    \begin{subproof}
                                        $1 < N$ or $1 = N$ by \cref{real_one_leq_naturalsplus}.
                                        \begin{byCase}
                                        \caseOf{$1 < N$.}
                                            Thus $0 < N$ by \cref{real_lt_trans}.
                                        \caseOf{$1 = N$.}
                                            Thus $0 < N$ by \cref{real_lt_subst_right}.
                                        \end{byCase}
                                    \end{subproof}
                                    Thus $N \neq 0$ by \cref{real_pos_nonzero}.
                                    Thus $\inv{N}\in\reals$ by \cref{real_mul_inverse}.
                                    $1 / N = 1 \rmul \inv{N}$ by \cref{real_div}.
                                    Thus $1 / N\in\reals$ by \cref{real_mul_closed}.
                                \end{subproof}
                                Show that $d((x,y)) < 1 / N$.
                                \begin{subproof}
                                    Thus $1 / n \leq 1 / N$.
                                    Thus $1 / n < 1 / N$ or $1 / n = 1 / N$ by \cref{real_leq_split}.
                                    \begin{byCase}
                                    \caseOf{$1 / n < 1 / N$.}
                                        Thus $d((x,y)) < 1 / N$ by \cref{real_lt_trans}.
                                    \caseOf{$1 / n = 1 / N$.}
                                        Thus $d((x,y)) < 1 / N$ by \cref{real_lt_subst_right}.
                                    \end{byCase}
                                \end{subproof}
                                Thus $d((x,y)) < \epsilon$ by \cref{real_lt_trans}.
                                Thus $y\in\metricBall{d}{X}{x}{\epsilon}$ by \cref{metric_ball}.
                            \end{subproof}
                                Thus $\metricBall{d}{X}{x}{1 / n}\subseteq \metricBall{d}{X}{x}{\epsilon}$ by \cref{subseteq}.
                            \end{subproof}
                            Thus $\apply{s}{\apply{u}{n}}\in\metricBall{d}{X}{x}{1 / n}$.
                            Thus $t(n) = \apply{s}{\apply{u}{n}}$.
                            Thus $t(n)\in\metricBall{d}{X}{x}{1 / n}$.
                            Thus $t(n)\in\metricBall{d}{X}{x}{\epsilon}$ by \cref{subseteq}.
                            Thus $t(n)\in U$ by \cref{subseteq}.
                        \end{subproof}
                        Thus there exists $N\in\naturalsPlus$ such that for all $n\in\naturalsPlus$ if $N\leq n$ then $t(n)\in U$.
                    \end{subproof}
                    Thus $t$ converges to $x$ in $X$ with respect to $T$ by \cref{converges_to}.
                \end{subproof}
                Thus there exist $t,x$ such that
                    $t$ is a sequence in $X$ and
                    $t$ is a subsequence of $s$ and
                    $x\in X$ and
                    $t$ converges to $x$ in $X$ with respect to $T$.
            \end{byCase}
        \end{subproof}
        Thus $X$ is sequentially compact with respect to $T$ by \cref{sequentially_compact}.
    \end{subproof}

    Show that if $X$ is compact with respect to $T$ then
        $X$ is limit point compact with respect to $T$ and
        $X$ is sequentially compact with respect to $T$.
    \begin{subproof}
        Assume $X$ is compact with respect to $T$.
        Thus $X$ is limit point compact with respect to $T$ by \cref{compact_implies_limit_point_compact}.
        Thus $X$ is sequentially compact with respect to $T$.
        Thus $X$ is limit point compact with respect to $T$ and
            $X$ is sequentially compact with respect to $T$.
    \end{subproof}

    Show that if $X$ is limit point compact with respect to $T$ and
        $X$ is sequentially compact with respect to $T$ then
        $X$ is compact with respect to $T$.
    \begin{subproof}
        Assume $X$ is limit point compact with respect to $T$ and
            $X$ is sequentially compact with respect to $T$.
        Thus $X$ is sequentially compact with respect to $T$.
        Thus $X$ is compact with respect to $T$.
    \end{subproof}
\end{proof}
    
% from §29 helper definition 1: compact subspace
\begin{definition}\label{compact_subspace}
    $C$ is a compact subspace of $X$ with respect to $T$ iff
        $C\subseteq X$ and
        $C$ is compact with respect to $\subspaceTopology{T}{C}$.
\end{definition}

% from §29 Item 1: locally compact at a point
\begin{definition}\label{locally_compact_at_point}
    $X$ is locally compact at $x$ with respect to $T$ iff
        $x\in X$ and
        there exists $C$ such that
            $C$ is a compact subspace of $X$ with respect to $T$ and
            there exists $U$ such that
                $U$ is a neighborhood of $x$ in $X$ with respect to $T$ and
                $U\subseteq C$.
\end{definition}

% from §29 Item 2: locally compact
\begin{definition}\label{locally_compact}
    $X$ is locally compact with respect to $T$ iff
        for all $x\in X$ we have $X$ is locally compact at $x$ with respect to $T$.
\end{definition}

% from §29 Item 3: compact implies locally compact
\begin{lemma}\label{compact_implies_locally_compact}
    Assume $T$ is a topology on $X$.
    Suppose $X$ is compact with respect to $T$.
    Then $X$ is locally compact with respect to $T$.
\end{lemma}
\begin{proof}
    Show that for all $x\in X$ we have $X$ is locally compact at $x$ with respect to $T$.
    \begin{subproof}
        Fix $x\in X$.
        Let $T_X = \subspaceTopology{T}{X}$.
        $X\subseteq X$ by \cref{subseteq_refl}.
        Thus $T_X$ is a topology on $X$ by \cref{subspace_topology_is_topology}.
        Show that $T_X = T$.
        \begin{subproof}
            Show that $T\subseteq T_X$.
            \begin{subproof}
                Show that for all $U$ such that $U\in T$ we have $U\in T_X$.
                \begin{subproof}
                    Fix $U$.
                    Assume $U\in T$.
                    $T\subseteq\pow{X}$ by \cref{topology_on}.
                    Thus $U\subseteq X$ by \cref{powerset_elim,subseteq}.
                    Thus $U = X\inter U$ by \cref{inter_absorb_supseteq_left}.
                    Thus $U\in T_X$ by \cref{subspace_topology}.
                \end{subproof}
                Thus $T\subseteq T_X$ by \cref{subseteq}.
            \end{subproof}
            Show that $T_X\subseteq T$.
            \begin{subproof}
                Show that for all $U$ such that $U\in T_X$ we have $U\in T$.
                \begin{subproof}
                    Fix $U$.
                    Assume $U\in T_X$.
                    Take $V\in T$ such that $U = X\inter V$ by \cref{subspace_topology}.
                    $T\subseteq\pow{X}$ by \cref{topology_on}.
                    Thus $V\subseteq X$ by \cref{powerset_elim,subseteq}.
                    Thus $U = V$ by \cref{inter_absorb_supseteq_left}.
                    Thus $U\in T$.
                \end{subproof}
                Thus $T_X\subseteq T$ by \cref{subseteq}.
            \end{subproof}
            Thus $T_X = T$ by \cref{subseteq_antisymmetric}.
        \end{subproof}
        Show that $X$ is compact with respect to $T_X$.
        \begin{subproof}
            Show that for all $A$ such that $A$ is an open covering of $X$ with respect to $T_X$
                there exists $B$ such that
                    $B\subseteq A$ and
                    $B$ is finite and
                    $B$ is a covering of $X$.
            \begin{subproof}
                Fix $A$.
                Assume $A$ is an open covering of $X$ with respect to $T_X$.
                Show that $A$ is an open covering of $X$ with respect to $T$.
                \begin{subproof}
                    Show that for all $U$ such that $U\in A$ we have $U$ is open in $X$ with respect to $T$.
                    \begin{subproof}
                        Fix $U$.
                        Assume $U\in A$.
                        $U$ is open in $X$ with respect to $T_X$ by \cref{open_covering}.
                        Thus $U\in T_X$ by \cref{open_in}.
                        Take $V\in T$ such that $U = X\inter V$ by \cref{subspace_topology}.
                        $T\subseteq\pow{X}$ by \cref{topology_on}.
                        Thus $V\subseteq X$ by \cref{powerset_elim,subseteq}.
                        Thus $U = V$ by \cref{inter_absorb_supseteq_left}.
                        Thus $U\in T$.
                        Thus $U$ is open in $X$ with respect to $T$ by \cref{open_in}.
                    \end{subproof}
                    Thus $A$ is an open covering of $X$ with respect to $T$ by \cref{open_covering}.
                \end{subproof}
                Thus there exists $B$ such that
                    $B\subseteq A$ and
                    $B$ is finite and
                    $B$ is a covering of $X$ by \cref{compact}.
            \end{subproof}
            Thus $X$ is compact with respect to $T_X$ by \cref{compact}.
        \end{subproof}
        $X\subseteq X$ by \cref{subseteq_refl}.
        Thus $X$ is a compact subspace of $X$ with respect to $T$ by \cref{compact_subspace}.
        $X\in T$ by \cref{topology_on}.
        Thus $X$ is open in $X$ with respect to $T$ by \cref{open_in}.
        Thus $X$ is a neighborhood of $x$ in $X$ with respect to $T$ by \cref{neighborhood}.
        Thus $X$ is locally compact at $x$ with respect to $T$ by \cref{locally_compact_at_point}.
    \end{subproof}
    Thus $X$ is locally compact with respect to $T$ by \cref{locally_compact}.
\end{proof}

% from §29 helper lemma 1: set not member of itself
\begin{lemma}\label{set_notin_self}
    For all $x$ we have $x\notin x$.
\end{lemma}
\begin{proof}
    Fix $x$.
    Suppose not.
    Then $x\in x$.
    $x\in\{x\}$ by \cref{singleton_intro}.
    Show that $\{x\}\neq\emptyset$.
    \begin{subproof}
        Suppose not.
        Then $\{x\} = \emptyset$.
        Thus $x\in\emptyset$.
        Contradiction by \cref{emptyset}.
    \end{subproof}
    Take $a\in\{x\}$ such that for all $y\in a$ we have $y\notin\{x\}$ by \cref{foundation}.
    $a = x$ by \cref{singleton_elim}.
    Thus for all $y\in x$ we have $y\notin\{x\}$.
    Thus $x\notin\{x\}$.
    Contradiction by \cref{singleton_intro}.
\end{proof}

% from §29 helper lemma 2: empty set is compact
\begin{lemma}\label{emptyset_compact}
    Assume $T$ is a topology on $X$.
    Then $\emptyset$ is compact with respect to $T$.
\end{lemma}
\begin{proof}
    Show that for all $A$ such that $A$ is an open covering of $\emptyset$ with respect to $T$
        there exists $B$ such that
            $B\subseteq A$ and
            $B$ is finite and
            $B$ is a covering of $\emptyset$.
    \begin{subproof}
        Fix $A$.
        Assume $A$ is an open covering of $\emptyset$ with respect to $T$.
        Let $B = \emptyset$.
        $B\subseteq A$ by \cref{emptyset_subseteq}.
        $B$ is finite by \cref{emptyset_is_finite}.
        $\emptyset\subseteq\unions{B}$ by \cref{emptyset_subseteq}.
        Thus $B$ is a covering of $\emptyset$ by \cref{covering}.
    \end{subproof}
    Thus $\emptyset$ is compact with respect to $T$ by \cref{compact}.
\end{proof}

% from §29 helper lemma 3: intersection of compact subspaces is compact
\begin{lemma}\label{compact_subspace_inters}
    Assume $T$ is a topology on $X$.
    Assume $X$ is Hausdorff with respect to $T$.
    Suppose $C\subseteq\pow{X}$.
    Suppose $C$ is inhabited.
    Suppose for all $K$ such that $K\in C$ we have $K$ is a compact subspace of $X$ with respect to $T$.
    Then $\inters{C}$ is a compact subspace of $X$ with respect to $T$.
\end{lemma}
\begin{proof}
    Take $K_0$ such that $K_0\in C$.
    Show that for all $K$ such that $K\in C$ we have $K$ is closed in $X$ with respect to $T$.
    \begin{subproof}
        Fix $K$.
        Assume $K\in C$.
        Thus $K$ is a compact subspace of $X$ with respect to $T$.
        Thus $K\subseteq X$ by \cref{compact_subspace}.
        Let $T_K = \subspaceTopology{T}{K}$.
        Thus $K$ is compact with respect to $T_K$ by \cref{compact_subspace}.
        Thus $K$ is closed in $X$ with respect to $T$ by \cref{compact_subspace_closed}.
    \end{subproof}
    Thus $\inters{C}$ is closed in $X$ with respect to $T$ by \cref{closed_inters}.
    Show that $\inters{C}\subseteq K_0$.
    \begin{subproof}
        Show that for all $x$ such that $x\in\inters{C}$ we have $x\in K_0$.
        \begin{subproof}
            Fix $x$.
            Assume $x\in\inters{C}$.
            Thus $x\in K_0$ by \cref{inters}.
        \end{subproof}
        Thus $\inters{C}\subseteq K_0$ by \cref{subseteq}.
    \end{subproof}
    Let $T_{K0} = \subspaceTopology{T}{K_0}$.
    Let $T_C = \subspaceTopology{T}{\inters{C}}$.
    $K_0$ is a compact subspace of $X$ with respect to $T$ by \cref{compact_subspace}.
    Thus $K_0$ is compact with respect to $T_{K0}$ by \cref{compact_subspace}.
    $K_0\subseteq X$ by \cref{compact_subspace}.
    Thus $K_0$ is a subspace of $X$ with respect to $T$ by \cref{subspace_of}.
    Thus $T_{K0}$ is a topology on $K_0$ by \cref{subspace_topology_is_topology}.
    Show that $\inters{C}$ is closed in $K_0$ with respect to $T_{K0}$.
    \begin{subproof}
        Let $D = \inters{C}$.
        $D$ is closed in $X$ with respect to $T$.
        Thus $D\inter K_0$ is closed in $K_0$ with respect to $T_{K0}$ by \cref{closed_in_subspace_iff}.
        Show that $D\inter K_0 = \inters{C}$.
        \begin{subproof}
            Show that $D\inter K_0\subseteq \inters{C}$.
            \begin{subproof}
                Show that for all $x$ such that $x\in D\inter K_0$ we have $x\in\inters{C}$.
                \begin{subproof}
                    Fix $x$.
                    Assume $x\in D\inter K_0$.
                    Thus $x\in D$ by \cref{inter}.
                    Thus $x\in\inters{C}$.
                \end{subproof}
                Thus $D\inter K_0\subseteq \inters{C}$ by \cref{subseteq}.
            \end{subproof}
            Show that $\inters{C}\subseteq D\inter K_0$.
            \begin{subproof}
                Show that for all $x$ such that $x\in\inters{C}$ we have $x\in D\inter K_0$.
                \begin{subproof}
                    Fix $x$.
                    Assume $x\in\inters{C}$.
                    Thus $x\in D$.
                    Thus $x\in K_0$ by \cref{inters}.
                    Thus $x\in D\inter K_0$ by \cref{inter_intro}.
                \end{subproof}
                Thus $\inters{C}\subseteq D\inter K_0$ by \cref{subseteq}.
            \end{subproof}
            Thus $D\inter K_0 = \inters{C}$ by \cref{subseteq_antisymmetric}.
        \end{subproof}
        Thus $\inters{C}$ is closed in $K_0$ with respect to $T_{K0}$ by \cref{closed_in_subspace_iff}.
    \end{subproof}
    Let $T_{KC} = \subspaceTopology{T_{K0}}{\inters{C}}$.
    Thus $\inters{C}$ is compact with respect to $T_{KC}$ by \cref{closed_subspace_compact}.
    $\inters{C}\subseteq K_0$.
    Thus $T_{KC} = T_C$ by \cref{subspace_subspace_topology}.
    Thus $\inters{C}$ is compact with respect to $T_C$.
    $K_0\subseteq X$ by \cref{compact_subspace}.
    Thus $\inters{C}\subseteq X$ by \cref{subseteq_transitive}.
    Thus $\inters{C}$ is a compact subspace of $X$ with respect to $T$ by \cref{compact_subspace}.
\end{proof}

% from §29 helper lemma 4: union of compact subspaces is compact
\begin{lemma}\label{compact_subspace_union}
    Assume $T$ is a topology on $X$.
    Suppose $A$ is a compact subspace of $X$ with respect to $T$.
    Suppose $B$ is a compact subspace of $X$ with respect to $T$.
    Let $C = A\union B$.
    Then $C$ is a compact subspace of $X$ with respect to $T$.
\end{lemma}
\begin{proof}
    $A\subseteq X$ and $B\subseteq X$ by \cref{compact_subspace}.
    Thus $C\subseteq X$ by \cref{union_subsets_is_subset}.
    Let $T_C = \subspaceTopology{T}{C}$.
    Show that $C$ is compact with respect to $T_C$.
    \begin{subproof}
        Show that for all $U$ such that $U$ is an open covering of $C$ with respect to $T_C$
            there exists $D$ such that
                $D\subseteq U$ and
                $D$ is finite and
                $D$ is a covering of $C$.
        \begin{subproof}
            Fix $U$.
            Assume $U$ is an open covering of $C$ with respect to $T_C$.
            Thus $U$ is a covering of $C$ by \cref{open_covering}.
            Let $V = \{W\in T \mid \text{there exists $U_0\in U$ such that $U_0 = C\inter W$}\}$.
            Show that $V$ is a covering of $A$.
            \begin{subproof}
                Show that $A\subseteq\unions{V}$.
                \begin{subproof}
                    Show that for all $a$ such that $a\in A$ we have $a\in\unions{V}$.
                    \begin{subproof}
                        Fix $a$.
                        Assume $a\in A$.
                        Thus $a\in C$ by \cref{union_intro_left}.
                        $C\subseteq\unions{U}$ by \cref{covering}.
                        Thus $a\in\unions{U}$ by \cref{elem_subseteq}.
                        Take $U_0\in U$ such that $a\in U_0$ by \cref{unions_iff}.
                        $U_0$ is open in $C$ with respect to $T_C$ by \cref{open_covering}.
                        Thus $U_0\in T_C$ by \cref{open_in}.
                        Take $V_0\in T$ such that $U_0 = C\inter V_0$ by \cref{subspace_topology}.
                        Show that $V_0\in V$.
                        \begin{subproof}
                            $V_0\in T$.
                            $U_0\in U$ and $U_0 = C\inter V_0$.
                            Thus there exists $U_1\in U$ such that $U_1 = C\inter V_0$.
                            Thus $V_0\in V$.
                        \end{subproof}
                        $U_0\subseteq V_0$ by \cref{inter_lower_right}.
                        Thus $a\in V_0$ by \cref{elem_subseteq}.
                        Thus $a\in\unions{V}$ by \cref{unions_iff}.
                    \end{subproof}
                    Thus $A\subseteq\unions{V}$ by \cref{subseteq}.
                \end{subproof}
                Thus $V$ is a covering of $A$ by \cref{covering}.
            \end{subproof}
            Show that for all $V_0$ such that $V_0\in V$ we have $V_0$ is open in $X$ with respect to $T$.
            \begin{subproof}
                Fix $V_0$.
                Assume $V_0\in V$.
                Thus $V_0\in T$.
                Thus $V_0$ is open in $X$ with respect to $T$ by \cref{open_in}.
            \end{subproof}
            $A$ is a subspace of $X$ with respect to $T$ by \cref{subspace_of,compact_subspace}.
            Let $T_A = \subspaceTopology{T}{A}$.
            $A$ is compact with respect to $T_A$ by \cref{compact_subspace}.
            Take $V_A$ such that
                $V_A\subseteq V$ and
                $V_A$ is finite and
                $V_A$ is a covering of $A$
                by \cref{compact_subspace_open_covering}.

            Show that $V$ is a covering of $B$.
            \begin{subproof}
                Show that $B\subseteq\unions{V}$.
                \begin{subproof}
                    Show that for all $b$ such that $b\in B$ we have $b\in\unions{V}$.
                    \begin{subproof}
                        Fix $b$.
                        Assume $b\in B$.
                        Thus $b\in C$ by \cref{union_intro_right}.
                        $C\subseteq\unions{U}$ by \cref{covering}.
                        Thus $b\in\unions{U}$ by \cref{elem_subseteq}.
                        Take $U_0\in U$ such that $b\in U_0$ by \cref{unions_iff}.
                        $U_0$ is open in $C$ with respect to $T_C$ by \cref{open_covering}.
                        Thus $U_0\in T_C$ by \cref{open_in}.
                        Take $V_0\in T$ such that $U_0 = C\inter V_0$ by \cref{subspace_topology}.
                        Show that $V_0\in V$.
                        \begin{subproof}
                            $V_0\in T$.
                            $U_0\in U$ and $U_0 = C\inter V_0$.
                            Thus there exists $U_1\in U$ such that $U_1 = C\inter V_0$.
                            Thus $V_0\in V$.
                        \end{subproof}
                        $U_0\subseteq V_0$ by \cref{inter_lower_right}.
                        Thus $b\in V_0$ by \cref{elem_subseteq}.
                        Thus $b\in\unions{V}$ by \cref{unions_iff}.
                    \end{subproof}
                    Thus $B\subseteq\unions{V}$ by \cref{subseteq}.
                \end{subproof}
                Thus $V$ is a covering of $B$ by \cref{covering}.
            \end{subproof}
            $B$ is a subspace of $X$ with respect to $T$ by \cref{subspace_of,compact_subspace}.
            Let $T_B = \subspaceTopology{T}{B}$.
            $B$ is compact with respect to $T_B$ by \cref{compact_subspace}.
            Take $V_B$ such that
                $V_B\subseteq V$ and
                $V_B$ is finite and
                $V_B$ is a covering of $B$
                by \cref{compact_subspace_open_covering}.

            Let $V_0 = V_A\union V_B$.
            $V_0$ is finite by \cref{union_of_finite_is_finite}.
            Show that $V_0$ is a covering of $C$.
            \begin{subproof}
                Show that $C\subseteq\unions{V_0}$.
                \begin{subproof}
                    Show that for all $x$ such that $x\in C$ we have $x\in\unions{V_0}$.
                    \begin{subproof}
                        Fix $x$.
                        Assume $x\in C$.
                        Thus $x\in A$ or $x\in B$ by \cref{union_iff}.
                        \begin{byCase}
                        \caseOf{$x\in A$.}
                            $A\subseteq\unions{V_A}$ by \cref{covering}.
                            Thus $x\in\unions{V_A}$ by \cref{elem_subseteq}.
                            Take $v\in V_A$ such that $x\in v$ by \cref{unions_iff}.
                            Thus $v\in V_0$ by \cref{union_intro_left}.
                            Thus $x\in\unions{V_0}$ by \cref{unions_iff}.
                        \caseOf{$x\in B$.}
                            $B\subseteq\unions{V_B}$ by \cref{covering}.
                            Thus $x\in\unions{V_B}$ by \cref{elem_subseteq}.
                            Take $v\in V_B$ such that $x\in v$ by \cref{unions_iff}.
                            Thus $v\in V_0$ by \cref{union_intro_right}.
                            Thus $x\in\unions{V_0}$ by \cref{unions_iff}.
                        \end{byCase}
                    \end{subproof}
                    Thus $C\subseteq\unions{V_0}$ by \cref{subseteq}.
                \end{subproof}
                Thus $V_0$ is a covering of $C$ by \cref{covering}.
            \end{subproof}

            Let $g = \{(v, C\inter v) \mid v\in V_0\}$.
            Show that $g$ is a function from $V_0$ to $\pow{C}$.
            \begin{subproof}
                Show that $g$ is a function.
                \begin{subproof}
                    Show that $g$ is a relation.
                    \begin{subproof}
                        Show that for all $w$ such that $w\in g$ there exist $x,y$ such that $w = (x,y)$.
                        \begin{subproof}
                            Fix $w$.
                            Assume $w\in g$.
                            Take $v\in V_0$ such that $w = (v,C\inter v)$.
                            Thus there exists $x,y$ such that $w = (x,y)$.
                        \end{subproof}
                        Thus $g$ is a relation by \cref{relation}.
                    \end{subproof}
                    Show that $g$ is right-unique.
                    \begin{subproof}
                        Show that for all $x,y,z$ such that $(x,y)\in g$ and $(x,z)\in g$ we have $y = z$.
                        \begin{subproof}
                            Fix $x$.
                            Fix $y$.
                            Fix $z$.
                            Assume $(x,y)\in g$ and $(x,z)\in g$.
                            Take $v_1\in V_0$ such that $(x,y) = (v_1,C\inter v_1)$.
                            Take $v_2\in V_0$ such that $(x,z) = (v_2,C\inter v_2)$.
                            Thus $x = v_1$ and $y = C\inter v_1$ by \cref{pair_eq_iff}.
                            Thus $x = v_2$ and $z = C\inter v_2$ by \cref{pair_eq_iff}.
                            Thus $v_1 = v_2$.
                            Thus $y = z$.
                        \end{subproof}
                        Thus $g$ is right-unique by \cref{rightunique}.
                    \end{subproof}
                    Thus $g$ is a function.
                \end{subproof}
                Show that $\dom{g} = V_0$.
                \begin{subproof}
                    Show that $V_0\subseteq\dom{g}$.
                    \begin{subproof}
                        Show that for all $v$ such that $v\in V_0$ we have $v\in\dom{g}$.
                        \begin{subproof}
                            Fix $v$.
                            Assume $v\in V_0$.
                            Thus $(v,C\inter v)\in g$.
                            Thus $v\in\dom{g}$ by \cref{dom_iff}.
                        \end{subproof}
                        Thus $V_0\subseteq\dom{g}$ by \cref{subseteq}.
                    \end{subproof}
                    Show that $\dom{g}\subseteq V_0$.
                    \begin{subproof}
                        Show that for all $v$ such that $v\in\dom{g}$ we have $v\in V_0$.
                        \begin{subproof}
                            Fix $v$.
                            Assume $v\in\dom{g}$.
                            Take $y$ such that $(v,y)\in g$ by \cref{dom_iff}.
                            Take $v_1\in V_0$ such that $(v,y) = (v_1,C\inter v_1)$.
                            Thus $v = v_1$ by \cref{pair_eq_iff}.
                            Thus $v\in V_0$.
                        \end{subproof}
                        Thus $\dom{g}\subseteq V_0$ by \cref{subseteq}.
                    \end{subproof}
                    Thus $\dom{g} = V_0$ by \cref{subseteq_antisymmetric}.
                \end{subproof}
                Show that $g$ is a function to $\pow{C}$.
                \begin{subproof}
                    Show that for all $v$ such that $v\in\dom{g}$ we have $g(v)\in\pow{C}$.
                    \begin{subproof}
                        Fix $v$.
                        Assume $v\in\dom{g}$.
                        Thus $v\in V_0$.
                        Thus $g(v) = C\inter v$ by \cref{function_apply_iff}.
                        Thus $g(v)\subseteq C$ by \cref{inter_lower_left}.
                        Thus $g(v)\in\pow{C}$ by \cref{powerset_intro}.
                    \end{subproof}
                    Thus $g$ is a function to $\pow{C}$.
                \end{subproof}
                Thus $g\in\funs{V_0}{\pow{C}}$ by \cref{funs_intro}.
            \end{subproof}
            Let $D = \{w \mid \exists v\in V_0. w = C\inter v\}$.
            $g\in\funs{V_0}{\pow{C}}$.
            Thus $g$ is a function by \cref{funs_is_function}.
            Thus $g$ is right-unique by \cref{function_rightunique}.
            Thus $g$ is a relation by \cref{funs_is_relation}.
            Thus $\dom{g} = V_0$ by \cref{funs}.
            Let $E = \img{g}{V_0}$.
            Thus $E$ is finite by \cref{finite_image_function}.
            Show that $D\subseteq E$.
            \begin{subproof}
                Show that for all $W$ such that $W\in D$ we have $W\in E$.
                \begin{subproof}
                    Fix $W$.
                    Assume $W\in D$.
                    Take $v\in V_0$ such that $W = C\inter v$.
                    Thus $(v,W)\in g$.
                    Thus $W\in E$ by \cref{img_iff}.
                \end{subproof}
                Thus $D\subseteq E$ by \cref{subseteq}.
            \end{subproof}
            Thus $D$ is finite by \cref{subseteq_finite_is_finite}.
            Show that $D\subseteq U$.
            \begin{subproof}
                Show that for all $W$ such that $W\in D$ we have $W\in U$.
                \begin{subproof}
                    Fix $W$.
                    Assume $W\in D$.
                    Take $v\in V_0$ such that $W = C\inter v$.
                    $V_A\subseteq V$ and $V_B\subseteq V$.
                    Thus $V_0\subseteq V$ by \cref{union_subsets_is_subset}.
                    Thus $v\in V$ by \cref{elem_subseteq}.
                    Take $U_0\in U$ such that $U_0 = C\inter v$.
                    Thus $W\in U$.
                \end{subproof}
                Thus $D\subseteq U$ by \cref{subseteq}.
            \end{subproof}
            Show that $D$ is a covering of $C$.
            \begin{subproof}
                Show that $C\subseteq\unions{D}$.
                \begin{subproof}
                    Show that for all $x$ such that $x\in C$ we have $x\in\unions{D}$.
                    \begin{subproof}
                        Fix $x$.
                        Assume $x\in C$.
                        $C\subseteq\unions{V_0}$ by \cref{covering}.
                        Thus $x\in\unions{V_0}$ by \cref{elem_subseteq}.
                        Take $v\in V_0$ such that $x\in v$ by \cref{unions_iff}.
                        Let $W = C\inter v$.
                        Thus $W\in D$.
                        Thus $x\in W$ by \cref{inter_intro}.
                        Thus $x\in\unions{D}$ by \cref{unions_iff}.
                    \end{subproof}
                    Thus $C\subseteq\unions{D}$ by \cref{subseteq}.
                \end{subproof}
                Thus $D$ is a covering of $C$ by \cref{covering}.
            \end{subproof}
            Thus there exists $D$ such that
                $D\subseteq U$ and
                $D$ is finite and
                $D$ is a covering of $C$.
        \end{subproof}
        Thus $C$ is compact with respect to $T_C$ by \cref{compact}.
    \end{subproof}
    Thus $C$ is a compact subspace of $X$ with respect to $T$ by \cref{compact_subspace}.
\end{proof}

% from §29 helper lemma 5: compact subspace minus open set is compact
\begin{lemma}\label{compact_subspace_minus_open}
    Assume $T$ is a topology on $X$.
    Suppose $C$ is a compact subspace of $X$ with respect to $T$.
    Suppose $U$ is open in $X$ with respect to $T$.
    Let $D = C\setminus U$.
    Then $D$ is a compact subspace of $X$ with respect to $T$.
\end{lemma}
\begin{proof}
    $C\subseteq X$ by \cref{compact_subspace}.
    Thus $D\subseteq X$ by \cref{setminus_subseteq,subseteq_transitive}.
    Let $T_C = \subspaceTopology{T}{C}$.
    Let $T_D = \subspaceTopology{T}{D}$.
    Let $T_{CD} = \subspaceTopology{T_C}{D}$.
    $C$ is a subspace of $X$ with respect to $T$ by \cref{subspace_of,compact_subspace}.
    $T$ is a topology on $X$ by \cref{subspace_of}.
    Show that $D$ is closed in $C$ with respect to $T_C$.
    \begin{subproof}
        $X\setminus U\subseteq X$ by \cref{setminus_subseteq}.
        $X\setminus (X\setminus U) = X\inter U$ by \cref{setminus_setminus}.
        $U\in T$ by \cref{open_in}.
        $T\subseteq\pow{X}$ by \cref{topology_on}.
        Thus $U\in\pow{X}$ by \cref{subseteq}.
        Thus $U\subseteq X$ by \cref{pow_iff}.
        Thus $X\inter U = U$ by \cref{inter_absorb_supseteq_left}.
        Thus $X\setminus (X\setminus U) = U$ by \cref{setminus_setminus,inter_absorb_supseteq_left}.
        Thus $X\setminus U$ is closed in $X$ with respect to $T$ by \cref{closed_in}.
        Let $E = X\setminus U$.
        $D = E\inter C$ by \cref{setminus_eq_inter_complement,inter_comm}.
        Thus $D$ is closed in $C$ with respect to $T_C$ by \cref{closed_in_subspace_iff}.
    \end{subproof}
    $C$ is compact with respect to $T_C$ by \cref{compact_subspace}.
    $D\subseteq C$ by \cref{setminus_subseteq}.
    $T_C$ is a topology on $C$ by \cref{subspace_topology_is_topology}.
    Thus $D$ is compact with respect to $T_{CD}$ by \cref{closed_subspace_compact}.
    Thus $T_{CD} = T_D$ by \cref{subspace_subspace_topology}.
    Thus $D$ is compact with respect to $T_D$.
    Thus $D$ is a compact subspace of $X$ with respect to $T$ by \cref{compact_subspace}.
\end{proof}

% from §29 helper definition 2: open sets from X
\begin{definition}\label{one_point_open_sets}
    $\onePointOpenSets{T}{X}{p}
        = \{U\in\pow{X\union\{p\}} \mid \text{$U\subseteq X$ and $U\in T$}\}$.
\end{definition}

% from §29 helper definition 3: complements of compact sets
\begin{definition}\label{one_point_compact_complements}
    $\onePointCompactComplements{T}{X}{p}
        = \{U\in\pow{X\union\{p\}} \mid \text{there exists $C$ such that
            $C$ is a compact subspace of $X$ with respect to $T$ and
            $U = (X\union\{p\})\setminus C$}\}$.
\end{definition}

% from §29 helper definition 4: one-point compactification topology
\begin{definition}\label{one_point_compactification_topology}
    $\onePointTopology{T}{X}{p}
        = \onePointOpenSets{T}{X}{p}\union \onePointCompactComplements{T}{X}{p}$.
\end{definition}

% from §29 helper lemma 6: one-point topology is contained in powerset
\begin{lemma}\label{one_point_topology_subseteq_pow}
    Assume $p\notin X$.
    Let $Y = X\union\{p\}$.
    Let $S = \onePointTopology{T}{X}{p}$.
    Then $S\subseteq\pow{Y}$.
\end{lemma}
\begin{proof}
    Let $S_1 = \onePointOpenSets{T}{X}{p}$.
    Let $S_2 = \onePointCompactComplements{T}{X}{p}$.
    $S = S_1\union S_2$ by \cref{one_point_compactification_topology}.
    Show that for all $U$ such that $U\in S$ we have $U\in\pow{Y}$.
    \begin{subproof}
        Fix $U$.
        Assume $U\in S$.
        Thus $U\in S_1$ or $U\in S_2$ by \cref{union_iff}.
        \begin{byCase}
        \caseOf{$U\in S_1$.}
            Thus $U\in\pow{Y}$.
        \caseOf{$U\in S_2$.}
            Thus $U\in\pow{Y}$.
        \end{byCase}
    \end{subproof}
    Thus $S\subseteq\pow{Y}$ by \cref{subseteq}.
\end{proof}

% from §29 helper lemma 7: open set from X is open in one-point topology
\begin{lemma}\label{one_point_topology_intro_left}
    Assume $p\notin X$.
    Let $Y = X\union\{p\}$.
    Let $S = \onePointTopology{T}{X}{p}$.
    Suppose $U\subseteq X$.
    Suppose $U\in T$.
    Then $U\in S$.
\end{lemma}
\begin{proof}
    Let $S_1 = \onePointOpenSets{T}{X}{p}$.
    Let $S_2 = \onePointCompactComplements{T}{X}{p}$.
    $S = S_1\union S_2$ by \cref{one_point_compactification_topology}.
    Show that $U\in S_1$.
    \begin{subproof}
        $X\subseteq Y$ by \cref{union_upper_left}.
        Thus $U\subseteq Y$ by \cref{subseteq_transitive}.
        Thus $U\in\pow{Y}$ by \cref{pow_iff}.
        Thus $U\in S_1$.
    \end{subproof}
    Thus $U\in S$ by \cref{union_intro_left}.
\end{proof}

% from §29 helper lemma 8: complement of compact set is open in one-point topology
\begin{lemma}\label{one_point_topology_intro_right}
    Assume $p\notin X$.
    Let $Y = X\union\{p\}$.
    Let $S = \onePointTopology{T}{X}{p}$.
    Suppose there exists $C$ such that
        $C$ is a compact subspace of $X$ with respect to $T$ and
        $U = Y\setminus C$.
    Then $U\in S$.
\end{lemma}
\begin{proof}
    Let $S_1 = \onePointOpenSets{T}{X}{p}$.
    Let $S_2 = \onePointCompactComplements{T}{X}{p}$.
    $S = S_1\union S_2$ by \cref{one_point_compactification_topology}.
    Take $C$ such that
        $C$ is a compact subspace of $X$ with respect to $T$ and
        $U = Y\setminus C$.
    Show that $U\in S_2$.
    \begin{subproof}
        $U\subseteq Y$ by \cref{setminus_subseteq}.
        Thus $U\in\pow{Y}$ by \cref{pow_iff}.
        Thus $U\in S_2$.
    \end{subproof}
    Thus $U\in S$ by \cref{union_intro_right}.
\end{proof}

% from §29 helper lemma 9: open set without added point
\begin{lemma}\label{one_point_topology_without_p}
    Assume $p\notin X$.
    Let $Y = X\union\{p\}$.
    Let $S = \onePointTopology{T}{X}{p}$.
    Suppose $U\in S$.
    Suppose $p\notin U$.
    Then $U\subseteq X$ and $U\in T$.
\end{lemma}
\begin{proof}
    Let $S_1 = \onePointOpenSets{T}{X}{p}$.
    Let $S_2 = \onePointCompactComplements{T}{X}{p}$.
    $S = S_1\union S_2$ by \cref{one_point_compactification_topology}.
    Thus $U\in S_1$ or $U\in S_2$ by \cref{union_iff}.
    Show that $U\in S_1$.
    \begin{subproof}
        Suppose not.
        Then $U\notin S_1$.
        Thus $U\in S_2$ by \cref{union_iff}.
        Take $C$ such that
            $C$ is a compact subspace of $X$ with respect to $T$ and
            $U = Y\setminus C$
            by \cref{one_point_compact_complements}.
        $C\subseteq X$ by \cref{compact_subspace}.
        Show that $p\notin C$.
        \begin{subproof}
            Suppose not.
            Then $p\in C$.
            Thus $p\in X$ by \cref{subseteq}.
            Contradiction.
        \end{subproof}
        $p\in\{p\}$ by \cref{singleton_intro}.
        Thus $p\in Y$ by \cref{union_intro_right}.
        Thus $p\in Y\setminus C$ by \cref{setminus_intro}.
        Thus $p\in U$.
        Contradiction.
    \end{subproof}
    Thus $U\in S_1$.
    Thus $U\subseteq X$ and $U\in T$.
\end{proof}

% from §29 helper lemma 10: open set containing added point
\begin{lemma}\label{one_point_topology_with_p}
    Assume $p\notin X$.
    Let $Y = X\union\{p\}$.
    Let $S = \onePointTopology{T}{X}{p}$.
    Suppose $U\in S$.
    Suppose $p\in U$.
    Then there exists $C$ such that $C$ is a compact subspace of $X$ with respect to $T$
        and $U = (X\union\{p\})\setminus C$.
\end{lemma}
\begin{proof}
    Let $S_1 = \onePointOpenSets{T}{X}{p}$.
    Let $S_2 = \onePointCompactComplements{T}{X}{p}$.
    $S = S_1\union S_2$ by \cref{one_point_compactification_topology}.
    Thus $U\in S_1$ or $U\in S_2$ by \cref{union_iff}.
    Show that $U\in S_2$.
    \begin{subproof}
        Suppose not.
        Then $U\notin S_2$.
        Thus $U\in S_1$ by \cref{union_iff}.
        Thus $U\subseteq X$ and $U\in T$.
        Thus $p\in X$ by \cref{subseteq}.
        Contradiction.
    \end{subproof}
    Thus $U\in S_2$.
    Thus there exists $C$ such that $C$ is a compact subspace of $X$ with respect to $T$
        and $U = Y\setminus C$
        by \cref{one_point_compact_complements}.
\end{proof}

% from §29 helper lemma 11: one-point topology is a topology
\begin{lemma}\label{one_point_topology_is_topology}
    Assume $T$ is a topology on $X$.
    Assume $X$ is Hausdorff with respect to $T$.
    Assume $p\notin X$.
    Let $Y = X\union\{p\}$.
    Let $S = \onePointTopology{T}{X}{p}$.
    Then $S$ is a topology on $Y$.
\end{lemma}
\begin{proof}
    Show that $\emptyset\in S$.
    \begin{subproof}
        $\emptyset\subseteq X$ by \cref{emptyset_subseteq}.
        $\emptyset\in T$ by \cref{topology_on}.
        Thus $\emptyset\in S$ by \cref{one_point_topology_intro_left}.
    \end{subproof}
    Show that $Y\in S$.
    \begin{subproof}
        $\emptyset\subseteq X$ by \cref{emptyset_subseteq}.
        Let $T_0 = \subspaceTopology{T}{\emptyset}$.
        $T_0$ is a topology on $\emptyset$ by \cref{subspace_topology_is_topology}.
        Thus $\emptyset$ is compact with respect to $T_0$ by \cref{emptyset_compact}.
        Thus $\emptyset$ is a compact subspace of $X$ with respect to $T$ by \cref{compact_subspace}.
        $Y = Y\setminus\emptyset$ by \cref{setminus_emptyset}.
        Thus there exists $C$ such that $C$ is a compact subspace of $X$ with respect to $T$
            and $Y = Y\setminus C$.
        Thus $Y\in S$ by \cref{one_point_topology_intro_right}.
    \end{subproof}
    Show that for all $A$ such that $A\subseteq S$ we have $\unions{A}\in S$.
    \begin{subproof}
        Fix $A$.
        Assume $A\subseteq S$.
        \begin{byCase}
        \caseOf{$p\notin\unions{A}$.}
            Show that for all $U$ such that $U\in A$ we have $U\in T$.
            \begin{subproof}
                Fix $U$.
                Assume $U\in A$.
                Thus $U\in S$ by \cref{subseteq}.
                Show that $p\notin U$.
                \begin{subproof}
                    Suppose not.
                    Then $p\in U$.
                    Thus $p\in\unions{A}$ by \cref{unions_iff}.
                    Contradiction.
                \end{subproof}
                Thus $U\subseteq X$ and $U\in T$ by \cref{one_point_topology_without_p}.
            \end{subproof}
            Thus $A\subseteq T$ by \cref{subseteq}.
            $T\subseteq\pow{X}$ by \cref{topology_on}.
            Thus $A\subseteq\pow{X}$ by \cref{subseteq_transitive}.
            Thus $\unions{A}\subseteq X$ by \cref{unions_subseteq_of_powerset_is_subseteq}.
            Thus $\unions{A}\in T$ by \cref{topology_on}.
            Thus $\unions{A}\in S$ by \cref{one_point_topology_intro_left}.
        \caseOf{$p\in\unions{A}$.}
            Take $U_0\in A$ such that $p\in U_0$ by \cref{unions_iff}.
            $U_0\in S$ by \cref{subseteq}.
            Take $C_0$ such that
                $C_0$ is a compact subspace of $X$ with respect to $T$ and
                $U_0 = (X\union\{p\})\setminus C_0$
                by \cref{one_point_topology_with_p}.
            Let $A_1 = \{U\in A \mid p\notin U\}$.
            Show that $\unions{A_1}\in T$.
            \begin{subproof}
                Show that for all $U$ such that $U\in A_1$ we have $U\in T$.
                \begin{subproof}
                    Fix $U$.
                    Assume $U\in A_1$.
                    Thus $U\in A$ and $p\notin U$.
                    Thus $U\in S$ by \cref{subseteq}.
                    Thus $U\in T$ by \cref{one_point_topology_without_p}.
                \end{subproof}
                Thus $A_1\subseteq T$ by \cref{subseteq}.
                Thus $\unions{A_1}\in T$ by \cref{topology_on}.
            \end{subproof}
            Let $U = \unions{A_1}$.
            $T\subseteq\pow{X}$ by \cref{topology_on}.
            Thus $U\in\pow{X}$ by \cref{subseteq}.
            Thus $U\subseteq X$ by \cref{pow_iff}.
            Let $A_2 = \{U\in A \mid p\in U\}$.
            Let $F = \{Y\setminus U \mid U\in A_2\}$.
            Show that $F$ is inhabited.
            \begin{subproof}
                $U_0\in A_2$.
                Thus $Y\setminus U_0\in F$.
                Thus $F$ is inhabited.
            \end{subproof}
            Show that for all $C$ such that $C\in F$ we have
                $C$ is a compact subspace of $X$ with respect to $T$.
            \begin{subproof}
                Fix $C$.
                Assume $C\in F$.
                Take $U_1\in A_2$ such that $C = Y\setminus U_1$.
                Thus $U_1\in A$ and $p\in U_1$.
                Thus $U_1\in S$ by \cref{subseteq}.
                Take $C_1$ such that
                    $C_1$ is a compact subspace of $X$ with respect to $T$ and
                    $U_1 = (X\union\{p\})\setminus C_1$
                    by \cref{one_point_topology_with_p}.
                $C_1\subseteq X$ by \cref{compact_subspace}.
                Show that $X\subseteq Y$.
                \begin{subproof}
                    Show that for all $x$ such that $x\in X$ we have $x\in Y$.
                    \begin{subproof}
                        Fix $x$.
                        Assume $x\in X$.
                        Thus $x\in X\union\{p\}$ by \cref{union_intro_left}.
                        Thus $x\in Y$.
                    \end{subproof}
                    Thus $X\subseteq Y$ by \cref{subseteq}.
                \end{subproof}
                Thus $C_1\subseteq Y$ by \cref{subseteq_transitive}.
                $Y\setminus U_1 = Y\setminus (Y\setminus C_1)$.
                Thus $Y\setminus U_1 = Y\inter C_1$ by \cref{setminus_setminus}.
                Thus $Y\inter C_1 = C_1$ by \cref{inter_absorb_supseteq_left}.
                Thus $C = C_1$.
                Thus $C$ is a compact subspace of $X$ with respect to $T$.
            \end{subproof}
            Show that $F\subseteq\pow{X}$.
            \begin{subproof}
                Show that for all $C$ such that $C\in F$ we have $C\in\pow{X}$.
                \begin{subproof}
                    Fix $C$.
                    Assume $C\in F$.
                    Thus $C$ is a compact subspace of $X$ with respect to $T$.
                    Thus $C\subseteq X$ by \cref{compact_subspace}.
                    Thus $C\in\pow{X}$ by \cref{pow_iff}.
                \end{subproof}
                Thus $F\subseteq\pow{X}$ by \cref{subseteq}.
            \end{subproof}
            Let $C = \inters{F}$.
            $C$ is a compact subspace of $X$ with respect to $T$ by \cref{compact_subspace_inters}.
            Show that $\unions{A_2} = (X\union\{p\})\setminus C$.
            \begin{subproof}
                Show that $\unions{A_2}\subseteq (X\union\{p\})\setminus C$.
                \begin{subproof}
                    Show that for all $x$ such that $x\in\unions{A_2}$ we have
                        $x\in (X\union\{p\})\setminus C$.
                    \begin{subproof}
                        Fix $x$.
                        Assume $x\in\unions{A_2}$.
                        Take $U_1\in A_2$ such that $x\in U_1$ by \cref{unions_iff}.
                        Let $C_1 = Y\setminus U_1$.
                        Thus $U_1\in A$.
                        Thus $U_1\in S$ by \cref{subseteq}.
                        Thus $C_1\in F$.
                        $S\subseteq\pow{Y}$ by \cref{one_point_topology_subseteq_pow}.
                        Thus $U_1\in\pow{Y}$ by \cref{subseteq}.
                        Thus $U_1\subseteq Y$ by \cref{pow_iff}.
                        $Y\setminus C_1 = Y\setminus (Y\setminus U_1)$.
                        Thus $Y\setminus C_1 = Y\inter U_1$ by \cref{setminus_setminus}.
                        Thus $Y\inter U_1 = U_1$ by \cref{inter_absorb_supseteq_left}.
                        Thus $x\notin C_1$ by \cref{setminus_elim_right}.
                        Show that $x\notin C$.
                        \begin{subproof}
                            Suppose not.
                            Then $x\in C$.
                            Thus $x\in C_1$ by \cref{inters_iff_forall}.
                            Contradiction.
                        \end{subproof}
                        Thus $x\notin C$.
                        $x\in X\union\{p\}$ by \cref{setminus_elim_left}.
                        Thus $x\in (X\union\{p\})\setminus C$ by \cref{setminus_intro}.
                    \end{subproof}
                    Thus $\unions{A_2}\subseteq (X\union\{p\})\setminus C$ by \cref{subseteq}.
                \end{subproof}
                Show that $(X\union\{p\})\setminus C\subseteq \unions{A_2}$.
                \begin{subproof}
                    Show that for all $x$ such that $x\in (X\union\{p\})\setminus C$ we have
                        $x\in\unions{A_2}$.
                    \begin{subproof}
                        Fix $x$.
                        Assume $x\in (X\union\{p\})\setminus C$.
                        Thus $x\notin C$ by \cref{setminus_elim_right}.
                        Suppose not.
                        Then $x\notin\unions{A_2}$.
                        Show that for all $C_1$ such that $C_1\in F$ we have $x\in C_1$.
                        \begin{subproof}
                            Fix $C_1$.
                            Assume $C_1\in F$.
                            Take $U_1\in A_2$ such that $C_1 = Y\setminus U_1$.
                            Thus $U_1\in A$.
                            Thus $U_1\in S$ by \cref{subseteq}.
                            $S\subseteq\pow{Y}$ by \cref{one_point_topology_subseteq_pow}.
                            Thus $U_1\in\pow{Y}$ by \cref{subseteq}.
                            Thus $U_1\subseteq Y$ by \cref{pow_iff}.
                            $Y\setminus C_1 = Y\setminus (Y\setminus U_1)$.
                            Thus $Y\setminus C_1 = Y\inter U_1$ by \cref{setminus_setminus}.
                            Thus $Y\inter U_1 = U_1$ by \cref{inter_absorb_supseteq_left}.
                            Suppose not.
                            Then $x\notin C_1$.
                            Thus $x\in X\union\{p\}$ by \cref{setminus_elim_left}.
                            Thus $x\in Y$.
                            Thus $x\in Y\setminus C_1$ by \cref{setminus_intro}.
                            $Y\setminus C_1 = Y\setminus (Y\setminus U_1)$.
                            Thus $Y\setminus C_1 = Y\inter U_1$ by \cref{setminus_setminus}.
                            Thus $Y\inter U_1 = U_1$ by \cref{inter_absorb_supseteq_left}.
                            Thus $x\in U_1$.
                            Thus $x\in\unions{A_2}$ by \cref{unions_iff}.
                            Contradiction.
                        \end{subproof}
                        Thus $x\in C$ by \cref{inters_iff_forall}.
                        Contradiction.
                        Thus $x\in\unions{A_2}$.
                    \end{subproof}
                    Thus $(X\union\{p\})\setminus C\subseteq \unions{A_2}$ by \cref{subseteq}.
                \end{subproof}
                Thus $\unions{A_2} = (X\union\{p\})\setminus C$ by \cref{subseteq_antisymmetric}.
            \end{subproof}
            Show that $\unions{A} = \unions{A_1}\union\unions{A_2}$.
            \begin{subproof}
                Show that $\unions{A}\subseteq \unions{A_1}\union\unions{A_2}$.
                \begin{subproof}
                    Show that for all $x$ such that $x\in\unions{A}$ we have
                        $x\in\unions{A_1}\union\unions{A_2}$.
                    \begin{subproof}
                        Fix $x$.
                        Assume $x\in\unions{A}$.
                        Take $U_1\in A$ such that $x\in U_1$ by \cref{unions_iff}.
                        \begin{byCase}
                        \caseOf{$p\in U_1$.}
                            Thus $U_1\in A_2$.
                            Thus $x\in\unions{A_2}$ by \cref{unions_iff}.
                            Thus $x\in\unions{A_1}\union\unions{A_2}$ by \cref{union_intro_right}.
                        \caseOf{$p\notin U_1$.}
                            Thus $U_1\in A_1$.
                            Thus $x\in\unions{A_1}$ by \cref{unions_iff}.
                            Thus $x\in\unions{A_1}\union\unions{A_2}$ by \cref{union_intro_left}.
                        \end{byCase}
                    \end{subproof}
                    Thus $\unions{A}\subseteq \unions{A_1}\union\unions{A_2}$ by \cref{subseteq}.
                \end{subproof}
                Show that $\unions{A_1}\union\unions{A_2}\subseteq \unions{A}$.
                \begin{subproof}
                    Show that for all $x$ such that $x\in\unions{A_1}\union\unions{A_2}$ we have
                        $x\in\unions{A}$.
                    \begin{subproof}
                        Fix $x$.
                        Assume $x\in\unions{A_1}\union\unions{A_2}$.
                        Thus $x\in\unions{A_1}$ or $x\in\unions{A_2}$ by \cref{union_iff}.
                        \begin{byCase}
                        \caseOf{$x\in\unions{A_1}$.}
                            Take $U_1\in A_1$ such that $x\in U_1$ by \cref{unions_iff}.
                            Thus $U_1\in A$.
                            Thus $x\in\unions{A}$ by \cref{unions_iff}.
                        \caseOf{$x\in\unions{A_2}$.}
                            Take $U_1\in A_2$ such that $x\in U_1$ by \cref{unions_iff}.
                            Thus $U_1\in A$.
                            Thus $x\in\unions{A}$ by \cref{unions_iff}.
                        \end{byCase}
                    \end{subproof}
                    Thus $\unions{A_1}\union\unions{A_2}\subseteq \unions{A}$ by \cref{subseteq}.
                \end{subproof}
                Thus $\unions{A} = \unions{A_1}\union\unions{A_2}$ by \cref{subseteq_antisymmetric}.
            \end{subproof}
            Thus $U$ is open in $X$ with respect to $T$ by \cref{open_in}.
            Let $D = C\setminus U$.
            $D$ is a compact subspace of $X$ with respect to $T$ by \cref{compact_subspace_minus_open}.
            Show that $\unions{A} = (X\union\{p\})\setminus D$.
            \begin{subproof}
                Show that $\unions{A}\subseteq (X\union\{p\})\setminus D$.
                \begin{subproof}
                    Show that for all $x$ such that $x\in\unions{A}$ we have
                        $x\in (X\union\{p\})\setminus D$.
                    \begin{subproof}
                        Fix $x$.
                        Assume $x\in\unions{A}$.
                        Thus $x\in\unions{A_1}\union\unions{A_2}$.
                        Thus $x\in\unions{A_1}$ or $x\in\unions{A_2}$ by \cref{union_iff}.
                        \begin{byCase}
                        \caseOf{$x\in\unions{A_1}$.}
                            Thus $x\in U$.
                            Show that $x\notin D$.
                            \begin{subproof}
                                Suppose not.
                                Then $x\in D$.
                                Thus $x\in C$ and $x\notin U$ by \cref{setminus_elim_left,setminus_elim_right}.
                                Contradiction.
                            \end{subproof}
                            Thus $x\in X$ by \cref{subseteq}.
                            Thus $x\in X\union\{p\}$ by \cref{union_intro_left}.
                            Thus $x\in (X\union\{p\})\setminus D$ by \cref{setminus_intro}.
                        \caseOf{$x\in\unions{A_2}$.}
                            Thus $x\in (X\union\{p\})\setminus C$.
                            Thus $x\notin C$ by \cref{setminus_elim_right}.
                            Show that $x\notin D$.
                            \begin{subproof}
                                Suppose not.
                                Then $x\in D$.
                                Thus $x\in C$ by \cref{setminus_elim_left}.
                                Contradiction.
                            \end{subproof}
                            $x\in X\union\{p\}$ by \cref{setminus_elim_left}.
                            Thus $x\in (X\union\{p\})\setminus D$ by \cref{setminus_intro}.
                        \end{byCase}
                    \end{subproof}
                    Thus $\unions{A}\subseteq (X\union\{p\})\setminus D$ by \cref{subseteq}.
                \end{subproof}
                Show that $(X\union\{p\})\setminus D\subseteq \unions{A}$.
                \begin{subproof}
                    Show that for all $x$ such that $x\in (X\union\{p\})\setminus D$ we have
                        $x\in\unions{A}$.
                    \begin{subproof}
                        Fix $x$.
                        Assume $x\in (X\union\{p\})\setminus D$.
                        Thus $x\notin D$ by \cref{setminus_elim_right}.
                        \begin{byCase}
                        \caseOf{$x\in U$.}
                            Thus $x\in\unions{A_1}$.
                            Take $U_1\in A_1$ such that $x\in U_1$ by \cref{unions_iff}.
                            Thus $U_1\in A$.
                            Thus $x\in\unions{A}$ by \cref{unions_iff}.
                        \caseOf{$x\notin U$.}
                            Show that $x\in (X\union\{p\})\setminus C$.
                            \begin{subproof}
                                Suppose not.
                                Then $x\notin (X\union\{p\})\setminus C$.
                                Thus $x\in X\union\{p\}$ by \cref{setminus_elim_left}.
                                Show that $x\in C$.
                                \begin{subproof}
                                    Suppose not.
                                    Then $x\notin C$.
                                    Thus $x\in (X\union\{p\})\setminus C$ by \cref{setminus_intro}.
                                    Contradiction.
                                \end{subproof}
                                Thus $x\in C$ and $x\notin U$.
                                Thus $x\in D$ by \cref{setminus_intro}.
                                Contradiction.
                            \end{subproof}
                            Thus $x\in (X\union\{p\})\setminus C$.
                            Thus $x\in\unions{A_2}$.
                            Take $U_1\in A_2$ such that $x\in U_1$ by \cref{unions_iff}.
                            Thus $U_1\in A$.
                            Thus $x\in\unions{A}$ by \cref{unions_iff}.
                        \end{byCase}
                    \end{subproof}
                    Thus $(X\union\{p\})\setminus D\subseteq \unions{A}$ by \cref{subseteq}.
                \end{subproof}
                Thus $\unions{A} = (X\union\{p\})\setminus D$ by \cref{subseteq_antisymmetric}.
            \end{subproof}
            Thus there exists $C$ such that $C$ is a compact subspace of $X$ with respect to $T$
                and $\unions{A} = Y\setminus C$.
            Thus $\unions{A}\in S$ by \cref{one_point_topology_intro_right}.
        \end{byCase}
    \end{subproof}
    Show that for all $U,V$ such that $U\in S$ and $V\in S$ we have $U\inter V\in S$.
    \begin{subproof}
        Fix $U$.
        Fix $V$.
        Assume $U\in S$ and $V\in S$.
        \begin{byCase}
        \caseOf{$p\notin U$ and $p\notin V$.}
            Thus $U\subseteq X$ and $U\in T$ by \cref{one_point_topology_without_p}.
            Thus $V\subseteq X$ and $V\in T$ by \cref{one_point_topology_without_p}.
            Thus $U\inter V\subseteq U$ by \cref{inter_lower_left}.
            Thus $U\inter V\subseteq X$ by \cref{subseteq_transitive}.
            Thus $U\inter V\in T$ by \cref{topology_on}.
            Thus $U\inter V\in S$ by \cref{one_point_topology_intro_left}.
        \caseOf{$p\in U$ and $p\in V$.}
            Take $C_1$ such that
                $C_1$ is a compact subspace of $X$ with respect to $T$ and
                $U = (X\union\{p\})\setminus C_1$
                by \cref{one_point_topology_with_p}.
            Take $C_2$ such that
                $C_2$ is a compact subspace of $X$ with respect to $T$ and
                $V = (X\union\{p\})\setminus C_2$
                by \cref{one_point_topology_with_p}.
            Let $C_3 = C_1\union C_2$.
            $C_3$ is a compact subspace of $X$ with respect to $T$ by \cref{compact_subspace_union}.
            $U\inter V = (X\union\{p\})\setminus C_3$ by \cref{setminus_union}.
            Thus there exists $C$ such that $C$ is a compact subspace of $X$ with respect to $T$
                and $U\inter V = Y\setminus C$.
            Thus $U\inter V\in S$ by \cref{one_point_topology_intro_right}.
        \caseOf{$p\in U$ and $p\notin V$.}
            Take $C_1$ such that
                $C_1$ is a compact subspace of $X$ with respect to $T$ and
                $U = (X\union\{p\})\setminus C_1$
                by \cref{one_point_topology_with_p}.
            Thus $V\subseteq X$ and $V\in T$ by \cref{one_point_topology_without_p}.
            $C_1\subseteq X$ by \cref{compact_subspace}.
            Let $T_1 = \subspaceTopology{T}{C_1}$.
            Thus $C_1$ is compact with respect to $T_1$ by \cref{compact_subspace}.
            $C_1$ is closed in $X$ with respect to $T$ by \cref{compact_subspace_closed}.
            Thus $X\setminus C_1$ is open in $X$ with respect to $T$ by \cref{closed_in}.
            $X\setminus C_1\subseteq X$ by \cref{setminus_subseteq}.
            Thus $X\setminus C_1\in T$ by \cref{open_in}.
            $V\setminus C_1 = V\inter (X\setminus C_1)$ by \cref{setminus_eq_inter_complement}.
            Thus $V\setminus C_1\in T$ by \cref{topology_on}.
            $V\subseteq X\union\{p\}$ by \cref{subseteq_transitive,union_upper_left}.
            $C_1\subseteq X\union\{p\}$ by \cref{subseteq_transitive,union_upper_left}.
            Thus $V\setminus C_1 = V\inter ((X\union\{p\})\setminus C_1)$ by \cref{setminus_eq_inter_complement}.
            Thus $V\setminus C_1 = V\inter U$.
            Thus $U\inter V = V\setminus C_1$ by \cref{inter_comm}.
            Thus $U\inter V\subseteq V$ by \cref{inter_lower_left}.
            Thus $U\inter V\subseteq X$ by \cref{subseteq_transitive}.
            Thus $U\inter V\in T$.
            Thus $U\inter V\in S$ by \cref{one_point_topology_intro_left}.
        \caseOf{$p\notin U$ and $p\in V$.}
            Take $C_1$ such that
                $C_1$ is a compact subspace of $X$ with respect to $T$ and
                $V = (X\union\{p\})\setminus C_1$
                by \cref{one_point_topology_with_p}.
            Thus $U\subseteq X$ and $U\in T$ by \cref{one_point_topology_without_p}.
            $C_1\subseteq X$ by \cref{compact_subspace}.
            Let $T_1 = \subspaceTopology{T}{C_1}$.
            Thus $C_1$ is compact with respect to $T_1$ by \cref{compact_subspace}.
            $C_1$ is closed in $X$ with respect to $T$ by \cref{compact_subspace_closed}.
            Thus $X\setminus C_1$ is open in $X$ with respect to $T$ by \cref{closed_in}.
            $X\setminus C_1\subseteq X$ by \cref{setminus_subseteq}.
            Thus $X\setminus C_1\in T$ by \cref{open_in}.
            $U\setminus C_1 = U\inter (X\setminus C_1)$ by \cref{setminus_eq_inter_complement}.
            Thus $U\setminus C_1\in T$ by \cref{topology_on}.
            $U\subseteq X\union\{p\}$ by \cref{subseteq_transitive,union_upper_left}.
            $C_1\subseteq X\union\{p\}$ by \cref{subseteq_transitive,union_upper_left}.
            Thus $U\setminus C_1 = U\inter ((X\union\{p\})\setminus C_1)$ by \cref{setminus_eq_inter_complement}.
            Thus $U\setminus C_1 = U\inter V$.
            Thus $U\inter V = U\setminus C_1$.
            Thus $U\inter V\subseteq U$ by \cref{inter_lower_left}.
            Thus $U\inter V\subseteq X$ by \cref{subseteq_transitive}.
            Thus $U\inter V\in T$.
            Thus $U\inter V\in S$ by \cref{one_point_topology_intro_left}.
        \end{byCase}
    \end{subproof}
    $S\subseteq\pow{Y}$ by \cref{one_point_topology_subseteq_pow}.
    Thus $S$ is a topology on $Y$ by \cref{topology_on}.
\end{proof}

% from §29 helper lemma 12: subspace topology from one-point topology
\begin{lemma}\label{one_point_topology_subspace}
    Assume $T$ is a topology on $X$.
    Assume $X$ is Hausdorff with respect to $T$.
    Assume $p\notin X$.
    Let $Y = X\union\{p\}$.
    Let $S = \onePointTopology{T}{X}{p}$.
    Let $T_X = \subspaceTopology{S}{X}$.
    Then $T_X = T$.
\end{lemma}
\begin{proof}
    Show that $T\subseteq T_X$.
    \begin{subproof}
        Show that for all $U$ such that $U\in T$ we have $U\in T_X$.
        \begin{subproof}
            Fix $U$.
            Assume $U\in T$.
            $T\subseteq\pow{X}$ by \cref{topology_on}.
            Thus $U\subseteq X$ by \cref{pow_iff,subseteq}.
            Thus $U\in S$ by \cref{one_point_topology_intro_left}.
            $U = X\inter U$ by \cref{inter_absorb_supseteq_left}.
            Thus $U\in T_X$ by \cref{subspace_topology}.
        \end{subproof}
        Thus $T\subseteq T_X$ by \cref{subseteq}.
    \end{subproof}
    Show that $T_X\subseteq T$.
    \begin{subproof}
        Show that for all $U$ such that $U\in T_X$ we have $U\in T$.
        \begin{subproof}
            Fix $U$.
            Assume $U\in T_X$.
            Take $V\in S$ such that $U = X\inter V$ by \cref{subspace_topology}.
            \begin{byCase}
            \caseOf{$p\notin V$.}
                Thus $V\subseteq X$ and $V\in T$ by \cref{one_point_topology_without_p}.
                Thus $U = V$ by \cref{inter_absorb_supseteq_left}.
                Thus $U\in T$.
            \caseOf{$p\in V$.}
                Take $C$ such that
                    $C$ is a compact subspace of $X$ with respect to $T$ and
                    $V = (X\union\{p\})\setminus C$
                    by \cref{one_point_topology_with_p}.
                $C\subseteq X$ by \cref{compact_subspace}.
                Let $T_1 = \subspaceTopology{T}{C}$.
                Thus $C$ is compact with respect to $T_1$ by \cref{compact_subspace}.
                $C$ is closed in $X$ with respect to $T$ by \cref{compact_subspace_closed}.
                Thus $X\setminus C$ is open in $X$ with respect to $T$ by \cref{closed_in}.
                $X\setminus C\subseteq X$ by \cref{setminus_subseteq}.
                Thus $X\setminus C\in T$ by \cref{open_in}.
                Show that $X\subseteq Y$.
                \begin{subproof}
                    Show that for all $x$ such that $x\in X$ we have $x\in Y$.
                    \begin{subproof}
                        Fix $x$.
                        Assume $x\in X$.
                        Thus $x\in X\union\{p\}$ by \cref{union_intro_left}.
                        Thus $x\in Y$.
                    \end{subproof}
                    Thus $X\subseteq Y$ by \cref{subseteq}.
                \end{subproof}
                Thus $C\subseteq Y$ by \cref{subseteq_transitive}.
                $V = Y\setminus C$.
                Thus $U = X\inter (Y\setminus C)$.
                Thus $U = X\setminus C$ by \cref{setminus_eq_inter_complement}.
                Thus $U\in T$.
            \end{byCase}
        \end{subproof}
        Thus $T_X\subseteq T$ by \cref{subseteq}.
    \end{subproof}
    Thus $T_X = T$ by \cref{subseteq_antisymmetric}.
\end{proof}

% from §29 helper lemma 13: one-point topology is compact
\begin{lemma}\label{one_point_topology_compact}
    Assume $T$ is a topology on $X$.
    Assume $X$ is Hausdorff with respect to $T$.
    Assume $p\notin X$.
    Let $Y = X\union\{p\}$.
    Let $S = \onePointTopology{T}{X}{p}$.
    Then $Y$ is compact with respect to $S$.
\end{lemma}
\begin{proof}
    Show that for all $A$ such that $A$ is an open covering of $Y$ with respect to $S$
        there exists $B$ such that $B\subseteq A$ and $B$ is finite and $B$ is a covering of $Y$.
    \begin{subproof}
        Fix $A$.
        Assume $A$ is an open covering of $Y$ with respect to $S$.
        Thus $A$ is a covering of $Y$ and
            for all $U$ such that $U\in A$ we have $U$ is open in $Y$ with respect to $S$
            by \cref{open_covering}.
        $p\in\{p\}$ by \cref{singleton_intro}.
        Thus $p\in Y$ by \cref{union_intro_right}.
        Thus $Y\subseteq\unions{A}$ by \cref{covering}.
        Thus $p\in\unions{A}$ by \cref{subseteq}.
        Take $U_0\in A$ such that $p\in U_0$ by \cref{unions_iff}.
        $U_0$ is open in $Y$ with respect to $S$ by \cref{open_covering}.
        Thus $U_0\in S$ by \cref{open_in}.
        $S\subseteq\pow{Y}$ by \cref{one_point_topology_subseteq_pow}.
        Thus $U_0\in\pow{Y}$ by \cref{subseteq}.
        Thus $U_0\subseteq Y$ by \cref{pow_iff}.
        Take $C_0$ such that
            $C_0$ is a compact subspace of $X$ with respect to $T$ and
            $U_0 = (X\union\{p\})\setminus C_0$
            by \cref{one_point_topology_with_p}.
        Let $A_0 = A\setminus\{U_0\}$.
        Show that $A_0$ is a covering of $C_0$.
        \begin{subproof}
            Show that $C_0\subseteq\unions{A_0}$.
            \begin{subproof}
                Show that for all $x$ such that $x\in C_0$ we have $x\in\unions{A_0}$.
                \begin{subproof}
                    Fix $x$.
                    Assume $x\in C_0$.
                    Thus $x\notin U_0$ by \cref{setminus_elim_right}.
                    $C_0\subseteq X$ by \cref{compact_subspace}.
                    Thus $x\in X$ by \cref{subseteq}.
                    Show that $X\subseteq Y$.
                    \begin{subproof}
                        Show that for all $x$ such that $x\in X$ we have $x\in Y$.
                        \begin{subproof}
                            Fix $x$.
                            Assume $x\in X$.
                            Thus $x\in X\union\{p\}$ by \cref{union_intro_left}.
                            Thus $x\in Y$.
                    \end{subproof}
                    Thus $X\subseteq Y$ by \cref{subseteq}.
                \end{subproof}
                Thus $x\in Y$ by \cref{subseteq}.
                Thus $Y\subseteq\unions{A}$ by \cref{covering}.
                Thus $x\in\unions{A}$ by \cref{subseteq}.
                Take $U_1\in A$ such that $x\in U_1$ by \cref{unions_iff}.
                    Show that $U_1\neq U_0$.
                    \begin{subproof}
                        Suppose not.
                        Then $U_1 = U_0$.
                        Thus $x\in U_0$.
                        Contradiction.
                    \end{subproof}
                    Thus $U_1\notin\{U_0\}$ by \cref{singleton_iff}.
                    Thus $U_1\in A_0$ by \cref{setminus_intro}.
                    Thus $x\in\unions{A_0}$ by \cref{unions_iff}.
                \end{subproof}
                Thus $C_0\subseteq\unions{A_0}$ by \cref{subseteq}.
            \end{subproof}
            Thus $A_0$ is a covering of $C_0$ by \cref{covering}.
        \end{subproof}
        Let $T_X = \subspaceTopology{S}{X}$.
        $T_X = T$ by \cref{one_point_topology_subspace}.
        Let $T_C = \subspaceTopology{T}{C_0}$.
        Let $T_D = \subspaceTopology{S}{C_0}$.
        $S$ is a topology on $Y$ by \cref{one_point_topology_is_topology}.
        $C_0\subseteq X$ by \cref{compact_subspace}.
        Show that $X\subseteq Y$.
        \begin{subproof}
            Show that for all $x$ such that $x\in X$ we have $x\in Y$.
            \begin{subproof}
                Fix $x$.
                Assume $x\in X$.
                Thus $x\in X\union\{p\}$ by \cref{union_intro_left}.
                Thus $x\in Y$.
            \end{subproof}
            Thus $X\subseteq Y$ by \cref{subseteq}.
        \end{subproof}
        Thus $T_D = T_C$ by \cref{subspace_subspace_topology}.
        $C_0$ is compact with respect to $T_C$ by \cref{compact_subspace}.
        Thus $C_0$ is compact with respect to $T_D$.
        Thus $C_0\subseteq Y$ by \cref{subseteq_transitive}.
        Thus $C_0$ is a subspace of $Y$ with respect to $S$ by \cref{subspace_of}.
        Show that for all $U$ such that $U\in A_0$ we have $U$ is open in $Y$ with respect to $S$.
        \begin{subproof}
            Fix $U$.
            Assume $U\in A_0$.
            Thus $U\in A$ by \cref{setminus_elim_left}.
            Thus $U$ is open in $Y$ with respect to $S$ by \cref{open_covering}.
        \end{subproof}
        Take $B$ such that
            $B\subseteq A_0$ and
            $B$ is finite and
            $B$ is a covering of $C_0$
            by \cref{compact_subspace_open_covering}.
        Let $G = \{U_0,U_0\}$.
        $G$ is finite by \cref{pair_is_finite}.
        $U_0\in G$ by \cref{upair_intro_left}.
        Thus $\{U_0\}\subseteq G$ by \cref{singleton_subset_intro}.
        Thus $\{U_0\}$ is finite by \cref{subseteq_finite_is_finite}.
        Let $B_0 = B\union\{U_0\}$.
        $B_0$ is finite by \cref{union_of_finite_is_finite}.
        $A_0\subseteq A$ by \cref{setminus_subseteq}.
        Thus $B\subseteq A$ by \cref{subseteq_transitive}.
        Thus $\{U_0\}\subseteq A$ by \cref{singleton_subset_intro}.
        Thus $B_0\subseteq A$ by \cref{union_subsets_is_subset}.
        Show that $B_0$ is a covering of $Y$.
        \begin{subproof}
            Show that $Y\subseteq\unions{B_0}$.
            \begin{subproof}
                Show that for all $x$ such that $x\in Y$ we have $x\in\unions{B_0}$.
                \begin{subproof}
                    Fix $x$.
                    Assume $x\in Y$.
                    \begin{byCase}
                    \caseOf{$x\in U_0$.}
                        $U_0\in\{U_0\}$ by \cref{singleton_intro}.
                        Thus $U_0\in B_0$ by \cref{union_intro_right}.
                        Thus $x\in\unions{B_0}$ by \cref{unions_iff}.
                    \caseOf{$x\notin U_0$.}
                        Thus $x\in Y\setminus U_0$ by \cref{setminus_intro}.
                        $Y\setminus U_0 = Y\setminus (Y\setminus C_0)$.
                        Thus $Y\setminus U_0 = Y\inter C_0$ by \cref{setminus_setminus}.
                        $C_0\subseteq X$ by \cref{compact_subspace}.
                        Thus $C_0\subseteq Y$ by \cref{subseteq_transitive}.
                        Thus $Y\inter C_0 = C_0$ by \cref{inter_absorb_supseteq_left}.
                        Thus $x\in C_0$.
                        Thus $C_0\subseteq\unions{B}$ by \cref{covering}.
                        Thus $x\in\unions{B}$ by \cref{subseteq}.
                        Thus $x\in\unions{B_0}$ by \cref{unions_iff,union_intro_left}.
                    \end{byCase}
                \end{subproof}
                Thus $Y\subseteq\unions{B_0}$ by \cref{subseteq}.
            \end{subproof}
            Thus $B_0$ is a covering of $Y$ by \cref{covering}.
        \end{subproof}
        Thus there exists $B_0$ such that
            $B_0\subseteq A$ and
            $B_0$ is finite and
            $B_0$ is a covering of $Y$.
    \end{subproof}
    Thus $Y$ is compact with respect to $S$ by \cref{compact}.
\end{proof}

% from §29 helper lemma 14: one-point topology is Hausdorff
\begin{lemma}\label{one_point_topology_hausdorff}
    Assume $T$ is a topology on $X$.
    Assume $X$ is Hausdorff with respect to $T$.
    Assume $X$ is locally compact with respect to $T$.
    Assume $p\notin X$.
    Let $Y = X\union\{p\}$.
    Let $S = \onePointTopology{T}{X}{p}$.
    Then $Y$ is Hausdorff with respect to $S$.
\end{lemma}
\begin{proof}
    $S$ is a topology on $Y$ by \cref{one_point_topology_is_topology}.
    Show that for all $x,y$ such that $x\in Y$ and $y\in Y$ and $x\neq y$ there exist $U,V$ such that
        $U$ is a neighborhood of $x$ in $Y$ with respect to $S$ and
        $V$ is a neighborhood of $y$ in $Y$ with respect to $S$ and
        $U$ is disjoint from $V$.
    \begin{subproof}
        Fix $x$.
        Fix $y$.
        Assume $x\in Y$ and $y\in Y$ and $x\neq y$.
        \begin{byCase}
        \caseOf{$x\in X$ and $y\in X$.}
            Take $U,V$ such that
                $U$ is a neighborhood of $x$ in $X$ with respect to $T$ and
                $V$ is a neighborhood of $y$ in $X$ with respect to $T$ and
                $U$ is disjoint from $V$
                by \cref{hausdorff}.
            $U$ is open in $X$ with respect to $T$ by \cref{neighborhood}.
            $V$ is open in $X$ with respect to $T$ by \cref{neighborhood}.
            Thus $U\in T$ by \cref{open_in}.
            Thus $V\in T$ by \cref{open_in}.
            $T\subseteq\pow{X}$ by \cref{topology_on}.
            Thus $U\in\pow{X}$ by \cref{subseteq}.
            Thus $V\in\pow{X}$ by \cref{subseteq}.
            Thus $U\subseteq X$ by \cref{pow_iff}.
            Thus $V\subseteq X$ by \cref{pow_iff}.
            Thus $U\in S$ by \cref{one_point_topology_intro_left}.
            Thus $V\in S$ by \cref{one_point_topology_intro_left}.
            Thus $U$ and $V$ are open in $Y$ with respect to $S$ by \cref{open_in}.
            $x\in U$ by \cref{neighborhood}.
            $y\in V$ by \cref{neighborhood}.
            Thus $U$ is a neighborhood of $x$ in $Y$ with respect to $S$ by \cref{neighborhood}.
            Thus $V$ is a neighborhood of $y$ in $Y$ with respect to $S$ by \cref{neighborhood}.
            Thus there exist $U,V$ such that
                $U$ is a neighborhood of $x$ in $Y$ with respect to $S$ and
                $V$ is a neighborhood of $y$ in $Y$ with respect to $S$ and
                $U$ is disjoint from $V$.
        \caseOf{$x\in X$ and $y\notin X$.}
            $p\in\{p\}$ by \cref{singleton_intro}.
            Thus $p\in Y$ by \cref{union_intro_right}.
            Thus $y\in X$ or $y\in\{p\}$ by \cref{union_iff}.
            Thus $y\in\{p\}$.
            Thus $y = p$ by \cref{singleton_iff}.
            $X$ is locally compact at $x$ with respect to $T$ by \cref{locally_compact}.
            Take $C$ such that
                $C$ is a compact subspace of $X$ with respect to $T$ and
                there exists $U$ such that
                    $U$ is a neighborhood of $x$ in $X$ with respect to $T$ and
                    $U\subseteq C$
                by \cref{locally_compact_at_point}.
            Take $U$ such that
                $U$ is a neighborhood of $x$ in $X$ with respect to $T$ and
                $U\subseteq C$.
            $U$ is open in $X$ with respect to $T$ by \cref{neighborhood}.
            Thus $U\in T$ by \cref{open_in}.
            $T\subseteq\pow{X}$ by \cref{topology_on}.
            Thus $U\in\pow{X}$ by \cref{subseteq}.
            Thus $U\subseteq X$ by \cref{pow_iff}.
            Thus $U\in S$ by \cref{one_point_topology_intro_left}.
            Let $V = (X\union\{p\})\setminus C$.
            Thus $V = Y\setminus C$.
            Thus there exists $C_1$ such that
                $C_1$ is a compact subspace of $X$ with respect to $T$ and
                $V = Y\setminus C_1$.
            Thus $V\in S$ by \cref{one_point_topology_intro_right}.
            Thus $U$ and $V$ are open in $Y$ with respect to $S$ by \cref{open_in}.
            $x\in U$ by \cref{neighborhood}.
            $C\subseteq X$ by \cref{compact_subspace}.
            Show that $y\notin C$.
            \begin{subproof}
                Suppose not.
                Then $y\in C$.
                Thus $y\in X$ by \cref{subseteq}.
                Contradiction.
            \end{subproof}
            Thus $y\in Y$.
            Thus $y\in V$ by \cref{setminus_intro}.
            Show that $U$ is disjoint from $V$.
            \begin{subproof}
                Suppose not.
                Then there exists $z$ such that $z\in U$ and $z\in V$.
                Thus $z\in C$ by \cref{subseteq}.
                Thus $z\notin C$ by \cref{setminus_elim_right}.
                Contradiction.
            \end{subproof}
            Thus $U$ is a neighborhood of $x$ in $Y$ with respect to $S$ by \cref{neighborhood}.
            Thus $V$ is a neighborhood of $y$ in $Y$ with respect to $S$ by \cref{neighborhood}.
            Thus there exist $U,V$ such that
                $U$ is a neighborhood of $x$ in $Y$ with respect to $S$ and
                $V$ is a neighborhood of $y$ in $Y$ with respect to $S$ and
                $U$ is disjoint from $V$.
        \caseOf{$x\notin X$ and $y\in X$.}
            $p\in\{p\}$ by \cref{singleton_intro}.
            Thus $p\in Y$ by \cref{union_intro_right}.
            Thus $x\in X$ or $x\in\{p\}$ by \cref{union_iff}.
            Thus $x\in\{p\}$.
            Thus $x = p$ by \cref{singleton_iff}.
            $X$ is locally compact at $y$ with respect to $T$ by \cref{locally_compact}.
            Take $C$ such that
                $C$ is a compact subspace of $X$ with respect to $T$ and
                there exists $U$ such that
                    $U$ is a neighborhood of $y$ in $X$ with respect to $T$ and
                    $U\subseteq C$
                by \cref{locally_compact_at_point}.
            Take $U$ such that
                $U$ is a neighborhood of $y$ in $X$ with respect to $T$ and
                $U\subseteq C$.
            $U$ is open in $X$ with respect to $T$ by \cref{neighborhood}.
            Thus $U\in T$ by \cref{open_in}.
            $T\subseteq\pow{X}$ by \cref{topology_on}.
            Thus $U\in\pow{X}$ by \cref{subseteq}.
            Thus $U\subseteq X$ by \cref{pow_iff}.
            Thus $U\in S$ by \cref{one_point_topology_intro_left}.
            Let $V = (X\union\{p\})\setminus C$.
            Thus $V = Y\setminus C$.
            Thus there exists $C_1$ such that
                $C_1$ is a compact subspace of $X$ with respect to $T$ and
                $V = Y\setminus C_1$.
            Thus $V\in S$ by \cref{one_point_topology_intro_right}.
            Thus $U$ and $V$ are open in $Y$ with respect to $S$ by \cref{open_in}.
            $y\in U$ by \cref{neighborhood}.
            $C\subseteq X$ by \cref{compact_subspace}.
            Show that $x\notin C$.
            \begin{subproof}
                Suppose not.
                Then $x\in C$.
                Thus $x\in X$ by \cref{subseteq}.
                Contradiction.
            \end{subproof}
            Thus $x\in Y$.
            Thus $x\in V$ by \cref{setminus_intro}.
            Show that $U$ is disjoint from $V$.
            \begin{subproof}
                Suppose not.
                Then there exists $z$ such that $z\in U$ and $z\in V$.
                Thus $z\in C$ by \cref{subseteq}.
                Thus $z\notin C$ by \cref{setminus_elim_right}.
                Contradiction.
            \end{subproof}
            Thus $U$ is a neighborhood of $y$ in $Y$ with respect to $S$ by \cref{neighborhood}.
            Thus $V$ is a neighborhood of $x$ in $Y$ with respect to $S$ by \cref{neighborhood}.
            Thus $V$ is disjoint from $U$ by \cref{disjoint_symmetric}.
            Let $U_1 = V$.
            Let $V_1 = U$.
            Thus $U_1$ is a neighborhood of $x$ in $Y$ with respect to $S$ and
                $V_1$ is a neighborhood of $y$ in $Y$ with respect to $S$ and
                $U_1$ is disjoint from $V_1$.
            Thus there exist $U,V$ such that
                $U$ is a neighborhood of $x$ in $Y$ with respect to $S$ and
                $V$ is a neighborhood of $y$ in $Y$ with respect to $S$ and
                $U$ is disjoint from $V$.
        \end{byCase}
    \end{subproof}
    Thus $Y$ is Hausdorff with respect to $S$ by \cref{hausdorff}.
\end{proof}

% from §29 Item 4: one-point compactification existence (forward)
\begin{theorem}\label{one_point_compactification_existence_forward}
    Assume $T$ is a topology on $X$.
    If $X$ is locally compact with respect to $T$ and $X$ is Hausdorff with respect to $T$, then
        there exist $Y,S,p$ such that
            $S$ is a topology on $Y$ and
            $X\subseteq Y$ and
            $T = \subspaceTopology{S}{X}$ and
            $Y\setminus X = \{p\}$ and
            $Y$ is compact with respect to $S$ and
            $Y$ is Hausdorff with respect to $S$.
\end{theorem}
\begin{proof}
    Assume $X$ is locally compact with respect to $T$ and $X$ is Hausdorff with respect to $T$.
    Let $p = X$.
    $p\notin X$ by \cref{set_notin_self}.
    Let $Y = X\union\{p\}$.
    Let $S = \onePointTopology{T}{X}{p}$.
    $S$ is a topology on $Y$ by \cref{one_point_topology_is_topology}.
    Show that $X\subseteq Y$.
    \begin{subproof}
        Show that for all $x$ such that $x\in X$ we have $x\in Y$.
        \begin{subproof}
            Fix $x$.
            Assume $x\in X$.
            Thus $x\in X\union\{p\}$ by \cref{union_intro_left}.
            Thus $x\in Y$.
        \end{subproof}
        Thus $X\subseteq Y$ by \cref{subseteq}.
    \end{subproof}
    $T = \subspaceTopology{S}{X}$ by \cref{one_point_topology_subspace}.
    Show that $Y\setminus X = \{p\}$.
    \begin{subproof}
        Show that $Y\setminus X\subseteq \{p\}$.
        \begin{subproof}
            Show that for all $x$ such that $x\in Y\setminus X$ we have $x\in\{p\}$.
            \begin{subproof}
                Fix $x$.
                Assume $x\in Y\setminus X$.
                Thus $x\in Y$ and $x\notin X$ by \cref{setminus_elim_left,setminus_elim_right}.
                Thus $x\in X$ or $x\in\{p\}$ by \cref{union_iff}.
                Thus $x\in\{p\}$.
            \end{subproof}
            Thus $Y\setminus X\subseteq \{p\}$ by \cref{subseteq}.
        \end{subproof}
        Show that $\{p\}\subseteq Y\setminus X$.
        \begin{subproof}
            Show that for all $x$ such that $x\in\{p\}$ we have $x\in Y\setminus X$.
            \begin{subproof}
                Fix $x$.
                Assume $x\in\{p\}$.
                Thus $x = p$ by \cref{singleton_iff}.
                Thus $x\notin X$.
                Thus $x\in Y$ by \cref{union_intro_right}.
                Thus $x\in Y\setminus X$ by \cref{setminus_intro}.
            \end{subproof}
            Thus $\{p\}\subseteq Y\setminus X$ by \cref{subseteq}.
        \end{subproof}
        Thus $Y\setminus X = \{p\}$ by \cref{subseteq_antisymmetric}.
    \end{subproof}
    $Y$ is compact with respect to $S$ by \cref{one_point_topology_compact}.
    $Y$ is Hausdorff with respect to $S$ by \cref{one_point_topology_hausdorff}.
    Thus there exist $Y,S,p$ such that
        $S$ is a topology on $Y$ and
        $X\subseteq Y$ and
        $T = \subspaceTopology{S}{X}$ and
        $Y\setminus X = \{p\}$ and
        $Y$ is compact with respect to $S$ and
        $Y$ is Hausdorff with respect to $S$.
\end{proof}

% from §29 Item 4: one-point compactification existence (reverse)
\begin{theorem}\label{one_point_compactification_existence_reverse}
    Assume $T$ is a topology on $X$.
    If there exist $Y,S,p$ such that
            $S$ is a topology on $Y$ and
            $X\subseteq Y$ and
            $T = \subspaceTopology{S}{X}$ and
            $Y\setminus X = \{p\}$ and
            $Y$ is compact with respect to $S$ and
            $Y$ is Hausdorff with respect to $S$
        then $X$ is locally compact with respect to $T$ and $X$ is Hausdorff with respect to $T$.
\end{theorem}
\begin{proof}
    Assume there exist $Y,S,p$ such that
        $S$ is a topology on $Y$ and
        $X\subseteq Y$ and
        $T = \subspaceTopology{S}{X}$ and
        $Y\setminus X = \{p\}$ and
        $Y$ is compact with respect to $S$ and
        $Y$ is Hausdorff with respect to $S$.
    Take $Y,S,p$ such that
        $S$ is a topology on $Y$ and
        $X\subseteq Y$ and
        $T = \subspaceTopology{S}{X}$ and
        $Y\setminus X = \{p\}$ and
        $Y$ is compact with respect to $S$ and
        $Y$ is Hausdorff with respect to $S$.
    Thus $Y = X\union (Y\setminus X)$ by \cref{setminus_partition}.
    Thus $Y = X\union\{p\}$.
    $X$ is a subspace of $Y$ with respect to $S$ by \cref{subspace_of}.
    Thus $X$ is Hausdorff with respect to $T$ by \cref{subspace_hausdorff}.
    Show that $X$ is locally compact with respect to $T$.
    \begin{subproof}
        Show that for all $x\in X$ we have $X$ is locally compact at $x$ with respect to $T$.
        \begin{subproof}
            Fix $x$.
            Assume $x\in X$.
            $p\in\{p\}$ by \cref{singleton_intro}.
            Thus $p\in Y\setminus X$.
            Thus $p\in Y$ by \cref{setminus_elim_left}.
            $x\in Y$ by \cref{subseteq}.
            Show that $x\neq p$.
            \begin{subproof}
                Suppose not.
                Then $x = p$.
                Thus $p\in X$.
                Thus $p\in Y\setminus X$.
                Thus $p\notin X$ by \cref{setminus_elim_right}.
                Contradiction.
            \end{subproof}
            Take $U,V$ such that
                $U$ is a neighborhood of $x$ in $Y$ with respect to $S$ and
                $V$ is a neighborhood of $p$ in $Y$ with respect to $S$ and
                $U$ is disjoint from $V$
                by \cref{hausdorff}.
            Let $C = Y\setminus V$.
            Show that $C$ is closed in $Y$ with respect to $S$.
            \begin{subproof}
                $C\subseteq Y$ by \cref{setminus_subseteq}.
                $V$ is open in $Y$ with respect to $S$ by \cref{neighborhood}.
                Thus $V\in S$ by \cref{open_in}.
                $S\subseteq\pow{Y}$ by \cref{topology_on}.
                Thus $V\in\pow{Y}$ by \cref{subseteq}.
                Thus $V\subseteq Y$ by \cref{pow_iff}.
                $Y\setminus C = Y\setminus (Y\setminus V)$.
                Thus $Y\setminus C = Y\inter V$ by \cref{setminus_setminus}.
                Thus $Y\inter V = V$ by \cref{inter_absorb_supseteq_left}.
                Thus $Y\setminus C = V$.
                Thus $Y\setminus C$ is open in $Y$ with respect to $S$ by \cref{neighborhood}.
                Thus $C$ is closed in $Y$ with respect to $S$ by \cref{closed_in}.
            \end{subproof}
            $C\subseteq Y$ by \cref{setminus_subseteq}.
            Thus $C$ is compact with respect to $\subspaceTopology{S}{C}$ by \cref{closed_subspace_compact}.
            $p\in V$ by \cref{neighborhood}.
            Thus $p\notin C$ by \cref{setminus_elim_right}.
            Show that $C\subseteq X$.
            \begin{subproof}
                Show that for all $y$ such that $y\in C$ we have $y\in X$.
                \begin{subproof}
                    Fix $y$.
                    Assume $y\in C$.
                    Thus $y\in Y$ by \cref{setminus_elim_left}.
                    Thus $y\in X$ or $y\in\{p\}$ by \cref{union_iff}.
                    Suppose not.
                    Then $y\notin X$.
                    Thus $y\in\{p\}$.
                    Then $y = p$ by \cref{singleton_iff}.
                    Thus $y\notin C$.
                    Contradiction.
                    Thus $y\in X$.
                \end{subproof}
                Thus $C\subseteq X$ by \cref{subseteq}.
            \end{subproof}
            Let $T_C = \subspaceTopology{T}{C}$.
            Thus $\subspaceTopology{S}{C} = T_C$ by \cref{subspace_subspace_topology}.
            Thus $C$ is compact with respect to $T_C$.
            $C$ is a compact subspace of $X$ with respect to $T$ by \cref{compact_subspace}.
            $U$ is open in $Y$ with respect to $S$ by \cref{neighborhood}.
            Thus $U\in S$ by \cref{open_in}.
            $S\subseteq\pow{Y}$ by \cref{topology_on}.
            Thus $U\in\pow{Y}$ by \cref{subseteq}.
            Thus $U\subseteq Y$ by \cref{pow_iff}.
            $U\inter X = X\inter U$ by \cref{inter_comm}.
            Thus $U\inter X\in \subspaceTopology{S}{X}$ by \cref{subspace_topology}.
            Thus $U\inter X\in T$.
            Thus $U\inter X$ is open in $X$ with respect to $T$ by \cref{open_in}.
            $x\in U$ by \cref{neighborhood}.
            Thus $x\in U\inter X$ by \cref{inter_intro}.
            Show that $U\inter X\subseteq C$.
            \begin{subproof}
                Show that for all $y$ such that $y\in U\inter X$ we have $y\in C$.
                \begin{subproof}
                    Fix $y$.
                    Assume $y\in U\inter X$.
                    Thus $y\in U$.
                    Suppose not.
                    Then $y\notin C$.
                    Then $y\in Y$ by \cref{subseteq}.
                    Thus $y\in Y\setminus C$ by \cref{setminus_intro}.
                    $Y\setminus C = Y\setminus (Y\setminus V)$.
                    Thus $Y\setminus C = Y\inter V$ by \cref{setminus_setminus}.
                    $V$ is open in $Y$ with respect to $S$ by \cref{neighborhood}.
                    Thus $V\in S$ by \cref{open_in}.
                    Thus $V\in\pow{Y}$ by \cref{subseteq}.
                    Thus $V\subseteq Y$ by \cref{pow_iff}.
                    Thus $Y\inter V = V$ by \cref{inter_absorb_supseteq_left}.
                    Thus $y\in V$.
                    Thus $y\in U\inter V$ by \cref{inter_intro}.
                    Contradiction by \cref{disjoint}.
                    Thus $y\in C$.
                \end{subproof}
                Thus $U\inter X\subseteq C$ by \cref{subseteq}.
            \end{subproof}
            $U\inter X$ is a neighborhood of $x$ in $X$ with respect to $T$ by \cref{neighborhood}.
            Thus $X$ is locally compact at $x$ with respect to $T$ by \cref{locally_compact_at_point}.
        \end{subproof}
        Thus $X$ is locally compact with respect to $T$ by \cref{locally_compact}.
    \end{subproof}
    Thus $X$ is locally compact with respect to $T$ and $X$ is Hausdorff with respect to $T$.
\end{proof}

 
% from §29 Item 5: one-point compactification uniqueness
\begin{theorem}\label{one_point_compactification_uniqueness}
    Assume $T$ is a topology on $X$.
    Assume $S$ is a topology on $Y$.
    Assume $S_1$ is a topology on $Y_1$.
    Suppose $X$ is a subspace of $Y$ with respect to $S$.
    Suppose $X$ is a subspace of $Y_1$ with respect to $S_1$.
    Suppose $T = \subspaceTopology{S}{X}$.
    Suppose $T = \subspaceTopology{S_1}{X}$.
    Suppose $Y\setminus X = \{p\}$.
    Suppose $Y_1\setminus X = \{q\}$.
    Suppose $Y$ is compact with respect to $S$ and $Y$ is Hausdorff with respect to $S$.
    Suppose $Y_1$ is compact with respect to $S_1$ and $Y_1$ is Hausdorff with respect to $S_1$.
    Then there exists $h$ such that
        $h$ is a homeomorphism between $Y$ and $Y_1$ with respect to $S$ and $S_1$ and
        $\restrl{h}{X} = \identity{X}$.
\end{theorem}
\begin{proof}
    Thus $X\subseteq Y$ by \cref{subspace_of}.
    Thus $X\subseteq Y_1$ by \cref{subspace_of}.
    Thus $Y = X\union\{p\}$ by \cref{setminus_partition}.
    Thus $Y_1 = X\union\{q\}$ by \cref{setminus_partition}.
    $p\in\{p\}$ by \cref{singleton_intro}.
    Thus $p\in Y\setminus X$.
    Thus $p\notin X$ by \cref{setminus_elim_right}.
    Thus $p\in Y$ by \cref{setminus_elim_left}.
    $q\in\{q\}$ by \cref{singleton_intro}.
    Thus $q\in Y_1\setminus X$.
    Thus $q\notin X$ by \cref{setminus_elim_right}.
    Thus $q\in Y_1$ by \cref{setminus_elim_left}.

    Let $f = \identity{X}$.
    Let $g = \{(p,q)\}$.
    Let $h = f\union g$.

    $f$ is a function by \cref{id_is_function}.
    Show that $g$ is a function.
    \begin{subproof}
        Show that $g$ is a relation.
        \begin{subproof}
            Show that for all $w$ such that $w\in g$ there exist $x,y$ such that $w = (x,y)$.
            \begin{subproof}
                Fix $w$.
                Assume $w\in g$.
                Thus $w = (p,q)$ by \cref{singleton_iff}.
                Thus there exist $x,y$ such that $w = (x,y)$.
            \end{subproof}
            Thus $g$ is a relation by \cref{relation}.
        \end{subproof}
        Show that $g$ is right-unique.
        \begin{subproof}
            Show that for all $x,y_1,y_2$ such that $(x,y_1)\in g$ and $(x,y_2)\in g$ we have $y_1 = y_2$.
            \begin{subproof}
                Fix $x$.
                Fix $y_1$.
                Fix $y_2$.
                Assume $(x,y_1)\in g$ and $(x,y_2)\in g$.
                Thus $(x,y_1) = (p,q)$ by \cref{singleton_iff}.
                Thus $(x,y_2) = (p,q)$ by \cref{singleton_iff}.
                Thus $y_1 = q$ by \cref{pair_eq_iff}.
                Thus $y_2 = q$ by \cref{pair_eq_iff}.
                Thus $y_1 = y_2$.
            \end{subproof}
            Thus $g$ is right-unique by \cref{rightunique}.
        \end{subproof}
        Thus $g$ is a function.
    \end{subproof}

    $\dom{f} = X$ by \cref{id_dom}.
    $g = \cons{(p,q)}{\emptyset}$.
    Thus $\dom{g} = \cons{p}{\dom{\emptyset}}$ by \cref{dom_cons}.
    Thus $\dom{g} = \cons{p}{\emptyset}$ by \cref{dom_emptyset}.
    Thus $\dom{g} = \{p\}$.
    Show that for all $x$ such that $x\in\dom{f}\inter\dom{g}$ we have $f(x) = g(x)$.
    \begin{subproof}
        Fix $x$.
        Assume $x\in\dom{f}\inter\dom{g}$.
        Suppose not.
        Then $f(x)\neq g(x)$.
        Thus $x\in\dom{f}$ and $x\in\dom{g}$ by \cref{inter}.
        Thus $x\in X$.
        Thus $x\in\{p\}$.
        Thus $x = p$ by \cref{singleton_iff}.
        Thus $x\notin X$.
        Contradiction.
        Thus $f(x) = g(x)$.
    \end{subproof}
    Thus $h$ is a function by \cref{union_compatible_functions}.

    $\dom{h} = \dom{f}\union\dom{g}$ by \cref{dom_union}.
    Thus $\dom{h} = X\union\{p\}$.
    Thus $\dom{h} = Y$.

    Show that for all $x$ such that $x\in\dom{h}$ we have $h(x)\in Y_1$.
    \begin{subproof}
        Fix $x$.
        Assume $x\in\dom{h}$.
        Thus $x\in Y$.
        Thus $x\in X$ or $x\in\{p\}$ by \cref{union_iff}.
        \begin{byCase}
        \caseOf{$x\in X$.}
            $(x,x)\in f$ by \cref{id_elem_intro}.
            Thus $(x,x)\in h$ by \cref{union_intro_left}.
            Thus $h(x) = x$ by \cref{function_apply_intro}.
            Thus $x\in Y_1$ by \cref{union_intro_left}.
            Thus $h(x)\in Y_1$.
        \caseOf{$x\in\{p\}$.}
            Thus $x = p$ by \cref{singleton_iff}.
            $(p,q)\in g$ by \cref{singleton_intro}.
            Thus $(p,q)\in h$ by \cref{union_intro_right}.
            Thus $h(p) = q$ by \cref{function_apply_intro}.
            Thus $q\in Y_1$ by \cref{union_intro_right}.
            Thus $h(x)\in Y_1$.
        \end{byCase}
    \end{subproof}
    Thus $h$ is a function to $Y_1$.
    Thus $h\in\funs{Y}{Y_1}$ by \cref{funs_intro}.

    Show that for all $y$ such that $y\in Y_1$ there exists $x$ such that $x\in Y$ and $h(x) = y$.
    \begin{subproof}
        Fix $y$.
        Assume $y\in Y_1$.
        Thus $y\in X$ or $y\in\{q\}$ by \cref{union_iff}.
        \begin{byCase}
        \caseOf{$y\in X$.}
            Let $x = y$.
            $x\in Y$ by \cref{union_intro_left}.
            $(x,x)\in f$ by \cref{id_elem_intro}.
            Thus $(x,x)\in h$ by \cref{union_intro_left}.
            Thus $h(x) = y$ by \cref{function_apply_intro}.
            Thus there exists $x$ such that $x\in Y$ and $h(x) = y$.
        \caseOf{$y\in\{q\}$.}
            Thus $y = q$ by \cref{singleton_iff}.
            Let $x = p$.
            $x\in Y$ by \cref{union_intro_right}.
            $(p,q)\in g$ by \cref{singleton_intro}.
            Thus $(p,q)\in h$ by \cref{union_intro_right}.
            Thus $h(x) = y$ by \cref{function_apply_intro}.
            Thus there exists $x$ such that $x\in Y$ and $h(x) = y$.
        \end{byCase}
    \end{subproof}
    Thus $h\in\surj{Y}{Y_1}$ by \cref{surj}.

    Show that for all $x_1,x_2$ such that $x_1\in\dom{h}$ and $x_2\in\dom{h}$ and $h(x_1) = h(x_2)$
        we have $x_1 = x_2$.
    \begin{subproof}
        Fix $x_1$.
        Fix $x_2$.
        Assume $x_1\in\dom{h}$ and $x_2\in\dom{h}$ and $h(x_1) = h(x_2)$.
        Thus $x_1\in Y$ and $x_2\in Y$.
        Thus $x_1\in X$ or $x_1\in\{p\}$ by \cref{union_iff}.
        Thus $x_2\in X$ or $x_2\in\{p\}$ by \cref{union_iff}.
        \begin{byCase}
        \caseOf{$x_1\in X$ and $x_2\in X$.}
            $(x_1,x_1)\in f$ by \cref{id_elem_intro}.
            Thus $(x_1,x_1)\in h$ by \cref{union_intro_left}.
            $(x_2,x_2)\in f$ by \cref{id_elem_intro}.
            Thus $(x_2,x_2)\in h$ by \cref{union_intro_left}.
            Thus $h(x_1) = x_1$ by \cref{function_apply_intro}.
            Thus $h(x_2) = x_2$ by \cref{function_apply_intro}.
            Thus $x_1 = x_2$.
        \caseOf{$x_1\in X$ and $x_2\in\{p\}$.}
            Suppose not.
            Then $x_1\neq x_2$.
            Thus $x_2 = p$ by \cref{singleton_iff}.
            $(x_1,x_1)\in f$ by \cref{id_elem_intro}.
            Thus $(x_1,x_1)\in h$ by \cref{union_intro_left}.
            Thus $h(x_1) = x_1$ by \cref{function_apply_intro}.
            $(p,q)\in g$ by \cref{singleton_intro}.
            Thus $(p,q)\in h$ by \cref{union_intro_right}.
            Thus $h(x_2) = q$ by \cref{function_apply_intro}.
            Thus $x_1 = q$.
            Thus $x_1\notin X$.
            Contradiction.
            Thus $x_1 = x_2$.
        \caseOf{$x_1\in\{p\}$ and $x_2\in X$.}
            Suppose not.
            Then $x_1\neq x_2$.
            Thus $x_1 = p$ by \cref{singleton_iff}.
            $(p,q)\in g$ by \cref{singleton_intro}.
            Thus $(p,q)\in h$ by \cref{union_intro_right}.
            Thus $h(x_1) = q$ by \cref{function_apply_intro}.
            $(x_2,x_2)\in f$ by \cref{id_elem_intro}.
            Thus $(x_2,x_2)\in h$ by \cref{union_intro_left}.
            Thus $h(x_2) = x_2$ by \cref{function_apply_intro}.
            Thus $x_2 = q$.
            Thus $x_2\notin X$.
            Contradiction.
            Thus $x_1 = x_2$.
        \caseOf{$x_1\in\{p\}$ and $x_2\in\{p\}$.}
            Thus $x_1 = p$ by \cref{singleton_iff}.
            Thus $x_2 = p$ by \cref{singleton_iff}.
            Thus $x_1 = x_2$.
        \end{byCase}
    \end{subproof}
    Thus $h$ is injective by \cref{injective_function}.
    Thus $h\in\bijections{Y}{Y_1}$ by \cref{bijections}.
    Thus $h$ is a bijection from $Y$ to $Y_1$.

    Show that for all $U$ such that $U$ is open in $Y_1$ with respect to $S_1$
        we have $\preimg{h}{U}$ is open in $Y$ with respect to $S$.
    \begin{subproof}
        Fix $U$.
        Assume $U$ is open in $Y_1$ with respect to $S_1$.
        \begin{byCase}
        \caseOf{$q\notin U$.}
            Show that $U\subseteq X$.
            \begin{subproof}
                Show that for all $y$ such that $y\in U$ we have $y\in X$.
                \begin{subproof}
                    Fix $y$.
                    Assume $y\in U$.
                    Thus $y\in Y_1$.
                    Thus $y\in X$ or $y\in\{q\}$ by \cref{union_iff}.
                    Suppose not.
                    Then $y\in\{q\}$.
                    Thus $y = q$ by \cref{singleton_iff}.
                    Thus $q\in U$.
                    Contradiction.
                    Thus $y\in X$.
                \end{subproof}
                Thus $U\subseteq X$ by \cref{subseteq}.
            \end{subproof}
            Show that $\preimg{h}{U} = U$.
            \begin{subproof}
                Show that $\preimg{h}{U}\subseteq U$.
                \begin{subproof}
                    Show that for all $y$ such that $y\in\preimg{h}{U}$ we have $y\in U$.
                    \begin{subproof}
                    Fix $y$.
                    Assume $y\in\preimg{h}{U}$.
                    Take $b\in U$ such that $y\mathrel{h} b$ by \cref{preimg_iff}.
                    Show that $y\notin\{p\}$.
                    \begin{subproof}
                        Suppose not.
                        Then $y\in\{p\}$.
                        Thus $y = p$ by \cref{singleton_iff}.
                        Thus $(y,b)\in h$.
                        Show that $(y,b)\notin f$.
                        \begin{subproof}
                            Suppose not.
                            Then $(y,b)\in f$.
                            Take $a\in X$ such that $(y,b) = (a,a)$ by \cref{id_elem_inspect}.
                            Thus $y = a$ by \cref{pair_eq_iff}.
                            Thus $y\in X$.
                            Thus $y = p$ by \cref{singleton_iff}.
                            Thus $p\in X$.
                            Contradiction.
                            Thus $(y,b)\notin f$.
                        \end{subproof}
                        Thus $(y,b)\in g$ by \cref{union_iff}.
                        Thus $b = q$ by \cref{singleton_iff,pair_eq_iff}.
                        Thus $q\in U$.
                        Contradiction.
                        Thus $y\notin\{p\}$.
                    \end{subproof}
                    Thus $\preimg{h}{U}\subseteq\dom{h}$ by \cref{preimg_subseteq_dom}.
                    Thus $y\in\dom{h}$ by \cref{subseteq}.
                    Thus $y\in Y$.
                    Thus $y\in X$ or $y\in\{p\}$ by \cref{union_iff}.
                    Thus $y\in X$.
                    $(y,y)\in f$ by \cref{id_elem_intro}.
                    Thus $(y,y)\in h$ by \cref{union_intro_left}.
                    Thus $h(y) = y$ by \cref{function_apply_intro}.
                        Thus $y = b$ by \cref{function_apply_intro}.
                        Thus $y\in U$.
                    \end{subproof}
                    Thus $\preimg{h}{U}\subseteq U$ by \cref{subseteq}.
                \end{subproof}
                Show that $U\subseteq\preimg{h}{U}$.
                \begin{subproof}
                    Show that for all $y$ such that $y\in U$ we have $y\in\preimg{h}{U}$.
                    \begin{subproof}
                        Fix $y$.
                        Assume $y\in U$.
                        Thus $y\in X$ by \cref{subseteq}.
                        $(y,y)\in f$ by \cref{id_elem_intro}.
                        Thus $(y,y)\in h$ by \cref{union_intro_left}.
                        Thus $h(y) = y$ by \cref{function_apply_intro}.
                        Thus $y\mathrel{h} y$.
                        Thus $y\in\preimg{h}{U}$ by \cref{preimg_iff}.
                    \end{subproof}
                    Thus $U\subseteq\preimg{h}{U}$ by \cref{subseteq}.
                \end{subproof}
                Thus $\preimg{h}{U} = U$ by \cref{subseteq_antisymmetric}.
            \end{subproof}
            $U\in S_1$ by \cref{open_in}.
            $U = X\inter U$ by \cref{inter_absorb_supseteq_left}.
            Thus $U\in \subspaceTopology{S_1}{X}$ by \cref{subspace_topology}.
            Thus $U\in T$.
            Thus $U$ is open in $X$ with respect to $T$ by \cref{open_in}.
            $\{p\}\subseteq Y$ by \cref{singleton_subset_intro}.
            Let $G = \{p,p\}$.
            $G$ is finite by \cref{pair_is_finite}.
            $p\in G$ by \cref{upair_intro_left}.
            Thus $\{p\}\subseteq G$ by \cref{singleton_subset_intro}.
            Thus $\{p\}$ is finite by \cref{subseteq_finite_is_finite}.
            Thus $\{p\}$ is closed in $Y$ with respect to $S$ by \cref{finite_point_set_closed_hausdorff}.
            Thus $Y\setminus\{p\}$ is open in $Y$ with respect to $S$ by \cref{closed_in}.
            Show that $X = Y\setminus\{p\}$.
            \begin{subproof}
                Show that $X\subseteq Y\setminus\{p\}$.
                \begin{subproof}
                    Show that for all $y$ such that $y\in X$ we have $y\in Y\setminus\{p\}$.
                    \begin{subproof}
                        Fix $y$.
                        Assume $y\in X$.
                        Thus $y\in Y$ by \cref{union_intro_left}.
                        Thus $y\neq p$.
                        Thus $y\notin\{p\}$ by \cref{singleton_iff}.
                        Thus $y\in Y\setminus\{p\}$ by \cref{setminus_intro}.
                    \end{subproof}
                    Thus $X\subseteq Y\setminus\{p\}$ by \cref{subseteq}.
                \end{subproof}
                Show that $Y\setminus\{p\}\subseteq X$.
                \begin{subproof}
                    Show that for all $y$ such that $y\in Y\setminus\{p\}$ we have $y\in X$.
                    \begin{subproof}
                        Fix $y$.
                        Assume $y\in Y\setminus\{p\}$.
                        Thus $y\in Y$ by \cref{setminus_elim_left}.
                        Thus $y\in X$ or $y\in\{p\}$ by \cref{union_iff}.
                        Suppose not.
                        Then $y\in\{p\}$.
                        Thus $y = p$ by \cref{singleton_iff}.
                        Thus $y\notin Y\setminus\{p\}$ by \cref{setminus_elim_right}.
                        Contradiction.
                        Thus $y\in X$.
                    \end{subproof}
                    Thus $Y\setminus\{p\}\subseteq X$ by \cref{subseteq}.
                \end{subproof}
                Thus $X = Y\setminus\{p\}$ by \cref{subseteq_antisymmetric}.
            \end{subproof}
            Thus $X$ is open in $Y$ with respect to $S$.
            Thus $U$ is open in $Y$ with respect to $S$ by \cref{open_in_subspace_open_in_space,subspace_of}.
            Thus $\preimg{h}{U}$ is open in $Y$ with respect to $S$.
        \caseOf{$q\in U$.}
            Let $C = Y_1\setminus U$.
            Show that $C$ is closed in $Y_1$ with respect to $S_1$.
            \begin{subproof}
                $C\subseteq Y_1$ by \cref{setminus_subseteq}.
                $U\in S_1$ by \cref{open_in}.
                $S_1\subseteq\pow{Y_1}$ by \cref{topology_on}.
                Thus $U\in\pow{Y_1}$ by \cref{subseteq}.
                Thus $U\subseteq Y_1$ by \cref{pow_iff}.
                $Y_1\setminus C = Y_1\setminus (Y_1\setminus U)$.
                Thus $Y_1\setminus C = Y_1\inter U$ by \cref{setminus_setminus}.
                Thus $Y_1\inter U = U$ by \cref{inter_absorb_supseteq_left}.
                Thus $Y_1\setminus C = U$.
                Thus $Y_1\setminus C$ is open in $Y_1$ with respect to $S_1$ by \cref{open_in}.
                Thus $C$ is closed in $Y_1$ with respect to $S_1$ by \cref{closed_in}.
            \end{subproof}
            Show that $C\subseteq X$.
            \begin{subproof}
                Show that for all $y$ such that $y\in C$ we have $y\in X$.
                \begin{subproof}
                    Fix $y$.
                    Assume $y\in C$.
                    Thus $y\in Y_1$ by \cref{setminus_elim_left}.
                    Thus $y\in X$ or $y\in\{q\}$ by \cref{union_iff}.
                    Suppose not.
                    Then $y\in\{q\}$.
                    Thus $y = q$ by \cref{singleton_iff}.
                    Thus $y\notin C$ by \cref{setminus_elim_right}.
                    Contradiction.
                    Thus $y\in X$.
                \end{subproof}
                Thus $C\subseteq X$ by \cref{subseteq}.
            \end{subproof}
            $C\subseteq Y_1$ by \cref{setminus_subseteq}.
            $Y_1$ is compact with respect to $S_1$.
            Thus $C$ is compact with respect to $\subspaceTopology{S_1}{C}$ by \cref{closed_subspace_compact}.
            Let $T_C = \subspaceTopology{T}{C}$.
            Thus $\subspaceTopology{S_1}{C} = T_C$ by \cref{subspace_subspace_topology}.
            Thus $C$ is compact with respect to $T_C$.
            Thus $C$ is a compact subspace of $X$ with respect to $T$ by \cref{compact_subspace}.
            Thus $\subspaceTopology{S}{C} = T_C$ by \cref{subspace_subspace_topology}.
            Thus $C$ is compact with respect to $\subspaceTopology{S}{C}$.
            $C\subseteq X$.
            Thus $C\subseteq Y$ by \cref{subseteq_transitive}.
            Thus $C$ is a compact subspace of $Y$ with respect to $S$ by \cref{compact_subspace}.
            Thus $C$ is closed in $Y$ with respect to $S$ by \cref{compact_subspace_closed}.
            Thus $Y\setminus C$ is open in $Y$ with respect to $S$ by \cref{closed_in}.
            Show that $\preimg{h}{U} = Y\setminus C$.
            \begin{subproof}
                Show that $\preimg{h}{U}\subseteq Y\setminus C$.
                \begin{subproof}
                    Show that for all $y$ such that $y\in\preimg{h}{U}$ we have $y\in Y\setminus C$.
                    \begin{subproof}
                    Fix $y$.
                    Assume $y\in\preimg{h}{U}$.
                    Thus $\preimg{h}{U}\subseteq\dom{h}$ by \cref{preimg_subseteq_dom}.
                    Thus $y\in\dom{h}$ by \cref{subseteq}.
                    Thus $y\in Y$.
                        Take $b\in U$ such that $y\mathrel{h} b$ by \cref{preimg_iff}.
                        Show that $y\notin C$.
                        \begin{subproof}
                            Suppose not.
                            Then $y\in C$.
                            Thus $y\in X$ by \cref{subseteq}.
                            $(y,y)\in f$ by \cref{id_elem_intro}.
                            Thus $(y,y)\in h$ by \cref{union_intro_left}.
                            Thus $h(y) = y$ by \cref{function_apply_intro}.
                            Thus $y = b$ by \cref{function_apply_intro}.
                            Thus $y\in U$.
                            Thus $y\notin C$ by \cref{setminus_elim_right}.
                            Contradiction.
                            Thus $y\notin C$.
                        \end{subproof}
                        Thus $y\in Y\setminus C$ by \cref{setminus_intro}.
                    \end{subproof}
                    Thus $\preimg{h}{U}\subseteq Y\setminus C$ by \cref{subseteq}.
                \end{subproof}
                Show that $Y\setminus C\subseteq\preimg{h}{U}$.
                \begin{subproof}
                    Show that for all $y$ such that $y\in Y\setminus C$ we have $y\in\preimg{h}{U}$.
                    \begin{subproof}
                        Fix $y$.
                        Assume $y\in Y\setminus C$.
                        Thus $y\in Y$ and $y\notin C$ by \cref{setminus_elim_left,setminus_elim_right}.
                        Thus $y\in X$ or $y\in\{p\}$ by \cref{union_iff}.
                        \begin{byCase}
                        \caseOf{$y\in X$.}
                            Thus $y\in Y_1$ by \cref{union_intro_left}.
                            Show that $y\in U$.
                            \begin{subproof}
                                Suppose not.
                                Then $y\notin U$.
                                Thus $y\in C$ by \cref{setminus_intro}.
                                Contradiction.
                                Thus $y\in U$.
                            \end{subproof}
                            $(y,y)\in f$ by \cref{id_elem_intro}.
                            Thus $(y,y)\in h$ by \cref{union_intro_left}.
                            Thus $h(y) = y$ by \cref{function_apply_intro}.
                            Thus $y\mathrel{h} y$.
                            Thus $y\in\preimg{h}{U}$ by \cref{preimg_iff}.
                        \caseOf{$y\in\{p\}$.}
                            Thus $y = p$ by \cref{singleton_iff}.
                            $(p,q)\in g$ by \cref{singleton_intro}.
                            Thus $(p,q)\in h$ by \cref{union_intro_right}.
                            Thus $h(y) = q$ by \cref{function_apply_intro}.
                            Thus $q\in U$.
                            Thus $y\mathrel{h} q$.
                            Thus $y\in\preimg{h}{U}$ by \cref{preimg_iff}.
                        \end{byCase}
                    \end{subproof}
                    Thus $Y\setminus C\subseteq\preimg{h}{U}$ by \cref{subseteq}.
                \end{subproof}
                Thus $\preimg{h}{U} = Y\setminus C$ by \cref{subseteq_antisymmetric}.
            \end{subproof}
            Thus $\preimg{h}{U}$ is open in $Y$ with respect to $S$ by \cref{subseteq}.
        \end{byCase}
    \end{subproof}
    Thus $h$ is continuous from $Y$ to $Y_1$ with respect to $S$ and $S_1$ by \cref{continuous}.

    Thus $h$ is a homeomorphism between $Y$ and $Y_1$ with respect to $S$ and $S_1$
        by \cref{compact_hausdorff_bijection_homeomorphism}.

    Show that $\restrl{h}{X} = \identity{X}$.
    \begin{subproof}
        Show that $\restrl{h}{X}\subseteq \identity{X}$.
        \begin{subproof}
            Show that for all $w$ such that $w\in\restrl{h}{X}$ we have $w\in\identity{X}$.
            \begin{subproof}
                Fix $w$.
                Assume $w\in\restrl{h}{X}$.
                Then there exist $x,y$ such that $w = (x,y)$ and $x\in X$ by \cref{restrl}.
                Thus $\restrl{h}{X}\subseteq h$ by \cref{restrl_subseteq}.
                Thus $w\in h$ by \cref{subseteq}.
                Thus $(x,y)\in h$.
                Thus $h(x) = y$ by \cref{function_apply_intro}.
                $(x,x)\in f$ by \cref{id_elem_intro}.
                Thus $(x,x)\in h$ by \cref{union_intro_left}.
                Thus $h(x) = x$ by \cref{function_apply_intro}.
                Thus $y = x$.
                Thus $w\in\identity{X}$ by \cref{id_elem_intro}.
            \end{subproof}
            Thus $\restrl{h}{X}\subseteq \identity{X}$ by \cref{subseteq}.
        \end{subproof}
        Show that $\identity{X}\subseteq\restrl{h}{X}$.
        \begin{subproof}
            Show that for all $w$ such that $w\in\identity{X}$ we have $w\in\restrl{h}{X}$.
            \begin{subproof}
                Fix $w$.
                Assume $w\in\identity{X}$.
                Then there exists $x$ such that $w = (x,x)$ and $x\in X$ by \cref{id_elem_inspect}.
                Thus $(x,x)\in f$ by \cref{id_elem_intro}.
                Thus $(x,x)\in h$ by \cref{union_intro_left}.
                Thus $w\in\restrl{h}{X}$ by \cref{restrl}.
            \end{subproof}
            Thus $\identity{X}\subseteq\restrl{h}{X}$ by \cref{subseteq}.
        \end{subproof}
        Thus $\restrl{h}{X} = \identity{X}$ by \cref{subseteq_antisymmetric}.
    \end{subproof}
    Thus there exists $h$ such that
        $h$ is a homeomorphism between $Y$ and $Y_1$ with respect to $S$ and $S_1$ and
        $\restrl{h}{X} = \identity{X}$.
\end{proof}


% from §29 Item 6: compactification
\begin{definition}\label{compactification}
    $Y$ is a compactification of $X$ with respect to $S$ and $T$ iff
        $S$ is a topology on $Y$ and
        $T$ is a topology on $X$ and
        $X$ is a subspace of $Y$ with respect to $S$ and
        $T = \subspaceTopology{S}{X}$ and
        $X\neq Y$ and
        $\closure{X}{Y}{S} = Y$ and
        $Y$ is compact with respect to $S$ and
        $Y$ is Hausdorff with respect to $S$.
\end{definition}

% from §29 Item 7: one-point compactification
\begin{definition}\label{one_point_compactification}
    $Y$ is a one-point compactification of $X$ with respect to $S$ and $T$ iff
        $Y$ is a compactification of $X$ with respect to $S$ and $T$ and
        there exists $p$ such that $Y\setminus X = \{p\}$.
\end{definition}

% from §29 Item 8: one-point compactification characterization
\begin{theorem}\label{one_point_compactification_characterization}
    Assume $T$ is a topology on $X$.
    Then there exist $Y,S$ such that $Y$ is a one-point compactification of $X$ with respect to $S$ and $T$
        iff $X$ is locally compact with respect to $T$ and $X$ is Hausdorff with respect to $T$
            and $X$ is not compact with respect to $T$.
\end{theorem}
\begin{proof}
    Show that if there exist $Y,S$ such that $Y$ is a one-point compactification of $X$ with respect to $S$ and $T$ then
        $X$ is locally compact with respect to $T$ and $X$ is Hausdorff with respect to $T$ and
            $X$ is not compact with respect to $T$.
    \begin{subproof}
        Assume there exist $Y,S$ such that $Y$ is a one-point compactification of $X$ with respect to $S$ and $T$.
        Take $Y,S$ such that $Y$ is a one-point compactification of $X$ with respect to $S$ and $T$.
        Take $p$ such that $Y\setminus X = \{p\}$ by \cref{one_point_compactification}.
        $Y$ is a compactification of $X$ with respect to $S$ and $T$ by \cref{one_point_compactification}.
        $T = \subspaceTopology{S}{X}$ by \cref{compactification}.
        $Y$ is compact with respect to $S$ by \cref{compactification}.
        $Y$ is Hausdorff with respect to $S$ by \cref{compactification}.
        $X\subseteq Y$ by \cref{subspace_of,compactification}.
        $S$ is a topology on $Y$ by \cref{compactification}.

        Thus there exist $Y_0,S_0,p_0$ such that
            $S_0$ is a topology on $Y_0$ and
            $X\subseteq Y_0$ and
            $T = \subspaceTopology{S_0}{X}$ and
            $Y_0\setminus X = \{p_0\}$ and
            $Y_0$ is compact with respect to $S_0$ and
            $Y_0$ is Hausdorff with respect to $S_0$.

        Thus $X$ is locally compact with respect to $T$ and $X$ is Hausdorff with respect to $T$
            by \cref{one_point_compactification_existence_reverse}.

        Show that $X$ is not compact with respect to $T$.
        \begin{subproof}
            Suppose not.
            Then $X$ is compact with respect to $T$.
            Thus $X$ is a compact subspace of $Y$ with respect to $S$ by \cref{compact_subspace,compactification}.
            Thus $X$ is closed in $Y$ with respect to $S$ by \cref{compact_subspace_closed}.
            $X\subseteq X$ by \cref{subseteq_refl}.
            Thus $\closure{X}{Y}{S}\subseteq X$ by \cref{closure_subseteq_closed}.
            $X\subseteq \closure{X}{Y}{S}$ by \cref{subseteq_closure}.
            Thus $\closure{X}{Y}{S} = X$ by \cref{subseteq_antisymmetric}.
            $\closure{X}{Y}{S} = Y$ by \cref{compactification}.
            Thus $X = Y$.
            Contradiction by \cref{compactification}.
        \end{subproof}
        Thus $X$ is locally compact with respect to $T$ and $X$ is Hausdorff with respect to $T$
            and $X$ is not compact with respect to $T$.
    \end{subproof}

    Show that if $X$ is locally compact with respect to $T$ and $X$ is Hausdorff with respect to $T$
        and $X$ is not compact with respect to $T$
        then there exist $Y,S$ such that $Y$ is a one-point compactification of $X$ with respect to $S$ and $T$.
    \begin{subproof}
        Assume $X$ is locally compact with respect to $T$ and $X$ is Hausdorff with respect to $T$
            and $X$ is not compact with respect to $T$.
        Take $Y,S,p$ such that
            $S$ is a topology on $Y$ and
            $X\subseteq Y$ and
            $T = \subspaceTopology{S}{X}$ and
            $Y\setminus X = \{p\}$ and
            $Y$ is compact with respect to $S$ and
            $Y$ is Hausdorff with respect to $S$
            by \cref{one_point_compactification_existence_forward}.
        Thus $Y = X\union (Y\setminus X)$ by \cref{setminus_partition}.
        Thus $Y = X\union\{p\}$.

        Show that $p\in \closure{X}{Y}{S}$.
        \begin{subproof}
            Suppose not.
            Then $p\notin \closure{X}{Y}{S}$.
            Show that $\closure{X}{Y}{S}\subseteq X$.
            \begin{subproof}
                Show that for all $y$ such that $y\in\closure{X}{Y}{S}$ we have $y\in X$.
                \begin{subproof}
                    Fix $y$.
                    Assume $y\in\closure{X}{Y}{S}$.
                    $\closure{X}{Y}{S}$ is closed in $Y$ with respect to $S$ by \cref{closure_closed_in}.
                    Thus $\closure{X}{Y}{S}\subseteq Y$ by \cref{closed_in}.
                    Thus $y\in Y$ by \cref{subseteq}.
                    Thus $y\in X$ or $y\in\{p\}$ by \cref{union_iff}.
                    Suppose not.
                    Then $y\in\{p\}$.
                    Thus $y = p$ by \cref{singleton_iff}.
                    Thus $y\notin \closure{X}{Y}{S}$.
                    Contradiction.
                    Thus $y\in X$.
                \end{subproof}
                Thus $\closure{X}{Y}{S}\subseteq X$ by \cref{subseteq}.
            \end{subproof}
            $X\subseteq \closure{X}{Y}{S}$ by \cref{subseteq_closure}.
            Thus $\closure{X}{Y}{S} = X$ by \cref{subseteq_antisymmetric}.
            Thus $X$ is closed in $Y$ with respect to $S$ by \cref{closure_closed_in}.
            Thus $X$ is compact with respect to $T$ by \cref{closed_subspace_compact}.
            Contradiction.
        \end{subproof}
        Thus $p\in \closure{X}{Y}{S}$.

        $\closure{X}{Y}{S}$ is closed in $Y$ with respect to $S$ by \cref{closure_closed_in}.
        Thus $\closure{X}{Y}{S}\subseteq Y$ by \cref{closed_in}.
        $X\subseteq \closure{X}{Y}{S}$ by \cref{subseteq_closure}.
        Thus $\{p\}\subseteq \closure{X}{Y}{S}$ by \cref{singleton_subset_intro}.
        Thus $X\union\{p\}\subseteq \closure{X}{Y}{S}$ by \cref{union_subsets_is_subset}.
        Thus $Y\subseteq \closure{X}{Y}{S}$.
        Thus $\closure{X}{Y}{S} = Y$ by \cref{subseteq_antisymmetric}.

        $p\in\{p\}$ by \cref{singleton_intro}.
        Thus $p\in Y\setminus X$.
        Thus $p\notin X$ by \cref{setminus_elim_right}.
        Thus $X\neq Y$.
        Thus $X$ is a subspace of $Y$ with respect to $S$ by \cref{subspace_of}.
        Thus $Y$ is a compactification of $X$ with respect to $S$ and $T$ by \cref{compactification}.
        Thus there exist $Y,S$ such that $Y$ is a one-point compactification of $X$ with respect to $S$ and $T$
            by \cref{one_point_compactification}.
    \end{subproof}
\end{proof}

% from §29 Item 9: local compactness via neighborhoods
\begin{theorem}\label{locally_compact_neighborhood_characterization}
    Assume $T$ is a topology on $X$.
    Assume $X$ is Hausdorff with respect to $T$.
    Then $X$ is locally compact with respect to $T$ iff
        for all $x\in X$ we have
            for all $U$ such that $U$ is a neighborhood of $x$ in $X$ with respect to $T$ we have
                there exists $V$ such that
                    $V$ is a neighborhood of $x$ in $X$ with respect to $T$ and
                    $\closure{V}{X}{T}$ is a compact subspace of $X$ with respect to $T$ and
                    $\closure{V}{X}{T}\subseteq U$.
\end{theorem}
\begin{proof}
    Show that if for all $x\in X$ we have
        for all $U$ such that $U$ is a neighborhood of $x$ in $X$ with respect to $T$ we have
            there exists $V$ such that
                $V$ is a neighborhood of $x$ in $X$ with respect to $T$ and
                $\closure{V}{X}{T}$ is a compact subspace of $X$ with respect to $T$ and
                $\closure{V}{X}{T}\subseteq U$
        then $X$ is locally compact with respect to $T$.
    \begin{subproof}
        Assume for all $x\in X$ we have
            for all $U$ such that $U$ is a neighborhood of $x$ in $X$ with respect to $T$ we have
                there exists $V$ such that
                    $V$ is a neighborhood of $x$ in $X$ with respect to $T$ and
                    $\closure{V}{X}{T}$ is a compact subspace of $X$ with respect to $T$ and
                    $\closure{V}{X}{T}\subseteq U$.
        Show that for all $x\in X$ we have $X$ is locally compact at $x$ with respect to $T$.
        \begin{subproof}
            Fix $x$.
            Assume $x\in X$.
            Let $U = X$.
            $U\in T$ by \cref{topology_on}.
            Thus $U$ is open in $X$ with respect to $T$ by \cref{open_in}.
            Thus $U$ is a neighborhood of $x$ in $X$ with respect to $T$ by \cref{neighborhood}.
            Take $V$ such that
                $V$ is a neighborhood of $x$ in $X$ with respect to $T$ and
                $\closure{V}{X}{T}$ is a compact subspace of $X$ with respect to $T$ and
                $\closure{V}{X}{T}\subseteq U$.
            $V$ is open in $X$ with respect to $T$ by \cref{neighborhood}.
            Thus $V\in T$ by \cref{open_in}.
            $T\subseteq\pow{X}$ by \cref{topology_on}.
            Thus $V\in\pow{X}$ by \cref{subseteq}.
            Thus $V\subseteq X$ by \cref{pow_iff}.
            $V\subseteq \closure{V}{X}{T}$ by \cref{subseteq_closure}.
            Thus $X$ is locally compact at $x$ with respect to $T$ by \cref{locally_compact_at_point}.
        \end{subproof}
        Thus $X$ is locally compact with respect to $T$ by \cref{locally_compact}.
    \end{subproof}

    Show that if $X$ is locally compact with respect to $T$ then
        for all $x\in X$ we have
            for all $U$ such that $U$ is a neighborhood of $x$ in $X$ with respect to $T$ we have
                there exists $V$ such that
                    $V$ is a neighborhood of $x$ in $X$ with respect to $T$ and
                    $\closure{V}{X}{T}$ is a compact subspace of $X$ with respect to $T$ and
                    $\closure{V}{X}{T}\subseteq U$.
    \begin{subproof}
        Assume $X$ is locally compact with respect to $T$.
        Let $p = X$.
        $p\notin X$ by \cref{set_notin_self}.
        Let $Y = X\union\{p\}$.
        Let $S = \onePointTopology{T}{X}{p}$.
        $S$ is a topology on $Y$ by \cref{one_point_topology_is_topology}.
        $Y$ is compact with respect to $S$ by \cref{one_point_topology_compact}.
        $Y$ is Hausdorff with respect to $S$ by \cref{one_point_topology_hausdorff}.
        $T = \subspaceTopology{S}{X}$ by \cref{one_point_topology_subspace}.
        $X\subseteq Y$ by \cref{union_upper_left}.
        Thus $X$ is a subspace of $Y$ with respect to $S$ by \cref{subspace_of}.

        Show that for all $x\in X$ we have
            for all $U$ such that $U$ is a neighborhood of $x$ in $X$ with respect to $T$ we have
                there exists $V$ such that
                    $V$ is a neighborhood of $x$ in $X$ with respect to $T$ and
                    $\closure{V}{X}{T}$ is a compact subspace of $X$ with respect to $T$ and
                    $\closure{V}{X}{T}\subseteq U$.
        \begin{subproof}
            Fix $x$.
            Assume $x\in X$.
            Fix $U$.
            Assume $U$ is a neighborhood of $x$ in $X$ with respect to $T$.
            $U$ is open in $X$ with respect to $T$ by \cref{neighborhood}.
            Thus $U\in T$ by \cref{open_in}.
            $T\subseteq\pow{X}$ by \cref{topology_on}.
            Thus $U\in\pow{X}$ by \cref{subseteq}.
            Thus $U\subseteq X$ by \cref{pow_iff}.

            $X\in T$ by \cref{topology_on}.
            $X\subseteq X$ by \cref{subseteq_refl}.
            Thus $X\in S$ by \cref{one_point_topology_intro_left}.
            Thus $X$ is open in $Y$ with respect to $S$ by \cref{open_in}.
            Thus $U$ is open in $Y$ with respect to $S$ by
                \cref{open_in_subspace_open_in_space,subspace_of}.

            Let $C = Y\setminus U$.
            $C\subseteq Y$ by \cref{setminus_subseteq}.
            Thus $U\in S$ by \cref{open_in}.
            $S\subseteq\pow{Y}$ by \cref{topology_on}.
            Thus $U\in\pow{Y}$ by \cref{subseteq}.
            Thus $U\subseteq Y$ by \cref{pow_iff}.
            $Y\setminus C = Y\setminus (Y\setminus U)$.
            Thus $Y\setminus C = Y\inter U$ by \cref{setminus_setminus}.
            Thus $Y\inter U = U$ by \cref{inter_absorb_supseteq_left}.
            Thus $Y\setminus C = U$.
            Thus $Y\setminus C$ is open in $Y$ with respect to $S$ by \cref{open_in}.
            Thus $C$ is closed in $Y$ with respect to $S$ by \cref{closed_in}.
            Thus $C$ is compact with respect to $\subspaceTopology{S}{C}$ by \cref{closed_subspace_compact}.

            $x\in U$ by \cref{neighborhood}.
            Thus $x\notin C$ by \cref{setminus_elim_right}.
            Take $V,W$ such that
                $V$ is open in $Y$ with respect to $S$ and
                $W$ is open in $Y$ with respect to $S$ and
                $x\in V$ and
                $C\subseteq W$ and
                $V$ is disjoint from $W$
                by \cref{compact_hausdorff_separation}.

            $p\notin X$ by \cref{set_notin_self}.
            Thus $p\notin U$ by \cref{subseteq}.
            $p\in\{p\}$ by \cref{singleton_intro}.
            Thus $p\in Y$ by \cref{union_intro_right}.
            Thus $p\in C$ by \cref{setminus_intro}.
            Thus $p\in W$ by \cref{subseteq}.
            Thus $p\notin V$ by \cref{disjoint}.

            Show that $V\subseteq X$.
            \begin{subproof}
                Show that for all $y$ such that $y\in V$ we have $y\in X$.
                \begin{subproof}
                    Fix $y$.
                    Assume $y\in V$.
                    $V\in S$ by \cref{open_in}.
                    $S\subseteq\pow{Y}$ by \cref{topology_on}.
                    Thus $V\in\pow{Y}$ by \cref{subseteq}.
                    Thus $y\in Y$ by \cref{pow_iff,subseteq}.
                    Thus $y\in X$ or $y\in\{p\}$ by \cref{union_iff}.
                    Suppose not.
                    Then $y\in\{p\}$.
                    Thus $y = p$ by \cref{singleton_iff}.
                    Thus $y\notin V$.
                    Contradiction.
                    Thus $y\in X$.
                \end{subproof}
                Thus $V\subseteq X$ by \cref{subseteq}.
            \end{subproof}

            $V = X\inter V$ by \cref{inter_absorb_supseteq_left}.
            $V\in S$ by \cref{open_in}.
            Thus $V\in T$ by \cref{subspace_topology}.
            Thus $V$ is open in $X$ with respect to $T$ by \cref{open_in}.
            Thus $V$ is a neighborhood of $x$ in $X$ with respect to $T$ by \cref{neighborhood}.
            Thus $V\subseteq Y$ by \cref{subseteq_transitive}.

            Let $B = \closure{V}{Y}{S}$.
            $B$ is closed in $Y$ with respect to $S$ by \cref{closure_closed_in}.
            Thus $B\subseteq Y$ by \cref{closed_in}.
            Thus $B$ is compact with respect to $\subspaceTopology{S}{B}$ by \cref{closed_subspace_compact}.

            $W$ is open in $Y$ with respect to $S$ by \cref{open_in}.
            $Y\setminus W\subseteq Y$ by \cref{setminus_subseteq}.
            $Y\setminus (Y\setminus W) = Y\inter W$ by \cref{setminus_setminus}.
            $W\in S$ by \cref{open_in}.
            $S\subseteq\pow{Y}$ by \cref{topology_on}.
            Thus $W\in\pow{Y}$ by \cref{subseteq}.
            Thus $W\subseteq Y$ by \cref{pow_iff}.
            Thus $Y\inter W = W$ by \cref{inter_absorb_supseteq_left}.
            Thus $Y\setminus (Y\setminus W) = W$.
            Thus $Y\setminus (Y\setminus W)$ is open in $Y$ with respect to $S$ by \cref{open_in}.
            Thus $Y\setminus W$ is closed in $Y$ with respect to $S$ by \cref{closed_in}.
            Show that $V\subseteq Y\setminus W$.
            \begin{subproof}
                Show that for all $y$ such that $y\in V$ we have $y\in Y\setminus W$.
                \begin{subproof}
                    Fix $y$.
                    Assume $y\in V$.
                    Thus $y\notin W$ by \cref{disjoint}.
                    $V\subseteq Y$ by \cref{subseteq}.
                    Thus $y\in Y$ by \cref{subseteq}.
                    Thus $y\in Y\setminus W$ by \cref{setminus_intro}.
                \end{subproof}
                Thus $V\subseteq Y\setminus W$ by \cref{subseteq}.
            \end{subproof}
            Thus $B\subseteq Y\setminus W$ by \cref{closure_subseteq_closed}.

            Show that $Y\setminus W\subseteq Y\setminus C$.
            \begin{subproof}
                Show that for all $y$ such that $y\in Y\setminus W$ we have $y\in Y\setminus C$.
                \begin{subproof}
                    Fix $y$.
                    Assume $y\in Y\setminus W$.
                    Thus $y\in Y$ and $y\notin W$ by \cref{setminus_elim_left,setminus_elim_right}.
                    Suppose not.
                    Then $y\in C$.
                    Thus $y\in W$ by \cref{subseteq}.
                    Contradiction.
                    Thus $y\notin C$.
                    Thus $y\in Y\setminus C$ by \cref{setminus_intro}.
                \end{subproof}
                Thus $Y\setminus W\subseteq Y\setminus C$ by \cref{subseteq}.
            \end{subproof}
            $Y\setminus C = Y\setminus (Y\setminus U)$.
            Thus $Y\setminus C = Y\inter U$ by \cref{setminus_setminus}.
            Thus $Y\inter U = U$ by \cref{inter_absorb_supseteq_left}.
            Thus $Y\setminus C = U$.
            Thus $Y\setminus W\subseteq U$ by \cref{subseteq}.
            Thus $B\subseteq U$ by \cref{subseteq_transitive}.
            Thus $B\subseteq X$ by \cref{subseteq_transitive}.

            Let $T_B = \subspaceTopology{T}{B}$.
            Thus $\subspaceTopology{S}{B} = T_B$ by \cref{subspace_subspace_topology}.
            Thus $B$ is compact with respect to $T_B$.
            Thus $B$ is a compact subspace of $X$ with respect to $T$ by \cref{compact_subspace}.

            $\closure{V}{X}{T} = \closure{V}{Y}{S}\inter X$ by \cref{closure_in_subspace,subspace_of}.
            Thus $\closure{V}{X}{T} = B\inter X$.
            Thus $\closure{V}{X}{T} = X\inter B$ by \cref{inter_comm}.
            Thus $\closure{V}{X}{T} = B$ by \cref{inter_absorb_supseteq_left}.
            Thus $\closure{V}{X}{T}$ is a compact subspace of $X$ with respect to $T$.
            Thus $\closure{V}{X}{T}\subseteq U$ by \cref{subseteq}.
        \end{subproof}
    \end{subproof}
\end{proof}

% from §29 Item 10: locally compact subspaces
\begin{corollary}\label{locally_compact_subspace}
    Assume $T$ is a topology on $X$.
    Assume $X$ is locally compact with respect to $T$.
    Assume $X$ is Hausdorff with respect to $T$.
    Suppose $A$ is a subspace of $X$ with respect to $T$.
    If $A$ is closed in $X$ with respect to $T$ or $A$ is open in $X$ with respect to $T$ then
        $A$ is locally compact with respect to $\subspaceTopology{T}{A}$.
\end{corollary}
\begin{proof}
    Let $T_A = \subspaceTopology{T}{A}$.
    Assume $A$ is closed in $X$ with respect to $T$ or $A$ is open in $X$ with respect to $T$.
    \begin{byCase}
    \caseOf{$A$ is closed in $X$ with respect to $T$.}
        Show that $A$ is locally compact with respect to $T_A$.
        \begin{subproof}
            Show that for all $x\in A$ we have $A$ is locally compact at $x$ with respect to $T_A$.
            \begin{subproof}
                Fix $x$.
                Assume $x\in A$.
                $A\subseteq X$ by \cref{subspace_of}.
                Thus $x\in X$ by \cref{subseteq}.
                $X$ is locally compact at $x$ with respect to $T$ by \cref{locally_compact}.
                Take $C$ such that
                    $C$ is a compact subspace of $X$ with respect to $T$ and
                    there exists $U$ such that
                        $U$ is a neighborhood of $x$ in $X$ with respect to $T$ and
                        $U\subseteq C$
                    by \cref{locally_compact_at_point}.
                Take $U$ such that
                    $U$ is a neighborhood of $x$ in $X$ with respect to $T$ and
                    $U\subseteq C$.
                Let $D = C\inter A$.
                $D\subseteq A$ by \cref{inter_lower_right}.

                Let $T_C = \subspaceTopology{T}{C}$.
                $C$ is compact with respect to $T_C$ by \cref{compact_subspace}.
                $C$ is a subspace of $X$ with respect to $T$ by \cref{subspace_of,compact_subspace}.
                Show that $D$ is closed in $C$ with respect to $T_C$.
                \begin{subproof}
                    $D = A\inter C$ by \cref{inter_comm}.
                    Thus $D$ is closed in $C$ with respect to $T_C$ by \cref{closed_in_subspace_iff}.
                \end{subproof}
                $D\subseteq C$ by \cref{inter_lower_left}.
                Thus $T_C$ is a topology on $C$ by \cref{subspace_topology_is_topology,subspace_of}.
                $C\subseteq X$ by \cref{compact_subspace}.
                Thus $D$ is compact with respect to $\subspaceTopology{T_C}{D}$ by \cref{closed_subspace_compact}.
                Let $T_D = \subspaceTopology{T}{D}$.
                Thus $\subspaceTopology{T_C}{D} = T_D$ by \cref{subspace_subspace_topology}.
                Let $T_{AD} = \subspaceTopology{T_A}{D}$.
                Thus $T_D = T_{AD}$ by \cref{subspace_subspace_topology}.
                Thus $D$ is compact with respect to $T_{AD}$.
                Thus $D$ is a compact subspace of $A$ with respect to $T_A$ by \cref{compact_subspace}.

                $U$ is open in $X$ with respect to $T$ by \cref{neighborhood}.
                Thus $U\in T$ by \cref{open_in}.
                $U\inter A = A\inter U$ by \cref{inter_comm}.
                Thus $U\inter A\in T_A$ by \cref{subspace_topology}.
                Thus $T_A$ is a topology on $A$ by \cref{subspace_topology_is_topology,subspace_of}.
                Thus $U\inter A$ is open in $A$ with respect to $T_A$ by \cref{open_in}.
                $x\in U$ by \cref{neighborhood}.
                Thus $x\in U\inter A$ by \cref{inter_intro}.
                Thus $U\inter A$ is a neighborhood of $x$ in $A$ with respect to $T_A$ by \cref{neighborhood}.
                Show that $U\inter A\subseteq C\inter A$.
                \begin{subproof}
                    $U\inter A\subseteq U$ by \cref{inter_lower_left}.
                    Thus $U\inter A\subseteq C$ by \cref{subseteq_transitive}.
                    $U\inter A\subseteq A$ by \cref{inter_lower_right}.
                    Thus $U\inter A\subseteq C\inter A$ by \cref{subseteq_inter_iff}.
                \end{subproof}
                Thus $A$ is locally compact at $x$ with respect to $T_A$ by \cref{locally_compact_at_point}.
            \end{subproof}
            Thus $A$ is locally compact with respect to $T_A$ by \cref{locally_compact}.
        \end{subproof}
    \caseOf{$A$ is open in $X$ with respect to $T$.}
        Show that $A$ is locally compact with respect to $T_A$.
        \begin{subproof}
            Show that for all $x\in A$ we have $A$ is locally compact at $x$ with respect to $T_A$.
            \begin{subproof}
                Fix $x$.
                Assume $x\in A$.
                $A\subseteq X$ by \cref{subspace_of}.
                Thus $x\in X$ by \cref{subseteq}.
                $A$ is a neighborhood of $x$ in $X$ with respect to $T$ by \cref{neighborhood}.
                Thus for all $y\in X$ we have
                    for all $U$ such that $U$ is a neighborhood of $y$ in $X$ with respect to $T$ we have
                        there exists $V$ such that
                            $V$ is a neighborhood of $y$ in $X$ with respect to $T$ and
                            $\closure{V}{X}{T}$ is a compact subspace of $X$ with respect to $T$ and
                            $\closure{V}{X}{T}\subseteq U$
                    by \cref{locally_compact_neighborhood_characterization}.
                Take $V$ such that
                    $V$ is a neighborhood of $x$ in $X$ with respect to $T$ and
                    $\closure{V}{X}{T}$ is a compact subspace of $X$ with respect to $T$ and
                    $\closure{V}{X}{T}\subseteq A$
                    by \cref{locally_compact_neighborhood_characterization}.

                $V$ is open in $X$ with respect to $T$ by \cref{neighborhood}.
                Thus $V\in T$ by \cref{open_in}.
                $T\subseteq\pow{X}$ by \cref{topology_on}.
                Thus $V\in\pow{X}$ by \cref{subseteq}.
                Thus $V\subseteq X$ by \cref{pow_iff}.
                $V\subseteq \closure{V}{X}{T}$ by \cref{subseteq_closure}.
                Thus $V\subseteq A$ by \cref{subseteq_transitive}.
                $V = A\inter V$ by \cref{inter_absorb_supseteq_left}.
                Thus $V\in T_A$ by \cref{subspace_topology}.
                Thus $T_A$ is a topology on $A$ by \cref{subspace_topology_is_topology,subspace_of}.
                Thus $V$ is open in $A$ with respect to $T_A$ by \cref{open_in}.
                Thus $V$ is a neighborhood of $x$ in $A$ with respect to $T_A$ by \cref{neighborhood}.

                Let $D = \closure{V}{X}{T}$.
                $D\subseteq A$ by \cref{subseteq}.
                Let $T_D = \subspaceTopology{T}{D}$.
                Let $T_{AD} = \subspaceTopology{T_A}{D}$.
                Thus $T_D = T_{AD}$ by \cref{subspace_subspace_topology}.
                Thus $D$ is compact with respect to $T_{AD}$ by \cref{compact_subspace}.
                Thus $D$ is a compact subspace of $A$ with respect to $T_A$ by \cref{compact_subspace}.
                Thus $A$ is locally compact at $x$ with respect to $T_A$ by \cref{locally_compact_at_point}.
            \end{subproof}
            Thus $A$ is locally compact with respect to $T_A$ by \cref{locally_compact}.
        \end{subproof}
    \end{byCase}
\end{proof}

% from §29 helper lemma 15: homeomorphism preserves local compactness
\begin{lemma}\label{homeomorphism_preserves_local_compactness}
    Assume $T$ is a topology on $X$.
    Assume $S$ is a topology on $Y$.
    Suppose $f$ is a homeomorphism between $X$ and $Y$ with respect to $T$ and $S$.
    If $Y$ is locally compact with respect to $S$ then $X$ is locally compact with respect to $T$.
\end{lemma}
\begin{proof}
    Assume $Y$ is locally compact with respect to $S$.
    Let $g = \converse{f}$.
    $f$ is a bijection from $X$ to $Y$ by \cref{homeomorphism}.
    Thus $f\in\bijections{X}{Y}$ by \cref{bijections}.
    Thus $f\in\funs{X}{Y}$ by \cref{bijections_to_funs}.
    Thus $f$ is a function by \cref{funs_is_function}.
    Thus $\dom{f} = X$ by \cref{funs}.
    Thus $g$ is a bijection from $Y$ to $X$ by \cref{bijection_converse_is_bijection}.
    Thus $g\in\bijections{Y}{X}$ by \cref{bijections}.
    Thus $\dom{g} = Y$ by \cref{bijections_dom}.
    Thus $g$ is a function by \cref{bijection_is_function}.
    $g$ is continuous from $Y$ to $X$ with respect to $S$ and $T$ by \cref{homeomorphism}.

    Show that for all $x\in X$ we have $X$ is locally compact at $x$ with respect to $T$.
    \begin{subproof}
        Fix $x$.
        Assume $x\in X$.
        Thus $x\in\dom{f}$.
        Let $y = f(x)$.
        $y\in Y$ by \cref{funs_type_apply}.
        $Y$ is locally compact at $y$ with respect to $S$ by \cref{locally_compact}.
        Take $C$ such that
            $C$ is a compact subspace of $Y$ with respect to $S$ and
            there exists $U$ such that
                $U$ is a neighborhood of $y$ in $Y$ with respect to $S$ and
                $U\subseteq C$
            by \cref{locally_compact_at_point}.
        Take $U$ such that
            $U$ is a neighborhood of $y$ in $Y$ with respect to $S$ and
            $U\subseteq C$.

        Let $g_C = \restrl{g}{C}$.
        $C\subseteq Y$ by \cref{compact_subspace}.
        Thus $Y\subseteq \dom{g}$ by \cref{subseteq_refl}.
        Thus $C\subseteq \dom{g}$ by \cref{subseteq_transitive}.
        Thus $C$ is a subspace of $Y$ with respect to $S$ by \cref{subspace_of,compact_subspace}.
        Let $T_C = \subspaceTopology{S}{C}$.
        $g_C$ is continuous from $C$ to $X$ with respect to $T_C$ and $T$ by \cref{continuous_restrict_domain}.
        $C$ is compact with respect to $T_C$ by \cref{compact_subspace}.
        Let $A = \img{g_C}{C}$.
        Let $T_A = \subspaceTopology{T}{A}$.
        $A$ is compact with respect to $T_A$ by \cref{continuous_image_compact}.
        $g_C\in\funs{C}{X}$ by \cref{continuous}.
        $\funs{C}{X}\subseteq\rels{C}{X}$ by \cref{funs_subseteq_rels}.
        Thus $g_C\in\rels{C}{X}$ by \cref{subseteq}.
        Thus $\ran{g_C}\subseteq X$ by \cref{rels_ran_subseteq}.
        $A\subseteq \ran{g_C}$ by \cref{img_subseteq_ran}.
        Thus $A\subseteq X$ by \cref{subseteq_transitive}.
        Thus $A$ is a compact subspace of $X$ with respect to $T$ by \cref{compact_subspace}.

        Let $V = \img{g}{U}$.
        $V = \preimg{f}{U}$ by \cref{preim_eq_img_of_converse}.
        $U$ is open in $Y$ with respect to $S$ by \cref{neighborhood}.
        $y\in U$ by \cref{neighborhood}.
        $f$ is continuous from $X$ to $Y$ with respect to $T$ and $S$ by \cref{homeomorphism}.
        Thus $V$ is open in $X$ with respect to $T$ by \cref{continuous}.
        $x\mathrel{f} y$ by \cref{function_apply_elim}.
        Thus $x\in V$ by \cref{preimg_iff}.
        Thus $V$ is a neighborhood of $x$ in $X$ with respect to $T$ by \cref{neighborhood}.

        $U\subseteq C$ by \cref{subseteq}.
        Thus $V\subseteq \img{g}{C}$ by \cref{img_subseteq}.
        $g_C\in\funs{C}{X}$ by \cref{continuous_restrict_domain}.
        Thus $\dom{g_C} = C$ by \cref{funs}.
        Thus $\img{g_C}{C} = \ran{g_C}$ by \cref{img_dom}.
        Thus $\img{g}{C} = \img{g_C}{C}$ by \cref{restrl_ran}.
        Thus $V\subseteq A$ by \cref{subseteq}.
        Thus $X$ is locally compact at $x$ with respect to $T$ by \cref{locally_compact_at_point}.
    \end{subproof}
    Thus $X$ is locally compact with respect to $T$ by \cref{locally_compact}.
\end{proof}

% from §29 Item 11: locally compact Hausdorff embeds as open subspace
\begin{corollary}\label{locally_compact_open_subspace_compact_hausdorff_forward}
    Assume $T$ is a topology on $X$.
    If $X$ is locally compact with respect to $T$ and $X$ is Hausdorff with respect to $T$ then
        there exist $Y,S,A,f$ such that
            $S$ is a topology on $Y$ and
            $Y$ is compact with respect to $S$ and
            $Y$ is Hausdorff with respect to $S$ and
            $A$ is open in $Y$ with respect to $S$ and
            $f$ is a homeomorphism between $X$ and $A$ with respect to $T$ and $\subspaceTopology{S}{A}$.
\end{corollary}
\begin{proof}
    Assume $X$ is locally compact with respect to $T$ and $X$ is Hausdorff with respect to $T$.
    Let $p = X$.
    $p\notin X$ by \cref{set_notin_self}.
    Let $Y = X\union\{p\}$.
    Let $S = \onePointTopology{T}{X}{p}$.
    $S$ is a topology on $Y$ by \cref{one_point_topology_is_topology}.
    $Y$ is compact with respect to $S$ by \cref{one_point_topology_compact}.
    $Y$ is Hausdorff with respect to $S$ by \cref{one_point_topology_hausdorff}.
    Let $A = X$.
    $A\in T$ by \cref{topology_on}.
    $A\subseteq X$ by \cref{subseteq}.
    Thus $A\in S$ by \cref{one_point_topology_intro_left}.
    Thus $A$ is open in $Y$ with respect to $S$ by \cref{open_in}.
    $T = \subspaceTopology{S}{A}$ by \cref{one_point_topology_subspace}.
    Let $f = \identity{A}$.
    $f$ is a bijection from $X$ to $A$ by \cref{id_elem_bijections}.
    $f$ is continuous from $X$ to $A$ with respect to $T$ and $\subspaceTopology{S}{A}$ by
        \cref{identity_continuous}.
    Show that $\converse{f}$ is continuous from $A$ to $X$ with respect to $\subspaceTopology{S}{A}$ and $T$.
    \begin{subproof}
        $\converse{f}\in\funs{A}{X}$ by \cref{bijective_converse_are_funs}.
        Thus $\converse{f}$ is a function from $A$ to $X$.
        $A\in S$ by \cref{open_in}.
        $S\subseteq \pow{Y}$ by \cref{topology_on}.
        Thus $A\in \pow{Y}$ by \cref{subseteq}.
        Thus $A\subseteq Y$ by \cref{pow_iff}.
        Thus $A$ is a subspace of $Y$ with respect to $S$ by \cref{subspace_of}.
        Thus $\subspaceTopology{S}{A}$ is a topology on $A$ by \cref{subspace_topology_is_topology}.
        Show that for all $U$ such that $U$ is open in $X$ with respect to $T$ we have
            $\preimg{\converse{f}}{U}$ is open in $A$ with respect to $\subspaceTopology{S}{A}$.
        \begin{subproof}
            Fix $U$.
            Assume $U$ is open in $X$ with respect to $T$.
            $U\in T$ by \cref{open_in}.
            $T\subseteq \pow{X}$ by \cref{topology_on}.
            Thus $U\in \pow{X}$ by \cref{subseteq}.
            Thus $U\subseteq X$ by \cref{pow_iff}.
            Thus $U\subseteq A$ by \cref{subseteq}.
            $f$ is a relation by \cref{id_is_relation}.
            $\preimg{\converse{f}}{U} = \img{\converse{\converse{f}}}{U}$ by \cref{preim_eq_img_of_converse}.
            Thus $\preimg{\converse{f}}{U} = \img{f}{U}$ by \cref{converse_converse_eq}.
            $\img{f}{U} = A\inter U$ by \cref{id_img}.
            Thus $\preimg{\converse{f}}{U} = U$ by \cref{inter_absorb_supseteq_left}.
            $U\in T$ by \cref{open_in}.
            Thus $U\in \subspaceTopology{S}{A}$ by \cref{subseteq}.
            Thus $\preimg{\converse{f}}{U}$ is open in $A$ with respect to $\subspaceTopology{S}{A}$ by
                \cref{open_in}.
        \end{subproof}
        Thus $\converse{f}$ is continuous from $A$ to $X$ with respect to $\subspaceTopology{S}{A}$ and $T$
            by \cref{continuous}.
    \end{subproof}
    Thus $f$ is a homeomorphism between $X$ and $A$ with respect to $T$ and $\subspaceTopology{S}{A}$
        by \cref{homeomorphism}.
    Thus there exist $Y,S,A,f$ such that
        $S$ is a topology on $Y$ and
        $Y$ is compact with respect to $S$ and
        $Y$ is Hausdorff with respect to $S$ and
        $A$ is open in $Y$ with respect to $S$ and
        $f$ is a homeomorphism between $X$ and $A$ with respect to $T$ and $\subspaceTopology{S}{A}$.
\end{proof}

% from §29 Item 12: open subspace of compact Hausdorff is locally compact Hausdorff
\begin{corollary}\label{open_subspace_compact_hausdorff_locally_compact_hausdorff}
    Assume $T$ is a topology on $X$.
    If there exist $Y,S,A,f$ such that
        $S$ is a topology on $Y$ and
        $Y$ is compact with respect to $S$ and
        $Y$ is Hausdorff with respect to $S$ and
        $A$ is open in $Y$ with respect to $S$ and
        $f$ is a homeomorphism between $X$ and $A$ with respect to $T$ and $\subspaceTopology{S}{A}$
    then $X$ is locally compact with respect to $T$ and $X$ is Hausdorff with respect to $T$.
\end{corollary}
\begin{proof}
    Assume there exist $Y,S,A,f$ such that
        $S$ is a topology on $Y$ and
        $Y$ is compact with respect to $S$ and
        $Y$ is Hausdorff with respect to $S$ and
        $A$ is open in $Y$ with respect to $S$ and
        $f$ is a homeomorphism between $X$ and $A$ with respect to $T$ and $\subspaceTopology{S}{A}$.
    Take $Y,S,A,f$ such that
        $S$ is a topology on $Y$ and
        $Y$ is compact with respect to $S$ and
        $Y$ is Hausdorff with respect to $S$ and
        $A$ is open in $Y$ with respect to $S$ and
    $f$ is a homeomorphism between $X$ and $A$ with respect to $T$ and $\subspaceTopology{S}{A}$.

    $Y$ is locally compact with respect to $S$ by \cref{compact_implies_locally_compact}.
    $A\in S$ by \cref{open_in}.
    $S\subseteq \pow{Y}$ by \cref{topology_on}.
    Thus $A\in \pow{Y}$ by \cref{subseteq}.
    Thus $A\subseteq Y$ by \cref{pow_iff}.
    Thus $A$ is a subspace of $Y$ with respect to $S$ by \cref{subspace_of}.
    Let $T_A = \subspaceTopology{S}{A}$.
    Thus $T_A$ is a topology on $A$ by \cref{subspace_topology_is_topology}.
    $A$ is Hausdorff with respect to $T_A$ by \cref{subspace_hausdorff}.
    $A$ is locally compact with respect to $T_A$ by \cref{locally_compact_subspace}.
    Thus $X$ is locally compact with respect to $T$ by \cref{homeomorphism_preserves_local_compactness}.

    Show that $X$ is Hausdorff with respect to $T$.
    \begin{subproof}
        Show that for all $x_1,x_2$ such that $x_1\in X$ and $x_2\in X$ and $x_1\neq x_2$ there exist $U_1,U_2$ such that
            $U_1$ is a neighborhood of $x_1$ in $X$ with respect to $T$ and
            $U_2$ is a neighborhood of $x_2$ in $X$ with respect to $T$ and
            $U_1$ is disjoint from $U_2$.
        \begin{subproof}
            Fix $x_1$.
            Fix $x_2$.
            Assume $x_1\in X$ and $x_2\in X$ and $x_1\neq x_2$.
            $f$ is a bijection from $X$ to $A$ by \cref{homeomorphism}.
            $f\in\funs{X}{A}$ by \cref{bijections_to_funs,homeomorphism}.
            $f$ is a function by \cref{bijection_is_function}.
            $f$ is injective by \cref{bijections,homeomorphism}.
            Let $y_1 = f(x_1)$.
            Let $y_2 = f(x_2)$.
            $y_1\in A$ by \cref{funs_type_apply,homeomorphism}.
            $y_2\in A$ by \cref{funs_type_apply,homeomorphism}.
            Show that $y_1\neq y_2$.
            \begin{subproof}
                Suppose not.
                Then $y_1 = y_2$.
                $\dom{f} = X$ by \cref{funs}.
                Thus $x_1\in\dom{f}$ and $x_2\in\dom{f}$.
                Thus $x_1 = x_2$ by \cref{injective_function}.
                Contradiction.
            \end{subproof}
            Take $V_1,V_2$ such that
                $V_1$ is a neighborhood of $y_1$ in $A$ with respect to $T_A$ and
                $V_2$ is a neighborhood of $y_2$ in $A$ with respect to $T_A$ and
                $V_1$ is disjoint from $V_2$
                by \cref{hausdorff}.
            $V_1$ is open in $A$ with respect to $T_A$ by \cref{neighborhood}.
            $V_2$ is open in $A$ with respect to $T_A$ by \cref{neighborhood}.
            Let $U_1 = \preimg{f}{V_1}$.
            Let $U_2 = \preimg{f}{V_2}$.
            $U_1$ is open in $X$ with respect to $T$ by \cref{continuous,homeomorphism}.
            $U_2$ is open in $X$ with respect to $T$ by \cref{continuous,homeomorphism}.
            $\dom{f} = X$ by \cref{funs}.
            Thus $x_1\in\dom{f}$ and $x_2\in\dom{f}$.
            $x_1\mathrel{f} y_1$ by \cref{function_apply_elim}.
            $x_2\mathrel{f} y_2$ by \cref{function_apply_elim}.
            $y_1\in V_1$ by \cref{neighborhood}.
            $y_2\in V_2$ by \cref{neighborhood}.
            Thus $x_1\in U_1$ by \cref{preimg_iff}.
            Thus $x_2\in U_2$ by \cref{preimg_iff}.
            Show that $U_1$ is disjoint from $U_2$.
            \begin{subproof}
                Suppose not.
                Then there exists $z$ such that $z\in U_1$ and $z\in U_2$.
                Take $w_1\in V_1$ such that $z\mathrel{f} w_1$ by \cref{preimg_iff}.
                Take $w_2\in V_2$ such that $z\mathrel{f} w_2$ by \cref{preimg_iff}.
                Thus $f(z) = w_1$ by \cref{function_apply_intro}.
                Thus $f(z) = w_2$ by \cref{function_apply_intro}.
                Thus $w_1 = w_2$.
                Thus $w_1\in V_1\inter V_2$ by \cref{inter_intro}.
                Contradiction by \cref{disjoint}.
            \end{subproof}
            Thus $U_1$ is a neighborhood of $x_1$ in $X$ with respect to $T$ by \cref{neighborhood}.
            Thus $U_2$ is a neighborhood of $x_2$ in $X$ with respect to $T$ by \cref{neighborhood}.
            Thus there exist $U_1,U_2$ such that
                $U_1$ is a neighborhood of $x_1$ in $X$ with respect to $T$ and
                $U_2$ is a neighborhood of $x_2$ in $X$ with respect to $T$ and
                $U_1$ is disjoint from $U_2$.
        \end{subproof}
        Thus $X$ is Hausdorff with respect to $T$ by \cref{hausdorff}.
    \end{subproof}
    Thus $X$ is locally compact with respect to $T$ and $X$ is Hausdorff with respect to $T$.
\end{proof}
% from §30 Item 1: sequence implies closure
\begin{lemma}\label{sequence_implies_closure_first_countable}
    Assume $T$ is a topology on $X$.
    Suppose $A\subseteq X$.
    Suppose $s$ is a sequence in $A$.
    Suppose $x\in X$.
    Suppose $s$ converges to $x$ in $X$ with respect to $T$.
    Then $x\in\closure{A}{X}{T}$.
\end{lemma}
\begin{proof}
    Follows by \cref{sequence_implies_closure}.
\end{proof}

% from §30 helper lemma 1: nested countable basis at a point
\begin{lemma}\label{nested_countable_basis}
    Suppose $b$ is a countable basis at $x$ in $X$ with respect to $T$.
    Then there exists $c$ such that
        $c$ is a countable basis at $x$ in $X$ with respect to $T$ and
        for all $m,n\in\naturalsPlus$ if $m\leq n$ then $c(n)\subseteq c(m)$.
\end{lemma}
\begin{proof}
    $b$ is a sequence in $\pow{X}$ by \cref{countable_basis_at_point}.
    $1\in\naturalsPlus$ by \cref{real_one_in_naturals}.
    $b(1)$ is a neighborhood of $x$ in $X$ with respect to $T$ by \cref{countable_basis_at_point}.
    Thus $b(1)$ is open in $X$ with respect to $T$ by \cref{neighborhood}.
    Thus $T$ is a topology on $X$ by \cref{open_in}.
    Let $F = \{(n,\img{\restrl{b}{\initialSegment{n}}}{\initialSegment{n}}) \mid n\in\naturalsPlus\}$.
    Let $C = \{(n,\inters{\apply{F}{n}}) \mid n\in\naturalsPlus\}$.

    Show that $F$ is a function on $\naturalsPlus$.
    \begin{subproof}
        Show that $F$ is a function.
        \begin{subproof}
            Show that for all $n,U_1,U_2$ such that $(n,U_1)\in F$ and $(n,U_2)\in F$ we have $U_1 = U_2$.
            \begin{subproof}
                Fix $n$.
                Fix $U_1$.
                Fix $U_2$.
                Assume $(n,U_1)\in F$ and $(n,U_2)\in F$.
                Thus $U_1 = \img{\restrl{b}{\initialSegment{n}}}{\initialSegment{n}}$ and
                    $U_2 = \img{\restrl{b}{\initialSegment{n}}}{\initialSegment{n}}$ by \cref{pair_eq_iff}.
                Thus $U_1 = U_2$.
            \end{subproof}
            Thus $F$ is right-unique by \cref{rightunique}.
            Thus $F$ is a function.
        \end{subproof}
        Show that $\dom{F} = \naturalsPlus$.
        \begin{subproof}
            Show that $\naturalsPlus\subseteq\dom{F}$.
            \begin{subproof}
                Show that for all $n$ such that $n\in\naturalsPlus$ we have $n\in\dom{F}$.
                \begin{subproof}
                    Fix $n$.
                    Assume $n\in\naturalsPlus$.
                    Let $U = \img{\restrl{b}{\initialSegment{n}}}{\initialSegment{n}}$.
                    $(n,U)\in F$ by \cref{pair_eq_iff}.
                    Thus $n\in\dom{F}$ by \cref{dom_intro}.
                \end{subproof}
                Thus $\naturalsPlus\subseteq\dom{F}$ by \cref{subseteq}.
            \end{subproof}
            Show that $\dom{F}\subseteq\naturalsPlus$.
            \begin{subproof}
                Show that for all $n$ such that $n\in\dom{F}$ we have $n\in\naturalsPlus$.
                \begin{subproof}
                    Fix $n$.
                    Assume $n\in\dom{F}$.
                    Take $U$ such that $(n,U)\in F$ by \cref{dom_iff}.
                    Thus $n\in\naturalsPlus$ by \cref{pair_eq_iff}.
                \end{subproof}
                Thus $\dom{F}\subseteq\naturalsPlus$ by \cref{subseteq}.
            \end{subproof}
            Thus $\dom{F} = \naturalsPlus$ by \cref{subseteq_antisymmetric}.
        \end{subproof}
        Thus $F$ is a function on $\naturalsPlus$.
    \end{subproof}

    Show that for all $n\in\naturalsPlus$ we have $F(n) = \img{\restrl{b}{\initialSegment{n}}}{\initialSegment{n}}$.
    \begin{subproof}
        Fix $n$.
        Assume $n\in\naturalsPlus$.
        Let $U = \img{\restrl{b}{\initialSegment{n}}}{\initialSegment{n}}$.
        $(n,U)\in F$ by \cref{pair_eq_iff}.
        Thus $F(n) = U$ by \cref{function_apply_intro}.
    \end{subproof}

    Show that $C$ is a function on $\naturalsPlus$.
    \begin{subproof}
        Show that $C$ is a function.
        \begin{subproof}
            Show that for all $n,U_1,U_2$ such that $(n,U_1)\in C$ and $(n,U_2)\in C$ we have $U_1 = U_2$.
            \begin{subproof}
                Fix $n$.
                Fix $U_1$.
                Fix $U_2$.
                Assume $(n,U_1)\in C$ and $(n,U_2)\in C$.
                Thus $U_1 = \inters{F(n)}$ and $U_2 = \inters{F(n)}$ by \cref{pair_eq_iff}.
                Thus $U_1 = U_2$.
            \end{subproof}
            Thus $C$ is right-unique by \cref{rightunique}.
            Thus $C$ is a function.
        \end{subproof}
        Show that $\dom{C} = \naturalsPlus$.
        \begin{subproof}
            Show that $\naturalsPlus\subseteq\dom{C}$.
            \begin{subproof}
                Show that for all $n$ such that $n\in\naturalsPlus$ we have $n\in\dom{C}$.
                \begin{subproof}
                    Fix $n$.
                    Assume $n\in\naturalsPlus$.
                    Let $U = \inters{F(n)}$.
                    $(n,U)\in C$ by \cref{pair_eq_iff}.
                    Thus $n\in\dom{C}$ by \cref{dom_intro}.
                \end{subproof}
                Thus $\naturalsPlus\subseteq\dom{C}$ by \cref{subseteq}.
            \end{subproof}
            Show that $\dom{C}\subseteq\naturalsPlus$.
            \begin{subproof}
                Show that for all $n$ such that $n\in\dom{C}$ we have $n\in\naturalsPlus$.
                \begin{subproof}
                    Fix $n$.
                    Assume $n\in\dom{C}$.
                    Take $U$ such that $(n,U)\in C$ by \cref{dom_iff}.
                    Thus $n\in\naturalsPlus$ by \cref{pair_eq_iff}.
                \end{subproof}
                Thus $\dom{C}\subseteq\naturalsPlus$ by \cref{subseteq}.
            \end{subproof}
            Thus $\dom{C} = \naturalsPlus$ by \cref{subseteq_antisymmetric}.
        \end{subproof}
        Thus $C$ is a function on $\naturalsPlus$.
    \end{subproof}

    Show that for all $n\in\naturalsPlus$ we have $C(n) = \inters{F(n)}$.
    \begin{subproof}
        Fix $n$.
        Assume $n\in\naturalsPlus$.
        Let $U = \inters{F(n)}$.
        $(n,U)\in C$ by \cref{pair_eq_iff}.
        Thus $C(n) = U$ by \cref{function_apply_intro}.
    \end{subproof}

    Show that for all $n\in\naturalsPlus$ we have $C(n)$ is a neighborhood of $x$ in $X$ with respect to $T$.
    \begin{subproof}
        Fix $n\in\naturalsPlus$.
        Show that $F(n)$ is finite and inhabited.
        \begin{subproof}
            $\initialSegment{n}$ is finite by \cref{initial_segment_finite}.
            $b\in\funs{\naturalsPlus}{\pow{X}}$ by \cref{sequence_in}.
            Thus $\dom{b} = \naturalsPlus$ by \cref{funs}.
            Show that $\initialSegment{n}\subseteq\dom{b}$.
            \begin{subproof}
                Show that for all $i$ such that $i\in\initialSegment{n}$ we have $i\in\dom{b}$.
                \begin{subproof}
                    Fix $i$.
                    Assume $i\in\initialSegment{n}$.
                    $i\in\naturalsPlus$ by \cref{initial_segment_naturalsplus}.
                    Thus $i\in\dom{b}$ by \cref{subseteq}.
                \end{subproof}
                Thus $\initialSegment{n}\subseteq\dom{b}$ by \cref{subseteq}.
            \end{subproof}
            Let $g = \restrl{b}{\initialSegment{n}}$.
            $b$ is a function by \cref{funs_is_function,sequence_in}.
            Thus $b$ is right-unique by \cref{function_rightunique}.
            Thus $b$ is a relation by \cref{funs_is_relation}.
            $g$ is a function by \cref{restrl_of_function_is_function}.
            Show that $\dom{g} = \initialSegment{n}$.
            \begin{subproof}
                Thus $\dom{g}\subseteq\initialSegment{n}$ by \cref{restrl_dom}.
                Show that $\initialSegment{n}\subseteq\dom{g}$.
                \begin{subproof}
                    Show that for all $i$ such that $i\in\initialSegment{n}$ we have $i\in\dom{g}$.
                    \begin{subproof}
                        Fix $i$.
                        Assume $i\in\initialSegment{n}$.
                        $i\in\naturalsPlus$ by \cref{initial_segment_naturalsplus}.
                        Thus $i\in\dom{b}$ by \cref{subseteq}.
                        Take $U$ such that $(i,U)\in b$ by \cref{dom_iff}.
                        Thus $(i,U)\in g$ by \cref{restrl_iff}.
                        Thus $i\in\dom{g}$ by \cref{dom_intro}.
                    \end{subproof}
                    Thus $\initialSegment{n}\subseteq\dom{g}$ by \cref{subseteq}.
                \end{subproof}
                Thus $\dom{g} = \initialSegment{n}$ by \cref{subseteq_antisymmetric}.
            \end{subproof}
            Thus $g$ is a function on $\initialSegment{n}$.
            Thus $F(n) = \img{g}{\initialSegment{n}}$.
            Thus $F(n)$ is finite by \cref{finite_image_function}.
            $1\in\naturalsPlus$ by \cref{real_one_in_naturals}.
            $1 \leq n$ by \cref{real_one_leq_naturalsplus}.
            Thus $1\in\initialSegment{n}$ by \cref{initial_segment_naturalsplus}.
            Take $U_1$ such that $1\mathrel{b} U_1$ by \cref{dom_iff,funs}.
            Thus $U_1\in F(n)$ by \cref{img_iff,restrl_iff}.
            Thus $F(n)$ is inhabited by \cref{nonempty_inhabited}.
        \end{subproof}
        Show that $F(n)\subseteq T$.
        \begin{subproof}
            Show that for all $U$ such that $U\in F(n)$ we have $U\in T$.
            \begin{subproof}
                Fix $U$.
                Assume $U\in F(n)$.
                Thus $U\in\img{\restrl{b}{\initialSegment{n}}}{\initialSegment{n}}$.
                Take $k\in\initialSegment{n}$ such that $k\mathrel{\restrl{b}{\initialSegment{n}}} U$ by \cref{img_iff}.
                Thus $k\mathrel{b} U$ by \cref{restrl_iff}.
                $b\in\funs{\naturalsPlus}{\pow{X}}$ by \cref{sequence_in}.
                Thus $\dom{b} = \naturalsPlus$ by \cref{funs}.
                $k\in\naturalsPlus$ by \cref{initial_segment_naturalsplus}.
                Thus $k\in\dom{b}$ by \cref{subseteq}.
                $b$ is a function by \cref{funs_is_function,sequence_in}.
                Thus $b(k) = U$ by \cref{function_apply_iff}.
                $k\in\naturalsPlus$ by \cref{initial_segment_naturalsplus}.
                Thus $U$ is a neighborhood of $x$ in $X$ with respect to $T$ by \cref{countable_basis_at_point}.
                Thus $U$ is open in $X$ with respect to $T$ by \cref{neighborhood}.
                Thus $U\in T$ by \cref{open_in}.
            \end{subproof}
            Thus $F(n)\subseteq T$ by \cref{subseteq}.
        \end{subproof}
        Thus $C(n)\in T$ by \cref{finite_intersections_open_in_topology}.
        Show that $x\in C(n)$.
        \begin{subproof}
            Show that for all $U$ such that $U\in F(n)$ we have $x\in U$.
            \begin{subproof}
                Fix $U$.
                Assume $U\in F(n)$.
                Thus $U\in\img{\restrl{b}{\initialSegment{n}}}{\initialSegment{n}}$.
                Take $k\in\initialSegment{n}$ such that $k\mathrel{\restrl{b}{\initialSegment{n}}} U$ by \cref{img_iff}.
                Thus $k\mathrel{b} U$ by \cref{restrl_iff}.
                $b\in\funs{\naturalsPlus}{\pow{X}}$ by \cref{sequence_in}.
                Thus $\dom{b} = \naturalsPlus$ by \cref{funs}.
                $k\in\naturalsPlus$ by \cref{initial_segment_naturalsplus}.
                Thus $k\in\dom{b}$ by \cref{subseteq}.
                $b$ is a function by \cref{funs_is_function,sequence_in}.
                Thus $b(k) = U$ by \cref{function_apply_iff}.
                $k\in\naturalsPlus$ by \cref{initial_segment_naturalsplus}.
                Thus $b(k)$ is a neighborhood of $x$ in $X$ with respect to $T$ by \cref{countable_basis_at_point}.
                Thus $x\in b(k)$ by \cref{neighborhood}.
                Thus $x\in U$.
            \end{subproof}
            Thus $x\in C(n)$ by \cref{inters_iff_forall}.
        \end{subproof}
        Thus $C(n)$ is a neighborhood of $x$ in $X$ with respect to $T$ by \cref{neighborhood,open_in}.
    \end{subproof}

    Show that $C$ is a countable basis at $x$ in $X$ with respect to $T$.
    \begin{subproof}
        Show that $C$ is a sequence in $\pow{X}$.
        \begin{subproof}
            Show that $C$ is a function to $\pow{X}$.
            \begin{subproof}
                Show that for all $n\in\dom{C}$ we have $C(n)\in\pow{X}$.
                \begin{subproof}
                    Fix $n$.
                    Assume $n\in\dom{C}$.
                    $\dom{C} = \naturalsPlus$.
                    Thus $n\in\naturalsPlus$ by \cref{subseteq}.
                    $C(n)$ is a neighborhood of $x$ in $X$ with respect to $T$.
                    Thus $C(n)$ is open in $X$ with respect to $T$ by \cref{neighborhood}.
                    Thus $C(n)\in T$ by \cref{open_in}.
                    $T\subseteq\pow{X}$ by \cref{topology_on}.
                    Thus $C(n)\in\pow{X}$ by \cref{subseteq}.
                \end{subproof}
                Thus $C$ is a function to $\pow{X}$.
            \end{subproof}
            Thus $C\in\funs{\naturalsPlus}{\pow{X}}$ by \cref{funs_intro}.
            Thus $C$ is a sequence in $\pow{X}$ by \cref{sequence_in}.
        \end{subproof}
        Show that for all $n\in\naturalsPlus$ we have $C(n)$ is a neighborhood of $x$ in $X$ with respect to $T$.
        \begin{subproof}
            Fix $n\in\naturalsPlus$.
            $C(n)$ is a neighborhood of $x$ in $X$ with respect to $T$.
        \end{subproof}
        Show that for all $U$ such that $U$ is a neighborhood of $x$ in $X$ with respect to $T$
            there exists $n\in\naturalsPlus$ such that $C(n)\subseteq U$.
        \begin{subproof}
            Fix $U$.
            Assume $U$ is a neighborhood of $x$ in $X$ with respect to $T$.
            Take $k\in\naturalsPlus$ such that $b(k)\subseteq U$ by \cref{countable_basis_at_point}.
            Show that $C(k)\subseteq U$.
            \begin{subproof}
                Show that $b(k)\in F(k)$.
                \begin{subproof}
                    $k\in\initialSegment{k}$ by \cref{initial_segment_naturalsplus}.
                    $b\in\funs{\naturalsPlus}{\pow{X}}$ by \cref{sequence_in}.
                    Thus $\dom{b} = \naturalsPlus$ by \cref{funs}.
                    $k\in\naturalsPlus$ by \cref{initial_segment_naturalsplus}.
                    Thus $k\in\dom{b}$ by \cref{subseteq}.
                    $b$ is a function by \cref{funs_is_function,sequence_in}.
                    Thus $(k,b(k))\in b$ by \cref{function_apply_iff}.
                    Thus $(k,b(k))\in\restrl{b}{\initialSegment{k}}$ by \cref{restrl_iff}.
                    Thus $b(k)\in F(k)$ by \cref{img_iff}.
                \end{subproof}
                Thus $C(k)\subseteq b(k)$ by \cref{inters_subseteq_elem}.
                Thus $C(k)\subseteq U$ by \cref{subseteq_transitive}.
            \end{subproof}
            Thus there exists $n\in\naturalsPlus$ such that $C(n)\subseteq U$.
        \end{subproof}
        Thus $C$ is a countable basis at $x$ in $X$ with respect to $T$ by \cref{countable_basis_at_point}.
    \end{subproof}

    Show that for all $m,n$ such that $m\in\naturalsPlus$ and $n\in\naturalsPlus$ and $m\leq n$ we have $C(n)\subseteq C(m)$.
    \begin{subproof}
        Fix $m$.
        Fix $n$.
        Assume $m\in\naturalsPlus$ and $n\in\naturalsPlus$ and $m\leq n$.
        Show that $C(n)\subseteq C(m)$.
        \begin{subproof}
            Show that for all $z$ such that $z\in C(n)$ we have $z\in C(m)$.
            \begin{subproof}
                Fix $z$.
                Assume $z\in C(n)$.
                Thus for all $U$ such that $U\in F(n)$ we have $z\in U$ by \cref{inters_iff_forall}.
                Show that for all $U$ such that $U\in F(m)$ we have $z\in U$.
                \begin{subproof}
                    Fix $U$.
                    Assume $U\in F(m)$.
                    Show that $\initialSegment{m}\subseteq\initialSegment{n}$.
                    \begin{subproof}
                        Show that for all $i$ such that $i\in\initialSegment{m}$ we have $i\in\initialSegment{n}$.
                        \begin{subproof}
                            Fix $i$.
                            Assume $i\in\initialSegment{m}$.
                            Thus $i\in\naturalsPlus$ and $i \leq m$ by \cref{initial_segment_naturalsplus}.
                            $i\in\reals$ by \cref{real_naturals_subset_reals,subseteq}.
                            $m\in\reals$ and $n\in\reals$ by \cref{real_naturals_subset_reals,subseteq}.
                            Thus $i \leq n$ by \cref{real_leq_trans}.
                            Thus $i\in\initialSegment{n}$ by \cref{initial_segment_naturalsplus}.
                        \end{subproof}
                        Thus $\initialSegment{m}\subseteq\initialSegment{n}$ by \cref{subseteq}.
                    \end{subproof}
                    Show that $\img{\restrl{b}{\initialSegment{m}}}{\initialSegment{m}}\subseteq
                        \img{\restrl{b}{\initialSegment{n}}}{\initialSegment{n}}$.
                    \begin{subproof}
                        Show that for all $y$ such that $y\in\img{\restrl{b}{\initialSegment{m}}}{\initialSegment{m}}$ we have
                            $y\in\img{\restrl{b}{\initialSegment{n}}}{\initialSegment{n}}$.
                        \begin{subproof}
                            Fix $y$.
                            Assume $y\in\img{\restrl{b}{\initialSegment{m}}}{\initialSegment{m}}$.
                            Take $a\in\initialSegment{m}$ such that $a\mathrel{\restrl{b}{\initialSegment{m}}} y$ by \cref{img_iff}.
                            Thus $a\mathrel{b} y$ and $a\in\initialSegment{m}$ by \cref{restrl_iff}.
                            Thus $a\in\initialSegment{n}$ by \cref{subseteq}.
                            Thus $a\mathrel{\restrl{b}{\initialSegment{n}}} y$ by \cref{restrl_iff}.
                            Thus $y\in\img{\restrl{b}{\initialSegment{n}}}{\initialSegment{n}}$ by \cref{img_iff}.
                        \end{subproof}
                        Thus $\img{\restrl{b}{\initialSegment{m}}}{\initialSegment{m}}\subseteq
                            \img{\restrl{b}{\initialSegment{n}}}{\initialSegment{n}}$ by \cref{subseteq}.
                    \end{subproof}
                    Thus $U\in F(n)$ by \cref{subseteq}.
                    Thus $z\in U$.
                \end{subproof}
                Show that $F(m)$ is inhabited.
                \begin{subproof}
                    $1\in\naturalsPlus$ by \cref{real_one_in_naturals}.
                    $1 \leq m$ by \cref{real_one_leq_naturalsplus}.
                    Thus $1\in\initialSegment{m}$ by \cref{initial_segment_naturalsplus}.
                    $b\in\funs{\naturalsPlus}{\pow{X}}$ by \cref{sequence_in}.
                    Thus $\dom{b} = \naturalsPlus$ by \cref{funs}.
                    Thus $1\in\dom{b}$ by \cref{subseteq}.
                    $b$ is a function by \cref{funs_is_function,sequence_in}.
                    Thus $(1,b(1))\in b$ by \cref{function_apply_iff}.
                    Thus $(1,b(1))\in\restrl{b}{\initialSegment{m}}$ by \cref{restrl_iff}.
                    Thus $b(1)\in F(m)$ by \cref{img_iff}.
                    Thus $F(m)$ is inhabited by \cref{nonempty_inhabited}.
                \end{subproof}
                Thus $z\in C(m)$ by \cref{inters_iff_forall}.
            \end{subproof}
            Thus $C(n)\subseteq C(m)$ by \cref{subseteq}.
        \end{subproof}
    \end{subproof}

    Thus there exists $c$ such that
        $c$ is a countable basis at $x$ in $X$ with respect to $T$ and
        for all $m,n\in\naturalsPlus$ if $m\leq n$ then $c(n)\subseteq c(m)$.
\end{proof}

% from §30 Item 2: closure implies sequence in a first-countable space
\begin{lemma}\label{closure_implies_sequence_first_countable}
    Assume $T$ is a topology on $X$.
    Suppose $A\subseteq X$.
    Suppose $x\in X$.
    Suppose $X$ satisfies the first countability axiom with respect to $T$.
    Suppose $x\in\closure{A}{X}{T}$.
    Then there exists $s$ such that $s$ is a sequence in $A$ and $s$ converges to $x$ in $X$ with respect to $T$.
\end{lemma}
\begin{proof}
    Take $b$ such that $b$ is a countable basis at $x$ in $X$ with respect to $T$ by \cref{first_countability_axiom}.
    Take $c$ such that
        $c$ is a countable basis at $x$ in $X$ with respect to $T$ and
        for all $m,n\in\naturalsPlus$ if $m\leq n$ then $c(n)\subseteq c(m)$ by \cref{nested_countable_basis}.

    Show that for all $n\in\naturalsPlus$ we have $c(n)$ intersects $A$.
    \begin{subproof}
        Fix $n\in\naturalsPlus$.
        $c(n)$ is a neighborhood of $x$ in $X$ with respect to $T$ by \cref{countable_basis_at_point}.
        Thus $c(n)$ is open in $X$ with respect to $T$ and $x\in c(n)$ by \cref{neighborhood}.
        Thus $c(n)$ intersects $A$ by \cref{closure_iff_intersects_open}.
    \end{subproof}

    Let $R = \{w\in\naturalsPlus\times X \mid \snd{w}\in A\inter \apply{c}{\fst{w}}\}$.
    Show that $R$ is a relation.
    \begin{subproof}
        Show that for all $w$ such that $w\in R$ there exist $n,y$ such that $w = (n,y)$.
        \begin{subproof}
            Fix $w$.
            Assume $w\in R$.
            Thus $w\in\naturalsPlus\times X$.
            Take $n,y$ such that $n\in\naturalsPlus$ and $y\in X$ and $w = (n,y)$ by \cref{times_elem_is_tuple}.
        \end{subproof}
        Thus $R$ is a relation by \cref{relation}.
    \end{subproof}
    Show that for all $n\in\naturalsPlus$ there exists $y$ such that $n\mathrel{R} y$.
    \begin{subproof}
        Fix $n\in\naturalsPlus$.
        $c(n)$ intersects $A$.
        Thus $c(n)\inter A \neq \emptyset$ by \cref{intersects}.
        Thus $A\inter c(n) = c(n)\inter A$ by \cref{inter_comm}.
        Thus $A\inter c(n) \neq \emptyset$.
        Take $y$ such that $y\in A\inter c(n)$ by \cref{nonempty_inhabited}.
        Thus $y\in A$ by \cref{inter_elim_left}.
        Thus $y\in X$ by \cref{subseteq}.
        Thus $(n,y)\in\naturalsPlus\times X$ by \cref{times_tuple_intro}.
        $\fst{(n,y)} = n$ by \cref{fst_eq}.
        $\snd{(n,y)} = y$ by \cref{snd_eq}.
        Thus $\snd{(n,y)}\in A\inter \apply{c}{\fst{(n,y)}}$ by \cref{inter_intro}.
        Thus $(n,y)\in R$.
        Thus $n\mathrel{R} y$ by \cref{pair_eq_iff}.
    \end{subproof}
    Take $s$ such that $s$ is a function on $\naturalsPlus$ and
        for all $n\in\naturalsPlus$ we have $n\mathrel{R} s(n)$ by \cref{axiom_of_choice}.

    Show that $s$ is a sequence in $A$.
    \begin{subproof}
        Show that $s$ is a function to $A$.
        \begin{subproof}
            Show that for all $n\in\dom{s}$ we have $s(n)\in A$.
            \begin{subproof}
                Fix $n$.
                Assume $n\in\dom{s}$.
                $\dom{s} = \naturalsPlus$.
                Thus $n\in\naturalsPlus$ by \cref{subseteq}.
                Thus $n\mathrel{R} s(n)$.
                Thus $(n,s(n))\in R$ by \cref{pair_eq_iff}.
                Thus $\snd{(n,s(n))}\in A\inter \apply{c}{\fst{(n,s(n))}}$.
                $\fst{(n,s(n))} = n$ by \cref{fst_eq}.
                $\snd{(n,s(n))} = s(n)$ by \cref{snd_eq}.
                Thus $s(n)\in A\inter c(n)$.
                Thus $s(n)\in A$ by \cref{inter}.
            \end{subproof}
            Thus $s$ is a function to $A$.
        \end{subproof}
        Thus $s\in\funs{\naturalsPlus}{A}$ by \cref{funs_intro}.
        Thus $s$ is a sequence in $A$ by \cref{sequence_in}.
    \end{subproof}

    Show that $s$ converges to $x$ in $X$ with respect to $T$.
    \begin{subproof}
        Show that $s$ is a sequence in $X$.
        \begin{subproof}
            $s$ is a function to $A$.
            Thus $s$ is a function to $X$ by \cref{function_on_weaken_codom,subseteq}.
            $\dom{s} = \naturalsPlus$ by \cref{funs_elim}.
            Thus $s\in\funs{\naturalsPlus}{X}$ by \cref{funs_intro}.
            Thus $s$ is a sequence in $X$ by \cref{sequence_in}.
        \end{subproof}
        Show that for all $U$ such that $U$ is a neighborhood of $x$ in $X$ with respect to $T$
            there exists $N\in\naturalsPlus$ such that for all $n\in\naturalsPlus$ if $N\leq n$ then $s(n)\in U$.
        \begin{subproof}
            Fix $U$.
            Assume $U$ is a neighborhood of $x$ in $X$ with respect to $T$.
            $c$ is a countable basis at $x$ in $X$ with respect to $T$.
            $1\in\naturalsPlus$ by \cref{real_one_in_naturals}.
            Thus $c(1)$ is a neighborhood of $x$ in $X$ with respect to $T$ and
                for all $V$ such that $V$ is a neighborhood of $x$ in $X$ with respect to $T$
                    there exists $n\in\naturalsPlus$ such that $c(n)\subseteq V$ by \cref{countable_basis_at_point}.
            Thus for all $V$ such that $V$ is a neighborhood of $x$ in $X$ with respect to $T$
                there exists $n\in\naturalsPlus$ such that $c(n)\subseteq V$.
            Thus there exists $N\in\naturalsPlus$ such that $c(N)\subseteq U$.
            Take $N\in\naturalsPlus$ such that $c(N)\subseteq U$.
            Show that for all $n\in\naturalsPlus$ if $N\leq n$ then $s(n)\in U$.
            \begin{subproof}
                Fix $n\in\naturalsPlus$.
                Assume $N\leq n$.
                $c(n)\subseteq c(N)$ by \cref{nested_countable_basis}.
                Thus $c(n)\subseteq U$ by \cref{subseteq_transitive}.
                $n\mathrel{R} s(n)$.
                Thus $(n,s(n))\in R$ by \cref{pair_eq_iff}.
                Thus $\snd{(n,s(n))}\in A\inter \apply{c}{\fst{(n,s(n))}}$.
                $\fst{(n,s(n))} = n$ by \cref{fst_eq}.
                $\snd{(n,s(n))} = s(n)$ by \cref{snd_eq}.
                Thus $s(n)\in c(n)$ by \cref{inter}.
                Thus $s(n)\in U$ by \cref{subseteq}.
            \end{subproof}
            Thus there exists $N\in\naturalsPlus$ such that for all $n\in\naturalsPlus$ if $N\leq n$ then $s(n)\in U$.
        \end{subproof}
        Thus $s$ converges to $x$ in $X$ with respect to $T$ by \cref{converges_to}.
    \end{subproof}

    Thus there exists $s$ such that $s$ is a sequence in $A$ and $s$ converges to $x$ in $X$ with respect to $T$.
\end{proof}

% from §30 Item 3: continuous implies sequential continuity
\begin{theorem}\label{continuous_implies_sequential_first_countable}
    Assume $T$ is a topology on $X$.
    Assume $T_1$ is a topology on $Y$.
    Suppose $f$ is a function from $X$ to $Y$.
    Suppose $f$ is continuous from $X$ to $Y$ with respect to $T$ and $T_1$.
    Then for all $s$ such that $s$ is a sequence in $X$ we have
        for all $x\in X$ if $s$ converges to $x$ in $X$ with respect to $T$ then
        $f\circ s$ converges to $f(x)$ in $Y$ with respect to $T_1$.
\end{theorem}
\begin{proof}
    Follows by \cref{continuous_implies_sequential}.
\end{proof}

% from §30 Item 4: sequential continuity implies continuity in a first-countable space
\begin{theorem}\label{sequential_implies_continuous_first_countable}
    Assume $T$ is a topology on $X$.
    Assume $T_1$ is a topology on $Y$.
    Suppose $f$ is a function from $X$ to $Y$.
    Suppose $X$ satisfies the first countability axiom with respect to $T$.
    Suppose for all $s$ such that $s$ is a sequence in $X$ we have
        for all $x\in X$ if $s$ converges to $x$ in $X$ with respect to $T$ then
        $f\circ s$ converges to $f(x)$ in $Y$ with respect to $T_1$.
    Then $f$ is continuous from $X$ to $Y$ with respect to $T$ and $T_1$.
\end{theorem}
\begin{proof}
    Show that for all $A$ such that $A\subseteq X$ we have
        $\img{f}{\closure{A}{X}{T}}\subseteq \closure{\img{f}{A}}{Y}{T_1}$.
    \begin{subproof}
        Fix $A$.
        Assume $A\subseteq X$.
        Show that for all $y$ such that $y\in\img{f}{\closure{A}{X}{T}}$ we have
            $y\in\closure{\img{f}{A}}{Y}{T_1}$.
        \begin{subproof}
            Fix $y$.
            Assume $y\in\img{f}{\closure{A}{X}{T}}$.
            $f$ is a function from $X$ to $Y$.
            Thus $f\in\funs{X}{Y}$.
            Thus $f$ is a function by \cref{funs_is_function}.
            Take $x$ such that $x\in\closure{A}{X}{T}$ and $x\mathrel{f} y$ by \cref{img_iff}.
            $X$ is closed in $X$ with respect to $T$ by \cref{closed_whole_space}.
            Thus $\closure{A}{X}{T}\subseteq X$ by \cref{closure_subseteq_closed}.
            Thus $x\in X$ by \cref{subseteq}.
            Thus $(x,y)\in f$ by \cref{pair_eq_iff}.
            Thus $f(x) = y$ by \cref{function_apply_iff}.
            Take $s$ such that $s$ is a sequence in $A$ and $s$ converges to $x$ in $X$ with respect to $T$ by
                \cref{closure_implies_sequence_first_countable}.
            Show that $s$ is a sequence in $X$.
            \begin{subproof}
                $s\in\funs{\naturalsPlus}{A}$ by \cref{sequence_in}.
                Thus $s$ is a function to $A$ and $\dom{s} = \naturalsPlus$ by \cref{funs_elim}.
                Thus $s$ is a function to $X$ by \cref{function_on_weaken_codom,subseteq}.
                Thus $s\in\funs{\naturalsPlus}{X}$ by \cref{funs_intro}.
                Thus $s$ is a sequence in $X$ by \cref{sequence_in}.
            \end{subproof}
            Thus $f\circ s$ converges to $f(x)$ in $Y$ with respect to $T_1$ by assumption.
            Thus $f\circ s$ converges to $y$ in $Y$ with respect to $T_1$.
            Show that $f\circ s$ is a sequence in $\img{f}{A}$.
            \begin{subproof}
                $f\in\funs{X}{Y}$.
                Thus $f\circ s\in\funs{\naturalsPlus}{Y}$ by \cref{funs_circ,sequence_in}.
                Thus $f\circ s$ is a function to $Y$ and $\dom{f\circ s} = \naturalsPlus$ by \cref{funs_elim}.
                Show that $f\circ s$ is a function to $\img{f}{A}$.
                \begin{subproof}
                    $f\circ s$ is a function by \cref{funs_is_function}.
                    Show that for all $n\in\dom{f\circ s}$ we have $(f\circ s)(n)\in\img{f}{A}$.
                    \begin{subproof}
                        Fix $n\in\dom{f\circ s}$.
                        $\dom{f\circ s} = \naturalsPlus$.
                        Thus $n\in\naturalsPlus$ by \cref{subseteq}.
                        $s\in\funs{\naturalsPlus}{A}$ by \cref{sequence_in}.
                        Thus $s$ is a function to $A$ and $\dom{s} = \naturalsPlus$ by \cref{funs_elim}.
                        Thus $n\in\dom{s}$ by \cref{subseteq}.
                        Thus $s$ is a function by \cref{funs_is_function}.
                        Thus $s\in\rels{\naturalsPlus}{A}$ by \cref{subseteq,funs_subseteq_rels}.
                        Thus $s\subseteq\naturalsPlus\times A$ by \cref{rels_elim}.
                        Thus $(n,s(n))\in s$ by \cref{function_apply_elim}.
                        Thus $(n,s(n))\in\naturalsPlus\times A$ by \cref{subseteq}.
                        Thus $\snd{(n,s(n))}\in A$ by \cref{times_proj_elim}.
                        $\snd{(n,s(n))} = s(n)$ by \cref{snd_eq}.
                        Thus $s(n)\in A$.
                        Thus $s(n)\in X$ by \cref{subseteq}.
                        Thus $f$ is a function to $Y$ and $\dom{f} = X$ by \cref{funs_elim}.
                        Thus $s(n)\in\dom{f}$.
                        Thus $(s(n),f(s(n)))\in f$ by \cref{function_apply_elim}.
                        Thus $s(n)\mathrel{f} f(s(n))$ by \cref{pair_eq_iff}.
                        Thus $f(s(n))\in\img{f}{A}$ by \cref{img_iff}.
                        $\ran{s}\subseteq A$ by \cref{function_to_ran}.
                        Thus $\ran{s}\subseteq X$ by \cref{subseteq_transitive}.
                        Thus $f$ is composable with $s$.
                        Thus $(f\circ s)(n) = f(s(n))$ by \cref{circ_apply}.
                        Thus $(f\circ s)(n)\in\img{f}{A}$.
                    \end{subproof}
                    Thus $f\circ s$ is a function to $\img{f}{A}$.
                \end{subproof}
                Thus $f\circ s\in\funs{\naturalsPlus}{\img{f}{A}}$ by \cref{funs_intro}.
                Thus $f\circ s$ is a sequence in $\img{f}{A}$ by \cref{sequence_in}.
            \end{subproof}
            $f$ is a function to $Y$ and $\dom{f} = X$ by \cref{funs_elim}.
            Thus $x\in\dom{f}$ by \cref{subseteq}.
            Thus $(x,f(x))\in f$ by \cref{function_apply_elim}.
            Thus $f(x)\in\ran{f}$ by \cref{ran_intro}.
            $\ran{f}\subseteq Y$ by \cref{function_to_ran}.
            Thus $f(x)\in Y$ by \cref{subseteq}.
            $\img{f}{A}\subseteq\ran{f}$ by \cref{img_subseteq_ran}.
            Thus $\img{f}{A}\subseteq Y$ by \cref{subseteq_transitive}.
            Thus $f(x)\in\closure{\img{f}{A}}{Y}{T_1}$ by \cref{sequence_implies_closure}.
            Thus $y\in\closure{\img{f}{A}}{Y}{T_1}$.
        \end{subproof}
        Thus $\img{f}{\closure{A}{X}{T}}\subseteq \closure{\img{f}{A}}{Y}{T_1}$ by \cref{subseteq}.
    \end{subproof}
    Thus for all $B$ such that $B$ is closed in $Y$ with respect to $T_1$
        we have $\preimg{f}{B}$ is closed in $X$ with respect to $T$ by
        \cref{closure_image_subset_implies_preimg_closed}.
    Thus $f$ is continuous from $X$ to $Y$ with respect to $T$ and $T_1$ by
        \cref{preimg_closed_implies_continuous}.
\end{proof}
% from §30 Item 5: second countability axiom
\begin{definition}\label{second_countability_axiom}
    $X$ satisfies the second countability axiom with respect to $T$ iff
        there exists $b$ such that
            $b$ is a sequence in $\pow{X}$ and
            for all $n\in\naturalsPlus$ we have $b(n)$ is open in $X$ with respect to $T$ and
            for all $x\in X$ we have
                for all $U$ such that $U$ is open in $X$ with respect to $T$ and $x\in U$
                    there exists $n\in\naturalsPlus$ such that $x\in b(n)$ and $b(n)\subseteq U$.
\end{definition}

% from §30 Item 6: second countability implies first countability
\begin{lemma}\label{second_countability_implies_first}
    Assume $T$ is a topology on $X$.
    Suppose $X$ satisfies the second countability axiom with respect to $T$.
    Then $X$ satisfies the first countability axiom with respect to $T$.
\end{lemma}
\begin{proof}
    Take $b$ such that
        $b$ is a sequence in $\pow{X}$ and
        for all $n\in\naturalsPlus$ we have $b(n)$ is open in $X$ with respect to $T$ and
        for all $x\in X$ we have
            for all $U$ such that $U$ is open in $X$ with respect to $T$ and $x\in U$
                there exists $n\in\naturalsPlus$ such that $x\in b(n)$ and $b(n)\subseteq U$
        by \cref{second_countability_axiom}.
    Show that for all $x\in X$ there exists $c$ such that
        $c$ is a countable basis at $x$ in $X$ with respect to $T$.
    \begin{subproof}
        Fix $x\in X$.
        Let $R = \{(n,y) \mid n\in\naturalsPlus \mid
            (x\in b(n)\land y = b(n))\lor(x\notin b(n)\land y = X)\}$.
        Show that $R$ is a relation.
        \begin{subproof}
            Show that for all $w$ such that $w\in R$ there exist $n,y$ such that $w = (n,y)$.
            \begin{subproof}
                Fix $w$.
                Assume $w\in R$.
                Thus there exist $n,y$ such that $w = (n,y)$ by \cref{pair_eq_iff}.
            \end{subproof}
            Thus $R$ is a relation by \cref{relation}.
        \end{subproof}
        Show that for all $n\in\naturalsPlus$ there exists $y$ such that $n\mathrel{R} y$.
        \begin{subproof}
            Fix $n\in\naturalsPlus$.
            \begin{byCase}
                \caseOf{$x\in b(n)$.}
                    Let $y = b(n)$.
                    Thus $(n,y)\in R$ by \cref{pair_eq_iff}.
                    Thus $n\mathrel{R} y$ by \cref{pair_eq_iff}.
                \caseOf{$x\notin b(n)$.}
                    Let $y = X$.
                    Thus $(n,y)\in R$ by \cref{pair_eq_iff}.
                    Thus $n\mathrel{R} y$ by \cref{pair_eq_iff}.
            \end{byCase}
        \end{subproof}
        Take $c$ such that $c$ is a function on $\naturalsPlus$ and for all $n\in\naturalsPlus$ we have
            $n\mathrel{R} c(n)$ by \cref{axiom_of_choice}.
        Show that $c$ is a countable basis at $x$ in $X$ with respect to $T$.
        \begin{subproof}
            Show that $c$ is a sequence in $\pow{X}$.
            \begin{subproof}
                Show that $c$ is a function to $\pow{X}$.
                \begin{subproof}
                    Show that for all $n\in\dom{c}$ we have $c(n)\in\pow{X}$.
                    \begin{subproof}
                        Fix $n\in\dom{c}$.
                        $\dom{c} = \naturalsPlus$.
                        Thus $n\in\naturalsPlus$ by \cref{subseteq}.
                        Thus $n\mathrel{R} c(n)$.
                        Thus $(n,c(n))\in R$ by \cref{pair_eq_iff}.
                        Thus $(x\in b(n)\land c(n) = b(n))\lor(x\notin b(n)\land c(n) = X)$ by \cref{pair_eq_iff}.
                        \begin{byCase}
                            \caseOf{$x\in b(n)\land c(n) = b(n)$.}
                                Show that $b(n)\in\pow{X}$.
                                \begin{subproof}
                                    $b\in\funs{\naturalsPlus}{\pow{X}}$ by \cref{sequence_in}.
                                    Thus $b$ is a function by \cref{funs_is_function}.
                                    Thus $b\in\rels{\naturalsPlus}{\pow{X}}$ by \cref{subseteq,funs_subseteq_rels}.
                                    Thus $b\subseteq\naturalsPlus\times\pow{X}$ by \cref{rels_elim}.
                                    $\dom{b} = \naturalsPlus$ by \cref{funs_elim}.
                                    Thus $n\in\dom{b}$ by \cref{subseteq}.
                                    Thus $(n,b(n))\in b$ by \cref{function_apply_elim}.
                                    Thus $(n,b(n))\in\naturalsPlus\times\pow{X}$ by \cref{subseteq}.
                                    Thus $\snd{(n,b(n))}\in\pow{X}$ by \cref{times_proj_elim}.
                                    $\snd{(n,b(n))} = b(n)$ by \cref{snd_eq}.
                                    Thus $b(n)\in\pow{X}$.
                                \end{subproof}
                                Thus $c(n)\in\pow{X}$.
                            \caseOf{$x\notin b(n)\land c(n) = X$.}
                                $X\subseteq X$ by \cref{subseteq_refl}.
                                Thus $X\in\pow{X}$ by \cref{powerset_intro}.
                                Thus $c(n)\in\pow{X}$.
                        \end{byCase}
                    \end{subproof}
                    Thus $c$ is a function to $\pow{X}$.
                \end{subproof}
                Thus $c\in\funs{\naturalsPlus}{\pow{X}}$ by \cref{funs_intro}.
                Thus $c$ is a sequence in $\pow{X}$ by \cref{sequence_in}.
            \end{subproof}
            Show that for all $n\in\naturalsPlus$ we have $c(n)$ is a neighborhood of $x$ in $X$ with respect to $T$.
            \begin{subproof}
                Fix $n\in\naturalsPlus$.
                Thus $n\mathrel{R} c(n)$.
                Thus $(n,c(n))\in R$ by \cref{pair_eq_iff}.
                Thus $(x\in b(n)\land c(n) = b(n))\lor(x\notin b(n)\land c(n) = X)$ by \cref{pair_eq_iff}.
                \begin{byCase}
                    \caseOf{$x\in b(n)\land c(n) = b(n)$.}
                        Thus $b(n)$ is open in $X$ with respect to $T$.
                        Thus $c(n)$ is open in $X$ with respect to $T$.
                        Thus $x\in c(n)$.
                        Thus $c(n)$ is a neighborhood of $x$ in $X$ with respect to $T$ by \cref{neighborhood}.
                    \caseOf{$x\notin b(n)\land c(n) = X$.}
                        $X\in T$ by \cref{topology_on}.
                        Thus $X$ is open in $X$ with respect to $T$ by \cref{open_in}.
                        Thus $x\in X$.
                        Thus $c(n)$ is a neighborhood of $x$ in $X$ with respect to $T$ by \cref{neighborhood}.
                \end{byCase}
            \end{subproof}
            Show that for all $U$ such that $U$ is a neighborhood of $x$ in $X$ with respect to $T$
                there exists $n\in\naturalsPlus$ such that $c(n)\subseteq U$.
            \begin{subproof}
                Fix $U$.
                Assume $U$ is a neighborhood of $x$ in $X$ with respect to $T$.
                Thus $U$ is open in $X$ with respect to $T$ and $x\in U$ by \cref{neighborhood}.
                Thus there exists $n\in\naturalsPlus$ such that $x\in b(n)$ and $b(n)\subseteq U$.
                Take $n\in\naturalsPlus$ such that $x\in b(n)$ and $b(n)\subseteq U$.
                Thus $n\mathrel{R} c(n)$.
                Thus $(n,c(n))\in R$ by \cref{pair_eq_iff}.
                Thus $x\in b(n)\land c(n) = b(n)$ by \cref{pair_eq_iff}.
                Thus $c(n)\subseteq U$ by \cref{subseteq}.
            \end{subproof}
            Thus $c$ is a countable basis at $x$ in $X$ with respect to $T$ by \cref{countable_basis_at_point}.
        \end{subproof}
    \end{subproof}
    Thus $X$ satisfies the first countability axiom with respect to $T$ by \cref{first_countability_axiom}.
\end{proof}
% from §30 Item 7: subspace of a first-countable space is first-countable
\begin{lemma}\label{subspace_first_countable}
    Assume $T$ is a topology on $X$.
    Suppose $A\subseteq X$.
    Suppose $X$ satisfies the first countability axiom with respect to $T$.
    Then $A$ satisfies the first countability axiom with respect to $\subspaceTopology{T}{A}$.
\end{lemma}
\begin{proof}
    Let $T_1 = \subspaceTopology{T}{A}$.
    $T_1$ is a topology on $A$ by \cref{subspace_topology_is_topology}.
    Show that for all $x\in A$ there exists $c$ such that
        $c$ is a countable basis at $x$ in $A$ with respect to $T_1$.
    \begin{subproof}
        Fix $x\in A$.
        Thus $x\in X$ by \cref{subseteq}.
        Take $b$ such that $b$ is a countable basis at $x$ in $X$ with respect to $T$
            by \cref{first_countability_axiom}.
        Let $c = \{(n, A\inter b(n)) \mid n\in\naturalsPlus\}$.
        Show that $c$ is a countable basis at $x$ in $A$ with respect to $T_1$.
        \begin{subproof}
            Show that $c$ is a sequence in $\pow{A}$.
            \begin{subproof}
                Show that $c$ is a function to $\pow{A}$.
                \begin{subproof}
                    Show that $c$ is a function.
                    \begin{subproof}
                        Show that for all $n,U_1,U_2$ such that $(n,U_1)\in c$ and $(n,U_2)\in c$
                            we have $U_1 = U_2$.
                        \begin{subproof}
                            Fix $n$.
                            Fix $U_1$.
                            Fix $U_2$.
                            Assume $(n,U_1)\in c$ and $(n,U_2)\in c$.
                            Thus $U_1 = A\inter b(n)$ and $U_2 = A\inter b(n)$ by \cref{pair_eq_iff}.
                            Thus $U_1 = U_2$.
                        \end{subproof}
                        Thus $c$ is right-unique by \cref{rightunique}.
                        Thus $c$ is a function.
                    \end{subproof}
                    Show that for all $n\in\dom{c}$ we have $c(n)\in\pow{A}$.
                    \begin{subproof}
                        Fix $n\in\dom{c}$.
                        Take $U$ such that $(n,U)\in c$ by \cref{dom_iff}.
                        Thus $U = A\inter b(n)$ by \cref{pair_eq_iff}.
                        Show that $A\inter b(n)\subseteq A$.
                        \begin{subproof}
                            Show that for all $z$ such that $z\in A\inter b(n)$ we have $z\in A$.
                            \begin{subproof}
                                Fix $z$.
                                Assume $z\in A\inter b(n)$.
                                Thus $z\in A$ by \cref{inter_elim_left}.
                            \end{subproof}
                            Thus $A\inter b(n)\subseteq A$ by \cref{subseteq}.
                        \end{subproof}
                        Thus $A\inter b(n)\in\pow{A}$ by \cref{powerset_intro}.
                        Thus $c(n)\in\pow{A}$ by \cref{function_apply_iff}.
                    \end{subproof}
                    Thus $c$ is a function to $\pow{A}$.
                \end{subproof}
                Show that $\dom{c} = \naturalsPlus$.
                \begin{subproof}
                    Show that $\naturalsPlus\subseteq\dom{c}$.
                    \begin{subproof}
                        Show that for all $n$ such that $n\in\naturalsPlus$ we have $n\in\dom{c}$.
                        \begin{subproof}
                            Fix $n$.
                            Assume $n\in\naturalsPlus$.
                            $(n, A\inter b(n))\in c$ by \cref{pair_eq_iff}.
                            Thus $n\in\dom{c}$ by \cref{dom_intro}.
                        \end{subproof}
                        Thus $\naturalsPlus\subseteq\dom{c}$ by \cref{subseteq}.
                    \end{subproof}
                    Show that $\dom{c}\subseteq\naturalsPlus$.
                    \begin{subproof}
                        Show that for all $n$ such that $n\in\dom{c}$ we have $n\in\naturalsPlus$.
                        \begin{subproof}
                            Fix $n$.
                            Assume $n\in\dom{c}$.
                            Take $U$ such that $(n,U)\in c$ by \cref{dom_iff}.
                            Thus $n\in\naturalsPlus$ by \cref{pair_eq_iff}.
                        \end{subproof}
                        Thus $\dom{c}\subseteq\naturalsPlus$ by \cref{subseteq}.
                    \end{subproof}
                    Thus $\dom{c} = \naturalsPlus$ by \cref{subseteq_antisymmetric}.
                \end{subproof}
                Thus $c\in\funs{\naturalsPlus}{\pow{A}}$ by \cref{funs_intro}.
                Thus $c$ is a sequence in $\pow{A}$ by \cref{sequence_in}.
            \end{subproof}
            Show that for all $n\in\naturalsPlus$ we have $c(n)$ is a neighborhood of $x$ in $A$ with respect to $T_1$.
            \begin{subproof}
                Fix $n\in\naturalsPlus$.
                $b(n)$ is a neighborhood of $x$ in $X$ with respect to $T$ by \cref{countable_basis_at_point}.
                Thus $b(n)$ is open in $X$ with respect to $T$ and $x\in b(n)$ by \cref{neighborhood}.
                Thus $b(n)\in T$ by \cref{open_in}.
                Thus $A\inter b(n)\in T_1$ by \cref{subspace_topology}.
                Thus $A\inter b(n)$ is open in $A$ with respect to $T_1$ by \cref{open_in}.
                $c\in\funs{\naturalsPlus}{\pow{A}}$ by \cref{sequence_in}.
                Thus $c$ is a function by \cref{funs_is_function}.
                Thus $\dom{c} = \naturalsPlus$ by \cref{funs_elim}.
                Thus $n\in\dom{c}$ by \cref{subseteq}.
                $(n, A\inter b(n))\in c$ by \cref{pair_eq_iff}.
                Thus $c(n) = A\inter b(n)$ by \cref{function_apply_iff}.
                $x\in A\inter b(n)$ by \cref{inter_intro}.
                Thus $x\in c(n)$.
                Thus $c(n)$ is a neighborhood of $x$ in $A$ with respect to $T_1$ by \cref{neighborhood}.
            \end{subproof}
            Show that for all $U$ such that $U$ is a neighborhood of $x$ in $A$ with respect to $T_1$
                there exists $n\in\naturalsPlus$ such that $c(n)\subseteq U$.
            \begin{subproof}
                Fix $U$.
                Assume $U$ is a neighborhood of $x$ in $A$ with respect to $T_1$.
                Thus $U$ is open in $A$ with respect to $T_1$ and $x\in U$ by \cref{neighborhood}.
                Thus $U\in T_1$ by \cref{open_in}.
                Take $V\in T$ such that $U = A\inter V$ by \cref{subspace_topology}.
                Thus $x\in A\inter V$.
                Thus $x\in V$ by \cref{inter_elim_right}.
                $V$ is open in $X$ with respect to $T$ by \cref{open_in}.
                Thus $V$ is a neighborhood of $x$ in $X$ with respect to $T$ by \cref{neighborhood}.
                $1\in\naturalsPlus$ by \cref{real_one_in_naturals}.
                Thus $b(1)$ is a neighborhood of $x$ in $X$ with respect to $T$ and
                    for all $U$ such that $U$ is a neighborhood of $x$ in $X$ with respect to $T$
                        there exists $n\in\naturalsPlus$ such that $b(n)\subseteq U$
                    by \cref{countable_basis_at_point}.
                Thus for all $U$ such that $U$ is a neighborhood of $x$ in $X$ with respect to $T$
                    there exists $n\in\naturalsPlus$ such that $b(n)\subseteq U$.
                Thus there exists $n\in\naturalsPlus$ such that $b(n)\subseteq V$.
                Take $n\in\naturalsPlus$ such that $b(n)\subseteq V$.
                Show that $A\inter b(n)\subseteq A\inter V$.
                \begin{subproof}
                    Show that for all $z$ such that $z\in A\inter b(n)$ we have $z\in A\inter V$.
                    \begin{subproof}
                        Fix $z$.
                        Assume $z\in A\inter b(n)$.
                        Thus $z\in A$ by \cref{inter_elim_left}.
                        Thus $z\in b(n)$ by \cref{inter_elim_right}.
                        Thus $z\in V$ by \cref{subseteq}.
                        Thus $z\in A\inter V$ by \cref{inter_intro}.
                    \end{subproof}
                    Thus $A\inter b(n)\subseteq A\inter V$ by \cref{subseteq}.
                \end{subproof}
                $c\in\funs{\naturalsPlus}{\pow{A}}$ by \cref{sequence_in}.
                Thus $c$ is a function by \cref{funs_is_function}.
                Thus $\dom{c} = \naturalsPlus$ by \cref{funs_elim}.
                Thus $n\in\dom{c}$ by \cref{subseteq}.
                $(n, A\inter b(n))\in c$ by \cref{pair_eq_iff}.
                Thus $c(n) = A\inter b(n)$ by \cref{function_apply_iff}.
                Thus $c(n)\subseteq A\inter V$ by \cref{subseteq}.
                Thus $c(n)\subseteq U$ by \cref{subseteq}.
            \end{subproof}
            Thus $c$ is a countable basis at $x$ in $A$ with respect to $T_1$ by \cref{countable_basis_at_point}.
        \end{subproof}
    \end{subproof}
    Thus $A$ satisfies the first countability axiom with respect to $T_1$ by \cref{first_countability_axiom}.
\end{proof}
% from §30 helper lemma 2: basis elements are neighborhoods
\begin{lemma}\label{countable_basis_elem_neighborhood}
    Assume $b$ is a countable basis at $x$ in $X$ with respect to $T$.
    Assume $n\in\naturalsPlus$.
    Then $b(n)$ is a neighborhood of $x$ in $X$ with respect to $T$.
\end{lemma}
\begin{proof}
    Thus for all $k\in\naturalsPlus$ we have $b(k)$ is a neighborhood of $x$ in $X$ with respect to $T$
        by \cref{countable_basis_at_point}.
    Thus $b(n)$ is a neighborhood of $x$ in $X$ with respect to $T$.
\end{proof}
% from §30 Item 8: countable product of first-countable spaces is first-countable
\begin{lemma}\label{countable_product_first_countable}
    Assume $X$ is a function.
    Assume $T$ is a function.
    Assume $\dom{X} = \naturalsPlus$.
    Assume $\dom{T} = \naturalsPlus$.
    Assume that for all $n\in\naturalsPlus$ we have $\apply{T}{n}$ is a topology on $\apply{X}{n}$.
    Suppose for all $n\in\naturalsPlus$ we have $\apply{X}{n}$ satisfies the first countability axiom with respect to $\apply{T}{n}$.
    Then $\productFamily{X}$ satisfies the first countability axiom with respect to $\productTopologyIndexed{T}{X}$.
\end{lemma}
\begin{proof}
    Let $T_0 = \productTopologyIndexed{T}{X}$.
    Show that for all $x\in\productFamily{X}$ there exists $d$ such that
        $d$ is a countable basis at $x$ in $\productFamily{X}$ with respect to $T_0$.
    \begin{subproof}
        Fix $x\in\productFamily{X}$.
        Let $R = \{(n,b) \mid n\in\naturalsPlus \mid
            \text{$b$ is a countable basis at $\apply{x}{n}$ in $\apply{X}{n}$ with respect to $\apply{T}{n}$}\}$.
        Show that $R$ is a relation.
        \begin{subproof}
            Show that for all $w$ such that $w\in R$ there exist $n,b$ such that $w = (n,b)$.
            \begin{subproof}
                Fix $w$.
                Assume $w\in R$.
                Thus there exist $n,b$ such that $w = (n,b)$ by \cref{pair_eq_iff}.
            \end{subproof}
            Thus $R$ is a relation by \cref{relation}.
        \end{subproof}
        Show that for all $n\in\naturalsPlus$ there exists $b$ such that $n\mathrel{R} b$.
        \begin{subproof}
            Fix $n\in\naturalsPlus$.
            $x\in\productFamily{X}$.
            Thus $x\in\funs{\dom{X}}{\unions{\ran{X}}}$ by \cref{indexed_product}.
            Thus $x$ is a function and $\dom{x} = \dom{X}$ by \cref{funs_elim}.
            Thus $n\in\dom{x}$ by \cref{subseteq}.
            Thus $\apply{x}{n}\in\apply{X}{n}$ by \cref{indexed_product}.
            $\apply{X}{n}$ satisfies the first countability axiom with respect to $\apply{T}{n}$.
            Thus there exists $b$ such that $b$ is a countable basis at $\apply{x}{n}$ in $\apply{X}{n}$ with respect to $\apply{T}{n}$.
            Take $b$ such that $b$ is a countable basis at $\apply{x}{n}$ in $\apply{X}{n}$ with respect to $\apply{T}{n}$.
            Thus $(n,b)\in R$ by \cref{pair_eq_iff}.
            Thus $n\mathrel{R} b$ by \cref{pair_eq_iff}.
        \end{subproof}
        Take $b$ such that $b$ is a function on $\naturalsPlus$ and for all $n\in\naturalsPlus$ we have
            $n\mathrel{R}\apply{b}{n}$ by \cref{choice_function}.
        Show that for all $n\in\naturalsPlus$ we have
            $\apply{b}{n}$ is a countable basis at $\apply{x}{n}$ in $\apply{X}{n}$ with respect to $\apply{T}{n}$.
        \begin{subproof}
            Fix $n\in\naturalsPlus$.
            Thus $n\mathrel{R}\apply{b}{n}$ by \cref{choice_function}.
            Thus $\apply{b}{n}$ is a countable basis at $\apply{x}{n}$ in $\apply{X}{n}$ with respect to $\apply{T}{n}$
                by \cref{pair_eq_iff}.
        \end{subproof}

        Let $R_1 = \{(n,c) \mid n\in\naturalsPlus \mid
            \text{$c$ is a countable basis at $\apply{x}{n}$ in $\apply{X}{n}$ with respect to $\apply{T}{n}$ and
            for all $m,k\in\naturalsPlus$ if $m\leq k$ then $c(k)\subseteq c(m)$}\}$.
        Show that $R_1$ is a relation.
        \begin{subproof}
            Show that for all $w$ such that $w\in R_1$ there exist $n,c$ such that $w = (n,c)$.
            \begin{subproof}
                Fix $w$.
                Assume $w\in R_1$.
                Thus there exist $n,c$ such that $w = (n,c)$ by \cref{pair_eq_iff}.
            \end{subproof}
            Thus $R_1$ is a relation by \cref{relation}.
        \end{subproof}
        Show that for all $n\in\naturalsPlus$ there exists $c$ such that $n\mathrel{R_1} c$.
        \begin{subproof}
            Fix $n\in\naturalsPlus$.
            $\apply{b}{n}$ is a countable basis at $\apply{x}{n}$ in $\apply{X}{n}$ with respect to $\apply{T}{n}$.
            Thus there exists $c$ such that
                $c$ is a countable basis at $\apply{x}{n}$ in $\apply{X}{n}$ with respect to $\apply{T}{n}$ and
                for all $m,k\in\naturalsPlus$ if $m\leq k$ then $c(k)\subseteq c(m)$ by \cref{nested_countable_basis}.
            Take $c$ such that
                $c$ is a countable basis at $\apply{x}{n}$ in $\apply{X}{n}$ with respect to $\apply{T}{n}$ and
                for all $m,k\in\naturalsPlus$ if $m\leq k$ then $c(k)\subseteq c(m)$.
            Thus $(n,c)\in R_1$ by \cref{pair_eq_iff}.
            Thus $n\mathrel{R_1} c$ by \cref{pair_eq_iff}.
        \end{subproof}
        Take $c$ such that $c$ is a function on $\naturalsPlus$ and for all $n\in\naturalsPlus$ we have
            $n\mathrel{R_1}c(n)$ by \cref{choice_function}.
        Show that for all $n\in\naturalsPlus$ we have
            $c(n)$ is a countable basis at $\apply{x}{n}$ in $\apply{X}{n}$ with respect to $\apply{T}{n}$.
        \begin{subproof}
            Fix $n\in\naturalsPlus$.
            Thus $n\mathrel{R_1}\apply{c}{n}$ by \cref{choice_function}.
            Thus $\apply{c}{n}$ is a countable basis at $\apply{x}{n}$ in $\apply{X}{n}$ with respect to $\apply{T}{n}$
                by \cref{pair_eq_iff}.
        \end{subproof}
        Show that for all $n\in\naturalsPlus$ we have for all $m,k\in\naturalsPlus$ if $m\leq k$ then
            $\apply{c(n)}{k}\subseteq \apply{c(n)}{m}$.
        \begin{subproof}
            Fix $n\in\naturalsPlus$.
            Thus $n\mathrel{R_1}\apply{c}{n}$ by \cref{choice_function}.
            Thus for all $m,k\in\naturalsPlus$ if $m\leq k$ then
                $\apply{c(n)}{k}\subseteq \apply{c(n)}{m}$ by \cref{pair_eq_iff}.
        \end{subproof}

        Let $R_d = \{(n,g) \mid n\in\naturalsPlus \mid
            \text{$g$ is a function on $\initialSegment{n}$ and for all $i\in\initialSegment{n}$ we have
            $g(i) = \preimg{\piIndex{X}{i}}{\apply{c(i)}{n}}$}\}$.
        Show that $R_d$ is a relation.
        \begin{subproof}
            Show that for all $w$ such that $w\in R_d$ there exist $n,g$ such that $w = (n,g)$.
            \begin{subproof}
                Fix $w$.
                Assume $w\in R_d$.
                Thus there exist $n,g$ such that $w = (n,g)$ by \cref{pair_eq_iff}.
            \end{subproof}
            Thus $R_d$ is a relation by \cref{relation}.
        \end{subproof}
        Show that for all $n\in\naturalsPlus$ there exists $g$ such that $n\mathrel{R_d} g$.
        \begin{subproof}
            Fix $n\in\naturalsPlus$.
            Let $g = \{(i, \preimg{\piIndex{X}{i}}{\apply{c(i)}{n}}) \mid i\in\initialSegment{n}\}$.
            Show that $g$ is a function on $\initialSegment{n}$.
            \begin{subproof}
                Show that $g$ is a function.
                \begin{subproof}
                    Show that for all $i,A_1,A_2$ such that $(i,A_1)\in g$ and $(i,A_2)\in g$ we have $A_1 = A_2$.
                    \begin{subproof}
                        Fix $i$.
                        Fix $A_1$.
                        Fix $A_2$.
                        Assume $(i,A_1)\in g$ and $(i,A_2)\in g$.
                        Thus $A_1 = \preimg{\piIndex{X}{i}}{\apply{c(i)}{n}}$ and
                            $A_2 = \preimg{\piIndex{X}{i}}{\apply{c(i)}{n}}$ by \cref{pair_eq_iff}.
                        Thus $A_1 = A_2$.
                    \end{subproof}
                    Thus $g$ is right-unique by \cref{rightunique}.
                    Show that $g$ is a relation.
                    \begin{subproof}
                        Show that for all $w$ such that $w\in g$ there exist $i,A$ such that $w = (i,A)$.
                        \begin{subproof}
                            Fix $w$.
                            Assume $w\in g$.
                            Thus there exist $i,A$ such that $w = (i,A)$ by \cref{pair_eq_iff}.
                        \end{subproof}
                        Thus $g$ is a relation by \cref{relation}.
                    \end{subproof}
                    Thus $g$ is a function.
                \end{subproof}
                Show that $\dom{g} = \initialSegment{n}$.
                \begin{subproof}
                    Show that $\initialSegment{n}\subseteq\dom{g}$.
                    \begin{subproof}
                        Show that for all $i$ such that $i\in\initialSegment{n}$ we have $i\in\dom{g}$.
                        \begin{subproof}
                            Fix $i$.
                            Assume $i\in\initialSegment{n}$.
                            $(i, \preimg{\piIndex{X}{i}}{\apply{c(i)}{n}})\in g$ by \cref{pair_eq_iff}.
                            Thus $i\in\dom{g}$ by \cref{dom_intro}.
                        \end{subproof}
                        Thus $\initialSegment{n}\subseteq\dom{g}$ by \cref{subseteq}.
                    \end{subproof}
                    Show that $\dom{g}\subseteq\initialSegment{n}$.
                    \begin{subproof}
                        Show that for all $i$ such that $i\in\dom{g}$ we have $i\in\initialSegment{n}$.
                        \begin{subproof}
                            Fix $i$.
                            Assume $i\in\dom{g}$.
                            Take $A$ such that $(i,A)\in g$ by \cref{dom_iff}.
                            Thus $i\in\initialSegment{n}$ by \cref{pair_eq_iff}.
                        \end{subproof}
                        Thus $\dom{g}\subseteq\initialSegment{n}$ by \cref{subseteq}.
                    \end{subproof}
                    Thus $\dom{g} = \initialSegment{n}$ by \cref{subseteq_antisymmetric}.
                \end{subproof}
                Thus $g$ is a function and $\dom{g} = \initialSegment{n}$.
                Thus $g$ is a function on $\initialSegment{n}$.
            \end{subproof}
            Show that for all $i\in\initialSegment{n}$ we have
                $\apply{g}{i} = \preimg{\piIndex{X}{i}}{\apply{c(i)}{n}}$.
            \begin{subproof}
                Fix $i\in\initialSegment{n}$.
                $(i, \preimg{\piIndex{X}{i}}{\apply{c(i)}{n}})\in g$ by \cref{pair_eq_iff}.
                Thus $i\in\dom{g}$ by \cref{dom_intro}.
                Thus $\apply{g}{i} = \preimg{\piIndex{X}{i}}{\apply{c(i)}{n}}$ by \cref{function_apply_iff}.
            \end{subproof}
            Thus $(n,g)\in R_d$ by \cref{pair_eq_iff}.
            Thus $n\mathrel{R_d} g$ by \cref{pair_eq_iff}.
        \end{subproof}
        Take $h$ such that $h$ is a function on $\naturalsPlus$ and for all $n\in\naturalsPlus$ we have
            $n\mathrel{R_d}\apply{h}{n}$ by \cref{choice_function}.
        Let $d = \{(n, \inters{\img{\apply{h}{n}}{\initialSegment{n}}}) \mid n\in\naturalsPlus\}$.

        Show that $d$ is a countable basis at $x$ in $\productFamily{X}$ with respect to $T_0$.
        \begin{subproof}
            Show that $d$ is a sequence in $\pow{\productFamily{X}}$.
            \begin{subproof}
                Show that $d$ is a function to $\pow{\productFamily{X}}$.
                \begin{subproof}
                    Show that $d$ is a function.
                    \begin{subproof}
                        Show that for all $n,U_1,U_2$ such that $(n,U_1)\in d$ and $(n,U_2)\in d$ we have $U_1 = U_2$.
                        \begin{subproof}
                            Fix $n$.
                            Fix $U_1$.
                            Fix $U_2$.
                            Assume $(n,U_1)\in d$ and $(n,U_2)\in d$.
                            Thus $U_1 = \inters{\img{\apply{h}{n}}{\initialSegment{n}}}$ and
                                $U_2 = \inters{\img{\apply{h}{n}}{\initialSegment{n}}}$
                                by \cref{pair_eq_iff}.
                            Thus $U_1 = U_2$.
                        \end{subproof}
                        Thus $d$ is right-unique by \cref{rightunique}.
                        Thus $d$ is a function.
                    \end{subproof}
                    Show that for all $n\in\dom{d}$ we have $d(n)\in\pow{\productFamily{X}}$.
                    \begin{subproof}
                        Fix $n\in\dom{d}$.
                        Take $U$ such that $(n,U)\in d$ by \cref{dom_iff}.
                        Thus $U = \inters{\img{\apply{h}{n}}{\initialSegment{n}}}$ by \cref{pair_eq_iff}.
                        Thus $n\in\naturalsPlus$ by \cref{pair_eq_iff}.
                        Let $F = \img{\apply{h}{n}}{\initialSegment{n}}$.
                        Show that $F\subseteq \productSubbasis{T}{X}$.
                        \begin{subproof}
                            Show that for all $A$ such that $A\in F$ we have $A\in\productSubbasis{T}{X}$.
                            \begin{subproof}
                                Fix $A$.
                                Assume $A\in F$.
                                Take $i$ such that $i\in\initialSegment{n}$ and $(i,A)\in\apply{h}{n}$ by \cref{img_iff}.
                                Thus $n\mathrel{R_d}\apply{h}{n}$ by \cref{choice_function}.
                                Thus $\apply{h}{n}$ is a function and $\dom{\apply{h}{n}} = \initialSegment{n}$ by \cref{pair_eq_iff}.
                                Thus $i\in\dom{\apply{h}{n}}$ by \cref{dom_intro}.
                                Thus $\apply{\apply{h}{n}}{i} = A$ by \cref{function_apply_iff}.
                                Thus for all $j\in\initialSegment{n}$ we have
                                    $\apply{\apply{h}{n}}{j} = \preimg{\piIndex{X}{j}}{\apply{c(j)}{n}}$ by \cref{pair_eq_iff}.
                                Thus $A = \preimg{\piIndex{X}{i}}{\apply{c(i)}{n}}$.
                                Thus $i\in\naturalsPlus$ by \cref{initial_segment_naturalsplus}.
                                $\dom{X} = \naturalsPlus$.
                                Thus $i\in\dom{X}$ by \cref{subseteq}.
                                $c(i)$ is a countable basis at $\apply{x}{i}$ in $\apply{X}{i}$ with respect to $\apply{T}{i}$.
                                Thus $c(i)$ is a sequence in $\pow{\apply{X}{i}}$ by \cref{countable_basis_at_point}.
                                Thus $c(i)\in\funs{\naturalsPlus}{\pow{\apply{X}{i}}}$ by \cref{sequence_in}.
                                Thus $c(i)$ is a function to $\pow{\apply{X}{i}}$ and $\dom{c(i)} = \naturalsPlus$ by \cref{funs_elim}.
                                Thus for all $z\in\dom{c(i)}$ we have $\apply{c(i)}{z}\in\pow{\apply{X}{i}}$.
                                Thus $n\in\dom{c(i)}$ by \cref{subseteq}.
                                Thus $\apply{c(i)}{n}\in\pow{\apply{X}{i}}$.
                                $\apply{c(i)}{n}$ is a neighborhood of $\apply{x}{i}$ in $\apply{X}{i}$ with respect to $\apply{T}{i}$
                                    by \cref{countable_basis_at_point}.
                                Thus $\apply{c(i)}{n}$ is open in $\apply{X}{i}$ with respect to $\apply{T}{i}$ by \cref{neighborhood}.
                                Thus $\apply{c(i)}{n}\in\apply{T}{i}$ by \cref{open_in}.
                                Thus $A\subseteq\dom{\piIndex{X}{i}}$ by \cref{preimg_subseteq_dom}.
                                Show that $\dom{\piIndex{X}{i}}\subseteq\productFamily{X}$.
                                \begin{subproof}
                                    Show that for all $z$ such that $z\in\dom{\piIndex{X}{i}}$ we have $z\in\productFamily{X}$.
                                    \begin{subproof}
                                        Fix $z$.
                                        Assume $z\in\dom{\piIndex{X}{i}}$.
                                        Take $y$ such that $(z,y)\in\piIndex{X}{i}$ by \cref{dom_iff}.
                                        Take $x\in\productFamily{X}$ such that $(z,y) = (x,\apply{x}{i})$ by \cref{projection_map_indexed}.
                                        Thus $z = x$ by \cref{pair_eq_iff}.
                                        Thus $z\in\productFamily{X}$.
                                    \end{subproof}
                                    Thus $\dom{\piIndex{X}{i}}\subseteq\productFamily{X}$ by \cref{subseteq}.
                                \end{subproof}
                                Thus $A\subseteq\productFamily{X}$ by \cref{subseteq_transitive}.
                                Thus $A\in\pow{\productFamily{X}}$ by \cref{powerset_intro}.
                                Thus $A\in\productSubbasisAt{T}{X}{i}$ by \cref{product_subbasis_indexed}.
                                Thus $A\in\productSubbasis{T}{X}$ by \cref{product_subbasis_union_indexed}.
                            \end{subproof}
                            Thus $F\subseteq\productSubbasis{T}{X}$ by \cref{subseteq}.
                        \end{subproof}
                        Show that $F$ is finite.
                        \begin{subproof}
                            Thus $n\mathrel{R_d}\apply{h}{n}$ by \cref{choice_function}.
                            Thus $\apply{h}{n}$ is a function on $\initialSegment{n}$ by \cref{pair_eq_iff}.
                            $\initialSegment{n}$ is finite by \cref{initial_segment_finite}.
                            Thus $F$ is finite by \cref{finite_image_function}.
                        \end{subproof}
                        Show that $F$ is inhabited.
                        \begin{subproof}
                            $n\in\initialSegment{n}$ by \cref{initial_segment_naturalsplus}.
                            Thus $n\mathrel{R_d}\apply{h}{n}$ by \cref{choice_function}.
                            Thus $\apply{h}{n}$ is a function and $\dom{\apply{h}{n}} = \initialSegment{n}$ by \cref{pair_eq_iff}.
                            Thus $n\in\dom{\apply{h}{n}}$ by \cref{subseteq}.
                            Thus for all $i\in\initialSegment{n}$ we have
                                $\apply{\apply{h}{n}}{i} = \preimg{\piIndex{X}{i}}{\apply{c(i)}{n}}$ by \cref{pair_eq_iff}.
                            Thus $\apply{\apply{h}{n}}{n} = \preimg{\piIndex{X}{n}}{\apply{c(n)}{n}}$.
                            Thus $(n, \preimg{\piIndex{X}{n}}{\apply{c(n)}{n}})\in\apply{h}{n}$ by \cref{function_apply_iff}.
                            Thus $\preimg{\piIndex{X}{n}}{\apply{c(n)}{n}}\in F$ by \cref{img_iff}.
                        \end{subproof}
                        Take $A$ such that $A\in F$.
                        Thus $A\in\productSubbasis{T}{X}$ by \cref{subseteq}.
                        Show that $A\subseteq\unions{\productSubbasis{T}{X}}$.
                        \begin{subproof}
                            Show that for all $z$ such that $z\in A$ we have $z\in\unions{\productSubbasis{T}{X}}$.
                            \begin{subproof}
                                Fix $z$.
                                Assume $z\in A$.
                                Thus $z\in\unions{\productSubbasis{T}{X}}$ by \cref{unions_intro}.
                            \end{subproof}
                            Thus $A\subseteq\unions{\productSubbasis{T}{X}}$ by \cref{subseteq}.
                        \end{subproof}
                        Thus $U\subseteq A$ by \cref{inters_subseteq_elem}.
                        Thus $U\subseteq\unions{\productSubbasis{T}{X}}$ by \cref{subseteq_transitive}.
                        Thus $U\in\pow{\unions{\productSubbasis{T}{X}}}$ by \cref{powerset_intro}.
                        Thus $U\in\finiteIntersections{\productSubbasis{T}{X}}$ by \cref{finite_intersections}.
                        Thus $U\subseteq\unions{\productSubbasis{T}{X}}$ by \cref{pow_iff}.
                        $\unions{\productSubbasis{T}{X}} = \productFamily{X}$ by \cref{product_subbasis_is_subbasis,subbasis_on}.
                        Thus $U\subseteq\productFamily{X}$ by \cref{subseteq}.
                        Thus $U\in\pow{\productFamily{X}}$ by \cref{powerset_intro}.
                        Thus $d(n)\in\pow{\productFamily{X}}$ by \cref{function_apply_iff}.
                    \end{subproof}
                    Thus $d$ is a function to $\pow{\productFamily{X}}$.
                \end{subproof}
                Show that $\dom{d} = \naturalsPlus$.
                \begin{subproof}
                    Show that $\naturalsPlus\subseteq\dom{d}$.
                    \begin{subproof}
                        Show that for all $n$ such that $n\in\naturalsPlus$ we have $n\in\dom{d}$.
                        \begin{subproof}
                            Fix $n$.
                            Assume $n\in\naturalsPlus$.
                            $(n, \inters{\img{\apply{h}{n}}{\initialSegment{n}}})\in d$ by \cref{pair_eq_iff}.
                            Thus $n\in\dom{d}$ by \cref{dom_intro}.
                        \end{subproof}
                        Thus $\naturalsPlus\subseteq\dom{d}$ by \cref{subseteq}.
                    \end{subproof}
                    Show that $\dom{d}\subseteq\naturalsPlus$.
                    \begin{subproof}
                        Show that for all $n$ such that $n\in\dom{d}$ we have $n\in\naturalsPlus$.
                        \begin{subproof}
                            Fix $n$.
                            Assume $n\in\dom{d}$.
                            Take $U$ such that $(n,U)\in d$ by \cref{dom_iff}.
                            Thus $n\in\naturalsPlus$ by \cref{pair_eq_iff}.
                        \end{subproof}
                        Thus $\dom{d}\subseteq\naturalsPlus$ by \cref{subseteq}.
                    \end{subproof}
                    Thus $\dom{d} = \naturalsPlus$ by \cref{subseteq_antisymmetric}.
                \end{subproof}
                Thus $d\in\funs{\naturalsPlus}{\pow{\productFamily{X}}}$ by \cref{funs_intro}.
                Thus $d$ is a sequence in $\pow{\productFamily{X}}$ by \cref{sequence_in}.
            \end{subproof}

            Show that for all $n\in\naturalsPlus$ we have $d(n)$ is a neighborhood of $x$ in $\productFamily{X}$ with respect to $T_0$.
            \begin{subproof}
                Fix $n\in\naturalsPlus$.
                Let $F = \img{\apply{h}{n}}{\initialSegment{n}}$.
                Show that $x\in\inters{F}$.
                \begin{subproof}
                    Show that for all $A$ such that $A\in F$ we have $x\in A$.
                    \begin{subproof}
                        Fix $A$.
                        Assume $A\in F$.
                        Take $i$ such that $i\in\initialSegment{n}$ and $(i,A)\in\apply{h}{n}$ by \cref{img_iff}.
                        Thus $n\mathrel{R_d}\apply{h}{n}$ by \cref{choice_function}.
                        Thus $\apply{h}{n}$ is a function and $\dom{\apply{h}{n}} = \initialSegment{n}$ by \cref{pair_eq_iff}.
                        Thus $i\in\dom{\apply{h}{n}}$ by \cref{dom_intro}.
                        Thus $\apply{\apply{h}{n}}{i} = A$ by \cref{function_apply_iff}.
                        Thus for all $j\in\initialSegment{n}$ we have
                            $\apply{\apply{h}{n}}{j} = \preimg{\piIndex{X}{j}}{\apply{c(j)}{n}}$ by \cref{pair_eq_iff}.
                        Thus $A = \preimg{\piIndex{X}{i}}{\apply{c(i)}{n}}$.
                        $c(i)$ is a countable basis at $\apply{x}{i}$ in $\apply{X}{i}$ with respect to $\apply{T}{i}$.
                        Thus $n\in\naturalsPlus$.
                        Thus $\apply{c(i)}{n}$ is a neighborhood of $\apply{x}{i}$ in $\apply{X}{i}$ with respect to $\apply{T}{i}$
                            by \cref{countable_basis_elem_neighborhood}.
                        Thus $\apply{x}{i}\in\apply{c(i)}{n}$ by \cref{neighborhood}.
                        Thus $x\in A$ by \cref{preimg_iff,projection_map_indexed}.
                    \end{subproof}
                    Show that $F$ is inhabited.
                    \begin{subproof}
                        $n\in\initialSegment{n}$ by \cref{initial_segment_naturalsplus}.
                        Thus $n\mathrel{R_d}\apply{h}{n}$ by \cref{choice_function}.
                        Thus $\apply{h}{n}$ is a function and $\dom{\apply{h}{n}} = \initialSegment{n}$ by \cref{pair_eq_iff}.
                        Thus $n\in\dom{\apply{h}{n}}$ by \cref{subseteq}.
                        Thus for all $j\in\initialSegment{n}$ we have
                            $\apply{\apply{h}{n}}{j} = \preimg{\piIndex{X}{j}}{\apply{c(j)}{n}}$ by \cref{pair_eq_iff}.
                        Thus $\apply{\apply{h}{n}}{n} = \preimg{\piIndex{X}{n}}{\apply{c(n)}{n}}$.
                        Thus $(n, \preimg{\piIndex{X}{n}}{\apply{c(n)}{n}})\in\apply{h}{n}$ by \cref{function_apply_iff}.
                        Thus $\preimg{\piIndex{X}{n}}{\apply{c(n)}{n}}\in F$ by \cref{img_iff}.
                    \end{subproof}
                    Thus $x\in\inters{F}$ by \cref{inters_intro}.
                \end{subproof}
                Show that $F\subseteq\productSubbasis{T}{X}$.
                \begin{subproof}
                    Show that for all $A$ such that $A\in F$ we have $A\in\productSubbasis{T}{X}$.
                    \begin{subproof}
                        Fix $A$.
                        Assume $A\in F$.
                        Take $i$ such that $i\in\initialSegment{n}$ and $(i,A)\in\apply{h}{n}$ by \cref{img_iff}.
                        Thus $n\mathrel{R_d}\apply{h}{n}$ by \cref{choice_function}.
                        Thus $\apply{h}{n}$ is a function and $\dom{\apply{h}{n}} = \initialSegment{n}$ by \cref{pair_eq_iff}.
                        Thus $i\in\dom{\apply{h}{n}}$ by \cref{dom_intro}.
                        Thus $\apply{\apply{h}{n}}{i} = A$ by \cref{function_apply_iff}.
                        Thus for all $j\in\initialSegment{n}$ we have
                            $\apply{\apply{h}{n}}{j} = \preimg{\piIndex{X}{j}}{\apply{c(j)}{n}}$ by \cref{pair_eq_iff}.
                        Thus $A = \preimg{\piIndex{X}{i}}{\apply{c(i)}{n}}$.
                        Thus $i\in\naturalsPlus$ by \cref{initial_segment_naturalsplus}.
                        $\dom{X} = \naturalsPlus$.
                        Thus $i\in\dom{X}$ by \cref{subseteq}.
                        $c(i)$ is a countable basis at $\apply{x}{i}$ in $\apply{X}{i}$ with respect to $\apply{T}{i}$.
                        Thus $n\in\naturalsPlus$.
                        Thus $\apply{c(i)}{n}$ is a neighborhood of $\apply{x}{i}$ in $\apply{X}{i}$ with respect to $\apply{T}{i}$
                            by \cref{countable_basis_elem_neighborhood}.
                        Thus $\apply{c(i)}{n}$ is open in $\apply{X}{i}$ with respect to $\apply{T}{i}$ by \cref{neighborhood}.
                        Thus $\apply{c(i)}{n}\in\apply{T}{i}$ by \cref{open_in}.
                        Thus $A\subseteq\dom{\piIndex{X}{i}}$ by \cref{preimg_subseteq_dom}.
                        Show that $\dom{\piIndex{X}{i}}\subseteq\productFamily{X}$.
                        \begin{subproof}
                            Show that for all $z$ such that $z\in\dom{\piIndex{X}{i}}$ we have $z\in\productFamily{X}$.
                            \begin{subproof}
                                Fix $z$.
                                Assume $z\in\dom{\piIndex{X}{i}}$.
                                Take $y$ such that $(z,y)\in\piIndex{X}{i}$ by \cref{dom_iff}.
                                Take $x\in\productFamily{X}$ such that $(z,y) = (x,\apply{x}{i})$ by \cref{projection_map_indexed}.
                                Thus $z = x$ by \cref{pair_eq_iff}.
                                Thus $z\in\productFamily{X}$.
                            \end{subproof}
                            Thus $\dom{\piIndex{X}{i}}\subseteq\productFamily{X}$ by \cref{subseteq}.
                        \end{subproof}
                        Thus $A\subseteq\productFamily{X}$ by \cref{subseteq_transitive}.
                        Thus $A\in\pow{\productFamily{X}}$ by \cref{powerset_intro}.
                        Thus $A\in\productSubbasisAt{T}{X}{i}$ by \cref{product_subbasis_indexed}.
                        Thus $A\in\productSubbasis{T}{X}$ by \cref{product_subbasis_union_indexed}.
                    \end{subproof}
                    Thus $F\subseteq\productSubbasis{T}{X}$ by \cref{subseteq}.
                \end{subproof}
                Show that $F$ is finite.
                \begin{subproof}
                    Thus $n\mathrel{R_d}\apply{h}{n}$ by \cref{choice_function}.
                    Thus $\apply{h}{n}$ is a function on $\initialSegment{n}$ by \cref{pair_eq_iff}.
                    $\initialSegment{n}$ is finite by \cref{initial_segment_finite}.
                    Thus $F$ is finite by \cref{finite_image_function}.
                \end{subproof}
                Show that $F$ is inhabited.
                \begin{subproof}
                    $n\in\initialSegment{n}$ by \cref{initial_segment_naturalsplus}.
                    Thus $n\mathrel{R_d}\apply{h}{n}$ by \cref{choice_function}.
                    Thus $\apply{h}{n}$ is a function and $\dom{\apply{h}{n}} = \initialSegment{n}$ by \cref{pair_eq_iff}.
                    Thus $n\in\dom{\apply{h}{n}}$ by \cref{subseteq}.
                    Thus for all $j\in\initialSegment{n}$ we have
                        $\apply{\apply{h}{n}}{j} = \preimg{\piIndex{X}{j}}{\apply{c(j)}{n}}$ by \cref{pair_eq_iff}.
                    Thus $\apply{\apply{h}{n}}{n} = \preimg{\piIndex{X}{n}}{\apply{c(n)}{n}}$.
                    Thus $(n, \preimg{\piIndex{X}{n}}{\apply{c(n)}{n}})\in\apply{h}{n}$ by \cref{function_apply_iff}.
                    Thus $\preimg{\piIndex{X}{n}}{\apply{c(n)}{n}}\in F$ by \cref{img_iff}.
                \end{subproof}
                Take $A$ such that $A\in F$.
                Thus $A\in\productSubbasis{T}{X}$ by \cref{subseteq}.
                Show that $A\subseteq\unions{\productSubbasis{T}{X}}$.
                \begin{subproof}
                    Show that for all $z$ such that $z\in A$ we have $z\in\unions{\productSubbasis{T}{X}}$.
                    \begin{subproof}
                        Fix $z$.
                        Assume $z\in A$.
                        Thus $z\in\unions{\productSubbasis{T}{X}}$ by \cref{unions_intro}.
                    \end{subproof}
                    Thus $A\subseteq\unions{\productSubbasis{T}{X}}$ by \cref{subseteq}.
                \end{subproof}
                Thus $\inters{F}\subseteq A$ by \cref{inters_subseteq_elem}.
                Thus $\inters{F}\subseteq\unions{\productSubbasis{T}{X}}$ by \cref{subseteq_transitive}.
                Thus $\inters{F}\in\pow{\unions{\productSubbasis{T}{X}}}$ by \cref{powerset_intro}.
                Thus $\inters{F}\in\finiteIntersections{\productSubbasis{T}{X}}$ by \cref{finite_intersections}.
                $\finiteIntersections{\productSubbasis{T}{X}}$ is a basis for $T_0$ on $\productFamily{X}$ by
                    \cref{product_subbasis_finite_intersections_basis,product_topology_indexed,basis_for_topology,topology_generated_by_subbasis}.
                Thus $T_0 = \basisTopology{\finiteIntersections{\productSubbasis{T}{X}}}{\productFamily{X}}$ by \cref{basis_for_topology}.
                Thus $\inters{F}\in\basisTopology{\finiteIntersections{\productSubbasis{T}{X}}}{\productFamily{X}}$
                    by \cref{basis_elem_open_in_generated,basis_for_topology}.
                Thus $\inters{F}\in T_0$ by \cref{basis_for_topology}.
                $\basisTopology{\finiteIntersections{\productSubbasis{T}{X}}}{\productFamily{X}}$ is a topology on $\productFamily{X}$
                    by \cref{basis_topology_is_topology,basis_for_topology}.
                Thus $T_0$ is a topology on $\productFamily{X}$ by \cref{basis_for_topology}.
                Thus $\inters{F}$ is open in $\productFamily{X}$ with respect to $T_0$ by \cref{open_in}.
                Thus $\inters{F}$ is a neighborhood of $x$ in $\productFamily{X}$ with respect to $T_0$ by \cref{neighborhood}.
                $d\in\funs{\naturalsPlus}{\pow{\productFamily{X}}}$ by \cref{sequence_in}.
                Thus $d$ is a function by \cref{funs_is_function}.
                Thus $\dom{d} = \naturalsPlus$ by \cref{funs_elim}.
                Thus $n\in\dom{d}$ by \cref{subseteq}.
                $(n, \inters{F})\in d$ by \cref{pair_eq_iff}.
                Thus $d(n) = \inters{F}$ by \cref{function_apply_iff}.
                Thus $d(n)$ is a neighborhood of $x$ in $\productFamily{X}$ with respect to $T_0$.
            \end{subproof}

            Show that for all $U$ such that $U$ is a neighborhood of $x$ in $\productFamily{X}$ with respect to $T_0$
                there exists $n\in\naturalsPlus$ such that $d(n)\subseteq U$.
            \begin{subproof}
                Fix $U$.
                Assume $U$ is a neighborhood of $x$ in $\productFamily{X}$ with respect to $T_0$.
                Thus $U$ is open in $\productFamily{X}$ with respect to $T_0$ and $x\in U$ by \cref{neighborhood}.
                Show that $\productSubbasis{T}{X}$ is a subbasis on $\productFamily{X}$.
                \begin{subproof}
                    $X$ is a function.
                    Thus $X$ is right-unique and $X$ is a relation.
                    $T$ is a function.
                    Thus $T$ is right-unique and $T$ is a relation.
                    $\dom{T} = \dom{X}$.
                    $1\in\naturalsPlus$ by \cref{real_one_in_naturals}.
                    $\dom{X} = \naturalsPlus$.
                    Thus $1\in\dom{X}$ by \cref{subseteq}.
                    Thus there exists $a$ such that $a\in\dom{X}$.
                    Thus for all $n\in\dom{X}$ we have $\apply{T}{n}$ is a topology on $\apply{X}{n}$.
                    Thus $\productSubbasis{T}{X}$ is a subbasis on $\productFamily{X}$ by \cref{product_subbasis_is_subbasis}.
                \end{subproof}
                Show that $\finiteIntersections{\productSubbasis{T}{X}}$ is a basis for $T_0$ on $\productFamily{X}$.
                \begin{subproof}
                    $X$ is a function.
                    $T$ is a function.
                    $\dom{T} = \dom{X}$.
                    $1\in\naturalsPlus$ by \cref{real_one_in_naturals}.
                    $\dom{X} = \naturalsPlus$.
                    Thus $1\in\dom{X}$ by \cref{subseteq}.
                    Thus $\dom{X}$ is inhabited.
                    Thus for all $n\in\dom{X}$ we have $\apply{T}{n}$ is a topology on $\apply{X}{n}$.
                    Thus $\finiteIntersections{\productSubbasis{T}{X}}$ is a basis for $\productTopologyIndexed{T}{X}$ on $\productFamily{X}$
                        by \cref{product_subbasis_finite_intersections_basis}.
                    $T_0 = \productTopologyIndexed{T}{X}$.
                    Thus $\finiteIntersections{\productSubbasis{T}{X}}$ is a basis for $T_0$ on $\productFamily{X}$.
                \end{subproof}
                $U$ is open in $\productFamily{X}$ with respect to $T_0$.
                Thus $U\in T_0$ by \cref{open_in}.
                Thus $T_0 = \basisTopology{\finiteIntersections{\productSubbasis{T}{X}}}{\productFamily{X}}$
                    by \cref{basis_for_topology}.
                Thus $U\in\basisTopology{\finiteIntersections{\productSubbasis{T}{X}}}{\productFamily{X}}$.
                Take $B\in\finiteIntersections{\productSubbasis{T}{X}}$ such that $x\in B$ and $B\subseteq U$
                    by \cref{topology_generated_by_basis}.
                Take $F\subseteq\productSubbasis{T}{X}$ such that $F$ is finite and inhabited and $B = \inters{F}$ by \cref{finite_intersections}.

                Let $R_2 = \{(A,i) \mid A\in F \mid
                    \text{$i\in\dom{X}$ and there exists $V\in\apply{T}{i}$ such that
                    $A = \preimg{\piIndex{X}{i}}{V}$}\}$.
                Show that $R_2$ is a relation.
                \begin{subproof}
                    Show that for all $w$ such that $w\in R_2$ there exist $A,i$ such that $w = (A,i)$.
                    \begin{subproof}
                        Fix $w$.
                        Assume $w\in R_2$.
                        Thus there exist $A,i$ such that $w = (A,i)$ by \cref{pair_eq_iff}.
                    \end{subproof}
                    Thus $R_2$ is a relation by \cref{relation}.
                \end{subproof}
                Show that for all $A\in F$ there exists $i$ such that $A\mathrel{R_2} i$.
                \begin{subproof}
                    Fix $A\in F$.
                    $A\in\productSubbasis{T}{X}$ by \cref{subseteq}.
                    Take $i$ such that $i\in\dom{X}$ and $A\in\productSubbasisAt{T}{X}{i}$
                        by \cref{product_subbasis_union_indexed}.
                    Take $V\in\apply{T}{i}$ such that $A = \preimg{\piIndex{X}{i}}{V}$ by \cref{product_subbasis_indexed}.
                    Thus $(A,i)\in R_2$ by \cref{pair_eq_iff}.
                    Thus $A\mathrel{R_2} i$ by \cref{pair_eq_iff}.
                \end{subproof}
                Take $p$ such that $p$ is a function on $F$ and for all $A\in F$ we have $A\mathrel{R_2}\apply{p}{A}$
                    by \cref{choice_function}.
                Let $J = \img{p}{F}$.
                $F$ is finite.
                Thus $J$ is finite by \cref{finite_image_function}.
                Show that $J$ is inhabited.
                \begin{subproof}
                    Take $A$ such that $A\in F$.
                    $p$ is a function on $F$.
                    Thus $p$ is a function and $\dom{p} = F$.
                    Thus $A\in\dom{p}$ by \cref{subseteq}.
                    Thus $(A,\apply{p}{A})\in p$ by \cref{function_apply_iff}.
                    Thus $\apply{p}{A}\in J$ by \cref{img_iff}.
                \end{subproof}

                Let $R_3 = \{(A,k) \mid A\in F \mid
                    \text{there exists $V\in\apply{T}{\apply{p}{A}}$ such that
                    $A = \preimg{\piIndex{X}{\apply{p}{A}}}{V}$ and
                    $k\in\naturalsPlus$ and $\apply{c(\apply{p}{A})}{k}\subseteq V$}\}$.
                Show that $R_3$ is a relation.
                \begin{subproof}
                    Show that for all $w$ such that $w\in R_3$ there exist $A,k$ such that $w = (A,k)$.
                    \begin{subproof}
                        Fix $w$.
                        Assume $w\in R_3$.
                        Thus there exist $A,k$ such that $w = (A,k)$ by \cref{pair_eq_iff}.
                    \end{subproof}
                    Thus $R_3$ is a relation by \cref{relation}.
                \end{subproof}
                Show that for all $A\in F$ there exists $k$ such that $A\mathrel{R_3} k$.
                \begin{subproof}
                    Fix $A\in F$.
                    Thus $A\mathrel{R_2}\apply{p}{A}$ by \cref{choice_function}.
                    Thus $\apply{p}{A}\in\dom{X}$ and there exists $V\in\apply{T}{\apply{p}{A}}$ such that
                        $A = \preimg{\piIndex{X}{\apply{p}{A}}}{V}$ by \cref{pair_eq_iff}.
                    Take $V\in\apply{T}{\apply{p}{A}}$ such that $A = \preimg{\piIndex{X}{\apply{p}{A}}}{V}$.
                    Thus $B\subseteq A$ by \cref{inters_subseteq_elem}.
                    $x\in B$.
                    Thus $x\in A$ by \cref{subseteq}.
                    Thus $x\in \preimg{\piIndex{X}{\apply{p}{A}}}{V}$ by \cref{subseteq}.
                    Take $y\in V$ such that $(x,y)\in\piIndex{X}{\apply{p}{A}}$ by \cref{preimg_iff}.
                    Take $z\in\productFamily{X}$ such that $(z,\apply{z}{\apply{p}{A}}) = (x,y)$ by \cref{projection_map_indexed}.
                    Thus $z = x$ and $\apply{z}{\apply{p}{A}} = y$ by \cref{pair_eq_iff}.
                    Thus $\apply{x}{\apply{p}{A}} = y$.
                    Thus $\apply{x}{\apply{p}{A}}\in V$.
                    Show that $\apply{T}{\apply{p}{A}}$ is a topology on $\apply{X}{\apply{p}{A}}$.
                    \begin{subproof}
                        $\dom{X} = \naturalsPlus$.
                        Thus $\apply{p}{A}\in\naturalsPlus$ by \cref{subseteq}.
                        Thus $\apply{T}{\apply{p}{A}}$ is a topology on $\apply{X}{\apply{p}{A}}$.
                    \end{subproof}
                    Thus $V$ is open in $\apply{X}{\apply{p}{A}}$ with respect to $\apply{T}{\apply{p}{A}}$ by \cref{open_in}.
                    Thus $V$ is a neighborhood of $\apply{x}{\apply{p}{A}}$ in $\apply{X}{\apply{p}{A}}$ with respect to $\apply{T}{\apply{p}{A}}$ by \cref{neighborhood}.
                    $c(\apply{p}{A})$ is a countable basis at $\apply{x}{\apply{p}{A}}$ in $\apply{X}{\apply{p}{A}}$ with respect to $\apply{T}{\apply{p}{A}}$.
                    Thus there exists $k\in\naturalsPlus$ such that $\apply{c(\apply{p}{A})}{k}\subseteq V$ by \cref{countable_basis_at_point}.
                    Take $k\in\naturalsPlus$ such that $\apply{c(\apply{p}{A})}{k}\subseteq V$.
                    Thus $(A,k)\in R_3$ by \cref{pair_eq_iff}.
                    Thus $A\mathrel{R_3} k$ by \cref{pair_eq_iff}.
                \end{subproof}
                Take $q$ such that $q$ is a function on $F$ and for all $A\in F$ we have $A\mathrel{R_3}\apply{q}{A}$
                    by \cref{choice_function}.
                Let $K = \img{q}{F}$.
                $F$ is finite.
                Thus $K$ is finite by \cref{finite_image_function}.
                Show that $K$ is inhabited.
                \begin{subproof}
                    Take $A$ such that $A\in F$.
                    $q$ is a function on $F$.
                    Thus $q$ is a function and $\dom{q} = F$.
                    Thus $A\in\dom{q}$ by \cref{subseteq}.
                    Thus $(A,\apply{q}{A})\in q$ by \cref{function_apply_iff}.
                    Thus $\apply{q}{A}\in K$ by \cref{img_iff}.
                \end{subproof}

                Show that $J\subseteq\naturalsPlus$.
                \begin{subproof}
                    Show that for all $j$ such that $j\in J$ we have $j\in\naturalsPlus$.
                    \begin{subproof}
                        Fix $j$.
                        Assume $j\in J$.
                        Take $A$ such that $A\in F$ and $(A,j)\in p$ by \cref{img_iff}.
                        $p$ is a function on $F$.
                        Thus $p$ is a function and $\dom{p} = F$.
                        Thus $A\in\dom{p}$ by \cref{subseteq}.
                        Thus $\apply{p}{A} = j$ by \cref{function_apply_iff}.
                        Thus $A\mathrel{R_2}\apply{p}{A}$ by \cref{choice_function}.
                        Thus $\apply{p}{A}\in\dom{X}$ by \cref{pair_eq_iff}.
                        $\dom{X} = \naturalsPlus$.
                        Thus $j\in\naturalsPlus$.
                    \end{subproof}
                    Thus $J\subseteq\naturalsPlus$ by \cref{subseteq}.
                \end{subproof}

                Show that $K\subseteq\naturalsPlus$.
                \begin{subproof}
                    Show that for all $k$ such that $k\in K$ we have $k\in\naturalsPlus$.
                    \begin{subproof}
                        Fix $k$.
                        Assume $k\in K$.
                        Take $A$ such that $A\in F$ and $(A,k)\in q$ by \cref{img_iff}.
                        $q$ is a function on $F$.
                        Thus $q$ is a function and $\dom{q} = F$.
                        Thus $A\in\dom{q}$ by \cref{subseteq}.
                        Thus $\apply{q}{A} = k$ by \cref{function_apply_iff}.
                        Thus $A\mathrel{R_3}\apply{q}{A}$ by \cref{choice_function}.
                        Thus $\apply{q}{A}\in\naturalsPlus$ by \cref{pair_eq_iff}.
                        Thus $k\in\naturalsPlus$.
                    \end{subproof}
                    Thus $K\subseteq\naturalsPlus$ by \cref{subseteq}.
                \end{subproof}

                Let $G = J\union K$.
                Thus $G\subseteq\naturalsPlus$ by \cref{union_subsets_is_subset}.
                $\naturalsPlus\subseteq\reals$ by \cref{real_naturals_subset_reals}.
                Thus $G\subseteq\reals$ by \cref{subseteq_transitive}.
                $G$ is finite by \cref{union_of_finite_is_finite}.
                $G$ is inhabited.
                Take $m$ such that $m$ is a largest element of $G$ with respect to $\realOrderRel$
                    by \cref{finite_inhabited_largest_element,real_order_relation_simple}.
                Thus $m\in G$ by \cref{largest_element}.
                Thus $m\in\naturalsPlus$ by \cref{subseteq}.

                Show that for all $A\in F$ we have $d(m)\subseteq A$.
                \begin{subproof}
                    Fix $A\in F$.
                    Thus $A\mathrel{R_2}\apply{p}{A}$ by \cref{choice_function}.
                    $p$ is a function on $F$.
                    Thus $p$ is a function and $\dom{p} = F$.
                    Thus $A\in\dom{p}$ by \cref{subseteq}.
                    Thus $(A,\apply{p}{A})\in p$ by \cref{function_apply_iff}.
                    Thus $\apply{p}{A}\in J$ by \cref{img_iff}.
                    Thus $\apply{p}{A}\in G$ by \cref{union_intro_left}.
                    Show that $\apply{p}{A}\leq m$.
                    \begin{subproof}
                        Thus $\apply{p}{A}\mathrel{\realOrderRel} m$ or $\apply{p}{A} = m$ by \cref{largest_element}.
                        \begin{byCase}
                        \caseOf{$\apply{p}{A}\mathrel{\realOrderRel} m$.}
                            Thus $\apply{p}{A} < m$ by \cref{real_order_relation_elim}.
                            Thus $\apply{p}{A}\leq m$.
                        \caseOf{$\apply{p}{A} = m$.}
                            Thus $\apply{p}{A}\leq m$.
                        \end{byCase}
                    \end{subproof}
                    Thus $\apply{p}{A}\in\dom{X}$ by \cref{pair_eq_iff}.
                    $\dom{X} = \naturalsPlus$.
                    Thus $\apply{p}{A}\in\naturalsPlus$.
                    Thus $\apply{p}{A}\in\initialSegment{m}$ by \cref{initial_segment_naturalsplus}.

                    $A\mathrel{R_3}\apply{q}{A}$ by \cref{choice_function}.
                    $q$ is a function on $F$.
                    Thus $q$ is a function and $\dom{q} = F$.
                    Thus $A\in\dom{q}$ by \cref{subseteq}.
                    Thus $(A,\apply{q}{A})\in q$ by \cref{function_apply_iff}.
                    Thus $\apply{q}{A}\in K$ by \cref{img_iff}.
                    Thus $\apply{q}{A}\in G$ by \cref{union_intro_right}.
                    Show that $\apply{q}{A}\leq m$.
                    \begin{subproof}
                        Thus $\apply{q}{A}\mathrel{\realOrderRel} m$ or $\apply{q}{A} = m$ by \cref{largest_element}.
                        \begin{byCase}
                        \caseOf{$\apply{q}{A}\mathrel{\realOrderRel} m$.}
                            Thus $\apply{q}{A} < m$ by \cref{real_order_relation_elim}.
                            Thus $\apply{q}{A}\leq m$.
                        \caseOf{$\apply{q}{A} = m$.}
                            Thus $\apply{q}{A}\leq m$.
                        \end{byCase}
                    \end{subproof}

                    Let $i = \apply{p}{A}$.
                    Thus $i\in\initialSegment{m}$.
                    Thus $A\mathrel{R_2} i$ by \cref{choice_function}.
                    Thus $i\in\dom{X}$ by \cref{pair_eq_iff}.
                    Let $k = \apply{q}{A}$.
                    Thus $A\mathrel{R_3} k$ by \cref{choice_function}.
                    Thus there exists $V\in\apply{T}{i}$ such that
                        $A = \preimg{\piIndex{X}{i}}{V}$ and $k\in\naturalsPlus$ and $\apply{c(i)}{k}\subseteq V$
                        by \cref{pair_eq_iff}.
                    Take $V\in\apply{T}{i}$ such that
                        $A = \preimg{\piIndex{X}{i}}{V}$ and $k\in\naturalsPlus$ and $\apply{c(i)}{k}\subseteq V$.

                    Thus $k\leq m$.
                    Thus $i\mathrel{R_1}\apply{c}{i}$ by \cref{choice_function}.
                    Thus for all $m_1,k_1\in\naturalsPlus$ if $m_1\leq k_1$ then
                        $\apply{c(i)}{k_1}\subseteq \apply{c(i)}{m_1}$ by \cref{pair_eq_iff}.
                    Thus $\apply{c(i)}{m}\subseteq \apply{c(i)}{k}$.
                    Thus $\apply{c(i)}{m}\subseteq V$ by \cref{subseteq}.

                    Let $F_m = \img{\apply{h}{m}}{\initialSegment{m}}$.
                    Thus $m\mathrel{R_d}\apply{h}{m}$ by \cref{choice_function}.
                    Thus $\apply{h}{m}$ is right-unique and $\apply{h}{m}$ is a relation and
                        $\initialSegment{m} = \dom{\apply{h}{m}}$ and
                        for all $j\in\initialSegment{m}$ we have
                        $\apply{\apply{h}{m}}{j} = \preimg{\piIndex{X}{j}}{\apply{c(j)}{m}}$ by \cref{pair_eq_iff}.
                    Thus $\apply{h}{m}$ is a function.
                    Thus $\apply{\apply{h}{m}}{i} = \preimg{\piIndex{X}{i}}{\apply{c(i)}{m}}$.
                    Thus $i\in\dom{\apply{h}{m}}$ by \cref{subseteq}.
                    Thus $(i,\apply{\apply{h}{m}}{i})\in\apply{h}{m}$ by \cref{function_apply_iff}.
                    Thus $(i,\preimg{\piIndex{X}{i}}{\apply{c(i)}{m}})\in\apply{h}{m}$.
                    Thus $\preimg{\piIndex{X}{i}}{\apply{c(i)}{m}}\in F_m$ by \cref{img_iff}.
                    $d\in\funs{\naturalsPlus}{\pow{\productFamily{X}}}$ by \cref{sequence_in}.
                    Thus $d$ is a function by \cref{funs_is_function}.
                    Thus $\dom{d} = \naturalsPlus$ by \cref{funs_elim}.
                    Thus $m\in\dom{d}$ by \cref{subseteq}.
                    $(m, \inters{F_m})\in d$ by \cref{pair_eq_iff}.
                    Thus $d(m) = \inters{F_m}$ by \cref{function_apply_iff}.
                    Thus $d(m)\subseteq \preimg{\piIndex{X}{i}}{\apply{c(i)}{m}}$ by \cref{inters_subseteq_elem}.
                    Thus $\preimg{\piIndex{X}{i}}{\apply{c(i)}{m}}\subseteq \preimg{\piIndex{X}{i}}{V}$ by \cref{preimg_subseteq}.
                    Thus $\preimg{\piIndex{X}{i}}{\apply{c(i)}{m}}\subseteq A$ by \cref{subseteq}.
                    Thus $d(m)\subseteq A$ by \cref{subseteq_transitive}.
                \end{subproof}
                Thus $d(m)\subseteq \inters{F}$ by \cref{subseteq_inters_iff}.
                Thus $d(m)\subseteq B$ by \cref{subseteq}.
                Thus $d(m)\subseteq U$ by \cref{subseteq_transitive}.
            \end{subproof}
            Thus $d$ is a countable basis at $x$ in $\productFamily{X}$ with respect to $T_0$ by \cref{countable_basis_at_point}.
        \end{subproof}
    \end{subproof}
    Thus $\productFamily{X}$ satisfies the first countability axiom with respect to $T_0$ by \cref{first_countability_axiom}.
\end{proof}

% from §30 Item 9: subspace of a second-countable space is second-countable
\begin{lemma}\label{subspace_second_countable}
    Assume $T$ is a topology on $X$.
    Suppose $A\subseteq X$.
    Suppose $X$ satisfies the second countability axiom with respect to $T$.
    Then $A$ satisfies the second countability axiom with respect to $\subspaceTopology{T}{A}$.
\end{lemma}
\begin{proof}
    Let $T_1 = \subspaceTopology{T}{A}$.
    $T_1$ is a topology on $A$ by \cref{subspace_topology_is_topology}.
    Take $b$ such that
        $b$ is a sequence in $\pow{X}$ and
        for all $n\in\naturalsPlus$ we have $b(n)$ is open in $X$ with respect to $T$ and
        for all $x\in X$ we have
            for all $U$ such that $U$ is open in $X$ with respect to $T$ and $x\in U$
                there exists $n\in\naturalsPlus$ such that $x\in b(n)$ and $b(n)\subseteq U$
        by \cref{second_countability_axiom}.
    Let $c = \{(n, A\inter b(n)) \mid n\in\naturalsPlus\}$.
    Show that $c$ is a sequence in $\pow{A}$.
    \begin{subproof}
        Show that $c$ is a function to $\pow{A}$.
        \begin{subproof}
            Show that $c$ is a function.
            \begin{subproof}
                Show that for all $n,U_1,U_2$ such that $(n,U_1)\in c$ and $(n,U_2)\in c$ we have $U_1 = U_2$.
                \begin{subproof}
                    Fix $n$.
                    Fix $U_1$.
                    Fix $U_2$.
                    Assume $(n,U_1)\in c$ and $(n,U_2)\in c$.
                    Thus $U_1 = A\inter b(n)$ and $U_2 = A\inter b(n)$ by \cref{pair_eq_iff}.
                    Thus $U_1 = U_2$.
                \end{subproof}
                Thus $c$ is right-unique by \cref{rightunique}.
                Show that $c$ is a relation.
                \begin{subproof}
                    Show that for all $w$ such that $w\in c$ there exist $n,U$ such that $w = (n,U)$.
                    \begin{subproof}
                        Fix $w$.
                        Assume $w\in c$.
                        Thus there exist $n,U$ such that $w = (n,U)$ by \cref{pair_eq_iff}.
                    \end{subproof}
                    Thus $c$ is a relation by \cref{relation}.
                \end{subproof}
                Thus $c$ is a function.
            \end{subproof}
            Show that for all $n\in\dom{c}$ we have $c(n)\in\pow{A}$.
            \begin{subproof}
                Fix $n\in\dom{c}$.
                Take $U$ such that $(n,U)\in c$ by \cref{dom_iff}.
                Thus $U = A\inter b(n)$ by \cref{pair_eq_iff}.
                Show that $A\inter b(n)\subseteq A$.
                \begin{subproof}
                    Show that for all $z$ such that $z\in A\inter b(n)$ we have $z\in A$.
                    \begin{subproof}
                        Fix $z$.
                        Assume $z\in A\inter b(n)$.
                        Thus $z\in A$ by \cref{inter_elim_left}.
                    \end{subproof}
                    Thus $A\inter b(n)\subseteq A$ by \cref{subseteq}.
                \end{subproof}
                Thus $A\inter b(n)\in\pow{A}$ by \cref{powerset_intro}.
                Thus $c(n)\in\pow{A}$ by \cref{function_apply_iff}.
            \end{subproof}
            Thus $c$ is a function to $\pow{A}$.
        \end{subproof}
        Show that $\dom{c} = \naturalsPlus$.
        \begin{subproof}
            Show that $\naturalsPlus\subseteq\dom{c}$.
            \begin{subproof}
                Show that for all $n$ such that $n\in\naturalsPlus$ we have $n\in\dom{c}$.
                \begin{subproof}
                    Fix $n$.
                    Assume $n\in\naturalsPlus$.
                    $(n, A\inter b(n))\in c$ by \cref{pair_eq_iff}.
                    Thus $n\in\dom{c}$ by \cref{dom_intro}.
                \end{subproof}
                Thus $\naturalsPlus\subseteq\dom{c}$ by \cref{subseteq}.
            \end{subproof}
            Show that $\dom{c}\subseteq\naturalsPlus$.
            \begin{subproof}
                Show that for all $n$ such that $n\in\dom{c}$ we have $n\in\naturalsPlus$.
                \begin{subproof}
                    Fix $n$.
                    Assume $n\in\dom{c}$.
                    Take $U$ such that $(n,U)\in c$ by \cref{dom_iff}.
                    Thus $n\in\naturalsPlus$ by \cref{pair_eq_iff}.
                \end{subproof}
                Thus $\dom{c}\subseteq\naturalsPlus$ by \cref{subseteq}.
            \end{subproof}
            Thus $\dom{c} = \naturalsPlus$ by \cref{subseteq_antisymmetric}.
        \end{subproof}
        Thus $c\in\funs{\naturalsPlus}{\pow{A}}$ by \cref{funs_intro}.
        Thus $c$ is a sequence in $\pow{A}$ by \cref{sequence_in}.
    \end{subproof}
    Show that for all $n\in\naturalsPlus$ we have $c(n)$ is open in $A$ with respect to $T_1$.
    \begin{subproof}
        Fix $n\in\naturalsPlus$.
        $b(n)$ is open in $X$ with respect to $T$.
        Thus $b(n)\in T$ by \cref{open_in}.
        Thus $A\inter b(n)\in T_1$ by \cref{subspace_topology}.
        Thus $A\inter b(n)$ is open in $A$ with respect to $T_1$ by \cref{open_in}.
        $c\in\funs{\naturalsPlus}{\pow{A}}$ by \cref{sequence_in}.
        Thus $c$ is a function by \cref{funs_is_function}.
        Thus $\dom{c} = \naturalsPlus$ by \cref{funs_elim}.
        Thus $n\in\dom{c}$ by \cref{subseteq}.
        $(n, A\inter b(n))\in c$ by \cref{pair_eq_iff}.
        Thus $c(n) = A\inter b(n)$ by \cref{function_apply_iff}.
        Thus $c(n)$ is open in $A$ with respect to $T_1$.
    \end{subproof}
    Show that for all $x\in A$ we have
        for all $U$ such that $U$ is open in $A$ with respect to $T_1$ and $x\in U$
            there exists $n\in\naturalsPlus$ such that $x\in c(n)$ and $c(n)\subseteq U$.
    \begin{subproof}
        Fix $x\in A$.
        Fix $U$.
        Assume $U$ is open in $A$ with respect to $T_1$ and $x\in U$.
        Thus $U\in T_1$ by \cref{open_in}.
        Take $V\in T$ such that $U = A\inter V$ by \cref{subspace_topology}.
        Thus $x\in V$ by \cref{inter_elim_right}.
        $V$ is open in $X$ with respect to $T$ by \cref{open_in}.
        Thus there exists $n\in\naturalsPlus$ such that $x\in b(n)$ and $b(n)\subseteq V$.
        Take $n\in\naturalsPlus$ such that $x\in b(n)$ and $b(n)\subseteq V$.
        $c\in\funs{\naturalsPlus}{\pow{A}}$ by \cref{sequence_in}.
        Thus $c$ is a function by \cref{funs_is_function}.
        Thus $\dom{c} = \naturalsPlus$ by \cref{funs_elim}.
        Thus $n\in\dom{c}$ by \cref{subseteq}.
        $(n, A\inter b(n))\in c$ by \cref{pair_eq_iff}.
        Thus $c(n) = A\inter b(n)$ by \cref{function_apply_iff}.
        $x\in A\inter b(n)$ by \cref{inter_intro}.
        Thus $x\in c(n)$.
        Show that $A\inter b(n)\subseteq A\inter V$.
        \begin{subproof}
            Show that for all $z$ such that $z\in A\inter b(n)$ we have $z\in A\inter V$.
            \begin{subproof}
                Fix $z$.
                Assume $z\in A\inter b(n)$.
                Thus $z\in A$ by \cref{inter_elim_left}.
                Thus $z\in b(n)$ by \cref{inter_elim_right}.
                Thus $z\in V$ by \cref{subseteq}.
                Thus $z\in A\inter V$ by \cref{inter_intro}.
            \end{subproof}
            Thus $A\inter b(n)\subseteq A\inter V$ by \cref{subseteq}.
        \end{subproof}
        Thus $c(n)\subseteq U$ by \cref{subseteq}.
    \end{subproof}
    Thus $A$ satisfies the second countability axiom with respect to $T_1$ by \cref{second_countability_axiom}.
\end{proof}

% from §30 helper lemma 3: naturalsPlus closed under addition
\begin{lemma}\label{naturalsplus_add_closed}
    Suppose $a\in\naturalsPlus$ and $b\in\naturalsPlus$.
    Then $a + b\in\naturalsPlus$.
\end{lemma}
\begin{proof}
    Let $A = \{b\in\naturalsPlus \mid a + b\in\naturalsPlus\}$.
    Show that for all $b\in\naturalsPlus$ we have $b\in A$.
    \begin{subproof}
        Show that $A\subseteq\naturalsPlus$.
        \begin{subproof}
            Show that for all $b$ such that $b\in A$ we have $b\in\naturalsPlus$.
            \begin{subproof}
                Fix $b$.
                Assume $b\in A$.
                Thus $b\in\naturalsPlus$.
            \end{subproof}
            Thus $A\subseteq\naturalsPlus$ by \cref{subseteq}.
        \end{subproof}
        $1\in\naturalsPlus$ by \cref{real_one_in_naturals}.
        $a + 1\in\naturalsPlus$ by \cref{real_naturals_closed_succ}.
        Thus $1\in A$.
        Show that for all $n\in A$ we have $n + 1\in A$.
        \begin{subproof}
            Fix $n\in A$.
            Thus $n\in\naturalsPlus$.
            Thus $n + 1\in\naturalsPlus$ by \cref{real_naturals_closed_succ}.
            $a\in\reals$ by \cref{real_naturals_subset_reals,subseteq}.
            $n\in\reals$ by \cref{real_naturals_subset_reals,subseteq}.
            $1\in\reals$ by \cref{real_naturals_subset_reals,real_one_in_naturals,subseteq}.
            $a + n\in\naturalsPlus$.
            Thus $(a + n) + 1\in\naturalsPlus$ by \cref{real_naturals_closed_succ}.
            $a + (n + 1) = (a + n) + 1$ by \cref{real_add_assoc}.
            Thus $a + (n + 1)\in\naturalsPlus$.
            Thus $n + 1\in A$.
        \end{subproof}
        Thus for all $b\in\naturalsPlus$ we have $b\in A$ by \cref{real_naturals_induction}.
    \end{subproof}
    Thus $b\in A$.
    Thus $a + b\in\naturalsPlus$.
\end{proof}

% from §30 helper lemma 4: naturalsPlus closed under multiplication
\begin{lemma}\label{naturalsplus_mul_closed}
    Suppose $a\in\naturalsPlus$ and $b\in\naturalsPlus$.
    Then $a \rmul b\in\naturalsPlus$.
\end{lemma}
\begin{proof}
    Let $A = \{b\in\naturalsPlus \mid a \rmul b\in\naturalsPlus\}$.
    Show that for all $b\in\naturalsPlus$ we have $b\in A$.
    \begin{subproof}
        Show that $A\subseteq\naturalsPlus$.
        \begin{subproof}
            Show that for all $b$ such that $b\in A$ we have $b\in\naturalsPlus$.
            \begin{subproof}
                Fix $b$.
                Assume $b\in A$.
                Thus $b\in\naturalsPlus$.
            \end{subproof}
            Thus $A\subseteq\naturalsPlus$ by \cref{subseteq}.
        \end{subproof}
        $1\in\naturalsPlus$ by \cref{real_one_in_naturals}.
        $a\in\reals$ by \cref{real_naturals_subset_reals,subseteq}.
        $1\in\reals$ by \cref{real_naturals_subset_reals,real_one_in_naturals,subseteq}.
        $a \rmul 1 = 1 \rmul a$ by \cref{real_mul_comm}.
        $1 \rmul a = a$ by \cref{real_mul_ident}.
        Thus $a \rmul 1\in\naturalsPlus$.
        Thus $1\in A$.
        Show that for all $n\in A$ we have $n + 1\in A$.
        \begin{subproof}
            Fix $n\in A$.
            Thus $n\in\naturalsPlus$.
            Thus $a \rmul n\in\naturalsPlus$.
            Thus $a + (a \rmul n)\in\naturalsPlus$ by \cref{naturalsplus_add_closed}.
            $a\in\reals$ by \cref{real_naturals_subset_reals,subseteq}.
            $n\in\reals$ by \cref{real_naturals_subset_reals,subseteq}.
            $1\in\reals$ by \cref{real_naturals_subset_reals,real_one_in_naturals,subseteq}.
            $a \rmul (n + 1) = (a \rmul n) + (a \rmul 1)$ by \cref{real_distrib}.
            $a \rmul 1 = a$ by \cref{real_mul_comm,real_mul_ident}.
            Thus $a \rmul (n + 1) = (a \rmul n) + a$.
            Thus $a \rmul (n + 1)\in\naturalsPlus$.
            Thus $n + 1\in A$.
        \end{subproof}
        Thus for all $b\in\naturalsPlus$ we have $b\in A$ by \cref{real_naturals_induction}.
    \end{subproof}
    Thus $b\in A$.
    Thus $a \rmul b\in\naturalsPlus$.
\end{proof}

% from §30 helper lemma 5: successor product below next square
\begin{lemma}\label{naturalsplus_mul_succ_lt_square}
    Suppose $s\in\naturalsPlus$.
    Then $s \rmul (s + 1) < (s + 1) \rmul (s + 1)$.
\end{lemma}
\begin{proof}
    $s\in\reals$ by \cref{real_naturals_subset_reals,subseteq}.
    $1\in\reals$ by \cref{real_naturals_subset_reals,real_one_in_naturals,subseteq}.
    $0\in\reals$ by \cref{real_add_ident}.
    $0 < 1$ by \cref{real_zero_lt_one}.
    $0 + s < 1 + s$ by \cref{real_lt_add}.
    $0 + s = s$ by \cref{real_add_ident}.
    $1 + s = s + 1$ by \cref{real_add_comm}.
    Thus $s < s + 1$.
    $s + 1\in\naturalsPlus$ by \cref{real_naturals_closed_succ}.
    $s + 1\in\reals$ by \cref{real_naturals_subset_reals,real_naturals_closed_succ,subseteq}.
    $1 < s + 1$ or $1 = s + 1$ by \cref{real_one_leq_naturalsplus}.
    Show that $0 < s + 1$.
    \begin{subproof}
        \begin{byCase}
        \caseOf{$1 < s + 1$.}
            Thus $0 < s + 1$ by \cref{real_lt_trans}.
        \caseOf{$1 = s + 1$.}
            Thus $0 < s + 1$ by \cref{real_zero_lt_one}.
        \end{byCase}
    \end{subproof}
    Thus $s \rmul (s + 1) < (s + 1) \rmul (s + 1)$ by \cref{real_lt_mul_pos}.
\end{proof}

% from §30 helper lemma 6: bound for square against successor product
\begin{lemma}\label{naturalsplus_mul_succ_lt_square_bound}
    Suppose $s\in\naturalsPlus$ and $t\in\naturalsPlus$ and $s < t$.
    Then $s \rmul (s + 1) < t \rmul t$.
\end{lemma}
\begin{proof}
    $t = 1$ or there exists $m\in\naturalsPlus$ such that $t = m + 1$ by \cref{real_naturals_predecessor}.
    \begin{byCase}
    \caseOf{$t = 1$.}
        Show that $s \rmul (s + 1) < t \rmul t$.
        \begin{subproof}
            Suppose not.
            $1 < s$ or $1 = s$ by \cref{real_one_leq_naturalsplus}.
            Thus $s < 1$.
            $1\in\reals$ by \cref{real_naturals_subset_reals,real_one_in_naturals,subseteq}.
            $s\in\reals$ by \cref{real_naturals_subset_reals,subseteq}.
            Show that it is wrong that $s < 1$.
            \begin{subproof}
                Suppose not.
                Thus $s < 1$.
                \begin{byCase}
                \caseOf{$1 < s$.}
                    Thus $1 < 1$ by \cref{real_lt_trans}.
                    Contradiction by \cref{real_lt_irreflexive}.
                \caseOf{$1 = s$.}
                    Thus $1 < 1$ by \cref{real_lt_subst_left}.
                    Contradiction by \cref{real_lt_irreflexive}.
                \end{byCase}
            \end{subproof}
            Contradiction.
        \end{subproof}
    \caseOf{there exists $m\in\naturalsPlus$ such that $t = m + 1$.}
        Take $m\in\naturalsPlus$ such that $t = m + 1$.
        Thus $s < m + 1$.
        Thus $s \leq m + 1$.
        Thus $s \neq m + 1$.
        Thus $s \leq m$ by \cref{real_naturals_leq_succ_elim}.
        $s < m$ or $s = m$.
        \begin{byCase}
        \caseOf{$s = m$.}
            Thus $t = s + 1$.
            Thus $s \rmul (s + 1) < (s + 1) \rmul (s + 1)$ by \cref{naturalsplus_mul_succ_lt_square}.
            Thus $s \rmul (s + 1) < t \rmul t$.
        \caseOf{$s < m$.}
            $s\in\reals$ by \cref{real_naturals_subset_reals,subseteq}.
            $t\in\reals$ by \cref{real_naturals_subset_reals,subseteq}.
            $s + 1\in\naturalsPlus$ by \cref{real_naturals_closed_succ}.
            $s + 1\in\reals$ by \cref{real_naturals_subset_reals,real_naturals_closed_succ,subseteq}.
            $1\in\reals$ by \cref{real_naturals_subset_reals,real_one_in_naturals,subseteq}.
            $0\in\reals$ by \cref{real_add_ident}.
            $0 < 1$ by \cref{real_zero_lt_one}.
            $1 < s + 1$ or $1 = s + 1$ by \cref{real_one_leq_naturalsplus}.
            Show that $0 < s + 1$.
            \begin{subproof}
                \begin{byCase}
                \caseOf{$1 < s + 1$.}
                    Thus $0 < s + 1$ by \cref{real_lt_trans}.
                \caseOf{$1 = s + 1$.}
                    Thus $0 < s + 1$ by \cref{real_lt_subst_right}.
                \end{byCase}
            \end{subproof}
            Thus $s \rmul (s + 1) < t \rmul (s + 1)$ by \cref{real_lt_mul_pos}.
            $1\in\reals$ by \cref{real_naturals_subset_reals,real_one_in_naturals,subseteq}.
            $0\in\reals$ by \cref{real_add_ident}.
            $0 < 1$ by \cref{real_zero_lt_one}.
            $m\in\reals$ by \cref{real_naturals_subset_reals,subseteq}.
            $s + 1 < m + 1$ by \cref{real_lt_add}.
            Thus $s + 1 < t$.
            $1 < t$ or $1 = t$ by \cref{real_one_leq_naturalsplus}.
            Show that $0 < t$.
            \begin{subproof}
                \begin{byCase}
                \caseOf{$1 < t$.}
                    $0 < t$ by \cref{real_lt_trans,real_zero_lt_one}.
                \caseOf{$1 = t$.}
                    $0 < t$ by \cref{real_zero_lt_one}.
                \end{byCase}
            \end{subproof}
            Thus $(s + 1) \rmul t < t \rmul t$ by \cref{real_lt_mul_pos}.
            $(s + 1) \rmul t = t \rmul (s + 1)$ by \cref{real_mul_comm}.
            Thus $t \rmul (s + 1) < t \rmul t$.
            $s \rmul (s + 1)\in\reals$ by \cref{real_mul_closed}.
            $t \rmul (s + 1)\in\reals$ by \cref{real_mul_closed}.
            $t \rmul t\in\reals$ by \cref{real_mul_closed}.
            Thus $s \rmul (s + 1) < t \rmul t$ by \cref{real_lt_trans}.
        \end{byCase}
    \end{byCase}
\end{proof}
% from §30 helper lemma 7: injective pairing of naturalsPlus
\begin{lemma}\label{naturalsplus_pairing_injection}
    There exists $f\in\inj{\naturalsPlus\times\naturalsPlus}{\naturalsPlus}$.
\end{lemma}
\begin{proof}
    Let $g = \{((i,j), (i + j) \rmul (i + j) + i) \mid i\in\naturalsPlus \mid j\in\naturalsPlus\}$.
    Show that $g\in\inj{\naturalsPlus\times\naturalsPlus}{\naturalsPlus}$.
    \begin{subproof}
        Show that $g\in\funs{\naturalsPlus\times\naturalsPlus}{\naturalsPlus}$.
        \begin{subproof}
            Show that $g$ is a function.
            \begin{subproof}
                Show that for all $p,a,b$ such that $(p,a)\in g$ and $(p,b)\in g$ we have $a = b$.
                \begin{subproof}
                Fix $p$.
                Fix $a$.
                Fix $b$.
                Assume $(p,a)\in g$ and $(p,b)\in g$.
                Thus there exist $i,j$ such that $i\in\naturalsPlus$ and $j\in\naturalsPlus$ and
                    $p = (i,j)$ and $a = (i + j) \rmul (i + j) + i$ by \cref{pair_eq_iff}.
                Take $i,j\in\naturalsPlus$ such that $p = (i,j)$ and $a = (i + j) \rmul (i + j) + i$.
                Thus there exist $i_1,j_1$ such that $i_1\in\naturalsPlus$ and $j_1\in\naturalsPlus$ and
                    $p = (i_1,j_1)$ and $b = (i_1 + j_1) \rmul (i_1 + j_1) + i_1$ by \cref{pair_eq_iff}.
                Take $i_1,j_1\in\naturalsPlus$ such that $p = (i_1,j_1)$ and $b = (i_1 + j_1) \rmul (i_1 + j_1) + i_1$.
                Thus $(i,j) = (i_1,j_1)$ by \cref{pair_eq_iff}.
                Thus $i = i_1$ and $j = j_1$ by \cref{pair_eq_iff}.
                Thus $a = b$.
                \end{subproof}
                Thus $g$ is right-unique by \cref{rightunique}.
                Show that $g$ is a relation.
                \begin{subproof}
                    Show that for all $w$ such that $w\in g$ there exist $p,a$ such that $w = (p,a)$.
                    \begin{subproof}
                        Fix $w$.
                        Assume $w\in g$.
                        Thus there exist $p,a$ such that $w = (p,a)$ by \cref{pair_eq_iff}.
                    \end{subproof}
                    Thus $g$ is a relation by \cref{relation}.
                \end{subproof}
                Thus $g$ is a function.
            \end{subproof}
            Show that for all $p\in\dom{g}$ we have $g(p)\in\naturalsPlus$.
            \begin{subproof}
                Fix $p\in\dom{g}$.
                Take $a$ such that $(p,a)\in g$ by \cref{dom_iff}.
                Thus there exist $i,j$ such that $i\in\naturalsPlus$ and $j\in\naturalsPlus$ and
                    $p = (i,j)$ and $a = (i + j) \rmul (i + j) + i$ by \cref{pair_eq_iff}.
                Take $i,j\in\naturalsPlus$ such that $p = (i,j)$ and $a = (i + j) \rmul (i + j) + i$.
                Thus $i + j\in\naturalsPlus$ by \cref{naturalsplus_add_closed}.
                Thus $(i + j) \rmul (i + j)\in\naturalsPlus$ by \cref{naturalsplus_mul_closed}.
                Thus $a\in\naturalsPlus$ by \cref{naturalsplus_add_closed}.
                Thus $g(p)\in\naturalsPlus$ by \cref{function_apply_iff}.
            \end{subproof}
            Show that $\dom{g} = \naturalsPlus\times\naturalsPlus$.
            \begin{subproof}
                Show that $\dom{g}\subseteq\naturalsPlus\times\naturalsPlus$.
                \begin{subproof}
                    Show that for all $p$ such that $p\in\dom{g}$ we have $p\in\naturalsPlus\times\naturalsPlus$.
                    \begin{subproof}
                        Fix $p$.
                        Assume $p\in\dom{g}$.
                        Take $a$ such that $(p,a)\in g$ by \cref{dom_iff}.
                        Thus there exist $i,j$ such that $i\in\naturalsPlus$ and $j\in\naturalsPlus$ and
                            $p = (i,j)$ and $a = (i + j) \rmul (i + j) + i$ by \cref{pair_eq_iff}.
                        Take $i,j\in\naturalsPlus$ such that $p = (i,j)$ and $a = (i + j) \rmul (i + j) + i$.
                        Thus $p\in\naturalsPlus\times\naturalsPlus$ by \cref{times_tuple_intro}.
                    \end{subproof}
                    Thus $\dom{g}\subseteq\naturalsPlus\times\naturalsPlus$ by \cref{subseteq}.
                \end{subproof}
                Show that $\naturalsPlus\times\naturalsPlus\subseteq\dom{g}$.
                \begin{subproof}
                    Show that for all $p$ such that $p\in\naturalsPlus\times\naturalsPlus$ we have $p\in\dom{g}$.
                    \begin{subproof}
                        Fix $p$.
                        Assume $p\in\naturalsPlus\times\naturalsPlus$.
                        Take $i,j$ such that $i\in\naturalsPlus$ and $j\in\naturalsPlus$ and $p = (i,j)$ by \cref{times}.
                        $(p,(i + j) \rmul (i + j) + i)\in g$ by \cref{pair_eq_iff}.
                        Thus $p\in\dom{g}$ by \cref{dom_intro}.
                    \end{subproof}
                    Thus $\naturalsPlus\times\naturalsPlus\subseteq\dom{g}$ by \cref{subseteq}.
                \end{subproof}
                Thus $\dom{g} = \naturalsPlus\times\naturalsPlus$ by \cref{subseteq_antisymmetric}.
            \end{subproof}
            Thus $g\in\funs{\naturalsPlus\times\naturalsPlus}{\naturalsPlus}$ by \cref{funs_intro}.
        \end{subproof}
        Show that for all $p_1,p_2$ such that $p_1\in\naturalsPlus\times\naturalsPlus$ and $p_2\in\naturalsPlus\times\naturalsPlus$ and $g(p_1) = g(p_2)$ we have $p_1 = p_2$.
        \begin{subproof}
            Fix $p_1$.
            Fix $p_2$.
            Assume $p_1\in\naturalsPlus\times\naturalsPlus$ and $p_2\in\naturalsPlus\times\naturalsPlus$ and $g(p_1) = g(p_2)$.
            $g\in\funs{\naturalsPlus\times\naturalsPlus}{\naturalsPlus}$.
            Thus $\dom{g} = \naturalsPlus\times\naturalsPlus$ by \cref{funs_elim}.
            Thus $p_1\in\dom{g}$ and $p_2\in\dom{g}$ by \cref{subseteq}.
            $g$ is a function by \cref{funs_elim}.
            Thus $(p_1,g(p_1))\in g$ and $(p_2,g(p_2))\in g$ by \cref{function_apply_iff}.
            Thus there exist $i_1,j_1$ such that $i_1\in\naturalsPlus$ and $j_1\in\naturalsPlus$ and
                $p_1 = (i_1,j_1)$ and $g(p_1) = (i_1 + j_1) \rmul (i_1 + j_1) + i_1$ by \cref{pair_eq_iff}.
            Take $i_1,j_1\in\naturalsPlus$ such that $p_1 = (i_1,j_1)$ and $g(p_1) = (i_1 + j_1) \rmul (i_1 + j_1) + i_1$.
            Thus there exist $i_2,j_2$ such that $i_2\in\naturalsPlus$ and $j_2\in\naturalsPlus$ and
                $p_2 = (i_2,j_2)$ and $g(p_2) = (i_2 + j_2) \rmul (i_2 + j_2) + i_2$ by \cref{pair_eq_iff}.
            Take $i_2,j_2\in\naturalsPlus$ such that $p_2 = (i_2,j_2)$ and $g(p_2) = (i_2 + j_2) \rmul (i_2 + j_2) + i_2$.
            Thus $(i_1 + j_1) \rmul (i_1 + j_1) + i_1 = (i_2 + j_2) \rmul (i_2 + j_2) + i_2$.
            Let $s = i_1 + j_1$.
            Let $t = i_2 + j_2$.
            Thus $s \rmul s + i_1 = t \rmul t + i_2$.
            \begin{byCase}
            \caseOf{$s = t$.}
            Thus $s\in\naturalsPlus$ by \cref{naturalsplus_add_closed}.
            Thus $t\in\naturalsPlus$ by \cref{naturalsplus_add_closed}.
            $i_1\in\reals$ by \cref{real_naturals_subset_reals,subseteq}.
            $j_1\in\reals$ by \cref{real_naturals_subset_reals,subseteq}.
            $i_2\in\reals$ by \cref{real_naturals_subset_reals,subseteq}.
            $j_2\in\reals$ by \cref{real_naturals_subset_reals,subseteq}.
            $s\in\reals$ by \cref{real_naturals_subset_reals,subseteq}.
            $t\in\reals$ by \cref{real_naturals_subset_reals,subseteq}.
            $t \rmul t\in\reals$ by \cref{real_mul_closed}.
            Thus $t \rmul t + i_1 = t \rmul t + i_2$.
            $t \rmul t + i_1 = i_1 + t \rmul t$ by \cref{real_add_comm}.
            $t \rmul t + i_2 = i_2 + t \rmul t$ by \cref{real_add_comm}.
            Thus $i_1 + t \rmul t = i_2 + t \rmul t$.
            Thus $i_1 = i_2$ by \cref{real_add_cancel}.
            Thus $i_1 + j_1 = i_2 + j_2$.
            Thus $i_1 + j_1 = i_1 + j_2$.
            $i_1 + j_1 = j_1 + i_1$ by \cref{real_add_comm}.
            $i_1 + j_2 = j_2 + i_1$ by \cref{real_add_comm}.
            Thus $j_1 + i_1 = j_2 + i_1$.
            Thus $j_1 = j_2$ by \cref{real_add_cancel}.
            Thus $p_1 = p_2$ by \cref{pair_eq_iff}.
            \caseOf{$s < t$.}
                Suppose not.
                $j_1\in\naturalsPlus$.
                $j_1\in\reals$ by \cref{real_naturals_subset_reals,subseteq}.
                $1\in\reals$ by \cref{real_naturals_subset_reals,real_one_in_naturals,subseteq}.
                $0\in\reals$ by \cref{real_add_ident}.
                $0 < 1$ by \cref{real_zero_lt_one}.
                $0 + j_1 < 1 + j_1$ by \cref{real_lt_add}.
                $0 + j_1 = j_1$ by \cref{real_add_ident}.
                $1 + j_1 = j_1 + 1$ by \cref{real_add_comm}.
                Thus $0 < j_1$.
                $i_1\in\reals$ by \cref{real_naturals_subset_reals,subseteq}.
                Thus $s\in\naturalsPlus$ by \cref{naturalsplus_add_closed}.
                Thus $s\in\reals$ by \cref{real_naturals_subset_reals,subseteq}.
                Thus $0 + i_1 < j_1 + i_1$ by \cref{real_lt_add}.
                $0 + i_1 = i_1$ by \cref{real_add_ident}.
                $j_1 + i_1 = i_1 + j_1$ by \cref{real_add_comm}.
                Thus $i_1 < i_1 + j_1$.
                Thus $i_1 < s$.
                $s \rmul s\in\reals$ by \cref{real_mul_closed}.
                $i_1 + (s \rmul s) < s + (s \rmul s)$ by \cref{real_lt_add}.
                $i_1 + (s \rmul s) = s \rmul s + i_1$ by \cref{real_add_comm}.
                $s + (s \rmul s) = s \rmul s + s$ by \cref{real_add_comm}.
                Thus $s \rmul s + i_1 < s \rmul s + s$.
                $s \rmul (s + 1) = s \rmul s + (s \rmul 1)$ by \cref{real_distrib}.
                $s \rmul 1 = s$ by \cref{real_mul_comm,real_mul_ident}.
                Thus $s \rmul s + s = s \rmul (s + 1)$.
                Thus $s \rmul s + i_1 < s \rmul (s + 1)$.
                $s \rmul s + i_1\in\reals$ by \cref{real_add_closed}.
                $s + 1\in\reals$ by \cref{real_add_closed}.
                $s \rmul (s + 1)\in\reals$ by \cref{real_mul_closed}.
                Thus $t\in\naturalsPlus$ by \cref{naturalsplus_add_closed}.
                $t\in\reals$ by \cref{real_naturals_subset_reals,subseteq}.
                $t \rmul t\in\reals$ by \cref{real_mul_closed}.
                Thus $s \rmul (s + 1) < t \rmul t$ by \cref{naturalsplus_mul_succ_lt_square_bound}.
                Thus $s \rmul s + i_1 < t \rmul t$ by \cref{real_lt_trans}.
                $i_2\in\naturalsPlus$.
                $i_2\in\reals$ by \cref{real_naturals_subset_reals,subseteq}.
                $0 + i_2 < 1 + i_2$ by \cref{real_lt_add,real_zero_lt_one,real_add_ident,real_add_comm}.
                Thus $0 < i_2$.
                $0\in\reals$ by \cref{real_add_ident}.
                $0 + t \rmul t < i_2 + t \rmul t$ by \cref{real_lt_add}.
                $0 + t \rmul t = t \rmul t$ by \cref{real_add_ident}.
                Thus $t \rmul t < i_2 + t \rmul t$ by \cref{real_lt_subst_left}.
                $i_2 + t \rmul t = t \rmul t + i_2$ by \cref{real_add_comm}.
                Thus $t \rmul t < t \rmul t + i_2$ by \cref{real_lt_subst_right}.
                Thus $s \rmul s + i_1 < t \rmul t + i_2$ by \cref{real_lt_trans}.
                Thus it is wrong that $s \rmul s + i_1 = t \rmul t + i_2$ by \cref{real_lt_irreflexive}.
                Contradiction.
            \caseOf{$t < s$.}
                Suppose not.
                $j_2\in\naturalsPlus$.
                $j_2\in\reals$ by \cref{real_naturals_subset_reals,subseteq}.
                $1\in\reals$ by \cref{real_naturals_subset_reals,real_one_in_naturals,subseteq}.
                $0\in\reals$ by \cref{real_add_ident}.
                $0 < 1$ by \cref{real_zero_lt_one}.
                $0 + j_2 < 1 + j_2$ by \cref{real_lt_add}.
                $0 + j_2 = j_2$ by \cref{real_add_ident}.
                $1 + j_2 = j_2 + 1$ by \cref{real_add_comm}.
                Thus $0 < j_2$.
                $i_2\in\reals$ by \cref{real_naturals_subset_reals,subseteq}.
                Thus $0 + i_2 < j_2 + i_2$ by \cref{real_lt_add}.
                $0 + i_2 = i_2$ by \cref{real_add_ident}.
                $j_2 + i_2 = i_2 + j_2$ by \cref{real_add_comm}.
                Thus $i_2 < i_2 + j_2$.
                Thus $i_2 < t$.
                Thus $t\in\naturalsPlus$ by \cref{naturalsplus_add_closed}.
                Thus $t\in\reals$ by \cref{real_naturals_subset_reals,subseteq}.
                $t \rmul t\in\reals$ by \cref{real_mul_closed}.
                $i_2 + (t \rmul t) < t + (t \rmul t)$ by \cref{real_lt_add}.
                $i_2 + (t \rmul t) = t \rmul t + i_2$ by \cref{real_add_comm}.
                $t + (t \rmul t) = t \rmul t + t$ by \cref{real_add_comm}.
                Thus $t \rmul t + i_2 < t \rmul t + t$.
                $t \rmul (t + 1) = t \rmul t + (t \rmul 1)$ by \cref{real_distrib}.
                $t \rmul 1 = t$ by \cref{real_mul_comm,real_mul_ident}.
                Thus $t \rmul t + t = t \rmul (t + 1)$.
                Thus $t \rmul t + i_2 < t \rmul (t + 1)$.
                $t \rmul t + i_2\in\reals$ by \cref{real_add_closed}.
                $t + 1\in\reals$ by \cref{real_add_closed}.
                $t \rmul (t + 1)\in\reals$ by \cref{real_mul_closed}.
                Thus $s\in\naturalsPlus$ by \cref{naturalsplus_add_closed}.
                Thus $s\in\reals$ by \cref{real_naturals_subset_reals,subseteq}.
                $s \rmul s\in\reals$ by \cref{real_mul_closed}.
                Thus $t \rmul (t + 1) < s \rmul s$ by \cref{naturalsplus_mul_succ_lt_square_bound}.
                Thus $t \rmul t + i_2 < s \rmul s$ by \cref{real_lt_trans}.
                $i_1\in\naturalsPlus$.
                $i_1\in\reals$ by \cref{real_naturals_subset_reals,subseteq}.
                $0 + i_1 < 1 + i_1$ by \cref{real_lt_add,real_zero_lt_one,real_add_ident,real_add_comm}.
                Thus $0 < i_1$.
                $0\in\reals$ by \cref{real_add_ident}.
                $0 + s \rmul s < i_1 + s \rmul s$ by \cref{real_lt_add}.
                $0 + s \rmul s = s \rmul s$ by \cref{real_add_ident}.
                Thus $s \rmul s < i_1 + s \rmul s$ by \cref{real_lt_subst_left}.
                $i_1 + s \rmul s = s \rmul s + i_1$ by \cref{real_add_comm}.
                Thus $s \rmul s < s \rmul s + i_1$ by \cref{real_lt_subst_right}.
                Thus $t \rmul t + i_2 < s \rmul s + i_1$ by \cref{real_lt_trans}.
                Thus it is wrong that $s \rmul s + i_1 = t \rmul t + i_2$ by \cref{real_lt_irreflexive}.
                Contradiction.
            \end{byCase}
        \end{subproof}
        Thus $g\in\inj{\naturalsPlus\times\naturalsPlus}{\naturalsPlus}$ by \cref{inj}.
    \end{subproof}
    Thus there exists $f\in\inj{\naturalsPlus\times\naturalsPlus}{\naturalsPlus}$.
\end{proof}

% from §30 helper lemma 8: function equality by pointwise equality
\begin{lemma}\label{functions_eq_iff}
    Suppose $f\in\funs{X}{Y}$ and $g\in\funs{X}{Y}$.
    Suppose for all $x\in X$ we have $f(x) = g(x)$.
    Then $f = g$.
\end{lemma}
\begin{proof}
    $f$ is a function and $\dom{f} = X$ by \cref{funs_elim}.
    $g$ is a function and $\dom{g} = X$ by \cref{funs_elim}.
    Thus for all $x\in\dom{f}$ we have $f(x) = g(x)$.
    Thus $f\subseteq g$ by \cref{fun_subseteq,subseteq_refl}.
    Thus for all $x\in\dom{g}$ we have $f(x) = g(x)$.
    Thus $g\subseteq f$ by \cref{fun_subseteq,subseteq_refl}.
    Thus $f = g$ by \cref{subseteq_antisymmetric}.
\end{proof}

% from §30 helper lemma 9: injective coding of finite sequences of naturalsPlus
\begin{lemma}\label{naturalsplus_finite_sequences_inj}
    Suppose $n\in\naturalsPlus$.
    Then there exists $f\in\inj{\funs{\initialSegment{n}}{\naturalsPlus}}{\naturalsPlus}$.
\end{lemma}
\begin{proof}
    Let $A = \{n\in\naturalsPlus \mid \exists f\in\inj{\funs{\initialSegment{n}}{\naturalsPlus}}{\naturalsPlus}. f = f\}$.
    Show that for all $n\in\naturalsPlus$ we have $n\in A$.
    \begin{subproof}
        Show that $A\subseteq\naturalsPlus$.
        \begin{subproof}
            Show that for all $n$ such that $n\in A$ we have $n\in\naturalsPlus$.
            \begin{subproof}
                Fix $n$.
                Assume $n\in A$.
                Thus $n\in\naturalsPlus$.
            \end{subproof}
            Thus $A\subseteq\naturalsPlus$ by \cref{subseteq}.
        \end{subproof}
        $1\in\naturalsPlus$ by \cref{real_one_in_naturals}.
        Let $f_1 = \{(u, \apply{u}{1}) \mid u\in\funs{\initialSegment{1}}{\naturalsPlus}\}$.
        Show that $f_1\in\inj{\funs{\initialSegment{1}}{\naturalsPlus}}{\naturalsPlus}$.
        \begin{subproof}
            Show that $f_1\in\funs{\funs{\initialSegment{1}}{\naturalsPlus}}{\naturalsPlus}$.
            \begin{subproof}
                Show that $f_1$ is a function.
                \begin{subproof}
                    Show that for all $u,a,b$ such that $(u,a)\in f_1$ and $(u,b)\in f_1$ we have $a = b$.
                    \begin{subproof}
                        Fix $u$.
                        Fix $a$.
                        Fix $b$.
                        Assume $(u,a)\in f_1$ and $(u,b)\in f_1$.
                        Thus $a = u(1)$ and $b = u(1)$ by \cref{pair_eq_iff}.
                        Thus $a = b$.
                    \end{subproof}
                    Thus $f_1$ is right-unique by \cref{rightunique}.
                    Show that $f_1$ is a relation.
                    \begin{subproof}
                        Show that for all $w$ such that $w\in f_1$ there exist $u,a$ such that $w = (u,a)$.
                        \begin{subproof}
                            Fix $w$.
                            Assume $w\in f_1$.
                            Thus there exist $u,a$ such that $w = (u,a)$ by \cref{pair_eq_iff}.
                        \end{subproof}
                        Thus $f_1$ is a relation by \cref{relation}.
                    \end{subproof}
                    Thus $f_1$ is a function.
                \end{subproof}
                Show that $\dom{f_1} = \funs{\initialSegment{1}}{\naturalsPlus}$.
                \begin{subproof}
                    Show that $\dom{f_1}\subseteq\funs{\initialSegment{1}}{\naturalsPlus}$.
                    \begin{subproof}
                        Show that for all $u$ such that $u\in\dom{f_1}$ we have $u\in\funs{\initialSegment{1}}{\naturalsPlus}$.
                        \begin{subproof}
                            Fix $u$.
                            Assume $u\in\dom{f_1}$.
                            Take $a$ such that $(u,a)\in f_1$ by \cref{dom_iff}.
                            Thus there exists $u_0$ such that $u_0\in\funs{\initialSegment{1}}{\naturalsPlus}$ and
                                $(u,a) = (u_0,\apply{u_0}{1})$ by \cref{pair_eq_iff}.
                            Take $u_0\in\funs{\initialSegment{1}}{\naturalsPlus}$ such that $(u,a) = (u_0,\apply{u_0}{1})$.
                            Thus $u = u_0$ by \cref{pair_eq_iff}.
                            Thus $u\in\funs{\initialSegment{1}}{\naturalsPlus}$.
                        \end{subproof}
                        Thus $\dom{f_1}\subseteq\funs{\initialSegment{1}}{\naturalsPlus}$ by \cref{subseteq}.
                    \end{subproof}
                    Show that $\funs{\initialSegment{1}}{\naturalsPlus}\subseteq\dom{f_1}$.
                    \begin{subproof}
                        Show that for all $u$ such that $u\in\funs{\initialSegment{1}}{\naturalsPlus}$ we have $u\in\dom{f_1}$.
                        \begin{subproof}
                            Fix $u$.
                            Assume $u\in\funs{\initialSegment{1}}{\naturalsPlus}$.
                            $(u,\apply{u}{1})\in f_1$ by \cref{pair_eq_iff}.
                            Thus $u\in\dom{f_1}$ by \cref{dom_intro}.
                        \end{subproof}
                        Thus $\funs{\initialSegment{1}}{\naturalsPlus}\subseteq\dom{f_1}$ by \cref{subseteq}.
                    \end{subproof}
                    Thus $\dom{f_1} = \funs{\initialSegment{1}}{\naturalsPlus}$ by \cref{subseteq_antisymmetric}.
                \end{subproof}
                Show that for all $u\in\dom{f_1}$ we have $f_1(u)\in\naturalsPlus$.
                \begin{subproof}
                    Fix $u\in\dom{f_1}$.
                    Take $a$ such that $(u,a)\in f_1$ by \cref{dom_iff}.
                    Thus there exists $u_0$ such that $u_0\in\funs{\initialSegment{1}}{\naturalsPlus}$ and
                        $(u,a) = (u_0,\apply{u_0}{1})$ by \cref{pair_eq_iff}.
                    Take $u_0\in\funs{\initialSegment{1}}{\naturalsPlus}$ such that $(u,a) = (u_0,\apply{u_0}{1})$.
                    Thus $u = u_0$ and $a = u_0(1)$ by \cref{pair_eq_iff}.
                    Thus $u\in\funs{\initialSegment{1}}{\naturalsPlus}$.
                    Thus $u$ is a function to $\naturalsPlus$ and $\dom{u} = \initialSegment{1}$ by \cref{funs_elim}.
                    $1\in\initialSegment{1}$ by \cref{initial_segment_naturalsplus,real_one_in_naturals}.
                    Thus $1\in\dom{u}$ by \cref{subseteq}.
                    $(1,\apply{u}{1})\in u$ by \cref{function_apply_iff}.
                    Thus $\apply{u}{1}\in\ran{u}$ by \cref{ran_intro}.
                    $\ran{u}\subseteq\naturalsPlus$ by \cref{function_to_ran}.
                    Thus $\apply{u}{1}\in\naturalsPlus$ by \cref{subseteq}.
                    Thus $f_1(u)\in\naturalsPlus$ by \cref{function_apply_iff}.
                \end{subproof}
                Thus $f_1\in\funs{\funs{\initialSegment{1}}{\naturalsPlus}}{\naturalsPlus}$ by \cref{funs_intro}.
            \end{subproof}
            Show that for all $u_1,u_2$ such that $u_1\in\funs{\initialSegment{1}}{\naturalsPlus}$ and $u_2\in\funs{\initialSegment{1}}{\naturalsPlus}$ and $f_1(u_1) = f_1(u_2)$ we have $u_1 = u_2$.
            \begin{subproof}
                Fix $u_1$.
                Fix $u_2$.
                Assume $u_1\in\funs{\initialSegment{1}}{\naturalsPlus}$ and $u_2\in\funs{\initialSegment{1}}{\naturalsPlus}$ and $f_1(u_1) = f_1(u_2)$.
                $(u_1,\apply{u_1}{1})\in f_1$ by \cref{pair_eq_iff}.
                $(u_2,\apply{u_2}{1})\in f_1$ by \cref{pair_eq_iff}.
                $f_1$ is a function by \cref{funs_elim}.
                Thus $f_1(u_1) = u_1(1)$ and $f_1(u_2) = u_2(1)$ by \cref{function_apply_iff}.
                Thus $u_1(1) = u_2(1)$.
                Show that for all $k\in\initialSegment{1}$ we have $u_1(k) = u_2(k)$.
                \begin{subproof}
                    Fix $k\in\initialSegment{1}$.
                    Thus $k\in\naturalsPlus$ and $k \leq 1$ by \cref{initial_segment_naturalsplus}.
                    $1\in\reals$ by \cref{real_naturals_subset_reals,real_one_in_naturals,subseteq}.
                    $k\in\reals$ by \cref{real_naturals_subset_reals,subseteq}.
                    Thus $1 \leq k$ by \cref{real_one_leq_naturalsplus}.
                    Thus $k = 1$ by \cref{real_leq_antisym}.
                    Thus $u_1(k) = u_2(k)$.
                \end{subproof}
                Thus $u_1 = u_2$ by \cref{functions_eq_iff}.
            \end{subproof}
            Thus $f_1\in\inj{\funs{\initialSegment{1}}{\naturalsPlus}}{\naturalsPlus}$ by \cref{inj}.
        \end{subproof}
        Thus $1\in A$.
        Show that for all $n\in A$ we have $n + 1\in A$.
        \begin{subproof}
            Fix $n\in A$.
            Thus there exists $f\in\inj{\funs{\initialSegment{n}}{\naturalsPlus}}{\naturalsPlus}$ such that $f = f$.
            Take $f\in\inj{\funs{\initialSegment{n}}{\naturalsPlus}}{\naturalsPlus}$ such that $f = f$.
            Thus there exists $g\in\inj{\naturalsPlus\times\naturalsPlus}{\naturalsPlus}$ such that $g = g$ by \cref{naturalsplus_pairing_injection}.
            Take $g\in\inj{\naturalsPlus\times\naturalsPlus}{\naturalsPlus}$ such that $g = g$.
            Let $F = \{(u, \apply{g}{(\apply{f}{\restrl{u}{\initialSegment{n}}}, \apply{u}{n + 1})}) \mid u\in\funs{\initialSegment{n + 1}}{\naturalsPlus}\}$.
            Show that $F\in\inj{\funs{\initialSegment{n + 1}}{\naturalsPlus}}{\naturalsPlus}$.
            \begin{subproof}
                Show that $F\in\funs{\funs{\initialSegment{n + 1}}{\naturalsPlus}}{\naturalsPlus}$.
                \begin{subproof}
                    Show that $F$ is a function.
                    \begin{subproof}
                        Show that for all $u,a,b$ such that $(u,a)\in F$ and $(u,b)\in F$ we have $a = b$.
                        \begin{subproof}
                            Fix $u$.
                            Fix $a$.
                            Fix $b$.
                            Assume $(u,a)\in F$ and $(u,b)\in F$.
                            Thus $a = \apply{g}{(\apply{f}{\restrl{u}{\initialSegment{n}}}, \apply{u}{n + 1})}$ and
                                $b = \apply{g}{(\apply{f}{\restrl{u}{\initialSegment{n}}}, \apply{u}{n + 1})}$ by \cref{pair_eq_iff}.
                            Thus $a = b$.
                        \end{subproof}
                        Thus $F$ is right-unique by \cref{rightunique}.
                        Show that $F$ is a relation.
                        \begin{subproof}
                            Show that for all $w$ such that $w\in F$ there exist $u,a$ such that $w = (u,a)$.
                            \begin{subproof}
                                Fix $w$.
                                Assume $w\in F$.
                                Thus there exist $u,a$ such that $w = (u,a)$ by \cref{pair_eq_iff}.
                            \end{subproof}
                            Thus $F$ is a relation by \cref{relation}.
                        \end{subproof}
                        Thus $F$ is a function.
                    \end{subproof}
                    Show that $\dom{F} = \funs{\initialSegment{n + 1}}{\naturalsPlus}$.
                    \begin{subproof}
                        Show that $\dom{F}\subseteq\funs{\initialSegment{n + 1}}{\naturalsPlus}$.
                        \begin{subproof}
                            Show that for all $u$ such that $u\in\dom{F}$ we have $u\in\funs{\initialSegment{n + 1}}{\naturalsPlus}$.
                            \begin{subproof}
                                Fix $u$.
                                Assume $u\in\dom{F}$.
                                Take $a$ such that $(u,a)\in F$ by \cref{dom_iff}.
                                Thus there exists $u_0$ such that $u_0\in\funs{\initialSegment{n + 1}}{\naturalsPlus}$ and
                                    $(u,a) = (u_0,\apply{g}{(\apply{f}{\restrl{u_0}{\initialSegment{n}}}, \apply{u_0}{n + 1})})$ by \cref{pair_eq_iff}.
                                Take $u_0\in\funs{\initialSegment{n + 1}}{\naturalsPlus}$ such that
                                    $(u,a) = (u_0,\apply{g}{(\apply{f}{\restrl{u_0}{\initialSegment{n}}}, \apply{u_0}{n + 1})})$.
                                Thus $u = u_0$ by \cref{pair_eq_iff}.
                                Thus $u\in\funs{\initialSegment{n + 1}}{\naturalsPlus}$.
                            \end{subproof}
                            Thus $\dom{F}\subseteq\funs{\initialSegment{n + 1}}{\naturalsPlus}$ by \cref{subseteq}.
                        \end{subproof}
                        Show that $\funs{\initialSegment{n + 1}}{\naturalsPlus}\subseteq\dom{F}$.
                        \begin{subproof}
                            Show that for all $u$ such that $u\in\funs{\initialSegment{n + 1}}{\naturalsPlus}$ we have $u\in\dom{F}$.
                            \begin{subproof}
                                Fix $u$.
                                Assume $u\in\funs{\initialSegment{n + 1}}{\naturalsPlus}$.
                                $(u,\apply{g}{(\apply{f}{\restrl{u}{\initialSegment{n}}}, \apply{u}{n + 1})})\in F$ by \cref{pair_eq_iff}.
                                Thus $u\in\dom{F}$ by \cref{dom_intro}.
                            \end{subproof}
                            Thus $\funs{\initialSegment{n + 1}}{\naturalsPlus}\subseteq\dom{F}$ by \cref{subseteq}.
                        \end{subproof}
                        Thus $\dom{F} = \funs{\initialSegment{n + 1}}{\naturalsPlus}$ by \cref{subseteq_antisymmetric}.
                    \end{subproof}
                    Show that for all $u\in\dom{F}$ we have $F(u)\in\naturalsPlus$.
                    \begin{subproof}
                        Fix $u\in\dom{F}$.
                        Take $a$ such that $(u,a)\in F$ by \cref{dom_iff}.
                        Thus there exists $u_0$ such that $u_0\in\funs{\initialSegment{n + 1}}{\naturalsPlus}$ and
                            $(u,a) = (u_0,\apply{g}{(\apply{f}{\restrl{u_0}{\initialSegment{n}}}, \apply{u_0}{n + 1})})$ by \cref{pair_eq_iff}.
                        Take $u_0\in\funs{\initialSegment{n + 1}}{\naturalsPlus}$ such that
                            $(u,a) = (u_0,\apply{g}{(\apply{f}{\restrl{u_0}{\initialSegment{n}}}, \apply{u_0}{n + 1})})$.
                        Thus $u = u_0$ and $a = \apply{g}{(\apply{f}{\restrl{u_0}{\initialSegment{n}}}, \apply{u_0}{n + 1})}$ by \cref{pair_eq_iff}.
                        Thus $u\in\funs{\initialSegment{n + 1}}{\naturalsPlus}$.
                        Thus $u$ is a function to $\naturalsPlus$ and $\dom{u} = \initialSegment{n + 1}$ by \cref{funs_elim}.
                        $\restrl{u}{\initialSegment{n}}$ is a function by \cref{restrl_of_function_is_function}.
                        $\dom{\restrl{u}{\initialSegment{n}}} = \dom{u}\inter \initialSegment{n}$ by \cref{dom_of_restrl_eq_inter}.
                        Show that $\dom{\restrl{u}{\initialSegment{n}}} = \initialSegment{n}$.
                        \begin{subproof}
                            Show that $\dom{\restrl{u}{\initialSegment{n}}}\subseteq\initialSegment{n}$.
                            \begin{subproof}
                                Thus $\dom{\restrl{u}{\initialSegment{n}}}\subseteq\initialSegment{n}$ by \cref{restrl_dom}.
                            \end{subproof}
                            Show that $\initialSegment{n}\subseteq\dom{\restrl{u}{\initialSegment{n}}}$.
                            \begin{subproof}
                                Show that for all $k$ such that $k\in\initialSegment{n}$ we have $k\in\dom{\restrl{u}{\initialSegment{n}}}$.
                                \begin{subproof}
                                    Fix $k$.
                                    Assume $k\in\initialSegment{n}$.
                                    Thus $k\in\naturalsPlus$ by \cref{initial_segment_naturalsplus}.
                                    $k\in\reals$ by \cref{real_naturals_subset_reals,subseteq}.
                                    Thus $k \leq n$.
                                    $n\in\naturalsPlus$ by \cref{subseteq}.
                                    $n\in\reals$ by \cref{real_naturals_subset_reals,subseteq}.
                                    $0\in\reals$ by \cref{real_add_ident}.
                                    $1\in\reals$ by \cref{real_naturals_subset_reals,real_one_in_naturals,subseteq}.
                                    $n + 1\in\reals$ by \cref{real_add_closed}.
                                    $0 < 1$ by \cref{real_zero_lt_one}.
                                    $0 + n < 1 + n$ by \cref{real_lt_add}.
                                    $0 + n = n$ by \cref{real_add_ident}.
                                    $1 + n = n + 1$ by \cref{real_add_comm}.
                                    Thus $n < n + 1$.
                                    Thus $k \leq n + 1$ by \cref{real_leq_trans}.
                                    Thus $k\in\initialSegment{n + 1}$ by \cref{initial_segment_naturalsplus}.
                                    Thus $k\in\dom{u}$ by \cref{subseteq}.
                                    $(k,\apply{u}{k})\in u$ by \cref{function_apply_iff}.
                                    Thus $(k,\apply{u}{k})\in\restrl{u}{\initialSegment{n}}$ by \cref{restrl_iff}.
                                    Thus $k\in\dom{\restrl{u}{\initialSegment{n}}}$ by \cref{dom_intro}.
                                \end{subproof}
                                Thus $\initialSegment{n}\subseteq\dom{\restrl{u}{\initialSegment{n}}}$ by \cref{subseteq}.
                            \end{subproof}
                            Thus $\dom{\restrl{u}{\initialSegment{n}}} = \initialSegment{n}$ by \cref{subseteq_antisymmetric}.
                        \end{subproof}
                        Show that for all $k\in\dom{\restrl{u}{\initialSegment{n}}}$ we have $\apply{\restrl{u}{\initialSegment{n}}}{k}\in\naturalsPlus$.
                        \begin{subproof}
                            Fix $k\in\dom{\restrl{u}{\initialSegment{n}}}$.
                            Thus $k\in\dom{u}$ by \cref{dom_of_restrl_eq_inter,inter_elim_left}.
                            Thus $u(k)\in\naturalsPlus$.
                            Thus $k\in\initialSegment{n}$ by \cref{subseteq}.
                            $(k,\apply{u}{k})\in u$ by \cref{function_apply_iff}.
                            Thus $(k,\apply{u}{k})\in\restrl{u}{\initialSegment{n}}$ by \cref{restrl_iff}.
                            Thus $\apply{\restrl{u}{\initialSegment{n}}}{k} = u(k)$ by \cref{function_apply_iff}.
                            Thus $\apply{\restrl{u}{\initialSegment{n}}}{k}\in\naturalsPlus$.
                        \end{subproof}
                        Thus $\restrl{u}{\initialSegment{n}}\in\funs{\initialSegment{n}}{\naturalsPlus}$ by \cref{funs_intro}.
                        $f\in\funs{\funs{\initialSegment{n}}{\naturalsPlus}}{\naturalsPlus}$ by \cref{inj}.
                        Thus $\dom{f} = \funs{\initialSegment{n}}{\naturalsPlus}$ by \cref{funs_elim}.
                        Thus $\restrl{u}{\initialSegment{n}}\in\dom{f}$ by \cref{subseteq}.
                        Thus $\apply{f}{\restrl{u}{\initialSegment{n}}}\in\naturalsPlus$ by \cref{funs_elim}.
                        $n + 1\in\naturalsPlus$ by \cref{real_naturals_closed_succ}.
                        Thus $n + 1\in\initialSegment{n + 1}$ by \cref{initial_segment_naturalsplus}.
                        Thus $n + 1\in\dom{u}$ by \cref{subseteq}.
                        $(n + 1,\apply{u}{n + 1})\in u$ by \cref{function_apply_iff}.
                        Thus $\apply{u}{n + 1}\in\ran{u}$ by \cref{ran_intro}.
                        $\ran{u}\subseteq\naturalsPlus$ by \cref{function_to_ran}.
                        Thus $\apply{u}{n + 1}\in\naturalsPlus$ by \cref{subseteq}.
                        Thus $(\apply{f}{\restrl{u}{\initialSegment{n}}}, \apply{u}{n + 1})\in\naturalsPlus\times\naturalsPlus$ by \cref{times_tuple_intro}.
                        $g\in\funs{\naturalsPlus\times\naturalsPlus}{\naturalsPlus}$ by \cref{inj}.
                        Thus $\dom{g} = \naturalsPlus\times\naturalsPlus$ by \cref{funs_elim}.
                        Thus $(\apply{f}{\restrl{u}{\initialSegment{n}}}, \apply{u}{n + 1})\in\dom{g}$ by \cref{subseteq}.
                        Thus $a\in\naturalsPlus$ by \cref{funs_elim}.
                        Thus $F(u)\in\naturalsPlus$ by \cref{function_apply_iff}.
                    \end{subproof}
                    Thus $F\in\funs{\funs{\initialSegment{n + 1}}{\naturalsPlus}}{\naturalsPlus}$ by \cref{funs_intro}.
                \end{subproof}
                Show that for all $u_1,u_2$ such that $u_1\in\funs{\initialSegment{n + 1}}{\naturalsPlus}$ and $u_2\in\funs{\initialSegment{n + 1}}{\naturalsPlus}$ and $F(u_1) = F(u_2)$ we have $u_1 = u_2$.
                \begin{subproof}
                    Fix $u_1$.
                    Fix $u_2$.
                    Assume $u_1\in\funs{\initialSegment{n + 1}}{\naturalsPlus}$ and $u_2\in\funs{\initialSegment{n + 1}}{\naturalsPlus}$ and $F(u_1) = F(u_2)$.
                    Show that for all $u\in\funs{\initialSegment{n + 1}}{\naturalsPlus}$ we have $\restrl{u}{\initialSegment{n}}\in\funs{\initialSegment{n}}{\naturalsPlus}$.
                    \begin{subproof}
                        Fix $u$.
                        Assume $u\in\funs{\initialSegment{n + 1}}{\naturalsPlus}$.
                        Thus $u$ is a function and $\dom{u} = \initialSegment{n + 1}$ by \cref{funs_elim}.
                        $\restrl{u}{\initialSegment{n}}$ is a function by \cref{restrl_of_function_is_function}.
                        $\dom{\restrl{u}{\initialSegment{n}}} = \dom{u}\inter \initialSegment{n}$ by \cref{dom_of_restrl_eq_inter}.
                        Show that $\dom{\restrl{u}{\initialSegment{n}}} = \initialSegment{n}$.
                        \begin{subproof}
                            Show that $\dom{\restrl{u}{\initialSegment{n}}}\subseteq\initialSegment{n}$.
                            \begin{subproof}
                                Thus $\dom{\restrl{u}{\initialSegment{n}}}\subseteq\initialSegment{n}$ by \cref{restrl_dom}.
                            \end{subproof}
                            Show that $\initialSegment{n}\subseteq\dom{\restrl{u}{\initialSegment{n}}}$.
                            \begin{subproof}
                                Show that for all $k$ such that $k\in\initialSegment{n}$ we have $k\in\dom{\restrl{u}{\initialSegment{n}}}$.
                                \begin{subproof}
                                    Fix $k$.
                                    Assume $k\in\initialSegment{n}$.
                                    Thus $k\in\naturalsPlus$ by \cref{initial_segment_naturalsplus}.
                                    $k\in\reals$ by \cref{real_naturals_subset_reals,subseteq}.
                                    Thus $k \leq n$.
                                    $n\in\naturalsPlus$ by \cref{subseteq}.
                                    $n\in\reals$ by \cref{real_naturals_subset_reals,subseteq}.
                                    $0\in\reals$ by \cref{real_add_ident}.
                                    $1\in\reals$ by \cref{real_naturals_subset_reals,real_one_in_naturals,subseteq}.
                                    $n + 1\in\reals$ by \cref{real_add_closed}.
                                    $0 < 1$ by \cref{real_zero_lt_one}.
                                    $0 + n < 1 + n$ by \cref{real_lt_add}.
                                    $0 + n = n$ by \cref{real_add_ident}.
                                    $1 + n = n + 1$ by \cref{real_add_comm}.
                                    Thus $n < n + 1$.
                                    Thus $k \leq n + 1$ by \cref{real_leq_trans}.
                                    Thus $k\in\initialSegment{n + 1}$ by \cref{initial_segment_naturalsplus}.
                                    Thus $k\in\dom{u}$ by \cref{subseteq}.
                                    $(k,\apply{u}{k})\in u$ by \cref{function_apply_iff}.
                                    Thus $(k,\apply{u}{k})\in\restrl{u}{\initialSegment{n}}$ by \cref{restrl_iff}.
                                    Thus $k\in\dom{\restrl{u}{\initialSegment{n}}}$ by \cref{dom_intro}.
                                \end{subproof}
                                Thus $\initialSegment{n}\subseteq\dom{\restrl{u}{\initialSegment{n}}}$ by \cref{subseteq}.
                            \end{subproof}
                            Thus $\dom{\restrl{u}{\initialSegment{n}}} = \initialSegment{n}$ by \cref{subseteq_antisymmetric}.
                        \end{subproof}
                        Show that for all $k\in\dom{\restrl{u}{\initialSegment{n}}}$ we have $\apply{\restrl{u}{\initialSegment{n}}}{k}\in\naturalsPlus$.
                        \begin{subproof}
                            Fix $k\in\dom{\restrl{u}{\initialSegment{n}}}$.
                            Thus $k\in\dom{u}$ by \cref{dom_of_restrl_eq_inter,inter_elim_left}.
                            Thus $u(k)\in\naturalsPlus$.
                            Thus $k\in\initialSegment{n}$ by \cref{subseteq}.
                            $(k,\apply{u}{k})\in u$ by \cref{function_apply_iff}.
                            Thus $(k,\apply{u}{k})\in\restrl{u}{\initialSegment{n}}$ by \cref{restrl_iff}.
                            Thus $\apply{\restrl{u}{\initialSegment{n}}}{k} = u(k)$ by \cref{function_apply_iff}.
                            Thus $\apply{\restrl{u}{\initialSegment{n}}}{k}\in\naturalsPlus$.
                        \end{subproof}
                        Thus $\restrl{u}{\initialSegment{n}}\in\funs{\initialSegment{n}}{\naturalsPlus}$ by \cref{funs_intro}.
                    \end{subproof}
                    Thus $\restrl{u_1}{\initialSegment{n}}\in\funs{\initialSegment{n}}{\naturalsPlus}$ and
                        $\restrl{u_2}{\initialSegment{n}}\in\funs{\initialSegment{n}}{\naturalsPlus}$.
                    $(u_1,\apply{g}{(\apply{f}{\restrl{u_1}{\initialSegment{n}}}, \apply{u_1}{n + 1})})\in F$ by \cref{pair_eq_iff}.
                    $(u_2,\apply{g}{(\apply{f}{\restrl{u_2}{\initialSegment{n}}}, \apply{u_2}{n + 1})})\in F$ by \cref{pair_eq_iff}.
                    $F$ is a function by \cref{funs_elim}.
                    Thus $F(u_1) = \apply{g}{(\apply{f}{\restrl{u_1}{\initialSegment{n}}}, \apply{u_1}{n + 1})}$ and
                        $F(u_2) = \apply{g}{(\apply{f}{\restrl{u_2}{\initialSegment{n}}}, \apply{u_2}{n + 1})}$ by \cref{function_apply_iff}.
                    Thus $\apply{g}{(\apply{f}{\restrl{u_1}{\initialSegment{n}}}, \apply{u_1}{n + 1})} = \apply{g}{(\apply{f}{\restrl{u_2}{\initialSegment{n}}}, \apply{u_2}{n + 1})}$.
                    $f\in\funs{\funs{\initialSegment{n}}{\naturalsPlus}}{\naturalsPlus}$ by \cref{inj}.
                    Thus $\dom{f} = \funs{\initialSegment{n}}{\naturalsPlus}$ by \cref{funs_elim}.
                    Thus $\restrl{u_1}{\initialSegment{n}}\in\dom{f}$ and $\restrl{u_2}{\initialSegment{n}}\in\dom{f}$ by \cref{subseteq}.
                    Thus $\apply{f}{\restrl{u_1}{\initialSegment{n}}}\in\naturalsPlus$ and
                        $\apply{f}{\restrl{u_2}{\initialSegment{n}}}\in\naturalsPlus$ by \cref{funs_elim}.
                    Thus $\dom{u_1} = \initialSegment{n + 1}$ and $\dom{u_2} = \initialSegment{n + 1}$ by \cref{funs_elim}.
                    $n + 1\in\naturalsPlus$ by \cref{real_naturals_closed_succ}.
                    Thus $n + 1\in\initialSegment{n + 1}$ by \cref{initial_segment_naturalsplus}.
                    Thus $n + 1\in\dom{u_1}$ and $n + 1\in\dom{u_2}$ by \cref{subseteq}.
                    Thus $\apply{u_1}{n + 1}\in\naturalsPlus$ and $\apply{u_2}{n + 1}\in\naturalsPlus$ by \cref{funs_elim}.
                    Thus $(\apply{f}{\restrl{u_1}{\initialSegment{n}}}, \apply{u_1}{n + 1})\in\naturalsPlus\times\naturalsPlus$ and
                        $(\apply{f}{\restrl{u_2}{\initialSegment{n}}}, \apply{u_2}{n + 1})\in\naturalsPlus\times\naturalsPlus$ by \cref{times_tuple_intro}.
                    Thus $(\apply{f}{\restrl{u_1}{\initialSegment{n}}}, \apply{u_1}{n + 1}) = (\apply{f}{\restrl{u_2}{\initialSegment{n}}}, \apply{u_2}{n + 1})$ by \cref{inj}.
                    Thus $\apply{f}{\restrl{u_1}{\initialSegment{n}}} = \apply{f}{\restrl{u_2}{\initialSegment{n}}}$ and
                        $\apply{u_1}{n + 1} = \apply{u_2}{n + 1}$ by \cref{pair_eq_iff}.
                    Thus $\restrl{u_1}{\initialSegment{n}} = \restrl{u_2}{\initialSegment{n}}$ by \cref{inj}.
                    $u_1$ is a function and $\dom{u_1} = \initialSegment{n + 1}$ by \cref{funs_elim}.
                    $u_2$ is a function and $\dom{u_2} = \initialSegment{n + 1}$ by \cref{funs_elim}.
                    Show that for all $k\in\initialSegment{n + 1}$ we have $u_1(k) = u_2(k)$.
                    \begin{subproof}
                        Fix $k\in\initialSegment{n + 1}$.
                        $k\in\naturalsPlus$ by \cref{initial_segment_naturalsplus}.
                        $k < n + 1$ or $k = n + 1$ by \cref{initial_segment_naturalsplus}.
                        \begin{byCase}
                        \caseOf{$k = n + 1$.}
                            Thus $u_1(k) = u_2(k)$.
                        \caseOf{$k < n + 1$.}
                            Thus $k \leq n + 1$.
                            Thus $k \neq n + 1$.
                            Thus $k \leq n$ by \cref{real_naturals_leq_succ_elim}.
                            Thus $k\in\initialSegment{n}$ by \cref{initial_segment_naturalsplus}.
                            Thus $\apply{\restrl{u_1}{\initialSegment{n}}}{k} = \apply{\restrl{u_2}{\initialSegment{n}}}{k}$.
                            $(k,\apply{u_1}{k})\in u_1$ by \cref{function_apply_iff}.
                            Thus $(k,\apply{u_1}{k})\in\restrl{u_1}{\initialSegment{n}}$ by \cref{restrl_iff}.
                            $\restrl{u_1}{\initialSegment{n}}$ is a function by \cref{restrl_of_function_is_function}.
                            Thus $\apply{\restrl{u_1}{\initialSegment{n}}}{k} = u_1(k)$ by \cref{function_apply_iff}.
                            $(k,\apply{u_2}{k})\in u_2$ by \cref{function_apply_iff}.
                            Thus $(k,\apply{u_2}{k})\in\restrl{u_2}{\initialSegment{n}}$ by \cref{restrl_iff}.
                            $\restrl{u_2}{\initialSegment{n}}$ is a function by \cref{restrl_of_function_is_function}.
                            Thus $\apply{\restrl{u_2}{\initialSegment{n}}}{k} = u_2(k)$ by \cref{function_apply_iff}.
                            Thus $u_1(k) = u_2(k)$.
                        \end{byCase}
                    \end{subproof}
                    Thus $u_1 = u_2$ by \cref{functions_eq_iff}.
                \end{subproof}
                Thus $F\in\inj{\funs{\initialSegment{n + 1}}{\naturalsPlus}}{\naturalsPlus}$ by \cref{inj}.
            \end{subproof}
            Thus $n + 1\in A$.
        \end{subproof}
        Thus for all $n\in\naturalsPlus$ we have $n\in A$ by \cref{real_naturals_induction}.
    \end{subproof}
    Thus $n\in A$.
    Thus there exists $f\in\inj{\funs{\initialSegment{n}}{\naturalsPlus}}{\naturalsPlus}$ by \cref{pair_eq_iff}.
\end{proof}

% from §30 helper lemma 10: injective coding of all finite sequences of naturalsPlus
\begin{lemma}\label{naturalsplus_finite_sequences_all_inj}
    Let $F = \{ \funs{\initialSegment{k}}{\naturalsPlus} \mid k\in\naturalsPlus \}$.
    Let $S = \unions{F}$.
    Then there exists $h\in\inj{S}{\naturalsPlus}$.
\end{lemma}
\begin{proof}
    Let $F = \{ \funs{\initialSegment{k}}{\naturalsPlus} \mid k\in\naturalsPlus \}$.
    Let $S = \unions{F}$.
    Let $R = \{(k,f) \mid k\in\naturalsPlus \mid f\in\inj{\funs{\initialSegment{k}}{\naturalsPlus}}{\naturalsPlus}\}$.
    Show that for all $k\in\naturalsPlus$ there exists $f$ such that $k\mathrel{R} f$.
    \begin{subproof}
        Fix $k\in\naturalsPlus$.
        Thus there exists $f\in\inj{\funs{\initialSegment{k}}{\naturalsPlus}}{\naturalsPlus}$ such that $f = f$ by \cref{naturalsplus_finite_sequences_inj}.
        Take $f\in\inj{\funs{\initialSegment{k}}{\naturalsPlus}}{\naturalsPlus}$ such that $f = f$.
        Thus $k\mathrel{R} f$ by \cref{pair_eq_iff}.
    \end{subproof}
    Show that $R$ is a relation.
    \begin{subproof}
        Show that for all $w$ such that $w\in R$ there exist $k,f$ such that $w = (k,f)$.
        \begin{subproof}
            Fix $w$.
            Assume $w\in R$.
            Thus there exist $k,f$ such that $k\in\naturalsPlus$ and $f\in\inj{\funs{\initialSegment{k}}{\naturalsPlus}}{\naturalsPlus}$ and
                $w = (k,f)$ by \cref{pair_eq_iff}.
        \end{subproof}
        Thus $R$ is a relation by \cref{relation}.
    \end{subproof}
    Take $F$ such that $F$ is a function on $\naturalsPlus$ and for all $k\in\naturalsPlus$ we have $k\mathrel{R} \apply{F}{k}$
        by \cref{choice_function}.
    Thus there exists $g\in\inj{\naturalsPlus\times\naturalsPlus}{\naturalsPlus}$ such that $g = g$ by \cref{naturalsplus_pairing_injection}.
    Take $g\in\inj{\naturalsPlus\times\naturalsPlus}{\naturalsPlus}$ such that $g = g$.
    Let $h = \{(u, \apply{g}{(k, \apply{\apply{F}{k}}{u})}) \mid k\in\naturalsPlus \mid u\in\funs{\initialSegment{k}}{\naturalsPlus}\}$.
    Show that $h\in\inj{S}{\naturalsPlus}$.
    \begin{subproof}
        Show that $h\in\funs{S}{\naturalsPlus}$.
        \begin{subproof}
            Show that $h$ is a function.
            \begin{subproof}
                Show that for all $u,a,b$ such that $(u,a)\in h$ and $(u,b)\in h$ we have $a = b$.
                \begin{subproof}
                    Fix $u$.
                    Fix $a$.
                    Fix $b$.
                    Assume $(u,a)\in h$ and $(u,b)\in h$.
                    Thus there exists $k_1\in\naturalsPlus$ such that $u\in\funs{\initialSegment{k_1}}{\naturalsPlus}$ and
                        $a = \apply{g}{(k_1, \apply{\apply{F}{k_1}}{u})}$ by \cref{pair_eq_iff}.
                    Thus there exists $k_2\in\naturalsPlus$ such that $u\in\funs{\initialSegment{k_2}}{\naturalsPlus}$ and
                        $b = \apply{g}{(k_2, \apply{\apply{F}{k_2}}{u})}$ by \cref{pair_eq_iff}.
                    Take $k_1\in\naturalsPlus$ such that $u\in\funs{\initialSegment{k_1}}{\naturalsPlus}$ and
                        $a = \apply{g}{(k_1, \apply{\apply{F}{k_1}}{u})}$.
                    Take $k_2\in\naturalsPlus$ such that $u\in\funs{\initialSegment{k_2}}{\naturalsPlus}$ and
                        $b = \apply{g}{(k_2, \apply{\apply{F}{k_2}}{u})}$.
                    $u$ is a function and $\dom{u} = \initialSegment{k_1}$ by \cref{funs_elim}.
                    $u$ is a function and $\dom{u} = \initialSegment{k_2}$ by \cref{funs_elim}.
                    Thus $\initialSegment{k_1} = \initialSegment{k_2}$ by \cref{subseteq_antisymmetric,subseteq}.
                    Thus $k_1 \leq k_1$.
                    Thus $k_1\in\initialSegment{k_1}$ by \cref{initial_segment_naturalsplus}.
                    Thus $k_2 \leq k_2$.
                    Thus $k_2\in\initialSegment{k_2}$ by \cref{initial_segment_naturalsplus}.
                    Thus $k_1\in\initialSegment{k_2}$ and $k_2\in\initialSegment{k_1}$ by \cref{subseteq}.
                    Thus $k_1 \leq k_2$ and $k_2 \leq k_1$ by \cref{initial_segment_naturalsplus}.
                    $k_1\in\reals$ by \cref{real_naturals_subset_reals,subseteq}.
                    $k_2\in\reals$ by \cref{real_naturals_subset_reals,subseteq}.
                    Thus $k_1 = k_2$ by \cref{real_leq_antisym}.
                    Thus $a = b$.
                \end{subproof}
                Thus $h$ is right-unique by \cref{rightunique}.
                Show that $h$ is a relation.
                \begin{subproof}
                    Show that for all $w$ such that $w\in h$ there exist $u,a$ such that $w = (u,a)$.
                    \begin{subproof}
                        Fix $w$.
                        Assume $w\in h$.
                        Thus there exist $u,a$ such that $w = (u,a)$ by \cref{pair_eq_iff}.
                    \end{subproof}
                    Thus $h$ is a relation by \cref{relation}.
                \end{subproof}
                Thus $h$ is a function.
            \end{subproof}
            Show that for all $u\in\dom{h}$ we have $h(u)\in\naturalsPlus$.
            \begin{subproof}
                Fix $u\in\dom{h}$.
                Take $a$ such that $(u,a)\in h$ by \cref{dom_iff}.
                Thus there exists $k\in\naturalsPlus$ such that $u\in\funs{\initialSegment{k}}{\naturalsPlus}$ and
                    $a = \apply{g}{(k, \apply{\apply{F}{k}}{u})}$ by \cref{pair_eq_iff}.
                Take $k\in\naturalsPlus$ such that $u\in\funs{\initialSegment{k}}{\naturalsPlus}$ and
                    $a = \apply{g}{(k, \apply{\apply{F}{k}}{u})}$.
                Thus $k\mathrel{R}\apply{F}{k}$.
                Thus $\apply{F}{k}\in\inj{\funs{\initialSegment{k}}{\naturalsPlus}}{\naturalsPlus}$ by \cref{pair_eq_iff}.
                $\apply{F}{k}\in\funs{\funs{\initialSegment{k}}{\naturalsPlus}}{\naturalsPlus}$ by \cref{inj}.
                Thus $\dom{\apply{F}{k}} = \funs{\initialSegment{k}}{\naturalsPlus}$ by \cref{funs_elim}.
                Thus $u\in\dom{\apply{F}{k}}$ by \cref{subseteq}.
                Thus $\apply{\apply{F}{k}}{u}\in\naturalsPlus$ by \cref{funs_elim}.
                Thus $(k, \apply{\apply{F}{k}}{u})\in\naturalsPlus\times\naturalsPlus$ by \cref{times_tuple_intro}.
                $g\in\funs{\naturalsPlus\times\naturalsPlus}{\naturalsPlus}$ by \cref{inj}.
                Thus $\dom{g} = \naturalsPlus\times\naturalsPlus$ by \cref{funs_elim}.
                Thus $(k, \apply{\apply{F}{k}}{u})\in\dom{g}$ by \cref{subseteq}.
                Thus $a\in\naturalsPlus$ by \cref{funs_elim}.
                Thus $h(u)\in\naturalsPlus$ by \cref{function_apply_iff}.
            \end{subproof}
            Show that $\dom{h} = S$.
            \begin{subproof}
                Show that $\dom{h}\subseteq S$.
                \begin{subproof}
                    Show that for all $u$ such that $u\in\dom{h}$ we have $u\in S$.
                    \begin{subproof}
                        Fix $u$.
                        Assume $u\in\dom{h}$.
                        Take $a$ such that $(u,a)\in h$ by \cref{dom_iff}.
                        Thus there exists $k\in\naturalsPlus$ such that $u\in\funs{\initialSegment{k}}{\naturalsPlus}$ and
                            $a = \apply{g}{(k, \apply{\apply{F}{k}}{u})}$ by \cref{pair_eq_iff}.
                        Take $k\in\naturalsPlus$ such that $u\in\funs{\initialSegment{k}}{\naturalsPlus}$ and
                            $a = \apply{g}{(k, \apply{\apply{F}{k}}{u})}$.
                        $\funs{\initialSegment{k}}{\naturalsPlus}\in F$ by \cref{pair_eq_iff}.
                        Thus $u\in S$ by \cref{unions_iff}.
                    \end{subproof}
                    Thus $\dom{h}\subseteq S$ by \cref{subseteq}.
                \end{subproof}
                Show that $S\subseteq\dom{h}$.
                \begin{subproof}
                    Show that for all $u$ such that $u\in S$ we have $u\in\dom{h}$.
                    \begin{subproof}
                        Fix $u$.
                        Assume $u\in S$.
                        Take $A$ such that $A\in F$ and $u\in A$ by \cref{unions_iff}.
                        Thus there exists $k\in\naturalsPlus$ such that $A = \funs{\initialSegment{k}}{\naturalsPlus}$ by \cref{pair_eq_iff}.
                        Take $k\in\naturalsPlus$ such that $A = \funs{\initialSegment{k}}{\naturalsPlus}$.
                        Thus $u\in\funs{\initialSegment{k}}{\naturalsPlus}$.
                        $(u, \apply{g}{(k, \apply{\apply{F}{k}}{u})})\in h$ by \cref{pair_eq_iff}.
                        Thus $u\in\dom{h}$ by \cref{dom_intro}.
                    \end{subproof}
                    Thus $S\subseteq\dom{h}$ by \cref{subseteq}.
                \end{subproof}
                Thus $\dom{h} = S$ by \cref{subseteq_antisymmetric}.
            \end{subproof}
            Thus $h\in\funs{S}{\naturalsPlus}$ by \cref{funs_intro}.
        \end{subproof}
        Show that for all $u_1,u_2$ such that $u_1\in S$ and $u_2\in S$ and $h(u_1) = h(u_2)$ we have $u_1 = u_2$.
        \begin{subproof}
            Fix $u_1$.
            Fix $u_2$.
            Assume $u_1\in S$ and $u_2\in S$ and $h(u_1) = h(u_2)$.
            $h$ is a function and $\dom{h} = S$ by \cref{funs_elim}.
            Thus $u_1\in\dom{h}$ and $u_2\in\dom{h}$ by \cref{subseteq}.
            $(u_1,\apply{h}{u_1})\in h$ and $(u_2,\apply{h}{u_2})\in h$ by \cref{function_apply_iff}.
            Thus there exists $k_1\in\naturalsPlus$ such that $u_1\in\funs{\initialSegment{k_1}}{\naturalsPlus}$ and
                $\apply{h}{u_1} = \apply{g}{(k_1, \apply{\apply{F}{k_1}}{u_1})}$ by \cref{pair_eq_iff}.
            Thus there exists $k_2\in\naturalsPlus$ such that $u_2\in\funs{\initialSegment{k_2}}{\naturalsPlus}$ and
                $\apply{h}{u_2} = \apply{g}{(k_2, \apply{\apply{F}{k_2}}{u_2})}$ by \cref{pair_eq_iff}.
            Take $k_1\in\naturalsPlus$ such that $u_1\in\funs{\initialSegment{k_1}}{\naturalsPlus}$ and
                $\apply{h}{u_1} = \apply{g}{(k_1, \apply{\apply{F}{k_1}}{u_1})}$.
            Take $k_2\in\naturalsPlus$ such that $u_2\in\funs{\initialSegment{k_2}}{\naturalsPlus}$ and
                $\apply{h}{u_2} = \apply{g}{(k_2, \apply{\apply{F}{k_2}}{u_2})}$.
            Thus $\apply{g}{(k_1, \apply{\apply{F}{k_1}}{u_1})} = \apply{g}{(k_2, \apply{\apply{F}{k_2}}{u_2})}$.
            Thus $k_1\mathrel{R}\apply{F}{k_1}$ and $k_2\mathrel{R}\apply{F}{k_2}$.
            Thus $\apply{F}{k_1}\in\inj{\funs{\initialSegment{k_1}}{\naturalsPlus}}{\naturalsPlus}$ and
                $\apply{F}{k_2}\in\inj{\funs{\initialSegment{k_2}}{\naturalsPlus}}{\naturalsPlus}$ by \cref{pair_eq_iff}.
            $\apply{F}{k_1}\in\funs{\funs{\initialSegment{k_1}}{\naturalsPlus}}{\naturalsPlus}$ and
                $\apply{F}{k_2}\in\funs{\funs{\initialSegment{k_2}}{\naturalsPlus}}{\naturalsPlus}$ by \cref{inj}.
            Thus $\dom{\apply{F}{k_1}} = \funs{\initialSegment{k_1}}{\naturalsPlus}$ and
                $\dom{\apply{F}{k_2}} = \funs{\initialSegment{k_2}}{\naturalsPlus}$ by \cref{funs_elim}.
            Thus $u_1\in\dom{\apply{F}{k_1}}$ and $u_2\in\dom{\apply{F}{k_2}}$ by \cref{subseteq}.
            Thus $\apply{\apply{F}{k_1}}{u_1}\in\naturalsPlus$ and $\apply{\apply{F}{k_2}}{u_2}\in\naturalsPlus$ by \cref{funs_elim}.
            Thus $(k_1, \apply{\apply{F}{k_1}}{u_1})\in\naturalsPlus\times\naturalsPlus$ and
                $(k_2, \apply{\apply{F}{k_2}}{u_2})\in\naturalsPlus\times\naturalsPlus$ by \cref{times_tuple_intro}.
            Thus $(k_1, \apply{\apply{F}{k_1}}{u_1}) = (k_2, \apply{\apply{F}{k_2}}{u_2})$ by \cref{inj}.
            Thus $k_1 = k_2$ and $\apply{\apply{F}{k_1}}{u_1} = \apply{\apply{F}{k_2}}{u_2}$ by \cref{pair_eq_iff}.
            Thus $\apply{\apply{F}{k_1}}{u_1} = \apply{\apply{F}{k_1}}{u_2}$.
            Thus $u_1 = u_2$ by \cref{inj}.
        \end{subproof}
        Thus $h\in\inj{S}{\naturalsPlus}$ by \cref{inj}.
    \end{subproof}
    Thus there exists $h\in\inj{S}{\naturalsPlus}$.
\end{proof}

% from §30 helper lemma 11: enumerate finite sequences of naturalsPlus
\begin{lemma}\label{naturalsplus_finite_sequence_enumeration}
    Let $F = \{ \funs{\initialSegment{k}}{\naturalsPlus} \mid k\in\naturalsPlus \}$.
    Let $S = \unions{F}$.
    There exists $q$ such that
        $q$ is a function on $\naturalsPlus$ and
        for all $n\in\naturalsPlus$ we have $q(n)\in S$ and
        for all $u$ such that $u\in S$
            there exists $n\in\naturalsPlus$ such that $q(n) = u$.
\end{lemma}
\begin{proof}
    Let $F = \{ \funs{\initialSegment{k}}{\naturalsPlus} \mid k\in\naturalsPlus \}$.
    Let $S = \unions{F}$.
    Thus there exists $h\in\inj{S}{\naturalsPlus}$ such that $h = h$ by \cref{naturalsplus_finite_sequences_all_inj}.
    Take $h\in\inj{S}{\naturalsPlus}$ such that $h = h$.
    Let $u_0 = \{(1,1)\}$.
    $1\in\naturalsPlus$ by \cref{real_one_in_naturals}.
    Show that $u_0$ is a function on $\initialSegment{1}$.
    \begin{subproof}
        Show that $u_0$ is a function.
        \begin{subproof}
            Show that $u_0$ is right-unique.
            \begin{subproof}
                Show that for all $x,y,y'$ such that $(x,y)\in u_0$ and $(x,y')\in u_0$ we have $y = y'$.
                \begin{subproof}
                    Fix $x$.
                    Fix $y$.
                    Fix $y'$.
                    Assume $(x,y)\in u_0$ and $(x,y')\in u_0$.
                    Thus $(x,y) = (1,1)$ and $(x,y') = (1,1)$ by \cref{singleton_elim}.
                    Thus $y = 1$ and $y' = 1$ by \cref{pair_eq_iff}.
                    Thus $y = y'$.
                \end{subproof}
                Thus $u_0$ is right-unique by \cref{rightunique}.
            \end{subproof}
            Show that $u_0$ is a relation.
            \begin{subproof}
                Show that for all $w\in u_0$ there exists $x,y$ such that $w = (x,y)$.
                \begin{subproof}
                    Fix $w\in u_0$.
                    Thus $w = (1,1)$ by \cref{singleton_elim}.
                    Let $x = 1$.
                    Let $y = 1$.
                    Thus $w = (x,y)$.
                \end{subproof}
                Thus $u_0$ is a relation by \cref{relation}.
            \end{subproof}
            Thus $u_0$ is a function.
        \end{subproof}
        Show that $\dom{u_0} = \initialSegment{1}$.
        \begin{subproof}
            Show that $\dom{u_0}\subseteq\initialSegment{1}$.
            \begin{subproof}
                Show that for all $x$ such that $x\in\dom{u_0}$ we have $x\in\initialSegment{1}$.
                \begin{subproof}
                    Fix $x$.
                    Assume $x\in\dom{u_0}$.
                    Take $y$ such that $(x,y)\in u_0$ by \cref{dom}.
                    Thus $(x,y) = (1,1)$ by \cref{singleton_elim}.
                    Thus $x = 1$ by \cref{pair_eq_iff}.
                    $1\in\naturalsPlus$ by \cref{real_one_in_naturals}.
                    Thus $x\in\initialSegment{1}$ by \cref{initial_segment_naturalsplus}.
                \end{subproof}
                Thus $\dom{u_0}\subseteq\initialSegment{1}$ by \cref{subseteq}.
            \end{subproof}
            Show that $\initialSegment{1}\subseteq\dom{u_0}$.
            \begin{subproof}
                Show that for all $x$ such that $x\in\initialSegment{1}$ we have $x\in\dom{u_0}$.
                \begin{subproof}
                    Fix $x$.
                    Assume $x\in\initialSegment{1}$.
                    Thus $x\in\naturalsPlus$ and $x \leq 1$ by \cref{initial_segment_naturalsplus}.
                    $1\in\reals$ by \cref{real_naturals_subset_reals,real_one_in_naturals,subseteq}.
                    $x\in\reals$ by \cref{real_naturals_subset_reals,subseteq}.
                    $1 < x$ or $1 = x$ by \cref{real_one_leq_naturalsplus}.
                    \begin{byCase}
                    \caseOf{$1 = x$.}
                        Thus $(x,1) = (1,1)$.
                        Thus $(x,1)\in u_0$ by \cref{singleton_intro}.
                        Thus $x\in\dom{u_0}$ by \cref{dom}.
                    \caseOf{$1 < x$.}
                        Suppose not.
                        Show that $1 < 1$.
                        \begin{subproof}
                            $x < 1$ or $x = 1$ by \cref{real_leq_split}.
                            \begin{byCase}
                            \caseOf{$x < 1$.}
                                Thus $1 < 1$ by \cref{real_lt_trans}.
                            \caseOf{$x = 1$.}
                                Thus $1 < 1$ by \cref{real_lt_subst_right}.
                            \end{byCase}
                        \end{subproof}
                        Contradiction by \cref{real_lt_irreflexive}.
                    \end{byCase}
                \end{subproof}
                Thus $\initialSegment{1}\subseteq\dom{u_0}$ by \cref{subseteq}.
            \end{subproof}
            Thus $\dom{u_0} = \initialSegment{1}$ by \cref{subseteq_antisymmetric}.
        \end{subproof}
        Thus $u_0$ is a function on $\initialSegment{1}$.
    \end{subproof}
    Thus $u_0$ is a function.
    Show that for all $x\in\dom{u_0}$ we have $u_0(x)\in\naturalsPlus$.
    \begin{subproof}
        Fix $x\in\dom{u_0}$.
        Take $y$ such that $(x,y)\in u_0$ by \cref{dom}.
        Thus $(x,y) = (1,1)$ by \cref{singleton_elim}.
        Thus $y = 1$ by \cref{pair_eq_iff}.
        Thus $u_0(x) = y$ by \cref{function_apply_intro}.
        Thus $u_0(x) = 1$.
        Thus $u_0(x)\in\naturalsPlus$ by \cref{real_one_in_naturals}.
    \end{subproof}
    Thus $u_0\in\funs{\initialSegment{1}}{\naturalsPlus}$ by \cref{funs_intro}.
    Thus $u_0\in S$ by \cref{unions_iff}.
    Let $R_a = \{(n,u) \mid n\in\ran{h} \mid u\in S \land h(u) = n\}$.
    Let $R_b = (\naturalsPlus\setminus\ran{h})\times\{u_0\}$.
    Let $R_1 = R_a\union R_b$.
    Show that $R_1$ is a relation.
    \begin{subproof}
        Show that $R_a$ is a relation.
        \begin{subproof}
            Show that for all $w\in R_a$ there exists $x,y$ such that $w = (x,y)$.
            \begin{subproof}
                Fix $w\in R_a$.
                Take $n,u$ such that $n\in\ran{h}$ and $u\in S$ and $h(u) = n$ and $w = (n,u)$.
                Let $x = n$.
                Let $y = u$.
                Thus $w = (x,y)$.
            \end{subproof}
            Thus $R_a$ is a relation by \cref{relation}.
        \end{subproof}
        Show that $R_b$ is a relation.
        \begin{subproof}
            Show that for all $w\in R_b$ there exists $x,y$ such that $w = (x,y)$.
            \begin{subproof}
                Fix $w\in R_b$.
                Take $x,y$ such that $w = (x,y)$ by \cref{times_elem_is_tuple}.
                Thus $w = (x,y)$.
            \end{subproof}
            Thus $R_b$ is a relation by \cref{relation}.
        \end{subproof}
        Thus $R_1$ is a relation by \cref{union_relations_is_relation}.
    \end{subproof}
    Show that for all $n\in\naturalsPlus$ there exists $u$ such that $n\mathrel{R_1} u$.
    \begin{subproof}
        Fix $n\in\naturalsPlus$.
        \begin{byCase}
        \caseOf{$n\in\ran{h}$.}
        $h\in\funs{S}{\naturalsPlus}$ by \cref{inj}.
        Thus $h$ is a function and $\dom{h} = S$ by \cref{funs_elim}.
        Take $u$ such that $u\mathrel{h} n$ by \cref{ran_iff}.
        Thus $(u,n)\in h$.
        Thus $u\in\dom{h}$ by \cref{dom}.
        Thus $u\in S$.
        Thus $h(u) = n$ by \cref{function_apply_intro}.
        Thus $(n,u)\in R_a$.
        Thus $n\mathrel{R_1} u$ by \cref{union_intro_left}.
        \caseOf{$n\notin\ran{h}$.}
        Thus $n\in\naturalsPlus\setminus\ran{h}$ by \cref{setminus_intro}.
        $u_0\in\{u_0\}$ by \cref{singleton_intro}.
        Thus $(n,u_0)\in R_b$ by \cref{times_tuple_intro}.
        Thus $n\mathrel{R_1} u_0$ by \cref{union_intro_right}.
        \end{byCase}
    \end{subproof}
    Take $q$ such that $q$ is a function on $\naturalsPlus$ and for all $n\in\naturalsPlus$ we have $n\mathrel{R_1} \apply{q}{n}$
        by \cref{choice_function}.
    Show that for all $n\in\naturalsPlus$ we have $q(n)\in S$.
    \begin{subproof}
        Fix $n\in\naturalsPlus$.
        Thus $n\mathrel{R_1} \apply{q}{n}$ by \cref{choice_function}.
        \begin{byCase}
        \caseOf{$n\in\ran{h}$.}
        Thus $(n,\apply{q}{n})\in R_1$.
        Thus $(n,\apply{q}{n})\in R_a$ or $(n,\apply{q}{n})\in R_b$ by \cref{union_iff}.
        \begin{byCase}
        \caseOf{$(n,\apply{q}{n})\in R_a$.}
            Take $n_1,u_1$ such that $n_1\in\ran{h}$ and $u_1\in S$ and $h(u_1) = n_1$ and $(n,\apply{q}{n}) = (n_1,u_1)$.
            Thus $n_1 = n$ and $u_1 = \apply{q}{n}$ by \cref{pair_eq_iff}.
            Thus $\apply{q}{n}\in S$.
        \caseOf{$(n,\apply{q}{n})\in R_b$.}
            Thus $\apply{q}{n}\in\{u_0\}$ by \cref{times_tuple_elim}.
            Thus $\apply{q}{n} = u_0$ by \cref{singleton_elim}.
            Thus $\apply{q}{n}\in S$.
        \end{byCase}
        \caseOf{$n\notin\ran{h}$.}
        Thus $(n,\apply{q}{n})\in R_1$ by \cref{choice_function}.
        Show that $(n,\apply{q}{n})\in R_b$.
        \begin{subproof}
            Suppose not.
            Thus $(n,\apply{q}{n})\in R_a$ or $(n,\apply{q}{n})\in R_b$ by \cref{union_iff}.
            Thus $(n,\apply{q}{n})\in R_a$.
            Take $n_1,u_1$ such that $n_1\in\ran{h}$ and $u_1\in S$ and $h(u_1) = n_1$ and $(n,\apply{q}{n}) = (n_1,u_1)$.
            Thus $n_1 = n$ by \cref{pair_eq_iff}.
            Thus $n\in\ran{h}$.
            Contradiction.
        \end{subproof}
        Thus $\apply{q}{n}\in\{u_0\}$ by \cref{times_tuple_elim}.
        Thus $\apply{q}{n} = u_0$ by \cref{singleton_elim}.
        Thus $\apply{q}{n}\in S$.
        \end{byCase}
    \end{subproof}
    Show that for all $u\in S$ there exists $n\in\naturalsPlus$ such that $q(n) = u$.
    \begin{subproof}
        Fix $u\in S$.
        $h\in\funs{S}{\naturalsPlus}$ by \cref{inj}.
        Thus $h$ is a function and $\dom{h} = S$ by \cref{funs_elim}.
        Thus $u\in\dom{h}$.
        Let $n = h(u)$.
        Thus $n\in\naturalsPlus$ by \cref{funs_elim}.
        Thus $(u,n)\in h$ by \cref{function_apply_elim}.
        Thus $n\in\ran{h}$ by \cref{ran_intro}.
        Thus $(n,\apply{q}{n})\in R_1$ by \cref{choice_function}.
        Thus $(n,\apply{q}{n})\in R_a$ or $(n,\apply{q}{n})\in R_b$ by \cref{union_iff}.
        \begin{byCase}
        \caseOf{$(n,\apply{q}{n})\in R_a$.}
            Take $n_1,u_1$ such that $n_1\in\ran{h}$ and $u_1\in S$ and $h(u_1) = n_1$ and $(n,\apply{q}{n}) = (n_1,u_1)$.
            Thus $n_1 = n$ and $u_1 = \apply{q}{n}$ by \cref{pair_eq_iff}.
            Thus $h(\apply{q}{n}) = n$.
            Thus $\apply{q}{n}\in S$.
            Thus $h(\apply{q}{n}) = h(u)$.
            Thus $\apply{q}{n} = u$ by \cref{inj}.
        \caseOf{$(n,\apply{q}{n})\in R_b$.}
            Suppose not.
            Thus $n\in\naturalsPlus\setminus\ran{h}$ by \cref{times_tuple_elim}.
            Thus $n\notin\ran{h}$ by \cref{setminus_elim_right}.
            Contradiction.
        \end{byCase}
    \end{subproof}
    Thus there exists $q$ such that
        $q$ is a function on $\naturalsPlus$ and
        for all $n\in\naturalsPlus$ we have $q(n)\in S$ and
        for all $u\in S$ there exists $n\in\naturalsPlus$ such that $q(n) = u$.
\end{proof}
% from §30 Item 10: countable product of second-countable spaces is second-countable
\begin{lemma}\label{countable_product_second_countable}
    Assume $X$ is a function.
    Assume $T$ is a function.
    Assume $\dom{X} = \naturalsPlus$.
    Assume $\dom{T} = \naturalsPlus$.
    Assume that for all $n\in\naturalsPlus$ we have $\apply{T}{n}$ is a topology on $\apply{X}{n}$.
    Suppose for all $n\in\naturalsPlus$ we have $\apply{X}{n}$ satisfies the second countability axiom with respect to $\apply{T}{n}$.
    Then $\productFamily{X}$ satisfies the second countability axiom with respect to $\productTopologyIndexed{T}{X}$.
\end{lemma}
\begin{proof}
    Let $T_0 = \productTopologyIndexed{T}{X}$.
    Let $R = \{(n,b) \mid n\in\naturalsPlus \mid
        \text{$b$ is a sequence in $\pow{\apply{X}{n}}$ and
        for all $k\in\naturalsPlus$ we have $b(k)$ is open in $\apply{X}{n}$ with respect to $\apply{T}{n}$ and
        for all $x\in\apply{X}{n}$ we have
            for all $U$ such that $U$ is open in $\apply{X}{n}$ with respect to $\apply{T}{n}$ and $x\in U$
                there exists $k\in\naturalsPlus$ such that $x\in b(k)$ and $b(k)\subseteq U$}\}$.
    Show that $R$ is a relation.
    \begin{subproof}
        Show that for all $w$ such that $w\in R$ there exist $n,b$ such that $w = (n,b)$.
        \begin{subproof}
            Fix $w$.
            Assume $w\in R$.
            Thus there exist $n,b$ such that $w = (n,b)$ by \cref{pair_eq_iff}.
        \end{subproof}
        Thus $R$ is a relation by \cref{relation}.
    \end{subproof}
    Show that for all $n\in\naturalsPlus$ there exists $b$ such that $n\mathrel{R} b$.
    \begin{subproof}
        Fix $n\in\naturalsPlus$.
        $\apply{X}{n}$ satisfies the second countability axiom with respect to $\apply{T}{n}$.
        Thus there exists $b$ such that
            $b$ is a sequence in $\pow{\apply{X}{n}}$ and
            for all $k\in\naturalsPlus$ we have $b(k)$ is open in $\apply{X}{n}$ with respect to $\apply{T}{n}$ and
            for all $x\in\apply{X}{n}$ we have
                for all $U$ such that $U$ is open in $\apply{X}{n}$ with respect to $\apply{T}{n}$ and $x\in U$
                    there exists $k\in\naturalsPlus$ such that $x\in b(k)$ and $b(k)\subseteq U$
            by \cref{second_countability_axiom}.
        Take $b$ such that
            $b$ is a sequence in $\pow{\apply{X}{n}}$ and
            for all $k\in\naturalsPlus$ we have $b(k)$ is open in $\apply{X}{n}$ with respect to $\apply{T}{n}$ and
            for all $x\in\apply{X}{n}$ we have
                for all $U$ such that $U$ is open in $\apply{X}{n}$ with respect to $\apply{T}{n}$ and $x\in U$
                    there exists $k\in\naturalsPlus$ such that $x\in b(k)$ and $b(k)\subseteq U$.
        Thus $(n,b)\in R$ by \cref{pair_eq_iff}.
        Thus $n\mathrel{R} b$ by \cref{pair_eq_iff}.
    \end{subproof}
    Take $b$ such that $b$ is a function on $\naturalsPlus$ and for all $n\in\naturalsPlus$ we have
        $n\mathrel{R}\apply{b}{n}$ by \cref{choice_function}.

    Let $F = \{ \funs{\initialSegment{k}}{\naturalsPlus} \mid k\in\naturalsPlus \}$.
    Let $S = \unions{F}$.

    Take $q$ such that
        $q$ is a function on $\naturalsPlus$ and
        for all $n\in\naturalsPlus$ we have $q(n)\in S$ and
        for all $u$ such that $u\in S$
            there exists $n\in\naturalsPlus$ such that $q(n) = u$
        by \cref{naturalsplus_finite_sequence_enumeration}.

    Let $R_d = \{(n,g) \mid n\in\naturalsPlus \mid \text{$g$ is a function on $\dom{\apply{q}{n}}$ and
        for all $i\in\dom{\apply{q}{n}}$ we have
            $(i, \preimg{\piIndex{X}{i}}{\apply{\apply{b}{i}}{\apply{\apply{q}{n}}{i}}})\in g$}\}$.
    Show that $R_d$ is a relation.
    \begin{subproof}
        Show that for all $w$ such that $w\in R_d$ there exist $n,g$ such that $w = (n,g)$.
        \begin{subproof}
            Fix $w$.
            Assume $w\in R_d$.
            Thus there exist $n,g$ such that $w = (n,g)$ by \cref{pair_eq_iff}.
        \end{subproof}
        Thus $R_d$ is a relation by \cref{relation}.
    \end{subproof}
    Show that for all $n\in\naturalsPlus$ there exists $g$ such that $n\mathrel{R_d} g$.
    \begin{subproof}
        Fix $n\in\naturalsPlus$.
        $q$ is a function on $\naturalsPlus$.
        Thus $q$ is a function and $\dom{q} = \naturalsPlus$.
        Thus $n\in\dom{q}$ by \cref{subseteq}.
        Thus $\apply{q}{n}\in S$ by \cref{naturalsplus_finite_sequence_enumeration}.
        Take $m\in\naturalsPlus$ such that $\apply{q}{n}\in\funs{\initialSegment{m}}{\naturalsPlus}$ by \cref{unions_iff}.
        Thus $\apply{q}{n}$ is a function and $\dom{\apply{q}{n}} = \initialSegment{m}$ by \cref{funs_elim}.
        Let $g = \{(i, \preimg{\piIndex{X}{i}}{\apply{\apply{b}{i}}{\apply{\apply{q}{n}}{i}}}) \mid i\in\dom{\apply{q}{n}}\}$.
        Show that $g$ is a function on $\dom{\apply{q}{n}}$.
        \begin{subproof}
            Show that $g$ is a function.
            \begin{subproof}
                Show that for all $i,A_1,A_2$ such that $(i,A_1)\in g$ and $(i,A_2)\in g$ we have $A_1 = A_2$.
                \begin{subproof}
                    Fix $i$.
                    Fix $A_1$.
                    Fix $A_2$.
                    Assume $(i,A_1)\in g$ and $(i,A_2)\in g$.
                    Thus $A_1 = \preimg{\piIndex{X}{i}}{\apply{\apply{b}{i}}{\apply{\apply{q}{n}}{i}}}$ and
                        $A_2 = \preimg{\piIndex{X}{i}}{\apply{\apply{b}{i}}{\apply{\apply{q}{n}}{i}}}$ by \cref{pair_eq_iff}.
                    Thus $A_1 = A_2$.
                \end{subproof}
                Thus $g$ is right-unique by \cref{rightunique}.
                Show that $g$ is a relation.
                \begin{subproof}
                    Show that for all $w$ such that $w\in g$ there exist $i,A$ such that $w = (i,A)$.
                    \begin{subproof}
                        Fix $w$.
                        Assume $w\in g$.
                        Thus there exist $i,A$ such that $w = (i,A)$ by \cref{pair_eq_iff}.
                    \end{subproof}
                    Thus $g$ is a relation by \cref{relation}.
                \end{subproof}
                Thus $g$ is a function.
            \end{subproof}
            Show that $\dom{g} = \dom{\apply{q}{n}}$.
            \begin{subproof}
                Show that $\dom{\apply{q}{n}}\subseteq\dom{g}$.
                \begin{subproof}
                    Show that for all $i$ such that $i\in\dom{\apply{q}{n}}$ we have $i\in\dom{g}$.
                    \begin{subproof}
                        Fix $i$.
                        Assume $i\in\dom{\apply{q}{n}}$.
                        $(i, \preimg{\piIndex{X}{i}}{\apply{\apply{b}{i}}{\apply{\apply{q}{n}}{i}}})\in g$ by \cref{pair_eq_iff}.
                        Thus $i\in\dom{g}$ by \cref{dom_intro}.
                    \end{subproof}
                    Thus $\dom{\apply{q}{n}}\subseteq\dom{g}$ by \cref{subseteq}.
                \end{subproof}
                Show that $\dom{g}\subseteq\dom{\apply{q}{n}}$.
                \begin{subproof}
                    Show that for all $i$ such that $i\in\dom{g}$ we have $i\in\dom{\apply{q}{n}}$.
                    \begin{subproof}
                        Fix $i$.
                        Assume $i\in\dom{g}$.
                        Take $A$ such that $(i,A)\in g$ by \cref{dom_iff}.
                        Thus $i\in\dom{\apply{q}{n}}$ by \cref{pair_eq_iff}.
                    \end{subproof}
                    Thus $\dom{g}\subseteq\dom{\apply{q}{n}}$ by \cref{subseteq}.
                \end{subproof}
                Thus $\dom{g} = \dom{\apply{q}{n}}$ by \cref{subseteq_antisymmetric}.
            \end{subproof}
            Thus $g$ is a function and $\dom{g} = \dom{\apply{q}{n}}$.
            Thus $g$ is a function on $\dom{\apply{q}{n}}$.
        \end{subproof}
        Show that for all $i\in\dom{\apply{q}{n}}$ we have
            $\apply{g}{i} = \preimg{\piIndex{X}{i}}{\apply{\apply{b}{i}}{\apply{\apply{q}{n}}{i}}}$.
        \begin{subproof}
            Fix $i\in\dom{\apply{q}{n}}$.
            $(i, \preimg{\piIndex{X}{i}}{\apply{\apply{b}{i}}{\apply{\apply{q}{n}}{i}}})\in g$ by \cref{pair_eq_iff}.
            Thus $i\in\dom{g}$ by \cref{dom_intro}.
            Thus $\apply{g}{i} = \preimg{\piIndex{X}{i}}{\apply{\apply{b}{i}}{\apply{\apply{q}{n}}{i}}}$ by \cref{function_apply_iff}.
        \end{subproof}
        Thus $(n,g)\in R_d$ by \cref{pair_eq_iff}.
        Thus $n\mathrel{R_d} g$ by \cref{pair_eq_iff}.
    \end{subproof}
    Take $h$ such that $h$ is a function on $\naturalsPlus$ and for all $n\in\naturalsPlus$ we have
        $n\mathrel{R_d}\apply{h}{n}$ by \cref{choice_function}.

    Let $d = \{(n, \inters{\img{\apply{h}{n}}{\dom{\apply{q}{n}}}}) \mid n\in\naturalsPlus\}$.

    Show that $\productSubbasis{T}{X}$ is a subbasis on $\productFamily{X}$.
    \begin{subproof}
        $X$ is a function.
        Thus $X$ is right-unique and $X$ is a relation.
        $T$ is a function.
        Thus $T$ is right-unique and $T$ is a relation.
        $\dom{T} = \dom{X}$.
        $1\in\naturalsPlus$ by \cref{real_one_in_naturals}.
        $\dom{X} = \naturalsPlus$.
        Thus $1\in\dom{X}$ by \cref{subseteq}.
        Thus there exists $a$ such that $a\in\dom{X}$.
        Thus for all $n\in\dom{X}$ we have $\apply{T}{n}$ is a topology on $\apply{X}{n}$.
        Thus $\productSubbasis{T}{X}$ is a subbasis on $\productFamily{X}$ by \cref{product_subbasis_is_subbasis}.
    \end{subproof}

    Show that $d$ is a sequence in $\pow{\productFamily{X}}$.
    \begin{subproof}
        Show that $d\in\funs{\naturalsPlus}{\pow{\productFamily{X}}}$.
        \begin{subproof}
            Show that $d$ is a function.
            \begin{subproof}
                Show that for all $n,U_1,U_2$ such that $(n,U_1)\in d$ and $(n,U_2)\in d$ we have $U_1 = U_2$.
                \begin{subproof}
                    Fix $n$.
                    Fix $U_1$.
                    Fix $U_2$.
                    Assume $(n,U_1)\in d$ and $(n,U_2)\in d$.
                    Thus $U_1 = \inters{\img{\apply{h}{n}}{\dom{\apply{q}{n}}}}$ and
                        $U_2 = \inters{\img{\apply{h}{n}}{\dom{\apply{q}{n}}}}$ by \cref{pair_eq_iff}.
                    Thus $U_1 = U_2$.
                \end{subproof}
                Thus $d$ is right-unique by \cref{rightunique}.
                Show that $d$ is a relation.
                \begin{subproof}
                    Show that for all $w$ such that $w\in d$ there exist $n,U$ such that $w = (n,U)$.
                    \begin{subproof}
                        Fix $w$.
                        Assume $w\in d$.
                        Thus there exist $n,U$ such that $w = (n,U)$ by \cref{pair_eq_iff}.
                    \end{subproof}
                    Thus $d$ is a relation by \cref{relation}.
                \end{subproof}
                Thus $d$ is a function.
            \end{subproof}
            Show that for all $n\in\naturalsPlus$ we have $d(n)\in\pow{\productFamily{X}}$.
            \begin{subproof}
                Fix $n\in\naturalsPlus$.
                Let $U = \inters{\img{\apply{h}{n}}{\dom{\apply{q}{n}}}}$.
                Thus $(n,U)\in d$ by \cref{pair_eq_iff}.
                Thus $n\in\dom{d}$ by \cref{dom_intro}.
                Thus $d(n) = U$ by \cref{function_apply_iff}.
                Let $F = \img{\apply{h}{n}}{\dom{\apply{q}{n}}}$.
                Show that $F\subseteq\productSubbasis{T}{X}$.
                \begin{subproof}
                    Show that for all $A$ such that $A\in F$ we have $A\in\productSubbasis{T}{X}$.
                    \begin{subproof}
                        Fix $A$.
                        Assume $A\in F$.
                        Take $i$ such that $i\in\dom{\apply{q}{n}}$ and $(i,A)\in\apply{h}{n}$ by \cref{img_iff}.
                        Thus $n\mathrel{R_d}\apply{h}{n}$ by \cref{choice_function}.
                        Thus $(n,\apply{h}{n})\in R_d$.
                        Take $n_1,g_1$ such that $n_1\in\naturalsPlus$ and
                            $g_1$ is a function on $\dom{\apply{q}{n_1}}$ and
                            for all $j\in\dom{\apply{q}{n_1}}$ we have
                                $(j, \preimg{\piIndex{X}{j}}{\apply{\apply{b}{j}}{\apply{\apply{q}{n_1}}{j}}})\in g_1$ and
                            $(n,\apply{h}{n}) = (n_1,g_1)$.
                        Thus $n_1 = n$ and $g_1 = \apply{h}{n}$ by \cref{pair_eq_iff}.
                        Thus $\apply{h}{n}$ is a function on $\dom{\apply{q}{n}}$ and
                            for all $j\in\dom{\apply{q}{n}}$ we have
                                $(j, \preimg{\piIndex{X}{j}}{\apply{\apply{b}{j}}{\apply{\apply{q}{n}}{j}}})\in\apply{h}{n}$.
                        Thus $i\in\dom{\apply{h}{n}}$ by \cref{dom_intro}.
                        Thus $\apply{\apply{h}{n}}{i} = A$ by \cref{function_apply_iff}.
                        Thus $A = \preimg{\piIndex{X}{i}}{\apply{\apply{b}{i}}{\apply{\apply{q}{n}}{i}}}$.
                        $\apply{q}{n}\in S$ by \cref{naturalsplus_finite_sequence_enumeration}.
                        Take $m\in\naturalsPlus$ such that $\apply{q}{n}\in\funs{\initialSegment{m}}{\naturalsPlus}$ by \cref{unions_iff}.
                        Thus $\apply{q}{n}$ is a function and $\dom{\apply{q}{n}} = \initialSegment{m}$ by \cref{funs_elim}.
                        Thus $i\in\initialSegment{m}$ by \cref{subseteq}.
                        Thus $i\in\naturalsPlus$ by \cref{initial_segment_naturalsplus}.
                        $\dom{X} = \naturalsPlus$.
                        Thus $i\in\dom{X}$ by \cref{subseteq}.
                        $i\mathrel{R}\apply{b}{i}$ by \cref{choice_function}.
                        Thus $\apply{b}{i}$ is a sequence in $\pow{\apply{X}{i}}$ and
                            for all $k\in\naturalsPlus$ we have $\apply{\apply{b}{i}}{k}$ is open in $\apply{X}{i}$ with respect to $\apply{T}{i}$
                            by \cref{pair_eq_iff}.
                        Thus $\apply{q}{n}$ is a function to $\naturalsPlus$ by \cref{funs_elim}.
                        Thus $\apply{q}{n}$ is a function.
                        Thus $(i,\apply{\apply{q}{n}}{i})\in\apply{q}{n}$ by \cref{function_apply_elim}.
                        Thus $\apply{\apply{q}{n}}{i}\in\ran{\apply{q}{n}}$ by \cref{ran_iff}.
                        Thus $\apply{\apply{q}{n}}{i}\in\naturalsPlus$ by \cref{function_to_ran,subseteq}.
                        Thus $\apply{\apply{b}{i}}{\apply{\apply{q}{n}}{i}}$ is open in $\apply{X}{i}$ with respect to $\apply{T}{i}$.
                        Thus $\apply{\apply{b}{i}}{\apply{\apply{q}{n}}{i}}\in\apply{T}{i}$ by \cref{open_in}.
                        $\piIndex{X}{i}\in\funs{\productFamily{X}}{\apply{X}{i}}$ by \cref{projection_map_indexed_function}.
                        Thus $\dom{\piIndex{X}{i}} = \productFamily{X}$ by \cref{funs_elim}.
                        Thus $A\subseteq\dom{\piIndex{X}{i}}$ by \cref{preimg_subseteq_dom}.
                        Thus $A\subseteq\productFamily{X}$.
                        Thus $A\in\pow{\productFamily{X}}$ by \cref{powerset_intro}.
                        Thus $A\in\productSubbasisAt{T}{X}{i}$ by \cref{product_subbasis_indexed}.
                        Thus $A\in\productSubbasis{T}{X}$ by \cref{product_subbasis_union_indexed}.
                    \end{subproof}
                    Thus $F\subseteq\productSubbasis{T}{X}$ by \cref{subseteq}.
                \end{subproof}
                Show that $F$ is finite.
                \begin{subproof}
                    $\apply{q}{n}\in S$ by \cref{naturalsplus_finite_sequence_enumeration}.
                    Take $m\in\naturalsPlus$ such that $\apply{q}{n}\in\funs{\initialSegment{m}}{\naturalsPlus}$ by \cref{unions_iff}.
                    Thus $\apply{q}{n}$ is a function and $\dom{\apply{q}{n}} = \initialSegment{m}$ by \cref{funs_elim}.
                    Thus $\dom{\apply{q}{n}}$ is finite by \cref{initial_segment_finite}.
                    Thus $n\mathrel{R_d}\apply{h}{n}$ by \cref{choice_function}.
                    Thus $(n,\apply{h}{n})\in R_d$.
                    Thus $\apply{h}{n}$ is right-unique and $\apply{h}{n}$ is a relation and
                        $\dom{\apply{q}{n}} = \dom{\apply{h}{n}}$ by \cref{pair_eq_iff}.
                    Thus $F$ is finite by \cref{finite_image_function}.
                \end{subproof}
                Show that $F$ is inhabited.
                \begin{subproof}
                    $\apply{q}{n}\in S$ by \cref{naturalsplus_finite_sequence_enumeration}.
                    Take $m\in\naturalsPlus$ such that $\apply{q}{n}\in\funs{\initialSegment{m}}{\naturalsPlus}$ by \cref{unions_iff}.
                    Thus $\apply{q}{n}$ is a function and $\dom{\apply{q}{n}} = \initialSegment{m}$ by \cref{funs_elim}.
                    $m\in\initialSegment{m}$ by \cref{initial_segment_naturalsplus}.
                    Thus $m\in\dom{\apply{q}{n}}$ by \cref{subseteq}.
                    Thus $n\mathrel{R_d}\apply{h}{n}$ by \cref{choice_function}.
                    Thus $\apply{h}{n}$ is a function and $\dom{\apply{h}{n}} = \dom{\apply{q}{n}}$ by \cref{pair_eq_iff}.
                    Thus $m\in\dom{\apply{h}{n}}$ by \cref{subseteq}.
                    Thus $(m,\apply{\apply{h}{n}}{m})\in\apply{h}{n}$ by \cref{function_apply_iff}.
                    Thus $\apply{\apply{h}{n}}{m}\in F$ by \cref{img_iff}.
                \end{subproof}
                Thus $U = \inters{F}$.
                Show that $U\subseteq\unions{F}$.
                \begin{subproof}
                    Show that for all $x$ such that $x\in U$ we have $x\in\unions{F}$.
                    \begin{subproof}
                        Fix $x$.
                        Assume $x\in U$.
                        Thus $x\in\inters{F}$.
                        Thus $x\in\unions{F}$ by \cref{inters}.
                    \end{subproof}
                    Thus $U\subseteq\unions{F}$ by \cref{subseteq}.
                \end{subproof}
                Show that $\unions{F}\subseteq\unions{\productSubbasis{T}{X}}$.
                \begin{subproof}
                    Show that for all $x$ such that $x\in\unions{F}$ we have $x\in\unions{\productSubbasis{T}{X}}$.
                    \begin{subproof}
                        Fix $x$.
                        Assume $x\in\unions{F}$.
                        Take $A$ such that $A\in F$ and $x\in A$ by \cref{unions_iff}.
                        Thus $A\in\productSubbasis{T}{X}$ by \cref{subseteq}.
                        Thus $x\in\unions{\productSubbasis{T}{X}}$ by \cref{unions_intro}.
                    \end{subproof}
                    Thus $\unions{F}\subseteq\unions{\productSubbasis{T}{X}}$ by \cref{subseteq}.
                \end{subproof}
                Thus $U\subseteq\unions{\productSubbasis{T}{X}}$ by \cref{subseteq_transitive}.
                Thus $U\in\pow{\unions{\productSubbasis{T}{X}}}$ by \cref{powerset_intro}.
                $\unions{\productSubbasis{T}{X}} = \productFamily{X}$ by \cref{subbasis_on}.
                Thus $U\in\pow{\productFamily{X}}$ by \cref{subseteq}.
                Thus $d(n)\in\pow{\productFamily{X}}$.
            \end{subproof}
            Show that $\dom{d} = \naturalsPlus$.
            \begin{subproof}
                Show that $\dom{d}\subseteq\naturalsPlus$.
                \begin{subproof}
                    Show that for all $x$ such that $x\in\dom{d}$ we have $x\in\naturalsPlus$.
                    \begin{subproof}
                        Fix $x$.
                        Assume $x\in\dom{d}$.
                        Take $U$ such that $(x,U)\in d$ by \cref{dom_iff}.
                        Take $n$ such that $n\in\naturalsPlus$ and $(x,U) = (n,\inters{\img{\apply{h}{n}}{\dom{\apply{q}{n}}}})$ by \cref{pair_eq_iff}.
                        Thus $x = n$ by \cref{pair_eq_iff}.
                        Thus $x\in\naturalsPlus$.
                    \end{subproof}
                    Thus $\dom{d}\subseteq\naturalsPlus$ by \cref{subseteq}.
                \end{subproof}
                Show that $\naturalsPlus\subseteq\dom{d}$.
                \begin{subproof}
                    Show that for all $x$ such that $x\in\naturalsPlus$ we have $x\in\dom{d}$.
                    \begin{subproof}
                        Fix $x$.
                        Assume $x\in\naturalsPlus$.
                        Thus $(x,\inters{\img{\apply{h}{x}}{\dom{\apply{q}{x}}}})\in d$ by \cref{pair_eq_iff}.
                        Thus $x\in\dom{d}$ by \cref{dom_intro}.
                    \end{subproof}
                    Thus $\naturalsPlus\subseteq\dom{d}$ by \cref{subseteq}.
                \end{subproof}
                Thus $\dom{d} = \naturalsPlus$ by \cref{subseteq_antisymmetric}.
            \end{subproof}
            Thus $\dom{d} = \naturalsPlus$.
            Thus $\naturalsPlus = \dom{d}$.
            Thus $d\in\funs{\naturalsPlus}{\pow{\productFamily{X}}}$ by \cref{funs_intro}.
        \end{subproof}
        Thus $d$ is a sequence in $\pow{\productFamily{X}}$ by \cref{sequence_in}.
    \end{subproof}

    Show that for all $n\in\naturalsPlus$ we have $d(n)$ is open in $\productFamily{X}$ with respect to $T_0$.
    \begin{subproof}
        Fix $n\in\naturalsPlus$.
        Let $F = \img{\apply{h}{n}}{\dom{\apply{q}{n}}}$.
        Show that $F\subseteq\productSubbasis{T}{X}$.
        \begin{subproof}
            Show that for all $A$ such that $A\in F$ we have $A\in\productSubbasis{T}{X}$.
            \begin{subproof}
                Fix $A$.
                Assume $A\in F$.
                Take $i$ such that $i\in\dom{\apply{q}{n}}$ and $(i,A)\in\apply{h}{n}$ by \cref{img_iff}.
                Thus $n\mathrel{R_d}\apply{h}{n}$ by \cref{choice_function}.
                Thus $(n,\apply{h}{n})\in R_d$.
                Take $n_1,g_1$ such that $n_1\in\naturalsPlus$ and
                    $g_1$ is a function on $\dom{\apply{q}{n_1}}$ and
                    for all $j\in\dom{\apply{q}{n_1}}$ we have
                        $(j,\preimg{\piIndex{X}{j}}{\apply{\apply{b}{j}}{\apply{\apply{q}{n_1}}{j}}})\in g_1$ and
                    $(n,\apply{h}{n}) = (n_1,g_1)$ by \cref{pair_eq_iff}.
                Thus $n_1 = n$ and $g_1 = \apply{h}{n}$ by \cref{pair_eq_iff}.
                Thus $\apply{h}{n}$ is a function on $\dom{\apply{q}{n}}$ and
                    for all $j\in\dom{\apply{q}{n}}$ we have
                        $(j,\preimg{\piIndex{X}{j}}{\apply{\apply{b}{j}}{\apply{\apply{q}{n}}{j}}})\in\apply{h}{n}$.
                Thus $\apply{h}{n}$ is a function and $\dom{\apply{h}{n}} = \dom{\apply{q}{n}}$.
                Show that for all $j\in\dom{\apply{q}{n}}$ we have
                    $\apply{\apply{h}{n}}{j} = \preimg{\piIndex{X}{j}}{\apply{\apply{b}{j}}{\apply{\apply{q}{n}}{j}}}$.
                \begin{subproof}
                    Fix $j\in\dom{\apply{q}{n}}$.
                    Thus $(j,\preimg{\piIndex{X}{j}}{\apply{\apply{b}{j}}{\apply{\apply{q}{n}}{j}}})\in\apply{h}{n}$.
                    Thus $j\in\dom{\apply{h}{n}}$ by \cref{dom_intro}.
                    Thus $\apply{\apply{h}{n}}{j} = \preimg{\piIndex{X}{j}}{\apply{\apply{b}{j}}{\apply{\apply{q}{n}}{j}}}$ by \cref{function_apply_iff}.
                \end{subproof}
                Thus $i\in\dom{\apply{h}{n}}$ by \cref{dom_intro}.
                Thus $\apply{\apply{h}{n}}{i} = A$ by \cref{function_apply_iff}.
                Thus $A = \preimg{\piIndex{X}{i}}{\apply{\apply{b}{i}}{\apply{\apply{q}{n}}{i}}}$.
                $\apply{q}{n}\in S$ by \cref{naturalsplus_finite_sequence_enumeration}.
                Take $m\in\naturalsPlus$ such that $\apply{q}{n}\in\funs{\initialSegment{m}}{\naturalsPlus}$ by \cref{unions_iff}.
                Thus $\apply{q}{n}$ is a function to $\naturalsPlus$ and $\dom{\apply{q}{n}} = \initialSegment{m}$ by \cref{funs_elim}.
                Thus $i\in\dom{\apply{q}{n}}$.
                Thus $i\in\initialSegment{m}$ by \cref{subseteq}.
                Thus $i\in\naturalsPlus$ by \cref{initial_segment_naturalsplus}.
                $\dom{X} = \naturalsPlus$.
                Thus $i\in\dom{X}$ by \cref{subseteq}.
                $i\mathrel{R}\apply{b}{i}$ by \cref{choice_function}.
                Thus $\apply{b}{i}$ is a sequence in $\pow{\apply{X}{i}}$ and
                    for all $k\in\naturalsPlus$ we have $\apply{\apply{b}{i}}{k}$ is open in $\apply{X}{i}$ with respect to $\apply{T}{i}$
                    by \cref{pair_eq_iff}.
                Thus $\ran{\apply{q}{n}}\subseteq\naturalsPlus$ by \cref{function_to_ran}.
                Thus $(i,\apply{\apply{q}{n}}{i})\in\apply{q}{n}$ by \cref{function_apply_elim}.
                Thus $\apply{\apply{q}{n}}{i}\in\ran{\apply{q}{n}}$ by \cref{ran_iff}.
                Thus $\apply{\apply{q}{n}}{i}\in\naturalsPlus$ by \cref{subseteq}.
                Thus $\apply{\apply{b}{i}}{\apply{\apply{q}{n}}{i}}$ is open in $\apply{X}{i}$ with respect to $\apply{T}{i}$.
                Thus $\apply{\apply{b}{i}}{\apply{\apply{q}{n}}{i}}\in\apply{T}{i}$ by \cref{open_in}.
                $\piIndex{X}{i}\in\funs{\productFamily{X}}{\apply{X}{i}}$ by \cref{projection_map_indexed_function}.
                Thus $\dom{\piIndex{X}{i}} = \productFamily{X}$ by \cref{funs_elim}.
                Thus $A\subseteq\dom{\piIndex{X}{i}}$ by \cref{preimg_subseteq_dom}.
                Thus $A\subseteq\productFamily{X}$.
                Thus $A\in\pow{\productFamily{X}}$ by \cref{powerset_intro}.
                Let $U = \apply{\apply{b}{i}}{\apply{\apply{q}{n}}{i}}$.
                Thus $U\in\apply{T}{i}$.
                Thus $A = \preimg{\piIndex{X}{i}}{U}$.
                Thus $A\in\productSubbasisAt{T}{X}{i}$ by \cref{product_subbasis_indexed}.
                Thus $A\in\productSubbasis{T}{X}$ by \cref{product_subbasis_union_indexed}.
            \end{subproof}
            Thus $F\subseteq\productSubbasis{T}{X}$ by \cref{subseteq}.
        \end{subproof}
        Show that $F$ is finite.
        \begin{subproof}
            $\apply{q}{n}\in S$ by \cref{naturalsplus_finite_sequence_enumeration}.
            Take $m\in\naturalsPlus$ such that $\apply{q}{n}\in\funs{\initialSegment{m}}{\naturalsPlus}$ by \cref{unions_iff}.
            Thus $\apply{q}{n}$ is a function and $\dom{\apply{q}{n}} = \initialSegment{m}$ by \cref{funs_elim}.
            $\initialSegment{m}$ is finite by \cref{initial_segment_finite}.
            Thus $n\mathrel{R_d}\apply{h}{n}$ by \cref{choice_function}.
            Thus $(n,\apply{h}{n})\in R_d$.
            Take $n_1,g_1$ such that $n_1\in\naturalsPlus$ and
                $g_1$ is a function on $\dom{\apply{q}{n_1}}$ and
                $(n,\apply{h}{n}) = (n_1,g_1)$ and
                for all $j\in\dom{\apply{q}{n_1}}$ we have
                    $(j,\preimg{\piIndex{X}{j}}{\apply{\apply{b}{j}}{\apply{\apply{q}{n_1}}{j}}})\in g_1$
                by \cref{pair_eq_iff}.
            Thus $n_1 = n$ by \cref{pair_eq_iff}.
            Thus $g_1 = \apply{h}{n}$ by \cref{pair_eq_iff}.
            Thus $\apply{h}{n}$ is a function and $\dom{\apply{h}{n}} = \dom{\apply{q}{n}}$.
            Thus $F$ is finite by \cref{finite_image_function}.
        \end{subproof}
        Show that $F$ is inhabited.
        \begin{subproof}
            $\apply{q}{n}\in S$ by \cref{naturalsplus_finite_sequence_enumeration}.
            Take $m\in\naturalsPlus$ such that $\apply{q}{n}\in\funs{\initialSegment{m}}{\naturalsPlus}$ by \cref{unions_iff}.
            Thus $\apply{q}{n}$ is a function and $\dom{\apply{q}{n}} = \initialSegment{m}$ by \cref{funs_elim}.
            $m\in\initialSegment{m}$ by \cref{initial_segment_naturalsplus}.
            Thus $m\in\dom{\apply{q}{n}}$ by \cref{subseteq}.
            Thus $n\mathrel{R_d}\apply{h}{n}$ by \cref{choice_function}.
            Thus $\apply{h}{n}$ is a function and $\dom{\apply{h}{n}} = \dom{\apply{q}{n}}$ by \cref{pair_eq_iff}.
            Thus $m\in\dom{\apply{h}{n}}$ by \cref{subseteq}.
            Thus $(m,\apply{\apply{h}{n}}{m})\in\apply{h}{n}$ by \cref{function_apply_iff}.
            Thus $\apply{\apply{h}{n}}{m}\in F$ by \cref{img_iff}.
        \end{subproof}
        Show that $\inters{F}\in\pow{\unions{\productSubbasis{T}{X}}}$.
        \begin{subproof}
            Show that $\inters{F}\subseteq\unions{\productSubbasis{T}{X}}$.
            \begin{subproof}
                Show that for all $x$ such that $x\in\inters{F}$ we have $x\in\unions{\productSubbasis{T}{X}}$.
                \begin{subproof}
                    Fix $x$.
                    Assume $x\in\inters{F}$.
                    Thus $x\in\unions{F}$ by \cref{inters}.
                    Take $A$ such that $A\in F$ and $x\in A$ by \cref{unions_iff}.
                    Thus $A\in\productSubbasis{T}{X}$ by \cref{subseteq}.
                    Thus $x\in\unions{\productSubbasis{T}{X}}$ by \cref{unions_intro}.
                \end{subproof}
                Thus $\inters{F}\subseteq\unions{\productSubbasis{T}{X}}$ by \cref{subseteq}.
            \end{subproof}
            Thus $\inters{F}\in\pow{\unions{\productSubbasis{T}{X}}}$ by \cref{powerset_intro}.
        \end{subproof}
        Thus $\inters{F}\in\finiteIntersections{\productSubbasis{T}{X}}$ by \cref{finite_intersections}.
        Thus $\finiteIntersections{\productSubbasis{T}{X}}$ is a basis for $T_0$ on $\productFamily{X}$
            by \cref{product_subbasis_finite_intersections_basis}.
        Thus $T_0 = \basisTopology{\finiteIntersections{\productSubbasis{T}{X}}}{\productFamily{X}}$
            by \cref{product_topology_indexed,topology_generated_by_subbasis,basis_for_topology}.
        Thus $\inters{F}\in\basisTopology{\finiteIntersections{\productSubbasis{T}{X}}}{\productFamily{X}}$
            by \cref{basis_elem_open_in_generated,basis_for_topology}.
        Thus $\inters{F}\in T_0$ by \cref{basis_for_topology}.
        Thus $T_0$ is a topology on $\productFamily{X}$ by \cref{basis_for_topology,basis_topology_is_topology}.
        $d\in\funs{\naturalsPlus}{\pow{\productFamily{X}}}$ by \cref{sequence_in}.
        Thus $d$ is a function by \cref{funs_is_function}.
        Thus $\dom{d} = \naturalsPlus$ by \cref{funs_elim}.
        Thus $n\in\dom{d}$ by \cref{subseteq}.
        $(n,\inters{F})\in d$ by \cref{pair_eq_iff}.
        Thus $d(n) = \inters{F}$ by \cref{function_apply_iff}.
        Thus $d(n)$ is open in $\productFamily{X}$ with respect to $T_0$ by \cref{open_in}.
    \end{subproof}

    Show that for all $x\in\productFamily{X}$ we have
        for all $U$ such that $U$ is open in $\productFamily{X}$ with respect to $T_0$ and $x\in U$
            there exists $n\in\naturalsPlus$ such that $x\in d(n)$ and $d(n)\subseteq U$.
    \begin{subproof}
        Fix $x\in\productFamily{X}$.
        Fix $U$.
        Assume $U$ is open in $\productFamily{X}$ with respect to $T_0$ and $x\in U$.
        $U\in T_0$ by \cref{open_in}.
        $X$ is a function.
        Thus $X$ is right-unique and $X$ is a relation.
        $T$ is a function.
        Thus $T$ is right-unique and $T$ is a relation.
        $\dom{T} = \dom{X}$.
        $1\in\naturalsPlus$ by \cref{real_one_in_naturals}.
        $\dom{X} = \naturalsPlus$.
        Thus $1\in\dom{X}$ by \cref{subseteq}.
        Thus there exists $a$ such that $a\in\dom{X}$.
        Thus for all $n\in\dom{X}$ we have $\apply{T}{n}$ is a topology on $\apply{X}{n}$.
        Thus $\finiteIntersections{\productSubbasis{T}{X}}$ is a basis for $T_0$ on $\productFamily{X}$
            by \cref{product_subbasis_finite_intersections_basis}.
        Thus $T_0 = \basisTopology{\finiteIntersections{\productSubbasis{T}{X}}}{\productFamily{X}}$
            by \cref{product_topology_indexed,topology_generated_by_subbasis,basis_for_topology}.
        Thus $U\in\basisTopology{\finiteIntersections{\productSubbasis{T}{X}}}{\productFamily{X}}$.
        Take $B\in\finiteIntersections{\productSubbasis{T}{X}}$ such that $x\in B$ and $B\subseteq U$
            by \cref{topology_generated_by_basis}.
        Take $F\subseteq\productSubbasis{T}{X}$ such that $F$ is finite and inhabited and $B = \inters{F}$ by \cref{finite_intersections}.

        Let $R_2 = \{(A,i) \mid A\in F \mid
            \text{$i\in\dom{X}$ and there exists $V\in\apply{T}{i}$ such that $A = \preimg{\piIndex{X}{i}}{V}$}\}$.
        Show that $R_2$ is a relation.
        \begin{subproof}
            Show that for all $w$ such that $w\in R_2$ there exist $A,i$ such that $w = (A,i)$.
            \begin{subproof}
                Fix $w$.
                Assume $w\in R_2$.
                Thus there exist $A,i$ such that $w = (A,i)$ by \cref{pair_eq_iff}.
            \end{subproof}
            Thus $R_2$ is a relation by \cref{relation}.
        \end{subproof}
        Show that for all $A\in F$ there exists $i$ such that $A\mathrel{R_2} i$.
        \begin{subproof}
            Fix $A\in F$.
            $A\in\productSubbasis{T}{X}$ by \cref{subseteq}.
            Take $i$ such that $i\in\dom{X}$ and $A\in\productSubbasisAt{T}{X}{i}$
                by \cref{product_subbasis_union_indexed}.
            Take $V\in\apply{T}{i}$ such that $A = \preimg{\piIndex{X}{i}}{V}$ by \cref{product_subbasis_indexed}.
            Thus $(A,i)\in R_2$ by \cref{pair_eq_iff}.
            Thus $A\mathrel{R_2} i$ by \cref{pair_eq_iff}.
        \end{subproof}
        Take $p$ such that $p$ is a function on $F$ and for all $A\in F$ we have $A\mathrel{R_2}\apply{p}{A}$
            by \cref{choice_function}.
        Let $J = \img{p}{F}$.
        $F$ is finite.
        Thus $J$ is finite by \cref{finite_image_function}.
        Show that $J$ is inhabited.
        \begin{subproof}
            Take $A$ such that $A\in F$.
            $p$ is a function on $F$.
            Thus $p$ is a function and $\dom{p} = F$.
            Thus $A\in\dom{p}$ by \cref{subseteq}.
            Thus $(A,\apply{p}{A})\in p$ by \cref{function_apply_iff}.
            Thus $\apply{p}{A}\in J$ by \cref{img_iff}.
        \end{subproof}
        Show that $J\subseteq\naturalsPlus$.
        \begin{subproof}
            Show that for all $j$ such that $j\in J$ we have $j\in\naturalsPlus$.
            \begin{subproof}
                Fix $j$.
                Assume $j\in J$.
                Take $A$ such that $A\in F$ and $(A,j)\in p$ by \cref{img_iff}.
                $p$ is a function on $F$.
                Thus $p$ is a function and $\dom{p} = F$.
                Thus $A\in\dom{p}$ by \cref{subseteq}.
                Thus $\apply{p}{A} = j$ by \cref{function_apply_iff}.
                Thus $A\mathrel{R_2}\apply{p}{A}$ by \cref{choice_function}.
                Thus $\apply{p}{A}\in\dom{X}$ by \cref{pair_eq_iff}.
                $\dom{X} = \naturalsPlus$.
                Thus $j\in\naturalsPlus$.
            \end{subproof}
            Thus $J\subseteq\naturalsPlus$ by \cref{subseteq}.
        \end{subproof}

        Let $R_3 = \{(A,V) \mid A\in F \mid
            \text{$V\in\apply{T}{\apply{p}{A}}$ and $A = \preimg{\piIndex{X}{\apply{p}{A}}}{V}$}\}$.
        Show that $R_3$ is a relation.
        \begin{subproof}
            Show that for all $w$ such that $w\in R_3$ there exist $A,V$ such that $w = (A,V)$.
            \begin{subproof}
                Fix $w$.
                Assume $w\in R_3$.
                Thus there exist $A,V$ such that $w = (A,V)$ by \cref{pair_eq_iff}.
            \end{subproof}
            Thus $R_3$ is a relation by \cref{relation}.
        \end{subproof}
        Show that for all $A\in F$ there exists $V$ such that $A\mathrel{R_3} V$.
        \begin{subproof}
            Fix $A\in F$.
            Thus $A\mathrel{R_2}\apply{p}{A}$ by \cref{choice_function}.
            Thus there exists $V\in\apply{T}{\apply{p}{A}}$ such that $A = \preimg{\piIndex{X}{\apply{p}{A}}}{V}$
                by \cref{pair_eq_iff}.
            Take $V\in\apply{T}{\apply{p}{A}}$ such that $A = \preimg{\piIndex{X}{\apply{p}{A}}}{V}$.
            Thus $(A,V)\in R_3$ by \cref{pair_eq_iff}.
            Thus $A\mathrel{R_3} V$ by \cref{pair_eq_iff}.
        \end{subproof}
        Take $r$ such that $r$ is a function on $F$ and for all $A\in F$ we have $A\mathrel{R_3}\apply{r}{A}$
            by \cref{choice_function}.

        Let $R_4 = \{(i,V) \mid i\in J \mid
            V = \inters{\img{r}{\preimg{p}{\{i\}}}}\}$.
        Show that $R_4$ is a relation.
        \begin{subproof}
            Show that for all $w$ such that $w\in R_4$ there exist $i,V$ such that $w = (i,V)$.
            \begin{subproof}
                Fix $w$.
                Assume $w\in R_4$.
                Thus there exist $i,V$ such that $w = (i,V)$ by \cref{pair_eq_iff}.
            \end{subproof}
            Thus $R_4$ is a relation by \cref{relation}.
        \end{subproof}
        Show that for all $i\in J$ there exists $V$ such that $i\mathrel{R_4} V$.
        \begin{subproof}
            Fix $i\in J$.
            Let $H = \img{r}{\preimg{p}{\{i\}}}$.
            Let $V = \inters{H}$.
            Thus $(i,V)\in R_4$ by \cref{pair_eq_iff}.
            Thus $i\mathrel{R_4} V$ by \cref{pair_eq_iff}.
        \end{subproof}
        Take $v$ such that $v$ is a function on $J$ and for all $i\in J$ we have $i\mathrel{R_4}\apply{v}{i}$
            by \cref{choice_function}.

        Show that for all $i\in J$ we have $\apply{v}{i}\in\apply{T}{i}$.
        \begin{subproof}
            Fix $i\in J$.
            Let $H = \img{r}{\preimg{p}{\{i\}}}$.
            Thus $i\mathrel{R_4}\apply{v}{i}$ by \cref{choice_function}.
            Thus $\apply{v}{i} = \inters{H}$ by \cref{pair_eq_iff}.
            Show that $H\subseteq\apply{T}{i}$.
            \begin{subproof}
                Show that for all $V$ such that $V\in H$ we have $V\in\apply{T}{i}$.
                \begin{subproof}
                    Fix $V$.
                    Assume $V\in H$.
                    Thus $V\in\img{r}{\preimg{p}{\{i\}}}$.
                    Take $A$ such that $A\in\preimg{p}{\{i\}}$ and $(A,V)\in r$ by \cref{img_iff}.
                    Take $y\in\{i\}$ such that $A\mathrel{p} y$ by \cref{preimg_iff}.
                    Thus $(A,y)\in p$.
                    $p$ is a function on $F$.
                    Thus $p$ is a function and $\dom{p} = F$.
                    Thus $A\in\dom{p}$ by \cref{dom_intro}.
                    Thus $A\in F$ by \cref{subseteq}.
                    Thus $y = i$ by \cref{singleton_elim}.
                    Thus $\apply{p}{A} = y$ by \cref{function_apply_iff}.
                    Thus $\apply{p}{A} = i$.
                    $r$ is a function on $F$.
                    Thus $r$ is a function and $\dom{r} = F$.
                    Thus $A\in\dom{r}$ by \cref{subseteq}.
                    Thus $\apply{r}{A} = V$ by \cref{function_apply_iff}.
                    Thus $A\mathrel{R_3}\apply{r}{A}$ by \cref{choice_function}.
                    Thus $\apply{r}{A}\in\apply{T}{\apply{p}{A}}$ by \cref{pair_eq_iff}.
                    Thus $V\in\apply{T}{i}$.
                \end{subproof}
                Thus $H\subseteq\apply{T}{i}$ by \cref{subseteq}.
            \end{subproof}
            Show that $H$ is finite.
            \begin{subproof}
                $F$ is finite.
                Thus $\img{r}{F}$ is finite by \cref{finite_image_function}.
                Show that $H\subseteq\img{r}{F}$.
                \begin{subproof}
                    Show that for all $V$ such that $V\in H$ we have $V\in\img{r}{F}$.
                    \begin{subproof}
                        Fix $V$.
                        Assume $V\in H$.
                        Thus $V\in\img{r}{\preimg{p}{\{i\}}}$.
                        Take $A$ such that $A\in\preimg{p}{\{i\}}$ and $(A,V)\in r$ by \cref{img_iff}.
                        Take $y\in\{i\}$ such that $A\mathrel{p} y$ by \cref{preimg_iff}.
                        Thus $(A,y)\in p$.
                        $p$ is a function on $F$.
                        Thus $p$ is a function and $\dom{p} = F$.
                        Thus $A\in\dom{p}$ by \cref{dom_intro}.
                        Thus $A\in F$ by \cref{subseteq}.
                        Thus $V\in\img{r}{F}$ by \cref{img_iff}.
                    \end{subproof}
                    Thus $H\subseteq\img{r}{F}$ by \cref{subseteq}.
                \end{subproof}
                Thus $H$ is finite by \cref{subseteq_finite_is_finite}.
            \end{subproof}
            Show that $H$ is inhabited.
            \begin{subproof}
                Take $A$ such that $A\in F$ and $(A,i)\in p$ by \cref{img_iff}.
                $p$ is a function on $F$.
                Thus $p$ is a function and $\dom{p} = F$.
                Thus $A\in\dom{p}$ by \cref{subseteq}.
                Thus $\apply{p}{A} = i$ by \cref{function_apply_iff}.
                $i\in\{i\}$ by \cref{singleton_intro}.
                Thus $A\in\preimg{p}{\{i\}}$ by \cref{preimg_iff}.
                $r$ is a function on $F$.
                Thus $r$ is a function and $\dom{r} = F$.
                Thus $A\in\dom{r}$ by \cref{subseteq}.
                Thus $(A,\apply{r}{A})\in r$ by \cref{function_apply_iff}.
                Thus $\apply{r}{A}\in H$ by \cref{img_iff}.
            \end{subproof}
            Show that $\apply{T}{i}$ is a topology on $\apply{X}{i}$.
            \begin{subproof}
                $\dom{X} = \naturalsPlus$.
                Thus $i\in\dom{X}$ by \cref{subseteq}.
                Thus $\apply{T}{i}$ is a topology on $\apply{X}{i}$.
            \end{subproof}
            Thus $\apply{v}{i}\in\apply{T}{i}$ by \cref{finite_intersections_open_in_topology}.
        \end{subproof}

        Show that for all $i\in J$ we have $\apply{x}{i}\in\apply{v}{i}$.
        \begin{subproof}
            Fix $i\in J$.
            Let $H = \img{r}{\preimg{p}{\{i\}}}$.
            Thus $i\mathrel{R_4}\apply{v}{i}$ by \cref{choice_function}.
            Thus $\apply{v}{i} = \inters{H}$ by \cref{pair_eq_iff}.
            Show that for all $V\in H$ we have $\apply{x}{i}\in V$.
            \begin{subproof}
                Fix $V\in H$.
                Thus $V\in\img{r}{\preimg{p}{\{i\}}}$.
                Take $A$ such that $A\in\preimg{p}{\{i\}}$ and $(A,V)\in r$ by \cref{img_iff}.
                Take $y\in\{i\}$ such that $A\mathrel{p} y$ by \cref{preimg_iff}.
                Thus $(A,y)\in p$.
                $p$ is a function on $F$.
                Thus $p$ is a function and $\dom{p} = F$.
                Thus $A\in\dom{p}$ by \cref{dom_intro}.
                Thus $A\in F$ by \cref{subseteq}.
                Thus $y = i$ by \cref{singleton_elim}.
                Thus $\apply{p}{A} = y$ by \cref{function_apply_iff}.
                Thus $\apply{p}{A} = i$.
                $r$ is a function on $F$.
                Thus $r$ is a function and $\dom{r} = F$.
                Thus $A\in\dom{r}$ by \cref{subseteq}.
                Thus $\apply{r}{A} = V$ by \cref{function_apply_iff}.
                Thus $A\mathrel{R_3}\apply{r}{A}$ by \cref{choice_function}.
                Thus $A = \preimg{\piIndex{X}{\apply{p}{A}}}{\apply{r}{A}}$ by \cref{pair_eq_iff}.
                Thus $A = \preimg{\piIndex{X}{i}}{V}$.
                Thus $B = \inters{F}$.
                Thus $B\subseteq A$ by \cref{inters_subseteq_elem}.
                Thus $x\in A$ by \cref{subseteq}.
                Take $y\in V$ such that $(x,y)\in\piIndex{X}{i}$ by \cref{preimg_iff}.
                Take $z\in\productFamily{X}$ such that $(z,\apply{z}{i}) = (x,y)$ by \cref{projection_map_indexed}.
                Thus $z = x$ and $\apply{z}{i} = y$ by \cref{pair_eq_iff}.
                Thus $\apply{x}{i} = y$.
                Thus $\apply{x}{i}\in V$.
            \end{subproof}
            Show that $H$ is inhabited.
            \begin{subproof}
                Take $A$ such that $A\in F$ and $(A,i)\in p$ by \cref{img_iff}.
                $p$ is a function on $F$.
                Thus $p$ is a function and $\dom{p} = F$.
                Thus $A\in\dom{p}$ by \cref{subseteq}.
                Thus $\apply{p}{A} = i$ by \cref{function_apply_iff}.
                $i\in\{i\}$ by \cref{singleton_intro}.
                Thus $A\in\preimg{p}{\{i\}}$ by \cref{preimg_iff}.
                $r$ is a function on $F$.
                Thus $r$ is a function and $\dom{r} = F$.
                Thus $A\in\dom{r}$ by \cref{subseteq}.
                Thus $(A,\apply{r}{A})\in r$ by \cref{function_apply_iff}.
                Thus $\apply{r}{A}\in H$ by \cref{img_iff}.
            \end{subproof}
            Thus $\apply{x}{i}\in\inters{H}$ by \cref{inters_intro}.
            Thus $\apply{x}{i}\in\apply{v}{i}$.
        \end{subproof}

        Let $R_5 = \{(i,k) \mid i\in J \mid
            \text{$k\in\naturalsPlus$ and $\apply{x}{i}\in\apply{\apply{b}{i}}{k}$ and $\apply{\apply{b}{i}}{k}\subseteq\apply{v}{i}$}\}$.
        Show that $R_5$ is a relation.
        \begin{subproof}
            Show that for all $w$ such that $w\in R_5$ there exist $i,k$ such that $w = (i,k)$.
            \begin{subproof}
                Fix $w$.
                Assume $w\in R_5$.
                Thus there exist $i,k$ such that $w = (i,k)$ by \cref{pair_eq_iff}.
            \end{subproof}
            Thus $R_5$ is a relation by \cref{relation}.
        \end{subproof}
        Show that for all $i\in J$ there exists $k$ such that $i\mathrel{R_5} k$.
        \begin{subproof}
            Fix $i\in J$.
            Show that $\apply{v}{i}$ is open in $\apply{X}{i}$ with respect to $\apply{T}{i}$.
            \begin{subproof}
                $\apply{v}{i}\in\apply{T}{i}$.
                Thus $\apply{T}{i}$ is a topology on $\apply{X}{i}$.
                Thus $\apply{v}{i}$ is open in $\apply{X}{i}$ with respect to $\apply{T}{i}$ by \cref{open_in}.
            \end{subproof}
            $\apply{x}{i}\in\apply{v}{i}$.
            $i\mathrel{R}\apply{b}{i}$ by \cref{choice_function}.
            Thus $\apply{b}{i}$ is a sequence in $\pow{\apply{X}{i}}$ and
                for all $k\in\naturalsPlus$ we have $\apply{\apply{b}{i}}{k}$ is open in $\apply{X}{i}$ with respect to $\apply{T}{i}$ and
                for all $y\in\apply{X}{i}$ we have
                    for all $U$ such that $U$ is open in $\apply{X}{i}$ with respect to $\apply{T}{i}$ and $y\in U$
                        there exists $k\in\naturalsPlus$ such that $y\in\apply{\apply{b}{i}}{k}$ and $\apply{\apply{b}{i}}{k}\subseteq U$
                by \cref{pair_eq_iff}.
            Thus there exists $k\in\naturalsPlus$ such that
                $\apply{x}{i}\in\apply{\apply{b}{i}}{k}$ and $\apply{\apply{b}{i}}{k}\subseteq\apply{v}{i}$.
            Take $k\in\naturalsPlus$ such that
                $\apply{x}{i}\in\apply{\apply{b}{i}}{k}$ and $\apply{\apply{b}{i}}{k}\subseteq\apply{v}{i}$.
            Thus $(i,k)\in R_5$ by \cref{pair_eq_iff}.
            Thus $i\mathrel{R_5} k$ by \cref{pair_eq_iff}.
        \end{subproof}
        Take $s$ such that $s$ is a function on $J$ and for all $i\in J$ we have $i\mathrel{R_5}\apply{s}{i}$
            by \cref{choice_function}.

        $\naturalsPlus\subseteq\reals$ by \cref{real_naturals_subset_reals}.
        Thus $J\subseteq\reals$ by \cref{subseteq_transitive}.
        Take $m$ such that $m$ is a largest element of $J$ with respect to $\realOrderRel$
            by \cref{finite_inhabited_largest_element,real_order_relation_simple}.
        Thus $m\in J$ by \cref{largest_element}.
        Thus $m\in\naturalsPlus$ by \cref{subseteq}.

        Let $R_6 = \{(i,k) \mid i\in\initialSegment{m} \mid
            \text{$k\in\naturalsPlus$ and
            $((i\in J \land \apply{x}{i}\in\apply{\apply{b}{i}}{k} \land \apply{\apply{b}{i}}{k}\subseteq\apply{v}{i})
            \lor (i\notin J \land \apply{x}{i}\in\apply{\apply{b}{i}}{k}))$}\}$.
        Show that $R_6$ is a relation.
        \begin{subproof}
            Show that for all $w$ such that $w\in R_6$ there exist $i,k$ such that $w = (i,k)$.
            \begin{subproof}
                Fix $w$.
                Assume $w\in R_6$.
                Thus there exist $i,k$ such that $w = (i,k)$ by \cref{pair_eq_iff}.
            \end{subproof}
            Thus $R_6$ is a relation by \cref{relation}.
        \end{subproof}
        Show that for all $i\in\initialSegment{m}$ there exists $k$ such that $i\mathrel{R_6} k$.
        \begin{subproof}
            Fix $i\in\initialSegment{m}$.
            \begin{byCase}
            \caseOf{$i\in J$.}
                Thus $i\mathrel{R_5}\apply{s}{i}$ by \cref{choice_function}.
                Thus $\apply{s}{i}\in\naturalsPlus$ and
                    $\apply{x}{i}\in\apply{\apply{b}{i}}{\apply{s}{i}}$ and
                    $\apply{\apply{b}{i}}{\apply{s}{i}}\subseteq\apply{v}{i}$ by \cref{pair_eq_iff}.
                Thus $(i,\apply{s}{i})\in R_6$ by \cref{pair_eq_iff}.
                Thus $i\mathrel{R_6}\apply{s}{i}$ by \cref{pair_eq_iff}.
            \caseOf{$i\notin J$.}
                Show that $\apply{X}{i}$ is open in $\apply{X}{i}$ with respect to $\apply{T}{i}$.
                \begin{subproof}
                    $\dom{X} = \naturalsPlus$.
                    Thus $i\in\dom{X}$ by \cref{subseteq}.
                    Thus $\apply{T}{i}$ is a topology on $\apply{X}{i}$.
                    Thus $\apply{X}{i}\in\apply{T}{i}$ by \cref{topology_on}.
                    Thus $\apply{X}{i}$ is open in $\apply{X}{i}$ with respect to $\apply{T}{i}$ by \cref{open_in}.
                \end{subproof}
                $x\in\productFamily{X}$.
                Thus $x\in\funs{\dom{X}}{\unions{\ran{X}}}$ by \cref{indexed_product}.
                Thus $x$ is a function and $\dom{x} = \dom{X}$ by \cref{funs_elim}.
                Thus $i\in\dom{x}$ by \cref{subseteq}.
                Thus $\apply{x}{i}\in\apply{X}{i}$ by \cref{indexed_product}.
                $i\mathrel{R}\apply{b}{i}$ by \cref{choice_function}.
                Thus $\apply{b}{i}$ is a sequence in $\pow{\apply{X}{i}}$ and
                    for all $k\in\naturalsPlus$ we have $\apply{\apply{b}{i}}{k}$ is open in $\apply{X}{i}$ with respect to $\apply{T}{i}$ and
                    for all $y\in\apply{X}{i}$ we have
                        for all $U$ such that $U$ is open in $\apply{X}{i}$ with respect to $\apply{T}{i}$ and $y\in U$
                            there exists $k\in\naturalsPlus$ such that $y\in\apply{\apply{b}{i}}{k}$ and $\apply{\apply{b}{i}}{k}\subseteq U$
                    by \cref{pair_eq_iff}.
                Thus there exists $k\in\naturalsPlus$ such that
                    $\apply{x}{i}\in\apply{\apply{b}{i}}{k}$ and $\apply{\apply{b}{i}}{k}\subseteq\apply{X}{i}$.
                Take $k\in\naturalsPlus$ such that $\apply{x}{i}\in\apply{\apply{b}{i}}{k}$ and $\apply{\apply{b}{i}}{k}\subseteq\apply{X}{i}$.
                Thus $(i,k)\in R_6$ by \cref{pair_eq_iff}.
                Thus $i\mathrel{R_6} k$ by \cref{pair_eq_iff}.
            \end{byCase}
        \end{subproof}
        Take $u$ such that $u$ is a function on $\initialSegment{m}$ and for all $i\in\initialSegment{m}$ we have
            $i\mathrel{R_6}\apply{u}{i}$ by \cref{choice_function}.
        Show that $u\in\funs{\initialSegment{m}}{\naturalsPlus}$.
        \begin{subproof}
            Show that for all $i\in\dom{u}$ we have $u(i)\in\naturalsPlus$.
            \begin{subproof}
                Fix $i\in\dom{u}$.
                Thus $i\in\initialSegment{m}$ by \cref{subseteq}.
                Thus $i\mathrel{R_6}\apply{u}{i}$ by \cref{choice_function}.
                Thus $\apply{u}{i}\in\naturalsPlus$ by \cref{pair_eq_iff}.
            \end{subproof}
            Thus $u\in\funs{\initialSegment{m}}{\naturalsPlus}$ by \cref{funs_intro}.
        \end{subproof}
        Thus $u\in S$ by \cref{unions_intro}.
        Take $n\in\naturalsPlus$ such that $\apply{q}{n} = u$ by \cref{naturalsplus_finite_sequence_enumeration}.

        Show that $x\in d(n)$.
        \begin{subproof}
            Let $F_n = \img{\apply{h}{n}}{\dom{\apply{q}{n}}}$.
            Thus $\apply{q}{n} = u$.
            Thus $\dom{\apply{q}{n}} = \initialSegment{m}$ by \cref{funs_elim}.
            Show that for all $i\in\initialSegment{m}$ we have
                $x\in\preimg{\piIndex{X}{i}}{\apply{\apply{b}{i}}{\apply{u}{i}}}$.
            \begin{subproof}
                Fix $i\in\initialSegment{m}$.
                Thus $i\mathrel{R_6}\apply{u}{i}$ by \cref{choice_function}.
                Thus $\apply{x}{i}\in\apply{\apply{b}{i}}{\apply{u}{i}}$ by \cref{pair_eq_iff}.
                Thus $x\in\preimg{\piIndex{X}{i}}{\apply{\apply{b}{i}}{\apply{u}{i}}}$ by \cref{preimg_iff,projection_map_indexed}.
            \end{subproof}
            Show that for all $A\in F_n$ we have $x\in A$.
            \begin{subproof}
                Fix $A\in F_n$.
                Take $i$ such that $i\in\dom{\apply{q}{n}}$ and $(i,A)\in\apply{h}{n}$ by \cref{img_iff}.
                Thus $i\in\initialSegment{m}$ by \cref{subseteq}.
                Thus $n\mathrel{R_d}\apply{h}{n}$ by \cref{choice_function}.
                Thus $(n,\apply{h}{n})\in R_d$.
                Take $n_1,g_1$ such that $n_1\in\naturalsPlus$ and
                    $g_1$ is a function on $\dom{\apply{q}{n_1}}$ and
                    for all $j\in\dom{\apply{q}{n_1}}$ we have
                        $(j,\preimg{\piIndex{X}{j}}{\apply{\apply{b}{j}}{\apply{\apply{q}{n_1}}{j}}})\in g_1$ and
                    $(n,\apply{h}{n}) = (n_1,g_1)$ by \cref{pair_eq_iff}.
                Thus $n_1 = n$ by \cref{pair_eq_iff}.
                Thus $g_1 = \apply{h}{n}$ by \cref{pair_eq_iff}.
                Thus $\apply{h}{n}$ is a function on $\dom{\apply{q}{n}}$ and
                    for all $j\in\dom{\apply{q}{n}}$ we have
                        $(j,\preimg{\piIndex{X}{j}}{\apply{\apply{b}{j}}{\apply{\apply{q}{n}}{j}}})\in\apply{h}{n}$.
                Thus $\apply{h}{n}$ is a function and $\dom{\apply{h}{n}} = \dom{\apply{q}{n}}$.
                Show that for all $j\in\dom{\apply{q}{n}}$ we have
                    $\apply{\apply{h}{n}}{j} = \preimg{\piIndex{X}{j}}{\apply{\apply{b}{j}}{\apply{\apply{q}{n}}{j}}}$.
                \begin{subproof}
                    Fix $j\in\dom{\apply{q}{n}}$.
                    Thus $(j,\preimg{\piIndex{X}{j}}{\apply{\apply{b}{j}}{\apply{\apply{q}{n}}{j}}})\in\apply{h}{n}$.
                    Thus $j\in\dom{\apply{h}{n}}$ by \cref{dom_intro}.
                    Thus $\apply{\apply{h}{n}}{j} = \preimg{\piIndex{X}{j}}{\apply{\apply{b}{j}}{\apply{\apply{q}{n}}{j}}}$ by \cref{function_apply_iff}.
                \end{subproof}
                Thus $i\in\dom{\apply{h}{n}}$ by \cref{dom_intro}.
                Thus $\apply{\apply{h}{n}}{i} = A$ by \cref{function_apply_iff}.
                Thus $A = \preimg{\piIndex{X}{i}}{\apply{\apply{b}{i}}{\apply{\apply{q}{n}}{i}}}$.
                Thus $\apply{\apply{q}{n}}{i} = \apply{u}{i}$.
                Thus $x\in A$ by \cref{subseteq}.
            \end{subproof}
            Show that $F_n$ is inhabited.
            \begin{subproof}
                $m\in\initialSegment{m}$ by \cref{initial_segment_naturalsplus}.
                Thus $m\in\dom{\apply{q}{n}}$ by \cref{subseteq}.
                Thus $n\mathrel{R_d}\apply{h}{n}$ by \cref{choice_function}.
                Thus $(n,\apply{h}{n})\in R_d$.
                Take $n_1,g_1$ such that $n_1\in\naturalsPlus$ and
                    $g_1$ is a function on $\dom{\apply{q}{n_1}}$ and
                    for all $j\in\dom{\apply{q}{n_1}}$ we have
                        $(j,\preimg{\piIndex{X}{j}}{\apply{\apply{b}{j}}{\apply{\apply{q}{n_1}}{j}}})\in g_1$ and
                    $(n,\apply{h}{n}) = (n_1,g_1)$ by \cref{pair_eq_iff}.
                Thus $n_1 = n$ by \cref{pair_eq_iff}.
                Thus $g_1 = \apply{h}{n}$ by \cref{pair_eq_iff}.
                Thus $\apply{h}{n}$ is a function and $\dom{\apply{h}{n}} = \dom{\apply{q}{n}}$.
                Thus $m\in\dom{\apply{h}{n}}$ by \cref{subseteq}.
                Thus $(m,\apply{\apply{h}{n}}{m})\in\apply{h}{n}$ by \cref{function_apply_iff}.
                Thus $\apply{\apply{h}{n}}{m}\in F_n$ by \cref{img_iff}.
            \end{subproof}
            Thus $x\in\inters{F_n}$ by \cref{inters_intro}.
            $d\in\funs{\naturalsPlus}{\pow{\productFamily{X}}}$ by \cref{sequence_in}.
            Thus $d$ is a function by \cref{funs_is_function}.
            Thus $\dom{d} = \naturalsPlus$ by \cref{funs_elim}.
            Thus $n\in\dom{d}$ by \cref{subseteq}.
            $(n,\inters{F_n})\in d$ by \cref{pair_eq_iff}.
            Thus $d(n) = \inters{F_n}$ by \cref{function_apply_iff}.
            Thus $x\in d(n)$.
        \end{subproof}

        Show that $d(n)\subseteq U$.
        \begin{subproof}
            Let $F_n = \img{\apply{h}{n}}{\dom{\apply{q}{n}}}$.
            Thus $\apply{q}{n} = u$.
            Thus $\dom{\apply{q}{n}} = \initialSegment{m}$ by \cref{funs_elim}.
            Show that for all $A\in F$ we have $d(n)\subseteq A$.
            \begin{subproof}
                Fix $A\in F$.
                Thus $A\mathrel{R_2}\apply{p}{A}$ by \cref{choice_function}.
                $p$ is a function on $F$.
                Thus $p$ is a function and $\dom{p} = F$.
                Thus $A\in\dom{p}$ by \cref{subseteq}.
                Thus $(A,\apply{p}{A})\in p$ by \cref{function_apply_elim}.
                Thus $\apply{p}{A}\in J$ by \cref{img_iff}.
                Let $i = \apply{p}{A}$.
                Thus $i\in J$.
                Thus $i\in\naturalsPlus$ by \cref{subseteq}.
                Show that $i\leq m$.
                \begin{subproof}
                    Thus $i\mathrel{\realOrderRel} m$ or $i = m$ by \cref{largest_element}.
                    \begin{byCase}
                    \caseOf{$i\mathrel{\realOrderRel} m$.}
                        Thus $i < m$ by \cref{real_order_relation_elim}.
                        Thus $i\leq m$.
                    \caseOf{$i = m$.}
                        Thus $i\leq m$.
                    \end{byCase}
                \end{subproof}
                Thus $i\in\initialSegment{m}$ by \cref{initial_segment_naturalsplus}.
                Thus $i\mathrel{R_6}\apply{u}{i}$ by \cref{choice_function}.
                Thus $\apply{u}{i}\in\naturalsPlus$ and
                    $\apply{x}{i}\in\apply{\apply{b}{i}}{\apply{u}{i}}$ and
                    $\apply{\apply{b}{i}}{\apply{u}{i}}\subseteq\apply{v}{i}$ by \cref{pair_eq_iff}.
                Thus $A\mathrel{R_3}\apply{r}{A}$ by \cref{choice_function}.
                Thus $A = \preimg{\piIndex{X}{i}}{\apply{r}{A}}$ by \cref{pair_eq_iff}.
                Let $H = \img{r}{\preimg{p}{\{i\}}}$.
                Thus $i\mathrel{R_4}\apply{v}{i}$ by \cref{choice_function}.
                Thus $\apply{v}{i} = \inters{H}$ by \cref{pair_eq_iff}.
                $p$ is a function on $F$.
                Thus $p$ is a function and $\dom{p} = F$.
                Thus $A\in\dom{p}$ by \cref{subseteq}.
                Thus $(A,\apply{p}{A})\in p$ by \cref{function_apply_elim}.
                Thus $(A,i)\in p$.
                $i\in\{i\}$ by \cref{singleton_intro}.
                Thus $A\in\preimg{p}{\{i\}}$ by \cref{preimg_iff}.
                $r$ is a function on $F$.
                Thus $r$ is a function and $\dom{r} = F$.
                Thus $A\in\dom{r}$ by \cref{subseteq}.
                Thus $(A,\apply{r}{A})\in r$ by \cref{function_apply_iff}.
                Thus $\apply{r}{A}\in H$ by \cref{img_iff}.
                Thus $\apply{v}{i}\subseteq\apply{r}{A}$ by \cref{inters_subseteq_elem}.
                Thus $\apply{\apply{b}{i}}{\apply{u}{i}}\subseteq\apply{r}{A}$ by \cref{subseteq_transitive}.
                Thus $\preimg{\piIndex{X}{i}}{\apply{\apply{b}{i}}{\apply{u}{i}}}\subseteq A$ by \cref{preimg_subseteq,subseteq}.
                Thus $n\mathrel{R_d}\apply{h}{n}$ by \cref{choice_function}.
                Thus $(n,\apply{h}{n})\in R_d$.
                Take $n_1,g_1$ such that $n_1\in\naturalsPlus$ and
                    $g_1$ is a function on $\dom{\apply{q}{n_1}}$ and
                    for all $j\in\dom{\apply{q}{n_1}}$ we have
                        $(j,\preimg{\piIndex{X}{j}}{\apply{\apply{b}{j}}{\apply{\apply{q}{n_1}}{j}}})\in g_1$ and
                    $(n,\apply{h}{n}) = (n_1,g_1)$ by \cref{pair_eq_iff}.
                Thus $n_1 = n$ by \cref{pair_eq_iff}.
                Thus $g_1 = \apply{h}{n}$ by \cref{pair_eq_iff}.
                Thus $\apply{h}{n}$ is a function on $\dom{\apply{q}{n}}$ and
                    for all $j\in\dom{\apply{q}{n}}$ we have
                        $(j,\preimg{\piIndex{X}{j}}{\apply{\apply{b}{j}}{\apply{\apply{q}{n}}{j}}})\in\apply{h}{n}$.
                Thus $\apply{h}{n}$ is a function and $\dom{\apply{h}{n}} = \dom{\apply{q}{n}}$.
                Show that for all $j\in\dom{\apply{q}{n}}$ we have
                    $\apply{\apply{h}{n}}{j} = \preimg{\piIndex{X}{j}}{\apply{\apply{b}{j}}{\apply{\apply{q}{n}}{j}}}$.
                \begin{subproof}
                    Fix $j\in\dom{\apply{q}{n}}$.
                    Thus $(j,\preimg{\piIndex{X}{j}}{\apply{\apply{b}{j}}{\apply{\apply{q}{n}}{j}}})\in\apply{h}{n}$.
                    Thus $j\in\dom{\apply{h}{n}}$ by \cref{dom_intro}.
                    Thus $\apply{\apply{h}{n}}{j} = \preimg{\piIndex{X}{j}}{\apply{\apply{b}{j}}{\apply{\apply{q}{n}}{j}}}$ by \cref{function_apply_iff}.
                \end{subproof}
                Thus $i\in\dom{\apply{h}{n}}$ by \cref{subseteq}.
                Thus $\apply{\apply{h}{n}}{i} = \preimg{\piIndex{X}{i}}{\apply{\apply{b}{i}}{\apply{\apply{q}{n}}{i}}}$.
                Thus $\apply{\apply{q}{n}}{i} = \apply{u}{i}$.
                Thus $\preimg{\piIndex{X}{i}}{\apply{\apply{b}{i}}{\apply{u}{i}}}\in F_n$ by \cref{img_iff}.
                $d\in\funs{\naturalsPlus}{\pow{\productFamily{X}}}$ by \cref{sequence_in}.
                Thus $d$ is a function by \cref{funs_is_function}.
                Thus $\dom{d} = \naturalsPlus$ by \cref{funs_elim}.
                Thus $n\in\dom{d}$ by \cref{subseteq}.
                $(n,\inters{F_n})\in d$ by \cref{pair_eq_iff}.
                Thus $d(n) = \inters{F_n}$ by \cref{function_apply_iff}.
                Thus $d(n)\subseteq \preimg{\piIndex{X}{i}}{\apply{\apply{b}{i}}{\apply{u}{i}}}$ by \cref{inters_subseteq_elem}.
                Thus $d(n)\subseteq A$ by \cref{subseteq_transitive}.
            \end{subproof}
            Thus $d(n)\subseteq \inters{F}$ by \cref{subseteq_inters_iff}.
            Thus $d(n)\subseteq B$ by \cref{subseteq}.
            Thus $d(n)\subseteq U$ by \cref{subseteq_transitive}.
        \end{subproof}
    \end{subproof}
    Thus $\productFamily{X}$ satisfies the second countability axiom with respect to $T_0$ by \cref{second_countability_axiom}.
\end{proof}

% from §30 Item 11: dense subset
\begin{definition}\label{dense_in}
    $A$ is dense in $X$ with respect to $T$ iff $\closure{A}{X}{T} = X$.
\end{definition}

% from §30 Item 12: countable subcover in a second-countable space
\begin{lemma}\label{second_countable_open_covering_countable_subcover}
    Assume $T$ is a topology on $X$.
    Suppose $X$ satisfies the second countability axiom with respect to $T$.
    Then for all $C$ such that $C$ is an open covering of $X$ with respect to $T$
        there exists $J$ such that
            $J\subseteq\naturalsPlus$ and
            there exists $f$ such that
                $f$ is a function on $J$ and
                for all $j\in J$ we have $\apply{f}{j}\in C$ and
                $X\subseteq\unions{\img{f}{J}}$.
\end{lemma}
\begin{proof}
    Take $b$ such that
        $b$ is a sequence in $\pow{X}$ and
        for all $n\in\naturalsPlus$ we have $b(n)$ is open in $X$ with respect to $T$ and
        for all $x\in X$ we have
            for all $U$ such that $U$ is open in $X$ with respect to $T$ and $x\in U$
                there exists $n\in\naturalsPlus$ such that $x\in b(n)$ and $b(n)\subseteq U$
        by \cref{second_countability_axiom}.
    Fix $C$.
    Assume $C$ is an open covering of $X$ with respect to $T$.
    Let $J = \{n\in\naturalsPlus \mid \exists A\in C. b(n)\subseteq A\}$.
    Let $R = \{(n,A) \mid n\in J \mid A\in C \land b(n)\subseteq A\}$.
    Show that $R$ is a relation.
    \begin{subproof}
        Show that for all $w$ such that $w\in R$ there exist $n,A$ such that $w = (n,A)$.
        \begin{subproof}
            Fix $w$.
            Assume $w\in R$.
            Thus there exist $n,A$ such that $w = (n,A)$ by \cref{pair_eq_iff}.
        \end{subproof}
        Thus $R$ is a relation by \cref{relation}.
    \end{subproof}
    Show that for all $n\in J$ there exists $A$ such that $n\mathrel{R} A$.
    \begin{subproof}
        Fix $n\in J$.
        Thus there exists $A\in C$ such that $b(n)\subseteq A$.
        Take $A\in C$ such that $b(n)\subseteq A$.
        Thus $(n,A)\in R$ by \cref{pair_eq_iff}.
        Thus $n\mathrel{R} A$ by \cref{pair_eq_iff}.
    \end{subproof}
    Take $f$ such that $f$ is a function on $J$ and for all $n\in J$ we have $n\mathrel{R}\apply{f}{n}$
        by \cref{choice_function}.
    Show that for all $j\in J$ we have $\apply{f}{j}\in C$.
    \begin{subproof}
        Fix $j\in J$.
        Thus $j\mathrel{R}\apply{f}{j}$ by \cref{choice_function}.
        Thus $\apply{f}{j}\in C$ by \cref{pair_eq_iff}.
    \end{subproof}
    Show that $X\subseteq\unions{\img{f}{J}}$.
    \begin{subproof}
        Show that for all $x$ such that $x\in X$ we have $x\in\unions{\img{f}{J}}$.
        \begin{subproof}
            Fix $x$.
            Assume $x\in X$.
            $C$ is a covering of $X$ by \cref{open_covering}.
            Thus $X\subseteq\unions{C}$ by \cref{covering}.
            Thus $x\in\unions{C}$ by \cref{subseteq}.
            Take $A_0$ such that $A_0\in C$ and $x\in A_0$ by \cref{unions_iff}.
            $A_0$ is open in $X$ with respect to $T$ by \cref{open_covering}.
            Thus there exists $n\in\naturalsPlus$ such that $x\in b(n)$ and $b(n)\subseteq A_0$.
            Take $n\in\naturalsPlus$ such that $x\in b(n)$ and $b(n)\subseteq A_0$.
            Thus $n\in J$.
            Thus $n\mathrel{R}\apply{f}{n}$ by \cref{choice_function}.
            Thus $\apply{f}{n}\in C$ and $b(n)\subseteq\apply{f}{n}$ by \cref{pair_eq_iff}.
            Thus $x\in\apply{f}{n}$ by \cref{subseteq}.
            $f$ is a function on $J$ by \cref{choice_function}.
            Thus $f$ is a function and $\dom{f} = J$.
            Thus $n\in\dom{f}$ by \cref{subseteq}.
            Thus $(n,\apply{f}{n})\in f$ by \cref{function_apply_iff}.
            Thus $\apply{f}{n}\in\img{f}{J}$ by \cref{img_iff}.
            Thus $x\in\unions{\img{f}{J}}$ by \cref{unions_intro}.
        \end{subproof}
        Thus $X\subseteq\unions{\img{f}{J}}$ by \cref{subseteq}.
    \end{subproof}
    Thus there exists $J$ such that
        $J\subseteq\naturalsPlus$ and
        there exists $f$ such that
            $f$ is a function on $J$ and
            for all $j\in J$ we have $\apply{f}{j}\in C$ and
            $X\subseteq\unions{\img{f}{J}}$.
\end{proof}

% from §30 Item 13: countable dense subset in a second-countable space
\begin{lemma}\label{second_countable_countable_dense_subset}
    Assume $T$ is a topology on $X$.
    Suppose $X$ satisfies the second countability axiom with respect to $T$.
    Then there exists $J$ such that
        $J\subseteq\naturalsPlus$ and
        there exists $f$ such that
            $f$ is a function on $J$ and
            for all $j\in J$ we have $\apply{f}{j}\in X$ and
            $\img{f}{J}$ is dense in $X$ with respect to $T$.
\end{lemma}
\begin{proof}
    Take $b$ such that
        $b$ is a sequence in $\pow{X}$ and
        for all $n\in\naturalsPlus$ we have $b(n)$ is open in $X$ with respect to $T$ and
        for all $x\in X$ we have
            for all $U$ such that $U$ is open in $X$ with respect to $T$ and $x\in U$
                there exists $n\in\naturalsPlus$ such that $x\in b(n)$ and $b(n)\subseteq U$
        by \cref{second_countability_axiom}.
    Let $J = \{n\in\naturalsPlus \mid b(n)\neq\emptyset\}$.
    Let $R = \{(n,x) \mid n\in J \mid x\in b(n)\}$.
    Show that $R$ is a relation.
    \begin{subproof}
        Show that for all $w$ such that $w\in R$ there exist $n,x$ such that $w = (n,x)$.
        \begin{subproof}
            Fix $w$.
            Assume $w\in R$.
            Thus there exist $n,x$ such that $w = (n,x)$ by \cref{pair_eq_iff}.
        \end{subproof}
        Thus $R$ is a relation by \cref{relation}.
    \end{subproof}
    Show that for all $n\in J$ there exists $x$ such that $n\mathrel{R} x$.
    \begin{subproof}
        Fix $n\in J$.
        Thus $b(n)\neq\emptyset$ by \cref{pair_eq_iff}.
        Thus $b(n)$ is inhabited by \cref{nonempty_inhabited}.
        Take $x$ such that $x\in b(n)$.
        Thus $(n,x)\in R$ by \cref{pair_eq_iff}.
        Thus $n\mathrel{R} x$ by \cref{pair_eq_iff}.
    \end{subproof}
    Take $f$ such that $f$ is a function on $J$ and for all $n\in J$ we have $n\mathrel{R}\apply{f}{n}$
        by \cref{choice_function}.
    Show that for all $j\in J$ we have $\apply{f}{j}\in X$.
    \begin{subproof}
        Fix $j\in J$.
        Thus $j\mathrel{R}\apply{f}{j}$ by \cref{choice_function}.
        Thus $\apply{f}{j}\in b(j)$ by \cref{pair_eq_iff}.
        $b\in\funs{\naturalsPlus}{\pow{X}}$ by \cref{sequence_in}.
        Thus $b$ is a function to $\pow{X}$ and $\dom{b} = \naturalsPlus$ by \cref{funs_elim}.
        Thus $j\in\dom{b}$ by \cref{subseteq}.
        Thus $b(j)\in\pow{X}$.
        Thus $b(j)\subseteq X$ by \cref{powerset_elim}.
        Thus $\apply{f}{j}\in X$ by \cref{subseteq}.
    \end{subproof}
    Show that $\img{f}{J}$ is dense in $X$ with respect to $T$.
    \begin{subproof}
        Show that $\img{f}{J}\subseteq X$.
        \begin{subproof}
            Show that for all $y$ such that $y\in\img{f}{J}$ we have $y\in X$.
            \begin{subproof}
                Fix $y$.
                Assume $y\in\img{f}{J}$.
                Take $j$ such that $j\in J$ and $(j,y)\in f$ by \cref{img_iff}.
                $f$ is a function on $J$ by \cref{choice_function}.
                Thus $f$ is a function and $\dom{f} = J$.
                Thus $j\in\dom{f}$ by \cref{subseteq}.
                Thus $\apply{f}{j} = y$ by \cref{function_apply_iff}.
                Thus $\apply{f}{j}\in X$.
                Thus $y\in X$.
            \end{subproof}
            Thus $\img{f}{J}\subseteq X$ by \cref{subseteq}.
        \end{subproof}
        Show that for all $x\in X$ we have $x\in\closure{\img{f}{J}}{X}{T}$.
        \begin{subproof}
            Fix $x\in X$.
            Show that for all $U$ such that $U$ is open in $X$ with respect to $T$ and $x\in U$
                we have $U$ intersects $\img{f}{J}$.
            \begin{subproof}
                Fix $U$.
                Assume $U$ is open in $X$ with respect to $T$ and $x\in U$.
                Thus there exists $n\in\naturalsPlus$ such that $x\in b(n)$ and $b(n)\subseteq U$.
                Take $n\in\naturalsPlus$ such that $x\in b(n)$ and $b(n)\subseteq U$.
                Thus $b(n)\neq\emptyset$ by \cref{emptyset}.
                Thus $n\in J$.
                Thus $n\mathrel{R}\apply{f}{n}$ by \cref{choice_function}.
                Thus $\apply{f}{n}\in b(n)$ by \cref{pair_eq_iff}.
                Thus $\apply{f}{n}\in U$ by \cref{subseteq}.
                $f$ is a function on $J$ by \cref{choice_function}.
                Thus $f$ is a function and $\dom{f} = J$.
                Thus $n\in\dom{f}$ by \cref{subseteq}.
                Thus $(n,\apply{f}{n})\in f$ by \cref{function_apply_iff}.
                Thus $\apply{f}{n}\in\img{f}{J}$ by \cref{img_iff}.
                Thus $\apply{f}{n}\in U\inter\img{f}{J}$ by \cref{inter_intro}.
                Thus $U\inter\img{f}{J}\neq\emptyset$ by \cref{emptyset}.
                Thus $U$ intersects $\img{f}{J}$ by \cref{intersects}.
            \end{subproof}
            Thus $x\in\closure{\img{f}{J}}{X}{T}$ by \cref{closure_iff_intersects_open}.
        \end{subproof}
        Thus $X\subseteq\closure{\img{f}{J}}{X}{T}$ by \cref{subseteq}.
        $X$ is closed in $X$ with respect to $T$ by \cref{closed_whole_space}.
        Thus $\closure{\img{f}{J}}{X}{T}\subseteq X$ by \cref{closure_subseteq_closed}.
        Thus $\closure{\img{f}{J}}{X}{T} = X$ by \cref{subseteq_antisymmetric}.
        Thus $\img{f}{J}$ is dense in $X$ with respect to $T$ by \cref{dense_in}.
    \end{subproof}
    Thus there exists $J$ such that
        $J\subseteq\naturalsPlus$ and
        there exists $f$ such that
            $f$ is a function on $J$ and
            for all $j\in J$ we have $\apply{f}{j}\in X$ and
            $\img{f}{J}$ is dense in $X$ with respect to $T$.
\end{proof}

% from §30 Item 14: Lindelöf space
\begin{definition}\label{lindelof_space}
    $X$ is Lindelöf with respect to $T$ iff
        for all $C$ such that $C$ is an open covering of $X$ with respect to $T$
            there exists $J$ such that
                $J\subseteq\naturalsPlus$ and
                there exists $f$ such that
                    $f$ is a function on $J$ and
                    for all $j\in J$ we have $\apply{f}{j}\in C$ and
                    $X\subseteq\unions{\img{f}{J}}$.
\end{definition}

% from §30 Item 15: separable space
\begin{definition}\label{separable_space}
    $X$ is separable with respect to $T$ iff
        there exists $J$ such that
            $J\subseteq\naturalsPlus$ and
            there exists $f$ such that
                $f$ is a function on $J$ and
                for all $j\in J$ we have $\apply{f}{j}\in X$ and
                $\img{f}{J}$ is dense in $X$ with respect to $T$.
\end{definition}


% from §31 helper definition 1: closed singletons
\begin{definition}\label{one_point_closed}
    $X$ has closed singletons with respect to $T$ iff
        for all $x$ such that $x\in X$ we have $\{x\}$ is closed in $X$ with respect to $T$.
\end{definition}

% from §31 Item 1: regular space
\begin{definition}\label{regular_space}
    $X$ is regular with respect to $T$ iff
        $T$ is a topology on $X$ and
        $X$ has closed singletons with respect to $T$ and
        for all $x,B$ such that $x\in X$ and $B$ is closed in $X$ with respect to $T$ and $x\notin B$
            there exist $U,V$ such that
                $U$ is open in $X$ with respect to $T$ and
                $V$ is open in $X$ with respect to $T$ and
                $x\in U$ and
                $B\subseteq V$ and
                $U$ is disjoint from $V$.
\end{definition}

% from §31 Item 2: normal space
\begin{definition}\label{normal_space}
    $X$ is normal with respect to $T$ iff
        $T$ is a topology on $X$ and
        $X$ has closed singletons with respect to $T$ and
        for all $A,B$ such that $A$ is closed in $X$ with respect to $T$ and $B$ is closed in $X$ with respect to $T$
            and $A$ is disjoint from $B$
            there exist $U,V$ such that
                $U$ is open in $X$ with respect to $T$ and
                $V$ is open in $X$ with respect to $T$ and
                $A\subseteq U$ and
                $B\subseteq V$ and
                $U$ is disjoint from $V$.
\end{definition}

% from §31 helper lemma 1: regular implies Hausdorff
\begin{lemma}\label{regular_implies_hausdorff}
    Assume $X$ is regular with respect to $T$.
    Then $X$ is Hausdorff with respect to $T$.
\end{lemma}
\begin{proof}
    $T$ is a topology on $X$ by \cref{regular_space}.
    Show that for all $x_1,x_2$ such that $x_1\in X$ and $x_2\in X$ and $x_1\neq x_2$
        there exist $U_1,U_2$ such that
            $U_1$ is a neighborhood of $x_1$ in $X$ with respect to $T$ and
            $U_2$ is a neighborhood of $x_2$ in $X$ with respect to $T$ and
            $U_1$ is disjoint from $U_2$.
    \begin{subproof}
        Fix $x_1$.
        Fix $x_2$.
        Assume $x_1\in X$ and $x_2\in X$ and $x_1\neq x_2$.
        $X$ has closed singletons with respect to $T$ by \cref{regular_space}.
        Thus $\{x_2\}$ is closed in $X$ with respect to $T$ by \cref{one_point_closed}.
        Thus $x_1\notin\{x_2\}$.
        Take $U,V$ such that
            $U$ is open in $X$ with respect to $T$ and
            $V$ is open in $X$ with respect to $T$ and
            $x_1\in U$ and
            $\{x_2\}\subseteq V$ and
            $U$ is disjoint from $V$
            by \cref{regular_space}.
        $x_2\in\{x_2\}$ by \cref{singleton_intro}.
        Thus $x_2\in V$ by \cref{subseteq}.
        $U$ is a neighborhood of $x_1$ in $X$ with respect to $T$ by \cref{neighborhood}.
        $V$ is a neighborhood of $x_2$ in $X$ with respect to $T$ by \cref{neighborhood}.
    \end{subproof}
    Thus $X$ is Hausdorff with respect to $T$ by \cref{hausdorff}.
\end{proof}

% from §31 helper lemma 2: normal implies regular
\begin{lemma}\label{normal_implies_regular}
    Assume $X$ is normal with respect to $T$.
    Then $X$ is regular with respect to $T$.
\end{lemma}
\begin{proof}
    $T$ is a topology on $X$ by \cref{normal_space}.
    Show that for all $x$ such that $x\in X$ we have $\{x\}$ is closed in $X$ with respect to $T$.
    \begin{subproof}
        Fix $x$.
        Assume $x\in X$.
        $X$ has closed singletons with respect to $T$ by \cref{normal_space}.
        Thus $\{x\}$ is closed in $X$ with respect to $T$ by \cref{one_point_closed}.
    \end{subproof}
    Show that for all $x,B$ such that $x\in X$ and $B$ is closed in $X$ with respect to $T$ and $x\notin B$
        there exist $U,V$ such that
            $U$ is open in $X$ with respect to $T$ and
            $V$ is open in $X$ with respect to $T$ and
            $x\in U$ and
            $B\subseteq V$ and
            $U$ is disjoint from $V$.
    \begin{subproof}
        Fix $x$.
        Fix $B$.
        Assume $x\in X$ and $B$ is closed in $X$ with respect to $T$ and $x\notin B$.
        $X$ has closed singletons with respect to $T$ by \cref{normal_space}.
        Thus $\{x\}$ is closed in $X$ with respect to $T$ by \cref{one_point_closed}.
        Show that $\{x\}$ is disjoint from $B$.
        \begin{subproof}
            Suppose not.
            Then there exists $y$ such that $y\in\{x\}$ and $y\in B$ by \cref{disjoint}.
            Thus $y = x$ by \cref{singleton_elim}.
            Thus $x\in B$.
            Contradiction.
        \end{subproof}
        Take $U,V$ such that
            $U$ is open in $X$ with respect to $T$ and
            $V$ is open in $X$ with respect to $T$ and
            $\{x\}\subseteq U$ and
            $B\subseteq V$ and
            $U$ is disjoint from $V$
            by \cref{normal_space}.
        $x\in\{x\}$ by \cref{singleton_intro}.
        Thus $x\in U$ by \cref{subseteq}.
    \end{subproof}
    $X$ has closed singletons with respect to $T$ by \cref{one_point_closed}.
    Thus $X$ is regular with respect to $T$ by \cref{regular_space}.
\end{proof}

% from §31 Item 3: regular neighborhood closure property
\begin{lemma}\label{regular_closure_neighborhood}
    Assume $X$ is regular with respect to $T$.
    Then for all $x,U$ such that $x\in X$ and $U$ is a neighborhood of $x$ in $X$ with respect to $T$
        there exists $V$ such that
            $V$ is a neighborhood of $x$ in $X$ with respect to $T$ and
            $\closure{V}{X}{T}\subseteq U$.
\end{lemma}
\begin{proof}
    $T$ is a topology on $X$ by \cref{regular_space}.
    Fix $x$.
    Fix $U$.
    Assume $x\in X$ and $U$ is a neighborhood of $x$ in $X$ with respect to $T$.
    Let $B = X\setminus U$.
    Show that $B$ is closed in $X$ with respect to $T$.
    \begin{subproof}
        $B\subseteq X$ by \cref{setminus_subseteq}.
        $U$ is open in $X$ with respect to $T$ by \cref{neighborhood}.
        Thus $U\in T$ by \cref{open_in}.
        Thus $T\subseteq\pow{X}$ by \cref{topology_on}.
        Thus $U\in\pow{X}$ by \cref{subseteq}.
        Thus $U\subseteq X$ by \cref{pow_iff}.
        Thus $X\setminus B = U$ by \cref{double_relative_complement,subseteq}.
        Thus $X\setminus B$ is open in $X$ with respect to $T$ by \cref{open_in}.
        Thus $B$ is closed in $X$ with respect to $T$ by \cref{closed_in}.
    \end{subproof}
    Show that $x\notin B$.
    \begin{subproof}
        Thus $x\in U$ by \cref{neighborhood}.
        Suppose not.
        Then $x\in B$.
        Thus $x\in X\setminus U$.
        Thus $x\notin U$ by \cref{setminus_elim_right}.
        Contradiction.
    \end{subproof}
    Take $U_1,V_1$ such that
        $U_1$ is open in $X$ with respect to $T$ and
        $V_1$ is open in $X$ with respect to $T$ and
        $x\in U_1$ and
        $B\subseteq V_1$ and
        $U_1$ is disjoint from $V_1$
        by \cref{regular_space}.
    Let $V = U_1$.
    Show that $V$ is a neighborhood of $x$ in $X$ with respect to $T$.
    \begin{subproof}
        $V$ is open in $X$ with respect to $T$ by \cref{open_in}.
        Thus $V$ is a neighborhood of $x$ in $X$ with respect to $T$ by \cref{neighborhood}.
    \end{subproof}
    Show that $\closure{V}{X}{T}\subseteq U$.
    \begin{subproof}
        Show that $V\subseteq X\setminus V_1$.
        \begin{subproof}
            Show that for all $y$ such that $y\in V$ we have $y\in X\setminus V_1$.
            \begin{subproof}
                Fix $y$.
                Assume $y\in V$.
                Show that $y\notin V_1$.
                \begin{subproof}
                    Suppose not.
                    Then $y\in V_1$.
                    Thus there exists $z$ such that $z\in V$ and $z\in V_1$.
                    Thus $V$ is not disjoint from $V_1$ by \cref{disjoint}.
                    Contradiction.
                \end{subproof}
                $V$ is open in $X$ with respect to $T$ by \cref{open_in}.
                Thus $V\in T$ by \cref{open_in}.
                Thus $T\subseteq\pow{X}$ by \cref{topology_on}.
                Thus $V\in\pow{X}$ by \cref{subseteq}.
                Thus $V\subseteq X$ by \cref{pow_iff}.
                Thus $y\in X$ by \cref{subseteq}.
                Thus $y\in X\setminus V_1$ by \cref{setminus_intro}.
            \end{subproof}
            Thus $V\subseteq X\setminus V_1$ by \cref{subseteq}.
        \end{subproof}
        $V_1$ is open in $X$ with respect to $T$ by \cref{open_in}.
        Show that $X\setminus V_1$ is closed in $X$ with respect to $T$.
        \begin{subproof}
            $X\setminus V_1\subseteq X$ by \cref{setminus_subseteq}.
            Thus $V_1\in T$ by \cref{open_in}.
            Thus $T\subseteq\pow{X}$ by \cref{topology_on}.
            Thus $V_1\in\pow{X}$ by \cref{subseteq}.
            Thus $V_1\subseteq X$ by \cref{pow_iff}.
            Thus $X\setminus (X\setminus V_1) = V_1$ by \cref{double_relative_complement,subseteq}.
            Thus $X\setminus (X\setminus V_1)$ is open in $X$ with respect to $T$ by \cref{open_in}.
            Thus $X\setminus V_1$ is closed in $X$ with respect to $T$ by \cref{closed_in}.
        \end{subproof}
        $X\setminus V_1\subseteq X$ by \cref{setminus_subseteq}.
        Thus $V\subseteq X$ by \cref{subseteq_transitive}.
        Thus $\closure{V}{X}{T}\subseteq X\setminus V_1$ by \cref{closure_subseteq_closed}.
        $B\subseteq V_1$ by \cref{subseteq}.
        Thus $X\setminus B \supseteq X\setminus V_1$ by \cref{subseteq_implies_setminus_supseteq}.
        Thus $X\setminus V_1\subseteq X\setminus B$ by \cref{subseteq}.
        $U$ is open in $X$ with respect to $T$ by \cref{neighborhood}.
        $U\in T$ by \cref{open_in}.
        Thus $T\subseteq\pow{X}$ by \cref{topology_on}.
        Thus $U\in\pow{X}$ by \cref{subseteq}.
        Thus $U\subseteq X$ by \cref{pow_iff}.
        Thus $X\setminus B = U$ by \cref{double_relative_complement,subseteq}.
        Thus $X\setminus V_1\subseteq U$ by \cref{subseteq_transitive}.
        Thus $\closure{V}{X}{T}\subseteq U$ by \cref{subseteq_transitive}.
    \end{subproof}
\end{proof}

% from §31 Item 4: regularity from neighborhood closure property
\begin{lemma}\label{regular_from_closure_neighborhood}
    Suppose $T$ is a topology on $X$.
    Suppose for all $x$ such that $x\in X$ we have $\{x\}$ is closed in $X$ with respect to $T$.
    Suppose for all $x,U$ such that $x\in X$ and $U$ is a neighborhood of $x$ in $X$ with respect to $T$
        there exists $V$ such that
            $V$ is a neighborhood of $x$ in $X$ with respect to $T$ and
            $\closure{V}{X}{T}\subseteq U$.
    Then $X$ is regular with respect to $T$.
\end{lemma}
\begin{proof}
    Show that for all $x,B$ such that $x\in X$ and $B$ is closed in $X$ with respect to $T$ and $x\notin B$
        there exist $U,V$ such that
            $U$ is open in $X$ with respect to $T$ and
            $V$ is open in $X$ with respect to $T$ and
            $x\in U$ and
            $B\subseteq V$ and
            $U$ is disjoint from $V$.
    \begin{subproof}
        Fix $x$.
        Fix $B$.
        Assume $x\in X$ and $B$ is closed in $X$ with respect to $T$ and $x\notin B$.
        Let $U = X\setminus B$.
        Show that $U$ is a neighborhood of $x$ in $X$ with respect to $T$.
        \begin{subproof}
            $U$ is open in $X$ with respect to $T$ by \cref{closed_in}.
            Thus $x\in U$ by \cref{setminus_intro}.
            Thus $U$ is a neighborhood of $x$ in $X$ with respect to $T$ by \cref{neighborhood}.
        \end{subproof}
        Take $V$ such that
            $V$ is a neighborhood of $x$ in $X$ with respect to $T$ and
            $\closure{V}{X}{T}\subseteq U$.
        Let $W = X\setminus \closure{V}{X}{T}$.
        Show that $W$ is open in $X$ with respect to $T$.
        \begin{subproof}
            $V$ is open in $X$ with respect to $T$ by \cref{neighborhood}.
            Thus $V\in T$ by \cref{open_in}.
            Thus $T\subseteq\pow{X}$ by \cref{topology_on}.
            Thus $V\in\pow{X}$ by \cref{subseteq}.
            Thus $V\subseteq X$ by \cref{pow_iff}.
            $\closure{V}{X}{T}$ is closed in $X$ with respect to $T$ by \cref{closure_closed_in}.
            Thus $W$ is open in $X$ with respect to $T$ by \cref{closed_in}.
        \end{subproof}
        Show that $B\subseteq W$.
        \begin{subproof}
            $B\subseteq X$ by \cref{closed_in}.
            Thus $X\setminus (X\setminus B) = B$ by \cref{double_relative_complement,subseteq}.
            Thus $\closure{V}{X}{T}\subseteq X\setminus B$ by \cref{subseteq}.
            Thus $X\setminus \closure{V}{X}{T}\supseteq X\setminus (X\setminus B)$
                by \cref{subseteq_implies_setminus_supseteq}.
            Thus $X\setminus \closure{V}{X}{T}\supseteq B$.
            Thus $B\subseteq W$ by \cref{subseteq}.
        \end{subproof}
        Show that $V$ is disjoint from $W$.
        \begin{subproof}
            Suppose not.
            Then there exists $y$ such that $y\in V$ and $y\in W$ by \cref{disjoint}.
            $V$ is open in $X$ with respect to $T$ by \cref{neighborhood}.
            Thus $V\in T$ by \cref{open_in}.
            Thus $T\subseteq\pow{X}$ by \cref{topology_on}.
            Thus $V\in\pow{X}$ by \cref{subseteq}.
            Thus $V\subseteq X$ by \cref{pow_iff}.
            Thus $V\subseteq \closure{V}{X}{T}$ by \cref{subseteq_closure}.
            Thus $y\in \closure{V}{X}{T}$ by \cref{subseteq}.
            Thus $y\notin \closure{V}{X}{T}$ by \cref{setminus_elim_right}.
            Contradiction.
        \end{subproof}
        $V$ is open in $X$ with respect to $T$ by \cref{neighborhood}.
        Thus $x\in V$ by \cref{neighborhood}.
        Take $U_1$ such that $U_1 = V$.
        Take $V_1$ such that $V_1 = W$.
    \end{subproof}
    $X$ has closed singletons with respect to $T$ by \cref{one_point_closed}.
    Thus $X$ is regular with respect to $T$ by \cref{regular_space}.
\end{proof}

% from §31 Item 5: normal neighborhood closure property
\begin{lemma}\label{normal_closure_neighborhood}
    Assume $X$ is normal with respect to $T$.
    Then for all $A,U$ such that $A$ is closed in $X$ with respect to $T$ and $U$ is open in $X$ with respect to $T$
        and $A\subseteq U$
        there exists $V$ such that
            $V$ is open in $X$ with respect to $T$ and
            $A\subseteq V$ and
            $\closure{V}{X}{T}\subseteq U$.
\end{lemma}
\begin{proof}
    Fix $A$.
    Fix $U$.
    Assume $A$ is closed in $X$ with respect to $T$ and $U$ is open in $X$ with respect to $T$ and $A\subseteq U$.
    Let $B = X\setminus U$.
    Show that $B$ is closed in $X$ with respect to $T$.
    \begin{subproof}
        $B\subseteq X$ by \cref{setminus_subseteq}.
        $U\in T$ by \cref{open_in}.
        $T$ is a topology on $X$ by \cref{normal_space}.
        Thus $T\subseteq\pow{X}$ by \cref{topology_on}.
        Thus $U\in\pow{X}$ by \cref{subseteq}.
        Thus $U\subseteq X$ by \cref{pow_iff}.
        Thus $X\setminus B = U$ by \cref{double_relative_complement,subseteq}.
        Thus $X\setminus B$ is open in $X$ with respect to $T$ by \cref{open_in}.
        Thus $B$ is closed in $X$ with respect to $T$ by \cref{closed_in}.
    \end{subproof}
    Show that $A$ is disjoint from $B$.
    \begin{subproof}
        Suppose not.
        Then there exists $y$ such that $y\in A$ and $y\in B$ by \cref{disjoint}.
        Thus $y\in U$ by \cref{subseteq}.
        Thus $y\notin U$ by \cref{setminus_elim_right}.
        Contradiction.
    \end{subproof}
    Thus $A$ is closed in $X$ with respect to $T$.
    Thus $B$ is closed in $X$ with respect to $T$.
    Thus $A$ is disjoint from $B$.
    $T$ is a topology on $X$ by \cref{normal_space}.
    Take $V,W$ such that
        $V$ is open in $X$ with respect to $T$ and
        $W$ is open in $X$ with respect to $T$ and
        $A\subseteq V$ and
        $B\subseteq W$ and
        $V$ is disjoint from $W$
        by \cref{normal_space}.

    Show that $\closure{V}{X}{T}\subseteq U$.
    \begin{subproof}
        $V\in T$ by \cref{open_in}.
        $T$ is a topology on $X$ by \cref{normal_space}.
        Thus $T\subseteq\pow{X}$ by \cref{topology_on}.
        Thus $V\in\pow{X}$ by \cref{subseteq}.
        Thus $V\subseteq X$ by \cref{pow_iff}.
        Show that $V\subseteq X\setminus W$.
        \begin{subproof}
            Show that for all $y$ such that $y\in V$ we have $y\in X\setminus W$.
            \begin{subproof}
                Fix $y$.
                Assume $y\in V$.
                Show that $y\notin W$.
                \begin{subproof}
                    Suppose not.
                    Then $y\in W$.
                    Thus there exists $z$ such that $z\in V$ and $z\in W$.
                    Thus $V$ is not disjoint from $W$ by \cref{disjoint}.
                    Contradiction.
                \end{subproof}
                Thus $y\in X$ by \cref{subseteq}.
                Thus $y\in X\setminus W$ by \cref{setminus_intro}.
            \end{subproof}
            Thus $V\subseteq X\setminus W$ by \cref{subseteq}.
        \end{subproof}
        $W\in T$ by \cref{open_in}.
        Show that $X\setminus W$ is closed in $X$ with respect to $T$.
        \begin{subproof}
            $X\setminus W\subseteq X$ by \cref{setminus_subseteq}.
            $T$ is a topology on $X$ by \cref{normal_space}.
            Thus $T\subseteq\pow{X}$ by \cref{topology_on}.
            Thus $W\in\pow{X}$ by \cref{subseteq}.
            Thus $W\subseteq X$ by \cref{pow_iff}.
            Thus $X\setminus (X\setminus W) = W$ by \cref{double_relative_complement,subseteq}.
            Thus $X\setminus (X\setminus W)$ is open in $X$ with respect to $T$ by \cref{open_in}.
            Thus $X\setminus W$ is closed in $X$ with respect to $T$ by \cref{closed_in}.
        \end{subproof}
        Thus $\closure{V}{X}{T}\subseteq X\setminus W$ by \cref{closure_subseteq_closed}.
        $B\subseteq W$ by \cref{subseteq}.
        Thus $X\setminus B \supseteq X\setminus W$ by \cref{subseteq_implies_setminus_supseteq}.
        Thus $X\setminus W\subseteq X\setminus B$ by \cref{subseteq}.
        $U\in T$ by \cref{open_in}.
        $T$ is a topology on $X$ by \cref{normal_space}.
        Thus $T\subseteq\pow{X}$ by \cref{topology_on}.
        Thus $U\in\pow{X}$ by \cref{subseteq}.
        Thus $U\subseteq X$ by \cref{pow_iff}.
        Thus $X\setminus B = U$ by \cref{double_relative_complement,subseteq}.
        Thus $\closure{V}{X}{T}\subseteq U$ by \cref{subseteq_transitive}.
    \end{subproof}
    Thus there exists $V$ such that
        $V$ is open in $X$ with respect to $T$ and
        $A\subseteq V$ and
        $\closure{V}{X}{T}\subseteq U$.
\end{proof}

% from §31 Item 6: normality from neighborhood closure property
\begin{lemma}\label{normal_from_closure_neighborhood}
    Suppose $T$ is a topology on $X$.
    Suppose for all $x$ such that $x\in X$ we have $\{x\}$ is closed in $X$ with respect to $T$.
    Suppose for all $A,U$ such that $A$ is closed in $X$ with respect to $T$ and $U$ is open in $X$ with respect to $T$
        and $A\subseteq U$
        there exists $V$ such that
            $V$ is open in $X$ with respect to $T$ and
            $A\subseteq V$ and
            $\closure{V}{X}{T}\subseteq U$.
    Then $X$ is normal with respect to $T$.
\end{lemma}
\begin{proof}
    Show that for all $A,B$ such that $A$ is closed in $X$ with respect to $T$ and $B$ is closed in $X$ with respect to $T$
        and $A$ is disjoint from $B$
        there exist $U,V$ such that
            $U$ is open in $X$ with respect to $T$ and
            $V$ is open in $X$ with respect to $T$ and
            $A\subseteq U$ and
            $B\subseteq V$ and
            $U$ is disjoint from $V$.
    \begin{subproof}
        Fix $A$.
        Fix $B$.
        Assume $A$ is closed in $X$ with respect to $T$ and $B$ is closed in $X$ with respect to $T$
            and $A$ is disjoint from $B$.
        Let $U = X\setminus B$.
        Show that $A\subseteq U$.
        \begin{subproof}
            Show that for all $y$ such that $y\in A$ we have $y\in U$.
            \begin{subproof}
                Fix $y$.
                Assume $y\in A$.
                Show that $y\notin B$.
                \begin{subproof}
                    Suppose not.
                    Then $y\in B$.
                    Thus there exists $z$ such that $z\in A$ and $z\in B$.
                    Thus $A$ is not disjoint from $B$ by \cref{disjoint}.
                    Contradiction.
                \end{subproof}
                $A\subseteq X$ by \cref{closed_in}.
                Thus $y\in X$ by \cref{subseteq}.
                Thus $y\in X\setminus B$ by \cref{setminus_intro}.
                Thus $y\in U$.
            \end{subproof}
            Thus $A\subseteq U$ by \cref{subseteq}.
        \end{subproof}
        $U$ is open in $X$ with respect to $T$ by \cref{closed_in}.
        Take $V$ such that
            $V$ is open in $X$ with respect to $T$ and
            $A\subseteq V$ and
            $\closure{V}{X}{T}\subseteq U$.
        Let $W = X\setminus \closure{V}{X}{T}$.
        Show that $W$ is open in $X$ with respect to $T$.
        \begin{subproof}
            $V\in T$ by \cref{open_in}.
            Thus $T\subseteq\pow{X}$ by \cref{topology_on}.
            Thus $V\in\pow{X}$ by \cref{subseteq}.
            Thus $V\subseteq X$ by \cref{pow_iff}.
            $\closure{V}{X}{T}$ is closed in $X$ with respect to $T$ by \cref{closure_closed_in}.
            Thus $W$ is open in $X$ with respect to $T$ by \cref{closed_in}.
        \end{subproof}
        Show that $B\subseteq W$.
        \begin{subproof}
            $B\subseteq X$ by \cref{closed_in}.
            Thus $X\setminus (X\setminus B) = B$ by \cref{double_relative_complement,subseteq}.
            Thus $\closure{V}{X}{T}\subseteq X\setminus B$ by \cref{subseteq}.
            Thus $X\setminus \closure{V}{X}{T}\supseteq X\setminus (X\setminus B)$
                by \cref{subseteq_implies_setminus_supseteq}.
            Thus $X\setminus \closure{V}{X}{T}\supseteq B$.
            Thus $B\subseteq W$ by \cref{subseteq}.
        \end{subproof}
        Show that $V$ is disjoint from $W$.
        \begin{subproof}
            Suppose not.
            Then there exists $y$ such that $y\in V$ and $y\in W$ by \cref{disjoint}.
            $V\in T$ by \cref{open_in}.
            Thus $T\subseteq\pow{X}$ by \cref{topology_on}.
            Thus $V\in\pow{X}$ by \cref{subseteq}.
            Thus $V\subseteq X$ by \cref{pow_iff}.
            Thus $V\subseteq \closure{V}{X}{T}$ by \cref{subseteq_closure}.
            Thus $y\in \closure{V}{X}{T}$ by \cref{subseteq}.
            Thus $y\notin \closure{V}{X}{T}$ by \cref{setminus_elim_right}.
            Contradiction.
        \end{subproof}
        Thus there exist $U_1,V_1$ such that
            $U_1$ is open in $X$ with respect to $T$ and
            $V_1$ is open in $X$ with respect to $T$ and
            $A\subseteq U_1$ and
            $B\subseteq V_1$ and
            $U_1$ is disjoint from $V_1$.
        Take $U_1$ such that $U_1 = V$.
        Take $V_1$ such that $V_1 = W$.
    \end{subproof}
    $X$ has closed singletons with respect to $T$ by \cref{one_point_closed}.
    Thus $X$ is normal with respect to $T$ by \cref{normal_space}.
\end{proof}

% from §31 Item 7: subspace of a Hausdorff space is Hausdorff
\begin{lemma}\label{subspace_hausdorff_31}
    Assume $X$ is Hausdorff with respect to $T$.
    Assume $Y$ is a subspace of $X$ with respect to $T$.
    Let $T_Y = \subspaceTopology{T}{Y}$.
    Then $Y$ is Hausdorff with respect to $T_Y$.
\end{lemma}
\begin{proof}
    Thus $Y\subseteq X$ by \cref{subspace_of}.
    Thus $Y$ is Hausdorff with respect to $T_Y$ by \cref{subspace_hausdorff}.
\end{proof}

% from §31 Item 8: product of Hausdorff spaces is Hausdorff
\begin{lemma}\label{product_hausdorff_31}
    Assume $X$ is a function.
    Assume $T$ is a function.
    Assume $\dom{T} = \dom{X}$.
    Assume $\dom{X}$ is inhabited.
    Suppose for all $\alpha\in\dom{X}$ we have $\apply{X}{\alpha}$ is Hausdorff with respect to $\apply{T}{\alpha}$.
    Then $\productFamily{X}$ is Hausdorff with respect to $\productTopologyIndexed{T}{X}$.
\end{lemma}
\begin{proof}
    Thus $\productFamily{X}$ is Hausdorff with respect to $\productTopologyIndexed{T}{X}$ by \cref{product_hausdorff_product}.
\end{proof}

% from §31 Item 9: subspace of a regular space is regular
\begin{lemma}\label{subspace_regular}
    Assume $X$ is regular with respect to $T$.
    Assume $Y$ is a subspace of $X$ with respect to $T$.
    Let $T_Y = \subspaceTopology{T}{Y}$.
    Then $Y$ is regular with respect to $T_Y$.
\end{lemma}
\begin{proof}
    $T$ is a topology on $X$ by \cref{regular_space}.
    Thus $Y\subseteq X$ by \cref{subspace_of}.
    Thus $T_Y$ is a topology on $Y$ by \cref{subspace_topology_is_topology}.
    Show that for all $y$ such that $y\in Y$ we have $\{y\}$ is closed in $Y$ with respect to $T_Y$.
    \begin{subproof}
        Fix $y$.
        Assume $y\in Y$.
        Thus $y\in X$ by \cref{elem_subseteq}.
        $X$ has closed singletons with respect to $T$ by \cref{regular_space}.
        Thus $\{y\}$ is closed in $X$ with respect to $T$ by \cref{one_point_closed}.
        $\{y\}\subseteq Y$ by \cref{singleton_subset_intro}.
        Thus $\{y\}\inter Y = \{y\}$ by \cref{inter_absorb_supseteq_right}.
        Thus $\{y\}$ is closed in $Y$ with respect to $T_Y$ by \cref{closed_in_subspace_iff,subspace_of}.
    \end{subproof}
    Thus $Y$ has closed singletons with respect to $T_Y$ by \cref{one_point_closed}.
    Show that for all $x,B$ such that $x\in Y$ and $B$ is closed in $Y$ with respect to $T_Y$ and $x\notin B$
        there exist $U,V$ such that
            $U$ is open in $Y$ with respect to $T_Y$ and
            $V$ is open in $Y$ with respect to $T_Y$ and
            $x\in U$ and
            $B\subseteq V$ and
            $U$ is disjoint from $V$.
    \begin{subproof}
        Fix $x$.
        Fix $B$.
        Assume $x\in Y$ and $B$ is closed in $Y$ with respect to $T_Y$ and $x\notin B$.
        Take $C$ such that $C$ is closed in $X$ with respect to $T$ and $B = C\inter Y$
            by \cref{closed_in_subspace_iff,subspace_of}.
        Thus $x\in X$ by \cref{elem_subseteq}.
        Show that $x\notin C$.
        \begin{subproof}
            Suppose not.
            Thus $x\in C$.
            Thus $x\in C$ and $x\in Y$.
            Thus $x\in C\inter Y$ by \cref{inter_intro}.
            Thus $C\inter Y\subseteq B$ by \cref{subseteq_refl}.
            Thus $x\in B$ by \cref{subseteq}.
            Contradiction.
        \end{subproof}
        Take $U,V$ such that
            $U$ is open in $X$ with respect to $T$ and
            $V$ is open in $X$ with respect to $T$ and
            $x\in U$ and
            $C\subseteq V$ and
            $U$ is disjoint from $V$
            by \cref{regular_space}.
        Let $U_1 = U\inter Y$.
        Let $V_1 = V\inter Y$.
        $U\in T$ by \cref{open_in}.
        $V\in T$ by \cref{open_in}.
        Thus $U_1 = Y\inter U$ by \cref{inter_comm}.
        Thus $V_1 = Y\inter V$ by \cref{inter_comm}.
        Thus $U_1\in T_Y$ by \cref{subspace_topology}.
        Thus $V_1\in T_Y$ by \cref{subspace_topology}.
        Thus $U_1$ is open in $Y$ with respect to $T_Y$ by \cref{open_in}.
        Thus $V_1$ is open in $Y$ with respect to $T_Y$ by \cref{open_in}.
        Thus $x\in U_1$ by \cref{inter_intro}.
        Show that $B\subseteq V_1$.
        \begin{subproof}
            Thus $B\subseteq C$ by \cref{inter_lower_left}.
            Thus $B\subseteq V$ by \cref{subseteq_transitive}.
            Thus $B\subseteq Y$ by \cref{inter_lower_right}.
            Thus $B\subseteq V\inter Y$ by \cref{subseteq_inter_iff}.
        \end{subproof}
        Show that $U_1$ is disjoint from $V_1$.
        \begin{subproof}
            Suppose not.
            Then there exists $y$ such that $y\in U_1$ and $y\in V_1$.
            Thus $y\in U$ and $y\in V$ by \cref{inter}.
            Thus there exists $a$ such that $a\in U$ and $a\in V$.
            Contradiction by \cref{disjoint}.
        \end{subproof}
    \end{subproof}
    Thus $Y$ is regular with respect to $T_Y$ by \cref{regular_space}.
\end{proof}

% from §31 Item 10: product of regular spaces is regular
\begin{lemma}\label{product_regular}
    Assume $X$ is a function.
    Assume $T$ is a function.
    Assume $\dom{T} = \dom{X}$.
    Assume $\dom{X}$ is inhabited.
    Suppose for all $\alpha\in\dom{X}$ we have $\apply{X}{\alpha}$ is regular with respect to $\apply{T}{\alpha}$.
    Then $\productFamily{X}$ is regular with respect to $\productTopologyIndexed{T}{X}$.
\end{lemma}
\begin{proof}
    Let $T_0 = \productTopologyIndexed{T}{X}$.
    Let $S = \productSubbasis{T}{X}$.
    Show that for all $\alpha\in\dom{X}$ we have $\apply{T}{\alpha}$ is a topology on $\apply{X}{\alpha}$.
    \begin{subproof}
        Fix $\alpha$.
        Assume $\alpha\in\dom{X}$.
        $\apply{X}{\alpha}$ is regular with respect to $\apply{T}{\alpha}$.
        Thus $\apply{T}{\alpha}$ is a topology on $\apply{X}{\alpha}$ by \cref{regular_space}.
    \end{subproof}
    Show that for all $\alpha\in\dom{X}$ we have $\apply{X}{\alpha}$ is Hausdorff with respect to $\apply{T}{\alpha}$.
    \begin{subproof}
        Fix $\alpha$.
        Assume $\alpha\in\dom{X}$.
        $\apply{X}{\alpha}$ is regular with respect to $\apply{T}{\alpha}$.
        Thus $\apply{X}{\alpha}$ is Hausdorff with respect to $\apply{T}{\alpha}$ by \cref{regular_implies_hausdorff}.
    \end{subproof}
    Thus $\productFamily{X}$ is Hausdorff with respect to $T_0$ by \cref{product_hausdorff_31}.
    $T_0$ is a topology on $\productFamily{X}$ by \cref{hausdorff}.
    Show that for all $x$ such that $x\in\productFamily{X}$ we have $\{x\}$ is closed in $\productFamily{X}$ with respect to $T_0$.
    \begin{subproof}
        Fix $x$.
        Assume $x\in\productFamily{X}$.
        Let $A = \{x\}$.
        $A\subseteq\productFamily{X}$ by \cref{singleton_subset_intro}.
        Let $G = \{x,x\}$.
        $G$ is finite by \cref{pair_is_finite}.
        $x\in G$ by \cref{upair_intro_left}.
        $A\subseteq G$ by \cref{singleton_subset_intro}.
        Thus $A$ is finite by \cref{subseteq_finite_is_finite}.
        Thus $A$ is closed in $\productFamily{X}$ with respect to $T_0$ by \cref{finite_point_set_closed_hausdorff}.
    \end{subproof}
    Show that for all $x,U$ such that $x\in\productFamily{X}$ and
        $U$ is a neighborhood of $x$ in $\productFamily{X}$ with respect to $T_0$
        there exists $V$ such that
            $V$ is a neighborhood of $x$ in $\productFamily{X}$ with respect to $T_0$ and
            $\closure{V}{\productFamily{X}}{T_0}\subseteq U$.
    \begin{subproof}
        Fix $x$.
        Fix $U$.
        Assume $x\in\productFamily{X}$ and
            $U$ is a neighborhood of $x$ in $\productFamily{X}$ with respect to $T_0$.
        $U$ is open in $\productFamily{X}$ with respect to $T_0$ by \cref{neighborhood}.
        Thus $x\in U$ by \cref{neighborhood}.
        Thus $U\in T_0$ by \cref{open_in}.
        Thus $T_0\subseteq\pow{\productFamily{X}}$ by \cref{topology_on}.
        Thus $U\in\pow{\productFamily{X}}$ by \cref{subseteq}.
        $\finiteIntersections{S}$ is a basis for $T_0$ on $\productFamily{X}$ by \cref{product_subbasis_finite_intersections_basis}.
        Thus $T_0 = \basisTopology{\finiteIntersections{S}}{\productFamily{X}}$
            by \cref{product_topology_indexed,topology_generated_by_subbasis,basis_for_topology}.
        Thus $U\in\basisTopology{\finiteIntersections{S}}{\productFamily{X}}$ by \cref{basis_for_topology}.
        Take $B_1\in\finiteIntersections{S}$ such that $x\in B_1$ and $B_1\subseteq U$ by \cref{topology_generated_by_basis}.
        Take $F\subseteq S$ such that $F$ is finite and inhabited and $B_1 = \inters{F}$ by \cref{finite_intersections}.
        Let $R = \{w\in F\times\pow{\productFamily{X}} \mid \exists \beta\in\dom{X}. \exists U_0\in\apply{T}{\beta}. \exists V_0\in\apply{T}{\beta}.
            \fst{w} = \preimg{\piIndex{X}{\beta}}{U_0} \land
            \apply{x}{\beta}\in V_0 \land
            \closure{V_0}{\apply{X}{\beta}}{\apply{T}{\beta}}\subseteq U_0 \land
            \snd{w} = \preimg{\piIndex{X}{\beta}}{V_0}\}$.
        Show that $R$ is a relation.
        \begin{subproof}
            Show that for all $w$ such that $w\in R$ there exist $y,z$ such that $w = (y,z)$.
            \begin{subproof}
                Fix $w$.
                Assume $w\in R$.
                Thus $w\in F\times\pow{\productFamily{X}}$.
                Take $y,z$ such that $y\in F$ and $z\in\pow{\productFamily{X}}$ and $w = (y,z)$ by \cref{times_elem_is_tuple}.
            \end{subproof}
            Thus $R$ is a relation by \cref{relation}.
        \end{subproof}
        Show that for all $A\in F$ there exists $B$ such that $A\mathrel{R} B$.
        \begin{subproof}
            Fix $A$.
            Assume $A\in F$.
            Thus $A\in S$ by \cref{subseteq}.
            Take $\beta\in\dom{X}$ such that $A\in\productSubbasisAt{T}{X}{\beta}$ by \cref{product_subbasis_union_indexed}.
            Take $U_0\in\apply{T}{\beta}$ such that $A = \preimg{\piIndex{X}{\beta}}{U_0}$ by \cref{product_subbasis_indexed}.
            Thus $\inters{F}\subseteq A$ by \cref{inters_subseteq_elem}.
            Thus $x\in A$ by \cref{subseteq}.
            Take $y\in U_0$ such that $x\mathrel{\piIndex{X}{\beta}} y$ by \cref{preimg_iff}.
            Take $x_0\in\productFamily{X}$ such that $(x,y) = (x_0,\apply{x_0}{\beta})$ by \cref{projection_map_indexed}.
            Thus $x = x_0$ and $y = \apply{x}{\beta}$ by \cref{pair_eq_iff}.
            Thus $\apply{x}{\beta}\in U_0$.
            $U_0$ is open in $\apply{X}{\beta}$ with respect to $\apply{T}{\beta}$ by \cref{open_in}.
            Thus $U_0$ is a neighborhood of $\apply{x}{\beta}$ in $\apply{X}{\beta}$ with respect to $\apply{T}{\beta}$ by \cref{neighborhood}.
            Thus $\apply{x}{\beta}\in\apply{X}{\beta}$ by \cref{indexed_product}.
            $\apply{X}{\beta}$ is regular with respect to $\apply{T}{\beta}$.
            Take $V_0$ such that
                $V_0$ is a neighborhood of $\apply{x}{\beta}$ in $\apply{X}{\beta}$ with respect to $\apply{T}{\beta}$ and
                $\closure{V_0}{\apply{X}{\beta}}{\apply{T}{\beta}}\subseteq U_0$
                by \cref{regular_closure_neighborhood}.
            $V_0$ is open in $\apply{X}{\beta}$ with respect to $\apply{T}{\beta}$ by \cref{neighborhood}.
            Thus $V_0\in\apply{T}{\beta}$ by \cref{open_in}.
            Thus $\apply{x}{\beta}\in V_0$ by \cref{neighborhood}.
            Let $B = \preimg{\piIndex{X}{\beta}}{V_0}$.
            Let $w = (A,B)$.
            Show that $B\in\pow{\productFamily{X}}$.
            \begin{subproof}
                Show that $B\subseteq\productFamily{X}$.
                \begin{subproof}
                    $\piIndex{X}{\beta}\in\funs{\productFamily{X}}{\apply{X}{\beta}}$ by \cref{projection_map_indexed_function}.
                    Thus $\piIndex{X}{\beta}$ is a relation by \cref{funs_is_relation}.
                    Thus $B\subseteq\dom{\piIndex{X}{\beta}}$ by \cref{preimg_subseteq_dom}.
                    $\dom{\piIndex{X}{\beta}} = \productFamily{X}$ by \cref{funs}.
                    Thus $B\subseteq\productFamily{X}$ by \cref{subseteq_transitive}.
                \end{subproof}
                Thus $B\in\pow{\productFamily{X}}$ by \cref{powerset_intro}.
            \end{subproof}
            Thus $w\in F\times\pow{\productFamily{X}}$ by \cref{times_tuple_intro}.
            $\fst{w} = A$ and $\snd{w} = B$ by \cref{fst_eq,snd_eq}.
            Thus $\fst{w} = \preimg{\piIndex{X}{\beta}}{U_0}$.
            Thus $\snd{w} = \preimg{\piIndex{X}{\beta}}{V_0}$.
            Thus $w\in R$.
            Thus $(A,B)\in R$ by \cref{pair_eq_iff}.
            Thus $A\mathrel{R}B$.
        \end{subproof}
        Take $g$ such that $g$ is a function on $F$ and for all $A\in F$ we have $A\mathrel{R}\apply{g}{A}$
            by \cref{choice_function}.
        Let $F_1 = \img{g}{F}$.
        Thus $F_1$ is finite by \cref{finite_image_function}.
        Show that $F_1$ is inhabited.
        \begin{subproof}
            Take $A$ such that $A\in F$.
            Thus $(A,\apply{g}{A})\in g$ by \cref{function_apply_elim}.
            Thus $\apply{g}{A}\in F_1$ by \cref{img_iff}.
            Thus $F_1\neq\emptyset$ by \cref{emptyset}.
            Thus $F_1$ is inhabited by \cref{nonempty_inhabited}.
        \end{subproof}
        Show that $F_1\subseteq T_0$.
        \begin{subproof}
            Show that for all $B$ such that $B\in F_1$ we have $B\in T_0$.
            \begin{subproof}
                Fix $B$.
                Assume $B\in F_1$.
                Take $A$ such that $A\in F$ and $(A,B)\in g$ by \cref{img_iff}.
                Thus $\apply{g}{A} = B$ by \cref{function_apply_iff}.
                Thus $A\mathrel{R}B$ by \cref{choice_function}.
                Take $\beta\in\dom{X}$ such that there exists $U_0\in\apply{T}{\beta}$ such that there exists $V_0\in\apply{T}{\beta}$ such that
                    $A = \preimg{\piIndex{X}{\beta}}{U_0}$ and
                    $\apply{x}{\beta}\in V_0$ and
                    $\closure{V_0}{\apply{X}{\beta}}{\apply{T}{\beta}}\subseteq U_0$ and
                    $B = \preimg{\piIndex{X}{\beta}}{V_0}$.
                Take $U_0\in\apply{T}{\beta}$ such that there exists $V_0\in\apply{T}{\beta}$ such that
                    $A = \preimg{\piIndex{X}{\beta}}{U_0}$ and
                    $\apply{x}{\beta}\in V_0$ and
                    $\closure{V_0}{\apply{X}{\beta}}{\apply{T}{\beta}}\subseteq U_0$ and
                    $B = \preimg{\piIndex{X}{\beta}}{V_0}$.
                Take $V_0\in\apply{T}{\beta}$ such that
                    $A = \preimg{\piIndex{X}{\beta}}{U_0}$ and
                    $\apply{x}{\beta}\in V_0$ and
                    $\closure{V_0}{\apply{X}{\beta}}{\apply{T}{\beta}}\subseteq U_0$ and
                    $B = \preimg{\piIndex{X}{\beta}}{V_0}$.
                $\apply{X}{\beta}$ is regular with respect to $\apply{T}{\beta}$.
                Thus $\apply{T}{\beta}$ is a topology on $\apply{X}{\beta}$ by \cref{regular_space}.
                $V_0$ is open in $\apply{X}{\beta}$ with respect to $\apply{T}{\beta}$ by \cref{open_in}.
                $\piIndex{X}{\beta}$ is continuous from $\productFamily{X}$ to $\apply{X}{\beta}$
                    with respect to $T_0$ and $\apply{T}{\beta}$ by \cref{projection_map_indexed_continuous}.
                Thus $B$ is open in $\productFamily{X}$ with respect to $T_0$ by \cref{continuous,open_in}.
                Thus $B\in T_0$ by \cref{open_in}.
            \end{subproof}
            Thus $F_1\subseteq T_0$ by \cref{subseteq}.
        \end{subproof}
        Let $V = \inters{F_1}$.
        Thus $V\in T_0$ by \cref{finite_intersections_open_in_topology}.
        Thus $V$ is open in $\productFamily{X}$ with respect to $T_0$ by \cref{open_in}.
        Show that $x\in V$.
        \begin{subproof}
            Show that for all $B$ such that $B\in F_1$ we have $x\in B$.
            \begin{subproof}
                Fix $B$.
                Assume $B\in F_1$.
                Take $A$ such that $A\in F$ and $(A,B)\in g$ by \cref{img_iff}.
                Thus $\apply{g}{A} = B$ by \cref{function_apply_iff}.
                Thus $A\mathrel{R}B$ by \cref{choice_function}.
                Take $\beta\in\dom{X}$ such that there exists $U_0$ such that there exists $V_0$ such that
                    $U_0\in\apply{T}{\beta}$ and
                    $V_0\in\apply{T}{\beta}$ and
                    $A = \preimg{\piIndex{X}{\beta}}{U_0}$ and
                    $\apply{x}{\beta}\in V_0$ and
                    $\closure{V_0}{\apply{X}{\beta}}{\apply{T}{\beta}}\subseteq U_0$ and
                    $B = \preimg{\piIndex{X}{\beta}}{V_0}$.
                Thus $(x,\apply{x}{\beta})\in\piIndex{X}{\beta}$ by \cref{projection_map_indexed}.
                Thus $x\mathrel{\piIndex{X}{\beta}}\apply{x}{\beta}$.
                Thus $x\in B$ by \cref{preimg_iff}.
            \end{subproof}
            Thus $x\in\inters{F_1}$ by \cref{inters_iff_forall}.
        \end{subproof}
        Thus $V$ is a neighborhood of $x$ in $\productFamily{X}$ with respect to $T_0$ by \cref{neighborhood}.
        Show that $\closure{V}{\productFamily{X}}{T_0}\subseteq U$.
        \begin{subproof}
            Show that for all $A\in F$ we have $\closure{V}{\productFamily{X}}{T_0}\subseteq A$.
            \begin{subproof}
                Fix $A$.
                Assume $A\in F$.
                Thus $A\mathrel{R}\apply{g}{A}$ by \cref{choice_function}.
                Let $B = \apply{g}{A}$.
                Take $\beta\in\dom{X}$ such that there exists $U_0\in\apply{T}{\beta}$ such that there exists $V_0\in\apply{T}{\beta}$ such that
                    $A = \preimg{\piIndex{X}{\beta}}{U_0}$ and
                    $\apply{x}{\beta}\in V_0$ and
                    $\closure{V_0}{\apply{X}{\beta}}{\apply{T}{\beta}}\subseteq U_0$ and
                    $B = \preimg{\piIndex{X}{\beta}}{V_0}$.
                Take $U_0\in\apply{T}{\beta}$ such that there exists $V_0\in\apply{T}{\beta}$ such that
                    $A = \preimg{\piIndex{X}{\beta}}{U_0}$ and
                    $\apply{x}{\beta}\in V_0$ and
                    $\closure{V_0}{\apply{X}{\beta}}{\apply{T}{\beta}}\subseteq U_0$ and
                    $B = \preimg{\piIndex{X}{\beta}}{V_0}$.
                Take $V_0\in\apply{T}{\beta}$ such that
                    $A = \preimg{\piIndex{X}{\beta}}{U_0}$ and
                    $\apply{x}{\beta}\in V_0$ and
                    $\closure{V_0}{\apply{X}{\beta}}{\apply{T}{\beta}}\subseteq U_0$ and
                    $B = \preimg{\piIndex{X}{\beta}}{V_0}$.
                Thus $\apply{g}{A}\in F_1$ by \cref{img_iff,function_apply_elim}.
                Thus $V\subseteq B$ by \cref{inters_subseteq_elem}.
                $\apply{X}{\beta}$ is regular with respect to $\apply{T}{\beta}$.
                Thus $\apply{T}{\beta}$ is a topology on $\apply{X}{\beta}$ by \cref{regular_space}.
                Thus $\apply{T}{\beta}\subseteq\pow{\apply{X}{\beta}}$ by \cref{topology_on}.
                Thus $V_0\in\pow{\apply{X}{\beta}}$ by \cref{subseteq}.
                Thus $V_0\subseteq\apply{X}{\beta}$ by \cref{powerset_elim}.
                $\closure{V_0}{\apply{X}{\beta}}{\apply{T}{\beta}}$ is closed in $\apply{X}{\beta}$ with respect to $\apply{T}{\beta}$ by \cref{closure_closed_in}.
                Let $D = \apply{X}{\beta}\setminus \closure{V_0}{\apply{X}{\beta}}{\apply{T}{\beta}}$.
                $D$ is open in $\apply{X}{\beta}$ with respect to $\apply{T}{\beta}$ by \cref{closed_in}.
                $\piIndex{X}{\beta}$ is continuous from $\productFamily{X}$ to $\apply{X}{\beta}$
                    with respect to $T_0$ and $\apply{T}{\beta}$ by \cref{projection_map_indexed_continuous}.
                Thus $\preimg{\piIndex{X}{\beta}}{D}$ is open in $\productFamily{X}$ with respect to $T_0$ by \cref{continuous}.
                $\piIndex{X}{\beta}\in\funs{\productFamily{X}}{\apply{X}{\beta}}$ by \cref{projection_map_indexed_function}.
                $\apply{X}{\beta}$ is closed in $\apply{X}{\beta}$ with respect to $\apply{T}{\beta}$ by \cref{closed_whole_space}.
                Thus $\closure{V_0}{\apply{X}{\beta}}{\apply{T}{\beta}}\subseteq \apply{X}{\beta}$ by \cref{closure_subseteq_closed}.
                Thus $\preimg{\piIndex{X}{\beta}}{D} = \productFamily{X}\setminus \preimg{\piIndex{X}{\beta}}{\closure{V_0}{\apply{X}{\beta}}{\apply{T}{\beta}}}$
                    by \cref{preimg_setminus_function}.
                Thus $\productFamily{X}\setminus \preimg{\piIndex{X}{\beta}}{\closure{V_0}{\apply{X}{\beta}}{\apply{T}{\beta}}}$
                    is open in $\productFamily{X}$ with respect to $T_0$ by \cref{open_in}.
                Thus $\piIndex{X}{\beta}$ is a relation by \cref{funs_is_relation}.
                Thus $\preimg{\piIndex{X}{\beta}}{\closure{V_0}{\apply{X}{\beta}}{\apply{T}{\beta}}}\subseteq \dom{\piIndex{X}{\beta}}$
                    by \cref{preimg_subseteq_dom}.
                $\dom{\piIndex{X}{\beta}} = \productFamily{X}$ by \cref{funs}.
                Thus $\preimg{\piIndex{X}{\beta}}{\closure{V_0}{\apply{X}{\beta}}{\apply{T}{\beta}}}\subseteq \productFamily{X}$
                    by \cref{subseteq_transitive}.
                Thus $\preimg{\piIndex{X}{\beta}}{\closure{V_0}{\apply{X}{\beta}}{\apply{T}{\beta}}}$
                    is closed in $\productFamily{X}$ with respect to $T_0$ by \cref{closed_in}.
                $V_0\subseteq \closure{V_0}{\apply{X}{\beta}}{\apply{T}{\beta}}$ by \cref{subseteq_closure}.
                Thus $B\subseteq \preimg{\piIndex{X}{\beta}}{\closure{V_0}{\apply{X}{\beta}}{\apply{T}{\beta}}}$
                    by \cref{preimg_subseteq}.
                Thus $V\subseteq \preimg{\piIndex{X}{\beta}}{\closure{V_0}{\apply{X}{\beta}}{\apply{T}{\beta}}}$
                    by \cref{subseteq_transitive}.
                $T_0$ is a topology on $\productFamily{X}$ by \cref{hausdorff}.
                Thus $T_0\subseteq\pow{\productFamily{X}}$ by \cref{topology_on}.
                Thus $V\in\pow{\productFamily{X}}$ by \cref{subseteq}.
                Thus $V\subseteq\productFamily{X}$ by \cref{powerset_elim}.
                Thus $\closure{V}{\productFamily{X}}{T_0}\subseteq
                    \preimg{\piIndex{X}{\beta}}{\closure{V_0}{\apply{X}{\beta}}{\apply{T}{\beta}}}$
                    by \cref{closure_subseteq_closed}.
                Thus $\closure{V}{\productFamily{X}}{T_0}\subseteq \preimg{\piIndex{X}{\beta}}{U_0}$
                    by \cref{preimg_subseteq,subseteq_transitive}.
                Thus $\closure{V}{\productFamily{X}}{T_0}\subseteq A$ by \cref{subseteq}.
            \end{subproof}
            Thus $\closure{V}{\productFamily{X}}{T_0}\subseteq \inters{F}$ by \cref{subseteq_inters_iff}.
            Thus $\closure{V}{\productFamily{X}}{T_0}\subseteq B_1$ by \cref{subseteq}.
            Thus $\closure{V}{\productFamily{X}}{T_0}\subseteq U$ by \cref{subseteq_transitive}.
        \end{subproof}
    \end{subproof}
    Thus $\productFamily{X}$ is regular with respect to $T_0$ by \cref{regular_from_closure_neighborhood}.
\end{proof}

% from §32 Item 1: regular space with countable basis is normal
\begin{theorem}\label{regular_second_countable_normal}
    Assume $X$ is regular with respect to $T$.
    Suppose $X$ satisfies the second countability axiom with respect to $T$.
    Then $X$ is normal with respect to $T$.
\end{theorem}
\begin{proof}
    $T$ is a topology on $X$ by \cref{regular_space}.
    Take $b$ such that
        $b$ is a sequence in $\pow{X}$ and
        for all $n\in\naturalsPlus$ we have $b(n)$ is open in $X$ with respect to $T$ and
        for all $x\in X$ we have
            for all $U$ such that $U$ is open in $X$ with respect to $T$ and $x\in U$
                there exists $n\in\naturalsPlus$ such that $x\in b(n)$ and $b(n)\subseteq U$
        by \cref{second_countability_axiom}.
    Show that for all $A,B$ such that $A$ is closed in $X$ with respect to $T$ and $B$ is closed in $X$ with respect to $T$
        and $A$ is disjoint from $B$
        there exist $U,V$ such that
            $U$ is open in $X$ with respect to $T$ and
            $V$ is open in $X$ with respect to $T$ and
            $A\subseteq U$ and
            $B\subseteq V$ and
            $U$ is disjoint from $V$.
    \begin{subproof}
        Fix $A$.
        Fix $B$.
        Assume $A$ is closed in $X$ with respect to $T$ and $B$ is closed in $X$ with respect to $T$
            and $A$ is disjoint from $B$.
        Let $P = \{n\in\naturalsPlus \mid b(n)\inter A \neq \emptyset \land \closure{b(n)}{X}{T}\inter B = \emptyset\}$.
        Let $Q = \{n\in\naturalsPlus \mid b(n)\inter B \neq \emptyset \land \closure{b(n)}{X}{T}\inter A = \emptyset\}$.
        Let $U_1 = \{w \mid \exists n\in P. w = (n,b(n))\}$.
        Let $U_2 = \{w \mid \exists n\in\naturalsPlus\setminus P. w = (n,\emptyset)\}$.
        Let $U = U_1 \union U_2$.
        Let $V_1 = \{w \mid \exists n\in Q. w = (n,b(n))\}$.
        Let $V_2 = \{w \mid \exists n\in\naturalsPlus\setminus Q. w = (n,\emptyset)\}$.
        Let $V = V_1 \union V_2$.

        Show that $U$ is a function on $\naturalsPlus$.
        \begin{subproof}
            Show that $U$ is a function.
            \begin{subproof}
                Show that for all $w$ such that $w\in U$ there exist $x,y$ such that $w = (x,y)$.
                \begin{subproof}
                    Fix $w$.
                    Assume $w\in U$.
                    Then $w\in U_1$ or
                        $w\in U_2$
                        by \cref{union_iff}.
                    \begin{byCase}
                    \caseOf{$w\in U_1$.}
                        Take $n\in P$ such that $w = (n,b(n))$.
                        Thus there exist $x,y$ such that $w = (x,y)$.
                    \caseOf{$w\in U_2$.}
                        Take $n\in\naturalsPlus\setminus P$ such that $w = (n,\emptyset)$.
                        Thus there exist $x,y$ such that $w = (x,y)$.
                    \end{byCase}
                \end{subproof}
                Thus $U$ is a relation by \cref{relation}.
                Show that for all $n,y_1,y_2$ such that $(n,y_1)\in U$ and $(n,y_2)\in U$ we have $y_1 = y_2$.
                \begin{subproof}
                    Fix $n$.
                    Fix $y_1$.
                    Fix $y_2$.
                    Assume $(n,y_1)\in U$ and $(n,y_2)\in U$.
                    \begin{byCase}
                    \caseOf{$n\in P$.}
                        Show that $(n,y_1)\in U_1$.
                        \begin{subproof}
                            Suppose not.
                            Then $(n,y_1)\in U_2$ by \cref{union_iff}.
                            Take $n_0\in\naturalsPlus\setminus P$ such that $(n,y_1) = (n_0,\emptyset)$.
                            Thus $n = n_0$ by \cref{pair_eq_iff}.
                            Thus $n\notin P$ by \cref{setminus_elim_right}.
                            Contradiction.
                        \end{subproof}
                        Take $n_1\in P$ such that $(n,y_1) = (n_1,b(n_1))$.
                        Thus $y_1 = b(n)$ by \cref{pair_eq_iff}.
                        Show that $(n,y_2)\in U_1$.
                        \begin{subproof}
                            Suppose not.
                            Then $(n,y_2)\in U_2$ by \cref{union_iff}.
                            Take $n_0\in\naturalsPlus\setminus P$ such that $(n,y_2) = (n_0,\emptyset)$.
                            Thus $n = n_0$ by \cref{pair_eq_iff}.
                            Thus $n\notin P$ by \cref{setminus_elim_right}.
                            Contradiction.
                        \end{subproof}
                        Take $n_2\in P$ such that $(n,y_2) = (n_2,b(n_2))$.
                        Thus $y_2 = b(n)$ by \cref{pair_eq_iff}.
                        Thus $y_1 = y_2$.
                    \caseOf{$n\notin P$.}
                        Show that $(n,y_1)\in U_2$.
                        \begin{subproof}
                            Suppose not.
                            Then $(n,y_1)\in U_1$ by \cref{union_iff}.
                            Take $n_0\in P$ such that $(n,y_1) = (n_0,b(n_0))$.
                            Thus $n = n_0$ by \cref{pair_eq_iff}.
                            Thus $n\in P$.
                            Contradiction.
                        \end{subproof}
                        Take $n_1\in\naturalsPlus\setminus P$ such that $(n,y_1) = (n_1,\emptyset)$.
                        Thus $y_1 = \emptyset$ by \cref{pair_eq_iff}.
                        Show that $(n,y_2)\in U_2$.
                        \begin{subproof}
                            Suppose not.
                            Then $(n,y_2)\in U_1$ by \cref{union_iff}.
                            Take $n_0\in P$ such that $(n,y_2) = (n_0,b(n_0))$.
                            Thus $n = n_0$ by \cref{pair_eq_iff}.
                            Thus $n\in P$.
                            Contradiction.
                        \end{subproof}
                        Take $n_2\in\naturalsPlus\setminus P$ such that $(n,y_2) = (n_2,\emptyset)$.
                        Thus $y_2 = \emptyset$ by \cref{pair_eq_iff}.
                        Thus $y_1 = y_2$.
                    \end{byCase}
                \end{subproof}
                Thus $U$ is right-unique by \cref{rightunique}.
            \end{subproof}
            Show that $\dom{U} = \naturalsPlus$.
            \begin{subproof}
                Show that for all $n\in\dom{U}$ we have $n\in\naturalsPlus$.
                \begin{subproof}
                    Fix $n\in\dom{U}$.
                    Take $y$ such that $(n,y)\in U$ by \cref{dom_iff}.
                    Then $(n,y)\in U_1$ or
                        $(n,y)\in U_2$
                        by \cref{union_iff}.
                    \begin{byCase}
                    \caseOf{$(n,y)\in U_1$.}
                        Take $n_0\in P$ such that $(n,y) = (n_0,b(n_0))$.
                        Thus $n = n_0$ by \cref{pair_eq_iff}.
                        Thus $n\in P$.
                        Thus $n\in\naturalsPlus$.
                    \caseOf{$(n,y)\in U_2$.}
                        Take $n_0\in\naturalsPlus\setminus P$ such that $(n,y) = (n_0,\emptyset)$.
                        Thus $n = n_0$ by \cref{pair_eq_iff}.
                        Thus $n\in\naturalsPlus\setminus P$ by \cref{setminus_elim_left}.
                        Thus $n\in\naturalsPlus$ by \cref{setminus_elim_left}.
                    \end{byCase}
                \end{subproof}
                Thus $\dom{U}\subseteq\naturalsPlus$ by \cref{subseteq}.
                Show that for all $n\in\naturalsPlus$ we have $n\in\dom{U}$.
                \begin{subproof}
                    Fix $n\in\naturalsPlus$.
                    \begin{byCase}
                    \caseOf{$n\in P$.}
                        Let $w = (n,b(n))$.
                        $w\in U_1$.
                        Thus $w\in U$ by \cref{union_iff}.
                        Thus $n\in\dom{U}$ by \cref{dom_intro}.
                    \caseOf{$n\notin P$.}
                        Let $w = (n,\emptyset)$.
                        $n\in\naturalsPlus\setminus P$ by \cref{setminus_intro}.
                        Thus $w\in U_2$.
                        Thus $w\in U$ by \cref{union_iff}.
                        Thus $n\in\dom{U}$ by \cref{dom_intro}.
                    \end{byCase}
                \end{subproof}
                Thus $\naturalsPlus\subseteq\dom{U}$ by \cref{subseteq}.
                Thus $\dom{U} = \naturalsPlus$ by \cref{subseteq_antisymmetric}.
            \end{subproof}
            Thus $U$ is a function on $\naturalsPlus$.
        \end{subproof}

        Show that $U\in\funs{\naturalsPlus}{\pow{X}}$.
        \begin{subproof}
            Show that for all $n\in\dom{U}$ we have $U(n)\in\pow{X}$.
            \begin{subproof}
                Fix $n\in\dom{U}$.
                $\dom{U} = \naturalsPlus$.
                Thus $n\in\naturalsPlus$ by \cref{subseteq}.
                \begin{byCase}
                \caseOf{$n\in P$.}
                    $(n,U(n))\in U$ by \cref{function_apply_elim}.
                    Show that $(n,U(n))\in U_1$.
                    \begin{subproof}
                        Suppose not.
                        Then $(n,U(n))\in U_2$ by \cref{union_iff}.
                        Take $n_1\in\naturalsPlus\setminus P$ such that $(n,U(n)) = (n_1,\emptyset)$.
                        Thus $n = n_1$ by \cref{pair_eq_iff}.
                        Thus $n\notin P$ by \cref{setminus_elim_right}.
                        Contradiction.
                    \end{subproof}
                    Take $n_0\in P$ such that $(n,U(n)) = (n_0,b(n_0))$.
                    Thus $n = n_0$ by \cref{pair_eq_iff}.
                    $b\in\funs{\naturalsPlus}{\pow{X}}$ by \cref{sequence_in}.
                    Thus $b$ is a function and $\dom{b} = \naturalsPlus$ and for all $m\in\dom{b}$ we have $b(m)\in\pow{X}$
                        by \cref{funs_elim}.
                    Thus $b(n)\in\pow{X}$.
                    Thus $U(n)\in\pow{X}$ by \cref{pair_eq_iff}.
                \caseOf{$n\notin P$.}
                    $(n,U(n))\in U$ by \cref{function_apply_elim}.
                    Show that $(n,U(n))\in U_2$.
                    \begin{subproof}
                        Suppose not.
                        Then $(n,U(n))\in U_1$ by \cref{union_iff}.
                        Take $n_1\in P$ such that $(n,U(n)) = (n_1,b(n_1))$.
                        Thus $n = n_1$ by \cref{pair_eq_iff}.
                        Thus $n\in P$.
                        Contradiction.
                    \end{subproof}
                    Take $n_0\in\naturalsPlus\setminus P$ such that $(n,U(n)) = (n_0,\emptyset)$.
                    Thus $n = n_0$ by \cref{pair_eq_iff}.
                    $\emptyset\subseteq X$ by \cref{emptyset_subseteq}.
                    Thus $\emptyset\in\pow{X}$ by \cref{powerset_intro}.
                    Thus $U(n)\in\pow{X}$ by \cref{pair_eq_iff}.
                \end{byCase}
            \end{subproof}
            Thus $U$ is a function to $\pow{X}$.
            Thus $U\in\funs{\naturalsPlus}{\pow{X}}$ by \cref{funs_intro}.
        \end{subproof}

        Show that $V$ is a function on $\naturalsPlus$.
        \begin{subproof}
            Show that $V$ is a function.
            \begin{subproof}
                Show that for all $w$ such that $w\in V$ there exist $x,y$ such that $w = (x,y)$.
                \begin{subproof}
                    Fix $w$.
                    Assume $w\in V$.
                    Then $w\in V_1$ or
                        $w\in V_2$
                        by \cref{union_iff}.
                    \begin{byCase}
                    \caseOf{$w\in V_1$.}
                        Take $n\in Q$ such that $w = (n,b(n))$.
                        Thus there exist $x,y$ such that $w = (x,y)$.
                    \caseOf{$w\in V_2$.}
                        Take $n\in\naturalsPlus\setminus Q$ such that $w = (n,\emptyset)$.
                        Thus there exist $x,y$ such that $w = (x,y)$.
                    \end{byCase}
                \end{subproof}
                Thus $V$ is a relation by \cref{relation}.
                Show that for all $n,y_1,y_2$ such that $(n,y_1)\in V$ and $(n,y_2)\in V$ we have $y_1 = y_2$.
                \begin{subproof}
                    Fix $n$.
                    Fix $y_1$.
                    Fix $y_2$.
                    Assume $(n,y_1)\in V$ and $(n,y_2)\in V$.
                    \begin{byCase}
                    \caseOf{$n\in Q$.}
                        Show that $(n,y_1)\in V_1$.
                        \begin{subproof}
                            Suppose not.
                            Then $(n,y_1)\in V_2$ by \cref{union_iff}.
                            Take $n_0\in\naturalsPlus\setminus Q$ such that $(n,y_1) = (n_0,\emptyset)$.
                            Thus $n = n_0$ by \cref{pair_eq_iff}.
                            Thus $n\notin Q$ by \cref{setminus_elim_right}.
                            Contradiction.
                        \end{subproof}
                        Take $n_1\in Q$ such that $(n,y_1) = (n_1,b(n_1))$.
                        Thus $y_1 = b(n)$ by \cref{pair_eq_iff}.
                        Show that $(n,y_2)\in V_1$.
                        \begin{subproof}
                            Suppose not.
                            Then $(n,y_2)\in V_2$ by \cref{union_iff}.
                            Take $n_0\in\naturalsPlus\setminus Q$ such that $(n,y_2) = (n_0,\emptyset)$.
                            Thus $n = n_0$ by \cref{pair_eq_iff}.
                            Thus $n\notin Q$ by \cref{setminus_elim_right}.
                            Contradiction.
                        \end{subproof}
                        Take $n_2\in Q$ such that $(n,y_2) = (n_2,b(n_2))$.
                        Thus $y_2 = b(n)$ by \cref{pair_eq_iff}.
                        Thus $y_1 = y_2$.
                    \caseOf{$n\notin Q$.}
                        Show that $(n,y_1)\in V_2$.
                        \begin{subproof}
                            Suppose not.
                            Then $(n,y_1)\in V_1$ by \cref{union_iff}.
                            Take $n_0\in Q$ such that $(n,y_1) = (n_0,b(n_0))$.
                            Thus $n = n_0$ by \cref{pair_eq_iff}.
                            Thus $n\in Q$.
                            Contradiction.
                        \end{subproof}
                        Take $n_1\in\naturalsPlus\setminus Q$ such that $(n,y_1) = (n_1,\emptyset)$.
                        Thus $y_1 = \emptyset$ by \cref{pair_eq_iff}.
                        Show that $(n,y_2)\in V_2$.
                        \begin{subproof}
                            Suppose not.
                            Then $(n,y_2)\in V_1$ by \cref{union_iff}.
                            Take $n_0\in Q$ such that $(n,y_2) = (n_0,b(n_0))$.
                            Thus $n = n_0$ by \cref{pair_eq_iff}.
                            Thus $n\in Q$.
                            Contradiction.
                        \end{subproof}
                        Take $n_2\in\naturalsPlus\setminus Q$ such that $(n,y_2) = (n_2,\emptyset)$.
                        Thus $y_2 = \emptyset$ by \cref{pair_eq_iff}.
                        Thus $y_1 = y_2$.
                    \end{byCase}
                \end{subproof}
                Thus $V$ is right-unique by \cref{rightunique}.
            \end{subproof}
            Show that $\dom{V} = \naturalsPlus$.
            \begin{subproof}
                Show that for all $n\in\dom{V}$ we have $n\in\naturalsPlus$.
                \begin{subproof}
                    Fix $n\in\dom{V}$.
                    Take $y$ such that $(n,y)\in V$ by \cref{dom_iff}.
                    Then $(n,y)\in V_1$ or
                        $(n,y)\in V_2$
                        by \cref{union_iff}.
                    \begin{byCase}
                    \caseOf{$(n,y)\in V_1$.}
                        Take $n_0\in Q$ such that $(n,y) = (n_0,b(n_0))$.
                        Thus $n = n_0$ by \cref{pair_eq_iff}.
                        Thus $n\in Q$.
                        Thus $n\in\naturalsPlus$.
                    \caseOf{$(n,y)\in V_2$.}
                        Take $n_0\in\naturalsPlus\setminus Q$ such that $(n,y) = (n_0,\emptyset)$.
                        Thus $n = n_0$ by \cref{pair_eq_iff}.
                        Thus $n\in\naturalsPlus\setminus Q$ by \cref{setminus_elim_left}.
                        Thus $n\in\naturalsPlus$ by \cref{setminus_elim_left}.
                    \end{byCase}
                \end{subproof}
                Thus $\dom{V}\subseteq\naturalsPlus$ by \cref{subseteq}.
                Show that for all $n\in\naturalsPlus$ we have $n\in\dom{V}$.
                \begin{subproof}
                    Fix $n\in\naturalsPlus$.
                    \begin{byCase}
                    \caseOf{$n\in Q$.}
                        Let $w = (n,b(n))$.
                        $w\in V_1$.
                        Thus $w\in V$ by \cref{union_iff}.
                        Thus $n\in\dom{V}$ by \cref{dom_intro}.
                    \caseOf{$n\notin Q$.}
                        Let $w = (n,\emptyset)$.
                        $n\in\naturalsPlus\setminus Q$ by \cref{setminus_intro}.
                        Thus $w\in V_2$.
                        Thus $w\in V$ by \cref{union_iff}.
                        Thus $n\in\dom{V}$ by \cref{dom_intro}.
                    \end{byCase}
                \end{subproof}
                Thus $\naturalsPlus\subseteq\dom{V}$ by \cref{subseteq}.
                Thus $\dom{V} = \naturalsPlus$ by \cref{subseteq_antisymmetric}.
            \end{subproof}
            Thus $V$ is a function on $\naturalsPlus$.
        \end{subproof}

        Show that $V\in\funs{\naturalsPlus}{\pow{X}}$.
        \begin{subproof}
            Show that for all $n\in\dom{V}$ we have $V(n)\in\pow{X}$.
            \begin{subproof}
                Fix $n\in\dom{V}$.
                $\dom{V} = \naturalsPlus$.
                Thus $n\in\naturalsPlus$ by \cref{subseteq}.
                \begin{byCase}
                \caseOf{$n\in Q$.}
                    $(n,V(n))\in V$ by \cref{function_apply_elim}.
                    Show that $(n,V(n))\in V_1$.
                    \begin{subproof}
                        Suppose not.
                        Then $(n,V(n))\in V_2$ by \cref{union_iff}.
                        Take $n_1\in\naturalsPlus\setminus Q$ such that $(n,V(n)) = (n_1,\emptyset)$.
                        Thus $n = n_1$ by \cref{pair_eq_iff}.
                        Thus $n\notin Q$ by \cref{setminus_elim_right}.
                        Contradiction.
                    \end{subproof}
                    Take $n_0\in Q$ such that $(n,V(n)) = (n_0,b(n_0))$.
                    Thus $n = n_0$ by \cref{pair_eq_iff}.
                    $b\in\funs{\naturalsPlus}{\pow{X}}$ by \cref{sequence_in}.
                    Thus $b$ is a function and $\dom{b} = \naturalsPlus$ and for all $m\in\dom{b}$ we have $b(m)\in\pow{X}$
                        by \cref{funs_elim}.
                    Thus $b(n)\in\pow{X}$.
                    Thus $V(n)\in\pow{X}$ by \cref{pair_eq_iff}.
                \caseOf{$n\notin Q$.}
                    $(n,V(n))\in V$ by \cref{function_apply_elim}.
                    Show that $(n,V(n))\in V_2$.
                    \begin{subproof}
                        Suppose not.
                        Then $(n,V(n))\in V_1$ by \cref{union_iff}.
                        Take $n_1\in Q$ such that $(n,V(n)) = (n_1,b(n_1))$.
                        Thus $n = n_1$ by \cref{pair_eq_iff}.
                        Thus $n\in Q$.
                        Contradiction.
                    \end{subproof}
                    Take $n_0\in\naturalsPlus\setminus Q$ such that $(n,V(n)) = (n_0,\emptyset)$.
                    Thus $n = n_0$ by \cref{pair_eq_iff}.
                    $\emptyset\subseteq X$ by \cref{emptyset_subseteq}.
                    Thus $\emptyset\in\pow{X}$ by \cref{powerset_intro}.
                    Thus $V(n)\in\pow{X}$ by \cref{pair_eq_iff}.
                \end{byCase}
            \end{subproof}
            Thus $V$ is a function to $\pow{X}$.
            Thus $V\in\funs{\naturalsPlus}{\pow{X}}$ by \cref{funs_intro}.
        \end{subproof}

        Show that for all $n\in\naturalsPlus$ we have $U(n)$ is open in $X$ with respect to $T$ and
            $\closure{U(n)}{X}{T}$ is disjoint from $B$.
        \begin{subproof}
            Fix $n\in\naturalsPlus$.
            \begin{byCase}
            \caseOf{$n\in P$.}
                $b\in\funs{\naturalsPlus}{\pow{X}}$ by \cref{sequence_in}.
                Thus $b$ is a function by \cref{funs_elim}.
                Thus $b(n)\in\pow{X}$ by \cref{funs_type_apply}.
                $b(n)$ is open in $X$ with respect to $T$ by \cref{second_countability_axiom}.
                Thus $(n,b(n))\in U_1$.
                Thus $(n,b(n))\in U$ by \cref{union_intro_left}.
                $(n,U(n))\in U$ by \cref{function_apply_elim}.
                $U$ is right-unique by \cref{function_rightunique}.
                Thus $U(n) = b(n)$ by \hyperref[rightunique]{right-uniqueness}.
                Thus $U(n)$ is open in $X$ with respect to $T$.
                Show that there exists no $x$ such that $x\in\closure{U(n)}{X}{T}$ and $x\in B$.
                \begin{subproof}
                    Suppose not.
                    Then there exists $x$ such that $x\in\closure{U(n)}{X}{T}$ and $x\in B$.
                    Thus $x\in\closure{b(n)}{X}{T}$.
                    Thus $x\in\closure{b(n)}{X}{T}\inter B$ by \cref{inter_intro}.
                    $\closure{b(n)}{X}{T}\inter B = \emptyset$.
                    Thus $x\in\emptyset$.
                    Contradiction by \cref{emptyset}.
                \end{subproof}
                Thus $\closure{U(n)}{X}{T}$ is disjoint from $B$ by \cref{disjoint}.
            \caseOf{$n\notin P$.}
                $n\in\naturalsPlus\setminus P$ by \cref{setminus_intro}.
                Thus $(n,\emptyset)\in U_2$.
                Thus $(n,\emptyset)\in U$ by \cref{union_intro_right}.
                $(n,U(n))\in U$ by \cref{function_apply_elim}.
                $U$ is right-unique by \cref{function_rightunique}.
                Thus $U(n) = \emptyset$ by \hyperref[rightunique]{right-uniqueness}.
                Thus $\emptyset\in T$ by \cref{topology_on}.
                Thus $U(n)$ is open in $X$ with respect to $T$ by \cref{open_in}.
                $\emptyset$ is closed in $X$ with respect to $T$ by \cref{closed_emptyset}.
                Thus $U(n)\subseteq X$ by \cref{emptyset_subseteq}.
                Thus $U(n)\subseteq\emptyset$ by \cref{subseteq_emptyset_iff}.
                Thus $\closure{U(n)}{X}{T}\subseteq\emptyset$ by \cref{closure_subseteq_closed}.
                Show that there exists no $x$ such that $x\in\closure{U(n)}{X}{T}$ and $x\in B$.
                \begin{subproof}
                    Suppose not.
                    Then there exists $x$ such that $x\in\closure{U(n)}{X}{T}$ and $x\in B$.
                    Thus $x\in\emptyset$ by \cref{subseteq}.
                    Contradiction by \cref{emptyset}.
                \end{subproof}
                Thus $\closure{U(n)}{X}{T}$ is disjoint from $B$ by \cref{disjoint}.
            \end{byCase}
        \end{subproof}

        Show that for all $n\in\naturalsPlus$ we have $V(n)$ is open in $X$ with respect to $T$ and
            $\closure{V(n)}{X}{T}$ is disjoint from $A$.
        \begin{subproof}
            Fix $n\in\naturalsPlus$.
            \begin{byCase}
            \caseOf{$n\in Q$.}
                $b\in\funs{\naturalsPlus}{\pow{X}}$ by \cref{sequence_in}.
                Thus $b$ is a function by \cref{funs_elim}.
                Thus $b(n)\in\pow{X}$ by \cref{funs_type_apply}.
                $b(n)$ is open in $X$ with respect to $T$ by \cref{second_countability_axiom}.
                Thus $(n,b(n))\in V_1$.
                Thus $(n,b(n))\in V$ by \cref{union_intro_left}.
                $(n,V(n))\in V$ by \cref{function_apply_elim}.
                $V$ is right-unique by \cref{function_rightunique}.
                Thus $V(n) = b(n)$ by \hyperref[rightunique]{right-uniqueness}.
                Thus $V(n)$ is open in $X$ with respect to $T$.
                Show that there exists no $x$ such that $x\in\closure{V(n)}{X}{T}$ and $x\in A$.
                \begin{subproof}
                    Suppose not.
                    Then there exists $x$ such that $x\in\closure{V(n)}{X}{T}$ and $x\in A$.
                    Thus $x\in\closure{b(n)}{X}{T}$.
                    Thus $x\in\closure{b(n)}{X}{T}\inter A$ by \cref{inter_intro}.
                    $\closure{b(n)}{X}{T}\inter A = \emptyset$.
                    Thus $x\in\emptyset$.
                    Contradiction by \cref{emptyset}.
                \end{subproof}
                Thus $\closure{V(n)}{X}{T}$ is disjoint from $A$ by \cref{disjoint}.
            \caseOf{$n\notin Q$.}
                $n\in\naturalsPlus\setminus Q$ by \cref{setminus_intro}.
                Thus $(n,\emptyset)\in V_2$.
                Thus $(n,\emptyset)\in V$ by \cref{union_intro_right}.
                $(n,V(n))\in V$ by \cref{function_apply_elim}.
                $V$ is right-unique by \cref{function_rightunique}.
                Thus $V(n) = \emptyset$ by \hyperref[rightunique]{right-uniqueness}.
                Thus $\emptyset\in T$ by \cref{topology_on}.
                Thus $V(n)$ is open in $X$ with respect to $T$ by \cref{open_in}.
                $\emptyset$ is closed in $X$ with respect to $T$ by \cref{closed_emptyset}.
                Thus $V(n)\subseteq X$ by \cref{emptyset_subseteq}.
                Thus $V(n)\subseteq\emptyset$ by \cref{subseteq_emptyset_iff}.
                Thus $\closure{V(n)}{X}{T}\subseteq\emptyset$ by \cref{closure_subseteq_closed}.
                Show that there exists no $x$ such that $x\in\closure{V(n)}{X}{T}$ and $x\in A$.
                \begin{subproof}
                    Suppose not.
                    Then there exists $x$ such that $x\in\closure{V(n)}{X}{T}$ and $x\in A$.
                    Thus $x\in\emptyset$ by \cref{subseteq}.
                    Contradiction by \cref{emptyset}.
                \end{subproof}
                Thus $\closure{V(n)}{X}{T}$ is disjoint from $A$ by \cref{disjoint}.
            \end{byCase}
        \end{subproof}

        Show that $A\subseteq\unions{\img{U}{\naturalsPlus}}$.
        \begin{subproof}
            Show that for all $x$ such that $x\in A$ we have $x\in\unions{\img{U}{\naturalsPlus}}$.
            \begin{subproof}
                Fix $x$.
                Assume $x\in A$.
                $A\subseteq X$ by \cref{closed_in}.
                Thus $x\in X$ by \cref{subseteq}.
                Show that $x\notin B$.
                \begin{subproof}
                    Suppose not.
                    Then $x\in B$.
                    Thus there exists $y$ such that $y\in A$ and $y\in B$ by \cref{disjoint}.
                    Contradiction.
                \end{subproof}
                Take $U_0,V_0$ such that
                    $U_0$ is open in $X$ with respect to $T$ and
                    $V_0$ is open in $X$ with respect to $T$ and
                    $x\in U_0$ and
                    $B\subseteq V_0$ and
                    $U_0$ is disjoint from $V_0$
                    by \cref{regular_space}.
                Thus there exists $n\in\naturalsPlus$ such that $x\in b(n)$ and $b(n)\subseteq U_0$.
                Take $n\in\naturalsPlus$ such that $x\in b(n)$ and $b(n)\subseteq U_0$.
                Show that $b(n)$ intersects $A$.
                \begin{subproof}
                    $x\in b(n)$ and $x\in A$.
                    Thus $x\in b(n)\inter A$ by \cref{inter_intro}.
                    Thus $b(n)\inter A\neq\emptyset$ by \cref{emptyset}.
                    Thus $b(n)$ intersects $A$ by \cref{intersects}.
                \end{subproof}
                Show that $\closure{b(n)}{X}{T}$ is disjoint from $B$.
                \begin{subproof}
                    $V_0$ is open in $X$ with respect to $T$ by \cref{open_in}.
                    $V_0\in T$ by \cref{open_in}.
                    $T\subseteq\pow{X}$ by \cref{topology_on}.
                    Thus $V_0\in\pow{X}$ by \cref{subseteq}.
                    Thus $V_0\subseteq X$ by \cref{powerset_elim}.
                    Thus $X\setminus V_0\subseteq X$ by \cref{setminus_subseteq}.
                    Thus $X\setminus (X\setminus V_0) = V_0$ by \cref{double_relative_complement}.
                    Thus $X\setminus (X\setminus V_0)$ is open in $X$ with respect to $T$.
                    Thus $X\setminus V_0$ is closed in $X$ with respect to $T$ by \cref{closed_in}.
                    $U_0$ is disjoint from $V_0$ by \cref{disjoint}.
                    Show that $U_0\subseteq X\setminus V_0$.
                    \begin{subproof}
                        Show that for all $y$ such that $y\in U_0$ we have $y\in X\setminus V_0$.
                        \begin{subproof}
                            Fix $y$.
                            Assume $y\in U_0$.
                            Show that $y\notin V_0$.
                            \begin{subproof}
                                Suppose not.
                                Then $y\in V_0$.
                                Thus there exists $z$ such that $z\in U_0$ and $z\in V_0$.
                                Contradiction by \cref{disjoint}.
                            \end{subproof}
                            $U_0\in T$ by \cref{open_in}.
                            $T\subseteq\pow{X}$ by \cref{topology_on}.
                            Thus $U_0\in\pow{X}$ by \cref{subseteq}.
                            Thus $U_0\subseteq X$ by \cref{powerset_elim}.
                            Thus $y\in X$ by \cref{subseteq}.
                            Thus $y\in X\setminus V_0$ by \cref{setminus_intro}.
                        \end{subproof}
                        Thus $U_0\subseteq X\setminus V_0$ by \cref{subseteq}.
                    \end{subproof}
                    $b(n)\subseteq U_0$.
                    Thus $b(n)\subseteq X\setminus V_0$ by \cref{subseteq_transitive}.
                    $b\in\funs{\naturalsPlus}{\pow{X}}$ by \cref{sequence_in}.
                    Thus $b$ is a function by \cref{funs_elim}.
                    Thus $b(n)\in\pow{X}$ by \cref{funs_type_apply}.
                    Thus $b(n)\subseteq X$ by \cref{powerset_elim}.
                    Thus $\closure{b(n)}{X}{T}\subseteq X\setminus V_0$ by \cref{closure_subseteq_closed}.
                    Show that $X\setminus V_0$ is disjoint from $B$.
                    \begin{subproof}
                        Suppose not.
                        Then there exists $y$ such that $y\in X\setminus V_0$ and $y\in B$ by \cref{disjoint}.
                        $B\subseteq V_0$ by \cref{subseteq}.
                        Thus $y\in V_0$.
                        Thus $y\notin V_0$ by \cref{setminus_elim_right}.
                        Contradiction.
                    \end{subproof}
                    Thus $\closure{b(n)}{X}{T}$ is disjoint from $B$ by \cref{disjoint,subseteq}.
                \end{subproof}
                $b(n)$ intersects $A$.
                Thus $b(n)\inter A\neq\emptyset$ by \cref{intersects}.
                Show that $\closure{b(n)}{X}{T}\inter B = \emptyset$.
                \begin{subproof}
                    Suppose not.
                    Then $\closure{b(n)}{X}{T}\inter B \neq \emptyset$.
                    Then $\closure{b(n)}{X}{T}\inter B$ is inhabited by \cref{nonempty_inhabited}.
                    Take $y$ such that $y\in\closure{b(n)}{X}{T}\inter B$.
                    Thus $y\in\closure{b(n)}{X}{T}$ and $y\in B$ by \cref{inter}.
                    Thus there exists $y$ such that $y\in\closure{b(n)}{X}{T}$ and $y\in B$.
                    Contradiction by \cref{disjoint}.
                \end{subproof}
                Thus $n\in P$.
                Thus $(n,b(n))\in U_1$.
                Thus $(n,b(n))\in U$ by \cref{union_intro_left}.
                $(n,U(n))\in U$ by \cref{function_apply_elim}.
                $U$ is right-unique by \cref{function_rightunique}.
                Thus $U(n) = b(n)$ by \hyperref[rightunique]{right-uniqueness}.
                Thus $x\in U(n)$.
                Thus $x\in\unions{\img{U}{\naturalsPlus}}$ by \cref{unions_intro,img_iff}.
            \end{subproof}
            Thus $A\subseteq\unions{\img{U}{\naturalsPlus}}$ by \cref{subseteq}.
        \end{subproof}

        Show that $B\subseteq\unions{\img{V}{\naturalsPlus}}$.
        \begin{subproof}
            Show that for all $x$ such that $x\in B$ we have $x\in\unions{\img{V}{\naturalsPlus}}$.
            \begin{subproof}
                Fix $x$.
                Assume $x\in B$.
                $B\subseteq X$ by \cref{closed_in}.
                Thus $x\in X$ by \cref{subseteq}.
                Show that $x\notin A$.
                \begin{subproof}
                    Suppose not.
                    Then $x\in A$.
                    Thus there exists $y$ such that $y\in A$ and $y\in B$ by \cref{disjoint}.
                    Contradiction.
                \end{subproof}
                Take $U_0,V_0$ such that
                    $U_0$ is open in $X$ with respect to $T$ and
                    $V_0$ is open in $X$ with respect to $T$ and
                    $x\in U_0$ and
                    $A\subseteq V_0$ and
                    $U_0$ is disjoint from $V_0$
                    by \cref{regular_space}.
                Thus there exists $n\in\naturalsPlus$ such that $x\in b(n)$ and $b(n)\subseteq U_0$.
                Take $n\in\naturalsPlus$ such that $x\in b(n)$ and $b(n)\subseteq U_0$.
                Show that $b(n)$ intersects $B$.
                \begin{subproof}
                    $x\in b(n)$ and $x\in B$.
                    Thus $x\in b(n)\inter B$ by \cref{inter_intro}.
                    Thus $b(n)\inter B\neq\emptyset$ by \cref{emptyset}.
                    Thus $b(n)$ intersects $B$ by \cref{intersects}.
                \end{subproof}
                Show that $\closure{b(n)}{X}{T}$ is disjoint from $A$.
                \begin{subproof}
                    $V_0$ is open in $X$ with respect to $T$ by \cref{open_in}.
                    $V_0\in T$ by \cref{open_in}.
                    $T\subseteq\pow{X}$ by \cref{topology_on}.
                    Thus $V_0\in\pow{X}$ by \cref{subseteq}.
                    Thus $V_0\subseteq X$ by \cref{powerset_elim}.
                    Thus $X\setminus V_0\subseteq X$ by \cref{setminus_subseteq}.
                    Thus $X\setminus (X\setminus V_0) = V_0$ by \cref{double_relative_complement}.
                    Thus $X\setminus (X\setminus V_0)$ is open in $X$ with respect to $T$.
                    Thus $X\setminus V_0$ is closed in $X$ with respect to $T$ by \cref{closed_in}.
                    $U_0$ is disjoint from $V_0$ by \cref{disjoint}.
                    Show that $U_0\subseteq X\setminus V_0$.
                    \begin{subproof}
                        Show that for all $y$ such that $y\in U_0$ we have $y\in X\setminus V_0$.
                        \begin{subproof}
                            Fix $y$.
                            Assume $y\in U_0$.
                            Show that $y\notin V_0$.
                            \begin{subproof}
                                Suppose not.
                                Then $y\in V_0$.
                                Thus there exists $z$ such that $z\in U_0$ and $z\in V_0$.
                                Contradiction by \cref{disjoint}.
                            \end{subproof}
                            $U_0\in T$ by \cref{open_in}.
                            $T\subseteq\pow{X}$ by \cref{topology_on}.
                            Thus $U_0\in\pow{X}$ by \cref{subseteq}.
                            Thus $U_0\subseteq X$ by \cref{powerset_elim}.
                            Thus $y\in X$ by \cref{subseteq}.
                            Thus $y\in X\setminus V_0$ by \cref{setminus_intro}.
                        \end{subproof}
                        Thus $U_0\subseteq X\setminus V_0$ by \cref{subseteq}.
                    \end{subproof}
                    $b(n)\subseteq U_0$.
                    Thus $b(n)\subseteq X\setminus V_0$ by \cref{subseteq_transitive}.
                    $b\in\funs{\naturalsPlus}{\pow{X}}$ by \cref{sequence_in}.
                    Thus $b$ is a function by \cref{funs_elim}.
                    Thus $b(n)\in\pow{X}$ by \cref{funs_type_apply}.
                    Thus $b(n)\subseteq X$ by \cref{powerset_elim}.
                    Thus $\closure{b(n)}{X}{T}\subseteq X\setminus V_0$ by \cref{closure_subseteq_closed}.
                    Show that $X\setminus V_0$ is disjoint from $A$.
                    \begin{subproof}
                        Suppose not.
                        Then there exists $y$ such that $y\in X\setminus V_0$ and $y\in A$ by \cref{disjoint}.
                        $A\subseteq V_0$ by \cref{subseteq}.
                        Thus $y\in V_0$.
                        Thus $y\notin V_0$ by \cref{setminus_elim_right}.
                        Contradiction.
                    \end{subproof}
                    Thus $\closure{b(n)}{X}{T}$ is disjoint from $A$ by \cref{disjoint,subseteq}.
                \end{subproof}
                $b(n)$ intersects $B$.
                Thus $b(n)\inter B\neq\emptyset$ by \cref{intersects}.
                Show that $\closure{b(n)}{X}{T}\inter A = \emptyset$.
                \begin{subproof}
                    Suppose not.
                    Then $\closure{b(n)}{X}{T}\inter A \neq \emptyset$.
                    Then $\closure{b(n)}{X}{T}\inter A$ is inhabited by \cref{nonempty_inhabited}.
                    Take $y$ such that $y\in\closure{b(n)}{X}{T}\inter A$.
                    Thus $y\in\closure{b(n)}{X}{T}$ and $y\in A$ by \cref{inter}.
                    Thus there exists $y$ such that $y\in\closure{b(n)}{X}{T}$ and $y\in A$.
                    Contradiction by \cref{disjoint}.
                \end{subproof}
                Thus $n\in Q$.
                Thus $(n,b(n))\in V_1$.
                Thus $(n,b(n))\in V$ by \cref{union_intro_left}.
                $(n,V(n))\in V$ by \cref{function_apply_elim}.
                $V$ is right-unique by \cref{function_rightunique}.
                Thus $V(n) = b(n)$ by \hyperref[rightunique]{right-uniqueness}.
                Thus $x\in V(n)$.
                Thus $x\in\unions{\img{V}{\naturalsPlus}}$ by \cref{unions_intro,img_iff}.
            \end{subproof}
            Thus $B\subseteq\unions{\img{V}{\naturalsPlus}}$ by \cref{subseteq}.
        \end{subproof}

        Let $c = \{w \mid \exists n\in\naturalsPlus. w = (n,\closure{V(n)}{X}{T})\}$.
        Let $d = \{w \mid \exists n\in\naturalsPlus. w = (n,\closure{U(n)}{X}{T})\}$.
        Show that $c$ is a function on $\naturalsPlus$.
        \begin{subproof}
            Show that $c$ is a function.
            \begin{subproof}
                Show that for all $w$ such that $w\in c$ there exist $x,y$ such that $w = (x,y)$.
                \begin{subproof}
                    Fix $w$.
                    Assume $w\in c$.
                    Take $n\in\naturalsPlus$ such that $w = (n,\closure{V(n)}{X}{T})$.
                    Thus there exist $x,y$ such that $w = (x,y)$.
                \end{subproof}
                Thus $c$ is a relation by \cref{relation}.
                Show that for all $n,y_1,y_2$ such that $(n,y_1)\in c$ and $(n,y_2)\in c$ we have $y_1 = y_2$.
                \begin{subproof}
                    Fix $n$.
                    Fix $y_1$.
                    Fix $y_2$.
                    Assume $(n,y_1)\in c$ and $(n,y_2)\in c$.
                    Take $n_1\in\naturalsPlus$ such that $(n,y_1) = (n_1,\closure{V(n_1)}{X}{T})$.
                    Take $n_2\in\naturalsPlus$ such that $(n,y_2) = (n_2,\closure{V(n_2)}{X}{T})$.
                    Thus $n = n_1$ and $n = n_2$ by \cref{pair_eq_iff}.
                    Thus $y_1 = \closure{V(n)}{X}{T}$ and $y_2 = \closure{V(n)}{X}{T}$ by \cref{pair_eq_iff}.
                    Thus $y_1 = y_2$.
                \end{subproof}
                Thus $c$ is right-unique by \cref{rightunique}.
            \end{subproof}
            Show that $\dom{c} = \naturalsPlus$.
            \begin{subproof}
                Show that for all $n\in\dom{c}$ we have $n\in\naturalsPlus$.
                \begin{subproof}
                    Fix $n\in\dom{c}$.
                    Take $y$ such that $(n,y)\in c$ by \cref{dom_iff}.
                    Take $n_0\in\naturalsPlus$ such that $(n,y) = (n_0,\closure{V(n_0)}{X}{T})$.
                    Thus $n = n_0$ by \cref{pair_eq_iff}.
                    Thus $n\in\naturalsPlus$.
                \end{subproof}
                Thus $\dom{c}\subseteq\naturalsPlus$ by \cref{subseteq}.
                Show that for all $n\in\naturalsPlus$ we have $n\in\dom{c}$.
                \begin{subproof}
                    Fix $n\in\naturalsPlus$.
                    Let $w = (n,\closure{V(n)}{X}{T})$.
                    $w\in c$.
                    Thus $n\in\dom{c}$ by \cref{dom_intro}.
                \end{subproof}
                Thus $\naturalsPlus\subseteq\dom{c}$ by \cref{subseteq}.
                Thus $\dom{c} = \naturalsPlus$ by \cref{subseteq_antisymmetric}.
            \end{subproof}
            Thus $c$ is a function on $\naturalsPlus$.
        \end{subproof}

        Show that $d$ is a function on $\naturalsPlus$.
        \begin{subproof}
            Show that $d$ is a function.
            \begin{subproof}
                Show that for all $w$ such that $w\in d$ there exist $x,y$ such that $w = (x,y)$.
                \begin{subproof}
                    Fix $w$.
                    Assume $w\in d$.
                    Take $n\in\naturalsPlus$ such that $w = (n,\closure{U(n)}{X}{T})$.
                    Thus there exist $x,y$ such that $w = (x,y)$.
                \end{subproof}
                Thus $d$ is a relation by \cref{relation}.
                Show that for all $n,y_1,y_2$ such that $(n,y_1)\in d$ and $(n,y_2)\in d$ we have $y_1 = y_2$.
                \begin{subproof}
                    Fix $n$.
                    Fix $y_1$.
                    Fix $y_2$.
                    Assume $(n,y_1)\in d$ and $(n,y_2)\in d$.
                    Take $n_1\in\naturalsPlus$ such that $(n,y_1) = (n_1,\closure{U(n_1)}{X}{T})$.
                    Take $n_2\in\naturalsPlus$ such that $(n,y_2) = (n_2,\closure{U(n_2)}{X}{T})$.
                    Thus $n = n_1$ and $n = n_2$ by \cref{pair_eq_iff}.
                    Thus $y_1 = \closure{U(n)}{X}{T}$ and $y_2 = \closure{U(n)}{X}{T}$ by \cref{pair_eq_iff}.
                    Thus $y_1 = y_2$.
                \end{subproof}
                Thus $d$ is right-unique by \cref{rightunique}.
            \end{subproof}
            Show that $\dom{d} = \naturalsPlus$.
            \begin{subproof}
                Show that for all $n\in\dom{d}$ we have $n\in\naturalsPlus$.
                \begin{subproof}
                    Fix $n\in\dom{d}$.
                    Take $y$ such that $(n,y)\in d$ by \cref{dom_iff}.
                    Take $n_0\in\naturalsPlus$ such that $(n,y) = (n_0,\closure{U(n_0)}{X}{T})$.
                    Thus $n = n_0$ by \cref{pair_eq_iff}.
                    Thus $n\in\naturalsPlus$.
                \end{subproof}
                Thus $\dom{d}\subseteq\naturalsPlus$ by \cref{subseteq}.
                Show that for all $n\in\naturalsPlus$ we have $n\in\dom{d}$.
                \begin{subproof}
                    Fix $n\in\naturalsPlus$.
                    Let $w = (n,\closure{U(n)}{X}{T})$.
                    $w\in d$.
                    Thus $n\in\dom{d}$ by \cref{dom_intro}.
                \end{subproof}
                Thus $\naturalsPlus\subseteq\dom{d}$ by \cref{subseteq}.
                Thus $\dom{d} = \naturalsPlus$ by \cref{subseteq_antisymmetric}.
            \end{subproof}
            Thus $d$ is a function on $\naturalsPlus$.
        \end{subproof}

        Let $F = \{W \mid \exists n\in\naturalsPlus. W = U(n)\setminus \unions{\img{\restrl{c}{\initialSegment{n}}}{\initialSegment{n}}}\}$.
        Let $G = \{W \mid \exists n\in\naturalsPlus. W = V(n)\setminus \unions{\img{\restrl{d}{\initialSegment{n}}}{\initialSegment{n}}}\}$.
        Show that $F\subseteq T$.
        \begin{subproof}
            Show that for all $W$ such that $W\in F$ we have $W\in T$.
            \begin{subproof}
                Fix $W$.
                Assume $W\in F$.
                Take $n\in\naturalsPlus$ such that $W = U(n)\setminus \unions{\img{\restrl{c}{\initialSegment{n}}}{\initialSegment{n}}}$.
                Show that $\unions{\img{\restrl{c}{\initialSegment{n}}}{\initialSegment{n}}}$ is closed in $X$ with respect to $T$.
                \begin{subproof}
                    $\restrl{c}{\initialSegment{n}}$ is a function by \cref{restrl_of_function_is_function}.
                    $\initialSegment{n}$ is finite by \cref{initial_segment_finite}.
                    $c$ is a function on $\naturalsPlus$.
                    Thus $\dom{c} = \naturalsPlus$.
                    Thus $\dom{\restrl{c}{\initialSegment{n}}} = \dom{c}\inter \initialSegment{n}$ by \cref{dom_of_restrl_eq_inter}.
                    Thus $\dom{\restrl{c}{\initialSegment{n}}} = \naturalsPlus\inter \initialSegment{n}$.
                    Show that $\initialSegment{n}\subseteq\naturalsPlus$.
                    \begin{subproof}
                        Show that for all $k$ such that $k\in\initialSegment{n}$ we have $k\in\naturalsPlus$.
                        \begin{subproof}
                            Fix $k$.
                            Assume $k\in\initialSegment{n}$.
                            Thus $k\in\naturalsPlus$ by \cref{initial_segment_naturalsplus}.
                        \end{subproof}
                        Thus $\initialSegment{n}\subseteq\naturalsPlus$ by \cref{subseteq}.
                    \end{subproof}
                    Thus $\dom{\restrl{c}{\initialSegment{n}}} = \initialSegment{n}$ by \cref{inter_absorb_supseteq_left}.
                    Thus $\img{\restrl{c}{\initialSegment{n}}}{\initialSegment{n}}$ is finite by \cref{finite_image_function}.
                    Show that for all $C$ such that $C\in \img{\restrl{c}{\initialSegment{n}}}{\initialSegment{n}}$ we have
                        $C$ is closed in $X$ with respect to $T$.
                    \begin{subproof}
                        Fix $C$.
                        Assume $C\in \img{\restrl{c}{\initialSegment{n}}}{\initialSegment{n}}$.
                        Take $k\in\initialSegment{n}$ such that $k\mathrel{\restrl{c}{\initialSegment{n}}} C$ by \cref{img_iff}.
                        Thus $(k,C)\in \restrl{c}{\initialSegment{n}}$ by \cref{pair_eq_iff}.
                        Thus $(k,C)\in c$ by \cref{restrl_iff}.
                        Take $k_0\in\naturalsPlus$ such that $(k,C) = (k_0,\closure{V(k_0)}{X}{T})$.
                        Thus $C = \closure{V(k_0)}{X}{T}$ by \cref{pair_eq_iff}.
                        $V\in\funs{\naturalsPlus}{\pow{X}}$.
                        Thus $V$ is a function by \cref{funs_elim}.
                        Thus $V(k_0)\in\pow{X}$ by \cref{funs_type_apply}.
                        Thus $V(k_0)\subseteq X$ by \cref{powerset_elim}.
                        Thus $C$ is closed in $X$ with respect to $T$ by \cref{closure_closed_in}.
                    \end{subproof}
                    Show that for all $C$ such that $C\in \img{\restrl{c}{\initialSegment{n}}}{\initialSegment{n}}$ we have
                        $C\in\pow{X}$.
                    \begin{subproof}
                        Fix $C$.
                        Assume $C\in \img{\restrl{c}{\initialSegment{n}}}{\initialSegment{n}}$.
                        Thus $C$ is closed in $X$ with respect to $T$.
                        Thus $C\subseteq X$ by \cref{closed_in}.
                        Thus $C\in\pow{X}$ by \cref{powerset_intro}.
                    \end{subproof}
                    Thus $\img{\restrl{c}{\initialSegment{n}}}{\initialSegment{n}}\subseteq\pow{X}$ by \cref{subseteq}.
                    Thus $\unions{\img{\restrl{c}{\initialSegment{n}}}{\initialSegment{n}}}$ is closed in $X$ with respect to $T$
                        by \cref{closed_finite_unions}.
                \end{subproof}
                $U(n)$ is open in $X$ with respect to $T$.
                $U(n)\in T$ by \cref{open_in}.
                $T\subseteq\pow{X}$ by \cref{topology_on}.
                Thus $U(n)\in\pow{X}$ by \cref{subseteq}.
                Thus $U(n)\subseteq X$ by \cref{powerset_elim}.
                $\unions{\img{\restrl{c}{\initialSegment{n}}}{\initialSegment{n}}}$ is closed in $X$ with respect to $T$.
                Thus $\unions{\img{\restrl{c}{\initialSegment{n}}}{\initialSegment{n}}}\subseteq X$ by \cref{closed_in}.
                Thus $W = U(n)\inter (X\setminus \unions{\img{\restrl{c}{\initialSegment{n}}}{\initialSegment{n}}})$
                    by \cref{setminus_eq_inter_complement}.
                $X\setminus \unions{\img{\restrl{c}{\initialSegment{n}}}{\initialSegment{n}}}$ is open in $X$ with respect to $T$
                    by \cref{closed_in}.
                Thus $X\setminus \unions{\img{\restrl{c}{\initialSegment{n}}}{\initialSegment{n}}}\in T$ by \cref{open_in}.
                Thus $W\in T$ by \cref{topology_on}.
                Thus $W$ is open in $X$ with respect to $T$ by \cref{open_in}.
            \end{subproof}
            Thus $F\subseteq T$ by \cref{subseteq}.
        \end{subproof}

        Show that $G\subseteq T$.
        \begin{subproof}
            Show that for all $W$ such that $W\in G$ we have $W\in T$.
            \begin{subproof}
                Fix $W$.
                Assume $W\in G$.
                Take $n\in\naturalsPlus$ such that $W = V(n)\setminus \unions{\img{\restrl{d}{\initialSegment{n}}}{\initialSegment{n}}}$.
                Show that $\unions{\img{\restrl{d}{\initialSegment{n}}}{\initialSegment{n}}}$ is closed in $X$ with respect to $T$.
                \begin{subproof}
                    $\restrl{d}{\initialSegment{n}}$ is a function by \cref{restrl_of_function_is_function}.
                    $\initialSegment{n}$ is finite by \cref{initial_segment_finite}.
                    $d$ is a function on $\naturalsPlus$.
                    Thus $\dom{d} = \naturalsPlus$.
                    Thus $\dom{\restrl{d}{\initialSegment{n}}} = \dom{d}\inter \initialSegment{n}$ by \cref{dom_of_restrl_eq_inter}.
                    Thus $\dom{\restrl{d}{\initialSegment{n}}} = \naturalsPlus\inter \initialSegment{n}$.
                    Show that $\initialSegment{n}\subseteq\naturalsPlus$.
                    \begin{subproof}
                        Show that for all $k$ such that $k\in\initialSegment{n}$ we have $k\in\naturalsPlus$.
                        \begin{subproof}
                            Fix $k$.
                            Assume $k\in\initialSegment{n}$.
                            Thus $k\in\naturalsPlus$ by \cref{initial_segment_naturalsplus}.
                        \end{subproof}
                        Thus $\initialSegment{n}\subseteq\naturalsPlus$ by \cref{subseteq}.
                    \end{subproof}
                    Thus $\dom{\restrl{d}{\initialSegment{n}}} = \initialSegment{n}$ by \cref{inter_absorb_supseteq_left}.
                    Thus $\img{\restrl{d}{\initialSegment{n}}}{\initialSegment{n}}$ is finite by \cref{finite_image_function}.
                    Show that for all $C$ such that $C\in \img{\restrl{d}{\initialSegment{n}}}{\initialSegment{n}}$ we have
                        $C$ is closed in $X$ with respect to $T$.
                    \begin{subproof}
                        Fix $C$.
                        Assume $C\in \img{\restrl{d}{\initialSegment{n}}}{\initialSegment{n}}$.
                        Take $k\in\initialSegment{n}$ such that $k\mathrel{\restrl{d}{\initialSegment{n}}} C$ by \cref{img_iff}.
                        Thus $(k,C)\in \restrl{d}{\initialSegment{n}}$ by \cref{pair_eq_iff}.
                        Thus $(k,C)\in d$ by \cref{restrl_iff}.
                        Take $k_0\in\naturalsPlus$ such that $(k,C) = (k_0,\closure{U(k_0)}{X}{T})$.
                        Thus $C = \closure{U(k_0)}{X}{T}$ by \cref{pair_eq_iff}.
                        $U\in\funs{\naturalsPlus}{\pow{X}}$.
                        Thus $U$ is a function by \cref{funs_elim}.
                        Thus $U(k_0)\in\pow{X}$ by \cref{funs_type_apply}.
                        Thus $U(k_0)\subseteq X$ by \cref{powerset_elim}.
                        Thus $C$ is closed in $X$ with respect to $T$ by \cref{closure_closed_in}.
                    \end{subproof}
                    Show that for all $C$ such that $C\in \img{\restrl{d}{\initialSegment{n}}}{\initialSegment{n}}$ we have
                        $C\in\pow{X}$.
                    \begin{subproof}
                        Fix $C$.
                        Assume $C\in \img{\restrl{d}{\initialSegment{n}}}{\initialSegment{n}}$.
                        Thus $C$ is closed in $X$ with respect to $T$.
                        Thus $C\subseteq X$ by \cref{closed_in}.
                        Thus $C\in\pow{X}$ by \cref{powerset_intro}.
                    \end{subproof}
                    Thus $\img{\restrl{d}{\initialSegment{n}}}{\initialSegment{n}}\subseteq\pow{X}$ by \cref{subseteq}.
                    Thus $\unions{\img{\restrl{d}{\initialSegment{n}}}{\initialSegment{n}}}$ is closed in $X$ with respect to $T$
                        by \cref{closed_finite_unions}.
                \end{subproof}
                $V(n)$ is open in $X$ with respect to $T$.
                $V(n)\in T$ by \cref{open_in}.
                $T\subseteq\pow{X}$ by \cref{topology_on}.
                Thus $V(n)\in\pow{X}$ by \cref{subseteq}.
                Thus $V(n)\subseteq X$ by \cref{powerset_elim}.
                $\unions{\img{\restrl{d}{\initialSegment{n}}}{\initialSegment{n}}}$ is closed in $X$ with respect to $T$.
                Thus $\unions{\img{\restrl{d}{\initialSegment{n}}}{\initialSegment{n}}}\subseteq X$ by \cref{closed_in}.
                Thus $W = V(n)\inter (X\setminus \unions{\img{\restrl{d}{\initialSegment{n}}}{\initialSegment{n}}})$
                    by \cref{setminus_eq_inter_complement}.
                $X\setminus \unions{\img{\restrl{d}{\initialSegment{n}}}{\initialSegment{n}}}$ is open in $X$ with respect to $T$
                    by \cref{closed_in}.
                Thus $X\setminus \unions{\img{\restrl{d}{\initialSegment{n}}}{\initialSegment{n}}}\in T$ by \cref{open_in}.
                Thus $W\in T$ by \cref{topology_on}.
                Thus $W$ is open in $X$ with respect to $T$ by \cref{open_in}.
            \end{subproof}
            Thus $G\subseteq T$ by \cref{subseteq}.
        \end{subproof}

        Let $U_1 = \unions{F}$.
        Let $V_1 = \unions{G}$.
        Show that $U_1$ is open in $X$ with respect to $T$.
        \begin{subproof}
            $T$ is closed under arbitrary unions by \cref{topology_on}.
            Thus $U_1\in T$ by \cref{topology_on,subseteq}.
            Thus $U_1$ is open in $X$ with respect to $T$ by \cref{open_in}.
        \end{subproof}
        Show that $V_1$ is open in $X$ with respect to $T$.
        \begin{subproof}
            $T$ is closed under arbitrary unions by \cref{topology_on}.
            Thus $V_1\in T$ by \cref{topology_on,subseteq}.
            Thus $V_1$ is open in $X$ with respect to $T$ by \cref{open_in}.
        \end{subproof}

        Show that $A\subseteq U_1$.
        \begin{subproof}
            Show that for all $x$ such that $x\in A$ we have $x\in U_1$.
            \begin{subproof}
                Fix $x$.
                Assume $x\in A$.
                Thus $x\in\unions{\img{U}{\naturalsPlus}}$ by \cref{subseteq}.
                Take $W$ such that $W\in \img{U}{\naturalsPlus}$ and $x\in W$ by \cref{unions_iff}.
                Take $n\in\naturalsPlus$ such that $n\mathrel{U} W$ by \cref{img_iff}.
                Thus $(n,W)\in U$ by \cref{pair_eq_iff}.
                $(n,U(n))\in U$ by \cref{function_apply_elim}.
                $U$ is right-unique by \cref{function_rightunique}.
                Thus $U(n) = W$ by \hyperref[rightunique]{right-uniqueness}.
                Thus $x\in U(n)$.
                Show that $x\in U(n)\setminus \unions{\img{\restrl{c}{\initialSegment{n}}}{\initialSegment{n}}}$.
                \begin{subproof}
                    Show that $x\notin \unions{\img{\restrl{c}{\initialSegment{n}}}{\initialSegment{n}}}$.
                    \begin{subproof}
                        Suppose not.
                        Then there exists $C$ such that $C\in \img{\restrl{c}{\initialSegment{n}}}{\initialSegment{n}}$ and $x\in C$ by \cref{unions_iff}.
                        Take $k\in\initialSegment{n}$ such that $k\mathrel{\restrl{c}{\initialSegment{n}}} C$ by \cref{img_iff}.
                        Thus $(k,C)\in \restrl{c}{\initialSegment{n}}$ by \cref{pair_eq_iff}.
                        Thus $(k,C)\in c$ by \cref{restrl_iff}.
                        Take $k_0\in\naturalsPlus$ such that $(k,C) = (k_0,\closure{V(k_0)}{X}{T})$.
                        Thus $C = \closure{V(k_0)}{X}{T}$ by \cref{pair_eq_iff}.
                        Thus $x\in\closure{V(k_0)}{X}{T}$.
                        $\closure{V(k_0)}{X}{T}$ is disjoint from $A$.
                        Thus there exists $y$ such that $y\in \closure{V(k_0)}{X}{T}$ and $y\in A$.
                        Contradiction by \cref{disjoint}.
                    \end{subproof}
                    Thus $x\in U(n)\setminus \unions{\img{\restrl{c}{\initialSegment{n}}}{\initialSegment{n}}}$ by \cref{setminus_intro}.
                \end{subproof}
                Thus $U(n)\setminus \unions{\img{\restrl{c}{\initialSegment{n}}}{\initialSegment{n}}}\in F$.
                Thus $x\in U_1$ by \cref{unions_intro}.
            \end{subproof}
            Thus $A\subseteq U_1$ by \cref{subseteq}.
        \end{subproof}

        Show that $B\subseteq V_1$.
        \begin{subproof}
            Show that for all $x$ such that $x\in B$ we have $x\in V_1$.
            \begin{subproof}
                Fix $x$.
                Assume $x\in B$.
                Thus $x\in\unions{\img{V}{\naturalsPlus}}$ by \cref{subseteq}.
                Take $W$ such that $W\in \img{V}{\naturalsPlus}$ and $x\in W$ by \cref{unions_iff}.
                Take $n\in\naturalsPlus$ such that $n\mathrel{V} W$ by \cref{img_iff}.
                Thus $(n,W)\in V$ by \cref{pair_eq_iff}.
                $(n,V(n))\in V$ by \cref{function_apply_elim}.
                $V$ is right-unique by \cref{function_rightunique}.
                Thus $V(n) = W$ by \hyperref[rightunique]{right-uniqueness}.
                Thus $x\in V(n)$.
                Show that $x\in V(n)\setminus \unions{\img{\restrl{d}{\initialSegment{n}}}{\initialSegment{n}}}$.
                \begin{subproof}
                    Show that $x\notin \unions{\img{\restrl{d}{\initialSegment{n}}}{\initialSegment{n}}}$.
                    \begin{subproof}
                        Suppose not.
                        Then there exists $C$ such that $C\in \img{\restrl{d}{\initialSegment{n}}}{\initialSegment{n}}$ and $x\in C$ by \cref{unions_iff}.
                        Take $k\in\initialSegment{n}$ such that $k\mathrel{\restrl{d}{\initialSegment{n}}} C$ by \cref{img_iff}.
                        Thus $(k,C)\in \restrl{d}{\initialSegment{n}}$ by \cref{pair_eq_iff}.
                        Thus $(k,C)\in d$ by \cref{restrl_iff}.
                        Take $k_0\in\naturalsPlus$ such that $(k,C) = (k_0,\closure{U(k_0)}{X}{T})$.
                        Thus $C = \closure{U(k_0)}{X}{T}$ by \cref{pair_eq_iff}.
                        Thus $x\in\closure{U(k_0)}{X}{T}$.
                        $\closure{U(k_0)}{X}{T}$ is disjoint from $B$.
                        Thus there exists $y$ such that $y\in \closure{U(k_0)}{X}{T}$ and $y\in B$.
                        Contradiction by \cref{disjoint}.
                    \end{subproof}
                    Thus $x\in V(n)\setminus \unions{\img{\restrl{d}{\initialSegment{n}}}{\initialSegment{n}}}$ by \cref{setminus_intro}.
                \end{subproof}
                Thus $V(n)\setminus \unions{\img{\restrl{d}{\initialSegment{n}}}{\initialSegment{n}}}\in G$.
                Thus $x\in V_1$ by \cref{unions_intro}.
            \end{subproof}
            Thus $B\subseteq V_1$ by \cref{subseteq}.
        \end{subproof}

        Show that $U_1$ is disjoint from $V_1$.
        \begin{subproof}
            Suppose not.
            Then there exists $x$ such that $x\in U_1$ and $x\in V_1$ by \cref{disjoint}.
            Take $W_1$ such that $W_1\in F$ and $x\in W_1$ by \cref{unions_iff}.
            Take $W_2$ such that $W_2\in G$ and $x\in W_2$ by \cref{unions_iff}.
            Take $j\in\naturalsPlus$ such that $W_1 = U(j)\setminus \unions{\img{\restrl{c}{\initialSegment{j}}}{\initialSegment{j}}}$.
            Take $k\in\naturalsPlus$ such that $W_2 = V(k)\setminus \unions{\img{\restrl{d}{\initialSegment{k}}}{\initialSegment{k}}}$.
            $x\in U(j)$ by \cref{setminus_elim_left}.
            $x\notin \unions{\img{\restrl{c}{\initialSegment{j}}}{\initialSegment{j}}}$ by \cref{setminus_elim_right}.
            $x\in V(k)$ by \cref{setminus_elim_left}.
            $x\notin \unions{\img{\restrl{d}{\initialSegment{k}}}{\initialSegment{k}}}$ by \cref{setminus_elim_right}.
            \begin{byCase}
            \caseOf{$j\leq k$.}
                $j\in\initialSegment{k}$ by \cref{initial_segment_naturalsplus}.
                Take $C$ such that $C = \closure{U(j)}{X}{T}$.
                $j\mathrel{\restrl{d}{\initialSegment{k}}} C$ by \cref{restrl_iff}.
                Thus $C\in \img{\restrl{d}{\initialSegment{k}}}{\initialSegment{k}}$ by \cref{img_iff}.
                $U\in\funs{\naturalsPlus}{\pow{X}}$.
                Thus $U$ is a function by \cref{funs_elim}.
                Thus $U(j)\in\pow{X}$ by \cref{funs_type_apply}.
                Thus $U(j)\subseteq X$ by \cref{powerset_elim}.
                Thus $U(j)\subseteq \closure{U(j)}{X}{T}$ by \cref{subseteq_closure}.
                Thus $x\in C$ by \cref{subseteq}.
                Thus $x\in \unions{\img{\restrl{d}{\initialSegment{k}}}{\initialSegment{k}}}$ by \cref{unions_intro}.
                Contradiction.
            \caseOf{$k\leq j$.}
                $k\in\initialSegment{j}$ by \cref{initial_segment_naturalsplus}.
                Take $C$ such that $C = \closure{V(k)}{X}{T}$.
                $k\mathrel{\restrl{c}{\initialSegment{j}}} C$ by \cref{restrl_iff}.
                Thus $C\in \img{\restrl{c}{\initialSegment{j}}}{\initialSegment{j}}$ by \cref{img_iff}.
                $V\in\funs{\naturalsPlus}{\pow{X}}$.
                Thus $V$ is a function by \cref{funs_elim}.
                Thus $V(k)\in\pow{X}$ by \cref{funs_type_apply}.
                Thus $V(k)\subseteq X$ by \cref{powerset_elim}.
                Thus $V(k)\subseteq \closure{V(k)}{X}{T}$ by \cref{subseteq_closure}.
                Thus $x\in C$ by \cref{subseteq}.
                Thus $x\in \unions{\img{\restrl{c}{\initialSegment{j}}}{\initialSegment{j}}}$ by \cref{unions_intro}.
                Contradiction.
            \end{byCase}
        \end{subproof}

        Thus there exist $U_1,V_1$ such that
            $U_1$ is open in $X$ with respect to $T$ and
            $V_1$ is open in $X$ with respect to $T$ and
            $A\subseteq U_1$ and
            $B\subseteq V_1$ and
            $U_1$ is disjoint from $V_1$.
    \end{subproof}
    $X$ has closed singletons with respect to $T$ by \cref{regular_space}.
    Thus $X$ is normal with respect to $T$ by \cref{normal_space}.
\end{proof}

% from §32 helper definition 1: metric ball family avoiding a set
\begin{definition}\label{metric_ball_family}
    $\metricBallFamily{d}{X}{A}{B} = \{W\in\pow{X} \mid \exists a\in A. \exists \epsilon\in\reals.
        0 < \epsilon \land W = \metricBall{d}{X}{a}{\epsilon / (1 + 1)} \land
        \metricBall{d}{X}{a}{\epsilon}\inter B = \emptyset\}$.
\end{definition}

% from §32 Item 2: metrizable space is normal
\begin{theorem}\label{metrizable_normal}
    Suppose $X$ is metrizable with respect to $T$.
    Then $X$ is normal with respect to $T$.
\end{theorem}
\begin{proof}
    Take $d$ such that $d$ is a metric on $X$ and $T = \metricTopology{d}{X}$ by \cref{metrizable}.
    $1\in\reals$ by \cref{real_naturals_subset_reals,real_one_in_naturals,subseteq}.
    $0\in\reals$ by \cref{real_add_ident}.
    Thus $1 + 1\in\reals$ by \cref{real_add_closed}.
    Show that $0 < 1 + 1$.
    \begin{subproof}
        $0 < 1$ by \cref{real_zero_lt_one}.
        $0 + 1 < 1 + 1$ by \cref{real_lt_add}.
        Thus $1 < 1 + 1$ by \cref{real_add_ident}.
        Thus $0 < 1 + 1$ by \cref{real_lt_trans}.
    \end{subproof}
    Thus $1 + 1\neq 0$ by \cref{real_pos_nonzero}.
    Thus $\inv{(1 + 1)}\in\reals$ by \cref{real_mul_inverse}.
    Show that for all $A,B$ such that $A$ is closed in $X$ with respect to $T$ and $B$ is closed in $X$ with respect to $T$
        and $A$ is disjoint from $B$
        there exist $U,V$ such that
            $U$ is open in $X$ with respect to $T$ and
            $V$ is open in $X$ with respect to $T$ and
            $A\subseteq U$ and
            $B\subseteq V$ and
            $U$ is disjoint from $V$.
    \begin{subproof}
        Fix $A$.
        Fix $B$.
        Assume $A$ is closed in $X$ with respect to $T$ and $B$ is closed in $X$ with respect to $T$
            and $A$ is disjoint from $B$.
        Let $F = \metricBallFamily{d}{X}{A}{B}$.
        Let $G = \metricBallFamily{d}{X}{B}{A}$.
        Let $U = \unions{F}$.
        Let $V = \unions{G}$.

        Show that $U$ is open in $X$ with respect to $T$.
        \begin{subproof}
            Show that $F\subseteq T$.
            \begin{subproof}
                Show that for all $W$ such that $W\in F$ we have $W\in T$.
                \begin{subproof}
                    Fix $W$.
                    Assume $W\in F$.
                    Take $a,\epsilon$ such that
                        $a\in A$ and $\epsilon\in\reals$ and $0 < \epsilon$ and
                        $W = \metricBall{d}{X}{a}{\epsilon / (1 + 1)}$ and
                        $\metricBall{d}{X}{a}{\epsilon}\inter B = \emptyset$
                        by \cref{metric_ball_family}.
                    $a\in X$ by \cref{closed_in,subseteq}.
                    $\epsilon / (1 + 1)\in\reals$ by \cref{real_div,real_mul_closed}.
                    Show that $0 < \epsilon / (1 + 1)$.
                    \begin{subproof}
                        $0 \rmul \inv{(1 + 1)} < \epsilon \rmul \inv{(1 + 1)}$ by \cref{real_lt_mul_pos,real_inv_pos}.
                        $0 \rmul \inv{(1 + 1)} = 0$ by \cref{real_mul_zero}.
                        Thus $0 < \epsilon \rmul \inv{(1 + 1)}$ by \cref{real_lt_subst_left}.
                        Thus $0 < \epsilon / (1 + 1)$ by \cref{real_div}.
                    \end{subproof}
                    Show that for all $y$ such that $y\in\metricBall{d}{X}{a}{\epsilon / (1 + 1)}$ we have $y\in X$.
                    \begin{subproof}
                        Fix $y$.
                        Assume $y\in\metricBall{d}{X}{a}{\epsilon / (1 + 1)}$.
                        Thus $y\in X$ by \cref{metric_ball}.
                    \end{subproof}
                    Thus $\metricBall{d}{X}{a}{\epsilon / (1 + 1)}\subseteq X$ by \cref{subseteq}.
                    Thus $\metricBall{d}{X}{a}{\epsilon / (1 + 1)}\in\pow{X}$ by \cref{powerset_intro}.
                    $\metricBall{d}{X}{a}{\epsilon / (1 + 1)}\in\metricBasis{d}{X}$ by \cref{metric_basis}.
                    $\metricBasis{d}{X}$ is a basis on $X$ by \cref{metric_basis_is_basis}.
                    Thus $W\in T$ by \cref{basis_elem_open_in_generated,metric_topology,basis_for_topology}.
                \end{subproof}
                Thus $F\subseteq T$ by \cref{subseteq}.
            \end{subproof}
            $T$ is a topology on $X$ by \cref{metrizable}.
            Thus $U\in T$ by \cref{topology_on,subseteq}.
            Thus $U$ is open in $X$ with respect to $T$ by \cref{open_in}.
        \end{subproof}

        Show that $V$ is open in $X$ with respect to $T$.
        \begin{subproof}
            Show that $G\subseteq T$.
            \begin{subproof}
                Show that for all $W$ such that $W\in G$ we have $W\in T$.
                \begin{subproof}
                    Fix $W$.
                    Assume $W\in G$.
                    Take $b,\epsilon$ such that
                        $b\in B$ and $\epsilon\in\reals$ and $0 < \epsilon$ and
                        $W = \metricBall{d}{X}{b}{\epsilon / (1 + 1)}$ and
                        $\metricBall{d}{X}{b}{\epsilon}\inter A = \emptyset$
                        by \cref{metric_ball_family}.
                    $b\in X$ by \cref{closed_in,subseteq}.
                    $\epsilon / (1 + 1)\in\reals$ by \cref{real_div,real_mul_closed}.
                    Show that $0 < \epsilon / (1 + 1)$.
                    \begin{subproof}
                        $0 \rmul \inv{(1 + 1)} < \epsilon \rmul \inv{(1 + 1)}$ by \cref{real_lt_mul_pos,real_inv_pos}.
                        $0 \rmul \inv{(1 + 1)} = 0$ by \cref{real_mul_zero}.
                        Thus $0 < \epsilon \rmul \inv{(1 + 1)}$ by \cref{real_lt_subst_left}.
                        Thus $0 < \epsilon / (1 + 1)$ by \cref{real_div}.
                    \end{subproof}
                    Show that for all $y$ such that $y\in\metricBall{d}{X}{b}{\epsilon / (1 + 1)}$ we have $y\in X$.
                    \begin{subproof}
                        Fix $y$.
                        Assume $y\in\metricBall{d}{X}{b}{\epsilon / (1 + 1)}$.
                        Thus $y\in X$ by \cref{metric_ball}.
                    \end{subproof}
                    Thus $\metricBall{d}{X}{b}{\epsilon / (1 + 1)}\subseteq X$ by \cref{subseteq}.
                    Thus $\metricBall{d}{X}{b}{\epsilon / (1 + 1)}\in\pow{X}$ by \cref{powerset_intro}.
                    $\metricBall{d}{X}{b}{\epsilon / (1 + 1)}\in\metricBasis{d}{X}$ by \cref{metric_basis}.
                    $\metricBasis{d}{X}$ is a basis on $X$ by \cref{metric_basis_is_basis}.
                    Thus $W\in T$ by \cref{basis_elem_open_in_generated,metric_topology,basis_for_topology}.
                \end{subproof}
                Thus $G\subseteq T$ by \cref{subseteq}.
            \end{subproof}
            $T$ is a topology on $X$ by \cref{metrizable}.
            Thus $V\in T$ by \cref{topology_on,subseteq}.
            Thus $V$ is open in $X$ with respect to $T$ by \cref{open_in}.
        \end{subproof}

        Show that $A\subseteq U$.
        \begin{subproof}
            Show that for all $a$ such that $a\in A$ we have $a\in U$.
            \begin{subproof}
                Fix $a$.
                Assume $a\in A$.
                $A\subseteq X$ by \cref{closed_in}.
                Thus $a\in X$ by \cref{subseteq}.
                $B\subseteq X$ by \cref{closed_in}.
                Thus $X\setminus B$ is open in $X$ with respect to $T$ by \cref{closed_in}.
                $X\setminus B\subseteq X$ by \cref{setminus_subseteq}.
                $a\in X\setminus B$ by \cref{setminus_intro,disjoint}.
                There exists $\epsilon\in\reals$ such that $0 < \epsilon$ and
                    $\metricBall{d}{X}{a}{\epsilon}\subseteq X\setminus B$ by \cref{metric_open_iff}.
                Take $\epsilon\in\reals$ such that $0 < \epsilon$ and
                    $\metricBall{d}{X}{a}{\epsilon}\subseteq X\setminus B$.
                Show that $\metricBall{d}{X}{a}{\epsilon}\inter B = \emptyset$.
                \begin{subproof}
                    Suppose not.
                    Then $\metricBall{d}{X}{a}{\epsilon}\inter B\neq\emptyset$.
                    Take $x$ such that $x\in\metricBall{d}{X}{a}{\epsilon}\inter B$ by \cref{nonempty_inhabited}.
                    Thus $x\in\metricBall{d}{X}{a}{\epsilon}$ and $x\in B$ by \cref{inter}.
                    Thus $x\in X\setminus B$ by \cref{subseteq}.
                    Thus $x\notin B$ by \cref{setminus_elim_right}.
                    Contradiction.
                \end{subproof}
                Show that $0 < \epsilon / (1 + 1)$.
                \begin{subproof}
                    $0 \rmul \inv{(1 + 1)} < \epsilon \rmul \inv{(1 + 1)}$ by \cref{real_lt_mul_pos,real_inv_pos}.
                    $0 \rmul \inv{(1 + 1)} = 0$ by \cref{real_mul_zero}.
                    Thus $0 < \epsilon \rmul \inv{(1 + 1)}$ by \cref{real_lt_subst_left}.
                    Thus $0 < \epsilon / (1 + 1)$ by \cref{real_div}.
                \end{subproof}
                $d$ is zero-characterized on $X$ by \cref{metric_on}.
                Thus $d((a,a)) = 0$ by \cref{metric_zero_characterization}.
                Thus $d((a,a)) < \epsilon / (1 + 1)$ by \cref{real_lt_subst_left}.
                Thus $a\in\metricBall{d}{X}{a}{\epsilon / (1 + 1)}$ by \cref{metric_ball}.
                Show that for all $y$ such that $y\in\metricBall{d}{X}{a}{\epsilon / (1 + 1)}$ we have $y\in X$.
                \begin{subproof}
                    Fix $y$.
                    Assume $y\in\metricBall{d}{X}{a}{\epsilon / (1 + 1)}$.
                    Thus $y\in X$ by \cref{metric_ball}.
                \end{subproof}
                Thus $\metricBall{d}{X}{a}{\epsilon / (1 + 1)}\subseteq X$ by \cref{subseteq}.
                Thus $\metricBall{d}{X}{a}{\epsilon / (1 + 1)}\in\pow{X}$ by \cref{powerset_intro}.
                Thus $\metricBall{d}{X}{a}{\epsilon / (1 + 1)}\in F$.
                Thus $a\in U$ by \cref{unions_intro}.
            \end{subproof}
            Thus $A\subseteq U$ by \cref{subseteq}.
        \end{subproof}

        Show that $B\subseteq V$.
        \begin{subproof}
            Show that for all $b$ such that $b\in B$ we have $b\in V$.
            \begin{subproof}
                Fix $b$.
                Assume $b\in B$.
                $B\subseteq X$ by \cref{closed_in}.
                Thus $b\in X$ by \cref{subseteq}.
                $A\subseteq X$ by \cref{closed_in}.
                Thus $X\setminus A$ is open in $X$ with respect to $T$ by \cref{closed_in}.
                $X\setminus A\subseteq X$ by \cref{setminus_subseteq}.
                $b\in X\setminus A$ by \cref{setminus_intro,disjoint}.
                There exists $\epsilon\in\reals$ such that $0 < \epsilon$ and
                    $\metricBall{d}{X}{b}{\epsilon}\subseteq X\setminus A$ by \cref{metric_open_iff}.
                Take $\epsilon\in\reals$ such that $0 < \epsilon$ and
                    $\metricBall{d}{X}{b}{\epsilon}\subseteq X\setminus A$.
                Show that $\metricBall{d}{X}{b}{\epsilon}\inter A = \emptyset$.
                \begin{subproof}
                    Suppose not.
                    Then $\metricBall{d}{X}{b}{\epsilon}\inter A\neq\emptyset$.
                    Take $x$ such that $x\in\metricBall{d}{X}{b}{\epsilon}\inter A$ by \cref{nonempty_inhabited}.
                    Thus $x\in\metricBall{d}{X}{b}{\epsilon}$ and $x\in A$ by \cref{inter}.
                    Thus $x\in X\setminus A$ by \cref{subseteq}.
                    Thus $x\notin A$ by \cref{setminus_elim_right}.
                    Contradiction.
                \end{subproof}
                Show that $0 < \epsilon / (1 + 1)$.
                \begin{subproof}
                    $0 \rmul \inv{(1 + 1)} < \epsilon \rmul \inv{(1 + 1)}$ by \cref{real_lt_mul_pos,real_inv_pos}.
                    $0 \rmul \inv{(1 + 1)} = 0$ by \cref{real_mul_zero}.
                    Thus $0 < \epsilon \rmul \inv{(1 + 1)}$ by \cref{real_lt_subst_left}.
                    Thus $0 < \epsilon / (1 + 1)$ by \cref{real_div}.
                \end{subproof}
                $d$ is zero-characterized on $X$ by \cref{metric_on}.
                Thus $d((b,b)) = 0$ by \cref{metric_zero_characterization}.
                Thus $d((b,b)) < \epsilon / (1 + 1)$ by \cref{real_lt_subst_left}.
                Thus $b\in\metricBall{d}{X}{b}{\epsilon / (1 + 1)}$ by \cref{metric_ball}.
                Show that for all $y$ such that $y\in\metricBall{d}{X}{b}{\epsilon / (1 + 1)}$ we have $y\in X$.
                \begin{subproof}
                    Fix $y$.
                    Assume $y\in\metricBall{d}{X}{b}{\epsilon / (1 + 1)}$.
                    Thus $y\in X$ by \cref{metric_ball}.
                \end{subproof}
                Thus $\metricBall{d}{X}{b}{\epsilon / (1 + 1)}\subseteq X$ by \cref{subseteq}.
                Thus $\metricBall{d}{X}{b}{\epsilon / (1 + 1)}\in\pow{X}$ by \cref{powerset_intro}.
                Thus $\metricBall{d}{X}{b}{\epsilon / (1 + 1)}\in G$.
                Thus $b\in V$ by \cref{unions_intro}.
            \end{subproof}
            Thus $B\subseteq V$ by \cref{subseteq}.
        \end{subproof}

        Show that $U$ is disjoint from $V$.
        \begin{subproof}
            Suppose not.
            Then there exists $z$ such that $z\in U$ and $z\in V$ by \cref{disjoint}.
            Take $W_1$ such that $W_1\in F$ and $z\in W_1$ by \cref{unions_iff}.
            Take $W_2$ such that $W_2\in G$ and $z\in W_2$ by \cref{unions_iff}.
            Take $a,\epsilon_a$ such that
                $a\in A$ and $\epsilon_a\in\reals$ and $0 < \epsilon_a$ and
                $W_1 = \metricBall{d}{X}{a}{\epsilon_a / (1 + 1)}$ and
                $\metricBall{d}{X}{a}{\epsilon_a}\inter B = \emptyset$
                by \cref{metric_ball_family}.
            Take $b,\epsilon_b$ such that
                $b\in B$ and $\epsilon_b\in\reals$ and $0 < \epsilon_b$ and
                $W_2 = \metricBall{d}{X}{b}{\epsilon_b / (1 + 1)}$ and
                $\metricBall{d}{X}{b}{\epsilon_b}\inter A = \emptyset$
                by \cref{metric_ball_family}.
            $z\in\metricBall{d}{X}{a}{\epsilon_a / (1 + 1)}$ and $z\in\metricBall{d}{X}{b}{\epsilon_b / (1 + 1)}$.
            Thus $d((a,z)) < \epsilon_a / (1 + 1)$ and $d((b,z)) < \epsilon_b / (1 + 1)$ by \cref{metric_ball}.
            Show that $d((a,b)) < \epsilon_a / (1 + 1) + \epsilon_b / (1 + 1)$.
            \begin{subproof}
                $d$ satisfies the triangle inequality on $X$ by \cref{metric_on}.
                $d$ is symmetric on $X$ by \cref{metric_symmetric,metric_on}.
                Thus $z\in X$ by \cref{metric_ball}.
                $A\subseteq X$ by \cref{closed_in}.
                Thus $a\in X$ by \cref{subseteq}.
                $B\subseteq X$ by \cref{closed_in}.
                Thus $b\in X$ by \cref{subseteq}.
                Thus $d((z,b)) = d((b,z))$ by \cref{metric_symmetric}.
                $(a,z)\in X\times X$ and $(z,b)\in X\times X$ by \cref{times_tuple_intro}.
                Thus $d((a,z))\in\reals$ and $d((z,b))\in\reals$ by \cref{funs_type_apply,metric_on}.
                $\epsilon_a / (1 + 1)\in\reals$ and $\epsilon_b / (1 + 1)\in\reals$ by \cref{real_div,real_mul_closed}.
                $(a,b)\in X\times X$ by \cref{times_tuple_intro}.
                Thus $d((a,b))\in\reals$ by \cref{funs_type_apply,metric_on}.
                Thus $d((a,z)) + d((z,b))\in\reals$ by \cref{real_add_closed}.
                Thus $\epsilon_a / (1 + 1) + \epsilon_b / (1 + 1)\in\reals$ by \cref{real_add_closed}.
                Thus $d((a,b)) < d((a,z)) + d((z,b))$ or $d((a,b)) = d((a,z)) + d((z,b))$
                    by \cref{metric_triangle_inequality}.
                Thus $d((z,b)) < \epsilon_b / (1 + 1)$ by \cref{real_lt_subst_left}.
                Thus $d((a,z)) + d((z,b)) < \epsilon_a / (1 + 1) + \epsilon_b / (1 + 1)$ by \cref{real_lt_add_two}.
                \begin{byCase}
                \caseOf{$d((a,b)) < d((a,z)) + d((z,b))$.}
                    Thus $d((a,b)) < \epsilon_a / (1 + 1) + \epsilon_b / (1 + 1)$ by \cref{real_lt_trans}.
                \caseOf{$d((a,b)) = d((a,z)) + d((z,b))$.}
                    Thus $d((a,b)) < \epsilon_a / (1 + 1) + \epsilon_b / (1 + 1)$ by \cref{real_lt_subst_left}.
                \end{byCase}
            \end{subproof}
            $A\subseteq X$ by \cref{closed_in}.
            Thus $a\in X$ by \cref{subseteq}.
            $B\subseteq X$ by \cref{closed_in}.
            Thus $b\in X$ by \cref{subseteq}.
            $(a,b)\in X\times X$ by \cref{times_tuple_intro}.
            Thus $d((a,b))\in\reals$ by \cref{funs_type_apply,metric_on}.
            \begin{byCase}
            \caseOf{$\epsilon_a < \epsilon_b$.}
                $\epsilon_a / (1 + 1)\in\reals$ and $\epsilon_b / (1 + 1)\in\reals$ by \cref{real_div,real_mul_closed}.
                $\epsilon_a / (1 + 1) < \epsilon_b / (1 + 1)$ by \cref{real_lt_mul_pos,real_inv_pos,real_div}.
                Thus $\epsilon_a / (1 + 1) + \epsilon_b / (1 + 1) < \epsilon_b / (1 + 1) + \epsilon_b / (1 + 1)$ by \cref{real_lt_add}.
                Thus $\epsilon_a / (1 + 1) + \epsilon_b / (1 + 1)\in\reals$ by \cref{real_add_closed}.
                Thus $\epsilon_b / (1 + 1) + \epsilon_b / (1 + 1)\in\reals$ by \cref{real_add_closed}.
                Show that $\epsilon_b / (1 + 1) + \epsilon_b / (1 + 1) = \epsilon_b$.
                \begin{subproof}
                    $\epsilon_b / (1 + 1) = \epsilon_b \rmul \inv{(1 + 1)}$ by \cref{real_div}.
                    Thus $\epsilon_b / (1 + 1) + \epsilon_b / (1 + 1) =
                        (\epsilon_b \rmul \inv{(1 + 1)}) + (\epsilon_b \rmul \inv{(1 + 1)})$.
                    $(\epsilon_b + \epsilon_b) \rmul \inv{(1 + 1)} =
                        (\epsilon_b \rmul \inv{(1 + 1)}) + (\epsilon_b \rmul \inv{(1 + 1)})$ by \cref{real_right_distrib}.
                    Thus $\epsilon_b / (1 + 1) + \epsilon_b / (1 + 1) =
                        (\epsilon_b + \epsilon_b) \rmul \inv{(1 + 1)}$.
                    $\epsilon_b \rmul (1 + 1) = (1 + 1) \rmul \epsilon_b$ by \cref{real_mul_comm}.
                    $(1 + 1) \rmul \epsilon_b = (1 \rmul \epsilon_b) + (1 \rmul \epsilon_b)$ by \cref{real_right_distrib}.
                    Thus $\epsilon_b \rmul (1 + 1) = (1 \rmul \epsilon_b) + (1 \rmul \epsilon_b)$.
                    Thus $\epsilon_b \rmul (1 + 1) = \epsilon_b + \epsilon_b$ by \cref{real_mul_ident}.
                    Thus $\epsilon_b + \epsilon_b = \epsilon_b \rmul (1 + 1)$.
                    Thus $(\epsilon_b + \epsilon_b) \rmul \inv{(1 + 1)} =
                        (\epsilon_b \rmul (1 + 1)) \rmul \inv{(1 + 1)}$.
                    Thus $(\epsilon_b + \epsilon_b) \rmul \inv{(1 + 1)} =
                        \epsilon_b \rmul ((1 + 1) \rmul \inv{(1 + 1)})$ by \cref{real_mul_assoc}.
                    Thus $(\epsilon_b + \epsilon_b) \rmul \inv{(1 + 1)} = \epsilon_b \rmul 1$ by \cref{real_mul_inverse}.
                    $\epsilon_b \rmul 1 = 1 \rmul \epsilon_b$ by \cref{real_mul_comm}.
                    Thus $\epsilon_b \rmul 1 = \epsilon_b$ by \cref{real_mul_ident}.
                    Thus $\epsilon_b / (1 + 1) + \epsilon_b / (1 + 1) = \epsilon_b$.
                \end{subproof}
                Thus $d((a,b)) < \epsilon_b / (1 + 1) + \epsilon_b / (1 + 1)$ by \cref{real_lt_trans}.
                Thus $d((a,b)) < \epsilon_b$ by \cref{real_lt_subst_right}.
                $d$ is symmetric on $X$ by \cref{metric_symmetric,metric_on}.
                Thus $d((b,a)) = d((a,b))$ by \cref{metric_symmetric}.
                Thus $d((b,a)) < \epsilon_b$ by \cref{real_lt_subst_left}.
                Thus $a\in\metricBall{d}{X}{b}{\epsilon_b}$ by \cref{metric_ball}.
                Thus $a\in\metricBall{d}{X}{b}{\epsilon_b}\inter A$ by \cref{inter_intro}.
                Thus $\metricBall{d}{X}{b}{\epsilon_b}\inter A\subseteq\emptyset$ by \cref{subseteq_emptyset_iff}.
                Thus $a\in\emptyset$ by \cref{subseteq}.
                Contradiction by \cref{emptyset}.
            \caseOf{$\epsilon_a = \epsilon_b$.}
                $\epsilon_a / (1 + 1)\in\reals$ and $\epsilon_b / (1 + 1)\in\reals$ by \cref{real_div,real_mul_closed}.
                Thus $\epsilon_a / (1 + 1) + \epsilon_b / (1 + 1)\in\reals$ by \cref{real_add_closed}.
                Show that $\epsilon_a / (1 + 1) + \epsilon_b / (1 + 1) = \epsilon_a / (1 + 1) + \epsilon_a / (1 + 1)$.
                \begin{subproof}
                    Thus $\epsilon_b = \epsilon_a$.
                    Thus $\epsilon_b / (1 + 1) = \epsilon_a / (1 + 1)$.
                    Thus $\epsilon_a / (1 + 1) + \epsilon_b / (1 + 1) = \epsilon_a / (1 + 1) + \epsilon_a / (1 + 1)$ by \cref{real_add_eq_subst}.
                \end{subproof}
                Show that $\epsilon_a / (1 + 1) + \epsilon_a / (1 + 1) = \epsilon_a$.
                \begin{subproof}
                    $\epsilon_a / (1 + 1) = \epsilon_a \rmul \inv{(1 + 1)}$ by \cref{real_div}.
                    Thus $\epsilon_a / (1 + 1) + \epsilon_a / (1 + 1) =
                        (\epsilon_a \rmul \inv{(1 + 1)}) + (\epsilon_a \rmul \inv{(1 + 1)})$.
                    $(\epsilon_a + \epsilon_a) \rmul \inv{(1 + 1)} =
                        (\epsilon_a \rmul \inv{(1 + 1)}) + (\epsilon_a \rmul \inv{(1 + 1)})$ by \cref{real_right_distrib}.
                    Thus $\epsilon_a / (1 + 1) + \epsilon_a / (1 + 1) =
                        (\epsilon_a + \epsilon_a) \rmul \inv{(1 + 1)}$.
                    $\epsilon_a \rmul (1 + 1) = (1 + 1) \rmul \epsilon_a$ by \cref{real_mul_comm}.
                    $(1 + 1) \rmul \epsilon_a = (1 \rmul \epsilon_a) + (1 \rmul \epsilon_a)$ by \cref{real_right_distrib}.
                    Thus $\epsilon_a \rmul (1 + 1) = (1 \rmul \epsilon_a) + (1 \rmul \epsilon_a)$.
                    Thus $\epsilon_a \rmul (1 + 1) = \epsilon_a + \epsilon_a$ by \cref{real_mul_ident}.
                    Thus $\epsilon_a + \epsilon_a = \epsilon_a \rmul (1 + 1)$.
                    Thus $(\epsilon_a + \epsilon_a) \rmul \inv{(1 + 1)} =
                        (\epsilon_a \rmul (1 + 1)) \rmul \inv{(1 + 1)}$.
                    Thus $(\epsilon_a + \epsilon_a) \rmul \inv{(1 + 1)} =
                        \epsilon_a \rmul ((1 + 1) \rmul \inv{(1 + 1)})$ by \cref{real_mul_assoc}.
                    Thus $(\epsilon_a + \epsilon_a) \rmul \inv{(1 + 1)} = \epsilon_a \rmul 1$ by \cref{real_mul_inverse}.
                    $\epsilon_a \rmul 1 = 1 \rmul \epsilon_a$ by \cref{real_mul_comm}.
                    Thus $\epsilon_a \rmul 1 = \epsilon_a$ by \cref{real_mul_ident}.
                    Thus $\epsilon_a / (1 + 1) + \epsilon_a / (1 + 1) = \epsilon_a$.
                \end{subproof}
                Thus $d((a,b)) < \epsilon_a / (1 + 1) + \epsilon_a / (1 + 1)$ by \cref{real_lt_subst_right}.
                Thus $d((a,b)) < \epsilon_a$ by \cref{real_lt_subst_right}.
                Thus $b\in\metricBall{d}{X}{a}{\epsilon_a}$ by \cref{metric_ball}.
                Thus $b\in\metricBall{d}{X}{a}{\epsilon_a}\inter B$ by \cref{inter_intro}.
                Thus $\metricBall{d}{X}{a}{\epsilon_a}\inter B\subseteq\emptyset$ by \cref{subseteq_emptyset_iff}.
                Thus $b\in\emptyset$ by \cref{subseteq}.
                Contradiction by \cref{emptyset}.
            \caseOf{$\epsilon_b < \epsilon_a$.}
                $\epsilon_a / (1 + 1)\in\reals$ and $\epsilon_b / (1 + 1)\in\reals$ by \cref{real_div,real_mul_closed}.
                $\epsilon_b / (1 + 1) < \epsilon_a / (1 + 1)$ by \cref{real_lt_mul_pos,real_inv_pos,real_div}.
                Thus $\epsilon_b / (1 + 1) + \epsilon_a / (1 + 1) < \epsilon_a / (1 + 1) + \epsilon_a / (1 + 1)$ by \cref{real_lt_add}.
                $\epsilon_a / (1 + 1) + \epsilon_b / (1 + 1) = \epsilon_b / (1 + 1) + \epsilon_a / (1 + 1)$ by \cref{real_add_comm}.
                Thus $\epsilon_a / (1 + 1) + \epsilon_b / (1 + 1) < \epsilon_a / (1 + 1) + \epsilon_a / (1 + 1)$ by \cref{real_lt_subst_left}.
                Thus $\epsilon_a / (1 + 1) + \epsilon_b / (1 + 1)\in\reals$ by \cref{real_add_closed}.
                Thus $\epsilon_a / (1 + 1) + \epsilon_a / (1 + 1)\in\reals$ by \cref{real_add_closed}.
                Show that $\epsilon_a / (1 + 1) + \epsilon_a / (1 + 1) = \epsilon_a$.
                \begin{subproof}
                    $\epsilon_a / (1 + 1) = \epsilon_a \rmul \inv{(1 + 1)}$ by \cref{real_div}.
                    Thus $\epsilon_a / (1 + 1) + \epsilon_a / (1 + 1) =
                        (\epsilon_a \rmul \inv{(1 + 1)}) + (\epsilon_a \rmul \inv{(1 + 1)})$.
                    $(\epsilon_a + \epsilon_a) \rmul \inv{(1 + 1)} =
                        (\epsilon_a \rmul \inv{(1 + 1)}) + (\epsilon_a \rmul \inv{(1 + 1)})$ by \cref{real_right_distrib}.
                    Thus $\epsilon_a / (1 + 1) + \epsilon_a / (1 + 1) =
                        (\epsilon_a + \epsilon_a) \rmul \inv{(1 + 1)}$.
                    $\epsilon_a \rmul (1 + 1) = (1 + 1) \rmul \epsilon_a$ by \cref{real_mul_comm}.
                    $(1 + 1) \rmul \epsilon_a = (1 \rmul \epsilon_a) + (1 \rmul \epsilon_a)$ by \cref{real_right_distrib}.
                    Thus $\epsilon_a \rmul (1 + 1) = (1 \rmul \epsilon_a) + (1 \rmul \epsilon_a)$.
                    Thus $\epsilon_a \rmul (1 + 1) = \epsilon_a + \epsilon_a$ by \cref{real_mul_ident}.
                    Thus $\epsilon_a + \epsilon_a = \epsilon_a \rmul (1 + 1)$.
                    Thus $(\epsilon_a + \epsilon_a) \rmul \inv{(1 + 1)} =
                        (\epsilon_a \rmul (1 + 1)) \rmul \inv{(1 + 1)}$.
                    Thus $(\epsilon_a + \epsilon_a) \rmul \inv{(1 + 1)} =
                        \epsilon_a \rmul ((1 + 1) \rmul \inv{(1 + 1)})$ by \cref{real_mul_assoc}.
                    Thus $(\epsilon_a + \epsilon_a) \rmul \inv{(1 + 1)} = \epsilon_a \rmul 1$ by \cref{real_mul_inverse}.
                    $\epsilon_a \rmul 1 = 1 \rmul \epsilon_a$ by \cref{real_mul_comm}.
                    Thus $\epsilon_a \rmul 1 = \epsilon_a$ by \cref{real_mul_ident}.
                    Thus $\epsilon_a / (1 + 1) + \epsilon_a / (1 + 1) = \epsilon_a$.
                \end{subproof}
                Thus $d((a,b)) < \epsilon_a / (1 + 1) + \epsilon_a / (1 + 1)$ by \cref{real_lt_trans}.
                Thus $d((a,b)) < \epsilon_a$ by \cref{real_lt_subst_right}.
                Thus $b\in\metricBall{d}{X}{a}{\epsilon_a}$ by \cref{metric_ball}.
                Thus $b\in\metricBall{d}{X}{a}{\epsilon_a}\inter B$ by \cref{inter_intro}.
                Thus $\metricBall{d}{X}{a}{\epsilon_a}\inter B\subseteq\emptyset$ by \cref{subseteq_emptyset_iff}.
                Thus $b\in\emptyset$ by \cref{subseteq}.
                Contradiction by \cref{emptyset}.
            \end{byCase}
        \end{subproof}
    \end{subproof}
    $T$ is a topology on $X$ by \cref{metrizable}.
    $X$ is Hausdorff with respect to $T$ by \cref{metric_space_hausdorff}.
    Show that for all $x\in X$ we have $\{x\}$ is closed in $X$ with respect to $T$.
    \begin{subproof}
        Fix $x\in X$.
        Thus $\{x\}\subseteq X$ by \cref{singleton_subset_intro}.
        Show that $\{x\}$ is finite.
        \begin{subproof}
            $\{x,x\}$ is finite by \cref{pair_is_finite}.
            Show that for all $y$ such that $y\in\{x\}$ we have $y\in\{x,x\}$.
            \begin{subproof}
                Fix $y$.
                Assume $y\in\{x\}$.
                Thus $y = x$ by \cref{singleton_iff}.
                Thus $y\in\{x,x\}$ by \cref{upair_intro_left}.
            \end{subproof}
            Thus $\{x\}\subseteq\{x,x\}$ by \cref{subseteq}.
            Thus $\{x\}$ is finite by \cref{subseteq_finite_is_finite}.
        \end{subproof}
        Thus $\{x\}$ is closed in $X$ with respect to $T$ by \cref{finite_point_set_closed_hausdorff}.
        \end{subproof}
        Thus $X$ has closed singletons with respect to $T$ by \cref{one_point_closed}.
        Thus $X$ is normal with respect to $T$ by \cref{normal_space}.
\end{proof}

% from §32 Item 3: compact Hausdorff space is normal
\begin{theorem}\label{compact_hausdorff_normal}
    Suppose $X$ is compact with respect to $T$.
    Suppose $X$ is Hausdorff with respect to $T$.
    Then $X$ is normal with respect to $T$.
\end{theorem}
\begin{proof}
    $T$ is a topology on $X$ by \cref{hausdorff}.
    Show that for all $A,B$ such that $A$ is closed in $X$ with respect to $T$ and $B$ is closed in $X$ with respect to $T$
        and $A$ is disjoint from $B$
        there exist $U,V$ such that
            $U$ is open in $X$ with respect to $T$ and
            $V$ is open in $X$ with respect to $T$ and
            $A\subseteq U$ and
            $B\subseteq V$ and
            $U$ is disjoint from $V$.
    \begin{subproof}
        Fix $A$.
        Fix $B$.
        Assume $A$ is closed in $X$ with respect to $T$ and $B$ is closed in $X$ with respect to $T$
            and $A$ is disjoint from $B$.
        \begin{byCase}
        \caseOf{$A = \emptyset$.}
            Let $U = \emptyset$.
            Let $V = X$.
            $U\in T$ by \cref{topology_on}.
            $U$ is open in $X$ with respect to $T$ by \cref{open_in}.
            $V\in T$ by \cref{topology_on}.
            $V$ is open in $X$ with respect to $T$ by \cref{open_in}.
            $A\subseteq U$ by \cref{emptyset_subseteq}.
            $B\subseteq V$ by \cref{closed_in,subseteq}.
            Show that $U$ is disjoint from $V$.
            \begin{subproof}
                Suppose not.
                Then there exists $z$ such that $z\in U$ and $z\in V$ by \cref{disjoint}.
                Thus $z\in\emptyset$.
                Contradiction by \cref{emptyset}.
            \end{subproof}
        \caseOf{$A\neq\emptyset$.}
            $A\subseteq X$ by \cref{closed_in}.
            $B\subseteq X$ by \cref{closed_in}.
            Let $T_A = \subspaceTopology{T}{A}$.
            Let $T_B = \subspaceTopology{T}{B}$.
            Thus $A$ is compact with respect to $T_A$ by \cref{closed_subspace_compact}.
            Thus $B$ is compact with respect to $T_B$ by \cref{closed_subspace_compact}.
            Let $F = \{U\in T \mid \text{there exists $a\in A$ such that there exists $V$ such that
                    $V$ is open in $X$ with respect to $T$ and
                    $a\in U$ and
                    $B\subseteq V$ and
                    $U$ is disjoint from $V$}\}$.
            Show that $F$ is a covering of $A$.
            \begin{subproof}
                Show that for all $a$ such that $a\in A$ we have $a\in\unions{F}$.
                \begin{subproof}
                    Fix $a$.
                    Assume $a\in A$.
                    Thus $a\in X$ by \cref{subseteq}.
                    $a\notin B$ by \cref{disjoint}.
                    Thus $a\in X\setminus B$ by \cref{setminus_intro}.
                    Take $U,V$ such that
                        $U$ is open in $X$ with respect to $T$ and
                        $V$ is open in $X$ with respect to $T$ and
                        $a\in U$ and
                        $B\subseteq V$ and
                        $U$ is disjoint from $V$
                        by \cref{compact_hausdorff_separation}.
                    $U\in T$ by \cref{open_in}.
                    Thus $U\in F$.
                    Thus $a\in\unions{F}$ by \cref{unions_intro}.
                \end{subproof}
                Thus $A\subseteq\unions{F}$ by \cref{subseteq}.
                Thus $F$ is a covering of $A$ by \cref{covering}.
            \end{subproof}
            Show that for all $U$ such that $U\in F$ we have $U$ is open in $X$ with respect to $T$.
            \begin{subproof}
                Fix $U$.
                Assume $U\in F$.
                Thus $U\in T$.
                Thus $U$ is open in $X$ with respect to $T$ by \cref{open_in}.
            \end{subproof}
            $A$ is a subspace of $X$ with respect to $T$ by \cref{subspace_of}.
            Thus there exists $G$ such that
                $G\subseteq F$ and $G$ is finite and $G$ is a covering of $A$
                by \cref{compact_subspace_open_covering,subseteq}.
            Take $G$ such that $G\subseteq F$ and $G$ is finite and $G$ is a covering of $A$.
            Show that $G$ is inhabited.
            \begin{subproof}
                Take $a$ such that $a\in A$ by \cref{nonempty_inhabited}.
                $A\subseteq\unions{G}$ by \cref{covering}.
                Thus $a\in\unions{G}$ by \cref{subseteq}.
                Take $U\in G$ such that $a\in U$ by \cref{unions_iff}.
                Thus $G$ is inhabited.
            \end{subproof}
            Let $R = \{w\in G\times T \mid \text{there exists $U\in G$ such that there exists $V\in T$ such that
                    $w = (U,V)$ and
                    $V$ is open in $X$ with respect to $T$ and
                    $B\subseteq V$ and
                    $U$ is disjoint from $V$}\}$.
            Show that $R$ is a relation.
            \begin{subproof}
                Show that for all $w$ such that $w\in R$ there exist $U,V$ such that $w = (U,V)$.
                \begin{subproof}
                    Fix $w$.
                    Assume $w\in R$.
                    Take $U\in G$ such that there exists $V\in T$ such that
                        $w = (U,V)$ and
                        $V$ is open in $X$ with respect to $T$ and
                        $B\subseteq V$ and
                        $U$ is disjoint from $V$.
                    Take $V\in T$ such that
                        $w = (U,V)$ and
                        $V$ is open in $X$ with respect to $T$ and
                        $B\subseteq V$ and
                        $U$ is disjoint from $V$.
                \end{subproof}
                Thus $R$ is a relation by \cref{relation}.
            \end{subproof}
            Show that for all $U\in G$ there exists $V$ such that $U\mathrel{R} V$.
            \begin{subproof}
                Fix $U$.
                Assume $U\in G$.
                Thus $U\in F$ by \cref{subseteq}.
                Take $a\in A$ such that there exists $V$ such that
                    $V$ is open in $X$ with respect to $T$ and
                    $a\in U$ and
                    $B\subseteq V$ and
                    $U$ is disjoint from $V$.
                Take $V$ such that
                    $V$ is open in $X$ with respect to $T$ and
                    $a\in U$ and
                    $B\subseteq V$ and
                    $U$ is disjoint from $V$.
                $V\in T$ by \cref{open_in}.
                Thus $(U,V)\in R$.
                Thus $U\mathrel{R} V$.
            \end{subproof}
            Take $g$ such that $g$ is a function on $G$ and for all $U\in G$ we have $U\mathrel{R}\apply{g}{U}$ by \cref{choice_function}.
            Let $H = \img{g}{G}$.
            Show that $H\subseteq T$.
            \begin{subproof}
                Show that for all $V$ such that $V\in H$ we have $V\in T$.
                \begin{subproof}
                    Fix $V$.
                    Assume $V\in H$.
                    Take $U\in G$ such that $(U,V)\in g$ by \cref{img_iff}.
                    Thus $U\mathrel{R} V$ by \cref{function_apply_iff}.
                    Take $U_1\in G$ such that there exists $V_1\in T$ such that
                        $(U,V) = (U_1,V_1)$ and
                        $V_1$ is open in $X$ with respect to $T$ and
                        $B\subseteq V_1$ and
                        $U_1$ is disjoint from $V_1$.
                    Take $V_1\in T$ such that
                        $(U,V) = (U_1,V_1)$ and
                        $V_1$ is open in $X$ with respect to $T$ and
                        $B\subseteq V_1$ and
                        $U_1$ is disjoint from $V_1$.
                    Thus $V = V_1$ by \cref{pair_eq_iff}.
                    Thus $V\in T$.
                \end{subproof}
                Thus $H\subseteq T$ by \cref{subseteq}.
            \end{subproof}
            $H$ is finite by \cref{finite_image_function}.
            Show that $H$ is inhabited.
            \begin{subproof}
                Take $U$ such that $U\in G$ by \cref{nonempty_inhabited,inhabited_nonempty}.
                Thus $(U,\apply{g}{U})\in g$ by \cref{function_apply_elim,subseteq}.
                Thus $\apply{g}{U}\in H$ by \cref{img_iff}.
                Thus $H$ is inhabited.
            \end{subproof}
            Show that $G\subseteq T$.
            \begin{subproof}
                Show that for all $U$ such that $U\in G$ we have $U\in T$.
                \begin{subproof}
                    Fix $U$.
                    Assume $U\in G$.
                    Thus $U\in F$ by \cref{subseteq}.
                    Thus $U\in T$.
                \end{subproof}
                Thus $G\subseteq T$ by \cref{subseteq}.
            \end{subproof}
            Let $U_0 = \unions{G}$.
            Let $V_0 = \inters{H}$.
            $U_0\in T$ by \cref{topology_on}.
            $U_0$ is open in $X$ with respect to $T$ by \cref{open_in}.
            $V_0\in T$ by \cref{finite_intersections_open_in_topology}.
            $V_0$ is open in $X$ with respect to $T$ by \cref{open_in}.
            $A\subseteq U_0$ by \cref{covering}.
            Show that $B\subseteq V_0$.
            \begin{subproof}
                Show that for all $V$ such that $V\in H$ we have $B\subseteq V$.
                \begin{subproof}
                    Fix $V$.
                    Assume $V\in H$.
                    Take $U\in G$ such that $(U,V)\in g$ by \cref{img_iff}.
                    Thus $U\mathrel{R} V$ by \cref{function_apply_iff}.
                    Take $U_1\in G$ such that there exists $V_1\in T$ such that
                        $(U,V) = (U_1,V_1)$ and
                        $V_1$ is open in $X$ with respect to $T$ and
                        $B\subseteq V_1$ and
                        $U_1$ is disjoint from $V_1$.
                    Take $V_1\in T$ such that
                        $(U,V) = (U_1,V_1)$ and
                        $V_1$ is open in $X$ with respect to $T$ and
                        $B\subseteq V_1$ and
                        $U_1$ is disjoint from $V_1$.
                    Thus $V = V_1$ by \cref{pair_eq_iff}.
                    Thus $B\subseteq V$.
                \end{subproof}
                Thus $B\subseteq V_0$ by \cref{inters_greatest}.
            \end{subproof}
            Show that $U_0$ is disjoint from $V_0$.
            \begin{subproof}
                Suppose not.
                Then there exists $z$ such that $z\in U_0$ and $z\in V_0$.
                Take $U\in G$ such that $z\in U$ by \cref{unions_iff}.
                Thus $(U,\apply{g}{U})\in g$ by \cref{function_apply_elim,subseteq}.
                Thus $\apply{g}{U}\in H$ by \cref{img_iff}.
                Thus $V_0\subseteq \apply{g}{U}$ by \cref{inters_subseteq_elem}.
                Thus $z\in \apply{g}{U}$ by \cref{subseteq}.
                Thus $U\mathrel{R}\apply{g}{U}$ by \cref{choice_function}.
                Take $U_1\in G$ such that there exists $V_1\in T$ such that
                    $(U,\apply{g}{U}) = (U_1,V_1)$ and
                    $V_1$ is open in $X$ with respect to $T$ and
                    $B\subseteq V_1$ and
                    $U_1$ is disjoint from $V_1$.
                Take $V_1\in T$ such that
                    $(U,\apply{g}{U}) = (U_1,V_1)$ and
                    $V_1$ is open in $X$ with respect to $T$ and
                    $B\subseteq V_1$ and
                    $U_1$ is disjoint from $V_1$.
                Thus $U = U_1$ by \cref{pair_eq_iff}.
                Thus $\apply{g}{U} = V_1$ by \cref{pair_eq_iff}.
                Thus $U$ is disjoint from $\apply{g}{U}$.
                Thus $z\notin \apply{g}{U}$ by \cref{disjoint}.
                Contradiction.
            \end{subproof}
        \end{byCase}
    \end{subproof}
    Show that for all $x\in X$ we have $\{x\}$ is closed in $X$ with respect to $T$.
    \begin{subproof}
        Fix $x\in X$.
        Thus $\{x\}\subseteq X$ by \cref{singleton_subset_intro}.
        Show that $\{x\}$ is finite.
        \begin{subproof}
            $\{x,x\}$ is finite by \cref{pair_is_finite}.
            Show that for all $y$ such that $y\in\{x\}$ we have $y\in\{x,x\}$.
            \begin{subproof}
                Fix $y$.
                Assume $y\in\{x\}$.
                Thus $y = x$ by \cref{singleton_iff}.
                Thus $y\in\{x,x\}$ by \cref{upair_intro_left}.
            \end{subproof}
            Thus $\{x\}\subseteq\{x,x\}$ by \cref{subseteq}.
            Thus $\{x\}$ is finite by \cref{subseteq_finite_is_finite}.
        \end{subproof}
        Thus $\{x\}$ is closed in $X$ with respect to $T$ by \cref{finite_point_set_closed_hausdorff}.
    \end{subproof}
    Thus $X$ has closed singletons with respect to $T$ by \cref{one_point_closed}.
    Thus $X$ is normal with respect to $T$ by \cref{normal_space}.
\end{proof}

% from §32 helper definition 2: well order relation
\begin{definition}\label{well_order_relation}
    $R$ is a well order relation on $X$ iff
        $R$ is a simple order relation on $X$ and
        for all $A$ such that $A\subseteq X$ and $A\neq\emptyset$
            there exists $a$ such that $a$ is a smallest element of $A$ with respect to $R$.
\end{definition}

% from §32 Item 4: well-ordered set is normal in the order topology
\begin{theorem}\label{well_ordered_normal}
    Assume $R$ is a well order relation on $X$.
    Assume $T$ is the order topology on $X$ with respect to $R$.
    Then $X$ is normal with respect to $T$.
\end{theorem}
\begin{proof}
    Omitted. % do not touch
\end{proof}

% from §33 helper definition 1: unit rational set
\begin{definition}\label{unit_rationals}
    $\unitRationals =
        \{\fst{q} \rmul \inv{\snd{q}} \mid q\in\naturalsPlus\times\naturalsPlus\}$.
\end{definition}

% from §33 helper definition 2: unit rational set with zero
\begin{definition}\label{unit_rationals_zero}
    $\unitRationalsZero = \unitRationals \union \{0\}$.
\end{definition}

% from §33 helper definition 3: Urysohn admissible functions
\begin{definition}\label{urysohn_admissible}
    $\urysohnAdmissible{s}{X}{T}{A}{B}{n} =
        \{f\in\funs{\{0,1\}\union \img{s}{\initialSegment{n}}}{\pow{X}} \mid
            \text{for all $p\in\dom{f}$ we have $f(p)\in T$ and
            $f(1) = X\setminus B$ and
            $A\subseteq f(0)$ and
            $\closure{f(0)}{X}{T}\subseteq f(1)$ and
            for all $p\in\dom{f}$ for all $q\in\dom{f}$ if $p < q$ then
                $\closure{f(p)}{X}{T}\subseteq f(q)$ and
            for all $p\in\dom{f}$ if $1 < p$ then $f(p) = X$}\}$.
\end{definition}

% from §33 helper definition 4: Urysohn good pair predicate
\begin{definition}\label{urysohn_good_pair}
    $\urysohnGoodPair{s}{X}{T}{A}{B}{k}{g}$ iff
        $k\in\naturalsPlus$ and $g\in\urysohnAdmissible{s}{X}{T}{A}{B}{k}$.
\end{definition}

% from §33 helper definition 5: Urysohn good pair set
\begin{definition}\label{urysohn_good_pairs}
    $\urysohnGoodPairs{s}{X}{T}{A}{B} =
        \{(k,g) \mid k\in\naturalsPlus,
            g\in\rels{\{0,1\}\union \img{s}{\naturalsPlus}}{\pow{X}} \mid
            \text{$\urysohnGoodPair{s}{X}{T}{A}{B}{k}{g}$}\}$.
\end{definition}

% from §33 helper definition 6: Urysohn level set
\begin{definition}\label{urysohn_level_set}
    $\urysohnLevelSet{U}{x} = \{q\in\unitRationalsZero \mid x\in U(q)\}$.
\end{definition}

% from §33 helper lemma 1: unit rationals are dense in nonnegative intervals
\begin{lemma}\label{unit_rationals_dense_nonneg}
    Suppose $a\in\reals$ and $b\in\reals$ and $0 \leq a$ and $a < b$.
    Then there exists $p$ such that $p\in\unitRationals$ and $p\in\intervalopen{a}{b}$.
\end{lemma}
\begin{proof}
    Take $n, N$ such that $n\in\naturalsPlus$ and $N\in\naturalsPlus$ and $n / N\in\intervalopen{a}{b}$
        by \cref{real_rationals_nonneg_interval}.
    Let $p = n / N$.
    Let $q = (n,N)$.
    Thus $q\in\naturalsPlus\times\naturalsPlus$ by \cref{times_tuple_intro}.
    Thus $p = n \rmul \inv{N}$ by \cref{real_div}.
    Thus $p = \fst{q} \rmul \inv{\snd{q}}$ by \cref{fst_eq,snd_eq}.
    Thus $p\in\unitRationals$ by \cref{unit_rationals}.
    Thus $p\in\intervalopen{a}{b}$.
\end{proof}

% from §33 helper lemma 2: sequence enumerating unit rationals
\begin{lemma}\label{unit_rationals_sequence}
    There exists $s$ such that
        $s$ is a sequence in $\unitRationals$ and
        for all $p$ such that $p\in\unitRationals$ there exists $n\in\naturalsPlus$ such that $s(n) = p$.
\end{lemma}
\begin{proof}
    Take $g$ such that $g\in\inj{\naturalsPlus\times\naturalsPlus}{\naturalsPlus}$
        by \cref{naturalsplus_pairing_injection}.
    Let $p_0 = (1,1)$.
    Let $R = \converse{g} \union ((\naturalsPlus\setminus\ran{g})\times\{p_0\})$.
    Show that $R$ is a function.
    \begin{subproof}
        Show that for all $x,y_1,y_2$ such that $(x,y_1)\in R$ and $(x,y_2)\in R$ we have $y_1 = y_2$.
        \begin{subproof}
            Fix $x$.
            Fix $y_1$.
            Fix $y_2$.
            Assume $(x,y_1)\in R$ and $(x,y_2)\in R$.
            \begin{byCase}
            \caseOf{$x\in\ran{g}$.}
                Show that $(x,y_1)\in\converse{g}$.
                \begin{subproof}
                    Suppose not.
                    Thus $(x,y_1)\notin\converse{g}$.
                    Thus $(x,y_1)\in\converse{g}$ or $(x,y_1)\in(\naturalsPlus\setminus\ran{g})\times\{p_0\}$
                        by \cref{union_iff}.
                    \begin{byCase}
                    \caseOf{$(x,y_1)\in\converse{g}$.}
                        Contradiction.
                    \caseOf{$(x,y_1)\in(\naturalsPlus\setminus\ran{g})\times\{p_0\}$.}
                        Thus $x\in\naturalsPlus\setminus\ran{g}$ by \cref{times_tuple_elim}.
                        Thus $x\notin\ran{g}$ by \cref{setminus_elim_right}.
                        Contradiction.
                    \end{byCase}
                \end{subproof}
                Show that $(x,y_2)\in\converse{g}$.
                \begin{subproof}
                    Suppose not.
                    Thus $(x,y_2)\notin\converse{g}$.
                    Thus $(x,y_2)\in\converse{g}$ or $(x,y_2)\in(\naturalsPlus\setminus\ran{g})\times\{p_0\}$
                        by \cref{union_iff}.
                    \begin{byCase}
                    \caseOf{$(x,y_2)\in\converse{g}$.}
                        Contradiction.
                    \caseOf{$(x,y_2)\in(\naturalsPlus\setminus\ran{g})\times\{p_0\}$.}
                        Thus $x\in\naturalsPlus\setminus\ran{g}$ by \cref{times_tuple_elim}.
                        Thus $x\notin\ran{g}$ by \cref{setminus_elim_right}.
                        Contradiction.
                    \end{byCase}
                \end{subproof}
                Thus $y_1\mathrel{g} x$ and $y_2\mathrel{g} x$ by \cref{converse_iff}.
                $g\in\funs{\naturalsPlus\times\naturalsPlus}{\naturalsPlus}$ by \cref{inj}.
                Thus $g$ is a function by \cref{funs_is_function}.
                Thus $g(y_1) = x$ and $g(y_2) = x$ by \cref{function_apply_intro}.
                Thus $g(y_1) = g(y_2)$.
                Thus $y_1\in\dom{g}$ and $y_2\in\dom{g}$ by \cref{dom_iff}.
                Thus $\dom{g} = \naturalsPlus\times\naturalsPlus$ by \cref{funs}.
                Thus $y_1\in\naturalsPlus\times\naturalsPlus$ and $y_2\in\naturalsPlus\times\naturalsPlus$
                    by \cref{subseteq}.
                Thus $y_1 = y_2$ by \cref{inj}.
            \caseOf{$x\notin\ran{g}$.}
                Show that $(x,y_1)\in(\naturalsPlus\setminus\ran{g})\times\{p_0\}$.
                \begin{subproof}
                    Suppose not.
                    Thus $(x,y_1)\notin(\naturalsPlus\setminus\ran{g})\times\{p_0\}$.
                    Thus $(x,y_1)\in\converse{g}$ or $(x,y_1)\in(\naturalsPlus\setminus\ran{g})\times\{p_0\}$
                        by \cref{union_iff}.
                    \begin{byCase}
                    \caseOf{$(x,y_1)\in\converse{g}$.}
                        Thus $y_1\mathrel{g} x$ by \cref{converse_iff}.
                        Thus $x\in\ran{g}$ by \cref{ran_iff}.
                        Contradiction.
                    \caseOf{$(x,y_1)\in(\naturalsPlus\setminus\ran{g})\times\{p_0\}$.}
                        Contradiction.
                    \end{byCase}
                \end{subproof}
                Show that $(x,y_2)\in(\naturalsPlus\setminus\ran{g})\times\{p_0\}$.
                \begin{subproof}
                    Suppose not.
                    Thus $(x,y_2)\notin(\naturalsPlus\setminus\ran{g})\times\{p_0\}$.
                    Thus $(x,y_2)\in\converse{g}$ or $(x,y_2)\in(\naturalsPlus\setminus\ran{g})\times\{p_0\}$
                        by \cref{union_iff}.
                    \begin{byCase}
                    \caseOf{$(x,y_2)\in\converse{g}$.}
                        Thus $y_2\mathrel{g} x$ by \cref{converse_iff}.
                        Thus $x\in\ran{g}$ by \cref{ran_iff}.
                        Contradiction.
                    \caseOf{$(x,y_2)\in(\naturalsPlus\setminus\ran{g})\times\{p_0\}$.}
                        Contradiction.
                    \end{byCase}
                \end{subproof}
                Thus $y_1\in\{p_0\}$ and $y_2\in\{p_0\}$ by \cref{times_tuple_elim}.
                Thus $y_1 = p_0$ and $y_2 = p_0$ by \cref{singleton_iff}.
                Thus $y_1 = y_2$.
            \end{byCase}
        \end{subproof}
        Thus $R$ is right-unique by \cref{rightunique}.
        Show that $R$ is a relation.
        \begin{subproof}
            Show that for all $w$ such that $w\in R$ there exist $x,y$ such that $w = (x,y)$.
            \begin{subproof}
                Fix $w$.
                Assume $w\in R$.
                Thus $w\in\converse{g}$ or $w\in(\naturalsPlus\setminus\ran{g})\times\{p_0\}$ by \cref{union_iff}.
                \begin{byCase}
                \caseOf{$w\in\converse{g}$.}
                    Take $w_0\in g$ such that there exist $x,y$ such that $w_0 = (x,y)$ and $w = (y,x)$
                        by \cref{converse_relation}.
                    Take $x,y$ such that $w_0 = (x,y)$ and $w = (y,x)$.
                    Thus there exist $x_1,y_1$ such that $w = (x_1,y_1)$.
                \caseOf{$w\in(\naturalsPlus\setminus\ran{g})\times\{p_0\}$.}
                    Thus there exist $x,y$ such that $w = (x,y)$ by \cref{times_elem_is_tuple}.
                \end{byCase}
            \end{subproof}
            Thus $R$ is a relation by \cref{relation}.
        \end{subproof}
        Thus $R$ is a function.
    \end{subproof}
    Show that $\dom{R} = \naturalsPlus$.
    \begin{subproof}
        Show that for all $x\in\naturalsPlus$ we have $x\in\dom{R}$.
        \begin{subproof}
            Fix $x$.
            Assume $x\in\naturalsPlus$.
            \begin{byCase}
            \caseOf{$x\in\ran{g}$.}
                Take $p$ such that $p\mathrel{g} x$ by \cref{ran_iff}.
                Thus $x\mathrel{\converse{g}} p$ by \cref{converse_iff}.
                Thus $(x,p)\in R$ by \cref{union_iff}.
                Thus $x\in\dom{R}$ by \cref{dom_iff}.
            \caseOf{$x\notin\ran{g}$.}
                Thus $x\in\naturalsPlus\setminus\ran{g}$ by \cref{setminus_intro}.
                Thus $p_0\in\{p_0\}$ by \cref{singleton_intro}.
                Thus $(x,p_0)\in(\naturalsPlus\setminus\ran{g})\times\{p_0\}$ by \cref{times_tuple_intro}.
                Thus $(x,p_0)\in R$ by \cref{union_iff}.
                Thus $x\in\dom{R}$ by \cref{dom_iff}.
            \end{byCase}
        \end{subproof}
        Thus $\naturalsPlus\subseteq\dom{R}$ by \cref{subseteq}.
        Show that for all $x\in\dom{R}$ we have $x\in\naturalsPlus$.
        \begin{subproof}
            Fix $x$.
            Assume $x\in\dom{R}$.
            Take $y$ such that $(x,y)\in R$ by \cref{dom_iff}.
            Thus $(x,y)\in\converse{g}$ or $(x,y)\in(\naturalsPlus\setminus\ran{g})\times\{p_0\}$ by \cref{union_iff}.
            \begin{byCase}
            \caseOf{$(x,y)\in\converse{g}$.}
                Thus $y\mathrel{g} x$ by \cref{converse_iff}.
                $g\in\funs{\naturalsPlus\times\naturalsPlus}{\naturalsPlus}$ by \cref{inj}.
                Thus $g$ is a function by \cref{funs_is_function}.
                Thus $y\in\dom{g}$ by \cref{dom_iff}.
                Thus $\dom{g} = \naturalsPlus\times\naturalsPlus$ by \cref{funs}.
                Thus $y\in\naturalsPlus\times\naturalsPlus$ by \cref{subseteq}.
                Thus $g(y)\in\naturalsPlus$ by \cref{funs_type_apply}.
                Thus $g(y) = x$ by \cref{function_apply_intro}.
                Thus $x\in\naturalsPlus$.
            \caseOf{$(x,y)\in(\naturalsPlus\setminus\ran{g})\times\{p_0\}$.}
                Thus $x\in\naturalsPlus\setminus\ran{g}$ by \cref{times_tuple_elim}.
                Thus $x\in\naturalsPlus$ by \cref{setminus_elim_left}.
            \end{byCase}
        \end{subproof}
        Thus $\dom{R}\subseteq\naturalsPlus$ by \cref{subseteq}.
        Thus $\dom{R} = \naturalsPlus$ by \cref{subseteq_antisymmetric}.
    \end{subproof}
    Let $h = \{(n, \fst{R(n)} \rmul \inv{\snd{R(n)}}) \mid n\in\naturalsPlus\}$.
    Show that $h$ is a function from $\naturalsPlus$ to $\unitRationals$.
    \begin{subproof}
        Show that $h$ is a function.
        \begin{subproof}
            Show that for all $n,y_1,y_2$ such that $(n,y_1)\in h$ and $(n,y_2)\in h$ we have $y_1 = y_2$.
            \begin{subproof}
                Fix $n$.
                Fix $y_1$.
                Fix $y_2$.
                Assume $(n,y_1)\in h$ and $(n,y_2)\in h$.
                Thus $y_1 = \fst{R(n)} \rmul \inv{\snd{R(n)}}$ and
                    $y_2 = \fst{R(n)} \rmul \inv{\snd{R(n)}}$ by \cref{pair_eq_iff}.
                Thus $y_1 = y_2$.
            \end{subproof}
            Thus $h$ is right-unique by \cref{rightunique}.
            Show that $h$ is a relation.
            \begin{subproof}
                Show that for all $w$ such that $w\in h$ there exist $x,y$ such that $w = (x,y)$.
                \begin{subproof}
                    Fix $w$.
                    Assume $w\in h$.
                    Thus there exist $x,y$ such that $w = (x,y)$ by \cref{pair_eq_iff}.
                \end{subproof}
                Thus $h$ is a relation by \cref{relation}.
            \end{subproof}
            Thus $h$ is a function.
        \end{subproof}
        Show that for all $n\in\dom{h}$ we have $h(n)\in\unitRationals$.
        \begin{subproof}
            Fix $n\in\dom{h}$.
            Take $y$ such that $(n,y)\in h$ by \cref{dom_iff}.
            Thus $y = \fst{R(n)} \rmul \inv{\snd{R(n)}}$ by \cref{pair_eq_iff}.
            Thus $n\in\naturalsPlus$ by \cref{pair_eq_iff}.
            Thus $n\in\dom{R}$ by \cref{subseteq}.
            $R$ is a function.
            Thus $(n,R(n))\in R$ by \cref{function_apply_elim}.
            Thus $(n,R(n))\in\converse{g}$ or $(n,R(n))\in(\naturalsPlus\setminus\ran{g})\times\{p_0\}$ by \cref{union_iff}.
            Show that $R(n)\in\naturalsPlus\times\naturalsPlus$.
            \begin{subproof}
                \begin{byCase}
                \caseOf{$(n,R(n))\in\converse{g}$.}
                    Thus $R(n)\mathrel{g} n$ by \cref{converse_iff}.
                    Thus $(R(n),n)\in g$.
                    Thus $R(n)\in\dom{g}$ by \cref{dom_iff}.
                    $g\in\funs{\naturalsPlus\times\naturalsPlus}{\naturalsPlus}$ by \cref{inj}.
                    Thus $\dom{g} = \naturalsPlus\times\naturalsPlus$ by \cref{funs}.
                    Thus $R(n)\in\naturalsPlus\times\naturalsPlus$ by \cref{subseteq}.
                \caseOf{$(n,R(n))\in(\naturalsPlus\setminus\ran{g})\times\{p_0\}$.}
                    Thus $R(n)\in\{p_0\}$ by \cref{times_tuple_elim}.
                    $1\in\naturalsPlus$ by \cref{real_one_in_naturals}.
                    Thus $p_0\in\naturalsPlus\times\naturalsPlus$ by \cref{times_tuple_intro}.
                    Thus $\{p_0\}\subseteq\naturalsPlus\times\naturalsPlus$ by \cref{singleton_subset_intro}.
                    Thus $R(n)\in\naturalsPlus\times\naturalsPlus$ by \cref{subseteq}.
            \end{byCase}
        \end{subproof}
        Thus $\fst{R(n)}\in\naturalsPlus$ by \cref{fst_type}.
        Thus $\snd{R(n)}\in\naturalsPlus$ by \cref{snd_type}.
        Thus $y\in\unitRationals$ by \cref{unit_rationals}.
        Thus $h(n)\in\unitRationals$ by \cref{function_apply_intro}.
        \end{subproof}
        Thus $h$ is a function to $\unitRationals$.
        Show that $\dom{h} = \naturalsPlus$.
        \begin{subproof}
            Show that for all $n\in\naturalsPlus$ we have $n\in\dom{h}$.
            \begin{subproof}
                Fix $n$.
                Assume $n\in\naturalsPlus$.
                Thus $(n,\fst{R(n)} \rmul \inv{\snd{R(n)}})\in h$ by \cref{pair_eq_iff}.
                Thus $n\in\dom{h}$ by \cref{dom_iff}.
            \end{subproof}
            Thus $\naturalsPlus\subseteq\dom{h}$ by \cref{subseteq}.
            Show that for all $n\in\dom{h}$ we have $n\in\naturalsPlus$.
            \begin{subproof}
                Fix $n$.
                Assume $n\in\dom{h}$.
                Take $y$ such that $(n,y)\in h$ by \cref{dom_iff}.
                Thus $n\in\naturalsPlus$ by \cref{pair_eq_iff}.
            \end{subproof}
            Thus $\dom{h}\subseteq\naturalsPlus$ by \cref{subseteq}.
            Thus $\dom{h} = \naturalsPlus$ by \cref{subseteq_antisymmetric}.
        \end{subproof}
        Thus $h\in\funs{\naturalsPlus}{\unitRationals}$ by \cref{funs_intro}.
        Thus $h$ is a sequence in $\unitRationals$ by \cref{sequence_in}.
    \end{subproof}
    Show that for all $p$ such that $p\in\unitRationals$ there exists $n\in\naturalsPlus$ such that $h(n) = p$.
    \begin{subproof}
        Fix $p$.
        Assume $p\in\unitRationals$.
        Take $q$ such that $q\in\naturalsPlus\times\naturalsPlus$ and $p = \fst{q} \rmul \inv{\snd{q}}$
            by \cref{unit_rationals}.
        $g\in\inj{\naturalsPlus\times\naturalsPlus}{\naturalsPlus}$ by \cref{naturalsplus_pairing_injection}.
        Thus $g\in\funs{\naturalsPlus\times\naturalsPlus}{\naturalsPlus}$ by \cref{inj}.
        Thus $g$ is a function by \cref{funs_is_function}.
        Thus $\dom{g} = \naturalsPlus\times\naturalsPlus$ by \cref{funs}.
        Thus $q\in\dom{g}$ by \cref{subseteq}.
        Thus $g(q)\in\naturalsPlus$ by \cref{funs_type_apply}.
        Let $m = g(q)$.
        Thus $m\in\naturalsPlus$.
        Thus $(q,m)\in g$ by \cref{function_apply_elim}.
        Thus $q\mathrel{g} m$.
        Thus $m\mathrel{\converse{g}} q$ by \cref{converse_iff}.
        Thus $(m,q)\in R$ by \cref{union_iff}.
        $R$ is a function and $m\in\dom{R}$ by \cref{dom_iff}.
        Thus $R(m) = q$ by \cref{function_apply_intro}.
        Thus $(m,\fst{R(m)} \rmul \inv{\snd{R(m)}})\in h$ by \cref{pair_eq_iff}.
        Thus $h$ is right-unique by \cref{funs}.
        Thus $h$ is a relation by \cref{funs_is_relation}.
        Thus $h(m) = \fst{R(m)} \rmul \inv{\snd{R(m)}}$ by \cref{function_apply_intro}.
        Thus $h(m) = \fst{q} \rmul \inv{\snd{q}}$.
        Thus $h(m) = p$.
    \end{subproof}
    Thus there exists $s$ such that
        $s$ is a sequence in $\unitRationals$ and
        for all $p$ such that $p\in\unitRationals$ there exists $n\in\naturalsPlus$ such that $s(n) = p$.
\end{proof}


% from §33 helper lemma 3: good pair decomposition
\begin{lemma}\label{urysohn_good_pair_decomp}
    Let $G = \urysohnGoodPairs{s}{X}{T}{A}{B}$.
    Suppose $z\in G$.
    Then there exist $n,f$ such that
        $n\in\naturalsPlus$ and
        $f\in\urysohnAdmissible{s}{X}{T}{A}{B}{n}$ and
        $z = (n,f)$.
\end{lemma}
\begin{proof}
    Take $n,f$ such that
        $n\in\naturalsPlus$ and
        $f\in\rels{\{0,1\}\union \img{s}{\naturalsPlus}}{\pow{X}}$ and
        $\urysohnGoodPair{s}{X}{T}{A}{B}{n}{f}$ and
        $z = (n,f)$.
    Thus $n\in\naturalsPlus$ and $f\in\urysohnAdmissible{s}{X}{T}{A}{B}{n}$
        by \cref{urysohn_good_pair}.
\end{proof}

% from §33 helper lemma 4: good pair introduction
\begin{lemma}\label{urysohn_good_pair_intro}
    Let $G = \urysohnGoodPairs{s}{X}{T}{A}{B}$.
    Suppose $n\in\naturalsPlus$ and $f\in\urysohnAdmissible{s}{X}{T}{A}{B}{n}$.
    Then $(n,f)\in G$.
\end{lemma}
\begin{proof}
    Thus $\urysohnGoodPair{s}{X}{T}{A}{B}{n}{f}$ by \cref{urysohn_good_pair}.
    Let $D = \{0,1\}\union \img{s}{\initialSegment{n}}$.
    Thus $f\in\funs{D}{\pow{X}}$ by \cref{urysohn_admissible}.
    Thus $f\in\rels{D}{\pow{X}}$ by \cref{funs_subseteq_rels,elem_subseteq}.
    Thus $f\subseteq D\times\pow{X}$ by \cref{rels,powerset_elim}.
    Show that $D\subseteq \{0,1\}\union \img{s}{\naturalsPlus}$.
    \begin{subproof}
        $\{0,1\}\subseteq \{0,1\}\union \img{s}{\naturalsPlus}$ by \cref{union_upper_left}.
        Show that for all $x$ such that $x\in\initialSegment{n}$ we have $x\in\naturalsPlus$.
        \begin{subproof}
            Fix $x$.
            Assume $x\in\initialSegment{n}$.
            Thus $x\in\naturalsPlus$ by \cref{initial_segment_naturalsplus}.
        \end{subproof}
        Thus $\initialSegment{n}\subseteq\naturalsPlus$ by \cref{subseteq}.
        Thus $\img{s}{\initialSegment{n}}\subseteq \img{s}{\naturalsPlus}$ by \cref{img_subseteq}.
        Thus $\img{s}{\naturalsPlus}\subseteq \{0,1\}\union \img{s}{\naturalsPlus}$ by \cref{union_upper_right}.
        Thus $\img{s}{\initialSegment{n}}\subseteq \{0,1\}\union \img{s}{\naturalsPlus}$
            by \cref{subseteq_transitive}.
        Thus $D\subseteq \{0,1\}\union \img{s}{\naturalsPlus}$ by \cref{union_subsets_is_subset}.
    \end{subproof}
    Thus $D\times\pow{X}\subseteq (\{0,1\}\union \img{s}{\naturalsPlus})\times\pow{X}$
        by \cref{times_subseteq_left}.
    Thus $f\subseteq (\{0,1\}\union \img{s}{\naturalsPlus})\times\pow{X}$ by \cref{subseteq_transitive}.
    Thus $f\in\rels{\{0,1\}\union \img{s}{\naturalsPlus}}{\pow{X}}$ by \cref{rels_intro}.
    Thus $(n,f)\in G$.
\end{proof}

% from §33 helper lemma 6: continuity at a point via open neighborhoods
\begin{lemma}\label{continuous_at_open}
    Assume $f$ is a function from $X$ to $Y$.
    Assume $T$ is a topology on $X$.
    Assume $T_1$ is a topology on $Y$.
    Assume $x\in X$.
    Suppose for all $V$ such that $V$ is open in $Y$ with respect to $T_1$ and $f(x)\in V$ there exists $U$ such that
        $U$ is open in $X$ with respect to $T$ and $x\in U$ and $\img{f}{U}\subseteq V$.
    Then $f$ is continuous at $x$ from $X$ to $Y$ with respect to $T$ and $T_1$.
\end{lemma}
\begin{proof}
    Show that for all $V$ such that $V$ is a neighborhood of $f(x)$ in $Y$ with respect to $T_1$ there exists $U$ such that
        $U$ is a neighborhood of $x$ in $X$ with respect to $T$ and $\img{f}{U}\subseteq V$.
    \begin{subproof}
        Fix $V$.
        Assume $V$ is a neighborhood of $f(x)$ in $Y$ with respect to $T_1$.
        Thus $V$ is open in $Y$ with respect to $T_1$ by \cref{neighborhood}.
        Thus $f(x)\in V$ by \cref{neighborhood}.
        Take $U$ such that $U$ is open in $X$ with respect to $T$ and $x\in U$ and $\img{f}{U}\subseteq V$.
        Thus $U$ is a neighborhood of $x$ in $X$ with respect to $T$ by \cref{neighborhood}.
        Thus $\img{f}{U}\subseteq V$.
    \end{subproof}
    Thus $f$ is continuous at $x$ from $X$ to $Y$ with respect to $T$ and $T_1$ by \cref{continuous_at}.
\end{proof}

% from §33 helper lemma 7: continuity to a subspace via range inclusion
\begin{lemma}\label{continuous_subspace_range}
    Assume $f\in\funs{X}{\reals}$.
    Assume $T$ is a topology on $X$.
    Assume $I_0\subseteq\reals$.
    Let $T_0 = \subspaceTopology{\standardTopologyR}{I_0}$.
    Assume $f$ is continuous from $X$ to $\reals$ with respect to $T$ and $\standardTopologyR$.
    Assume $\img{f}{X}\subseteq I_0$.
    Then $f$ is continuous from $X$ to $I_0$ with respect to $T$ and $T_0$.
\end{lemma}
\begin{proof}
    $\standardBasisR$ is a basis on $\reals$ by \cref{standard_basis_is_basis}.
    Thus $\basisTopology{\standardBasisR}{\reals}$ is a topology on $\reals$
        by \cref{basis_topology_is_topology}.
    Thus $\standardTopologyR$ is a topology on $\reals$ by \cref{standard_topology_reals}.
    Thus $I_0$ is a subspace of $\reals$ with respect to $\standardTopologyR$ by \cref{subspace_of}.
    Thus $T_0$ is a topology on $I_0$ by \cref{subspace_topology_is_topology,subspace_of}.
    Show that for all $B$ such that $B$ is open in $I_0$ with respect to $T_0$
        we have $\preimg{f}{B}$ is open in $X$ with respect to $T$.
    \begin{subproof}
        Fix $B$.
        Assume $B$ is open in $I_0$ with respect to $T_0$.
        Take $U\in\standardTopologyR$ such that $B = I_0\inter U$ by \cref{subspace_topology,open_in}.
        $U$ is open in $\reals$ with respect to $\standardTopologyR$ by \cref{open_in}.
        $\preimg{f}{U}$ is open in $X$ with respect to $T$ by \cref{continuous}.
        Show that $\preimg{f}{B} = \preimg{f}{U}$.
        \begin{subproof}
            Show that for all $x$ such that $x\in\preimg{f}{B}$ we have $x\in\preimg{f}{U}$.
            \begin{subproof}
                Fix $x$.
                Assume $x\in\preimg{f}{B}$.
                Take $y\in B$ such that $x\mathrel{f} y$ by \cref{preimg_iff}.
                Thus $y\in U$ by \cref{inter_elim_right}.
                Thus $x\in\preimg{f}{U}$ by \cref{preimg_iff}.
            \end{subproof}
            Thus $\preimg{f}{B}\subseteq \preimg{f}{U}$ by \cref{subseteq}.
            Show that for all $x$ such that $x\in\preimg{f}{U}$ we have $x\in\preimg{f}{B}$.
            \begin{subproof}
                Fix $x$.
                Assume $x\in\preimg{f}{U}$.
                Take $y\in U$ such that $x\mathrel{f} y$ by \cref{preimg_iff}.
                $\preimg{f}{U}\subseteq \dom{f}$ by \cref{preimg_subseteq_dom}.
                Thus $x\in\dom{f}$ by \cref{subseteq}.
                Thus $x\in X$ by \cref{funs}.
                Thus $f$ is a function by \cref{funs_is_function}.
                Thus $f$ is right-unique by \cref{function_rightunique}.
                Thus $f$ is a relation by \cref{funs_is_relation}.
                Thus $y = f(x)$ by \cref{function_apply_iff}.
                Thus $y\in\img{f}{X}$ by \cref{img_of_function_intro,inter_intro}.
                Thus $y\in I_0$ by \cref{subseteq}.
                Thus $y\in I_0\inter U$ by \cref{inter_intro}.
                Thus $y\in B$.
                Thus $x\in\preimg{f}{B}$ by \cref{preimg_iff}.
            \end{subproof}
            Thus $\preimg{f}{U}\subseteq \preimg{f}{B}$ by \cref{subseteq}.
            Thus $\preimg{f}{B} = \preimg{f}{U}$ by \cref{subseteq_antisymmetric}.
        \end{subproof}
        Thus $\preimg{f}{B}$ is open in $X$ with respect to $T$.
    \end{subproof}
    Thus for all $B$ such that $B$ is open in $I_0$ with respect to $T_0$
        we have $\preimg{f}{B}$ is open in $X$ with respect to $T$.
    Thus $f$ is a function by \cref{funs_is_function}.
    Thus $f$ is right-unique by \cref{function_rightunique}.
    Thus $f$ is a relation by \cref{funs_is_relation}.
    Thus $\dom{f} = X$ by \cref{funs_elim}.
    Thus $\img{f}{\dom{f}} = \ran{f}$ by \cref{img_dom}.
    Thus $\ran{f} = \img{f}{X}$ by \cref{funs}.
    Thus $\ran{f}\subseteq I_0$ by \cref{subseteq}.
    Show that for all $x\in\dom{f}$ we have $f(x)\in I_0$.
    \begin{subproof}
        Fix $x\in\dom{f}$.
        Thus $(x,f(x))\in f$ by \cref{function_apply_elim}.
        Thus $f(x)\in\ran{f}$ by \cref{ran_iff}.
        Thus $f(x)\in I_0$ by \cref{subseteq}.
    \end{subproof}
    Thus $f\in\funs{X}{I_0}$ by \cref{funs_intro}.
    Thus $f$ is continuous from $X$ to $I_0$ with respect to $T$ and $T_0$
        by \cref{continuous,funs_intro,subspace_topology_is_topology,subspace_of}.
\end{proof}

% from §33 helper lemma 8: initial good pair when $r_1 = 1$
\begin{lemma}\label{urysohn_u1_eq1}
    Assume $U_1 = X\setminus B$.
    Assume $U_1$ is open in $X$ with respect to $T$.
    Assume $U_0$ is open in $X$ with respect to $T$.
    Assume $A\subseteq U_0$.
    Assume $\closure{U_0}{X}{T}\subseteq U_1$.
    Assume $r_1 = 1$.
    Assume $\img{s}{\initialSegment{1}} = \{r_1\}$.
    Let $G = \urysohnGoodPairs{s}{X}{T}{A}{B}$.
    Then there exists $u_1$ such that $(1,u_1)\in G$.
\end{lemma}
\begin{proof}
    Let $u_1 = \{(0,U_0),(1,U_1)\}$.
    Show that $(1,u_1)\in G$.
    \begin{subproof}
        Show that $u_1\in\urysohnAdmissible{s}{X}{T}{A}{B}{1}$.
        \begin{subproof}
            Show that $u_1\in\funs{\{0,1\}\union \img{s}{\initialSegment{1}}}{\pow{X}}$.
            \begin{subproof}
                Show that $u_1$ is right-unique.
                \begin{subproof}
                    Show that for all $a,b,b'$ such that $(a,b)\in u_1$ and $(a,b')\in u_1$ we have $b = b'$.
                    \begin{subproof}
                        Fix $a$.
                        Fix $b$.
                        Fix $b'$.
                        Assume $(a,b)\in u_1$ and $(a,b')\in u_1$.
                        Thus $(a,b) = (0,U_0)$ or $(a,b) = (1,U_1)$ by \cref{upair_elim}.
                        \begin{byCase}
                        \caseOf{$(a,b) = (0,U_0)$.}
                            Thus $a = 0$ by \cref{pair_eq_iff}.
                            Show that $b = b'$.
                            \begin{subproof}
                                Suppose not.
                                Thus $b \neq b'$.
                                Thus $(a,b') = (0,U_0)$ or $(a,b') = (1,U_1)$ by \cref{upair_elim}.
                                \begin{byCase}
                                \caseOf{$(a,b') = (0,U_0)$.}
                                    Thus $b = b'$ by \cref{pair_eq_iff}.
                                    Contradiction.
                                \caseOf{$(a,b') = (1,U_1)$.}
                                    Thus $a = 1$ by \cref{pair_eq_iff}.
                                    Contradiction.
                                \end{byCase}
                            \end{subproof}
                        \caseOf{$(a,b) = (1,U_1)$.}
                            Thus $a = 1$ by \cref{pair_eq_iff}.
                            Show that $b = b'$.
                            \begin{subproof}
                                Suppose not.
                                Thus $b \neq b'$.
                                Thus $(a,b') = (0,U_0)$ or $(a,b') = (1,U_1)$ by \cref{upair_elim}.
                                \begin{byCase}
                                \caseOf{$(a,b') = (1,U_1)$.}
                                    Thus $b = b'$ by \cref{pair_eq_iff}.
                                    Contradiction.
                                \caseOf{$(a,b') = (0,U_0)$.}
                                    Thus $a = 0$ by \cref{pair_eq_iff}.
                                    Contradiction.
                                \end{byCase}
                            \end{subproof}
                        \end{byCase}
                    \end{subproof}
                    Thus $u_1$ is right-unique by \cref{rightunique}.
                \end{subproof}
                Show that $u_1$ is a relation.
                \begin{subproof}
                    Show that for all $w$ such that $w\in u_1$ there exist $x,y$ such that $w = (x,y)$.
                    \begin{subproof}
                        Fix $w$.
                        Assume $w\in u_1$.
                        Thus $w = (0,U_0)$ or $w = (1,U_1)$ by \cref{upair_elim}.
                        Thus there exist $x,y$ such that $w = (x,y)$.
                    \end{subproof}
                    Thus $u_1$ is a relation by \cref{relation}.
                \end{subproof}
                Thus $u_1$ is a function.
                Show that $\dom{u_1} = \{0,1\}\union \img{s}{\initialSegment{1}}$.
                \begin{subproof}
                    Show that $\dom{u_1} = \{0,1\}$.
                    \begin{subproof}
                        Show that $\dom{u_1}\subseteq \{0,1\}$.
                        \begin{subproof}
                            Show that for all $x\in\dom{u_1}$ we have $x\in\{0,1\}$.
                            \begin{subproof}
                                Fix $x$.
                                Assume $x\in\dom{u_1}$.
                                Take $y$ such that $(x,y)\in u_1$ by \cref{dom_iff}.
                                Thus $(x,y) = (0,U_0)$ or $(x,y) = (1,U_1)$ by \cref{upair_elim}.
                                Thus $x = 0$ or $x = 1$ by \cref{pair_eq_iff}.
                                Thus $x\in\{0,1\}$ by \cref{upair_intro_left,upair_intro_right}.
                            \end{subproof}
                            Thus $\dom{u_1}\subseteq \{0,1\}$ by \cref{subseteq}.
                        \end{subproof}
                        Show that $\{0,1\}\subseteq \dom{u_1}$.
                        \begin{subproof}
                            Show that for all $x\in\{0,1\}$ we have $x\in\dom{u_1}$.
                            \begin{subproof}
                                Fix $x$.
                                Assume $x\in\{0,1\}$.
                                Thus $x = 0$ or $x = 1$ by \cref{upair_elim}.
                                \begin{byCase}
                                \caseOf{$x = 0$.}
                                    Thus $(0,U_0)\in u_1$ by \cref{upair_intro_left}.
                                    Thus $x\in\dom{u_1}$ by \cref{dom_intro}.
                                \caseOf{$x = 1$.}
                                    Thus $(1,U_1)\in u_1$ by \cref{upair_intro_right}.
                                    Thus $x\in\dom{u_1}$ by \cref{dom_intro}.
                                \end{byCase}
                            \end{subproof}
                            Thus $\{0,1\}\subseteq \dom{u_1}$ by \cref{subseteq}.
                        \end{subproof}
                        Thus $\dom{u_1} = \{0,1\}$ by \cref{subseteq_antisymmetric}.
                    \end{subproof}
                    Show that $\{0,1\}\union \img{s}{\initialSegment{1}} = \{0,1\}$.
                    \begin{subproof}
                        Thus $\img{s}{\initialSegment{1}} = \{1\}$.
                        Thus $\img{s}{\initialSegment{1}}\subseteq \{1\}$ by \cref{subseteq_refl}.
                        $1\in\{0,1\}$ by \cref{upair_intro_right}.
                        Thus $\{1\}\subseteq \{0,1\}$ by \cref{singleton_subset_intro}.
                        Thus $\img{s}{\initialSegment{1}}\subseteq \{0,1\}$ by \cref{subseteq_transitive}.
                        Thus $\{0,1\}\union \img{s}{\initialSegment{1}}\subseteq \{0,1\}$ by \cref{union_subsets_is_subset,subseteq_refl}.
                        Thus $\{0,1\}\subseteq \{0,1\}\union \img{s}{\initialSegment{1}}$ by \cref{union_intro_left,subseteq}.
                        Thus $\{0,1\}\union \img{s}{\initialSegment{1}} = \{0,1\}$ by \cref{subseteq_antisymmetric}.
                    \end{subproof}
                    Thus $\dom{u_1} = \{0,1\}\union \img{s}{\initialSegment{1}}$ by \cref{subseteq_antisymmetric}.
                \end{subproof}
                Show that for all $p\in\dom{u_1}$ we have $u_1(p)\in\pow{X}$.
                \begin{subproof}
                    Fix $p$.
                    Assume $p\in\dom{u_1}$.
                    Thus $p\in\{0,1\}$.
                    Thus $p = 0$ or $p = 1$ by \cref{upair_elim}.
                    \begin{byCase}
                    \caseOf{$p = 0$.}
                        Thus $(0,U_0)\in u_1$ by \cref{upair_intro_left}.
                        Thus $u_1(p) = U_0$ by \cref{function_apply_intro}.
                        $U_0$ is open in $X$ with respect to $T$.
                        Thus $T$ is a topology on $X$ and $U_0\in T$ by \cref{open_in}.
                        Thus $T\subseteq\pow{X}$ by \cref{topology_on}.
                        Thus $u_1(p)\in\pow{X}$ by \cref{elem_subseteq}.
                    \caseOf{$p = 1$.}
                        Thus $(1,U_1)\in u_1$ by \cref{upair_intro_right}.
                        Thus $u_1(p) = U_1$ by \cref{function_apply_intro}.
                        $U_1$ is open in $X$ with respect to $T$.
                        Thus $T$ is a topology on $X$ and $U_1\in T$ by \cref{open_in}.
                        Thus $T\subseteq\pow{X}$ by \cref{topology_on}.
                        Thus $u_1(p)\in\pow{X}$ by \cref{elem_subseteq}.
                    \end{byCase}
                \end{subproof}
                Thus $u_1\in\funs{\{0,1\}\union \img{s}{\initialSegment{1}}}{\pow{X}}$ by \cref{funs_intro}.
            \end{subproof}
            Show that for all $p\in\dom{u_1}$ we have
                $u_1(p)\in T$ and
                $u_1(1) = X\setminus B$ and
                $A\subseteq u_1(0)$ and
                $\closure{u_1(0)}{X}{T}\subseteq u_1(1)$ and
                for all $p\in\dom{u_1}$ for all $q\in\dom{u_1}$ if $p < q$ then
                    $\closure{u_1(p)}{X}{T}\subseteq u_1(q)$ and
                for all $p\in\dom{u_1}$ if $1 < p$ then $u_1(p) = X$.
            \begin{subproof}
                Fix $p$.
                Assume $p\in\dom{u_1}$.
                Thus $u_1$ is a function by \cref{funs_is_function}.
                Thus $u_1$ is right-unique and $u_1$ is a relation.
                Show that $u_1(p)\in T$.
                \begin{subproof}
                    Thus $p\in\{0,1\}$.
                    Thus $p = 0$ or $p = 1$ by \cref{upair_elim}.
                    \begin{byCase}
                    \caseOf{$p = 0$.}
                        Thus $(0,U_0)\in u_1$ by \cref{upair_intro_left}.
                        Thus $u_1(p) = U_0$ by \cref{function_apply_intro}.
                        $U_0$ is open in $X$ with respect to $T$.
                        Thus $u_1(p)\in T$ by \cref{open_in}.
                    \caseOf{$p = 1$.}
                        Thus $(1,U_1)\in u_1$ by \cref{upair_intro_right}.
                        Thus $u_1(p) = U_1$ by \cref{function_apply_intro}.
                        $U_1$ is open in $X$ with respect to $T$.
                        Thus $u_1(p)\in T$ by \cref{open_in}.
                    \end{byCase}
                \end{subproof}
                Show that $u_1(1) = X\setminus B$.
                \begin{subproof}
                    Thus $(1,U_1)\in u_1$ by \cref{upair_intro_right}.
                    Thus $u_1(1) = U_1$ by \cref{function_apply_intro}.
                    Thus $u_1(1) = X\setminus B$.
                \end{subproof}
                Show that $A\subseteq u_1(0)$.
                \begin{subproof}
                    Thus $(0,U_0)\in u_1$ by \cref{upair_intro_left}.
                    Thus $u_1(0) = U_0$ by \cref{function_apply_intro}.
                    Thus $A\subseteq u_1(0)$.
                \end{subproof}
                Show that $\closure{u_1(0)}{X}{T}\subseteq u_1(1)$.
                \begin{subproof}
                    Thus $(0,U_0)\in u_1$ by \cref{upair_intro_left}.
                    Thus $u_1(0) = U_0$ by \cref{function_apply_intro}.
                    Thus $(1,U_1)\in u_1$ by \cref{upair_intro_right}.
                    Thus $u_1(1) = U_1$ by \cref{function_apply_intro}.
                    Thus $\closure{u_1(0)}{X}{T}\subseteq u_1(1)$ by \cref{subseteq}.
                \end{subproof}
                Show that for all $p\in\dom{u_1}$ for all $q\in\dom{u_1}$ if $p < q$ then
                    $\closure{u_1(p)}{X}{T}\subseteq u_1(q)$.
                \begin{subproof}
                    Fix $p\in\dom{u_1}$.
                    Fix $q\in\dom{u_1}$.
                    Assume $p < q$.
                    Thus $p\in\{0,1\}$ and $q\in\{0,1\}$.
                    Show that $p = 0$.
                    \begin{subproof}
                        Suppose not.
                        Thus $p = 1$ by \cref{upair_elim}.
                        Thus $q = 0$ or $q = 1$ by \cref{upair_elim}.
                        \begin{byCase}
                        \caseOf{$q = 0$.}
                            $0\in\reals$ and $1\in\reals$ by \cref{real_add_ident,real_naturals_subset_reals,real_one_in_naturals,subseteq}.
                            $0 < 1$ by \cref{real_zero_lt_one}.
                            Thus $1 < q$ by \cref{real_lt_subst_left}.
                            Thus $1 < 0$ by \cref{real_lt_subst_right}.
                            Thus $1 < 1$ by \cref{real_lt_trans}.
                            Contradiction by \cref{real_lt_irreflexive}.
                        \caseOf{$q = 1$.}
                            Thus $1 < q$ by \cref{real_lt_subst_left}.
                            Thus $1 < 1$ by \cref{real_lt_subst_right}.
                            Thus $1\in\reals$ by \cref{real_naturals_subset_reals,real_one_in_naturals,subseteq}.
                            Contradiction by \cref{real_lt_irreflexive}.
                        \end{byCase}
                    \end{subproof}
                    Show that $q = 1$.
                    \begin{subproof}
                        Suppose not.
                        Thus $q = 0$ by \cref{upair_elim}.
                        Thus $p = 0$.
                        Thus $0\in\reals$ by \cref{real_add_ident}.
                        Thus $p < 0$ by \cref{real_lt_subst_right}.
                        Thus $0 < 0$ by \cref{real_lt_subst_left}.
                        Contradiction by \cref{real_lt_irreflexive}.
                    \end{subproof}
                    Thus $(p,U_0) = (0,U_0)$ by \cref{pair_eq_iff}.
                    Thus $(q,U_1) = (1,U_1)$ by \cref{pair_eq_iff}.
                    Thus $(p,U_0)\in u_1$ and $(q,U_1)\in u_1$ by \cref{upair_intro_left,upair_intro_right}.
                    Thus $u_1$ is a function.
                    Thus $u_1$ is right-unique and $u_1$ is a relation.
                    Thus $u_1(p) = U_0$ and $u_1(q) = U_1$ by \cref{function_apply_intro}.
                    Thus $\closure{u_1(p)}{X}{T}\subseteq u_1(q)$ by \cref{subseteq}.
                \end{subproof}
                Show that for all $p\in\dom{u_1}$ if $1 < p$ then $u_1(p) = X$.
                \begin{subproof}
                    Fix $p$.
                    Assume $p\in\dom{u_1}$.
                    Show that if $1 < p$ then $u_1(p) = X$.
                    \begin{subproof}
                        Assume $1 < p$.
                        Suppose not.
                        Thus $u_1(p) \neq X$.
                        Thus $p\in\{0,1\}$.
                        Then $p = 0$ or $p = 1$ by \cref{upair_elim}.
                        \begin{byCase}
                        \caseOf{$p = 0$.}
                            $0\in\reals$ and $1\in\reals$ by \cref{real_add_ident,real_naturals_subset_reals,real_one_in_naturals,subseteq}.
                            $0 < 1$ by \cref{real_zero_lt_one}.
                            Thus $1 < 0$ by \cref{real_lt_subst_right}.
                            Thus $1 < 1$ by \cref{real_lt_trans}.
                            Contradiction by \cref{real_lt_irreflexive}.
                        \caseOf{$p = 1$.}
                            Thus $1 < 1$ by \cref{real_lt_subst_right}.
                            Thus $1\in\reals$ by \cref{real_naturals_subset_reals,real_one_in_naturals,subseteq}.
                            Contradiction by \cref{real_lt_irreflexive}.
                        \end{byCase}
                    \end{subproof}
                \end{subproof}
            \end{subproof}
            Thus $u_1\in\urysohnAdmissible{s}{X}{T}{A}{B}{1}$ by \cref{urysohn_admissible}.
        \end{subproof}
        $1\in\naturalsPlus$ by \cref{real_one_in_naturals}.
        $1\in\naturalsPlus$ by \cref{real_one_in_naturals}.
        Thus $(1,u_1)\in G$ by \cref{urysohn_good_pair_intro}.
    \end{subproof}
    Thus there exists $u_1$ such that $(1,u_1)\in G$.
\end{proof}
% from §33 helper lemma 9: initial good pair when $1 < r_1$
\begin{lemma}\label{urysohn_u1_gt1}
    Assume $U_1 = X\setminus B$.
    Assume $U_1$ is open in $X$ with respect to $T$.
    Assume $U_0$ is open in $X$ with respect to $T$.
    Assume $A\subseteq U_0$.
    Assume $\closure{U_0}{X}{T}\subseteq U_1$.
    Assume $r_1\in\reals$.
    Assume $1 < r_1$.
    Assume $\img{s}{\initialSegment{1}} = \{r_1\}$.
    Let $G = \urysohnGoodPairs{s}{X}{T}{A}{B}$.
    Then there exists $u_1$ such that $(1,u_1)\in G$.
\end{lemma}
\begin{proof}
    Let $U_r = X$.
    Let $u_1 = \{(0,U_0),(1,U_1)\}\union\{(r_1,U_r)\}$.
    Show that $(1,u_1)\in G$.
    \begin{subproof}
        Show that $u_1\in\urysohnAdmissible{s}{X}{T}{A}{B}{1}$.
        \begin{subproof}
            Show that $u_1\in\funs{\{0,1\}\union \img{s}{\initialSegment{1}}}{\pow{X}}$.
            \begin{subproof}
                Show that $r_1 > 0$.
                \begin{subproof}
                    Suppose not.
                    Thus $r_1 = 0$.
                    Thus $1 < 0$ by \cref{real_lt_subst_right}.
                    $0 < 1$ by \cref{real_zero_lt_one}.
                    $0\in\reals$ by \cref{real_add_ident}.
                    $1\in\reals$ by \cref{real_mul_ident}.
                    Thus $1 < 1$ by \cref{real_lt_trans}.
                    Contradiction by \cref{real_lt_irreflexive}.
                \end{subproof}
                Show that $u_1$ is right-unique.
                \begin{subproof}
                    Show that for all $a,b,b'$ such that $(a,b)\in u_1$ and $(a,b')\in u_1$ we have $b = b'$.
                    \begin{subproof}
                        Fix $a$.
                        Fix $b$.
                        Fix $b'$.
                        Assume $(a,b)\in u_1$ and $(a,b')\in u_1$.
                        Thus $(a,b)\in\{(0,U_0),(1,U_1)\}$ or $(a,b)\in\{(r_1,U_r)\}$ by \cref{union_iff}.
                        Thus $(a,b) = (0,U_0)$ or $(a,b) = (1,U_1)$ or $(a,b) = (r_1,U_r)$
                            by \cref{upair_elim,singleton_elim}.
                        \begin{byCase}
                        \caseOf{$(a,b) = (0,U_0)$.}
                            Thus $a = 0$ by \cref{pair_eq_iff}.
                            Show that $b = b'$.
                            \begin{subproof}
                                Suppose not.
                                Thus $b \neq b'$.
                                Thus $(a,b')\in\{(0,U_0),(1,U_1)\}$ or $(a,b')\in\{(r_1,U_r)\}$ by \cref{union_iff}.
                                \begin{byCase}
                                \caseOf{$(a,b')\in\{(0,U_0),(1,U_1)\}$.}
                                    Thus $(a,b') = (0,U_0)$ or $(a,b') = (1,U_1)$ by \cref{upair_elim}.
                                    \begin{byCase}
                                    \caseOf{$(a,b') = (0,U_0)$.}
                                        Thus $b = b'$ by \cref{pair_eq_iff}.
                                        Contradiction.
                                    \caseOf{$(a,b') = (1,U_1)$.}
                                        Thus $a = 1$ by \cref{pair_eq_iff}.
                                        Thus $1\neq 0$ by \cref{real_mul_ident}.
                                        Contradiction.
                                    \end{byCase}
                                \caseOf{$(a,b')\in\{(r_1,U_r)\}$.}
                                    Thus $(a,b') = (r_1,U_r)$ by \cref{singleton_elim}.
                                    Thus $a = r_1$ by \cref{pair_eq_iff}.
                                    Thus $r_1 \neq 0$.
                                    Contradiction.
                                \end{byCase}
                            \end{subproof}
                        \caseOf{$(a,b) = (1,U_1)$.}
                            Thus $a = 1$ by \cref{pair_eq_iff}.
                            Show that $b = b'$.
                            \begin{subproof}
                                Suppose not.
                                Thus $b \neq b'$.
                                Thus $(a,b')\in\{(0,U_0),(1,U_1)\}$ or $(a,b')\in\{(r_1,U_r)\}$ by \cref{union_iff}.
                                \begin{byCase}
                                \caseOf{$(a,b')\in\{(0,U_0),(1,U_1)\}$.}
                                    Thus $(a,b') = (0,U_0)$ or $(a,b') = (1,U_1)$ by \cref{upair_elim}.
                                    \begin{byCase}
                                    \caseOf{$(a,b') = (1,U_1)$.}
                                        Thus $b = b'$ by \cref{pair_eq_iff}.
                                        Contradiction.
                                    \caseOf{$(a,b') = (0,U_0)$.}
                                        Thus $a = 0$ by \cref{pair_eq_iff}.
                                        Thus $1\neq 0$ by \cref{real_mul_ident}.
                                        Contradiction.
                                    \end{byCase}
                                \caseOf{$(a,b')\in\{(r_1,U_r)\}$.}
                                    Thus $(a,b') = (r_1,U_r)$ by \cref{singleton_elim}.
                                    Thus $a = r_1$ by \cref{pair_eq_iff}.
                                    Thus $r_1 \neq 1$.
                                    Contradiction.
                                \end{byCase}
                            \end{subproof}
                        \caseOf{$(a,b) = (r_1,U_r)$.}
                            Thus $a = r_1$ by \cref{pair_eq_iff}.
                            Show that $b = b'$.
                            \begin{subproof}
                                Suppose not.
                                Thus $b \neq b'$.
                                Thus $(a,b')\in\{(0,U_0),(1,U_1)\}$ or $(a,b')\in\{(r_1,U_r)\}$ by \cref{union_iff}.
                                \begin{byCase}
                                \caseOf{$(a,b')\in\{(r_1,U_r)\}$.}
                                    Thus $(a,b') = (r_1,U_r)$ by \cref{singleton_elim}.
                                    Thus $b = b'$ by \cref{pair_eq_iff}.
                                    Contradiction.
                                \caseOf{$(a,b')\in\{(0,U_0),(1,U_1)\}$.}
                                    Thus $(a,b') = (0,U_0)$ or $(a,b') = (1,U_1)$ by \cref{upair_elim}.
                                    \begin{byCase}
                                    \caseOf{$(a,b') = (0,U_0)$.}
                                        Thus $a = 0$ by \cref{pair_eq_iff}.
                                        Thus $r_1 \neq 0$.
                                        Contradiction.
                                    \caseOf{$(a,b') = (1,U_1)$.}
                                        Thus $a = 1$ by \cref{pair_eq_iff}.
                                        Thus $r_1 \neq 1$.
                                        Contradiction.
                                    \end{byCase}
                                \end{byCase}
                            \end{subproof}
                        \end{byCase}
                    \end{subproof}
                    Thus $u_1$ is right-unique by \cref{rightunique}.
                \end{subproof}
                Show that $u_1$ is a relation.
                \begin{subproof}
                    Show that for all $w$ such that $w\in u_1$ there exist $x,y$ such that $w = (x,y)$.
                    \begin{subproof}
                        Fix $w$.
                        Assume $w\in u_1$.
                        Thus $w\in\{(0,U_0),(1,U_1)\}$ or $w\in\{(r_1,U_r)\}$ by \cref{union_iff}.
                        Thus $w = (0,U_0)$ or $w = (1,U_1)$ or $w = (r_1,U_r)$
                            by \cref{upair_elim,singleton_elim}.
                        Thus there exist $x,y$ such that $w = (x,y)$.
                    \end{subproof}
                    Thus $u_1$ is a relation by \cref{relation}.
                \end{subproof}
                Thus $u_1$ is a function.
                Show that $\dom{u_1} = \{0,1\}\union \img{s}{\initialSegment{1}}$.
                \begin{subproof}
                    Show that $\dom{u_1} = \{0,1\}\union\{r_1\}$.
                    \begin{subproof}
                        Show that $\dom{u_1}\subseteq \{0,1\}\union\{r_1\}$.
                        \begin{subproof}
                            Show that for all $x\in\dom{u_1}$ we have $x\in\{0,1\}\union\{r_1\}$.
                            \begin{subproof}
                                Fix $x$.
                                Assume $x\in\dom{u_1}$.
                                Take $y$ such that $(x,y)\in u_1$ by \cref{dom_iff}.
                                Thus $(x,y)\in\{(0,U_0),(1,U_1)\}$ or $(x,y)\in\{(r_1,U_r)\}$ by \cref{union_iff}.
                                \begin{byCase}
                                \caseOf{$(x,y)\in\{(0,U_0),(1,U_1)\}$.}
                                    Thus $(x,y) = (0,U_0)$ or $(x,y) = (1,U_1)$ by \cref{upair_elim}.
                                    Thus $x = 0$ or $x = 1$ by \cref{pair_eq_iff}.
                                    Thus $x\in\{0,1\}$ by \cref{upair_intro_left,upair_intro_right}.
                                    Thus $x\in\{0,1\}\union\{r_1\}$ by \cref{union_intro_left}.
                                \caseOf{$(x,y)\in\{(r_1,U_r)\}$.}
                                    Thus $(x,y) = (r_1,U_r)$ by \cref{singleton_elim}.
                                    Thus $x = r_1$ by \cref{pair_eq_iff}.
                                    Thus $x\in\{r_1\}$ by \cref{singleton_intro}.
                                    Thus $x\in\{0,1\}\union\{r_1\}$ by \cref{union_intro_right}.
                                \end{byCase}
                            \end{subproof}
                            Thus $\dom{u_1}\subseteq \{0,1\}\union\{r_1\}$ by \cref{subseteq}.
                        \end{subproof}
                        Show that $\{0,1\}\union\{r_1\}\subseteq \dom{u_1}$.
                        \begin{subproof}
                            Show that for all $x\in\{0,1\}\union\{r_1\}$ we have $x\in\dom{u_1}$.
                            \begin{subproof}
                                Fix $x$.
                                Assume $x\in\{0,1\}\union\{r_1\}$.
                                Thus $x\in\{0,1\}$ or $x\in\{r_1\}$ by \cref{union_iff}.
                                \begin{byCase}
                                \caseOf{$x\in\{0,1\}$.}
                                    Thus $x = 0$ or $x = 1$ by \cref{upair_elim}.
                                    \begin{byCase}
                                    \caseOf{$x = 0$.}
                                        Thus $(0,U_0)\in\{(0,U_0),(1,U_1)\}$ by \cref{upair_intro_left}.
                                        Thus $(0,U_0)\in u_1$ by \cref{union_intro_left}.
                                        Thus $x\in\dom{u_1}$ by \cref{dom_intro}.
                                    \caseOf{$x = 1$.}
                                        Thus $(1,U_1)\in\{(0,U_0),(1,U_1)\}$ by \cref{upair_intro_right}.
                                        Thus $(1,U_1)\in u_1$ by \cref{union_intro_left}.
                                        Thus $x\in\dom{u_1}$ by \cref{dom_intro}.
                                    \end{byCase}
                                \caseOf{$x\in\{r_1\}$.}
                                    Thus $x = r_1$ by \cref{singleton_elim}.
                                    Thus $(r_1,U_r)\in\{(r_1,U_r)\}$ by \cref{singleton_intro}.
                                    Thus $(r_1,U_r)\in u_1$ by \cref{union_intro_right}.
                                    Thus $x\in\dom{u_1}$ by \cref{dom_intro}.
                                \end{byCase}
                            \end{subproof}
                            Thus $\{0,1\}\union\{r_1\}\subseteq \dom{u_1}$ by \cref{subseteq}.
                        \end{subproof}
                        Thus $\dom{u_1} = \{0,1\}\union\{r_1\}$ by \cref{subseteq_antisymmetric}.
                    \end{subproof}
                    Show that $\{0,1\}\union \img{s}{\initialSegment{1}} = \{0,1\}\union\{r_1\}$.
                    \begin{subproof}
                        Thus $\img{s}{\initialSegment{1}} = \{r_1\}$.
                        Thus $\img{s}{\initialSegment{1}}\subseteq \{r_1\}$ by \cref{subseteq_refl}.
                        $r_1\in\{r_1\}$ by \cref{singleton_intro}.
                        Thus $\{r_1\}\subseteq \{0,1\}\union\{r_1\}$ by \cref{union_intro_right,subseteq}.
                        Thus $\img{s}{\initialSegment{1}}\subseteq \{0,1\}\union\{r_1\}$ by \cref{subseteq_transitive}.
                        Thus $\{0,1\}\union \img{s}{\initialSegment{1}}\subseteq \{0,1\}\union\{r_1\}$ by \cref{union_subsets_is_subset,subseteq_refl}.
                        Thus $\{0,1\}\subseteq \{0,1\}\union \img{s}{\initialSegment{1}}$ by \cref{union_intro_left,subseteq}.
                        Thus $\{r_1\}\subseteq \{0,1\}\union \img{s}{\initialSegment{1}}$ by \cref{union_intro_right,subseteq}.
                        Thus $\{0,1\}\union\{r_1\}\subseteq \{0,1\}\union \img{s}{\initialSegment{1}}$ by \cref{union_subsets_is_subset}.
                        Thus $\{0,1\}\union \img{s}{\initialSegment{1}} = \{0,1\}\union\{r_1\}$ by \cref{subseteq_antisymmetric}.
                    \end{subproof}
                    Thus $\dom{u_1} = \{0,1\}\union \img{s}{\initialSegment{1}}$ by \cref{subseteq_antisymmetric}.
                \end{subproof}
                Show that for all $p\in\dom{u_1}$ we have $u_1(p)\in\pow{X}$.
                \begin{subproof}
                    Fix $p$.
                    Assume $p\in\dom{u_1}$.
                    Thus $p\in\{0,1\}\union\{r_1\}$.
                    Thus $p\in\{0,1\}$ or $p\in\{r_1\}$ by \cref{union_iff}.
                    \begin{byCase}
                    \caseOf{$p\in\{0,1\}$.}
                        Thus $p = 0$ or $p = 1$ by \cref{upair_elim}.
                        \begin{byCase}
                        \caseOf{$p = 0$.}
                            Thus $(0,U_0)\in\{(0,U_0),(1,U_1)\}$ by \cref{upair_intro_left}.
                            Thus $(0,U_0)\in u_1$ by \cref{union_intro_left}.
                            Thus $u_1(p) = U_0$ by \cref{function_apply_intro}.
                            $U_0$ is open in $X$ with respect to $T$.
                            Thus $T$ is a topology on $X$ and $U_0\in T$ by \cref{open_in}.
                            Thus $T\subseteq\pow{X}$ by \cref{topology_on}.
                            Thus $u_1(p)\in\pow{X}$ by \cref{elem_subseteq}.
                        \caseOf{$p = 1$.}
                            Thus $(1,U_1)\in\{(0,U_0),(1,U_1)\}$ by \cref{upair_intro_right}.
                            Thus $(1,U_1)\in u_1$ by \cref{union_intro_left}.
                            Thus $u_1(p) = U_1$ by \cref{function_apply_intro}.
                            $U_1$ is open in $X$ with respect to $T$.
                            Thus $T$ is a topology on $X$ and $U_1\in T$ by \cref{open_in}.
                            Thus $T\subseteq\pow{X}$ by \cref{topology_on}.
                            Thus $u_1(p)\in\pow{X}$ by \cref{elem_subseteq}.
                        \end{byCase}
                    \caseOf{$p\in\{r_1\}$.}
                        Thus $p = r_1$ by \cref{singleton_elim}.
                        Thus $(r_1,U_r)\in\{(r_1,U_r)\}$ by \cref{singleton_intro}.
                        Thus $(r_1,U_r)\in u_1$ by \cref{union_intro_right}.
                        Thus $u_1(p) = U_r$ by \cref{function_apply_intro}.
                        $U_r = X$.
                        $U_0$ is open in $X$ with respect to $T$.
                        Thus $T$ is a topology on $X$ by \cref{open_in}.
                        Thus $X\in T$ by \cref{topology_on}.
                        Thus $u_1(p)\in T$.
                        Thus $T\subseteq\pow{X}$ by \cref{topology_on}.
                        Thus $u_1(p)\in\pow{X}$ by \cref{elem_subseteq}.
                    \end{byCase}
                \end{subproof}
                Thus $u_1\in\funs{\{0,1\}\union \img{s}{\initialSegment{1}}}{\pow{X}}$ by \cref{funs_intro}.
            \end{subproof}
            Show that for all $p\in\dom{u_1}$ we have
                $u_1(p)\in T$ and
                $u_1(1) = X\setminus B$ and
                $A\subseteq u_1(0)$ and
                $\closure{u_1(0)}{X}{T}\subseteq u_1(1)$ and
                for all $p\in\dom{u_1}$ for all $q\in\dom{u_1}$ if $p < q$ then
                    $\closure{u_1(p)}{X}{T}\subseteq u_1(q)$ and
                for all $p\in\dom{u_1}$ if $1 < p$ then $u_1(p) = X$.
            \begin{subproof}
                Fix $p$.
                Assume $p\in\dom{u_1}$.
                Thus $u_1$ is a function by \cref{funs_is_function}.
                Thus $u_1$ is right-unique and $u_1$ is a relation.
                Show that $u_1(p)\in T$.
                \begin{subproof}
                    Thus $p\in\{0,1\}\union\{r_1\}$.
                    Thus $p\in\{0,1\}$ or $p\in\{r_1\}$ by \cref{union_iff}.
                    \begin{byCase}
                    \caseOf{$p\in\{0,1\}$.}
                        Thus $p = 0$ or $p = 1$ by \cref{upair_elim}.
                        \begin{byCase}
                        \caseOf{$p = 0$.}
                            Thus $(0,U_0)\in\{(0,U_0),(1,U_1)\}$ by \cref{upair_intro_left}.
                            Thus $(0,U_0)\in u_1$ by \cref{union_intro_left}.
                            Thus $u_1(p) = U_0$ by \cref{function_apply_intro}.
                            $U_0$ is open in $X$ with respect to $T$.
                            Thus $u_1(p)\in T$ by \cref{open_in}.
                        \caseOf{$p = 1$.}
                            Thus $(1,U_1)\in\{(0,U_0),(1,U_1)\}$ by \cref{upair_intro_right}.
                            Thus $(1,U_1)\in u_1$ by \cref{union_intro_left}.
                            Thus $u_1(p) = U_1$ by \cref{function_apply_intro}.
                            $U_1$ is open in $X$ with respect to $T$.
                            Thus $u_1(p)\in T$ by \cref{open_in}.
                        \end{byCase}
                    \caseOf{$p\in\{r_1\}$.}
                        Thus $p = r_1$ by \cref{singleton_elim}.
                        Thus $(r_1,U_r)\in\{(r_1,U_r)\}$ by \cref{singleton_intro}.
                        Thus $(r_1,U_r)\in u_1$ by \cref{union_intro_right}.
                        Thus $u_1(p) = U_r$ by \cref{function_apply_intro}.
                        $U_r = X$.
                        $U_0$ is open in $X$ with respect to $T$.
                        Thus $T$ is a topology on $X$ by \cref{open_in}.
                        Thus $X\in T$ by \cref{topology_on}.
                        Thus $u_1(p)\in T$.
                    \end{byCase}
                \end{subproof}
                Show that $u_1(1) = X\setminus B$.
                \begin{subproof}
                    Thus $(1,U_1)\in\{(0,U_0),(1,U_1)\}$ by \cref{upair_intro_right}.
                    Thus $(1,U_1)\in u_1$ by \cref{union_intro_left}.
                    Thus $u_1(1) = U_1$ by \cref{function_apply_intro}.
                    Thus $u_1(1) = X\setminus B$.
                \end{subproof}
                Show that $A\subseteq u_1(0)$.
                \begin{subproof}
                    Thus $(0,U_0)\in\{(0,U_0),(1,U_1)\}$ by \cref{upair_intro_left}.
                    Thus $(0,U_0)\in u_1$ by \cref{union_intro_left}.
                    Thus $u_1(0) = U_0$ by \cref{function_apply_intro}.
                    Thus $A\subseteq u_1(0)$.
                \end{subproof}
                Show that $\closure{u_1(0)}{X}{T}\subseteq u_1(1)$.
                \begin{subproof}
                    Thus $(0,U_0)\in\{(0,U_0),(1,U_1)\}$ by \cref{upair_intro_left}.
                    Thus $(0,U_0)\in u_1$ by \cref{union_intro_left}.
                    Thus $u_1(0) = U_0$ by \cref{function_apply_intro}.
                    Thus $(1,U_1)\in\{(0,U_0),(1,U_1)\}$ by \cref{upair_intro_right}.
                    Thus $(1,U_1)\in u_1$ by \cref{union_intro_left}.
                    Thus $u_1(1) = U_1$ by \cref{function_apply_intro}.
                    Thus $\closure{u_1(0)}{X}{T}\subseteq u_1(1)$ by \cref{subseteq}.
                \end{subproof}
                Show that for all $p\in\dom{u_1}$ for all $q\in\dom{u_1}$ if $p < q$ then
                    $\closure{u_1(p)}{X}{T}\subseteq u_1(q)$.
                \begin{subproof}
                    Fix $p\in\dom{u_1}$.
                    Fix $q\in\dom{u_1}$.
                    Assume $p < q$.
                    Thus $p\in\{0,1\}\union\{r_1\}$ and $q\in\{0,1\}\union\{r_1\}$.
                    Show that $q \neq 0$.
                    \begin{subproof}
                        Suppose not.
                        Thus $q = 0$.
                        Thus $p < 0$ by \cref{real_lt_subst_right}.
                        Thus $p\in\{0,1\}\union\{r_1\}$.
                        Thus $p\in\{0,1\}$ or $p\in\{r_1\}$ by \cref{union_iff}.
                        \begin{byCase}
                        \caseOf{$p\in\{0,1\}$.}
                            Thus $p = 0$ or $p = 1$ by \cref{upair_elim}.
                            \begin{byCase}
                            \caseOf{$p = 0$.}
                                Thus $0 < 0$ by \cref{real_lt_subst_left}.
                                $0\in\reals$ by \cref{real_add_ident}.
                                Contradiction by \cref{real_lt_irreflexive}.
                            \caseOf{$p = 1$.}
                                Thus $1 < 0$ by \cref{real_lt_subst_left}.
                                $0\in\reals$ by \cref{real_add_ident}.
                                $1\in\reals$ by \cref{real_mul_ident}.
                                $0 < 1$ by \cref{real_zero_lt_one}.
                                Thus $1 < 1$ by \cref{real_lt_trans}.
                                Contradiction by \cref{real_lt_irreflexive}.
                            \end{byCase}
                        \caseOf{$p\in\{r_1\}$.}
                            Thus $p = r_1$ by \cref{singleton_elim}.
                            Thus $r_1 < 0$ by \cref{real_lt_subst_left}.
                            $0\in\reals$ by \cref{real_add_ident}.
                            $1\in\reals$ by \cref{real_mul_ident}.
                            Thus $1 < 0$ by \cref{real_lt_trans}.
                            $0 < 1$ by \cref{real_zero_lt_one}.
                            Thus $1 < 1$ by \cref{real_lt_trans}.
                            Contradiction by \cref{real_lt_irreflexive}.
                        \end{byCase}
                    \end{subproof}
                    Thus $q\in\{1\}\union\{r_1\}$.
                    Thus $q\in\{1\}$ or $q\in\{r_1\}$ by \cref{union_iff}.
                    \begin{byCase}
                    \caseOf{$q\in\{r_1\}$.}
                        Thus $q = r_1$ by \cref{singleton_elim}.
                        Thus $(q,U_r)\in\{(r_1,U_r)\}$ by \cref{singleton_intro}.
                        Thus $(q,U_r)\in u_1$ by \cref{union_intro_right}.
                        Thus $u_1(q) = U_r$ by \cref{function_apply_intro}.
                        $U_r = X$.
                        Show that $\closure{u_1(p)}{X}{T}\subseteq u_1(q)$.
                        \begin{subproof}
                            $U_0$ is open in $X$ with respect to $T$.
                            Thus $T$ is a topology on $X$ by \cref{open_in}.
                            Thus $\emptyset\in T$ by \cref{topology_on}.
                            Thus $X\setminus X = \emptyset$ by \cref{setminus_self}.
                            Thus $X\subseteq X$ by \cref{subseteq_refl}.
                            Thus $\emptyset$ is open in $X$ with respect to $T$ by \cref{open_in}.
                            Thus $X$ is closed in $X$ with respect to $T$ by \cref{closed_in}.
                            Thus $T\subseteq\pow{X}$ by \cref{topology_on}.
                            Thus $u_1(p)\in\pow{X}$ by \cref{elem_subseteq}.
                            Thus $u_1(p)\subseteq X$ by \cref{pow_iff}.
                            Thus $\closure{u_1(p)}{X}{T}\subseteq X$ by \cref{closure_subseteq_closed}.
                            Thus $\closure{u_1(p)}{X}{T}\subseteq u_1(q)$ by \cref{subseteq}.
                        \end{subproof}
                    \caseOf{$q\in\{1\}$.}
                        Thus $q = 1$ by \cref{singleton_elim}.
                        Show that $p = 0$.
                        \begin{subproof}
                            Suppose not.
                            Thus $p\in\{0,1\}\union\{r_1\}$.
                            Thus $p\in\{0,1\}$ or $p\in\{r_1\}$ by \cref{union_iff}.
                            \begin{byCase}
                            \caseOf{$p\in\{0,1\}$.}
                                Thus $p = 0$ or $p = 1$ by \cref{upair_elim}.
                                \begin{byCase}
                                \caseOf{$p = 1$.}
                                    Thus $1 < 1$ by \cref{real_lt_subst_left}.
                                    $1\in\reals$ by \cref{real_mul_ident}.
                                    Contradiction by \cref{real_lt_irreflexive}.
                                \caseOf{$p = 0$.}
                                    Contradiction.
                                \end{byCase}
                            \caseOf{$p\in\{r_1\}$.}
                                Thus $p = r_1$ by \cref{singleton_elim}.
                                Thus $r_1 < 1$ by \cref{real_lt_subst_left}.
                                Contradiction.
                            \end{byCase}
                        \end{subproof}
                        Thus $(p,U_0) = (0,U_0)$ by \cref{pair_eq_iff}.
                        Thus $(q,U_1) = (1,U_1)$ by \cref{pair_eq_iff}.
                        Thus $(p,U_0)\in u_1$ by \cref{union_intro_left,upair_intro_left}.
                        Thus $(q,U_1)\in u_1$ by \cref{union_intro_left,upair_intro_right}.
                        Thus $u_1(p) = U_0$ and $u_1(q) = U_1$ by \cref{function_apply_intro}.
                        Thus $\closure{u_1(p)}{X}{T}\subseteq u_1(q)$ by \cref{subseteq}.
                    \end{byCase}
                \end{subproof}
                Show that for all $p\in\dom{u_1}$ if $1 < p$ then $u_1(p) = X$.
                \begin{subproof}
                    Fix $p\in\dom{u_1}$.
                    Show that if $1 < p$ then $u_1(p) = X$.
                    \begin{subproof}
                        Assume $1 < p$.
                        Suppose not.
                        Thus $u_1(p) \neq X$.
                        Thus $p\in\{0,1\}\union\{r_1\}$.
                        Thus $p\in\{0,1\}$ or $p\in\{r_1\}$ by \cref{union_iff}.
                        \begin{byCase}
                        \caseOf{$p\in\{0,1\}$.}
                            Thus $p = 0$ or $p = 1$ by \cref{upair_elim}.
                            \begin{byCase}
                            \caseOf{$p = 0$.}
                                Thus $1 < 0$ by \cref{real_lt_subst_left}.
                                $0\in\reals$ by \cref{real_add_ident}.
                                $1\in\reals$ by \cref{real_mul_ident}.
                                $0 < 1$ by \cref{real_zero_lt_one}.
                                Thus $1 < 1$ by \cref{real_lt_trans}.
                                Contradiction by \cref{real_lt_irreflexive}.
                            \caseOf{$p = 1$.}
                                Thus $1 < 1$ by \cref{real_lt_subst_left}.
                                $1\in\reals$ by \cref{real_mul_ident}.
                                Contradiction by \cref{real_lt_irreflexive}.
                            \end{byCase}
                        \caseOf{$p\in\{r_1\}$.}
                            Thus $p = r_1$ by \cref{singleton_elim}.
                            Thus $1 < r_1$ by \cref{real_lt_subst_left}.
                            Contradiction.
                        \end{byCase}
                    \end{subproof}
                \end{subproof}
            \end{subproof}
            Thus $u_1\in\urysohnAdmissible{s}{X}{T}{A}{B}{1}$ by \cref{urysohn_admissible}.
        \end{subproof}
        $1\in\naturalsPlus$ by \cref{real_one_in_naturals}.
        Thus $(1,u_1)\in G$ by \cref{urysohn_good_pair_intro}.
    \end{subproof}
\end{proof}

% from §33 helper lemma 12: unit rational is positive
\begin{lemma}\label{unit_rational_pos}
    Assume $r_1\in\unitRationals$.
    Then $0 < r_1$.
\end{lemma}
\begin{proof}
    Take $q$ such that
        $q\in\naturalsPlus\times\naturalsPlus$ and
        $r_1 = \fst{q} \rmul \inv{\snd{q}}$
        by \cref{unit_rationals}.
    Take $m,n$ such that $q = (m,n)$ and $m\in\naturalsPlus$ and $n\in\naturalsPlus$
        by \cref{times}.
    Thus $\fst{q} = m$ and $\snd{q} = n$ by \cref{fst_eq,snd_eq}.
    Thus $m\in\reals$ and $n\in\reals$ by \cref{real_naturals_subset_reals,subseteq}.
    Show that $m > 0$.
    \begin{subproof}
        $1\in\reals$ by \cref{real_mul_ident}.
        Thus $1 < m$ or $1 = m$ by \cref{real_one_leq_naturalsplus}.
        \begin{byCase}
        \caseOf{$1 < m$.}
            $0\in\reals$ by \cref{real_add_ident}.
            $0 < 1$ by \cref{real_zero_lt_one}.
            Thus $0 < m$ by \cref{real_lt_trans}.
        \caseOf{$1 = m$.}
            $0 < 1$ by \cref{real_zero_lt_one}.
            Thus $0 < m$ by \cref{real_lt_subst_right}.
        \end{byCase}
    \end{subproof}
    Show that $n > 0$.
    \begin{subproof}
        $1\in\reals$ by \cref{real_mul_ident}.
        Thus $1 < n$ or $1 = n$ by \cref{real_one_leq_naturalsplus}.
        \begin{byCase}
        \caseOf{$1 < n$.}
            $0\in\reals$ by \cref{real_add_ident}.
            $0 < 1$ by \cref{real_zero_lt_one}.
            Thus $0 < n$ by \cref{real_lt_trans}.
        \caseOf{$1 = n$.}
            $0 < 1$ by \cref{real_zero_lt_one}.
            Thus $0 < n$ by \cref{real_lt_subst_right}.
        \end{byCase}
    \end{subproof}
    Thus $n \neq 0$ by \cref{real_pos_nonzero}.
    Thus $\inv{n}\in\reals$ by \cref{real_mul_inverse}.
    Thus $\inv{n} > 0$ by \cref{real_inv_pos}.
    Thus $r_1 = m \rmul \inv{n}$ by \cref{fst_eq,snd_eq}.
    Thus $m \rmul \inv{n} > 0$ by \cref{real_mul_pos_pos}.
    Thus $r_1 > 0$ by \cref{real_lt_subst_right}.
\end{proof}

% from §33 helper lemma 10: admissible good pair for $r_1 < 1$ given $U_r$
\begin{lemma}\label{urysohn_u1_lt1_admissible}
    Assume $U_1 = X\setminus B$.
    Assume $U_1$ is open in $X$ with respect to $T$.
    Assume $U_0$ is open in $X$ with respect to $T$.
    Assume $A\subseteq U_0$.
    Assume $\closure{U_0}{X}{T}\subseteq U_1$.
    Assume $r_1\in\reals$.
    Assume $r_1\in\unitRationals$.
    Assume $r_1 < 1$.
    Assume $\img{s}{\initialSegment{1}} = \{r_1\}$.
    Assume $U_r$ is open in $X$ with respect to $T$.
    Assume $\closure{U_0}{X}{T}\subseteq U_r$.
    Assume $\closure{U_r}{X}{T}\subseteq U_1$.
    Let $G = \urysohnGoodPairs{s}{X}{T}{A}{B}$.
    Then there exists $u_1$ such that $(1,u_1)\in G$.
\end{lemma}
\begin{proof}
    Let $u_1 = \{(0,U_0),(1,U_1)\}\union\{(r_1,U_r)\}$.
    Show that $(1,u_1)\in G$.
    \begin{subproof}
        Show that $u_1\in\urysohnAdmissible{s}{X}{T}{A}{B}{1}$.
        \begin{subproof}
            Show that $u_1\in\funs{\{0,1\}\union \img{s}{\initialSegment{1}}}{\pow{X}}$.
            \begin{subproof}
                Thus $0 < r_1$ by \cref{unit_rational_pos}.
                Thus $r_1 \neq 0$ by \cref{real_pos_nonzero}.
                Show that $u_1$ is right-unique.
                \begin{subproof}
                    Show that for all $a,b,b'$ such that $(a,b)\in u_1$ and $(a,b')\in u_1$ we have $b = b'$.
                    \begin{subproof}
                        Fix $a$.
                        Fix $b$.
                        Fix $b'$.
                        Assume $(a,b)\in u_1$ and $(a,b')\in u_1$.
                        Thus $(a,b)\in\{(0,U_0),(1,U_1)\}$ or $(a,b)\in\{(r_1,U_r)\}$ by \cref{union_iff}.
                        Thus $(a,b) = (0,U_0)$ or $(a,b) = (1,U_1)$ or $(a,b) = (r_1,U_r)$
                            by \cref{upair_elim,singleton_elim}.
                        \begin{byCase}
                        \caseOf{$(a,b) = (0,U_0)$.}
                            Thus $a = 0$ by \cref{pair_eq_iff}.
                            Show that $b = b'$.
                            \begin{subproof}
                                Suppose not.
                                Thus $b \neq b'$.
                                Thus $(a,b')\in\{(0,U_0),(1,U_1)\}$ or $(a,b')\in\{(r_1,U_r)\}$ by \cref{union_iff}.
                                \begin{byCase}
                                \caseOf{$(a,b')\in\{(0,U_0),(1,U_1)\}$.}
                                    Thus $(a,b') = (0,U_0)$ or $(a,b') = (1,U_1)$ by \cref{upair_elim}.
                                    \begin{byCase}
                                    \caseOf{$(a,b') = (0,U_0)$.}
                                        Thus $b = b'$ by \cref{pair_eq_iff}.
                                        Contradiction.
                                    \caseOf{$(a,b') = (1,U_1)$.}
                                        Thus $a = 1$ by \cref{pair_eq_iff}.
                                        Thus $1\neq 0$ by \cref{real_mul_ident}.
                                        Contradiction.
                                    \end{byCase}
                                \caseOf{$(a,b')\in\{(r_1,U_r)\}$.}
                                    Thus $(a,b') = (r_1,U_r)$ by \cref{singleton_elim}.
                                    Thus $a = r_1$ by \cref{pair_eq_iff}.
                                    Thus $r_1 \neq 0$.
                                    Contradiction.
                                \end{byCase}
                            \end{subproof}
                        \caseOf{$(a,b) = (1,U_1)$.}
                            Thus $a = 1$ by \cref{pair_eq_iff}.
                            Show that $b = b'$.
                            \begin{subproof}
                                Suppose not.
                                Thus $b \neq b'$.
                                Thus $(a,b')\in\{(0,U_0),(1,U_1)\}$ or $(a,b')\in\{(r_1,U_r)\}$ by \cref{union_iff}.
                                \begin{byCase}
                                \caseOf{$(a,b')\in\{(0,U_0),(1,U_1)\}$.}
                                    Thus $(a,b') = (0,U_0)$ or $(a,b') = (1,U_1)$ by \cref{upair_elim}.
                                    \begin{byCase}
                                    \caseOf{$(a,b') = (1,U_1)$.}
                                        Thus $b = b'$ by \cref{pair_eq_iff}.
                                        Contradiction.
                                    \caseOf{$(a,b') = (0,U_0)$.}
                                        Thus $a = 0$ by \cref{pair_eq_iff}.
                                        Thus $1\neq 0$ by \cref{real_mul_ident}.
                                        Contradiction.
                                    \end{byCase}
                                \caseOf{$(a,b')\in\{(r_1,U_r)\}$.}
                                    Thus $(a,b') = (r_1,U_r)$ by \cref{singleton_elim}.
                                    Thus $a = r_1$ by \cref{pair_eq_iff}.
                                    Thus $r_1 \neq 1$.
                                    Contradiction.
                                \end{byCase}
                            \end{subproof}
                        \caseOf{$(a,b) = (r_1,U_r)$.}
                            Thus $a = r_1$ by \cref{pair_eq_iff}.
                            Show that $b = b'$.
                            \begin{subproof}
                                Suppose not.
                                Thus $b \neq b'$.
                                Thus $(a,b')\in\{(0,U_0),(1,U_1)\}$ or $(a,b')\in\{(r_1,U_r)\}$ by \cref{union_iff}.
                                \begin{byCase}
                                \caseOf{$(a,b')\in\{(r_1,U_r)\}$.}
                                    Thus $(a,b') = (r_1,U_r)$ by \cref{singleton_elim}.
                                    Thus $b = b'$ by \cref{pair_eq_iff}.
                                    Contradiction.
                                \caseOf{$(a,b')\in\{(0,U_0),(1,U_1)\}$.}
                                    Thus $(a,b') = (0,U_0)$ or $(a,b') = (1,U_1)$ by \cref{upair_elim}.
                                    \begin{byCase}
                                    \caseOf{$(a,b') = (0,U_0)$.}
                                        Thus $a = 0$ by \cref{pair_eq_iff}.
                                        Thus $r_1 \neq 0$.
                                        Contradiction.
                                    \caseOf{$(a,b') = (1,U_1)$.}
                                        Thus $a = 1$ by \cref{pair_eq_iff}.
                                        Thus $r_1 \neq 1$.
                                        Contradiction.
                                    \end{byCase}
                                \end{byCase}
                            \end{subproof}
                        \end{byCase}
                    \end{subproof}
                    Thus $u_1$ is right-unique by \cref{rightunique}.
                \end{subproof}
                Show that $u_1$ is a relation.
                \begin{subproof}
                    Show that for all $w$ such that $w\in u_1$ there exist $x,y$ such that $w = (x,y)$.
                    \begin{subproof}
                        Fix $w$.
                        Assume $w\in u_1$.
                        Thus $w\in\{(0,U_0),(1,U_1)\}$ or $w\in\{(r_1,U_r)\}$ by \cref{union_iff}.
                        Thus $w = (0,U_0)$ or $w = (1,U_1)$ or $w = (r_1,U_r)$
                            by \cref{upair_elim,singleton_elim}.
                        Thus there exist $x,y$ such that $w = (x,y)$.
                    \end{subproof}
                    Thus $u_1$ is a relation by \cref{relation}.
                \end{subproof}
                Thus $u_1$ is a function.
                Show that $\dom{u_1} = \{0,1\}\union \img{s}{\initialSegment{1}}$.
                \begin{subproof}
                    Show that $\dom{u_1} = \{0,1\}\union\{r_1\}$.
                    \begin{subproof}
                        Show that $\dom{u_1}\subseteq \{0,1\}\union\{r_1\}$.
                        \begin{subproof}
                            Show that for all $x\in\dom{u_1}$ we have $x\in\{0,1\}\union\{r_1\}$.
                            \begin{subproof}
                                Fix $x$.
                                Assume $x\in\dom{u_1}$.
                                Take $y$ such that $(x,y)\in u_1$ by \cref{dom_iff}.
                                Thus $(x,y)\in\{(0,U_0),(1,U_1)\}$ or $(x,y)\in\{(r_1,U_r)\}$ by \cref{union_iff}.
                                \begin{byCase}
                                \caseOf{$(x,y)\in\{(0,U_0),(1,U_1)\}$.}
                                    Thus $(x,y) = (0,U_0)$ or $(x,y) = (1,U_1)$ by \cref{upair_elim}.
                                    Thus $x = 0$ or $x = 1$ by \cref{pair_eq_iff}.
                                    Thus $x\in\{0,1\}$ by \cref{upair_intro_left,upair_intro_right}.
                                    Thus $x\in\{0,1\}\union\{r_1\}$ by \cref{union_intro_left}.
                                \caseOf{$(x,y)\in\{(r_1,U_r)\}$.}
                                    Thus $(x,y) = (r_1,U_r)$ by \cref{singleton_elim}.
                                    Thus $x = r_1$ by \cref{pair_eq_iff}.
                                    Thus $x\in\{r_1\}$ by \cref{singleton_intro}.
                                    Thus $x\in\{0,1\}\union\{r_1\}$ by \cref{union_intro_right}.
                                \end{byCase}
                            \end{subproof}
                            Thus $\dom{u_1}\subseteq \{0,1\}\union\{r_1\}$ by \cref{subseteq}.
                        \end{subproof}
                        Show that $\{0,1\}\union\{r_1\}\subseteq \dom{u_1}$.
                        \begin{subproof}
                            Show that for all $x\in\{0,1\}\union\{r_1\}$ we have $x\in\dom{u_1}$.
                            \begin{subproof}
                                Fix $x$.
                                Assume $x\in\{0,1\}\union\{r_1\}$.
                                Thus $x\in\{0,1\}$ or $x\in\{r_1\}$ by \cref{union_iff}.
                                \begin{byCase}
                                \caseOf{$x\in\{0,1\}$.}
                                    Thus $x = 0$ or $x = 1$ by \cref{upair_elim}.
                                    \begin{byCase}
                                    \caseOf{$x = 0$.}
                                        Thus $(0,U_0)\in\{(0,U_0),(1,U_1)\}$ by \cref{upair_intro_left}.
                                        Thus $(0,U_0)\in u_1$ by \cref{union_intro_left}.
                                        Thus $x\in\dom{u_1}$ by \cref{dom_intro}.
                                    \caseOf{$x = 1$.}
                                        Thus $(1,U_1)\in\{(0,U_0),(1,U_1)\}$ by \cref{upair_intro_right}.
                                        Thus $(1,U_1)\in u_1$ by \cref{union_intro_left}.
                                        Thus $x\in\dom{u_1}$ by \cref{dom_intro}.
                                    \end{byCase}
                                \caseOf{$x\in\{r_1\}$.}
                                    Thus $x = r_1$ by \cref{singleton_elim}.
                                    Thus $(r_1,U_r)\in\{(r_1,U_r)\}$ by \cref{singleton_intro}.
                                    Thus $(r_1,U_r)\in u_1$ by \cref{union_intro_right}.
                                    Thus $x\in\dom{u_1}$ by \cref{dom_intro}.
                                \end{byCase}
                            \end{subproof}
                            Thus $\{0,1\}\union\{r_1\}\subseteq \dom{u_1}$ by \cref{subseteq}.
                        \end{subproof}
                        Thus $\dom{u_1} = \{0,1\}\union\{r_1\}$ by \cref{subseteq_antisymmetric}.
                    \end{subproof}
                    Show that $\{0,1\}\union \img{s}{\initialSegment{1}} = \{0,1\}\union\{r_1\}$.
                    \begin{subproof}
                        Thus $\img{s}{\initialSegment{1}} = \{r_1\}$.
                        Thus $\img{s}{\initialSegment{1}}\subseteq \{r_1\}$ by \cref{subseteq_refl}.
                        $r_1\in\{r_1\}$ by \cref{singleton_intro}.
                        Thus $\{r_1\}\subseteq \{0,1\}\union\{r_1\}$ by \cref{union_intro_right,subseteq}.
                        Thus $\img{s}{\initialSegment{1}}\subseteq \{0,1\}\union\{r_1\}$ by \cref{subseteq_transitive}.
                        Thus $\{0,1\}\union \img{s}{\initialSegment{1}}\subseteq \{0,1\}\union\{r_1\}$ by \cref{union_subsets_is_subset,subseteq_refl}.
                        Thus $\{0,1\}\subseteq \{0,1\}\union \img{s}{\initialSegment{1}}$ by \cref{union_intro_left,subseteq}.
                        Thus $\{r_1\}\subseteq \{0,1\}\union \img{s}{\initialSegment{1}}$ by \cref{union_intro_right,subseteq}.
                        Thus $\{0,1\}\union\{r_1\}\subseteq \{0,1\}\union \img{s}{\initialSegment{1}}$ by \cref{union_subsets_is_subset}.
                        Thus $\{0,1\}\union \img{s}{\initialSegment{1}} = \{0,1\}\union\{r_1\}$ by \cref{subseteq_antisymmetric}.
                    \end{subproof}
                    Thus $\dom{u_1} = \{0,1\}\union \img{s}{\initialSegment{1}}$ by \cref{subseteq_antisymmetric}.
                \end{subproof}
                Show that for all $p\in\dom{u_1}$ we have $u_1(p)\in\pow{X}$.
                \begin{subproof}
                    Fix $p$.
                    Assume $p\in\dom{u_1}$.
                    Thus $p\in\{0,1\}\union\{r_1\}$.
                    Thus $p\in\{0,1\}$ or $p\in\{r_1\}$ by \cref{union_iff}.
                    \begin{byCase}
                    \caseOf{$p\in\{0,1\}$.}
                        Thus $p = 0$ or $p = 1$ by \cref{upair_elim}.
                        \begin{byCase}
                        \caseOf{$p = 0$.}
                            Thus $(0,U_0)\in\{(0,U_0),(1,U_1)\}$ by \cref{upair_intro_left}.
                            Thus $(0,U_0)\in u_1$ by \cref{union_intro_left}.
                            Thus $u_1(p) = U_0$ by \cref{function_apply_intro}.
                            $U_0$ is open in $X$ with respect to $T$.
                            Thus $T$ is a topology on $X$ and $U_0\in T$ by \cref{open_in}.
                            Thus $T\subseteq\pow{X}$ by \cref{topology_on}.
                            Thus $u_1(p)\in\pow{X}$ by \cref{elem_subseteq}.
                        \caseOf{$p = 1$.}
                            Thus $(1,U_1)\in\{(0,U_0),(1,U_1)\}$ by \cref{upair_intro_right}.
                            Thus $(1,U_1)\in u_1$ by \cref{union_intro_left}.
                            Thus $u_1(p) = U_1$ by \cref{function_apply_intro}.
                            $U_1$ is open in $X$ with respect to $T$.
                            Thus $T$ is a topology on $X$ and $U_1\in T$ by \cref{open_in}.
                            Thus $T\subseteq\pow{X}$ by \cref{topology_on}.
                            Thus $u_1(p)\in\pow{X}$ by \cref{elem_subseteq}.
                        \end{byCase}
                    \caseOf{$p\in\{r_1\}$.}
                        Thus $p = r_1$ by \cref{singleton_elim}.
                        Thus $(r_1,U_r)\in\{(r_1,U_r)\}$ by \cref{singleton_intro}.
                        Thus $(r_1,U_r)\in u_1$ by \cref{union_intro_right}.
                        Thus $u_1(p) = U_r$ by \cref{function_apply_intro}.
                        $U_r$ is open in $X$ with respect to $T$.
                        Thus $T$ is a topology on $X$ and $U_r\in T$ by \cref{open_in}.
                        Thus $T\subseteq\pow{X}$ by \cref{topology_on}.
                        Thus $u_1(p)\in\pow{X}$ by \cref{elem_subseteq}.
                    \end{byCase}
                \end{subproof}
                Thus $u_1\in\funs{\{0,1\}\union \img{s}{\initialSegment{1}}}{\pow{X}}$ by \cref{funs_intro}.
            \end{subproof}
            Show that for all $p\in\dom{u_1}$ we have
                $u_1(p)\in T$ and
                $u_1(1) = X\setminus B$ and
                $A\subseteq u_1(0)$ and
                $\closure{u_1(0)}{X}{T}\subseteq u_1(1)$ and
                for all $p\in\dom{u_1}$ for all $q\in\dom{u_1}$ if $p < q$ then
                    $\closure{u_1(p)}{X}{T}\subseteq u_1(q)$ and
                for all $p\in\dom{u_1}$ if $1 < p$ then $u_1(p) = X$.
            \begin{subproof}
                Fix $p$.
                Assume $p\in\dom{u_1}$.
                Thus $u_1$ is a function by \cref{funs_is_function}.
                Thus $u_1$ is right-unique and $u_1$ is a relation.
                Show that $u_1(p)\in T$.
                \begin{subproof}
                    Thus $p\in\{0,1\}\union\{r_1\}$.
                    Thus $p\in\{0,1\}$ or $p\in\{r_1\}$ by \cref{union_iff}.
                    \begin{byCase}
                    \caseOf{$p\in\{0,1\}$.}
                        Thus $p = 0$ or $p = 1$ by \cref{upair_elim}.
                        \begin{byCase}
                        \caseOf{$p = 0$.}
                            Thus $(0,U_0)\in\{(0,U_0),(1,U_1)\}$ by \cref{upair_intro_left}.
                            Thus $(0,U_0)\in u_1$ by \cref{union_intro_left}.
                            Thus $u_1(p) = U_0$ by \cref{function_apply_intro}.
                            $U_0$ is open in $X$ with respect to $T$.
                            Thus $u_1(p)\in T$ by \cref{open_in}.
                        \caseOf{$p = 1$.}
                            Thus $(1,U_1)\in\{(0,U_0),(1,U_1)\}$ by \cref{upair_intro_right}.
                            Thus $(1,U_1)\in u_1$ by \cref{union_intro_left}.
                            Thus $u_1(p) = U_1$ by \cref{function_apply_intro}.
                            $U_1$ is open in $X$ with respect to $T$.
                            Thus $u_1(p)\in T$ by \cref{open_in}.
                        \end{byCase}
                    \caseOf{$p\in\{r_1\}$.}
                        Thus $p = r_1$ by \cref{singleton_elim}.
                        Thus $(r_1,U_r)\in\{(r_1,U_r)\}$ by \cref{singleton_intro}.
                        Thus $(r_1,U_r)\in u_1$ by \cref{union_intro_right}.
                        Thus $u_1(p) = U_r$ by \cref{function_apply_intro}.
                        $U_r$ is open in $X$ with respect to $T$.
                        Thus $u_1(p)\in T$ by \cref{open_in}.
                    \end{byCase}
                \end{subproof}
                Show that $u_1(1) = X\setminus B$.
                \begin{subproof}
                    Thus $(1,U_1)\in\{(0,U_0),(1,U_1)\}$ by \cref{upair_intro_right}.
                    Thus $(1,U_1)\in u_1$ by \cref{union_intro_left}.
                    Thus $u_1(1) = U_1$ by \cref{function_apply_intro}.
                    Thus $u_1(1) = X\setminus B$.
                \end{subproof}
                Show that $A\subseteq u_1(0)$.
                \begin{subproof}
                    Thus $(0,U_0)\in\{(0,U_0),(1,U_1)\}$ by \cref{upair_intro_left}.
                    Thus $(0,U_0)\in u_1$ by \cref{union_intro_left}.
                    Thus $u_1(0) = U_0$ by \cref{function_apply_intro}.
                    Thus $A\subseteq u_1(0)$.
                \end{subproof}
                Show that $\closure{u_1(0)}{X}{T}\subseteq u_1(1)$.
                \begin{subproof}
                    Thus $(0,U_0)\in\{(0,U_0),(1,U_1)\}$ by \cref{upair_intro_left}.
                    Thus $(0,U_0)\in u_1$ by \cref{union_intro_left}.
                    Thus $u_1(0) = U_0$ by \cref{function_apply_intro}.
                    Thus $(1,U_1)\in\{(0,U_0),(1,U_1)\}$ by \cref{upair_intro_right}.
                    Thus $(1,U_1)\in u_1$ by \cref{union_intro_left}.
                    Thus $u_1(1) = U_1$ by \cref{function_apply_intro}.
                    Thus $\closure{u_1(0)}{X}{T}\subseteq u_1(1)$ by \cref{subseteq}.
                \end{subproof}
                Show that for all $p\in\dom{u_1}$ for all $q\in\dom{u_1}$ if $p < q$ then
                    $\closure{u_1(p)}{X}{T}\subseteq u_1(q)$.
                \begin{subproof}
                    Fix $p\in\dom{u_1}$.
                    Fix $q\in\dom{u_1}$.
                    Assume $p < q$.
                    Thus $p\in\{0,1\}\union\{r_1\}$ and $q\in\{0,1\}\union\{r_1\}$.
                    Show that $q \neq 0$.
                    \begin{subproof}
                        Suppose not.
                        Thus $q = 0$.
                        Thus $p < 0$ by \cref{real_lt_subst_right}.
                        Thus $p\in\{0,1\}\union\{r_1\}$.
                        Thus $p\in\{0,1\}$ or $p\in\{r_1\}$ by \cref{union_iff}.
                        \begin{byCase}
                        \caseOf{$p\in\{0,1\}$.}
                            Thus $p = 0$ or $p = 1$ by \cref{upair_elim}.
                            \begin{byCase}
                            \caseOf{$p = 0$.}
                                Thus $0 < 0$ by \cref{real_lt_subst_left}.
                                $0\in\reals$ by \cref{real_add_ident}.
                                Contradiction by \cref{real_lt_irreflexive}.
                            \caseOf{$p = 1$.}
                                Thus $1 < 0$ by \cref{real_lt_subst_left}.
                                $0\in\reals$ by \cref{real_add_ident}.
                                $1\in\reals$ by \cref{real_mul_ident}.
                                $0 < 1$ by \cref{real_zero_lt_one}.
                                Thus $1 < 1$ by \cref{real_lt_trans}.
                                Contradiction by \cref{real_lt_irreflexive}.
                            \end{byCase}
                        \caseOf{$p\in\{r_1\}$.}
                            Thus $p = r_1$ by \cref{singleton_elim}.
                            Thus $r_1 < 0$ by \cref{real_lt_subst_left}.
                            $0\in\reals$ by \cref{real_add_ident}.
                            Thus $0 < r_1$ by \cref{unit_rational_pos}.
                            Thus $0 < 0$ by \cref{real_lt_trans}.
                            Contradiction by \cref{real_lt_irreflexive}.
                        \end{byCase}
                    \end{subproof}
                    Thus $q\in\{1\}\union\{r_1\}$.
                    Thus $q\in\{1\}$ or $q\in\{r_1\}$ by \cref{union_iff}.
                    \begin{byCase}
                    \caseOf{$q\in\{r_1\}$.}
                        Thus $q = r_1$ by \cref{singleton_elim}.
                        Show that $p = 0$.
                        \begin{subproof}
                            Suppose not.
                            Thus $p\in\{0,1\}\union\{r_1\}$.
                            Thus $p\in\{0,1\}$ or $p\in\{r_1\}$ by \cref{union_iff}.
                            \begin{byCase}
                            \caseOf{$p\in\{0,1\}$.}
                                Thus $p = 0$ or $p = 1$ by \cref{upair_elim}.
                                \begin{byCase}
                                \caseOf{$p = 1$.}
                                    Thus $1 < r_1$ by \cref{real_lt_subst_left}.
                                    Contradiction.
                                \caseOf{$p = 0$.}
                                    Contradiction.
                                \end{byCase}
                            \caseOf{$p\in\{r_1\}$.}
                                Thus $p = r_1$ by \cref{singleton_elim}.
                                Thus $r_1 < r_1$ by \cref{real_lt_subst_left}.
                                Contradiction by \cref{real_lt_irreflexive}.
                            \end{byCase}
                        \end{subproof}
                        Thus $(p,U_0) = (0,U_0)$ by \cref{pair_eq_iff}.
                        Thus $(q,U_r) = (r_1,U_r)$ by \cref{pair_eq_iff}.
                        Thus $(p,U_0)\in u_1$ by \cref{union_intro_left,upair_intro_left}.
                        Thus $(q,U_r)\in u_1$ by \cref{union_intro_right,singleton_intro}.
                        Thus $u_1(p) = U_0$ and $u_1(q) = U_r$ by \cref{function_apply_intro}.
                        Thus $\closure{u_1(p)}{X}{T}\subseteq u_1(q)$ by \cref{subseteq}.
                    \caseOf{$q\in\{1\}$.}
                        Thus $q = 1$ by \cref{singleton_elim}.
                        Thus $p\in\{0,1\}\union\{r_1\}$.
                        Thus $p\in\{0,1\}$ or $p\in\{r_1\}$ by \cref{union_iff}.
                        \begin{byCase}
                        \caseOf{$p\in\{0,1\}$.}
                            Show that $p = 0$.
                            \begin{subproof}
                                Suppose not.
                                Thus $p = 1$ by \cref{upair_elim}.
                                Thus $1 < 1$ by \cref{real_lt_subst_left}.
                                $1\in\reals$ by \cref{real_mul_ident}.
                                Contradiction by \cref{real_lt_irreflexive}.
                            \end{subproof}
                            Thus $(p,U_0) = (0,U_0)$ by \cref{pair_eq_iff}.
                            Thus $(q,U_1) = (1,U_1)$ by \cref{pair_eq_iff}.
                            Thus $(p,U_0)\in u_1$ by \cref{union_intro_left,upair_intro_left}.
                            Thus $(q,U_1)\in u_1$ by \cref{union_intro_left,upair_intro_right}.
                            Thus $u_1(p) = U_0$ and $u_1(q) = U_1$ by \cref{function_apply_intro}.
                            Thus $\closure{u_1(p)}{X}{T}\subseteq u_1(q)$ by \cref{subseteq}.
                        \caseOf{$p\in\{r_1\}$.}
                            Thus $p = r_1$ by \cref{singleton_elim}.
                            Thus $(p,U_r) = (r_1,U_r)$ by \cref{pair_eq_iff}.
                            Thus $(q,U_1) = (1,U_1)$ by \cref{pair_eq_iff}.
                            Thus $(p,U_r)\in u_1$ by \cref{union_intro_right,singleton_intro}.
                            Thus $(q,U_1)\in u_1$ by \cref{union_intro_left,upair_intro_right}.
                            Thus $u_1(p) = U_r$ and $u_1(q) = U_1$ by \cref{function_apply_intro}.
                            Thus $\closure{u_1(p)}{X}{T}\subseteq u_1(q)$ by \cref{subseteq}.
                        \end{byCase}
                    \end{byCase}
                \end{subproof}
                Show that for all $p\in\dom{u_1}$ if $1 < p$ then $u_1(p) = X$.
                \begin{subproof}
                    Fix $p\in\dom{u_1}$.
                    Assume $1 < p$.
                    Suppose not.
                    Thus $u_1(p) \neq X$.
                    Thus $p\in\{0,1\}\union\{r_1\}$.
                    Thus $p\in\{0,1\}$ or $p\in\{r_1\}$ by \cref{union_iff}.
                    \begin{byCase}
                    \caseOf{$p\in\{0,1\}$.}
                        Thus $p = 0$ or $p = 1$ by \cref{upair_elim}.
                        \begin{byCase}
                        \caseOf{$p = 0$.}
                            Thus $1 < 0$ by \cref{real_lt_subst_left}.
                            $0\in\reals$ by \cref{real_add_ident}.
                            $1\in\reals$ by \cref{real_mul_ident}.
                            $0 < 1$ by \cref{real_zero_lt_one}.
                            Thus $1 < 1$ by \cref{real_lt_trans}.
                            Contradiction by \cref{real_lt_irreflexive}.
                        \caseOf{$p = 1$.}
                            Thus $1 < 1$ by \cref{real_lt_subst_left}.
                            $1\in\reals$ by \cref{real_mul_ident}.
                            Contradiction by \cref{real_lt_irreflexive}.
                        \end{byCase}
                    \caseOf{$p\in\{r_1\}$.}
                        Thus $p = r_1$ by \cref{singleton_elim}.
                        Thus $1 < r_1$ by \cref{real_lt_subst_left}.
                        Contradiction.
                    \end{byCase}
                \end{subproof}
            \end{subproof}
            Thus $u_1\in\urysohnAdmissible{s}{X}{T}{A}{B}{1}$ by \cref{urysohn_admissible}.
        \end{subproof}
        $1\in\naturalsPlus$ by \cref{real_one_in_naturals}.
        Thus $(1,u_1)\in G$ by \cref{urysohn_good_pair_intro}.
    \end{subproof}
\end{proof}

% from §33 helper lemma 11: initial good pair when $r_1 < 1$
\begin{lemma}\label{urysohn_u1_lt1}
    Assume $X$ is normal with respect to $T$.
    Assume $U_1 = X\setminus B$.
    Assume $U_1$ is open in $X$ with respect to $T$.
    Assume $U_0$ is open in $X$ with respect to $T$.
    Assume $A\subseteq U_0$.
    Assume $\closure{U_0}{X}{T}\subseteq U_1$.
    Assume $r_1\in\reals$.
    Assume $r_1\in\unitRationals$.
    Assume $r_1 < 1$.
    Assume $\img{s}{\initialSegment{1}} = \{r_1\}$.
    Let $G = \urysohnGoodPairs{s}{X}{T}{A}{B}$.
    Then there exists $u_1$ such that $(1,u_1)\in G$.
\end{lemma}
\begin{proof}
    $U_1$ is open in $X$ with respect to $T$.
    Thus $\closure{U_0}{X}{T}\subseteq U_1$.
    $U_0$ is open in $X$ with respect to $T$.
    Thus $T$ is a topology on $X$ and $U_0\in T$ by \cref{open_in}.
    Thus $T\subseteq\pow{X}$ by \cref{topology_on}.
    Thus $U_0\in\pow{X}$ by \cref{elem_subseteq}.
    Thus $U_0\subseteq X$ by \cref{pow_iff}.
    Thus $\closure{U_0}{X}{T}$ is closed in $X$ with respect to $T$ by \cref{closure_closed_in}.
    Thus $0 < r_1$ by \cref{unit_rational_pos}.
    Take $U_r$ such that
        $U_r$ is open in $X$ with respect to $T$ and
        $\closure{U_0}{X}{T}\subseteq U_r$ and
        $\closure{U_r}{X}{T}\subseteq U_1$
        by \cref{normal_closure_neighborhood}.
    Thus there exists $u_1$ such that $(1,u_1)\in G$ by \cref{urysohn_u1_lt1_admissible}.
\end{proof}

% from §33 helper lemma 13: admissible chain is linearly ordered by inclusion
\begin{lemma}\label{urysohn_unit_interval_chain}
    Assume $s$ is a sequence in $\unitRationals$.
    Let $G = \urysohnGoodPairs{s}{X}{T}{A}{B}$.
    Let $R = \{(z,w) \mid z\in G, w\in G \mid
        \fst{w} = \fst{z} + 1 \land \snd{z}\subseteq\snd{w}\}$.
    Assume $c$ is a function on $G$.
    Assume for all $z\in G$ we have $z\mathrel{R} \apply{c}{z}$.
    Assume $u$ is a sequence in $G$.
    Assume for all $n\in\naturalsPlus$ we have $u(n + 1) = \apply{c}{u(n)}$.
    Let $F = \{\snd{u(n)} \mid n\in\naturalsPlus\}$.
    Then for all $f,g$ such that $f\in F$ and $g\in F$ we have $f\subseteq g$ or $g\subseteq f$.
\end{lemma}
\begin{proof}
    Show that for all $f,g$ such that $f\in F$ and $g\in F$ we have $f\subseteq g$ or $g\subseteq f$.
    \begin{subproof}
        Show that for all $n\in\naturalsPlus$ we have $\snd{u(n)}\subseteq \snd{u(n + 1)}$.
        \begin{subproof}
            Fix $n\in\naturalsPlus$.
            Thus $u$ is a sequence in $G$.
            Thus $u\in\funs{\naturalsPlus}{G}$ by \cref{sequence_in}.
            Thus $u(n)\in G$ by \cref{funs_type_apply}.
            Thus $u(n)\mathrel{R} \apply{c}{u(n)}$.
            Thus $u(n + 1) = \apply{c}{u(n)}$.
            Thus $u(n)\mathrel{R} u(n + 1)$.
            Thus $(u(n),u(n + 1))\in R$.
            Thus there exist $z,w$ such that
                $z\in G$ and
                $w\in G$ and
                $(z,w) = (u(n),u(n + 1))$ and
                $\fst{w} = \fst{z} + 1$ and
                $\snd{z}\subseteq \snd{w}$.
            Thus $\snd{u(n)}\subseteq \snd{u(n + 1)}$ by \cref{pair_eq_iff}.
        \end{subproof}
        Let $A_0 = \{m\in\naturalsPlus \mid \text{for all $n\in\naturalsPlus$ if $n \leq m$ then
            $\snd{u(n)}\subseteq \snd{u(m)}$}\}$.
        $A_0\subseteq\naturalsPlus$ by \cref{subseteq}.
        Show that $1\in A_0$.
        \begin{subproof}
            $1\in\naturalsPlus$ by \cref{real_one_in_naturals}.
            Show that for all $n\in\naturalsPlus$ if $n \leq 1$ then $\snd{u(n)}\subseteq \snd{u(1)}$.
            \begin{subproof}
                Fix $n\in\naturalsPlus$.
                Assume $n \leq 1$.
                $1\in\reals$ and $n\in\reals$ by \cref{real_naturals_subset_reals,subseteq}.
                Thus $1 < n$ or $1 = n$ by \cref{real_one_leq_naturalsplus}.
                Thus $1 \leq n$.
                Thus $n = 1$ by \cref{real_leq_antisym}.
                Thus $\snd{u(n)} = \snd{u(1)}$.
                Thus $\snd{u(n)}\subseteq \snd{u(1)}$ by \cref{subseteq_refl}.
            \end{subproof}
            Thus $1\in A_0$.
        \end{subproof}
        Show that for all $k\in A_0$ we have $k + 1\in A_0$.
        \begin{subproof}
            Fix $k$.
            Assume $k\in A_0$.
            Thus $k\in\naturalsPlus$ by \cref{subseteq}.
            Show that for all $n\in\naturalsPlus$ if $n \leq k + 1$ then $\snd{u(n)}\subseteq \snd{u(k + 1)}$.
            \begin{subproof}
                Fix $n\in\naturalsPlus$.
                Assume $n \leq k + 1$.
                \begin{byCase}
                \caseOf{$n = k + 1$.}
                    Thus $\snd{u(n)} = \snd{u(k + 1)}$.
                    Thus $\snd{u(n)}\subseteq \snd{u(k + 1)}$ by \cref{subseteq_refl}.
                \caseOf{$n \neq k + 1$.}
                    Thus $n \leq k$ by \cref{real_naturals_leq_succ_elim}.
                    Thus $\snd{u(n)}\subseteq \snd{u(k)}$.
                    Thus $\snd{u(k)}\subseteq \snd{u(k + 1)}$.
                    Thus $\snd{u(n)}\subseteq \snd{u(k + 1)}$ by \cref{subseteq_transitive}.
                \end{byCase}
            \end{subproof}
            Thus $k + 1\in A_0$.
        \end{subproof}
        Thus for all $m\in\naturalsPlus$ we have $m\in A_0$ by \cref{real_naturals_induction}.
        Fix $f$.
        Fix $g$.
        Assume $f\in F$ and $g\in F$.
        Take $n$ such that $n\in\naturalsPlus$ and $f = \snd{u(n)}$.
        Take $m$ such that $m\in\naturalsPlus$ and $g = \snd{u(m)}$.
        $n\in\reals$ and $m\in\reals$ by \cref{real_naturals_subset_reals,subseteq}.
        Thus $n < m$ or $n = m$ or $m < n$ by \cref{real_lt_trichotomy}.
        \begin{byCase}
        \caseOf{$n = m$.}
            Thus $g = \snd{u(n)}$.
            Thus $f = g$.
            Thus $f\subseteq g$ by \cref{subseteq_refl}.
        \caseOf{$n < m$.}
            Thus $n \leq m$.
            Thus $\snd{u(n)}\subseteq \snd{u(m)}$.
            Thus $f\subseteq g$.
        \caseOf{$m < n$.}
            Thus $m \leq n$.
            Thus $\snd{u(m)}\subseteq \snd{u(n)}$.
            Thus $g\subseteq f$.
        \end{byCase}
    \end{subproof}
\end{proof}

% from §33 helper lemma 14: unions of the admissible chain are a function
\begin{lemma}\label{urysohn_unit_interval_union_function}
    Assume $s$ is a sequence in $\unitRationals$.
    Let $G = \urysohnGoodPairs{s}{X}{T}{A}{B}$.
    Let $R = \{(z,w) \mid z\in G, w\in G \mid
        \fst{w} = \fst{z} + 1 \land \snd{z}\subseteq\snd{w}\}$.
    Assume $c$ is a function on $G$.
    Assume for all $z\in G$ we have $z\mathrel{R} \apply{c}{z}$.
    Assume $u$ is a sequence in $G$.
    Assume for all $n\in\naturalsPlus$ we have $u(n + 1) = \apply{c}{u(n)}$.
    Assume for all $n\in\naturalsPlus$ we have $\fst{u(n)} = n$.
    Let $F = \{\snd{u(n)} \mid n\in\naturalsPlus\}$.
    Let $U = \unions{F}$.
    Then $U$ is a function.
\end{lemma}
\begin{proof}
    Show that for all $g$ such that $g\in F$ we have $g$ is a function.
    \begin{subproof}
        Fix $g$.
        Assume $g\in F$.
        Take $n$ such that $n\in\naturalsPlus$ and $g = \snd{u(n)}$.
        Thus $u$ is a sequence in $G$.
        Thus $u\in\funs{\naturalsPlus}{G}$ by \cref{sequence_in}.
        Thus $u(n)\in G$ by \cref{funs_type_apply}.
        Thus there exist $k,f$ such that
            $k\in\naturalsPlus$ and
            $f\in\urysohnAdmissible{s}{X}{T}{A}{B}{k}$ and
            $u(n) = (k,f)$
            by \cref{urysohn_good_pair_decomp}.
        Thus $\snd{u(n)} = f$ by \cref{snd_eq}.
        Thus $g = f$.
        Thus $f\in\funs{\{0,1\}\union \img{s}{\initialSegment{k}}}{\pow{X}}$ by \cref{urysohn_admissible}.
        Thus $f$ is a function by \cref{funs_is_function}.
        Thus $g$ is a function.
    \end{subproof}
    Thus $F$ is a family of functions.
    Thus for all $f,g$ such that $f\in F$ and $g\in F$ we have $f\subseteq g$ or $g\subseteq f$
        by \cref{urysohn_unit_interval_chain}.
    Thus $U$ is a function by \cref{unions_of_compatible_family_of_function_is_function}.
\end{proof}

% from §33 helper lemma 15: Urysohn function is a function to reals
\begin{lemma}\label{urysohn_unit_interval_f_funs}
    Assume $s$ is a sequence in $\unitRationals$.
    Assume for all $p$ such that $p\in\unitRationals$ there exists $n\in\naturalsPlus$ such that $s(n) = p$.
    Let $G = \urysohnGoodPairs{s}{X}{T}{A}{B}$.
    Assume $u$ is a sequence in $G$.
    Assume for all $n\in\naturalsPlus$ we have $\fst{u(n)} = n$.
    Let $F = \{\snd{u(n)} \mid n\in\naturalsPlus\}$.
    Let $U = \unions{F}$.
    Assume $U$ is a function.
    Let $f = \{(x,y) \mid x\in X, y\in\reals \mid
        \text{$y$ is an infimum of $\urysohnLevelSet{U}{x}$}\}$.
    Then $f$ is a function from $X$ to $\reals$.
\end{lemma}
\begin{proof}
    Show that $f$ is a function from $X$ to $\reals$.
    \begin{subproof}
        Show that $f$ is a function.
        \begin{subproof}
            Show that for all $x,y_1,y_2$ such that $(x,y_1)\in f$ and $(x,y_2)\in f$ we have $y_1 = y_2$.
            \begin{subproof}
                Fix $x$.
                Let $Q_x = \urysohnLevelSet{U}{x}$.
                Fix $y_1$.
                Fix $y_2$.
                Assume $(x,y_1)\in f$ and $(x,y_2)\in f$.
                Thus $y_1$ and $y_2$ are infimums of $Q_x$.
                Thus $y_1 = y_2$ by \cref{real_infimum_unique}.
            \end{subproof}
            Thus $f$ is right-unique by \cref{rightunique}.
            Show that for all $w$ such that $w\in f$ there exist $x,y$ such that $w = (x,y)$.
            \begin{subproof}
                Fix $w$.
                Assume $w\in f$.
                Take $x,y$ such that
                    $x\in X$ and $y\in\reals$ and $w = (x,y)$ and
                    $y$ is an infimum of $\urysohnLevelSet{U}{x}$.
            \end{subproof}
            Thus $f$ is a relation by \cref{relation}.
        \end{subproof}
        Thus $f$ is a function.
        Show that $X\subseteq\dom{f}$.
        \begin{subproof}
            Show that for all $x\in X$ we have $x\in\dom{f}$.
            \begin{subproof}
                Fix $x$.
                Assume $x\in X$.
                Let $Q_x = \urysohnLevelSet{U}{x}$.
                Show that $Q_x\neq \emptyset$.
                \begin{subproof}
                    $1\in\reals$ by \cref{real_mul_ident}.
                    $0\in\reals$ by \cref{real_add_ident}.
                    $0 < 1$ by \cref{real_zero_lt_one}.
                    Thus $0 \leq 1$ by \cref{real_lt_imp_leq}.
                    Let $b = 1 + 1$.
                    Thus $b\in\reals$ by \cref{real_add_closed}.
                    $0 + 1 < 1 + 1$ by \cref{real_lt_add}.
                    $0 + 1 = 1$ by \cref{real_add_ident}.
                    Thus $1 < b$ by \cref{real_lt_subst_left}.
                    Take $p$ such that $p\in\unitRationals$ and $p\in\intervalopen{1}{b}$
                        by \cref{unit_rationals_dense_nonneg}.
                    Thus $1 < p$ by \cref{intervalopen}.
                    Thus there exists $k\in\naturalsPlus$ such that $s(k) = p$ by assumption.
                    Take $k$ such that $k\in\naturalsPlus$ and $s(k) = p$.
                    Thus $u$ is a sequence in $G$.
                    Thus $u\in\funs{\naturalsPlus}{G}$ by \cref{sequence_in}.
                    Thus $u(k)\in G$ by \cref{funs_type_apply}.
                    Take $n,t$ such that
                        $n\in\naturalsPlus$ and
                        $t\in\urysohnAdmissible{s}{X}{T}{A}{B}{n}$ and
                        $u(k) = (n,t)$
                        by \cref{urysohn_good_pair_decomp}.
                    Thus $\fst{u(k)} = n$ by \cref{fst_eq}.
                    Thus $k = n$.
                    Thus $u(k) = (k,t)$ by \cref{pair_eq_iff}.
                    Thus $\snd{u(k)} = t$ by \cref{snd_eq}.
                    Thus $t\in F$.
                    Thus $t\in\funs{\{0,1\}\union \img{s}{\initialSegment{k}}}{\pow{X}}$
                        by \cref{urysohn_admissible}.
                    Thus $t$ is a function by \cref{funs_is_function}.
                    Thus $t$ is right-unique and $t$ is a relation.
                    Thus $\dom{t} = \{0,1\}\union \img{s}{\initialSegment{k}}$ by \cref{funs_elim}.
                    Show that $p\in\dom{t}$.
                    \begin{subproof}
                        Thus $k \leq k$.
                        Thus $k\in\initialSegment{k}$ by \cref{initial_segment_naturalsplus}.
                        Let $y = k$.
                        Thus $y\in\initialSegment{k}$.
                        Thus $s(y) = s(k)$.
                        Thus $s$ is a sequence in $\unitRationals$.
                        Thus $s\in\funs{\naturalsPlus}{\unitRationals}$ by \cref{sequence_in}.
                        Thus $s$ is a function by \cref{funs_is_function}.
                        Thus $\dom{s} = \naturalsPlus$ by \cref{funs_elim}.
                        Thus $y\in\dom{s}$ by \cref{subseteq}.
                        Thus $(y,s(y))\in s$ by \cref{function_apply_elim}.
                        Thus $(y,s(k))\in s$.
                        Thus $s(k)\in\img{s}{\initialSegment{k}}$ by \cref{img_iff}.
                        Thus $p\in\img{s}{\initialSegment{k}}$.
                        Thus $p\in\dom{t}$ by \cref{union_intro_right}.
                    \end{subproof}
                    Show that for all $q\in\dom{t}$ if $1 < q$ then $t(q) = X$.
                    \begin{subproof}
                        Show that $0\in\dom{t}$.
                        \begin{subproof}
                            $0\in\{0,1\}$ by \cref{upair_intro_left}.
                            Thus $0\in\dom{t}$ by \cref{union_intro_left,subseteq}.
                        \end{subproof}
                        Show that $1\in\dom{t}$.
                        \begin{subproof}
                            $1\in\{0,1\}$ by \cref{upair_intro_right}.
                            Thus $1\in\dom{t}$ by \cref{union_intro_left,subseteq}.
                        \end{subproof}
                        $0\in\reals$ by \cref{real_add_ident}.
                        $1\in\reals$ by \cref{real_mul_ident}.
                        $0 < 1$ by \cref{real_zero_lt_one}.
                        Thus $\closure{t(0)}{X}{T}\subseteq t(1)$ and
                            for all $q\in\dom{t}$ if $1 < q$ then $t(q) = X$
                            by \cref{urysohn_admissible}.
                        Thus for all $q\in\dom{t}$ if $1 < q$ then $t(q) = X$.
                    \end{subproof}
                    Thus $t(p) = X$.
                    Thus $(p,t(p))\in t$ by \cref{function_apply_elim}.
                    Thus $(p,t(p))\in U$ by \cref{unions_intro}.
                    Thus $U(p) = t(p)$ by \cref{function_apply_intro}.
                    Thus $x\in X$.
                    Thus $x\in U(p)$.
                    Thus $p\in Q_x$.
                    Thus $Q_x\neq \emptyset$.
                \end{subproof}
                Show that $Q_x\subseteq\reals$.
                \begin{subproof}
                    Show that for all $p\in Q_x$ we have $p\in\reals$.
                    \begin{subproof}
                        Fix $p$.
                        Assume $p\in Q_x$.
                        Thus $p\in\unitRationalsZero$ by \cref{urysohn_level_set}.
                        Thus $p\in\unitRationals$ or $p\in\{0\}$ by \cref{unit_rationals_zero,union_iff}.
                        \begin{byCase}
                        \caseOf{$p\in\unitRationals$.}
                            Take $q$ such that $q\in\naturalsPlus\times\naturalsPlus$ and
                                $p = \fst{q} \rmul \inv{\snd{q}}$
                                by \cref{unit_rationals}.
                            Thus $\fst{q}\in\naturalsPlus$ and $\snd{q}\in\naturalsPlus$ by \cref{fst_type,snd_type}.
                            Thus $\fst{q}\in\reals$ and $\snd{q}\in\reals$ by \cref{real_naturals_subset_reals,subseteq}.
                            $1\in\reals$ by \cref{real_mul_ident}.
                            $0\in\reals$ by \cref{real_add_ident}.
                            $0 < 1$ by \cref{real_zero_lt_one}.
                            Show that $0 < \snd{q}$.
                            \begin{subproof}
                                Thus $1 < \snd{q}$ or $1 = \snd{q}$ by \cref{real_one_leq_naturalsplus}.
                                \begin{byCase}
                                \caseOf{$1 < \snd{q}$.}
                                    Thus $0 < \snd{q}$ by \cref{real_lt_trans}.
                                \caseOf{$1 = \snd{q}$.}
                                    Thus $0 < \snd{q}$ by \cref{real_lt_subst_right}.
                                \end{byCase}
                            \end{subproof}
                            Thus $\snd{q}\neq 0$ by \cref{real_pos_nonzero}.
                            Thus $\inv{\snd{q}}\in\reals$ by \cref{real_mul_inverse}.
                            Thus $p\in\reals$ by \cref{real_mul_closed}.
                        \caseOf{$p\in\{0\}$.}
                            Thus $p = 0$ by \cref{singleton_elim}.
                            $0\in\reals$ by \cref{real_add_ident}.
                            Thus $p\in\reals$.
                        \end{byCase}
                    \end{subproof}
                    Thus $Q_x\subseteq\reals$ by \cref{subseteq}.
                \end{subproof}
                Show that $Q_x$ is bounded below.
                \begin{subproof}
                    Show that $0$ is a lower bound of $Q_x$.
                    \begin{subproof}
                        $0\in\reals$ by \cref{real_add_ident}.
                        Show that for all $p\in Q_x$ we have $p\in\reals$ and $0 \leq p$.
                        \begin{subproof}
                            Fix $p$.
                            Assume $p\in Q_x$.
                            Thus $p\in\unitRationalsZero$.
                            Thus $p\in\unitRationals$ or $p\in\{0\}$ by \cref{unit_rationals_zero,union_iff}.
                            \begin{byCase}
                            \caseOf{$p\in\{0\}$.}
                                Thus $p = 0$ by \cref{singleton_elim}.
                                $0\in\reals$ by \cref{real_add_ident}.
                                Thus $p\in\reals$.
                                Thus $0 \leq p$.
                            \caseOf{$p\in\unitRationals$.}
                                Take $q$ such that $q\in\naturalsPlus\times\naturalsPlus$ and
                                    $p = \fst{q} \rmul \inv{\snd{q}}$
                                    by \cref{unit_rationals}.
                                Thus $\fst{q}\in\naturalsPlus$ and $\snd{q}\in\naturalsPlus$ by \cref{fst_type,snd_type}.
                                Thus $\fst{q}\in\reals$ and $\snd{q}\in\reals$ by \cref{real_naturals_subset_reals,subseteq}.
                                $1\in\reals$ by \cref{real_mul_ident}.
                                Thus $1 < \fst{q}$ or $1 = \fst{q}$ by \cref{real_one_leq_naturalsplus}.
                                Thus $1 < \snd{q}$ or $1 = \snd{q}$ by \cref{real_one_leq_naturalsplus}.
                                $0\in\reals$ by \cref{real_add_ident}.
                                $0 < 1$ by \cref{real_zero_lt_one}.
                                Show that $0 < \fst{q}$.
                                \begin{subproof}
                                    \begin{byCase}
                                    \caseOf{$1 < \fst{q}$.}
                                        Thus $0 < \fst{q}$ by \cref{real_lt_trans}.
                                    \caseOf{$1 = \fst{q}$.}
                                        Thus $0 < \fst{q}$ by \cref{real_lt_subst_right}.
                                    \end{byCase}
                                \end{subproof}
                                Show that $0 < \snd{q}$.
                                \begin{subproof}
                                    \begin{byCase}
                                    \caseOf{$1 < \snd{q}$.}
                                        Thus $0 < \snd{q}$ by \cref{real_lt_trans}.
                                    \caseOf{$1 = \snd{q}$.}
                                        Thus $0 < \snd{q}$ by \cref{real_lt_subst_right}.
                                    \end{byCase}
                                \end{subproof}
                                Thus $\snd{q}\neq 0$ by \cref{real_pos_nonzero}.
                                Thus $\inv{\snd{q}}\in\reals$ by \cref{real_mul_inverse}.
                                Thus $p\in\reals$ by \cref{real_mul_closed}.
                                Thus $\inv{\snd{q}}$ is positive by \cref{real_inv_pos}.
                                Thus $0 < p$ by \cref{real_mul_pos_pos_gt}.
                                Thus $0 \leq p$.
                            \end{byCase}
                        \end{subproof}
                        Thus for all $p\in Q_x$ we have $p\in\reals$.
                        Thus $Q_x\subseteq\reals$ by \cref{subseteq}.
                        Thus for all $p\in Q_x$ we have $0 \leq p$.
                        Thus $0$ is a lower bound of $Q_x$ by \cref{real_lower_bound}.
                    \end{subproof}
                    Thus $Q_x$ is bounded below.
                \end{subproof}
                Take $p$ such that $p$ is an infimum of $Q_x$ by \cref{real_infimum_exists}.
                Thus $p$ is a lower bound of $Q_x$ by \cref{real_infimum}.
                Thus $p\in\reals$ by \cref{real_lower_bound}.
                Thus $(x,p)\in f$.
                Thus $x\in\dom{f}$ by \cref{dom_intro}.
            \end{subproof}
            Thus $X\subseteq\dom{f}$ by \cref{subseteq}.
        \end{subproof}
        Show that for all $x\in\dom{f}$ we have $x\in X$.
        \begin{subproof}
            Fix $x$.
            Assume $x\in\dom{f}$.
            Take $p$ such that $(x,p)\in f$ by \cref{dom_iff}.
            Take $x_0,y_0$ such that
                $x_0\in X$ and
                $y_0\in\reals$ and
                $(x,p) = (x_0,y_0)$ and
                $y_0$ is an infimum of $\urysohnLevelSet{U}{x_0}$.
            Thus $x = x_0$ by \cref{pair_eq_iff}.
            Thus $x\in X$.
        \end{subproof}
        Thus $\dom{f}\subseteq X$ by \cref{subseteq}.
        Thus $\dom{f} = X$ by \cref{subseteq_antisymmetric}.
        Show that for all $x\in\dom{f}$ we have $f(x)\in\reals$.
        \begin{subproof}
            Fix $x\in\dom{f}$.
            Let $Q_x = \urysohnLevelSet{U}{x}$.
            Take $p$ such that $(x,p)\in f$ by \cref{dom_iff}.
            Thus $p$ is an infimum of $Q_x$.
            Thus $p$ is a lower bound of $Q_x$ by \cref{real_infimum}.
            Thus $p\in\reals$ by \cref{real_lower_bound}.
            Thus $f(x) = p$ by \cref{function_apply_intro}.
            Thus $f(x)\in\reals$.
        \end{subproof}
        Thus $f\in\funs{X}{\reals}$ by \cref{funs_intro}.
    \end{subproof}
\end{proof}

% from §33 helper lemma 16: existence of a first good pair
\begin{lemma}\label{urysohn_unit_interval_u1}
    Assume $X$ is normal with respect to $T$.
    Assume $A$ is closed in $X$ with respect to $T$.
    Assume $B$ is closed in $X$ with respect to $T$.
    Assume $A$ is disjoint from $B$.
    Assume $s$ is a sequence in $\unitRationals$.
    Let $G = \urysohnGoodPairs{s}{X}{T}{A}{B}$.
    Then there exists $u_1$ such that $(1,u_1)\in G$.
\end{lemma}
\begin{proof}
    Show that there exists $u_1$ such that $(1,u_1)\in G$.
    \begin{subproof}
        Let $U_1 = X\setminus B$.
        $U_1$ is open in $X$ with respect to $T$ by \cref{closed_in}.
        Show that $A\subseteq U_1$.
        \begin{subproof}
            Show that for all $x\in A$ we have $x\in U_1$.
            \begin{subproof}
                Fix $x$.
                Assume $x\in A$.
                Thus $x\notin B$ by \cref{disjoint}.
                Thus $A\subseteq X$ by \cref{closed_in}.
                Thus $x\in X$ by \cref{subseteq}.
                Thus $x\in X\setminus B$ by \cref{setminus_intro}.
                Thus $x\in U_1$.
            \end{subproof}
            Thus $A\subseteq U_1$ by \cref{subseteq}.
        \end{subproof}
        Take $U_0$ such that
            $U_0$ is open in $X$ with respect to $T$ and
            $A\subseteq U_0$ and
            $\closure{U_0}{X}{T}\subseteq U_1$
            by \cref{normal_closure_neighborhood}.
        Let $r_1 = s(1)$.
        $s\in\funs{\naturalsPlus}{\unitRationals}$ by \cref{sequence_in}.
        $1\in\naturalsPlus$ by \cref{real_one_in_naturals}.
        Thus $r_1\in\unitRationals$ by \cref{funs_type_apply}.
        Show that $\img{s}{\initialSegment{1}} = \{r_1\}$.
        \begin{subproof}
            Show that $\img{s}{\initialSegment{1}}\subseteq \{r_1\}$.
            \begin{subproof}
                Show that for all $y\in \img{s}{\initialSegment{1}}$ we have $y\in\{r_1\}$.
                \begin{subproof}
                    Fix $y$.
                    Assume $y\in \img{s}{\initialSegment{1}}$.
                    Thus there exists $k$ such that $k\in\initialSegment{1}$ and $(k,y)\in s$ by \cref{img_iff}.
                    Take $k$ such that $k\in\initialSegment{1}$ and $(k,y)\in s$.
                    Thus $k\in\naturalsPlus$ and $k \leq 1$ by \cref{initial_segment_naturalsplus}.
                    Thus $1 \leq k$ by \cref{real_one_leq_naturalsplus}.
                    $k\in\reals$ and $1\in\reals$ by \cref{real_naturals_subset_reals,subseteq}.
                    Thus $k = 1$ by \cref{real_leq_antisym}.
                    $s$ is a function by \cref{funs_is_function}.
                    Thus $s(1) = y$ by \cref{function_apply_intro}.
                    Thus $y = r_1$.
                    Thus $y\in\{r_1\}$ by \cref{singleton_intro}.
                \end{subproof}
                Thus $\img{s}{\initialSegment{1}}\subseteq \{r_1\}$ by \cref{subseteq}.
            \end{subproof}
            Show that $\{r_1\}\subseteq \img{s}{\initialSegment{1}}$.
            \begin{subproof}
                $1\leq 1$.
                Thus $1\in\initialSegment{1}$ by \cref{initial_segment_naturalsplus,real_one_in_naturals}.
                $s$ is a function by \cref{funs_is_function}.
                Thus $\dom{s} = \naturalsPlus$ by \cref{funs_elim}.
                Thus $1\in\dom{s}$.
                Thus $(1,r_1)\in s$ by \cref{function_apply_elim}.
                Thus $r_1\in \img{s}{\initialSegment{1}}$ by \cref{img_iff}.
                Thus $\{r_1\}\subseteq \img{s}{\initialSegment{1}}$ by \cref{singleton_subset_intro}.
            \end{subproof}
            Thus $\img{s}{\initialSegment{1}} = \{r_1\}$ by \cref{subseteq_antisymmetric}.
        \end{subproof}
        \begin{byCase}
        \caseOf{$r_1 = 1$.}
            Take $u_1$ such that $(1,u_1)\in G$ by \cref{urysohn_u1_eq1}.
        \caseOf{$r_1 \neq 1$.}
            Show that $r_1\in\reals$.
            \begin{subproof}
                Take $q$ such that
                    $q\in\naturalsPlus\times\naturalsPlus$ and
                    $r_1 = \fst{q} \rmul \inv{\snd{q}}$
                    by \cref{unit_rationals}.
                Take $m,n$ such that $q = (m,n)$ and $m\in\naturalsPlus$ and $n\in\naturalsPlus$
                    by \cref{times}.
                Thus $\fst{q} = m$ and $\snd{q} = n$ by \cref{fst_eq,snd_eq}.
                Thus $m\in\reals$ and $n\in\reals$ by \cref{real_naturals_subset_reals,subseteq}.
                Show that $n \neq 0$.
                \begin{subproof}
                    $1\in\reals$ by \cref{real_mul_ident}.
                    Thus $1 < n$ or $1 = n$ by \cref{real_one_leq_naturalsplus}.
                    \begin{byCase}
                    \caseOf{$1 < n$.}
                        Suppose not.
                        Thus $n = 0$.
                        Thus $1 < 0$ by \cref{real_lt_subst_right}.
                        $0 < 1$ by \cref{real_zero_lt_one}.
                        Thus $1 < 1$ by \cref{real_lt_trans}.
                        Contradiction by \cref{real_lt_irreflexive}.
                    \caseOf{$1 = n$.}
                        Thus $n \neq 0$ by \cref{real_mul_ident}.
                    \end{byCase}
                \end{subproof}
                Thus $\inv{n}\in\reals$ by \cref{real_mul_inverse}.
                Thus $r_1\in\reals$ by \cref{real_mul_closed}.
            \end{subproof}
            $1\in\reals$ by \cref{real_mul_ident}.
            Thus $r_1 < 1$ or $1 < r_1$ by \cref{real_lt_trichotomy}.
            \begin{byCase}
            \caseOf{$1 < r_1$.}
                Take $u_1$ such that $(1,u_1)\in G$ by \cref{urysohn_u1_gt1}.
            \caseOf{$r_1 < 1$.}
                Take $u_1$ such that $(1,u_1)\in G$ by \cref{urysohn_u1_lt1}.
            \end{byCase}
        \end{byCase}
    \end{subproof}
    Thus there exists $u_1$ such that $(1,u_1)\in G$.
\end{proof}

% from §33 helper lemma 17: indices of the iteration sequence
\begin{lemma}\label{urysohn_unit_interval_fst_u}
    Assume $s$ is a sequence in $\unitRationals$.
    Let $G = \urysohnGoodPairs{s}{X}{T}{A}{B}$.
    Let $R = \{(z,w) \mid z\in G, w\in G \mid
        \fst{w} = \fst{z} + 1 \land \snd{z}\subseteq\snd{w}\}$.
    Assume $c$ is a function on $G$.
    Assume for all $z\in G$ we have $z\mathrel{R} \apply{c}{z}$.
    Assume $u$ is a sequence in $G$.
    Assume $u(1) = a_0$.
    Assume $a_0 = (1,u_1)$.
    Assume for all $n\in\naturalsPlus$ we have $u(n + 1) = \apply{c}{u(n)}$.
    Then for all $n\in\naturalsPlus$ we have $\fst{u(n)} = n$.
\end{lemma}
\begin{proof}
    Show that for all $n\in\naturalsPlus$ we have $\fst{u(n)} = n$.
    \begin{subproof}
        Let $A_0 = \{n\in\naturalsPlus \mid \fst{u(n)} = n\}$.
        $A_0\subseteq\naturalsPlus$ by \cref{subseteq}.
        Show that $1\in A_0$.
        \begin{subproof}
            $1\in\naturalsPlus$ by \cref{real_one_in_naturals}.
            Thus $u(1) = a_0$.
            Thus $u(1) = (1,u_1)$.
            Thus $\fst{u(1)} = 1$ by \cref{fst_eq}.
            Thus $1\in A_0$.
        \end{subproof}
        Show that for all $k\in A_0$ we have $k + 1\in A_0$.
        \begin{subproof}
            Fix $k$.
            Assume $k\in A_0$.
            Thus $k\in\naturalsPlus$ by \cref{subseteq}.
            Thus $u$ is a sequence in $G$.
            Thus $u\in\funs{\naturalsPlus}{G}$ by \cref{sequence_in}.
            Thus $u(k)\in G$ by \cref{funs_type_apply}.
            Thus $u(k)\mathrel{R} u(k + 1)$.
            Thus $(u(k),u(k + 1))\in R$.
            Thus there exist $z,w$ such that
                $z\in G$ and
                $w\in G$ and
                $(z,w) = (u(k),u(k + 1))$ and
                $\fst{w} = \fst{z} + 1$.
            Thus $\fst{u(k + 1)} = \fst{u(k)} + 1$ by \cref{pair_eq_iff}.
            Thus $\fst{u(k + 1)} = k + 1$.
            Thus $k + 1\in A_0$.
        \end{subproof}
        Thus for all $n\in\naturalsPlus$ we have $n\in A_0$ by \cref{real_naturals_induction}.
        Thus for all $n\in\naturalsPlus$ we have $\fst{u(n)} = n$.
    \end{subproof}
\end{proof}

% from §33 helper lemma 18: continuous subspace range
\begin{lemma}\label{urysohn_unit_interval_continuous_subspace}
    Assume $f\in\funs{X}{\reals}$.
    Assume $T$ is a topology on $X$.
    Assume $I_0\subseteq\reals$.
    Let $T_0 = \subspaceTopology{\standardTopologyR}{I_0}$.
    Assume $f$ is continuous from $X$ to $\reals$ with respect to $T$ and $\standardTopologyR$.
    Assume $\img{f}{X}\subseteq I_0$.
    Then $f$ is continuous from $X$ to $I_0$ with respect to $T$ and $T_0$.
\end{lemma}
\begin{proof}
    Thus $f$ is continuous from $X$ to $I_0$ with respect to $T$ and $T_0$
        by \cref{continuous_subspace_range}.
\end{proof}


% from §33 helper lemma 19: continuity at zero
\begin{lemma}\label{urysohn_unit_interval_continuous_zero}
    Assume $X$ is normal with respect to $T$ and $A$ is closed in $X$ with respect to $T$ and
        $B$ is closed in $X$ with respect to $T$ and $A$ is disjoint from $B$.
    Assume $s$ is a sequence in $\unitRationals$ and for all $p$ such that $p\in\unitRationals$
        there exists $n\in\naturalsPlus$ such that $s(n) = p$.
    Let $G = \urysohnGoodPairs{s}{X}{T}{A}{B}$.
    Let $R = \{(z,w) \mid z\in G, w\in G \mid
        \fst{w} = \fst{z} + 1 \land \snd{z}\subseteq\snd{w}\}$.
    Assume $c$ is a function on $G$ and for all $z\in G$ we have $z\mathrel{R} \apply{c}{z}$.
    Assume $u$ is a sequence in $G$ and $u(1) = a_0$ and $a_0 = (1,u_1)$ and
        for all $n\in\naturalsPlus$ we have $u(n + 1) = \apply{c}{u(n)}$ and $\fst{u(n)} = n$.
    Let $F = \{\snd{u(n)} \mid n\in\naturalsPlus\}$.
    Let $U = \unions{F}$.
    Assume $U$ is a function.
    Let $I_0 = \intervalclosed{0}{1}$.
    Let $f = \{(x,y) \mid x\in X, y\in\reals \mid
        \text{$y$ is an infimum of $\urysohnLevelSet{U}{x}$}\}$.
    Assume $f\in\funs{X}{\reals}$.
    Assume $\img{f}{X}\subseteq I_0$.
    Assume $x\in X$.
    Assume $V$ is open in $\reals$ with respect to $\standardTopologyR$.
    Assume $a\in\reals$ and $b\in\reals$ and $a < b$.
    Assume $f(x)\in\intervalopen{a}{b}$.
    Assume $\intervalopen{a}{b}\subseteq V$.
    Assume $f(x) = 0$.
    Then there exists $U$ such that
        $U$ is open in $X$ with respect to $T$ and $x\in U$ and $\img{f}{U}\subseteq V$.
\end{lemma}
\begin{proof}
    Let $Q_x = \urysohnLevelSet{U}{x}$.
    Thus $f(x) < b$ by \cref{intervalopen}.
    Thus $0 < b$ by \cref{real_lt_subst_left}.
    $0\in\reals$ by \cref{real_add_ident}.
    Take $q$ such that $q\in\unitRationals$ and $q\in\intervalopen{0}{b}$
        by \cref{unit_rationals_dense_nonneg}.
    Thus $q\in\reals$ by \cref{intervalopen}.
    Thus $0 < q$ and $q < b$ by \cref{intervalopen}.
    Thus $q\in\unitRationalsZero$ by \cref{unit_rationals_zero,union_intro_left}.
    Let $U = U(q)$.
    Show that $U$ is open in $X$ with respect to $T$ and $x\in U$ and
        $\img{f}{U}\subseteq V$.
    \begin{subproof}
        Show that $U$ is a neighborhood of $x$ in $X$ with respect to $T$.
        \begin{subproof}
            Show that $U$ is open in $X$ with respect to $T$.
            \begin{subproof}
                Show that $U(q)\in T$.
                \begin{subproof}
                    Thus $q\in\unitRationals$.
                    Take $k$ such that $k\in\naturalsPlus$ and $s(k) = q$
                        by \cref{unit_rationals_sequence}.
                    Thus $u$ is a sequence in $G$.
                    Thus $u\in\funs{\naturalsPlus}{G}$ by \cref{sequence_in}.
                    Thus $u(k)\in G$ by \cref{funs_type_apply}.
                    Take $n,t$ such that
                        $n\in\naturalsPlus$ and
                        $t\in\urysohnAdmissible{s}{X}{T}{A}{B}{n}$ and
                        $u(k) = (n,t)$
                        by \cref{urysohn_good_pair_decomp}.
                    Thus $\fst{u(k)} = n$ by \cref{fst_eq}.
                    Thus $k = n$.
                    Thus $u(k) = (k,t)$ by \cref{pair_eq_iff}.
                    Thus $\snd{u(k)} = t$ by \cref{snd_eq}.
                    Thus $t\in F$.
                    Thus $t\in\funs{\{0,1\}\union \img{s}{\initialSegment{k}}}{\pow{X}}$
                        by \cref{urysohn_admissible}.
                    Thus $t$ is a function by \cref{funs_is_function}.
                    Thus $t$ is right-unique and $t$ is a relation.
                    Thus $\dom{t} = \{0,1\}\union \img{s}{\initialSegment{k}}$
                        by \cref{funs_elim}.
                    Show that $q\in\dom{t}$.
                    \begin{subproof}
                        Thus $k \leq k$.
                        Thus $k\in\initialSegment{k}$ by \cref{initial_segment_naturalsplus}.
                        Let $y = k$.
                        Thus $y\in\initialSegment{k}$.
                        Thus $s(y) = s(k)$.
                        Thus $s$ is a sequence in $\unitRationals$.
                        Thus $s\in\funs{\naturalsPlus}{\unitRationals}$ by \cref{sequence_in}.
                        Thus $s$ is a function by \cref{funs_is_function}.
                        Thus $\dom{s} = \naturalsPlus$ by \cref{funs_elim}.
                        Thus $y\in\dom{s}$ by \cref{subseteq}.
                        Thus $(y,s(y))\in s$ by \cref{function_apply_elim}.
                        Thus $(y,s(k))\in s$.
                        Thus $s(k)\in\img{s}{\initialSegment{k}}$ by \cref{img_iff}.
                        Thus $q\in\img{s}{\initialSegment{k}}$.
                        Thus $q\in\dom{t}$ by \cref{union_intro_right}.
                    \end{subproof}
                    Thus $t(q)\in T$ by \cref{urysohn_admissible,subseteq}.
                    Thus $(q,t(q))\in t$ by \cref{function_apply_elim}.
                    Thus $(q,t(q))\in U$ by \cref{unions_intro}.
                    Thus $U(q) = t(q)$ by \cref{function_apply_intro}.
                    Thus $U(q)\in T$.
                \end{subproof}
                Thus $U\in T$.
                $T$ is a topology on $X$ by \cref{normal_space}.
                Thus $U$ is open in $X$ with respect to $T$ by \cref{open_in}.
            \end{subproof}
            Show that $x\in U$.
            \begin{subproof}
                Show that $x\in U(q)$.
                \begin{subproof}
                    Thus $0 < q$.
                    Thus $f(x) < q$ by \cref{real_lt_subst_right}.
                    Show that $q$ is not a lower bound of $Q_x$.
                    \begin{subproof}
                        Show that $f(x)$ is an infimum of $Q_x$.
                        \begin{subproof}
                            Thus $f$ is a function by \cref{funs_is_function}.
                            Thus $\dom{f} = X$ by \cref{funs_elim}.
                            Thus $x\in\dom{f}$ by \cref{subseteq}.
                            Thus $(x,f(x))\in f$ by \cref{function_apply_elim}.
                            Take $x_0,y_0$ such that
                                $x_0\in X$ and
                                $y_0\in\reals$ and
                                $(x,f(x)) = (x_0,y_0)$ and
                                $y_0$ is an infimum of $\urysohnLevelSet{U}{x_0}$.
                            Thus $x = x_0$ and $f(x) = y_0$ by \cref{pair_eq_iff}.
                            Thus $y_0$ is an infimum of $Q_x$.
                            Thus $f(x)$ is an infimum of $Q_x$.
                        \end{subproof}
                        Thus $f(x)$ is a lower bound of $Q_x$ by \cref{real_infimum}.
                        Thus $f(x) \leq q$ by \cref{real_lower_bound}.
                        Thus $f(x) < q$ by \cref{real_lt_subst_left}.
                        Thus it is wrong that $q$ is a lower bound of $Q_x$.
                    \end{subproof}
                    Show that there exists $r$ such that $r\in Q_x$ and $r < q$.
                    \begin{subproof}
                        Suppose not.
                        Show that $q$ is a lower bound of $Q_x$.
                        \begin{subproof}
                            Thus $q\in\reals$.
                            Show that $Q_x\subseteq\reals$.
                            \begin{subproof}
                                Show that for all $r\in Q_x$ we have $r\in\reals$.
                                \begin{subproof}
                                    Fix $r$.
                                    Assume $r\in Q_x$.
                                    Thus $r\in\unitRationalsZero$ by \cref{urysohn_level_set}.
                                    Thus $r\in\unitRationals$ or $r\in\{0\}$
                                        by \cref{unit_rationals_zero,union_iff}.
                                    \begin{byCase}
                                    \caseOf{$r\in\unitRationals$.}
                                        Take $t$ such that $t\in\naturalsPlus\times\naturalsPlus$ and
                                            $r = \fst{t} \rmul \inv{\snd{t}}$
                                            by \cref{unit_rationals}.
                                        Thus $\fst{t}\in\naturalsPlus$ and $\snd{t}\in\naturalsPlus$
                                            by \cref{fst_type,snd_type}.
                                        Thus $\fst{t}\in\reals$ and $\snd{t}\in\reals$
                                            by \cref{real_naturals_subset_reals,subseteq}.
                                        Show that $\snd{t}\neq 0$.
                                        \begin{subproof}
                                            $1\in\reals$ by \cref{real_mul_ident}.
                                            Thus $1 < \snd{t}$ or $1 = \snd{t}$
                                                by \cref{real_one_leq_naturalsplus}.
                                            $0\in\reals$ by \cref{real_add_ident}.
                                            $0 < 1$ by \cref{real_zero_lt_one}.
                                            \begin{byCase}
                                            \caseOf{$1 < \snd{t}$.}
                                                Thus $0 < \snd{t}$ by \cref{real_lt_trans}.
                                                Thus $\snd{t}\neq 0$ by \cref{real_pos_nonzero}.
                                            \caseOf{$1 = \snd{t}$.}
                                                Thus $0 < \snd{t}$ by \cref{real_lt_subst_right}.
                                                Thus $\snd{t}\neq 0$ by \cref{real_pos_nonzero}.
                                            \end{byCase}
                                        \end{subproof}
                                        Thus $\inv{\snd{t}}\in\reals$ by \cref{real_mul_inverse}.
                                        Thus $r\in\reals$ by \cref{real_mul_closed}.
                                    \caseOf{$r\in\{0\}$.}
                                        Thus $r = 0$ by \cref{singleton_elim}.
                                        $0\in\reals$ by \cref{real_add_ident}.
                                        Thus $r\in\reals$.
                                    \end{byCase}
                                \end{subproof}
                                Thus $Q_x\subseteq\reals$ by \cref{subseteq}.
                            \end{subproof}
                            Show that for all $r\in Q_x$ we have $q < r$ or $q = r$.
                            \begin{subproof}
                                Fix $r$.
                                Assume $r\in Q_x$.
                                Thus $r\in\reals$ by \cref{subseteq}.
                                Show that $q < r$ or $q = r$.
                                \begin{subproof}
                                    Suppose not.
                                    Thus $r < q$ or $r = q$ or $q < r$
                                        by \cref{real_lt_trichotomy}.
                                    Thus $r < q$.
                                    Thus there exists $r$ such that $r\in Q_x$ and $r < q$.
                                    Contradiction.
                                \end{subproof}
                            \end{subproof}
                            Thus $q$ is a lower bound of $Q_x$ by \cref{real_lower_bound}.
                        \end{subproof}
                        Contradiction.
                    \end{subproof}
                    Take $r$ such that $r\in Q_x$ and $r < q$.
                    Thus $x\in U(r)$.
                    Show that $x\in U(q)$.
                    \begin{subproof}
                        Thus $r\in\unitRationalsZero$ by \cref{urysohn_level_set}.
                        Thus $r\in\unitRationals$ or $r\in\{0\}$
                            by \cref{unit_rationals_zero,union_iff}.
                        Thus $q\in\unitRationals$.
                        Take $k_q$ such that $k_q\in\naturalsPlus$ and $s(k_q) = q$
                            by \cref{unit_rationals_sequence}.
                        Thus $u$ is a sequence in $G$.
                        Thus $u\in\funs{\naturalsPlus}{G}$ by \cref{sequence_in}.
                        Thus $u(k_q)\in G$ by \cref{funs_type_apply}.
                        Take $n_q,t_q$ such that
                            $n_q\in\naturalsPlus$ and
                            $t_q\in\urysohnAdmissible{s}{X}{T}{A}{B}{n_q}$ and
                            $u(k_q) = (n_q,t_q)$
                            by \cref{urysohn_good_pair_decomp}.
                        Thus $\fst{u(k_q)} = n_q$ by \cref{fst_eq}.
                        Thus $\fst{u(k_q)} = k_q$.
                        Thus $k_q = n_q$.
                        Thus $u(k_q) = (k_q,t_q)$ by \cref{pair_eq_iff}.
                        Thus $\snd{u(k_q)} = t_q$ by \cref{snd_eq}.
                        Thus $t_q\in F$.
                        Thus $t_q\in\funs{\{0,1\}\union \img{s}{\initialSegment{k_q}}}{\pow{X}}$
                            by \cref{urysohn_admissible}.
                        Thus $t_q$ is a function by \cref{funs_is_function}.
                        Thus $t_q$ is right-unique and $t_q$ is a relation.
                        Thus $\dom{t_q} = \{0,1\}\union \img{s}{\initialSegment{k_q}}$
                            by \cref{funs_elim}.
                        Show that $q\in\dom{t_q}$.
                        \begin{subproof}
                            Thus $k_q \leq k_q$.
                            Thus $k_q\in\initialSegment{k_q}$ by \cref{initial_segment_naturalsplus}.
                            Let $y = k_q$.
                            Thus $y\in\initialSegment{k_q}$.
                            Thus $s(y) = s(k_q)$.
                            Thus $s$ is a sequence in $\unitRationals$.
                            Thus $s\in\funs{\naturalsPlus}{\unitRationals}$ by \cref{sequence_in}.
                            Thus $s$ is a function by \cref{funs_is_function}.
                            Thus $\dom{s} = \naturalsPlus$ by \cref{funs_elim}.
                            Thus $y\in\dom{s}$ by \cref{subseteq}.
                            Thus $(y,s(y))\in s$ by \cref{function_apply_elim}.
                            Thus $(y,s(k_q))\in s$.
                            Thus $s(k_q)\in\img{s}{\initialSegment{k_q}}$ by \cref{img_iff}.
                            Thus $q\in\img{s}{\initialSegment{k_q}}$.
                            Thus $q\in\dom{t_q}$ by \cref{union_intro_right}.
                        \end{subproof}
                        Thus $(q,t_q(q))\in t_q$ by \cref{function_apply_elim}.
                        Thus $(q,t_q(q))\in U$ by \cref{unions_intro}.
                        Thus $U(q) = t_q(q)$ by \cref{function_apply_intro}.
                        \begin{byCase}
                        \caseOf{$r\in\unitRationals$.}
                            Show that $x\in U(q)$.
                            \begin{subproof}
                                Take $k_r$ such that $k_r\in\naturalsPlus$ and $s(k_r) = r$
                                    by \cref{unit_rationals_sequence}.
                                Thus $u$ is a sequence in $G$.
                                Thus $u\in\funs{\naturalsPlus}{G}$ by \cref{sequence_in}.
                                Thus $u(k_r)\in G$ by \cref{funs_type_apply}.
                                Take $n_r,t_r$ such that
                                    $n_r\in\naturalsPlus$ and
                                    $t_r\in\urysohnAdmissible{s}{X}{T}{A}{B}{n_r}$ and
                                    $u(k_r) = (n_r,t_r)$
                                    by \cref{urysohn_good_pair_decomp}.
                                Thus $\fst{u(k_r)} = n_r$ by \cref{fst_eq}.
                                Thus $\fst{u(k_r)} = k_r$.
                                Thus $k_r = n_r$.
                                Thus $u(k_r) = (k_r,t_r)$ by \cref{pair_eq_iff}.
                                Thus $\snd{u(k_r)} = t_r$ by \cref{snd_eq}.
                                Thus $t_r\in F$.
                                Thus $t_r\in\funs{\{0,1\}\union \img{s}{\initialSegment{k_r}}}{\pow{X}}$
                                    by \cref{urysohn_admissible}.
                                Thus $t_r$ is a function by \cref{funs_is_function}.
                                Thus $t_r$ is right-unique and $t_r$ is a relation.
                                Thus $\dom{t_r} = \{0,1\}\union \img{s}{\initialSegment{k_r}}$
                                    by \cref{funs_elim}.
                                Show that $r\in\dom{t_r}$.
                                \begin{subproof}
                                    Thus $k_r \leq k_r$.
                                    Thus $k_r\in\initialSegment{k_r}$
                                        by \cref{initial_segment_naturalsplus}.
                                    Let $y = k_r$.
                                    Thus $y\in\initialSegment{k_r}$.
                                    Thus $s(y) = s(k_r)$.
                                    Thus $s$ is a sequence in $\unitRationals$.
                                    Thus $s\in\funs{\naturalsPlus}{\unitRationals}$ by \cref{sequence_in}.
                                    Thus $s$ is a function by \cref{funs_is_function}.
                                    Thus $\dom{s} = \naturalsPlus$ by \cref{funs_elim}.
                                    Thus $y\in\dom{s}$ by \cref{subseteq}.
                                    Thus $(y,s(y))\in s$ by \cref{function_apply_elim}.
                                    Thus $(y,s(k_r))\in s$.
                                    Thus $s(k_r)\in\img{s}{\initialSegment{k_r}}$ by \cref{img_iff}.
                                    Thus $r\in\img{s}{\initialSegment{k_r}}$.
                                    Thus $r\in\dom{t_r}$ by \cref{union_intro_right}.
                            \end{subproof}
                            Thus $(r,t_r(r))\in t_r$ by \cref{function_apply_elim}.
                            Thus $t_r\subseteq t_q$ or $t_q\subseteq t_r$
                                by \cref{urysohn_unit_interval_chain}.
                            \begin{byCase}
                            \caseOf{$t_r\subseteq t_q$.}
                                    Thus $(r,t_r(r))\in t_q$ by \cref{subseteq}.
                                    Thus $r\in\dom{t_q}$ by \cref{dom_intro}.
                                    Thus $\closure{t_q(r)}{X}{T}\subseteq t_q(q)$
                                        by \cref{urysohn_admissible}.
                                    Thus $(r,t_q(r))\in t_q$ by \cref{function_apply_elim}.
                                    Thus $(r,t_q(r))\in U$ by \cref{unions_intro}.
                                    Thus $U(r) = t_q(r)$ by \cref{function_apply_intro}.
                                    Thus $t_q(r)\in\pow{X}$ by \cref{funs_type_apply}.
                                    Thus $t_q(r)\subseteq X$ by \cref{powerset_elim}.
                                    Thus $U(r)\subseteq X$.
                                    Thus $\closure{U(r)}{X}{T}\subseteq U(q)$.
                                    Thus $U(r)\subseteq\closure{U(r)}{X}{T}$
                                        by \cref{subseteq_closure,normal_space}.
                                    Thus $x\in\closure{U(r)}{X}{T}$ by \cref{subseteq}.
                                    Thus $x\in U(q)$ by \cref{subseteq}.
                            \caseOf{$t_q\subseteq t_r$.}
                                    Thus $(q,t_q(q))\in t_r$ by \cref{subseteq}.
                                    Thus $q\in\dom{t_r}$ by \cref{dom_intro}.
                                    Thus $\closure{t_r(r)}{X}{T}\subseteq t_r(q)$
                                        by \cref{urysohn_admissible}.
                                    Thus $(r,t_r(r))\in t_r$ by \cref{function_apply_elim}.
                                    Thus $(r,t_r(r))\in U$ by \cref{unions_intro}.
                                    Thus $U(r) = t_r(r)$ by \cref{function_apply_intro}.
                                    Thus $t_r(r)\in\pow{X}$ by \cref{funs_type_apply}.
                                    Thus $t_r(r)\subseteq X$ by \cref{powerset_elim}.
                                    Thus $U(r)\subseteq X$.
                                    Thus $(q,t_r(q))\in t_r$ by \cref{function_apply_elim}.
                                    Thus $(q,t_r(q))\in U$ by \cref{unions_intro}.
                                    Thus $U(q) = t_r(q)$ by \cref{function_apply_intro}.
                                    Thus $\closure{U(r)}{X}{T}\subseteq U(q)$.
                                    Thus $U(r)\subseteq\closure{U(r)}{X}{T}$
                                        by \cref{subseteq_closure,normal_space}.
                                    Thus $x\in\closure{U(r)}{X}{T}$ by \cref{subseteq}.
                                    Thus $x\in U(q)$ by \cref{subseteq}.
                            \end{byCase}
                        \end{subproof}
                        \caseOf{$r\in\{0\}$.}
                            Thus $r = 0$ by \cref{singleton_elim}.
                            $0\in\{0,1\}$ by \cref{upair_intro_left}.
                            Thus $0\in\dom{t_q}$ by \cref{union_intro_left,subseteq}.
                            Thus $\closure{t_q(0)}{X}{T}\subseteq t_q(q)$
                                by \cref{urysohn_admissible}.
                            Thus $(0,t_q(0))\in t_q$ by \cref{function_apply_elim}.
                            Thus $(0,t_q(0))\in U$ by \cref{unions_intro}.
                            Thus $U(r) = t_q(0)$ by \cref{function_apply_intro}.
                            Thus $t_q(0)\in\pow{X}$ by \cref{funs_type_apply}.
                            Thus $t_q(0)\subseteq X$ by \cref{powerset_elim}.
                            Thus $U(r)\subseteq X$.
                            Thus $\closure{U(r)}{X}{T}\subseteq U(q)$.
                            Thus $U(r)\subseteq\closure{U(r)}{X}{T}$
                                by \cref{subseteq_closure,normal_space}.
                            Thus $x\in\closure{U(r)}{X}{T}$ by \cref{subseteq}.
                            Thus $x\in U(q)$ by \cref{subseteq}.
                        \end{byCase}
                    \end{subproof}
                \end{subproof}
                Thus $x\in U$.
            \end{subproof}
            Thus $U$ is a neighborhood of $x$ in $X$ with respect to $T$
                by \cref{neighborhood}.
        \end{subproof}
        Thus $U$ is open in $X$ with respect to $T$ by \cref{neighborhood}.
        Thus $x\in U$ by \cref{neighborhood}.
        Show that $\img{f}{U}\subseteq V$.
        \begin{subproof}
            Show that for all $y\in U$ we have $f(y)\in V$.
            \begin{subproof}
                Fix $y$.
                Assume $y\in U$.
                Let $Q_y = \urysohnLevelSet{U}{y}$.
                Thus $y\in U(q)$.
                Thus $q\in\unitRationalsZero$
                    by \cref{unit_rationals_zero,union_intro_left}.
                Thus $q\in Q_y$ by \cref{urysohn_level_set}.
                Show that $f(y) \leq q$.
                \begin{subproof}
                    Show that $f(y)$ is a lower bound of $Q_y$.
                    \begin{subproof}
                        Show that $f(y)$ is an infimum of $Q_y$.
                        \begin{subproof}
                            Thus $f$ is a function by \cref{funs_is_function}.
                            Thus $\dom{f} = X$ by \cref{funs_elim}.
                            Thus $U$ is open in $X$ with respect to $T$
                                by \cref{neighborhood}.
                            Thus $T$ is a topology on $X$ by \cref{open_in}.
                            Thus $U\in T$ by \cref{open_in}.
                            Thus $T\subseteq\pow{X}$ by \cref{topology_on}.
                            Thus $U\in\pow{X}$ by \cref{subseteq}.
                            Thus $U\subseteq X$ by \cref{powerset_elim}.
                            Thus $y\in X$ by \cref{subseteq}.
                            Thus $y\in\dom{f}$ by \cref{subseteq}.
                            Thus $(y,f(y))\in f$ by \cref{function_apply_elim}.
                            Take $x_0,y_0$ such that
                                $x_0\in X$ and
                                $y_0\in\reals$ and
                                $(y,f(y)) = (x_0,y_0)$ and
                                $y_0$ is an infimum of $\urysohnLevelSet{U}{x_0}$.
                            Thus $y = x_0$ and $f(y) = y_0$ by \cref{pair_eq_iff}.
                            Thus $f(y)$ is an infimum of $Q_y$.
                        \end{subproof}
                        Thus $f(y)$ is a lower bound of $Q_y$ by \cref{real_infimum}.
                    \end{subproof}
                    Thus $f(y) \leq q$ by \cref{real_lower_bound}.
                \end{subproof}
                Thus $U$ is open in $X$ with respect to $T$
                    by \cref{neighborhood}.
                Thus $T$ is a topology on $X$ by \cref{open_in}.
                Thus $U\in T$ by \cref{open_in}.
                Thus $T\subseteq\pow{X}$ by \cref{topology_on}.
                Thus $U\in\pow{X}$ by \cref{subseteq}.
                Thus $U\subseteq X$ by \cref{powerset_elim}.
                Thus $y\in X$ by \cref{subseteq}.
                Thus $f(y)\in\reals$ by \cref{funs_type_apply}.
                Thus $f$ is a function by \cref{funs_is_function}.
                Thus $\dom{f} = X$ by \cref{funs_elim}.
                Thus $y\in\dom{f}$ by \cref{subseteq}.
                Thus $(y,f(y))\in f$ by \cref{function_apply_elim}.
                Thus $f(y)\in\img{f}{X}$ by \cref{img_iff}.
                Thus $f(y)\in I_0$ by \cref{subseteq}.
                Thus $0 \leq f(y)$ by \cref{intervalclosed}.
                Thus $a < f(x)$ by \cref{intervalopen}.
                Thus $a < 0$ by \cref{real_lt_subst_right}.
                Show that $a < f(y)$.
                \begin{subproof}
                    Thus $0 < f(y)$ or $0 = f(y)$.
                    \begin{byCase}
                    \caseOf{$0 < f(y)$.}
                            Thus $a < f(y)$ by \cref{real_lt_trans}.
                    \caseOf{$0 = f(y)$.}
                            Thus $a < f(y)$ by \cref{real_lt_subst_right}.
                    \end{byCase}
                \end{subproof}
                Thus $f(y) < b$ by \cref{real_lt_trans}.
                Thus $f(y)\in\intervalopen{a}{b}$ by \cref{intervalopen}.
                Thus $f(y)\in V$ by \cref{subseteq}.
            \end{subproof}
            Show that for all $z\in\img{f}{U}$ we have $z\in V$.
            \begin{subproof}
                Fix $z$.
                Assume $z\in\img{f}{U}$.
                Take $y$ such that $y\in U$ and $y\mathrel{f} z$ by \cref{img_iff}.
                Thus $f$ is a function by \cref{funs_is_function}.
                Thus $f(y) = z$ by \cref{function_apply_intro}.
                Thus $f(y)\in V$.
                Thus $z\in V$.
            \end{subproof}
            Thus $\img{f}{U}\subseteq V$ by \cref{subseteq}.
        \end{subproof}
    \end{subproof}
    Thus there exists $U$ such that
        $U$ is open in $X$ with respect to $T$ and $x\in U$ and $\img{f}{U}\subseteq V$.
\end{proof}

% from §33 helper lemma 20: positive case produces open neighborhood
\begin{lemma}\label{urysohn_unit_interval_continuous_pos_open}
    Assume $X$ is normal with respect to $T$.
    Assume $A$ is closed in $X$ with respect to $T$.
    Assume $B$ is closed in $X$ with respect to $T$.
    Assume $A$ is disjoint from $B$.
    Assume $s$ is a sequence in $\unitRationals$.
    Assume for all $p$ such that $p\in\unitRationals$ there exists $n\in\naturalsPlus$ such that $s(n) = p$.
    Let $G = \urysohnGoodPairs{s}{X}{T}{A}{B}$.
    Let $R = \{(z,w) \mid z\in G, w\in G \mid
        \fst{w} = \fst{z} + 1 \land \snd{z}\subseteq\snd{w}\}$.
    Assume $c$ is a function on $G$.
    Assume for all $z\in G$ we have $z\mathrel{R} \apply{c}{z}$.
    Assume $u$ is a sequence in $G$.
    Assume $u(1) = a_0$.
    Assume $a_0 = (1,u_1)$.
    Assume for all $n\in\naturalsPlus$ we have $u(n + 1) = \apply{c}{u(n)}$.
    Assume for all $n\in\naturalsPlus$ we have $\fst{u(n)} = n$.
    Let $F = \{\snd{u(n)} \mid n\in\naturalsPlus\}$.
    Let $U = \unions{F}$.
    Assume $U$ is a function.
    Let $I_0 = \intervalclosed{0}{1}$.
    Let $f = \{(x,y) \mid x\in X, y\in\reals \mid
        \text{$y$ is an infimum of $\urysohnLevelSet{U}{x}$}\}$.
    Assume $f\in\funs{X}{\reals}$.
    Assume $\img{f}{X}\subseteq I_0$.
    Assume $x\in X$.
    Assume $V$ is open in $\reals$ with respect to $\standardTopologyR$.
    Assume $a\in\reals$ and $b\in\reals$ and $a < b$.
    Assume $f(x)\in\intervalopen{a}{b}$.
    Assume $\intervalopen{a}{b}\subseteq V$.
    Assume $0 < f(x)$.
    Then there exist $p,q,r_0,W$ such that
        $p\in\unitRationals$ and
        $p\in\intervalopen{a}{f(x)}$ and
        $p\in\intervalopen{0}{f(x)}$ and
        $q\in\unitRationals$ and
        $q\in\intervalopen{f(x)}{b}$ and
        $r_0\in\unitRationals$ and
        $r_0\in\intervalopen{p}{f(x)}$ and
        $p\in\unitRationalsZero$ and
        $q\in\unitRationalsZero$ and
        $r_0\in\unitRationalsZero$ and
        $W = U(q)\setminus \closure{U(p)}{X}{T}$ and
        $W$ is open in $X$ with respect to $T$.
\end{lemma}
\begin{proof}
    Show that there exists $p$ such that
        $p\in\unitRationals$ and
        $p\in\intervalopen{a}{f(x)}$ and
        $p\in\intervalopen{0}{f(x)}$.
    \begin{subproof}
        $0\in\reals$ by \cref{real_add_ident}.
        Thus $f(x)\in\reals$ by \cref{funs_type_apply}.
        Thus $a < 0$ or $a = 0$ or $0 < a$ by \cref{real_lt_trichotomy}.
        \begin{byCase}
        \caseOf{$a < 0$.}
            Take $p_0$ such that $p_0\in\unitRationals$ and $p_0\in\intervalopen{0}{f(x)}$
                by \cref{unit_rationals_dense_nonneg}.
            Thus $p_0\in\reals$ by \cref{intervalopen}.
            Thus $0 < p_0$ and $p_0 < f(x)$ by \cref{intervalopen}.
            Thus $a < p_0$ by \cref{real_lt_trans}.
            Thus $p_0\in\intervalopen{a}{f(x)}$ by \cref{intervalopen}.
            Thus there exists $p$ such that
                $p\in\unitRationals$ and
                $p\in\intervalopen{a}{f(x)}$ and
                $p\in\intervalopen{0}{f(x)}$.
        \caseOf{$a = 0$.}
            Take $p_0$ such that $p_0\in\unitRationals$ and $p_0\in\intervalopen{0}{f(x)}$
                by \cref{unit_rationals_dense_nonneg}.
            Thus $p_0\in\reals$ by \cref{intervalopen}.
            Thus $0 < p_0$ and $p_0 < f(x)$ by \cref{intervalopen}.
            Thus $a < p_0$ by \cref{real_lt_subst_left}.
            Thus $p_0\in\intervalopen{a}{f(x)}$ by \cref{intervalopen}.
            Thus there exists $p$ such that
                $p\in\unitRationals$ and
                $p\in\intervalopen{a}{f(x)}$ and
                $p\in\intervalopen{0}{f(x)}$.
        \caseOf{$0 < a$.}
            Thus $a < f(x)$ by \cref{intervalopen}.
            Take $p_0$ such that $p_0\in\unitRationals$ and $p_0\in\intervalopen{a}{f(x)}$
                by \cref{unit_rationals_dense_nonneg}.
            Thus $p_0\in\reals$ by \cref{intervalopen}.
            Thus $a < p_0$ and $p_0 < f(x)$ by \cref{intervalopen}.
            Thus $0 < p_0$ by \cref{real_lt_trans}.
            Thus $p_0\in\intervalopen{0}{f(x)}$ by \cref{intervalopen}.
            Thus there exists $p$ such that
                $p\in\unitRationals$ and
                $p\in\intervalopen{a}{f(x)}$ and
                $p\in\intervalopen{0}{f(x)}$.
        \end{byCase}
    \end{subproof}
    Take $p$ such that
        $p\in\unitRationals$ and
        $p\in\intervalopen{a}{f(x)}$ and
        $p\in\intervalopen{0}{f(x)}$.
    Thus $f(x)\in\reals$ by \cref{funs_type_apply}.
    Thus $f(x) < b$ by \cref{intervalopen}.
    Take $q$ such that $q\in\unitRationals$ and $q\in\intervalopen{f(x)}{b}$
        by \cref{unit_rationals_dense_nonneg}.
    Thus $p\in\reals$ and $q\in\reals$ by \cref{intervalopen}.
    Thus $a < p$ and $p < f(x)$ by \cref{intervalopen}.
    Thus $0 < p$ by \cref{intervalopen}.
    Thus $0 \leq p$ by \cref{real_lt_imp_leq}.
    Take $r_0$ such that $r_0\in\unitRationals$ and $r_0\in\intervalopen{p}{f(x)}$
        by \cref{unit_rationals_dense_nonneg}.
    Thus $r_0\in\reals$ by \cref{intervalopen}.
    Thus $p < r_0$ and $r_0 < f(x)$ by \cref{intervalopen}.
    Thus $r_0\in\unitRationalsZero$ by \cref{unit_rationals_zero,union_intro_left}.
    Thus $f(x) < q$ and $q < b$ by \cref{intervalopen}.
    Thus $p\in\unitRationalsZero$ and $q\in\unitRationalsZero$
        by \cref{unit_rationals_zero,union_intro_left}.
    Let $W = U(q)\setminus \closure{U(p)}{X}{T}$.
        Show that $W$ is open in $X$ with respect to $T$.
            \begin{subproof}
                Show that $U(q)\in T$.
                \begin{subproof}
                    Thus $q\in\unitRationals$.
                    Take $k$ such that $k\in\naturalsPlus$ and $s(k) = q$
                        by \cref{unit_rationals_sequence}.
                    Thus $u$ is a sequence in $G$.
                    Thus $u\in\funs{\naturalsPlus}{G}$ by \cref{sequence_in}.
                    Thus $u(k)\in G$ by \cref{funs_type_apply}.
                    Take $n,t$ such that
                        $n\in\naturalsPlus$ and
                        $t\in\urysohnAdmissible{s}{X}{T}{A}{B}{n}$ and
                        $u(k) = (n,t)$
                        by \cref{urysohn_good_pair_decomp}.
                    Thus $\fst{u(k)} = n$ by \cref{fst_eq}.
                    Thus $k = n$.
                    Thus $u(k) = (k,t)$ by \cref{pair_eq_iff}.
                    Thus $\snd{u(k)} = t$ by \cref{snd_eq}.
                    Thus $t\in F$.
                    Thus $t\in\funs{\{0,1\}\union \img{s}{\initialSegment{k}}}{\pow{X}}$
                        by \cref{urysohn_admissible}.
                    Thus $t$ is a function by \cref{funs_is_function}.
                    Thus $t$ is right-unique and $t$ is a relation.
                    Thus $\dom{t} = \{0,1\}\union \img{s}{\initialSegment{k}}$
                        by \cref{funs_elim}.
                    Show that $q\in\dom{t}$.
                    \begin{subproof}
                        Thus $k \leq k$.
                        Thus $k\in\initialSegment{k}$ by \cref{initial_segment_naturalsplus}.
                        Let $y = k$.
                        Thus $y\in\initialSegment{k}$.
                        Thus $s(y) = s(k)$.
                        Thus $s$ is a sequence in $\unitRationals$.
                        Thus $s\in\funs{\naturalsPlus}{\unitRationals}$ by \cref{sequence_in}.
                        Thus $s$ is a function by \cref{funs_is_function}.
                        Thus $\dom{s} = \naturalsPlus$ by \cref{funs_elim}.
                        Thus $y\in\dom{s}$ by \cref{subseteq}.
                        Thus $(y,s(y))\in s$ by \cref{function_apply_elim}.
                        Thus $(y,s(k))\in s$.
                        Thus $s(k)\in\img{s}{\initialSegment{k}}$ by \cref{img_iff}.
                        Thus $q\in\img{s}{\initialSegment{k}}$.
                        Thus $q\in\dom{t}$ by \cref{union_intro_right}.
                    \end{subproof}
                    Thus $t(q)\in T$ by \cref{urysohn_admissible,subseteq}.
                    Thus $(q,t(q))\in t$ by \cref{function_apply_elim}.
                    Thus $(q,t(q))\in U$ by \cref{unions_intro}.
                    Thus $U(q) = t(q)$ by \cref{function_apply_intro}.
                    Thus $U(q)\in T$.
                \end{subproof}
                Show that $U(p)\subseteq X$.
                \begin{subproof}
                    Thus $p\in\unitRationals$.
                    Take $k_p$ such that $k_p\in\naturalsPlus$ and $s(k_p) = p$
                        by \cref{unit_rationals_sequence}.
                    Thus $u$ is a sequence in $G$.
                    Thus $u\in\funs{\naturalsPlus}{G}$ by \cref{sequence_in}.
                    Thus $u(k_p)\in G$ by \cref{funs_type_apply}.
                    Take $n_p,t_p$ such that
                        $n_p\in\naturalsPlus$ and
                        $t_p\in\urysohnAdmissible{s}{X}{T}{A}{B}{n_p}$ and
                        $u(k_p) = (n_p,t_p)$
                        by \cref{urysohn_good_pair_decomp}.
                    Thus $\fst{u(k_p)} = n_p$ by \cref{fst_eq}.
                    Thus $k_p = n_p$.
                    Thus $u(k_p) = (k_p,t_p)$ by \cref{pair_eq_iff}.
                    Thus $\snd{u(k_p)} = t_p$ by \cref{snd_eq}.
                    Thus $t_p\in F$.
                    Thus $t_p\in\funs{\{0,1\}\union \img{s}{\initialSegment{k_p}}}{\pow{X}}$
                        by \cref{urysohn_admissible}.
                    Thus $t_p$ is a function by \cref{funs_is_function}.
                    Thus $t_p$ is right-unique and $t_p$ is a relation.
                    Thus $\dom{t_p} = \{0,1\}\union \img{s}{\initialSegment{k_p}}$
                        by \cref{funs_elim}.
                    Show that $p\in\dom{t_p}$.
                    \begin{subproof}
                        Thus $k_p \leq k_p$.
                        Thus $k_p\in\initialSegment{k_p}$
                            by \cref{initial_segment_naturalsplus}.
                        Let $y = k_p$.
                        Thus $y\in\initialSegment{k_p}$.
                        Thus $s(y) = s(k_p)$.
                        Thus $s$ is a sequence in $\unitRationals$.
                        Thus $s\in\funs{\naturalsPlus}{\unitRationals}$ by \cref{sequence_in}.
                        Thus $s$ is a function by \cref{funs_is_function}.
                        Thus $\dom{s} = \naturalsPlus$ by \cref{funs_elim}.
                        Thus $y\in\dom{s}$ by \cref{subseteq}.
                        Thus $(y,s(y))\in s$ by \cref{function_apply_elim}.
                        Thus $(y,s(k_p))\in s$.
                        Thus $s(k_p)\in\img{s}{\initialSegment{k_p}}$ by \cref{img_iff}.
                        Thus $p\in\img{s}{\initialSegment{k_p}}$.
                        Thus $p\in\dom{t_p}$ by \cref{union_intro_right}.
                    \end{subproof}
                    Thus $t_p(p)\in\pow{X}$ by \cref{funs_type_apply}.
                    Thus $t_p(p)\subseteq X$ by \cref{powerset_elim}.
                    Thus $(p,t_p(p))\in t_p$ by \cref{function_apply_elim}.
                    Thus $(p,t_p(p))\in U$ by \cref{unions_intro}.
                    Thus $U(p) = t_p(p)$ by \cref{function_apply_intro}.
                    Thus $U(p)\subseteq X$.
                \end{subproof}
                Thus $T$ is a topology on $X$ by \cref{normal_space}.
                Thus $\closure{U(p)}{X}{T}$ is closed in $X$ with respect to $T$
                    by \cref{closure_closed_in}.
                Thus $X\setminus \closure{U(p)}{X}{T}$ is open in $X$ with respect to $T$
                    by \cref{closed_in}.
                Thus $X\setminus \closure{U(p)}{X}{T}\in T$ by \cref{open_in}.
                Thus $U(q)\inter (X\setminus \closure{U(p)}{X}{T})\in T$
                    by \cref{topology_on}.
                Thus $T\subseteq\pow{X}$ by \cref{topology_on}.
                Thus $U(q)\in\pow{X}$ by \cref{subseteq}.
                Thus $U(q)\subseteq X$ by \cref{powerset_elim}.
                Show that $W = U(q)\inter (X\setminus \closure{U(p)}{X}{T})$.
                \begin{subproof}
                    Show that for all $z\in W$ we have
                        $z\in U(q)\inter (X\setminus \closure{U(p)}{X}{T})$.
                    \begin{subproof}
                        Fix $z$.
                        Assume $z\in W$.
                        Thus $z\in U(q)$ by \cref{setminus_elim_left}.
                        Thus $z\notin\closure{U(p)}{X}{T}$ by \cref{setminus_elim_right}.
                        Thus $z\in X$ by \cref{subseteq}.
                        Thus $z\in X\setminus \closure{U(p)}{X}{T}$ by \cref{setminus_intro}.
                        Thus $z\in U(q)\inter (X\setminus \closure{U(p)}{X}{T})$
                            by \cref{inter_intro}.
                    \end{subproof}
                    Show that for all $z\in U(q)\inter (X\setminus \closure{U(p)}{X}{T})$
                        we have $z\in W$.
                    \begin{subproof}
                        Fix $z$.
                        Assume $z\in U(q)\inter (X\setminus \closure{U(p)}{X}{T})$.
                        Thus $z\in U(q)$ by \cref{inter_elim_left}.
                        Thus $z\in X\setminus \closure{U(p)}{X}{T}$ by \cref{inter_elim_right}.
                        Thus $z\notin\closure{U(p)}{X}{T}$ by \cref{setminus_elim_right}.
                        Thus $z\in W$ by \cref{setminus_intro}.
                    \end{subproof}
                    Thus $W\subseteq U(q)\inter (X\setminus \closure{U(p)}{X}{T})$
                        by \cref{subseteq}.
                    Thus $U(q)\inter (X\setminus \closure{U(p)}{X}{T})\subseteq W$
                        by \cref{subseteq}.
                    Thus $W = U(q)\inter (X\setminus \closure{U(p)}{X}{T})$
                        by \cref{subseteq_antisymmetric}.
                \end{subproof}
                Thus $W\in T$.
                Thus $W$ is open in $X$ with respect to $T$ by \cref{open_in}.
            \end{subproof}
            Thus $W$ is open in $X$ with respect to $T$.
    Thus there exist $p,q,r_0,W$ such that
        $p\in\unitRationals$ and
        $p\in\intervalopen{a}{f(x)}$ and
        $p\in\intervalopen{0}{f(x)}$ and
        $q\in\unitRationals$ and
        $q\in\intervalopen{f(x)}{b}$ and
        $r_0\in\unitRationals$ and
        $r_0\in\intervalopen{p}{f(x)}$ and
        $p\in\unitRationalsZero$ and
        $q\in\unitRationalsZero$ and
        $r_0\in\unitRationalsZero$ and
        $W = U(q)\setminus \closure{U(p)}{X}{T}$ and
        $W$ is open in $X$ with respect to $T$.
\end{proof}

% from §33 helper lemma 21: positive case point membership
\begin{lemma}\label{urysohn_unit_interval_continuous_pos_x_in_W}
    Assume $X$ is normal with respect to $T$.
    Assume $A$ is closed in $X$ with respect to $T$.
    Assume $B$ is closed in $X$ with respect to $T$.
    Assume $A$ is disjoint from $B$.
    Assume $s$ is a sequence in $\unitRationals$.
    Assume for all $p$ such that $p\in\unitRationals$ there exists $n\in\naturalsPlus$ such that $s(n) = p$.
    Let $G = \urysohnGoodPairs{s}{X}{T}{A}{B}$.
    Let $R = \{(z,w) \mid z\in G, w\in G \mid
        \fst{w} = \fst{z} + 1 \land \snd{z}\subseteq\snd{w}\}$.
    Assume $c$ is a function on $G$.
    Assume for all $z\in G$ we have $z\mathrel{R} \apply{c}{z}$.
    Assume $u$ is a sequence in $G$.
    Assume $u(1) = a_0$.
    Assume $a_0 = (1,u_1)$.
    Assume for all $n\in\naturalsPlus$ we have $u(n + 1) = \apply{c}{u(n)}$.
    Assume for all $n\in\naturalsPlus$ we have $\fst{u(n)} = n$.
    Let $F = \{\snd{u(n)} \mid n\in\naturalsPlus\}$.
    Let $U = \unions{F}$.
    Assume $U$ is a function.
    Let $I_0 = \intervalclosed{0}{1}$.
    Let $f = \{(x,y) \mid x\in X, y\in\reals \mid
        \text{$y$ is an infimum of $\urysohnLevelSet{U}{x}$}\}$.
    Assume $f\in\funs{X}{\reals}$ and $\img{f}{X}\subseteq I_0$.
    Assume $x\in X$ and $V$ is open in $\reals$ with respect to $\standardTopologyR$ and
        $a\in\reals$ and $b\in\reals$ and $a < b$ and
        $f(x)\in\intervalopen{a}{b}$ and $\intervalopen{a}{b}\subseteq V$ and $0 < f(x)$.
    Assume $p\in\unitRationals$ and $p\in\intervalopen{a}{f(x)}$ and $p\in\intervalopen{0}{f(x)}$ and
        $q\in\unitRationals$ and $q\in\intervalopen{f(x)}{b}$ and
        $r_0\in\unitRationals$ and $r_0\in\intervalopen{p}{f(x)}$ and
        $p\in\unitRationalsZero$ and $q\in\unitRationalsZero$ and $r_0\in\unitRationalsZero$.
    Let $W = U(q)\setminus \closure{U(p)}{X}{T}$.
    Then $x\in W$.
\end{lemma}
\begin{proof}
    Let $Q_x = \urysohnLevelSet{U}{x}$.
    Show that $f(x)$ is an infimum of $Q_x$.
    \begin{subproof}
        Thus $f$ is a function by \cref{funs_is_function}.
        Thus $\dom{f} = X$ by \cref{funs_elim}.
        Thus $x\in\dom{f}$ by \cref{subseteq}.
        Thus $(x,f(x))\in f$ by \cref{function_apply_elim}.
        Take $x_0,y_0$ such that
            $x_0\in X$ and
            $y_0\in\reals$ and
            $(x,f(x)) = (x_0,y_0)$ and
            $y_0$ is an infimum of $\urysohnLevelSet{U}{x_0}$.
        Thus $x = x_0$ and $f(x) = y_0$ by \cref{pair_eq_iff}.
        Thus $f(x)$ is an infimum of $Q_x$.
    \end{subproof}
    Thus $f(x)$ is a lower bound of $Q_x$ by \cref{real_infimum}.
    Show that for all $r\in Q_x$ we have $f(x) \leq r$.
    \begin{subproof}
        Fix $r$.
        Assume $r\in Q_x$.
        Thus $f(x) \leq r$ by \cref{real_lower_bound}.
    \end{subproof}
    Show that $x\in U(q)$.
        \begin{subproof}
            Thus $f(x) < q$.
            Show that $f(x)$ is an infimum of $Q_x$.
            \begin{subproof}
                Thus $f$ is a function by \cref{funs_is_function}.
                Thus $\dom{f} = X$ by \cref{funs_elim}.
                Thus $x\in\dom{f}$ by \cref{subseteq}.
                Thus $(x,f(x))\in f$ by \cref{function_apply_elim}.
                Take $x_0,y_0$ such that
                    $x_0\in X$ and
                    $y_0\in\reals$ and
                    $(x,f(x)) = (x_0,y_0)$ and
                    $y_0$ is an infimum of $\urysohnLevelSet{U}{x_0}$.
                Thus $x = x_0$ and $f(x) = y_0$ by \cref{pair_eq_iff}.
                Thus $f(x)$ is an infimum of $Q_x$.
            \end{subproof}
            Thus $f(x)$ is a lower bound of $Q_x$ by \cref{real_infimum}.
            Thus $f(x) \leq q$ by \cref{real_lower_bound}.
            Thus $f(x) < q$ by \cref{real_lt_subst_left}.
            Thus it is wrong that $q$ is a lower bound of $Q_x$.
            Show that there exists $r$ such that $r\in Q_x$ and $r < q$.
            \begin{subproof}
                Suppose not.
                Show that $q$ is a lower bound of $Q_x$.
                \begin{subproof}
                    Thus $q\in\reals$.
                    Show that $Q_x\subseteq\reals$.
                    \begin{subproof}
                        Show that for all $r\in Q_x$ we have $r\in\reals$.
                        \begin{subproof}
                            Fix $r$.
                            Assume $r\in Q_x$.
                            Thus $r\in\unitRationalsZero$ by \cref{urysohn_level_set}.
                            Thus $r\in\unitRationals$ or $r\in\{0\}$
                                by \cref{unit_rationals_zero,union_iff}.
                            \begin{byCase}
                            \caseOf{$r\in\unitRationals$.}
                                Take $t$ such that $t\in\naturalsPlus\times\naturalsPlus$ and
                                    $r = \fst{t} \rmul \inv{\snd{t}}$
                                    by \cref{unit_rationals}.
                                Thus $\fst{t}\in\naturalsPlus$ and $\snd{t}\in\naturalsPlus$
                                    by \cref{fst_type,snd_type}.
                                Thus $\fst{t}\in\reals$ and $\snd{t}\in\reals$
                                    by \cref{real_naturals_subset_reals,subseteq}.
                                Show that $\snd{t}\neq 0$.
                                \begin{subproof}
                                    $1\in\reals$ by \cref{real_mul_ident}.
                                    Thus $1 < \snd{t}$ or $1 = \snd{t}$
                                        by \cref{real_one_leq_naturalsplus}.
                                    $0\in\reals$ by \cref{real_add_ident}.
                                    $0 < 1$ by \cref{real_zero_lt_one}.
                                    \begin{byCase}
                                    \caseOf{$1 < \snd{t}$.}
                                        Thus $0 < \snd{t}$ by \cref{real_lt_trans}.
                                        Thus $\snd{t}\neq 0$ by \cref{real_pos_nonzero}.
                                    \caseOf{$1 = \snd{t}$.}
                                        Thus $0 < \snd{t}$ by \cref{real_lt_subst_right}.
                                        Thus $\snd{t}\neq 0$ by \cref{real_pos_nonzero}.
                                    \end{byCase}
                                \end{subproof}
                                Thus $\inv{\snd{t}}\in\reals$ by \cref{real_mul_inverse}.
                                Thus $r\in\reals$ by \cref{real_mul_closed}.
                            \caseOf{$r\in\{0\}$.}
                                Thus $r = 0$ by \cref{singleton_elim}.
                                $0\in\reals$ by \cref{real_add_ident}.
                                Thus $r\in\reals$.
                            \end{byCase}
                        \end{subproof}
                        Thus $Q_x\subseteq\reals$ by \cref{subseteq}.
                    \end{subproof}
                    Show that for all $r\in Q_x$ we have $q < r$ or $q = r$.
                    \begin{subproof}
                        Fix $r$.
                        Assume $r\in Q_x$.
                        Thus $r\in\reals$ by \cref{subseteq}.
                        Show that $q < r$ or $q = r$.
                        \begin{subproof}
                            Suppose not.
                            Thus $r < q$ or $r = q$ or $q < r$
                                by \cref{real_lt_trichotomy}.
                            Thus $r < q$.
                            Thus there exists $r$ such that $r\in Q_x$ and $r < q$.
                            Contradiction.
                        \end{subproof}
                    \end{subproof}
                    Thus $q$ is a lower bound of $Q_x$ by \cref{real_lower_bound}.
                \end{subproof}
                Contradiction.
            \end{subproof}
            Take $r$ such that $r\in Q_x$ and $r < q$.
            Thus $x\in U(r)$.
            Show that $x\in U(q)$.
            \begin{subproof}
                Thus $r\in\unitRationalsZero$ by \cref{urysohn_level_set}.
                Thus $r\in\unitRationals$ or $r\in\{0\}$
                    by \cref{unit_rationals_zero,union_iff}.
                Thus $q\in\unitRationals$.
                Take $k_q$ such that $k_q\in\naturalsPlus$ and $s(k_q) = q$
                    by \cref{unit_rationals_sequence}.
                Thus $u$ is a sequence in $G$.
                Thus $u\in\funs{\naturalsPlus}{G}$ by \cref{sequence_in}.
                Thus $u(k_q)\in G$ by \cref{funs_type_apply}.
                Take $n_q,t_q$ such that
                    $n_q\in\naturalsPlus$ and
                    $t_q\in\urysohnAdmissible{s}{X}{T}{A}{B}{n_q}$ and
                    $u(k_q) = (n_q,t_q)$
                    by \cref{urysohn_good_pair_decomp}.
                Thus $\fst{u(k_q)} = n_q$ by \cref{fst_eq}.
                Thus $\fst{u(k_q)} = k_q$.
                Thus $k_q = n_q$.
                Thus $u(k_q) = (k_q,t_q)$ by \cref{pair_eq_iff}.
                Thus $\snd{u(k_q)} = t_q$ by \cref{snd_eq}.
                Thus $t_q\in F$.
                Thus $t_q\in\funs{\{0,1\}\union \img{s}{\initialSegment{k_q}}}{\pow{X}}$
                    by \cref{urysohn_admissible}.
                Thus $t_q$ is a function by \cref{funs_is_function}.
                Thus $t_q$ is right-unique and $t_q$ is a relation.
                Thus $\dom{t_q} = \{0,1\}\union \img{s}{\initialSegment{k_q}}$
                    by \cref{funs_elim}.
                Show that $q\in\dom{t_q}$.
                \begin{subproof}
                    Thus $k_q \leq k_q$.
                    Thus $k_q\in\initialSegment{k_q}$ by \cref{initial_segment_naturalsplus}.
                    Let $y = k_q$.
                    Thus $y\in\initialSegment{k_q}$.
                    Thus $s(y) = s(k_q)$.
                    Thus $s$ is a sequence in $\unitRationals$.
                    Thus $s\in\funs{\naturalsPlus}{\unitRationals}$ by \cref{sequence_in}.
                    Thus $s$ is a function by \cref{funs_is_function}.
                    Thus $\dom{s} = \naturalsPlus$ by \cref{funs_elim}.
                    Thus $y\in\dom{s}$ by \cref{subseteq}.
                    Thus $(y,s(y))\in s$ by \cref{function_apply_elim}.
                    Thus $(y,s(k_q))\in s$.
                    Thus $s(k_q)\in\img{s}{\initialSegment{k_q}}$ by \cref{img_iff}.
                    Thus $q\in\img{s}{\initialSegment{k_q}}$.
                    Thus $q\in\dom{t_q}$ by \cref{union_intro_right}.
                \end{subproof}
                Thus $(q,t_q(q))\in t_q$ by \cref{function_apply_elim}.
                Thus $(q,t_q(q))\in U$ by \cref{unions_intro}.
                Thus $U(q) = t_q(q)$ by \cref{function_apply_intro}.
                \begin{byCase}
                \caseOf{$r\in\unitRationals$.}
                    Show that $x\in U(q)$.
                    \begin{subproof}
                        Take $k_r$ such that $k_r\in\naturalsPlus$ and $s(k_r) = r$
                            by \cref{unit_rationals_sequence}.
                        Thus $u$ is a sequence in $G$.
                        Thus $u\in\funs{\naturalsPlus}{G}$ by \cref{sequence_in}.
                        Thus $u(k_r)\in G$ by \cref{funs_type_apply}.
                        Take $n_r,t_r$ such that
                            $n_r\in\naturalsPlus$ and
                            $t_r\in\urysohnAdmissible{s}{X}{T}{A}{B}{n_r}$ and
                            $u(k_r) = (n_r,t_r)$
                            by \cref{urysohn_good_pair_decomp}.
                        Thus $\fst{u(k_r)} = n_r$ by \cref{fst_eq}.
                        Thus $\fst{u(k_r)} = k_r$.
                        Thus $k_r = n_r$.
                        Thus $u(k_r) = (k_r,t_r)$ by \cref{pair_eq_iff}.
                        Thus $\snd{u(k_r)} = t_r$ by \cref{snd_eq}.
                        Thus $t_r\in F$.
                        Thus $t_r\in\funs{\{0,1\}\union \img{s}{\initialSegment{k_r}}}{\pow{X}}$
                            by \cref{urysohn_admissible}.
                        Thus $t_r$ is a function by \cref{funs_is_function}.
                        Thus $t_r$ is right-unique and $t_r$ is a relation.
                        Thus $\dom{t_r} = \{0,1\}\union \img{s}{\initialSegment{k_r}}$
                            by \cref{funs_elim}.
                        Show that $r\in\dom{t_r}$.
                        \begin{subproof}
                            Thus $k_r \leq k_r$.
                            Thus $k_r\in\initialSegment{k_r}$
                                by \cref{initial_segment_naturalsplus}.
                            Let $y = k_r$.
                            Thus $y\in\initialSegment{k_r}$.
                            Thus $s(y) = s(k_r)$.
                            Thus $s$ is a sequence in $\unitRationals$.
                            Thus $s\in\funs{\naturalsPlus}{\unitRationals}$ by \cref{sequence_in}.
                            Thus $s$ is a function by \cref{funs_is_function}.
                            Thus $\dom{s} = \naturalsPlus$ by \cref{funs_elim}.
                            Thus $y\in\dom{s}$ by \cref{subseteq}.
                            Thus $(y,s(y))\in s$ by \cref{function_apply_elim}.
                            Thus $(y,s(k_r))\in s$.
                            Thus $s(k_r)\in\img{s}{\initialSegment{k_r}}$ by \cref{img_iff}.
                            Thus $r\in\img{s}{\initialSegment{k_r}}$.
                            Thus $r\in\dom{t_r}$ by \cref{union_intro_right}.
                    \end{subproof}
                    Thus $(r,t_r(r))\in t_r$ by \cref{function_apply_elim}.
                    Thus $t_r\subseteq t_q$ or $t_q\subseteq t_r$
                        by \cref{urysohn_unit_interval_chain}.
                    \begin{byCase}
                    \caseOf{$t_r\subseteq t_q$.}
                            Thus $(r,t_r(r))\in t_q$ by \cref{subseteq}.
                            Thus $r\in\dom{t_q}$ by \cref{dom_intro}.
                            Thus $\closure{t_q(r)}{X}{T}\subseteq t_q(q)$
                                by \cref{urysohn_admissible}.
                            Thus $(r,t_q(r))\in t_q$ by \cref{function_apply_elim}.
                            Thus $(r,t_q(r))\in U$ by \cref{unions_intro}.
                            Thus $U(r) = t_q(r)$ by \cref{function_apply_intro}.
                            Thus $t_q(r)\in\pow{X}$ by \cref{funs_type_apply}.
                            Thus $t_q(r)\subseteq X$ by \cref{powerset_elim}.
                            Thus $U(r)\subseteq X$.
                            Thus $\closure{U(r)}{X}{T}\subseteq U(q)$.
                            Thus $U(r)\subseteq\closure{U(r)}{X}{T}$
                                by \cref{subseteq_closure,normal_space}.
                            Thus $x\in\closure{U(r)}{X}{T}$ by \cref{subseteq}.
                            Thus $x\in U(q)$ by \cref{subseteq}.
                    \caseOf{$t_q\subseteq t_r$.}
                            Thus $(q,t_q(q))\in t_r$ by \cref{subseteq}.
                            Thus $q\in\dom{t_r}$ by \cref{dom_intro}.
                            Thus $\closure{t_r(r)}{X}{T}\subseteq t_r(q)$
                                by \cref{urysohn_admissible}.
                            Thus $(r,t_r(r))\in t_r$ by \cref{function_apply_elim}.
                            Thus $(r,t_r(r))\in U$ by \cref{unions_intro}.
                            Thus $U(r) = t_r(r)$ by \cref{function_apply_intro}.
                            Thus $t_r(r)\in\pow{X}$ by \cref{funs_type_apply}.
                            Thus $t_r(r)\subseteq X$ by \cref{powerset_elim}.
                            Thus $U(r)\subseteq X$.
                            Thus $(q,t_r(q))\in t_r$ by \cref{function_apply_elim}.
                            Thus $(q,t_r(q))\in U$ by \cref{unions_intro}.
                            Thus $U(q) = t_r(q)$ by \cref{function_apply_intro}.
                            Thus $\closure{U(r)}{X}{T}\subseteq U(q)$.
                            Thus $U(r)\subseteq\closure{U(r)}{X}{T}$
                                by \cref{subseteq_closure,normal_space}.
                            Thus $x\in\closure{U(r)}{X}{T}$ by \cref{subseteq}.
                            Thus $x\in U(q)$ by \cref{subseteq}.
                    \end{byCase}
                \end{subproof}
                \caseOf{$r\in\{0\}$.}
                    Thus $r = 0$ by \cref{singleton_elim}.
                    $0\in\{0,1\}$ by \cref{upair_intro_left}.
                    Thus $0\in\dom{t_q}$ by \cref{union_intro_left,subseteq}.
                    Thus $\closure{t_q(0)}{X}{T}\subseteq t_q(q)$
                        by \cref{urysohn_admissible}.
                    Thus $(0,t_q(0))\in t_q$ by \cref{function_apply_elim}.
                    Thus $(0,t_q(0))\in U$ by \cref{unions_intro}.
                    Thus $U(r) = t_q(0)$ by \cref{function_apply_intro}.
                    Thus $t_q(0)\in\pow{X}$ by \cref{funs_type_apply}.
                    Thus $t_q(0)\subseteq X$ by \cref{powerset_elim}.
                    Thus $U(r)\subseteq X$.
                    Thus $\closure{U(r)}{X}{T}\subseteq U(q)$.
                    Thus $U(r)\subseteq\closure{U(r)}{X}{T}$
                        by \cref{subseteq_closure,normal_space}.
                    Thus $x\in\closure{U(r)}{X}{T}$ by \cref{subseteq}.
                    Thus $x\in U(q)$ by \cref{subseteq}.
                \end{byCase}
        \end{subproof}
        \end{subproof}
        Show that $x\notin\closure{U(p)}{X}{T}$.
        \begin{subproof}
            Show that $\closure{U(p)}{X}{T}\subseteq U(r_0)$.
            \begin{subproof}
                Thus $p\in\unitRationals$.
                Thus $r_0\in\unitRationals$.
                Take $k_p$ such that $k_p\in\naturalsPlus$ and $s(k_p) = p$
                    by \cref{unit_rationals_sequence}.
                Take $k_r$ such that $k_r\in\naturalsPlus$ and $s(k_r) = r_0$
                    by \cref{unit_rationals_sequence}.
                Thus $u$ is a sequence in $G$.
                Thus $u\in\funs{\naturalsPlus}{G}$ by \cref{sequence_in}.
                Thus $u(k_p)\in G$ by \cref{funs_type_apply}.
                Thus $u(k_r)\in G$ by \cref{funs_type_apply}.
                Take $n_p,t_p$ such that
                    $n_p\in\naturalsPlus$ and
                    $t_p\in\urysohnAdmissible{s}{X}{T}{A}{B}{n_p}$ and
                    $u(k_p) = (n_p,t_p)$
                    by \cref{urysohn_good_pair_decomp}.
                Take $n_r,t_r$ such that
                    $n_r\in\naturalsPlus$ and
                    $t_r\in\urysohnAdmissible{s}{X}{T}{A}{B}{n_r}$ and
                    $u(k_r) = (n_r,t_r)$
                    by \cref{urysohn_good_pair_decomp}.
                Thus $\fst{u(k_p)} = n_p$ by \cref{fst_eq}.
                Thus $\fst{u(k_r)} = n_r$ by \cref{fst_eq}.
                Thus $k_p = n_p$.
                Thus $k_r = n_r$.
                Thus $u(k_p) = (k_p,t_p)$ by \cref{pair_eq_iff}.
                Thus $u(k_r) = (k_r,t_r)$ by \cref{pair_eq_iff}.
                Thus $\snd{u(k_p)} = t_p$ by \cref{snd_eq}.
                Thus $\snd{u(k_r)} = t_r$ by \cref{snd_eq}.
                Thus $t_p\in F$.
                Thus $t_r\in F$.
                Thus $t_p\in\funs{\{0,1\}\union \img{s}{\initialSegment{k_p}}}{\pow{X}}$
                    by \cref{urysohn_admissible}.
                Thus $t_r\in\funs{\{0,1\}\union \img{s}{\initialSegment{k_r}}}{\pow{X}}$
                    by \cref{urysohn_admissible}.
                Thus $t_p$ is a function by \cref{funs_is_function}.
                Thus $t_r$ is a function by \cref{funs_is_function}.
                Thus $t_p$ is right-unique and $t_p$ is a relation.
                Thus $t_r$ is right-unique and $t_r$ is a relation.
                Thus $\dom{t_p} = \{0,1\}\union \img{s}{\initialSegment{k_p}}$
                    by \cref{funs_elim}.
                Thus $\dom{t_r} = \{0,1\}\union \img{s}{\initialSegment{k_r}}$
                    by \cref{funs_elim}.
                Show that $p\in\dom{t_p}$.
                \begin{subproof}
                    Thus $k_p \leq k_p$.
                    Thus $k_p\in\initialSegment{k_p}$
                        by \cref{initial_segment_naturalsplus}.
                    Let $y = k_p$.
                    Thus $y\in\initialSegment{k_p}$.
                    Thus $s(y) = s(k_p)$.
                    Thus $s$ is a sequence in $\unitRationals$.
                    Thus $s\in\funs{\naturalsPlus}{\unitRationals}$ by \cref{sequence_in}.
                    Thus $s$ is a function by \cref{funs_is_function}.
                    Thus $\dom{s} = \naturalsPlus$ by \cref{funs_elim}.
                    Thus $y\in\dom{s}$ by \cref{subseteq}.
                    Thus $(y,s(y))\in s$ by \cref{function_apply_elim}.
                    Thus $(y,s(k_p))\in s$.
                    Thus $s(k_p)\in\img{s}{\initialSegment{k_p}}$ by \cref{img_iff}.
                    Thus $p\in\img{s}{\initialSegment{k_p}}$.
                    Thus $p\in\dom{t_p}$ by \cref{union_intro_right}.
                \end{subproof}
                Show that $r_0\in\dom{t_r}$.
                \begin{subproof}
                    Thus $k_r \leq k_r$.
                    Thus $k_r\in\initialSegment{k_r}$
                        by \cref{initial_segment_naturalsplus}.
                    Let $y = k_r$.
                    Thus $y\in\initialSegment{k_r}$.
                    Thus $s(y) = s(k_r)$.
                    Thus $s$ is a sequence in $\unitRationals$.
                    Thus $s\in\funs{\naturalsPlus}{\unitRationals}$ by \cref{sequence_in}.
                    Thus $s$ is a function by \cref{funs_is_function}.
                    Thus $\dom{s} = \naturalsPlus$ by \cref{funs_elim}.
                    Thus $y\in\dom{s}$ by \cref{subseteq}.
                    Thus $(y,s(y))\in s$ by \cref{function_apply_elim}.
                    Thus $(y,s(k_r))\in s$.
                    Thus $s(k_r)\in\img{s}{\initialSegment{k_r}}$ by \cref{img_iff}.
                    Thus $r_0\in\img{s}{\initialSegment{k_r}}$.
                    Thus $r_0\in\dom{t_r}$ by \cref{union_intro_right}.
                \end{subproof}
                Thus $(p,t_p(p))\in t_p$ by \cref{function_apply_elim}.
                Thus $(r_0,t_r(r_0))\in t_r$ by \cref{function_apply_elim}.
                Thus $(p,t_p(p))\in U$ by \cref{unions_intro}.
                Thus $(r_0,t_r(r_0))\in U$ by \cref{unions_intro}.
                Thus $U(p) = t_p(p)$ by \cref{function_apply_intro}.
                Thus $U(r_0) = t_r(r_0)$ by \cref{function_apply_intro}.
                Thus $t_p\subseteq t_r$ or $t_r\subseteq t_p$
                    by \cref{urysohn_unit_interval_chain}.
                \begin{byCase}
                \caseOf{$t_p\subseteq t_r$.}
                        Thus $(p,t_p(p))\in t_r$ by \cref{subseteq}.
                        Thus $t_r(p) = t_p(p)$ by \cref{function_apply_intro}.
                        Thus $U(p) = t_r(p)$.
                        Thus $p\in\dom{t_r}$ by \cref{dom_intro}.
                        Thus $r_0\in\dom{t_r}$ by \cref{dom_intro}.
                        Thus $p < r_0$.
                        Thus $\closure{t_r(p)}{X}{T}\subseteq t_r(r_0)$
                            by \cref{urysohn_admissible}.
                        Thus $\closure{U(p)}{X}{T}\subseteq U(r_0)$.
                \caseOf{$t_r\subseteq t_p$.}
                        Thus $(r_0,t_r(r_0))\in t_p$ by \cref{subseteq}.
                        Thus $t_p(r_0) = t_r(r_0)$ by \cref{function_apply_intro}.
                        Thus $U(r_0) = t_p(r_0)$.
                        Thus $p\in\dom{t_p}$ by \cref{dom_intro}.
                        Thus $r_0\in\dom{t_p}$ by \cref{dom_intro}.
                        Thus $p < r_0$.
                        Thus $\closure{t_p(p)}{X}{T}\subseteq t_p(r_0)$
                            by \cref{urysohn_admissible}.
                        Thus $\closure{U(p)}{X}{T}\subseteq U(r_0)$.
                \end{byCase}
            \end{subproof}
            Show that $x\notin U(r_0)$.
            \begin{subproof}
                Suppose not.
                Thus $x\in U(r_0)$.
                Thus $r_0\in Q_x$ by \cref{urysohn_level_set}.
                Thus $f(x) \leq r_0$ by \cref{real_lower_bound}.
                Thus $f(x)\in\reals$ by \cref{funs_type_apply}.
                Thus $r_0\in\reals$ by \cref{intervalopen}.
                Thus $f(x) < r_0$ or $f(x) = r_0$ by \cref{real_leq_split}.
                Thus $r_0 < f(x)$ by \cref{intervalopen}.
                \begin{byCase}
                \caseOf{$f(x) < r_0$.}
                        Thus $f(x) < f(x)$ by \cref{real_lt_trans}.
                        Contradiction by \cref{real_lt_irreflexive}.
                \caseOf{$f(x) = r_0$.}
                        Thus $f(x) < f(x)$ by \cref{real_lt_subst_left}.
                        Contradiction by \cref{real_lt_irreflexive}.
                \end{byCase}
            \end{subproof}
                Thus $x\notin\closure{U(p)}{X}{T}$ by \cref{not_in_subseteq}.
        \end{subproof}
        Thus $x\notin\closure{U(p)}{X}{T}$.
        Thus $x\in W$ by \cref{setminus_intro}.
\end{proof}

% from §33 helper lemma 22: positive case image control
\begin{lemma}\label{urysohn_unit_interval_continuous_pos_img}
    Assume $X$ is normal with respect to $T$ and $A$ is closed in $X$ with respect to $T$ and
        $B$ is closed in $X$ with respect to $T$ and $A$ is disjoint from $B$.
    Assume $s$ is a sequence in $\unitRationals$ and for all $p$ such that $p\in\unitRationals$
        there exists $n\in\naturalsPlus$ such that $s(n) = p$.
    Let $G = \urysohnGoodPairs{s}{X}{T}{A}{B}$.
    Let $R = \{(z,w) \mid z\in G, w\in G \mid
        \fst{w} = \fst{z} + 1 \land \snd{z}\subseteq\snd{w}\}$.
    Assume $c$ is a function on $G$ and for all $z\in G$ we have $z\mathrel{R} \apply{c}{z}$.
    Assume $u$ is a sequence in $G$ and $u(1) = a_0$ and $a_0 = (1,u_1)$ and
        for all $n\in\naturalsPlus$ we have $u(n + 1) = \apply{c}{u(n)}$ and $\fst{u(n)} = n$.
    Let $F = \{\snd{u(n)} \mid n\in\naturalsPlus\}$.
    Let $U = \unions{F}$.
    Assume $U$ is a function.
    Let $I_0 = \intervalclosed{0}{1}$.
    Let $f = \{(x,y) \mid x\in X, y\in\reals \mid
        \text{$y$ is an infimum of $\urysohnLevelSet{U}{x}$}\}$.
    Assume $f\in\funs{X}{\reals}$ and $\img{f}{X}\subseteq I_0$.
    Assume $x\in X$ and $V$ is open in $\reals$ with respect to $\standardTopologyR$ and
        $a\in\reals$ and $b\in\reals$ and $a < b$ and
        $f(x)\in\intervalopen{a}{b}$ and $\intervalopen{a}{b}\subseteq V$ and $0 < f(x)$.
    Assume $p\in\unitRationals$ and $p\in\intervalopen{a}{f(x)}$ and $p\in\intervalopen{0}{f(x)}$ and
        $q\in\unitRationals$ and $q\in\intervalopen{f(x)}{b}$ and
        $r_0\in\unitRationals$ and $r_0\in\intervalopen{p}{f(x)}$ and
        $p\in\unitRationalsZero$ and $q\in\unitRationalsZero$ and $r_0\in\unitRationalsZero$.
    Let $W = U(q)\setminus \closure{U(p)}{X}{T}$.
    Assume $W\subseteq X$.
    Then $\img{f}{W}\subseteq V$.
\end{lemma}
\begin{proof}
    Show that for all $y\in W$ we have $f(y)\in V$.
        \begin{subproof}
            Fix $y$.
            Assume $y\in W$.
            Thus $y\in X$ and $f(y)\in\reals$ by \cref{subseteq,funs_type_apply}.
            Let $Q_y = \urysohnLevelSet{U}{y}$.
            Thus $y\in U(q)$ by \cref{setminus_elim_left}.
            Thus $y\notin\closure{U(p)}{X}{T}$ by \cref{setminus_elim_right}.
            Thus $q\in Q_y$ by \cref{urysohn_level_set}.
            Show that $f(y) \leq q$.
            \begin{subproof}
                Show that $f(y)$ is a lower bound of $Q_y$.
                \begin{subproof}
                    Show that $f(y)$ is an infimum of $Q_y$.
                    \begin{subproof}
                        Thus $(y,f(y))\in f$ by \cref{funs_is_function,funs_elim,subseteq,function_apply_elim}.
                        Take $x_0,y_0$ such that
                            $x_0\in X$ and
                            $y_0\in\reals$ and
                            $(y,f(y)) = (x_0,y_0)$ and
                            $y_0$ is an infimum of $\urysohnLevelSet{U}{x_0}$.
                        Thus $y = x_0$ and $f(y) = y_0$ by \cref{pair_eq_iff}.
                        Thus $f(y)$ is an infimum of $Q_y$.
                    \end{subproof}
                    Thus $f(y)$ is a lower bound of $Q_y$ by \cref{real_infimum}.
                \end{subproof}
                Thus $f(y) \leq q$ by \cref{real_lower_bound}.
            \end{subproof}
            Show that $p \leq f(y)$.
            \begin{subproof}
                Show that $p$ is a lower bound of $Q_y$.
                \begin{subproof}
                    Show that $Q_y\subseteq\reals$.
                    \begin{subproof}
                        Show that for all $r\in Q_y$ we have $r\in\reals$.
                        \begin{subproof}
                            Fix $r$.
                            Assume $r\in Q_y$.
                            Thus $r\in\unitRationalsZero$ by \cref{urysohn_level_set}.
                            Thus $r\in\unitRationals$ or $r\in\{0\}$
                                by \cref{unit_rationals_zero,union_iff}.
                            \begin{byCase}
                            \caseOf{$r\in\unitRationals$.}
                                Take $q$ such that $q\in\naturalsPlus\times\naturalsPlus$ and
                                    $r = \fst{q} \rmul \inv{\snd{q}}$
                                    by \cref{unit_rationals}.
                                Thus $\fst{q}\in\naturalsPlus$ and $\snd{q}\in\naturalsPlus$
                                    by \cref{fst_type,snd_type}.
                                Thus $\fst{q}\in\reals$ and $\snd{q}\in\reals$
                                    by \cref{real_naturals_subset_reals,subseteq}.
                                Show that $\snd{q}\neq 0$.
                                \begin{subproof}
                                    $1\in\reals$ by \cref{real_mul_ident}.
                                    Thus $1 < \snd{q}$ or $1 = \snd{q}$
                                        by \cref{real_one_leq_naturalsplus}.
                                    $0\in\reals$ by \cref{real_add_ident}.
                                    $0 < 1$ by \cref{real_zero_lt_one}.
                                    \begin{byCase}
                                    \caseOf{$1 < \snd{q}$.}
                                        Thus $0 < \snd{q}$ by \cref{real_lt_trans}.
                                        Thus $\snd{q}\neq 0$ by \cref{real_pos_nonzero}.
                                    \caseOf{$1 = \snd{q}$.}
                                        Thus $0 < \snd{q}$ by \cref{real_lt_subst_right}.
                                        Thus $\snd{q}\neq 0$ by \cref{real_pos_nonzero}.
                                    \end{byCase}
                                \end{subproof}
                                Thus $\inv{\snd{q}}\in\reals$ by \cref{real_mul_inverse}.
                                Thus $r\in\reals$ by \cref{real_mul_closed}.
                            \caseOf{$r\in\{0\}$.}
                                Thus $r = 0$ by \cref{singleton_elim}.
                                $0\in\reals$ by \cref{real_add_ident}.
                                Thus $r\in\reals$.
                            \end{byCase}
                        \end{subproof}
                        Thus $Q_y\subseteq\reals$ by \cref{subseteq}.
                    \end{subproof}
                    Thus $p\in\reals$ by \cref{intervalopen}.
                    Show that for all $r\in Q_y$ we have $p \leq r$.
                    \begin{subproof}
                        Fix $r$.
                        Assume $r\in Q_y$.
                        Thus $r\in\reals$ by \cref{subseteq}.
                        Thus $y\in U(r)$ by \cref{urysohn_level_set}.
                        Show that $p \neq r$.
                        \begin{subproof}
                            Suppose not.
                            Thus $p = r$.
                            Thus $y\in U(p)$.
                            Show that $U(p)\subseteq X$.
                            \begin{subproof}
                                Thus $p\in\unitRationals$.
                                Take $k_p$ such that $k_p\in\naturalsPlus$ and $s(k_p) = p$
                                    by \cref{unit_rationals_sequence}.
                                Thus $u$ is a sequence in $G$.
                                Thus $u\in\funs{\naturalsPlus}{G}$ by \cref{sequence_in}.
                                Thus $u(k_p)\in G$ by \cref{funs_type_apply}.
                                Take $n_p,t_p$ such that
                                    $n_p\in\naturalsPlus$ and
                                    $t_p\in\urysohnAdmissible{s}{X}{T}{A}{B}{n_p}$ and
                                    $u(k_p) = (n_p,t_p)$
                                    by \cref{urysohn_good_pair_decomp}.
                                Thus $\fst{u(k_p)} = n_p$ by \cref{fst_eq}.
                                Thus $k_p = n_p$.
                                Thus $u(k_p) = (k_p,t_p)$ by \cref{pair_eq_iff}.
                                Thus $\snd{u(k_p)} = t_p$ by \cref{snd_eq}.
                                Thus $t_p\in F$.
                                Thus $t_p\in\funs{\{0,1\}\union \img{s}{\initialSegment{k_p}}}{\pow{X}}$
                                    by \cref{urysohn_admissible}.
                                Thus $t_p$ is a function by \cref{funs_is_function}.
                                Thus $t_p$ is right-unique and $t_p$ is a relation.
                                Thus $\dom{t_p} = \{0,1\}\union \img{s}{\initialSegment{k_p}}$
                                    by \cref{funs_elim}.
                                Show that $p\in\dom{t_p}$.
                                \begin{subproof}
                                    Thus $k_p \leq k_p$.
                                    Thus $k_p\in\initialSegment{k_p}$
                                        by \cref{initial_segment_naturalsplus}.
                                    Let $y_0 = k_p$.
                                    Thus $y_0\in\initialSegment{k_p}$.
                                    Thus $s(y_0) = s(k_p)$.
                                    Thus $s\in\funs{\naturalsPlus}{\unitRationals}$ by \cref{sequence_in}.
                                    Thus $(y_0,s(y_0))\in s$ by \cref{funs_is_function,funs_elim,subseteq,function_apply_elim}.
                                    Thus $(y_0,s(k_p))\in s$.
                                    Thus $s(k_p)\in\img{s}{\initialSegment{k_p}}$ by \cref{img_iff}.
                                    Thus $p\in\img{s}{\initialSegment{k_p}}$.
                                    Thus $p\in\dom{t_p}$ by \cref{union_intro_right}.
                                \end{subproof}
                                Thus $t_p(p)\subseteq X$ by \cref{funs_type_apply,powerset_elim}.
                                Thus $(p,t_p(p))\in t_p$ by \cref{function_apply_elim}.
                                Thus $(p,t_p(p))\in U$ by \cref{unions_intro}.
                                Thus $U(p) = t_p(p)$ by \cref{function_apply_intro}.
                                Thus $U(p)\subseteq X$.
                            \end{subproof}
                            Thus $U(p)\subseteq\closure{U(p)}{X}{T}$
                                by \cref{subseteq_closure,normal_space}.
                            Thus $y\in\closure{U(p)}{X}{T}$ by \cref{subseteq}.
                            Contradiction.
                        \end{subproof}
                        Thus $p \neq r$.
                        Thus $p < r$ or $r < p$ by \cref{real_lt_trichotomy}.
                        \begin{byCase}
                        \caseOf{$p < r$.}
                            Thus $p \leq r$ by \cref{real_lt_imp_leq}.
                        \caseOf{$r < p$.}
                            Show that $p \leq r$.
                            \begin{subproof}
                                Suppose not.
                                Show that $\closure{U(r)}{X}{T}\subseteq U(p)$.
                                \begin{subproof}
                                    Thus $p\in\unitRationals$.
                                    Take $k_p$ such that $k_p\in\naturalsPlus$ and $s(k_p) = p$
                                        by \cref{unit_rationals_sequence}.
                                    Thus $u$ is a sequence in $G$.
                                    Thus $u\in\funs{\naturalsPlus}{G}$ by \cref{sequence_in}.
                                    Thus $u(k_p)\in G$ by \cref{funs_type_apply}.
                                    Take $n_p,t_p$ such that
                                        $n_p\in\naturalsPlus$ and
                                        $t_p\in\urysohnAdmissible{s}{X}{T}{A}{B}{n_p}$ and
                                        $u(k_p) = (n_p,t_p)$
                                        by \cref{urysohn_good_pair_decomp}.
                                    Thus $\fst{u(k_p)} = n_p$ by \cref{fst_eq}.
                                    Thus $k_p = n_p$.
                                    Thus $u(k_p) = (k_p,t_p)$ by \cref{pair_eq_iff}.
                                    Thus $\snd{u(k_p)} = t_p$ by \cref{snd_eq}.
                                    Thus $t_p\in F$.
                                    Thus $t_p\in\funs{\{0,1\}\union \img{s}{\initialSegment{k_p}}}{\pow{X}}$
                                        by \cref{urysohn_admissible}.
                                    Thus $t_p$ is a function by \cref{funs_is_function}.
                                    Thus $t_p$ is right-unique and $t_p$ is a relation.
                                    Thus $\dom{t_p} = \{0,1\}\union \img{s}{\initialSegment{k_p}}$
                                        by \cref{funs_elim}.
                                    Show that $p\in\dom{t_p}$.
                                    \begin{subproof}
                                        Thus $k_p \leq k_p$.
                                        Thus $k_p\in\initialSegment{k_p}$
                                            by \cref{initial_segment_naturalsplus}.
                                        Let $y_p = k_p$.
                                        Thus $y_p\in\initialSegment{k_p}$.
                                        Thus $s(y_p) = s(k_p)$.
                                        Thus $s\in\funs{\naturalsPlus}{\unitRationals}$ by \cref{sequence_in}.
                                        Thus $(y_p,s(y_p))\in s$ by \cref{funs_is_function,funs_elim,subseteq,function_apply_elim}.
                                        Thus $(y_p,s(k_p))\in s$.
                                        Thus $s(k_p)\in\img{s}{\initialSegment{k_p}}$ by \cref{img_iff}.
                                        Thus $p\in\img{s}{\initialSegment{k_p}}$.
                                        Thus $p\in\dom{t_p}$ by \cref{union_intro_right}.
                                    \end{subproof}
                                    Thus $(p,t_p(p))\in t_p$ by \cref{function_apply_elim}.
                                    Thus $(p,t_p(p))\in U$ by \cref{unions_intro}.
                                    Thus $U(p) = t_p(p)$ by \cref{function_apply_intro}.
                                    Thus $r\in\unitRationalsZero$ by \cref{urysohn_level_set}.
                                    Thus $r\in\unitRationals$ or $r\in\{0\}$
                                        by \cref{unit_rationals_zero,union_iff}.
                                    \begin{byCase}
                                    \caseOf{$r\in\unitRationals$.}
                                        Take $k_r$ such that $k_r\in\naturalsPlus$ and $s(k_r) = r$
                                            by \cref{unit_rationals_sequence}.
                                        Thus $u(k_r)\in G$ by \cref{funs_type_apply}.
                                        Take $n_r,t_r$ such that
                                            $n_r\in\naturalsPlus$ and
                                            $t_r\in\urysohnAdmissible{s}{X}{T}{A}{B}{n_r}$ and
                                            $u(k_r) = (n_r,t_r)$
                                            by \cref{urysohn_good_pair_decomp}.
                                        Thus $\fst{u(k_r)} = n_r$ by \cref{fst_eq}.
                                        Thus $k_r = n_r$.
                                        Thus $u(k_r) = (k_r,t_r)$ by \cref{pair_eq_iff}.
                                        Thus $\snd{u(k_r)} = t_r$ by \cref{snd_eq}.
                                        Thus $t_r\in F$.
                                        Thus $t_r\in\funs{\{0,1\}\union \img{s}{\initialSegment{k_r}}}{\pow{X}}$
                                            by \cref{urysohn_admissible}.
                                        Thus $t_r$ is a function by \cref{funs_is_function}.
                                        Thus $t_r$ is right-unique and $t_r$ is a relation.
                                        Thus $\dom{t_r} = \{0,1\}\union \img{s}{\initialSegment{k_r}}$
                                            by \cref{funs_elim}.
                                        Show that $r\in\dom{t_r}$.
                                        \begin{subproof}
                                            Thus $k_r \leq k_r$.
                                            Thus $k_r\in\initialSegment{k_r}$
                                                by \cref{initial_segment_naturalsplus}.
                                            Let $y_r = k_r$.
                                            Thus $y_r\in\initialSegment{k_r}$.
                                            Thus $s(y_r) = s(k_r)$.
                                            Thus $s$ is a sequence in $\unitRationals$.
                                            Thus $s\in\funs{\naturalsPlus}{\unitRationals}$ by \cref{sequence_in}.
                                            Thus $(y_r,s(y_r))\in s$ by \cref{funs_is_function,funs_elim,subseteq,function_apply_elim}.
                                            Thus $(y_r,s(k_r))\in s$.
                                            Thus $s(k_r)\in\img{s}{\initialSegment{k_r}}$ by \cref{img_iff}.
                                            Thus $r\in\img{s}{\initialSegment{k_r}}$.
                                            Thus $r\in\dom{t_r}$ by \cref{union_intro_right}.
                                        \end{subproof}
                                        Thus $(r,t_r(r))\in t_r$ by \cref{function_apply_elim}.
                                        Thus $(r,t_r(r))\in U$ by \cref{unions_intro}.
                                        Thus $U(r) = t_r(r)$ by \cref{function_apply_intro}.
                                        Thus $t_r\subseteq t_p$ or $t_p\subseteq t_r$
                                            by \cref{urysohn_unit_interval_chain}.
                                        \begin{byCase}
                                        \caseOf{$t_r\subseteq t_p$.}
                                            Thus $(r,t_r(r))\in t_p$ by \cref{subseteq}.
                                            Thus $t_p(r) = t_r(r)$ by \cref{function_apply_intro}.
                                            Thus $U(r) = t_p(r)$.
                                            Thus $r\in\dom{t_p}$ by \cref{dom_intro}.
                                            Thus $p\in\dom{t_p}$ by \cref{dom_intro}.
                                            Thus $r < p$.
                                            Thus $\closure{t_p(r)}{X}{T}\subseteq t_p(p)$
                                                by \cref{urysohn_admissible}.
                                            Thus $\closure{U(r)}{X}{T}\subseteq U(p)$.
                                        \caseOf{$t_p\subseteq t_r$.}
                                            Thus $(p,t_p(p))\in t_r$ by \cref{subseteq}.
                                            Thus $t_r(p) = t_p(p)$ by \cref{function_apply_intro}.
                                            Thus $U(p) = t_r(p)$.
                                            Thus $r\in\dom{t_r}$ by \cref{dom_intro}.
                                            Thus $p\in\dom{t_r}$ by \cref{dom_intro}.
                                            Thus $r < p$.
                                            Thus $\closure{t_r(r)}{X}{T}\subseteq t_r(p)$
                                                by \cref{urysohn_admissible}.
                                            Thus $\closure{U(r)}{X}{T}\subseteq U(p)$.
                                        \end{byCase}
                                    \caseOf{$r\in\{0\}$.}
                                        Thus $r = 0$ by \cref{singleton_elim}.
                                        Thus $r\in\{0,1\}$ by \cref{upair_intro_left}.
                                        Thus $r\in\{0,1\}\union \img{s}{\initialSegment{k_p}}$
                                            by \cref{union_intro_left}.
                                        Thus $r\in\dom{t_p}$ by \cref{funs_elim}.
                                        Thus $(r,t_p(r))\in t_p$ by \cref{function_apply_elim}.
                                        Thus $(r,t_p(r))\in U$ by \cref{unions_intro}.
                                        Thus $U(r) = t_p(r)$ by \cref{function_apply_intro}.
                                        Thus $r\in\dom{t_p}$ by \cref{dom_intro}.
                                        Thus $p\in\dom{t_p}$ by \cref{dom_intro}.
                                        Thus $r < p$.
                                        Thus $\closure{t_p(r)}{X}{T}\subseteq t_p(p)$
                                            by \cref{urysohn_admissible}.
                                        Thus $\closure{U(r)}{X}{T}\subseteq U(p)$.
                                    \end{byCase}
                                \end{subproof}
                                Show that $U(r)\subseteq X$.
                                \begin{subproof}
                                    Thus $r\in\unitRationalsZero$ by \cref{urysohn_level_set}.
                                    Thus $r\in\unitRationals$ or $r\in\{0\}$
                                        by \cref{unit_rationals_zero,union_iff}.
                                    \begin{byCase}
                                    \caseOf{$r\in\unitRationals$.}
                                        Take $k_r$ such that $k_r\in\naturalsPlus$ and $s(k_r) = r$
                                            by \cref{unit_rationals_sequence}.
                                        Thus $u$ is a sequence in $G$.
                                        Thus $u\in\funs{\naturalsPlus}{G}$ by \cref{sequence_in}.
                                        Thus $u(k_r)\in G$ by \cref{funs_type_apply}.
                                        Take $n_r,t_r$ such that
                                            $n_r\in\naturalsPlus$ and
                                            $t_r\in\urysohnAdmissible{s}{X}{T}{A}{B}{n_r}$ and
                                            $u(k_r) = (n_r,t_r)$
                                            by \cref{urysohn_good_pair_decomp}.
                                        Thus $\fst{u(k_r)} = n_r$ by \cref{fst_eq}.
                                        Thus $k_r = n_r$.
                                        Thus $u(k_r) = (k_r,t_r)$ by \cref{pair_eq_iff}.
                                        Thus $\snd{u(k_r)} = t_r$ by \cref{snd_eq}.
                                        Thus $t_r\in F$.
                                        Thus $t_r\in\funs{\{0,1\}\union \img{s}{\initialSegment{k_r}}}{\pow{X}}$
                                            by \cref{urysohn_admissible}.
                                        Thus $t_r$ is a function by \cref{funs_is_function}.
                                        Thus $t_r$ is right-unique and $t_r$ is a relation.
                                        Thus $\dom{t_r} = \{0,1\}\union \img{s}{\initialSegment{k_r}}$
                                            by \cref{funs_elim}.
                                        Show that $r\in\dom{t_r}$.
                                        \begin{subproof}
                                            Thus $k_r \leq k_r$.
                                            Thus $k_r\in\initialSegment{k_r}$
                                                by \cref{initial_segment_naturalsplus}.
                                            Let $y_r = k_r$.
                                            Thus $y_r\in\initialSegment{k_r}$.
                                            Thus $s(y_r) = s(k_r)$.
                                            Thus $s$ is a sequence in $\unitRationals$.
                                            Thus $s\in\funs{\naturalsPlus}{\unitRationals}$ by \cref{sequence_in}.
                                            Thus $s$ is a function by \cref{funs_is_function}.
                                            Thus $\dom{s} = \naturalsPlus$ by \cref{funs_elim}.
                                            Thus $y_r\in\dom{s}$ by \cref{subseteq}.
                                            Thus $(y_r,s(y_r))\in s$ by \cref{function_apply_elim}.
                                            Thus $(y_r,s(k_r))\in s$.
                                            Thus $s(k_r)\in\img{s}{\initialSegment{k_r}}$ by \cref{img_iff}.
                                            Thus $r\in\img{s}{\initialSegment{k_r}}$.
                                            Thus $r\in\dom{t_r}$ by \cref{union_intro_right}.
                                        \end{subproof}
                                        Thus $(r,t_r(r))\in t_r$ by \cref{function_apply_elim}.
                                        Thus $(r,t_r(r))\in U$ by \cref{unions_intro}.
                                        Thus $U(r) = t_r(r)$ by \cref{function_apply_intro}.
                                        Thus $t_r(r)\subseteq X$ by \cref{funs_type_apply,powerset_elim}.
                                        Thus $U(r)\subseteq X$.
                                    \caseOf{$r\in\{0\}$.}
                                        Thus $r = 0$ by \cref{singleton_elim}.
                                        Thus $p\in\unitRationals$.
                                        Take $k_p$ such that $k_p\in\naturalsPlus$ and $s(k_p) = p$
                                            by \cref{unit_rationals_sequence}.
                                        Thus $u$ is a sequence in $G$.
                                        Thus $u\in\funs{\naturalsPlus}{G}$ by \cref{sequence_in}.
                                        Thus $u(k_p)\in G$ by \cref{funs_type_apply}.
                                        Take $n_p,t_p$ such that
                                            $n_p\in\naturalsPlus$ and
                                            $t_p\in\urysohnAdmissible{s}{X}{T}{A}{B}{n_p}$ and
                                            $u(k_p) = (n_p,t_p)$
                                            by \cref{urysohn_good_pair_decomp}.
                                        Thus $\fst{u(k_p)} = n_p$ by \cref{fst_eq}.
                                        Thus $k_p = n_p$.
                                        Thus $u(k_p) = (k_p,t_p)$ by \cref{pair_eq_iff}.
                                        Thus $\snd{u(k_p)} = t_p$ by \cref{snd_eq}.
                                        Thus $t_p\in F$.
                                        Thus $t_p\in\funs{\{0,1\}\union \img{s}{\initialSegment{k_p}}}{\pow{X}}$
                                            by \cref{urysohn_admissible}.
                                        Thus $t_p$ is a function by \cref{funs_is_function}.
                                        Thus $t_p$ is right-unique and $t_p$ is a relation.
                                        Thus $\dom{t_p} = \{0,1\}\union \img{s}{\initialSegment{k_p}}$
                                            by \cref{funs_elim}.
                                        Thus $r\in\{0,1\}$ by \cref{upair_intro_left}.
                                        Thus $r\in\{0,1\}\union \img{s}{\initialSegment{k_p}}$
                                            by \cref{union_intro_left}.
                                        Thus $r\in\dom{t_p}$ by \cref{funs_elim}.
                                        Thus $(r,t_p(r))\in t_p$ by \cref{function_apply_elim}.
                                        Thus $(r,t_p(r))\in U$ by \cref{unions_intro}.
                                        Thus $U(r) = t_p(r)$ by \cref{function_apply_intro}.
                                        Thus $t_p(r)\subseteq X$ by \cref{funs_type_apply,powerset_elim}.
                                        Thus $U(r)\subseteq X$.
                                    \end{byCase}
                                \end{subproof}
                                Thus $U(r)\subseteq\closure{U(r)}{X}{T}$
                                    by \cref{subseteq_closure,normal_space}.
                                Thus $y\in\closure{U(r)}{X}{T}$ by \cref{subseteq}.
                                Thus $y\in U(p)$ by \cref{subseteq}.
                                Show that $U(p)\subseteq X$.
                                \begin{subproof}
                                    Thus $p\in\unitRationals$.
                                    Take $k_p$ such that $k_p\in\naturalsPlus$ and $s(k_p) = p$
                                        by \cref{unit_rationals_sequence}.
                                    Thus $u$ is a sequence in $G$.
                                    Thus $u\in\funs{\naturalsPlus}{G}$ by \cref{sequence_in}.
                                    Thus $u(k_p)\in G$ by \cref{funs_type_apply}.
                                    Take $n_p,t_p$ such that
                                        $n_p\in\naturalsPlus$ and
                                        $t_p\in\urysohnAdmissible{s}{X}{T}{A}{B}{n_p}$ and
                                        $u(k_p) = (n_p,t_p)$
                                        by \cref{urysohn_good_pair_decomp}.
                                    Thus $\fst{u(k_p)} = n_p$ by \cref{fst_eq}.
                                    Thus $k_p = n_p$.
                                    Thus $u(k_p) = (k_p,t_p)$ by \cref{pair_eq_iff}.
                                    Thus $\snd{u(k_p)} = t_p$ by \cref{snd_eq}.
                                    Thus $t_p\in F$.
                                    Thus $t_p\in\funs{\{0,1\}\union \img{s}{\initialSegment{k_p}}}{\pow{X}}$
                                        by \cref{urysohn_admissible}.
                                    Thus $t_p$ is a function by \cref{funs_is_function}.
                                    Thus $t_p$ is right-unique and $t_p$ is a relation.
                                    Thus $\dom{t_p} = \{0,1\}\union \img{s}{\initialSegment{k_p}}$
                                        by \cref{funs_elim}.
                                    Show that $p\in\dom{t_p}$.
                                    \begin{subproof}
                                        Thus $k_p \leq k_p$.
                                        Thus $k_p\in\initialSegment{k_p}$
                                            by \cref{initial_segment_naturalsplus}.
                                        Let $y_p = k_p$.
                                        Thus $y_p\in\initialSegment{k_p}$.
                                        Thus $s(y_p) = s(k_p)$.
                                        Thus $s$ is a sequence in $\unitRationals$.
                                        Thus $s\in\funs{\naturalsPlus}{\unitRationals}$ by \cref{sequence_in}.
                                        Thus $(y_p,s(y_p))\in s$ by \cref{funs_is_function,funs_elim,subseteq,function_apply_elim}.
                                        Thus $(y_p,s(k_p))\in s$.
                                        Thus $s(k_p)\in\img{s}{\initialSegment{k_p}}$ by \cref{img_iff}.
                                        Thus $p\in\img{s}{\initialSegment{k_p}}$.
                                        Thus $p\in\dom{t_p}$ by \cref{union_intro_right}.
                                    \end{subproof}
                                    Thus $t_p(p)\subseteq X$ by \cref{funs_type_apply,powerset_elim}.
                                    Thus $(p,t_p(p))\in t_p$ by \cref{function_apply_elim}.
                                    Thus $(p,t_p(p))\in U$ by \cref{unions_intro}.
                                    Thus $U(p) = t_p(p)$ by \cref{function_apply_intro}.
                                    Thus $U(p)\subseteq X$.
                                \end{subproof}
                                    Thus $U(p)\subseteq\closure{U(p)}{X}{T}$
                                        by \cref{subseteq_closure,normal_space}.
                                Thus $y\in\closure{U(p)}{X}{T}$ by \cref{subseteq}.
                                Contradiction.
                            \end{subproof}
                        \end{byCase}
                    \end{subproof}
                    Thus $p$ is a lower bound of $Q_y$ by \cref{real_lower_bound}.
                \end{subproof}
                Show that $f(y)$ is an infimum of $Q_y$.
                \begin{subproof}
                    Thus $(y,f(y))\in f$ by \cref{funs_is_function,funs_elim,subseteq,function_apply_elim}.
                    Take $x_0,y_0$ such that
                        $x_0\in X$ and
                        $y_0\in\reals$ and
                        $(y,f(y)) = (x_0,y_0)$ and
                        $y_0$ is an infimum of $\urysohnLevelSet{U}{x_0}$.
                    Thus $y = x_0$ and $f(y) = y_0$ by \cref{pair_eq_iff}.
                    Thus $f(y)$ is an infimum of $Q_y$.
                \end{subproof}
                Thus $f(y) \geq p$ by \cref{real_infimum}.
                Thus $p \leq f(y)$.
            \end{subproof}
            Show that $a < f(y)$.
            \begin{subproof}
                Thus $a < p$ and $p\in\reals$ by \cref{intervalopen}.
                Thus $p < f(y)$ or $p = f(y)$ by \cref{real_leq_split}.
                \begin{byCase}
                \caseOf{$p < f(y)$.}
                    Thus $a < f(y)$ by \cref{real_lt_trans}.
                \caseOf{$p = f(y)$.}
                    Thus $a < f(y)$ by \cref{real_lt_subst_right}.
                \end{byCase}
            \end{subproof}
            Show that $f(y) < b$.
            \begin{subproof}
                Thus $q < b$ and $q\in\reals$ by \cref{intervalopen}.
                Thus $f(y) < q$ or $f(y) = q$ by \cref{real_leq_split}.
                \begin{byCase}
                \caseOf{$f(y) < q$.}
                    Thus $f(y) < b$ by \cref{real_lt_trans}.
                \caseOf{$f(y) = q$.}
                    Thus $f(y) < b$ by \cref{real_lt_subst_left}.
                \end{byCase}
            \end{subproof}
            Thus $f(y)\in V$ by \cref{intervalopen,subseteq}.
        \end{subproof}
                                    Show that $\img{f}{W}\subseteq V$.
                                    \begin{subproof}
    Show that for all $z\in\img{f}{W}$ we have $z\in V$.
    \begin{subproof}
        Fix $z$.
        Assume $z\in\img{f}{W}$.
        Take $y$ such that $y\in W$ and $(y,z)\in f$ by \cref{img_iff}.
        Thus $f(y) = z$ by \cref{funs_is_function,function_apply_intro}.
        Thus $f(y)\in V$.
        Thus $z\in V$.
    \end{subproof}
    Thus $\img{f}{W}\subseteq V$ by \cref{subseteq}.
                                    \end{subproof}
\end{proof}

% from §33 helper lemma 23: continuity at a point
\begin{lemma}\label{urysohn_unit_interval_continuous_at_point}
    Assume $X$ is normal with respect to $T$.
    Assume $A$ is closed in $X$ with respect to $T$.
    Assume $B$ is closed in $X$ with respect to $T$.
    Assume $A$ is disjoint from $B$.
    Assume $s$ is a sequence in $\unitRationals$.
    Assume for all $p$ such that $p\in\unitRationals$ there exists $n\in\naturalsPlus$ such that $s(n) = p$.
    Let $G = \urysohnGoodPairs{s}{X}{T}{A}{B}$.
    Let $R = \{(z,w) \mid z\in G, w\in G \mid
        \fst{w} = \fst{z} + 1 \land \snd{z}\subseteq\snd{w}\}$.
    Assume $c$ is a function on $G$.
    Assume for all $z\in G$ we have $z\mathrel{R} \apply{c}{z}$.
    Assume $u$ is a sequence in $G$.
    Assume $u(1) = a_0$.
    Assume $a_0 = (1,u_1)$.
    Assume for all $n\in\naturalsPlus$ we have $u(n + 1) = \apply{c}{u(n)}$.
    Assume for all $n\in\naturalsPlus$ we have $\fst{u(n)} = n$.
    Let $F = \{\snd{u(n)} \mid n\in\naturalsPlus\}$.
    Let $U = \unions{F}$.
    Assume $U$ is a function.
    Let $I_0 = \intervalclosed{0}{1}$.
    Let $f = \{(x,y) \mid x\in X, y\in\reals \mid
        \text{$y$ is an infimum of $\urysohnLevelSet{U}{x}$}\}$.
    Assume $f\in\funs{X}{\reals}$.
    Assume $\img{f}{X}\subseteq I_0$.
    Assume $x\in X$.
    Assume $V$ is open in $\reals$ with respect to $\standardTopologyR$.
    Assume $f(x)\in V$.
    Then there exists $U_0$ such that
        $U_0$ is open in $X$ with respect to $T$ and $x\in U_0$ and $\img{f}{U_0}\subseteq V$.
\end{lemma}
\begin{proof}
    Thus $V\in\standardTopologyR$ by \cref{open_in}.
    Thus $V\in\basisTopology{\standardBasisR}{\reals}$ by \cref{standard_topology_reals}.
    Take $B_1$ such that $B_1\in\standardBasisR$ and $f(x)\in B_1$ and $B_1\subseteq V$
        by \cref{topology_generated_by_basis}.
    Take $a,b\in\reals$ such that $a < b$ and $B_1 = \intervalopen{a}{b}$
        by \cref{standard_basis_reals}.
    Thus $a < f(x)$ and $f(x) < b$ by \cref{intervalopen}.
    Show that $f(x)\in I_0$.
    \begin{subproof}
        Thus $f$ is a function by \cref{funs_is_function}.
        Thus $\dom{f} = X$ by \cref{funs_elim}.
        Thus $x\in\dom{f}$ by \cref{subseteq}.
        Thus $(x,f(x))\in f$ by \cref{function_apply_elim}.
        Thus $f(x)\in\img{f}{X}$ by \cref{img_iff}.
        Thus $\img{f}{X}\subseteq I_0$.
        Thus $f(x)\in I_0$ by \cref{subseteq}.
    \end{subproof}
    Thus $0 \leq f(x)$ by \cref{intervalclosed}.
    Thus $f(x)\in\reals$ by \cref{funs_type_apply}.
    $0\in\reals$ by \cref{real_add_ident}.
    Thus $0 < f(x)$ or $0 = f(x)$.
    \begin{byCase}
    \caseOf{$0 = f(x)$.}
        Thus there exists $U_0$ such that
            $U_0$ is open in $X$ with respect to $T$ and $x\in U_0$ and $\img{f}{U_0}\subseteq V$
            by \cref{urysohn_unit_interval_continuous_zero}.
    \caseOf{$0 < f(x)$.}
        Take $p,q,r_0,W$ such that
            $p\in\unitRationals$ and
            $p\in\intervalopen{a}{f(x)}$ and
            $p\in\intervalopen{0}{f(x)}$ and
            $q\in\unitRationals$ and
            $q\in\intervalopen{f(x)}{b}$ and
            $r_0\in\unitRationals$ and
            $r_0\in\intervalopen{p}{f(x)}$ and
            $p\in\unitRationalsZero$ and
            $q\in\unitRationalsZero$ and
            $r_0\in\unitRationalsZero$ and
            $W = U(q)\setminus \closure{U(p)}{X}{T}$ and
            $W$ is open in $X$ with respect to $T$
            by \cref{urysohn_unit_interval_continuous_pos_open}.
        Thus $x\in W$ by \cref{urysohn_unit_interval_continuous_pos_x_in_W}.
        Thus $W\in T$ by \cref{open_in}.
        $T$ is a topology on $X$ by \cref{normal_space}.
        Thus $T\subseteq\pow{X}$ by \cref{topology_on}.
        Thus $W\in\pow{X}$ by \cref{subseteq}.
        Thus $W\subseteq X$ by \cref{powerset_elim}.
        Thus $\img{f}{W}\subseteq V$ by \cref{urysohn_unit_interval_continuous_pos_img}.
        Thus there exists $U_0$ such that
            $U_0$ is open in $X$ with respect to $T$ and $x\in U_0$ and $\img{f}{U_0}\subseteq V$.
    \end{byCase}
\end{proof}

% from §33 helper lemma 24: Urysohn function is continuous to reals
\begin{lemma}\label{urysohn_unit_interval_continuous_reals}
    Assume $X$ is normal with respect to $T$.
    Assume $A$ is closed in $X$ with respect to $T$.
    Assume $B$ is closed in $X$ with respect to $T$.
    Assume $A$ is disjoint from $B$.
    Assume $s$ is a sequence in $\unitRationals$.
    Assume for all $p$ such that $p\in\unitRationals$ there exists $n\in\naturalsPlus$ such that $s(n) = p$.
    Let $G = \urysohnGoodPairs{s}{X}{T}{A}{B}$.
    Let $R = \{(z,w) \mid z\in G, w\in G \mid
        \fst{w} = \fst{z} + 1 \land \snd{z}\subseteq\snd{w}\}$.
    Assume $c$ is a function on $G$.
    Assume for all $z\in G$ we have $z\mathrel{R} \apply{c}{z}$.
    Assume $u$ is a sequence in $G$.
    Assume $u(1) = a_0$.
    Assume $a_0 = (1,u_1)$.
    Assume for all $n\in\naturalsPlus$ we have $u(n + 1) = \apply{c}{u(n)}$.
    Assume for all $n\in\naturalsPlus$ we have $\fst{u(n)} = n$.
    Let $F = \{\snd{u(n)} \mid n\in\naturalsPlus\}$.
    Let $U = \unions{F}$.
    Assume $U$ is a function.
    Let $I_0 = \intervalclosed{0}{1}$.
    Let $f = \{(x,y) \mid x\in X, y\in\reals \mid
        \text{$y$ is an infimum of $\urysohnLevelSet{U}{x}$}\}$.
    Assume $f\in\funs{X}{\reals}$.
    Assume $\img{f}{X}\subseteq I_0$.
    Then $f$ is continuous from $X$ to $\reals$ with respect to $T$ and $\standardTopologyR$.
\end{lemma}
\begin{proof}
    Show that for all $x\in X$ we have
        $f$ is continuous at $x$ from $X$ to $\reals$ with respect to $T$ and $\standardTopologyR$.
    \begin{subproof}
        Fix $x\in X$.
        Show that for all $V$ such that $V$ is open in $\reals$ with respect to $\standardTopologyR$ and
            $f(x)\in V$ there exists $U_0$ such that
                $U_0$ is open in $X$ with respect to $T$ and $x\in U_0$ and $\img{f}{U_0}\subseteq V$.
        \begin{subproof}
            Fix $V$.
            Assume $V$ is open in $\reals$ with respect to $\standardTopologyR$.
            Assume $f(x)\in V$.
            Thus there exists $U_0$ such that
                $U_0$ is open in $X$ with respect to $T$ and $x\in U_0$ and $\img{f}{U_0}\subseteq V$
                by \cref{urysohn_unit_interval_continuous_at_point}.
        \end{subproof}
        Show that $\standardTopologyR$ is a topology on $\reals$.
        \begin{subproof}
            $\standardBasisR$ is a basis on $\reals$ by \cref{standard_basis_is_basis}.
            Thus $\basisTopology{\standardBasisR}{\reals}$ is a topology on $\reals$
                by \cref{basis_topology_is_topology}.
            Thus $\standardTopologyR$ is a topology on $\reals$ by \cref{standard_topology_reals}.
        \end{subproof}
        $T$ is a topology on $X$ by \cref{normal_space}.
        Thus $f$ is continuous at $x$ from $X$ to $\reals$ with respect to $T$ and $\standardTopologyR$
            by \cref{continuous_at_open}.
    \end{subproof}
    $\standardBasisR$ is a basis on $\reals$ by \cref{standard_basis_is_basis}.
    Thus $\standardTopologyR$ is a topology on $\reals$
        by \cref{basis_topology_is_topology,standard_topology_reals}.
    $T$ is a topology on $X$ by \cref{normal_space}.
    Thus $f$ is continuous from $X$ to $\reals$ with respect to $T$ and $\standardTopologyR$
        by \cref{continuous_at_implies_continuous}.
\end{proof}

% from §33 helper lemma 34: admissible values above $1$
\begin{lemma}\label{urysohn_admissible_above_one}
    Assume $f\in\urysohnAdmissible{s}{X}{T}{A}{B}{n}$.
    Assume $p\in\dom{f}$.
    Assume $1 < p$.
    Then $f(p) = X$.
\end{lemma}
\begin{proof}
    Thus $f\in\funs{\{0,1\}\union \img{s}{\initialSegment{n}}}{\pow{X}}$
        by \cref{urysohn_admissible}.
    Thus $\dom{f} = \{0,1\}\union \img{s}{\initialSegment{n}}$ by \cref{funs_elim}.
    Thus $0\in\{0,1\}\union \img{s}{\initialSegment{n}}$ by \cref{upair_intro_left,union_intro_left}.
    Thus $0\in\dom{f}$ by \cref{funs_elim}.
    Thus $1\in\{0,1\}\union \img{s}{\initialSegment{n}}$ by \cref{upair_intro_right,union_intro_left}.
    Thus $1\in\dom{f}$ by \cref{funs_elim}.
    Let $u = 0$.
    Let $v = 1$.
    $u < v$ by \cref{real_zero_lt_one}.
    Thus $u\in\dom{f}$ and $v\in\dom{f}$.
    Thus for all $q\in\dom{f}$ if $1 < q$ then $f(q) = X$ by \cref{urysohn_admissible}.
    Thus $f(p) = X$.
\end{proof}

% from §33 helper lemma 35: unit rationals are real
\begin{lemma}\label{unit_rational_real}
    Assume $r\in\unitRationals$.
    Then $r\in\reals$.
\end{lemma}
\begin{proof}
    Take $q$ such that
        $q\in\naturalsPlus\times\naturalsPlus$ and
        $r = \fst{q} \rmul \inv{\snd{q}}$
        by \cref{unit_rationals}.
    Take $m,n$ such that $q = (m,n)$ and $m\in\naturalsPlus$ and $n\in\naturalsPlus$
        by \cref{times}.
    Thus $\fst{q} = m$ and $\snd{q} = n$ by \cref{fst_eq,snd_eq}.
    Thus $m\in\reals$ and $n\in\reals$ by \cref{real_naturals_subset_reals,subseteq}.
    Show that $n \neq 0$.
    \begin{subproof}
        Suppose not.
        Thus $n = 0$.
        $1\in\reals$ by \cref{real_mul_ident}.
        Thus $1 < n$ or $1 = n$ by \cref{real_one_leq_naturalsplus}.
        \begin{byCase}
        \caseOf{$1 < n$.}
            $0\in\reals$ by \cref{real_add_ident}.
            $0 < 1$ by \cref{real_zero_lt_one}.
            Thus $0 < n$ by \cref{real_lt_trans}.
            Thus $n \neq 0$ by \cref{real_pos_nonzero}.
            Contradiction.
        \caseOf{$1 = n$.}
            $0 < 1$ by \cref{real_zero_lt_one}.
            Thus $0 < n$ by \cref{real_lt_subst_right}.
            Thus $n \neq 0$ by \cref{real_pos_nonzero}.
            Contradiction.
        \end{byCase}
    \end{subproof}
    Thus $\inv{n}\in\reals$ and $n \rmul \inv{n} = 1$ by \cref{real_mul_inverse}.
    Thus $r\in\reals$ by \cref{real_mul_closed}.
\end{proof}

% from §33 helper lemma 25: extension when $r\in\dom{f}$
\begin{lemma}\label{urysohn_unit_interval_extend_in_dom}
    Assume $X$ is normal with respect to $T$.
    Assume $A$ is closed in $X$ with respect to $T$.
    Assume $B$ is closed in $X$ with respect to $T$.
    Assume $A$ is disjoint from $B$.
    Assume $s$ is a sequence in $\unitRationals$.
    Let $G = \urysohnGoodPairs{s}{X}{T}{A}{B}$.
    Let $R = \{(z,w) \mid z\in G, w\in G \mid
        \fst{w} = \fst{z} + 1 \land \snd{z}\subseteq\snd{w}\}$.
    Suppose $n\in\naturalsPlus$ and $f\in\urysohnAdmissible{s}{X}{T}{A}{B}{n}$ and $z = (n,f)$.
    Let $r = s(n + 1)$.
    Assume $r\in\dom{f}$.
    Then there exists $w$ such that $z\mathrel{R} w$.
\end{lemma}
\begin{proof}
    Thus $f\in\funs{\{0,1\}\union \img{s}{\initialSegment{n}}}{\pow{X}}$ and
        for all $p\in\dom{f}$ we have $f(p)\in T$ and
        $f(1) = X\setminus B$ and
        $A\subseteq f(0)$ and
        $\closure{f(0)}{X}{T}\subseteq f(1)$ and
        for all $p\in\dom{f}$ for all $q\in\dom{f}$ if $p < q$ then
            $\closure{f(p)}{X}{T}\subseteq f(q)$ and
        for all $p\in\dom{f}$ if $1 < p$ then $f(p) = X$
        by \cref{urysohn_admissible}.
    Show that $\dom{f} = \{0,1\}\union \img{s}{\initialSegment{n + 1}}$.
    \begin{subproof}
        Show that $\img{s}{\initialSegment{n}}\subseteq \img{s}{\initialSegment{n + 1}}$.
        \begin{subproof}
            Show that for all $y\in \img{s}{\initialSegment{n}}$ we have
                $y\in \img{s}{\initialSegment{n + 1}}$.
            \begin{subproof}
                Fix $y$.
                Assume $y\in \img{s}{\initialSegment{n}}$.
                Take $k$ such that $k\in\initialSegment{n}$ and $(k,y)\in s$ by \cref{img_iff}.
                $k\in\naturalsPlus$ and $k \leq n$ by \cref{initial_segment_naturalsplus}.
                $n\in\naturalsPlus$.
                $n\in\reals$ and $k\in\reals$ by \cref{real_naturals_subset_reals,subseteq}.
                $1\in\reals$ by \cref{real_naturals_subset_reals,real_one_in_naturals,subseteq}.
                $0\in\reals$ by \cref{real_add_ident}.
                $0 < 1$ by \cref{real_zero_lt_one}.
                $0 + n < 1 + n$ by \cref{real_lt_add}.
                $0 + n = n$ by \cref{real_add_ident}.
                $1 + n = n + 1$ by \cref{real_add_comm}.
                $n + 1\in\reals$ by \cref{real_add_closed}.
                Thus $n < n + 1$ by \cref{real_lt_subst_left}.
                Thus $n \leq n + 1$ by \cref{real_lt_imp_leq}.
                Thus $k \leq n + 1$ by \cref{real_leq_trans}.
                Thus $k\in\initialSegment{n + 1}$ by \cref{initial_segment_naturalsplus}.
                Thus $y\in \img{s}{\initialSegment{n + 1}}$ by \cref{img_iff}.
            \end{subproof}
            Thus $\img{s}{\initialSegment{n}}\subseteq \img{s}{\initialSegment{n + 1}}$
                by \cref{subseteq}.
        \end{subproof}
        Show that $\{0,1\}\union \img{s}{\initialSegment{n}}\subseteq \{0,1\}\union \img{s}{\initialSegment{n + 1}}$.
        \begin{subproof}
            Show that for all $y\in \{0,1\}\union \img{s}{\initialSegment{n}}$ we have
                $y\in \{0,1\}\union \img{s}{\initialSegment{n + 1}}$.
            \begin{subproof}
                Fix $y$.
                Assume $y\in \{0,1\}\union \img{s}{\initialSegment{n}}$.
                Then $y\in \{0,1\}$ or $y\in \img{s}{\initialSegment{n}}$ by \cref{union_iff}.
                \begin{byCase}
                \caseOf{$y\in \{0,1\}$.}
                    Thus $y\in \{0,1\}\union \img{s}{\initialSegment{n + 1}}$ by \cref{union_intro_left}.
                \caseOf{$y\in \img{s}{\initialSegment{n}}$.}
                    Thus $y\in \img{s}{\initialSegment{n + 1}}$ by \cref{subseteq}.
                    Thus $y\in \{0,1\}\union \img{s}{\initialSegment{n + 1}}$ by \cref{union_intro_right}.
                \end{byCase}
            \end{subproof}
            Thus $\{0,1\}\union \img{s}{\initialSegment{n}}\subseteq \{0,1\}\union \img{s}{\initialSegment{n + 1}}$
                by \cref{subseteq}.
        \end{subproof}
        Show that for all $y\in \img{s}{\initialSegment{n + 1}}$ we have $y\in\dom{f}$.
        \begin{subproof}
            Fix $y$.
            Assume $y\in \img{s}{\initialSegment{n + 1}}$.
            Thus there exists $k$ such that $k\in\initialSegment{n + 1}$ and $(k,y)\in s$ by \cref{img_iff}.
            Take $k$ such that $k\in\initialSegment{n + 1}$ and $(k,y)\in s$.
            Thus $s\in\funs{\naturalsPlus}{\unitRationals}$ by \cref{sequence_in}.
            Thus $s$ is a function by \cref{funs_is_function}.
            Thus $s(k) = y$ by \cref{function_apply_intro}.
            \begin{byCase}
            \caseOf{$k\in\initialSegment{n}$.}
                Thus $y\in \img{s}{\initialSegment{n}}$ by \cref{img_iff}.
                Thus $y\in \{0,1\}\union \img{s}{\initialSegment{n}}$ by \cref{union_intro_right}.
                Thus $\dom{f} = \{0,1\}\union \img{s}{\initialSegment{n}}$ by \cref{funs_elim}.
                Thus $y\in\dom{f}$.
            \caseOf{$k\notin\initialSegment{n}$.}
                $k\in\naturalsPlus$ and $k\leq n + 1$ by \cref{initial_segment_naturalsplus}.
                Show that $k = n + 1$.
                \begin{subproof}
                    Suppose not.
                    Thus $k \neq n + 1$.
                    Thus $k \leq n$ by \cref{real_naturals_leq_succ_elim}.
                    Thus $k\in\initialSegment{n}$ by \cref{initial_segment_naturalsplus}.
                    Contradiction.
                \end{subproof}
                Thus $y = r$.
                Thus $y\in\dom{f}$.
            \end{byCase}
        \end{subproof}
        Thus $\img{s}{\initialSegment{n + 1}}\subseteq \dom{f}$ by \cref{subseteq}.
        Show that $\{0,1\}\union \img{s}{\initialSegment{n + 1}}\subseteq \dom{f}$.
        \begin{subproof}
            Show that for all $y\in \{0,1\}\union \img{s}{\initialSegment{n + 1}}$ we have $y\in\dom{f}$.
            \begin{subproof}
                Fix $y$.
                Assume $y\in \{0,1\}\union \img{s}{\initialSegment{n + 1}}$.
                Then $y\in \{0,1\}$ or $y\in \img{s}{\initialSegment{n + 1}}$ by \cref{union_iff}.
                \begin{byCase}
                \caseOf{$y\in \{0,1\}$.}
                    $\dom{f} = \{0,1\}\union \img{s}{\initialSegment{n}}$ by \cref{funs_elim}.
                    Thus $y\in\dom{f}$ by \cref{union_intro_left,subseteq}.
                \caseOf{$y\in \img{s}{\initialSegment{n + 1}}$.}
                    Thus $y\in\dom{f}$ by \cref{subseteq}.
                \end{byCase}
            \end{subproof}
            Thus $\{0,1\}\union \img{s}{\initialSegment{n + 1}}\subseteq \dom{f}$ by \cref{subseteq}.
        \end{subproof}
        Show that $\dom{f}\subseteq \{0,1\}\union \img{s}{\initialSegment{n + 1}}$.
        \begin{subproof}
            Show that for all $y\in\dom{f}$ we have $y\in \{0,1\}\union \img{s}{\initialSegment{n + 1}}$.
            \begin{subproof}
                Fix $y$.
                Assume $y\in\dom{f}$.
                $\dom{f} = \{0,1\}\union \img{s}{\initialSegment{n}}$ by \cref{funs_elim}.
                Thus $y\in \{0,1\}\union \img{s}{\initialSegment{n}}$.
                Then $y\in \{0,1\}$ or $y\in \img{s}{\initialSegment{n}}$ by \cref{union_iff}.
                \begin{byCase}
                \caseOf{$y\in \{0,1\}$.}
                    Thus $y\in \{0,1\}\union \img{s}{\initialSegment{n + 1}}$ by \cref{union_intro_left}.
                \caseOf{$y\in \img{s}{\initialSegment{n}}$.}
                    Thus $y\in \img{s}{\initialSegment{n + 1}}$ by \cref{subseteq}.
                    Thus $y\in \{0,1\}\union \img{s}{\initialSegment{n + 1}}$ by \cref{union_intro_right}.
                \end{byCase}
            \end{subproof}
            Thus $\dom{f}\subseteq \{0,1\}\union \img{s}{\initialSegment{n + 1}}$ by \cref{subseteq}.
        \end{subproof}
        Thus $\dom{f} = \{0,1\}\union \img{s}{\initialSegment{n + 1}}$ by \cref{subseteq_antisymmetric}.
    \end{subproof}
    Let $w = (n + 1,f)$.
    Show that $w\in G$.
    \begin{subproof}
        Thus $n + 1\in\naturalsPlus$ by \cref{naturalsplus_add_closed,real_one_in_naturals}.
        Show that $f\in\urysohnAdmissible{s}{X}{T}{A}{B}{n + 1}$.
        \begin{subproof}
            Thus $f$ is a function to $\pow{X}$ by \cref{funs_elim}.
            Thus $f\in\funs{\{0,1\}\union \img{s}{\initialSegment{n + 1}}}{\pow{X}}$ by \cref{funs_intro}.
            Thus $f\in\urysohnAdmissible{s}{X}{T}{A}{B}{n + 1}$ by \cref{urysohn_admissible}.
        \end{subproof}
        Thus $w\in G$.
    \end{subproof}
    Show that $z\mathrel{R} w$.
    \begin{subproof}
        $z\in G$.
        $w\in G$.
        $\fst{z} = n$ and $\snd{z} = f$ by \cref{fst_eq,snd_eq}.
        $\fst{w} = n + 1$ and $\snd{w} = f$ by \cref{fst_eq,snd_eq}.
        Thus $\fst{w} = \fst{z} + 1$.
        Thus $\snd{z}\subseteq \snd{w}$ by \cref{subseteq_refl}.
        Thus $(z,w)\in R$.
        Thus $z\mathrel{R} w$.
    \end{subproof}
\end{proof}

% from §33 helper lemma 26: extension for $1 < r$ (setup)
\begin{lemma}\label{urysohn_unit_interval_extend_gt1_setup}
    Assume $X$ is normal with respect to $T$.
    Assume $A$ is closed in $X$ with respect to $T$.
    Assume $B$ is closed in $X$ with respect to $T$.
    Assume $A$ is disjoint from $B$.
    Assume $s$ is a sequence in $\unitRationals$.
    Suppose $n\in\naturalsPlus$ and $f\in\urysohnAdmissible{s}{X}{T}{A}{B}{n}$ and $z = (n,f)$.
    Let $r = s(n + 1)$.
    Assume $r\notin\dom{f}$.
    Assume $1 < r$.
    Let $U_r = X$.
    Let $h = \{(r,U_r)\}$.
    Let $g = f\union h$.
    Then $g$ is a function.
\end{lemma}
\begin{proof}
    Show that for all $x$ such that $x\in\dom{f}\inter\dom{h}$ we have $f(x) = h(x)$.
    \begin{subproof}
        Suppose not.
        Then there exists $x$ such that
            $x\in\dom{f}\inter\dom{h}$ and
            $f(x) \neq h(x)$.
        Take $x$ such that $x\in\dom{f}\inter\dom{h}$ and $f(x) \neq h(x)$.
        $x\in\dom{h}$ by \cref{inter_elim_right}.
        Take $y$ such that $(x,y)\in h$ by \cref{dom_iff}.
        Thus $(x,y) = (r,U_r)$ by \cref{singleton_elim}.
        Thus $x = r$ by \cref{pair_eq_iff}.
        Thus $r\in\dom{f}$.
        Contradiction.
    \end{subproof}
    Thus $f\in\funs{\{0,1\}\union \img{s}{\initialSegment{n}}}{\pow{X}}$
        by \cref{urysohn_admissible}.
    Thus $f$ is a function by \cref{funs_is_function}.
    Thus $f$ is right-unique and $f$ is a relation.
    Show that $h$ is right-unique.
    \begin{subproof}
        Show that for all $a,b,b'$ such that $(a,b)\in\{(r,U_r)\}$ and $(a,b')\in\{(r,U_r)\}$ we have
            $b = b'$.
        \begin{subproof}
            Fix $a$.
            Fix $b$.
            Fix $b'$.
            Assume $(a,b)\in\{(r,U_r)\}$ and $(a,b')\in\{(r,U_r)\}$.
            Thus $(a,b) = (r,U_r)$ and $(a,b') = (r,U_r)$ by \cref{singleton_elim}.
            Thus $b = b'$ by \cref{pair_eq_iff}.
        \end{subproof}
        Thus $h$ is right-unique by \cref{rightunique}.
    \end{subproof}
    Show that $h$ is a relation.
    \begin{subproof}
        Show that for all $w\in\{(r,U_r)\}$ there exist $x, y$ such that $w = (x,y)$.
        \begin{subproof}
            Fix $w$.
            Assume $w\in\{(r,U_r)\}$.
            Thus $w = (r,U_r)$ by \cref{singleton_elim}.
            Thus there exist $x, y$ such that $w = (x,y)$.
        \end{subproof}
        Thus $h$ is a relation by \cref{relation}.
    \end{subproof}
    Thus $h$ is a function.
    Thus $f\union h$ is a function by \cref{union_compatible_functions}.
    Thus $g$ is a function.
\end{proof}

% from §33 helper lemma 27: extension for $1 < r$
\begin{lemma}\label{urysohn_unit_interval_extend_gt1}
    Assume $X$ is normal with respect to $T$.
    Assume $A$ is closed in $X$ with respect to $T$.
    Assume $B$ is closed in $X$ with respect to $T$.
    Assume $A$ is disjoint from $B$.
    Assume $s$ is a sequence in $\unitRationals$.
    Let $G = \urysohnGoodPairs{s}{X}{T}{A}{B}$.
    Let $R = \{(z,w) \mid z\in G, w\in G \mid
        \fst{w} = \fst{z} + 1 \land \snd{z}\subseteq\snd{w}\}$.
    Suppose $n\in\naturalsPlus$ and $f\in\urysohnAdmissible{s}{X}{T}{A}{B}{n}$ and $z = (n,f)$.
    Let $r = s(n + 1)$.
    Assume $r\notin\dom{f}$.
    Assume $1 < r$.
    Then there exists $w$ such that $z\mathrel{R} w$.
\end{lemma}
\begin{proof}
    Thus $f\in\funs{\{0,1\}\union \img{s}{\initialSegment{n}}}{\pow{X}}$ and
        for all $p\in\dom{f}$ we have $f(p)\in T$ and
        $f(1) = X\setminus B$ and
        $A\subseteq f(0)$ and
        $\closure{f(0)}{X}{T}\subseteq f(1)$ and
        for all $p\in\dom{f}$ for all $q\in\dom{f}$ if $p < q$ then
            $\closure{f(p)}{X}{T}\subseteq f(q)$ and
        for all $p\in\dom{f}$ if $1 < p$ then $f(p) = X$
        by \cref{urysohn_admissible}.
    Let $U_r = X$.
    Let $h = \{(r,U_r)\}$.
    Let $g = f\union h$.
    Thus $g$ is a function by \cref{urysohn_unit_interval_extend_gt1_setup}.
    Show that $\dom{g} = \{0,1\}\union \img{s}{\initialSegment{n + 1}}$.
    \begin{subproof}
        $\dom{g} = \dom{f}\union\dom{h}$ by \cref{dom_union}.
        Show that $\dom{h} = \{r\}$.
        \begin{subproof}
            Show that for all $x\in\dom{h}$ we have $x\in\{r\}$.
            \begin{subproof}
                Fix $x$.
                Assume $x\in\dom{h}$.
                Take $y$ such that $(x,y)\in h$ by \cref{dom_iff}.
                Thus $(x,y) = (r,U_r)$ by \cref{singleton_elim}.
                Thus $x = r$ by \cref{pair_eq_iff}.
                Thus $x\in\{r\}$ by \cref{singleton_intro}.
            \end{subproof}
            Thus $\dom{h}\subseteq \{r\}$ by \cref{subseteq}.
            $r\in\dom{h}$ by \cref{singleton_intro,dom_intro}.
            Show that for all $x\in\{r\}$ we have $x\in\dom{h}$.
            \begin{subproof}
                Fix $x$.
                Assume $x\in\{r\}$.
                Thus $x = r$ by \cref{singleton_elim}.
                Thus $x\in\dom{h}$.
            \end{subproof}
            Thus $\{r\}\subseteq\dom{h}$ by \cref{subseteq}.
            Thus $\dom{h} = \{r\}$ by \cref{subseteq_antisymmetric}.
        \end{subproof}
        Thus $\dom{g} = \dom{f}\union\{r\}$.
        $\dom{f} = \{0,1\}\union \img{s}{\initialSegment{n}}$ by \cref{funs_elim}.
        Show that $\dom{g}\subseteq \{0,1\}\union \img{s}{\initialSegment{n + 1}}$.
        \begin{subproof}
            Show that for all $y\in\dom{g}$ we have $y\in \{0,1\}\union \img{s}{\initialSegment{n + 1}}$.
            \begin{subproof}
                Fix $y$.
                Assume $y\in\dom{g}$.
                Thus $y\in\dom{f}\union\{r\}$.
                Then $y\in\dom{f}$ or $y\in\{r\}$ by \cref{union_iff}.
                \begin{byCase}
                \caseOf{$y\in\dom{f}$.}
                    Thus $y\in \{0,1\}\union \img{s}{\initialSegment{n}}$.
                    Then $y\in \{0,1\}$ or $y\in \img{s}{\initialSegment{n}}$ by \cref{union_iff}.
                    \begin{byCase}
                    \caseOf{$y\in \{0,1\}$.}
                        Thus $y\in \{0,1\}\union \img{s}{\initialSegment{n + 1}}$ by \cref{union_intro_left}.
                    \caseOf{$y\in \img{s}{\initialSegment{n}}$.}
                        Take $k$ such that $k\in\initialSegment{n}$ and $(k,y)\in s$ by \cref{img_iff}.
                        $k\in\naturalsPlus$ and $k \leq n$ by \cref{initial_segment_naturalsplus}.
                        $n\in\naturalsPlus$.
                        $n\in\reals$ and $k\in\reals$ by \cref{real_naturals_subset_reals,subseteq}.
                        $1\in\reals$ by \cref{real_naturals_subset_reals,real_one_in_naturals,subseteq}.
                        $0\in\reals$ by \cref{real_add_ident}.
                        $0 < 1$ by \cref{real_zero_lt_one}.
                        $0 + n < 1 + n$ by \cref{real_lt_add}.
                        $0 + n = n$ by \cref{real_add_ident}.
                        $1 + n = n + 1$ by \cref{real_add_comm}.
                        $n + 1\in\reals$ by \cref{real_add_closed}.
                        Thus $n < n + 1$ by \cref{real_lt_subst_left}.
                        Thus $n \leq n + 1$ by \cref{real_lt_imp_leq}.
                        Thus $k \leq n + 1$ by \cref{real_leq_trans}.
                        Thus $k\in\initialSegment{n + 1}$ by \cref{initial_segment_naturalsplus}.
                        Thus $y\in \img{s}{\initialSegment{n + 1}}$ by \cref{img_iff}.
                        Thus $y\in \{0,1\}\union \img{s}{\initialSegment{n + 1}}$ by \cref{union_intro_right}.
                    \end{byCase}
                \caseOf{$y\in\{r\}$.}
                    Thus $y = r$ by \cref{singleton_elim}.
                    Thus $y = s(n + 1)$.
                    $n + 1\in\naturalsPlus$ by \cref{naturalsplus_add_closed,real_one_in_naturals}.
                    Thus $n + 1 \leq n + 1$.
                    Thus $n + 1\in\initialSegment{n + 1}$ by \cref{initial_segment_naturalsplus}.
                    $s\in\funs{\naturalsPlus}{\unitRationals}$ by \cref{sequence_in}.
                    Thus $\dom{s} = \naturalsPlus$ by \cref{funs_elim}.
                    Thus $n + 1\in\dom{s}$.
                    $s$ is a function by \cref{funs_is_function}.
                    Thus $(n + 1,y)\in s$ by \cref{function_apply_elim}.
                    Thus $y\in \img{s}{\initialSegment{n + 1}}$ by \cref{img_iff}.
                    Thus $y\in \{0,1\}\union \img{s}{\initialSegment{n + 1}}$ by \cref{union_intro_right}.
                \end{byCase}
            \end{subproof}
            Thus $\dom{g}\subseteq \{0,1\}\union \img{s}{\initialSegment{n + 1}}$ by \cref{subseteq}.
        \end{subproof}
        Show that $\{0,1\}\union \img{s}{\initialSegment{n + 1}}\subseteq \dom{g}$.
        \begin{subproof}
            Show that for all $y\in \{0,1\}\union \img{s}{\initialSegment{n + 1}}$ we have $y\in\dom{g}$.
            \begin{subproof}
                Fix $y$.
                Assume $y\in \{0,1\}\union \img{s}{\initialSegment{n + 1}}$.
                Then $y\in \{0,1\}$ or $y\in \img{s}{\initialSegment{n + 1}}$ by \cref{union_iff}.
                \begin{byCase}
                \caseOf{$y\in \{0,1\}$.}
                    Thus $y\in \{0,1\}\union \img{s}{\initialSegment{n}}$ by \cref{union_intro_left}.
                    Thus $y\in\dom{f}$.
                    Thus $y\in\dom{g}$ by \cref{union_intro_left}.
                \caseOf{$y\in \img{s}{\initialSegment{n + 1}}$.}
                    Thus there exists $k$ such that $k\in\initialSegment{n + 1}$ and $(k,y)\in s$ by \cref{img_iff}.
                    Take $k$ such that $k\in\initialSegment{n + 1}$ and $(k,y)\in s$.
                    \begin{byCase}
                    \caseOf{$k\in\initialSegment{n}$.}
                        Thus $y\in \img{s}{\initialSegment{n}}$ by \cref{img_iff}.
                        Thus $y\in \{0,1\}\union \img{s}{\initialSegment{n}}$ by \cref{union_intro_right}.
                        Thus $y\in\dom{f}$.
                        Thus $y\in\dom{g}$ by \cref{union_intro_left}.
                    \caseOf{$k\notin\initialSegment{n}$.}
                        $k\in\naturalsPlus$ and $k\leq n + 1$ by \cref{initial_segment_naturalsplus}.
                        Show that $k = n + 1$.
                        \begin{subproof}
                            Suppose not.
                            Thus $k \neq n + 1$.
                            Thus $k \leq n$ by \cref{real_naturals_leq_succ_elim}.
                            Thus $k\in\initialSegment{n}$ by \cref{initial_segment_naturalsplus}.
                            Contradiction.
                        \end{subproof}
                        Thus $s\in\funs{\naturalsPlus}{\unitRationals}$ by \cref{sequence_in}.
                        Thus $s$ is a function by \cref{funs_is_function}.
                        Thus $s(k) = y$ by \cref{function_apply_intro}.
                        Thus $r = s(n + 1)$.
                        Thus $r = s(k)$.
                        Thus $r = y$.
                        Thus $y = r$.
                        Thus $y\in\{r\}$ by \cref{singleton_intro}.
                        Thus $y\in\dom{g}$ by \cref{union_intro_right}.
                    \end{byCase}
                \end{byCase}
            \end{subproof}
            Thus $\{0,1\}\union \img{s}{\initialSegment{n + 1}}\subseteq \dom{g}$ by \cref{subseteq}.
        \end{subproof}
        Thus $\dom{g} = \{0,1\}\union \img{s}{\initialSegment{n + 1}}$ by \cref{subseteq_antisymmetric}.
    \end{subproof}
    Show that for all $p\in\dom{g}$ we have $g(p)\in\pow{X}$.
    \begin{subproof}
        Fix $p$.
        Assume $p\in\dom{g}$.
        Thus $p\in\dom{f}\union\{r\}$.
        Then $p\in\dom{f}$ or $p\in\{r\}$ by \cref{union_iff}.
        \begin{byCase}
        \caseOf{$p\in\dom{f}$.}
            Thus $f\in\funs{\{0,1\}\union \img{s}{\initialSegment{n}}}{\pow{X}}$
                by \cref{urysohn_admissible}.
            Thus $\dom{f} = \{0,1\}\union \img{s}{\initialSegment{n}}$ by \cref{funs_elim}.
            Thus $p\in \{0,1\}\union \img{s}{\initialSegment{n}}$.
            Thus $f(p)\in\pow{X}$ by \cref{funs_type_apply}.
            Thus $f$ is a function by \cref{funs_is_function}.
            Thus $f$ is right-unique and $f$ is a relation.
            Thus $(p,f(p))\in f$ by \cref{function_apply_elim}.
            Thus $(p,f(p))\in g$ by \cref{union_intro_left}.
            Thus $g(p) = f(p)$ by \cref{function_apply_intro}.
            Thus $g(p)\in\pow{X}$.
        \caseOf{$p\in\{r\}$.}
            Thus $p = r$ by \cref{singleton_elim}.
            Thus $(r,U_r)\in h$ by \cref{singleton_intro}.
            Thus $(r,U_r)\in g$ by \cref{union_intro_right}.
            Thus $(p,U_r)\in g$.
            Thus $g(p) = U_r$ by \cref{function_apply_intro}.
            $T$ is a topology on $X$ by \cref{normal_space}.
            Thus $X\in T$ by \cref{topology_on}.
            Thus $U_r\in T$.
            Thus $T\subseteq\pow{X}$ by \cref{topology_on}.
            Thus $U_r\in\pow{X}$ by \cref{subseteq}.
            Thus $g(p)\in\pow{X}$.
        \end{byCase}
    \end{subproof}
    Let $w = (n + 1,g)$.
    Show that $w\in G$.
    \begin{subproof}
        Thus $n + 1\in\naturalsPlus$ by \cref{naturalsplus_add_closed,real_one_in_naturals}.
        Show that $g\in\urysohnAdmissible{s}{X}{T}{A}{B}{n + 1}$.
        \begin{subproof}
            Show that $g\in\funs{\{0,1\}\union \img{s}{\initialSegment{n + 1}}}{\pow{X}}$.
            \begin{subproof}
                Thus for all $p\in\dom{g}$ we have $g(p)\in\pow{X}$.
                Thus $g$ is a function by \cref{urysohn_unit_interval_extend_gt1_setup}.
                Thus $g$ is a function to $\pow{X}$.
                Thus $\dom{g} = \{0,1\}\union \img{s}{\initialSegment{n + 1}}$.
                Thus $g\in\funs{\{0,1\}\union \img{s}{\initialSegment{n + 1}}}{\pow{X}}$ by \cref{funs_intro}.
            \end{subproof}
            Show that for all $p\in\dom{g}$ we have
                $g(p)\in T$ and
                $g(1) = X\setminus B$ and
                $A\subseteq g(0)$ and
                $\closure{g(0)}{X}{T}\subseteq g(1)$ and
                for all $p\in\dom{g}$ for all $q\in\dom{g}$ if $p < q$ then
                    $\closure{g(p)}{X}{T}\subseteq g(q)$ and
                for all $p\in\dom{g}$ if $1 < p$ then $g(p) = X$.
            \begin{subproof}
                Fix $p$.
                Assume $p\in\dom{g}$.
                Show that $g(p)\in T$.
                \begin{subproof}
                    Thus $\dom{g} = \dom{f}\union\dom{h}$ by \cref{dom_union}.
                    Thus $p\in\dom{f}\union\dom{h}$.
                    Then $p\in\dom{f}$ or $p\in\dom{h}$ by \cref{union_iff}.
                    \begin{byCase}
                    \caseOf{$p\in\dom{f}$.}
                        Thus $f(p)\in T$ by \cref{urysohn_admissible}.
                        Thus $f$ is a function by \cref{funs_is_function}.
                        Thus $(p,f(p))\in f$ by \cref{function_apply_elim}.
                        Thus $(p,f(p))\in g$ by \cref{union_intro_left}.
                        Thus $g(p) = f(p)$ by \cref{function_apply_intro}.
                        Thus $g(p)\in T$.
                    \caseOf{$p\in\dom{h}$.}
                        Take $y$ such that $(p,y)\in h$ by \cref{dom_iff}.
                        Thus $(p,y) = (r,U_r)$ by \cref{singleton_elim}.
                        Thus $p = r$ by \cref{pair_eq_iff}.
                        Thus $(r,U_r)\in g$ by \cref{union_intro_right,singleton_intro}.
                        Thus $g(p) = U_r$ by \cref{function_apply_intro}.
                        $T$ is a topology on $X$ by \cref{normal_space}.
                        Thus $X\in T$ by \cref{topology_on}.
                        Thus $U_r\in T$.
                        Thus $g(p)\in T$.
                    \end{byCase}
                \end{subproof}
                Thus $T$ is a topology on $X$ by \cref{normal_space}.
                Thus $g(p)$ is open in $X$ with respect to $T$ by \cref{open_in}.
                Show that $g(1) = X\setminus B$.
                \begin{subproof}
                    Thus $f\in\funs{\{0,1\}\union \img{s}{\initialSegment{n}}}{\pow{X}}$
                        by \cref{urysohn_admissible}.
                    Thus $1\in\dom{f}$ by \cref{funs_elim,upair_intro_right,union_intro_left}.
                    $f$ is a function by \cref{funs_is_function}.
                    Thus $(1,f(1))\in f$ by \cref{function_apply_elim}.
                    Thus $(1,f(1))\in g$ by \cref{union_intro_left}.
                    Thus $g(1) = f(1)$ by \cref{function_apply_intro}.
                    Thus $f(1) = X\setminus B$.
                    Thus $g(1) = X\setminus B$.
                \end{subproof}
                Show that $A\subseteq g(0)$.
                \begin{subproof}
                    Thus $f\in\funs{\{0,1\}\union \img{s}{\initialSegment{n}}}{\pow{X}}$
                        by \cref{urysohn_admissible}.
                    Thus $0\in\dom{f}$ by \cref{funs_elim,upair_intro_left,union_intro_left}.
                    $f$ is a function by \cref{funs_is_function}.
                    Thus $(0,f(0))\in f$ by \cref{function_apply_elim}.
                    Thus $(0,f(0))\in g$ by \cref{union_intro_left}.
                    Thus $g(0) = f(0)$ by \cref{function_apply_intro}.
                    Thus $A\subseteq f(0)$.
                    Thus $A\subseteq g(0)$.
                \end{subproof}
                Show that $\closure{g(0)}{X}{T}\subseteq g(1)$.
                \begin{subproof}
                    Thus $f\in\funs{\{0,1\}\union \img{s}{\initialSegment{n}}}{\pow{X}}$
                        by \cref{urysohn_admissible}.
                    Thus $0\in\dom{f}$ by \cref{funs_elim,upair_intro_left,union_intro_left}.
                    $f$ is a function by \cref{funs_is_function}.
                    Thus $(0,f(0))\in f$ by \cref{function_apply_elim}.
                    Thus $(0,f(0))\in g$ by \cref{union_intro_left}.
                    Thus $g(0) = f(0)$ by \cref{function_apply_intro}.
                    Thus $1\in\dom{f}$ by \cref{funs_elim,upair_intro_right,union_intro_left}.
                    Thus $(1,f(1))\in f$ by \cref{function_apply_elim}.
                    Thus $(1,f(1))\in g$ by \cref{union_intro_left}.
                    Thus $g(1) = f(1)$ by \cref{function_apply_intro}.
                    Thus $\closure{f(0)}{X}{T}\subseteq f(1)$.
                    Thus $\closure{g(0)}{X}{T}\subseteq g(1)$ by \cref{subseteq}.
                \end{subproof}
                Show that for all $p\in\dom{g}$ for all $q\in\dom{g}$ if $p < q$ then
                    $\closure{g(p)}{X}{T}\subseteq g(q)$.
                \begin{subproof}
                    Fix $p$.
                    Assume $p\in\dom{g}$.
                    Fix $q$.
                    Assume $q\in\dom{g}$.
                    Assume $p < q$.
                    Thus $\dom{g} = \dom{f}\union\dom{h}$ by \cref{dom_union}.
                    Then $p\in\dom{f}$ or $p\in\dom{h}$ by \cref{union_iff}.
                    Then $q\in\dom{f}$ or $q\in\dom{h}$ by \cref{union_iff}.
                    \begin{byCase}
                    \caseOf{$q\in\dom{h}$.}
                        Take $y$ such that $(q,y)\in h$ by \cref{dom_iff}.
                        Thus $(q,y) = (r,U_r)$ by \cref{singleton_elim}.
                        Thus $q = r$ by \cref{pair_eq_iff}.
                        Thus $(r,U_r)\in g$ by \cref{union_intro_right,singleton_intro}.
                        Thus $g(q) = U_r$ by \cref{function_apply_intro}.
                        Thus $g(q) = X$.
                        $X$ is closed in $X$ with respect to $T$ by \cref{closed_whole_space}.
                        Thus $g\in\funs{\{0,1\}\union \img{s}{\initialSegment{n + 1}}}{\pow{X}}$.
                        Thus $\dom{g} = \{0,1\}\union \img{s}{\initialSegment{n + 1}}$ by \cref{funs_elim}.
                        Thus $p\in \{0,1\}\union \img{s}{\initialSegment{n + 1}}$.
                        Thus $g(p)\in\pow{X}$ by \cref{funs_type_apply}.
                        Thus $g(p)\subseteq X$ by \cref{powerset_elim}.
                        Thus $\closure{g(p)}{X}{T}\subseteq X$ by \cref{closure_subseteq_closed}.
                        Thus $\closure{g(p)}{X}{T}\subseteq g(q)$ by \cref{subseteq}.
                    \caseOf{$q\in\dom{f}$.}
                        \begin{byCase}
                        \caseOf{$p\in\dom{f}$.}
                            Thus $\closure{f(p)}{X}{T}\subseteq f(q)$ by \cref{urysohn_admissible}.
                            Thus $f$ is a function by \cref{funs_is_function}.
                            Thus $(p,f(p))\in f$ and $(q,f(q))\in f$ by \cref{function_apply_elim}.
                            Thus $(p,f(p))\in g$ and $(q,f(q))\in g$ by \cref{union_intro_left}.
                            Thus $g(p) = f(p)$ and $g(q) = f(q)$ by \cref{function_apply_intro}.
                            Thus $\closure{g(p)}{X}{T}\subseteq g(q)$ by \cref{subseteq}.
                        \caseOf{$p\in\dom{h}$.}
                            Take $y$ such that $(p,y)\in h$ by \cref{dom_iff}.
                            Thus $(p,y) = (r,U_r)$ by \cref{singleton_elim}.
                            Thus $p = r$ by \cref{pair_eq_iff}.
                            Thus $r < q$ by \cref{real_lt_subst_left}.
                            Show that $r\in\reals$.
                            \begin{subproof}
                                Thus $s\in\funs{\naturalsPlus}{\unitRationals}$ by \cref{sequence_in}.
                                Thus $n + 1\in\naturalsPlus$ by \cref{naturalsplus_add_closed,real_one_in_naturals}.
                                Thus $r\in\unitRationals$ by \cref{funs_type_apply}.
                                Take $t$ such that
                                    $t\in\naturalsPlus\times\naturalsPlus$ and
                                    $r = \fst{t} \rmul \inv{\snd{t}}$
                                    by \cref{unit_rationals}.
                                Thus $\fst{t}\in\naturalsPlus$ and $\snd{t}\in\naturalsPlus$
                                    by \cref{fst_type,snd_type}.
                                Thus $\fst{t}\in\reals$ and $\snd{t}\in\reals$
                                    by \cref{real_naturals_subset_reals,subseteq}.
                                Show that $\snd{t}\neq 0$.
                                \begin{subproof}
                                    $1\in\reals$ by \cref{real_mul_ident}.
                                    Thus $1 < \snd{t}$ or $1 = \snd{t}$
                                        by \cref{real_one_leq_naturalsplus}.
                                    $0\in\reals$ by \cref{real_add_ident}.
                                    $0 < 1$ by \cref{real_zero_lt_one}.
                                    \begin{byCase}
                                    \caseOf{$1 < \snd{t}$.}
                                        Thus $0 < \snd{t}$ by \cref{real_lt_trans}.
                                        Thus $\snd{t}\neq 0$ by \cref{real_pos_nonzero}.
                                    \caseOf{$1 = \snd{t}$.}
                                        Thus $0 < \snd{t}$ by \cref{real_lt_subst_right}.
                                        Thus $\snd{t}\neq 0$ by \cref{real_pos_nonzero}.
                                    \end{byCase}
                                \end{subproof}
                                Thus $\inv{\snd{t}}\in\reals$ by \cref{real_mul_inverse}.
                                Thus $r\in\reals$ by \cref{real_mul_closed}.
                            \end{subproof}
                            Show that $q\in\reals$.
                            \begin{subproof}
                                Thus $f\in\funs{\{0,1\}\union \img{s}{\initialSegment{n}}}{\pow{X}}$
                                    by \cref{urysohn_admissible}.
                                Thus $\dom{f} = \{0,1\}\union \img{s}{\initialSegment{n}}$ by \cref{funs_elim}.
                                Thus $q\in \{0,1\}\union \img{s}{\initialSegment{n}}$.
                                Then $q\in \{0,1\}$ or $q\in \img{s}{\initialSegment{n}}$ by \cref{union_iff}.
                                \begin{byCase}
                                \caseOf{$q\in \{0,1\}$.}
                                    Thus $q = 0$ or $q = 1$ by \cref{upair_elim}.
                                    \begin{byCase}
                                    \caseOf{$q = 0$.}
                                        $0\in\reals$ by \cref{real_add_ident}.
                                        Thus $q\in\reals$.
                                    \caseOf{$q = 1$.}
                                        $1\in\reals$ by \cref{real_mul_ident}.
                                        Thus $q\in\reals$.
                                    \end{byCase}
                                \caseOf{$q\in \img{s}{\initialSegment{n}}$.}
                                    Take $k$ such that $k\in\initialSegment{n}$ and $(k,q)\in s$ by \cref{img_iff}.
                                    Thus $s\in\funs{\naturalsPlus}{\unitRationals}$ by \cref{sequence_in}.
                                    Thus $k\in\naturalsPlus$ by \cref{initial_segment_naturalsplus}.
                                    Thus $s$ is a function by \cref{funs_is_function}.
                                    Thus $s(k) = q$ by \cref{function_apply_intro}.
                                    Thus $s(k)\in\unitRationals$ by \cref{funs_type_apply}.
                                    Thus $q\in\unitRationals$.
                                    Take $t$ such that
                                        $t\in\naturalsPlus\times\naturalsPlus$ and
                                        $q = \fst{t} \rmul \inv{\snd{t}}$
                                        by \cref{unit_rationals}.
                                    Thus $\fst{t}\in\naturalsPlus$ and $\snd{t}\in\naturalsPlus$
                                        by \cref{fst_type,snd_type}.
                                    Thus $\fst{t}\in\reals$ and $\snd{t}\in\reals$
                                        by \cref{real_naturals_subset_reals,subseteq}.
                                    Show that $\snd{t}\neq 0$.
                                    \begin{subproof}
                                        $1\in\reals$ by \cref{real_mul_ident}.
                                        Thus $1 < \snd{t}$ or $1 = \snd{t}$
                                            by \cref{real_one_leq_naturalsplus}.
                                        $0\in\reals$ by \cref{real_add_ident}.
                                        $0 < 1$ by \cref{real_zero_lt_one}.
                                        \begin{byCase}
                                        \caseOf{$1 < \snd{t}$.}
                                            Thus $0 < \snd{t}$ by \cref{real_lt_trans}.
                                            Thus $\snd{t}\neq 0$ by \cref{real_pos_nonzero}.
                                        \caseOf{$1 = \snd{t}$.}
                                            Thus $0 < \snd{t}$ by \cref{real_lt_subst_right}.
                                            Thus $\snd{t}\neq 0$ by \cref{real_pos_nonzero}.
                                        \end{byCase}
                                    \end{subproof}
                                    Thus $\inv{\snd{t}}\in\reals$ by \cref{real_mul_inverse}.
                                    Thus $q\in\reals$ by \cref{real_mul_closed}.
                                \end{byCase}
                            \end{subproof}
                            $1\in\reals$ by \cref{real_mul_ident}.
                            Thus $1 < q$ by \cref{real_lt_trans}.
                            Thus $q\in\dom{f}$.
                            Thus $f(q) = X$ by \cref{urysohn_admissible_above_one}.
                            Thus $f$ is a function by \cref{funs_is_function}.
                            Thus $(q,f(q))\in f$ by \cref{function_apply_elim}.
                            Thus $(q,f(q))\in g$ by \cref{union_intro_left}.
                            Thus $g(q) = f(q)$ by \cref{function_apply_intro}.
                            Thus $(r,U_r)\in g$ by \cref{union_intro_right,singleton_intro}.
                            Thus $g(p) = U_r$ by \cref{function_apply_intro}.
                            Thus $g(p) = X$.
                            $X$ is closed in $X$ with respect to $T$ by \cref{closed_whole_space}.
                            Thus $g\in\funs{\{0,1\}\union \img{s}{\initialSegment{n + 1}}}{\pow{X}}$.
                            Thus $\dom{g} = \{0,1\}\union \img{s}{\initialSegment{n + 1}}$ by \cref{funs_elim}.
                            Thus $p\in \{0,1\}\union \img{s}{\initialSegment{n + 1}}$.
                            Thus $g(p)\in\pow{X}$ by \cref{funs_type_apply}.
                            Thus $g(p)\subseteq X$ by \cref{powerset_elim}.
                            Thus $\closure{g(p)}{X}{T}\subseteq X$ by \cref{closure_subseteq_closed}.
                            Thus $\closure{g(p)}{X}{T}\subseteq g(q)$ by \cref{subseteq}.
                        \end{byCase}
                    \end{byCase}
                \end{subproof}
                Show that for all $p\in\dom{g}$ if $1 < p$ then $g(p) = X$.
                \begin{subproof}
                    Fix $p$.
                    Assume $p\in\dom{g}$.
                    Assume $1 < p$.
                    \begin{byCase}
                    \caseOf{$p = r$.}
                        Thus $(r,U_r)\in g$ by \cref{union_intro_right,singleton_intro}.
                        Thus $g(p) = U_r$ by \cref{function_apply_intro}.
                        Thus $g(p) = X$.
                    \caseOf{$p\neq r$.}
                        Show that $p\in\dom{f}$.
                        \begin{subproof}
                            Take $y$ such that $(p,y)\in g$ by \cref{dom_iff}.
                            Thus $(p,y)\in f$ or $(p,y)\in h$ by \cref{union_iff}.
                            Show that $(p,y)\in f$.
                            \begin{subproof}
                                Suppose not.
                                Thus $(p,y)\notin f$.
                                Thus $(p,y)\in h$ by \cref{union_iff}.
                                Thus $(p,y) = (r,U_r)$ by \cref{singleton_elim}.
                                Thus $p = r$ by \cref{pair_eq_iff}.
                                Contradiction.
                            \end{subproof}
                            Thus $p\in\dom{f}$ by \cref{dom_intro}.
                        \end{subproof}
                        $f$ is a function by \cref{funs_is_function}.
                        Thus $(p,f(p))\in f$ by \cref{function_apply_elim}.
                        Thus $(p,f(p))\in g$ by \cref{union_intro_left}.
                        Thus $g(p) = f(p)$ by \cref{function_apply_intro}.
                        Show that $f(p) = X$.
                        \begin{subproof}
                            Show that $0\in\dom{f}$.
                            \begin{subproof}
                                $\dom{f} = \{0,1\}\union \img{s}{\initialSegment{n}}$ by \cref{funs_elim}.
                                $0\in\{0,1\}$ by \cref{upair_intro_left}.
                                Thus $0\in\dom{f}$.
                            \end{subproof}
                            Thus $f(0)\in T$ and $f(1) = X\setminus B$ and $A\subseteq f(0)$ and
                                $\closure{f(0)}{X}{T}\subseteq f(1)$ and
                                for all $p\in\dom{f}$ for all $q\in\dom{f}$ if $p < q$ then
                                    $\closure{f(p)}{X}{T}\subseteq f(q)$ and
                                for all $p\in\dom{f}$ if $1 < p$ then $f(p) = X$
                                by \cref{urysohn_admissible}.
                            Thus if $1 < p$ then $f(p) = X$.
                            Thus $f(p) = X$.
                        \end{subproof}
                        Thus $g(p) = X$.
                    \end{byCase}
                \end{subproof}
            \end{subproof}
            Thus $g\in\urysohnAdmissible{s}{X}{T}{A}{B}{n + 1}$ by \cref{urysohn_admissible}.
        \end{subproof}
        Thus $(n + 1,g)\in G$ by \cref{urysohn_good_pair_intro}.
        Thus $w\in G$.
    \end{subproof}
    Show that $z\mathrel{R} w$.
    \begin{subproof}
        Thus $z\in G$.
        Thus $w\in G$.
        Thus $\fst{z} = n$ by \cref{fst_eq}.
        Thus $\fst{w} = n + 1$ by \cref{fst_eq}.
        Thus $\fst{w} = \fst{z} + 1$.
        Thus $\snd{z} = f$ by \cref{snd_eq}.
        Thus $\snd{w} = g$ by \cref{snd_eq}.
        $f\subseteq g$ by \cref{union_upper_left}.
        Thus $\snd{z}\subseteq \snd{w}$.
        Thus $(z,w)\in R$.
    \end{subproof}
    Thus $z\mathrel{R} w$.
\end{proof}

% from §33 helper lemma 28: extension for $r < 1$ (setup)
\begin{lemma}\label{urysohn_unit_interval_extend_lt1_setup}
    Assume $X$ is normal with respect to $T$.
    Assume $A$ is closed in $X$ with respect to $T$.
    Assume $B$ is closed in $X$ with respect to $T$.
    Assume $A$ is disjoint from $B$.
    Assume $s$ is a sequence in $\unitRationals$.
    Let $G = \urysohnGoodPairs{s}{X}{T}{A}{B}$.
    Let $R = \{(z,w) \mid z\in G, w\in G \mid
        \fst{w} = \fst{z} + 1 \land \snd{z}\subseteq\snd{w}\}$.
    Suppose $n\in\naturalsPlus$ and $f\in\urysohnAdmissible{s}{X}{T}{A}{B}{n}$ and $z = (n,f)$.
    Let $r = s(n + 1)$.
    Assume $r\notin\dom{f}$.
    Assume $r < 1$.
    Let $S = \{p\in\dom{f} \mid p < r\}$.
    Let $T_1 = \{q\in\dom{f} \mid r < q\}$.
    Then there exist $p_0,q_0,U_r,h,g$ such that
        $p_0$ is a largest element of $S$ with respect to $\realOrderRel$ and
        $q_0$ is a smallest element of $T_1$ with respect to $\realOrderRel$ and
        $U_r$ is open in $X$ with respect to $T$ and
        $\closure{f(p_0)}{X}{T}\subseteq U_r$ and
        $\closure{U_r}{X}{T}\subseteq f(q_0)$ and
        $h = \{(r,U_r)\}$ and
        $g = f\union h$ and
        $g$ is a function.
\end{lemma}
\begin{proof}
    Thus $f\in\funs{\{0,1\}\union \img{s}{\initialSegment{n}}}{\pow{X}}$ and
        for all $p\in\dom{f}$ we have $f(p)\in T$ and
        $f(1) = X\setminus B$ and
        $A\subseteq f(0)$ and
        $\closure{f(0)}{X}{T}\subseteq f(1)$ and
        for all $p\in\dom{f}$ for all $q\in\dom{f}$ if $p < q$ then
            $\closure{f(p)}{X}{T}\subseteq f(q)$ and
        for all $p\in\dom{f}$ if $1 < p$ then $f(p) = X$
        by \cref{urysohn_admissible}.
    Show that $S$ is inhabited and $T_1$ is inhabited.
    \begin{subproof}
        $\dom{f} = \{0,1\}\union \img{s}{\initialSegment{n}}$ by \cref{funs_elim}.
        Show that $0\in\dom{f}$.
        \begin{subproof}
            $0\in\{0,1\}$ by \cref{upair_intro_left}.
            Thus $0\in\dom{f}$.
        \end{subproof}
        Show that $1\in\dom{f}$.
        \begin{subproof}
            $1\in\{0,1\}$ by \cref{upair_intro_right}.
            Thus $1\in\dom{f}$.
        \end{subproof}
        Thus $s\in\funs{\naturalsPlus}{\unitRationals}$ by \cref{sequence_in}.
        Thus $n + 1\in\naturalsPlus$ by \cref{naturalsplus_add_closed,real_one_in_naturals}.
        Thus $r\in\unitRationals$ by \cref{funs_type_apply}.
        Thus $0 < r$ by \cref{unit_rational_pos}.
        Thus $0\in S$.
        Thus $r < 1$.
        Thus $1\in T_1$.
    \end{subproof}
    Show that $S$ is finite and $T_1$ is finite.
    \begin{subproof}
        Let $D = \{0,1\}\union \img{s}{\initialSegment{n}}$.
        $\dom{f} = D$ by \cref{funs_elim}.
        $\{0,1\}$ is finite by \cref{pair_is_finite}.
        Show that $\img{s}{\initialSegment{n}}$ is finite.
        \begin{subproof}
            $\initialSegment{n}$ is finite by \cref{initial_segment_finite}.
            $s\in\funs{\naturalsPlus}{\unitRationals}$ by \cref{sequence_in}.
            $s$ is a function by \cref{funs_is_function}.
            Let $g_0 = \restrl{s}{\initialSegment{n}}$.
            $g_0$ is a function by \cref{restrl_of_function_is_function}.
            $\dom{s} = \naturalsPlus$ by \cref{funs_elim}.
            $\dom{g_0} = \dom{s}\inter \initialSegment{n}$ by \cref{dom_of_restrl_eq_inter}.
            Show that $\dom{g_0} = \initialSegment{n}$.
            \begin{subproof}
                Thus $\dom{g_0}\subseteq\initialSegment{n}$ by \cref{restrl_dom}.
                Show that for all $x$ such that $x\in\initialSegment{n}$ we have $x\in\dom{g_0}$.
                \begin{subproof}
                    Fix $x$.
                    Assume $x\in\initialSegment{n}$.
                    Thus $x\in\naturalsPlus$ by \cref{initial_segment_naturalsplus}.
                    Thus $x\in\dom{s}$.
                    Thus $x\in\dom{s}\inter \initialSegment{n}$ by \cref{inter_intro}.
                    Thus $x\in\dom{g_0}$.
                \end{subproof}
                Thus $\initialSegment{n}\subseteq\dom{g_0}$ by \cref{subseteq}.
                Thus $\dom{g_0} = \initialSegment{n}$ by \cref{subseteq_antisymmetric}.
            \end{subproof}
            Thus $g_0$ is a function on $\initialSegment{n}$.
            Thus $\img{g_0}{\initialSegment{n}}$ is finite by \cref{finite_image_function}.
            Show that $\initialSegment{n}\subseteq\dom{s}$.
            \begin{subproof}
                Show that for all $x$ such that $x\in\initialSegment{n}$ we have $x\in\dom{s}$.
                \begin{subproof}
                    Fix $x$.
                    Assume $x\in\initialSegment{n}$.
                    Thus $x\in\naturalsPlus$ by \cref{initial_segment_naturalsplus}.
                    Thus $x\in\dom{s}$.
                \end{subproof}
                Thus $\initialSegment{n}\subseteq\dom{s}$ by \cref{subseteq}.
            \end{subproof}
            Thus $\img{g_0}{\initialSegment{n}} = \img{s}{\initialSegment{n}\inter \initialSegment{n}}$
                by \cref{restrl_img}.
            Thus $\img{g_0}{\initialSegment{n}} = \img{s}{\initialSegment{n}}$
                by \cref{inter_idempotent}.
            Thus $\img{s}{\initialSegment{n}}$ is finite.
        \end{subproof}
        Thus $D$ is finite by \cref{union_of_finite_is_finite}.
        Thus $\dom{f}$ is finite.
        Show that $S\subseteq\dom{f}$.
        \begin{subproof}
            Show that for all $x$ such that $x\in S$ we have $x\in\dom{f}$.
            \begin{subproof}
                Fix $x$.
                Assume $x\in S$.
                Thus $x\in\dom{f}$.
            \end{subproof}
            Thus $S\subseteq\dom{f}$ by \cref{subseteq}.
        \end{subproof}
        Show that $T_1\subseteq\dom{f}$.
        \begin{subproof}
            Show that for all $x$ such that $x\in T_1$ we have $x\in\dom{f}$.
            \begin{subproof}
                Fix $x$.
                Assume $x\in T_1$.
                Thus $x\in\dom{f}$.
            \end{subproof}
            Thus $T_1\subseteq\dom{f}$ by \cref{subseteq}.
        \end{subproof}
        Thus $S$ is finite by \cref{subseteq_finite_is_finite}.
        Thus $T_1$ is finite by \cref{subseteq_finite_is_finite}.
    \end{subproof}
    Show that $S\subseteq\reals$ and $T_1\subseteq\reals$.
    \begin{subproof}
        Show that $\dom{f}\subseteq\reals$.
        \begin{subproof}
            $\dom{f} = \{0,1\}\union \img{s}{\initialSegment{n}}$ by \cref{funs_elim}.
            Show that $\{0,1\}\subseteq\reals$.
            \begin{subproof}
                Show that for all $x$ such that $x\in\{0,1\}$ we have $x\in\reals$.
                \begin{subproof}
                    Fix $x$.
                    Assume $x\in\{0,1\}$.
                    $x = 0$ or $x = 1$ by \cref{upair_elim}.
                    \begin{byCase}
                    \caseOf{$x = 0$.}
                        $0\in\reals$ by \cref{real_add_ident}.
                        Thus $x\in\reals$.
                    \caseOf{$x = 1$.}
                        $1\in\reals$ by \cref{real_naturals_subset_reals,real_one_in_naturals,subseteq}.
                        Thus $x\in\reals$.
                    \end{byCase}
                \end{subproof}
                Thus $\{0,1\}\subseteq\reals$ by \cref{subseteq}.
            \end{subproof}
            Show that $\img{s}{\initialSegment{n}}\subseteq\reals$.
            \begin{subproof}
                Show that for all $y$ such that $y\in\img{s}{\initialSegment{n}}$ we have $y\in\reals$.
                \begin{subproof}
                    Fix $y$.
                    Assume $y\in\img{s}{\initialSegment{n}}$.
                    Take $k$ such that $k\in\initialSegment{n}$ and $(k,y)\in s$ by \cref{img_iff}.
                    $k\in\naturalsPlus$ by \cref{initial_segment_naturalsplus}.
                    $s\in\funs{\naturalsPlus}{\unitRationals}$ by \cref{sequence_in}.
                    $s$ is a function by \cref{funs_is_function}.
                    Thus $s(k) = y$ by \cref{function_apply_intro}.
                    Thus $s(k)\in\unitRationals$ by \cref{funs_type_apply}.
                    Thus $y\in\unitRationals$.
                    Show that $y\in\reals$.
                    \begin{subproof}
                        Take $q$ such that $q\in\naturalsPlus\times\naturalsPlus$ and
                            $y = \fst{q} \rmul \inv{\snd{q}}$ by \cref{unit_rationals}.
                        Thus $\fst{q}\in\naturalsPlus$ by \cref{fst_type}.
                        Thus $\snd{q}\in\naturalsPlus$ by \cref{snd_type}.
                        $\fst{q}\in\reals$ and $\snd{q}\in\reals$
                            by \cref{real_naturals_subset_reals,subseteq}.
                        $1 < \snd{q}$ or $1 = \snd{q}$ by \cref{real_one_leq_naturalsplus}.
                        $0 < 1$ by \cref{real_zero_lt_one}.
                        Show that $0 < \snd{q}$.
                        \begin{subproof}
                            \begin{byCase}
                            \caseOf{$1 = \snd{q}$.}
                                Thus $0 < \snd{q}$ by \cref{real_zero_lt_one}.
                            \caseOf{$1 < \snd{q}$.}
                                $0\in\reals$ by \cref{real_add_ident}.
                                $1\in\reals$ by \cref{real_naturals_subset_reals,real_one_in_naturals,subseteq}.
                                Thus $0 < \snd{q}$ by \cref{real_lt_trans}.
                            \end{byCase}
                        \end{subproof}
                        Show that $\snd{q}\neq 0$.
                        \begin{subproof}
                            Suppose not.
                            Thus $\snd{q} = 0$.
                            Thus $0 < 0$.
                            Contradiction by \cref{real_lt_irreflexive}.
                        \end{subproof}
                        Thus $\inv{\snd{q}}\in\reals$ and
                            $\snd{q} \rmul \inv{\snd{q}} = 1$ by \cref{real_mul_inverse}.
                        Thus $y\in\reals$ by \cref{real_mul_closed}.
                    \end{subproof}
                \end{subproof}
                Thus $\img{s}{\initialSegment{n}}\subseteq\reals$ by \cref{subseteq}.
            \end{subproof}
            Thus $\dom{f}\subseteq\reals$ by \cref{union_subsets_is_subset}.
        \end{subproof}
        Show that $S\subseteq\dom{f}$.
        \begin{subproof}
            Show that for all $x$ such that $x\in S$ we have $x\in\dom{f}$.
            \begin{subproof}
                Fix $x$.
                Assume $x\in S$.
                Thus $x\in\dom{f}$.
            \end{subproof}
            Thus $S\subseteq\dom{f}$ by \cref{subseteq}.
        \end{subproof}
        Show that $T_1\subseteq\dom{f}$.
        \begin{subproof}
            Show that for all $x$ such that $x\in T_1$ we have $x\in\dom{f}$.
            \begin{subproof}
                Fix $x$.
                Assume $x\in T_1$.
                Thus $x\in\dom{f}$.
            \end{subproof}
            Thus $T_1\subseteq\dom{f}$ by \cref{subseteq}.
        \end{subproof}
        Thus $S\subseteq\reals$ by \cref{subseteq_transitive}.
        Thus $T_1\subseteq\reals$ by \cref{subseteq_transitive}.
    \end{subproof}
    $\realOrderRel$ is a simple order relation on $\reals$ by \cref{real_order_relation_simple}.
    Take $p_0$ such that $p_0$ is a largest element of $S$ with respect to $\realOrderRel$
        by \cref{finite_inhabited_largest_element}.
    Take $q_0$ such that $q_0$ is a smallest element of $T_1$ with respect to $\realOrderRel$
        by \cref{finite_inhabited_smallest_element}.
    Thus $p_0 < r$ and $r < q_0$ by \cref{largest_element,smallest_element}.
    Show that $\closure{f(p_0)}{X}{T}\subseteq f(q_0)$.
    \begin{subproof}
        Thus $p_0\in S$ and $q_0\in T_1$ by \cref{largest_element,smallest_element}.
        Thus $p_0\in\dom{f}$ and $p_0 < r$.
        Thus $q_0\in\dom{f}$ and $r < q_0$.
        Thus $p_0\in\reals$ by \cref{subseteq}.
        Thus $q_0\in\reals$ by \cref{subseteq}.
        Show that $r\in\reals$.
        \begin{subproof}
            Thus $s\in\funs{\naturalsPlus}{\unitRationals}$ by \cref{sequence_in}.
            Thus $n + 1\in\naturalsPlus$ by \cref{naturalsplus_add_closed,real_one_in_naturals}.
            Thus $r\in\unitRationals$ by \cref{funs_type_apply}.
            Take $t$ such that
                $t\in\naturalsPlus\times\naturalsPlus$ and
                $r = \fst{t} \rmul \inv{\snd{t}}$
                by \cref{unit_rationals}.
            Thus $\fst{t}\in\naturalsPlus$ and $\snd{t}\in\naturalsPlus$
                by \cref{fst_type,snd_type}.
            Thus $\fst{t}\in\reals$ and $\snd{t}\in\reals$
                by \cref{real_naturals_subset_reals,subseteq}.
            Show that $\snd{t}\neq 0$.
            \begin{subproof}
                $1\in\reals$ by \cref{real_mul_ident}.
                Thus $1 < \snd{t}$ or $1 = \snd{t}$
                    by \cref{real_one_leq_naturalsplus}.
                $0\in\reals$ by \cref{real_add_ident}.
                $0 < 1$ by \cref{real_zero_lt_one}.
                \begin{byCase}
                \caseOf{$1 < \snd{t}$.}
                    Thus $0 < \snd{t}$ by \cref{real_lt_trans}.
                    Thus $\snd{t}\neq 0$ by \cref{real_pos_nonzero}.
                \caseOf{$1 = \snd{t}$.}
                    Thus $0 < \snd{t}$ by \cref{real_lt_subst_right}.
                    Thus $\snd{t}\neq 0$ by \cref{real_pos_nonzero}.
                \end{byCase}
            \end{subproof}
            Thus $\inv{\snd{t}}\in\reals$ by \cref{real_mul_inverse}.
            Thus $r\in\reals$ by \cref{real_mul_closed}.
        \end{subproof}
        Thus $p_0 < q_0$ by \cref{real_lt_trans}.
        Thus $\closure{f(p_0)}{X}{T}\subseteq f(q_0)$ by \cref{urysohn_admissible}.
    \end{subproof}
    Show that $\closure{f(p_0)}{X}{T}$ is closed in $X$ with respect to $T$.
    \begin{subproof}
        $T$ is a topology on $X$ by \cref{normal_space}.
        Thus $p_0\in\dom{f}$.
        Thus $f(p_0)\in T$ by \cref{urysohn_admissible}.
        Thus $T\subseteq\pow{X}$ by \cref{topology_on}.
        Thus $f(p_0)\in\pow{X}$ by \cref{elem_subseteq}.
        Thus $f(p_0)\subseteq X$ by \cref{powerset_elim}.
        Thus $\closure{f(p_0)}{X}{T}$ is closed in $X$ with respect to $T$
            by \cref{closure_closed_in}.
    \end{subproof}
    Show that $f(q_0)$ is open in $X$ with respect to $T$.
    \begin{subproof}
        $T$ is a topology on $X$ by \cref{normal_space}.
        Thus $q_0\in\dom{f}$.
        Thus $f(q_0)\in T$ by \cref{urysohn_admissible}.
        Thus $f(q_0)$ is open in $X$ with respect to $T$ by \cref{open_in}.
    \end{subproof}
    Take $U_r$ such that
        $U_r$ is open in $X$ with respect to $T$ and
        $\closure{f(p_0)}{X}{T}\subseteq U_r$ and
        $\closure{U_r}{X}{T}\subseteq f(q_0)$
        by \cref{normal_closure_neighborhood}.
    Let $h = \{(r,U_r)\}$.
    Let $g = f\union h$.
    Show that $g$ is a function.
    \begin{subproof}
        Show that for all $x$ such that $x\in\dom{f}\inter\dom{h}$ we have $f(x) = h(x)$.
        \begin{subproof}
            Suppose not.
            Then there exists $x$ such that
                $x\in\dom{f}\inter\dom{h}$ and
                $f(x) \neq h(x)$.
            Take $x$ such that $x\in\dom{f}\inter\dom{h}$ and $f(x) \neq h(x)$.
            $x\in\dom{h}$ by \cref{inter_elim_right}.
            Take $y$ such that $(x,y)\in h$ by \cref{dom_iff}.
            Thus $(x,y) = (r,U_r)$ by \cref{singleton_elim}.
            Thus $x = r$ by \cref{pair_eq_iff}.
            Thus $r\in\dom{f}$.
            Contradiction.
        \end{subproof}
        $f$ is a function by \cref{funs_is_function}.
        Thus $f$ is right-unique and $f$ is a relation.
        Show that $h$ is right-unique.
        \begin{subproof}
            Show that for all $a,b,b'$ such that $(a,b)\in\{(r,U_r)\}$ and $(a,b')\in\{(r,U_r)\}$ we have
                $b = b'$.
            \begin{subproof}
                Fix $a$.
                Fix $b$.
                Fix $b'$.
                Assume $(a,b)\in\{(r,U_r)\}$ and $(a,b')\in\{(r,U_r)\}$.
                Thus $(a,b) = (r,U_r)$ and $(a,b') = (r,U_r)$ by \cref{singleton_elim}.
                Thus $b = b'$ by \cref{pair_eq_iff}.
            \end{subproof}
            Thus $h$ is right-unique by \cref{rightunique}.
        \end{subproof}
        Show that $h$ is a relation.
        \begin{subproof}
            Show that for all $w\in\{(r,U_r)\}$ there exist $x, y$ such that $w = (x,y)$.
            \begin{subproof}
                Fix $w$.
                Assume $w\in\{(r,U_r)\}$.
                Thus $w = (r,U_r)$ by \cref{singleton_elim}.
                Thus there exist $x, y$ such that $w = (x,y)$.
            \end{subproof}
            Thus $h$ is a relation by \cref{relation}.
        \end{subproof}
        Thus $h$ is a function.
        Thus $f\union h$ is a function by \cref{union_compatible_functions}.
        Thus $g$ is a function.
    \end{subproof}
    Thus there exist $p_0,q_0,U_r,h,g$ such that
        $p_0$ is a largest element of $S$ with respect to $\realOrderRel$ and
        $q_0$ is a smallest element of $T_1$ with respect to $\realOrderRel$ and
        $U_r$ is open in $X$ with respect to $T$ and
        $\closure{f(p_0)}{X}{T}\subseteq U_r$ and
        $\closure{U_r}{X}{T}\subseteq f(q_0)$ and
        $h = \{(r,U_r)\}$ and
        $g = f\union h$ and
        $g$ is a function.
\end{proof}

% from §33 helper lemma 29: extension for $r < 1$
\begin{lemma}\label{urysohn_unit_interval_extend_lt1}
    Assume $X$ is normal with respect to $T$.
    Assume $A$ is closed in $X$ with respect to $T$.
    Assume $B$ is closed in $X$ with respect to $T$.
    Assume $A$ is disjoint from $B$.
    Assume $s$ is a sequence in $\unitRationals$.
    Let $G = \urysohnGoodPairs{s}{X}{T}{A}{B}$.
    Let $R = \{(z,w) \mid z\in G, w\in G \mid
        \fst{w} = \fst{z} + 1 \land \snd{z}\subseteq\snd{w}\}$.
    Suppose $n\in\naturalsPlus$ and $f\in\urysohnAdmissible{s}{X}{T}{A}{B}{n}$ and $z = (n,f)$.
    Let $r = s(n + 1)$.
    Assume $r\notin\dom{f}$.
    Assume $r < 1$.
    Let $S = \{p\in\dom{f} \mid p < r\}$.
    Let $T_1 = \{q\in\dom{f} \mid r < q\}$.
    Then there exists $w$ such that $z\mathrel{R} w$.
\end{lemma}
\begin{proof}
    Thus $f\in\funs{\{0,1\}\union \img{s}{\initialSegment{n}}}{\pow{X}}$ and
        for all $p\in\dom{f}$ we have $f(p)\in T$ and
        $f(1) = X\setminus B$ and
        $A\subseteq f(0)$ and
        $\closure{f(0)}{X}{T}\subseteq f(1)$ and
        for all $p\in\dom{f}$ for all $q\in\dom{f}$ if $p < q$ then
            $\closure{f(p)}{X}{T}\subseteq f(q)$ and
        for all $p\in\dom{f}$ if $1 < p$ then $f(p) = X$
        by \cref{urysohn_admissible}.
    Take $p_0,q_0,U_r,h,g$ such that
        $p_0$ is a largest element of $S$ with respect to $\realOrderRel$ and
        $q_0$ is a smallest element of $T_1$ with respect to $\realOrderRel$ and
        $U_r$ is open in $X$ with respect to $T$ and
        $\closure{f(p_0)}{X}{T}\subseteq U_r$ and
        $\closure{U_r}{X}{T}\subseteq f(q_0)$ and
        $h = \{(r,U_r)\}$ and
        $g = f\union h$ and
        $g$ is a function
        by \cref{urysohn_unit_interval_extend_lt1_setup}.
    Show that $\dom{g} = \{0,1\}\union \img{s}{\initialSegment{n + 1}}$.
    \begin{subproof}
        $\dom{g} = \dom{f}\union\dom{h}$ by \cref{dom_union}.
        Show that $\dom{h} = \{r\}$.
        \begin{subproof}
            Show that for all $x\in\dom{h}$ we have $x\in\{r\}$.
            \begin{subproof}
                Fix $x$.
                Assume $x\in\dom{h}$.
                Take $y$ such that $(x,y)\in h$ by \cref{dom_iff}.
                Thus $(x,y) = (r,U_r)$ by \cref{singleton_elim}.
                Thus $x = r$ by \cref{pair_eq_iff}.
                Thus $x\in\{r\}$ by \cref{singleton_intro}.
            \end{subproof}
            Thus $\dom{h}\subseteq \{r\}$ by \cref{subseteq}.
            $r\in\dom{h}$ by \cref{singleton_intro,dom_intro}.
            Show that for all $x\in\{r\}$ we have $x\in\dom{h}$.
            \begin{subproof}
                Fix $x$.
                Assume $x\in\{r\}$.
                Thus $x = r$ by \cref{singleton_elim}.
                Thus $x\in\dom{h}$.
            \end{subproof}
            Thus $\{r\}\subseteq\dom{h}$ by \cref{subseteq}.
            Thus $\dom{h} = \{r\}$ by \cref{subseteq_antisymmetric}.
        \end{subproof}
        Thus $\dom{g} = \dom{f}\union\{r\}$.
        $\dom{f} = \{0,1\}\union \img{s}{\initialSegment{n}}$ by \cref{funs_elim}.
        Show that $\dom{g}\subseteq \{0,1\}\union \img{s}{\initialSegment{n + 1}}$.
        \begin{subproof}
            Show that for all $y\in\dom{g}$ we have $y\in \{0,1\}\union \img{s}{\initialSegment{n + 1}}$.
            \begin{subproof}
                Fix $y$.
                Assume $y\in\dom{g}$.
                Thus $y\in\dom{f}\union\{r\}$.
                Then $y\in\dom{f}$ or $y\in\{r\}$ by \cref{union_iff}.
                \begin{byCase}
                \caseOf{$y\in\dom{f}$.}
                    Thus $y\in \{0,1\}\union \img{s}{\initialSegment{n}}$.
                    Then $y\in \{0,1\}$ or $y\in \img{s}{\initialSegment{n}}$ by \cref{union_iff}.
                    \begin{byCase}
                    \caseOf{$y\in \{0,1\}$.}
                        Thus $y\in \{0,1\}\union \img{s}{\initialSegment{n + 1}}$ by \cref{union_intro_left}.
                    \caseOf{$y\in \img{s}{\initialSegment{n}}$.}
                        Take $k$ such that $k\in\initialSegment{n}$ and $(k,y)\in s$ by \cref{img_iff}.
                        $k\in\naturalsPlus$ and $k \leq n$ by \cref{initial_segment_naturalsplus}.
                        $n\in\naturalsPlus$.
                        $n\in\reals$ and $k\in\reals$ by \cref{real_naturals_subset_reals,subseteq}.
                        $1\in\reals$ by \cref{real_naturals_subset_reals,real_one_in_naturals,subseteq}.
                        $0\in\reals$ by \cref{real_add_ident}.
                        $0 < 1$ by \cref{real_zero_lt_one}.
                        $0 + n < 1 + n$ by \cref{real_lt_add}.
                        $0 + n = n$ by \cref{real_add_ident}.
                        $1 + n = n + 1$ by \cref{real_add_comm}.
                        $n + 1\in\reals$ by \cref{real_add_closed}.
                        Thus $n < n + 1$ by \cref{real_lt_subst_left}.
                        Thus $n \leq n + 1$ by \cref{real_lt_imp_leq}.
                        Thus $k \leq n + 1$ by \cref{real_leq_trans}.
                        Thus $k\in\initialSegment{n + 1}$ by \cref{initial_segment_naturalsplus}.
                        Thus $y\in \img{s}{\initialSegment{n + 1}}$ by \cref{img_iff}.
                        Thus $y\in \{0,1\}\union \img{s}{\initialSegment{n + 1}}$ by \cref{union_intro_right}.
                    \end{byCase}
                \caseOf{$y\in\{r\}$.}
                    Thus $y = r$ by \cref{singleton_elim}.
                    Thus $y = s(n + 1)$.
                    $n + 1\in\naturalsPlus$ by \cref{naturalsplus_add_closed,real_one_in_naturals}.
                    Thus $n + 1 \leq n + 1$.
                    Thus $n + 1\in\initialSegment{n + 1}$ by \cref{initial_segment_naturalsplus}.
                    $s\in\funs{\naturalsPlus}{\unitRationals}$ by \cref{sequence_in}.
                    Thus $\dom{s} = \naturalsPlus$ by \cref{funs_elim}.
                    Thus $n + 1\in\dom{s}$.
                    $s$ is a function by \cref{funs_is_function}.
                    Thus $(n + 1,y)\in s$ by \cref{function_apply_elim}.
                    Thus $y\in \img{s}{\initialSegment{n + 1}}$ by \cref{img_iff}.
                    Thus $y\in \{0,1\}\union \img{s}{\initialSegment{n + 1}}$ by \cref{union_intro_right}.
                \end{byCase}
            \end{subproof}
            Thus $\dom{g}\subseteq \{0,1\}\union \img{s}{\initialSegment{n + 1}}$ by \cref{subseteq}.
        \end{subproof}
        Show that $\{0,1\}\union \img{s}{\initialSegment{n + 1}}\subseteq \dom{g}$.
        \begin{subproof}
            Show that for all $y\in \{0,1\}\union \img{s}{\initialSegment{n + 1}}$ we have $y\in\dom{g}$.
            \begin{subproof}
                Fix $y$.
                Assume $y\in \{0,1\}\union \img{s}{\initialSegment{n + 1}}$.
                Then $y\in \{0,1\}$ or $y\in \img{s}{\initialSegment{n + 1}}$ by \cref{union_iff}.
                \begin{byCase}
                \caseOf{$y\in \{0,1\}$.}
                    Thus $y\in \{0,1\}\union \img{s}{\initialSegment{n}}$ by \cref{union_intro_left}.
                    Thus $y\in\dom{f}$.
                    Thus $y\in\dom{g}$ by \cref{union_intro_left}.
                \caseOf{$y\in \img{s}{\initialSegment{n + 1}}$.}
                    Thus there exists $k$ such that $k\in\initialSegment{n + 1}$ and $(k,y)\in s$ by \cref{img_iff}.
                    Take $k$ such that $k\in\initialSegment{n + 1}$ and $(k,y)\in s$.
                    \begin{byCase}
                    \caseOf{$k\in\initialSegment{n}$.}
                        Thus $y\in \img{s}{\initialSegment{n}}$ by \cref{img_iff}.
                        Thus $y\in \{0,1\}\union \img{s}{\initialSegment{n}}$ by \cref{union_intro_right}.
                        Thus $y\in\dom{f}$.
                        Thus $y\in\dom{g}$ by \cref{union_intro_left}.
                    \caseOf{$k\notin\initialSegment{n}$.}
                        $k\in\naturalsPlus$ and $k\leq n + 1$ by \cref{initial_segment_naturalsplus}.
                        Show that $k = n + 1$.
                        \begin{subproof}
                            Suppose not.
                            Thus $k \neq n + 1$.
                            Thus $k \leq n$ by \cref{real_naturals_leq_succ_elim}.
                            Thus $k\in\initialSegment{n}$ by \cref{initial_segment_naturalsplus}.
                            Contradiction.
                        \end{subproof}
                        Thus $y = r$.
                        Thus $y\in\{r\}$ by \cref{singleton_intro}.
                        Thus $y\in\dom{g}$ by \cref{union_intro_right}.
                    \end{byCase}
                \end{byCase}
            \end{subproof}
            Thus $\{0,1\}\union \img{s}{\initialSegment{n + 1}}\subseteq \dom{g}$ by \cref{subseteq}.
        \end{subproof}
        Thus $\dom{g} = \{0,1\}\union \img{s}{\initialSegment{n + 1}}$ by \cref{subseteq_antisymmetric}.
    \end{subproof}
    Show that $\dom{g} = \dom{f}\union\{r\}$.
    \begin{subproof}
        $\dom{g} = \dom{f}\union\dom{h}$ by \cref{dom_union}.
        Show that $\dom{h} = \{r\}$.
        \begin{subproof}
            Show that for all $x\in\dom{h}$ we have $x\in\{r\}$.
            \begin{subproof}
                Fix $x$.
                Assume $x\in\dom{h}$.
                Take $y$ such that $(x,y)\in h$ by \cref{dom_iff}.
                Thus $(x,y) = (r,U_r)$ by \cref{singleton_elim}.
                Thus $x = r$ by \cref{pair_eq_iff}.
                Thus $x\in\{r\}$ by \cref{singleton_intro}.
            \end{subproof}
            Thus $\dom{h}\subseteq \{r\}$ by \cref{subseteq}.
            $r\in\dom{h}$ by \cref{singleton_intro,dom_intro}.
            Thus $\{r\}\subseteq\dom{h}$ by \cref{singleton_subset_intro}.
            Thus $\dom{h} = \{r\}$ by \cref{subseteq_antisymmetric}.
        \end{subproof}
        Thus $\dom{g} = \dom{f}\union\{r\}$.
    \end{subproof}
    Show that for all $p\in\dom{g}$ we have $g(p)$ is open in $X$ with respect to $T$.
    \begin{subproof}
        Fix $p$.
        Assume $p\in\dom{g}$.
        Thus $p\in\dom{f}\union\{r\}$.
        Then $p\in\dom{f}$ or $p\in\{r\}$ by \cref{union_iff}.
        \begin{byCase}
        \caseOf{$p\in\dom{f}$.}
            Thus $p\in\dom{f}$.
            Thus $f(p)\in T$ by \cref{urysohn_admissible}.
            Thus $f(p)$ is open in $X$ with respect to $T$ by \cref{open_in,topology_on,normal_space}.
            Thus $f(p)\in T$.
            $f$ is a function by \cref{funs_is_function}.
            Thus $f$ is right-unique and $f$ is a relation.
            Thus $(p,f(p))\in f$ by \cref{function_apply_elim}.
            Thus $(p,f(p))\in g$ by \cref{union_intro_left}.
            Thus $g(p) = f(p)$ by \cref{function_apply_intro}.
            Thus $g(p)$ is open in $X$ with respect to $T$.
        \caseOf{$p\in\{r\}$.}
            Thus $p = r$ by \cref{singleton_elim}.
            Thus $(r,U_r)\in h$ by \cref{singleton_intro}.
            Thus $(r,U_r)\in g$ by \cref{union_intro_right}.
            Thus $(p,U_r) = (r,U_r)$ by \cref{pair_eq_iff}.
            Thus $(p,U_r)\in g$.
            Thus $g(p) = U_r$ by \cref{function_apply_intro}.
            Thus $g(p)$ is open in $X$ with respect to $T$.
        \end{byCase}
    \end{subproof}
    Show that for all $p,q$ such that $p\in\dom{g}$ and $q\in\dom{g}$ and $p < q$ we have
        $\closure{g(p)}{X}{T}\subseteq g(q)$.
    \begin{subproof}
        Fix $p$.
        Fix $q$.
        Assume $p\in\dom{g}$ and $q\in\dom{g}$ and $p < q$.
        Thus $p\in\dom{f}\union\{r\}$ and $q\in\dom{f}\union\{r\}$.
        Then $p\in\dom{f}$ or $p\in\{r\}$ by \cref{union_iff}.
        Then $q\in\dom{f}$ or $q\in\{r\}$ by \cref{union_iff}.
        \begin{byCase}
        \caseOf{$p\in\dom{f}$.}
            \begin{byCase}
            \caseOf{$q\in\dom{f}$.}
                Thus $\closure{f(p)}{X}{T}\subseteq f(q)$ by \cref{urysohn_admissible}.
                $f$ is a function by \cref{funs_is_function}.
                Thus $(p,f(p))\in f$ and $(q,f(q))\in f$ by \cref{function_apply_elim}.
                Thus $(p,f(p))\in g$ by \cref{union_intro_left}.
                Thus $(q,f(q))\in g$ by \cref{union_intro_left}.
                Thus $g(p) = f(p)$ and $g(q) = f(q)$ by \cref{function_apply_intro}.
                Thus $\closure{g(p)}{X}{T}\subseteq g(q)$.
            \caseOf{$q\in\{r\}$.}
                Thus $q = r$ by \cref{singleton_elim}.
                Thus $q\notin\dom{f}$.
                Thus $p\neq q$.
                Thus $p < r$.
                $p\in\dom{f}$.
                Show that $\closure{g(p)}{X}{T}\subseteq g(q)$.
                \begin{subproof}
                    $f$ is a function by \cref{funs_is_function}.
                    Thus $(p,f(p))\in f$ by \cref{function_apply_elim}.
                    Thus $(p,f(p))\in g$ by \cref{union_intro_left}.
                    Thus $g(p) = f(p)$ by \cref{function_apply_intro}.
                    Show that $\closure{f(p)}{X}{T}\subseteq U_r$.
                    \begin{subproof}
                        Thus $p\in S$.
                        Thus $p\mathrel{\realOrderRel} p_0$ or $p = p_0$ by \cref{largest_element}.
                        \begin{byCase}
                        \caseOf{$p\mathrel{\realOrderRel} p_0$.}
                            Thus $p < p_0$ by \cref{real_order_relation_elim}.
                            Thus $p_0\in S$ by \cref{largest_element}.
                            Thus $p_0\in\dom{f}$.
                            Thus $\closure{f(p)}{X}{T}\subseteq f(p_0)$ by \cref{urysohn_admissible}.
                            $T$ is a topology on $X$ by \cref{normal_space}.
                            Thus $T\subseteq\pow{X}$ by \cref{topology_on}.
                            Thus $f(p_0)\in\pow{X}$ by \cref{elem_subseteq}.
                            Thus $f(p_0)\subseteq X$ by \cref{powerset_elim}.
                            $f(p_0)\subseteq \closure{f(p_0)}{X}{T}$ by \cref{subseteq_closure}.
                            Thus $\closure{f(p)}{X}{T}\subseteq \closure{f(p_0)}{X}{T}$ by \cref{subseteq_transitive}.
                            Thus $\closure{f(p)}{X}{T}\subseteq U_r$ by \cref{subseteq_transitive}.
                        \caseOf{$p = p_0$.}
                            Thus $\closure{f(p)}{X}{T}\subseteq U_r$.
                        \end{byCase}
                    \end{subproof}
                    Thus $(r,U_r)\in h$ by \cref{singleton_intro}.
                    Thus $(r,U_r)\in g$ by \cref{union_intro_right}.
                    Thus $g(q) = U_r$ by \cref{function_apply_intro}.
                    Thus $\closure{g(p)}{X}{T}\subseteq U_r$.
                    Thus $\closure{g(p)}{X}{T}\subseteq g(q)$.
                \end{subproof}
            \end{byCase}
        \caseOf{$p\in\{r\}$.}
            \begin{byCase}
            \caseOf{$q\in\dom{f}$.}
                Thus $p = r$ by \cref{singleton_elim}.
                Thus $q\neq r$.
                Thus $r < q$.
                Show that $\closure{g(p)}{X}{T}\subseteq g(q)$.
                \begin{subproof}
                    Thus $(r,U_r)\in h$ by \cref{singleton_intro}.
                    Thus $(r,U_r)\in g$ by \cref{union_intro_right}.
                    Thus $g(p) = U_r$ by \cref{function_apply_intro}.
                    $f$ is a function by \cref{funs_is_function}.
                    Thus $(q,f(q))\in f$ by \cref{function_apply_elim}.
                    Thus $(q,f(q))\in g$ by \cref{union_intro_left}.
                    Thus $g(q) = f(q)$ by \cref{function_apply_intro}.
                    Show that $\closure{U_r}{X}{T}\subseteq f(q)$.
                    \begin{subproof}
                        Thus $q\in T_1$.
                        Thus $q_0\mathrel{\realOrderRel} q$ or $q_0 = q$ by \cref{smallest_element}.
                        \begin{byCase}
                        \caseOf{$q_0\mathrel{\realOrderRel} q$.}
                            Thus $q_0 < q$ by \cref{real_order_relation_elim}.
                            Thus $q_0\in T_1$ by \cref{smallest_element}.
                            Thus $q_0\in\dom{f}$.
                            Thus $\closure{f(q_0)}{X}{T}\subseteq f(q)$ by \cref{urysohn_admissible}.
                            Thus $f(q_0)\in T$ by \cref{urysohn_admissible}.
                            $T$ is a topology on $X$ by \cref{normal_space}.
                            Thus $T\subseteq\pow{X}$ by \cref{topology_on}.
                            Thus $f(q_0)\in\pow{X}$ by \cref{elem_subseteq}.
                            Thus $f(q_0)\subseteq X$ by \cref{powerset_elim}.
                            $f(q_0)\subseteq \closure{f(q_0)}{X}{T}$ by \cref{subseteq_closure}.
                            Thus $\closure{U_r}{X}{T}\subseteq \closure{f(q_0)}{X}{T}$ by \cref{subseteq_transitive}.
                            Thus $\closure{U_r}{X}{T}\subseteq f(q)$ by \cref{subseteq_transitive}.
                        \caseOf{$q_0 = q$.}
                            Thus $\closure{U_r}{X}{T}\subseteq f(q)$.
                        \end{byCase}
                    \end{subproof}
                    Thus $\closure{g(p)}{X}{T}\subseteq g(q)$.
                \end{subproof}
            \caseOf{$q\in\{r\}$.}
                Show that $\closure{g(p)}{X}{T}\subseteq g(q)$.
                \begin{subproof}
                    Suppose not.
                    Thus $q = r$ by \cref{singleton_elim}.
                    Thus $p = r$ by \cref{singleton_elim}.
                    Thus $p < q$.
                    Thus $p < r$ by \cref{real_lt_subst_right}.
                    Thus $r = p$.
                    Thus $s\in\funs{\naturalsPlus}{\unitRationals}$ by \cref{sequence_in}.
                    Thus $n + 1\in\naturalsPlus$ by \cref{naturalsplus_add_closed,real_one_in_naturals}.
                    Thus $r\in\unitRationals$ by \cref{funs_type_apply}.
                    Thus $r\in\reals$ by \cref{unit_rational_real}.
                    Thus $p\in\reals$.
                    Thus $p < p$ by \cref{real_lt_subst_right}.
                    Contradiction by \cref{real_lt_irreflexive}.
                \end{subproof}
            \end{byCase}
        \end{byCase}
    \end{subproof}
    Show that for all $p\in\dom{g}$ if $1 < p$ then $g(p) = X$.
    \begin{subproof}
        Fix $p$.
        Assume $p\in\dom{g}$.
        Assume $1 < p$.
        Thus $p\in\dom{f}\union\{r\}$.
        Then $p\in\dom{f}$ or $p\in\{r\}$ by \cref{union_iff}.
        \begin{byCase}
        \caseOf{$p\in\dom{f}$.}
            Thus $f(p) = X$ by \cref{urysohn_admissible_above_one}.
            $f$ is a function by \cref{funs_is_function}.
            Thus $(p,f(p))\in f$ by \cref{function_apply_elim}.
            Thus $(p,f(p))\in g$ by \cref{union_intro_left}.
            Thus $g(p) = f(p)$ by \cref{function_apply_intro}.
            Thus $g(p) = X$.
        \caseOf{$p\in\{r\}$.}
            Show that $g(p) = X$.
            \begin{subproof}
                Suppose not.
                Thus $p = r$ by \cref{singleton_elim}.
                Thus $r < 1$.
                Thus $p < 1$ by \cref{real_lt_subst_left}.
                Thus $s\in\funs{\naturalsPlus}{\unitRationals}$ by \cref{sequence_in}.
                Thus $n + 1\in\naturalsPlus$ by \cref{naturalsplus_add_closed,real_one_in_naturals}.
                Thus $r\in\unitRationals$ by \cref{funs_type_apply}.
                Thus $r\in\reals$ by \cref{unit_rational_real}.
                Thus $p\in\reals$.
                $1\in\reals$ by \cref{real_naturals_subset_reals,real_one_in_naturals,subseteq}.
                Thus $1 < p$.
                Thus $p < p$ by \cref{real_lt_trans}.
                Contradiction by \cref{real_lt_irreflexive}.
            \end{subproof}
        \end{byCase}
    \end{subproof}
    Show that $g\in\funs{\{0,1\}\union \img{s}{\initialSegment{n + 1}}}{\pow{X}}$.
    \begin{subproof}
        Show that $g$ is a function to $\pow{X}$.
        \begin{subproof}
            Show that for all $p\in\dom{g}$ we have $g(p)\in\pow{X}$.
            \begin{subproof}
                Fix $p$.
                Assume $p\in\dom{g}$.
                Thus $g(p)$ is open in $X$ with respect to $T$.
                Thus $g(p)\in T$ by \cref{open_in}.
                $T$ is a topology on $X$ by \cref{normal_space}.
                Thus $T\subseteq\pow{X}$ by \cref{topology_on}.
                Thus $g(p)\in\pow{X}$ by \cref{elem_subseteq}.
            \end{subproof}
            Thus $g$ is a function to $\pow{X}$.
        \end{subproof}
        Thus $g\in\funs{\{0,1\}\union \img{s}{\initialSegment{n + 1}}}{\pow{X}}$ by \cref{funs_intro}.
    \end{subproof}
    Show that $g(1) = X\setminus B$ and $A\subseteq g(0)$ and
        $\closure{g(0)}{X}{T}\subseteq g(1)$.
    \begin{subproof}
        $\dom{f} = \{0,1\}\union \img{s}{\initialSegment{n}}$ by \cref{funs_elim}.
        $0\in\{0,1\}$ by \cref{upair_intro_left}.
        $1\in\{0,1\}$ by \cref{upair_intro_right}.
        Thus $0\in\dom{f}$ and $1\in\dom{f}$ by \cref{union_intro_left}.
        $f$ is a function by \cref{funs_is_function}.
        Thus $(0,f(0))\in f$ and $(1,f(1))\in f$ by \cref{function_apply_elim}.
        Thus $(0,f(0))\in g$ and $(1,f(1))\in g$ by \cref{union_intro_left}.
        Thus $g(0) = f(0)$ and $g(1) = f(1)$ by \cref{function_apply_intro}.
        Thus $g(1) = X\setminus B$ and $A\subseteq g(0)$ and
            $\closure{g(0)}{X}{T}\subseteq g(1)$ by \cref{urysohn_admissible}.
    \end{subproof}
    Show that for all $p\in\dom{g}$ we have $g(p)\in T$.
    \begin{subproof}
        Fix $p$.
        Assume $p\in\dom{g}$.
        Thus $g(p)$ is open in $X$ with respect to $T$.
        Thus $g(p)\in T$ by \cref{open_in}.
    \end{subproof}
    Thus $g\in\urysohnAdmissible{s}{X}{T}{A}{B}{n + 1}$ by \cref{urysohn_admissible}.
    Let $w = (n + 1,g)$.
    Show that $w\in G$.
    \begin{subproof}
        Thus $n + 1\in\naturalsPlus$ by \cref{naturalsplus_add_closed,real_one_in_naturals}.
        Thus $w\in G$ by \cref{urysohn_good_pair_intro}.
    \end{subproof}
    Show that $z\mathrel{R} w$.
    \begin{subproof}
        Thus $z\in G$.
        Thus $w\in G$.
        Thus $\fst{z} = n$ by \cref{fst_eq}.
        Thus $\fst{w} = n + 1$ by \cref{fst_eq}.
        Thus $\fst{w} = \fst{z} + 1$.
        Thus $\snd{z} = f$ by \cref{snd_eq}.
        Thus $\snd{w} = g$ by \cref{snd_eq}.
        $f\subseteq g$ by \cref{union_upper_left}.
        Thus $\snd{z}\subseteq \snd{w}$.
        Thus $(z,w)\in R$.
    \end{subproof}
    Thus $z\mathrel{R} w$.
\end{proof}

% from §33 helper lemma 30: extension relation on good pairs
\begin{lemma}\label{urysohn_unit_interval_extend}
    Assume $X$ is normal with respect to $T$.
    Assume $A$ is closed in $X$ with respect to $T$.
    Assume $B$ is closed in $X$ with respect to $T$.
    Assume $A$ is disjoint from $B$.
    Assume $s$ is a sequence in $\unitRationals$.
    Assume for all $p$ such that $p\in\unitRationals$ there exists $n\in\naturalsPlus$ such that $s(n) = p$.
    Let $G = \urysohnGoodPairs{s}{X}{T}{A}{B}$.
    Let $R = \{(z,w) \mid z\in G, w\in G \mid
        \fst{w} = \fst{z} + 1 \land \snd{z}\subseteq\snd{w}\}$.
    Then for all $z\in G$ there exists $w$ such that $z\mathrel{R} w$.
\end{lemma}
\begin{proof}
    Fix $z$.
    Assume $z\in G$.
    Take $n,f$ such that
        $n\in\naturalsPlus$ and
        $f\in\urysohnAdmissible{s}{X}{T}{A}{B}{n}$ and
        $z = (n,f)$
        by \cref{urysohn_good_pair_decomp}.
    Thus $f\in\funs{\{0,1\}\union \img{s}{\initialSegment{n}}}{\pow{X}}$ and
        for all $p\in\dom{f}$ we have $f(p)\in T$ and
        $f(1) = X\setminus B$ and
        $A\subseteq f(0)$ and
        $\closure{f(0)}{X}{T}\subseteq f(1)$ and
        for all $p\in\dom{f}$ for all $q\in\dom{f}$ if $p < q$ then
            $\closure{f(p)}{X}{T}\subseteq f(q)$ and
        for all $p\in\dom{f}$ if $1 < p$ then $f(p) = X$
        by \cref{urysohn_admissible}.
    Let $r = s(n + 1)$.
    $s\in\funs{\naturalsPlus}{\unitRationals}$ by \cref{sequence_in}.
    Thus $r\in\unitRationals$ by \cref{funs_type_apply,naturalsplus_add_closed,real_one_in_naturals}.
    Show that $r\in\reals$ and $0 < r$.
    \begin{subproof}
        Take $q$ such that $q\in\naturalsPlus\times\naturalsPlus$ and $r = \fst{q} \rmul \inv{\snd{q}}$
            by \cref{unit_rationals}.
        Thus $\fst{q}\in\naturalsPlus$ by \cref{fst_type}.
        Thus $\snd{q}\in\naturalsPlus$ by \cref{snd_type}.
        $\fst{q}\in\reals$ and $\snd{q}\in\reals$ by \cref{real_naturals_subset_reals,subseteq}.
        $1 < \fst{q}$ or $1 = \fst{q}$ by \cref{real_one_leq_naturalsplus}.
        $1 < \snd{q}$ or $1 = \snd{q}$ by \cref{real_one_leq_naturalsplus}.
        $0 < 1$ by \cref{real_zero_lt_one}.
        $0\in\reals$ by \cref{real_add_ident}.
        $1\in\reals$ by \cref{real_naturals_subset_reals,real_one_in_naturals,subseteq}.
        Show that $0 < \fst{q}$.
        \begin{subproof}
            \begin{byCase}
            \caseOf{$1 < \fst{q}$.} Thus $0 < \fst{q}$ by \cref{real_lt_trans}.
            \caseOf{$1 = \fst{q}$.} Thus $0 < \fst{q}$ by \cref{real_zero_lt_one}.
            \end{byCase}
        \end{subproof}
        Show that $0 < \snd{q}$.
        \begin{subproof}
            \begin{byCase}
            \caseOf{$1 < \snd{q}$.} Thus $0 < \snd{q}$ by \cref{real_lt_trans}.
            \caseOf{$1 = \snd{q}$.} Thus $0 < \snd{q}$ by \cref{real_zero_lt_one}.
            \end{byCase}
        \end{subproof}
        Thus $\snd{q} \neq 0$ by \cref{real_pos_nonzero}.
        Thus $\inv{\snd{q}}\in\reals$ by \cref{real_mul_inverse}.
        Thus $0 < \inv{\snd{q}}$ by \cref{real_inv_pos}.
        Thus $r\in\reals$ by \cref{real_mul_closed}.
        Thus $r > 0$ by \cref{real_mul_pos_pos_gt}.
    \end{subproof}
    \begin{byCase}
    \caseOf{$r\in\dom{f}$.}
        Thus there exists $w$ such that $z\mathrel{R} w$ by \cref{urysohn_unit_interval_extend_in_dom}.
    \caseOf{$r\notin\dom{f}$.}
        \begin{byCase}
        \caseOf{$1 < r$.}
            Thus there exists $w$ such that $z\mathrel{R} w$ by \cref{urysohn_unit_interval_extend_gt1}.
        \caseOf{$r < 1$.}
            Let $S = \{p\in\dom{f} \mid p < r\}$.
            Let $T_1 = \{q\in\dom{f} \mid r < q\}$.
            Thus there exists $w$ such that $z\mathrel{R} w$ by \cref{urysohn_unit_interval_extend_lt1}.
        \end{byCase}
    \end{byCase}
\end{proof}

% from §33 helper lemma 31: Urysohn function vanishes on $A$
\begin{lemma}\label{urysohn_unit_interval_value_A}
    Assume $X$ is normal with respect to $T$.
    Assume $A$ is closed in $X$ with respect to $T$.
    Assume $B$ is closed in $X$ with respect to $T$.
    Assume $A$ is disjoint from $B$.
    Assume $s$ is a sequence in $\unitRationals$.
    Assume for all $p$ such that $p\in\unitRationals$ there exists $n\in\naturalsPlus$ such that $s(n) = p$.
    Let $G = \urysohnGoodPairs{s}{X}{T}{A}{B}$.
    Let $R = \{(z,w) \mid z\in G, w\in G \mid
        \fst{w} = \fst{z} + 1 \land \snd{z}\subseteq\snd{w}\}$.
    Assume $c$ is a function on $G$.
    Assume for all $z\in G$ we have $z\mathrel{R} \apply{c}{z}$.
    Assume $u$ is a sequence in $G$.
    Assume $u(1) = a_0$.
    Assume $a_0 = (1,u_1)$.
    Assume for all $n\in\naturalsPlus$ we have $u(n + 1) = \apply{c}{u(n)}$.
    Assume for all $n\in\naturalsPlus$ we have $\fst{u(n)} = n$.
    Let $F = \{\snd{u(n)} \mid n\in\naturalsPlus\}$.
    Let $U = \unions{F}$.
    Assume $U$ is a function.
    Let $I_0 = \intervalclosed{0}{1}$.
    Let $f = \{(x,y) \mid x\in X, y\in\reals \mid
        \text{$y$ is an infimum of $\urysohnLevelSet{U}{x}$}\}$.
    Assume $f\in\funs{X}{\reals}$.
    Assume $\img{f}{X}\subseteq I_0$.
    Then for all $x\in A$ we have $f(x) = 0$.
\end{lemma}
\begin{proof}
    Fix $x\in A$.
    Let $Q_x = \urysohnLevelSet{U}{x}$.
    Show that $x\in U(0)$.
    \begin{subproof}
        $1\in\naturalsPlus$ by \cref{real_one_in_naturals}.
        Thus $u$ is a sequence in $G$.
        Thus $u\in\funs{\naturalsPlus}{G}$ by \cref{sequence_in}.
        Thus $u(1)\in G$ by \cref{funs_type_apply}.
        Take $n,t$ such that
            $n\in\naturalsPlus$ and
            $t\in\urysohnAdmissible{s}{X}{T}{A}{B}{n}$ and
            $u(1) = (n,t)$
            by \cref{urysohn_good_pair_decomp}.
        Thus $t\in\funs{\{0,1\}\union \img{s}{\initialSegment{n}}}{\pow{X}}$
            by \cref{urysohn_admissible}.
        Thus $t$ is a function by \cref{funs_is_function}.
        Thus $t$ is right-unique and $t$ is a relation.
        Thus $\dom{t} = \{0,1\}\union \img{s}{\initialSegment{n}}$ by \cref{funs_elim}.
        $0\in\{0,1\}$ by \cref{upair_intro_left}.
        Thus $0\in\dom{t}$ by \cref{union_intro_left,subseteq}.
        Thus $(0,t(0))\in t$ by \cref{function_apply_elim}.
        Thus $\snd{u(1)} = t$ by \cref{snd_eq}.
        Thus $t\in F$.
        Thus $(0,t(0))\in U$ by \cref{unions_intro}.
        Thus $U(0) = t(0)$ by \cref{function_apply_intro}.
        Thus $A\subseteq t(0)$ by \cref{urysohn_admissible}.
        Thus $x\in t(0)$ by \cref{subseteq}.
        Thus $x\in U(0)$.
    \end{subproof}
    Thus $0\in Q_x$.
    Show that $f(x)$ is an infimum of $Q_x$.
    \begin{subproof}
        Thus $f$ is a function by \cref{funs_is_function}.
        Thus $\dom{f} = X$ by \cref{funs_elim}.
        Thus $A\subseteq X$ by \cref{closed_in}.
        Thus $x\in X$ by \cref{subseteq}.
        Thus $x\in\dom{f}$ by \cref{subseteq}.
        Thus $(x,f(x))\in f$ by \cref{function_apply_elim}.
        Take $x_0,y_0$ such that
            $x_0\in X$ and
            $y_0\in\reals$ and
            $(x,f(x)) = (x_0,y_0)$ and
            $y_0$ is an infimum of $\urysohnLevelSet{U}{x_0}$.
        Thus $x = x_0$ and $f(x) = y_0$ by \cref{pair_eq_iff}.
        Thus $f(x)$ is an infimum of $Q_x$.
    \end{subproof}
    Thus $f(x)$ is a lower bound of $Q_x$ by \cref{real_infimum}.
    Thus $f(x) \leq 0$ by \cref{real_lower_bound}.
    Show that $f(x)\in I_0$.
    \begin{subproof}
        Thus $f$ is a function by \cref{funs_is_function}.
        Thus $\dom{f} = X$ by \cref{funs_elim}.
        Thus $A\subseteq X$ by \cref{closed_in}.
        Thus $x\in X$ by \cref{subseteq}.
        Thus $x\in\dom{f}$ by \cref{subseteq}.
        Thus $(x,f(x))\in f$ by \cref{function_apply_elim}.
        Thus $f(x)\in\img{f}{X}$ by \cref{img_iff}.
        Thus $f(x)\in I_0$ by \cref{subseteq}.
    \end{subproof}
    Thus $0 \leq f(x)$ by \cref{intervalclosed}.
    $0\in\reals$ by \cref{real_add_ident}.
    Thus $A\subseteq X$ by \cref{closed_in}.
    Thus $x\in X$ by \cref{subseteq}.
    Thus $f(x)\in\reals$ by \cref{funs_type_apply}.
    Thus $f(x) = 0$ by \cref{real_leq_antisym}.
\end{proof}

% from §33 helper lemma 32: Urysohn function equals $1$ on $B$
\begin{lemma}\label{urysohn_unit_interval_value_B}
    Assume $X$ is normal with respect to $T$.
    Assume $A$ is closed in $X$ with respect to $T$.
    Assume $B$ is closed in $X$ with respect to $T$.
    Assume $A$ is disjoint from $B$.
    Assume $s$ is a sequence in $\unitRationals$.
    Assume for all $p$ such that $p\in\unitRationals$ there exists $n\in\naturalsPlus$ such that $s(n) = p$.
    Let $G = \urysohnGoodPairs{s}{X}{T}{A}{B}$.
    Let $R = \{(z,w) \mid z\in G, w\in G \mid
        \fst{w} = \fst{z} + 1 \land \snd{z}\subseteq\snd{w}\}$.
    Assume $c$ is a function on $G$.
    Assume for all $z\in G$ we have $z\mathrel{R} \apply{c}{z}$.
    Assume $u$ is a sequence in $G$.
    Assume $u(1) = a_0$.
    Assume $a_0 = (1,u_1)$.
    Assume for all $n\in\naturalsPlus$ we have $u(n + 1) = \apply{c}{u(n)}$.
    Assume for all $n\in\naturalsPlus$ we have $\fst{u(n)} = n$.
    Let $F = \{\snd{u(n)} \mid n\in\naturalsPlus\}$.
    Let $U = \unions{F}$.
    Assume $U$ is a function.
    Let $I_0 = \intervalclosed{0}{1}$.
    Let $f = \{(x,y) \mid x\in X, y\in\reals \mid
        \text{$y$ is an infimum of $\urysohnLevelSet{U}{x}$}\}$.
    Assume $f\in\funs{X}{\reals}$.
    Assume $\img{f}{X}\subseteq I_0$.
    Then for all $x\in B$ we have $f(x) = 1$.
\end{lemma}
\begin{proof}
    Fix $x\in B$.
    Let $Q_x = \urysohnLevelSet{U}{x}$.
    Show that $x\notin U(1)$.
    \begin{subproof}
        Show that $U(1) = X\setminus B$.
        \begin{subproof}
            $1\in\naturalsPlus$ by \cref{real_one_in_naturals}.
            Thus $u$ is a sequence in $G$.
            Thus $u\in\funs{\naturalsPlus}{G}$ by \cref{sequence_in}.
            Thus $u(1)\in G$ by \cref{funs_type_apply}.
            Take $n,t$ such that
                $n\in\naturalsPlus$ and
                $t\in\urysohnAdmissible{s}{X}{T}{A}{B}{n}$ and
                $u(1) = (n,t)$
                by \cref{urysohn_good_pair_decomp}.
            Thus $t\in\funs{\{0,1\}\union \img{s}{\initialSegment{n}}}{\pow{X}}$
                by \cref{urysohn_admissible}.
            Thus $t$ is a function by \cref{funs_is_function}.
            Thus $t$ is right-unique and $t$ is a relation.
            Thus $\dom{t} = \{0,1\}\union \img{s}{\initialSegment{n}}$ by \cref{funs_elim}.
            $1\in\{0,1\}$ by \cref{upair_intro_right}.
            Thus $1\in\dom{t}$ by \cref{union_intro_left,subseteq}.
            Thus $(1,t(1))\in t$ by \cref{function_apply_elim}.
            Thus $\snd{u(1)} = t$ by \cref{snd_eq}.
            Thus $t\in F$.
            Thus $(1,t(1))\in U$ by \cref{unions_intro}.
            Thus $U(1) = t(1)$ by \cref{function_apply_intro}.
            Thus $t(1) = X\setminus B$ by \cref{urysohn_admissible}.
            Thus $U(1) = X\setminus B$.
        \end{subproof}
        Suppose not.
        Thus $x\in U(1)$.
        Thus $U(1)\subseteq X\setminus B$ by \cref{subseteq_refl}.
        Thus $x\in X\setminus B$ by \cref{subseteq}.
        Thus $x\notin B$ by \cref{setminus_elim_right}.
        Contradiction.
    \end{subproof}
    Show that $f(x)$ is an infimum of $Q_x$.
    \begin{subproof}
        Thus $f$ is a function by \cref{funs_is_function}.
        Thus $\dom{f} = X$ by \cref{funs_elim}.
        Thus $B\subseteq X$ by \cref{closed_in}.
        Thus $x\in X$ by \cref{subseteq}.
        Thus $x\in\dom{f}$ by \cref{subseteq}.
        Thus $(x,f(x))\in f$ by \cref{function_apply_elim}.
        Take $x_0,y_0$ such that
            $x_0\in X$ and
            $y_0\in\reals$ and
            $(x,f(x)) = (x_0,y_0)$ and
            $y_0$ is an infimum of $\urysohnLevelSet{U}{x_0}$.
        Thus $x = x_0$ and $f(x) = y_0$ by \cref{pair_eq_iff}.
        Thus $f(x)$ is an infimum of $Q_x$.
    \end{subproof}
    Thus $f(x)$ is a lower bound of $Q_x$ by \cref{real_infimum}.
    Show that $1$ is a lower bound of $Q_x$.
    \begin{subproof}
        Show that for all $q\in Q_x$ we have $1 \leq q$.
        \begin{subproof}
            Fix $q$.
            Assume $q\in Q_x$.
            Thus $q\in\unitRationalsZero$ and $x\in U(q)$ by \cref{urysohn_level_set}.
            Show that $1 \leq q$.
            \begin{subproof}
                Suppose not.
                Thus $q < 1$.
                Thus $x\in U(q)$.
                Thus $x\in\unions{\img{U}{\{q\}}}$ by \cref{apply}.
                Then there exists $y$ such that $y\in\img{U}{\{q\}}$ and $x\in y$ by \cref{unions_iff}.
                Take $y$ such that $y\in\img{U}{\{q\}}$ and $x\in y$.
                Take $a$ such that $a\in\{q\}$ and $(a,y)\in U$ by \cref{img_iff}.
                Thus $a = q$ by \cref{singleton_elim}.
                Thus $(q,y)\in U$.
                Thus $q\in\dom{U}$ by \cref{dom_intro}.
                Thus $(q,U(q))\in U$ by \cref{function_apply_elim}.
                Thus there exists $t_q$ such that $t_q\in F$ and $(q,U(q))\in t_q$ by \cref{unions_iff}.
                Thus $q\in\dom{t_q}$ by \cref{dom_intro}.
                Thus there exists $n_q$ such that $n_q\in\naturalsPlus$ and $t_q = \snd{u(n_q)}$.
                $u$ is a sequence in $G$.
                Thus $u\in\funs{\naturalsPlus}{G}$ by \cref{sequence_in}.
                Thus $u(n_q)\in G$ by \cref{funs_type_apply}.
                Take $m_q,v_q$ such that
                    $m_q\in\naturalsPlus$ and
                    $v_q\in\urysohnAdmissible{s}{X}{T}{A}{B}{m_q}$ and
                    $u(n_q) = (m_q,v_q)$
                    by \cref{urysohn_good_pair_decomp}.
                Thus $\snd{u(n_q)} = v_q$ by \cref{snd_eq}.
                Thus $t_q = v_q$.
                Thus $t_q\in\urysohnAdmissible{s}{X}{T}{A}{B}{m_q}$.
                Thus $\snd{u(1)} = u_1$ by \cref{snd_eq}.
                Thus $u_1\in F$.
                $1\in\naturalsPlus$ by \cref{real_one_in_naturals}.
                Thus $u(1)\in G$ by \cref{funs_type_apply}.
                Take $n_1,t_1$ such that
                    $n_1\in\naturalsPlus$ and
                    $t_1\in\urysohnAdmissible{s}{X}{T}{A}{B}{n_1}$ and
                    $u(1) = (n_1,t_1)$
                    by \cref{urysohn_good_pair_decomp}.
                Thus $t_1 = u_1$ by \cref{snd_eq}.
                Thus $u_1\in\urysohnAdmissible{s}{X}{T}{A}{B}{n_1}$.
                Thus $t_q\subseteq u_1$ or $u_1\subseteq t_q$ by \cref{urysohn_unit_interval_chain}.
                \begin{byCase}
                \caseOf{$t_q\subseteq u_1$.}
                    Thus $(q,U(q))\in u_1$ by \cref{subseteq}.
                    Thus $q\in\dom{u_1}$ by \cref{dom_intro}.
                    Thus $u_1\in\funs{\{0,1\}\union \img{s}{\initialSegment{n_1}}}{\pow{X}}$
                        by \cref{urysohn_admissible}.
                    Thus $u_1$ is a function by \cref{funs_is_function}.
                    Thus $u_1$ is right-unique and $u_1$ is a relation.
                    Thus $U(q) = u_1(q)$ by \cref{function_apply_intro}.
                    $1\in\{0,1\}$ by \cref{upair_intro_right}.
                    Thus $1\in\dom{u_1}$ by \cref{funs_elim,union_intro_left}.
                    Thus $q < 1$.
                    Thus $\closure{u_1(q)}{X}{T}\subseteq u_1(1)$ by \cref{urysohn_admissible}.
                    Thus $u_1(q)\in T$ by \cref{urysohn_admissible}.
                    $T$ is a topology on $X$ by \cref{normal_space}.
                    Thus $T\subseteq\pow{X}$ by \cref{topology_on}.
                    Thus $u_1(q)\in\pow{X}$ by \cref{elem_subseteq}.
                    Thus $u_1(q)\subseteq X$ by \cref{powerset_elim}.
                    Thus $u_1(q)\subseteq \closure{u_1(q)}{X}{T}$ by \cref{subseteq_closure}.
                    Thus $U(q)\subseteq u_1(1)$ by \cref{subseteq_transitive}.
                    Thus $(1,u_1(1))\in u_1$ by \cref{function_apply_elim}.
                    Thus $(1,u_1(1))\in U$ by \cref{unions_intro}.
                    Thus $U(1) = u_1(1)$ by \cref{function_apply_intro}.
                    Thus $x\in U(1)$ by \cref{subseteq}.
                    Contradiction.
                \caseOf{$u_1\subseteq t_q$.}
                    Thus $u_1\in\funs{\{0,1\}\union \img{s}{\initialSegment{n_1}}}{\pow{X}}$
                        by \cref{urysohn_admissible}.
                    Thus $u_1$ is a function by \cref{funs_is_function}.
                    Thus $\dom{u_1} = \{0,1\}\union \img{s}{\initialSegment{n_1}}$ by \cref{funs_elim}.
                    $1\in\{0,1\}$ by \cref{upair_intro_right}.
                    Thus $1\in\dom{u_1}$ by \cref{union_intro_left}.
                    Thus $(1,u_1(1))\in u_1$ by \cref{function_apply_elim}.
                    Thus $(1,u_1(1))\in t_q$ by \cref{subseteq}.
                    Thus $1\in\dom{t_q}$ by \cref{dom_intro}.
                    Thus $\closure{t_q(q)}{X}{T}\subseteq t_q(1)$ by \cref{urysohn_admissible}.
                    Thus $t_q(q)\in T$ by \cref{urysohn_admissible}.
                    $T$ is a topology on $X$ by \cref{normal_space}.
                    Thus $T\subseteq\pow{X}$ by \cref{topology_on}.
                    Thus $t_q(q)\in\pow{X}$ by \cref{elem_subseteq}.
                    Thus $t_q(q)\subseteq X$ by \cref{powerset_elim}.
                    Thus $t_q(q)\subseteq \closure{t_q(q)}{X}{T}$ by \cref{subseteq_closure}.
                    Thus $t_q\in\funs{\{0,1\}\union \img{s}{\initialSegment{m_q}}}{\pow{X}}$
                        by \cref{urysohn_admissible}.
                    Thus $t_q$ is a function by \cref{funs_is_function}.
                    Thus $t_q$ is right-unique and $t_q$ is a relation.
                    Thus $U(q) = t_q(q)$ by \cref{function_apply_intro}.
                    Thus $U(q)\subseteq t_q(1)$ by \cref{subseteq_transitive}.
                    Thus $(1,t_q(1))\in t_q$ by \cref{function_apply_elim}.
                    Thus $(1,t_q(1))\in U$ by \cref{unions_intro}.
                    Thus $U(1) = t_q(1)$ by \cref{function_apply_intro}.
                    Thus $x\in U(1)$ by \cref{subseteq}.
                    Contradiction.
                \end{byCase}
            \end{subproof}
        \end{subproof}
        $1\in\reals$ by \cref{real_mul_ident}.
        Show that $Q_x\subseteq\reals$.
        \begin{subproof}
            Show that for all $q\in Q_x$ we have $q\in\reals$.
            \begin{subproof}
                Fix $q$.
                Assume $q\in Q_x$.
                Thus $q\in\unitRationalsZero$ by \cref{urysohn_level_set}.
                Then $q\in\unitRationals$ or $q\in\{0\}$ by \cref{unit_rationals_zero,union_iff}.
                \begin{byCase}
                \caseOf{$q\in\unitRationals$.}
                    Thus $q\in\reals$ by \cref{unit_rational_real}.
                \caseOf{$q\in\{0\}$.}
                    Thus $q = 0$ by \cref{singleton_elim}.
                    Thus $q\in\reals$ by \cref{real_add_ident}.
                \end{byCase}
            \end{subproof}
            Thus $Q_x\subseteq\reals$ by \cref{subseteq}.
        \end{subproof}
        Thus $1$ is a lower bound of $Q_x$ by \cref{real_lower_bound}.
    \end{subproof}
    Thus $1 \leq f(x)$ by \cref{real_infimum}.
    Show that $f(x)\in I_0$.
    \begin{subproof}
        Thus $f$ is a function by \cref{funs_is_function}.
        Thus $\dom{f} = X$ by \cref{funs_elim}.
        Thus $B\subseteq X$ by \cref{closed_in}.
        Thus $x\in X$ by \cref{subseteq}.
        Thus $x\in\dom{f}$ by \cref{subseteq}.
        Thus $(x,f(x))\in f$ by \cref{function_apply_elim}.
        Thus $f(x)\in\img{f}{X}$ by \cref{img_iff}.
        Thus $f(x)\in I_0$ by \cref{subseteq}.
    \end{subproof}
    Thus $f(x) \leq 1$ by \cref{intervalclosed}.
    $1\in\reals$ by \cref{real_mul_ident}.
    Thus $B\subseteq X$ by \cref{closed_in}.
    Thus $x\in X$ by \cref{subseteq}.
    Thus $f(x)\in\reals$ by \cref{funs_type_apply}.
    Thus $f(x) = 1$ by \cref{real_leq_antisym}.
\end{proof}

% from §33 helper lemma 33: Urysohn lemma for the unit interval
\begin{lemma}\label{urysohn_unit_interval}
    Assume $X$ is normal with respect to $T$.
    Assume $A$ is closed in $X$ with respect to $T$.
    Assume $B$ is closed in $X$ with respect to $T$.
    Assume $A$ is disjoint from $B$.
    Let $I_0 = \intervalclosed{0}{1}$.
    Let $T_0 = \subspaceTopology{\standardTopologyR}{I_0}$.
    Then there exists $f$ such that
        $f$ is continuous from $X$ to $I_0$ with respect to $T$ and $T_0$ and
        $\img{f}{A}\subseteq \{0\}$ and
        $\img{f}{B}\subseteq \{1\}$.
\end{lemma}
\begin{proof}
    Take $s$ such that
        $s$ is a sequence in $\unitRationals$ and
        for all $p$ such that $p\in\unitRationals$ there exists $n\in\naturalsPlus$ such that $s(n) = p$
        by \cref{unit_rationals_sequence}.
    Let $P_n = \{0,1\}\union \img{s}{\initialSegment{n}}$.

    Let $G = \urysohnGoodPairs{s}{X}{T}{A}{B}$.

    Let $R = \{(z,w) \mid z\in G, w\in G \mid
        \fst{w} = \fst{z} + 1 \land \snd{z}\subseteq\snd{w}\}$.

    Thus for all $z\in G$ there exists $w$ such that $z\mathrel{R} w$
        by \cref{urysohn_unit_interval_extend}.

    Show that $R$ is a relation.
    \begin{subproof}
        Show that for all $v$ such that $v\in R$ there exist $x,y$ such that $v = (x,y)$.
        \begin{subproof}
            Fix $v$.
            Assume $v\in R$.
            Take $z,w$ such that $z\in G$ and $w\in G$ and $v = (z,w)$.
        \end{subproof}
        Thus $R$ is a relation by \cref{relation}.
    \end{subproof}

    Take $c$ such that $c$ is a function on $G$ and for all $z\in G$ we have $z\mathrel{R} \apply{c}{z}$
        by \cref{choice_function}.

    Take $u_1$ such that $(1,u_1)\in G$ by \cref{urysohn_unit_interval_u1}.
    Let $a_0 = (1,u_1)$.

    Show that $c\in\funs{G}{G}$.
    \begin{subproof}
        Show that for all $x\in\dom{c}$ we have $c(x)\in G$.
        \begin{subproof}
            Fix $x$.
            Assume $x\in\dom{c}$.
            Thus $x\in G$.
            Thus $(x,c(x))\in R$.
            Thus there exist $z,w$ such that
                $z\in G$ and
                $w\in G$ and
                $(z,w) = (x,c(x))$.
            Thus $c(x)\in G$ by \cref{pair_eq_iff}.
        \end{subproof}
        Thus $c$ is a function.
        Thus $c$ is a function to $G$.
        Thus $G = \dom{c}$.
        Thus $c\in\funs{G}{G}$ by \cref{funs_intro}.
    \end{subproof}
    Take $u$ such that
        $u$ is a sequence in $G$ and
        $u(1) = a_0$ and
        for all $n\in\naturalsPlus$ we have $u(n + 1) = \apply{c}{u(n)}$
        by \cref{iteration_sequence_naturalsplus}.
    Thus for all $n\in\naturalsPlus$ we have $\fst{u(n)} = n$
        by \cref{urysohn_unit_interval_fst_u}.

    Let $F = \{\snd{u(n)} \mid n\in\naturalsPlus\}$.
    Let $U = \unions{F}$.
    Thus $U$ is a function by \cref{urysohn_unit_interval_union_function}.

    Let $f = \{(x,y) \mid x\in X, y\in\reals \mid
        \text{$y$ is an infimum of $\urysohnLevelSet{U}{x}$}\}$.
    Thus $f$ is a function from $X$ to $\reals$ by \cref{urysohn_unit_interval_f_funs}.

    Show that for all $x\in X$ we have $f(x)\in I_0$.
    \begin{subproof}
        Fix $x\in X$.
        Let $Q_x = \urysohnLevelSet{U}{x}$.
        Show that $0 \leq f(x)$.
        \begin{subproof}
            $0\in\reals$ by \cref{real_add_ident}.
            Show that $0$ is a lower bound of $Q_x$.
            \begin{subproof}
                Show that for all $p\in Q_x$ we have $p\in\reals$ and $0 \leq p$.
                \begin{subproof}
                    Fix $p$.
                    Assume $p\in Q_x$.
                    Thus $p\in\unitRationalsZero$ by \cref{urysohn_level_set}.
                    Thus $p\in\unitRationals$ or $p\in\{0\}$ by \cref{unit_rationals_zero,union_iff}.
                    \begin{byCase}
                    \caseOf{$p\in\unitRationals$.}
                        Take $q$ such that $q\in\naturalsPlus\times\naturalsPlus$ and
                            $p = \fst{q} \rmul \inv{\snd{q}}$
                            by \cref{unit_rationals}.
                        Thus $\fst{q}\in\naturalsPlus$ and $\snd{q}\in\naturalsPlus$ by \cref{fst_type,snd_type}.
                        Thus $\fst{q}\in\reals$ and $\snd{q}\in\reals$
                            by \cref{real_naturals_subset_reals,subseteq}.
                        $0\in\reals$ by \cref{real_add_ident}.
                        $1\in\reals$ by \cref{real_mul_ident}.
                        $0 < 1$ by \cref{real_zero_lt_one}.
                        Show that $0 < \fst{q}$.
                        \begin{subproof}
                            $1 < \fst{q}$ or $1 = \fst{q}$ by \cref{real_one_leq_naturalsplus}.
                            \begin{byCase}
                            \caseOf{$1 < \fst{q}$.}
                                Thus $0 < \fst{q}$ by \cref{real_lt_trans}.
                            \caseOf{$1 = \fst{q}$.}
                                Thus $0 < \fst{q}$ by \cref{real_lt_subst_right}.
                            \end{byCase}
                        \end{subproof}
                        Show that $0 < \snd{q}$.
                        \begin{subproof}
                            $1 < \snd{q}$ or $1 = \snd{q}$ by \cref{real_one_leq_naturalsplus}.
                            \begin{byCase}
                            \caseOf{$1 < \snd{q}$.}
                                Thus $0 < \snd{q}$ by \cref{real_lt_trans}.
                            \caseOf{$1 = \snd{q}$.}
                                Thus $0 < \snd{q}$ by \cref{real_lt_subst_right}.
                            \end{byCase}
                        \end{subproof}
                        Thus $\snd{q}\neq 0$ by \cref{real_pos_nonzero}.
                        Thus $\inv{\snd{q}}\in\reals$ by \cref{real_mul_inverse}.
                        Thus $\inv{\snd{q}}$ is positive by \cref{real_inv_pos}.
                        Thus $0 < \inv{\snd{q}}$.
                        Thus $0 < p$ by \cref{real_mul_pos_pos_gt}.
                        Thus $p\in\reals$ by \cref{real_mul_closed}.
                        Thus $0 \leq p$.
                    \caseOf{$p\in\{0\}$.}
                        Thus $p = 0$ by \cref{singleton_elim}.
                        Thus $p\in\reals$ by \cref{real_add_ident}.
                        Thus $0 \leq p$.
                    \end{byCase}
                \end{subproof}
                Show that $Q_x\subseteq\reals$.
                \begin{subproof}
                    Show that for all $p\in Q_x$ we have $p\in\reals$.
                    \begin{subproof}
                        Fix $p$.
                        Assume $p\in Q_x$.
                        Thus $p\in\reals$ and $0 \leq p$.
                        Thus $p\in\reals$.
                    \end{subproof}
                    Thus $Q_x\subseteq\reals$ by \cref{subseteq}.
                \end{subproof}
                Show that for all $p\in Q_x$ we have $0 \leq p$.
                \begin{subproof}
                    Fix $p$.
                    Assume $p\in Q_x$.
                    Thus $p\in\reals$ and $0 \leq p$.
                    Thus $0 \leq p$.
                \end{subproof}
                Thus $0$ is a lower bound of $Q_x$ by \cref{real_lower_bound}.
            \end{subproof}
            Show that $f(x)$ is an infimum of $Q_x$.
            \begin{subproof}
                Thus $f$ is a function by \cref{funs_is_function}.
                Thus $\dom{f} = X$ by \cref{funs_elim}.
                Thus $x\in\dom{f}$ by \cref{subseteq}.
                Thus $(x,f(x))\in f$ by \cref{function_apply_elim}.
                Take $x_0,y_0$ such that
                    $x_0\in X$ and
                    $y_0\in\reals$ and
                    $(x,f(x)) = (x_0,y_0)$ and
                    $y_0$ is an infimum of $\urysohnLevelSet{U}{x_0}$.
                Thus $x = x_0$ and $f(x) = y_0$ by \cref{pair_eq_iff}.
                Thus $f(x)$ is an infimum of $Q_x$.
            \end{subproof}
            Thus $0 < f(x)$ or $0 = f(x)$ by \cref{real_infimum}.
            Thus $0 \leq f(x)$.
        \end{subproof}
        Show that $f(x) \leq 1$.
        \begin{subproof}
            Suppose not.
            Thus $1 < f(x)$.
            $1\in\reals$ by \cref{real_mul_ident}.
            $0 < 1$ by \cref{real_zero_lt_one}.
            Thus $0 \leq 1$ by \cref{real_lt_imp_leq}.
            Thus $f(x)\in\reals$ by \cref{funs_type_apply}.
            Take $p$ such that $p\in\unitRationals$ and $p\in\intervalopen{1}{f(x)}$
                by \cref{unit_rationals_dense_nonneg}.
            Thus $p\in\reals$ by \cref{intervalopen}.
            Thus $p > 1$.
            Thus $p < f(x)$ by \cref{intervalopen}.
            Thus $p\in\unitRationalsZero$ by \cref{unit_rationals_zero,union_iff}.
            Show that $x\in U(p)$.
            \begin{subproof}
                Show that $U(p) = X$.
                \begin{subproof}
                    Take $k$ such that $k\in\naturalsPlus$ and $s(k) = p$ by \cref{unit_rationals_sequence}.
                    Thus $u$ is a sequence in $G$.
                    Thus $u\in\funs{\naturalsPlus}{G}$ by \cref{sequence_in}.
                    Thus $u(k)\in G$ by \cref{funs_type_apply}.
                    Take $n,t$ such that
                        $n\in\naturalsPlus$ and
                        $t\in\urysohnAdmissible{s}{X}{T}{A}{B}{n}$ and
                        $u(k) = (n,t)$
                        by \cref{urysohn_good_pair_decomp}.
                    Thus $\fst{u(k)} = n$ by \cref{fst_eq}.
                    Thus $n = k$.
                    Thus $u(k) = (k,t)$ by \cref{pair_eq_iff}.
                    Thus $\snd{u(k)} = t$ by \cref{snd_eq}.
                    Thus $t\in F$.
                    Thus $t\in\funs{\{0,1\}\union \img{s}{\initialSegment{k}}}{\pow{X}}$
                        by \cref{urysohn_admissible}.
                    Thus $\dom{t} = \{0,1\}\union \img{s}{\initialSegment{k}}$ by \cref{funs_elim}.
                    Show that $p\in\dom{t}$.
                    \begin{subproof}
                        Thus $k \leq k$.
                        Thus $k\in\initialSegment{k}$ by \cref{initial_segment_naturalsplus}.
                        Let $y = k$.
                        Thus $y\in\initialSegment{k}$.
                        Thus $s(y) = s(k)$.
                        Thus $s$ is a sequence in $\unitRationals$.
                        Thus $s\in\funs{\naturalsPlus}{\unitRationals}$ by \cref{sequence_in}.
                        Thus $s$ is a function by \cref{funs_is_function}.
                        Thus $\dom{s} = \naturalsPlus$ by \cref{funs_elim}.
                        Thus $y\in\dom{s}$ by \cref{subseteq}.
                        Thus $(y,s(y))\in s$ by \cref{function_apply_elim}.
                        Thus $(y,s(k))\in s$.
                        Thus $s(k)\in\img{s}{\initialSegment{k}}$ by \cref{img_iff}.
                        Thus $p\in\img{s}{\initialSegment{k}}$.
                        Thus $p\in\dom{t}$ by \cref{union_intro_right}.
                    \end{subproof}
                    Show that for all $q\in\dom{t}$ if $1 < q$ then $t(q) = X$.
                    \begin{subproof}
                        Show that $0\in\dom{t}$.
                        \begin{subproof}
                            $0\in\{0,1\}$ by \cref{upair_intro_left}.
                            Thus $0\in\dom{t}$ by \cref{union_intro_left,subseteq}.
                        \end{subproof}
                        Show that $1\in\dom{t}$.
                        \begin{subproof}
                            $1\in\{0,1\}$ by \cref{upair_intro_right}.
                            Thus $1\in\dom{t}$ by \cref{union_intro_left,subseteq}.
                        \end{subproof}
                        $0\in\reals$ by \cref{real_add_ident}.
                        $1\in\reals$ by \cref{real_mul_ident}.
                        $0 < 1$ by \cref{real_zero_lt_one}.
                        Thus $\closure{t(0)}{X}{T}\subseteq t(1)$ and
                            for all $q\in\dom{t}$ if $1 < q$ then $t(q) = X$
                            by \cref{urysohn_admissible}.
                        Thus for all $q\in\dom{t}$ if $1 < q$ then $t(q) = X$.
                    \end{subproof}
                    Thus $1 < p$.
                    Thus $t(p) = X$.
                    $t$ is a function by \cref{funs_is_function}.
                    Thus $t$ is right-unique and $t$ is a relation.
                    Thus $(p,t(p))\in t$ by \cref{function_apply_elim}.
                    Thus $(p,t(p))\in U$ by \cref{unions_intro}.
                    Thus $U(p) = t(p)$ by \cref{function_apply_intro}.
                    Thus $U(p) = X$.
                \end{subproof}
                Thus $x\in U(p)$.
            \end{subproof}
            Thus $p\in Q_x$.
            Show that $f(x)$ is an infimum of $Q_x$.
            \begin{subproof}
                Thus $f$ is a function by \cref{funs_is_function}.
                Thus $\dom{f} = X$ by \cref{funs_elim}.
                Thus $x\in\dom{f}$ by \cref{subseteq}.
                Thus $(x,f(x))\in f$ by \cref{function_apply_elim}.
                Take $x_0,y_0$ such that
                    $x_0\in X$ and
                    $y_0\in\reals$ and
                    $(x,f(x)) = (x_0,y_0)$ and
                    $y_0$ is an infimum of $\urysohnLevelSet{U}{x_0}$.
                Thus $x = x_0$ and $f(x) = y_0$ by \cref{pair_eq_iff}.
                Thus $f(x)$ is an infimum of $Q_x$.
            \end{subproof}
            Thus $f(x)$ is a lower bound of $Q_x$ by \cref{real_infimum}.
            Thus $f(x) < p$ or $f(x) = p$ by \cref{real_lower_bound}.
            \begin{byCase}
            \caseOf{$f(x) < p$.}
                Thus $f(x) < f(x)$ by \cref{real_lt_trans}.
                Contradiction by \cref{real_lt_irreflexive}.
            \caseOf{$f(x) = p$.}
                Thus $f(x) < f(x)$ by \cref{real_lt_subst_left}.
                Contradiction by \cref{real_lt_irreflexive}.
            \end{byCase}
        \end{subproof}
        Thus $f(x)\in\reals$ by \cref{funs_type_apply}.
        Thus $f(x)\in I_0$ by \cref{intervalclosed}.
    \end{subproof}
    Show that $\img{f}{X}\subseteq I_0$.
    \begin{subproof}
        Show that for all $y\in\img{f}{X}$ we have $y\in I_0$.
        \begin{subproof}
            Fix $y$.
            Assume $y\in\img{f}{X}$.
            Take $x$ such that $x\in X$ and $(x,y)\in f$ by \cref{img_iff}.
            Thus $f$ is a function by \cref{funs_is_function}.
            Thus $y = f(x)$ by \cref{function_apply_intro}.
            Thus $f(x)\in I_0$.
            Thus $y\in I_0$.
        \end{subproof}
        Thus $\img{f}{X}\subseteq I_0$ by \cref{subseteq}.
    \end{subproof}

    Thus for all $x\in A$ we have $f(x) = 0$
        by \cref{urysohn_unit_interval_value_A}.

    Thus for all $x\in B$ we have $f(x) = 1$
        by \cref{urysohn_unit_interval_value_B}.

    Show that $\img{f}{A}\subseteq \{0\}$.
    \begin{subproof}
        Show that for all $y\in\img{f}{A}$ we have $y\in\{0\}$.
        \begin{subproof}
            Fix $y$.
            Assume $y\in\img{f}{A}$.
            Take $x$ such that $x\in A$ and $(x,y)\in f$ by \cref{img_iff}.
            Thus $f$ is a function by \cref{funs_is_function}.
            Thus $y = f(x)$ by \cref{function_apply_iff}.
            Thus $f(x) = 0$.
            Thus $y = 0$.
            Thus $y\in\{0\}$ by \cref{singleton_intro}.
        \end{subproof}
        Thus $\img{f}{A}\subseteq \{0\}$ by \cref{subseteq}.
    \end{subproof}
    Show that $\img{f}{B}\subseteq \{1\}$.
    \begin{subproof}
        Show that for all $y\in\img{f}{B}$ we have $y\in\{1\}$.
        \begin{subproof}
            Fix $y$.
            Assume $y\in\img{f}{B}$.
            Take $x$ such that $x\in B$ and $(x,y)\in f$ by \cref{img_iff}.
            Thus $f$ is a function by \cref{funs_is_function}.
            Thus $y = f(x)$ by \cref{function_apply_iff}.
            Thus $f(x) = 1$.
            Thus $y = 1$.
            Thus $y\in\{1\}$ by \cref{singleton_intro}.
        \end{subproof}
        Thus $\img{f}{B}\subseteq \{1\}$ by \cref{subseteq}.
    \end{subproof}

    $I_0\subseteq\reals$ by \cref{intervalclosed,subseteq}.
    Show that $f$ is continuous from $X$ to $I_0$ with respect to $T$ and $T_0$.
    \begin{subproof}
        Thus $f$ is continuous from $X$ to $\reals$ with respect to $T$ and $\standardTopologyR$
            by \cref{urysohn_unit_interval_continuous_reals}.
        $T$ is a topology on $X$ by \cref{normal_space}.
        Thus $T_0 = \subspaceTopology{\standardTopologyR}{I_0}$.
        Thus $f\in\funs{X}{\reals}$ by \cref{urysohn_unit_interval_f_funs}.
        Thus $f$ is continuous from $X$ to $I_0$ with respect to $T$ and $T_0$
            by \cref{urysohn_unit_interval_continuous_subspace}.
    \end{subproof}
    Thus there exists $f$ such that
        $f$ is continuous from $X$ to $I_0$ with respect to $T$ and $T_0$ and
        $\img{f}{A}\subseteq \{0\}$ and
        $\img{f}{B}\subseteq \{1\}$.
\end{proof}

% from §33 Item 1: Urysohn lemma
\begin{theorem}\label{urysohn_lemma}
    Assume $X$ is normal with respect to $T$.
    Assume $A$ is closed in $X$ with respect to $T$.
    Assume $B$ is closed in $X$ with respect to $T$.
    Assume $A$ is disjoint from $B$.
    Assume $a\in\reals$ and $b\in\reals$ and $a < b$.
    Let $I = \intervalclosed{a}{b}$.
    Let $T_I = \subspaceTopology{\standardTopologyR}{I}$.
    Then there exists $f$ such that
        $f$ is continuous from $X$ to $I$ with respect to $T$ and $T_I$ and
        for all $x\in A$ we have $f(x) = a$ and
        for all $x\in B$ we have $f(x) = b$.
\end{theorem}
\begin{proof}
    Let $I_0 = \intervalclosed{0}{1}$.
    Let $T_0 = \subspaceTopology{\standardTopologyR}{I_0}$.
    Take $g$ such that
        $g$ is continuous from $X$ to $I_0$ with respect to $T$ and $T_0$ and
        $\img{g}{A}\subseteq \{0\}$ and
        $\img{g}{B}\subseteq \{1\}$
        by \cref{urysohn_unit_interval}.
    Let $c = b - a$.
    Let $j = \identity{I_0}$.
    Let $g_1 = j\circ g$.
    Show that $I_0$ is a subspace of $\reals$ with respect to $\standardTopologyR$.
    \begin{subproof}
        $\standardBasisR$ is a basis on $\reals$ by \cref{standard_basis_is_basis}.
        Thus $\basisTopology{\standardBasisR}{\reals}$ is a topology on $\reals$
            by \cref{basis_topology_is_topology}.
        Thus $\standardTopologyR$ is a topology on $\reals$ by \cref{standard_topology_reals}.
        Show that $I_0\subseteq\reals$.
        \begin{subproof}
            Show that for all $y\in I_0$ we have $y\in\reals$.
            \begin{subproof}
                Fix $y$.
                Assume $y\in I_0$.
                Thus $y\in\reals$ by \cref{intervalclosed}.
            \end{subproof}
            Thus $I_0\subseteq\reals$ by \cref{subseteq}.
        \end{subproof}
        Thus $I_0$ is a subspace of $\reals$ with respect to $\standardTopologyR$ by \cref{subspace_of}.
    \end{subproof}
    $g_1$ is continuous from $X$ to $\reals$ with respect to $T$ and $\standardTopologyR$
        by \cref{continuous_expand_range}.
    Let $k = \{(x,a) \mid x\in X\}$.
    Let $h = \{(x,c) \mid x\in X\}$.
    $T$ is a topology on $X$ by \cref{normal_space}.
    $a\in\reals$ by assumption.
    $\neg{a}\in\reals$ by \cref{real_add_inverse}.
    Thus $b + \neg{a}\in\reals$ by \cref{real_add_closed}.
    Thus $b - a\in\reals$ by \cref{real_sub}.
    Thus $c\in\reals$.
    $\standardBasisR$ is a basis on $\reals$ by \cref{standard_basis_is_basis}.
    Thus $\basisTopology{\standardBasisR}{\reals}$ is a topology on $\reals$
        by \cref{basis_topology_is_topology}.
    Thus $\standardTopologyR$ is a topology on $\reals$ by \cref{standard_topology_reals}.
    Show that $I\subseteq\reals$.
    \begin{subproof}
        Show that for all $y\in I$ we have $y\in\reals$.
        \begin{subproof}
            Fix $y$.
            Assume $y\in I$.
            Thus $y\in\reals$ by \cref{intervalclosed}.
        \end{subproof}
        Thus $I\subseteq\reals$ by \cref{subseteq}.
    \end{subproof}
    Thus $I$ is a subspace of $\reals$ with respect to $\standardTopologyR$ by \cref{subspace_of}.
    $T_I$ is a topology on $I$ by \cref{subspace_topology_is_topology}.
    Show that $I_0\subseteq\reals$.
    \begin{subproof}
        Show that for all $y\in I_0$ we have $y\in\reals$.
        \begin{subproof}
            Fix $y$.
            Assume $y\in I_0$.
            Thus $y\in\reals$ by \cref{intervalclosed}.
        \end{subproof}
        Thus $I_0\subseteq\reals$ by \cref{subseteq}.
    \end{subproof}
    $T_0$ is a topology on $I_0$ by \cref{subspace_topology_is_topology}.
    $k$ is a function from $X$ to $\reals$ by \cref{constant_map_function}.
    Show that for all $x\in X$ we have $k(x) = a$.
    \begin{subproof}
        Fix $x$.
        Assume $x\in X$.
        Thus $(x,a)\in k$.
        $k$ is a function.
        Thus $k(x) = a$ by \cref{function_apply_iff}.
    \end{subproof}
    $k$ is continuous from $X$ to $\reals$ with respect to $T$ and $\standardTopologyR$
        by \cref{constant_continuous}.
    $h$ is a function from $X$ to $\reals$ by \cref{constant_map_function}.
    Show that for all $x\in X$ we have $h(x) = c$.
    \begin{subproof}
        Fix $x$.
        Assume $x\in X$.
        Thus $(x,c)\in h$.
        $h$ is a function.
        Thus $h(x) = c$ by \cref{function_apply_iff}.
    \end{subproof}
    $h$ is continuous from $X$ to $\reals$ with respect to $T$ and $\standardTopologyR$
        by \cref{constant_continuous}.
    Let $m = \{(x, h(x) \rmul g_1(x)) \mid x\in X\}$.
    Let $f_0 = \{(x, k(x) + m(x)) \mid x\in X\}$.
    $m$ is continuous from $X$ to $\reals$ with respect to $T$ and $\standardTopologyR$
        by \cref{real_function_multiplication_continuous}.
    $f_0$ is continuous from $X$ to $\reals$ with respect to $T$ and $\standardTopologyR$
        by \cref{real_function_addition_continuous}.

    Show that $\img{f_0}{X}\subseteq I$.
    \begin{subproof}
        Show that for all $x\in X$ we have $f_0(x)\in I$.
        \begin{subproof}
            Fix $x\in X$.
            $g\in\funs{X}{I_0}$ by \cref{continuous}.
            Thus $X = \dom{g}$ by \cref{funs_elim}.
            Thus $x\in\dom{g}$.
            Thus $g(x)\in I_0$ by \cref{funs_elim}.
            Thus $g$ is a function by \cref{funs_is_function}.
            Thus $g$ is right-unique by \cref{function_rightunique}.
            Thus $g$ is a relation by \cref{funs_is_relation}.
            Thus $g\in\rels{X}{I_0}$ by \cref{funs_subseteq_rels,subseteq}.
            Thus $\ran{g}\subseteq I_0$ by \cref{rels_ran_subseteq}.
            Thus $\dom{j} = I_0$ by \cref{id_dom}.
            Thus $\ran{g}\subseteq\dom{j}$.
            $j$ is right-unique by \cref{id_rightunique}.
            $j$ is a relation by \cref{id_is_relation}.
            Thus $g_1(x) = j(g(x))$ by \cref{circ_apply}.
            Thus $g_1(x) = g(x)$ by \cref{id_apply}.
            Thus $g_1(x)\in I_0$.
            Thus $g_1(x)\in\reals$ by \cref{intervalclosed}.
            Thus $0 \leq g_1(x)$ and $g_1(x) \leq 1$ by \cref{intervalclosed}.
            Thus $c\in\reals$.
            $m\in\funs{X}{\reals}$ by \cref{continuous}.
            Thus $m$ is a function by \cref{funs_is_function}.
            Thus $(x, h(x) \rmul g_1(x))\in m$.
            Thus $m(x) = h(x) \rmul g_1(x)$ by \cref{function_apply_iff}.
            Thus $h(x) = c$.
            Thus $m(x) = c \rmul g_1(x)$.
            Show that $0 \leq c \rmul g_1(x)$.
            \begin{subproof}
                \begin{byCase}
                \caseOf{$g_1(x) = 0$.}
                    Thus $g_1(x) \rmul c = 0$ by \cref{real_mul_zero}.
                    Thus $c \rmul g_1(x) = g_1(x) \rmul c$ by \cref{real_mul_comm}.
                    Thus $c \rmul g_1(x) = 0$.
                    Thus $0 \leq c \rmul g_1(x)$.
                \caseOf{$g_1(x) \neq 0$.}
                    Thus $0 < g_1(x)$ or $g_1(x) < 0$ by \cref{real_lt_trichotomy}.
                    Thus $0 < g_1(x)$.
                    Thus $0 < c$ by \cref{real_lt_imp_pos_diff}.
                    Thus $0 < c \rmul g_1(x)$ by \cref{real_mul_pos_pos_gt}.
                    Thus $0 \leq c \rmul g_1(x)$.
                \end{byCase}
            \end{subproof}
            Thus $0 \leq m(x)$.
            Show that $c \rmul g_1(x) \leq c$.
            \begin{subproof}
                \begin{byCase}
                \caseOf{$g_1(x) = 1$.}
                    Thus $c \rmul g_1(x) = g_1(x) \rmul c$ by \cref{real_mul_comm}.
                    Thus $c \rmul g_1(x) = 1 \rmul c$.
                    Thus $c \rmul g_1(x) = c$ by \cref{real_mul_ident}.
                    Thus $c \rmul g_1(x) \leq c$.
                \caseOf{$g_1(x) \neq 1$.}
                    Thus $g_1(x) < 1$ by \cref{real_lt_trichotomy}.
                    Thus $0 < c$ by \cref{real_lt_imp_pos_diff}.
                    $1\in\reals$ by \cref{real_naturals_subset_reals,real_one_in_naturals,subseteq}.
                    Thus $g_1(x) \rmul c < 1 \rmul c$ by \cref{real_lt_mul_pos}.
                    Thus $c \rmul g_1(x) = g_1(x) \rmul c$ by \cref{real_mul_comm}.
                    Thus $c \rmul 1 = 1 \rmul c$ by \cref{real_mul_comm}.
                    Thus $c \rmul g_1(x) < c \rmul 1$.
                    Thus $c \rmul g_1(x) < c$ by \cref{real_mul_ident}.
                    Thus $c \rmul g_1(x) \leq c$.
                \end{byCase}
            \end{subproof}
            Thus $m(x) \leq c$.
            Thus $c \rmul g_1(x)\in\reals$ by \cref{real_mul_closed}.
            Thus $m(x)\in\reals$.
            $f_0\in\funs{X}{\reals}$ by \cref{continuous}.
            Thus $f_0$ is a function by \cref{funs_is_function}.
            Thus $(x, k(x) + m(x))\in f_0$.
            Thus $f_0(x) = k(x) + m(x)$ by \cref{function_apply_iff}.
            Thus $k(x) = a$.
            Thus $f_0(x) = a + m(x)$ by \cref{real_add_eq_subst}.
            Show that $a + m(x)\in I$.
            \begin{subproof}
                Thus $a + m(x)\in\reals$ by \cref{real_add_closed}.
                $0\in\reals$ by \cref{real_add_ident}.
                Thus $0 \leq m(x)$.
                Thus $m(x) \leq c$.
                Thus $0 + a \leq m(x) + a$ by \cref{real_leq_add}.
                Thus $a \leq m(x) + a$ by \cref{real_add_ident}.
                Thus $a \leq a + m(x)$ by \cref{real_add_comm}.
                Thus $m(x) + a \leq c + a$ by \cref{real_leq_add}.
                Thus $a + m(x) = m(x) + a$ by \cref{real_add_comm}.
                Thus $a + c = c + a$ by \cref{real_add_comm}.
                Thus $a + m(x) \leq a + c$.
                Show that $a + c = b$.
                \begin{subproof}
                    Thus $c = b - a$.
                    Thus $a + c = a + (b - a)$.
                    Thus $a + (b - a) = a + (b + \neg{a})$ by \cref{real_sub}.
                    Thus $a + (b + \neg{a}) = (a + b) + \neg{a}$ by \cref{real_add_assoc}.
                    Thus $(a + b) + \neg{a} = (b + a) + \neg{a}$ by \cref{real_add_comm}.
                    Thus $(b + a) + \neg{a} = b + (a + \neg{a})$ by \cref{real_add_assoc}.
                    Thus $a + \neg{a} = 0$ by \cref{real_add_inverse}.
                    Thus $b + (a + \neg{a}) = b + 0$.
                    Thus $b + 0 = 0 + b$ by \cref{real_add_comm}.
                    Thus $0 + b = b$ by \cref{real_add_ident}.
                    Thus $a + c = b$.
                \end{subproof}
                Thus $a + m(x) \leq b$.
                Thus $a + m(x)\in I$ by \cref{intervalclosed}.
            \end{subproof}
            Thus $f_0(x)\in I$ by \cref{function_apply_intro}.
        \end{subproof}
        Show that for all $y\in\img{f_0}{X}$ we have $y\in I$.
        \begin{subproof}
            Fix $y$.
            Assume $y\in\img{f_0}{X}$.
            Take $x$ such that $x\in X$ and $(x,y)\in f_0$ by \cref{img_iff}.
            $f_0\in\funs{X}{\reals}$ by \cref{continuous}.
            Thus $f_0$ is a function by \cref{funs_is_function}.
            Thus $y = f_0(x)$ by \cref{function_apply_iff}.
            Thus $y\in I$.
        \end{subproof}
        Thus $\img{f_0}{X}\subseteq I$ by \cref{subseteq}.
    \end{subproof}

    $f_0$ is continuous from $X$ to $I$ with respect to $T$ and $T_I$
        by \cref{continuous_restrict_range}.

    Show that for all $x\in A$ we have $f_0(x) = a$.
    \begin{subproof}
        Fix $x\in A$.
        Thus $A\subseteq X$ by \cref{closed_in}.
        Thus $x\in X$ by \cref{subseteq}.
        $g\in\funs{X}{I_0}$ by \cref{continuous}.
        Thus $g$ is a function by \cref{funs_is_function}.
        Thus $X = \dom{g}$ by \cref{funs_elim}.
        Thus $x\in\dom{g}$.
        Thus $(x,g(x))\in g$ by \cref{function_apply_elim}.
        Thus $g(x)\in\img{g}{A}$ by \cref{img_iff}.
        Thus $g(x)\in\{0\}$ by \cref{subseteq}.
        Thus $g(x) = 0$ by \cref{singleton_elim}.
        Thus $g(x)\in I_0$ by \cref{funs_elim}.
        Thus $g$ is right-unique by \cref{function_rightunique}.
        Thus $g$ is a relation by \cref{funs_is_relation}.
        Thus $g\in\rels{X}{I_0}$ by \cref{funs_subseteq_rels,subseteq}.
        Thus $\ran{g}\subseteq I_0$ by \cref{rels_ran_subseteq}.
        Thus $\dom{j} = I_0$ by \cref{id_dom}.
        Thus $\ran{g}\subseteq\dom{j}$.
        $j$ is right-unique by \cref{id_rightunique}.
        $j$ is a relation by \cref{id_is_relation}.
        Thus $g_1(x) = j(g(x))$ by \cref{circ_apply}.
        Thus $g_1(x) = g(x)$ by \cref{id_apply}.
        Thus $g_1(x)\in I_0$.
        Thus $g_1(x)\in\reals$ by \cref{intervalclosed}.
        Thus $g_1(x) = 0$.
        $m\in\funs{X}{\reals}$ by \cref{continuous}.
        Thus $m$ is a function by \cref{funs_is_function}.
        Thus $(x, h(x) \rmul g_1(x))\in m$.
        Thus $m(x) = h(x) \rmul g_1(x)$ by \cref{function_apply_iff}.
        Thus $h(x) = c$.
        Thus $m(x) = c \rmul g_1(x)$.
        Thus $c \rmul g_1(x) = g_1(x) \rmul c$ by \cref{real_mul_comm}.
        Thus $c \rmul g_1(x) = 0$ by \cref{real_mul_zero}.
        Thus $m(x) = 0$.
        $f_0\in\funs{X}{\reals}$ by \cref{continuous}.
        Thus $f_0$ is a function by \cref{funs_is_function}.
        Thus $(x, k(x) + m(x))\in f_0$.
        Thus $f_0(x) = k(x) + m(x)$ by \cref{function_apply_iff}.
        Thus $k(x) = a$.
        Thus $f_0(x) = a + 0$ by \cref{real_add_eq_subst}.
        Thus $a + 0 = 0 + a$ by \cref{real_add_comm}.
        Thus $f_0(x) = a$ by \cref{real_add_ident}.
    \end{subproof}
    Show that for all $x\in B$ we have $f_0(x) = b$.
    \begin{subproof}
        Fix $x\in B$.
        Thus $B\subseteq X$ by \cref{closed_in}.
        Thus $x\in X$ by \cref{subseteq}.
        $g\in\funs{X}{I_0}$ by \cref{continuous}.
        Thus $g$ is a function by \cref{funs_is_function}.
        Thus $X = \dom{g}$ by \cref{funs_elim}.
        Thus $x\in\dom{g}$.
        Thus $(x,g(x))\in g$ by \cref{function_apply_elim}.
        Thus $g(x)\in\img{g}{B}$ by \cref{img_iff}.
        Thus $g(x)\in\{1\}$ by \cref{subseteq}.
        Thus $g(x) = 1$ by \cref{singleton_elim}.
        Thus $g(x)\in I_0$ by \cref{funs_elim}.
        Thus $g$ is right-unique by \cref{function_rightunique}.
        Thus $g$ is a relation by \cref{funs_is_relation}.
        Thus $g\in\rels{X}{I_0}$ by \cref{funs_subseteq_rels,subseteq}.
        Thus $\ran{g}\subseteq I_0$ by \cref{rels_ran_subseteq}.
        Thus $\dom{j} = I_0$ by \cref{id_dom}.
        Thus $\ran{g}\subseteq\dom{j}$.
        $j$ is right-unique by \cref{id_rightunique}.
        $j$ is a relation by \cref{id_is_relation}.
        Thus $g_1(x) = j(g(x))$ by \cref{circ_apply}.
        Thus $g_1(x) = g(x)$ by \cref{id_apply}.
        Thus $g_1(x)\in I_0$.
        Thus $g_1(x)\in\reals$ by \cref{intervalclosed}.
        Thus $g_1(x) = 1$.
        $m\in\funs{X}{\reals}$ by \cref{continuous}.
        Thus $m$ is a function by \cref{funs_is_function}.
        Thus $(x, h(x) \rmul g_1(x))\in m$.
        Thus $m(x) = h(x) \rmul g_1(x)$ by \cref{function_apply_iff}.
        Thus $h(x) = c$.
        Thus $m(x) = c \rmul g_1(x)$.
        Thus $c \rmul g_1(x) = g_1(x) \rmul c$ by \cref{real_mul_comm}.
        Thus $m(x) = 1 \rmul c$.
        Thus $m(x) = c$ by \cref{real_mul_ident}.
        $f_0\in\funs{X}{\reals}$ by \cref{continuous}.
        Thus $f_0$ is a function by \cref{funs_is_function}.
        Thus $(x, k(x) + m(x))\in f_0$.
        Thus $f_0(x) = k(x) + m(x)$ by \cref{function_apply_iff}.
        Thus $k(x) = a$.
        Thus $f_0(x) = a + c$ by \cref{real_add_eq_subst}.
        Show that $a + c = b$.
        \begin{subproof}
            Thus $c = b - a$.
            Thus $a + c = a + (b - a)$.
            Thus $a + (b - a) = a + (b + \neg{a})$ by \cref{real_sub}.
            Thus $a + (b + \neg{a}) = (a + b) + \neg{a}$ by \cref{real_add_assoc}.
            Thus $(a + b) + \neg{a} = (b + a) + \neg{a}$ by \cref{real_add_comm}.
            Thus $(b + a) + \neg{a} = b + (a + \neg{a})$ by \cref{real_add_assoc}.
            Thus $a + \neg{a} = 0$ by \cref{real_add_inverse}.
            Thus $b + (a + \neg{a}) = b + 0$.
            $0\in\reals$ by \cref{real_add_ident}.
            $b\in\reals$ by assumption.
            Thus $b + 0 = 0 + b$ by \cref{real_add_comm}.
            Thus $0 + b = b$ by \cref{real_add_ident}.
            Thus $a + c = b$.
        \end{subproof}
        Thus $f_0(x) = b$.
    \end{subproof}
    Thus there exists $f$ such that
        $f$ is continuous from $X$ to $I$ with respect to $T$ and $T_I$ and
        for all $x\in A$ we have $f(x) = a$ and
        for all $x\in B$ we have $f(x) = b$.
\end{proof}

